\documentclass[11pt,a4paper]{article}
\usepackage[T1]{fontenc}
\usepackage[utf8]{inputenc}
\usepackage{lmodern,amsmath,amssymb,amsthm,mathtools,microtype,mathrsfs}
\usepackage[margin=25mm,top=23mm,bottom=25mm]{geometry}
\usepackage{booktabs,array,longtable,enumitem,xcolor,fancyhdr}
\usepackage[colorlinks=true,linkcolor=blue!45!black,urlcolor=blue!45!black,citecolor=blue!45!black]{hyperref}
\hypersetup{pdftitle={Reading the Alpoge--Claude RH and Dirichlet-GRH proof: working version 0.35},pdfauthor={ChatGPT},pdfsubject={Partial human-readable reconstruction with explicit arithmetic inputs},pdfkeywords={Riemann hypothesis, Lean, theta lifts, prime sums, Arakelov slopes}}
\setlength{\parindent}{0pt}
\setlength{\parskip}{5pt plus 1pt}
\setlist{topsep=3pt,itemsep=3pt,leftmargin=*}
\setlength{\headheight}{14pt}
\pagestyle{fancy}\fancyhf{}
\fancyhead[L]{\small RH and Dirichlet GRH: a reconstruction}
\fancyhead[R]{\small Working version 0.35}
\fancyfoot[C]{\thepage}
\renewcommand{\headrulewidth}{0.3pt}
\numberwithin{equation}{section}
\newtheorem{theorem}{Theorem}[section]
\newtheorem{proposition}[theorem]{Proposition}
\newtheorem{lemma}[theorem]{Lemma}
\newtheorem{corollary}[theorem]{Corollary}
\theoremstyle{remark}\newtheorem{remark}[theorem]{Remark}
\newcommand{\Rea}{\operatorname{Re}}
\newcommand{\HS}{\mathrm{HS}}
\newcommand{\tr}{\operatorname{tr}}
\newcommand{\codim}{\operatorname{codim}}
\newcommand{\covol}{\operatorname{covol}}
\newcommand{\OK}{\mathcal O_K}
\newcommand{\src}[2]{\href{https://github.com/CocoSgt/riemann-hypothesis-lean/blob/7cbf7a0c32fca1badf05931879a37c3f8a2b03b4/src/#1}{\texttt{#2}}}
\newcommand{\status}[1]{\textbf{Status:} #1}
\begin{document}
\begin{center}
{\LARGE\bfseries Reading the Alp\"oge--Claude RH/GRH proof}\\[7pt]
{\large Completed global assembly, original lowest types, and auxiliary primes}\\[-1pt]
{\large A refined completion audit with 71 new proof claims}\\[8pt]
Working notes for Benjamin Matschke\quad |\quad 22 November 2026
\end{center}

\begingroup\setlength{\fboxsep}{9pt}
\noindent\fcolorbox{blue!35!black}{blue!3}{\begin{minipage}{0.95\linewidth}
\textbf{What version 0.35 achieves.}
Both original losses and the exact same-family tied-ledger and final
\texttt{Closed} assembly are now reconstructed. The selected global
assembly leaf is proved on its original inputs. The original full
lowest compact type is smooth and irreducible with multiplicity one;
the exact extremal coefficient and Hilbert-map stages are supplied.
Count II has its complete original auxiliary-prime supplier under
unchanged \texttt{CohP}, and residual Hecke values for every supplied
integral type.

The selected register has 34 open coupled leaves: 23 C, 9 B and 2 D.
Common geometry, compatible choices, local/spectral estimates,
simultaneous ioP and global Sobolev completion, ordinary/control and
same-ideal automorphic support remain. RH retains I/II/V;
unrestricted Dirichlet GRH also retains III.

Every historical mathematical proof, source array and dependency edge
is retained. The current audit governs older frontier statements.
We accept the underlying Lean proof as fully checked and correct.
This remains \textbf{a partial human reconstruction}; no new Lean
replay is claimed.
\end{minipage}}\endgroup

\section{The route and the present boundary}

The released Lean proof has a short mathematical route at its top level:
\begin{enumerate}
\item Construct a suitably centered, smoothed prime sum that detects a fixed target $L$-function.
\item Build iterated theta-period maps with a common integral target. Compare finite-place divisibility with their archimedean sizes and overlaps.
\item Obtain a one-sided square-root bound for the prime sum throughout arbitrarily late intervals $[Y,Y^{3/2}]$.
\item Show that any zero to the right of $1/2$ would force a violation on every sufficiently late such interval.
\end{enumerate}
The difficulty is concentrated in the quantitative arithmetic in steps 2--3. Steps 1 and 4, the numerical end of step 3, and several algebraic and slope arguments inside step 2 now have compact proofs below. This division follows actual source interfaces, rather than the chronology of the discovery.

The source is the CocoSgt mirror at commit
\begin{center}\small\texttt{7cbf7a0c32fca1badf05931879a37c3f8a2b03b4}.\end{center}
The main statement uses Mathlib's definition of RH; its generalization concerns Dirichlet $L$-functions. Sutherland's verification page, in the 20 November 2026, 22:05 UTC snapshot inspected for version 0.1, reports compilation of all 175,187 modules, standard-axiom dependencies for the headline results, and independent checker runs. No full build or kernel replay was run for this reconstruction. Seven pivotal files were also byte-identical in the Deicyde mirror at revision \texttt{e7c8972c701d3396947808ab20cc16dc14623d33}.\cite{verification,main}

\paragraph{How to read the provenance.}
Tags such as \textsf{[LC2]} identify claims in \texttt{provenance.json} and its readable companion \texttt{provenance.txt}. Every tag states whether it is a source translation, a derived simplification, or an input still awaiting a human proof. The ledger records the original filename, declaration, line range, pinned link, and full-file SHA-256. The accompanying \texttt{backtrack.py} can verify those hashes and print a cited passage from a local checkout. The more technical notes preserve calculations omitted from the main text.


\section{Current dependency audit: completed global assembly and original lowest types}
\label{sec:dependency-audit-v35}

\paragraph{Status and exact scope.}
The final-global assembly implication is now reconstructed on its original
inputs. This continuation adds 71 human claims. It also proves
full original lowest-type multiplicity one, the extremal coefficient and
Hilbert-map stages, the complete original auxiliary-prime supplier, and
residual Hecke values for every supplied integral type.
The underlying Alpoge--Claude Lean proof is accepted as fully checked
and correct. This is still a \textbf{partial human reconstruction}:
the common geometric, analytic, local, spectral and ordinary arithmetic
producers listed below remain to be written out.

The selected sufficient graph has 121 nodes and
34 open coupled leaves: 23 C,
9 B and 2 D.
The entire conditional \texttt{TASK\_\allowbreak GL\_\allowbreak FINAL\_\allowbreak ASSEMBLY} leaf is
closed; the theta and CountII leaves are narrowed. All 245 historical
compound targets in 107 scope groups are retained. These counts are
neither completion percentages nor a theorem of absolute minimality.
Optional routes remain outside the selected AND requirements.

\begin{center}\small
\begin{tabular}{cp{0.80\textwidth}}
A & A published theorem applied with every original hypothesis and normalization checked.\\
B & An elementary or moderately difficult argument to reconstruct independently.\\
C & A substantial or deep reconstruction obligation.\\
D & A possible gap in our human understanding of the formal source.\\
\end{tabular}
\end{center}
The two D leaves remain \texttt{AD\_\allowbreak VANISH\_\allowbreak K\_\allowbreak MATCH} and
\texttt{HOST\_\allowbreak RIGID\_\allowbreak S}; they do not allege an error in the checked proof.

\paragraph{The complete original global consumer.}
LFS1--LFS17 derive both original losses from the same chosen tops,
bad-prime bound, local gains and chosen realized degree pair.
Four prime classes yield eleven terms, with order constant 12,
leading coefficient $2\log(4)C$ and the thirteen-error remainder.
The gain and prime-logarithmic dictionaries retain every finite-place
weight. The asked-sum identity holds for every permitted reading,
arithmetic-record equality and cutoff. The original $a_v,\allowbreak{} b_v,\allowbreak{} i_1$
remain unchanged; no \texttt{Cover\allowbreak All} premise is added.

GFI1--GFI17 construct the actual common ledger, all ties, one character
and cutoff, every ledger inequality and every error, size, growth and
cost field. The original chosen-data conversion and the final
\texttt{Closed} order are proved. Its epsilon cutoff and joined
transport tuple precede the family and $C_H$; the eventual threshold
precedes all members, base points, words, places and readings.
The selected constants remain $C_{\rm abs}=2,\allowbreak{} C_0=3$, column constant
$4$, $e_D=2,\allowbreak{} e_L=0$. The column constant and chosen quadratic parameter
$C_L$ have their distinct original roles. The final zero-free endpoint
is still $1/2+9\epsilon$.
Compatible family/selector existence, duties and the actual
\texttt{Link\allowbreak G/\allowbreak Link\allowbreak S/\allowbreak Link\allowbreak IRR/\allowbreak Link\allowbreak DIV} producers remain upstream.

\paragraph{The original lowest type and Hilbert stage.}
HCU reconstructs the signed frame domain, smooth triangular factor,
connectedness, local/global holomorphic uniqueness and the full Borel
dictionary. Finite compact-weight arguments and the explicit Vogan-norm
inequality give lowering on the full original lowest subspace.
Coefficient uniqueness proves that every lowest-type summand is the
same smooth irreducible seed. Thus \texttt{Low\allowbreak Sm/\allowbreak Low\allowbreak Irr}, including
multiplicity one of the original supremum, is supplied. The ordered
base model at its original coefficient embedding is checked.

On a bounded original \texttt{Hol\allowbreak Form} with the supplied same-target
\texttt{Fin\allowbreak Pair/\allowbreak Is\allowbreak Reg}, holomorphic integration gives the exact extremal
coefficient identity. The orbit Gram identity gives the full original
Hilbert multilinear map and finite eigenlaws. What remains is one
$io_i$ valid for every parameter, actual boundedness/finite volume,
smooth representatives and their derivatives, Sobolev point evaluation
and uniform growth, central finiteness and cusp closure. These are
needed for the pointwise smooth \texttt{MMap} on the full original
\texttt{Wall}; geometric readings, rational tests and base rows are
separate contracts.

\paragraph{The original auxiliary primes and residual values.}
CAP1--CAP14 prove full \texttt{Ar\allowbreak Pl} under unchanged \texttt{Coh\allowbreak P}:
the exact residual subgroup and its normality, the $\mathrm{PL}$-relative
Frattini tower, split degree-one Chebotarev primes for every list,
the actual enlarged levels, complete closed kernel and inverse-Frobenius
tame relation. The single bound
\[
 |\operatorname{Hom}_\ell(p_N^{-1}H)|
 \le |\operatorname{Hom}_\ell(H)|\ell^{q[\Gamma:H]}
\]
is independent of $N$. Unconditional finite-Galois Chebotarev is the
one new published input, with its hypotheses and convention checked.
CAP15--CAP17 construct original \texttt{Is\allowbreak Res} values for any supplied
integral type at every \texttt{QData} level.
Ordinary/control, the type and \texttt{neat\allowbreak U/\allowbreak TTheta\allowbreak Ex}, four further
Hida endpoints, same-ideal all-level \texttt{Has\allowbreak Aut}, \texttt{Choice},
first-order/obstruction bounds and \texttt{Raw} remain producers.

\paragraph{Final cut and next priority.}
The exact final ten-input interface and every shared common-object
requirement are unchanged. RH and two-primary targets retain I/II/V;
unrestricted Dirichlet targets additionally retain III, SeedIV and
the chain arguments. The original window exponent 401, floor
$x^{400}$ and $1/2+9\epsilon$ estimate retain their scopes.
Re-audit the common original construction and the remaining theta/ordinary producers against the final ten-input cut. The next bounded analytic producer is the simultaneous JointHarm.exists\_\allowbreak ioP: first construct its parameter-free exists\_\allowbreak io\_\allowbreak poly on the actual CoefMapTie, then prove the original compact Weil-character laws for every parameter t using the completed full lowest type. In parallel reconstruct the original one-point NiceI Choice at h=0 and a=n=4, including its finite detector and full coefficient-extension transport; prove the proposed zero-dimensional Kap specialization before using it, and only then assess its possible B2 consequence. These are work priorities, not additional closure claims. Keep bounded HolForm, the actual FinPair/IsReg target, smooth representatives, Sobolev bounds, central finiteness and cusp closure explicit, together with original ordinary/control and same-J HasAut. The common Bundle/readings/rational tests/local installation and compatible second layer still dominate several final inputs; prioritize any concrete producer that supplies several of those interfaces. Do not repeat the closed global assembly or substitute a compatible model for the original common objects.

\paragraph{Every remaining selected obligation.}
Each following paragraph is the full current contract from the graph.
Categories and completion status are separate; no conditional consumer
is counted as its missing producer.

\begin{description}[style=nextline]
\item[1. GLOBAL\_\allowbreak DENOMINATOR --- C] Actual LinkN local table and binding producer. Closed(Des $\mathcal C$ 3) LinkN: the actual etaSum is bounded with exact dZ0 and same-host vvTerm, subject only to the literal ShaFinOdd and summability guards.

\item[2. GLOBAL\_\allowbreak RETAINED\_\allowbreak DEFECT --- C] Original retained-row defect bound LinkIRR. Rows.DelpNonneg and delp$\le$exp(-x\textasciicircum{}(ez/7)) uniformly for every late actual row datum.

\item[3. GLOBAL\_\allowbreak PERIOD\_\allowbreak HEIGHT --- C] Full actual period and height comparison LinkDIVt. Closed conditional TuSize -> DGamSize -> (LinkEta -> DIVinN -> LinkDIV), with every original LinkDIV\_\allowbreak at clause.

\item[4. GLOBAL\_\allowbreak LOCAL\_\allowbreak GAIN --- C] Original global LinkG installation. One uniform LinkG with actual local lattice lengths, served-orbit coverage and MeanKappaBound; delta(x)->0 before varying actual indices.

\item[5. GLOBAL\_\allowbreak SPECTRAL --- C] Original normalized kernel/Gram estimates LinkS. Closed TuSize -> LinkS with exact Gamma normalization, off-diagonal Hilbert-Schmidt bound, diagonal error and Dall/Dstar estimate, errors exp(-x\textasciicircum{}(epsilon/8)).

\item[6. TASK\_\allowbreak PROP\_\allowbreak GK\_\allowbreak INSTALL --- B] Install OneChar and the congruence integer into PropGK. Complete the finite orbit/\allowbreak Frobenius/\allowbreak valuation argument producing one nonzero integer before every later coefficient prime, with the original e,\allowbreak{} N,\allowbreak{} kernel,\allowbreak{} p and fixed-K constants.

\item[7. TASK\_\allowbreak ONE\_\allowbreak CHAR\_\allowbreak PROOF --- C] Reconstruct the original one-character argument. From exact InpZ and AutomorphyUnramifiedCoefficients, prove the original OneChar trace/length estimate on the supplied coefficient datum and GD input.

\item[8. TASK\_\allowbreak INP\_\allowbreak Z\_\allowbreak ASSEMBLY --- B] Install the exact Kato input record. Build the original CohInp/InpZ record using exactly lemM,\allowbreak{} lemMprime,\allowbreak{} lemP,\allowbreak{} RigII,\allowbreak{} AdVanishK and number on the same arithmetic objects. No rank-four replacement is permitted.

\item[9. RIG\_\allowbreak II --- C] Type-II polarized adjoint rigidity. Marked original datum, actual type-II point and modularity of all source extensions imply vanishing for its exact adRep.

\item[10. AD\_\allowbreak VANISH\_\allowbreak K\_\allowbreak MATCH --- D] Exact rank-three polarized adjoint matching. AdVanishK uses the full negative-transpose complex-conjugation action and irreducibility on the infinite cyclotomic kernel. RX does not transfer to every rank by name.

\item[11. AUT\_\allowbreak UNRAM\_\allowbreak COEF --- C] Actual unramified-coefficient automorphy. For each actual Dat with coefficient ramification one and fixed iota, choose one marked g1 supplying FinA, AutB and LemW together.

\item[12. CNR3 --- C] Uniform rank-three and rank-two congruence integer. Fixed quadratic K -> c,\allowbreak{} A -> e,\allowbreak{} N,\allowbreak{} p,\allowbreak{} w,\allowbreak{} finite Galois-stable roots-of-unity orbit -> nonzero integer NN -> every later coefficient prime, closure, embedding, Frobenius and finite admissible collection.

\item[13. COUNT\_\allowbreak II\_\allowbreak SUPPORT --- C] Two actual rank-four Links with common threshold and same-ideal Hida/P1. For each of the two original QX rank-four Links, at fixed iota choose one common coefficient threshold before every later coefficient DVR and actual universal B/Lambda; produce one ideal J with ring,\allowbreak{} dense,\allowbreak{} pt,\allowbreak{} end1,\allowbreak{} end2,\allowbreak{} twist AND P1core. JSP supplies the exact common ideal and conditional support transport. CIP supplies every-map split-CM reciprocity with the positive cyclotomic norm sign and the full universal NiceI Norm with original D0, integral representation, ordered flags and determinant law. CAP now supplies the full original ArPl under unchanged CohP: PhiCond, the exact polarized residual H, the PL-relative Frattini tower, every prescribed list and all actual TWGrp/TWPl levels, inverse arithmetic Frobenius, complete closed kernel and one uniform Hom bound. CAP also constructs the original IsRes values for any supplied integral type at every QData level. Remaining C producers: original ordinary/control/arithmetic-fiber data and integral type/neatU/TThetaEx, four further Hida fields on this J, all-level same-J automorphic modules/HasAut, Choice, first-order and obstruction bounds, and Raw, retaining original CohP and all stated side conditions. Both Links, ell=3, fixed iota and common threshold order remain. No generic NiceE Norm or ordinary/control existence follows from the completed consumers.

\item[14. TASK\_\allowbreak CASE3\_\allowbreak AUT\_\allowbreak JOIN --- B] Complete the original Case-III automorphy/image join. Install the actual Case-III ProModAlt, SeedIVAlt and shared PMExAlt with the inherited CohomologyHyp/Ferrero-Washington/HFit/SeedII facts into the original whole Case-III automorphy/image conclusion.

\item[15. TASK\_\allowbreak CASE3\_\allowbreak PROMOD\_\allowbreak JOIN --- C] Reconstruct the actual Case-III chain-to-promodularity proof. Use the original tensor-factor tower, ChainRankFour, ChainDesc4, seeded rank-four chain reduction, actual PMExAlt and SeedIVAlt to construct ProModAlt on every original Ctx3 star point at the same fixed iota, preserving all finite-coefficient, inertia and support maps.

\item[16. SEED\_\allowbreak IV\_\allowbreak ALT --- C] Actual Case-III seed with prechosen residual datum. Base Ctx3/Four tower, iota and Amb -> w,\allowbreak{} s before every later Dies tower -> compatible finite-coefficient seed with exact residual support/root entries.

\item[17. TENSOR\_\allowbreak FACTOR\_\allowbreak TOWER --- B] Actual tensor-factor tower and characters. Original finite-field, tame, factor and character constraints; source-level construction is sizeable and has not been reconstructed here.

\item[18. CHAIN\_\allowbreak RANK\_\allowbreak FOUR --- C] Actual SomeIn-to-AllIn propagation. Same source Space and PM ideal, residual point Good4, original bound 30(C+1)\textasciicircum{}2 <= d; geometry must preserve all identities.

\item[19. CHAIN\_\allowbreak DESC4 --- B] Original residual chain descent. Finite residual descent retaining source Chain4 compatibility.

\item[20. RANK4\_\allowbreak CHAIN\_\allowbreak REDUCTION --- C] Seeded reduction to the original residual chain. Select T1 before every later Deep tower; retain one compatible seed and Chain4 for every original residual lift over Q.

\item[21. TASK\_\allowbreak COUNT\_\allowbreak V\_\allowbreak INSTALL --- B] Complete the original CountV finite assembly. From exactly the listed source host/\allowbreak cohomology inputs, complete countV\_\allowbreak of\_\allowbreak finiteness\_\allowbreak duality\_\allowbreak euler\_\allowbreak host, including original host/frame formation, normalized length, the CaseV guard and R1<=2\textasciicircum{}c2.

\item[22. THM\_\allowbreak S6\_\allowbreak S --- C] Original unitary-three ordinary support and automorphy. Original host/frame and fixed iota -> one marked g1 -> every universal R,\allowbreak{} u -> weight-zero finiteness and every Type-I BPrime\_\allowbreak s point.

\item[23. HOST\_\allowbreak FRAME\_\allowbreak LENGTH --- B] Original Case-V host/frame and normalized-length installation. Actual HostData, frame, local-condition equality, abstract count and normalized original Selmer length must all refer to the same objects.

\item[24. HOST\_\allowbreak RIGID\_\allowbreak S --- D] Rigidity on the original CountV frame. Same original host/frame, marked point and eta-kernel; published adjoint vanishing is usable only after exact representation and hypothesis matching.

\item[25. HOST\_\allowbreak NUMBER\_\allowbreak S --- C] Uniform original-host rank-three catalogue. One exponent before all closures, coefficient rings, levels, characters, frames, embeddings and marked points; one separating trace function controls all later finite pairwise nonisomorphic collections.

\item[26. TASK\_\allowbreak HOST\_\allowbreak VAN\_\allowbreak E\_\allowbreak INSTALL --- B] Transport K/duality to literal HostVanE. Use the already supplied LemmaK/local-global duality and finiteness on the actual CountV host/frame and exact original coefficient modules.

\item[27. TASK\_\allowbreak HOST\_\allowbreak XPP\_\allowbreak INSTALL --- B] Transport CorX to literal HostXpp. Preserve Psi\textasciicircum{}-2 eta, the prime-to-ell quotient, actual local conditions and the original CountV host while installing BR7.

\item[28. DESIGN\_\allowbreak CHOICES\_\allowbreak COMMON --- C] Construct the original common design family and bundle. For the one target HostExW, construct the original compatible DesignChoices on Cdes=WithTw(WithSm Cprime), where Cprime=WithChN(A intersect target class), A is the one Built construction class, and cT(eps,\allowbreak{} y)=cTmin(y). Retain every original field, tuple, word, metric and common bundle duty. DSZ10 and RMP11 prove the size/rank properties of members of this class; they do not produce a member. The RecExP/RecAtG and common object/row route remains a coupled construction.

\item[29. LOCAL\_\allowbreak NBUILT\_\allowbreak COMMON --- C] Actual local specialization on the common selected family. Closed on the same Des(C,3) class: one eventual family of actual NBuilt certificates for the original Fm.D(x,\allowbreak{} i), including NdSpec, original iV/Hecke actions, every required maximal ideal, $\Psi$/range and projection square. Keep the constructor certificate before forgetting it to TSel.

\item[30. SECOND\_\allowbreak LAYER\_\allowbreak COMPATIBLE\_\allowbreak SELECTION --- C] Original second-layer selector and class duties. For the same A, supply Duty on A and a Closed common original TSel family whose selected records are the original true-column rec7 records (or have the entire dependent tuple transported), with their real witness and original coefficient-field root data. Existence/compatibility remains arithmetic; merely abstractly isomorphic fields or independently selected columns do not suffice.

\item[31. BASE\_\allowbreak THETA\_\allowbreak READINGS --- C] Actual base reading on the common bundle. For every family in the original class, eventually on the SAME selected Bundle and for each i<m, supply the actual R,\allowbreak{} F,\allowbreak{} Read,\allowbreak{} ReadL, weak models/Weil data and phase-one pins. Identify their LiftTie.theta, mtRead, mjRead and ThDef3 basis with that Bundle.P.A3, P.Zs, P.S0 and its stored maps. Exact field, embeddings, metrics, measures, row family and weak choices are retained. This is one part of the coupled BTH12 producer, not a freely replaceable isomorphic model.

\item[32. BASE\_\allowbreak THETA\_\allowbreak CUSPIDAL\_\allowbreak EXTENSION --- C] Extend the original convergent base integrals. On the SAME selected Bundle and SAME reading as BASE\_\allowbreak THETA\_\allowbreak READINGS, extend the exact lowest-type theta integral for each finite basis vector of HV(i,\allowbreak{} t0) to the original full Wall\_\allowbreak i, with one family threshold before i and basis vectors. BTH/FIA/CTC/TAE prove convergence, finite cancellation, genuine cuspidality, actual-action smoothness, one growth exponent for every derivative word, stored-vacuum transfer and exact lowering. HCU now proves the actual signed-frame holomorphic dictionary and uniqueness, full LowSm/LowIrr including multiplicity one of the original lowKSub, the exact base-model ordered-parameter installation, and holomorphicity of the extremal cusp packet. Under bounded HolForm, FinPair and IsReg on the SAME target, it proves the full original extremal coefficient identity and BackForm.exists\_\allowbreak map Hilbert multilinear realization with finite eigenlaws. Remaining C subcontracts: one original io\_\allowbreak i satisfying IsIoW for EVERY parameter t; actual boundedness and finite-volume regular target; pointwise smooth representatives and every derivative of the constructed Hilbert image, Sobolev point-evaluation estimates and uniform growth, central finiteness, and cusp preservation under closed orbit span, producing the literal Sobolev-continuous MMap agreeing with the original integral. Preserve all fields, adjoints, levels and common choices. Separate geometric readings, rational tests and actual base rows remain outside this analytic packet.

\item[33. BASE\_\allowbreak THETA\_\allowbreak RATIONAL\_\allowbreak TESTS --- C] Simultaneous original base rational tests. On the SAME original Bundle/reading/theta map, for every i<m supply one finite separating complex-linear test family. Its values must lie in the fixed embedded coefficient field on theta of EVERY original rational source vector and on EVERY stored target-coordinate basis vector. BTH8 and BTH11 use this one base LiftTestB family for rational adjoint image and matrix coordinates. Finite dimensionality alone does not supply these rationality identities.

\item[34. SLR\_\allowbreak BASE\_\allowbreak ROWS\_\allowbreak PRODUCER --- C] Actual base row identities on that same bundle. For the SAME selected Bundle, common original fields and readings used by the other three BTH12 inputs, supply actual RowsHold0\_\allowbreak 3 uniformly eventually for every family. Its original primitive row realization, spectrum/type and basis identifications must be retained. BTH supplies a theta consumer only; it does not construct the base row identity. No separate witness may be chosen merely to satisfy this row requirement.

\end{description}

\paragraph{Validation and continuity.}
Every predecessor proof, claim field, source array and dependency edge
is retained. Exact new source and proof ranges are bound to the pinned
commit and independently reviewed. The release checks all source ranges,
graph reachability, every displayed remaining contract and review hashes,
then checks TeX syntax and references in draft mode without a PDF.
These checks are not a new Lean replay. Historical RF16 remains governed
by DEX11. The current audit and additive scope updates govern older
frontier statements without silently rewriting historical mathematics.


\paragraph{Historical version-0.34 audit.}
The following audit is preserved as a historical checkpoint.
Section~\ref{sec:dependency-audit-v35} and the version-0.35 scope
updates give the current remaining obligations. The historical
conditional arguments retain their stated hypotheses.

\section{Current dependency audit: original global consumers, theta derivatives, and Count-II normalization}
\label{sec:dependency-audit-v34}

\status{Version 0.34 accepts the underlying Alp\"oge--Claude formal
proof as fully checked and correct. It adds 55 human proof claims:
GVC1--18, GFA1--12, TAE1--13 and CIP1--12. The human reconstruction
remains incomplete. Completed implications retain their original
producer premises; category D concerns our understanding of the
formal source. This current audit supersedes historical frontier
statements without changing any historical mathematical proof.}

\subsection{What this continuation supplies}

\paragraph{The exact global numerical consumers.}
GVC reconstructs the degree products and their finite Vandermonde
comparison on the original words. The large top--bottom factors
cancel in the cross-word ratio because the actual multiplicities
and shifted coordinates match. The remaining factors are controlled
uniformly by the original parameter windows. Together with the
previous mass and index identities, this proves the literal
\texttt{Count\allowbreak F}. It also installs all sixteen numerical and eight
weight clauses of \texttt{Nums\allowbreak L}/\texttt{Wts\allowbreak L}, with one threshold
before the word and family index.

Direct sum estimates applied to the original lower and maximal
column slopes give the exact \texttt{Conclusion5}/\texttt{Conclusion6}
and \texttt{F5}/\texttt{F6} consumers, and then
\texttt{Slopes\allowbreak MF}. Every original determinant shift is retained.
These sufficient arguments do not automatically prove every stronger
intermediate statement called \texttt{Part\allowbreak Five} or \texttt{Prop\allowbreak Z}
in the source. GVC18 also checks that the needed bound on
\(L.\mathrm{logDerivativeAtOne}\) is the \texttt{Dl\allowbreak Size} projection
of the original \texttt{Fam\allowbreak P.\allowbreak sizes}; it is not a new independent
producer for these consumers. Existence of that original family
remains a separate obligation.

\paragraph{Quotient slopes with the actual type-prime cost.}
For the original lattice \(N_u=I_uN_u^0\), GFA1 proves the precise
congruence cost
\[
 \frac1{r_u}\sum_{w\mid T_u}2\ell_we_w
                  L.\mathrm{lenMN}(u,\allowbreak{} w)
 \le36|TT|m(a/k+b')\log T_u.                         \tag{AU34.1}
\]
A finite subset of the places is allowed. The weight calculation
first gives constant nine, and the original retained-rank ratio is
at least one half. The source's factor two in the squared metric
then gives exactly 36. The original local length and ramified rank
identities, not substitute lattices, enter this calculation.

For the correction away from that prime, a useful direct argument
uses only what the consumer requires. Put
\[
 b=e^{x^\epsilon},\quad a=xb,\quad
 X'=e^{x^{\epsilon/3}},\quad S=e^{x^{\epsilon/2}}.
\]
The actual definitions of the local type-I/type-II correction and
ceilings give \(\tau_p\le aS^{20}\), with \(\tau_p=0\) for
\(p\ge X'\). The finite-place degree identity and a crude count of
primes then give
\[
 T_a=\sum_w2\ell_we_w\tau_{p(w)}
       \le aS^{24}\le b^2                              \tag{AU34.2}
\]
for all sufficiently large \(x\), uniformly before every index.
This bypasses the stronger weighted-prime estimate solely for this
consumer. It does not assert that stronger estimate has been proved.
The exact ideal shift, integral local quotient lengths, nonnegative
rounding error and the zero contribution at the exceptional prime
now give \texttt{Places\allowbreak F} and \texttt{Slopes\allowbreak NF} in the source's
normalization, from the original local-gain and metric premises.

\paragraph{The frame inequality has its full height quantifier.}
GFA6--9 check positive target rank, the algebraic dual with its
original metric, the actual rational Bessel map and every Gram
normalization. Let \(s\) be the weighted mean source slope minus
height, let \(u\ge0\) be the maximal excess above \(s\), let
\(\delta\) be the target discard, and let \(d,\allowbreak{} e_0\) be the
original source-discard and diagonal errors. The reconstructed
frame estimate gives, for \emph{every} \(H>0\) with
\(s-\mu_{\min}(RR^\vee)\le H\),
\[
 \mu(RR^\vee)\ge s-
 \frac{4uH/(u+H)+(5\sigma+4d+4\sqrt{e_0})H}{1-\delta}.
                                                               \tag{AU34.3}
\]
For \(u>0\), \(H/(u+H)\le\log(1+H/u)\); for \(u=0\), both
terms multiplied by \(u\) are zero. This proves the literal
\texttt{Frame\_\allowbreak at} with its logarithmic term. No upper restriction
on \(H\) is added. One pair of error bounds is chosen before the
family index and yields \texttt{Frame\allowbreak F} with the original \(C_H\).

GFA10--12 check the actual node codes and intervals, all seventeen
\texttt{Win\_\allowbreak at} clauses, and all four \texttt{Char\_\allowbreak at} clauses.
The primitive-character choices agree at all positive integers,
including ramified ones. Uniqueness uses CRT, multiplicativity and
the minimum conductor at each prime; it needs no theorem asserting
a prime in each arithmetic progression. The three quadratic
readings use the original host's stored \texttt{fl\allowbreak PS} field, with
the correct interchange of the two sign coordinates. Thin-character
and quadratic estimates remain the original chosen inputs.

The bounded residual global task is now explicit: prove the
original \(\texttt{Loss\allowbreak FS}(\mathcal F,0)\), and then carry the
same selected values and characters through the actual
\texttt{link\allowbreak Core}/\texttt{Chain\allowbreak Tied} and final \texttt{Closed}
wrapper. The constants remain
\(C_{\rm abs}=2,\allowbreak{} C_0=3,\allowbreak{} C_5=4,\allowbreak{} e_D=2,\allowbreak{} e_L=0\).
The compatible second-layer selector and local/spectral premises
are not supplied by these consumer proofs.

\paragraph{Real derivatives and one growth exponent for the original theta integral.}
TAE first differentiates the reconstructed Gaussian oscillator
formula in every Heisenberg smooth seminorm. Polynomial Gaussian
domination justifies the derivatives at every order. The actual
full action agrees with that oscillator construction on the
original theta space, with its continuous scalar character.
For an arbitrary adelic left factor, exact Haar \(L^2\) unitarity
transports the derivative; continuity converts equality almost
everywhere to equality at every point. No unproved supremum-norm
boundedness of that left factor is assumed.

A common finite Heisenberg level thickens the real cube into an
adelic set of positive measure. Compactness bounds its rational
translate multiplicity, and unitarity bounds its integral.
Iterating a one-dimensional evaluation inequality then proves
\[
 |R_{\mathbf X}\Phi_L(xg)|\le C_{\mathbf X,\allowbreak{} x}
                       H_T(g)^N,
 \qquad N=\dim_{\mathbb R}\mathrm{RW}(T,\allowbreak{} S),             \tag{AU34.4}
\]
with the \emph{same} \(N\) for every derivative word. The constants
may depend on the word and translate. The source has the original
compact rational representatives from BTH3; a finite coefficient
basis supplies a dominating function and justifies differentiation
under the same integral and Haar measure.

TAE also reconstructs the stored-vacuum passage, using the original
universal finite place-model condition. Adding an algebraically
extended annihilator to the place-model map leaves its defining
finite intertwiners and initial vacuum unchanged. Universality
therefore forces the stored vector to satisfy the real vacuum
equations. The exact commutator computation then identifies the
complex lowering generator with the original mixed Laplacian,
including coefficient \(-4\pi i/s_\tau^2\). Positive and negative
minors in the literal \(\Delta\) use disjoint source columns, so
that Laplacian kills \(\Delta\).

Under the explicitly retained lowest-type smoothness and
irreducibility premises, the source and compact target actions
propagate lowering from \(\Delta\) to every coefficient value.
Combined with FIA/CTC cuspidality and the original finite-level
identities, this gives all eight fields of
\texttt{Is\allowbreak Hol\allowbreak Cusp\allowbreak KMap} for the exact lowest-type integral.
It is not yet the required Hilbert multiplicity map.
The remaining C subcontracts are: one \(io_i\) valid for every
parameter \(t\); \texttt{Low\allowbreak Sm}/\texttt{Low\allowbreak Irr} for the full
original lowest compact type, including multiplicity one;
boundedness and target finite volume; and the original extremal
coefficient/Hilbert realization with Sobolev continuity. Same-Bundle
readings, rational tests and actual base rows remain separate.

\paragraph{Exact reciprocity and the required Count-II normalization.}
CIP matches classical local/global reciprocity to every source
predicate map in the required split-CM scope. Arithmetic reciprocity
sends a uniformizer to arithmetic Frobenius, whereas the source
requires the inverse valuation sign. The proof explicitly inverts
that convention. Global-map uniqueness is checked in every finite
quotient. At a split CM place \(w\ne c(w)\), a local unit in the
global reciprocity kernel is in the finite-idelic closure of the
global units. For any global unit \(e\), \(e/c(e)\) is a root of
unity; its finite possibilities and the coordinates away from
\(w,\allowbreak{} c(w)\) force the original local unit to be one. This proves
the required equality for \emph{every} \texttt{Is\allowbreak Local\allowbreak Artin},
not only a constructed witness. Norm compatibility and the checked
inversion give the source's positive cyclotomic norm sign.
The three category-A application certificates cite Milne,
\emph{Class Field Theory}, version 4.03, with exact theorem locators
and checked number-field/local-field hypotheses. The convention,
uniqueness and character matching are proved separately in CIP.

For the irreducible-point normalization, put
\(A=K[[t]]\) in odd characteristic. An integral unimodular symmetric
form has a vector of unit self-pairing: otherwise polarization,
using the invertibility of two, would contradict unimodularity.
Integral Gram--Schmidt diagonalizes it over \(A\). A finite extension
of the constant field supplies the finitely many residue square
roots, and coefficient recursion supplies square roots of the
resulting unit power series. This gives \(P=1\) by an integral
basis change while keeping \(t\) unchanged.

The transformed representation remains integral. Absolute
irreducibility gives full matrix span; intersections with the
original invariant flag give saturated integral submodules in the
same order. The proof checks the original characters, polynomial
law, determinant, N3 and N4, including the \emph{old} \(D_0\).
Thus the entire universal \texttt{Nice\allowbreak I} instance of \texttt{Norm}
required by JSP11 is supplied, including characteristic three.
The general \texttt{Nice\allowbreak E} normalization is not asserted.
The ordinary arithmetic/control system, four further Hida fields,
same-ideal all-level \texttt{Has\allowbreak Aut}, and the other single-point
support inputs remain C producers. Both original rank-four Links
still use one coefficient threshold at fixed \(\iota\).

\subsection{The selected complete-proof cut and next priorities}

The final interface is unchanged: one compatible design and nine
original \texttt{Closed} predicates. The cutoff and transport tuple
are joined before choosing the common family; eventual thresholds
come before its indices, words and places. RH and two-primary
targets require I/II/V; unrestricted Dirichlet GRH additionally
requires III, SeedIV and the original chains. The two fixed fields
in Counts I and V remain distinct. Keep the exact host window
\(\mathrm{VWinAt}(\chi,401)\), cutoff \(\lfloor x^{400}\rfloor\),
and final numerical exponent \(1/2+9\epsilon\). The AQ alternative
retains its own arithmetic envelope.

The selected graph still has 35 open coupled leaves: 23 C, 10 B
and 2 D. This unchanged count is compatible with substantial
progress inside the three refined leaves. The nested Count-II
reciprocity-understanding issue is now closed; it was not an
additional top-level D leaf. The 121-node graph, 245 historical
targets in 107 groups, thirteen shared refinements and eight
optional routes remain available. None of these counts is a
percentage of proof completion or an absolute minimality claim.

The next bounded global work is the original loss estimate, then
its exact tied-ledger and Closed installation. In parallel, the
common \(ioP\), lowest-type multiplicity-one and Hilbert realization
serve several analytic branches; the ordinary rank-four
arithmetic/control and same-ideal automorphic producer remains an
independent deep obstruction already needed for RH. The recursive
common construction, local \texttt{NBuilt}/\texttt{Nd\allowbreak Spec}, and the
original denominator, retained defect, period-height, local-gain
and spectral producers must stay visible as shared dependencies.

A means a direct published application with hypotheses checked;
B an elementary or moderately difficult independent reconstruction;
C a substantial deep reconstruction; D a source-understanding
issue in the human reconstruction. Completed A applications remain
in the registry. The following register gives every remaining
selected contract verbatim, without replacing common witnesses
by unrelated existential choices.

\begin{description}[style=nextline]
\item[1. GLOBAL\_\allowbreak DENOMINATOR --- C] Actual LinkN local table and binding producer. Closed(Des $\mathcal C$ 3) LinkN: the actual etaSum is bounded with exact dZ0 and same-host vvTerm, subject only to the literal ShaFinOdd and summability guards.

\item[2. GLOBAL\_\allowbreak RETAINED\_\allowbreak DEFECT --- C] Original retained-row defect bound LinkIRR. Rows.DelpNonneg and delp$\le$exp(-x\textasciicircum{}(ez/7)) uniformly for every late actual row datum.

\item[3. GLOBAL\_\allowbreak PERIOD\_\allowbreak HEIGHT --- C] Full actual period and height comparison LinkDIVt. Closed conditional TuSize -> DGamSize -> (LinkEta -> DIVinN -> LinkDIV), with every original LinkDIV\_\allowbreak at clause.

\item[4. GLOBAL\_\allowbreak LOCAL\_\allowbreak GAIN --- C] Original global LinkG installation. One uniform LinkG with actual local lattice lengths, served-orbit coverage and MeanKappaBound; delta(x)->0 before varying actual indices.

\item[5. GLOBAL\_\allowbreak SPECTRAL --- C] Original normalized kernel/Gram estimates LinkS. Closed TuSize -> LinkS with exact Gamma normalization, off-diagonal Hilbert-Schmidt bound, diagonal error and Dall/Dstar estimate, errors exp(-x\textasciicircum{}(epsilon/8)).

\item[6. TASK\_\allowbreak PROP\_\allowbreak GK\_\allowbreak INSTALL --- B] Install OneChar and the congruence integer into PropGK. Complete the finite orbit/\allowbreak Frobenius/\allowbreak valuation argument producing one nonzero integer before every later coefficient prime, with the original e,\allowbreak{} N,\allowbreak{} kernel,\allowbreak{} p and fixed-K constants.

\item[7. TASK\_\allowbreak ONE\_\allowbreak CHAR\_\allowbreak PROOF --- C] Reconstruct the original one-character argument. From exact InpZ and AutomorphyUnramifiedCoefficients, prove the original OneChar trace/length estimate on the supplied coefficient datum and GD input.

\item[8. TASK\_\allowbreak INP\_\allowbreak Z\_\allowbreak ASSEMBLY --- B] Install the exact Kato input record. Build the original CohInp/InpZ record using exactly lemM,\allowbreak{} lemMprime,\allowbreak{} lemP,\allowbreak{} RigII,\allowbreak{} AdVanishK and number on the same arithmetic objects. No rank-four replacement is permitted.

\item[9. RIG\_\allowbreak II --- C] Type-II polarized adjoint rigidity. Marked original datum, actual type-II point and modularity of all source extensions imply vanishing for its exact adRep.

\item[10. AD\_\allowbreak VANISH\_\allowbreak K\_\allowbreak MATCH --- D] Exact rank-three polarized adjoint matching. AdVanishK uses the full negative-transpose complex-conjugation action and irreducibility on the infinite cyclotomic kernel. RX does not transfer to every rank by name.

\item[11. AUT\_\allowbreak UNRAM\_\allowbreak COEF --- C] Actual unramified-coefficient automorphy. For each actual Dat with coefficient ramification one and fixed iota, choose one marked g1 supplying FinA, AutB and LemW together.

\item[12. CNR3 --- C] Uniform rank-three and rank-two congruence integer. Fixed quadratic K -> c,\allowbreak{} A -> e,\allowbreak{} N,\allowbreak{} p,\allowbreak{} w,\allowbreak{} finite Galois-stable roots-of-unity orbit -> nonzero integer NN -> every later coefficient prime, closure, embedding, Frobenius and finite admissible collection.

\item[13. COUNT\_\allowbreak II\_\allowbreak SUPPORT --- C] Two actual rank-four Links with common threshold and same-ideal Hida/P1. For each of the two original QX rank-four Links, at fixed iota choose one common coefficient threshold before every later coefficient DVR and actual universal B/Lambda; produce one ideal J with ring,\allowbreak{} dense,\allowbreak{} pt,\allowbreak{} end1,\allowbreak{} end2,\allowbreak{} twist AND P1core. JSP constructs the exact common Hecke-kernel/annihilator ideal and conditional RingProps/Dense/TWBase and the one-point NiceI transport. CIP now proves exact source reciprocity at the original split CM places for every IsLocalArtin and every admissible global map, with the positive cyclotomic norm sign, and the full universal NiceI instance of Norm used by that transport, preserving the original D0, integral representation, ordered flags and determinant law. Remaining C producers: original ordinary/control/arithmetic-fiber data, the four further Hida fields on this J, all-level same-J automorphic modules/HasAut, and Choice, ArPl, first-order and obstruction bounds, and Raw, with original CohP and stated side conditions. Both Links, ell=3 and common threshold order remain. Generic NiceE Norm is not claimed; the nested reciprocity-understanding obligation is closed in its full required scope.

\item[14. TASK\_\allowbreak CASE3\_\allowbreak AUT\_\allowbreak JOIN --- B] Complete the original Case-III automorphy/image join. Install the actual Case-III ProModAlt, SeedIVAlt and shared PMExAlt with the inherited CohomologyHyp/Ferrero-Washington/HFit/SeedII facts into the original whole Case-III automorphy/image conclusion.

\item[15. TASK\_\allowbreak CASE3\_\allowbreak PROMOD\_\allowbreak JOIN --- C] Reconstruct the actual Case-III chain-to-promodularity proof. Use the original tensor-factor tower, ChainRankFour, ChainDesc4, seeded rank-four chain reduction, actual PMExAlt and SeedIVAlt to construct ProModAlt on every original Ctx3 star point at the same fixed iota, preserving all finite-coefficient, inertia and support maps.

\item[16. SEED\_\allowbreak IV\_\allowbreak ALT --- C] Actual Case-III seed with prechosen residual datum. Base Ctx3/Four tower, iota and Amb -> w,\allowbreak{} s before every later Dies tower -> compatible finite-coefficient seed with exact residual support/root entries.

\item[17. TENSOR\_\allowbreak FACTOR\_\allowbreak TOWER --- B] Actual tensor-factor tower and characters. Original finite-field, tame, factor and character constraints; source-level construction is sizeable and has not been reconstructed here.

\item[18. CHAIN\_\allowbreak RANK\_\allowbreak FOUR --- C] Actual SomeIn-to-AllIn propagation. Same source Space and PM ideal, residual point Good4, original bound 30(C+1)\textasciicircum{}2 <= d; geometry must preserve all identities.

\item[19. CHAIN\_\allowbreak DESC4 --- B] Original residual chain descent. Finite residual descent retaining source Chain4 compatibility.

\item[20. RANK4\_\allowbreak CHAIN\_\allowbreak REDUCTION --- C] Seeded reduction to the original residual chain. Select T1 before every later Deep tower; retain one compatible seed and Chain4 for every original residual lift over Q.

\item[21. TASK\_\allowbreak COUNT\_\allowbreak V\_\allowbreak INSTALL --- B] Complete the original CountV finite assembly. From exactly the listed source host/\allowbreak cohomology inputs, complete countV\_\allowbreak of\_\allowbreak finiteness\_\allowbreak duality\_\allowbreak euler\_\allowbreak host, including original host/frame formation, normalized length, the CaseV guard and R1<=2\textasciicircum{}c2.

\item[22. THM\_\allowbreak S6\_\allowbreak S --- C] Original unitary-three ordinary support and automorphy. Original host/frame and fixed iota -> one marked g1 -> every universal R,\allowbreak{} u -> weight-zero finiteness and every Type-I BPrime\_\allowbreak s point.

\item[23. HOST\_\allowbreak FRAME\_\allowbreak LENGTH --- B] Original Case-V host/frame and normalized-length installation. Actual HostData, frame, local-condition equality, abstract count and normalized original Selmer length must all refer to the same objects.

\item[24. HOST\_\allowbreak RIGID\_\allowbreak S --- D] Rigidity on the original CountV frame. Same original host/frame, marked point and eta-kernel; published adjoint vanishing is usable only after exact representation and hypothesis matching.

\item[25. HOST\_\allowbreak NUMBER\_\allowbreak S --- C] Uniform original-host rank-three catalogue. One exponent before all closures, coefficient rings, levels, characters, frames, embeddings and marked points; one separating trace function controls all later finite pairwise nonisomorphic collections.

\item[26. TASK\_\allowbreak HOST\_\allowbreak VAN\_\allowbreak E\_\allowbreak INSTALL --- B] Transport K/duality to literal HostVanE. Use the already supplied LemmaK/local-global duality and finiteness on the actual CountV host/frame and exact original coefficient modules.

\item[27. TASK\_\allowbreak HOST\_\allowbreak XPP\_\allowbreak INSTALL --- B] Transport CorX to literal HostXpp. Preserve Psi\textasciicircum{}-2 eta, the prime-to-ell quotient, actual local conditions and the original CountV host while installing BR7.

\item[28. DESIGN\_\allowbreak CHOICES\_\allowbreak COMMON --- C] Construct the original common design family and bundle. For the one target HostExW, construct the original compatible DesignChoices on Cdes=WithTw(WithSm Cprime), where Cprime=WithChN(A intersect target class), A is the one Built construction class, and cT(eps,\allowbreak{} y)=cTmin(y). Retain every original field, tuple, word, metric and common bundle duty. DSZ10 and RMP11 prove the size/rank properties of members of this class; they do not produce a member. The RecExP/RecAtG and common object/row route remains a coupled construction.

\item[29. LOCAL\_\allowbreak NBUILT\_\allowbreak COMMON --- C] Actual local specialization on the common selected family. Closed on the same Des(C,3) class: one eventual family of actual NBuilt certificates for the original Fm.D(x,\allowbreak{} i), including NdSpec, original iV/Hecke actions, every required maximal ideal, $\Psi$/range and projection square. Keep the constructor certificate before forgetting it to TSel.

\item[30. SECOND\_\allowbreak LAYER\_\allowbreak COMPATIBLE\_\allowbreak SELECTION --- C] Original second-layer selector and class duties. For the same A, supply Duty on A and a Closed common original TSel family whose selected records are the original true-column rec7 records (or have the entire dependent tuple transported), with their real witness and original coefficient-field root data. Existence/compatibility remains arithmetic; merely abstractly isomorphic fields or independently selected columns do not suffice.

\item[31. TASK\_\allowbreak GL\_\allowbreak FINAL\_\allowbreak ASSEMBLY --- B] Full final global numerical and lattice assembly. On the SAME original supplied family, compatible second-layer selector and Duty, GLS13 gives the full eleven-field FamGLmin(D3(3,\allowbreak{} A),4). GVC now supplies literal CountF and SlopesMF, using the logarithmic-derivative size already projected from FamP.sizes; GFA supplies LevelTFamily with constant36, PlacesF/SlopesNF with original metrics/support, positive RRRankF/LatF, the literal all-H FrameF, NodesF, every Win\_\allowbreak at clause and exact CharF. Remaining B work: prove original LossFS(Fm,0) from the chosen small shares and weights, then install the actual same-family linkCore/ChainTied ledger and final Closed wrapper with Cabs=2,\allowbreak{} C0=3,\allowbreak{} C5=4,\allowbreak{} eD=2,\allowbreak{} eL=0. Keep original top/bottom values, characters, all ledger fields and thresholds. The epsilon cutoff and TransportedAlt tuple precede Fm,\allowbreak{} CH; the eventual x threshold precedes i,\allowbreak{} u and places. Compatible family/selector existence and LinkG/LinkS/LinkIRR/LinkDIV remain upstream producers.

\item[32. BASE\_\allowbreak THETA\_\allowbreak READINGS --- C] Actual base reading on the common bundle. For every family in the original class, eventually on the SAME selected Bundle and for each i<m, supply the actual R,\allowbreak{} F,\allowbreak{} Read,\allowbreak{} ReadL, weak models/Weil data and phase-one pins. Identify their LiftTie.theta, mtRead, mjRead and ThDef3 basis with that Bundle.P.A3, P.Zs, P.S0 and its stored maps. Exact field, embeddings, metrics, measures, row family and weak choices are retained. This is one part of the coupled BTH12 producer, not a freely replaceable isomorphic model.

\item[33. BASE\_\allowbreak THETA\_\allowbreak CUSPIDAL\_\allowbreak EXTENSION --- C] Extend the original convergent base integrals. On the SAME selected Bundle and SAME reading as BASE\_\allowbreak THETA\_\allowbreak READINGS, extend the exact lowest-type theta integral for each finite basis vector of HV(i,\allowbreak{} t0) to the original Wall\_\allowbreak i, with one family threshold before i and basis vectors. BTH3/FIA/CTC prove convergence, finite cancellation and genuine cuspidality. TAE1–13 now prove actual-action smoothness, uniform growth with one exponent for every derivative word, dominated differentiation, stored-vacuum transfer, exact lowering, finite-level and compact-equivariance installation; given the following producer premises this gives the literal IsHolCuspKMap. Remaining C subcontracts: one original io\_\allowbreak i satisfying IsIoW for EVERY parameter t; LowSm and LowIrr for the full original lowKSub (including multiplicity one); boundedness and finite-volume target input; the exact extremal coefficient identity and Sobolev-continuous Hilbert multiplicity-map realization agreeing with the given integral. Keep all original fields, adjoints, levels and common choices. Separate geometric readings, rational tests and actual base rows are not supplied by this analytic packet.

\item[34. BASE\_\allowbreak THETA\_\allowbreak RATIONAL\_\allowbreak TESTS --- C] Simultaneous original base rational tests. On the SAME original Bundle/reading/theta map, for every i<m supply one finite separating complex-linear test family. Its values must lie in the fixed embedded coefficient field on theta of EVERY original rational source vector and on EVERY stored target-coordinate basis vector. BTH8 and BTH11 use this one base LiftTestB family for rational adjoint image and matrix coordinates. Finite dimensionality alone does not supply these rationality identities.

\item[35. SLR\_\allowbreak BASE\_\allowbreak ROWS\_\allowbreak PRODUCER --- C] Actual base row identities on that same bundle. For the SAME selected Bundle, common original fields and readings used by the other three BTH12 inputs, supply actual RowsHold0\_\allowbreak 3 uniformly eventually for every family. Its original primitive row realization, spectrum/type and basis identifications must be retained. BTH supplies a theta consumer only; it does not construct the base row identity. No separate witness may be chosen merely to satisfy this row requirement.

\end{description}

\paragraph{Validation and continuity.}
The release preserves predecessor claim fields, source arrays and
mathematical dependency edges. New source and proof ranges are bound
to the pinned commit and independently reviewed. The release
validator checks all inherited source ranges, graph connectivity,
exact remaining contracts, review hashes and draft TeX references.
These checks are not a new Lean replay. Historical RF16 remains
governed by its explicit DEX11 correction; no historical proof
bytes are silently rewritten. The current audit and additive scope
updates govern superseded historical frontier statements.


\paragraph{Historical version-0.33 audit.}
The following audit is preserved as a historical checkpoint.
Section~\ref{sec:dependency-audit-v34} and the version-0.34 scope
updates give the current remaining obligations. The historical
conditional mathematics retains its original hypotheses.

\section{Current dependency audit: original cuspidality, column slopes, and the same Hecke ideal}
\label{sec:dependency-audit-v33}

\status{Version 0.33 accepts the underlying Alp\"oge--Claude formal
proof as fully checked and correct. The new human work proves the
original finite isotropic averaging argument and genuine theta
cuspidality, all three GL4 implications and their FamGLmin assembly,
and the exact Hecke-kernel and single-point support transports.
The human reconstruction is still incomplete. The current register
below states every remaining selected obligation, with its original
objects and quantifiers; category D concerns our understanding of
the formal source, not a claim that the checked proof has a gap.}

\subsection{What this continuation actually supplies}

\paragraph{One finite average on the original kernels.}
FIA1--FIA8 use the original real datum, character, finite frames,
weighted lattice, residue maps, Fock seed and Weil covariance.
Choose once a prime $\mathfrak P$ in the stored nonempty good-prime
set and put $N=\#k(\mathfrak P)>0$. For each original index and
parameter, adelic pair $(h,\allowbreak{} g)$ and nonzero rational vector $\xi$,
there is a list $a_0,\ldots,\allowbreak{} a_{N-1}$ such that
\[
 \sum_{j=0}^{N-1}H_{\xi a_j}\omega(h,\allowbreak{} g)\operatorname{pt}(P)=0
                         \quad\text{for every polynomial }P.
 \tag{AU33.1}
\]
The period is fixed for the datum, and the list is chosen before
$P$. The list may depend on the other parameters. The finite
argument requires only $\xi\ne0$; isotropy is needed in the
constant-term argument which follows. The local phase uses the full
oriented symplectic form, while a raw Heisenberg operator uses its
half. FIA checks the sign and both half-datum guards explicitly.

\paragraph{Convergent constant terms and actual cuspidality.}
For an isotropic $\xi$ write $Z_\xi$ for the original rank-one
transvection group. CTC1 identifies it with $\mathbb A_F$, with
rational subgroup $F$, using $1+(t\xi_0)\xi\xi^*G_T$.
CTC2 supplies the uniform bound on compact Weil orbits from a
common finite level, exact Haar $L^2$ unitarity and an iterated
one-variable Sobolev estimate. CTC3--CTC7 reconstruct the local
character duality, algebraic split-half separation and positive
action on the along coinvariants. For a nonzero rational row the
positive definite source supplies a nonzero Hermitian norm; a
transvection then gives a unit phase different from one times a
positive real scalar. That product cannot be one.

Fourier coefficients on the compact along quotient therefore show
that the exact transvection constant-term functional satisfies
\[
 \mathcal C_\xi(H_{\xi a}f)=\mathcal C_\xi(f),\qquad
 \mathcal C_\xi(f)=\int_{Z_\xi(F)\backslash Z_\xi(\mathbb A)}
                                     \omega(1,\allowbreak{} z)f(0)\,\allowbreak{} dz.
 \tag{AU33.2}
\]
Applying it to (AU33.1) gives $N\mathcal C_\xi(f)=0$.
The original definite source has compact rational representatives
from its original finite class presentation. A finite coefficient
basis makes the two-variable kernel continuous and bounded there.
Thus Fubini is justified and the actual theta integral has zero
transvection constants.

CTC11 writes each original isotropic-frame unipotent group in Witt
coordinates as a two-step group of matrices $n(A,\allowbreak{} B)$, with $B$
skew-Hermitian. Its quotient by rational points is compact, by
successively reducing $A$ and the central $B$ coordinate in adelic
vector quotients. A transvection along any frame column is central.
Haar invariance and Fubini now force the full unipotent constant
term to vanish. Hence CTC13 proves genuine cuspidality of the
original base integral. It does not use the formal convention
assigning zero to a nonintegrable Bochner integral.

The remaining base extension is smaller but still substantial:
smoothness and uniform derivative growth, lowering and prescribed
lowest type, original finite-level/equivariance installation, and
bounded holomorphic realization in the specified Hilbert
multiplicity model. Separately, one original $\mathrm{ioP}$ must retain all
its parameter conditions; this is not the smooth Hilbert extension itself. Same-Bundle geometric readings,
simultaneous rational tests and actual base rows remain separate
coupled producers.

\paragraph{The three original GL4 slots.}
GLS1--GLS3 construct the literal column-value record from the same
supplied forms: original class lattices, conjugate coefficient
embeddings, local polynomial membership at every finite place,
and the exact infinite-place norm. For every supplied ColVal,
GLS6 proves
\[
                     \mu(Q_{c,\tau})=-\frac{\log T}{k}.
 \tag{AU33.3}
\]
This uses the original squared Hermitian slope, its factor two,
the discriminant index $T^{2f}$, and scalar extension without a
hidden finite-index change. GLS4 proves the exact highest-vector
Fock norm by a finite normal-order recurrence; GLS5 uses normalized
Weyl/coefficient lattice duality and Cauchy--Binet for its integral
determinant estimate. These give the original lower slope bound.

For the upper slope, GLS8 applies Banaszczyk's published Euclidean
transference theorem, Theorem 2.1 in \emph{Mathematische Annalen}
296 (1993), 625--635, on the actual positive trace lattice.
Its full rank, dimension, metric and trace/codifferent dual are
checked. The subsequent original class-map and coefficient-norm
arguments give the literal constant $4$, not an unspecified
constant. The primary source and complete hypothesis check are in
\texttt{global\_\allowbreak slope/\allowbreak published\_\allowbreak applications.\allowbreak json}.

GLS13 assembles all eleven original fields
\[
 h_{2b},\allowbreak{} h_{2v},\allowbreak{} h_3,\allowbreak{} h_{Za},\allowbreak{} h_{Zb},\allowbreak{} t_{CV},\allowbreak{} h_W,\allowbreak{} h_\mu,\allowbreak{} h_I,\allowbreak{} h_R,\allowbreak{} h_B
\]
of $\mathrm{FamGLmin}(D_3(3,\mathcal A),4)$ on the same supplied
second-layer selector and Duty. Earlier RSL/GLE/DSZ/RMP arguments
provide the other fields in their exact scopes. The common second
layer still has to be produced. The final numerical/lattice
implication through \texttt{link\allowbreak GLOBAL7\_\allowbreak min/\allowbreak items} and
\texttt{link8} still has to be written. It retains
$C_{\rm abs}=2,\allowbreak{} C_0=3,\allowbreak{} C_5=4,\allowbreak{} e_D=2,\allowbreak{} e_L=0$.

\paragraph{Count II keeps one ideal throughout.}
For each of the two original rank-four Links, JSP constructs the
map from the given universal ordinary ring to the original Hecke
algebra. It writes explicit preimages of every spherical and
ordinary Hecke generator, proving surjectivity. Therefore
\[
          J=\ker(B\longrightarrow\mathbb T)
                        =\operatorname{Ann}_B(M),\qquad
                         B/J\simeq\mathbb T
 \tag{AU33.4}
\]
in the nonzero-module case; if $M=0$ the exact source alternative
is $J=B$. JSP verifies all eight RingProps fields using coefficientwise
polynomial-law tests at separating arithmetic fibers. It proves
the Dense implication with the exact weight and reciprocity
dictionary kept explicit. It also constructs the actual TWBase on
this same module and ideal.

The characteristic-prime argument is the source's one-point NiceI
case. Injective local coefficient extensions preserve kernels;
the normalized full matrix span supplies a nonzero trace
discriminant in every characteristic. Under HasAut and the six
retained arithmetic/patching contracts, the source one-point
support conclusion then gives P1core for the same $J$.
This reconstructs the transport, not the substantial patching
input itself. It works for the allowed coefficient prime $3$.

Still to produce are the original ordinary arithmetic/control data
and separating fibers, the precise reciprocity matching, the four
further Hida endpoints, all-level automorphic modules/HasAut, and
the retained single-point support inputs. The fixed $\iota$
comes first; one coefficient threshold for both Links precedes
every later DVR and every original universal $B,\Lambda$.
The seven properties must hold on each Link's same $J$.
The exact Artin-map/\texttt{Is\allowbreak Wt\allowbreak Part} matching is a category D
subobligation within the coupled category C producer. The count
of two D leaves below refers to the selected graph granularity;
it does not suppress this nested source-understanding task.

\subsection{The selected complete-proof cut}

The final interface remains one compatible DesignExists and the
nine original Closed predicates: LinkN, LinkEta, LinkIRR,
SmallLevelRankEq, the TuSize/DGamSize implication to LinkDIVt,
the ChosenE implication to NuFullF and MuRank, LinkG, the TuSize
implication to LinkS, and the TuSize/DGamSize implication to
LinkGLOBAL. The Closed epsilon cutoffs and transport tuple are joined before
selecting one family. Eventual x thresholds may depend on that
supplied family and its class witness, but precede the indices,
places and other variables over which the bound is uniform. The construction-class TuSize and RankM proofs from
version 0.32 remain available; they do not create that class's
compatible member.

RH and the two-primary specialization still use I, II and V.
Unrestricted Dirichlet GRH additionally uses III, SeedIV and the
original chain. The Count-I and Count-V fields and prime scopes
remain distinct; the Count-V twist exponent is not the bounded
host dyadic exponent. The host still uses the original
$\mathrm{VWinAt}(\chi,401)$ and cutoff $\lfloor x^{400}\rfloor$;
the final reconstructed numerical exponent remains
$1/2+9\varepsilon$. The AQ route is an alternative conditional
route with its own arithmetic envelope.

Each remaining selected leaf is classified A, B, C or D. A means
a direct published application with hypotheses checked. Completed
A applications are retained in their registry and are not counted
as unresolved leaves. B means an elementary or moderately difficult
independent reconstruction. C means a substantial deep reconstruction.
D means a possible gap in our human understanding of the source.
A completed conditional argument can have open C producers;
its completion does not assert their existence.

The register now has 35 open leaves: 23 C, 10 B and 2 D.
The three GL4 implication leaves are closed; the theta and Count-II
leaves are narrowed. These are coupled sufficient contracts, not
equal-sized tasks or a percentage of proof completion. All 245
inherited compound targets in 107 groups remain dispositioned.
The 121-node selected graph is rooted in RH and Dirichlet GRH;
13 shared work refinements and eight optional routes remain
separate from its logical AND rules.

\subsection{Priority by contribution to the final theorem}

The next bounded global task is the residual numerical/lattice
assembly from the now complete same-selector FamGLmin package.
It should be reconstructed directly on the source's existing
frame, HN and loss variables, with all constants displayed.
The main shared analytic task is the smooth-growth/lowering and
Hilbert realization of the now cuspidal original theta integral.
This serves the common base reading, small-rank and later
period/spectral branches. The main independent arithmetic task
is the original rank-four ordinary control/arithmetic producer
and all-level same-ideal support construction. JSP identifies
which algebraic transports no longer need to be repeated.

The common recursive/geometric construction, local NBuilt/NdSpec,
the original denominator, retained defect, period, local-gain and
spectral producers remain on the path to both roots. Keep their
bindings visible while working on a bounded subtask. The next
branch should be selected for its effect on these root
dependencies, not for its position at the end of the latest chapter.

\subsection{Every remaining selected obligation}
The following contracts are reproduced from the machine-readable
current register. Their order groups the source's interfaces; it
does not assign equal cost or independent witnesses to them.

\begin{description}[style=nextline]
\item[1. GLOBAL\_\allowbreak DENOMINATOR --- C] Actual LinkN local table and binding producer. Closed(Des $\mathcal C$ 3) LinkN: the actual etaSum is bounded with exact dZ0 and same-host vvTerm, subject only to the literal ShaFinOdd and summability guards.

\item[2. GLOBAL\_\allowbreak RETAINED\_\allowbreak DEFECT --- C] Original retained-row defect bound LinkIRR. Rows.DelpNonneg and delp$\le$exp(-x\textasciicircum{}(ez/7)) uniformly for every late actual row datum.

\item[3. GLOBAL\_\allowbreak PERIOD\_\allowbreak HEIGHT --- C] Full actual period and height comparison LinkDIVt. Closed conditional TuSize -> DGamSize -> (LinkEta -> DIVinN -> LinkDIV), with every original LinkDIV\_\allowbreak at clause.

\item[4. GLOBAL\_\allowbreak LOCAL\_\allowbreak GAIN --- C] Original global LinkG installation. One uniform LinkG with actual local lattice lengths, served-orbit coverage and MeanKappaBound; delta(x)->0 before varying actual indices.

\item[5. GLOBAL\_\allowbreak SPECTRAL --- C] Original normalized kernel/Gram estimates LinkS. Closed TuSize -> LinkS with exact Gamma normalization, off-diagonal Hilbert-Schmidt bound, diagonal error and Dall/Dstar estimate, errors exp(-x\textasciicircum{}(epsilon/8)).

\item[6. TASK\_\allowbreak PROP\_\allowbreak GK\_\allowbreak INSTALL --- B] Install OneChar and the congruence integer into PropGK. Complete the finite orbit/\allowbreak Frobenius/\allowbreak valuation argument producing one nonzero integer before every later coefficient prime, with the original e,\allowbreak{} N,\allowbreak{} kernel,\allowbreak{} p and fixed-K constants.

\item[7. TASK\_\allowbreak ONE\_\allowbreak CHAR\_\allowbreak PROOF --- C] Reconstruct the original one-character argument. From exact InpZ and AutomorphyUnramifiedCoefficients, prove the original OneChar trace/length estimate on the supplied coefficient datum and GD input.

\item[8. TASK\_\allowbreak INP\_\allowbreak Z\_\allowbreak ASSEMBLY --- B] Install the exact Kato input record. Build the original CohInp/InpZ record using exactly lemM,\allowbreak{} lemMprime,\allowbreak{} lemP,\allowbreak{} RigII,\allowbreak{} AdVanishK and number on the same arithmetic objects. No rank-four replacement is permitted.

\item[9. RIG\_\allowbreak II --- C] Type-II polarized adjoint rigidity. Marked original datum, actual type-II point and modularity of all source extensions imply vanishing for its exact adRep.

\item[10. AD\_\allowbreak VANISH\_\allowbreak K\_\allowbreak MATCH --- D] Exact rank-three polarized adjoint matching. AdVanishK uses the full negative-transpose complex-conjugation action and irreducibility on the infinite cyclotomic kernel. RX does not transfer to every rank by name.

\item[11. AUT\_\allowbreak UNRAM\_\allowbreak COEF --- C] Actual unramified-coefficient automorphy. For each actual Dat with coefficient ramification one and fixed iota, choose one marked g1 supplying FinA, AutB and LemW together.

\item[12. CNR3 --- C] Uniform rank-three and rank-two congruence integer. Fixed quadratic K -> c,\allowbreak{} A -> e,\allowbreak{} N,\allowbreak{} p,\allowbreak{} w,\allowbreak{} finite Galois-stable roots-of-unity orbit -> nonzero integer NN -> every later coefficient prime, closure, embedding, Frobenius and finite admissible collection.

\item[13. COUNT\_\allowbreak II\_\allowbreak SUPPORT --- C] Two actual rank-four Links with common threshold and same-ideal Hida/P1. For each of the two original QX rank-four Links, at fixed iota choose one common coefficient threshold before every later coefficient DVR and every actual universal B/Lambda; produce one ideal J with ring,\allowbreak{} dense,\allowbreak{} pt,\allowbreak{} end1,\allowbreak{} end2,\allowbreak{} twist AND P1core. JSP constructs J=ker(B$\to$T)=Ann\_\allowbreak B(M), proves ring/dense and the same-ideal TWBase from explicitly supplied original ordinary Hecke arithmetic data, and proves the one-point NiceI$\to$P1 implication from HasAut and six exact arithmetic/patching contracts. Remaining producers are that ordinary/control/arithmetic-fiber system including its exact reciprocity dictionary, the four further Hida fields on this J, the all-level automorphic modules/HasAut, and the retained single-point support inputs. The exact artOf/every-map IsWtPart reciprocity matching is a nested D source-understanding obligation inside this coupled C producer; KX3 alone does not supply it. The common ideal, ell=3 scope, both Links and threshold order are unchanged.

\item[14. TASK\_\allowbreak CASE3\_\allowbreak AUT\_\allowbreak JOIN --- B] Complete the original Case-III automorphy/image join. Install the actual Case-III ProModAlt, SeedIVAlt and shared PMExAlt with the inherited CohomologyHyp/Ferrero-Washington/HFit/SeedII facts into the original whole Case-III automorphy/image conclusion.

\item[15. TASK\_\allowbreak CASE3\_\allowbreak PROMOD\_\allowbreak JOIN --- C] Reconstruct the actual Case-III chain-to-promodularity proof. Use the original tensor-factor tower, ChainRankFour, ChainDesc4, seeded rank-four chain reduction, actual PMExAlt and SeedIVAlt to construct ProModAlt on every original Ctx3 star point at the same fixed iota, preserving all finite-coefficient, inertia and support maps.

\item[16. SEED\_\allowbreak IV\_\allowbreak ALT --- C] Actual Case-III seed with prechosen residual datum. Base Ctx3/Four tower, iota and Amb -> w,\allowbreak{} s before every later Dies tower -> compatible finite-coefficient seed with exact residual support/root entries.

\item[17. TENSOR\_\allowbreak FACTOR\_\allowbreak TOWER --- B] Actual tensor-factor tower and characters. Original finite-field, tame, factor and character constraints; source-level construction is sizeable and has not been reconstructed here.

\item[18. CHAIN\_\allowbreak RANK\_\allowbreak FOUR --- C] Actual SomeIn-to-AllIn propagation. Same source Space and PM ideal, residual point Good4, original bound 30(C+1)\textasciicircum{}2 <= d; geometry must preserve all identities.

\item[19. CHAIN\_\allowbreak DESC4 --- B] Original residual chain descent. Finite residual descent retaining source Chain4 compatibility.

\item[20. RANK4\_\allowbreak CHAIN\_\allowbreak REDUCTION --- C] Seeded reduction to the original residual chain. Select T1 before every later Deep tower; retain one compatible seed and Chain4 for every original residual lift over Q.

\item[21. TASK\_\allowbreak COUNT\_\allowbreak V\_\allowbreak INSTALL --- B] Complete the original CountV finite assembly. From exactly the listed source host/\allowbreak cohomology inputs, complete countV\_\allowbreak of\_\allowbreak finiteness\_\allowbreak duality\_\allowbreak euler\_\allowbreak host, including original host/frame formation, normalized length, the CaseV guard and R1<=2\textasciicircum{}c2.

\item[22. THM\_\allowbreak S6\_\allowbreak S --- C] Original unitary-three ordinary support and automorphy. Original host/frame and fixed iota -> one marked g1 -> every universal R,\allowbreak{} u -> weight-zero finiteness and every Type-I BPrime\_\allowbreak s point.

\item[23. HOST\_\allowbreak FRAME\_\allowbreak LENGTH --- B] Original Case-V host/frame and normalized-length installation. Actual HostData, frame, local-condition equality, abstract count and normalized original Selmer length must all refer to the same objects.

\item[24. HOST\_\allowbreak RIGID\_\allowbreak S --- D] Rigidity on the original CountV frame. Same original host/frame, marked point and eta-kernel; published adjoint vanishing is usable only after exact representation and hypothesis matching.

\item[25. HOST\_\allowbreak NUMBER\_\allowbreak S --- C] Uniform original-host rank-three catalogue. One exponent before all closures, coefficient rings, levels, characters, frames, embeddings and marked points; one separating trace function controls all later finite pairwise nonisomorphic collections.

\item[26. TASK\_\allowbreak HOST\_\allowbreak VAN\_\allowbreak E\_\allowbreak INSTALL --- B] Transport K/duality to literal HostVanE. Use the already supplied LemmaK/local-global duality and finiteness on the actual CountV host/frame and exact original coefficient modules.

\item[27. TASK\_\allowbreak HOST\_\allowbreak XPP\_\allowbreak INSTALL --- B] Transport CorX to literal HostXpp. Preserve Psi\textasciicircum{}-2 eta, the prime-to-ell quotient, actual local conditions and the original CountV host while installing BR7.

\item[28. DESIGN\_\allowbreak CHOICES\_\allowbreak COMMON --- C] Construct the original common design family and bundle. For the one target HostExW, construct the original compatible DesignChoices on Cdes=WithTw(WithSm Cprime), where Cprime=WithChN(A intersect target class), A is the one Built construction class, and cT(eps,\allowbreak{} y)=cTmin(y). Retain every original field, tuple, word, metric and common bundle duty. DSZ10 and RMP11 prove the size/rank properties of members of this class; they do not produce a member. The RecExP/RecAtG and common object/row route remains a coupled construction.

\item[29. LOCAL\_\allowbreak NBUILT\_\allowbreak COMMON --- C] Actual local specialization on the common selected family. Closed on the same Des(C,3) class: one eventual family of actual NBuilt certificates for the original Fm.D(x,\allowbreak{} i), including NdSpec, original iV/Hecke actions, every required maximal ideal, $\Psi$/range and projection square. Keep the constructor certificate before forgetting it to TSel.

\item[30. SECOND\_\allowbreak LAYER\_\allowbreak COMPATIBLE\_\allowbreak SELECTION --- C] Original second-layer selector and class duties. For the same A, supply Duty on A and a Closed common original TSel family whose selected records are the original true-column rec7 records (or have the entire dependent tuple transported), with their real witness and original coefficient-field root data. Existence/compatibility remains arithmetic; merely abstractly isomorphic fields or independently selected columns do not suffice.

\item[31. TASK\_\allowbreak GL\_\allowbreak FINAL\_\allowbreak ASSEMBLY --- B] Full final global numerical and lattice assembly. GLS13 assembles all eleven original FamGLmin(D3(3,\allowbreak{} A),4) fields on the SAME supplied second-layer selector and Duty, using the proved GL4/GL6 inputs. What remains is the human numerical/lattice implication linkGLOBAL7\_\allowbreak min/items and link8 on that same class, from these fields, original congruence length, ChosenWall/Duty and TuSize, with Cabs=2,\allowbreak{} C0=3,\allowbreak{} C5=4,\allowbreak{} eD=2,\allowbreak{} eL=0. Preserve the actual frame/HN and loss inequalities. The Closed epsilon cutoff and TransportedAlt tuple precede the supplied Fm,\allowbreak{} CH; the eventual x threshold may depend on that Fm,\allowbreak{} CH but precedes i,\allowbreak{} u and every place. FamGLmin existence is conditional on the common upstream second layer.

\item[32. BASE\_\allowbreak THETA\_\allowbreak READINGS --- C] Actual base reading on the common bundle. For every family in the original class, eventually on the SAME selected Bundle and for each i<m, supply the actual R,\allowbreak{} F,\allowbreak{} Read,\allowbreak{} ReadL, weak models/Weil data and phase-one pins. Identify their LiftTie.theta, mtRead, mjRead and ThDef3 basis with that Bundle.P.A3, P.Zs, P.S0 and its stored maps. Exact field, embeddings, metrics, measures, row family and weak choices are retained. This is one part of the coupled BTH12 producer, not a freely replaceable isomorphic model.

\item[33. BASE\_\allowbreak THETA\_\allowbreak CUSPIDAL\_\allowbreak EXTENSION --- C] Extend the original convergent base integrals. On the SAME selected Bundle and SAME reading as BASE\_\allowbreak THETA\_\allowbreak READINGS, extend the exact lowest-type theta integral for each finite basis vector of HV(i,\allowbreak{} t0) to the original Wall\_\allowbreak i with one family threshold before i and the basis vectors. BTH3 proves convergence; FIA8 and CTC1–CTC13 now prove the actual finite average, ZCusp and genuinely convergent full cuspidality. The remaining work is original smoothness and uniform derivative-growth estimates, lowering and prescribed lowest type, finite-level/equivariance installation and Sobolev-continuous bounded holomorphic/Hilbert multiplicity realization. Retain one original ioP with its universal parameter conditions; base-only existence is insufficient.

\item[34. BASE\_\allowbreak THETA\_\allowbreak RATIONAL\_\allowbreak TESTS --- C] Simultaneous original base rational tests. On the SAME original Bundle/reading/theta map, for every i<m supply one finite separating complex-linear test family. Its values must lie in the fixed embedded coefficient field on theta of EVERY original rational source vector and on EVERY stored target-coordinate basis vector. BTH8 and BTH11 use this one base LiftTestB family for rational adjoint image and matrix coordinates. Finite dimensionality alone does not supply these rationality identities.

\item[35. SLR\_\allowbreak BASE\_\allowbreak ROWS\_\allowbreak PRODUCER --- C] Actual base row identities on that same bundle. For the SAME selected Bundle, common original fields and readings used by the other three BTH12 inputs, supply actual RowsHold0\_\allowbreak 3 uniformly eventually for every family. Its original primitive row realization, spectrum/type and basis identifications must be retained. BTH supplies a theta consumer only; it does not construct the base row identity. No separate witness may be chosen merely to satisfy this row requirement.

\end{description}

\paragraph{Validation and continuity.}
The release preserves predecessor claim fields, source arrays and
mathematical dependency edges. New source and proof ranges are bound
to the pinned commit and independently reviewed. The release
validator checks all inherited source ranges, graph connectivity,
exact remaining contracts, review hashes and draft TeX references.
These checks are not a new Lean replay. Historical RF16 remains
governed by its explicit DEX11 correction; no historical proof
bytes are silently rewritten. The current audit and additive scope
updates govern superseded historical frontier statements.


\paragraph{Historical version-0.32 audit.}
The following audit is retained as a historical checkpoint.
Section~\ref{sec:dependency-audit-v33} and the version-0.33 additive
scope updates specify the current open obligations. The historical
conditional arguments retain their original hypotheses.

\section{Current dependency audit: actual size and rank, local operators, and the base theta frontier}
\label{sec:dependency-audit-v32}

\status{Version 0.32 accepts the underlying Alp\"oge--Claude Lean
proof as fully checked and correct. The new work closes the original
construction-class size and rank obligations and the specified global
local calculations. It also narrows the actual base-theta construction.
The human reconstruction remains incomplete. Every remaining selected
obligation is classified below; an unresolved human reading is not a
claim of a gap in the formal proof.}

\subsection{What is now proved on the original objects}

The most useful change is in design existence. DEX7 had reduced it
to three compatible inputs on one class: original \texttt{Design\allowbreak Choices},
the closed \texttt{Tu\allowbreak Size} implication, and the closed original
\texttt{Rank\allowbreak M} formula. DSZ10 and RMP11 now prove the latter two
on the exact construction class. Thus only the actual common-family
construction remains inside this selected design cut. This is a
statement about the dependencies of the human proof, not an assertion
that its most difficult construction has become elementary.

The class is essential. Let
\[
 \mathcal A=\texttt{cls}(Q,\mathrm{Lk},\allowbreak{} E),\quad
 \mathcal C'=\texttt{With\allowbreak Ch\allowbreak N}(\mathcal A\cap\mathcal G),\quad
 \mathcal C_{\rm des}=\texttt{With\allowbreak Tw}(\texttt{With\allowbreak Sm}(\mathcal C')),
                                                               \tag{AU32.1}
\]
where $\mathcal G$ is the actual target class. Membership in
$\mathcal A$ contains the original \texttt{Link} and \texttt{QOf}
property. Its positivity field permits \texttt{cls\_\allowbreak at}; the
source's \texttt{pick} then gives the very Bundle whose $R$ field
supplies the columns. The same \texttt{Built} witness contains the
actual real-place data. DSZ and RMP use these projections on the
same $D$, level, field, weights, metrics and stored degree variables.
They prove properties of every member of $\mathcal A$.
They do not prove that this class has a member satisfying all the
other required duties.

\paragraph{The degree and prime-size calculation.}
The new finite product comparison is
\[
 0<d_n(\lambda)\le d_\tau(\lambda)
         \left(\max_{a,\allowbreak{} b}|\lambda_a-\lambda_b|\right)^{pq}.
                                                               \tag{AU32.2}
\]
DSZ2 checks that the source's quantitative gap exceeds the actual
general construction and rigidity threshold in DC1--DC13. DC18--DC19,
with the previously matched Harish--Chandra theorem and the exact
compact-dual Haar normalization, therefore identify the stored row
degree with $d_r(\lambda)$ using the literal \texttt{dsig\allowbreak Tie}.
This closes the original $h_3$ slot on every supplied family.

For prime size, use the original family cloud, not a freshly chosen
equally spaced model. Its interior coordinates satisfy
$u_i=v_{i-1}-1/2-\xi_i\Delta$, $\xi_i\in\{0,1\}$.
At distance $d$ the row-to-column factor is at most
$4d/(4d-1)$. Hence the interior logarithmic loss is at most
$\frac{k}{3}(1+\log k)$. The remaining factors give
\[
 \log\frac{V_r(v,-b)}{V_k(u)}\le5kmL.
                                                               \tag{AU32.3}
\]
The factorial factors cancel exactly against the mass ratio.
Combining this upper estimate with QW6's proved lower estimate gives,
for $X_u=(fk)^{-1}\log(D_*/d_*(u))$ and
$C_0=\max(5,\allowbreak{} C_{YL}+1)$,
\[
 15mL\le X_u\le(7+3C_0)mL
                                                               \tag{AU32.4}
\]
on one late-family threshold. The original $\Gamma_u$ correction
and original prime window then yield
\[
 B^{m/2}\le T_u,\qquad
 \log T_u\le(8+3C_0)mL.
                                                               \tag{AU32.5}
\]
The source's larger constant
$c_T(y)=\max(11,16+3\max(5,\allowbreak{} y+1))$ is unchanged.
DSZ10 supplies the needed original volume identity on the construction
class from \texttt{Built}/\texttt{Real\allowbreak Tied} and RSL1--RSL3.
This proof replaces only the upper comparison needed for
\texttt{Tu\allowbreak Size}; it does not assert the stronger full Vandermonde
variation or counting package.

\paragraph{The original class number and rank.}
RMP begins with the actual column lattice $L$:
$L\subset L^\#\subset T^{-1}L$ and $|L^\#/L|=T^{2f}$.
At every one of the $f$ split places over $T$, the original nonzero
prime-order functional gives one nonzero discriminant component.
Duality gives equal lengths to the two components; their total
cardinality forces each length to be one. Choosing the half singled
out by that functional gives the literal self-dual \texttt{over\allowbreak Lat}.
The local level fixes a specified nonzero functional, giving
$T^k-1$, rather than the projective-line index.

The actual auxiliary levels and the original genus mass now give
\[
 |\mathrm{cls}_0|=\mu_k\,\mathrm{idxPi}\,(T^k-1)^f.
                                                               \tag{AU32.6}
\]
Two distinct auxiliary primes force the finite rational level
stabilizers to be trivial. The original \texttt{Forms0} value maps
therefore identify the dimension of the actual $\mathrm{Mu}_0$
with this class count times the exact Weyl dimension product.
The stored \texttt{Fact\allowbreak Du\allowbreak Star} gives the literal rank identity
\[
                    d_M(u)=d_*(u)(T_u^k-1)^f.
                                                               \tag{AU32.7}
\]
The class-number proof is stronger than its source-facing contract:
it does not need \texttt{GBig}. The rank proof does not use a
theta square, finite-character assertion or metric comparison.
The existing construction-class Bundle supplies its needed columns.

\paragraph{The original local Hecke and congruence calculations.}
GLE1--GLE2 prove every clause of the eventual
\texttt{Hida\allowbreak Setting} from the actual parameter windows and the
lower bound in \texttt{Tu\allowbreak Size}. The remaining GLE proofs retain
the same original \texttt{Forms0}, finite field-valued character
and stabilizer identifications. They give the exact lattice
coordinates and
\[
 U_v(B_{w,\allowbreak{} v})\subset S_w,\qquad D_v(S_w)\subset B_{w,\allowbreak{} v},
 \qquad U_vD_v=T^{k-1}\quad\hbox{on }M^{\rm ram}.
                                                               \tag{AU32.8}
\]
The finite two-step move calculation includes the line contribution,
the inverse diamond convention, and the character-one comparison
between equal endpoints. The second branch is also exact:
\[
 T^{b_v}S_w\subset B_{w,\allowbreak{} v},\qquad
 \ell(B_{w,\allowbreak{} v}/T^{b_v}S_w)
      =\frac{e_wd_S}{k}(|\lambda|+kb_v),
                                                               \tag{AU32.9}
\]
where $|\lambda|$ denotes the signed sum of its entries.
The determinant shift is retained when computing each elementary
divisor. Exact sequences give the one-operator bound, then the
source's original ordered product and its $m$th power give the
complete congruence length bound. No interchange of differently
ordered factors is assumed. Empty unseen-place products are treated
as the identity.

These close five explicit GL6 implication tasks. Together with the
actual second-layer supplier and D3's built-in \texttt{Tu\allowbreak Size},
they give the remaining original $N_\lambda\subset S^{\rm ram}$
inclusion used in RSL fullness. The second-layer selector remains
a shared producer. For the global-slope branch, the original three
GL4 slots and the final numerical/lattice assembly remain unwritten.

\subsection{The smaller actual base-theta frontier}

BTH1 proves smooth-topology cyclicity of nonzero smooth vectors in
the irreducible source models. BTH2 then proves that the multiplicity
map is determined by its lowest-type values. These steps make the
continuity needed for uniqueness explicit.

BTH3 proves absolute convergence of the actual base lowest-type theta
integral. The original positive-definite real groups are compact;
the stored finite double-coset representatives give a compact set of
adelic representatives. Unfolding bounds the measure of the prescribed
fundamental domain by that of this compact set. The actual weak Weil
continuity and automorphy fields make the integrand bounded there.
This does not by itself prove a smooth cuspidal target extension.

Once the requisite base extensions and exact source readings are
supplied, BTH assembles the linear map, activates the original weak
choices and constructs the good-projected target map. Its Hilbert
reading uses an explicit tensor witness in the geometric theta law.
One base rational test family supplies both the rational adjoint
image and rational coordinates in the stored target basis.
The finite unitary argument kills the orthogonal complement of the
original good cut without replacing the source idempotent by a
self-adjoint one. This gives the actual base theta square.

BTH13--BTH15 also prove the exact normalization for the source's
already stored Weil action. A continuous adelic unitary character
trivial on rational and split-supported points is trivial; the
Heisenberg scalar lemma and positive split normalization then show
that two full Weil data with the same characters agree on the theta
space. Comparing the original action with the canonical construction,
with its determinant-character twist, gives its required finite
lattice-stabilizer character. The canonical action is therefore
identified with the actual source action where the covariance proof
uses it; this is a proved step, not an additional producer premise.

The exact remaining same-bundle inputs are listed in the register
below. Their principal mathematical parts are the actual geometric
reading, the smooth cuspidal extension of the now-convergent basis
integrals, the simultaneous rational tests, and the original base
rows. They must hold on one selected Bundle with its exact field,
embeddings, Hilbert and rational metrics, level, basis and chosen
maps. Independent models with matching dimensions do not suffice.
The stronger nonbase theta package remains relevant only to its
other consumers, not automatically to this base rank argument.

\subsection{The final theorem interface and the count branches}

The final sufficient interface is unchanged: one compatible
\texttt{Design\allowbreak Exists} and nine \texttt{Closed} predicates,
\[
\begin{gathered}
 \texttt{Link\allowbreak N},\quad\texttt{Link\allowbreak Eta},\quad\texttt{Link\allowbreak IRR},\quad
 \texttt{Small\allowbreak Level\allowbreak Rank\allowbreak Eq},\\
 \texttt{Tu\allowbreak Size}\Rightarrow\texttt{DGam\allowbreak Size}\Rightarrow\texttt{Link\allowbreak DIVt},
 \qquad\texttt{Chosen\allowbreak E}\Rightarrow
       (\texttt{Nu\allowbreak Full\allowbreak F}\wedge\texttt{Mu\allowbreak Rank}),\\
 \texttt{Link\allowbreak G},\qquad\texttt{Tu\allowbreak Size}\Rightarrow\texttt{Link\allowbreak S},\qquad
 \texttt{Tu\allowbreak Size}\Rightarrow\texttt{DGam\allowbreak Size}\Rightarrow\texttt{Link\allowbreak GLOBAL}.
                                                               \tag{AU32.10}
\end{gathered}
\]
All outer thresholds are joined before one family is selected.
The constants $C_{\rm abs}=2$, $C_0=3$, $e_D=2$, $e_L=0$
and the ZW exponent $1/2+9\epsilon$ are unchanged. HX still
uses the original \texttt{VWin\allowbreak At}$(\chi,401)$ and cutoff
$\lfloor x^{400}\rfloor$.

For RH and targets of power-of-two character order, the selected
arithmetic count route requires I, II and V. General Dirichlet
targets additionally require the actual Case-III/SeedIV/chain
construction. CountII retains all six original \texttt{Hida\allowbreak Props\allowbreak Alt}
fields and \texttt{P1core} on the same ideal of each of its two
original rank-four Links, after one coefficient threshold.
CountI's fixed field and CountV's Gaussian/quartic host are distinct.
The Case-V exponent $e=1$ is the original twist exponent, not the
unbounded dyadic host exponent. These distinctions survive every
geometric simplification in the present version.

\subsection{Classification and priority}

\begin{description}[style=nextline]
\item[A: matched published application.]
An identified published theorem is applied to the original object,
with all consumed hypotheses checked. The retained HC-Deg, MF and
GV applications remain recorded in their earlier scopes. They do
not become new unproved local claims in this version.
\item[B: independent elementary or moderate reconstruction.]
Some B arguments are complete certificates; the register below
contains only those B proof tasks still unwritten. A theorem name
in the Lean source is not counted as its human proof.
\item[C: substantial construction or proof.]
The exact original target is identified, but its substantial
geometric, analytic, representation-theoretic or arithmetic
producer remains to be reconstructed.
\item[D: incomplete human source understanding.]
The original instance or hypothesis bridge still needs a precise
human identification. This does not question the checked Lean proof.
\end{description}

Prioritize the common original row/theta reading, smooth cuspidal
extension and base rational tests: their construction is shared
with several other period, spectral and specialization consumers.
The next bounded global branch is the actual \texttt{Col\allowbreak Val}
dictionary and its two GL4 slope estimates, followed by the final
\texttt{Fam\allowbreak GLmin} assembly. In parallel, the same-ideal rank-four
ordinary Hecke/support producer remains the main independent count
priority. Repeating the now proved size, rank or local length
arguments would contribute less to the final theorem.

The audit records selected sufficient dependencies. Its AND rules
retain the proof of each implication and every actual producer
premise. Shared refinements are navigation links rather than extra
logical implications. All 245 inherited compound continuation records
and their 107 groups remain, with precise current subclaim updates.
Neither their counts nor the selected leaf count is a percentage of
the mathematical work, and no absolute minimality theorem is asserted.

\subsection{Complete current selected obligation register}

The following is the complete open-leaf register for this selected
route. Completed certificates and matched A applications are in the
provenance and published-input registries. Coupled hypotheses below
must be supplied on the same original objects.

\begin{description}[style=nextline]
\item[1. C: Actual LinkN local table and binding producer.]
\texttt{GLOBAL\_\allowbreak DENOMINATOR}. Closed(Des \(\mathcal C\) 3) LinkN: the actual etaSum is bounded with exact dZ0 and same-host vvTerm, subject only to the literal ShaFinOdd and summability guards.

\item[2. C: Original retained-row defect bound LinkIRR.]
\texttt{GLOBAL\_\allowbreak RETAINED\_\allowbreak DEFECT}. Rows.DelpNonneg and delp\(\le\)exp(-x\textasciicircum{}(ez/7)) uniformly for every late actual row datum.

\item[3. C: Full actual period and height comparison LinkDIVt.]
\texttt{GLOBAL\_\allowbreak PERIOD\_\allowbreak HEIGHT}. Closed conditional TuSize -> DGamSize -> (LinkEta -> DIVinN -> LinkDIV), with every original LinkDIV\_\allowbreak at clause.

\item[4. C: Original global LinkG installation.]
\texttt{GLOBAL\_\allowbreak LOCAL\_\allowbreak GAIN}. One uniform LinkG with actual local lattice lengths, served-orbit coverage and MeanKappaBound; delta(x)->0 before varying actual indices.

\item[5. C: Original normalized kernel/Gram estimates LinkS.]
\texttt{GLOBAL\_\allowbreak SPECTRAL}. Closed TuSize -> LinkS with exact Gamma normalization, off-diagonal Hilbert-Schmidt bound, diagonal error and Dall/Dstar estimate, errors exp(-x\textasciicircum{}(epsilon/8)).

\item[6. B: Install OneChar and the congruence integer into PropGK.]
\texttt{TASK\_\allowbreak PROP\_\allowbreak GK\_\allowbreak INSTALL}. Complete the finite orbit/Frobenius/valuation argument producing one nonzero integer before every later coefficient prime, with the original e,\allowbreak{} N,\allowbreak{} kernel,\allowbreak{} p and fixed-K constants.

\item[7. C: Reconstruct the original one-character argument.]
\texttt{TASK\_\allowbreak ONE\_\allowbreak CHAR\_\allowbreak PROOF}. From exact InpZ and AutomorphyUnramifiedCoefficients, prove the original OneChar trace/length estimate on the supplied coefficient datum and GD input.

\item[8. B: Install the exact Kato input record.]
\texttt{TASK\_\allowbreak INP\_\allowbreak Z\_\allowbreak ASSEMBLY}. Build the original CohInp/InpZ record using exactly lemM,\allowbreak{} lemMprime,\allowbreak{} lemP,\allowbreak{} RigII,\allowbreak{} AdVanishK and number on the same arithmetic objects. No rank-four replacement is permitted.

\item[9. C: Type-II polarized adjoint rigidity.]
\texttt{RIG\_\allowbreak II}. Marked original datum, actual type-II point and modularity of all source extensions imply vanishing for its exact adRep.

\item[10. D: Exact rank-three polarized adjoint matching.]
\texttt{AD\_\allowbreak VANISH\_\allowbreak K\_\allowbreak MATCH}. AdVanishK uses the full negative-transpose complex-conjugation action and irreducibility on the infinite cyclotomic kernel. RX does not transfer to every rank by name.

\item[11. C: Actual unramified-coefficient automorphy.]
\texttt{AUT\_\allowbreak UNRAM\_\allowbreak COEF}. For each actual Dat with coefficient ramification one and fixed iota, choose one marked g1 supplying FinA, AutB and LemW together.

\item[12. C: Uniform rank-three and rank-two congruence integer.]
\texttt{CNR3}. Fixed quadratic K -> c,\allowbreak{} A -> e,\allowbreak{} N,\allowbreak{} p,\allowbreak{} w,\allowbreak{} finite Galois-stable roots-of-unity orbit -> nonzero integer NN -> every later coefficient prime, closure, embedding, Frobenius and finite admissible collection.

\item[13. C: Two actual rank-four Links with common threshold and same-ideal Hida/P1.]
\texttt{COUNT\_\allowbreak II\_\allowbreak SUPPORT}. For each of the two original QX Links, at the already fixed iota, choose one common coefficient threshold before every later coefficient DVR and every actual universal B/Lambda; then construct one ideal J0 with ring,\allowbreak{} dense,\allowbreak{} pt,\allowbreak{} end1,\allowbreak{} end2,\allowbreak{} twist AND the original P1core for that very J0.

\item[14. B: Complete the original Case-III automorphy/image join.]
\texttt{TASK\_\allowbreak CASE3\_\allowbreak AUT\_\allowbreak JOIN}. Install the actual Case-III ProModAlt, SeedIVAlt and shared PMExAlt with the inherited CohomologyHyp/Ferrero-Washington/HFit/SeedII facts into the original whole Case-III automorphy/image conclusion.

\item[15. C: Reconstruct the actual Case-III chain-to-promodularity proof.]
\texttt{TASK\_\allowbreak CASE3\_\allowbreak PROMOD\_\allowbreak JOIN}. Use the original tensor-factor tower, ChainRankFour, ChainDesc4, seeded rank-four chain reduction, actual PMExAlt and SeedIVAlt to construct ProModAlt on every original Ctx3 star point at the same fixed iota, preserving all finite-coefficient, inertia and support maps.

\item[16. C: Actual Case-III seed with prechosen residual datum.]
\texttt{SEED\_\allowbreak IV\_\allowbreak ALT}. Base Ctx3/Four tower, iota and Amb -> w,\allowbreak{} s before every later Dies tower -> compatible finite-coefficient seed with exact residual support/root entries.

\item[17. B: Actual tensor-factor tower and characters.]
\texttt{TENSOR\_\allowbreak FACTOR\_\allowbreak TOWER}. Original finite-field, tame, factor and character constraints; source-level construction is sizeable and has not been reconstructed here.

\item[18. C: Actual SomeIn-to-AllIn propagation.]
\texttt{CHAIN\_\allowbreak RANK\_\allowbreak FOUR}. Same source Space and PM ideal, residual point Good4, original bound 30(C+1)\textasciicircum{}2 <= d; geometry must preserve all identities.

\item[19. B: Original residual chain descent.]
\texttt{CHAIN\_\allowbreak DESC4}. Finite residual descent retaining source Chain4 compatibility.

\item[20. C: Seeded reduction to the original residual chain.]
\texttt{RANK4\_\allowbreak CHAIN\_\allowbreak REDUCTION}. Select T1 before every later Deep tower; retain one compatible seed and Chain4 for every original residual lift over Q.

\item[21. B: Complete the original CountV finite assembly.]
\texttt{TASK\_\allowbreak COUNT\_\allowbreak V\_\allowbreak INSTALL}. From exactly the listed source host/cohomology inputs, complete countV\_\allowbreak of\_\allowbreak finiteness\_\allowbreak duality\_\allowbreak euler\_\allowbreak host, including original host/frame formation, normalized length, the CaseV guard and R1<=2\textasciicircum{}c2.

\item[22. C: Original unitary-three ordinary support and automorphy.]
\texttt{THM\_\allowbreak S6\_\allowbreak S}. Original host/frame and fixed iota -> one marked g1 -> every universal R,\allowbreak{} u -> weight-zero finiteness and every Type-I BPrime\_\allowbreak s point.

\item[23. B: Original Case-V host/frame and normalized-length installation.]
\texttt{HOST\_\allowbreak FRAME\_\allowbreak LENGTH}. Actual HostData, frame, local-condition equality, abstract count and normalized original Selmer length must all refer to the same objects.

\item[24. D: Rigidity on the original CountV frame.]
\texttt{HOST\_\allowbreak RIGID\_\allowbreak S}. Same original host/frame, marked point and eta-kernel; published adjoint vanishing is usable only after exact representation and hypothesis matching.

\item[25. C: Uniform original-host rank-three catalogue.]
\texttt{HOST\_\allowbreak NUMBER\_\allowbreak S}. One exponent before all closures, coefficient rings, levels, characters, frames, embeddings and marked points; one separating trace function controls all later finite pairwise nonisomorphic collections.

\item[26. B: Transport K/duality to literal HostVanE.]
\texttt{TASK\_\allowbreak HOST\_\allowbreak VAN\_\allowbreak E\_\allowbreak INSTALL}. Use the already supplied LemmaK/local-global duality and finiteness on the actual CountV host/frame and exact original coefficient modules.

\item[27. B: Transport CorX to literal HostXpp.]
\texttt{TASK\_\allowbreak HOST\_\allowbreak XPP\_\allowbreak INSTALL}. Preserve Psi\textasciicircum{}-2 eta, the prime-to-ell quotient, actual local conditions and the original CountV host while installing BR7.

\item[28. C: Construct the original common design family and bundle.]
\texttt{DESIGN\_\allowbreak CHOICES\_\allowbreak COMMON}. For the one target HostExW, construct the original compatible DesignChoices on Cdes=WithTw(WithSm Cprime), where Cprime=WithChN(A intersect target class), A is the one Built construction class, and cT(eps,\allowbreak{} y)=cTmin(y). Retain every original field, tuple, word, metric and common bundle duty. DSZ10 and RMP11 prove the size/rank properties of members of this class; they do not produce a member. The RecExP/RecAtG and common object/row route remains a coupled construction.

\item[29. C: Actual local specialization on the common selected family.]
\texttt{LOCAL\_\allowbreak NBUILT\_\allowbreak COMMON}. Closed on the same Des(C,3) class: one eventual family of actual NBuilt certificates for the original Fm.D(x,\allowbreak{} i), including NdSpec, original iV/Hecke actions, every required maximal ideal, \(\Psi\)/range and projection square. Keep the constructor certificate before forgetting it to TSel.

\item[30. C: Original second-layer selector and class duties.]
\texttt{SECOND\_\allowbreak LAYER\_\allowbreak COMPATIBLE\_\allowbreak SELECTION}. For the same A, supply Duty on A and a Closed common original TSel family whose selected records are the original true-column rec7 records (or have the entire dependent tuple transported), with their real witness and original coefficient-field root data. Existence/compatibility remains arithmetic; merely abstractly isomorphic fields or independently selected columns do not suffice.

\item[31. B: Original column-lattice lower slope.]
\texttt{TASK\_\allowbreak GL4\_\allowbreak SLOPE\_\allowbreak LOWER}. Prove hZa=Closed FamilySlopeLowerOpenPart using the actual ColVal.Q, exact base-change lattice and metric: column\_\allowbreak lattice\_\allowbreak exact\_\allowbreak slope followed by the standard Schur slope and hook estimate.

\item[32. B: Original class-map upper slope.]
\texttt{TASK\_\allowbreak GL4\_\allowbreak SLOPE\_\allowbreak MAX}. Prove hZb=Closed FamilySlopeMaxOpenPart(4) using class\_\allowbreak map\_\allowbreak archimedean on original class bases and metrics and the constant-four hook bound of PStmtFree.

\item[33. B: Actual ColVal dictionary.]
\texttt{TASK\_\allowbreak GL4\_\allowbreak COLUMN\_\allowbreak VALUES}. Prove tCV: eventual Nonempty ColVal(D,\allowbreak{} u) on the same family. Supply original ColTieW polynomial realization, exact local Q\_\allowbreak (c,\allowbreak{} tau) membership and archimedean norm formula, then colValOfTiedS\_\allowbreak holds. An abstract slice model alone is insufficient.

\item[34. B: Full final global numerical and lattice assembly.]
\texttt{TASK\_\allowbreak GL\_\allowbreak FINAL\_\allowbreak ASSEMBLY}. Given precisely the eleven original FamGLmin fields (h2b,\allowbreak{} h2v,\allowbreak{} h3,\allowbreak{} hZa,\allowbreak{} hZb,\allowbreak{} tCV,\allowbreak{} hW,\allowbreak{} hmu,\allowbreak{} hI,\allowbreak{} hR,\allowbreak{} hB), the reconstructed congruence length, original Duty/ChosenWall and TuSize, complete linkGLOBAL7\_\allowbreak min/items and link8 on the same class, with Cabs=2,\allowbreak{} C0=3,\allowbreak{} C5=4,\allowbreak{} eD=2,\allowbreak{} eL=0. FR1/GF8 give frame/HN consequences only in their stated scopes; the remaining inequalities and all actual slot hypotheses must be written.

\item[35. C: Actual base reading on the common bundle.]
\texttt{BASE\_\allowbreak THETA\_\allowbreak READINGS}. For every family in the original class, eventually on the SAME selected Bundle and for each i<m, supply the actual R,\allowbreak{} F,\allowbreak{} Read,\allowbreak{} ReadL, weak models/Weil data and phase-one pins. Identify their LiftTie.theta, mtRead, mjRead and ThDef3 basis with that Bundle.P.A3, P.Zs, P.S0 and its stored maps. Exact field, embeddings, metrics, measures, row family and weak choices are retained. This is one part of the coupled BTH12 producer, not a freely replaceable isomorphic model.

\item[36. C: Extend the original convergent base integrals.]
\texttt{BASE\_\allowbreak THETA\_\allowbreak CUSPIDAL\_\allowbreak EXTENSION}. On the SAME selected Bundle and the SAME reading as BASE\_\allowbreak THETA\_\allowbreak READINGS, for each i<m extend the exact lowest-type theta integral for every vector of one finite basis of HV(i,\allowbreak{} t0) to the original Wall\_\allowbreak i. Absolute convergence is already BTH3. The remaining assertion includes smoothness, equivariance, cuspidality, prescribed finite level and Sobolev-continuous multiplicity realization. Supply these with one family threshold before i and the basis vectors.

\item[37. C: Simultaneous original base rational tests.]
\texttt{BASE\_\allowbreak THETA\_\allowbreak RATIONAL\_\allowbreak TESTS}. On the SAME original Bundle/reading/theta map, for every i<m supply one finite separating complex-linear test family. Its values must lie in the fixed embedded coefficient field on theta of EVERY original rational source vector and on EVERY stored target-coordinate basis vector. BTH8 and BTH11 use this one base LiftTestB family for rational adjoint image and matrix coordinates. Finite dimensionality alone does not supply these rationality identities.

\item[38. C: Actual base row identities on that same bundle.]
\texttt{SLR\_\allowbreak BASE\_\allowbreak ROWS\_\allowbreak PRODUCER}. For the SAME selected Bundle, common original fields and readings used by the other three BTH12 inputs, supply actual RowsHold0\_\allowbreak 3 uniformly eventually for every family. Its original primitive row realization, spectrum/type and basis identifications must be retained. BTH supplies a theta consumer only; it does not construct the base row identity. No separate witness may be chosen merely to satisfy this row requirement.

\end{description}


\paragraph{Historical version-0.31 audit.}
The following audit is retained as a historical checkpoint.
Section~\ref{sec:dependency-audit-v32} and the version-0.32 additive
scope updates specify the current remaining producer obligations.
Its previous conditional proofs retain their original hypotheses.

\section{Current dependency audit: compatible producers and completed consumers}
\label{sec:dependency-audit-v31}

\status{Version 0.31 accepts the underlying Lean proof as fully checked
and correct. It adds 41 source-linked human proof claims in four
components. The full human proof is still unfinished. The advance is
that several apparently large obligations now have small, exact
producer interfaces, while the finite arguments consuming them are
written out. The selected route remains ordinary RH and Dirichlet
GRH; no assertion about a broader class of automorphic $L$-functions
is added.}

\subsection{How to read the audit}

An obligation has both a category and a completion status.
\begin{description}[style=nextline]
\item[A: a matched published theorem.]
The theorem is identified and its hypotheses are checked on the
original object. A name or a nearby result is insufficient.
\item[B: an elementary or moderately difficult argument.]
The argument can be reconstructed independently. Some B items below
are complete; others remain explicit proof tasks.
\item[C: a substantial mathematical producer.]
Its exact target is understood, but its human proof or compatible
construction still requires substantial reconstruction.
\item[D: an unresolved human source identification.]
The exact interpretation, original instance or quantifier bridge
has not yet been settled. This describes our understanding and
does not allege a defect in the checked formal proof.
\end{description}

The directed audit records selected sufficient implications.
An AND edge includes every premise of its displayed rule, including
any still-unwritten proof of that rule. A completed conditional
implication does not supply its antecedent. Shared work is recorded
separately and is not silently promoted to a sufficient conjunction.
All classes, selectors, fields, coefficient embeddings, lattices,
metrics and original operators are retained.

The earlier 245 continuation records and their 107 groups remain
available with their complete compound scopes. Each now carries
explicit version-0.31 subclaim reductions where applicable.
Their historical A/B/C/D counts are navigation statistics; they are
not a completion percentage or a count of independent mathematical
obligations. The current remaining-leaf register classifies every
open task in the selected logical graph.

\subsection{The unchanged final interface}

For a fixed primitive target character $\chi$, the selected final
construction still has exactly ten input contracts:
\[
\begin{gathered}
\texttt{Design\allowbreak Exists},\quad
\texttt{Link\allowbreak N},\quad\texttt{Link\allowbreak Eta},\quad\texttt{Link\allowbreak IRR},\quad
\texttt{Small\allowbreak Level\allowbreak Rank\allowbreak Eq},\\
\texttt{Tu\allowbreak Size}\Rightarrow\texttt{DGam\allowbreak Size}\Rightarrow\texttt{Link\allowbreak DIVt},
\quad
\texttt{Chosen\allowbreak E}\Rightarrow(\texttt{Nu\allowbreak Full\allowbreak F}\wedge\texttt{Mu\allowbreak Rank}),\\
\texttt{Link\allowbreak G},\quad
\texttt{Tu\allowbreak Size}\Rightarrow\texttt{Link\allowbreak S},\quad
\texttt{Tu\allowbreak Size}\Rightarrow\texttt{DGam\allowbreak Size}\Rightarrow\texttt{Link\allowbreak GLOBAL}.
\end{gathered}                                                   \tag{AU31.1}
\]
The nine predicates after design existence are \texttt{Closed} on
the same design class. Their transported bounds are joined before
one family is selected. The three class projections and the target
\texttt{Hyp\allowbreak V} are unchanged. The constants are
$C_{\rm abs}=2$, $C_0=3$, $e_D=2$, $e_L=0$.
The reconstructed ZW argument therefore gives the original window
exponent $1/2+9\epsilon$ from the literal tied family.

The host is already constructed from the original
\texttt{VWin\allowbreak At}$(\chi,401)$ by HX, with cutoff
$\lfloor x^{400}\rfloor$. The count input for a target of
power-of-two order, including the principal target for RH, retains
I, II and V. KA5--KA6 remove Case III only in this specialization.
An arbitrary Dirichlet target still requires Case III.
Nothing in the present geometric reductions changes those count
producer boundaries.

\subsection{Design existence: a proved three-input sufficient rule}

DEX7 proves the following rule on one common class $\mathcal C$,
with one common function $c_T$:
\[
\begin{gathered}
\texttt{Design\allowbreak Choices}(\mathcal C,3,c_T),\\
\texttt{Closed}\bigl(\mathcal C,\,
c_T(\epsilon,C_{YL})\le C_{Tu}\Rightarrow\texttt{Tu\allowbreak Size}\bigr),\\
\texttt{Closed}\bigl(\mathcal C,\,
\hbox{eventual original }\texttt{Rank\allowbreak M}\bigr)
\quad\Longrightarrow\quad
\texttt{Design\allowbreak Exists}(\mathcal C,3).
\end{gathered}                                                   \tag{AU31.2}
\]
All three premises are still actual producer obligations, classified
C until their original constructions have been written out.
The implication itself is B and complete.

The numerical ingredients are now explicit. DEX1--DEX2 prove
\[
\texttt{Boxes\allowbreak F}\wedge\texttt{QDRows\allowbreak F}
                 \Longrightarrow\texttt{Chosen\allowbreak Thin}(1)
                                                               \tag{AU31.3}
\]
on every supplied original family. The real-character argument uses
the actual conductor divisibility and a bound for the odd part of
the modulus. It does not bound the dyadic part.
DEX3 checks the published rational Minkowski theorem for the exact
integer \texttt{mink\allowbreak M}; rank zero is handled separately.
DEX4 gives
\[
 \log\operatorname{minkM}(2z)\le12z\log z\quad(z\ge2),
\]
so the actual \texttt{DGam\allowbreak Dvd\allowbreak F} supplies \texttt{DGam\allowbreak Size} with
constant $12$.

For the rank proviso, DEX5 uses the exact formula
$d_M(u)=d_\star(u)(T_u^k-1)^f$ and the original narrow prime window.
After the original $T_\star$, $\Gamma$ and $D_\star$ identities,
\[
\frac{d_M(u)\Gamma^2}{D_\star}
 =\left(\frac{T_u}{T_\star}\right)^{fk}(1-T_u^{-k})^f
 \in[1/2,e].
                                                               \tag{AU31.4}
\]
This is stronger than the required $[1/3,3]$ interval.
It proves the numerical proviso from the actual rank formula;
it does not create that formula from an unrelated dimension count.

DEX8--DEX11 also reconstruct the whole finite assembly from the
literal recursive premise \texttt{Rec\allowbreak Ex\allowbreak P(with\allowbreak QD QOf)} and the same
target \texttt{Host\allowbreak Ex\allowbreak W} to compatible design choices.
The universal eventual row predicates apply to every linked bundle;
one bundle supplied by the existence premise therefore satisfies
their finite conjunction. The dependent tuple equality is retained.
The thresholds, large-base intersections, injective word code,
kit coverage, fallback real indices and original object identities
are all checked.

The recursive construction remains conditional. QW9 and MF20
already close its historical logarithmic-window premise; the
remaining \texttt{Rec\allowbreak At\allowbreak G} core and eventual common bundle/row
property must still be supplied. Reconstructing the window again
would contribute nothing to the current frontier.

\subsection{Small-level rank: the base tuple suffices}

SLR1--SLR9 remove the former D uncertainty about the original
level projection. On every supplied original bundle,
\[
                     \mathrm{lPrj}\,\mathrm{lZ0}=1.
                                                               \tag{AU31.5}
\]
The adjoint computation retains the original $|P|$ scaling of the
inner product. It proves the required projection and surjectivity;
it does not assert equality of two integral lattices when only
their inclusions are known.

The smaller remaining producer is one common bundle selector with
the following two properties at one eventual threshold:
\[
\texttt{Rows\allowbreak Hold0\_\allowbreak 3}(P)
\qquad\hbox{and}\qquad
\text{the original base-tuple }\texttt{th\allowbreak Theta}\text{ square on }P.
                                                               \tag{AU31.6}
\]
The finite type decomposition uses generalized primary spaces,
including their nilpotent parts. The actual stored bases identify
the row dimension with $D_0$ directly. The base theta square and
(AU31.5) then make the original theta map surjective; its adjoint
has the same rank, giving the exact \texttt{Small\allowbreak Level\allowbreak Rank\allowbreak Eq}.
No nonbase theta square, \texttt{Tup\allowbreak All\allowbreak Used}, full pairing part of
\texttt{Div\allowbreak Rows3}, or \texttt{Rows\allowbreak Hold\allowbreak Lam3} is needed for this
consumer.

The finite consumer is B and complete. The actual compatible
row and theta producer in (AU31.6) is C. In particular, the Hilbert
reading, lowest-$K$-type comparison and rational-adjoint reading
must refer to this bundle's actual maps. This is now a sharply
bounded target for the next substantial reconstruction.

\subsection{Local coverage: finite algebra after actual specialization}

LTD3--LTD5 prove
\[
\begin{gathered}
\text{actual }\texttt{Nd\allowbreak Spec}
 \Longrightarrow\texttt{Locs\allowbreak Cover\allowbreak Tw}\text{ on the original local installation},\\
\text{eventual actual }\texttt{NBuilt}\text{ on }F.D(x,i)
 \Longrightarrow\texttt{Link\allowbreak Eta}(F).
\end{gathered}                                                   \tag{AU31.7}
\]
The algebra is finite. Successively refine invariant generalized
eigenspaces for the commuting split integral Hecke operators.
Each strict refinement increases the number of nonzero blocks, so
the process stops in at most the dimension. On a final block every
generator is a scalar plus a commuting nilpotent. The scalar
character is well defined, and grouping these characters by their
original maximal ideals proves the coverage and joint projection
injectivity.

The specialization maps, ranges and projection squares in
\texttt{Nd\allowbreak Spec} are essential. The original
\texttt{NBuilt} construction carries them; a generic \texttt{TSel}
record does not. Producing this compatible local data is C, while
its coverage consumer is now B and complete.

LTD1--LTD2 and LTD6--LTD8 also settle the literal selection,
localization, operator and field transports. The rational row
operators of rank $r$ remain distinct from the ambient operators
of rank $n_V$. The prescribed $c_{\rm Hk}^k$ normalization is
retained. Field membership is checked inside the actual embedded
coefficient field; existence of some abstract number field is
insufficient. For bottom periods, \texttt{Sums\allowbreak In} already assumes
the normalized monomial-integral rationality subject to
integrability. The separate analytic hypotheses discharge those
guards; integrating field-valued kernels alone would not establish it.

\subsection{The second layer: mass, metrics and the smaller fullness rule}

RSL1--RSL3 match the exact universal mass and volume inputs on
$W.h_T,W.L_0$. They use the existing, explicitly normalized
GV Tamagawa and MF Hasse--Minkowski applications. They close the
same-object interpretation of \texttt{h\allowbreak U} and \texttt{h\allowbreak G}, including
the reciprocal orientation and half power in the volume ratio.
The original \texttt{Vol\allowbreak Tie} remains essential.

RSL4--RSL6 show that the true-column additive characters already
supply the needed $g_0$th and $g_1$th roots in the original
coefficient field. Every supplied symmetric-form column on that
record has level-zero character of exponent dividing $g_0g_1$.
Thus one keeps the selector's actual columns and real witness.
There is no need to choose a new column system and then hope
it matches the theta data. TG12 had already closed the earlier
terminal-column existence implication in its actual supplied-data
scope; that older frontier is not reopened.

RSL7--RSL10 construct the actual integral and isometric model.
Two auxiliary primes kill the stabilizers. Restriction to free
orbit representatives yields an original rational slice $W$ and
an integral decomposition at every coefficient place. Finite
character values are roots of unity in the fixed field, hence have
modulus one at every embedding; the exact norm identity proves
the isometry at every archimedean place.

With $q=T-1$, let
$A_0=\ker(\mathcal O_K^q\xrightarrow{\sum}\mathcal O_K)$.
The original full and ramified lattices are respectively
\[
 W\otimes(\mathcal O_K^q)^{\otimes f},
 \qquad W\otimes A_0^{\otimes f},
                                                               \tag{AU31.8}
\]
integrally and with their original metrics. The Gram determinant
of the basis $e_j-e_q$ is $q$, giving
\[
 \frac{d_u}{d_{u,S}}
   =\left(1-\frac1{T-1}\right)^f,\qquad
 \mu_{\rm ram}=\mu_{\rm full}
                     -\frac{f\log(T-1)}{T-2},\qquad
 \mu_{\max,\rm ram}\le\mu_{\max,\rm full}.             \tag{AU31.9}
\]
This supplies the actual transport previously retained in LC7.

RSL11--RSL12 prove a smaller fullness interface:
\[
\begin{gathered}
\text{actual second layer},\qquad\texttt{Tu\allowbreak Size},\qquad
\text{eventual }N_\lambda(D,u,w)\subset S_{\rm ram}(D,u,w)\\
\Longrightarrow\quad\texttt{Nu\allowbreak Full\allowbreak F}\wedge\texttt{Mu\allowbreak Rank}.
\end{gathered}                                                   \tag{AU31.10}
\]
Local inclusion, the original \texttt{nu\allowbreak Def}, and one eligible
place away from $T_u$ give
$cR\subset M\subset R$ with $c\ne0$. The actual
\texttt{nprime\allowbreak Def} expresses the final lattice as the original
nonzero fractional ideal times $M$. It is consequently finitely
generated and spans the original division space. The nonzero
orbit slice gives positive rank.

The original prime has $T\equiv1\pmod4$, so primality gives $T\ge3$
pointwise. This supplies the positivity used above without the full
\texttt{Hida\allowbreak Setting} theorem. The displayed \texttt{Tu\allowbreak Size} premise is
therefore harmless retained source-facing data; on the final class
$D3(3,\mathcal A)$ it is already a projection of membership.
The bare-inclusion fullness argument needs no size-to-setting
installation, GL4 slope estimate or final global-slope assembly.

\subsection{Exactly what the global slope still needs}

The eleven slots of \texttt{Fam\allowbreak GLmin} now have the following
current disposition:
\begin{longtable}{p{0.12\textwidth}p{0.46\textwidth}p{0.30\textwidth}}
\textbf{Slot}&\textbf{Original content}&\textbf{Current status}\\\hline
$h_{2b}$&Adjacent genus-mass ratio&MF19 complete.\\
$h_{2v}$&Actual base volume equals mass&RSL1--RSL3 complete from real tie.\\
$h_3$&Two-sided degree bound from the gap&B: finite product and original-row installation remain; DC supplies the deep normalized producer.\\
$h_{Za}$&Actual column lower slope&B: exact lattice/metric formula and hook estimate remain.\\
$h_{Zb}$&Actual class-map maximal slope&B: class-map degree and constant-four estimate remain.\\
$tCV$&Actual polynomial/local/metric column dictionary&B: full \texttt{Col\allowbreak Tie\allowbreak W}/\texttt{Col\allowbreak Val} installation remains.\\
$h_W$&Free ramified action model&RSL7--RSL10 complete from second layer.\\
$h_\mu$&Actual ramified mean slope&RSL9--RSL10 complete from second layer.\\
$h_I$&Two integral level-zero Hecke maps&B: original local coset calculation remains.\\
$h_R$&Ramified stability and inverse composite&B: exact operator calculation remains.\\
$h_B$&Second-branch inclusion and length&B: original DVR calculation remains.\\
\end{longtable}

In addition, the original \texttt{Tu\allowbreak Size}-to-eventual-\texttt{Hida\allowbreak Setting}
installation, finite composite-length estimate and full final
numerical/lattice assembly remain explicit B proof tasks.
The size installation is a real guard: \texttt{Fam\allowbreak U} means
eventually $\texttt{Hida\allowbreak Setting}\Rightarrow P$, not eventually $P$.
The selected strong operator route must therefore supply all of its
level, alcove and $b_{\rm Top}\ge0$ conditions before using those
three conditional contracts.
The three GL4 slots are an exact sufficient cut for those slots,
not a replacement for all eleven. The three full GL6 contracts
supply the weaker inclusion of (AU31.10) through the completed
RSL12 argument, but their stronger length and inverse-composite
parts remain necessary on the selected global-slope route.

The remaining C obligations elsewhere are the common recursive
geometric construction, actual row/theta/local specialization,
original denominator and period-height arithmetic, retained
defect/local gain/spectral estimates, original rank producer and
the I/II/V count producers. Case III is additional only for
unrestricted target order. Exact count-side D matches include the
original adjoint representation and host rigidity; their scopes
are unchanged from version 0.30.

\subsection{A corrected historical intermediate estimate}

Independent review found that RF.45--RF.46 omitted the
$\epsilon/100$ term in the literal upper window for $f$.
DEX11 supplies the corrected proof:
\[
 f\le4x^{2+8\epsilon+\epsilon/100},\qquad
 m^5fx^\epsilon\le128x^{9/2+9\epsilon},\qquad
 k\ge x^{9/2+10\epsilon}.
                                                               \tag{AU31.11}
\]
The positive exponent gap is still $\epsilon$. It gives the same
uniform $\log\#\mathrm{Word}\le k/4$ bound and the same
$\#\mathrm{Word}\le2^{k-3}$ conclusion in RF.48.
The historical bytes are preserved for provenance; the current
audit and additive RF16 update explicitly supersede that
intermediate estimate. This is a correction to the human text.

\subsection{Priority for the next substantial work}

The highest-value shared producer is now the original base-row
and base-theta square on a common bundle. The immediate target
is the exact Hilbert/lowest-$K$-type and rational-adjoint reading,
joined to the actual \texttt{Rows\allowbreak Hold0\_\allowbreak 3} producer, then carried
through that bundle's construction. This would finish the
arithmetic input of the narrowed small-rank consumer and clarify
interfaces shared with local coverage, denominators, period height
and the spectrum.

In parallel, the three actual GL6 operator and length calculations
are bounded B tasks with a clear payoff: their weaker parts feed
the already complete fullness rule, and their full statements
feed global slope. The original \texttt{Rank\allowbreak M} supplier is another
focused target, using the already available mass and class-space
work while preserving the actual class-number and level-index
identities. The same-ideal rank-four \texttt{Hida\allowbreak Props\allowbreak Alt} plus
\texttt{P1core} producer remains an independent substantial count
priority. Its coefficient threshold and both original Links must
be retained.

These priorities follow the shared dependencies and the explicit
remaining cuts. They are qualitative judgments about contribution,
not measured runtimes or a claimed percentage of the proof.
The formal source is not recompiled in this round. Source/proof
hashes, inherited-file continuity, logical reachability, independent
mathematical reviews and draft-mode TeX references are checked
separately.


\paragraph{Historical version-0.30 audit.}
The next audit is retained as a historical checkpoint. Its frontier
statements are superseded precisely by Section~\ref{sec:dependency-audit-v31}
and the version-0.31 scope updates. Its conditional proofs and source
identities retain their stated hypotheses.

\paragraph{Current scope and historical chapters.}
The version-0.30 audit below records the preceding boundary of this exposition.
Earlier version-specific summaries and conditional arguments are retained
to preserve their provenance.  Their statements that a host, metric,
numerical gap, or Case-III premise is still open must be read with the
precise supersessions in this audit.  Closed implications retain their
actual arithmetic hypotheses.  The published applications are checked
in their stated scopes; inherited certificates are not described as new
independent reviews.  This round performs no Lean build or kernel replay.

\section{A connected minimal sufficient dependency audit}
\label{sec:dependency-audit-v30}

\status{The stipulated correctness of the checked Lean proof is accepted.
This audit concerns only completion of its human reconstruction.
It connects the final theorem to the actual common-family and host-count
obligations, separates ordinary RH from unrestricted Dirichlet GRH, and
classifies every retained continuation record.  Minimality here is
relative to the explicitly selected source route after removing proved
implications, duplicate producers and stronger optional interfaces.  It
is not a proof that no mathematically shorter proof exists, nor a
transitive enumeration of every Lean declaration.}

\subsection{How to read the audit}

The classification of an argument and its completion status are separate.
A deep published theorem may be a completed application of class A; an
elementary assertion may be an open class-B task if its original objects
have not yet been identified.  In particular, a proved implication
$P\Rightarrow Q$ is not a proof of $Q$ while $P$ is still a producer
obligation.

\begin{center}
\begin{tabular}{p{.08\linewidth}p{.83\linewidth}}
A&Direct application of a published theorem, with the exact consumed
hypotheses checked.  The certificate records its theorem, normalization,
source objects and parameter range.\\
B&Elementary or moderately difficult argument that can be reconstructed
independently, including finite algebra, numerical inequalities and
precise assembly on the original objects.\\
C&Genuinely substantial mathematical reconstruction of the actual
arithmetic, automorphic, geometric or spectral producer.\\
D&A possible gap in our understanding of the source: an unresolved
definition, object identity, normalization, theorem match or quantifier
transport.  This is not an assertion of a gap in the formal proof.\\
\end{tabular}
\end{center}

The new machine-readable map has explicit logical edges for the chosen
source implications.  Work-sharing and possible-reuse edges are separately
marked: identifying a common local dictionary is useful to several
producers, but the dictionary alone is not asserted to prove any of them.
The complete records and definitions remain attached to the frontier.

This repairs a navigation defect in version 0.29.  Following only its
logical \texttt{requires}/alternative edges from the two final roots
reached 28 of its 151 nodes; the count/cohomology branch was attached
through explanatory \texttt{related\_\allowbreak nodes}, not a logical host edge.
The earlier map was a sufficient-interface guide, but could not be used
as a connected completion audit.  The earlier map and every proof are
preserved as historical records.

\subsection{The exact conditional route to the final theorem}

Fix a primitive target $\chi$ first.  The selected route is:
\[
 \begin{gathered}
 \text{original count suppliers}
       \Longrightarrow \mathrm{VWinAt}(\chi,401)
       \mathrel{\mathop{\Longrightarrow}^{\rm HX4}}
                 \mathrm{HostExW}(3/2,(\mathrm{GChi}\chi).\mathrm{req}),\\
 \text{same-target host}+
       \text{one compatible design and its closed links}
       \Longrightarrow \mathrm{FamilyClassConnected},\\
 \mathrm{FamilyClassConnected}
       \mathrel{\mathop{\Longrightarrow}^{\rm ZW4,WB8,AN3}}
                L(s,\chi)\ne0\quad(\Re s>1/2),
 \end{gathered}                                                    \tag{DA30.1}
\]
with the principal pole treated as in the exported statement.
Every arrow is quantified over the original source models and preserves
the target; the host window alone does not construct the compatible
geometric family.  Applying this route at $\chi=1$ gives ordinary RH.
Applying it for every primitive target, and then GH6's imprimitive-factor
and functional-equation argument, gives Dirichlet GRH.

The common-family interface is exactly the argument list of
\texttt{top\allowbreak Des\allowbreak Win}.  It asks for a compatible design and the following
nine closed predicates on the same class, together with its class
inclusions and \texttt{Hyp\allowbreak V}:
\[
\begin{gathered}
 \mathrm{LinkN},\ \mathrm{LinkEta},\ \mathrm{LinkIRR},\
 \mathrm{SmallLevelRankEq},\
 (\mathrm{TuSize}\Rightarrow\mathrm{DGamSize}\Rightarrow\mathrm{LinkDIVt}),\\
 (\mathrm{ChosenE}\Rightarrow\mathrm{NuFullF}\wedge\mathrm{MuRank}),\
 \mathrm{LinkG},\
 (\mathrm{TuSize}\Rightarrow\mathrm{LinkS}),\\
 (\mathrm{TuSize}\Rightarrow\mathrm{DGamSize}\Rightarrow\mathrm{LinkGLOBAL}).
\end{gathered}                                                    \tag{DA30.2}
\]
Thus there are exactly ten first-level producer interfaces: the design
and these nine closed predicates.  In the selected \texttt{shape\allowbreak Win}
caller, $(C_{\rm abs},C_0,e_D,e_L)=(2,3,2,0)$; design existence is on
$\mathrm{WithTw}(\mathrm{WithSm}\,\mathcal C')$, the nine predicates
are closed on $\mathrm{Des}(\mathcal C,3)$, and the additional guards
are $\mathcal C'\subseteq\mathcal C\cap\mathcal E$ and
$\mathcal E\Rightarrow\mathrm{HypV}$.  Here ``closed'' means that for
every sufficiently small positive $\epsilon$ one threshold is chosen
before all families and height parameters in that class; the predicate
then holds whenever the family is large at that threshold.  Design
existence must work for every such transported threshold.

The finite threshold maximum is taken before design existence selects
one family.  It would be invalid to supply each predicate on an
independently chosen family.  The target class gives both the required
$\mathrm{EAtH}$ and the same cohomological $\mathrm{HypV}$.
DA1--DA2 prove the assembly; MT7 now discharges its previously retained
metric premise.  The thresholds, coefficient fields, word sets, lattices,
local flags, maps and metrics keep their literal bindings.

\subsection{The smaller counting cut for ordinary RH}

For an unrestricted target, KA4 and the existing YX/HV/TW/OB/VW chain
give the sufficient cohomological frontier
\[
 \mathrm{CorGK}_{K_\chi},\quad \mathrm{PMExAlt},\quad
 \mathrm{CountIII}_\chi,\quad \mathrm{CountV}_\chi,                \tag{DA30.3}
\]
relative to their recorded classical and prime-pair inputs.
Here $K_\chi=K_*(p_K(\chi))$ is chosen once for the fixed target;
$\mathrm{PropGK}_{K_\chi}$ is a sufficient stronger input for its
first entry.  The quantifiers over later hosts, coefficient DVRs,
closures, finite prime sets and layer functions remain intact.

KA5 proves a genuine target specialization.  Every original host
character $\psi$ satisfies
$\operatorname{ord}(\psi)\mid2^{s+2}\operatorname{ord}(\chi)$.
Case III requires $\operatorname{ord}(\psi)$ to be odd and
$\psi\ne1$.  If $\operatorname{ord}(\chi)$ is a power of two,
these conditions contradict one another.  KA6 checks that each consumer
through \texttt{VWin\allowbreak At} uses only that same fixed target.
Consequently (DA30.3) reduces to
\[
 \boxed{\quad
 \mathrm{CorGK}_{K_\chi},\quad\mathrm{PMExAlt},\quad
 \mathrm{CountV}_\chi\qquad
        (\operatorname{ord}(\chi)\text{ a power of two}).\quad}    \tag{DA30.4}
\]
In particular, the separate Case-III seed and promodular construction
do not lie on this selected route to ordinary RH.  They remain on the
unrestricted Dirichlet route.  No universal \texttt{Count\allowbreak III} theorem
has been inferred from a single-target empty case.

For $\chi=1$, $K_\chi\cong\mathbb Q(\sqrt{-7})$ and the exceptional
set is $\{3,7\}$; KA gives the reduced-form proof of its class number.
Count V still concerns unbounded dyadic conductors at those two primes.
KA's concrete examples show that Cases I, II and V remain nonempty,
including characteristic three.  A theorem requiring $\ell>3$ therefore
cannot automatically close the entire principal support branch.

\subsection{What this version actually proves}

\begin{center}
\begin{longtable}{p{.12\linewidth}p{.53\linewidth}p{.24\linewidth}}
Claims&Completed argument&Boundary retained\\\hline
MT1--MT7&The original ideal relation, exact slope change, all floor-error
support costs and one uniform constant give literal \texttt{Link\allowbreak Mu}.&
Actual \texttt{Link\allowbreak DIV}, fullness and positive rank; the family itself.\\
HX1--HX4&All old host inputs are installed on the original field and
prime pair; the cohomological rescaling preserves $401\to400$.&
$\mathrm{VWinAt}(\chi,401)$; exact classical prime-pair leaves.\\
KA1--KA4&All Case-I \texttt{GD} fields, two-prime valuation aggregation,
finite kernel-orbit mass and original coefficient return give
$\mathrm{PropGK}\Rightarrow\mathrm{CorGK}\Rightarrow\mathrm{CountI}$.&
The actual uniform congruence integer/one-character producer.\\
KA5--KA6&Case-III vacuity at a power-of-two-order target propagates through
every height and window consumer.&Three remaining cohomological inputs
in (DA30.4); the separate global family.\\
ZW1--ZW4&The full weighted-slope and numerical error argument maps literal
\texttt{Chain\allowbreak Tied} into the exact $1/2+(e_D+e_L+7)\epsilon$ window.&
Production of the arithmetic ledger upstream.\\
AQ1--AQ5&An arbitrarily late positive-part integral bound is enough for
zero exclusion; integer values suffice for pointwise windows after
elementary interpolation.&No new arithmetic supplier of either weaker
contract is claimed.\\
\end{longtable}
\end{center}

The metric proof uses the actual identity $N_u=I_uM_u$ already stored
in \texttt{Fam\allowbreak P}, with
$I_u=\prod_w\mathfrak p_w^{-\lfloor e_ws_{u,w}\rfloor}$.
The comparison metric is unchanged by definition, so
\[
 \mu(N_u)-\mu(M_u)
       -2\sum_w\ell_we_ws_{u,w}
      =2\sum_w\ell_w(\lfloor e_ws_{u,w}\rfloor-e_ws_{u,w}).
\]
After every actual exceptional prime family is included, MT proves the
uniform constant $2(\log4+5+10|C_\Delta|)$ on the original $X''$ scale.
This removes an entire independent interface; it does not supply the
upstream divisibility or spectral theorem.

ZW similarly removes a numerical obligation without inventing an
arithmetic shortcut.  If $E$ is the original error bracket and
$R=1+|K|+|c_{HH}|+|e_K|$, its proof gives
\[
 \mathrm{ErrX5at}\le200R^3E\,x^{(e_D+e_L+2)\epsilon}
                  \le E\,x^{(e_D+e_L+3)\epsilon}
\]
after one uniform threshold.  WB then gives the additional four epsilon
units.  The geometric slope inequality is a supplied ledger field;
the finite numerical deduction from that field is now fully written.

\subsection{The remaining global obligations}

The table classifies the still-unsupplied original interfaces; many
conditional sublemmas are already complete.  Shared dependencies listed
in the right column are work-sharing information, not automatically
sufficient hypotheses for the entire left-hand contract.
The dictionary and row-system entries refine the ten interfaces in
(DA30.2); they are not extra independent first-level inputs.

\begin{center}
\begin{longtable}{p{.22\linewidth}p{.08\linewidth}p{.59\linewidth}}
Original obligation&Class&Exact remaining work\\\hline
Compatible design&C&One actual Core/layout, coefficient embeddings,
$\mathrm{Over/RecOK/LedOK}$, rows and lattices, with all requested
theta/local data.  Initial zero theta or empty toric data are not the
required final installation.\\
Local/theta/operator dictionary&D&Identify the actual maps, local models,
nonzero theta lifts and critical toric ratios used by the same
$\mathrm{NTh/NChoice/ChoiceSide/LocAll}$ family.  Their remaining
nonvanishing, convergence, separation and comparison arguments are C.\\
Actual rows and spectrum&C&Original $\mathrm{St/ColPin/R3.Read}$,
rational integral realization, retained rows, and original operator
actions.  Canonical representation irreducibility alone does not
construct these.\\
$\mathrm{LinkN}$&C&Actual local Galois denominator table and its exact
type/theta/rank bindings.  LN proves the full conditional family chain,
including the restoration cost $d_j\log|G|$.\\
$\mathrm{LinkEta}$&C&Compatible twist coverage at every actual local
place; nonempty places alone are insufficient.\\
$\mathrm{LinkIRR}$&C&One original row system with injective label count,
all joins and row count, yielding the precise retained defect.\\
Small-level rank&D&Locate and match the actual projection-surjectivity
supplier $z_L$ and original $\mu_0$; the subsequent finite-dimensional
argument is B.\\
$\mathrm{LinkDIVt}$&C&Original rational period, specified isometries,
heights, ideal shifts, support and word-height variation, all on the
same target and operators.\\
Fullness/rank and global slope&D/B/C&Match the common real-place second
layer, mass ratio and genus-volume normalization.  The same source
\texttt{link\_\allowbreak gl} feeds both interfaces; count this arithmetic
matching once.  Object identification is D, finite installation is B,
and a substantive comparison not already supplied remains C.\\
$\mathrm{LinkG}$&C&Simultaneous original lattice lengths, orbit coverage,
mean-kappa bound and uniform gain; previously proved local statements
must use the one actual global instance.\\
$\mathrm{LinkS}$&C&Normalized finite-word, Poisson, image/norm, kernel,
mass and Gram identities with the exact $\Gamma$, $D_*$, $D_{\rm all}$
and exponential error terms.\\
\end{longtable}
\end{center}

The original class-inclusion and target-class definitions are B
identifications, recorded separately from the C construction of their
objects.  These definitions and the complete argument lists occur in
the pinned source-anchor inventory.  No large package is assigned A
merely because its name contains ``literature'' or a familiar theorem.

\subsection{The remaining counting obligations}

\begin{center}
\begin{longtable}{p{.20\linewidth}p{.08\linewidth}p{.61\linewidth}}
Branch&Class&Original unresolved producer and smaller substeps\\\hline
Count I&C&\texttt{One\allowbreak Char} requires actual unramified-coefficient
automorphy and the Kato input.  The rank-three/rank-two congruence theorem
must choose one nonzero integer before every later coefficient prime.
Its final orbit/Frobenius/valuation installation into \texttt{Prop\allowbreak GK}
is B.  Exact Kato and polarized-adjoint literature matching is D until
the original representations and hypotheses are certified.  The
already matched QX adjoint application does not itself certify this
different rank-three use.\\
Count II&C&On each of QX's two actual Links, one and the same ideal must
satisfy all six \texttt{Hida\allowbreak Props\allowbreak Alt} fields and \texttt{P1core}, above
one common coefficient threshold.  QX and YX already give
$\mathrm{PMExAlt}\Rightarrow\mathrm{CountII}$.
Matching a candidate published support theorem is D.\\
Count III&C&For unrestricted targets: actual \texttt{Seed\allowbreak IVAlt},
tensor-factor tower, \texttt{Chain\allowbreak Rank\allowbreak Four}, residual descent and seeded
chain reduction, then the full original Case-III promodular witness.
Finite tower/descent portions are B.  This separate branch is absent
from (DA30.4), while the common support input remains.\\
Count V&C/D&Original rank-three support/automorphy \texttt{Thm\allowbreak S6\_\allowbreak s},
host rigidity and uniform host catalogue.  The original frame/length
installation, \texttt{Host\allowbreak Van\allowbreak E}, and \texttt{Host\allowbreak Xpp} are B after the
already supplied cohomology and CorX inputs.  Rank-four statements
cannot be imported without a checked D bridge.\\
\end{longtable}
\end{center}

The detailed count inventory records the exact order of all existential
and universal choices.  Some of its edges are explicitly only proposed
reuse paths: for example, a separate proof of the rank-three catalogue
does not require a rank-four-to-rank-three bridge.  Likewise separate
existence of a Hida ideal and a support ideal does not give their
required common witness.

\subsection{Published inputs and stronger optional routes}

Seven applications were freshly checked against primary sources in this
round: the quadratic analytic class-number normalization; Bombieri's
exceptional-factor density theorem; Heath--Brown's hybrid density
specialization; Thorne's exact polarized adjoint vanishing; Caraiani's
conjugate-self-dual local compatibility and temperedness; Hida's
rank-two fixed-central control application; and complete-local
connectedness after the dimension bounds.  HX1 adds the precisely
specialized least-prime theorem.  Their applications are A in these
scopes, not broader arithmetic existence theorems.

The literature inventory also retains 101 entries from 22 earlier
classical-input registries.  These were inventoried, not all independently
rechecked in this round.  Their original written hypothesis certificates
remain the authority for the exact scope.  An empty convenience-summary
field in the new inventory does not erase the detailed original record.
Where an old certificate is reused for a different rank, polarization,
local condition or coefficient-prime range, that new match is D.

In particular, the rank-two Hida input used for the Case-II seed does not
supply the rank-four six-field Hecke ideal.  Allen--Newton--Thorne's
residually reducible automorphy theorem is a candidate support source,
but its residual-Schur, constituent and local hypotheses and its
$\ell>3$ range have not been matched to both actual Links.  It is not
an automatic A closure of \texttt{P1core}.

The backlog keeps optional stronger routes separate: universal
Harish--Chandra interfaces beyond the selected large-gap route,
internal proofs of already matched named classical theorems, all-rank
counting packages, unchanged-level versions, and the stronger one-level
symmetric predicate.  Their absence does not become an additional
requirement of the selected sufficient route.  Every one of the 245
inherited continuation records has an explicit disposition, including
compound records whose closed and still-open subscopes differ.

The complete disposition contains 245 records in 107 scope groups.
Their primary classes are A: 3, B: 10, C: 118 and D: 114; independently,
215 records remain applicable to the combined RH/Dirichlet-GRH project,
25 concern optional stronger routes, and 5 are closed in their stated
scope.  These are bookkeeping counts of inherited records, with overlap
and compound scopes, not counts of independent mathematical obligations
or measures of proof completion.  In particular, ``applicable'' does not
mean required separately for ordinary RH: the Case-III specialization
above still applies.  The three class-A records in this continuation
inventory are not the same inventory as the seven freshly rechecked
published applications plus HX1.

\subsection{A weaker terminal contract worth keeping available}

AQ1--AQ4 show that pointwise square-root control is stronger than the
analytic endpoint needs.  For each fixed target and each $\theta$ in
a sequence decreasing to $1/2$, it is enough to fix one $\delta>0$
such that for every $\eta>0$ and every $Y_0$ there is
$Y\ge\max(2,Y_0)$ with
\[
 \int_Y^{Y^{1+\delta}}\max(H_\chi(x),0)x^{-\theta-1}\,dx
                        \le Y^\eta.                              \tag{DA30.5}
\]
The positive part is essential; a signed integral could hide a positive
peak through cancellation.  The family and $\delta$ precede $\eta$
and the choice of window.  AQ's gamma localization and nearest-pole
argument prove the full implication, including possible real poles.

AQ5 also shows that integer upper bounds on a fixed-stretch window
are enough: on $[n,n+1]$ the affine interpolation $L_n$ satisfies
$0\le H_\chi-L_n\le1/(8n)$, because
$H_\chi''(x)=-1/x$ there.  Slightly shrinking the stretch keeps the
endpoints inside the integer window.

Neither observation currently supplies an arithmetic producer.  The
audit therefore records (DA30.5) as an alternative research contract,
not as a completed shortcut around the common-family obligations.
The older ``direct gap'' contract is treated the same way: bundling
the unresolved arithmetic into one predicate does not reduce its work.

\subsection{The next branch by contribution to completion}

The next default investigation should be the \emph{actual shared
local/theta/operator installation}, with the shared real-place
second-layer/mass normalization investigated in parallel when possible.
These inputs are used by design existence, denominator comparisons,
period/height control, local gains and the spectral theorem.  Resolving
their exact source dictionary is a concrete D task and determines which
large C proofs really remain.  Its output must be a literal common
instance and an exact downstream consumer; another conditional theorem
assuming that instance would not retire the producer.

Among bounded follow-ups, the actual level-projection surjectivity
$z_L$ is a promising target once its supplier is located.  The original
Count-V frame/length and vanishing installations, and the
\texttt{One\allowbreak Char} plus rank-three-number to \texttt{Prop\allowbreak GK} join, are
other finite B chains.  Their value is to expose or remove whole
consumer interfaces, not to increase the number of lemma tags.

The rank-four ordinary support construction remains a substantial
high-value task on both RH and GRH routes, but it is no longer the
automatic next choice solely because version 0.29 ended there.
The separate Case-III seed/promodular route is a priority for completing
unrestricted Dirichlet GRH after taking account of its absence from
the selected ordinary-RH route.

The full human-readable proof is still unfinished.  This version gives a
connected, classified route to that goal, completes five useful groups
of arguments, and makes the remaining original-object obligations
reviewable without changing the assumption about the checked formal
proof.  Source and text checks in this round do not constitute a new
Lean replay.


\section{The prime sum and the exact missing arithmetic theorem}

Write $\Lambda$ for the von Mangoldt function and $\gamma$ for Euler's constant. For $x\ge1$, define
\begin{equation}\label{eq:G}
G(x)=\sum_{n\le x}\Lambda(n)\left(\frac{x}{n}-1\right)
       -x\log x+(1+\gamma)x.
\end{equation}
For a nonprincipal primitive Dirichlet character $\chi$, define
\begin{equation}\label{eq:Gchi}
G_\chi(x)=\sum_{n\le x}\Lambda(n)\chi(n)\left(\frac{x}{n}-1\right)
                  +x\frac{L'}{L}(1,\chi).
\end{equation}
The plus sign in the last term is essential. For the primitive principal character, write $G_1=G$. A principal character of larger modulus must first be replaced by its primitive inducer, or its missing Euler factors must be included explicitly. These normalizations agree with the code.\cite{charfamily,level}

If a fixed primitive character $\psi$ has order $o$, let $\mathcal C_\psi$ be the set of distinct primitive characters inducing its powers. The source's cyclic target is
\begin{equation}\label{eq:Hpsi}
H_\psi(x)=G(x)+\frac{1}{3o+1}
                   \sum_{\substack{\xi\in\mathcal C_\psi\\\xi\ne1}}\Rea G_\xi(x).
\end{equation}
Every weight is positive, and the target $\psi$ occurs. The family stays fixed while the auxiliary arithmetic data and the windows vary. For ordinary RH, $o=1$, $\mathcal C_\psi=\{1\}$, and $H_\psi=G$.

\paragraph{Arithmetic window theorem (source result; proof not yet reconstructed here).}
For each fixed $\psi$, there are constants $A_\psi>0$ and $\varepsilon_0>0$ such that for every $0<\varepsilon<\varepsilon_0$ and every $Y_0$, there exists $Y\ge\max(2,Y_0)$ with
\begin{equation}\label{eq:window-arith}
H_\psi(x)\le x^{1/2+A_\psi\varepsilon}
                    \qquad(Y\le x\le Y^{3/2}).
\end{equation}
The source obtains the exponent $1/2+(e_D+e_L+7)\varepsilon$. One may enlarge its fixed coefficient to a positive $A_\psi$. The existence quantifiers for $e_D,e_L$ precede the choice of $\varepsilon$; this order is indispensable.\cite{coregm,topwin,conclusion}

Calling \eqref{eq:window-arith} a lemma does not make its proof easy. It contains the main achievement of the release. Sections~\ref{sec:scalar}--\ref{sec:remaining} identify its constituent arguments and separate the deductions proved here from the arithmetic inputs. \textsf{[AN1, SC2]}

\section{A complete readable analytic endpoint}\label{sec:analytic}

\status{Proved here from standard complex analysis and elementary analytic properties of zeta. The arithmetic window theorem is a hypothesis, not a consequence of this section.}

\begin{theorem}\label{thm:ordinary}
Suppose that, for every $\eta>0$, there are arbitrarily large $Y$ such that
\begin{equation}\label{eq:ordinary-window}
G(x)\le x^{1/2+\eta}\qquad(Y\le x\le Y^{3/2}).
\end{equation}
Then RH holds.
\end{theorem}

The source proves a more general weighted-character version.\cite{zerofree,landau} The proof below specializes it to ordinary RH, removing the real-zero case. \textsf{[AN2]}

\subsection{The Mellin transform records every off-line zero}

For $\Rea s>1$, direct integration gives
\[
\int_1^\infty\left(\frac{x}{n}-1\right)_+x^{-s-1}\,dx
       =\frac{n^{-s}}{s(s-1)}.
\]
Termwise integration of the absolutely convergent series therefore gives, with $g(u)=G(e^u)$,
\begin{equation}\label{eq:Phi}
\begin{split}
\Phi(s)&:=\int_0^\infty g(u)e^{-su}\,du
        =\int_1^\infty G(x)x^{-s-1}\,dx\\
&=-\frac{\zeta'/\zeta(s)}{s(s-1)}
  -\frac{1}{(s-1)^2}+\frac{1+\gamma}{s-1}.
\end{split}
\end{equation}
Since $\zeta'/\zeta(s)=-1/(s-1)+\gamma+O(s-1)$, the singularity at $s=1$ cancels. In $\Rea s>0$, the remaining poles are precisely the zeta zeros. A zero $\rho$ of multiplicity $m_\rho$ gives a simple pole of $\Phi$ with residue
\begin{equation}\label{eq:res}
\frac{m_\rho}{\rho(1-\rho)}\ne0.
\end{equation}
There are no positive real poles: $\zeta(s)>0$ for real $s>1$, while $\zeta(s)<0$ for $0<s<1$, by the alternating eta-series formula. Standard continuation and the Euler product locate all poles under consideration in $0<\Rea s\le1$.

The crude bound $|g(u)|\ll e^{2u}$ is enough for what follows. For example, $\Lambda(n)\le\log n$ bounds the uncentered sum by $O(x\log^2(2x))$, which is $O(x^2)$.

\subsection{Localize to one logarithmic window}

For $\Rea z>2$, differentiation under the integral yields
\begin{equation}\label{eq:Ak}
A_k(z):=\frac{(-1)^k}{k!}\Phi^{(k)}(z)
       =\int_0^\infty g(u)\frac{u^k}{k!}e^{-zu}\,du.
\end{equation}
Fix $\sigma=336$ and set $k=\lceil420U\rceil$ for $U\ge1$. The elementary estimate
\begin{equation}\label{eq:tails}
\int_{(0,\infty)\setminus[U,3U/2]}
\frac{u^k}{k!}e^{(2-\sigma)u}\,du\ll\sigma^{-k}
\end{equation}
is proved at the end of this section. Thus the derivative integral is concentrated in the logarithm of $[Y,Y^{3/2}]$, where $U=\log Y$.

Fix $1/2<\theta<1$, and use \eqref{eq:ordinary-window} with $\eta=\theta-1/2$. On its arbitrarily late logarithmic windows, $g(u)\le e^{\theta u}$. Because $g$ is real and $|g|+g=2g_+$, for every real $t$ we have
\[
|A_k(\sigma+it)|+A_k(\sigma)
\le2\int_0^\infty g_+(u)\frac{u^k}{k!}e^{-\sigma u}\,du.
\]
Inside the window use $u^ke^{-(\sigma-\theta)u}/k!\le(\sigma-\theta)^{-k}$, which follows from the exponential series. Outside use \eqref{eq:tails} and the crude growth bound. This proves
\begin{equation}\label{eq:onesided}
|A_k(\sigma+it)|+A_k(\sigma)
       \ll(U+1)(\sigma-\theta)^{-k}.
\end{equation}
The extra real derivative $A_k(\sigma)$ is what permits a one-sided upper bound. A bound on $|G|$ is unnecessary.

\subsection{A nearest pole contradicts the localized estimate}

Suppose that a pole has real part greater than $\theta$, and let $\mathcal Z$ denote the pole set. Set
\[
D=\min\{\sigma-\theta,\ \inf_{\rho\in\mathcal Z}|\sigma-\rho|\}.
\]
There is a pole $\rho_1$ whose horizontal distance $\sigma-\Rea\rho_1$ from the vertical line $\Rea s=\sigma$ is less than $D$. If $D=\sigma-\theta$, use the hypothesized pole. Otherwise a nearest pole to $\sigma$ exists in a compact disk; it is nonreal, so its horizontal distance is strictly less than its distance $D$.

Choose $t$ near $\Im\rho_1$ so that the nearest pole $\rho$ to $z=\sigma+it$ is unique, at distance $d<D$. Only finitely many poles matter in a fixed neighborhood. Each pair excludes at most one value of $t$ at which their distances agree; pairs with equal imaginary parts have different horizontal distances since both real parts lie left of $\sigma$. Hence such a choice is possible.

Choose $d<r<D$ with no other pole in the closed radius-$r$ disk about $z$. Subtract the principal part $c/(s-\rho)$ from $\Phi$ and apply Cauchy's derivative estimate on that disk. At the real center, the radius-$r$ disk is pole-free. Consequently
\[
A_k(z)=\frac{c}{(z-\rho)^{k+1}}+O(r^{-k}),\qquad
A_k(\sigma)=O(r^{-k}),\qquad c\ne0.
\]
Combining with \eqref{eq:onesided} gives
\[
\frac{|c|}{d^{k+1}}\ll r^{-k}+(U+1)(\sigma-\theta)^{-k}.
\]
This is impossible along the arbitrarily late windows: $k\asymp U\to\infty$ and $d<r<D\le\sigma-\theta$. Thus there is no zero to the right of $\theta$. Letting $\theta\downarrow1/2$ and using the functional equation proves Theorem~\ref{thm:ordinary}.

\paragraph{Proof of the kernel estimate \eqref{eq:tails}.}
On the upper tail multiply the integrand by $e^{14(u-3U/2)}\ge1$, and on the lower tail by $e^{14(U-u)}\ge1$. Extending both integrals to $(0,\infty)$ gives
\[
\int_{(0,\infty)\setminus[U,3U/2]}
\frac{u^k}{k!}e^{-334u}\,du
\le\frac{e^{-21U}}{320^{k+1}}+\frac{e^{14U}}{348^{k+1}}.
\]
After multiplying by $336^k$, the first term is at most $e^{1/20}/320$, using $\log(336/320)\le1/20$ and $k\le420U+1$. The second is at most $1/348$, using $\log(348/336)\ge1/29$ and $k\ge420U$. This proves the estimate with an absolute constant.\hfill$\square$

\subsection{The extension to positive character families}

For $\chi\ne1$, the same calculation gives
\begin{equation}\label{eq:Mchi}
M_\chi(s):=\int_1^\infty G_\chi(x)x^{-s-1}\,dx
=-\frac{L'/L(s,\chi)-sL'/L(1,\chi)}{s(s-1)}.
\end{equation}
The apparent pole at one is removable. If $H=\sum a_\chi\Rea G_\chi$ with $a_\chi>0$, then its transform is
\[
\Phi_H(s)=\frac12\sum_\chi a_\chi\{M_\chi(s)+M_{\bar\chi}(s)\}.
\]
At a zero $\rho$ of any constituent it has residue
\begin{equation}\label{eq:weightedres}
\frac{b_\rho}{\rho(1-\rho)},\qquad
b_\rho=\frac12\sum_\chi a_\chi
             \{m_\chi(\rho)+m_{\bar\chi}(\rho)\}>0.
\end{equation}
Coincident zeros cannot cancel. \textsf{[AN3]} Real zeros $\beta\in(0,1)$ give positive residues.\cite{poles}

The preceding nearest-pole proof still works if some pole has horizontal distance less than $D$. Otherwise every pole has horizontal distance at least $D$. A pole to the right of $\theta$ forces $D<\sigma-\theta$; a pole at distance $D$ from $\sigma$ must then be the unique nearest pole, which is real: $\beta=\sigma-D$. Its residue $c$ is positive. Choose $D<r<\sigma-\theta$ small enough that $\beta$ is the only pole in the closed radius-$r$ disk about $\sigma$. Then
\[
A_k(\sigma)=cD^{-k-1}+O(r^{-k})\qquad(D<r<\sigma-\theta).
\]
Taking $t=0$ in \eqref{eq:onesided} gives the same contradiction. Therefore a bound $H(x)\le Cx^\theta$ on arbitrarily late windows excludes every constituent zero with real part greater than $\theta$.

Apply this to \eqref{eq:Hpsi} and \eqref{eq:window-arith}, choosing $\varepsilon$ after fixing a proposed zero $\Rea\rho=\beta>1/2$, with $A_\psi\varepsilon<\beta-1/2$. This excludes all such zeros. Functional equations reflect left-hand zeros into the excluded right half-plane. For an imprimitive character, the additional factors $1-\chi^*(p)p^{-s}$ are nonzero for $\Rea s>0$ because $|\chi^*(p)p^{-s}|<1$. Hence the primitive result gives Dirichlet GRH. These final transitions also appear explicitly in the source.\cite{critline}

\section{Why the exponent is exactly a square root}\label{sec:scalar}

\status{The following scalar lemma and proof are complete. The fact that arithmetic constructions satisfy its hypothesis is still an input. \textsf{[SC1]}}

\begin{proposition}\label{prop:scalar}
Let $x\ge3$, $\varepsilon>0$, $C\ge0$; let $n,n_T,n_2,m$ be positive integers; let $b,k>0$ and $0\le a\le2xb$. Put
\[
d_1=\frac{n_T}{n},\quad d_2=\frac{n_2}{n},\quad Q=n_Tmb,
\quad D=\frac{Qd_2}{2}.
\]
Assume
\begin{gather}\label{eq:scales}
d_1\ge\tfrac12x^{-\varepsilon/8},\quad x^{-\varepsilon}\le d_2\le1,\\
x^{1/2-\varepsilon}\le m\le x^{1/2+\varepsilon/2},\qquad
xm\le k\le n^2x^{1/2-\varepsilon}.
\end{gather}
Define
\begin{equation}\label{eq:B}
B=n_Ta\left(1+\frac{m^2}{k}\right)
+m\left(2nmbx^\varepsilon+
\frac{n_2kbx^\varepsilon}{n(nx^{-\varepsilon})}\right).
\end{equation}
If a real $V$ satisfies
\begin{equation}\label{eq:scalarhyp}
D\bigl(V-4x^{1/2+\varepsilon}\bigr)
       \le x^{C\varepsilon}B+\frac{2Q}{x},
\end{equation}
then
\begin{equation}\label{eq:scalarcon}
V\le26x^{1/2+(C+21/8)\varepsilon}.
\end{equation}
If also $26\le x^{3\varepsilon/8}$, then $V\le x^{1/2+(C+3)\varepsilon}$.
\end{proposition}

\begin{proof}
Set $Z=x^{1/2+13\varepsilon/8}$. Dividing the four terms of $B$ by $Q$ gives
\[
\frac{B}{Q}=\frac{a}{mb}+\frac{am}{bk}
                  +\frac{2mx^\varepsilon}{d_1}
                  +\frac{d_2kx^{2\varepsilon}}{d_1n^2}.
\]
The respective bounds are
\begin{align*}
\frac{a}{mb}&\le 2x/m\le2x^{1/2+\varepsilon}\le2Z,\\
\frac{am}{bk}&\le2xm/k\le2\le2Z,\\
\frac{2mx^\varepsilon}{d_1}&\le4x^{1/2+\varepsilon/2+\varepsilon+\varepsilon/8}=4Z,\\
\frac{d_2kx^{2\varepsilon}}{d_1n^2}
&\le2x^{1/2-\varepsilon+2\varepsilon+\varepsilon/8}\le2Z.
\end{align*}
Thus $B\le10QZ$. Divide \eqref{eq:scalarhyp} by $Q/2$ to get
\[
d_2\bigl(V-4x^{1/2+\varepsilon}\bigr)
       \le20x^{C\varepsilon}Z+4/x\le22x^{C\varepsilon}Z.
\]
Since $d_2\ge x^{-\varepsilon}$, this gives
\[
V\le4x^{1/2+\varepsilon}
       +22x^{1/2+(C+21/8)\varepsilon}
\le26x^{1/2+(C+21/8)\varepsilon}.
\]
Finally absorb $26$ using $26\le x^{3\varepsilon/8}$.
\end{proof}

This is the numerical content of \texttt{proposition\_size\_bound}.\cite{size}
In particular, the exact normalized error is
\begin{equation}\label{eq:balance}
\boxed{\frac{B}{D}=\frac{2a}{bd_2}\left(\frac1m+\frac{m}{k}\right)
  +\frac{4mx^\varepsilon}{d_1d_2}
  +\frac{2kx^{2\varepsilon}}{d_1n^2}.}
\end{equation}
In the actual construction, $a\ge xb$ and $d_1,d_2\le1$. The $1/m$ part of the first term is at least $2x/m$, while the second term is at least $4m$. The opposing costs $x/m$ and $m$ explain the square-root chain length. The other parameter bounds ensure that the remaining contributions are no larger than $x^{1/2+O(\varepsilon)}$. This calculation does not establish the divisibility or norm estimates that make \eqref{eq:scalarhyp} true.

\subsection{How the scalar bound becomes the fixed target bound}

The source first proves an arithmetic gap bound for every translate of a selected finite set. Its exact Fourier identity is \texttt{ch\_balance}; averaging over the subgroup produces the positive coefficient $D=n_Tmbd_2/2$ in \eqref{eq:scalarhyp}.\cite{balance,chcore}

At this stage $V=G_{\rm mA}(x)-\Sigma(x)$, where $G_{\rm mA}$ is a weighted sum of retained character ramps and $\Sigma$ is a sum of absolute values of omitted auxiliary ramps. The allowance $4x^{1/2+\varepsilon}$ pays for shifts of the cutoffs, approximate Fourier weights, and the linear terms at $s=1$. The source then identifies $G_{\rm mA}=H_\psi$ and proves
\[
\Sigma(x)\le4o\,x^{1/2+\varepsilon}.
\]
With $C=e_D+e_L+3$, where the construction gives $e_D+e_L\ge0$, Proposition~\ref{prop:scalar} yields
\begin{align*}
H_\psi(x)&\le x^{1/2+(e_D+e_L+6)\varepsilon}
                +4o\,x^{1/2+\varepsilon}\\
&\le(1+4o)x^{1/2+(e_D+e_L+6)\varepsilon}
 \le x^{1/2+(e_D+e_L+7)\varepsilon}
\end{align*}
once $x^\varepsilon\ge1+4o$. These identifications and estimates are explicit in \texttt{GmAEq}, \texttt{SigALe}, and \texttt{CoreGm}.\cite{coregm,gm,sigma}

% Standalone section fragment. Uses only macros/environments already in v0.1.
% Source claim identifiers are resolved in period_lattice_provenance.json.
\section{From finite periods to one integral target}\label{sec:period-lattice}

\status{The descent, gluing, integrality, and norm arguments below are proved here. Their inputs concerning the actual theta kernels and local divisibility are identified explicitly. For coefficient maps, ``rational'' means defined over the coefficient number field $K$.}

\subsection{The objects behind the maps}

The finite models are spaces of equivariant functions on decorated hermitian lattices. A point $d$ consists of a unimodular lattice in a positive definite hermitian space over a CM field, together with specified nonzero residue functionals and residue frames. A rational unitary group, together with a finite level subgroup, acts on these points. Once finiteness of the orbit set and stabilizers has been established, choose representatives $d_c$ and put $\Gamma_c=\operatorname{Stab}(d_c)$. For a coefficient representation $(V,\rho)$,
\[
 \mathcal A=\{f:f(gd)=\rho(g)f(d)\}
       \simeq\bigoplus_c V^{\Gamma_c},\qquad
 \mathcal L_{\Delta}=\{f\in\mathcal A:f(d)\in\Delta(d)\text{ for every }d\}.
\]
Here $\Delta(d)$ and its larger companion $\nabla(d)$ are explicit integral highest-weight modules; their polynomial definitions are recorded in the source notes. The initial source $C_u$ is the joint kernel of the residue-unit averaging operators
\[
 \operatorname{Nm}_P=\sum_{a\in\kappa(P)^\times}\langle a\rangle,
 \qquad C_u=\bigcap_{P\in S_{T,u}}\ker\operatorname{Nm}_P,
 \qquad M_u^{\rm ram}=C_u\cap\mathcal L_{\Delta}.
\]
Here $S_{T,u}$ is the set of primes dividing the chosen element $\pi_T(u)$, rather than necessarily every prime above the rational type prime $T_u$. The working lattice $N_u$ is a further arithmetic choice in $C_u$; it must not be identified with $M_u^{\rm ram}$. In the final divisibility bridge it is exactly the coordinate pullback of the record's $\texttt{NuLat}$. Its required local conditions remain inputs here. \textsf{[PL1, PL2, PL4, PL12]}

A kernel acts by the finite formula
\[
 (T_kf)(d')=\sum_c\frac{k(d_c,d')f(d_c)}{|\Gamma_c|}.
\]
Let $\mathcal T_{u,t}$ denote the chain of these operators, including the final evaluation or coupled-kernel step. The period in the code is
\begin{equation}\label{eq:pl-period}
 P_{u,t}(\phi,F)=\sum_c\frac{
 b_{\rm alg}\bigl(\operatorname{br}_u(\mathcal T_{u,t}\phi)(d_c),F(d_c)\bigr)}{|\Gamma_c|}.
\end{equation}
The coefficient pairing is bilinear: for polynomial coefficients it is a fixed normalization of $\sum_\alpha\alpha!f_\alpha g_\alpha$. Thus \eqref{eq:pl-period} is a $K$-bilinear finite sum, with visible stabilizer denominators. Constructing and estimating the individual theta kernels is additional work. \textsf{[PL3, PL15]}

On the row side, finite-group character and level projectors select a subspace $W_t$ of the analogous function space. Its lattice is
$M_t=W_t\cap\mathcal L_{\nabla}$. The selected theta-adjoint image gives an orthogonal quotient $\pi_t:W_t\twoheadrightarrow Q_t$, with its quotient metric. Its lattice is the \emph{image} $\pi_tM_t$. Rationality and common-image theorems for the theta maps produce isomorphisms
\[
 \iota_t:Q_t\overset\sim\longrightarrow Q:=Q_{t_0},
 \qquad \iota_{t_0}=1.
\]
Concretely, if $J_t$ is the injective complex theta map on $Q_t$ and all $J_t$ have the same image, then $\iota_t=J_{t_0}^{-1}J_t$; the assertion that this matrix lies in $K$ is a substantive input. \textsf{[PL5, PL6]}

\subsection{Canonical descent and comparison scalars}

\begin{proposition}[Descent of the period]\label{prop:pl-descent}
Suppose $P_{u,t_0}(\phi,F)=0$ whenever $\pi_{t_0}F=0$. There is a unique $K$-bilinear map $B_u:C_u\times Q\to K$ satisfying
\begin{equation}\label{eq:pl-base}
 B_u(\phi,\pi_{t_0}F)=P_{u,t_0}(\phi,F).
\end{equation}
It can be computed from any $K$-linear right inverse $r$ of $\pi_{t_0}$ by $B_u(\phi,y)=P_{u,t_0}(\phi,ry)$.

Fix an embedding $\sigma:K\hookrightarrow\mathbb C$. Suppose nonzero scalars $\xi_{u,t}\in\mathbb C$, with $\xi_{u,t_0}=1$, satisfy
\begin{equation}\label{eq:pl-complex}
 \xi_{u,t}\,\sigma(P_{u,t}(\phi,F))
     =\sigma\bigl(B_u(\phi,\iota_t\pi_tF)\bigr).
\end{equation}
Suppose also that whenever $\xi_{u,\mu(w)}=\sigma(a)$ for $a\in K$, one has $\operatorname{ord}_w(a)=0$. Then there are $\eta_{u,t}\in K^\times$, with $\eta_{u,t_0}=1$, such that
\begin{equation}\label{eq:pl-factor}
 P_{u,t}(\phi,F)=\eta_{u,t}B_u(\phi,\iota_t\pi_tF),
 \qquad \operatorname{ord}_w\eta_{u,\mu(w)}=0.
\end{equation}
\end{proposition}

\begin{proof}
Since $F-r\pi_{t_0}F$ is in the kernel of $\pi_{t_0}$, the displayed definition gives \eqref{eq:pl-base}. The same observation proves independence of $r$, and surjectivity of $\pi_{t_0}$ proves uniqueness.

If $\xi_{u,t}=\sigma(a)$, choose $\eta_{u,t}=a^{-1}$. Otherwise every value of $P_{u,t}$ must be zero: a nonzero value in \eqref{eq:pl-complex} would express $\xi_{u,t}$ as the image under $\sigma$ of the quotient of two elements of $K$. The right side is then zero too, so choose $\eta_{u,t}=1$. These choices give all the claims, including the finite-place unit condition. No nonvanishing assumption on $B_u$ was used.
\end{proof}

This is the source's actual right-inverse construction and scalar descent. The scalar is inverted when passing to \eqref{eq:pl-factor}. The kernel-vanishing input follows from the theta identity once $\ker\pi_t\subseteq\ker\Theta_t$ is known. The underlying linear algebra is short: if $\Theta_tW_t\subseteq E_0$, $\Theta_t^*E_0\subseteq W_t$, and $F\in W_t$ is orthogonal to $\Theta_t^*E_0$, then
$\|\Theta_tF\|^2=\langle F,\Theta_t^*\Theta_tF\rangle=0$.
Rationality and the theta identity still have to be supplied for the actual models. \textsf{[PL7, PL8, PL18]}

\subsection{Gluing the common lattice and using divisibility}

\begin{theorem}[One integral target for every word]\label{thm:pl-integral}
Let $M_t\subset W_t$ and $N_u\subset C_u$ be full $\OK$-lattices, and let $\pi_t$ and $\iota_t$ be as above. For each finite place $w$, choose a row parameter $\mu(w)$, equal to $t_0$ outside a finite set. Put
\begin{equation}\label{eq:pl-R}
 R=\bigcap_w\bigl(\iota_{\mu(w)}\pi_{\mu(w)}M_{\mu(w)}\bigr)_w
       \subset Q,
\end{equation}
where the subscript means localization over $\mathcal O_{K,w}$, inside the $K$-vector space. Then $R$ is a full lattice, uniquely determined by
$R_w=(\iota_{\mu(w)}\pi_{\mu(w)}M_{\mu(w)})_w$.

Assume \eqref{eq:pl-factor}. Let $\sigma_u(w)$ be finitely supported real numbers such that, for $\phi\in N_u$ and $F\in M_{\mu(w)}$,
\begin{equation}\label{eq:pl-val}
 P_{u,\mu(w)}(\phi,F)=0
 \quad\text{or}\quad
 \operatorname{ord}_wP_{u,\mu(w)}(\phi,F)\ge\sigma_u(w).
\end{equation}
Define
\begin{equation}\label{eq:pl-shift}
 k_u(w)=\lfloor\sigma_u(w)\rfloor,\qquad
 I_u=\prod_w\mathfrak p_w^{-k_u(w)},\qquad N'_u=I_uN_u.
\end{equation}
Then
\begin{equation}\label{eq:pl-integral}
 B_u(N'_u,R)\subseteq\OK,\qquad
 B_u:N'_u\longrightarrow R^\vee:=\operatorname{Hom}_{\OK}(R,\OK).
\end{equation}
The same $R$ works for every $u$.
\end{theorem}

\begin{proof}
Write $L_w=\iota_{\mu(w)}\pi_{\mu(w)}M_{\mu(w)}$ and $L_0=\pi_{t_0}M_{t_0}$. Images under surjective linear maps are full lattices, so all these lattices are full. They equal $L_0$ outside finitely many $w$.

Choose $0\ne d\in\OK$ with $dL_0\subseteq L_w\subseteq d^{-1}L_0$ for the exceptional $w$, and with $d$ in each exceptional prime. The finite torsion module $A=d^{-1}L_0/dL_0$ is the direct sum of its prime-primary components. In the component at each prime dividing $d$, prescribe the submodule given by $(L_w)_w/(dL_0)_w$. Pull back their direct sum to $d^{-1}L_0$. The resulting full lattice has exactly the required localizations; at primes not dividing $d$, both bounds localize to $L_0$.

A lattice is the intersection of its localizations: if $x$ fails to belong, the proper ideal $\{a\in\OK:ax\in L\}$ lies in a maximal ideal $\mathfrak p$, which prevents $x\in L_{\mathfrak p}$. This proves that the constructed lattice is \eqref{eq:pl-R}, as well as uniqueness.

For integrality, \eqref{eq:pl-factor} and the unit condition transfer \eqref{eq:pl-val} to $B_u$ on $N_u\times L_w$. The same bound holds after localization, since the new denominators are $w$-units. As valuations are integers, the bound is at least $\lceil\sigma_u(w)\rceil$. Multiplication by $I_u$ subtracts $k_u(w)$ from this bound. Thus $B_u(N'_u,R)$ is integral at every finite place, hence belongs to $\OK=\bigcap_w\mathcal O_{K,w}$. Finally, \eqref{eq:pl-R} contains no $u$.
\end{proof}

The proof also allows any finitely supported integer $k_u(w)\le\lceil\sigma_u(w)\rceil$; the released construction uses the floor. Its $\sigma_u(w)$ is the ramification index times the normalized gain recorded at $w$. A negative gain shrinks the lattice in \eqref{eq:pl-shift}, so ``shifted'' is more accurate than ``enlarged.'' \textsf{[PL9--PL12]}

\subsection{Why the quotient also gives the right norm bound}

At an infinite place $v$, write the base period as
$P_{u,t_0}(\phi,F)=P_{\rm end}(T_r\cdots T_1\phi,F)$.
Give $Q_{t_0}$ the quotient metric from $W_{t_0}$. Suppose
\[
 |P_{\rm end}(z,F)|\le p_v\|z\|\|F\|,\quad
 \|T_j\|\le a_{j,v},\quad
 \|\iota_{t_0}q\|_Q^2\ge s_v\|q\|_{Q_{t_0}}^2,
 \qquad p_v,a_{j,v},s_v>0.
\]
Here the metric on the common target $Q$ can differ from the quotient metric, even though $\iota_{t_0}$ is the identity in the chosen coordinates. For $y\in Q_v$, take the orthogonal, minimum-norm lift $F$ of $\iota_{t_0}^{-1}y$. Substituting this $F$ in \eqref{eq:pl-base} proves
\begin{equation}\label{eq:pl-norm}
 \|B_u\|_v^2\le\frac{p_v^2\prod_j a_{j,v}^2}{s_v}.
\end{equation}
There is no loss from descent through the quotient itself: a functional vanishing on its kernel has exactly the same norm as the descended functional.

In the source, $p$ ranges over a finite nonempty set of spectral pieces,
$\kappa_v>0$, and every $\lambda_{j,p,v}>0$. The eigenvalue comparison gives
$s_v=\kappa_v\min_p\prod_j\lambda_{j,p,v}>0$.
Consequently, if $B_u\ne0$, its squared-norm height satisfies
\begin{equation}\label{eq:pl-height}
 h(B_u)\le\operatorname{mean}_v\!\left(
 \log p_v^2+\sum_j\log a_{j,v}^2-\log\kappa_v
       -\min_p\sum_j\log\lambda_{j,p,v}\right).
\end{equation}
The mean is over all infinite places, with the source's equal weights for totally complex $K$. Formula \eqref{eq:pl-norm} also applies when $B_u=0$; then one uses positive majorants without taking $\log0$. \textsf{[PL11, PL20]}

What remains is now localized: prove the rational common-image and theta comparison identities; establish \eqref{eq:pl-val} for the chosen $N_u$; and bound the chain norms and eigenvalues at every infinite place. This section supplies the complete algebraic passage from those inputs to integral maps with a common target and the required form of height estimate. It does not yet establish their quantitative arithmetic inputs.

\section{The three coefficient lattices and the normalized neighbor map}
\label{sec:edge-lattices}

\status{The integral coefficient-lattice comparisons retained in [JR12] and [LW4] are proved here for the source's polynomial model. This includes full rank, change of lattice coordinates, and the neighbor-map inclusion. An elementary polynomial argument replaces the source's semistandard-basis route to the weight estimate. The pointwise local tests and selected-symbol rank are reconstructed below; Sections~\ref{sec:frame-geometry}--\ref{sec:edge-assembly} now supply their compatible data and the full Type-II edge theorem.}

\subsection{An integral polynomial model and its scaling}

Let $W$ be a DVR with fraction field $Q$, uniformizer $p$, and residue field of size $q=p^{n'}$. Let $\mathcal O$ be a DVR with fraction field $K$, containing the images of $W$ under embeddings $\tau_s=\iota\sigma^s$, $0\le s<n'$. Assume $w\in\mathcal O$ satisfies $w^2=p\ne0$. Put $k=m+1$. At each $s$ use the decreasing integer weight
\[
 \lambda_s=
 \begin{cases}
   \lambda^0_s,&s\text{ even},\\
   (-\lambda^0_{s,k},\ldots,-\lambda^0_{s,1}),&s\text{ odd},
 \end{cases}
 \qquad
 A=\sum_s\sum_{i=1}^{r}\lambda_{s,i}.
\]
Here $0\le r\le k$, and $A$ is an integer of either sign. These are exactly the source's \texttt{lamAt} and \texttt{Ar}; the reversal at odd $s$ is part of the definition. \textsf{[EL1]}\cite{edge-coefficients,edge-operators}

For a decreasing weight $\lambda$, write $e=\lambda_k$ and $\nu_i=\lambda_i-e\ge0$. Let $\mathcal D_\nu(R)$ consist of polynomials in a $k\times k$ matrix $X$, with column multidegree $\nu$, invariant under every substitution
$C_j\mapsto C_j+uC_i$ with $i<j$. Invariance means a polynomial identity in $u$, over the coefficient ring $R$. Put $T(g)f(X)=f(gX)$, so $T(g)T(h)=T(hg)$, and set
\[
 \rho_\lambda(g)=\det(g)^{-e}T(g^{-1}),\qquad
 H_\lambda(g\mathcal O^k)=\rho_\lambda(g)\mathcal D_\nu(\mathcal O)
 \subset \mathcal D_\nu(K).
 \tag{EL.1}
\]
The column conditions are finitely many linear equations on the finite free module of polynomials of column degree $\nu$. Thus $\mathcal D_\nu(\mathcal O)$ is finite free, and clearing denominators gives
$K\mathcal D_\nu(\mathcal O)=\mathcal D_\nu(K)$.
Row substitutions commute with column additions. Hence $\rho_\lambda(U)$ preserves $\mathcal D_\nu(\mathcal O)$ for $U\in\operatorname{GL}_k(\mathcal O)$, including when $e<0$, since $\det U$ is a unit. This proves that (EL.1) is independent of the basis $g$ and is a full lattice. It also proves
$H_\lambda(a\Lambda)=\rho_\lambda(a)H_\lambda(\Lambda)$.
\textsf{[EL2]}

\begin{lemma}[The needed weight bound without tableaux]
\label{lem:edge-weight}
For $f\in\mathcal D_\nu(K)$, every monomial of $f$ has total degree at most
$N_r=\sum_{i=1}^r\nu_i$ in the first $r$ rows. Consequently, for
$D_w=\operatorname{diag}(w^{-1}I_r,I_{k-r})$,
\[
 w^{\,\sum_{i=1}^r\lambda_i}H_\lambda(g\mathcal O^k)
 \subseteq H_\lambda(gD_w\mathcal O^k).
 \tag{EL.2}
\]
\end{lemma}
\begin{proof}
The cases $r=0,k$ are immediate. Otherwise write the generic matrix as
$X=\left(\begin{smallmatrix}B&C\\ E&F\end{smallmatrix}\right)$, with $B$ of size $r$.
In the polynomial ring localized at $\det B$, column invariance gives
\[
 f\!\left(\begin{pmatrix}zI_r&0\\0&I\end{pmatrix}X\right)
 =
 f\!\left(\begin{pmatrix}zB&0\\ E&F-EB^{-1}C\end{pmatrix}\right).
 \tag{EL.3}
\]
Indeed, multiply on the right by
$\left(\begin{smallmatrix}I&-B^{-1}C\\0&I\end{smallmatrix}\right)$.
Its coefficients are independent of $z$. On the right, every entry involving $z$ lies in the first $r$ columns, whose total prescribed degree is $N_r$.
Localization is injective, so the degree bound holds before localization as well. For integral $f$, it follows term by term that
$w^{N_r}T(D_w)f\in\mathcal D_\nu(\mathcal O)$.
Finally $\det D_w=w^{-r}$ and $\sum_{i\le r}\lambda_i=N_r+re$.
Applying the membership criterion
$h\in H_\lambda(g\mathcal O^k)\iff\det(g)^eT(g)h\in\mathcal D_\nu(\mathcal O)$
therefore gives (EL.2), with no sign restriction on $e$ or on $\sum_{i\le r}\lambda_i$.
\textsf{[EL3--EL4]}
\end{proof}

\subsection{Edges, midpoint lattices, and integral polars}

An edge $e=(L^-,L^+)$ consists of full $W$-lattices with
$pL^+\subseteq L^-\subseteq L^+$ and $|L^+/L^-|=q^r$.
Smith normal form gives, after reordering and absorbing units into basis vectors,
\[
 L^-=gW^k,\qquad
 L^+=g\operatorname{diag}(p^{-1}I_r,I)W^k.
 \tag{EL.4}
\]
To see the count, the elementary divisors divide the uniformizer $p$, so each is a unit or a unit times $p$. The quotient cardinality forces exactly $r$ of the second kind.
Writing $g_s=\tau_s(g)$, the extended source, midpoint, and target lattices are
\[
 g_s\mathcal O^k,\qquad
 \Lambda_s(e):=\tau_s(L^-)\mathcal O+w\,\tau_s(L^+)\mathcal O
      =g_sD_w\mathcal O^k,\qquad
 g_sD_w^2\mathcal O^k .
 \tag{EL.5}
\]
This is the source's \texttt{LeAt}: its factor $w^{-1}p$ equals $w$.
\textsf{[EL5]}

Use separate polynomial variables at different $s$ and the ambient space
$V=\bigotimes_s\mathcal D_{\nu_s}(K)$, where $\nu_{s,i}=\lambda_{s,i}-\lambda_{s,k}$.
This is the source's literal \texttt{Vsp}, with full lattices
\[
 H(L)=\bigotimes_s H_{\lambda_s}(\tau_s(L)\mathcal O),
 \qquad
 H_e=\bigotimes_s H_{\lambda_s}(\Lambda_s(e)).
\]
Tensoring (EL.2) and then applying it a second time gives
\[
 w^A H(L^-)\subseteq H_e,\qquad
 w^A H_e\subseteq H(L^+).
 \tag{EL.6}
\]
For a full lattice $M\subset V$, let
$M^\vee=\{\phi\in V^*:\phi(M)\subseteq\mathcal O\}$; for a lattice in $V^*$ use the analogous polar in $V$.
An $\mathcal O$-basis and its coordinate functionals show
$(M^\vee)^\vee=M$.
Set
\[
 R_e=w^{-A}H_e,\qquad R'_e=w^{-A}H_e^\vee.
\]
Equation (EL.6) gives
$H(L^-)\subseteq R_e$ and $H(L^+)^\vee\subseteq R'_e$.
Moreover the pairing
\[
 \beta_e(x,\phi)=\phi(p^A x),\qquad
 \beta_e(w^{-A}b_i,w^{-A}b_j^*)=\delta_{ij},
 \tag{EL.7}
\]
is perfect over $\mathcal O$. Here the sole scalar calculation is
$w^{-A}p^Aw^{-A}=1$, valid for every integer $A$ because $w^2=p\ne0$.
In particular,
$p^A R_e=(R'_e)^\vee$.
\textsf{[EL6]}

\subsection{The three slots and the neighbor operator}

For an integer $t\ge0$, define the ball slot, the target test lattice, and the dual slot by
\[
 B_t(L)=w^{-t}H(L)\cap\bigcap_{e:\,e^-=L}R_e,\qquad
 T_t(L)=w^{-t}H(L)^\vee\cap\bigcap_{e:\,e^+=L}R'_e,\qquad
 P_t(L)=T_t(L)^\vee .
 \tag{EL.8}
\]
Thus the three coefficient slots are $H(L),B_t(L),P_t(L)$;
$T_t(L)$ itself is a lattice in the dual vector space.
The preceding inclusions imply
\[
 \boxed{\quad
 w^tH(L)\subseteq P_t(L)\subseteq H(L)
       \subseteq B_t(L)\subseteq w^{-t}H(L).
 \quad}
 \tag{EL.9}
\]
Indeed $H(L)^\vee\subseteq T_t(L)\subseteq w^{-t}H(L)^\vee$, and taking polars gives the first two inclusions; the other two follow directly from (EL.8).
Every slot is full and finite free, of the same rank $\dim_K V$: each is sandwiched between two full lattices over the DVR.
This proves all four comparisons in \texttt{RestLocChain}, including the two multiplication-by-$w^t$ statements. \textsf{[EL7]}

For any edge $e=(L^-,L^+)$, (EL.7)--(EL.8) give
\[
 \boxed{\qquad p^A B_t(L^-)\subseteq P_t(L^+).\qquad}
 \tag{EL.10}
\]
The proof is only evaluation: $x\in B_t(L^-)$ lies in $R_e$, while every $\phi\in T_t(L^+)$ lies in $R'_e$, so $\phi(p^Ax)\in\mathcal O$.
The same argument gives
$p^A(w^cR_e)\subseteq w^cP_t(L^+)$ for $c\ge0$.
It therefore also transports the prescribed errors at unwanted neighboring targets, as required in the joint-rank cancellation. \textsf{[EL10]}

For clarity, these constructions respect the group action used on sections.
The identity after (EL.1), applied to all $\tau_s(a)$, sends $H(L)$ and $H_e$ to $H(aL)$ and $H_{ae}$.
On dual vectors use the contragredient action $\phi\mapsto\phi\circ\rho(a)^{-1}$.
It follows by taking intersections and polars that all three slots transform by $\rho(a)$.
Scalar character values in $\mathcal O^\times$ preserve them as well. \textsf{[EL8]}

Finally, on a tuple of lattices $\mathbf L$, the source's neighbor operator is precisely
\[
 (Y_vF)(\mathbf L)
  =p^{A_v}\sum_{L'\,\to\, L_v}
       F(L_1,\ldots,L_{v-1},L',L_{v+1},\ldots).
 \tag{EL.11}
\]
The sum is finite: its possible $L'$ lie between $pL_v$ and $L_v$ and hence correspond to subspaces of the finite vector space $L_v/pL_v$.
Tensor (EL.10) with the unchanged slots at the other places and add the terms. This proves
$Y_v(\text{ball at }v)\subseteq(\text{dual at }v)$.
For equivariant sections, multiplication by $a$ bijects the incoming edges, so reindexing (EL.11) proves equivariance of the result. Sections extended by zero off lattice tuples remain zero there: either there is no edge at $v$, or the unchanged bad coordinate makes every summand vanish.
These observations prove \texttt{RestNat} and its assembly with the coefficient comparisons. They require neither positivity of $A_v$ nor integrality of $p^{A_v}$ by itself. \textsf{[EL9]}

The integral polynomial component of [JR12]/[LW4] is thus discharged. What remains elsewhere is the quantitative production and separation of the local test vectors, the finite-orbit and stabilizer input for the global section space, and the other arithmetic and analytic steps of the RH/GRH argument.

\section{The polynomial that isolates one edge}\label{sec:symbols}

\status{The localizer, its elementary keys, the finite correction process, and the degree accounting are proved here. Section~\ref{sec:path-generators} now proves the integral path-generator span and its key-symbol image. Together these establish realization for the stated symbol family. Its quantitative dimension is supplied by Sections~\ref{sec:coarse-lift}--\ref{sec:rank-closure}; the local parameter construction is now supplied in Sections~\ref{sec:frame-geometry}--\ref{sec:edge-assembly}; Sections~\ref{sec:initial-membership}--\ref{sec:word-transport} now prove incoming membership.} \textsf{[SY1--SY8]}

Work in the Type-II setting: $p\ge13$, $2\le r<k$, an even Frobenius period $n'$, and a coefficient DVR $O$ with uniformizer $w$ and $w^2=p$. A frame consists of disjoint increasing blocks of even positions $a_{j,1}<\cdots<a_{j,r}=z_j$, periodic under $j\mapsto j+J$ by addition of $n'$. Put $Q_j=p^{z_{j+1}-z_j}$ and let $Q=p^{G'}$ bound all Frobenius exponents used below. The full frame and size conditions are those in the ledger; in particular, they are stronger than just periodicity.

The chart substitution $\Phi_g$ is an integral invertible change of coordinates followed by multiplication of the normal variables by $w$; let $\Pi(g)$ be its rational plane. At an even position the normal variables have column $b\ge r$; at an odd position they have $b<r$. Its inverse exists over $\operatorname{Frac}(O)$. Write $\mathcal T$ for the integral polynomials in the remaining, tangential variables. All variables here have positions in the selected set $TT$.

\subsection{A finite-field telescoping identity}

Let $M_{j,I}$ be the $r$ by $r$ minor, in columns $I$, of the rows obtained by raising row $a_{j,l}$ to the power $p^{z_j-a_{j,l}}$. Let $I_0=\{0,\ldots,r-1\}$, $\Delta_j=M_{j,I_0}$, and let $\mathcal I$ consist of the $r(k-r)$ sets obtained by replacing one column of $I_0$. Define
\begin{align}
R_{\rm loc}&=\Delta_0\prod_{I\in\mathcal I}
 \left(\prod_{j=0}^{J-1}\Delta_j^{Q_j-1}
             -\prod_{j=0}^{J-1}M_{j,I}^{Q_j-1}\right),\label{eq:sy-localizer}\\
R_0&=\prod_{j=0}^{J-1}\Delta_j^{\alpha_j},\qquad
\alpha_j=\mathbf1_{j=0}+r(k-r)(Q_j-1).
\end{align}
These are exactly the source's \texttt{Rloc} and \texttt{R0}.\cite{symbol-source,symbol-elementary}

\begin{lemma}[Plane isolation]\label{lem:sy-isolation}
For every other rational $r$-plane chart $g'$, all coefficients of $\Phi_{g'}R_{\rm loc}$ are divisible by $w$. In its own chart, $\Phi_gR_{\rm loc}\equiv\Phi_gR_0\pmod w$.
\end{lemma}
\begin{proof}
Translate all relevant lifted positions by a common multiple of $n'$ to make them nonnegative; this changes neither differences nor coordinates. Reduce modulo $w$, and put $q_0=p^{n'}$. In the other chart the rows, after their indicated Frobenius powers, have the form $T_jH^{(p^{z_j})}$, with $H\in\operatorname{Mat}_{r,k}(\mathbf F_{q_0})$. Hence
\[
M_{j,I}=\det(T_j)\det(H_I)^{p^{z_j}},\qquad
\sum_j p^{z_j}(Q_j-1)=p^{z_0}(q_0-1).
\]
For every nonzero $\det(H_I)$ the product over $j$ in \eqref{eq:sy-localizer} is consequently the same polynomial $\prod_j\det(T_j)^{Q_j-1}$. If $\det(H_{I_0})=0$, the initial factor $\Delta_0$ vanishes. Otherwise the right block of $H$ is nonzero, since the two planes differ. Cramer's rule gives a nonzero exchanged-column minor. The corresponding difference factor vanishes. The signs from ordering exchanged columns disappear because $Q_j-1$ is even. In the original chart each exchanged minor contains a normal variable, so only the $R_0$ term survives modulo $w$.
\end{proof}

We need slightly more of the last assertion. Expanding the product by nonempty subsets $S\subseteq\mathcal I$, each error term loses $|S|(Q_j-1)$ copies of $\Delta_j$. In the original chart, expanding its exchanged minors writes it as
\begin{equation}\label{eq:sy-local-error}
w^e X^{\nu}\prod_j(\Phi_g\Delta_j)^{\alpha_j-b_j}H,
\quad e\ge1,\quad |\nu|\le e,\quad 0\le b_j\le e,\quad H\in\mathcal T.
\end{equation}
Indeed, each exchanged determinant uses a normal column; after substitution a normal variable raised to $p^a$ contributes $w^{p^a}X^{p^a}$. Thus its normal degree is paid for by its $w$-power, and $e\ge\sum_j|S|(Q_j-1)\ge b_j$. This proves the source's more precise localizer expansion, not merely its reduction modulo $w$. \textsf{[SY2--SY3]}

\subsection{First-order keys and why top determinants occur}

For each position $s$, take the $r$ most recent even selected positions $a_l<s$, and let $T_s$ have entries $X_{a_l,b}^{p^{s-a_l}}$, $b<r$. Set $D_s=\det T_s$. A \emph{key} of order $d$ is an integral polynomial $K$ with a normal-homogeneous leading term $K_0$ of degree $d\le2Q$, such that
\[
K_0\in w^dO[X],\quad \Phi_gK-K_0\in w^{d+1}O[X],\quad
\Phi_{g'}K\in w^dO[X]\quad\hbox{for every chart }g'.
\]

For each normal variable $X_{s,b}$ there are order-one keys and tangential integral coefficients $H_l$ satisfying
\begin{equation}\label{eq:sy-key-identity}
wX_{s,b}D_s=\sum_lH_lK_{0,l}.
\end{equation}
Here is the construction. If $s$ is even, use the bordered $(r+1)$-minor with first row at $s$, the older $r$ rows raised to $p^{s-a_l}$, and columns $I_0\cup\{b\}$. Its leading term is $(-1)^rwX_{s,b}D_s$; all other terms contain $w^{p^{s-a_l}}\in w^2O$. If $s$ is odd, pair each older primal row, raised to that power, with the dual row at $s$. Their leading terms form $wT_s(X_{s,0},\ldots,X_{s,r-1})^{\mathsf t}$; multiplying by $\operatorname{adj}(T_s)$ proves \eqref{eq:sy-key-identity}. On any residue-plane chart the first keys vanish because $r+1$ vectors lie in an $r$-space, and the second because a vector and a covector annihilating that space pair to zero. This checks all key conditions. \textsf{[SY4]}

If a label monomial is $X^m$, the source's digit condition implies $m_{s,b}\le Q$. There are at most $kQ$ normal factors at each $s$. Multiplying \eqref{eq:sy-key-identity} for those factors therefore expresses
\[
w^{\deg_N m}X^m\prod_{s\in TT}D_s^{kQ}
\]
as a sum of products of key leading terms, of total key order $\deg_Nm$, with coefficients in $\mathcal T$. The unused determinant powers are nonnegative. This proves the residual-monomial assertion \texttt{RestMon}; no division by $D_s$ is used.

\subsection{Finite correction, with the exact determinant budget}

A term of level $\ell=i+\sum_a d_a$ has formal symbol and integral raw lift
\begin{align*}
S_{\mu,\lambda}&=w^i\prod_aK_{0,a}\,(\Phi_gR_0)^c
            \prod_j(\Phi_g\Delta_j)^\lambda\prod_sD_s^\mu H,\\
U_{\mu,\lambda}&=w^i\prod_aK_a\,R_{\rm loc}^c
            \prod_j\Delta_j^\lambda\prod_s\widehat D_s^\mu\widehat H.
\end{align*}
The hats denote integral polynomials before substitution, with $\Phi_g\widehat D_s=D_s$ and $\Phi_g\widehat H=H$. Thus $S_{\mu,\lambda}\in w^\ell O[X]$, while $\Phi_{g'}U_{\mu,\lambda}\in w^{\ell+c}O[X]$ on other planes.

Fix the desired precision $L=o+c$. Expanding $\Phi_gU-S$ and using \eqref{eq:sy-local-error} and \eqref{eq:sy-key-identity} gives, modulo $w^L$, finitely many new symbols at levels $\ell+E$, with
\begin{equation}\label{eq:sy-spares}
1\le E<L-\ell,\qquad
\mu'=\mu-(2Q+1)E,\qquad \lambda'=\lambda-E.
\end{equation}
For completeness, the coefficient estimate behind the budget is this: in a key error of order $d$, a monomial of normal degree $v$ has coefficient divisible by $w^x$, where $x=\max(d+1,v)$. Put $E=x-d$. Then $v\le(2Q+1)E$. Replacing its $v$ normal factors by \eqref{eq:sy-key-identity} uses at most that many copies of each $D_s$. A localizer error spends at most $E$ frame determinants. Products add both the gained levels and these costs; surplus determinants go into $H$. Terms already of level $L$ may be discarded. The coefficient bound uses integrality of $K$ before substitution, not just the displayed leading congruence. \textsf{[SY5]}

Starting at $\ell\ge o$ with $\mu=(2Q+3)c$, $\lambda=c$ leaves enough determinants at every step: the invariants are $\mu+(2Q+1)\ell$ and $\lambda+\ell$, and every correction has level at most $L-1$. Downward induction on $L-\ell$ now produces an integral $\widetilde U$ with
\[
\Phi_g\widetilde U\equiv S\pmod {w^{o+c}},\qquad
\Phi_{g'}\widetilde U\equiv0\pmod {w^{o+c}}\quad(\Pi(g')\ne\Pi(g)).
\]
Explicitly, replace each child symbol by its already corrected lift and subtract those lifts from $U$; the new error is in $w^L$. The process is finite, so it requires no completion or infinite limit. Projection to a fixed multidegree preserves all congruences because substitutions preserve positions.

\subsection{The actual symbol family and its degree budget}

The source's fixed multiplier is
\begin{equation}\label{eq:sy-sptrue}
\mathrm{SPtrue}=R_0^c\prod_j\Delta_j^c
                    \prod_{s\in TT}\widehat D_s^{(2Q+3)c+kQ}.
\end{equation}
Its substitution is integral and has normal degree zero. For an allowed path-generator polynomial $P$ of normal degree at most $T$, the two generator assertions proved in Section~\ref{sec:path-generators} put its segment factors into the algebra generated by key symbols. The label factor is handled by SY4. Apply the finite correction to each normal-homogeneous component $\sigma_o$ of $P\Phi_g(\mathrm{SPtrue})$, $0\le o\le T$. It yields integral polynomials $F_o$ of the required multidegree satisfying
\begin{gather*}
\Phi_gF_o\equiv w^o\sigma_o\pmod {w^{o+c}},\qquad
\Phi_{g'}F_o\equiv0\pmod {w^{o+c}}\quad(\Pi(g')\ne\Pi(g)),\\
\sum_{o=0}^T\sigma_o=P\Phi_g(\mathrm{SPtrue}).
\end{gather*}
These are the actual separated realizations. The source's intermediate \texttt{Fs} are instead rational preimages of $\sigma_o$ under $\Phi_g$; distinguishing the two avoids an unwarranted assertion that those intermediate preimages are integral. \textsf{[SY6]}

Finally let $E'_t$ be the degree of \eqref{eq:sy-sptrue} at position $t$. Frame rows occur once per period, and at most $N_1+1$ top determinants involve a fixed position. Consequently
\begin{align*}
E'_t&\le\{c(1+rkQ)+c\}Q
 +(N_1+1)\{(2Q+3)c+kQ\}rQ\\
&\le15c(N_1+rk)Q^2=:E_0
\end{align*}
when $k\le c$, $N_1\le k$, $r\ge2$ and $Q\ge3$. Also $E'_t=0$ off $TT$. The source's size condition $a\ge Ck^3r^2NQ^5cY$, $C\ge30$, makes $E_0\le a$, so subtracting this degree budget from the prescribed degree is legitimate. Together with Section~\ref{sec:path-generators}, this completes the realization and degree part of \texttt{SymbolRealisableBound}. The quantitative selected-symbol dimension is now proved in Section~\ref{sec:rank-closure}. Sections~\ref{sec:integral-weyl}--\ref{sec:digit-tests} now supply the pointwise dual-test constructions and incoming word closure. Sections~\ref{sec:selected-dimension}--\ref{sec:rank-closure} prove the supplied-data rank and degree theorem; Sections~\ref{sec:frame-geometry}--\ref{sec:edge-assembly} now supply the full local Type-II parameter matching. Global abundance of the required orbits remains separate. \textsf{[SY7--SY8]}

\section{Integral path generators and Frobenius columns}
\label{sec:path-generators}

\status{Both generator assertions retained in [SY6] are proved below. The same integral filtered basis supplies the \texttt{SPAN} input to the local Gram estimate. The detailed companion proves the combinatorial and Frobenius calculations with the literal source conventions. No integral decomposition into rational isotypic summands is asserted.}

\subsection{One integral basis for all prefix conditions}

Give letters $x_0,\ldots,x_{N-1}$ in dimension $n$ vector or covector signs and degrees $d_i$. A signed Pieri path is a sequence of decreasing integer tuples $\lambda^0=0,\ldots,\lambda^N$: a vector step increases by a horizontal strip of size $d_i$, and a covector step decreases by a horizontal strip of size $d_i$. Set $h_i(c)=\#\{a:c\le\lambda^i_a\}$. Each column starts at $1_{\varnothing}$; when $c>0$, a height increase applies exterior multiplication by $x_i$, a decrease applies contraction. For $c\le0$ interchange those operations. Use the source's right-wedge signs. Call the resulting component $K_c(I_c)$. Choose $|I_c|=h_N(c)$ for $c>0$, and $|I_c|=n-h_N(c)$ otherwise. Complement the negative-column subsets. Require the increasing lists in these modified columns to be weakly increasing along rows. Put
\[
 \Psi_{\lambda,I}=\prod_cK_c(I_c).
 \tag{IP.1}
\]
Only finitely many factors differ from $1$. These are exactly the source's \texttt{Kcol}, \texttt{IsSt}, and \texttt{Psi}. \textsf{[IP1]}\cite{path-basis,path-symbol}

\begin{lemma}[Integral filtered path basis]
The polynomials (IP.1) form a basis over $\mathbb Z$. Over any integral domain $R$, the simultaneous conditions
\[
 v_i(F)\ge\theta_i,\qquad M_i(F)\le U_i
 \tag{IP.2}
\]
are equivalent to an $R$-linear combination of the same basis elements satisfying $\lambda^i_n\ge\theta_i$ and $\lambda^i_1\le U_i$ for every prefix $i$. Here $v_i,M_i$ are the minimum and maximum weights of the generic conjugated rank-one torus on that prefix.
\end{lemma}
\begin{proof}
For canonical terminal subsets, order the last letter's exponents with coordinate $n$ first for vectors, coordinate $1$ first for covectors. The leading component is
\[
 \pm\Psi_{\text{preceding path, canonical}}
       \prod_a x_{N-1,a}^{|\lambda^N_a-\lambda^{N-1}_a|}.
\]
For fixed endpoint its exponent determines the preceding endpoint; induction proves canonical-path independence over every field. A generic matrix separates the other terminal choices: its endpoint-weight coefficient is the canonical product times a product of minors. Jacobi's complementary-minor identity turns this into a common determinant power times flag minors. Their diagonal leading monomials recover the weakly ordered tableau rows, so they are distinct. Taking a maximal endpoint in dominance proves independence of all (IP.1), in every characteristic. \textsf{[IP2]}

Row insertion of a weakly increasing word gives
\[
 \sum_{\mu/\lambda\text{ horizontal strip of size }d}s_n(\mu)
   =s_n(\lambda)\binom{n+d-1}{d},
\]
where $s_n$ counts standard terminal choices. Determinant shifts permit negative weights; complementing and reversing columns gives the downward version. Iteration counts exactly the ordinary monomials. The square integer coefficient matrix is thus nonsingular modulo every prime, so its determinant is $\pm1$. \textsf{[IP3]}

For the full prefix, every transformed column product has coordinate weights between $\lambda_n$ and $\lambda_1$. Conversely, a maximal bad endpoint in a basis expansion survives in its generic endpoint-weight component by the preceding independence argument. A good endpoint cannot dominate it. For an earlier prefix, cut each exterior recurrence there: this proves containment for good paths. Iterating the same counting identity from that endpoint gives exactly the dimension of the constrained prefix space times the tail monomial dimension. Thus containment is equality. Uniqueness of basis expansion handles all prefixes simultaneously. Passing to the fraction field and back retains integral coefficients. \textsf{[IP4]}
\end{proof}

Each column is multilinear exactly in its used letters, and letter $i$ occurs in $d_i$ columns. A positive column's height is at most its preceding vector occurrences. Cutting at prefix $i$ also gives
\[
 v_i(K_c(I_c))\ge q(c,h_i(c)),\quad
 \lambda^i_n=\sum_cq(c,h_i(c)),\quad
 q(c,h)=\mathbf1_{c>0,\,h=n}-\mathbf1_{c\le0,\,h<n}.
 \tag{IP.3}
\]
Indeed the intermediate exterior weights are $0,1$ for positive columns and $0,-1$ for negative ones. These facts are precisely the remaining \texttt{SPAN} clauses. \textsf{[IP5]}

\subsection{The actual generators and their keys}

For the source frame segment, an \texttt{MGen0} generator is $AG$ in blocks of dimensions $r,K$. Apply the lemma separately. The defining zero-gap inequalities give, for every surviving path pair $(\nu,\kappa)$,
\[
 \kappa^{i+1}_1\le\nu^i_r\quad(\text{vector}),\qquad
 \kappa^i_1\le\nu^{i+1}_r\quad(\text{covector}).
\]
On the initial preceding-frame block, $G$ has degree zero and $A$ uses only vectors. These same inequalities follow there. They are exactly the missing cross-block strip inequalities, so $(\nu,\kappa)$ is a genuine joined path. Multiplying integral expansions proves \texttt{GeneratorSpanInclusion}. \textsf{[IP6]}

For a joined path put $S_c=K_c^{\nu}(I'_c)K_c^{\kappa}(I''_c)$. The literal split-product identity gives $P=\prod_cS_c$. Let selected lifted positions be $s_0<\cdots<s_{N-1}\le z$, with $p^{z-s_0}\le Q=p^{G'}$. A step is costly when its height rises from $h\ge r$ or falls from $h\le r$. Set
\[
 a_c=\sum_{i\text{ costly}}p^{z-s_i}<2Q.
 \tag{IP.4}
\]
Every monomial of $S_c$ has this weighted normal degree. The source bound $N+r\le K$ keeps negative columns strictly above height $r$. \textsf{[IP7]}

Work over the source DVR $\mathcal O$, with $w^2=p$, $p\ge13$. Construct $F_c^g$ using the same exterior recurrence in the adapted chart-$g$ letters, but raise every component to $p^{s-t}$ between successive used positions, then to $p^{z-t}$ at the end. This is the integral polynomial \texttt{KKc}. Its degree at used position $s_i$ is $p^{z-s_i}$. \textsf{[IP8]}

Here is the essential all-chart calculation. In a moving integral exterior basis, write $VL_s=w^aZ$. Witt Frobenius gives, for $P=p^{t-s}$,
\[
 V^{[P]}L_t=w^{aP}(Z^{[P]}+pR),\qquad R\text{ integral}.
 \tag{IP.5}
\]
Hence an extra factor $w$ becomes two extra factors across a strict gap. The nominal subsets are first-block subsets below height $r$, subsets containing that block above it, and last-block subsets for negative columns. Inductively every component has the predicted order; nonnominal components have one extra. At a costly step, normal terms acquire $w$, while tangential terms require a nonnominal predecessor and retain two extra powers. This proves order $a_c$ in every chart. \textsf{[IP9]}

In its own chart, divide by the predicted power and reduce modulo $w$. Exactly the prescribed block operation survives, with its original exterior sign. Frobenius distributes in the residue field, giving
\[
 \Phi_gF_c^g\equiv w^{a_c}\operatorname{phSeg}(S_c)
                 \pmod{w^{a_c+1}}.
 \tag{IP.6}
\]
Thus each factor is a literal key, with the required selected-position support. For $o=\sum_ca_c$, their leading terms multiply to $w^o\operatorname{phSeg}(P)$, whose normal degree is $o$. This is exactly \texttt{Sy}, proving \texttt{GeneratorImageIsSymbol}. Zero factors cause no exception. \textsf{[IP10--IP11]}

\section{Integral Weyl comparison from an explicit averaging operator}
\label{sec:integral-weyl}

\status{The following proves the integral Fock comparison used in CL.5--CL.8, and supplies the quantitative isotypic-projector lemma needed by the constrained Gram argument. The background is finite-dimensional rational representation theory of split $\mathrm{GL}_k$ in characteristic zero: complete reducibility, highest weights and the elementary PBW weight bound. The integral claims below are proved separately.}

\subsection{The denominator, and why the operator projects}

Let $\mathcal O$ be a DVR of mixed characteristic $(0,p)$ with fraction field $K$.
Suppose operators $E_{ab}$ satisfy the $\mathfrak{gl}_k$ commutator relations, and let $h\ge0$ be an integer. Write
\[
 D_p(k,h)=\binom{k}{2}\lfloor\log_p(h+k-1)\rfloor,
 \qquad E_{ab}^{[j]}=E_{ab}^j/j!;
 \tag{IW.1}
\]
the cases $k=1$ or $h=0$ below permit exponent zero.
For $z\notin S$, $s=|S|$, define
\[
 \Omega_0=1,\quad \Omega_{q+1}=\sum_{i\in S}E_{zi}\Omega_qE_{iz},\quad
 P_{S,z}=\sum_{q=0}^{h}\frac{(-1)^qs!}{(q+s)!q!}\Omega_q.
\]
Compose these steps while adjoining the $k$ indices successively; call the result $I_h$.
If $M$ is an $\mathcal O$-module stable under every off-diagonal divided power, then
\[
 p^{D_p(k,h)}I_hM\subseteq M. \tag{IW.2}
\]
Indeed, commuting left and right multiplications expands the $q$th summand into divided-power words with coefficients
$(-1)^qs!\prod_i\alpha_i!/(q+s)!$, where $\sum_i\alpha_i=q$.
Its denominator is a product of $s$ binomial coefficients with upper entries at most $h+s$.
Legendre's formula gives $v_p\binom nt\le\lfloor\log_p n\rfloor$.
Thus a step costs at most $s\lfloor\log_p(h+k-1)\rfloor$; sum $s=0,\ldots,k-1$. \textsf{[IW1]}\cite{weyl-integral,weyl-form}

Suppose all diagonal eigenvalues on the relevant module lie in $[-h,h]$.
If $E_{ab}u=0$ for $a,b\in S$ and $E_{zz}u=0$, the commutator relations give
\[
 E_{jz}\Omega_{q+1}u=\Omega_{q+1}E_{jz}u
 +(q+1)(q+s+1)\Omega_qE_{jz}u\quad(j\in S).
\]
The displayed coefficients telescope to $c_h\Omega_hE_{jz}u$, where $c_h=(-1)^hs!/((h+s)!h!)$.
Each right-hand string inside this expression reaches $z$-weight $-h-1$, so it vanishes.
The opposite root also kills $P_{S,z}u$: for a vector $y$ of both diagonal weights zero with $E_{jz}y=0$,
$E_{jz}E_{zj}^{m+1}y=-m(m+1)E_{zj}^my$; nilpotence and downward induction give $E_{zj}y=0$.
Induction over $S$ proves that $I_h$ sends weight zero to invariants.
Every correction lies in a root-operator image, and the construction commutes with intertwiners.
Consequently on $\operatorname{End}(U)$, for irreducible $U$, it sends a weight-zero $T$ to $(\operatorname{tr}T/\dim U)1$. \textsf{[IW2]}

\subsection{The actual polynomial lattice and its normalized pairing}

Use a square matrix $X=(X_{ij})$ and a partition $\nu_1\ge\cdots\ge\nu_k=0$.
Let $\nabla$ be the $\mathcal O$-module of polynomials of column multidegree $\nu$ invariant under $X\mapsto XU$ for upper unitriangular $U$. Put
\[
 H=\prod_{j=1}^{k}\det(X_{[1,j],[1,j]})^{\nu_j-\nu_{j+1}},
\]
where $\nu_{k+1}=0$. Let $\Delta$ be the span of divided-power words applied to $H$.
Put $h=\nu_1$ and $V=K\nabla$.
Gaussian decomposition $X=LDU$ shows that the highest line is exactly $KH$: a highest vector is also invariant under lower unitriangular left multiplication.
Complete reducibility gives irreducibility of $V$; its weights have every coordinate between $0$ and $h$.
The coefficient lattice $\nabla$ is graded and divided-power stable, with highest line $\mathcal OH$ since $H(I)=1$.
The highest-weight projection $C$ therefore maps $\nabla$ into $\Delta$.
The Hom lattice $\operatorname{Hom}_{\mathcal O}(\nabla,\Delta)$ is divided-power stable, since its $j$th divided action is
$\sum_{a+b=j}(-1)^b E_{uv}^{[a]}T E_{uv}^{[b]}$.
Apply (IW.2) to $C$. Its invariant projection is $(\dim V)^{-1}1$; multiplication by the integer $\dim V$ proves
\[
 p^{D_p(k,h)}\nabla\subseteq\Delta\subseteq\nabla. \tag{IW.3}
\]
If $h=0$, both lattices equal $\mathcal O$. \textsf{[IW3]}

Define the \emph{symmetric bilinear} Fock form and its exact normalization by
\[
 F(P,Q)=\sum_\alpha P_\alpha Q_\alpha\alpha!,\qquad
 N=F(H,H)\in\mathbb Z_{>0},\qquad B=F/N.
\]
Write $R(g)P(X)=P(gX)$, the source's substitution convention.
Monomial calculation gives transpose adjointness; highest-weight projection and $H(I)=1$ give
\[
 B(R(g^t)H,Q)=B(H,R(g)Q)=Q(g)\qquad(Q\in V). \tag{IW.4}
\]
Thus $B(\Delta,\nabla)\subseteq\mathcal O$.
Conversely, if $B(\Delta,Q)\subseteq\mathcal O$, substitute a generic lower unitriangular $L$ in (IW.4).
The coefficients of $R(L^t)H$ are integral combinations of divided-power words, so $Q(L)$ is integral.
Gaussian decomposition and column homogeneity recover $Q(X)$ after inverting principal minors.
These minors are primitive polynomials; the DVR Gauss lemma removes their denominators.
Hence $Q\in\nabla$. Dualizing full lattices gives
\[
 \Delta^\#=\nabla,\qquad \nabla^\#=\Delta,
 \qquad p^{D_p(k,h)}\nabla\subseteq\nabla^\#\subseteq\nabla. \tag{IW.5}
\]
The expanded note proves the Gaussian recovery and nondegeneracy explicitly. \textsf{[IW5]}

\subsection{Tensor transport and the shared projector bound}

For a dominant integer weight $\lambda$, use $\nu_i=\lambda_i-\lambda_k$ and the determinant twist $\lambda_k$.
For a basis matrix $A$, the source convention is $L(A)=\det(A)^{-\lambda_k}R(A^{-1})\nabla$; for tuples use the tensor product of these lattices.
Transpose adjointness cancels the determinant twists at $A$ and $A^{-t}$.
Tensoring (IW.5), and identifying the linear dual with the opposite representation by the normalized Fock forms, proves
\[
 \begin{aligned}
 p^{\kappa}\,L(A^{-t})&\subseteq\Psi(L(A)^*)\subseteq L(A^{-t}),\\
 \kappa&=\sum_{s:\,h_s>0}\binom{k}{2}\lfloor\log_p(h_s+k-1)\rfloor.
 \end{aligned}
 \tag{IW.6}
\]
This is the actual scalar-normalized tensor pairing, not a choice of arbitrary perfect form. \textsf{[IW6]}
At an edge, the midpoint matrices are $A_s=\sigma^s(g)\operatorname{diag}(w^{-1}I_r,I)$ with $w^2=p$.
The opposite edge gives exactly $A_s^{-t}$; the common $w^{-A_r}$ factors in $R_e$ and $R'_e$ preserve (IW.6).

For completeness, (IW.2) also controls every isotypic projector $e_U$ on a graded, divided-power-stable full lattice $M$ in a split semisimple rational representation.
Here $h$ must bound $|\mu_a-\eta_a|$ for \emph{two} weights of that representation.
For an integral basis $u_i$ of $M\cap U$, let $\mathcal A$ be $I_h$ after weight-zero projection on $\operatorname{Hom}(U,V)$.
Then
\[
 \begin{aligned}
 e_U(y)&=(\dim U)\sum_i\mathcal A(u_i^*\otimes y)(u_i),\\
 p^{D_p(k,h)}e_U(M)&\subseteq M.
 \end{aligned}\tag{IW.7}
\]
Schur's trace identity proves the first formula on each irreducible summand; (IW.2) proves the second. \textsf{[IW4]}
Finally, $h_s+k\le a^2$ gives each exponent at most $k(k-1)\log_p a$.
For the two selected families this implies the source's stated $\kappa_0+\kappa_f\le4Nk^2\log_p a$. \textsf{[IW7]}

\section{The integral complementary basis and its exact count}
\label{sec:complementary-basis}

\status{This supplies the integral basis, complementary-family count and Cartan multiplication used in CL1--CL5. The proofs work over every integral domain, including residue fields; they use polynomial identities, not evaluations at the points of a finite field.}

\paragraph{CB1. A leading-term calculation.}\cite{complement-basis}
Let $R$ be a domain and $X$ an $r\times n$ matrix of indeterminates.
Write $\nabla_\nu(R)$ for polynomials of column multidegree $\nu$ invariant under every generic operation $C_j\mapsto C_j+tC_l$, $l<j$.
Order variables column first, then row, with earlier variables larger.
If $X^e$ leads a nonzero member, then
\[
 \sum_a e_{a,t}=\nu_t,\qquad
 \sum_{b\le a}e_{b,j}\le\sum_{b<a}e_{b,j-1}. \tag{CB.1}
\]
Here is the substantive step. For two columns $x,y$, the more general bound
$\deg_t f(x,y+tx)\le s$ implies
$\sum_{i\le a}b_i\le s+\sum_{i<a}a_i$ for their leading exponents.
First retain the homogeneous component fixing each pair's total degree and every other variable's exponent; this preserves the leading monomial and the bound.
At the first row write
\[
 f=\sum_{u=0}^{K}x^uy^{e+K-u}f_u,\qquad f_K\ne0.
\]
After substitution this is $(y+tx)^eQ$, with
$\deg_tQ\le s-e$ and $e\le s$.
Setting $x=1,y=-t$ in $Q$ leaves $f_K(x',y'+tx')$; since $\deg_yQ\le K$, its $t$-degree is at most $s-e+K$.
Induction on rows proves the assertion. Generic invariance gives $s=0$ and hence (CB.1).

\paragraph{CB2. A monic integral basis.}
For an array satisfying (CB.1), make tableau row $t$ from $e_{a,t}$ copies of $a$, in increasing order.
The inequalities say exactly that entries increase strictly down columns.
For each tableau column take the minor using its entries as matrix rows and the first matrix columns; multiply these minors to obtain $b_e$.
The diagonal term leads each minor with coefficient one, so $b_e$ has leading monomial $X^e$.
Every $b_e$ belongs to $\nabla_\nu(R)$, and distinct $e$ give independence.
For any $f$, subtract its leading coefficient times the corresponding $b_e$.
The leading monomial decreases; the fixed multidegree admits finitely many monomials, so this terminates.
Thus the $b_e$ are an integral basis, homogeneous for row weights.
The same integer polynomials give bases after passage to any domain, including reduction modulo a prime.
Products in disjoint blocks likewise give a tensor basis.

\paragraph{CB3. Counting without a dimension theorem.}
For a decreasing integer tuple $\lambda$ of length $k$, put
\[
 D_k(\lambda)=\prod_{i<j}\frac{\lambda_i-\lambda_j+j-i}{j-i},\qquad D_0=D_1=1.
\]
Tableaux of shape $\nu=\lambda-\lambda_k$ with entries at most $k$ have this count.
Indeed, deleting the entries $k$ leaves a shape $\mu$ satisfying
$\nu_i\ge\mu_i\ge\nu_{i+1}$.
Induct on $k$, using $x_i=\nu_i+k-i$ and $y_i=\mu_i+k-1-i$ and the identity
\[
 \sum_{x_{i+1}\le y_i<x_i}\prod_{i<j}(y_i-y_j)
 =\frac{1}{(k-1)!}\prod_{i<j}(x_i-x_j). \tag{CB.2}
\]
For completeness, write the left Vandermonde as a determinant and sum each row.
Choose $Q_j$ with $Q_j(t+1)-Q_j(t)=t^{k-1-j}$, of leading coefficient $1/(k-j)$.
The resulting determinant is $\det[Q_j(x_i)-Q_j(x_{i+1})]$.
Append the constant column and undo successive row differences: the determinant has degrees $k-1,\ldots,0$ and leading-coefficient product $1/(k-1)!$, proving (CB.2).

\paragraph{CB4--CB5. The actual complement.}
For a selected weight $\lambda^0$, put
\[
 \ell=(\lambda^0_2,\lambda^0_2,\lambda^0_3,\ldots,\lambda^0_k).
\]
Use the effective weight $\ell$ at even positions and $-\operatorname{rev}\ell$ at odd positions, and subtract its last entry to obtain $\nu$.
At an even position impose $e_{1,1}=\nu_1$: deleting the first row and column gives shape $(\nu_2,\ldots,\nu_k)$.
At an odd position impose $e_{1,t}=0$: since $\nu_k=0$, delete the first row and last column.
The inverse operations insert the indicated zeros and, in the even case, $\nu_1$.
Constant shifts and negative reversal preserve $D$, so both counts are $D_{k-1}(\lambda^0_2,\ldots,\lambda^0_k)$.
Consequently
\[
 n_2=\prod_{s\in T}D_{k-1}(\lambda^0_{s,2},\ldots,\lambda^0_{s,k})
       \prod_{s\notin T}D_k(\lambda^0_s). \tag{CB.3}
\]
The tensor basis also gives $\dim V=\prod_sD_k(\lambda^0_s)$, the source's \texttt{FewBadOrbits} interface.
If $d_s=\lambda^0_{s,1}-\lambda^0_{s,2}\ge a>0$ and
$0\le\lambda^0_{s,2}-\lambda^0_{s,j}\le c_0$, then
\[
\begin{aligned}
 \frac{\binom{d_s+k-1}{k-1}D_{k-1}(\lambda^0_{s,2},\ldots,\lambda^0_{s,k})}{D_k(\lambda^0_s)}
 &=\prod_{j=2}^k\frac{d_s+j-1}{d_s+\lambda^0_{s,2}-\lambda^0_{s,j}+j-1}\\
 &\ge1-\frac{(k-1)c_0}{a}.
\end{aligned}
\]
Multiplication yields the source's weaker factor $(1-kc_0/a)^{|T|}$ when $kc_0/a\le1$.

\paragraph{CB6--CB7. Integral multiplication and the source lattices.}\cite{complement-cartan}
A product of minors of shape $\nu+\mu$ has
$(\nu_j-\nu_{j+1})+(\mu_j-\mu_{j+1})$ factors of height $j$.
Partition them accordingly. CB2 therefore proves
\[
 \nabla_{\nu+\mu}(R)=\operatorname{span}_R
 \{fg:f\in\nabla_\nu(R),\ g\in\nabla_\mu(R)\}. \tag{CB.4}
\]
Over a DVR primitive factors have primitive product: reduce their coefficient polynomials modulo the uniformizer.

For the actual basis matrix $M$, row substitution $\mathcal R(M)f(X)=f(MX)$ is an anti-action. The source lattice is exactly
\[
 L_\lambda(M)=\det(M)^{-\lambda_k}\mathcal R(M^{-1})\nabla_{\lambda-\lambda_k}(\mathcal O).
\]
Changing $M$ to $MU$, $U\in\mathrm{GL}_k(\mathcal O)$, preserves this lattice; hence it agrees with the source's intersection over all bases.
For the midpoint matrices $M_s$ of EL, the actual edge-coordinate isomorphism is
\[
 \mathcal C_\lambda(F)=w^{A_r(\lambda)}
 \left(\prod_s\det(M_s)^{\lambda^{\rm eff}_{s,k}}\right)
 \operatorname{blockAct}(M_s)F.
\]
It takes $R_e(\lambda)$ onto the integral tensor coefficient lattice.
All exponents add with weights, so
$\mathcal C_{\lambda+\mu}(FG)=\mathcal C_\lambda(F)\mathcal C_\mu(G)$.
Thus (CB.4) is the integral edge statement \texttt{OntoAll}, and applying
$\mathcal C_{\rm rest}^{-1}$ to the selected basis products constructs the exact complementary family.
No denominator or characteristic restriction is introduced.
This supplies the basis, count and Cartan inputs retained in CL.
The distinct incoming argument is IN/WT. CG/LA/GP/RS/RK supply selected-symbol density; arithmetic windows remain separate.

\section{The complementary weight family and primitive dual tests}
\label{sec:cloud}

\status{This section expands the source's Type-II \texttt{CloudStep}: we prove the complementary dimension comparison, the reduction modulo the uniformizer for every primitive combination, the contraction argument, and the valuation bookkeeping. The normalized Weyl comparison is proved in Section~\ref{sec:integral-weyl}; Sections~\ref{sec:local-gram}--\ref{sec:digit-tests} prove the constrained Gram and admissible-test calculations. Section~\ref{sec:complementary-basis} proves the complementary standard-monomial/count and Cartan identifications; Sections~\ref{sec:initial-membership}--\ref{sec:word-transport} prove incoming membership. The selected-symbol rank is proved in RK; FG/OR/CP/EA below supply the complete local frame, outgoing realization and numerical compatibility. The expanded CL note proves the characteristic-zero operator bound used in contraction.}

\subsection{A complementary family of the right size}

Use the edge lattices and perfect pairings of Section~\ref{sec:edge-lattices}. Write $k=m+1$, $q=p^{n'}$, $N=|\mathcal T|$, and
\[
 d_s=\begin{cases}\lambda^0_{s,1}-\lambda^0_{s,2},&s\in\mathcal T,\\0,&s\notin\mathcal T,\end{cases}
 \qquad \lambda^{\rm rem}_s=\lambda^0_s-(d_s,0,\ldots,0).
\]
Thus the first two remainder weights agree at every selected twist. Let $D_j(\mu)$ denote the Weyl dimension
$\prod_{i<l}(\mu_i-\mu_l+l-i)/(l-i)$ for a decreasing $j$-tuple. The complementary family has the exact cardinality
\[
 n_2=\prod_{s\in\mathcal T}D_{k-1}(\lambda^0_{s,2},\ldots,\lambda^0_{s,k})
       \prod_{s\notin\mathcal T}D_k(\lambda^0_s).
 \tag{CL.1}\label{eq:cl-count}
\]
Here is the construction behind this formula. Use the integral monic standard-monomial basis of the remainder representation, indexed by interlaced exponent arrays $e(a,t)$, with column sums $\nu_t$; arrays use indices $0,\ldots,k-1$. At a selected even twist retain those with $e(0,0)=\nu_0$. At a selected odd twist retain those with $e(0,t)=0$ for every $t$. At unselected twists retain the entire basis. In the even case deleting the first row and column identifies the selected arrays with the $(k-1)$-variable arrays of tail shape. In the odd case delete the first row and last column; the last column has sum zero after the prescribed dual reversal. The latter shape is the dual of the tail, up to a determinant twist, so it has the same dimension. Tensor these families and transport them to $R_e(\lambda^{\rm rem})$. This gives $h_1,\ldots,h_{n_2}$. \textsf{[CL1]}\cite{cloud-construction,cloud-flat}

Suppose $d_s\ge a>0$ and the spread of the selected tail is at most $c_0$, with $kc_0/a\le1$. Weyl's formula, after subtracting the last weight, gives
\[
 D_k(d+\mu_0,\mu)
 =D_{k-1}(\mu)\prod_{j=1}^{k-1}\frac{d+\mu_0-\mu_{j-1}+j}{j}
 \le D_{k-1}(\mu)(1+c_0/a)^{k-1}\binom{d+k-1}{k-1}.
\]
For $0\le x\le1/k$, $(1-kx)(1+x)^{k-1}\le1$: its logarithmic derivative is negative on $x<1/k$, and the endpoint is immediate. Multiplying the resulting inequalities proves
\[
 (1-kc_0/a)^N\dim V
 \le n_2\prod_{s\in\mathcal T}\binom{d_s+k-1}{k-1}.
 \tag{CL.2}\label{eq:cl-weyl}
\]
Consequently, if the symbol construction supplies a saturated free family $B_1,\ldots,B_b$ of size
$b\ge(1-C_rNr^2k/Y)q^{-r(k-r)}\prod_s\binom{d_s+k-1}{k-1}$, then
\[
 bn_2\ge\frac{\rho\dim V}{q^{r(k-r)}},\qquad
 \rho=(1-C_rNr^2k/Y)(1-kc_0/a)^N.
 \tag{CL.3}
\]
If the first factor is negative this last bound is automatic; otherwise multiply (CL.2) by that factor. The symbol dimension estimate itself is a separate input. \textsf{[CL2]}

\subsection{Why every primitive combination survives}

Let $P$ be the integral polynomial obtained from the source's determinant factor, with $\bar P\ne0$ modulo $w$. Suppose the flat lifts satisfy $\operatorname{subst}(F_b)=B_bP$ in the appropriate renamed variables, and put $S_{b,i}=\Theta(F_b)h_i$. Then
\[
 \exists(b,i):c_{b,i}\in\mathcal O^\times
 \quad\Longrightarrow\quad
 U:=\sum_{b,i}c_{b,i}S_{b,i}\notin wR_e(\lambda^0).
 \tag{CL.4}\label{eq:cl-primitive}
\]
To prove this, first observe why the polynomial products used above are independent over the residue field. At an even twist the flat generators are $X_{c,0}$; the selected remainder's first column is fixed, so the first column of a product's leading exponent recovers every flat exponent $\alpha_c$. At an odd twist they are the signed maximal minors omitting row $c$. The diagonal entries of their leading exponent are $\sum_{c>t}\alpha_c$, $0\le t<k-1$. The selected remainder has zero diagonal: its zero first row and interlacing force $e(a,t)=0$ for $a\le t$. These tail sums, together with $\sum_c\alpha_c=d_s$, recover $\alpha$. The remainder exponent is then recovered by subtraction. Thus all leading monomials are distinct; their leading coefficients are units. Tensoring over the twists preserves independence.

The coordinate change defining $R_e$ sends its elements to integral remainder polynomials, and sends the displayed products to these flat-generator products. Its scalar factors multiply because both the determinant weight and $A$ are additive in the weight. If $U\in wR_e$, reduction modulo $w$ therefore gives
\[
 \sum_i\Theta_0\!\left(\bar P\sum_b\bar c_{b,i}\bar B_b\right)\bar h_i=0.
\]
The independence just proved makes each polynomial in parentheses zero. Cancel $\bar P$ in the residue polynomial ring. Saturation of $\operatorname{span}_{\mathcal O}(B_b)$ and linear independence of the $B_b$ imply $c_{b,i}\in w\mathcal O$ for every pair, a contradiction. This is a test of arbitrary primitive combinations, not just of individual family members. \textsf{[CL3]}

\subsection{Contraction reduces the dual test to a flat representation}

First work on the selected twists alone. Section~\ref{sec:integral-weyl} supplies the actual normalized maps $\Psi_0,\Psi_f$ from dual to opposite-edge representations satisfying
\[
 p^{\kappa_i}\bar R_i\subseteq\Psi_i(R'_i)\subseteq\bar R_i,
 \qquad i=0,f,
 \tag{CL.5}\label{eq:cl-sandwich}
\]
and their analogous vertex-lattice comparisons. Cartan multiplication spans $\bar R_0$ by $y\beta$, with $y\in\bar R_f$, $\beta\in\bar R_{\rm rem}$. In the standard-monomial model this spanning follows by splitting each product of flag minors into the two prescribed shapes. Set $\kappa=\kappa_0+\kappa_f$ and
$\Lambda_\beta=p^{\kappa_0}\Psi_0^{-1}M_\beta\Psi_f$, where $M_\beta(y)=y\beta$. Define $\rho_\beta$ by
\[
 \beta_f(\rho_\beta u,T)=\beta_0(u,\Lambda_\beta T).
 \tag{CL.6}
\]
Then $\rho_\beta R_0\subseteq R_f$. If all $\rho_\beta u$ lie in $w^{2\kappa+1}R_f$, take any $Z\in R'_0$, express $\Psi_0Z=\sum y\beta$, and use $T=\Psi_f^{-1}(p^{\kappa_f}y)\in R'_f$. Equation (CL.6) makes $p^\kappa\beta_0(u,\Psi_0^{-1}(y\beta))$ divisible by $w^{2\kappa+1}$. Since $p^\kappa=w^{2\kappa}$, summing gives $\beta_0(u,Z)\in w\mathcal O$. Perfectness implies $u\in wR_0$. \textsf{[CL5]}

Contracting the nonselected tensor factors against an integral dual basis preserves this detection statement: if every contraction were divisible by $w$, the original tensor would be divisible by $w$. For the primitive $U$ in (CL.4), therefore choose a contraction and write
$\rho_\beta U=w^\nu\Theta(F_1)$, with $0\le\nu\le2\kappa$ and $\Theta(F_1)$ primitive in $R_f$. The Fock contraction identity places $F_1$ in the span of words in $X_{s,i}\partial_{s,j}$ applied to the $F_b$, using at most $c_0$ letters at each selected twist. A characteristic-zero proof of this assertion, including its highest-weight argument, is given in the expanded note. It makes no assertion that these operator coefficients are integral. \textsf{[CL4]}

Put $\ell=rNc_0+2rn'c''$. The same lattice comparisons at the other edges and at the target give a dual map $\Lambda$ with loss at most $\ell$ in $w$-exponents and a scalar $c\in\mathcal O$ such that
\[
 c\,\beta_0(U,\Lambda T)=p^\kappa\beta_f(\Theta F_1,T).
 \tag{CL.7}\label{eq:cl-transfer}
\]
Explicitly, $\Lambda=w^{2(A_f-A_0)}\rho_\beta^*$ and $c=w^{2\kappa-\nu}$; hence $c$ need not be a unit. The source's nonselected tensor factors can be contracted by an integral dual basis without adding to $\kappa$; they contribute the $2rn'c''$ height loss. For the selected factors,
\[
 \kappa_i=\sum_{s\in\mathcal T,\,h_{i,s}>0}
 \binom{k}{2}\lfloor\log_p(h_{i,s}+k-1)\rfloor,
 \qquad h_{i,s}+k\le a^2,
 \quad\kappa\le4Nk^2\log_p a.
 \tag{CL.8}
\]
Here $h_{i,s}$ is the spread of the corresponding top or flat weight. These are actual lattice comparison exponents, not dimensions of the representations. \textsf{[CL5--CL6]}

\subsection{The flat valuation bound and what remains}

Write $L=Nk^2\log_p a\ge1$ and normalize $v(p)=1$, so $v(w)=1/2$. For a flat polynomial of multidegree $d$, put
\[
 \mathfrak f=\prod_s d_s!,\qquad
 \kappa_{\rm fl}=\sum_s k\lfloor\log_p d_s\rfloor,
 \quad \lfloor\log_p0\rfloor:=0.
\]
The normalized Fock pairing $p^{\kappa_{\rm fl}}\langle G,H\rangle/\mathfrak f$ is integral for integral $G,H$: Legendre's formula bounds the valuation of each relevant multinomial coefficient by $\kappa_{\rm fl}$. The constrained construction in Sections~\ref{sec:local-gram}--\ref{sec:digit-tests}, with incoming \texttt{MemAt} now proved in IN/WT and the stated frame/size data, gives an admissible test polynomial $H$ with
\[
 \langle G,H\rangle=M+E,\qquad
 v(M)\le v(\mathfrak f)+29L,\qquad
 v(E)\ge c_P-20L+v(\mathfrak f)-\kappa_{\rm fl},
 \tag{CL.9}\label{eq:cl-flat-error}
\]
where $\kappa_{\rm fl}\le20L$ and $c_P>69L$. The strict inequality forces $v(E)>v(M)$, so
\[
 y=p^{\kappa_{\rm fl}}\langle G,H\rangle/\mathfrak f\in\mathcal O,
 \qquad y\mid p^{D_X},\qquad D_X=\lceil49L\rceil\le50L.
 \tag{CL.10}
\]
The expanded CL note derives (CL.9) from the selection and Gram estimates. Those pointwise constructions are now proved in PG and DT. The Gram remainder itself can be zero, although the digit-normal approximation still contributes error; the conservative constant is unchanged. \textsf{[CL7--CL9]}

The localizer realization writes $H$ as a sum of allowed symbols and realizes its pairing as $y+w^{c_P}z$. If $yb=p^{D_X}$ and $2D_X<c_P$, then
\[
 (y+w^{c_P}z)b=p^{D_X}\bigl(1+w^{c_P-2D_X}zb\bigr);
\]
the parenthesized factor is a unit, so the realized flat value still divides $p^{D_X}$. Applying (CL.7), an integral full value $x=\beta_0(U,\Lambda T)$ satisfies $cx=p^\kappa y'$, where $y'b'=p^{D_X}$. Thus $x(cb')=p^{\kappa+D_X}$. We obtain, for \emph{every} primitive coefficient array,
\[
 D_1=\kappa+D_X\le54Nk^2\log_p a,
 \qquad t_T\le t_0+Nc_0+\ell,
 \qquad \Lambda T\in w^{c_*}R'_{e'}\quad(e'\ne e),
 \tag{CL.11}\label{eq:cl-final}
\]
using $c_P=c_*+\ell$. The source hypotheses include $D_1\le D_p$ and $c_*\ge8(D_p+1)$ once its large constant exceeds the displayed bounds. They ensure the needed strict perturbation inequalities. This extracts the usable constant $54$ from the pinned witness choices. IW, PG and DT now discharge the integral Weyl and pointwise constrained-test inputs. CB and IN/WT now discharge the complementary basis/Cartan and incoming \texttt{MemAt} inputs. RK supplies quantitative selected-symbol rank. FG constructs the actual frames, OR supplies the additional outgoing realization, CP matches all pointwise budgets, and EA proves the full Type-II conclusion below. \textsf{[CL10--CL12]}

\section{Initial symbols satisfy the incoming filtration}
\label{sec:initial-membership}

\paragraph{IN1--IN2: the exact statement.}\cite{initial-blocks}
Assume the archive's \texttt{DataII} hypotheses with $C\geq1$, and put
$\mathbb K=\operatorname{Frac}(O)$, $k=r+K$.  For every actual symbol
$S=\mathrm{SPR}\,B_b$, every physical digit block, every selected normal
joint projector $e_\kappa$, and every normal coefficient, we prove
\[
 H=[Z^\nu](1\otimes e_\kappa)\operatorname{lev}_{\rho,d_P,d_N}(S)
 \quad\Longrightarrow\quad H\in X_D(0,1/8).
 \tag{IN1}
\]
This is precisely \texttt{Lemma5aAt}, the initial half of incoming
\texttt{MemAt}.  Here $v_T$ denotes the minimum exponent of a
generic-conjugated rank-one torus acting on the signed letters in $T$;
$v_T(0)=+\infty$.  The required conditions are plane multidegree $d_P$,
chain bounds $\theta_{j,i}$ after every frame-digit extraction, and link
bounds $\theta_{j,u}-s_{j,u}/8$.

Expand each segment factor into its defining plane--normal generators
$A_jG_j$, and expand $\mathrm{SPR}$ into monomials $X^\gamma$.
Each $\gamma$ uses only plane coordinates at even, hence vector, letters:
every determinant defining $\mathrm{SPR}$ has this property.
Retain the actual total degree
$d^{\rm tot}=\mathrm{G2d}+\deg\gamma$ through the compatibility condition
\[
 \sum_c\rho(x,c)+q_x\bigl(d_P(x)+d_N(x)\bigr)
       =d^{\rm tot}(\operatorname{pos}x),
 \qquad q_x=p^{z_j-\operatorname{pos}x}.
 \tag{IN2}
\]
Individual expanded summands need not have that degree.  We first prove
membership only on blocks satisfying (IN2), extend by linearity, and
then use the actual symbol's homogeneity: all incompatible blocks vanish.

\paragraph{IN3--IN4: the normal projector supplies the threshold.}
For one product summand, put $b=\lfloor(l+\gamma)_{\rm plane}/q\rfloor$.
Its digit block is zero unless $\rho\equiv l+\gamma\pmod q$; otherwise it is
$A_*\otimes G_*$.  The normal factor is the prescribed homogeneous
component of the product of renamed $G_j$'s.  Normal labels disappear
because their exponents are $<q$, and $G_j$ has no preceding-frame variables.
Consequently
\[
 [Z^\nu](1\otimes e_\kappa)(A_*\otimes G_*)
             =[Z^\nu](e_\kappa G_*)\,A_*.
 \tag{IN3}
\]
If $e_\kappa G_*\neq0$, every selected prefix type is present.
Indeed canonical isotypic projectors commute for nested or disjoint
letter sets, so their product is the unique joint projector.
The generic maximum weight of a nonzero highest-weight type $\lambda$
is $\lambda_1$; homogeneous extraction cannot increase it.
Thus the original normal-prefix upper bounds also bound the corresponding
$\kappa_{j,u,1}$.

Write $x_{j,u}$ for the ordered nonpreceding-frame letters,
$B_{j,u}=\mathrm{Bo}(j,x_{j,u})$, and
$s_{j,u}=\lfloor B_{j,u}/(Yk^2)\rfloor$.  The defining pair inequalities give
\[
 v_{<x_{j,u}}(A_j)\geq\theta_{j,u},\qquad
 \theta_{j,u}=\max\left\{
 \begin{array}{l}
 B_{j,u}/2,\\
 \kappa_{j,u+1,1}+m_*+s_{j,u}\quad\text{if }x_{j,u}\text{ is even},\\
 \kappa_{j,u-1,1}+m_*+s_{j,u}\quad\text{if }u>0,
                         \ x_{j,u-1}\text{ is odd}
 \end{array}\right\}.
 \tag{IN4}
\]
Absent terms are omitted.  For the last line, the supported sets
$\{x\leq x_{j,u-1}\}$ and $\{x<x_{j,u}\}$ coincide, and an odd
predecessor leaves the oldest preceding vector, hence $B$ and $s$, unchanged.

\paragraph{IN5--IN6: frame digits cost no chain order.}
Give each segment independent intermediate letters and current-frame
labels, and give its preceding frame independent variables $w_{j-1,m}$.
Let $\chi_j$ rename $A_j$ into these variables and
$\Pi=\prod_j\chi_jA_j$.  The substitution $\phi$ fixes intermediate letters
and sends
\[
 \mathrm{lab}_{j,m,a}\mapsto X_{j,m,a},\qquad
 w_{j,m,a}\mapsto X_{j,m,a}^{P_j},\qquad
 P_j=p^{z_{j+1}-z_j}.
\]
Exactly $A_*=[d_P](X^b\phi\Pi)=X^b\phi\Pi'$, where $\Pi'$ keeps
the model multidegrees mapping to $d_P$.  Such filters commute with
the generic torus.  On the model chain consisting of $w_{j-1,m}$ and
the first $i$ intermediate letters, only $\chi_jA_j$ contributes; therefore
$v_{C_{j,i}}(\Pi')\geq\theta_{j,i}$.

Expand $\Pi'=\sum_M\mathrm{lab}^M Q_M$.  Labels lie outside every chain,
so each $Q_M$ retains its chain bound and uses only intermediate and $w$
variables.  Direct division of monomial exponents gives
\[
 \operatorname{Fdig}_{\sigma}(X^\beta\phi Q_M)=
 \begin{cases}
 X^{\operatorname{lift}\lfloor\beta/P\rfloor}Q_M,
                   &\beta\bmod P=\sigma,\\
 0,&\text{otherwise}.
 \end{cases}
 \tag{IN5}
\]
Here $P=1$ on intermediate letters.  With $\beta=b+\phi\!\operatorname{exp}(M)$,
the extra monomial has no covector variables inside the chain: labels
are frame vectors, and $b$ has no odd intermediate exponents.
Its order is nonnegative.  Summation proves the exact chain bound.

\paragraph{IN7--IN8: the sole loss is a preceding-frame degree.}
For a link ending at $x_{j,u}$, where $u=n_j+r-1-m$ and $m+2\leq r$,
remove the preceding frame $W_j$ from the original prefix.  If a
homogeneous piece has total $W_j$-degree $d_W$, this lowers order by
at most $d_W$: apply the inverse vector twist on $W_j$, whose Laurent
exponents are at least $-d_W$, then extract its coefficients.
Remaining letters rename into the link; all other factors contribute
only vector variables there.  A piece surviving $[d_P]$ satisfies
$P_{j-1}\deg_{W_{j,m'}}\leq d_P(j-1,m')-\deg_{j-1,m'}b$.  Hence
\[
 v_{T_{j,m}}(A_*)\geq\theta_{j,u}-\Delta_j,\qquad
 \Delta_j=\sum_{m'=0}^{r-1}
 \frac{\bigl(d_P(j-1,m')-\deg_{j-1,m'}b\bigr)_+}{p^{g_j}},
 \quad g_j=z_j-z_{j-1}.
 \tag{IN6}
\]
The fixed-degree relation (IN2) gives, for distinct integers
$t_{m'}=z_{j-1}-\operatorname{pos}(j-1,m')\geq0$,
$p^{t_{m'}}(d_P-\deg b)_+\leq\mathrm{G2d}\leq2a$.
Thus $\Delta_j\leq3a/p^{g_j}$, since $p\geq3$.
If $D=z_j-\operatorname{oldest}(j,x_{j,u})$, frame separation gives
$D+L_1+1\leq g_j$.  The source size hypotheses with $C\geq1$ give
\[
 a\leq B_{j,u}p^D,\qquad 48Yk^2\leq p^{L_1+1},\qquad
 1\leq B_{j,u}/(Yk^2)\leq2s_{j,u}.
\]
Consequently
\[
 \Delta_j\leq\frac{3a}{p^{g_j}}
 \leq\frac{3B_{j,u}}{p^{L_1+1}}
 \leq\frac{B_{j,u}}{16Yk^2}\leq\frac{s_{j,u}}8.
 \tag{IN7}
\]

\paragraph{IN9: assembly.}
Equations (IN3)--(IN7), plane homogeneity, and linearity prove (IN1).
This closes initial \texttt{MemAtG(0,1/8)} for the actual symbol family.
The following word-transport argument spends the separate operator slack.

\section{Transport through the allowed operator words}
\label{sec:word-transport}

\status{This section proves the operator-word half of incoming \texttt{MemAt}. Its input is the actual initial-symbol assertion \texttt{MemAtG}$(0,1/8)$, proved in IN9. Exact digit identities, a normal-type calculation and a numerical prefix budget prove \texttt{WordConditionAt}; linearity then gives the assertion for the whole word span. The existence of the frame and initial symbols is supplied separately.}

\subsection{The prefix budget}

Put $k=r+K$. A valid periodic frame has increasing even positions
$f_{j,0}<\cdots<f_{j,r-1}=z_j$, with
$z_{j-1}<f_{j,0}$ and $f_{j+J,i}=f_{j,i}+n'$.
Let $\mathcal W_j$ be the selected positions between $f_{j-1,0}$ and $z_j$.
The selected positions in $(z_{j-1},z_j]$, ordered as $x_{j,u}$,
are the internal letters followed by the current frame. Write $N_{j,u}$
for the first $u$ of these letters. These true letters represent the selected
residues bijectively. Define
\[
 \delta_{j,u}(s)=
 \mathbf1_{\{\exists x\in\mathcal W_j:\ x\equiv s\pmod {n'},\ x\le x_{j,u}\}}.
 \tag{WT.1}
\]
The weak inequality includes every letter affecting either adjacent
normal prefix or the plane cutoff; it also includes the previous frame
whose carry variables enter the next segment's chains. \textsf{[WT1]}\cite{incoming-words}

For normal joint highest weights $\kappa$, let $T_{j,u}(\kappa)$ denote
the source's \texttt{thStar}. It is the maximum of $B_{j,u}/2$ and the
active terms
$\kappa_{j,u+1,1}+m_*+s_{j,u}$ when $x_{j,u}$ is even, and
$\kappa_{j,u-1,1}+m_*+s_{j,u}$ when $u\ge1$ and $x_{j,u-1}$ is odd.
Here $s_{j,u}=\lfloor B_{j,u}/(Yk^2)\rfloor$.

Write $\mathcal M_{\alpha,\beta}(L;F)$ for the following precise universal
condition: after every actual digit/block extraction and normal joint
projection, every plane coefficient has chain orders at least
$T-\alpha s-L$ and link orders at least $T-\beta s-L$.
Positive link thresholds below these bounds are allowed; this equivalent
form permits changing one block degree during differentiation.
The input is $\mathcal M_{0,1/8}(0;F)$. \textsf{[WT1,WT6]}

\subsection{Digit extraction and the signed twist}

At one physical coordinate let $q=p^{z_j-x}$ and let $\mathcal L_\rho$
retain exponents $\rho+qb$ and replace them by $b$ in a coarse variable $y$.
For $0\le\rho<q$, direct differentiation of monomials gives
\[
\begin{aligned}
 \mathcal L_\rho(\partial_X F)
 &=\begin{cases}
 ((\rho+1)+qy\partial_y)\mathcal L_{\rho+1}F,&\rho<q-1,\\
 q\partial_y\mathcal L_0F,&\rho=q-1,
 \end{cases}\\
 \mathcal L_\rho(XF)
 &=\begin{cases}
 \mathcal L_{\rho-1}F,&\rho>0,\\
 y\mathcal L_{q-1}F,&\rho=0.
 \end{cases}
\end{aligned}
\tag{WT.2}
\]
Other residues are unchanged. After block projection, the derivative
carry term uses input degree $d+1$, and the multiplication carry term
uses $d-1$; the latter vanishes if $d=0$.
The same identities apply to the second, frame-to-carry digit extraction.
Thus the Frobenius exponents introduce scalar factors $q$, with no
amplification of the order loss. \textsf{[WT3]}

The signed-twist matrix is
$A(t)=g\operatorname{diag}(1,\ldots,t,\ldots,1)g^{-1}$.
Its entries have exponents $0,1$, and those of $A(t)^{-1}$ have exponents
$0,-1$. A vector coordinate transforms by $A(t)$ and a covector by the
inverse; the derivative transforms by the opposite matrix. Consequently
\[
\begin{array}{c|cc|cc}
 &\multicolumn{2}{c|}{\text{minimum-order loss}}
 &\multicolumn{2}{c}{\text{maximum-order increase}}\\
 &\text{vector}&\text{covector}&\text{vector}&\text{covector}\\ \hline
 y\cdot&0&1&1&0\\
 \partial_y&1&0&0&1\\
 y\partial_y&1&1&1&1
\end{array}
\tag{WT.3}
\]
Every entry is zero when the letter is outside the prefix. In particular,
Euler costs one, since one of its two factors has zero cost.
Together with (WT.2), this proves the plane derivative, multiplication
and Euler steps, including every second-digit chain. \textsf{[WT2,WT4]}

\subsection{The normal projection and one full operator}

Normal prefix projectors commute: their letter sets are nested within a
segment and disjoint across segments. On each finite multidegree space
they give a finite joint decomposition $\sum_\kappa e_\kappa=1$.
The maximum generic-twist exponent on type $\kappa_{j,u}$ is
$\kappa_{j,u,1}$. Hence for
$D\in\{\partial_y,y\cdot,y\partial_y\}$,
\[
 e_{\kappa'}D e_\kappa\ne0
 \quad\Longrightarrow\quad
 \kappa'_{j,u,1}\le
 \kappa_{j,u,1}+\mathbf1_{\{x\in N_{j,u}\}}
 \quad\text{at every cut}.                         \tag{WT.4}
\]
Indeed (WT.3) bounds the maximum before projecting, and an excluded type
has zero projection. Only these finitely many compatible input types
contribute to an output coefficient. \textsf{[WT5]}

Both adjacent indicators in $T_{j,u}$ are bounded by
$\delta_{j,u}(s)$. Taking maxima therefore gives
$T_{j,u}(\kappa')\le T_{j,u}(\kappa)+\delta_{j,u}(s)$.
For a plane operation, raise the input threshold by this same indicator
and use (WT.3). For a normal operation, use (WT.4) and the displayed
comparison. Formula (WT.2) and linearity cover all terms. Thus
\[
 \mathcal M_{0,1/8}(L;F)
 \Longrightarrow
 \mathcal M_{0,1/8}
   (L+2\delta(s);X_{s,a}\partial_{s,b}F).
 \tag{WT.5}
\]
The universal block quantifiers justify every intermediate degree change;
the complete operator preserves physical multidegree. \textsf{[WT6]}

\subsection{Absorbing the entire word}

For an allowed word $\ell$, at most $c_0$ operators occur at each selected
residue. Its accumulated loss satisfies
\[
 L_\ell(j,u)=2\sum_{(s,a,b)\in\ell}\delta_{j,u}(s)
 \le2\#\mathcal W_jc_0\le2N_1c_0.
\]
The actual frame bound gives $B_{j,u}\ge a/p^{G'}$.
The source's numerical hypothesis, for $C\ge16$, implies
$a\ge16N_1c_0Yk^2p^{G'}$; hence the integer floor satisfies
$s_{j,u}\ge16N_1c_0$.
Therefore $L_\ell\le s/8$. Induction in (WT.5), followed by weakening
the cutoffs, proves
\[
 \operatorname{MemAtG}(0,1/8;F)
 \Longrightarrow
 \operatorname{MemAtG}(1/8,3/8;\operatorname{opWord}_\ell F).
 \tag{WT.6}
\]
This is \texttt{WordConditionAt}. Each defining condition is linear in
$F$, so it extends from the initial symbols to their span and then to
the literal \texttt{OpCl} word span. Together with IN9 this proves
incoming \texttt{MemAt} for the actual family. \textsf{[WT7--WT9]}

\section{The constrained local Gram test}
\label{sec:local-gram}

\status{This section reconstructs the projection--derivative argument behind CL9. It uses the explicitly identified signed-path basis, integral-projector and factorial-unit lemmas of the companion IP, IW and DT notes. The construction gives the source constants $(26,2,24,60)$ and, more strongly, permits the remainder $R=0$. The expanded note proves the intermediate assertions and records their exact source dependencies.}

\subsection{The conditions and the two orders}

Work over the source DVR with fraction field $K$, prime $p\ge5$, and $v(p)=1$; arbitrary ramification is allowed. Each letter has $r$ coordinates and is a vector or covector. Write
\[
 \langle F,H\rangle=\sum_\alpha F_\alpha H_\alpha\alpha!,\quad
 b(F)=\min_{F_\alpha\ne0}\left(v(F_\alpha)+\tfrac12v_p(\alpha!)\right),\quad
 n_0=\tfrac12v_p\!\left(\prod_xd_x!\right).
\]
For a set $S$ of letters, $\operatorname{ord}_S F$ is the least exponent of $t$ after applying a generic conjugate of $\operatorname{diag}(1,\ldots,1,t)$ to those letters, with the dual action on covectors. It is distinct from the coefficient valuation $v$.

Segments $j\in\mathbb Z/J\mathbb Z$ contain signed base letters $L_j$ and vector letters $z_{j,m}$, $0\le m<r$. Put $P_{j,m}=p^{g_{j,m}}$. The enlarged polynomial ring replaces each $z$ by a label $a$, carry $u$, and tangent $\tau$, with
\[
 \phi(a_{j,m,c})=z_{j,m,c},\qquad
 \phi(u_{j,m,c})=z_{j,m,c}^{P_{j,m}},\qquad
 \phi(\tau_{j,m,c})=z_{j,m,c};
\]
tangents with $m=0$ are excluded. A \emph{good} monomial has label exponents $<P_{j,m}$, the prescribed degrees after $\phi$, and tangent degree at $(j,m)$ at most $d_{L,j}:=\sum_{x\in L_j}d_x$.

The chain sets $S_{j,i}$ consist of a nested subset of $L_j$ and all carries from $j-1$. The link sets are
$Y_{j,m}=L_j\cup\{u_{j-1,*}\}\cup\{\tau_{j,m'}:m'>m\}$.
The input $F$ has $b(F)=0$, prescribed multidegree, chain orders $\theta_{j,i}$ in every digit component, and link orders $\theta^L_{j,m}$ before splitting. Let $\Gamma$ multiply coefficient $\alpha$ by the prescribed unit $\gamma(\alpha)$, whose two signs admit integral binomial-polynomial approximations modulo $p^h$ of degree $\le6h$.

Set
\[
 \kappa=\sum_x r\bigl(\log_p(d_x+r)+1\bigr),\quad
 e=r^2\bigl(\log_p(2\sum_xd_x+r)+1\bigr),\quad
 D_m=e\sum_j n_j,\quad D_l=eJ(r-1).
\]
Here $\kappa$ is the source's \texttt{kp}. Assume $c\ge1$,
\[
 c>4(D_m+D_l+3\kappa),\qquad
 \frac{r d_{L,j}}{P_{j,m}}+26rc\le s_{j+1,i}.
 \tag{PG.1}
\]
Write $h=\lceil c\rceil\le2c$. Thus the last inequality holds with $13rh$ in place of $26rc$. \textsf{[PG1]}\cite{gram-construction}

\subsection{Constructing a detector and imposing the chains}

The weighted lattice $\{b\ge0\}$ is stable under divided powers: a root operator of order $q$ changes the weighted coefficient valuation by
\[
 \tfrac12v_p\binom{\alpha_b}{q}
 +\tfrac12v_p\binom{\alpha_a+q}{q}\ge0.
\]
The integral-projector lemma therefore gives each order-cutoff projector a loss at most $e$ in $b$. Projectors for nested or disjoint letter sets commute and are Fock self-adjoint. Choose $\alpha$ attaining $b(F)=0$, and project $x^\alpha$ onto all original link conditions. The pairing with $F$ is unchanged. Approximate $\Gamma^{-1}$ modulo $p^h$, then clear its coefficient denominators by $p^{\lceil n_0+D_l-v_p(\alpha!)/2\rceil}$. This gives an integral detector $H$ with link loss $6h$ and
\[
 v\langle F,\Gamma H\rangle\le n_0+D_l+1.
 \tag{PG.2}
\]
The signed-path basis lifts $H$ to an integral good polynomial $Z$ with $\phi Z=H$ and the same link bounds. Explicitly, in each path-column product move the frame letters $m>0$ to tangents exactly when that column contains a base letter; each tangent then occurs at most $d_{L,j}$ times. Split the remaining labels into residue and carry exponents. The column order bounds prove that this preserves the required links. \textsf{[PG3--PG5]}

Let $\Pi$ be the product of the chain projectors at levels $\theta-s$. It preserves goodness and the link conditions. On a good monomial its coefficient denominators cost at most $D_m+\kappa/2$. Factorial splitting supplies a unit function
$\gamma_*(\zeta)=\gamma(A)\gamma_{\rm dig}(\zeta)$, where $\phi(x^\zeta)=x^{A+M}$ and $M$ is the tangent exponent. Choose integral diagonal approximations $g_\pm\equiv\gamma_*^{\pm1}\pmod{p^h}$, each of order loss at most $12h$, and set
\[
 \Omega=\operatorname{diag}(g_-)\Pi\operatorname{diag}(g_+),
 \qquad G=p^{N_2}\Omega Z,\qquad N_2=\lceil D_m+\kappa/2\rceil.
 \tag{PG.3}
\]
Then $G$ is integral and good, with chain loss $12h$ and total link loss $30h$. \textsf{[PG6,PG8]}

\subsection{Why the constrained projection preserves detection}

The decisive derivative estimate is elementary. For $P=p^g$,
\[
 a_l=[z^M](1+z)^\rho\bigl((1+z)^P-1\bigr)^l,
 \quad
 \binom{\rho+P\beta}{M}=\sum_l a_l\binom\beta l,
 \quad p^{\max(0,l-\lfloor M/P\rfloor)}\mid a_l.
\]
Indeed $(1+z)^P-1=z^P+pQ(z)$: a term contributing degree $M$ can use at most $\lfloor M/P\rfloor$ copies of $z^P$. Tensoring and discarding coefficients divisible by $p^h$ replaces divided differentiation of order $M$ by carry differentiation of order at most $|M|/P+h$. Together with the two $6h$ unit corrections, (PG.1) proves that every digit component of $\Gamma^{-1}\partial^{[M]}\Gamma F$ satisfies the weakened chains, up to coefficients of valuation $\ge h-n_0$. \textsf{[PG2,PG7]}

Put $\Lambda(W)=\langle F,\Gamma\phi W\rangle$. The factorial identity on each fixed label/tangent block rewrites $\Lambda$ as a constant times a Fock pairing with that conjugated derivative. Its chain-compatible part is fixed by $\Pi$, hence cancels in $\Lambda(\Omega Z)-\Lambda(Z)$. The explicit error calculation gives
\[
 v\bigl(\Lambda(\Omega Z)-\Lambda(Z)\bigr)
 \ge h+n_0-2\kappa-D_m>n_0+D_l+1.
 \tag{PG.4}
\]
Thus (PG.2) survives. Since $\kappa\ge1$, integrality and factorial bounds give
\[
 b(\phi G)\ge n_0-\kappa,\qquad
 v\Lambda(G)\le n_0+D_m+D_l+3\kappa.
 \tag{PG.5}
\]
Finally $12h\le24c$ and $30h\le60c$. These are exactly the required orders and Gram bound, with $E=D_m+D_l+2\kappa$ and \emph{$R=0$}. The report supplies the block cancellation in (PG.4); this is a construction of the constrained test, not an assumption that a suitable nonzero pairing exists. \textsf{[PG9--PG11]}

\section{Digit factorials, normal tests, and admissible assembly}
\label{sec:digit-tests}

\status{We prove the factorial-unit approximation, construct the normal test, and perform the coefficient selection with its exact scalar. We then prove outgoing symbol membership under the stated frame, degree and margin data. The incoming assertion \texttt{MemAt}, placing the whole operator-word family in the signed-chain space, is proved in IN/WT under the same supplied data. The expanded DT note gives the segment-by-segment membership calculation.}

\subsection{A shorter factorial-unit approximation}

Let $p$ be any prime, $P_v=p^{g_v}$, $0\le\rho_v<P_v$, and $\mathbf b\in\mathbb N^V$, for a finite set $V$. Put
\[
 \Gamma(\mathbf b)=\prod_{v\in V}
 \frac{(\rho_v+P_vb_v)!}{\rho_v!\,b_v!\,(P_v!)^{b_v}}.
 \tag{DT.1}
\]
This is a $p$-adic unit. For every integer $c\ge1$ and $\varepsilon\in\{1,-1\}$, there is a polynomial
\[
 Q(\mathbf b)=\sum_{|\mathbf j|\le2(c-1)}q_{\mathbf j}
       \prod_v\binom{b_v}{j_v},\qquad q_{\mathbf j}\in\mathbb Z_{(p)},
 \qquad v_p(\Gamma(\mathbf b)^\varepsilon-Q(\mathbf b))\ge c.
 \tag{DT.2}
\]
The bound is independent of $|V|$ and holds also at $p=2,3$. It improves the source interfaces $4c$ and $6c$; later numerical constants are left unchanged.

Here is a finite-product proof. If $t_i\ge1$ and $T=\sum_i t_i$, the ordered sum
\[
 S_{\mathbf t}(b)=\sum_{0\le u_1<\cdots<u_m<b}u_1^{t_1}\cdots u_m^{t_m}
\]
is an integer-valued polynomial of degree at most $T+m\le2T$. Induct using discrete summation and $\sum_{u<b}\binom uj=\binom b{j+1}$; Newton coefficients are integer finite differences. Thus a factor $F(u)=1+\sum_{t\ge1}a_tu^t$ with $a_t\in p^t\mathbb Z_{(p)}$ has
$\prod_{u<b}F(u)$ approximated modulo $p^c$ by keeping total grade $T<c$. The contribution of grade $T$ is an integral binomial polynomial of degree at most $2T$, with coefficients divisible by $p^T$. Products over several variables retain this property by adding grades. \textsf{[DT1--DT2]}

Write $U(n)=\prod_{1\le i\le n,\ p\nmid i}i$. Removing multiples of $p$ repeatedly gives
\[
 \frac{(\rho+Pb)!}{\rho!b!(P!)^b}
 =\prod_{i=1}^{\rho}\left(1+\frac{P}{i}b\right)
  \prod_{\ell=1}^{g}\prod_{u=0}^{b-1}
       \prod_{\substack{1\le a\le p^\ell\\p\nmid a}}
            \left(1+\frac{p^\ell}{a}u\right).
 \tag{DT.3}
\]
Every nonconstant linear coefficient is in $p\mathbb Z_{(p)}$. Apply the graded product argument; for the inverse use the finite geometric expansion modulo $p^c$. This proves (DT.2), including uniformity in all natural $\mathbf b$, without an infinite-product convergence argument. \textsf{[DT3--DT4]}\cite{digit-approximation}

\subsection{Normal projection with its exact scalar}

The coefficient multiplier $\prod_v\binom{\alpha_v}{j_v}$ is the operator
$X^{\mathbf j}\partial^{[\mathbf j]}$. Its conjugation module has all coordinate weights at most $|\mathbf j|$, so it increases a prefix's largest highest-weight coordinate by at most that amount. It also preserves each letter degree. This proves the weight cost of (DT.2), without an integral isotypic decomposition. \textsf{[DT5--DT6]}

In the actual local setting take $w^2=p$. On the normal polynomial space let $e$ be the joint projector onto the chosen prefix types, $D$ its summed IW4 denominator, and $\nu$ a monomial of normal multidegree $d^{\rm no}$. Define
\[
 F_{\rm no}=v_p\!\left(\prod_xd_x^{\rm no}!\right),\quad U=v_p(\nu!),\quad
 c_Y=w^{F_{\rm no}}p^Dw^{-U},\qquad Y=c_Ye(X^\nu).
 \tag{DT.4}
\]
The polynomial $w^{-U}X^\nu$ belongs to the weighted Fock lattice. Projector integrality followed by the factor $w^{F_{\rm no}}$ makes $Y$ coefficient-integral. Apply (DT.2) to the inverse normal factorial unit and set $N_0=\operatorname{diag}(Q)Y$. Then
\[
 N_0-\operatorname{diag}(\Gamma_{\rm no}^{-1})Y
       \in p^c\mathcal O[X],\qquad
 v(c_Y)=F_{\rm no}/2+D-U/2.
 \tag{DT.5}
\]
Its prefix first weights increase by at most $2(c-1)\le4c$, exactly the allowed \texttt{NormTest} slack. Thus the substantive \texttt{NoAt} construction is explicit. \textsf{[DT7]}

\subsection{Selecting the plane coefficient and preserving detection}

Start with a coefficient-integral polynomial $G$ having a unit coefficient at $\alpha_0$. Split its true exponents uniquely as $\rho+P\beta$, and separate plane and normal carries. Put
$V_P=v_p(\beta_{P,0}!)$ and $V_N=v_p(\beta_{N,0}!)$.
For the resulting normal coefficient polynomial $y$, $b(y)\le V_N/2$.
Since the finite joint projectors sum to the identity, one satisfies $b(ey)\le b(y)$.
Choose a normal monomial $\nu$ attaining this norm, and let $P_c$ be its plane coefficient after projection. Then
\[
 b(P_c)+v_p(\nu!)/2\le(V_P+V_N)/2.
\]
Choose $z$ of valuation $b(P_c)$ from an attaining coefficient and the appropriate power of $w$; set $F_2=z^{-1}P_c$, so $b(F_2)=0$.
The incoming \texttt{MemAt} condition, now supplied by IN/WT, is precisely what puts $F_2$ in the required signed-chain space. \textsf{[DT8]}

The exact factorial identity on this digit block is
\[
 \alpha!=C_\rho\,\beta_P!\,\beta_N!\,
       \Gamma_{\rm pl}(\beta_P)\Gamma_{\rm no}(\beta_N),\qquad
 C_\rho=\prod_v\rho_v!\prod_x(P_x!)^{d_x^{\rm pl}+d_x^{\rm no}}.
\]
Fock self-adjointness therefore gives the ideal pairing scalar
$k_0=C_\rho c_Y\nu!z$. With $F_{\rm tot}=v_p(\prod_xd_x!)$ and
$F_{\rm pl}=v_p(\prod_xd_x^{\rm pl}!)$, the explicit size condition
\[
 D+(F_{\rm pl}-V_P)/2+(F_{\rm no}-V_N)/2\le D_n
\]
implies
\[
 v(k_0)\le F_{\rm tot}-F_{\rm pl}/2+D_n.
 \tag{DT.6}
\]
Indeed $v(C_\rho)+V_P+V_N=v_p(\alpha_0!)\le F_{\rm tot}$. Replacing the ideal normal factor by $N_0$ introduces pairing error of valuation at least $c+F_{\rm tot}-\kappa_{\rm fl}$. These are the exact \texttt{PickAt} estimates, and CL8 already gives $D_n\le9Nk^2\log_p a$.

\subsection{Why the output is an allowed symbol}

This last step concerns \texttt{MembAt}, distinct from \texttt{MemAt}. Under the source's \texttt{MFH} frame and degree conditions, each relevant slack satisfies $s\ge128c$.
At a cut put $T=\texttt{thStar}$ and $\tau=T-s/2-60c$.
PG's chain and link bounds give plane order at least $\tau$; the normal first weight is at most $a+4c$.
For an even position the definition of $T$ gives $T\ge a_{\rm after}+m+s$, hence $\tau\ge a_{\rm after}+m+4c$. At an odd position use the next cut with $a_{\rm before}$. Also $T\ge B/2$ gives $\tau\ge B/2-s$. These are exactly the two separation bounds defining the segment symbol space. \textsf{[DT10]}\cite{test-membership}

All segment constraints cut out rational subspaces. Their intersections with the coefficient lattices are saturated and therefore direct summands over the DVR. Extending their bases proves that simultaneous constraints on disjoint segments give an integral sum of products of constrained segment factors. No integral decomposition into irreducible representations is needed. \textsf{[DT9]}
The carry identity $p^{z_j-x}p^{z_{j+1}-z_j}=p^{z_{j+1}-x}$ verifies the actual Frobenius substitution on variables. Goodness and \texttt{BlockOK} give the physical degrees; $\rho<P$ gives the label bounds. The normal order is fixed, with plane/normal roles exchanged at odd positions. Thus the assembled test is in the permitted integral symbol span. The DT report proves all four interfaces \texttt{DecAt}, \texttt{AlgAt}, \texttt{DegAt}, and \texttt{LabAt} pointwise. \textsf{[DT10]}

The initial-symbol and operator-word proof of incoming \texttt{MemAt} is now supplied by IN/WT. Its initial route retains the fixed total-degree blocks. The stronger unblocked assertion is unnecessary; IN records the exact source interface. The rank theorem is now proved in Section~\ref{sec:rank-closure}. FG/OR/CP/EA below supply the full local Type-II frame construction, realization and numerical compatibility. \textsf{[DT11, IN1--IN9, WT1--WT9]}

\section{Primitive symbol selection and the dimension ingredients}\label{sec:selected-dimension}

\status{The integral selection, nonzero multiplier, and normal-degree bound are proved here. The following elementary ingredients feed the complete quantitative count in Sections~\ref{sec:coarse-lift}--\ref{sec:rank-closure}.}

\paragraph{A residue basis gives an actual integral summand.}
Let $O$ be a DVR, $\mathbb f$ its residue field, and $G$ a set of polynomials supported on a fixed finite set $S$ of monomials. Choose $g_1,\ldots,g_d\in G$ whose reductions form a basis of $\operatorname{span}_{\mathbb f}\overline G$. Their coefficient matrix has a $d\times d$ minor $A$ with unit determinant. If $\pi_S$ selects those coefficient positions, then
\[
 R=A^{-1}\pi_S:O[\mathbf X]\longrightarrow O^d,
 \qquad R(g_i)=e_i.
\]
Thus $O[\mathbf X]=\bigoplus_i Og_i\oplus\ker R$. In particular, the selected summand has a basis of \emph{actual generators}, and its rank is exactly the residue-field dimension. This proves \texttt{RestSel}, by a shorter argument than successive division by a uniformizer. It does not estimate that dimension. \textsf{[SD1]}\cite{symbol-selection}

A finite saturated polynomial lattice reduces injectively modulo the uniformizer: if $\sum c_i b_i=\varpi P$, saturation puts $P$ in the lattice, and basis coordinates force every $c_i$ to be divisible by $\varpi$. Intersections with monomial-support conditions preserve saturation. For segment products, independence also survives a monomial substitution whenever the map
\[
 (f,(m_j)_j)\longmapsto \ell_f+\sum_j E_jm_j
\]
is injective on the allowed supports: it sends distinct tensor-product monomials to distinct monomials. In the source, this injectivity is obtained by successive division in
$n=\ell+P(u+Qv)$, with $0\le\ell<P$ and $0\le u<Q$; a missing old-frame term has $v=0$. CG1 and LA now prove the full geometry, saturation and exact-degree identifications for the actual segment spaces. \textsf{[SD2]}

\paragraph{The determinant multiplier survives reduction.}
For each late letter, take the last $r$ even selected letters preceding it. They occupy distinct positions modulo $n'$. Indeed, every interval of length $n'$ contains the $r$ periodic translates of one frame; if the chosen $r$ predecessors spanned a whole period, those frame letters and the last predecessor would give at least $r+1$ predecessors after the first. The substituted determinant of these rows is nonzero over every field: setting its diagonal variables to $1$ and other entries to $0$ makes it $1$, even after the prescribed positive powers. Their product $\Delta$ is therefore nonzero modulo $\varpi$. \textsf{[SD3]}

\paragraph{A direct normal-degree bound.}
Put $N=|TT|$, $k=r+K$, $\mathcal L=C_{\rm La}N_s^2\log a$, and
$A=\mathcal L a/(pk)$, $B=ap^{-L_1}$. Count the last $K$ coordinates at even positions and the first $r$ at odd positions. Old-frame letters contribute zero: they are even and have only plane coordinates. At a late letter $x$ of segment $j$, Frobenius substitution multiplies the allowed normal order by $p^{\operatorname{dist}(j,x)}$, giving at most
\[
 p^{\operatorname{dist}(j,x)}o(j,x)\le A+10B.
\]
The late letters across one period partition $TT$; hence segment factors contribute at most $N(A+10B)$. Each label exponent is below a relevant power at most $Q$, so labels contribute at most $NkQ$. Products add these degrees, and sums cannot enlarge support. Consequently every selected-span polynomial satisfies the slightly sharper bound
\[
 \deg_{\rm normal}P\le N(A+10B+kQ)
                 \le10N(A+B+kQ).
\]
This proves the degree part of \texttt{Conclusion2}, with $C_{\rm deg}=10$, under its stated frame and size hypotheses. \textsf{[SD4]}

\paragraph{The exact rank-transfer step.}
Let $\overline\Lambda$ be the span of the source's coarse generators and $\overline N$ the span of the reduced actual symbols. If every coarse generator $u$ satisfies $\overline{\Delta u}=\overline g$ for some actual generator $g$, multiplication by $\overline\Delta\ne0$ injects $\overline\Lambda$ into $\overline N$. SD1 then gives an integral summand of rank at least $\dim\overline\Lambda$. The required congruence is proved separately in CG8, with an actual generator as its witness. \textsf{[SD5]}

\paragraph{Two elementary pieces of the count.}
For $r\ge2$, $Q>2b$, and any $c,\rho$, the number of vectors $0\le u_i<Q-2b$ with $c+\sum_i u_i\equiv\rho\pmod Q$ is at least
\[
 Q^{r-1}\left(1-\frac{2rb}{Q}\right).
\]
There are $Q^{r-1}$ unrestricted solutions; for each coordinate, its $2b$ forbidden values account for $2bQ^{r-2}$ solutions. Subtract their union. A terminal positive path step of degree $d$ satisfies
$d\le\mu'_1-\mu_k\le2b$ when its new first weight is $b$ and earlier total negative degree is at most $b$. Thus, if $4Yb\le Q$, applying the count to each of $r$ rows gives the factor
\[
 \left[Q^{r-1}\max\!\left(0,1-\frac r{2Y}\right)\right]^r
\]
used in \texttt{RestMon}. GP10 proves the numerical inequality $4Yb\le Q$ for the actual parameters. \textsf{[SD6]}

If $J\le N$, $r\ge2$, $k\ge5$, $Y\ge1$, and $1-1/(32Y)\le\alpha\le1$, elementary multiplication and Bernoulli's inequality give
\[
 1-\frac{Nr^2k}{Y}\le
 \alpha^J\left[(1-5/Y)(1-1/Y)(1-r/(2Y))\right]^N
\]
whenever the left side is positive. RS now proves all preceding loss bounds and the actual summation, so this gives the final source coefficient $C_{\rm rk}=1$. \textsf{[SD7]}

\paragraph{The quantitative conclusion, now proved below.}
The source chooses $C_{\rm La},C_k\ge1088$, $C_{\rm rk}=1$, $C_{\rm deg}=10$ and proves
\[
 \operatorname{rank}S_g\ge q^{-rK}
 \left(1-\frac{Nr^2k}{Y}\right)
 \prod_{t\in TT}\binom{\deg_t+k-1}{k-1},\qquad q=p^{n'}.
\]
CG proves the coarse-generator congruence; LA proves the actual segment-product injection and exact degree matching; GP proves good-path membership, independence and abundance; RS proves degree padding and the complete residue-family summation. RK assembles these with SD1, SD3 and SD4 into the full local theorem. SD8 is retained as the stable identifier of this formerly open frontier, now discharged under the supplied-data hypotheses. \textsf{[SD8]}\cite{symbol-density}

\section{Lifting the coarse generators modulo the uniformizer}
\label{sec:coarse-lift}

\status{The actual \texttt{GeneratorModuloUniformiser} congruence is proved here, under the source's frame and numerical data with $C_{\rm La},C_k\ge32$. The integral path basis IP3--IP4 supplies the needed parts of \texttt{NestedTypeBound}. The coarse-family dimension is a separate counting question.}

\paragraph{A common integral basis for the actual segment spaces.}
List the segment letters $\mathcal W_j$ increasingly as $x_0,\ldots,x_{N-1}$, including
the preceding frame. Write $\mathsf X_R(m,\theta,o)$ for the source's
$\mathrm{XX}'$ over a domain $R$. It is the span of separated plane/normal
products $AG$ satisfying the prefix gaps, intersected with the prescribed
normal-degree support bound. Normal variables are absent from the
preceding frame.

Apply IP3--IP4 separately to the plane and normal variables. The paired
polynomials $\Psi_{\nu,I'}\Psi_{\kappa,I''}$ form an integral basis;
the allowed subset is exactly
\[
\begin{array}{ll}
 b''_i=0,&x_i\text{ in the preceding frame},\\
 \kappa^{i+1}_1+m_{x_i}\le\nu^i_r,\quad b''_i\le o_{x_i},
       &x_i\text{ late and even},\\
 \kappa^i_1+m_{x_i}\le\nu^{i+1}_r,\quad b'_i\le o_{x_i},
       &x_i\text{ late and odd},\\
 \theta_{x_i}\le\nu^i_r,&x_i\text{ late}.
\end{array}                                                    \tag{CG.1}
\]
Indeed, expand $A,G$ using their respective lower and upper prefix
bounds. Their tensor products satisfy the gap inequalities. Projection
onto the normal-degree support condition keeps or kills each whole
paired basis vector, since its plane and normal letter degrees are fixed.
Conversely an allowed pair satisfies the defining bounds by choosing
$M_{x_i}=\kappa^{i+1}_1$ at an even letter and $\kappa^i_1$ at an odd one.
This proves the claimed basis identification.

The same integral polynomials give the basis over $O$ and its residue
field $\mathbb f$. Thus the actual $\mathsf X_O$ is saturated, reduction
maps it onto $\mathsf X_{\mathbb f}$, and projection to a total letter
multidegree preserves it. This proves the first, third and sixth
\texttt{NestedTypeBound} clauses used below and by the segment-product
argument. No integral isotypic splitting is required. \textsf{[CG1]}

\paragraph{The coarse scale and its margin.}
Put $k=r+K$, $z_j=f_{j,r-1}$ and $d_j(x)=z_j-x$. For a late letter $x$,
let $y=y_j(x)$ be its oldest of $r$ even predecessors, and set
$B_x=\deg_{y\bmod n'}/p^{d_j(y)}$, $s_x=\lfloor B_x/(k^2Y)\rfloor$.
Let $D_j$ be the product of the corresponding plane determinants to
powers $s_x$, and $\Delta=\prod_j\phi_j(D_j)$, where
$\phi_j(X_{x,c})=Z_{x\bmod n',c}^{p^{d_j(x)}}$.
Its degree is $E$, and the coarse target degree is $a'=d-E$, where
$d_t=\deg_t-E'_t$ on $TT$ and is zero elsewhere.
The source chooses
\[
 T_j=\frac{a}{8Yk^2N_s p^{d_j(f_{j-1,0})}},\qquad
 P_j=p^{R_j}\le T_j<pP_j,\qquad m_0=kN_s+1.           \tag{CG.2}
\]
The proved degree and scale estimates give
\[
 a'_t\ge0.9\deg_t\ (t\in TT),\quad E\le d,\quad P_j>m_*,\quad
 8Yk^2N_s p^{d_j(y)}P_j\le a\quad(y\in\mathcal W_j),  \tag{CG.3}
\]
where here $m_*=\texttt{mStar}$. \textsf{[CG3,RS2--RS3]}

Put $\mathcal L=C_{\rm La}N_s^2\log a$. For a coarse segment of letter degree $b$, put
$q_y=a'_{y\bmod n'}/(p^{d_j(y)}P_j)$. Its threshold and order bound are
\[
\begin{aligned}
 \theta^c_x&=\frac23
 \begin{cases}
  \lfloor q_y\rfloor-(k+r+1)-m_0,&y\text{ in the preceding frame},\\
  b(y),&y\text{ intermediate},
 \end{cases}\\
 o^c_x&=\frac{\mathcal L a}{pkp^{d_j(x)}P_j}
       +\mathbf1_{x\text{ even}}\frac{4a}{p^{d_j(y)}P_j}.
\end{aligned}                                                    \tag{CG.4}
\]
These define the coarse $\mathsf X(m_0,\theta^c,o^c)$; the actual
$M_j(b)$ also fixes degree $b$ and bounds current-frame plane exponents.
A coarse generator is
$u=Z^\ell\prod_j\phi_j(\operatorname{expand}_{P_j}y_j)$, with
$y_j\in M_j(b_j)$, physical degree $a'$, and label exponents
$\ell_{t,c}<p^{d_j(x)}P_j$ at the unique late representatives.

For a nonzero coarse generator, an intermediate physical row occurs
in only one segment. Its degree identity and its label bound give
$b(y)>q_y-k$. Set $c_y=\deg_{y\bmod n'}/(p^{d_j(y)}P_j)$.
Then $q_y\ge0.9c_y$ and $c_y\ge40kN_s$ by (CG.3).
Substitution into either case of (CG.4), using
$\lfloor q_y\rfloor>q_y-1$, proves
\[
 P_j(\theta^c_x-kN_s)\ge B_x/2,
 \qquad P_j o^c_x+kP_j\le o_x.                       \tag{CG.5}
\]
For the second inequality, the frame gap gives
$d_j(y)\ge d_j(x)+L_1+1$; hence both $kP_j$ and
$4a/p^{d_j(y)}$ are at most $a/p^{L_1+d_j(x)}$.
This retains the source's even/odd coefficients $10,1$ in $o_x$.
The companion gives both floor cases explicitly. \textsf{[CG4--CG5]}

\paragraph{The local lift.}
Split a coarse label at its unique late representative:
\[
 \ell_{t,c}=\ell^0_{t,c}+p^{d_j(x)}e_{j,x,c},\qquad
 0\le\ell^0_{t,c}<p^{d_j(x)},\quad0\le e_{j,x,c}<P_j. \tag{CG.6}
\]
For a coarse paired basis product $A_0G_0$, over $\mathbb f$ set
$A=A_0^{P_j}\ell_P D_j$ and $G=G_0^{P_j}\ell_N$, where
$\ell_P\ell_N=X^{e_j}$.
The canonical path coefficients lie in the prime field. Frobenius gives
\[
 \overline{D_j\operatorname{expand}_{P_j}(A_0G_0)X^{e_j}}=AG.
\]
Expansion fixes coefficients and is applied linearly to a general
coarse polynomial; no inverse coefficient Frobenius is assumed.

Let $d_S$ count odd plane label exponents and $d_V$ even normal label
exponents. Then $d_S+d_V\le kN_sP_j$. Take
$(\nu_*,\kappa_*)=(\nu^i_r,\kappa^{i+1}_1)$ at an even $x_i$ and
$(\nu^{i+1}_r,\kappa^i_1)$ at an odd one. At the respective prefixes,
\[
 v(A)\ge P_j\nu_*-d_S+s_x,\qquad
 M(G)\le P_j\kappa_*+d_V.                            \tag{CG.7}
\]
The determinant belonging to $x$ contributes $s_x$ because all its
rows precede $x$; every other determinant has nonnegative order.
The coarse gap $\nu_*-\kappa_*\ge kN_s+1$ and $P_j>m_*$
give the fine gap $m_*+s_x$. The always-required minimum follows from
the first inequality in (CG.5). The normal degree is at most
$P_jo^c_x+kP_j\le o_x$. Supports are unchanged. Thus $AG$ lies in the
fine $\mathsf X_{\mathbb f}$, and CG1 lifts it integrally.
Linearity handles an arbitrary coarse $y_j$; projecting its lift to
the letter degree of $D_j\operatorname{expand}_{P_j}(y_j)X^{e_j}$
preserves both membership and congruence. \textsf{[CG2,CG6--CG7]}

\paragraph{One allowed global generator.}
Call these homogeneous lifts $x_j$. The digit identity gives
\[
 \Delta u=Z^{\ell^0}\prod_j\phi_j
       (D_j\operatorname{expand}_{P_j}(y_j)X^{e_j}),
 \qquad g=Z^{\ell^0}\prod_j\phi_j(x_j).               \tag{CG.8}
\]
The segment replacements preserve all physical degrees, so $g$ has
degree $E+a'=d$, and its label and factors are allowed. Consequently
$g\in\mathrm{Gen}$ and $\overline{\Delta u}=\overline g$.
For $u=0$ take $g=0$. This closes the coarse-generator congruence
required in SD5 and SD8. \textsf{[CG8--CG9]}


\noindent\textit{Primary source route.}\cite{coarse-lift}

\section{Assembling the actual segment spaces}
\label{sec:segment-assembly}

\status{This closes the actual segment-product injection and degree matching retained in SD2. With the shared integral path identification CG1 and the numerical bound RS3, it proves \texttt{RestLam} under the source's \texttt{Dat.Hyp}, with $C_{\rm La},C_k\ge1088$. Counting sufficiently many segment paths is a separate argument.}

\paragraph{LA1--LA2. The frame geometry and its scales.}
Write $k=r+K$, $f_{j,i}=\mathrm{fr}(j,i)$, $z_j=f_{j,r-1}$ and $s_j=f_{j-1,0}$.
Frame validity says that the $r$ selected frame letters are ordered, consecutive frames are separated, there are no selected letters between adjacent entries of a frame, and
$f_{j+J,i}=f_{j,i}+n'$.
Consequently segment $j$ has precisely the old-frame letters $f_{j-1,i}$ and the late selected letters
\[
 L_j=TT_{\mathbb Z}\cap(z_{j-1},z_j].
\]
For $0\le j<J$, these late sets partition the selected residue positions: their union lies in the half-open interval $(z_{-1},z_{-1}+n']$.
Old-frame positions likewise occur at most once per period. An old/late overlap is exactly a current-frame letter, shared with the next segment.

The apparently complicated coarse powers have a simple exact form. Put
\[
 H=\log_p\left\lfloor\frac{a}{8Yk^2N_s}\right\rfloor,
 \qquad A_j=\max\{z_j,H+s_j\}.
\]
Here $\log_p$ denotes the integer logarithm, with $\log_p0=0$.
The identity $\log_p(N/p^d)=\max(\log_pN-d,0)$, with integer division, gives
\[
 R_j=\max\{H-(z_j-s_j),0\},\qquad
 P_j(x)=p^{A_j-x}\ (x\le z_j),\qquad Q_j=p^{A_{j+1}-A_j}. \tag{LA.1}
\]
Both $z_j$ and $s_j$ increase strictly, so these exponents have the stated signs.
Also $A_{j+J}=A_j+n'$. Thus
\[
 P_{j+1}(x)=P_j(x)Q_j\quad(x\le z_j),
 \qquad \prod_{j=0}^{J-1}Q_j=p^{n'}. \tag{LA.2}
\]
These facts prove every condition of the source's \texttt{LamGeo}, including the overlap at the period boundary.

\paragraph{LA3. Saturation for the actual restricted spaces.}
Let $O$ be the source DVR with residue field $\mathbb f$.
Its segment lattice is exactly
\[
 M_j(b)=XX'_j(m_0,\theta_b,o_c)
 \cap\{\text{row degrees }b\}
 \cap\{m(f_{j,i},c)<Q_j\text{ for }c<r\},
 \qquad m_0=kN_s+1.
\]
CG1 identifies the first module with the span of a subset of the same integral paired path basis over $O$ and $\mathbb f$; hence it is saturated.
Its hypotheses hold because $m_0\ge0$ and
$\#\mathrm{letters}_j+r\le N_s+r\le K$.
The two remaining restrictions are coefficient-support conditions, so their intersections preserve saturation.
Fixed row degrees and finite letter support make $M_j(b)$ finite free.
An $O$-basis stays independent after reduction: a lifted relation
$\sum c_i b_i=\varpi P$ puts $P$ back in $M_j(b)$ by saturation, and basis coordinates imply $\varpi\mid c_i$ for every $i$.
This establishes the actual \texttt{RestLamA} input, rather than assuming that reduction preserves an arbitrary lattice basis.

\paragraph{LA4. Recovering every label and segment monomial.}
An allowed label $\ell$ satisfies, at the unique late representative $x\in L_j$ of $t$,
\[
 0\le\ell_{t,c}<P_j(x),\qquad
 \sum_c\ell_{t,c}\le a'_t,\qquad
 \sum_c\ell_{t,c}\equiv a'_t\pmod{P_j(x)}.
\]
Write $n_j(x)=(a'_t-\sum_c\ell_{t,c})/P_j(x)$.
The terminal parameters satisfy $0\le e_{j,i}<rQ_j$ and
$e_{j,i}\equiv n_j(f_{j,i})\pmod{Q_j}$.
For allowed segment monomials $m_j$, the final exponent at a frame plane coordinate has the form
\[
 N_{t,c}=\ell_{t,c}+P_j(x)\bigl(m_j(x,c)+Q_jm_{j+1}(x,c)\bigr),
 \qquad m_j(x,c)<Q_j. \tag{LA.3}
\]
Indices and letters are translated periodically when necessary.
Reduction modulo $P_j(x)$ recovers $\ell$; division and reduction modulo $Q_j$ recover the current exponent and then the old exponent.
At a nonframe position, or a normal coordinate $c\ge r$, the old term is zero, so division alone suffices.
Outside $TT$ all exponents vanish.
Finally the current-frame row sums recover $e_{j,i}$.
Thus the exponent map is injective across \emph{all} label and terminal-parameter families, proving \texttt{RestLamInj}.

\paragraph{LA5. Exact row degrees, including natural subtraction.}
The prescribed degrees $b_j=b_{\ell,e,j}$ are $n_j(x)$ at intermediate letters, $e_{j,i}$ at current-frame letters, and
$(n_{j-1}(x)-e_{j-1,i})/Q_{j-1}$ at old-frame letters.
The last numerator is nonnegative: RS3 gives
\[
 \frac{a'_t}{P_j(x)Q_j}\ge\frac{36}{5}Yk^2N_s\ge k+r+1.
\]
Since $\sum_c\ell_{t,c}<kP_j(x)$, it follows that
$n_j(x)>Q_j(k+r+1)-k\ge rQ_j>e_{j,i}$.
The defining congruence therefore gives the exact identity
\[
 e_{j,i}+Q_j\frac{n_j(x)-e_{j,i}}{Q_j}=n_j(x).
\]
Summing (LA.3) over coordinates now gives $a'_t$ at every frame position; the intermediate positions give the same result directly.
This proves \texttt{RestLamDeg}, without importing the source's full \texttt{RestPad} argument.

\paragraph{LA6. The resulting dimension inequality.}
Choose the residue-independent bases from LA3 for every $(\ell,e,j)$.
LA4 sends distinct tensor-product monomials to distinct final monomials, so the corresponding product map is injective on the direct sum over $(\ell,e)$.
LA5 puts every product in the actual coarse generator set.
Consequently
\[
 \boxed{\displaystyle
 \sum_{\ell\in LA}\sum_{e\in EE(\ell)}
 \prod_{j=0}^{J-1}\operatorname{rank}_O M_j(b_{\ell,e,j})
 \le\dim_{\mathbb f}\overline\Lambda.} \tag{LA.4}
\]
This is the complete \texttt{RestLam} assembly. Its proof does not estimate the summands; their abundance belongs to the separate segment and residue-family counts.


\noindent\textit{Primary source route.}\cite{segment-assembly}

\section{Good paths give the required segment ranks}
\label{sec:good-paths}

\status{The actual \texttt{RestMem}, \texttt{RestInd},
\texttt{GoodPathWeight}, \texttt{BetaRatioBound} and \texttt{RestSeg}
are proved under \texttt{Dat.Hyp} and $C_{\mathcal L},C_k\geq1088$.
The proof uses the earlier IP integral path basis, CB dimension formula,
SD6 residue count and the RS2--RS3 arithmetic. FG/CP/EA below supply the actual local frames and compatible parameters;
global abundance of the required orbit data remains a separate step.}

\paragraph{GP1--GP2: the actual paths and their total mass.}
Assume \texttt{Dat.Hyp} and $C_{\mathcal L},C_k\geq1088$.
In one segment write $x_0<\cdots<x_{N-1}$, $n_0=N-r$, $k=r+K$.
The first $r$ letters are the old frame; the last $r$ are the current
frame $\tau_t$.  Put $P_i=P_z(j,x_i)$, $Q=Q_{c,j}$, $m=kN_s+1$, and
\[
 B_i=\begin{cases}(B_{w,i}-m)_+,&i<r,\\b(x_i),&r\leq i<n_0,\end{cases}
 \qquad \beta=\sum_{i<n_0}B_i,\qquad
 R=\prod_{i<n_0}\binom{B_i+k-1}{k-1}.
 \tag{GP1}
\]
For counting, $b$ need only satisfy the source's intermediate-degree
condition \texttt{AdmLL}; membership below additionally requires
\texttt{BOK}.
The RS bounds give $k\leq B_i\leq2a/P_i$ and
$B_{i'}\geq(p^2/5)B_i$ for earlier letters $i<i'<n_0$ of the same
parity.  Also $N\leq N_s$, $N\geq3r$, and
$h:=K-2N\geq k/2$.  For $i\geq r$, let $o(i)$ be the $r$-th most
recent vector letter; it lies before $n_0$.

Take the literal signed horizontal-strip paths $\mu^0=0,\ldots,\mu^N$:
their first $n_0$ strip sizes are $B_i$, and their terminal conditions are
$\mu^{n_0+t}_l=\beta$ for $l<t$.  Let $\nu,\kappa$ be the first $r$
and last $K$ rows.  A path has weight $w(\mu)=d_K(\kappa^N)$, the
number of its standard normal arrays.  Signed Pieri at the fixed-size
steps and GT branching at the terminal steps give
\[
 \sum_{\mu}w(\mu)=R,\qquad
 \#\{\mu\}\leq\exp(2N^2L),\qquad L=\log(1+2aN).
 \tag{GP2}
\]
For the second bound, at each stage only the first $N$ and last $N$
coordinates can be nonzero; each has at most $1+\beta\leq1+2aN$
possible values.  Terminal branching is exact: deleting the forced
$\beta$ entries turns every terminal strip into one interlacing row.

\paragraph{GP3--GP6: why the bad paths have small mass.}
At stage $i$ delete its $(i-n_0)_+$ forced heads, obtaining $\eta^i$.
Let $T_+(i,u)$ be the sum of the positive parts after coordinate $u$,
and $T_-(i,u)$ the same sum for the reversed negative parts.
Let $S_\pm(i,u)$ be the suffix sums of the reverse-chronological lists
of earlier preterminal vector/covector degrees.
Horizontal-strip interlacing gives
\[
 T_\pm(i,u)\leq S_\pm(i,u).
 \tag{GP3}
\]
Indeed an active step prepends its degree to the reference list and
satisfies $T(i+1,u+1)\leq T(i,u)$ and
$T(i+1,0)\leq T(i,0)+B_i$; an opposite-sign or terminal step decreases
the relevant tails.  Thus a deficit $X$ at a suffix persists to the
endpoint, shifting its index once for each later active step.

Set $f_h(x)=\sum_{t=1}^h\log(1+x/t)$.  For a reference list
$b_u\geq h$, $b_u\geq4b_{u+1}$ and nonnegative entries $a_u$ satisfying the
suffix bounds, concavity and summation by parts give
\[
 \begin{split}
 \sum_u\bigl(f_h(a_u)-f_h(b_u)\bigr)
 &\leq\sum_u(g_u-g_{u-1})\Delta_u,\\
 g_u=f'_h(b_u),\quad g_{-1}&=0,\quad
 \Delta_u=\sum_{l\geq u}(a_l-b_l)\leq0.
 \end{split}
 \tag{GP4}
\]
Since $h/(x+h)\leq f'_h(x)\leq h/(x+1)$, the coefficients satisfy
$g_u-g_{u-1}\geq h/(4b_u)$.  A deficit $X$ therefore contributes
at most $-Xh/(4b_u)$.

The endpoint $\lambda=\kappa^N$ has at most $N$ positive and $N$
negative entries.  In its Weyl product, pairs of nonzero entries cost
at most $2N^2L$.  For each nonzero entry, its interactions with zero
entries are bounded by $f_h(|\lambda_l|)+2NL$: the distances beyond
$1,\ldots,h$ number at most $2N$.  Hence
\[
 \begin{aligned}
 \log w(\mu)&\leq
 \sum_l f_h((\lambda_l)_+)+\sum_l f_h((-\lambda_l)_+)+6N^2L,\\
 \log R&\geq\sum_{i<n_0}f_h(B_i).
 \end{aligned}
 \tag{GP5}
\]

Put $\Xi_i=\mathcal L a/(pkP_i)$.  Call a path typical if a covector
step loses at most $\Xi_i$ positive mass, a vector step $i\geq r$
loses at most $\Xi_i$ negative mass, and
$\mu^i_{r-1}\geq2B_{o(i)}/3$.  Failure of a mass condition creates
a suffix deficit $\Xi_i$.  Its reference entry is an earlier
$B_l\leq4a/(pP_i)$, so (GP4) costs at least $\mathcal L h/(16k)$.
Failure of the last condition creates a deficit $B_{o(i)}/3$ at that
reference entry, costing $h/12$.  Combining (GP2)--(GP5),
\[
 \sum_{\text{atypical }\mu}w(\mu)
 \leq R\exp\!\left(8N^2L-\frac{\min(\mathcal L,k)}{32}\right)
 \leq R a^{-N_s^2}.
 \tag{GP6}
\]
Here $L\leq2\log a$ and
$\min(\mathcal L,k)\geq1088N_s^2\log a$; the estimate is stronger
than the source's exponent with $11N^2L$.

\paragraph{GP7--GP9: good paths yield actual independent generators.}
Tracing row $r-1$ backwards through its $r-1$ later vector steps and
summing the geometric degree separation gives
$\mu^i_{r-1}\leq(1+5/(p^2-5))B_{o(i)}<2B_{o(i)}$.
Interlacing then bounds the normal-degree increase at a vector step
by $\Xi_i+2B_{o(i)}$, and the plane-degree decrease at a covector
step by $\Xi_i$.  These are precisely the source's
$\mathrm{thc}=2B_{o(i)}/3$ and
$\mathrm{oc}_i=\Xi_i+\mathbf1_{x_i\text{ even}}4a/P_{o(i)}$
bounds.  Thus the good-path/normal-array set $T$ satisfies
$|T|\geq(1-a^{-N_s^2})R$.

For $(\mu,I)\in T$, take the actual auxiliary exponents
$u_{tc}+2\beta<Q$ with terminal row degrees
$d_{n_0+t}+\sum_cu_{tc}=b(\tau_t)$.  The generators are exactly
\[
 \begin{aligned}
 F_{\mu,I,u}&=C_b\prod_{t,c}X_{\tau_t,c}^{u_{tc}}
 \Psi_{\rm pl}(\nu,I_{\rm can})\Psi_{\rm nor}(\kappa,I),\\
 C_b&=\det(X_{W})^m\prod_{x\in W}X_{x,0}^{\,b_x-B_{w,x}}.
 \end{aligned}
 \tag{GP7}
\]
The IP prefix bounds and the crossing inequalities
$\kappa^{i+1}_0\leq\nu^i_{r-1}$ for vectors,
$\kappa^i_0\leq\nu^{i+1}_{r-1}$ for covectors prove the literal
$XX'$ conditions: $\det(X_W)^m$ supplies the additional $m$.
The normal polynomial has no old-frame variables.  Row degrees add
to $b$, and SD6 gives $d_{n_0+t}\leq2\beta$, so every terminal plane
exponent is $<Q$.  Hence $F_{\mu,I,u}\in M_j(b)$.

For independence, retain only the antidiagonal variable
$Y_t=X_{\tau_t,r-1-t}$ in each terminal plane row of the path polynomial.
Column expansion makes it a signed monomial in the $Y_t$ times a
truncated-path standard polynomial.  Its retained column sets record
the time of every terminal column rise, hence recover the full path.
The IP integral basis and disjoint terminal variables prove independence
of these specialized products, including the auxiliary monomials.
In a relation before specialization, scale the killed variables by $T$
and take the least auxiliary $T$-degree with nonzero coefficient;
this gives a forbidden relation among the specialized products.
Finally cancel the common nonzero $C_b$.  This proves
\texttt{RestMem} and \texttt{RestInd} over the actual DVR.

\paragraph{GP10--GP12: the residue bound and \texttt{RestSeg}.}
Let $D=\operatorname{dist}(j,\tau_0)$, $g=z_{j+1}-z_j$.
The preceding gap and distinct distances give
$\beta\leq4a/p^{R_j+D+L_1+1}$.  RS3 gives
$Q=p^{g+R_{j+1}-R_j}>a/(8pYk^2N_s p^{D+R_j})$.  Therefore
\[
 \frac\beta Q<\frac{32Yk^2N_s}{p^{L_1}}
 \leq\frac1{4pY}\leq\frac1{4Y}.
 \tag{GP8}
\]
The elementary SD6 residue count now supplies, for every terminal row
and every desired residue, at least
$Q^{r-1}\max(0,1-r/(2Y))$ choices of $u$.
Their resulting degrees $e_t=d_{n_0+t}+\sum_cu_{tc}$ satisfy
$0\leq e_t<rQ$.  Set $G_j(b)=|\mathrm{Idx}(j,\mathrm{restr}\,b)|$.
This index set equals $\mathrm{Idx}(j,b)$: old-frame inputs use
$B_w-m$, independently of the old values of $b$.
Membership, independence and finite-dimensionality give
$G_j(b)\leq\operatorname{rank}_O M_j(b)$ for every \texttt{BOK} $b$;
the count gives
\[
 \sum_{\substack{0\leq e_t<rQ\\e_t\equiv\rho_t\pmod Q}}
 G_j(b[e])
 \geq(1-a^{-N_s^2})R
       \left[Q^{r-1}\max\!\left(0,1-\frac r{2Y}\right)\right]^r.
 \tag{GP9}
\]
The function depends only on late row degrees.  These are exactly
the clauses of \texttt{RestSeg 1088}; no path-abundance or numerical
ratio hypothesis remains unproved.


\noindent\textit{Primary source route.}\cite{good-paths}

\section{Summing all residue families with the exact loss}\label{sec:residue-summation}

\status{The actual degree padding and all numerical and summation inputs to
\texttt{RestSum} are proved here. The segment functions are supplied by GP;
the injection into the coarse residue space is supplied by LA.}

Write $k=r+K$, $N=|TT|$, $B(n)=\binom{n+k-1}{k-1}$, and
$f_{j,i}=\mathrm{fr}(j,i)$. Set $z_j=f_{j,r-1}$ and $s_j=f_{j-1,0}$.
The late letters of segment $j$ are the selected lifts in $(z_{j-1},z_j]$;
they split into $LL_j$ and the current frame. Over one period they tile
$TT$ bijectively. A full segment lies in $(z_{j-2},z_j]$, so a residue
occurs at most twice in the full segment sets. In particular $J\le N$.
These statements follow directly from ordered periodic frames. \textsf{[RS1]}

\paragraph{The actual numerical margins.}
Each determinant row contributes at most $2a/(k^2Y)$ to its residue's
degree. There are at most $N_s$ determinants per segment; the preceding
two-occurrence bound gives
\[
 E_t\le\frac{4N_sa}{k^2Y},\qquad
 a'_t=\deg_t-E'_t-E_t\ge\frac9{10}\deg_t.
 \tag{RS.1}
\]
For the second inequality, the source hypotheses give
$a\ge32$, $k\ge5$, $N_s\ge1$, $k\ge32N_s$, $E_0\le a/32$,
and hence $E_t\le a/40$. They also give the sharper total loss
$E'_t+E_t\le5\deg_t/(kY)$. Thus all these natural subtractions are
ordinary nonnegative differences. \textsf{[RS2]}

Put $U=8Yk^2N_s$, $T_j=a/(Up^{z_j-s_j})$,
$R_j=\log_p\lfloor T_j\rfloor$, and $P_j(x)=p^{z_j-x+R_j}$.
The size assumption gives
\[
 T_j\ge4p(E_0+m_*+1),\quad
 p^{R_j}\le T_j<p^{R_j+1},\quad
 \frac{a'_{\operatorname{pos}x}}{P_j(x)}
 \ge\frac{36}{5}Yk^2N_s. \tag{RS.2}
\]
The scale relation has a short exact proof. With
$H=\log_p\lfloor a/U\rfloor$,
\[
 A_j=z_j+R_j=\max(z_j,H+s_j),\qquad
 P_j(x)=p^{A_j-x},\qquad Q_j=p^{A_{j+1}-A_j}.
\]
Both sequences inside the maximum strictly increase. Since
$A_{j+J}=A_j+n'$, this proves
\[
 P_j(x)Q_j=P_{j+1}(x),\qquad \prod_{j=0}^{J-1}Q_j=p^{n'}=q.
 \tag{RS.3}
\]
It also proves periodicity and consistency at equal-residue late
occurrences. \textsf{[RS3; LA2]}

\paragraph{Degree padding and digit partitions.}
For a label $\lambda_t$ let $L_t=\sum_c\lambda_{t,c}$, with
$0\le\lambda_{t,c}<P$, $L_t\le a'_t$ and $L_t\equiv a'_t\pmod P$.
Then $n_j(x)=(a'_t-L_t)/P_j(x)$ is an integer in
$(a'_t/P_j(x)-k,a'_t/P_j(x)]$. Terminal degrees satisfy
$0\le e_{j,i}<rQ_j$ and $e_{j,i}\equiv n_j(f_{j,i})\pmod{Q_j}$.
Equation (RS.2) in segment $j+1$ gives $n_j(f_{j,i})>e_{j,i}$.
The old-frame degree is therefore exactly $(n_j-e_{j,i})/Q_j$.

More generally, if $s<kP$ and $e<rQ$, successive floors give
\[
 \left\lfloor\frac{(\lfloor(n-s)_+/P\rfloor-e)_+}{Q}\right\rfloor
 \ge\left\lfloor\frac n{PQ}\right\rfloor-k-r.
\]
This proves the source's weaker old-frame margin $k+r+1$ and hence
\texttt{RestPad}. The exact decomposition $w_c=\lambda_c+Pv_c$ gives
\[
 B(n)=\sum_{\substack{0\le\lambda_c<P,\ L\le n\\L\equiv n\pmod P}}
 B\!\left(\frac{n-L}{P}\right).
 \tag{RS.4}
\]
If $kP\le n$, there are exactly $P^{k-1}$ labels satisfying that
congruence: any $k-1$ coordinates determine the last. \textsf{[RS4--RS5]}

\paragraph{Two binomial comparisons.}
The product formula for $B$ and Bernoulli's inequality show
\[
\begin{split}
 B(a'_t)&\ge(1-5/Y)B(\deg_t),\\
 (P')^{k-1}B(b)&\ge(1-1/Y)B(a'_t),
 \quad b=\bigl(\lfloor a'_t/P'\rfloor-(k+r+1)-(kN_s+1)\bigr)_+,
\end{split}\tag{RS.5}
\]
where the second line uses $P'=P_j(f_{j,i})Q_j$.
Indeed (RS.1) loses at most $5\deg_t/(kY)$, while
$P'b\ge a'_t-4kN_sP'\ge(1-1/(kY))a'_t$ by (RS.2).
Compare each of the $k-1$ factors and use $(1-\theta)^{k-1}\ge1-(k-1)\theta$.
\textsf{[RS6]}

\paragraph{The exact summation.}
Let $\alpha=1-a^{-N_s^2}$ and $\tau=\max(0,1-r/(2Y))$.
Suppose $G_j(b)\ge0$ depends only on late degrees and obeys the source's
terminal-residue lower bound
\[
 \sum_{e_j}G_j(\operatorname{setTau}_j(b,e_j))
 \ge\alpha W_j[Q_j^{r-1}\tau]^r\prod_{x\in LL_j}B(b(x)),
\]
where $W_j=\prod_i B((\mathrm{Bw}_j(f_{j-1,i})-kN_s-1)_+)$.
For a fixed label, the choices $e_j$ are independent; finite
distributivity makes the sum of products a product of sums.
Now sum the labels using (RS.4). At an $LL$ residue this supplies
$B(a'_t)$; at a frame residue it supplies $P_j(f_{j,i})^{k-1}$.
Periodicity shifts $W_j$ to $W_{j+1}$. Multiplication by $q^{rK}$
supplies $Q_j^K$ at each of the $r$ frame entries, and
\[
 Q_j^{r-1}P_j(f_{j,i})^{k-1}Q_j^K
   =(P_j(f_{j,i})Q_j)^{k-1}.
\]
Thus (RS.5), the tiling, and $\eta=(1-5/Y)(1-1/Y)(1-r/(2Y))$ give
\[
 q^{rK}\sum_{\lambda\in LA}\sum_{e\in EE(\lambda)}
       \prod_jG_j(b_j(\lambda,e))
 \ge\alpha^J\eta^N\prod_{t\in TT}B(\deg_t).
 \tag{RS.6}
\]
If $1-Nr^2k/Y>0$, all factors used above are nonnegative,
$\alpha\ge1-1/(32Y)$, and SD7 gives
$\alpha^J\eta^N\ge1-Nr^2k/Y$.
A nonpositive target follows directly from the nonnegative sum.
This proves \texttt{RestSum 1088 1} with every actual summation and
numerical premise discharged. \textsf{[RS7--RS8]}


\noindent\textit{Primary source route.}\cite{residue-sum}

\section{The quantitative selected-symbol theorem}\label{sec:rank-closure}

\status{The source's rank and degree conclusion is now proved for supplied
valid frames and numerical data. This closes the density frontier SD8;
FG/OR/CP/EA below construct and match the data needed by the full local Type-II theorem.}

\begin{theorem}[Selected symbols, with the source constants]
Assume the exact frame, DVR, degree and size hypotheses of
\texttt{Conclusion2}, with $C_{\rm La},C_k\ge1088$.
Put $k=r+K$, $N=|TT|$, $q=p^{n'}$ and
\[
 \mathcal B=\left(1-\frac{Nr^2k}{Y}\right)
       \prod_{t\in TT}\binom{\deg_t+k-1}{k-1}.
\]
The span $\mathcal N$ of the actual allowed integral generators contains
an integral direct summand $S_g$ of the ambient polynomial module,
with a basis consisting of actual generators and
\[
 \operatorname{rank}_O S_g\ge q^{-rK}\mathcal B.
 \tag{RK.1}
\]
For every finite field $\mathbb F_q$, the same rank is at least
$\mathcal B/\#\operatorname{Gr}(r,k)(\mathbb F_q)$.
Every polynomial in $\mathcal N$ satisfies
\[
 \deg_{\rm normal}P\le
 10N\left(\frac{\mathcal L a}{pk}+ap^{-L_1}+kQ\right),
 \qquad \mathcal L=C_{\rm La}N_s^2\log a.
 \tag{RK.2}
\]
Thus the literal witnesses are $C_0=1088$, $C_{\rm rk}=1$ and
$C_{\rm deg}=10$.
\end{theorem}

\begin{proof}
Let $\overline\Lambda$ be the coarse-generator span and $\overline N$
the reduced actual-generator span. GP supplies nonnegative segment
functions $G_j$ with the required terminal-residue sums. RS4 makes each
family degree admissible, so GP also bounds $G_j$ by the corresponding
segment rank. RS8 and LA6 give the full chain
\[
\begin{aligned}
 q^{-rK}\mathcal B
 &\le\sum_{\ell\in LA}\sum_{e\in EE(\ell)}\prod_jG_j(b_{\ell,e,j})\\
 &\le\sum_{\ell\in LA}\sum_{e\in EE(\ell)}
                  \prod_j\operatorname{rank}_O M_j(b_{\ell,e,j})
 \le\dim_{\mathbb f}\overline\Lambda.
\end{aligned}\tag{RK.3}
\]
Every multiplication in the middle inequality is between nonnegative
factors. This is the complete \texttt{RestCount} conclusion. \textsf{[RK1]}

CG8 gives, for every coarse generator $u$, an actual generator $g$ with
$\overline{\Delta u}=\overline g$. SD3 proves $\overline\Delta\ne0$.
Multiplication by $\overline\Delta$ is therefore an injective linear map
from $\overline\Lambda$ into $\overline N$. Hence
$\dim\overline\Lambda\le\dim\overline N$. \textsf{[RK2]}

All generators have one fixed finite row multidegree. SD1 chooses actual
generators with a unit coefficient minor, producing an ambient direct
summand $S_g$ of rank $\dim\overline N$. This proves (RK.1).
The graphs of the $q^{rK}$ linear maps $\mathbb F_q^r\to\mathbb F_q^K$
are distinct $r$-planes. Thus $\#\operatorname{Gr}(r,k)(\mathbb F_q)
\ge q^{rK}$; division gives the asserted plane-normalized bound when
$\mathcal B\ge0$, and a negative numerator gives a trivial bound.
Finally (RK.2) is SD4, valid on the entire span $\mathcal N$.
The standard quadratic DVR setting supplies the uniformized DVR used
above because its uniformizer $w$ satisfies $w^2=p$.
This proves every rank, basis, complement and degree clause of the
source's \texttt{Conclusion2}. \textsf{[RK3]}
\end{proof}

The following FG/OR/CP/EA sections now match these data to the full
local Type-II hypotheses. Global abundance and simultaneous parameter
existence, the separate $r=1$ construction, and the arithmetic windows for
RH and Dirichlet GRH remain as stated in Section~\ref{sec:remaining}.


\noindent\textit{Primary source route.}\cite{rank-closure}

\section{Constructing the periodic frames from marked tiles}
\label{sec:frame-geometry}

\status{The actual \texttt{FramesOK} conclusion is proved from
\texttt{OrbitHyp}, including its strict segment span and stronger $L_2$
gaps. The construction below is the source's frame, with one common
marked tile for all coordinates. Abundance of the good orbits remains
a separate question.}

\paragraph{FG1: the three concrete orbit conditions.}
Write $n=n'>0$ for the even period and
$T\subset\{0,\ldots,n-1\}$ for the selected residues. Put
\[
\begin{gathered}
 u=\lceil1/\delta\rceil_+,\qquad \ell=(6r+1)u,\qquad
 M=\lceil(h+3)10^{6r+1}\log x\rceil_+,\\
 G'=(M+5)\ell,\qquad
 N_1=\left\lceil\max\{6eG'\delta,(h+3)\log_2x\}\right\rceil_++1 .
\end{gathered}                                                    \tag{FG.1}
\]
Here $\lceil\cdot\rceil_+$ is the nonnegative ceiling, $\delta=\texttt{de}$,
$r\ge2$, $L_1\ge1$, and the actual hypotheses include
$2L_1\le L_2<u$ and $3G'\le n$. Thus $u\ge1$, and $\delta>0$ follows.

A cyclic window of length $L$ is
$\{i,i+1,\ldots,i+L-1\}\pmod n$. The source requires at most $r-1$
selected residues in every $L_2$ window and at most $N_1$ in every $G'$
window. Its third condition concerns the $t=\lfloor n/\ell\rfloor$
complete tiles. Tile $s$ has cells
\[
 C(s,c)=[s\ell+(c-1)u,s\ell+cu),\qquad 1\le c\le6r+1.
\]
It is \emph{marked} when its odd cells contain no selected point and
each even cell contains exactly one selected point, itself even.
Every cyclic block of $M+1$ tile indices must contain a marked tile.
These are precisely \texttt{CondA}, \texttt{CondIntervals} and
\texttt{CondTiles}; no frame existence is hidden in them.

\paragraph{FG2--FG4: enumerate the tiles and take their middle block.}
List all marked indices as $s_0<\cdots<s_{J-1}$.
They exist by the tile condition; moreover $t\ge3(M+5)>M+1$.
For every $j=qJ+a$ with $q\in\mathbb Z$ and $0\le a<J$, set
\[
 B_j=qn+s_a\ell.
\]
The next marked index has positive cyclic distance at most $M+1$:
apply the tile condition immediately after the current index.
Writing $n=t\ell+R$, $0\le R<\ell$, the successive physical gap is
that cyclic distance times $\ell$, with an additional $R$ exactly
at the wrap. Hence
\[
 B_{j+J}=B_j+n,\qquad
 \ell\le B_{j+1}-B_j<(M+2)\ell.                         \tag{FG.2}
\]
The remainder $R$ is retained.

In the lifted marked tile starting at $B_j$, let $a_{j,v}$ be its
selected point in cell $2v+2$, $0\le v<3r$. Then
\[
 B_j+(2v+1)u\le a_{j,v}<B_j+(2v+2)u,\qquad
 f_{j,i}:=a_{j,r+i}\quad(0\le i<r).                    \tag{FG.3}
\]
These are the actual \texttt{bf\_tp} and \texttt{bf\_F} choices.
Because $n$ is even, all lifted selected points remain even.
Uniqueness gives $f_{j+J,i}=f_{j,i}+n$.
The frames are increasing and ordered, and no selected point lies
between two consecutive entries of one frame. For a fixed $j$,
the same index $s_a$ places every $f_{j,i}\pmod n$ in cell
$2r+2+2i$. Every marked tile occurs.

The empty odd cells also prove a stronger local fact: every selected
$y<a_{j,v}$ satisfies $y+u<a_{j,v}$. Inside the tile, use the previous
even cell; outside it, use its first empty cell. Thus this statement
does not assume that intervening unmarked tiles are empty.

\paragraph{FG5: exact endpoint and span conditions.}
Let $\widetilde T$ be the periodic lift of $T$, and set
\[
 W_j=\widetilde T\cap[f_{j-1,0},f_{j,r-1}],\qquad
 L_j=\widetilde T\cap(f_{j-1,r-1},f_{j,0}).
\]
The first $r$ members of $L_j$ are $a_{j-1,2r},\ldots,a_{j-1,3r-1}$;
the last $r$ are $a_{j,0},\ldots,a_{j,r-1}$.
All are even, and each has gap $>u>L_2$ from its preceding selected
point. The same holds for $f_{j,0}$. These two lists are disjoint,
so $|L_j|\ge2r$, which includes the required $|L_j|\ge r$.
This checks the literal first/last-$r$ endpoint predicate, including
its boundary inequalities.

The exact cell positions and (FG.2) give
\[
 f_{j,r-1}-f_{j-1,0}
 <B_j-B_{j-1}+(2r-1)u
 <(M+3)\ell\le G'.                                    \tag{FG.4}
\]
The source's strict bound is preserved.

\paragraph{FG6--FG7: counts and the oldest vector.}
Reduction modulo $n$ is injective on every integer half-open interval
of length at most $n$: two distinct congruent points differ by at
least $n$. Since $G'\le n/3$, (FG.4) and the cyclic-window condition imply
$|W_j|\le N_1$.

For a selected $z\in W_j$ after the old frame, all $r$ old-frame
entries are even predecessors. Thus the source's $r$-th most recent
even predecessor $y=\operatorname{oldest}(j,z)$ exists; its
default-zero branch is unused. If $z-y\le L_1$, the $r$ most recent
even predecessors lie in $[y,z)\subseteq[y,y+L_2)$.
This contradicts the $r-1$ window bound. Therefore
\[
 z-\operatorname{oldest}(j,z)>L_1.                    \tag{FG.5}
\]
Together, these facts prove every clause of
$\texttt{Valid}(L_1,N_1)$, the strict span (FG.4), and both stronger
$L_2$ gap conditions. With $F.TT=TT$, this is the exact
\texttt{FramesOK}. The common-tile and coverage properties above also
prove the stronger \texttt{RestFrames} conclusion.

\paragraph{FG8: frame count and scope.}
The marked complete tiles are disjoint and each contains $3r$ selected
points, so $3rJ\le|TT|$. In particular the source's bound
$Jr\le|TT|$ holds. More generally that weaker bound follows for any
valid periodic frame by injecting its $Jr$ entries in one period into
$TT$. The same interval injection gives $|W_j|\le|TT|$; the existing
choice $N_s=\min(N_1,2|TT|)$ is retained. No numerical or probabilistic
existence of the input orbit has been inferred from this construction.
\noindent\textit{Primary source route.}\cite{frames-actual}


\section{The outgoing one-sided symbols really have integral lifts}
\label{sec:outgoing-realization}

\status{This proves the additional outgoing realization used by the cloud
test. Its polynomial family is \texttt{monPh}, which needs a separate
determinant extraction before the earlier SY correction applies. The
actual numerical hypotheses with $C\ge100$ suffice.}

\paragraph{OR1--OR2: minimum weight forces a determinant.}
Work first over any field. For signed letters, let $v_S(P)$ be the least
$t$-exponent after a generic conjugate of
$\operatorname{diag}(1,\ldots,1,t)$ acts on the letters in $S$.
Vectors transform directly and covectors inversely. Choose $r$ distinct
vector letters $A\subset S$, put $D_A=\det(X_A)$, and let
$B$ be the total degree in the other vector letters of $S$.
For a letter-multihomogeneous polynomial,
\[
 v_S(P)\ge\theta>B
 \quad\Longrightarrow\quad D_A^{\theta-B}\mid P
 \qquad(\theta\in\mathbb Z).                              \tag{OR.1}
\]
Indeed, modulo the prime polynomial $D_A$, the generic selected matrix has
rank $r-1$. Choose a basis whose last direction it annihilates. The selected
vectors are then fixed by the cocharacter; the other vectors contribute
degree at most $B$ in $t$, and covectors only nonpositive powers.
The lower bound $\theta>B$ forces the specialized polynomial to vanish.
Thus $D_A\mid P$.

For any other subset $S'$, transforming only its selected rows is a
rank-one matrix update. Its determinant is affine in $t$, with minimum
exponent at most $1$. Consequently
\[
 v_{S'}(D_A^mH)\le v_{S'}(H)+m.                           \tag{OR.2}
\]
Iterating proves (OR.1), since division decreases only the selected
letters' degrees. Over an integral domain the quotient remains integral:
$D_A^m$ has unit leading coefficient, so polynomial division over the
domain gives the same zero remainder as over its fraction field.

\paragraph{OR3--OR5: extracting the actual segment divisor.}
Write $k=r+K$, $Q=p^{G'}$, $c=c_P$, $N=|TT|$, and
$\alpha_*=\max_{0\le j<J}\alpha_j$. The exact exponents are
\[
\begin{aligned}
m_1&=(2Q+3)c+kQ+c,&m_2&=c\alpha_*+2c,\\
m_*&=(N_1+1)\bigl((2Q+3)c+kQ+c\alpha_*+3c\bigr).
\end{aligned}                                                   \tag{OR.3}
\]
For each late segment letter $x$, let $A_x$ be its $r$ most recent even
selected predecessors, and let $D_x$ be their plane determinant.
Let $D_{{\rm old},j}$ be the preceding-frame determinant. Extract
\[
                 Dv_j=\prod_{x\ {\rm late}}D_x^{m_1}D_{{\rm old},j}^{m_2}.
                                                               \tag{OR.4}
\]
The total extraction cost is at most $|{\rm letters}_j|m_1+m_2\le m_*$.
Formula (OR.2) controls \emph{every} prefix after each extraction; the
other-vector degrees can only decrease. Thus a starting minimum
$B_x+m_*$ suffices for all tasks. The $D_{{\rm old},j}$ task uses the prefix before
the first late letter, containing exactly the preceding frame, so its
other-vector degree is zero.

The outgoing degree assumption is only $d_y\le1.01\,\operatorname{levDeg}(y)$
on the preceding frame and intermediate letters. An even predecessor
outside $A_x$ lies in these sets and is $2h$ positions before the oldest
member of $A_x$, for a distinct integer $h\ge1$. Hence, writing $B_o$ for
that oldest member's normalized degree,
\[
 B_x\le2.02B_o\sum_{h\ge1}p^{-2h}
       =\frac{2.02B_o}{p^2-1}\le\frac{B_o}{80}
       \qquad(p\ge13).                                    \tag{OR.5}
\]
There is no lower approximation assumption on $d_y$.

The actual hypotheses give $r,k\ge2$, $Y\ge r$, $N\ge1$, $c\ge1$,
$N_1\le k$, $\alpha_*\le1+rkQ$, and
$a\ge Ck^3r^2NQ^5cY$. The frame span gives $B_o\ge a/Q$. Therefore
\[
 m_*\le8rk^2Qc\le B_o/100,\qquad
 s=\lfloor B_o/(Yk^2)\rfloor\le B_o/8
 \quad(C\ge100).
\]
The prescribed minimum $B_o/2-s$ exceeds $B_x+m_*$.
Division thus leaves exactly the zero-gap plane/normal conditions
\texttt{MGen0}.

\paragraph{OR6: passing from the entire constrained space.}
An element of \texttt{XX'} with fixed whole-letter degrees is a finite
integral sum of the plane/normal products just considered. To see this,
project its defining product span to those degrees: the generic
substitutions preserve each letter degree, so every component retains all
minimum/maximum inequalities. The final order-support intersection can
be forgotten for this inclusion. After division, IP6 and IP11 give
\[
 \xi_{ij}=Dv_j\,\xi'_{ij},\qquad
 \phi_j(\xi'_{ij})\in\mathcal S,
\]
where $\mathcal S$ is the algebra of finite sums of key symbols with their
prescribed normal orders.

\paragraph{OR7: the exact extra multiplier.}
Use $D_s,\Delta_j,R_0$ for the already substituted top determinants,
frame determinants and $R_0=\prod_j\Delta_j^{\alpha_j}$.
Set
\[
 {\cal F}_c=R_0^c\prod_j\Delta_j^c\prod_{s\in TT}D_s^{(2Q+3)c},
 \qquad {\cal D}=\prod_{s\in TT}D_s^{kQ}.
\]
Entrywise Frobenius gives, inside the integral tangential algebra,
\[
 \phi_j(D_x)=D_{\bar x}^{p^{z_j-x}}+pE_x,\qquad
 \phi_j(D_{{\rm old},j})=\Delta_{j-1}^{p^{z_j-z_{j-1}}}+pE_j.
\]
The binomial theorem gives
$(D^P+pE)^m\in D^{m-c}\mathcal T+p^c\mathcal T$ when $P\ge1,m\ge c$.
The late intervals partition one full period, so each $s\in TT$ occurs
once. Combining the factors, and absorbing
$\prod_j\Delta_j^{c(\alpha_*-\alpha_j)}$, proves the literal
\texttt{RestPh} identity
\[
                    \prod_j\phi_j(Dv_j)={\cal F}_c{\cal D}H+p^cY,
 \qquad H\in\mathcal T,\quad Y\text{ integral}.             \tag{OR.6}
\]

\paragraph{OR8: the outgoing lift and its precision.}
For the actual fine label $X^\mu$, put
$G=\sum_i\prod_j\phi_j(\xi'_{ij})$. Then
\[
 \sigma=\operatorname{monPh}
       =(X^\mu{\cal D}GH){\cal F}_c+p^cX^\mu GY.            \tag{OR.7}
\]
Here $G\in\mathcal S$ is integral, and SY4 puts
$X^\mu{\cal D}$ in $\mathcal S$. Take the normal-degree-$o$ component.
SY5's finite correction applied to the first term gives an integral
$F_o$ such that
\[
 \Phi_gF_o\equiv w^o\sigma\pmod{w^{o+c}},\qquad
 \Phi_{g'}F_o\equiv0\pmod{w^{o+c}}
 \quad(\Pi(g')\ne\Pi(g)).                                 \tag{OR.8}
\]
The second term contributes $w^op^c=w^{o+2c}$, so it respects this
precision. Projection preserves any prescribed position multidegree.
These are all clauses of the actual \texttt{Realisable} conclusion.
The result closes the outgoing one-sided application under the supplied
Type-II hypotheses; it asserts no global supply of those hypotheses.
\noindent\textit{Primary source route.}\cite{outgoing-actual}


\section{Matching the actual cloud data and denominator budgets}
\label{sec:cloud-parameters}

\status{The numerical data for the actual signed-chain test are now
constructed from the Type-II hypotheses. The floor in the slack is retained.
The outgoing realization used at the end is the separate OR argument;
ordinary source-symbol realization alone would not suffice.}

\paragraph{CP1: one set of budgets, with the correct floor.}
Write $k=m+1$, $N=|TT|$, $Q=p^{G'}$, and
\[
 L=Nk^2\log_p a,\qquad
 \ell=rNc_0+2rn'c'',\qquad c_P=c_*+\ell.
\]
Besides the actual orbit, weight and frame conditions, the hypotheses give
\[
\begin{gathered}
 r\ge2,\quad r\le Y,\quad N_1+3r\le k,\quad
 256k^2NY^2\le p^{L_1-1},\\
 CL\le D_p,\quad 8(D_p+1)\le c_*,\quad
 Ck^3r^2NQ^5c_PY\le a .
\end{gathered}                                                     \tag{CP.1}
\]
The last inequality, with $C\ge100$, gives
$a\ge\max(2,p,k,4N,r)$ and $L\ge1$; also $c_P>8CL$.
For $B>0$ put $u=B/(Yk^2)$ and $s(B)=\lfloor u\rfloor$.
If $B\ge a/Q^2$, then
\[
 u\ge Ckr^2NQ^3c_P\ge2,\qquad s(B)\ge u/2.
\]
The gap budget in (CP.1) gives
$p^{L_1+1}\ge48rk^2Y$, whence
\[
 3rp^{-L_1-1}B\le u/16\le s(B)/8.                     \tag{CP.2}
\]
If $B\ge a/Q$, then $u\ge Ckr^2NQ^4c_P\ge100r^2c_P$.
The last target is an integer, so
\[
                         s(B)\ge100r^2c_P.            \tag{CP.3}
\]
This proves directly the two numerical comparisons used downstream;
no rounding term is omitted.

\paragraph{CP2: the literal chain and outgoing margins.}
Let $z_j=f_{j,r-1}$, $D_j=z_j-f_{j,0}$,
$g_j=z_{j+1}-z_j$, and $P_j=p^{g_j}$.
For a cut in segment $j+1$, its oldest normalized degree $B$ satisfies
\[
 B\ge a/p^{D_j+g_j}\ge a/Q.
\]
The periodic wrap uses a common translation by $n'$ and leaves
$D_j+g_j$ unchanged. The stronger frame gap places every intermediate
letter of segment $j$ more than $L_2$ before $f_{j,0}$.
For its digit-block plane degree $d_x$, the exact degree identity gives
$p^{z_j-x}d_x\le2a$. Summing over distinct exponents and using $p\ge3$,
\[
 d_{LL,j}\le3ap^{-D_j-L_1-1},\qquad
 r d_{LL,j}/P_j\le s(B)/8.
\]
Since $26r\le25r^2$ for $r\ge2$, (CP.3) now proves
\[
 r d_{LL,j}/P_j+26rc_P\le3s(B)/8.                      \tag{CP.4}
\]
This is exactly the source's \texttt{SignedChainData.Hyp 26}.
For every non-old-frame cut the same oldest-degree bound gives
$s(B)\ge100r^2c_P\ge128c_P$. Together with the prescribed selected
degrees and zero degrees off $TT$, this is the complete
\texttt{MFH} condition required by DT10.

\paragraph{CP3: all logarithmic costs come from actual counts.}
The late letters in a single period inject into $TT$. Thus there are
at most $N$ true letters, at most $N$ chains, at most $N$ counted
links, and at most $2N$ normal projector cuts. Their block degrees are
at most $2a$. Put $h=\log_p a$. Since $a\ge p,4N,r,2$,
\[
 1+\log_p(d_x+r)\le4h,\qquad
 1+\log_p(2\sum_xd_x+r)\le4h.
\]
Consequently the literal definitions give
\[
 D_m\le4L,\quad D_l\le4L,\quad
 0\le k_p\le4L,\quad M_D:=D_m+D_l+3k_p\le20L.         \tag{CP.5}
\]
Also $\kappa_{\rm fl}=\sum_{t\in TT}k\operatorname{ilog}_p(d_t)
\le2Nkh\le20L$.
For each normal cut,
$2\sum_{x\in S}d_x+K-1\le4Na+k\le a^3$ and
$\binom K2\le k^2$. Hence its projector cost sums to $D_{\rm no}\le6L$.
Legendre's formula bounds the two multinomial deficits by
$3Nrh$ and $3NKh$. Their halves plus $D_{\rm no}$ are at most $9L$.
These are the actual \texttt{NumStep 20} and \texttt{SizeAt 9}
budgets. In particular $4M_D<c_P$.

\paragraph{CP4: the signed-chain object actually exists.}
Choose a unit coefficient of the primitive input polynomial and take
its digit class and plane/normal block degrees, as in DT8. Enumerate
the intermediate letters followed by the current frame. At a cut
write $B_u$ for the oldest degree, $s_u=s(B_u)$ and
\[
 T_u=\max\bigl(B_u/2,\ t_{1,u},\ t_{2,u}\bigr),
\]
where $t_{1,u}=a_{u+1}+m_*+s_u$ at an even letter and $B_u/2$
otherwise; $t_{2,u}=a_{u-1}+m_*+s_u$ after an odd letter and $B_u/2$
otherwise. Here $a_u$ is the first coordinate of the normal joint
highest weight. Use all old-frame carries and the first $i$
intermediate letters as chain $i$. Set its threshold to
$T_i-s_i/8$, its slack to $3s_i/8$, and each required link threshold
to $T_u-3s_u/8$.

These chains are strictly increasing. Link thresholds are positive:
$3s_u/8<B_u/2\le T_u$. To avoid any circular choice, first use zero
normal weights, select the joint projector by IW/DT, and then update
the thresholds to its actual weights. The projector depends on the
letter sets and degrees, not on these thresholds. DT supplies the
unit digit factor and its finite binomial approximations. The result
is literally \texttt{mkDt}, with every \texttt{Tied} identity satisfied.
IN/WT supply its incoming condition for the whole permitted word
span, since $C\ge100\ge16$. Thus the PG test applies to this datum.

\paragraph{CP5: the value bound is now unconditional on extra local budgets.}
With $f=\prod_td_t!$, DT's coefficient selection gives
$v_p(k_0)\le v_p(f)-n_0+9L$.
PG, with $(C_H,C_E,C_6,C_{18})=(26,2,24,60)$, gives a main pairing
of valuation at most $v_p(f)+29L$. The retained approximation error
has valuation at least $c_P-20L+v_p(f)-\kappa_{\rm fl}$.
Since $c_P>800L>69L$, its valuation is strictly higher than the main
term's. The normalized integral pairing therefore divides
\[
 p^{D_X},\qquad D_X=\lceil49L\rceil_+\le50L.           \tag{CP.6}
\]
The PG simplification $R=0$ is allowed, but does not remove the normal
digit-approximation error. DT10 now applies with the proved MFH
margins and decomposes the test into the actual permitted one-sided
symbols. This supplies \texttt{Proposition5At} and \texttt{NSStep}.

\paragraph{CP6: a uniform height, with constant $10$.}
By RK and the compatibility in EA, the initial symbols have normal
degree at most
\[
 X=10N\left(1088\min(N_1,2N)^2\log(a)\frac{a}{pk}
              +ap^{-L_1}+kQ\right).
\]
The fixed localizer multiplier has normal degree zero. Conjugating
an elementary operator by the chart substitution gives a linear
combination of operators which replace one variable by one variable;
each increases normal degree by at most one. The permitted word uses
at most $c_0$ operators at each of the $N$ positions. Hence
\[
 t_F=\lfloor X\rfloor_++Nc_0\le t_0(C,10)+Nc_0
\]
for $C\ge1088$. Every symbol component with normal degree $>t_F$
pairs to zero with the input, so discard it. The remaining components
have a common height bound independent of the primitive input.

\paragraph{CP7: realization and the normalized dual pairing.}
Apply OR to these one-sided symbols on the dual edge chart. Its
integral lift $F_o$ satisfies
$\Phi F_o=w^o\sigma+ w^{o+c_P}E$ and has order at least $o+c_P$
at every other plane. The opposite-edge Fock identification of IW
sends $w^{-o}F_o$ to an integral edge dual test after normalization by
$p^{\kappa_{\rm fl}}/f$. Legendre's inequality
$f\mid p^{\kappa_{\rm fl}}\alpha!$ makes its error integral.
Summing the tests gives exactly
\[
 \beta_{\rm fl}(\Theta F_1,T)=y+w^{c_P}z,\quad z\in O,
 \qquad yb=p^{D_X}.
\]
The lift has target height $\le t_F$ and other-edge order $c_P$.
Because $D_X\le D_p$ and $c_P\ge8(D_p+1)>2D_X$,
the value $y+w^{c_P}z$ still divides $p^{D_X}$: its product with $b$
is $p^{D_X}(1+w^{c_P-2D_X}bz)$, whose parenthesis is a unit.
This proves the actual \texttt{FlatAt} assertion, including its
integral error and every other-edge condition.

\paragraph{CP8: full cloud test and quantifier order.}
CB and the primitive-array argument CL3 apply to every coefficient
array having a unit entry. The normalized contraction of CL/IW has
cost $\kappa\le4L$ and shifts height and other-edge order by $\ell$.
If its scalar is $c$, the exact identity is
$c\beta(U,\Lambda T)=p^\kappa\beta_{\rm fl}(\Theta F_1,T)$.
The scalar need not be a unit. If the latter flat value times $b'$
equals $p^{D_X}$, then
\[
 \beta(U,\Lambda T)(cb')=p^{\kappa+D_X},\qquad
 D_1:=\kappa+D_X\le54L.                               \tag{CP.7}
\]
Other-edge order becomes $c_P-\ell=c_*$ and target height is at most
$t_0(C,10)+Nc_0+\ell$. Thus the exact cloud interface has
$(C_{t_0},C_{D_1})=(10,54)$ for every $C\ge4129$, since OR's actual
realization threshold is $100$. All constants are chosen before the local
field, edge, basis and coefficient array; the test itself may depend
on the array. The cloud result is a local theorem under the stated
Type-II hypotheses, not the later global supply of those hypotheses.
\noindent\textit{Primary source route.}\cite{cloud-actual}


\section{From selected symbols to controlled edges}\label{sec:edge-assembly}

\status{The assembly and its numerical comparisons are proved here. The
uniform cloud theorem is an explicit input, supplied separately by the
cloud-compatibility argument.}

Write $k=r+K$, $N=|TT|$, $N_1=\texttt{orbitWindowCount}$,
$Q=p^{G'}$, $q=p^{n'}$, and
\[
 \ell=rNc_0+2rn'c'',\qquad c'=c_*+\ell,\qquad
 X_C=\frac{CN_1^2\log(a)\,a}{pk}+ap^{-L_1}+kQ.
 \tag{EA.1}
\]
Use the literal \texttt{Hyps}, \texttt{OrbitHyp}, \texttt{WeightHyp} and
coefficient setting of CL. In particular $w^2=p$,
\[
\begin{gathered}
 N_1^2\le rN,\quad k\ge N_1+3r,\quad Y\ge r\ge2,
 \quad CN_1^2\log a\le k,\\
 256k^2NY^2\le p^{L_1-1},\quad CNk^2\log_p a\le D_p,
 \quad8(D_p+1)\le c_*,\\
 Ck^3r^2NQ^5c'Y\le a.
\end{gathered}\tag{EA.2}
\]
The other orbit and weight clauses remain in force.

\begin{theorem}[Uniform edge assembly]
Fix RK's $C_{\rm La}=C_k=1088$, $C_{\rm rk}=1$, $C_{\rm deg}=10$.
Suppose the actual \texttt{CloudStep}$(C,A,B,1088)$ holds for every
$C\ge C_{\rm cloud}$, with fixed $A,B>0$. Then the source's
\texttt{ConclusionE} holds for
\[
 C\ge C_0:=\max(4129,C_{\rm cloud},A,B).
 \tag{EA.3}
\]
Concretely, for every integer
$t\ge CrN(X_C+c_0)+2rn'c''$, it gives the literal
\texttt{LPprime} with denominator exponent $D_p$, error depth $c_*$ and
rank loss $CNr^2k/Y$, at every target lattice.
\end{theorem}

\begin{proof}
\emph{Match the selected-symbol data.}
FG gives the full \texttt{FramesOK}. Segment lengths are $<G'\le n'/3$,
so their positions inject modulo $n'$; with $N_s=\min(N_1,2N)$ this gives
$|\mathrm{letters}_j|\le N_s$ and $p^{\mathrm{dist}_j(x)}\le Q$.
The remaining RK inequalities follow from $N_s\le N_1,2N$ and (EA.2).
Here is the substantial size check. The actual reserve and degree are
\[
\begin{aligned}
 E_0&=15c'(N_1+rk)Q^2,\\
 m_*&=(N_1+1)\bigl((2Q+3)c'+kQ+c'\alpha_{\max}+3c'\bigr),
 \qquad\alpha_{\max}\le1+rkQ.
\end{aligned}
\]
Put $Z=c'rkQ$. Then $E_0+m_*+1\le Z(30Q+11(N_1+1)+1)$.
Since $N_sN_1\le rN$ and $N_s\le2N$,
\[
\begin{split}
 N_s(30Q+11(N_1+1)+1)&\le60NQ+11rN+24N\le95rNQ,\\
 32pk^2N_sQ^2(E_0+m_*+1)&\le3040k^3r^2NQ^5c'.
\end{split}\tag{EA.4}
\]
Multiplying by $Y$ verifies RK's final hypothesis. The source threshold
$3041+|C_k|=4129$ suffices for every comparison. \textsf{[EA1--EA2]}

For $C\ge30$, SY/IP's actual determinant multiplier has degree
$E'(s)\le E_0\le a\le d_s$ on $TT$ and zero elsewhere. Thus its finite
correction gives \texttt{SourcesRealisable} with these same parameters.
RK supplies an ambient direct summand with an actual-generator basis
$B_1,\ldots,B_b$. It is finite because its row multidegree is fixed.
Let $A'=\max(A,10)$ and
\[
 T_0=\left\lfloor10N
 \left(\frac{1088N_s^2\log(a)\,a}{pk}+ap^{-L_1}+kQ\right)\right\rfloor
 \le A'NX_C.
 \tag{EA.5}
\]
The report verifies the determinant budget and all natural subtractions.
\textsf{[EA2--EA3]}

\emph{Realise the flat family in the actual edge.}
Choose a basis $H$ of its source. Neighboring targets have edge bases
$Hg^{-T}$ with $g\in\operatorname{GL}_k(W)$; their rational planes are
the spans of the first $r$ rows of $\bar g^{-1}$. This matches the source's
\texttt{ratPlane} convention. The standard complete unramified DVR/Witt
identification respects Frobenius. Dividing a level-$o$ integral raw lift
by $w^o$, then summing $0\le o\le T_0$, gives flat vectors $V_i,S_i$ with
\[
\begin{gathered}
 V_i\in w^{-T_0}H_{\rm flat}(L^-),\quad S_i\in R_e(\mathrm{flat}),
 \quad V_i-S_i\in w^{c'}R_e(\mathrm{flat}),\\
 V_i\in w^{c'}R_{e'}(\mathrm{flat})
 \quad(e'\ne e\text{ with the same source}).
\end{gathered}\tag{EA.6}
\]
This uses the cloud's exact flat-coordinate identification; it does not
assert that rational preimages are integral. \textsf{[EA3--EA4]}

\emph{Multiply and preserve every primitive test.}
The cloud supplies $n_2$ complementary vectors
$h_j\in R_e(\mathrm{rest})\cap w^{-o_S}H_{\rm rest}(L^-)$,
with $o_S\le\ell$. CB's multiplication and EL's inclusions give
$f_{ij}=V_ih_j$ source height $T_0+o_S$, while its errors at other outgoing
edges have order $c'-o_S\ge c_*$. Its own-edge difference from $S_ih_j$
has order $c'$. The pairing is exactly
$\beta(F,G)=G(p^{A_\lambda}F)$, with possibly negative $A_\lambda$.

For any primitive coefficient array the cloud supplies a dual test with
integral value $a_0\mid p^{D_1}$. We have
$D_1\le BNk^2\log_p a\le D_p$. The realised value differs by
$w^{c'}z\in p^{D_1+1}\mathcal O$. If $a_0b_0=p^{D_1}$, then
\[
 a_0+p^{D_1+1}z'=a_0(1+pb_0z'),
 \tag{EA.7}
\]
whose parenthesis is a unit. The valuation bound therefore survives for
every primitive array. \textsf{[EA5]}

\emph{Count and absorb the costs.}
Put $x=Nr^2k/Y$, $u=(1-kc_0/a)^N$ and
$\Pi=\prod_{s\in TT}\binom{d_s+k-1}{k-1}$. RK and the complementary
count give $b\ge q^{-rK}(1-x)\Pi$ and $u\dim V\le\Pi n_2$.
Since the size bound implies $Nkc_0Y\le a$, we have $0\le u\le1$.
Hence
\[
 bn_2\ge q^{-rK}(1-x)u\dim V.
 \tag{EA.8}
\]
When $1-x<0$ this is simply nonnegativity; no inequality is multiplied by
a negative factor. The source and target heights are at most
$A'NX_C+\ell$ and $A'NX_C+Nc_0+\ell$. Both are $\le t$, since
$A'\le Cr$ and $r+1\le Cr$. Bernoulli gives $u\ge1-x$; splitting
$x\le1$ and $x>1$ yields $(1-x)u\ge1-2x\ge1-Cx$.
\textsf{[EA6--EA7]}

The kernels of $(u,v)\mapsto u+Mv$ for
$M\in\operatorname{Mat}_{r,K}(\mathbb F_q)$ give $q^{rK}$ different
edges into $W^k$. Summing (EA.8) yields at least $(1-Cx)\dim V$ vectors.
For a primitive combination across all edges, use a test from an edge
with a unit coefficient. All other edges contribute only
$p^{D_p+1}\mathcal O$, so (EA.7) applies again.

Let $J\subset\mathcal O^M$ be the image of all target tests under their
pairings with these $M$ vectors. Smith normal form and primitive detection
give $p^{D_p}\mathcal O^M\subset J$. Choosing preimages of
$p^{D_p}e_i$ gives exactly \texttt{GoodRankK}, with identity coefficient
matrix. Tensor polynomial covariance and the contragredient dual action
transport this witness to every target lattice, preserving
$G(p^{A_\lambda}F)$. This proves the required \texttt{LPprime}, including
its $4(D_p+1)\le c_*$ condition. All thresholds were fixed before the local
data, proving \texttt{ConclusionE}. \textsf{[EA8--EA9]}
\end{proof}

The \texttt{RestCloudX}, \texttt{RestCloudY} and split-cloud variants feed
the same implication once their actual uniform cloud premise is established.
The separate Type-I construction and the arithmetic windows remain outside
this edge-assembly theorem.

\begin{corollary}[The actual local Type-II theorem]\label{cor:typeII}
Under the literal Type-II coefficient, orbit and weight hypotheses,
with the numerical inequalities (EA.2), every $C\ge4129$ gives the
source's full \texttt{ConclusionE}. In its \texttt{LPprime} output,
the rank loss is at most $CNr^2k/Y$, the denominator exponent is $D_p$,
and other-edge errors have order $c_*$. One may use every integer height
\[
 t\ge CrN\left(\frac{CN_1^2\log(a)\,a}{pk}
                  +ap^{-L_1}+kQ+c_0\right)+2rn'c''.
\]
\end{corollary}
\begin{proof}
CP8 proves the uniform \texttt{CloudStep}$(C,10,54,1088)$ for
$C\ge4129$, using OR's proved threshold $100$. Substitute
$(C_{\rm cloud},A,B)=(4129,10,54)$ in the edge-assembly theorem.
Every threshold is fixed before choosing the local field, orbit, weights,
target lattice or primitive coefficient array. This discharges the
cloud premise within this reconstruction. \textsf{[CP8, EA9]}
\end{proof}
The result constructs the local objects from valid input data. It leaves
the global selection and abundance of such data to the arithmetic part
of the RH/GRH argument.
\noindent\textit{Primary source route.}\cite{edges-actual}


\section{The separate rank-one construction: a polynomial averaging mechanism}
\label{sec:type-one}

\status{The Type-I block lattices, their primitive-test interpretation, the
actual gamma degree and leading-form identities, logarithmic integrality,
and the polynomial-matrix averaging lemma are proved below.
Sections~\ref{sec:graded-gram}--\ref{sec:paired-lifts} now supply its
actual graded Gram, simultaneous-partner and cleaned-lift inputs.
The later RU/RD/HM/NP/TC sections supply the remaining local inputs and prove the full rank-one theorem in Section~\ref{sec:type-one-closure}.
The earlier $r\geq2$ segment theorem is not used for this branch.}

\paragraph{TI1--TI2: the actual polynomials and integral lattices.}
The edge rank is one, but its segment polynomials have $r=m$ coordinates,
where the ambient dimension is $m+1$.  A tie joins adjacent vector and
covector rows.  With scalar variables $Z_i,S_i$ and an auxiliary row $\xi$,
the source uses
\[
 E_i=Z_iS_i+\sum_c x_{a_i,c}x_{a_i+1,c},\qquad
 E_+=Z_+S_++\sum_c x_{u,c}\xi_c,
\]
and the integral span of
\[
 \Bigl(\prod_iE_i^{j_i}Z_i^{z_i}\Bigr)E_+^{j_+}Z_+^{z_+}f(x)g(\xi).
 \tag{TI1.1}
\]
Here the levels fit both tied degrees, $f$ belongs to the actual integral
prefix space of residual degrees, and $\deg g=B-j_+$.
The block lattice is its rational span intersected with
\[
 v_p(F_\alpha)+\tfrac12v_p(\alpha!)+\tfrac12v_p(s_b)\geq0,
 \qquad s_b=\prod_{i<\sigma}(A_i-w_i)!\,\mathrm{sg}_i!.
 \tag{TI1.2}
\]
This is a full finite free lattice: it is a rational subspace intersected
with a finite free weighted-monomial lattice.  Its form $s_b\,\mathrm{Fock}$
is integral term by term.  Smith bases show that its largest saturated
subspace in which every primitive vector has a partner of valuation at
most $D$ has rank exactly $\#\{i:v_p(d_i)\leq D\}$.
Indeed the indicated Smith vectors work; a larger saturated subspace
contains a primitive vector with all those coordinates zero.

\paragraph{TI3: the contraction polynomial.}
Let a signed chain have $N+1$ rows of degrees $d_i$, total $S$, and last
vector row $w$ of degree $o$. Fix $\rho$, a nonnegative integer shift $z$,
and append a covector $y$; all row degrees are nonnegative integers.
The source substitution $\mathcal F$ is
\[
\begin{array}{lll}
\text{vector }x:&x_\rho\mapsto\sum_cx_cy_c,&x_c\mapsto y_\rho x_c\ (c\ne\rho),\\
\text{covector }x:&x_\rho\mapsto y_\rho x_\rho,
 &x_c\mapsto y_\rho x_c-x_\rho y_c\ (c\ne\rho).
\end{array}
\]
Put $\operatorname{fr}_n(f)=y_\rho^{n-S}\mathcal F(f)$ for $n\geq S$,
$E=\sum_cw_cy_c$, $F=\sum_c\partial_{w_c}\partial_{y_c}$ and
$A=o+n+r+z-1$.  With $B$ denoting the bilinear Fock form, set
\[
 \Phi_{o,n,z}=\sum_{j=0}^o\frac{(-1)^jE^jF^j}{j!^2\binom Aj},\qquad
 \Gamma_{f,g}(n)=
 \frac{\binom Ao B(\Phi_{o,n,z}\operatorname{fr}_n f,
                         \operatorname{fr}_n g)}
 {(n-S)!S!\prod_i d_i!}.
 \tag{TI3.1}
\]
If nonzero $f,g$ have respective orders $e,e'$ in $w_\rho$, then
\[
 \deg\Gamma_{f,g}\leq S+o-\max(e,e'),\qquad
 [T^{S+o-e}]\Gamma_{f,g}
 =\frac{e!\binom{z+e}e B(\operatorname{in}_ef,\operatorname{in}_eg)}
 {o!S!\prod_i d_i!}\quad(e=e').
 \tag{TI3.2}
\]
Here $\operatorname{in}_e$ retains precisely coordinate degree $e$.

For a short proof, $[F,E]=r+\mathrm{Euler}_w+\mathrm{Euler}_y$ gives
the rational harmonic decomposition.  A finite binomial identity shows
that $\Phi$ acts on $E^l\ker F$ by $\binom{z+l}l/\binom Al$; hence
\[
 \Phi_{o,n,z}E^e=
 \frac{(z+1)^{\overline e}}{A^{\underline e}}
 E^e\Phi_{o-e,n-e,z+e}.
 \tag{TI3.3}
\]
Write $f=w_\rho^eh$, $q=\mathcal F(h)$ and $v=\mathcal F(g)$.
With $m=n-S$, the Leibniz formula
\[
 F^j(y_\rho^mq)=\sum_{u=0}^j\binom ju m^{\underline u}
 y_\rho^{m-u}\partial_{w_\rho}^uF^{j-u}q
\]
shows that $B(E^{e+j}F^j(y_\rho^mq),y_\rho^mv)/m!$ has degree at
most $S+j$.  Its leading coefficient is
$B(w_\rho^{e+j}\partial_{w_\rho}^jh,g)$.
The remaining coefficient in (TI3.3) is a polynomial of degree $o-e-j$.
At the top degree its sum reduces to
\[
 \sum_j\frac{(-w_\rho)^j}{j!}\partial_{w_\rho}^jh
 =h\big|_{w_\rho=0}.
\]
This proves both formulas (TI3.2), with symmetry supplying $\max(e,e')$.

\paragraph{TI4: integrality costs only logarithmically many powers of $p$.}
For integral $f,g$, let $L_0=\lfloor\log_p(S+o+r)\rfloor$ and
$L_3=\lfloor\log_p(3S)\rfloor$ when $S>0$. Then $p^K\Gamma$ has
integral binomial coefficients provided
\[
 K\geq(2r+1)L_0+(N+2)(r-1)L_3.
 \tag{TI4.1}
\]
To prove it, a multinomial coefficient with $r$ parts totaling $R$ has
valuation at most $(r-1)\lfloor\log_pR\rfloor$, by Legendre's formula.
Consequently $F^j\operatorname{fr}_n(f)/j!$ has coefficient valuations
at least $v_p(j!)-(r-1)L_0$.
Use self-adjointness in (TI3.1) and
$\binom Ao/\binom Aj=\binom{A-j}{o-j}/\binom oj$.
After accounting for the factorials in every row, the remaining factorial
valuation is exactly
\[
 v_p\binom{n-j}{S-j}-2v_p\binom oj-v_p\binom Sj.
\]
Evaluate only at $n=S,\ldots,2S+o\leq3S$.
This proves integrality with (TI4.1) at degree plus one consecutive integers;
Newton differences give integral coefficients in $\binom{T-S}j$, and
integer translation gives them in $\binom Tj$.
The case $S=0$ is constant.  This also proves the source's
$\texttt{GamInt}$ with $K=10(N+3)r^2(L_0+1)$.

\paragraph{TI5--TI6: why averaging removes most bad factors.}
Let $M(T)$ be an $n\times n$ matrix of integral binomial polynomials of
degree at most $d$, with nonzero determinant of degree $\delta$ and
leading coefficient of valuation $\lambda$.
On an integer progression $t_j=t_0+uj$, $p\nmid u$, $0\leq j<Y$, assume
$h\in\mathbb N$, $1\leq Y\leq p^h$ and $d\leq p^h$.
If $\beta(t)$ counts Smith factors of valuation at
least $E>0$, then
\[
 \lambda+\frac\delta{p-1}\leq nL
 \quad\Longrightarrow\quad
 \frac1Y\sum_j\beta(t_j)
 \leq n\left(\frac dY+\frac{L+dh/Y}{E}\right).
 \tag{TI5.1}
\]
Indeed $M(t+p^hX)$ is coefficientwise integral.  After Smith reduction at
$t$, its determinant coefficient of degree $k<\beta(t)$ has valuation
at least $E(\beta(t)-k)$.
Factoring the determinant over a splitting field gives
\[
 E\beta(t)\leq\lambda+\sum_\zeta\min\{v_p(t-\zeta),h\}
       +E\#\{\zeta:v_p(t-\zeta)\geq h\}.
\]
For each root the first sum averages to at most $1/(p-1)+h/Y$,
by counting residue classes modulo $p^i$; the small discs are disjoint,
so its last count sums to at most one. This proves (TI5.1).

For the actual matrix $p^{\lceil\kappa\rceil}\Gamma$, suppose the basis is
adapted to coordinate order and each order-$e$ initial Gram block, normalized
by $\prod_i d_i!$, has determinant valuation at most $\kappa_4$ times its rank.
Equal-order blocks alone supply the determinant's top coefficient in (TI3.2).
If $S,z\leq2a$, $r\geq2$, $\ell=\log_p(8a+r)+1$, and
$\kappa\geq(N+3)r^2\ell$, (TI4.1) supplies integral binomial coefficients.
Legendre's formula and (TI3.2) give
\[
 \frac1Y\sum_j\beta(n_j)
 \leq n\left\{\frac{2S}Y+
 \frac{\lceil\kappa\rceil+\kappa_4+3\ell+2S\ell/Y}{E}\right\},
 \tag{TI6.1}
\]
for $n_j=s+uj$, $Y,u\geq1$, $p\nmid u$, $S\leq s$ and $n_{Y-1}\leq2a$.
This keeps the integral normalization and the rounding term explicit.
GG below now proves the retained graded Gram premise; TP gives the
sharper source normalization and removes the rounding loss in its exact
one-tie estimate.

\paragraph{TI7--TI9: the bridge to the final exceptional function.}
The source next weights these estimates by varying highest-weight
multiplicities, supplies simultaneous integral partners, and constructs
paired maps into the cleaned rank-one edge lattices.
Those representation and lifting inputs are now proved in GG, TP
and RL below. SW proves the form and weighted-sum mechanisms, retaining
the actual coarse rank inequality and explicitly stated numerical inputs.
Once their dimension-weighted inequality and numerical density bounds are
supplied, summing simultaneous degree translations gives
$\operatorname{mean}k_F\leq2\operatorname{mean}k_x$.
The actual final function has a lower-shift term (set to one at the boundary),
two upper-shift terms and $12Nkc/b'$.  Boundary and last-term costs at most
$1/Y$ each therefore give
\[
 \operatorname{mean}\kappa
 \leq6(C_A+C_s)NN_1/Y+2/Y
 \leq[6(C_A+C_s)+8]NN_1/Y.
 \tag{TI8.1}
\]
This elementary transfer is proved in the companion. RU/RD now
prove the actual rank and density estimates; HM identifies the harmonic
objects and cloud maps; NP supplies every numerical field. TC below
assembles them into the complete local conclusion.
\noindent\textit{Primary source route.}\cite{type-one-actual}


\section{The actual graded Gram bound for Type I}
\label{sec:graded-gram}

\status{This supplies the previously retained graded initial-Gram premise
of TI6, together with the full-word \texttt{HypP}/\texttt{HypW} inputs.
The argument controls the actual coefficient-image lattice.
It uses the proved IW averaging theorem and retains no additional local
integral representation-theory premise.}

\paragraph{GG1: the lattice and its normalization.}
Let a signed word have $K+1$ letters of degrees $d_i$, total degree
$D=\sum_i d_i$, and final vector letter $W$. Work over $\mathbb Q_p$,
with integral coefficient lattice $\mathbb Z_p[X]$.
Assume $r\ge2$, put $\rho=r-1$, and $z=W_\rho$.
For finitely many prefix bounds $(J_\ell,M_\ell)$, let $Y$ be the
multihomogeneous space whose types at prefix $J_\ell$ have first
highest weight at most $M_\ell$; $M_\ell=\infty$ is allowed.
Let $M_w$ be its integral lattice of full weight $w$, highest for
$GL_{\{0,\ldots,\rho-1\}}$ acting on all letters. The literal
\texttt{grLat} is
\[
 \begin{gathered}
 {\cal G}_e=\operatorname{in}_e(M_w\cap z^e\mathbb Q_p[X]),\qquad
 \operatorname{in}_e\text{ retains }z^e,\\
 B_d(X^\alpha,X^\beta)=
 \delta_{\alpha\beta}\frac{\prod_{i,a}\alpha_{ia}!}{\prod_i d_i!}.
 \end{gathered}
                                                               \tag{GG.1}
\]
The diagonal entries are reciprocals of integer multinomials.

\paragraph{GG2: all required integral projectors.}
Put
\[
 e_1=\binom r2\operatorname{Nat.log}_p(D+r-1),\qquad
 D_4=\binom r2\operatorname{Nat.log}_p(4D+r-1).                  \tag{GG.2}
\]
Any two prefix weights differ at one coordinate by at most $D$:
each lies between minus the covector degree and the vector degree.
Elementary substitutions show that the coefficient lattice is
weight-graded and stable under root divided powers. IW4 therefore
makes every prefix type projector integral after multiplication by
$p^{e_1}$. They are Fock self-adjoint and commute for nested prefixes,
because the smaller projector commutes with the larger diagonal action.

On a fixed full weight $w$, the smaller-group projector onto type
$w|_{\{0,\ldots,\rho-1\}}$ is precisely the highest-vector projector:
the top-weight space consists of highest vectors, and every other
type contributes lowering images. This proves \texttt{HypE} with the
same cost.

For an isotypic lattice $M$, let $\Delta$ be the divided-power span
of its integral highest vectors. Its integral top-weight projector
$C$ satisfies, by IW2 on $\operatorname{Hom}(M,\Delta)$,
$I(C)=(\dim U)^{-1}\operatorname{Id}$, where $U$ is one irreducible
factor. Multiplying by the integer $\dim U$ gives
$p^{e_1}M\subset\Delta$. The multiplicity space is untouched.
This proves \texttt{HypW} without integral semisimplicity.

\paragraph{GG3--GG4: the actual restriction-image lattice.}
Allow the final degree to be $o$. Let $X_o$ consist of $z$-free
polynomials satisfying the earlier prefix bounds. Each full-prefix
bound becomes both the earlier-word $GL_r$ bound and the full-word
$GL_{r-1}$ bound. These are exactly \texttt{XSp}.
If all full bounds are at least $D$, restriction lifts trivially.
Otherwise let $c<D$ be their finite minimum and set
\[
 {\cal W}=\bigoplus_oX_o,\qquad
 {\cal L}_{\rm im}=\sum_o
       \{h|_{z=0}:h\in Y_o\cap\mathbb Z_p[X]\}.                 \tag{GG.3}
\]
Restriction lands in $X_o$: extracting $W$-coefficients tensors
with $\operatorname{Sym}^o(V)^*$, whose coordinate weights are
nonpositive, and hence preserves the earlier first-weight bound.

Write $A_{ab}$ for the earlier action,
$L_{ab}=A_{ab}+W_a\partial_b$, and $E_W=\sum W_i\partial_i$.
On the remaining coordinates define
\[
 U_a=W_a(c-E_W)-\sum_bW_bA_{ab},\qquad
 E_{**}=c-E_W,\quad E_{*b}=\partial_b,\quad
 E_{a*}=U_a,\quad E_{ab}=L_{ab}.                               \tag{GG.4}
\]
Direct commutators give $\mathfrak{gl}_r$.
The family $U_a=cW_a-\sum_bW_bL_{ab}$ is covariant.
On coordinate-$0$ weight $c$, $U_0$ vanishes; the other $U_a$
do not raise that weight. Thus it preserves the first-weight cutoff.
Derivatives form a covector family and preserve it as well.
This proves stability directly, without the source's intermediate
rational-surjectivity theorem.

Every nonzero $X_o$ has $o\le c+D_{\rm early}$: a generic coordinate
change makes its $W_0^o$ coefficient nonzero, whose earlier signed
weight is at least $-D_{\rm early}$. The auxiliary weights therefore
lie in $[-D_{\rm early},c+2D_{\rm early}]$, of width at most $4D$.

Divided powers are integral. For $U_a$, put $q=1+tW_a$.
Its exponential is $q^c$ times substitution
\[
 W_j\mapsto W_j/q,\quad
 x_b\mapsto x_b-tW_bx_a/q,\quad
 y_a\mapsto y_a+t\sum_bW_by_b,                                 \tag{GG.5}
\]
on earlier vector and covector letters; other covector coordinates
are fixed. Integer powers of $q$ have integral series coefficients.
The same formula before restriction proves that $ {\cal L}_{\rm im}$
is divided-power stable.

An auxiliary highest vector is killed by every $\partial_b$, so has
$W$-degree zero and lifts to itself. IW4 first isolates each type of
an integral vector in $ {\cal W}$, at cost $D_4$; GG2's highest-vector
generation costs another $D_4$. Their generators lie in
$ {\cal L}_{\rm im}$. Taking the original degree component proves
\[
 p^{2D_4}(X_{d_K}\cap\mathbb Z_p[X])
       \subset\{h|_{z=0}:h\in Y_d\cap\mathbb Z_p[X]\}.           \tag{GG.6}
\]

\paragraph{GG5--GG6: order, saturation and projection.}
Dividing by $z^e$ reduces the final degree and full weight by $e$,
and subtracts $e$ from precisely the full-prefix bounds.
This is an equivalence: after a generic group change, multiplication
by $z^e$ adds $e$ to the largest first-coordinate Laurent exponent.
Project the lift in (GG.6) to its full weight and smaller-group
highest vectors, costing $e_1$. If $V_e=\mathbb Q_p{\cal G}_e$,
\[
 p^b(V_e\cap\mathbb Z_p[X])\subset{\cal G}_e,\qquad
 b=e_1+2D_4.                                                   \tag{GG.7}
\]
This establishes saturation control for the actual image.

For $K\ge1$, first project integrally to $z$-degree zero. The rational
projection onto the reduced highest restriction space then uses
$K+2$ nested type projectors. Those at prefix
lengths $0,1$ are identity or zero. The remaining $K$ cost $Ke_1$;
highest projection costs $e_1$. Strip $z^e$, apply this product and
multiply back. It gives a Fock-orthogonal projection $P$ onto $V_e$
with integral cost $A=(K+1)e_1$. The identity
$\operatorname{Fock}(z^ef,z^eg)=e!\operatorname{Fock}(f,g)$
preserves orthogonality.

\paragraph{GG7--GG8: determinant and the precise segment constant.}
For primitive $u\in{\cal G}_e$, its smallest coefficient valuation
is at most $b$ by (GG.7). Choose its corresponding monomial $m$.
Then $h=p^{A+b}P(m)\in{\cal G}_e$, and (GG.1) gives
$v_pB_d(u,h)\le A+2b$.
Smith form bounds every elementary-divisor valuation by this number.
Thus for any basis of rank $k$,
\[
 v_p\det B_d\le\bigl((K+3)e_1+4D_4\bigr)k
 \le (K+7)\frac{r^2}{2}\ell\,k,                                \tag{GG.8}
\]
when both logarithms in (GG.2) are at most $\ell$.
For $K=0$ the possible graded lattice is a monomial line with Gram
entry $1/\binom{d_0}{e}$, so its valuation is at most zero.

In the actual segment let $\ell=1+\log_p(8a+r)$, $a\ge1$, $p\ge3$.
The literal size hypothesis gives full total at most $8a+r$ and
$2D_i+1\le8a+r$ for each tie-prefix total.
Both required logarithms are therefore at most $\ell$.
Since $K=\operatorname{pos}_i<2s+1$, (GG.8) proves exactly
\[
 \kappa_4=(s+4)r^2\ell.                                       \tag{GG.9}
\]
GG2 also supplies full-word \texttt{HypP}/\texttt{HypW} with the
permitted larger exponent $\lfloor r^2\ell\rfloor$ required by
the partner construction.
\noindent\textit{Primary source route.}\cite{graded-gram-actual}


\section{Compatible tie partners and their exact weighted mean}
\label{sec:tie-partners}

\status{The actual simultaneous-partner construction and the source's
dimension-weighted multiple-tie mean are proved here. This uses GG's
quantitative prefix projections, integral highest-vector generation and
restricted graded Gram estimate. The sharper gamma integrality constant
is proved explicitly. Segment and digit-class transfers and the actual
edge lifts are separate steps.}

\paragraph{TP1--TP3: the exact spaces and the sharper one-tie estimate.}
Work over $\mathcal O=\mathbb Z_p$, $K=\mathbb Q_p$, with $p\ne2$ and
$r\ge2$. A signed word has $N$ letters of fixed degrees. Let $Y$ be its
polynomial space subject to the actual prefix conditions
$\lambda_1\le M_k$ at lengths $J_k$, and put
$Y_{\mathcal O}=Y\cap\mathcal O[X]$. Its form is
$B_L=\mathrm{Fock}/\prod_jd_j!$.
A tie $i$ joins vector $a_i$ to covector $a_i+1$, whose degree is $n_i$;
the positions increase, the pairs are disjoint, and $z_i\in\mathbb N$.
Write $S_i$ for the preceding total degree and require $n_i\ge S_i$.

Before that covector, let $M_{i,w}$ be the space of full torus weight $w$
that is highest for the coordinates other than $\rho=r-1$, with all
cutoffs $J_k\le a_i+1$. TI3's explicit freeze map identifies
$M_{i,w}\cap\mathcal O[X]$ with the entire integral highest-vector lattice
of weight $w-n_i e_\rho$ after adjoining the covector. Indeed, on
$y_\rho\ne0$, its substitution is a common unitriangular change of vectors
and inverse transpose change of covectors. Highest-vector invariance
recovers a polynomial from its specialization $y=e_\rho$, which is the
integral inverse. This also preserves the old prefix conditions.
The order filtration in $x_{a_i,\rho}$ is saturated; lifting bases of its
successive coefficient images gives an integral adapted basis.

For a prefix of $K+1$ old letters, total $S$ and last vector degree $o$,
the exact strengthening of TI4 is
\[
 \begin{gathered}
 p^{c_I}\Gamma_{f,g}(n)=\sum_{j=0}^{S+o}c_j\binom nj,
 \qquad c_j\in\mathcal O,\\
 c_I=((K+4)r+6)\lfloor\log_p(S+o+r)\rfloor.
 \end{gathered}
 \tag{TP.1}
\]
Here $n\ge S$ and $z\in\mathbb N$. A short proof expands the contraction
as $D_\rho+D_\perp$. A paired monomial term becomes an integral-binomial
polynomial times a scalar. The two differentiated coefficients and the
$K+2$ row-factorial groupings lose at most $(K+4)rB$, where
$B=\lfloor\log_p(S+o+r)\rfloor$. The remaining factorials cancel exactly to
\[
 -v_p\binom S{j-a}-2v_p\binom oj-v_p\binom j{a'},
 \qquad 0\le a,a'\le j\le o.
\]
This loses at most $4B$, within the source's $6B$ allowance. The expanded
companion gives every coefficient and the binomial translation.

Let $b$ be the adapted-basis size and suppose the equal-order initial Gram
determinants have valuation at most $\kappa_4$ times their sizes, as GG
proves. For $A(n)=p^{\lceil\kappa\rceil}\Gamma(n)$, TI3 gives degree
$\sum_u(S+o-e_u)$ and leading coefficient
\[
 \prod_e\left[
 \left(\frac{p^{\lceil\kappa\rceil}e!\binom{z+e}e}{o!S!}\right)^{|I_e|}
 \det G_e\right].
\]
Only equal-order permutations attain the top degree. If
$2S,z+o,Y,r\le X$,
$\ell\ge\lfloor\log_pX\rfloor+1$ and
$\kappa\ge((K+4)r+6)\ell$, (TP.1) gives integral binomial coefficients.
Legendre's formula, $\lceil\kappa\rceil<\kappa+1$ and
$v_p\binom{z+e}e\le\ell-1$ give the precise leading-coefficient budget
$b(\kappa+3\ell+\kappa_4)$. TI5, with the integer
$h=\lfloor\log_pX\rfloor+1$, therefore yields
\[
 \underset{0\le j<Y}{\operatorname{Avg}}
 \#\{\text{Smith factors of }A(s+dj):v_p>E\}
 \le b\left[\frac{2S}{Y}
 +\frac{\kappa+3\ell+\kappa_4+2S\ell/Y}{E}\right],
 \tag{TP.2}
\]
for $E\ge1$, $s\ge S$, $d>0$ and $p\nmid d$.

\paragraph{TP4--TP7: one system of integral partners for all ties.}
Use GG's common projection and highest-vector-generation exponent $e_1$.
For one tie, Smith bases give integral highest vectors $h_x,h''_x$ and
scalars $d_x$ such that
$c_*\mathrm{Fock}(\phi_i h_x,h''_y)=\delta_{xy}d_x$ within each weight.
The normalization is
\[
 c_*=
 \frac{p^{\lceil\kappa\rceil}\binom{o+n_i+r+z_i-1}o}
 {(n_i-S)!S!\prod_{j\le a_i}d_j!}.
\]
Row-factorial grouping also makes $c_*\mathrm{Fock}(h_x,h_y)$ integral.
Call $x$ good when $v_p(d_x)\le E$, and set $e=\lfloor E\rfloor$.
Define on highest vectors
\[
 \varpi h_x=1_{x\,\mathrm{good}}h_x,\qquad
 \tau h_y=
 \sum_{\substack{x\,\mathrm{good}\\\lambda_x=\lambda_y}}
 \frac{p^e}{d_x}\,c_*\mathrm{Fock}(h_x,h_y)h''_x.
 \tag{TP.3}
\]
These coefficients are integral. Extend the maps equivariantly over $K$,
then coefficientwise over later variables. The projection onto the earlier
cutoffs costs $(a_i+2)e_1$; the pure-type projection and highest-vector
generation cost another $2e_1$. Thus the resulting $T_i,\pi_i$ satisfy
\[
 \begin{gathered}
 p^{K_i}T_i(Y_{\mathcal O})\subseteq Y_{\mathcal O},
 \qquad K_i=e_1(a_i+4)\le e_1(N+2),\\
 B_L(\phi_i u,T_i v)=p^eB_L(\pi_i u,v),
 \qquad \pi_i\phi_j=\phi_j\pi_i\quad(j>i).
 \end{gathered}
 \tag{TP.4}
\]
The pairing identity permits $u$ with only the cutoffs
$J_k\le a_i+1$; $v$ satisfies all cutoffs. It follows on highest vectors
from (TP.3), then on their cyclic modules by Fock contravariance.

Let $D_{i,w}(n)$ be the dimension of the actual continuation space
\texttt{YYsp}: multiply a nonzero frozen highest vector's cyclic module by
later polynomials and impose all cutoffs. If $b_{i,w}(n)$ is its bad Smith
count, the saturated lattice
\[
 L_1=Y_{\mathcal O}\cap\bigcap_i\ker(\pi_i-1)
 \quad\text{satisfies}\quad
 \operatorname{codim}L_1\le\sum_{i,w}b_{i,w}(n)D_{i,w}(n).
 \tag{TP.5}
\]
A primitive $u\in L_1$ is primitive in $Y_{\mathcal O}$ and has a unit
coefficient. Project its corresponding monomial onto all prefix cutoffs
and multiply by $p^{e_1(N+1)}$. This gives $z\in Y_{\mathcal O}$ with
$|B_L(u,z)|\ge p^{-e_1(N+1)}$: the normalized monomial pairing is the
inverse of a product of integer multinomials. Define
\[
 u'=(p^{K_0}T_0)\cdots(p^{K_{t-1}}T_{t-1})z,
 \qquad D=\lfloor t(E+\kappa)+\kappa\rfloor.
\]
If $e_1(N+2)\le\kappa$, peeling (TP.4) in this order gives
$|B_L((\phi_0\cdots\phi_{t-1})u,u')|\ge p^{-D}$.
Later ties preserve the needed earlier cutoffs; the partial inverses need
not commute. Smith form now supplies $\operatorname{rank}L_1$ integral
pairs with matrix $p^D I$. Thus the source's \texttt{bad} count is bounded
by (TP.5).

\paragraph{TP8--TP11: retaining the multiplicities in the average.}
Write $m_{i,w}=\dim M_{i,w}$. Rational highest-weight decomposition and the
freeze isomorphism give the exact identity
\[
 \sum_w m_{i,w}D_{i,w}(n)=\dim Y(n).
 \tag{TP.6}
\]
Each rational cyclic copy is preserved by the later-prefix projectors;
copies of the same highest weight have isomorphic continuation spaces.
No integral direct-sum assertion enters this identity.

Also $D_{i,w}$ increases with $n_i$. For an increase $u\ge0$, the frozen
highest vector changes from $h$ to $y_\rho^u h$. Project multiplication by
$y_\rho^u$ onto this Cartan component. The resulting map commutes with
raising operators, shifts weights and is nonzero on the highest line,
so its kernel is zero. It preserves later cutoffs because covector weights
are nonpositive and the Cartan projector commutes with the larger-prefix
action. Earlier cutoffs follow from freezing. This proves an injection of
the actual continuation spaces.

On the grid $n_i=s_i+d_i j$, $0\le j<Y_i$, let
$n_i^+=s_i+d_i(Y_i-1)$ and $S_i^+=S_i(n^+)\le s_i$.
Apply (TP.2) holding the other coordinates fixed, replace $D_{i,w}(n)$
by its endpoint value, and use (TP.6). Finite summation gives
\[
 \boxed{\begin{gathered}
 \operatorname{Avg}_n\operatorname{bad}(n;D)
 \le\sum_i\epsilon_i\operatorname{Avg}_n\dim Y(n^{(i,+)}),\\
 \epsilon_i=\frac{2S_i^+}{Y_i}
 +\frac{2\kappa+2S_i^+\ell/Y_i}{E}.
 \end{gathered}}
 \tag{TP.7}
\]
Here $n^{(i,+)}$ changes only coordinate $i$ to its endpoint. The general
conditions are
$\kappa\ge\tfrac32(N+2)r^2\ell$, $e_1(N+2)\le\kappa$ and
$3\ell+\kappa_4\le\kappa$, with the degree bounds used above.
In the literal alternating chain $N=2s+1$, the source size inequalities
$\operatorname{total}\le8a+r$,
$o_i+n_i^++r+z_i+1\le8a+r$, $Y_i\le8a+r$ and
$2S_i^++1\le8a+r$ allow exactly
\[
 \begin{gathered}
 \ell=\log_p(8a+r)+1,\quad e_1=\lfloor r^2\ell\rfloor,
 \\\kappa_4=(s+4)r^2\ell,\quad
 \kappa=3(s+2)r^2\ell.
 \end{gathered}
 \tag{TP.8}
\]
GG supplies the projection, highest-vector and Gram bounds at these constants.
Equations (TP.7)--(TP.8), with $a,E\ge1$
and the literal admissible cutoffs, are precisely \texttt{TheoremMV7}.
They retain the endpoint dimensions with their actual multiplicities.
\textup{[TP1--TP11]}
\noindent\textit{Primary source route.}\cite{tie-partners-actual}


\section{From the Type-I level forms to weighted segment bounds}
\label{sec:segment-weighted}

\status{The actual integral level comparison, block-lattice conversion,
terminal Gram window and weighted summation are proved here.
The final digit-class estimate uses TP's proved multiple-tie mean and
explicitly retains the coarse rank bound (SW8.1) and the stated numerical
conditions on the parameter box. No global dimension-density theorem is
inferred from the summation alone.}

\paragraph{SW1: levels are a direct sum, and their integral span is saturated.}
Use TI1's variables, and write
$j=(j_1,\ldots,j_t,j_+)$ for a level.  Before multiplying by
$\prod E_i^{j_i}E_+^{j_+}\prod Z_i^{z_i}Z_+^{z_+}$, its lattice is exactly
\[
 L_j=Y_{\mathbb Z_p}(d(j))\otimes
           \mathbb Z_p[\xi]_{B-j_+},\qquad
 \operatorname{rank}L_j=\dim Y(d(j))\binom{B-j_++r-1}{r-1}.
 \tag{SW1.1}
\]
The tensor product is embedded by multiplication in disjoint variables.
Its integral description follows coefficient by coefficient in $\xi$.
For a finite relation between different levels, take the lexicographically
largest $S$-degree $j$. Its coefficient is the nonzero monomial
$\prod Z_i^{z_i+j_i}Z_+^{z_++j_+}$ times that level's coefficient polynomial;
no smaller level contributes. Thus the relation is zero level by level.
The same argument proves saturation: if the total polynomial is integral,
extracting this coefficient makes its chain coefficients integral and hence
members of $Y_{\mathbb Z_p}$, subtract the corresponding integral level,
and descend. Formation of $Y$ commutes with field extension: in its finite
multidegree space, each prefix condition is the common kernel of its
forbidden-weight coefficient maps after words in that prefix's own
generators. Each decreasing kernel sequence stabilizes; intersect the
finitely many resulting kernels. Thus the actual block
rank, also over the ramified coefficient field, equals the sum in (SW1.1).
This argument uses the actual $S$ and $Z$ variables, not a dimension count
for a substituted ambient space.

\paragraph{SW2: the integral harmonic change of levels.}
For one pair of residual degrees $(o,n)$ and shift $z$, put
$\widetilde E=ZS+\sum_c w_ce_c$ and
$\widetilde F=\partial_Z\partial_S+\sum_c\partial_{w_c}\partial_{e_c}$.
For a polynomial $f$ independent of $Z,S$, define
\[
 P_{o,n,z}f=\sum_{\ell=0}^{\min(o,n)}
 \frac{(-1)^\ell\widetilde E^\ell Z^zF^\ell f}
 {\ell!^2\binom{o+n+r+z-1}{\ell}}.
 \tag{SW2.1}
\]
The commutator $[\widetilde F,\widetilde E]
=r+1+\mathrm{Euler}$ gives
$\widetilde F P_{o,n,z}f=0$. Its inverse expansion is
\[
 Z^zf=\sum_{\ell=0}^{\min(o,n)}
 \widetilde E^\ell P_{o-\ell,n-\ell,z}
 \left(\frac{F^\ell f}
 {\ell!^2\binom{o+n+r+z-\ell}{\ell}}\right).
 \tag{SW2.2}
\]
These finite identities follow by collecting $\widetilde E^\ell F^\ell$
and the falling-factorial binomial identity; the companion gives the
coefficient cancellation explicitly.

Both expansions cost at most
$u=r\lfloor\log_p(2A)\rfloor+\lfloor\log_p(6A+r)\rfloor$
when the relevant row degrees and shifts are at most $2A$.
Indeed $p^{r\lfloor\log_p(2A)\rfloor}F^\ell/\ell!^2$ preserves the
appropriate integral residual lattice, by multinomial carries, and the
remaining inverse binomial costs at most the second logarithm.
The prefix conditions are preserved: a vector-row derivative and its
coordinate companions form a stable space whose first weight does not
increase; a complete vector--covector contraction commutes with the prefix
action and preserves weight. If the prefix cuts the pair, only its vector
row is inside, so the first observation applies.
This proves the source's \texttt{HypJoin} and \texttt{HypU}.
Different pairs have disjoint variables. Applying (SW2.1)--(SW2.2) successively
therefore proves, for the actual original and harmonic level lattices
$SS_{\mathbb Z_p},PS_{\mathbb Z_p}$,
\[
 p^U PS_{\mathbb Z_p}\subset SS_{\mathbb Z_p},\qquad
 p^U SS_{\mathbb Z_p}\subset PS_{\mathbb Z_p},\qquad
 U=(t+1)\left\lceil(r+1)(\log_p(8A+r)+1)\right\rceil.
 \tag{SW2.3}
\]
Their rational spans agree. Since $\widetilde E$ and $\widetilde F$ are
Fock adjoints, different harmonic levels are orthogonal. On level $j$,
the unnormalized harmonic map $\Theta_j$ has pairing
\[
 B(\Theta_jf,\Theta_jg)=c_j\,B(\Phi_{\rm all}f,g),\quad
 c_j=\prod_{\nu=1}^{t+1}
 j_\nu!^2 z_\nu!\binom{o_\nu+n_\nu+r+z_\nu+j_\nu}{j_\nu}.
 \tag{SW2.4}
\]
Here $o_\nu,n_\nu$ are the residual degrees at that pair.
The norm identity follows from
$\widetilde F\widetilde E^\ell h
=\ell(d_h+r+\ell)\widetilde E^{\ell-1}h$
for harmonic $h$, and
$B(Pf,Pg)=z!B(f,\Phi_{o,n,z}g)$.

\paragraph{SW3: convert the good level factors into the actual block lattice.}
Let $d_j!=\prod_a d_a(j)!$ include the terminal $\xi$ degree, and let
$b_j(D)$ be the bad Smith count of
$B(\Phi_{\rm all}\,\cdot,\cdot)/d_j!$ on $L_j$ at depth $D\in\mathbb Z$.
Choose dual families for all good factors, move them by $p^U\Theta_j$,
and multiply the partner on level $j$ by
$p^M/(c_jd_j!)$, where $M=\max_j v_p(c_jd_j!)$.
The multiplier is integral; (SW2.3)--(SW2.4) give one common pairing
$p^e\delta_{ab}$, with $e=D+2U+M$, and at least
$\operatorname{rank}SS_{\mathbb Z_p}-\sum_jb_j(D)$ pairs.

Let $\Lambda=\log_p(8a+r)+1$, $A=(3s+2)a$, $t\le s$, $r\ge4$ and
$p\ge(r+1)^2$, with $2s+2\le r$. Legendre's formula, applied to the
actual row multidegrees, gives for every supported monomial $\mu$ and level $j$
\[
 v_p(c_jd_j!)\le v_p(\mu!)+
       \big((N+1)r+3(t+1)\big)\lfloor\log_p(8A+r)\rfloor,
 \qquad N=2s+1.
 \tag{SW3.1}
\]
All levels have the same extended row multidegrees: each $Z$ belongs to
its vector row and each $S$ to its covector row. There are
$(N+1)r+2(t+1)$ variables; bounding multinomial carries in these rows
gives that many logarithms. The remaining $t+1$ logarithms bound the
binomials in $c_j$. The factorial part is bounded by the row factorial
product because $j!^2z!o!n!$ divides $(o+z+j)!(n+j)!$.
Since $\log_p(8A+r)\le\Lambda$, it follows that
\[
 e\le v_p(\mu!)+D+\ell_{\rm lev},\qquad
 \ell_{\rm lev}=10(s+1)(r+1)\Lambda+\kappa_0,\quad \kappa_0\ge0.
 \tag{SW3.2}
\]

Here is the exact ramified-lattice conversion. Write $s_b$ for TI1's
factorial constant, $w^2=p$, $F_0=\lfloor D+\ell_{\rm lev}\rfloor$, and
$q=F_0-e-v_p(s_b)$. Multiply both families by $w^q$.
For each coefficient, (SW3.2) gives
$q+v_p(\mu!)+v_p(s_b)\ge0$, precisely TI1's weighted integrality condition.
Their block pairing has valuation
$v_p(s_b)+q+e=F_0$. Saturating the first-family span only improves the
primitive-test bound: a primitive vector of its saturation needs at
worst division by a power of $w$, which lowers the valuation of a suitable
pairing. Thus the source's actual \texttt{BlockBoundUniformiser} holds,
with rank loss $\sum_jb_j(D)$ and depth $D+\ell_{\rm lev}$.

\paragraph{SW4: tensoring and diagonalizing the intermediate form.}
Before the terminal contraction the form is
$b_0=B(\Phi_{\rm ties}\,\cdot,\cdot)/d_j!$ on (SW1.1).
If $B(\Phi_{\rm ties}x_i,y_k)/\prod_a d_a!=p^D\delta_{ik}$,
then
$x_i\xi^\alpha$ and
$\big((B-j_+)!/\alpha!\big)y_i\xi^\alpha$
are integral dual families for $b_0$.
Their pairing is $p^D\delta_{ik}\delta_{\alpha\beta}$.
Consequently
\[
 \operatorname{bad}(b_0,L_j,D)
 \le\operatorname{bad}(b_Y,Y_{\mathbb Z_p},D)
          \binom{B-j_++r-1}{r-1}.
 \tag{SW4.1}
\]
The disjoint tie operators are self-adjoint and commute, so $b_0$ is
symmetric. Over $\mathbb Z_p$ with $p\ne2$, every finite free lattice
with symmetric rational form has an orthogonal integral basis.
Choose a vector with least nonzero pairing valuation. If a minimal
entry is off the diagonal, the sum of its two basis vectors has that
valuation since $2$ is a unit. Pairing divided by its diagonal entry is
integral, giving an orthogonal direct summand; iterate, then retain the
radical. The good basis vectors are exactly those whose diagonal
valuations are at most $D$. This justifies the orthogonal-family step
without assuming a self-dual lattice.

\paragraph{SW5: the terminal Gram window is another polynomial average.}
Fix an orthogonal good family $f_1,\ldots,f_n$ for $b_0$, with nonzero
diagonal entries $d_i$ and $v_p(d_i)\le D$. Let $o,n_\xi$ be the remaining
degrees, $J=\min(o,n_\xi)$, $q=o+n_\xi+r-1$, and vary the terminal shift $z$.
For
$B_z(f,g)=B(\Phi_{o,n_\xi,z}\Phi_{\rm ties}f,g)/d_j!$, set
\[
 \begin{gathered}
 M(z)=p^{x_0}\binom{q+z}{J}\big(B_z(f_i,f_k)\big)_{i,k},\\
 x_0=\big((t+1)(r+1)+(N+1)r+2(t+1)\big)L,\quad
 L=\lfloor\log_p(8A+r)\rfloor+1.
 \end{gathered}
 \tag{SW5.1}
\]
It is an integral binomial-polynomial matrix of degree at most $J$.
Its leading matrix is $p^{x_0}\operatorname{diag}(d_i)/J!$.
To see this, multiply the projector formula by $\binom{q+z}J$:
its $\ell$th term becomes
\[
 (-1)^\ell\binom{q+z-\ell}{J-\ell}
 \frac{B(E^\ell F^\ell\Phi_{\rm ties}f_i,f_k)}
      {\ell!^2\binom J\ell\,d_j!}.
 \tag{SW5.2}
\]
The first $t+1$ operations cost $(t+1)(r+1)L$ powers of $p$;
normalized Fock pairing costs at most $((N+1)r+2(t+1))L$ by row
multinomial carries. Vandermonde's identity gives integral binomial
coefficients in $z$. Only $\ell=0$ has degree $J$, proving the leading formula.
Its determinant has degree $Jn$, and Legendre gives
$v_p(\operatorname{lc}\det M)+Jn/(p-1)\le n(x_0+D+L)$.
The scalar in (SW5.1) is nonzero and integral for $z\ge0$; multiplying
by it can only increase bad Smith counts. TI5 therefore yields
\[
 \frac1W\sum_{i=0}^{W-1}\operatorname{bad}(B_{z_0+i},\langle f\rangle,k)
 \le n\left\{\frac JW+
       \frac{x_0+D+L+Jh/W}{k+1}\right\},
 \quad W,J+1\le p^h,\quad k\ge0.
 \tag{SW5.3}
\]

In the actual window put
$\kappa=3(s+2)r^2\Lambda$,
$D_1=t(Y\kappa+C_M\kappa)+C_M\kappa$, $C_M\ge1$,
$D=\lfloor D_1\rfloor$, $k=\lfloor YD_1\rfloor$.
The literal numerical inputs are
$W\ge a/((r+1)Y)$, $W\le a+1$,
$J\le5b'/2$, $pb'<a$, $p\ge(r+1)^2Y^2$, $Y\ge1$, and $o_+\le a$.
They give $J/W\le1/(2Y)$.
Choose $h=\lceil\log_p(a+1)\rceil$; then $h\le\kappa$ and
$x_0+L\le\kappa\le D_1$. Since $k+1>YD_1$, (SW5.3) is at most
$3n/Y$, hence at most the source's $5n/Y$.
Add the discarded orthogonal directions and use (SW4.1). This proves
the actual \texttt{RestPlus} terminal average.

\paragraph{SW6: fixed digit classes have a fixed index set.}
In a compatible digit class, the coarse degree at site $u$ satisfies
$P_u C_u=a_u-\sum_c\mathrm{digit}_{u,c}$.
The literal segment inequalities give
$(r+1)P_u\le b'$, $\sum_c\mathrm{digit}_{u,c}<b'$,
$b_u\le5b'/2$, $a>25b'$, and $a_u\ge a$.
Thus $C_u\ge A_u$, where $A_u$ is the other coarse degree.
Every kept tie order satisfies $j_i\le o_i\le A_u$,
so its extra condition $j_i\le C_u$ is automatic throughout the box.
The early orders also remain legal: their bound
$\mathrm{sg}_u\le2\lambda A_0/(r+1)$, together with
$4\lambda<r+1$, $A_0\le a/Q$ and $P_u\le Q$,
gives $P_u\mathrm{sg}_u\le a/2<P_uC_u$.
The other kept predicates are independent of $a_u$.
The actual set of kept blocks and legal levels is therefore constant.

Let $u_i(a)$ replace the $i$th tied degree by its largest compatible
endpoint. It stays in the same box. Only one factor in the ambient
dimension $m(a)$ changes, from $\binom{x+r}r$ to $\binom{y+r}r$.
Here $x\ge1$, $y\ge x$, and the source bounds
$xP\ge3a/4$, $yP-xP\le a/((r+1)Y)+x_a$,
$rx_a\le a/8$ give $r(y-x)\le3x$. Hence
\[
 \frac{m(u_i(a))}{m(a)}
 =\prod_{h=1}^r\frac{y+h}{x+h}
 \le(y/x)^r\le\exp\!\big(r(y-x)/x\big)\le e^3.
 \tag{SW6.1}
\]

\paragraph{SW7: the numerical loss in the multiple-tie average.}
TP's exact mean theorem gives, for tie $i$ with progression length $Y_i$,
\[
 \epsilon_i=\frac{2S_i}{Y_i}
       +\frac{2\kappa+2S_i\Lambda/Y_i}{Y\kappa}.
 \tag{SW7.1}
\]
The input is the prefix total $S_i$ with the other tied coordinates at
their upper endpoints, not at the current point.
For the actual box the source reduces the check to
\[
 S_iP_i\le\tfrac52b'+(2s+1)\frac{2a}{p},\quad
 Y_iP_i\ge\frac{a}{2(r+1)Y},\quad
 (r+1)Y^2b'\le a/2,\quad (r+1)^2Y^2\le p.
 \tag{SW7.2}
\]
Thus $S_i/Y_i\le7/Y$. Since $24\Lambda\le\kappa$,
$\epsilon_i\le17/Y\le(30+C_M)/Y$.
Compatible residues form progressions of step one in coarse degree:
division by $P_i$ removes the original spacing $P_i$.
Their initial point obeys
$n_{0,i}P_i\ge a-(r+1)P_i-5b'/2$.
Together with the first bound in (SW7.2), $pb'<a$,
$(r+1)P_i\le b'$ and $p\ge4s+8$, this gives $S_i\le n_{0,i}$.
These checks explain the source's \texttt{RestTMSize}/\texttt{RestTMEps}
interfaces; the other degree and logarithm bounds of TP remain part of
the explicitly stated admissible-input package.

\paragraph{SW8: retain the dimension weights when summing.}
Let $A$ be the compatible top-degree box, $W$ the terminal window,
$I$ the fixed set of block--level indices, and $d(z)$ the terminal
monomial dimension. Write $y(a,z)=\dim Y(a,z)$ and $q(a,z)$ for its
bad count. The required coarse rank input is the literal inequality
\[
 \sum_{z\in I}y(a,z)d(z)\le m(a).
 \tag{SW8.1}
\]
SW1 identifies its left side with the sum of the actual block ranks.
The right-side estimate is the upper-rank clause of
\texttt{flipOneThree.adm11}; it is retained here as a separate input.
Suppose, as proved above and in TP under their matching inputs,
\[
 \begin{aligned}
 \sum_{w\in W}b(a,w,z)&\le |W|
       \{q(a,z)d(z)+5y(a,z)d(z)/Y\},\\
 \sum_{a\in A}q(a,z)&\le\frac{30+C_M}{Y}
       \sum_{i=1}^t\sum_{a\in A}y(u_i(a),z).
 \end{aligned}
 \tag{SW8.2}
\]
Sum first in $z$, use (SW8.1) at $u_i(a)$, then (SW6.1):
\[
 \boxed{\quad
 \sum_{a,w,z}b(a,w,z)
 \le \frac{(30+C_M)t e^3+5}{Y}
                |W|\sum_{a\in A}m(a).\quad}
 \tag{SW8.3}
\]
No injectivity of the endpoint replacement is needed: each updated mass
is bounded separately by $e^3m(a)$. This is why the varying multiplicity
weights cause only a constant loss.

\paragraph{SW9: depth arithmetic.}
For $N_1\ge1$, $s\le N_1$ and $Y\ge1$,
$\kappa\le9N_1(r+1)^2\Lambda$.
It follows directly that
\[
 YD_1+\ell_{\rm lev}
 \le C Y^2N_1^2(r+1)^2\Lambda+\kappa_0
 \quad\text{if }C\ge18(1+C_M)+20.
 \tag{SW9.1}
\]
Taking also
$C\ge2((30+C_M)e^3+5)e^3$
matches the source's \texttt{SegmentWeightedCountingBound}.
With TP's literal constant $C_M=1$, these are fixed constants, independent
of the field, digits and polynomial degrees.

\paragraph{SW10: the coarse-degree gap used in the raw lifts.}
Here $k=r+1$, $\lambda=\mathrm{La}$, $\epsilon=N_1\lambda/k$,
$\delta=\mathrm{de}$, and $\beta=\mathrm{be}$ have their literal
\texttt{Std} meanings. Put $Q_j=Pp^d$, $A_0=\lfloor(a-E_0)/Q_j\rfloor$,
and $A=(g-\sum\text{digits})/P$. Then
$g-kP<AP\le g$, while
$b'(1-\beta)\le g\le b'(1+\beta)$.
The precise standard inequalities used are
\[
 40\epsilon\le\delta<1/6,\quad
 0\le\beta\le\delta/2-18\epsilon,\quad
 E_0+Q\le\epsilon a,\quad
 kP\le kQ\le\epsilon b',\quad 0\le E_0,\quad Q_j\le Q.
 \tag{SW10.1}
\]
The distance $d$ is odd. The forbidden band about $a/b'$ and the
definition of the split index $\sigma$ imply
$p^d<(a/b')(1-2\delta)$ for $i<\sigma$, and
$p^d>(a/b')(1+2\delta)$ for $i\ge\sigma$.
The elementary inequalities
\[
 (1-2\delta)(1+\beta)(1+\delta)\le1-\epsilon,\qquad
 (1-\beta-\epsilon)(1-\delta)(1+2\delta)\ge1
\]
now give, respectively,
\[
 A(1+\delta)<A_0\quad(i<\sigma),\qquad
 A_0<A(1-\delta)\quad(i\ge\sigma).
 \tag{SW10.2}
\]
For the first, multiply by $Q_j$ and compare with
$a(1-\epsilon)<A_0Q_j$, retaining the strict floor inequality.
For the second use $AP>b'(1-\beta-\epsilon)$ and $A_0Q_j\le a$.
Finally $kQ\le\epsilon b'$ follows from $b'\ge k^3Q$,
$\epsilon\ge1/k$ and $k\ge1$.
Thus the exact coarse gap required by RL follows from literal
\texttt{Std}, including its excluded-band hypothesis.

\paragraph{SW11: scope and the completed application.}
This completes the integral-form and summation mechanisms.
RU/RD below supply both coarse rank bounds, HM the actual harmonic
and cloud comparisons, and NP the uniform numerical package and
shifted height estimate. RL and SW12 construct the cleaned lifts
under literal \texttt{Std}. TC proves the full local exceptional-function
theorem; the global obligations are listed later.
\noindent\textit{Primary source route.}\cite{segment-weighted-actual}


\subsection{SW12: the numerical interfaces for the actual kept blocks}
\label{sec:segment-numerics}

\status{The literal standard hypotheses imply $\mathrm{BLk}(0)$,
$o_+\le a_+$, and the source total-order bound with $C_t\ge30$.
NP below supplies those hypotheses uniformly and proves the final shifted-height comparison.}

Assume $\mathrm{Std}(C_0,C_\delta,C_{E1},C_{E2},1,X)$.
Write $k=m+1$, $\lambda=X.\mathrm{La}(C_0)$, $\delta=X.\mathrm{de}$,
$\beta=\mathrm{be}(X,C_0)$ and $\epsilon=N_1\lambda/k$.
The clauses used below give
\[
 \begin{gathered}
 p\ge5,\quad k\ge5,\quad p\ge k^2,\quad
 40\epsilon\le\delta<1/6,\quad
 0\le\beta\le\delta/2-18\epsilon,\\
 \lambda/k\le1/4,\quad \epsilon\ge1/k,\quad
 b'\le a/p,\quad k^3Q\le b'.
 \end{gathered}
\]
In one actual segment, $G\ge2$ is even and its interior distances
$0<d_0<\cdots<d_{s-1}<G$ are odd. Put $P_i=p^{G-d_i}$,
$g_i=b_T(\mathrm{pos}(j,i))$, $A_0=\lfloor(a-E_0)/p^G\rfloor$,
and let $A_i$ be the coarse degree after the $k$ compatible digits.
Thus $g_i-kP_i<A_iP_i\le g_i$ and
$b'(1-\beta)\le g_i\le b'(1+\beta)\le2b'$.
The split is $i<\sigma\iff p^{d_i}<a/b'$.
For the effective order $e_i=w_i+\mathrm{sg}_i$ before the split and
$e_i=o_i$ after it, the signed integer cap is
$\mathrm{cap}_l=A_0+\sum_{i<l}(A_i-e_i)$.
No positivity of $A_0$ or of a cap is assumed. All cap comparisons below
are upper bounds; the nonnegative effective orders can therefore be
discarded when obtaining them. The remaining window inequalities are
literal clauses of $\mathrm{Std}$.

\paragraph{The late cutoff.}
Odd spacing and $E_0\ge0$ give
\[
 P_i\mathrm{cap}_i
 \le\frac a{p^{d_i}}+2b'\sum_{h<i}p^{d_h-d_i}
 \le\frac a{p^{d_i}}+\frac{2b'}p.
 \tag{SW.12a}
\]
Indeed the geometric sum is at most $(p^2-1)^{-1}\le p^{-1}$.
Also $2/p\le\epsilon$ and $kP_i\le kQ\le\epsilon b'$.
For $i\ge\sigma$, the excluded-band clause of $\mathrm{Std}$ strengthens
the failed split test to $p^{d_i}>(a/b')(1+2\delta)$.
Set $t=a/p^{d_i}$, $u=1-\delta/2+17\epsilon$ and
$v=1-\delta/4-\epsilon$. Then
\[
 t<\frac{b'}{1+2\delta},\qquad
 P_i\mathrm{cap}_i\le t+\epsilon b',\qquad A_iP_i>b'u.
\]
Both $u,v$ are positive. The explicit remainders
\[
\begin{aligned}
 uv-\epsilon-(1-\delta/2)(1-\delta/4)
   &=\epsilon(15-15\delta/4-17\epsilon)\ge0,\\
 (1+2\delta)(1-\delta/2)(1-\delta/4)-1
   &=\delta(5/4-11\delta/8+\delta^2/4)\ge0
\end{aligned}
\]
show $P_i\mathrm{cap}_i<A_iP_i v$.
Adding $(\lambda/k)A_iP_i\le\epsilon A_iP_i$ proves
\[
 \mathrm{cap}_i+(\lambda/k)A_i\le A_i(1-\delta/4).
 \tag{SW.12b}
\]
The kept-block upper bound for $o_i$ now gives the strengthened late
clause of $\mathrm{BLk}(0)$; all its early clauses are already in
the definition of a kept block.

\paragraph{The terminal extra order.}
The distinct positive exponents $G-d_i$ imply
$\sum_iP_i^{-1}\le(p-1)^{-1}\le5/(4p)$.
Consequently $\mathrm{cap}_s\le a/p^2+3b'/p$ and
$\sum_i2a/P_i\le4a/p$. The literal kept terminal bound is
\[
 o_+\le\mathrm{cap}_s+
  (\lambda/k+4/p^2)\sum_i2a/P_i+\lambda b_+/k,
 \qquad b_+\le2b'.
\]
It follows that
\[
\begin{aligned}
 o_+&\le L_{\rm cap}:=
  \frac a{p^2}+\frac{3b'}p+
  \left(\frac\lambda k+\frac4{p^2}\right)\frac{4a}p
  +\frac{2\lambda b'}k\\
 &\le\left[\frac1{25}+\frac3{25}
       +\left(\frac14+\frac4{25}\right)\frac45+\frac1{10}\right]a
   =\frac{147}{250}a\le a_+.
\end{aligned}
\tag{SW.12c}
\]
Here $L_{\rm cap}$ is exactly \texttt{localOrderCap}, and the terminal
window gives $a_+\ge a$. This proves the remaining $\mathrm{BLk}(0)$
clauses and the conclusion required by \texttt{OPlusAt\_of} whenever
its explicit hypothesis supplying $\mathrm{Std}$ is available.

\paragraph{The full order, with its minimum term.}
Put
\[
 R=\frac a{p^2}+\frac{b'}p+\frac{\lambda a}{pk}
     +\frac{\lambda b'}k+kQ,
 \qquad M_i=\min\{g_i,a_{\rm start}/p^{d_i}\}.
\]
The early excluded-band inequality gives
$g_i\le a_{\rm start}/p^{d_i}$; the late one gives
$a/p^{d_i}\le M_i$.
Multiplying the respective kept inequalities
$e_i\le A_i(1+4/p^2)+(\lambda/k)\mathrm{cap}_i$ and
$e_i\le\mathrm{cap}_i+(\lambda/k)A_i$ by $P_i$, and using
(SW.12a), therefore proves
\[
 P_i e_i\le M_i+T,\qquad
 T=\frac{2b'}p+\frac{2\lambda b'}k+\frac{8b'}{p^2}
       +\frac{\lambda a}{pk}+\frac{2\lambda b'}{pk}.
\]
The digit contribution is at most $2kQ$ per frame. Using $P_i\ge p$
in the terminal bound, and $1/p\le1/5$, gives the following coefficients
of the five summands of $R$ in the complete error:
\[
 (1+2s),\quad6s,\quad3s,\quad4s+2,\quad2s.
\]
Each is at most $30(s+1)$. Thus the actual natural-number order obeys
\[
 \mathrm{blockOrd}_j\le\sum_iM_{j,i}+30(s_j+1)R.
 \tag{SW.12d}
\]
The cyclic segment interiors are disjoint; each contributes its selected
odd points and one distinct even start. Hence
$\sum_j(s_j+1)\le|TT|$. Summing (SW.12d), for $C_t\ge30$ and $c\ge0$,
\[
 \sum_j\mathrm{blockOrd}_j
 \le\sum_{j,i}M_{j,i}+C_t|TT|(R+Qc)=X.t_0(C_0,C_t).
 \tag{SW.12e}
\]
This is the literal \texttt{ord\_all} endpoint. Its proof also works for
$\mathrm{Std}$ with arbitrary $H$, since this last order estimate uses
neither of the clauses containing $H$.
\noindent\textit{Primary source route.}\cite{segment-numerics-actual}


\section{Actual paired lifts and finite cleanup}
\label{sec:paired-lifts}

We now construct the two polynomials used in the Type-I pairing. Work
in the literal unramified quadratic setting, with uniformizer \(w^2=p\),
\(v_p(p)=1\), \(k=m+1\), and \(Q=p^{G'}\). Assume the actual
\(\mathrm{Std}\) hypotheses, \(C_{E1},C_{E2}\ge17\), \(C_m\ge17\).
The kept-block, terminal-degree and order bounds are those of SW12;
the coarse-degree alternative is SW10. The endpoint height below is
\(t_0\); its subsequent comparison with \(t_I\) is separate.

\paragraph{RL1--RL3: correction of the actual segment form.}
For a block \(b\), the physical monomial substitution is
\[
 I_b(\alpha)=\rho_b+\sum_v P_v\alpha_v e_{\iota(v)}.
\]
The map \(\iota\) is injective, \(P_v=p^{g_v}\) at odd positions and
\(P_v=1\) at the terminal position. Balanced degrees imply that the
total auxiliary degree \(c_\ell\) at each physical letter, and the total
degree at each radix \(P_v\), are fixed. Put
\[
 C_b=\prod_\ell c_\ell!,\quad N_b=C_b\,\mathrm{sp}_b,\quad
 v_b=v_p(N_b),\qquad
 L_j=\sum_\ell\lfloor\log_p d_\ell\rfloor,\quad
 z_j=\frac{p^{kL_j}}{\prod_\ell d_\ell!}.
\]
Here \(\mathrm{sp}_b\) is the source's integer \(\mathrm{specConst}\);
zero degrees contribute zero to \(L_j\).
DT's inverse factorial-unit approximation supplies
\[
 g_b(\alpha)=\sum_{|\nu|\le4c}q_\nu\binom{\alpha}{\nu},\qquad
 I_b(\alpha)!\,g_b(\alpha)
   =\Phi_b\alpha!(1+p^ch_\alpha),\quad q_\nu,h_\alpha\in O.
 \tag{RL.1}
\]
Only odd variables occur in \(\nu\). Explicitly,
\(\Phi_b=R_b\prod_v\rho_v!(P_v!)^{\alpha^0_v}\), where \(R_b\)
contains the fixed factorials outside \(\iota\), and \(\alpha^0\)
puts each normal-letter degree in its first coordinate.

The scalar
\[
 \kappa_b=\frac{z_jp^{v_b}\Phi_b}{\mathrm{sp}_b}
 \quad\text{satisfies}\quad
 0\le v_p(\kappa_b)\le(k+1)L_j.                  \tag{RL.2}
\]
Indeed \(C_b/\alpha^0!\) is a product of binomial coefficients, whereas
\(\prod d_\ell!/I_b(\alpha^0)!\) is a product of \(k\)-part multinomials.
Their valuations are at most \(L_j\) and \((k-1)L_j\), respectively.
Legendre's formula makes \(v_p(\kappa_b)\) equal to \(kL_j\) plus the
first valuation minus the second.

Replace each binomial in \(g_b\) by the diagonal operator
\(x^\nu\partial^{[\nu]}\), obtaining \(G_b\), and set
\[
 s_b(u)=w^{v_b}G_bu,\qquad s'_b(u)=w^{v_b}u.       \tag{RL.3}
\]
The weighted coefficient condition is
\(v_p(u_\alpha)+\tfrac12v_p(\alpha!\mathrm{sp}_b)\ge0\).
Since \(\alpha!\mathrm{sp}_b\mid N_b\), both outputs are integral.
These operators reduce the SSZ margin by at most \(4c\): expand the
coefficient of \(t^n\) under \(x\mapsto(1+t)x\), and in each tied factor
use \(ZS=E-q\), where \(E=ZS+q\) and \(q=\sum_aW_aE_a\).
Each changed coordinate or separated pair costs at most one prefix
weight. All binomial coefficients are integral; SW1 supplies saturation.
Thus the maps land in \(\mathrm{SSZ}(m_0(13))\), because
\[
 m_0(C_m)-4c\ge m_0(13),\qquad
 m_0(C)=CkN_1Qc+12c .
\]
Equation (RL.1), summed coefficientwise, gives
\[
 z_j\operatorname{Fock}(\operatorname{io}_bs_bu,
                        \operatorname{io}_bs'_bv)
 =\kappa_b\bigl(\mathrm{sp}_b\operatorname{Fock}(u,v)+w^c\eta\bigr).
 \tag{RL.4}
\]
Distinct blocks are exactly orthogonal: reduce physical exponents modulo
\(P_i\), recover the early coordinate-zero degrees, then recover the
remaining block orders from the balanced degrees. Segments use disjoint
physical letters. Moreover \(z_{\rm Ten}\alpha!\in O\) for every physical
monomial of the fixed degree, by the same multinomial valuation bound.
Consequently (RL.4) factors over segments and is unchanged modulo \(w^c\)
when either symbol is changed modulo \(w^c\).

\paragraph{RL4--RL7: actual raw polynomials and their tail spaces.}
Put \(A=13kN_1Q\), \(\mu=mm+nA\), \(E=(c-n)A\), where
\(mm\in\{0,12c\}\). The numerical wrappers admit explicit short estimates.
With \(\epsilon=N_1\mathrm{La}/k\), the early coarse alternative and
\(\mathrm{Std}\) give
\[
 \mu+N_1E\le4\epsilon A_0,\quad
 \sum_{i<\sigma}\mathrm{sg}_i\le2\epsilon A_0+N_1E,\quad
 \sum_{\ell\le i}w_\ell\le A_i(1+\epsilon),\quad
 A_i(1+\delta)<A_0 .
\]
Since \(40\epsilon\le\delta<1/6\), these prove the literal
floor inequalities \(\mathrm{ColumnNumericsAt}\).
IP's joined one-plane/normal paths therefore give actual Frobenius
column polynomials, of letter degree at most \(Q\).

Multiply these columns by the Frobenius minors and contractions
\[
 K_{W,a}=X_{\mathrm{odd},+,a}X_{\mathrm{top},+,0}^{p^d}
          -X_{\mathrm{odd},+,0}X_{\mathrm{top},+,a}^{p^d},
 \qquad K_S=\sum_aX_{\mathrm{top},+,a}^{p^d}X_{\mathrm{odd},-,a},
\]
with the recorded digit and level multiplicities. Include
\[
 \left[X_{\mathrm{top}_{j_0},+,0}
 \prod_{a>0}\left\{\prod_jX_{\mathrm{top}_j,+,0}^{p^{G_j}-1}
                    -\prod_jX_{\mathrm{top}_j,+,a}^{p^{G_j}-1}
             \right\}\right]^c.                         \tag{RL.5}
\]
This localizer vanishes modulo \(w\) at every other line, by cyclic
Frobenius telescoping. A residual monomial makes the own leading term
exactly the prescribed symbol. Its exponents are nonnegative because
\[
 \mathrm{usedX}\le a-E_0+(N_1+c)kQ,\qquad
 E_0\ge\mathrm{localOrderCap}+cA+(N_1+c)kQ .
\]

For tails removing at most \(e\) factors, divide each new exponent by
\(P_i\), keeping both quotient and remainder. The remainder defines
the new digits. The quotients give a new block, and
\(Z^zS^y=Z^{z-y}(E-q)^y\) gives its exact integral expansion.
The required \(z\ge y\) follows from
\[
 \delta A_i/4\ge c(A+2Q+3k)+4.
\]
Deleting columns costs \(N_1e\) in the prefix bound; the new monomial
and tied factors cost at most \((N_1+1)(k+Qe)\). The cap decreases by
at most \(N_1(3k+2Qe)\), and
\[
 N_1e+(N_1+1)(k+Qe)+N_1(3k+2Qe)\le eA.          \tag{RL.6}
\]
Taking floors preserves this integer inequality. This proves the
actual \(\mathrm{NumCoreAt}\) and tail-space membership, including
the natural-number subtractions.

\paragraph{RL8--RL11: finite cleanup and the genuine paired endpoint.}
Write the product as \(f_R\prod_af_a^{m_a}\), with exact rigid
chart factor and own leading terms \(w^{o_a}\sigma_a\).
For \(t_a=w^{-o_a}\Phi f_a-\sigma_a\),
\[
 v_w([X^M]t_a)\ge\max(1,\operatorname{nd}(M)-o_a).
\]
Expansion and (RL.6) give a raw lift of height \(o\) satisfying
\[
 \Phi(\mathrm{raw})=\sigma+\sum_{h=1}^{n-1}w^h\tau_h\pmod{w^n},
 \quad \tau_h\in\mathrm{SSg}(mm+(n-h)A,(c-n+h)A,o+h).
\]
Inductively lift \(\tau_h\) to precision \(n-h\) and subtract
\(\sum_hw^hF_h\). Multiplication by \(w^h\) cancels its additional
height \(h\); all other-edge conditions remain divisible by \(w^c\).
At \(n=c\) this produces actual cleaned lifts of height
\(o\le\lfloor t_0(C_0,C_t)\rfloor\), \(C_t\ge30\).

For the edge basis \(b\), write
\[
 b'=b^{-T}\operatorname{diag}(p,1,\ldots,1),\qquad
 D_w=\operatorname{diag}(w^{-1},1,\ldots,1).
\]
The reversed dual midpoint satisfies \(b'D_w=(bD_w)^{-T}\).
The actual flat maps and normalized Fock map
\(\Xi\) therefore satisfy
\[
 \mathrm{beHH}(\Theta_bF,\Xi\Theta_{b'}F')
          =z_{\rm Ten}\operatorname{Fock}(\Phi_1F,\Phi_1F').
\]
Normalized integrality sends the dual cleaned lift into
\(\mathrm{CCe}'\). The resulting actual vectors satisfy
\[
 \begin{gathered}
 F_\beta(u)\in\mathrm{CCe},\quad G_\beta(u)\in\mathrm{CCe}',\\
 \mathrm{beHH}(F_\beta(u),G_{\beta'}(u'))
 =\mathbf1_{\beta=\beta'}\kappa_\beta
       \prod_j\mathrm{segForm}_j(u_j,u'_j)+w^c\nu ,\\
 v_p(\kappa_\beta)\le4\kappa_0.
 \end{gathered}                                  \tag{RL.7}
\]
Nonretained block components are zero. This proves the literal
\(\mathrm{PairLiftAt}\) conclusion; the separate digit-class rank
and subsequent global parameter choices are not inferred from it.
\noindent\textit{Primary source route.}\cite{paired-lifts-actual}



\paragraph{The paired-lift endpoint and its application.}
RL constructs the actual paired maps into $\mathrm{CCe}$ and
$\mathrm{CCe}'$ at height $\lfloor t_0\rfloor$, with scalar
valuation at most $4\kappa_0$. The following sections prove the
remaining rank, harmonic, numerical and general-weight steps,
culminating in TC6's complete local Type-I theorem.

\section{The actual coarse upper rank on each digit class}
\label{sec:rank-upper}

\status{This proves the upper-rank premise retained by SW8 for the
literal block lattices, under \texttt{Std} with $H\ge1$. The integral
path basis and its prefix adaptation are IP3--IP4; their exact dimension
count is CB3. No additional matching or representation premise remains.
The uniform supply of \texttt{Std} is a separate numerical result.
RU9--RU11 also prove the actual segment-to-HH assembly, using RD9's
lower count and RL11's pair lifts.}

\paragraph{RU1--RU3: ranks are weighted path counts.}
Work over a characteristic-zero DVR $O$ with fraction field $K$.
Put $r=X.m$, $k=r+1$ and
\[
 h_r(d)=\binom{d+r-1}{r-1},\qquad
 D_r(\lambda)=\prod_{a<b}
       \frac{\lambda_a-\lambda_b+b-a}{b-a}.                 \tag{RU.1}
\]
For a signed word, a path $\kappa^0=0,\ldots,\kappa^N$ adds a
horizontal strip of the prescribed degree at a vector letter and
removes one at a covector letter. IP4 selects exactly those paths
whose first row obeys every prescribed prefix cutoff. CB3 counts
the terminal standard choices. Thus the actual constrained space is
\[
 \dim_K Y=\sum_{\kappa\ {\rm allowed}}D_r(\kappa^N).
                                                               \tag{RU.2}
\]
The source defines its cutoff as the largest Lie-stable subspace
with the corresponding coordinate-weight bound. In characteristic
zero this is precisely IP4's generic-group bound: exponentiate the
nilpotent root operators and use the diagonal weight decomposition.

For a block at level $(j,j_+)$, SW1's literal multiplication map is
injective: the largest monomial in the independent $S$ variables
recovers that level's coefficient. Its two coefficient factors use
disjoint chain and terminal $\xi$ variables. Clearing denominators
therefore gives
\[
 \operatorname{rank}(\text{block})
   =\sum_{\text{legal }(j,j_+)}
       \dim_KY_{j,j_+}\,h_r(B-j_+).                         \tag{RU.3}
\]
This is also the rank of the actual factorial-weighted block lattice:
it lies in the same rational space and contains a full integral
coefficient lattice.

Two exact counting identities will evaluate the upper bound:
\[
 \begin{aligned}
 \sum_{\kappa\ {\rm unrestricted}}D_r(\kappa^N)
       &=\prod_\ell h_r(d_\ell),\\
 \sum_{\nu_{a+1}\le\lambda_a\le\nu_a}D_r(\lambda)
       &=D_{r+1}(\nu),&
 \sum_{j=0}^{B}h_r(B-j)&=h_{r+1}(B).
 \end{aligned}                                             \tag{RU.4}
\]
The first counts the ordinary monomial basis using IP3. The second
is CB3's tableau branching, valid also for signed dominant weights
after a common determinant shift. The last counts monomials by the
exponent of one distinguished variable. All path sets are finite:
each coordinate is bounded by the total signed degrees.

\paragraph{The actual segment dictionary.}
Fix compatible digits $\rho_i^z,\rho_i^y$ at the $s$ interior
positions and write
\[
 P_i=p^{G-d_i},\quad
 A_i=\frac{b_i-\sum_c\rho^z_{ic}}{P_i},\quad
 C_i=\frac{a_i-\sum_c\rho^y_{ic}}{P_i},\quad
 A_*=\left\lfloor\frac{a-E_0}{p^G}\right\rfloor,\quad B=b_+.
                                                               \tag{RU.5}
\]
Compatibility makes $A_i,C_i$ nonnegative integers.
Split positions into early $i<\sigma$ and late $i\ge\sigma$.
A kept block has fields $(o_i,w_i,s_i,o_+)$, where $s_i$ denotes
the source's \texttt{sg}. Early positions have
$o_i=0$, $w_i\le A_i$, $s_i\le C_i$; late positions have
$o_i\le A_i$, $w_i=s_i=0$. Put
\[
 e_i=\begin{cases}w_i+s_i&i<\sigma,\\o_i&i\ge\sigma,\end{cases}
 \qquad c_n=A_*+\sum_{i<n}(A_i-e_i).                         \tag{RU.6}
\]
The residual word is $V_0,E_0,\ldots,V_{s-1},E_{s-1},U$,
alternating vectors and covectors and ending with a vector. At a
legal level its degrees are
\[
 \begin{array}{c|cc}
 &d_{2i}&d_{2i+1}\\ \hline
 i<\sigma&w_i&C_i-s_i\\
 i\ge\sigma&o_i-j_i&C_i-j_i
 \end{array},
 \qquad d_{2s}=o_+-j_+.                                    \tag{RU.7}
\]
Here $j_i\le\min(o_i,C_i)$ and $j_+\le\min(o_+,B)$.
The actual cutoffs are $\lfloor c_i-M\rfloor$ after each late
$V_i$ and $\lfloor c_s-M\rfloor$ at the endpoint, with $M\ge0$.

\paragraph{RU4--RU5: the missing early row bounds.}
Let $\lambda=\mathrm{La}$, $\epsilon=N_1\lambda/k$,
$\delta=\mathrm{de}$ and $\beta=\mathrm{be}$.
The literal \texttt{Std} inequalities and SW10 give
\[
 \begin{gathered}
 k\ge5,\quad p\ge k^2,\quad s\le N_1,\quad
 40\epsilon\le\delta<1/6,\quad \lambda,N_1\ge1,\\
 0\le\beta<1/12,\qquad
 A_*\ge a/(2Q)\ge k,\qquad A_i(1+\delta)<A_*\quad(i<\sigma).
 \end{gathered}                                             \tag{RU.8}
\]
For $l<i$, increasing physical distances give $P_l\ge pP_i$.
Digit compatibility gives $A_lP_l\le b_l$ and
$b_i-kP_i<A_iP_i$. Using
$b'(1-\beta)\le b_i\le b'(1+\beta)$ yields
$A_l\le2(A_i+k)/p$.
Also \texttt{Kept} gives $s_l\le2\lambda A_*/k$ early.
Consequently
\[
 \sum_{l\le i}(A_l+s_l)
 \le A_i+s\left(\frac{2(A_i+k)}p+\frac{2\lambda A_*}k\right)
 \le A_i+3\epsilon A_*\le A_*.                              \tag{RU.9}
\]
For the last inequalities, $2(A_i+k)/p\le4A_*/k^2
\le\lambda A_*/k$, while
$(A_*-A_i)(1+\delta)>A_*\delta$ gives
$A_*-A_i>34\epsilon A_*$.
Since $w_l\le A_l$, for every $n\le i+1$,
\[
 \sum_{l\le i}w_l-c_n
 \le\sum_{l\le i}(A_l+s_l)-A_*\le0.                         \tag{RU.10}
\]
This proves both neighboring cap bounds without assuming early
caps are monotone.

A vector strip increases the first row by at most its degree;
a covector strip cannot increase it. Hence an early residual path
has $\kappa^{2i+1}_0\le\sum_{l\le i}w_l$.
At late positions the actual cutoff gives
$\kappa^{2i+1}_0\le c_i\le c_{i+1}$.
The terminal cutoff gives $\kappa^{2s+1}_0\le c_s$.
Thus, at every position,
\[
 \kappa^{2i+1}_0\le\min(c_i,c_{i+1}),\qquad
 \kappa^{2s+1}_0\le c_s.                                   \tag{RU.11}
\]

\paragraph{RU6: insert one row and recover all the data.}
Define a path with $r+1$ coordinates by
\[
 \nu^0=0,\quad
 \nu^{2i+1}=(c_i,\kappa^{2i}),\quad
 \nu^{2i+2}=(c_i+A_i-d_{2i},\kappa^{2i+1}).                  \tag{RU.12}
\]
The odd formula includes $i=s$.
It is a path for the unrestricted word
\[
                    +A_*,+A_0,-C_0,\ldots,+A_{s-1},-C_{s-1}.
\]
Indeed, adjoining first entries $f\le f'$ to a horizontal strip
$\alpha/\gamma$ is legal exactly when $\alpha_0\le f$;
its size increases by $f'-f$.
Equation (RU.11) supplies this inequality, and
\[
 c_i+A_i-d_{2i}-c_{i+1}
   =\begin{cases}s_i&i<\sigma,\\j_i&i\ge\sigma,\end{cases}
 \qquad
 c_i+A_i-d_{2i}-c_{i+1}+d_{2i+1}=C_i.                      \tag{RU.13}
\]
The final vector strip and endpoint bound say that
$\kappa^{2s+1}$ interlaces $\nu^{2s+1}$.

Send a source datum to
$(\nu,\kappa^{2s+1},j_+)$.
This preserves its weight
$D_r(\kappa^{2s+1})h_r(B-j_+)$.
It is injective: the last $r$ rows of $\nu$ recover every earlier
$\kappa$, and the stored endpoint recovers the last one.
Strip row-sum differences recover all $d_\ell$. Then early
$(w_i,s_i)=(d_{2i},C_i-d_{2i+1})$; late
$(j_i,o_i)=(C_i-d_{2i+1},d_{2i}+j_i)$; and
$o_+=d_{2s}+j_+$. The fixed digits and zero fields finish recovery.

\paragraph{RU7--RU8: sum the target weights.}
Enlarge the image to all unrestricted target paths, all interlacing
endpoints and $0\le j_+\le B$. Their weights are nonnegative.
Summing in that order with (RU.4), and using (RU.2)--(RU.3), proves
\[
 \boxed{\operatorname{rankOnDigits}
 \le h_{r+1}(A_*)\prod_i h_{r+1}(A_i)h_{r+1}(C_i)
                         h_{r+1}(B)=\texttt{mmB}.}         \tag{RU.14}
\]
This holds for every real $M\ge0$, with no lower-bound margin
restriction. At $M=0$ it is exactly SW8's missing upper inequality.
The numerical companion's uniform implication to \texttt{Std}
therefore supplies the literal \texttt{UpAt} and \texttt{RestUp}
package. Its thresholds may be taken as
$T_0=T_1=1$ and
$T_2=1000(CE_1+CE_2+1)(C_c+1)$, in the source's original order of
constant choices.

\paragraph{RU9: restore the full tensor dimension.}
Here use RD9, under $C_0\ge880$ and \texttt{Std} with $H=4$.
Write $q_j=p^{G_j}$, $q=p^{n'}$, $J=\#\{\text{segments}\}$,
$N=|TT|$, and $R_j=\operatorname{segRank}_j(M)$. Define
\[
 \mathcal B=e^{-\lambda/8}
       +\frac{kQ(M+N_1k)}{b'}+\frac1Y+\frac{kE_0}{a},
 \qquad \kappa_P=C_\kappa N\mathcal B,\quad C_\kappa\ge3.
                                                               \tag{RU.15}
\]
Coordinatewise division $e=\rho+Pu$ gives the exact identity
\[
 \sum_{\substack{0\le\rho_c<P,\ \sum\rho_c\le g\\
                  P\mid g-\sum\rho_c}}
      h_k\!\left(\frac{g-\sum\rho_c}{P}\right)=h_k(g).
                                                               \tag{RU.16}
\]
Summing RD9 over all interior digit rows therefore gives
\[
 R_j\ge L_j h_k(A_{*,j})
       \left(\prod_{i\text{ interior }j}h_k(a_i)h_k(b_i)\right)
       h_k(b_{\rm next}),
 \quad
 L_j=(1-e^{-\lambda/8})(1-e^{-\lambda/3})
              \left(1-\frac{r\lceil M\rceil}{A_{*,j}}\right),
                                                               \tag{RU.17}
\]
when $\lceil M\rceil\le A_{*,j}N_1\lambda/k$.
Kept blocks form a finite set: their digits, $o_i,w_i,s_i$ are
bounded, and $c_s\le A_{*,j}+\sum_i b_i$ bounds the terminal
$o_+$ through its literal kept inequality.

Suppose $\mathcal B<1/6$. Put
$S=kQ/a$, $t=kQ(M+N_1k)/b'$ and
$\delta_f=2/(kY)+(E_0+Q)/a$.
Since $t\ge S(M+5)$ and $A_{*,j}\ge a/(2Q)$,
\[
 \frac{k\lceil M\rceil}{A_{*,j}}<2S(M+1)\le2t<1/3,
 \qquad k\delta_f\le2\mathcal B<1/3.                       \tag{RU.18}
\]
Thus the required margin condition holds. With
$\theta=1-\delta_f$, Bernoulli and the product inequality give
$L_j\theta^r\ge1-3\mathcal B$: the sum of the four losses is at
most
$2e^{-\lambda/8}+S(2M+3)+2/Y+kE_0/a\le3\mathcal B$.
For the actual head degree $g=a_{\rm start}$, its window and the
floor defining $A_{*,j}$ give, for $1\le u\le r$,
\[
 \theta(g+u)\le a-E_0+u-Q\le a-E_0+u-1
                  <q_j(A_{*,j}+u).
\]
Multiplication yields
$\theta^r h_k(g)\le q_j^rh_k(A_{*,j})$.
Hence $(1-3\mathcal B)T_j\le q_j^rR_j$, where
\[
 T_j=h_k(a_{\rm start})\bigl(\prod_{i\text{ interior }j}
                 h_k(a_i)h_k(b_i)\bigr)h_k(b_{\rm next}).
\]

The cyclic intervals partition the selected odd positions, and
the successor permutes their even endpoints. Consequently
$\prod_jT_j=\dim HH$ and $\prod_jq_j=q$.
This is the literal polynomial tensor dimension: CB3 gives
$h_k(d)$ for both vector shape $(d,0,\ldots,0)$ and normalized
covector shape $(d,\ldots,d,0)$; their disjoint-variable tensor
bases multiply. Since $J\le N$,
\[
 \boxed{(1-\kappa_P)\dim HH
       \le\operatorname{gaussian}(q,k)\prod_jR_j.}         \tag{RU.19}
\]
Indeed multiply the segment bounds and use Bernoulli and
$\operatorname{gaussian}(q,k)=1+q+\cdots+q^r\ge q^r$.
If $\kappa_P\ge1$, the claim is immediate; otherwise $N\ge2$
implies $\mathcal B<1/6$. Thus (RU.19) holds for every $M\ge0$.

\paragraph{RU10--RU11: produce the actual HH certificate.}
The weighted block lattice is finite free: in its finite monomial
support the coefficient of $X^\alpha$ lies in
$w^{-v_p(\alpha!)-v_p(\mathrm{specConst})}O$, where $w^2=p$.
Denominator clearing gives its full rational rank. Summing over
the finite kept set proves $\operatorname{rank}_O L_j=R_j$.
Moreover $\mathrm{specConst}\cdot\mathrm{fock}$ is integral on
each block, by adding the two coefficient valuation bounds.

Apply Smith form separately on each block. If a saturated good
submodule $L_{1,j}$ has rank $g_j$, at least $g_j$ diagonal
entries have $v_p\le D_j$. Otherwise its map to those good
coordinates has a nonzero kernel. Divide a kernel vector by its
largest common uniformizer power; saturation retains it and makes
it primitive. Every pairing against $L_j$ then has valuation
$>D_j$, contradicting its good-partner property. The selected
pairs retain single-block support.

For supplied good shares $g_j\ge(1-\kappa_j)R_j$, RL11 (with
$CE_1,CE_2,C_m\ge17$ and $C_t\ge30$) tensors
these choices into actual $F_I\in\mathrm{CC}_e$ and
$G_I\in\mathrm{CC}'_e$ at height $\lfloor t_0\rfloor$, with
\[
 \beta_{HH}(F_I,G_J)=\delta_I\mathbf1_{I=J}+w^c\nu_{IJ},
 \quad 0\ne\delta_I\in O,\quad
 v_p(\delta_I)\le D_*:=\sum_jD_j+4\kappa_0.              \tag{RU.20}
\]
Here $D_j,\kappa_j\ge0$, $c>4D_*$, and the literal normalization
is $\beta_{HH}(F,G)=p^{A_{HH}}G(F)$.
Factor the matrix as $\Delta(I+E)$; all $E_{IJ}$ lie in $wO$
because $v_w(\delta_I)\le2D_*<c$. The adjugate inverts $I+E$
over $O$. If $\delta_{\max}$ has maximal valuation, multiplication
on the right by
\[
 (I+E)^{-1}\operatorname{diag}(\delta_{\max}/\delta_I)
                                                               \tag{RU.21}
\]
gives $\delta_{\max}I$. These are integral column operations,
so they preserve every lattice and precision condition on the
$G_I$, with no height loss. This is the literal
\texttt{FactorsLE} certificate of rank $\prod_jg_j$ and depth $D_*$.

Finally split off $\kappa_P+\sum_j\kappa_j\ge1$, when the desired
mass is nonpositive. Otherwise every remaining factor is nonnegative;
$\prod_j(1-\kappa_j)\ge1-\sum_j\kappa_j$ together with (RU.19)
gives exactly
\[
 \operatorname{RhoGE}_{HH}\!\left(
 D_*,\frac{(1-\kappa_P-\sum_j\kappa_j)\dim HH}
                 {\operatorname{gaussian}(q,k)}\right).  \tag{RU.22}
\]
It is the source's \texttt{aux\_3}/\texttt{HerrExtensionClass}
for the actual $\mathrm{CC}_e,\mathrm{CC}'_e$.
The separate weighted-partner and numerical results supply the
particular $D_j,\kappa_j$ and compare their totals with $D_0,k_x$;
HM then transports this certificate to the harmonic model.
\noindent\textit{Primary source route.}\cite{coarse-rank-upper}


\section{The actual coarse lower rank and density}
\label{sec:rank-lower}

\status{The weighted lower count, the actual kept-block injection and the
three source loss factors are proved under literal $\mathrm{Std}(H=4)$
and $C_0\ge880$. The shared identities RU1--RU3 identify this count with
the actual lattice rank. With RU8 this gives the source two-sided
half-density window. Uniform parameter supply is a separate input.}

Write $r=X.m$, $k=r+1$, $\lambda=C_0N_1^2\log a$,
$\epsilon=N_1\lambda/k$, and $\delta=X.\mathrm{de}$.
Fix an actual segment and compatible digits, and put
\[
 P_i=p^{G-d_i},\qquad A_*=\left\lfloor\frac{a-E_0}{p^G}\right\rfloor,
 \qquad
 A_i=\frac{b_i-\sum_c r^z_{ic}}{P_i},\quad
 C_i=\frac{a_i^{\rm top}-\sum_c r^y_{ic}}{P_i},\quad B=b_+.
\]
The quotients are nonnegative integers by compatibility.
Here $0<d_0<\cdots<d_{s-1}<G$ are odd and $G$ is even.
Let
$S_d(g)=\binom{g+d-1}{d-1}$ and
\[
 \mathcal D(A)=S_k(A)\prod_{i<s}S_k(A_i)S_k(C_i),\qquad
 \mathcal B=\mathcal D(A_*)S_k(B)=\mathrm{mmB}.
\]
For a real margin $h\ge0$, put $M=\lceil h\rceil$ and suppose
$M\le\epsilon A_*$. The claimed exact lower bound is
\[
 \mathrm{rankOnDigits}(h)\ \ge\
 (1-e^{-\lambda/8})(1-e^{-\lambda/3})
 \left(1-\frac{rM}{A_*}\right)\mathcal B.
 \tag{RD.1}
\]

\paragraph{The actual pair distribution.}
Put $A=A_*-M$. A pair consists of a signed Pieri chain
$\mu^0,\ldots,\mu^{2s+1}$ in $k$ coordinates, starting at zero,
with initial vector degree $A$ and subsequent pairs of signed degrees
$(+A_i,-C_i)_{0\le i<s}$, and an $r$-coordinate weight $\ell$
interlacing the endpoint.
The weight of a pair is the Weyl dimension $w_r(\ell)$.
The signed path identity and the interlacing branching identity,
proved in RU2--RU3, give
\[
                       \sum_{\text{pairs}}w_r(\ell)=\mathcal D(A).
 \tag{RD.2}
\]
Delete the first coordinate, writing $K_t^n=\mu_{t+1}^n$, and set
\[
\begin{gathered}
 c_i=\mu_0^{2i+1},\quad \rho_i=K_0^{2i+2},\quad
 W_i=\sum K^{2i+2}-\sum K^{2i+1},\\
 O_i=c_i+A_i-c_{i+1},\quad S_i=O_i-W_i,\quad
 u_i=O_i-\rho_i,\quad n_i=-\ell_{r-s+i}.
\end{gathered}
\]
Horizontal interlacing gives $0\le W_i\le A_i$, $0\le S_i\le C_i$
and $\rho_i\le c_i,c_{i+1}$. A pair is \emph{kept} precisely when
\[
 u_i+\frac{4(A_i+r)}{p^2}
       \le\frac\lambda k\max(c_i,A_i),\qquad
 C_i(1-\lambda/k)\le n_i\quad(i<s).
 \tag{RD.3}
\]

\paragraph{Why the bad weighted mass is small.}
The standard hypotheses and odd spacing imply
\[
\begin{gathered}
 k\ge10,\quad 4\lambda\le k,\quad p\ge k^2,\quad
 A,A_i,C_i\ge k^3,\quad 2s+1\le r,\\
 \sum_{l<i}A_l\le\frac{4A_i}{p^2},\qquad
 \sum_{l<i}C_l\le\frac{4C_i}{p^2}.
\end{gathered}
\]
For the geometric bounds, the consecutive degrees satisfy
$p^2A_i\le2A_{i+1}$ and $p^2C_i\le2C_{i+1}$: multiply by
$P_{i+1}$, use $P_i\ge p^2P_{i+1}$, and use the degree windows and
$kP_{i+1}\le b'/25$.
Following the upper and lower interlacing boundaries yields
\[
\begin{gathered}
 0\le\rho_i\le A_i(1+4/p^2),\quad u_i\ge-4A_i/p^2,
 \quad0\le n_i\le C_i(1+4/p^2),\\
 0\le\ell_0\le c_s,\quad
 0\le\ell_t\le\rho_{s-t}\ (1\le t\le s),\quad
 \ell_t=0\ (s<t<r-s).
\end{gathered}
\]
For example, negative mass of a tail grows by at most $C_i$ at a
covector step and decreases at a vector step; following its last
coordinate bounds $n_i$ by $\sum_{l\le i}C_l$.

Here is the quantitative core. Put $q=4/p^2$, $\eta=k^{-2}$,
$\tau=\lambda/(4k)$ and
\[
 \widehat c_i=c_i+r,\quad \widehat A_i=(A_i+r)(1+q),\quad
 \widehat\rho_i=\rho_i+r,\quad\widehat u_i=u_i+q(A_i+r).
\]
Then $\widehat c_{i+1}=\widehat c_i+\widehat A_i-
\widehat u_i-\widehat\rho_i$, with
$0\le\widehat\rho_i\le\min(\widehat c_i,\widehat A_i)$ and
$\widehat u_i\ge0$. The identity
\[
 \rho(c+a-u-\rho)=ca-(c-\rho)(a-\rho)-\rho u
 \tag{RD.4}
\]
gives $\widehat\rho_i\widehat c_{i+1}\le
\widehat A_i\widehat c_i$. If the first inequality in (RD.3) fails,
then $\widehat u_i>2\tau\max(\widehat c_i,\widehat A_i)$,
and the bound improves by $1-\tau$: split at
$\widehat\rho_i=\min(\widehat c_i,\widehat A_i)/2$ in (RD.4).
If the second fails, the factor $n_i+r$ instead improves by
$1-\lambda/k+\eta\le e^{-\tau}$.
Telescoping the quadratic inequalities therefore gives
\[
 (\ell_0+r)\prod_i(\rho_i+r)(n_i+r)
 \le\theta A\prod_iA_iC_i,\qquad
 \theta^{r-1}\le
 \exp\!\left(\frac{(k-2)(4s+1)}{k^2}
             -\frac{(k-2)\lambda}{4k}\right)
 \le e^{2-\lambda/5}.
 \tag{RD.5}
\]
The factors lost in regularizing are at most $e^\eta$ initially and
$e^{2\eta}$ per interior degree.

One signed Pieri path to any dominant weight proves
$w_r(\ell)\le\prod_tS_r(|\ell_t|)$: add positive coordinates from
top to bottom, then subtract negative ones from bottom to top, and
apply the total path-dimension identity. The bounds
$y^{r-1}/(r-1)!\le S_r(y)\le(y+r)^{r-1}/(r-1)!$ and (RD.5)
give $w_r(\ell)\le e^{2-\lambda/5}\mathcal D(A)$ for each bad pair.
There are at most $(8a)^{4(N_1+1)^2}$ pairs: every state is determined
by its first $s+1$ and last $s$ coordinates, all bounded by
$A+(2s+1)\lfloor2a/p\rfloor\le2a$.
Even retaining the source's extra envelope
$(8a+k)^{(2N_1+2)^2}$, the total bad weight is at most
\[
 e^{2-\lambda/5+64N_1^2\log a}\mathcal D(A)
 \le e^{-\lambda/8}\mathcal D(A),\qquad C_0\ge880.
 \tag{RD.6}
\]
Thus the kept pairs retain at least $1-e^{-\lambda/8}$ of (RD.2).

\paragraph{Why these pairs give the actual filtered blocks.}
The coarse alternative SW10 says
$A_i(1+\delta)<A_*$ before the split $\sigma$, and
$A_*<A_i(1-\delta)$ after it. From
$c_i\le A+\sum_{l<i}A_l$, the late case gives $c_i\le A_i$ and
\[
                 O_i\le c_i+\lambda A_i/k\le A_i.
\]
In the early case, induction in (RD.3) gives
\[
 c_i\ge A_*(1-e),\quad
 e=4N_1\lambda/k+12N_1/p^2\le9\delta/80,
 \quad e(1+\delta)\le\delta.
\]
Hence $A_i\le c_i$, and
$O_i\le A_i(1+4/p^2)+\lambda c_i/k$,
$S_i\le2\lambda A_*/k$.
Finally, terminal interlacing telescopes to
\[
 D:=\sum\ell-\sum K^{2s+1}
 \le c_s+\sum_iC_i-\sum_i n_i
 \le c_s+(\lambda/k)\sum_iC_i.
\]

For $0\le j_+\le J:=\lfloor\lambda B/k\rfloor$, define
\[
\begin{array}{c|ccc}
 &o_i&w_i&\mathrm{sg}_i\\ \hline
 i<\sigma&0&W_i&S_i\\
 i\ge\sigma&O_i&0&0
\end{array}
 \qquad o_+=D+j_+.
\]
Keep the original digits, choose the late levels $j_i=S_i$, and use
terminal level $j_+$. The preceding inequalities are exactly the
source kept-block inequalities; the terminal one follows also from
$C_i\le2a/P_i$. Moreover
\[
 \mathrm{cap}_l=A_*+\sum_{i<l}(A_i-O_i)=M+c_l.
 \tag{RD.7}
\]
The actual level chain is $K^{n+1}$, followed by $\ell$.
Its degrees are $W_i,C_i-S_i,D$, exactly the level-adjusted block
degrees. Levels are legal because $S_i\le O_i,C_i$ and $j_+\le D+j_+,B$.
At each tested prefix its first coordinate is at most $c_i$, or $c_s$
terminally; (RD.7) and $M\ge h$ give the required
$\lfloor\mathrm{cap}-h\rfloor$ cutoff.
The map is injective: the tail chain and its fixed degree sums recover
every omitted first coordinate; the terminal state recovers $\ell$,
and its level recovers $j_+$. Its weight is
$w_r(\ell)S_r(B-j_+)$. This proves
\[
 \mathrm{chainCount}(h)\ge
 (1-e^{-\lambda/8})\mathcal D(A_*-M)
                  \sum_{j=0}^{J}S_r(B-j).
 \tag{RD.8}
\]
All sums are finite: the kept inequalities bound every block order,
then every level and chain coordinate.

\paragraph{The remaining factors and the half-density window.}
The hockey-stick identity and the product formula for binomials give
\[
\begin{aligned}
 \sum_{j=0}^{J}S_r(B-j)&=S_k(B)-S_k(B-J-1)
                    \ge(1-e^{-\lambda/3})S_k(B),\\
 S_k(A_*-M)&\ge(1-M/A_*)^rS_k(A_*)
                    \ge(1-rM/A_*)S_k(A_*).
\end{aligned}
\]
For the first bound use $J+1>\lambda B/k$ and
$k(B+r)\le3rB$, valid since $r\ge4$ and $B\ge k$.
Applying these to (RD.8), and the actual rank identity RU1--RU2,
proves (RD.1). Its last factor may be negative; that binomial comparison
is used last, with a nonnegative coefficient.

In the source window $0\le h\le b'/(kQY)$,
$MkQ\le2b'$ and $A_*Q\ge a/2\ge60b'$ give
$rM/A_*\le1/30$ and the required $M\le\epsilon A_*$.
Since $\lambda\ge24$, the two exponential factors are at least
$3/4$ and $8/9$. Together with RU8,
\[
       \frac12\,\mathrm{mmB}
       \le\mathrm{rankOnDigits}(h)\le\mathrm{mmB},
       \qquad 0\le h\le\frac{b'}{kQY}.
 \tag{RD.9}
\]
Indeed the proved lower factor is at least $29/45>1/2$.
This is the actual one-segment comparison, including its digit and
rounding conventions. The expanded report proves each intervening
inequality and identifies the exact source declarations.

\subsection{RD11--RD15: harmonic density and the final window mean}
\label{sec:rank-density-mean}

\status{The actual density and column-box numerical records are proved
with the source thresholds. A finite telescoping argument strengthens
the harmonic density and shortens its proof. With HM7 and the completed
local mean and margin bounds, it gives the literal final exceptional
mean $[6(C_A+C_s)+8]NN_1/Y$.}

\paragraph{A finite harmonic calculation.}
In the HM1 convention put $P(t)=\binom{t+r}{r}$, $P(-1)=0$, and
$H(\alpha,\beta)=P(\alpha)P(\beta)-P(\alpha-1)P(\beta-1)$.
Let $R_j=P(\alpha)P(\beta)/[P(\alpha+j)P(\beta+j)]$.
Monotonicity of $t/(t+r)$ gives the exact comparison
\[
 \frac{H(\alpha,\beta)}{P(\alpha+j)P(\beta+j)}
 =R_j\left(1-\frac{\alpha\beta}{(\alpha+r)(\beta+r)}\right)
 \ge R_j-R_{j+1}.
 \tag{RD.10}
\]
Thus its sum for $0\le j\le J$ is at least $1-R_{J+1}$.
Writing $L=J+1$ and discarding the first polynomial ratio,
\[
 R_L\le\prod_{t=1}^{L}\left(1-\frac r{\beta+r+t}\right)
 \le e^{-rL/(\beta+r+L)}\le e^{-\lambda'/3}
 \tag{RD.11}
\]
whenever $r\ge2$, $\beta+L+r\le2b'$ and
$L\ge\lambda'b'/(r+1)$. This improves the source lower density by
removing its extra factor $1-rJ_0/\alpha$; it works also at zero degrees.

\paragraph{The literal enlarged windows.}
Use the final data $X.\mathrm{Hyp}(K.C_0,C_H)$ and assume
$C_H\ge\max(T_{\rm LP},400)$, where NP4 proves the base estimates and
\[
 T_{\rm LP}=\max\{108,8K.C_A+93,1,K.C_0^{-2}\}.
\]
Write $a=\lceil a_{\mathbb R}\rceil$,
$\lambda_0=\mathrm{LaI}(a_{\mathbb R})$,
$\lambda' =\mathrm{LaI}(a)$, $Z=\lambda_0b'/k$,
$J_0=\lambda'b'/k$, $J=\lfloor J_0\rfloor$, $c=c_{\rm cl}$,
and $e=J+c$. The source windows are
\[
\begin{aligned}
 W_a(\tau)&=[a+\tau,\lfloor a_{\mathbb R}+a_{\mathbb R}/(kY)\rfloor+\tau],\\
 W_b(\tau)&=[\lceil b_l\rceil+\tau,\lfloor b_u\rfloor+\tau],
\end{aligned}
\]
intersected with the integers. Their enlarged versions $W_a^+,W_b^+$
raise only the upper endpoints by $e$.
The literal hypotheses and NP4 rounding give
\[
\begin{gathered}
 \lambda_0\le\lambda'\le2\lambda_0,\quad
 \lambda_0\ge8\log(NY),\quad c\le Z,\quad J_0\le2Z,\\
 b'\ge3200k,\quad800Z\le b',\quad Z\ge1,\qquad
 w_a,w_b\ge Y_bNZ/2>0,
\end{gathered}
\]
where $w_a=|W_a(0)|$, $w_b=|W_b(0)|$ and $Y_b\ge400$.
For the window bounds, rounding loses at most one from the real widths;
those widths are at least $Y_bN\lambda_0b'$ and $Y_bNZ$.
For $\tau\le c$ and $\beta\in W_b(\tau)$,
$\beta+J+1+r\le3b'/2+3Z+k<2b'$, so (RD.11) applies uniformly.
Choose $d=1-e^{-\lambda'/3}$. Then
\[
 Ne^{-\lambda'/3}\le\frac13,
 \qquad \frac{Ne}{w_a},\frac{Ne}{w_b}\le\frac3{200},
 \qquad d^N\ge\frac23.
\]
The first inequality follows by cubing:
$N^3e^{-\lambda'}\le N^{-5}\le1/32<1/27$.
Bernoulli also gives
$ (1+u)^N\le(1-Nu)^{-1}$ for $Nu<1$.
Consequently
\[
 |W_a^+|^N|W_b^+|^N
 \le(200/197)^2w_a^Nw_b^N
 \le2d^Nw_a^Nw_b^N.
 \tag{RD.12}
\]
Together with (RD.10)--(RD.11), this proves every field of the actual
$\mathrm{DensityNumericalFacts}$. HM7 therefore gives the two translated
averages needed below.

\paragraph{Boundary loss and the exceptional mean.}
Already $C_H\ge T_{\rm LP}$ implies
$w_a,w_b\ge2YNc$ and $12Nkc/b'\le1/Y$.
Indeed the size clause gives $C_Hk^3NYc\le b'$,
while the window widths satisfy
$w_a\ge216b'-1$ and $w_b\ge216b'/k-1$; also $b'\ge k$.
Bernoulli now bounds the fraction below a shifted lower endpoint by
\[
 1-\left(1-\frac c{w_a}\right)^N
       \left(1-\frac c{w_b}\right)^N\le\frac1Y.
 \tag{RD.13}
\]
These are exactly the two nontrivial fields of
$\mathrm{ColumnSpanNumericalFacts}$.

SW8--SW9, RU8 and NP3--NP6 give each actual segment mean at most
$C_A(s_j+1)N_1/Y$; NP9 gives $\kappa_P\le C_sNN_1/Y$.
The sufficient fixed coefficients are
\[
 C_A\ge\max\{56,2(31e^3+5)e^3,1\},\qquad
 C_s\ge C_\kappa(C_m+4C_{E1}+C_{E2}+15).
\]
Each assignment on the $s_j+1$ distinct segment coordinates has
$|W_a^+|^{N-s_j-1}$ extensions. Finite marginalization, followed by
$\sum_j(s_j+1)\le N$, therefore proves
$\operatorname{mean}k_x\le(C_A+C_s)NN_1/Y$.
No independence between different segments is needed.

The source $k_{\rm text}$ is the inside term $k_F(A-c,B-c)$
(replaced by $1$ outside), plus $2k_F(A+c,B+c)+12Nkc/b'$.
Subtraction is injective on the inside box; addition is a bijection
onto the translated box. Since $k_F\ge0$, HM7 and (RD.13) give
\[
 \operatorname{mean}k_{\rm text}
 \le6\operatorname{mean}k_x+\frac2Y
 \le[6(C_A+C_s)+8]\frac{NN_1}{Y}.
 \tag{RD.14}
\]
The coefficient is exactly $K.C_{\rm mean}$. The expanded report records
the natural-to-signed window identification and all constant thresholds.
\noindent\textit{Primary source route.}\cite{coarse-rank-lower}


\section{Harmonic models, density, and the actual cloud maps}
\label{sec:harmonic-model}

\status{HM1--HM12 identify the literal polynomial lattices, prove the
bounded harmonic comparison and its density consequence, and construct
the complementary-middle cloud maps and coefficient-ring transport.
The HH rank input is RU11; its numerical supply and the final applications
of the cloud maps are separate. No integral semisimplicity is assumed.}

\paragraph{HM1--HM2: an explicit harmonic quotient.}
Work in the source coefficient setting $w^2=p$, with $v_p(w)=1/2$.
Put $k=m+1$, $S_d=\binom{d+k-1}{k-1}$ and $S_{-1}=0$.
For a $k\times k$ matrix $Z$, let
$c_i(Z)=(-1)^i\det Z_{\widehat i,\{0,\ldots,k-2\}}$.
The actual two-letter map sends $x_i$ to the first-column entry
of one matrix and $y_i$ to $c_i$ of an independent matrix.
Its leading monomials recover every exponent of $x,y$, so it is
injective and saturated over $O$. Identifying the matrices gives
\[
 \beta_{ab}:P_{ab}\longrightarrow H_{(a,0,\ldots,-b)},\qquad
 \ker\beta_{ab}=hP_{a-1,b-1},\quad h=\sum_i x_i y_i,
 \qquad H_{ab}=S_aS_b-S_{a-1}S_{b-1}.                    \tag{HM.1}
\]
Indeed CB6 gives surjectivity, the repeated-column determinant gives
$\beta(h)=0$, and CB3's dimension formula gives equality of kernels.
The cofactor identity
$c(g^{-1}Z)=\det(g)^{-1}g^t c(Z)$ checks the actual determinant twists
for every lattice $gO^k$.

Let $\Delta=\sum_i\partial_{x_i}\partial_{y_i}$.
The commutator $[\Delta,h]F=(a+b+k)F$ gives
\[
 \pi_{AB}=\sum_{\ell=0}^{\min(A,B)}
 \frac{(-1)^\ell}{\binom{A+B+k-2}{\ell}}
 \frac{h^\ell}{\ell!}\frac{\Delta^\ell}{\ell!},\qquad
 E_j=\binom{A+B+k-j-1}{j}^{-1}
 \frac{h^j}{j!}\pi_{A-j,B-j}\frac{\Delta^j}{j!}.        \tag{HM.2}
\]
These project onto $h^j\ker\Delta$ and sum to the identity.
Both divided operators preserve the lattice with basis
$x^\alpha y^\beta/\alpha!$. Legendre's formula gives
$v_p\binom nr\le\lfloor\log_p n\rfloor$; comparison with the
ordinary coefficient lattice costs at most $(k-1)\lfloor\log_p(A+B+k)\rfloor$.
Thus both actual maps $\operatorname{pr}_j,\operatorname{io}_j$
obtained from (HM.2) have denominator at most
\[
 U=\left\lceil(m+2)\bigl(\log_p(8a+m)+1\bigr)\right\rceil,
 \qquad \sum_j\operatorname{io}_j\operatorname{pr}_j=1,
                                                               \tag{HM.3}
\]
when $A+B+1\le8a$. Division by $h^j$ costs nothing: it has a
monic leading term. Equivariance gives the same bound at every lattice.

\paragraph{HM3--HM5: normalization and rank splitting.}
Tensor over $N=|TT|$ disjoint letters and put $u=NU$.
Write $A_H=\mathrm{AHH}$, $A_j=\mathrm{Ar}(\lambda_j)$,
$d_j=A_H-A_j\ge0$. For cleaned vectors $x,\Phi$ the literal maps are
\[
 x_j=p^u w^{d_j}\operatorname{pr}_jx,\qquad
 \Phi_j=p^{2u+d_j}\Phi\circ(w^{-d_j}\operatorname{io}_j),\qquad
 \mathrm{be}_j(x_j,\Phi_j)
   =p^{3u}\mathrm{be}_H(\operatorname{io}_j\operatorname{pr}_jx,\Phi).
                                                               \tag{HM.4}
\]
Applying the lattice bounds to the source, midpoint, target and every
competing edge verifies all cleaned conditions, with unchanged height
and conductor.

If integral matrices satisfy $\sum_jB_j=zI_n$, the total number of
their Smith factors with valuation at most $v_p(z)$ is at least $n$.
Otherwise their small-factor coordinate rows have a common primitive
kernel vector, contradicting $\sum_jB_jx=zx$.
Consequently an HH certificate at depth $D$ splits among the actual
harmonic components at depth $D+3u\le D+4u$.

There are $N_e=(p^{n'k}-1)/(p^{n'}-1)$ outgoing rank-one edges.
Transport by $\det(g)^{-\lambda_{\rm last}}\operatorname{rowAct}(g^{-1})$
preserves the actual lattices and pairing. A certificate at one edge,
with primal family in $\mathrm{CCL}$ and dual family merely in $R'_e$,
therefore gives one at every edge. If $2D<c$, dividing each diagonal
block by its scalar makes the joint pairing matrix the identity modulo
$w$: different-edge pairings lie in $w^cO$. Hence
\[
                         n\le\dim V(\lambda)/N_e.       \tag{HM.5}
\]
This proves the one-sided \texttt{nLeD} as well as its cleaned-dual version.

\paragraph{HM6--HM7: the actual density inequality.}
Use the fixed-weight numerical bounds supplied by NP6, set
$D_1=D_0+4NU$, and apply RU11's HH mass
\[
                         (1-k_x(A,B))\dim HH/N_e.
\]
For $J=\lfloor J_0\rfloor$, decreasing all degrees by at most $J$
lowers the segment-minimum height by at most $NJ_0$.
Thus $t_*(\alpha,\beta)=\lceil t_0(\alpha,\beta)+NJ_0\rceil$
admits every small-shift certificate. Combining (HM.5), the defining
infimum $k_F$, and the telescoping identity
$\sum_j\prod_sH_{A_s-j_s,B_s-j_s}=\prod_sS_{A_s}S_{B_s}$ gives
\[
 \sum_{0\le j_s\le J} k_F(A-j,B-j)
       \prod_s H_{A_s-j_s,B_s-j_s}
 \le k_x(A,B)\prod_sS_{A_s}S_{B_s}.                    \tag{HM.6}
\]
For $R_j=S_\alpha S_\beta/(S_{\alpha+j}S_{\beta+j})$,
monotonicity of $t/(t+k-1)$ gives
\[
 \sum_{j=0}^{J}\frac{H_{\alpha\beta}}
                         {S_{\alpha+j}S_{\beta+j}}
 \ge1-R_{J+1}
 \ge1-\exp\!\left(-\frac{(k-1)(J+1)}{\beta+J+k}\right).\tag{HM.7}
\]
Summing (HM.6) under $(\alpha,\beta,j)\mapsto(\alpha+j,\beta+j,j)$
therefore proves the source \texttt{Avg}: if $\delta$ bounds (HM.7)
uniformly and the extended product box has size at most $2\delta^N$
times the original box, then $\operatorname{mean}k_F\le2\operatorname{mean}k_x$.
RD11--RD15 supply these numerical inequalities.

\paragraph{HM8: actual duality and disjoint tensor factors.}
The transported row-homogeneous CB2 basis $x_i$ has target exponents
$\pi_i$ and orders $o_i=A-\pi_i\ge0$. Thus
\[
 R_e=\bigoplus_iOw^{-o_i}x_i,\qquad
 R'_e=\bigoplus_iOw^{-o_i}p^{-\pi_i}x_i^*,\qquad
 p^Aw^{-2o_i}p^{-\pi_i}=1.                            \tag{HM.8}
\]
The order inequality counts entries in the first edge rows of each
tableau. This proves the exact dual-height and perfect-pairing interfaces.
Splitting letter blocks into $TT$ and its complement gives actual
\texttt{RestTensor}: multiply the disjoint polynomial factors, multiply
their basis evaluations on the dual side, and multiply basis cardinalities.
Transported product bases prove both lattice equalities, including midpoints.

\paragraph{HM9--HM11: complementary pairing and cloud maps.}
For complementary middle shapes $\alpha'_i=M-\alpha_{r-1-i}$, put
$e=\min\alpha$, $\gamma=\alpha-e$, $T=M-e$ and let $J_r$ reverse rows.
Jacobi's complementary-minor identity proves that
\[
 \tau f(Y)=\det(Y)^T f(Y^{-t}J_r)                      \tag{HM.9}
\]
is an integral isomorphism from shape $\alpha'$ to $\gamma$.
IW3--IW5 bound the normalized Fock dual lattice by
$p^\kappa M_\gamma\subset M_\gamma^\#\subset M_\gamma$,
$\kappa=\binom r2\lfloor\log_p(\gamma_0+r-1)\rfloor$.
Take $\kappa=0$ for the one-dimensional case.
Smith reduction of $p^\kappa B_0$, followed by rounding its valuations
upward, gives integral families with pairings $\delta_{ab}p^{\kappa_b}$,
$0\le\kappa_b\le\kappa$.
The pairing $L(f,f')=p^\kappa B_0(f/\det^e,\tau f')$
has simultaneous row covariance $\det(g)^M$.

The actual contraction is coefficientwise
\[
 \Lambda(v,t)(Z)=L_{Y,Y'}\bigl(v(Z\operatorname{diag}(1,Y,1))
                 t(Z\operatorname{diag}(1,Y',1))\bigr).\tag{HM.10}
\]
Lift each middle tableau minor by adjoining row and column zero,
and add the required power of $Z_{00}$. The resulting integral
$s_b,s'_a$ have row-zero degree $2c$ and last-row degree zero.
The highest-line calculation, evaluated at $Z=I$, gives exactly
$\Lambda(Fs_b,s'_a)=\delta_{ab}p^{\kappa_b}H_\zeta F$.
Tensoring produces the source cloud maps with $\#b=\dim\mathrm{Cl}$.
Multiplication by the primitive $H_\zeta$ has exact integral preimage,
by Gauss valuation. The actual dual map is
$g\mapsto p^{Nc}g\circ L_a$; testing against other target edges
costs precisely $2Nc$ powers of $w$.
This proves full \texttt{RestCloudMaps}, with conductor $c_2-2Nc$.

\paragraph{HM12: coefficient-ring independence.}
Hensel-lift a common residue generator. The powers
\[
                          1,\theta,\ldots,\theta^{n'-1}
\]
form a $\mathbb Z_p$-basis of each unramified $W$.
A common quadratic residue lift $\eta$ gives the basis
$1,\eta,w,\eta w$ of $O$ over $W$; successive reduction modulo $w$
proves independence, and completeness proves spanning.
Matching these bases preserves their defining equations and sends $w$ to $w$.
Uniqueness of the Hensel lifts makes the isomorphism commute with Frobenius.
Coefficient transport then preserves every actual polynomial lattice,
dual, normalized pairing and integral rank certificate. Thus the literal
\texttt{LPprime} statement is independent of the chosen RS setting.
\noindent\textit{Primary source route.}\cite{harmonic-cloud-actual}


\section{Uniform parameters and the actual Type-I height}
\label{sec:numerical-parameters}

\status{NP1--NP9 supply the literal standard, segment, rounded-weight and
projection numerical records, including the shifted comparison with $t_I$.
These finite inequalities do not assert the existence of the global
arithmetic data to which the local theorem is to be applied.}

\paragraph{NP1--NP2: a uniform standard package.}
Write $k=m+1$, $N=|TT|$, $Q=p^{G'}$, $\lambda=C_0N_1^2\log a$,
$\epsilon=N_1\lambda/k$, and
\[
 E_0=C_{E1}\left(\frac a{p^2}+\frac {b'}p+
              \frac{\lambda N_1a}{pk}+\frac{\lambda b'}k\right)
       +C_{E2}k^2N_1Q^2c.
\]
For $H\ge1$, $C_c\ge0$, and $C_0,C_{E1},C_{E2}\ge1$, set
\[
 T_2=1000(C_{E1}+C_{E2}+1)(C_c+1)H.                 \tag{NP.1}
\]
The literal $\mathrm{D1num}$, extended-window hypothesis $\mathrm{HIext}$,
and $\mathrm{CBound}(C_c)$ imply
\[
                  \mathrm{Std}(C_0,C_\delta,C_{E1},C_{E2},H,X)
\]
whenever
each of $C_{Yb},C_\delta,C_k,C_D,C_b$ is at least $T_2$.
The report lists every field of these records. The substantive estimates
are as follows.

Put $F=c_*+Nc_{\rm cl}+n'c_{\rm cl}$. The depth hypothesis gives
$Nk\log_pa\le D_{\rm LP}$, whence $c\le2C_cF$.
The size hypothesis and $C_b\ge2C_cH,H$ give
\[
 Hk^3NQ^4cY\le b',\qquad Hk^3NQ^4Y\le b'.        \tag{NP.2}
\]
Set $A=N_1\lambda/k^2$. Each of
\[
 a/p^2,\quad b'/p,\quad\lambda N_1a/(pk),\quad
 \lambda b'/k,\quad k^2N_1Q^2c,\quad Q
\]
is at most $aA$: use $p\ge k^2$, $kb'\le a$,
$k^3Q^2c\le b'$, $k^3Q\le b'$, and $N_1,\lambda\ge1$.
Since $k\ge4C_{E1}+C_{E2}+1$, this proves
\[
                         (E_0+Q)/a\le\epsilon.    \tag{NP.3}
\]
If $\eta=(b_h-b_l)/b'$, then
$\delta=C_\delta\epsilon+2\eta$ and
$\beta=\eta+2\lambda/k$. Thus
$40\epsilon\le\delta<1/6$ and
$0\le\beta\le\delta/2-18\epsilon$.
Finally $2J_0=2\lambda b'/k\le a/(kY)$ enlarges the original windows
to exactly the standard windows. All other fields follow directly or
by removing factors at least one.
The quantifier order is preserved: choose $C_c$, then $C_0\ge1$,
then $C_{E1},C_{E2},C_m\ge1$, and only then the five constants
bounded below by (NP.1). In particular $H=4$, needed by RD, is uniform.

\paragraph{NP3: the actual segment satisfies the size conditions.}
For a segment put $r=m$, $s\le N_1$, and $P_i=p^{G-d_i}$.
The extended hypotheses give the source record $\mathrm{SegFacts}$:
\[
 \begin{gathered}
 p\ge(r+1)^2Y^2,\quad 2s+2\le r,\quad
 p\le P_i\le Q_j\le Q,\quad (r+1)P_i\le b',\\
 4\lambda<r+1,\quad (r+1)b'\le a/4,\quad
 b_T\le\tfrac52b',\quad 0\le x_a\le2J_0.
 \end{gathered}                                      \tag{NP.4}
\]
Here $C_{Yb}\ge2,C_b\ge1,C_D,C_{E1},C_{E2}\ge0$ suffice.
Indeed $p\ge k^2\ge4k$ gives $kb'\le a/4$, while the size
inequality gives $kQ\le b'$ and $2Y_bN\lambda<k$ gives $4\lambda<k$.
For the multiple-tie model of TP, a nonterminal letter degree $d_l$
satisfies $d_lP_i\le2a$ and the terminal degree is at most $2a$.
Consequently all letter and tie degrees are at most $2a$ and the
sum of letter degrees is at most $4a\le8a$.
The two bounds
\[
 S_iP_i\le\tfrac52b'+(2s+1)2a/p,\qquad
 n_{0i}P_i\ge a-kP_i-\tfrac52b'
\]
give $S_i\le n_{0i}$ because $p\ge4s+8$ and $pb'<a$.
They apply when all tie coordinates are simultaneously at their
compatible endpoints, by SW6.
Under $C_0\ge1$ and the NP2 thresholds, the remaining bounds in
SW5--SW7 follow from $kY^2b'\le a/2$, $p\ge k^2Y^2$,
$a/(kY)\le|W|\le a+1$, and SW12.
Thus TP11 applies to the actual model with $C_M=1$, including its
logarithmic degree hypotheses; SW5--SW9 need no additional numerical
assumption. The coarse rank input is supplied separately by RU.

\paragraph{NP4: rounding the radius without changing the hypotheses.}
For the final Type-I statement distinguish the real radius $a_r$ from
$a=\lceil a_r\rceil$, and put $L=\log_pa_r$.
For a fixed positive constant tuple $K$, define
\[
 \begin{split}
 T_{\rm LP}&=\max\{108,\,8C_A+93,\,1,\,C_0^{-2}\},\\
 T_M&=\max\{T_{\rm LP},\,4K.CH+1,\,2C_m+24\}.
 \end{split}                                         \tag{NP.5}
\]
Assume the literal $\mathrm{HypI}(C_0,C_H)$, with $C_H\ge T_M$.
Its elementary consequences are $k\ge5,p\ge25,b'\ge8,a_r\ge200$,
$N,N_1^2\le k$, and $L,\log a_r,\lambda(a_r)\ge1$.
Since $a_r\le a<a_r+1$,
\[
 0\le\log a-\log a_r\le1/200,\quad
 \lambda(a)\le2\lambda(a_r),\quad
 \lambda(a)\log a\le2\lambda(a_r)\log a_r.          \tag{NP.6}
\]
Replace $a_r$ by $a$, $Y_b$ by $Y_b/2$, and the constant $C_H$ by
$K.CH$. The resulting data still satisfy $\mathrm{HypI}$.
For the delicate clause, write $\delta_r,\delta_a$ for the
old and new excluded-band widths. Then
$\delta_a+1/k\le\delta_r$, while $k\le a_r$ and $\delta_a<1/6$.
These imply
\[
 [a(1-2\delta_a),a(1+2\delta_a)]
 \subseteq[a_r(1-2\delta_r),a_r(1+2\delta_r)].        \tag{NP.7}
\]
The lower endpoint uses monotonicity. For the upper endpoint,
the rounding error is at most $1+2\delta_a<4/3$, whereas the
additional allowance is $2a_r/k\ge2$.
The other clauses use (NP.6), with factors two or four accounted for
by (NP.5). In particular the forbidden odd powers remain forbidden.

\paragraph{NP5--NP6: depths, windows and multiplicity.}
Put $\ell=1+\log_p(8a+m)$, $u_0=\lceil(k+1)\ell\rceil$,
$\kappa_1=\lceil(k-1)^2\ell\rceil$, $n_s=\mathrm{numSeg}(TT)$,
and $v=\log_pa$. The source depths are
\[
 \begin{split}
 D_0&=n_s(C_AY^2N_1^2k^2\ell+4Nkv)+8Nkv,\\
 D_1&=D_0+4Nu_0.
 \end{split}
\]
Here $\ell\le3L$, $v\le2L$. With $P=NY^2k^3L$, the four contributions
to $D_1$ are bounded by $3C_AP,8P,16P,16P$, and
$N\kappa_1\le4P$. Hence
\[
 2D_1+N\kappa_1\le(6C_A+84)P\le D_p,\qquad
 4D_1<c_4:=c_*+2Nc_{\rm cl}+2n'c_{\rm cl}.           \tag{NP.8}
\]
The strict inequality uses $8(D_p+1)\le c_*$.
Thus $4D_0<c_4$ and $2D_1<c_4$ as required.

For $J_0=\lambda(a)b'/k$, $J=\lfloor J_0\rfloor$, the hypotheses give
$3\le J\le J_0\le b'/216$ and $c_{\rm cl}\le J_0$.
Let $u_a=\lfloor a_r+a_r/(kY)\rfloor$, $u_b=\lfloor b_u\rfloor$,
$l_b=\lceil b_l\rceil$. Then
\[
 \begin{gathered}
 J\le a,l_b,\qquad
 a+a/(kY)\le u_a+J+c_{\rm cl}\le a+a/(kY)+2J_0,\\
 u_b+J+c_{\rm cl}\le b_u+2J_0,\qquad
 (u_a+J+c_{\rm cl})+(u_b+J+c_{\rm cl})+1\le8a.
 \end{gathered}                                      \tag{NP.9}
\]
Each follows by substituting the floor/ceiling bounds; the report gives
the inequalities explicitly. Moreover $c_4\le2(c_*+Nc_{\rm cl}+n'c_{\rm cl})$
gives $\mathrm{CBoundI}(4)$, and
\[
 m_0=C_mkN_1Qc_4+12c_4\le b'/(kQY).                 \tag{NP.10}
\]
For (NP.10), multiply by $kQY$ and dominate the result by
$(2C_m+24)k^2N_1Q^2Y(c_4/2)$; the final size hypothesis applies.
Equations (NP.6)--(NP.10) supply exactly both
$\mathrm{FixedWeightNumericalFacts}$ and $\mathrm{MeanKappaNumericalFacts}$,
uniformly in all data with $C_H\ge T_M$.

\paragraph{NP7--NP8: the true final height.}
Every odd marked point $y$ belongs to exactly one open cyclic interval
between consecutive even marked points. Its segment start is
$y_{\rm prev}$ and its segment distance is $d_{\rm prev}$.
This is a bijection, including the interval crossing zero, so the
segment order sum equals
\[
 S=\sum_{y\ {\rm odd}}\min\{b_y,a_{y_{\rm prev}}/p^{d_{\rm prev}}\}.
\]
For $c\ge0,P\ge1$,
$\min(b+c,(a+c)/P)\le\min(b,a/P)+c$.
Thus shifting both degree families by $c_{\rm cl}$ changes $S$ to
$S'\le S+Nc_{\rm cl}$.
The actual integer height is $t_*=\lceil t_0+NJ_0\rceil$,
so it costs at most one more than $t_0+NJ_0$.
In comparing the local error terms at $a$ with the final ones at $a_r$,
their respective factors are $2,1,4,2$; the extra $NJ_0$ uses another
factor two on $N\lambda(a_r)b'/k$.
The allowance $CNkQ$ absorbs the rounding unit, since $NkQ\ge1$,
and $CNc_{\rm cl}$ absorbs $S'-S$. Therefore
\[
 C\ge4C_t+4\quad\Longrightarrow\quad
 t_*(a_w+c_{\rm cl},b_w+c_{\rm cl})+2n'c_{\rm cl}
       \le t_I(C_0,C;a_w,b_w).                      \tag{NP.11}
\]
No floor or additional shift has been suppressed.
Finally the size inequality also gives
$12Nkc_{\rm cl}/b'\le1/9$, $2c_{\rm cl}\le8a$,
and $b'/3\le a,b_l$. Together with (NP.8), these are precisely
$\mathrm{FactsLP}(K,C)$, supplied with thresholds
$(C_H,C)\ge(T_{\rm LP},4C_t+4)$.
All constant choices precede the variable data. Combining several such
packages only takes a maximum of their thresholds, and therefore does
not reverse any of the source quantifiers.

\paragraph{NP9: the explicit coarse loss constant.}
The actual loss is
\[
 \kappa_P=C_\kappa N\left(e^{-\lambda/8}
       +\frac{kQ(m_0+N_1k)}{b'}+\frac1Y+\frac{kE_0}a\right).
\]
Under NP2, set $\alpha=a/Y$ and $\beta=b'/Y$.
The standard size bounds give
$k^3Q^2c\le\beta$ and $k^3Q\le\beta$, while
$\lambda\le k$, $kb'\le a$, and $\lambda Yb'\le a$.
Since $\lambda\ge8\log(NY)$, the exponential term is at most $N_1/Y$.
Expanding $m_0=C_mkN_1Qc+12c$ bounds the second term by
$(C_m+13)N_1/Y$.
After multiplication by $k$, the five summands in $E_0$ are bounded
by $\alpha,\alpha,N_1\alpha,\alpha,N_1\alpha$, respectively.
Thus
\[
 \boxed{\quad\kappa_P\le C_s\,\frac{NN_1}{Y},\qquad
 C_s=C_\kappa(C_m+4C_{E1}+C_{E2}+15).\quad}          \tag{NP.12}
\]
The constant is chosen after $C_\kappa$ but before any variable data,
as required by the source's $\mathrm{FamilyFourthBound}$.


\noindent\textit{Primary source route.}\cite{type-one-numerical}


\section{The complete local Type-I theorem}
\label{sec:type-one-closure}

\status{TC1--TC6 complete the source's supplied-data
\texttt{Conclusion2}. The proof uses the actual maps, ranks and
numerical bounds reconstructed above. In particular, no unproved
\texttt{AdmB} or \texttt{AdmT} package is assumed. The existence of
the simultaneous global arithmetic data remains separate.}

\paragraph{TC1: the cloud count is a finite matrix argument.}
Write $G=(p^{n'k}-1)/(p^{n'}-1)$ and
$d(\lambda)=\dim V(\lambda)/G$. HM11 constructs the actual maps
\[
 i:V(H)\longrightarrow V(H^+),\quad
 U_b:V(H)\longrightarrow V(\lambda),\quad
 L_a:V(\lambda)\longrightarrow V(H^+)
\]
with
\[
 L_aU_b=\delta_{ab}p^{\kappa_b}i,\qquad
 0\le\kappa_b\le K,\qquad \#\{b\}=B.                    \tag{TC.1}
\]
The exact midpoint-lattice preimage makes $i$ a primitive
embedding. Extending an integral basis gives a retraction that
preserves the midpoint lattices. Perfect normalized duality
transfers its dual with the factor $p^{A_H-A_{H^+}}$.

Take $n_1,n_2$ good factors for $H,H^+$ at depth $D$, and
absorb the integral column operations so their diagonal scalars
are $z_1,z_2$, of valuation at most $D$. Smith-reduce the cross
pairing matrix $N$ between the image of the first family and
the second dual family. Let $r$ count its invariant factors
dividing $z=z_1z_2$, and put $s=n_1-r$. On the remaining rows
the cross matrix is $zE$ with $E\in wO$.

The combined pairing on these $s$ rows and the second family is
\[
 M_0=\begin{pmatrix}z_1I_s&zE\\ B_0&z_2I_{n_2}\end{pmatrix}.
\]
Both $U=I-EB_0$ and $V=I-B_0E$ are invertible over $O$, and
\[
 M_0
 \begin{pmatrix}
 z_2U^{-1}&-zEV^{-1}\\ -B_0U^{-1}&z_1V^{-1}
 \end{pmatrix}
       =zI.                                             \tag{TC.2}
\]
These are $s+n_2$ factors at depth $2D$ with a cleaned source
and the full midpoint dual. HM5's one-sided upper bound, using
$4D<c_2$, gives $s+n_2\le d(H^+)$. Thus
$r\ge n_1+n_2-d(H^+)$.
The good rows of $N$, followed by (TC.1), give
\[
 Br\ \hbox{factors for }\lambda,\qquad
 r\ge n_1+n_2-d(H^+),\qquad
 \hbox{depth }2D+K.                                    \tag{TC.3}
\]
The normalized pullbacks through $L_a$ have the conductor
$c_2-2Nc$ proved in HM11. Multiplying each by
$p^{K-\kappa_a}$ makes the common scalar exactly $p^Kz$.
This proves the source's cloud count without an additional
intersection-dimension assumption.

\paragraph{TC2: compare the dimensions factor by factor.}
At each selected position let
\[
 \lambda_s=(A_s+c,\beta_{s,1},\ldots,\beta_{s,k-2},-B_s-c),
 \qquad |\beta_{s,i}|\le c,\qquad A_s,B_s\ge b'/3.
\]
Let $H$ have extremes $(A_s,-B_s)$ and zero middle weights;
$H^+$ has extremes $(A_s+2c,-B_s-2c)$. Let $\mathrm{Cl}$ be
the product of the middle Weyl dimensions. The Weyl product CB3
gives
\[
 \mathrm{Cl}\dim H\le\dim\lambda\le E\mathrm{Cl}\dim H,
 \qquad \dim H^+\le E\dim H,\qquad
 E=(1+6c/b')^{N(2k-3)}.                                \tag{TC.4}
\]
Indeed each remaining numerator increases by between zero and
$2c$, or by $4c$ in the extreme-to-extreme factor for $H^+$.
Its corresponding denominator is at least $A_s$, $B_s$ or
$A_s+B_s$.
For $\varepsilon=12Nkc/b'\le1/9$,
$E\le e^\varepsilon$, $(1-\varepsilon)E\le1$ and $E\le2$.
Consequently TC1 converts flat retained fractions
$1-\kappa_1,1-\kappa_2$ into the general fraction
\[
 1-\kappa_1-2\kappa_2-\varepsilon,\qquad
 \hbox{depth }2D+N\lceil\kappa_{\rm one}\rceil,\qquad
 \hbox{conductor }c_2-2Nc.                             \tag{TC.5}
\]
Here $0\le\kappa_1\le1$, $\kappa_2\ge0$.
The sign check is essential:
$(1-\kappa_1)\mathrm{Cl}\,d(H)
\ge(1-\kappa_1)(1-\varepsilon)d(\lambda)$.

\paragraph{TC3: restore the weights outside the selected set.}
Write $\lambda=\lambda_T+\lambda_O$ on disjoint letter sets.
Multiplication in disjoint polynomial variables is an integral
tensor isomorphism for every lattice family. The zero-weight
factors are constants, so dimensions multiply and the normalized
pairing is the tensor product. If all coordinates of $\lambda_O$
are bounded in absolute value by $c$, its first-minus-last weight
sum is at most $h=2n'c$. In an elementary-divisor frame, the
tableau row degrees and the signed determinant twist give
\[
 R_{\lambda_O,e}\subseteq w^{-h}H_{\lambda_O,\rm src},
 \qquad R'_{\lambda_O,e}
       \subseteq w^{-h}H_{\lambda_O,\rm tgt}^{*}.        \tag{TC.6}
\]
Tensor with the exact dual midpoint bases of HM8. The height
increases by $h$, the conductor decreases by $h$, and the depth
is unchanged. Rank multiplies by $\dim V(\lambda_O)$.
Thus the retained \emph{fraction} is unchanged. This proves
the source's \texttt{ResidueClassCount}.

\paragraph{TC4: the edge factors give the actual integral rank.}
There are exactly $G$ incoming rank-one edges at any target
lattice. Choose the local factors on each. Their combined pairing
matrix has diagonal blocks $z_eI$, with $z_e\mid p^D$, and
off-diagonal entries in $w^cO$. After multiplying column $e$
by $p^D/z_e$, the matrix is
\[
                         p^D(I+E),\qquad E\in wO,      \tag{TC.7}
\]
provided $c>2D$. Its remaining inverse is integral. Summing
the local ranks therefore gives at least
$(1-\kappa)\dim V(\lambda)$ factors dividing $p^D$.
The vectors retain their individual other-edge vanishing.
This is the full matrix certificate \texttt{GoodRankK}, hence
\texttt{LPprime}, with the source's stronger conductor hypotheses
$8D<c$ and $4(D+1)\le c$ retained.

\paragraph{TC5: attain the infimum and apply the two shifts.}
Let $\kappa_F(A,B)$ be the least nonnegative flat deficit at
height $t_*(A,B)$, conductor $c_4$ and depth $D_1$.
It lies in $[0,1]$, and its infimum is admissible.
To see this at any edge, put $T=(1-\kappa_F)d(H)$.
If $T>0$, a sufficiently close admissible deficit forces its
integer rank to exceed $\lceil T\rceil-1$, hence to be at least
$T$; if $T\le0$, rank zero suffices.

For $c=c_{\rm cl}$ the actual final deficit is
\[
 \begin{split}
 \kappa_{\rm text}(a_w,b_w)
  ={}&\begin{cases}
      \kappa_F(a_w-c,b_w-c),&
          a_w-c\ge a,\ b_w-c\ge b_l,\\
      1,&\text{otherwise},
      \end{cases}\\
     &+2\kappa_F(a_w+c,b_w+c)+12Nkc/b'.
 \end{split}                                           \tag{TC.8}
\]
All inequalities and shifts are coordinatewise. If this deficit
is at least one, use the zero certificate. Otherwise the first
branch holds. Enlarge its height to $t_*(a_w+c,b_w+c)$ and
apply TC2 to the two flat weights. Apply TC3 next. The resulting
conductor is exactly
\[
 c_4-2Nc-2n'c=c_*.
\]
NP5 and NP7 give respectively
\[
 2D_1+N\lceil\kappa_{\rm one}\rceil\le D_p,\qquad
 t_*(a_w+c,b_w+c)+2n'c\le t_I.                         \tag{TC.9}
\]
TC4 now proves \texttt{MidLP} for every requested integer
height $t\ge t_I$.

\paragraph{TC6: a completely explicit local theorem.}
The following sufficient constants keep all previous estimates:
\[
 \begin{gathered}
 C_0\ge880,\quad C_{E1}=C_{E2}=C_m=17,\quad
 C_\kappa=3,\quad C_t=30,\quad C_c=4,\\
 C_A=50000,\qquad C_s=351,\qquad C_{H,\rm inner}=700000.
 \end{gathered}                                         \tag{TC.10}
\]
NP2 gives the last threshold as $1000\cdot35\cdot5\cdot4$.
The SW mean coefficient is at most
$2(31e^3+5)e^3<45468$ using $e<3$, and NP9 gives
$C_s=3(17+4\cdot17+17+15)=351$.

There is no remaining pointwise good-share assumption.
RU10 makes each segment lattice finite free; the ranks of its
saturated good sublattices form a nonempty finite set. Choose
one of maximal rank. The literal definition of $\kappa_{TI}$
then supplies its good share, including the rank-zero case.
Moreover
\[
 \sum_jD_j+4\kappa_0=D_0,\qquad
 \kappa_P+\sum_j\kappa_{TI,j}=\kappa_x.                 \tag{TC.11}
\]
RU11 therefore supplies exactly HM6's aggregate certificate.
HM6--HM7 and RD11--RD15 give
\[
 \operatorname{mean}\kappa_{\rm text}
     \le \{6(C_A+C_s)+8\}\frac{NN_1}{Y}
     =302114\,\frac{NN_1}{Y}.                          \tag{TC.12}
\]
The final numerical thresholds are
$T_{\rm LP}=400093$, $T_M=2800001$ and
$T_{\rm density}=400093$; the height coefficient only needs
$C\ge124$. Thus we have proved:

\begin{quote}
\textbf{Local Type-I theorem.}
For $C_0\ge880$, $C_H\ge2800001$ and $C=302114$,
the source's literal \texttt{HypI} implies
\texttt{Conclusion2}.
There is a nonnegative exceptional function of the degree
families, with mean at most $CNN_1/Y$, which works for every
allowed coefficient-ring setting and every dominant
\texttt{WeightI}. At every target lattice and every integer
$t\ge t_I$, it gives at least
$(1-\kappa)\dim V(\lambda)$ integral factors dividing $p^{D_p}$,
with the required other-edge vanishing to order $c_*$.
\end{quote}

For uniformity in the coefficient ring, choose one setting and
its $\kappa_{\rm text}$; HM12 transports its \texttt{LPprime}
certificate to every other setting. Equality of two independently
computed exceptional functions is unnecessary. If no setting
exists, use $\kappa=0$ and the conditional assertion is vacuous.
All constants precede the local data. The expanded companion
states every field of \texttt{HypI} and the exact $t_I$ formula.

Together with EA9, this completes both local Type-I and Type-II
constructions. It does not supply the global abundance,
theta/spectral identities, metric comparison, cohomological
windows or simultaneous prime choices still needed for RH
and Dirichlet GRH.
\noindent\textit{Primary source route.}\cite{type-one-complete}


\section{Good-orbit counts and uniform Type-II data}
\label{sec:type-two-family}

\status{We prove the actual family Type-II hypotheses and their local
application. Separately, a finite probability argument produces a subset
with small weighted bad-orbit mass. Identifying that subset with the one
used by the arithmetic host remains an explicit construction step.}

Fix $0<\epsilon<1/52$, $\zeta=\epsilon/J>0$, and the nonnegative global
constant tuple $\mathbf C$. The literal parameters are
\[
\begin{gathered}
 \delta=x^{-\epsilon/8},\quad h=2+8\epsilon,\quad
 K_0=\lceil8(h+2)/\epsilon\rceil,\quad r=K_0+1,\\
 u=\lceil\delta^{-1}\rceil,\quad \ell=(6r+1)u,\quad
 P_0=10^{-(6r+1)},\quad M=\lceil(h+3)\log x/P_0\rceil,\\
 G=(M+5)\ell,\qquad
 N_1=\lceil\max(6eG\delta,(h+3)\log_2x)\rceil+1.
\end{gathered}\tag{GII.1}
\]
Thus $G=\texttt{Gp}$ and $N_1=\texttt{N1}$, including their ceilings.
All assertions below use one sufficiently large $x$ threshold after the
fixed constants have been chosen.

\subsection{The finite counting input}

\paragraph{GII1: enough tops inside each good orbit.}
Let $T\subset\mathbb Z/n'\mathbb Z$ satisfy the actual good-orbit
conditions, with $3G\le n'$. A marked complete tile of length $\ell$
has one selected even point in each even-numbered cell of length $u$,
and no selected points in odd-numbered cells. Every $M+1$ successive
complete tiles contain a marked one. List the marked tile starts
periodically on $\mathbb Z$. Their successive gaps are less than
$(M+2)\ell$: the wrap contributes the remainder of $n'$ upon division
by $\ell$, which is less than $\ell$. From any position, the next
marked tile's second cell therefore supplies an even top at positive
distance at most $G$.

The map $(t,d)\mapsto t-d$, for $t\in T$ and $1\le d\le G$,
covers the orbit. Conversely, double-counting all cyclic $G$-windows,
each of which contains at most $N_1$ tops, gives
\[
 n'\le |T|G\le n'N_1.\tag{GII.2}
\]
The upper bound improves the source's $|T|G\le N_1(n'+G)$; we retain
its downstream constants. The least positive even successor has gap
$\le G$ and at most $N_1$ interior tops. It is distinct because $G<n'$.
Hence there are at least two even tops, and their literal cyclic
segments satisfy $\texttt{segG}\le G$ and $\texttt{segS}\le N_1$.
This proves \texttt{OrbitSegs}. When $L_1\le L_2$ and the short-window
cap is $r-1$, fewer than $r$ even tops occur in the preceding $L_1$
positions. These are the remaining positional clauses of
\texttt{AuxBlock8.adm1}.

\paragraph{GII2--GII3: one uniform constant and the actual orbit map.}
For $x\ge3$, put
\[
 C_G=2(6r+1)\bigl((h+3)10^{6r+1}+6\bigr),\qquad C_e=20C_G.
\]
The inequalities $1\le u\delta\le2$, $M< (h+3)\log x/P_0+1$ and
$\log x\ge1$ give $G\delta\le C_G\log x$. Moreover $G\delta\ge5$ and
$(h+3)\log_2x\le2M\le2G\delta$, so the first argument of the maximum
in (GII.1) dominates. Thus $N_1<6eG\delta+2\le16.8G\delta$. Consequently
\[
 G\le\frac{C_e\log x}{\delta},\quad N_1\le C_e\log x,\qquad
 \frac{n'\delta}{C_e\log x}\le N:=|T|\le40n'\delta.
 \tag{GII.3}
\]
For an inert prime, $y\mapsto t_0+yh_p$ is the actual orbit bijection,
and the representative identities identify $N$ with \texttt{Nv}.
The predicate \texttt{IsV} means that the unique index $y$ is even.
Transporting the preceding counts and distances proves
\texttt{adm1} for every $C_{31}\ge C_e$.

In a supplied family, $r,h,C_e$ depend only on $\epsilon$. For
$x\ge\max(3,x_1+1)$ these statements hold for every member. If
$\mathbf C.\mathrm{CN}\ge C_e$, they also prove the exact
\texttt{FamilyWindowLogBound}. Nonnegative substitution gives
$\rho\le\rho_B\le x^{-2}$, using the family's stated bound on $\rho_B$.
No current family index is selected; empty index sets cause no difficulty.

\subsection{Supplying every Type-II inequality}

Write
\[
\begin{gathered}
 a=xe^{x^\epsilon},\quad X'=e^{x^{\epsilon/3}},\quad b'=a/X',
 \quad p_0=k^2x^{26},\quad Q=p^G,\\
 L_1(p)=\lceil\log_p((kx^{26})^4)\rceil,\qquad
 L_2=2\sup_{p\ \mathrm{inert}}L_1(p),\quad c=4\mathrm{Sc},\quad c_0=2b'+2c.
\end{gathered}\tag{GII.4}
\]
The relevant \texttt{SizesAt} clauses give
$x^{9/2}\le k\le x^5\le x^6$, $n\le x^3$ and $\mathrm{Sc}\le b'/x$.
The actual parameter-fill record gives
$\Lambda:=C_{\rm La}N_1^2\log a=\mathrm{LaG}$.
Eventually $c\le b'$, $c_0\le5b'$ and $b'\ge1$.

\paragraph{GII4: levels, prime size and weights.}
Since $\log(kx^{26})\le31\log x$ and $\log p\ge\log2>0.69$,
$L_2\le514\log x<u$ eventually; its definition gives $2L_1\le L_2$.
Also $3G\le x^{\epsilon/4}$ eventually. The host's small primes split
up to $\lfloor x^{\zeta/5}\rfloor$; inert classification excludes them
and $2$. Hence $p>x^{\zeta/5}$, so $p\ge13$ eventually, and
\[
 L_2\le L_b:=2\lceil620/\zeta\rceil.\tag{GII.5}
\]
The latter constant will also control the probability argument.

Let $I_b=2n\delta x^\epsilon\Lambda b'/k$. If $\rho\le1/7$, its
nonnegative summands give $I_b\le b'$. The actual integer weights lie in
$[a,a+a/(kx^{13})]$ and $[b',b'+I_b]$, precisely the global windows.
For $x\ge24$, $c\le b'/6$ and $2b'\le a$. Therefore, with
$c''=\lfloor c\rfloor$,
\[
 (a-c)+c''\le a_t\le2(a-c)-c'',\qquad
 b_t+2c''\le\lfloor c_0\rfloor.\tag{GII.6}
\]
For the upper $a_t$ bound, use $a_t\le3a/2$ and $3c\le a/4$.
Thus all shifted and rounded weight clauses are supplied.

\paragraph{GII5: the lower count supplies the rank margin.}
Assume $p<p_0$, $x^{\epsilon/4}\le n'\le n$ and $N\le n'$.
Put $y=x^{\epsilon/24}$. Choose $x$ so that both $C_e\log x$ and
$C_{31}\log x$ are at most $y$. The lower count gives
$Ny\ge n'\delta\ge y^3$, whence
\[
 N_1^2\le y^2\le N\le rN.\tag{GII.7}
\]
Also $r,N_1\le x$ eventually, so $N_1+3r\le4x\le k$ and
$r\le x^{26}$. Since $\log a\le2x$,
$C_k'N_1^2\log a\le2C_k'x^3\le k$ eventually. Finally the ceiling in
$L_1$ gives $p^{L_1}\ge(kx^{26})^4$. Dividing by $p<p_0$ yields
\[
 256k^2N(x^{26})^2\le k^2x^{78}\le p^{L_1-1},\tag{GII.8}
\]
using $N\le x^3$ and $x\ge500$.

\paragraph{GII6: the full exponential size budget.}
Let $N_+=40n'\delta$ and use the actual denominator
\[
 D=\left\lceil C_DN_+(N_1^3+r^5+k^2)\log_pa
                 +2rn'c+rN_+c_0\right\rceil.\tag{GII.9}
\]
Because $p\le x^{38}$, (GII.3) gives
$\log Q\le38C_G(\log x)^2x^{\epsilon/8}\le x^{7\epsilon/48}$
eventually. Put $E=\exp(5x^{7\epsilon/48})$, so $Q^5\le E$.
For $x\ge\max(500,480C_D,C_s)$, $\log_pa\le4x$ and
\[
\begin{aligned}
 D&\le x^{17}+202x^4b'+1\le x^{18}b',\\
 8(D+1)+rNc_0+2rn'c&\le9x^{18}b',\\
 C_sk^3r^2NQ^5[8(D+1)+rNc_0+2rn'c]x^{26}
   &\le9x^{68}Eb'\le x^{69}Eb'.
\end{aligned}\tag{GII.10}
\]
These estimates include the ceiling's extra unit. Since
$7\epsilon/48<\epsilon/3$, choose $x$ so that
$70\log x+5x^{7\epsilon/48}\le x^{\epsilon/3}$. Then
$x^{69}E+1\le x^{70}E\le X'$, and the last line of (GII.10) is
$\le a-b'\le a-c$. This proves every field of \texttt{NumII}.

\paragraph{GII7--GII8: the literal family application.}
A candidate prime has $n'=\operatorname{ord}(h_p)\ge x^{\epsilon/4}$;
its order divides $n$. For each good representative, (GII.3) supplies
the count used above. We have now proved all of \texttt{HypIIAt},
uniformly in the prime, orbit and family member. In a \texttt{FamP},
\texttt{sizes}, \texttt{fills} and the bound on $\rho_B$ supply exactly
the records used here. Thus
\[
 \exists x_0\ \forall x\ge x_0\ \forall i\in\mathrm{Ix}(x),
 \qquad\texttt{HypIIAt}(D_{x,i},L_{x,i},X_{x,i},\mathbf C).
 \tag{GII.11}
\]
The constant tuple and $\epsilon,J$ precede $x_0$; $i,p$ and the orbit
follow it. If no family with the given $\epsilon$ exists the universal
claim is vacuous, rather than requiring an invalid division by $\epsilon$.

For completeness, use $c\ge0$ and $a-c\ge1$, already proved above.
Take a local constant $C\ge4129$ and require
$C\le C_k,C_t,C_{\rm La},C_D,C_s,C_k'$. Apply EA9 at
\[
 m=k-1,\quad a_0=a-c,\quad Y=x^{26},\quad
 c''=\lfloor c\rfloor,\quad c_{0,\mathrm{int}}=\lfloor c_0\rfloor,
 \quad D_p=D,\quad c_*=8(D+1).
\]
Equations (GII.6--10) verify its hypotheses. In particular
$CNk^2\log_pa_0\le D$. Monotonicity of every nonnegative height term
gives the exact global height
\[
 t_{\rm II}=C_t rN\left(\frac{\Lambda a}{pk}+ap^{-L_1}+c_0+kQ\right)+2rn'c.
\]
Hence the local witness works at $\lceil t_{\rm II}\rceil$ and loss
$\kappa_{\rm II}=\min(1,C_kNr^2kx^{-26})$. The truncation at $1$ uses
the empty witness when its rank requirement is nonpositive. This proves
the actual \texttt{adm6}, at the supplied standard coefficient setting;
the finite extension to an actual place is separate.

\subsection{Probability and a simultaneous choice of tops}

\paragraph{GII9: condition on a selected point correctly.}
Select the ambient points independently with probability $\delta$ and
condition on a fixed point being selected. Its orbit has one forced
index $i_0$; the other indices remain independent. At most $L_2$ short
windows contain $i_0$. Failure of their cap $K_0$ therefore has probability
at most
\[
 L_2\binom{L_2}{K_0}\delta^{K_0}
 +n'\binom{L_2}{K_0+1}\delta^{K_0+1}.
\]
Since $\delta^{K_0}\le x^{-(h+2)}$, $n'\le n\le x^{h+\epsilon/8}$
and $L_2\le L_b$, this is at most $(L_b+1)2^{L_b}x^{-2}$.
The additional $\delta$ in the second term is essential.

For $\delta\le1/10$, $(1-\delta)^u\ge1/4$ and
$\lfloor u/2\rfloor\delta\ge9/20$. Thus an empty cell has probability
$\ge1/4$, and exactly one selected even point has probability $\ge9/80$.
A complete tile is marked with probability at least $P_0$.
Among any $M+1$ distinct consecutive tiles, discard the one containing
$i_0$, if any. The other $M$ marking events are independent, so a union
bound gives tile failure at most
$n'(1-P_0)^M\le n'x^{-(h+3)}\le x^{-2}$.

Put $B=N_1-1$. For a long window, the conditional failure probability
is bounded by the binomial tail
\[
 \binom GB\delta^B\le(eG\delta/B)^B\le2^{-B}\le x^{-(h+3)}.
\]
If $B>G$ the tail is zero. A union bound over $n'$ windows contributes
at most $x^{-2}$. Therefore
\[
 \Pr(\text{bad orbit}\mid\text{point selected})\le C_bx^{-2},
 \qquad C_b=(L_b+1)2^{L_b}+2.\tag{GII.12}
\]
Marginalizing over points outside the orbit multiplies by $1$; removing
the conditioning multiplies by $\delta$. This is the literal
\texttt{BadProb}, uniform over all sufficiently long steps and all points.

\paragraph{GII10: one subset works for every prime in the finite sum.}
Let $\mu=\delta|H|\ge16$. Chebyshev gives size between $\mu/2$ and
$2\mu$ with probability at least $3/4$. For nonnegative weights
$\omega_p$, the weighted bad count $Z$ has expectation at most
$B_0=C_bx^{-2}\mu\sum_p\omega_p$. Markov gives $Z\le4B_0$ with
probability at least $3/4$; if $B_0=0$, then $Z=0$ almost surely.
The events intersect, proving the exact existential \texttt{TopsOK}.

Take $\omega_p=\log p/p$ over the source's primes $p<X'$ of large
orbit order. For $n=\lfloor y\rfloor\ge1$,
\[
 \frac n2\sum_{p\le n}\frac{\log p}{p}
 \le\sum_{p\le n}\lfloor n/p\rfloor\log p
 \le\log(n!)\le n\log n.
\]
Hence $\sum_{p<y}\log p/p\le2\log y$, including the empty small cases.
The selected subset consequently satisfies
\[
 Z\le16C_b|T|x^{-2+\epsilon/3}
   \le |T|x^{-2+\epsilon}\tag{GII.13}
\]
eventually. This is a simultaneous weighted abundance statement.
BO and BS below extend the construction to the host's biased probabilities
and coset shares. HC identifies the selected set with the actual host
subset and proves the exact \texttt{BadOK}/\texttt{ChosenBad} transfer.
\noindent\textit{Primary source route.}\cite{type-two-family}


\section{Supplying the actual Type-I hypotheses along a family}
\label{sec:type-one-family}

\status{GI1--GI8 prove the actual numerical hypothesis and its
uniform family version, then connect TC6 to the source's global
time and exceptional-share mean. GII supplies the good-orbit
count and segment geometry; GJ supplies the logarithmic size
bounds and the canonical exceptional function. The arithmetic
family and its stated size assumptions remain supplied data.}

\paragraph{GI1--GI2: the scales and one uniform threshold.}
Fix $0<\varepsilon<1/52$ and the nonnegative coefficient vector
$K$. Write $k=D.R.k$, $n=|HH(D)|$, and use the source definitions
\[
\begin{gathered}
 a=xe^{x^\varepsilon},\quad
 X'=e^{x^{\varepsilon/3}},\quad
 b'=a/X',\quad \delta_1=x^{-\varepsilon/8},\quad
 Y=x^{13},\quad Y_b=x^\varepsilon/20,\\
 \lambda=K.\mathrm{La}\,N_1^2\log a,\quad
 N_*=2n\delta_1,\quad
 c=4L.\mathrm{Sc}=16k\Delta,\quad
 I=\frac{N_*x^\varepsilon\lambda b'}k,\\
 \rho=\frac{4N_*x^\varepsilon\lambda}k
          +\frac{K.\rho\,N_1\lambda}k.
\end{gathered}                                                    \tag{GI.1}
\]
For a candidate prime $p$ and an actual good representative,
let $n'=\operatorname{ord}(h_p)$, let $T$ be its selected
positions, and put
\[
\begin{gathered}
 N=|T|,\qquad N_+=40n'\delta_1,\qquad Q=p^{G'},\\
 D_f=\left\lceil K.\mathrm{DI}\,N_+Y^2
                    (N_1^5+k^3)\log_p a\right\rceil_+.
\end{gathered}                                                    \tag{GI.2}
\]
The order of an element is at most the group order, and the
orbit parametrization is injective before its period. Thus
$N\le n'\le n$. The literal family size assumptions give
\[
 x^{9/2}\le k\le x^6,\qquad n\le x^3,\qquad
 \log\Delta\le\tfrac12x^{\varepsilon/4},\qquad
 c\le\lambda b'/k.                                      \tag{GI.3}
\]
GJ2, at $\eta=\varepsilon/48$, gives
\[
 x^{47\varepsilon/48}\le\lambda\le x^{49\varepsilon/48},
 \qquad \log Q\le x^{23\varepsilon/48}\quad(p<X').        \tag{GI.4}
\]
GII2 gives a fixed $B=B(\varepsilon)>0$ with
$G'\le B\log x\,x^{\varepsilon/8}$ and $N_1\le B\log x$.
One exact choice is
\[
 B=40\bigl((h+3)10^{6r+1}+6\bigr)(6r+1),\quad
 h=2+8\varepsilon,\quad
 r=\lceil8(h+2)/\varepsilon\rceil_++1.
\]
Enlarge a threshold depending only on $\varepsilon,K$ so that
\[
\begin{gathered}
 x\ge\max\{e,7,320K.\mathrm{DI}+1,K.\mathrm{kI}\},\\
 \max\{3B,514\}\log x\le x^{\varepsilon/8},\qquad
 (2B+3)\log x\le x,\qquad
 128\log x\le x^{47\varepsilon/48},\\
 20K.\mathrm{kI}\le x^\varepsilon,\qquad
 39\log x\le x^{\varepsilon/3},\qquad
 K.\mathrm{kI}\,x^4\le x^{9/2},\\
 340\log x\le x^\varepsilon,\quad
 16x^{23\varepsilon/48}\le x^\varepsilon,\quad
 2x^{\varepsilon/4}\le x^\varepsilon,\quad
 4x^{\varepsilon/3}\le x^\varepsilon.
\end{gathered}                                                    \tag{GI.5}
\]
Each inequality follows from a strict gap between exponents;
for example, putting $t=\log x$ and using
$e^{ut}\ge u^2t^2/2$ proves $C\log x\le x^u$ eventually.
The threshold precedes every prime, orbit and family index.

The candidate condition $n'\ge x^{\varepsilon/4}$ now gives
$3G'\le n'$. Also, for every prime $q$,
\[
 L_1(q)=\lceil\log_q((kx^{26})^4)\rceil_+
       \le\lceil256\log x\rceil_+,
\]
because $\log q\ge1/2$. Taking the bounded natural supremum
in the definition of $L_2$ gives
\[
 L_2<512\log x+2\le514\log x
       \le x^{\varepsilon/8}\le U=\lceil x^{\varepsilon/8}\rceil_+.
                                                               \tag{GI.6}
\]
The empty defining set gives $L_2=0$. These are exactly the two
fields of the source's larger \texttt{NumMean} record needed
here. GII1--GII3 now give the actual segment geometry and
$N\le N_+$.

\paragraph{GI3: verify the small numerical fields.}
Assume $p\ge p_0=k^2x^{26}$, $p<X'$ and $\rho\le x^{-2}$.
The first thirteen fields of \texttt{AuxBlock8.HypI} are
\[
\begin{gathered}
 k^2Y^2\le p,\quad p<a/b',\quad \rho<1/6,
 \quad 2N_1+2\le k-1,\quad 8\log(NY)\le\lambda,\\
 Y_bN\lambda b'/k\le I,\quad 2I\le b',\quad
 K.\mathrm{kI}\le Y_b,\quad
 kY^2Y_bN\lambda\log a\le X',\\
 K.\mathrm{kI}\,N\lambda^2\le k,\quad
 2Y_bN\lambda<k,\quad c\le\lambda b'/k,\quad N\le N_+.
\end{gathered}                                                    \tag{GI.7}
\]
Here is a direct check. Equations (GI.3)--(GI.5) give
$1\le\lambda\le x$, $\lambda^2\le x$,
$x^\varepsilon\lambda\le x$, $N_1\le x$ and
$\log a\le2x$. The rank margin follows from
$2N_1+3\le(2B+3)\log x\le x\le k$; the logarithmic bound
follows from $NY\le x^{16}$. The actual good orbits have $N>0$.
Moreover $N\le40n\delta_1$ gives $Y_bN\le N_*x^\varepsilon$.
The nonnegative summands of $\rho\le1/7$ give
$2I\le b'/14$, which is stronger than required.
The remaining product estimates are
\[
\begin{aligned}
 kY^2Y_bN\lambda\log a&\le2x^{38}\le x^{39}\le X',\\
 K.\mathrm{kI}\,N\lambda^2&\le K.\mathrm{kI}\,x^4\le k,\\
 2Y_bN\lambda&\le x^4/10<k.
\end{aligned}                                                    \tag{GI.8}
\]
The prime, conductor and top-count clauses are precisely their
stated hypotheses. Thus no one of the thirteen fields is an
additional asymptotic assumption.

\paragraph{GI4: the full depth and exponential budget.}
The fourteenth field is
\[
 K.\mathrm{kI}\,k^3NQ^4[8(D_f+1)+(N+n')c]Y\le b'.       \tag{GI.9}
\]
Since $N_+\le40x^3$, $N_1^5+k^3\le2x^{18}$ and
$\log_p a\le4x$, its actual rounded depth satisfies
\[
 D_f\le320K.\mathrm{DI}\,x^{48}+1\le x^{49}.
\]
Set $E_1=e^{x^{\varepsilon/4}/2}$ and
$E_2=e^{4x^{23\varepsilon/48}}$. Then $c\le16x^6E_1$,
$Q^4\le E_2$, and, for $x\ge7$,
\[
 8(D_f+1)+(N+n')c\le48x^{49}E_1\le x^{51}E_1.
\]
The left side of (GI.9) is at most $x^{86}E_1E_2$.
The final four inequalities of (GI.5) give
\[
 85\log x+\tfrac12x^{\varepsilon/4}
       +4x^{23\varepsilon/48}+x^{\varepsilon/3}
       \le x^\varepsilon.
\]
Exponentiating proves $x^{86}E_1E_2\le
xe^{x^\varepsilon-x^{\varepsilon/3}}=b'$ and completes
the literal fourteen-part hypothesis.

\paragraph{GI5: the actual family statement.}
For a supplied \texttt{FamP}, the window hypothesis is the
eventual inequality $N_1\le K.\mathrm{CN}\log x$.
Substitution into the nonnegative second summand of $\rho$ gives
$\rho\le\rho_B\le x^{-2}$, where the last bound is exactly
\texttt{FamP.rho\_le}. This uses no prime sum over the excluded
band. The family supplies (GI.3), and its filling identity
$\lambda=L.\mathrm{LaG}$ makes the conductor inequality and
the chosen lower-degree window exact.

Indeed, at every selected position the actual integer degrees
lie in
\[
 W_a=[\lceil a\rceil,\lfloor a+a/(kY)\rfloor]\cap\mathbb Z,
 \qquad
 W_b=[\lceil b'\rceil,\lfloor b'+I\rfloor]\cap\mathbb Z.  \tag{GI.10}
\]
This follows from \texttt{aT\_win}, \texttt{bT\_ge} and
\texttt{SizeBot}, followed by the definition of the orbit
restriction. Combining these memberships with (GI.6)--(GI.9)
proves the exact \texttt{Joints.HypIAt}: eventually, for
\emph{every} family index, candidate Type-I prime and good
representative, the two spacing bounds, numerical hypothesis
and actual weight windows hold. A single maximum of the fixed
thresholds suffices; no choice depending on the prime or index
has been made. This does not assert existence of a family.

\paragraph{GI6: convert to the completed local theorem.}
Take
\[
\begin{gathered}
 C_0=K.\mathrm{La},\quad
 0\le C_H\le\min\{K.\mathrm{kI},K.\mathrm{DI}\},\quad
 2C_H\le K.\rho,\\
 m=k-1,\quad c_{\rm cl}=\lfloor c\rfloor_+,\quad
 D_p=D_f,\quad c_*=8(D_f+1),\quad
 b_l=b',\quad b_u=b'+I.
\end{gathered}                                                    \tag{GI.11}
\]
GII1 gives an even selected point, at least two segments, and
their bounds $\mathrm{segG}\le G'$, $\mathrm{segS}\le N_1$.
The only delicate normalization is the local band radius:
\[
 \delta_I=\frac{C_HN_1\lambda}k+\frac{2I}{b'},\qquad
 2\delta_I\le\rho<1/6.                                  \tag{GI.12}
\]
Hence the local forbidden interval
$[X'(1-2\delta_I),X'(1+2\delta_I)]$ is contained in the
global interval with radius $\rho$. Exclusion of
\texttt{InBand0} therefore excludes every odd prime power
required locally. An actual candidate excludes the larger
$\rho_B$ interval, so supplies this premise.

All remaining local fields follow from (GI.7)--(GI.9):
$m+1=k$, $N_1\ge1$, $I\le b'/2$, $c_{\rm cl}\le c$,
and
\[
 C_HNY^2(N_1^5+k^3)\log_p a\le D_f
\]
by $N\le N_+$ and $C_H\le K.\mathrm{DI}$.
The last size product only decreases on replacing
$K.\mathrm{kI},c$ by $C_H,c_{\rm cl}$. This proves the
literal local \texttt{HypI}, with its original depth and
uniformizer normalization.

\paragraph{GI7: compare the exact time bounds.}
For each odd $y\in T$, its nearest preceding even selected
point has a unique least backward distance $1\le d<n'$.
The forward and backward descriptions are equivalent modulo
$n'$. Consequently the local \texttt{dPrev,yPrev} and global
\texttt{dMinus} define the same leading sum
\[
 S=\sum_{y\in T,\ y\text{ odd}}
        \min\{b_y,a_{y-d}/p^d\}.
\]
This identity holds for arbitrary signed integer degrees.
Put $E=a/p^2+\lambda a/(pk)$ and
\[
\begin{aligned}
 c_{\rm loc}&=8(D_f+1)+CNk\log_p a+4Nc_{\rm cl}+2n'c_{\rm cl},\\
 c_6&=8(D_f+1)+K.\kappa_2N_+k\log_p a+4N_+c+2n'c,
\end{aligned}
\]
where $K.\kappa_1,K.\kappa_2$ abbreviate the two global
\texttt{kappaThreshold} coefficients. The actual definitions are
\[
\begin{aligned}
 t_I={}&S+CNN_1(E+b'/p+\lambda b'/k)\\
       &+CN(Qc_{\rm loc}+kQ+k\log_p a+c_{\rm cl})+2n'c_{\rm cl},\\
 t_{If}={}&S+K.\kappa_1NN_1E\\
       &+K.\kappa_2N(Qc_6+kQ+k\log_p a+c)+2n'c.
\end{aligned}                                                    \tag{GI.13}
\]
If $2C\le K.\kappa_1$ and $C\le K.\kappa_2$, then
$b'\le a/p$, $c_{\rm loc}\le c_6$ and $c_{\rm cl}\le c$
prove term by term
$t_I\le t_{If}\le\lceil t_{If}\rceil_+$.
The common signed sum $S$ is unchanged.

\paragraph{GI8: the actual exceptional mean.}
TC6 supplies one function $\kappa$ on $W_a^T\times W_b^T$,
chosen before the coefficient rings, with mean at most
$CNN_1/Y$. GI7 makes all its certificates valid at the exact
global time $\lceil t_{If}\rceil_+$. Leastness of the source's
canonical exceptional function therefore gives
\[
 0\le\kappa_I(\mathbf a,\mathbf b)
       \le\kappa(\mathbf a|_T,\mathbf b|_T).
\]
Restriction to $T$ is a bijection from the source's orbit box
to $W^T$; its inverse extends by zero. Thus both double sums
and their denominators agree, proving the literal
\texttt{adm8} bound
\[
 \operatorname{meanLoc}(\kappa_I)\le C\,NN_1/Y.           \tag{GI.14}
\]
GJ1 supplies attainment and unpacks the actual family field
\texttt{ex.kapI\_def}; an arbitrary instance must still satisfy
\texttt{KapIDef}. Empty boxes give zero under the source convention.

The following fixed vector is sufficient for TC6 and every
comparison above:
\[
\begin{gathered}
 K.\mathrm{La}\ge880,\qquad
 K.\mathrm{kI},K.\mathrm{DI}\ge2800001,\qquad
 K.\rho\ge5600002,\\
 K.\kappa_1\ge604228,\qquad K.\kappa_2\ge302114,
 \qquad C=302114.
\end{gathered}                                                    \tag{GI.15}
\]
The mean coefficient may be any $C_m\ge302114$.
These choices precede the scale, arithmetic object, prime,
weights and coefficient rings. Global prime averaging and the
existence of the simultaneous arithmetic family remain separate.
\noindent\textit{Primary source route.}\cite{type-one-family}


\section{Finite coefficient extension and the actual standard subring}
\label{sec:finite-base-change}

\status{BC1--BC8 construct the source's standard coefficient ring at a
specified finite place and prove the actual \texttt{LPBaseChange} theorem.
The parameters $t,D,\kappa,c_*$ are unchanged. This closes the finite
extension step that HM12's ring-isomorphism argument did not address.
For a literal \texttt{recC} datum the standard subring is automatic;
existence of the full globally compatible family remains separate.}

\paragraph{BC1: what the standard ring actually is.}
Let $p\ge5$ be prime, $f\ge2$ even and $q=p^f$. With $W(k)$ denoting
Witt vectors, set
\[
 W=W(\mathbb F_q),\quad W_2=W(\mathbb F_{q^2}),\quad
 O_s=W_2[\omega]/(\omega^2-p),\quad K_s=\operatorname{Frac}(O_s).
                                                               \tag{BC.1}
\]
The field $Q=\operatorname{Frac}(W)$ embeds in $K_s$. Its Frobenius
$\sigma$ reduces to $x\mapsto x^p$ and satisfies $\sigma^f=1$.
The extension $W_2/W$ is unramified of degree two, whereas adjoining
$\omega$ is ramified of degree two. Thus $O_s$ is free of rank four
over $W$, has uniformizer $\omega$ and residue field $\mathbb F_{q^2}$.
These are the source's \texttt{WQ, OO, KO, jO, iota, wU}.

To check the assertions directly, lift a residue-field basis of
$\mathbb F_{q^2}/\mathbb F_q$ and successively subtract its residue
expansion, dividing the remainder by $p$. Completeness gives a spanning
basis of $W_2$ over $W$; reduction of a relation after removing its
smallest $p$-power proves independence. The polynomial $X^2-p$ is
irreducible since a square in $\operatorname{Frac}(W_2)$ has even
$p$-valuation. For $a+b\omega\ne0$ the candidate orders
$2v_p(a)$ and $2v_p(b)+1$ differ. Factoring the smaller power of
$\omega$ leaves a unit. This proves the DVR and residue assertions;
$\omega^2=p$ also proves completeness.

Every integer is a square in $O_s$. Write $z=p^m u$ with $p\nmid u$.
Every nonzero element of $\mathbb F_q$ is a square in
$\mathbb F_{q^2}$, since $q-1$ divides $(q^2-1)/2$.
Hensel iteration lifts a square root of $\bar u$ to $y\in W_2$
(the derivative $2y$ is a unit), and $(\omega^m y)^2=z$.
The case $z=0$ is immediate. This checks the full standard quadratic
setting, including its finite structural map from $W$.

\paragraph{BC2: the finite extension at the chosen place exists.}
We need two elementary facts, with proofs recorded here because they
supply the actual coefficient ring.

First, if $R$ is a complete DVR with characteristic-zero fraction field
$K$ and $E/K$ is finite, its integral closure $S$ in $E$ is a finite
free complete DVR. Choose an integral primitive element and the full
order $A=R[\theta]$. Separability and integrality of traces give
\[
 A\subset S\subset A^\#,
 \qquad A^\#=\{x\in E:\operatorname{Tr}_{E/K}(xA)\subset R\}.
                                                               \tag{BC.2}
\]
The trace form is nondegenerate, so $A^\#$ is a finite fractional
lattice. Hence $S$ is finite and torsion free, therefore free over $R$.
Clearing polynomial coefficients shows $\operatorname{Frac}(S)=E$.
The ring $S$ is complete in the uniformizer-adic topology. If it had
several maximal ideals, the finite Artinian ring $S/\pi S$ would
have a nontrivial idempotent. Elements of $1+\pi S$ are units by the
convergent geometric series; Newton iteration for $X^2-X$ therefore lifts it
to $S$, contradicting that $S$ is a domain. Thus $S$ is local.
It is normal and one-dimensional, hence a DVR; explicitly, for
$0\ne a\in\mathfrak m$, the equality $\sqrt{(a)}=\mathfrak m$
and finite generation give a power $\mathfrak m^n\subset(a)$. Choose the least $n$ with
$\mathfrak m^n\subset(a)$ and $b\in\mathfrak m^{n-1}\setminus(a)$.
Then $(b/a)\mathfrak m\subset S$. If this lay in $\mathfrak m$,
the determinant trick would make $b/a$ integral. Otherwise
$(b/a)x$ is a unit for some $x\in\mathfrak m$, forcing
$\mathfrak m=(x)$. The two adic topologies agree, proving completeness.

Second, let $L$ be the source number field and $v$ a specified prime
above $p$. Inside $\Omega=\overline{K_s}$, each embedding $\tau:L\to\Omega$
defines
\[
 E_\tau=K_s\tau(L),\qquad
 S_\tau=\text{integral closure of }O_s\text{ in }E_\tau.  \tag{BC.3}
\]
This extension is finite, and its maximal ideal induces a prime
$P_\tau$ of $L$ above $p$. Suppose none is $v$. The Chinese remainder
theorem gives an integer $a\in L$ that is $0$ modulo $v$ and $1$
modulo every distinct $P_\tau$. All $\tau(a)^{-1}$ are integral over
$O_s$, so their product $N_{L/\mathbb Q}(a)^{-1}$ is integral.
But $p\mid N_{L/\mathbb Q}(a)\ne0$, since the norm is an integer
multiple of $a$ in the ring of integers. This makes $p^{-1}$ integral
over the integrally closed ring $O_s$, a contradiction.
Thus some $\tau$ induces $v$, and
\[
 \tau(x)\in S_\tau\quad\Longleftrightarrow\quad v(x)\ge0.
                                                               \tag{BC.4}
\]
For the reverse implication use a local fraction with denominator
outside $v$; for the converse use $x^{-1}$ in the local maximal ideal.
These are precisely the objects \texttt{KfRec, ORec, embRec, extRec}.

\paragraph{BC3: exact lattice identities and ramification.}
Let $\phi:R\hookrightarrow S$ be a finite inclusion of DVRs, with
fraction fields $K\hookrightarrow K'$. The finite torsion-free
$R$-module $S$ is free. Moreover $1$ is primitive: if $\pi$ were
invertible in $S$, then $\pi^{-1}\in K$ would be integral over $R$.
Choose an $R$-basis $b_1=1,b_2,\ldots,b_d$ of $S$. For a
$K$-space $V$ and any $R$-submodule $M\subset V$,
\[
 SM=\bigoplus_{a=1}^d b_aM\subset K'\otimes_KV.
 \quad\text{Consequently}\quad
 S\Bigl(\bigcap_iM_i\Bigr)=\bigcap_iSM_i,
 \qquad SM\cap V=M.                                     \tag{BC.5}
\]
The intersection may be infinite: finite freeness, not flatness alone,
justifies the coordinate argument. Kernels, invariant submodules defined
by universal coefficient equations, saturation and ranks also commute
with this extension. These statements give faithful descent of lattice
inclusions and equalities.

The dual map is the semilinear map
$\delta(G)(a\otimes x)=a\phi_K(G(x))$. A lattice basis and its
algebraic dual basis prove
\[
 S\delta(M^\vee)=(SM)^\vee.                              \tag{BC.6}
\]
The pairing is $K'$-valued evaluation. No trace-dual or different is
involved. If $\phi(\pi_R)=u\pi_S^e$, then
\[
 \operatorname{ord}_S\phi_K(x)=e\operatorname{ord}_R(x),
 \qquad
 \frac{\operatorname{ord}_S\phi_K(x)}{\operatorname{ord}_S(p)}
 =\frac{\operatorname{ord}_R(x)}{\operatorname{ord}_R(p)}.
                                                               \tag{BC.7}
\]
Lengths multiply by $e$, as follows by applying extension to each
Smith factor $R/(\pi_R^a)$. Normalized lengths and determinant
valuations are unchanged. In particular the image of $\omega$ has
integer valuation $e$ and $p$-normalized valuation $1/2$; it need
not be a uniformizer.

\paragraph{BC4--BC5: identify the actual polynomial lattices.}
CB2 supplies the same monic integral tableau basis for the literal
dual-Weyl module $\nabla_\mu$ over every domain. Hence
$S\otimes_R\nabla_\mu(R)=\nabla_\mu(S)$, including the generic
unipotent-invariance identities. For decreasing integral $\lambda$,
put $\mu_i=\lambda_i-\lambda_k$ and $T(g)f(X)=f(gX)$. The actual
lattice at $\Lambda=gR^k$ is
\[
 H_\lambda(\Lambda)
 =\det(g)^{-\lambda_k}T(g^{-1})\nabla_\mu(R).             \tag{BC.8}
\]
Replacing $g$ by $gh$, $h\in\mathrm{GL}_k(R)$, leaves it unchanged:
$T((gh)^{-1})=T(g^{-1})T(h^{-1})$, the inner factor preserves
$\nabla_\mu(R)$, and $\det(h)^{-\lambda_k}$ is a unit.
Thus the source's intersection over all lattice bases equals
(BC.8). Determinants, inverses and substitution commute with coefficient
maps, so (BC.8) commutes with extension. Products of tableau bases in
disjoint variable blocks give the same assertion for the actual tensor.

Here is the full dictionary. For $s=0,\ldots,f-1$, use
$\lambda_s=\lambda_s^0$ when $s$ is even and
$\lambda_{s,i}=-\lambda^0_{s,k+1-i}$ when $s$ is odd. Set
$A_r=\sum_s\sum_{i\le r}\lambda_{s,i}$ and
\[
\begin{aligned}
 \Lambda_s(L)&=O\,\iota(\sigma^sL),&
 \Lambda_s(e)&=\Lambda_s(e.\mathrm{src})
                   +\omega^{-1}\Lambda_s(p\,e.\mathrm{tgt}),\\
 H(L)&=\bigotimes_s H_{\lambda_s}(\Lambda_s(L)),&
 H(e)&=\bigotimes_s H_{\lambda_s}(\Lambda_s(e)),\\
 R_e&=\omega^{-A_r}H(e),&R'_e&=\omega^{-A_r}H(e)^\vee,
 \qquad b(F,G)=p^{A_r}G(F).
\end{aligned}                                             \tag{BC.9}
\]
An adapted $W$-basis of the edge has target
$g\operatorname{diag}(p^{-1}I_r,I)W^k$ and source $gW^k$;
therefore its midpoint lattice is
$\iota(\sigma^sg)\operatorname{diag}(\omega^{-1}I_r,I)O^k$.
These are genuine full lattices, with no implicit choice of representation.

Let $\nu$ map polynomial coefficients and $\delta$ map duals. Their
common tableau basis identifies the two representation spaces by scalar
extension, with equal dimension. The exact clean lattices are
\[
\begin{aligned}
 N_{\rm src}(L)&=\omega^{-t}H(L)\cap
                           \bigcap_{e.\mathrm{src}=L}R_e,\\
 N_{\rm tgt}(L)&=\omega^{-t}H(L)^\vee\cap
                           \bigcap_{e.\mathrm{tgt}=L}R'_e.
\end{aligned}                                             \tag{BC.10}
\]
Equations (BC.5)--(BC.9) identify their scalar extensions with the
actual target clean lattices. They also preserve every extra
$\omega^{c_*}R_{e'}$ condition. Crucially,
$b'(\nu F,\delta G)=\phi_K(b(F,G))$, even if $A_r<0$.

\paragraph{BC6--BC7: transport the rank certificate without loss.}
The source's \texttt{GoodRankK} certificate is an integral map
$u:O^n\to N_{\rm tgt}(L)$ and an integral matrix $M$ with
\[
 \sum_i M_{ji}b(F_i,u(x))=p^Dx_j\qquad(x\in O^n).
\]
For $y\in S^n$ put $u'(y)=\sum_\ell y_\ell\delta(u(\epsilon_\ell))$
and $M'=\phi(M)$. The previous lattice identities give the correct
codomain, and
\[
 \sum_iM'_{ji}b'(\nu F_i,u'(y))
 =\sum_\ell y_\ell\phi_K\!\left(
         \sum_iM_{ji}b(F_i,u(\epsilon_\ell))\right)=p^Dy_j.
                                                               \tag{BC.11}
\]
Keep the same $W$-edges and send $F_i$ to $\nu F_i$. All clean
and alternative-edge divisibilities hold, and the rank inequality
$n\ge(1-\kappa)\dim V$ is unchanged. So is $4(D+1)\le c_*$.
This proves
\[
 \texttt{LPprime}_{O}(t,D,\kappa,c_*)
 \ \Longrightarrow\
 \texttt{LPprime}_{S}(t,D,\kappa,c_*).                    \tag{BC.12}
\]
If a source weight is not dominant, the defined representation has a
zero tensor factor; the empty list and $n=0$ supply the literal
exceptional branch. The larger coefficient residue field never replaces
$q=\#(W/pW)$ in the edge data. The expanded report also proves descent
for a fixed family of tests by fractional Smith form; no arbitrary newly
chosen family of vectors is claimed to descend.

\paragraph{BC8: match the actual global-family interface.}
The source \texttt{StdSubAt} asks for the standard ring and a finite map
satisfying $\phi_K\iota_s=\iota$, $\phi j_s=j$ and
$\phi(\omega_s)=\omega$. In \texttt{locPackC/recC}, the rings are
literally $WQ,QQ,ORec,KfRec$, and all three maps are the composites
constructed in BC1--BC2. Thus \texttt{ringTied\_recC} uses identity
isomorphisms and the interface follows. For a compatibly isomorphic
presentation, compose the field/ring isomorphisms and their inverses;
finiteness and the three displayed equations persist. The actual
Frobenius on $Q$ stays fixed, using its supplied period and reduction
properties.

This proves \texttt{StdSubAt} pointwise for every well-formed
\texttt{recC}, and eventually for a supplied \texttt{IsRecCEv} family.
The latter is one sufficient identification, not a necessary route:
an actual \texttt{GroupCompatibilityFamily} also supplies the required
ring identifications directly, and BC applies to its ring field.
It does not assume that an arbitrary Hida coefficient DVR already
contains the doubled residue field: that field was deliberately built
into \texttt{recC}. Supplying the remaining global construction and its
simultaneous compatibility is the precise boundary retained here.
\noindent\textit{Primary source route.}\cite{finite-base-change}


\section{Applying the local theorems to the arithmetic family}
\label{sec:family-join}

\noindent\textbf{Status.}
The following arguments identify the actual canonical deficit, cyclic
coordinates and finite-place application. The eventual numerical suppliers
are proved in GI/GII, and the coefficient extension in BC. Existence of
the full global arithmetic construction retains the boundaries stated
below. The source's global chain length $m_{\rm ch}=D.R.m$ is distinct
from the local parameter $m_{\rm loc}=k-1$.
Here $\lceil z\rceil_+=\max(0,\lceil z\rceil)$ denotes the source's
natural-valued ceiling.

\paragraph{GJ1: the universal infimum is attained.}
Fix the orbit data and degree families. Let $S$ be the nonnegative
deficits for which the required certificate holds for every permitted
coefficient setting and dominant weight, at the source height
$\lceil t_{If}\rceil_+$, depth $D_{If}$ and conductor $8(D_{If}+1)$.
The zero certificate gives $1\in S$, so $q_0=\inf S\in[0,1]$.
Fix any one of those settings, weights and target lattices, and put
$T=(1-q_0)d$, where $d$ is its representation dimension.
For $T>0$ choose $q\in S$ sufficiently close to $q_0$ that
$(1-q)d>\lceil T\rceil-1$. Its integer certificate rank is at least
$\lceil T\rceil\ge T$. For $T\le0$ use rank zero. All other conditions
are unchanged. The fixed objects were arbitrary, so $q_0\in S$
with the original universal quantifiers. Thus the actual
$X.\texttt{kapI}=q_0$ supplies \texttt{adm7}, and is bounded above by
each admissible deficit, as required by GI's mean comparison.
This is the actual family choice: for $x$ beyond the family start,
the field \texttt{(Fm.ex x hx i).kapI\_def} says precisely that
\texttt{Fm.Gx.kapI} is this infimum. No replacement of the family
or its degree assignments is required.

\paragraph{GJ2: uniform scale estimates.}
Write $\varepsilon$ for the fixed family exponent and recall
\[
 a=x e^{x^\varepsilon},\quad X'=e^{x^{\varepsilon/3}},\quad
 b'=a/X',\quad \delta_1=x^{-\varepsilon/8},\quad
 p_0=k^2x^{26},\quad Y=x^{13}.
\]
GII2 gives $G'\le C_e\log x/\delta_1$ and
$1\le N_1\le C_e\log x$, uniformly in the family index.
Let $c_\Lambda=K.\texttt{La}>0$ and
$\Lambda=c_\Lambda N_1^2\log a$. For any $\eta>0$ choose $x\ge3$ with
\[
 \left(2c_\Lambda C_e^2+38C_e+c_\Lambda^{-1}\right)(\log x)^3
 \le x^\eta.
\]
This is an explicit uniform choice: $\log x\ge
\max(1,24M/\eta^4)$ suffices for the coefficient $M$ on the left,
by the fourth term of the exponential series. Since $k\le x^6$,
$\log a=\log x+x^\varepsilon$ and $p_0\le x^{38}$, we obtain
\begin{align*}
 x^{\varepsilon-\eta}&\le\Lambda\le x^{\varepsilon+\eta},\\
 p<p_0&\ \Longrightarrow\ \log(p^{G'})\le x^{\varepsilon/8+\eta},\\
 p<X'&\ \Longrightarrow\ \log(p^{G'})\le x^{11\varepsilon/24+\eta}.
\end{align*}
For the last two bounds use $\log p\le38\log x$ or
$\log p<x^{\varepsilon/3}$ respectively. These prove the actual
\texttt{adm5} fields; later we use $\eta=\varepsilon/48$.

\paragraph{GJ3: exact orbit identification.}
For a class-zero prime let $h=h_P$, $n'=\operatorname{ord}(h)$,
and let $t_0$ be its chosen orbit representative. The selector identities
$b_h(t+h)=b_h(t)$ and $t\in b_h(t)+\mathbb Nh$ imply
$b_h(b_h(t))=b_h(t)$. Hence
\[
 \phi(i)=t_0+ih\quad(0\le i<n')
\]
is a bijection from the cyclic positions onto this orbit. It identifies
$T=\{i:\phi(i)\in TT\}$ with $\{t\in TT:b_h(t)=t_0\}$, and therefore
\[
 N_v=|T|,\qquad a_{\rm loc}=a_L\circ\phi,\quad
 b_{\rm loc}=b_L\circ\phi,\qquad
 \sum_{t_0\in\mathrm{goodReps}}N_v\le|TT|. \tag{GJ.1}
\]
The last inequality is disjointness of the representative fibres.
The actual \texttt{GatTies}/\texttt{LatTied} equations identify these
positions and extremes with the coefficient representation at the place.
Its integer bound on the remaining coordinates is at most
$\lfloor c_G\rfloor$. Thus the local base-weight and dominance conditions
hold for precisely that representation.

\paragraph{GJ4: the two prime ranges give the exact orbit certificate.}
Assume \texttt{HypIAt}, \texttt{HypIIAt}, the local interfaces
\texttt{adm6}/\texttt{adm7}, $\rho\le\rho_B$, and the compatible
standard coefficient subring supplied by BC. Fix a served orbit at an
actual finite place. It is a good orbit of a candidate prime, so (GJ.1)
applies. Work first over the standard subring.

If $p<p_0$, apply the Type-II interface using \texttt{HypIIAt}; the
actual rank parameter is $r_G$. If $p\ge p_0$, the rank parameter is
one; \texttt{HypIAt} and GJ1 give the Type-I certificate. The candidate
condition excludes the band of width $\rho_B$, hence also the contained
band of width $\rho$. In both cases the definitions give exactly
\[
 t=\lceil t_v\rceil_+,\qquad D=D_p,\qquad
 \kappa=\kappa_v,\qquad c_*=8(D_p+1).                 \tag{GJ.2}
\]
BC transports the certificate to the actual coefficient DVR with
these same quantities. The image of the chosen square root of $p$
may have larger valuation than a uniformizer; the $p$-normalized
depth and the powers defining height and vanishing remain unchanged.
This proves the literal \texttt{OrbitLatticeBounds}. Empty served sets
require no certificate and supply no assertion of abundance.

\paragraph{GJ5: the uniform family consequence.}
The GI/GII bounds are eventual uniformly in the family index. Taking
the maximum of their finitely many thresholds and the family's start
therefore applies GJ4 to all members at once, whenever the actual
standard-subring and lattice identifications hold. One sufficient
Type-I coefficient vector is
\[
\begin{gathered}
 K.\texttt{La}\ge880,\quad K.\texttt{kI},K.\texttt{DI}\ge2800001,
 \quad K.\texttt{rho}\ge5600002,\\
 K.\texttt{kappaThresholdOne}\ge604228,\qquad
 K.\texttt{kappaThresholdTwo}\ge302114,\qquad C_m=302114.
\end{gathered}                                                   \tag{GJ.3}
\]
These constants precede all family members. The Type-II coefficient
conditions are recorded in GII. The source construction still has to
realize the selected subset, degree pair, local comparison data and the
remaining global objects together. GII proves uniform Bernoulli
selection; BO/BS now prove the biased selection and exact
\texttt{Frame.Shares}, and HC installs them in the host.
GW proves the degree choice, and HW follows its actual preselection
and family transport. LI constructs local-place, completion and
integral-character identifications for the supplied global data.
BC's ring identifications
come directly from supplied \texttt{GroupCompatibilityFamily} data;
\texttt{IsRecCEv} is an alternative sufficient interface. None of these
implications supplies the entire global construction by itself.
\noindent\textit{Primary source route.}\cite{family-orbit-join}


\section{One degree assignment controlling all candidate primes}
\label{sec:global-weight-choice}

\noindent\textbf{Status.}
GW1--GW6 prove the local mean, global finite averaging and exact
weighted-loss choice, including their transport through the family.
The additional \texttt{CoverAll} hypothesis
realizes the selected degrees in a family member. Existence of a
family with that coverage is separate. HW follows the source's active
construction, which selects one pair before forming the final family.

Fix $0<\varepsilon<1/52$ and use the scales in GJ2. All thresholds
below are uniform in the family index. The chain length is $m_{\rm ch}$,
whereas the coefficient representation has rank parameter $k$.

\paragraph{GW1: the actual degree boxes and small mean coefficients.}
The size bounds give
\[
 k\le x^6,\qquad m_{\rm ch}\le x^{1/2+\varepsilon/100},\qquad
 |TT|\le2x^{2+8\varepsilon+\varepsilon/100},\qquad
 N_1\le C_e\log x.
\]
Since $r_G$, $K.\texttt{kII}$ and $C_m=302114$ are fixed, eventually
\[
 K.\texttt{kII}\,r_G^2 kx^{-26}\le x^{-12},\qquad
 C_mN_1/x^{13}\le x^{-12}.                         \tag{GW.1}
\]
For example, $K.\texttt{kII}\,r_G^2\le x^8$ and
$C_m C_e\log x\le x$ suffice. The two integer windows are
\begin{align*}
 I_a&=[\lceil a\rceil,\lfloor a+a/(kx^{13})\rfloor]\cap\mathbb Z,\\
 I_b&=[\lceil b'\rceil,
       \lfloor b'+N_*x^\varepsilon\Lambda b'/k\rfloor]\cap\mathbb Z,
 \qquad N_*=2n\delta_1.
\end{align*}
They are nonempty eventually. Indeed $a\ge x^{19}$,
$b'\ge x^6$, $N_*\ge1$ and $\Lambda\ge1$ eventually, so both real
window widths are at least one. The elementary implication
$v-u\ge1\Rightarrow\lceil u\rceil\le\lfloor v\rfloor$ handles rounding.

\paragraph{GW2: the mean loss for each good orbit.}
Let $p$ be a candidate prime and $v$ a good orbit with $N_v$ selected
positions. For $p<p_0$ its deficit is independent of the degrees:
\[
 \kappa_v=\min(1,K.\texttt{kII}\,N_vr_G^2kx^{-26})
 \le N_vx^{-12}.
\]
For $p\ge p_0$, GJ1 and GI8 give nonnegativity and
\[
 \mathbb E_{\rm local}\kappa_v
 \le C_mN_vN_1/x^{13}\le N_vx^{-12}.                 \tag{GW.2}
\]
Here GI5 supplies the exact numerical and band hypotheses; the
average is over the literal two degree boxes on that orbit. Thus
both assertions in the source's \texttt{MeanKappaBound} hold, including
the pointwise small-prime bound.

\paragraph{GW3: pass from local averages to one common probability space.}
Choose degrees independently and uniformly from $I_a$ and $I_b$ at
each point of $TT$, and assign zero outside $TT$. Write $\mathcal B$
for this finite nonempty product. For a fixed orbit the injection in
GJ3 identifies its $N_v$ positions with a subset of $TT$. Every local
degree pair has exactly
\[
 |I_a|^{|TT|-N_v}|I_b|^{|TT|-N_v}
\]
global extensions. Consequently its marginal average is precisely
the average in (GW.2). Since the good-orbit fibres are disjoint,
\[
 \Sigma_p=m_{\rm ch}\sum_{v\ \mathrm{good}}\kappa_v\ge0,
 \qquad
 \mathbb E_{\mathcal B}\Sigma_p
 \le m_{\rm ch}|TT|x^{-12}.                         \tag{GW.3}
\]
No independence between different primes is needed.

\paragraph{GW4: the total weighted exceptional loss is small.}
Put $t=x^{\varepsilon/3}$ and $X'=e^t$. A harmonic-sum estimate gives
\[
 \sum_{p\ \mathrm{candidate}}\frac{\log p}{p}
 \le t\sum_{j=1}^{\lceil e^t\rceil-1}\frac1j
 \le t(1+t)\le2t^2.                                \tag{GW.4}
\]
This uses only the inclusion of candidate primes in the integers
below $X'$. For $s\ge0$,
$\mathbf1_{\{s>x^{-\varepsilon}\}}\le x^\varepsilon s$.
Define
\[
 Z=\sum_{p\ \mathrm{candidate}}\frac{\log p}{p}
   \bigl(\mathbf1_{\{\Sigma_p>x^{-\varepsilon}\}}+\Sigma_p\bigr).
\]
Equations (GW.3)--(GW.4) and the size bounds imply
\[
 \mathbb E Z
 \le8x^{-19/2+(1453/150)\varepsilon}
 =8x^{-9-\gamma},\qquad
 \gamma=\tfrac12-\tfrac{1453}{150}\varepsilon
 >\tfrac{2447}{7800}>0.                            \tag{GW.5}
\]
This retains the source's stronger intermediate harmonic bound and
uses a slightly simpler, weaker set of size estimates; a prime number
theorem or independence across primes is unnecessary.

\paragraph{GW5: select the single degree pair.}
For $32x^{-\gamma}\le1$, some pair in $\mathcal B$ satisfies
$Z\le x^{-9}/4$, because otherwise its finite average would be larger.
For that same pair, simultaneously over all candidate primes,
\[
 \boxed{\quad
 \sum_{p\ \mathrm{candidate}}\frac{\log p}{p}
  \bigl(2\mathbf1_{\{\Sigma_p>x^{-\varepsilon}\}}+4\Sigma_p\bigr)
 \le4Z\le x^{-9}.
 \quad}                                                   \tag{GW.6}
\]
This is the literal \texttt{WReq}, hence \texttt{WeightVectorExists};
nonnegativity and the bound on $Z$ also give \texttt{MeanEx}.
It controls the weighted exceptional set, not every prime separately.

\paragraph{GW6: transfer the choice through the actual family.}
At a fixed $x\in\texttt{Fm.Win}$ choose a member $i_0$ and apply GW5.
Transfer its selected degrees to the common integer labels
$\texttt{Fm.TTn}(x)$. The field \texttt{bTop\_eq}, together with
$\texttt{LaG}=\Lambda$ from \texttt{fills}, identifies its degree
windows exactly with the family's windows.

Every member at this $x$ has the same frame, selected integer labels,
prime classes, Frobenius steps, scale, orbit selector and canonical
deficit function. These are the explicit \texttt{FamCore}/\texttt{FamP}
equalities. Relabel by the frame identification. The definitions of
candidates, good representatives, pulled-back degrees and $\Sigma_p$
then agree term by term. Thus every member carrying the chosen degree
pair satisfies (GW.6). This proves \texttt{ChosenW} directly from
GW1--GW5, without importing the larger \texttt{LinkG} package.

If \texttt{Fm.CoverAll} also holds, it supplies a member carrying
any admissible pair, in particular this one. This proves
\texttt{ChosenWS}; \texttt{coverS} can then move its distinguished
position while preserving the degrees. \texttt{CoverAll} is a
separate property, not a field of \texttt{FamP}. It is an optional
route: HW below reconstructs the actual preselected-pair transport
without requiring it. Neither route infers full arithmetic-family
existence from finite averaging alone.
\noindent\textit{Primary source route.}\cite{global-degree-choice}


\section{Realizing the chosen degrees in the actual family}
\label{sec:host-weight-join}

\status{We prove the degree-realization part of the actual
\texttt{HostWW--KitS--ChosenWall} route. Arithmetic hosts, the point
constructor for every admissible degree pair, and its numerical records
remain supplied inputs. The stronger property \texttt{CoverAll} is not
needed on this route.}

Fix the scale $s$, ranks $R$, host $H$, constants $C,c_D,c_h$ with
\texttt{Std} $C$ and $C.\texttt{GK.Nonneg}$, and an
admissible word code. A supplied constructed point $P$ includes its actual
ties and extra data. Assume \texttt{Val}, \texttt{Pt.From} and
\texttt{CanonInst}, as in the source. Write $T=H.\texttt{TT}$ and
$\Delta=\lceil x^{c_D}\rceil_{\mathbb N}$,
$h_0=2\lceil x^{c_h}\rceil_{\mathbb N}$, and $w''_n=b'/(xn)$.
For integer functions $a,b$ on $R.A$, let
\[
 \mathcal E_H(a,b)=
 \sum_{p\in\mathcal C_H}\frac{\log p}{p}
 \left(2\mathbf1_{\{\Sigma_{H,p}(a,b)>x^{-\epsilon}\}}
       +4\Sigma_{H,p}(a,b)\right).                 \tag{HW.1}
\]
Here $\mathcal C_H$ and $\Sigma_{H,p}$ are the literal host candidate
set and canonical-deficit sum. The source's \texttt{WOK} is
$\mathcal E_H(a,b)\le x^{-9}$; \texttt{WReq} is the point version.

\paragraph{HW1: exact host--point transport.}
The value map on $R.A$ is injective. Thus \texttt{Pt.From} identifies
the selected sets, Frobenius steps, $q,d$, and prime classifications,
not just their cardinalities. The actual numerical bases agree because
\[
 \mathrm{Sc}=4k\Delta,\qquad \eta=mh_0.             \tag{HW.2}
\]
The two definitions of the universal admissible-deficit set
$\mathcal K(p,n',T',a',b')$ are identical after these substitutions.
The point's extra record gives
$\kappa_I=\inf\mathcal K$; separately, \texttt{CanonInst} gives the
canonical orbit selector $\mathrm{bs}=\mathrm{bs}_H$.
Consequently the representatives, local deficits, and every summand in
(HW.1) agree. No choice of a minimizing certificate is involved.

If GW5 supplies \texttt{WeightVectorExists} at $P$, transfer its chosen
functions along this identification. The box has value zero outside
$T$. Its integer endpoints are exactly the inequalities
\[
 a\le a_t\le a+\frac{a}{kx^{13}},\qquad
 b'\le b_t\le b'+\frac{2n\delta x^\epsilon\Lambda b'}k
 \quad(t\in T).                                  \tag{HW.3}
\]
Hence it supplies the host's \texttt{WinOK} and \texttt{WOK}.
Conversely, if these host functions are the degrees in
\texttt{Pt.From}, their zero extensions are the point's actual degrees,
and host \texttt{WOK} gives actual \texttt{WReq} with the same constants.

\paragraph{HW2: choose the pair, then reconstruct every position.}
Assume the actual point constructor is available for \emph{each}
admissible pair and every $u_0<n$, uniformly beyond its supplied
threshold. First use any admissible pair to build temporary points and
kits. With the supplied \texttt{NumAt} records, the family constructor
produces a temporary family to which GW5 applies. Its window membership
gives an index at every sufficiently large allowed scale. HW1 transfers
the selected pair from that point back to its original host. This is the
degree-selection assertion of \texttt{HostWW}.

Keep this new pair fixed and invoke the point constructor again, for
each $u_0\in\{0,\ldots,n-1\}$. These points form an actual kit.
For completeness, its less immediate equality fields can be checked
directly. Their canonical $\kappa_I$ agree because $k,s,\mathrm{Sc}$
determine the same admissible-deficit sets; for every $p\in\mathbb N$
their exponents are the same
$\lceil C_e f k^2\log_p a\rceil_{\mathbb N}$, with the source's
value zero at $p=0,1$. At a lifted level $1\le j\le m$,
the filling equation and the common Hodge profile give
\[
 \operatorname{off}(j,(0,t))=
 \begin{cases}\lfloor b'/x\rfloor,&t\in T,\\
               \lceil w''_n\rceil,&t\notin T
 \end{cases}
 +(j-1)h_0+\frac{h_0}{2}.                         \tag{HW.4}
\]
Outside the lifted levels \texttt{CanonInst} prescribes $b'/x$ or
$w''_n$ directly. The remaining host equalities come from
\texttt{Pt.From}; the common code identifies the word labels. Indexing
by $u_0$ gives \texttt{Kit.cover}, while the fixed pair gives degree
equality at every common value label. HW1 and (HW.3) give
\texttt{WallPt} at \emph{every} reconstructed point.

All constants and construction thresholds are fixed before the scale;
the new pair may depend on that scale and host. The points for all
$u_0$ are chosen only after the pair. Thus the argument does not require
an already fixed family to contain every admissible degree pair.

\paragraph{HW3: actual family transport and an existing carrier.}
Construct \texttt{famOfKits} from these kits and their supplied
\texttt{NumAt} records. It copies each point's $D,L,G_x$, ties and extra
data; its common fields are the kit equalities just proved. Therefore
\texttt{WallPt} holds at every family index. The identity
$2\mathbf1_E=(2\text{ if }E\text{ else }0)$ converts its loss to the
literal \texttt{ChW\_at}; beyond the family threshold this proves
\texttt{ChosenWall}.

For an allowed $x$, choose an existing index $i_0$ using
\texttt{Ix\_win}. Transfer its actual degrees to the common integer
labels $T_n(x)$, assigning zero elsewhere. The point windows and the
family's $b_T\le b_{\rm top}$ give the required admissibility. The
index $i_0$ itself carries this pair, and every other carrier satisfies
\texttt{ChW\_at} because every family member does. This is exactly
\texttt{ChosenWS}. Empty index sets outside the scale window cause no
problem. BO/BS/HC below supply the biased tops and
\texttt{Frame.Shares}, and expand the conditional host assembly.
The underlying field and analytic inputs, together with the actual
point/numerical construction, remain separate obligations.
\noindent\textit{Primary source route.}\cite{host-degree-realization}


\section{Good-orbit estimates for the actual biased selection}
\label{sec:biased-orbits}

\status{This proves the nonconstant Bernoulli estimate needed by the
arithmetic top-set construction. It allows an arbitrary inclusion profile
between $3\delta/4$ and $\delta$, keeps the actual parity and forced-point
conditions, and supplies the weighted expectation used in the subsequent
finite selection. It does not construct an arithmetic host.}

Fix $\epsilon>0$ and real constants $C_o,L_b$. Put
\[
\begin{gathered}
 \theta=\epsilon/8,\quad h=2+8\epsilon,\quad
 K=\lceil(h+2)/\theta\rceil_{\mathbb N},\quad r=K+1,
 \quad\delta=x^{-\theta},\\
 u=\lceil\delta^{-1}\rceil_{\mathbb N},\quad
 \ell=(6r+1)u,\quad\pi_0=10^{-(6r+1)},\\
 M=\lceil(h+3)\log x/\pi_0\rceil_{\mathbb N},\quad
 G=(M+5)\ell,\\
 N=\left\lceil\max\{6eG\delta,(h+3)\log_2x\}\right\rceil_{\mathbb N}+1.
\end{gathered}                                                   \tag{BO.1}
\]
These are the literal \texttt{KE}, \texttt{rE}, \texttt{goodU},
\texttt{orbitSmallMultiplicity}, \texttt{goodG'} and
\texttt{orbitWindowCount}. Choose each coordinate independently, with
probability $q_i\in[3\delta/4,\delta]$.

\paragraph{BO1: the finite product law and count tails.}
On a finite set $S$, a subset $U$ has weight
$\prod_{i\in U}q_i\prod_{i\in S\setminus U}(1-q_i)$.
Summing over any unused coordinates multiplies by
$\prod_i(q_i+1-q_i)=1$. Consequently conditioning on a selected
coordinate $i_0$ leaves the other coordinates independent with their
original probabilities. Events on disjoint sets still factor.

For count tails, compare each coordinate with a Bernoulli($\delta$)
coordinate using joint states $(0,0),(0,1),(1,1)$ of probabilities
$1-\delta,\delta-q_i,q_i$. Padding to $g$ coordinates gives, for
$|S|\le g$,
\[
 \Pr(|U|\ge b)\le \binom{g}{b}\delta^b,
 \qquad
 \Pr(|U|\ge b)\le
 B(g,\delta,b):=\sum_{j=b}^{g}\binom{g}{j}\delta^j(1-\delta)^{g-j}.
                                                               \tag{BO.2}
\]
The first bound also follows by counting selected $b$-subsets. This
comparison is used only for increasing count events, not for marked tiles.

\paragraph{BO2: the biased marked-tile probability.}
A complete tile has $6r+1$ consecutive cells of length $u$. Its
odd-numbered cells must be empty; every even-numbered cell must contain
exactly one selected point, whose \emph{coordinate} is even.
Any such cell contains at least $\lfloor u/2\rfloor$ even coordinates.
The corresponding event probabilities are therefore at least
\[
 A_0=(1-\delta)^u,\qquad
 A_1=\lfloor u/2\rfloor\frac{3\delta}{4}(1-\delta)^{u-1}.
\]
The singleton events are disjoint, and all cells are disjoint, so the
actual marked probability is at least the source's
\[
 \pi=A_0^{3r+1}A_1^{3r}=\texttt{markedLow}(\delta,r,u).
                                                               \tag{BO.3}
\]
If $0<\delta\le1/100$, then $1\le u\delta\le1+\delta$ and
$\log(1-\delta)\ge-\delta/(1-\delta)$ give
$A_0\ge e^{-103/100}>3/10$. Also
$\lfloor u/2\rfloor\delta\ge99/200$, whence
$A_1\ge891/8000>1/10$. Thus
\[
 0\le\pi_0\le\pi\le1.                                        \tag{BO.4}
\]
This is the additional calculation needed for the biased law.

\paragraph{BO3: all three failures, with one point forced.}
Use the source's cyclic windows
$W(i,L)=\{(i+j)\bmod n:0\le j<L\}$. A set is good when all
$L_2$-windows contain at most $K$ points, every $M+1$ successive
complete tiles contains a marked tile, and all $G$-windows contain at
most $N$ points. Complete tile indices are read modulo
$t=\lfloor n/\ell\rfloor$.

Assume $(M+1)\ell\le n$ and force $i_0<n$ to be selected. A failed
short window still contains at least $K$ unforced selected points.
For each starting tile, its $M+1$ tiles are distinct; remove the at most
one containing $i_0$. At least $M$ independent, unaffected tiles remain.
If $i_0$ is in the incomplete remainder, none is removed. A failed
long window implies at least $N-1$ unforced selected points. Union bounds
over the $n$ starts, using $t\le n$, prove the exact source bound
\[
 \Pr(\mathrm{bad}\mid i_0\text{ selected})\le
 n\binom{L_2}{K}\delta^K+n(1-\pi)^M+nB(G,\delta,N-1).
                                                               \tag{BO.5}
\]
The long-window threshold is deliberately weakened to the source's
$N-1$. No cyclic rotation of the parity condition is used.

\paragraph{BO4: one explicit uniform threshold.}
Set
\[
\begin{gathered}
 C_1=(h+3)/\pi_0+2,\quad C_*=2(6r+1)C_1,\\
 x_0=\exp\max\{1,\,4C_*/\theta^2,\,\log(100)/\theta\},\qquad
 C_b=(|C_o|+1)\bigl((\max\{L_b,1\})^K+2\bigr).
\end{gathered}                                                   \tag{BO.6}
\]
For $x\ge x_0$, writing $y=\log x$ gives
$C_*(y+1)\le2C_*y\le(\theta y)^2/2\le e^{\theta y}$.
The exact ceiling bounds now yield
\[
 (M+1)\ell\le C_*(\log x+1)x^\theta
             \le x^{2\theta}.                                  \tag{BO.7}
\]
Hence the tile condition holds whenever
$x^{\epsilon/4}\le n\le C_ox^h$. If $L_2\le L_b$, the three terms
in (BO.5) are respectively at most
\[
 (|C_o|+1)(\max\{L_b,1\})^Kx^{-2},\qquad
 (|C_o|+1)x^{-3},\qquad (|C_o|+1)x^{-3}.                         \tag{BO.8}
\]
Indeed $\theta K\ge h+2$ and $\pi_0M\ge(h+3)\log x$.
For the last term put $b=N-1$. The factorial estimate
$b!\ge(b/e)^b$ and (BO.2) give
$B(G,\delta,b)\le6^{-b}\le e^{-b}$, since $b\ge6eG\delta$;
also $b\ge(h+3)\log_2x\ge(h+3)\log x$. The empty tail is harmless.
Consequently the conditional probability is at most $C_bx^{-2}$,
uniformly over every allowed length, profile, base and forced point.

\paragraph{BO5: the literal ambient-orbit and weighted statement.}
Let $A$ be a finite abelian group and choose its points independently
with probabilities $q(a)\in[3\delta/4,\delta]$. Fix a step $v$ of
order $n$, a point $a$, and any base $b$ with $a=b+i_0v$.
The map $i\mapsto b+iv$, $0\le i<n$, is injective. Marginalizing
outside its image and then factoring out $q(a)$ proves exactly
\[
 \Pr(a\in T\text{ and its based orbit is bad})
 =q(a)\Pr(\mathrm{bad}\mid i_0\text{ selected})
 \le\delta C_bx^{-2}.                                          \tag{BO.9}
\]
The base may depend on both the step and $a$, as the source's
$\mathrm{bs}(v,a)$ does. For a finite set of steps $v_p$ satisfying
the length bounds and arbitrary nonnegative weights $\omega_p$,
linearity gives
\[
 \mathbb E\sum_p\omega_p
   \#\{a\in T:\text{the prescribed based orbit is bad}\}
 \le \delta |A|C_bx^{-2}\sum_p\omega_p.                         \tag{BO.10}
\]
No independence between the different bad-orbit events is required.
This supplies the probability input to the simultaneous coset selection;
that selection and arithmetic realization are separate steps.
\noindent\textit{Primary source route.}\cite{biased-orbit-probability}


\section{One selected set with biased coset shares and small bad mass}
\label{sec:balanced-tops}

\status{BS1--BS6, using BO's biased orbit estimate, prove the literal
\texttt{topsGroup} conclusion. They produce one set satisfying size,
coset-share and weighted badness requirements together. BS7 explains
the character weights and the ordinary-RH specialization. Arithmetic
host and point construction beyond these finite-selection properties
remain separate.}

\paragraph{BS1: all coset counts share one variance budget.}
Let $A$ be a finite set, $c:A\to C$ a map to a finite set of classes,
and choose its points independently with probabilities
$0\le p_a\le\delta\le1$. Write $T$ for the random subset and
\[
 Y_j=|T\cap c^{-1}(j)|-\sum_{a\in c^{-1}(j)}p_a.
\]
Independence gives, for arbitrary real $g_a,b$,
\[
 \mathbb E\left(b+\sum_a g_a\mathbf1_{\{a\in T\}}\right)^2
 =\left(b+\sum_a g_ap_a\right)^2+
       \sum_a g_a^2p_a(1-p_a).
                                                        \tag{BS.1}
\]
Consequently
$\mathbb E\sum_jY_j^2=\sum_ap_a(1-p_a)\le\delta|A|$.
There is no factor equal to the number of classes.

\paragraph{BS2: meet the count and badness conditions simultaneously.}
Let $I$ be finite, $\omega_i\ge0$, and let
$\mathrm{Bad}(i,a,T)$ be any predicates. Assume
\[
 \mathbb P(a\in T,\ \mathrm{Bad}(i,a,T))\le\beta\delta
 \quad(i\in I,\ a\in A),\qquad \beta\ge0.
\]
Put
\[
 B(T)=\sum_{i\in I}\omega_i
         |\{a\in T:\mathrm{Bad}(i,a,T)\}|,\qquad
 M=\beta\delta|A|\sum_{i\in I}\omega_i.
\]
Then $\mathbb EB\le M$. If $D>0$ and $D^2\ge2\delta|A|$,
some positive-probability subset satisfies simultaneously
\[
 |Y_j|\le D\quad(j\in C),\qquad B(T)\le4M.             \tag{BS.2}
\]
For $M>0$, average the nonnegative quantity
$F(T)=D^{-2}\sum_jY_j^2+B(T)/(4M)$.
Its expectation is at most $3/4$, so a value at most $3/4$
exists in the probability support and gives (BS.2).
If $M=0$, nonnegativity forces $B=0$ on that support; choose
there a point with $\sum_jY_j^2\le D^2/2$.
No independence among the different badness predicates is used.

\paragraph{BS3: convert biased counts into the actual shares.}
Now let $A$ be a finite abelian group, $A_0\le A$,
$o=[A:A_0]$, $n_0=|A_0|$ and $N=on_0$.
Use probabilities $\delta$ on $A_0$ and $3\delta/4$ outside it.
Each coset has size $n_0$. If its count $C_\gamma$ differs from its
mean by at most $D$, then, for $S=|T|$,
\[
 \left|S-\frac{3o+1}{4}\delta n_0\right|\le oD.
                                                        \tag{BS.3}
\]
Assume $0<\delta$, $0\le\eta\le1$ and
$4oD\le\eta\delta n_0$. Equation (BS.3) gives the slightly
stronger size bounds
\[
 \frac34\delta N\le S\le\frac54\delta N.
                                                        \tag{BS.4}
\]
Set $\lambda_\gamma=(4\text{ if }\gamma\in A_0
\text{ else }3)/(3o+1)$. Since
$\lambda_\gamma(3o+1)\delta n_0/4=\mathbb EC_\gamma$,
\[
 |C_\gamma-\lambda_\gamma S|
 \le D+\lambda_\gamma oD\le3D
 \le\frac{3\eta\delta n_0}{4o}
 \le\eta S.                                            \tag{BS.5}
\]
These are exactly the source's coset shares, with its weaker
$\delta N/2\le S\le2\delta N$ also supplied.

\paragraph{BS4: one explicit scale condition suffices.}
Fix $\varepsilon>0,C_o,L_b$ before the group and let
\[
 h=2+8\varepsilon,\quad
 \delta=x^{-\varepsilon/8},\quad
 \tfrac12x^h\le N\le C_ox^h,\quad o\le C_o,\quad
 D=\sqrt{2\delta N}.
\]
For $x\ge1$ with $x^\varepsilon\ge64C_o^4$,
\[
 \delta N\ge\tfrac12x^{2+63\varepsilon/8}
 \ge32o^4x^2,
 \qquad
 4oD\le x^{-1}\delta n_0.                            \tag{BS.6}
\]
The last inequality follows by squaring positive quantities and
using $N=on_0$. Thus BS3 applies with $\eta=x^{-1}$,
which is stronger than the required $x^{-1+\varepsilon}$.
If $C_o<1$, the index assumption admits no such group.

\paragraph{BS5: the literal biased selection theorem.}
Fix $L_2\le L_b$, a finite prime list $P$, elements $h_p\in A$
with $\operatorname{ord}(h_p)\ge x^{\varepsilon/4}$, and
nonnegative weights $\omega_p$. The supplied selector satisfies
$b(h,t)+ih=t$ for some $i\ge0$. It is fixed before selecting $T$.
Define
\[
 B_p(T)=\{t\in T:\text{the selected orbit based at }b(h_p,t)
                       \text{ is not good}\}.
\]
The precise good-orbit parameters are BO's
$L_2,K,r,\delta,h,x$. BO5 proves
\[
 \mathbb P(t\in B_p(T))\le C_bx^{-2}\delta,
\]
uniformly in $p,t$ and the indicated group data. In particular, the
base may depend on $t$; no rotation-invariance of the even-index
condition is assumed.
Apply BS2 with $\beta=C_bx^{-2}$ and then BS3--BS4. One set $T$,
with $S=|T|$, satisfies
\[
 \boxed{\begin{gathered}
  \delta N/2\le S\le2\delta N,\\
  |C_\gamma-\lambda_\gamma S|\le
                   x^{-1+\varepsilon}S\quad(\gamma\in A),\\
  \sum_{p\in P}\omega_p|B_p(T)|
       \le4C_bx^{-2}\delta N\sum_{p\in P}\omega_p .
 \end{gathered}}                                      \tag{BS.7}
\]
The constants $C_b,x_0$ depend only on the previously fixed
$\varepsilon,C_o,L_b$. One may take the maximum of BO's threshold,
$1$, and $(64C_o^4)^{1/\varepsilon}$.
All group, selector, prime-list and weight choices follow that
threshold. This is the literal \texttt{topsGroup} statement; the
construction actually retains the stronger error $x^{-1}|T|$.

\paragraph{BS6: the prime-weighted badness allowance.}
For $z\ge2$, put $n=\lfloor z\rfloor$.
The factorization of $n!$ and
$\lfloor n/p\rfloor\ge n/(2p)$ give
\[
 \frac n2\sum_{p\le n}\frac{\log p}{p}
 \le\log(n!)\le n\log n.
\]
Thus $\sum_{p<z}\log p/p\le2\log z$; the sum is empty
when $1<z<2$. For a candidate set below
$X'=\exp(x^{\varepsilon/3})$, use
$\omega_p=\log p/p$ in (BS.7) and $\,\delta N\le2|T|$:
\[
 \sum_{p\in P}\frac{\log p}{p}|B_p(T)|
 \le16C_b|T|x^{-2+\varepsilon/3}
 \le |T|x^{-2+\varepsilon},                          \tag{BS.8}
\]
once $x^{2\varepsilon/3}\ge16C_b$. This keeps the same
set's coset shares and size bounds.

\paragraph{BS7: why the bias produces the GRH character weights.}
On $Q=A/A_0$, the target probability measure is
\[
 \mu=\frac{3o}{3o+1}\,u_Q+\frac1{3o+1}\,\delta_0,
                                                        \tag{BS.9}
\]
where $u_Q$ is uniform. For a character $\chi:Q\to\mathbb C^\times$,
translation of its finite sum shows $\sum_{q\in Q}\chi(q)=0$
unless $\chi=1$. Therefore
\[
 \sum_{q\in Q}\mu(q)\chi(q)=
 \begin{cases}1,&\chi=1,\\(3o+1)^{-1},&\chi\ne1.\end{cases}
                                                        \tag{BS.10}
\]
Under the shares bound with error $\eta$, the empirical normalized
character average differs from this value by at most $o\eta$:
sum the $o$ count errors and use $|\chi(q)|=1$.
For a cyclic character level, the nontrivial coefficients are precisely
\texttt{Level.wt}; \texttt{Level.Gm} supplies the trivial term separately
with coefficient one. Its \texttt{Level.ofSub} has no exceptional
distinguished character.
This explains the weights, while the arithmetic character
identification and explicit formula remain their own arguments.
For ordinary RH, $o=1$ gives $A_0=A$; the shares error is identically
zero and the selection is uniform Bernoulli.
\noindent\textit{Primary source route.}\cite{balanced-tops}


\section{From the selected tops to the actual arithmetic host}
\label{sec:host-tops}

\status{HC1--HC10 prove the finite, numerical and algebraic joins below.
The exact analytic inputs \texttt{ShortNorm}, \texttt{SmallOne} and
\texttt{ThinH}, and the supplied character field and pair window, remain
explicit hypotheses. In particular, constructing a host from its supplied
\texttt{HostAx} data does not construct those field data.}

\paragraph{HC1--HC2: the actual top-set choice.}
Write $n=|A|$, $o=[A:A_0]$, $n_0=|A_0|$, so $n=on_0$.
Write $Q=QS$ and $N=Qqd$ for the host moduli, and
$r=k+m-1$ for its rank parameter.
Fix $0<\varepsilon<1/52$ and $C_0\ge1$, with
$C_o=\max(C_0,2)$, before the scale or field. The actual size data give
\[
 \begin{gathered}
 \delta=x^{-\varepsilon/8},\quad h=2+8\varepsilon,
 \quad x^h\le n\le2x^h,\quad o\le C_0,\\
 k\le2x^{9/2+10\varepsilon},\quad
 x^{1/J}\le Y\le x,\quad 1\le J\le C_0.
 \end{gathered}                                      \tag{HC.1}
\]
The supplied \texttt{PairRow} says every prime
$p\le Y^{\varepsilon/5}$ splits in $\mathbb Q(\sqrt{-d})$.
An inert-class prime therefore satisfies
$p>x^{\varepsilon/(5C_o)}$. This also excludes $2$, which belongs
to the host's designated small-prime set. For $x\ge16$,
\[
 \begin{gathered}
 \log((kx^{26})^4)
 \le4\log2+(122+40\varepsilon)\log x<124\log x,\\
 L_2=2\sup_{p\,\mathrm{inert}}
       \left\lceil\log_p((kx^{26})^4)\right\rceil
 \le2\left(\frac{620C_o}{\varepsilon}+1\right).
 \end{gathered}                                      \tag{HC.2}
\]
The supremum is over natural numbers; the displayed common upper
bound proves boundedness and also covers an empty indexing set.

Apply BS5 with this last upper bound. Use the actual Frobenius steps
$h_p$ and the actual finite candidate set
\[
 P=\{p<X':p\text{ prime, inert-class},\
             \operatorname{ord}(h_p)\ge x^{\varepsilon/4}\},
 \qquad X'=\exp(x^{\varepsilon/3}).
\]
For each $t$, retain the prescribed base $b(h_p,t)$ satisfying
$b(h_p,t)+ih_p=t$ for some $i\ge0$. The badness predicate uses
the indices $0,\ldots,\operatorname{ord}(h_p)-1$ and their actual
even parity. Its base may depend on $t$. On these inert candidates
the source's \texttt{stepP} is exactly $h_p$. Thus BS5 supplies one
set $T=TT$ such that
\[
 \begin{gathered}
 \delta n/2\le |T|\le2\delta n,\\
 \bigl||T\cap(\gamma+A_0)|-\lambda_\gamma|T|\bigr|
       \le x^{-1+\varepsilon}|T|,\\
 \sum_{p\in P}\frac{\log p}{p}|B_p(T)|
       \le4C_bx^{-2}\delta n\sum_{p\in P}\frac{\log p}{p},\\
 \lambda_\gamma=
 \frac{4\mathbf1_{A_0}(\gamma)+3\mathbf1_{A\setminus A_0}(\gamma)}{3o+1}.
 \end{gathered}                                      \tag{HC.3}
\]
These are literally the size, \texttt{Frame.Shares} and
\texttt{BadOK} predicates. The source update $H[T]$ changes only
$H.TT$. Its field, \texttt{HostAx}, Frobenius law, prime classes,
$L_2$ and $P$ remain identical. Hence the data used to choose $T$
remain the data of the resulting host. The constants $C_b,x_0$
precede all these host choices.

\paragraph{HC3: transport preserves the actual bad sum.}
A point with \texttt{Pt.From} identifies its ranks, field data,
prime classes, and every Frobenius step with those of the host.
Equality of the $TT$ coordinate images, together with the
injectivity of those coordinates, identifies the actual sets.
With the canonical selector, each bad-top predicate is therefore
identical, term by term. This proves \texttt{chBad\_of\_from};
the size and share statements transfer by the same identification.
Every supplied family assembled from such points inherits
\texttt{ChosenTT} and \texttt{ChosenBad} with one threshold.
This does not supply its point or family construction.
For the same set, BS6 also gives
\[
 \sum_{p\in P}\frac{\log p}{p}|B_p(T)|
 \le16C_b|T|x^{-2+\varepsilon/3}
 \le |T|x^{-2+\varepsilon},                         \tag{HC.4}
\]
once $x^{2\varepsilon/3}\ge16C_b$.

\paragraph{HC4: what the host witness actually needs.}
A supplied \texttt{HostAx} already includes a nonempty CM-type and
an ideal $\mathfrak b$ satisfying
\[
 \mathfrak D_{E/K}=\mathfrak b\overline{\mathfrak b},
 \qquad \mathfrak b+\overline{\mathfrak b}=\mathcal O_E,
 \qquad \sigma\mathfrak b=\mathfrak b\quad
                   (\operatorname{gal}(\sigma)\in\Psi).       \tag{HC.5}
\]
Choose an embedding in that CM-type, the displayed ideal, and the
square root of $-q$ furnished by \texttt{FieldRows}. Copy them,
the supplied choices, \texttt{FieldRows}, and $q,d>4tr+1$ into
\texttt{HostNF}; its definition gives a \texttt{Host} with exactly
the original \texttt{HostAx} and choices. This proves the entire
source predicate named \texttt{admitted}. That name is not a
Lean placeholder. The substantive field facts (HC.5) remain
part of the supplied input. For the actual construction choose
$\mathfrak c_K=\mathcal O_K$: its norm is $1$, and it is primitive
because an integer $m>1$ cannot be a unit in $\mathcal O_K$.

\paragraph{HC5: the actual auxiliary primes.}
The analytic input \texttt{ShortNorm} gives constants $A,c,c'>0$
and a lower bound
\[
 \#\{p\in(X,X(1+u)]:p\text{ prime},\ \operatorname{NormOne}(d,M,p)\}
 \ge\frac{c'uX}{(4dM^2)^A\log X},                    \tag{HC.6}
\]
under its explicit hypotheses
$d\ge100$ prime, $d\equiv3\pmod4$, $M\ge3$, $X\ge1$,
$\log(4dM^2)\le c\sqrt{\log X}$ and
$e^{-c\sqrt{\log X}}\le u\le1$.
Here $\operatorname{NormOne}$ means that integers $a,b_1$ satisfy
\[
 a^2+db_1^2=4p,\qquad M\mid b_1,\qquad 2M\mid a-2-b_1.
                                                               \tag{HC.7}
\]
It is not replaced by a mere splitting condition.
If $d\mid N$ and $\log N\le3x^{\varepsilon/5}\log^2x$, the
large-$x$ consequences of (HC.6), uniformly for
$\log Z\ge8x^\varepsilon$, give at least $Z^{9/10}$ such
primes in $(Z,2Z]$. Indeed, with $v=x^{\varepsilon/5}$,
$\log(4dN^2)\le10v^2$ and $\log Z\ge8v^5$;
the denominator in (HC.6), with $u=1$, is at most
$c'Z^{1/10}$ once $x$ is large.

For the literal exceptional set
\[
 BB(r,g)=\{g\}\cup
 \{\ell:\exists\,1\le j\le r,\ \ell\text{ prime},
                      g^j<\ell^{4v_\ell(g^j-1)}\},
\]
there are at most $3r+1$ elements. Four exceptional prime-power
factors for one $j$ would already have product greater than $g^j$.
For fixed $\ell,j$, take the minimum valuation $w$ among candidates
in $[Z,2Z]$. Then $Z<(\ell^w)^4$; there are at most $2j$
roots of $z^j=1$ modulo $\ell^w$, with at most
$Z^{3/4}+1$ candidates per residue. Consequently
\[
 \begin{gathered}
 \#\{g\in[Z,2Z]:BB(r,g)\cap BB(r,g_0)\ne\varnothing\}
 \\
 \le(3r+1)[1+2r^2(Z^{3/4}+1)]\le20r^3Z^{3/4}.
 \end{gathered}                                      \tag{HC.8}
\]
Take $Z=(kb)^8$, $b=e^{x^\varepsilon}$, and $r\le x^5$.
Eventually $20r^3<Z^{3/20}$, so two of the norm primes have
disjoint $BB$ sets.

The separate analytic input \texttt{SmallOne} gives an exponent
$a_0\ge1$ and a prime $1\bmod M$ below $M^{a_0}$ for every
$M\ge2$. Apply it to $15N$, then $15Np$, then $Npp'$.
This gives distinct primes $p,p',q_0$, with
\[
 \begin{gathered}
 p,p'\equiv1\pmod{15N},\quad q_0\equiv1\pmod N,\qquad
 p,p',q_0\le N^{C_1},\\
 C_1=2a_0^3+4a_0^2+a_0.
 \end{gathered}                                      \tag{HC.9}
\]
Eventually $N^{C_1}<b<(kb)^8$, so these also differ from the
two norm primes. The five required level primes have now been
chosen with their exact separation conditions.

\paragraph{HC6: the second set is a different finite construction.}
For $L\ge2$, there is $S\subset\mathbb N$ with $4Ls<n_0$,
no relation
$j(s_2-s_1)=i(s_3-s_1)$ among distinct members with
$0<|i|,|j|\le L$ and $i\ne j$, and
\[
 |S|\ge \frac{n_0}{L}
             e^{-32\sqrt{\log n_0\log L}}.           \tag{HC.10}
\]
Here is the construction. Put $x_0=\log n_0$, $y=\log L$,
$s=\sqrt{x_0y}$, $d_0=\lfloor s/y\rfloor$,
$a=(x_0-y-3)/d_0$ and
$d=\lfloor e^a/(4L)\rfloor$ when the right side exceeds $1$;
otherwise $\{0\}$ suffices for $n_0>0$.
A sphere in $\{0,\ldots,d-1\}^{d_0}$ contains at least
$d^{d_0}/(d_0d^2)$ vectors. Encode them in base $4Ld$.
Any indicated relation has coefficients of absolute value below
the base, hence holds coordinatewise; three distinct collinear
points cannot lie on a sphere. The inequalities
$d_0\ge32$, $y+7\le a\le2s$ and
$3+d_0(7+y)+d_0+2a\le32s$ give (HC.10).

The actual $L=L_w$ satisfies
$2\le L_w\le(4084+60/\varepsilon)\log x$.
Since $n_0$ grows like $x^{2+8\varepsilon}$ up to fixed constants,
putting $E=e^{32\sqrt{\log n_0\log L_w}}$ gives, uniformly,
$E\ge2$ and $4C_oL_wE\le x^\varepsilon$ for large $x$.
The lower bound before erasure is then at least
$4x^{-\varepsilon}n\ge4$. Erase $0$; losing at most one element
gives the actual second set $T_2$ with
\[
 \begin{gathered}
 |T_2|\ge\frac{n_0}{L_w}
       e^{-64\sqrt{\log n_0\log L_w}},\qquad
 |T_2|\ge x^{-\varepsilon}n,\qquad
 4L_w|T_2|\le n_0.
 \end{gathered}                                      \tag{HC.11}
\]
The last inequality counts positive integers below $n_0/(4L_w)$.
Thus the source's absolute constant $C_{T_2}$ can be $64$.

\paragraph{HC7--HC8: the remaining elementary host joins.}
The exact rank formulas imply $r\le x^5$ and $t\le x^{400}$
for large $x$. For example, setting $y=x^{1/100}\ge2$ gives
\[
 \begin{gathered}
 f\le4y^{216},\quad k\le2y^{470},\quad r\le4y^{470},
 \quad m\ge y^{49},\quad L^2\le y^4,\\
 t\le10708y^{39967}\le y^{40000}.
 \end{gathered}
\]
The pair window gives $\log q>Y^{\varepsilon/5}/2$ and $q\mid d+1$.
These imply $q,d>4tr+1$ and $d\ge100$ once
$814\log x\le x^{\varepsilon/(5C_0)}$.
If $Q\le C_0n$, $\ell=\log Y\ge\max(2,41000C_0/\varepsilon^2)$
and $z=Y^{\varepsilon/5}$, then
\[
 \begin{gathered}
 820C_0\ell\le z,\qquad
 \log Q\le2C_0+\tfrac94\log x\le z(\ell^2-1),\\
 \log(Qqd)\le2z\ell^2.
 \end{gathered}                                      \tag{HC.12}
\]
The prescribed splitting congruences give \texttt{PairRow}, and
(HC.12) gives exactly the modulus bound needed for HC5.

Writing $b_1=Nu$, $a-2-b_1=2Nv$ in (HC.7) gives
\[
 \begin{gathered}
 \pi=\frac{a+b_1\sqrt{-d}}2
    =1+N\left(v+u\frac{1+\sqrt{-d}}2\right)\in\mathcal O_K,\\
 \pi\bar\pi=g,\qquad \pi\equiv1\pmod N,
 \qquad K=\mathbb Q(\sqrt{-d}).
 \end{gathered}                                      \tag{HC.13}
\]
Expanding the norm also gives $g\equiv1\pmod N$.
For any chosen level prime $l\equiv1\pmod N$, the supplied
Artin/Frobenius law gives trivial Frobenius and inertia. Its
decomposition group is trivial, so it splits into $[E:\mathbb Q]$
primes; quadratic reciprocity also places it in \texttt{splitK}.

The analytic \texttt{PairGood} input bounds the nontrivial twists
$\xi(\cdot/q)^i(\cdot/d)^j$ and gives their zero-free rectangles.
It yields the three \texttt{QDq} characters, the six
\texttt{Boxq} characters, and $|dlogH|\le4\log x$ directly.
For the other real characters, CRT and the unit groups at prime
powers give the precise decomposition
\[
 \begin{gathered}
 \chi=\rho(\cdot/q)^i(\cdot/d)^j,\qquad
 \operatorname{cond}(\chi)=\operatorname{cond}(\rho)q^id^j,\\
 \operatorname{cond}(\rho)\mid8Q_{\rm odd}.
 \end{gathered}                                      \tag{HC.14}
\]
Indeed a real character kills the odd-order principal units and
the squares among the dyadic units. Since
$Q\mid\operatorname{lcm}(2^{e+2},M)$ implies $Q_{\rm odd}\mid M$,
\texttt{PairGood.boxes} applies to (HC.14). Whenever $q$ or $d$
divides the conductor, $(i,j)\ne(0,0)$, giving exactly
\texttt{BoxH} and \texttt{RealBox}.

\paragraph{HC9: the exact conditional host theorem.}
Retain \texttt{ThinH}: under $Q_{\rm odd}\le C_0$, (HC.12) and
\texttt{RealBox}, every relevant primitive host character obeys
\[
 |g_\chi(y)|\le y e^{-c_6x^{\varepsilon/8}}
 \quad\left(y\ge e^{x^{\varepsilon/3}}/3\right).       \tag{HC.15}
\]
Given the three analytic inputs, choose
$C_{T_2}=64$, $C_8=\max(C_0,C_1)$, $C_5=4$, and the constants
$c_6,C_b$ before the host. Take one maximum of the preceding
thresholds. For every supplied \texttt{HostAx}, ranks and pair
window satisfying (HC.1), $Q\le C_0n$,
$0<M\le C_0$, $Q\mid\operatorname{lcm}(2^{e+2},M)$, and
\texttt{FieldRows} at $\lfloor Y^{\varepsilon/5}\rfloor$,
HC4--HC8 construct all choices and the host. Install HC2's $TT$.
The resulting host has the original \texttt{HostAx}, the exact
five-prime level set, integral norm elements with their congruences,
both second-set bounds, \texttt{QDq}, \texttt{BoxW},
$|dlogH|\le x^\varepsilon$, \texttt{ThinOK} and \texttt{BadOK}.
These are all conjuncts of
\texttt{HostGood}$(s,C_8,4,c_6,C_b,H,q_0)$.
Its shares hold separately. The expanded proof lists every
\texttt{FieldRows} assumption and checks every conjunct.

\paragraph{HC10: character windows and the ordinary specialization.}
For a primitive character $\psi$, the source field construction
supplies, for $e\ge e_0$, a host field with
$n=c2^e$, $o=\operatorname{ord}(\psi)$,
$Q\mid\operatorname{lcm}(2^{e+2},M_\psi)$,
$Q\le4M_\psi n$, and the exact kept-character, kernel and twist
data. This construction, the uniform rank constructor, the analytic
\texttt{PairEx}, and the cohomological \texttt{VWinAt}$(\psi,401)$
remain inputs here.

Their elementary window join is explicit. Put $\kappa=c^{1/h}$.
The supplied cohomological interval $[W,W^2]$, with
$W\ge\max(\kappa^3,\kappa^{400})$, gives intervals
$[Z,Z^{3/2}]$ for $Z=\kappa W$. Writing $x=\kappa x'$ gives
\[
 x^h\le c2^{e(x')}\le2x^h,\qquad
 x^{400}\le(x')^{401},\qquad W\le x'\le W^2.        \tag{HC.16}
\]
Thus the dyadic field size and the required cohomological depth
hold together. Take $C_0=J+\operatorname{ord}(\psi)+8M_\psi$ and
a threshold forcing $Y^{\varepsilon/5}\ge5+4M_\psi$ and $e\ge e_0$.
HC9 then supplies the host on every allowed scale, with one fixed
set of constants. The kernel and twist identities identify every
bounded, unkept character with a nontrivial $q,d$ twist, so
\texttt{PairGood} supplies the remaining bounds in \texttt{EAtH}.
Since $t\le x^{400}$, the supplied cohomological bound restricts
to the actual rank. This proves the literal conditional
\texttt{HostExW}$(3/2)$ endpoint of \texttt{supplyOfAdmOf}.

For $o=1$, $A_0=A$ and the shares error is zero. The source's
ordinary host-only specialization takes $M=1$, $J=1$, $Y=x$,
and uses the supplied trivial-character field and rank constructors.
It does not need \texttt{VWinAt} to obtain that host-only statement;
it also does not supply the additional cohomological bound.
The full arithmetic field and point existence questions remain
distinct from these completed selection and transport arguments.
\noindent\textit{Primary source route.}\cite{actual-host-tops}


\section{Frobenius orbits as actual local places}
\label{sec:local-identifications}

\status{LI1--LI8 construct the literal local-place identifications,
completion maps, lattice frames and integral away character for supplied
host and level data. The position identity uses the original coefficient
embedding and every prime above $p$. The remaining global comparison
lattice and class-rank theorems are separate inputs.}

\paragraph{The actual arithmetic setting.}
The host has
$G=\operatorname{Gal}(E/\mathbb Q)=(\{\pm1\}^2)\times A$,
$c=((-1,-1),1)$, $F=E^{\langle c\rangle}$ and
$\Psi=\{((\epsilon,1),a)\}$, so $G=\Psi\times\langle c\rangle$.
Write $\sigma_g$ for its action on $E$ and
$\phi_g=\phi_0\sigma_g:E\to\mathbb C$.
Assume $p$ is unramified in $E$ and its arithmetic Frobenius at every
prime above $p$ is
\[
 f=((-1,-1),h)=zc,\qquad z=((1,1),h)\in\Psi,
 \qquad n=\operatorname{ord}(h)\ge2\text{ is even}.       \tag{LI.1}
\]
Put $W_n=W(\mathbb F_{p^n})$, $Q_n=\operatorname{Frac}(W_n)$,
and let $\operatorname{Fr}$ reduce to $x\mapsto x^p$.
LI7 below checks these hypotheses from the literal \texttt{GIdx}
and the supplied host Frobenius law.

\paragraph{LI1--LI2: construct and classify the local embeddings.}
For an unramified prime $P$ of $E$, its completed integer ring has
uniformizer $p$ and residue field $\mathbb F_{p^m}$. Here is the needed
explicit identification with $W(\mathbb F_{p^m})$.
Choose a residue generator $\alpha$, lift its monic minimal polynomial
to $P_0\in\mathbb Z_p[T]$, and lift its simple root by
$b\mapsto b-P_0(b)/P_0'(b)$. The error squares at each step.
Successive residue subtraction and division by $p$ show that
$1,b,\ldots,b^{m-1}$ span the completed ring over $\mathbb Z_p$.
Dividing a hypothetical relation by its smallest $p$-power and reducing
modulo $p$ proves independence. The same construction in the Witt ring
gives the required isomorphism and its uniqueness for a prescribed
residue map. Consequently there is an embedding $i:E\to Q_m$ with
\[
 i(\mathcal O_E)\subset W_m,\qquad
 \overline i:\mathcal O_E\twoheadrightarrow\mathbb F_{p^m},
 \qquad\ker\overline i=P.                               \tag{LI.2}
\]
These are exactly the source's \texttt{Good} conditions.

Two such embeddings inducing $P$ differ by a power of Frobenius:
their residue maps do, and simple-root uniqueness lifts that equality.
Normality of $E/\mathbb Q$ also shows
$\operatorname{Fr}i=i\sigma_f$. Thus $f^m=1$, whereas residue
surjectivity and $f^n=1$ give $\operatorname{Fr}^n=1$ on
$\mathbb F_{p^m}$. The finite-field Frobenius has exact order $m$,
so $m=n$. Fix such an $i$ and put $i_g=i\sigma_g$ and
$P_g=\ker(\overline{i_g})$. We have
\[
 P_{g'}=P_g\ \Longleftrightarrow\ g'=gf^s\text{ for some }s.
                                                               \tag{LI.3}
\]
Every prime above $p$ occurs: construct its embedding by (LI.2),
then use normality to compare it with $i$.

The completion comparison is exact. For any embedding satisfying
(LI.2), $i(x)\in W_n$ if and only if $x$ is integral at its prime:
local denominators outside the prime map to units; for a nonintegral
$x$, its inverse maps to the maximal ideal. Residue lifting approximates
every Witt vector by global integers modulo $p^r$. Global integers are
also dense in the local completion, since a denominator outside the
prime is invertible modulo each prime power. Applying the integrality
criterion to $x/p^r$ identifies the quotients at every precision.
Their inverse limits, and then fraction fields, are therefore isomorphic.

\paragraph{LI3: the orbits label all real-subfield primes.}
The evenness in (LI.1) gives $c\notin\langle f\rangle$: an equality
$f^s=c$ would force both $n\mid s$ and $s$ odd. Hence
$P_g\ne P_{gc}$. Let $v_g=P_g\cap\mathcal O_F$.
Given any residue at $P_g$, choose by CRT an integer $b$ attaining it
and vanishing at $P_{gc}$. Then $b+c(b)\in\mathcal O_F$ has that
same residue. Thus $i_g|_F$ still satisfies (LI.2), and
\[
 Q_n\simeq F_{v_g},\qquad W_n\simeq\mathcal O_{F_{v_g}}. \tag{LI.4}
\]
In particular these primes split in $E/F$.

If $v_{g'}=v_g$ and $y\in P_{g'}$, then $yc(y)\in v_g$,
so $y\in P_g\cup P_{gc}$. An ideal contained in a union of two
ideals lies in one; maximality gives $P_{g'}=P_g$ or $P_{gc}$.
Together with (LI.3), this proves
\[
 \boxed{\ \Psi/\langle z\rangle\ \simeq\
          \{v\subset\mathcal O_F:v\mid p\}\ },
 \qquad\#\{v\mid p\}=|\Psi|/n.                         \tag{LI.5}
\]
Surjectivity uses lying over and the already proved surjectivity for
$E$-primes. Use the literal representatives $a_v=\texttt{orbRep}(z,v)$;
the source positions are $a_vz^s$, $0\le s<n$.

\paragraph{LI4: align the original complex and coefficient embeddings.}
Let $L$ be the coefficient number field, $e_C:L\to\mathbb C$, and
$e:L\to K=\texttt{KfRec}$ the prescribed-place embedding from BC2.
The finite extension $K/\mathbb Q_p$ embeds in
$\overline{\mathbb Q}_p\subset\mathbb C_p$; choose $i_K$.
Integrality of monic roots and their unit inverses shows
$i_K^{-1}(\{|x|\le1\})=\texttt{ORec}$.
Extend $i_Ke$ to an embedding
$\lambda:\overline{\mathbb Q}_{\mathbb C}\to\overline{\mathbb Q}_p$,
so $i_K e=\lambda e_C$ on $L$.
With $\iota:Q_n\to K$, normality of $E$ lets us replace $i$ by
one Galois conjugate so that $i_K\iota i=\lambda\phi_0$.
The preceding construction then uses this adjusted $i$. For
$i_v=i\sigma_{a_v}$ it gives
\[
 i_K\iota\operatorname{Fr}^{s}(i_v(x))
   =\lambda\phi_{a_vz^sc^s}(x).                         \tag{LI.6}
\]
This is the literal \texttt{PosCompat}. The factor $c^s$ is essential:
odd positions conjugate $x\in E$. Every argument of $\lambda$ is
algebraic; no extension to transcendental complex numbers is used.

\paragraph{LI5: the splitting and frame use the given lattice.}
Let $e_v:Q_n\to F_v$ be (LI.4), and $j=e_vi_v$. Choose
$0\ne b\in E$ with $c(b)=-b$, as supplied by the host's
$\sqrt{-d}$. Since $1,b$ is an $F$-basis,
\[
 \Phi_v(x\otimes t)=(j(x)t,j(c(x))t)                    \tag{LI.7}
\]
is an isomorphism $E\otimes_FF_v\simeq F_v\times F_v$:
its inverse sends $(u,v)$ to
$1\otimes(u+v)/2+b\otimes(u-v)/(2j(b))$.
It sends scalars diagonally and conjugation to coordinate interchange.

Write $\Lambda=\texttt{Lc.Lb}\subset E^k$ for the supplied finite
full integer lattice. The first projection of the literal localized
lattice is exactly
\[
 \pi_1(\texttt{locL}_v\Lambda)
   =M_v:=\mathcal O_{F_v}\text{-span}\,j(\Lambda).        \tag{LI.8}
\]
Indeed the source integral tensor subring projects into
$\mathcal O_{F_v}$ by integrality, and contains all its diagonal
scalars. The two inclusions in (LI.8) follow from the defining spans.
The module $M_v$ is finite, torsion free and full: an $E$-basis chosen
from $\Lambda$ retains nonzero determinant under $j$.
Choose a DVR basis as the columns of $B_v\in\mathrm{GL}_k(F_v)$.
Then $t\in M_v$ exactly when $B_v^{-1}t\in\mathcal O_{F_v}^k$.
This proves the exact \texttt{fr\_lat} clause. Equations (LI.4)--(LI.8)
give every field of \texttt{LocId}, with its recovered
$\texttt{iw}(v)=i_v$.

\paragraph{LI6: the same positions induce the prescribed valuations.}
Let $\theta_g:E\to L$ be the actual \texttt{tauLL} embedding with
$e_C\theta_g=\phi_g$. Injectivity of $i_K$ in (LI.6) gives
\[
 e\theta_{a_vz^s}(c^sx)=\iota\operatorname{Fr}^s(i_v(x)).
                                                               \tag{LI.9}
\]
For $0\ne x\in F$, BC3 therefore gives the exact normalization
\[
 \frac{\nu_w(\theta_{a_vz^s}(x))}{\nu_w(p)}=v_p(i_v(x)),
 \qquad
 x\in v\ \Longleftrightarrow\
 \nu_w(\theta_{a_vz^s}(x))>0\quad(x\in\mathcal O_F).
                                                               \tag{LI.10}
\]
The last equivalence includes $x=0$, whose valuation is $+\infty$.
Thus ramified coefficient extension changes the integer valuation scale
by the stated factor; it does not alter which prime is represented.

\paragraph{LI7--LI8: the actual application and integral character.}
For \texttt{GIdx}, the class \texttt{inertClass} excludes the conductor.
The supplied host law gives trivial inertia, hence unramifiedness, and
Frobenius $((-1,-1),h)$ at every prime. Its \texttt{HNeZero} clause
makes $\operatorname{ord}(h)$ even and positive. These are precisely
(LI.1), with $n=\texttt{frobOrder}$. LI1--LI6 therefore prove the actual
\texttt{exists\_locId} conclusion, including position compatibility
for every choice of away character.

Here is that character under the literal \texttt{Joint0} and
\texttt{LevelZeroCharacterRational} assumptions. The local subgroup
equality comes from \texttt{Joint0}; the matching Gram identity is
\texttt{Cmp.gramN\_zero}. For an away-level
element $a$, set its components above $p$ to $1$ and retain all other
components, obtaining $K(a)$ in the original level subgroup $U$.
Only finitely many components change; the unitary equations and level
conditions hold componentwise. Injectivity into the local components
makes $K$ a homomorphism. At level zero, the character
$\chi:U\to\mathbb C^\times$ is a product of additive characters of
finite residue groups. Their values are roots of unity. Rationality
and injectivity of $e_C$ give a multiplicative $r:U\to L^\times$
with $e_Cr=\chi$; both $r(\kappa)$ and its inverse are algebraic
integers. Hence
\[
 \psi(a)=e(r(K(a)))^{-1}\in\texttt{ORec}^{\times},
 \qquad \psi(a)e(r(K(a)))=1.                            \tag{LI.11}
\]
Uniqueness of the componentwise truncation proves this for every
$\kappa$ allowed in \texttt{PsiTied}, establishing \texttt{exists\_psi}.

These choices can be made for all supplied words and places. The
following SF, CR and CS sections now prove the common-frame algebra,
class-rank equality and integral comparison span under their exact
compatible-forms hypotheses. The simultaneous global construction
supplying those hypotheses remains distinct from these local
identifications; the existence of the host is not inferred from them.
\noindent\textit{Primary source route.}\cite{actual-local-places}


\section{One global frame for all the local comparisons}
\label{sec:standard-global-frame}

\status{SF1--SF7 prove the common-frame and companion-lattice arguments,
including their specialization to the actual positive spaces and carried
base lattices. They replace independent local basis choices by base changes
of a single global matrix. The complete standard-column objects are still
supplied with their stated carrier properties; their arithmetic existence
is not inferred from these linear-algebra calculations.}

\paragraph{The mathematical issue.}
LI5 constructed a frame separately at every local place. The comparison
theorem needs more: one matrix over the original CM field must generate
every such frame. The source achieves this using a common self-dual lattice
and an explicit companion identity. The argument is elementary once the
actual base lattices are identified.

Write $E$ for the CM field, $\mathcal O=\mathcal O_E$, and
$\bar x$ for its CM involution. For a matrix $M$, put
$M^*=\bar M^{\mathsf T}$. In coordinates let
\[
 h_J(x,y)=\bar x^{\mathsf T}Jy,\qquad
 \Lambda^\#=\{y:h_J(x,y)\in\mathcal O\text{ for every }x\in\Lambda\}.
                                                               \tag{SF.1}
\]
The dual in (SF.1) is an $\mathcal O$-module: addition and multiplication
in the second argument preserve the integrality condition. Thus the
source's $\mathcal O$-span of \texttt{dualSet} is exactly this module.
It is a Hermitian dual, with no trace or different ideal.

\paragraph{SF1: the global matrix and the exact integer coordinates.}
Let $V$ be an $n$-dimensional $E$-space, $c:V\simeq E^n$ a coordinate
isomorphism, and $b_1,\ldots,b_n$ an orthonormal $E$-basis for $h$.
Assume $h(v,w)=h_J(c(v),c(w))$ and
$L_0=\bigoplus_j\mathcal O b_j$. Define
\[
 B=[\,c(b_1)\ \cdots\ c(b_n)\,]\in\mathrm{GL}_n(E).
\]
Then
\[
 \boxed{\ B^*JB=I,\qquad c(L_0)=B\mathcal O^n,\qquad
 x\in c(L_0)\Longleftrightarrow B^{-1}x\in\mathcal O^n.\ }       \tag{SF.2}
\]
Indeed the $(i,j)$ entry of $B^*JB$ is $h(b_i,b_j)$, while
$B^{-1}c(v)$ is the unique coefficient vector of $v$ in the basis $b$.
This proves both assertions. In the source, the first-level equality
$\operatorname{rank}(X_0)=R.k$ only reindexes this basis: apply the same
permutation to rows, columns and coefficient vectors. It neither changes
the Gram matrix nor introduces a denominator. This is the full argument
for \texttt{stdB\_core} and \texttt{exists\_stdB\_of\_tied}.

\paragraph{SF2: the actual spaces already have standard coordinates.}
For the supplied true-column record $I$ with \texttt{I.Ok}, its
\texttt{spaces} field identifies every $X_i$ with the positive space
\[
 E^{N(i)},\qquad h(x,y)=\sum_j\bar x_jy_j,
 \qquad L_{0,i}=\mathcal O^{N(i)}.
                                                               \tag{SF.3}
\]
The space $X'$ is the same construction in dimension $N(m)-1$.
The designated frames are the coordinate bases, not bases obtained by
an existence theorem. Their pairings are $\delta_{ij}$ and their
integer spans are the displayed lattices. Consequently
\texttt{StdTied} holds, $\texttt{gramP}=I$,
$\texttt{gramN}(i)=I$, and membership in
$\texttt{selfDualLatticeCoords}(i)$ is coordinatewise integrality.
In these literal coordinates one can choose $B=I$ in (SF.2).

The \texttt{Pins.iX} clause sends the $j$th basis vector of $X'$ to
the $j$th basis vector of $X_m$. By expanding a vector in its basis,
the inclusion is therefore
\[
 (x_1,\ldots,x_{N(m)-1})\longmapsto
 (x_1,\ldots,x_{N(m)-1},0).                              \tag{SF.4}
\]
There is no change of the underlying class lattices. By definition,
$\texttt{latAt}(u,a,i,c)$ is the original representative lattice pulled
back through the inverse coordinate map. Mapping forward gives exactly
the original representative lattice. The frame set is identified by the
identity map, so its determinant functions are unchanged. These observations
prove all the coordinate assertions of
\texttt{spaces\_in\_coordinates}, under its literal
\texttt{Ok, Pins, Laws, Later} premises.

\paragraph{SF3: the companion identity is a one-coordinate calculation.}
For $0\ne\pi\in\mathcal O$ and $n\ge1$, put
\[
 L_\pi=\mathcal O e_1\oplus\cdots\oplus\mathcal O e_{n-1}
                         \oplus\pi\mathcal O e_n.
\]
Testing (SF.1) against $e_1,\ldots,e_{n-1},\pi e_n$ gives
\[
 L_\pi^\#=
 \mathcal O e_1\oplus\cdots\oplus\mathcal O e_{n-1}
                         \oplus\bar\pi^{-1}\mathcal O e_n.
\]
The same test also proves sufficiency. Hence
\[
 \boxed{\ L_\pi+\bar\pi L_\pi^\#=\mathcal O^n.\ }       \tag{SF.5}
\]
For an explicit decomposition of an integer vector $x$, take
$x'=x-x_ne_n\in L_\pi$ and
$y=(x_n/\bar\pi)e_n\in L_\pi^\#$; then
$x=x'+\bar\pi y$. Conversely each summand in (SF.5) is integral.
The case $n=0$ is the zero module and needs no nonzero-$\pi$ argument.
If $\pi\bar\pi=T\in\mathbb Z$, the same coordinate descriptions give
\[
 L_\pi\subset\mathcal O^n\subset L_\pi^\#,
 \qquad T L_\pi^\#\subset L_\pi.                       \tag{SF.6}
\]
The conjugate on $\pi$ in (SF.5) is required by the convention
that the first argument of $h$ is conjugate linear.

\paragraph{SF4: identify the actual carried lattices.}
The source defines $\texttt{baseL}(n,\pi)$ by
\[
 x_j\in\mathcal O\quad\text{for every }j,
 \qquad x_n\in\pi\mathcal O,
                                                               \tag{SF.7}
\]
using the equivalent zero-indexed condition $j+1=n$ for the last
coordinate. Thus it is precisely $L_\pi$.

Here are the actual joins, rather than an assumption that the companion
formula happens to hold. Supply $D=I.\texttt{rec7}(l,Z)$ and standard
forms $F_u$, Hermitian columns $H$, and level-prime data $A_u$, with
\texttt{StdF}, \texttt{StdLvC}, \texttt{CarH} and \texttt{CarT}
as in the source. At level zero, \texttt{StdF.ends} supplies
\texttt{Reads.lb}, which identifies the forms-side lattice with
\texttt{baseL}(N(0),X.\texttt{piT}(u)). The \texttt{CarH.car0}
identity pulls it back along the exact rank reindexing. Hence the
Hermitian level-zero lattice is $L_{\pi_u}$, with
$\pi_u=X.\texttt{piT}(u)$.

At $1\le i<m$ use \texttt{StdF.mid}, \texttt{Reads.lb} and
\texttt{CarH.carC}. At $i=m$ use \texttt{StdF.ends} and
\texttt{CarH.carM}; the terminal chart set is a singleton.
These give the same description for every chart in its own dimension.
The equalities \texttt{CarT.piU} and \texttt{StdF.piU} identify
the multiplier in every case with the same $\pi_u$.
It is nonzero because $\pi_u\bar\pi_u=T_u$ and $T_u$ is prime.
SF2 identifies the Gram and self-dual lattices. Applying (SF.5) now
proves the complete lattice assertion
\texttt{SymmetricBaseCompatibility}, the chart-indexed version of
\texttt{BaseTied}, including its higher-level clauses. In particular, with
$\mathfrak a_u=(\bar\pi_u)$,
\[
 L_u+\mathfrak a_u L_u^\#=L_0                         \tag{SF.8}
\]
is the actual level-zero \texttt{Companion} identity, for the same
columns used in the comparison. No basis change depending on $u$ has
been introduced.

\paragraph{SF5: a localization lemma with an explicit inverse.}
Let $v$ be a split finite place of the real subfield $F$,
$K=F_v$, $R=\mathcal O_{F_v}$, and let
$j:E\to K$ be the chosen first projection of
$E\otimes_FK\simeq K\times K$ from LI5.
Suppose $L\subset B\mathcal O^n$ and
\[
 N B e_i\in L\quad(1\le i\le n),\qquad N\in R^\times,
                                                               \tag{SF.9}
\]
for an integer $N$. Put $M_v=R\text{-span}\,j(L)$.
Then
\[
 M_v=j(B)R^n,\qquad
 t\in M_v\Longleftrightarrow j(B)^{-1}t\in R^n.         \tag{SF.10}
\]
For the forward inclusion, $B^{-1}x$ is integral for $x\in L$,
and $j(\mathcal O)\subset R$. For the reverse, if
$u=j(B)^{-1}t\in R^n$, then
\[
 t=\sum_i (u_i/N)\,j(NBe_i).                           \tag{SF.11}
\]
Every coefficient belongs to $R$ and every generator comes from $L$.
LI5 identifies $M_v$ with the literal first projection of
$\texttt{locL}_vL$. One can also lift (SF.11) inside that lattice:
the tensor scalars $1\otimes(u_i/N)$ belong to its integral tensor
subring and multiply the images of $NBe_i$. Thus (SF.10) proves the
source's exact existential \texttt{fr\_lat} clause, not just equality
of abstract lattices.

\paragraph{SF6: the same $B$ works at every admissible place.}
For a supplied column level $L=L_u$, integrality of its pairing gives
$L\subset L^\#$, and its \texttt{Lb\_dual} condition gives
$T_u L^\#\subset L$; this extends from \texttt{dualSet} to its span
by linearity. From (SF.8),
\[
 L\subset L_0=B\mathcal O^n\subset L^\#,
 \qquad T_uBe_i\in L.                                  \tag{SF.12}
\]
At a prime admitted by \texttt{GIdx}, $p\nmid T_u$. For every
$v\mid p$, $T_u$ is a unit of $R$: integers $a,b$ with
$aT_u+bp=1$ show that its residue is nonzero. Apply SF5 with $N=T_u$.
Replace only the frame in the already constructed local identification by
\[
 B_v=j_v(B).                                             \tag{SF.13}
\]
Its lattice axiom follows from (SF.10); all other local-identification
fields are unchanged. In particular the completion, splitting map,
prime bijection, embedding $\texttt{iw}$, original orbit representatives
and position map from LI1--LI7 are identical. The alignment identity
with $c^s$ at odd positions and the prescribed-place valuation identity
therefore remain literally the same. This proves
\texttt{LocId.withFr}, \texttt{locIdStd},
\texttt{locIdStd\_fr} and \texttt{locIdStd\_iw}.

\paragraph{SF7: the common-frame conclusion and its boundary.}
Fix one supplied true-column construction and the standard forms and
carrier identifications in SF4. SF2 supplies its standard self-dual
spaces, SF4 supplies its companion identities, and SF1 chooses $B$
once from the common level-zero space. For every word $u$ and every
\texttt{GIdx} place, LI supplies a local identification and SF6 replaces
its frame by (SF.13). This gives simultaneously
\[
 B^*J_0B=I,\qquad
 \pi_{1,v}(\operatorname{locL}_v L_u)=j_v(B)R_v^n
 \quad\text{for all such }u,v.                          \tag{SF.14}
\]
For the actual positive-space coordinates $B=I$ is available. The source
uses the more flexible existence formulation, and both choices satisfy
the same exact matrix and integrality conditions.

This discharges the common-frame \emph{algebra}; it does not construct
the entire supplied standard-column package. Its retained source supplier
is \texttt{colsTiedSym\_std}: besides the lattice identities proved here,
that theorem constructs the forms, levels, characters, residues, class
maps and their simultaneous compatibilities. Its hypotheses
\texttt{Ok, Pins, Laws, Later}, and any earlier host/point existence
needed to supply them, keep their stated status. The class-rank equality,
global comparison-lattice span, spectral/metric estimates and arithmetic
windows require their respective arguments; (SF.14) alone proves none
of those conclusions.


\section{The actual class-rank identity}
\label{sec:class-rank}

\status{CR1--CR6 prove the literal \texttt{locPack\_d\_eq} and
\texttt{pack\_d\_eq} conclusions from supplied \texttt{D7},
\texttt{LevelData}, \texttt{Forms0}, and the local package of LI.
The class correspondence and the trivial-stabilizer argument are proved
below. The construction of those global records remains a separate task.
An equality of ranks does not establish an integral or metric comparison.}

\paragraph{The quantity to be compared.}
Let $k$ be the coefficient rank and put $m_{\mathrm{loc}}=k-1$;
this local polynomial parameter is distinct from the global chain
length $D.R.m$. Let $\mathcal O\subset K$ be the coefficient DVR and
its fraction field. At a prime $p$, the source has places $v$ and
Frobenius positions $s=0,\ldots,n_v-1$, with a decreasing integer
weight $\lambda^0_{v,s}$ of length $k$. A tuple $\mathbf L=(L_v)_v$
of full $W_v$-lattices gives a polynomial lattice
\[
 H(\mathbf L)=\bigotimes_{v,s}
       D_{\lambda_{v,s}}(\Lambda_{v,s}),\qquad
 \Lambda_{v,s}=\mathcal O\,\iota_v(\sigma_v^s L_v),
 \quad
 \lambda_{v,s}=
 \begin{cases}
 \lambda^0_{v,s},&s\text{ even},\\
 -\operatorname{rev}\lambda^0_{v,s},&s\text{ odd}.
 \end{cases}                                                    \tag{CR.1}
\]
The tensor is realized as products in disjoint polynomial-variable
blocks; it is the literal \texttt{HtupW}.
For each away component $i$, a group $\Gamma_i$ acts on lattice tuples
and polynomial fibers, with an $\mathcal O^{\times}$-valued character
$\chi_i$. The module $M$ consists of functions satisfying
\[
 f_i(\gamma\mathbf L)=\chi_i(\gamma)\rho_i(\gamma)f_i(\mathbf L),
 \qquad f_i(\mathbf L)\in H(\mathbf L).
                                                               \tag{CR.2}
\]
As in the source, these functions are extended by zero to tuples which
are not full lattices. Its rank $d=\operatorname{rank}_{\mathcal O}M$
is the number appearing in the local conclusions.

\paragraph{CR1. The literal Weyl fibers have the required rank.}
Here $D_\lambda$ uses the integral \emph{Weyl} lattice $\Delta$,
not the dual-Weyl coefficient lattice $\nabla$ of CB2.
For $\nu_j=\lambda_j-\lambda_k$, let $\Delta_\nu(\mathbb Z)$ be the
integer span of the divided-power words on its highest polynomial.
All these words have bounded multidegree. Thus this is a finite free
subgroup of a finite free group of integer polynomials. Choose any
integer basis $b_1,\ldots,b_q$. The characteristic-zero argument of IW3
identifies its rational span with that of $\nabla_\nu$; CB2--CB3 give
\[
 q=D_k(\lambda):=
       \prod_{a<b}\frac{\lambda_a-\lambda_b+b-a}{b-a}.
                                                               \tag{CR.3}
\]
Divided-power polynomial operators have integer coefficients and commute
with coefficient extension. Every generator over $\mathcal O$ is
therefore an $\mathcal O$-linear combination of the same $b_j$.
A nonzero integer minor of their coefficient matrix remains nonzero in
$K$, so these images are still independent. This proves the exact
base-change assertion for $\Delta$, without asserting
$\Delta=\nabla$ integrally.

For $\Lambda=g\mathcal O^k$, the source's \texttt{Dl0} is
\[
 D_\lambda(\Lambda)
   =\det(g)^{-\lambda_k}R(g^{-1})\Delta_\nu(\mathcal O),
       \qquad R(g)P(X)=P(gX).                            \tag{CR.4}
\]
This expression is independent of the chosen basis $g$:
changing it to $gh$, $h\in\mathrm{GL}_k(\mathcal O)$, preserves
the Weyl lattice and introduces only an integral determinant unit.
The displayed invertible $K$-linear transformation sends the actual
integer basis to an $\mathcal O$-basis of (CR.4).
Products of these bases in disjoint blocks prove
\[
 \operatorname{rank}_{\mathcal O}H(\mathbf L)
    =q_0:=\prod_{v,s}D_k(\lambda^0_{v,s}).                \tag{CR.5}
\]
Indeed, constant weight shifts and negative reversal preserve the
product in (CR.3). Neither the basis matrices nor their determinant
twists need be integral units.

The required characteristic zero follows from the actual Hida setting.
If $W_v$ had prime characteristic $\ell$, then $p=0$ when $\ell=p$,
contradicting that $p$ is a uniformizer; if $\ell\ne p$, then $p$ is a
unit in the prime subfield, again a contradiction. Thus $W_v$, its
fraction field, and its embedded extension $K$ have characteristic zero.

\paragraph{CR2. Evaluate on one representative of each orbit.}
Suppose $\Gamma_i$ acts freely on the full lattice tuples and has finitely
many orbits. Choose one representative $\mathbf L_x$ of each orbit.
The determinant-twisted row action in (CR.4) takes
$H(\mathbf L)$ onto $H(\gamma\mathbf L)$: apply the formula to a
basis matrix and then apply $\gamma^{-1}$ for the reverse inclusion.
The scalar $\chi_i(\gamma)$ is a unit. Consequently evaluation is an
\emph{integral} linear isomorphism
\[
 M_i\ \longrightarrow\ \prod_{x\in\Gamma_i\backslash\mathrm{Lat}}
                                  H(\mathbf L_x).       \tag{CR.6}
\]
For its inverse, prescribe $a_x\in H(\mathbf L_x)$ and set
$f_i(\gamma\mathbf L_x)=\chi_i(\gamma)\rho_i(\gamma)a_x$.
The orbit determines $x$ and freeness determines $\gamma$, so this
is well-defined. The action and character laws give (CR.2), and the
fiber transport gives its integrality. Conversely, (CR.2) determines
every value from the representatives. This proof also explains why
nontrivial stabilizers would require invariant subspaces rather than
arbitrary fibers.

\paragraph{CR3. The actual class labels are the local lattice orbits.}
Let $\mathcal G=\mathrm U(J)(\mathbb A_{F,f})$, let
$\mathcal H=\mathrm U(J)(F)$ denote its rational subgroup, and let
$U$ be the supplied level. Remove its conditions above $p$ to obtain
$U'$. The local identifications of LI supply a projection and section
\[
 \pi:\mathcal G\longrightarrow
       G_p:=\prod_{v\mid p}\mathrm{GL}_k(Q_v),\qquad
 j:G_p\longrightarrow\mathcal G,
 \qquad \pi j=\operatorname{id}_{G_p},\quad j(G_p)\subset U'.
\]
Here $j$ inserts a split unitary element above $p$ and the identity
elsewhere; its second split component is the adjoint inverse of the
first. In the frames of LI5,
\[
 U=\{a\in U':\pi(a)\in K_p\},
       \qquad K_p=\prod_{v\mid p}\mathrm{GL}_k(W_v).
                                                               \tag{CR.7}
\]
To check this equality for the actual level, $\texttt{GIdx}$ excludes
$p$ from the two auxiliary primes and the integer $T$ controlling
$\Lambda\subset\Lambda^{\#}\subset T^{-1}\Lambda$.
Localization therefore makes $\Lambda$ self-dual. In the integral
splitting $E\otimes_F F_v=F_v\times F_v$, its two projected lattices
are mutually dual. A unitary element preserves both exactly when its
first component preserves the first projected lattice. LI5 identifies
that condition with integrality of the framed matrix and its inverse,
which is precisely membership in $\mathrm{GL}_k(W_v)$.

Choose representatives $y_i$ for $\mathcal H\backslash\mathcal G/U'$
with $\pi(y_i)=1$. They can always be so normalized: replace any
representative $y$ by $yj(\pi(y))^{-1}$. Define
\[
 \Gamma_i=\{\gamma\in\mathcal H:y_i^{-1}\gamma y_i\in U'\}.
\]
Every full lattice tuple has the form $g\mathbf W$, where
$\mathbf W=(W_v^k)_v$, and two matrices give the same tuple exactly
when they differ on the right by $K_p$. Define its class by
\[
 (i,g\mathbf W)\longmapsto[\,y_i j(g)\,]
                        \in\mathcal H\backslash\mathcal G/U.
                                                               \tag{CR.8}
\]
Changing $g$ by $gk$, $k\in K_p$, does not change this class, by
(CR.7). More precisely, multiplication in the double quotient gives
\[
 [y_i j(g)]=[y_{i'}j(g')]
 \quad\Longleftrightarrow\quad
 i=i'\ \text{ and }\
 \exists\gamma\in\Gamma_i:\quad
       g'\mathbf W=\pi(\gamma)g\mathbf W.               \tag{CR.9}
\]
For the reverse implication the needed right-level element is
$j(g)^{-1}y_i^{-1}\gamma^{-1}y_i j(g')$; its membership follows
from (CR.7) and
$g'^{-1}\pi(\gamma)g\in K_p$.
For the forward implication, first pass to the $U'$-quotient to get
$i=i'$, and then project the equality to $G_p$.

There are no missing classes. If $x=\gamma y_i u'$ with $u'\in U'$,
then $j(\pi(u'))^{-1}u'\in U$ by (CR.7). Thus $[x]$ is the class
of $y_i j(\pi(u'))$. Composing (CR.8) with the supplied bijection
from the level's double quotient to its finite class set $C_0$ gives
the literal map \texttt{clsR}, its equality criterion, and its
surjectivity. In particular,
\[
 \coprod_i\Gamma_i\backslash\mathrm{Lat}\ \simeq\ C_0.
                                                               \tag{CR.10}
\]

\paragraph{CR4. Why the actual stabilizers are trivial.}
First, every rational stabilizer of a level class is finite.
For an adelic representative $x$, intersect the transformed local
lattices $x_v\Lambda_v$ with $E^k$. Clearing the finitely many adelic
denominators sandwiches this intersection between
$a\mathcal O_E^k$ and $a^{-1}\mathcal O_E^k$ for a nonzero integer
$a$, so it is a full integer lattice. The stabilizer preserves it
and the Hermitian form, which is definite at every archimedean place.
In the Minkowski embedding, each of its basis vectors can therefore
have only finitely many possible images: a discrete lattice meets
a bounded set in finitely many points. Hence the stabilizer is finite.

At each of the source's two distinct auxiliary rational primes
$g_0,g_1$, a conjugate of a stabilizer element is integral and reduces
to an upper-unitriangular matrix. This is the literal
$\texttt{iwU}$ condition, weaker than being the identity modulo the
prime. We need the following elementary fact. If $A\in\mathrm{GL}_k(R)$
over a complete DVR of residue characteristic $g$ is upper
unitriangular modulo its maximal ideal $\mathfrak m$ and
$A^N=1$ with $(N,g)=1$, then $A=1$. Indeed its reduction has
$g$-power order, so it is already the identity. If
$A=1+X$ with $X\in M_k(\mathfrak m^r)$, then
\[
 A^N\equiv1+NX\pmod{\mathfrak m^{r+1}}
       \qquad(r\ge1).
\]
As $N$ is a unit, induction puts $X$ in every power of
$\mathfrak m$, whence $X=0$.

An element of finite order $N=g^a b$, $(b,g)=1$, in this local
subgroup consequently satisfies $A^{g^a}=1$. Thus its order is a
power of $g$. The maps from the rational field into either split
completion are injective; the global stabilizer element has both
$g_0$-power and $g_1$-power order, and is the identity.
Finally, if $\gamma\in\Gamma_i$ fixes $g\mathbf W$, then (CR.7)
puts $(y_i j(g))^{-1}\gamma(y_i j(g))$ in $U$.
The rational-stabilizer result therefore proves exactly the freeness
required in CR2, as well as triviality of the level-zero class
stabilizers.

\paragraph{CR5--CR6. Match the two dimensions.}
Combining (CR.5), (CR.6) and (CR.10) gives the general rank formula
\[
 \boxed{\ d=|C_0|\prod_{v,s}D_k(\lambda^0_{v,s}).\ }      \tag{CR.11}
\]
The supplied \texttt{Forms0} identifies the complex level-zero
automorphic forms with their coefficient values on $C_0$, subject
only to the class-stabilizer invariance conditions. CR4 removes
those conditions. Each coefficient fiber is the tensor of the
characteristic-zero Weyl representations of weights $\lambda_\tau$.
Consequently
\[
 \dim_{\mathbb C}\mathrm{Mu}_0
       =|C_0|\prod_{\tau\in\mathrm{Tau}}D_k(\lambda_\tau).
                                                               \tag{CR.12}
\]
The actual position map is
$\texttt{gatRec\_pos}=\texttt{eTau}^{-1}\circ\texttt{posG}$,
that is, $(v,s)\mapsto\texttt{eTau}^{-1}(\texttt{posG}(v,s))$;
it is a bijection. The definition \texttt{lamRec} identifies
$\lambda^0_{v,s}$ with precisely the corresponding $\lambda_\tau$.
The rank reindexing uses $m_{\mathrm{loc}}+1=k=N(0)$, while negative reversal at
odd positions was already accounted for in (CR.5). Thus
\[
 \boxed{\ (X.\texttt{gat}\ r\ e).d
          =\dim_{\mathbb C}(D.\texttt{mem}\ u).A.\mathrm{Mu}_0\ }
             \qquad(r,e\in\mathbb N).                    \tag{CR.13}
\]
This is the source's \texttt{locPack\_d\_eq}; substituting the
constructed \texttt{pack} gives \texttt{pack\_d\_eq}.

The exported wrapper takes \texttt{Joint0} as an argument but its
proof does not use it. This is mathematically exact: the two supplied
level descriptions already use the same finite set $C_0$ and the
same weights, which suffice for (CR.11)--(CR.12). Their actual
lattices, local subgroups and representatives need not be identified
to compare these dimensions. \texttt{Joint0} retains its separate
role in the character and comparison constructions. No integral
identification, Gram identity, metric transport, discarded-target
estimate, or global arithmetic window follows from (CR.13) alone.


\section{The exact global comparison lattice}
\label{sec:comparison-span}

\status{CS1--CS9 prove the integral descent, local comparison and
global span arguments for the supplied compatible forms. In particular,
the source's \texttt{ValIntRec} and comparison-lattice equality are
deduced below; they are not new assumptions. The conditional comparison proof below retains its stated
\texttt{ValEx} premise; VE1--VE9 in Section~\ref{sec:value-existence}
now construct those polynomial values from the exact compatible global
form record. FC/CV and the adelic-geometry sections reconstruct the
specified portions of that record. Their remaining arithmetic and
geometric suppliers are recorded separately.}

\paragraph{Setting and notation.}
Keep the CM field $E/F$, the same global frame $B$ from SF, and the
same embeddings and positions from LI. Write $L$ for the coefficient
number field, $e:L\hookrightarrow K$ for the prescribed-place embedding,
$O\subset K$ for the local DVR and
$R=L\cap e^{-1}(O)=\mathcal O_{L,w}$.
Let $V$ be the global space of level-zero forms and
$\mathcal L\subset V$ its specified full integer lattice. The comparison
map $C:V\to\mathcal F_K$ is semilinear for $e$; its values are polynomial
functions of the local lattice tuples. The target $M\subset\mathcal F_K$
is the literal module \texttt{G.M}: its values lie in the integral
\emph{Weyl} tensor lattices and satisfy the prescribed level-character
equivariance. The Weyl lattice $\Delta$ is generated by divided-power
words on the highest polynomial. It is not replaced by the coefficient
lattice $\nabla$ of IW.

The intended conclusion is an equality of integral modules:
\[
 \boxed{\quad O\text{-span}\,C(\mathcal L)=M.\quad}       \tag{CS.1}
\]
Equality after taking the $K$-span is only one step in its proof.

\paragraph{CS1: exact descent through a valuation ring.}
Let $L\subset K$ be any field extension, $O$ a valuation ring of $K$,
and $R=O\cap L$. For polynomials $q_i,Q$ over $L$,
\[
 Q\in O\text{-span}\{q_i\}\quad\Longleftrightarrow\quad
 Q\in R\text{-span}\{q_i\}.                            \tag{CS.2}
\]
Only finitely many $q_i$ occur in a given expression. If they are
linearly independent over $L$, an invertible coefficient minor shows
that their coefficients in an expression for $Q$ lie in $L$; hence
they lie in $R$. Otherwise choose a nonzero relation
$\sum_i a_iq_i=0$. The valuation-ring property permits a pivot $j$
with $a_j\ne0$ and $a_i/a_j\in R$ for every $i$.
In $Q=\sum_i c_iq_i$, replace
\[
 c_i\longmapsto c_i-c_j a_i/a_j.
\]
All new coefficients remain in $O$ and the $j$th becomes zero.
Induction proves (CS.2); the reverse inclusion is immediate.
This is the source's \texttt{descVal}. It does not require $K/L$ to
be finite or $O$ to be a finite $R$-module.

\paragraph{CS2: the integral Weyl comparison, including odd positions.}
For a decreasing integral weight $\lambda=(\lambda_1,\ldots,\lambda_n)$,
put $\nu_i=\lambda_i-\lambda_n$ and
$\lambda^*=(-\lambda_n,\ldots,-\lambda_1)$. Let $R_n$ be the
column-reversal matrix and $\epsilon=\det R_n\in\{1,-1\}$.
On invertible matrices define
\[
 (\mathcal D_\lambda f)(X)=
 \epsilon^{-\lambda_1}\det(X)^{\lambda_1-\lambda_n}
                      f(X^{-\mathsf T}R_n).             \tag{CS.3}
\]
It restricts to an integral isomorphism
$\Delta_{\nu(\lambda^*)}(\mathbb Z)\simeq
\Delta_{\nu(\lambda)}(\mathbb Z)$.
Here is a direct proof. Jacobi's complementary-minor identity sends
each initial minor of $X^{-\mathsf T}R_n$ to a signed complementary
minor of $X$ divided by $\det X$. In the highest polynomial the total
denominator exponent is $\lambda_1-\lambda_n$, so (CS.3) sends it
to a signed row-reversal of the highest polynomial for $\lambda$.
For $a\ne b$, substitution by $I+tE_{ab}$ intertwines with substitution
by $I-tE_{ba}$. Comparing powers of the indeterminate $t$ intertwines
their divided powers integrally, with signs. Row permutations preserve
the lattice: signed transpositions are products of integral elementary
matrices, and diagonal sign changes act integrally on every weight.
Diagonal divided powers act by integer binomial coefficients on weight
vectors. Thus all generating words have polynomial integral images.
Applying (CS.3) twice, with $\lambda$ and $\lambda^*$, is the identity;
this proves both inclusions, including all signs.

Equivalently, if $\widetilde f=\mathcal D_\lambda f$, then
\[
 \det(Z)^{\lambda_n}\widetilde f(Z)=
 \det(Z^{-\mathsf T}R_n)^{-\lambda_1}
                                    f(Z^{-\mathsf T}R_n).       \tag{CS.4}
\]
This is the literal \texttt{OddId}/\texttt{OddFact} identity. Negative
determinant twists are allowed. They are evaluated only on invertible
matrices; no determinant denominator is inserted into $\Delta$.

At a lattice $A O^n$, the integral factor used here is exactly
\[
 D_\lambda(AO^n)=
 \det(A)^{-\lambda_n}\,R(A^{-1})\Delta_\nu(O),\qquad
 (R(g)f)(X)=f(gX).                                      \tag{CS.5}
\]
The divided-power generators have integer coefficients and commute
with coefficient extension. Hence (CS.5) and every tensor product of
these factors are the $O$-spans of their corresponding integer
generators after the displayed substitutions. In even positions use
the basis matrix $N_\tau$; in odd positions use $N_\tau^{-\mathsf T}$
and (CS.4), with the source's exact position reindexing: $\lambda_\tau=\lambda^0_{v,s}$,
while the local factor has weight $\lambda^0_{v,s}$ for even $s$ and
$(\lambda^0_{v,s})^*$ for odd $s$. This is \texttt{WtTie}. CS1 descends their coefficients to $R$.
This proves \texttt{WeylLatticeDescent}: membership of a coefficient-field
polynomial in the local Weyl tensor lattice is equivalent to an
$R$-linear combination of the global integer Weyl tensors in the
corresponding frame. Equality of the functions on all invertible complex
frames implies equality of the polynomials: substitute $X+tI$ in each
block and use that only finitely many $t$ give a singular block.

\paragraph{CS3: the odd-position lattice has no hidden index.}
Fix a class lattice $L_c$, its adelic representative $y_c$, and a
matrix $B'$ which is a basis of $L_c^\#$ at all primes over $p$.
Such a matrix exists. The ring obtained from $\mathcal O_E$ by
inverting integers prime to $p$ is a semilocal Dedekind domain;
each of its ideals is principal, by prescribing the finitely many
valuations using CRT. Its finite torsion-free modules are free.
Applied to the full Hermitian dual lattice, this gives exactly
\texttt{IsLocBasis}, whose denominators are integers prime to $p$.

At $v\mid p$, $p\nmid T$ and the source inclusions
$L_b\subset L_b^\#$, $T L_b^\#\subset L_b$ show that $L_b$ is
locally self-dual. The unitary $y_c$ carries it to $L_c$, so the same
holds for $L_c$ and for the lattice spanned locally by $B'$.
Consequently its Gram matrix
\[
 G'=B'^*JB',\qquad X^*=\bar X^{\mathsf T},
\]
and $G'^{-1}$ have integral entries at these places. The first assertion
tests pairs of basis vectors; the second tests the dual basis and uses
self-duality. The \texttt{Joint0} identity identifies the level lattice
with the one framed by $B$, so this argument concerns the same local
identification as SF6.

For a comparison point $\gamma y_c\kappa$ with $\gamma^*J\gamma=J$
and $\kappa$ in the level group, put
\[
 N=B^{-1}\gamma B'.
\]
The point identity and $\kappa$'s preservation of the base lattice give
$L_v=\operatorname{colspan}_{W_v}i_v(N)$.
Indeed the change-of-basis matrix between the two descriptions and its
inverse preserve the integral coordinate module; both are integral,
so their column spans agree. Now $B^*JB=I$ gives
\[
 N^*N=G',\qquad (N^{-1})^*=NG'^{-1}.                    \tag{CS.6}
\]
For an even position, LI's compatibility identifies the two matrices
entry by entry. For an odd position it identifies the coefficient
matrix $e(t_\tau(N)^{-\mathsf T})$ with the Frobenius image of
$(N^{-1})^*$. By (CS.6) that differs from the Frobenius image of $N$
by right multiplication by the integral invertible matrix $G'^{-1}$.
Thus its $O$-column span is identical. This proves the literal
\texttt{RestLatW}, \texttt{RestGramW} and \texttt{RestLatO} clauses.
The $c^s$ in LI is essential to (CS.6); deleting it would give the
wrong odd-position matrix.

\paragraph{CS4: comparison preserves exactly the specified integrality.}
Supply the source's polynomial-value condition \texttt{ValEx}.
It says that the global form at a comparison point is represented by
one coefficient-field polynomial $Q$, after the specified frame
substitution and level character. CS2 and CS3 identify the membership
condition for $e(Q)$ in the local Weyl tensor lattice with the global
integer-Weyl expression in the basis $B'$.

The character factor does not change this condition. Under
\texttt{LevelZeroCharacterRational}, its factors have representatives
$r\in L$ with $r^a=1$ for a positive integer $a$: they come from
additive characters of finite residue fields. Hence
$a\,v_w(r)=0$, so both $r$ and $r^{-1}$ belong to $R$.
Multiplication by the character, or its inverse, preserves the
$R$-span of integer Weyl tensors. Translating the frame by
$\gamma$ and then by $\gamma^{-1}$ proves both directions of
\texttt{valIntAt\_of}; using one local basis in the reverse direction
and every local basis in the forward direction proves its exact
quantifiers. In the actual record this is \texttt{ValIntRec}.

\paragraph{CS5: the localization criterion is an equality.}
The global lattice $\mathcal L$ is defined by the specified Weyl
integrality conditions at every finite coefficient place. Suppose
$x\in V$ satisfies them at $w$. Choose $0\ne d\in\mathcal O_L$
with $dx\in\mathcal L$, using fullness and finite generation.
By CRT choose $s\in\mathcal O_L$ with
\[
 s\equiv1\pmod w,\qquad
 v_{w'}(s)\ge v_{w'}(d)\quad(w'\ne w).                 \tag{CS.7}
\]
Only primes dividing $d$ impose nontrivial conditions. At $w$, $sx$
is integral because $s$ is; elsewhere $sx=(s/d)(dx)$ is integral.
Thus $sx\in\mathcal L$, while $s^{-1}\in R$. This proves
$x\in R\mathcal L$. The converse follows by linearity of every
local condition. CS4 and the comparison's point/class surjectivity
therefore give
\[
 \boxed{\ C(x)\in M\quad\Longleftrightarrow\quad
                         x\in R\mathcal L.\ }          \tag{CS.8}
\]
The level-equivariance part follows from the same point identity and
the matching inverse away character from LI8. Polynomial uniqueness
makes the comparison independent of a chosen point representative.

\paragraph{CS6: one global matrix realizes every local lattice tuple.}
For the finitely many distinct places $v\mid p$, take arbitrary
$g_v\in\mathrm{GL}_n(F_v)$. Choose $a$ so large that
$p^a g_v^{-1}$ is integral for every $v$. Simultaneous approximation
gives $A\in\operatorname{Mat}_n(F)$ with
\[
 A_v-g_v\in p^{a+1}\operatorname{Mat}_n(\mathcal O_{F_v}).
\]
For completeness, approximate each entry separately in its completion,
clear their finitely many global denominators with one integer, and
apply CRT to the distinct prime-power ideals with correspondingly
increased precision. Dividing back gives the displayed approximation.
Now $g_v^{-1}A_v=I+pD_v$, which is invertible over
$\mathcal O_{F_v}$ by reduction modulo $p$. Hence
\[
 A_v\mathcal O_{F_v}^{n}=g_v\mathcal O_{F_v}^{n}.    \tag{CS.9}
\]
Its determinant is nonzero locally, so $A\in\mathrm{GL}_n(F)$.
The empty tuple permits $A=I$. Via LI's completion maps this proves
\texttt{WeakApprox} and \texttt{RatTup}: the approximation produces
\emph{exact} lattice equality, not an approximation to the desired
integral lattice.

\paragraph{CS7: rational integral functions span the local module.}
Call a local polynomial or function rational if every coefficient
belongs to $e(L)$. By CS6 the tuple lattices can be represented by one
global matrix. LI's positions show that all its Frobenius-twisted
matrix entries belong to $e(L)$. The integer divided-power generators
in CS2 therefore show that every local Weyl tensor lattice is the
$O$-span of rational elements.

This also holds for its stabilizer-invariant part. Here is a proof
that does not divide by a stabilizer order. Choose finitely many
rational generators and let $D_R$ be their $R$-span in the finite
$L$-space $U_L$. It is finite free. Let $U_0\subset U_L$ be the
subspace cut out by the rational level-character equations. These
equations are rational because the matrices are global and the
characters have the coefficient-field representatives from LI8.
The intersection $D_R\cap U_0$ is saturated in $D_R$, so it is a
direct summand over the DVR $R$. Clearing a common $R$-denominator
shows that its $L$-span is exactly $U_0\cap LD_R$. Choose a basis
adapted to this summand. The same basis over $O$ proves
\[
 (O D_R)\cap(K U_0)=O(D_R\cap U_0).                    \tag{CS.10}
\]
The common kernel of even infinitely many rational equations commutes
with this field extension: their rows span a finite-dimensional space,
so a finite subset suffices. Thus (CS.10) is exactly the invariant
claim, without an averaging denominator.

There are finitely many component/lattice orbits. Evaluate a function
at one representative per orbit and extend each value equivariantly,
setting it zero on the other orbits. A stabilizer-invariant value makes
this well defined. The sum of these finitely many extensions recovers
the original function. Rational values extend to rational functions
because the character-twisted action is rational. It follows that
\[
 M=O\text{-span}\{F\in M:F\text{ is rational}\}.        \tag{CS.11}
\]
This is the source's \texttt{MRatSpan}.

\paragraph{CS8: rank equality gives rational surjectivity.}
The comparison is injective: every global class has a comparison point,
and if its polynomial value vanishes, the value identity forces the
global form to vanish on every invertible complex frame. The global
realization of forms is itself injective. Its values have rational
coefficients by construction.

An $L$-independent family in $V$ stays $K$-independent after comparison.
Indeed, expand the coefficients of a finite $K$-relation in an
$L$-basis of their span. Taking each basis coordinate in every
polynomial coefficient produces an $L$-relation. Injectivity kills
each relation separately. The supplied rational-forms record identifies
the complex images of a basis of $V$ with a basis of \texttt{Fi(0)}
(its \texttt{emb\_basis} clause), and \texttt{Fi\_zero} identifies
that space with $\mathcal M_{u,0}$. Thus $\dim_LV=\dim_{\mathbb C}\mathcal M_{u,0}$.
CR6 now gives $\dim_LV=\operatorname{rank}_O M$; this is the complete
\texttt{card\_eq\_d} bridge. CS5 shows the basis images
belong to $KM$, because $\mathcal L$ spans $V$. Finiteness of $M$
then shows these images form a $K$-basis of $KM$.

If $F\in M$ is rational, write $F=\sum_i a_iC(b_i)$ with $a_i\in K$.
Choose an $L$-linear retraction $\pi:K\to L$ with $\pi(1)=1$,
and apply it coefficient by coefficient. Rationality gives
\[
 F=C\Bigl(\sum_i\pi(a_i)b_i\Bigr).                     \tag{CS.12}
\]
This proves \texttt{surjRat\_pf}. It is stronger than rational-span
equality in the precise way required next, but it does not yet assert
integral surjectivity.

\paragraph{CS9: the integral span and its exact hypotheses.}
For $x\in\mathcal L$, CS5 gives $C(x)\in M$, proving the easy
inclusion in (CS.1). Conversely, by CS7 it suffices to consider a
rational $F\in M$. CS8 writes $F=C(x)$, and CS5 gives
$x\in R\mathcal L$. Since $e(R)\subset O$, semilinearity implies
$C(x)\in O\text{-span}\,C(\mathcal L)$. This proves (CS.1).

For the literal source data, the theorem just proved is
\texttt{comparison\_lattice\_span\_eq} from
\texttt{CmpCore}, \texttt{GatIs} and \texttt{ValIntRec}.
CS1--CS4 supply the last premise from the same frames and positions,
supplied \texttt{Forms0} and \texttt{Joint0}, coefficient-field character
rationality, and \texttt{CmpCore.ex}.
\texttt{GatIs} identifies the target with the actual local group
model, supplying its finite class set, fraction field and finite
Weyl lattice. CR supplies the exact rank equality. VE now proves \texttt{ValEx} from the same global record and
supplies every field of the actual comparison hypothesis. In particular,
the existence of a representing polynomial is no longer an additional
assumption. The remaining input is the compatible global record with its
actual arithmetic and geometric construction; FC, CV and the adelic
geometry below specify which parts have been supplied. For supplied
\texttt{Joint0}, \texttt{SameUloc} is its literal local-subgroup equality
and needs no further theorem.

Nothing here identifies Hermitian norms or covolumes. The exact
integral module equality is a necessary input to the later metric
transport and ramified restrictions; those metric and arithmetic
window statements retain their separate status.


\section{Patching the selected local lattices into one global lattice}
\label{sec:lattice-patching}

\status{LP1--LP5 prove the finite-place patching and its exact application
to the selected source lattice. The global geometric datum, the local
comparison and its commensurability bounds are supplied. This does not
construct the geometric point or the stronger numerical record
\texttt{NumAt}.}

Write $R=\mathcal O_K$ for the integer ring of a number field, and let
$V$ be a $K$-vector space. For an $R$-submodule $M\subset V$ and a
nonzero prime ideal $\mathfrak p$, identify its localization with
\[
 M_{\mathfrak p}=\{x\in V:\text{some }s\in R\setminus\mathfrak p
                                      \text{ satisfies }sx\in M\}.
\]
All intersections below are inside $V$. A module $A_{\mathfrak p}$ is
called local at $\mathfrak p$ when $sx\in A_{\mathfrak p}$ and
$s\notin\mathfrak p$ imply $x\in A_{\mathfrak p}$. This is closure
under the units inverted in $R_{\mathfrak p}$; it does not mean
saturation with respect to every nonzero element of $R$.

\paragraph{LP1: clear every denominator except at one chosen prime.}
For $0\ne c\in R$ and a fixed $\mathfrak p$, there is
$t\notin\mathfrak p$ such that, for every $\mathfrak q\ne\mathfrak p$,
some $s\notin\mathfrak q$ satisfies $c\mid st$.
Indeed factor
\[
                 (c)=\mathfrak p^a\mathfrak a,
                 \qquad\mathfrak a+\mathfrak p=R.
\]
Choose $t\in\mathfrak a\setminus\mathfrak p$. If
$\mathfrak q\ne\mathfrak p$, then
$\mathfrak p^a\not\subset\mathfrak q$: for $a>0$ the contrary
would imply $\mathfrak p\subset\mathfrak q$ and hence equality;
for $a=0$ it would say $R\subset\mathfrak q$.
Choose $s\in\mathfrak p^a\setminus\mathfrak q$. Then
$st\in(c)$, as required. No principal generator for
$\mathfrak p$ is chosen.

\paragraph{LP2: localization commutes with this particular intersection.}
Suppose $B_{\mathfrak q}\subset V$ is local at $\mathfrak q$ for
every nonzero prime, and that every $x\in V$ has a common nonzero
denominator:
\[
 \forall x\in V\quad\exists,0\ne c\in R\quad
                   cx\in B_{\mathfrak q}\ \text{for every }\mathfrak q.
                                                               \tag{LP.1}
\]
Then
\[
 N:=\bigcap_{\mathfrak q}B_{\mathfrak q}
       \quad\Longrightarrow\quad N_{\mathfrak p}=B_{\mathfrak p}.
                                                               \tag{LP.2}
\]
One inclusion follows from the local property. For the other, take
$x\in B_{\mathfrak p}$ and choose $c$ in (LP.1). LP1 gives
$t\notin\mathfrak p$. At $\mathfrak p$ we have
$tx\in B_{\mathfrak p}$ directly. At every other
$\mathfrak q$, write $st=cd$ with $s\notin\mathfrak q$.
Then $s(tx)=d(cx)\in B_{\mathfrak q}$, and the local property
gives $tx\in B_{\mathfrak q}$. Thus $tx\in N$, proving (LP.2).
The denominator hypothesis is the reason this infinite intersection
may be localized; no unrestricted commutation rule is being used.

\paragraph{LP3: finite modification of a full reference module.}
Let $S$ be a finite set of primes and suppose that $M$ is full in the
explicit sense that every $x\in V$ has a nonzero integer multiple in
$M$. For $\mathfrak p\in S$, prescribe a local module
$A_{\mathfrak p}$ and a positive integer $n_{\mathfrak p}$ with
\[
 n_{\mathfrak p}M\subset A_{\mathfrak p},\qquad
 n_{\mathfrak p}A_{\mathfrak p}\subset M_{\mathfrak p}.
                                                               \tag{LP.3}
\]
Set $B_{\mathfrak p}=A_{\mathfrak p}$ on $S$ and
$B_{\mathfrak p}=M_{\mathfrak p}$ off $S$. If $mx\in M$ and
$d=\prod_{\mathfrak p\in S}n_{\mathfrak p}$, then $dmx$ belongs
to every $B_{\mathfrak p}$. Hence LP2 constructs the exact patch
\[
 \boxed{\quad
 N_{\mathfrak p}=A_{\mathfrak p}\ (\mathfrak p\in S),\qquad
 N_{\mathfrak p}=M_{\mathfrak p}\ (\mathfrak p\notin S).
 \quad}                                                       \tag{LP.4}
\]
There is also the useful global bound
\[
                        dM\subset N\subset d^{-1}M.            \tag{LP.5}
\]
For the right inclusion, $dx\in M_{\mathfrak p}$ at every prime
when $x\in N$. To justify the resulting global membership, observe
that $M=\bigcap_{\mathfrak p}M_{\mathfrak p}$ for any submodule:
if $x\notin M$, the proper ideal $\{r:rx\in M\}$ is contained
in a maximal ideal, which contradicts $x$ belonging to that
localization. In the present full situation this ideal contains a
nonzero integer, so the maximal ideal is nonzero. The same argument
proves uniqueness of $N$ from (LP.4).
If $M$ is a finite full lattice, (LP.5) and Noetherianity make $N$
a finite full lattice too. The intersection proof of existence
itself does not require finite generation.

\paragraph{LP4: the literal served-prime and type-prime patch.}
Fix a supplied word $u$ and put
$M=M_{\rm ram}(u)\subset V_u$, the retained reference lattice of
PL4. This is the lattice called \texttt{DV.MramL}; it is not the
whole unretained source lattice. Let $T_u$ be its own type prime.
For a good place $w$ at a served rational prime, prescribe
\[
 A_w=\{x\in V_u:\operatorname{cmp}_{u,w}(\nu_u x)
                                    \in\mathcal N_{u,w}\},       \tag{LP.6}
\]
where $\mathcal N_{u,w}$ is the actual selected local module
\texttt{NP}, and served means the literal set \texttt{VV} is
nonempty. At a place above $T_u$, prescribe instead
\[
                  A_w=\{x:\nu_u x\in L^\lambda_{u,w}\}.         \tag{LP.7}
\]
The valuation compatibility of the comparison identifies every
scalar integral at $w$ with a scalar of the coefficient DVR.
The modules in (LP.6)--(LP.7) are therefore closed under these
scalars, hence local at $w$. This proves the local saturation
condition used in LP2 directly, including for denominators outside
$w$.

The supplied served-place conditions say
$A_w\subset M_w$ and $n_wM\subset A_w$ for some $n_w>0$.
They imply both clauses of (LP.3). The supplied type-place conditions
say exactly $n_wM\subset A_w$ and $n_wA_w\subset M_w$.
Let $S$ consist of these served places and the places above $T_u$.
It is finite: a served prime satisfies $p<X'$, and a number field
has finitely many places above each rational prime. Moreover the
good-place condition excludes $p\mid T_u$, so (LP.6) and (LP.7)
never compete. LP3 now gives one $N_u$ satisfying both prescriptions,
with
\[
 \begin{array}{ll}
 (N_u)_w\subset M_w,&p_w\ne T_u,\\
 (N_u)_w=M_w,&p_w\ne T_u\text{ and }\texttt{VV}(p_w)=\varnothing.
 \end{array}                                                   \tag{LP.8}
\]
No containment in $M_w$ is imposed at the type prime. In particular
(LP.8) cannot be strengthened to $N_u\subset M$ without an
additional type-place assumption. Nor is $N_u$ generally saturated:
already the prescription $A_{(2)}=2\mathbb Z_{(2)}$, with reference
$\mathbb Z$ and no other modifications, gives $N=2\mathbb Z$.

\paragraph{LP5: install the patch without losing the valuation scale.}
Once the canonical refilled size/comparison datum is supplied, replace
its \texttt{NuLat} field by the modules $N_u$ above. The source
\texttt{setLat} sets
\[
 \operatorname{lenMN}(u,w)
  =\frac{\operatorname{toNat}\bigl(\operatorname{lenLoc}_w
                                      (N_u,M_{\rm ram}(u))\bigr)}{e_w},
              \qquad e_w>0.                                  \tag{LP.9}
\]
Here \texttt{lenLoc} is the source's extended-natural local length
and \texttt{toNat} its conversion; (LP.9) does not assert finiteness
at an arbitrary place. Whenever this local length is a finite integer
$\ell$, substitution gives the exact required identity
\[
                   \ell=e_w\operatorname{lenMN}(u,w).           \tag{LP.10}
\]
At a non-type unserved place the two localized modules agree by
(LP.8), so the length is zero. At served places the actual integral
comparison is exactly (LP.6), not merely an abstract isomorphism.
All other fields of the size datum are retained. Consequently its
scale, selected degrees, rank and host bounds, and canonical orbit
selector are unchanged; the instance is transported by keeping its
same field values. These are the lattice and length portions of
\texttt{exists\_ledPart}.

\paragraph{Updated boundary of the point construction.}
LP patches the actual selected local lattices and installs their
ramification-normalized lengths. RF constructs canonical refill and
conditional point installation. PN/SZ complete the numerical package;
SZ constructs initial size and rank data; AC gives coded columns and
the conditional level datum. AH now constructs the actual field/host
under its literal $q,d$ and splitting hypotheses. The remaining
existence inputs include the analytic prime and top/logarithmic windows
and the geometric \texttt{Lay/RecOK/LedOK/RecExP} with its mass/window
estimates. The new local period/theta/index/comparison arguments do
not themselves construct those geometric points.

\paragraph{Version-0.29 scope of the historical conditional proofs.}
KX supplies \texttt{LevEx}, \texttt{LocCFTAll}, \texttt{LamEx},
\texttt{PhiAll} and \texttt{ProModularRingExists}, retaining the
last statement's original nonempty-pair antecedent.  QX supplies
the literal adjoint lifting theorem and the complete downstream
two-Link/support join.  Its actual rank-four ordinary Hecke ideals
and \texttt{P1core} producer remain explicit.

SX8 supplies \texttt{SeedIIAlt}.  UX12 supplies every clause of
the original case-II \texttt{ProModAlt} witness from that seed and
\texttt{PMExAlt}.  YX then applies the existing original-point,
rigidity, catalogue and count machinery to prove
\texttt{PMExAlt}$\Rightarrow$\texttt{CountII}.
It makes only pointwise case-II substitutions in the older proofs;
their full two-case statements and original hypotheses remain.
The added sufficient height cut is \texttt{CountI},
\texttt{CountIII}, \texttt{CountV} and \texttt{PMExAlt}.

The compound historical seed and promodularity targets retain
their unsupplied case-IV and case-III scope.  Current supplier
replacements are additive and recorded in
Section~\ref{sec:dependency-audit-v29}.  The separate actual
family, gap, theta/local, denominator and metric/spectral
obligations remain substantive.

\paragraph{Version-0.28 scope of the historical conditional proofs.}
BX now supplies the entire \texttt{BlochKatoLocal},
\texttt{LocalPublishedTheorems} and \texttt{CohomologyHyp} records,
with named classical local theorems and explicit model/cochain/lattice
transport.  RX proves the actual \texttt{RigiditySymplecticFour}
conclusion by a direct cyclotomic-irreducibility and adjoint-Selmer
argument; it bypasses, without promoting, the larger original
ordinary automorphy-lifting packages.  CX supplies the entire SP9
catalogue and \texttt{NumberSp4Alt} on the original Galois level.
Its auxiliary label level $M^7$ is included in a fixed polynomial bound.

JX1 joins these conclusions to AX/DX/MX/PX/SP and proves
\texttt{SymplecticFourImageDense}$\Rightarrow$
\texttt{CountII}$\wedge$\texttt{CountIII}, on the same original
objects and with normalized coefficient return.  JX2 consequently
reduces the selected height/window cut to \texttt{CountI},
\texttt{CountV}, and image density, retaining the window's prime-pair
inputs.  IX closes actual valued-point, nonsplit-block, purity and
promodular transport chains; it retains the actual modular
quotient/link/support and seed producers.  Neither image density
nor the remaining counts are asserted unconditionally.

Historical proofs keep their original hypotheses and source ranges.
Current supplier replacements are additive and are recorded in
Section~\ref{sec:dependency-audit-v28}.  The separate common-family,
arithmetic-gap, genuine theta/local, actual denominator and exact
metric/global/spectral obligations remain substantive.

\paragraph{Version-0.27 scope of the historical conditional proofs.}
AX now constructs the original auxiliary symplectic data for both
cases and proves the ordinary, level, dual and safe unlinking
properties.  It uses the given coefficient ring when the required
dyadic roots are present, or the permitted finite extension with
SP8/RV1 normalized return otherwise.  DX/MX/PX keep the ring of
this actual datum throughout.  Its G1 conclusion retains exactly
\texttt{SymplecticFourImageDense}.  DX supplies the original
\texttt{PhiCond}, \texttt{Nums} and \texttt{LocBase}; MX uses them
to construct a literal \texttt{Model} for every original good datum.
The charts are characteristic-zero charts with their integral
lattices retained, not assertions of integral formal smoothness.

PX proves the internal G2 comparison through the actual
\texttt{D.Ineq} under the same datum's \texttt{Fin0},
\texttt{NoEnd} and \texttt{AllIrr RigidAt}.  Image-density supplies
the first two predicates; actual irreducible rigidity and the exact
SP9 catalogue remain.  The source \texttt{rigSp4\_alt} producer,
including its local Hodge/adjoint joins, is still a separate target.
The full \texttt{BlochKatoLocal} package is not required for the
new model construction, but no unused part of it is promoted.

The actual count cut remains \texttt{CountI/II/III/V}.  HV/TW/OB/VW
already supplies every subsequent height/window assembly step.
The common arithmetic family, genuine theta/local state, actual
denominator tables, arithmetic gap and exact metric/global/spectral
comparisons remain independent substantive obligations.  Historical
proofs retain their original hypotheses; the current frontier is in
Section~\ref{sec:dependency-audit-v27} and the version-0.27 audit.

\paragraph{Version-0.26 scope of the historical conditional proofs.}
The historical proofs and their original hypotheses are preserved.
RV now supplies literal \texttt{SoFin}, \texttt{SoV} and coefficient
independence; BR supplies the entire \texttt{CorX} branch; SP supplies
\texttt{LemmaM'} and \texttt{LemmaM}, and FP supplies the universal
\texttt{LemmaP} with its exact original hypotheses.  The latter's
arithmetic installation is still required.  SP's paired-count
endpoint explicitly retains actual data/admissible-class production,
the weighted model and point-orbit comparison, and the finite
catalogue.  These are not supplied by naming generic packages.

HV supplies \texttt{CountVs}, \texttt{CountIV} and every original-field
character, case, layer and finite-prime summation step to
\texttt{PointHV}; CZ already supplies \texttt{CountZ}.  The remaining
actual count cut is \texttt{CountI/II/III/V}.  Together with TW8,
OB8 and VW this is a complete implication to the same cohomological
window and fixed-pair endpoints.  No separate \texttt{PointHV}
assembly obligation remains.  The common family/gap, genuine
theta/local state, actual denominator tables and exact metric/global/
spectral comparisons retain their original roles.  The current
frontier is in Section~\ref{sec:dependency-audit-v26} and the
version-0.26 audit.

\paragraph{Version-0.25 scope of the historical conditional proofs.}
Historical mathematical proofs and their original hypotheses are preserved.
LK1--LK4 now supply the three exact native local-K interfaces. TB1--TB10
prove the all-prime torsion supplier, and KB1--KB8 construct the actual
\texttt{KBridge} and supply \texttt{LemmaK 31}. TW1--TW8 supply
the literal Ferrero--Washington rank specialization and the full
$\exists K\,\forall C_t\ge1,\ \texttt{ThmBab}\ K\ C_t$.
Thus the simultaneous two-primary/reference-layer input of VW3--VW4,
QP25 and the older cohomological budget is supplied, alongside the
earlier odd-minus theorem OB8. CZ1--CZ6 independently supply
\texttt{PropZ} and \texttt{CountZ}, uniformly with exponent zero.
The selected major cohomological frontier is now \texttt{PointHV};
its other counts and remaining assembly retain their stated inputs.
The actual common family/gap, genuine theta/local data, actual local
denominator tables and global metric/spectral comparisons remain.
The current frontier is stated in Sections~\ref{sec:dependency-audit-v25}
and~\ref{sec:remaining}, and in the version-0.25 audit.

\paragraph{Version-0.24 scope of the historical conditional proofs.}
Historical mathematical proofs and their original hypotheses are preserved.
LD1--LD4 now supply the stated finite-duality and unramified-reciprocity
interfaces relative to exact classical theorems. OS1--OS9 prove the
literal \texttt{PsiBound} supplier, and OB1--OB8 prove
\texttt{ThmBcFree} with constant $240$. Thus the odd-minus input
in VW3--VW4, QP25 and the earlier cohomological budget is now supplied;
older statements that it remains to be written are historical.
The remaining bespoke cohomological cut is \texttt{ThmBab} plus
\texttt{PointHV}, on the same original fields and simultaneous windows.
LN1--LN10 prove the conditional integral-cut, congruence-length,
prime-summation and original-family \texttt{LinkN} chain. Their
actual local Galois tables and type/theta/rank compatibility inputs
remain open; the stronger arithmetic cut is not inferred from a
good cut. The full arithmetic family/gap, genuine theta/local data,
global spectral estimates and exact metric transport remain required.
The current frontier is stated in Sections~\ref{sec:dependency-audit-v24}
and~\ref{sec:remaining}; the standalone version-0.24 audit records
the same-object and quantifier obligations explicitly.

\paragraph{Version-0.23 scope of the historical conditional proofs.}
Historical proof paragraphs and their original hypotheses are preserved.
HP1--HP6 now supply the literal characteristic-zero highest-weight,
irreducibility, classification and descent premises of MR5/MR19 and
AR13--AR14, including the full original-field weight endpoint.
VW1--VW4 give the complete cohomological-window assembly, conditional
on the exact deep suppliers \texttt{ThmBab}, \texttt{PointHV} and
\texttt{ThmBcFree}. WB1--WB8 give the full arithmetic-gap-to-character-window
implication, including ramification and uniform moving cutoffs, and
join the already complete AN2--AN3/GH6 analytic endpoint. DA1--DA2
make the common-family quantifiers and conditional global-link assembly
explicit. Thus older statements that these particular downstream
assemblies are still future work are historical. The actual arithmetic
gap/family, deep cohomological estimates, genuine theta/local data and
global metric/spectral comparisons remain open in this reconstruction.
The current frontier is in Section~\ref{sec:remaining} and the semantic
audit in Section~\ref{sec:dependency-audit}.

\paragraph{Version-0.22 scope of the historical conditional proofs.}
The historical proofs retain their original text and hypotheses.
MR now supplies actual fixed-state form realization and original-DVR
lattices/signatures, but its initial SideB has zero theta/empty toric
fields and requires a supplied NTh. QD supplies the quantitative
moved-vacuum dictionary and common scalar units on the original choices.
DC supplies the sufficient large-gap discrete-series range and exact
formal degree using the named Harish--Chandra theorem, and restricted
RealTied under GV's Tamagawa input. QP supplies relative progressions
and full simultaneous PairEx under the named classical CN/BD/HD inputs;
RealBox/ThinH hold on a compatible supplied host, and winPairs retains
VWinAt. Older descriptions of these specific suppliers as future work
are historical. Genuine same-object theta/local installation, critical
ratios, remaining Core/layout and global/cohomological/spectral joins
remain as stated in Section~\ref{sec:remaining}.

\paragraph{Version-0.21 scope of the historical conditional proofs.}
The historical proofs below retain their original text and hypotheses.
NR now constructs the actual all-prime coefficient family, inducing
readings and conditional bounded integral row representations. CW
constructs the full finite and real Weil action, with exact rational
pullback automorphy and bottom suppliers. GV proves the exact measure
normalizations and direct definite quotient-to-mass identity, and uses
the explicitly named classical unitary Tamagawa theorem for the
indefinite volume and true-volume stage consequences. CC proves the
Dirichlet analytic inputs, exact centered integrals and uniform
late-cutoff thinning under the original zero-exclusion and
\texttt{RealBox} hypotheses.
Older descriptions of those specific suppliers as future work are
historical. Genuine form/local realization, the remaining exact Core
and discrete-series/formal-degree inputs, quantitative theta and
critical-ratio joins, and the independent global branches remain
as stated in Section~\ref{sec:remaining}.

\paragraph{Version-0.20 scope of the historical conditional proofs.}
The historical proofs below retain their original premises and text.
GA now constructs normalized adic characters from the stipulated global
Artin map and supplies exact local character laws and conditional
installed-character identifications. WA proves qualitative finite Weil
arithmetic, the normalized jet field and period fields for supplied
bottom triangles. MF proves the exact unramified CM mass ratio, with
classical Hasse--Minkowski stated explicitly. QR constructs the actual
ray-class family and proves its uniform weighted interval estimate;
the earlier reductions therefore supply the required norm primes.
Older descriptions of these particular suppliers as future work are
historical. Compatible nonempty arithmetic realizations and the full
geometric core, new full bottom-Weil existence, critical ratio fields,
quantitative dictionaries and exact embedding joins, and the four
global branches remain as stated in Section~\ref{sec:remaining}.

\paragraph{Version-0.19 scope of the historical conditional proofs.}
The historical proofs below retain their original premises and text.
HF now constructs the actual host-character frame and continues the
same-base initial coefficient, jet and bare-stage joins. AR expands
actual local arithmetic level/character comparisons and conditional
coefficient-weight and installed series realizations. TT supplies
primitive finite/local theta-test calculations under their exact field
and analytic interfaces. QW proves the original quantitative mass window
from its stated mass law and the exact prime-count/scale transports
from their uniform analytic/arithmetic inputs. The bounded-modulus
real-zero margin is also proved. Older statements of those reductions
as future work are historical. Compatible nonempty realizations and
the full core, genuine global mass and weighted interval suppliers,
remaining theta identities and the four global branches remain as
stated in Section~\ref{sec:remaining}.

\paragraph{Version-0.18 scope of the historical conditional proofs.}
The historical proofs below retain their original premises and text.
RC now expands actual rank, character and local-level mechanics.
LS constructs the actual local period rings and the bare Galois stage
on the same adapted rows. IG constructs an original sized preliminary
input and the exact chosen-volume empty-jet stage. TD supplies rational
theta cuts, common-image and metric deductions and common-test
normalization under explicit primitive inputs. Older statements of
these steps as future work are historical. Compatible nonempty
realizations, actual frame/core and quantitative windows, primitive
theta tests/operator identities, and the four global branches remain
as stated in Section~\ref{sec:remaining}.

\paragraph{Version-0.17 scope of the historical conditional proofs.}
The historical proofs below retain their original premises and text.
VI and CA now construct the actual vacuum, normalized indicator,
covariance and conditional canonical analytic join. NC constructs the
choice/tuple records under their exact rank, character and support ties.
SE supplies restricted lattices, several bare stage records and the pins
normal form. Older references to these steps as future work are historical;
the current boundary is stated in Section~\ref{sec:remaining} and the
new sections. Full \texttt{NStage}/\texttt{LocStage}, arithmetic ties,
initial geometric and analytic-window suppliers, and the four independent
global branches remain distinct.

\paragraph{Version-0.16 scope of the historical conditional proofs.}
The following previously reconstructed proofs retain their stated inputs
and their original text. AH in Section~\ref{sec:arithmetic-hosts} now
constructs the conditional arithmetic host and character field; PA in
Section~\ref{sec:period-arithmetic} supplies its periods and trace
isometries; TA in Section~\ref{sec:theta-arithmetic} supplies the
one-dimensional normalized theta arithmetic; CI in
Section~\ref{sec:canonical-indices} supplies the local-index and
nonvanishing theorems; and TX in Section~\ref{sec:twisted-comparison}
supplies the standard all-twist comparison. References below to those
tasks as future work describe the historical stage, not the present
boundary. The current remaining hypotheses are stated in the new
sections and Section~\ref{sec:remaining}. In particular the actual
multivariable vacuum/indicator/covariance, normalized arithmetic
\texttt{NChoice}/\texttt{TupTied}, \texttt{Pins.iX}, \texttt{Later},
and the analytic/geometric/global branches are still distinguished.

\paragraph{Version-0.15 scope of the following conditional constructions.}
The following AD/DR/FC/CV/NB/VE/RF sections preserve their version-0.14
proofs and the hypotheses stated there. The additional law premise is now
supplied independently by TL13 in Section~\ref{sec:column-laws}, from
the literal \texttt{TC.Ok}. TG in Section~\ref{sec:terminal-geometry}
supplies the complete reference symmetric and terminal constructions.
PN and SZ complete every numerical parameter and size family, and AC
with SZ proves the three conditional arithmetic constructors isolated
by RF10. References below to those tasks as remaining describe their
historical stage; the current boundaries are stated in the new sections
and Section~\ref{sec:remaining}. Actual host, \texttt{Ok/Pins/Later},
analytic prime and geometric point suppliers are still retained where
specified.

\section{Explicit adelic base points and supported modifications}
\label{sec:adelic-geometry}

\status{AD1 constructs the actual diagonal base point and its last-row
residues. AD2--AD3 prove the local transition and supported-splice
arguments from supplied realizers. AD4 installs prescribed rational
frames without changing the carried lattice or the other decorations.
The initial realization of a self-dual decorated lattice remains a
separate geometric input, as do the identities for the particular
up/down neighbors. No rational strong-approximation assertion is used.}

\paragraph{Conventions.}
Let $E/F$ be the supplied CM extension, write $\bar x$ for its
involution, and put $\mathbb A=E\otimes_F\mathbb A_{F,f}$.
Use its canonical identification with $\mathbb A_{E,f}$, so an
adelic equality can be checked at every finite prime $P$ of $E$.
Write $\widehat{\mathcal O}=\prod_P\mathcal O_{E_P}$ for its
integral subring. Conjugation interchanges the $P$ and $\bar P$
components, using the induced isomorphism of their completions.
A matrix $g\in\mathrm{GL}_n(\mathbb A)$ is unitary when
$\bar g^{\mathsf T}g=1$. It carries the standard lattice to $N$
when $g\widehat{\mathcal O}^{n}=\widehat N$.
The full rational lattice is recovered from these local modules
by LP2--LP3 and the completion argument proved in AD4 below.

For a finite prime $P$, the characteristic function $e_P$ of its
component is an integral finite adele. Thus
\[
 e_P^2=e_P,\qquad e_Pe_Q=0\ (P\ne Q),\qquad
 \overline{e_P}=e_{\bar P}.                         \tag{AD.1}
\]
These are literal elements of the restricted product: only one
component differs from zero. They give finite componentwise
constructions inside the same adelic algebra. The source's
\texttt{idem} is the same idempotent, written as the idele with
component $2$ at $P$ and $1$ elsewhere, minus $1$.

\paragraph{AD1: a diagonal unitary base point with the prescribed residue.}
Let $0\ne\pi\in\mathcal O_E$, put $q=\bar\pi$, and assume
$(\pi,q)=\mathcal O_E$. The finite set $S$ of primes dividing
$\pi$ is disjoint from $\bar S$: otherwise one prime would
contain both terms in a Bezout identity for $\pi,q$.
Define
\[
 a_P=
 \begin{cases}
 q^{-1},&P\in S,\\
 \pi,&\bar P\in S,\\
 1,&P\notin S\cup\bar S.
 \end{cases}                                         \tag{AD.2}
\]
All displayed field elements are mapped to their indicated
completions. This is a finite modification of $1$, hence a finite
adele. At a prime in $S$, $q$ is a local unit; at a prime in
$\bar S$, $\pi$ is a local unit. Consequently both $a$ and
$a^{-1}$ belong to $\widehat{\mathcal O}$. At $P\in S$,
$\bar a_P=q$, while at $P\in\bar S$, $\bar a_P=\pi^{-1}$.
Checking the three cases proves
\[
                 a\bar a=1,\qquad a^{-1}=\bar a.      \tag{AD.3}
\]
In the idempotent notation, (AD.2) is also the finite formula
\[
 a=1-\sum_{P\in S}(e_P+\bar e_P)
          +\sum_{P\in S}e_Pq^{-1}
          +\sum_{P\in S}\bar e_P\pi.
\]
This proves the literal \texttt{exists\_\allowbreak unit}, including the
integrality of both inverses and $a_P^{-1}=q$ at each selected prime.

For $n>0$ put
\[
 k=\operatorname{diag}(1,\ldots,1,a),\qquad
 k^{-1}=\operatorname{diag}(1,\ldots,1,\bar a).        \tag{AD.4}
\]
Equation (AD.3) proves that $k$ is unitary. Both matrices are
integral, so $k\widehat{\mathcal O}^{n}=
\widehat{\mathcal O}^{n}$. It therefore realizes the actual
standard lattice, with empty type, auxiliary and twist decoration.
If $\pi q=T\ne0$ is a rational integer, its components are $1$
at every prime not dividing $T$. All $E$-components over an
$F$-prime $v\nmid T$ are therefore $1$, so the projection of
$k$ to $E\otimes_FF_v$ is exactly the identity. At a selected
prime $P\mid\pi$, the last row of $k^{-1}$ is
\[
              (0,\ldots,0,q),                       \tag{AD.5}
\]
already before reduction. Thus its entries have the precise
\texttt{HasRes} values required in \texttt{exists\_\allowbreak base} and
\texttt{exists\_\allowbreak kStd}. In dimension zero take the unique empty
identity matrix; the last-row assertions have no instances.
The actual source hypotheses $T=T_u$, $\pi=\pi_u$ give this
construction for every word and every dimension, simultaneously.

\paragraph{AD2: local lattice inclusion controls transition matrices.}
Suppose $a,b\in\mathrm{GL}_n(\mathbb A)$ carry the standard
lattice to full lattices $N,N'$. If $N'_P\subseteq N_P$, then
\[
                      a_P^{-1}b_P\in M_n(\mathcal O_{E_P}).
                                                               \tag{AD.6}
\]
Indeed every column $b_Pe_j$ belongs to $N'_P$ and hence to
$a_P\mathcal O_{E_P}^n$; its coordinates under $a_P^{-1}$
are integral. Here a localization inclusion is first extended
to the completion. The extension is justified by clearing a
denominator on each element of a finite generating set.
Equivalently, for $x\in N'$ choose $s\notin P$ with
$sx\in N$; $s$ is a local unit, so integrality of
$a_P^{-1}(sx)$ implies integrality of $a_P^{-1}x$.
Finite sums and integral scalar multiples give the completed
module statement. If $N'_P=N_P$, apply the same reasoning in
the opposite direction to conclude
\[
                      a_P^{-1}b_P\in\mathrm{GL}_n(\mathcal O_{E_P}).
                                                               \tag{AD.7}
\]
This is the source's \texttt{valued\_\allowbreak inv\_\allowbreak mul\_\allowbreak le}, with the
inverse assertion available when the lattices agree.

One residue cancellation will also be needed. If $s\in\mathcal O_E
\setminus P$, $z\in E_P$, and $sz$ has residue $\bar s c$,
then $z$ is integral and has residue $c$. The first assertion
uses that $s$ is a unit; the second cancels its nonzero residue
in $\mathcal O_E/P$. In the source's explicit \texttt{HasRes}
notation, choose $t\in\mathcal O_E$ and $r\in P$ with
$ts+r=1$. If $sz-b$ has positive valuation, then
\[
             z-tb=t(sz-b)+rz                         \tag{AD.8}
\]
has positive valuation too, and the residue of $tb$ is $c$.
This proves \texttt{hasRes\_\allowbreak of\_\allowbreak mul}. It extends a residue
identity on a global lattice to any vector in its localization,
by multiplying the vector by a denominator outside $P$.

\paragraph{AD3: confine a realized modification to one conjugate pair.}
Let $X,X_0$ be the source's decorated lattices with type, auxiliary
and frame-prime sets $S_T,S_G,Q$. They are pairwise disjoint,
and their union contains no conjugate pair. Choose $P\in S_T$.
Suppose $h$ realizes $X$ and $h_0$ realizes $X_0$, in the full
\texttt{IsPtOf} sense. Assume that their localized lattices agree
away from $P,\bar P$, their type functionals agree at every
other prime of $S_T$, and their auxiliary reductions and rational
frames agree. These are the actual hypotheses used in the splice.

Put $e=e_P+e_{\bar P}$. It is an idempotent fixed by conjugation.
Define, entrywise,
\[
 g=(1-e)h+eh_0,\qquad
 g^{-1}=(1-e)h^{-1}+eh_0^{-1},\qquad
 \varepsilon=h^{-1}g.                                \tag{AD.9}
\]
The first two matrices are inverse because $e(1-e)=0$.
Since $\bar e=e$ and $h,h_0$ are unitary, the conjugate
transpose of the first is the second. Thus $g$ and
$\varepsilon$ are unitary. Their components satisfy
\[
 g_w=\begin{cases}(h_0)_w,&w=P,\bar P,\\h_w,&\text{otherwise},
       \end{cases}
 \qquad \varepsilon_w=1\quad(w\ne P,\bar P).          \tag{AD.10}
\]
There is no restriction on the size of the changed components;
only finitely many are changed.

We check all five fields of \texttt{IsPtOf} for $g,X_0$.
Unitarity was proved above. For the lattice field,
$g_w^{-1}(h_0)_w$ and its inverse are identity at the two
changed primes and are integral elsewhere by (AD.7).
Therefore $g^{-1}h_0\in\mathrm{GL}_n(\widehat{\mathcal O})$,
and $g\widehat{\mathcal O}^n=h_0\widehat{\mathcal O}^n
=\widehat{N_0}$ exactly.
At $P$ the type residue is that of $h_0$. At any other
$P'\in S_T$, its component is that of $h$; a vector in
$N_0$ belongs to $N_{P'}$, so multiply it by $s\notin P'$
to put it in $N$. The old residue identity and equality
of the prescribed type functionals give the new identity
after cancellation by (AD.8).
The identical argument handles every auxiliary residue vector.
The no-conjugate-pair and disjointness assumptions ensure that
these primes lie outside $\{P,\bar P\}$. Finally all primes
in $Q$ are outside this pair, where $g=h$; the exact rational
frame matrices are unchanged by hypothesis. Thus the frame
field also holds.

Since $h\varepsilon=g$, this proves the complete supported
modification conclusion of \texttt{exists\_\allowbreak splice}, from
the supplied realizer $h_0$. The source obtains $h_0$ from
\texttt{exists\_\allowbreak isPtOf\_\allowbreak cm}; that geometric supplier is
kept explicit. For an actual up/down move one must still
prove its stated equality of local lattices and decorations.
The support conclusion alone does not prove those premises.

\paragraph{AD4: install the exact rational frames at finitely many primes.}
First justify the source's lattice-recovery step. For any
$x\in\mathrm{GL}_n(\mathbb A)$ put
$M_P=x_P\mathcal O_{E_P}^n$. Both $x$ and $x^{-1}$ are integral
outside a finite set. Choose a nonzero rational integer $d$ clearing
their entries at the exceptional primes; then, at every $P$,
\[
 d\mathcal O_{E_P}^n\subseteq M_P\subseteq d^{-1}\mathcal O_{E_P}^n.
                                                               \tag{AD.13}
\]
The rational modules $B_P=M_P\cap E^n$ are local at $P$ in the
sense of LP2 and equal $\mathcal O_{E,(P)}^n$ outside that finite
set. They satisfy the bounds required by LP3. Thus
$N=\bigcap_P B_P$ is a finite full lattice with
$N_{(P)}=B_P$ and $d\mathcal O_E^n\subseteq N\subseteq
d^{-1}\mathcal O_E^n$.
For each fixed $P$, use density of $E$ in its \emph{single}
completion $E_P$ to choose rational columns $v_j$ sufficiently
close to $x_Pe_j$ that
\[
             x_P^{-1}(v_1,\ldots,v_n)
                    \in 1+\mathfrak m_P M_n(\mathcal O_{E_P})
                    \subset\mathrm{GL}_n(\mathcal O_{E_P}).
                                                               \tag{AD.14}
\]
The last inclusion follows by reducing the determinant modulo
$\mathfrak m_P$; it also exhibits the needed openness.
The $v_j$ belong to $B_P$ and form an $\mathcal O_{E_P}$-basis
of $M_P$, so the completion of $N_{(P)}$ is exactly $M_P$.
For a finite lattice between $d\mathcal O_E^n$ and
$d^{-1}\mathcal O_E^n$, its integral adelic span is the product
of these completed local modules: multiply by $d$ and apply the
finite Chinese remainder theorem to the submodule modulo $d^2$.
Therefore $\widehat N=x\widehat{\mathcal O}^n$.
The intersection defining $N$ is precisely
$\operatorname{globLat}(x)$, so this proves the literal
\texttt{Carries} assertion used by \texttt{carries\_\allowbreak globLat}.
The approximation was separate at each single prime, not a
rational strong-approximation assertion. For $n=0$ the statement
is immediate.

Let $x$ be unitary and carry the standard lattice to $N$.
Let $Q$ be a finite set containing no conjugate pair.
At every $P\in Q$ supply a rational matrix $M_P\in
\mathrm{GL}_n(E)$ which is a local basis of $N$ there.
Define a new adelic matrix by
\[
 g_P=M_P\quad(P\in Q),\qquad
 g_{\bar P}=\overline{M_P}^{\,-\mathsf T}\quad(P\in Q),
 \qquad g_w=x_w\quad(w\notin Q\cup\bar Q).           \tag{AD.11}
\]
At the conjugate prime the field matrix is transported to that
completion. Formula (AD.1) constructs each entry, so $g$ is
an actual finite adele. The inverse is obtained by the same
formulas with each selected matrix inverted. The two selected
components are conjugate inverse transposes; hence $g$ is
unitary, including the unchanged components.

Put $b=x^{-1}g$, also unitary. At $P\in Q$, $x_P$ and
$M_P$ are bases of the same local lattice, so both $b_P$
and $b_P^{-1}$ are integral by (AD.7). Unitarity gives
\[
            b_{\bar P}=\overline{b_P^{-1}}^{\,\mathsf T}.
                                                               \tag{AD.12}
\]
It follows that $b_{\bar P}$ and its inverse are integral
as well. Away from $Q\cup\bar Q$ they are identity.
Therefore $b\in\mathrm{GL}_n(\widehat{\mathcal O})$ and
$g$ carries the very same lattice $N$. Its selected component
is exactly $M_P$, rather than merely congruent to it.
Outside $Q\cup\bar Q$, $g^{-1}=x^{-1}$, so all pullback
coordinate functions remain identical. This proves
\texttt{exists\_\allowbreak modify} with all its equality clauses.

For the actual decorated-point application, the type and auxiliary
primes avoid $Q\cup\bar Q$. Their residues are consequently
preserved, while the specified rational frames are installed
exactly. Thus the frame portion of the general realization
theorem is now explicit. Its remaining input is a unitary
adele already realizing the self-dual lattice and the finite
type/auxiliary reductions. The source supplies it using
\texttt{unimodGenus}, \texttt{decoPtSurj\_\allowbreak of} and
\texttt{reduction\_\allowbreak surjective\_\allowbreak bis}; those arithmetic
realization arguments are not inferred from (AD.11).


\section{Realizing the decorated lattices by actual unitary adeles}
\label{sec:decorated-realization}

\status{DR1 proves the needed local Hermitian normal form, including
residue characteristic two, from the literal unramified CM-field
condition. DR2 constructs a unitary finite-adelic matrix carrying the
standard lattice to any supplied self-dual lattice. DR3 installs the
prescribed residue functionals, auxiliary reduction bases and rational
frames, and proves every field of the source's \texttt{IsPtOf}.
Thus the initial decorated-point realizability used in FC and NB is
now supplied under the original CM-field and valid-decoration hypotheses.
The existence of those arithmetic host and column-law data remains a
separate input.}

\paragraph{Conventions and exact arithmetic input.}
Let $E/F$ be the supplied quadratic CM extension, with involution
$a\mapsto\bar a$ and fixed field $F$. Put $\mathcal O=\mathcal O_E$
and use $h(x,y)=\bar x^{\mathsf T}y$ on $E^n$.
The literal \texttt{CMData.cmUnram} condition includes
\[
 \text{for every maximal }P\subset\mathcal O\text{ there is }
 \theta\in\mathcal O\text{ with }\theta-\bar\theta\notin P.
 \tag{DR.1}
\]
This is the local arithmetic assertion used below. It is part of the
supplied CM data; we do not infer it from the name of that record.
Its other clauses identify the involution with complex conjugation
under the embeddings. We use \texttt{CMData.fixed} and
$[E:F]=2$ to identify its fixed field exactly with the given $F$.

For a finite prime $v$ of $F$, write
\[
 K=F_v,\quad R=\mathcal O_{F_v},\quad
 B=E\otimes_F K,\quad
 A=\prod_{P\mid v}\mathcal O_{E_P}\subset B.
 \tag{DR.2}
\]
Thus $A$ is the integral algebra in the completed CM algebra, with
the same involution, not an arbitrarily chosen quadratic extension.
The completion decomposes over the primes above $v$ as in the local
identification construction. Its fixed algebra is $K$, and its fixed
integral subring is $R$. For the first assertion, take any nonfixed
$\theta\in E$: the basis $1,\theta$ remains a $K$-basis after
extension, and comparing the two coefficients of a fixed element
gives the assertion. The second follows by intersecting with the
integral ring, since a scalar in $K$ is integral in all these finite
extensions exactly when its $v$-valuation is nonnegative.

\paragraph{DR1: local orthonormal bases, including the dyadic case.}
First derive the precise local algebra from (DR.1). Choose a prime
$P\mid v$ and its $\theta$ from that condition. If there is another
prime above $v$, it is $\bar P$, since $E/F$ is quadratic Galois.
The element $\delta=\theta-\bar\theta$ is a unit there as well:
$\bar\delta=-\delta$, so membership at $\bar P$ would imply
membership at $P$. Consequently $\delta\in A^\times$.
For every $a\in A$, the elements
\[
 b=\frac{a-\bar a}{\theta-\bar\theta},\qquad
 c=a-b\theta                                             \tag{DR.3}
\]
are integral and fixed by conjugation, and hence belong to $R$.
Conversely $R+R\theta\subseteq A$. Therefore
\[
 A=R\oplus R\theta,
 \qquad A\simeq R[X]/\bigl(X^2-(\theta+\bar\theta)X
                                  +\theta\bar\theta\bigr).
 \tag{DR.4}
\]
The coefficients lie in $R$ and the discriminant is
$\delta^2\in R^\times$. Reduction modulo a uniformizer $\varpi$
of $R$ is consequently a separable quadratic algebra over the finite
residue field $k$. It is either $k\times k$, with conjugation
exchanging the factors, or the quadratic field $k_2/k$, with its
nontrivial field involution.

In the first case the two simple residue roots lift to roots in $R$.
Explicitly, if $f(a_m)\in\varpi^mR$ and $f'(a_m)$ is a unit, choose
$t_m\in R$ so that
$f(a_m)/\varpi^m+f'(a_m)t_m\equiv0\pmod\varpi$ and put
$a_{m+1}=a_m+\varpi^m t_m$. The errors improve by one power and
the sequence converges to a root. The two lifted roots have unit
difference, so the Chinese remainder theorem in (DR.4) gives
$A\simeq R\times R$. In the second case $A/\varpi A=k_2$ is a
field. An element with nonzero residue is a unit: lift its residue
inverse and invert the resulting $1+\varpi z$ by a convergent
geometric series. Every nonzero element is then a power of
$\varpi$ times a unit, using the finite free $R$-module (DR.4).
Thus $A$ is the unramified quadratic discrete valuation ring over $R$.
This establishes the split/unramified alternatives from the literal
source hypothesis, without assuming them separately.

We need the norm on units in the nonsplit case. Let $|k|=q$.
The finite-field norm $k_2^\times\to k^\times$ is onto, since
it is $z\mapsto z^{q+1}$ on the cyclic group of order $q^2-1$.
The finite-field trace is onto as well. Indeed
$z\mapsto z+z^q$ is a $k$-linear map into the one-dimensional
space $k$, and it is not zero: the nonzero polynomial $X+X^q$
has degree $q<q^2$ and cannot vanish at every element of $k_2$.
This proof works when $\operatorname{char}k=2$.

Let $u\in R^\times$. Choose $b_1\in A^\times$ with
$b_1\bar b_1\equiv u\pmod\varpi$ using the finite-field norm.
Suppose $b_m\bar b_m-u=\varpi^m c_m$. Choose $z_m\in A$
whose residue trace is $-c_m/u$, and set
\[
                    b_{m+1}=b_m(1+\varpi^m z_m).
 \tag{DR.5}
\]
For $m\ge1$ the identity
\[
 (1+\varpi^m z)(1+\varpi^m\bar z)
       \equiv1+\varpi^m(z+\bar z)\pmod{\varpi^{m+1}}
 \tag{DR.6}
\]
holds because $2m\ge m+1$. Hence
$b_{m+1}\bar b_{m+1}\equiv u\pmod{\varpi^{m+1}}$.
The sequence is Cauchy and remains in the same nonzero residue class.
Its limit $b\in A^\times$ satisfies $b\bar b=u$.
No division by $2$ or assumption of odd residue characteristic is
used. In the split case the norm is $(a,b)\mapsto ab$ and is
plainly onto the units of $R$.

Now let $M$ be a free rank-$n$ lattice over the nonsplit ring $A$
with a perfect integral Hermitian pairing, conjugate linear in the
first argument. Its Gram determinant in an integral basis is a unit,
so its reduction is a nondegenerate Hermitian form over $k_2/k$.
If $n>0$, this reduced form has a vector of nonzero self-pairing.
For otherwise, for every $x,y$ and every $a\in k_2$,
\[
 0=h(x+ay,x+ay)=a\,h(x,y)+\bar a\,\overline{h(x,y)}.
 \tag{DR.7}
\]
Nondegeneracy supplies $x,y$ with $h(x,y)\ne0$. As $a$ varies,
the right side is the trace of every element of $k_2$, contradicting
its surjectivity. In particular this argument remains valid in
characteristic two, where the trace of $1$ itself can be zero.

Lift such a vector to $v\in M$. Its norm $h(v,v)=u$ is a unit
of $R$. By (DR.5) choose $b\in A^\times$ with $b\bar b=u$,
and replace $v$ by $b^{-1}v$. It now has norm one. The projection
$x\mapsto v h(v,x)$ gives the integral direct sum
\[
             M=Av\oplus M_1,\qquad M_1=v^\perp\cap M.
 \tag{DR.8}
\]
The complement is free of rank $n-1$, since $A$ is a discrete
valuation ring. It is again perfect: a vector of its fraction-field
span pairing integrally with $M_1$ pairs integrally with all of $M$
and is therefore in $M$, hence in $M_1$. Induction gives an integral
orthonormal basis of $M$. The rank-zero case is immediate.

For the split ring $A=R\times R$, write $M=M_+\oplus M_-$.
The Hermitian pairing is a perfect duality between the two free
$R$-lattices. Choose a basis $e_i$ of $M_+$ and its dual basis
$f_i$ of $M_-$. The vectors $(e_i,f_i)$ form an $A$-basis of $M$
with Gram matrix $I$. Thus the same orthonormal-basis conclusion
holds in both local alternatives. This proves the normal-form
ingredient used by \texttt{unimodGenus}, with the actual completed
CM algebra and its actual involution.

\paragraph{DR2: an actual finite-adelic matrix carrying a self-dual lattice.}
Let $N\subset E^n$ be a supplied full $\mathcal O$-lattice with
$N=N^\#$ for the standard Hermitian pairing. Choose a nonzero
rational integer $d$ with
\[
                 d\mathcal O^n\subseteq N
                         \subseteq d^{-1}\mathcal O^n.
 \tag{DR.9}
\]
Such a $d$ follows by clearing the finitely many coordinates of
generators and of an $E$-basis contained in $N$. Thus $N_P$ is
standard outside the finitely many primes dividing $d$.

At every finite prime $v$ of $F$, let $\widehat N_v$ be the product
of its completed local lattices above $v$. It is a free rank-$n$
module over the integral algebra $A$ in (DR.2). Hermitian duality
commutes with these localizations and completions: in a rational
local basis its defining inequalities are that the coordinates
after the conjugate-transpose Gram matrix are integral. Passing
from the local ring to its completion does not change this dual
matrix formula. At a conjugate pair it pairs the two different
components, exactly as in NB1; at a fixed prime it is the ordinary
unramified Hermitian dual. Consequently $\widehat N_v$ is perfect.

If the prime is nonsplit, DR1 gives an orthonormal integral basis
of $\widehat N_v$. Let $g_v\in\mathrm{GL}_n(B)$ be its column
matrix. Then $g_v^*g_v=I$ and $g_v A^n=\widehat N_v$.
At a split pair $P,\bar P$, one can make the choice even more
explicit. Choose any completed basis matrix $B_P$ of $N_P$ and set
\[
              g_P=B_P,\qquad
              g_{\bar P}=\overline{B_P}^{\,-\mathsf T}.
 \tag{DR.10}
\]
The bar transports entries to the conjugate completion.
The equality $N=N^\#$ says exactly that the second matrix is
a basis of the lattice at $\bar P$. Formula (DR.10) is unitary.
At every place where $N$ is standard, choose the identity matrix.

These choices define a matrix $g$ over the actual finite adelic
algebra $E\otimes_F\mathbb A_{F,f}$. To see this directly, expand
every $g_v$ and its inverse in a fixed $F$-basis $1,\theta$ of $E$.
Each coefficient is its identity-matrix coefficient outside finitely
many $v$, so it is an actual finite adele of $F$. Taking the corresponding
two tensor terms constructs every matrix entry in the stated algebra.
Both the entries and those of the inverse are integral outside the
finite set in (DR.9), where they are the identity.
The inverse is the conjugate transpose
componentwise, so $g$ is unitary. The chosen local bases give
\[
          g\widehat{\mathcal O}^{n}
                    =\prod_P\widehat N_P=\widehat N.
 \tag{DR.11}
\]
Here the last equality is the integral adelic-span identity proved
in AD4 using the finite sandwich and Chinese remaindering. It holds
for adelic vectors, not just rational ones. Thus (DR.11) is the
literal \texttt{Carries} assertion.
For rational vectors, intersection with $E^n$ recovers $N$: membership
in a completed local lattice is the nonnegativity of the valuations
of its rational coefficients in a local basis, hence is local
membership; LP1 then applies. The source's
$\operatorname{globLat}(g)$ is therefore precisely $N$.
This proves the actual \texttt{Formalism.L2.Col.UnimodGenus}
existence needed in \texttt{exists\_\allowbreak isPtOf\_\allowbreak cm}: that predicate is
the adelic $\operatorname{globLat}$ existence assertion just proved.
The similarly named \texttt{UnimodGenusLoc} and
\texttt{PI.UnimodGenus} ask for rational isometries of semilocal
lattices and are not asserted here. The direct adelic construction
bypasses that part of the source route and invokes no rational
approximation by unitary matrices.

\paragraph{DR3: realize every prescribed residue and rational frame.}
Let $p$ be a valid framed decorated lattice of rank $n$, with
underlying self-dual lattice $N$. Its selected type primes $S_T$,
auxiliary primes $S_G$, and rational-frame primes $Q$ are pairwise
disjoint finite sets. Assume the source's exact orientation condition
\[
 P\in S_T\cup S_G\cup Q
       \quad\Longrightarrow\quad
 \bar P\notin S_T\cup S_G\cup Q.                    \tag{DR.12}
\]
In particular each selected prime lies in a split conjugate pair;
neither component is selected twice. The valid decoration supplies
a nonzero residue-linear $f_P:N_P\to\mathcal O/P$ at $P\in S_T$,
a residue-basis map $\mathrm{be}_P:N_P\to(\mathcal O/P)^n$ at
$P\in S_G$, and an exact rational local basis matrix at $P\in Q$.

Choose a rational local basis $B_P$ of $N_P$ for each $P$ in the
union. Such a basis exists because a full lattice over the local
discrete valuation ring is free; its vectors lie in $E^n$ and form
an $E$-basis. Write $k_P=\mathcal O/P$.
The functional $f_P$ descends to a nonzero row
$r_P:k_P^n\to k_P$ in the coordinates $B_P$.
Residue-linearity over $\mathcal O_{(P)}$ follows from its
$\mathcal O$-linearity by clearing a denominator outside $P$.
Choose a basis of $\ker r_P$ and then a vector on which $r_P=1$.
Its column matrix $U_P\in\mathrm{GL}_n(k_P)$ satisfies
\[
                         r_PU_P=e_n^{\mathsf T}.    \tag{DR.13}
\]
Lift its entries through the actual quotient map
$\mathcal O\twoheadrightarrow k_P$ to a rational matrix
$\widetilde U_P$. Its determinant has nonzero residue, so the
matrix and its inverse are integral at $P$. Replacing $B_P$ by
$B_P\widetilde U_P$ preserves its local lattice and makes the last
coordinate of its inverse reduce exactly to $f_P$.

At an auxiliary prime, the supplied map $\mathrm{be}_P$ induces
an isomorphism $N_P/PN_P\simeq k_P^n$.
For clarity its global kernel and surjectivity clauses suffice:
clear a denominator $s\notin P$ to put a vector in $N$;
zero residue then puts $sx$ in $PN$, hence $x$ in $PN_P$.
Surjectivity already has global witnesses. Relative to $B_P$ this
isomorphism is a matrix $M_P\in\mathrm{GL}_n(k_P)$.
Lift $M_P^{-1}$ as above and replace $B_P$ by that lifted
right multiple. Its inverse-coordinate reduction is now the exact
given map $\mathrm{be}_P$, with every coordinate retained.
At a prime of $Q$ use the supplied rational frame matrix itself;
the defining framed-decoration condition says that it is a local
basis of $N_P$. Disjointness makes these prescriptions compatible.

Apply DR2 to obtain an initial unitary matrix carrying the standard
lattice to $N$. Apply AD4 once, with its modification set equal to
$S_T\cup S_G\cup Q$ and the just chosen rational local matrices.
Condition (DR.12) is exactly its no-conjugate-pair premise.
It gives a unitary $g$ carrying the very same $N$, whose component
at each selected $P$ is exactly $B_P$, and whose conjugate component
is its conjugate inverse transpose. Thus its pullback on a global
$x\in N$ has at $P$ the integral coordinates $B_P^{-1}x$.

The residue statements just proved are the source's literal
\texttt{HasRes} statements. To check this last point explicitly,
write any rational integral coordinate as $a/s$ with $s\notin P$.
Choose $r\in\mathcal O$ with
$\rho_P(r)=\rho_P(a)\rho_P(s)^{-1}$.
Then $a/s-r\in P\mathcal O_{(P)}$, so its image in the completion
has multiplicative valuation strictly less than one. This is precisely
$\texttt{HasRes}\,P\,(a/s)\,\rho_P(r)$, using the actual quotient
map. We obtain
\[
\begin{array}{ll}
 P\in S_T:&\operatorname{res}_P((g^{-1}x)_n)=f_P(x),\\
 P\in S_G:&\operatorname{res}_P((g^{-1}x)_j)
                                      =\mathrm{be}_P(x)_j
                                      \quad\text{for every }j,\\
 P\in Q:&g_Pe_j=p.\operatorname{fr}(P)(e_j)
                                      \quad\text{for every }j.
\end{array}                                                     \tag{DR.14}
\]
Together with unitarity and (DR.11), these are all five fields
\texttt{uni}, \texttt{lat}, \texttt{f}, \texttt{be}, and \texttt{fr}
of \texttt{IsPtOf}. If $n=0$, a nonzero functional on the zero
module is impossible, so a valid decoration forces $S_T=\varnothing$.
The empty identity matrix gives the conclusion directly; all auxiliary
and frame coordinate clauses are vacuous.
This proves the full \texttt{exists\_\allowbreak isPtOf\_\allowbreak cm} assertion.

For the true-column application, the source's all-twist wrapper
introduces no further geometry. A type prime contains $T_u$, while
the auxiliary and twist primes do not; these sets and their conjugates
are therefore disjoint. Two selected type primes cannot be conjugate
because $(\pi_u,\bar\pi_u)=\mathcal O$. The supplied twist-level
split condition handles the remaining auxiliary/twist union and its
orientation, while its disjointness clause separates those two sets.
These are exactly the checks in \texttt{exists\_\allowbreak ptOf\_\allowbreak twists}.
Hence DR3 supplies its actual points, the \texttt{yStd} points of FC3,
and the initial realizers used in NB6. AD3 and NB then supply the
supported moves as already proved. No additional point-realizability
hypothesis is needed beyond the literal supplied CM and valid-decoration
data. The full-column terminal incidence and special symmetric-level
identifications, and the upstream existence of the arithmetic host,
remain distinct from this completed realization construction.


\paragraph{Using the realization theorem in the following constructions.}
The FC and NB arguments below state the initial-realizer premise at each
place where it is used. For the supplied CM and valid-decoration data,
DR3 now supplies that premise. Those conditional intermediate statements
therefore introduce no further realizability assumption in the actual
comparison. VE11 records their combined application with the original
arithmetic host and true-column law data.

\section{Constructing charts, residue functionals and Hermitian carriers}
\label{sec:global-forms}

\status{FC1--FC18 prove the arithmetic and linear-algebra constructions
specified below, including actual finite charts from supplied adelic
representatives, replacement of their residue functionals without changing
their level groups, exact residue projections, the Fock metric and period
identities, the common auxiliary-prime data, and the Hermitian carrier
conversion. FC19 joins these with CV9 and NB7--NB8 to construct the
actual step-zero \texttt{Forms0}, its matching level carrier, and all
the global input records needed by VE9. This construction retains the
initial decorated-point realizability input. The complete simultaneous
\texttt{FormsTiedSym} and \texttt{ColsTiedSym} constructors still require
the terminal incidence construction and the specified symmetric
auxiliary-prime level descriptions. The original \texttt{TC.Ok},
\texttt{Pins}, \texttt{Laws} and \texttt{Later} data remain hypotheses.}

\paragraph{The data used here.}
Write $E/F$ for the supplied CM extension, $\mathcal O=\mathcal O_E$,
and $x\mapsto\bar x$ for its involution. Fix a word $u$, a step $i\le m$,
and $n=N_i>0$. The exact coordinate identifications of SF2 put the
Hermitian form in the form $h(x,y)=\bar x^{\mathsf T}y$ and preserve the
original class lattice and frame matrices. Put $p=T_u$ and
$\pi=\pi_u$, so $p$ is a rational prime and
$\pi\bar\pi=p$. The prime data include
$(\pi,\bar\pi)=\mathcal O$, the selected set
$S=\{P:P\mid\pi\}$, $|S|=[F:\mathbb Q]$, and
$|\mathcal O/P|=p$ for $P\in S$.

The source's \texttt{Laws} are substantive supplied data. They include
finite orbit and stabilizer statements for the actual column actions,
the level and period action laws, and the polynomial lattices and their
pairings. They do not mean merely that some unnamed compatible objects
exist. Whenever the proof below uses adelic realizability or a supported
move not yet reconstructed, that use is stated separately.

\paragraph{FC1: the level-prime index and its real residue field.}
First, $(\pi)=\prod_{P\in S}P$. Indeed, conjugation preserves absolute
ideal norm, and $[E:\mathbb Q]=2[F:\mathbb Q]=2f$, so
$\mathrm N((\pi))^2=p^{2f}$ and $\mathrm N((\pi))=p^f$.
The $f$ distinct primes in $S$ already contribute $p^f$; none can have
a larger exponent. Since distinct maximal ideals are comaximal,
\[
 r\in\bigcap_{P\in S}P\quad\Longleftrightarrow\quad\pi\mid r.
 \tag{FC.1}
\]
This is the \texttt{BasePi.dvd} assertion, with its exact squarefree
content, rather than merely a statement about the support of $(\pi)$.

Restriction gives a bijection from $S$ to the primes $v$ of $F$ above
$p$. For surjectivity, choose a prime $Q$ of $E$ over $v$. Since
$p=\pi\bar\pi\in Q$, either $\pi\in Q$, or
$\pi\in\bar Q$; conjugate primes have the same restriction. For
injectivity, the primes over a fixed $v$ in this quadratic Galois
extension are one conjugate orbit. Two selected primes cannot be
conjugate, since that would put both $\pi$ and $\bar\pi$ in one
prime, contrary to their coprimality. Thus the map is the literal
\texttt{underT}; its inverse does not select a different orientation.

There is also an explicit residue-field descent. Let
$\rho:\mathcal O\twoheadrightarrow k$ be the chosen quotient at $P$,
where $\rho(\pi)=0$ and $\rho(\bar\pi)\ne0$, and let
$v=P\cap\mathcal O_F$. For $y\in\mathcal O$, the elements
\[
 t=\pi+\bar\pi,\qquad z=y\bar\pi+\bar y\pi
 \tag{FC.2}
\]
are integral and fixed by conjugation, hence belong to $\mathcal O_F$.
The residue of $t$ is nonzero. Choose $s\in\mathcal O_F$ with
$st\equiv1\pmod v$. Then $\rho(sz)=\rho(y)$. Consequently the
restriction of $\rho$ to $\mathcal O_F$ is onto, has kernel $v$, and
induces the exact isomorphism
\[
 k\simeq\mathcal O_F/v,
 \qquad \rho(a)\longmapsto a\bmod v\quad(a\in\mathcal O_F).
 \tag{FC.3}
\]
No choice of an abstract field of the same cardinality is substituted
for this map. Together with the supplied embedding/prime correspondence,
these are the prime-index and residue fields of \texttt{atTTied} and
\texttt{CarH.fv}. The coefficient-place labels \texttt{tauOf} and
\texttt{seenP} still use that supplied arithmetic correspondence.

\paragraph{FC2: coupling data are determined by their reductions, and
every prescribed reduction can be lifted.}
Here the reference twist-prime set is empty, as it is for $\mu_0$ in
the source. For a decorated lattice $D$, write $D_P$ for its localized
lattice and $b_P:D_P\to k_P^n$ for its supplied reduction frame.
At an auxiliary prime $P$, the rational integer $q_P$ lies in
$P\setminus P^2$; thus it is a uniformizer in $\mathcal O_P$ and
\[
 b_P(x)=0\quad\Longleftrightarrow\quad x\in q_P D_P
 \quad(x\in D_P).                                      \tag{FC.4}
\]
This follows by writing $x$ in the local frame: zero reduction means
that every coordinate is divisible by the uniformizer.

At the coupling step $i=m-1$, the additional datum is a coset
$y+M$, where $q_Py\in D_P$ and $M$ imposes $x\in D_P$ at each
auxiliary prime. Its reduction at $P$ is the singleton
$\{b_P(q_Py)\}$. If two such cosets have the same reductions,
(FC.4) gives $y-y'\in D_P$ at every relevant $P$, hence
$y-y'\in M$ and the cosets agree. At all other steps the additional
set is empty. Thus a column point is determined by its decorated
lattice and the family of coupling reductions. This proves the actual
\texttt{eq\_\allowbreak of\_\allowbreak deco\_\allowbreak bar} assertion in the reference-tuple case.

For existence, prescribe vectors $v_P\in k_P^n$ for an arbitrary
decorated lattice $D$. Its global reduction frames supply
$x_P\in D$ with $b_P(x_P)=v_P$. Choose $e_P\in\mathcal O$ with
$e_P\equiv1\pmod P$ and such that $e_P/q_P$ is integral at every
other auxiliary prime. Such elements follow from the Chinese remainder
theorem, using sufficiently large powers at the other primes. Then
\[
 y=\sum_P\frac{e_P}{q_P}x_P                            \tag{FC.5}
\]
satisfies $q_Qy\in D_Q$ and $b_Q(q_Qy)=v_Q$ for every $Q$:
the $Q$-summand gives $v_Q$, and each other summand is a multiple of
$q_Q$ in $D_Q$. The coset $y+M$ is therefore the required point.
This proves \texttt{exists\_\allowbreak ptAt\_\allowbreak bar\_\allowbreak eq}; no adelic approximation
theorem is being used in this finite residue construction.

\paragraph{FC3: the finite chart fibers and their exact representatives.}
Let $\mathscr C_i$ be the supplied finite column-class set. A class
$c$ has an index $b(c)$ and a preferred point $r(c)$; let
$C(c)$ denote its coupling-reduction family. Define
\[
 \operatorname{Ch}_i=\operatorname{im}
       \bigl(c\mapsto(b(c),C(c))\bigr),\qquad
 \operatorname{chOf}(c)=(b(c),C(c)).                     \tag{FC.6}
\]
This is a finite image, not a new finite-class assumption. Choose a
representative $c_\beta$ of each image point. If $H_b$ is the original
product of finite reduction groups, put
$H_{b,C}=H_b\cap\operatorname{Stab}(C)$, component by component.

The fiber over $\beta=(b,C)$ is exactly the decorated-lattice orbit
set for $\mathrm U(E)\times H_{b,C}$. To prove injectivity, an element
of this group relating the decorated lattices also preserves $C$;
FC2 then says that it relates the whole points. They represent the same
original class. For surjectivity, lift any decorated lattice with the
prescribed $C$ by (FC.5), and take its preferred representative. Preferred
representatives normalize the coupling coset by choosing one value from
its orbit. Starting with the same $C$ therefore gives the same normalized
value: the two sets of attainable values are identical. The moving
element consequently lies in $H_{b,C}$. This proves both directions,
including the normalization step which would be missing from an arbitrary
partition into fibers.

For the passage to adelic charts we retain one precise geometric input:
each actual decorated lattice has a realizing unitary adelic point.
AD1 constructs the diagonal base point $k_u$ with the lattice and
last-row residues specified in \texttt{kStd\_\allowbreak spec}. Choose the actual source representative
$y_c=y_{\rm Std}(u,i,c)$ so that $y_ck_u$ realizes $r(c)$.
Pulling a vector back by $y_c^{-1}$ gives the chart base lattice
$L_\pi=\mathcal O^{n-1}\oplus\pi\mathcal O$. Equality of the
decorated points says exactly that a quotient matrix preserves this
lattice and has reduction in $H_{b,C}$: the lattice equality checks
integrality of the matrix and its inverse, and equality of the finite
decorations checks the reduction condition. Rational changes of frame
act on the left. Thus the fiber bijection just proved is also the
adelic double-quotient bijection required of a \texttt{LevelChart}.
The supplied injective stabilizer map has exactly the rational elements
in $y_cUy_c^{-1}$ as its image, by the same equality-of-points criterion.

The chart character is the restriction of the original finite character,
composed with reduction. More explicitly, right multiplication by an
adelic $x\in U$ acts on the decoration by the inverse of its reduction.
The definition of the level character contains this same inverse, so
the two inversions cancel. Its value is precisely the original column
character of the moving element. This proves the character clause of
\texttt{ReadsK}, in addition to its lattice, level and stabilizer clauses.
The matrices remain $y_{\rm Std}$, not representatives chosen anew
inside each chart. At $i=0,m$ the column index is unique and the coupling
set is empty; the nonempty image (FC.6) is therefore a singleton.
This gives the exact endpoint \texttt{Unique} assertions.

\paragraph{FC4: normalize the residue functional without changing the
level.}
The following local calculation supplies the substantial part of
\texttt{charts\_\allowbreak funKilled}. Let $A$ be the integral algebra
$\mathcal O_E\otimes_{\mathcal O_F}\mathcal O_{F_v}$ inside
$E\otimes_FF_v$, with its involution; write $q=\bar\pi$.
At a selected level prime, (FC.1) and FC1 identify both $A/\pi A$
and $A/qA$ with fields of order $p$. Put
\[
 L=A^{n-1}\oplus\pi A,\quad
 L^\#=A^{n-1}\oplus q^{-1}A,\quad
 L^\#/L\simeq A/pA,\quad [x]\longmapsto[qx_n].         \tag{FC.7}
\]
The middle identity is SF3, after base change. Suppose $g$ is unitary
and preserves $L$. It also preserves $L^\#$. Applying $g$ to
$\pi e_j\in L$ shows that its last row is integral; applying it to
$e_j\in L^\#$ shows that all other rows are integral. Moreover
\[
 g_{nj}\in\pi A\ (j<n),\qquad
 g_{jn}\in qA\ (j<n).                                 \tag{FC.8}
\]
On the quotient (FC.7), $g$ therefore acts by multiplication by
$d=g_{nn}$. The $(n,n)$ entry of $g^*g=1$ gives
\[
 d\bar d\equiv1\pmod{pA}.                             \tag{FC.9}
\]

Let $f:L^\#\to M$ be any nonzero $A$-linear functional vanishing on
$L$, where the additive group of $M$ has prime order $p$. Multiplication
by $\pi$ on $M$ is either zero or onto. In the second case
$qM=0$, since $q\pi=p$ and $pM=0$. Thus $f$ factors nontrivially
through either $A/\pi A$ or $A/qA$. The resulting map to $M$ is
an isomorphism: both additive groups have prime order. It follows that
$f(gx)=f(x)$ for every $x$ exactly when $d\equiv1$ modulo the
chosen one of these ideals. By (FC.9) and conjugation, the two conditions
are equivalent. In particular they are also equivalent to $g$ acting
trivially on all of $L^\#/L$.

Replace $f$ by the canonical quotient
\[
 f_0:L^\#\longrightarrow A/\pi A,\qquad
 f_0(x)=qx_n\bmod\pi A.                               \tag{FC.10}
\]
It is nonzero, kills $L$, has a value group of order $p$, and
$f_0(\pi x)=0$. The intersection of its stabilizer with the unitary
stabilizer of $L$ is exactly the old intersection. Therefore one can
replace only the chart's value module and functional: its base lattice,
every local level group, global level, character, adelic representatives
and stabilizer maps stay literally unchanged. Its \texttt{ReadsK}
proof stays valid. Doing this independently at each level prime and
each chart yields the actual \texttt{FunKilled} property. For $n=1$
the off-diagonal statements are empty and the same proof applies.

\paragraph{FC5: the quotient projections are constructed, not assumed.}
Let $\mathfrak p\subset\mathcal O$ be maximal and let two lattices
$A_0,M_0\subset E^n$ satisfy $(A_0)_{\mathfrak p}=(M_0)_{\mathfrak p}$.
For $x\in A_0$, choose $s\notin\mathfrak p$ with $sx\in M_0$ and set
\[
 \operatorname{pr}(x)=\bar s^{-1}[sx]
                 \in M_0/\mathfrak pM_0.              \tag{FC.11}
\]
Using a common denominator proves independence of $s$, additivity and
$\operatorname{pr}(ax)=\bar a\operatorname{pr}(x)$.
If $x\in M_0$, this is its usual residue. Given $[y]\in M_0/\mathfrak pM_0$,
choose $s\notin\mathfrak p$ with $sy\in A_0$ and $t$ with
$ts\equiv1\pmod{\mathfrak p}$; then $tsy\in A_0$ maps to $[y]$.
Thus $\operatorname{pr}$ is onto. Its kernel is exactly
$\mathfrak pA_0$: if its value is zero, clearing a second denominator
gives $rx\in\mathfrak pA_0$ for some $r\notin\mathfrak p$;
choose $a$ with $ar\equiv1\pmod{\mathfrak p}$ and write
$x=arx+(1-ar)x$.

For the actual decoration, its defining local laws give
$(L^\#)_{\mathfrak p}=N_{\mathfrak p}$ and
$(L+\mathfrak pL^\#)_{\mathfrak p}=(\operatorname{dn}N)_{\mathfrak p}$.
Apply (FC.11) to these two pairs. They are the source's
\texttt{pr} and \texttt{prk}, with the required surjectivity and kernels.
The decoration functional satisfies
\[
 f_{\mathfrak p}(z)=0\quad\Longleftrightarrow\quad
 z\in L+\mathfrak pL^\#\qquad(z\in L^\#).              \tag{FC.12}
\]
For the nontrivial implication, $f_{\mathfrak p}(z)=0$ puts $z$ in
$L_{\mathfrak p}$; clear a denominator $s$ and choose
$ts+r=1$, $r\in\mathfrak p$, to write
$z=t(sz)+rz\in L+\mathfrak pL^\#$. The converse follows by
linearity. Finally, if $f_{\mathfrak p}(z)=f_{\mathfrak p}(e)=1$,
then $z-e\in L_{\mathfrak p}$ and
$pz-pe\in\mathfrak p(\operatorname{dn}N)_{\mathfrak p}$.
Their projections agree. This is the special \texttt{te} compatibility,
and completes \texttt{exists\_\allowbreak level\_\allowbreak projections} with the literal
\texttt{IdL} identities on global elements.

\paragraph{FC6: the residue comparison has an explicit formula.}
Let $f:L^\#\to M$ now be a nonzero normalized functional, so
$f(\pi x)=0$. Its factorization in FC4 gives a unique additive
isomorphism $e:M\simeq k_P$ satisfying
\[
 e(f(x))=\rho_P(\bar\pi x_n).                          \tag{FC.13}
\]
Indeed, the map $x\mapsto\rho_P(\bar\pi x_n)$ is onto, and its
kernel is killed by $f$: subtract
$\pi\bar\pi^{-1}ae_n$ when $\bar\pi x_n=\pi a$; the
remainder lies in $L$. The induced nonzero map $k_P\to M$ is an
isomorphism between groups of order $p$.

This normalization is compatible with the actual decorated-class
functional, not merely with an abstract residue field. If $y_ck_u$
realizes the point, its inverse is $k_u^{-1}y_c^{-1}$. The defining
last-row residue of $k_u^{-1}$ is
$(0,\ldots,0,\bar\pi)$. Reduction of its product with
$y_c^{-1}z$ therefore gives
\[
 f_P(sz)=\bar s\,\rho_P\!\left(
        \bar\pi (y_{c,v}^{-1}z)_n\right)               \tag{FC.14}
\]
whenever $sz$ is in the original decorated lattice $N$.
All the products in this reduction are integral at $P$: the first
$n-1$ coordinates of $y_{c,v}^{-1}z$ are integral, and its last is
integral at $P$ because $\bar\pi$ is a unit there. This checks the
integrality needed when multiplying residue identities.

Choose $s\notin P$ with $sz\in N$, as permitted by FC5.
The defining rule (FC.11), (FC.13) and (FC.14) then agree after
multiplication by the nonzero scalar $\bar s$; cancel it.
This proves \texttt{FvTieG}. For a rational base-coordinate vector
$x$, it also proves the precise terminal formula
$e(f(x))=\operatorname{res}(r)$ when $r=\bar\pi x_n\in\mathcal O$.
Combined with $h(\pi e_n,x)=\bar\pi x_n$, this is the
\texttt{TilBaseFv} calculation once the terminal vector and its
incidence data have been supplied.

\paragraph{FC7: the metric comparison is a finite polynomial calculation.}
For polynomials in all frame entries write
$P=\sum_\alpha p_\alpha X^\alpha$ and define the normalized Fock form
\[
 N_\lambda(P,Q)=c_\lambda^{-1}
       \sum_\alpha\alpha!p_\alpha\overline{q_\alpha},
 \qquad c_\lambda>0,                                  \tag{FC.15}
\]
where $c_\lambda$ is the product of the Fock squared norms (the
self-pairings) of the source's highest polynomials. This convention is linear in $P$ and conjugate
linear in $Q$. It is positive definite. For completeness, the unnormalized
sum equals the Gaussian integral
$\pi^{-d}\int_{\mathbb C^d}P(z)\overline{Q(z)}e^{-|z|^2}\,dz$:
angular integration kills unequal exponents and radial integration gives
$\alpha!$. A unitary change of variables preserves this integral and
therefore the form. Blockwise left multiplication of the frame matrices
by unitary matrices is such a change of variables. Adding a common
integer to every entry of a highest weight changes neither its
normalized weight nor $c_\lambda$, hence does not change (FC.15).

In the actual standard coordinates, an orthonormal frame $o$ satisfies
$o_\tau^*o_\tau=I$, so $|\det o_\tau|=1$. Inserting the
determinant-twisted polynomial definitions gives
\[
 \operatorname{ipV}(f_P,f_Q)=N_\lambda(Q,P),
 \qquad f_P(z)=\operatorname{fnOf}(o^{-1},P)(z).
 \tag{FC.16}
\]
The order $Q,P$ is essential: $\operatorname{ipV}$ is conjugate linear
in its first argument. No reversal of the Hermitian convention is
being made.

The global Gram matrix is defined by the finite class sum with weights
$a_c=|\operatorname{Stab}(c)|^{-1}>0$. Expanding a vector in its
chosen basis and distributing this finite sum proves, at an embedding
$s$ of the coefficient field,
\[
 \operatorname{maj}_s(x,x)
   =\sum_c a_c\,N_\lambda(sP_c,sP_c)
   =\sum_c a_c\prod_\tau|\det o_\tau|^{2\delta_\tau}
       N_\lambda\bigl(P_c^s(oX),P_c^s(oX)\bigr).
 \tag{FC.17}
\]
Taking real parts is the source's \texttt{NormTieG}. The following
descent argument justifies the embedded Gram expansion even when the
coefficient field $L$ is larger than a CM field. Define
\[
 L^{\mathrm c}=\{a\in L:
    s(a^*)=\overline{s(a)}\text{ for every }s:L\hookrightarrow\mathbb C\}.
 \tag{FC.17a}
\]
It is a subfield, since the displayed equality is preserved by field
operations. It contains the subfield
$L_0=\mathbb Q(\tau(E),\chi(g):\tau,g)$ generated by the embedded
CM-field images and the character values in the stabilizer actions.
For the first generators this follows from the supplied
\texttt{FOk.emb\_\allowbreak star} and CM conjugation laws. For a root of unity
$\zeta$, its distinguished embedding has modulus one, and
\texttt{FOk.sig\_\allowbreak star} identifies star there with complex conjugation;
injectivity
of that embedding gives $\zeta^*=\zeta^{-1}$. Every other embedding
also maps it to a root of unity, hence
$s(\zeta^*)=s(\zeta)^{-1}=\overline{s(\zeta)}$.
These are the generators actually used by the polynomial action.

The integer Weyl-polynomial basis identifies the stabilizer equations
with linear equations over $L^{\mathrm c}$. Their invariant space has
an $L$-basis with coordinates in $L^{\mathrm c}$. One direct proof,
valid simultaneously for all equations, chooses an
$L^{\mathrm c}$-basis $a_j$ of $L$ and applies its coordinate
projections to every coefficient of an invariant vector. Each projection
commutes with every action matrix, so gives another invariant vector
with coefficients in $L^{\mathrm c}$. The original vector is
$\sum_j a_j$ times these projected vectors. The latter therefore span
over $L$, and a basis can be extracted from them.

Crucially, this is the basis used by the actual source: the definition
\texttt{CSD.basis} selects a basis with these coordinates whenever
\texttt{HasRatBasis} holds, as just proved. It is not an arbitrary
$L$-basis followed by an unwarranted interchange of embedding and
involution. For two vectors of this selected basis, every polynomial
coefficient belongs to $L^{\mathrm c}$; thus (FC.17a) proves that
their coefficient pairing, after embedding, is exactly (FC.15) with
the Hermitian order reversed as in (FC.16). Expand a general vector
$x=\sum_j a_jb_j$ only after those basis-pairing identities are proved.
Both sides are then the same complex quadratic form in the scalars
$s(a_j)$, with the usual complex conjugates. No equality
$s(a_j^*)=\overline{s(a_j)}$ is required for arbitrary coefficients
$a_j\in L$. This also explains the statement in any other coordinate
basis: transport the already obtained complex quadratic form by the
complex change-of-basis matrix.

The combined matrix of all basis values is injective. It has a left
inverse over $L$, whose embedded matrix remains a left inverse; hence
a nonzero complex coefficient vector has a nonzero class value.
The positive weighted Fock sum therefore makes the embedded Gram
matrix positive definite. The source's standard Hermitian frame
realizes precisely this Gram matrix. Polynomial injectivity (CV1)
identifies any proposed family $P_c$ in the norm statement with the
actual coefficient polynomials. Thus (FC.16)--(FC.17) prove the literal
\texttt{ipV\_\allowbreak norm\_\allowbreak tie}, rather than assuming a metric comparison.
The same argument in dimension $N_m-1$, with the terminal constant
weight shift, supplies \texttt{ipP} and \texttt{NormPTie} for the
terminal value construction of CV10. It proves no independent spectral
or metric estimate elsewhere in the reconstruction.

\paragraph{FC8: the period map is coefficient conjugation.}
Let $\mathcal C P$ be the frame function obtained by conjugating the
coefficients of $P$ and retaining the prescribed determinant twist.
It is conjugate linear and injective by CV1; it maps onto the prescribed
complex polynomial value space, because that space is defined over the
real integer highest-polynomial generators. For a rational unitary
matrix $\gamma$ and a character value $\chi$ of finite order,
$\gamma^{-1}=\gamma^*$ and $\overline{\chi^{-1}}=\chi$.
Substitution into the polynomial action gives the exact intertwining
identity
\[
 \mathcal C\bigl(\chi^{-1}\rho(\gamma)P\bigr)
          =\operatorname{tr}_{\gamma,\chi}(\mathcal C P).
 \tag{FC.18}
\]
Here $\mathcal C$ conjugates coefficients; it does not replace the
frame variable by its complex conjugate.

Evaluate an equivariant polynomial section on each preferred terminal
class representative and apply $\mathcal C$; call the resulting map
$j_M$. The source's \texttt{PerIn.H} is the whole group. Values on
preferred representatives therefore determine a section, proving
injectivity. Conversely, lift each prescribed class value uniquely by
$\mathcal C$ and extend by $\chi^{-1}\rho(\gamma)$ along its orbit.
Two possible moving elements differ by a stabilizer. Equation (FC.18),
the invariance of the class value and injectivity of $\mathcal C$
show that the lift is fixed by that stabilizer. The extension is
well-defined and equivariant. Thus $j_M$ maps onto exactly
the source's \texttt{ptC} space.

Conjugating the second polynomial in (FC.15) turns its Hermitian
coefficient pairing into the symmetric bilinear coefficient pairing.
Term by term in the weighted class sum this proves
\[
 \operatorname{period}(w,F)=\operatorname{ipP}(j_MF,w).
 \tag{FC.19}
\]
For nonzero $z\in\texttt{ptC}$, at least one class polynomial is
nonzero. Its Fock norm is positive, all class weights are positive,
and every other summand is nonnegative. Hence
$\operatorname{Re}\operatorname{ipP}(z,z)>0$. These statements prove
all five assertions of the literal \texttt{period\_\allowbreak conj} supplier
under the supplied column/period laws; they do not assert existence of
those laws for a new arithmetic host.

\paragraph{FC9: convert the chart levels and classes to Hermitian
carriers.}
This construction uses the auxiliary-prime descriptions
$\texttt{GLevH}$ of the same charts, constructed in FC14 below. At an auxiliary prime, a matrix
congruent to the identity lies in the chart level because its reduction
is the identity of its finite subgroup. Away from the auxiliary and
level primes the chart already has its stated full lattice stabilizer.
At a level prime, FC4 proves the remaining Hermitian condition exactly:
\[
 U_v=\{g\in\mathrm U(h):gL=L,\quad
                    (g-1)L^\#\subset L\}.            \tag{FC.20}
\]
In fact the old definition is the intersection with the stabilizer of
the nonzero prime-order functional, and FC4 shows that this is equivalent
to acting trivially on the whole discriminant module. This also explains
why one cannot simply replace a functional by an arbitrary zero map.

Use the coordinate basis of SF2 to carry the lattice, local groups and
global group back to the actual Hermitian space. Transport along the
single rank-index bijection takes addition, scalar multiplication,
duals, localizations and matrix conjugation to themselves. For each
class, use the very same $y_c$ and stabilizer homomorphism in these
coordinates. Their quotient bijection, lattice-transport condition,
injectivity and range formula are unchanged. This proves
\texttt{classTie\_\allowbreak of\_\allowbreak chart}, including every equality in
\texttt{SameCar}, rather than just an isomorphism of class sets.

At step zero the same transport produces \texttt{ColLevel}. Its
type-prime splitting and frames are the supplied $G$ and the same
\texttt{glev0.fr}. The additional claim that a rational stabilizer
preserves its class lattice follows as follows. Conjugating it by $y_c$
puts it in $U$, so it preserves every completed class lattice.
A full $\mathcal O$-lattice is the intersection of its finite localizations
inside $E^n$, by the denominator argument of LP1. It therefore preserves
the global lattice; apply the same argument to its inverse.

For $1\le i<m$, apply this construction to each chart fiber. At
$i=m$, the chart type is a singleton by FC3, so restriction to its
fiber and the inverse inclusion are exact inverse class maps. There
is no empty endpoint fiber. Use the residue equivalence (FC.3) and the
same place bijection for the remaining fields. These choices construct
the entire \texttt{CarH} output of \texttt{hermColsCh\_\allowbreak of\_\allowbreak forms},
conditional on the given standard forms and their level descriptions.
No separate Hermitian carrier existence assumption is needed after
those forms are available.

\paragraph{FC10: the step-zero value module is killed by the actual
uniformizer.}
The previous normalization proves this for the charts constructed here.
There is also a useful intrinsic proof for any supplied standard forms,
which is the source's \texttt{pi\_\allowbreak kills}. Choose a class, a vector in
its nonzero residue-functional image, and a lift $z$ using the
surjectivity of its projection. By \texttt{FvTie}, the corresponding
value $m=f(y_c^{-1}z)$ in the chart value group is nonzero.
Multiplication of $z$ by $\pi$ multiplies its projected value by
$\operatorname{res}(\pi)=0$. Linearity and \texttt{FvTie} therefore
give $\pi m=0$. An additive group of prime order is generated by
any nonzero element, so $\pi M=0$.

If a previously chosen \texttt{ColLevel} has the correct carriers but
different residue-module fields, transport the forms' functional along
the exact rank reindexing and use its value module. Nonzeroness,
$p$-torsion, order $p$, vanishing on the base lattice and its stabilizer
formula transport with it. This changes only those functional fields;
the lattice, local and global groups, class representatives, stabilizer
maps, splitting and frames remain identical. It proves the literal
\texttt{stepZeroHermitianLevel\_\allowbreak with\_\allowbreak formsFunctional} and
\texttt{stepZeroHermitianLevel\_\allowbreak uniformiser\_\allowbreak kills\_\allowbreak values} assertions.

\paragraph{FC11: what the naturality calculation does once a supported
move exists.}
This paragraph proves the algebraic part of the move comparison, with
its geometric premise exposed. Suppose an actual up/down/diamond move
keeps the chart label, has moving column-group element $g=(\gamma,h)$,
and has an adelic realization $y_c\delta k$ with $\delta$ equal to
the identity away from the specified level place. Let $c'$ be the
new preferred class. By \texttt{ReadsK},
\[
 \kappa=y_{c'}^{-1}\gamma^{-1}y_c\delta\in U,\qquad
 y_c\delta=\gamma y_{c'}\kappa,
 \qquad \chi_U(\kappa)=\chi_{\rm col}(g).              \tag{FC.21}
\]
The reason for membership is exact: $\gamma^{-1}$ moves the realized
point to the preferred point by the finite component $h$, which
preserves the chart coupling data. The level-reading equivalence of
FC3 applies, and its character-reading equality gives the last formula.
For frames with $z=\gamma z'$, the transported function is
$\chi_{\rm col}(g)f(\gamma^{-1}z)=\chi_U(\kappa)f(z')$.
This proves the character and frame parts of \texttt{upNatTie} and
\texttt{dnNatTieG\_\allowbreak cols}; the inverse $\kappa^{-1}$ in the source's
alternative factorization introduces exactly the displayed inverse
character, with no new scalar.

The elementary local-basis correction used in an up move is also exact.
If $a$ is a uniformizer at $P$ and $x-z\in PA$, then
$a^{-1}(x-z)\in A_P$, whence
\[
 (A+\mathcal Oa^{-1}x)_P=(A+\mathcal Oa^{-1}z)_P.
 \tag{FC.22}
\]
For a down-move functional, an identity known on the original global
lattice extends to its localization: choose $s\notin P$ making the
vector integral, use the residue-linearity of both sides, and cancel
$\bar s\ne0$. This is the cancellation in
\texttt{dn\_\allowbreak level\_\allowbreak functional}. These arguments do not construct
$\delta$, prove that the source's actual neighbor has the required
local lattice, or establish that a move preserves the coupling label.
Those are distinct geometric assertions retained below.

\paragraph{FC12: a common finite orbit carrier for the penultimate
charts.}
For a group $J$ and a finite family of finite-index subgroups
$U_\beta\le J$, set
\[
 Y=\coprod_\beta J/U_\beta,\qquad y_\beta=(\beta,U_\beta).
 \tag{FC.23}
\]
This is a finite $J$-set. Every point belongs to exactly one orbit
$J y_\beta$, and the stabilizer of $y_\beta$ is exactly $U_\beta$.
If $\chi_\beta$ is the restriction of a common character $\chi_J$,
its stabilizer character is also the required one. These assertions
follow directly by multiplying cosets and retaining their disjoint-union
label. They prove the actual \texttt{exists\_\allowbreak orbits} part of
\texttt{stepBefore\_\allowbreak oneLevel}.

For the actual penultimate charts, the supplied index law
\texttt{last\_\allowbreak uniq} makes the finite column-group index independent
of the class. Choose a point whose coupling coset contains zero.
Such a point exists by taking the zero coset of the prescribed coupling
lattice. Membership of zero is preserved by every linear column action,
so its preferred representative still has zero coupling reduction;
let $\beta_0$ be its chart. All zero-containing coupling cosets have
the same reduction family: translating their defining lattice by zero
leaves that lattice, and a reduction frame sends it to the prescribed
standard lattice. In the reference case these reductions at the auxiliary
primes are the singleton zero.

Here is the exact group calculation. Let $J_0$ be the unitary integral
adelic group whose reduction fixes the last row at every selected level
prime, and let $r$ be its product of auxiliary-prime reduction maps.
The AD1 base point normalizes $J_0$: if
$e_n^{\mathsf T}k^{-1}=b e_n^{\mathsf T}$ after reduction, then
$b$ is a unit (multiply by $k$), and the equation
$e_n^{\mathsf T}x=e_n^{\mathsf T}$ is equivalent to the same
equation for $k^{-1}xk$. At auxiliary primes $k=1$, so conjugation
does not change $r(x)$. The point-reading argument of FC3 consequently
gives
\[
 U_\beta=\{x\in J_0:r(x)^{-1}\in H,
                   \ r(x)^{-1}C_\beta=C_\beta\},
 \qquad \chi_\beta(x)=\theta(r(x)^{-1}),             \tag{FC.23a}
\]
where $H$ and $\theta$ are the same finite column group and character
for all $\beta$. The zero reduction $C_{\beta_0}$ is fixed by $H$.
Thus choosing $J$ to be the already constructed chart level at
$\beta_0$ gives $U_\beta\subset J.U$, equality of all base
lattices, and restriction of the character of $J$ to $U_\beta$.
Moreover $\ker r\subset U_\beta$ and $r$ has finite image, so
each $U_\beta$ has finite index in $J.U$. Applying (FC.23) now
constructs the actual finite orbit carrier with its exact stabilizers.

This proves the common level, inclusions, character restrictions and
orbit part of \texttt{stepBefore\_\allowbreak oneLevel}. What remains there is
the specified identification of its reduction group with
$\texttt{HSym}$ and its character with the displayed
$\texttt{lamSym}$ formula. The generic reduction-image construction
of FC14 alone does not prove that particular matrix shape and character
formula.

\paragraph{FC13: construct the one common auxiliary-prime package.}
Write $q_j=g_j$ for the two distinct supplied auxiliary rational primes,
and $a_j=\texttt{piG}(j)$, with $a_j\bar a_j=q_j$.
The supplied complete splitting at these primes gives, for each
$v\mid q_j$, a splitting
$E\otimes_FF_v\simeq F_v\times F_v$ which carries conjugation
to exchange of factors and scalars to the diagonal. Its residue field
is the prime field $\mathbb F_{q_j}$. Orient the splitting so that the
first component of $\bar a_j$ has valuation less than one.
This is always possible: the two components of $a_j$ have product
$q_j$, whose valuation is less than one. If the required component is
not small, the other is; exchange the two factors. Both displayed
splitting identities survive that exchange. The equality
$a_j\bar a_j=q_j$ and the chosen $a_j$ themselves are unchanged.

Choose one selected auxiliary prime $w_j$ of $E$ above $q_j$; its
existence follows from the supplied positive cardinality of the selected
set. The rings $k_v$ and $\mathcal O/w_j$ both have order $q_j$.
The unique unital map from $\mathbb F_{q_j}$ to either is an
isomorphism: its kernel is zero and domain and codomain have the same
finite cardinality. Let $e_v:\mathbb F_{q_j}\simeq k_v$ and
$e_w:\mathbb F_{q_j}\simeq\mathcal O/w$ denote these maps.
For the coefficient-field embedding $\sigma_0:L\hookrightarrow\mathbb C$,
define
\[
 \psi_{j,v}(x)=\sigma_0\!\left(
       \psi^{\rm col}_{w_j}(e_{w_j}e_v^{-1}x)\right).
 \tag{FC.24}
\]
It is an additive character. If it were trivial, surjectivity of the
field isomorphisms and injectivity of $\sigma_0$ would make the supplied
column character trivial. Thus it is nontrivial. The column-character
law says that $\psi^{\rm col}_w$ has the same value on every integer
at every selected $w\mid q_j$. Since every element of these prime
fields is an integer residue, (FC.24) agrees with the required character
at every such $w$, not only at $w_j$.

The supplied multiplicative-character law gives an injective homomorphism
$\omega_j:\mathbb F_{q_j}^{\times}\to L^{\times}$ for which
the $s$th column character is $\omega_j^{2(s+1)}$.
Extend $\sigma_0\omega_j$ to a multiplicative character on the field,
and transport its powers by $e_v^{-1}$:
\[
 \chi_{j,v,b}(x)
       =\bigl(\sigma_0\omega_j\bigr)^{2(b+1)}(e_v^{-1}x).
 \tag{FC.25}
\]
This character has the prescribed values on every nonzero integer
residue. Injectivity gives the primitive character
$\sigma_0\omega_j$ order $q_j-1$, since
$\mathbb F_{q_j}^{\times}$ is cyclic of that order. If the powers
in (FC.25) for $a,b<m$ were equal, transporting them to $w_j$ would
make the corresponding supplied column characters equal. Their
injectivity law gives $a=b$. Equations (FC.24)--(FC.25) also show
agreement across all $v\mid q_j$ on integer residues and exhibit the
single primitive multiplicative character required by
\texttt{CharacterCompatible}.

The oriented splittings, the original $a_j$, and these characters are
the complete $G=\texttt{GLoc}$ object. They prove all three fields of
\texttt{IdG} and the complete \texttt{CharacterCompatible} assertion
in \texttt{exists\_\allowbreak gLoc\_\allowbreak idG\_\allowbreak chiTied}. Choose this $G$ once,
independently of the word $u$. Since $G.\texttt{piG}=X.\texttt{piG}$,
the required membership of its uniformizers in the designated subfield
is exactly the supplied $X.\texttt{piG\_\allowbreak K}$ clause. Complete splitting,
the original character laws and that subfield clause remain part of
the host data, not consequences of this transport argument.

\paragraph{FC14: construct the chart reduction groups and their
characters.}
Fix any chart $S$ from FC3 and any $v\mid q_j$.
Because $q_j\ne p$, the element $p$ is a unit at $v$. Its base
lattice satisfies
$p\mathcal O^n\subset L_\pi\subset\mathcal O^n$, so the
completed lattice is the standard integral lattice. This is the explicit
denominator calculation of SF5. The first split projection therefore has
the identity matrix as a frame, including the required existential
lifting of every integral first-projection vector.

Let $V_v$ be the unitary stabilizer of this completed lattice and let
$r_v:V_v\to\mathrm{GL}_n(k_v)$ be reduction of the first split
matrix. Its second split matrix is the inverse transpose of the first,
because the Hermitian Gram matrix is $I$. Hence $r_v(x)=I$ implies
that both split matrices of $x$ and $x^{-1}$ are congruent to $I$.
Embed $x$ at the single place $v$, with identity elsewhere. It preserves
the completed lattice and every reduction frame of the decorated point:
at $v$ this follows from the congruence, and elsewhere it is the identity.
The base point $k_u$ is also the identity at $v$.
Thus right multiplication by this embedded element leaves the same
decorated point and its coupling reductions. The exact \texttt{ReadsK}
level and character clauses give
\[
 \ker r_v\subset U_v,\qquad
 \chi_U(\iota_v x)=1\quad(x\in\ker r_v),             \tag{FC.26}
\]
where $\iota_v$ denotes the supported embedding.

Define $H_v=r_v(U_v)$ and
$\lambda_v(r_vx)=\chi_U(\iota_vx)^{-1}$ for $x\in U_v$.
Equation (FC.26) makes this independent of the lift and multiplicative.
Moreover
\[
 U_v=r_v^{-1}(H_v)\quad\hbox{inside }V_v.             \tag{FC.27}
\]
Indeed, if $r_vx=r_vy$ with $y\in U_v$, then
$y^{-1}x\in\ker r_v\subset U_v$.
Only finitely many places divide the two distinct nonzero $q_j$.
For $x\in U$, remove its components at these places using the
$\iota_v$. The remaining element is the identity at every bad place;
the given \texttt{StepLevel} character is trivial on that subgroup.
Consequently
\[
 \chi_U(x)=\prod_{j=0}^1\prod_{v\mid q_j}
                  \lambda_v(r_vx_v)^{-1}.            \tag{FC.28}
\]
Each $H_v$ is finite, so each $\lambda_v$ has finite image and all
its values have absolute value one. Equations (FC.27)--(FC.28), with
the identity frame, supply every field of \texttt{GLevH}.
This proves the actual \texttt{chart\_\allowbreak gLevH} supplier from the
constructed charts. No further auxiliary-prime description is assumed
for those charts.

\paragraph{FC15: the explicit step-zero level and character.}
At $i=0$, $n=N_0=k$, the coupling set is empty, and the supplied
index law identifies the original finite column subgroup at each
selected auxiliary prime with the upper unitriangular group
$\mathrm U_n(k_v)$. This identification is part of
$\texttt{idx\_\allowbreak ok.zero\_\allowbreak G}$; its matrix content is simply
$A_{ab}=0$ for $b<a$ and $A_{aa}=1$. The source's
$\texttt{InJ}(k,t,N_0,\mathrm{false})$ has exactly this content:
all its block labels are singleton indices, while its extra fixed-column
and last-row conditions have no applicable index.

The oriented first split factor in FC13 corresponds to the unique
selected prime $w\mid\bar a_j$ above $v$. Its residue map
$\rho_w:\mathcal O/w\simeq k_v$ is the map induced by that
first projection. It carries every integer to itself. Reading the
same point after a matrix supported at $v$ changes its reduction frame
by the inverse matrix at $w$ and leaves all other selected reductions
unchanged. Since $\mathrm U_n$ is a group, the inverse has the
required shape exactly when the matrix itself does. The absence of
coupling and twist-prime data at step zero, together with \texttt{ReadsK},
therefore proves
\[
 x\in U_v\quad\Longleftrightarrow\quad
 r_v(x)\in\mathrm U_n(k_v)\qquad(x\in V_v).          \tag{FC.29}
\]
This proves the specified group identity, rather than merely choosing
the image subgroup as in FC14.

For $A\in\mathrm U_n$, put
$s(A)=\sum_{a=1}^{n-1}A_{a,a+1}$. Matrix multiplication gives
$s(AB)=s(A)+s(B)$: in the $(a,a+1)$ entry only the two diagonal
terms survive. Thus $s(A^{-1})=-s(A)$.
At step zero the column character is
$\prod_w\psi_w^{\rm col}(s(A_w))$; there are no determinant
block factors or twist-prime factors. The chart reads
$A_w=r_w(x)^{-1}$. Applying (FC.24), the actual residue maps
$\rho_w$, and the last identity gives
\[
 \chi_U(x)=\prod_{j=0}^1\prod_{v\mid q_j}
               \psi_{j,v}\bigl(s(r_vx_v)\bigr)^{-1}. \tag{FC.30}
\]
This is exactly the source's \texttt{lamJ} formula with
$(a,nb,skip,fc,zl)=(k,0,k,N_0,\mathrm{false})$.
The finite-product reindexing is legitimate: restriction is a bijection
on each selected auxiliary-prime set, by the same conjugate-prime and
coprimality argument as FC1, and the two sets are disjoint because
$q_0\ne q_1$. The additive-character values have finite order and
absolute value one. The identity frame, (FC.29), (FC.30), and this norm
calculation prove all fields of \texttt{glev\_\allowbreak zero}, including
\texttt{fr0}$=I$.

\paragraph{FC16: the exact joins to the comparison hypotheses.}
The level-zero coefficient-field condition has a short proof which is
stronger than requiring the particular construction (FC.24).
For any additive character $\psi:k_v\to\mathbb C$ at $v\mid q_j$,
\[
 \psi(x)^{q_j}=\psi(q_jx)=\psi(0)=1.                 \tag{FC.31}
\]
The supplied \texttt{Base.roots\_\allowbreak Ks} contains a primitive
$q_j(q_j-1)$st root $\zeta$ in $L$, with its specified complex
embedding. Every complex $q_j$th root of unity is a power of
$\sigma_0(\zeta)$: it is a root of
$X^{q_j(q_j-1)}-1$, whose distinct roots are precisely those powers.
Hence $\psi(x)=\sigma_0(\zeta^a)$ for some integer $a$.
Applying this to any $G.\texttt{psi0}$ proves the literal
\texttt{LevelZeroCharacterRational} condition. This proves the actual
\texttt{additiveCharacter\_\allowbreak typePrimes\_\allowbreak values\_\allowbreak coefficientField}
supplier without using its \texttt{IdG} argument, and without another
global forms assumption.

The step-zero carrier construction in FC9 uses only the step-zero
chart, its lattice/stabilizer data and the level of FC15. Call this
level $L_0(u)$. Any completion of our partial forms data to
$T_u:\texttt{Forms0}$, keeping these specified fields, satisfies
\texttt{Joint0} with $L_0$ directly: its five assertions are equality
of the reindexed base lattice, local groups, global group, class matrices
and stabilizer matrices, exactly the \texttt{car0} identities in FC9.
They are unaffected by FC10's functional replacement.

This also removes a need to supply separate companion ideals for
these carriers. The FC3 base-lattice equality and FC9's exact reindexing
give $L_0(u)=L_{\pi_u}$ in the standard space. SF2 supplies
\texttt{StdTied}, and SF3 gives
\[
 L_0(u)+(\bar\pi_u)L_0(u)^\#=\mathcal O^{N_0}.
 \tag{FC.32}
\]
Thus the comparison ideals are the explicit $(\bar\pi_u)$ for these
same level-zero carriers. Their construction needs no terminal
\texttt{CarT} data. Once the remaining fields of \texttt{Forms0}
are supplied, VE9's \texttt{StdTied}, \texttt{Joint0},
\texttt{LevelZeroCharacterRational} and companion premises are all
discharged by these exact joins. FC19 supplies the full
\texttt{Forms0} directly, using NB7--NB8 for its naturality fields.

\paragraph{FC17: the distinguished-block lattice comparison.}
Fix a coefficient prime $w$ above $p$ which is seen, and a selected
level prime $P$. For a class decoration, put $A=L^\#$ and let
$K_P\subset A$ be the kernel of its $P$-residue functional.
The decorated-lattice laws and the denominator argument of FC5 give
\[
 (K_P)_P=(L+PA)_P,\qquad (K_P)_Q=A_Q
             \quad(Q\ne P\text{ selected}).          \tag{FC.33}
\]
For the second equality, choose $a\in P\setminus Q$.
Multiplication by $a$ kills the $P$-residue of every vector in $A$,
while $a$ becomes a unit at $Q$; this proves the nontrivial inclusion.
For the first, zero residue is exactly local membership in $L$,
and the kernel after clearing denominators is $L+PA$, as in (FC.12).

Choose rational $p$-local basis matrices $B,B'$ for $A$ and
$L+PA$ respectively. Such bases exist because both are full lattices;
for the second lattice use $L\subset L+PA\subset A$.
The rational $p$-local basis condition implies the prime-local basis
condition at every $Q\mid p$, by CV6. Indeed a rational integer
prime to $p$ cannot belong to $Q$, and a product of the finitely many
coordinate denominators clears the converse implication.
At a seen $w$ the supplied place correspondence is exact:
the place below $w$ through the block embedding labelled
$\tau_Q=\texttt{tauAt}(w,Q)$ is $Q$.
Thus (FC.33) gives a basis of $K_P$ in every block, namely
\[
 B_\tau=\begin{cases}B',&\tau=\tau_P,\\B,&\tau\ne\tau_P.
            \end{cases}                              \tag{FC.34}
\]
Apply the exact localized Weyl-lattice transport of CV7 to this family.
Membership in the $w$-localization of the polynomial lattice attached
to $K_P$ is equivalent to one integral-span polynomial evaluated with
$B_\tau^{-1}$ in each block. Formula (FC.34) is precisely the
piecewise inverse-frame formula in \texttt{BwTieG}. CV8 takes the
product of the classwise denominators to turn these conditions into
membership in the global $\texttt{Bw}$ localization. All frame maps,
embeddings and the determinant shift are the same identities as CV9.
The source quantifies over every pair $B,B'$; the forward implication
works for every such pair, and the reverse uses the existence just proved,
so it is not vacuous. This proves the full actual \texttt{bw\_\allowbreak tie}
supplier. These local matrices describe the class lattice and do not
replace the common global comparison frame of SF1--SF6. The assertion
is restricted to seen $w$ of residue characteristic $p$, as in the
source; it makes no claim about unseen coefficient primes.

\paragraph{FC18: the actual diamond move from its defining functional
equation.}
Let $s\in k_P^\times$, choose an integer $a$ representing $s$,
and let $\delta$ satisfy the source's \texttt{IsDia} condition at
$v=P\cap\mathcal O_F$. This means that $\delta$ is the identity
away from $v$, is unitary, preserves the base lattice, and satisfies
\[
             a f(\delta x)=f(x)\qquad(x\in L^\#).
 \tag{FC.35}
\]
Its inverse preserves the same lattices. FC4 therefore proves that
both matrices have integral entries. Substitute $\delta^{-1}x$
in (FC.35), and use FC6's exact functional identification, to obtain
\[
 \rho_P\bigl(\bar\pi(\delta^{-1}x)_n\bigr)
                =s\,\rho_P(\bar\pi x_n).
 \tag{FC.36}
\]
Every coordinate vector belongs to $L^\#$, and $\bar\pi$ is a
unit at $P$. Testing (FC.36) on them gives
\[
        (\delta^{-1})_{nj}\equiv s\,\delta_{nj}\pmod P.
 \tag{FC.37}
\]
Here the symbol on the right is the Kronecker delta.
Thus the source's integral-support-last-row conclusion follows directly
from the normalized functional; no choice of a special diamond matrix
is needed.

Suppose $y_ck_u$ realizes the preferred decorated point. AD1 gives
$k_u^{-1}$ last row $\bar\pi e_n^{\mathsf T}$ modulo $P$,
and $k_u$ preserves the standard integral adelic lattice.
For a vector in the original decorated lattice, $y_c^{-1}z$ is
therefore integral. Reducing the product
$k_u^{-1}\delta^{-1}y_c^{-1}z$ is legitimate. Equations (FC.37)
and the last-row formula multiply its old functional by $s$.
The underlying lattice is unchanged, because $\delta$ and its inverse
are integral. Every auxiliary reduction, every other selected level-prime
functional, and every twist-prime frame is unchanged, since its place
is outside $v$. No selected prime is conjugate to $P$.
Consequently $y_c\delta k_u$ realizes exactly the source's
\texttt{dia1} decoration.

This operation changes only the $P$-functional. In particular its
coupling reductions are unchanged. The supplied level law
\texttt{dia\_\allowbreak act} says that the operation commutes with the column
group action, so the sets of attainable coupling reductions are also
unchanged. The preferred-normalization argument of FC3 proves that the
new preferred class has the same chart label. Apply FC11 to the
realizer just constructed: it gives the exact adelic factorization,
character scalar and frame transport. This proves all clauses of the
actual \texttt{dilTied} supplier, under the literal supplied column
laws and the chart realizability input. It also supplies the dilation
field required when CV9's chart values are made into an endpoint
\texttt{ValTied} record.

\paragraph{FC19: direct construction of Forms0 and its matching comparison data.}
Fix the actual $D=I.\texttt{rec7}\,l\,Z$ and the supplied
\texttt{Ok}, \texttt{Pins}, \texttt{Laws} and \texttt{Later} data.
Retain the explicit initial realization hypothesis used in FC3 and
NB6: every valid framed decorated point needed there has a unitary
finite-adelic realizer in the full \texttt{IsPtOf} sense.
We now construct, for every word $u$, an actual
$T_u:\texttt{Forms0}\,D\,u$, level data $L_0$, and all of
VE9's global joins. The construction uses only step-zero form data.

\emph{First construct the form record.}
Choose the common $G$ of FC13 once for all words. At $i=0$, FC3
gives a unique chart label $\beta_u$ and
$\operatorname{chOf}(c)=\beta_u$ for every class $c$. FC4 normalizes
its functional without changing its level, class matrices or stabilizer
matrices. Define $\operatorname{cls0}_u$ by retaining this level and
setting
\[
 y_u(c)=K_u(\beta_u).y(c,\operatorname{chOf}(c)=\beta_u),\qquad
 s_u(c,\gamma)=K_u(\beta_u).\operatorname{stabI}
                      (c,\operatorname{chOf}(c)=\beta_u)(\gamma).
 \tag{FC.38}
\]
The map from the original class set to this chart fiber is a bijection
with inverse the first projection. Thus the class bijection, carried
lattice identity, and injective stabilizer with its exact image are
the same assertions as FC3. This is \texttt{clsOfChart} itself,
with no replacement representatives.

Take $\operatorname{atT}_u$ from FC1 and the frame equivalence $e_Y$
from SF2. The determinant witness in \texttt{ValTied.de} is $C=I$:
the Gram matrix is $I$, the integral-coordinate lattice is
$\mathcal O^{N_0}$, and the original determinant label equals
$\det(e_Y z)$ at every embedding. CV9 supplies the complex-form,
coefficient-field and localized integral-value identities.
Because the chart set is a singleton, a function satisfies its
chartwise adelic condition if and only if its restriction through
the bijection in (FC.38) does. Surjectivity of that bijection proves
both implications by evaluation on every class; no extension is lost.

The global integral field of \texttt{ValTied} also follows from
these localized assertions. For any $\mathcal O_L$-submodule $M$
of an $L$-vector space and any $x$ in every localization $M_w$,
put $\mathfrak b_x=\{a\in\mathcal O_L:ax\in M\}$.
If $x\notin M$, this ideal is proper and lies in a maximal ideal
$w$. Membership $x\in M_w$ supplies an $a\notin w$ with
$ax\in M$, contradicting $\mathfrak b_x\subseteq w$.
Every maximal ideal of $\mathcal O_L$ is a coefficient prime, so
\[
                         M=\bigcap_w M_w.           \tag{FC.39}
\]
Here localization membership supplies such a denominator by taking
the product of the finitely many denominators in an expression for
$x$ in the localized span. This argument does not require an
unproved freeness statement for $M$.
Apply it to the actual form lattice, using the full prime correspondence
and CV9 at every coefficient prime. FC7 supplies the exact inner-product
field and FC18 supplies both dilation fields. These data are precisely
\texttt{valOfG}, and give $\operatorname{val0}_u$ with unchanged $e_Y$.

The remaining fields of \texttt{Forms0} are now specified individually.
Its \texttt{gl} is the same $G$, and \texttt{gl\_\allowbreak K} follows from
$G.\texttt{piG}=X.\texttt{piG}$ and the supplied
$X.\texttt{piG\_\allowbreak K}$. Its \texttt{glev0} is the actual level of FC15,
with parameters
$(a,t,nb,skip,fc,zl)=(k,t,0,k,N_0,\mathrm{false})$ and frame $I$.
FC5 gives \texttt{lv0} by retaining its quotient projections;
FC6 gives the additive residue isomorphisms \texttt{fMe0} and
their exact \texttt{fv0} identity. NB7 and NB8 give
\texttt{up\_\allowbreak natural\_\allowbreak level\_\allowbreak zero} and
\texttt{down\_\allowbreak natural\_\allowbreak level\_\allowbreak zero}: substitute the unique chart
label, discard the now automatic chart-equality clause, and retain
the same $\gamma,\eta,\kappa$, lattices, functionals and frame scalar.
These are the source's \texttt{upOfG} and \texttt{dnOfG} conversions.
FC7 gives \texttt{norm0}, and FC17 gives \texttt{bw0}; their endpoint
conversions merely use the identical $e_Y$. We have therefore filled
every field of \texttt{Forms0} directly. Neither the period map nor
the penultimate or terminal form records are needed for this step.

\emph{Next construct the matching level carrier.}
Let $r:k=N_0$ be the fixed \texttt{rank\_\allowbreak eq\_\allowbreak first\_\allowbreak level}
identity. Apply the reindexing construction in FC9 to
$\operatorname{cls0}_u.\operatorname{lev}$, its actual class and
stabilizer matrices, $G$, and the just constructed
$\operatorname{glev0}_u$. This construction is the proof of
\texttt{exists\_\allowbreak colLevel} with only its step-zero arguments:
every occurrence of its apparent full form parameter is through
\texttt{cls0}, \texttt{gl}, or \texttt{glev0}.
Transport the lattice, local groups, global group, functional,
class matrices and stabilizer matrices by $r$; transport every
relation by the same coordinate permutation. The upper-unitriangular
description in FC15 gives the required auxiliary level clause for
\texttt{ColLevel}; the finite prime of $F$ above $q_0$ required in
that record exists by lying over. The stabilizer preserves each actual global
class lattice because it preserves every completed lattice. For a rational
vector, membership in the completion of a prime-local lattice is the same
as membership in that localization: write it in a local basis and test
that each rational coefficient has nonnegative valuation. LP1's local
membership criterion then recovers the global lattice.
Thus these data satisfy all \texttt{ColLevel} fields. Call the result
$L_0(u)$ and choose it for each word.

In particular its five carried objects satisfy exactly
\[
\begin{split}
 \operatorname{reL}_{r}(L_0(u).L_b)&=
                         \operatorname{cls0}_u.\operatorname{lev}.L_b,\\
 \operatorname{reGL}_{r}(L_0(u).U_{\rm loc,v})&=
                         \operatorname{cls0}_u.\operatorname{lev}.U_{\rm loc,v},\\
 \operatorname{reGL}_{r}(L_0(u).U)&=
                         \operatorname{cls0}_u.\operatorname{lev}.U,\\
 \operatorname{reGL}_{r}(L_0(u).y_c)&=y_u(c),\qquad
 \operatorname{reGL}_{r}(L_0(u).s_{c,\gamma})=s_u(c,\gamma).
\end{split}
 \tag{FC.40}
\]
For groups the notation denotes their image under the reindexing
homomorphism. Equation (FC.40) is the literal
$\texttt{Joint0}\,L_0\,u\,T_u$, not merely an abstract isomorphism
of class spaces. The splitting used in $L_0(u)$ is $G.\texttt{sp}$,
and its local frame reindexes to the same identity frame of
$\operatorname{glev0}_u$. Functional normalization has not changed
any term of (FC.40).

Finally SF2 proves \texttt{StdTied} for this actual $D$, including
all of its positive spaces and its terminal positive space; it is
independent of the missing terminal incidence construction. FC16 gives
\texttt{LevelZeroCharacterRational} for $T_u.\texttt{gl}=G$.
The transported base lattice is exactly
$L_{\pi_u}=\mathcal O^{k-1}\oplus\pi_u\mathcal O$ because the
single rank reindexing retains the last coordinate. Put
$\mathfrak a_u=(\bar\pi_u)$. SF3 then proves
\[
 L_0(u).L_b+\mathfrak a_u\bigl(L_0(u).L_b\bigr)^\#
       =\mathcal O^k
       =D.\operatorname{sp}.\operatorname{crdLat0}
                         \bigl((D.\operatorname{sp}.X_0).L_0\bigr).
 \tag{FC.41}
\]
The common standard comparison frame can be taken to be $B=I$.
Consequently the \texttt{Forms0}, \texttt{Joint0}, \texttt{StdTied},
level-zero character-rationality and companion hypotheses in VE9 all
have actual simultaneous suppliers here, under the stated true-column
and initial realizability hypotheses. No full \texttt{CarT},
\texttt{FormsTiedSym}, or \texttt{ColsTiedSym} has been assumed in
order to construct these comparison inputs.

\emph{The separate full-column boundary.}
Ordinary choice over the word/step sets also assembles the other
pointwise fields proved in FC1--FC18. Use the same $G$ everywhere;
unused steps $i>m$ may have singleton chart type. CV9 and FC7 give
the chart values and norm ties, FC8 gives the period tie, and CV10
and FC7 give terminal row values and norms once its row chart is
supplied. The full simultaneous \texttt{FormsTiedSym} route still
requires the particular penultimate $\texttt{HSym}/\texttt{lamSym}$
identification (FC12 already gives the actual common level and finite
carrier), the terminal row level and \texttt{ReadsP}, and the terminal
vector/incidence and adelic evaluation assertion \texttt{til\_\allowbreak ties}.
The corresponding full \texttt{CarT} record must be tied to the same
objects before \texttt{colsTiedSym\_\allowbreak std} is closed. Initial decorated
realizability and the existence of the original arithmetic host,
characters, column laws and later-stage data remain explicit inputs.


\section{The actual complex, rational and integral chart values}
\label{sec:chart-values}

\status{CV1--CV9 prove the source's
\texttt{ValuationCompatibleChart} for the true column family, at every
coefficient prime, from its supplied compatible charts. CV10 gives the
same value statements for the smaller terminal row space. The chart
and carrier constructions are treated in FC; the present argument
constructs the rational forms and proves exact integral membership.
Its conclusion has no norm, Petersson, period or dilation clause.}

\paragraph{Notation and the actual data.}
Keep the true column datum \(I:\mathrm{TC}(B,X)\), its
\(\mathrm{Ok}\), \(\mathrm{Pins}\), \(\mathrm{Laws}\) and
\(\mathrm{Later}\) hypotheses, and the resulting record
\(D=I.\mathrm{rec7}(l,Z)\). Fix a word \(u\), a twist \(a\), and
a level \(i\le m\). Write \(E/F\) for the CM extension,
\(L=B.\mathrm{LL}\) for the coefficient number field,
\(\sigma:L\hookrightarrow\mathbb C\) for its prescribed embedding,
and \(j_\tau:E\hookrightarrow L\) for the prescribed finite family
of embeddings. Thus the actual complex embeddings are
\(\iota_\tau=\sigma\circ j_\tau\). Put \(n=N_i\) and let
\(\mathcal C\) be the source's finite class set
\((I.\mathrm{stp}(u,a,i)).\mathrm{Cl}\).
The member-forms record is denoted by \(\mathcal M\), its
\(L\)-space by \(V_{\mathrm{form}}\), and its integral lattice
by \(\mathcal L\).

Each class has a coordinate lattice \(H_c\subset E^n\).
SF gives the literal standard Hermitian form
\(h(x,y)=\overline{x}^{\,t}y\) in these coordinates, and
identifies \(H_c\) with the source's decorated lattice \(L_c\).
Its Hermitian dual is
\[
 H_c^\#=\{y:\overline{x}^{\,t}y\in\mathcal O_E
                           \text{ for every }x\in H_c\}.
\]
The decorated lattice called \texttt{Lsh} is exactly \(H_c^\#\).
The local basis matrices below describe these class lattices in
the fixed global coordinates; they do not replace the common
global frame of SF.

\paragraph{CV1: frame functions retain the determinant twist.}
Let \(\mu_\tau=\mathcal M.\mathrm{lam}(i,\tau)\), and put
\[
 c_i=\frac{i(t+1)}2,\qquad
 \lambda_\tau=\mu_\tau-c_i(1,\ldots,1),\qquad
 d_\tau=\operatorname{detShift}(\lambda_\tau).
                                                               \tag{CV.1}
\]
Here \(t\) is odd, so \(c_i=i((t+1)/2)\) is an integer with
no rounding discrepancy. This identifies the source's two
opposite conventions for \texttt{shW}, hence its two literal
definitions of \(\lambda\). Their polynomial weights
\(\nu_\tau=\operatorname{natW}(\lambda_\tau)
=\operatorname{natW}(\mu_\tau)\) agree.

A frame \(z=(z_\tau)_\tau\) consists of invertible complex
matrices. It need not be orthonormal. For invertible matrices
\(A=(A_\tau)_\tau\), define
\[
 \mathcal E_{\lambda,A}(P)(z)
   =\left(\prod_\tau\det(A_\tau z_\tau)^{d_\tau}\right)
                      P((A_\tau z_\tau)_\tau).
                                                               \tag{CV.2}
\]
This is exactly \texttt{fnOf}; negative determinant powers are
evaluated only on invertible matrices. Evaluation in (CV.2) is
injective. Indeed its determinant factor is nonzero, so a
polynomial in its kernel vanishes on every invertible tuple.
For an arbitrary tuple \(X\), substitute
\(X_\tau+s_\tau I\). Each determinant excludes only finitely
many values of \(s_\tau\). A polynomial vanishing for the
remaining independent values is zero; evaluate at all
\(s_\tau=0\). This also proves injectivity after mapping
coefficients through \(\sigma\).

Use the source's polynomial action
\[
 \rho_\lambda(g)P(X)
   =\left(\prod_\tau\det(g_\tau)^{d_\tau}\right)
                              P((g_\tau^tX_\tau)_\tau).
 \qquad
 \mathcal E_{\lambda,1}(\sigma(\rho_\lambda(g)P))
       =\mathcal E_{\lambda,\sigma(g)^t}(\sigma P).
                                                               \tag{CV.3}
\]
Substitution and determinant multiplication prove this identity
and \(\rho_\lambda(gh)=\rho_\lambda(g)\rho_\lambda(h)\).
Left translating a frame function by an invertible tuple
therefore preserves its polynomial representation space.
The coordinate equivalence \(e_Y\) used below has
\((e_Yz)_\tau=z_\tau\) exactly; it introduces neither a scalar
nor a change of coefficient embedding.

\paragraph{CV2: the rational and integral polynomial spaces are exact.}
For each \(\tau\), let
\(\Delta_{\nu_\tau}(\mathbb Z)\) be the lattice generated by
divided-power words on the highest polynomial, as in CS2.
Let \(\Delta_\lambda^{\mathrm{std}}\subset L[X]\) be the
\(\mathcal O_L\)-span of the products
\(\prod_\tau f_\tau(X_\tau)\), with
\(f_\tau\in\Delta_{\nu_\tau}(\mathbb Z)\).
It is the source's \texttt{delStd}.
Let \(V_\lambda\) be the source's coefficient space
\texttt{VV}. In characteristic zero, IW3 and CS2 give
\[
       L\Delta_\lambda^{\mathrm{std}}=V_\lambda.
\]
This is a field-span equality. The integral Weyl lattice
\(\Delta\) is retained throughout.

At a coefficient prime \(w\), put
\(R_w=(\mathcal O_L)_w\).
The prescribed valuation identifies \texttt{intW} with
\(R_w=\{r:|r|_w\le1\}\). Directly from the generators,
\[
 \begin{split}
 \operatorname{RatSpan}(\mu)
       &=\{\sigma P:P\in V_\lambda\},\\
 \operatorname{IntSpan}_w(\mu)
       &=\{\sigma P:P\in
                      (\Delta_\lambda^{\mathrm{std}})_w\}.
 \end{split}                                                   \tag{CV.4}
\]
For the first equality, expand each Weyl factor in its integral
divided-power generators and distribute the finite tensor
products. Their \(L\)-span is \(V_\lambda\).
For the second, a finite sum with coefficients in \(R_w\)
has a common denominator outside \(w\), so it lies in the
localization of \(\Delta_\lambda^{\mathrm{std}}\).
Conversely, if \(sP\) belongs to that lattice and \(s\notin w\),
expand \(sP\) in the same integral tensors and divide their
coefficients by \(s\). All resulting coefficients remain in
\(R_w\). This proves the literal additive-closure descriptions
of both spans, with no integral index discarded.
The weight shift in (CV.1) leaves the generators unchanged;
its determinant factor remains in (CV.2).

\paragraph{CV3: rational complex functions give actual rational forms.}
For the moment take any source step with class set
\(\mathcal C\), stabilizers \(\Gamma_c\), matrices
\(\gamma_s\in\mathrm{GL}_n(E)\), and coefficient-field
characters \(\eta_s\in L^\times\). Its complex member space
consists of tuples \(\phi_c\in
\mathcal E_{\lambda,1}(V_{\lambda,\mathbb C})\) satisfying
\[
 \phi_c(z)=\sigma(\eta_s)\,
                 \phi_c(\iota(\gamma_s^{-1})z)
                       \quad(s\in\Gamma_c).                  \tag{CV.5}
\]
Its rational forms are tuples \(P_c\in V_\lambda\) fixed by
\(\eta_s\rho_\lambda(j(\gamma_s)^{-t})\), and their
embedding is \(\phi_c=\mathcal E_{\lambda,1}(\sigma P_c)\).
Equation (CV.3) intertwines these two actions. Consequently
\[
 \begin{split}
 \phi\text{ is the image of a rational form}
 \quad\Longleftrightarrow\quad&
 \phi\text{ belongs to the complex member space, and}\\[-2pt]
 &\phi_c=\mathcal E_{\lambda,1}(Q_c)
     \text{ for some }Q_c\in\operatorname{RatSpan}(\mu)
     \text{ for every }c.
 \end{split}                                                   \tag{CV.6}
\]
To prove the reverse implication, lift \(Q_c\) by (CV.4).
The stabilizer equation and (CV.3), followed by injectivity
in CV1, force each lifted polynomial to satisfy the rational
fixed-vector equation. Thus the tuple is the required
rational form. This constructs a preimage; no averaging
over stabilizers is used.

\paragraph{CV4: extend from class representatives to all adelic points.}
Let \(G(\mathbb A_f)\) be the unitary adelic group,
\(\Gamma=G(F)\), \(U\) a supplied level subgroup, and
\(\chi:U\to\mathbb C^\times\) its character. Suppose
\(c\mapsto[y_c]\) is a bijection
\(\mathcal C\simeq\Gamma\backslash G(\mathbb A_f)/U\).
A tuple of polynomial frame functions \(\phi_c\) extends
to an adelic function with
\[
 f(\gamma y\kappa,\iota(\gamma)z)=\chi(\kappa)f(y,z)
                                                               \tag{CV.7}
\]
if and only if
\[
 y_c=\gamma y_c\kappa,\quad z=\iota(\gamma)z'
       \quad\Longrightarrow\quad
             \phi_c(z)=\chi(\kappa)\phi_c(z').                  \tag{CV.8}
\]
The necessary implication is restriction. For sufficiency,
write \(y=\gamma y_c\kappa\) and set
\[
             f(y,z)=\chi(\kappa)
                          \phi_c(\iota(\gamma^{-1})z).         \tag{CV.9}
\]
If two decompositions are used, bijectivity first gives the
same \(c\). With
\(\delta=\gamma_2^{-1}\gamma_1\) and
\(v=\kappa_1\kappa_2^{-1}\), their equality gives
\(y_c=\delta y_cv\). Equation (CV.8) compares the two
frame arguments and proves that (CV.9) is independent of
the decomposition. It also proves (CV.7).
CV1 shows that every \(f(y,\cdot)\) is still a polynomial
frame function of the required weight. This is the
substantive content of \texttt{mem\_\allowbreak fiC\_\allowbreak iff}.

\paragraph{CV5: the actual stabilizer characters agree.}
The chart reading \texttt{ReadsK} gives a matrix \(k_\beta\)
for each chart and says that \(y_ck_{\beta(c)}\) represents
the actual decorated point \(p_c\). Representation here
retains five conditions: unitarity, the carried global
lattice, the last-coordinate residues at the type primes,
the full prescribed residues at the auxiliary primes, and
the exact rational frames at the twist primes.

These five conditions are preserved by rational unitary
translation:
\[
       g\text{ represents }p
       \quad\Longrightarrow\quad
       \gamma g\text{ represents }(\gamma,1,1)p.
                                                               \tag{CV.10}
\]
Indeed
\((\gamma g)^{-1}x=g^{-1}\gamma^{-1}x\).
This identity transports the lattice and both residue
conditions. The prescribed frame matrices become
\(\gamma\) times the old matrices, which is exactly their
decorated-point action. Unitarity is preserved by
multiplication. Restricting from the common twist-prime set
\(\mathrm{SQ}\) to the particular set \(Q_a\) preserves
these conditions: \(Q_a\subset\mathrm{SQ}\) and
\(Q_a\) is disjoint from the auxiliary-prime set. The latter
follows from \(I.\mathrm{Ok}\), since the selected good
twist primes avoid the set containing the auxiliary primes.

Write a point stabilizer as \(s=(\gamma_s,\ell_s)\).
Since \(sp_c=p_c\),
\[
       (1,\ell_s)p_c=(\gamma_s^{-1},1)p_c.
\]
If \(y_c=\gamma_sy_c\kappa\), then
\[
       \gamma_s^{-1}(y_ck_{\beta(c)})
                     =y_c\kappa k_{\beta(c)}.
                                                               \tag{CV.11}
\]
By (CV.10), the right side represents
\((1,\ell_s)p_c\). Apply the character-reading clause of
\texttt{ReadsK} to this purely level action. The step
character ignores the rational unitary component, so it gives
\[
                    \chi(\kappa)=\sigma(\eta_s).              \tag{CV.12}
\]
The chart's stabilizer-range identity says precisely that
every rational \(\gamma\) with
\(\gamma\in y_cUy_c^{-1}\) comes from a point stabilizer,
with its same rational matrix. Since
\(y_c=\gamma y_c\kappa\) is equivalent to
\(\gamma=y_c\kappa^{-1}y_c^{-1}\), (CV.12) identifies
(CV.5) with (CV.8), in both directions. Applying CV4 on
each chart proves the actual complex member-space clause.

\paragraph{CV6: local bases, scalar extension and the two duals.}
For a rational prime \(p\), write \(\mathcal O_{E,(p)}\)
for the localization obtained by inverting rational integers
prime to \(p\). A matrix \(b\in\mathrm{GL}_n(E)\) is the
source's \(\operatorname{IsLocBasis}(p,H)\) exactly when
\[
                  H_{(p)}=b\mathcal O_{E,(p)}^n.              \tag{CV.13}
\]
If \(\mathfrak q\mid p\), every integer prime to \(p\) is a
unit at \(\mathfrak q\), by Bezout's identity. Localizing
(CV.13) therefore gives
\(H_{\mathfrak q}=b(\mathcal O_E)_{\mathfrak q}^n\).
Conversely, common denominators for the finitely many
coordinates justify this equality in the source's explicit
localization convention. A matrix as in (CV.13) exists for
every full lattice \(H\): over the semilocal Dedekind ring
\(\mathcal O_{E,(p)}\), a finite torsion-free module is free.
One can choose generators of its ideal summands by imposing
their valuations simultaneously at the finitely many primes
over \(p\), using CRT. This is the same rational semilocal
basis construction used in CS3.

For a coefficient prime \(w\), let
\(\mathfrak q_\tau=j_\tau^{-1}(w)\).
It lies over the same rational prime \(p_w\), since
\(j_\tau(p_w)=p_w\). Put
\[
 M_\tau=\mathcal O_L\text{-span}\,j_\tau(H),\qquad
 M_\tau^\vee=\{\ell:\ell^tx\in\mathcal O_L
                                  \text{ for every }x\in M_\tau\}.
\]
This last dual is the ordinary linear dual, with no
conjugation. If \(b\) satisfies (CV.13) for \(p_w\), then
\[
 (M_\tau)_w=j_\tau(b)R_w^n,\qquad
 (M_\tau^\vee)_w=j_\tau(b)^{-t}R_w^n.                          \tag{CV.14}
\]
For the first identity, a denominator outside
\(\mathfrak q_\tau\) maps outside \(w\), so both localized
inclusions survive scalar extension. For the second, the
pairing with the columns of \(j_\tau(b)\) identifies the
local dual. To pass from local pairings to the localized
global dual, clear the finitely many pairing denominators
on a finite generating set of \(M_\tau\).
Apply this construction to \(H=H_c^\#\).
The Hermitian dual used to obtain \(H_c^\#\) and the
ordinary dual in (CV.14) perform different operations;
both are retained.

\paragraph{CV7: localization of the transported Weyl lattice is exact.}
For any full \(H\subset E^n\), the source's
\(\mathcal D_\lambda(H)=\texttt{delAtr}(j,\lambda,H)\)
is the global polynomial lattice obtained from the ordinary
duals \(M_\tau^\vee\). At a prime \(w\), choose bases
\(g_{\tau,w}\) of these duals, and set
\[
 A_w=\rho_\lambda(g_w)
                      (\Delta_\lambda^{\mathrm{std}})_w.
\]
This does not depend on the bases: two choices differ by
matrices in \(\mathrm{GL}_n(R_w)\), which preserve the
integral Weyl lattice and its determinant twist in both
directions. To see the integrality directly, eliminate
an invertible matrix over the local ring using unit pivots.
Its elementary factors act by finite divided-power
expansions; diagonal units and permutations preserve the
integral weight decomposition. Their determinant powers
are units. This is integral stability, not a comparison
of dimensions.

By definition \(\mathcal D_\lambda(H)=\bigcap_w A_w\)
inside \(V_\lambda\), and
\[
                   \mathcal D_\lambda(H)_w=A_w.              \tag{CV.15}
\]
Here is a denominator proof of this localization step.
All the lattices \(M_\tau^\vee\) agree with the standard
ones away from finitely many primes: choose one nonzero
integer bounding each of these finitely generated full
lattices above and below the standard lattice.
Consequently \(A_w=(\Delta_\lambda^{\mathrm{std}})_w\)
away from a finite set. At each of its remaining primes,
\(A_w\) is a full \(R_w\)-lattice. Clear its finitely many
basis denominators using multipliers in \(\mathcal O_L\);
these are integral at all other primes and therefore
preserve their local modules. Their finite product gives
a nonzero common multiplier.
Thus every element of \(V_\lambda\) has one nonzero
multiple belonging to all \(A_w\). Each \(A_w\) is local
at \(w\), so LP2 proves (CV.15).

For the basis \(b\) of CV6, (CV.14)--(CV.15) give
\[
 \begin{split}
 P\in\mathcal D_\lambda(H)_w
 \quad\Longleftrightarrow\quad&
 P=\rho_\lambda\bigl((j_\tau(b)^{-t})_\tau\bigr)P_0\\
 &\text{for some }
             P_0\in(\Delta_\lambda^{\mathrm{std}})_w.
 \end{split}                                                   \tag{CV.16}
\]
Using CV1 and CV2 converts this into the precise value
criterion
\[
 \begin{split}
 P\in\mathcal D_\lambda(H)_w
 \quad\Longleftrightarrow\quad&
 \mathcal E_{\lambda,1}(\sigma P)(z)
      =\mathcal E_{\lambda,\iota(b^{-1})}(Q)(z)
                       \quad\text{for every frame }z,\\
 & Q\in\operatorname{IntSpan}_w(\mu)
                  \text{ for some }Q .
 \end{split}                                                   \tag{CV.17}
\]
The reverse direction uses polynomial injectivity from CV1.
The transpose in (CV.16) becomes exactly \(b^{-1}\) in
(CV.17), and its determinant factor is included in
\(\mathcal E\). No conjugation or determinant factor is
dropped.

\paragraph{CV8: one denominator works for all classes.}
The source's actual lattice is
\[
 \mathcal L=
 \{x\in V_{\mathrm{form}}:
                x_c\in\mathcal D_\lambda(H_c^\#)
                               \text{ for every }c\}.
                                                               \tag{CV.18}
\]
Since \(\mathcal C\) is finite,
\[
 x\in\mathcal L_w
 \quad\Longleftrightarrow\quad
 x_c\in\mathcal D_\lambda(H_c^\#)_w
                              \text{ for every }c.            \tag{CV.19}
\]
The forward implication is immediate. For the reverse,
choose \(s_c\notin w\) clearing the denominator for class
\(c\), and take \(s=\prod_c s_c\).
Then \(s\notin w\) and \(sx\in\mathcal L\).
Scalar multiplication preserves the stabilizer equations,
so this is a denominator within the rational forms space.
The source's valuation module \(\mathcal M.S(i,w)\)
is exactly \(\mathcal L_w\), because its \texttt{intW}
is the same \(R_w\). Also
\(\mathcal L=\bigcap_w\mathcal L_w\): a vector outside
\(\mathcal L\) has a proper denominator ideal contained
in a prime, contradicting membership in that localization.

\paragraph{CV9: the actual valuation-compatible chart.}
Supply the source charts \(K_\beta\), the class map
\(\beta:\mathcal C\to\mathrm{Ch}\), their \texttt{ReadsK}
identity, and \(e_Y\) satisfying \texttt{IdY}.
The source hypotheses are exactly
\[
 I.\mathrm{Ok},\quad\mathrm{Pins}(B,X),\quad
 l:I.\mathrm{Laws},\quad
 Z:\mathrm{Later}(B,I.\mathrm{cols}(l)),\quad i\le m.
                                                               \tag{CV.20}
\]
The bad-place set may be the one supplied to these
charts; no extra exclusion on \(w\) is needed.
The standard-coordinate identification of SF gives
\(\operatorname{gramN}=I\),
\(\operatorname{latCrd}(I.\mathrm{latAt}(c))=H_c\),
and \(\operatorname{dualSet}(H_c)=H_c^\#=\texttt{Lsh}_c\).
The embedding identity is the literal
\(\operatorname{embT}(D)_\tau=\sigma\circ j_\tau\),
proved from the base datum's prescribed embedding
identity, and (CV.1) identifies the weights.

CV4--CV5 give
\[
 \phi\in\mathcal M.\mathrm{Fi}(i)
 \ \Longleftrightarrow\
 \bigl(c,z\mapsto\phi_c(e_Y^{-1}z)\bigr)
       \in\operatorname{FiC}(K_\beta)
                   \quad\text{on every chart }\beta.         \tag{CV.21}
\]
CV3 gives (CV.6) with the actual
\(\mathcal M.\mathrm{emb}(i)\), \(e_Y\), and
\(\operatorname{RatSpan}(D,\mathcal M.\mathrm{lam}(i))\).
Finally CV7--CV8 give, for every coefficient prime \(w\),
\[
 \begin{split}
 x\in\mathcal M.S(i,w)
 \quad\Longleftrightarrow\quad
 &\text{for every }c\text{ and every }
    b\in\operatorname{IsLocBasis}(p_w,H_c^\#),\\
 &\text{there is }Q\in\operatorname{IntSpan}_w(\mu)
    \text{ such that, for every }z,\\[-2pt]
 &\mathcal M.\operatorname{emb}(i)(x)_c(z)
       =\mathcal E_{\lambda,\iota(b^{-1})}(Q)(e_Yz).
 \end{split}                                                   \tag{CV.22}
\]
The reverse implication is not vacuous: CV6 supplies at
least one such \(b\) for each class and residue prime.
These are precisely the three clauses of
\texttt{ValuationCompatibleChart}, with
\(\texttt{okW}\equiv\mathrm{True}\).
They prove \texttt{TC.val\_\allowbreak tied} and its exported
\texttt{hermitianChart\_\allowbreak complex\_\allowbreak rational\_\allowbreak integral\_\allowbreak forms}
statement. The all-prime integral criterion is now proved
for the supplied charts, rather than placed among their
assumed properties.

\paragraph{CV10: the smaller terminal row space uses the same proof.}
For \(\texttt{clsP\_\allowbreak tie}\), take \(n'=N_m-1\), the actual
row step \(I.\mathrm{stX}(\mu_0)\), and its supplied row
class chart with \texttt{ReadsP}. Its untwisted coefficient
weight is \(\mu'=\texttt{sgW}\), and its frame weight is
\[
                  \lambda'=\mu'-\frac{m(t+1)}2(1,\ldots,1).
                                                               \tag{CV.23}
\]
The equality of the two source definitions follows from
the same oddness calculation as (CV.1). Thus both
\(V_{\lambda'}=V_{\mu'}\) and
\(\Delta_{\lambda'}^{\mathrm{std}}
 =\Delta_{\mu'}^{\mathrm{std}}\), while the frame
determinant exponent stays that of \(\lambda'\).
At this row step the type-prime decoration set is empty.
The decorated lattice \texttt{Lsh} therefore equals its
carried lattice \(N_c\) exactly.

CV4--CV5 prove the terminal \texttt{ptC} clause using the
row chart's same stabilizer matrices and character.
CV3 proves its \texttt{rat} clause.
CV6--CV8, followed by intersection over all coefficient
primes, prove its \texttt{lat} clause with \(N_c\) in place
of \(H_c^\#\). The frame identification is the identity
on matrix entries, as in \texttt{clsP\_\allowbreak tie}.
The supplied terminal chart and its remaining construction are tracked in FC; this
corollary supplies its value clauses. Its
\texttt{ipP} and \texttt{NormPTie} clauses require the
separate normalized Fock and Hermitian metric argument
and are not consequences of (CV.21)--(CV.22).

\paragraph{Boundary of this construction.}
The proof constructs rational preimages and proves the
exact integral coefficient criterion for the actual
step lattices. It retains the supplied true column datum,
its \texttt{Ok/Pins/Laws/Later} witnesses, and the chart
reading identities used in CV5. FC describes their
simultaneous chart and carrier assembly. CV does not
infer those inputs from their names, identify an
integral Weyl lattice with its dual-Weyl enlargement,
or infer a metric identity from the lattice equality.


\section{Actual up and down neighbors and their chart naturality}
\label{sec:neighbor-geometry}

\status{NB1--NB5 construct and identify the actual neighboring decorated
lattices, their residue maps and their chart labels. NB6--NB8 prove the
full up/down naturality statements using AD3's supported splice. The
initial adelic realization of a valid decorated point remains an explicit
input. No equality of integral lattices is deduced merely from equal rank.}

\paragraph{Conventions and exact inputs.}
Put $\mathcal O=\mathcal O_E$, $V=E^n$, and
$h(x,y)=\bar x^{\mathsf T}y$. For a full $\mathcal O$-lattice $M$ write
\[
 M^\#=\{y\in V:h(M,y)\subseteq\mathcal O\},\qquad
 M_P=\mathcal O_{(P)}M\subset V.
\]
Here $M_P$ is localization inside the rational vector space; completed
lattices appear only when applying AD3. All sums and ideal products
below are sums and products of the actual integral submodules of $V$.
Let $d=(N,(f_Q)_{Q\in S_T},(\mathrm{be}_R)_{R\in S_G})$ be a valid
source decoration. Thus $N=N^\#$ is full, each $f_Q$ is a nonzero
residue-linear map $N_Q\to k_Q=\mathcal O/Q$, extended by zero outside
$N_Q$, and each $\mathrm{be}_R$ gives the prescribed residue basis
isomorphism at its auxiliary prime. Put
\[
 L=\{x\in N:f_Q(x)=0\text{ for all }Q\in S_T\},\qquad H=L^\#.
\]
Fix $P\in S_T$ such that $P\ne\bar P$, no member of $S_T$ is $\bar P$,
and all auxiliary and twist-frame primes avoid $P,\bar P$. These are
the precise \texttt{MoveHyp} and frame-disjointness requirements.
For an up move take $q\in P\setminus P^2$. In the true-column application
$q=T_u$: FC1 gives $v_P(\pi_u)=1$, $v_P(\bar\pi_u)=0$, and
$T_u=\pi_u\bar\pi_u$, so this is the actual rational integer $T_u$,
not a replacement uniformizer.

For the actual data the separation at the moved prime follows from
the stated prime and twist laws. No selected type prime can be the
conjugate of another, because $(\pi_u,\bar\pi_u)=\mathcal O$.
The auxiliary rational primes are different from $T_u$. A twist
prime has residue characteristic outside the set $S$ containing
$T_u$, so it cannot contain $T_u$; restriction to $\mathcal O_F$
proves the same assertion for every prime above the moved base-field
prime. Thus all auxiliary and twist primes avoid $P,\bar P$.
When using AD3 we retain its full structural hypotheses on these
actual sets: pairwise disjointness and absence of a conjugate pair
in their union. These are supplied by the same prime/selection/twist
data used in the chart construction, rather than by any neighbor
existence assertion.

The application in NB5--NB8 retains the supplied
\texttt{TC.Ok}, \texttt{TC.Pins}, $l:\texttt{I.Laws}$ and
$Z:\texttt{TC.Later}$, and the exact chart data
\texttt{ReadsK}, \texttt{IdY}, \texttt{IdV}, \texttt{IdL}.
FC constructs the indicated chart and projection data under its stated
geometric inputs. NB uses neither the final Forms0 assembly nor its
up/down clauses as a premise.

\paragraph{NB1: localization of the Hermitian dual uses the conjugate prime.}
For every full lattice $M$ one has
\[
 (M^\#)_P=\{y\in V:h(x,y)\in\mathcal O_{(P)}
                         \text{ for all }x\in M_{\bar P}\}.
 \tag{NB.1}
\]
Indeed, if $sy\in M^\#$ with $s\notin P$, and $tx\in M$ with
$t\notin\bar P$, then
$h(x,y)=h(tx,sy)/(\bar t s)$ is integral at $P$.
Conversely, take finitely many $\mathcal O$-generators $m_j$ of $M$.
The proposed condition puts every $h(m_j,y)$ in $\mathcal O_{(P)}$.
A product $s\notin P$ of their denominators gives
$h(M,sy)\subseteq\mathcal O$, and hence $sy\in M^\#$.

The pairing between $\mathcal O_{(\bar P)}^n$ and
$\mathcal O_{(P)}^n$ is perfect. In local bases its dual is given by
the inverse conjugate transpose of the basis matrix; taking this
operation twice returns the original local lattice. Thus
$M^{\#\#}=M$, first locally by this matrix calculation and then globally
by LP1. The same calculation proves that $M^\#$ is a full lattice:
clearing the finitely many basis-matrix denominators places it between
two nonzero scalar multiples of $\mathcal O^n$.
Consequently equality of the localized lattices at every finite prime
proves equality of the global integral lattices. Conjugation has been
used as an isomorphism between the two local rings; it is not being
treated as an automorphism of $\mathcal O_{(P)}$ when $P\ne\bar P$.
Equality $M_P=M'_P$ also gives the source's \texttt{SameAt} witnesses:
clear denominators outside $P$ on finitely many generators and multiply
them to obtain one $a\notin P$ with $aM\subseteq M'$, and repeat in
the opposite direction.

\paragraph{NB2: the actual up lattice and its residue kernel.}
Choose $x\in N$ with $f_P(x)=1$ and define the source's actual lattices
\[
 A=N+P^{-1}x,\qquad U=A^\#+\bar P A,
 \qquad P^{-1}=\{a\in E:aP\subseteq\mathcal O\}.
 \tag{NB.2}
\]
Here $P^{-1}x$ denotes its $\mathcal O$-span. Since
$N\subseteq A\subseteq P^{-1}N$, $A$ is full and
$\bar P N\subseteq A^\#\subseteq N$. Localizing the invertible ideal
$P^{-1}$, and using (NB.1), gives the complete table
\[
\begin{array}{c|ccc}
 &P&\bar P&R\notin\{P,\bar P\}\\ \hline
 A&N_P+q^{-1}\mathcal O_{(P)}x&N_{\bar P}&N_R\\
 A^\#&N_P&(A^\#)_{\bar P}&N_R\\
 U&A_P&(A^\#)_{\bar P}&N_R.
\end{array}                                                    \tag{NB.3}
\]
For example $\bar P$ is a unit ideal at $P$, while
$\bar P N_{\bar P}\subseteq(A^\#)_{\bar P}$; these facts justify
both exceptional entries in the last row. At $P$ and $\bar P$ these
entries are dual to each other by (NB.1); at every other prime they are
the entries of $N=N^\#$. Therefore $U=U^\#$ globally. The definition
also gives $U\subseteq A$ and $\bar P A\subseteq U$. This proves all
the lattice assertions of \texttt{upN\_\allowbreak spec}, including actual equality
away from the conjugate pair.

Update only the selected type functional, setting
\[
 f_P^+(y)=f_P(qy)\quad(y\in U_P),\qquad
 f_P^+(y)=0\quad(y\notin U_P).                       \tag{NB.4}
\]
Every $y\in U_P$ has the form $n+q^{-1}a x$ with
$n\in N_P$ and $a\in\mathcal O_{(P)}$. Residue-linearity extends to
this local ring by clearing a denominator outside $P$, so
$f_P^+(y)=a\bmod P$. Hence
\[
 \ker(f_P^+|_{U_P})=N_P,\qquad U_P/N_P\simeq k_P.     \tag{NB.5}
\]
In particular the new functional is nonzero on the global lattice:
multiply $q^{-1}x\in U_P$ by a denominator $s\notin P$ placing it in
$U$; its new value is the nonzero residue of $s$.

Keep every other $f_Q$ and every $\mathrm{be}_R$ literally unchanged.
Their local domains agree by (NB.3), so their local laws persist.
The global nonzero and surjectivity witnesses are obtained by clearing
a denominator outside their own prime and rescaling by its inverse
residue. The auxiliary kernel assertion persists as well: if
$z\in U$ has zero auxiliary residue at $R$, then $z\in R U_R$;
choose $s\notin R$ with $sz\in R U$ and $t$ with
$ts\equiv1\pmod R$. The equality
$z=t(sz)+(1-ts)z$ puts $z$ in $RU$. Thus this is a valid decoration,
the literal \texttt{d.up P q x}, with no fallback case.

For completeness the finite residue conditions give
\[
 L_P=\ker(f_P|_{N_P}),\qquad L_R=N_R\ (R\notin S_T),
 \qquad H_P=N_P.                                     \tag{NB.6}
\]
To obtain the first equality, clear a denominator into $N$, then
multiply by elements of each $Q\in S_T\setminus\{P\}$ which lie in
$Q$ but not in $P$; the other conditions vanish without changing the
localization at $P$. The second equality has the same proof. For the
third, $\bar P\notin S_T$ and (NB.1) apply. Writing $L^+,H^+$ for the
new decoration, (NB.3)--(NB.6) yield
\[
 N\subseteq H,\qquad (L^+)_P=N_P=H_P,\qquad
 (H^+)_P=U_P=(H+\mathcal Oq^{-1}x)_P.                \tag{NB.7}
\]
Together with (NB.4), these are the complete
\texttt{MoveUpLat.lattices\_\allowbreak at\_\allowbreak up} assertions.

\paragraph{NB3: the actual down lattice and its extended functional.}
Put
\[
 K_P^{\rm gl}=\{x\in N:f_P(x)=0\},\qquad
 D=K_P^{\rm gl}+P(K_P^{\rm gl})^\# .                 \tag{NB.8}
\]
The notation $K_P^{\rm gl}$ denotes a global lattice, not a localization.
It lies between $PN$ and $N$, so it is full. At every $R\ne P$ it has
localization $N_R$: multiply an element of $N$ by an element of
$P\setminus R$ to kill its residue. Its localization at $P$ is
$\ker(f_P|_{N_P})$. The paired-dual calculation in NB1 now gives
\[
 D_P=(K_P^{\rm gl})_P=L_P,\qquad
 D_{\bar P}=((K_P^{\rm gl})^\#)_{\bar P},\qquad
 D_R=N_R\quad(R\notin\{P,\bar P\}).                 \tag{NB.9}
\]
Indeed $((K_P^{\rm gl})^\#)_P=N_P$ and
$PN_P\subseteq(K_P^{\rm gl})_P$; at $\bar P$, the ideal $P$ is a unit.
The two exceptional entries of (NB.9) are dual to one another.
Thus $D=D^\#$, which proves the actual \texttt{dnN\_\allowbreak spec}.
Also $N_P/D_P\simeq k_P$, by the nonzero functional $f_P$.
These are local indices; $D$ and $N$ need not be comparable globally.

Let $g:D/PD\to k_P$ be any nonzero $k_P$-linear functional. For
$y\in D_P$, choose $s\notin P$ with $sy\in D$ and define
\[
 f_P^-(y)=\rho_P(s)^{-1}g([sy]),                     \tag{NB.10}
\]
where $\rho_P:\mathcal O\to k_P$ is the actual quotient map;
put $f_P^-(y)=0$ outside $D_P$. If $t$ is another denominator,
linearity gives $\rho_P(t)g([sy])=\rho_P(s)g([ty])$.
This proves independence, additivity and residue-linearity. On $D$
the formula is exactly $g\circ[\cdot]$, so it is nonzero there.
This is precisely both \texttt{TDat.extF} and \texttt{modP.ext};
no different quotient functional has been chosen.

Update $f_P$ to (NB.10), leave all other functionals and auxiliary maps
unchanged, and use the same witness and kernel arguments as in NB2.
The result is the actual valid \texttt{d.dn P (TDat.extF D P g)}.
If $L^-,H^-$ are its kernel lattice and Hermitian dual, then
\[
                         (H^-)_P=D_P=L_P.           \tag{NB.11}
\]
This follows from (NB.6), applied to the new valid decoration, and is
the substantive local identity in \texttt{dn\_\allowbreak locP\_\allowbreak lat}.

\paragraph{NB4: the actual moves preserve frames, coupling and chart label.}
The source's framed point move \texttt{PtY.mv} returns its input if the
proposed new lattice invalidates a frame or its coupling set. That
alternative does not occur here. Every twist frame is a basis of the
same localized lattice before and after (NB.3) or (NB.9). The actual
coupling sets depend only on the auxiliary and twist localizations:
at the penultimate step they are cosets defined from those local
lattices and the prescribed uniformizer powers, and at other steps
the coupling set is empty. FC2 described the reference case with no
twist primes. These localizations are all
unchanged. Thus the actual moved point has exactly the new decoration
of NB2 or NB3, the old rational frames, and the old coupling subset
$Y\subset V$. This proves the full content of \texttt{mv\_\allowbreak up} and
\texttt{mv\_\allowbreak dn}, including validity.

Here is why the same statement survives passage to the preferred
class representatives. Two framed points will be called equal on
the coupling data when they have the same twist frames, the same
$Y$, and the same auxiliary local lattices and maps
$\mathrm{be}_R$. Their chosen auxiliary lift-frames are then equal,
since those choices depend on the same lattice and reduction map.
Consequently their \texttt{cosetIn} labels are equal. Applying the
same column-group element to both points preserves this relation:
localization commutes with its $E$-linear rational part, the two
coupling sets have the same image, and the finite components transform
both frame and residue maps by the same formulas. This proves the
source's \texttt{SameG.smul} assertion directly.

At the penultimate step the preferred representative is chosen to
have the selected value of \texttt{cosetIn} from its nonempty set of
attainable values along the orbit. The old point and its raw neighbor
have the same value after every common group action; hence their two
sets of attainable values are identical. The fixed choice on this
set gives the same preferred value. At every other step both labels
are empty. The column index $b$ is unchanged, and \texttt{ReadsK.key}
identifies a chart label exactly by $b$ and this preferred coupling
value. Therefore both actual up and down moves keep \texttt{chOf}.
This proves \texttt{up\_\allowbreak keeps\_\allowbreak chart} and \texttt{dn\_\allowbreak keeps\_\allowbreak chart};
equality of the raw labels alone would not have established this.

We shall also use invariance under a pure rational unitary action.
At a twist prime the rational frame changes from $F_R$ to $\gamma F_R$,
so its coordinates of $\gamma y$ are identical to the old coordinates
of $y$. At an auxiliary prime, let $F_R,F'_R$ be the chosen old and
new lift-frames. Then
\[
 B_R=(F'_R)^{-1}\gamma F_R\in\mathrm{GL}_n(\mathcal O_{(R)}),
 \qquad B_R\equiv I\pmod R.                          \tag{NB.12}
\]
The reduction identity follows because both frames lift the same
transported auxiliary basis. The actual coupling law gives
$q_R F_R^{-1}y\in\mathcal O_{(R)}^n$, where the auxiliary rational
prime $q_R$ is a uniformizer. Thus
$(B_R-I)F_R^{-1}y\in\mathcal O_{(R)}^n$. The new coordinates and old
coordinates define exactly the same coupling coset. At twist primes
they were equal already, including any prescribed higher congruence
power. This proves \texttt{cosetIn\_\allowbreak smul\_\allowbreak unitary} with the actual
chosen lift-frames and supplies the rational-action condition needed
for the level-reading argument below.

\paragraph{NB5: transport by the actual moving element and compare projections.}
Let $c'$ be the class of either raw neighbor and let
$g_0=(\gamma,h_0)$ be the source's chosen moving element, so
\[
                      g_0\,r(c')=p_{\rm raw}.        \tag{NB.13}
\]
Its rational part $\gamma$ is unitary. The finite part changes only
the finite decoration; hence $\gamma$ carries the preferred new
$N,L,H$ to the corresponding raw $N,L,H$. For $L$ this follows by
pulling back all type functionals, and for $H$ it follows from
$h(\gamma x,\gamma y)=h(x,y)$. The rational matrix is therefore
exactly \texttt{gam(mv)}, not an independently chosen lattice isometry.

For an up move start with $z\in H_c$ satisfying
$\varphi_c(\operatorname{pr}_c z)=1$, where $\varphi_c$ is the
functional on $N_c/PN_c$. Choose the source's quotient representative
$x_0\in N_c$ of $\operatorname{pr}_c z$. FC5 gives
$z-x_0\in PH_c$ and $f_P(x_0)=1$. Since $q$ is a uniformizer,
$q^{-1}(z-x_0)\in(H_c)_P$. Thus (NB.7) and (NB.13) give
\[
 (\gamma L_{c'})_P=(H_c)_P,\qquad
 (\gamma H_{c'})_P=(H_c+\mathcal Oq^{-1}z)_P.         \tag{NB.14}
\]
For $x\in H_{c'}$ with $q\gamma x\in H_c$, choose a common
$s\notin P$ with $sx\in N_{c'}$ and $sq\gamma x\in N_c$.
Such a denominator exists by $(H_c)_P=(N_c)_P$ and its new analogue.
On these global elements (NB.4) and (NB.13) give
$f_{c',P}(sx)=f_{c,P}(sq\gamma x)$. FC5 identifies the two global
reductions with the projections. Their residue-linearity and
$\rho_P(s)\ne0$ then give the exact identity
\[
 \varphi_{c'}(\operatorname{pr}_{c'}x)
       =\varphi_c(\operatorname{pr}_c(q\gamma x)).    \tag{NB.15}
\]

For a down move, let $g:D_c/PD_c\to k_P$ be its specified nonzero
functional and put $K_c=L_c+PH_c$. Its localization is $(D_c)_P$
by (NB.6) and (NB.9). Equations (NB.11) and (NB.13) give
\[
                         (\gamma H_{c'})_P=(L_c)_P. \tag{NB.16}
\]
If $x\in H_{c'}$ and $\gamma x\in K_c$, choose $s\notin P$ with
$sx\in N_{c'}$. Equation (NB.13) places $s\gamma x$ in the actual
global lattice $D_c$. By (NB.10) its new residue is $g([s\gamma x])$.
FC5 identifies $[sx]$ and $[s\gamma x]$ with
$\operatorname{pr}_{c'}(sx)$ and $\operatorname{prk}_c(s\gamma x)$.
Cancel the same nonzero residue $\rho_P(s)$ to obtain
\[
 \varphi_{c'}(\operatorname{pr}_{c'}x)
                   =g(\operatorname{prk}_c(\gamma x)).           \tag{NB.17}
\]
This proves \texttt{dn\_\allowbreak level\_\allowbreak functional} for the exact source
projection maps and chosen moving matrix. No assertion that $K_c=D_c$
globally has been used; only their equality at $P$ is required.

\paragraph{NB6: a supported adelic factorization with the exact character.}
Assume the retained realization input: each valid framed decorated
point under consideration has a unitary adelic realizer in the full
\texttt{IsPtOf} sense. Take $h=y_ck$ realizing $r(c)$ and a realizer
of the actual raw neighbor. NB2--NB4 verify every comparison premise
of AD3: the lattices agree off $P,\bar P$, every other type functional
is the same, and all auxiliary maps and rational frames are unchanged.
AD3 supplies a unitary $\varepsilon$, equal to $1$ off that pair,
such that $y_ck\varepsilon$ realizes the raw neighbor. Set
\[
                         \eta=k\varepsilon k^{-1}. \tag{NB.18}
\]
Then $y_c\eta k=y_ck\varepsilon$. At every $F$-prime other than
$v_0=P\cap\mathcal O_F$, all its $E$-components are the identity,
so its projection to $E\otimes_FF_v$ is exactly $1$.
Conjugating by the same base point $k$ preserves this support and
unitarity. NB4 ensures that the new preferred class has the same
chart and therefore the very same base point $k$.

For the actual $g_0=(\gamma,h_0)$ of (NB.13), put
\[
 \kappa_0=y_{c'}^{-1}\operatorname{rat}(\gamma)^{-1}y_c\eta.
                                                               \tag{NB.19}
\]
The point realized by $y_{c'}\kappa_0k$ is
$\gamma^{-1}p_{\rm raw}=h_0r(c')$. Its coupling label equals that
of $r(c')$: NB4 proves equality of the raw and preferred labels,
and also invariance under the rational unitary $\gamma^{-1}$.
The exact level-reading equivalence in \texttt{ReadsK} therefore
places $\kappa_0$ in the chart level $U_\beta$; its character-reading
identity gives
\[
 \chi_\beta(\kappa_0)=\sigma(\chi_{\rm col}(g_0)).  \tag{NB.20}
\]
Here $\sigma$ is the supplied embedding of the coefficient field,
with the same identification used in CV and FC11. Set
$\kappa=\kappa_0^{-1}$. Rearranging (NB.19) and using the character
homomorphism proves
\[
 \operatorname{rat}(\gamma)y_{c'}=y_c\eta\kappa,
 \qquad \chi_\beta(\kappa)^{-1}
                   =\sigma(\chi_{\rm col}(g_0)).     \tag{NB.21}
\]
This one direct-splice proof supplies both the up factorization and
the conclusion of \texttt{move\_\allowbreak down\_\allowbreak adelic}, with the latter's
exact moving matrix. The source's down proof uses a reverse splice;
under the explicit realization input the direct splice above gives
the same required existential conclusion without changing its scalar.

\paragraph{NB7: the literal up naturality statement is now proved.}
Take the exact true-column data listed in the conventions, $i\le m$,
a chart class $c$, its selected prime $v$, and
$z\in\operatorname{LdG}(\operatorname{lat}c)$ with
$\mathrm{fv}_c(\operatorname{pr}_c z)=1$.
The coordinate identities of SF and FC identify
$\operatorname{lat}c=L_c$, $\operatorname{LdG}(\operatorname{lat}c)=H_c$,
and $\operatorname{PvG}(v)=P$. The identity \texttt{IdV} identifies
the base-field prime $(e_Vv)$ with $P\cap\mathcal O_F$, with its prescribed
orientation. Apply NB2 and NB4 to the quotient representative of
$\operatorname{pr}_c z$. They produce the source's actual up class
$c'=\mathrm{up}(v,c,\operatorname{pr}_c z)$ in the same chart.
Take its actual $\gamma$ from (NB.13), and $\eta,\kappa$ from NB6.
The support and adelic clauses are (NB.18)--(NB.21); the two lattice
clauses are (NB.14) with $q=T_u$; the functional clause is (NB.15).

Finally, for arbitrary $h$ and frames $s,s'$ related by
$e_Y(s)=\operatorname{atInf}(\gamma)e_Y(s')$, the literal
\texttt{IdY} and embedding identities give $\gamma^{-1}s=s'$
componentwise. The source's transport formula is consequently
\[
 \mathrm{upT}(v,c,\operatorname{pr}_c z)h(s)
       =\sigma(\chi_{\rm col}(g_0))h(s')
       =\chi_\beta(\kappa)^{-1}h(s').                \tag{NB.22}
\]
This checks every clause of \texttt{UpNatTieG}, and proves the actual
\texttt{UpChart.upNatTie}. The identity on frames uses
$\iota_\tau=\sigma\circ j_\tau$ at the supplied embeddings; no
complex conjugation, determinant factor or scalar normalization is
inserted into (NB.22).

\paragraph{NB8: the literal down naturality statement is now proved.}
For the same exact data, take any nonzero
$g\in\operatorname{Hom}_{k_P}(D_c/PD_c,k_P)$, which is the source's
\texttt{lk} dual under \texttt{IdL}. NB3 constructs the actual
$\mathrm{dnPt}$ using this $g$ and (NB.10), and NB4 puts its class
$c'=\mathrm{dn}(v,c,g)$ in the old chart. Use the actual $\gamma$
from (NB.13) and the supported $\eta,\kappa$ from NB6.
Equation (NB.21) supplies the exact adelic relation, (NB.16) supplies
the transported dual-lattice equality, and (NB.17) supplies the
functional comparison on every $x\in H_{c'}$ with
$\gamma x\in L_c+PH_c$. The same frame calculation as in NB7 yields
\[
                  \mathrm{dnT}(v,c,g)h(s)
                      =\chi_\beta(\kappa)^{-1}h(s'). \tag{NB.23}
\]
These are all the clauses of \texttt{DnNatTieG}, hence the actual
\texttt{DnMove.dnNatTieG\_\allowbreak cols}.

The conclusion closes the neighbor and coupling assertions formerly
retained in the Forms0 construction. Its scope remains conditional on
the literal supplied true-column laws and on initial adelic realization.
The latter is not supplied by a name for a record, by an equal-rank
argument, or by the local modifications just proved. NB makes no
additional claim about terminal-row incidence or the upstream
construction of \texttt{TC.Ok}, \texttt{Pins}, \texttt{Laws}, or
\texttt{Later}.


\section{Constructing the comparison values over the coefficient field}
\label{sec:value-existence}

\status{VE1--VE9 remove the separate polynomial-value existence premise
from CS. They prove block rationality, adelic point existence and
representative independence, then construct the literal \texttt{ValEx}
and concrete comparison under supplied global forms, standard-frame,
level and character data. FC19 now supplies those data from the actual
true-column and initial-realization hypotheses, and VE10 obtains their
concrete comparison. DR1--DR3 then construct that initial realization
from the supplied CM and decoration data; VE11 discharges this remaining
geometric premise. The independent metric, spectral, cohomological
and arithmetic-window inputs retain their separate scope.}

\paragraph{The data and the assertion.}
Let $E/F$ be the CM extension, $L$ the coefficient field with its
specified embedding $\jmath:L\hookrightarrow\mathbb C$, and
$t_\tau:E\hookrightarrow L$ the specified lifts of the complex
embeddings. Keep the same frame $B$, place identification and position
map as LI and SF. The supplied level-zero forms give functions
$f_{x,c}(z)$ for $x\in V$, a class $c$ and an invertible complex
matrix frame $z=(z_\tau)_\tau$. Here the source's
\texttt{Frames} consists of all invertible matrix frames: its Gram
label does not require the matrices themselves to be unitary.

For an away component $i$ and a tuple of full local lattices
$\mathbf L$, we will construct a unique polynomial
$Q_{x,i,\mathbf L}$ over $L$ in the prescribed tensor representation.
For every basis $g$ of $\mathbf L$ and every presentation of its
associated adelic point as $a=\gamma y_c\kappa$, its evaluation must
equal
\[
 \operatorname{ev}_C(\jmath Q_{x,i,\mathbf L})(z)
       =\chi(\kappa)f_{x,c}(z'),\qquad z=\gamma_\infty z'.
                                                        \tag{VE.1}
\]
Thus one polynomial must work for \emph{all} these presentations.

\paragraph{VE1: explicit rational polynomials for each block.}
Write $k=m_{\rm loc}+1$, $\lambda=(\lambda_1,\ldots,\lambda_k)$,
$e=\lambda_k$, $\lambda^*=(-\lambda_k,\ldots,-\lambda_1)$ and
$e^*=-\lambda_1$. Put
\[
 T(Y)=Y^{-\mathsf T}R_k,
\]
where $R_k$ reverses columns. The integral duality of CS2 supplies,
for every $f\in\Delta_{\nu(\lambda)}(\mathbb Z)$, a polynomial
$f^*\in\Delta_{\nu(\lambda^*)}(\mathbb Z)$ satisfying
\[
 \det(Y)^e f(Y)=\det(T(Y))^{e^*}f^*(T(Y))
       \quad(Y\in\mathrm{GL}_k).
                                                        \tag{VE.2}
\]
Its sign normalization is exactly the one in CS2; explicitly,
$f^*(X)=\epsilon^{\lambda_k}\det(X)^{\lambda_1-\lambda_k}
f(X^{-\mathsf T}R_k)$ on invertible matrices, with
$\epsilon=\det R_k$. Complementary minors show that this is an
integer polynomial on the specified Weyl space.

Suppose the block coordinate matrix is $M=\jmath(M_0)$ with
$M_0\in\mathrm{GL}_k(L)$. Define
\[
 F_{\rm ev}(X)=\det(M_0)^{-e}f(M_0^{-1}X),\qquad
 F_{\rm odd}(X)=\det(M_0)^{e^*}f^*(M_0^{\mathsf T}X).
                                                        \tag{VE.3}
\]
All coefficients lie in $L$. These polynomials belong to the
required dual-Weyl block spaces. Indeed integral Weyl generators lie
in the column-homogeneous, upper-column-unipotent invariant space;
row substitutions commute with the column action and preserve
column degrees. Scalar multiplication preserves this space too.
No equality of the integral lattices $\Delta$ and $\nabla$ is used.

If $Z=MY$, the even expression gives
$\det Z^e\,\jmath(F_{\rm ev})(Z)=\det Y^e f(Y)$.
If $Z=T(MY)=M^{-\mathsf T}T(Y)$, then
$M^{\mathsf T}Z=T(Y)$, and (VE.2) gives
\[
 \det Z^{e^*}\,\jmath(F_{\rm odd})(Z)=\det Y^e f(Y).
                                                        \tag{VE.4}
\]
The source's \texttt{WtTie} assigns precisely these two weights to
even and odd positions. Equations (VE.3)--(VE.4) therefore prove
\texttt{BlockRationalityField}. For the actual frame coordinates,
$M_0=t_\tau(B^{-1})$; hence they also prove
\texttt{BlockRationalityAtPlace}, with the original embeddings.

\paragraph{VE2: tensor the blocks and recover representative values.}
The rationality field of the supplied \texttt{Forms0} says that each
$f_{x,c}$ is a finite sum of terms
\[
 \jmath(r)\prod_\tau\det(z_\tau)^{\lambda_{\tau,k}}
                                  f_\tau(z_\tau),
 \qquad r\in L,\quad f_\tau\in\Delta_{\nu(\lambda_\tau)}(\mathbb Z).
                                                        \tag{VE.5}
\]
For one term choose the block polynomials from VE1, multiply them
in disjoint variable blocks, and multiply by $r$. The result belongs
to the source's $L$-tensor space \texttt{VtupR}. Evaluating its product
gives (VE.5), since the position map is a bijection. Add the resulting
polynomials for a finite sum, and use their negatives for additive
inverses. This proves \texttt{evRat\_\allowbreak gen}, \texttt{evRat\_\allowbreak of\_\allowbreak blocks}
and the actual \texttt{EvRatRep}: every representative value has a
polynomial $Q_c$ over $L$.

\paragraph{VE3: rational transport from one class representative.}
For a rational unitary $\gamma$, put $H=B^{-1}\gamma B$.
The same standard-frame identity as SF gives $H^*H=I$. In position
$(v,s)$ define the coefficient-field matrix
\[
 h_{v,s}=
 \begin{cases}
 t_\tau(H),&s\text{ even},\\
 t_\tau(\overline H),&s\text{ odd},
 \end{cases}
 \qquad\tau=\operatorname{pos}(v,s).
                                                        \tag{VE.6}
\]
Because $\overline H=(H^{-1})^{\mathsf T}$, these are the exact
even and odd frame-transport matrices, including the conjugation
already required by LI. The natural action is
\[
 (\rho(h)Q)(X)=
    \left(\prod_{v,s}\det(h_{v,s})^{-e_{v,s}}\right)
             Q\bigl((h_{v,s}^{-1}X_{v,s})_{v,s}\bigr).
                                                        \tag{VE.7}
\]
It preserves \texttt{VtupR}: each row substitution preserves its
dual-Weyl block, tensor generators map to tensor generators, and
linearity handles their span. This proves the source's
\texttt{VtupStab} directly.

If $z=\gamma_\infty z'$, then the block coordinates satisfy
$Z(z)=hZ(z')$. Substitution in (VE.7), including its determinant
factor, proves
$\operatorname{ev}_C(\jmath\rho(h)Q)(z)
=\operatorname{ev}_C(\jmath Q)(z')$.
The supplied level-zero character rationality gives
$\chi(\kappa)=\jmath(r_\kappa)$: the level character is a finite
product of the supplied character values and their inverses, and
the higher-block factors are absent at level zero. Thus
\[
 Q=r_\kappa\rho(h)Q_c
                                                        \tag{VE.8}
\]
is a rational polynomial giving the right side of (VE.1) for this
presentation. This proves \texttt{EvRatPt}; the rational action and
character lift have been constructed rather than postulated.

\paragraph{VE4: construct the required adelic points.}
At each place above $p$, LI's split unitary group is isomorphic to
$\mathrm{GL}_k(F_v)$ by projection to its first component; the
second component is determined by the Hermitian equation. Given
$g=(g_v)_v$, insert the first components $B_vg_vB_v^{-1}$ at these
finitely many places and the identity at all others. The resulting
section $j_p(g)$ is a finite adele in the unitary group.
The away representative $y_i$ is the identity above $p$, so
\[
 a=y_i j_p(g)
                                                        \tag{VE.9}
\]
has the exact prescribed local components and equals $y_i$ away
from $p$. The supplied class representatives exhaust the double
quotient; hence $a=\gamma y_c\kappa$ for a rational unitary
$\gamma$ and $\kappa\in U$. This is \texttt{PtEx}.

The same finite-place section also proves the two companion point
properties used by CS. For any unitary adele $b$, divide it by the
section of its components above $p$. The quotient is the identity
above $p$ and equals $b$ elsewhere. If $b$ is in the level away
from $p$, this quotient lies in the full level. Applying this to
the away double-coset decomposition of $y_c$ gives \texttt{PtAll}.
Applying it to $y_i^{-1}\gamma y_i$ for an away stabilizer gives
\texttt{AwayEx}. This uses finite adelic insertion and the supplied
class decompositions; it does not require a new strong-approximation
theorem for rational unitary points.

\paragraph{VE5: the ambiguity is in the kernel of the character.}
Suppose $a,a'$ describe the same $i,\mathbf L$, with bases $g,g'$.
Then $\kappa_0=a^{-1}a'$ is the identity away from $p$. Above $p$,
its first split component in the common frame is $g_v^{-1}g'_v$.
Equality of the lattice spans says that this matrix and its inverse
are integral. The actual hyperspecial description of the level,
proved in LI/SF and CR3, puts $\kappa_0$ in every local level group,
hence in $U$. The good-prime condition excludes the two auxiliary
character primes. Since $\kappa_0$ is the identity there, the
literal level-character support condition gives
\[
 a'=a\kappa_0,\qquad \chi(\kappa_0)=1.
                                                        \tag{VE.10}
\]
This proves \texttt{PtOne}. The equality of local level groups here
is \texttt{Joint0.Uloc\_\allowbreak eq}; it is not inferred from rank equality.

\paragraph{VE6: one polynomial works for every presentation.}
The supplied \texttt{Forms0} also gives an adelic extension
$F_x(a,z)$ of $f_{x,c}(z)$, with
\[
 F_x(\gamma y_c\kappa,z)
       =\chi(\kappa)f_{x,c}(z')\quad(z=\gamma_\infty z').
                                                        \tag{VE.11}
\]
This is \texttt{AdExt}: membership in the actual form space and its
adelic realization are fields of the supplied global form datum.
Choose one basis and one presentation by VE4 and its polynomial
by VE3. Its evaluation is $F_x(a,z)$ for every $z$.
For another basis and presentation, VE5 gives $a'=a\kappa_0$ with
$\chi(\kappa_0)=1$, so $F_x(a',z)=F_x(a,z)$. Applying (VE.11)
to the second presentation proves (VE.1) for the same polynomial.
This proves \texttt{valEx\_\allowbreak of\_\allowbreak rep'} and the actual
\texttt{ex\_\allowbreak C\_\allowbreak blk}. All basis and representative choices disappear
from the required evaluation identity.

\paragraph{VE7: uniqueness and linearity require no extra existence.}
The block-coordinate map from frames onto tuples of invertible
matrices is surjective: the even inverse is $Y=M^{-1}Z$, and the
odd inverse is
$Y=M^{-1}(ZR_k)^{-\mathsf T}$.
The determinant factors in evaluation are nonzero. Consequently a
polynomial with zero frame evaluations vanishes on all invertible
block tuples. To see that it is zero everywhere, let $\mathfrak D$
be the product of their generic determinants. The polynomial
$Q\mathfrak D$ vanishes on every tuple: either every block is
invertible or one factor of $\mathfrak D$ is zero. Polynomial
functions over the infinite field $\mathbb C$ separate polynomials,
and $\mathfrak D\ne0$ by evaluation at the identity tuple.
Thus $Q=0$.

If two polynomials satisfy (VE.1), choose one presentation by VE4
and subtract their evaluations. The preceding argument and the
injectivity of $\jmath$ prove equality over $L$. The unique value
$Q_{x,i,\mathbf L}$ is additive and $L$-linear in $x$, since sums
and scalar multiples satisfy the same identity. After mapping
coefficients by $e:L\hookrightarrow K$, and extending by zero on
non-lattice tuples as in the source, this constructs the literal
semilinear comparison \texttt{cmpL}.

\paragraph{VE8: supply all fields of the concrete comparison.}
For the actual recursion, the input data are
\texttt{Forms0}, the common standard $B$, its matching local frames,
\texttt{Joint0}, the coefficient-field level character, and LI's
original position/valuation compatibility. VE1--VE2 supply block
and representative rationality; VE4--VE5 give the point properties;
VE6--VE7 give existence, uniqueness and evaluation injectivity.
LI8 supplies the matching inverse away character. CS1--CS4 now apply
with a \emph{proved} \texttt{ValEx}, so supply \texttt{ValIntRec};
CR and CS5--CS9 give the exact integral span. Hence every field of
\texttt{CmpHyp0P} is supplied:
\[
 \texttt{ex},\ \texttt{evInj},\ \texttt{ptEx},\
 \texttt{ptAll},\ \texttt{chi},\ \texttt{latM}.
                                                        \tag{VE.12}
\]
In particular this reconstructs \texttt{comparison\_\allowbreak statements}
and \texttt{exists\_\allowbreak isConcrete}. The temporary use of CS under a
\texttt{ValEx} premise is now discharged without circularity: VE1--VE7
do not use CS's span or integrality conclusions, only its independent
integral duality CS2.

\paragraph{VE9: the exact conditional existence theorem.}
Suppose a supplied $D7$ has level data $L_0$, forms $T_u$, a common
standard-frame datum, \texttt{Joint0} for each $u$, coefficient-field
level-zero characters, and companion ideals
\[
 L_{b,u}+\mathfrak a_u L_{b,u}^{\#}=\Lambda_0.
                                                        \tag{VE.13}
\]
For every supplied scale $\mathrm{sc}$, construct the local
identifications and their original coefficient embeddings by LI;
construct the inverse away characters by LI8; replace only the frames
using SF and (VE.13); and supply the ranks by CR. The resulting
record uses $C_e=1$, exactly as in
\texttt{exists\_\allowbreak concrete\_\allowbreak of\_\allowbreak record}. Its forms are transported
unchanged to the recursion, so the joint level and character
identities are retained. VE8 then gives a concrete comparison $c$.
For every good place its actual image satisfies
\[
 \boxed{\quad\mathcal O\text{-span}\,c(\mathcal L)=M.\quad}
                                                        \tag{VE.14}
\]
This closes the former extra block-rationality and polynomial-value
premises. The remaining construction problem is to supply the
global forms, level/character and standard companion data from
the same true-column geometry. The chart and form constructions
in FC and CV address that problem. None of
(VE.12)--(VE.14) constructs spectral estimates, metric transport,
cohomological windows, or a global prime-sum estimate.

\paragraph{VE10: the actual true-column comparison follows from the direct form construction.}
Take the actual true-column record $I:\texttt{TC}(B,X)$ with its
specified witnesses $I.\texttt{Ok}$, $\texttt{Pins}$,
$l:I.\texttt{Laws}$ and $Z:\texttt{Later}(B,I.\texttt{cols}(l))$.
Put $D=I.\texttt{rec7}\,l\,Z$. Retain the explicit geometric
hypothesis that the valid framed decorated points used in the chart
and neighbor constructions have unitary finite-adelic realizers.
Then for every supplied scale the actual concrete comparison exists,
with $C_e=1$ and with its original coefficient embeddings and positions.
More precisely, the output is the recursion
$\operatorname{recC}(D,\mathrm{sc},1,L_0,I_{\rm loc},\psi,h_d)$
constructed by \texttt{exists\_\allowbreak concrete\_\allowbreak of\_\allowbreak record}, with its transported
forms and concrete comparison; it is not a change to the definition of $D$.

Indeed FC19 constructs $T_u:\texttt{Forms0}\,D\,u$ for every
word, using the actual CV values, NB up/down naturality and FC residue,
norm, distinguished-block and dilation fields. It simultaneously
constructs $L_0(u)$ with the exact five \texttt{Joint0} carrier
equalities. Its \texttt{StdTied} data use the same standard coordinates;
FC16 gives the coefficient-field level-zero character values; and
the companion ideal is the explicit $(\bar\pi_u)$, with the exact
identity (FC.41). These are every global-record premise of VE9.
Apply VE9 to these objects, rather than to a separately supplied
form record. At each good place $w$ the resulting map therefore obeys
\[
 \mathcal O_w\text{-span}\,
          c_{u,w}(\mathcal L_u)=M_{u,w},\qquad C_e=1. \tag{VE.15}
\]
Here $\mathcal L_u$ is the actual global forms lattice and $M_{u,w}$
is the actual target coefficient lattice of the local model. The
assertion is an integral span equality; no index is suppressed.
Existence, uniqueness, integrality and the rank identification are
the already proved VE and CS/CR statements on these same objects.

In this intermediate formulation the initial realizer is an explicit
geometric supplier, constructed in DR and applied in VE11 below. AD1
constructs its diagonal base point, AD4 supplies global-lattice
recovery and exact frame correction, and AD3 supplies its supported
modifications; none of these manufactures the initial decorated
realizer. The full terminal \texttt{FormsTiedSym}/\texttt{ColsTiedSym}
construction is a separate remaining task and is not used in this
direct level-zero comparison. The true-column arithmetic laws and
the later metric, spectral, cohomological and prime-window arguments
retain their stated boundaries.

\paragraph{VE11: discharge the initial-realizer premise on the actual column data.}
Keep exactly the true-column data and scale in VE10. Its CM-field
record includes the literal unramified condition (DR.1), and its
valid framed decorations have the self-duality, nonzero type
functionals, auxiliary residue bases and rational local frames used
in DR3. The type, auxiliary and twist prime sets satisfy (DR.12):
the type primes divide the selected $T_u$, the others do not;
the split-level data orient the auxiliary/twist union; and
$(\pi_u,\bar\pi_u)=\mathcal O$ orients the type primes.
The disjointness checks are the ones in
\texttt{exists\_\allowbreak ptOf\_\allowbreak twists}. Thus DR3 supplies a unitary finite-adelic
realizer for every valid framed decorated point occurring in the
charts and neighbor moves.

There is no stronger rational-genus hypothesis at this step.
DR2 proves the source's actual adelic \texttt{Col.UnimodGenus}
by constructing the matrix and its carried lattice. DR3 installs
all the residue and frame clauses of \texttt{IsPtOf}. The rational
semilocal predicates with similar genus names are bypassed by this
construction. In particular no rational unitary approximation or
global integral-basis assumption is inserted.

Apply VE10 with these supplied realizers. For every supplied scale
$\mathrm{sc}$ there exist the required local identifications, levels,
characters and rank data for
\[
 \operatorname{recC}(I.\texttt{rec7}\,l\,Z,
                  \mathrm{sc},1,L_0,I_{\rm loc},\psi,h_d),
                                                        \tag{VE.16}
\]
and its concrete comparison satisfies (VE.15) at every good place.
The form spaces, coefficient embeddings, odd-position involutions
and normalized valuations are precisely the ones retained in VE10.
This removes the independent decorated-point-realization premise
from the actual integral comparison theorem.

The dependency is one-way: DR1 uses local unramified Hermitian algebra;
DR2 uses only local duality, lattice recovery and finite adelic assembly;
DR3 uses exact frame modification. None uses the value polynomials,
comparison maps or their integral-span conclusion. The supplied
arithmetic host and actual $I.\texttt{Ok}/\texttt{Pins}/
\texttt{Laws}/\texttt{Later}$ remain hypotheses. The full terminal
symmetric-column, metric, spectral, cohomological and prime-window
constructions remain the separate tasks already identified.


\section{Canonical refill and the conditional arithmetic point}
\label{sec:canonical-refill}

\status{RF1--RF8 construct the literal canonical refill, complete the
\texttt{exists\_\allowbreak ledPart} implication using LP, and install the resulting
ideal-scaled lattice and height in the point. RF9 proves the two scale
comparisons needed by that construction. RF10 proves the uniform point join
from the explicitly retained recursive, host-column, level and initial-size
suppliers. The geometric \texttt{Lay}/\texttt{RecOK}/\texttt{LedOK} package,
the analytic supply of level primes and the stronger \texttt{NumAt} record
are not inferred from the refill. RF11--RF15 additionally prove the
quantitative rank, thin-size, gap and band estimates and construct every
field of the block-length record \texttt{ParamsT}; RF16 proves the
actual upper bound on word cardinality. The complete
\texttt{ParameterChoice'} and \texttt{SizesAt} suppliers remain explicit.}

Fix the actual datum $D$, a size datum $L_0$, the local constants $K$,
and an error datum $\mathcal E$. Write $p_w$ and $e_w>0$ for the residue
characteristic and ramification index of a finite place. The set of words
$u$ is finite. Here $m=D.R.m$ is the global chain length, $k=D.R.k$,
$n=|A|$, and $J=\{1,\ldots,m\}$. Put
\[
 b=e^{x^\epsilon},\quad a=xb,\quad X'=e^{x^{\epsilon/3}},\quad
 b'=a/X',\quad w''=b'/(xn),\quad X''=a/w''=X'xn.       \tag{RF.1}
\]
The top degree functions $a_T,b_T$, the weights $w_j(\tau)$, and all
geometric and representation data are the supplied ones. In particular,
this construction will retain $L_0$'s scale and its exact selected degrees.

\paragraph{RF1: choose the same representative throughout an orbit.}
The finite abelian group $A$ comes with its source enumeration
$\operatorname{ofNat}$ and integer labels. For $h,t\in A$, define
$\operatorname{bs}(h,t)$ to be $\operatorname{ofNat}(v)$ for the least
nonnegative $v$ such that
\[
                  \operatorname{ofNat}(v)+ih=t
                  \quad\hbox{for some }i\geq0.                 \tag{RF.2}
\]
Such a $v$ exists by taking the label of $t$. The same candidate labels
work for $t$ and $t+h$: add one to $i$ in one direction, and add
$\operatorname{ord}(h)-1$ in the other. Thus
\[
 \operatorname{bs}(h,t)+ih=t\ \hbox{for some }i\geq0,
 \qquad\operatorname{bs}(h,t+h)=\operatorname{bs}(h,t).          \tag{RF.3}
\]
This is the literal \texttt{Arith.bsH}; its choice uses the given labels,
so it remains compatible when points are compared in those labels.

\paragraph{RF2: fill the actual weights, including the unused indices.}
Label the archimedean positions by $z=(\delta,t)\in\{0,1\}\times A$,
using the original bijection $\tau(z)$. A top is $(0,t)$ with
$t\in TT$; a thick position has first coordinate $1$ and second
coordinate in the translated second set; all other positions are thin.
These are the actual top and thick sets: injectivity of the position map
shows that their label descriptions agree with the geometric descriptions.
Extend $b_T$ by zero away from $TT$ and set
\[
 o_j(z)=
 \begin{cases}
 w_j(\tau(z))-\mathbf1_{\mathrm{top}}(z)b_T(t),&j\in J,\\
 b'/x,&j\notin J,\ z\text{ top},\\
 b,&j\notin J,\ z\text{ thick},\\
 w'',&j\notin J,\ z\text{ thin}.
 \end{cases}                                                  \tag{RF.4}
\]
The three supplied weight windows give, respectively,
\[
 \frac{b'}x\leq o_j(z)\leq\frac{b'}x+\eta,\qquad
 b\leq o_j(z)\leq b+\eta+1,\qquad
 w''\leq o_j(z)\leq w''+\eta.                                 \tag{RF.5}
\]
They apply for every $j$, because (RF.4) gives their lower endpoints
outside $J$. Define
\[
 \operatorname{fill}_j(z)=\mathbf1_{\mathrm{top}}(z)b_T(t)+o_j(z).
\]
Then $\operatorname{fill}_j(z)=w_j(\tau(z))$ for $j\in J$.
If $w''\geq1$, every fill is at least $1$: thin positions use $w''$,
thick positions use $b\geq1$, and tops use $b_T\geq b'>0$ and
$b'/x=nw''\geq1$.

Use (RF.3) for the orbit selector, (RF.4) for its offsets, and set
\[
 e_p=\left\lceil C_e f k^2\log_p a\right\rceil_+,
 \qquad \kappa_I=\inf\{\kappa:\kappa\text{ is an admissible
                            universal Type-I loss}\}.         \tag{RF.6}
\]
The admissible set is the actual \texttt{kapISet}; its nonemptiness and
infimum attainment were proved in GJ1. Defining (RF.6) gives
\texttt{EpDef} and \texttt{KapIDef}; it does not by itself identify $e_p$
with the geometric stabilizer sandwich exponent. That distinct assertion
is \texttt{EpTie}. Equations (RF.3)--(RF.4) prove \texttt{CanonInst},
including its prescribed first-coordinate-zero offsets outside $J$.

\paragraph{RF3: one finite set supports every numerical field.}
Assume that the twist set
$S_{\rm tw}=\{w:\mu_w\ne\mu_0\}$ is finite and that the positive
integers $D_{\Gamma,u}$ are supplied. Let $T_u$ denote the actual type
prime of word $u$, and choose
\[
 B_S=\lceil X''\rceil_++g_0+g_1+d+
                           \sum_u(T_u+D_{\Gamma,u}),\qquad
 S=\{w:p_w\leq B_S\}\cup S_{\rm tw}.                          \tag{RF.7}
\]
This set is finite. Indeed, a place over a rational prime $p$ determines
a prime ideal dividing the nonzero principal ideal $(p)$ in the number
ring; unique ideal factorization makes that set finite. There are only
finitely many $p\leq B_S$.
If $w\notin S$, then
\[
 \begin{gathered}
 p_w>X'',\quad p_w\notin\{g_0,g_1,d\},\quad
 p_w\ne T_u,\quad p_w\nmid D_{\Gamma,u}\quad(\text{every }u),\\
                         \mu_w=\mu_0.                         \tag{RF.8}
 \end{gathered}
\]
For the divisibility assertion use $D_{\Gamma,u}>0$: a positive integer
divisible by $p_w$ is at least $p_w$, contradicting $p_w>B_S$.
In particular $S$ contains every place of residue characteristic at most
$X''$, as well as every nontrivially twisted place.

\paragraph{RF4: the literal gain, forcing, error and deficit values.}
We spell out the outer cases, since they are what makes the refill
possible. Let $\mathcal R_p$ be the orbit representatives from RF1 and
$\mathcal V_p=\texttt{VV}(p)$ the actual served representatives. Write
$g^0_{jv},g^U_{jv},P_{jv}$ for the source local gain scores and correction,
$t_v,D_p$ for the local depth and degree budgets, and
$\Sigma_p=m\sum_{v\in\mathrm{goodReps}(p)}\kappa_v$.
All are evaluated at the same selected $a_T,b_T$, scale and instance.
The explicit formulas for these unchanged scalar scores are recorded in
the expanded report. Set
\[
 S_v=\left\lceil\sum_{j\in J}
        \max\{0,g^U_{jv}-P_{jv}-\max(g^0_{jv},0)\}\right\rceil,
                                                               \tag{RF.9}
\]
and define
\[
 \begin{split}
 \tau_p&=\begin{cases}0,&\mathcal V_p=\varnothing,\\
 e_p+m\sum_{v\in\mathcal V_p}(\lceil t_v\rceil_++D_p+1),
                                      &\mathcal V_p\ne\varnothing,
                         \end{cases}\\
 \iota_p&=\begin{cases}0,&\mathcal V_p=\varnothing,\\
 2\Sigma_p\max_{v\in\mathcal V_p}S_v,&\mathcal V_p\ne\varnothing.
                         \end{cases}                          
 \end{split}
 \tag{RF.10}
\]
The gain is zero when $p>X''$ or when $p$ is of the source inert class
and $h_p=0$. Otherwise it is
\[
 G_{jp}=\sum_{v\in\mathcal V_p}
       \max\{g^U_{jv}-P_{jv},\max(g^0_{jv},0)\}
          +\sum_{v\in\mathcal R_p\setminus\mathcal V_p}
                                      \max(g^0_{jv},0).         \tag{RF.11}
\]
At a place $w$, install $G_{j,p_w}$ only for the inert or signal class,
and install zero for the other classes. In particular it vanishes at
$g_0,g_1$: \texttt{HostArith} places these level primes in the split class.

The forcing value is, in priority order,
\[
 F_{jp}=\begin{cases}
 F^G_j,&p\in\{g_0,g_1\},\\
 F^d_j,&\operatorname{cls}(p)=\mathrm{ramD},\\
 F^{\rm in}_{jp},&p\leq X'',\ p\text{ inert or signal},\\
 0,&\text{otherwise}.
 \end{cases}                                                  \tag{RF.12}
\]
Here $F^G_j=-\operatorname{fockDeg}_j/2$,
$F^d_j=-(k-1)\sum_t\min(w_j(0,t),w_j(1,t))/2$, and
$F^{\rm in}_{jp}$ is the explicit finite orbit/linear-program expression
in the report. In the error formula the own type prime has first priority:
\[
 E_{ujw}=\begin{cases}
 E^T_{uw},&p_w=T_u,\\
 v_{p_w}(D_{\Gamma,u})+E^G_{0,ujw},&p_w\ne T_u,\ p_w=g_0,\\
 v_{p_w}(D_{\Gamma,u})+E^G_{1,ujw},&p_w\ne T_u,g_0,\ p_w=g_1,\\
 v_{p_w}(D_{\Gamma,u})+E^d_j,&\text{next, if }\operatorname{cls}(p_w)=\mathrm{ramD},\\
 v_{p_w}(D_{\Gamma,u})+E^C_{jp_w},&\text{next, if }p_w\leq X'',\ p_w\text{ inert},\ h_{p_w}=0,\\
 v_{p_w}(D_{\Gamma,u})+E^I_{jp_w},&\text{next, if }p_w\leq X'',\ p_w\text{ inert or signal},\\
 v_{p_w}(D_{\Gamma,u}),&\text{otherwise}.
 \end{cases}                                                  \tag{RF.13}
\]
The report records every term of $E^T,E^G,E^d,E^C,E^I$, including the
ramification division in $E^G$ and the odd-order correction in $E^I$.
No sign or size estimate for these quantities is asserted merely by
installing their values.

Every field in (RF.10)--(RF.13) is zero off $S$. For $F$ and $E$, use
(RF.8), the implication $\operatorname{cls}(p)=\mathrm{ramD}\Rightarrow
p=d$, and $v_p(D_{\Gamma,u})=0$ when $p\nmid D_{\Gamma,u}$.
For $G$ use its $p>X''$ cutoff. A served prime must satisfy $p<X'$;
since $X'\leq X''$ by (RF.1), no place outside $S$ is served, giving
$\tau=\iota=0$ there. Conjugate places have the same rational residue
characteristic, so (RF.12) also gives $F_{j,\bar w}=F_{jw}$.

\paragraph{RF5: refill the record without changing its parameters.}
Replace the five numerical fields of $L_0$ by RF4 and its support by
$S$. Define
\[
 \begin{split}
 \mathrm{drop}&=\{p:\operatorname{NoGain}(p)\text{ and }p\ne T_u
                                             \text{ for every }u\},\\
 TT_p&=\begin{cases}TT,&p\text{ signal or }p\geq X',\\
 \{t\in TT:\operatorname{bs}(\operatorname{step}_p,t)\in\mathcal V_p\},
                                                       &\text{otherwise},
 \end{cases}\\
 w^{\rm part}_{j,t,p,e}&=\operatorname{fill}_j
             \bigl(\operatorname{partner}(p,t,e)\bigr).        
 \end{split}
 \tag{RF.14}
\]
Store $TT_p$ as the corresponding subset of the subtype $TT$. This
subtype introduces no lost labels because $TT_p\subset TT$; forgetting
the subtype gives exactly $TT_p$. Its full-set clause follows from the
first branch in (RF.14), and $w^{\rm part}\geq1$ follows from RF2.
The set $\mathrm{drop}$ is generally infinite. Finite support is asserted
for the five numerical fields, not for that set.

The remaining essential supplied condition is
\[
 \mu_w=\mu_0\quad\Longrightarrow\quad F_{jw}=0\quad(j\in J).
                                                               \tag{RF.15}
\]
Together with $L_0$'s untwisted vanishing of $A_j(w)$, this proves the
refilled untwisted field. At a place of neither inert nor signal class,
the gain is zero by definition and $\mathcal V_p=\varnothing$ by the
candidate condition; hence both deficits vanish. Type primes are excluded
from $\mathrm{drop}$ by definition.

All other fields are copied. The formulas used in RF1--RF4 depend on
the copied fields and on the fixed $D$, never on the fields being replaced.
Reconstructing the canonical instance over the refilled datum therefore
gives exactly the same selector, offsets, $\kappa_I$ and $e_p$.
The weight identity in RF2, partner identity and subset identity in
(RF.14), and the definitions of gain and deficits prove every clause of
\texttt{ParametersFill}. Its final clause $\mathrm{La}(K,s)=L.\mathrm{LaG}$
is the retained \texttt{Val.La\_\allowbreak eq}. Thus we obtain the literal
\texttt{exists\_\allowbreak refill} conclusion, preserving
\[
           L.sc=L_0.sc,\qquad L.a_T=L_0.a_T,\qquad L.b_T=L_0.b_T,
                                                               \tag{RF.16}
\]
and \texttt{Val}, together with \texttt{CanonInst}, \texttt{EpDef},
\texttt{KapIDef}, \texttt{FDef}, \texttt{EDef}, both support containments,
all-place \texttt{ParametersFill}, and the exact drop-set equivalence.

\paragraph{RF6: a served prime is never discarded.}
Nonempty $\mathcal V_p$ forces the literal candidate conditions
\[
 p\text{ prime},\quad p<X',\quad p\text{ inert},\quad
 x^{\epsilon/4}\leq\operatorname{ord}(h_p),\quad
 p\geq p_0\Longrightarrow p\notin\mathrm{Band},\quad
                         \Sigma_p\leq x^{-\epsilon}.           \tag{RF.17}
\]
The band's power condition is
$|p^i-X'|\leq\rho_B X'$ for some $i\geq1$, with $p_0\leq p<X'$.
Consequently (RF.17) excludes the small-order, power-band and excess-loss
alternatives in \texttt{NoGain}, as well as its $p>X''$ and
non-inert/non-signal alternatives. The only remaining alternative is
\texttt{IsCs}: inert class with nonzero odd-order $h_p$.

Here is the needed group argument. The frame has a homomorphism from
$A$ into a $2$-group, injective on a cyclic subgroup $A_0$, and
$[A:A_0]=o\leq x^{\epsilon/100}$. An odd-order element has zero image
in the $2$-group. Its multiple by $o$ lies in $A_0$ and still has zero
image; injectivity gives $oh_p=0$. Thus its order divides $o$, contrary to
\[
 x^{\epsilon/4}\leq\operatorname{ord}(h_p)\leq o
                 \leq x^{\epsilon/100}<x^{\epsilon/4}.
\]
We have proved
\[
 \mathcal V_p\ne\varnothing\ \Longrightarrow\
 \neg\operatorname{NoGain}(p),\qquad
 p\in\mathrm{drop}\ \Longrightarrow\ \mathcal V_p=\varnothing.
                                                               \tag{RF.18}
\]
This uses the actual small odd-part bound; a numerical lower bound on the
order alone would not imply even order in an arbitrary finite group.

\paragraph{RF7: complete the supplied-data \texttt{LedPart} construction.}
Assume the literal \texttt{Lay}, \texttt{RecOK} and \texttt{LedOK} data,
\texttt{Std}, $c_g\geq10$, an initial $L_0$ with the required
\texttt{Val} and degree identities, \texttt{HostArith}, and $w''\geq1$.
Retain also the source's $X'\leq T_u$ assumption. The relevant content of
\texttt{RecOK} is finite twist support, positive $D_{\Gamma,u}$, fullness
of each actual $M_{\rm ram}(u)$, and the two-sided commensurability at
its own type prime. The quantified \texttt{LedOK} supplies (RF.15) and,
for every canonical instance satisfying the defining numerical and fill
identities, the actual \texttt{LatTied}, \texttt{EpTie} and served-place
containment/commensurability predicates of LP4.

Apply \texttt{LedOK} first to $L_0$ to get (RF.15), then apply RF5.
Because (RF.16) and \texttt{Val} survive, apply \texttt{LedOK} again to
this refilled datum. This second application is essential: one cannot
transfer the geometric local predicates from an unrelated instance.
LP4 now patches the prescribed served and own-type local modules into
one $N_u$ for every word. LP5 installs these modules and
\[
 \operatorname{lenMN}(u,w)
       =\frac{\operatorname{toNat}(\operatorname{lenLoc}_w
                          (N_u,M_{\rm ram}(u)))}{e_w}.          \tag{RF.19}
\]
At a non-type place where $p\in\mathrm{drop}$, (RF.18) makes the place
unserved, so LP gives $N_{u,w}=M_{{\rm ram},u,w}$. The same is true
if $p$ is of neither inert nor signal class. At every non-type place
LP gives the required local containment. These are the first two clauses
of \texttt{NuDef}. Its third clause is the exact finite-length equation
$\ell=e_w\operatorname{lenMN}(u,w)$ from (RF.19).
Equality of localized modules makes their localized quotient zero, so
every non-type unserved length is zero. This proves \texttt{TiesAt.lats0}
and \texttt{NuUnserved}, not just the patch's existence.

The served and type identities are precisely \texttt{LatDefs}; all-place
filling and both support containments come from RF5; \texttt{LatTied} and
\texttt{EpTie} come from the second \texttt{LedOK} application. The
supplied level/error data and (RF.13) give \texttt{ErrDefs}; its three
constant inequalities follow from the standard fixed constants and
$c_g\geq10$. We have therefore constructed the entire literal
\texttt{LedPart}, with the same scale, selected degrees, \texttt{Val}
and canonical instance. No global containment $N_u\subset M_{\rm ram}(u)$
is added at a type prime, and finiteness of an arbitrary local length is
not asserted by (RF.19).

\paragraph{RF8: install the actual final lattice and height.}
For the constructed $L$ put
\[
 s_{uw}=\sum_{j\in J}\bigl(G^{\rm kept}_{jw}+F_{jw}-E_{ujw}\bigr),
 \qquad
 I_u=\prod_{w\in S}\mathfrak p_w^{-\lfloor e_ws_{uw}\rfloor},
 \qquad N'_u=I_uN_u.                                         \tag{RF.20}
\]
Here $G^{\rm kept}_{jw}=0$ on the drop set and is the installed gain
otherwise. The fractional-ideal product is finite by RF3--RF5; its powers
are integral and may have either sign. In particular the floor is taken
\emph{after} ramification scaling. The product $I_uN_u$ means the integer-ring
span of all $ay$ with $a\in I_u$, $y\in N_u$, exactly the source
\texttt{nPrime}.

Replace the source-lattice field of $D$ by $N'_u$, set its height field to
\[
 h_u=\texttt{hExpl6}(D,L,Y.X,C_{\rm sz},u),
\]
and set its common height allowance to
\[
                  H'=\max\{1,\operatorname{sFrame}(N',h)
                                      +\mu_{\max}(RR)\}.       \tag{RF.21}
\]
The vector spaces, forms, host, weights, local comparisons and local
selected modules are retained. Hence the defining $N'$ and $h$ fields
now hold as identities. The numerical functions $\mathcal V_p,G,\tau,
\iota,TT_p$ and the local module $NP$ are unchanged, since none reads the
replaced output-lattice or height field. The auxiliary row-side maps are
unchanged as well: induction on their defining chain leaves each map,
its image and its adjoint composition identical. Thus every supplied
comparison, kernel and auxiliary datum transports to the new record.
This constructs the literal point \texttt{LedPart.pt} and preserves
its scale, degree identities, \texttt{Val} and \texttt{CanonInst}.
If the supplied recursive datum has the tame, Minkowski and local
quadratic-twisting properties, these transport too. Installing their
records does not prove their geometric or metric premises.

\paragraph{RF9: the two eventual scale inequalities have elementary proofs.}
For fixed $\epsilon>0$, \texttt{RankVal} gives
$n\leq2x^{2+8\epsilon}\leq x^3$ when $x\geq4$, because
$\epsilon<1/52$. Also
\[
                         w''=\frac{e^{x^\epsilon-x^{\epsilon/3}}}{n}.
\]
For $x^{2\epsilon/3}\geq2$ and $\log x\geq12/\epsilon^2$,
$x^\epsilon-x^{\epsilon/3}\geq x^\epsilon/2\geq3\log x$.
The last inequality follows from
$e^{\epsilon\log x}\geq(\epsilon\log x)^2/2$.
Thus $w''\geq x^3/n\geq1$, with an explicit threshold depending only
on $\epsilon$.
Next, if $c_s>0$, $m\geq1$, $T>0$, and
$\log T\geq c_smx^\epsilon$, then
\[
 x^{2\epsilon/3}\geq c_s^{-1}
       \quad\Longrightarrow\quad \log T\geq x^{\epsilon/3}
       \quad\Longrightarrow\quad X'\leq T.                   \tag{RF.22}
\]
These prove the two threshold lemmas used in the point join without
changing any previously fixed local constant.

\paragraph{RF10: the uniform point join and its exact remaining suppliers.}
The recursive premise \texttt{RecExP} has a useful order of choices.
After the external constants it supplies $c_s>0$ and a threshold.
For each sufficiently large admissible column datum it chooses positive
$d_u,D_*$ and lower bounds on the logarithms of the desired type-prime
centers. It then supplies $D,Y$, \texttt{Over}, \texttt{RecOK},
\texttt{LedOK}, and the twisting data \emph{for every} level record
with those $d_u,D_*$ and the prescribed $\Gamma$. Its requested property
$Q$ holds for every compatible completed \texttt{LedPart}.
The universal quantifier over the level record permits the level choice
to be made after $d_u,D_*$ have been fixed.

Retain the following three substantive suppliers with their literal
uniform hypotheses: \texttt{chGood\_\allowbreak of\_\allowbreak host}, which constructs the
compatible weight columns from the host and selected degrees;
\texttt{exists\_\allowbreak levCh\_\allowbreak of\_\allowbreak levelPrime}, which consumes
\texttt{LevelPrime} and chooses the prescribed norm/congruence primes;
and \texttt{exists\_\allowbreak sizeData\_\allowbreak over}, which constructs $L_0$ and the
same-host/word-label identifications from the recursive datum. These
are not consequences of the refill. With them, the point join is now
explicit. Enlarge the middle constant thresholds to include $c_g\geq10$,
replace the lower $c_D$ threshold by its maximum with $0$, and the lower
$c_h$ threshold by its maximum with $402$. Take the maximum of the
recursive, three supplier and two RF9 thresholds. For each admissible
host, selected degree pair, word code and distinguished position:
construct the columns; obtain $d_u,D_*$ recursively; choose the levels;
apply the universal recursive assertion; obtain $L_0$; apply RF7; and
install the point by RF8. The recursive property $Q$ applies to this
very \texttt{LedPart}, and the retained tame, Minkowski and twisting
properties give the accompanying duties.

The threshold order is exactly
\[
 \epsilon_0;\quad(\epsilon,C_8,C_5,c_6,C_b);\quad K_0;\quad C;
 \quad c_{D,0};\quad c_D;\quad c_{h,0};\quad c_h;\quad x_0;
 \quad(s,R,H,q_0,a_T,b_T,\mathrm{code},u_0).                    \tag{RF.23}
\]
No threshold may depend on the final degree pair or distinguished
position. The resulting point retains the original host field, top and
second sets, primes, Frobenius labels, word code and $u_0$ through the
supplied \texttt{Over}/initial-size identifications. This is the
\texttt{objCP\_\allowbreak and\_\allowbreak of\_\allowbreak recExP} join, with its three construction
suppliers still explicit. It does not assert \texttt{RecExP},
\texttt{LevelPrime}, or \texttt{NumAt}. The exact next-source list
records those remaining constructions rather than treating the record
type \texttt{ObjCP} as a proof of their existence.

\subsection{Actual numerical estimates beyond the refill}

These estimates have the precise \texttt{SzRow} order of quantifiers:
for each $0<\epsilon<1/52$ and fixed standard constant record $C$, choose
$c_{D,0}$; for every $c_D\geq c_{D,0}$ choose $c_{h,0}$; for every
$c_h\geq c_{h,0}$ choose $x_0$; then the conclusions hold for every
supplied $D,L$ with $\operatorname{Val}(C,c_D,c_h,L)$,
$L.\mathrm{sc}.\epsilon=\epsilon$, and $x\geq x_0$.
There is no choice of a different $D$ or $L$ in these assertions.
We write $\mathcal L=x^\epsilon$, $\delta=x^{-\epsilon/8}$,
$r=k+m-1$, $f=2n$, and $N_m=k+m(t+1)$.
All terminating decimals below denote exact rational numbers.

\paragraph{RF11: an explicit comparison principle for all the scales.}
If $a<b$, $p$ is a nonnegative integer and $A$ is fixed, then
\[
             A x^a(\log x)^p=o(x^b)\quad(x\longrightarrow\infty).
                                                               \tag{RF.24}
\]
Indeed, with $y=\log x\geq1$ and $d=b-a>0$, the positive series for
$e^{dy}$ gives
\[
 \frac{|A|y^p}{e^{dy}}
       \leq \frac{|A|(p+1)!}{d^{p+1}y}\longrightarrow0.
\]
This supplies an actual threshold for any requested positive upper
constant. Taking the maximum of three thresholds proves, for
$0\leq a<b$ and $\kappa>0$,
\[
 c\log x+A x^a(\log x)^p+A\leq\kappa x^b                 \tag{RF.25}
\]
eventually, with absolute values used for negative coefficients.
Consequently, for any fixed real $c$ and $\kappa>0$, all the following
hold at one threshold depending only on $\epsilon,c,\kappa$:
\begin{gather*}
 x^c\leq e^{x^{\epsilon/8}},\qquad
 x^c\leq (X')^\kappa,\qquad x^cX'\leq b^\kappa,
 \qquad b'x^c\leq b,\\
 b^{1-\kappa}x^c\leq b',\qquad
 b^{1-\kappa}x^c\leq w''\leq b'
       \quad(1\leq n\leq x^3),\\
 e^{c x^{\epsilon/5}(\log x)^2}x^c\leq(X')^\kappa.
                                                               \tag{RF.26}
\end{gather*}
For example the lower bound for $w''$ reduces, using
$\log n\leq3\log x$, to
$(c+3)\log x+x^{\epsilon/3}\leq\kappa x^\epsilon$.
The other inequalities follow by taking logarithms in (RF.1).
Thus this comparison principle proves the actual scale-polynomial
supplier used below, including its uniformity in $n$.

\paragraph{RF12: ranks, the large order, and the seven rank-parameter bounds.}
Put
\[
 T=\lceil x^{377}\rceil_+,
 \quad W=10r\left(f+\frac{f^2(k-1)^3\mathcal L^2}{m}\right),
 \quad n_W=8\lceil W/8\rceil_+ .                            \tag{RF.27}
\]
The supplied rank condition says $r+t=8fr+4+Tn_W$.
For $x\geq1000^{1000}$, constants at most $1000$ can be absorbed into
an additional exponent $0.001$: if $A\leq1000$ and $u+0.001\leq v$,
then $Ax^u\leq x^v$. The three defining rank windows therefore give
\[
\begin{array}{c|ccccccc}
\text{quantity}&f&r&k&\mathcal L&m&m^{-1}&f\binom r2\\ \hline
\text{upper power of }x&2.156&4.694&4.694&0.020&0.501&-0.498&11.544
\end{array}                                                   \tag{RF.28}
\]
For instance $r\leq2k\leq4x^{9/2+10\epsilon}$ and
$m\geq x^{1/2-\epsilon/16}\geq x^{0.498}$; also
$f\binom r2\leq fr^2$.
The second term of $W$ has exponent at most
$4.694+2(2.156)+3(4.694)+2(0.020)-0.498=22.630$.
The first has exponent at most $6.850$. Since $n_W<W+8$,
\[
 n_W\leq28x^{22.7}\leq x^{22.9},\quad
 x^{377}\leq T\leq x^{378},\quad
 x^{377}\leq t\leq13x^{400.9}\leq x^{401}.                  \tag{RF.29}
\]
The lower bound for $t$ uses $n_W\geq1$ and $8fr\geq r$.
Finally $N_m=k+m(t+1)\leq3x^{401.6}\leq x^{402}$.
This proves all seventeen conclusions of the source's
\texttt{rank\_\allowbreak basic\_\allowbreak bounds}: in addition to (RF.28)--(RF.29),
\[
\begin{gathered}
 f\leq x^3,\quad m\leq x,\quad r\leq x^5,\quad k\leq r\leq2k,
 \quad1\leq\mathcal L\leq x,\\
 x^{6.5}\leq fr\mathcal L\leq x^{6.87},\quad
 f\binom r2\leq x^{12},\quad kf/m\leq x^7,\\
 x^{1/2-\epsilon/16}\leq m\leq2x^{1/2-\epsilon/16}.
\end{gathered}                                                 \tag{RF.30}
\]
The lower bound for $fr\mathcal L$ follows from
$f\geq x^2$, $r\geq k\geq x^{9/2}$ and $\mathcal L\geq1$;
the upper bound and $kf/m$ follow by adding the exponents in (RF.28).

Let $R_0=nx^{-\epsilon}$ and let $C_k,C_6,C_f$ be the actual fixed
constants. At a further threshold the seven required inequalities are
\[
\begin{gathered}
 C_km^5f\mathcal L x^\epsilon\leq k,\qquad
 k\leq x^{1/2}nR_0,\qquad \sqrt{k/m}\leq R_0,\\
 m^2\mathcal L\leq k,\quad mx^{1/2}\log x\leq k,\quad
 C_6m^6\mathcal L\leq k\log k,\quad C_f\mathcal L\leq f\log k.
\end{gathered}                                                 \tag{RF.31}
\]
For the first three, the ratios of the left sides to the available
lower bounds on the right sides are bounded respectively by
$128|C_k|x^{-5\epsilon/16}$,
$2x^{-5\epsilon}$, and, after squaring the third inequality,
$2x^{-63\epsilon/16}$. They tend to zero.
For the last four put $y=x^{1/2}$ and require
$y\geq\max(8,64C_6,C_f)$. Then
$m\leq2y$, $k\geq y^9$, $f\geq y^4$,
$\mathcal L\leq y$, $0\leq\log x\leq y^2$, and $\log k\geq1$.
The four left sides are at most $4y^3$, $2y^4$, $y^8$, and $y^2$,
respectively. This proves the literal rank-parameter families, not only
their eventual conjunction.

\paragraph{RF13: thin-size domination and every prescribed column gap.}
From (RF.26) with $c=5|C_w|$ and $\kappa=1/10$, and $1\leq k\leq x^5$,
\[
                    k^{C_w}\leq w'',\qquad b^{9/10}\leq w''.
                                                               \tag{RF.32}
\]
Indeed both quantities are at most the product
$b^{9/10}x^{5|C_w|}$, which is at most $w''$.
This retains the source's fixed value $c_w=9/10$.

For the gap conditions let
$\Delta_{\mathrm{cl}}=\lceil x^{c_D}\rceil_+$,
$H_0=2\lceil x^{c_h}\rceil_+$ and
$G=C_{\mathrm{gap}}km^4+N_m m^{10}$.
Take $c_D,c_h\geq430$, $x\geq\max(2,C_{\mathrm{gap}})$,
and enlarge the scale threshold until $b,b',w''\geq x^{430}$.
By (RF.30),
$G\leq x^{10}+x^{412}<x^{430}$.
The actual \texttt{HCVal} columns are
\[
 v_i=(k-2-i)4\Delta_{\mathrm{cl}},\qquad b_j=b_1+jH_0,
       \qquad e_j=b_j+H_0/2.                                 \tag{RF.33}
\]
At a top, $b_1=b_T+\lfloor b'/x\rfloor\geq b'$; at a thick
position $b_1\geq\lceil b\rceil$; at a thin position
$b_1=\lceil w''\rceil$. Therefore $b_1\geq x^{430}$ in all cases.
The three gaps are exactly $4\Delta_{\mathrm{cl}}$, $v_{k-2}+b_1=b_1$, and $H_0$.
Each is strictly greater than $G$. This proves all three clauses of
\texttt{szRow\_\allowbreak sig\_\allowbreak gaps} with the actual columns.

\paragraph{RF14: the full band, lambda and host-prime bounds.}
Here the integer orbit-count construction is also expanded. Set
$h=2+8\epsilon$, $r_G=\lceil8(h+2)/\epsilon\rceil_++1$,
$q=6r_G+1$, and $\ell=\log x\geq1$. The source definitions are
\[
 U=\lceil1/\delta\rceil_+,\quad
 M=\lceil(h+3)10^q\ell\rceil_+,\quad
 G'=(M+5)qU,
 \quad N_1=\left\lceil\max(6eG'\delta,(h+3)\log_2x)\right\rceil_++1.
                                                               \tag{RF.34}
\]
Since $\delta U\leq2$,
$M+5\leq((h+3)10^q+6)\ell$, and $\log_2x\leq2\ell$, the
explicit constant
\[
 B_N=6e\bigl((h+3)10^q+6\bigr)2q+2(h+3)+2
\]
satisfies $1\leq N_1\leq B_N\log x$.
For $K_{\Lambda}\geq1$, put $B=K_{\Lambda}B_N^2$. The actual
$\Lambda=K_{\Lambda}N_1^2\log a$ consequently satisfies
\[
              x^\epsilon\leq\Lambda
                   \leq2B x^\epsilon(\log x)^2              \tag{RF.35}
\]
once $\log x\leq x^\epsilon$.

For any fixed real $c$, set
\[
 \rho(c)=\frac{8n\delta x^\epsilon\Lambda
                         +c\log x\,\Lambda}{k}.
\]
The rank windows, $\delta\leq1$ and (RF.35) bound its two numerators by
$32B x^{2+10\epsilon}(\log x)^2$ and
$2|c|B x^\epsilon(\log x)^3$.
By (RF.24), each is at most
$\tfrac12x^{5/2+10\epsilon}$ at a uniform threshold.
Dividing by $k\geq x^{9/2+10\epsilon}$ gives
\[
                         \rho(c)\leq x^{-2}.                \tag{RF.36}
\]
The choice $c=K_\rho K_N$ is exactly $\rho_B$; the choice
$c=C_\rho$ is exactly the parameter condition $\rho_P$.
Together with the supplied host top-count and $b_T$ bounds, this proves
the whole top/band/bottom parameter family as well.

Write $Q_c=QS\,q_0d$ for the host conductor product, using $q_0$ to
distinguish the host prime from the integer $q$ in (RF.34).
In this numerical argument $C_8=C.P.C_{p\mathrm{de}}$ is the constant
in the supplied \texttt{HostVal}.
The supplied \texttt{PairRow} at some $Y\leq x$ gives
\[
 \log Q_c\leq2x^{\epsilon/5}(\log x)^2,\qquad
 \log p_{\mathrm{de}},\ \log p'_{\mathrm{de}}
            \leq2C_8x^{\epsilon/5}(\log x)^2.
                                                               \tag{RF.37}
\]
Here the second statement follows from the actual inequalities
$p_{\mathrm{de}},p'_{\mathrm{de}}\leq Q_c^{C_8}$.
RF11 makes the last bound strictly smaller than $x^{\epsilon/3}$;
hence both auxiliary primes are strictly below $X'$.
For the already supplied level primes set $Y_g=(kb)^8$.
The host condition says $Y_g\leq g_i\leq2Y_g$. Since $k\leq x^5$,
RF11 gives
\[
 g_i\leq e^{9\mathcal L},\qquad \log Y_g\leq9\mathcal L,
 \qquad N_m^4Q_c^{C_Y}\leq Y_g^{1/4}.                       \tag{RF.38}
\]
For the last inequality take logarithms and use
$1608\log x+2C_Yx^{\epsilon/5}(\log x)^2\leq2x^\epsilon$;
the right logarithm is $2\log k+2x^\epsilon$.
These are the complete type-prime scale inequalities with the fixed
$C_{YL}=9$. They concern the supplied $g_i$ and do not choose any prime.

Finally, the lower and upper top-count bounds are precisely
\texttt{ChTT\_\allowbreak at}. The identity
$\#\mathrm{Word}=\lceil3\Gamma_L^2x^{C_F}\rceil_+$ implies
\texttt{ChCode\_\allowbreak at} with constants $(3,C_F)$ by the defining ceiling
inequality; it is a lower bound, not the separate exponential upper
bound on the word count. Since $1\leq J\leq C_8$,
$\epsilon/J\geq\epsilon/C_8$. Thus
$x\geq11^{5C_8/\epsilon}$ also gives
$11\leq\lfloor x^{(\epsilon/J)/5}\rfloor_+\leq\mathrm{splitBound}$.
This proves the full source \texttt{szRow\_\allowbreak band\_\allowbreak lambda\_\allowbreak primes}
and the additional split-bound clause, with no host existence inference.

\paragraph{RF15: construct every field of the actual block-length record.}
Suppose a compatible parameter record $P$ is supplied, so that
$P.C_{nW}=C.P.C_{nW}=10$. Set
\[
 \mathsf{TO}=T,\quad \mathsf n_W=n_W,\quad
 \mathsf n_P=8fr+4,\quad Z=fr\mathcal L,\quad
 \mathsf P=C_PZ^{12},\quad \mathsf D=f\binom r2,
 \quad s=r+6.                                                \tag{RF.39}
\]
Copy the seven constants from $C.T$. Standardness fixes
$e_P=12$, $C_p=1$, $c_{TO}=377$, gives $C_P\geq1000$, and gives the
other required lower bounds. Choose $c_D\geq800$, $c_h\geq1000$,
and then $x$ large enough that $C_P\leq x^{0.001}$,
$C_{hT},C_{bN}\leq x$ and RF11--RF14 apply.
Since $Z\leq x^{6.87}$ and $Z\geq x^{6.5}$,
\[
 1000x^{78}\leq\mathsf P\leq x^{82.441}\leq x^{85},
 \qquad \mathsf D\leq x^{11.544}\leq x^{12}.                \tag{RF.40}
\]
The defining second-set density gives
$\delta_2\geq x^{-\epsilon}$, hence
$\delta_2^{-1}\leq x^{0.020}$.
The strongest lower bound on $\mathsf{TO}$ now has an explicit exponent:
\[
 \frac{\mathsf P^4\mathsf D^3kf x}{m\delta_2}
 \leq x^{4(82.441)+3(11.544)+4.694+2.156+1-0.498+0.020}
 =x^{371.768}\leq x^{377}\leq\mathsf{TO}.                  \tag{RF.41}
\]
In particular this uses the exact order scale $377$; replacing the
inputs prematurely by much coarser bounds would not prove the claim.

The count interval is $W\leq\mathsf n_W<W+8$ and
$8\mid\mathsf n_W$. Since $m\leq k-1$, $f\geq2$ and
$\mathcal L\geq1$,
$W\geq10fr(1+(k-1)^2)\geq fr(r+1)$.
Also $(k-1)f+2r\leq2fr\leq W$ and $8fr+4\leq10fr\leq W$.
These prove all three lower-count conditions.
Moreover
$(2fm)^4\mathsf n_W\leq16x^{33.528}\leq\mathsf P$ and
$\mathsf n_W\leq\mathsf P$.
The exact splitting identity is the supplied rank identity in (RF.27).
Because $k,m$ are even, $r$ and $s=r+6$ are odd, and
$r+5<s\leq r+7$.

For completeness the remaining polynomial tests admit the following
explicit upper bounds. Increase $x_0$ to absorb the fixed coefficients.
\[
\begin{array}{l|r}
\text{quantity bounded above}&\text{power of }x\\ \hline
(\mathsf D+2)\mathsf n_W+2r&36\\
40(\mathsf D+2)&13\\
11\mathsf D^2&25\\
r+2(2\cdot5^4r)(\mathsf D+1)&18\\
\mathsf D k/m&17
\end{array}
\qquad
\begin{array}{l|r}
\text{quantity bounded above}&\text{power of }x\\ \hline
\mathsf{TO}^2&756\\
(\mathsf{TO}+\mathsf D+2)
 (\mathsf{TO}+\mathsf D+2+2N_m)&782\\
4\mathsf{TO}^2+2t&757\\
\mathsf P\mathsf D f m^2N_m\mathsf{TO}&882\\
\mathsf P(fmk)^3k/m&117\\
C_{hT}\mathsf{TO}N_m f^2m^2k\mathsf D\mathsf P&891
\end{array}                                                   \tag{RF.42}
\]
Every entry of the left table is at most $\mathsf{TO}$.
The first two entries of the right table are at most
$\Delta_{\mathrm{cl}}=\lceil x^{c_D}\rceil_+$; the last four are at most
$h_0=L.\operatorname{etaSz}/m=2\lceil x^{c_h}\rceil_+$.
For the product of two sums, use
$\mathsf{TO}+\mathsf D+2\leq x^{379}$ and the second factor
$\leq x^{403}$. For the first height entry use
$4x^{756}+2x^{401}\leq6x^{756}\leq x^{757}$.
The other height exponents are obtained by multiplying the explicit
bounds $\mathsf P\leq x^{85}$, $\mathsf D\leq x^{12}$,
$f\leq x^3$, $m\leq x$, $k\leq x^5$, $N_m\leq x^{402}$,
$\mathsf{TO}\leq x^{378}$; in the cubic entry use $m\geq1$.

The tests against $b$ and $\sqrt b$ are not additional assumptions:
\[
\begin{array}{l|r}
\mathsf{TO}\mathsf P x^{C_b}&x^{463+C_b}\\
\mathsf{TO}\mathsf P&x^{463}\\
\mathsf P^3\mathsf{TO}^2x^{C_b}&x^{1011+C_b}\\
\mathsf P^3\mathsf{TO}^2&x^{1011}\\
\mathsf{TO}\mathsf D^2\log b\,x^{C_b}&x^{403+C_b}\\
\mathsf{TO}\mathsf D^2\log b&x^{403}\\
C_{bN}rmk\mathsf P^2\mathsf D^3&x^{218}
\end{array}                                                   \tag{RF.43}
\]
RF11 puts the first, third, fifth and seventh rows below $b$, and the
second, fourth and sixth below $b^{1/2}$.
The three polynomial-size bounds are immediate from
$\mathsf P\leq x^{85}$, $\mathsf D\leq x^{12}$,
$\mathsf{TO}\leq x^{378}\leq x^{401}$.
This checks every defining field of \texttt{ParamsT}, including its
constant list and the identity $e_P=C.T.e_P$.
It proves the entire \texttt{szRow\_\allowbreak blockLengthParams} implication,
uniformly for every supplied compatible $P$.

\paragraph{The exact numerical supplier boundary.}
Including RF16 below, seven of the twenty-one parameter families in
\texttt{explicitConstants\_\allowbreak parameterConditions\_\allowbreak large} have now been
proved here: the three rank-$k$ bounds; the four other rank bounds;
top-count/band/bottom bounds; thin-size bounds; all signature gaps;
the type-prime scale bounds; and the upper word-cardinality bound.
RF15 proves the full block-length
supplier, and RF14 proves the full band/lambda/prime supplier.
The fourteen other parameter families and the complete
\texttt{SizesAt} package are still retained, with their exact source
targets recorded in the continuation files.

If those remaining suppliers give a compatible $P$ and
\texttt{SizesAt} for the same $D,L$, RF15 gives \texttt{ParamsT} and
RF14 gives every nontrivial remaining inequality in \texttt{NumAt}.
Its other fields are the literal identities in \texttt{Val}:
the constants, $b_T$, the top ideal $cc_K$, the logarithmic depth,
the common cloud vector, the rank windows, and the formula for $L_w$
using $C_t=401$. Thus \texttt{NumAt}, \texttt{ChTT\_\allowbreak at} and
\texttt{ChCode\_\allowbreak at} follow conditionally on precisely these two
remaining packages. Neither those packages nor the simultaneous
geometric construction is asserted by this numerical extension.

\paragraph{RF16: the literal upper bound on the number of words.}
This last estimate uses the actual exponential value of $\Gamma_L$.
Write $G=C.C_\Gamma\geq1$, $F=C.P.C_F\geq1$, and
$H=m^5f x^\epsilon$. The exact \texttt{Val} identities give
\[
 \Gamma_L=e^{GH},\qquad
 \#\mathrm{Word}=\left\lceil3e^{2GH}x^F\right\rceil_+.
                                                               \tag{RF.44}
\]
The rank windows, with the full saving in the exponent of $m$, imply
\[
 H\leq128x^{9/2+139\epsilon/16},\qquad
 k\geq x^{9/2+160\epsilon/16}.                              \tag{RF.45}
\]
The difference is $21\epsilon/16>0$. Fix a threshold at which
$k\geq6$,
\[
 256Gx^{9/2+139\epsilon/16}
                  \leq\tfrac18x^{9/2+160\epsilon/16},
 \qquad F\log x+\log6
                  \leq\tfrac18x^{9/2+160\epsilon/16}.        \tag{RF.46}
\]
For the first inequality it suffices that
$x^{21\epsilon/16}\geq2048G$; RF11 supplies the second.
Both thresholds depend only on $\epsilon,C$.
Put $A=3e^{2GH}x^F\geq3$. Since
$\lceil A\rceil_+<A+1\leq2A$,
\[
 \log\#\mathrm{Word}
       \leq\log6+2GH+F\log x\leq k/4
       \leq(k-3)\log2.                                     \tag{RF.47}
\]
The last inequality uses $k-3\geq k/2$ and $\log2>1/2$.
Exponentiating, and identifying the natural subtraction with the
integer $k-3$ because $k\geq6$, proves
\[
                         \#\mathrm{Word}\leq2^{k-3}.         \tag{RF.48}
\]
This is the complete \texttt{szRow\_\allowbreak card\_\allowbreak words\_\allowbreak le} family.
It has $c_{D,0}=c_{h,0}=0$, so it combines with the preceding
families without changing their quantifier order. The number of words
is not treated as polynomial: its logarithm is compared to the
larger rank $k$ with the explicit positive gap in (RF.45).


\section{Deriving the column laws and isolating canonical evaluation inputs}
\label{sec:column-laws}

\status{TL1--TL13 reconstruct the actual implication
\texttt{TC.Ok}$\Rightarrow$\texttt{TC.Laws}, including its lattice
finiteness, coefficient, level, and period suppliers. These proofs use the
literal arithmetic and local-index clauses of \texttt{Ok}; they do not
use a previously assembled \texttt{Forms0} or a terminal comparison to
prove its own law hypothesis. TL14--TL21 reconstruct the finite-field,
orientation, and finite-generation parts of the canonical evaluation
constructor, with the remaining local nonvanishing and analytic
algebraicity suppliers identified explicitly. No assertion here produces
the five independent stages of \texttt{Later}, and no Lean build or
kernel verification is asserted.}

\paragraph{Conventions.}
Write $E$ for the fixed number field with involution, $\mathcal O$ for
its integers, and $h(x,y)=\bar x^{\mathsf T}y$ on $E^n$. The field law
used below is
\[
 \tau(\bar a)=\overline{\tau(a)}
 \quad\text{for every embedding }\tau:E\hookrightarrow\mathbb C.
 \tag{TL.1}
\]
For coefficient field $K$, embeddings $e_\tau:E\to K$, and
$\sigma:K\to\mathbb C$, \texttt{FOk} also supplies
$e_\tau(\bar a)=e_\tau(a)^*$ and
$\sigma(b^*)=\overline{\sigma(b)}$. It does not assert the latter
identity for every embedding of an arbitrary larger $K$.
The rank, prime, character, and index data called \texttt{TC.Ok}
throughout this section are the actual source data, including
\texttt{ChIdx.Ok}. Their existence is a separate question, treated
only to the extent stated in TL14--TL21.

\paragraph{TL1: finitely many self-dual Hermitian lattice classes.}
Put $d=[E:\mathbb Q]$, $D=nd$, and $q=|\operatorname{disc}(E)|$.
For a full self-dual $\mathcal O$-lattice $N=N^\#$, its underlying
integer lattice has the positive integral scalar product
\[
 B(x,y)=\operatorname{Tr}_{E/\mathbb Q}h(x,y).
 \tag{TL.2}
\]
Symmetry follows by conjugating the trace, and positivity follows by
writing the trace as the sum of $|\tau(x_i)|^2$. Let $N^*$ denote the
dual for $B$. If $y\in N^*$, then, for all $a\in\mathcal O$ and
$x\in N$, $\operatorname{Tr}(a h(x,y))$ is integral: apply the defining
condition to $\bar a x$. Thus $h(x,y)$ belongs to the codifferent.
The inverse of an integral trace matrix has denominator dividing its
determinant, so $q\mathfrak D^{-1}\subseteq\mathcal O$. Consequently
\[
 N\subseteq N^*,\qquad qN^*\subseteq N,\qquad
 1\le\det B|_N=[N^*:N]\le q^D.                 \tag{TL.3}
\]
In particular nonzero lattice vectors have length at least one and the
covolume is at most $C=q^{D/2}$.

Here is an elementary uniform bound for independent short vectors.
Let $v_D$ be the volume of the Euclidean unit ball and choose
$R\ge1$ with $v_D R>3^{D-1}C$. If the vectors of $N$ of length at
most $2R$ spanned a subspace of dimension $r<D$, then all the lattice
points in any translate of a radius-$R$ ball would lie in one affine
$r$-plane: differences of two such points have length at most $2R$.
Their disjoint radius-$1/2$ balls in that plane give at most
$(2R+1)^r\le3^{D-1}R^{D-1}$ points. On the other hand, averaging the
number of lattice points over a fundamental parallelepiped gives
$v_D R^D/\operatorname{covol}(N)\ge v_D R^D/C$, a contradiction.
There are therefore $D$ independent vectors $v_i$ of length at most
$2R$.

Set $N_0=\sum_i\mathbb Z v_i$. Its integral Gram entries are bounded
by $4R^2$, and $b=[N:N_0]\le(2R)^D$, since
$\operatorname{covol}(N)\ge1$. Moreover
$N_0\subseteq N\subseteq b^{-1}N_0$. For each of the finitely many
Gram matrices and integers $b$, there are finitely many intermediate
groups, as they correspond to subgroups of
$b^{-1}N_0/N_0$.

One must also retain the $\mathcal O$-action, rather than just the
integer quadratic form. Fix an integral basis $a_s$ of $\mathcal O$.
Multiplication by $a_s$ has operator norm at most
$\max_\tau|\tau(a_s)|$ for (TL.2). In the $v_i$ coordinates its
entries have denominator dividing $b$; they are bounded because the
Gram entries are bounded and the inverse Gram matrix is bounded by
the cofactor formula and $\det\operatorname{Gram}(v_i)\ge1$.
There are consequently only finitely many possible multiplication
matrices. Identical Gram, intermediate-group, and multiplication data
give an integer isometry commuting with $\mathcal O$, hence an
$E$-linear isometry. It preserves $h$ because all the identities
$\operatorname{Tr}(a h(x,y))$ agree and the field trace pairing is
nondegenerate. This proves finiteness up to the actual rational
unitary group. For $n=0$ there is just the zero lattice.
Finally the automorphism group of any one lattice is finite: images
of a fixed integer basis lie in finitely many bounded lattice sets.

\paragraph{TL2: decorations, local frames, and coupling cosets preserve
finiteness.}
A valid decoration consists of $N=N^\#$, a nonzero residue-linear
functional at each type prime, and a residue basis at each auxiliary
prime. Values on finitely many $\mathcal O$-generators of $N$
determine these maps on $N$. They determine the localized maps by
clearing a denominator outside the relevant prime; the prescribed
extension is zero outside that localization. All their possible
values lie in finite residue fields. Thus there are finitely many
decorations over each fixed $N$, and TL1 gives finitely many decorated
orbits.

For a fixed decoration, rational frames for $N_P$ form a torsor under
$\operatorname{GL}_n(\mathcal O_{(P)})$. Restricting the permitted
right-frame transformations to a subgroup of finite index leaves
finitely many frame classes. Taking the product over the finitely
many twist primes preserves finiteness. At the coupling step the
additional coordinate is a coset subject to finitely many bounds
$q_P^{a_P}N_P$ modulo $q_P^{b_P}N_P$. Its possible local cosets are
finite: if $a_P<b_P$, this is a quotient of a free local module by
a nonzero power of its maximal ideal, and otherwise there is only
one coset. The map from global coupling cosets to their family of
local cosets is injective by the definition of the simultaneous
local-condition module. Zero supplies a nonempty fiber. At other
steps the coupling coordinate is the fixed empty set.

The elementary orbit argument being used is this: an equivariant
map with finitely many base orbits and finite fibers above chosen
base representatives has finitely many source orbits. Move a source
point over a chosen representative, then choose within that finite
fiber. No finite cardinality of the full rational matrix group is
assumed.

A stabilizer maps injectively into the rational unitary automorphisms
of its decorated lattice. Indeed, once that rational matrix is known,
surjectivity of each auxiliary residue basis determines the auxiliary
matrix, and equality of each rational local frame determines its
right-frame matrix. Coupling coordinates add a condition, not a new
kernel. Stabilizers are therefore finite. Their orders divide the
corresponding lattice-automorphism orders. The finite set of those
orders, over the finite rank set in \texttt{gamRanks}, has a positive
least common multiple \texttt{DGam}.

\paragraph{TL3: the actual twist levels have finite index and finite
characters.}
For any nonzero integer $c$, the quotient
$\mathcal O_{(P)}/c\mathcal O_{(P)}$ is finite: it is a localization
of the finite ring $\mathcal O/c\mathcal O$. A subgroup of
$\operatorname{GL}_n(\mathcal O_{(P)})$ containing the principal
congruence kernel modulo $c$ therefore has finite index.

For the actual intermediate index matrix $h$ at step $j$, the
\texttt{boundQ} clause of \texttt{ChIdx.Ok} says
\[
 h^{-1}\Lambda\subseteq\Lambda,
 \qquad q^j h\Lambda\subseteq\Lambda.             \tag{TL.4}
\]
If $\kappa\equiv1\pmod{q^{j+1}}$ on the first coordinate block,
put $a=h^{-1}\operatorname{diag}(\kappa,1)h$. For $x\in\Lambda$,
the first coordinates of $q^jh x$ are integral, so
$(\kappa-1)\operatorname{ini}(h x)\in q\mathcal O_{(P)}^n$.
The first inclusion in (TL.4), applied to the embedded first block,
gives $(a-1)\Lambda\subseteq q\Lambda$. The same calculation for
$\kappa^{-1}$ proves $a^{-1}\Lambda\subseteq\Lambda$ and hence
$a\Lambda=\Lambda$. Its reduction is the identity on every
specified residue flag. This is precisely membership in the source's
\texttt{Jat}, whose first-block projection defines \texttt{levQ}.
Thus the actual level contains the congruence kernel of depth $j+1$.
At the top step the level is generated by an extended row level and
the congruence kernel of its specified depth plus one, so the same
conclusion holds. The row level itself is a preimage of a subgroup
of the finite residue general linear group.

The level characters are products of finite residue multiplicative
characters, finite residue additive characters, and the prescribed
local additive character on $q^{-D}\mathcal O_{(P)}$ modulo
$\mathcal O_{(P)}$. Every value is a root of unity: these additive
quotients and residue multiplicative groups are finite. Products and
inverses retain that property. The local-index character formula in
\texttt{ChIdx.Ok} identifies these products with the actual level
characters. This argument proves their finite order after the index
data have been supplied; it does not construct those index data.

Combining TL2 and these facts gives all seven clauses of
\texttt{GStep.Ok}: finite orbits, finite stabilizers, unitary rational
matrices, root-of-unity characters, exact transport of $L^\#$, fullness
of that lattice, and stabilizer injectivity. Replacing the preferred
orbit representative changes none of these seven assertions.

\paragraph{TL4: the zero step is free and all its twists have the same
classes.}
At step zero the auxiliary level has unipotent reduction at a prime
above $g_0$, and the base prime bound is
\[
 [E:\mathbb Q]\,k=2fk<g_0-1.                     \tag{TL.5}
\]
Let $\gamma$ stabilize a decorated point. By TL2 it has finite order.
Its prime-to-$g_0$ part is congruent to the identity: a unipotent
matrix in residue characteristic $g_0$ has $g_0$-power order. A
finite-order matrix $u=1+X$ with order $r$ prime to $g_0$ and
$X\in P$ is the identity. Indeed $(1+X)^r=1$ first puts $X$ in
$P^2$, and iteration puts it in every $P^a$, whose intersection is
zero. Hence $\gamma$ has $g_0$-power order. If it were nontrivial,
a power would have order $g_0$ and its rational characteristic
polynomial on the underlying $[E:\mathbb Q]k$-dimensional space
would contain $\Phi_{g_0}$, of degree $g_0-1$. This contradicts
(TL.5). Stabilizer injectivity now kills the auxiliary and local
frame matrices as well. The zero-step action is free.

The actual index at zero is unique by \texttt{zero\_\allowbreak uniq}; its
auxiliary level is the same \texttt{colLev 0} for every twist, its
twist-frame level is the full local general linear group, and its
coupling coordinate is empty. Forgetting those local frames
therefore gives a bijection with the same decorated orbit set for
every twist. To see injectivity directly, any two frames over an
identical local lattice differ by an allowed integral invertible
matrix. Adding frames proves surjectivity. The prescribed preferred
representative minimizes the same decorated-orbit datum, so the
preferred decorations agree literally. This proves the actual
\texttt{cl0} bijections and their decoration equality.
Finally \texttt{tup mu0}=1 and the injective finite-idele reading
identity imply $\xi_{\mu_0}=1$. Every other read character has finite
order, hence has complex absolute value one. The defining products
for $\rho$ and $\rho'$ are therefore of modulus one. TL2's positive
\texttt{DGam}, its divisibility, and the trivial zero stabilizers
give the complete chain laws.

\paragraph{TL5: covariance and isometry of the coefficient family.}
This paragraph starts from a generic \texttt{GStep.Ok}, independently
of the final column-law record. Let $\lambda$ be any supplied
dominant tensor weight, $V_\lambda$ its polynomial coefficient space,
and
\[
 R(g)=\chi(g)\,R_\lambda(e_\tau(\gamma(g))) .       \tag{TL.6}
\]
Exact lattice transport says
$L^\#(gd)=\gamma(g)L^\#(d)$. Extension along $e_\tau$ carries this
to equality of coefficient lattices; taking coordinate duals changes
the frame by inverse transpose. The defining minor polynomials and
their divided-power translates then give exactly
\[
 \nabla(gd)=R(g)\nabla(d),\qquad
 \Delta(gd)=R(g)\Delta(d).                         \tag{TL.7}
\]
Here the character scalar causes no integral error. If
$\chi(g)^r=1$, both $\chi(g)$ and its inverse are algebraic integers,
so multiplication by it preserves every $\mathcal O_K$-module.

The rational matrix is unitary. Under every $e_\tau$ this is the
unitary identity for the coefficient pairing, since \texttt{FOk}
commutes with conjugation on $E$. The normalized tensor Fock form
is invariant under unitary changes of variables; this follows on
monomials from the tensor symmetric-power scalar product and then
on the Weyl subspace by restriction. At the distinguished embedding,
a root of unity has absolute value one. Injectivity of $\sigma$
and \texttt{FOk.sig\_\allowbreak star} give
$\chi(g)^*\chi(g)=1$. Thus (TL.6) also preserves the pairing.
Together with the finite orbit and stabilizer clauses these are all
the assertions of the actual \texttt{fam.Ok}.

\paragraph{TL6: rational fixed vectors give a basis of the complex
forms.}
Choose one representative of each class $c$ and write $\Gamma_c$
for its finite stabilizer. The rational forms are precisely
$\prod_c V_\lambda^{\Gamma_c}$. The matrix coefficient realization
$\operatorname{ev}:V_\lambda\to\operatorname{Val}_{\mathbb C}$ is
injective and intertwines $R(g)$ with
$F(Y)\mapsto\sigma(\chi(g))F(\gamma(g)^{-1}Y)$; evaluation of the
minor generators proves this identity. In a fixed polynomial basis
its columns remain independent after any field embedding: an
injective finite matrix has a left inverse, which can be embedded.
Thus rationally independent form vectors have complexly independent
evaluations. The defined complex polynomial value space is the
complex span of the evaluated standard coefficient basis (equivalently
the complex matrix-coefficient realization of the same Weyl module).
Thus evaluation after scalar extension is onto this space as well
as injective; no surjectivity onto a larger function space is assumed.

For spanning, the projector
\[
 e_c=|\Gamma_c|^{-1}\sum_{g\in\Gamma_c}R(g)        \tag{TL.8}
\]
is a matrix over $K$. Its image is exactly the invariant subspace:
translation permutes the summands, and it fixes an invariant vector.
The same identities hold after scalar extension. Hence every complex
stabilizer-fixed value is a complex linear combination of rational
fixed values. Apply this separately at the finitely many classes.
The actual chosen rational basis therefore evaluates to a basis of
\texttt{Fi}, proving \texttt{LawBasis} and finite dimensionality.
Division in (TL.8) takes place only in the characteristic-zero
field; it is not an assertion of integral averaging.

\paragraph{TL7: the full integral forms lattice and the actual
all-place Hermitian metric.}
For a full lattice $M\subset E^n$, its extension along each embedding
and its coordinate dual are full. The corresponding integral Weyl
lattices are full and finitely generated. One direct verification
puts $M$ between two nonzero scalar multiples of the standard
lattice. The finite degrees of the minors and determinant twists
put its Weyl lattice between two scalar multiples of the standard
Weyl lattice, whose finite divided-power construction spans the
coefficient representation. The same argument applies to the
$\nabla$ lattice and its dual $\Delta$.

For finitely many classes and full coefficient lattices $M_c$, put
\[
 \mathcal L=\{x\in\textstyle\prod_cV_\lambda^{\Gamma_c}:
                         x_c\in M_c\ \forall c\}. \tag{TL.9}
\]
Its inclusion in $\prod_c M_c$ makes it finitely generated, since
$\mathcal O_K$ is Noetherian. For any rational form $x$, choose
nonzero $a_c\in\mathcal O_K$ with $a_cx_c\in M_c$. The product
$a=\prod_c a_c$ puts $ax$ in $\mathcal L$ and is nonzero.
Thus $K\mathcal L$ is the whole forms space.

The all-place positivity needs care about the chosen basis. Define
\[
 K^{\rm cm}=\{b\in K:\phi(b^*)=\overline{\phi(b)}
                       \text{ for every }\phi:K\to\mathbb C\}.
 \tag{TL.10}
\]
This is a subfield. It contains every $e_\tau(E)$ by (TL.1), and
every character root of unity: first use \texttt{sig\_\allowbreak star} and
injectivity at $\sigma$ to get $b^*=b^{-1}$, then use $|\phi(b)|=1$
at any $\phi$. All representation matrices (TL.6) therefore have
entries in $K^{\rm cm}$. Row reduction of the finite invariant
equations produces a basis with coefficients in that subfield.
The source's \texttt{CSD.basis} selects precisely such a basis when
this rational-basis condition holds. This is not an assertion that
an arbitrary $K$-basis has conjugation-compatible coefficients.

Let $A_c$ be the matrix of these basis values at $c$, and $G$ the
standard normalized tensor Fock Gram matrix. The actual class-sum
matrix is $H=\sum_c |\Gamma_c|^{-1}A_c^*GA_c$. Because only the
basis coefficients and the fixed representation pairing enter,
(TL.10) gives for every embedding $\phi$
\[
 \phi(H)=\sum_c |\Gamma_c|^{-1}
               \phi(A_c)^*\phi(G)\phi(A_c).       \tag{TL.11}
\]
The matrix $\phi(G)$ is positive definite: the standard Fock
monomial scalar product is positive and restriction to the Weyl
space is positive. The stacked matrices $A_c$ have trivial kernel;
a left inverse shows that they still do so after $\phi$.
Consequently every nonzero vector has a strictly positive summand
in (TL.11). This proves positivity at every infinite place and,
together with (TL.9), the literal \texttt{LawGood}.

At $\sigma$ the evaluation isometry identifies that class sum with
\[
 \|x\|^2=\sum_c\frac{1}{|\Gamma_c|}
               \operatorname{ipVal}(\operatorname{ev}x_c,
                                     \operatorname{ev}x_c).
 \tag{TL.12}
\]
The normalizing denominator is the product of the Fock
self-pairings of the highest vectors, hence of their squared
norms. It is not the product of their norms. The source's infinite
place representing $\sigma$ may choose its conjugate embedding;
the real value of the self-pairing is unchanged. The positive square
root frame of $\phi(H)$ consequently gives exactly the source's
\texttt{majAt} and \texttt{LawNorm}, with its original normalization.

\paragraph{TL8: the residue ranks, distinguished vector, and dilation
laws.}
For a full rank-$n$ lattice $M$, localization at $P$ is free of rank
$n$ over the discrete valuation ring $\mathcal O_{(P)}$. The natural
map $M/PM\to M_P/PM_P$ is an isomorphism: a denominator outside
$P$ is invertible modulo $P$, and if $sx\in PM$ with $s\notin P$,
choose $t$ with $st\equiv1\pmod P$ to recover $x\in PM$.
Both residue modules in \texttt{TDat} therefore have dimension $n$.
The supplied additive, residue-linear $f$ kills $PN$, so it factors
uniquely to a nonzero linear functional $\mathrm{fv}:N/PN\to k_P$.
Choose $e\in N$ with $f(e)=1$, by scaling a nonzero value by a lift
of its residue inverse. The clause \texttt{te\_\allowbreak mem} gives
$Te\in\mathrm{dn}N\setminus P\mathrm{dn}N$, so its residue class
$\mathrm{te}$ is nonzero. TL11 verifies that clause for the actual
decorated lattices.

Dilation by $a\in k_P^\times$ replaces only $f_P$ by $af_P$.
Its inverse is dilation by $a^{-1}$, it commutes with the acting
group, and it preserves $N$ and $L^\#$. It therefore induces a
group action on classes. The raw stabilizer is unchanged; changing
to the preferred representative conjugates it, so its cardinality
is also unchanged. If $g$ takes the preferred representative of the
dilated class to the dilation of the old representative, its
rational matrix induces an isomorphism
$e:N_{c_a}/PN_{c_a}\to N_c/PN_c$ and
\[
 \mathrm{fv}_{c_a}(v)=a\,\mathrm{fv}_c(e(v)).      \tag{TL.13}
\]
This proves all the literal \texttt{TDat.Laws}, with the factor
$a$ in its correct direction.

\paragraph{TL9: the diamond action preserves complex and integral
zero-step forms.}
Since the zero-step action is free, there is a unique group element
$g_d$ carrying the preferred representative of $[d]$ to $d$.
It satisfies $g_{hd}=hg_d$ and $g_{\operatorname{rep}c}=1$.
Extend a class value $F_c$ to points by
$\widetilde F(d)=\operatorname{tr}(g_d)F_{[d]}$.
Commutation of dilation with the group action proves directly that
pullback along dilation has identity and multiplication laws. The
residue-unit group is commutative, so the order reversal for
pullbacks causes no change to the source's formula. The evaluation
intertwining in TL6 shows that it preserves the polynomial value
space; stabilizer invariance is automatic at this free step.

For integral values, put
$y_c=R(g_{\operatorname{dia}_a\operatorname{rep}c})
 x_{[\operatorname{dia}_a\operatorname{rep}c]}$.
Equation (TL.7), and the fact that dilation leaves $L^\#$ unchanged,
put $y_c$ in the actual coefficient lattice at $c$. It is a rational
form by freeness. Its evaluation is the diamond transform of the
evaluation of $x$. This proves the precise existential
\texttt{dia\_\allowbreak lat} statement, in addition to complex stability and
the representation laws.

\paragraph{TL10: spanning at every coefficient place, and
\texttt{laws\_\allowbreak all}.}
The same proof as (TL.9) applies to the coefficient lattice built
from \texttt{ker} at any step and selected prime: \texttt{TDat.Ok}
supplies its fullness. Thus the actual \texttt{BLat} spans the
rational forms space. For any coefficient place $w$, its localization
$\operatorname{loc}_w(\texttt{BLat})$ contains \texttt{BLat}, so it
has the same full $K$-span. This holds for every place in the supplied
\texttt{Places} record; no seen-place or residue-characteristic
restriction is introduced.

Now start with arbitrary number fields and a generic chain
$I:\texttt{ChainIn}$, level datum $J:\texttt{AtIn}$, period datum
$Q:\texttt{PerIn}$, coefficient embeddings, dominant weights, and
places. Under exactly \texttt{FOk}, \texttt{I.Ok}, \texttt{J.Ok},
and \texttt{Q.Ok}, TL5--TL7 apply at every chain step and also to
$Q.S'$ with its shifted terminal weight. TL8 supplies every level
law, TL9 supplies the zero-step diamond laws and integral stability,
and the preceding denominator argument supplies every
\texttt{B\_\allowbreak span}. These are precisely all fields of
\texttt{J.Laws} and \texttt{FmLaws}. In particular the genuine
proof of \texttt{laws\_\allowbreak all} requires no terminal incidence or
polynomial-value comparison premise.

\paragraph{TL11: the actual level datum is valid and tied to the
decorated lattices.}
We now return to \texttt{TC.Ok}. Its base prime laws give a rational
prime $T$, selected primes $S_T$ with $|S_T|=f$ and $|k_P|=T$,
no conjugate pair in $S_T$, and $T\in P\setminus P^2$.
Auxiliary and twist-frame primes avoid $P,\bar P$: the former have
different rational residue characteristics, and the tuple excludes
ramification at $T$. Thus the source's \texttt{MoveHyp} holds.

Here is the local lattice calculation needed for the actual data.
The elementary duality formula NB1 depends only on lattices and
LP1, and is used in that form. If $N=N^\#$ and
$K_P^{\rm gl}=\{x\in N:f_P(x)=0\}$, then the down lattice
\[
 D=K_P^{\rm gl}+P(K_P^{\rm gl})^\#                \tag{TL.14}
\]
has localizations $D_P=\ker(f_P|_{N_P})$,
$D_{\bar P}=((K_P^{\rm gl})^\#)_{\bar P}$, and $D_R=N_R$ elsewhere.
The exceptional two are dual to each other, proving $D=D^\#$.
Choose a local basis with $f_P$ its last residue coordinate. Then
$D_P=\mathcal O_{(P)}^{n-1}\oplus T\mathcal O_{(P)}$ in that basis.
If $f_P(e)=1$, the last coordinate of $Te$ is a unit times $T$;
it lies in $D_P$ but not in $PD_P$. At $\bar P$ the down lattice
contains $N_{\bar P}$, and at all other primes it equals $N_R$.
Hence $Te\in D\setminus PD$ globally. This proves the actual
\texttt{te\_\allowbreak mem}.

For the up operation with $f_P(x)=1$, put
$A=N+P^{-1}x$ and $U=A^\#+\bar P A$. Its localizations are
$U_P=N_P+T^{-1}\mathcal O_{(P)}x$,
$U_{\bar P}=(A^\#)_{\bar P}$, and $U_R=N_R$ elsewhere.
Again the two exceptional lattices are dual and $U=U^\#$.
The updated functional is $f_P^+(y)=f_P(Ty)$ on $U_P$; in the
displayed basis its kernel is $N_P$. Given the down move's specified
nonzero residue-linear map $g:D/PD\to k_P$, its actual functional is
$f_P^-(y)=\rho_P(s)^{-1}g([sy])$ for $y\in D_P$, where
$s\notin P$ and $sy\in D$, and is zero outside $D_P$.
More generally, any ambient down functional giving a valid decoration
descends to such a map on $D/PD$, and its localization and off-lattice
clauses force this same formula.
Replacing $s$ by another denominator does not change the value:
multiply the two numerator equations and use residue-linearity.
For down after the specified up, (TL.14) returns $D=N$ and the
specified $g$ is the original residue map induced by $f_P$.
The displayed formula is then $\rho_P(s)^{-1}f_P(sy)=f_P(y)$.
Thus both the lattice and its functional return literally to the
old ones. If $x-x'\in PN$, then
$P^{-1}(x-x')\subset P^{-1}(PN)\subset N$, so
$N+P^{-1}x=N+P^{-1}x'$ globally. Their raw up lattices and
functionals are therefore equal even for invalid inputs. Thus both
the validity and fallback branches agree, and up depends only on
the stated residue class without an additional validity assumption.

All remaining decorations and local frames are unchanged, since
their primes avoid the exceptional pair. The coupling coordinate is
defined by local conditions at those other primes, so its allowed
set is unchanged as well. Thus these are literal operations on the
source points, not only operations on lattice isomorphism classes.
Applying a rational unitary matrix transports sums, duals, and
residue maps; hence the up/down action identities hold exactly,
including their frame and coupling coordinates. Dilation simply
rescales one functional, proving all its raw identities. Fullness
of the level kernel follows, for example, by placing it between a
nonzero integer multiple of $L^\#$ and $L^\#$; the finite residue
conditions provide such an integer.

These calculations establish all \texttt{TDat.Ok} and
\texttt{TDat.Tied} clauses for \texttt{tdatAt}. Rechoosing preferred
representatives changes none of them. The Weil additive character
is nontrivial because the specified $\xi_0$ is a unit at the selected
prime; its value on a generator of the inverse-prime quotient has
nontrivial residue phase. Together with the base prime/cardinality
laws, TL4's zero-step freeness, and the positive constants already
in the base data, this gives the literal \texttt{atIn.Ok} for every
word and every twist. This proof does not use NB7/NB8 or the final
\texttt{Forms0} constructor.

\paragraph{TL12: the actual period datum and its level-prime overlap.}
Let $n=N_m$, $n'=n-1$, and $\pi=\pi_T$ with
$\pi\bar\pi=T$. The row step has no type primes. TL2--TL3 apply to
its residue row level and its reduced twist level, giving
\texttt{stepX.Ok}. Extend its rational, auxiliary, and twist acting
matrices by a last coordinate equal to one. This gives an
injective acting-group homomorphism: its first block recovers the
original element. Its rational matrix is
$\operatorname{diag}(\gamma,1)$, so the first coordinates transform
by $\gamma$ and the last coordinate stays fixed.
The point's twist frame, however, extends as
$\operatorname{diag}(\mathrm{fr},\pi)$. Its exact equivariance is
\[
 \operatorname{diag}(\gamma,1)
 \operatorname{diag}(\mathrm{fr},\pi)
 \operatorname{diag}(q,1)^{-1}
 =\operatorname{diag}(\gamma\mathrm{fr}q^{-1},\pi).
\]

The top auxiliary level is exactly the image of the row level under
this diagonal embedding. At a twist prime the top level is its
extended row level times the specified principal congruence group.
The character is pulled back from reduction, is trivial on that
congruence group, and on $\operatorname{diag}(g,1)$ equals the row
character: every old block determinant is unchanged and the added
block determinant is one. These facts prove the exact
\texttt{chr\_\allowbreak ext} identity and the equivariance of the extended
framed decoration. Since $T$ is excluded at every twist prime,
$\pi$ is a unit there, so the specified extended frames are valid.

For a row lattice $N=N^\#$ the extended decoration has
\[
 \widetilde N=N\oplus\mathcal O,\qquad
 \widetilde f_P(x,z)=\rho_P(\bar\pi z),\qquad
 \widetilde L=N\oplus\pi\mathcal O,
 \qquad\widetilde L^\#=N\oplus\bar\pi^{-1}\mathcal O.
 \tag{TL.15}
\]
To check the third equality, the selected primes are exactly the
prime divisors of $(\pi)$ with exponent one: its norm is $T^f$ and
the $f$ selected primes each have norm $T$. Their intersection is
therefore $(\pi)$, and $\bar\pi$ is a unit at each one. The last
equality follows by taking the Hermitian dual; in particular the
denominator is $\bar\pi$, not $\pi$.
With $t_z=\bar\pi$, (TL.15) says exactly
\[
 (x,z)\in\widetilde L^\#\iff x\in L^\#\text{ and }t_zz\in\mathcal O,
 \quad (x,z)\in\widetilde N\iff x\in L^\#\text{ and }z\in\mathcal O,
 \quad \bar t_z t_z=T.                            \tag{TL.16}
\]

The terminal preferred-representative predicate asks to lie in the
image of this extension. Every extended point itself satisfies it,
so the representative selected in its orbit also lies in that image.
This proves the actual \texttt{rep\_\allowbreak til}; it is not a claim that
all terminal orbits lie in the image. The subgroup $H$ in this
period datum is the full row acting group, so its finite index is
one and its rational-image condition is immediate.
The preceding identities prove every field of \texttt{perIn.Ok}
and \texttt{perIn.Over(atIn)}, for every word and twist, directly
from \texttt{TC.Ok}. No terminal comparison from TG is used.

\paragraph{TL13: the complete implication \texttt{Ok}$\Rightarrow$
\texttt{Laws}, and its exact scope.}
For any actual base $B$, prechoices $X$, and true-column data
$I:\texttt{TC}(B,X)$, assume only $I.\texttt{Ok}$ in addition to the
data already built into $B,X,I$. TL3--TL4 give all \texttt{ch} fields
and nonempty zero indices; TL11 gives \texttt{lev}; TL12 gives
\texttt{per} and \texttt{ovr}; TL10 gives \texttt{atL} and
\texttt{fmL} at every word and twist. The base fields and primes
give \texttt{pr}, TL4 gives \texttt{xi0} and \texttt{bij0}, and
every decorated level kernel is a full lattice by the finite-residue
argument in TL11, giving \texttt{lat}. These are all eleven fields
of the literal \texttt{TC.Laws} record.

Thus the law premise in the previously conditional constructions
can now be supplied by this proof whenever \texttt{TC.Ok} is
supplied. The quantifiers are all words, all character indices,
all natural-number chain steps, all classes, and, in
\texttt{B\_\allowbreak span}, all coefficient places. No new \texttt{Pins},
\texttt{Later}, chosen adelic realizer, or terminal comparison is
needed to derive these laws. This implication alone does not supply
\texttt{TC.Ok}; in particular its local index nonvanishing clauses
remain present in that input.

\paragraph{TL14: one number field contains all the normalized frame
character values.}
We give the substantive arithmetic argument behind
\texttt{frame\_\allowbreak values}. Let $\xi$ be a continuous Hecke character
of a CM field, with infinity type $n_\tau\in\mathbb Z$, and let
$r\in\mathbb Z$ satisfy $n_\tau-r\in2\mathbb Z$ for every CM
embedding. Put
\[
 a_P=\xi(P)\,\mathrm N(P)^{-r/2}
 \quad\text{at unramified primes }P.              \tag{TL.17}
\]
The complex target has a neighborhood of one containing no
nontrivial subgroup. Continuity on the finite ideles therefore
gives a finite set $S$ outside which local units are killed and,
at each $P\in S$, an open principal-unit subgroup in the kernel.
One can see this without a global character theorem: choose a
basic open subgroup inside the inverse image of that neighborhood.
Its image is a subgroup contained in the neighborhood, hence is
trivial. The quotient of local units by the resulting congruence
subgroup is finite. Thus the local character values on units at
$S$ have bounded root-of-unity orders.

Let $M$ contain all conjugates of $E$ and those finitely many roots
of unity. The ordinary ideal class group is finite. For clarity,
its usual lattice proof only needs a uniform short-vector bound:
for an integral ideal $I$, averaging a ball of volume exceeding
its Minkowski covolume gives a translated ball containing two
lattice points. Their difference yields $0\ne a\in I$ with
$|\mathrm N(a)|\le C_E\mathrm N(I)$. Hence
$(a)I^{-1}$ is integral of bounded norm. There are finitely many
ideals of bounded norm, since an ideal of norm $b$ contains
$b\mathcal O$ and corresponds to a subgroup of the finite group
$\mathcal O/b\mathcal O$. This proves class-group finiteness.

Let its order be $h$. If unramified primes $P,P'$ have the same
class, then $P(P')^{h-1}=(a)$ for a nonzero algebraic integer $a$ which is
a unit at every prime of $S$. The Hecke product formula gives
\[
 a_Pa_{P'}^{h-1}
  =\left(\prod_{Q\in S}\xi_Q(a)\right)^{-1}
    \prod_\tau \tau(a)^{-n_\tau}
       \bigl(\tau(a)\overline{\tau(a)}\bigr)^{(n_\tau-r)/2}
  \ \in M.                                         \tag{TL.18}
\]
Indeed $\xi_\infty(a)=\prod_\tau(\tau(a)/|\tau(a)|)^{n_\tau}$,
and $\mathrm N((a))=\prod_\tau|\tau(a)|^2$; the parity assumption
is exactly what turns the remaining powers of absolute values
into the integral powers in (TL.18).
Choose one unramified prime in each class which occurs. Setting
$P=P'$ in (TL.18) makes its chosen value algebraic over $M$.
Adjoin these finitely many nonzero values. Equation (TL.18) now
places every $a_P$ in this one number field. This proves the
literal \texttt{ValuesIn(r/2)} assertion, including its uniformity
over all unramified primes, rather than just individual
algebraicity.

For the supplied frame, $\mu_j$ has even infinity exponents,
$\eta$ has exponent one and uses $r=1$, and $\chi_\Psi$ has
exponent two and uses $r=0$. The finite-order $\delta,\delta'$
and the finite auxiliary characters require only cyclotomic
extensions. A finite compositum, then its normal closure, gives
one finite Galois field containing all five character families,
all field embeddings, and the auxiliary roots, exactly as in
\texttt{frame\_\allowbreak values}. A primitive generator may be required
inside the fixed coefficient field in \texttt{large\_\allowbreak 2}; this
does not amount to enlarging that fixed field automatically.

\paragraph{TL15: finite root data, actual residue characters, and
finite-idele readings.}
For a finite-order character $\xi$ of order dividing $N>0$, its
values are powers of a primitive $N$th root. At each of its
finitely many ramified primes with norm $q$, adjoin
$\sqrt q$ and a primitive $q^{2m}$th root. Taking the union for
a finite family $\Xi$ gives finitely many algebraic generators
which imply the exact \texttt{Holds} condition for every member.
This includes the whole character image, not just its values at
uniformizers. A field containing one primitive $N$th root
contains every $N$th root by its cyclic list of powers.

At an auxiliary prime above $g_i$, identify its residue field
with $\mathbb F_{g_i}$ by the prescribed prime-field isomorphism.
Lift the primitive multiplicative character $\omega_i$ and the
additive character $b\mapsto\exp(2\pi\mathrm i b/g_i)$ through
the injective coefficient embedding. Define
\[
 \psi_{0,P}(b)=\psi_i(\overline{\xi_0}\,b),\qquad
 \chi_{j,P}(b)=\omega_i(b)^{2(j+1)}.               \tag{TL.19}
\]
Here $\overline{\xi_0}$ means reduction at $P$, not complex
conjugation. Since $\xi_0$ is a local unit, $\psi_{0,P}$ is
nontrivial. The primitive $\omega_i$ has order $g_i-1$; the
inequality $2m<g_i-1$ makes the powers in (TL.19) distinct.
The powered characters themselves need not be primitive.
The reductions of $\xi_0$ at the selected primes over a given
$g_i$ agree under these prime-field maps: $\xi_0$ lies in the
specified quadratic subfield, and all these primes lie over
the one selected prime of that subfield. Thus the additive and
multiplicative compatibility clauses hold with their actual
residue maps. Tameness makes a local unit character depend only
on its residue; compare a unit to the natural integer with the
same residue and use the prescribed formula for $\mu_j$ on that
integer. This proves the inverse-character tie in
\texttt{Chars.Tied}. Lift the normalized $\mu_j$ and $\eta$
prime values using TL14. The exponential in (TL.19) gives
\texttt{WeilTied} literally. These arguments construct all of
\texttt{Chars.Ok}, \texttt{Tied}, and \texttt{WeilTied} from the
stated frame and root data.

For the reading step, a finite-order character is trivial on
the infinite ideles: every complex unit has an $N$th root, so
its value under an order-$N$ character is one. Its finite-idele
restriction therefore determines it. If all its values lie in
$\sigma(K)$, pull them back uniquely to $K^\times$. Multiplication
and every triviality identity descend through the injective
$\sigma$. Continuity provides the open congruence kernel as in
TL14. The supplied \texttt{IsTwist} conditions say precisely
finite order, triviality on the base-field ideles, tameness, and
ramification only at the selected split, degree-one unramified
primes outside the excluded set. They give the tuple's wild-unit,
inert-place, and base-field kernel clauses. For an infinite
product of local units the assertion follows first on the
finite-support subgroup and then by its density and continuity.
Conversely the tuple's open kernel gives a continuous finite-idele
character, and the rational-unit identity descends it to the
idele class group with trivial infinite component. These two
constructions are inverse. In particular the read tuple is
unique and the tuple for the trivial character is one.

\paragraph{TL16: orienting the auxiliary factors and obtaining the
quotient part of \texttt{Pins}.}
Suppose an integral $a$ lies in the specified imaginary quadratic
subfield and satisfies $a\bar a=g$, with $g$ a rational prime.
Neither $a$ nor $\bar a$ is a unit: conjugation would otherwise
make both units. Their principal ideals in the quadratic field
have norms greater than one with product $g^2$, so each has norm
$g$ and is a maximal ideal. A prescribed prime above $g$ contains
one of them. Choose $b=a$ or $b=\bar a$ so that the ideal of
$\bar b$ is that prescribed prime. A prime of $E$ contains
$\bar b$ if and only if its restriction is that prime. Norm and
the required congruence to one persist under conjugation. Doing
this independently at $g_0,g_1$ proves the actual
\texttt{exists\_\allowbreak oriented} statement.

Under the given $\texttt{Ctx.Over}$ relation replace the type-prime
factors by the context's own $\pi_T(u)$, preserving the now-oriented
auxiliary factors. Their prescribed norm and congruence laws make
this another valid \texttt{PreX}. The context defines
$c(u)=\pi_T(u_0)/\pi_T(u)$, and its prime norm makes the denominator
nonzero. Hence
\[
 c(u)\pi_T(u)=\pi_T(u_0)\quad\text{for every word }u. \tag{TL.20}
\]
This supplies \texttt{Pins.cof}. The other field,
\texttt{Pins.iX}, is the literal agreement of the supplied row
frame with the first columns of the full frame; it is still an
explicit antecedent in \texttt{evLg\_\allowbreak canonical}. No change of
coordinates here proves an arbitrary supplied frame inclusion.

\paragraph{TL17: finite principalization and the precise period-field
boundary.}
For any nonzero ideal $I$ of a number field, class-group finiteness
gives $I^h=(\beta)$ for some $h>0$. Adjoin a root $x^h=\beta$.
It is integral, and in the larger integer ring
$(x)^h=(I\mathcal O')^h$. Unique ideal factorization then gives
$(x)=I\mathcal O'$ by comparing exponents at every prime.
Fractional ideals are treated after clearing a denominator.
Thus finitely many rank-one fractional lattices become principal
over one finite extension, with algebraic generators. Applied to
the coordinate-axis sections of the module generated by the
conjugate vectors, this supplies the algebraic derivative scalars
$c_D$. Their sections commute with this extension: over a
Dedekind domain the larger integer ring is torsion-free, hence
flat, so taking the kernel of the projection to the other
coordinates commutes with extension. This proves the finite
principalization content of \texttt{exists\_\allowbreak dgen}.

For the period field, the host identity
$\mathfrak D_{E/K_0}=\mathfrak b\bar{\mathfrak b}$ gives
$\mathfrak b^{-1}\mathfrak D_{E/K_0}=\bar{\mathfrak b}$.
Hence each source period-line lattice is the extension of a
conjugate of $\bar{\mathfrak b}$. The preceding root construction
principalizes every one simultaneously. Adjoining also
$j(\mathbb C/\mathfrak c)$ therefore gives the actual
\texttt{FieldOK}, provided this CM $j$-value is algebraic.
The source proves the stronger integrality assertion using the
modular polynomial/CM-lattice theorem. That analytic modular
input is retained here; it has not been inferred from ideal
principalization. The finite extension and its generator are
constructed once that one algebraicity input is supplied.

Similarly, a full \texttt{AmbIn} includes the chosen period pair,
the quadratic trace isometry and its complex linear extension,
and the local correction data, as well as $c_D$ and a discriminant
twelfth root. TL17 closes its finite principalization part, not
the separate supplier of these period and trace-isometry data.
These are recorded as exact inputs in TL21.

\paragraph{TL18: reducing normalized-jet algebraicity to the primitive
theta inputs.}
There is a useful elementary step in the analytic supplier which
can be reconstructed without presuming its conclusion. Consider
an algebra of smooth germs over $\overline{\mathbb Q}$, generated
by germs whose derivatives in algebraic coordinate directions
stay in the same algebra and whose values at zero are algebraic.
The sum, product, and chain rules prove by induction that every
iterated derivative of every germ in the algebra is algebraic
at zero. The same holds after a linear substitution with algebraic
coefficients.

Suppose the supplied four-copy identity places
$f(z_1)f(z_2)f(z_3)f(z_4)$ in this algebra. Apply the same
multi-derivative separately in each copy and set all four
variables to zero. The result is $(D^\alpha f(0))^4$, which is
algebraic. Hence $D^\alpha f(0)$ is algebraic too. This argument
works even when $f(0)=0$; it never divides by a value of $f$.
If the available algebraic directions are a spanning family,
choose independent columns and invert their nonsingular minor
over $\overline{\mathbb Q}$ to express any algebraic direction
with algebraic coefficients. This proves the source's direction
descent in the four-copy calculation.

For the actual vacuum, the four-copy identity is a finite sum
of products of one-dimensional torsion translates of the
normalized even theta function. Its coefficients are the
algebraic correction factor $\kappa$, signs, and inverse square
roots of finite indices. Each factor of $\kappa$ satisfies its
stated $2h_{K_0}$th-power equation with an algebraic right-hand
side, so it is algebraic; only finitely many nontrivial factors
occur. Thus, \emph{given} that four-copy identity and the
normalized CM torsion-translate Taylor theorem, the preceding
argument proves algebraicity of the standard vacuum's Taylor
coefficients.

The remaining transport is also explicit. Two rational
Heisenberg translations compose with a phase
$\exp(2\pi\mathrm i q)$, $q\in\mathbb Q$, so their extra factor
is algebraic. The intertwining operator is a finite sum of such
translations with coefficients equal to signs times an inverse
square root of an index. It preserves algebraicity of all
torsion-translate Taylor coefficients. The source's normalized
indicator theta function is, by its explicit finite presentation,
a finite algebraic linear combination of rational translations
of this moved vacuum. The multi-derivative and coordinate
reindexing formulas then prove \texttt{JetAlg}, and the supplied
entirety identities give \texttt{ESmooth}.

The remaining primitive statements are therefore precise: the
normalized CM theta Taylor theorem, its torsion-translation
algebraicity theorem, and the vacuum/four-copy and normalized
indicator identities. This paragraph proves the closure and
transport steps conditional on those identities. It does not
claim to have proved those analytic identities from formal
definitions alone.

\paragraph{TL19: finite generators for all bounded-degree normalized
jets.}
Assume the actual smoothness, unitary covariance, normalization
character law, and \texttt{JetAlg} of the theta family in question;
also assume its $c_D$ are algebraic, as in TL17. Fix positive
modulus $M$, a finite prime set, two ranks, and a total degree
$d$. TL1 gives finitely many self-dual lattice representatives
in the two ranks. There are finitely many divisors $c\mid M$.
For fixed representatives and $c$, the permitted rational shift
is bounded in a fixed fractional Hom lattice, while the theta
function is unchanged by adding $c$ times its local integral
Hom lattice. The resulting quotient is finite: localize at the
finitely many relevant primes, where a bounded lattice modulo
a deeper full lattice has finite residue cardinality, and use
the injective map to those local quotients as in TL2.
There are also finitely many multi-indices of degree at most
$d$.

Take all normalized derivatives for these finite choices.
They are algebraic by \texttt{JetAlg}. Adjoin them, the finitely
many $c_D$, and generators for all embedded copies of $E$.
For arbitrary self-dual lattices choose rational unitary maps
to the fixed representatives. The actual covariance identity
expresses the theta function as the representative function
composed with its specified linear change of variables, whose
coefficients are embedded elements of $E$. Repeated chain rule
expresses each derivative of degree $d$ as a finite sum of
representative derivatives of that degree times products of
these coefficients. The normalization-character change is
the prescribed product of embedded determinants, also in the
same field. Finally the rescaling by $c_D$ uses algebraic
factors already adjoined. This proves \texttt{FinAlg(jetSet)}:
one finite algebraic extension contains the whole infinite
set of these jets.

For canonical column data the exact modulus is
\[
 M(u,S_G)=Q^{m+1}\bigl(T_u Q^{m+1}\bigr),
 \qquad Q=\prod_{P\in S_G}q_P.                    \tag{TL.21}
\]
All prime sets used by the steps are subsets of the finite
$\texttt{badP}(B,\Xi)$ for the finite character family $\Xi$.
There are finitely many words, steps $i<m$, and such pairs of
prime sets. The union of their \texttt{FinAlg} sets is again
\texttt{FinAlg}, by taking the finite union of generators.
The source's \texttt{allSet} includes only indices satisfying
\texttt{FinAlg}; its finiteness proof alone does not show that
an actual step belongs to it. The analytic inputs and the
preceding argument supply exactly that missing membership.
The actual multidegrees have total degree \texttt{totDeg},
so containment of the chosen generators in the coefficient
field proves the literal \texttt{JetsIn} for those canonical
data. This is a finite-generation reduction, not a proof of
analytic algebraicity by selecting an empty set.

\paragraph{TL20: what the local index and canonicalization really
require.}
At an ordinary twist prime the initial index construction selects
finite representatives recursively. Starting with $F_0=\{1\}$,
retain at the next step exactly the admitted double cosets which
have a nonzero first-unfolding coefficient from a member of $F_i$.
Given finite admitted level-coset representatives, normal-form
reduction, and invariance of nonvanishing under the indicated
unipotent/double-coset moves, choose one normalized representative
in each retained coset. This produces a finite $F_{i+1}$, with
the exact predecessor, uniqueness, and exhaustivity properties.
At $m-1$ set $F_{m-1}=\{1\}$ and above it set $F_j=\varnothing$,
as stipulated by the source. Finite products over the local
primes give the global index sets. A level character descends
to its projected group when it is trivial on the kernel:
define its value using a lift, and compare two lifts by their
quotient in that kernel. This explains the finite selection
and character-descent parts of \texttt{exists\_\allowbreak chIdx}; the
needed local coset, normal-form, kernel-character, and
first-unfolding theorems remain genuine inputs here.

For an already supplied $I.\texttt{Ok}$ of depth at least $2m$,
much of the canonical arithmetic can be extracted without an
existence argument. Its square-root clause gives $\sqrt{q_P}$.
Its reference local-character formula gives the required
coefficient-field values of $\mu_j(q_P)^{-1}$. If the old depth
is $D=2m+a$, evaluate its Weil identity at
$\xi_0^{-1}/q_P^{2m}$, using the integer $q_P^a$ in that
identity. The local-unit clause makes this an allowed argument
and gives exactly $\exp(2\pi\mathrm i/q_P^{2m})$ in the
coefficient field. Thus the primitive roots required for the
canonical depth are already present. The base congruences give
$q_P\equiv1\pmod4$ at auxiliary primes and the tuple separates
twist primes from them.

Replacing the constants and indices by \texttt{twCst} and
\texttt{twIdx} now produces \texttt{IsTw} definitionally,
and leaves the host, character family, tuple readings, selection,
and ambient theta data unchanged. To assert that the replacement
still satisfies \texttt{Ok}, however, the source invokes the
actual twisted first-unfolding nonvanishing and its double-coset
invariance, together with the corresponding character
nonvanishing theorem. These do not follow from the root and
square-root extraction just proved. They remain the exact
\texttt{twIdx\_\allowbreak ok\_\allowbreak closed} boundary. Canonical replacement
therefore does not make the local-index hypotheses disappear.

\paragraph{TL21: the quantified canonical evaluation assembly and
the remaining construction boundary.}
The phrase \texttt{EvLg} has two nested meanings. Its inner
\texttt{large\_\allowbreak 2} assertion chooses finite algebraic generator
sets $S_1$ and $S_2$, and finite prime/avoidance conditions,
before quantifying over a base $B$ lying over the fixed context,
a frame $f$, and a twist datum $T$. The antecedent requires
$S_1$ in the \emph{fixed} coefficient field and $S_2$ in the full
coefficient field. It does not assert the result for an arbitrary
unchanged coefficient field. In this constructor the additional
prime and avoidance sets can be empty. The twist family
\[
 \Xi=\{1,\delta,\delta'\}\cup
       \{(\delta'\text{ if }\operatorname{reschar}(w)=p_\delta
                      \text{ else }\delta)\psi_w:w\}
 \tag{TL.22}
\]
is finite: outside the finite set of supported rational primes
$\psi_w=1$, and there are only finitely many coefficient places
over each supported prime.

Choose $S_1$ to contain a primitive generator for TL14's common
frame-value field. Choose $S_2$ to contain the period-field
generators of TL17, the roots and square roots of TL15, and the
jet generators of TL19. Their finite union is algebraic under
the primitive inputs explicitly retained in TL17--TL20.
With a nonempty \texttt{PreX}, TL16 gives the required auxiliary
orientation and type-prime quotient identity. Under the stated
frame-inclusion antecedent this gives both fields of
\texttt{Pins}.

The finite character set $\Xi$ itself is the index type. TL15
gives its unique tuples, the three distinguished indices, the
trivial reference tuple, and the actual residue character data.
The local-index suppliers in TL20, applied with initial
$M=A=B=m$ and twist depth $2m$, then give the actual initial
\texttt{TC.Ok}; the supplied ambient period/trace-isometry
constructor supplies its specified \texttt{AmbIn}. The selection,
frame, excluded set, and range of $\xi$ are exactly the prescribed
ones, because they were used as these record fields. TL20's
canonical index supplier preserves \texttt{Ok} and gives
\texttt{IsTw}; TL19 gives \texttt{JetsIn}. TL14 supplies
\texttt{ValLL0} for every later assembly with the same base and
frame, and TL17 supplies \texttt{FieldOK}. Finally TL13 now
supplies \texttt{Laws} without a further premise.

The outer \texttt{Ev} quantifies a positive $\varepsilon_0$, then
$\varepsilon,C_8$, a transported lower-bound datum, constants,
and sufficiently large contexts. The inner construction just
described works for every such context, so the outer thresholds
can be chosen trivially (for example $\varepsilon_0=1$ and
zero size thresholds) while retaining the literal inner
coefficient-field requirements. This explains the source's
\texttt{ev\_\allowbreak of\_\allowbreak forall} step; it neither constructs an arithmetic
host nor removes the existing context hypotheses.

Accordingly, the remaining sources for the full canonical
existence conclusion are the actual local-index/nonvanishing
theorems; the period/trace-isometry ambient datum and CM
$j$-algebraicity; and the primitive normalized-theta identities
specified in TL18 together with their unitary covariance law.
The finite fields, coefficient roots, tuple reading, orientation,
ideal principalization, finite derivative reduction, and the
entire implication \texttt{Ok}$\Rightarrow$\texttt{Laws} have
been reconstructed above. The record \texttt{Later} still
consists of five separately supplied stages. A universal
statement about all such later assemblies is not an existence
proof of those stages, of a primitive host, or of RH/GRH.



\section{The symmetric terminal levels and their Hermitian geometry}
\label{sec:terminal-geometry}

\status{TG1--TG12 construct the previously missing symmetric
penultimate level, terminal row chart, full and row auxiliary levels,
terminal incidence and Hermitian carriers. They finish the actual
\texttt{FormsTiedSym} and \texttt{ColsTiedSym} constructions.
They use the actual arithmetic
host and the stated \texttt{TC.Ok}, \texttt{Pins}, \texttt{Laws} and
\texttt{Later} data. Initial decorated-point realizers are supplied by
DR3, not retained as an additional hypothesis. Integral lattices,
specified residue maps, preferred class representatives and the common
global coordinate frame are preserved throughout.}

\paragraph{Conventions and supplied data.}
Keep the true column datum $I:\mathrm{TC}(B,X)$ and write
$D=I.\mathrm{rec7}(l,Z)$. The words $u$ range over the original word
set; all constructions below use the one auxiliary-prime datum $G$ of
FC13. Fix $u$ temporarily, put $p=T_u$, $\pi=\pi_u$ and
$\mathcal O=\mathcal O_E$, and use the standard coordinates of SF2:
\[
 h(x,y)=\bar x^{\mathsf T}y,\qquad
 \pi\bar\pi=p,\qquad(\pi,\bar\pi)=\mathcal O.
 \tag{TG.1}
\]
The selected type primes are $S=\{P:P\mid\pi\}$. Their residue
maps are the original maps $\operatorname{res}_P$; the selected
auxiliary primes lie over $g_0,g_1$, are disjoint from $S\cup\bar S$,
and satisfy $\pi\equiv1$ there. The reference twist has no $Q$-frame
primes. The supplied rank identities give
\[
 s=m/2-1,\quad r=k+2s+1=k+m-1,\quad
 N_{m-1}=r+(m-1)t,\quad n=N_m-1=r+mt.
 \tag{TG.2}
\]
Here $m\ge2$ is even and $t>0$ is odd. The row space is $E^n$,
the full terminal space is $E^{n+1}$, and the inclusion $\iota$ is
$x\mapsto(x,0)$. The distinguished vector spanning its orthogonal
line is not silently identified with the last basis vector.

\paragraph{TG1: the penultimate symmetric group is the actual residue group.}
For any relevant matrix size $d$, put
\[
 \rho(a)=
 \begin{cases}s+a&a<k,\\a-k&k\le a<k+s,\\a&a\ge k+s,
 \end{cases}
 \qquad
 \ell(a)=
 \begin{cases}\rho(a)&a<r,\\r+\lfloor(a-r)/t\rfloor&a\ge r.
 \end{cases}
 \tag{TG.3}
\]
The group $J_{\rm sym}(d)$ consists of the invertible residue
matrices $M$ such that both $M$ and $M^{-1}$ satisfy
\[
 M_{ab}=0\quad\hbox{if }\ell(b)<\ell(a),\qquad
 M_{aa}=1\ (a<r),\qquad M_{a,k+s}=0\ (a<k).
 \tag{TG.4}
\]
This is literally \texttt{JSymG} with its final Boolean equal to
\texttt{false}, hence the source's \texttt{HSym}. At step $m-1$
the actual finite group $\mathrm{levG}$, transported through the
specified auxiliary residue-field isomorphism, equals
$J_{\rm sym}(N_{m-1})$.

Here is a proof of both inclusions. At an auxiliary prime $P$ the
integer $q=g_a$ is a uniformizer in $\mathcal O_{E,P}$, since the
prime is completely split and has residue cardinality $g_a$.
The extra-space parameter at step $m-1$ is zero. Thus the actual
flag lattice at a head cut $b<r$, or a block cut $b=r+jt$, is
\[
 F_b=\{x:\ x_a\in q\mathcal O_{E,P}
                \text{ if }\rho(a)\ge d(b),\quad
                x_a\in\mathcal O_{E,P}\text{ otherwise}\},
 \quad d(b)=\begin{cases}s&b=k+s,\\b&b\ne k+s.\end{cases}
 \tag{TG.5}
\]
The marked vector for the head cut $b$ is $e_{\rho^{-1}(b)}$.
The local group preserves these lattices and satisfies
$(g-1)e_{\rho^{-1}(b)}\in F_b$. Taking $b=\rho(a)$ forces the
head diagonal to be one and every entry below it in the ordering
$\ell$ to vanish modulo $P$. At $b=k+s$ the cut begins at $s$;
its marked-vector condition additionally forces
$M_{a,k+s}=0$ for $a<k$. The block cuts force all lower off-block
entries to vanish. Applying the same argument to $g^{-1}$ proves
all of (TG.4) for the inverse as well: preservation of $F_b$ turns
$(g-1)e\in F_b$ into $(g^{-1}-1)e=-g^{-1}(g-1)e\in F_b$.

Conversely, lift $M$ entry by entry to $\mathcal O_{E,P}$.
Its determinant is a unit, so the lift $g$ and its inverse are
integral and reduce to $M,M^{-1}$. If a row lies on the scaled
side of a cut (TG.5) and a column does not, then
$\ell(\text{column})<\ell(\text{row})$. Therefore the corresponding
entry of $g$ is divisible by $q$. These divisibilities prove
$gF_b\subset F_b$, and the inverse conditions prove equality.
For a marked head vector, every scaled coordinate of
$(g-1)e_{\rho^{-1}(b)}$ vanishes modulo $P$: use the diagonal and
lower-entry conditions, and at the exceptional cut use precisely
the extra zero block in (TG.4). Division by the uniformizer now
gives the required membership in $F_b$. Hence the lift belongs
to the actual local symmetric flag group. The supplied
\texttt{grpG} and \texttt{last\_\allowbreak h} laws identify its reduction
image with $\mathrm{levG}$. This proves equality, including
surjectivity onto the stated finite group; it is not just a
containment of possible reductions.

\paragraph{TG2: the exact penultimate character and common level.}
For $b$ complete $t$-blocks after the head, define
\[
 \Lambda_b(M)=
 \psi\!\left(\sum_{\substack{0\le a<r-1\\a+1\ne k}}
                     M_{a,a+1}\right)
 \prod_{j=0}^{b-1}\chi_j\!\left(
       \det M_{[r+jt,r+(j+1)t)}\right).
 \tag{TG.6}
\]
This is exactly \texttt{lamSym}; the omitted edge is the entry
$(k-1,k)$ in the original coordinate order. The character of the
actual penultimate finite level is $\Lambda_{m-1}$, after the
same residue-field transport as in TG1.

Indeed, with extra-space parameter zero, the coordinate matrix
in the flag character is the original matrix. Its additive
argument is the two superdiagonal sums separated by the omitted
edge; its other factors are the displayed determinants. Reduction
commutes with these sums and determinants. The extra inverse
character has argument zero and contributes one. Every block
determinant is nonzero on the group, so the source's extension
of a multiplicative character by zero contributes its ordinary
unit value. FC13 identifies $\psi$ and every $\chi_j$ under the
prescribed residue embeddings. Agreement there on natural
integer residues suffices because each auxiliary residue field
is the prime field of order $g_a$. This proves (TG.6) on every
group element and also proves that it is a character: it is the
transport of the actual level character.

Take the common level $J$ and finite $J$-set of FC12, obtained
from the zero-coupling chart. TG1 identifies its auxiliary
reduction groups with $J_{\rm sym}(N_{m-1})$. FC14 then gives
the actual \texttt{GLevH}, with frame equal to the identity at
every auxiliary place and local character (TG.6), $b=m-1$.
The global character retains the source's inverse convention:
$\chi_J(x)$ is the product of $\Lambda_{m-1}$ evaluated at the
inverse reductions of $x$. None of the common-level construction
is changed. In particular
\[
 Y=\coprod_\beta J/U_\beta,\qquad
 \operatorname{Stab}_J([1]_\beta)=U_\beta,
 \qquad \chi_J|_{U_\beta}=\chi_\beta.
 \tag{TG.7}
\]
The finite reduction kernel lies in every $U_\beta$, which makes
all these coset sets finite. This proves the full
\texttt{stepBefore\_\allowbreak oneLevel} conclusion, including its specified
$\mathrm{HSym}$ and $\mathrm{lamSym}$ clauses.

\paragraph{TG3: construct the terminal row chart.}
The raw row action has empty type-prime set and, at the reference
twist, empty $Q$-frame set. Its valid decoration has a full
self-dual lattice $N_c\subset E^n$, the prescribed auxiliary
residue bases, and the actual finite groups
$\Lambda_P=\mathrm{rowTyG}_P$. Define the row base lattice to be
$\mathcal O^n$. At auxiliary places impose reduction in
$\Lambda_P$; elsewhere take the full integral unitary stabilizer.
Let $U_P$ be this adelic level and give it character
\[
 \chi_P(x)=\sigma\!\left(
       \prod_{Q\in S_G}\lambda_Q(\operatorname{red}_Q(x)^{-1})
                         \right).
 \tag{TG.8}
\]
Each preferred row decoration has an actual adelic representative
$y_P(c)$ by DR3, with these exact residue bases. Thus
$y_P(c)\widehat{\mathcal O}^{n}=\widehat N_c$.

The classes of the original row action identify bijectively with
the double quotient by the rational unitary group and $U_P$.
To check this assertion, write the group action as its rational
factor $\gamma$ and its auxiliary finite factor $a$. A relation
between two decorations gives, by the exact carried-lattice and
residue-basis equations, a matrix $x$ with
$\operatorname{rat}(\gamma)y_P(c')=y_P(c)x$ and $x\in U_P$.
Conversely this adelic equality recovers the lattice by rational
intersection (AD4), and recovers the finite factor by reduction;
it is the original decoration relation. For surjectivity, start
with any unitary finite adele, intersect its carried lattice
with $E^n$, and use its reduced inverse coordinates as the
auxiliary bases. At each place its lattice is a unitary image of
the self-dual standard lattice; the paired-dual calculation of
NB1 and the local-global intersection of LP1 make the global
lattice self-dual. Its reduced inverse coordinates are an actual
residue basis, so the resulting decoration is valid. Its preferred
representative is in the same double coset.

The stabilizer map is the actual rational factor $s\mapsto\gamma_s$.
It is injective: when $\gamma_s=1$, preservation of the auxiliary
bases forces every finite factor to be one; there are no other
factors because the reference $Q$-set is empty. Its range is
exactly the rational unitary elements in
$y_P(c)U_Py_P(c)^{-1}$. Equation (TG.8) is the corresponding
character relation, with the inverse already accounted for.
These statements prove every \texttt{ReadsP} and
\texttt{ClsPTied} field, not merely the existence of a name for
such a chart.

Apply CV10 to this actual chart, using its unchanged frame
identification. It gives the complex, rational and exact integral
row-value clauses. The constant shift of the terminal weight by
$m(t+1)/2$ leaves its polynomial weight, integral $\Delta$ lattice
and Fock pairing unchanged; the determinant factor stays in the
frame evaluation. FC7 therefore supplies the Petersson and
\texttt{NormPTie} clauses with the specified embeddings and
rational basis. We obtain the source's full
\texttt{clsP\_\allowbreak tie}, with its stipulated identity map on row
frame functions. No equality between integral $\Delta$ and
dual-Weyl lattices is used.

\paragraph{TG4: the full and row terminal levels and their frames.}
On $E^n$, the actual row finite group is the parabolic with
symmetric head and $m$ blocks of size $t$. By reading its
head and block conditions, it is $J_{\rm sym}(n)$ of (TG.4).
On $E^{n+1}$ the terminal finite group is the same parabolic
with one final block of size one, and with the final column
fixed. The lower block-zero condition forces the final row to
vanish away from its diagonal. Consequently it is exactly
$J_{\rm sym}^{+}(n+1)$, where both the final row and final
column are those of the identity, for the matrix and its inverse.
This is \texttt{JSymG true}. Conversely these conditions give
the original parabolic and fixed-column conditions, so both
identifications are equalities.

The character in each case is (TG.6) with $b=m$. The same
skipped-superdiagonal and block-determinant calculation proves
this, including the coefficient embedding. The full terminal
base lattice is
\[
 L_\pi=\mathcal O^n\oplus\pi\mathcal O,
 \tag{TG.9}
\]
whereas the row base is $\mathcal O^n$. At an auxiliary place
$v$, $\pi$ is a unit. For its prescribed split projection
$\phi_v$, choose
\[
 F_P(v)=I_n,\qquad
 F_M(v)=\operatorname{diag}(1,\ldots,1,\phi_v(\pi)).
 \tag{TG.10}
\]
The latter spans the same completed integral lattice as the
identity and commutes with the full residue group, because that
group fixes the last row and column. Thus the level description
and character are unchanged. FC14 supplies both actual
\texttt{GLevH} records in these frames. Their first columns
satisfy $F_M\iota=\iota F_P$, and their final column is
$\phi_v(\pi e_{n+1})$. These are exactly the two
\texttt{FrExtFr} equations. They hold for every auxiliary place
and use the one $G$ shared across all words.

\paragraph{TG5: the terminal vector, lattices and residue functional.}
Put $z_0=\pi e_{n+1}$. For the actual raw terminal extension of a
row decoration with self-dual lattice $N_c$, its three lattices
are distinct:
\[
 \widetilde N_c=\iota N_c\oplus\mathcal Oe_{n+1},\qquad
 \widetilde L_c=\iota N_c\oplus\pi\mathcal Oe_{n+1},\qquad
 \widetilde L_c^{\#}=\iota N_c\oplus\bar\pi^{-1}
                                      \mathcal Oe_{n+1}.
 \tag{TG.11}
\]
The raw type functional is
$\widetilde f_P(x)=\operatorname{res}_P(\bar\pi x_{n+1})$.
These identities also follow directly from the definitions:
the row type set is empty, so its $N,L,L^{\#}$ coincide;
the extra $N$ coordinate is integral; the simultaneous kernels
of $\operatorname{res}_P(\bar\pi x_{n+1})$ require
$x_{n+1}\in\bigcap_{P\in S}P=\pi\mathcal O$ by FC1;
and Hermitian duality sends $\pi\mathcal O$ to
$\bar\pi^{-1}\mathcal O$. Thus the normalization is
$\bar\pi$, rather than $\pi$.

The source's chosen movement $m_c$ takes the preferred full
representative of $\widetilde c$ to this raw extension. Let
$\gamma_c$ be its rational unitary factor. The finite factors
leave the type kernel lattice unchanged. Consequently
\[
 \gamma_c L_{\widetilde c}
      =\iota N_c+\mathcal O z_0,
 \qquad h(z_0,z_0)=\bar\pi\pi=p.
 \tag{TG.12}
\]
The vector $z_0$ spans the specified orthogonal line. Indeed,
any nonzero vector $z_c$ orthogonal to $\iota E^n$ has zero first
$n$ coordinates, so $z_0=(\pi/(z_c)_{n+1})z_c$. Nonvanishing
of the last coordinate follows from $z_c\ne0$. This proves the
required span statement without changing the distinguished
$z_c$. Applied to the standard row base, (TG.12) is also the
exact \texttt{TilBaseLat} identity.

For the functional clause, let $x\in L_{\widetilde c}^{\#}$ and
suppose $a=h(z_0,\gamma_cx)\in\mathcal O$. Set $b=\bar\pi$.
The dual-lattice identities at the selected and conjugate primes
give $bL_{\widetilde c}^{\#}\subset N_{\widetilde c}$.
On this latter lattice, FC5 identifies the quotient projection
with the original decoration projection, while the period
extension has exactly the functional in (TG.11). Hence
\[
 \operatorname{res}_P(b)\,
   f_P(\operatorname{pr}_P x)
   =f_P(\operatorname{pr}_P(bx))
   =\operatorname{res}_P(ba).
 \tag{TG.13}
\]
Here the rational movement is the same $m_c$ as in (TG.12);
its finite factors do not rescale the type functional.
Since $\operatorname{res}_P(b)\ne0$, cancellation gives
\[
 f_P(\operatorname{pr}_P x)
       =\operatorname{res}_P\bigl(h(z_0,\gamma_cx)\bigr).
 \tag{TG.14}
\]
This is the full \texttt{TilTied} functional clause on the
Hermitian dual, not just on its sublattice $N$. The identical
calculation on the base lattice, with FC6's exact completed
functional, gives \texttt{TilBaseFv}.

\paragraph{TG6: the actual terminal adelic incidence and evaluation.}
For a row unitary matrix $y$, let
$\operatorname{ext}(y)=\operatorname{diag}(y,1)$. By (TG.11)
and the integral adelic-span identity of AD4, this extension
carries $\widehat{\mathcal O}^{n+1}$ to
$\widehat{\widetilde N}_c$ whenever $y=y_P(c)$.
Let $k$ be the actual basepoint matrix of AD1. Then
$\operatorname{ext}(y_P(c))k$ realizes the raw terminal
decoration, including its exact residue bases. At a type prime
the inverse last row of $k$ has residue
$\bar\pi e_{n+1}$, which gives the functional in (TG.11).
At auxiliary primes $k=1$ and $\pi\equiv1$, so the new last
residue coordinate $x_{n+1}/\pi$ equals $x_{n+1}$ modulo the
prime; the other coordinates are the original row residue basis.
There are no reference $Q$-frames to check. This proves all
\texttt{IsPtOf} clauses of the extension.

Write $y_M(\widetilde c)$ for the chosen full chart representative.
The exact level-reading relation, applied to the finite part of
the same movement $m_c$ used in TG5, gives
\[
 x_c=y_M(\widetilde c)^{-1}
       \operatorname{rat}(\gamma_c)^{-1}
       \operatorname{ext}(y_P(c))\in U_M,
 \qquad \chi_M(x_c)=\sigma(\chi_{\rm col}(m_c)).
 \tag{TG.15}
\]
The terminal coupling set is trivial, so this reading introduces
no further chart condition. Set $\kappa_c=x_c^{-1}$. Then
\[
 \operatorname{rat}(\gamma_c)y_M(\widetilde c)
       =\operatorname{ext}(y_P(c))\kappa_c,
 \qquad
 \chi_M(\kappa_c)^{-1}=\sigma(\chi_{\rm col}(m_c)).
 \tag{TG.16}
\]
These equations fix both inverse conventions.

They prove the literal \texttt{TilAd} statement for every matrix
$g$ satisfying the stated extension conditions. In fact those
conditions force $g=\operatorname{ext}(y_P(c))$: the equation
$g\iota=\iota y_P(c)$ determines the first $n$ columns, and
$g z_0=z_0$ determines the last column because the image of
the nonzero field element $\pi$ is a unit in every coefficient
algebra under consideration. Thus no alternative extension is
being chosen in \texttt{TilAd}.

Finally, on frame functions the source's terminal restriction is
defined using $m_c$ as
\[
 (\operatorname{ev}_{z_0}\varphi)(c,z)
  =\sigma(\chi_{\rm col}(m_c))\,
           \varphi(\widetilde c,\gamma_c^{-1}z).
 \tag{TG.17}
\]
With the unchanged frame identification and $eW_z=\mathrm{id}$,
if $z=\gamma_c z'$ then (TG.16)--(TG.17) give
$(eW_z\operatorname{ev}_{z_0}\varphi)(c,z)
=\chi_M(\kappa_c)^{-1}\varphi(\widetilde c,z')$.
This proves the source's exact \texttt{EvzTie}, together with
the prescribed identity $eW_z=\mathrm{id}$. TG5--TG6 therefore
prove \texttt{til\_\allowbreak ties} with its actual chosen $\gamma_c$,
vector and lattices, not an abstract substitute incidence.

\paragraph{TG7: construct the full simultaneous symmetric forms.}
Choose $G$ once as in FC13. At every step and every word, DR3
supplies the decorated realizers used in FC3. Construct those
charts and apply FC4's functional normalization, which preserves
their lattice, local and global levels, characters, representatives
and stabilizer maps. Keep their original chart labels. The finite
class law and surjectivity of the label map make each label set
finite; the terminal label set is the one-element set, because
its coupling data are empty. At the endpoints reindex this
one-element chart as a \texttt{ClsTied}.

We now assign every part of \texttt{FormsTiedSym} to a construction
already proved. FC1 supplies the exact type-prime data. CV9 and
FC7 give the endpoint and middle complex, rational, integral,
Petersson and norm clauses. FC5--FC6 give their quotient maps
and normalized completed residue functionals; FC17 gives their
integral basis comparison. NB7--NB8, with DR3 discharging their
initial realizer premise, give all up and down naturality clauses.
FC18 gives endpoint and middle dilation, including preservation
of the chart label. FC14 gives the middle reduction groups and
characters, and FC15 gives the prescribed step-zero group.
TG1--TG2 give the common penultimate group, character, frame and
finite orbit carrier. TG3 gives the actual row class/value/norm
objects. TG4 gives the terminal full and row levels and their
frame extension. TG5--TG6 give every terminal lattice, functional,
adelic-incidence and evaluation clause. FC8 supplies the period
map and its metric compatibility with these same frame functions.

All frame identifications are the fixed coordinate identifications
of SF2; thus the required \texttt{DGam} equality and the identity
$eW_z$ are literal equalities. The quotient maps retain the
prescribed formulas on $N_c$ and on the down kernel lattice,
so \texttt{StdLvC} holds, including its identity penultimate
frame. The step-zero and row auxiliary frames are identity for
every word, and every $G$ is the one chosen above. Its
\texttt{CharacterCompatible} law is FC13. These assignments
construct the exact conclusion of
\texttt{formsTiedSym\_\allowbreak std}: a simultaneous family of full
\texttt{FormsTiedSym}, \texttt{StdF}, \texttt{StdLvC}, finite
chart sets, a unique terminal chart, a constant $G$, and the
stated character compatibility. No terminal field remains an
independent existence premise.

\paragraph{TG8: the discriminant, valuation orientation and projections.}
For every valid class lattice $L_c$, put $H_c=L_c^{\#}$.
Then
\[
                     H_c/L_c\simeq_{\mathcal O}\mathcal O/(p).
 \tag{TG.18}
\]
To prove this, use the self-dual companion $N_c$ and choose a
local basis at a selected prime $P$ adapted to its nonzero type
functional. Then $L_{c,P}$ is obtained from $N_{c,P}$ by
replacing one summand $\mathcal O_P$ by $P\mathcal O_P$.
At $\bar P$, $L_{c,\bar P}=N_{c,\bar P}$. Hermitian duality
between the paired primes therefore gives
$H_{c,P}=N_{c,P}$, and at $\bar P$ exactly one enlarged dual
summand. Consequently $(H_c/L_c)_P$ and
$(H_c/L_c)_{\bar P}$ are one-dimensional residue modules;
the quotient vanishes at every other prime. The same calculation
shows $pH_c\subset L_c$. FC1 gives
$(p)=\prod_{P\in S}P\bar P$ with no repetitions. The Chinese
remainder idempotents decompose the finite $\mathcal O/(p)$-module
$H_c/L_c$ into these one-dimensional factors. Choosing one
nonzero vector in each factor gives (TG.18). This proves the
source's \texttt{DiscCl}; it is a module equivalence, with no
claim of an isometry of discriminant pairings.

The source's \texttt{SeenCl} has the original valuation
orientation. For a coefficient place $w$, pull its valuation back
along the prescribed embedding
$E\xrightarrow{\sigma_{\tau(w,v)}}E\xrightarrow{j_E}L$;
call it $\nu$. The \texttt{tauOf} identity says that the positive
ideal of its restriction to $\mathcal O_F$ is the prescribed
$v$. The positive ideal on $\mathcal O_E$ is a nonzero prime:
all integral elements have nonnegative valuation, and $p$ has
positive valuation. It is therefore a maximal prime over $v$.
By FC1 there are exactly the paired primes $P,\bar P$ over $v$.
The supplied \texttt{seenP} identity gives $\nu(\pi)>0$ in
the seen case, hence selects $P$; in the other case it gives
$\nu(\bar\pi)>0$ and selects $\bar P$. Thus
\[
 \nu(a)>0\quad\Longleftrightarrow\quad
 \begin{cases}a\in P&\text{if seen},\\a\in\bar P&\text{otherwise}
 \end{cases}\qquad(a\in\mathcal O).
 \tag{TG.19}
\]
Injectivity of the distinguished coefficient embedding identifies
the $L$-element appearing in \texttt{SeenCl} with this precise
pullback. This rules out changing either the residue prime or
the embedding in (TG.19).

Let $C_i$ denote the fixed coordinate isometry from the original
Hermitian space to $E^{N_i}$. Directly from the pairing,
$C_i(L^{\#})=(C_iL)^{\#}$. Linearity also gives
$C_i(P L)=P C_i(L)$, commutation with sums, and commutation with
localization: a witness $az\in L$, $a\notin P$, is sent to the
same witness $aC_i z\in C_iL$. Therefore the Hermitian
projections may be defined by
\[
 \operatorname{pr}^{H}_P(z)=\operatorname{pr}_P(C_i z),
 \qquad
 \operatorname{prk}^{H}_P(z)=\operatorname{prk}_P(C_i z).
 \tag{TG.20}
\]
They have exactly the surjectivity, residue-linearity and kernels
of FC5. In particular the two kernels are $PH_c$ and
$P(L_c+PH_c)$, respectively; the kernel of the composed type
functional is $L_c+PH_c$. The multiplication-by-$p$ relation
between the two quotients is preserved by $C_i$.

At a middle step the standalone and chart projections are equal
on the whole required domains. They agree by definition on
$N_c$ and on the original down lattice. For arbitrary
$z\in H_c$, multiply by $b=\bar\pi$; then $bz\in N_c$.
For $z\in L_c+PH_c$, the same multiplication lies in the
original down lattice, by its local kernel description at $P$
and the unchanged conditions at other primes. Apply the known
projection equality there and cancel the nonzero scalar
$\operatorname{res}_P(b)$. This proves both equalities in
\texttt{pr\_\allowbreak lvC\_\allowbreak eq\_\allowbreak lv}, including their integral domains.

\paragraph{TG9: Hermitian row carriers and every level-prime move.}
Pull each row lattice of TG3 back by $C_P=\operatorname{crdP}$,
and transport its unitary level and rational stabilizer matrices
through the same rank reindexing. The row base becomes the
original standard lattice $X'.L_0$, since its orthonormal frame
has integral coordinates. This gives the row class carrier,
with identical representatives, double quotient and stabilizer
range. Its level is hyperspecial away from the two auxiliary
primes. It contains the required principal congruence subgroup
at an auxiliary place: if $(x-1)$ sends the standard integral
lattice into $g_a$ times itself, testing on coordinate vectors
gives a matrix $x=I+g_a Z$ with $Z$ integral. Its residue is
$I$, which lies in the row residue group of TG4. Thus the
literal \texttt{StepLevel0}, \texttt{ClassTie} and
\texttt{SameCar} conditions all hold. The lattice identity is
$C_P(\mathrm{lattice}_2(c))=N_c$.

Transport the actual NB up and down matrices by
$\Gamma=C_i^{-1}\gamma C_i$. Write $L=L_c$, $H=L^{\#}$,
$K=L+PH$, and let primes and quotient vectors keep their
original meanings. For an up vector $z\in H$ with
$f_P(\operatorname{pr}_Pz)=1$, NB5 gives
\[
 (\gamma L_{c'})_P=H_P,\qquad
 (\gamma H_{c'})_P=(H+\mathcal O p^{-1}z)_P,
 \quad f'_{P}(\operatorname{pr}'_P x)
      =f_P(\operatorname{pr}_P(p\gamma x)).
 \tag{TG.21}
\]
The last identity is asserted on the exact domain
$x\in H_{c'}$, $p\gamma x\in H$. For a down functional
$g\ne0$, the corresponding formulas are
\[
 (\gamma H_{c'})_P=L_P,\qquad
 f'_{P}(\operatorname{pr}'_P x)
             =g(\operatorname{prk}_P(\gamma x))
       \quad(x\in H_{c'},\ \gamma x\in K).
 \tag{TG.22}
\]
NB6 gives in both cases the same moving matrix and a factorization
\[
       \operatorname{rat}(\gamma)y_{c'}=y_c\eta\kappa,
       \quad\kappa\in U_\beta,
       \quad\eta_v=1\text{ for }v\ne P\cap\mathcal O_F.
 \tag{TG.23}
\]
DR3 supplies its realizer hypothesis. At any other place,
(TG.23) and the level condition show that $\gamma$ carries the
new completed lattice to the old one and preserves its
normalized functional. Rational intersection then gives equality
of localizations off $P,\bar P$; FC1 identifies these as the
only two primes above the support place. The other type
functionals agree by the same completed equation and FC6.
Using (TG.20), and cancelling the coordinate isometry, proves
all clauses of the literal Hermitian \texttt{UpCl} and
\texttt{DnCl}.

The diamond case is equally constructive. For $a\in k_P^\times$,
choose an integer representing $a$. At the split base place,
take a unit whose residue is $a^{-1}$ and multiply the last
coordinate in the selected split factor by that unit and in
the other factor by its inverse. Set the adelic matrix $\delta$
to be the identity elsewhere. It is unitary, preserves the base
lattice, and satisfies $a f(\delta x)=f(x)$; thus it is an
actual \texttt{IsDia} element. FC18 gives the exact diamond
class and a factorization (TG.23). Both $\delta$ and $\kappa$
stabilize the completed base lattice, so AD4 gives
\[
 \gamma L_{c'}=L_c,\qquad
 f'_{P}(\operatorname{pr}'_P x)
        =a f_P(\operatorname{pr}_P(\gamma x)),
 \tag{TG.24}
\]
with equality of every other type functional. These are the
\texttt{DilCl} clauses after coordinate transport.

NB4 and FC18 preserve the middle chart label. At the terminal
step its uniqueness supplies the same assertion. Equation
(TG.23), reindexed without changing its factors, is consequently
the required \texttt{Ad0S} or \texttt{AdIS} relation for each
move. This proves all six clauses of
\texttt{SymmetricAdelicCompatibility}. For indices beyond $m$,
or move parameters not satisfying their defining nonzero
condition, choose the identity isometry; the source asserts no
move identity in those cases.

\paragraph{TG10: the terminal Hermitian carrier and stabilizer map.}
Define the Hermitian type-prime datum $A_u$ using the original
$\pi$, primes, residue maps and coefficient-place orientation
from FC1 and TG8; use projections (TG.20), the move isometries
of TG9, and the row carrier just constructed. Put
\[
 z_0^H=C_m^{-1}z_0,\qquad
 \widetilde\Gamma_c=C_m^{-1}\gamma_c C_m.
 \tag{TG.25}
\]
Equations (TG.12)--(TG.14), together with
$C_m\iota_H=\iota C_P$, prove the Hermitian terminal lattice
and functional clauses. Equation (TG.18) gives its discriminant
clause and (TG.19) its seen-prime clause. Coordinate transport
proves all remaining quotient and multiplication clauses as in
TG8. Thus $A_u$ is the actual \texttt{AtT}, not just selected
maps without their laws.

Every one of the nine \texttt{CarT} requirements is now fixed:
$\pi$ and the residue map are unchanged; the row carrier and
row-lattice identity are TG9; the two endpoint projections are
(TG.20); the middle projection is TG8's cancellation identity;
and the vector and matrix are (TG.25). This also proves the
source's \texttt{hermitianSide\_\allowbreak classes\_\allowbreak eq\_\allowbreak formClasses}
constructor. The same coordinate construction applies to any
supplied $H$ with the literal \texttt{CarH} equalities, so the
\texttt{hermitian\_\allowbreak data\_\allowbreak at\_\allowbreak level\_\allowbreak prime} conclusion does
not require an additional choice of compatible carriers.

For the stabilizer clause, let $s$ stabilize the preferred row
class $c$. Its raw extension has rational matrix
$\operatorname{diag}(\gamma_s,1)$, so it intertwines $\iota$
and fixes $z_0$. The source's stabilizer map is conjugation by
the actual movement $m_c$:
\[
 \operatorname{stabMap}_c(s)=m_c^{-1}\operatorname{ext}(s)m_c.
 \tag{TG.26}
\]
Taking rational factors yields
\[
 \gamma_c\,\gamma_{\operatorname{stabMap}_c(s)}\,
                      \gamma_c^{-1}
       =\operatorname{diag}(\gamma_s,1).
 \tag{TG.27}
\]
Transport (TG.27) through $C_m,C_P$. On the hyperplane it is
the row stabilizer action, and on $z_0^H$ it is the identity.
These are exactly the two
\texttt{SymmetricStabilizerMapCompatibility} equations.

For the source's formulation with an arbitrary supplied
\texttt{FormsTiedSym} satisfying \texttt{StdF}, the same
$\gamma_c$ is forced, rather than assumed. Its evaluation
clause has $eW_z=\mathrm{id}$ and the standard frame map.
Test that clause on the indicator of one invertible frame.
The nonzero character scalar forces
$\operatorname{frMul}(\gamma_{m_c}^{-1}\gamma_c)z=z$.
Cancel the invertible frame and use injectivity of one prescribed
field embedding to obtain $\gamma_c=\gamma_{m_c}$. Thus
(TG.27) holds in that formulation too. With TG9, this proves
the full source conclusion: an $A$ with \texttt{CarT},
symmetric adelic compatibility and symmetric stabilizer
compatibility, under the stated supplied data.

\paragraph{TG11: normalize the level-zero functional and all residue ties.}
There is one functional choice in passing from forms to a
Hermitian level carrier that \texttt{SameCar} alone does not
specify. We make it explicit. At level zero replace the
Hermitian target module and functional by the corresponding
forms target module and functional, pulled back through the
fixed rank reindexing. This preserves its base lattice, every
local and global level, representatives, stabilizer maps,
auxiliary splittings and auxiliary frames. The stabilizer of
the functional is preserved because the forms level has
exactly that stabilizer, and \texttt{SameCar} identifies its
local group. Nonvanishing, cardinality, the $p$-annihilation
law and vanishing on the base lattice are transported as well.
Thus this is a valid replacement of a \texttt{ColLevel},
with all the equalities required by
\texttt{stepZeroHermitianLevel\_\allowbreak with\_\allowbreak formsFunctional}.

The new target module $M_v$ is annihilated by the actual
local image of $\pi$. To check it without assuming an extra
linearity property of the additive equivalence $M_v\simeq k_P$,
choose a class and a dual-lattice vector whose projected type
functional is nonzero. Such a vector exists by surjectivity of
the projection and nonvanishing of that functional. Let
$m_0\ne0$ be the corresponding value of the completed base
functional. Multiplying the vector by $\pi$ multiplies its
projection by $\operatorname{res}_P(\pi)=0$, so the same
functional comparison gives $\pi m_0=0$. The additive group
of $M_v$ has prime order $p$; every nonzero element generates
it. Since scalar multiplication commutes with integer
multiplication, $\pi M_v=0$. This is precisely the source's
\texttt{pi\_\allowbreak kills}, and proves the replacement's additional
annihilation conclusion.

For clarity, the residue comparison can also be constructed for
any carrier satisfying that annihilation conclusion, not only
the replacement just chosen. At $v\mid p$ write
$S_v=\operatorname{intV}(E,v)$. The prescribed split orientation
identifies it with the product of the two complete local integer
rings. The selected component of $\pi$ is a uniformizer and
the other is a unit. Consequently
\[
 S_v/(\pi)\simeq\mathbb F_p,\qquad
 L_{\pi,v}^{\#}/L_{\pi,v}\simeq S_v/(p),
 \tag{TG.28}
\]
where the second quotient is generated over $S_v$ by the
class of $\bar\pi^{-1}e_{\mathrm{last}}$. Both nonzero completed
functionals in question vanish on $L_{\pi,v}$, take values
in groups of order $p$, and are killed by $\pi$ on the target.
They therefore factor through the one-dimensional
$S_v/(\pi)$-quotient of (TG.28). Each induced map is an
isomorphism of additive groups. Composing one with the inverse
of the other gives the unique additive equivalence of their
targets which identifies the two functionals. This proves the
source's cyclic-kernel argument directly, with its selected
prime orientation; an arbitrary abstract group of order $p$
would not determine that orientation by itself.

The resulting $fMe_v$ gives the exact level-zero residue identity
on every class: if $y_c t=C_0z$ locally, then
\[
        f_P(\operatorname{pr}^{H}_P z)=fMe_v(fT_v(t)).
 \tag{TG.29}
\]
For intermediate charts define the residue functional on their
completed Hermitian dual by first applying the coordinate
reindexing, then the forms functional, then its prescribed
$fMe$. The equation $y_ct=C_i z$ is equivalent, after this
reindexing, to $t=y_c^{-1}C_i z$. FC6 and TG8 therefore prove
the identical relation (TG.29) there. At the terminal step the
unique chart identifies this with its endpoint functional.
These are all the \texttt{SymmetricResidueFieldLevels}
clauses; no linear-algebra existence theorem for an unspecified
functional is being substituted for their pointwise equations.

Finally set the level-zero auxiliary generators to the original
$X.\mathrm{piG}$. Their designated-subfield membership, norms,
split orientation and scalar-diagonal splitting laws are FC13.
The frame-lattice equivalence is FC15 transported by the same
rank reindexing, and the last two fields are the annihilation
law and (TG.29). This constructs the literal \texttt{ColLevelX}
for every word, hence \texttt{SymmetricLevelZeroData}, with
the exact identity $\mathrm{piG}=X.\mathrm{piG}$ for every word.

\paragraph{TG12: the complete symmetric column comparison.}
Construct the forms family by TG7. FC9 transports its charts to
a Hermitian carrier $H$, and TG8--TG10 construct $A$ with
\texttt{CarT} and both symmetric compatibility laws. Now apply
TG11's level-zero functional replacement. Every \texttt{CarH}
field survives, as do $\mathrm{CarT}$ and the two compatibility
laws: their level-zero clauses use only the unchanged group,
representatives and stabilizer maps. The type-prime indexing,
residue fields, other levels and class data are also unchanged,
so the same $A$ remains well typed. This order makes clear that
the functional normalization is not an input to its own
Hermitian carrier construction.

Two further displayed identities finish the Hermitian column
record. The supplied \texttt{Pins.cof} equation, with the
unchanged $\pi_u$, is exactly
\[
                 c_u\pi_u=\pi_{u_0}.
 \tag{TG.30}
\]
For each endpoint or chart base lattice, SF2 and the standard
base-lattice formula give $L_\pi$ as in (TG.9), in the
appropriate dimension. Its dual is
$\mathcal O^{d-1}\oplus\bar\pi^{-1}\mathcal O$, so
coordinatewise
\[
               L_\pi+\bar\pi L_\pi^{\#}=\mathcal O^d.
 \tag{TG.31}
\]
This is SF3's companion identity, and the right side is the
coordinate image of the original standard lattice. Equation
(TG.31) is precisely \texttt{SymmetricBaseCompatibility},
including the version with the principal ideal
$(\bar\pi)$ multiplying the dual. Together with TG11, these
prove the actual remaining conclusions of
\texttt{herm\_\allowbreak rest\_\allowbreak tied}: cofactor compatibility, base
compatibility, all residue-field levels and level-zero data
with the prescribed auxiliary generators.

Take these objects as the fields of \texttt{ColsHermS}, with
SF2's \texttt{StdTied}, and use the TG7 forms as its forms
family. The carrier equalities are FC9 and TG10; all projection,
vector and rational-matrix equalities are (TG.20), (TG.25) and
their stated chart versions. The splitting equality and
step-zero frame equality come from the same FC9 transport.
All auxiliary generators agree because TG11 used
$X.\mathrm{piG}$. All $G$ objects agree because FC13 was chosen
once. Their character compatibility is FC13. The step-zero
frames are identity by FC15, and the row frames are identity
by TG4, so both required frame-constancy assertions hold for
every pair of words.

This constructs the source's exact endpoint
\[
 \exists\,C:\mathrm{ColsTiedSym}(D),\qquad
       I.\mathrm{StdC}(l,Z,C)\ \land\
       \forall u\ I.\mathrm{StdLvC}(l,Z,u,C.\mathrm{forms}(u)).
 \tag{TG.32}
\]
It is the conclusion of \texttt{colsTiedSym\_\allowbreak std}. Each
substantive supplier in this final join has been given above
or proved in the cited earlier claims. In particular neither
\texttt{FormsTiedSym}, \texttt{CarT}, initial decorated-point
realizability, nor terminal incidence remains an independent
input to (TG.32). The conclusion still concerns the supplied
arithmetic host and true-column data with their stated laws
and \texttt{Later}. It does not construct that arithmetic
host or the later analytic package, and does not by itself
assert RH or GRH. The stronger \texttt{OneLevelSymmetricColumns}
assertions about all earlier chart groups and singleton chart
sets are separate from (TG.32).



\section{Completion of the actual parameter record}
\label{sec:parameter-completion}

\status{PN2--PN15 prove the fourteen parameter families retained after RF16.
Together with RF12--RF14 and RF16, PN16 constructs every field of the
literal \texttt{ParameterChoice'} record, with its prescribed constant
list, under the supplied \texttt{Val} and \texttt{HostArith} hypotheses.
The threshold is uniform in the actual datum and size datum. This closes
\texttt{explicitConstants\_\allowbreak parameterConditions\_\allowbreak large}; the construction
of the host, columns and initial size datum remains a separate input.
The \texttt{SizesAt} and \texttt{NumAt} joins are treated in the following
component.}

\paragraph{PN1: common data, exact scales, and the order of thresholds.}
Fix $0<\epsilon<1/52$ and a standard constant record $C$. All assertions
in this section have the source's \texttt{SzRow} order
\[
 \exists c_{D,0}\ \forall c_D\geq c_{D,0}\
 \exists c_{h,0}\ \forall c_h\geq c_{h,0}\
 \exists x_0\ \forall D,L:\quad
 \left.
 \begin{array}{c}
 L.\mathrm{sc}.\epsilon=\epsilon,\quad L.\mathrm{sc}.x\geq x_0,\\
 \operatorname{Val}(C,c_D,c_h,L)
 \end{array}\right\}\Longrightarrow\hbox{the stated conclusion}.
                                                               \tag{PN.1}
\]
Here $L$ is an actual \texttt{SizeData\textsubscript{2}} over the same
actual $D$. Its windows and record identities are available; no new
geometric object is chosen in a numerical proof. The eventual threshold
$x_0$ may depend on $\epsilon,C,c_D,c_h$, but not on $D,L$ or on the
auxiliary scale parameter $J$.

Use the notation
\begin{gather*}
 n=D.R.n,\quad f=2n,\quad r=k+m-1,\quad
 N=N_m=k+m(t+1),\quad \mathcal L=x^\epsilon,\quad\ell=\log x,\\
 b=e^{\mathcal L},\quad a=xb,\quad X'=e^{x^{\epsilon/3}},\quad
 b'=\frac{xb}{X'},\quad w=\frac{b}{nX'},\quad X''=xnX'=\frac{a}{w},\\
 \Delta_{\rm cl}=\lceil x^{c_D}\rceil_+,\quad
 H=2\lceil x^{c_h}\rceil_+,\quad
 S=4k\Delta_{\rm cl},\quad\eta=mH,\quad h_0=H,\quad
 d_2=\frac{|T_2|}{n},\quad\lambda=L.\mathrm{LaG}.
                                                               \tag{PN.2}
\end{gather*}
Thus $S=L.\mathrm{Sc}$ and $\eta=L.\mathrm{etaSz}$, by the actual
\texttt{Sc\_\allowbreak eq}, \texttt{Del\_\allowbreak eq} and \texttt{eta\_\allowbreak eq} identities.
The scalar $\Delta_{\rm cl}$ is unrelated to an identification of the
integral Weyl operator $\Delta$ with its dual $\nabla$.
Set
\begin{align*}
 A_*&=nt^2(N+t)\mathcal L^2+kn^2(S+t+\mathcal L),\\
 \ell_*&=\frac{mb}{A_*},\qquad
 \ell_A=\frac{k\mathcal L}{md_2},\qquad
 r_S=\frac m2-1,\qquad M=r+mt=N-1.
                                                               \tag{PN.3}
\end{align*}
These are the literal \texttt{lStar}, \texttt{lA}, \texttt{rS} and
\texttt{Mrt}; $m$ is even and at least two. Also
$\texttt{qSig}=m$. All denominators in (PN.3) are positive.

RF12 gives, after a threshold depending only on $\epsilon$,
\[
 1\leq n\leq f\leq x^3,\quad
 2\leq m\leq x,\quad x^{1/4}\leq m,\quad
 k\leq r\leq x^5,\quad
 1\leq t\leq x^{401},\quad N\leq x^{402},\quad
 1\leq\mathcal L\leq x.                                    \tag{PN.4}
\]
The weaker lower bound $m\geq x^{1/4}$ follows from RF12's exact
$m\geq x^{1/2-\epsilon/16}$. The supplied second-set bound gives
$d_2\geq x^{-\epsilon}\geq x^{-1}$, while $d_2\leq1$ is already a
size-datum field. Consequently $\ell_A\leq x^7$. RF14 gives
$\mathcal L\leq\lambda\leq B\mathcal L\ell^2$ for a fixed $B>0$;
RF11 therefore gives $1\leq\lambda\leq x$ eventually.
For $c_D,c_h\geq0$ and $x\geq8$, rounding gives
\[
 x^{c_D}\leq\Delta_{\rm cl}\leq2x^{c_D},\qquad
 x^{c_h}\leq H\leq4x^{c_h},\qquad
 S\leq x^{c_D+6},\qquad \eta\leq x^{c_h+3}.             \tag{PN.5}
\]

We repeatedly use the already proved RF11 comparison: for every fixed
real exponent $v$, every positive fixed coefficient, and
$1\leq n\leq x^3$, a sufficiently large common threshold makes any
prescribed polynomial in $x$ smaller than $b,b',w$; it also gives
$b'x^v\leq b$ and $x^vX'\leq b$. The coefficient is absorbed using a
positive exponent margin. These are uniform statements in the actual
$n,D,L$, not estimates obtained by specializing to one datum.

\paragraph{PN2: both host primes and the exact conductor-product bound.}
This proves family [1]. Put $C_8=C.P.C_{\rm pde}\geq1$,
$\zeta=\epsilon/J$, and $Q_c=QS\,q\,d=\operatorname{condE}(D)$.
The supplied \texttt{HostVal} gives $1\leq J\leq C_8$ and a real $Y$
with $x^{1/J}\leq Y\leq x$. Its \texttt{PairRow} clauses say
\begin{gather*}
 e^{Y^{\epsilon/5}/2}<q\leq e^{Y^{\epsilon/5}},\qquad
 e^{Y^{\epsilon/5}(\log Y)^2/2}<d
                 \leq e^{Y^{\epsilon/5}(\log Y)^2},\\
 \log Q_c\leq2Y^{\epsilon/5}(\log Y)^2.
\end{gather*}
Taking logarithms, and using $\ell/J\leq\log Y\leq\ell$, proves
\begin{gather*}
 \tfrac12x^{\zeta/5}<\log q\leq x^{\epsilon/5},\qquad
 \tfrac12x^{\zeta/5}(\ell/J)^2<\log d
                                  \leq x^{\epsilon/5}\ell^2,\\
 \log Q_c\leq2x^{\epsilon/5}\ell^2.                         \tag{PN.6}
\end{gather*}
For $\ell\geq C_8$, the lower bound for $\log d$ is at least the
lower bound displayed for $\log q$. Since
$\zeta\geq\epsilon/C_8$, RF11 makes
\[
 (2C.C_t+5)\ell<\tfrac12x^{\epsilon/(5C_8)}
                                      \leq\tfrac12x^{\zeta/5}.
\]
Thus $x^{2L.C_t+5}<q,d$, using $L.C_t=C.C_t=401$.
Requiring $x^{\epsilon/(30C_8)}\geq2$ also gives
\[
 x^{\zeta/6}\leq\tfrac12x^{\zeta/5}<\log q.            \tag{PN.7}
\]
All thresholds are uniform in $J$ because its upper bound is fixed.

By (PN.4),
\[
 \log(Q_c^{36}t^{150})
 \leq72x^{\epsilon/5}\ell^2+60150\ell
 \leq\tfrac12x^\epsilon                                  \tag{PN.8}
\]
eventually, by the positive exponent gaps in RF11. Exponentiating
gives $Q_c^{36}t^{150}\leq b^{1/2}$. The host inequalities
$p_{\rm de},p'_{\rm de}\leq Q_c^{C_8}$ similarly give
\[
 \log p_{\rm de},\ \log p'_{\rm de}
       \leq2C_8x^{\epsilon/5}\ell^2\leq\mathcal L.       \tag{PN.9}
\]
Finally $|L.\mathrm{logDerivativeAtOne}|\leq x^\epsilon$ is precisely
the supplied \texttt{Val.dlog} field. This completes the five-part
family. Here \texttt{condE} is the defined product $QSqd$; the proof
does not construct a host field or prove an analytic logarithmic-
derivative estimate from the numerical inequalities.

\paragraph{PN3: the four logarithmic and block-window inequalities.}
This proves family [4]. For $m\geq4$, evenness gives
$r_S=m/2-1\geq m/4$. RF11 and (PN.4) imply, eventually,
\[
 0\leq\log f\leq3\ell,\qquad
 9\ell^2\leq\tfrac14x^{1/4}\leq\tfrac m4\leq r_S.
                                                               \tag{PN.10}
\]
Hence $\log f\leq\sqrt{r_S}$ and
$(\log f)^2\leq r_S\leq m\leq m^2\mathcal L$.
The standard constants give $c_{LL}=C.P.c_{LL}>0$.
Since $r\geq k\geq x$, requiring $x\geq c_{LL}^{-1}$ gives
$c_{LL}r\geq1$ and $\log r\geq\ell$. Therefore
\[
 \log\log b=\epsilon\ell\leq\ell
                    \leq c_{LL}rf\log r.                    \tag{PN.11}
\]
The actual block expression, with $C.P.C_{nW}=10$, is
\[
 n_W(L,10)=10r\left(f+\frac{f^2(k-1)^3\mathcal L^2}{m}\right)
             \geq10rf\geq(4f+10)r,                         \tag{PN.12}
\]
because $f\geq2$. These are all four fields of
\texttt{szRow\_\allowbreak logf}; no conductor-field identity is needed for the
source field named \texttt{log\_\allowbreak sq\_\allowbreak conductor\_\allowbreak le}.

\paragraph{PN4: all five bounds involving the height scale.}
This proves family [5]. Choose $c_h\geq401$. Then
$t\leq x^{401}\leq x^{c_h}\leq\eta$.
The upper bound in (PN.5) and the rank bounds give
\begin{align*}
 \eta t^3fk^3&\leq x^{c_h+3+1203+3+15}=x^{c_h+1224},\\
 xt\eta&\leq x^{c_h+405},\\
 m(t+1)\eta&\leq x^{c_h+406},
\end{align*}
where $t+1\leq x^{402}$ after $x\geq2$.
RF11 makes $x^{c_h+1224}\leq b,w$, and also
$(c_h+3)\ell\leq x^{\epsilon/4}$. Since $\eta>0$,
\[
 \boxed{\quad
 \eta\leq\frac{b}{t^3fk^3},\quad
 \log\eta\leq x^{\epsilon/4},\quad
 t\leq\eta,\quad t\eta\leq\frac wx,\quad
 m(t+1)\eta\leq w.\quad}                                  \tag{PN.13}
\]
The height exponent was fixed before these exponential comparisons
were imposed. Increasing the eventual $x$ threshold does not change it.

\paragraph{PN5: the four coupled height and cloud bounds.}
This proves family [6]. For $c_D,c_h\geq0$ put
$Q=x^{c_D+c_h+402}$ and take $x\geq8$.
Then $k,m,r,N,\mathcal L\leq Q$, $\Delta_{\rm cl}\leq2Q$,
$H\leq4Q$, and consequently
\[
 S\leq8Q^2\leq Q^3,\qquad
 \eta\leq4Q^2\leq Q^3,\qquad d_2Q\geq1.
\]
RF11 supplies $w\geq Q^{10}$. Since $m\geq1$,
\[
 d_2mw\geq Q^9,\qquad w\geq1.
\]
The four desired comparisons now follow by ordinary multiplication:
\begin{align*}
 m(S+\eta)&\leq2Q^4\leq Q^{10}\leq w\mathcal L,\\
 N\eta&\leq Q^4\leq Q^9\leq d_2mw,\\
 \eta^2r\mathcal L&\leq Q^8\leq Q^9\leq d_2mw^2,\\
 kS&\leq Q^4\leq w.
                                                               \tag{PN.14}
\end{align*}
The first three are also exactly the corresponding coupled
\texttt{SizeEta} clauses; the fourth is the parameter field
\texttt{kS\_\allowbreak le}.

\paragraph{PN6: the cloud upper bounds and its logarithm.}
This proves family [8]. Using $k\leq x^5$ and
$\Delta_{\rm cl}\leq2x^{c_D}$ gives
\[
 4kS\leq32x^{c_D+10},\qquad
 4C.P.C_{\Delta G}k^2\Delta_{\rm cl}
                       \leq8C.P.C_{\Delta G}x^{c_D+10}.
\]
Take $x\geq\max(32,8C.P.C_{\Delta G})$ and use RF11 to obtain
$x^{c_D+11}\leq b'$. Because $\lambda\geq1$, both right-hand
sides are at most $\lambda b'$. Dividing the first by the positive
$k$ yields
\[
 4S\leq\frac{\lambda b'}k,\qquad
 4C.P.C_{\Delta G}k^2\Delta_{\rm cl}\leq\lambda b'.       \tag{PN.15}
\]
Also
$\log\Delta_{\rm cl}\leq\log2+c_D\ell
\leq\tfrac12x^{\epsilon/4}$ eventually, by RF11.
This checks the literal half factor in \texttt{logDel}.

\paragraph{PN7: actual integer window length and logarithmic depth.}
This proves family [9]. The supplied identities and standard constants
are
\[
 L_w=10\lceil405\ell\rceil_++2\lceil30/\epsilon\rceil_++22,
 \quad C.P.C_{Lw}=4455,\quad C.P.C_{j0}=405,
 \quad j_0=\left\lceil\log_5\frac{nN}{m}\right\rceil_+.
                                                               \tag{PN.16}
\]
For every natural $l\geq11$, choose the actual integer
$j=\lceil7\ell\rceil_+$. Since $l-1\geq10>e$,
\[
 \ell_A\leq x^7=e^{7\ell}\leq e^j\leq(l-1)^j.
\]
Moreover $j\leq\lceil405\ell\rceil_+$, so the exact displayed
formula for $L_w$ gives $2j+2\leq L_w$. The ceiling inequalities give
\[
 L_w<4050\ell+60/\epsilon+34\leq4455\ell
\]
once $\ell\geq60/\epsilon+34$.

The rank inequality $k+m-1+5\leq t$ implies
$N=k+m(t+1)\leq2mt$. Thus $1<nN/m\leq2x^{404}$ by (PN.4).
Using $\log5>1$ and $\log2\leq\log5$, for $\ell\geq2$,
\[
 j_0<\frac{\log2+404\ell}{\log5}+1
         \leq404\ell+2\leq405\ell.                         \tag{PN.17}
\]
The elementary inequality $\log5>1$ follows, for example, from the
exponential series bound $e<3<5$. This proves the full integer-witness,
window-length and depth family, with all rounding preserved.

\paragraph{PN8: the complete exact column and arithmetic tuple.}
This proves family [10]. The supplied \texttt{HCVal} gives integers
$B_\tau$ with
\[
 v_\tau(i)=(k-2-i)4\Delta_{\rm cl},\quad
 b_\tau(j)=B_\tau+jH,\quad e_\tau(j)=b_\tau(j)+H/2.      \tag{PN.18}
\]
At a top corresponding to $t_T\in TT$,
$B_\tau=b_T(t_T)+\lfloor b'/x\rfloor$, which gives the exact
\texttt{topPlace\_\allowbreak first\_\allowbreak bound}. Off the top set, \texttt{HCVal}
also gives
\[
 u_\tau(0)=(k-2)4\Delta_{\rm cl}+k4\Delta_{\rm cl}+\tfrac12,
 \qquad u_\tau(k-1)=-\tfrac12.
\]
Consequently
\[
 u_\tau(0)-v_\tau(0)=S+\tfrac12\leq2S,
 \qquad u_\tau(k-1)=v_\tau(k-2)-\tfrac12.                 \tag{PN.19}
\]
Here $S\geq4$ since the natural ceiling $\Delta_{\rm cl}$ is at
least one. Adjacent $b$ coordinates differ by exactly $H=h_0$, and
every difference $e_\tau(j)-b_\tau(j)$ is exactly $H/2$.
Thus the adjacent-gap interval $[h_0,h_0+1]$ and both centering
errors, including the last coordinate, hold with error zero.

For each $j$, the last \texttt{HCVal} clause supplies
$\operatorname{MUQ}(\varphi_\tau,e_\tau(j))$ for the actual embeddings
$\varphi_\tau=D.\mathrm{host}.\mathrm{emb}(D.eTau(\tau))$.
By definition it means that there are integers $a_\tau$ with
$e_\tau(j)=a_\tau$ and, for every root of unity $\xi$ in the host field,
\[
                    \prod_\tau\varphi_\tau(\xi)^{a_\tau}=1.
                                                               \tag{PN.20}
\]
An integer considered as a rational has numerator that integer.
Replacing $a_\tau$ by $\operatorname{num}(e_\tau(j))$ therefore gives
exactly the host's \texttt{MuCond} field \texttt{e\_\allowbreak mu}; no parity
condition is silently substituted for this root-of-unity character
condition.

The remaining exact fields are as follows. The equality
$\eta_F=6x^{-C.P.C_F}$ is \texttt{Val.etaF\_\allowbreak eq}, and the standard
$C.P.C_{\eta F}=6$ gives its required coefficient. The two congruences
$p_{\rm de},p'_{\rm de}\equiv1\pmod{15}$ and their inequalities with
both $g_i$ are the supplied \texttt{Host.pde\_\allowbreak mod} and
\texttt{Host.pde\_\allowbreak ne\_\allowbreak g}. Finally \texttt{HostVal} says
$cc_K=\mathcal O_K$, so its absolute ideal norm is one:
\[
       \log\operatorname{Nm}(cc_K)=0\leq\mathcal L,
       \qquad \operatorname{Nm}(cc_K)=1\leq C.P.C_{ccK}\,d,
                                                               \tag{PN.21}
\]
using $C.P.C_{ccK}=1$ and the supplied positive prime $d$.
This supplies all twelve parts of \texttt{szRow\_\allowbreak exact}. It transports
the given character and prime conditions; it is not their existence proof.

\paragraph{PN9: every low node is quantitatively far from every midpoint.}
This proves family [11]. Choose $c_D\geq402$, $c_h\geq7c_D$,
and $Q=x^{c_D}\geq28$. From (PN.4), each of $n,t,N,\mathcal L,k$
is at most $Q$, and $\Delta_{\rm cl}\leq2Q$. Thus
\[
 A_*\leq2Q^6+8Q^5+2Q^4\leq12Q^6,
 \qquad 2A_*+t/2+3\leq28Q^6\leq Q^7.                   \tag{PN.22}
\]
Since $b/\ell_*=A_*/m\leq A_*$, and $H/2=\lceil x^{c_h}\rceil_+$
is at least $Q^7$,
\[
                 2b/\ell_*+t/2+3\leq H/2.
\]
For integer indices $i,j$, (PN.18) gives
\[
 |b_\tau(i)-e_\tau(j)|=|i-j-\tfrac12|H\geq H/2,
                                                               \tag{PN.23}
\]
because an integer has distance at least $1/2$ from $1/2$.
This proves the required inequality simultaneously for every
$\tau$ and every pair of allowed indices.

\paragraph{PN10: all cross-weight ratios and the final-bottom bound.}
This proves family [12]. RF11 lets us require
\[
               \eta+1\leq w\leq b'\leq b,
               \qquad 8x^6b'\leq b,\qquad x\geq8.        \tag{PN.24}
\]
Indeed $\eta\leq x^{c_h+3}$, while $w$ dominates every fixed
power; $w=b'/(xn)$; and $b'x^7\leq b$ suffices for the second
inequality. The supplied top upper bound, $\delta\leq1$,
$\mathcal L,\lambda\leq x$, $n\leq x^3\leq x^4$ and $k\geq1$
give
\[
              b_T(t_T)\leq b'+2x^6b'.                       \tag{PN.25}
\]
Apply the three actual \texttt{SizeData\textsubscript{2}} windows.
At a top their lower endpoint is at least $b'$, and their upper
endpoint is
\[
 b_T(t_T)+b'/x+\eta\leq3b'+2x^6b'\leq5x^6b'\leq2b.
\]
At a thick place the window is $[b,b+\eta+1]\subseteq[w,2b]$;
at a thin place it is $[w,w+\eta]\subseteq[w,2b]$.
The three cases exhaust the actual positions. Therefore
\[
          w\leq w_j(\tau)\leq2b
                  \quad(1\leq j\leq m).                   \tag{PN.26}
\]
Since $X''w=xb$ and $x\geq2$, this implies
$w_j(\tau)\leq X''w_{j'}(\tau')$ for every permitted pair.

For the second assertion of the family, \texttt{HCVal} gives
$B_\tau\geq w$ in each case: at a top,
$B_\tau=b_T+\lfloor b'/x\rfloor\geq b'$; at a thick place,
$B_\tau\geq\lceil b\rceil$; at a thin place,
$B_\tau=\lceil w\rceil$. Therefore
\[
 b_\tau(m-1)=B_\tau+(m-1)H
                \leq B_\tau+\eta\leq2B_\tau=2b_\tau(0).
                                                               \tag{PN.27}
\]
This proof uses the already supplied weight windows. It does not use
the desired parameter record to create an initial size datum.

\paragraph{PN11: all cloud-unit and count lower bounds.}
This proves family [14]. Write
$c=C.P.C_\Delta$, $F=C.P.C_F$ and $K=C.P.C_{\rm Count}$;
standardness gives $c,F,K\geq1$. Choose
\[
                 c_D\geq5c+Fc+406K+412,
                 \qquad x\geq\max(1,c).                    \tag{PN.28}
\]
The five literal left-hand sides are bounded as follows:
\[
\begin{array}{c|c}
 \text{left-hand side}&\text{upper bound}\\ \hline
 k^c(cm)^3/\eta_F^c&x^{5c+6+Fc}\\
 cmN&x^{404}\\
 k^c(cm^4\mathcal L)^4&x^{5c+24}\\
 Nm^{10}&x^{412}\\
 (Nf\mathcal L)^K&x^{406K}
\end{array}                                                   \tag{PN.29}
\]
For the first row, use $k^c\leq x^{5c}$,
$(cm)^3\leq x^6$ and
$\eta_F^c=(6x^{-F})^c\geq x^{-Fc}$.
For the third, $cm^4\mathcal L\leq x^6$.
The other three rows follow directly from (PN.4).
Every displayed exponent is at most the right side of (PN.28),
and hence every row is at most $x^{c_D}\leq\Delta_{\rm cl}$.
This proves all four conjuncts of \texttt{Del\_\allowbreak ge} and the separate
\texttt{Del\_\allowbreak count} field with their actual constants and real powers.

\paragraph{PN12: all ten node-height bounds, with a stronger first row.}
This proves family [15]. The following calculation also permits the
coefficients in the two marked rows to vary with $x$: assume
$0\leq A_r,A_k\leq x$, $x\geq200$, and put $Y=x^{c_D}\geq1$.
Choose
\[
                  c_h\geq\max(1300+c_D,406K),
                  \qquad K=C.P.C_{\rm Count}.                \tag{PN.30}
\]
Then $H\geq x^{1300}Y$. The two parts of $A_*$ satisfy
\[
 nt^2(N+t)\mathcal L^2\leq2x^{1209},\qquad
 kn^2(S+t+\mathcal L)\leq3x^{412}Y,
\]
because $S\leq x^6Y$ and $t,\mathcal L\leq x^{401}Y$.
In particular $A_*\leq5x^{1299}Y$, so
\[
 12A_*+2t+8\leq70x^{1299}Y\leq x^{1300}Y\leq H.
                                                               \tag{PN.31}
\]
This is stronger than the literal first row,
$12b/\ell_*+2t+8=12A_*/m+2t+8\leq H$.
The remaining nine left-hand sides have the explicit bounds
\[
\begin{array}{c|c}
 kt\mathcal L&x^{407}\\
 A_r r^{19}f^5\mathcal L&x^{112}\\
 (Nf\mathcal L)^K&x^{406K}\\
 4N^3+2t+6&12x^{1206}\leq x^{1207}\\
 A_k k m^{9/2}&x^{11}\\
 2M^3+t+3&6x^{1206}\leq x^{1207}\\
 120r\,q_{\rm sig}^{3/2}&120x^7\leq x^8\\
 Nm^{10}&x^{412}\\
 16q_{\rm sig}N&16x^{403}\leq x^{404}.
\end{array}                                                   \tag{PN.32}
\]
Here $M=N-1\leq N$, $q_{\rm sig}=m$,
$m^{9/2}\leq m^5$ and $m^{3/2}\leq m^2$ because $m\geq1$.
Every row is at most $H$ by (PN.30). To obtain the exact parameter
family, take $A_r=C.P.C_{h0r}$ and $A_k=C.P.C_{h0k}$, enlarging the
$x$ threshold to dominate these fixed constants. For the size-condition
application, the same calculation allows
$A_r=A_k=x^{\epsilon/100}\leq x$ without changing the exponent budget.
No inequality involving the desired parameter record has been used.

\paragraph{PN13: the additional cubic-rank and degree-power bounds.}
This proves family [16]. Let $K_N=C.P.C_{h0N}\geq1$ and choose
\[
                   c_h\geq\max(1210,3+402K_N).
                                                               \tag{PN.33}
\]
For $x\geq4$, the rank bounds imply
\[
 k+m(t+r+7)\leq x^5+3x^{402}\leq4x^{402}\leq x^{403}.
\]
It follows that
\[
 t+2\{k+m(t+(r+6)+1)\}^3+1
       \leq4x^{1209}\leq x^{1210}\leq H.                   \tag{PN.34}
\]
Also $fN^{K_N}\leq x^{3+402K_N}\leq H$.
These are the two parameter fields
\texttt{heightBound\_\allowbreak ge\_\allowbreak rank\_\allowbreak cubic} and
\texttt{heightBound\_\allowbreak ge\_\allowbreak degree\_\allowbreak power}; they are distinct from
RF15's four \texttt{ParamsT} height fields.

\paragraph{PN14: the full rank-bottom-cloud inequality.}
This proves family [18]. Use RF11 to require
\[
              b'x^{14}\leq b,
              \qquad x^{c_D+20}\leq b.
\]
Since $S\leq x^{c_D+6}$, both $b'$ and $S$ are at most $b/x^{14}$.
The uniform bounds $r\leq x^5$, $\ell\leq x$, $\ell_A\leq x^7$
give
\[
 r(b'+S)\ell\ell_A
       \leq x^{13}\frac{2b}{x^{14}}
       =\frac{2b}{x}\leq d_2mb,                           \tag{PN.35}
\]
where the last inequality uses $d_2\geq x^{-1}$ and $m\geq2$.
This proves the actual \texttt{rbS} field. The same inequality is the
last clause of \texttt{SizeEta}, since its \texttt{szlA} is $\ell_A$.

\paragraph{PN15: both strict lower-scale inequalities.}
This proves family [19]. Set
\[
 Q=x^{c_D+c_h+402},\qquad P=Q^7,
\]
and take $x\geq100$. Each of $n,t,N,\mathcal L,k,m$ is at most $Q$,
$\Delta_{\rm cl}\leq2Q$ and $H\leq4Q$. As in (PN.22),
\[
 0<A_*\leq2Q^6+8Q^5+2Q^4\leq P,\quad
 S\leq8Q^2\leq P,\quad\eta\leq4Q^2\leq P,\quad t\leq P.
\]
RF11 gives both $P^2\leq b'$ and $b'P^2\leq b$ at one threshold.
Consequently
\[
 \frac b{\ell_*}=\frac{A_*}{m}\leq P,
 \qquad \ell_*=\frac{mb}{A_*}\geq\frac bP
                  \geq b'P\geq100b'.                      \tag{PN.36}
\]
Because $P\leq b'$ and $b'>0$,
\begin{align*}
 b'+S+\eta+b/\ell_*&\leq b'+3P\leq4b'<100b'\leq\ell_*,\\
 2\eta+t+8b/\ell_*+7&\leq11P+7\leq12P
                                      \leq12b'<100b'\leq\ell_*.
                                                               \tag{PN.37}
\end{align*}
This preserves the strict signs in both source fields.
Since \texttt{szlStar} and \texttt{lStar} have the same definition,
these are also precisely the two \texttt{SizeLstar} inequalities.

\paragraph{PN16: construct all 111 fields of the compatible parameter record.}
Write $c=C.P.C_\Delta$, $F=C.P.C_F$, $K=C.P.C_{\rm Count}$ and
$K_N=C.P.C_{h0N}$. One sufficient common initial choice is
\begin{align*}
 c_{D,0}&=\max\{800,\ 5c+Fc+406K+412\},\\
 c_{h,0}(c_D)&=\max\{1000,\ 7c_D,\ 1300+c_D,\
                                 406K,\ 3+402K_N\}.
                                                               \tag{PN.38}
\end{align*}
For every $c_D\geq c_{D,0}$ and $c_h\geq c_{h,0}(c_D)$, take the
maximum of the finitely many eventual thresholds in PN1--PN15 and
RF12--RF14, RF16. This is allowed by (PN.1): every threshold precedes
the choice of $D,L$. The same chosen $D,L$ then satisfies all twenty-one
families. Their allocation is exact:
\[
\begin{array}{c|c@{\qquad}c|c}
 \text{family}&\text{proof}&\text{family}&\text{proof}\\ \hline
 1&\mathrm{PN2}&2,3&\mathrm{RF12}\\
 4&\mathrm{PN3}&5&\mathrm{PN4}\\
 6&\mathrm{PN5}&7&\mathrm{RF14}\\
 8&\mathrm{PN6}&9&\mathrm{PN7}\\
 10&\mathrm{PN8}&11&\mathrm{PN9}\\
 12&\mathrm{PN10}&13&\mathrm{RF13}\\
 14&\mathrm{PN11}&15&\mathrm{PN12}\\
 16&\mathrm{PN13}&17&\mathrm{RF13}\\
 18&\mathrm{PN14}&19&\mathrm{PN15}\\
 20&\mathrm{RF16}&21&\mathrm{RF14}.
\end{array}                                                   \tag{PN.39}
\]

Copy each of the following 23 real constants from $C.P$ into $P$:
\begin{gather*}
 C_{kLow},C_{k6},C_{fL},c_{LL},C_{nW},C_\rho,C_{\Delta G},
 C_{Lw},C_{j0},C_w,c_w,C_\Delta,\\
 C_{\rm Count},C_{h0r},C_{h0k},C_{h0N},C_{\rm gap},C_{\eta F},
 C_Y,C_{YL},C_{\rm pde},C_{ccK},C_F.
                                                               \tag{PN.40}
\end{gather*}
All are positive. For $c_{LL}$ this follows from
$1\leq c_{LL}^{-1}$ in the standard last-constant list; the other
nonfixed entries are at least one there or in the middle list.
The remaining fixed entries are
$C_{nW}=10$, $C_{Lw}=4455$, $C_{j0}=405$, $c_w=9/10$,
$C_{\eta F}=6$, $C_{YL}=9$ and $C_{ccK}=1$, with
$C_{\rm pde}\geq1$ supplied explicitly by standardness.
Thus the record's 23 positivity fields are also supplied.

The twenty-one families provide exactly 61 further named fields,
including fields that contain several quantified inequalities.
The remaining four fields are exactly the four clauses of the
explicit hypothesis $\operatorname{HostArith}(D.\mathrm{host})$:
\begin{enumerate}
\item for each $i$, an integral $\pi\in K$ with
      $\pi\overline\pi=g_i$;
\item $g_i\equiv1\pmod{QSqd}$ for each $i$;
\item the exact full splitting cardinalities above both auxiliary primes;
\item every supplied level prime has class $\mathrm{splitK}$.
\end{enumerate}
They become \texttt{g\_\allowbreak norm}, \texttt{g\_\allowbreak mod},
\texttt{pde\_\allowbreak split} and \texttt{level\_\allowbreak split}, respectively.
PN8's numerator conversion supplies the literal \texttt{e\_\allowbreak mu} field.
The ledger's field map checks every name in
\texttt{ParameterChoice'} exactly once:
\[
              23+23+61+4=111.                              \tag{PN.41}
\]
This is a concrete record, with the actual geometric data, embeddings,
weights and primes left in place. Since its constants were copied in
the order (PN.40), both required equalities hold exactly:
\[
                  \operatorname{parConsts}(P)=C.PK,
                  \qquad P.C_{nW}=C.P.C_{nW}.                \tag{PN.42}
\]
We have therefore proved the entire source conclusion
\[
 \operatorname{SzRow}\!\left(
 C,D,L\longmapsto
 \operatorname{HostArith}(D.\mathrm{host})\Longrightarrow
 \exists P:\operatorname{ParameterChoice}'(D,L),\
 \operatorname{parConsts}(P)=C.PK\ \land\ P.C_{nW}=C.P.C_{nW}
 \right).                                                    \tag{PN.43}
\]

\paragraph{The remaining existence boundary.}
No parameter family is retained after PN16. RF15 can now be applied to
this actual compatible $P$, with its already proved 54 \texttt{ParamsT}
fields; the numerical join with \texttt{SizesAt} is handled separately.
The hypothesis here is still an actual initial $D,L$ satisfying
\texttt{Val} and \texttt{HostArith}, including its rank windows,
host-prime windows, column identities, root-of-unity condition and
weight windows. Numerical estimates do not supply that host or those
columns. The four independent global theta/spectral, metric/dimension,
cohomological and final prime-window branches keep their earlier status.


\section{All size conditions and the initial arithmetic size datum}
\label{sec:size-conditions}

\status{SZ1--SZ8 prove the complete literal \texttt{SizesAt}
package, with the nested uniform thresholds of \texttt{SzRow}.
Together with PN and RF they construct every field of \texttt{NumAt}
in SZ9. SZ10--SZ12 independently construct the initial
\texttt{SizeData\textsubscript{2}} from the stated arithmetic and
Harish--Chandra data, and transport it through the actual
\texttt{Over} construction with its exact host and word-code identities.
They neither choose an arithmetic host nor construct
\texttt{ChOK}, \texttt{Over}, or \texttt{RecOK}.
The initial zero numerical fields are subsequently replaced by RF's
canonical refill.}

\paragraph{Notation and the order of the choices.}
In SZ1--SZ9 fix $0<\epsilon<1/52$, a standard constant record $C$,
and the same actual $D,L$ as in PN and RF, with
$\operatorname{Val}(C,c_D,c_h,L)$. To avoid confusing the size record
with its logarithmic scale, write
\[
 \begin{gathered}
 \mathcal L=x^\epsilon,\quad \delta=x^{-\epsilon/8},\quad
 \theta=\epsilon/100,\quad z=\epsilon/J,\quad
 N=N(m),\quad r=k+m-1,\quad f=2n,\\
 \Delta=\lceil x^{c_D}\rceil_+,\quad
 h=2\lceil x^{c_h}\rceil_+,\quad S=4k\Delta,\quad \eta=mh,\\
 \Lambda=L.\mathrm{LaG},\quad d_2=|T_2|/n,\quad
 N_* =2n\delta,\quad B=n x^{-\epsilon},\\
 Q=Q_Sqd,\quad \ell_A=\frac{kx^\epsilon}{md_2},\quad
 D_*=nt^2(N+t)\mathcal L^2+kn^2(S+t+\mathcal L),\quad
 \ell_* =\frac{mb}{D_*}.
 \end{gathered}                                                   \tag{SZ.1}
\]
Here $b,a,b',w'',X',X''$ are exactly (RF.1). In particular,
\emph{$\ell_*$ has the numerator $mb$}, so that
$b/\ell_*=D_*/m$. All natural quantities in real inequalities are
coerced to $\mathbb R$.

Every assertion called uniform below has the order
\[
 \forall\epsilon\ \forall C\ \exists c_{D,0}\
 \forall c_D\geq c_{D,0}\ \exists c_{h,0}\
 \forall c_h\geq c_{h,0}\ \exists x_0\
 \forall D,L.                                                    \tag{SZ.2}
\]
The last quantifiers require the fixed $\epsilon$, $x\geq x_0$ and
\texttt{Val}. In particular $x_0$ cannot depend on $D,L$, the supplied
primes, $J$, or a chosen position. We take finite maxima of RF11's
power/exponential thresholds and those in PN. Since
$1\leq J\leq C_8:=C.P.Cpde$, every occurrence of $z$ in a threshold
can be bounded below by $\epsilon/C_8$. This keeps (SZ.2) intact.

\paragraph{SZ1: the seven basic rank, density and length conditions.}
Put $\alpha=1/2-\epsilon/16$, $\beta=2+8\epsilon$, and
$\gamma=9/2+10\epsilon$. The literal rank windows say
\[
 x^\beta\leq n\leq2x^\beta,\qquad
 x^\alpha\leq m\leq2x^\alpha,\qquad
 x^\gamma\leq k\leq2x^\gamma.                                  \tag{SZ.3}
\]
Thus $x^\theta\geq2$ proves \texttt{SizeN}, \texttt{SizeM}, and
\texttt{SizeK}; the square root in the middle window is exactly
$x^\alpha$. The rank axiom and $L.\mathrm{t\_win}$ give
$1\leq t\leq x^{C_t}\leq x^{C_t+\theta}$, proving
\texttt{SizeT}. The host's two bounds
$\delta n/2\leq|TT|\leq2\delta n$ give
\[
       x^{-\theta}\delta n\leq |TT|\leq N_*.
\]
This is \texttt{SizeTT}. Its other density clause says
$x^{-\epsilon}n\leq |T_2|$; divide by $n>0$ and use the actual
\texttt{typeDensity\_\allowbreak eq\_\allowbreak card\_\allowbreak div} to obtain
\texttt{SizeDe2}. Finally,
\[
 L_w=10\lceil405\log x\rceil_+
             +2\lceil30/\epsilon\rceil_++22
       \leq4050\log x+60/\epsilon+34.
\]
For example $\log x\geq5000+100/\epsilon$ makes the right side
at most $(\log x)^2$. This proves \texttt{SizeLw}, with no rounding
term discarded.

\paragraph{SZ2: all fourteen coupled rank conditions.}
We use $\log\mathrm{rd}_E\leq\mathcal L$, established independently
in SZ3 below. Set $\zeta=\epsilon/1000$. For large $x$, (SZ.3)
implies
\[
 m\leq x^{\alpha+\zeta},\quad f\leq x^{\beta+\zeta},\quad
 k\leq x^{\gamma+\zeta},\quad
 100+8|C_\Gamma|\leq x^\zeta.                                 \tag{SZ.4}
\]
The following table verifies the first eight clauses of
\texttt{SizeKRel}. A listed exponent bounds the whole left side
above by that power of $x$; the right exponent is a lower bound for
the right side. For the fourth row the displayed inequality is the
square of the source inequality. All numbers in the table are exact
terminating rational decimals.
\begin{center}\small
\begin{tabular}{p{0.47\textwidth}cc}
Required inequality & upper exponent & lower exponent\\\hline
$x^\theta m^5f\mathcal L x^\epsilon\leq k$
 & $4.5+9.7035\epsilon$ & $4.5+10\epsilon$\\
$8C_\Gamma m^5f\mathcal L\leq k$
 & $4.5+8.6945\epsilon$ & $4.5+10\epsilon$\\
$k\leq\sqrt{x}\,nB$
 & $4.5+10.001\epsilon$ & $4.5+15\epsilon$\\
$k\leq mB^2$
 & $4.5+10.001\epsilon$ & $4.5+13.9375\epsilon$\\
$m^2\mathcal L\leq k$
 & $1+0.877\epsilon$ & $4.5+10\epsilon$\\
$m\sqrt{x}\log x\leq k$
 & $2-0.0615\epsilon$ & $4.5+10\epsilon$\\
$x^\theta m^6\mathcal L\leq k\log k$
 & $3+0.641\epsilon$ & $4.5+10\epsilon$\\
$x^\theta m^3\mathcal L\log\mathrm{rd}_E
                  \leq k(\log k)^2$
 & $1.5+1.8255\epsilon$ & $4.5+10\epsilon$
\end{tabular}
\end{center}
The sixth row uses $\log x\leq x$, and the last two use
$\log k\geq1$. Each exponent gap is positive for the entire allowed
range of $\epsilon$. Since $B>0$, the fourth row is equivalent to
$\sqrt{k/m}\leq B$.

For the remaining six clauses, $m$ is even and
$r_S=m/2-1$. The elementary logarithmic comparison from RF11 gives
$\log f\leq x^{1/5}$, whereas (SZ.3) eventually gives
$m\geq4x^{2/5}$ and hence $r_S\geq x^{2/5}$.
Therefore
\[
 \log f\leq\sqrt{r_S},\qquad
 (\log f)^2\leq r_S\leq m^2\mathcal L.
\]
Next $\log\log b=\log\mathcal L\leq\mathcal L$ and
$x^\theta\mathcal L\leq f$, by comparing exponents in (SZ.3).
As $r\geq k\geq5$, $\log r,\log k\geq1$, proving
\[
 x^\theta\log\log b\leq rf\log r,
       \qquad x^\theta\mathcal L\leq f\log k.
\]
Lastly $f=2n\leq m^2n$ because $m\geq2$, and eventually
$f\leq k\leq r\leq r^2$ by (SZ.3). These are the literal last two
clauses. Thus no clause of \texttt{SizeKRel} is omitted.

\paragraph{SZ3: prime, conductor and root-discriminant conditions.}
Choose the $Y$ already supplied by \texttt{HostVal}; no prime is
chosen here. Its window gives $x^{1/J}\leq Y\leq x$, and
\texttt{PairRow} gives
\[
 \begin{split}
 \tfrac12 x^{z/5}&<\log q\leq x^{\epsilon/5},\\
 \tfrac12 x^{z/5}(\log x/J)^2
             &<\log d\leq x^{\epsilon/5}(\log x)^2,\\
 \log Q&\leq2x^{\epsilon/5}(\log x)^2.
 \end{split}                                                    \tag{SZ.5}
\]
At a uniform threshold $\log x\geq C_8$ and
$x^{\epsilon/(100C_8)}\geq2$. Thus $\log x/J\geq1$ and
$x^{z/100}\geq2$. The two lower bounds in (SZ.5) give
$x^{19z/100}\leq\log q,\log d$, and
$z/6<19z/100$ gives the additional $q$ lower bound.
Absorbing $(\log x)^2$ into $x^\theta$ gives both upper bounds
$x^{\epsilon/5+\theta}$.

The supplied host's actual derived discriminant facts, rather than a
new assertion about the existence of that host, give
\[
 1\leq\mathrm{rd}_E\leq\frac{Q_S}{2}\sqrt{qd}\leq Q.
                                                                    \tag{SZ.6}
\]
For completeness the lower host bound is
$2^{e+1-2^{-e}}\sqrt d$, which is at least one because $e\geq0$
and $d$ is a positive prime. The upper bound in (SZ.6) is exactly
$H.\mathrm{der.rd\_le}$. Consequently
\[
 0\leq\log\mathrm{rd}_E\leq2x^{\epsilon/5}(\log x)^2
                                     \leq\mathcal L.           \tag{SZ.7}
\]
Together with \texttt{Val.dlog}, these prove all eight clauses of
\texttt{SizeQD}. The host also gives
$p_{\mathrm{de}},p'_{\mathrm{de}}\leq Q^{C_8}$. Hence
\[
 \log p_{\mathrm{de}},\log p'_{\mathrm{de}}
      \leq2C_8x^{\epsilon/5}(\log x)^2\leq\mathcal L,
\]
which proves \texttt{SizePde}.

RF12 gives $t\leq x^{401}$. Take logarithms in the required
conductor inequality. The sufficient inequality is
\[
 36\log Q+150\log t
    \leq72x^{\epsilon/5}(\log x)^2+60150\log x
    \leq\tfrac12 x^\epsilon=\log\sqrt b.                       \tag{SZ.8}
\]
RF11 proves its last step with a uniform threshold. Exponentiation
proves $Q^{36}t^{150}\leq\sqrt b$, the complete
\texttt{SizeCond}. Here $Q$ is the literal conductor expression
in \texttt{szCond}; no further conductor identity is presumed.

\paragraph{SZ4: the exact lambda, rho and bottom windows.}
RF14 proves, from the rounded orbit count, constants $A_C>0$ and
uniform thresholds such that
\[
        x^\epsilon\leq\Lambda
                    \leq A_Cx^\epsilon(\log x)^2.
\]
Absorb $A_C(\log x)^2$ into $x^\theta$; this proves
\texttt{SizeLa}. The coefficient $x^\theta$ in
\texttt{SizeRho} is variable, so one must not substitute it into a
bound proved only for a fixed coefficient. Instead use (SZ.3) directly:
\[
 4N_*x^\epsilon\Lambda
       \leq16x^{2+9.885\epsilon},\qquad
 x^\theta\log x\Lambda
       \leq x^{1.02\epsilon}\log x.                            \tag{SZ.9}
\]
Each is at most $\frac12x^{2.5+10\epsilon}$ for all sufficiently
large $x$. Their sum is at most $k x^{-2}$ by (SZ.3).
Division by $k>0$ gives precisely \texttt{SizeRho}.
Finally \texttt{Val.bT\_\allowbreak le}, with $N_*=2n\delta$, is exactly
\texttt{SizeBot}; no enlargement of the bottom window is needed.

\paragraph{SZ5: the eta, node and strict-star conditions on the same datum.}
PN4 proves the five eta inequalities, PN5 the four coupled eta/cloud
inequalities, and PN14 the rank-bottom-cloud inequality. In the
literal order of \texttt{SizeEta}, the supplied entries are
\[
 (\mathrm{PN4}_3,\mathrm{PN4}_1,\mathrm{PN4}_2,
   \mathrm{PN4}_4,\mathrm{PN4}_5,
   \mathrm{PN5}_1,\mathrm{PN5}_2,\mathrm{PN5}_3,\mathrm{PN14}).
                                                                    \tag{SZ.10}
\]
Thus the ninth clause includes the exact $\ell_A$, $d_2$ and
$b'+S$ factors; it is not inferred just from a small-$\eta$ estimate.

PN12 proves the following strengthened bounds on $h=\eta/m$:
\[
 \begin{gathered}
 12D_*+2t+8\leq h,\quad kt\mathcal L\leq h,\quad
 A_r r^{19}f^5\mathcal L\leq h,\quad
 4N^3+2t+6\leq h,\\
 A_k k m^{9/2}\leq h,\quad
 2(r+mt)^3+t+3\leq h,\quad Nm^{10}\leq h,
 \qquad 0\leq A_r,A_k\leq x.
 \end{gathered}                                                    \tag{SZ.11}
\]
Take $A_r=A_k=x^\theta$, since $0<\theta<1$ and $x\geq1$.
Since $m\geq1$ and $D_*>0$,
$12b/\ell_*=12D_*/m\leq12D_*$. This proves all seven
clauses of \texttt{SizeNodes}, including the nonintegral power
$m^{9/2}$. Finally PN15 proves
\[
 b'+S+\eta+b/\ell_*<\ell_*,\qquad
                2\eta+t+8b/\ell_*+7<\ell_*.
\]
These are exactly \texttt{SizeLstar}. Each application keeps the
same $D,L,c_D,c_h$ and uses the maximum of the finitely many
thresholds; it is not a selection of different parameters for different
clauses.

\paragraph{SZ6: all eight spacing conditions.}
Take $c_D\geq412$ and $x\geq32$. RF12 gives
$m\leq x$, $k\leq x^5$, $N\leq x^{402}$; also
$\theta<1$ and $\Delta\geq x^{c_D}$. Therefore
\[
 Nm^{10}\leq x^{412}\leq\Delta,\qquad
 x^\theta mN\leq x^{404}\leq\Delta.
\]
The ceiling upper bound $\Delta\leq2x^{c_D}$ and RF11 imply
\[
        \log\Delta\leq\log2+c_D\log x
                                   \leq\tfrac12x^{\epsilon/4}.
                                                                    \tag{SZ.12}
\]
At the scale-comparison threshold for the fixed exponent $c_D+12$,
$x^{c_D+12}\leq w''\leq b'$. Hence
\[
             16xk^2\Delta
                    \leq32x^{c_D+11}\leq x^{c_D+12}\leq w''.
                                                                    \tag{SZ.13}
\]
As $m\leq k$, $x^\theta\leq x$, $S=4k\Delta$ and
$\Lambda\geq1$, (SZ.13) proves the four required inequalities
\[
 kS\leq w'',\quad x^\theta mk\Delta\leq w'',\quad
          4S\leq\Lambda b'/k,\quad
          4x^\theta k^2\Delta\leq\Lambda b'.
\]
For the last clause, the host's literal $v$ is
$v=\mathrm{rd}_E/2>0$. Thus $\log v\leq\log\mathrm{rd}_E
\leq\mathcal L\leq x$ by SZ3. Therefore
$x^\theta m\log v\leq x^3\leq k$. This remains true if
$\log v$ is negative; multiplication uses only positive factors.
All eight clauses of \texttt{SizeDel} are now proved.

\paragraph{SZ7: the actual weight envelope, prime scales and word count.}
For fixed $c_h$, RF11's exponential comparison gives, for
$c=\max(c_h,0)+10$,
\[
 \sqrt b\,x^c\leq w'',\qquad b'x^c\leq b,
       \qquad\eta+1\leq w''\leq b'.                           \tag{SZ.14}
\]
The last inequality follows also directly from
$\eta=2m\lceil x^{c_h}\rceil_+\leq4x^{1+\max(c_h,0)}$,
after increasing the threshold. The first part gives
\texttt{SizeW}. The proof of SZ4 shows that the bottom surplus
$2n\delta x^\epsilon\Lambda/k$ tends to zero; alternatively
$\Lambda\leq x^2$, $f=2n\leq x^3$ and $k\geq1$ bound it by
$x^6$. It follows from (SZ.14) that
\[
          b_T+b'/x\leq(x^6+2)b'\leq b.
\]
Use the three \emph{actual} size-record weight windows. At a top,
$b_T+b'/x\leq w_j\leq b_T+b'/x+\eta$; at a thick position,
$b\leq w_j\leq b+\eta+1$; at a thin position,
$w''\leq w_j\leq w''+\eta$. They give the common envelope
\[
               w''\leq w_j(\tau)\leq2b
             \quad(1\leq j\leq m,\ \hbox{every }\tau).       \tag{SZ.15}
\]
Since $X''w''=xb$ and $x\geq2$, (SZ.15) proves every
cross-position, cross-index clause of \texttt{SizeXpp}.

Write $\beta_\tau=(D.hc\,\tau).b_0$. The \texttt{HCVal}
formulas give $b_{m-1}=\beta_\tau+(m-1)h$ and
$\beta_\tau\geq w''$: top, thick and thin positions give,
respectively, $\beta_\tau\geq b'$, $\beta_\tau\geq b$, and
$\beta_\tau=\lceil w''\rceil$. Thus
\[
       b_{m-1}\leq\beta_\tau+\eta\leq2\beta_\tau,
\]
which is exactly \texttt{SizeBm}.

The host's interval $(kb)^8\leq g_i\leq2(kb)^8$ proves both
the lower bound and $g_i\leq2g_j$ in \texttt{SizeG}. Moreover
\[
 \log g_i\leq\log2+8\log k+8\mathcal L
          \leq40\log x+9\mathcal L+9
          \leq x^\theta\mathcal L.
\]
The last inequality is another fixed-logarithm comparison. RF16 gives
the complete \texttt{SizeWords} bound
$|\mathrm{Word}|\leq2^{k-3}$, and RF14 gives
$11\leq\mathrm{splitBound}$. These statements concern the supplied
primes and words; no new prime or code has been selected.

\paragraph{SZ8: assemble the complete literal \texttt{SizesAt}.}
The definition has twenty-three named size clauses and one final
splitting bound. The following exhaustive map records their suppliers:
\begin{center}\small
\begin{tabular}{p{0.68\textwidth}l}
Clauses & supplier\\\hline
$N,M,K,T,TT,De2,Lw$ & SZ1\\
$KRel$ & SZ2\\
$Cond,QD,Pde$ & SZ3\\
$La,Rho,Bot$ & SZ4\\
$Eta,Nodes,Lstar$ & SZ5\\
$Del$ & SZ6\\
$W,Bm,Xpp,G,Words$ and $11\leq\mathrm{splitBound}$ & SZ7
\end{tabular}
\end{center}
Take the maximum of the $c_D$ thresholds, then, for that $c_D$,
the maximum of the $c_h$ thresholds, and finally the maximum of the
$x$ thresholds. This is precisely the quantifier order (SZ.2).
All conjunctions refer to the same objects. The resulting statement is
the full \texttt{szRow\_\allowbreak sizesAt}. The separately defined
\texttt{SizeTu} does not occur in \texttt{SizesAt}; it has not
silently been added to this theorem.

\paragraph{SZ9: every field of \texttt{NumAt} and the density/code pair.}
Assume additionally the literal \texttt{HostArith}. PN16 constructs
an actual $P:\texttt{ParameterChoice'}\,D\,L$ with
\[
          \operatorname{parConsts}(P)=C.PK,
                       \qquad P.CnW=C.P.CnW=10.
\]
Apply RF15 to \emph{this same $P$}. It constructs an actual
$T:\texttt{ParamsT}\,D\,L\,P$ with its prescribed seven constants
and $T.eP=C.T.eP$. SZ8 supplies the complete \texttt{sizes} field.
The remaining fields of \texttt{NumAt} are supplied as follows:
\begin{center}\small
\begin{tabular}{p{0.44\textwidth}p{0.48\textwidth}}
Field or group of fields & proof\\\hline
$\rho_B\leq x^{-2}$; $x^\epsilon\leq\Lambda$;
$g_i\leq e^{9\mathcal L}$; both $p_{\mathrm{de}}<X'$ bounds
 & RF14, on the same $D,L$\\
$C_\Gamma,C_c,C_t,CT_2$ equalities; bottom upper bound;
exact log-five depth & the corresponding fields of \texttt{Val}\\
$ccK=\top$ & \texttt{HostVal}\\
$L_w=10\lceil(C.Ct+4)\log x\rceil_+
                 +2\lceil30/\epsilon\rceil_++22$
 & $C.Ct=401$ and the literal $L_w$ equality\\
All positions have the same $v$ column
 & $v_i=(k-2-i)4\Delta$, independent of the position\\
The three raw rank windows & \texttt{Val.ranks}, exactly (SZ.3)
\end{tabular}
\end{center}
Together with the six $P,T$ data and constant fields and the
\texttt{sizes} field, this accounts for all twenty-one fields of
the actual record. RF14 also gives \texttt{ChTT\_\allowbreak at} and
$\texttt{ChCode\_\allowbreak at}(L,3,C.P.CF)$. Taking the finite maxima once
more proves
\[
 \texttt{HostArith}(D.host)\Longrightarrow
 \operatorname{Nonempty}(\texttt{NumAt}(C,L))
 \ \land\ \texttt{ChTT\_\allowbreak at}(L)
 \ \land\ \texttt{ChCode\_\allowbreak at}(L,3,C.P.CF)                    \tag{SZ.16}
\]
as an actual \texttt{SzRow} theorem. Neither full
\texttt{ParameterChoice'} nor \texttt{SizesAt} remains an
unproved numerical premise of this implication. Existence of the
arithmetic/geometric $D$ is a different assertion.

\paragraph{SZ10: a torsion estimate and all thirteen initial-window clauses.}
We now prove the initial windows \emph{without assuming an existing
size datum or PN16}. Fix a scale $s$, actual $D$, top bottom-degrees
$b_T$, and the literal \texttt{RankVal}, \texttt{HostVal}, and
\texttt{HCVal}, with
\[
 c_D\geq0,\quad c_h\geq402,\qquad
                  b'\leq b_T\leq2b'.                         \tag{SZ.17}
\]
All thresholds here depend only on $\epsilon,c_D,c_h,C_8$.

First let $F$ be a number field of degree $d_F$. A root of unity
of exact order $a$ has cyclotomic degree $\varphi(a)\leq d_F$.
For $a=\prod p^{e_p}$,
\[
  \frac{\varphi(a)^2}{a}
          =\prod_{p\mid a}p^{e_p-2}(p-1)^2\geq\tfrac12.
                                                                    \tag{SZ.18}
\]
Each odd-prime factor is at least one; the factor at $2$, if present,
is at least $1/2$. Thus $a\leq2d_F^2$. This first proves that all
roots of unity form a finite group, since their orders are bounded.
A finite subgroup of a field's multiplicative group is cyclic, so its
generator shows that its cardinality is also at most $2d_F^2$.
For the supplied host $[E:\mathbb Q]=2f$, giving
\[
             |\mu(E)|\leq8f^2\leq x^7\leq h/2             \tag{SZ.19}
\]
after RF12 and $x\geq8$. This supplies the actual torsion bound used
by the thick-position interval.

RF11--RF12, without a size record, give at a common threshold
\[
 k+5\leq b'/x,\quad \eta\leq b,\quad b'\leq b,\quad
 6\leq b,\quad1\leq t\leq x^{401},\quad nN/m\leq x^{405}.
                                                                    \tag{SZ.20}
\]
Here the crude $f\leq x^3$, $m\leq x$, $N\leq x^{402}$
bounds are used only where needed. By \texttt{HCVal}, writing
$\beta_\tau$ for its integral base value,
\[
 b_j=\beta_\tau+jh,\quad e_j=\beta_\tau+jh+h/2,
       \qquad w_j=e_{j-1}\ (1\leq j\leq m).                  \tag{SZ.21}
\]
In particular
$\beta_\tau+1\leq w_j\leq\beta_\tau+\eta-h/2$.
We check all thirteen clauses of the source window theorem:
\begin{enumerate}
\item At a top,
$b_0-b_T=\lfloor b'/x\rfloor>(k-1)/2+3$ by (SZ.20).
\item For every position and $0\leq j<m$, $b_j\leq b^2$.
Indeed (SZ.17) and the three possible base values imply
$\beta_\tau\leq3b+2$: tops use $2b'+b'/x$, thick positions
use $\lceil b\rceil+|\mu(E)|$, and thin positions use
$\lceil w''\rceil$ with $w''\leq b'$. Equations
(SZ.19)--(SZ.21) give $b_j\leq4b+2\leq b^2$ for $b\geq6$.
\item At a thick position the integral base satisfies
$\lceil b\rceil\leq\beta_\tau<\lceil b\rceil+|\mu(E)|$.
Since both endpoints are integers and $|\mu(E)|\leq h/2$,
$b\leq\beta_\tau\leq b+h/2$. Consequently
$b\leq w_j\leq b+\eta+1$, the required thick window.
\item At a top,
$\beta_\tau=b_T+\lfloor b'/x\rfloor$. The bounds
$b'/x-1<\lfloor b'/x\rfloor\leq b'/x$ and $h/2\geq1$
give $b_T+b'/x\leq w_j\leq b_T+b'/x+\eta$.
\item At a thin position $\beta_\tau=\lceil w''\rceil$.
The same $h/2\geq1$ gives $w''\leq w_j\leq w''+\eta$.
\item Off the top set, $u_0-v_0=S+1/2$.
Since $S=4k\Delta\geq1$, it lies in $[S,2S]$.
\item For every $i<k-1$,
$v_i=(k-2-i)4\Delta$ lies in $[0,S]$.
\item For $i,j<m$,
$|b_i-e_j|=|(i-j)h-h/2|\leq\eta$.
\item The rank axiom gives $1\leq t$.
\item RF12 gives $t\leq x^{401}$.
\item Set $q_0=\lceil\log_5(nN/m)\rceil_+$. Since
$\log5\geq1$, (SZ.20) gives
$q_0\leq\lceil405\log x\rceil_+$. Thus
$10q_0+20\leq L_w$ by its exact formula.
\item $|T_2|>0$ follows from $|T_2|\geq x^{-\epsilon}n>0$.
\item $|T_2|\leq n$ follows from the host's actual
\texttt{second\_\allowbreak set\_\allowbreak range}: its integer labels lie in the range
of length $n_0\leq n$ (the condition $L_w\geq2$ follows from
the exact formula). No density assumption stronger than the source
one is substituted.
\end{enumerate}
These give the complete \texttt{windows\_\allowbreak of\_\allowbreak explicitParameters}
conclusion. The multiplicative-unit condition in \texttt{HCVal}
is retained for subsequent parameter construction; none of these
thirteen window estimates needs to manufacture it.

\paragraph{SZ11: construct the initial size record field by field.}
Assume the literal inputs to \texttt{exists\_\allowbreak sizeData\_\allowbreak explicit}:
\texttt{Std}, $c_D\geq0$, $c_h\geq402$, the arithmetic and
column data of SZ10, the bounded logarithmic derivative, the prescribed
$\Delta$, $CT_2$, word count and $\eta_F$, and
\[
 \log\Gamma_u=C_\Gamma m^5f\mathcal L,
 \qquad
 b'\leq b_T\leq b'+2n\delta x^\epsilon\Lambda b'/k,           \tag{SZ.22}
\]
where now $\Lambda=G.\mathrm{La}(C.GK,s)$. Also retain the source's
$a_T$ window, the two top $u$ endpoints, and the untwisted equality
$A_j=0$. These are actual constructor inputs, not conclusions of
this step.

There is no circular use of SZ4's existing-size hypothesis here.
RF14's independent orbit-count bound gives
$\Lambda\leq A_Cx^\epsilon(\log x)^2$ directly for the scale.
Using $\delta\leq1$ and the raw rank windows,
\[
 \frac{2n\delta x^\epsilon\Lambda}{k}
       \leq4A_C x^{-5/2}(\log x)^2\leq1.                    \tag{SZ.23}
\]
Thus (SZ.22) implies (SZ.17), and SZ10 applies. Increase the
threshold so that $C_8\leq x^\theta$. Define the initial record by
\[
 \begin{array}{c|l}
 \text{fields}&\text{chosen values}\\\hline
 \mathrm{sc},a_T,b_T&s,a_T,b_T\\
 \mathrm{supp},\mathrm{drop}&\varnothing,\varnothing\\
 \mathrm{gain},F,E,\tau,\mathrm{io}&0,0,0,0,0\\
 \mathrm{NuLat},\mathrm{lenMN}&\perp,0\\
 j_0&\lceil\log_5(nN/m)\rceil_+\\
 d_2,\Lambda,S,\eta&|T_2|/n,G.\mathrm{La}(C.GK,s),4k\Delta,mh\\
 C_c,C_t,C_\Gamma&C.Cc,C.Ct,C.CGam\\
 TT_w,\mathrm{wPart}&TT,1
 \end{array}                                                     \tag{SZ.24}
\]
Take the logarithmic-derivative field to be the actual
$\operatorname{dlogL}(D,1)=\operatorname{dlogH}(D.host)$.
The $a_T,b_T$ and top-$u$ fields are the supplied ones. SZ10 proves
the top-bottom separation, upper $b_j$ bound, three weight windows,
cloud gap, $v$ bounds, $b/e$ closeness, depth and $t$ bounds.
For each word's middle $u_i$, strict interlacing gives
$v_i<u_i<v_{i-1}$; since these $v$ columns are unchanged by the
word, $0\leq u_i\leq S$. This supplies \texttt{cloud\_\allowbreak u} on
its exact range $1\leq i$, $i+1<k$.

All support, conjugacy, non-inert and above-band implications for
the chosen zero fields are immediate equalities. The untwisted field
uses the supplied $A_j=0$ and $F=0$. The empty drop set excludes
every type prime, and $TT_w=TT$ proves its fullness condition.
The splitting tie is the host field. The other two ties follow from
\[
 o\leq C_8\leq x^\theta,\qquad
             Q_S\leq C_8n\leq4nx^\theta.
\]
SZ10 gives $0<|T_2|\leq n$, proving $0<d_2\leq1$ and its
defining density equality. The positive $S,\eta,C_\Gamma$ and
unit lower bound on $\mathrm{wPart}$ are immediate from the chosen
values and \texttt{Std}. Finally the exact
\texttt{Gam\_\allowbreak eq} is the supplied equality in (SZ.22).
This accounts for every data and proof field of
\texttt{SizeData\textsubscript{2}}, not merely its numerical subset.

The fifteen fields of \texttt{Val} now follow from the supplied rank,
host, column, derivative, spacing, $CT_2$, word and $\eta_F$ data,
together with the defining equalities (SZ.24) and the bottom window
(SZ.22). Thus there is an actual $L$ with
\[
      L.\mathrm{sc}=s,\quad L.a_T=a_T,\quad L.b_T=b_T,
                         \operatorname{Val}(C,c_D,c_h,L).       \tag{SZ.25}
\]
This is an initial record: its zero gains and lattices are legal
values for the constructor, and RF5 later replaces them by the
canonical refill. No final forcing or lattice estimate is asserted
from the zeros.

\paragraph{SZ12: transport through \texttt{Over} with exact labels and code.}
Supply the literal \texttt{ChOK} datum $P$, its level record $Lv$,
\[
 Lv.\Gamma=\Gamma_S(C,s,P.R),\qquad
             \texttt{Over}(D,P,Lv),\qquad\texttt{RecOK}(s,C,c_D,D).
                                                                    \tag{SZ.26}
\]
The existential decomposition of \texttt{Over} identifies $D$ with
its actual staged construction. Its host has the same rank, $TT$
labels, second set, $q,d,g,p_{\mathrm{de}},p'_{\mathrm{de}}$ and
distinguished translate as $P.H$; its level-prime set is the old
one united with the prescribed level primes. In particular the old
set is contained in the new one, rather than necessarily equal to it.
The \texttt{Pre.ofCh} construction proves that every member of this
union is in class \texttt{splitK}. The other three
\texttt{HostArith} clauses are exactly the arithmetic part of
\texttt{ChOK}. Hence \texttt{HostArith}(D.host) is supplied.

Use SZ11 with $a_T(t)=a(t.\mathrm{val})$ and
$b_T(t)=b(t.\mathrm{val})$. The rank, host, column, derivative,
$CT_2$, word, degree-window and top-$u$ inputs come from
\texttt{ChOK}; the spacing, $\eta_F$ and untwisted inputs come
from \texttt{RecOK}. The staged identity
$\Gamma_u=Lv.\Gamma$ and (SZ.26) give
$\log\Gamma_u=C_\Gamma m^5f\mathcal L$ by $\log(e^y)=y$.
Thus (SZ.25) yields $L.\mathrm{sc}=s$,
\texttt{Val}, and the exact \texttt{WtEq}: injectivity of the
frame-label map identifies any two top positions having the same
label, so their two degrees are the prescribed $a,b$ values.

For clarity the simultaneous identities proved by this construction
are
\[
 \begin{gathered}
 D.R=P.R,\quad
 D.host.TT.\mathrm{image}(\mathrm{val})
          =P.H.TT.\mathrm{image}(\mathrm{val}),\quad
 D.host.T_2=P.H.T_2,\\
 (q,d,g,p_{\mathrm{de}},p'_{\mathrm{de}})_D
          =(q,d,g,p_{\mathrm{de}},p'_{\mathrm{de}})_P,\quad
 P.H.\mathrm{levelPrimes}\subseteq D.host.\mathrm{levelPrimes},\\
 D.host.u_0.\mathrm{val}=P.H.u_0.\mathrm{val},\qquad
 (G.h_P(D,l)).\mathrm{val}=(h_{P.H}(l)).\mathrm{val}.
 \end{gathered}                                                    \tag{SZ.27}
\]
The word types are identified by an actual equivalence $e_W$, which
is the identity after unfolding the stages. For $l<k-2$, compare
the source's two expressions for the same $u_l$ entry: the
\texttt{u\_\allowbreak letter} formula and the \texttt{ChOK.code} formula.
Cancel their identical base terms to obtain
\[
 \Delta\,\mathbf1_{\mathrm{letter}(u,l)}
             =\Delta\,\mathbf1_{c(e_Wu,l)}.
\]
Since $\Delta>0$, the two Boolean bits agree. For $l\geq k-2$
both extended bits are false by definition. Hence, for every $u,l$,
\[
 \operatorname{wordBits}(D,u,l)=
 \begin{cases}c(e_Wu,l),&l<k-2,\\ \mathrm{false},&l\geq k-2.
 \end{cases}                                                      \tag{SZ.28}
\]
This proves the complete literal \texttt{exists\_\allowbreak sizeData\_\allowbreak over},
including the labels and code, and discharges the initial-size
supplier retained in RF10. The arithmetic column constructor
\texttt{chGood\_\allowbreak of\_\allowbreak host}, the level-prime constructor and the
geometric \texttt{RecExP}/\texttt{RecOK}/\texttt{LedOK} supply
remain distinct tasks. SZ9 can now be applied to this $L$ when its
uniform thresholds are included; RF's refill preserves
\texttt{Val}, so it can also be applied to that refilled datum.

\paragraph{SZ13: construct the actual ranks over the prescribed frame.}
The numerical theorems above assume an actual rank record. We can also
construct that record over any supplied frame $\phi$ whose order $n$
satisfies the required window
\[
             x^{2+8\epsilon}\leq n\leq2x^{2+8\epsilon}.
                                                                    \tag{SZ.29}
\]
This preserves the given group, labels, projection and distinguished
generator; no replacement of the frame is made. The uniform threshold
$x\geq16$ suffices. Put
\[
 A=\sqrt{x\delta}=x^{1/2-\epsilon/16},\qquad
 K=x^{9/2+10\epsilon},\qquad
          m=2\lceil A/2\rceil_+,\quad k=2\lceil K/2\rceil_+.
                                                                    \tag{SZ.30}
\]
For any $y\geq2$, $y\leq2\lceil y/2\rceil_+<y+2\leq2y$.
Here $2\leq A\leq x$, $x^4\leq K$, and
$x^2\leq n\leq2x^3$, using $0<\epsilon<1/52$.
Thus $k,m$ are even, their two rank windows hold, and
\[
       2\leq m,\quad m+3\leq k,\quad m\leq2n\leq k.
                                                                    \tag{SZ.31}
\]
For example $m\leq2x$, $2n\geq2x^2$, and
$2n\leq4x^3\leq x^4\leq k$ prove the last two comparisons;
$2x+3\leq x^4$ proves the remaining gap.

Set $f=2n$, $r=k+m-1$, and define
\[
 \begin{split}
 T&=\lceil x^{377}\rceil_+,\\
 W&=\left\lceil\frac{10r}{8}
       \left(f+\frac{f^2(k-1)^3\mathcal L^2}{m}\right)
                                                    \right\rceil_+,\\
 t&=(8f-1)r+4+8TW.
 \end{split}                                                    \tag{SZ.32}
\]
The last expression is a nonnegative integer and equals the source's
natural subtraction $8fr+4+8TW-r$, because $f\geq1$.
Since $r$ is odd and $8f-1$ is odd, $t$ is odd. Moreover
\[
             t\geq7r+4\geq r+5,\qquad 3m<t,
\]
where $r=k+m-1\geq m+4$. In particular $m\leq t+1$, so
$m^2\leq k+m(t+1)$. These are all the additional inequalities
in the literal \texttt{Ranks} record. Together with the unchanged
$\phi$ and (SZ.31), they therefore construct an actual $R$.

The definitions of $T$ and $8W$ are precisely $\mathrm{TOv}(s)$ and
$\mathrm{nWv}(R,s)$. There is no circular choice of $t$: the latter
expression depends on $n,f,k,m,r$ and the scale, all fixed before
$t$ is defined. Equation (SZ.32) gives
\[
             r+t=8fr+4+\mathrm{TOv}(s)\mathrm{nWv}(R,s).
                                                                    \tag{SZ.33}
\]
This, (SZ.29), and (SZ.30) are exactly
$\mathrm{RankVal}(s,R)$, with $R.\mathrm{toFrame}=\phi$.
Thus the literal \texttt{exists\_\allowbreak ranks\_\allowbreak rankVal\_\allowbreak frame} is
proved, uniformly over all such prescribed frames.

There is also an explicit cyclic-frame specialization. Given $B\geq1$,
choose the least integer $e\geq0$ with $2^e\geq B$. Minimality
gives $2^e\leq2B$. For $B=x^{2+8\epsilon}$, (SZ.29) therefore
holds with $n=2^e$, and $e\geq2$. The cyclic frame uses
$A=\mathbb Z/2^e\mathbb Z$, generator $1$, and the homomorphism
\[
 A\longrightarrow
  (\mathbb Z/2^{e+2}\mathbb Z)^\times/\langle-1\rangle,
                            \qquad j\longmapsto[5^j].       \tag{SZ.34}
\]
It is an isomorphism. Indeed
$v_2(5^{2^j}-1)=j+2$, by induction from
$(u^2-1)=(u-1)(u+1)$ and $u\equiv1\pmod4$.
Thus $5$ has order $2^e$ modulo $2^{e+2}$, and its powers never
equal $-1$ there because they are $1$ modulo $4$. The quotient
also has order $2^e$, so (SZ.34) is bijective.
Transport this additive group to the source's labels
$\mathrm{Fin}(2^e)$ by the residue-class bijection. Its projection
is surjective, and a multiple of its distinguished generator whose
projection is zero is itself zero, because the projection is
injective. These are exactly the frame axioms.
The chosen isomorphism is the source's \texttt{cycIso}: a homomorphism
from a cyclic group is determined by its value at $1$.
Consequently this is the literal \texttt{cyclicFrame}(e), and
(SZ.30)--(SZ.33) prove \texttt{exists\_\allowbreak ranks\_\allowbreak cyclic\_\allowbreak frame}.
This constructs rank and frame data only; the number field and
arithmetic host with those labels remain separate existence problems.


\section{Construct the arithmetic columns and the level datum}
\label{sec:arithmetic-constructors}

This section proves the two conditional constructors retained in RF10:
\texttt{chGood\_\allowbreak of\_\allowbreak host} and \texttt{exists\_\allowbreak levCh\_\allowbreak of\_\allowbreak levelPrime}.
The first starts with the actual arithmetic host, its prescribed top weights,
and an injective Boolean code. The second starts with the analytic
\texttt{LevelPrime} assertion and the prescribed positive quantities
$\mathrm{du}_u,D_*$ satisfying their logarithmic window. We construct the
columns and the complete level record, including their integrality and
root-of-unity conditions. The existence of those input hosts, good top
weights, and analytic primes is not asserted here.

Write $n=R.n$, $f=2n$, $r=k+m-1$ for the literal rank record. For a scale
$s$ put
\[
 x=s.x,\quad \epsilon=s.\mathrm{eps},\quad
 \mathcal L=x^\epsilon,\quad b=\exp(\mathcal L),\quad
 X'=\exp(x^{\epsilon/3}),\quad
 a=xb,\quad b'=\frac{xb}{X'},\quad w=\frac{b}{nX'}.
                                                               \tag{AC.1}
\]
The letter $b_j$ below is a column coordinate, whereas $b$ in (AC.1)
is a positive real scale. In the arithmetic host let
$E=H.E$, $s_d=H.\mathrm{sqd}$, and
$K=\mathbb Q(s_d)$, where $s_d^2=-d$.
The literal host hypotheses give a prime $d\equiv7\pmod {16}$,
an isomorphism
\[
 \operatorname{Gal}(E/\mathbb Q)
       \simeq\{\pm1\}^2\times A,
 \qquad |A|=n,
\]
and the second sign is the action on $s_d$.
These statements, rather than an arbitrary choice of a CM type, are
used in AC3.

\paragraph{AC1: correct all roots of unity with one bounded shift.}
Let $E$ be a number field, $F$ any field, and
$\varphi_i:E\hookrightarrow F$ a finite family of field embeddings.
Fix one index $i_0$ and arbitrary integers $a_i$. Put
$w_E=|\mu(E)|$, the number of roots of unity in $E$.
There exists an integer $q$ with $0\leq q<w_E$ such that
\[
       \prod_i\varphi_i(z)^{a_i+q\,1_{i=i_0}}=1
                  \qquad(z\in\mu(E)).                       \tag{AC.2}
\]
Here $w_E$ is the literal \texttt{Units.torsionOrder}.
For completeness, roots of unity are finite: an element of order $h$
has cyclotomic degree $\phi(h)\leq[E:\mathbb Q]$, while
$\phi(h)^2\geq h/2$ by the prime factor calculation in SZ10.
There are thus finitely many possible orders and finitely many roots
for each order. A finite subgroup of a field's multiplicative group is
cyclic. Indeed, combining elements of maximal prime-power orders gives
an element whose order is the exponent $e$ of the subgroup; all its
elements solve $X^e-1=0$, so its order is at most $e$, and it is cyclic.
Choose a generator $\zeta$ of $\mu(E)$.

The element
$U=\prod_i\varphi_i(\zeta)^{a_i}$ satisfies $U^{w_E}=1$.
The $w_E$ distinct powers of $\varphi_{i_0}(\zeta)$ are all the roots
of $X^{w_E}-1$ in $F$. Hence
$\varphi_{i_0}(\zeta)^q=U^{-1}$ for one $0\leq q<w_E$.
This proves (AC.2) at $\zeta$. Every $z\in\mu(E)$ is $\zeta^j$,
and the left side at $z$ is the $j$th power of its value at $\zeta$.
This proves the full assertion for every root of unity, including
negative exponents. It is the substance of \texttt{exists\_\allowbreak shift}.

\paragraph{AC2: the quadratic-subfield norm of every root is one.}
We prove the following elementary statement, also present as
\texttt{norm\_\allowbreak root\_\allowbreak of\_\allowbreak unity\_\allowbreak quadratic\_\allowbreak subfield} in the source:
if $d>3$ is prime, $s_d^2=-d$ in a number field $E$, and
$4\mid[E:\mathbb Q]$, then
\[
                  N_{E/K}(z)=1\qquad(z\in\mu(E)),
            \qquad K=\mathbb Q(s_d).                         \tag{AC.3}
\]
This will establish the needed host norm property rather than simply
reuse the field \texttt{HostDer.norm\_\allowbreak mu}.

First, $[K:\mathbb Q]=2$. If $y=a+c s_d\in K$ has rational square,
the coefficient of $s_d$ in $y^2$ gives $2ac=0$. Thus a rational number
$u$ which is a square in $K$ is either a rational square or $-d$ times
a rational square. It follows that
\[
             -1,-2,-3,2\quad\hbox{are not squares in }K.
                                                               \tag{AC.4}
\]
For $-j$ with $1\leq j\leq3$, the only possible second case would be
$j=d c^2$ with $c\in\mathbb Q$. The exponent of the prime $d$ on the
right is odd and on the left is zero, a contradiction. The first
case is excluded by negativity. For $2$, the second case is negative
and the first contradicts the elementary irrationality of $\sqrt2$.

The only root in $K$ of odd order is $1$. Otherwise a prime divisor
$p$ of that order gives a primitive $p$th root in $K$. Its degree
$p-1$ is at most two, so $p=3$. Such a root $v$ obeys
$v^2+v+1=0$, whence $(2v+1)^2=-3$, contradicting (AC.4).

Next let $y^{2^a}=1$. If $y^2=1$, then $y=\pm1$ and
$N_{E/K}(y)=1$, since $[E:K]=[E:\mathbb Q]/2$ is even.
If $y^2\ne1$, some power $i$ of $y$ has $i^2=-1$.
Put $K_1=K(i)$, which is quadratic over $K$ by (AC.4), and put
$v=N_{E/K_1}(y)$. The order of $v$ is a power of two.
It cannot exceed four: if it did, a power $u$ of $v$ would satisfy
$u^4=-1$, giving
\[
                         (u-u^3)^2=2.
\]
An element of $K_1$ whose square lies in $K$ has square either a
square of an element of $K$ or its negative, by the same coefficient
argument in the basis $1,i$. We would obtain a square root of $2$
or $-2$ in $K$, again contradicting (AC.4).
Thus $v\in\{1,-1,i,-i\}$. Each of these four elements has norm one
from $K_1$ to $K$, since the nonidentity $K$-automorphism sends
$i$ to $-i$. Norm transitivity proves $N_{E/K}(y)=1$.

Finally write an exponent killing $z$ as $2^a o$ with $o$ odd.
The result for $y=z^o$ gives $N_{E/K}(z)^o=1$.
The odd-order result in $K$ gives (AC.3).
This proof requires neither a class-number assertion nor an
unramifiedness or prime-distribution assertion.

\paragraph{AC3: identify the CM-type product and permit parallel shifts.}
For the literal host, the Galois isomorphism above and the Galois
degree formula give $[E:\mathbb Q]=4n$. Thus AC2 applies.
Let $\Psi\subset\{\pm1\}^2\times A$ be the kernel of the second
sign, and let $e_T:T\simeq\Psi$ be any finite reindexing.
For $g\in\Psi$, the automorphism $\sigma_g=H.\mathrm{gal}^{-1}(g)$
fixes $s_d$ and hence $K$. Conversely an automorphism fixing $K$
fixes $s_d\ne0$, so its second sign is $+1$.
Consequently $g\mapsto\sigma_g$ is a bijection from $\Psi$ to
$\operatorname{Gal}(E/K)$. The ordinary norm formula gives
\[
       \prod_{\tau\in T}\sigma_{e_T(\tau)}(z)=N_{E/K}(z)
                     \quad\hbox{inside }E.                  \tag{AC.5}
\]
Compose with the fixed host embedding $\varphi_0:E\hookrightarrow
\mathbb C$. Its compositions with these automorphisms are precisely
$H.\mathrm{emb}(e_T(\tau))$. By AC2,
\[
             \prod_{\tau\in T}H.\mathrm{emb}(e_T(\tau))(z)=1
                          \qquad(z\in\mu(E)).               \tag{AC.6}
\]

Write $\mathrm{MU}(a)$ for the assertion that this product with
integer exponents $a_\tau$ is one for every root $z$.
Equation (AC.6) proves $\mathrm{MU}(1)$; raising it to any integer
power proves $\mathrm{MU}(c)$ for every constant integer $c$.
Multiplying the two root identities proves
\[
                 \mathrm{MU}(a)\Longrightarrow
                      \mathrm{MU}(a+c).                     \tag{AC.7}
\]
If rational weights equal $a_\tau+c$ with these integer exponents,
they satisfy the literal $\mathrm{MUQ}$ definition: the required
integral lift is exactly $a+c$. This proves
\texttt{muq\_\allowbreak of\_\allowbreak parallel}, including its integrality assertion.
The condition is a product identity for all roots of unity; parity
of the sum alone would not suffice.

\paragraph{AC4: construct every field of a coded column.}
Let $D,h\in\mathbb N$ satisfy
\[
        D\geq1,\qquad h\geq(t+3)/2,
\]
and let $U_0,U_L,B\in\mathbb Z$ satisfy
\[
             4D(k-2)\leq U_0,\qquad U_L<0,
                         -B\leq U_L.                        \tag{AC.8}
\]
For a Boolean word $c=(c_0,\ldots,c_{k-3})$, put
\[
 \begin{aligned}
 v_i&=4D(k-2-i) &&(0\leq i\leq k-2),\\
 u_0&=U_0+\tfrac12,\qquad
 u_{k-1}=U_L+\tfrac12,\\
 u_{\ell+1}&=v_\ell-\tfrac12-D\,1_{c_\ell=\mathrm{true}}
                                     &&(0\leq\ell\leq k-3),\\
 b_j&=B+2hj,\qquad e_j=B+(2j+1)h &&(0\leq j<m).
 \end{aligned}                                                 \tag{AC.9}
\]
These are the four actual data fields of an $\mathrm{HCParam}(R)$.
We verify its nine proof fields.

For each interior index, $v_\ell-u_{\ell+1}$ is $1/2$ or
$D+1/2$, while
$u_{\ell+1}-v_{\ell+1}$ is $4D-1/2$ or $3D-1/2$;
all four numbers are positive. At the first endpoint,
$u_0>v_0$ by (AC.8); at the last, $u_{k-1}\leq-1/2<0=v_{k-2}$.
These prove both strict interlacing fields. Also
$u_{k-1}=U_L+1/2>-B=-b_0$.
The gaps $e_j-b_j$ and $b_{j+1}-e_j$ both equal $h$, so both
required lower bounds are valid. Each $b_j,e_j,v_i$ is integral,
and each $u_i$ is an integer plus $1/2$.
This proves all fields, including the strict bottom inequality.
It proves \texttt{exists\_\allowbreak hcParam\_\allowbreak of\_\allowbreak code} by an explicit formula.

\paragraph{AC5: uniform numerical prerequisites for the column formula.}
Fix $0<\epsilon<1/52$, $c_D\geq0$, and $c_h\geq402$.
There is $x_0$, depending only on these three fixed numbers, such that
for every scale with this $\epsilon$, every $x\geq x_0$, and every
actual $R$ satisfying $\mathrm{RankVal}(s,R)$, set
\[
                  D=\lceil x^{c_D}\rceil_+,
             h=\lceil x^{c_h}\rceil_+,
                    h_0=2h.                                \tag{AC.10}
\]
Then
\[
 D\geq1,\quad (t+3)/2\leq h,\quad 4kD\leq a,
       \quad k\leq b'/x,\qquad b',b,w>0.                    \tag{AC.11}
\]
Indeed, use RF12 to obtain $k\leq x^5$ and $t\leq x^{401}$,
and require $x\geq8$. Since $D\leq2x^{c_D}$,
\[
       4kD\leq8x^{c_D+5}\leq x^{c_D+6}.
\]
RF11 lets us require $x^{c_D+6}X'\leq b$.
Since $X'\geq1$ and $a=xb$, this proves $4kD\leq a$.
It also proves $k\leq x^5\leq b/X'=b'/x$.
The estimate $h\geq x^{402}\geq(t+3)/2$ follows from
$t\leq x^{401}$ and $x\geq8$.
All positivity assertions follow from (AC.1) and $n\geq1$.
The threshold precedes the choices of $s,R$, the host, the code,
and the top weights. This is the exact uniformity of
\texttt{hc\_\allowbreak sizes}.

\paragraph{AC6: construct the whole column family with exact root congruences.}
Retain AC5, supply an actual host $H$ with
$\mathrm{HostVal}(s,C_8,H)$, and supply integers $a_T(t),b_T(t)$
on the host's top set satisfying
\[
                         a\leq a_T(t),\qquad b'\leq b_T(t).
                                                               \tag{AC.12}
\]
Let $e_T:T\simeq\Psi$ be any finite reindexing. We construct
\emph{simultaneously for every Boolean word} a family
$p_{c,\tau}\in\mathrm{HCParam}(R)$.

The top places are $((1,1),\iota(t))$, while thick places are
$((-1,1),\iota(s\,a_0+H.u_0))$ for $s\in H.T_2$.
Both lie in $\Psi$ and their first signs make them disjoint.
The inequality $|H.T_2|\geq x^{-\epsilon}n>0$ in
$\mathrm{HostVal}$ makes the thick set nonempty.
Choose a thick place $\tau_0\in T$.
For top places the index $t$ is unique, because $\iota$ is injective.
Define the integer base
\[
 B^{\mathrm{raw}}_\tau=
 \begin{cases}
 b_T(t)+\lfloor b'/x\rfloor,&e_T(\tau)=\tau_1(t),\\
 \lceil b\rceil,&e_T(\tau)\text{ is thick},\\
 \lceil w\rceil,&e_T(\tau)\text{ is thin}.
 \end{cases}                                                 \tag{AC.13}
\]
AC1 supplies $0\leq q<w_E$ such that
\[
             B_\tau=B^{\mathrm{raw}}_\tau+q\,1_{\tau=\tau_0}
                          \quad\hbox{satisfies }\mathrm{MU}(B).
                                                               \tag{AC.14}
\]
The correction changes no top or thin coordinate. On every thick
place it gives exactly
\[
                 \lceil b\rceil\leq B_\tau
                            <\lceil b\rceil+w_E.             \tag{AC.15}
\]
Thus the source's strict torsion bound is preserved.

Write $k=2\kappa$, which is possible because $R.k$ is even.
Choose endpoint integers
\[
 (U_0(\tau),U_L(\tau))=
 \begin{cases}
 (a_T(t)+\kappa-1,-b_T(t)-\kappa),
                                  &e_T(\tau)=\tau_1(t),\\
 (4D(k-2)+4kD,-1),                &\text{otherwise}.
 \end{cases}                                                 \tag{AC.16}
\]
At a top place, $a_T\geq4kD$ gives the first condition in (AC.8).
Also $b_T>0$ and
$\lfloor b'/x\rfloor\geq k\geq\kappa$, so $U_L<0$ and
$-B\leq-b_T-\kappa=U_L$.
Off the top set, $B\geq1$ by positivity of $b,w$ and $q\geq0$;
(AC.8) follows directly. Apply AC4 using these same $D,h,B,U_0,U_L$
for every word.

The resulting family has
\[
 \begin{split}
 v_i(\tau)&=4D(k-2-i),\\
 b_j(\tau)&=B_\tau+jh_0,\qquad
 e_j(\tau)=B_\tau+(j+\tfrac12)h_0,\\
 u_{\ell+1}(\tau)&=v_\ell(\tau)-\tfrac12
                         -D\,1_{c_\ell=\mathrm{true}}.
 \end{split}                                                 \tag{AC.17}
\]
Since $h_0=2h$, the weight $e_j$ is the integral weight $B+(2j+1)h$.
Equations (AC.14) and (AC.7) prove $\mathrm{MUQ}(e_j)$ for every $j$.
This supplies the exact root condition in $\mathrm{HCVal}$.
All its other fields are (AC.13), (AC.15), (AC.16), and (AC.17).
In particular, off the top set the endpoints are
$4D(k-2)+4kD+1/2$ and $-1/2$, and on a top place they are
\[
         u_0=a_T(t)+(k-1)/2,\qquad
         u_{k-1}=-b_T(t)-(k-1)/2.                            \tag{AC.18}
\]
The functions $v,b,e$ and both endpoint functions are independent
of the word. This proves every conclusion of
\texttt{exists\_\allowbreak hc\_\allowbreak family}, including the five compatibility
identities between any two words and their prescribed top values.
No unproved global simultaneous choice of these column families
remains under the stated inputs.

\paragraph{AC7: the literal \texttt{chGood\_\allowbreak of\_\allowbreak host} constructor.}
Fix $\epsilon,C_5,c_6,C_b,c_D,c_h,C$ as in that theorem:
$0<\epsilon<1/52$, $\mathrm{Std}(C)$, $c_D\geq0$, $c_h\geq402$.
Then a uniform threshold works for all the following supplied data:
an actual $R,H$, $\mathrm{RankVal}(s,R)$,
$H.CT_2=C.CT_2$, the literal
\[
 \mathrm{HostGood}(s,C.P.C_{\mathrm{pde}},C_5,c_6,C_b,
                                      H,q_0(H)),
                                                               \tag{AC.19}
\]
prescribed weights $a_T,b_T$ with $\mathrm{WinOK}$, a code
$c:\mathrm{Fin}(M)\hookrightarrow\{\mathrm{false},\mathrm{true}\}^{k-2}$,
and
\[
 M=\left\lceil3\Gamma^2x^{C.P.C_F}\right\rceil_+,
       \qquad \Gamma=\exp(C.C_\Gamma m^5f\mathcal L).
                                                               \tag{AC.20}
\]
An arbitrary natural $u$ may also be prescribed as the host's
unipotent translate.

Put $N_c=H.Q_S H.q H.d$ and $C_8=C.P.C_{\mathrm{pde}}\geq1$.
The pair-row part of (AC.19) gives
\[
              \log N_c\leq2x^{\epsilon/5}(\log x)^2.
\]
For every $z\leq N_c^{C_8}$, require the uniform elementary
inequality
$2C_8(\log x)^2< x^{2\epsilon/15}$.
Then $z<X'$. The host bounds therefore give
\[
                q_0(H),p_{\mathrm{de}},p'_{\mathrm{de}}<X'.
                                                               \tag{AC.21}
\]
On the other hand, each $g_i\geq(kb)^8\geq b\geq X'$.
Thus $q_0(H)\ne g_i$, with no need to assume this extra inequality.

The host's level-prime set is exactly
$\{g_0,g_1,p_{\mathrm{de}},p'_{\mathrm{de}},q_0(H)\}$ by
(AC.19). Every member is prime and congruent to $1$ modulo $N_c$:
for the two $g_i$ this is $\mathrm{HostArith}$; for the two
$p_{\mathrm{de}}$ it follows by reducing their stated congruence
modulo $\operatorname{lcm}(15,N_c)$; for $q_0(H)$ it is already
in (AC.19). Here $N_c>1$.
The same input supplies $H.cc_K=\mathcal O_K$ and the disjointness
of the two literal $\mathrm{BB}$ sets.

There is also a stronger congruence in \texttt{ChOK.small}.
The source defines $r_0(H)$ to be the least prime not dividing
$N_c(p_{\mathrm{de}}-1)(p'_{\mathrm{de}}-1)N(cc_K)$ and defines
$q_0(H)$ as the least prime $p$ with
\[
              p\equiv1\pmod {r_0(H)N_c},
                           \qquad p\nmid N(cc_K).
\]
Because $q_0(H)$ is already prime by (AC.19), this defining set is
nonempty: its natural infimum would be zero if it were empty.
The least element of a nonempty set of natural numbers belongs to
that set. Thus the stronger congruence follows directly from the
literal definition. Since $cc_K=\mathcal O_K$ has norm one,
$r_0(H)$ is exactly \texttt{leastOff}
$\bigl(N_c(p_{\mathrm{de}}-1)(p'_{\mathrm{de}}-1)\bigr)$.
This argument does not add a new prime-existence hypothesis.

Replace only $H.u_0$ by $R.\mathrm{ofNat}(u)$.
The definitions of $\mathrm{HostGood}$ and $\mathrm{WinOK}$ do not
read that field, so (AC.19) and all prescribed weights persist.
The thick set can change; construct the AC6 family after this
replacement, so its thick-place correction uses the correct new set.
The lower halves of $\mathrm{WinOK}$ give (AC.12).
Since the positive expression in (AC.20) has positive ceiling,
$M>0$. Choose the base word $0\in\mathrm{Fin}(M)$ and put
\[
                         \mathrm{hcW}(v,\tau)=p_{c(v),\tau}.
\]
Together with the supplied $R$, the translated $H$, and the preceding
level and $\mathrm{BB}$ proofs, these data fill the actual
\texttt{PreCh} record. Its five word-compatibility fields are exactly
those proved in AC6.

For precision, the sixteen fields of $\mathrm{ChOK}$ are supplied as
follows:
\[
\begin{array}{c|l}
\text{fields}&\text{proof or supplied value}\\ \hline
\mathrm{rank},\mathrm{ct2}&\mathrm{RankVal},\ H.CT_2=C.CT_2\\
\mathrm{val},\mathrm{arith},\mathrm{thick},\mathrm{dlog}
      &\text{the corresponding clauses of (AC.19)}\\
\mathrm{gnorm},\mathrm{pde}&\text{the integral norm and congruence clauses of (AC.19)}\\
\mathrm{small}&\text{level membership, (AC.21), and the stronger congruence above}\\
\mathrm{lev}&\text{the five-element level set and (AC.21)}\\
\mathrm{win},\mathrm{card},\mathrm{inj}
      &\mathrm{WinOK},\ \text{(AC.20)},\ \text{the supplied injective code}\\
\mathrm{hc},\mathrm{code},\mathrm{top}&\text{AC6, (AC.17), and (AC.18)}
\end{array}                                                   \tag{AC.22}
\]
In the \texttt{small} row, inequality with the two $p_{\mathrm{de}}$
is a clause of (AC.19), while inequality with the $g_i$ was proved
above. The retained $\mathrm{HostGood}$ is the second field of
$\mathrm{ChGood}$. This proves the literal endpoint with its
prescribed translate, weights and code. Its threshold depends only
on the fixed displayed constants, and is chosen before $R,H,a_T,b_T,M,c,u$.

\paragraph{AC8: turn an integral norm-form solution into the required element.}
For $d$ prime with $d\equiv3\pmod4$, write $s_d^2=-d$ and let
conjugation send $s_d$ to $-s_d$. The literal assertion
$\mathrm{NormOne}(d,M,p)$ means that there are $A,B\in\mathbb Z$
with
\[
 A^2+dB^2=4p,\qquad M\mid B,\qquad 2M\mid A-2-B.
                                                               \tag{AC.23}
\]
The pair-based and two-integer source definitions are identical.
Write $B=M u$, $A=2+Mu+2Mv$, and put $e=(d+1)/4\in\mathbb Z$.
The element $\omega=(1+s_d)/2$ is integral because
$\omega^2-\omega+e=0$. Define
\[
       \pi=1+M(v+u\omega)=\frac{A+B s_d}{2}\in\mathcal O_E.
                                                               \tag{AC.24}
\]
It lies in $K$, satisfies $\pi\overline\pi=p$, and
$\pi-1\in M\mathcal O_E$. Moreover expansion gives
\[
     p-1=M\left(u+2v+M(eu^2+uv+v^2)\right).
                                                               \tag{AC.25}
\]
Thus $p\equiv1\pmod M$. This proves every conclusion of
\texttt{normOne\_\allowbreak host}, with the actual element (AC.24), not just
an abstract splitting assertion.

The converse used in \texttt{normOne\_\allowbreak iff\_\allowbreak of\_\allowbreak sqrt} can also be
proved directly. Suppose $M>0$, $\pi\in\mathcal O_E\cap K$,
$\pi\overline\pi=p$, and $\pi-1=M y$ with $y\in\mathcal O_E$.
Then $y\in K$, so write $y=a+c s_d$ with $a,c\in\mathbb Q$.
Both its trace $2a$ and norm $a^2+dc^2$ are rational algebraic
integers and hence integers, say $t$ and $v_0$.
The identity $d(2c)^2=4v_0-t^2\in\mathbb Z$ implies
$2c\in\mathbb Z$: in lowest terms its denominator squared divides
the squarefree integer $d$. Write $2c=u$.
Reducing $t^2+du^2=4v_0$ modulo four and using $d\equiv3\pmod4$
shows that $t,u$ have the same parity. Set
$A=2+Mt$ and $B=Mu$. Then (AC.23) follows by taking the norm of
$\pi=1+My$, and $2M\mid A-2-B$ follows from the parity.
For $M=0$, both sides of the equivalence force $\pi=1$ and $p=1$;
use $A=2,B=0$. Finally, replacing $M$ by any divisor $M'$ preserves
(AC.23), so the divisor version of the congruence follows as well.

\paragraph{AC9: construct the full level record from its prime solutions.}
Supply an actual $P:\mathrm{PreCh}$, positive integers $\mathrm{du}_v$,
positive $D_*$, and $\Gamma>1$. Write
\[
 U_v=\left(\frac{D_*}{\mathrm{du}_v}\right)^{1/(fk)}
                         \Gamma^{-2/(fk)},\qquad
 M_c=N_c(g_0g_1)^2.                                         \tag{AC.26}
\]
Suppose for every word $v$ we have a prime $T_v$ and an integral
$\pi_v\in K$ with
\[
 \begin{gathered}
 \pi_v\overline{\pi_v}=T_v,\quad
 \pi_v-1\in M_c\mathcal O_E,\quad
 T_v\ne p_{\mathrm{de}},p'_{\mathrm{de}},\\
       U_v\leq T_v\leq(1+1/(fk))U_v.
 \end{gathered}                                              \tag{AC.27}
\]
These data construct the actual $\mathrm{LevCh}(P)$.
The sole additional arithmetic field to check is
$T_v\equiv1\pmod {N_c}$.
Since $N_c\mid M_c$, write $\pi_v=1+N_c\alpha$ with
$\alpha\in\mathcal O_E$. Then
\[
 \frac{T_v-1}{N_c}
       =\alpha+\overline\alpha+N_c\alpha\overline\alpha
\]
is both rational and an algebraic integer, so it is an integer.
Since $N_c>1$, the literal remainder is $1$.

Set the five data fields $T,\pi T,\mathrm{duStar},Dstar,GamP$
to $T_v,\pi_v,\mathrm{du}_v,D_*,\Gamma$, respectively.
The ten proof fields are: primality, the congruence just proved,
$\pi T\in K$, its norm identity, its congruence modulo $M_c$,
avoidance of the two specified primes, positivity of $\mathrm{duStar}$,
positivity of $Dstar$, $1<GamP$, and the two-sided window.
They are exactly (AC.27) and the stipulated positive values.
No pairwise distinctness of the $T_v$ is required by this record.
This proves \texttt{LevelChoice.exists\_\allowbreak levCh}.

The associated host has level set
$H.\mathrm{levelPrimes}\cup\{T_v\}$. Its old membership fields persist.
Each new prime is coprime to $Q_S$, since it is $1$ modulo the
multiple $N_c$ of $Q_S$. The $cc_K$ coprimality fields follow from
$cc_K=\mathcal O_K$. Thus this is also the literal
\texttt{LevCh.host} used by the later joins, with the same field,
rank, frame labels, and column data.

\paragraph{AC10: use exactly the retained analytic level-prime assertion.}
Assume the literal $\mathrm{LevelPrime}$ statement. For every fixed
$0<\epsilon<1/52$, $c_s>0$, and standard $C$, there is a threshold
$x_0$ with the following property, uniformly over \emph{all} $c_D,c_h$
and all subsequently supplied $s,P,q_1,a_T,b_T,c$.
Suppose $s.\mathrm{eps}=\epsilon$, $x\geq x_0$, and
$\mathrm{ChOK}(s,C,c_D,c_h,P,q_1,a_T,b_T,c)$.
Supply positive $\mathrm{du}_v,D_*$ with
\[
       c_s m\mathcal L\leq\log U_v\quad\text{for every }v,
       \qquad
       \Gamma=\exp(C.C_\Gamma m^5 f\mathcal L),               \tag{AC.28}
\]
where $U_v$ is (AC.26) for this $\Gamma$.
We construct $L:\mathrm{LevCh}(P)$ with the exact equalities
\[
              L.\mathrm{duStar}=\mathrm{du},\quad
              L.Dstar=D_*,\quad L.GamP=\Gamma.              \tag{AC.29}
\]

Here is the precise analytic input and how its hypotheses are met.
$\mathrm{LevelPrime}$ supplies, once its threshold depending on
$\epsilon,c_s$ is met, a prime $T$ in
$[U,(1+1/(fk))U]$ and
$\mathrm{NormOne}(d,N_c(g_0g_1)^2,T)$, provided
\[
 \begin{gathered}
 \mathrm{RankVal}(s,R),\quad d\text{ prime},\quad d\equiv3\pmod4,
       \quad d\geq100,\quad 0<N_c,\quad d\mid N_c,\\
 \log N_c\leq3x^{\epsilon/5}(\log x)^2,\qquad
 1\leq g_i\leq b^9,\quad U>0,\quad c_s m\mathcal L\leq\log U.
 \end{gathered}                                              \tag{AC.30}
\]
The conditions on $d$, $N_c$, and $\log N_c$ are the literal
$\mathrm{ModRow}(s,N_c,d)$; $\mathrm{RankVal}$ is a separate input.
The host supplies primality and $d\equiv7\pmod {16}$.
Its bound $4tr+1<d$, together with $r\geq6,t\geq11$, gives
$d>265>100$. The other two elementary conductor conditions are
immediate from $N_c=Q_Sqd$. The pair row inside $\mathrm{ChOK.val}$
has $1\leq Y\leq x$ and
$\log N_c\leq2Y^{\epsilon/5}(\log Y)^2$; monotonicity gives
the required stronger bound with coefficient two, hence (AC.30).

Require the RF11--RF12 thresholds, $x\geq2$, and
$x\geq(1/c_s)^4$. Then
\[
 k\leq x^5,\quad x^{41}\leq b,\quad m\geq x^{1/4},
        \qquad 2k^8\leq2x^{40}\leq x^{41}\leq b.
\]
The host's upper bound $g_i\leq2(kb)^8$ now gives $g_i\leq b^9$.
Its primality gives $g_i\geq1$.
The same thresholds imply $c_s m\geq1$, and therefore
\[
              X'\leq\exp(c_s m\mathcal L)\leq U_v.
                                                               \tag{AC.31}
\]
All $U_v$ are positive; $\Gamma>1$ because $C.C_\Gamma\geq1$
and $m,f,\mathcal L$ are positive.
For $p_{\mathrm{de}}$ and $p'_{\mathrm{de}}$, the field
$\mathrm{ChOK.lev}$ gives either equality with some $g_i$ or a
strict bound by $X'$. The equality is excluded by the literal host
axiom $p_{\mathrm{de}},p'_{\mathrm{de}}\ne g_i$.
Thus both specified primes are less than $X'$.

Apply $\mathrm{LevelPrime}$ for each word with $U=U_v$.
It supplies all the prime windows and integer norm solutions at
one common threshold: its quantifier over $U$ occurs after its
threshold, so neither the word nor $\mathrm{du},D_*$ changes $x_0$.
Use AC8 to produce $\pi_v$ with congruence modulo $M_c$.
Equation (AC.31) gives $T_v\geq U_v\geq X'$, excluding both
$p_{\mathrm{de}}$ primes. AC9 constructs the desired record and
(AC.29). Taking the maximum of these finitely many thresholds
proves exactly \texttt{exists\_\allowbreak levCh\_\allowbreak of\_\allowbreak levelPrime}.

Consequently RF10's two arithmetic constructor arrows are now proved.
They remain conditional on the same actual inputs: the supplied
$\mathrm{HostGood}$, top windows and injective code for AC7; and
$\mathrm{LevelPrime}$ plus the positive discriminant data and their
logarithmic window for AC10. The construction of recursive global
forms, the joint $\mathrm{RecExP}$ data, and the independent analytic,
spectral and pairing branches are not consequences of these arrows.


\section{Constructing the character field and its arithmetic host}
\label{sec:arithmetic-hosts}

\status{AH1--AH18 reconstruct the finite character construction, its actual
cyclotomic fixed field, every field of \texttt{HostDer}, the half-different
ideal and its trace-self-dual inverse, and the literal \texttt{fieldChi}
conclusion. The hypotheses on the primes $q,d$, their splitting congruences,
and their eventual analytic windows remain hypotheses. These arguments
construct arithmetic data once those hypotheses hold; they do not select
the analytic windows.}

Throughout this section $U_a=(\mathbb Z/a\mathbb Z)^\times$ for $a>0$,
$T_e=2^{e+2}$, and $H_e=U_{T_e}/\langle-1\rangle$.
The symbol $\varphi(a)$ denotes Euler's totient, whereas $\phi$ denotes a
frame. Write $n=\phi.\mathrm{ord}=|A|$, $n_0=|A_0|$ and
$o=n/n_0$. We use multiplicative notation for the finite group $A$;
transporting it to an additive group on $\{0,\ldots,n-1\}$ gives the
literal frame. Put
\[
 G=\{\pm1\}\times\{\pm1\}\times A,\qquad
 c=(-1,-1,1),\quad \iota_d=(1,-1,1),\quad
 \Psi=\{(s,t,a):t=1\}.                         \tag{AH.1}
\]
The exact prime hypothesis is
\[
 q,d\text{ prime},\qquad q\equiv3\pmod8,
 \quad d\equiv7\pmod{16},\qquad
 \left(\frac{-d}{q}\right)=1.                 \tag{AH.2}
\]
It implies that $q,d$ are distinct odd primes. In particular no choice of
such primes is included in the arguments below.

\paragraph{AH1: finite character duality and the two-adic quotient.}
The quotient $H_e$ is cyclic of order $2^e$. Indeed,
$v_2(5^{2^j}-1)=j+2$: the case $j=0$ is immediate, and the induction step
uses $5^{2^j}+1\equiv2\pmod8$ for $j\ge1$, with the first step checked
directly. Thus $5$ has order $2^e$ modulo $T_e$. Its powers are $1$ modulo
$4$, so $-1$ is outside this subgroup, and counting the $2^{e+1}$ odd
residue classes proves $U_{T_e}=\langle-1\rangle\times\langle5\rangle$.

Let $M>0$ and let $\Xi$ be a subgroup of the complex characters of $U_M$.
Set
\[
 B=U_M/\operatorname{ann}(\Xi),\qquad
 \operatorname{ann}(\Xi)=\{u:\chi(u)=1\text{ for every }\chi\in\Xi\}.
                                                               \tag{AH.3}
\]
Then $|B|=|\Xi|$, and the characters of $B$ are exactly the characters in
$\Xi$. To see both assertions, sum $\chi(u)$ over $u\in U_M$ and
$\chi\in\Xi$. The sum of a nontrivial character is zero, since translating
the sum by a point where the character is not $1$ multiplies it by that
value. Summing first over $\Xi$ gives
$|\operatorname{ann}(\Xi)||\Xi|=|U_M|$. Every character in $\Xi$ factors
through $B$, and a finite abelian group has exactly its order many complex
characters: on each cyclic factor of order $a$ one chooses any $a$th root
of unity. The resulting cardinality equality proves the assertion.
This argument includes groups with odd-order factors.

We also use the following elementary reduction fact. If $a\mid b$,
the map $U_b\to U_a$ is onto. Lift a unit modulo $a$ to an integer and
use CRT to make it $1$ at every prime dividing $b$ but not $a$.
The resulting integer is coprime to $b$ and has the required image.

\paragraph{AH2: the minimal modulus and the actual mixed frame.}
Fix $e\ge0$, put $W_e=\operatorname{lcm}(T_e,M)$, and let
$\xi_W:U_{W_e}\to B$ be reduction modulo $M$ followed by the quotient.
Choose the least positive integer $Q=Q_e$ satisfying
\[
 Q\mid W_e,\quad T_e\mid Q,\quad
 u\equiv1\pmod Q\Longrightarrow\xi_W(u)=1.       \tag{AH.4}
\]
The set is nonempty because $W_e$ belongs to it. Reduction $U_{W_e}\to U_Q$
is surjective; its kernel lies in $\ker\xi_W$, so there is a unique
surjective homomorphism $\xi:U_Q\to B$ with
$\xi(u\bmod Q)=\xi_W(u)$. Define
\[
 \operatorname{cl}:U_Q\longrightarrow H_e,
 \qquad \eta(u)=(\operatorname{cl}(u),\xi(u)),
 \qquad A=\eta(U_Q)\subset H_e\times B.          \tag{AH.5}
\]
The first projection $p_1:A\to H_e$ is onto. On $\ker p_2$ it is
injective, since an element with both coordinates $1$ is the identity.
Hence $\ker p_2$ is cyclic. Choose a generator $a_0$ and set
$A_0=\langle a_0\rangle=\ker p_2$. If $p_1(a_0^j)=1$, then $a_0^j=1$;
this is precisely the generator condition of \texttt{Frame}.
Relabeling the finite group by $\operatorname{Fin}(|A|)$ gives the literal
\texttt{ModData.frame}, with its supplied group isomorphism, not merely a
group of the same order. Let $a:U_Q\to A$ denote $\eta$ with this relabeling,
and let $\xi_A:A\to B$ denote $p_2$. Then
\[
 p_1a=\operatorname{cl},\quad \xi_Aa=\xi,
 \quad A_0=\ker\xi_A,
 \quad o=[A:A_0]=|B|=|\Xi|,
 \quad 2^e\mid n.                              \tag{AH.6}
\]
If every element of $B$ has order dividing $2^j$, then $A$ is a $2$-group,
because it is a subgroup of $H_e\times B$. No such assumption is made
for the construction itself.

\paragraph{AH3: CRT constructs the Artin map, including its sign twist.}
More generally take any $Q>0$, $T_e\mid Q$, and any homomorphism
$\xi:U_Q\to B$, and form $A$ as in AH2. Assume (AH.2),
$\gcd(Q,qd)=1$, $e\ge2$, and
$(-d/\ell)=1$ for every odd prime $\ell\mid Q$.
Set $N=Qqd$ and
\[
 \alpha(u)=\left(\frac{u}{q}\right),\qquad
 \beta(u)=\left(\frac{u}{d}\right),\qquad
 w(s)=a(s\bmod Q),\quad s\in\{\pm1\}.
\]
The literal homomorphism is
\[
 \operatorname{art}(u)=
 \bigl(\alpha(u),\beta(u),a(u\bmod Q)w(\alpha(u))\bigr).
                                                               \tag{AH.7}
\]
For a prescribed $(s,t,b)\in G$, choose $v\in U_Q$ with
$a(v)=b\,w(s)$. Choose a residue of Legendre sign $s$ modulo $q$ and
one of sign $t$ modulo $d$; one may use $1$ or $-1$ because both primes
are $3$ modulo $4$. CRT gives a unit $u\in U_N$ with these three
residues. Since $w(s)^2=1$, (AH.7) maps it to $(s,t,b)$.
Thus the Artin map is surjective.

The remaining identities follow from the same formula:
\[
 \begin{aligned}
 \operatorname{art}(-1)&=c,\\
 p_1(\operatorname{art}(u)_A)&=[u\bmod T_e],\\
 u\equiv1\pmod{Qq}&\Longrightarrow
              \operatorname{art}(u)\in\{1,\iota_d\},\\
 u\equiv1\pmod{Qd}&\Longrightarrow
              \operatorname{art}(u)^2=1.
 \end{aligned}                                                   \tag{AH.8}
\]
For the first identity the $A$ coordinate is $a(-1)w(-1)=1$;
for the second, $p_1w(s)=1$ because $s=\pm1$ dies in $H_e$.
These calculations prove every Artin field of
\texttt{ModData.charData}. The factor $w(\alpha(u))$ is indispensable
when the original characters are not even.

\paragraph{AH4: the cyclotomic fixed field and the actual quadratic roots.}
Take any \texttt{CharData}, with its surjection
$\operatorname{art}:U_N\to G$ and the identities (AH.8), and put
\[
 L=\mathbb Q(\zeta_N),\qquad
 U=\ker(\operatorname{art}),\qquad E=L^U.        \tag{AH.9}
\]
Under the usual cyclotomic identification
$\tau_u(\zeta_N)=\zeta_N^u$, the fundamental theorem of Galois theory gives
\[
 \operatorname{Gal}(E/\mathbb Q)\simeq U_N/U\simeq G,
 \qquad [E:\mathbb Q]=4n.                       \tag{AH.10}
\]
The isomorphism is characterized by
$\operatorname{gal}(\tau_u|_E)=\operatorname{art}(u)$, so the inclusion
$E\hookrightarrow L$ satisfies the literal \texttt{art\_\allowbreak spec}.

For completeness, the quadratic elements can be constructed inside this
same cyclotomic field. For a prime $p\equiv3\pmod4$ dividing $N$, choose
$\zeta_p=\zeta_N^{N/p}$ and set
$g_p=\sum_{a\bmod p}(a/p)\zeta_p^a$. Changing $a$ to $u^{-1}a$ gives
$\tau_u g_p=(u/p)g_p$. Also $\overline{g_p}=-g_p$ and
\[
 g_p\overline{g_p}
 =\sum_{t\ne0}(t/p)\sum_{b\ne0}\zeta_p^{(t-1)b}=p.
\]
Here the inner sum is $p-1$ when $t=1$ and $-1$ otherwise, and
$\sum_{t\ne0}(t/p)=0$. Hence $g_p^2=-p$.
If $(u/p)$ is a character of $\operatorname{art}(u)$, its value is $1$
on $U$, so $g_p\in E$ with the prescribed Galois action. In particular
$s_d=g_d$ realizes the required second sign. For the Artin map of AH3,
$s_q=g_q$ realizes the first sign as well. This proves
\texttt{CharData.toHostRaw}, \texttt{ModData.host} and
\texttt{AnyGrp.host}, including their specified roots, embeddings and
Galois maps.

\paragraph{AH5: exactness of the modulus and the primitive conductor.}
Return to the minimal $Q_e$ of AH2 and assume $e\ge1$.
For every prime $\ell\mid Q$ there is a unit $v\in U_Q$ with
\[
 v\equiv1\pmod{Q/\ell},\qquad \eta(v)\ne1.      \tag{AH.11}
\]
If $T_e\mid Q/\ell$ and no such $v$ existed, every unit modulo $W_e$
congruent to $1$ modulo $Q/\ell$ would have trivial $\xi_W$-value.
Thus $Q/\ell$ would satisfy (AH.4), contradicting minimality.
Otherwise $\ell=2$ and $v_2(Q)=e+2$. The unit $v=1+Q/2$ is $1$ modulo
$Q/2$ but is neither $1$ nor $-1$ modulo $T_e$; the latter assertion uses
$e\ge1$ and $Q/2\equiv2^{e+1}\pmod{T_e}$. Its $H_e$ coordinate is
therefore nontrivial. Lift $v$ by CRT to a unit $u$ which is $1$ modulo
$q$ and $d$. Then $u\equiv1\pmod{N/\ell}$ and
$\operatorname{art}(u)\ne1$. This is every clause of
\texttt{HostRaw.Exact}, for the actual modulus $Q$.

When $\Xi=\langle\psi\rangle$ for a primitive character of conductor $M$,
\[
                    M\mid Q,\qquad Q=W_e.       \tag{AH.12}
\]
Indeed, let $g=\gcd(M,Q)$ and let $v\in U_M$ be $1$ modulo $g$.
Generalized CRT gives $z\equiv v\pmod M$ and $z\equiv1\pmod Q$.
It is a unit modulo $W_e$, because $T_e\mid Q$; (AH.4) implies
$\psi(v)=1$. Thus $\psi$ factors through $U_g$. Primitivity forces
$g=M$, proving $M\mid Q$; together with $T_e\mid Q\mid W_e$ this gives
the displayed equality. The literal source only needs its divisibility
consequence.

\paragraph{AH6: one eventually constant multiplier in the frame order.}
For the fixed $M,\Xi$, define
\[
 J_e=\xi(\ker\operatorname{cl})\subset B.
\]
Projection of $A\subset H_e\times B$ onto $H_e$ has kernel
$\{1\}\times J_e$. Consequently
\[
              n_e=2^e|J_e|,\qquad 1\le|J_e|\le|\Xi|.             \tag{AH.13}
\]
These subgroups form a descending sequence. To check this without assuming
$Q_e\mid Q_{e+1}$, lift a representative at level $Q_{e+1}$ to a unit
modulo $W_{e+1}$, reduce it modulo $W_e$, then modulo $Q_e$.
Its two-adic residue, originally $\pm1$ modulo $T_{e+1}$, is $\pm1$
modulo $T_e$, and its $B$ coordinate is unchanged. Hence
$J_{e+1}\subset J_e$. The nonincreasing positive integers $|J_e|$ attain
a minimum. Choose an index $e_1$ attaining it and write that minimum as
$c_0$. For all $e\ge e_1$,
\[
 1\le c_0\le|\Xi|,\qquad n_e=c_0\,2^e.          \tag{AH.14}
\]
The constants precede $q,d$ and every analytic scale. This is the exact
quantifier order of \texttt{ordAny}.

\paragraph{AH7: the prescribed characters are actually kept.}
For $\chi\in\Xi$, let $\chi_B:B\to\mathbb C^\times$ be its descent and
write $\theta_A=\chi_B\xi_A$. Define the character of the actual group $G$
by
\[
 \theta_\chi(s,t,b)=\theta_A\bigl(b\,w(s)\bigr). \tag{AH.15}
\]
It is trivial on $(1,1,A_0)$ and on $\iota_d$, so belongs to the literal
$X^+$ character set. Substitution in (AH.7) cancels $w(\alpha)^2$ and gives
\[
       \theta_\chi(\operatorname{art}(u))=\chi(u)                 \tag{AH.16}
\]
whenever both sides are read on units of a common multiple of $N$ and $M$.
Thus the primitive Dirichlet character induced by $\theta_\chi$ is the
primitive character induced by $\chi$. This uses uniqueness of the
primitive inducing character, which follows by factoring both characters
through the common modulus and then through the gcd of their conductors.
Accordingly every primitive character in \texttt{Level.ofSub} is kept in
$X^+$.

The first-sign character is not among them. Use CRT to choose
$u\equiv1\pmod Q$, $u\equiv1\pmod d$ and $(u/q)=-1$.
For every character (AH.15), its value on $\operatorname{art}(u)$ is $1$,
whereas the first sign is $-1$. Equal primitive inducing characters would
have equal values on all units modulo $N$, a contradiction. Since
\texttt{Level.ofSub.psK=none}, this is the entire
\texttt{ExceptionalCharacterKept} predicate. It does not assert that all
characters in $X^+$ are kept.

\paragraph{AH8: the exact kernel and every allowed quadratic twist.}
Let $\psi$ be primitive of conductor $M$, and $\Xi=\langle\psi\rangle$.
For every integer $a$ coprime to $N$, (AH.7) and (AH.12) give
\[
 \begin{split}
 \operatorname{art}(a)=1\quad\Longleftrightarrow\quad&
 a\equiv\pm1\pmod{T_e},\quad\psi(a)=1,\\
 & (a/q)=(a/d)=1. 
 \end{split}                                                     \tag{AH.17}
\]
Indeed, the first two coordinates must be $1$; then $w(\alpha)=1$ and
the remaining coordinate is trivial exactly when both components of
$\eta(a)$ are trivial. The annihilator of $\langle\psi\rangle$ is
$\ker\psi$. This proves \texttt{KerChi}, including $M\mid Q$.

Now let $\theta:G\to\mathbb C^\times$ be trivial on $(1,1,A_0)$.
Its restriction to $A$ factors through $A/A_0\simeq B$, hence by AH1 is
some $\chi\in\Xi$. The restrictions to
$s\mapsto(s,1,w(s))$ and $t\mapsto(1,t,1)$ have values $s^i,t^j$
for some $i,j\in\{0,1\}$. This follows because the image of the generator
of a group of order $2$ has square $1$, hence is $1$ or $-1$.
The factorization
\[
 \operatorname{art}(a)
  =(\alpha,1,w(\alpha))(1,\beta,1)(1,1,a(a\bmod Q))
\]
therefore gives
\[
 \theta(\operatorname{art}(a))
   =\chi(a)(a/q)^i(a/d)^j,\qquad
 \theta(\iota_d)=(-1)^j.                        \tag{AH.18}
\]
This is \texttt{TwistForm} with its prescribed $\iota_d$ sign. It includes
all characters trivial on the designated cyclic subgroup, with no
restriction to even characters or to $2$-groups.

\paragraph{AH9: cyclotomic reduction gives the complete local Galois laws.}
Let $L=\mathbb Q(\zeta_N)$ and let $P_L$ lie above a rational prime $p$.
Write $N=p^v m$, with $p\nmid m$. In the residue field, the reduction of
$\zeta_N$ has exact order $m$. Its order divides $m$ because a root of
$p$-power order reduces to $1$. A root of order prime to $p$ keeps its
order: if $\eta\ne1$, $\eta^a=1$, $p\nmid a$, and $\eta\equiv1\pmod{P_L}$,
then $1+\eta+\cdots+\eta^{a-1}=0$ would reduce to $a=0$, a contradiction.
Applying this to powers of $\zeta_N^{p^v}$ gives exactness. Thus
\[
 \zeta_N^r\equiv\zeta_N^s\pmod{P_L}
       \quad\Longleftrightarrow\quad r\equiv s\pmod m.          \tag{AH.19}
\]

Here one may test integral residue maps on $\zeta_N$, since
$\mathcal O_L=\mathbb Z[\zeta_N]$. The relevant integral statement has
a direct local proof. For each $p$, the prime-to-$p$ polynomial
$X^m-1$ has nonzero derivative at all its roots modulo $p$.
Hensel factorization therefore makes $\mathbb Z_p[\zeta_m]$ a product
of unramified discrete valuation rings. On each factor,
$\Phi_{p^v}(1+X)$ is Eisenstein if $v>0$: its constant coefficient is
$p$, its leading coefficient is $1$, and the others are divisible by $p$.
Adjoining $\zeta_{p^v}-1$ consequently gives a discrete valuation ring,
with the power basis of this Eisenstein polynomial. Since
$\mathbb Z_p[\zeta_N]=\mathbb Z_p[\zeta_m,\zeta_{p^v}]$, this is the
full integral closure locally. The finite index between
$\mathbb Z[\zeta_N]$ and $\mathcal O_L$ is therefore divisible by no prime,
so is $1$. For $v=0$ the unramified argument already suffices.

It follows from (AH.19), on the integral generator, that
\[
 \begin{aligned}
 \tau_u\in I(P_L)&\Longleftrightarrow u\equiv1\pmod m,\\
 \tau_u\text{ induces arithmetic Frobenius at }P_L
      &\Longleftrightarrow u\equiv p\pmod m.
 \end{aligned}                                                   \tag{AH.20}
\]
For the second assertion the two ring maps on the residue field agree on
the generator precisely when $\overline{\zeta_N}^u=
\overline{\zeta_N}^{p}$. A unit $u$ in this congruence class exists by
CRT, choosing its $p$-part to be $1$. The decomposition group is the
preimage of the cyclic residue Galois group generated by this Frobenius,
so
\[
 \tau_u\in D(P_L)
       \quad\Longleftrightarrow\quad
       u\equiv p^j\pmod m\text{ for some }j\ge0.                \tag{AH.21}
\]

These formulas descend to every Galois subfield $E$ of $L$ under its
actual restriction map. More explicitly, restriction maps decomposition
and inertia groups onto the respective groups at $P=P_L\cap\mathcal O_E$.
For decomposition, lift an automorphism of $E$ to $L$ and adjust the
lift by an element of $\operatorname{Gal}(L/E)$ so that it fixes $P_L$;
this is possible because that group is transitive on the primes above
$P$. For inertia, the resulting residue action fixes $\kappa(P)$.
The decomposition group over $E$ maps onto
$\operatorname{Gal}(\kappa(P_L)/\kappa(P))$; cancel this residue action
by an element over $E$. This proves the claimed surjectivity, including
the inertia direction that restriction alone would not prove.
The residue surjectivity used here is the usual finite Galois prime
theorem: the stabilizer acts on the residue field with kernel inertia
and quotient the finite-field Galois group. It follows by localizing
the normal extension at the prime and lifting the embeddings of its
maximal unramified subextension.

Consequently, for the cyclotomic realization of any raw host, writing
$N_p=N/p^{v_p(N)}$, one has the literal identities
\[
 \begin{aligned}
 \sigma\in I(P)&\Longleftrightarrow
     \exists u\in U_N: u\equiv1\pmod{N_p},\quad
                       \operatorname{gal}(\sigma)=\operatorname{art}(u),\\
 \sigma\text{ is arithmetic Frobenius at }P&\Longleftrightarrow
     \exists u\in U_N: u\equiv p\pmod{N_p},\quad
                       \operatorname{gal}(\sigma)=\operatorname{art}(u),\\
 \sigma\in D(P)&\Longleftrightarrow
     \exists u\in U_N,\ j\ge0: u\equiv p^j\pmod{N_p},\quad
                       \operatorname{gal}(\sigma)=\operatorname{art}(u).
 \end{aligned}                                                   \tag{AH.22}
\]
For the Frobenius converse, two lifts induce the same residue Frobenius
exactly when their quotient is in inertia. This supplies the necessary
lift without imposing that a general lift itself be Frobenius. If $p\nmid N$,
the inertia group is trivial and the unique arithmetic Frobenius has
image $\operatorname{art}(p\bmod N)$.

\paragraph{AH10: the CM field, the exact CM-type and the sign characters.}
Now let $H$ be an arbitrary \texttt{HostRaw}, not necessarily one just
constructed in AH4. Its specified embedding in $\mathbb Q(\zeta_N)$
and \texttt{art\_\allowbreak spec} imply that $\operatorname{art}$ is onto: extend
an automorphism of $E$ to the cyclotomic field and read its exponent on
$\zeta_N$. They also identify the image of $E$ with the fixed field of
$\ker\operatorname{art}$. Thus AH9 applies to this actual raw host.

Complex conjugation on $\zeta_N$ is $\tau_{-1}$, so its restriction to $E$
is $\operatorname{gal}^{-1}(c)$; denote it by $\overline{\phantom{x}}$.
Every embedding of $E$ into $\mathbb C$ extends to the cyclotomic field,
and hence
$\psi(\overline{x})=\overline{\psi(x)}$ for every embedding $\psi$.
Because $s_d^2=-d$ and $d>0$, no embedding of $E$ is real. The fixed field
$F=E^{\langle\overline{\phantom{x}}\rangle}$ is totally real, is the
maximal real subfield, and $E/F$ is quadratic. Therefore
\[
 [E:\mathbb Q]=4n,\qquad [F:\mathbb Q]=2n,\qquad
 x\in F\Longleftrightarrow\overline{x}=x.        \tag{AH.23}
\]
An automorphism induces conjugation at even one complex embedding exactly
when it is this automorphism, by injectivity of that embedding. This
proves the full \texttt{conj\_\allowbreak iff} predicate.

The actual CM-type consists of the embeddings with
$\psi(s_d)=i\sqrt d$. Each embedding sends $s_d$ to one of
$\pm i\sqrt d$, and conjugation exchanges the two possibilities.
Thus the type has $2n$ elements, selects exactly one from each conjugate
pair, and, for $\psi$ in the type,
\[
 \psi\sigma\text{ is in the type}
       \quad\Longleftrightarrow\quad
       \operatorname{gal}(\sigma)\in\Psi.        \tag{AH.24}
\]
Also $K=\mathbb Q(s_d)$ is exactly the fixed field of $\Psi$.

There is a nonzero first-sign eigenvector even when a raw host does not
assert \texttt{PsQ}. Take the quadratic fixed field of the kernel of the
first-sign character, choose $x$ moved by conjugation there, and use
$a=x-\overline{x}\ne0$. Automorphisms of first sign $1$ fix it and those
of first sign $-1$ negate it. This proves \texttt{ps\_\allowbreak root}; it does not
replace the stronger square-root-of-$-q$ supplied in AH4 for our constructed
host.

Finally compare the image of the raw element $s_d$ with the Gauss sum
$g_d$ of AH4. Their squares are both $-d$, so they differ by sign. The
Gauss action therefore forces
\[
 \operatorname{art}(u)_d=(u/d).                  \tag{AH.25}
\]
For $p\nmid N$, AH9, (AH.25), quadratic reciprocity and $d\equiv3\pmod4$
give every coordinate required by \texttt{frob}: the first sign is the
defined $\mathrm{psAt}(p)$, the second is $(p/d)=(-d/p)$, and the $H_e$
coordinate is the class of $p$ modulo $T_e$.

\paragraph{AH11: ramified decomposition and the unramified CM extension.}
The quadratic integral calculation in AC8 gives
\[
 \mathcal O_K=\mathbb Z[\omega],\quad
 \omega=(1+s_d)/2,\quad
 f_\omega(X)=X^2-X+(d+1)/4,quad \operatorname{disc}(K)=-d.
                                                               \tag{AH.26}
\]
At $2$ this polynomial is $X(X-1)$ modulo $2$, since $d\equiv7\pmod{16}$.
At an odd prime different from $d$ its roots are distinct, and it splits
exactly when $(-d/p)=1$; at $d$ the prime ramifies. Thus every prime
dividing $Qq$ splits in $K$. Restriction of a decomposition automorphism
at such a prime fixes its prime in $K$. In a split quadratic extension
that stabilizer is trivial, so
\[
              p\mid Qq\quad\Longrightarrow\quad D(P)\subset\Psi.
                                                               \tag{AH.27}
\]
This proves \texttt{decomp\_\allowbreak typeI}, including the prime $2$.

At $d$, (AH.22) and $N_d=Qq$ give
\[
                     I(P)=\{1,\iota_d\}.        \tag{AH.28}
\]
The containment follows from \texttt{art\_\allowbreak d}. For the reverse inclusion,
CRT gives $u\equiv1\pmod{Qq}$ and $u\equiv-1\pmod d$.
Its second sign is $-1$ by (AH.25), so its image is $\iota_d$.

We need a further first-sign fact at $d$. The integer $-d$ is a square
modulo $Qq$. At every odd prime factor, start with the prescribed
nonzero square root and lift by solving the linear congruence for its
next digit; $2s$ is invertible. At $2$, start with $s=1$ modulo $8$.
If $s^2\equiv-d\pmod{2^a}$, $a\ge3$, replacing $s$ by
$s+2^{a-1}$ changes its square by $2^a$ modulo $2^{a+1}$, and therefore
corrects the next bit when necessary. CRT combines the prime powers.
The first sign of $\operatorname{art}$ factors through $U_{Qq}$ by
\texttt{art\_\allowbreak d}. If $s^2\equiv-d\pmod{Qq}$, then $d\equiv-s^2$, whence
\[
 u\equiv d\pmod{Qq}\quad\Longrightarrow\quad
              \operatorname{art}(u)_q=-1.       \tag{AH.29}
\]
This proof uses the first sign of $\operatorname{art}(-1)$ and the square
of a sign; it does not assume \texttt{PsQ} for a general raw host.

Conjugation fixes no prime above a divisor of $N$. For a divisor of $Qq$
this follows from (AH.27), since $c\notin\Psi$. At $d$, choose a
Frobenius lift $F$ with exponent $u\equiv d\pmod{Qq}$. If conjugation
were in the decomposition group, then
$c=\rho\,\operatorname{gal}(F)^j$ with
$\rho\in\{1,\iota_d\}$ and $j\ge0$. The $A$ coordinate and
\texttt{art\_\allowbreak h} would give $d^j\equiv\pm1\pmod{T_e}$.
Reducing modulo $16$, the congruence $d\equiv7$ forces $j$ even, because
$7^j$ is $7$ for odd $j$ and $1$ for even $j$.
But the first sign, using (AH.29), gives $-1=(-1)^j$, a contradiction.
Thus \texttt{decompConj} holds. At primes outside $N$ inertia is trivial;
at primes dividing $N$ conjugation is not even in the decomposition group.
It follows that conjugation is in no finite inertia group, proving the
literal \texttt{unramEF}: $E/F$ is unramified at every finite prime.

\paragraph{AH12: splitting, ordinary places and the exact $\eta$ clauses.}
The factorization in (AH.26) gives
\[
 \#\{\mathfrak p\subset\mathcal O_K:\mathfrak p\mid p\}=2
 \quad\Longleftrightarrow\quad
 p=2\ \text{or}\ \bigl(p\ne2\text{ and }(-d/p)=1\bigr).        \tag{AH.30}
\]
For a prime $P$ of $E$, the source's ordinary condition says
$\sigma\tau\in\Psi\Longleftrightarrow\sigma\in\Psi$ for every
$\tau\in D(P)$ and every $\sigma$. Taking $\sigma=1$ shows this is
equivalent to $D(P)\subset\Psi$; the converse follows because $\Psi$ is
a subgroup. Since restriction maps $D(P)$ onto the decomposition group
in $K$, this is equivalent to splitting in $K$. This proves the full
\texttt{ordinary} equivalence for all rational primes, including $d$.

Since $E/F$ is quadratic and unramified at finite primes, every nonzero
prime of $F$ either has two primes above it or has one prime of residue
degree $2$. The latter case occurs exactly when conjugation fixes a prime
$P$ above it. For $p\nmid N$, the decomposition group is generated by
the arithmetic Frobenius $\sigma$, so this is equivalent to
$c\in\langle\operatorname{gal}(\sigma)\rangle$.
For $p\mid N$, AH11 shows it never occurs. With $\eta_F=1$ in the split
case and $-1$ in the inert case, these are precisely all three clauses of
\texttt{HostDer.eta}.

\paragraph{AH13: constructing the actual half-different ideal.}
Let $\mathfrak D=\mathfrak D_{E/K}$.
The relative inertia group for $E/K$ is $I(P)\cap\Psi$. At a prime above
$d$ it is trivial by (AH.28); outside $N$ it is trivial by AH9.
In a finite extension of number fields, a prime divides the different
exactly when it is ramified. This follows locally from the trace dual
over discrete valuation rings, or from the derivative of a local integral
generator; in the unramified case that derivative is a unit.
Therefore
\[
 P\mid\mathfrak D\quad\Longrightarrow\quad P\mid Qq.           \tag{AH.31}
\]
Every such rational prime splits in $K$. Hence, for
$\mathfrak p=P\cap\mathcal O_K$,
$\overline{\mathfrak p}\ne\mathfrak p$.

The different is invariant under every $\sigma\in\operatorname{Gal}(E/\mathbb Q)$.
Indeed, $K$ is stable because the Galois group is abelian, and
\[
 \operatorname{Tr}_{E/K}(\sigma z)
       =\sigma|_K\bigl(\operatorname{Tr}_{E/K}(z)\bigr).
\]
This preserves integrality in $\mathcal O_K$, hence preserves the
codifferent and its inverse $\mathfrak D$. In particular the exponents
of $P$ and $\overline P$ in $\mathfrak D$ agree.
Order the finite set of primes of $K$ below the support of $\mathfrak D$,
and define
\[
 \mathfrak b=
 \prod_{\substack{P\mid\mathfrak D\\
              P\cap\mathcal O_K<\overline P\cap\mathcal O_K}}
                       P^{v_P(\mathfrak D)}.                  \tag{AH.32}
\]
Exactly one prime of each conjugate pair contributes. Unique factorization
of nonzero ideals gives
\[
 \mathfrak D=\mathfrak b\overline{\mathfrak b},\qquad
 \mathfrak b+\overline{\mathfrak b}=\mathcal O_E.                 \tag{AH.33}
\]
An element of $\operatorname{Gal}(E/K)=\Psi$ preserves contractions to
$K$, commutes with conjugation, and preserves the exponents of the
different. It therefore preserves the selected product:
$\sigma\mathfrak b=\mathfrak b$ for every $\sigma\in\Psi$.
This proves every clause of the source theorem named
\texttt{KcDiff.principal}. Its conclusion is a selected integral ideal
$\mathfrak b$. It does not assert that this ideal is generated by one
element and uses no class-number hypothesis.

\paragraph{AH14: the inverse ideal is self-dual for the conjugated trace.}
The codifferent is
\[
 \mathfrak D^{-1}
   =\{z\in E:\operatorname{Tr}_{E/K}(z\mathcal O_E)
                                      \subset\mathcal O_K\}.
\]
For any nonzero fractional $\mathcal O_E$-ideal $J$, its trace dual satisfies
\[
 J^*=\{x:\operatorname{Tr}_{E/K}(xJ)\subset\mathcal O_K\}
       =\mathfrak D^{-1}J^{-1}.                  \tag{AH.34}
\]
To prove the equality, the trace condition is equivalent to
$xJ\subset\mathfrak D^{-1}$: for $y\in J$ and $a\in\mathcal O_E$,
the element $ay$ is still in $J$. Since a nonzero ideal in a Dedekind domain
is invertible, the latter inclusion is equivalent to
$x\in\mathfrak D^{-1}J^{-1}$.

Take any integral ideal $\mathfrak b$ satisfying (AH.33), not necessarily
the particular choice (AH.32). It is nonzero because the different is
nonzero. Conjugation maps $\mathfrak b^{-1}$ bijectively onto
$(\overline{\mathfrak b})^{-1}$: multiply the defining inclusion
$y\mathfrak b\subset\mathcal O_E$ by conjugation in both directions.
Apply (AH.34) with $J=(\overline{\mathfrak b})^{-1}$ to obtain
\[
 J^*=(\mathfrak b\overline{\mathfrak b})^{-1}
             \overline{\mathfrak b}=\mathfrak b^{-1}.
\]
Consequently, for every $x\in E$,
\[
 x\in\mathfrak b^{-1}
 \quad\Longleftrightarrow\quad
 \operatorname{Tr}_{E/K}(x\overline y)\in\mathcal O_K
                 \text{ for every }y\in\mathfrak b^{-1}.        \tag{AH.35}
\]
This proves \texttt{KcDiffSmulMem.unimodular} in both directions, with
the literal $\mathcal O_K$-valued trace pairing and the literal inverse
fractional ideal. No principal generator or metric comparison has been
substituted for this pairing.

\paragraph{AH15: a finite-kernel formula for the root discriminant.}
Let $E\subset L=\mathbb Q(\zeta_N)$ be any subfield, and let
$U\subset U_N$ be its fixing subgroup. For $f\mid N$ define
\[
 K_f=\ker(U_N\longrightarrow U_f),\qquad
 \rho_U(f)=\frac{|K_f\setminus U|}{|K_f|},\qquad
 S(U)=\sum_{a\mid N}\Lambda(a)\rho_U(N/a),        \tag{AH.36}
\]
where $\Lambda(p^j)=\log p$ for $j\ge1$ and is zero on other positive
integers. Then
\[
           \log\operatorname{rd}(E)=S(U),\qquad
 \operatorname{rd}(E)=|\operatorname{disc}E|^{1/[E:\mathbb Q]}.
                                                               \tag{AH.37}
\]
Here is a derivation, so this formula introduces no analytic input.
By AH9, $\mathcal O_L=\mathbb Z[\zeta_N]=\mathcal O_E[\zeta_N]$.
If $f_E$ is the minimal polynomial of $\zeta_N$ over $E$, the relative
different is $(f_E'(\zeta_N))$. The monogenic different formula follows
directly from Lagrange interpolation: for a monic separable polynomial
of degree $a$, the sums
$\sum_{f_E(z)=0}z^j/f_E'(z)$ are $0$ for $j<a-1$ and $1$ for $j=a-1$.
Reducing higher powers by $f_E$, these identities give a trace-dual basis
of $\mathcal O_E[\zeta_N]$ in
$f_E'(\zeta_N)^{-1}\mathcal O_E[\zeta_N]$, with a triangular change of
basis of determinant $\pm1$. This proves the ideal formula in both
directions. Also
\[
 f_E'(\zeta_N)=
       \prod_{\substack{u\in U\\u\ne1}}(\zeta_N-\zeta_N^u).
\]
The discriminant tower identity, obtained from the determinant of the
trace pairing in a tower, is
\[
 |\operatorname{disc}L|
   =|\operatorname{disc}E|^{[L:E]}
      N_{L/\mathbb Q}(\mathfrak D_{L/E}).
\]
Subtract its logarithm from the corresponding product for $E=\mathbb Q$.
This gives
\[
 [L:E]\log|\operatorname{disc}E|
       =\sum_{u\notin U}\log|N_{L/\mathbb Q}(\zeta_N-\zeta_N^u)|.
                                                               \tag{AH.38}
\]
If $h$ is the additive order of $u-1$ modulo $N$, then $h>1$ and
$\zeta_N-\zeta_N^u$ is a root of unity times $1-\eta$ for a primitive
$h$th root $\eta$. Its norm has logarithm
$\varphi(N)\Lambda(h)/\varphi(h)$. Indeed,
$\Phi_h(1)=p$ for $h=p^j$ and $1$ otherwise. The last evaluation follows
by induction from
$\prod_{a\mid h,\,a>1}\Phi_a(1)=h$, obtained by evaluating
$(X^h-1)/(X-1)$ at $1$. Taking the norm first in $\mathbb Q(\eta)$ and
then in $L$ proves the stated logarithm.

For each divisor $h>1$ of $N$, the elementary identity
\[
 \frac{\Lambda(h)}{\varphi(h)}
   =\sum_{\substack{a\mid N\\h\mid a}}
             \Lambda(a)\frac{\varphi(N/a)}{\varphi(N)}           \tag{AH.39}
\]
holds. If $h$ is not a prime power both sides vanish. If $h=p^j$ and
$N=p^v m$, $p\nmid m$, the right side is the geometric totient sum
$\log p\,\sum_{i=j}^v\varphi(p^{v-i})/\varphi(p^v)
=\log p/\varphi(p^j)$.
The condition $h\mid a$ says $u\equiv1\pmod{N/a}$, and
$|K_f|=\varphi(N)/\varphi(f)$ by surjective reduction.
Interchange the two finite sums in (AH.38), use (AH.39), and divide by
$[L:E][E:\mathbb Q]=\varphi(N)$. This proves (AH.37).

\paragraph{AH16: all three numerical discriminant bounds and its support.}
First, for any subgroup $U\subset U_N$,
\[
                  [U_N:U]\le\exp(2S(U)).        \tag{AH.40}
\]
Choose the least divisor $c_1$ of $N$ for which $K_{c_1}\subset U$.
If $K_a,K_b\subset U$, then $K_{\gcd(a,b)}\subset U$:
for $u\equiv1\pmod{\gcd(a,b)}$, CRT gives a unit $v$ with
$v\equiv1\pmod a$, $v\equiv u\pmod b$, and then
$u=v(v^{-1}u)\in K_aK_b$. Lift from the lcm modulus to $N$ using AH1.
It follows by minimality that $c_1\mid f$ whenever $K_f\subset U$.
If $c_1\nmid f$, the proper subgroup $K_f\cap U$ has index at least $2$
in $K_f$, so $\rho_U(f)\ge1/2$. Counting prime powers shows
\[
 \sum_{\substack{a\mid N\\c_1\nmid N/a}}\Lambda(a)=\log c_1:
\]
for $p^v\Vert N$ and $p^b\Vert c_1$ the sum contains precisely the last
$b$ powers of $p$. Hence $2S(U)\ge\log c_1$, while
$[U_N:U]\le\varphi(c_1)\le c_1$. This proves (AH.40).

Apply this to $U=\ker\operatorname{art}$ of a raw host.
Surjectivity and $|G|=4n$ give
\[
                         4n\le\operatorname{rd}(E)^2.           \tag{AH.41}
\]
For the upper bound, $\rho_U(f)\le1-\varphi(f)/\varphi(N)$ because
$1\in K_f\cap U$. At $f=N/d=Qq$, the image of $K_f$ has order $2$ by
(AH.28), so $\rho_U(f)=1/2$. At $f=N/q=Qd$, the group $K_f$ is cyclic,
being isomorphic to $U_q$ by CRT, and \texttt{art\_\allowbreak q} makes its image
have exponent dividing $2$. Thus its image has order at most $2$, and
$\rho_U(N/q)\le1/2$. In the remaining terms of (AH.36), retain just the
improvement at powers of $2$ over the bound $\rho\le1$.
Writing $N=2^v m$ with $m$ odd, their total improvement is
\[
 \log2\sum_{j=1}^v\frac{\varphi(N/2^j)}{\varphi(N)}
       =\log2.
\]
Since $q,d$ are distinct odd primes, these corrections are disjoint.
Using $\sum_{a\mid N}\Lambda(a)=\log N$ gives
\[
 \begin{split}
 \log\operatorname{rd}(E)
 &\le\log N-\log2-\tfrac12\log q-\tfrac12\log d,\\
 \operatorname{rd}(E)&\le\frac Q2\sqrt{qd}.       
 \end{split}
 \tag{AH.42}
\]

For the lower bound, \texttt{art\_\allowbreak h} implies that every $u\in U$ reduces
to $1$ or $-1$ modulo $T_e$. Write again $N=2^v m$, $m$ odd; here
$v\ge e+2$. For $0\le i\le e+1$, the reduction of $K_{2^i m}$ in
$U_{T_e}$ has order $2^{e+1}/\varphi(2^i)$ and equal-size fibers.
Among its elements, at most two of $\{1,-1\}$ are possible if $i=0,1$,
and at most one if $i\ge2$. Therefore
\[
 \rho_U(2^i m)\ge
     1-\frac{\varphi(2^i)+\mathbf1_{i<2}}{2^{e+1}}.
\]
Use in (AH.36) the distinct prime-power terms $a=2^{v-i}$ for these
indices and the term $a=d$. All other terms are nonnegative.
Since $\sum_{i=0}^{e+1}\varphi(2^i)=2^{e+1}$ and
$\rho_U(N/d)=1/2$, this proves
\[
 \begin{split}
 \log\operatorname{rd}(E)
   &\ge(e+1-2^{-e})\log2+\tfrac12\log d,\\
 2^{e+1-2^{-e}}\sqrt d&\le\operatorname{rd}(E).  
 \end{split}
 \tag{AH.43}
\]
The counts also justify the stated reductions when $i=0$, since
$U_1$ is the trivial group.

No rational prime outside $N$ divides $\operatorname{disc}E$, by AH9 and
the ramification criterion for the different. The prime $d$ does divide
it because its inertia group is nontrivial. The same is true of $2$:
CRT makes $K_m$ surject onto $U_{T_e}$, and \texttt{art\_\allowbreak h} makes its
image surject onto the nontrivial group $H_e$; thus $I_2\ne1$.
Finally $E/F$ is unramified at every finite prime, so the discriminant
tower identity gives $|\operatorname{disc}E|=|\operatorname{disc}F|^2$.
There are $2n$ conjugate pairs of complex embeddings of $E$, an even
number, so its discriminant has positive sign. Thus
\[
 \operatorname{disc}E=(\operatorname{disc}F)^2,\qquad
 2d\mid\operatorname{disc}E,\qquad
 p\mid\operatorname{disc}E\Longrightarrow p\mid N
 \quad(p\text{ prime}).                         \tag{AH.44}
\]
These are the literal \texttt{discr\_\allowbreak sq}, \texttt{discr\_\allowbreak primes},
\texttt{rd\_\allowbreak ge}, \texttt{rd\_\allowbreak le} and \texttt{rd\_\allowbreak sq\_\allowbreak ge} fields.

\paragraph{AH17: the full raw-to-derived host theorem.}
The preceding arguments prove \texttt{HostRaw.der} for every raw host,
with no new arithmetic premise. The exact coverage of its thirty fields
is as follows:
\[
\begin{array}{p{0.72\linewidth}l}
\texttt{art\_\allowbreak surj, cm, conj\_\allowbreak iff, finrank\_\allowbreak E, finrank\_\allowbreak F, mem\_\allowbreak F,
cmType, ps\_\allowbreak root, art\_\allowbreak ch}&\mathrm{AH10}\\[2pt]
\texttt{frob, frob\_\allowbreak any, inertia, decomp}&\mathrm{AH9,AH10}\\[2pt]
\texttt{ps\_\allowbreak d, inertia\_\allowbreak d, decomp\_\allowbreak typeI, decompConj, unramEF}
 &\mathrm{AH11}\\[2pt]
\texttt{splitK, ordinary, eta}&\mathrm{AH12}\\[2pt]
\texttt{diffSupport, principal}&\mathrm{AH13}\\[2pt]
\texttt{unimodular}&\mathrm{AH14}\\[2pt]
\texttt{discr\_\allowbreak sq, discr\_\allowbreak primes, rd\_\allowbreak ge, rd\_\allowbreak le, rd\_\allowbreak sq\_\allowbreak ge}
 &\mathrm{AH15,AH16}\\[2pt]
\texttt{norm\_\allowbreak mu}&\mathrm{AC2,AC3}.
\end{array}                                                     \tag{AH.45}
\]
For the last line, the raw hypotheses already give
$[E:\mathbb Q]=4n$, $d>3$ prime, and $K=\mathbb Q(\sqrt{-d})$;
AC2--AC3 prove $N_{E/K}(\zeta)=1$ for every root of unity, including its
$2$-primary part. Thus that result introduces no dependency on
\texttt{HostDer.norm\_\allowbreak mu} itself. Appending these proved fields to $H$
gives the literal \texttt{HostAx.ofRaw H}, whose raw projection is $H$.
In particular the arithmetic input to HC4's host witness is now supplied.

\paragraph{AH18: the complete prescribed-character field theorem.}
Let $\psi$ be any primitive complex Dirichlet character of conductor $M>0$,
and set $\Xi=\langle\psi\rangle$, $o=|\Xi|$. Choose $c_0,e_1$ by AH6
and set
\[
                   e_0=\max(e_1,2,4M,64o).       \tag{AH.46}
\]
For every $e\ge e_0$, every $q,d$ satisfying (AH.2), and every integer
$b\ge4M$, assume additionally
\[
 q\ne3,\qquad \gcd(M,qd)=1,\qquad
 (-d/p)=1\text{ for every odd prime }p\le b.     \tag{AH.47}
\]
Use the frame from AH2, the Artin map from AH3, the fixed field from AH4,
and AH17 to obtain an actual $H:\texttt{HostAx}\,\phi$.
Then $\phi.e=e$, $n=c_0\,2^e$, $H.q=q$, $H.d=d$, and
$\phi.o=o=(\texttt{Level.cyc}\,\psi).o$.

Here are all the remaining quantitative rows. Since
$Q\mid W_e\mid4M2^e$ and $n\ge2^e$,
\[
 Q\mid\operatorname{lcm}(T_e,M),\qquad Q\le4Mn.
\]
Also $4M\le e<2^e$ implies
$Q\le4M2^e\le2^{2e}\le n^2$. The bound $o\le2^o$ gives
$o^{64}\le2^{64o}\le2^e$. If an odd prime $\ell$ divides $Q$, it divides
$M$; hence $4\ell\le4M\le b$, and the splitting hypothesis in (AH.47)
applies. Thus the eight fields of \texttt{FieldRows} are all proved:
the first-sign Legendre formula and the actual square root of $-q$ come
from AH3--AH4, exactness from AH5, $Q\le n^2$ and $o^{64}\le2^e$ from
these inequalities, $q\ne3$ and splitting below $b$ from (AH.47), and
$4\ell\le b$ for odd $\ell\mid Q$ from the divisor argument.
AH7 supplies \texttt{KeptInXp} and
\texttt{ExceptionalCharacterKept}; AH8 supplies \texttt{KerChi} and
\texttt{TwistForm}. This is every conjunct of the literal
\texttt{UpperTheory.Top.fieldChi} theorem, with $c_0\ge1$ and with
$c_0,e_0$ chosen before $e,q,d,b$.

The result does not select primes satisfying the simultaneous analytic
pair window, does not supply \texttt{VVle}, and does not prove a top-window
estimate. The original $q,d$ congruence and splitting hypotheses are
preserved exactly. HC10 can now use this constructed field, and SZ13 can
construct ranks over its actual frame when the required order window is
supplied. Those joins retain their analytic inputs and the separate
\texttt{HostGood}, \texttt{WinOK}, \texttt{LevelPrime},
$d_u$--$D_s$, \texttt{Over} and \texttt{RecOK} interfaces.



\section{Constructing the ambient periods and their arithmetic}
\label{sec:period-arithmetic}

We now supply the ambient-period and CM-$j$ inputs retained in TL17
and TL21. The input is the literal arithmetic host $H$: its number
fields $K=\mathbb Q(\sqrt{-d})\subset E$, chosen embeddings, conjugation,
and inverse half-different $A=H.\mathrm{AA0}$ are fixed. In particular,
the integral trace dual of $A$ is $A$ itself. The construction of the
host is a separate arithmetic question. None of the arguments below
uses the existence of column laws, comparison maps, or a later-stage
geometric point.

Write $\iota:K\hookrightarrow\mathbb C$ for the distinguished embedding,
$\theta=\sqrt{-d}$, and $\sigma$ for conjugation. The actual hypotheses
give $d\equiv7\pmod {16}$ and $\iota(\theta)=i\sqrt d$.
The trace form is linear in its first variable:
\[
 h(x,y)=\operatorname{Tr}_{E/K}(x\bar y).
                                                               \tag{PA.1}
\]
This convention fixes all transposes and conjugates below.

\paragraph{PA1: ideal lattices, area, and the discriminant.}
Let $I$ be a nonzero integral ideal of $\mathcal O_K$, and let
$\lambda\in\mathbb C^\times$. Then
$\Lambda=\lambda\iota(I)$ is a genuine rank-two complex lattice.
Indeed, $I$ is a free abelian group of rank two: it has finite index
in $\mathcal O_K$, since for $0\ne a\in I$ the full-rank multiplication
matrix of $a$ has nonzero integral determinant. For such an $a$,
$a$ and $\theta a$ belong to the rational span of $I$ and have
real-linearly independent complex images. Consequently any integral
basis of $I$ gives a real basis of $\mathbb C$. Its image of
$\mathbb Z^2$ is discrete. The additive-subgroup closure in the source
is already this additive subgroup; it is not a topological closure.

For an integral basis $(\omega_1,\omega_2)$ of $\Lambda$, the source
area, defined as an infimum of positive areas of lattice pairs, is
\[
 A_\Lambda=|\operatorname{Im}(\omega_1\bar\omega_2)|>0.
                                                               \tag{PA.2}
\]
Every other pair has area an integral determinant times (PA.2),
and the basis pair attains the least positive value. This proves
the exact area dictionary, including the absolute value.

The lattice sums
$g_2=60\sum_{0\ne\omega\in\Lambda}\omega^{-4}$ and
$g_3=140\sum_{0\ne\omega\in\Lambda}\omega^{-6}$ converge absolutely:
the number of lattice points in a disk of radius $R$ is $O((1+R)^2)$.
They therefore agree with the period-pair sums by reindexing.
Their scaling laws are $g_2(\lambda\Lambda)=\lambda^{-4}g_2(\Lambda)$
and $g_3(\lambda\Lambda)=\lambda^{-6}g_3(\Lambda)$.
Also
\[
 \Delta_\Lambda=g_2^3-27g_3^2\ne0,\qquad
 \Delta_{\lambda\Lambda}=\lambda^{-12}\Delta_\Lambda.
                                                               \tag{PA.3}
\]
Here is the standard analytic justification of nonvanishing. The
normally convergent Weierstrass series defines an even elliptic
function $\wp$ with one double pole per fundamental parallelogram.
Comparing Laurent series and cancelling poles gives
$\wp'^2=4\wp^3-g_2\wp-g_3$. The three nonzero half-period classes
are zeros of $\wp'$, by oddness and periodicity. Their $\wp$-values
are distinct: otherwise $\wp-e$ would have two distinct zeros of
order at least two, despite having just one pole of order two.
The equality of numbers of zeros and poles follows by integrating
the logarithmic derivative around a translated fundamental
parallelogram, whose opposite edges cancel. Thus the cubic has
three distinct roots, proving (PA.3). This argument also fixes the
source's normalization of $g_2,g_3,\Delta$.

\paragraph{PA2: the specified normalized period pair and half-period.}
Take $\Lambda_0=\iota(\mathcal O_K)$. Choose
$\delta^{12}=\Delta_{\Lambda_0}$, so $\delta\ne0$, and use PA1
to choose a period pair with lattice $\Lambda=\delta\Lambda_0$.
Then $\Delta_\Lambda=1$. Set $\omega_0=\delta$ and $t_0=\delta/2$.
The exact three clauses required in
\texttt{IsTwoTorsionPointUnkilled} are
\[
 2t_0\in\Lambda,\qquad
 \frac{1+i\sqrt d}{2}t_0\notin\Lambda,\qquad
 \frac{1-i\sqrt d}{2}t_0\notin\Lambda.                         \tag{PA.4}
\]
The first uses $1\in\mathcal O_K$. For either other clause,
membership and cancellation of $\delta$ would make
$(1\pm i\sqrt d)/4$ an algebraic integer. Its sum with its complex
conjugate is $1/2$, which is not an algebraic integer: a rational
algebraic integer is an integer, by the rational-root argument for
a monic integral polynomial. This proves (PA.4). In particular
$t_0\notin\Lambda$, since $(1+\theta)/2\in\mathcal O_K$; that
integrality follows from
$X^2-X+(d+1)/4$.
These are the actual period and point in
\texttt{exists\_\allowbreak periodPair\_\allowbreak lattice\_\allowbreak eq\_\allowbreak lamOf}, with no theta
nonvanishing input.

\paragraph{PA3: four copies of the trace form have a rational isometry.}
The host's embeddings indexed by $\tau$ are exactly the $K$-embeddings
of $E$. Conjugation commutes with them, since both composites become
ordinary complex conjugation after a complex embedding. Hence
$h(x,x)$ is fixed by conjugation in $K$, and is rational. If $x\ne0$,
\[
 \iota(h(x,x))=\sum_\tau |\tau(x)|^2>0.                       \tag{PA.5}
\]
Gram--Schmidt over $K$ therefore gives a basis $v_1,\ldots,v_n$
with $h(v_i,v_j)=0$ for $i\ne j$ and
$h(v_i,v_i)=a_i\in\mathbb Q_{>0}$. At each step subtract
$h(x,v_i)a_i^{-1}v_i$; positivity ensures that the norm of every
nonzero new vector is nonzero. No square root in $K$ is needed.

For each positive rational $a$, the four-squares theorem gives
$p,q,r,s\in\mathbb Q$ with $p^2+q^2+r^2+s^2=a^{-1}$.
One obtains this rational version directly by writing $a=A/B>0$,
representing $AB$ as a sum of four integer squares, and dividing
the four integers by $A$. The matrix
\[
 C=\begin{pmatrix}
 p&q&r&s\\-q&p&-s&r\\-r&s&p&-q\\-s&-r&q&p
 \end{pmatrix}
 \quad\hbox{satisfies}\quad aCC^{\mathsf T}=aC^{\mathsf T}C=1.
                                                               \tag{PA.6}
\]
Its four rows, applied to four copies of each $v_i$, give $4n$
orthonormal vectors in $E^4$. They are a basis, so their coordinate
map is a $K$-linear isometry
\[
 \gamma_0:E^4\longrightarrow K^{\{1,2,3,4\}\times\{1,\ldots,n\}},
 \qquad
 \sum_w\gamma_0(x)_w\overline{\gamma_0(y)_w}
       =\sum_{a=1}^4\operatorname{Tr}_{E/K}(x_a\bar y_a).
                                                               \tag{PA.7}
\]
This is precisely \texttt{IsoKG}; the source's equality after
embedding $K$ into $E$ is equivalent to (PA.7). The only classical
integer theorem used in (PA.6) is Lagrange's four-squares theorem,
the same \texttt{Nat.sum\_\allowbreak four\_\allowbreak squares} used in the pinned proof.

\paragraph{PA4: the complex extension is explicit, unique, and algebraic.}
Let $b_w=\gamma_0^{-1}(e_w)$ and put
$B_{(a,\tau),w}=\tau((b_w)_a)$. The trace identity gives
\[
                         B^*B=1.
                                                               \tag{PA.8}
\]
Indeed, its $(w',w)$ entry is the embedded trace pairing of
$b_w$ and $b_{w'}$. Since $B$ is square, also $BB^*=1$.
Expanding $x=\sum_w\gamma_0(x)_w b_w$ gives
\[
 (\tau(x_a))_{a,\tau}=B(\iota(\gamma_0(x)_w))_w.
                                                               \tag{PA.9}
\]
Thus $g_{\mathbb C}=B^*$ is the required complex-linear extension.
The vectors on the left of (PA.9), as $x$ varies, span the entire
complex space, because they include the columns of $B$; this also
proves uniqueness. Every entry of $B$ is algebraic, so the entries
of $g_{\mathbb C}$ are algebraic. Consequently this specific map
and its inverse preserve vectors with algebraic coordinates.
It is not necessary to assume that an arbitrary complex-linear
extension preserves such vectors.

\paragraph{PA5: separate the two primes above a split rational prime.}
Suppose $p$ is prime and either $p=2$ or $(-d/p)=1$.
There are a prime $\mathfrak p\mid p$ and $\alpha,\beta\in\mathcal O_K$
such that
\[
 \alpha\in\mathfrak p,\qquad \beta\notin\mathfrak p,
 \qquad \beta=\bar\alpha.                                    \tag{PA.10}
\]
For $p=2$ use $a=(1+\theta)/2$. Its trace is one, and its norm
$(d+1)/4$ is even. For odd $p$, choose $r\not\equiv0\pmod p$
with $r^2\equiv-d\pmod p$ and use $a=r+\theta$. Its norm is
divisible by $p$ and its trace $2r$ is prime to $p$.
In either case choose a prime ideal above $p$. Exactly one of
$a,\bar a$ belongs to it: at least one does by their product,
and both cannot by their trace. Relabel them to obtain (PA.10).
The existence of a prime above $p$ follows by taking a maximal
ideal of the nonzero finite ring $\mathcal O_K/p\mathcal O_K$;
that ring is nonzero since $\mathcal O_K$ is a free integral
module of rank two. This proof includes the dyadic case.

\paragraph{PA6: a local basis with a denominator outside one prime.}
Let $L\subset K^r$ be a full finitely generated
$\mathcal O_K$-lattice. Localize at $\mathfrak p$. The localization
is a discrete valuation ring, and $L_{\mathfrak p}$ is free of
rank $r$: successive choice of a coordinate with least valuation
and elementary elimination proves the usual basis theorem over
a discrete valuation ring. Its basis can be multiplied, vector
by vector, by elements outside $\mathfrak p$ to give vectors
$b_j\in L$ which remain a local basis. Let $g$ have these columns.
Clearing the coordinates of a finite generating family of $L$
by one $t\notin\mathfrak p$ gives
\[
       t g^{-1}L\subset\mathcal O_K^r,
       \qquad t g\mathcal O_K^r\subset L.                      \tag{PA.11}
\]
The second containment already holds with $t=1$, before taking
the common denominator. This is the exact two-sided basis
approximation consumed by the ambient construction. It does not
assert that the global lattice is a free $\mathcal O_K$-module.

\paragraph{PA7: unitary approximation at the selected split prime.}
Here is an explicit algebraic replacement for a named approximation
theorem. Use (PA.10) and put $q=\alpha/\bar\alpha$ and
$e_N=(1+q^N)^{-1}$. In the $\mathfrak p$-adic topology,
\[
       e_N\longrightarrow1,\qquad
       \bar e_N=\frac{q^N}{1+q^N}\longrightarrow0.              \tag{PA.12}
\]
In particular these expressions are defined for $N>0$.
For a matrix $A$ with $1+A$ invertible, set
$c(A)=(1+A)^{-1}(1-A)$. Direct multiplication gives
$c(c(A))=A$ and $1+c(A)=2(1+A)^{-1}$.
For any matrix $S$ set
\[
      T_N=e_NS-\bar e_NS^*,\qquad T_N^*=-T_N.
                                                               \tag{PA.13}
\]
When $1+T_N$ is invertible, $c(T_N)$ is unitary, since
$(1-T_N)^*(1-T_N)=(1+T_N)^*(1+T_N)$ and both factors commute
with their inverses. If $S=c(A)$, (PA.12) gives
$c(T_N)\to A$. Invertibility holds for all sufficiently large
$N$, by the nonzero determinant at the limit.

To approximate an arbitrary $g\in\operatorname{GL}_r(K)$,
choose a scalar $a\ne0,-1$ with $\det(a1+g)\ne0$; only finitely
many scalars are excluded. Apply the preceding construction to
$a1$ and $a^{-1}g$, and multiply the resulting unitary matrices.
Their product $u_N$ tends to $g$. Thus, for some $N$,
$g^{-1}u_N$ and $u_N^{-1}g$ have entries in
$\mathcal O_{K,\mathfrak p}$. Clear their finitely many denominators
by one $t'\notin\mathfrak p$. With $u=u_N$ we have
\[
 u^*u=1,\qquad
 t'g^{-1}u,\ t'u^{-1}g\in M_r(\mathcal O_K).                   \tag{PA.14}
\]
The argument uses a single prime topology. Complex conjugation
is not assumed continuous for that topology; (PA.12), rather
than such a false continuity assertion, is the essential step.

\paragraph{PA8: the rational denominator and the actual AtP condition.}
Suppose $L$ is self-dual for the standard Hermitian pairing.
Use PA6--PA7 and replace $L$ by $L'=u^{-1}L$. Combining
(PA.11) and (PA.14) gives, for some $t\notin\mathfrak p$,
\[
                   tL'\subset\mathcal O_K^r,
                 \qquad t\mathcal O_K^r\subset L'.            \tag{PA.15}
\]
Self-duality also gives both containments with $\bar t$.
For example, pair $x\in L'$ with $te_j$ to obtain
$\bar t x_j\in\mathcal O_K$; conversely pair $\bar t e_j$
with every $y\in L'$ and use $ty_j\in\mathcal O_K$.
Now put $a=\bar\alpha t+\alpha\bar t$. It is integral and
fixed by conjugation, hence an integer. Moreover
$a\equiv\bar\alpha t\not\equiv0\pmod{\mathfrak p}$.
Both containments in (PA.15) therefore hold with the positive
integer $s=|a|$, and $p\nmid s$. This is slightly shorter than
the generic source proof, which uses $|a\bar a|$ and only
assumes integral norms. It is valid for the actual quadratic $K$.

Apply this to
$L=\gamma_0(A^4)\subset K^{4n}$. It is finitely generated over
$\mathcal O_K$ because $\mathcal O_E$ is finite over
$\mathcal O_K$. It spans: for $0\ne b\in A$, every $x/b\in E$
has a nonzero integer multiple integral over $\mathbb Z$, so
every $x$ has an integer multiple in $A$. Finally $A^4$ is
self-dual for the sum of trace forms; testing one coordinate
at a time proves the converse implication. Isometry (PA.7)
therefore makes $L$ self-dual. For $\gamma_p=u^{-1}\gamma_0$,
\[
 s\gamma_p(A^4)\subset\mathcal O_K^{4n},\qquad
 s\gamma_p^{-1}(\mathcal O_K^{4n})\subset A^4,
 \qquad (s,p)=1.                                               \tag{PA.16}
\]
This proves both literal clauses of \texttt{AtP}, with one
denominator working for every vector. It supplies more than the
pointwise denominators required by the record.

\paragraph{PA9: all ambient record fields are now supplied.}
Let $h_K$ be the class number of $K$. For each relevant split prime
choose $\mathfrak p_p$ as in PA5, a generator $\pi_p$ of
$\mathfrak p_p^{h_K}$, and $\gamma_p$ as in PA8. The generator
exists because the ideal class has order dividing $h_K$; it is
nonzero because the ideal is nonzero. Choose
\[
 r_p^{2h_K}=\frac{\iota(\pi_p)}{\iota(\bar\pi_p)}.
                                                               \tag{PA.17}
\]
For indices not satisfying the prime/splitting condition choose
arbitrary fallback values, using the prime above two and
$\gamma_0$. This yields all eight fields of \texttt{KapIn}
including the four universally quantified specifications.
PA2, PA3, PA4, and this choice yield every field of
\texttt{VacIn}: $L$, its literal lattice equality, $\omega_0$,
the three half-period clauses, $\gamma_0$, $g_{\mathbb C}$, its
exact evaluation formula, and the correction data.

The derivative generators $c_D$ and their one common coefficient
number field are supplied by the ideal-principalization and flat
base-change argument of TL17. Use those generators together with
the \emph{same} period and vacuum data. Then
\[
                \exists J:\texttt{AmbIn}(H),\quad J.\texttt{Ok}
                                                               \tag{PA.18}
\]
holds: the two \texttt{Ok} fields are exactly
$\delta^{12}=\Delta_{\Lambda_0}$ and \texttt{IsDgen}.
No normalized theta Taylor theorem is a premise of (PA.18).
Its existence does not yet assert \texttt{StdOK}, a vacuum
four-copy identity, or any of the five \texttt{Later} stages.

\paragraph{PA10: the correction product has genuine finite support.}
The elementary determinant argument also explains the correction
factor in later vacuum formulas. Throughout this paragraph $p$
ranges over the rational primes satisfying the literal
$\texttt{SplitK}(H,p)$ predicate; the product in (PA.21) is over
that same subtype. The arbitrary fallback values of \texttt{KapIn}
outside this subtype play no role. For a fixed rational isometry
$\gamma_0$, choose a positive integer $N$ with
\[
 N\gamma_0(A^4)\subset\mathcal O_K^{4n},\qquad
 N\gamma_0^{-1}(\mathcal O_K^{4n})\subset A^4.                  \tag{PA.19}
\]
Finite generation and rational spanning give such a common $N$.
At a rational prime $p\nmid N$, both $\gamma_0$ and $\gamma_p$
identify the same localized lattice with the standard one.
Their transition and its inverse are integral at each prime of
$K$ above $p$, by (PA.16). Its determinant is therefore a unit
there. Thus the source exponent
\[
 y_p=\operatorname{ord}_{\mathfrak p_p}
                  \det(\gamma_p\gamma_0^{-1})
                                                               \tag{PA.20}
\]
vanishes unless $p\mid N$. For a fixed prime, replacing
$\gamma_p$ by another isometry satisfying \texttt{AtP} changes
the determinant by a local unit, so $y_p$ is unchanged.
This proves both \texttt{YFinite} and \texttt{YIndep}.
Each $r_p$ in (PA.17) is nonzero and algebraic, and hence
\[
                     \kappa=\prod_p r_p^{-y_p}
                     \in\overline{\mathbb Q}^{\times}.         \tag{PA.21}
\]
This is a finite product for the actual data. No inference about
support is made merely from Lean's totalized \texttt{finprod}.
The assertion concerns independence from the local isometries;
it does not identify arbitrary independent choices of the roots
$r_p$ in (PA.17).

\paragraph{PA11: a modular function with an integral unit leading term.}
We record the classical modular calculation needed for the exact
CM integrality input. Use the usual normalization
\[
 j(\tau)=\frac{E_4(\tau)^3}{\Delta_{\rm mod}(\tau)},\qquad
 q=e^{2\pi i\tau},\qquad j=q^{-1}+744+\cdots\in q^{-1}\mathbb Z[[q]].
                                                               \tag{PA.22}
\]
The source calls this function \texttt{modularJ}. It is
$1728g_2^3/(g_2^3-27g_3^2)$ for the lattice
$\mathbb Z+\mathbb Z\tau$, not $g_2^3/\Delta_\Lambda$.
For completeness, the integral expansion follows directly from
$E_4=1+240\sum_{n\ge1}\sigma_3(n)q^n$,
$E_6=1-504\sum_{n\ge1}\sigma_5(n)q^n$, and
$\Delta_{\rm mod}=(E_4^3-E_6^2)/1728=q+O(q^2)$.
The integer divisibility needed for the quotient is
$12\mid5\sigma_3(n)+7\sigma_5(n)$, obtained term by term from
$12\mid5a^3+7a^5$. Expansion then gives an integral power
series for $\Delta_{\rm mod}/q$ with constant term one; its
formal inverse is integral by recursive coefficient comparison.
The standard Eisenstein-series transformation and Fourier
formulas justify these convergent series. They follow by
reindexing their absolutely convergent lattice sums and the
partial-fraction expansion of $\pi\cot\pi z$, respectively.
Nonvanishing of $\Delta_{\rm mod}$ on the upper half-plane
follows equally from PA1 and the nonzero normalization constant.

Let $F$ be holomorphic on the upper half-plane, invariant under
$\operatorname{SL}_2(\mathbb Z)$, and suppose it has a convergent
Laurent expansion at the cusp with finitely many negative powers,
all coefficients algebraic integers, and leading coefficient an
algebraic unit $c$. Then
\[
                       F=P(j),\quad P\in\overline{\mathbb Z}[X],
                       \operatorname{lc}(P)=c.                 \tag{PA.23}
\]
To prove this, subtract $c j^m$ to cancel the leading pole
$c q^{-m}$ and continue by decreasing $m$. The remainder is
invariant, holomorphic, and bounded at the cusp, and therefore
constant: on the compactified standard fundamental domain it
extends across the cusp and across the finite elliptic quotient
charts, so the maximum principle applies. The final constant is
its integral constant Fourier coefficient. This proves (PA.23)
with all coefficient assertions. If $F(\tau_0)=0$, division by
the algebraic unit $c$ gives a monic polynomial with algebraic
integer coefficients vanishing at $j(\tau_0)$. Integrality is
transitive, so $j(\tau_0)$ is integral over $\mathbb Z$.

The same conclusion holds if the initial expansion is in
$q^{1/n}$ with an algebraic-unit leading coefficient. Raising
$F$ to the $n$th power makes its leading exponent integral;
periodicity under $\tau\mapsto\tau+1$ forces every nonintegral
Fourier exponent to have zero coefficient. All remaining
coefficients are still algebraic integers, and the leading
coefficient is $c^n$. This is the exact unit-leading-expansion
argument in the source.

\paragraph{PA12: an explicit nonsquare determinant proves CM integrality.}
Suppose $a\tau^2+b\tau+c=0$, where $a,b,c\in\mathbb Z$,
$a\ne0$, and $\operatorname{Im}\tau>0$. Put
$D=4ac-b^2>0$, $k=D+|b|$, and
\[
 M=\begin{pmatrix}k-b&-c\\a&k\end{pmatrix},\qquad
 n=\det M=k^2-bk+ac.
                                                               \tag{PA.24}
\]
Then $M\tau=\tau$ and
$4n=(2k-b)^2+D$. In particular $n>0$ and is not a square.
Indeed $m=2k-b\ge D>0$, so
$m^2<m^2+D<(m+1)^2$; the right side of $4n=m^2+D$
cannot be an integer square.

Consider the finite set of holomorphic functions $j(Az)$ as
$A$ ranges over all integral matrices of determinant $n$.
It is finite because left integral row operations reduce $A$
to a triangular matrix
$\left(\begin{smallmatrix}u&v\\0&w\end{smallmatrix}\right)$
with $u,w>0$, $uw=n$, and $0\le v<w$. Left multiplication by
$\operatorname{SL}_2(\mathbb Z)$ does not change the function.
Right multiplication permutes this finite set. Consequently
\[
             F(z)=\prod_{G\in\{j(A\,\cdot)\}}(j(z)-G(z))
                                                               \tag{PA.25}
\]
is invariant and holomorphic. Each factor has an integral
Laurent expansion in $q^{1/n}$ with an algebraic-unit leading
coefficient. In fact
$j((uz+v)/w)=\zeta_w^{-v}q^{-u/w}+O(1)$, and
$u/w\ne1$ because $n$ is nonsquare. The leading term of the
difference is therefore either $q^{-1}$ or
$-\zeta_w^{-v}q^{-u/w}$, with no leading cancellation.
All remaining coefficients lie in a cyclotomic integer ring.
The product is nonzero as a function and has a unit leading
term. The factor supplied by $M$ vanishes at $\tau$.
PA11 now proves that $j(\tau)$ is an algebraic integer.
Using the set of distinct functions, as the source does, avoids
any unjustified multiplicity count.

If a lattice $\mathbb Z+\mathbb Z\tau$ is stable under a
scalar $\alpha\notin\mathbb Z$, write
$\alpha\tau=p\tau+q$ and $\alpha=r\tau+s$ with integers
$p,q,r,s$. Then $r\ne0$ and
$r\tau^2+(s-p)\tau-q=0$. Thus the preceding argument applies.
An action by $\mathcal O_K$ supplies such an $\alpha$, since
$\mathcal O_K$ has integral rank two. This proves the literal
\texttt{isIntegral\_\allowbreak modularJ\_\allowbreak of\_\allowbreak ringOfIntegers} assertion.

\paragraph{PA13: the actual base $j$-value and the finite FieldOK extension.}
The base lattice in the source is $\iota(H.\mathrm{ccK})$,
where $H.\mathrm{ccK}\ne0$. By PA1 it has a period pair.
After homothety and a possible change of oriented basis it is
$\mathbb Z+\mathbb Z\tau$ with $\operatorname{Im}\tau>0$,
and is stable under $\mathcal O_K$. PA12 and the invariant
dictionary therefore give
\[
 \texttt{jLat}(\texttt{baseLat}(H))
       =1728\,g_2^3/\Delta\quad\hbox{integral over }\mathbb Z.
                                                               \tag{PA.26}
\]
This proves the exact \texttt{isIntegral\_\allowbreak jLat} target.

The remaining ideal in \texttt{FieldOK} is also explicit.
From $\mathfrak D_{E/K}=\mathfrak b\bar{\mathfrak b}$ and
$A=\mathfrak b^{-1}$ we have
$A\mathfrak D_{E/K}=\bar{\mathfrak b}$. Hence each
\texttt{ellLat} is the extension of a Galois conjugate of
$\bar{\mathfrak b}$. Choose $h>0$ and $\beta$ with
$\bar{\mathfrak b}^{h}=(\beta)$, and adjoin one root
$z_g^h=\varphi_0(\sigma_g\beta)$ for each of the finitely many
$g$. Also adjoin (PA.26) to the prescribed starting number field
$M$. The resulting $M'$ is finite over $\mathbb Q$ and contains
$M$. In any compatible coefficient number field containing
$M'$, the extended ideal is $(z_g)$: their $h$th powers agree,
and equality of nonzero ideals follows prime by prime from
$h v_{\mathfrak q}(I)=h v_{\mathfrak q}(J)$.
This supplies both \texttt{FieldOK} clauses with their literal
embedding compatibility and universal quantifier over later
coefficient fields. The finite ideal extension is the argument
already reconstructed in TL17; PA26 supplies its previously
retained CM-$j$ input. A fixed coefficient field may only be
used when it contains this prescribed $M'$.

\paragraph{PA14: normalized invariants and torsion coordinates are algebraic.}
Let $\Lambda=\delta\iota(I)$ be any actual curve-ideal lattice
with $\delta^{12}=\Delta_{\iota(I)}$. PA1 gives
$\Delta_\Lambda=1$, and PA12 gives integral $j(\Lambda)$.
Therefore
\[
 g_2(\Lambda)^3=j(\Lambda)/1728,\qquad
 g_3(\Lambda)^2=(g_2(\Lambda)^3-1)/27                         \tag{PA.27}
\]
show $g_2,g_3\in\overline{\mathbb Q}$. The other entries of
\texttt{DataIn} are embeddings of elements of the number field
$E$, hence algebraic as well.

If $mt\in\Lambda$, $m>0$, and $t\notin\Lambda$, then
\[
                        \wp_\Lambda(t),\ \wp'_\Lambda(t)
                              \in\overline{\mathbb Q}.         \tag{PA.28}
\]
One elementary proof keeps the curve and its normalizations
visible. The map $z\mapsto(\wp(z),\wp'(z)/2)$ identifies
$\mathbb C/\Lambda$ with the nonsingular curve
$y^2=x^3-g_2x/4-g_3/4$, carrying addition to the rational
chord-and-tangent law. The differential equation in PA1 and
the elliptic addition formula prove this identification; the
degree-two function $\wp$ is surjective, and its fibers are
the pairs $\{z,-z\}$. The $m$-torsion has exactly $m^2$ complex
points, since it is $(m^{-1}\Lambda)/\Lambda$.
Its equations and addition law are defined over
$\overline{\mathbb Q}$. If one coordinate were transcendental,
a field embedding fixing $\overline{\mathbb Q}$ could send it
to infinitely many values: send a transcendental $x$ to $x+n$,
extend to the algebraic closure of $\overline{\mathbb Q}(x)$,
and then, if desired, extend a transcendence basis to all of
$\mathbb C$. Each image would still be an $m$-torsion point,
contradicting finiteness. The second coordinate is algebraic
over the first by the displayed cubic. This proves (PA.28)
without assuming a transcendental period is algebraic.
Together, (PA.27)--(PA.28) reconstruct the actual
\texttt{admJ\_\allowbreak and\_\allowbreak torsion\_\allowbreak data\_\allowbreak algebraic} supplier.
They do not yet prove algebraicity of normalized sigma values
or Taylor coefficients at those points.

\paragraph{PA15: exact torsion-translation transfer to the even theta function.}
This last result has an explicit remaining premise. Suppose
$\Lambda$ is a genuine lattice, $s=\mathrm{sgStar}_\Lambda$
is its normalized odd theta function, and every rational-lattice
translate $U_xs$ is entire with algebraic Taylor coefficients at
zero. Here
\[
 U_xf(z)=f(z+x)\exp\!\left(-\frac{\pi z\bar x}{A_\Lambda}
                         -\frac{\pi x\bar x}{2A_\Lambda}\right).
                                                               \tag{PA.29}
\]
The ordinary analytic laws of $s$ require no arithmetic premise.
To see them, start from the normally convergent sigma product
$\sigma_\Lambda(z)=z\prod_{0\ne\omega\in\Lambda}
 (1-z/\omega)\exp(z/\omega+z^2/(2\omega^2))$.
It is odd, has derivative one at zero, and has simple zeros
exactly at the lattice. Its quasi-periods satisfy the Legendre
relation; this follows by integrating its logarithmic derivative
around a fundamental parallelogram containing its one simple
zero. For a positively oriented basis the relation is
$\eta_1\omega_2-\eta_2\omega_1=2\pi i$. It implies that there
is a complex $a$ with
$\eta_\omega=a\omega+\pi\bar\omega/A_\Lambda$ on both basis
periods, and hence on every period. Multiplying the sigma
product by $\exp(-az^2/2)$ gives the normalized $s$ and the
law $U_\omega s=\eta(\omega)s$, where the sign is positive
exactly on $2\Lambda$. These are the zero, derivative, and
translation laws used here. Their analytic normalization makes
no assertion that $a$ or a translated value is algebraic.
To identify this product with the literal chosen \texttt{sgStar},
note that any odd normalized entire function with these laws
vanishes at zero and therefore at every lattice point. Dividing
it by the product removes those possible simple poles; the
quotient is entire and periodic, hence constant on the compact
torus. The derivative-one normalization makes that constant one.

For rational-lattice points $x,y$ define
$\epsilon(x,y)=\exp(\pi(y\bar x-x\bar y)/(2A_\Lambda))$.
Expansion of exponents gives
\[
 U_xU_y=\epsilon(x,y)U_{x+y}.                                  \tag{PA.30}
\]
If $mx,ny\in\Lambda$, the signed area of this pair is an
integer multiple of $A_\Lambda$. Thus
$\epsilon(x,y)^{2mn}=1$, so these phases are nonzero algebraic
numbers. This retains the sign and the factor two of the source.

For $t\notin\Lambda$ with $2t\in\Lambda$, the source's even
theta function is $\vartheta_t=U_ts/(U_ts)(0)$. The denominator
is nonzero because the zeros of $s$ are precisely $\Lambda$;
its algebraicity follows from the stated translate premise.
Equation (PA.30) gives the identity of entire functions
\[
 U_x\vartheta_t=
       \frac{\epsilon(x,t)}{(U_ts)(0)}\,U_{x+t}s.                \tag{PA.31}
\]
Its Taylor coefficients are algebraic. For $\omega\in\Lambda$
the odd-theta law is
$U_\omega s=\eta(\omega)s$, where
$\eta(\omega)=1$ if $\omega/2\in\Lambda$ and $-1$ otherwise.
Combining it with (PA.30)--(PA.31) gives
\[
 U_{x+\omega}\vartheta_t=
 \frac{\epsilon(x+\omega,t)}{\epsilon(x,t)}
 \frac{\eta(\omega)}{\epsilon(x+t,\omega)}\,U_x\vartheta_t.
                                                               \tag{PA.32}
\]
The multiplier is a nonzero algebraic number. Entirety follows
from translation and multiplication by an exponential, and
Taylor coefficients commute with multiplication by the constant.
This is the complete stated theorem
\texttt{algebraicTaylor\_\allowbreak torsionTranslate\_\allowbreak evenTheta\_\allowbreak alt}; its
odd-theta translate premise remains explicit. PA14 supplies
algebraic elliptic invariants and torsion coordinates, which are
inputs to the sigma/Taylor reciprocity argument reconstructed
separately in TA; it does not replace that argument.

\paragraph{Boundary after the period construction.}
The ambient carrier, its actual local corrections, and the
CM-$j$/coefficient-field extension are now constructed from the
host. The primitive normalized odd-theta Taylor/reciprocity theorem
is supplied in the following TA section. The vacuum four-copy
identity and the exact normalized indicator and covariance
identities still need their own proofs.
The complete reference and twisted column constructions remain
conditional on the expressly stated host, index and later-stage
inputs. Nothing in PA1--PA15 supplies the independent global
spectral, metric, cohomological, or simultaneous-prime estimates.


\section{The arithmetic of normalized theta coefficients}
\label{sec:theta-arithmetic}

The arithmetic premise retained in PA15 can now be supplied. We prove
the actual one-dimensional statements \texttt{galSg}, \texttt{gal1},
and \texttt{use1}, including their coefficient-field and torsion-point
conditions. This section uses the literal host and curve datum of PA,
not a column law, a vacuum identity, or a later-stage geometric point.
It does not yet prove the multidimensional vacuum or covariance
identities. The ordinary uniformization theorem for a nonsingular
complex elliptic curve is a named classical input here. We explain
its precise uses, and the intervening algebraic descent arguments,
instead of incorporating that whole classical theory as a new axiom
with the desired theta conclusion.

Write $\Lambda=\lambda\iota(I)$ for the curve lattice, with
$\Delta_\Lambda=1$, and $A=\operatorname{area}(\Lambda)>0$. Let
$s_\Lambda=\mathrm{sgStar}(\Lambda)$ be the normalized odd theta
function described in PA15. Its derivative at zero is one, its
zeros are exactly $\Lambda$, and
\[
 s_\Lambda(z+\omega)=\eta(\omega)
  \exp\!\left(\frac{\pi}{A}(z+\omega/2)\bar\omega\right)s_\Lambda(z),
 \quad
 \eta(\omega)=\begin{cases}1&\omega\in2\Lambda,\\-1&\omega\notin2\Lambda.
 \end{cases}                                                    \tag{TA.1}
\]
The translation operator is the actual operator
\[
 U_xs(z)=s(z+x)\exp\!\left(-\frac{\pi}{A}z\bar x
                           -\frac{\pi}{2A}x\bar x\right).       \tag{TA.2}
\]
All statements about coefficients below use $[z^n]f=f^{(n)}(0)/n!$.
There is no claim that a general field embedding commutes with limits,
complex differentiation, or the exponential function.

\paragraph{TA14: analytic normalization without an arithmetic premise.}
We isolate this preliminary claim so that the dependency graph does
not use the arithmetic assumption of PA15 to prove that assumption.
For a genuine lattice the canonical product
\[
 \sigma(z)=z\prod_{0\ne\omega\in\Lambda}
       (1-z/\omega)\exp(z/\omega+z^2/(2\omega^2))
\]
converges normally on compact sets: the logarithm of each factor
is $O(|z/\omega|^3)$ away from finitely many factors, and the
lattice sum of $|\omega|^{-3}$ converges. It has simple zeros
exactly at $\Lambda$, is odd by pairing $\omega,-\omega$, and
has derivative one at zero. Differentiating its logarithm
gives $(\log\sigma)''=-\wp$. Therefore the increments of its
logarithmic derivative along periods are constants
$\eta_\omega$, additive in $\omega$. Integrating around a
positively oriented fundamental parallelogram, with the one
residue of the logarithmic derivative, gives Legendre's
identity $\eta_1\omega_2-\eta_2\omega_1=2\pi i$.
For $\operatorname{Im}(\bar\omega_1\omega_2)=A>0$, this says
$\eta_\omega=c\omega+\pi\bar\omega/A$ for one $c\in\mathbb C$.
The sign in the integrated period law is $-1$ on each basis
period, by oddness at its half-period. Composing the two laws
and using Legendre gives the sign
$(-1)^{r+s+rs}$ at $r\omega_1+s\omega_2$, which is one
exactly when both integers are even. Hence
$f(z)=e^{-cz^2/2}\sigma(z)$ has precisely (TA.1), with its
stated zeros, parity and derivative.

If an entire $G$ has those same translation laws, $G/f$ is
elliptic and has at most a single simple pole on the torus.
The residue theorem on a fundamental parallelogram forces
that residue to vanish, so it has no pole. The maximum
principle makes it constant; its value is $G'(0)$.
This proves uniqueness and identifies the source's chosen
$\mathrm{sgStar}$ with $f$, without any field-membership
assertion. The same uniqueness gives the scaling identity
$s_{a\Lambda}(az)=a s_\Lambda(z)$ for $a\ne0$, since
$\operatorname{area}(a\Lambda)=|a|^2A$.

Finally, for $v\notin\Lambda$ the function
$-f(u+v)f(u-v)/(f(u)^2f(v)^2)$ is elliptic in $u$,
has just the double pole $u^{-2}$ at zero, and vanishes
at $u=v$. Subtracting $\wp(u)$ and cancelling poles shows
that it is $\wp(u)-\wp(v)$. This is the exact sigma addition
identity used in TA3 and TA7. All analytic facts used below
are therefore available independently of odd-theta
coefficient algebraicity.

\paragraph{TA1: the actual conjugate subgroup is a period lattice.}
The definition of $\mathrm{lamConj}(M,\rho,\Lambda)$ selects an
additive subgroup with invariants $\rho(g_2),\rho(g_3)$. One must
justify that this is the correct lattice, rather than assume that
an arbitrary chosen subgroup is discrete. Here is the missing
classification argument. If either Eisenstein sum of orders four
or six is nonzero, it is summable: in the source the totalized sum
of a nonsummable family is zero. Summability implies that only
finitely many nonzero subgroup elements can have absolute value
at most one. Otherwise the corresponding inverse powers would
not tend to zero on the cofinite filter. Removing those finitely
many points leaves a neighborhood of zero containing only zero.
Thus the subgroup is discrete.

A discrete subgroup of the real plane has rank at most two.
For completeness, choose a shortest nonzero vector; if the group
is not contained in its real line, reduce along that vector and
choose a least positive perpendicular height. Subtracting integral
multiples of the two vectors places any point in a bounded
fundamental parallelogram. Discreteness and the minimal choices
then give an integral basis. The one-dimensional case is cyclic.
For a cyclic subgroup $a\mathbb Z$ the absolutely convergent sums give
\[
 g_2=\frac{4\pi^4}{3a^4},\qquad
 g_3=\frac{8\pi^6}{27a^6},\qquad g_2^3-27g_3^2=0;
\]
the zero subgroup also has discriminant zero. Hence a subgroup
with nonzero discriminant must be a rank-two lattice.

For two rank-two lattices with the same $g_2,g_3$, their Weierstrass
functions have the same Laurent series at zero. One can obtain
uniqueness directly from the coefficient comparison in TA2 below,
applied to $z^2\wp(z)$. The identity theorem extends equality as
meromorphic functions on $\mathbb C$, and their pole sets are the
two lattices. Thus the lattices are equal.

Existence for every pair $a,b$ with $a^3-27b^2\ne0$ follows from
classical elliptic uniformization: the smooth projective cubic
$y^2=x^3-ax/4-b/4$, with differential $dx/(2y)$, is
$\mathbb C/\Lambda'$ with $x=\wp_{\Lambda'}$ and
$y=\wp'_{\Lambda'}/2$. The differential normalization fixes the
scale, so its invariants are exactly $a,b$. Applying this with
$a=\rho(g_2)$ and $b=\rho(g_3)$ is legitimate because $\rho$ is
injective and the discriminant is one. Consequently the selected
\texttt{lamConj} is this unique genuine lattice $\Lambda'$.
In particular, if $\rho$ fixes both invariants, $\Lambda'=\Lambda$.
This also supplies the arbitrary-subgroup uniqueness clause of
the invariant dictionary, which PA1 deliberately did not assume.

\paragraph{TA2: a formal uniqueness theorem with all initial data.}
For a formal series $F$ over a characteristic-zero field set
\[
 W_s(F)=F'^2-FF''-sF^2.
\]
The cleared Weierstrass equation is
\[
 (W_s(F)'F-2W_s(F)F')^2
     =4W_s(F)^3-aW_s(F)F^4-bF^6.                               \tag{TA.3}
\]
For fixed $s,a,b$ there is at most one solution with
$F(0)=0$, $[z]F=1$, and $[z^2]F=0$. To prove this, write $F=zV$,
$V(0)=1$, and put $P=W_s(F)V^{-2}$. Cancelling $V^6$ in (TA.3)
gives
\[
 (zP'-2P)^2=4P^3-az^4P-bz^6,\qquad P(0)=1.                  \tag{TA.4}
\]
If two solutions $P,Q$ first differ in degree $n$, write
$P-Q=z^nE$, $E(0)\ne0$. Subtract (TA.4), cancel $z^n$,
and take constant terms. The two sides are
$-4(n-2)E(0)$ and $12E(0)$. Thus $(n+1)E(0)=0$, a
contradiction. This proves uniqueness of $P$.

For two original solutions $F=zV$ and $G=zT$, equality of $P$
gives equality of the derivatives of $V'/V$ and $T'/T$.
Their constant terms vanish because $[z^2]F=[z^2]G=0$.
In characteristic zero a formal series with derivative zero is
constant, so $V'/V=T'/T$; consequently $(V/T)'=0$ and
$V/T=1$. This proves $F=G$. The degree-two coefficient of
(TA.3), using the three initial conditions, also gives
\[
                         s=-2[z^3]F.                         \tag{TA.5}
\]
These are purely formal operations. In particular they commute
with every field embedding, with no convergence hypothesis on
the coefficientwise image series.

\paragraph{TA3: the differential equation and a CM polynomial.}
The purely analytic canonical product construction in TA14 gives
$s_\Lambda(z)=e^{-cz^2/2}\sigma_\Lambda(z)$ for the ordinary
Weierstrass sigma function. Since
$(\log\sigma_\Lambda)''=-\wp_\Lambda$, one obtains
\[
 W_c(s_\Lambda)=\wp_\Lambda s_\Lambda^2.                     \tag{TA.6}
\]
Equivalently (TA.6) follows by differentiating the sigma addition
formula twice at zero, so it uses no arithmetic property of $c$.
Together with $\wp'^2=4\wp^3-g_2\wp-g_3$, it gives (TA.3)
for the Taylor series of $s_\Lambda$. Formula (TA.5) identifies
the constant $c$ from that series.

Let $\alpha=\iota(\sqrt{-d})$, so
$\alpha\bar\alpha=d=:N$, $\alpha\Lambda\subset\Lambda$,
$\alpha\ne0$, and $\alpha^2=-d\ne d$. By (TA.1),
\[
 R(z)=\frac{s_\Lambda(\alpha z)^2}{s_\Lambda(z)^{2N}}         \tag{TA.7}
\]
is even and $\Lambda$-periodic: its exponential multipliers
cancel because $\alpha\bar\alpha=N$, and its signs have been
squared. Its only poles modulo $\Lambda$ occur at zero, with
order at most $2N-2$. An even elliptic function with poles only
there is a polynomial in $\wp$: successively subtract a multiple
of $\wp^j$ to cancel the highest even pole, then apply the
maximum principle to the holomorphic remainder on the compact
torus. Thus $R=Q(\wp)$ for a polynomial $Q$.

Suppose now $M\subset\mathbb C$ contains $\alpha,g_2,g_3$ and
all Taylor coefficients of $s_\Lambda$. By (TA.5), $c\in M$;
by (TA.6), the Laurent series of $\wp$ lies in $M((z))$.
The Laurent series of (TA.7) also lies in $M((z))$. Since
$\wp^j=z^{-2j}+\cdots$, the same descending pole cancellation
shows that every coefficient of $Q$ belongs to $M$.
This supplies the polynomial descent required before applying
$\rho$; it is not an assumption that a complex polynomial
appearing in an elliptic-function argument is already defined
over the coefficient field. The source uses an $M$-linear
retraction for this descent; distinct pole orders give the
shorter argument here.

\paragraph{TA4: CM eliminates the possible quadratic exponential.}
Let $\rho:M\hookrightarrow\mathbb C$ fix $\iota(K)$, with
$g_2,g_3$ and all coefficients of $s_\Lambda$ in $M$.
Let $\Lambda'$ be TA1's conjugate lattice and $s'=s_{\Lambda'}$.
Write $c'$ for its constant in (TA.6). Multiplication of a
function by $e^{uz^2}$ changes its constant from $c'$ to
$c'-2u$; direct differentiation verifies
\[
 W_{c'-2u}(e^{uz^2}s')=e^{2uz^2}W_{c'}(s').                  \tag{TA.8}
\]
Take $u=(c'-\rho(c))/2$. The coefficientwise image under $\rho$
of the original Taylor series and the actual Taylor series of
$e^{uz^2}s'$ have identical initial data and satisfy (TA.3) with
the same $\rho(c),\rho(g_2),\rho(g_3)$. TA2 makes the series
equal. In particular the coefficientwise image has now been
identified with a convergent series by a formal uniqueness
argument, rather than by an illicit continuity assertion.

Clear the denominator in (TA.7), substitute (TA.6), and apply
$\rho$ coefficientwise. TA3 shows that every coefficient being
transported is in $M$. Taylor injectivity for entire functions
then gives the analytic identity
\[
 s'(\alpha z)^2e^{\kappa z^2}
       =s'(z)^{2N}\rho(Q)(\wp_{\Lambda'}(z)),\qquad
             \kappa=2u(\alpha^2-N).                          \tag{TA.9}
\]
The identity is first read off the lattice and then extended
after clearing all poles. At a period $\omega\in\Lambda'$,
the right side vanishes: its possible pole has order at most
$2N-2$, while $s'^{2N}$ has order $2N$. Hence
$\alpha\omega\in\Lambda'$. The squared period law for $s'$
can therefore be used on both sides of (TA.9). Translating by
$\omega$ and cancelling on the nonempty open set where the
factors do not vanish gives
\[
                  e^{\kappa(z+\omega)^2}=e^{\kappa z^2}.
\]
Both sides are entire, so this holds for every $z$. For a
nonzero period, differentiation gives $2\kappa\omega=0$;
thus $\kappa=0$, and $\alpha^2\ne N$ forces $u=0$.
We have proved the exact statement \texttt{galSg}:
\[
 \rho([z^n]s_\Lambda)=[z^n]s_{\Lambda'}\quad(n\ge0),          \tag{TA.10}
\]
under its original assumption that all coefficients on the
left belong to $M$. The next arguments will supply algebraicity;
it was not silently assumed in this step.

\paragraph{TA5: torsion-point transport and the exact CM action.}
Under Weierstrass uniformization the class of $t$ is the origin
if $t\in\Lambda$, and otherwise is
$(\wp_\Lambda(t),\wp'_\Lambda(t)/2)$ on
$y^2=x^3-g_2x/4-g_3/4$. The source relation \texttt{StarPt}
says precisely that the second point is the coordinatewise
$\rho$-image of the first, with both coordinates in $M$ for
a point other than the origin, or that both points are the origin.
The ordinate can be zero, in particular at nonzero $2$-torsion.
This interpretation
proves uniqueness modulo $\Lambda'$, invariance under replacing
$t$ by a period translate, compatibility with addition, and
preservation of the exact annihilating integers. The group law
uses rational chord-and-tangent formulas, including the origin
and vertical-tangent cases, and therefore commutes with field
embeddings. Existence follows by uniformizing the image point.

We also need the action of every $a\in\mathcal O_K$, not just
integer multiplication. The map $z\mapsto\iota(a)z$ on
$\mathbb C/\Lambda$ is an algebraic endomorphism of its cubic.
Indeed its coordinate functions are elliptic meromorphic
functions, hence rational expressions in $\wp,\wp'$; the
identities giving a morphism and its group law are algebraic.
For a complex field automorphism $\sigma$, conjugate these
expressions and their identities. The resulting morphism on
the conjugate cubic is holomorphic; its lift through the
covering $\mathbb C\to\mathbb C/\Lambda'$ is $z\mapsto\beta z$.
The pullback of $dx/(2y)$ shows
$\beta=\sigma(\iota(a))$: this is an algebraic identity of
differentials, so conjugation of its coefficients is legitimate.
If $\sigma$ fixes $\iota(a)$, the transported action is exactly
the same scalar, with no unspecified isogeny or change of scale.

The output coordinates of this CM action must also belong to
$M$. Apply the preceding calculation to every automorphism
fixing all of $M$. It fixes the invariants, the input point,
and $\iota(a)$; hence it fixes the output point. The fixed
field of $\operatorname{Aut}(\mathbb C/M)$ is $M$: an algebraic
element outside $M$ has a distinct conjugate in characteristic
zero, and an element transcendental over $M$ can be shifted
in a transcendence basis; either embedding extends to the
algebraically closed field $\mathbb C$. Thus each nonzero
output coordinate is in $M$. This supplies the field-membership
part of the literal \texttt{StarPt} conclusion as well as its
coordinate identity.

For an arbitrary embedding $\rho:M\to\mathbb C$ it suffices
to extend its restriction on the countable field generated by
the invariants, both the input coordinates and the output
CM coordinates just proved to belong to $M$, and the
specified CM scalars. Such an embedding extends to an
automorphism of $\mathbb C$: extend it first to algebraic
closures, choose transcendence bases of the complements
(both have cardinality $|\mathbb R|$), identify the bases,
and extend through algebraic closures again. We do not assert
extension of an arbitrary embedding on an uncountable $M$.
All quantities whose images matter were included in the
countable subfield, so the preceding calculation proves the
literal CM clause for $\rho$.

If every nonzero $n$-torsion point has its two coordinates in
$M$, use one such extension simultaneously on their finite
set. There are exactly $n^2$ points, identified with
$\Lambda/n\Lambda$; their images are distinct $n$-torsion
points and exhaust the target's $n^2$ points. This proves the
surjectivity clause with its exact extra hypothesis. We have
now accounted for all seven clauses of \texttt{galStar}.
In particular $\Lambda'$ is stable under $\iota(\mathcal O_K)$:
apply the CM clause to the origin, represented by $0$ and any
$\omega'\in\Lambda'$.

\paragraph{TA6: conjugate curve ideals with the required coprimality.}
For later use we supply \texttt{galCurve}, including its
arbitrary prescribed positive integer $n$. Let
$\Lambda'=\mathbb Z\omega_1+\mathbb Z\omega_2$ be the lattice
of TA1, stable under $\mathcal O_K$ by TA5. For a nonreal
integral $a\in K$, write
$\iota(a)\omega_1=r\omega_1+s\omega_2$ with $r,s\in\mathbb Z$.
Here $s\ne0$, since otherwise $\iota(a)=r$. Put
$\delta=\omega_1/s$ and
$I=\{b\in\mathcal O_K:\delta\iota(b)\in\Lambda'\}$.
Stability makes $I$ an ideal, and $s\in I$ makes it nonzero.
Both periods are $\delta$ times integral elements, namely
$s$ and $a-r$, so $\Lambda'=\delta\iota(I)$.

Every nonzero ideal class of a Dedekind domain has a
representative coprime to a given nonzero ideal $J$.
Here is the elementary form used. Choose $a\in I$ whose
valuation at each prime dividing $J$ is exactly that of $I$,
by the finite Chinese remainder theorem in $I/\mathfrak p I$.
Then $(a)=I A$ with $A+J=\mathcal O_K$. Choose $a'\in A$
with the same exact valuations at the primes of $J$ and write
$(a')=A B$, again with $B+J=\mathcal O_K$.
It follows that $(a)B=(a')I$, and
\[
 \Lambda'=\delta\,\iota(a/a')\,\iota(B).                     \tag{TA.11}
\]
Apply this with $J=(2dn)$. Then $B$ has norm prime to $2dn$:
if a rational prime dividing $2dn$ divided its norm, some
prime above it would divide $B$ and $J$, a contradiction.
Its norm is therefore prime both to $2d$, as required of a
curve ideal, and to the prescribed $n$.

The discriminant of $\Lambda'$ is one, since its two invariants
are conjugates of those of $\Lambda$. The scaling identity in
PA1 therefore gives the exact twelfth-root equation for the
new factor $\delta'=\delta\iota(a/a')$. Thus $(B,\delta')$
is a literal curve datum, not merely a homothetic torus.
This proves \texttt{galCurve} without using theta coefficients.

\paragraph{TA7: division quotients are determined over the point field.}
For $u\notin\Lambda$ and $j\ge0$ define the analytic germ
\[
 \Psi_j(z)=\frac{s_\Lambda(j(z+u))}
                 {s_\Lambda(z+u)^{j^2}}.                    \tag{TA.12}
\]
The denominator is nonzero near zero. The sigma addition
formula and its derivative give
\[
 \Psi_0=0,\quad\Psi_1=1,\quad
 \Psi_2=-\wp'_\Lambda(z+u),\quad
 (j+1)^2\Psi_j\Psi_{j+2}
       =\Psi_{j+1}\Psi_{j+1}''-(\Psi_{j+1}')^2.              \tag{TA.13}
\]
For the recurrence, differentiate
$s(v+w)s(v-w)=\operatorname{wr}(s)(v)s(w)^2
-\operatorname{wr}(s)(w)s(v)^2$, where
$\operatorname{wr}(f)=ff''-f'^2$, and substitute
$v=(j+1)(z+u)$ and $w=z+u$. The product rule
$\operatorname{wr}(fg)=\operatorname{wr}(f)g^2+f^2\operatorname{wr}(g)$
and the corresponding power rule cancel the common powers
of $s(z+u)$ and give (TA.13).

The system $p'=q$, $q'=6p^2-g_2/2$, with
$p(0)=\wp(u)$ and $q(0)=\wp'(u)$, determines all derivatives
of $p,q$ recursively. Consequently their Taylor coefficients
belong to the field generated by these three data and commute
with embeddings of that field. Induct using (TA.13), starting
with $\Psi_1,\Psi_2$. At each step $\Psi_j$ for $j>0$ is a
nonzero germ, although its value at zero can vanish: otherwise
$s(j(z+u))$ would vanish on an open set and hence identically.
Division takes place in the Laurent-series field. If the
quotient is known analytically to be holomorphic, its Laurent
series has no negative terms and its Taylor coefficients lie
in the same field. This proves coefficient descent and
conjugacy for every $\Psi_j$, including at torsion points
where dividing numerical values at zero would be invalid.

We will also use a formal root observation. If $F(0)\ne0$,
all coefficients of $F^r$ and $F(0)$ lie in a characteristic-zero
field, and $r>0$, then all coefficients of $F$ lie there.
The coefficient of degree $k$ in $F^r$ is
$rF(0)^{r-1}[z^k]F$ plus a polynomial in earlier coefficients.
The same recursion proves conjugacy of two such roots when
their constant terms are paired. The constant term must be
specified; a power identity alone leaves a root of unity.

\paragraph{TA8: the constant that resolves the torsion root ambiguity.}
Fix $m>0$, $u\notin\Lambda$, and $m^2u\in\Lambda$.
Put $k=m^2-1$ and
\[
 K(z)=(U_us_\Lambda(z))^{4m^2},\qquad K(0)\ne0.
\]
Direct cancellation of the exponential factors in (TA.2)
gives
\begin{align}
 U_{2mu}s_\Lambda(2mz)&=K(z)\Psi_{2m}(z),                    \tag{TA.14}\\
 s_\Lambda(2m^2z)&=K(z)^{m^2}\Psi_{2m^2}(z).                \tag{TA.15}
\end{align}
In the second identity $2m^2u\in2\Lambda$, so (TA.1) makes
its translation operator the identity.
There is also the essential constant identity
\[
                         K(0)=-\frac{\Psi_k(0)}{\Psi_{k+2}(0)}.
                                                                  \tag{TA.16}
\]
Indeed $(k+1)u=m^2u$ is a period. Apply (TA.1) to the
points $u$ and $-u$ with that common period, and use oddness.
The quotient of the resulting values is
$-\exp(-2\pi m^2|u|^2/A)$. The powers of $s(u)$ in the
quotient of (TA.12) contribute $s(u)^{4m^2}$, giving
(TA.16). Its denominator is nonzero because
$(k+2)u\equiv u\pmod\Lambda$. Also $m=1$ is impossible
in this case, so $k>0$ and there is no exceptional zero
numerator. This is the specific root normalization required
by the source, not an arbitrary choice of an algebraic root.

\paragraph{TA9: the literal Galois identity for every rational translate.}
Take precisely the hypotheses of \texttt{gal1}: $\rho$ fixes
$\iota(K)$, the two lattice invariants belong to $M$, $m>0$,
$m(x/2)\in\Lambda$, and
\[
 \mathrm{StarPt}(M,\rho,\Lambda;x/(2m),x'/(2m));              \tag{TA.17}
\]
all coefficients of $G=U_xs_\Lambda$ are assumed in $M$.
Write $u=x/(2m)$, $u'=x'/(2m)$. TA5 gives
$m^2u'\in\Lambda'$. If $u\in\Lambda$, then $u'\in\Lambda'$
and both $x/2,x'/2$ are periods. The translations are the
identity by (TA.1), so TA4 applies directly.

Otherwise TA7 pairs all division quotients for $u,u'$ under
$\rho$. First apply its field-membership version to the
original lattice alone. Equation (TA.14), the assumed
coefficients of $G$, and Laurent division by the nonzero
$\Psi_{2m}$ show that the coefficients of $K$ lie in $M$.
Equation (TA.15) then shows that all coefficients of
$s_\Lambda$ lie in $M$, since rescaling the variable by the
nonzero integer $2m^2$ is invertible coefficientwise.
We can therefore use TA4; we have supplied, rather than
assumed, its additional coefficient-field premise.

Transport (TA.15), using TA4 and TA7, and cancel the nonzero
division-quotient germ. The series of $K^{m^2}$ and
$(K')^{m^2}$ are paired. Formula (TA.16) pairs their constant
terms $K(0),K'(0)$; TA7's root recursion pairs $K,K'$.
Finally (TA.14), followed by inverse integer rescaling,
gives the exact conclusion
\[
 \rho([z^n]U_xs_\Lambda)=[z^n]U_{x'}s_{\Lambda'}\quad(n\ge0).
                                                                  \tag{TA.18}
\]
This proves \texttt{galois\_\allowbreak translate\_\allowbreak sgStar}, with the same
integer $m$ and the literal $2m$-division condition. In
particular no root-of-unity factor has been lost.

\paragraph{TA10: algebraic coefficients for all actual torsion translates.}
PA14 supplies algebraic $g_2,g_3$ and algebraic coordinates
for every torsion point of the normalized CM lattice. Let
$x$ be any rational-lattice point and choose $r>0$ with
$rx\in\Lambda$. For every automorphism
$\sigma\in\operatorname{Aut}(\mathbb C/\overline{\mathbb Q})$
apply TA9 with $M=\mathbb C$, $\rho=\sigma$, $m=2r$, and
$x'=x$. Its hypotheses are all now justified: the coefficient
membership is automatic in $\mathbb C$, $\sigma$ fixes $K$
and the invariants, and the coordinates of $x/(4r)$ are
algebraic torsion coordinates. TA1 identifies the conjugate
lattice with the original lattice. Hence $\sigma$ fixes
every coefficient of $U_xs_\Lambda$.

An element of $\mathbb C$ fixed by all these automorphisms
belongs to $\overline{\mathbb Q}$. To see this directly, a
transcendental element can be included in a transcendence
basis over $\overline{\mathbb Q}$; replacing it by itself
plus one defines an automorphism of the rational-function
field, which extends through its algebraic closure to
$\mathbb C$. It does not fix that element. Thus
\[
 [z^n]U_xs_\Lambda\in\overline{\mathbb Q}
       \quad(x\in\mathbb Q\Lambda,\ n\ge0).                 \tag{TA.19}
\]
The functions are entire by (TA.2). This supplies the full
odd-theta premise retained in PA15. Applying PA15 now gives
algebraic Taylor coefficients for every rational translate
of the normalized even theta function, and the exact
nonzero algebraic lattice-shift multipliers. Its normalization
denominator is nonzero because the chosen half-point is
outside the lattice.

\paragraph{TA11: the distinguished half-point is unique modulo the lattice.}
Put $w=(1+i\sqrt d)/2$. The hypotheses give
$w(1-w)=(d+1)/4=2k$ for an integer $k$.
On $V=(\tfrac12\Lambda)/\Lambda\simeq\mathbb F_2^2$, the
endomorphisms $w$ and $1-w$ are complementary idempotents:
their product is zero and their sum is the identity.
If a class $t$ has both images nonzero, neither idempotent
has rank zero; consequently each has rank one. There is
exactly one element with nonzero components in both
one-dimensional summands. It follows that any two points
satisfying the literal three unkilled clauses differ by
a period. This proves the source's
\texttt{two\_\allowbreak torsion\_\allowbreak unique} without enumerating parity cases.

The normalized even function
$\vartheta_{\Lambda,t}=U_ts_\Lambda/U_ts_\Lambda(0)$ is unchanged
when $t$ is replaced by a period translate: the translation
composition law in PA15 gives the same nonzero scalar in
numerator and denominator. Thus it depends only on this
unique distinguished class.

\paragraph{TA12: even-theta conjugacy, including embeddings not fixing $K$.}
First let $\rho:\mathbb C\hookrightarrow\mathbb C$ fix
$\iota(K)$. TA5 transports the $16$-torsion point $t/8$
to some $u$ modulo $\Lambda'$. Apply TA9 with $x=t$,
$x'=8u$, and $m=4$: its condition $4(t/2)=2t\in\Lambda$
and its $8$-division point are exactly the ones just chosen.
Addition, CM multiplication, and exact-order preservation
in TA5 show that $8u$ satisfies the three unkilled clauses.
By TA6, $\Lambda'$ is an actual $\mathcal O_K$-stable curve
lattice. TA11 makes $8u$ congruent to any specified
unkilled point $t'$ on it. Divide (TA.18) by its degree-zero
identity. We obtain
\[
 \rho([z^n]\vartheta_{\Lambda,t})
                   =[z^n]\vartheta_{\Lambda',t'}.            \tag{TA.20}
\]

For ordinary complex conjugation, the canonical product,
the area identity $A_{\bar\Lambda}=A_\Lambda$, and (TA.2)
give
\[
 s_{\bar\Lambda}(z)=\overline{s_\Lambda(\bar z)},\qquad
 \vartheta_{\bar\Lambda,\bar t}(z)
                 =\overline{\vartheta_{\Lambda,t}(\bar z)}.   \tag{TA.21}
\]
One may alternatively use uniqueness of the normalized odd
theta period laws; the derivative at zero remains one.
Taking Taylor coefficients in (TA.21) is legitimate for
this specific continuous conjugation map. Conjugation
interchanges $w$ and $1-w$, so it preserves the two
nonvanishing clauses of the distinguished half-point.

Every embedding of $\mathbb C$ restricts on the quadratic
field $K$ either to the identity or to conjugation. In the
second case compose it with complex conjugation, apply the
fixed-$K$ result, and use (TA.21). TA1 identifies the new
conjugate lattice as $\bar\Lambda'$, since its two invariants
are the complex conjugates of those of $\Lambda'$. This
proves (TA.20) for every embedding of $\mathbb C$.

Finally let the original $\rho:M\hookrightarrow\mathbb C$
have only the literal coefficient assumptions in
\texttt{use1}. TA10 makes all even-theta coefficients
algebraic, and PA14 makes both invariants algebraic.
Extend the restriction of $\rho$ to
$M\cap\overline{\mathbb Q}$ to an automorphism of
$\mathbb C$, by the countable-field argument of TA5.
It agrees on both invariants and all coefficients to which
the statement applies. TA1 identifies the same
\texttt{lamConj}; therefore (TA.20) descends to the original
embedding. We have proved all three actual \texttt{use1}
clauses: algebraicity of every odd-theta rational translate,
algebraicity of every distinguished even theta coefficient,
and its prescribed conjugacy under the original $M,\rho$
and half-points. The quantifier order and coefficient-field
membership conditions are unchanged.

\paragraph{TA13: the updated analytic boundary.}
PA14, TA10 and PA15 now supply the primitive normalized
one-dimensional theta coefficients and all rational
translations; TA12 supplies the conjugacy conclusion of
\texttt{cmTheta\_\allowbreak taylorCoeff\_\allowbreak algebraic\_\allowbreak and\_\allowbreak conjugate}.
The generalized formal root argument already used in TL18
can therefore invoke these exact suppliers. PA4 also
supplies algebraic complex substitution, and PA10 supplies
the finite algebraic determinant correction $\kappa$.
This is a concrete reduction of the remaining vacuum
problem. It does not construct the actual
\texttt{StdOK}/four-copy vacuum identity, the normalized
indicator/covariance theorem, \texttt{Pins.iX}, or any of
the five later-stage existence assertions. Those remain
separately recorded targets. All proofs here are human
reconstructions with named classical inputs; no Lean
kernel replay or full self-contained RH/GRH proof is
claimed by this section.


\section{Constructing the local indices and their canonical replacement}
\label{sec:canonical-indices}

\status{CI1--CI18 reconstruct the finite local calculations left open in
TL20: finite admitted cosets, lattice normal forms, preservation of first
unfolding nonvanishing, descent of the actual projected character, the
auxiliary-prime one-coset calculation, and the resulting initial and
canonical local-index records. CI19 identifies the exact contribution
to the construction of \texttt{TC.Ok}. Arithmetic host, character, tuple,
and ambient theta data remain hypotheses where stated; no analytic
theta identity or stage of \texttt{Later} is inferred from these local
calculations. These are mathematical reconstructions from the pinned
source, without a Lean build or kernel-verification claim.}

\paragraph{Local conventions.}
Let $P$ be a nonzero prime of the integers of the number field $E$,
$O=\mathcal O_{E,P}$, and $q=|\mathcal O_E/P|$. Except in the finite
coset argument CI3, assume the literal \texttt{DegOneUnram} condition:
$q$ is a rational prime and $q\notin P^2$. Thus $q$ is a uniformizer
of the discrete valuation ring $O$, and $O/qO=\mathbb F_q$. The
coefficient field $K$ has characteristic zero. All local lattices below
are in rational $E$-vector spaces, rather than completions.

At step $j$, put $n_j=k+j(t+1)$ and $c_j=m-1-j$, for
$0\le j\le m-1$. The source rank assumptions include $m\ge4$ even
and $k\ge m+3$. Write
\[
 Y_j=X_j\oplus\bigoplus_{a=0}^{c_j-1}Ee_a,
 \qquad X_j=E^{n_j}.
 \tag{CI.1}
\]
The group $R_j$ preserves $X_j+Ee_0+\cdots+Ee_{a-1}$,
fixes $e_0$, and sends $e_a$ to $e_a$ plus an earlier vector.
Its action on $X_j$ may be any invertible linear map; calling
all of $R_j$ unipotent would be incorrect. Its kernel on $X_j$ is
$N_j$. Set
\[
 a(r)=\sum_{a=0}^{c_j-2}(r e_{a+1})_{e_a}.
 \tag{CI.2}
\]
Then $a(rs)=a(r)+a(s)$, as follows by multiplying these triangular
matrices. At a twist prime the basic lattice is $O^{n_j+c_j}$;
at an auxiliary, or type, prime it is
$\Lambda_j=O^{n_j}\oplus\bigoplus_a q^{a+1}Oe_a$.
The words ``nonzero unfolding'' always mean the source predicate
\texttt{FiniteSumNonzeroLargeDepth}: all sufficiently deep box sums
are nonzero as functions of the retained matrix variable. It is not
an assertion about an unproved analytic integral.

\paragraph{CI1: finite quotients, stable box sums, and affine transport.}
For every $s\ge1$, the integers $0,\ldots,q^s-1$ represent $O/q^sO$
exactly once. Indeed they represent $O/qO$; after removing that first
digit and dividing by $q$, induction gives existence. Uniqueness
follows because the $P$-valuation of a nonzero rational integer is
its usual $q$-adic valuation. Consequently
\[
 \mathcal B_{A,B}=\{0,1,\ldots,q^{A+B}-1\}/q^A
 \tag{CI.3}
\]
represents $q^{-A}O/q^BO$.

Suppose $f:E^d\to K$ is supported in $q^{-A_0}O^d$ and is invariant
under $q^{B_0}O^d$. For $A\ge A_0$, $B\ge B_0$, define
\[
 I(f)=q^{-Bd}\sum_{u\in\mathcal B_{A,B}^d}f(u).
 \tag{CI.4}
\]
Increasing $A$ only adds zero support cosets; increasing $B$ by $b$
splits each old period coset into $q^{bd}$ equal terms. The factor
$q^{-Bd}$ cancels this multiplicity. Thus (CI.4) is independent of
both sufficiently large depths. The same proof applies uniformly
to a family with common support and period bounds.

For an invertible $E$-linear map $T$ and $b\in E^d$, the functions
$f$ and $f(Tu+b)$ have bounded support and a full period lattice.
Their stable sums vanish simultaneously. To see this without a
measure theorem, choose an $O$-lattice $L$ which is a period lattice
for both functions. The change of variable identifies the sum over
$L$-cosets for $f(Tu+b)$ with the sum over $TL$-cosets for $f$.
Use the common sublattice $H=L\cap TL$, which has finite index in
each, to compare either sum with the sum over $H$-cosets. The two
comparisons multiply by the positive integers $[L:H]$ and $[TL:H]$.
Every box normalization is also a nonzero rational number.
Characteristic zero therefore proves the asserted equivalence.
Bounds on the entries of $T,T^{-1},b$ choose the needed box depths;
there is no unstated convergence argument.

\paragraph{CI2: a nonzero finite sum forces an affine eigenvalue to be one.}
Here an even more useful conclusion holds. Let $V$ be an additive
group, $L\subset V$ a subgroup, and $f:V\to K$ be $L$-periodic,
with only finitely many nonzero $L$-cosets and nonzero sum over them.
Suppose $T$ is an additive automorphism and
\[
 f(Tv+b)=\chi f(v)\quad(v\in V).
 \tag{CI.5}
\]
Then $\chi=1$, even if $T$ does not preserve the initially chosen
period subgroup $L$.

First $f\ne0$ and the surjectivity of $T+b$ give $\chi\ne0$.
Let $M=\{u:f(v+u)=f(v)\ \text{for all }v\}$ be the full period
group. Equation (CI.5) and its inverse imply $T(M)=M$. Fix
$v_0$ with $f(v_0)\ne0$. Distinct classes in $M/L$ give distinct
nonzero support classes $v_0+u+L$, so $[M:L]$ is finite.
Each nonzero $M$-coset contributes exactly $[M:L]$ equal terms to
the sum over $L$-cosets. Hence the sum over $M$-cosets is nonzero.
But $T+b$ permutes these cosets and their finite nonzero support.
Summing (CI.5) gives $S=\chi S$ with $S\ne0$, proving the claim.
This is the actual finite-sum rigidity used in the projected-kernel
argument; root-of-unity containment alone would not prove it.

\paragraph{CI3: finitely many admitted local level cosets.}
Let $\Lambda=\bigoplus_l Oq^{e_l}e_l$, and call $h$ admitted at
depth $j$ if
\[
 h^{-1}\Lambda\subseteq\Lambda,
 \qquad q^j h\Lambda\subseteq\Lambda.
 \tag{CI.6}
\]
Color $h$ by the matrix of $q^jh$ in this lattice basis modulo
$q^{j+1}$. There are finitely many colors. More literally, one may
lift every local integral entry to $\mathcal O_E$ modulo
$q^{j+1}O$ and use its class in the finite global quotient
$\mathcal O_E/(q^{j+1})$, exactly as the source does. Such lifts
exist because a denominator outside $P$ is invertible modulo a
sufficient power of $P$; alternatively localization of this finite
quotient is a quotient of it.

Equal colors give $(h'-h)\Lambda\subseteq q\Lambda$.
Applying $h^{-1}$ and using (CI.6) gives
$(h^{-1}h'-1)\Lambda\subseteq q\Lambda$; interchanging $h,h'$
gives the same assertion for its inverse. Thus $h^{-1}h'$ is an
integral basis change congruent to one modulo $q$.

Every local group $J_j$ in this construction contains that
principal congruence subgroup. At a twist prime it preserves the
two lattices defining $J_Q$; at a type prime it preserves the
inverse images of the residue flags and fixes their prescribed
residue basis vectors. These are direct consequences of being
the identity on $\Lambda/q\Lambda$. Hence equal colors imply
$h^{-1}h'\in J_j$, proving finiteness of admitted right cosets and
therefore of admitted $R_j\backslash\mathrm{GL}(Y_j)/J_j$
classes. The source's type-prime admission condition requires
stabilizing $\Lambda_j$, and therefore implies (CI.6) automatically.
This proves the actual
\texttt{exists\_\allowbreak finite\_\allowbreak index\_\allowbreak of\_\allowbreak adm\_\allowbreak cosets} supplier.

\paragraph{CI4: a rational normal form with integral projected group.}
At a twist prime let $M=hO^{n+c}$ with
$O^{n+c}\subseteq M\subseteq q^{-j}O^{n+c}$. Put
$W_a=X+Ee_0+\cdots+Ee_a$. The image of $M\cap W_a$ under its
$e_a$-coordinate is a nonzero fractional ideal of $O$, hence
$q^{\nu_a}O$. Choose $v_a\in M\cap W_a$ with that coordinate
$q^{\nu_a}$, and put $u_a=q^{-\nu_a}v_a$.
Successively eliminating the last coordinate proves
\[
 M=(M\cap X)\oplus\bigoplus_{a=0}^{c-1}Oq^{\nu_a}u_a,
 \qquad -j\le\nu_a\le0.
 \tag{CI.7}
\]
The inequalities follow by comparing the same coordinate image
with $O$ and $q^{-j}O$. Choose a unitriangular map which fixes
$X,e_0$ and carries $u_a$ to $e_a$ for $a>0$; it exists by
successive elimination. Write $u_0=e_0+z_0$. Since
$\nu_0\le0$, $u_0\in M$, and since $e_0\in M$, one has
$z_0\in M\cap X$. Now choose $g\in\mathrm{GL}(X)$ carrying
the free full $O$-lattice $M\cap X$ onto $O^n$, and extend it
by the identity on the tail. The resulting $\rho\in R_j$ gives
\[
 \rho M=O^n\oplus Oq^{\nu_0}(e_0+z)
          \oplus\bigoplus_{a=1}^{c-1}Oq^{\nu_a}e_a,
 \qquad z\in O^n.
 \tag{CI.8}
\]
If $c=0$, only the first summand occurs. Formula (CI.8) again lies
between $O^{n+c}$ and $q^{-j}O^{n+c}$, so the normalized matrix
is still admitted. It is in the same double coset, with right
factor equal to one.

The intersection of (CI.8) with $X$ is exactly $O^n$.
If $r\in R_j\cap h'J_jh'^{-1}$, then $r$ preserves $h'O^{n+c}$
and $X$, and its inverse does too. Its restriction to $X$ therefore
preserves $O^n$ in both directions. This proves the actual
$\operatorname{pr}_X C_{h'}\subseteq\mathrm{GL}_n(O)$ clause,
as well as \texttt{TwNF}, \texttt{TwBdd}, and \texttt{TwNormRep}.
No global free basis for an $\mathcal O_E$-lattice has been assumed.

\paragraph{CI5: the unfolding fibers are affine spaces, with the stated phase.}
For a step with source tail $f_0,\ldots,f_c$ and target tail
$e_0,\ldots,e_{c-1}$, the matrix $x(t)=\texttt{xx1}(x,t)$ is
described by its columns:
\[
 \begin{array}{c|c}
 \text{source column}&\text{image}\\ \hline
 X'&x(X')\subseteq X\\
 f_0&e_0\\
 f_a\ (1\le a<c)&e_a+\text{an arbitrary vector in }
             X+Ee_0+\cdots+Ee_{a-1}\\
 f_c&\text{an arbitrary vector in }Y.
 \end{array}
 \tag{CI.9}
\]
The fixed row $f_0\mapsto e_0$ applies only when $c>0$.
When $c=0$, the sole source tail column is $f_0=f_c$ and is
the free last column; the target has no tail. The entries
called arbitrary are the coordinates of $t$.
For $r\in R_j$, with restriction $g$ to $X$, column substitution
gives an invertible affine transformation $T_rt+b_r$ such that
\[
 r x(t)=\texttt{xx1}(gx,T_rt+b_r),\qquad
 \operatorname{chArg}(r x(t))=
       \operatorname{chArg}(x(t))+a(r).
 \tag{CI.10}
\]
The inverse comes from $r^{-1}$. For the phase identity, the
$e_a$-coordinate of the image of $f_{a+1}$ acquires precisely
$(r e_{a+1})_{e_a}$ for $a<c-1$; the last column acquires no
constant. All lower terms have zero $e_a$-coordinate.
Right composition by a source $r'$ has the corresponding formula
with retained variable $xg'$ and phase shift $a(r')$. Indeed
each new source column is itself plus earlier columns, so only
the immediately preceding tail column contributes to the phase.
These calculations prove both left and right affine transports,
including the invertibility and surjectivity on the retained
matrix variable.

\paragraph{CI6: twist-prime test functions, covariance, and uniform bounds.}
Reorder coordinates as the first $k$ coordinates of $X$, then
the $c$ tail coordinates, then the remaining coordinates of $X$.
In these coordinates $J_Q$ consists of the matrices in
$\mathrm{GL}(O)$ whose reduction preserves the first $r=k+j+c$
coordinates. Thus its lower-left residue block is zero.
The actual $\Phi_Q$ is supported on integral rectangular matrices
with that lower-left block zero modulo $q$, and on this support is
\[
 \Phi_Q(y)=\xi\bigl(\det\overline{y}_{[0,r),[0,r)}\bigr),
 \tag{CI.11}
\]
where a multiplicative character is extended by zero at zero.
An element $\kappa\in J_Q$ and its inverse preserve this support.
Block multiplication gives
$\overline{(\kappa^{-1}y)}_{11}
=\overline{\kappa^{-1}}_{11}\bar y_{11}$ and hence
\[
 \Phi_Q(\kappa^{-1}y)=
   \xi(\det\overline{\kappa^{-1}}_{11})\Phi_Q(y).
 \tag{CI.12}
\]
The scalar is nonzero, since the leading block and its inverse
are inverse modulo $q$. Right multiplication has the same
determinant calculation. Factoring the full determinant into
the leading and complementary diagonal blocks identifies these
scalars with the exact \texttt{lamRow} and \texttt{muAt} factors
used by the source. In particular \texttt{lamRow} is a unit-valued
character of $J_Q$, and is one on its principal congruence subgroup.

Write the translated datum as
$G_{a,h}(y)=s(a,h)\Phi_Q(h^{-1}ya)$. Its scalar $s(a,h)$ is the
displayed product of local determinant characters and square-root
powers in \texttt{Loc.om}. It is nonzero when the prescribed
uniformizer character value and square root are nonzero:
every nonzero element is uniquely $q^v u$ with $u\in O^\times$,
so its local character is a nonzero power times $\xi(\bar u)$.
If $a$ is admitted at step $i$ and $h$ at step $i+1$, then
\[
 G_{a,h}(y)\ne0\Longrightarrow q^{i+1}y\text{ is integral},
 \qquad G_{a,h}(y+q^{i+1}z)=G_{a,h}(y)
 \tag{CI.13}
\]
for integral $z$. For support, use $y=h(h^{-1}ya)a^{-1}$.
For periodicity, $h^{-1}q^{i+1}za$ is divisible by $q$, because
$h^{-1}$ is integral and $q^ia$ is integral. This preserves every
residue entry in (CI.11), including the zero-support cases.
The phase $\operatorname{chArg}$ is a sum of entries, so its value
on this support lies in $q^{-(i+1)}O$. These are uniform bounds,
not assumptions about the chosen matrix $y$.

\paragraph{CI7: actual twisted double-coset nonvanishing.}
Suppose $h'=r h\kappa$ with $r\in R_j$, $\kappa\in J_Q$, and
$h,h'$ admitted. By (CI.12) and the determinant factors,
\[
 G_{a,h'}(y)=C\,G_{a,h}(r^{-1}y),\qquad C\ne0.
 \tag{CI.14}
\]
Assume the depth-$D$ character $\psi$ is multiplicative under
addition on $q^{-D}O$ and is one on $O$. It is nonzero throughout
$q^{-D}O$, since $\psi(u)\psi(-u)=1$. Take $D\ge j$.
If the translated test function has empty support, both
unfoldings vanish. Otherwise choose a support point.
Where either side of (CI.14) is nonzero, CI6 bounds both phase
arguments in $q^{-D}O$. Their difference, the constant in
(CI.10), is therefore in that group too. Consequently the two
unfolding integrands are related by an invertible affine change
of $t$ and a nonzero scalar, including their zero-support terms.
CI1 proves that their stable sums vanish simultaneously. The
invertible retained-variable change $x\mapsto g^{-1}x$ then gives
\[
 \texttt{FiniteSumNonzeroLargeDepth}(a,h)
 \quad\Longleftrightarrow\quad
 \texttt{FiniteSumNonzeroLargeDepth}(a,h').
 \tag{CI.15}
\]
The same reasoning gives source-side transport whenever the
source covariance law and corresponding support bounds hold.
It uses only finite sums over commensurable lattices. This proves
the twist-prime content of \texttt{TwistedEisensteinNonvanishing}
and \texttt{unfoldingNonvanishingDoubleCoset\_\allowbreak holds}, including
the actual normalized representative from CI4. The source's
analytic-sounding name does not add an analytic hypothesis here.

\paragraph{CI8: the actual projected-kernel character is trivial.}
Put $C_h=R_j\cap hJ_jh^{-1}$ and
\[
 \chi_h(r)=\psi(a(r))^{-1}\lambda_j(h^{-1}rh).
 \tag{CI.16}
\]
At a twist prime, $q^j r$ is integral for $r\in C_h$: use
$r=h\kappa h^{-1}$ and admission of $h$. Thus $a(r)\in q^{-j}O$.
For $D\ge j$, (CI.2), the character law of CI6, and additivity
of $\psi$ prove that (CI.16) is a nonzero multiplicative character.

Suppose a nonzero unfolding enters $h$ from some admitted $a$.
If $r\in C_h$ acts identically on $X_j$, then it is triangular
with every diagonal entry one, so $\det r=1$. The target
covariance formula therefore gives
$G_{a,h}(r y)=\lambda_j(h^{-1}rh)G_{a,h}(y)$ without an extra
determinant factor. Fix a retained $x$ at which the stable
unfolding sum is nonzero. By CI5 its integrand satisfies
\[
 f_x(T_rt+b_r)=\chi_h(r)f_x(t).
 \tag{CI.17}
\]
The support and period bounds of CI6 make it a finite sum on a
quotient of the free-coordinate lattice. CI2 proves
$\chi_h(r)=1$. Thus \texttt{ChiTriv} is proved for the actual
kernel, not merely for an assumed quotient character.

Define $\bar\chi_h(g)=\chi_h(r)$ using any lift of
$g\in\operatorname{pr}_X C_h$. Two lifts differ by the kernel,
so this is well defined and multiplicative. CI4 places the group
inside $\mathrm{GL}_n(O)$. This is the precise \texttt{TwChrOk}
and \texttt{characterNonvanishing\_\allowbreak holds} conclusion. At the last
index $c_j=0$ the projection has trivial kernel. At index zero
$h=1$, the twist character is one: $a(r)\in O$ and the
complementary determinant block in \texttt{lamQB} is empty,
since $n_0=k$.

\paragraph{CI9: recursive finite index sets with exact quantifiers.}
Fix one local prime. Start with $F_0=\{1\}$. Given finite $F_i$,
consider admitted matrices $h$ at step $i+1$ for which
\emph{some} $a\in F_i$ has nonzero unfolding into $h$. CI3 gives
finitely many possible double cosets. At a twist prime choose
the normalized representative of CI4 in each retained class;
CI7 preserves its nonzero predecessor. At a type prime retain
any representative in the class. Choose exactly one per class.
This constructs $F_{i+1}$, and proves, separately, its predecessor
property, uniqueness modulo double cosets, and coverage of every
admitted matrix having a nonzero predecessor. The construction
is for $i<m-2$. Set $F_{m-1}=\{1\}$ and $F_j=\varnothing$ for
$j>m-1$, exactly as in the source.

At the finitely many local primes take
$\beta_j=\prod_P F_{P,j}$. For a prescribed component in a
nonempty $\beta_{i+1}$, predecessor choices at all primes form a
member of $\beta_i$; hence the source's predecessor quantifier
is satisfied. Similarly choose the covering representative
componentwise. This proves the actual \texttt{exists\_\allowbreak step} and
\texttt{exists\_\allowbreak sets} construction. Intermediate index sets are
not asserted nonempty merely because the initial set is nonempty.
Their exact nonemptiness requirements are handled by the
nonvanishing arguments, not by changing an empty product silently.

\paragraph{CI10: the actual symmetric residue pattern and its characters.}
Put $s=m/2-1$, $r=k+m-1$, and use the reordered coordinates of
CI6. For a square matrix of size $N=r+jt$, the allowed off-diagonal
positions $(p,p')$ are precisely
\[
\begin{array}{c|l}
 \text{column }p'&\text{allowed row }p\\ \hline
 p'<k&p<p'\ \text{or}\ k\le p<k+s\\
 k\le p'<r&k\le p<p'\ \text{or}\ (p'>k+s\ \text{and}\ p<k)\\
 p'\ge r&p<r\ \text{or}\
       \lfloor(p-r)/t\rfloor\le\lfloor(p'-r)/t\rfloor.
\end{array}
\tag{CI.18}
\]
Only indices within the matrix are used. When $t=0$ the last row
of the table has no columns. Let $\mathcal A_N$ be the matrices
whose off-diagonal support satisfies (CI.18). The relation
``equal or allowed'' is transitive: for a column in the first
two intervals, its allowed predecessors stay in those intervals
and a second predecessor satisfies the same inequality; for a
column in the last interval, block numbers can only decrease.
The one missing corner, from column $k+s$ to a row $p<k$, cannot
arise through an intermediate index, since that column has only
earlier tail predecessors. Thus $\mathcal A_N$ is an algebra.

On the first $r$ indices the strict allowed relation has no
cycles. For example it strictly increases the ordering function
which first lists $k,\ldots,k+s-1$, then $0,\ldots,k-1$, then
$k+s,\ldots,r-1$. Hence
$(AB)_{pp}=A_{pp}B_{pp}$ for $p<r$ in this algebra. An invertible
$A\in\mathcal A_N$ has inverse in the algebra: multiplication
by $A$ is injective on the finite-dimensional algebra, hence
surjective, so its inverse is the element mapping to one.
It follows that
\[
 J_N^{\rm res}=\{A\in\mathcal A_N^\times:A_{pp}=1\ (p<r)\}
 \tag{CI.19}
\]
is a group. Conjugating the rational matrices by the diagonal
basis of $\Lambda_j$ and reducing modulo $q$, the defining
\texttt{JGsym} flag and residue-vector clauses are exactly
(CI.19). In fact each flag is the inverse image of its span of
earlier reordered basis vectors. Its exceptional boundary
\texttt{bd5} imposes precisely the missing corner just described.
This proves both directions of the entry description.

Define
\[
 \ell(A)=\sum_{a=0}^{k-2}A_{a,a+1}
          +\sum_{a=k}^{r-2}A_{a,a+1},\qquad
 \lambda_j(A)=\psi_0(\ell(A))
       \prod_{b=0}^{j-1}\chi_b(\det A_{[r+bt,r+(b+1)t)}).
 \tag{CI.20}
\]
For each displayed adjacent pair, the matrix product has only
the two possible intermediate indices equal to its endpoints:
the table rules out all others. The omitted pair $(k-1,k)$ is
zero throughout the pattern. Thus $\ell(AB)=\ell(A)+\ell(B)$.
Each diagonal $t$-block multiplies as a block, since all
intermediate indices outside it contribute zero. Its determinant
is nonzero for a group member, by the same statement for its
inverse. Therefore (CI.20) is a unit-valued character and is one
on a matrix reducing to one. This reconstructs the actual
symmetric \texttt{lamRow} character law.

The rectangular test function has $a=r+(j-1)t$ columns and
$a+t$ rows. On its pattern support it is
\[
 \Phi_S(Y)=\psi_0\!\left(\sum_{b=0}^{r-2}Y_{b,b+1}\right)
 \frac{\prod_{b=0}^{j-2}\chi_b(\det Y_{[r+bt,r+(b+1)t)})}
      {\chi_{j-1}(\det Y_{[0,a)})},
 \tag{CI.21}
\]
with the same zero extension when a determinant vanishes.
The pattern is stable under left and right multiplication by
the corresponding groups (CI.19). To verify its covariance,
multiply blocks: each old diagonal block determinant gains the
corresponding determinant of the acting matrix, the leading
$a$-determinant gains its leading determinant, and the adjacent
entry sum gains $\ell$ of that matrix. For target multiplication
the full determinant is the leading $a$-determinant times the
new last $t$-block determinant. These identities give, with the
source's exact local determinant character $\mu$,
\[
 \Phi(\kappa y)\mu(\det\kappa)
    =\lambda_j(\kappa)\Phi(y),\qquad
 \Phi(y\kappa')\mu(\det\kappa')
    =\lambda_{j-1}(\kappa')\Phi(y).
 \tag{CI.22}
\]
The identities also hold off support, because the inverse
matrices preserve support. Reduction of the scaled rational
entries gives the literal \texttt{PhiGS} source and target
laws. This supplies the type-prime covariance needed below,
rather than assuming it as part of a later column-law record.

\paragraph{CI11: a nonzero residue unfolding has the identity double coset.}
In CI11--CI13 assume $j+1<m$, hence $c=c_j\ge1$; this is
the source's one-coset range. The zero-tail projected kernel
is handled separately in CI8.
Let $h\in\mathrm{GL}_{a+t}(\mathbb F_q)$. If the finite
unfolding of $\Phi_S(h^{-1}\texttt{xx1}(x,t))$ is nonzero,
some summand has $Y=h^{-1}\texttt{xx1}(x,t)$ on nonzero support.
Its leading $a\times a$ block is invertible by (CI.21).
Complete the $a$ columns of $Y$ by the last $t$ standard columns.
The resulting square matrix $\kappa$ is invertible, has determinant
that leading determinant, and belongs to $J_{a+t}^{\rm res}$
by (CI.18).

From (CI.9) choose $\nu\in N_j$ whose tail columns are the first
$c$ tail columns of $\texttt{xx1}(x,t)$; its triangular diagonal
is one. Then $g=\nu^{-1}h\kappa$ fixes every tail basis vector.
Moreover it sends into $X$ each reordered source column outside
the tail interval and its one free last column. Let $P_X,P_A$
be the projections onto $X$ and the tail, and put
$B=P_AgP_X$. Since $B$ kills the tail and takes values in it,
$B^2=0$. The matrix $\kappa'=1+B$ is invertible with inverse
$1-B$. The only possible nonzero entries of $B$ have a tail row
and a column at or after the omitted free column; all these
positions are allowed by (CI.18). Thus $\kappa'\in J_{a+t}^{\rm res}$.
For $\gamma=g(1-B)$ the tail component of $\gamma v$ equals that
of $v$, while $\gamma$ fixes the tail and maps $X$ into $X$.
Consequently $\gamma\in R_j$. We obtain
\[
 h=\nu\gamma\kappa'\kappa^{-1}\in R_jJ_{a+t}^{\rm res}.
 \tag{CI.23}
\]
This proves the support part of \texttt{OneCosetSharp} by matrix
completion and elimination; cancellation in a Fourier sum has
not been assumed to identify its support.

\paragraph{CI12: an explicit nonzero residue unfolding at the identity.}
Choose $x_0:E^{n'}\to E^n$ on coordinate vectors by
$x_0e_b=e_b$ for $b<k$ and $x_0e_b=e_{b+1}$ for $b\ge k$.
Use this same matrix over the residue field. For any free
coordinates $t$, the reordered columns of
$Y=\texttt{xx1}(x_0,t)$ outside $k<p'\le k+c$ are standard
columns. If $Y$ satisfies the support conditions of (CI.21), its
leading block is upper triangular with diagonal one, and its
old $t$-blocks are identity matrices. Indeed in the remaining
columns (CI.18) permits only rows strictly earlier than the
column; all other columns were already standard. Hence all
multiplicative-character factors in (CI.21) are one.
The only possibly nonzero terms of its adjacent-entry sum are
the $c$ terms of $\operatorname{chArg}(Y)$. Therefore
\[
 \Phi_S(\texttt{xx1}(x_0,t))
       \psi_0(-\operatorname{chArg}(\texttt{xx1}(x_0,t)))
   =\begin{cases}1&\text{if the support pattern holds},\\
                  0&\text{otherwise}.
     \end{cases}
 \tag{CI.24}
\]
There is at least one supporting choice: all free entries zero
except the entry sending the last source tail vector $f_c$ to
the target vector $e_k$ in $X$, which is one. In reordered
coordinates this is the rectangular identity matrix. The finite
unfolding at $x_0$ is thus the positive integer counting these
supporting choices. It is nonzero in characteristic zero.
This proves the nonvanishing half of \texttt{OneCosetSharp};
even a nontriviality hypothesis on $\psi_0$ is unnecessary for
this particular finite-matrix calculation.

\paragraph{CI13: reduction, lifting, and the actual type-prime one-coset theorem.}
Let $h$ stabilize $\Lambda_j$. In its scaled basis, both $h$ and
$h^{-1}$ are integral and reduce to inverse residue matrices.
Its determinant is a unit, so the scalar in the translated
$\Phi_{GS}$ is nonzero. The scaled entries of an unfolding
matrix are again of the form (CI.9). More precisely, for a free
entry in row $l$ and source column $b$, put
\[
 d_f=(b-n'+1)-\begin{cases}0&l<n,\\l-n+1&l\ge n.
                       \end{cases}
 \quad 1\le d_f\le c+1.
 \tag{CI.25}
\]
Its scaled entry is $q^{d_f}t_f$. The retained $x$ entries
remain unscaled. If one is nonintegral the unfolding is zero,
since every nonzero test-function value requires integral
scaled entries.

For integral $x$, reduce it and every $q^{d_f}t_f$ modulo $q$.
In a box (CI.3) with $A\ge c+1$, $B\ge1$, every residue value
has exactly $q^{B+d_f-1}$ admissible representatives in the
$f$th free coordinate. This follows by writing its numerator
as $q^{A-d_f}(u+qv)$, where $u$ is the chosen residue digit.
The phase satisfies
$\overline{q\operatorname{chArg}(\texttt{xx1}(x,t))}
=\operatorname{chArg}(\texttt{xx1}(\bar x,\bar t))$.
The literal type-prime tie
$\psi(a/q)=\psi_0(\bar a)$ consequently gives the exact identity
\[
 \sigma(E^{A,B}_{1,h}(x))
   =\sigma(s(1,h))\,
        q^{\sum_f(d_f-1)} E^{\rm res}_{1,\bar h}(\bar x).
 \tag{CI.26}
\]
Here the residue function is (CI.21), and $\sigma:K\hookrightarrow
\mathbb C$ is the specified embedding. All factors in front are
nonzero. Every residue $x$ lifts entrywise to an integral rational
$x$. Thus eventual nonzero unfolding is equivalent to nonzero
residue unfolding, in both directions and with the stated
uniform depths.

Apply CI11 to $\bar h=\bar\rho\bar\kappa$. Lift the entries of
$\bar\rho$ to $O$, preserving its zero entries and its entries
equal to one. Its leading $X$ block has unit determinant; its
triangular tail has diagonal one. Hence the lift and its inverse
are integral. Conjugating back by the scaled lattice basis gives
$\rho\in R_j\cap\operatorname{Stab}(\Lambda_j)$ with that
reduction. The matrix $\rho^{-1}h$ has residue $\bar\kappa$ and
therefore belongs to $J_j$ by CI10. Consequently
\[
 E_{1,h}\ne0\ \Longrightarrow\ h\in R_jJ_j,
 \qquad E_{1,1}\ne0.
 \tag{CI.27}
\]
The second assertion follows from CI12 and (CI.26).

For completeness these statements transport in both source and
target double cosets, with the stabilizer hypotheses present
in the source. All source and target matrices then preserve the
scaled lattices, so their moved Hom lattice is the same; it has
bounded coordinates and a full period lattice. In that Hom
lattice $q\operatorname{chArg}$ is integral, because successive
tail exponents differ by one. Formula (CI.22), CI5, and the finite
affine change of CI1 prove the two-sided transport exactly as in
CI7, now with depth one. In particular, if the source matrix
$a$ is in the identity double coset and both matrices stabilize
their prescribed lattices, then
\[
 E_{a,h}\ne0\quad\Longleftrightarrow\quad h\in R_jJ_j.
 \tag{CI.28}
\]
This is the actual auxiliary-prime input to \texttt{twSurjG},
\texttt{twNonVanG}, and the reference-chart index matching.
It has been proved independently of \texttt{TC.Laws}.

\paragraph{CI14: the projected symmetric levels and their exact characters.}
At the identity representative, restricting $r\in R_j\cap J_j$
to $X_j$ gives an element of $\mathrm{GL}_{n_j}(O)$, since
$\Lambda_j\cap X_j=O^{n_j}$. Its reduction $A$ has diagonal one
for $p<k+j$; its allowed off-diagonal positions are
\[
\begin{array}{c|l}
 p'<k&p<p'\ \text{or}\ (k\le p\ \text{and}\ p+s+2\le k+j)\\
 k\le p'<k+j&k\le p<p'\ \text{or}\
                     (p<k\ \text{and}\ k+j<p'+s+1)\\
 p'\ge k+j&p<k+j\ \text{or}\
       \lfloor(p-(k+j))/t\rfloor\le
                 \lfloor(p'-(k+j))/t\rfloor.
\end{array}
\tag{CI.29}
\]
Indeed deleting the $c_j$ tail positions from (CI.18) replaces
each remaining index $p\ge k$ by $p+c_j$, and
$c_j+j=2s+1$. Substitution in the three inequalities gives
exactly (CI.29). Conversely lift a residue matrix satisfying
(CI.29) entrywise, with determinant a local unit, and extend it
by the identity on the tail. The same substitutions show that
this extension belongs to $R_j\cap J_j$. This proves the exact
projected group, including surjectivity onto the residue group.

For such an extension of $g$ with residue $A$, split the
superdiagonal sum (CI.20) into the initial $X$ part, the tail
part, the tail-to-$X$ boundary, and the remaining $X$ part.
The boundary entry is zero because $rX\subseteq X$; the tail
part is $q a(r)$ in scaled coordinates. The old block determinants
are precisely the $X$ block determinants. Using
$\psi(a(r))=\psi_0(\overline{q a(r)})$, (CI.16) therefore becomes
\[
 \bar\chi_j(A)=
 \psi_0\!\left(\sum_{b=0}^{k-2}A_{b,b+1}
             +\sum_{b=0}^{j-2}A_{k+b,k+b+1}\right)
 \prod_{b=0}^{j-1}
       \chi_b(\det A_{[k+j+bt,k+j+(b+1)t)}).
 \tag{CI.30}
\]
The two copies of the tail phase cancel, with the inverse
on $\psi$ exactly as in (CI.16). Formulas (CI.29)--(CI.30) are
the source's \texttt{colLevSym} group and character; at $j=0$
they give \texttt{colLev Ch 0}. Thus the actual
\texttt{ocbGrpGLoc}, \texttt{ocbChrGLoc}, and
\texttt{closedLevels\_\allowbreak by\_\allowbreak entries} assertions are proved.

The same calculation also makes double-coset character
transport explicit. If $h'=\rho h\kappa$, then
$C_{h'}=\rho C_h\rho^{-1}$. Its projected group is conjugated
by $g=\rho|_X$. By (CI.2), $a(\rho r\rho^{-1})=a(r)$, while
$\lambda(\kappa^{-1}(h^{-1}rh)\kappa)=\lambda(h^{-1}rh)$.
Hence the projected characters agree under that conjugation.
At type primes the matrices involved preserve $\Lambda_j$, so
$g\in\mathrm{GL}_{n_j}(O)$ and this is also the actual residue
level conjugation used in reference-chart matching.

\paragraph{CI15: constructing the full initial local-index record.}
Assume the literal inputs of \texttt{exists\_\allowbreak chIdx\_\allowbreak of\_\allowbreak cst}:
every row prime is degree one and unramified; twist and type
prime sets are disjoint; $\texttt{sharp}=\texttt{true}$;
\texttt{LocCst.Ok}, $m\le M\le A,B$, the specified unit $\xi_0$
and Weil tie, the exact local determinant-character values and
their type-prime residue tie, and
\[
 D_P=1\quad(P\text{ type}),\qquad
 D_P\ge2M,\quad D_P\ge2(m-1)+1\quad(P\text{ twist}).
 \tag{CI.31}
\]
These are explicit arithmetic and constant hypotheses, rather
than an already supplied index record. CI9 constructs its finite
sets and representatives. At type primes all chosen representatives
are in the identity double coset: induct from $F_0=\{1\}$ using
CI13. That local set is nonempty at every constructed step,
because transport of $E_{1,1}\ne0$ supplies an eligible next
representative. CI14 therefore constructs its residue level and
character by conjugating the identity level. At twist primes use
CI8's projected group and descended character. At zero choose
the literal top twist level and \texttt{colLev Ch 0}; at the last
step the zero tail makes kernel descent immediate.

Here are the remaining bounds required by the record. They are
included to avoid concealing them in a record-assembly assertion.
For type primes write $\mathcal H_i=\operatorname{Hom}_O
(\Lambda_i,\Lambda_{i+1})$. The datum is supported in
$\mathcal H_i$ and invariant under $q\mathcal H_i$, by its
integral-entry and residue definition. Admitted source and target
matrices preserve this same Hom lattice. Its coordinate exponents
are differences of numbers between $0$ and $m-1$; hence its
support is contained in $q^{-m}$ times the ordinary integral Hom
lattice, and its periods contain $q^m$ times that lattice.
Also $q\operatorname{chArg}$ is integral on its support.
At twist primes CI6 gives the sharper bound $i+1\le m$.
Multiplication by the phase preserves these uniform period
bounds, because the change of phase argument is integral.
CI1 now proves the literal \texttt{deep} equality for all
$A'\ge A$, $B'\ge B$.

The retained $X$ entries of \texttt{xx1} are actual matrix
entries. A nonzero summand therefore bounds them by $q^{-A}$;
adding a matrix in $q^B\operatorname{Hom}(O^{n_i},O^{n_{i+1}})$
changes every summand by a period and leaves its phase unchanged.
Multiplication by the fixed \texttt{volN} factor in
\texttt{XiAt} preserves both statements. At the terminal
\texttt{joinCol} step the source tail has one coordinate and
the target tail is zero. In the scaled bases its extra column
is $q y$. Thus at a type prime support requires $q y$ integral,
and $y\mapsto y+z$ with $z\in O^n$ changes that column by a
$q$-period. The $x_0$ columns are integral and have $q$-periods.
At a twist prime the same argument with (CI.13) gives support
$q^{-M}$ and periods $q^M$ for $y$, and the weaker $q^{-A}$
support and $q^B$ periods for $x_0$. These prove exactly
\texttt{psiSupp}, \texttt{psiInv}, \texttt{xiSupp},
\texttt{xiInv}, \texttt{psiSuppX}, and \texttt{psiInvX}.

The projected integrality field follows from CI4 at twist
primes and $\Lambda_j\cap X_j=O^{n_j}$ at type primes.
The four group/character fields follow from the constructions
above: at type primes a residue matrix and any integral lift
give the same value by (CI.30), transported by CI14; at twist
primes any two projected lifts give the same value by CI8.
The zero and last uniqueness and matrix fields, predecessor,
injectivity, and coverage fields are precisely CI9. The
arithmetic character, depth, and Weil fields are the stated
inputs (CI.31). This checks every field of
\texttt{ChIdx.Ok} and proves the actual
\texttt{exists\_\allowbreak chIdx\_\allowbreak of\_\allowbreak cst}, including its twist normal-form
conclusion. No \texttt{Forms0}, \texttt{TC.Laws}, theta-algebraicity
or \texttt{Later} record has entered this proof.

\paragraph{CI16: the actual local constants, including the auxiliary Gauss sum.}
The constant hypotheses used in CI15 can themselves be
constructed from the root and local-character data already
isolated in TL14--TL15. Suppose $\xi_0\in O^\times$, and the
coefficient field contains a primitive $q^D$th root. Adjust its
power so that its specified complex embedding is
$\zeta=\exp(2\pi\mathrm i/q^D)$; two primitive roots differ by
a power coprime to $q^D$. For $x\in q^{-D}O$, let $b$ satisfy
$\xi_0x-b/q^D\in O$, and define
\[
 \psi(x)=\zeta_K^b;
 \qquad\psi(x)=1\text{ outside }q^{-D}O.
 \tag{CI.32}
\]
CI1 gives $b$, uniquely modulo $q^D$. Therefore the definition
is independent of its representative, is additive into
multiplication on $q^{-D}O$, is one on $O$, and is not one at
$q^{-1}$, since $\xi_0$ is a local unit. The stated Weil formula
holds for every allowed integer representative $b$, not just
the chosen one. At a type prime take $D=1$; the supplied
\texttt{Ch.WeilTied} formula identifies (CI.32) with
$\psi_0(\overline{qx})$. Thus its \texttt{psi0} clause is also exact.

For the canonical type-prime square root the source uses the
quadratic Gauss sum, not an arbitrary root. Let $\epsilon$ be
the quadratic character of $\mathbb F_q^\times$, extended by
zero, and $G=\sum_x\epsilon(x)\psi_0(x)$. Here $q\equiv1\pmod4$
and $\psi_0\ne1$. Translation of the complete additive sum
shows $\sum_x\psi_0(ax)=0$ for $a\ne0$, and multiplication by
a nonsquare shows $\sum_x\epsilon(x)=0$. With $y=xt$ in $G^2$,
\[
 G^2=\sum_t\epsilon(t)\sum_{x\ne0}\psi_0(x(1+t))
     =\epsilon(-1)(q-1)-\sum_{t\ne-1}\epsilon(t)=q.
 \tag{CI.33}
\]
The last equality uses $\epsilon(-1)=1$, since $4\mid q-1$.
This proves the actual \texttt{ciSq} square identity over $K$.
At twist primes a supplied square root witnesses the predicate
defining \texttt{twSqrt}. Its chosen value therefore squares to
$q$ and is nonzero. No sign or positive-embedding normalization
is part of that definition.

Choose $\mu_q$ to be the specified inverse reference local
character value times the inverse twist value at type primes;
at twist primes use the prescribed product of the two inverse
coefficient-field units. Their nonvanishing and all \texttt{muq}
equalities follow directly from those values and injectivity
of the embedding. The residue-unit tie is the literal
\texttt{Ch.Tied} clause. This constructs
\texttt{exists\_\allowbreak locCst\_\allowbreak of\_\allowbreak values}.

The tied data with integral weights also supply the required
reference values at ramified primes. Here is the additional
argument behind \texttt{locK\_\allowbreak mem\_\allowbreak of\_\allowbreak tied}. For a fixed local
element $x\in E^\times$ at $P$, choose $a\in E^\times$ with
$a/x\equiv1\pmod P$ and $a\equiv1$ at the other finitely many
ramified primes. This is elementary approximation: clear the
denominators by an integer $b$, solve for an algebraic integer
$y$ modulo the requisite powers of the distinct prime ideals
by the Chinese remainder theorem, and set $a=y/b$; choose a
nonzero solution. Tameness gives $\mu_P(a)=\mu_P(x)$ and makes
the other ramified factors one. At every unramified prime the
local value at $a$ is an integral power of its tied uniformizer
value in $K$, and only finitely many such powers are nonzero.
The infinity type has exponents $-2z_\tau$ with $z_\tau\in\mathbb Z$,
so its value is a product of
$(\tau(a)/\overline{\tau(a)})^{-z_\tau}$, also in the embedded
coefficient field. The Hecke product formula now places
$\mu_P(x)$ in that field. This proves the reference-value
supplier from the literal tied integral-weight data, and gives
the full \texttt{exists\_\allowbreak locCst} used for the initial indices.

For initial indices take $M=A=B=m$, $D_Q=2m$ and
$D_G=1$. For canonical indices the actual definition is
$M=m$, $A=B=m+1$, $D_Q=2m$, $D_G=1$, with
\texttt{sharp=true}. The nine fields of \texttt{LocCst.Ok}
now follow from (CI.32), (CI.33), the supplied twist square
roots, and these inequalities. In particular the larger canonical
box depths do not change $M$ to $m+1$.

\paragraph{CI17: the specific canonical indices satisfy their full record.}
Consider the source's actual \texttt{twCst} and \texttt{twIdx},
not an arbitrary existence witness. Assume $X.\texttt{Ok}$,
$Pr.\texttt{Ok}$, $Ch.\texttt{Ok}$, the given character and
Weil ties, type-prime inverse reference values in $K$,
$q_P\equiv1\pmod4$ at type primes, degree-one unramified twist
primes disjoint from them, and the root, square-root, and
local-unit clauses of \texttt{TwCstQ}.

The canonical step selects a class at each twist prime among
admitted matrices having a nonzero predecessor in the preceding
global index set. CI3 makes this class set finite, so its
``finite if finite, empty otherwise'' definition takes the
finite branch. CI4 and CI7 provide a member of every such class
satisfying all predicates in the choice of \texttt{twRep}:
admission, normal form, projected integrality, and the original
nonzero-predecessor condition. Thus the chosen representative
has those properties, regardless of which eligible witness
the choice operation returns.

The canonical matrix is one at every type prime. The definition
of \texttt{twNext} retains its explicit condition that every
type prime have a predecessor into that identity matrix.
CI13 proves the identity nonzero assertion and the one-coset
support statement needed for this condition. For coverage,
start with any tuple of admitted matrices having the required
predecessors: its twist classes choose the product component,
and its type components are in the identity double coset by
CI13. Transport to the identity supplies the defining condition
of \texttt{twNext}. Conversely this condition and the twist
candidate condition give every requested predecessor. Equality
of all local double cosets gives equality of the selected
product of classes, hence equality of indices. This proves the
actual predecessor, injectivity and coverage fields, without
assuming that a candidate set is nonempty.

At positive twist steps CI8 proves the actual character used
by \texttt{twLam}; at the last step its projection kernel is
zero, and at zero it is the trivial top level. At all type
steps CI14 identifies the group and character with the actual
\texttt{colLevSym} definition. The support, periodicity and
depth fields follow from the uniform calculations in CI15,
now with the canonical constants of CI16. The base and terminal
index types are the one-element type stipulated by the source,
and the types above the terminal step are empty. The arithmetic
ties come from CI16 and the present literal hypotheses.
This proves \texttt{twIdx\_\allowbreak ok\_\allowbreak closed}; its formerly separate
\texttt{TwFinQ}, \texttt{TwNormDC}, \texttt{TwChrQ},
\texttt{TwZeroQ}, and \texttt{TwAnQ} suppliers have all been
proved by CI1--CI16.

\paragraph{CI18: canonical replacement really preserves \texttt{Ok}.}
Let $I$ be an actual supplied column datum with
$I.\texttt{Ok}$ and old twist depths at least $2m$, as in the
source's \texttt{Deep} hypothesis. TL20 already extracted its
primitive $q^{2m}$th roots, square roots, and inverse reference
local values. To recall the root step precisely: if the old
depth is $D=2m+b$, insert
$x=\xi_0^{-1}/q^{2m}$ and integer $q^b$ in its Weil identity.
The local-unit clause verifies both allowed denominators, and
the resulting value has complex image
$\exp(2\pi\mathrm i/q^{2m})$. Its order is exactly $q^{2m}$
by injectivity. This extracts the arithmetic premises only;
the nonvanishing and kernel arguments are supplied by CI7,
CI8, and CI10--CI14, not by this extraction.

Apply CI17 for each tuple in $I$. Replacing its constants and
indices by the defined \texttt{twCst} and \texttt{twIdx} gives
their actual \texttt{ChIdx.Ok}. Every other field of $I$ is
unchanged. The coordinate, host, character, tuple-reading,
selection, and ambient-data fields of \texttt{TC.Ok} are
therefore inherited literally. This proves
\texttt{twIdx\_\allowbreak ok\_\allowbreak of\_\allowbreak ok} and \texttt{canon\_\allowbreak ok};
\texttt{IsTw} follows from the definitions. TL13 then derives
the column laws for the canonical datum. No use of those laws
was made in proving its local-index predicates.

\paragraph{CI19: the resulting canonical evaluation construction boundary.}
For the initial constructor, choose the finite character family,
tuples, residue characters and coefficient fields of TL14--TL16.
Supply the literal field/host/selection data and the ambient
period datum of the constructor. CI16 gives local constants
with $M=A=B=m$ and twist depth $2m$; CI15 gives their actual
indices. Installing those records therefore supplies the
local-index fields of the initial \texttt{TC.Ok}. The other
fields are precisely the supplied raw constructor data, and
must be checked by their own arithmetic and ambient suppliers.
CI18 now gives a canonical datum with \texttt{Ok} and
\texttt{IsTw}, and TL13 gives \texttt{Laws}.

Accordingly the local-coset, normal-form, projected-character,
and first-unfolding boundary in TL20--TL21 is discharged. The
remaining hypotheses in that \texttt{EvLg} construction include
the actual host and raw arithmetic data and the normalized-theta
vacuum/indicator identities, smoothness and covariance used by
TL18--TL19 for \texttt{JetsIn}. In this version PA9 separately
supplies the literal ambient period/trace-isometry record,
and PA13 supplies its CM-field conclusion. They can be combined
with the present construction using the same ambient datum;
neither was a premise of CI1--CI18. AH17--AH18 separately construct
the arithmetic host under their exact prime and splitting hypotheses,
and TA10/TA12 supply the one-dimensional normalized theta arithmetic.
These do not by themselves identify the actual vacuum and indicator
family or prove its covariance. The coefficient-field containment
antecedents and their original order of quantifiers remain
unchanged. None of this produces the five independent stages
of \texttt{Later}, or proves RH or GRH from the local index
calculation alone.


\section{Variable twists and the complete twisted column comparison}
\label{sec:twisted-comparison}

The reference comparison TG12 extends to every member of the actual
twist family. The construction below proves the local chart statements
before assembling their global comparisons. In particular it retains
the prescribed frame depth, the inverse in the level character, and the
actual moving element in the terminal period map.

We use the supplied arithmetic host, the true column datum \(I\), its
literal \texttt{Ok} and \texttt{Pins}, and its \texttt{Later} data.
The laws may be taken from TL13. Fix a word \(u\) and a twist \(a\).
Write \(Q_E(a)\) for its selected primes of \(E\), and \(Q_F(a)\)
for their underlying primes of \(F\). The selection contains no conjugate
pair; these primes avoid the type and auxiliary primes. At every
\(P\in Q_E(a)\), the original \(\pi=\pi_u\) is a unit.
Completions are interpreted through the actual product
\(E\otimes_F F_v=\prod_{P\mid v}E_P\).
The notation \(\Delta\) always denotes the integral Weyl lattice used
in CV and CR. We never replace it by its dual Weyl lattice.

\paragraph{TX1: completing a rational level and its character.}
Let \(P\) be a finite prime, let \(H\subset\operatorname{GL}_n(O_{E,(P)})\),
and let \(\lambda:H\to K^\times\) be a character. Suppose that the
principal congruence subgroup of depth \(N\) is contained in \(H\)
and that \(\lambda\) is trivial on it. Let \(d\ne0\) have valuation at
least \(N\). In the completed integral general linear group put
\[
 K_P(d)=\{z:z-I\in dM_n(O_{E_P})\},\qquad
 \widehat H=H K_P(d).
\]
Here the rational matrices in \(H\) are embedded by the prescribed
completion map. Conjugating \(z-I\) by an integral matrix with integral
inverse preserves \(dM_n(O_{E_P})\). Thus \(K_P(d)\) is normal in
the completed integral general linear group, and \(\widehat H\)
is a subgroup.

There is a unique character
\[
 \widehat\lambda(gz)=\lambda(g)\quad(g\in H,\ z\in K_P(d)).
 \tag{TX1.1}
\]
Indeed two such expressions imply that the rational matrix
\(g^{-1}g'\) is congruent to \(I\) to depth \(N\). It therefore
belongs to the principal subgroup on which \(\lambda\) is trivial.
Normality proves multiplicativity. The same argument gives the exact
intersection
\[
 \widehat H\cap\operatorname{GL}_n(E)=H
 \quad\text{among the rational matrices integral with integral inverse.}
 \tag{TX1.2}
\]
It also shows that any exponent annihilating \(\lambda\) annihilates
\(\widehat\lambda\). The subgroup is open since it contains \(K_P(d)\).
These are the source constructions \texttt{qbU} and \texttt{qbChr};
the rational intersection and character restriction are part of the
proof, rather than assumptions on an unspecified completed level.

\paragraph{TX2: the actual terminal depth and the paired-prime diagonal.}
Let \(c=\texttt{topDep}(R,I.\texttt{cst}(a).M)\).
The terminal rational level is generated by the block extension of
the row level and the principal subgroup defined by divisibility
of the matrix discrepancy by \(q_P^{\,c+1}\).
The latter is normal in the integral general linear group, so this
generated subgroup is their product. Its character is the row
character on the first factor and is trivial on the principal factor.
Consequently the precise source choice
\[
 d_{\rm top}(a,P)=q_P^{\,c+1}
 \tag{TX2.1}
\]
satisfies the hypotheses of TX1. This uses the source valuation of
\(q_P\): in additive \(P\)-valuation its principal depth is
\(N_P=v_P(q_P^{\,c+1})\). One must not identify \(N_P\) with
\(c+1\) without a degree-one unramified hypothesis. Nor do we
replace the prescribed exponent \(c+1\) by an unspecified large
integer.

Put \(D_\pi=\operatorname{diag}(1,\ldots,1,\pi)\). At each
selected \(P\) prescribe the component \(D_\pi\); at its conjugate
prescribe its inverse conjugate transpose; prescribe identity at all
other primes. Because no selected prime is conjugate to another, these
prescriptions are compatible. The resulting finite adele \(t_a\)
is unitary. Both paired components are integral with integral inverse,
since \(\pi\) is a unit at these primes. All prescribed type and
auxiliary reductions are identity. Thus
\[
 t_a\in K^1_{ST,SG},\qquad
 \rho_P(t_a)=D_\pi\ (P\in Q_E(a)),\qquad
 (t_a)_v=1\ (v\notin Q_F(a)).
 \tag{TX2.2}
\]
This constructs the required element directly. The rational matrix
\(D_\pi\) itself is not being asserted to be unitary: unitarity is
the paired-prime condition on the entire adele.

\paragraph{TX3: the framed level dictionary, including its inverse.}
Let \(K_0\) be the reference terminal chart, with level \(U_0\),
character \(\chi_0\), and its actual standard companion \(k\).
The reference reading gives
\[
 x\in U_0\ \Longleftrightarrow\ k^{-1}xk\in\operatorname{levU}(\Lambda),
 \qquad
 \chi_0(x)=\sigma(\chi_\Lambda(h))
 \tag{TX3.1}
\]
when the auxiliary reduction of \(k^{-1}xk\) is \(h^{-1}\).
At \(Q_E(a)\), \(k\) is identity. Complete each actual terminal
frame level \((\Theta_P,\lambda_P)\) by TX1 and TX2, and set
\[
 U_a=U_0\cap\bigcap_{P\in Q_E(a)}
 \{x:\rho_P(t_a^{-1}xt_a)\in\widehat\Theta_P\},
 \quad
 \chi_a(x)=\chi_0(x)
 \left(\prod_{P\in Q_E(a)}
       \sigma\widehat\lambda_P(\rho_P(t_a^{-1}xt_a))\right)^{-1}.
 \tag{TX3.2}
\]

Suppose \(ykt_a\) and \(yxkt_a\) realize framed points \(p_0,p_1\).
Put \(z=(kt_a)^{-1}x(kt_a)\). Since \(t_a\) lies in the congruence
kernel of the reference level, it normalizes that level and has
trivial auxiliary reduction. Thus the underlying decoration and
auxiliary conditions on \(z\) are exactly the reference ones.
If \(A_P,A'_P\) are the two rational frames, equality of their
underlying local lattices gives
\[
 \rho_P(z)=A_P^{-1}A'_P
       \in\operatorname{GL}_n(O_{E,(P)}).
 \tag{TX3.3}
\]
By (TX1.2), the new level condition is precisely
\(\theta_P=(A_P^{-1}A'_P)^{-1}\in\Theta_P\).
The frame action is \(A_P\mapsto A_P\theta_P^{-1}\).
Together with the reference auxiliary action this proves
\[
 x\in U_a\quad\Longleftrightarrow\quad
 p_1=(1,h,\theta)\,p_0
 \quad\text{for some actual finite-level element }(h,\theta).
 \tag{TX3.4}
\]
The local completion character in (TX3.2) is evaluated at
\(\theta_P^{-1}\). Its outer inverse therefore gives
\(\lambda_P(\theta_P)\), and (TX3.1) gives
\[
 \chi_a(x)=\sigma\bigl(\chi_\Lambda(h)\chi_\Theta(\theta)\bigr).
 \tag{TX3.5}
\]
In particular \(\chi_a(x)=\chi_0(x)\) for the same matrix \(x\)
whose components at \(Q_F(a)\) are identity. This is a pointwise
equality of characters on a specified subgroup, not an identification
of characters on unrelated isomorphic groups.

\paragraph{TX4: constructing the framed chart, rather than postulating it.}
The level \(U_a\) is open, lies in the prescribed lattice stabilizer,
and contains a principal congruence subgroup at every relevant place.
It therefore defines the required step level with the same base
lattice \(L_\pi\). Outside \(Q_F(a)\) its local group and its lattice
are literally those of \(K_0\).

For each preferred framed-point class \(c\), DR3 supplies a full
\texttt{IsPtOf} realizer \(g_c\) of its extension to the actual \(SQ\)
point. Choose
\[
 y_c=g_c(kt_a)^{-1}.
 \tag{TX4.1}
\]
The point reading is then true for these very representatives.
The level dictionary (TX3.4) proves injectivity of the associated
map from point orbits to unitary double cosets.

Here is the surjectivity argument, including the rational-frame issue.
Start with an arbitrary unitary adele \(y\). At the finitely many
selected primes approximate the columns of \(ykt_a\), which give
completed bases of full lattices. Rational
matrices are dense in these completions, and the conditions of being
a basis, preserving each prescribed congruence, and having discrepancy
in \(K_P(d_{\rm top})\) are open. Choose rational matrices in these
neighbourhoods. Prescribe the corresponding small correction at each
\(P\), its inverse conjugate transpose at \(\bar P\), and identity
elsewhere. The paired-prime splicing of AD4 gives a correction
\(z\in K^1_{ST,SG}\) with precisely these components. The element
\(x=kt_a z(kt_a)^{-1}\) lies in \(U_a\), by (TX3.2), and
\[
 yxkt_a=ykt_a z
\]
has the chosen rational frames. The intersection construction for its
lattices and reductions gives a
valid framed point, hence one of the point classes. This proves
surjectivity.

The same dictionary gives the stabilizer identity
\[
 \operatorname{im}(\Gamma_c\longrightarrow U(E))
       =U(E)\cap y_cU_a y_c^{-1}.
 \tag{TX4.2}
\]
Injectivity follows from the actual point action: once its rational
component and the point are fixed, its auxiliary reductions and frame
components are forced. Pulling back the standard completed lattices
through \(y_c\) recovers the source lattice assigned to \(c\).
Thus the orbit bijection, the integral lattices, and the stabilizers
are all attached to (TX4.1). These are the substantive clauses of
\texttt{exists\_\allowbreak clsG\_\allowbreak top}.

\paragraph{TX5: terminal readings, roots, base points, and extension.}
At the last step the index set is a singleton and the coupling space
is empty. TX4 consequently gives the actual single chart of
\texttt{chart\_\allowbreak twm}. Its \texttt{Reads} clauses are as follows.
Its base lattice is \(L_\pi\); its companion is \(kt_a\); its
components outside the type and twist primes are unchanged; its
inverse has the required \(\bar\pi\) last row at type primes; and
its representatives, local move dictionary, character dictionary,
and stabilizers are exactly (TX4.1), (TX3.4), (TX3.5), and (TX4.2).
This proves \texttt{reads\_\allowbreak top}.

The reference character has finite exponent \(e_0>0\). Each actual
frame character has finite exponent; a common multiple \(e_Q>0\)
annihilates their product and its inverse. Thus \(e_0e_Q\)
annihilates \(\chi_a\), proving \texttt{root\_\allowbreak top}.
If \(gkt_a\) and \(ykt_a\) realize the same preferred point,
apply (TX3.4) with the identity finite-level element. It gives
\[
 g^{-1}y\in U_a,\qquad \chi_a(g^{-1}y)=1.
 \tag{TX5.1}
\]
In particular it gives \texttt{top\_\allowbreak base}, with its exact order
of multiplication.

Finally let \(g'\) realize an actual row point \(d\), and suppose
\(\operatorname{IsExt}(z_0,g',g)\), where
\(z_0=\pi e_{\rm last}\). Since \(\pi\ne0\), the extension
equations force \(g=\operatorname{diag}(g',1)\). The underlying
decoration of \(gk\) is the terminal extension of \(d\), by TG6:
\[
 N=N_{\rm row}\oplus O_E,\quad
 L=N_{\rm row}\oplus\pi O_E,\quad
 L^\sharp=N_{\rm row}\oplus\bar\pi^{-1}O_E.
 \tag{TX5.2}
\]
Its functional is the reduction of \(\bar\pi\) times the last
coordinate. The auxiliary reading uses the prescribed congruence
\(\pi\equiv1\) there. Right multiplication by \(t_a\in K^1_{ST,SG}\)
preserves all these data. At \(P\in Q_E(a)\), however, it supplies
the required point frame:
\[
 \operatorname{diag}(A_P,1)\,D_\pi
       =\operatorname{diag}(A_P,\pi).
 \tag{TX5.3}
\]
Thus \(gkt_a\) realizes \(\widetilde F(d)\), including every twist
frame. Acting matrices are \(\operatorname{diag}(\gamma,1)\);
point frames are (TX5.3). This proves the full extension clause
of \texttt{chart\_\allowbreak twm}, and hence its \texttt{TopBase} implication.
No separate terminal point or realization premise remains.

\paragraph{TX6: middle charts with the actual coupling stabilizers.}
For \(1\le i<m\), a point has an index \(b\), an auxiliary coupling
class \(\bar C\), and at each twist prime a coupling coset \(C_P\).
The latter is empty away from the penultimate stage. At that stage
it has the form
\[
 C_P=z_P+q_P^{M}O_{E,(P)}^n,
 \qquad z_P\in q_P^{-M}O_{E,(P)}^n,
\]
as follows by applying the actual rational frame to the two
\texttt{locLat} bounds in \texttt{ybSet}. Define a chart fibre by
fixing these data. There are finitely many fibres: the index set and
auxiliary residue spaces are finite, and the twist cosets are indexed
by the finite quotient
\(q_P^{-M}O_{E,(P)}^n/q_P^{M}O_{E,(P)}^n\).
Its additive depth is \(2M v_P(q_P)\), with the actual valuation.

In the fibre \(\beta\), cut the auxiliary level by the stabilizer
of \(\bar C\), and cut the rational twist level by the stabilizers
of the actual \(C_P\). The uncut twist level is a subgroup of
\(\operatorname{GL}_n(O_{E,(P)})\): this is part of the definition
of \texttt{LevI}, also before the penultimate stage. Its character
is trivial on a sufficiently deep principal subgroup, by the actual
\texttt{ThetaOpen} statement \texttt{lvQ\_\allowbreak open}. Here is the
underlying reason. Bound the denominators of the index matrix
\(h\) and \(h^{-1}\). If \(g\) is sufficiently close to identity,
then \(h^{-1}\operatorname{diag}(g,I)h\) and its inverse satisfy
the defining twist-level congruences, with identity residue matrix.
The zero extension is integral, so the additive character is trivial
on its phase; the residue row character is also \(1\).
The actual group and character readings in \texttt{Ok} then put
\(g\) in the twist level with character \(1\).

Increase this depth until
\((g-I)z_P\in q_P^{M}O_{E,(P)}^n\).
Integral \(g,g^{-1}\) already preserve the direction lattice, so
this additionally stabilizes \(C_P\); the empty coset gives no
extra restriction. Thus the cut level and its character satisfy
TX1. Character triviality here comes from the proved congruence
statement, rather than from finite order alone.

Choose an admissible companion \(k_\beta\), with identity components
at \(Q_E(a)\), as supplied in TX7. Apply DR3 to the actual preferred
points and use representatives \(y_c=g_c k_\beta^{-1}\).
The proof of TX4 now applies with these stabilizer cuts: a local
move is allowed precisely when it preserves the fixed coupling
data. The criterion is necessary by applying the move to the
point, and sufficient because the auxiliary and twist reductions
recover the coupling class. Rational approximation inside the
cut open subgroups proves double-coset surjectivity. This gives
finite charts, actual readings, lattices and stabilizers at every
middle step, with their actual finite-order characters.

The subsequent \texttt{FunKilled} normalization changes only the
value module and base functional by the FC4 construction on
\(L_\pi^\sharp/L_\pi\).
It changes neither the local groups nor their characters, the
representatives, the coupling stabilizers, or the companion
\(k_\beta\). Hence all readings persist. FC9 then transports
these same charts to the Hermitian description, giving
\texttt{SameCar}, rather than choosing unrelated Hermitian charts.

\paragraph{TX7: matching a middle chart to the reference away from \(Q\).}
We justify the admissible companions and the reference comparison
used in TX6. At an auxiliary prime, a nonzero actual index
representative lies in the identity double coset. The induction
starts with the zero-step identity, uses nonvanishing of the
preceding kernel and double-coset transport at each successor,
and ends with the prescribed final identity. The finite-sum
factorization and identity nonvanishing needed for this argument
are the local statements of CI7 and CI13. Thus, for a twisted
representative \(h_a\) and a reference representative \(h_0\),
write \(h_a=r_a\kappa_a\), \(h_0=r_0\kappa_0\) with \(r\) in
the left level and \(\kappa\) in the right level. The equality
\[
 h_a=(r_ar_0^{-1})h_0(\kappa_0^{-1}\kappa_a)
 \tag{TX7.1}
\]
conjugates the projected level groups and their characters.
It follows by substituting in the defining stabilizer equation;
the kernel eigencharacter is unchanged under the corresponding
affine change of summation variables. At the penultimate stage
both representatives are literally identity, so this conjugator
can be chosen to be identity there.

At that exceptional stage one must also match the coupling class.
Choose a vector \(y\) representing the given auxiliary coupling
residues, and form in the reference datum the coset
\[
 Y_0=y+\operatorname{locLat}_{SG}(N;0).
 \tag{TX7.2}
\]
There is no reference \(Q\)-condition. The actual reduction
formula shows that \(Y_0\) has exactly the given auxiliary coupling
vector. Move this point to its preferred representative. If the
auxiliary component of this move is \(g_0\), use \(g_0^{-1}\)
as the additional conjugator. The equality
\(\bar C_0=g_0^{-1}\bar C\) shows directly that the two coupling
stabilizers are conjugate. Before the penultimate stage the
coupling spaces are empty and this correction is unnecessary.

Lift the resulting finite auxiliary matrices to a unitary adele
\(s_\beta\), identity at type primes and outside the auxiliary
places, using paired-prime integral lifting. This is the AD4
construction; no global surjectivity of
\(\operatorname{GL}_n(O_E)\to\operatorname{GL}_n(O_E/P)\)
is needed. Put
\[
 k_\beta=k_{0,g(\beta)}s_\beta.
 \tag{TX7.3}
\]
The conjugators preserve the standard lattices, so this is an
admissible companion and is identity at \(Q_E(a)\).
For an adele whose \(Q_F(a)\)-components are identity, its frame
move is identity. Equations (TX7.1)--(TX7.3) and the coupling
calculation then give the same membership condition and the same
character value in the chart and its reference chart. Taking
adeles supported at a single place proves the literal equality
\[
 U_{\beta,v}=U_{0,g(\beta),v}\quad(v\notin Q_F(a)),
 \qquad
 \chi_\beta(x)=\chi_{0,g(\beta)}(x)
       \quad(x_v=1\text{ for }v\in Q_F(a)).
 \tag{TX7.4}
\]
The base lattices are also the same \(L_\pi\). The reference
twist itself uses its reference charts unchanged.

\paragraph{TX8: a character on the common penultimate local group.}
A character on each chart need not extend to the group generated
by all charts. The following argument supplies that missing step.
Fix \(v\in Q_F(a)\) and the penultimate index, whose representative
is identity. For each selected \(P\mid v\), use the positive
integer \(D(P)\) in the actual local constants. Let
\(\mathcal N_v\) consist of local matrices which have a rational
approximation \(g_P\) to this depth satisfying the actual rational
middle-level relation, with character value \(z_P\).
This is exactly the source \texttt{NearSet}; approximation and
congruence are measured through its prescribed completion maps.

Two witnesses \(g_P,g'_P\) for the same matrix differ by an
integral matrix congruent to identity to depth \(D(P)\).
Both matrices and their inverses are integral: at the penultimate
stage the projected rational level lies in
\(\operatorname{GL}_n(O_{E,(P)})\). At auxiliary primes its
character is the prescribed residue character, and at twist
primes it is the actual \(\lambda_P\) on the rational level.
At this penultimate step the tail is empty, so the additive
\texttt{n1Arg} phase in \texttt{chiInv} is absent. The remaining
\texttt{lamRow} character factors through the actual residue
matrices. The positive prescribed exponent \(D(P)\) makes a
matrix congruent to identity modulo \(q_P^{D(P)}\) act trivially
in these residue matrices. Both characters are therefore trivial
on the stipulated congruence subgroup, and
\[
 z_P=z'_P.
 \tag{TX8.1}
\]
This proves independence of every rational-lift choice.

Products and inverses of witnesses are again witnesses. For
example the discrepancy identity
\[
 (AB)^{-1}(CD)
   =(B^{-1}D)\bigl(D^{-1}(A^{-1}C)D\bigr)
 \tag{TX8.2}
\]
preserves the same congruence depth because \(D,D^{-1}\) are
integral; positivity of the depth ensures the required integral
neighbourhood. The rational middle-level character is
multiplicative. Thus \(\mathcal N_v\) is a group, and
\[
 \operatorname{val}_v(x)
       =\sigma\!\left(\prod_{P\mid v}z_P\right),
 \qquad
 \psi_v(x)=\operatorname{val}_v(x)^{-1}
 \tag{TX8.3}
\]
are well-defined characters. Factors at primes not over \(v\)
are \(1\), as in the literal source definition.

Each penultimate chart group at \(v\) lies in \(\mathcal N_v\).
Indeed approximate its completed frames by rational frames while
preserving the finitely many open level and coupling conditions.
The actual local-move dictionary then supplies a rational
middle-level witness; \(k_\beta=1\) at \(P\) identifies its
congruence with the stipulated \(D(P)\)-congruence. Its character
is the inverse product of these rational values, by the same
inverse-reduction calculation as TX3. Consequently
\[
 \chi_\beta(\operatorname{emb}_v x)=\psi_v(x)
       \quad(x\in U_{\beta,v}).
 \tag{TX8.4}
\]
Different close witnesses give the same value by (TX8.1).
Let \(W_v\) be the subgroup generated by all \(U_{\beta,v}\).
Then \(W_v\subset\mathcal N_v\), and the restriction of \(\psi_v\)
is the required common character on \(W_v\).

The penultimate chart family is nonempty. Choose the index supplied
by \texttt{last\_\allowbreak uniq}, the base decoration, and the zero-centred
actual \texttt{locLat} coupling coset. It is allowed since
\(0\) belongs to every lattice involved. This gives a point class
and hence a chart. Its principal congruence subgroup lies in
\(W_v\), proving that \(W_v\) is open.

Each chart group has finite index in \(W_v\). All groups lie
in the same local lattice stabilizer, and each chart contains a
principal congruence subgroup. The latter is the kernel of the
action on a finite module \(L/P^rL\); its containing group
therefore has finite index. This argument establishes both the
common character and the finiteness needed below.

\paragraph{TX9: the common penultimate level and its finite coupling carrier.}
Let \(H_0\) be the reference common penultimate level of FC12 and
TG2, with its identity auxiliary frames. Define \(J\) by replacing
its local group at \(v\in Q_F(a)\) by \(W_v\), keeping all other
local groups and the base lattice unchanged. The product
decomposition of the finite adele ring gives a homomorphism
\(\operatorname{off}_Q:J\to H_0\) which replaces the
\(Q_F(a)\)-components by identity. Put
\[
 \chi_J(\kappa)=
       \chi_{H_0}(\operatorname{off}_Q\kappa)
       \prod_{v\in Q_F(a)}\psi_v(\kappa_v).
 \tag{TX9.1}
\]
Decompose a chart element into its part supported on \(Q_F(a)\)
and its part identity there. TX8 identifies the first character;
TX7 identifies the second. Thus \(\chi_J\) restricts to the
actual character of every penultimate chart. It equals the
reference common character on the same elements identity at \(Q\).

For completeness, the finite carrier can be constructed explicitly.
Let \(\rho:J\to G_{\rm aux}\) be the actual auxiliary reduction
map, and let \(w_\beta\) be the auxiliary coupling vector of
chart \(\beta\). If \(K_{\beta,v}=U_{\beta,v}\subset W_v\),
the full chart stabilizer in \(J\) is
\[
 H_\beta=\{\kappa:
       \rho(\kappa)w_\beta=w_\beta,\quad
       \kappa_v\in K_{\beta,v}\ (v\in Q_F(a))\}.
 \tag{TX9.2}
\]
Take the disjoint union of the \(J\)-orbits of
\((w_\beta,(K_{\beta,v})_v)\). Record the chart label in the
second coordinate, so distinct chart orbits remain distinct.
Equivalently this carrier is \(\coprod_\beta J/H_\beta\), with
the first map \(y_G\) recording the translated auxiliary vector
and the second map \(y_Q\) recording the labelled local cosets.
Equation (TX9.2) proves that the pair \((y_G,y_Q)\) is injective.
It also proves the required stabilizers and equivariance:
\(y_G(\kappa y)=\rho(\kappa)y_G(y)\), and \(y_Q\) is unchanged
when all \(\kappa_v\), \(v\in Q_F(a)\), are identity.
The carrier is finite by TX8 and finiteness of the auxiliary
residue spaces.

At auxiliary primes \(J\) is the exact symmetric group
\(H_{\rm Sym}\) of TG2, with its actual reduction map and
identity frame. Away from the auxiliary and twist primes it
is the reference local lattice group. This proves all the
\texttt{TwLevJ} and \texttt{SymmetricTwistedJointForms} clauses,
including the common character; no unsupported extension of
characters to a generated group is used.

\paragraph{TX10: constructing \texttt{Law0} with the integral lattice.}
At step zero the auxiliary level is independent of the twist,
every twist-frame level is the whole integral general linear
group, the coupling coset is empty, and the index is unique.
Forgetting the twist frames therefore gives a bijection
\[
 e_a:\mathrm{Cl}(u,\mu_0,0)\longrightarrow\mathrm{Cl}(u,a,0).
 \tag{TX10.1}
\]
Surjectivity follows by choosing rational bases of the localized
underlying lattice at the finitely many frame primes. If two
underlying decorations differ by a rational unitary element,
their frame discrepancy is integral with integral inverse and
hence belongs to the full zero-step frame level. This proves
injectivity on point orbits.

Preferred representatives require a further check. Their
underlying decoration is the canonical representative for the
common action on decorations. Taking the canonical representative
again has no effect. Thus the preferred decorations corresponding
under (TX10.1) are literally equal, and in particular so are
their lattices \(L^\sharp\), functionals, and auxiliary maps.
The zero-step stabilizers are trivial by the actual free-action
law proved in TL. Consequently coefficient functions on the
finite class set satisfy the equivariance condition automatically.
Pullback along \(e_a^{-1}\) gives the \(L\)-linear equivalence
\[
 E_a f(c)=f(e_a^{-1}c)
 \tag{TX10.2}
\]
between the actual rational form spaces.

Membership in the global integral lattice is pointwise
membership in the same integral \(\Delta\)-lattice attached to
the preferred \(L^\sharp\). Its equality at corresponding
classes proves, in both directions,
\[
 f\in\mathcal L_{u,\mu_0,0}
       \quad\Longleftrightarrow\quad
 E_af\in\mathcal L_{u,a,0}.
 \tag{TX10.3}
\]
The coefficient embeddings commute with the same class
reindexing, proving the literal embedding clause of
\texttt{Law0}. At \(a=\mu_0\), the canonical class equivalence
is identity; injectivity of the form embedding forces the
chosen comparison to be the specified \(e_{\mu_0}\).
More generally, the embedding clause and injectivity of the
form embedding force every permissible \(E_a\) to equal
(TX10.2). The conclusions therefore hold for any supplied
\(\ell_0:\texttt{Law0}\), not only for the witness just
constructed.
Thus \texttt{Law0} is derived from the given \texttt{Ok} and laws.
It is not an additional arithmetic or integral-comparison premise.

\paragraph{TX11: all-twist chains, values and the same reference object.}
Use TG12 for the reference columns and retain that object throughout.
For each \((u,a)\), choose the middle charts of TX6--TX7, the
terminal chart of TX5, and the common penultimate level of
TX8--TX9. Make the terminal choice first and use that same choice
when collecting the steps. At \(a=\mu_0\) use the original
reference charts. At step zero use (TX10.1); outside the operative
range the chart index can be a singleton, since its reading
conditions have no instances.

Set the class lattice equal to the actual \(I.\texttt{latAt}\),
and use the identity map on frame coordinates. The coordinate
maps to the reference Hermitian spaces have the same underlying
frames by \texttt{StdF}; thus \texttt{IdY} holds literally.
The \texttt{TwChTied} data now use these chart sets, these
lattices, the actual point stabilizers, and the finite carrier
of TX9. Its level containment and stabilizer clauses are
(TX9.2). The reference comparison map \(g(\beta)\), the local
equalities (TX7.4), and prime disjointness give
\texttt{TwChJointS}. TX9 gives its symmetric joint forms.

The type-prime comparison can be chosen with both the prime
identification and residue-field identification equal to identity.
For the diamond operator at step zero, the two moving elements
act between equal underlying preferred decorations by TX10.
Their quotient fixes that decoration. Add the uniquely determined
frame correction in the full zero-step frame level; freeness
forces both their rational and auxiliary components to agree.
The zero-step frame character is trivial, so their full
transport factors also agree. This proves the diamond comparison.
The zero-step complex form spaces are related by (TX10.2),
since freeness removes the stabilizer condition. All other
type-prime parameters are the actual unchanged parameters.
This constructs \texttt{TwAtT}, with the stipulated residue
identity, rather than assuming it.

Apply CV9 to each actual reading and \texttt{IdY}. This supplies
the complex image, rational image, and the place-by-place
integral \(\Delta\)-lattice assertion. Restricting the latter
to the required good coefficient places gives exactly the
\texttt{ValuationCompatibleChart} field of \texttt{TwFormsTiedS}.
FC18 applies to the normalized charts because TX6 retained
\texttt{FunKilled}, all readings, and the same type-prime map.
It supplies the diamond clause. FC7 applies to these same frame
functions and actual integral form space, supplying its metric
and norm clauses. Equations (TX7.4) and (TX9.1) supply every
off-\(Q\) character clause.

The reference frame equivalence is identity, and TX10 identifies
the chosen zero-step equivalence with \(e_{\mu_0}\). The standard
data therefore have exactly the source \texttt{StdW} fields:
the actual \(e_{\rm Cl,0}\), identity frame and prime maps,
\(I.\texttt{latAt}\), the prescribed residues, all readings,
the terminal base and extension clauses of TX5, and
\texttt{FunKilled}. This proves \texttt{tw\_\allowbreak chains\_\allowbreak std} and
\texttt{tw\_\allowbreak ties\_\allowbreak with\_\allowbreak points}, with \texttt{Law0} supplied by
TX10 and no separately assumed chain or base-point object.

\paragraph{TX12: the actual level maps and zero/reference compatibility.}
On each class use the quotient projections of FC5, with their
actual kernels, for the level and kernel-level modules.
The type-prime residue map is identity under TX11's
identification. Thus FC6 identifies the transported functional.
NB7--NB8 apply to each actual up/down move: DR3 supplies its
initial full realizer; the supported splice preserves every
other type functional, auxiliary reduction, twist frame and
coupling label. The original moving element therefore gives
the required transport factor. This proves the up and down
naturality clauses for these very projections. The integral
basis clause is FC17 applied to CV9. Hence the resulting level
data are \texttt{SymmetricTwistedLevelData}, and their
\texttt{StdLv} projections are the literal quotient maps.

At step zero TX10 identifies the preferred decorations and
all their local kernels. Raw up lattices depend only on the
residue class: if \(x-x'\in PN\), then
\(P^{-1}(x-x')\subset N\), so
\(N+P^{-1}x=N+P^{-1}x'\). This also identifies the invalid
branches of the total move definition. The down kernel
construction and the diamond rescaling likewise depend only
on the same underlying decoration. Their preferred
representatives agree by canonical normalization. As in TX11,
freeness equates the rational and auxiliary components of the
moving elements; their frame characters are trivial. Thus
the class maps and the transport maps agree, not merely the
unframed neighbours.

The \(B_w\) assertion at step zero is also integral: (TX10.2)
preserves the pointwise \(\Delta\)-lattice of each equal
\(L^\sharp_{\rm ker}\), and localization preserves this
equivalence in both directions. At the reference twist all
frame maps are identity. The residue-field map is identity
because it agrees on the images of the surjective quotient
\(O_E\to O_E/P\). Every quotient projection, functional,
move, transport, diamond operator and \(B_w\)-condition then
agrees. These arguments prove the full \texttt{TwLvSameS},
including both its step-zero and reference-twist parts.

\paragraph{TX13: the actual twisted terminal row chart.}
The row point datum has unique index and empty coupling.
Its base lattice is \(O_E^{\,n}\), its type-prime set is empty,
and its twist-frame level is the inverse image of the actual
finite row level under reduction at \(P\). It therefore
contains the depth-one principal subgroup, on which its
character is trivial. Apply TX1 with the exact choices
\[
 k=1,\qquad t=1,\qquad d_{\rm row}(P)=q_P.
 \tag{TX13.1}
\]
Thus the row scalar uses the first power of \(q_P\), whereas
the full terminal scalar in TX2 uses power \(c+1\). A
nonuniformizer \(q_P\) may have a larger additive valuation;
its positive valuation still places its congruence subgroup
inside the depth-one reduction kernel.

DR3 realizes each actual preferred row point. Use those
realizers themselves as \(y'_c\); no companion factor occurs
in (TX13.1). Equality of rational local frames gives the
integral discrepancy used in (TX3.3), and the completed level
intersects the rational group in exactly the row frame level.
The same argument as TX4 proves the orbit bijection,
stabilizer identification and lattice pullback. It proves
\texttt{ReadsP} for the resulting \texttt{ClassesAtPrimes}.
The character is the actual auxiliary character times the
inverse product of completed row characters, so it has a
finite exponent. The base lattice agrees with the reference
row base lattice. Outside \(Q_F(a)\) the local groups agree
with the reference row groups; on the same adele identity at
\(Q_F(a)\), the characters agree. This gives the stronger
off-\(Q\) equality, and hence the required auxiliary-place
\texttt{TwPAtGS} equality.

\paragraph{TX14: row values, norms and the conjugate period map.}
The actual row weight has the prescribed common determinant
shift. Since \(t\) is odd, write \(t=2s+1\); then
\[
 m(t+1)/2=m(s+1),\qquad (t+1)/2=s+1.
 \tag{TX14.1}
\]
This identifies the two source shift conventions entry by
entry. Adding the same integer to all weight entries does
not change the polynomial natural weight, the integral
\(\Delta\)-lattice, or its Fock pairing; the determinant
factor is retained in the frame-function map.

Apply CV10 to the actual reading of TX13, with that shift.
The frame map \(e_W\) is identity on underlying frame
functions. Polynomial equivariance plus the chart dictionary
proves its complex image. The rational image is characterized
by values in the prescribed coefficient field in rational
frames. For the integral assertion, choose finitely many
class representatives and use the finite rational-coordinate
description of \(\Delta\). Clearing denominators gives the
global lattice intersection, and localizing that intersection
gives the actual local \(\Delta\)-lattice. The multipliers
are chosen outside the fixed coefficient prime at which one
is localizing: its local membership already gives the required
bound there, while finite prime avoidance, or the Chinese
remainder theorem, supplies an element with the required
orders at the other finitely many bad primes and nonzero
residue at the fixed prime. Taking this multiplier integral
also preserves every already-good place. This proves the
localization inclusion without introducing another bad place.
This is the same two-inclusion argument as CV7--CV10, now
with TX13's row representatives. It proves the literal
\texttt{ValuationAtPrimes} assertion at all places specified
by \texttt{okPT}.

In orthonormal coordinates its pairing is
\[
 \langle z,z'\rangle_P
 =\sum_c |\Gamma_c|^{-1}
     \operatorname{nfAll}_{\mathrm{sg}}(Q_c,P_c),
 \tag{TX14.2}
\]
where \(P_c,Q_c\) are the coefficient polynomials of \(z,z'\).
The order \(Q_c,P_c\) is the source Hermitian convention.
Equation (TX14.1) identifies the shifted pairing exactly.
FC7's rational-basis argument in the actual conjugation-stable
coefficient field therefore proves \texttt{NormPTieG}.
It does not interchange embeddings on arbitrary complex
coefficients.

Finally use coefficient conjugation on the actual polynomial
model to define the semilinear map \(j_M\) of FC8. Polynomial
evaluation is injective; conjugating twice gives the inverse
on the actual row-form image. Conjugating the equivariance
equations gives precisely the row automorphy equations, so
\(j_M\) maps onto that image and only that image. Expanding
the finite class sum with the same weights as (TX14.2) gives
\[
 \operatorname{period}(w,F)=\langle j_MF,w\rangle_P.
 \tag{TX14.3}
\]
This constructs the row value, norm and period clauses for
every twist; it does not replace the row chart by the
reference chart.

\paragraph{TX15: the terminal vector, moving element and evaluation for all twists.}
Fix a row class \(c\). Let \(m_c\) be the actual moving
element taking the preferred representative of its terminal
class \(\widetilde c\) to \(\widetilde F(\operatorname{rep}c)\),
and let \(\gamma_c\) be its rational unitary component.
Take \(z_0=\pi e_{\rm last}\). Then
\(\langle z_0,z_0\rangle=p\), and it spans the prescribed
orthogonal line. The raw lattices are (TX5.2), so the action
of this same moving element gives
\[
 \gamma_c L_{\widetilde c}
      =\iota(N_c)+O_Ez_0.
 \tag{TX15.1}
\]
For \(x\in L_{\widetilde c}^\sharp\) and
\(r=\langle z_0,\gamma_cx\rangle\in O_E\), the functional is
\[
 f_v(\operatorname{pr}_v x)=\operatorname{res}_v(r).
 \tag{TX15.2}
\]
To verify it on the full dual, put \(s=\bar\pi\).
Then \(sx\) lies in the intermediate self-dual lattice.
The raw last-coordinate formula and the projection reading
give \(f_v(\operatorname{pr}_v(sx))=\operatorname{res}_v(sr)\).
Both sides are \(\operatorname{res}_v(s)\) times the
corresponding side of (TX15.2). At the selected type prime
\(s\) is a unit, so cancellation proves (TX15.2).
This is a full integral-domain argument; no equality is
inferred from equality after tensoring with a field.

Let \(y_M\) and \(y'_c\) be the already chosen full and row
representatives, and let \(k\) denote the actual terminal
companion \(kt_a\) of TX5. Put
\[
 x_c=y_M^{-1}\operatorname{rat}(\gamma_c)^{-1}
                 \operatorname{ext}(y'_c),\qquad
 \kappa_c=x_c^{-1}.
 \tag{TX15.3}
\]
TX5 shows that \(\operatorname{ext}(y'_c)k\) realizes
\(\widetilde F(\operatorname{rep}c)\). Pulling back by
\(\gamma_c^{-1}\) leaves a move with rational component
identity and with the auxiliary/frame components of \(m_c\).
The terminal coupling is empty, so TX3 applies and gives
\[
 x_c\in U_M,\qquad
 \chi_M(\kappa_c)^{-1}=\sigma(\chi(m_c)),\qquad
 \operatorname{rat}(\gamma_c)y_M
         =\operatorname{ext}(y'_c)\kappa_c.
 \tag{TX15.4}
\]
This fixes both inverse orientations. Any adelic matrix
extending \(y'_c\) and fixing \(z_0\) is this block extension,
by \(\pi\ne0\); hence the last equality gives the universal
\texttt{til\_\allowbreak ad} clause, not only one chosen instance.

A row stabilizer \(s\) maps to
\(m_c^{-1}\operatorname{ext}(s)m_c\). Conjugating its rational
matrix by \(\gamma_c\) cancels the two moving elements and
leaves \(\operatorname{diag}(\gamma_s,1)\). This proves the
literal \texttt{IsExt} clause for \texttt{stabMap}.
Finally the definition of the column transport evaluates
at the inverse-translated frame and multiplies by
\(\sigma(\chi(m_c))\). With (TX15.4), for every frame
function \(\phi\), and every related pair of frames, it gives
\[
 e_{W,z}(\operatorname{ev}_z\phi)(c,s)
       =\chi_M(\kappa_c)^{-1}\phi(\widetilde c,s'),
 \quad
 \operatorname{IsTr}(\gamma_c)s'=s.
 \tag{TX15.5}
\]
Here \(e_{W,z}\) is identity on the underlying frame functions.
Equation (TX15.5) holds for arbitrary \(\phi\), not just for
an invariant form. Equations (TX15.1)--(TX15.5), together
with TX14, prove the full \texttt{til\_\allowbreak tw},
\texttt{TwPTiedS}, \texttt{TwPEvzS}, and \texttt{StdP}
conclusions for the same charts and moving elements.

\paragraph{TX16: the complete standard all-twist comparison and its exact scope.}
Let \(C\) be the reference \texttt{ColsTiedSym} constructed
in TG12. Its carriers give \texttt{CarH} on its same
Hermitian chains, and its standard level data give the
identity penultimate frames required in TX9. TX10 supplies
\(\ell_0:\texttt{Law0}\); alternatively start with any
supplied \(\ell_0\), whose chosen maps are the same by the
uniqueness proved in TX10. Apply TX11 to these objects,
TX12 to its resulting charts, and TX13--TX15 to its
resulting full terminal companions. Thus every record in
the final join has already been constructed on the same
class sets, lattices, frame maps, residue maps and characters.

To state the literal source endpoint, retain a supplied
\(\texttt{nch}:\texttt{NChoice}(D)\) and
\(\texttt{tp}:\texttt{TupTied}(D,I.\texttt{cf},\texttt{nch})\)
with the literal equality \(\texttt{tp.xi}=I.\texttt{xi}\).
The actual character reading and selected-prime law give
\[
 v\in Q_F(a)
 \ \Longleftrightarrow\
 \exists P\mid v,\ \xi_a\text{ is ramified at }P
 \ \Longleftrightarrow\
 \texttt{tp.InQ'}(a,v).
 \tag{TX16.1}
\]
The forward implication uses the nontrivial local character
at a selected twist prime; the reverse uses the selection
of the prime above every underlying ramification place.
This identifies the prime predicate itself, rather than
merely comparing cardinalities of the two exceptional sets.

Use TX11 for \(\texttt{ch},\texttt{chJ},\texttt{chM},
\texttt{forms}\), TX12 for \(\texttt{lv},\texttt{same}\),
TX9 for \(\texttt{lvJ}\), TX13--TX15 for
\(\texttt{pt},\texttt{ptG},\texttt{ptE}\).
These are exactly the ten fields of
\(\texttt{SymmetricTwistedColumnsAtPrimes}\).
Their standard clauses are TX11's \texttt{StdW},
TX12's \texttt{StdLv}, and TX15's \texttt{StdP}.
This proves
\[
 \exists\,t:
 \texttt{SymmetricTwistedColumnsAtPrimes}
       (\texttt{tp},C,I.\texttt{tw}(\ell,\ell_0,Z)),
 \qquad I.\texttt{StdTw}(\ell,Z,\ell_0,t),
 \tag{TX16.2}
\]
the actual \texttt{colsTwPTiedS\_\allowbreak std} conclusion.

The geometry and comparison through TX15 need neither
\texttt{NChoice} nor \texttt{TupTied}. Their appearance in
(TX16.2) is the exact original arithmetic interface:
\texttt{NChoice} includes finite discrete-valuation-ring
extensions, valuation embeddings, and basis and integral
lattice comparisons for the normalized \(M_t\)-spaces;
\texttt{TupTied} also prescribes the distinguished twists
and the characters at those same selected prime sets.
TL15's finite-idele character readings alone do not supply
these stronger arithmetic data. They remain explicit,
as do the supplied arithmetic host, \texttt{Ok},
\texttt{Pins}, and \texttt{Later}. Conversely no independent
\texttt{Law0}, realizability, chart, common-character,
twisted-chain, terminal-vector or period-comparison premise
has been introduced. The stronger
\texttt{OneLevelSymmetricColumns} assertion is not inferred
from (TX16.2); it has its own additional level restrictions.


\section{The actual vacuum identity and algebraic germs}
\label{sec:vacuum-identity}

This section supplies the vacuum and four-copy inputs retained in
TL18. Fix the literal arithmetic datum \(D\), a curve ideal
\(\mathfrak c\), its prescribed discriminant root \(\delta\), and
an actual \(\mathscr B:\texttt{Base}(D,\mathfrak c,\delta)\).
We assume the source's \(\texttt{Crv}\) conditions. Every period,
half-point, isometry and correction root below is the one in
\(\mathscr B\); no replacement base is chosen.

Write \(K=\mathbb Q(\sqrt{-d})\subset E\), \(A=D.\mathrm{host}.\mathrm{AA0}\),
\(n=|\mathrm{Tau}|\), \(\iota:K\hookrightarrow\mathbb C\), and
\(\operatorname{emb}(b)=(\tau(b))_\tau\). Thus
\(\vartheta=(1+i\sqrt d)/2\), \(d\equiv7\pmod {16}\), and
\[
 \Lambda=\delta\iota(\mathfrak c),\qquad
 a_\Lambda=\operatorname{area}(\Lambda),\qquad
 H(z,w)=a_\Lambda^{-1}\sum_\tau z_\tau\overline{w_\tau},
 \qquad E_H(z,w)=\operatorname{Im}H(z,w).
 \tag{VI.1}
\]
The form is linear in its first variable. The supplied half-point is
\(t=\omega_0/2\); it satisfies the three actual unkilled clauses.
The arithmetic host, including its trace-dual fractional ideal, remains
an input. The one-variable analytic and CM assertions used below are
the reconstructed PA and TA assertions, including TA14.

\paragraph{VI1: the same arithmetic lattices, complex map, and correction.}
The additive subgroup called \(\texttt{lamM}\) is exactly
\[
             \mathcal M=\delta\operatorname{emb}(\mathfrak c A).
 \tag{VI1.1}
\]
Indeed its defining generators are
\(\delta\iota(c)\operatorname{emb}(b)=
\delta\operatorname{emb}(cb)\), with \(c\in\mathfrak c\),
\(b\in A\); their finite sums are the ideal product
\(\mathfrak c A\). The closure in this definition is additive
subgroup generation. It is not a topological closure.

Let \(\gamma_0\) and \(g=\mathscr B.g_{\mathbb C}\) be the supplied
maps. PA4 proves that the vectors
\(\operatorname{emb}_4(x)=(\tau(x_a))_{a,\tau}\) span
\(\mathbb C^{4n}\), and that the unique complex extension of
\(\gamma_0\) has an algebraic matrix and is an isometry for the
sum Hermitian form. Since the supplied \(g\) has exactly that
reading on these spanning vectors, it is this same map. In particular
both \(g\) and \(g^{-1}\) preserve algebraic coordinate vectors.

Put
\[
 \mathcal S=\Lambda^{4n},\qquad
 \mathcal L=g(\mathcal M^4),\qquad
 L_K=\gamma_0(A^4)\subset K^{4n}.
 \tag{VI1.2}
\]
The same reading and \(K\)-linearity give
\(\mathcal L=\delta\iota(\mathfrak c L_K)\), coordinate by coordinate.
The denominator argument of PA10 gives an integer \(N>0\) with
\(NL_K\subset\mathcal O_K^{4n}\) and
\(N\mathcal O_K^{4n}\subset L_K\). Multiplying these inclusions
by \(\mathfrak c\) and applying \(\delta\iota\) proves
\[
                N\mathcal L\subset\mathcal S,\qquad
                N\mathcal S\subset\mathcal L.
 \tag{VI1.3}
\]
Both are full, finitely generated lattices, and both relative
quotients by their intersection are finite. Thus the source
indices \(\texttt{idx}_1,\texttt{idx}_2\) are positive. The first
inclusion also proves that each coordinate of each \(y\in\mathcal L\)
is a rational-lattice point of \(\Lambda\).

Finally PA10 applies to the actual \(\mathscr B.kc\). Its determinant
exponents vanish outside finitely many rational primes, and its
chosen roots satisfy their nonzero algebraic \(2h_K\)-power
equations. Hence the actual correction is
\[
                  \kappa=\prod_p r_p^{-y_p}
                         \in\overline{\mathbb Q}^{\,\times}.
 \tag{VI1.4}
\]
This is a genuinely finite product on the prescribed prime subtype.
It retains the chosen roots, without asserting that different
independent root choices give the same correction.

\paragraph{VI2: the full period lattices are self-dual with the stated multiplier.}
Here is the integral calculation behind the self-duality input.
Let \(L\subset E^N\), \(L'\subset E^{N'}\) be full self-dual
\(\mathcal O_E\)-lattices and put
\[
 P=\operatorname{Hom}_{\mathcal O_E}(L,L')
       \subset M_{N',N}(E).
 \tag{VI2.1}
\]
The pairing on matrices is \(h(T,U)=\operatorname{tr}(TU^*)\).
On their complex embeddings, (VI.1) is extended by summing
over the row and column indices as well as \(\tau\).
Its \(\mathcal O_E\)-dual on \(P\) is \(P\) itself. To check
this, localize at each finite prime. Choose bases for the two
modules and, where conjugation changes the prime, for the modules
at the conjugate prime. Self-duality identifies the corresponding
dual bases. The matrix units are then dual for the trace pairing.
This argument requires perfect integral pairings, not orthonormal
integral bases.

The host's trace-dual identity for \(A\) is equivalently
\(A\bar A=\mathfrak D_{E/K}^{-1}\). The trace-dual ideal formula,
applied to the preceding matrix units, therefore gives
\[
 (AP)^{\#_K}=AP,\qquad
 (\mathfrak c AP)^{\#_K}=\bar{\mathfrak c}^{-1}AP,\qquad
 (\operatorname N\mathfrak c)\,
       (\mathfrak c AP)^{\#_K}=\mathfrak c AP.
 \tag{VI2.2}
\]
Here \(\#_K\) tests \(\operatorname{Tr}_{E/K}h\) for
\(\mathcal O_K\)-valuedness. The last equality uses the actual
quadratic ideal identity
\(\mathfrak c\bar{\mathfrak c}=
(\operatorname N\mathfrak c)\mathcal O_K\).

The source lattice is precisely
\[
 \mathcal A(L,L')=\texttt{La}(L,L')
                 =\delta\operatorname{emb}(\mathfrak c AP).
 \tag{VI2.3}
\]
Expanding its defining generators proves this equality in both
directions. PA1 gives
\(a_\Lambda=|\delta|^2(\operatorname N\mathfrak c)\sqrt d/2\).
For \(b,c\) in the rational matrix space, the resulting alternating
pairing is
\[
 E_H(\delta\operatorname{emb}b,\delta\operatorname{emb}c)
 =\frac{1}{\operatorname N\mathfrak c}
    \operatorname{Tr}_{K/\mathbb Q}
       \left(\frac{\operatorname{Tr}_{E/K}h(b,c)}
                         {\sqrt{-d}}\right).
 \tag{VI2.4}
\]
The codifferent of \(K/\mathbb Q\) is
\((\sqrt{-d})^{-1}\mathcal O_K\). Testing (VI2.4) against all
\(\mathcal O_K\)-multiples of \(c\), and using (VI2.2), proves
that the alternating integral dual of \(\mathcal A(L,L')\)
is itself. This test also handles a vector initially given in
the full real space: in an integral lattice basis its pairings
are integral and its coordinates are obtained by inverting a
nonsingular rational pairing matrix. It therefore lies in the
rational span, where (VI2.2) applies.

These are full lattices. Any two of the modules in (VI2.3) have
the same rational span and admit common integer denominator
bounds, so the resulting period lattices are commensurable.
This proves the self-duality and commensurability clauses of the
actual \(\texttt{LatAll}\). In particular \(\mathcal M\) is
self-dual for \(E_H\); its fourth power is self-dual, and the
isometry \(g\) makes \(\mathcal L\) self-dual. The coordinate
period lattice \(\mathcal S\) is self-dual as well: on a period
basis the alternating matrix has off-diagonal entries \(1,-1\),
up to orientation, and one tests coordinates separately.

Every lattice just described is stable under \(\vartheta\).
For its vectors \(x,y\), write
\(e=E_H(x,y)\in\mathbb Z\) and
\(k=E_H(\vartheta x,y)\in\mathbb Z\). Since
\[
 \sqrt d\,\operatorname{Re}H(x,y)=2k-e,
 \tag{VI2.5}
\]
the quadratic expression
\(q(x)=\sqrt d\,\operatorname{Re}H(x,x)/2\) is an integer,
and
\[
 q(x+y)-q(x)-q(y)-E_H(x,y)\in2\mathbb Z.
 \tag{VI2.6}
\]
Consequently \(\alpha(x)=\exp(\pi i q(x))\) takes values in
\(\{1,-1\}\) and is exactly the required semicharacter:
\[
 \alpha(x+y)=\alpha(x)\alpha(y)\exp(\pi iE_H(x,y)).
 \tag{VI2.7}
\]
On every actual \(\texttt{La}\), the source's
\((-1)^{\operatorname{round}(q(x))}\) equals this \(\alpha(x)\),
because \(q(x)\) is integral. Thus this computation supplies the
literal multiplier, including its rounding convention.

\paragraph{VI3: the chosen half-point supplies the standard character and StdOK.}
For the one-dimensional period lattice let
\[
 \eta(w)=
 \begin{cases}1&w\in2\Lambda,\\-1&w\notin2\Lambda,\end{cases}
 \qquad
 \chi_0(w)=\eta(w)\exp(\pi iE_H(w,\omega_0)).
 \tag{VI3.1}
\]
The unkilled conditions imply that
\(\omega_0,\vartheta\omega_0,\omega_0+\vartheta\omega_0\)
are all nonzero in \(\Lambda/2\Lambda\). Thus the first two
form a basis of this two-dimensional \(\mathbb F_2\)-space.
Its alternating form is nondegenerate, so
\[
 q(\omega_0)=E_H(\vartheta\omega_0,\omega_0)
                         \quad\hbox{is odd}.
 \tag{VI3.2}
\]
Also \(q(\vartheta\omega_0)=
\frac{1+d}{4}q(\omega_0)\) is even, since \(d\equiv7\pmod8\).
It follows directly from (VI3.1) that \(\alpha\) and \(\chi_0\)
have the same values on both basis vectors: respectively \(-1\)
and \(1\). Their product is an ordinary character trivial on
\(2\Lambda\), because their semicharacter factors cancel.
They therefore agree on all of \(\Lambda\).
Taking products over the coordinates gives
\[
       \alpha|_{\mathcal S}=\texttt{stdChar}(L,\omega_0).
 \tag{VI3.3}
\]

We can now verify all ten fields of \(\texttt{StdData}\), hence
the actual \(\mathscr B.\texttt{StdOK}\).
The period and nonzero half-period are supplied by the actual
unkilled point. The isometry is VI1; \(\alpha(0)=1\);
the standard character is (VI3.3); and VI2 proves the image
lattice's semicharacter and self-duality. For every four-block
vector, isometry and the block sum give the exact identity
\[
             \alpha(gz)=\prod_{a=0}^3\alpha(z_a).
 \tag{VI3.4}
\]
The two positive relative indices are VI1. These verifications
prove \(\texttt{stdOK}\) and \(\texttt{stdOK\_\allowbreak of\_\allowbreak crv}\);
they do not leave \(\texttt{StdOK}\) as an independent premise.

\paragraph{VI4: the Heisenberg formula and a well-defined finite average.}
For a Hermitian form linear in its first variable define
\[
 U_xF(z)=F(z+x)
          \exp\!\left(-\pi H(z,x)-\frac{\pi}{2}H(x,x)\right).
 \tag{VI4.1}
\]
Expanding the exponent, with
\(H(x,y)-H(y,x)=2iE_H(x,y)\), gives
\[
 U_xU_y=e^{-\pi iE_H(x,y)}U_{x+y},\qquad
 U_xU_y=e^{-2\pi iE_H(x,y)}U_yU_x.
 \tag{VI4.2}
\]
For a lattice with semicharacter \(\alpha\), put
\(B_x=\alpha(x)U_x\). Equations (VI2.6)--(VI2.7) and integrality
of \(E_H\) show \(B_xB_y=B_{x+y}\) for vectors of that lattice.
In particular the theta condition is precisely \(B_xF=F\).

Let \(\Lambda_1,\Lambda_2\) be commensurable lattices carrying
the same semicharacter, and let \(F\) satisfy the theta condition
for \(\Lambda_1\). If \(y\in\Lambda_2\) and
\(d\in\Lambda_1\cap\Lambda_2\), then
\(B_{y+d}F=B_yB_dF=B_yF\). Thus
\[
 I_{12}F=
 \frac1{\sqrt{r_{12}}}
 \sum_{[y]\in\Lambda_2/(\Lambda_1\cap\Lambda_2)}B_yF,
 \qquad r_{12}=[\Lambda_2:\Lambda_1\cap\Lambda_2],
 \tag{VI4.3}
\]
is independent of every representative choice. The square root
is the positive real square root. Translation by any
\(x\in\Lambda_2\) permutes the summands, proving
\(I_{12}F\) is fixed by \(\Lambda_2\). Every summand is entire
when \(F\) is entire, so the finite sum is entire.
This is exactly the source's \(\texttt{inter}\), whose
\texttt{finsum} is an ordinary finite sum here.

\paragraph{VI5: finite Fourier inversion, nonvanishing, and the positive normalization.}
Suppose both lattices in VI4 are self-dual. The pairing
\[
 (\Lambda_1/(\Lambda_1\cap\Lambda_2))\times
 (\Lambda_2/(\Lambda_1\cap\Lambda_2))
 \longrightarrow\mathbb C^\times,\qquad
 (x,y)\longmapsto e^{-2\pi iE_H(x,y)}
 \tag{VI5.1}
\]
is well-defined and nondegenerate. For example, if its character
is trivial for every \(x\in\Lambda_1\), then
\(y\in\Lambda_1^\#=\Lambda_1\), so its quotient class is zero.
The other variable is identical.

A nontrivial character on a finite group has sum zero: translate
the sum by a point where its value is not \(1\). A trivial
character has sum equal to the group order. Summing (VI5.1)
in both orders consequently proves
\[
                   r_{12}=r_{21}=:r.
 \tag{VI5.2}
\]
Now expand \(I_{21}I_{12}F\). Move \(B_x\), \(x\in\Lambda_1\),
past \(B_y\) by (VI4.2), and use \(B_xF=F\).
The sum over \(x\) vanishes unless \(y\) is the zero quotient
class, where it equals \(r\). The two normalization factors
are \(r^{-1/2}\), so
\[
                         I_{21}I_{12}F=F.
 \tag{VI5.3}
\]
This proves both inversion and nonvanishing on a nonzero
theta vector. It also transports a one-dimensional theta
space to a one-dimensional theta space, in both directions.

For precision, the same positive normalization preserves the
Gaussian-weighted torus mean. On the common finite cover by
\(\Lambda_1\cap\Lambda_2\), use the normalized Haar mean of
\(F\overline G e^{-\pi H(z,z)}\). Directly from (VI4.1),
\[
 |U_xF(z)|^2e^{-\pi H(z,z)}
                 =|F(z+x)|^2e^{-\pi H(z+x,z+x)}.
 \tag{VI5.4}
\]
Thus \(B_x\) is unitary for this mean. The translates \(B_yF\)
belong to distinct character spaces for
\(\Lambda_1/(\Lambda_1\cap\Lambda_2)\), by (VI5.1), and hence
are pairwise orthogonal: applying an element with distinct
eigenvalues to their inner product proves it is zero.
All their squared norms equal that of \(F\). Equation (VI4.3)
therefore preserves the mean squared norm. The normalized mean
of a function periodic for either original lattice is unchanged
on this finite cover. This proves the stated isometry without
introducing an unspecified Fourier scalar or an analytic
nonvanishing assumption. The source's \(\texttt{normSq}\) is
exactly this normalized average over its chosen fundamental
domain. Fullness provides a bounded fundamental domain with
positive finite volume; periodicity makes the average independent
of that choice. Thus the argument gives the literal
\(\texttt{IopIso}\) norm equality.

\paragraph{VI6: the standard product theta space is exactly one line.}
Let \(s_\Lambda\) be the actual normalized odd theta function and
put
\[
 \theta_0(z)=\frac{U_t s_\Lambda(z)}{(U_t s_\Lambda)(0)},
 \qquad t=\omega_0/2.
 \tag{VI6.1}
\]
TA14 and PA15 give the source's function, its entirety and
simple zeros at \(-t+\Lambda\). The denominator is nonzero
because \(t\notin\Lambda\); \(\theta_0(0)=1\). Its multiplier
is precisely (VI3.1). This also identifies the period-pair
and \(\texttt{GL3Theta.th0}\) definitions on the same lattice.
TA12 keeps this normalized function on the prescribed
distinguished half-point class under coefficient conjugacy.
Thus the arithmetic input used later is attached to this
function, rather than to an unspecified generator of its line.

An entire one-variable function \(F\) with this multiplier
is a scalar multiple of \(\theta_0\). Indeed \(F/\theta_0\)
is elliptic and has at most one possible simple pole on the
torus. The sum of residues on a fundamental parallelogram is
zero, by cancellation of opposite edges. That sole possible
residue is therefore zero, so the pole is removable.
The quotient is entire and periodic, hence constant.
Its value at zero is \(F(0)\).

For finitely many coordinates put
\(F_0(z)=\prod_j\theta_0(z_j)\). The exponential in (VI4.1)
splits as a product, so \(F_0\) has the required standard
multiplier and \(F_0(0)=1\). Conversely fix all but one
coordinate of a standard theta function. Its remaining
one-variable slice has the preceding multiplier, so replacing
that coordinate by zero extracts the factor \(\theta_0\).
Repeating over all coordinates proves
\[
               F(z)=F(0)\prod_j\theta_0(z_j).
 \tag{VI6.2}
\]
This proves the actual standard theta-line theorem, together
with its normalized nonzero generator.

\paragraph{VI7: constructing the four-copy root and the actual chosen vacuum.}
Use the lattices and maps in VI1 and the multiplier in VI3.
By VI5--VI6 the function
\[
 K_4(z)=
 \frac1{\sqrt r}
 \sum_{[y]\in\mathcal L/(\mathcal L\cap\mathcal S)}
       \alpha(y)\,(U_yF_0)(gz)
 \tag{VI7.1}
\]
is nonzero. It spans the theta space for \(\mathcal M^4\),
the four-block form, and the product multiplier; the exact
change of variables is \(g\), by (VI3.4).
Set \(P=\kappa K_4\). It too is a nonzero generator.

Choose a point \(z^*=(z_0^*,z_1^*,z_2^*,z_3^*)\) with
\(P(z^*)\ne0\), and define
\(f_1(v)=P(v,z_1^*,z_2^*,z_3^*)\).
The block form and \(\alpha(0)=1\) show directly that \(f_1\)
is in the one-copy theta space. It is nonzero at \(z_0^*\).
Its tensor product \(f_1^{\otimes4}\) is in the four-copy line,
so \(f_1^{\otimes4}=cP\). Here \(c\ne0\): evaluate the left
side with all four variables equal to \(z_0^*\).
Choose \(b\in\mathbb C\) with \(b^4=c^{-1}\). Then
\(f=bf_1\) satisfies
\[
                   f^{\otimes4}=\kappa K_4.
 \tag{VI7.2}
\]
This takes a fourth root of one nonzero complex scalar; no
algebraicity over \(\mathbb Q\) is asserted for that scalar
or its chosen root, and no branch of a pointwise analytic
fourth root has been assumed.

The one-copy theta space is also a line. Given another vector
\(F\) in it, the mixed tensor \(F\otimes f\otimes f\otimes f\)
lies in the four-copy line. Evaluate its last three variables
at a point where \(f\ne0\) to obtain \(F=cf\).
If two such vectors satisfy (VI7.2), their ratio consequently
has fourth power \(1\).

We have proved existence, nonvanishing, spanning and the exact
normalization for every admissible choice. The literal
\(\mathscr B.\texttt{thMM}\) is the choice from this nonempty
predicate, so it satisfies (VI7.2) and is nonzero.
This proves both occurrences of \(\texttt{thMM\_\allowbreak spec}\), as
well as \(\texttt{thMM\_\allowbreak ne\_\allowbreak zero}\), with no independent
vacuum-existence premise.

\paragraph{VI8: the finite four-copy expression is an algebraic germ.}
Write \(f=\mathscr B.\texttt{thMM}\). For the coordinate form,
direct expansion of (VI4.1) gives
\[
                 (U_yF_0)(w)
                    =\prod_j (U_{y_j}\theta_0)(w_j).
 \tag{VI8.1}
\]
Every \(y_j\) appearing in (VI7.1) is a rational-lattice
point by (VI1.3). TA10, PA15 and the same half-point
normalization therefore give all Taylor coefficients of
each \(U_{y_j}\theta_0\) in \(\overline{\mathbb Q}\).
All its derivatives at zero are algebraic as well.
Each coordinate linear form \(z\mapsto(gz)_j\) has algebraic
coefficients by VI1. Thus every factor in (VI8.1), after this
substitution, is a literal \(\texttt{Gen.lin}\) generator for
the algebraic-coordinate directions.

The remaining coefficients in (VI7.1)--(VI7.2) are
\(\kappa\in\overline{\mathbb Q}^{\times}\), signs
\(\alpha(y)\), and \(r^{-1/2}\in\overline{\mathbb Q}^{\times}\).
The quotient is finite, and changing its representatives does
not change the function by VI4. The \(\texttt{Gen}\) rules for
constants, finite sums and products now prove
\[
 \texttt{Gen}\left(z\longmapsto\prod_{a=0}^3 f(z_a)\right).
 \tag{VI8.2}
\]
The function \(f\) is entire by VI7. Hence (VI8.2) is exactly
the source \(\texttt{Root4}(f)\), proving
\(\texttt{thetaVector\_\allowbreak fourthPower\_\allowbreak algebraicGerm}\).
There is no need to infer algebraic coefficients for an
unspecified complex-linear map from rational spanning alone:
VI1 supplied the actual algebraic matrix.

\paragraph{VI9: the same identity handles every rational Heisenberg translate.}
The preceding argument can be strengthened in a way useful for
the later indicator calculation. Let \(x\in\mathbb Q\mathcal M\)
and set \(f_x=U_xf\). Put \(x^{(4)}=(x,x,x,x)\) and
\(v=gx^{(4)}\). Some \(m>0\) has \(mv\in\mathcal L\).
Applying \(U_{x^{(4)}}\) to (VI7.2), using isometry of \(g\)
and the first identity in (VI4.2), gives
\[
 f_x^{\otimes4}(z)=
 \frac{\kappa}{\sqrt r}
 \sum_{[y]\in\mathcal L/(\mathcal L\cap\mathcal S)}
  \alpha(y)e^{-\pi iE_H(v,y)}
       \prod_j (U_{v_j+y_j}\theta_0)((gz)_j).
 \tag{VI9.1}
\]
The sign and the factor \(\pi\) in the phase are those of
the actual Heisenberg multiplication formula.
Since \(\mathcal L\) is integral,
\(mE_H(v,y)=E_H(mv,y)\in\mathbb Z\). Thus every displayed
phase is a root of unity and is algebraic. By (VI1.3),
every \(v_j+y_j\) is a rational-lattice point of \(\Lambda\).
The same TA10 argument as in VI8 proves
\(\texttt{Root4}(f_x)\).

For any multi-derivative \(D^\nu\) in algebraic coordinate
directions, differentiate (VI9.1) by that same operator
separately in each of the four blocks and then set all
variables to zero. TL18's product and chain rules show that
the resulting number is algebraic. Its left side is
\((D^\nu f_x(0))^4\), so \(D^\nu f_x(0)\) is algebraic.
This remains valid when \(f_x(0)=0\); no value of \(f_x\)
has been divided out. We have proved algebraicity of all
Taylor coefficients of every rational translate of this
actual vacuum, without invoking a theorem about arbitrary
theta functions with algebraic coefficients.

\paragraph{VI10: rectangular vacua and the actual moved functions.}
For arbitrary \(N,N'\), the standard matrix lattice satisfies
the exact integral equality
\[
 \mathcal A(\mathcal O_E^N,\mathcal O_E^{N'})
       =\mathcal M^{\,N'\times N}.
 \tag{VI10.1}
\]
An integral homomorphism between the standard modules is a
matrix over \(\mathcal O_E\). Multiplying by \(A\) gives
matrices with entries in \(A\); conversely every such
matrix is the sum of its individual matrix entries.
Applying the defining period generators proves (VI10.1)
in both directions.

The rectangular vacuum is exactly
\[
 f_{N,N'}(z)=\mathscr B.\texttt{thZero}(N,N')(z)
              =\prod_{i,j} f((z_{\tau ij})_\tau).
 \tag{VI10.2}
\]
The Hermitian form and its multiplier split over these
columns, so VI7 proves that it is entire, satisfies the
literal theta condition, and is nonzero. The one-copy
spanning assertion of VI7, applied successively to the
columns, proves that it spans the full standard theta line.
This also covers an empty matrix index set: the product is
the constant \(1\) on the zero-dimensional space.

If \(X\in\mathbb Q\mathcal A(\mathcal O_E^N,\mathcal O_E^{N'})\),
the Heisenberg operator on (VI10.2) is the product of the
operators on its columns. Every column belongs to
\(\mathbb Q\mathcal M\). VI9 therefore proves that all Taylor
coefficients of \(U_Xf_{N,N'}\) are algebraic. At \(X=0\) this
is the literal \(\texttt{thZero\_\allowbreak coefIn}\); division by the
multi-index factorials preserves algebraicity.

For any actual self-dual pair \(L,L'\), VI2 gives finite
commensurability with the standard period lattice and the
same multiplier. Both actual moving operators defining
\(\texttt{thMoved}\) are consequently the positive finite
sums (VI4.3). They preserve the theta condition and do not
kill the vacuum, by VI5. All their translating vectors lie
in the rational span of the standard period lattice.
Their coefficients are signs times inverse square roots
of positive integers.

This finite presentation is stable under a further rational
translation. For \(X,Y\) in the rational span, choose
integers \(m,n>0\) with \(mX,nY\) in the standard integral
period lattice. Then \(mnE_H(X,Y)\in\mathbb Z\), so the
phase in \(U_XU_Y=e^{-\pi iE_H(X,Y)}U_{X+Y}\) is algebraic.
Finite sums and VI9--VI10 now prove the actual assertion
\(\texttt{ta\_\allowbreak thMoved}\): the moved function is entire and
every rational translate has algebraic Taylor coefficients.
All sums here are genuinely finite for these lattices;
no support assertion is inferred from a totalized sum.

\paragraph{VI11: normalized indicator jets and smoothness on the same ambient data.}
Retain the literal seam \(s\), fields \(X\), character \(\chi\),
and constants \(c_D\) of \(\texttt{ambOf}(s,\mathscr B,\chi,c_D)\).
Assume exactly the source's \(X.\texttt{Ok}\), \(s.\texttt{Ok}\),
\(\texttt{Crv}\), \(\texttt{IsChi}\), and the stated even ranks.
The independent finite-presentation assertion CA2 gives, for
each allowed actual indicator,
\[
 \chi_F^{-1}\,\texttt{thInd}(z)
       =\sum_{\ell\in S} r_\ell\,
           U_{y_\ell}\bigl(\mathscr B.\texttt{thMoved}(L,L')\bigr)
                    (\operatorname{reindex}_s z),
 \qquad r_\ell\in\overline{\mathbb Q},
 \quad y_\ell\in\mathbb Q\mathcal A(L,L').
 \tag{VI11.1}
\]
CA2 proves this formula from the actual finite indicator;
it does not assume algebraic jets. VI2 identifies the
rational span on the right with that in VI10. Hence VI10
proves algebraicity of every coefficient of the right-hand
side. The seam reindexing has entries \(0,1\), and its
multi-derivative dictionary only introduces the explicit
integer factorials. Thus (VI11.1) proves the literal
\(\texttt{JetAlg}\) for every finite prime set and every
depth in the source quantifier order.

The ambient entire function is a constant multiple of the
same moved theta function after this coordinate reindexing.
VI10 proves its complex smoothness of every order, giving
the literal \(\texttt{ESmooth}\). Multiplication by a complex
constant preserves smoothness; no algebraicity of the
unnormalized factor \(\chi_F\) is silently assumed.

Consequently \(\texttt{esmooth\_\allowbreak jetAlg}\) is supplied by the
actual vacuum construction, the one-variable PA/TA input,
and CA2's independently proved finite identity. The dependency
direction is VI1--VI10 and CA2 into VI11; none of those
primitive assertions uses the final \(\texttt{FinAlg}\) or
\(\texttt{EvLg}\) construction. The arithmetic host and the
named classical one-variable analytic facts remain explicit;
no full RH or GRH conclusion is asserted here.


\section{The normalized indicators, covariance, and canonical finite jet field}
\label{sec:canonical-analytic}

\status{The claims in this section reconstruct the actual analytic and finite
Fourier identities used in \texttt{finAlg\_\allowbreak step}, and then join them to the
canonical local-index construction on the same data. The finite operator
claims CA1--CA4 do not assume a theta vacuum. The actual vacuum and its
algebraic translates are supplied by VI. The normalization in covariance
is proved, including its determinant character and its inert-prime sign;
it is not deduced from theta-line uniqueness alone. The final construction
retains the literal frame-inclusion and raw \texttt{PreX} inputs, and does
not construct a later geometric stage. All statements are human
reconstructions from the pinned source, without a Lean build or a
kernel-verification claim.}

\paragraph{Conventions.}
Fix the source datum $D$, its CM field $E$, imaginary quadratic subfield
$K=\mathbb Q(\sqrt{-d})$, induced CM type $\Phi$, and a curve ideal
$\mathfrak c$ with discriminant normalization $\delta$. Write
$V=\prod_{\tau\in\Phi}\operatorname{Mat}_{b\times a}(\mathbb C)$ and
$H$ for the source Hermitian form, linear in its first argument:
\[
 H(z,z')=\frac1{\operatorname{area}(\delta\mathfrak c)}
       \sum_{\tau,i,j}z_{\tau ij}\overline{z'_{\tau ij}},\qquad
 E_H(z,z')=\operatorname{Im}H(z,z'),\qquad
 Q(z)=\frac{\sqrt d}{2}H(z,z).
 \tag{CA.1}
\]
The last quantity is real. Rational points mean the image
$X=\delta\operatorname{Emb}_{\Phi}\operatorname{Mat}_{b\times a}(E)$,
not arbitrary complex points. For self-dual Hermitian lattices $L,L'$,
write $\Lambda(L,L')$ for the source period lattice \texttt{La}; under
this rational identification it is
$\mathfrak c\,\mathfrak A\operatorname{Hom}_{\mathcal O_E}(L,L')$,
where $\mathfrak A$ is the source trace-dual ideal. Every assertion about
this identification uses that same ideal and that same $\delta$.
For additive lattices $A,B$, set $[B:A\cap B]=n(A,B)$; this is the
source relative-index convention.

\paragraph{CA1: the exact rational Heisenberg operators.}
Define
\[
 (U_x f)(z)=e^{-\pi H(z,x)-\pi H(x,x)/2}f(z+x),\qquad
 \alpha(x)=e^{\pi iQ(x)},\qquad W_x=\alpha(x)U_x.
 \tag{CA.2}
\]
Expanding $H(z+x,y)$ and $H(x+y,x+y)$, using
$H(y,x)=\overline{H(x,y)}$, gives
\[
 U_xU_y=e^{-\pi iE_H(x,y)}U_{x+y},\qquad
 W_xW_y=e^{-2\pi i c(x,y)}W_{x+y},\qquad
 c(x,y)=\tfrac12\bigl(\sqrt d\,\operatorname{Re}H(x,y)
                         +E_H(x,y)\bigr).
 \tag{CA.3}
\]
These formulas hold on every function, so they also hold termwise in
finite sums and under differentiation of entire functions. In particular
$W_xW_y=e^{-2\pi iE_H(x,y)}W_yW_x$.

The integer $w=(1+\sqrt{-d})/2\in\mathcal O_K\subset\mathcal O_E$
has $w+\bar w=1$; all embeddings in $\Phi$ send it to
$(1+i\sqrt d)/2$. Thus the last cocycle in (CA.3) is
$\gamma(wx,y)$, where
\[
 \gamma(x,y)=e^{-2\pi iE_H(x,y)}.
 \tag{CA.4}
\]
It is alternating and biadditive, and
$\gamma(rx,y)=\gamma(x,\bar r y)$ for $r\in\mathcal O_E$.
The existence and integrality of $w$ are the actual quadratic-order
facts in the host, including the prime $2$; no division by $2$ is
performed in a residue field.

For rational $x,y$, both $E_H(x,y)$ and $\sqrt d\,
\operatorname{Re}H(x,y)$ are rational. Indeed (CA.1) becomes the
trace of an element of $E$ to the induced quadratic field, divided
by the rational ideal norm and by $\sqrt d/2$; its imaginary and
$\sqrt d$ times its real part are rational. Consequently the phases
in (CA.2)--(CA.4) are roots of unity. All their algebraic conjugates
have modulus one. This is the exact rationality used below, not
algebraicity of an arbitrary complex Gaussian factor.

\paragraph{CA2: a finite presentation of the normalized indicator.}
Use the literal \texttt{ambOf} datum obtained from a seam $s$, with
\texttt{Fields.Ok}, \texttt{Seam.Ok}, the above curve data, self-dual
$L,L'$, an arbitrary finite prime set $S$, and an integer $c>0$.
Let $\Lambda=\Lambda(L,L')$, $A_c=\texttt{AC}$,
$B_c=\texttt{BC}$, and choose the actual representative sets
\[
 R_A\subset\Lambda\cap A_c\quad\bmod B_c,
 \qquad R_B\subset B_c\quad\bmod\Lambda.
 \tag{CA.5}
\]
Here ``modulo'' means modulo the intersection when neither lattice
contains the other. Let $s_x$ be the source's chosen rational lift of
$x$, let $v=\texttt{volC}$, and put $k=\texttt{chiF}(L,L')$.
Writing $G=\texttt{thMoved}(L,L')$ in host coordinates, the actual
definition of the indicator expands to
\[
 k^{-1}\operatorname{thInd}_{S,L,L',c,x}
       =\sum_{y\in R_B}r_yU_{y-s_x}G,
 \tag{CA.6}
\]
where
\[
\begin{split}
 r_y={}&(k^{-1}k)\,\frac{v}{(\sqrt v)^2}\,
 \left(\frac1{|R_A|}\sum_{t\in R_A}
       {\bf1}_{B_c}(t)e^{2\pi iE_H(t,y)}\right)\\
 &\hspace{12mm}\cdot\alpha(s_x)\alpha(y)
                      e^{\pi iE_H(s_x,y)}.
\end{split}
 \tag{CA.7}
\]
If a degenerate definition gives $v\le0$, the two inverse square
roots in the source make their combined factor zero; in the present
full-lattice situation $v>0$ and the displayed factor is one. Likewise
$k^{-1}k$ is zero or one, with no algebraicity assumption on $k$.
For an actual character satisfying \texttt{IsChi}, it is one.

To check the normalization, expand the source finite Fourier
transform \texttt{hatS}, then its two factors $(\sqrt v)^{-1}$.
The translation identity
\[
 (U_yG)(z-s_x)e^{-\pi H(z,-s_x)-\pi H(-s_x,-s_x)/2}
       =e^{\pi iE_H(s_x,y)}(U_{y-s_x}G)(z)
\]
is CA1. It gives precisely (CA.7); in particular the inverse character
has canceled before any coefficient is called algebraic.

Every $r_y$ is algebraic and every algebraic embedding sends it to a
number of modulus at most one: the inner expression is an average
of zeros and roots of unity, and the remaining nonzero factors have
modulus one under every embedding. Every $y-s_x$ is rational with
respect to $\Lambda$.

The representative sets exist and satisfy the concrete estimate
\[
 |R_B|\le c^{[E:\mathbb Q]ab}.
 \tag{CA.8}
\]
For completeness, the local exponents defining either half lattice
lie between $-1$ and $1$. Hence $c\Lambda\subset B_c$ and
$cB_c\subset\Lambda$. The map $y\mapsto cy$ injects
$B_c/(B_c\cap\Lambda)$ into $\Lambda/c\Lambda$. A full rational
lattice has a $\mathbb Z$-basis of size $[E:\mathbb Q]ab$, so the
latter quotient has $c^{[E:\mathbb Q]ab}$ elements. The same
inclusions prove finiteness for $R_A$. This argument also covers
$a=0$ or $b=0$, when the quotients have one element.

Thus (CA.6) is the literal source predicate \texttt{Pres}, with the
sum of the absolute values of the conjugated coefficients bounded
by (CA.8). It holds for every $S$ here; the stronger no-conjugate-pair
hypothesis needed for \texttt{LaHyp} and covariance is not inserted
into this finite-presentation assertion.

\paragraph{CA3: the finite Hermitian norm Gauss sum, including characteristic two.}
Let $k/k_0$ be a finite field with nontrivial involution $\sigma$,
$|k_0|=q$, and let $\eta:k_0\to\mathbb C^\times$ be a nontrivial
additive character. Then
\[
 \sum_{t\in k}\eta(t\sigma(t))=-q=-\sqrt{|k|}.
 \tag{CA.9}
\]
Indeed $k^\times$ is cyclic and the norm $t\mapsto t^{q+1}$ maps
onto $k_0^\times$, with every nonzero fiber of size $q+1$.
Translation by a point at which $\eta\ne1$ shows
$\sum_{u\in k_0}\eta(u)=0$. The left side is consequently
$1+(q+1)(-1)=-q$.

The precise form needed for the Heisenberg cocycle is as follows.
Let $\chi:k\to\mathbb C^\times$ be nontrivial, additive, and
$\sigma$-invariant, and suppose $w+\sigma(w)=1$.
Then $\chi=\eta\circ\operatorname{Tr}_{k/k_0}$ for a nontrivial
$\eta$, and
\[
 \sum_{t\in k}\chi(wt\sigma(t))=-\sqrt{|k|}.
 \tag{CA.10}
\]
The kernel of the trace is the image of $\sigma-1$, on which
$\chi$ is trivial. The trace is onto: $X+X^q$ is a nonzero polynomial
of degree $q<q^2$, even in characteristic two. This proves the
factorization and shows its character is nontrivial. Since
$t\sigma(t)\in k_0$ and $\operatorname{Tr}(w)=1$, (CA.10) is
(CA.9). There is no missing dyadic exception.

\paragraph{CA4: finite intertwiners with their correct signs compose exactly.}
Let $R=\mathcal O_E$ act on an additive group $X$, let $W_x$ be
linear operators satisfying (CA.3) in the form
$W_xW_y=\gamma(wx,y)W_{x+y}$, and let $A,B,C$ be mutually
commensurable self-dual $R$-submodules for $\gamma$. On an $A$-fixed
vector define
\[
 S_{A,B}v=\sum_{x\in B/(A\cap B)}W_xv,\qquad
 J_{A,B}v=\lambda(n(A,B))^{-1}S_{A,B}v.
 \tag{CA.11}
\]
The representatives do not matter: changing $x$ by an element of
$A\cap B$ has phase one. The vector $S_{A,B}v$ is $B$-fixed.

Choose a multiplicative $\lambda:\mathbb N_{>0}\to\mathbb C^\times$
such that, at every residue field $R/P$,
\[
 \lambda(|R/P|)^2=|R/P|,
 \qquad \bar P=P\Longrightarrow
              \lambda(|R/P|)=-\sqrt{|R/P|}.
 \tag{CA.12}
\]
The involution on a fixed residue field must be nontrivial. These
conditions are exactly \texttt{GoodNorm}. For the actual host they
hold: $E/\mathbb Q$ is abelian and $E/E^+$ is unramified. Over a
rational prime $p$ all primes have the same residue degree $f_p$,
and conjugation fixes either all of them or none, since it commutes
with every Galois automorphism. Set
\[
 \lambda(p)=\begin{cases}
 \sqrt p\,e^{\pi i/f_p},&\text{conjugation fixes the primes over }p,\\
 \sqrt p,&\text{otherwise},
 \end{cases}
\]
and extend multiplicatively. The fixed residue involution is
nontrivial by the actual unramified CM condition, as checked in
DR1/AH. This proves (CA.12) without choosing inconsistent square
roots at different primes.

With this normalization,
\[
 J_{B,C}J_{A,B}v=J_{A,C}v.
 \tag{CA.13}
\]
Here is the finite proof. Self-duality gives perfect pairings between
the two finite quotients by each intersection. Averaging a nontrivial
character gives zero, so
$S_{B,A}S_{A,B}v=n(A,B)v$. If $A,B$ are adjacent, write
$B/(A\cap B)=R/P$ and $q_P=|R/P|$. There are three cases.
If $B=(A\cap B)+(B\cap C)$, direct representatives give
$n(A,C)=q_Pn(B,C)$ and $S_{B,C}S_{A,B}=S_{A,C}$.
If $A=(A\cap B)+(A\cap C)$, the reverse calculation gives
$n(B,C)=q_Pn(A,C)$ and $S_{B,C}S_{A,B}=q_PS_{A,C}$.
The factors in (CA.12) cancel in both cases.

Otherwise $A\cap C=B\cap C$. Let $x$ generate
$B/(A\cap B)$. The perfect pairing supplies $x=a+c_0$ with
$a\in A,c_0\in C$. Put $\chi(t)=\gamma(a,tc_0)$. It descends to
a nontrivial additive character on $R/P$; isotropy of $B$ gives
$\chi(\bar t)=\chi(t)$. Nontriviality forces $\bar P=P$:
if some $r\in P$ had $\bar r\notin P$, the inverse of $\bar r$
modulo $P$ and $\chi(\bar t)=\chi(t)$ would make $\chi$ trivial.
Expanding $W_{tx}v$ by the cocycle, then using the $C$-invariance of
$S_{A,C}v$, gives
\[
 S_{B,C}S_{A,B}v
   =\left(\sum_{t\in R/P}\chi(wt\bar t)\right)S_{A,C}v
   =\lambda(q_P)S_{A,C}v.
\]
Here $n(B,C)=n(A,C)$, so (CA.13) again follows. This is the case
where replacing every signed square root by a positive square root
would lose a sign.

Finally reduce general $A,B$ to adjacency. Choose a simple nonzero
$R$-submodule $Rx$ of $B/(A\cap B)$ and set
$A_1=(A\cap(Rx)^\perp)+Rx$. The perfect pairing proves
$A_1=A_1^\perp$; its quotient over $A\cap A_1$ is $R/P$ and
$A_1=(A_1\cap A)+(A_1\cap B)$. Therefore
$n(A,B)=|R/P|n(A_1,B)$, which strictly decreases the positive
integer index. Induction, the adjacent calculation, and the
intermediate-lattice sum formula establish (CA.13) in finitely many
steps. The argument uses finite modules; it makes no assumption
that a global ideal or lattice is free.

\paragraph{CA5: even rectangular ranks identify the signed and actual intertwiners.}
Use VI2's full, self-dual period lattices in the rational space $X$.
The action of $R=\mathcal O_E$ is entrywise through the embeddings;
CA1 gives the adjoint identity and the integral element $w$ needed
by CA4. On a period lattice $A$, $Q(x)$ is integral and
$\alpha(x)=(-1)^{Q(x)}=\texttt{al}(x)$. Thus its theta functions
are exactly the entire $A$-fixed vectors for $W$, and the source
operator is
\[
 I_{A,B}=n(A,B)^{-1/2}S_{A,B}.
 \tag{CA.14}
\]

Suppose $a,b$ are even. If only $L$ changes in $\Lambda(L,L')$,
then at a prime $P$ the length of the quotient by the intersection
is a multiple of $b$. To verify this, choose an
$\mathcal O_{E,P}$-basis of $L'_P$ and diagonalize the two source
lattices over that discrete valuation ring. The Hom lattice is a
direct sum of $b$ copies of the corresponding source dual lattice.
Tensoring by the common locally principal ideal
$\mathfrak c\mathfrak A$ does not alter its relative index.
If only $L'$ changes, the analogous length is a multiple of $a$.
These are local bases; no global basis is required.

The only nonpositive square-root choices in CA4 occur over the
conjugation-fixed primes. Their signs are therefore
$(-1)^{b\ell_P}$ or $(-1)^{a\ell_P}$, both one. At every other
prime the normalization is already positive. Hence
\[
 J_{\Lambda(L,L'),\Lambda(M,L')}=I_{\Lambda(L,L'),\Lambda(M,L')},
 \qquad
 J_{\Lambda(L,L'),\Lambda(L,M')}=I_{\Lambda(L,L'),\Lambda(L,M')}.
 \tag{CA.15}
\]
In particular, if $A_0=\Lambda(\mathcal O_E^a,\mathcal O_E^b)$
and $\theta_0=\texttt{thZero}$, then the actual definition satisfies
\[
 \texttt{thMoved}(L,L')
   =I_{\Lambda(L,\mathcal O_E^b),\Lambda(L,L')}
      I_{A_0,\Lambda(L,\mathcal O_E^b)}\theta_0
   =J_{A_0,\Lambda(L,L')}\theta_0.
 \tag{CA.16}
\]
CA4 now proves both one-variable transport laws and independence of
the order of these two moves. The equality does not claim that
every positive-normalized intertwiner between arbitrary lattices
composes without a sign.

Rational unitary changes of coordinates commute with these $J$
operators in their natural form. If $T$ preserves $H$, $Q$, and
$X$, then it carries representatives bijectively, preserves the
index, and
$P_TW_x=W_{Tx}P_T$, where $P_Tf(z)=f(T^{-1}z)$. Therefore
\[
 P_TJ_{A,B}=J_{TA,TB}P_T.
 \tag{CA.17}
\]
This is a direct finite-sum identity, independent of a separately
constructed adelic Weil representation.

\paragraph{CA6: computing the scalar by a Gaussian trace.}
We give the scalar calculation used for the actual rational action.
Let $A$ be a full self-dual period lattice, let $u$ be unitary for
$H$, put $a_0=u^{-1}$, and suppose $\det_{\mathbb C}(1-a_0)\ne0$.
Assume $A$ and $uA$ are commensurable and choose $m>0$ such that
\[
 mA\subset(1-a_0)(A\cap a_0^{-1}A).
 \tag{CA.18}
\]
Rationality of $u$ supplies such an $m$ by clearing denominators.
Put $C=\tfrac12(1+a_0)(1-a_0)^{-1}$. For nonzero
$\theta\in\operatorname{Th}(A)$, theta-line uniqueness and the
positive-Iop isometry (VI5) give
\[
 P_u\theta=c_u I_{A,uA}\theta,\qquad |c_u|=1.
\]
The additional information needed to determine its phase is
\[
 c_u=r\,\det_{\mathbb C}(1-a_0)^{-1}
       \sum_{x\in A/mA}\alpha(x)e^{-\pi H(Cx,x)},
 \qquad r>0.
 \tag{CA.19}
\]

Here is a derivation, including the analytic step. Define the kernel
\[
 K_A(z,t)=\sum_{x\in A}\alpha(x)
  \exp\!\left(\pi H(t+x,z)-\pi H(t,x)-\frac\pi2H(x,x)\right).
 \tag{CA.20}
\]
Positive definiteness bounds each summand, uniformly for $z,t$ in
a compact set, by $C_1e^{-C_2\|x\|^2+C_3\|x\|}$ with $C_2>0$.
The number of lattice points in a radius-$R$ ball is $O((1+R)^{2\dim V})$,
by disjoint small balls around lattice points. The bound is summable,
also after every fixed derivative. Thus the kernel is entire in $t$,
and shifting $x$ proves that $K_A(z,\cdot)$ lies in the theta line.
The scalar equation for $\theta$ therefore holds for this kernel.

Take a fundamental domain $F$ for
$A_m=(1-a_0)^{-1}(mA)\subset A\cap a_0^{-1}A$.
For fixed $z$, the kernel equation is
$K_A(z,a_0t)=c_u n(A,uA)^{-1/2}
\sum_y\alpha(y)U_y(K_A(z,\cdot))(t)$, with $y$ representing
$uA/(A\cap uA)$. Set $t=z$ and multiply by
$e^{-\pi H(z,z)}$. For each $y$, expansion of (CA.20) gives
\[
 U_y(K_A(z,\cdot))(z)e^{-\pi H(z,z)}
 =\sum_{x\in A}b(x,y)e^{2\pi iE_H(x+y,z)},
\]
where
$b(x,y)=\alpha(x)\exp(-\pi H(x,y)-\pi H(x,x)/2
-\pi H(y,y)/2+2\pi iE_H(x,y))$.
Each character on the right is periodic modulo $A_m$, because
$x\in A$, $y\in uA$, and $A_m\subset A\cap uA$.
Its integral on $F$ vanishes unless $x+y=0$: a nonzero real
linear functional $E_H(x+y,\cdot)$ gives a nontrivial torus
character, whose Haar integral is zero by translation.
If $y\notin A$ there is no such $x$. If $y\in A$, the single
remaining coefficient is $b(-y,y)=\alpha(-y)$; multiplying by
the outer $\alpha(y)$ gives one. Thus only the zero class of
$uA/(A\cap uA)$ remains, with integral $\operatorname{vol}(F)$.
The finite-Iop side is exactly
$c_u n(A,uA)^{-1/2}\operatorname{vol}(F)$, whereas the other
side is the integral of $K_A(z,a_0z)e^{-\pi H(z,z)}$.

On the other side write each $x\in A$ uniquely as
$s-(1-a_0)y$, with $s\in A/mA$ and $y\in A_m$.
The semicharacter identity converts this shift into $z\mapsto z+y$,
so the integral unfolds to a finite sum of integrals over $V$.
Set $B=(1-a_0)^{-1}$. The integral for $s$ is
\[
 \frac{\kappa}{\det_{\mathbb C}(1-a_0)}
       \alpha(s)e^{-\pi i\operatorname{Im}H(Bs,s)},\qquad \kappa>0.
 \tag{CA.21}
\]
To check (CA.21), diagonalize the unitary $a_0$ in an $H$-orthonormal
basis. Each eigenvalue $\zeta\ne1$ has $\operatorname{Re}(1-\zeta)>0$.
The elementary complex Gaussian integral with linear terms is
\[
 \int_{\mathbb C}e^{-\pi h(1-\zeta)|z|^2+\pi(pz+q\bar z)}\,dz
   =\frac1{h(1-\zeta)}
      e^{\pi pq/(h(1-\zeta))}.
\]
It follows from the two real Gaussian integrals, first for real
positive coefficients and then by holomorphic continuation in the
right half-plane; the dominating Gaussian on compact parameter
sets justifies that continuation and differentiation. Multiplying
the one-coordinate formulas and substituting
$(1-a_0)Bs=s$ leaves precisely the phase in (CA.21).
Equivalently, the remaining exponent identity is
$-H(s,s)/2-H(Bs,a_0^{-1}s)=-i\operatorname{Im}H(Bs,s)$.
Since $C=B-\tfrac12$ and $C$ is skew-adjoint,
$H(Cs,s)=i\operatorname{Im}H(Bs,s)$. Dividing by the first
positive volume factor proves (CA.19). All unfolding interchanges
are justified by the Gaussian bound above and the finite number
of cosets.

\paragraph{CA7: the finite Gauss factor is nonnegative in the even-rank action.}
Now $A=A_0$ is the standard rectangular period lattice and $u$
acts on just one matrix factor. The other factor has even dimension.
Let $P$ project onto its first half and let $S$ interchange the two
halves. They preserve $A$, commute with $a_0,C$, and satisfy
$P^2=P$, $S^2=1$, $SP=(1-P)S$. Both halves are $H$-orthogonal.
For
\[
 \varphi(x)=\alpha(x)e^{-\pi H(Cx,x)}
           =\exp\!\left(\pi i\left(Q(x)-\operatorname{Im}H(Cx,x)\right)\right),
 \tag{CA.22}
\]
the semicharacter identity, (CA.18), and integrality of $E_H$ on
$A$ show $\varphi(x+mz)=\varphi(x)$. These same identities give
$\varphi(x)=\varphi(Px)\varphi((1-P)x)$ and
$\varphi(Sx)=\varphi(x)$.

There is an integral $\ell\in\mathcal O_E$ with
\[
 \ell\bar\ell=-1+4m b_0,\qquad b_0\in\mathcal O_{E^+}.
 \tag{CA.23}
\]
At a split local pair choose components $1,-1$; at a nonsplit
unramified place, DR1's norm surjectivity gives norm $-1$ modulo
$4m$. The Chinese remainder theorem combines these finitely many
choices into $\ell$. The quotient $b_0$ is integral and fixed by
conjugation. Thus (CA.23) includes all primes dividing $2m$.

Let $L_0$ be entrywise multiplication by $\ell$, and $B_0$ by $b_0$.
They preserve $A$ and commute with $P,a_0,C$; moreover
\[
 H(L_0x,L_0y)=-H(x,y)+4mH(B_0x,y).
\]
It follows that $L_0$ induces a bijection of $A/mA$: if
$L_0x\in mA$, the displayed pairing identity makes
$E_H(x,y)$ divisible by $m$ for every $y\in A$, and self-duality
gives $x\in mA$. A direct exponent check gives
$\varphi(L_0x)=\overline{\varphi(x)}$. In detail, write
$mx=y-a_0y$ as in (CA.18), so $2mCx=y+a_0y\in A$.
The difference between the two exponents, divided by $\pi i$, is
\[
 4m\left(\tfrac{\sqrt d}{2}\operatorname{Re}H(B_0x,x)
                 -\operatorname{Im}H(B_0Cx,x)\right)\in2\mathbb Z.
\]
The first term is $2m$ times the integer
$Q(B_0x+x)-Q(B_0x)-Q(x)$; the second is twice the integral
pairing of $B_0(y+a_0y)$ with $x$.

The finite quotient $A/mA$ is the direct sum of the image and
kernel of $P$. Its two Gauss sums are equal by $S$. The first is
real because the bijection $L_0$ conjugates its summands and
preserves that half. Therefore their product, the sum in (CA.19),
is the square of a real number and is nonnegative. It cannot be
zero, since $|c_u|=1$. Equation (CA.19) consequently gives
\[
 c_u=\frac{|\det_{\mathbb C}(1-u^{-1})|}
               {\det_{\mathbb C}(1-u^{-1})}.
 \tag{CA.24}
\]
This positivity calculation is the step that fixes the phase; a
unit-modulus assertion alone would not give (CA.24).

\paragraph{CA8: the exact determinant powers for every rational unitary matrix.}
For $\gamma\in U(a,E)$ and $\gamma'\in U(b,E)$ put
\[
 (T_{\gamma,\gamma'}z)_\tau
       =\tau(\gamma')z_\tau\tau(\gamma)^{-1},\qquad
 P_{\gamma,\gamma'}f(z)=f(T_{\gamma,\gamma'}^{-1}z).
 \tag{CA.25}
\]
This is exactly the source \texttt{pull}: the inverse on the left
is the conjugate transpose, whereas the factor on the right is
$\gamma$. It maps $\Lambda(L,L')$ to
$\Lambda(\gamma L,\gamma'L')$.

First assume $\det(1-\gamma)\ne0$. For the action on the source
factor, $\det_{\mathbb C}(1-T_{\gamma,1}^{-1})
=\prod_\tau\tau(\det(1-\gamma))^b$. Since $a$ is even,
\[
 \det\gamma\,\overline{\det(1-\gamma)}=\det(1-\gamma).
\]
If $b=2r$ this identity gives
$|w^{2r}|/w^{2r}=\tau(\det\gamma)^{-r}$ for
$w=\tau(\det(1-\gamma))$. Equations (CA.15) and (CA.24) prove
\[
 P_{\gamma,1}\theta_0
   =\left(\prod_\tau\tau(\det\gamma)^{-b/2}\right)
       J_{A_0,T_{\gamma,1}A_0}\theta_0.
 \tag{CA.26}
\]
For the target factor use $\gamma'^{-1}=\bar\gamma'^{t}$.
The conjugate determinant identity has the opposite phase, yielding
\[
 P_{1,\gamma'}\theta_0
   =\left(\prod_\tau\tau(\det\gamma')^{a/2}\right)
       J_{A_0,T_{1,\gamma'}A_0}\theta_0.
 \tag{CA.27}
\]
If one rectangular dimension is zero, these are identities on
constant functions with empty products; the Gaussian argument is
unnecessary.

The restriction excluding eigenvalue one can be removed exactly.
There are infinitely many rational norm-one elements
$z=(1+t\sqrt{-d})/(1-t\sqrt{-d})$, $t\in\mathbb Q$.
Choose $z\ne1$ outside the finitely many inverse eigenvalues of
$\gamma$. Then $\gamma=(z\gamma)(z^{-1}1)$ and each factor has
$1-$factor invertible. Apply (CA.26) to each factor, use
(CA.17), then compose the two $J$ operators by CA4. Determinants
multiply, so (CA.26) holds for every $\gamma$. The same argument
applies to (CA.27). Combining the two commuting actions gives
\[
 P_{\gamma,\gamma'}\theta_0
    =\rho(\gamma,\gamma')J_{A_0,T_{\gamma,\gamma'}A_0}\theta_0,
 \quad
 \rho(\gamma,\gamma')=
       \prod_\tau\tau(\det\gamma)^{-b/2}
                       \tau(\det\gamma')^{a/2}.
 \tag{CA.28}
\]

\paragraph{CA9: all three equivariance laws for the actual normalized family.}
Let $\chi$ satisfy the literal \texttt{IsChi} laws. For self-dual
$L,L'$ define the source scalar and actual family
\[
 d_\chi(L,L')=
  \chi(\operatorname{detId}(\mathcal O_E^a,L))^{b/2}
  \chi(\operatorname{detId}(\mathcal O_E^b,L'))^{-a/2},\qquad
 e(L,L')=d_\chi(L,L')\,\texttt{thMoved}(L,L').
 \tag{CA.29}
\]
These determinant ideals have norm one. Multiplicativity and the
principal norm-one law imply that their character values are
nonzero, and
\[
 d_\chi(\gamma L,\gamma'L')=
       \rho(\gamma,\gamma')d_\chi(L,L').
 \tag{CA.30}
\]
Indeed the determinant ideal of an image is the old ideal times
$(\det\gamma)$, and \texttt{IsChi} assigns to this principal ideal
$\prod_\tau\tau(\det\gamma)^{-1}$. This determines both signs of
the exponents in (CA.30).

Using the actual equality (CA.16), naturality (CA.17), and (CA.28),
\[
\begin{split}
 P_{\gamma,\gamma'}e(L,L')
  &=d_\chi(L,L')J_{TA_0,T\Lambda(L,L')}
                        P_{\gamma,\gamma'}\theta_0\\
  &=\rho(\gamma,\gamma')d_\chi(L,L')
                        J_{A_0,T\Lambda(L,L')}\theta_0
   =e(\gamma L,\gamma'L').
\end{split}
 \tag{CA.31}
\]
For a change in just one lattice, (CA.13), (CA.15), and
$\operatorname{detId}(L,M)\operatorname{detId}(M,N)=\operatorname{detId}(L,N)$ similarly give
\[
\begin{split}
 e(L,L'')&=\chi(\operatorname{detId}(L',L''))^{-a/2}
            I_{\Lambda(L,L'),\Lambda(L,L'')}e(L,L'),\\
 e(M,L')&=\chi(\operatorname{detId}(L,M))^{b/2}
            I_{\Lambda(L,L'),\Lambda(M,L')}e(L,L').
\end{split}
 \tag{CA.32}
\]
Equations (CA.31)--(CA.32) are all three clauses of
\texttt{Base.EqP}, hence of the actual
\texttt{normalisedFamily\_\allowbreak satisfies\_\allowbreak equivarianceLaws} theorem.
The construction starts with the supplied Base's own
$\theta_0=\texttt{thZero}$, furnished by VI10. It neither chooses
a different family nor assumes the sought covariance of this one.

\paragraph{CA10: lift independence and covariance of the actual indicator.}
Assume the literal \texttt{LaHyp}, $\mathfrak A\ne0$, and $M>0$.
Thus $S$ contains no conjugate pair, the period ideal is a unit at
each $P\in S$, and $e(L,L')$ is fixed by the period lattice.
Use rational coordinates and put
$\mathcal P=\mathfrak c\mathfrak A\operatorname{Hom}(L,L')$.
For $c>0$ the two half lattices have the following local readings:
\[
\begin{array}{c|ccc}
 &P\in S&\bar P\in S&P,\bar P\notin S\\ \hline
 (A_c)_P&c^{-1}\mathcal P_P&c\mathcal P_P&\mathcal P_P\\
 (B_c)_P&c\mathcal P_P&c^{-1}\mathcal P_P&\mathcal P_P.
\end{array}
 \tag{CA.33}
\]
These localizations determine the global modules by LP1. Opposite
exponents cancel in the conjugate pairing, so both half lattices
are integral for the cocycle $c(x,y)$ of (CA.3). This follows
also by applying the integral pairing to $wx,y$, since $w$ is
integral and preserves each lattice.

Let $x\in M^{-1}\operatorname{Hom}_S(L,L')$. There exists
$t\in A_{Mc}$ with $t-x\in c\operatorname{Hom}_S(L,L')$.
To see this without an extra approximation hypothesis, clear a
common denominator for $x$ and an integral basis of the full
period module. Chinese remainders in this finitely generated
module give $y$ congruent to $(Mc)^{-1}x$ modulo $\mathcal P_P$
at the finitely many $P\in S$, and integral in $\mathcal P_Q$
at every other prime. Equivalently one may perform the CRT in
a common integral denominator module; it imposes only finitely
many residue conditions. Then $t=Mc\,y$ has the stated properties
by the three cases of (CA.33), using
$\mathcal P_P=\operatorname{Hom}(L_P,L'_P)$ on $S$.

Take any such $t$ and any representatives
$R\subset B_c$ modulo $B_c\cap\mathcal P$. Put
\[
 n_c=[\mathcal P\cap A_c:B_c\cap\mathcal P\cap A_c]>0.
\]
The literal finite Fourier definition reduces to
\[
 \operatorname{thInd}_{S,L,L',c,x}
    =\frac1{n_c}W_{-t}\sum_{y\in R}W_y e(L,L').
 \tag{CA.34}
\]
Indeed its inner representative set is
$(\mathcal P\cap A_c)/(B_c\cap\mathcal P\cap A_c)$.
The indicator of $B_c$ selects only its zero class. At that
class its Fourier phase is one, since $y\in B_c$ and the
cocycle is integral on $B_c$. Thus the Fourier coefficient is
$v/(n_c\sqrt v)$; the additional source factor $1/\sqrt v$
cancels $v>0$, giving exactly $1/n_c$.

The resulting formula is independent of the chosen lift and
representatives. The sum $F=\sum_{y\in R}W_y e$ is $B_c$-fixed:
translation permutes its quotient representatives and the period
and cocycle phases are one. Let $s$ be the source's original
lift in $A_c$. At $P\in S$, $t-s\in c\mathcal P_P$; at the
conjugate primes $s\in c\mathcal P_P$ and
$t\in Mc\mathcal P_P$; elsewhere both are integral.
Consequently $t\in A_c$ and $t-s\in B_c$.
The integral cocycle on $A_c$ gives
$W_{-t}W_{t-s}=W_{-s}$, so $W_{-t}F=W_{-s}F$.
This proves (CA.34) for the actual chosen-lift definition, not
just for a conveniently redefined indicator.

Now let $\gamma,\gamma'$ be rational unitary matrices.
The rational map $\phi(x)=\gamma'x\gamma^{-1}$ carries
$\mathcal P$, $A_c$, $B_c$, $M^{-1}\operatorname{Hom}_S$,
and $c\operatorname{Hom}_S$ to the corresponding modules for
$\gamma L,\gamma'L'$. This follows directly from the defining
mapping property and from localization; scalar denominators
commute with $\phi$. It preserves $n_c$, sends $t,R$ to valid
lift and representative data, and its complex map satisfies
$P_TW_y=W_{\phi(y)}P_T$. Applying CA9 to (CA.34) proves
\[
 \operatorname{thInd}_{S,\gamma L,\gamma'L',c,\phi(x)}(z)
       =\operatorname{thInd}_{S,L,L',c,x}(T^{-1}z).
 \tag{CA.35}
\]
This is exactly \texttt{ThIndEq}, with the source inverse on
the target factor. Equation (CA.30) is separately the literal
\texttt{ChiLaw}. Neither statement discards the determinant
factor before normalizing the jet.

\paragraph{CA11: the actual step has a finite algebraic jet field.}
Let $B$ be the given base, $X:\texttt{PreX}(B)$, and
$I:\texttt{ColIn}(B,X)$ satisfy its literal \texttt{Ok}.
Fix an index $a_0$, a word $u$, arbitrary steps $i,j$, and
$M>0,n\ge0$. Then
\[
 \operatorname{FinAlg}\!\left(
  \operatorname{jetSet}(I.A,N_i,N_j,I.S_b(a_0,u),M,n)\right).
 \tag{CA.36}
\]
The source also supplies \texttt{HasLetters}; we retain that
antecedent when matching \texttt{finAlg\_\allowbreak step}, although the
analytic proof below does not need an auxiliary stage.

First verify its actual period hypotheses. Every $N_i$ is even.
The selected set $S=I.S_b(a_0,u)$ is the union of the type-prime
set and the auxiliary/tuple set. A type prime and its conjugate
cannot both occur, because the two factors of its prescribed
norm element are coprime. The same applies to each auxiliary
factor. Tuple primes are chosen on only one side of a split
pair; primes of different supported rational types cannot collide.
This proves that $S$ has no conjugate pair. Their exact local
conditions say $q_P=\operatorname N P$ is prime and
$q_P\notin P^2$: for the prescribed norm elements this follows
by factoring their product and using the degree and norm count;
for tuple primes it is the literal split, degree-one, unramified
condition. These are the cases in \texttt{bad\_\allowbreak primes}.

At such a $P$ the relative different of $E/K$ is a unit.
In fact $q_P\notin P^2$ says that the ramification index over
$\mathbb Q_{q_P}$ is one; the residue extension is separable,
so the completed integral extension is unramified and its trace
pairing is perfect. The same holds over the intermediate field
$K$. The actual integral half-different $\mathfrak b$ satisfies
$\mathfrak b\bar{\mathfrak b}=\mathfrak D_{E/K}$.
Its two nonnegative local valuations must therefore be zero.
Thus $\mathfrak A=\mathfrak b^{-1}$ localizes to
$\mathcal O_{E,P}$. Since the curve ideal here is exactly
$\mathfrak c=\mathcal O_K$, this is the required
\texttt{AAUnit}; it is not inferred from self-duality alone.
VI2 supplies full period lattices, their self-duality, integral
quadratic multiplier, and finite generation. VI10 supplies the
actual fixed theta vector. Together these prove every
\texttt{LaHyp} clause, and $\mathfrak A\ne0$.

VI11, using only the independent finite presentation CA2, gives
\texttt{ESmooth} and \texttt{JetAlg}. CA9--CA10 give the actual
\texttt{VacEq}, \texttt{ThIndEq}, and \texttt{ChiLaw}.
The numbers $c_D(\tau)$ are algebraic: \texttt{AmbIn.Ok.dgen}
places each one in its specified number-field image.

Here is the finite-field argument, making the degree assertion
explicit. TL1 gives finitely many rational-unitary classes of
self-dual lattices in each rank. Choose representatives for
$L,L'$, finitely many positive divisors $c$ of $M$, and then
representatives for
\[
 M^{-1}\operatorname{Hom}_S(L,L')/
                    c\operatorname{Hom}_S(L,L').
 \tag{CA.37}
\]
The last quotient is finite by localization of a full lattice
and CRT; only the primes of $S$ dividing $Mc$ contribute.
Changing $x$ by its denominator module leaves the actual lift
predicate, hence the chosen indicator, unchanged.
For each of these finitely many functions retain its normalized
derivatives at zero of every total order at most $n$.
VI11 makes this a finite algebraic set. Adjoin also the finitely
many $c_D(\tau)$ and the images under every $\tau$ of a fixed
$\mathbb Q$-basis of $E$. The resulting field $K_0$ is generated
by finitely many algebraic numbers and contains every embedded
rational-unitary matrix entry and determinant.

For arbitrary $L,L',x$, choose unitary matrices carrying the
chosen lattice representatives to them. Equation (CA.35)
expresses their indicator by linear substitution in a selected
one. Repeated differentiation of a linear substitution is a
finite sum of derivatives of the same total order, multiplied
by products of the substitution's matrix entries. All these
entries lie in $K_0$. Equation (CA.30) shows that the ratio of
normalizing characters is in $K_0^\times$. Finally
\texttt{jetC} multiplies by the prescribed product of
$c_D(\tau)$ powers. Thus every normalized order-$n$ jet lies
in $K_0$, proving (CA.36). The polynomial here is the finite
Taylor component of a fixed degree, or its linear substitution;
no theta function is asserted to be a polynomial.

There is no hidden \texttt{LocStage}/\texttt{NStage} assumption
in this argument. The source proof of \texttt{finAlg\_\allowbreak step}
introduces arbitrary complete stages to access notation indexed
by $D7$. Its host-level \texttt{VacIn} already contains exactly
the period pair, half-point, rational isometry, complex reading,
and correction roots used above. It defines the same
\texttt{thMM}, \texttt{thZero}, \texttt{thMoved}, and
\texttt{eNat}; the source \texttt{toRec} and \texttt{*\_\allowbreak rec}
identities simply copy these fields. Every CA1--CA10 and VI
argument is expressed in those fields, the host, and its
embeddings. Applying those arguments directly to $I.J.V$ proves
(CA.36) for the literal $I.A=I.J.\texttt{amb}(B,I.\chi)$.
No unused complete-stage witness is required or claimed.

\paragraph{CA12: one finite generator set gives all canonical jets.}
Fix a preliminary datum $P$, character $\chi$, and finite twist
family $\Xi$. The actual finite bad-prime set is
$\mathcal B=\texttt{badP}(P,\Xi)$. For each word $u$, step
$i<m$, and subsets $S,R\subset\mathcal B$, put
\[
 q_R=\prod_{Q\in R}\operatorname NQ,\qquad
 M_{u,R}=q_R^{m+1}\bigl(T_uq_R^{m+1}\bigr).
 \tag{CA.38}
\]
The source \texttt{stepSet} uses this exact modulus, ranks
$N_i,N_{i+1}$, degree $\operatorname{totDeg}(P,u,i)$, and the
specific ambient datum $\texttt{amb0}(P.\mathrm{host})$.
Its \texttt{allSet} is the union only over those indices for
which \texttt{FinAlg(stepSet)} holds. There are finitely many
indices: the word and step sets are finite, and $S,R$ range over
two finite power sets. Taking the union of the finite algebraic
generator sets proves \texttt{FinAlg(allSet)}. This argument
does not silently assert finiteness for every arbitrarily chosen
prime set; the restriction in the actual definition is retained.

Let $Z=\texttt{gens}(P,\chi,\Xi)$ be the actual chosen finite
generator set. Every element is algebraic, and if a subfield
$K_0\subset\mathbb C$ contains $Z$, then it contains
\texttt{allSet}, by the defining universal property of field
generation. The fallback value $Z=\varnothing$ is never used
when $\Xi$ is finite.

Now let $I$ satisfy \texttt{Ok}, \texttt{IsTw}, and
$I.J=\texttt{amb0}(B.\mathrm{host})$, and assume
\[
 Z(B.\mathrm{toPre},I.C.\chi_\Psi,\operatorname{range}I.\xi)
                   \subset B.\mathrm{embLL}(B.LL).
 \tag{CA.39}
\]
Canonicality says that both exponents in $I.\mathrm{cst}$
are $m+1$, so its literal modulus is exactly (CA.38) for
$R=\texttt{rowSG}$; the type factor $T_u$ and all $q_Q$ are
positive. Its $S_b$ and $\texttt{rowSG}$ are subsets of
$\mathcal B$. Crucially, CA11 proves \texttt{FinAlg} for this
actual indexed step, so it satisfies the restriction in
\texttt{allSet}. The weight-to-degree identity of TL19 gives
$e.\mathrm{degree}=\operatorname{totDeg}(P,u,i)$ for every
multi-index tested by \texttt{JetsIn}. Therefore each such jet
lies in \texttt{allSet}, hence in the field image in (CA.39).
This proves the actual \texttt{jetsIn\_\allowbreak of\_\allowbreak isTw}.

\paragraph{CA13: the same-data canonical evaluation theorem.}
The preceding arguments close the analytic premises retained in
CI19/TL21. We spell out the quantifiers of the resulting
\texttt{EvLg} statement. Fix a source context $K$. Its
\texttt{large\_\allowbreak 2} witnesses may be chosen as follows.
The fixed coefficient-field requirement contains a generator
of TL14's common frame-value field. The full coefficient-field
requirement contains the PA13 period-field generators, TL15's
finite roots and square roots for the prescribed finite twist
family, the auxiliary-prime square roots, and CA12's $Z$.
Take the union of these finite sets; all its entries are
algebraic. The extra finite prime and avoidance requirements
can be empty. These sets are chosen before the later base and
field-containment antecedents, exactly as \texttt{large\_\allowbreak 2}
requires. The twist family is finite by TL21's support argument.

Suppose a base $B$, frame $f$, and twist datum $T$ lie over
$K$ and satisfy those field requirements. Retain the literal
frame-inclusion antecedent
\[
 i_X'(\operatorname{frame}_{X'}(a))
       =\operatorname{frame}_{X_m}(b)
       \quad\text{whenever the finite indices have equal values},
 \tag{CA.40}
\]
and suppose $\texttt{PreX}(B.\mathrm{toPre})$ is nonempty.
For every prescribed selection $P_v$ satisfying \texttt{IsSel},
the following construction uses this same $B,f,T,P_v$.

First orient the auxiliary factors of a supplied \texttt{PreX}
by TL16, choosing each factor or its conjugate according to
$f.P$. Replace only its type-prime elements by the context's
actual $\pi_T$. Since $c_u=\pi_T(u_0)/\pi_T(u)$, cancellation
of the nonzero denominator gives
$c_u\pi_T(u)=\pi_T(u_0)$; auxiliary orientation is unchanged.
Together with (CA.40), this is exactly both fields of
\texttt{Pins}. It asserts no new frame identity for arbitrary
raw frames.

Choose the actual $J=\texttt{amb0}(B.\mathrm{host})$;
PA9 proves its \texttt{AmbIn.Ok}. TL14--TL15 and the fixed/full
field containments supply the actual character values, tuples,
selected primes, and finite root data. CI15--CI17 construct
the local index records and all their nonvanishing and character
clauses, with initial $M=A=B=m$ and depth $2m$ off the auxiliary
primes. The fieldwise join of CI19 therefore gives a literal
$I_0:\texttt{ColIn}$ with
\[
 I_0.C=f,\quad I_0.\mathrm{sel}=P_v,\quad
 I_0.S=\{p\in\texttt{offSet}(B):p\ne p_\delta,p_\delta'\},
 \quad\operatorname{range}I_0.\xi=T.\mathrm{set}(f),\quad I_0.J=J.
 \tag{CA.41}
\]
The depth $2m$ satisfies the exact \texttt{Deep} bound needed
by CI18. Hence the specified replacement
$I=\texttt{canon}(I_0)$ satisfies \texttt{Ok} and
\texttt{IsTw}, preserving every field in (CA.41).
The same full-field requirement contains CA12's generators,
so CA12 supplies \texttt{JetsIn} for this $I$, with no change
of $J$, frame, or twist range. TL13 then also supplies
\texttt{Laws}, although the literal \texttt{evLg\_\allowbreak canonical}
conclusion does not require it as an extra conjunct.

TL14 supplies \texttt{ValLL0} for every later assembly
$C,W,V,S,N,G$ on this same base and frame; PA13 supplies the
same \texttt{FieldOK}. These are the final two conjuncts of
the source conclusion. A universal assertion over later
assemblies is not their existence assertion.

The construction works for every context $K$ already allowed
by the source. Thus the outer \texttt{Ev} wrapper uses
$\varepsilon_0=1$ and zero additional size thresholds, as in
\texttt{ev\_\allowbreak of\_\allowbreak forall}; all its original context and inner
field-containment quantifiers remain. This proves the actual
\texttt{evLg\_\allowbreak canonical} implication with its stated antecedents.
It neither creates an arbitrary preliminary datum, removes
(CA.40), nor constructs \texttt{Later}. The SE component
separately describes its stage suppliers and its still-retained
\texttt{LocStage}/\texttt{NStage} boundary. No inference to an
unconditional proof of RH or GRH is made here.


\section{Finite valuation extensions and the actual normalized arithmetic choice}
\label{sec:arithmetic-choice}

This section constructs the arithmetic comparison required by
\texttt{NChoice}, and then its exact character and tuple ties.
Throughout, \(D\) is a supplied \texttt{D7}. Write
\[
 L=D.\texttt{what}_2.\texttt{LL},\qquad L_0=D.\texttt{LL0},
 \qquad m=D.R.m,\qquad
 W_{\rm tw}=\{w:\mu_{\rm At}(w)\ne\mu_0\}.
\]
Thus \(L\) is a number field Galois over \(\mathbb Q\), equipped
with its specified complex embedding and conjugation. Its places
are the specified prime ideals of \(\mathcal O_L\).
The additive valuation has the convention \(v_w(0)=+\infty\);
\(v_w(x)\geq0\) means integrality, and \(v_w(x)>0\) means
membership in the maximal ideal.

The rings, spaces and lattices already carried by \(D\) are
retained as supplied objects. In particular, this section does
not construct \texttt{NTw}, \texttt{NExp}, or the row model.
It proves the finite DVR extension and integral comparison
needed to assemble their \texttt{NChoice}. The equality of the
two ranks, named \texttt{RankTie}, is a separate premise.
The resulting semilinear maps preserve the specified lattice
spans; their definition makes no assertion about Hecke operators
or Hermitian forms.

\paragraph{NC1: finite extensions of a complete DVR in characteristic zero.}
Let \(A\) be a complete DVR, with uniformizer \(\pi\) and fraction
field \(K\) of characteristic zero. Let \(K'/K\) be finite and put
\(B=\{x\in K':x\text{ is integral over }A\}\).
Then \(B\) is a complete DVR, finite and free over \(A\);
its fraction field is \(K'\), and \(A\longrightarrow B\) is an
injective local homomorphism.

Here is a direct proof of the finiteness that is needed in this
application. The finite extension is separable. Choose a primitive
element, multiply it by a nonzero element of \(A\) to make it
integral, and denote the result by \(\alpha\).
The order \(C=A[\alpha]\) is free of rank \([K':K]\), with
basis \(1,\alpha,\ldots,\alpha^{d-1}\).
For \(x\in B\), the numbers
\(\operatorname{Tr}_{K'/K}(x\alpha^i)\) belong to \(A\).
Indeed, they are integral over \(A\), as sums of conjugates of
integral elements, and they belong to \(K\); the DVR \(A\) is
integrally closed. The trace pairing is nondegenerate, so
\[
 C^\vee=\{x\in K':\operatorname{Tr}_{K'/K}(xC)\subset A\}
\]
is a free \(A\)-module of rank \(d\): its coordinate matrix is
the inverse of the nonsingular trace matrix of the displayed
basis. Thus \(C\subset B\subset C^\vee\).
A submodule of a finite free module over a DVR is finite free.
It follows that \(B\) is finite over \(A\), and that
\(B[1/\pi]=K'\).

Every maximal ideal of the integral finite \(A\)-algebra \(B\)
lies over \((\pi)\). The algebra \(B/\pi B\) is finite dimensional
over \(A/\pi A\), hence Artinian. If it had two local factors,
it would have a nontrivial idempotent \(\bar e\).
Such an idempotent lifts to \(B\). To see the lifting explicitly,
start with any lift \(e_0\); when \(e_r^2-e_r\in\pi^a B\),
\((2e_r-1)^2=1+4(e_r^2-e_r)\) is a unit and
\[
 e_{r+1}=e_r-(e_r^2-e_r)(2e_r-1)^{-1}
 \quad\text{satisfies}\quad
 e_{r+1}^2-e_{r+1}\in\pi^{2a}B.
 \tag{NC.1}
\]
The inverse of \(1+\pi b\) is its convergent geometric series.
Completeness here follows from the finite free \(A\)-module
structure of \(B\). The sequence in (NC.1) converges to an
idempotent with residue \(\bar e\), contradicting that \(B\)
is a domain. Thus \(B\) is local.

For completeness, the usual DVR criterion can also be made
explicit. A nonzero prime of \(B\) contracts to a nonzero prime
of \(A\), because \(B[1/\pi]\) is a field; therefore \(B\) has
dimension one. It is Noetherian and integrally closed.
If \(\mathfrak n\) is its maximal ideal and \(0\ne a\in\mathfrak n\),
the Artinian local ring \(B/aB\) has a nonzero socle. Choose a
representative \(u\notin aB\) with
\(\mathfrak n u\subset aB\), and put \(t=u/a\).
Then \(t\notin B\) but \(t\mathfrak n\subset B\).
The inclusion \(t\mathfrak n\subset\mathfrak n\) would make \(t\)
integral over \(B\), by applying the determinant trick to a finite
generating set of the nonzero ideal \(\mathfrak n\).
Consequently \(t\mathfrak n=B\), so \(\mathfrak n=t^{-1}B\)
is principal. Since \(\mathfrak n^c\subset\pi B\) for some \(c\),
the intersection of all its powers is zero. Repeated division
by its generator expresses every nonzero element as a unit times
a unique nonnegative power. This proves the DVR assertion.
The inclusions
\(\pi B\subset\mathfrak n\) and \(\mathfrak n^c\subset\pi B\)
also show that the two adic topologies agree, so \(B\) is complete
for its maximal ideal. Its maximal ideal contracts to \((\pi)\),
which proves locality of \(A\longrightarrow B\).

This proves the characteristic-zero instance of the integral
closure supplier used in the pinned source. That instance is
enough: every coefficient fraction field below contains \(L_0\).

\paragraph{NC2: an embedding reads a place exactly when it reads its maximal ideal.}
Let \(F\) be a number field, \(B\) a DVR with fraction field \(K'\),
and \(e:F\longrightarrow K'\) a field embedding.
Every algebraic integer of \(F\) maps into \(B\): it satisfies a
monic polynomial over \(\mathbb Z\subset B\), and \(B\) is
integrally closed. This gives a unique ring map
\(\phi:\mathcal O_F\longrightarrow B\).
Suppose a nonzero prime \(\mathfrak p\subset\mathcal O_F\)
satisfies
\[
 \phi^{-1}(\mathfrak m_B)=\mathfrak p.
 \tag{NC.2}
\]
Then
\[
 v_{\mathfrak p}(x)\geq0
 \quad\Longleftrightarrow\quad e(x)\in B
 \qquad(x\in F).
 \tag{NC.3}
\]
Indeed, \(\mathcal O_{F,\mathfrak p}\) is the valuation ring at
\(\mathfrak p\). If \(x=a/s\) with \(s\notin\mathfrak p\),
then \(\phi(s)\) is a unit, giving the forward implication.
Conversely, if \(v_{\mathfrak p}(x)<0\), the element \(x^{-1}\)
belongs to its maximal ideal. Write \(x^{-1}=a/s\) with
\(a\in\mathfrak p\), \(s\notin\mathfrak p\).
Its image lies in \(\mathfrak m_B\), so its inverse cannot
belong to \(B\). This proves (NC.3), including \(x=0\).
For nonzero \(x\), the same argument gives
\(v_{\mathfrak p}(x)>0\) precisely when \(e(x)\in\mathfrak m_B\).

We also need transitivity on primes, with the direction of
conjugation specified. If \(F/F_0\) is finite Galois and
\(\mathfrak p,\mathfrak q\) lie over the same prime of
\(\mathcal O_{F_0}\), there is
\(\sigma\in\operatorname{Gal}(F/F_0)\) with
\(\sigma^{-1}\mathfrak q=\mathfrak p\).
If not, the Chinese remainder theorem gives an integer \(a\)
which is zero modulo \(\mathfrak q\) and one modulo every
conjugate of \(\mathfrak p\). Its norm
\(\prod_\tau\tau(a)\) lies in
\(\mathfrak q\cap\mathcal O_{F_0}\), and hence in
\(\mathfrak p\), although every factor avoids \(\mathfrak p\).
This contradicts primality. This argument proves precisely
the prime alignment used next; it needs no class-number premise.

\paragraph{NC3: the finite coefficient extension at a prescribed prime.}
Let \(L/\mathbb Q\) be a finite Galois extension, \(L_0\subset L\)
an intermediate field, \(A\) a complete DVR and
\(i:L_0\longrightarrow K=\operatorname{Frac}(A)\) an embedding.
Let \(\mathfrak p\) be a nonzero prime of \(\mathcal O_L\).
Assume that the maximal ideal pulled back through \(i\) on
\(\mathcal O_{L_0}\) is exactly
\(\mathfrak p\cap\mathcal O_{L_0}\).
There exist a finite complete DVR extension \(A\subset B\)
and an embedding \(e:L\longrightarrow\operatorname{Frac}(B)\)
such that
\[
 e|_{L_0}=i,\qquad
 v_{\mathfrak p}(x)\geq0\ \Longleftrightarrow\ e(x)\in B.
 \tag{NC.4}
\]
Here \(i\) in the first equality is followed by the fraction-field
embedding of \(A\subset B\). In particular, for \(x\in L_0\)
and \(y,z\in A\),
\[
 i(x)z=y\quad\Longrightarrow\quad e(x)z=y
 \quad\text{in }\operatorname{Frac}(B).
 \tag{NC.5}
\]

Choose a primitive element for \(L/L_0\). A root of its minimal
polynomial, with coefficients mapped by \(i\), in an algebraic
closure of \(K\) gives an \(L_0\)-embedding of \(L\) into a
finite extension \(K'/K\). Equivalently, adjoin the images of
a finite \(L_0\)-basis; every such image is algebraic over \(K\).
Apply NC1 to \(K'\), and call the resulting embedding \(e_1\).
The inverse image of \(\mathfrak m_B\) on \(\mathcal O_L\)
is a prime \(\mathfrak q\) above
\(\mathfrak p\cap\mathcal O_{L_0}\). It is nonzero because this
contracted prime contains a rational prime. Since \(L/L_0\)
is Galois, NC2 supplies \(\sigma\) with
\(\sigma^{-1}\mathfrak q=\mathfrak p\).
Set \(e=e_1\circ\sigma\). This changes the prime read by the
embedding and fixes its restriction to \(L_0\).
NC2 now proves (NC.4), and restriction to \(L_0\) proves (NC.5).

The literal source first adjoins a square root as well. It calls
that construction with \(a=1\) here, so this extra generator is
unnecessary; adjoining it would still give the same finite
extension argument. The assertion concerns valuation rings,
not an equality of numerical normalizations of two valuations.
A further ramification index may occur in \(B\).

\paragraph{NC4: actual readings, the residue characteristic, and the cover implication.}
For a twist index \(i\) and side \(b\in\{\mathrm{false},\mathrm{true}\}\),
put
\[
 A_0=O_0(i),\qquad A_k=O_k(i),\qquad
 s_{i,b}=\begin{cases}\mathrm{emb}_0(i)&b=\mathrm{true},\\
                       \mathrm{emb}_{0b}(i)&b=\mathrm{false}.
          \end{cases}
\]
The condition \(\texttt{ReadsAt}(i,b,w)\) says that for every
\(x\in\mathcal O_{L_0}\),
\[
 s_{i,b}(x)\in\mathfrak m_{A_0}
       \quad\Longleftrightarrow\quad v_w(x)>0.
 \tag{NC.6}
\]
The image on the left is interpreted through the unique integral
map of NC2. This is an equality of the prime being read.

The map \(A_0\to O(i)\) is local. If a nonunit had unit image,
write it as a unit times a positive power of a uniformizer;
the uniformizer would have unit image. The nonzero integer
\(\ell(i)\) is also a unit times a power of that uniformizer,
so its image would be a unit. This contradicts the supplied
condition \(\ell(i)\in\mathfrak m_{O(i)}\).
The map \(O(i)\to A_k\) is injective and finite, hence integral
and local. Moreover \(A_k\) is complete: it is finite torsion
free over the complete DVR \(O(i)\), hence a finite free
complete module; its maximal-ideal and base-adic topologies
agree because the corresponding quotient is Artinian.
Consequently
\(\phi_k:A_0\to A_k\) is injective and local.
It therefore extends to an embedding of fraction fields.

Apply NC3 to the composite
\[
 L_0\xrightarrow{s_{i,b}}\operatorname{Frac}(A_0)
       \longrightarrow\operatorname{Frac}(A_k)
\]
and the specified prime of \(L\) at \(w\).
Locality and (NC.6) verify its prime hypothesis.
This constructs, for every \(k\geq0\), a
\(\rho:\texttt{Reading}(i,k,b,w)\) consisting of a DVR \(B\),
an injective finite local \(A_k\)-algebra structure, and
\(e_\rho:L\to\operatorname{Frac}(B)\), such that
\[
 v_w(x)\geq0\Longleftrightarrow e_\rho(x)\in B,\qquad
 s_{i,b}(x)z=y\Longrightarrow
 e_\rho(x)\phi_k(z)=\phi_k(y).
 \tag{NC.7}
\]
The second equality is followed into \(B\) and its fraction field.
It holds even for \(z=0\); the later \texttt{NChTie} field
requires it only for \(z\ne0\).

Equation (NC.6) applied to \(\ell(i)\) also proves
\(\ell(i)=\operatorname{resChar}(w)\).
If two distinct rational primes belonged to the same prime
ideal, their integer Bezout combination would place \(1\)
there, which is impossible.
Finally, the supplied \texttt{TwCover} reads every prime of
\(L_0\) with side \(\mathrm{true}\). Restricting \(w\) to \(L_0\)
therefore supplies \texttt{ReadsAt} at every \(w\), and in
particular the exact off-twist condition
\(\texttt{ReadElsewhere}(D)\).
This proves the implication from the cover; it does not
construct \texttt{TwCover}.

\paragraph{NC5: the finite rational generator set of a complex form.}
Let \(V\) be an \(L\)-vector space, \(M\) a finite-dimensional
complex vector space, and \(j:V\to M\) semilinear for an
embedding \(\iota:L\to\mathbb C\).
Suppose an \(L\)-basis \((b_a)\) maps to a complex basis.
Let \(T\subset M\) be finite, span \(M\), and lie in \(j(V)\).
Then \(V\) is finite dimensional, \(j\) is injective, and
\[
 S=\{v\in V:j(v)\in T\}
 \quad\text{is finite with}\quad \operatorname{span}_L S=V.
 \tag{NC.8}
\]
The basis index is finite because its image is independent in
the finite-dimensional space \(M\).
In basis coordinates \(j\) applies \(\iota\) to every coefficient;
thus it is injective, and \(S\) has at most \(|T|\) elements.
If \(S\) failed to span, a nonzero \(L\)-linear functional
\(\lambda\) would vanish on it. Define a complex linear
functional by
\(\lambda_{\mathbb C}(j(b_a))=\iota(\lambda(b_a))\).
Semilinearity gives
\(\lambda_{\mathbb C}(j(v))=\iota(\lambda(v))\).
It vanishes on \(T\), hence on \(M\), and then on every
\(j(b_a)\), contradicting \(\lambda\ne0\).

For the literal row objects take
\[
 V=M_{t,L}(\mu,j),\quad M=M_t(\mu,j),\quad
 T=\texttt{MtLat}(\mu,j).\texttt{gen}.
 \tag{NC.9}
\]
The premises are exactly \texttt{jMt\_\allowbreak basis}, the finite
generator and spanning fields of the supplied lattice, and
\texttt{MtLat\_\allowbreak sub}. Thus the preimage of the actual displayed
generator set is finite and spans over \(L\).
This proves rational descent of those generators; it makes
no new claim that the row model itself exists.

\paragraph{NC6: scalar extension preserves the rank and the full lattice.}
Let \(A\hookrightarrow B\) be an injective map of domains,
write \(K=\operatorname{Frac}(A)\), \(K'=\operatorname{Frac}(B)\),
and let \(V\) be a \(K\)-vector space.
If \(\Lambda\subset V\) is a finitely generated \(A\)-submodule
spanning \(V\) over \(K\), put
\[
 V'=K'\otimes_A V,\qquad
 \Lambda'=\operatorname{span}_B\{1\otimes v:v\in\Lambda\}.
 \tag{NC.10}
\]
Then
\[
 \dim_{K'}V'=\dim_K V,\qquad
 \Lambda'\text{ is finite and spans }V'\text{ over }K'.
 \tag{NC.11}
\]
The injection extends to \(K\hookrightarrow K'\), and
\(K'\otimes_A V\cong K'\otimes_K V\): every nonzero element
of \(A\) already acts invertibly on both modules.
A \(K\)-basis therefore becomes a \(K'\)-basis, proving the
dimension statement. If \(\lambda_1,\ldots,\lambda_s\)
generate \(\Lambda\), then \(1\otimes\lambda_i\) generate
\(\Lambda'\) over \(B\).
For any \(v\in V\), clear denominators in a finite expression
of \(v\) in the \(\lambda_i\). There is \(0\ne d\in A\)
with \(dv\in\Lambda\), and
\[
 a\otimes v=(a/d)(1\otimes dv).
 \tag{NC.12}
\]
This proves the spanning assertion for every pure tensor
and hence for all of \(V'\).
In particular, if \(B\) is a DVR, \(\Lambda'\) is a full
finite free \(B\)-lattice.

\paragraph{NC7: matching full DVR lattices gives the required semilinear map.}
Let \(B\) be a DVR with fraction field \(K'\),
\(e:L\hookrightarrow K'\), and let \(V\) be a finite-dimensional
\(L\)-space. Suppose \(S\subset V\) is finite and spans \(V\).
Let \(U\) be a \(K'\)-space with full finite \(B\)-lattice
\(\Lambda_U\), and suppose
\(\dim_LV=\dim_{K'}U=r\).
There is an \(e\)-semilinear map \(f:V\to U\) such that
\[
 \operatorname{span}_B f(S)=\Lambda_U,
 \qquad
 \{f(b_a)\}_a\text{ is a }K'\text{-basis of }U
 \text{ for every }L\text{-basis }(b_a)\text{ of }V.
 \tag{NC.13}
\]

Choose one \(L\)-basis and let \(f_0:V\to(K')^r\) apply \(e\)
to its coordinates. The module
\(A_S=\operatorname{span}_B f_0(S)\) is finitely generated,
torsion free, and spans \((K')^r\). A finite torsion-free
module over a DVR is free; clearing denominators in a relation
shows that a \(B\)-basis is also a \(K'\)-basis. Thus both
\(A_S\) and \(\Lambda_U\) have \(B\)-bases of cardinality \(r\).
Map these bases to one another and extend \(K'\)-linearly,
obtaining \(g:(K')^r\overset{\sim}{\longrightarrow}U\).
Then \(f=g\circ f_0\) has the asserted integral span.
For any other \(L\)-basis, its change-of-basis matrix has
nonzero determinant. Applying the field embedding to that
determinant keeps it nonzero, proving the assertion for
every basis.

One elementary proof of the freeness used here is to take
a generating family of smallest cardinality. In a nontrivial
linear relation, choose a nonzero coefficient of smallest
valuation; divide all coefficients by it. Torsion freeness
allows cancellation, and the chosen coefficient becomes a
unit. One generator is then a combination of the others,
contradicting minimality. This proves freeness without an
assumed integral comparison map.
The induced map \(K'\otimes_{L,e}V\to U\) is an isomorphism;
\(f\) itself need not be surjective as a map of sets.

\paragraph{NC8: the literal \texttt{exists\_\allowbreak formsIso} conclusion.}
Let \(A\hookrightarrow B\), \(V\), and \(\Lambda\) be as in
NC6, with \(B\) a DVR. Let \(e:L\to\operatorname{Frac}(B)\).
For a tuple \(\mu\), index \(j\), and any proposition \(c\),
assume
\[
 c\Longrightarrow
 \dim_L M_{t,L}(\mu,j)=\dim_{\operatorname{Frac}(A)}V.
 \tag{NC.14}
\]
Then there is an \(e\)-semilinear map
\[
 f:M_{t,L}(\mu,j)\longrightarrow
       \operatorname{Frac}(B)\otimes_A V
 \tag{NC.15}
\]
which, whenever \(c\) holds, carries every \(L\)-basis to a
basis after extension of scalars and satisfies
\[
 \operatorname{span}_B
 f\bigl(\{x:j_{Mt}(x)\in\texttt{MtLat}(\mu,j).\texttt{gen}\}\bigr)
 =
 \operatorname{span}_B\{1\otimes\lambda:\lambda\in\Lambda\}.
 \tag{NC.16}
\]
If \(c\) is false, take \(f=0\). Otherwise NC5 gives the
finite spanning generator set on the left, NC6 gives the
full target lattice and its rank, and NC7 constructs \(f\).
This proves both conclusions simultaneously on the actual
generator sets. The only equality of dimensions used is
(NC.14); it is not produced by the construction.

\paragraph{NC9: all fields of the normalized \texttt{NChoice}.}
A twist packet is
\[
 p_w=(i_w,Q_w,a_w,\chi_w),\qquad
 \chi_w:\mathrm{RCG}(i_w,Q_w)\longrightarrow O_{a_w}(i_w)^\times.
\]
The source convention is
\[
 \chi^{[b]}=
 \begin{cases}\chi&b=\mathrm{true},\\
                \chi^{-1}&b=\mathrm{false},
 \end{cases}
 \qquad(\chi^{[b]})^{[b]}=\chi.
 \tag{NC.17}
\]
Suppose, for every \(w\), we have a side \(b_w\), a reading
\(\rho_w:\texttt{Reading}(i_w,a_w,b_w,w)\), and
\(\texttt{ReadsAt}(i_w,b_w,w)\).
For \(w\in W_{\rm tw}\), assume the exact rank equality
\[
 \begin{split}
 \dim_L M_{t,L}(\mu_{\rm Tw}(w),j)
 &=
 \dim_{\operatorname{Frac}(O_{a_w}(i_w))}
 M_t^{\rm loc}\bigl(i_w,j,b_w,a_w,Q_w,
                 \chi_w^{[b_w]}\circ\operatorname{art}_{i_w,Q_w}\bigr),
 \\
 &\hspace{100pt}1\leq j\leq m.
 \end{split}
 \tag{NC.18}
\]
This is exactly \(\texttt{RankTie}(p_w,b_w,w)\).

Set \(B_w=\rho_w.O_e\), \(K_w=\operatorname{Frac}(B_w)\)
and \(e_w=\rho_w.\mathrm{emb}\).
For each \(w,j\), apply NC8 to the actual local
\(\texttt{MtLat}\) and to the condition
\[
 c_{w,j}=(w\in W_{\rm tw})\ \wedge\ (1\leq j\leq m).
 \tag{NC.19}
\]
Its finite generation and spanning hypotheses are the supplied
\texttt{NExp.MtLat\_\allowbreak fg} and \texttt{MtLat\_\allowbreak span}; its rank
hypothesis is (NC.18). Choose the resulting map as
\(\texttt{eMt}(w,j)\).
Outside (NC.19) it may be the zero map, exactly as in NC8.

The full record has the following 21 fields. This list includes
the typeclass fields, whose content matters for the construction.
\[
\begin{array}{ll}
 \texttt{twI},\texttt{twI\_\allowbreak l},\texttt{sd}
     & i_w,\ \ell(i_w)=\operatorname{resChar}(w),\ b_w,\\
 \texttt{Q},\texttt{k'},\texttt{psi'}
     & Q_w,\ a_w,\ \chi_w^{[b_w]},\\
 \texttt{Oe},\texttt{crOe},\texttt{domOe},\texttt{dvrOe}
     & B_w\text{ and its ring, domain and DVR structures},\\
 \texttt{algOe},\texttt{algOOe},\texttt{towOe}
     & O_{a_w}(i_w)\to B_w,\ O(i_w)\to B_w,\text{ their tower},\\
 \texttt{Oe\_\allowbreak inj},\texttt{Oe\_\allowbreak fin},\texttt{Oe\_\allowbreak loc}
     & \text{injectivity, finiteness and locality from }\rho_w,\\
 \texttt{embW},\texttt{embW\_\allowbreak w}
     & e_w\text{ and the valuation-ring equivalence (NC.7)},\\
 \texttt{eMt},\texttt{eMt\_\allowbreak basis},\texttt{eMt\_\allowbreak lat}
     & \text{the maps and the two conclusions of NC8}.
\end{array}
\tag{NC.20}
\]
The algebra \(O(i_w)\to B_w\) is the actual composite through
\(O_{a_w}(i_w)\), so the scalar-tower field holds as an equality
of those maps. NC4 supplies \texttt{twI\_\allowbreak l}.
For every \(w\in W_{\rm tw}\) and \(1\leq j\leq m\),
\texttt{eMt\_\allowbreak basis} is quantified over every chosen-index basis,
and \texttt{eMt\_\allowbreak lat} is precisely (NC.16) for
\(\mu=\mu_{\rm Tw}(w)\). All 21 fields are therefore proved.

Call the result \(C=\texttt{choiceOfPack}(p,b,\rho)\).
The packet extracted from a choice applies the side convention
once more to \(\texttt{psi'}\). Equation (NC.17) gives the
literal equality
\[
 \texttt{NChoicePack}(C,w)=p_w.
 \tag{NC.21}
\]
Thus the packet, including its selected primes and character,
is recovered exactly; it is not recovered only up to inversion
or isomorphism.

\paragraph{NC10: descending finite-order characters and conjugating a reading.}
Let \(e:L\hookrightarrow K'\) be a field embedding.
If \(L\) contains a primitive \(n\)-th root \(\zeta\), \(n>0\),
and \(\chi:\Gamma\to(K')^\times\) satisfies \(\chi(g)^n=1\),
then there is a unique character \(\psi:\Gamma\to L^\times\)
with \(e(\psi(g))=\chi(g)\).
The \(n\) distinct powers of \(e(\zeta)\) are all the roots of
\(X^n-1\) in \(K'\): a polynomial of degree \(n\) has at most
\(n\) roots. Thus each \(\chi(g)\) has a unique inverse image
among the powers of \(\zeta\).
Injectivity of \(e\) proves that these inverse images preserve
products and the identity.

For a packet \(p=(i,Q,a,\chi)\), the condition
\(\texttt{ValuesIn}(p)\) is exactly the existence of such an
\(n\) and primitive root in \(L\), with \(\chi(\delta)^n=1\)
for every \(\delta\in\mathrm{RCG}(i,Q)\).
Given a reading \(\rho\), apply the preceding construction
to
\(\chi^{[b]}\circ\operatorname{art}_{i,Q}\), followed into
\(\operatorname{Frac}(\rho.O_e)^\times\).
This gives \(\psi_0:G_E\to L^\times\) satisfying
\[
 e_\rho(\psi_0(g))
 =\chi^{[b]}(\operatorname{art}_{i,Q}(g))
 \quad\text{in }\operatorname{Frac}(\rho.O_e).
 \tag{NC.22}
\]
Here and below \(G_E\) is the exact Galois group \(G(E,D.nd.Om)\).
Equation (NC.22) is the literal condition
\(\rho.\texttt{Pulls}(\chi^{[b]},\psi_0)\).

To transport a reading across conjugate places, additionally
assume \(L_0/\mathbb Q\) is Galois. Complex conjugation then
preserves \(L_0\) and \(\mathcal O_{L_0}\).
The supplied two coefficient embeddings satisfy
\[
 s_{i,!b}(x)=s_{i,b}(\bar x),\qquad
 v_{\bar w}(x)=v_w(\bar x).
 \tag{NC.23}
\]
Keep the same DVR and put \(e_{\rho^c}=e_\rho\circ\bar{\phantom{x}}\).
Equations (NC.7) and (NC.23) prove every field of
\(\rho^c:\texttt{Reading}(i,a,!b,\bar w)\).
For a root of unity \(u\in L\), its specified complex image
has absolute value one, so \(\bar u=u^{-1}\).
If (NC.22) holds, the character \(\psi_0\) has finite order
pointwise: the target comes from the finite group
\(\mathrm{RCG}(i,Q)\), and \(e_\rho\) is injective.
Therefore
\[
 e_{\rho^c}(\psi_0(g))
 =e_\rho(\psi_0(g)^{-1})
 =\bigl(\chi^{[b]}(\operatorname{art}(g))\bigr)^{-1}
 =\chi^{[!b]}(\operatorname{art}(g)).
 \tag{NC.24}
\]
The conjugate reading pulls the \emph{same} \(L\)-valued
character, with the side reversed.

\paragraph{NC11: a coherent common character exists under the exact value hypotheses.}
Suppose packets \(t_w\) and sides \(b_w\) are supplied for
\(w\in W_{\rm tw}\), and satisfy the reading and conjugation
parts of \texttt{GoodData}:
\[
 \texttt{ReadsAt}(t_w.i,b_w,w),\quad
 t_{\bar w}=t_w,\quad
 \bar w\ne w\Longrightarrow b_{\bar w}=!b_w.
 \tag{NC.25}
\]
Assume \(L_0/\mathbb Q\) is Galois and
\(\texttt{ValuesIn}(t_w)\) for every such \(w\).
Then there are characters \(\psi_{0,w}:G_E\to L^\times\)
such that
\[
 \psi_{0,\bar w}=\psi_{0,w},\qquad
 w\notin W_{\rm tw}\Longrightarrow\psi_{0,w}=1,
 \tag{NC.26}
\]
and each \(w\in W_{\rm tw}\) has a reading pulling
\(\psi_{0,w}\) from the prescribed packet and side.
This is the exact \texttt{IsCharOf}(t,b,\(\psi_0\)) assertion.

Choose one representative of each orbit of the involution
\(w\mapsto\bar w\). On a twisted representative use NC4
and NC10 to choose a reading and descended character.
On the other member of a two-element orbit use (NC.24);
(NC.25) gives exactly its packet and side.
On a one-element orbit the original choice already suffices:
there is no side-reversal requirement there.
Put the character equal to one on the remaining orbits.
This proves all of (NC.26) and all the required readings.
The additional Hecke-goodness clauses of \texttt{GoodData}
are not used in this character construction.

The result chooses a common character; it does not identify it
with an independently prescribed local Galois character.
That exact compatibility is the next hypothesis,
\texttt{LocsCarry}, and is kept on these same objects.

\paragraph{NC12: the actual tied choice, with all five \texttt{NChTie} fields.}
Retain \(\texttt{ReadElsewhere}(D)\), meaning that every
\(w\notin W_{\rm tw}\) has some \(\texttt{ReadsAt}(i,b,w)\).
On \(W_{\rm tw}\) suppose packets \(t_w\), sides \(b_w\),
the readings in (NC.25), and \(\texttt{RankTie}(t_w,b_w,w)\)
are supplied. Suppose also that a \emph{fixed} family
\(\psi_{0,w}\) satisfies \texttt{IsCharOf}. For each existing
local object \(a:D.nd.Loc\), put \(w_a=D.nd.w(a)\);
the record supplies \(w_a\in W_{\rm tw}\).
The condition \texttt{LocsCarry} requires
\[
 \begin{gathered}
 D.ntw.twOf(a)=t_{w_a}.i,\qquad
 D.ntw.twSide(a)=b_{w_a},\qquad
 D.nd.Q(a)=t_{w_a}.Q,\\
 D.nd.embK(a)(\psi_{0,w_a}(g))
   =D.ntw.psiL(a)(g)
       \quad\text{in }\operatorname{Frac}(D.nd.O(a))
       \quad(g\in G_E).
 \end{gathered}
 \tag{NC.27}
\]
The last line includes the specified injection of the unit
on the right into that fraction field.

For \(w\in W_{\rm tw}\), use the reading supplied by
\texttt{IsCharOf}, so that it pulls this very \(\psi_{0,w}\).
For \(w\notin W_{\rm tw}\), choose a reading index and side
from \texttt{ReadElsewhere}, set
\[
 p_w=(i_w,\varnothing,0,1),
 \tag{NC.28}
\]
and use NC4 with \(k=0\). It pulls \(\psi_{0,w}=1\).
Apply NC9 to this family. Rank equality is needed only
on \(W_{\rm tw}\), where it is the prescribed \texttt{RankTie}.
The resulting \(C:\texttt{NChoice}(D)\) has
\[
 \texttt{NChoicePack}(C,w)=t_w,\quad C.sd(w)=b_w
 \qquad(w\in W_{\rm tw}).
 \tag{NC.29}
\]
This is exactly \texttt{Records}(C,t,b).

For \(\texttt{NChTie}(D,C)\), its first three fields
\texttt{twOf\_\allowbreak eq}, \texttt{twSide\_\allowbreak eq}, and \texttt{Q\_\allowbreak eq}
follow from (NC.27) and (NC.29).
Its fourth field \texttt{psi\_\allowbreak eq} uses the witness
\(\psi_0\) already fixed above: conjugation invariance is
(NC.26), the chosen readings prove its pullback equality
at every \(w\), and the last line of (NC.27) gives its equality
at every existing local object.
Its fifth field \texttt{embedding\_\allowbreak baseField\_\allowbreak condition}
is the denominator equality in (NC.7), restricted to \(z\ne0\).
The composite \(A_0\to O(i_w)\to O_{a_w}(i_w)\to B_w\)
is the same one used in (NC.20). This proves
\[
 \exists\, C:\texttt{NChoice}(D),\qquad
 \texttt{NChTie}(D,C)\ \wedge\ \texttt{Records}(C,t,b).
 \tag{NC.30}
\]
No \texttt{formsIso} or reading-extension existence is
left as an extra premise in (NC.30).

If the full \texttt{GoodData}(I,t,b) is supplied, (NC.29)
also gives \texttt{TwConjAt}(D,C) and
\texttt{TwGoodAt}(D,C,I). For the former, substitute the
identical packets at \(w,\bar w\) and their opposite sides
when distinct. For the latter, substitute the identical
packet in the \texttt{GoodPack} assertion; (NC.17) restores
its normalized character. This transports supplied
Hecke-goodness and does not establish the analytic
\texttt{GoodTw} supplier.

\paragraph{NC13: the exact rank and reading requirements are necessary as well.}
For every supplied \(C:\texttt{NChoice}(D)\) and
\(w\in W_{\rm tw}\),
\[
 \texttt{RankTie}(\texttt{NChoicePack}(C,w),C.sd(w),w)
 \tag{NC.31}
\]
holds. Choose an \(L\)-basis of the relevant \(M_{t,L}\).
The field \texttt{eMt\_\allowbreak basis} makes its images a basis after
extension to \(\operatorname{Frac}(C.Oe(w))\), so the two
dimensions there equal the cardinality of that basis.
NC6 identifies the target dimension with the original
local fraction-field dimension; (NC.17) removes the two
side inversions. This proves (NC.31).

If \(C\) is tied, it also gives
\(\texttt{ReadsAt}(C.twI(w),C.sd(w),w)\) for every \(w\).
For \(x\in\mathcal O_{L_0}\), let \(y\in A_0\) be its
integral image through the side embedding, as in NC2.
The tied denominator identity with denominator \(1\)
identifies \(C.embW(w)(x)\) with the image of \(y\) in \(B_w\).
The two coefficient maps are local, so this image is a
nonunit exactly when \(y\) is a nonunit.
For \(x\ne0\), \texttt{embW\_\allowbreak w} applied to \(x^{-1}\)
identifies nonunit status with \(v_w(x)>0\);
for \(x=0\) both statements hold directly.
Thus the precise maximal-ideal reading (NC.6) follows.
In particular \texttt{ReadElsewhere} follows from a tied
choice. Two tied choices have identical \(i,b,Q\) at each
existing local object, since each equals the corresponding
entry in (NC.27).

These converse facts explain the retained hypotheses.
They cannot be used to prove \texttt{RankTie} before the
choice exists, nor do they assert uniqueness of its DVR,
embedding, character away from the constrained objects,
or chosen lattice isomorphisms.

\paragraph{NC14: finite-order tame characters and nontrivial distinguished twists.}
Now fix the exact \texttt{Stages} \(X\), character frame \(f\),
and \(\texttt{TwDat}\ T\) used in the tuple endpoint.
Write \(E=X.B.host.E\), \(F=X.B.host.F\), and
\[
 \xi_{\rm Tw}(w)
  =\delta_{\operatorname{resChar}(w)}\,T.\psi(w),\qquad
 \delta_\ell=
 \begin{cases}f.\delta'&\ell=X.B.host.pde,\\
               f.\delta&\text{otherwise}.
 \end{cases}
 \tag{NC.32}
\]
Let
\(\Xi_T=\{1,f.\delta,f.\delta'\}\cup
              \{\xi_{\rm Tw}(w):w\}\).
Every member has finite order, is trivial on \(F\) in the
specified finite and infinite restriction sense, and is
tame at every finite place.
Indeed these assertions hold for both distinguished
characters by their \texttt{IsDelta} fields, and for
\(T.\psi(w)\) by its \texttt{IsTwist} field.
For a product, exponents \(n,n'>0\) can be replaced by
\(nn'\); restrictions and local values multiply.
This proves the assertions for (NC.32), as well as for
the identity character.

Every \(\xi_{\rm Tw}(w)\) is nontrivial.
For either distinguished character supported above its
prime \(p\), choose a prime \(P\) of \(E\) above \(p\).
Such a prime exists: \(\mathcal O_E/p\mathcal O_E\) is a
nonzero finite ring, so it has a maximal ideal.
It is nonzero because \(1/p\) is not an algebraic integer.
The twist character \(T.\psi(w)\) is unramified at \(P\).
Otherwise its \texttt{IsTwist} ramification clause says
that \(N(P)\) is a rational prime outside \texttt{offSet}.
As \(P\mid p\), that prime is \(p\), but both distinguished
primes belong to \texttt{offSet}, a contradiction.
The \texttt{IsDelta.ine} clause says that the restriction
of the distinguished character to local units has order
exactly \(15\): exponent \(n\) kills it if and only if
\(15\mid n\). In particular it is nontrivial.
An unramified twist is trivial on those units, so their
product cannot be the identity.

There are only finitely many places of \(L\) above a finite
set of rational primes: each is an ideal of the finite ring
\(\mathcal O_L/p\mathcal O_L\).
Consequently
\[
 \{w:\operatorname{resChar}(w)\in T.pl\}
 \quad\text{and}\quad \Xi_T
 \quad\text{are finite}.
 \tag{NC.33}
\]
For the second assertion, outside the first set
\(T.\psi(w)=1\), so (NC.32) equals one of the two
distinguished characters. Conjugation preserves residue
characteristic, and \(T.\psi(\bar w)=T.\psi(w)\), so
\(\xi_{\rm Tw}(\bar w)=\xi_{\rm Tw}(w)\) as well.

\paragraph{NC15: \texttt{MuIsW} identifies the exact finite support.}
Retain the supplied \(\texttt{MuIsW}(X,f,T,W,\xi)\).
Its ten clauses are: injectivity and range \(\Xi_T\) of
\(\xi\); its specified values \(1,f.\delta,f.\delta'\)
at \(\mu_0,\mu_\star,\mu_{\star'}\); the identity
\(\xi(\mu_{\rm Tw}(w))=\xi_{\rm Tw}(w)\); the assignments
\(\mu_{\rm At}(w)=\mu_{\rm Tw}(w)\) on \(W\) and
\(\mu_{\rm At}(w)=\mu_0\) off \(W\); the inclusion
of \(W\) into the first set of (NC.33); and
\(T.\psi(w)=1\) off \(W\).
Then
\[
 \mu_{\rm At}(w)\ne\mu_0\quad\Longleftrightarrow\quad w\in W,
 \qquad W\text{ finite},\qquad
 w\in W\Longrightarrow\bar w\in W.
 \tag{NC.34}
\]
If \(w\notin W\), its off-support assignment gives
\(\mu_{\rm At}(w)=\mu_0\).
If \(w\in W\) and this equality held, applying \(\xi\)
would give \(\xi_{\rm Tw}(w)=1\), contrary to NC14.
This proves the equivalence. Finiteness follows from
the supplied inclusion and (NC.33).
The row field \(\mu_{\rm At}(\bar w)=\mu_{\rm At}(w)\)
then proves conjugation stability.
In particular \(W\) is the exact support, although it
need not be the set of all places above \(T.pl\).

\paragraph{NC16: all twelve fields of \texttt{TupTied} on the same characters.}
Under the hypotheses of NC15, let
\(C:\texttt{NChoice}(X.toD7)\) satisfy
\[
 C.Q(w)=T.Q(w)\qquad(w\in W).
 \tag{NC.35}
\]
There is a
\(\mathrm{tp}:\texttt{TupTied}(X.toD7,\texttt{Fr.at}(X,f),C)\)
with the literal equalities
\[
 \mathrm{tp}.\xi=\xi,\qquad
 \mathrm{tp}.\psi=T.\psi.
 \tag{NC.36}
\]
Its twelve fields are proved as follows.
The fields \texttt{xi} and \texttt{psi} are those exact
functions. \texttt{xi\_\allowbreak inj} is the \texttt{MuIsW} injection.
For \texttt{xi\_\allowbreak fin}, \texttt{xi\_\allowbreak triv}, and
\texttt{xi\_\allowbreak tame}, the range identity puts every
\(\xi(\mu)\) in \(\Xi_T\), and NC14 gives the three assertions.
The field
\texttt{trueColumnChoices\_\allowbreak heckeChar\_\allowbreak reference\_\allowbreak eq\_\allowbreak one}
and the two fields \texttt{xi\_\allowbreak star}, \texttt{xi\_\allowbreak star'}
are exactly the three specified values in NC15.
The field \texttt{xi\_\allowbreak tw} is (NC.32) and the supplied
twisted-tuple identity, since \texttt{Fr.at} uses the same
distinguished characters.
For \texttt{psi\_\allowbreak is}, a nonreference place belongs to \(W\)
by (NC.34); substitute (NC.35) in the actual
\(T.\texttt{psi\_\allowbreak is}(w)\) assertion.
For \texttt{psi\_\allowbreak un}, a reference place is outside \(W\),
and the support clause of \texttt{MuIsW} gives \(T.\psi(w)=1\).
This proves all twelve fields and both equalities (NC.36).

No equality of the coefficient indices \(C.k'(w)\) and
\(T.k(w)\) is required by this endpoint. Nor does it identify
the Galois character \(\psi_{0,w}\) of NC12 with the Hecke
character \(T.\psi(w)\) without a supplied reciprocity
compatibility. They are different types of data.

\paragraph{NC17: the same-object join and the precise remaining supplier boundary.}
Combine NC12 with NC15--NC16 for the same \(D=X.toD7\).
Suppose the twisted packet has
\[
 t_w.Q=T.Q(w)\qquad(w\in W=W_{\rm tw}).
 \tag{NC.37}
\]
Taking the \(Q\)-component of (NC.29) proves (NC.35).
Thus NC12 and NC16 construct simultaneously
\[
 \begin{gathered}
 C:\texttt{NChoice}(D),\qquad \texttt{NChTie}(D,C),\\
 \mathrm{tp}:\texttt{TupTied}(D,\texttt{Fr.at}(X,f),C),
 \qquad \mathrm{tp}.\xi=\xi,\quad \mathrm{tp}.\psi=T.\psi.
 \end{gathered}
 \tag{NC.38}
\]
The component entries \(C.Q,C.k',C.sd,C.twI\) on \(W\)
are the ones in \(t,b\). Off \(W\), the empty-prime,
exponent-zero, trivial-character packet (NC.28) supplies
the unconstrained entries.

For an instance of the previously reconstructed TX16 endpoint,
also retain the literal frame equality
\(I.\texttt{cf}=\texttt{Fr.at}(X,f)\) and
\(\xi=I.\texttt{xi}\), together with its already stated
geometric and stage premises. Equation (NC.36) then gives
\(\mathrm{tp}.\xi=I.\texttt{xi}\) on the same objects,
so (NC.38) supplies its \texttt{NChoice}/\texttt{TupTied}
arguments. TX16's geometric conclusion is therefore
available under these more explicit arithmetic premises.
No new comparison of independently chosen frames or
character families is silently inserted.

The complete conditional suppliers proved here are
\texttt{exists\_\allowbreak coeff\_\allowbreak ext},
\texttt{ReadsAt.nonempty\_\allowbreak reading},
\texttt{exists\_\allowbreak formsIso},
\texttt{choiceOf}/\texttt{choiceOfPack},
\texttt{exists\_\allowbreak choice},
\texttt{exists\_\allowbreak tied\_\allowbreak choice\_\allowbreak of\_\allowbreak reads}, and
\texttt{tupTied\_\allowbreak of\_\allowbreak muIsW}, including their finite-support
and normalization consequences. For \texttt{exists\_\allowbreak choice}
alone, \texttt{IsCharOf} and \texttt{LocsCarry} are unnecessary:
use arbitrary readings from NC4 before NC9.

The remaining simultaneous supply is explicit:
the actual \texttt{D7} row and local models;
twisted packets and sides which satisfy \texttt{ReadsAt}
and \texttt{RankTie}; \texttt{ReadElsewhere} or its
\texttt{TwCover} supplier; a common family satisfying
\texttt{IsCharOf} and the same-object \texttt{LocsCarry};
and \texttt{MuIsW} with the exact prime equality (NC.37)
and, where TX16 is used, the stated frame and \(\xi\)
identifications.
NC11 can supply \texttt{IsCharOf} under its primitive-root,
conjugation, and \(L_0/\mathbb Q\)-Galois hypotheses,
but it does not supply \texttt{LocsCarry} for that new family.
All constructions are relative to these data; this section
does not prove a full \texttt{Later}, \texttt{Pins}, host
window, or final RH/GRH existence theorem.


\section{Restricted lattices, auxiliary stages, and coordinate pins}
\label{sec:stage-existence}

\status{SE1--SE16 prove the scalar-descent and restricted-lattice
construction of the literal divisibility stage, the finite-dimensional
row model, the ordinary local and spectral-stage model, the compact-dual
Haar normalization and the jet-stage model, and coordinate normalization
of the two pins. The precise remaining condition in the assembled
\texttt{Later} constructor is an \texttt{NStage} on the same adapted row
object. Its local period-ring and Galois suppliers are not declared
proved here. These existence statements do not assert the quantitative
spectral, metric, cohomological or prime-window properties needed later.}

In SE1--SE6, $K$ is a number field when integral lattices or infinite
places are used, $\sigma:K\hookrightarrow\mathbb C$ is fixed, and a
$\sigma$-semilinear map satisfies $e(ax)=\sigma(a)e(x)$.
For a Hermitian lattice $A=(V,b,X,L)$, our convention is
\[
 \|x\|_{A,v}=\|X_v^{-1}\sigma_v(b^{-1}x)\|_2.
 \tag{SE.1}
\]
The word ``good'' means that every $X_v$ is invertible, $L$ is a finitely
generated $\mathcal O_K$-module, and $KL=V$. It does not assert that $L$
is free over $\mathcal O_K$.

\paragraph{SE1: scalar extension preserves every independent family.}
Let $b_1,\ldots,b_n$ be a $K$-basis of $V$, and suppose
$e(b_1),\ldots,e(b_n)$ are complex-linearly independent in $M$.
Put $M_0=\operatorname{span}_{\mathbb C}e(V)$. Every semilinear map
$g:V\to N$ extends to a complex-linear map on $M$: prescribe
$e(b_i)\mapsto g(b_i)$ on $M_0$ and extend by zero on a vector-space
complement. This gives $F(e(x))=g(x)$ because
$e(x)=\sum_i\sigma(b_i^*(x))e(b_i)$.

In particular, $e$ is injective. If $z_1,\ldots,z_t$ are independent over
$K$, extend their coordinate functionals from their span to $V$ and put
$g(x)=(\sigma(z_1^*(x)),\ldots,\sigma(z_t^*(x)))$.
An extension $F$ sends $e(z_j)$ to the $j$th standard vector. Applying
$F$ to a complex linear relation proves that the $e(z_j)$ are independent
over $\mathbb C$. This argument does not require $\sigma(K)$ to be
topologically dense, or $e(V)$ to be a complex subspace.

\paragraph{SE2: a finite common kernel descends to the coefficient field.}
Let $T_s\in\operatorname{End}_{\mathbb C}(M)$ for $s$ in a finite set
$S$, and assume $T_s(e(V))\subseteq e(V)$ for every $s$. Then
\[
 M_0\cap\bigcap_{s\in S}\ker T_s
 =\operatorname{span}_{\mathbb C}
 e\left(\{x\in V:T_s e(x)=0\text{ for all }s\}\right).
 \tag{SE.2}
\]
Injectivity in SE1 gives unique $K$-linear $N_s$ with $eN_s=T_se$.
Combine these into $N:V\to V^S$, and choose a right inverse of
$N:V\to\operatorname{im}N$. Extend that right inverse to a linear
$r:V^S\to V$. The map $p=1-rN$ takes values in $\ker N$ and is the
identity there. If $\iota_s:V\to V^S$ is the $s$th coordinate inclusion,
\[
 1_V=p+\sum_s r\iota_sN_s.
 \tag{SE.3}
\]
Extend $ep$ and each $er\iota_s$ by SE1 to complex-linear maps $P,R_s$.
Equation (SE.3), first on the $e(b_i)$ and then on their complex span,
gives $1_{M_0}=P+\sum_sR_sT_s$. Also $P(M_0)$ is contained in the
right side of (SE.2). Thus a common-kernel vector $z\in M_0$ equals
$Pz$ and belongs to that right side. The reverse inclusion follows from
complex linearity of the $T_s$. The finite index set is essential to
this displayed finite-sum proof; no assertion about arbitrary infinite
intersections is being substituted for it.

\paragraph{SE3: the restricted metric has explicit frames.}
Let $U\subseteq V$ be a $K$-subspace with basis $c_1,\ldots,c_t$. For
each infinite place $v$, let $E_v$ be the matrix whose columns are
$\sigma_v(b^{-1}c_j)$. By SE1, $E_v$ has full column rank. Hence
\[
 A_v=(X_v^{-1}E_v)^*(X_v^{-1}E_v)
 \tag{SE.4}
\]
is positive definite. A positive definite Hermitian matrix is unitarily
diagonalizable with positive eigenvalues; taking their positive square
roots gives an invertible $B_v$ with $B_v^*B_v=A_v$. Set $Y_v=B_v^{-1}$.
For $x=\sum_j a_jc_j$,
\[
 \|Y_v^{-1}\sigma_v(a)\|_2^2
 =\sigma_v(a)^*A_v\sigma_v(a)
 =\|X_v^{-1}\sigma_v(b^{-1}x)\|_2^2.
 \tag{SE.5}
\]
Taking nonnegative square roots proves exact equality of norms, at
every infinite place. The empty matrix treats $U=0$: its determinant
is $1$ and both norms are zero. No numerical approximation to a matrix
square root enters the statement.

\paragraph{SE4: intersection with a rational subspace is a full lattice.}
Give $U$ the lattice
$L_U=\{x\in U:x\in L\}$ and the frames of SE3. Its image in $V$ is
exactly $L\cap U$. Since $\mathcal O_K$ is Noetherian, $L\cap U$ is a
finitely generated submodule of the finitely generated module $L$.
For $x\in U$, write $x$ as a finite $K$-linear combination of vectors
of $L$. Clear all denominators by some
$0\ne a\in\mathcal O_K$. Then $ax\in L\cap U$ and
$x=a^{-1}(ax)$ belongs to $KL_U$. Thus $L_U$ spans $U$, and together
with SE3 it is a good Hermitian lattice. This proof uses neither a
principal-ideal assumption nor a global integral basis for $L_U$.
The standard lattice $\mathcal O_K^d\subseteq K^d$ with identity
frames is good for the same reasons, including $d=0$.

\paragraph{SE5: the ramified subspace is the extension of its actual form.}
Fix the source's member forms $M$ and column data $A$. Its embedding
$e=M.\mathrm{emb}(0)$ has complex-independent basis image spanning
$A.\mathrm{Mu0}$, by \texttt{emb\_\allowbreak basis} and \texttt{Fi\_\allowbreak zero}.
The operators $\rho_v(a)$ preserve $e(V)$: they preserve the embedded
integral lattice by \texttt{rhoT\_\allowbreak lat}, and every vector is a
$K$-linear combination of lattice vectors. Their finite sums
$\mathrm{Nm}_v=\sum_a\rho_v(a)$ therefore also preserve $e(V)$.

Define
\[
 U=\{x\in V:e(x)\in A.\mathrm{Mram}\},\qquad
 A.\mathrm{Mram}=A.\mathrm{Mu0}\cap\bigcap_v\ker\mathrm{Nm}_v.
 \tag{SE.6}
\]
The local-prime index set is finite. SE2 proves that $e(U)$ spans this
exact complex ramified subspace. SE1 proves that every $K$-basis of $U$
maps to a complex-independent family there. Consequently the source's
\texttt{ramEmb} is a scalar-extension presentation of the actual
\texttt{Mram}; equality of dimensions alone is not used as an
identification. Its lattice is the actual intersection
$M.\mathrm{ramL}=M.\mathrm{lat}(0).L\cap U$ from SE4.

\paragraph{SE6: the literal divisibility-stage record exists.}
For every $B,C$ and every $W:\texttt{RowsRest}(B,C)$ there is a
$V:\texttt{DivStage}(B,C,W)$ whose lattices are all good and such that
\[
 \nu_u\bigl(\mathrm{Nu}(u).L\bigr)=(C.\mathrm{fm}(u)).\mathrm{ramL}.
 \tag{SE.7}
\]
Take $\mathrm{RR}=K^{D_{\rm rows}}$ with its standard lattice and
identity frames. For each $u$, take $\mathrm{Nu}(u)$ to be precisely
the subspace, intersection lattice and restricted frames of SE3--SE5.
Let $j_{EE}$ map standard coordinates through $W.\mathrm{onb}$, let
$\nu_u$ be the subspace inclusion, and set
\[
 j_{N,u}=(W.e\mathrm{Mu}_u)^{-1}\circ\mathrm{ramEmb}_u.
 \tag{SE.8}
\]
The full field check is as follows. The fields \texttt{div\_\allowbreak good},
\texttt{jEE\_\allowbreak norm} and \texttt{nuIn\_\allowbreak norm} follow respectively from
invertible frames, orthonormal coordinates and (SE.5).
\texttt{jEE\_\allowbreak basis} is the orthonormal basis assertion;
\texttt{jN\_\allowbreak basis} is SE5 transported through the displayed linear
equivalence. The fields \texttt{nuIn}, \texttt{nuIn\_\allowbreak inj} and
\texttt{nuIn\_\allowbreak range} are the inclusion, its injectivity and its range.
\texttt{nuIn\_\allowbreak jN} is (SE.8) followed by $W.e\mathrm{Mu}_u$.
Use exactly $W.D_{\rm all}$ and $W.\mathrm{delp}$ for
\texttt{Dall\_\allowbreak eq} and \texttt{delP\_\allowbreak eq}. The \texttt{div} data are
completed by
\[
 B_k=0,\quad h=0,\quad\eta_F=0,\quad
 HH=s\mathrm{Frame}(\mathrm{Nu},0)+\mathrm{RR}.\mu_{\max}.
 \tag{SE.9}
\]
Then \texttt{HH\_\allowbreak ge} is equality, and the bound on \texttt{delP}
is the bound already in $W$. SE4 gives \texttt{AllGood} and (SE.7).
This accounts for all fourteen \texttt{DivStage} fields. The zero
$B_k$ in this existence construction does not provide a later estimate
for independently prescribed nonzero maps.

\paragraph{SE7: the actual finite class data supply \texttt{Fits}.}
For an actual period datum $Q$ with $Q.\mathrm{Ok}$ and $Q.H(b)$ the
full group, evaluate an equivariant form on one representative of every
class. This defines a linear map
\[
 Q.M_t\longrightarrow\prod_{c\in Q.S'.\mathrm{Cl}} VV_c.
 \tag{SE.10}
\]
The class set and the polynomial representation spaces $VV_c$ are
finite. If all evaluations vanish, equivariance and the specified
movement from the representative to any point force all values to
vanish; thus (SE.10) is injective and $Q.M_t$ is finite dimensional.

Now fix the actual $I$ and $l:I.\mathrm{Laws}$, and put $C=I.\mathrm{cols}(l)$.
The source period spaces are independent of the word $u$ at fixed
twist $\mu$, with their stated common polynomial type. Set
$d(\mu)=\dim_{\mathbb C}(I.\mathrm{perIn}(u_0,\mu).M_t)$.
Choose coordinate linear equivalences to $\mathbb C^{d(\mu)}$;
these give both \texttt{mt0} and \texttt{mtTw} on this same $C$.
By TL7 (or the weighted positive class-sum calculation FC7),
\[
 \operatorname{ip}_u(x,y)=
 \sum_c w_c\langle x(c),y(c)\rangle_c
 \tag{SE.11}
\]
is conjugate-linear in $x$, linear in $y$, Hermitian, and positive
definite on $\mathrm{Mu0}$: every $w_c>0$, and a nonzero equivariant
form has a nonzero value on some representative. Restriction to
$\mathrm{Mram}$ gives \texttt{herm}, \texttt{pos}, \texttt{defn}.
All five fields of $\texttt{Fits}(C,d)$ are supplied. We retain the
literal $I.\mathrm{Laws}$ and finite-class hypotheses.

\paragraph{SE8: the row record has an explicit finite-dimensional model.}
For every $C$,
\[
 \texttt{Nonempty}\,\texttt{RowsRest}(B,C)
 \quad\Longleftrightarrow\quad
 \exists d:\ C.\mathrm{Tup}\to\mathbb N,\ \texttt{Fits}(C,d).
 \tag{SE.12}
\]
For the forward implication use $d(\mu)=\dim W.\mathrm{Amb}(\mu)$,
the two equivalences $W.eM_{t,0},W.eM_{t,\mu}$, and the isometry
$W.e\mathrm{Mu}$ to recover the five fields of \texttt{Fits}.

For the reverse implication choose arbitrary nonnegative $n,n_0,a$.
Use $EE=\mathbb C^n$, $EE_0=\mathbb C^{n_0}$ with their standard
orthonormal bases, $D_{\rm all}=2n$, $D_{\rm rows}=n$,
$\mathrm{delp}=1/2$, and a singleton $FT$. Set
$\mathrm{Amb}(\mu)=\mathbb C^{d(\mu)}$ and
$\mathrm{Mult}(\mu,0)=\mathbb C^n$, with all positive-level
$\mathrm{Mult}(\mu,j)=\mathbb C^{d(\mu)}$.
Every transition $io$ is coordinate truncation or extension by zero;
the top identification uses $m\ge2$.
Set $\mu\mathrm{At}=\mu_0$ and $\mu\mathrm{Tw}=\mu_\star$.

Take $\mathrm{Mu}(u)$ to be the actual $\mathrm{Mram}$, endowed with
(SE.11), and $e\mathrm{Mu}$ the identity. Transfer the actual base
period $P\widetilde{\mathrm{Col}}$ along the chosen \texttt{mt0}
equivalence; for other twists transfer the actual twisted period along
\texttt{mtTw}, including the given $e0\mathrm{Tw}$ identification.
Define $B_u$ by the base period after coordinate truncation
$\mathbb C^n\to\mathbb C^{d(\mu_0)}$. Its two linearity fields
hold because the finite column chain is a composite of linear maps
and $P\widetilde{\mathrm{Col}}$ is bilinear. This is a complex
bilinear period pairing, whereas (SE.11) is Hermitian.

For the remaining fields take $M_j=M_t=\mathbb C^a$,
$M_{t,K}=K^a$, the coordinate integral lattices, and
$jM_t(x)=(\sigma(x_i))_i$. SE1, together with the spanning standard
vectors, proves \texttt{jMt\_\allowbreak basis} for every $K$-basis; the generators
are images of the standard $K$-vectors, proving \texttt{MtLat\_\allowbreak sub}.
Put $e\mathrm{Op}=eM_t=\theta_P=1$, $\mathrm{isot}=\top$,
$\lambda_P=\mathrm{ze}=1$, $\kappa_{\rm big}=1$,
$\kappa_{\rm small}=-1$, $I_j=jw=\eta_J=0$, and
$\mathrm{shaLen}=1$. Idempotence, self-adjointness, coefficient-field
preservation and nonnegativity now follow directly. The base/twist
period ties hold by the definitions above. These assignments account
for every scalar, operator, lattice, period and metric field of the
record; choosing $n=n_0=a=1$ proves the required existence.

\paragraph{SE9: the ordinary local data and the specified split places.}
Let $E/F$ be the quadratic CM extension in $B$, with its nontrivial
involution $c$, and let $v$ be a finite place of $F$ of residue
characteristic $\ell$. Take
\[
 F_v=\widehat F_v,\qquad E_v=F_v\otimes_F E,
 \quad c_v=1\otimes c,
 \quad O_v=\{x\in E_v:x\text{ integral over }\mathcal O_{F_v}\}.
 \tag{SE.13}
\]
The completion is a characteristic-zero local field: an integral basis
and the $v$-adic norm make it a finite extension of $\mathbb Q_\ell$.
The original field is dense, its normalized valuation is the given
one, and tensoring a basis of $E/F$ gives the stated free rank two and
spanning property. Give $E_v$ its finite-dimensional vector topology
and Borel sigma algebra. The standard character of $\mathbb Q_\ell$
is $x\mapsto\exp(2\pi i\{x\}_\ell)$, where the principal part
$\{x\}_\ell\in\mathbb Z[1/\ell]/\mathbb Z$ is determined by the
finitely many negative digits. Composing with
$\operatorname{Tr}_{E_v/\mathbb Q_\ell}$ gives exactly the character
and principal-part specification in \texttt{LocAt}. These choices
supply all its fields before the optional split datum.

For clarity, the local-field facts just used are ordinary completion
and tensor-base-change facts: $F_v\otimes_FE\simeq\prod_{w\mid v}E_w$,
and $[E_w:\mathbb Q_\ell]=e_wf_w$. They follow by completing an
integral basis, decomposing the completed integral algebra into its
local factors by CRT, and taking fraction fields. If $\ell$ splits
completely in $E$, the equality
$\sum_{w\mid\ell}e_wf_w=[E:\mathbb Q]$, together with the maximal
number of distinct primes, forces every $e_w=f_w=1$. Hence
$F_v=\mathbb Q_\ell$ and $E_v\simeq F_v\times F_v$; $c_v$ swaps
the two factors. This is the required split algebra and involution,
and $\ell$ is its specified uniformizer. Equivalently, for
$E=F(\alpha)$ with $\alpha^2=a$ and $c\alpha=-\alpha$, the two local
roots $\pm\beta$ give $x+y\alpha\mapsto(x+y\beta,x-y\beta)$.

Use \texttt{some} of this split datum exactly when
$\ell\in\{g_0,g_1\}$, and \texttt{none} elsewhere. The host's full
splitting hypotheses supply the required implication at these two
primes. Thus the source's biconditional for the optional datum is
satisfied. It makes no claim that other places fail to split.
This ordinary completion construction is distinct from
\texttt{LocStage}, which also requires Fontaine period-ring data.

\paragraph{SE10: the literal spectral-stage record exists.}
Suppose $B$ has its given positive gap $\Delta$, injective Boolean
letter map and exact $u$- and $v$-weight identities
(\texttt{HasLetters}); suppose (SE.11) is positive definite on
each $\mathrm{Mu0}$; and let $W$ satisfy
$D_{\rm rows}\le D_{\rm all}$. Use $HH=\mathbb C^{D_{\rm all}}$,
and give $Cu(u)=\mathrm{Mu0}(u)$ precisely the metric (SE.11).
If $b_i=W.\mathrm{onb}_i$, the map
\[
 x\longmapsto\sum_{i<D_{\rm rows}}\langle x,b_i\rangle e_i
 \tag{SE.14}
\]
is conjugate-linear under the source convention and norm preserving:
its squared norm is $\sum_i|\langle x,b_i\rangle|^2=\|x\|^2$.
It supplies the field \texttt{sEE}, including the correct semilinearity.

Choose the local data of SE9, $\mathrm{dsig}=\mathrm{cdScale}
=\mathrm{volN}=\mathrm{volH}=1$, singleton \texttt{Term} and
\texttt{Pair}, and the one-block point partition $[(1,r_S+1)]$.
Take $P_u,a\mathrm{Term},G\mathrm{term},s_0,p\mathrm{val},
\mathrm{meanCount}$ to be zero, $c\mathrm{Term}=\mathrm{massPair}=1$,
the predicates \texttt{isId}, \texttt{inS}, \texttt{relevant},
\texttt{FeqF}, \texttt{scalarPair} to be true, and
\texttt{BBne}, \texttt{generic} to be false. Use the actual
$\Delta$, letter map and weight identities; they are not replaced
by zero data. Then \texttt{s\_\allowbreak split} is SE9, \texttt{s\_\allowbreak HH} is its
dimension, \texttt{s\_\allowbreak Cu} and \texttt{s\_\allowbreak Cu\_\allowbreak isom} are the identity
with the chosen metric, \texttt{s\_\allowbreak PlF} is the identity on finite
places, and \texttt{dsig\_\allowbreak pos} is $1>0$. Together with (SE.14) and
the displayed \texttt{s} this checks all nine \texttt{SStage} fields.
In particular both term types are nonempty. This construction imposes
no nonzero spectral contribution.

\paragraph{SE11: compact-dual Haar normalization is constructible.}
Put $J=\operatorname{diag}(1_P,-1_Q)$ and
$\mathfrak u_{P,Q}=\{X:X^*J+JX=0\}$. The real-linear map
\[
 \mathfrak u_{P,Q}\times\mathfrak u_{P,Q}\longrightarrow M_{P+Q}(\mathbb C),
 \qquad(X,Y)\longmapsto X+iY
 \tag{SE.15}
\]
is an isomorphism. Its inverse sends $Z$ to
\[
 X=\tfrac12(Z-JZ^*J),\qquad
 Y=-\tfrac i2(Z+JZ^*J).
 \tag{SE.16}
\]
Both satisfy the defining Lie-algebra equation, and $X+iY=Z$.
Start with any Lebesgue measure $\lambda_0$ on this real vector space.
Its product pushed through (SE.15) is $a$ times entrywise complex
Lebesgue measure for some finite $a>0$. Thus
$\lambda=a^{-1/2}\lambda_0$ satisfies the literal
\texttt{IsOmegaLeb} product-measure normalization.

Here is the density fact needed to normalize the group measures. The
matrix exponential maps a neighborhood of $0$ in
$\mathfrak u_{P,Q}$ analytically onto a neighborhood of $1$ in
$U(P,Q)$ and has derivative the identity. One can see this without a
global structural theorem: the convergent local matrix logarithm is
inverse to the exponential, and $g^*Jg=J$ implies
$(\log g)^*=-J(\log g)J$. Translate a positive volume density at the
identity by left multiplication. It defines a smooth left-invariant
measure, finite on compact sets and positive on open sets. Haar
uniqueness identifies every Haar measure with a positive constant
multiple of it. In exponential coordinates its density relative to
$\lambda$ is continuous and has a finite positive value $c$ at $0$.
For the exact entrywise-norm balls
$B_r=\{X:\sum_{ij}|X_{ij}|<r\}$, these coordinates give
\[
 \frac{\mu(\exp B_r)}{\lambda(B_r)}\longrightarrow c\in(0,\infty).
 \tag{SE.17}
\]
Indeed, for any $\epsilon>0$ the continuous density lies between
$c-\epsilon$ and $c+\epsilon$ on all sufficiently small $B_r$;
integration and division by their finite positive volumes prove the
limit. The matrix logarithm also gives injectivity on those balls.
For $P+Q=0$, use the unit point masses; the ratio is identically $1$.

The compact group $U(N)$ is a closed bounded matrix set. Normalize its
Haar measure $p_N$ to total mass $1$, and let $c_N>0$ be its density
in (SE.17). For $Q>0$, multiply any Haar measure on $U(P,Q)$, of
density $d_{P,Q}>0$, by $c_{P+Q}/d_{P,Q}$. For $Q=0$ keep exactly
$p_P$. This yields all four \texttt{CDHaar} fields: a measure at
each signature, Haar invariance, probability in the compact case,
and the same finite nonzero Omega density for $(P,Q)$ and $(P+Q,0)$.
The ordinary finite-dimensional change-of-variables and Haar
existence/uniqueness theorems used in this proof are explicit classical
background; the common-density conclusion is proved here.

\paragraph{SE12: a jet-stage model with all prescribed ranks.}
Given a positive volume function, a \texttt{CDHaar}, a
\texttt{ThetaJet}, the positive \texttt{GlobalExtra} scalars, and
nonempty adjacent class sets for each of its situations, the literal
\texttt{GJStage} exists. For each situation choose the word and its
two classes. A prime of $\mathcal O_K$ above the situation's rational
prime exists by lying over, and its residue characteristic is that
prime. A fixed element of the nonempty $\mathrm{Tau}$ supplies each
required archimedean label.

For every required $u,w$, let
$d=\dim_{\mathbb C}(C.\mathrm{mem}(u)).A.\mathrm{Mu0}$ and
$\ell=\mathrm{resChar}(w)$. Use $\mathbb Z_\ell\subset\mathbb Q_\ell$
as the coefficient DVR, no variable places, the one-element level
group acting trivially, the trivial character, and $d$ components.
A polynomial in no variables is a constant. Thus the integral
automorphic module in each component is $\mathbb Z_\ell$ and the
whole module is $\mathbb Z_\ell^d$; its constant-coefficient map is
an isomorphism. The class set is a singleton in each component, so
\texttt{FinClasses} holds, and extension to $\mathbb Q_\ell$ gives
exact rank $d$. This proves the source's prescribed-rank
\texttt{GAtP} supplier, including rank zero. These constructions
fill \texttt{gat}, \texttt{gat\_\allowbreak d}, all jet-place and class fields;
the other fields are the supplied volumes, Haar and extra scalars.

In particular choose volume $1$, all extra scalars $1$, both situation
types empty, and the \texttt{CDHaar} of SE11. No adjacent classes or
inert-prime choices are then required. This proves
$\texttt{Nonempty}\,\texttt{GJStage}(B,C)$ for every $B,C$.
The optional source conclusion using its specific
$\mathrm{baseVol}=\mathrm{trueVol}$ additionally needs the positivity
and identification of that specific volume function; it is not
inferred from the constant-volume model. Likewise nonempty jets and
actual lift estimates require their stated additional data.

\paragraph{SE13: zero idempotents and the exact adapted row object.}
Let $r\ge2$. At each selected prime, work over its finite residue field
$\mathbb F_{g_k}$. Among monic degree-$r$ polynomials over that field
with nonzero constant term, distinguish $f_0=(X-1)^r$.
On any complex vector space define
\[
 e_f=\begin{cases}1&f=f_0,\\0&f\ne f_0.\end{cases}
 \tag{SE.18}
\]
These are central orthogonal idempotents, their support is finite,
and their sum is $1$. A good polynomial in the source's good tuple
is squarefree; $f_0$ is not squarefree since $r\ge2$. Hence every
good-polynomial idempotent is zero. The set of chosen split primes
is nonempty ($f>0$), so the product for a good tuple is zero, and
the sum $e_0$ over good tuples is zero. This proves the exact
idempotent calculation used by
\texttt{idempotent\_\allowbreak good\_\allowbreak tuples\_\allowbreak eq\_\allowbreak zero}.

The source's \texttt{adapt} is a change of row object:
\[
 e\mathrm{Op}=0,\quad eM_t=0,\quad\eta_J=0,
 \quad\mathrm{shaLen}(p,e)=
 \begin{cases}(\operatorname{shaLen}(L_p,e).\mathrm{toNat}:\mathbb Q)
                       &p\text{ prime},\\0&\text{otherwise}.
 \end{cases}
 \tag{SE.19}
\]
It preserves the other fields. Zero is self-adjoint, idempotent and
preserves the coefficient-field image; the new Sha scalar is
nonnegative by its definition. Thus these updates preserve the
\texttt{RowsRest} type. When the source's trivial central actions are
used, (SE.18) makes the \texttt{IdempotentMatchesAction} condition
hold for these zero operators. No finiteness of the Sha group follows
from (SE.19): extended-natural infinity is also sent to $0$ by
\texttt{toNat}. This lemma proves the idempotent and adaptation steps;
it does not construct the full \texttt{NStage} or its local period rings.

\paragraph{SE14: the five-stage join retains one precise input.}
Fix $I$, $l:I.\mathrm{Laws}$ and \texttt{HasLetters} for the same
$B$. Let
\[
 C=I.\mathrm{cols}(l),\qquad
 W_0=\mathrm{RowsRestModel.model}(\mathrm{fits}(l),1,1,1),
 \qquad W=\mathrm{adapt}(W_0).
 \tag{SE.20}
\]
Here the Sha scalar used in \texttt{adapt} is supplied as in SE13;
its defining cyclotomic/Galois data remain explicit if needed.
If $N:\texttt{NStage}(B,C,W)$ is supplied, then
$\texttt{Later}(B,C)$ is inhabited. Indeed, SE7--SE8 supply $W_0$,
SE13 supplies $W$ with $D_{\rm rows}=1,D_{\rm all}=2$,
SE6 supplies $V:\texttt{DivStage}(B,C,W)$, and SE10 supplies
$S:\texttt{SStage}(B,C,W)$ using $1\le2$ and the same class metric.
SE11--SE12 supply $G:\texttt{GJStage}(B,C)$. The tuple
$(W,V,S,N,G)$ is exactly the five-field \texttt{Later} record.

The dependency on this same $W$ matters. An $N$ for a different row
object cannot be substituted by a dimension comparison. The source's
remaining \texttt{exists\_\allowbreak nStage\_\allowbreak adapt} supplies such an $N$ from
\texttt{LocStage} and its other constructions; only the idempotent
part of that supplier has been proved above. In particular the
\texttt{LocStage} fields include actual $A_{\inf}$, $B_{\rm cris}$,
$B_{\rm dR}$, actions, filtration, and tied comparison data, far
beyond SE9's ordinary completions. These remain named targets.

There is a concrete reason to keep the later quantitative hypotheses
separate: in (SE.20), $M_t=\mathbb C$ but $eM_t=0$, so
$\operatorname{rank}(eM_t)=0$ on a one-dimensional space. Thus the
bare row record and its adaptation cannot imply any positive lower
bound for that projector rank. A later theorem with such a conclusion
must use further premises. The same caution applies to the zero
spectral terms in SE10 and the empty jets in SE12.

There is also a concrete arithmetic compatibility obstruction for this
particular bare-stage model. The source's \texttt{nTw} sets
$\mathrm{Tw}=\varnothing$. But the \texttt{NChoice} record requires
\[
       \mathrm{twI}:B.\mathrm{places}.\mathrm{Pl}\longrightarrow
                         \mathrm{Tw}
 \tag{SE.24}
\]
at every coefficient place, including places where $\mu\mathrm{At}=\mu_0$.
The domain is nonempty: lying over supplies a prime of the coefficient
number field above $2$, as in SE12. Thus (SE.24) cannot exist for this
specific empty-twist model. Equivalently its \texttt{ReadElsewhere}
condition fails even though its nonreference support is empty. This
does not rule out other compatible \texttt{NStage}/\texttt{Later}
constructions; it proves that the displayed bare construction cannot
by itself supply NC's arithmetic choice interface.

\paragraph{SE15: diagonal coordinates supply the auxiliary frames.}
Suppose the source ranks are fixed, $E$ has its CM involution, $T$ is
prime and $p_z\in E$ satisfies $\overline{p_z}p_z=T$. Take
$X_j=E^{N_j}$ and $X'=E^{\mathrm{rank}X'}$ with positive standard
forms and coordinate inclusions. For $j<m$ put
$h=\lfloor(m-1-j)/2\rfloor$. The rank and positive index of $Y_j$
are $N_j+(m-1-j)$ and $N_j+h$. Its first $N_j$ coordinates are
$X_j$; the next $h$ positive and $h$ negative coordinates give
\[
 u_i=e_i+f_i,\qquad v_i=(e_i-f_i)/2,
 \quad\langle u_i,u_k\rangle=\langle v_i,v_k\rangle=0,
 \quad\langle u_i,v_k\rangle=\delta_{ik}.
 \tag{SE.21}
\]
They are orthogonal to $X_j$. If an extra negative coordinate is
present, put $p_z$ there: its norm is $-T$ and it is orthogonal to all
these vectors. At the top, put $z=p_z e_{\mathrm{rank}Y_m}$ in the
extra positive coordinate; it has norm $T$ and is orthogonal to the
embedded $Y_m$. All vectors and orthogonality identities are obtained
by multiplying diagonal entries. This proves the actual
\texttt{spacesNorm\_\allowbreak hasFrame}, coordinate-form identities, top norm,
signatures, $V\mathrm{neg}=\mathrm{false}$, and $T$ preservation.
For completeness the remaining coordinate spaces are
$W_j=E^r$ with diagonal signature $(k+j-1,m-j)$ for $1\le j\le m$,
and $V_j=E^{n_V}$ with the positive diagonal form. Their bases and
inclusions are the standard coordinate ones, completing the
\texttt{spacesNorm\_\allowbreak coord} assertions.

For any $n,N$ and matching coordinate indices $a,b$, extension by
zero satisfies $\mathrm{extTo}_{n,N}(e_a)=e_b$: evaluate each
coordinate, where both sides are $1$ at the matching index and $0$
elsewhere. Applying this to $X'\hookrightarrow X_m$ is precisely
the \texttt{Pins.iX} equation, with its fixed frame indices.

\paragraph{SE16: normalizing an existing input supplies both pins.}
Start with the source's already supplied $P$, $X:\texttt{PreX}(P)$,
oriented frame data, letter identities and the required size/
$du_\star$ properties. Write $\pi_u=X.\pi_T(u)$; then
$\pi_u\overline{\pi_u}=T_u$ and every $\pi_u\ne0$, since $T_u$
is prime. Apply SE15 with $T=T_{u_0}$ and $p_z=\pi_{u_0}$, and
replace the space fields by these diagonal spaces. Define
\[
 c(u)=\pi_{u_0}/\pi_u.
 \tag{SE.22}
\]
It is nonzero, satisfies
$c(u)\overline{c(u)}=T_{u_0}/T_u$, and, by direct cancellation,
\[
 c(u)\pi_u=\pi_{u_0}.
 \tag{SE.23}
\]
Thus (SE.23) is \texttt{Pins.cof}; the coordinate computation of
SE15 is \texttt{Pins.iX}. The preserved number-field, prime, weight
and word fields transport $X$, its orientation, the letter identities,
and the stated $du_\star$ equalities exactly as in
\texttt{ExistsOkPins.exists\_\allowbreak pre}; all changed space fields have
been checked in SE15. The scalar norm relation above is the changed
\texttt{cOf\_\allowbreak eq} field. Subsequent coefficient-field and base
extension can carry these coordinate equalities unchanged.

This is a normal-form construction for an existing arithmetic and
geometric input. The existence of the initial sized
\texttt{PreX}/oriented input, its analytic prime windows, and the
later choice of compatible $I.\mathrm{Ok}$ are not consequences of
the two cancellation and coordinate equations. Accordingly the
ledger closes their normalization steps while retaining those exact
existence suppliers.

\medskip
\noindent\textbf{Source and scope record.}
The companion provenance ledger gives declaration-level pinned Git
ranges and hashes for every SE claim. SE1--SE6 follow the substantive
\texttt{LatticeEmbeddingMemUnitGroup} proof; SE7--SE8 reconstruct
\texttt{fits} and \texttt{ModelLamPDef}; SE9--SE10 supply
\texttt{ExistsInnerProductSpace}; SE11 proves the four-field
\texttt{CDHaar} conclusion by an ordinary Lie-coordinate argument;
SE12 supplies the exact bare \texttt{GJStage}; SE13--SE14 leave the
full \texttt{NStage}/\texttt{LocStage} supplier open; SE15--SE16
reconstruct the coordinate and scalar pin normalization. None of
these constructions identifies an auxiliary stage with an externally
specified geometric or spectral realization without the relevant tie.


\section{Arithmetic compatibility of the actual choices}
\label{sec:arithmetic-compatibility}

This section supplies the rank and character comparisons upstream of NC.
Throughout, $D$ is one fixed datum, $E=D.\mathrm{host}.E$,
$L=D.\mathrm{what}_2.LL$, $L_0=D.LL0$, and
$\iota_L:L\hookrightarrow\mathbb C$ is its specified embedding.
The notation $\bar v$ denotes the specified conjugation on coefficient
places, while conjugation of primes of $E$ is written $\bar P$.
The actual coefficient rings, period modules and realization maps of $D$
are retained. In particular, a new auxiliary row model cannot be substituted
for the given $D$ when applying a rank comparison.

Write $G=G_E$, $G_f=U(E)(\mathbb A_f)$, and $\Gamma=U(E)(E)$ in the
source's standard positive definite unitary realization. For a twist index
$i$, put $K_{i,k}=\operatorname{Frac}(O_{i,k})$. A reading $\rho$ has an
injective map $\rho_L:L\to A_\rho=\operatorname{Frac}(O_\rho)$ and an
injective map $d_\rho:O_{i,k}\to A_\rho$. The expression
$\rho.\mathrm{Pulls}(\psi,\psi_0)$ means the literal equality
\[
 \rho_L(\psi_0(g))=d_\rho(\psi(\operatorname{art}_{i,Q}(g)))
 \qquad(g\in G).                                      \tag{RC.1}
\]
The normalized packet is
$p(i,b,Q,k,\psi)=(i,Q,k,\psi^{\epsilon_b})$, with
$\epsilon_{\mathrm{true}}=1$ and
$\epsilon_{\mathrm{false}}=-1$. Applying this side operation twice gives
$\psi$ itself. All Frobenius conventions below are those of the source:
an arithmetic Frobenius $\sigma_P$ has
$\chi(P)=r(\sigma_P)^{-1}$.

\paragraph{RC1: finite evaluations prove scalar extension for the form space.}
Let $K\hookrightarrow\mathbb C$ be an embedding $e$. Let $W_K,W_C$ be
finite-dimensional representations of $\Gamma$, with bases on the same
finite index set, such that every matrix of $\rho_C(\gamma)$ is the
entrywise image under $e$ of the corresponding matrix of $\rho_K(\gamma)$.
This is the exact predicate \texttt{WtMatch}; equality of dimensions alone
is insufficient. Let $J\le G_f$ and $\chi:J\to K^\times$. Assume that
$J$ contains a positive congruence subgroup on which $\chi=1$.
For
\[
 \mathcal A_K=\{F:G_f\to W_K:
 F(\gamma g)=\rho_K(\gamma)F(g),\
 F(g\kappa)=\chi(\kappa)F(g)\},
\]
assume $\dim_K\mathcal A_K<\infty$. Then the same equations over $\mathbb C$,
with character $e\chi$, define precisely
$\mathbb C\otimes_{K,e}\mathcal A_K$. In particular,
\[
 \dim_K\mathcal A_K=\dim_{\mathbb C}\mathcal A_C.          \tag{RC.2}
\]

First, any finite-dimensional subspace $V$ of functions $S\to W$ is
detected by finitely many evaluations. Starting with $V$, if the current
common kernel is nonzero, choose a nonzero function in it and a point where
it is nonzero. Adding that evaluation strictly decreases the kernel
dimension. After at most $\dim V$ steps the kernel is zero.
Choose such a finite set $S_0$ for $\mathcal A_K$.

The common weight bases identify the coefficientwise map
$W_K\to W_C$ and intertwine every $\Gamma$-action. Images of a $K$-basis
of $\mathcal A_K$ are $\mathbb C$-independent: evaluation on $S_0$, followed
by the finite weight-coordinate map, gives a finite matrix of full column
rank over $K$; a nonzero maximal minor remains nonzero under $e$.

For spanning, choose a basis $(\beta_a)$ of $\mathbb C$ over $e(K)$.
For $F\in\mathcal A_C$, expand each weight coordinate in this basis.
Extracting the coefficient of $\beta_a$ gives a $K$-valued function $F_a$.
Because the representation matrices and the level character are defined
over $e(K)$, each $F_a$ satisfies both defining equations and lies in
$\mathcal A_K$. The finitely many coordinates of $F|_{S_0}$ involve only
finitely many $\beta_a$, say those in $A_0$. If $a\notin A_0$, $F_a$ vanishes
on $S_0$, hence everywhere. Thus
$F=\sum_{a\in A_0}\beta_a e(F_a)$ globally, and each $F_a$ is a finite
$K$-combination of the chosen form basis. This proves spanning.
The congruence-kernel hypothesis supplies the smoothness required in
\texttt{NG.TypeSp}: every function with the stated right transformation
law is fixed by that same congruence subgroup. No commutation of scalar
extension with an arbitrary infinite product has been assumed.

\paragraph{RC2: the literal level and weight comparison supplies RankTie.}
Let $m:\texttt{MtReal}(D,\mu,j')$ and $X:\texttt{NExpTied}(D)$.
For an admissible twist $(i,b,k,Q,\psi)$ and $j\in\operatorname{Fin}(D.R.m)$,
put $c=X.\mathrm{mtV}(i,j+1,b,k,Q,\psi)$.
Suppose $e:K_{i,k}\hookrightarrow\mathbb C$ satisfies
\[
 \texttt{LevEq}(m.J,m.\chi,c.J,e\,c.\chi),\qquad
 \texttt{WtMatch}(e,c.\rho,m.\mathrm{wt}.\rho).
                                                               \tag{RC.3}
\]
Then
\[
 \dim_L \mathrm{MtL}(\mu,j')
 =\dim_{K_{i,k}}\mathrm{MtV}(i,j+1,b,k,Q,\psi).             \tag{RC.4}
\]
The local lattice is finitely generated and spans the actual local
fraction-field module, so that module is finite-dimensional.
The map $c.\mathrm{real}$ identifies it with its $K_{i,k}$-form space.
RC1 changes coefficients to $\mathbb C$. The first predicate in (RC.3)
is equality of subgroups and of their character values, so it identifies
the resulting complex form space with the range of $m.\mathrm{real}$.
The latter map is injective and its range is exactly that form space;
therefore its dimension is $\dim_{\mathbb C}\mathrm{Mt}(\mu,j')$.
Finally, the actual \texttt{jMt\_\allowbreak basis} field says that the images of an
$L$-basis under $\mathrm{jMt}$ form a complex basis. This gives (RC.4).

Apply this to $\mu=\mu_{\mathrm{Tw}}(v)$ and $j'=j+1$ for every
$1\le j'\le D.R.m$. It gives \texttt{RankEq} for the unnormalized
character and \texttt{RankTie} for its normalized packet, because the
two applications of $\psi\mapsto\psi^{\epsilon_b}$ cancel.
Neither an abstract dimension estimate nor the existence of an unrelated
row model supplies (RC.3).

\paragraph{RC3: finite-order Hecke characters have the required coefficient-valued Artin character.}
Retain the exact published global reciprocity input: a surjective open
quotient homomorphism
\[
 \operatorname{Art}:\mathbb A_E^\times\longrightarrow G_E^{\mathrm{ab}},
\]
whose kernel is the closure of principal ideles times the connected
infinite component, and whose inertia, uniformizer and decomposition
images are the ones stated in \texttt{IsGlobalArtin}.
Let $\chi$ be a continuous Hecke character with $\chi^n=1$, $n>0$,
and suppose $L$ contains a primitive $n$th root of unity.
There is a character $\psi_0:G_E\to L^\times$, with open kernel and
$\psi_0^n=1$, whose complex image corresponds to $\chi$ under this
specific Artin map. Thus the literal \texttt{GalOf} predicate holds.

The image of $\chi$ is contained in the finite set $\mu_n(\mathbb C)$.
Its kernel is both open and closed. Principal ideles are killed by the
Hecke-character law, and the connected infinite component has connected
image in a finite set, hence is also killed. Closedness of the kernel
therefore kills the entire Artin kernel. Quotient descent gives a
homomorphism $G_E^{\mathrm{ab}}\to\mu_n(\mathbb C)$; the open quotient
property proves its continuity. Compose with abelianization.
The map $\mu_n(L)\to\mu_n(\mathbb C)$ induced by $\iota_L$ is bijective:
both groups consist of powers of the specified primitive root and its
image. Pulling back the character through this bijection gives $\psi_0$,
preserves its kernel, and proves all the assertions.

This construction can be chosen as a function of the pair $(\chi,n)$
alone. That choice is used below to obtain exact equalities of common
characters at places with the same data. It does not require independent
choices of a Galois character at every coefficient place.

\paragraph{RC4: Frobenius detection follows from the stated reciprocity map and weak approximation.}
Let $\varphi:G_E\to B$ be a continuous homomorphism to a Hausdorff
commutative topological group. Suppose that, outside a finite set $S$,
$\varphi(\sigma)=1$ for every valuation above $P$ and every arithmetic
Frobenius $\sigma$ at it. Then $\varphi=1$.

Descend $\varphi$ through topological abelianization and compose with
$\operatorname{Art}$, obtaining a continuous idele character $\theta$.
It is trivial on principal ideles. The uniformizer-image clause of
\texttt{IsGlobalArtin}, with its inverse-Frobenius convention, says that
$\theta$ kills every local uniformizer at each $P\notin S$.
If $\pi$ is one such uniformizer and $x\in E_P^\times$, then
$x\pi^{1-v_P(x)}$ is another; hence $\theta(x)=1$.
Thus every local multiplicative group outside $S$ is killed.

The resulting character depends only on the factors at $S$ and the
infinite places. To justify this assertion without an infinite product
calculation, approximate an idele that is $1$ at those factors by its
finite-coordinate truncations; these converge in the restricted-product
topology, and all are killed. Weak approximation makes the diagonal
$E^\times$ dense in the finite product
$\prod_{P\in S}E_P^\times\times E_\infty^\times$.
Indeed, additive weak approximation applies simultaneously to a finite
set of nonzero coordinates, and the unit locus is open, with continuous
inversion. The induced character on this product is therefore $1$.
Surjectivity of $\operatorname{Art}$ gives $\varphi=1$.
This argument uses the precise global reciprocity input of RC3; it does
not add a Chebotarev hypothesis.

\paragraph{RC5: the common character has exactly the required local behavior.}
Assume $\chi$ is an \texttt{IsPsi} character for $\ell$ and $Q$, and
$\psi_0$ satisfies \texttt{GalOf} as in RC3. Then $\psi_0$ is
anticyclotomic, is ramified exactly at $Q$, and is trivial on the entire
decomposition group above each prime over $\ell$.

Here the source host supplies the quadratic extension $E/F$, its
involution, and the unramified normalization
$v_P(a)=v_{P\cap F}(a)$ for $a\in F^\times$.
The condition $\chi|_{\mathbb A_F^\times}=1$ gives, at every pair of
unramified conjugate primes,
\[
                  \chi(P)\chi(\bar P)=1.                \tag{RC.5}
\]
At a split prime this is the product of the two local values on a
uniformizer of $F$; at a fixed prime the single local value is itself
$1$. The stated valuation normalization is what permits this uniformizer
calculation. Conjugation carries a valuation above $P$ to one above
$\bar P$, preserves the residue cardinality, and sends an arithmetic
Frobenius to an arithmetic Frobenius there. Therefore the continuous
character
$g\mapsto \iota_L(\psi_0(g)\psi_0(cgc^{-1}))$
is $1$ on every Frobenius outside a finite conjugation-stable set.
Its continuity follows from the open kernel of $\psi_0$ and continuity
of conjugation on the profinite Galois group; no continuity of a
coefficient-field isomorphism is needed. RC4 and injectivity of
$\iota_L$ give $\psi_0(cgc^{-1})=\psi_0(g)^{-1}$.

The inertia-image clause of reciprocity says that $\chi$ is unramified
at $P$ if and only if $\psi_0$ kills inertia there. The exact
\texttt{IsPsi.ram} biconditional therefore gives
\texttt{RamExactly}$(\psi_0,Q)$, in both directions.
At $P\mid\ell$, \texttt{IsPsi.atl} kills the whole local
$E_P^\times$. The closure of its Artin image is the abelianized
decomposition image. Since the kernel of the complex character is
closed, it kills that closure; injectivity then gives triviality of
$\psi_0$ on the decomposition group. Merely knowing unramifiedness at
$\ell$ would not prove this last statement.

\paragraph{RC6: actual ring-class descent gives all seven TwOK clauses.}
Fix a reading $\rho$ of $(i,b,k)$ at $v$ and a character $\psi_0$ with
open kernel, $\ell_i^k$-power exponent, anticyclotomic symmetry, exact
ramification $Q$, and decomposition triviality above $\ell_i$.
Assume that $Q$ consists of split places and is disjoint from $S_0(i)$.
Use the specified $\ell_i^k$th root $z_{i,k}\in O_{i,k}^\times$.
For every $g$, $\rho_L(\psi_0(g))$ is a power of $d_\rho(z_{i,k})$.
There is consequently a unique value
$\widetilde\psi(g)\in\langle z_{i,k}\rangle$ with
$d_\rho(\widetilde\psi(g))=\rho_L(\psi_0(g))$.
Injectivity of $d_\rho$ proves multiplicativity and equality of kernels.

Retain the actual \texttt{NTw.art\_\allowbreak univ} field. It factors every
open-kernel, anticyclotomic, finite $\ell_i$-primary character unramified
outside $Q$ through the specified finite group
$\operatorname{RCG}(i,Q)$, when $Q$ avoids primes over $\ell_i$.
That avoidance follows from $Q\cap S_0(i)=\varnothing$, since $S_0(i)$
contains all those primes. Applying the universal property to the finite
cyclic group $\langle z_{i,k}\rangle$ gives
\[
 \psi:\operatorname{RCG}(i,Q)\to O_{i,k}^{\times},
 \qquad \widetilde\psi=\psi\circ\operatorname{art}_{i,Q}.
                                                               \tag{RC.6}
\]
This is the actual descent, not an assumption of (RC.6).

Every descent satisfying (RC.1), not only the one just chosen, has the
seven \texttt{TwOK} properties: split support, disjointness from $S_0$,
open kernel, anticyclotomic symmetry, image in
$\langle z_{i,k}\rangle$, exact ramification $Q$, and triviality on
decomposition above $\ell_i$. The last five follow by applying the two
injective coefficient maps to the corresponding equalities for $\psi_0$.
The root condition follows because all roots of $X^{\ell_i^k}-1$ in
the fraction field are powers of $z_{i,k}$. These clauses imply the
literal \texttt{AdmTw}: $S_0$ also contains the type and auxiliary
exceptional primes, while the exponent and unramifiedness clauses follow
from the displayed cyclic image and exact ramification.
Existence of the ring-class object and its universal map remains part
of the retained \texttt{NTw} input.

\paragraph{RC7: two readings of the common character differ by an allowable power.}
Suppose two readings, possibly on different sides and at different
coefficient places, have the same $i,Q,k$, and their descents $\psi,\psi'$
read the same $\psi_0$. Suppose $L$ contains a primitive
$n=\ell_i^k$th root $z$ and $\psi_0^n=1$. Then
\[
 \psi'=\psi^a\quad\text{for some }a\in\mathbb Z
 \text{ with }\gcd(|a|,\ell_i)=1.                        \tag{RC.7}
\]
For $k=0$, both characters are $1$ by surjectivity of
$\operatorname{art}_{i,Q}$, and $a=1$ works.
For $k>0$, each image of $z$ and each image of $z_{i,k}$ is a primitive
$n$th root in the relevant reading fraction field. Thus changing from
one chosen primitive generator to the other is raising to an exponent
invertible modulo $n$. Write $\psi_0(g)=z^{t_g}$.
Comparing these powers in the two fraction fields gives
$\psi'(\operatorname{art}(g))
 =\psi(\operatorname{art}(g))^a$, with $a$ a product of two such
invertible exponents. Surjectivity of $\operatorname{art}$ gives (RC.7)
on the whole ring-class group.

For normalized packets the exponent is
$\epsilon_{b'}\epsilon_b a$; changing one side negates $a$.
It is still prime to $\ell_i$. This sign is essential: power ambiguity
is a statement about the normalized packet as well as about its
unnormalized character.

\paragraph{RC8: a reading supplies the exact side embedding and a complex extension.}
A reading determines
\[
 \rho_L(x)=\operatorname{Frac}(O_0(i)\to O_\rho)
                  (\operatorname{emb}_{i,b}(x))
       \quad(x\in L_0),                                 \tag{RC.8}
\]
and identifies the nonnegative-valuation predicate at $v$ with
integrality of $\operatorname{emb}_{i,b}(x)$ in $O_0(i)$.
To prove (RC.8), write the fractional image of $x$ with nonzero
denominator in $O_0(i)$ and apply the reading's literal
denominator identity. Fraction-field injectivity clears that denominator.
For the valuation assertion use the following reflection principle:
an injective local map from a valuation ring $A$ to a domain $B$
reflects integrality under its map of fraction fields.
If $t\notin A$, then $t^{-1}$ lies in the maximal ideal of $A$.
If the image of $t$ lay in $B$, the image of $t^{-1}$ would be a unit,
contrary to locality. The local maps needed here are supplied by NC4.
This proves the actual \texttt{SidePl} and \texttt{emb\_\allowbreak side} interfaces.

For a supplied \texttt{TwFrameC}, there is also an embedding
$\iota_\rho:A_\rho\hookrightarrow\mathbb C$ satisfying
$\iota_\rho\rho_L=\iota_L$ on all of $L$.
Here is the cardinality argument. A first-countable Hausdorff completion
of a dense set $S$ has cardinality at most $|S|^{\aleph_0}$: select one
convergent sequence for each point, and uniqueness of limits makes that
selection injective. Hence $|\mathbb Q_\ell|\le\mathfrak c$,
its algebraic closure has cardinality at most $\mathfrak c$, and its
completion $\mathbb C_\ell$ has cardinality at most
$\mathfrak c^{\aleph_0}=\mathfrak c$. The frame embeds $O_{i,k}$ into
$\mathbb C_\ell$. Finite generation of $O_\rho$ over $O_{i,k}$ and
formation of fractions therefore give $|A_\rho|\le\mathfrak c$.

Adjoin $\mathfrak c$ independent variables to $A_\rho$ and take an
algebraic closure. The resulting algebraically closed field has
characteristic zero and transcendence degree $\mathfrak c$ over the
number field $\rho_L(L)$. The same is true of $\mathbb C$ over
$\iota_L(L)$: an algebraic extension of a countable field on $\kappa$
indeterminates has cardinality $\max(\aleph_0,\kappa)$, so its
uncountable cardinality forces the stated transcendence degree.
Choose a bijection of transcendence bases extending
$\rho_L(x)\mapsto\iota_L(x)$, then extend the rational-function-field
isomorphism to algebraic closures and restrict to $A_\rho$.
This is an algebraic embedding; no analytic continuity is asserted.

\paragraph{RC9: coherent readings also give exact complex Hecke readings.}
Suppose families $s(v),Q(v),k(v),\psi_0(v)$ are unchanged under
$v\mapsto\bar v$, each $s(v)$ reads at $v$ on at least one side, and
every reading admits a descent of $\psi_0(v)$.
Assume $L_0/\mathbb Q$ is Galois, as required for the conjugate-reading
construction NC10. There are simultaneous choices
$b(v),\rho_v,\psi_v$ satisfying (RC.1), such that
\[
 b(\bar v)=!b(v)\quad(\bar v\ne v),\qquad
 p(s(\bar v),b(\bar v),Q(\bar v),k(\bar v),\psi_{\bar v})
 =p(s(v),b(v),Q(v),k(v),\psi_v).                          \tag{RC.9}
\]
Choose one representative of each orbit of the involution. At that
representative choose a reading and a descent. On a two-element orbit,
use the conjugate reading and the same normalized packet; unnormalize
with the opposite side. NC10 proves the character-pulling identity,
including the inversion needed for the change of side.
At a fixed place keep the original choice; there is no requirement
that a Boolean side equal its own negation. The common families ensure
that the packet remains in precisely the same $i,Q,k$ fiber.
Unpacking it into a descent uses $\operatorname{psiAt}_b(\operatorname{psiAt}_b(\psi))=\psi$.

If $\operatorname{GalOf}(\chi(v),\psi_0(v))$ and
$\iota_v\rho_v=\iota_L$, set
$e_v=\iota_v\circ\operatorname{Frac}(O_{s(v),k(v)}\to O_{\rho_v})$.
Equation (RC.8) proves \texttt{SideRead} for $e_v$.
Applying $\iota_v$ to (RC.1) proves
\texttt{ReadsHecke}$(e_v,\psi_v\circ\operatorname{art},\chi(v))$
using the same Artin map. These two conclusions concern the same
embedding $e_v$, as required in the later rank comparison.

\paragraph{RC10: Frobenius comparison supplies the nonempty-local character tie.}
This criterion does not assume that $D.nd.\mathrm{Loc}$ is empty.
Let $\psi_0(v)$ have \texttt{GalOf}$(\chi(v),\psi_0(v))$ and
\texttt{IsPsi}$(\ell_v,Q(v),\chi(v))$.
For every local index $a$, suppose the already specified character
$\psi_{L,a}=D.ntw.\mathrm{psiL}(a)$ satisfies
\[
 \texttt{IsGalOf}\bigl(D,a,\mathrm{univ},
               \psi_{L,a},\chi(w(a)),0\bigr).            \tag{RC.10}
\]
Then for every $g\in G_E$,
\[
 \operatorname{embK}_a(\psi_0(w(a))(g))
       =\psi_{L,a}(g)\quad\hbox{in }\operatorname{Frac}(O_a),
 \qquad D.nd.Q(a)=Q(w(a)).                              \tag{RC.11}
\]

Let $j_a=\mathrm{iotaQC}_a\circ\mathrm{iotaQ}_a$ on
$\operatorname{Frac}(O_a)$. Its specified \texttt{iota\_\allowbreak LL} law is
$j_a\operatorname{embK}_a=\iota_L$.
The two complex characters
$j_a(\psi_{L,a})$ and $\iota_L(\psi_0(w(a)))$ have open kernels:
this is \texttt{NTw.psiL\_\allowbreak open} for the first and \texttt{GalOf} for
the second. Consequently they are locally constant and continuous,
even though $j_a$ need not be continuous on the fraction field.
Outside the finite set consisting of primes above $\ell_a$ and the
ramification of $\chi(w(a))$, equation (RC.10) with exponent $0$
and the Artin convention give the same value
$\chi(w(a))(P)^{-1}$ at every Frobenius. Apply RC4 to their quotient.
Injectivity of $j_a$ and \texttt{iota\_\allowbreak LL} give the first equality
in (RC.11).

For any prime $P$ and valuation above it, an inertia element is killed
by one side of that equality if and only if it is killed by the other.
The literal \texttt{NTw.psiL\_\allowbreak ram} and RC5 identify the sets where an
inertia element with nontrivial value exists. Their exact
biconditionals therefore give the conductor-set equality in (RC.11).
Now let the normalized packet at $w(a)$ have index $s(w(a))$, side
$b(w(a))$, and support $Q(w(a))$. If
\[
        \mathrm{twOf}(a)=s(w(a)),\qquad
        \mathrm{twSide}(a)=b(w(a)),                       \tag{RC.12}
\]
then (RC.11)--(RC.12) prove every clause of \texttt{LocsCarry}.
This proves the character and conductor portions from the weaker
Frobenius-value interface. It does not construct (RC.10), or choose
the local index and side compatibly with an independently built stage.

\paragraph{RC11: the exact simultaneous choice exists under its stated supplier hypotheses.}
Retain one $D$, a character frame $cf$, $M:\texttt{MtFam}(D)$,
$X:\texttt{NExpTied}(D)$, $F:\texttt{TwFrameC}(D)$, and the global
reciprocity input of RC3. Assume $L_0/\mathbb Q$ is Galois and
$D.nd.\mathrm{Loc}$ is empty. Supply functions
$\xi:\mathrm{Tup}\to\mathrm{HChar}(E)$,
$\chi(v)$, $Q(v)$, $k(v)$ and $s(v)$ with the following exact properties.

\begin{enumerate}
\item Each $s(v)$ has a reading at $v$ on some side, and
$s(\bar v)=s(v)$.
\item Every \texttt{NChoice} $C$ on this same $D$ with $C.Q=Q$
has a \texttt{TupTied} $tp$ with $tp.\xi=\xi$ and $tp.\psi=\chi$.
\item $\chi(v)$ is \texttt{IsPsi} for $\ell_v,Q(v)$, has order dividing
$\ell_v^{k(v)}$, and $L$ contains a primitive root of that order.
\item $Q(v)=\varnothing$ off the nonreference support and
$Q(v)\cap S_0(s(v))=\varnothing$.
\item Equality of the nonnegative-valuation predicates on $L_0$
at $v,v'$ implies $k(v)=k(v')$ and $\chi(v)=\chi(v')$.
Also $\chi(\bar v)=\chi(v)$ and $k(\bar v)=k(v)$.
\item For every $v,b,\psi$, every proof that $(s(v),k(v),Q(v),\psi)$
is admissible, and every $j\in\operatorname{Fin}(D.R.m)$, assume the
following implication. If $v$ is nonreference, $\ell_{s(v)}=\ell_v$,
and \texttt{TwOK} holds, then \emph{for every} $e$ satisfying both
\texttt{SideRead} and \texttt{ReadsHecke} for $\chi(v)$,
the exact predicates \texttt{LevEq} and \texttt{CellWt} of RC2 hold
for $M(\mu_{\mathrm{Tw}}(v),j+1)$ and this cell.
\end{enumerate}

Then there are $C:\texttt{NChoice}(D)$ and
$tp:\texttt{TupTied}(D,cf,C)$ with
\[
 tp.\xi=\xi,\quad tp.\psi=\chi,\quad C.Q=Q,\quad C.k'=k,
 \quad C.\mathrm{twI}=s,                                \tag{RC.13}
\]
all choices have \texttt{TwOK} for their unnormalized character
$C.\psi'\circ\operatorname{art}$, and all six fields of
\texttt{ChoiceSide} hold.

To prove this, a reading first gives $\ell_{s(v)}=\ell_v$.
The exact ramification condition makes $Q(v)$ a function of $\chi(v)$:
if $\chi(v)=\chi(v')$, their support sets are equal. Choose the character
$\psi_0(v)$ of RC3 as a function of
$(\chi(v),\ell_{s(v)}^{k(v)})$ alone. Thus all four families needed
in RC9 are conjugation-invariant. RC5--RC6 give a descent for every
reading, and RC9 chooses coherent $\rho_v,b(v),\psi_v$.
RC8 gives $\iota_v$ and RC9 gives both reading predicates for the
same $e_v$. Hypothesis 6 now applies, and RC2 supplies the exact rank
equalities. NC9 constructs the actual finite DVR extensions,
integral bases and normalized \texttt{NChoice}; it does not assume the
desired integral form isomorphism. Hypothesis 2 gives $tp$ on this $C$.

Here are the six fields of \texttt{ChoiceSide}, including the ones that
remain nonvacuous when the local index set is empty.
The \texttt{NChTie} local index, side and support equalities are vacuous;
its global common-character, conjugation and denominator identities
are (RC.1), (RC.9) and the reading's denominator law.
The six subfields of \texttt{ChoiceC} are index conjugacy,
\texttt{SidePl}, the fraction-field side equality, opposite sides on
two-element orbits, prime-to-$\ell$ power comparison at a common index,
and equality of conjugate packets. All but the power comparison follow
directly from RC8--RC9.

For that comparison, suppose $s(v')=s(v)$. If the sides agree, RC8
identifies the two nonnegative-valuation predicates on $L_0$.
If they differ, conjugate the reading at $v'$; the side then agrees,
and hypothesis 5 together with its conjugation clause again identifies
$k$ and $\chi$. The chosen $\psi_0$ and $Q$ therefore agree.
RC7 gives exactly the required power relation on normalized packets.
The \texttt{off} field is hypothesis 4. The \texttt{adm} field is RC6.

For \texttt{PsiChoiceCompatible}, if
$\iota_L(y)=\chi(v)(P)^{-1}$ at an allowed Frobenius $\sigma_P$,
the Artin convention gives $y=\psi_0(v)(\sigma_P)$.
Equation (RC.1) is then precisely the required comparison in
$\operatorname{Frac}(C.Oe(v))$.
Finally \texttt{eMt\_\allowbreak level} uses the very $\iota_v$ chosen above and
the subgroup/character equalities supplied by hypothesis 6.
Its literal conclusion is equality of levels and character values
under that complex extension. It does not assert that an arbitrary
integral basis isomorphism constructed by NC9 intertwines every Hecke
operator. Thus no extra equivariance is inferred from rank equality.

\paragraph{RC12: prescribed tuple assignments give an actual MuW row datum.}
Fix $B,f,T$ and one supplied $w:\texttt{ColW}(B,f,T)$.
Let $W_t$ be a conjugation-stable set of coefficient places.
Assume an existing $\texttt{RowsRest}(B,w.C)$; SE8 supplies one under
its stated column hypotheses. The exact range law of $w.I.\xi$ contains
every character $T.\xi_{\mathrm{Tw}}(f,v)$, so choose
\[
 w.I.\xi(m_t(v))=T.\xi_{\mathrm{Tw}}(f,v),\qquad
 m_a(v)=
 \begin{cases}m_t(v)&v\in W_t,\\w.C.\mu_0&v\notin W_t.\end{cases}
                                                               \tag{RC.14}
\]
The right side of the first equality is unchanged by $v\mapsto\bar v$:
the residue characteristic and $T.\psi$ are unchanged. Injectivity
of $w.I.\xi$ therefore gives $m_t(\bar v)=m_t(v)$.
Conjugation-stability of $W_t$ gives the same identity for $m_a$.
Moreover $m_a(v)$ is either $\mu_0$ or $m_t(v)$.
Replace precisely the five tuple-assignment fields
$(\mathrm{muAt},\mathrm{muAt}_{\mathrm{conj}},\mathrm{muTw},\mathrm{muTw}_{\mathrm{conj}},\mathrm{muAt}_{\mathrm{or}})$
of the existing row datum by these values. These are the only row-record
laws involving those assignments, so all other fields are retained.
The three \texttt{MuW} equations are exactly (RC.14) and its two cases.

The source's fuller \texttt{reset} also replaces the two weight arrays
by their prescribed integral values and $\mathrm{shaLen}$ by
$\mathrm{shaTrue}$. Those additional interfaces are separate from this
tuple proof. In particular, changing tuple assignments in an existing
row record does not by itself prove that an already supplied geometric
\texttt{RowPk}, its numerical data, or a local arithmetic stage remains
tied to that changed record.

\paragraph{RC13: the actual row lift preserves MuW and supplies all ten MuIsW fields.}
Now retain the actual $w:\texttt{ColW}$,
$P:\texttt{Pre2}(K,w)$ and $x:\texttt{St}(P)$.
The source row lift $q=P.q$ has $W_0$ with \texttt{MuW} and defines
$W=q.p.\mathrm{what}_2(q.e,W_0,q.n,\eta,h_\eta)$.
Its defining equations are
$W.\mathrm{muAt}=W_0.\mathrm{muAt}$ and $W.\mathrm{muTw}=W_0.\mathrm{muTw}$.
The stage $w.\mathrm{st}(x.o.Z)$ uses this exact $W$.
Consequently it satisfies \texttt{MuIsW} with the original $w.I.\xi$,
$T$ and $P.W_t$.

For completeness, the ten fields have the following individual sources.
Injectivity is $w.ok.\xi_{\mathrm{inj}}$; the range is $w.\xi_{\mathrm{range}}$;
the reference value is $w.l.\xi_0$.
The two distinguished values are $w.ok.\xi_{\mathrm{star}}$ and
$w.ok.\xi_{\mathrm{star'}}$, transported through the literal identity
$w.C_{\mathrm{eq}}$ of the reference frame.
The twist, on-support and off-support clauses are $q.\mu$,
preserved by the two defining equations of $W$ above.
The support-residue and off-support character clauses are $P.\mathrm{sub}$
and $P.\mathrm{supp}$. This verifies all ten fields on the actual stage;
no auxiliary zero-form model is used.

The actual frame equality
$cfOf(x)=\operatorname{Fr.at}(w.\mathrm{st}(x.o.Z),f)$
also follows from $w.C_{\mathrm{eq}}$.
The conductor function $QOf(P)$ equals $T.Q$ at every place:
on $P.W_t$ this is its definition; off it, $T.\psi=1$, so the exact
\texttt{IsPsi.ram} law forces $T.Q=\varnothing$.
Thus any $C$ on this same stage with $C.Q=QOf(P)$ satisfies the
support condition of NC16, which supplies a tuple tie with exactly
$tp.\xi=w.I.\xi$ and $tp.\psi=T.\psi$.
The literal \texttt{ColW}, \texttt{Pre2}, \texttt{RowPk} and
\texttt{St} records remain simultaneous inputs here; this transport is
not an existence proof for all their geometric and arithmetic fields.

\paragraph{RC14: a rank-one conjugation detects the parahoric scalar-action subgroup.}
At a split real-field prime, put $K=F_v$ and let $\mathcal O$ be its
valuation ring. The unitary group there is $\operatorname{GL}_n(K)$:
under $E\otimes_F K=K\times K$, a unitary matrix is the pair
$(h,h^{-T})$. Thus every $h$ has a unique adelic lift supported at $v$.
It is integral exactly when $h\in\operatorname{GL}_n(\mathcal O)$.
This lift is an actual finite adele: its two components at $v$ are
prescribed, it is the identity elsewhere, and expansion in a fixed
$F$-basis of $E$ gives finite-adele coefficients and inverses.

Let $P_r\le\operatorname{GL}_n(\mathcal O)$ be the inverse image of
the upper block parabolic of block sizes $r,n-r$.
Suppose a nonzero function $F$ on $G_f$ has the prescribed scalar
action on $P_r$, and its scalar character has the same kernel as
$u\mapsto\xi(\det u_{22})$. Suppose $\xi(a)\ne1$ for some
$a\in\mathcal O^\times$. Every integral element supported at $v$ that
acts by a scalar on $F$ belongs to $P_r$.

Indeed, write its local matrix as $h$. If $h\notin P_r$, there are
$i\ge r$, $j<r$ with $h_{ij}\in\mathcal O^\times$.
Set $t=(a-1)/h_{ij}\in\mathcal O$, and define
\[
 A=I+t(he_j)e_i^T,\qquad
 B=I+t e_j(e_i^Th)=h^{-1}Ah.                            \tag{RC.15}
\]
The identity $\det(I+uv^T)=1+v^Tu$ gives
$\det A=\det B=a$. Both matrices and their inverses are integral:
their entries are integral, their determinant is a unit, and the
adjugate formula gives integral inverses.
Matrix $A$ changes only column $i\ge r$; matrix $B$ changes only row
$j<r$. Both therefore lie in $P_r$. Their lower-right block
determinants are respectively $a$ and $1$.
Their adelic lifts $u,w$ satisfy $w=\kappa^{-1}u\kappa$.
The scalar for $u$ on $F$ is nontrivial and that for $w$ is $1$.
The scalar for $\kappa$ is nonzero, since right translation is invertible
and $F\ne0$. Applying the two equal products $\kappa w=u\kappa$
to the nonzero line spanned by $F$ forces those two scalars to agree,
a contradiction. Only equality of character kernels was used, so this
also applies to two complex readings of one coefficient character.

\paragraph{RC15: an automorphism cannot send one parahoric into both opposite parahorics.}
Let $0<r<n$. At a fixed split prime, the two possible embeddings
$E\hookrightarrow F_v$ are conjugate. One local matrix is consequently
the inverse transpose of the other. In the compact group
$K_v=\operatorname{GL}_n(\mathcal O)$, their parahorics are the inverse
images $P_+,P_-$ of opposite block parabolics in
$\operatorname{GL}_n(\kappa_v)$.
They have equal finite index, since inverse transpose interchanges
them. Their intersection has strictly larger index than either one:
$I+E_{0r}$ belongs to $P_+$ and not $P_-$.

Suppose $\theta$ is the source's \texttt{FrameAut}: it preserves the
global integral group and preserves support at each prime, in both
directions. Its restriction to the subgroup supported at $v$ is
therefore an automorphism of $K_v$. If
\[
                  \theta(P_\pm)\subseteq P_+\cap P_-,
                                                               \tag{RC.16}
\]
then preservation of finite subgroup index gives
$[K_v:P_\pm]\ge[K_v:P_+\cap P_-]>[K_v:P_\pm]$, impossible.
Consequently, if the image under $\theta$ of a parahoric for some
split embedding lies in the parahorics for two split embeddings at
the same prime, those two embeddings are equal.
There are only two possibilities because $E/F$ is quadratic:
write $E=F\oplus Fy$ with $\bar y=-y$ and $y\ne0$; the image of $y$
is determined by $y^2\in F$ up to sign. This proves the literal
split-place equality used in \texttt{LevCmp.levEq}, including its
orientation.

\paragraph{RC16: the common scalar action recovers the additive level character.}
Let the residue characteristic at the split prime be $p$.
Let $U_r\le K_v$ be the inverse image of the subgroup
\[
 \left\{\begin{pmatrix}u&B\\0&D\end{pmatrix}:
 u\text{ upper unitriangular of size }r,\
 D\in\operatorname{GL}_{n-r}(\kappa_v)\right\},
 \qquad 2\le r\le n.
\]
Its index in $K_v$ is finite and prime to $p$.
For a direct proof, it contains the inverse image of the full upper
unitriangular subgroup $U_n(\kappa_v)$. Reduction
$\operatorname{GL}_n(\mathcal O)\to\operatorname{GL}_n(\kappa_v)$
is surjective by entrywise lifting and the unit-determinant test.
The index of that smaller subgroup is
\[
 \frac{|\operatorname{GL}_n(\kappa_v)|}{|U_n(\kappa_v)|}
       =\prod_{a=1}^{n}(q_v^a-1),                       \tag{RC.17}
\]
since an ordered basis has
$\prod_{a=0}^{n-1}(q_v^n-q_v^a)$ choices and $U_n$ has
$q_v^{n(n-1)/2}$ elements. Every factor on the right is prime to $p$,
and the index of $U_r$ divides it.

Let nonzero functions $F,F'$ transform on $U_r$ by characters whose
additive factors are $\alpha,\alpha'$ and whose lower-block determinant
factors are arbitrary. Assume $\alpha,\alpha'$ are trivial on the
maximal ideal. Let $\theta$ preserve $K_v$ and suppose that, for every
$u\in U_r$, $\theta(u)$ acts by scalars $c,c'$ on $F,F'$, with
$c^m=(c')^m$, where $\gcd(m,p)=1$.
Then $\alpha(y)=\alpha'(y)$ for every $y\in\mathcal O$.

To prove this, let $H\le K_v$ consist of elements acting by a scalar
on both functions. A nonzero function has a unique nonzero scalar
for any such element; these scalars define homomorphisms
$\lambda,\lambda':H\to\mathbb C^\times$.
Set $\eta=(\lambda/\lambda')^m$.
The hypotheses give $\theta(U_r)\le\ker\eta\le H\le K_v$.
Thus $d=[H:\ker\eta]$ is finite and divides
$[K_v:\theta(U_r)]=[K_v:U_r]$, so $\gcd(d,p)=1$.
For the transvection $x_y=I+yE_{01}$, $y\in\mathcal O$,
the lower-block determinant is $1$ and the top superdiagonal sum is $y$.
It lies in both original level groups and in $H$, and hence
\[
       (\alpha(y)/\alpha'(y))^{md}=1.
                                                               \tag{RC.18}
\]
Also $p y$ belongs to the maximal ideal, so
$\alpha(y)^p=\alpha'(y)^p=1$. Bezout for $p$ and $md$ now gives
$\alpha(y)/\alpha'(y)=1$. The source uses $m=p-1$.
The conclusion concerns values on integral $y$, which are exactly the
superdiagonal values occurring in the level-character formula.

\paragraph{RC17: one explicit integral frame gives a trivial stabilizer class.}
Let $n\ge2$, $r\ge2$, and let a split prime $P\mid p$ satisfy
$p\ge2n$ and $p\notin P^2$. Then an adelic element $g$ can be chosen
such that
\[
 \gamma\in U_n(E),\quad
 g^{-1}\gamma g\in U_{\max},\quad
 (g^{-1}\gamma g)_P\in U_r
             \quad\Longrightarrow\quad\gamma=1.         \tag{RC.19}
\]
Use the integral matrix $A$ with first column $(1,\ldots,1)^T$,
second column $(0,1,\ldots,n-1)^T$, and remaining columns the corresponding
standard basis columns. It is lower triangular with diagonal entries
$1$, so $A\in\operatorname{GL}_n(\mathbb Z)$. Lift it at the chosen
split prime and use the identity at all other primes. This gives
$g\in U_{\max}$ by the explicit split lift of RC14.

Since $g$ is integral, the hypotheses make $\gamma$ an everywhere
integral rational unitary matrix. Each of its entries is an algebraic
integer of absolute value at most $1$ under every complex embedding:
the sum of squared absolute values in its row is $1$.
Such a nonzero entry is a root of unity. One elementary proof is to
consider its powers: their conjugates remain bounded, so the integral
coefficients of their monic characteristic polynomials lie in a fixed
finite set. Only finitely many roots of those finitely many polynomials
occur, forcing two powers to agree.
At least one row entry is nonzero; its absolute value is $1$, so every
other entry in that row is zero. Thus row $i$ has one nonzero entry
$\zeta_i$, in column $\sigma(i)$, where $\zeta_i$ is a root of unity.

Reduction of $\gamma A=AH$, with $H\in U_r$, on the first column gives
$\zeta_i\equiv1\pmod P$.
A torsion unit congruent to $1$ at this unramified odd prime is $1$.
Here is the exact elementary argument. If a nontrivial such unit has
finite order, a power $u=1+z$, $z\in P$, has prime order $a$.
The geometric sum $1+u+\cdots+u^{a-1}=0$ shows $a\in P$,
so $a=p$. Expanding it modulo $P^2$ gives
\[
       0=p+\frac{p(p-1)}2 z+z^2 t
\]
for some integral $t$. Since $p$ is odd and $p,z\in P$, this forces
$p\in P^2$, contrary to the hypothesis. Therefore all $\zeta_i=1$.

The second column now gives
$\sigma(i)-i\equiv H_{01}\pmod p$ for every $i$.
The difference between any two integers $\sigma(i)-i$ has absolute
value at most $2(n-1)<p$, so divisibility by $p$ forces exact equality.
Their common value is nonnegative from $i=0$ and nonpositive from
$i=n-1$, hence is zero. Thus $\sigma(i)=i$ and $\gamma=1$.

For the actual $D$, the source hypotheses supply these inequalities
and unramifiedness. Put $n=D.R.nV=r+t$.
The congruence $g_k\equiv1\pmod{Q_Sqd}$ and primality of $g_k$ imply
$g_k\ge Q_Sqd\ge q>4tr+1\ge2(r+t)$.
The specified count of primes above $g_k$ is $[E:\mathbb Q]$.
In the equality $\sum_{P\mid g_k}e_Pf_P=[E:\mathbb Q]$, every summand
is at least $1$, so all $e_P=f_P=1$. In particular $g_k\notin P^2$.
This uses the actual type-prime fields; it does not assert existence
of those primes without the host/type input.

\paragraph{RC18: the stabilizer class gives a nonzero form for every nonzero weight.}
Let $J$ have the source's product-level shape: integral parahoric
conditions at the ramified split primes, integral unipotent conditions
at the nonempty set of type primes, and the integral maximal compact
condition outside those sets. A product level permits the component
of any $\kappa\in J$ at one prime to be isolated inside $J$.
Each isolated component is integral by its stated local condition;
the complement condition treats all other primes. Consequently
$J\subseteq U_{\max}$, and every element of $J$ satisfies the specified
type-prime unipotent condition.

Choose the class $g$ of RC17 at one actual type prime. Its rational
stabilizer $\Gamma\cap gJg^{-1}$ is trivial.
For a nonzero $w$ in any coefficient representation $W$ and any
$J$-character $\chi$, define
\[
 F(\gamma g\kappa)=\chi(\kappa)\rho(\gamma)w
       \quad(\gamma\in\Gamma,\ \kappa\in J),
 \qquad F=0\text{ off }\Gamma gJ.                        \tag{RC.20}
\]
The expression is unique: if
$\gamma g\kappa=\gamma' g\kappa'$, then
$\gamma'^{-1}\gamma\in\Gamma\cap gJg^{-1}$, hence both factors agree.
Multiplication on the left by $\Gamma$ and on the right by $J$
preserves the double coset, so (RC.20) satisfies both transformation
laws on it and on its complement. Also $F(g)=w\ne0$.
For an irreducible weight the space is nonzero by the definition of
irreducibility. This proves the nonzero-form conclusion used for both
the complex \texttt{MtReal} model and the coefficient-field cell.
The congruence-kernel field in \texttt{MtReal} converts the former
function into its smooth form-space range, and the cell's actual
\texttt{real} isomorphism supplies its preimage in \texttt{MtV}.

\paragraph{RC19: the local arguments give the literal global LevEq.}
Retain the exact hypotheses of \texttt{LevCmp.levEq} on
$m:\texttt{MtReal}$ and the cell $c$ of $X:\texttt{NExpTied}$:
the actual type-prime data \texttt{GPlaces};
product, ramified-parahoric, type-unipotent and outside-compact
descriptions for $m$; the ramified and type descriptions for $c$;
$\mathrm{badV}\subseteq gPl\cup qPl$;
nontriviality of the common ramified component on some local unit;
and the exact \texttt{Joint} datum.
The cell's type-character description must hold for every complex
reading extending the specified side. Put $S=gPl\cup qPl$.
Then
\[
             \texttt{LevEq}(m.J,m.\chi,c.J,e\,c.\chi).
                                                               \tag{RC.21}
\]

We spell out the role of \texttt{Joint}. It gives an automorphism
$\theta$ preserving integrality and each prime support, a subgroup $J_1$,
another side embedding $e'$, and an injective $e'$-semilinear map
$\omega$ from the cell's weight space to a complex vector space.
The subgroup contains specified parahorics over $qPl$ and specified
unipotent groups at all type primes. Every element of $J_1$ acts,
after applying $\theta$, by scalars on all the realized $m$-forms and
on all cell forms after $\omega$; at a type prime their
$(g_k-1)$st powers agree. This is substantially more than equality
of ranks or an arbitrary linear isomorphism of two form spaces.

The cell is a product level directly from its \texttt{J\_\allowbreak mem} formula.
Indeed, isolate one adelic component using the idempotent of that
prime; every listed matrix condition becomes either the original
condition or the identity condition. Integrality is preserved.
Its outside-compact description follows from its \texttt{hyp} field,
the bad-set inclusion, and injectivity of $e$, which preserves the
character kernel. RC18 supplies nonzero forms on both sides.
Their preimages under the actual realization maps, and injectivity
of $\omega$, give two nonzero complex-valued weight functions to
which RC14--RC16 apply.

At a type prime, two local multiplicative components representing the
same Hecke character are equal. To see uniqueness without a continuity
shortcut, both components kill sufficiently deep principal units.
Approximate an arbitrary local unit by a nonzero element of $F$ so that
the quotient lies in both kernels. Their prescribed values on that
rational element agree, and therefore so do their values on the local
unit. This works for every element of $F_v^\times$, with the depth of
approximation adjusted to its valuation. RC16, with $p=g_k$ and
$m=g_k-1$, identifies the additive factors on integral elements.
The groups have the same unipotent matrix condition. The lower block
of an integral parahoric matrix has unit determinant: reduction is
block upper triangular and the full determinant is nonzero.
Its superdiagonal sum is integral. The character formulas consequently
agree on every supported group element.

At a ramified prime, \texttt{Joint} provides a parahoric at some split
embedding. RC14, applied to both nonzero functions, puts its image
under $\theta$ inside both proposed parahorics. The change from $e$
to $e'$ has no effect on the coefficient character's kernel.
RC15 forces the two split embeddings to be identical, and uniqueness
of the common Hecke component gives equality of their determinant
characters. Thus the two levels and characters agree on each
single-prime subgroup over $S$.

Finally, $S$ is finite: the type set lies over two rational primes,
and a Hecke character is ramified at only finitely many finite primes.
For the latter assertion, continuity and a complex neighborhood
containing no nontrivial subgroup force the character to kill the
unit factors outside a finite set.
For $\kappa\in m.J$, isolate its components at the finitely many primes
of $S$, successively divide them off, and leave a factor supported
outside $S$. The isolated factors belong to $c.J$ with equal character
values by the local comparison; the remaining factor belongs by the
common outside-compact rule and has character $1$ on both sides.
Multiplying proves membership and character equality globally.
Reverse the roles of the two product levels to obtain equality of
subgroups. This proves (RC.21), and RC2 can then be applied with the
separate exact \texttt{CellWt} input.

\paragraph{RC20: an actual coherent index and side selector on the twisted places.}
Assume $L_0/\mathbb Q$ is Galois, the literal
\texttt{TwCoverPair} holds, and an index $i_0\in D.\mathrm{ntw}.Tw$
is supplied. Write $v\sim v'$ if the two valuations have the same
positive-valuation predicate on $\mathcal O_{L_0}$. Write $v\approx v'$
if $v\sim v'$ or $\bar v\sim v'$.
Conjugation preserves $\sim$: normality of $L_0$ lets complex conjugation
act on $L_0$, and after clearing an integer denominator it acts on
its ring of integers as well. Thus $\approx$ is an equivalence relation,
and $v\approx\bar v$.
The coverage hypothesis says precisely that the prime of
$\mathcal O_{L_0}$ below every \emph{twisted} coefficient place is
the center of one of the two side embeddings of some actual twist index.
Consequently every twisted $v$ admits some \texttt{ReadsAt}$(i,b,v)$.

For each $\approx$-class choose an index satisfying
\[
 \exists u,b,\quad v\approx u\quad\hbox{and}\quad
                  \texttt{ReadsAt}(i,b,u),                    \tag{RC.22}
\]
if one exists, using $i_0$ otherwise. Equivalence makes the set of
eligible indices exactly the same throughout a class, so a fixed
choice function gives $s(v')=s(v)$ whenever $v\approx v'$.
At a twisted $v$, (RC.22) is nonempty. A reading at $u$ transports
along $\sim$ without changing its side; if conjugation is needed,
conjugating the reading reverses the side. Hence choose a preliminary
$b(v)$ with $\texttt{ReadsAt}(s(v),b(v),v)$ at every twisted $v$.
Outside that set its value may be arbitrary.

Choose one representative $o(v)\in\{v,\bar v\}$ of each conjugation
orbit. Define the final side by $b(v)$ at the representative and
by the opposite of the representative's side at the other place.
The reading conjugation identity and $s(\bar v)=s(v)$ show that the
final side still reads the selected index. At a two-element orbit
the sides are opposite; at a fixed place no impossible sign reversal
is required. This is exactly \texttt{exists\_\allowbreak settings}.
Its reading conclusion holds on the twisted places only.
It does not give the all-place hypothesis \texttt{hs} of RC11,
nor does it supply a reading for an empty index type.

\paragraph{RC21: finite integral bases and prime avoidance give the actual ExistsData.}
Fix the exact expansion datum $\mathcal E$, the function $N_0$ on
twist indices, and the following five source hypotheses:
$\texttt{aux\_\allowbreak 3}(D,\mathrm{bk0},\mathcal E,N_0)$,
\texttt{TwistedPairSettingHyps}, \texttt{FullRank},
\texttt{adm7}, and \texttt{TwCoverPair}, together with
$L_0/\mathbb Q$ Galois. The first is the actual twisting theorem
at every nonzero index: given its setting, any positive error budget,
the full-rank theta datum, any finite integral theta basis, and any
finite forbidden set, it supplies a packet satisfying
\texttt{GoodTw}. We use that hypothesis with its full meaning,
including the lattice-index inequalities, the finite-extension
quantifiers in \texttt{IdxAt}, and the three arithmetic parts
\texttt{partAT}, \texttt{partC}, \texttt{partD}.
We do not prove these inequalities by choosing an arbitrary character.
Under these hypotheses the literal conclusion is
\[
 \exists S\ \ \forall A:Tw\to\operatorname{Finset}(Pl(E))\
 \exists(t,sd)\quad \texttt{GoodData}(S,t,sd),             \tag{RC.23}
\]
with $Q(t(v))\cap A(s(v))=\varnothing$, with $t$ constant on
$\approx$-classes of twisted places, and, if the twisted-place
set is finite, with the rational-prime supports of packets from
different $\approx$-classes pairwise disjoint.

Here is the integral basis construction, which occurs \emph{before}
the forbidden sets $A$ are specified. Fix an active $i$ and
$1\le j\le D.R.m$, and write $O=O_i$.
At the trivial twist $k=0,Q=\varnothing,\psi=1$, all seven
\texttt{TwOK} clauses hold. The hypothesis \texttt{adm7}
(literally \texttt{admT4}) places the scalar extensions
$1\otimes\theta(a)$, for $a\in\mathcal E.\mathrm{Dat0}(i,j,b)$,
inside the image of the given finitely generated module
$\mathrm{MtLat}$ under $\mathrm{thP}$.
The finite injective local extension $O\to O_{i,0}$ is flat, since it
is torsion-free over the DVR $O$, and faithfully flat, since it is local.
Thus $x\mapsto1\otimes x$ is injective. The containing image is a
finitely generated $O$-module by finiteness of $O_{i,0}/O$.
Noetherianity of $O$ and injectivity show that the original theta span
is finitely generated.

Its image under the actual integral averaging operator
$AO=|G_0|e_O$ is also finitely generated and torsion-free.
It has a basis chosen from the vectors $AO\,\theta(a)$ themselves.
Indeed, first express finitely many module generators as finite
combinations of this generating family; their combined supports give
a finite generating subset. Choose one with minimal cardinality.
If it has a nontrivial relation, choose a nonzero coefficient of
least DVR valuation. It divides every other coefficient. Divide the
relation by it using torsion-freeness. The chosen coefficient is now
$1$, so its vector is a combination of the others, contrary to
minimality. Hence the chosen subset is a basis. Together with the
retained \texttt{FullRank} identity this gives all fields of
\texttt{TwIn.Ok}. Use empty subsets at inactive indices or out-of-range
$j$. Choice over $(i,j,b)$ gives one family $S$ independent of $A$.

We next prove the required finite avoidance step. For a finite set
$B$ of primes of $E$, let $B^+$ be the union, over $P'\in B$,
of all primes containing the nonzero integer $\mathrm N(P')$.
This set is finite, by factorization of nonzero ideals in
$\mathcal O_E$. If $P'$ lies over the rational prime $p$, its norm is
$p^f$ for some $f>0$; every prime $P$ over $p$ therefore contains
$\mathrm N(P')$. It follows that $Q\cap B^+=\varnothing$ implies
that no prime in $Q$ lies over the same rational prime as a prime in
$B$. This is the exact symmetric relation \texttt{Apart}.

For each active $i$ and each finite $B$, apply \texttt{aux\_\allowbreak 3}
with the fixed basis $S_i$, the actual positive budget
$\mathrm{twEps}(D,i)$, its supplied setting, and forbidden set
\[
             \mathrm{AvRec}(D)\ \cup\ A(i)\ \cup\ B^+ .
                                                               \tag{RC.24}
\]
The resulting $(Q,k,\psi)$ has the same \texttt{GoodTw} properties
after forgetting the extra forbidden sets, is disjoint from $A(i)$,
and is \texttt{Apart} from $B$.
At an inactive index choose $(\varnothing,0,1)$: the required
goodness assertion is conditional on activity and hence vacuous.

An elementary induction now selects good packets at all indices
which are pairwise \texttt{Apart} on any prescribed finite subset
$T\subseteq Tw$. Start with a good packet at every index.
When adding $a$ to a previously treated finite set $T$, keep all
other packets and choose the packet at $a$ with $B$ equal to the
union of their conductor sets. The preceding construction and
symmetry of \texttt{Apart} give every new pairwise condition.
All old conditions are unchanged.

If there are no twisted places, (RC.23) is immediate, with empty
packet function and arbitrary side. Otherwise a single covered
twisted place supplies the index $i_0$ needed in RC20.
Use its selector $s,sd$. If the twisted-place set is finite, take
$T$ to be its finite image under $s$; otherwise take $T=\varnothing$,
as no pairwise-apart conclusion is then required.
Select packets $f(i)$ by the preceding induction and set
$t(v)=(s(v),f(s(v)))$.
A reading at a twisted place makes its index active, so every
required \texttt{GoodPack} clause holds there.
RC20 gives reading, side reversal, conjugation invariance, and
constancy of $t$ on $\approx$-classes.
Finally, if two twisted places have the same selected index,
their readings imply $v\approx v'$: equal sides have the same
center on $\mathcal O_{L_0}$, and opposite sides differ by conjugation.
Consequently distinct $\approx$-classes use distinct indices, and
the finite selection supplies the promised \texttt{Apart} property.
This proves every clause and the quantifier order of
\texttt{ExistsData}.

\medskip
\noindent\textbf{Exact boundary.}
The actual rank, character descent, finite reading choices, local
level comparison and tuple compatibility have now been proved
under the displayed inputs. The simultaneous construction still
requires the supplied global Artin reciprocity map; the actual
ring-class universal property and twist coverage; the exact
geometric and representation-theoretic realization data producing
\texttt{Joint}, local shape, and \texttt{CellWt}; and the genuine
twisting/full-rank/integral-inclusion hypotheses used in RC21.
For nonempty local data, RC10 gives a sufficient exact
Frobenius-character criterion, while the local index and side
identifications remain necessary.
For tuples, the actual common \texttt{ColW}, \texttt{Pre2},
\texttt{RowPk} and \texttt{St} objects are retained; resetting tuple
maps in an unrelated row model does not supply them.
The bare stage with $Tw=Loc=\varnothing$ supplies none of the
all-place coverage required for RC11.
The final theorem \texttt{choices\_\allowbreak exist\_\allowbreak for} still depends on these
simultaneous suppliers and their stated analytic hypotheses.


\section{The actual local period rings and the bare Galois stage}
\label{sec:local-stages}

\status{This section reconstructs the local-ring and invariant calculations
needed by \texttt{LocStage}, and the actual empty-local-index construction
of \texttt{NStage}. Its period rings are the specified Witt, divided-power,
and adic constructions. The final stage is attached to the same base,
columns, and adapted row record. Its two index types are empty; it supplies
bare stage existence, without an arithmetic twist reading, a nonzero
projector, or a comparison theorem for an arbitrary geometric Galois
representation. These are human mathematical proofs from the pinned
source, without a Lean build or a kernel-verification claim.}

\paragraph{Conventions.}
Throughout LS1--LS17, $p$ is any rational prime, including $2$,
$\Omega=\overline{\mathbb Q}_p$, $C=\widehat\Omega$, $O=\mathcal O_C$,
and $G=\operatorname{Gal}(\Omega/\mathbb Q_p)$. We normalize the
absolute value by $|p|=p^{-1}$. For a finite extension
$K\subset\Omega$, put $G_K=\operatorname{Gal}(\Omega/K)$ and
$K_0=K\cap\mathbb Q_p^{\mathrm{ur}}$, where
$\mathbb Q_p^{\mathrm{ur}}$ is the field generated by roots of unity of
order prime to $p$. Completion, finite-extension valuation theory,
finite-field Galois theory, and the defining Witt identities are used in
their ordinary mathematical senses. NC1--NC2 already provide the
complete-DVR and normalized-valuation facts. All further period-ring
properties used below are proved here.

\paragraph{LS1: an explicit Ax approximation and completed fixed fields.}
If $H\le G$, $a\in\Omega$, and $|ha-a|\le r$ for every $h\in H$,
then there is $b\in\Omega^H$ such that
\[
                     |a-b|\le16r.                 \tag{LS.1}
\]
Here is the finite-orbit proof, including the uniform constant. If the
orbit has at most $2^k$ elements, let
\[
 C_k=2^{\,4-(2k+4)/2^k},\qquad
 d_k=2^{(k+1)/2^k};\quad C_0=1,\quad C_{k+1}=C_kd_k\le16.
 \tag{LS.2}
\]
Induct on $k$. The one-point case is immediate. Otherwise let $T$ be
the orbit, $n=|T|$, with $2^k<n\le2^{k+1}$, and form
$f(X)=\prod_{z\in T}(X-z)$. Its Hasse derivative
$P=\partial^{[n-2^k]}f$ is $H$-fixed, has degree $2^k$, and leading
coefficient $\binom n{n-2^k}$. Expansion at $a$ gives
\[
 |P(a)|\le r^{2^k},\qquad
 \left|\binom n{n-2^k}\right|\ge n^{-1}.          \tag{LS.3}
\]
The second estimate follows from Legendre's factorial formula: the
$p$-valuation of a binomial coefficient is a sum of zeros and ones,
with no terms above $\lfloor\log_p n\rfloor$.
Factor $P$ over $\Omega$. Some root $z$ satisfies
$|z-a|^{2^k}\le n r^{2^k}$, hence $|z-a|\le d_kr$.
Every $H$-conjugate of $z$ is a root of $P$ and
$|hz-z|\le d_kr$ by the ultrametric inequality. Induction applied to
$z$ gives an $H$-fixed $b$ with
$|z-b|\le C_kd_kr$. The same inequality bounds $|a-b|$.
Every algebraic element has a finite orbit, proving (LS.1).

Consequently
\[
                         C^H=\overline{\Omega^H}.   \tag{LS.4}
\]
The inclusion from right to left follows from isometry of the action.
For the other inclusion, approximate $x\in C^H$ by $a\in\Omega$
within $\delta$; then $|ha-a|\le\delta$ for every $h\in H$.
Equation (LS.1) gives a fixed algebraic point within $16\delta$ of
$x$. In particular $C^{G_K}=K$, since a finite extension of
$\mathbb Q_p$ is complete and therefore closed in $C$.

\paragraph{LS2: a uniform cyclotomic projection bound.}
Choose compatible primitive roots $\zeta_n$ of order $p^n$, with
$\zeta_0=1$, $\zeta_{n+1}^p=\zeta_n$, and put
$K_n=\mathbb Q_p(\zeta_n)$. For $n\ge1$, $m\ge0$, $N=p^m$, and
$z=\zeta_{n+m}$, every $y\in K_{n+m}$ has a unique expansion
\[
 y=\sum_{i=0}^{N-1}a_i z^i,\quad a_i\in K_n,
 \qquad |p|\,|a_0|\le |y|.                         \tag{LS.5}
\]
We give the bound without a trace-integrality assumption. For a
primitive $p^r$-th root $z$, all primitive conjugates have the same
distance from $1$: the geometric-sum identity gives the inequality
in each direction. Evaluating the cyclotomic product at $1$ yields
\[
              |z-1|^{e_r}=p^{-1},\qquad
              e_r=p^{r-1}(p-1).                    \tag{LS.6}
\]
For $Q\in\mathbb Q_p[X]$ of degree $<e_r$, the nonzero terms of
$Q(z-1)$ have distinct valuations modulo $\mathbb Z$. Thus its norm
is the largest term norm, and
$p^{-1}|Q_j|\le|Q(z-1)|$ for every coefficient. Taylor substitution
$X\mapsto X+1$ and its inverse have integral binomial coefficients;
therefore the same estimate holds for the coefficients of a
polynomial evaluated at $z$. This also proves that
$1,z,\ldots,z^{e_r-1}$ are independent, so the cyclotomic degree is
exactly $e_r$.

Write each $a_i=P_i(\zeta_n)$, with $\deg P_i<e_n$, and set
$P(X)=\sum_{i<N}P_i(X^N)X^i$. Its degree is $<Ne_n=e_{n+m}$.
Its coefficients at powers divisible by $N$ are exactly those of
$P_0$. Applying the preceding estimate and then the ultrametric
inequality to $P_0(\zeta_n)$ proves (LS.5).
The extension $K_{n+m}/K_n$ has degree $N$, and averaging its roots
shows that $a_0=N^{-1}\operatorname{Tr}_{K_{n+m}/K_n}(y)$.
These projections are compatible as $m$ varies, commute with $G$,
and have norm at most $p$, independently of $m$ and $n$.
They extend to maps
\[
 R_n:\widehat{\bigcup_{r\ge1}K_r}\longrightarrow K_n,
 \qquad R_nx\longrightarrow x.                    \tag{LS.7}
\]
Indeed they are eventually the identity on the dense union, and the
uniform bound extends this convergence to its completion. This proof
also covers $p=2$, $n=1$, where $K_1=\mathbb Q_2$.

\paragraph{LS3: every nonzero Tate twist has zero invariants.}
Let $\chi:G\to\mathbb Z_p^\times$ be the character determined by
$g\zeta_n=\zeta_n^{\chi(g)\bmod p^n}$. For any finite $K$ and any
nonzero integer $i$,
\[
 \bigl\{x\in C:g(x)\chi(g)^i=x\quad(g\in G_K)\bigr\}=0.
 \tag{LS.8}
\]
First take $K=\mathbb Q_p$. The kernel of $\chi$ fixes precisely
the cyclotomic union in $\Omega$, so LS1 puts such an $x$ in the
completion in (LS.7). Equivariance gives
$g(R_nx)\chi(g)^i=R_nx$. For each $n$ there is
$g\in G_{K_n}$ with $\chi(g)^i\ne1$. Otherwise the image of this
group under $\chi$ would consist of roots of $X^{|i|}-1$, hence be
finite, contradicting the unbounded degrees
$[K_{n+m}:K_n]=p^m$ proved in LS2. This $g$ fixes $R_nx\in K_n$,
so $R_nx=0$ for every $n$. Taking the limit proves the assertion.

For general $K$, let $d=[G:G_K]$ and set
$f(g)=g(x)\chi(g)^i$. The hypothesis makes $f$ constant on cosets
$gG_K$. Thus $Y=\prod_{gG_K}f(g)$ satisfies
$h(Y)\chi(h)^{id}=Y$ for every $h\in G$, by permutation of the
cosets. The already proved case gives $Y=0$. A factor is zero, and
all character factors are units, so $x=0$. No estimate involving
$[K:\mathbb Q_p]^{-1}$ is needed.

\paragraph{LS4: the actual Fontaine map and its principal kernel.}
Set
\[
 R=\varprojlim_{x\mapsto x^p}O/pO,\quad A=W(R),\quad
 \theta\left(\sum_{j\ge0}p^j[x_j]\right)
       =\sum_{j\ge0}p^j x_j^\sharp .                \tag{LS.9}
\]
The sharp map is multiplicative; its reduction modulo $p$ is the
zeroth inverse-limit coordinate. It is not asserted to be additive.
The Witt identities make (LS.9) a ring homomorphism. Every element $y\in O$ has a $p$-th root in $O$:
approximate $y$ by algebraic elements $y_n$ and successively
choose roots $x_n^p=y_n$ with
$|x_{n+1}-x_n|\le|y_{n+1}-y_n|^{1/p}$, using the factorization
of $X^p-y_{n+1}$ at $x_n$. The ultrametric inequality makes
$x_n$ Cauchy when $y_n$ is Cauchy, and its limit is the required
root. Iterating gives compatible $p$-power roots of $y$, hence
a preimage under sharp. Thus $\theta$ is onto. The ring $R$ is a perfect valuation domain:
its norm is $|x|=|x^\sharp|$, and if $|y|\le|x|$, the ratios of
the exact sharp-root coordinates of $y$ and $x$ give an element of
$R$ whose product with $x$ is $y$. A zero sharp value forces every
coordinate to be zero.

Choose $\varpi\in R$ with $\varpi^\sharp=p$. Then
\[
                 \ker\theta=(p-[\varpi]).          \tag{LS.10}
\]
To prove this, write $a=[a_0]+pb$. If $\theta a=0$, then
$|a_0|\le|p|$, so $a_0=\varpi c$ in the valuation ring $R$.
Hence
$a=(p-[\varpi])(-[c])+p([c]+b)$, and applying $\theta$ shows
$[c]+b$ is again in the kernel. Iterating and summing the
$p$-adically convergent coefficients proves divisibility. The
reverse inclusion is immediate. The nonzero zeroth Witt coefficient
of $p-[\varpi]$ also makes this element a nonzerodivisor.

For the compatible cyclotomic system of LS2 let
$\varepsilon\in R$ have sharp-root coordinates $\zeta_n$, and put
\[
 \mu=[\varepsilon]-1,\quad
 \nu=[\varepsilon^{1/p}]-1,\quad
 \xi=\sum_{j=0}^{p-1}[\varepsilon^{1/p}]^j,
 \qquad \mu=\xi\nu.                                \tag{LS.11}
\]
Then $\theta\xi=0$, and $\xi$ is a unit multiple of the generator
in (LS.10). To check the unit, reduce a factorization
$\xi=(p-[\varpi])u$ modulo $p$. The zeroth coefficient of $\xi$
is $(\varepsilon^{1/p}-1)^{p-1}$ and has norm $|p|$, so $|u_0|=1$;
a Witt vector with unit zeroth coefficient is a unit. For odd $p$,
the norm assertion follows from (LS.6), since the error between
sharp of a difference and the difference of sharp values is in
$pO$, of strictly smaller norm. For $p=2$ take one more inverse
Frobenius root: $\zeta_2^2=-1$, so
$|(\zeta_2-1)^2|=|2|$, whereas $|\zeta_2-1|>|2|$.
The same reduction argument followed by squaring gives
$|\varepsilon^{1/2}-1|=|2|$. This checks the dyadic case explicitly.

The quotient $A/(\xi^n)$ is $p$-torsion-free and $p$-adically
complete for every $n$. Indeed $\xi\bmod p\ne0$ in the domain
$R$, so $p^k\mid\xi^n b$ implies $p^k\mid b$. This makes
$(\xi^n)$ closed in the complete ring $A$: a Cauchy sequence of
multiples of $\xi^n$ has Cauchy quotients. Compatible quotient
classes modulo $p^k$ admit successively adjusted lifts with
successive differences in $p^kA$, proving completeness. These
arguments also show $p$-saturation of every $(\xi^n)$.

\paragraph{LS5: the actual de Rham DVR, filtration, and algebraic section.}
Put $A_P=A[1/p]$ and
\[
 B_{\mathrm{dR}}^+=\varprojlim_n A_P/(\xi^n),\qquad
 B_{\mathrm{dR}}=B_{\mathrm{dR}}^+[1/\xi],\qquad
 \operatorname{Fil}^aB_{\mathrm{dR}}=\xi^aB_{\mathrm{dR}}^+.
 \tag{LS.12}
\]
These agree with the source's completion and localization, since its
localized theta kernel is $(\xi)$, and every other generator differs
by a unit. The map $A_P\to B_{\mathrm{dR}}^+$ is injective.
In fact $\xi^n\mid a$ for every $n$ forces $a_0=0$ modulo $p$,
because $|\xi_0|<1$; after division by $p$, saturation permits the
same argument repeatedly, and $p$-adic separation gives $a=0$.
Clearing a power of $p$ proves the assertion on $A_P$.

The completion is separated and complete for $(\xi)$, with residue
field $C$, and multiplication by $\xi$ is injective. For the last
claim, lift an element modulo $\xi^{n+1}$ and cancel $\xi$ in
$A_P$; the resulting class modulo $\xi^n$ vanishes. An element of
the completion with nonzero residue is a unit, by the geometric
series lifting of its inverse. Every nonzero element is therefore
$\xi^r$ times a unit, for a unique $r\ge0$. This proves that
$B_{\mathrm{dR}}^+$ is a complete DVR and $B_{\mathrm{dR}}$ is its
fraction field. The filtration in (LS.12) is decreasing,
multiplicative, exhaustive and separated. The Galois action on
$C$, its valuation ring, the inverse-limit coordinates, the Witt
vectors, and then the completion preserves $\ker\theta$ and all
its powers. Thus it preserves every filtration step and commutes
with the residue map $\theta_D:B_{\mathrm{dR}}^+\to C$.

There is a unique $\mathbb Q_p$-algebra map
\[
 o:\Omega\longrightarrow B_{\mathrm{dR}}^+,
 \qquad\theta_D(o(a))=a,                            \tag{LS.13}
\]
where $\mathbb Q_p$ enters by the actual map $\mathbb Z_p\to A$.
For each algebraic $a$, lift the simple residue root $a$ of its
minimal polynomial by Newton iteration in the complete DVR.
The derivative has nonzero residue and is a unit, so the lift is
unique. On each finite separable subextension this is a field
embedding; uniqueness makes the embeddings compatible, additive,
multiplicative, and $G$-equivariant. In particular it agrees with
the prescribed constants from $\mathbb Z_p$. This construction does
not identify the residue map with a coefficient section on all of
$C$; the section asserted here has the algebraic domain $\Omega$.

\paragraph{LS6: the divided-power completion embeds in the actual de Rham ring.}
Inside $A_P$, set
\[
 D_0=A[\gamma_n:n\ge0],\qquad \gamma_n=\xi^n/n!,
 \qquad D=\widehat{D_0}^{\,(p)}.                    \tag{LS.14}
\]
Since $\ker\theta=(\xi)$, this is the source's divided-power
envelope $A_{\mathrm{cris},0}$, and $D=A_{\mathrm{cris}}$.
The multiplication identity
$\gamma_i\gamma_j=\binom{i+j}{i}\gamma_{i+j}$ shows that $D_0$
is the $A$-span of the displayed generators. It is $p$-torsion-free.
Its $p$-adic completion remains $p$-torsion-free, has dense image
of $D_0$, and satisfies
\[
 p^r\mid s\text{ in }D\quad\Longleftrightarrow\quad
 p^r\mid s\text{ in }D_0\qquad(s\in D_0).          \tag{LS.15}
\]
These assertions follow by projecting the inverse limit to
$D_0/p^rD_0$; to cancel $p$, lift one level farther and cancel in
the torsion-free ring $D_0$.

Modulo $\xi^nA_P$, the image $M_n$ of $D_0$ contains
$A/(\xi^n)$ and satisfies
\[
 p^n M_n\subset A/(\xi^n)\subset M_n.              \tag{LS.16}
\]
Indeed the terms $\gamma_j$ with $j\ge n$ disappear, while
$p^n/j!\in\mathbb Z_p$ for $j<n$. By LS4 the smaller module is
$p$-adically complete; the sandwich makes $M_n$ complete and
separated for the same topology. More explicitly, a Cauchy sequence
in $M_n$, after removal of a fixed initial term, has a tail in the
smaller module, and the two induced topologies differ by at most
$n$ powers of $p$. Extend $D_0\to M_n$ to its completion.
Uniqueness makes these maps compatible in $n$, giving
\[
 j:D\longrightarrow B_{\mathrm{dR}}^+,
 \quad j|_{D_0}\text{ is the actual inclusion},\quad
 p^n j(x)\bmod\xi^n\in\operatorname{im}(A).        \tag{LS.17}
\]
This is the literal boundedness field \texttt{IsCrisP.bdd}, with
$N=n$. Any map with the stated extension and boundedness is unique:
modulo $\xi^n$ its values lie in a fixed $p$-power enlargement of
$A/(\xi^n)$; dense approximations and separation there force equality.

We prove injectivity, rather than inferring it from density. Let
$L_N=\sum_{j<N}A\gamma_j$. The elementary relation
$\xi\gamma_j=(j+1)\gamma_{j+1}$ gives
\[
 L_N\cap\xi^N A_P\subset A\,(N\gamma_N).          \tag{LS.18}
\]
For induction, write $s=a\gamma_N+y$ with $y\in L_N$ and assume
$\xi^{N+1}\mid s$. The preceding induction step writes
$y=cN\gamma_N$. Multiplication by $N!$ and cancellation of
$\xi^N$ show that $\xi\mid a+Nc$ in $A_P$, hence in $A$ by
$p$-saturation. The displayed relation proves the next step.

Suppose $j(x)=0$, and write $x=s+pz$, $s\in D_0$.
For every $n$, (LS.17) and a sufficiently accurate dense approximation
give another expression $x=u+pz'$ with
$u\in D_0\cap\xi^nA_P$. Here is the adjustment explicitly:
write $x=s'+p^{n+1}v$ and choose $a\in A$ with
$p^n j(v)\equiv a\pmod{\xi^n}$; then $u=s'+pa$ has zero image
modulo $\xi^n$, and $x-u$ is divisible by $p$ in $D$.
Equation (LS.15) gives $s=u+pw$, $w\in D_0$.
Choose $M$ with $s\in L_M$, take $n=pM$, and split $w$ into
its terms of index $<n$ and its remaining terms. Thus
$s-pw_{<n}\in L_n\cap\xi^nA_P$. By (LS.18) this difference
is $a n\gamma_n$, divisible by $p$ in $D_0$. Hence $p\mid s$,
so $x=py$. As $p$ is a unit in $B_{\mathrm{dR}}^+$, also
$j(y)=0$. Iteration puts $x$ in every $p^rD$, whence $x=0$.

The actual Galois action preserves $D_0$ and extends to $D$.
The Witt Frobenius $\varphi$ does so as well: if
$\theta a=0$, then $a^p\in pD_0$ and
$\varphi(a)=a^p+pb\in pD_0$, so
$\varphi(a)^n/n!\in D_0$. The resulting actions commute, and
$j$ is Galois equivariant by its bounded uniqueness. No continuity
from the $p$-adic topology on $D$ to the $\xi$-adic topology on the
whole de Rham ring is presumed.

\paragraph{LS7: the logarithmic period, with integral convergence at $2$.}
For any $e\in R$ with $e^\sharp=1$, let $m=[e]-1$.
There is $w\in D_0$ such that
\[
                    pw=m^{p-1},\qquad w^p\in pD_0. \tag{LS.19}
\]
First suppose the first inverse-Frobenius sharp value is a nontrivial
$p$-th root. Put $X=[e^{1/p}]$ and $\eta=\sum_{i<p}X^i$.
The characteristic-$p$ polynomial identity gives
$m^{p-1}=\eta^p+pz$ with $z\in A$; applying $\theta$ shows
$\theta z=0$. Thus $w=\eta^p/p+z\in D_0$.
The divided powers of $\eta$ and $z$ show $w^p\in pD_0$:
the term $\eta^{p^2}/p^p$ is integral and divisible by $p$, since
$v_p((p^2)!)=p+1$, and the other terms follow from the binomial
coefficients and $z^p\in pD_0$. If the first nontrivial sharp root
occurs later, descend to it and use
$[e]-1=\xi(e)([e^{1/p}]-1)$ repeatedly. If there is no such root,
all coordinates of $e$ equal $1$, and (LS.19) holds with $w=0$.

Define $c_k=m^k/(k+1)$ in $A_P$. These elements lie in $D_0$.
For $q=\lfloor k/(p-1)\rfloor$, $r=k-q(p-1)$ and
$v=v_p(k+1)$, we have $v\le q$ and
\[
 c_k=(\text{a }p\text{-adic unit})\,p^{q-v}m^r w^q,
 \qquad p^{\lfloor q/p\rfloor}\mid c_k.            \tag{LS.20}
\]
The inequality $v(p-1)+1\le p^v\le k+1$ proves $v\le q$;
the divisibility follows from (LS.19). For $k>0$, also
$c_k\in\sqrt{pD_0}$. Therefore
\[
 u_e=\sum_{k\ge0}(-1)^k c_k\in D^\times,
 \qquad t_e=m u_e=\log[e]\in D.                   \tag{LS.21}
\]
Convergence follows from (LS.20). Modulo $p$, the difference
$u_e-1$ is a finite sum of nilpotents, hence nilpotent; lifting its
inverse through the $p$-adic completion proves that $u_e$ is a unit.
This argument uses no odd-prime exception. Multiplying a partial
sum by its factorial denominator gives exactly the finite sums in
the source's \texttt{IsLogA}, so (LS.21) has its stated meaning.

In $B_{\mathrm{dR}}^+$, $j(t_e)$ is the formal logarithm of
$1+j(m)$: modulo $\xi^n$, its terms of degree at least $n$
vanish, and (LS.17) transports the integral limiting relation.
Formal power-series identities and injectivity of $j$ give
$t_{e^a}=a t_e$ for integers $a$, and $\varphi(t_e)=pt_e$.
For the fixed primitive system $\varepsilon$, put $t=t_\varepsilon$.
It is nonzero and
\[
       \varphi(t)=pt,\qquad g(t)=\chi(g)t,
       j(t)\in\operatorname{Fil}^1\setminus\operatorname{Fil}^2.
 \tag{LS.22}
\]
For the Galois formula, choose integers $a_n\equiv\chi(g)\pmod{p^n}$.
The element $g\varepsilon\,\varepsilon^{-a_n}$ has its first
$n$ sharp-root coordinates equal to $1$; its $p^n$-th root $e_n$
therefore has sharp value $1$. Additivity of formal logarithms gives
$g(t)-a_nt=p^nt_{e_n}\in p^nD$. Completeness and separation give
the asserted limit. Finally $\mu=\xi\nu$ and
$\theta_D\nu=\zeta_1-1\ne0$; together with $u_\varepsilon$ a unit,
this proves the filtration assertion.

\paragraph{LS8: the actual crystalline core and unramified constants.}
For the chosen $\varepsilon$ take $B=D[1/\mu]$.
Since $t=\mu u_\varepsilon$, this is $D[1/t]$ as well;
the relation $pw=\mu^{p-1}$ makes $p$ a unit in $B$.
The injection $j$ extends uniquely to
\[
                    \iota:B\hookrightarrow B_{\mathrm{dR}}. \tag{LS.23}
\]
LS13 below checks the identification with all localization generators
in the source, so this notation does not substitute a formal field
model for its actual ring.

For a root $c\in\Omega$ of order $m$ prime to $p$, choose the unique
compatible $p$-roots within its finite group of $m$-th roots. They
define $c^\flat\in R$ with $(c^\flat)^m=1$ and
$(c^\flat)^\sharp=c$. In the complete DVR, both
$j([c^\flat])$ and $o(c)$ are roots of $X^m-1$ with residue $c$.
Their difference times the geometric sum is zero, and the latter
has nonzero residue $mc^{m-1}$. Thus
\[
                       j([c^\flat])=o(c).          \tag{LS.24}
\]
The subring of $B$ generated by $\mathbb Q_p$ and these lifts is a
field: every nonzero element is algebraic over $\mathbb Q_p$, so
its inverse is a polynomial in it, by its minimal polynomial with
nonzero constant term. Hence (LS.24) extends to a unique embedding
$u:\mathbb Q_p^{\mathrm{ur}}\to B$ with $\iota u=o$.

Reduction is injective on roots of unity of order prime to $p$.
The residue field of $\Omega$ is $\overline{\mathbb F}_p$; Hensel
lifting $X^{p^d}-X$ identifies each finite residue field with its
unramified extension. Its residue Frobenius therefore lifts to
compatible automorphisms, giving
$\sigma:\mathbb Q_p^{\mathrm{ur}}\to\mathbb Q_p^{\mathrm{ur}}$
with $\sigma(c)=c^p$ on these roots. Every finite intermediate
extension is cyclic unramified, so $\sigma$ preserves $K_0$ and
is bijective there. On its valuation ring it reduces to Frobenius.
Uniqueness on the generating roots gives
\[
 \varphi(u(c))=u(\sigma c),\qquad
 g(u(c))=u(gc),\qquad\iota(gb)=g(\iota b).         \tag{LS.25}
\]
The compatible roots, the tilt, $o$, $j$, $\iota$, $u$, $t$,
$\sigma$, and $\chi$ are now all on these same rings. The actions
extend to $B$, since (LS.22) makes the images of $t$ units there.
These give all maps and laws required by \texttt{Core}; LS13 below
identifies this localization with the source's precise period ring
and transports the construction. The Galois and Frobenius maps
are induced on coordinates and completions, and the coefficient
embeddings are the ones just constructed.

\paragraph{LS9: Frobenius is injective after divided-power completion.}
This needs a completion argument; injectivity on $A$ alone would
not suffice. Write $\xi=w+pa$ with $w=[\xi_0]$ a Teichm\"uller
element. Then $\varphi(w)=w^p$, $w\bmod p\ne0$, and
\[
 D_0\subset M=A[w/p]=A[\xi/p]\subset A_P,
 \qquad p^nM\bmod\xi^n\subset\operatorname{im}(A).
 \tag{LS.26}
\]
The first inclusion follows from $p^j/j!\in\mathbb Z_p$; the
last follows by expanding a polynomial in $\xi/p$ and separating
its terms of degree less than $n$.

We record the finite coefficient calculation used with this larger
ring. If $s\in M$ and, for every $r$, there are $u,v\in M$ with
\[
              \varphi(s+p^k u)=p^r v,
 \quad\text{then}\quad s\in p^kM.                 \tag{LS.27}
\]
Here $\varphi$ on $A_P$ is the extension of the bijective Witt
Frobenius. For completeness, put $I=(p,w)$. Regularity of $w$
modulo $p$ gives
$px\in I^{N+1}\Rightarrow x\in I^N$: expand
$I^{N+1}=Aw^{N+1}+pI^N$, reduce modulo $p$, then cancel $p$.
Every $v\in M$ has a denominator-clearing expression
$p^Nv=c$, $c\in I^N$.
For a polynomial $h=\sum h_iX^i$ define
\[
 F_0=0,\quad F_{N+1}=\varphi(h_N)w^{pN}+pF_N,
 \quad
 pF_N=p^N\varphi\!\left(\sum_{i<N}h_i(w/p)^i\right).
 \tag{LS.28}
\]
The precise induction used here is the following, with any integer
$l\ge1$ and recurrence $F_{N+1}=a_Nw^{lN}+pF_N$:
\[
 b w^{lN}+F_N\in I^{N+r},\quad i<N,\quad
 l(i+1)\le i+r+1
 \quad\Longrightarrow\quad a_i\in(p,w^l).             \tag{LS.28a}
\]
Induct on $N$. For the step from $N$ to $N+1$, put
$D=N+r+1$ and $A'=bw^l+a_N$. The hypothesis is
$A'w^{lN}+pF_N\in I^D$. If $D\le lN$, remove the first
term, which already belongs to $I^D$, and apply $p$-saturation
to obtain $F_N\in I^{N+r}$. The induction hypothesis with
$b=0$ handles $i<N$; the terminal $i=N$ is impossible since
its required inequality would give $l(N+1)\le D\le lN$.
If $D>lN$, put $e=D-lN>0$. Expand
$I^D=Aw^D+pI^{D-1}$ to write
\[
 A'w^{lN}+pF_N=w^Dc+pG,\qquad G\in I^{N+r}.
\]
Reduction modulo $p$ and cancellation of the regular element
$w^{lN}$ give $A'-w^ec=pd$ for some $d\in A$.
Cancel $p$ in the displayed equation to obtain
$d w^{lN}+F_N=G$. The induction hypothesis with $b=d$
handles $i<N$. At $i=N$, the required inequality is exactly
$e\ge l$, and
$a_N=pd+w^ec-bw^l\in(p,w^l)$. This proves (LS.28a)
with all its range conditions.

Now suppose $\varphi(h(w/p))=p^{p(i+1)+1}v$ with $v\in M$.
Write $p^{N_0}v=c\in I^{N_0}$ and take
$N=N_0+\deg(h)+i+1$ (the zero polynomial is immediate).
The last equality in (LS.28) gives
\[
 F_N=c\,p^{\deg(h)+i+1+p(i+1)}\in I^{N+p(i+1)}.
\]
Apply (LS.28a) with $l=p$, $r=p(i+1)$, $b=0$, and
$a_j=\varphi(h_j)$. Its inequalities hold, so
$\varphi(h_i)\in(p,w^p)$. Bijectivity of Witt Frobenius,
$\varphi(p)=p$, and $\varphi(w)=w^p$ then give
$h_i\in(p,w)$.

For $k=1$, apply this calculation to $h=f+p g$, where
$s=f(w/p)$ and $u=g(w/p)$ in the hypothesis of (LS.27), choosing
$r=p(i+1)+1$ separately for each coefficient $f_i$.
It follows that every $f_i\in(p,w)$, and substitution
$w=p(w/p)$ gives $s\in pM$. Repeated cancellation of $p$
proves (LS.27) for every $k$.

Now suppose $x\in D$ and $\varphi x=0$. Write
$x=s+p^kz$ with $s\in D_0$. Approximate $z$ arbitrarily closely
by elements of $D_0$. Applying Frobenius, and using (LS.15), gives
the hypothesis of (LS.27), so $s\in p^kM$.
For fixed $n$ choose $k=n+r$. By (LS.26) and (LS.17),
\[
 j(x)\bmod\xi^n
   =j(s)+p^{n+r}j(z)\bmod\xi^n
     \in p^r\bigl(A/(\xi^n)\bigr).
\]
This holds for every $r$; separation in LS4 makes the class zero.
Thus $j(x)=0$ and LS6 gives $x=0$. Frobenius on $D$ is injective.
It is injective on $B$ as well: clear a power of $t$, apply
Frobenius, and cancel the nonzero denominator using $D\hookrightarrow B$.

\paragraph{LS10: the integral Frobenius fixed ring is $\mathbb Z_p$.}
One has
\[
                         D^{\varphi=1}=\mathbb Z_p. \tag{LS.29}
\]
First $\varphi(D_0)\subset A+pD_0$. On a positive divided power
$a^n/n!$, $\theta a=0$, this follows from
$\varphi a\in pD_0$ and $v_p(n!)\le n-1$.
Approximation modulo $p$ extends the inclusion to $D$.
If $x=\varphi x$, choose $a\in A$ with $x\equiv a\pmod{pD}$.
Then $a^p-a\in pD$. It lies in $D_0$, so (LS.15) and
$\theta(D_0)\subset O$ imply
$(\theta a)^p-\theta a\in pO$. Its residue in the residue field
of $O$ therefore belongs to $\mathbb F_p$.
Choose $c\in\mathbb Z$ with $|\theta(a-c)|<1$.
Some power $\theta(a-c)^N$ is divisible by $p$ in $O$.
Lift its quotient by $\theta$, giving $b\in A$ with
$(a-c)^N-pb\in\ker\theta$. Its $p$-th power belongs to $pD_0$,
so $\bar a-c$ is nilpotent in $D/pD$. It also satisfies
$(\bar a-c)^p=\bar a-c$. A nilpotent satisfying this equation is
zero: multiply by the inverse of $1-(\bar a-c)^{p-1}$.
Hence $x\equiv c\pmod{pD}$.

As $p$ is regular and fixed by $\varphi$, the quotient $(x-c)/p$
is again fixed. Iterating with digits $0\le c_j<p$ gives
$x\equiv\sum_{j<n}p^jc_j\pmod{p^nD}$.
The digits converge to $z\in\mathbb Z_p$, and separation gives
$x=z$. Conversely Frobenius fixes $\mathbb Z_p$ by its Witt
definition (or by density of the integers and $p$-adic separation).

\paragraph{LS11: simultaneous theta vanishing and the Frobenius image.}
Two integral identities supply the fundamental intersection calculation:
\[
 \begin{split}
 &\theta(\varphi^na)=0\ (n\ge0),\ a\in A
                              \quad\Longrightarrow\quad \mu\mid a;\\
 &\varphi(D)\subset A+wD,
                 \qquad pw=\mu^{p-1}.
 \end{split}                                                   \tag{LS.30}
\]
For the first, let $\varepsilon_n=\varepsilon^{1/p^n}$,
$\xi_n=\sum_{i<p}[\varepsilon_{n+1}]^i$, and
$q_n=\prod_{j<n}\xi_j$. Then
$\mu=q_n([\varepsilon_n]-1)$.
Induction using $\ker\theta=(\xi)$ shows $q_n\mid a$:
after the first $n$ factors have been removed, apply $\theta\varphi^n$;
the removed factors have nonzero images $p$, so the next factor
must divide. Modulo $p$ this gives
\[
 |a_0|\le|\varepsilon-1|^{\,1-p^{-n}}\quad(n\ge0).
\]
Letting $n$ increase, the valuation-ring criterion implies
$(\varepsilon-1)\mid a_0$. Choose $b\in R$ with
$a_0=(\varepsilon-1)b$. Then $a=\mu[b]+pa'$.
Every $\theta\varphi^n(a')$ is zero, since the corresponding
values of $a$ and $\mu$ are zero and $p$ is regular in $O$.
Iteration and $p$-adic completeness prove $\mu\mid a$.

For the second identity, the geometric-sum polynomial gives
$\varphi\xi=pc+\mu^{p-1}=p(c+w)$, $c\in A$.
Thus
$\varphi(\xi^n/n!)=(p^n/n!)(c+w)^n\in A+wD$.
It follows for $D_0$ and then for $D$: successively write a dense
approximation to $y$, obtaining
$\varphi y=a_0+wz_0+p\varphi y_1$.
Iteration gives the convergent coefficients
$\sum p^ja_j\in A$ and $\sum p^jz_j\in D$, proving the claim.

\paragraph{LS12: integral Frobenius eigenvectors in the matching filtration.}
For $r\ge0$,
\[
 \varphi x=p^r x,\quad j(x)\in\operatorname{Fil}^r
 \quad\Longrightarrow\quad
 p^k x=z t^r\text{ for some }k\ge0, z\in\mathbb Z_p.
 \tag{LS.31}
\]
The case $r=0$ is LS10. For the induction step write
$\varphi x=p^{r+1}x=a+wz'$ by LS11. Every
$\theta_Dj(\varphi^n x)$ vanishes, as does
$\theta_Dj(\varphi^n w)$, the latter by
$p w=\mu^{p-1}$ and $\theta\varphi^n\mu=0$.
Thus $\theta\varphi^n a=0$ for every $n$, so $a=\mu b$.
With $t=\mu u$ as in LS7, set
\[
 x'=u^{-1}(pb+\mu^{p-2}z'),\qquad
                         p^{r+2}x=t x'.             \tag{LS.32}
\]
The exponent $p-2$ is nonnegative, also at $p=2$.
Apply Frobenius to the equality and cancel $pt$ in the domain $D$:
$\varphi x'=p^r x'$. Since $j(t)$ is a uniformizer and $p$ a
unit in $B_{\mathrm{dR}}^+$, (LS.32) gives
$j(x')\in\operatorname{Fil}^r$. Induction proves (LS.31), after
multiplication by one more power of $p$.

\paragraph{LS13: the actual crystalline localization and its fundamental intersection.}
Every generator $\mu_e=[e]-1$, with $e^\sharp=1$ and $e\ne1$,
is a unit in $D[1/t]$. Indeed LS7 gives
$t_e=\mu_eu_e\ne0$, $u_e\in D^\times$,
$\varphi t_e=pt_e$, and $j(t_e)\in\operatorname{Fil}^1$.
LS12 yields $p^kt_e=z t$ with $0\ne z\in\mathbb Z_p$.
Because $p$ is already a unit in $D[1/t]$, so is $z$; hence
$t_e$, and then $\mu_e$, is invertible. Conversely the source
localization inverts $\mu_\varepsilon$, and thus $t$.
The two localization maps are inverse and fix $D$. This identifies
the ring constructed here with the literal $B_{\mathrm{cris}}$
and the maps in LS8 with the source's maps. In particular the
constructed \texttt{Core} exists on the specified rings.

Now
\[
 B_{\mathrm{cris}}^{\varphi=1}\cap
        \iota^{-1}(B_{\mathrm{dR}}^+)=\mathbb Q_p.    \tag{LS.33}
\]
If $b$ belongs to the left side, write $bt^r=x\in D$.
Then $\varphi x=p^r x$ and $j(x)\in\operatorname{Fil}^r$.
Equation (LS.31) gives $p^kx=z t^r$, so $b=z/p^k$.
The reverse inclusion follows from the coefficient maps.
This proves exactly the \texttt{PeriodData.fund} field; a stronger
surjective fundamental exact sequence is not being asserted.

\paragraph{LS14: the de Rham invariants are exactly the finite base field.}
For every finite $K\subset\Omega$,
\[
                        B_{\mathrm{dR}}^{G_K}=o(K).       \tag{LS.34}
\]
First an invariant $z\in B_{\mathrm{dR}}^+$ with residue zero
vanishes. If $z=t^n y$, invariance and $g(t)=\chi(g)t$
give $g(y)=\chi(g)^{-n}y$. For $n>0$, its residue is a
nonzero Tate-twist invariant and is zero by LS3. Thus $t\mid y$.
Starting with $t\mid z$, induction makes $z$ divisible by every
power of $t$, and separation gives $z=0$.
More generally, if $g(y)=\chi(g)^n y$, $n\ge0$, then
\[
                y=t^n w,\qquad w\in(B_{\mathrm{dR}}^+)^{G_K}.
                                                               \tag{LS.35}
\]
For $n>0$, LS3 kills the residue of $y$; divide by $t$ and
induct with exponent $n-1$. The case $n=0$ is tautological.

Now clear the denominator of an invariant $b\in B_{\mathrm{dR}}$:
$y=t^n b\in B_{\mathrm{dR}}^+$ for some $n\ge0$.
Equation (LS.35) gives $y=t^n w$ with $w$ invariant, so $b=w$.
By LS1 its residue belongs to $K$. Subtract its unique algebraic
lift $o(k)$ from LS5. The difference is invariant with zero
residue, hence is zero by the first paragraph. Conversely $o(k)$
is invariant for $k\in K$, by equivariance of the algebraic
section. This proves (LS.34), including the exclusion of negative
filtration invariant elements.

\paragraph{LS15: the crystalline invariants are exactly the unramified base field.}
For the same $K$,
\[
                         B_{\mathrm{cris}}^{G_K}=u(K_0).
                                                               \tag{LS.36}
\]
Let $b$ be invariant. Frobenius commutes with $G_K$, and LS14
therefore gives elements $y_i\in K$ with
$\iota(\varphi^i b)=o(y_i)$. If $q\in\mathbb Q_p[X]$
annihilates $y_0$, injectivity of $\iota$ gives $q(b)=0$;
Frobenius fixes $\mathbb Q_p$, so $q(\varphi^i b)=0$ for
all $i$. The nonzero minimal polynomial of $y_0$ has finitely
many roots. Thus two iterates of $b$ agree, and injectivity of
Frobenius (LS9) permits cancellation of their common initial
iterate. There is an integer $N>0$ with $\varphi^N b=b$.

Choose a primitive root $\zeta\in\mathbb Q_p^{\mathrm{ur}}$
of order $p^N-1$. Its Frobenius orbit
$\zeta,\sigma\zeta,\ldots,\sigma^{N-1}\zeta$ consists of
$N$ distinct points: $p^i-1<p^N-1$ for $0<i<N$.
When $N=1$, this assertion has no distinctness condition, including
$p=2$, where $\zeta=1$. For every integer $n\ge0$, set
\[
 S_n=\sum_{i=0}^{N-1}u((\sigma^i\zeta)^n)\,\varphi^i b.
                                                               \tag{LS.37}
\]
Frobenius permutes the summands. Their images under $\iota$
are algebraic coefficients in $B_{\mathrm{dR}}^+$.
The fundamental intersection LS13 therefore gives
$S_n\in\mathbb Q_p$. Consequently every polynomial-weighted
sum in (LS.37), with coefficients in
$\mathbb Q_p^{\mathrm{ur}}$, belongs to $u(\mathbb Q_p^{\mathrm{ur}})$.
Take
$P(X)=\prod_{j=1}^{N-1}(X-\sigma^j\zeta)$.
All terms except the one with $i=0$ vanish, so
$u(P(\zeta))b\in u(\mathbb Q_p^{\mathrm{ur}})$.
Its nonzero coefficient is invertible, whence $b=u(c)$ for some
$c\in\mathbb Q_p^{\mathrm{ur}}$. Comparing with
$\iota b=o(y_0)$ and using injectivity of $o$ gives $c=y_0\in K$.
Thus $c\in K_0$. Conversely $u(K_0)$ is fixed by $G_K$ by
LS8. This proves the exact invariant-field assertion without
assuming a crystalline comparison theorem for representations.

\paragraph{LS16: all fields of the actual finite-extension period datum.}
For every finite $K\subset\Omega$, the preceding construction
gives the literal datum
\[
 \operatorname{PeriodData}
   (G_K,\mathbb Q_p,K_0,K,B_{\mathrm{cris}},B_{\mathrm{dR}})
                                                               \tag{LS.38}
\]
on the same two period rings, with no new choice of a model for
each $K$. We spell out the full record verification. Both $K$
and $K_0$ are finite-dimensional over $\mathbb Q_p$; the latter
is a subspace of $K$. The residue Frobenius $\sigma$ of LS8
restricts to a bijective $\mathbb Q_p$-map on $K_0$.
The restricted $G_K$-actions on $B_{\mathrm{cris}}$ and
$B_{\mathrm{dR}}$ are respectively $K_0$- and $K$-linear,
because those constants are fixed. Their inverses are the actions
of $g^{-1}$, and the group laws hold on coordinates and hence
on completions and localizations. Their invariant rings are
exactly the coefficient fields by LS14--LS15.

Frobenius is an injective $\mathbb Q_p$-algebra map by LS9;
it acts as $\sigma$ on $K_0$ and commutes with the group action
by LS8 and the commuting Witt operations. The map $\iota$ is an
injective $K_0$-algebra map and is equivariant. For each integer
$a$, take the actual $\operatorname{Fil}^a=\xi^aB_{\mathrm{dR}}^+$.
This is a $K$-subspace: $o(K)\subset B_{\mathrm{dR}}^+$ and
multiplication preserves the indicated power. It contains $1$
in degree zero, is decreasing, exhaustive, separated,
multiplicative and $G_K$-stable by LS5.
Take the restriction of $\chi$, valued in $\mathbb Q_p^\times$,
and the same $t$ of LS7. It is a unit in $B_{\mathrm{cris}}$,
$g(t)=\chi(g)t$, $\varphi(t)=pt$, and its image lies in
$\operatorname{Fil}^1\setminus\operatorname{Fil}^2$ because
it is a uniformizer. Finally LS13 gives the precise fundamental
intersection field. These are all fields of the source record.

The equalities in \texttt{IsRealPD} are literal identifications:
the prime is $p$; the two actions, $\varphi$, $\iota$, $t$,
$\chi$ and every filtration subset are those just specified;
and the restricted $\sigma$ is the same unramified Frobenius.
In particular, changing $K$ changes the coefficient subfields and
restricted group, but does not change the filtration subsets in
$B_{\mathrm{dR}}$.

\paragraph{LS17: the complete local prime, Ainf tie, and all-prime stage.}
Fix one construction of LS1--LS16 for each rational prime $p$.
It yields simultaneously
\[
 \exists(P:\texttt{LocPrime}\ p)\ (AI:\texttt{LocAinf}\ p\ P),
                 \texttt{LocTied}\ p\ P\ AI,\qquad
                 \texttt{LocStage}.                         \tag{LS.39}
\]
Here is the remaining field audit, including the ties to actual
rings. Take the underlying rings of $P$ to be the actual
$B_{\mathrm{cris}}$ and $B_{\mathrm{dR}}$, with their
$\mathbb Q_p$, unramified, and algebraic coefficient structures
from $u$ and $o$. If the source presents these as renamed types,
transport every operation along the identity ring equivalences
$e_B,e_D$. The datum for every finite $K$ is LS16.
The global actions, Frobenius, inclusion, $t$, and $\chi$ are
those of LS8; therefore every restriction equality in
\texttt{LocPrime} holds. Its equality of filtration subsets for
any two finite $K,K'$ follows from LS16. The actions on
unramified and algebraic constants are the ones in LS8 and LS5.
Frobenius on the valuation ring of $K_0$ reduces to $x\mapsto x^p$.

The cyclotomic clause includes its numerical approximation:
for $g\in G$ and $m\ge0$, choose the least nonnegative integer
$a$ representing $\chi(g)\bmod p^m$. Then
\[
 |\chi(g)-a|\le p^{-m},\qquad g(\zeta)=\zeta^a
                   \quad(\zeta^{p^m}=1).                  \tag{LS.40}
\]
The compatible primitive system $\zeta_m$ is the one chosen in
LS2. For every finite extension $L/\mathbb Q_p$, use its unique
extended absolute value and put $v(x)=-\log_p|x|$ for $x\ne0$,
$v(0)=\infty$. Algebraicity makes the nonzero values rational,
and $v(p)=1$. A finite complete DVR extension, as in NC1--NC2,
has $v(x)\ge0$ exactly on the integral closure of $\mathbb Z_p$.
This verifies both normalized valuation fields of $P$.

For $AI$, extend $G$ isometrically from $\Omega$ to $C$ and
act coordinatewise on $R=\varprojlim O/p$. The coordinate
formula holds because each action preserves $pO$ and commutes
with Frobenius. Take the same $\varepsilon$, its exact sharp-root
system, and the actual map $A\hookrightarrow B_{\mathrm{cris}}$.
This map is injective by LS5--LS6, is equivariant, and commutes
with Witt Frobenius by its construction. Its image under
$\iota$ is in filtration zero. For every finite $K$,
\[
 \iota(a)\in\operatorname{Fil}^1
      \quad\Longleftrightarrow\quad\theta(a)=0
      \quad\Longleftrightarrow\quad\xi\mid a\text{ in }A.
                                                               \tag{LS.41}
\]
This is the kernel calculation LS4, not merely divisibility in
a period field. For every prime-to-$p$ root $c$ and every tilt
representative satisfying the specified root-coordinate condition,
the Teichm\"uller image is $u(c)$ by LS8: uniqueness of the
root in the complete DVR proves the equality for all such choices.
These verify all \texttt{LocAinf} fields.

The tie uses $e_B,e_D$ just specified. It sends the prescribed
$\mathbb Z_p$ constants to the actual constants because $o$
and $u$ agree on them. Each algebraic $x\in\Omega$ has the
lift $o(x)$ with the required residue. Actions and Frobenius
are their actual coordinate constructions. The transported
inclusion satisfies \texttt{IsCrisDR}: its restriction to $D$
is the actual $j$, with the bounds proved in LS6, and its
extension is the localization map LS8--LS13. Filtrations agree
for every $K$ and degree. The same $t$ maps to $j(t)$, whose
formal logarithm identity was proved in LS7. Finally the
$A_{\inf}$ inclusion is the actual one. Thus every field of
\texttt{LocTied}, including its integral and logarithmic clauses,
has been verified. Choosing these triples for all primes gives
\texttt{LocStage}; this uses the constructions at every prime,
even when a later application has an empty local index type.

\paragraph{LS18: the global cyclotomic, residue, and prime-label suppliers.}
For the remainder fix the source $B:\texttt{Base}$, its host $H$,
$C:\texttt{Cols}(B)$, and a row object $W:\texttt{RowsRest}(B,C)$.
Write $r=B.R.r\ge2$ and $f=B.R.f>0$.
The nonempty pieces of the bare \texttt{NData} and \texttt{NTypes}
construction have the following elementary suppliers.

For each rational prime $p$, let $\Omega_H=\overline{H.E}$ and
$S_p$ be the finite set of primes of $H.E$ above $p$.
Finiteness follows from the finite factorization of the nonzero
ideal $(p)$. Choose a valuation above each such prime.
The cyclotomic character on
$\operatorname{Gal}(\Omega_H/H.E)$ is unramified outside $S_p$.
Indeed inertia at residue characteristic $q\ne p$ fixes the
residue of each $p^n$-th root of unity. The quotient of two such
roots with the same residue is $1$: if $z^m=1$, $q\nmid m$,
and $z\equiv1$, cancellation of $z-1$ in the geometric sum
would force its residue $m$ to be zero unless $z=1$.
Thus inertia fixes every such root, and the inverse-limit
character is trivial on inertia. It descends through the quotient
unramified outside $S_p$. Together with the chosen valuations,
this is the actual \texttt{LData} used in the source. Set
\[
 s_H(p,e)=
 \begin{cases}
  (\operatorname{shaLen}(L_p,e).\mathrm{toNat}:\mathbb Q),
                                      &p\text{ prime},\\
  0,&p\text{ not prime}.
 \end{cases}                                                \tag{LS.42}
\]
If the extended-natural length is the finite integer $n$,
then $s_H(p,e)=n$. This proves the exact conditional
\texttt{sha\_\allowbreak eq} field. It neither proves finiteness nor rules
out infinite length, whose \texttt{toNat} is zero.

There is a complex-conjugation lift on $\Omega_H$. Embed this
algebraic closure into $\mathbb C$ over $\mathbb Q$.
Its image is the field of algebraic numbers, stable under complex
conjugation. Conjugate through the embedding; injectivity gives
an involution, and the host's CM embedding identity identifies
its restriction to $H.E$ with the prescribed conjugation.

For each rational prime $g$, there is a multiplicative map
$\overline{\mathbb F}_g\to\mathbb C$ injective on nonzero
elements. To construct it directly, exhaust the algebraic closure
by $\mathbb F_{g^{n!}}$. Their multiplicative groups are cyclic
and nested. Embed each into the complex roots of unity, extending
the previous embedding: a prescribed primitive $d$-th root
extends to a primitive $N$-th root whenever $d\mid N$, by the
surjectivity of
$(\mathbb Z/N)^\times\to(\mathbb Z/d)^\times$, proved prime
by prime and by the Chinese remainder theorem. Put the value at
zero equal to zero. This supplies the exact multiplicative
\texttt{teichC} interface; additivity is not required.
For any prime $\ell$, choose a reduction from the ring of all
algebraic integers in $\mathbb C$ to
$\overline{\mathbb F}_\ell$. Such a map exists by taking a
maximal ideal above $(\ell)$ in this integral extension of
$\mathbb Z$ and embedding its algebraic residue field in the
algebraic closure. Every nonzero value of the preceding map is
a root of unity, hence integral. Composing gives the required
multiplicative map between residue algebraic closures. It is
injective on elements of order prime to $\ell$: their quotient
reducing to $1$ is $1$ by the same geometric-sum argument.
This proves exactly the restricted injectivity in
\texttt{iota\_\allowbreak inj}, without asserting an embedding of fields of
different characteristics.

Finally the source congruences $B.g\_cond$ make each $g_k$
split completely in $H.E$, by the host splitting calculation
already reconstructed in AH. The Galois group acts transitively
on primes above $g_k$; equality of their number with its order
makes this action free. Host conjugation is a nonidentity
involution, so its action on these primes has no fixed point.
There are $2f$ primes. Choose one prime from each conjugate pair,
giving an injective map $\operatorname{Fin}(f)$ into the primes
above $g_k$, with no two selected primes conjugate. The quadratic
field $H.K$ has discriminant $-H.d$; its identity extension over
$H.K$ is abelian, its ramified rational primes divide the source
conductor, and the same congruence makes $g_k$ split there.
These verify the selected-prime and ray-field clauses without
requiring a nontrivial class-field extension.

\paragraph{LS19: an actual tame inertia generator at every selected place.}
Let $V$ be a valuation of $\Omega_H$ above a finite place of the
number field $H.E$, and $I$ its inertia subgroup. There is
$\gamma\in I$ satisfying the source's \texttt{TameGenAt}:
\[
 I\subset\overline{\langle\gamma,P_V\rangle},\qquad
 P_V=\{\tau\in I:|\tau(y)/y-1|<1\text{ for every }y\ne0\}.
                                                               \tag{LS.43}
\]
Here $P_V$ is the explicitly defined wild subgroup, so (LS.43)
does not presuppose a tame-inertia classification theorem.

We use a compactness lemma. Suppose a compact group $I$ has
continuous maps $\theta_n$ to finite cyclic groups, with nested
kernels. The set of elements whose $\theta_n$-image generates
$\theta_n(I)$ is nonempty and closed. These sets are nested:
a generator maps to a generator in every quotient.
Compactness gives an element $\gamma$ in their intersection.
For each $\sigma\in I$, the sets
\[
 \overline{\langle\gamma\rangle}\cap
                       \theta_n^{-1}(\theta_n\sigma)
\]
are again nonempty, closed and nested, because a power of
$\gamma$ has any prescribed image in the finite cyclic group.
Their intersection contains an $h$ with
$h^{-1}\sigma\in\bigcap_n\ker\theta_n$.
Thus it suffices to identify this common kernel with $P_V$.

The restriction of the valuation to $\mathbb Q$ is controlled
by a single prime $q$: all other rational primes have norm $1$.
For any algebraic $y\ne0$, its polynomial equation has at least
two terms of largest norm; otherwise their sum could not be zero.
Equating those two norms gives
$|y|^a=|q|^b$ for some $a>0$, $b\in\mathbb Z$.
Choose $r_n\in\Omega_H$ with $r_n^{n!}=q$.
For $\sigma\in I$, take the residue class of
$\sigma(r_n)/r_n$ in the group of norm-one elements modulo
principal units, and call it $\theta_n(\sigma)$.
It is a homomorphism, because inertia acts trivially on residues.
Its image is contained in the reduction of the finite group of
$n!$-th roots of unity, hence is finite cyclic. Its fibers are
closed: the action on the fixed algebraic element $r_n$ factors
through a finite Galois extension. The kernels are nested.
Indeed $r_{n+1}^{n+1}/r_n$ has norm $1$, and the action of
inertia on this quotient changes it only by a principal unit.

If $\tau$ belongs to every kernel and $y\ne0$, choose $a,b$
as above and $n$ with $a\mid n!$. The element
$r_n^{b n!/a}$ has the same norm as $y$, and its inertia
coboundary is a principal unit. The quotient of $y$ by this
element is a unit, whose inertia coboundary also has residue $1$.
Thus $|\tau(y)/y-1|<1$, so $\tau\in P_V$.
The reverse inclusion follows by testing $y=r_n$.
Inertia is a closed subgroup of the profinite absolute Galois
group and is compact. The lemma therefore proves (LS.43).
Applying it to every selected prime from LS18 supplies the
literal \texttt{tg}, \texttt{tg\_\allowbreak mem}, and \texttt{tg\_\allowbreak gen}
fields; no selected-place generator is omitted by an empty-index
argument.

\paragraph{LS20: the full empty-index algebraic and automorphic records.}
Suppose
\[
 W.\mathrm{shaLen}=s_H,\qquad
 W.\eta_J(j,w)=0\ (1\le j\le m),\qquad
 W.e\mathrm{Op}=0,\quad W.eM_t=0.                    \tag{LS.44}
\]
The source constructs the following records on exactly these
$B,C,W$. Take $\texttt{NData.Om}=\Omega_H$, the chosen
valuations and $L_p$ from LS18, and put
\[
 \mathrm{Loc}=\varnothing,\quad\mathrm{locs}(j,w)=\varnothing,
 \quad\mathrm{PlG}(k)=\operatorname{Fin}(f),\quad
 \mathrm{lab}=(X-1)^r,\quad\chi_T=1.                  \tag{LS.45}
\]
The split-place and cardinality fields come from LS18 and the
\texttt{RowsRest} finite-group cardinal formula. The polynomial
is monic of degree $r$ with nonzero constant term. Its Sha clause
is precisely the finite-value assertion in LS18. Every field
indexed by $i\in\mathrm{Loc}$, including its local DVR,
representation, eigenvector, auxiliary-prime, and local-generator
clauses, is a function on the empty type. This supplies all of
\texttt{NData}, without introducing a local representation.

For clarity the nonempty part of \texttt{NAut} is also explicit.
At each place let a local parameter be a finite multiset of
nonzero complex numbers. Its dimension is cardinality; zero
is the empty multiset; addition is multiset sum. Twisting by a
real $s$ multiplies every entry by $\mathrm{Norm}(w)^s$;
character twisting multiplies by the product of the entries of
the character parameter. These scalars are nonzero. Set
\texttt{unr} and \texttt{mono0} identically true,
\texttt{frobEig} equal to that multiset, and temperedness equal
to the condition that all entries have absolute value $1$.
Set the inertia polynomial and \texttt{items2} equal to
$(X-1)^{\mathrm{dimension}}$. Multiset addition and scaling
verify the dimension, sum, twist, Frobenius, and inertia laws;
$(X-1)^{a+b}=(X-1)^a(X-1)^b$ gives their polynomial identity.

A global parameter of rank $a$ is an arbitrary family of these
local multisets of size $a$, together with one complex multiset
of size $a$ for each embedding of $H.E$ into $\mathbb C$.
Conjugate self-duality means inversion of the local multiset at
the conjugate place. Regular algebraicity means that each
infinite multiset has no repetitions and consists of elements
$n+(a-1)/2$, $n\in\mathbb Z$. The global twist scales the
finite entries and adds the sum of the twisting infinite entries.
Inversion of a product proves preservation of self-duality;
cardinality and the one-element case give the remaining stated
twist laws. Choose both \texttt{etaT} and its inverse to have
finite multiset $\{-1\}$ and infinite multiset $\{0\}$.
The inert-place and inverse clauses hold, and at a split pair
$(-1)(-1)=1$. Use the actual selected primes of LS18 for
\texttt{plGw}. All remaining \texttt{NAut} clauses have an
index in $\mathrm{Loc}$ and are vacuous. These are literal record
models; they are not asserted to classify automorphic representations.

In \texttt{NBC} take both representation and weight indexing
types empty. In \texttt{NTw} take the actual CM-conjugation
lift from LS18 and $\mathrm{Tw}=\varnothing$. Every other
field is indexed by $\mathrm{Tw}$ or $\mathrm{Loc}$, including
coefficient-DVR, character, Galois, and local-global ties.
For \texttt{NExp} use the zero distinguished subspace and zero
coefficient function, and use the empty functions for its twist
clauses. All \texttt{NTh} fields are twist-indexed and are
supplied by the same empty functions. This verifies the actual
constructors \texttt{nBC}, \texttt{nTw}, \texttt{nExp}, and
\texttt{nTh} with their exact common \texttt{NData}; no
existence assertion over a nonempty twist set is used.

\paragraph{LS21: every type, projector, bridge, and row tie in the bare NStage.}
Take the model of LS20 and an arbitrary supplied
$L:\texttt{LocStage}$. In \texttt{NTypes}, set the ray field
to $H.K$ with its identity extension, and use the conductor,
splitting, and disjoint \texttt{BB} data of $B$, checked in LS18.
Use the block multiset $\{(r,1)\}$ and the trivial residual
representation. Its kernel is the whole Galois group and is open.
Use the actual tame generators of LS19 and the multiplicative
residue maps of LS18. The inertia characteristic polynomial is
correct because the roots of $(X-1)^a$ are $a$ copies of $1$,
and every multiplicative map sends $1$ to $1$.

Define, in the relevant finite group algebra,
\[
 e_f=\begin{cases}1,&f=(X-1)^r,\\0,&\text{otherwise}.
       \end{cases}                                      \tag{LS.46}
\]
These elements are central, idempotent, and pairwise orthogonal.
The finite set of monic degree-$r$ polynomials with nonzero
constant coefficient contains $(X-1)^r$, and summing the $e_f$
over it gives $1$. Central actions on the ambient and middle
spaces are the augmentation map followed by scalar identity
endomorphisms; this is an algebra homomorphism. Take the type
index to be a singleton, every type label to be $(X-1)^r$,
each local type multiset to consist of $r$ copies of $1$, and
both type subspaces to be the whole space. These choices verify
the remaining nonempty \texttt{NTypes} clauses. Its two ties
quantified by a local index are vacuous.

A good polynomial is squarefree. Since $r\ge2$, it cannot be
$(X-1)^r$, so its $e_f$ is zero. There is a selected prime
because $f>0$. Each product indexed by a good tuple therefore
has a zero factor, and the sum $e_0$ over good tuples is zero.
This is the full calculation underlying SE13. By (LS.44), the
augmentation image of $e_0$ equals the required zero operators
on both ambient and middle spaces. Every coefficient of $e_0$
is zero, so the required coefficient-field equality with
\texttt{NExp.eLL0} also holds.

Put all these records together and retain $L$ literally as the
local-stage field. The bridge fields \texttt{nbr},
\texttt{nbr\_\allowbreak ind}, and \texttt{self\_\allowbreak good} have local indices;
\texttt{ntwbr} and \texttt{ntwbr\_\allowbreak ind} have twist indices.
Their types are empty. The local, base-change, residual,
middle-parameter, $Z$, and $Z\mathrm{good}$ ties likewise
have local indices. Finally, for $1\le j\le m$,
\[
 W.\eta_J(j,w)=0
       =\sum_{i\in\mathrm{locs}(j,w)}\eta_{\mathrm{Loc}}(i).
                                                               \tag{LS.47}
\]
Thus no field remains unchecked, and we obtain the literal endpoint
\[
 \exists N:\texttt{NStage}(B,C,W),\qquad N.\mathrm{toLocStage}=L.
                                                               \tag{LS.48}
\]
In particular LS17 supplies such an $L$. For arbitrary $W_0$,
SE13's exact update $W=\mathrm{adapt}(W_0)$ satisfies (LS.44).
This proves \texttt{exists\_\allowbreak nStage\_\allowbreak of\_\allowbreak rows},
\texttt{exists\_\allowbreak nStage\_\allowbreak adapt}, and their existential row-object
version, including preservation of the same prescribed $L$.

\paragraph{LS22: the bare Later record now exists, with its exact boundary.}
Let $I$ be the supplied common column input, let $l:I.\mathrm{Laws}$,
and suppose \texttt{HasLetters} holds for the same $B$.
Set $C=I.\mathrm{cols}(l)$, and take exactly the row model
$W_0$ and update $W=\mathrm{adapt}(W_0)$ of SE14.
By LS17 and LS21, there is an
$N:\texttt{NStage}(B,C,W)$. SE6--SE12 supply the other
three auxiliary stages on this same $B,C,W$, and SE14 now gives
\[
                         \operatorname{Nonempty}
                              (\texttt{Later}(B,C)).       \tag{LS.49}
\]
This discharges the full bare \texttt{LocStage}/\texttt{NStage}
existence input left open there, including the specified integral
period rings and their ties. It neither replaces an already fixed
later stage by a different one nor supplies extra properties that
are absent from the bare record.

The distinction from arithmetic choice is forced by the exact
index types. The constructed stage has both
$\mathrm{Loc}=\varnothing$ and $\mathrm{Tw}=\varnothing$.
The arithmetic \texttt{NChoice} needs a function from every
coefficient-field prime to $\mathrm{Tw}$, even at a reference
weight. A coefficient number field has a prime above $2$, so
this domain is nonempty and such a function cannot exist for
this particular stage. The empty-local-index hypotheses used in
RC do not remove that requirement: its additional twist-reading
hypothesis must still be supplied. No crystalline, de Rham, or
Wach comparison for a nonempty family of geometric representations
has been deduced merely from the period-ring construction.
Similarly the model's zero projectors do not furnish positive rank
bounds or the ties \texttt{RankTie}, \texttt{MuIsW}, and the
same-object arithmetic data needed by later quantitative arguments.
These are the precise continuing boundaries after (LS.49).


\section{The original sized input and its compatible preliminary geometry}
\label{sec:initial-geometry}

\status{IG1--IG11 construct an original preliminary datum, its integral
norm factors, diagonal spaces, letter identities and exact positive
integer $d_u$. They prove the literal \texttt{exists\_\allowbreak pre\_\allowbreak sized}
conclusion and supply both coordinate/scalar pins after the prescribed
orientation. IG12 constructs the specified \texttt{baseVol} empty-jet
stage. Its chosen-volume definition is distinguished from an assertion
that this number is the geometric covolume. The final paragraphs
reconstruct the recursive-existence implication with its mass-window
and actual recursive-data premises retained.}

The first prime selection below is qualitative. Its purpose is to
construct one actual preliminary input. It gives no prime in a
prescribed short interval, no uniform bound for that prime, and no
large-scale \texttt{ChGood} assertion. In particular its use in the
first eleven paragraphs does not discharge \texttt{LevelPrime}.
All fields, inclusions, frames and norms are those of the displayed
construction. A later coefficient-field extension may preserve them;
it is not identified with an arbitrary previously chosen base.

\paragraph{IG1: partial ray zeta functions from a shifted plane lattice.}
Let $K$ be an imaginary quadratic field, $\mathcal O=\mathcal O_K$,
and $M\geq2$ an integer. The ray group $G$ modulo $M\mathcal O$ is
the group of fractional ideals prime to $M$, modulo principal ideals
generated by elements congruent to $1$ at the primes of $M$.
It is finite. Indeed the ordinary ideal class group is finite by the
short-vector argument in TL14; the kernel of the ray-to-ordinary
class map is a quotient of the finite group
$(\mathcal O/M\mathcal O)^\times$. For this latter assertion, a
principal ideal prime to $M$ is classified, up to a generator
congruent to $1$, by the residue of its generator, modulo residues
of global units. A denominator prime to $M$ can be inverted in
this finite ring. CRT supplies representatives prime to $M$ in
every ordinary ideal class. These observations also identify the
source's finite quotient of the prime-to-$M$ ideal monoid with $G$;
thus the argument applies to any literal \texttt{RayData}, through
its stated quotient isomorphism.

Fix a complex embedding $\sigma:K\hookrightarrow\mathbb C$, and
write $A_0>0$ for the Euclidean covolume of $\sigma(\mathcal O)$.
For a class $c\in G$ choose a nonzero integral ideal $\mathfrak a$,
prime to $M$, in class $c^{-1}$. Since
$\mathfrak a+M\mathcal O=\mathcal O$, choose
$\alpha_0\in\mathfrak a$ with $\alpha_0\equiv1\pmod {M\mathcal O}$.
Set
\[
 \Lambda=\sigma(M\mathfrak a),\qquad z=\sigma(\alpha_0),
 \qquad A=\operatorname{covol}(\Lambda)
       =M^2\mathrm N(\mathfrak a)A_0.                  \tag{IG.1}
\]
The elements of $\mathfrak a$ congruent to $1$ modulo $M$ form
$\alpha_0+M\mathfrak a$: coprimality gives
$\mathfrak a\cap M\mathcal O=M\mathfrak a$.
This coset excludes zero, since $M\mathcal O$ is a proper ideal.
The map
\[
 \alpha\longmapsto(\alpha)\mathfrak a^{-1}
                                                               \tag{IG.2}
\]
maps that coset onto the integral ideals in ray class $c$.
Its fibers are the orbits of
$U_M=\{u\in\mathcal O^\times:u\equiv1\pmod M\}$.
This group is finite, since $|\sigma(u)|=1$ and a discrete lattice
meets the unit circle in finitely many points. Its action is free;
put $w_M=|U_M|>0$. Conversely, if two elements in (IG.2) give the
same ideal their quotient is a unit, and their local congruences
show that this unit belongs to $U_M$. Thus the multiplicity is
exactly $w_M$, not the number of all units unless these coincide.

For $\Re s>1$, counting lattice points in disks gives absolute
convergence and the exact identity
\[
 \zeta_c(s):=\sum_{\substack{\mathfrak b\subset\mathcal O\\
                       (\mathfrak b,M)=1,\ [\mathfrak b]=c}}
                \mathrm N(\mathfrak b)^{-s}
 =\frac{\mathrm N(\mathfrak a)^s}{w_M}
       \sum_{\lambda\in z+\Lambda}|\lambda|^{-2s}.    \tag{IG.3}
\]
For completeness the required counting bound follows by putting
disjoint disks of a fixed sufficiently small radius about lattice
points. All disks with centers in a disk of radius $R$ lie in a
disk of radius $R+O(1)$, so their number is $O(1+R^2)$.
Dyadic annuli then prove convergence in (IG.3), locally uniformly
in $\Re s>1$.

Put $\theta(t)=\sum_{\lambda\in z+\Lambda}e^{-\pi t|\lambda|^2}$.
The Gaussian Fourier transform and Fourier expansion of its
periodization give
\[
 \theta(t)=\frac1{At}
       \sum_{\nu\in\Lambda^*}
          e^{-\pi|\nu|^2/t}e^{2\pi\mathrm i\langle\nu,z\rangle}.
                                                               \tag{IG.4}
\]
Here $\Lambda^*$ is the Euclidean dual lattice for the pairing
appearing in the phase. To justify the formula directly, periodize
$e^{-\pi t|x|^2}$ on a fundamental parallelogram. Its $\nu$th
Fourier coefficient is $A^{-1}t^{-1}e^{-\pi|\nu|^2/t}$, by the
one-dimensional Gaussian integral in each orthogonal coordinate.
Both lattice sums and their derivatives converge uniformly on
compact sets; Fourier summation therefore gives (IG.4).
Because $0\notin z+\Lambda$, $\theta(t)$ decays exponentially as
$t\to\infty$. The nonzero terms of the dual sum give
$\theta(t)-(At)^{-1}=O(t^{-1}e^{-c/t})$ as $t\downarrow0$.
These estimates hold after multiplication by every fixed power
of $t$ and $\log t$.

Consequently
\[
 E_c(s)=\int_1^\infty\theta(t)t^{s-1}\,dt
       +\int_0^1\bigl(\theta(t)-(At)^{-1}\bigr)t^{s-1}\,dt
\]
is entire, by locally uniform differentiated integration.
Mellin integration of each Gaussian in the half-plane of absolute
convergence yields the continuation
\[
 \zeta_c(s)=
  \frac{\mathrm N(\mathfrak a)^s\pi^s}{w_M\Gamma(s)}
          \left(E_c(s)+\frac1{A(s-1)}\right).           \tag{IG.5}
\]
The reciprocal gamma function is entire with simple zeros at the
nonpositive integers. Hence (IG.5) has no pole except a simple
pole at $1$, whose residue is the same positive number for every
class:
\[
          \rho=\frac{\pi}{w_M M^2 A_0}>0.
                                                               \tag{IG.6}
\]
In particular $Z_c(s)=(s-1)\zeta_c(s)$ is entire,
$Z_c(1)=\rho$ and $Z_c(-1)=0$. This proves the actual partial-zeta
continuation assertion, including its common residue and its
negative-integer value. No nonvanishing theorem or prime existence
has been used in this argument.

\paragraph{IG2: nonvanishing at one and prime norms in every ray class.}
Write $h=|G|$. For each complex character $\chi$ of $G$, put
$L_\chi(s)=\sum_c\chi(c)\zeta_c(s)$ initially in $\Re s>1$.
Unique factorization of nonzero ideals gives
\[
 L_\chi(s)=\prod_{\mathfrak p\nmid M}
       (1-\chi([\mathfrak p])\mathrm N(\mathfrak p)^{-s})^{-1}.
                                                               \tag{IG.7}
\]
Absolute convergence follows from (IG.3), summed over the finite
class group. Character orthogonality and (IG.6) show that every
nontrivial $L_\chi$ is entire. The principal $L_1$ has just one
pole, at $1$, with residue $h\rho$, and $L_1(-1)=0$.

Here is the nonvanishing argument with its analytic domain made
explicit. Suppose $L_{\chi_0}(1)=0$ for some $\chi_0\ne1$.
Then $F(s)=\prod_\chi L_\chi(s)$ is entire, since this zero
cancels the only possible pole. If $d_{\mathfrak p}$ is the order
of $[\mathfrak p]$ in $G$, finite character duality gives
\[
 F(s)=\prod_{\mathfrak p\nmid M}
    (1-\mathrm N(\mathfrak p)^{-d_{\mathfrak p}s})^{-h/d_{\mathfrak p}}
      =\sum_{n\ge1}b_n n^{-s}\quad(\Re s>1),
 \qquad b_n\ge0,\quad b_1=1.                          \tag{IG.8}
\]
The coefficient expansion is absolutely convergent in that stated
half-plane. Fix a real $\sigma_0>1$. Termwise differentiation there
is justified because $(\log n)^r\ll_{r,\delta}n^\delta$ for every
$\delta>0$; choose $\delta<\sigma_0-1$. For any real
$\sigma<\sigma_0$, the entire Taylor series of $F$ at $\sigma_0$
has nonnegative terms after inserting these derivatives:
\[
 \begin{split}
 F(\sigma)
 &=\sum_{r\ge0}\frac{(\sigma_0-\sigma)^r}{r!}
           \sum_{n\ge1}b_n(\log n)^r n^{-\sigma_0}\\
 &=\sum_{n\ge1}b_n n^{-\sigma}\ \geq1.               
 \end{split}
 \tag{IG.9}
\]
Tonelli justifies the interchange; finiteness of the first sum
comes from its equality with the entire Taylor series. At
$\sigma=-1$ this contradicts $F(-1)=0$, forced by the principal
factor. Thus every nontrivial $L_\chi(1)$ is nonzero. This proof
uses the entire continuation obtained under the supposed zero;
it assumes neither an abscissa bound for (IG.8) nor the desired
nonvanishing assertion.

For real $s>1$, expand the logarithm in (IG.7). The terms of
degree at least two have a uniformly bounded total as $s\downarrow1$:
their absolute value is at most a constant times
$\sum_{\mathfrak p}\mathrm N(\mathfrak p)^{-2}<\infty$.
For nonprincipal $\chi$, the full logarithm stays bounded near
$1$. Indeed $L_\chi$ is nonzero on a small disk about $1$ and
has an analytic logarithm there; its difference from the Euler
logarithm on the real interval is a constant multiple of
$2\pi\mathrm i$. For the principal character its logarithm is
$-\log(s-1)+O(1)$ by the positive residue. Orthogonality now gives,
for each $c\in G$,
\[
 \sum_{[\mathfrak p]=c}\mathrm N(\mathfrak p)^{-s}
                   =-\frac1h\log(s-1)+O(1).           \tag{IG.10}
\]
In a quadratic field any non-prime prime-ideal norm is $p^2$.
Their contribution at $s\ge1$ is bounded by
$2\sum_{p}p^{-2}$; there are also only finitely many ideals of
bounded norm. The divergence in (IG.10) therefore leaves, in
every ray class and above any given bound, a prime ideal whose
norm is a rational prime. This is exactly the prime-in-class
conclusion used in \texttt{ExistsNormPrime}.

\paragraph{IG3: integral norm factors and rational progression primes.}
Let $M\ne0$, let $a\in\mathcal O_K$ be coprime to $M$, and let
$N$ be a prescribed bound. For $M\ge2$, apply IG2 to the class
of $(a)$ and choose $\mathfrak p$ with prime norm $p>N$.
Ray equivalence supplies nonzero $x,y\equiv1\pmod M$ with
$x\mathfrak p=y(a)$. The containment
$ya\in x\mathfrak p\subseteq x\mathcal O_K$ gives
$ya=x\pi$ for some $\pi\in\mathcal O_K$.
Cancellation of the nonzero ideal $(x)$ gives
$\mathfrak p=(\pi)$, and
\[
       \pi-a=(y-1)a-(x-1)\pi\in M\mathcal O_K.
                                                               \tag{IG.11}
\]
The field norm is $|\sigma(\pi)|^2>0$, so
$N_{K/\mathbb Q}(\pi)=p$. If $M=1$, use $M=2,a=1$
and discard its extra congruence.

When $K$ is the source host's specified imaginary quadratic
subfield of $E$, map this actual integral element into $\mathcal O_E$.
Conjugation on $E$ restricts to the nontrivial automorphism of $K$;
hence $\pi\bar\pi=p$ in $E$. For a rational integer $a$, write
$\pi=a+M\beta$. Then
\[
 \frac{p-a^2}{M}=a(\beta+\bar\beta)+M\beta\bar\beta
                                                               \tag{IG.12}
\]
is both rational and an algebraic integer, and therefore an integer.
Thus $p\equiv a^2\pmod M$, with the precise integral factor and
the precise congruence required by \texttt{exists\_\allowbreak norm\_\allowbreak prime}.

We also need ordinary rational primes in a reduced residue class;
they can be obtained from the same argument without a new
nonvanishing input. Given $M>0$ and $r$ coprime to $M$, choose a
prime $\ell\equiv3\pmod4$ not dividing $M$. Such primes exist
outside every finite set: $4\prod\ell_i-1$ has an odd number of
prime factors congruent to $3$ modulo $4$ and none of the listed
ones. The quadratic field $K_0=\mathbb Q(\sqrt{-\ell})$ has
discriminant $-\ell$, and is unramified at every prime dividing
$M$, including $2$. At those primes its integral quadratic algebra
is split or unramified. The norm on its units is onto: reduce to
the surjective finite-field norm and lift by the trace correction
in (DR.5)--(DR.6), which works also at $2$; for a split algebra use
$(r,1)$. Choose such local units with norm $r$ and apply CRT to
their residues modulo the prime powers in $M$. This gives
$a\in\mathcal O_{K_0}$ with $N(a)\equiv r\pmod M$ and
$(a,M)=1$. Applying (IG.11) in $K_0$ gives prime norms
$p>N$ with $p\equiv N(a)\equiv r\pmod M$.
Thus all progression-prime selections below follow from IG1--IG3;
no assertion about primes in a shrinking interval has entered.

\paragraph{IG4: an actual cyclic arithmetic host with the required labels.}
The source provides the particularly simple rank record
$R=\texttt{Rt}(0)$: its frame is the cyclic frame of exponent $2$,
\[
       k=8,\qquad m=2,\qquad t=15,\qquad n=4,\qquad f=8.
                                                               \tag{IG.13}
\]
These values satisfy every parity and rank inequality directly:
$r=k+m-1=9$, $r+5\le t$, $3m<t$, $m\le f\le k$ and
$m^2\le k+m(t+1)$. Its map $h$ is the isomorphism
$j\mapsto[5^j]$ from SZ13, and $o=1$, so $o^{64}\le2^e$
with $e=2$.

More generally, for the same cyclic frame at any $e\ge2$, and
given finite bounds $B,N$, put
$P_B=\prod_{0\le i<B}(2i+1)$, with empty product $1$.
It is odd and
$\gcd(8P_B-1,16P_B)=1$. IG3 gives a prime
\[
 d>N+4tr+3,\qquad d\equiv8P_B-1\pmod {16P_B}.
                                                               \tag{IG.14}
\]
Then $d\equiv7\pmod {16}$ and $-d\equiv1\pmod p$ for
every odd prime $p\le B$. The residue $1+6d$ is coprime to
$8d$ and is $3$ modulo $8$. Choose a prime
\[
 q>N+4tr+3,\qquad q\equiv1+6d\pmod {8d}.             \tag{IG.15}
\]
We have $q\equiv1\pmod d$ and $q,d\equiv3\pmod4$.
Quadratic reciprocity therefore gives $(-d/q)=1$:
the signs in $(-1/q)$ and the reciprocity law for $(d/q)$
cancel. These are exactly \texttt{HostHyp}, with $q\ne3$.

Use AH1--AH5 with the trivial prescribed-character group,
$Q_S=2^{e+2}$, and this actual cyclic frame. Concretely the field
is the fixed field of the kernel of
\[
 (\mathbb Z/Q_Sqd\mathbb Z)^\times\longrightarrow
       \{\pm1\}^2\times\mathbb Z/2^e\mathbb Z,
 \quad u\longmapsto\bigl((u/q),(u/d),[u\bmod Q_S]\bigr),
                                                               \tag{IG.16}
\]
where the last class is read through the same $[5]$ generator.
The sign twist of AH3 is trivial here because $[-1]=1$ in the
last quotient. CRT proves surjectivity. AH4 supplies the actual
square roots of $-q,-d$; AH17 supplies all fields of
\texttt{HostDer} for this raw host. Exactness and
$Q_S\le(2^e)^2$ follow from AH5 and $e\ge2$.

Take $\mathfrak c_K=\mathcal O_K$ and the half-different ideal of
AH13. A rational integer $a>1$ is not a unit in $\mathcal O_K$,
as its norm is $a^2>1$; this proves the required primitivity of
$\mathfrak c_K$. Its norm is one, so all its coprimality clauses
hold. The small-splitting assertion is (IG.14); no odd prime
divides $Q_S$. Choose an embedding in the nonempty CM type.
For the initial, subsequently replaced level, the literal choices
$\{31,61,67,71\}$, auxiliary pair $(67,71)$ and distinguished
pair $(31,61)$ satisfy the prime, distinctness, modulo-$15$, and
$Q_S$-coprimality clauses. Set the prescribed finite thick set
to $TT$, the second set to $\{0\}$, the unipotent translate to
$0$, $L_w=2$ and $CT_2=0$. The second-set range holds because
the distinguished cyclic generator has positive order, and its
no-relation assertion has no pair of distinct elements to check.
This constructs an actual \texttt{Host} with the same cyclic
labels and with $q,d>N$; it does not assume an inhabitant of that
record. We use $B=0$ and $TT=\{0\}$ for (IG.13).

\paragraph{IG5: separated auxiliary primes with their integral norm factors.}
Fix the host of IG4 and write $c=Q_Sqd>1$. Choose
$\pi_0\in\mathcal O_E\cap K$ by IG3 with
$\pi_0\bar\pi_0=g_0>N$ prime and $\pi_0\equiv1\pmod c$.
To construct the second prime, define
\[
 P=g_0\prod_{1\le d\le r}(g_0^d-1),\quad
       M=(cP)^{2r+1},\quad
       N'=M^8+1+2M+P.                                 \tag{IG.17}
\]
Both $P$ and $M$ are positive, and $c\mid M$.
Apply IG3 to residue $1+M$ modulo $M^2$ and a bound larger
than $N'+g_0+N$. It gives
\[
 \pi_1\bar\pi_1=g_1,\qquad
 \pi_1\equiv1+M\pmod {M^2},\qquad
 g_1\equiv(1+M)^2\pmod {M^2}.                         \tag{IG.18}
\]
Thus $g_1\ne g_0$, both $g_i>N$, and
$\pi_i\equiv1\pmod c$, $g_i\equiv1\pmod c$.

We check the stronger separation of the actual bad-prime sets.
The source defines
\[
 \mathrm{BB}(r,g)=\{g\}\cup
 \{\ell:\ell\text{ prime},\ \exists\,1\le d\le r,
                   g^d<\ell^{4v_\ell(g^d-1)}\}.       \tag{IG.19}
\]
Every member $\ell$ of $\mathrm{BB}(r,g_0)$ divides $P$.
Put $a=2r+1$; then $\ell^a\mid M$ and $2d<\ell^a$ for
$1\le d\le r$. From the bound and congruence in (IG.18) write
$g_1=1+2M+M^2b$, $b\ge0$, and set
$x=M(2+Mb)$ and $S=\sum_{i=0}^{d-1}(1+x)^i$.
The geometric-sum identity and reduction modulo $M$ give
\[
 g_1^d-1=M(2+Mb)S,
       \qquad (2+Mb)S\equiv2d\pmod M.                 \tag{IG.20}
\]
Since $0<2d<\ell^a$, $\ell^a$ does not divide $(2+Mb)S$.
Consequently
\[
 \ell^{4v_\ell(g_1^d-1)}
 \le (\ell^a)^4(\ell^{v_\ell(M)})^4
 \le M^8\le g_1\le g_1^d.                            \tag{IG.21}
\]
This excludes its exceptional-prime clause in (IG.19), and
$g_1>P$ excludes $\ell=g_1$. Hence the two \texttt{BB} sets
are disjoint, with the literal fourth-power valuation threshold.

Finally use IG3's rational progression conclusion to choose
$p_\delta>g_1$ and $p'_\delta>p_\delta$, both
$1\pmod {15c}$. Replace the level by
$\{g_0,g_1,p_\delta,p'_\delta\}$ and use these ordered pairs
as the auxiliary and distinguished primes. Congruence to $1$
modulo $c$ proves coprimality with $Q_S$; all four primes are
distinct as required. The field, Artin map, half-different,
thick set and rank labels are unchanged. This proves the actual
releveling conclusion and all the integral factor congruences.
For IG4's fixed rank choose $N=2fk+1$, which gives
$2fk<g_i-1$.

\paragraph{IG6: two distinct coded columns and a common spread bound.}
Choose the two constant Boolean words on $\{0,\ldots,k-3\}$,
one true and one false; $k\ge m+3$ makes this set nonempty.
Put $\Delta=1$, $h=t+3$, $u_F=4(k-2)$, $u_L=-1$ and $b_0=1$.
For either code $c$, define
\[
 \begin{aligned}
 v_i&=4(k-2-i) &&(0\le i<k-1),\\
 u_0&=4(k-2)+\tfrac12,\qquad u_{k-1}=-\tfrac12,\\
 u_{i+1}&=v_i-\tfrac12-c_i &&(0\le i<k-2),\\
 b_j&=1+2jh,\qquad e_j=b_j+h &&(0\le j<m).
 \end{aligned}                                                   \tag{IG.22}
\]
These formulas are independent of the CM embedding. Every
$u_i$ is a half-integer; every $v_i,b_j,e_j$ is integral.
At an interior position, $v_i-u_{i+1}$ is $1/2$ or $3/2$,
and $u_{i+1}-v_{i+1}$ is $7/2$ or $5/2$.
The endpoints have gaps $1/2$. Also
$u_{k-1}=-1/2>-b_0=-1$ and
$e_j-b_j=b_{j+1}-e_j=h$, which satisfies the source's
gap requirement because $2h\ge t+3$.
These are all inequality and integrality fields of
\texttt{HCParam}. The common $v,b,e$, first and last values
give the five \texttt{PreCh} compatibility equations.
The two columns are distinct because their value at $u_1$
differs by one. Their letters obey exactly
\[
       v_i-v_{i+1}=4\Delta,\qquad
       u_{i+1}=v_i-\tfrac12-\Delta c_i.                \tag{IG.23}
\]
For each $j$, the exponent $e_j=1+(2j+1)(t+3)$ is constant
over the CM type. AC2--AC3 show that its product on every
root of unity is one. This is the literal \texttt{MUQ}
condition in \texttt{PreX.e\_\allowbreak even}; that field's name does
not impose an additional parity assumption on $e_j$.

There are finitely many triples consisting of a word, an index
$0\le j\le m$, and a CM embedding. Let $S$ be the maximum of
the real spreads of their actual weights $\lambda_j$.
Choose an integer $a\ge1$ sufficiently large that
\[
                 S\le2a,\qquad 2fN_m+1<4a.           \tag{IG.24}
\]
For example an integer strictly above
$\max(1,S/2,(2fN_m+1)/4)$ works. This proves all spread
inequalities simultaneously, with one $a$ fixed before the
type prime is chosen. It is a bound on the actual finite weight
family, and does not replace its integral Weyl lattices.

\paragraph{IG7: a common type prime with an exact initial window.}
Apply IG3 with modulus
\[
                  M_T=Q_Sqd(g_0g_1)^2
\]
and residue $1$, choosing a prime $T$ larger than
$4a+p_\delta+p'_\delta$. Its integral factor $\pi_T\in K$
satisfies $\pi_T\bar\pi_T=T$ and
$\pi_T-1\in M_T\mathcal O_E$.
Use this same $T,\pi_T$ for both words. It is distinct from
both distinguished primes; it is $1$ modulo $Q_Sqd$ by (IG.12).
Set initially
\[
 d_u=1,\qquad D_*=4T^{fk},\qquad\Gamma_P=2.
                                                               \tag{IG.25}
\]
All are positive and $\Gamma_P>1$. Since $T>0$ and $fk>0$,
the real-power rules give exactly
\[
     (D_*/d_u)^{1/(fk)}\Gamma_P^{-2/(fk)}=T.
                                                               \tag{IG.26}
\]
Thus the entire two-sided window is
$T\le T\le(1+1/(fk))T$.
Together with primality, field membership, norm, congruence
and avoidance, these are all fifteen fields of \texttt{LevCh}.
Extend the existing level by the image of this constant type-prime
map; the host field and its labels remain unchanged.
The scalar in \texttt{Pre.ofCh} is the literal quotient
$c(u)=\pi_T(u_0)/\pi_T(u)=1$. In a nonconstant variant the
same quotient obeys
$c(u)\overline{c(u)}T_u=T_{u_0}$ by cancellation of its
nonzero denominator. No new rational norm equation is assumed.

\paragraph{IG8: construct the spaces and every field of the original PreX.}
Use the actual CM extension and involution of IG4, and the norm
factor $p_z=\pi_T$ of IG7. The coordinate construction of SE15
gives $X_j=E^{N_j}$ and $X'=E^{\operatorname{rank}Y_m}$ with
positive diagonal forms and zero-extension inclusions. In $Y_j$
put a positive sign on its first
$N_j+\lfloor(m-1-j)/2\rfloor$ coordinates and a negative
sign on the remainder. The hyperbolic pairs are exactly
$e_i+f_i,(e_i-f_i)/2$ in the matched positive/negative slots.
Their pairings are $0,0,1$, respectively. An extra negative
coordinate, when present, contains $p_z$ and has norm $-T$.
At the top the extra positive coordinate contains $p_z$ and
has norm $T$. All the asserted orthogonalities follow because
their coordinate supports are disjoint. Take
$W_j$ with signature $(k+j-1,m-j)$ and every $V_j$ positive;
their coordinate lattices are $\mathcal O_E$ in each slot.

There is also a finite-adelic isometry field in the actual
\texttt{Spaces} record. It must be supplied, not inferred from
an archimedean signature. AH11 gives the literal
$\mathrm{cmUnram}$ of this extension. DR1 proves that at every
finite base-field prime the integral CM algebra is split or
unramified, and its unit norm is onto. Choose a unit $y_v$
with $\bar y_vy_v=-1$ there. These choices are integral units
at every finite place and therefore form a finite idele $y$,
with $\bar yy=-1$. To change one diagonal sign multiply that
coordinate by $y$; to keep a sign multiply by $1$.
The inverse multiplier is $-\bar y$, and entrywise calculation
gives $\bar g^{\mathsf T}J_sg=J_t$ for any two of these sign
patterns. This supplies the literal
\texttt{baseSpace\_\allowbreak isometric\_\allowbreak graded} on the same finite CM
algebra. Together with SE15 it verifies all the remaining
space fields and \texttt{DiagSpaces}, including
$V\mathrm{neg}=\mathrm{false}$ for every index.

Construct $P=\texttt{Pre.ofCh}$ with these spaces and the two
columns of IG6. Its base word is the first code. Now define
$X:\texttt{PreX}(P)$ using precisely $\pi_0,\pi_1$ from IG5,
$\pi_T$ from IG7, and $a$ from (IG.24).
The six auxiliary/type factor membership, norm and congruence
fields are already proved. The distinguished primes are
$1$ modulo the host conductor $Q_Sqd$, so the cyclotomic
Frobenius formula AH9 gives complete splitting in $E$.
Thus each has exactly $[E:\mathbb Q]$ primes above it, which
is the exact \texttt{pde\_\allowbreak split} cardinality clause.
The field \texttt{ccK\_\allowbreak top} is IG4; the \texttt{MUQ} clause
and injectivity of the weight family are IG6.
Finally $2fk<g_i-1$ is IG5, $4a<T_u$ is IG7,
and the two remaining size/spread inequalities are (IG.24).
These exhaust the fields of \texttt{PreX}.
The word set has cardinal two, $TT=\{0\}$ is nonempty, and
(IG.23) gives the injective two-letter family. This reconstructs
\texttt{exists\_\allowbreak pre\_\allowbreak two\_\allowbreak words} from its underlying arithmetic
and coordinate geometry.

\paragraph{IG9: the exact mass-index expression is a positive integer.}
We next prove the integrality needed before resetting $d_u$.
For a prime $g$ let $U_k(\mathbb F_g)$ be the upper
unitriangular group and write
\[
 i_g=[\mathrm{GL}_k(\mathbb F_g):U_k(\mathbb F_g)]>0,
 \qquad |\mathrm{GL}_k(\mathbb F_g)|
                         =g^{k(k-1)/2}i_g.             \tag{IG.27}
\]
This is exactly the real quantity \texttt{idxColLoc}, since
there are $k(k-1)/2$ independently chosen strict-upper entries.

Let $L$ be a full positive definite Hermitian lattice of rank
$k$ over $\mathcal O_E$, and let $\Gamma$ be its unitary
automorphism group. It is finite by CR4's bounded-image argument.
Assume $g>[E:\mathbb Q]k+1$ and that a prime of $E$ over $g$
has residue field $\mathbb F_g$. Reduction of $L$ there has
dimension $k$: localize its torsion-free finite module at this
DVR, choose a basis, and reduce that basis. Reduction embeds
$\Gamma$ into $\mathrm{GL}_k(\mathbb F_g)$.
Indeed any finite subgroup of the principal congruence kernel
is a $g$-group. The proof is the binomial induction in CR4:
an element of order prime to $g$ congruent to $1$ is congruent
to $1$ at every depth and hence is $1$.
The exponent of this finite-order element is a unit in the local
DVR, so the induction works with arbitrary ramification; it does
not identify the rational prime with a local uniformizer.
But an element of order $g$ would have a primitive $g$th root
as an eigenvalue on the rational space of dimension
$[E:\mathbb Q]k$, forcing $g-1\le[E:\mathbb Q]k$.
Cauchy's theorem therefore excludes a nontrivial kernel,
and also proves $g\nmid|\Gamma|$.
Lagrange and (IG.27) now give $|\Gamma|\mid i_g$.

There are finitely many classes in the genus of $L$.
One direct justification uses the trace form: it is a positive
definite integral quadratic form of fixed rank and determinant.
Successive-minimum reduction, with its positive integral minimum,
gives an integral basis with bounded Gram entries; there are
finitely many such Gram matrices. For each one the matrices
giving multiplication by a fixed integral basis of $\mathcal O_E$
are bounded in the trace norm (their archimedean operator norms
are fixed), so only finitely many integral matrices occur.
The Hermitian form is recovered from the trace pairings with
these multiplications by the nondegeneracy of the field trace.
Thus the Hermitian, not merely quadratic, isometry classes are
finite. This is the finite-genus input to the source's mass.
Consequently
\[
       \operatorname{mass}(L)i_g
          =\sum_{[L']\in\operatorname{genus}(L)}
                         \frac{i_g}{|\operatorname{Aut}(L')|}
               \in\mathbb Z_{>0}.                    \tag{IG.28}
\]
The genus contains $L$ itself, so the sum is nonempty. This
uses the finite sum defining the definite mass, not an analytic
Siegel mass formula.

For a strictly decreasing half-integral tuple $u_1,\ldots,u_k$,
the source's dimension factor is
\[
              D(u)=\frac{\prod_{a<b}(u_a-u_b)}
                                   {\prod_{j=1}^{k-1}j!}
                       \in\mathbb Z_{>0}.             \tag{IG.29}
\]
Here is an integral proof. Subtract the common half to get
integers $z_i$, and use the integer-valued polynomials
$\binom{X}{j}$, $0\le j<k$. Their evaluation matrix at the
integers $z_i$ is integral, and its determinant is their
Vandermonde divided by $\prod_{j=0}^{k-1}j!$.
The sign is fixed by the decreasing order, so (IG.29) is
positive. This argument includes negative $z_i$, for which
the binomial polynomial still has integral values.

Apply (IG.28) to the standard positive form and standard lattice
over the actual $E$. Congruence $g_0\equiv1\pmod {Q_Sqd}$
gives a residue-degree-one prime by AH9, and
$[E:\mathbb Q]=2f$ makes IG5's size bound sufficient.
Thus, for every supplied HC column,
\[
 \begin{split}
 d^\mathrm{true}(u)
 &=\operatorname{mass}(\mathcal O_E^k)
                       i_{g_0}^{\,f}i_{g_1}^{\,f}
                        \prod_\tau D(u_\tau)\\
 &=(\operatorname{mass}(\mathcal O_E^k)i_{g_0})
                       i_{g_0}^{\,f-1}i_{g_1}^{\,f}
                        \prod_\tau D(u_\tau)
                         \in\mathbb Z_{>0}.           
 \end{split}
 \tag{IG.30}
\]
The first line is literally \texttt{duStarOf}. Since $f\ge1$,
every factor in the second line is a positive integer. This
proves \texttt{duStarOf\_\allowbreak nat} without replacing a dimension
by an integral-lattice equality.

\paragraph{IG10: the original sized preliminary input and its exact window.}
Take the actual two-word $P,X$ of IG8 and keep its base word
$u_0$. Let $d=d^\mathrm{true}(u_0)\in\mathbb Z_{>0}$ from
IG9. Restrict the word type to a singleton, keeping its old
base column, type prime, integral norm factor, host, spaces,
coordinate frames, and the same $a$. Retain $\Gamma_P$ and
set
\[
      d_u=d,\qquad D_* =T_{u_0}^{fk}\Gamma_P^2d.
                                                               \tag{IG.31}
\]
All quantities in the real-power calculation are positive.
Consequently
\[
 (D_*/d)^{1/(fk)}\Gamma_P^{-2/(fk)}
  =T_{u_0}\Gamma_P^{2/(fk)}\Gamma_P^{-2/(fk)}=T_{u_0}.
                                                               \tag{IG.32}
\]
Thus both original window inequalities hold, including their
unchanged upper factor $1+1/(fk)$. The inherited norm
quotient scalar and its identities are unchanged at the
base word. Restrict its letter function too; a function from
a singleton is injective, and the two equations (IG.23)
still hold. The fields of \texttt{PreX} independent of words
are copied; its type factors, $T$ bound and spreads are the
old fields at $u_0$. Injectivity of the weight family is now
automatic.

We have therefore constructed, with no supplied preliminary
input, an actual $P$ satisfying
\[
 \mathrm{HasLetters}(P),\quad
       \mathrm{PreX}(P)\ne\varnothing,\quad
       \mathrm{DiagSpaces}(P.\mathrm{sp}),\quad
       d_u=d^\mathrm{true}(u)\quad\text{for every word }u.
                                                               \tag{IG.33}
\]
This is the literal \texttt{JetsInhab.exists\_\allowbreak pre\_\allowbreak sized}
conclusion. The source intentionally reduces the two-word
family to a singleton here; (IG.33) does not claim two
distinct words with a shared newly normalized $D_*$.
The latter would require simultaneous, externally prescribed
windows and is a different statement.

\paragraph{IG11: orientation and both pins on this same preliminary datum.}
Let a prescribed \texttt{PiBar} choose a prime
$\mathfrak p_i$ of the imaginary quadratic subfield over
each $g_i$. The factor $\pi_i$ of IG5 and its conjugate
generate maximal ideals of norm $g_i$ in that subfield.
The chosen prime contains one of them. Replace $\pi_i$ by
itself or its conjugate so that
$(\bar\pi_i)=\mathfrak p_i$ there. For every prime $w$ of
$E$,
\[
     \bar\pi_i\in w\quad\Longleftrightarrow\quad
                         w|_K=\mathfrak p_i.           \tag{IG.34}
\]
This follows by contraction and maximality; it is the exact
orientation condition. The norm is unchanged and conjugation
preserves the congruence to one modulo the rational integer
$Q_Sqd$. Only these auxiliary factor choices change. All
type factors, weights, windows and coordinate frames remain
the ones in IG10.

The source scalar $c(u)=\pi_T(u_0)/\pi_T(u)$ satisfies
$c(u)\pi_T(u)=\pi_T(u_0)$, since its denominator has nonzero
prime norm. The coordinate inclusion $X'\hookrightarrow X_m$
takes the $a$th standard frame vector to the vector with the
same coordinate index; evaluate every coordinate to check the
identity. These are exactly \texttt{Pins.cof} and
\texttt{Pins.iX}. Thus IG10 supplies the previously retained
original input to SE15--SE16, and TL16 applies to any prescribed
orientation. No arbitrary rational frame or previously fixed
type-prime factor is identified with this choice.

The existence of a full Hecke-character frame is a separate
question: orientation here holds for every supplied
\texttt{PiBar} (in particular one coming from such a frame).
It does not manufacture the frame's $\mu_j,\eta,\chi_\Psi,
\delta,\delta'$ by renaming their defining record. The next-source
record identifies their actual supplier if that additional
unconditional join is required.

\paragraph{IG12: the exact chosen-volume empty-jet stage.}
Fix an actual base $B$, its columns $C$, and a compact-dual Haar
system $\mathrm{cd}$, constructed for example by SE11.
For a Hermitian space $W$, the source first defines
$\mathrm{VolIs}(W,x)$ using all Haar measures normalized by
the finite lattice stabilizer and the specified archimedean
compact-dual measures, and all their rational fundamental
domains. Its actual definition of \texttt{trueVol} is
\[
 v(W)=
 \begin{cases}
 \text{a chosen }x\text{ with }x>0\text{ and }\mathrm{VolIs}(W,x),
                    &\text{if such an }x\text{ exists},\\
 1,                 &\text{otherwise}.
 \end{cases}                                                   \tag{IG.35}
\]
Classical choice in the first branch retains its two stated
properties; the second branch is positive as well. Therefore
$v(W)>0$ for every $W$, without an assumption of finite
geometric volume. If a positive \texttt{VolIs} witness has
separately been proved, the first branch applies and
$\mathrm{VolIs}(W,v(W))$ follows. The positivity argument alone
does not give this extra conclusion.

The literal $\mathrm{baseVol}(B,\mathrm{cd},W)$ is (IG.35) for
$B$'s actual CM embedding family and the supplied same
$\mathrm{cd}$. Apply the model construction of SE12 with this
positive function and with both jet index types empty.
Its two adjacent-rank antecedents are then empty universal
statements. This gives an actual $J:\texttt{GJStage}(B,C)$
whose fields satisfy
\[
  J.\mathrm{cd}=\mathrm{cd},\qquad
  J.\mathrm{vol}=\mathrm{baseVol}(B,J.\mathrm{cd}),\qquad
  J.\mathrm{jet}=\mathrm{noJet}(B),
                                                               \tag{IG.36}
\]
and every prescribed \texttt{gat} has finite classes.
The volume equality uses the same Haar datum after substitution;
it is not merely equality with a constant positive model.
This proves the exact \texttt{GJInhab.exists\_\allowbreak noJet} endpoint.
It also gives \texttt{nonempty\_\allowbreak stage} after choosing SE11's
Haar system. For \texttt{exists\_\allowbreak oneJet} the identical argument
retains its two nonempty adjacent-class inputs and its specified
prime/rank situation. The chosen-volume construction supplies
neither those classes nor the geometric \texttt{VolTieG} theorem.

\paragraph{IG13: the exact recursive sizes before imposing their window.}
For an arbitrary supplied $\texttt{PreCh}$, write $d_P(u)$ for
the expression in (IG.30), and let $D_P$ be the source's
\texttt{DsP}:
\[
 D_P=\operatorname{mass}(\mathcal O_E^r)
        \bigl(j_{g_0}^{\,f}j_{g_1}^{\,f}\bigr)
        \prod_\tau d_\infty(\mathrm{hc}(u_0,\tau)),
 \qquad r=k+m-1.                                      \tag{IG.37}
\]
Here $j_g$ is the literal row-level index and
$d_\infty=\texttt{trueDsig}$ is the defined formal-degree
polynomial. The identity matrix belongs to the finite row-level
group, so its order and the order of $\mathrm{GL}_r(\mathbb F_g)$
are both positive. Their quotient $j_g$ is positive. The
concatenation
\[
       (v_0,\ldots,v_{k-2},-b_0,\ldots,-b_{m-1})
\]
is strictly decreasing: the interlacing gives decreasing $v$,
the two $b,e$ gaps give increasing $b$, and
$v_{k-2}>u_{k-1}>-b_0$ gives the intervening strict gap.
The source's $\lambda_S$ is exactly this ordered concatenation,
through the finite-index equality $r=(k-1)+m$. Its
Vandermonde divided by the positive factorial product is
$d_\infty$. Thus each factor in (IG.37) is positive, by the
finite-genus argument of IG9, and
\[
 D_P>0,\qquad
 (2/3)^f\operatorname{mass}(\mathcal O_E^r)
       j_{g_0}^{\,f}j_{g_1}^{\,f}
       \prod_\tau D(v_\tau,-b_\tau)\ \le D_P.           \tag{IG.38}
\]
In fact the product following $(2/3)^f$ equals $D_P$ exactly.
This proves \texttt{DsP\_\allowbreak pos\_\allowbreak and\_\allowbreak lower}; it uses the defined
formal-degree polynomial, and asserts no realization of a
discrete series or its analytic formal-degree theorem.

The integer assertion for $d_P(u)$ is also uniform over the
large-scale inputs occurring in the recursive arrow. Indeed
the already established rank and exponential-versus-polynomial
bounds give, at one threshold depending only on
$0<\epsilon<1/52$,
\[
       f\le x^3,\qquad b\ge x^4,\qquad x\ge10.
\]
If the supplied \texttt{ChOK} has $g_i\ge(kb)^8$, then
$g_i\ge kb\ge10kx^3>2fk+1$, using $k\ge1$.
Its arithmetic clauses give $g_i\equiv1\pmod {Q_Sqd}$.
Apply IG9 separately to every actual word. This gives positive
integers $d_u$ with $d_u=d_P(u)$, at a threshold independent of
the word, $c_D,c_h$, coefficient field or later level choice.
This is the literal \texttt{exists\_\allowbreak du\_\allowbreak nat} conclusion.
The separate fixed-rank existence construction of IG10 is not
being used to replace these supplied large-scale inputs.

\paragraph{IG14: the recursive-existence join and its exact remaining inputs.}
We reconstruct the quantifier and same-object step in
\texttt{recExP\_\allowbreak of\_\allowbreak recAtG}. Its original proof invokes the
following quantitative theorem \texttt{window\_\allowbreak lower}; in
this component that theorem remains an explicit input:
for fixed $\epsilon,C,c_D\ge0$ in its stated range, once $c_h$
and then $x$ exceed its thresholds, every supplied
\texttt{ChOK} and the exact choices $d_u=d_P(u)$,
$D_*=D_P$, $\Gamma=\mathrm{GamS}(C,s,R)$ satisfy
\[
 m\mathcal L\ \le\
  \log\left((D_*/d_u)^{1/(fk)}\Gamma^{-2/(fk)}\right)
                       \quad\text{for every }u.       \tag{IG.39}
\]
The source derives this from \texttt{SiegelMassRatioCM},
the local-index ratio and an averaged Vandermonde lower bound.
IG13 supplies its positivity and size-normalization inputs.
Neither IG7's freely chosen window nor (IG.38) proves (IG.39).
In particular, our reconstruction of the arrow below is
conditional on this remaining quantitative theorem; it is
not a claim to have removed it from the original source proof.

Suppose the supplied $\mathrm{RecAtG}$ gives, over every
sufficiently large good context $K$, an actual $D,Y$ with
\[
 \mathrm{Over}(D,K.P,K.Lv),\quad
 \mathrm{RecOK},\quad\mathrm{LedOK},\quad\mathrm{TWAt},
 \quad\mathrm{FromCoreAt}(Q',K,D,Y),\quad
 \mathrm{TupUsedAt},\quad\mathrm{TwAllAt},\quad
 \mathrm{MtStableAt}.                                  \tag{IG.40}
\]
Also suppose the supplied $\mathrm{EvP}$ asserts $Q_O(\kappa,\beta.D,\beta.O)$
for every sufficiently large actual $\beta:\mathrm{BPt}(K)$
whose core satisfies $Q'$, under $\kappa.\mathrm{Pos}$.
We explain all threshold choices, since changing their order
would change the theorem.

Take the minimum of the two positive upper bounds for $\epsilon$.
For fixed $\epsilon,C_8,C_5,c_6,C_b$, take the supremum of their
two transported lower-bound records; its components are maxima,
so a standard $C$ large for this record is large for each.
For this $C$, take
\[
 c_1=\max(c_1^{\mathrm{rec}},c_1^{Q},0).
\]
After $c_D\ge c_1$ is fixed, take the maximum of the two
$c_h$ thresholds and the threshold of (IG.39).
After $c_h$ is fixed, take the maximum of the four $x$ thresholds
for (IG.40), the $Q_O$ conclusion, (IG.39), and IG13's integral
$d_P$. If $\kappa.\mathrm{Pos}$ is false, use threshold $0$
for the second of these four: the asserted implication is
then vacuous. This conditional choice occurs only after all
components of $\kappa$ have been fixed.

Now let $s,P,a,b,c$ satisfy the resulting good-scale antecedents.
Choose $d_u=d_P(u)\in\mathbb Z_{>0}$ by IG13 and choose
$D_*=D_P>0$. Equation (IG.39) supplies the source's logarithmic
window with the literal constant $c_s=1$.
These choices precede the next universal quantifier.
For \emph{every} subsequently supplied
$Lv:\mathrm{LevCh}(P)$ with
$Lv.d=d$, $Lv.D_*=D_P$ and $Lv.\Gamma=\mathrm{GamS}$, form
\[
 K=(s,C,c_D,c_h,P,q_1,a,b,c,\mathrm{Std},\mathrm{ChOK},
                                        Lv,d_P,D_P,\mathrm{GamS}).
                                                               \tag{IG.41}
\]
The last three entries mean the exact defining equalities of
the context; $q_1$ is the one in the supplied good datum.
Thus this is the actual context over the same $P,Lv$, with
no newly selected prime or word relabeling. Apply (IG.40)
to obtain its $D,Y$.

Fix any permitted $\mathrm{LedPart}$ of $Y$, with the supplied
scale, weight, numerical-value and canonical-instance equations.
The $\mathrm{FromCoreAt}$ field is a witness $X,f,T,\mathfrak c$
and an equality of dependent pairs
\[
             (D,\mathrm{layObj}(Y))
                     =(X.\mathrm{toD7},\mathfrak c.O).
                                                               \tag{IG.42}
\]
Transport first along this literal equality. The fixed
$\mathrm{LedPart}$, its scale/weight/value/canonical equations,
the good datum, $\mathrm{RecOK}$ and $\mathrm{LedOK}$, and the
two used/twist fields of (IG.40) now form exactly a
$\beta:\mathrm{BPt}(K)$. Its $Q'$ condition is the same core
condition in (IG.42). The supplied $\mathrm{EvP}$ therefore
gives $Q_O$ for its actual point and object whenever
$\kappa.\mathrm{Pos}$ holds.

The target also requires $\mathrm{QD}$ after the point's
\texttt{setDiv} operation. Here
$\mathrm{QD}=\mathrm{TupUsedAt}\wedge\mathrm{TwAllAt}
\wedge\mathrm{MtStableAt}$. Changing only the ramified
divisor lattices and their height parameters leaves the
tuple labels, their use, the coefficient module
$\widetilde M$, and its integral scalar action unchanged.
For the twist clause, keep the original conjugacy witness,
the integers $N_0$, and the same finite selection sets.
The definitions of the automorphic record and the twist
epsilon used in every $\mathrm{TwGoodOne}$ are unchanged;
the paired local membership statements transport by their
specified identity maps. Thus the same witnesses prove
the new $\mathrm{TwGoodOne}$ and $\mathrm{QD}$.
This is the actual \texttt{qd\_\allowbreak setDiv} implication, not a
fresh arithmetic choice of twist data.

We have proved the complete \texttt{RecExP} quantifier and
transport conclusion under (IG.39), (IG.40), and the stated
$\mathrm{EvP}$ premise, including its $\mathrm{withQD}$
conjunct. The original source arrow has only
$\mathrm{SiegelMassRatioCM}$ as its extra numerical premise
because it separately invokes \texttt{window\_\allowbreak lower}.
That quantitative deduction, and the supply of actual
recursive cores, remain explicit continuation targets here.

For precision the upstream \texttt{LayOfCore.lay} does more
than join bare auxiliary stages. It starts with an actual
$\mathrm{V4.Core}$ and $\mathrm{VolumeEqMass}_2$ on the same
$X.\mathrm{toD7}$. Its core specifies the true scalar
degrees, the real metric/volume ties, actual local row and
column laws, the common comparison maps, the ramified local
lattice bounds, and the selected arithmetic character data.
The two size equalities of the resulting $\mathrm{Lay}$
come from \texttt{sizes\_\allowbreak eq\_\allowbreak count}; they are not forced by
positive volume or by a bare finite-dimensional model.
Its $\mathrm{RecOK}$ has nine literal duties: the specified
$\Delta$, $\eta_F$, untwisted row vanishing, finite twist
support, positive $D_\Gamma$, fullness of the ramified
integral lattice, tame distinctness, the actual Minkowski
$D_\Gamma$ choice, and the two-sided type-prime local bounds.
Its $\mathrm{LedOK}$ additionally compares the
actual integral $\mathrm{NP}$ lattices for every admissible
canonical instance. These are retained at their actual
producer ranges. IG10--IG12 do not silently supply them.


\section{Rational theta images, level cuts, and the common target metric}\label{sec:theta-descent}

\status{This section expands the linear algebra behind the remaining theta comparison. It constructs the rational adjoint cuts, common row and lift matrices, and their exact metrics from explicit tests, equivariance and positive operator identities. It also proves the normalization that turns the source's finite period tests into the required rational tests. Existence of the actual convergent tests, their period formula, and the positive theta operator identities remains a separate analytic and arithmetic input. No common-image assertion is inferred from dimensions alone.}

Fix the literal coefficient field $K$, its distinguished embedding
$\sigma:K\hookrightarrow\mathbb C$, the source tuple set $T$ and base
$t_0$. All scalar extensions use this embedding unless another infinite
place is displayed. Inner products are conjugate-linear in their first
argument. The notation $j:F\to V$ denotes a $\sigma$-semilinear
coefficient form: some $K$-basis of $F$ maps to a complex basis of the
finite-dimensional Hermitian space $V$. It follows, by coordinates,
that $j$ is injective and its image spans $V$. This condition is stronger
than injectivity of $j$ alone.

\paragraph{TD1: finite separating tests recover rational coordinates.}
Let $f:\mathbb C^d\to E$ be injective and complex-linear, and let
$\psi_1,\ldots,\psi_N:E\to\mathbb C$ be complex-linear maps with
$\bigcap_a\ker\psi_a=0$. Suppose all
$\psi_a(f(e_i))$ lie in $\sigma(K)$. If $f(c)$ also has all its test
values in $\sigma(K)$, then $c\in\sigma(K)^d$.
Indeed choose $A\in M_{N,d}(K)$ and $b\in K^N$ with
\[
 \sigma(A)_{ai}=\psi_a(f(e_i)),\qquad
 \sigma(b)_a=\psi_a(f(c)).
\]
Linearity gives $\sigma(A)c=\sigma(b)$. If $Ax=0$ for $x\in K^d$,
then all tests annihilate $f(\sigma(x))$, so $x=0$. Choose a
$K$-linear left inverse $B$ of $A$, by extending a basis of its image.
Then $BA=1$ and
\[
 c=\sigma(B)\sigma(A)c=\sigma(Bb).
\]
This proves every coordinate assertion, including the empty-dimensional
case. It uses neither continuity of $\sigma$ nor any conjugacy of analytic
functions. The separating tests and their values on the \emph{same}
basis are indispensable. This is the exact mechanism of
\texttt{coord\_\allowbreak rat}.

\paragraph{TD2: the tested theta kernel descends over the fixed field.}
Let $\Theta:V\to E$ be complex-linear and suppose a separating test
family on $E$ has
$\psi_a(\Theta(jx))\in\sigma(K)$ for every $x\in F$.
In the chosen basis $b_i$ of $F$, write
$\psi_a(\Theta(jb_i))=\sigma(A_{ai})$. The common complex kernel of
the tests after $\Theta$ is $\ker\Theta$, so its coordinate equation is
$\sigma(A)c=0$. Gaussian elimination over $K$ gives a basis of
$\ker A$; the same pivot columns and equations remain valid after the
injective field map $\sigma$. Therefore
\[
 \ker\Theta
 =\operatorname{span}_{\mathbb C}
   \{jx:x\in F,\ \Theta(jx)=0\}.
\]
No uniform degree bound on the coefficients was needed. The statement
uses actual membership in $\sigma(K)$, rather than algebraicity over
$\mathbb Q$. The same elimination also proves preservation of rank,
intersection and finite sums under this scalar extension whenever the
maps or subspaces are defined over $K$.

\paragraph{TD3: the actual adjoint image is rational.}
Equip $E$ with its supplied positive Hermitian inner product.
Suppose in addition that $F$ has a $K$-valued sesquilinear form $h$,
with its given involution $*$, satisfying
\[
 \langle jx,jy\rangle=\sigma(h(x,y)).
\]
If $F\ne0$, choose $x$ with $jx\ne0$ and compare this identity for
$(ax,x)$ and $(x,x)$. Cancellation of $\|jx\|^2$ proves
$\sigma(a^*)=\overline{\sigma(a)}$ for every $a\in K$.
Thus star compatibility is proved from the supplied metric, not added as
an unrecorded premise. If $F=0$, the adjoint-image assertion is immediate.

The finite separating test family embeds $E$ into $\mathbb C^N$, so
$E$ is finite-dimensional. For each test $\psi_a$, let $e_a\in E$ be its Riesz vector:
$\psi_a(e)=\langle e_a,e\rangle$. These vectors span $E$, since a
vector orthogonal to all of them is killed by all tests. For
$y_a=\Theta^*e_a$, the Gram tests on $V$ give
\[
 \langle jb_i,y_a\rangle
 =\langle\Theta(jb_i),e_a\rangle
 =\overline{\psi_a(\Theta(jb_i))}\in\sigma(K).
\]
The matrix $G_{ik}=h(b_i,b_k)$ is invertible: its embedded matrix is the
positive Gram matrix of the complex basis $jb_i$. Solving the displayed
linear equations with $G^{-1}$ shows $y_a=jx_a$ for some $x_a\in F$.
Every such $y_a$ belongs to $\operatorname{im}\Theta^*$; their complex
span is the whole image. Consequently, for the actual preimage
$Z=j^{-1}(\operatorname{im}\Theta^*)$,
\[
 \operatorname{span}_{\mathbb C}j(Z)=\operatorname{im}\Theta^*.
 \tag{TD.1}
\]
This proves \texttt{range\_\allowbreak adjoint\_\allowbreak rat}. Applying the argument to each
literal $\Theta_t$ and each single lift $\theta'_{j,t}$ gives both
clauses of \texttt{RatZ} from \texttt{RowTest}, \texttt{LiftTest},
the coefficient-basis conditions and their metric identities.

\paragraph{TD4: adjoint closure makes the cut stable.}
For every linear map between finite-dimensional Hermitian spaces,
\[
 \operatorname{im}\Theta^*=(\ker\Theta)^\perp.
 \tag{TD.2}
\]
One inclusion follows from the adjoint identity; the other follows by
orthogonal decomposition, or by rank--nullity. In particular,
$\Theta$ is injective on $\operatorname{im}\Theta^*$: if
$y=\Theta^*z$ and $\Theta y=0$, then
$\|y\|^2=\langle z,\Theta y\rangle=0$.
Suppose an operator $A$ has adjoint $A'$ and $A'$ preserves
$\ker\Theta$. For $y\in\operatorname{im}\Theta^*$ and
$k\in\ker\Theta$,
$\langle k,Ay\rangle=\langle A'k,y\rangle=0$.
Hence $Ay\in\operatorname{im}\Theta^*$.

An operator on $F$ extends uniquely to $V$ by its values on $jb_i$;
expanding $x$ in the basis proves that extension sends $jx$ to $j(Ax)$.
Thus, when the source operator family is closed under adjoints and
preserves $\ker\Theta$, its rational preimage $Z$ is stable under every
operator. Applying this separately to the row and lift families proves
\texttt{cutStable}. An arbitrary invariant kernel without adjoint
closure would not suffice for this step.

\paragraph{TD5: finite group actions and their adjoints extend exactly.}
Let $G$ be finite, $r$ a $K$-representation on $F$, $\rho(g)$ the
complex extensions on $V$, and $\tau$ a unitary representation on $E$.
Assume the rational metric is $r(G)$-invariant and
\[
 \Theta(j(r(g)x))=\tau(g)\Theta(jx).
\]
Because $j(F)$ spans, these equations prove
$\rho(gh)=\rho(g)\rho(h)$, $\rho(1)=1$,
$\Theta\rho(g)=\tau(g)\Theta$, and unitarity of $\rho(g)$.
For unitarity, first extend the inner-product equation in its second
variable by complex linearity, then in its first by conjugate linearity.
Taking adjoints and using the inverse group element yields
\[
 \rho(g)\Theta^*=\Theta^*\tau(g).
 \tag{TD.3}
\]
For $a=\sum_g a_g g\in K[G]$, define
$r(a)=\sum_g a_g r(g)$ and
$\rho_\sigma(a)=\sum_g\sigma(a_g)\rho(g)$, and likewise for $\tau$.
Finite sums give
\[
 j r(a)=\rho_\sigma(a)j,\quad
 \Theta\rho_\sigma(a)=\tau_\sigma(a)\Theta,\quad
 \rho_\sigma(a)\Theta^*=\Theta^*\tau_\sigma(a).
 \tag{TD.4}
\]
The last formula is extended from (TD.3); it is \emph{not} obtained by
mistakenly taking an adjoint without conjugating the coefficients.
These are exactly the five \texttt{LevelOperators.Hyp} fields and their
finite-group-algebra consequences.

\paragraph{TD6: rationality survives an actual idempotent level cut.}
Let $e\in K[G]$ act idempotently on $F$. Then $P=\rho_\sigma(e)$
is idempotent on $V$ and
$\operatorname{im}P=\operatorname{span}_{\mathbb C}j(\operatorname{im}r(e))$.
Put $R=\operatorname{im}\Theta^*$ and retain (TD.1).
Equation (TD.4) shows $P(R)\subset R$. If $v\in R\cap\operatorname{im}P$,
write $v=\sum_i c_i jz_i$ with $z_i\in Z$. Since $Pv=v$,
\[
 v=\sum_i c_i j(r(e)z_i),\qquad
 r(e)z_i\in Z\cap\operatorname{im}r(e).
\]
The reverse inclusion is immediate, proving the exact \texttt{RatLev}
intersection formula.

Moreover every $u\in\operatorname{im}P$ has the same theta image as
some $u'\in R\cap\operatorname{im}P$. Decompose $u=p+k$ with
$p\in R$, $k\in\ker\Theta$, and set $u'=Pp$. Then
$Pu=u$, $Pk\in\ker\Theta$ by (TD.4), and
$\Theta u'=\Theta u$. This proves the actual \texttt{Hyp.dn}
statement. Self-adjointness of $P$ is unnecessary; its idempotence,
coefficient-field definition and the two invariant subspaces are the
properties used.

\paragraph{TD7: the level character and its inverse are kept distinct.}
For a subgroup $H\le G$ and a character $\chi:H\to K^\times$, the
level projector is
\[
 e_{H,\chi}=\frac1{|H|}\sum_{h\in H}\chi(h)^{-1}h.
\]
Its square is itself: the coefficient of $k\in H$ in the square is
$|H|^{-2}\sum_h\chi(h)^{-1}\chi(h^{-1}k)^{-1}
=|H|^{-1}\chi(k)^{-1}$. Reindexing also proves
$h_0 e_{H,\chi}=\chi(h_0)e_{H,\chi}$.
A vector has character $\sigma\chi$ precisely when this projector
fixes it. If $e_X$ is the actual central idempotent for the supplied
finite set of absolutely irreducible characters, then
$e_Xe_{H,\chi}$ is idempotent and its image is the intersection of the
selected block and character subspaces. Here the finite-character
idempotent assertion is the usual finite-group character orthogonality
statement; the present argument uses its explicit centrality and
idempotence, already present in the source's \texttt{goodAt} hypotheses.

If $v\in E$ is fixed by $\tau_\sigma(e_X)$ and has character
$\sigma\chi$, (TD.4) implies
$\Theta^*v\in\operatorname{im}\rho_\sigma(e_Xe_{H,\chi})$.
This is \texttt{Hyp.up}. A character equation on rational vectors
extends to their complex span and then through $\Theta$, which proves
\texttt{Hyp.type}. The literal definitions use $\chi=\psi$ on the row
side and $\chi=\psi^{-1}$ on \texttt{lvlA}, the lift side. These two
choices are not interchanged during transport. The principal level
$H=\{1\}$ gives $e_{H,1}=1$ exactly.

\paragraph{TD8: positive theta operators give the common image and its metric.}
For each tuple $t$, let $S_t\subset V_t$ be the actual complex level
space and let
$U_t:\mathbb C^{d_t}\to V_t$ be the extension of the supplied rational
row coordinates. Assume, as in \texttt{CutRows},
\[
 \operatorname{im}U_t=\operatorname{im}\Theta_t^*\cap S_t,
 \quad \Theta_t S_t\subset E_0,
 \quad \Theta_t^*E_0\subset S_t,
 \tag{TD.5}
\]
and $U_t$ is injective with Gram matrix $n_P H_{Q,t}$, where $n_P>0$
and $H_{Q,t}$ is the actual positive quotient metric. Suppose the fixed
finite orthogonal family $E_p$ spans $E_0$ and
\[
 \Theta_t\Theta_t^*e=c_{t,p}e\quad(e\in E_p),\qquad
 c_{t,p}=\kappa_{v_0}\prod_j\lambda_{j,t,p,v_0}>0.
 \tag{TD.6}
\]
The map $I_t=\Theta_t U_t$ is injective by TD4. It takes values in
$E_0$. Conversely, for $e\in E_p$, (TD.5) puts $\Theta_t^*e$ in
$\operatorname{im}U_t$. If $U_t a=\Theta_t^*e$, then
$I_t(c_{t,p}^{-1}a)=e$. Since the $E_p$ span, every $I_t$ has image
\emph{exactly} $E_0$. This proves all clauses of \texttt{H2} and
$d_t=\dim_{\mathbb C}E_0$ for every $t$.

For the metric, let $x=U_t a$ and suppose $\Theta_t x\in E_p$.
Both $\Theta_t^*\Theta_t x$ and $c_{t,p}x$ lie in
$\operatorname{im}\Theta_t^*$, and (TD.6) gives the same theta image
for them. TD4 implies equality. Consequently, for every $b$,
\[
 \langle I_t a,I_t b\rangle
   =n_P c_{t,p}\langle a,b\rangle_{Q,t}.
 \tag{TD.7}
\]
The scalar is real, so conjugate linearity causes no extra phase.
This proves the literal \texttt{H4} identity, not only a norm bound.
If all row spaces are zero the argument remains valid; no positive
row dimension is deduced. In the actual implementation (TD.5) is
\texttt{RatLev}, \texttt{DnLev}, \texttt{UpLev}, and (TD.6) is
\texttt{ThmBop} with \texttt{PartsE}. Their primitive theta-theoretic
suppliers remain explicit.

\paragraph{TD9: the common row identifications are rational and coherent.}
Retain TD8 and one \emph{common} \texttt{RowTest} on the target $E$,
with rational values on every $\Theta_t(j_t x)$. The columns of $I_t$
therefore have rational test values. Apply TD1 with $f=I_{t_0}$ to the
unique preimage of each column of $I_t$. The resulting matrix
$M_t\in M_{d_0,d_t}(K)$ satisfies
\[
 I_{t_0}\sigma(M_t)=I_t. \tag{TD.8}
\]
It is exactly the source's \texttt{rowMat}: the \texttt{solve} choice
is unique because $I_{t_0}$ is injective, and the entrywise
\texttt{kOf} fallback is never used after rationality is proved.
The common image and injectivity make $\sigma(M_t)$ bijective;
rank invariance from TD2 makes $M_t$ bijective over $K$.
At the base tuple, uniqueness gives $M_{t_0}=1$.

For any two tuples the matrix $M_s^{-1}M_t$ is the unique rational
comparison through the same target. Thus
$(M_r^{-1}M_s)(M_s^{-1}M_t)=M_r^{-1}M_t$ exactly; no choice of a new
isomorphism is involved. In particular the rational common-image
hypothesis in PL6 is now derived from the specified tests and the
operator/type inputs of TD8. Equal dimensions alone would not give
(TD.8), and unrelated rational test systems for different tuples
would not prove it over this fixed $K$.

\paragraph{TD10: the single lifts have the same exact descent.}
For a fixed step $j$, replace $\Theta_t$ by $\theta'_{j,t}$,
$S_t$ by the actual lift level, and $E_0$ by
$\operatorname{im}w_{0,j}$. The \texttt{CutLifts} assumptions give
\[
 \theta'_{j,t}(\theta'_{j,t})^*w=\lambda_{j,t,p,v_0}w
 \quad\text{on }W_{j,p}\cap\operatorname{im}w_{0,j},
\]
with positive scalar, and these subspaces span that image. The proof of
TD8 gives injectivity of $T_{j,t}$, equality
$\operatorname{im}T_{j,t}=\operatorname{im}w_{0,j}$, and
\[
 \langle T_{j,t}a,T_{j,t}b\rangle
   =n_P\lambda_{j,t,p,v_0}\langle a,b\rangle_{T,j,t}
\]
when $T_{j,t}a$ lies in the $p$-part. These are \texttt{H5} and
\texttt{ThmA}. The principal map $w_{0,j}=e_{W,j}\operatorname{lin}(b_{W,j})$
is injective: $b_{W,j}$ is a complex basis and
$\langle e_{W,j}x,e_{W,j}y\rangle=n_P\langle x,y\rangle$ with $n_P>0$.

The \texttt{LiftTestB} premise supplies one separating test family
whose values are rational both on all $T_{j,t}$ columns and on the
\emph{chosen principal basis} $w_{0,j}e_i$. TD1 now proves that the
actual \texttt{liftMat} $L_{j,t}$ is over $K$, with
\[
 w_{0,j}\sigma(L_{j,t})=T_{j,t},
\]
and that it is bijective. The extra principal-basis clause in
\texttt{LiftTestB} cannot be deleted. This constructs the matrix; it
does not identify its image lattice integrally without the separate
lattice statements.

\paragraph{TD11: the base-lift computation needs only image containment.}
The source's theorem
\texttt{baseLevelLiftMatrix\_\allowbreak solves\_\allowbreak relation\_\allowbreak extendedData}
uses a weaker local argument than a full common-image theorem. For the
same $R,F,j$ at the base tuple, it suffices that each column
$T_{j,t_0}e_k$ belongs to $\operatorname{im}w_{0,j}$, that $w_{0,j}$
is injective, and that the \texttt{LiftTestB} values on both column
families are rational. TD1 then proves the chosen solution matrix
rational, and linearity gives
\[
 w_{0,j}\sigma(L_{j,t_0}x)=T_{j,t_0}\sigma(x)
 \quad(x\in K^{d_{j,t_0}}). \tag{TD.9}
\]
No surjectivity of $T_{j,t_0}$ is used here.

For completeness, the source's \texttt{tC\_\allowbreak single\_\allowbreak mem} obtains that
containment as follows. The lower-level source is cut by
$e_Xe_{H,\psi^{-1}}$, whose image lies in that of $e_X$. TD6 replaces
its vectors, without changing the lift, by vectors in the rational
adjoint image at principal level. Those coordinates lie in the
principal forms span, defined as the span of the corresponding principal
images over all tuples. TD7 supplies their lower-level character.
The supplied \texttt{PLevW} identity identifies the intersection of
these two conditions with $\operatorname{im}w_{0,j}$. This is precisely
the \texttt{DnLevL} containment used in the source proof.

Thus \texttt{WeakCharacterisation.base.models} and the specified
\texttt{HoldsL} fields supply (TD.9) on the original realization. The
large outer statement also lists published representation-theory and
base-enlargement inputs, but its local proof does not use them after
these model fields and \texttt{HoldsL} are supplied. This observation
removes an unused detour from the human proof; it neither constructs
\texttt{HoldsL} nor asserts its analytic prerequisites.

\paragraph{TD12: distinct spectral parts are orthogonal at every place.}
Here are the exact linear-algebra requirements behind
\texttt{majAt\_\allowbreak eq\_\allowbreak zero\_\allowbreak of\_\allowbreak mem\_\allowbreak Qp\_\allowbreak of\_\allowbreak ne}. At an infinite
place $v$, use the actual positive matrix $H_v$ on the finite class
space. The rational Hecke and inertia operators commute, are closed
under $H_v$-adjoints, and commute with the finite level action. The
source obtains these assertions by transporting its \texttt{OkH} and
\texttt{InertRationalAdjoint} matrix identities. Their entries, and the
original metric entries, lie in the specified maximal CM subfield;
conjugate transpose therefore transports correctly to every $v$.
This is a place-by-place algebraic statement, not continuity of an
arbitrary field embedding.

Every operator in the commuting adjoint-closed family is normal. A
normal operator has orthogonal eigenspaces: if $Ax=\alpha x$, then
$A^*x=\bar\alpha x$, since the squared norms of
$(A-\alpha)x$ and $(A^*-\bar\alpha)x$ agree. If also $Ay=\beta y$,
the adjoint identity gives
$(\beta-\alpha)\langle x,y\rangle=0$.
Existence of a complex eigenvalue and induction on the invariant
orthogonal complement give diagonalizability. Repeating on joint
eigenspaces gives simultaneous diagonalization; an infinite indexing
set causes no problem, since a finite subset spans its finite-dimensional
operator algebra. Hence the generalized eigenspaces occurring in the
source equal the ordinary eigenspaces.

The actual invariant level and adjoint-image subspaces are stable
under the whole family and its adjoints, so their restrictions remain
normal. Disjoint conditions on the joint eigenvalue systems therefore
give orthogonal sums at $v$. The source's residue labels $p\ne q$
have disjoint systems by \texttt{diffAt\_\allowbreak of\_\allowbreak mem\_\allowbreak ixFin3} and
\texttt{excl3}. The passage to the source's quotient metric needs a further check.
Let $P$ denote the actual surjective coordinate map $Pq$, with $PzE=1$.
The identity $PA=\bar A P$ for every operator and its adjoint makes
$\ker P$ invariant under the whole adjoint-closed family. Its orthogonal
complement is therefore invariant too. The quotient metric is exactly
the metric of the unique orthogonal lift: in matrix coordinates it is
$(PH_v^{-1}P^*)^{-1}$, and the lift of $y$ is
$H_v^{-1}P^*(PH_v^{-1}P^*)^{-1}y$. This follows by direct multiplication
and the condition of orthogonality to $\ker P$. Orthogonal projection
onto $(\ker P)^\perp$ commutes with the operators, so it preserves every
joint eigenspace. The supplied $zE x$ and $zE y$ lie in their respective
spectral sums; replacing them by their orthogonal lifts preserves those
sums. Their quotient inner product consequently vanishes. Thus
$\langle x,y\rangle_{Q,v}=0$ for $x\in Q_p$, $y\in Q_q$.
This proves \texttt{PartsOrth} for all places (and thus for the
source's required places $v\ne v_0$), with the stated rational-adjoint
and label-separation suppliers retained.

\paragraph{TD13: construct the exact weighted target metric.}
Let $Q=\bigoplus_p Q_p$ be the actual rational direct sum and let
$P_p$ be its coordinate projections. Directness gives
$P_pP_q=0$ for $p\ne q$ and $\sum_pP_p=1$. At $v\ne v_0$ set
\[
 c_{p,v}=\kappa_v\prod_j\lambda_{j,t_0,p,v}>0,\qquad
 H_{E,v}=\sum_p c_{p,v}\,v(P_p)^*H_{Q,v}v(P_p).
 \tag{TD.10}
\]
The matrix is positive definite: its quadratic form is a sum of
nonnegative terms, and if all vanish, each $v(P_p)z=0$, whence $z=0$.
Choose the positive square root and frame $X_{E,v}=H_{E,v}^{-1/2}$;
then $(X_{E,v}^{-1})^*X_{E,v}^{-1}=H_{E,v}$. This is the metric
specified by the source's \texttt{weightFrame}; uniqueness of a
particular matrix square-root choice is unnecessary for equality of
its norm.

At the distinguished place set
\[
 H_{E,v_0}=n_P^{-1}\operatorname{Gram}(I_{t_0}e_i),
 \tag{TD.11}
\]
which is positive definite by TD8. Splice this frame into (TD.10)
exactly at $v_0$, as \texttt{rowXE} does. TD7's level conventions and
TD8's factor $n_P$ give (TD.7) at the base; at other places TD12 and
$\sum_qP_q=1$ give
\[
 \langle x,y\rangle_{E,v}=c_{p,v}\langle x,y\rangle_{Q,v}
 \quad(x\in Q_p).
 \tag{TD.12}
\]
If the finite part index is nonempty, orthogonality implies the exact
lower bound needed by PL11:
\[
 \|z\|_{E,v}^2\ge
  \left(\kappa_v\min_p\prod_j\lambda_{j,t_0,p,v}\right)
  \|z\|_{Q,v}^2.
 \tag{TD.13}
\]
The corresponding base-place bound uses the actual spectral-part
coordinate identification (\texttt{H3}) and (TD.7). If the index set is
empty its spanning condition forces $Q=0$; do not assign a minimum to
an empty set. No quantitative lower bound on the positive
$\lambda$'s, or bound on their logarithms, has been proved here.

\paragraph{TD14: these matrices give the literal principal lattice.}
Let $M_t$ be the source's full projected row lattice and let
$\mu(w)=t_0$ outside the finite twist set. With the matrices of TD9,
put $L_w=M_{\mu(w)}^{\rm cmp}M_{\mu(w)}$, where the superscript
$\mathrm{cmp}$ distinguishes the comparison matrix $M_t^{\rm cmp}$
from the lattice. Equivalently $L_w$ is the actual
\texttt{matLat(rowMat(mu(w)),abelian2(mu(w)))}.
The comparison matrices are bijective, so these are full lattices;
$M_{t_0}^{\rm cmp}=1$ makes them equal to the base lattice off the
finite twist set. PL9 therefore constructs exactly
\[
 R=\bigcap_w(L_w)_w,\qquad R_w=(L_w)_w.
\]
For any $y\in K^{d_0}$, (TD.8) and injectivity give the stronger
identification
\[
 y\in M_t^{\rm cmp}M_t
 \quad\Longleftrightarrow\quad
 I_{t_0}\sigma(y)\in\{I_t\sigma(z):z\in M_t\}.
\]
Membership in $R$ is thus precisely the source's condition that, for
every $w$, some $d\in\mathcal O_K\setminus w$ makes
$d y$ belong to this image lattice. The same target carries the frames
of TD13, hence is the literal \texttt{principalHermitianLattice} with
its \texttt{Good} fields. This step does not infer local divisibility
of periods or an integral isometry from mere rational bijectivity.
It supplies the exact object to which PL9--PL11 apply.

\paragraph{TD15: one normalization makes all row tests rational.}
The source's row-test construction starts with complex-linear tests
$p_a$ on its Hilbert model $E$, separating all of $E$. The integral
definitions become linear only after their stated convergence
hypothesis \texttt{LinK}; retain that hypothesis. Let the actual finite
period identity be
\[
 \sigma(z_t)c_t\,p_a(\Theta_t j_t x)
       =\kappa\,w_{t,x,a},
 \tag{TD.14}
\]
for vectors in the selected good block, with $z_t\in K$ a root of
unity, $c_t\ne0$, $w_{t,x,a}\in\sigma(K)$ and
$c_t/c_{t_0}\in\sigma(K)$. The same identity extends to all source
vectors after replacing $x$ by its actual central-block projection:
TD5 and the identity action of $e_X$ on the Hilbert target show that
this replacement preserves $\Theta_tj_tx$. The rational right side
may change to $w_{t,x',a}$; no equality of the two finite sums is
needed.

If $\kappa\ne0$, use the \emph{single} scalar
$k=c_{t_0}\kappa^{-1}$. Writing
$w_{t,x',a}=\sigma(b)$ and $c_t/c_{t_0}=\sigma(r)$, (TD.14) gives
\[
 k\,p_a(\Theta_tj_tx)=\sigma\!\left(\frac{b}{z_t r}\right).
 \tag{TD.15}
\]
The denominators are nonzero because $z_t$ is a root of unity and
both $c_t,c_{t_0}$ are nonzero. Multiplying every $p_a$ by $k$ retains
separation. If $\kappa=0$, the nonzero left prefactor in (TD.14) forces
all $p_a(\Theta_tj_tx)=0$; choose $k=1$. These tests are still
separating on $E$ and have rational values on every theta image,
which in this branch is zero. Thus no unjustified nonvanishing premise
for the global scalar is introduced.

In the literal source $c_t=\prod_{i,j}c_{i,j,t}$ and
$c_t/c_{t_0}=\prod_{i,j}\texttt{sqRatio}_{i,j,t}$.
\texttt{ConstK} supplies each nonzero factor, \texttt{RatioK} each
ratio in the same field, \texttt{Formula2K} (TD.14), and
\texttt{SumRatK} the rational right side. This proves
\texttt{rowTest\_\allowbreak of\_\allowbreak choice} from those exact premises. Finite sums
and products preserve the coefficient field; their convergence,
period formula and actual coefficient-field assertions are not proved
merely by this closure observation.

\paragraph{TD16: conjugation, uniformity and the remaining test boundary.}
The Hilbert model and the complex side may differ by the source's
Boolean conjugation twist. Write $c_s$ for either the identity or
complex conjugation, and give $\operatorname{Tw}_s E$ the scalar action
for which its map to $E$ sends $a\cdot x$ to $c_s(a)x$.
If $\sigma'=c_s\circ\sigma$, transport a test by
\[
 \psi^{(s)}(x)=c_s\bigl(\psi(\operatorname{out}x)\bigr).
\]
It is complex-linear, because the two applications of $c_s$ cancel
on scalars. The test values transform from $\sigma(K)$ to
$\sigma'(K)$ and separation is preserved. Applying the same operation
twice is the inverse. This proves \texttt{test\_\allowbreak tw\_\allowbreak iff} and the
literal \texttt{rowTest\_\allowbreak toCx3}, using $R.\texttt{emb0}$ for the
specified embeddings; no arbitrary automorphism is treated as
continuous.

Each global test input has the order
\[
 \exists C_{F,0}\in\mathbb R\quad
 \forall\varepsilon\in\mathbb R\quad \forall C_F\ge C_{F,0}\quad
 \exists x_0\in\mathbb R\quad
 \forall D,\mathrm{large}_2,R,F,\mathrm{ch},\ldots .
\]
The implication for these final variables retains
$\mathrm{large}_2.\mathrm{sc}.\varepsilon=\varepsilon$,
$x_0\le\mathrm{large}_2.\mathrm{sc}.x$, the exact
\texttt{HidaSetting}, $\mathrm{ch}.\mathrm{Is}$ and
$\mathrm{ch}.\mathrm{AtScale}(C_F,\mathrm{large}_2.\mathrm{sc}.b)$.
For finitely many inputs take their maximum $C_{F,0}$ first, then,
\emph{after} fixing $\varepsilon,C_F$, the maximum of their $x_0$'s.
TD15 and the conjugation transport then prove the conditional
\texttt{rowTest\_\allowbreak realisation} with the same $D,L,R,F,\mathrm{ch}$ and
all original setting predicates. The proof of this assembly does not
use \texttt{SumOKK} once \texttt{SumRatK} is supplied; the former is
retained by the literal outer theorem and may be needed to establish
the latter.

There is a sharp obstruction to dropping the common fixed-field input.
In $E=\mathbb C$ and $K=\mathbb Q$, no separating family of
complex-linear tests takes rational values on both $1$ and $i$:
a test has $\psi(i)=i\psi(1)$, so rationality of both values forces
$\psi(1)=0$; all tests would annihilate $1$. Both input numbers are
algebraic. Hence the algebraic vacuum and jet results of VI/CA alone
do not imply the prescribed \texttt{RowTest}. The exact remaining
suppliers are the compatible test choice and \texttt{FieldOK},
convergent separating functionals, the unfolded period formula and
rational normalization ratios, together with the actual theta
operator/type identities and their positive eigenvalues in TD8--TD10.
The spectral kernel bounds, dimension balance, congruence denominators
and global prime estimates remain independent frontiers.


\section{Arithmetic realization of the local level and weight comparisons}
\label{sec:arithmetic-realization}

The datum $D$ remains fixed throughout this section. Put $E=D.\mathrm{host}.E$,
$F=E^+$, $n=D.R.nV=r+t$, and $K_{i,k}=\operatorname{Frac}(O_{i,k})$.
The index $j$ always lies in $\operatorname{Fin}(D.R.m)$; its cell is at
position $j+1$. The group of finite adelic unitary points is $G_f$, and its
rational subgroup is $\Gamma$. We use multiplicative valuations, so
membership in the maximal ideal means valuation strictly less than $1$.
At a split place $s$, write $v=s.v$, $P=s.w$, and
$\operatorname{loc}_s:G_f\to\mathrm{GL}_n(F_v)$ for the selected component.
The selection matters: the conjugate component is its inverse transpose.

Fix an actual $\texttt{NExpTied}$ object $X$, an admissible twist
$(i,b,k,Q,\psi)$, and its actual cell $M=X.\mathrm{mtV}(i,j+1,b,k,Q,\psi)$.
The cell comes with its coefficient module $W$, representation $\rho$,
level $J$, character $\chi$, realization $\mathrm{real}$, and selected
prime sets $Q_P,D_P$. These are retained data, not constructions from a
dimension formula. For a supplied stage $(B,f,T,w,P_0,x,l)$, write
\[
 \eta=\mu_{0,j}\,\xi(\mu_{\mathrm{Tw}}(v_0))
      =\mu_{0,j}\,\delta_\ell\,\psi_{\mathrm H},
 \qquad
 \delta_\ell=\begin{cases}\delta'&\ell=p_\delta,\\
                          \delta&\ell\ne p_\delta.\end{cases}       \tag{AR.1}
\]
Here $\ell=\operatorname{resChar}(v_0)=D.\mathrm{ntw}.l(i)$, and the
second equality is the actual tuple identity, not an identification of
unrelated characters. The complex characters supplied by HF satisfy the
needed character laws when those HF hypotheses hold. They do not, by
themselves, supply normalized Galois-character compatibility or the
nonempty twist coverage used below.

\paragraph{AR1: exact ramification gives the required bad-place bound.}
Assume the actual $\texttt{IsMu0}$, $\texttt{IsDelta}$,
$\texttt{IsPsi}$ and $\texttt{IsTwist}$ laws in (AR.1).
Outside the two type primes $\mu_{0,j}$ is unramified. The twist character
$\psi_{\mathrm H}$ is ramified exactly at $Q$ and avoids the two auxiliary
rational primes. The selected auxiliary character is unramified outside
the primes over the opposite auxiliary prime, and its inertial image has
order $15$ at every such prime. Then
\[
 \mathrm{badV}(M)\subset g\mathrm{Pl}(D)\cup q\mathrm{Pl}(D,\eta).
                                                               \tag{AR.2}
\]
Moreover the extra finite-order factor $\delta_\ell\psi_{\mathrm H}$ is
unramified at each type prime.

If a prime ideal contains distinct rational primes $p,q$, a Bezout identity
$ap+bq=1$ puts $1$ in that ideal, which is impossible. In the source's
$\texttt{IsTwist}$ condition a ramified prime has prime absolute norm
outside the forbidden rational set. If it also lies above a forbidden
$p$, the norm and $p$ both belong to the prime ideal, so their primality
forces equality and contradicts that exclusion. This proves the claimed
unramifiedness at the auxiliary places, with the literal norm condition.

Now take a bad place and a prime of $E$ above it witnessing badness.
A type prime is already in $g\mathrm{Pl}$. At a prime of $Q$, both
$\mu_{0,j}$ and $\delta_\ell$ are unramified. If their product with
$\psi_{\mathrm H}$ were unramified, cancellation on local units would
make $\psi_{\mathrm H}$ unramified, contrary to its exact ramification.
At an auxiliary prime, $\mu_{0,j}$ and $\psi_{\mathrm H}$ are
unramified; cancellation would make $\delta_\ell$ unramified. Its
inertial image would then have order $1$, contradicting order $15$.
At a type prime, $Q$ is excluded by admissibility and the auxiliary
primes have different rational residue characteristics. Hence both
extra factors are unramified there. These arguments prove (AR.2)
without estimating, enlarging or replacing $Q$.

\paragraph{AR2: selected halves and split-place separation are exact.}
Suppose $P_q$ selects a prime over each prime of $F$. On split primes
assume that the series choice set is exactly $\operatorname{range}P_q$.
For a finite conjugation-stable set $S$ of split primes,
the series half $\operatorname{half}(\operatorname{range}P_q,S)$ equals
the cell half $\operatorname{qpOf}(P_q,S)$.

For a split prime $P$ above $v$, the only primes above $v$ are $P$ and
$\bar P$. Membership in the series half means $P\in S$ and
$P=P_q(v)$. The cell condition is $P\in S$ and either this equality
or exclusion of the other prime from $S$. The second alternative,
when $P_q(v)=\bar P$, is impossible because $S$ is conjugation-stable.
Conversely $P=P_q(v)$ immediately satisfies the cell condition.
This also explains why an arbitrary half cannot be substituted.

The actual sets $Q_P,D_P$ are disjoint, contain no conjugate pair, and
avoid type primes. If selected split primes in $Q_P\cup D_P$ have the
same underlying $v$, they are equal or conjugate. A conjugate pair
within one of the sets is excluded by its defining condition. A
conjugate pair between the sets would put a prime of $Q$ over the
forbidden auxiliary rational prime. A coincidence with a type place
would similarly put the corresponding $g_a$ in a conductor or
auxiliary prime. Thus, at a selected conductor place, it is the only
member of $Q_P\cup D_P$ above that place and no type prime lies above it.
Finally, if $v\in q\mathrm{Pl}(D,\eta)$, choose a prime above it where
$\eta$ is ramified. AR1 shows that it lies in $Q$ or the auxiliary set.
The covering fields of $Q_P,D_P$ select it or its conjugate, producing
the required member of $Q_P\cup D_P$ over precisely $v$.

\paragraph{AR3: tame residue values agree with the specified Artin map.}
Retain the source's published global Artin reciprocity, with its
inertia and inverse-arithmetic-Frobenius convention. Let $\varphi:G_E\to A$
be a homomorphism into a commutative group with open kernel. At a prime
$P$ of residue characteristic $p$, suppose every inertial value has
order dividing a power of an integer $m$ prime to $p$.
Its value on a tame representative depends only on the residue-unit
and the stated exponent. In particular,
\[
 u\equiv1\pmod P\ \Longrightarrow\
 \operatorname{tameVal}(\varphi,P,a,u)=1.                 \tag{AR.3}
\]
If a complex embedding reads $\varphi$ as a Hecke character $\chi_{\mathrm H}$
under the specified Artin map, then for $u\in\mathcal O_E\setminus P$,
\[
 e\!\left(\operatorname{tameVal}(\varphi,P,1,u)\right)
       =\chi_{\mathrm H,P}(u).                          \tag{AR.4}
\]
The notation in (AR.4) includes the literal inclusion of coefficient
units into the fraction field.

The open kernel contains the fixing subgroup of a finite Galois
extension. In a finite local Galois extension the first ramification
quotient acts on a uniformizer through a residue-field unit; each
successive quotient embeds in the additive residue field by the
leading coefficient of $\sigma(\pi)-\pi$. The latter quotients are
$p$-groups and the filtration terminates. Consequently the wild inertia
image in a finite quotient is a $p$-group. Its image under $\varphi$ is
also of $m$-power order, so it is trivial. Tame inertia is read by the
ratios $\sigma(\pi_N)/\pi_N$, for prime-to-$p$ roots of a uniformizer.
Frobenius conjugates its tame generator to its $q$th power, where
$q=\#k_P$. Since the target is commutative, its image has exponent
dividing $q-1$. Thus the ratio for a $(q-1)$st root already determines
the value. Changing the valuation above $P$ conjugates inertia and
does not change a commutative character. This proves independence of
the choices used in $\operatorname{tameVal}$ and proves (AR.3).
These are the usual finite ramification and Hensel facts for the
completed number field; no assertion that wild inertia is trivial
in the absolute Galois group is being made.

The sign in (AR.4) can be checked without a choice of a tame generator.
Choose $\varpi\in P\setminus P^2$, $\pi^{q-1}=\varpi$, and
$t^{q-1}=u$. Arithmetic Frobenius lifts fixing respectively $\pi$ and
$\pi t$ give an inertia element $\tau'^{-1}\tau$. Its ratio on $\pi$
reduces to $u$. Their images under the stated Artin map differ by the
local unit $u$, since the two uniformizer norms are
$(-1)^q\varpi$ and $(-1)^qu\varpi$. For an arbitrary local unit $x$
with the same residue, Hensel lifting gives
$x/u=z^{q-1}$ with $z\equiv1$ modulo the maximal ideal. Lift $z$ to
inertia using the inertia-image clause and multiply by that lift to
the power $q-1$. It changes the Artin image to that of $x$ and leaves
the tame ratio unchanged. Evaluation of the reading identity proves
(AR.4). Replacing the inertia element by its $a$th power gives the
stated exponent-$a$ version.

\paragraph{AR4: the auxiliary side is read with its actual inverse.}
Let $\epsilon_b=1$ for $b=\mathrm{true}$ and $-1$ otherwise.
Assume the literal $\texttt{DeltaRead}(D,cf)$ and
$\texttt{SideRead}(i,b,k,e)$. Then the complex character obtained by
reading $\delta_i^{\epsilon_b}$ is the Hecke character
$cf.\mathrm{delAt}(\ell)$. Hence at $P\in D.\mathrm{ntw}.pD(i)$,
\[
 e\!\left(\operatorname{tameVal}(\delta_i,P,\epsilon_b,u)\right)
       =cf.\mathrm{delAt}(\ell)_P(u).                   \tag{AR.5}
\]

The Hecke character has finite order. RC3 descends it through the
specified Artin quotient to a continuous complex Galois character.
The read character also has open kernel, so is continuous even though
the field embedding $e$ itself is not assumed continuous.
Outside finitely many primes, $\texttt{DeltaRead}$ supplies an element
$a\in L_0$ with
$\iota(a)=cf.\mathrm{delAt}(\ell)(P)^{-1}$ and the prescribed
coefficient equality for $\delta_i(\mathrm{Frob}_P)$.
On the true side, $\texttt{SideRead}$ reads this equality directly.
On the false side the coefficient involution sends $a^{-1}$ to $a$:
the complex values are unitary, so conjugation equals inversion.
The additional inverse in $\delta_i^{\epsilon_b}$ cancels this change.
Thus both sides give the same inverse Frobenius value. RC4 proves
equality of the characters.

The inertia image of $\delta_i$ has order $15$. The auxiliary norm is
$1$ modulo $15$, so that norm is coprime to $15$. Since the norm
is a power of the residue characteristic, that characteristic
does not divide $15$. Apply AR3 with $m=15$, preserving $\epsilon_b$.
On $Q_P$ the analogous application uses the $\ell$-power character
$\psi$ and $\ell\notin P$. AR1 cancels the unramified factors in
(AR.1), so these two formulas read exactly the local component of
$\eta$, not just the separate twist factor.

\paragraph{AR5: the conductor-place subgroup and character are parahoric.}
Let $P\in Q_P\cup D_P$, $s=M.\mathrm{sp}(P)$, and
$\kappa\in G_f$ be trivial outside $s.v$. Then
\[
 \kappa\in M.J
 \quad\Longleftrightarrow\quad
 \operatorname{loc}_s(\kappa)\in\operatorname{ParaRT}(r).
                                                               \tag{AR.6}
\]
Here $\operatorname{ParaRT}(r)$ consists of integral invertible matrices
whose lower-left $t$-by-$r$ block lies in the maximal ideal.
There is a local component $\xi:F_{s.v}^{\times}\to\mathbb C^\times$
of $\eta$ such that, for every allowed reading $e$,
\[
 e(\chi(\kappa))=\xi(\det D_\kappa),                     \tag{AR.7}
\]
where $D_\kappa$ is the lower-right block. This is the literal
$\texttt{LevelPara}$ assertion.

Unfold $M.J$: it imposes global maximal-compact membership, the type
unipotent-block conditions, and parahoric conditions at every prime
of $Q_P,D_P$. AR2 and the support condition make every other local
matrix the identity. The sole remaining condition is the one in
(AR.6). Global maximal-compact membership is equivalent to local
integrality here: the unselected conjugate component is the inverse
transpose, so both a matrix and its inverse are integral exactly when
the selected component is in $\mathrm{GL}_n(\mathcal O_v)$.

Because $\eta$ is tame, compose its $E_P$ component with the split
completion isomorphism $F_v\simeq E_P$ to obtain $\xi$. It is trivial
on $1+\mathfrak m_v$ and gives the required values on $E^\times$.
The block-triangular reduction of an integral invertible matrix has
invertible diagonal blocks, so $\det D_\kappa$ is a unit. Choose its
integral residue representative in $\mathcal O_E$.
The literal product formula for $\chi$ has type factors and one
factor at each member of $Q_P,D_P$. The type factors are $1$ at the
identity. Every other conductor factor has residue $1$, hence equals
$1$ by (AR.3). The unique remaining factor equals
$\xi(\det D_\kappa)$ by (AR.4) or (AR.5).
This proves both the full subgroup equality and its character law.
Together with AR2 it proves $\texttt{paraShape\_\allowbreak export}$.

\paragraph{AR6: a prescribed nontrivial residue character extends to the completion.}
Let $\mathcal O_v\subset F_v$ be the valuation ring and
$\mathrm{red}:\mathcal O_v\to\mathbb F_p$ the specified residue map.
For every nontrivial additive character
$a:\mathbb F_p\to\mathbb C^\times$, there is an additive character
$a_v:F_v\to\mathbb C^\times$ such that
\[
 a_v|_{\mathcal O_v}=a\circ\mathrm{red},\quad
 a_v|_{\mathfrak m_v}=1,\quad
 a_v|_{\mathcal O_v}\ne1.                               \tag{AR.8}
\]

A homomorphism from a subgroup $H$ of an abelian group to
$\mathbb C^\times$ extends to the whole group. Order partial
extensions by inclusion and take a maximal one. If $x$ lies outside
its domain, either no positive multiple of $x$ lies in the domain,
or the least such multiple is $N x$. In the first case give $x$ any
nonzero value. In the second choose an $N$th root of the known value
at $Nx$. The formula on $H+\mathbb Zx$ is well-defined and extends
the homomorphism, contradicting maximality. Apply this to the
additive groups $\mathcal O_v\subset F_v$.
The maximal ideal maps to zero in the residue field. Conversely a
unit cannot map to zero, since ring homomorphisms preserve units;
thus the kernel is exactly the maximal ideal. A residue with
nontrivial $a$ has an integral lift, giving the last assertion.
The character is locally constant because its kernel contains the
open subgroup $\mathfrak m_v$.

\paragraph{AR7: the type character is the actual local component of the base character.}
Retain the exact $\texttt{SideB}$ normalized-character identities
$\texttt{chiD}$, $\texttt{chi}$, $\texttt{eta\_\allowbreak eq}$,
$\texttt{ringSpec\_\allowbreak normalisedCharacter\_\allowbreak eq}$ and
$\texttt{del\_\allowbreak eq}$, as well as $\texttt{ValLL0}$ and
$\texttt{TauTrue}$. Write $\tau_{j+1,a}$ for the single $L_0$-valued
type character supplied by $\texttt{TauTrue}$. Then, for a unit
$z\in\mathbb F_{g_a}^{\times}$ and an integral representative $u$ at
the type prime $P_a$,
\[
 \iota\bigl(\tau_{j+1,a}(z)\bigr)=\mu_{0,j,P_a}(u).
                                                               \tag{AR.9}
\]
This is an equality for this $\tau$, including its orientation.

We retain the algebraic-Hecke-character-to-adic-character existence
input used by $\texttt{GalChar.exists\_\allowbreak isGal}$; finite-order RC3 alone
does not apply to the normalized algebraic block character.
Under that input, construct an adic character for
$\mu_{0,j}\delta_\ell\eta_T^{-1}$ with the source's norm shift.
The normalized compatibility field gives its Frobenius values,
outside a finite set, as exactly those of the actual $\chi_1$.
Use the same reading of $L_0$ on both sides and RC4 to conclude
equality of these characters. This proves the actual
$\texttt{firstTwistCharacter\_\allowbreak eq\_\allowbreak ladic\_\allowbreak heckeCharacter}$ comparison.
No continuity of an isomorphism $\overline{\mathbb Q}_\ell\simeq\mathbb C$
is used; continuity is on the adic side before applying the algebraic
reading.

To apply it at $P_a$, choose a rational prime $p>g_a$ and a coefficient
place above $p$. The \emph{all-place} field $\texttt{Bs.cover}$ supplies
a twist index with $\ell=p\ne g_a$. This step requires a nonempty
twist family. The prime $P_a$ is unramified and of degree one over
$g_a$, so $g_a\in P_a\setminus P_a^2$ is a uniformizer. Apply the
Artin-unit construction of AR3 to a $(g_a-1)$st root of $g_a$.
The actual $\texttt{tau\_\allowbreak galois\_\allowbreak first\_\allowbreak character}$ law identifies the
result with $\tau_{j+1,a}(z)$. In the normalized block character the
auxiliary factor, $\eta_T$, and the norm factor are all trivial on
these local units. What remains is precisely $\mu_{0,j,P_a}(u)$.
This proves (AR.9). Any unramified factor may be multiplied into
$\mu_{0,j}$ without changing (AR.9).

\paragraph{AR8: the type-place subgroup and its nondegenerate character are explicit.}
Fix a type place $s$ above $P_a$. There are characters $a_s,\xi_s$,
chosen \emph{before} the complex reading $e$, such that $a_s$ satisfies
(AR.8), $\xi_s$ is a local component of $\eta$, and, for every
$e$ satisfying $\texttt{SideRead}$,
\[
 \begin{split}
  \kappa\in M.J&\ \Longleftrightarrow\
       \operatorname{loc}_s(\kappa)\in\operatorname{UniRT}(r),\\
  e(\chi(\kappa))&=
       a_s\!\left(\sum_{a=0}^{r-2} A_{\kappa,a,a+1}\right)
       \xi_s(\det D_\kappa).
 \end{split}                                                   \tag{AR.10}
\]
The formulas apply to $\kappa$ supported at $s.v$, with the character
formula restricted to the level. In $\operatorname{UniRT}(r)$ the
parahoric condition is supplemented by upper-unitriangular reduction
of the first $r$-by-$r$ block $A_\kappa$.

Define the residue additive character by
$z\mapsto\iota(\theta_a(\xi_{\mathrm{red},a}z))$.
Its value at $1$ is not $1$ when $r\ge2$. Insert the element with a
single $1$ on the first superdiagonal at this type place and identities
elsewhere into $\texttt{Bs.psiU}$. The formula for $\psi_U$ becomes
$\theta_a(\xi_{\mathrm{red},a})$, whereas
$X.\mathrm{psiU\_gen}$ gives the nontrivial generator value of its
nondegenerate additive character. Injectivity of the coefficient map
proves the assertion. AR6 extends the prescribed character.
For the formally possible $r<2$ case, choose any nontrivial residue
character; the superdiagonal sum is empty and no prescribed value
is used.

AR7 and AR1 give $\xi_s$ whose unit values are the fixed
$\tau_{j+1,a}$ residue values. Neither choice involves $e$.
AR2 makes all conductor and other type matrices the identity.
Unfolding $M.J$ therefore leaves exactly $\operatorname{UniRT}(r)$.
Its upper block reduces to an upper-unitriangular matrix, whose
determinant is $1$; the lower block has unit determinant by invertibility
of the full reduced matrix. The literal $\mathrm{gVal}$ expression
is the product of the additive superdiagonal character and the
$\tau$-character of that determinant. Reduction commutes with the
finite superdiagonal sum and with the determinant. The
$\texttt{SideRead}$ equality, $\texttt{Bs.psiU}$ and
$\texttt{TauTrue}$ now give (AR.10). This proves
$\texttt{levelUni\_\allowbreak export}$ with its original $\exists a_s,\xi_s\;\forall e$
quantifier order.

\paragraph{AR9: an identity frame identifies the two integral reduction operations.}
Let $P_V$ be the supplied $\texttt{IdFrame}$, with its exact rank
identification $\operatorname{rank}V=n$ and $D.\mathrm{host}.F=E^+$.
Let $\Phi=\mathrm{frameEquiv}(P_V)$. Then
$\kappa\in U_{\max}$ if and only if its framed matrix lies in the
stabilizer of the standard integral adelic lattice. If $\kappa$ is
integral, its framed residue entry is $\bar u$ exactly when the
corresponding selected local entry differs from $s.i(u)$ by an
element of valuation less than $1$.

The frame equation identifies each entry with the same finite adele
under the specified tensor interchange and rank reindexing. An
invertible matrix belongs to the integral general linear group
exactly when its entries and those of its inverse are integral:
the two matrices over the integral ring multiply to $1$, as is
checked after its injective inclusion in the adelic ring. This proves
the stabilizer assertion in both directions.
For the reduction assertion choose an integral $b$ representing the
global adelic residue. By construction the projected adelic entry
differs from $s.i(b)$ by an element of the maximal ideal. It differs
from $s.i(u)$ by such an element if and only if
$s.i(u-b)$ is in that ideal, by the ultrametric triangle inequality
in both directions. This last condition is exactly $u-b\in P$.
Thus the two residue operations coincide. Entrywise, this proves
the parabolic-block and upper-unitriangular equivalences used above
and below; there is no additional freedom to conjugate the frame.

\paragraph{AR10: a common unipotent datum gives the exact power comparison.}
Consider two geometric level data $L,L'$ in this same identity frame.
Assume the actual $\texttt{UniDat}$ data: their leading rank is $r$,
their blocks are $(r,t)$, their type-prime sets are the indexed $P_a$,
their conductor primes are split, lie in $\operatorname{range}P_q$
and avoid the type rational primes, and their residue-group
isomorphisms identify both level subgroups with the same upper
unitriangular group $U$. Finally assume their unipotent characters
are the same fixed function $f:U\to\mathbb C^\times$.
Then both levels contain every supported parahoric element at a
selected conductor place outside the type primes and every supported
$\operatorname{UniRT}(r)$ element at a type place. For an element
$\kappa$ in their intersection supported at $P_a$,
\[
 \chi_L(\kappa)^{g_a-1}
       =\chi_{L'}(\kappa)^{g_a-1}.                       \tag{AR.11}
\]

The containment statements follow from AR9. At other local places
the reductions are identity matrices. At the selected place the
block condition is exactly the required one; membership in $U$ is
detected after the specified residue-group isomorphism.
For a general element in the intersection, pull the two residue
tuples back to $U$. The entrywise equations for those isomorphisms
show that both preimages have the same matrix entries, hence are
the same $u\in U$. Both unipotent factors are therefore $f(u)$.
For an element supported at $P_a$, every remaining character factor
is $1$, except the diagonal block factors at $P_a$. Those factors
are characters of $\mathbb F_{g_a}^{\times}$ and their
$(g_a-1)$st powers are $1$. Raising the two products to that power
leaves the same $f(u)^{g_a-1}$, proving (AR.11).

\paragraph{AR11: the actual Joint is an intersection of two levels.}
Retain the precise input objects of $\texttt{joint\_\allowbreak of}$:
$Y:\texttt{R3.Read}$, $M_{\rm glob}:\texttt{MtFam}$,
$\texttt{MtAd}(Y,M_{\rm glob})$, the rational datum
$\texttt{RatAd}(Y,c_G,k_{K_s},I)$, its $\texttt{SideB}$, the
series $\texttt{SerAt}$ on that same datum and its same identity
frames, $\texttt{Keep}$, the actual level-character law,
$\texttt{Elem}$, and the selected-prime condition
$R.Q(t)\subset\operatorname{range}P_q$.
For $\mu$ satisfying $\texttt{Used0}(D,\mu)$ assume the exact equality
$\eta=R.C.\mu_0(j)\,R.\xi_H(Y.\mathrm{tupT}(\mu))$.
These data construct $\texttt{Joint}(\eta,M_{\rm glob}(\mu,j+1),X,b,h_A,j)$.

Let $L$ be the geometric level $\mathrm{datVP}$ with its literal
false-side level, and $L'$ the series level
$(\rho.B(i,j,b)).\mathrm{levS}(\rho.\mathrm{tw}(i,b,k,Q,\psi))$.
The residue isomorphisms come from $\texttt{Keep.grp}$, with
inverse $c_G$ on the geometric side. The rational law
$\texttt{RatAd.lvl}$ identifies its subgroup with $U$.
Its character law is
$\psi_U=(\psi_0)^{-1}$; the conjugated level in false-side
$\mathrm{datVP}$ contributes the other inverse. Thus its actual
unipotent factor is $f(u)=\iota(\psi_U(u))$.
On the series side the $\texttt{CellTie.U\_\allowbreak eq}$,
$\texttt{psiU\_\allowbreak eq}$, $\texttt{Bs.psiU}$ and $\rho.\mathrm{e\_LL0}$
fields give that same $f$. The selected-prime and block fields give
all remaining $\texttt{UniDat}$ clauses of AR10.

Take the frame automorphism to be the identity and set
\[
 J_1=\bigl(\mathrm{Lev.finE}\circ P_V.\psi\bigr)^{-1}
                   \bigl(\operatorname{lev}(L)\cap\operatorname{lev}(L')\bigr).
                                                               \tag{AR.12}
\]
Take $e=\rho.\mathrm{ek}(i,b,k)$, the polynomial target of the
series, and the injective semilinear map $\omega$ supplied by
$\texttt{SerAt.join}$. Its scalar restriction is the required side
embedding: equality on the coefficient ring extends to its fraction
field by writing each element as a fraction.
The exact join identity is
\[
 \mathrm{rdV}(f)(\Phi(g))=\omega((M.\mathrm{real}\,f)(g)). \tag{AR.13}
\]
AR10 proves both local containment fields for $J_1$. For
$\kappa\in J_1$ take the two scalars to be
$\chi_L(P_V.\psi(\kappa))$ and $\chi_{L'}(P_V.\psi(\kappa))$.
The $\texttt{MtAd.read}$ identity and injectivity of its coefficient
map transfer the first right-translation law from the actual
$Y.\mathrm{IL.jMt}$ form. For the second, (AR.13) transfers the
right-translation law of the series form. That law holds first on
the actual integral $\mathrm{MtLat}$ and then on its full fraction-field
span by linearity; it is not inferred merely from injectivity of
$\mathrm{rdV}$. AR10 proves the required equality of powers of the
two scalars. These are every field of $\texttt{Joint}$.

For the supplied stage, $\mu_{\mathrm{Tw}}(v_0)=\mu_{\mathrm{At}}(v_0)$
on the twisted support proves $\texttt{Used0}$. The literal
$\texttt{lev.xiH}$, $\texttt{tup}$ and $\texttt{tupTie}$ equations
give the displayed equality of $\eta$. Thus this also proves
$\texttt{joint}$ for that same stage. AR21 below constructs the
nontrivial reading used in (AR.13) for an installed series; the other
series and global realization fields are not created by (AR.12).

\paragraph{AR12: matrix intertwiners descend, and dual polynomiality is preserved.}
Let $e:K\hookrightarrow\mathbb C$, and let $A_\gamma,B_\gamma$ be
finite matrices over $K$, indexed by an arbitrary group. If a nonzero
complex matrix satisfies
\[
 X\,e(A_\gamma)=e(B_\gamma)\,X\qquad\hbox{for every }\gamma,
                                                               \tag{AR.14}
\]
then a nonzero matrix over $K$ satisfies the corresponding equations.
Choose a nonzero entry of $X$ and a $K$-linear functional
$\lambda:\mathbb C\to K$ that does not kill it; extend that entry
to a $K$-basis to obtain such a functional. Apply $\lambda$ entrywise.
It commutes with sums and multiplication by coefficients in $K$,
so every equation (AR.14) survives. The chosen entry proves that the
resulting matrix is nonzero. This works for infinitely many
$\gamma$ without an unjustified finite intersection argument.

There is also an explicit polynomiality statement needed in its
application. Suppose the dual of a finite-dimensional $K$-representation
of the rational unitary group has a basis in which its matrices are
polynomials in the entries of the embeddings of rational matrices.
Choose a basis of the original space and use its dual basis. If $A,B$
are the two constant change-of-basis matrices, then the original
matrix at $\gamma$ is the transpose of
$A\,\rho^\vee(\gamma^{-1})B$. On a unitary matrix $g$,
\[
 (g^{-1})_{ab}=\bar g_{ba}.
 \tag{AR.15}
\]
Replace each variable $(\sigma,a,b)$ in the polynomial matrix for
the dual by $(\sigma\circ c,b,a)$, where $c$ is the actual CM
involution, and perform these constant matrix multiplications and
the transpose. This gives polynomial matrices for the original
representation. Mapping their coefficients through $e$ proves
the literal $\texttt{IsAlgWeight}$ condition after complexification.
The source's $\texttt{HWAtF}$ includes the required polynomial
matrix presentation; existence of a weight vector alone would not
suffice for this argument.

\paragraph{AR13: the polynomial coefficient representation descends over the actual field.}
We state precisely the representation-theoretic inputs retained here.
First, an irreducible algebraic representation of the rational definite
unitary group, for a complete CM type, is isomorphic over $\mathbb C$
to the canonical tensor coefficient module $V_{\mathbb C}(\lambda)$
for dominant integral weights $\lambda$. Second, the dual of that
canonical module has exactly the highest weights specified in
$\texttt{contragredient\_\allowbreak polynomialWeight\_\allowbreak unique\_\allowbreak highestWeight}$,
with the source's conjugate-CM-type convention. These are the named
classification and canonical highest-weight inputs, not conclusions
of finite-dimensionality.

Under these two inputs, let $\rho$ be an irreducible finite-dimensional
representation over $K$ and suppose its dual has the literal
$\texttt{HWAtF}$ weights $\lambda_0(\tau)$ for each embedding in a
complete CM type. Then there are dominant weights $\lambda$ and
a $K$-linear equivariant isomorphism
\[
 T:W\xrightarrow{\ \sim\ }V_K(\lambda),\qquad
             \lambda(\tau)=\lambda_0(\tau).             \tag{AR.16}
\]
In particular the field of definition is $K$, not a newly chosen
splitting field.

Choose a $K$-basis of $W$ and complexify its matrices using $e$.
AR12 proves algebraicity. A nonzero complex representation contains
an irreducible subrepresentation: choose a nonzero invariant subspace
of minimal dimension. Algebraicity passes to this subspace: express
its inclusion by a constant matrix, choose a constant left inverse,
and multiply the ambient polynomial matrix by these two matrices.
By the stated classification this subspace is
$V_{\mathbb C}(\lambda)$ for some $\lambda$. Its inclusion in
$W_{\mathbb C}$ is a nonzero complex intertwiner. The canonical
coefficient modules have bases $b_K,b_{\mathbb C}$ on the same finite
index set, with $b_{\mathbb C}=e(b_K)$, and their action matrices
commute with coefficient extension. AR12 descends the inclusion to
a nonzero $K$-intertwiner
$\alpha:V_K(\lambda)\to W$.
Its image is a nonzero invariant subspace, so irreducibility of
$W$ makes $\alpha$ surjective. On the other hand the original
complex inclusion gives
$\dim V_K(\lambda)=\dim V_{\mathbb C}(\lambda)\le\dim_K W$.
Surjectivity gives the reverse inequality. Thus $\alpha$ is an
isomorphism; there is no need to assume that complexification of an
arbitrary irreducible representation stays irreducible.
Transfer the prescribed dual highest-weight equations across
$\alpha$ and then across the common coefficient bases. The second
named input gives $\lambda_0(\tau)=\lambda(\tau)$.
This proves $\texttt{irreducible\_\allowbreak weight\_\allowbreak descends\_\allowbreak to\_\allowbreak polynomial\_\allowbreak module}$
from the two precise complex inputs.

\paragraph{AR14: the actual infinitesimal parameters give CellWt.}
For antitone integer sequences $a,b$ of length $n$, equality of
their source infinitesimal-character multisets implies $a=b$.
Indeed the entries are
$a_i-i+(n-1)/2$. The sequence $a_i-i$ is strictly decreasing, and
so is $b_i-i$. Removing the common shift and comparing successively
the largest entries of the two finite multisets proves equality
at each index. This proves $\texttt{eq\_\allowbreak of\_\allowbreak infOf\_\allowbreak eq}$, including
the case $n=0$.

Apply the complex classification input of AR13 to the actual
$m.\mathrm{wt}$ in $\texttt{MtReal}$, obtaining weights
$\lambda_C$. Its supplied $\texttt{Mt}$ weight field, transferred
through the canonical duality with the conjugate CM type, says
\[
 \operatorname{infOf}(\lambda_C(\tau))
           =\operatorname{paramV}(R,hc_\tau,j,1/2).
                                                               \tag{AR.17}
\]
For the coefficient representation $X.\rho_W(i,j+1,b)$, its
$\texttt{rhoW\_\allowbreak hw}$ field supplies $\lambda_0$.
The exact $\texttt{CharSide}$ assertion
$\operatorname{infG}(\eta_T)=\{1/2\}$ and
$\texttt{rhoW\_\allowbreak wt}$ give the same right side of (AR.17).
AR13 supplies dominant $\lambda_K=\lambda_0$ and an isomorphism
over $K_0=\operatorname{Frac}(O_i)$. The decreasing-sequence argument
therefore gives $\lambda_K=\lambda_C$.
The equality
\[
 (e\circ\mathrm{toKk})\circ\mathrm{kEmb}_\tau
                      =\mathrm{embT}_\tau
                                                               \tag{AR.18}
\]
is exactly $\texttt{SideRead}$ applied to the specified embedded
host element. The actual CM type is injective, disjoint from its
conjugate and covers the complex embeddings, as required in AR13.
The host conjugation is the CM involution because their images
under the specified complex embedding agree.

Pulling back the common canonical polynomial bases through the two
isomorphisms now gives $\texttt{WtMatch}$ over $e\circ\mathrm{toKk}$.
For clarity, if $b_K$ is a canonical basis, then
\[
 \operatorname{coord}_{e(b_K)}(e(y))
                         =e(\operatorname{coord}_{b_K}(y)),
                                                               \tag{AR.19}
\]
by expanding $y$ in that basis. The polynomial action commutes with
coefficient extension, so (AR.19) proves entrywise equality of
every representation matrix.

Finally the actual cell field $\texttt{bc}$ supplies the base-change
matrices from $K_0$ to $K_{i,k}$. Its basis need not be the one just
chosen. If $P,Q$ are inverse change-of-basis matrices over $K_0$,
their images remain inverse over $K_{i,k}$; change the cell basis
by $P$ and use $QAP$ for the action matrices. Applying $e$ gives
the same conjugated complex matrices. Thus
\[
 \texttt{WtMatch}(e,M.\rho,m.\mathrm{wt}.\rho)
                      =\texttt{CellWt}(m,X,b,h_A,j,e)            \tag{AR.20}
\]
holds for every allowed $e$. This proves $\texttt{cellWt\_\allowbreak all}$
under the actual weight fields and the two named classification
inputs; it does not construct those weight fields or $\texttt{MtReal}$.

\paragraph{AR15: the coefficient module is the span of its extremal orbit over a field.}
Here $V_K(\lambda)$ is the tensor of the source's \emph{dual Weyl}
coefficient modules, with determinant shifts. We retain the standard
characteristic-zero equality of the divided-power Weyl module and
the dual Weyl module over $K$, the literal
$\texttt{dualWeyl\_\allowbreak eq\_\allowbreak weyl}$. No integral equality is inferred.
Under that input,
\[
 V_K(\lambda)
   =\operatorname{span}_K
       \{\rho_{\mathrm{poly}}(g)\Delta_\lambda:
                         g\in\prod_\tau\mathrm{GL}_n(K)\}.
                                                               \tag{AR.21}
\]

The nontrivial direction is elementary interpolation. A subspace
stable under row substitutions and containing a polynomial $F$
contains every coefficient of its substitution
$F(X(I+zE_{ba}))$. For each nonnegative integer $z$, the matrix is
invertible: if $a\ne b$ it is a transvection, and if $a=b$ its
nontrivial diagonal entry is $1+z\ne0$ in characteristic zero.
Taking more distinct such $z$ than the degree and inverting the
Vandermonde matrix expresses every coefficient as a linear
combination of these substitutions. Those coefficients are exactly
the row divided powers. By iteration the orbit span contains every
divided-power word applied to the highest polynomial.
The retained characteristic-zero equality makes these span the
dual Weyl module. Tensoring the argument over the embeddings and
absorbing the nonzero determinant characters proves (AR.21).
The reverse inclusion follows directly from stability of that
coefficient module under the polynomial action.
This assertion concerns fields of characteristic zero. It does
not identify the integral lattices $\Delta$ and $\nabla$ in CS.

\paragraph{AR16: complementary minors give the harmonic extremal value and its full image.}
Let $w$ be the actual good weight, with positive and negative
lengths $p_\tau,q_\tau$ satisfying $p_\tau+q_\tau\le m_F$ and
$p_\tau+q_\tau\le n$. Here $m_F$ is the Fock row rank; it is not the
global chain length. Write $a_\tau^+,a_\tau^-$ for its decreasing
nonnegative row lengths and put
$d_{\tau,i}^{\pm}=a_{\tau,i}^{\pm}-a_{\tau,i+1}^{\pm}$, extending
the last row length by zero. Let $\Delta_F=w.\delta(m_F)$, and let
$\mathrm{ev}_y$ evaluate at the source's $\mathrm{fockPt}(w.p,y)$.
For the leading principal minor $L_k(y)$,
\[
 \begin{split}
 \mathrm{ev}_y(\mathrm{minorPlus}_{\tau,k})&=L_k(y_\tau),\\
 \mathrm{ev}_y(\mathrm{minorMinus}_{\tau,k})
    &=(-1)^{k(k-1)/2}\,
                L_{n-k}(y_\tau)\det(y_\tau)^{-1}.
 \end{split}                                                   \tag{AR.22}
\]

The first equality is substitution into a determinant. For the
second, substitution produces the bottom principal minor of the
inverse matrix, with one row reversal. The complementary-minor
identity changes that minor to $L_{n-k}/\det$, while the reversal
has sign $(-1)^{k(k-1)/2}$. The identity holds on all invertible
matrices: multiply by the determinant and use the adjugate
cofactor identity, so no invertibility of an intermediate leading
minor is required.

Multiplying (AR.22), with exponents $d^\pm$, gives
\[
 \mathrm{harmFn}_{B}(\Delta_F)(y)
       =s_w\,\mathrm{coefFn}_{\lambda}(\Delta_\lambda)(B^{-1}y),
 \quad
 s_w=\prod_{\tau}\prod_{i=0}^{q_\tau-1}
             (-1)^{i(i+1)d^-_{\tau,i}/2}\in\{1,-1\}.
                                                               \tag{AR.23}
\]
The exponents agree because consecutive drops telescope:
$\sum_i d^-_{\tau,i}=a^-_{\tau,0}$, and the exponent of each
proper leading minor is
$\lambda_{\tau,k}-\lambda_{\tau,k+1}$; the remaining determinant
exponent is $\lambda_{\tau,n-1}$. This also covers an empty positive
or negative block. It verifies the exact sign, not just a nonzero
proportionality.

The space $\mathrm{harm}(m_F,w,B)$ is, by definition, the span of
the Fock-group orbit of $\Delta_F$. Direct substitution gives
$\mathrm{harmFn}(g\cdot P)(y)=\mathrm{harmFn}(P)(g^{-1}y)$.
The coefficient action has the same law after the inverse-transpose
change $g\mapsto g^{-T}$. As $g$ varies this is a bijection, and
$s_w$ is nonzero. AR15 and (AR.23) therefore show
\[
 \mathrm{harmFn}_{B}(\mathrm{harm}(m_F,w,B))
              =\mathrm{coefFn}_{\lambda}(V_{\mathbb C}(\lambda)).
                                                               \tag{AR.24}
\]
This proves the harmonic image identity, conditional only on the
explicit characteristic-zero input in AR15.

\paragraph{AR17: evaluation is injective on the harmonic orbit space.}
The harmonic evaluation map in (AR.24) is injective. This assertion
has a direct positive-inner-product proof and does not require
the classification input of AR13.
On the polynomial algebra use the Fock pairing
\[
 \langle X^\alpha,X^\beta\rangle_F
                       =\delta_{\alpha\beta}\,\alpha!.
                                                               \tag{AR.25}
\]
It is positive definite. For the mixed Fock substitution $T(g)$,
the plus variable with column index $a$ maps to
$\sum_b(g^{-1})_{ab}X_b$, whereas the minus variable maps to
$\sum_b g_{ba}X_b$, with the row and embedding indices fixed.
Applying the multinomial
formula to monomials proves
$T(g)^*=T(g^*)$, where $g^*$ is conjugate transpose in each factor.
Both sides stay in the finite homogeneous polynomial space in use.

Let $e_0=\mathrm{ev}_1(\Delta_F)$ and
$A=\langle\Delta_F,\Delta_F\rangle_F$.
Formula (AR.22) at $1$ gives $e_0=s_w\ne0$.
For every $g$,
\[
 e_0\langle\Delta_F,T(g)\Delta_F\rangle_F
                         =A\,\mathrm{ev}_1(T(g)\Delta_F).       \tag{AR.26}
\]
To verify it first take $g=ul$, with $u$ upper unitriangular
and $l$ lower triangular. Each leading flag minor shows that
$T(l)\Delta_F=c(l)\Delta_F$, with $c(l)=1$ when $l$ is lower
unitriangular. In particular $T(u^*)\Delta_F=\Delta_F$, so
adjointness gives
$\langle\Delta_F,T(u)\Delta_F\rangle_F=A$.
Evaluating the minors at the upper unitriangular $u^{-1}$ gives
$\mathrm{ev}_1(T(u)\Delta_F)=e_0$.
Multiplication by $c(l)$ proves (AR.26) on this big cell.
Every invertible matrix can be approached by matrices in that cell:
for $(1-z)g+zI$, each required trailing principal minor is a
polynomial in $z$ that is nonzero at $z=1$. Avoid their finitely many
zeros and let $z\to0$. Substitution and inverse matrices are
continuous at $g$. Thus (AR.26) holds everywhere.

By linearity, (AR.26) holds for every polynomial in the orbit span.
If $P$ is in that span and its harmonic function is zero, then
$\mathrm{ev}_1(T(g)P)=0$ for every $g$; the determinant prefactor
in $\mathrm{harmFn}$ is never zero. Equation (AR.26) gives
$\langle\Delta_F,T(g)P\rangle_F=0$.
Adjointness then says
$\langle T(g)\Delta_F,P\rangle_F=0$ for every $g$.
Because $P$ itself is in their span, $\langle P,P\rangle_F=0$,
and positivity proves $P=0$.

The coefficient-function map on $V_{\mathbb C}(\lambda)$ is also
injective: remove its nonzero determinant factor and obtain a
polynomial vanishing on the product of general linear groups.
That set is a nonempty Zariski-open dense subset of the matrix
space, so the polynomial is zero; alternatively apply the
one-variable finite-zero argument successively to its entries.
Together with AR16 this gives a unique linear isomorphism
\[
 \Psi:V_{\mathbb C}(\lambda)\xrightarrow{\ \sim\ }
                      \mathrm{harm}(m_F,w,B),
 \qquad \mathrm{harmFn}\circ\Psi=\mathrm{coefFn}.        \tag{AR.27}
\]
Its equivariance follows by applying the injective harmonic
evaluation map to the two candidate translates.

\paragraph{AR18: the coefficient embedding has all five WtMap properties.}
Retain a good geometric base $B$, its identity frame, the actual
rank equality, a coefficient ring $R$ with fraction field $K$,
$e:K\hookrightarrow\mathbb C$, $f:K_0\hookrightarrow K$, and
$s:L_0\hookrightarrow K_0$ with $efs=\iota$.
Assume the literal $\texttt{IsCoefRing}$, the common weight
$\lambda=B.w.\lambda$, an equivariant
$T_0:W_0\simeq V_{K_0}(\lambda)$, and the actual base-change
matrices for $W$ over $K$. Then there is an injective
$e$-semilinear map $\omega:W\to\mathrm{PolF}$ satisfying every
field of $\texttt{WtMap}$.

Choose common base-change bases $b_0,b_1$, set $c_0=T_0b_0$, and define
\[
 \Omega(w)=\sum_a e((b_1^*w)_a)\,
                      (ef)(c_{0,a}),\qquad \omega=\Psi\Omega.    \tag{AR.28}
\]
Images of the $c_{0,a}$ form a complex basis of the canonical module:
change from its canonical coefficient basis and apply the field
embedding to the inverse change-of-basis matrices. Thus $\Omega$,
and hence $\omega$, is injective. Its values are harmonic by (AR.27).
The same matrix equation used in AR14 shows that it intertwines
the rational unitary actions. The identity frame identifies the
rational matrices on both sides, so this is the literal
geometric action in $\texttt{WtMap}$.

Two lattice assertions remain, and both are essential. Let
$O=e(R)\subset\mathbb C$ and let $\Lambda_h$ be the geometric
coefficient lattice at the actual point $h$. By definition it
consists of harmonic polynomials whose harmonic functions belong
to the $O$-span of the coefficient functions of
$\nabla_{\!L_0}(\mathrm{latAt}(h))$. Each such generator lifts through
$s$ to $V_{K_0}(\lambda)$, and therefore through (AR.28) to $W$.
Its $O$-linear combinations also lift, using $R$-scalars.
Injectivity of coefficient evaluation shows that the resulting
image is exactly the desired harmonic polynomial. Thus
$\Lambda_h\subset\operatorname{range}\omega$.

Conversely, for every $w$ and $h$ there is a nonzero $c\in R$ with
$\omega(cw)\in\Lambda_h$. The lattice
$\nabla_{\!L_0}(\mathrm{latAt}(h))$ spans the coefficient module.
One way to check this is to sandwich the underlying full lattice
between two nonzero integral multiples of a standard lattice,
apply the finite tensor polynomial construction, and clear the
determinant shifts. Its coefficient-field span is unchanged.
Express the finitely many $c_{0,a}$ in the image of that spanning
lattice. Their finitely many coefficients, multiplied by the
coordinates of $w$, lie in $K=\operatorname{Frac}(R)$.
Clear their denominators and take the product of the nonzero
denominators. This gives the required $c$.
These are precisely the range and denominator fields of
$\texttt{WtMap}$, in addition to injectivity, harmonic membership
and equivariance. No surjectivity onto all complex harmonic
polynomials over the smaller field is asserted.

\paragraph{AR19: the global level equality includes every character factor.}
Assume the actual setup $\rho$ is good, pinned to the same
$R,c_G$, has the prescribed level of $D.\mathrm{nexp}$, and has
choice set $\operatorname{range}P_q$ at split primes.
Retain $\texttt{Keep}$, $\texttt{TauTrue}$, $\texttt{ValLL0}$,
the actual base-character law and the stated split auxiliary
places. Then, for every actual $\texttt{TwOK}$ twist,
\[
 g\in M.J\ \Longleftrightarrow\ \Phi(g)\in B.\mathrm{levJ}(t),
 \qquad
 e(\chi(g))=B.\mathrm{levLam}(t)(\Phi(g)).                \tag{AR.29}
\]
This is $\texttt{LevMatch}$, not merely equality on elements
supported at a single place.

AR2 identifies the two selected conductor sets exactly, and the
pinning gives the same type primes, residue isomorphisms and
identity frame. AR9 converts the maximal-compact and every
block-reduction condition. The entrywise residue-group equation
identifies membership of the type tuple in $U$, proving the first
part of (AR.29) for an arbitrary $g$.

For its character, use the actual product formulas on both sides.
The unipotent character is a homomorphism on the product of the
type groups. Decomposing a tuple into its single-place tuples
expresses its value as the product of the corresponding values;
the upper-unitriangular hypothesis ensures each tuple is in its
domain. At each type place, AR7 reads the diagonal residue
character, and the common $\psi_U$ identity reads the unipotent
factor. At $Q_P$, AR3 reads the residue character of $\psi$; at
$D_P$, AR4 reads the residue character of $\delta_i^{\epsilon_b}$.
For any integral invertible matrix satisfying the block condition,
reduction of the lower-right determinant equals the determinant
of the reduced lower block. Both are units, as in AR5.
Hence the same integral representative can be used in every
corresponding factor. Multiplication proves (AR.29).
On the false side the auxiliary Hecke character is first the
character of $\delta_i^{-1}$ and the residue exponent is $-1$;
these operations must be kept together. The calculation gives
the actual cell factor in (AR.5), rather than an extra inverse.

\paragraph{AR20: the actual form reading is pointwise and its integral preimage spans.}
Here is the complete linear-algebra construction used by
$\texttt{read\_\allowbreak cell}$. Suppose $\Phi:G\to G'$ is a group isomorphism,
the rational subgroups correspond in both directions, and
the weight and level comparisons are those of AR18 and AR19.
Let
$\mathrm{real}:V\simeq\mathcal A_K(G,\Gamma,J,\rho,\chi)$
be the supplied actual realization. Assume finitely many
representatives $y_a$ for the geometric double cosets and the
actual covariance of the coefficient lattices. Define
\[
 r(f)(h)=\omega((\mathrm{real}\,f)(\Phi^{-1}h)).          \tag{AR.30}
\]
Then $r$ is injective and semilinear, every integral geometric
form is in its range, and
\[
 \operatorname{span}_K\{f:r(f)\in B.\mathrm{mtLat}\}=V.  \tag{AR.31}
\]

The two transformation laws for $r(f)$ follow by substitution in
(AR.30), rational-subgroup correspondence, and (AR.29).
Injectivity follows pointwise from injectivity of $\omega$ and
$\mathrm{real}$. For an integral geometric form $F$, AR18 gives
a unique preimage $u(g)\in W$ of each $F(\Phi(g))$.
Apply $\omega$ to the required left and right covariance
equations for $u$. They become the known equations for $F$, so
injectivity proves them. Thus $u$ lies in the source form space
and $\mathrm{real}^{-1}(u)$ is the desired preimage.

For an arbitrary $f\in V$, choose a nonzero denominator $c_a\in R$
at each representative $y_a$ using the fifth $\texttt{WtMap}$ field.
Put $c=\prod_a c_a\ne0$. At $y_a$, the value of $r(cf)$ is
$(c/c_a)$ times the integral value provided by $c_a$. Here
$c/c_a=\prod_{b\ne a}c_b$ belongs to $R$, so this value is in the
coefficient lattice. Every point is $\gamma y_a\kappa$.
The geometric action transports that lattice to the required
one, and the right level scalar lies in $e(R)$. To see the last
assertion, the level scalar is a product of characters of finite
residue groups, hence a root of unity. Its preimage under $e$ is
integral over $R$, being a root of $X^N-1$. The actual coefficient
ring is integrally closed, so that preimage belongs to $R$.
Thus $r(cf)$ is integral everywhere. Finally $f=c^{-1}(cf)$ proves
(AR.31). This proof does not require the representatives to have
trivial stabilizer, since the pre-existing form already has its
stabilizer covariance.

\paragraph{AR21: the installed reading supplies the nontrivial series join.}
Under the full input list of
$\texttt{ExistsReadingRealSide.exists\_\allowbreak reading}$, AR13 and AR14
identify the coefficient weight with $B.w.\lambda$:
the setup's good-base weight field and the pinning give the
literal lift weight, and equality of its infinitesimal parameter
uses the same decreasing-sequence argument.
AR18 constructs $\omega$, AR19 gives $\texttt{LevMatch}$, and AR20
constructs $r=\mathrm{rdV}$ with all three required properties:
injectivity, range containing the integral forms, and full span
of their preimage. It also gives, by its definition,
\[
 r(f)(\Phi(g))=\omega((M.\mathrm{real}\,f)(g)),
                                                               \tag{AR.32}
\]
the exact join equation.

The choices are made once for each $(i,j,b,k,Q,\psi)$. If
$\texttt{TwOK}$ holds, make the construction above; if it fails,
choose the zero map, because the conclusions are all conditional
on $\texttt{TwOK}$. This preserves the source's quantifiers and
does not require all possible twists to be admissible. Proof
arguments witnessing admissibility do not alter the cell data.
When $\texttt{exists\_\allowbreak tie\_\allowbreak withSeries}$ installs this reading, its
explicit equality $\mathrm{CellTie.rdV}=r$ converts (AR.32) into
$\texttt{SerAt.join}$. Changing between two identity frames does
not change the finite adelic map: their entrywise matrix-preservation
equations are the same, and the ambient embeddings are injective.
This is exactly the final frame replacement in
$\texttt{exists\_\allowbreak preTie}$.

Thus the join is supplied by an actual coefficient map and
denominator argument for the installed series. We have not proved
the whole $\texttt{exists\_\allowbreak preTie}$ theorem from arithmetic data:
its good pinned setup, scalar normalization, faithful reading,
theta/class-sum installation, Hecke tie and universal-character
fields remain their separately stated inputs. Nor may an arbitrary
previous $\mathrm{rdV}$ be silently replaced by (AR.30).

\paragraph{AR22: the precise rank and choice boundary after these constructions.}
For the same supplied stage, AR1, AR5 and AR8 provide the missing
bad-place and local-shape inputs in RC19. AR11 supplies its actual
$\texttt{Joint}$, with AR21 available for the installed reading.
The existing product-level and nonzero-form arguments in RC14--RC19
therefore give the literal $\texttt{LevEq}$. Independently AR14
gives $\texttt{CellWt}$. RC1--RC2 now give the required dimension
equality and normalized $\texttt{RankTie}$ for each twist under
these hypotheses. This completes the local arithmetic-comparison
part of $\texttt{StatementCOfCasselmanWallach}$ with its specified
global realizations and coefficient structures.

This conclusion does not manufacture $\texttt{LocsCarry}$.
For every actual local index its full
$\texttt{IsGalOf}(D,a,\mathrm{univ},\psi_{L,a},\chi_{w(a)},0)$,
its selected $\mathrm{twOf}$ and its side identities are still needed,
as in RC10. With these, the nonvacuous all-place twist readings,
and the remaining exact choice hypotheses, NC and RC11 construct
the normalized $\texttt{NChoice}$. The separate
$\texttt{MuIsW}$ join still needs the same $\xi,\psi,Q$ and the
actual $\texttt{ColW}$, $\texttt{Pre2}$, $\texttt{RowPk}$ and
$\texttt{St}$ identities of RC13.
The bare LS stage has empty twist and local index sets and cannot
supply $\texttt{Bs.cover}$ in AR7 or the all-place twist readings
required by this construction.

The remaining substantive suppliers are therefore explicit:
global reciprocity and the normalized algebraic adic-character
constructor; the two complex highest-weight inputs and the stated
characteristic-zero coefficient-module identity; compatible
nonempty twist/local data with their normalized character and
weight fields; and the common global realization/pinning/series
installation. The complex HFrame character constructor addresses
its own part of these data. No analytic host/window existence,
local-to-global realization, or RH/GRH conclusion follows merely
from the conditional rank equality proved here.


\section{The actual host character frame and the original column input}
\label{sec:host-character-frame}

\status{HF1--HF13 construct the five actual host character families,
with their prescribed infinity types, restrictions, tame local
characters and inertia orders. HF14--HF15 construct their determinant
quotient and relative determinant laws. HF16--HF19 join these
constructions to the previously reconstructed original geometry,
coefficient fields, column laws, jets and later stages on the same
base. The final paragraph identifies precisely why these existence
results do not yet construct a quantitative \texttt{V4.Core}.}

We use the literal host of the source. Write $E$ for its CM field,
$F=E^+$ for its maximal real subfield, $K=\mathbb Q(\sqrt{-d})$
for its distinguished imaginary quadratic subfield, and
$c:x\mapsto\bar x$ for its specified conjugation. The CM-type is
$\Phi=(\phi_\tau)_{\tau\in\Psi}$, where
$\phi_\tau(\sqrt{-d})=i\sqrt d$. AH9--AH12 establish these
identifications and the fact that $E/F$ is unramified at every
finite prime. AC2--AC3 establish the host's root-of-unity
identities. A general CM field used in intermediate arguments
will be denoted by $L$.

Let $\mathbb I_L$ be the full idele group and
\[
 U_L=L_\infty^\times\times\prod_P\mathcal O_{L_P}^{\times}
       \subset\mathbb I_L,\qquad
 \operatorname{ph}(z)=z/|z|\quad(z\in\mathbb C^\times).
 \tag{HF.1}
\]
This $U_L$ is the source's \texttt{unitIdl}; its infinite
components are unrestricted. A Hecke character is a continuous
homomorphism $\mathbb I_L/L^\times\to\mathbb C^\times$.
Unitary means that every value has absolute value one.
Our finite valuations $\operatorname{ord}_P$ are additive with a
uniformizer of order one; where the source uses multiplicative
valuations, a principal unit satisfies $|u-1|_P<1$.
The exact infinity convention is
\[
 \chi_\infty(k)=\prod_\tau
       \operatorname{ph}(\phi_\tau(k))^{n_\tau}.
 \quad
 \texttt{HasInfQ}(a)\ \text{means}\ n_\tau=2a_\tau\in\mathbb Z.
 \tag{HF.2}
\]
In particular the factor two in the definition of
\texttt{IsMu0.inf} is retained throughout.

\paragraph{HF1: roots of unity and the compatibility needed for extension.}
For a CM field $L/F$, every embedding carries $c$ to ordinary
complex conjugation. If $x\bar x=1$ and
$\operatorname{ord}_P(x)=0$ at every finite prime, then $x$ is
an algebraic integer and every complex conjugate of $x$ has
absolute value one. To see directly that it is a root of unity,
consider the monic characteristic polynomials over $\mathbb Q$
of $x,x^2,x^3,\ldots$ acting on $L$. Their coefficients are
integers, and their elementary-symmetric formulas give uniform
bounds depending only on $[L:\mathbb Q]$. There are only finitely
many such polynomials and finitely many roots of all of them.
Two positive powers of $x$ coincide, so $x$ has finite order.
Conversely a root of unity is an algebraic-integer unit.

For the actual host, AC2--AC3 show that the specified element
$s_g\in\Psi$ is a nonidentity involution and
$s_g(\zeta)=\zeta^{-1}$ for every root of unity $\zeta\in E$.
The map $\tau\mapsto\tau s_g$ pairs the CM-type without fixed
points. The host group is abelian, so
\[
 \phi_{\tau s_g}(\zeta)=\phi_\tau(\zeta)^{-1},\qquad
 \prod_{\tau\in\Psi}\phi_\tau(\zeta)=1,\qquad
 N_{E/K}(\zeta)=1.                                    \tag{HF.3}
\]
The final equality also follows before applying the distinguished
complex embedding: $\Psi$ is exactly $\operatorname{Gal}(E/K)$,
and the same pairing cancels its norm product.

For an integral weight $e:\Psi\to\mathbb Z$, the source's
$\mathrm{MU}(\Phi,e)$ is the assertion
$\prod_\tau\phi_\tau(\zeta)^{e_\tau}=1$ for every such $\zeta$.
The predicate $\mathrm{MUQ}$ includes a choice of these integral
weights equal to the specified rational weights. Thus it gives
the same assertion with exponent $-e$ by inversion. We will use
the actual $\mathrm{MUQ}$ hypothesis for nonconstant weights;
(HF.3) proves it for constant integral weights.

\paragraph{HF2: continuous circle characters with prescribed values.}
Let $G$ be an abelian topological group, $U\subset G$ an open
subgroup, $\psi:U\to S^1$ continuous, $\alpha:A\to G$ a
homomorphism from an abelian group, and $a:A\to S^1$ another
homomorphism. Suppose
\[
       \alpha(y)\in U\quad\Longrightarrow\quad
                    a(y)=\psi(\alpha(y)).             \tag{HF.4}
\]
On $U\alpha(A)$ set
$\gamma(u\alpha(y))=\psi(u)a(y)$.
If two expressions represent the same element, their quotient
is an element of $\alpha(A)\cap U$, and (HF.4) proves equality
of their values. This is a homomorphism, continuous because its
restriction to the open subgroup $U$ is $\psi$.

Every homomorphism from a subgroup of an abelian group to $S^1$
extends to the entire group. Here is the algebraic argument
used in this application. Order partial extensions by inclusion
and take a maximal one. For an element $g$ outside its domain,
if no positive power of $g$ lies in the domain, choose any value
of absolute value one and extend on the subgroup generated by
$g$. Otherwise the integers $n$ with $g^n$ in the domain form
$n_0\mathbb Z$. Choose an $n_0$th root in $S^1$ of the known
value at $g^{n_0}$, and define the extension on $h g^j$
multiplicatively. The minimality of $n_0$ makes it well defined.
Both cases contradict maximality unless the domain was all of
$G$. The resulting extension $\nu:G\to S^1$ is continuous:
continuity at one is already known on the open subgroup $U$,
and translations give continuity everywhere.
Consequently
\[
 \nu|_U=\psi,\qquad \nu(\alpha(y))=a(y).
 \tag{HF.5}
\]
This proves the substantive open-subgroup extension used by
both character producers. It uses divisibility of the circle,
not a global reciprocity existence theorem.

\paragraph{HF3: the finite-unit model and its exact local readings.}
For a finite prime $P$ of $L$, choose a uniformizer $\pi_P$ in
$L_P$, and for $x\in L_P^\times$ put
\[
       x^\circ=x\pi_P^{-\operatorname{ord}_P(x)}.
 \tag{HF.6}
\]
This is a multiplicative map into the units of the local
valuation ring. If $\theta_P:\kappa(P)^\times\to S^1$ is a
residue character, then
$\widetilde\theta_P(x)=\theta_P(\overline{x^\circ})$
is a continuous unitary character of $L_P^\times$.
It kills $1+P\mathcal O_{L_P}$ and agrees with $\theta_P$
on units. A residue of a local unit may be computed from any
global algebraic integer within distance strictly less than
one; two such approximations differ by an element of $P$.
Thus this description agrees with the source's
\texttt{resPart}, including its readings on rational integers.
Every nonzero residue class is reached by an algebraic integer
which is a unit at $P$, so these readings also detect the exact
order of $\theta_P$.

For a CM-type $\phi_i$ selecting each infinite place once,
integers $n_i$, and a finite set $S$ of finite primes, define
\[
 \psi(z)=\prod_i\operatorname{ph}(z_{\phi_i})^{n_i}
                \prod_{P\in S}\widetilde\theta_P(z_P).
 \tag{HF.7}
\]
The completion at $\phi_i$ is identified with $\mathbb C$ by
that specified embedding. Finite products and the local
continuity just proved make $\psi$ continuous and unitary.
At an infinite component it has exactly the stated phase;
at a finite component it is $\widetilde\theta_P$ if $P\in S$,
and is one otherwise. In particular
\[
 \psi(\operatorname{diag}k)
   =\prod_i\operatorname{ph}(\phi_i(k))^{n_i}
      \prod_{P\in S}\widetilde\theta_P(k).
 \tag{HF.8}
\]
The subgroup $U_L$ is open: it is the product of the infinite
group with the standard open compact unit subgroup of the
finite ideles. Imposing finitely many principal-unit
congruences remains open.

Conjugation extends isometrically from $L$ to
$L_P\to L_{\bar P}$ and to all archimedean completions.
It therefore defines a continuous involution $C$ of the
idele group, preserves $U_L$, takes a local component at $P$
to the conjugate component at $\bar P$, and fixes every
idele coming from $F$. These assertions follow coordinate by
coordinate; no identification of different selected primes is
made. They are the actual \texttt{IsConjI} laws used below.

\paragraph{HF4: anticyclotomic characters with a prescribed tame part.}
Let $L/F$ be CM with CM-type $\phi_i$. Let $e_i\in\mathbb Z$,
let $S$ contain no pair $P,\bar P$, including no fixed prime,
and prescribe residue characters $\theta_P$ at $P\in S$.
Assume the exact root-of-unity compatibility
\[
 \prod_i\phi_i(\zeta)^{-e_i}
       \prod_{P\in S}\theta_P(\bar\zeta_P)=1
                 \quad\text{for every }\zeta\in\mu(L).
 \tag{HF.9}
\]
Use (HF.7) with $n_i=-e_i$. Let $A$ be the subgroup of
$L^\times$ consisting of $x$ with $x\bar x=1$, map it
diagonally to the ideles, and prescribe value one on $A$.
If its diagonal lies in $U_L$, HF1 makes $x$ a root of
unity, so (HF.9) is exactly the compatibility (HF.4).
HF2 gives a continuous unitary $\nu$ agreeing with $\psi$
on $U_L$ and killing these norm-one principal ideles. Set
\[
                    \chi(z)=\nu(z/Cz).               \tag{HF.10}
\]
For $k\in L^\times$, $k/\bar k$ has norm one. Thus
$\chi(\operatorname{diag}k)=1$, and $\chi$ is a Hecke
character. It is unitary and continuous. It is trivial
on the ideles of $F$ because $C$ fixes them.

For an infinite component,
$\operatorname{ph}(z/\bar z)=\operatorname{ph}(z)^2$;
therefore its infinity exponents are exactly $-2e_i$.
For a finite local unit $u$ at $P$, the idele
$u/Cu$ lies in $U_L$, so it can be evaluated using $\psi$.
The absence of conjugate pairs in $S$ gives
\[
 \begin{array}{ll}
  \chi_P(u)=\theta_P(\bar u),
                         &P\in S,\\
  \chi_{\bar P}(u)=
       \theta_P(\overline{c(u)})^{-1},
                         &P\in S,\\
  \chi_Q(u)=1,             &Q\notin S\cup\bar S.
 \end{array}                                                \tag{HF.11}
\]
Every finite local character is consequently tame.
For every nonnegative integer $n$,
\[
 \bigl[\chi_P(u)^n=1\text{ for every local unit }u\bigr]
   \Longleftrightarrow \theta_P^n=1
   \Longleftrightarrow
 \bigl[\chi_{\bar P}(u)^n=1\text{ for every local unit }u\bigr].
 \tag{HF.12}
\]
Surjectivity of reduction proves the forward implications;
conjugation is a local-field isomorphism, and inversion does
not change the kernel of a power. This is the exact local-order
part of \texttt{exists\_\allowbreak hecke\_\allowbreak trivOn}, as well as its
unitarity, real-subfield restriction and infinity-type clauses.

\paragraph{HF5: finite order when the infinity type is trivial.}
Suppose a continuous Hecke character $\chi$ is tame at a
finite set $S$, unramified away from $S$, and has trivial
infinity type. Its infinite restriction is identically one:
the diagonal $L$ is dense in $L\otimes_{\mathbb Q}\mathbb R$
by a rational basis and density of $\mathbb Q$ in $\mathbb R$,
so the prescribed identity on nonzero diagonal elements
extends by continuity to the archimedean unit group.

Choose a positive integer $M$ killing all the residue
characters at $S$, for example the product of
$|\kappa(P)^\times|$. Then $\chi^M$ kills every finite
local unit and the full infinite group. Finite-support unit
ideles are dense in the product of finite local unit groups;
hence continuity shows that $\chi^M$ kills all of $U_L$.
The finite valuations give
\[
       \mathbb I_L/(U_L L^\times)\simeq
                                  \operatorname{Cl}(\mathcal O_L).
 \tag{HF.13}
\]
Indeed their map onto fractional ideals is onto by products
of local uniformizers, its kernel is $U_L$, and diagonal
elements give precisely principal ideals. The class group is
finite by the lattice proof in TL14. If its order is $h$, then
$\chi^{Mh}=1$. If every local unit character already has order
dividing $15$, we may take $M=15$.
This proves the finite-order supplier required below; exact
order $15$ on inertia does not assert exact global order $15$.

\paragraph{HF6: selected split primes and the residue norm product.}
If a rational prime $p$ has $[E:\mathbb Q]$ distinct primes
above it, the identity $\sum_{P\mid p}e_Pf_P=[E:\mathbb Q]$
forces every $e_P=f_P=1$. Thus
$\kappa(P)=\mathbb F_p$ by the unique prime-field map,
$p\notin P^2$, and the decomposition group of every such
prime is trivial. Conjugation is nontrivial, so
$P\ne\bar P$. The two primes of $K$ above $p$ are exchanged
by conjugation.

Choose one prime $\mathfrak p$ of $K$ above $p$, and let
$S(\mathfrak p)$ be all primes of $E$ above it.
The group $\operatorname{Gal}(E/K)$ acts simply transitively
on this set. Fix $P_0\in S(\mathfrak p)$ and, for every $P$,
choose the unique automorphism carrying $P$ to $P_0$.
Since all residue maps on rational integers are the same,
reducing the norm product at $P_0$ gives
\[
       \prod_{P\in S(\mathfrak p)}\bar u_P
                    =\overline{N_{E/K}(u)}
                       \quad\text{in }\mathbb F_p^\times
 \tag{HF.14}
\]
for every $u\in\mathcal O_E^\times$. Therefore, if a fixed
character $\omega:\mathbb F_p^\times\to S^1$ is transported
to each of these residue fields and $N_{E/K}(u)=1$,
the product of the resulting character values is one.
This proves the exact residue-product supplier used for both
$\mu_0$ and $\delta$, including its uniformity over every
prime above the selected quadratic prime.

For each of the two auxiliary primes $g_k$, lying-over chooses
such a prime $\mathfrak p_k$ of $K$. The finite multiplicative
group $\mathbb F_{g_k}^\times$ is cyclic: a finite subgroup of
a field has at most $a$ roots of $X^a-1$, which forces its
exponent to equal its order. Choose a generator and send it
to a primitive complex $(g_k-1)$st root. This constructs
the injective character $\omega_k$ and hence the literal
$P:\texttt{PiBar}$ and $W:\texttt{Omg}$.
More generally, if $r\mid p-1$, sending a generator to a
primitive $r$th root constructs a character $\omega$ with
$\omega^n=1$ exactly when $r\mid n$.

\paragraph{HF7: the actual $\mu_{0,j}$ characters.}
Fix the chosen $P,W$ and an integral weight
$e_j:\Psi\to\mathbb Z$ satisfying $\mathrm{MU}(\Phi,e_j)$.
The two distinct auxiliary primes $g_0,g_1$ have disjoint
sets of primes above them. Put
$S=\bigcup_k S(\mathfrak p_k)$ and, on each selected prime
above $g_k$, prescribe
\[
             \theta_P=\omega_k^{-2(j+1)}.             \tag{HF.15}
\]
No conjugate pair lies in $S$. For a root of unity $\zeta$,
the infinite part of (HF.9) is one by $\mathrm{MU}$, and
each finite product is one by (HF.3) and (HF.14).
HF4 therefore constructs a character $\mu_{0,j}$ which
is unitary, is trivial on $F$, has infinity exponents
$-2e_{j,\tau}$, is tame everywhere, and is unramified
away from $g_0,g_1$. On every $P\mid\mathfrak p_k$ and
every natural integer $a$ prime to $g_k$ it satisfies
\[
       \mu_{0,j,P}(a)
             =\omega_k(a\bmod g_k)^{-2(j+1)}.          \tag{HF.16}
\]
These are all six fields of the literal
$\texttt{IsMu0}(P,W,e_j,j)$.
In particular its \texttt{HasInfQ} weight is $-e_j$,
not $-2e_j$. If the weights were supplied rationally,
$\mathrm{MUQ}$ supplies their integral representatives.
Doing this for every $j\in\{0,\ldots,m-1\}$ uses the same
$P,W$ throughout.

\paragraph{HF8: the unramified CM-type character $\chi_\Psi$.}
In HF4 take $S=\varnothing$ and $e_\tau=-1$ for every
$\tau$. Condition (HF.9) is exactly the middle equality
of (HF.3). The resulting $\chi_\Psi$ is unitary, trivial
on the ideles of $F$, unramified at every finite prime,
and has infinity exponent $2$ at every member of the
actual CM-type. This proves all four fields of
$\texttt{IsChiPsi}(H,\chi_\Psi)$.
The construction is on the actual host; it does not replace
the host by an imaginary quadratic field or by the synthetic
character types of a bare local-stage model.

\paragraph{HF9: the actual quadratic product formula on the real subfield.}
For a finite prime $v$ of $F$, put
\[
 \eta_v(a)=
 \begin{cases}
     (-1)^{\operatorname{ord}_v(a)},&v\text{ inert in }E/F,\\
     1,&v\text{ split in }E/F,
 \end{cases}
 \qquad
 \eta_w(a)=\operatorname{sgn}(w(a))\quad(w\mid\infty).
 \tag{HF.17}
\]
There is no ramified finite case by AH11. We prove, for this
actual host, the product formula
\[
             \prod_v\eta_v(a)\prod_{w\mid\infty}\eta_w(a)=1
                         \qquad(a\in F^\times).        \tag{HF.18}
\]
All finite products here have finite support.

First, at a prime $v\nmid d$ of residue cardinality $p^f$,
the sign distinguishing inert and split is
\[
             \eta_v(\pi_v)=\left(\frac{p^f}{d}\right).
 \tag{HF.19}
\]
For odd $p$, reduce the integral element $s_d=\sqrt{-d}$
at a prime $P$ over $v$. It is nonzero. The arithmetic
Frobenius of the unramified quadratic extension acts as
identity in the split case and as conjugation in the inert
case. On this element it acts by
$s_d^{p^f}=((-d/p))^f s_d$, by Euler's criterion and
iteration. Since $p\ne2$, the two signs are distinct.
Quadratic reciprocity and $d\equiv3\pmod4$ identify
$(-d/p)$ with $(p/d)$. At $p=2$, AH12 gives splitting,
and $d\equiv7\pmod{16}$ gives $(2/d)=1$, so (HF.19)
also holds there.

Let $y\in\mathcal O_F$ be nonzero and a unit at every
prime over $d$. Factoring $(y)$ and using (HF.19) gives
\[
 \prod_v\eta_v(y)=
       \left(\frac{|N_{F/\mathbb Q}(y)|}{d}\right),\qquad
 \prod_{w\mid\infty}\eta_w(y)
                    =\operatorname{sgn}N_{F/\mathbb Q}(y).
 \tag{HF.20}
\]
As $(-1/d)=-1$, their product is
$(N_{F/\mathbb Q}(y)/d)$.
We now show that this Jacobi symbol is one using an actual
quadratic subfield of $F$, rather than assuming the product
formula. The specified elements $\sqrt{-q}$ and $\sqrt{-d}$
give $\sqrt{qd}\in F$. The integer $qd$ is not a rational
square since $q,d$ are distinct primes. Put
$F_0=\mathbb Q(\sqrt{qd})$ and
$z=N_{F/F_0}(y)$. This is integral, so its trace $A$ and norm
$N=N_{F/\mathbb Q}(y)$ are integers. Writing
$2z=A+B\sqrt{qd}$ with $B\in\mathbb Q$ yields
\[
                       A^2-4N=qdB^2.                 \tag{HF.21}
\]
The right side is an integer. Its $d$-adic order is
$1+2\operatorname{ord}_d(B)$, so integrality forces
$\operatorname{ord}_d(B)\ge0$ and hence $d\mid A^2-4N$.
The norm $N$ is prime to $d$: an inverse to $y$ modulo
$d\mathcal O_F$ exists, and determinants of multiplication
show that $N$ is a unit modulo $d$. Thus
$N\equiv(A/2)^2\pmod d$ is a nonzero square, proving
(HF.18) for $y$ prime to $d$.

To remove this restriction, let $y\in\mathcal O_F$ be arbitrary
and nonzero. Every prime $v\mid d$ splits in $E/F$ by AH12.
For each pair $P,\bar P$ above such a $v$, use CRT modulo
sufficiently large powers of these distinct primes to choose
$b\in\mathcal O_E$ with
\[
 \operatorname{ord}_P(b)=\operatorname{ord}_v(y),
                    \qquad\operatorname{ord}_{\bar P}(b)=0.
 \tag{HF.22}
\]
The exact order $n$ is imposed by congruence to a local
uniformizer power modulo $P^{n+1}$. Let $a=b\bar b\in F$.
Unramifiedness gives
$\operatorname{ord}_v(a)=\operatorname{ord}_v(y)$ at the
chosen bad primes. Moreover the product in (HF.18) is one
for $a$: all real embeddings are positive, and at an inert
prime its order is twice the order of $b$.

Consequently $x=y/a$ is a unit at every prime above $d$.
Its denominator ideal is coprime to their product. CRT
therefore supplies $t\in\mathcal O_F$, a unit at those
primes, with $tx\in\mathcal O_F$; both $t$ and $tx$ are
units there. The already proved case applies to each of
them, hence to $x$, and then to $y=ax$. Finally write an
arbitrary nonzero element of $F$ as a quotient of two
algebraic integers. Multiplicativity proves (HF.18) in
full. This is the substantive host-specific proof of
\texttt{quadraticCharProductFormulaCMHost}.

\paragraph{HF10: an unramified character with the quadratic restriction.}
We next prove the extension principle actually used for
$\eta$. Let $L/F$ be a CM quadratic extension unramified at
finite primes, with CM-type $\phi_i$ and odd integers $n_i$.
Assume (HF.18) for its finite and infinite quadratic signs.
Suppose that for every root of unity $\zeta\in L$ one can
choose $\delta\in L^\times$ satisfying
\[
 \delta=\zeta\bar\delta,\qquad
 \operatorname{ord}_P(\delta)\in2\mathbb Z
       \quad\text{when }P\text{ lies over an inert prime of }F,
 \qquad
 \prod_i\operatorname{ph}(\phi_i(\delta))^{n_i}=1.
 \tag{HF.23}
\]
The character to be extended on $U_L$ is
$\psi(z)=\prod_i\operatorname{ph}(z_{\phi_i})^{n_i}$,
with no finite factors.

For a finite-support vector $m=(m_v)_v$ on the primes of
$F$, let $\beta(m)$ be the product of the local $F$
uniformizer ideles of exponent $m_v$, mapped to all
primes of $L$ above $v$. Set
$b(m)=\prod_v\eta_v(\pi_v)^{m_v}$.
We prescribe
\[
 \chi(\operatorname{diag}k\,\beta(m))=b(m),
          \qquad k\in L^\times.                      \tag{HF.24}
\]
It suffices, by HF2, to verify compatibility when this
idele lies in $U_L$. For $P\mid v$, its order gives
\[
                    \operatorname{ord}_P(k)=-m_v,
 \tag{HF.25}
\]
because the relative ramification index is one. Thus
$k/\bar k$ is a norm-one algebraic unit, hence a root
of unity by HF1. Choose $\delta$ for this root using
(HF.23). Then $k/\delta$ is fixed by conjugation,
so equals some $a\in F^\times$.

At an inert prime, (HF.23) and (HF.25) give
$m_v\equiv-\operatorname{ord}_v(a)\pmod2$.
The split factors are one. Hence (HF.18) implies
\[
 b(m)=\prod_v\eta_v(a)^{-1}
                =\prod_{w\mid\infty}\eta_w(a)
                =\prod_i\operatorname{ph}(\phi_i(a))^{n_i}.
 \tag{HF.26}
\]
The last equality uses oddness of every $n_i$ and the
bijection from the CM-type to the real places of $F$.
Multiplying by the last equality in (HF.23) shows
$b(m)=\psi(\operatorname{diag}k\,\beta(m))$.
This is exactly the required compatibility; HF2 now
constructs the continuous unitary character (HF.24).
Taking $m=0$ shows it kills all principal ideles.

It is unramified at every finite prime because it agrees
with $\psi$ on $U_L$. Its infinity type is $n_i$.
For a local $F$-component, divide by the corresponding
uniformizer idele of exponent $\operatorname{ord}_v(a)$;
the quotient is a finite unit idele, on which the value is
one. Equation (HF.24) then gives the finite restriction
$\eta_v(a)$. At infinity the odd phase gives
$\eta_w(a)$. Thus all clauses of the literal
\texttt{ResEta} are obtained. This proves the actual
unramified extension theorem from (HF.23), not merely
the existence of an abstract character on a quotient.

\paragraph{HF11: the host phase lemma and the actual $\eta$ character.}
We verify (HF.23) on the host with $n_\tau=1$.
Let $N=Q_Sqd$ be its specified cyclotomic conductor.
AC2--AC3 show that every host root of unity has order
dividing $N$. If $\zeta\ne-1$, set $\delta=1+\zeta$.
Then $\delta=\zeta\bar\delta$.
The involution $\tau\mapsto\tau s_g$ of HF1 takes
$\phi_\tau(\delta)$ to its complex conjugate, so their
normalized phases multiply to one. It follows that
\[
                 \prod_\tau\operatorname{ph}
                         (\phi_\tau(1+\zeta))=1.       \tag{HF.27}
\]
If $P\nmid N$ and $1+\zeta$ had positive order at $P$,
then $-\zeta$ would reduce to one. Its order divides
$N$, which is prime to the residue characteristic.
Reduction is injective on roots of such order: if
$u\ne1$, $u^a=1$ and $u\equiv1\pmod P$, the identity
$1+u+\cdots+u^{a-1}=0$ reduces to $a=0$.
This is impossible. Thus $\delta$ is a unit at every
$P\nmid N$.

For $\zeta=-1$, take $\delta=s_d=\sqrt{-d}$.
It satisfies $\delta=-\bar\delta$ and is a unit outside
$d$. At every CM-type embedding its phase is $i$.
The host rank hypotheses imply $4\mid|\Psi|$, so its
total phase is $i^{|\Psi|}=1$.
Every prime dividing $N$ splits in $E/F$ by AH12.
Consequently any prime over an inert prime of $F$
is outside $N$, and both choices of $\delta$ have
order zero there. This verifies every clause of
(HF.23), with no parity exception at residue characteristic
two.

HF9 supplies the product formula, and HF10 therefore
constructs a unitary Hecke character $\eta$, unramified
at all finite primes, with restriction $\eta_{E/F}$
on $F$ and infinity exponent one at every member of
$\Phi$. These are all four fields of
$\texttt{IsEtaT}(H,\eta)$.
This $\eta$ has the actual complex Hecke-character
domain and restriction. It is not identified with the
formal two-element character type used to inhabit a
bare local-stage record.

\paragraph{HF12: characters of exact inertial order fifteen.}
Let $p$ be a rational prime with $p\equiv1\pmod{15}$
which splits completely in $E$. Choose a prime
$\mathfrak p$ of $K$ above $p$. By HF6 there is a
character $\omega:\mathbb F_p^\times\to S^1$
of exact order $15$. Transport it to every residue
field at $P\in S(\mathfrak p)$.
Use HF4 with $e_\tau=0$ and these residue characters.
Its root compatibility follows from (HF.3) and
(HF.14). We obtain $\delta_p$, trivial on $F$,
tame at every finite prime, unramified away from $p$,
and with trivial infinity type.

Every $P\mid p$ lies either over $\mathfrak p$ or over
its conjugate. Thus (HF.12) says, literally for each
such $P$ and each $n\in\mathbb N$,
\[
 \left[\delta_{p,P}(u)^n=1
                 \text{ for every }u\in\mathcal O_{E_P}^\times\right]
                      \Longleftrightarrow 15\mid n.   \tag{HF.28}
\]
At all other primes it is unramified. HF5 gives a
positive finite global order, in fact an order dividing
$15\,|\operatorname{Cl}(\mathcal O_E)|$.
These are exactly all five fields of
$\texttt{IsDelta}(p,\delta_p)$.
Equation (HF.28) also proves $\delta_p\ne1$: take
$n=1$ at any prime above $p$.

Apply this construction at the host's distinct
$p_\delta,p_\delta'$ to obtain $\delta,\delta'$.
Their ramification supports are disjoint and nonempty,
so they are distinct from one and from each other.
If a set of forbidden rational primes omits $p$,
the same character satisfies \texttt{IsTwist} for that
set: complete splitting gives the split-prime,
prime-residue-norm and $p\notin P^2$ clauses at every
ramified prime. No \texttt{IsPsi} prime-power-order
condition is inferred; that is a different constructor.

\paragraph{HF13: the full frame on the unchanged host and weights.}
Let $H$ be the supplied host and
$e:\operatorname{Fin}(m)\to\Psi\to\mathbb Q$.
Assume that both $g_k$ and both $p_\delta,p_\delta'$
split completely in $E$, that each $e_{j,\tau}$ is
integral, and that $\mathrm{MUQ}(\Phi,e_j)$ holds for
each $j$. HF6 chooses one $P,W$; HF7 constructs all
$\mu_{0,j}$ for these same choices; HF8 and HF11
construct $\chi_\Psi$ and $\eta$; and HF12 constructs
$\delta,\delta'$. The resulting data
\[
          (P,W,\mu_0,\eta,\chi_\Psi,\delta,\delta')
 \tag{HF.29}
\]
with the proved certificates give the literal
$\texttt{HFrame}(H,e)$ of
\texttt{NonemptyHostHeckeCharacterFrame.nonempty\_\allowbreak hFrame}.
Each producer was proved above, including its local
and archimedean restrictions.

For a preliminary datum $P_0$ and its supplied
$\texttt{PreX}$, the required weights are
$e_{j,\tau}=(P_0.\mathrm{hc}_\tau).e_j$;
its \texttt{e\_\allowbreak even} is exactly the required
$\mathrm{MUQ}$ statement. Its level primes are
congruent to one modulo $N$. AH9 implies that their
inertia is trivial and their arithmetic Frobenius is
the identity; hence their decomposition groups are
trivial and they split completely. Thus HF13 gives
the actual $\texttt{Top.Fr}(P_0)$ on that same
preliminary datum. The finite field of character
values will be chosen after this frame, and orientation
of the auxiliary prime factors will be made using
this frame's $P$.

\paragraph{HF14: the actual determinant quotient and rational unitary map.}
Let $K_D$ be the subgroup generated by the following
ideles: those in $U_E$ congruent to one at every prime
above $g_0,g_1$, and the finite local ideles of elements
of $F^\times$. This is exactly the algebraic
\texttt{Subgroup.closure} in the source's \texttt{KD}.
The first set is already an open subgroup, so $K_D$
is open and closed; it consequently also equals its
topological closure. These local generators include
an element of $F^\times$ of order one at each specified
finite prime, which supplies the uniformizers used
below. The subgroup contains every infinite idele.
It also contains every diagonal
$a\in F^\times$. To prove the latter, take all real-subfield
primes below the finite set where $a$ has nonzero order or
where an auxiliary congruence is imposed. Multiply the
corresponding local $F$-ideles with component $a$.
The quotient of $\operatorname{diag}a$ by this finite product
is a unit idele equal to one at the congruence primes.
Both factors therefore lie in $K_D$.

For any character trivial on $F$, tame everywhere and
unramified away from $g_0,g_1$, its finite-idele
restriction kills $K_D$. On the first generating subgroup
this follows from tameness at the auxiliary primes and
unramifiedness elsewhere. To pass from finite products
of unit components to every unit idele, use their density
and continuity. It kills the local $F$ generators by
the exact restriction law. Its kernel is closed, so
it kills the closure. Thus the actual characters
$\mu_{0,j}$ and $\chi_\Psi$ descend to
\[
              D_T=\mathbb I_E/K_D,\qquad
              U_O=\operatorname{image}(U_E\to D_T).
 \tag{HF.30}
\]
Since $\chi_\Psi$ is unramified everywhere, its descended
finite restriction $\chi_D$ kills $U_O$.

For $z\in E^\times$ with $z\bar z=1$, an explicit
Hilbert-90 solution is
\[
 k=1+z\quad(z\ne-1),\qquad
 k=a-\bar a\quad(z=-1,\ a\ne\bar a).
 \tag{HF.31}
\]
In both cases $k\ne0$ and $z\bar k=k$.
Any two such $k$ differ multiplicatively by an element
of $F^\times$, whose diagonal lies in $K_D$.
Consequently
$e_1(z)=[\operatorname{diag}k]\in D_T$
is well defined. Products of solutions solve the
product equation, so it is a homomorphism. It has
precisely the source's rational-unitary determinant
mapping property.

The Hecke product identity on $k$ and (HF.2) for
$\chi_\Psi$ give
\[
 \chi_D(e_1(z))
   =\chi_{\Psi,\infty}(k)^{-1}
   =\prod_\tau\operatorname{ph}(\phi_\tau(k))^{-2}
   =\prod_\tau\phi_\tau(z)^{-1}.
 \tag{HF.32}
\]
For the last equality, apply $\phi_\tau$ to
$z\bar k=k$: the quotient $k/\bar k$ is the square
of its normalized complex phase. Equations
(HF.30)--(HF.32) prove the quotient, kernel,
rational-point and character laws used in the
actual determinant package.

\paragraph{HF15: relative lattice determinants and the full determinant package.}
An idele is in $K_DU_E$ if and only if its additive
orders at $P$ and $\bar P$ are equal for every conjugate
pair. One direction holds for each generator and then
for its closure, since the local order-equality conditions
are closed. Conversely, a balanced idele has only
finitely many nonzero orders. At each underlying prime
of $F$, subtract that common order using a local
$F$-uniformizer idele. Relative unramifiedness makes
its order one at every prime above the given prime.
The resulting quotient is in $U_E$.
It follows in particular that two ideles have the same
class in $D_T/U_O$ whenever all their differences
$\operatorname{ord}_P-\operatorname{ord}_{\bar P}$
are equal.

For a full self-dual lattice $M\subset E^n$, let $I_M$
be its determinant fractional ideal and write
$d_P(M)=\operatorname{ord}_P(I_M)$.
These orders have finite support. The paired localization
of Hermitian duality, proved in NB1, gives
\[
                 d_{\bar P}(M)=-d_P(M).              \tag{HF.33}
\]
Explicitly, a local basis matrix at $\bar P$ has dual
basis $(\bar B^t)^{-1}$ at $P$; self-duality identifies
this dual module with $M_P$. Its determinant has
the negative conjugate order. This uses local bases,
not a global freeness assumption.

Put $\Delta_P=d_P(M')-d_P(M)$. The source representative
for the relative determinant is
\[
 a(M,M')=\prod_P
       \operatorname{loc}_P(\pi_P)^{\max(\Delta_P,0)}.
 \tag{HF.34}
\]
The nonnegative exponent is deliberate: it is the
source's integer \texttt{toNat}. By (HF.33),
$\Delta_{\bar P}=-\Delta_P$, so
\[
 \operatorname{ord}_P a(M,M')
       -\operatorname{ord}_{\bar P}a(M,M')=\Delta_P.
 \tag{HF.35}
\]
At a fixed prime (HF.33) forces $\Delta_P=0$.
For three lattices the differences $\Delta_P$
telescope. The characterization of $K_DU_E$ therefore
proves
\[
 \operatorname{relDet}(M,M)=1,\qquad
 \operatorname{relDet}(M,M')
   \operatorname{relDet}(M',M'')
                       =\operatorname{relDet}(M,M'').
 \tag{HF.36}
\]
If $M'=\gamma M$ for a rational unitary matrix, then
$\Delta_P=\operatorname{ord}_P(\det\gamma)$.
Choose $k$ by (HF.31) for $z=\det\gamma$.
The equality $z=k/\bar k$ gives
$\Delta_P=\operatorname{ord}_P(k)
-\operatorname{ord}_{\bar P}(k)$, so (HF.35) proves
\[
            \operatorname{relDet}(M,\gamma M)
                 =[\,e_1(\det\gamma)\,]\in D_T/U_O.
 \tag{HF.37}
\]
Rank zero causes no exception: the determinant ideal
is the unit ideal and every displayed order is zero.

Take $\mu_{\mathrm{til},j+1}$ to be the descended
finite character $\mu_{0,j}$ in the source's range,
and one outside it. Together with $D_T,U_O,\chi_D,e_1$
and (HF.34), HF14 and (HF.36)--(HF.37) give every
field of the actual \texttt{trueDtPack}. Its relative
determinant is the source's positive-part representative,
not a replacement by a signed product with different
conventions.

The columns use the finer quotient obtained by imposing
principal-unit congruences also at a prescribed finite
set $Q$ and its conjugates. Call its subgroup $K_{D,Q}$.
The proofs in HF14 apply verbatim with this larger finite
congruence set: it is open, contains the infinite ideles
and the diagonal of $F^\times$, and the same finite
characters kill it. Moreover
$K_{D,Q}U_E=K_DU_E$, since both are the balanced-order
subgroup just computed. Thus
$D_{T,Q}/U_{O,Q}\simeq D_T/U_O$, with the identity on
idele representatives. The positive-part formula
(HF.34), (HF.36), and (HF.37) therefore give exactly
\texttt{trueDtPackQ} for every finite $Q$.
This is the determinant package actually inserted in
the column record.

\paragraph{HF16: a common field for the values of this actual frame.}
Apply TL14 to the frame just constructed. For
$\mu_{0,j}$ the infinity exponents are even, so the
normalization parameter is $r=0$; for $\eta$ it is
$r=1$; for $\chi_\Psi$ it is $r=0$. The two finite-order
characters and the finite residue characters require
only cyclotomic extensions. The finite compositum and
its normal closure therefore give a number field
$M\subset\mathbb C$, Galois over $\mathbb Q$, containing
all host embeddings, all $\omega_k$ values, and
\[
 \begin{array}{c|ccccc}
 \text{character}&\mu_{0,j}&\eta&\chi_\Psi&\delta&\delta'\\ \hline
 \text{normalization }s&0&1/2&0&0&0
 \end{array}
 \qquad
       \chi(P)\,N(P)^{-s}\in M
                    \quad\text{whenever }\chi_P\text{ is unramified}.
 \tag{HF.38}
\]
Here $N(P)^{-s}$ is the positive real normalization
specified in \texttt{valTw}. The assertion is uniform
over all unramified primes; it is not only separate
algebraicity for each prime. TL14 proves the needed
uniformity using finitely many ideal-class representatives
and the parity of the infinity exponents.

Choose a primitive generator $\alpha$ of $M$.
If an actual coefficient embedding $\sigma:L_0\to\mathbb C$
contains $\alpha$ in its image, then it contains $M$:
all elements of $M=\mathbb Q(\alpha)$ are rational
functions of $\alpha$ with nonzero denominators.
Equation (HF.38) then gives the literal
$\texttt{ValLL0}$ when $L_0$ is the fixed coefficient
subfield, and \texttt{ValLL} by its inclusion into the
full coefficient field. For every later assembly the
frame \texttt{Fr.at} uses these same characters and
weights, so the same witness proves the uniform
assertion over those assemblies. Containment of
$\alpha$ is a field choice or a field-containment
hypothesis; it does not follow from the record name
\texttt{HFrame}.

\paragraph{HF17: extending coefficients without changing the preliminary datum.}
We give the finite construction behind
\texttt{BaseExt.exists\_\allowbreak base}, since the original
existence join needs an actual base, not just a possible
field of character values. Fix a preliminary datum
$P_0$ and three finite sets of algebraic complex numbers
$S_0,S_1,S_2$ to be included respectively in $K_s,L_0,L$.
Adjoin their elements, all embeddings of a rational
basis of $E$, the finitely many required roots of
unity, and the finite principalization sets described
below. Take the splitting field over $\mathbb Q$
of the product of their minimal polynomials and
$X^2+1$. This is one finite Galois field $L\subset\mathbb C$,
closed under conjugation and containing $i$.

For ordinary ideal capitulation, choose one integral
ideal $\mathfrak c$ in each ideal class of $E$ and
$\mathfrak c^h=(\gamma_\mathfrak c)$, where $h>0$
is the class number. Adjoin a root
$b_\mathfrak c^h=\phi_0(\gamma_\mathfrak c)$.
Comparison of prime-ideal exponents after extension
shows $\phi_0(\mathfrak c)\mathcal O_L=(b_\mathfrak c)$.
Every other ideal in this class differs by a principal
fractional ideal, so these finitely many roots
principalize every extended ideal. The zero ideal
has generator zero.

The source also needs type-norm generators with a
specified absolute value at \emph{every} embedding of
$L$. For a nonnegative exponent function $r_\tau$ on
all complex embeddings of $E$, suppose
$r_\tau+r_{\bar\tau}=n$ is constant, and put
\[
 \Phi_r(\mathfrak a)=\prod_\tau
           (\tau\mathfrak a\,\mathcal O_L)^{r_\tau},
 \qquad
 N_r(x)=\prod_\tau\tau(x)^{r_\tau}.
 \tag{HF.39}
\]
Adjoin for every class representative a root
$b_\mathfrak c^h=N_r(\gamma_\mathfrak c)$.
Unique factorization of fractional ideals gives
$\Phi_r(\mathfrak c)=(b_\mathfrak c)$.
For any embedding $\sigma:L\to\mathbb C$, its action
permutes the embeddings of $E$ and commutes with
conjugation on that CM field. Pairing conjugate
embeddings in the norm product gives
\[
        |\sigma N_r(x)|^2=|N_{E/\mathbb Q}(x)|^n.
 \tag{HF.40}
\]
It follows by taking the positive $h$th root that
$|\sigma b_\mathfrak c|^2=N(\mathfrak c)^n$.
If $(x)\mathfrak a=(y)\mathfrak c$, then
\[
 \beta_\mathfrak a=N_r(x)^{-1}N_r(y)b_\mathfrak c,
 \qquad
 (\beta_\mathfrak a)=\Phi_r(\mathfrak a),\qquad
 |\sigma\beta_\mathfrak a|^2=N(\mathfrak a)^n.
 \tag{HF.41}
\]
The last equality follows by taking ideal norms of
the relation and using (HF.40). No assertion that
an arbitrary embedding of $L$ commutes with its
conjugation is needed: the conjugation identity
used in (HF.40) is only on the embedded CM field $E$.

Apply this to the finitely many source exponents
$r_\tau=p_j$ on the CM-type and $q_j$ on its conjugate,
$0\le j\le m$. This supplies the exact
\texttt{kPsi} generators and their norm equations.
The ordinary roots supply \texttt{capit}.
Now set the actual coefficient field to $L$, and
choose $K_s=L_0=L$ as subfields of it, with the
inclusion into $\mathbb C$ and the restricted complex
conjugation. Insert the specified host embedding,
roots, capitulation and type-norm generators.
Every field of \texttt{Base} is thereby supplied.
The preliminary part is \emph{equal} to $P_0$:
none of its host, spaces, words, weights, type
primes, masses or windows is changed.
This proves the literal three-set extension
statement and makes the subsequent finite-field
requirements simultaneous.

\paragraph{HF18: original canonical columns with pins and algebraic jets.}
Start with the sized preliminary datum and its actual
nonempty decoration supplied by IG10. In particular it
has \texttt{HasLetters}, diagonal spaces, the exact
$\operatorname{duStarOf}$ values, and the integral
weight compatibility \texttt{MUQ}. HF13 constructs a
frame $f$ for this host and these weights.
Orient each auxiliary factor by IG11 using the
\emph{chosen} primes $f.P$: replacing a factor by
its conjugate preserves its norm and congruences,
and makes its conjugate vanish exactly above the
chosen quadratic prime. Apply the normalized-space
construction of IG11 and SE15. It retains the host,
weights, type primes and letters, and supplies
\[
 i_X'(e_a)=e_a,\qquad
       c_u\,\pi_{T,u}=\pi_{T,u_0}.                  \tag{HF.42}
\]
The auxiliary orientation remains unchanged.
Write $P_0$ for this now fixed preliminary datum,
$X$ for its decoration, and $f$ for the same frame.
Only after this normalization do we extend the
coefficient field.

Let
\[
 \Xi=\{1,\delta,\delta'\},\qquad
 S=\{p\in\operatorname{offSet}(P_0):
                     p\ne p_\delta,\ p\ne p_{\delta'}\}.
 \tag{HF.43}
\]
The type primes, $g_0,g_1$, and every rational prime
dividing the host conductor belong to $S$.
For the conductor primes this follows directly
from $p_\delta,p_{\delta'}\equiv1$ modulo the
conductor; the other exclusions are part of the
initial prime separation. HF12 shows that all three
characters in $\Xi$ are \texttt{IsTwist} for $S$.
In particular this assertion concerns the source's
finite tame twists, not its stronger prime-power
\texttt{IsPsi} predicate.

Choose the following finite algebraic sets before
forming the base. Include the primitive generator
$\alpha$ of HF16 in the fixed coefficient field.
Include the PA13 generator for \texttt{Per.FieldOK},
square roots of $g_0,g_1$, and the TL15
\texttt{Holds} generators for $\Xi$ with exponent
$2m$. Finally include the actual finite set
\[
 Z=\texttt{TwIdxEx.gens}(P_0,f.\chi_\Psi,\Xi)
                                                               \tag{HF.44}
\]
from CA12. Its definition depends on this preliminary
datum, frame character and finite family, not on the
base about to be chosen. HF17 now supplies one base
$B$ with $B.\mathrm{toPre}=P_0$ containing all these
sets. Thus the value-field, roots, period field and
jet-field requirements hold simultaneously. There
is no later enlargement that could change the
column input under discussion.

Choose a prime selector over $F$, and use the
specific ambient datum $J=\texttt{amb0}(B.\mathrm{host})$,
whose full laws were proved in PA9. The TL15
character construction gives \texttt{Chars.Ok},
\texttt{Tied} and \texttt{WeilTied} for this $f$.
Here the orientation of $\pi_G$ is essential: it
identifies the source's selected set $v_G$ with
$f.P.\mathrm{Div}$, so HF7's inverse residue powers
are the local characters required by the constructor.
The reference element $\xi_0$ is a unit at these
primes because $g_k\ne d$ and the host different
formula has no denominator there.

For each $\xi\in\Xi$, its finite character values
lie in the chosen coefficient field. The tied
tuple construction of TL15 reads this same complex
character. If a tuple prime $P$ lies in $Q_E$,
the local reading law gives a ramified prime
$P'$ of $\xi$ with $N(P)=N(P')$. The split
degree-one property makes this norm a rational
prime; it is outside $S$, so it is different from
$d$. Hence $\xi_0$ is a unit at $P$ as well.
The square-root and primitive-root requirements
are precisely those already included by
\texttt{Holds}. CI15--CI17 then give the actual
local indices and constants, with
\[
 M=A=B=m,\qquad D(P)=2m
                  \quad(P\notin v_{G_0}\cup v_{G_1}). \tag{HF.45}
\]
Thus every field of the literal raw
\texttt{InEx.exists\_\allowbreak colIn} is supplied: its index
set is $\Xi$, its tuple is the coefficient-field
reading just constructed, its selected characters
are $1,\delta,\delta'$, and its frame, selector and
ambient datum are the fixed $f$, selector and $J$.
The remaining \texttt{Ok} fields are the already
fixed diagonal-space, prime-containment and
$\operatorname{duStarOf}$ identities.

Apply the actual canonical replacement of CI18
to this raw input. The depth hypothesis needed
there is exactly (HF.45). It preserves the base,
decoration, frame, character family, selector,
ambient datum and \texttt{Ok}; it gives
\texttt{IsTw} and changes the two modulus exponents
to $A=B=m+1$. CA12, with the very same set (HF.44),
now proves \texttt{JetsIn}. Its modulus is
\[
 \left(\prod_{Q\in\mathrm{rowSG}}NQ\right)^{m+1}
 \left(T_u\left(\prod_{Q\in\mathrm{rowSG}}NQ\right)^{m+1}\right),
 \tag{HF.46}
\]
not the raw modulus with exponent $m$.
The two equalities (HF.42) give the literal
\texttt{Pins}, independent of this replacement.

We have therefore constructed one original tuple
$(B,X,I)$ with
\[
 \begin{gathered}
 \texttt{HasLetters}(B.\mathrm{toPre}),\quad
 I.\texttt{Ok},\quad\texttt{Pins}(B,X),\quad
 I.\texttt{IsTw},\quad I.\texttt{JetsIn},\\
 I.C=f,\qquad I.S=S,\qquad
 \operatorname{range}(I.\xi)=\Xi,\qquad
 I.J=\texttt{amb0}(B.\mathrm{host}).                  \tag{HF.47}
 \end{gathered}
\]
HF16 also gives \texttt{ValLL0} for every assembly
over this base, and PA13 gives its actual
\texttt{Per.FieldOK}. Finally $\mu_\star\ne\mu_0$:
otherwise the reading law would give $\delta=1$;
at any prime above $p_\delta$, HF12 with exponent
one would then imply $15\mid1$. The same argument
works for $\delta'$, and the disjoint ramification
distinguishes the two nontrivial twists.

\paragraph{HF19: laws, later stages, twisted forms and all lattice jets on that base.}
Apply TL13 to the \texttt{Ok} input in (HF.47).
This constructs its actual \texttt{Laws}; there
is no additional law-existence premise. The
step-zero class bijection preserves the entire
decorated lattice, and the level-zero stabilizers
are trivial. Consequently a coefficient-field
form at step zero is simply a function on the
finite class set with values in the common
polynomial representation. Pullback along this
class bijection is a linear isomorphism. It
preserves the integral lattice because the
lattice at each matched class is equal; its
complex embedding is exactly the source's
\texttt{e0Tw}. This is the \texttt{Law0} proved
in TX10, with the present input and every twist.

The determinant package needed by
$C=I.\mathrm{cols}(l)$ is supplied by HF15 for the
actual finite set $I.SQ$. Its determinant
characters and lattice-determinant representatives
are consequently the ones used in this record.
The class sums give its Hermitian positive
metrics, and the terminal class spaces have
finite dimension, as in FC7 and SE7--SE8.
LS22 can therefore be applied with precisely
this $B,C$ and its retained \texttt{HasLetters}.
It constructs
\[
 Z=(W,V,S_{\mathrm{st}},N,G):
                  \texttt{Later}(B,C),\qquad
 D=I.\mathrm{rec7}(l,Z).
 \tag{HF.48}
\]
All five stage types in this statement are the
actual types over $B,C$; no pre-existing unrelated
stage tuple is substituted.

For these stages define the twisted forms by the
actual chain and metric forms for each $I.\iota$.
Use the preceding class-pullback isomorphism as
the degree-zero map, and the identity as the
map at the reference tuple $\mu_0$. The lattice
and embedding clauses are exactly the two
statements just proved for \texttt{Law0}, or
identities at $\mu_0$. This fills every field of
$I.\texttt{tw}(l,\mathrm{Law0},Z)$ and proves
\[
           \operatorname{Nonempty}(\texttt{TwForms}(D)).
 \tag{HF.49}
\]
In particular the existential
\texttt{TrueColsTest.exists\_\allowbreak all} is now realized
with the additional canonicality and jet clauses
in (HF.47).

The same input supplies the ordinary kernel
\texttt{Hyp} over the canonical seam of
\emph{any} stages above $B$. Indeed its definite
field and prime data are the \texttt{Ok} data;
the ambient curve and semicharacter are PA9 with
the actual $\chi_\Psi$ of HF8; the character
values follow from HF16; the selector and
forbidden-prime inclusions are those of HF18;
and the local constants and their $\mu_q$ tie
are the CI15--CI17 constants. These are all the
clauses of \texttt{hyp\_\allowbreak of\_\allowbreak colIn}.

Here is the precise passage from \texttt{JetsIn}
to the source's \texttt{JetsAll} for self-dual
lattices. Put $H=\operatorname{homAt}(S,N,N')$.
A coefficient function of the required range
is supported on $c_A^{-1}H$ and invariant under
$c_BH$. If it is nonzero at $x$, then $c_Ax\in H$,
and hence $(c_Ac_B)x\in H$. The global jet
statement in (HF.47) applies to this $x$ with
translation denominator $c_B$; $c_B>0$ follows
from the positive type prime and positive local
norms in (HF.46). The multi-index degree it tests
is
\[
 \deg_\tau=\sum_k(\nu_{u,i+1,\tau})_k,
 \tag{HF.50}
\]
which is exactly the required degree by TL19.
Conversely the indicator of $c_A^{-1}H$ is
invariant under $c_BH$, since $c_BH\subset
c_A^{-1}H$, so testing all such coefficient
functions recovers each prescribed point jet.
This proves the literal range/degree equivalence
and then \texttt{jetsAll\_\allowbreak of\_\allowbreak colIn}, for every
$u,a$ and $i<m$. It supplies ordinary lattice
jets over the same ambient datum and seam;
it does not assert the global trace or spectral
identities used by later theta-test theorems.

\paragraph{HF20: the exact remaining global-core join.}
The conclusions (HF.47)--(HF.49) remove the former
existence boundary at the actual host frame and
at the original \texttt{Ok}/\texttt{Pins}/jets/laws/
\texttt{Later} join. They do not supply all fields
of \texttt{V4.Core}. That record refers to one
specified context $K$, one stage tuple $X$,
one frame $f$, and one nontrivial twisting datum
$T$ on their common base. Its additional
requirements include the common object
\texttt{Obj2}, the column and geometric lattice
layouts, the actual row layout and metric
stability, the theta input/output and gauge
identities, the real and signature ties, the
level-character generators, and the arithmetic
choice ties at the prescribed primes.
None can be inferred solely from nonemptiness
of the component record types.

For clarity, the prime-forcing field already
gives a concrete obstruction to using the
untwisted empty model. Its prescribed finite
set can be chosen explicitly as
\[
 \{g_0,g_1,d\}\ \cup\
       \{0,1,\ldots,\lfloor X''_K(n)\rfloor_{\ge0}\}.
 \tag{HF.51}
\]
Every prime at most $X''_K(n)$, and every
distinguished prime required by \texttt{forced},
lies in this set. But for
$T=\texttt{TwDat.untw}(B,\varnothing)$ the set
$T.\mathrm{pl}$ is empty. The \texttt{forced}
clause of a proposed core would put the prime
$g_0$ into that empty set, a contradiction.
Thus the original three-character column
construction, or LS22's bare empty local/twist
stage, cannot by itself be promoted to a core
for that empty twisting datum.

The exact constructive reduction is instead
the source's \texttt{coreOfLay}. On one fixed
$X,f,T$, retain $K.\texttt{Over}(X.B.\mathrm{toPre})$
and supply its actual layout $Y$,
\texttt{RecOK}, the universally quantified
\texttt{LedOK}, \texttt{MuIsW}, \texttt{forced},
the numerical gate identity, the untwisted
row-kernel identity, metric stability, real
tie, true signature, primitive level-character
generators, and the specified arithmetic-choice
conductor-support and exponent equalities
$Y.\mathrm{nch}.Q(w)=T.Q(w)$ and
$Y.\mathrm{nch}.k'(w)=T.k(w)$. Then its common object
is $\operatorname{layObj}(Y)$; the seven column
fields come from $Y$ and \texttt{RecOK}; the
geometric filling and ledger fields come from
$Y$ and \texttt{LedOK}; its row layout is the
given untwisted row identity; and its character,
gauge and arithmetic ties are the stated
same-object hypotheses. In particular the
off-twist value clause uses the \emph{same}
\texttt{MuIsW.off} equality inside the
\texttt{LedOK} conclusion, rather than a
newly chosen tuple map. This fills the literal
core when those remaining suppliers are
available. HF20 records this exact conditional
join and the empty-model obstruction; it makes
no unconditional core, quantitative window,
theta-test, RH or GRH assertion.


\section{Arithmetic of the actual common theta tests}
\label{sec:theta-test-inputs}

\status{TT1--TT10 reconstruct the finite local calculations, the finite
unfolding support, the expansion-independent global sum, and its
coefficient-field descent. They prove the actual \texttt{SumOKK} and
\texttt{SumRatK} conclusions; the latter retains its literal
\texttt{FieldOK} premise. TT11--TT14 construct the elementary parts of
that field and reduce its jet part to a specific comparison of the
same moved vacuum. TT15--TT19 give the compact-period arithmetic,
normalization-ratio argument and compatible finite-field join, with
their remaining analytic and finite-Weil inputs displayed. They do
not infer \texttt{LinK}, \texttt{InjK}, \texttt{Formula2K}, a tower
formula or positive spectral eigenvalues from finite-sum arithmetic.}

Throughout, $\mathcal C$ denotes one literal \texttt{Choice R}, with
its fixed grid $G$, Weil datum $W$, row measures, normal datum $Z$,
frames, representatives and reference twist $t_0=\mu_0$. The
hypotheses called \texttt{HidaSetting} include the actual
\texttt{LevCh.Is}, the arithmetic host and prime exclusions, full
Weil types on both ends, standard additive character, standard
source lattice, and $\beta_X=1$. We retain these hypotheses whenever
we specialize a local calculation to $\mathcal C$. Put
$K=\operatorname{range}(\mathrm{embLL})\subset\mathbb C$.
Membership in $K$ is a field statement. No paragraph identifies an
integral Weyl lattice with its dual or replaces an integral span by
its rational span.

\paragraph{TT1: locally constant integrals are finite arithmetic sums.}
Let $k$ be a nonarchimedean local field, $V=k^d$, and let $\mu$ be
additive Haar measure. Suppose a compact open set $H$ contains zero
and $\mu(H)=1$. If $f:V\to K$ is locally constant with compact
support, then its complex integral belongs to $K$.

Here is the normalization argument, including the case where $H$
is a set rather than a subgroup. The nonzero fibers of $f$ form a
finite collection of compact open sets: compactness gives a finite
cover by neighborhoods on which $f$ is constant. For each such fiber,
a finite cover by translates of a sufficiently small additive open
subgroup supplies a subgroup preserving that fiber. Intersect these
subgroups, their negatives, and the corresponding subgroup for
$1_H$. The result is an open subgroup $U$ preserving $f$ and $1_H$.
Because $0\in H$, $U\subset H$, so $U$ is compact. The support of
$f$ and $H$ are finite unions of $U$-cosets. If $H$ has $N$ such
cosets, then $N>0$ and
\[
 \mu(U)=N^{-1},\qquad
 \int_V f\,d\mu=N^{-1}\sum_{a\in A}f(a),                 \tag{TT.1}
\]
where $A$ is any finite set of support representatives. Every
summand is in $K$, and $N^{-1}\in\mathbb Q\subset K$.
This proves integrability as well as field membership. In
particular it applies to the source's finite product
\texttt{Vars}, with each factor normalized by
$\mu_k(\mathcal O_k)=1$. No topological closure of $K$ is used.

\paragraph{TT2: singular matrix balls and rational multiplicative volumes.}
Let $\nu:k^\times\to\mathbb C^\times$ be ramified, meaning that
$\nu|_{\mathcal O_k^\times}$ is nontrivial. On a matrix ball
$B=X+cM_n(\mathcal O_k)$ containing a singular matrix $X$, the
integral
\[
 \int_{B\cap\mathrm{GL}_n(k)}
       \nu(\det g)|\det g|^s\,dg=0                         \tag{TT.2}
\]
holds with the source's left Haar measure and totalized integral.
If $n=0$ there is no singular $X$, so this case does not arise.
For $n>0$ choose a nonzero left-kernel row $w_0X=0$, scale it to
$w\in\mathcal O_k^n$ with some $w_i=1$, and choose a unit $u$
with $\nu(u)\ne1$. Then
\[
 M=I+(u-1)e_iw,\quad \det M=u,\quad MX=X,\quad
 M^{-1}=I-u^{-1}(u-1)e_iw.                                 \tag{TT.3}
\]
Both $M$ and $M^{-1}$ are integral. Thus left multiplication by
$M$ preserves $B\cap\mathrm{GL}_n(k)$ and its measure, while
multiplying the integrand by $\nu(u)$. Subtracting proves (TT.2).
For a genuinely convergent integral this is ordinary change of
variables. The same equality is valid for the totalized integral;
it is not an assertion of convergence.

Conversely, suppose every matrix in a compact open ball $B$ is
invertible. Let $J\subset\mathrm{GL}_n(\mathcal O_k)$ be open
with Haar volume one, and assume $\nu$ is trivial on
$1+\pi^e\mathcal O_k$. There is a compact open subgroup
\[
 U\subset J\cap(1+\pi^eM_n(\mathcal O_k))
\]
whose right action preserves $B$. To obtain it, cover the compact
set $B$ by finitely many matrix neighborhoods contained in $B$ and
use continuity of multiplication; perform the same construction
for the complement near the compact set in question. Equivalently,
a sufficiently deep principal subgroup preserves each ball by the
bound on $B^{-1}$. The determinant character and determinant norm
are constant on each right $U$-coset. The left cosets $gU$ have
volume $\mu(U)=[J:U]^{-1}\in\mathbb Q$. Hence, whenever
$\nu(\det g)|\det g|^s\in K$, the integral over $B$ is a finite
sum of elements of $K$ times this rational volume. Notice that the
normalization concerns $J$, not all of $\mathrm{GL}_n(k)$.

\paragraph{TT3: the actual frame-orbit integral has values in the field.}
Let $\Psi$ be supported on the matrix lattice $\operatorname{Mat}
(\Lambda)$ and invariant under translation by
$\pi^e\operatorname{Mat}(\Lambda)$, where $\Lambda$ is a diagonal
full lattice and $e\geq1$. Suppose its values, the values of the
ramified character $\nu$, and every half-integral determinant norm
occurring below lie in $K$. For the fixed full-rank matrix
$Y=Y_0h$ in \texttt{smallF}, consider
\[
 \mathcal F(Y)=\int_{\mathrm{GL}_a(k)}
       \Psi(gY)\nu(\det g)|\det g|^{D/2}\,dg.              \tag{TT.4}
\]
Its value is in $K$, with the source's normalization
$\operatorname{vol}(J)=1$ and literal conductor bound
$\operatorname{CondLE}(\pi,e,\nu)$.

Indeed $Y$ has a right inverse. The implication
$\Psi(gY)\ne0\Rightarrow gY\in\operatorname{Mat}(\Lambda)$
therefore bounds every entry of $g$. Multiplication by $Y$ also
shows that $g\mapsto\Psi(gY)$ is invariant under addition of
$\pi^NM_a(\mathcal O_k)$ for some $N$. Thus its support is a
finite union of such matrix balls, with a constant coefficient
$\Psi(gY)$ on each. A ball containing a singular matrix contributes
zero by TT2. On every remaining ball TT2 gives a value in $K$.
Summing proves (TT.4) when the original integrand is integrable.
If it is not integrable, its source integral is defined to be zero,
which is also in $K$. This last branch is essential: TT3 proves the
arithmetic of \texttt{smallF}, not its analytic absolute convergence.

\paragraph{TT4: finite Fourier flips preserve the specified arithmetic.}
Write $q=|\mathcal O_k/\pi|$. For the source's block lattice $L$
and complementary block lattice $L'$, finite Fourier exchange has
coefficients
\[
 q^{-e r}\,\psi(\langle X,Y\rangle),                       \tag{TT.5}
\]
where $r$ is the number of transformed entries and the representatives
range over the prescribed finite quotient. The inverse transform uses
the opposite character and its prescribed normalization. If $K$
contains the required additive-character values, both transforms
preserve $K$-valued functions. The identities governing their support
and period follow from
\[
 \sum_{x\in A}\chi(x)=
 \begin{cases}|A|&\chi=1,\\0&\chi\ne1,\end{cases}          \tag{TT.6}
\]
for a finite additive group $A$. In the second case translate by an
$x_0$ with $\chi(x_0)\ne1$ and cancel. Applying (TT.6) entry by
entry proves the exact exchange of a lattice and its annihilator;
there is no analytic Fourier-limit step.

The middle operation \texttt{phiA} is zero off its prescribed
primitive-frame lattice. On that lattice it is a determinant
character/norm factor times (TT.4), so it too is $K$-valued. It has
the claimed depth for the following explicit reason. If $X$ is a
primitive frame and $Z\in\pi^e\operatorname{Mat}(L')$, extend
$X$ to an integral basis of $L'$. In that basis the rows of $Z$ are
$\pi^e$-divisible; fill the remaining rows by zero. The resulting
matrix $I+A$, with $A\in\pi^eM(\mathcal O_k)$, is invertible
and gives $X+Z=X(I+A)$. Transporting back gives a lattice
congruence element $\kappa$. Its inverse transpose preserves the
dual lattice, its determinant is in $1+\pi^e\mathcal O_k$, and
the defining change of frame leaves \texttt{smallF} unchanged.
The conductor bound kills the determinant character and its norm
is one. Nonprimitive frames remain nonprimitive after this small
translation. These arguments prove both the support and period
parts of \texttt{hasDepth\_\allowbreak phiA}.

A square flip is precisely forward Fourier exchange, this middle
operation, and inverse Fourier exchange. The other tensor variables
are unchanged. Induction along the actual ordered list
\texttt{flipOrder} therefore preserves both the three clauses of
\texttt{IsTensor} and field membership. At a new vertex its level
is the full integral group, so the required open-subgroup and
integrality hypotheses persist. The original seed has the same
properties because it is the prescribed product of lattice
indicators and determinant characters. This proves arithmetic for
the actual \texttt{Tflip}, not a substitute transform.

Here are the exact field hypotheses in the global application.
The character is
$\nu_{ij,t}=\mu_{0,j}\xi_t/\chi_{m-i}$. If $x=u\pi^a$ then
\[
 |x|^{z/2}=(\sqrt q)^{-az}.                                \tag{TT.7}
\]
Thus \texttt{NumbersIn.chars}, \texttt{sqrt} and \texttt{roots}
supply all coefficients. In particular $q\in\pi\mathcal O_k$,
so $q^e$ kills $\pi^{-e}\mathcal O_k/\mathcal O_k$ and
$\psi(x)^{q^e}=1$ there. This does not identify the rational
integer $q$ with $\pi$, or its valuation with one. The same
calculation, including the normalized half-twist of $\eta$, proves
that the source's \texttt{moved} tensor is $K$-valued after its
specified inverse character factor and frame substitution.

\paragraph{TT5: flat integrals, their phase bounds and their local values.}
For fixed column and row variables $x,y$, the flat integrand is
\[
 e\longmapsto T_g(\operatorname{fill}(x,y,e))
                       \psi(\operatorname{phase}(x,e)),     \tag{TT.8}
\]
where $T_g$ is the actual moved opposite-path tensor. It is locally
constant. It has compact support in the free variables $e$: each
free variable is an entry of a specified \texttt{stepPt} matrix,
and the \texttt{downMap} frame is invertible. If the tensor is
nonzero, its lattice support bounds the entries of every filled
edge after conjugating by fixed matrices. Solving these entry
identities bounds each free variable. A closed bounded box in a
finite-dimensional local vector space is compact. This proves
actual integrability of (TT.8), in contrast to the totalized
multiplicative integral in TT3.

We also record the uniform bound needed for the phase field. For a
fixed moved tensor, all column and row entries visible in its edge
matrices are bounded on the support of the flat integral. The one
remaining last-row direction is bounded by the same translation
argument as TT2: translate the relevant free variables by a small
additive direction. Tensor periodicity leaves the first factor
unchanged, while the second factor is multiplied by
$\psi(a t)$, where $a$ is the last-row entry. A nonzero integral
forces $\psi(a t)=1$ for every $t$ in the chosen small ball.
The annihilator of a local lattice is a lattice; hence $a$ lies in
a fixed bounded ball. This supplies the omitted last-row bound.
Each term of \texttt{phase} is either a matrix entry or a product
of two such bounded entries. The ultrametric triangle inequality
gives, for some integer $N$, the exact implication
\[
 \operatorname{flat}(x,y)\ne0,\quad
 T_g(\operatorname{fill}(x,y,e))\ne0
 \ \Longrightarrow\ \operatorname{phase}(x,e)
                          \in\pi^{-N}\mathcal O_k.          \tag{TT.9}
\]
No bound on $e$ or $x,y$ is asserted when these nonvanishing
conditions fail.

If \texttt{PhasesIn} holds, the integrand is $K$-valued wherever
its tensor factor is nonzero, and it is zero elsewhere. TT1 proves
that every local \texttt{term.datum} value is in $K$. The auxiliary
AA vacuum is a lattice indicator times the standard conductor
phase, and the same compact-box argument proves its integrability.
Its separately prescribed scalar is a product of powers
$q^{n/2}$; it belongs to $K$ by (TT.7) at the AA primes.

\paragraph{TT6: nonzero unfolding terms have finitely many coset labels.}
We give the local bounded-representative argument, since mere local
constancy would not imply finite support on an infinite double
quotient. Write $\operatorname{Bd}(\gamma,M)$ when all entries of
$M$ have valuation at most $\gamma$, and take $\gamma\geq1$.
The elementary block identity behind the argument is this. If
\[
 M=\begin{pmatrix}A&B\\0&D\end{pmatrix},\qquad
 M^{-1}=\begin{pmatrix}A'&B'\\0&D'\end{pmatrix},
\]
and $D,B',D'$ have bound $\gamma$, then
\[
 M=\begin{pmatrix}A&0\\0&I\end{pmatrix}
       \begin{pmatrix}I&-B'D\\0&D\end{pmatrix},\qquad
 \begin{pmatrix}I&-B'D\\0&D\end{pmatrix}^{-1}
       =\begin{pmatrix}I&B'\\0&D'\end{pmatrix}.             \tag{TT.10}
\]
The two right-hand matrices have bound $\gamma^2$. This is exact
block multiplication, using $AB'+BD'=0$ and $D'D=I$.

At one step of the unfolding, perform lattice Gaussian elimination
to write the next matrix as $h=b\kappa$, where $\kappa$ is
integral with integral inverse and $b$ is block upper triangular
for the source's \texttt{midLev}/\texttt{lev} flag. Suppose the
previous matrix and its inverse have a fixed entry bound. A nonzero
flat integral has a point with nonzero integrand, so tensor support
bounds $h^{-1}\operatorname{stepPt}(x,e)$ after multiplying by
the previous matrix. For a free tail direction use
$d=bE_{\mathrm{tail}\,i,\,\mathrm{tail}\,(i+1)}t$, with
zero-based tail indices. This $E$ is rectangular: its tail has
$c$ rows and $c+1$ columns, exactly as in \texttt{stepPt}.
Thus for the last row $i=c-1$ its column $i+1=c$ still exists.
For $c=1$ the same translation uses row zero and column one,
and bounds the sole tail diagonal; no endpoint is omitted.
Write $\rho=\lfloor c/2\rfloor$ and $p=\rho+1$. The exact free pattern is: no core--core entries;
tail row $l$ and a core column when $l+1<p$; a core row and tail
column $j$ when $p<j+1$; and tail row $l$, tail column $j$ when
$l<j$. The auxiliary translation uses only the free entries,
not the extra \texttt{Kept} entries. Its sole possible nonzero
column is tail $i+1$. Upper triangularity for
$\operatorname{midLev}(\rho)$, which assigns level $2\rho+1$
to the core and $2(l+1)$ to tail $l$, forces a nonzero tail row
to have $l\leq i$, and a nonzero core row to have $i\geq\rho$.
These are exactly the required free inequalities in column $i+1$.
All core columns are zero. Moreover,
\[
 h^{-1}d=\kappa^{-1}E_{\mathrm{tail}\,i,\,\mathrm{tail}\,(i+1)}t.
\]
Since $\kappa^{-1}$ is integral, a fixed ball for $t$ makes this
a tensor period after multiplication by the bounded previous
matrix. The ball is independent of the possibly unbounded $b$.
The source's character reads the sum of entries
$(\mathrm{tail}\,l,\mathrm{tail}\,(l+1))$, so on $d$ it is
exactly $b_{ii}t$. Zero core columns keep the additional row
terms in the phase unchanged. Translating the integral then forces
\[
                    \psi(b_{ii}t)=1                         \tag{TT.11}
\]
throughout that ball. The annihilator bound gives a common bound
for every tail diagonal entry of $b$. The fixed entries of
\texttt{stepPt} likewise bound the tail columns of the inverse
of the diagonal block part of $b$. Removing its flag-unipotent
part and applying (TT.10) yields
\[
                  h=r w,\qquad r\in R_j,\qquad
     \operatorname{Bd}(C,w),\quad\operatorname{Bd}(C,w^{-1}), \tag{TT.12}
\]
where $R_j$ is exactly the core general-linear group times the
specified symmetric unipotent group. No enlargement to all
unipotent matrices has been made. This argument also covers an
empty core; a zero tail has no new tail equations.

The set of $w$ in (TT.12) is compact. To see closedness without
assuming inversion extends to singular matrices, view it as the
projection of the compact set of pairs $(w,v)$ in the two bounded
matrix boxes satisfying $wv=vw=I$. If $U$ is the prescribed open
right level, its image in $R_j\backslash\mathrm{GL}/U$ is finite:
the open right cosets cover that compact set and a finite subcover
suffices. Induct on the vertices. There is only the fixed initial
label; finitely many previous representatives have one common
entry bound; (TT.12) then leaves finitely many next labels. The
finite product of these finite label sets contains every nonzero
unfolding term. Taking the product over the finitely many active
places proves the literal \texttt{S.supp} is finite.

\paragraph{TT7: pure-coset expansion and fixed-depth uniqueness.}
For a fixed nonzero unfolding label, TT5 gives bounded support and
a sufficiently small common translation period for its flat datum.
After clearing the finitely many frame and half-lattice denominators,
choose one positive integer $k$ so that each relevant half-point
set is partitioned into finitely many cosets
\[
                  x_0+k\Lambda_{\rm half},\qquad
                       kx_0\in\Lambda.                    \tag{TT.13}
\]
Here the source's idempotent is fixed, and
$\Lambda_{\rm half}$ means its \texttt{IsHalf} subgroup; it is
not an arbitrary choice of a half of the ambient vector space.
Existence of the finite representatives is the finite quotient of
a bounded fractional lattice by a full integral sublattice. CRT
and localization identify this quotient with the actual local
cosets in the given frames. There is no uniform bound for $k$
claimed in this paragraph.

For each tuple of representatives take the product of its top and
row coset indicators and assign to it the value of the flat datum
at the tuple of centers. Only finitely many coefficients are
nonzero. Periodicity makes this a pointwise expansion on half
points; off the support both sides vanish. With the prescribed
prime set, idempotent and $z=1$, these are actual
\texttt{FinDatum}s and give \texttt{IsCosetExpansion}. In
particular some positive depth exists, so the source's least depth
\texttt{expDepth} exists and is positive.

The following stronger uniqueness is useful. Two depth-$k$ coset
indicators with the same idempotent that are nonzero at a common
point have the same coset. With the same prime set and $z=1$ they
are the same \texttt{FinDatum}. Thus the product indicators of
different pure data have disjoint supports. Evaluating any
expansion at a tuple of centers of a nonzero term $d$ gives
\[
       \tau(d)=\prod_w
        \operatorname{datum}_w(x_w,y_w).                   \tag{TT.14}
\]
It also forces that same $d$ to occur in every other expansion at
the same depth, with the same coefficient. Consequently the union
of the supports of \emph{all} depth-$k$ expansions is contained
in one chosen finite support. This proves the literal
\texttt{finite\_\allowbreak pureDat\_\allowbreak of\_\allowbreak depth}; it does not use the false
inference that a union of arbitrary finite sets is finite.

\paragraph{TT8: expansion independence and the actual SumOK supplier.}
The class pairing is independent even of expansions not restricted
to the special disjoint basis (TT.13). At fixed idempotent, each
finite theta-vector construction is linear in its finite datum
function: apply the prescribed finite Heisenberg sum to that
function. The theta-vector characterization and its normalization
identify this explicit vector with the chosen $Z.\mathrm{vec}$.
Thus the kernel operator is linear in each datum function; after
composition down the top and row chains, the final class pairing
is multilinear in the $2m$ datum slots.

To justify the tensor argument without an infinite-dimensional
assumption, collect the finitely many functions occurring in the
two expansions in each slot and take their finite-dimensional
spans. Point evaluations separate each such span: if every
function in a subspace vanished at every point, that subspace
would be zero. Choose finitely many evaluations giving an
injective coordinate map by successively reducing the kernel.
Their tensor product is injective. Therefore a finite sum of
pure tensors whose associated product function vanishes pointwise
is the zero tensor. Every multilinear functional annihilates it.
Apply this to the difference of the two expansions. Both vanish
off the specified half points, so their equality on half points
is sufficient. Summing over classes preserves the equality.

The remaining class conditions are actual consequences of the
geometric setting. Every top and row space used here is positive
definite at all archimedean places. Its level is open compact:
finite local projected levels are open and bounded, and the
outside lattice stabilizers are compact. The lattice argument of
CR3 gives finitely many classes; the automorphism group of a
positive definite integral lattice is finite, so every stabilizer
order is a positive integer. Accordingly
\[
             \operatorname{invGam}(c)
                    =\prod_j|\Gamma_{c,j}|^{-1}\in\mathbb Q^\times.
                                                                  \tag{TT.15}
\]
Trivial stabilizers are not required. The displayed weight lists
in the grid give the indicated interlacing and common determinant
shift by substitution. Reading the row coefficient space uses its
actual \texttt{funEquiv}; it is an isomorphism of the stated
polynomial representations, not an assertion about equality of
two integral lattices.

TT4--TT7, the preceding multilinearity, and these class facts
supply, respectively, all eleven fields of \texttt{SumOK}:
\texttt{flat}, \texttt{flatA}, \texttt{fin}, \texttt{ex},
\texttt{vec}, \texttt{indep}, \texttt{cls}, \texttt{gam},
\texttt{wts}, \texttt{br}, and \texttt{read}. They hold for every
literal choice at the given Hida setting. Thus \texttt{SumOKK}
holds with both its additional largeness thresholds equal to zero.
The source's \texttt{rowSum\_\allowbreak of\_\allowbreak hypotheses} has a nominal
\texttt{FunPr.Hypotheses2} argument, but its proof does not use it;
the proof above likewise uses none of that additional package.

\paragraph{TT9: rational polynomial coordinates, operators and branching.}
Let $b_1,\ldots,b_r$ be the prescribed bideterminant basis of a
coefficient space $VV_\lambda$. Every $b_i$ has integer
coefficients. A polynomial in its complex span has all coefficients
in $K$ if and only if its $b_i$-coordinates lie in $K$. For the
nontrivial implication choose finitely many monomials such that
the corresponding coefficient matrix of the $b_i$ has a nonzero
$r\times r$ minor. That minor has rational entries and inverse
over $\mathbb Q$. Reading those polynomial coefficients recovers
all $b_i$-coordinates in $K$. This is finite-dimensional field
descent, with no assertion about unimodularity of that minor.

A rational frame acts by linear substitution in the variables and
an integral power of a nonzero determinant. All its coefficients
lie in $K$ after the fixed embedding of $E$; inverse transpose
uses the inverse of that same frame. The matrices $Z.\mathrm{kk}$
are in $K$ under \texttt{JetsIn}. Hence the exact top and row
chains preserve $K$-polynomial coefficients. The row chain starts
with the source's \texttt{funOf} and the same inverse rational
frame, in that order.

The final branching pairing also preserves $K$. The Fock pairing
has rational diagonal coefficients (products of factorials) in
the monomial basis. Restriction, adjugate substitution and
bideterminant operations have rational coefficients. Orthogonal
projection onto the indicated rational polynomial subspace is
computed by inverting its rational nondegenerate Fock Gram
matrix, and is therefore rational. The maps \texttt{covOf} and
\texttt{funOf} are inverse rational maps on their stated spaces.
Finally the extremal normalization \texttt{extNormB} is the
nonzero rational Fock self-pairing of the fixed extremal
polynomial. Dividing by it proves
\[
       U,V\hbox{ with coefficients in }K
       \quad\Longrightarrow\quad\operatorname{brPair}(U,V)\in K.
                                                                  \tag{TT.16}
\]
These statements concern the bilinear Fock/branching conventions
in the source, including their specified conjugate identification;
no field is silently assumed stable under arbitrary complex
conjugation. All maps used in (TT.16) are the rational polynomial
maps after that identification.

For the actual source vector $x$ in its rational form, VE5--VE11
and the given model identity express its value at a finite point
as a rational coefficient vector multiplied by the actual level
character. That character is a product of the type characters and
inverse twist characters; the initial flag character is one.
Their values lie in $K$. Applying the fixed rational coefficient
identification gives $\texttt{formC}(j_Xx)$ with polynomial
coefficients in $K$. This checks the same-object input to the row
chain; a rank equality alone would not suffice.

\paragraph{TT10: the actual global sum is rational over the given field.}
Choose a depth-$\mathrm{expDepth}$ expansion $\tau_\beta$ for
each nonzero unfolding label. TT8 permits this choice inside the
literal definition of $v_K$. That definition is
\[
\begin{split}
 v_K(t,x,a)={}&s_{AA}
 \sum_{\beta\in\operatorname{supp}}
       \left(\prod_w c_w(\beta_w)\right)
 \sum_{c\in\operatorname{Classes}_\beta}
       \operatorname{invGam}(c)\\[-2mm]
 &\hspace{6mm}\cdot\sum_{d\in\operatorname{supp}\tau_\beta}
       \tau_\beta(d)\,
       \operatorname{brPair}\bigl(
        \operatorname{colMap}_d\Psi_a(c_{\rm top}),
        \operatorname{rowMap}_d\operatorname{formC}(j_Xx)(c_{\rm row})
                              \bigr).                       
\end{split}
 \tag{TT.17}
\]
Every displayed sum is finite by TT6--TT8. The local coefficients
$c_w(\beta_w)$ are the prescribed rational quotient volumes.
TT5 gives $s_{AA}\in K$, and (TT.14) and TT5 give
$\tau_\beta(d)\in K$. Equation (TT.15) treats the stabilizers.
The test column $\Psi_a$ is the scalar extension of its actual
coefficient-field vector, so TT9 applies on both sides of the
pairing. Every term of (TT.17) is consequently in $K$.

This proves the literal \texttt{SumRatK}, for the original choice
and the original reference field, under its stated
\texttt{FieldOK}. It uses TT8 for \texttt{SumOKK}. No new
largeness threshold is needed, and no global period-integral
convergence or identity with (TT.17) is inferred from this
arithmetic calculation.

\paragraph{TT11: the character, root and phase portions of NumbersIn.}
There is one number field containing all local values of each
algebraic Hecke character used here. A direct argument avoids
assuming that infinitely many algebraic numbers generate a finite
extension. Choose an open unit subgroup on which the finite
character is trivial, and let $\mathfrak f$ be its modulus.
The relevant ray class group is finite. For a finite idele, divide
by a chosen representative of its ray class and a principal idele;
the remaining unit lies in that open subgroup. The principal
value is the inverse infinity-type monomial, hence belongs to the
normal closure of the number field defining the embeddings. The
finitely many chosen ray-class values are algebraic: raise each
to the order of its ray class, and use that same principal-value
formula. Adjoin these finitely many roots. This proves the claimed
single field. Finite-order characters are simpler: all their
values belong to one cyclotomic field.

Apply this to $\mu_0$, the grid characters $\chi_i$,
$\chi_\psi$, and the finite-order twists and discriminant
characters. The actual integral infinity exponents follow from
\texttt{LevCh.Is}; the half-twist of $\eta$ is kept in its
\texttt{ValuesIn}(1/2) convention, supplied by the same-HFrame construction HF16 and TL14. We do not adjoin $\sqrt q$ for infinitely
many rational primes. Instead, \texttt{NumbersIn.sqrt} concerns
only the finite union of AA and active local primes. Adjoin those
finitely many positive square roots, the eighth roots of unity,
and the roots of order $q_w^{e_w}$ at the active places. This
constructs the literal \texttt{CharsIn} and \texttt{RootCl}.

For phases, (TT.9) and $q_w\in\pi_w\mathcal O_w$ give
$\psi_w(z)^{q_w^{N_w}}=1$ on the required support. There are
finitely many twists, finite grid data, unfolding labels and active
places, by the actual finite quotient defining \texttt{G.Datum}
and TT6. Take the product of their positive orders and adjoin
its roots of unity. This proves \texttt{PhaseCl}, including its
two nonvanishing conditions. The fields just constructed depend
on the fixed actual data; their compositum is finite over
$\mathbb Q$. These are the actual \texttt{NFChars},
\texttt{NFRoot} and \texttt{NFPhase} suppliers, with zero
additional largeness thresholds.

\paragraph{TT12: the moved lattices and vacuum use the same rational frame.}
At a definite node let $B$ be the normal datum's rational frame,
$B^*JB=I$, let $u$ satisfy the actual lattice identity
$\widehat\Lambda=uB\widehat{\mathcal O}_E^n$, and let $x$
be the finite point. In $B$-coordinates the
actual moved lattice is
\[
 L_x=E^n\cap a\widehat{\mathcal O}_E^{\,n},\qquad
             a=B^{-1}(x u)B.                               \tag{TT.18}
\]
The order $xu$, and the inverse on the left of $B$, are forced by
$Bv\in\operatorname{globLat}(x,\Lambda)$: its defining condition
is $x^{-1}Bv\in\widehat\Lambda$, and
$\widehat\Lambda=uB\widehat{\mathcal O}_E^n$.
Both $x$ and $u$ preserve $J$, so direct multiplication gives
$a^*a=I$. Localization and the paired-prime duality of NB1 show
that each localization of (TT.18) is self-dual. The finite
intersection/localization argument LP2 then gives a full global
self-dual lattice. This proves the actual
\texttt{srcLat}/\texttt{tgtLat} statements with the source's
swap and sign at a row edge; it does not choose new frames.

On the completed Hom lattice, the action is
$X\mapsto g_T Xg_S^{-1}$, with the source's transpose convention
when written as row coordinates. Substitution of the two lattice
conditions proves inclusion in the new Hom lattice. Apply the
inverse action for the reverse inclusion. Consequently the actual
Weil image of the old vacuum has the new lattice translation and
half-translation laws, by Heisenberg covariance. At infinity the
acting point is the identity, so its annihilation operators commute
with this action; the Fock vacuum equations persist. Its norm is
one because the supplied Weil action is unitary for the same
fundamental-domain measure. These prove the full moved-vacuum
conditions, including normalization, from the given Weil datum.
They do not yet prove that the scalar comparing two normalized
vacuum vectors is algebraic.

\paragraph{TT13: normalized jets from finite arithmetic on the same vacuum.}
We isolate a narrower remaining arithmetic assertion than
\texttt{NFJet}. For every actual definite edge $p$ and finite
pair $g$ used in the class sums, require
\[
 W_p(g)\operatorname{vac}_p
       \in\operatorname{span}_{\overline{\mathbb Q}}
       \{\operatorname{Heis}(v)\operatorname{vac}_p:
                                      v_\infty=0\}.         \tag{TT.19}
\]
The span means a finite linear combination. It concerns this
particular $W_p$, its given splitting characters, and its given
unit vacuum. TT12 by itself does not imply (TT.19): an unspecified
complex unit scalar on a line need not be algebraic. We retain
(TT.19) exactly where it is used below.

The normalized base entire function $e_{0,p}$ in the normal datum
has the source's \texttt{ThetaVectorUpToUnits} description: an
algebraic unit times a finite positive intertwiner applied to the
product one-variable theta function in a rational orthonormal
basis. TA10, VI4 and VI9--VI10 imply that every rational
Heisenberg translate of this function has algebraic Taylor
coefficients. This is the direct finite-translate argument. It
does not require the stronger general assertion that an arbitrary
algebraic-coefficient entire function remains so after every
translation.

For clarity, finite-adelic translations have exactly the required
rational-coordinate interpretation. If $v_\infty=0$, choose
$n>0$ with $nv$ in the completed Hom lattice. Additive approximation
and CRT choose a rational vector in the prescribed finite coset.
Subtract it and a lattice vector from $v$; what remains is its
opposite rational archimedean coordinate. Heisenberg composition
and the vacuum's half-lattice law introduce values of the standard
character on rational finite adeles. Multiplication by an integer
puts each argument in the character's conductor lattice, so these
values are roots of unity. The exact Fock derivative formula is
therefore a root of unity times the derivative of one rational
translate of $e_{0,p}$, multiplied by the original period factors.
This also proves the source's finite-translation jet statement
without changing the normalized half-point or $\delta$.

A coset datum with $z=1$ is an actual finite algebraic Heisenberg
combination of the vacuum. To see this, take a common finite
lattice quotient on which its indicator and character are
constant, and apply the finite Fourier identity (TT.6). The
coefficients are rational multiples of roots of unity; the
standard theta-vector characterization identifies this finite sum
with $Z.\mathrm{vec}(d)$. Heisenberg covariance and (TT.19) give
the same conclusion after the actual finite point acts on that
vector. It follows that its raw jet table has the form
\[
 J(m)=c_{0,p}\left(\prod_\tau
        \operatorname{opp}_p(\operatorname{period}(\delta,\tau))
                      ^{d_\tau(m)}\right)A(m),
                   \qquad A(m)\in\overline{\mathbb Q}.     \tag{TT.20}
\]
Here $\operatorname{opp}_p$ is precisely identity or conjugation
for that oriented edge, and $d_\tau(m)$ is the total degree at
$\tau$. Both preserve algebraicity. No arbitrary field embedding
is interchanged with a limit, derivative or integral.

The jet operator is homogeneous of the specified Fock degree
$f_{p,\tau}$ at each place. Thus (TT.20) factors out
$c_{0,p}\prod_\tau\operatorname{opp}_p
(\operatorname{period}(\delta,\tau))^{f_{p,\tau}}$.
The actual scalar $Z.\mathrm{scal}(1,p)$ is its inverse.
All periods and $c_{0,p}$ are nonzero in the normal datum. The
factors cancel exactly. The remaining operator is obtained from
$A(m)$ by rational Fock differentiation, rational projection and
bideterminant coordinates (TT9). Hence every actual normalized
entry $Z.\mathrm{kk}(d,x,y)_{ab}$ is algebraic. Alternatively,
the literal \texttt{DictAt} formula supplies (TT.20) directly:
its $\kappa$, $r_i$ and rational translates are algebraic, and
its same $c_0,\delta$ cancel in precisely this calculation. The
extra norm, denominator and all-embedding estimates in
\texttt{DictAt} are not needed for this qualitative conclusion.

\paragraph{TT14: finitely many jet entries give one field.}
For the fixed choice, consider all twists, all finite grid data,
all nonzero unfolding labels, all classes, all top and row edges,
and all their bideterminant matrix indices. These are finite by
the actual finite datum quotient and TT6--TT8. For each label use
the least depth \texttt{expDepth}. TT7 shows that the pure data
appearing in \emph{any} expansion at that depth form a finite
set. Thus the set $E_{\rm jet}$ of all entries required by
\texttt{JetsIn} is finite, despite its universal quantifier over
expansions. Under (TT.19), or the stated qualitative dictionary,
TT13 shows every member of $E_{\rm jet}$ is algebraic. Then
\[
                  M_{\rm jet}=\mathbb Q(E_{\rm jet})        \tag{TT.21}
\]
is a number field and proves the literal \texttt{JetCl}. In
particular a same-data supplier of (TT.19) would close
\texttt{NFJet}; no uniform degree bound is claimed here. This
argument supplies the finite-generation step that pointwise
algebraicity alone does not supply.

\paragraph{TT15: algebraic compact averaging and one surviving degree.}
Let $C$ be a compact group with probability Haar measure, and
$\rho:C\to\mathrm{GL}(V)$ a continuous finite-dimensional
complex representation. Suppose a dense subset of $C$ has
matrices over $\overline{\mathbb Q}$ in the chosen basis. Then
\[
                         P=\int_C\rho(g)\,dg               \tag{TT.22}
\]
has algebraic entries. Here is an algebraic proof of this fact,
which does not presume that integrals of algebraic-valued
functions are algebraic. Put $F=\overline{\mathbb Q}$, and form
over $F$ the span $D_F$ of all columns of $\rho(\gamma)-I$ and
the common kernel $I_F$ of those matrices, for the dense set.
Finite-dimensionality reduces both constructions to finitely many
matrices. Their complex scalar extensions are respectively
$D=\operatorname{span}_{\mathbb C}(\rho(g)-I)V$ and $V^C$:
for $D$, its finite-dimensional span is closed, so density puts
all $g$ in it; for $V^C$, use continuity of $g\mapsto\rho(g)v$.
Haar invariance shows that $P$ is the identity on $V^C$ and zero
on $D$. Also $v-Pv$ lies in $D$, by integrating in this closed
finite-dimensional subspace. Therefore
$V=V^C\oplus D$. The same direct-sum decomposition holds over
$F$, by ranks of finite matrices. Its projection has coefficients
in $F$, proving (TT.22). The zero-dimensional case is included.

For a supplied actual bottom triangle, the real group on the
source side is a product of compact unitary groups. In its
standard frames the Fock vacuum character is
\[
            \chi(g)=\prod_\tau\det(g_\tau)^{\,N-k_\tau},
       \qquad b_\tau=k_\tau-N/2,\quad k_\tau\in\mathbb Z,
                                                                  \tag{TT.23}
\]
where $N$ is the target rank and $b_\tau$ is the infinity
exponent of the specified splitting-character ratio. The parity
restriction of that ratio gives the displayed integrality.
The sign is the negative source-space sign: the positive
metaplectic half-trace and the splitting-character exponent
combine to $N/2-b_\tau=N-k_\tau$. More explicitly, diagonalize a
unitary source matrix with eigenvalues $e^{\mathrm i\theta_j}$.
For one oscillator,
$H=(-d^2/dx^2+x^2)/2$ satisfies
$H(e^{-x^2/2})=\tfrac12e^{-x^2/2}$, since the second derivative
is $(x^2-1)e^{-x^2/2}$. In this negative-source polarization each
of the $N$ Gaussian coordinates in source direction $j$ therefore
contributes the rotation factor $e^{\mathrm i\theta_j/2}$; the
prescribed splitting character contributes
$e^{-\mathrm i b_\tau\theta_j}$.
Their product is $e^{\mathrm i(N/2-b_\tau)\sum_j\theta_j}$.
The vacuum is unchanged by a unitary change of the diagonalizing
basis, and a creation monomial is carried to the corresponding
linear substitution. The integer exponent makes the result
independent of the angles chosen. This gives (TT.23) on the
whole compact unitary group, including its special-unitary part. A
finite-order twist is trivial on its connected real subgroup, so
the same $\chi$ works for every twist $t$.

On monomials with place-degree vector $d$, the resulting finite
matrix representation is $\chi(g)\operatorname{subst}_d(g)$.
Algebraic-entry unitary matrices are dense: use the Cayley
transform on skew-Hermitian matrices with algebraic entries,
which is dense in the open chart without eigenvalue $-1$.
Its matrices are algebraic, including negative determinant
powers in (TT.23). Thus TT15's first argument applies to every
degree block. At a central scalar $z$ in one real place, the
representation is multiplication by $\chi(z)z^{d_\tau}$.
Choose $z=e^{\mathrm i}$, which has infinite order. Haar
averaging vanishes unless
$\chi(z)z^{d_\tau}=1$ at every place. There is at most one
nonnegative degree vector $d_0$ satisfying these equations. If
there is none, all averages are zero and any fixed $d_0$ may be
used with zero coefficients. Otherwise every averaged kernel is
an algebraic linear combination of the finitely many kernels of
degree $d_0$. This is the actual fixed-degree conclusion used in
\texttt{avg\_\allowbreak kernel\_\allowbreak alg}, on the same vacuum and the same
probability Haar measure.

\paragraph{TT16: normalized bottom periods and their finite field.}
Fix a literal \texttt{IsBottomTri} with its actual Haar measures
and finite-order twists. Retain the following finite-Weil input:
for both source sides, its standard theta datum and its images
under finite adelic pairs are finite algebraic combinations of
finite Heisenberg translates of the corresponding fixed vacuum,
as in (TT.19). This is the content whose source producer is
\texttt{moved\_\allowbreak translate\_\allowbreak algComb}; its existence is not a
consequence of compact averaging alone.

For a fixed place degree $d$, VI9--VI11 and the finite-translation
argument of TT13 give a common nonzero period
\[
       \Omega_d=\overline{c_0\prod_\tau
                      \operatorname{period}(\delta,\tau)^{d_\tau}}
                                                                  \tag{TT.24}
\]
for the negatively oriented source; use the corresponding
unconjugated expression on the other source side. The unspecified
unit relating the actual normalized vacuum to its entire model
is included in $c_0$, so it is not asserted algebraic. Nonzero
vacuum norm and nonzero periods imply $\Omega_d\ne0$.
The retained finite-Weil assertion and the explicit finite
Fourier expansion of the standard theta datum show that every
finite-point monomial kernel of degree $d$, divided by
$\Omega_d$, is algebraic.

There is a common compact open finite level for these finitely
many datum functions and twists. Lattice stabilizers preserve the
vacuum line, and their action on it is a continuous character.
A continuous character of a profinite group into the circle has
finite image: its compact image is a closed subgroup of the
circle, and a quotient of a profinite group is totally
disconnected; the only such closed subgroups of the circle are
finite. Shrink the congruence subgroup to kill this character,
and to preserve the finitely many Heisenberg translations and
their conductor phases. This supplies a level fixing all the
finite datum functions. Polynomial creation operators commute
with finite-place action, so the level does not depend on the
monomial.

Let $z_l$ be the finite set of definite adelic class
representatives, let $w_l>0$ be their finite stabilizer orders,
and put $v=\mu(J)>0$, where $J$ is the product of that finite
level and the full compact real group. With probability Haar
measure $dg$ on the latter, the exact quotient formula is
\[
 I_t(P,h)=v\sum_l\frac1{w_l}
                    \int_C\operatorname{ker}_t(P,h)(z_lg)\,dg.
                                                                  \tag{TT.25}
\]
It follows by decomposing the adelic quotient into these classes;
the finite stabilizer makes each fiber a quotient of multiplicity
$w_l$, giving its reciprocal. Compactness gives convergence for
each continuous kernel. The number $v$ is the actual measure
normalization and need not be rational. TT15 reduces the integral
to the single degree $d_0$ with algebraic coefficients. Thus
\[
                         I_t(P,h)/(v\Omega_{d_0})
                                      \in\overline{\mathbb Q}     \tag{TT.26}
\]
for monomials $P$ and finite points $h$. This produces one
nonzero period for all twists and all degrees, including the
case where every period integral is zero.

We must still pass from algebraicity to a number field. Left
rational invariance, the common finite level, and the finite
definite class set reduce $h$ to finitely many representatives,
up to an archimedean point inverse to the archimedean component
of a rational unitary element. Its matrix coefficients lie in
the normal closure of $E$. Fock substitution and the integer
determinant powers in (TT.23) therefore have coefficients in one
fixed number field. Only finitely many monomials can have a
nonzero integral: central real scalars force the degree vector
of TT15, or the same argument applied before averaging for each
of the finitely many twists. Consequently all integrals are
linear combinations, over a fixed number field $M_0$, of finitely
many values $I_t(P_b,h_a)$. Adjoin their normalized algebraic
values from (TT.26). The resulting finite extension contains
all $I_t(P,h)/(v\Omega_{d_0})$. Repeat independently for $I'_t$;
its period and its Haar volume may differ. The compositum proves
the exact \texttt{SumsIn} assertion.

This proves the substantive definite-period implication from the
specified finite-Weil arithmetic. To obtain \texttt{NFBot} for
all squares, one also needs an actual bottom triangle at each
square. The source builds standard positive spaces of ranks
\texttt{botBig}, \texttt{botSmall}, and their sum, then chooses
Weil representations with its displayed parity-twisted characters.
The lattice and rank construction is elementary; the existence
of those additional full adelic Weil representations is retained
here. Their existence is not supplied merely by the grid's
already chosen $W$.

\paragraph{TT17: common rational periods force the bottom scalar ratio.}
We prove the precise algebraic comparison used after TT16. Suppose
there is one complex vector space $\Pi$, complex-linear maps
$p:P\twoheadrightarrow\Pi$, $p':P'\to\Pi$ with nonzero image,
and injective maps $J_t:\Pi\to\{\text{functions of }g\}$ such
that
\[
 I_t(P)(g)=\lambda_tJ_t(pP)(g),\qquad
 I'_t(P')(g)=J_t(p'P')(g),\qquad\lambda_t\ne0.              \tag{TT.27}
\]
These are the substantive factorization hypotheses; they are
kept on the same bottom triangle. Assume its monomial periods
are rational over $K$ after division by nonzero $\Omega,\Omega'$,
and finite-point vanishing of a difference in (TT.27) implies
vanishing at every point. Then
\[
                       \lambda_t/\lambda_{t_0}\in K^\times.\tag{TT.28}
\]

Choose $P'$ with $p'P'\ne0$ and $P$ with $pP=p'P'$. There is a
finite point $g_0$ where $I'_{t_0}(P')(g_0)\ne0$, by injectivity
and the finite-point determination hypothesis. Choose a
$K$-linear functional $\varphi:\mathbb C\to K$ nonzero on this
value, by extending a basis of its one-dimensional $K$-span.
No continuity of $\varphi$ is assumed or needed. For a polynomial
$Q=\sum a_mX^m$, put
$\operatorname{rat}_{\varphi,w}(Q)=\sum\varphi(wa_m)X^m$,
viewing its coefficients in $\mathbb C$ through $K\subset\mathbb C$.
If a complex-linear functional $\ell$ has
$\ell(X^m)\in\Omega K$ for every $m$, direct finite expansion gives
\[
 \ell(\operatorname{rat}_{\varphi,w\Omega}(Q))
                         =\Omega\varphi(w\ell(Q)).          \tag{TT.29}
\]
Set
$Q=\operatorname{rat}_{\varphi,\lambda_{t_0}^{-1}\Omega}(P)$
and $Q'=\operatorname{rat}_{\varphi,\Omega'}(P')$.
At $t_0$, (TT.29), finite-point determination and injectivity yield
\[
             \Omega'\lambda_{t_0}pQ=\Omega p'Q',
                              \qquad p'Q'\ne0.              \tag{TT.30}
\]
The second assertion follows by evaluating at $g_0$ and using
$\varphi(I'_{t_0}(P')(g_0))\ne0$.
For each $t$ choose a finite point with $I'_t(Q')\ne0$ and write
$I_t(Q)=\Omega z$, $I'_t(Q')=\Omega'z'$ with $z,z'\in K$,
$z'\ne0$. Apply $J_t$ to (TT.30) and multiply by $\lambda_t$.
Cancelling the nonzero periods gives
$\lambda_{t_0}z=\lambda_tz'$, hence (TT.28).

For the actual definite continuous automorphic kernels,
finite-point determination has a direct justification. Rational
unitary points are dense in the compact real group: the Cayley
chart is rational on the vector space of skew-Hermitian matrices,
and weak approximation approximates their entries at all real
places. For a fixed finite component $g_f$ and rational points
$\gamma_n$ with $\gamma_{n,\infty}\to g_\infty$, left
invariance gives
$F(\gamma_{n,\infty},g_f)=F(1,\gamma_{n,f}^{-1}g_f)=0$.
Continuity gives $F(g_\infty,g_f)=0$. No convergence of the finite
components is required. Compactness of the definite quotient
also gives the integrability needed to use the monomial period
identities. Thus TT17 proves the actual
\texttt{Factor.ratIn}/\texttt{ratIn\_\allowbreak of\_\allowbreak definite} implication;
it does not construct the common factorization $\Pi,p,p',J_t$.

\paragraph{TT18: the exact arithmetic consequence of a matching tower.}
Fix a square $(i,j)$ and retain its actual nonzero square
constants $c_t$, bottom factorization (TT.27), and the source's
tower identity. Write $p=\texttt{pp}(j)$,
$q=\texttt{qq}(j)$, $e=\texttt{eTop}(i,j)$, and let
$\Lambda_t$ be the completed Hecke functions of the actual
characters. The identity is
\[
 \frac{c_t}{c_{t_0}}\prod_{a=0}^{q-1}
       \frac{\Lambda_t(e+a+1)}{\Lambda_{t_0}(e+a+1)}
 =\mathrm{embLL}(\zeta)\left(\prod_{b<p+q}u_b\right)
       \frac{\lambda_t}{\lambda_{t_0}}
       \prod_{a=0}^{p-1}
       \frac{\Lambda_t(e+q-a)}{\Lambda_{t_0}(e+q-a)}.        \tag{TT.31}
\]
Every $u_b$ is the literal \texttt{StepRatio} with final argument
one, and hence lies in $K$ by its displayed coefficient-field
factorization. The away-unit $\zeta$ is an element of the same
coefficient field. Under \texttt{NumbersIn.ratios}, all displayed
$L$-ratios lie in $K$. The exact tower-number inequalities put
their arguments between $1$ and \texttt{halfRank}, and hence
within the upper bound \texttt{bigN}+1 required by that field
clause. The supplied \texttt{LNe} is stronger: it gives
nonvanishing at every integer from $1-\texttt{e0}$ through
\texttt{halfRank}. It therefore makes every denominator and
every factor on the left of (TT.31) nonzero.

Use the bottom triangle selected by \texttt{NumbersIn.bottom},
not an unrelated one. The separate existence input supplies a
witness that this selection is allowed; its factorization is
assumed for every valid bottom triangle. TT17 then gives
$\lambda_t/\lambda_{t_0}\in K^\times$ for the selected one.
Divide (TT.31) by its nonzero left product. This proves the actual
$\texttt{sqRatio}(i,j,t)\in K$. Taking the maxima of the original
constant, tower and nonvanishing thresholds gives the literal
\texttt{RatioK} implication, with all original inputs retained.
In particular a merely algebraic numerator is insufficient, and
qualitative nonvanishing at $s=1$ for an unrelated finite-order
character does not supply this critical interval.

\paragraph{TT19: a finite coefficient field and its same-choice boundary.}
For a fixed actual choice, TT11 supplies its character, root,
square-root and phase fields. Under the exact finite-Weil input
(TT.19), TT14 supplies its jet field. Under the bottom existence
and finite-Weil inputs of TT16, it supplies its bottom period
fields. If the finitely indexed critical ratios in
\texttt{LRatiosIn} are algebraic, adjoining them supplies the
remaining ratio field. The functions $\Lambda$ in that clause
are not separately arbitrary choices: the characterization
\texttt{IsHeckeL} fixes their completed functions; its uniqueness
identifies any two witnesses. There are only finitely many twists,
squares and integer arguments $1\le s\le\texttt{bigN}+1$.
Thus adjoining these finitely many algebraic values is legitimate.
For arguments with zero reference denominator the clause imposes
no requirement. No algebraicity of these critical values is
proved by adjoining formal symbols.

The compositum of these finitely many number fields is finite and
satisfies all eight fields of \texttt{NumbersIn}: characters,
eighth roots, local depth roots, phases, the specified square
roots, jets, critical ratios, and bottom sums. Field inclusion
preserves each clause. For the existential bottom clause use the
same chosen bottom triangle and the same two periods after
inclusion. For jets use the same $Z.\mathrm{kk}$ entries and
least-depth expansions. For phases use the same local integral,
frames and representatives. This proves the actual finite-field
join, rather than constructing eight unrelated choices.

To conclude literal \texttt{FieldOK}, this compositum must be
embedded in the range of the \emph{same} \texttt{embLL} used by
$R$ and $\mathcal C$. If that embedding is supplied, all eight
clauses transfer immediately and TT10 applies. If the coefficient
field is to be enlarged, the source's \texttt{LAgree},
\texttt{GAgree}, \texttt{WAgree}, \texttt{SAgree} and flip
agreements must be provided. They preserve the preliminary datum,
the actual HFrame and reference twist, every space/lattice and
frame, the character functions, theta realization, finite point
and jet data, and the actual local Fourier integrals. Merely
knowing that a finite extension exists does not prove these
same-object agreements or construct \texttt{Choice.Is}; in
particular \texttt{RealOK} is the additional separation condition
in that record. The producers \texttt{exField\_\allowbreak of\_\allowbreak keyFin\_\allowbreak exKey}
and \texttt{dat\_\allowbreak choice\_\allowbreak exists} remain conditional on these
compatible supplies.

The remaining primitive analytic tasks are consequently precise.
\texttt{LinK} still needs convergence of the actual iterated
cuspidal theta/Bessel tests; \texttt{InjK} still needs their
separation theorem. \texttt{Formula2K} still needs the same-data
projection, local square-flip matching, absolutely justified
unfolding, and the final period-to-jet identity. TT8 and TT10
supply its finite-sum side, not these identities. \texttt{RatioK}
still needs the tower and common bottom factorization and its
critical nonvanishing, as displayed in TT17--TT18. The stronger
\texttt{ThmBop}/\texttt{ThmAop} and
\texttt{PartsE}/\texttt{PartsW} include positive scalar operator
identities and spanning assertions; none is inferred here from
rationality or finite support. The source ledger below locates
actual producers for these tasks, and records which smaller
arithmetic subclaims have been discharged in this section.


\section{The quantitative mass window and its prime-counting boundary}
\label{sec:quantitative-windows}

\status{QW1--QW9 prove the literal conditional \texttt{window\_\allowbreak lower}
from \texttt{SiegelMassRatioCM}, including the Euler factors, both local
indices, the Vandermonde comparison and the uniform passage to the actual
host. QW10--QW13 prove the weighted-to-unweighted norm-prime reduction
and the two scale consequences of the explicit short-interval input.
QW14--QW15 prove classical nonvanishing at one and the finite-conductor
real-zero margin. The global mass-ratio theorem, the compatible arithmetic
\texttt{LFam}, and its uniform weighted short-interval lower bound remain
separate inputs. No arbitrary short interval is supplied by qualitative
prime existence or by the finite-conductor margin.}

In this section $B=e^L$, $L=x^\epsilon$, $a=xB$, and $r=k+m-1$.
The letter $B$ here is the source scale \texttt{s.b}; in QW10--QW13
the distinct symbol $\mathcal B=4dM^2$ denotes the analytic conductor
bound. The finite host has $[E:\mathbb Q]=2f$. All genus masses, local
indices and Harish--Chandra coordinates are the literal ones on the
same supplied host and preliminary datum.

\paragraph{QW1: absolute logarithmic control of the actual Euler product.}
Let $E$ satisfy the source \texttt{IsCMUnram} condition and put
$f=[E:\mathbb Q]/2$. For a maximal ideal $\mathfrak p$ above a rational
prime $p$, write $N\mathfrak p=p^{d_{\mathfrak p}}$. The actual factor is
\[
 E_{\mathfrak p}(M)=
 \begin{cases}
 (1-(-1)^M(N\mathfrak p)^{-M/2})^{-1},&
       \overline{\mathfrak p}=\mathfrak p,\\
 (1-(N\mathfrak p)^{-M})^{-1/2},&
       \overline{\mathfrak p}\ne\mathfrak p.
 \end{cases}                                                   \tag{QW.1}
\]
If $\overline{\mathfrak p}=\mathfrak p$, the unramified CM condition
provides an integer $z$ with $z-\bar z\notin\mathfrak p$. Thus residue
conjugation is nontrivial. It would be the identity on a residue field
$\mathbb F_p$, so $d_{\mathfrak p}\ge2$. In all cases, using
$\log(1+t)\le-\log(1-t)$ for $0\le t<1$,
\[
 E_{\mathfrak p}(M)>0,\qquad
 |\log E_{\mathfrak p}(M)|
 \le\frac{d_{\mathfrak p}}2[-\log(1-p^{-M})].                    \tag{QW.2}
\]
Indeed the fixed-prime factor has absolute logarithm at most
$-\log(1-p^{-M})$; the other has at most half this bound.
The ideal factorization identity
$\sum_{\mathfrak p\mid p}e_{\mathfrak p}d_{\mathfrak p}=2f$ gives
$\sum_{\mathfrak p\mid p}d_{\mathfrak p}/2\le f$.
For $M\ge2$, unique factorization of positive integers gives
\[
 \sum_p-\log(1-p^{-M})=\log\zeta(M)<\infty.
\]
This last equality can be obtained by increasing finite Euler products:
their positive expansions enumerate the integers supported on those
primes, and their limit is $\sum_{n\ge1}n^{-M}$.
Therefore the sum of the absolute logarithms in (QW.2) converges.
Exponentiating its signed sum constructs the actual positive product
$L_E(M)=\prod_{\mathfrak p}E_{\mathfrak p}(M)$ and proves
\[
               |\log L_E(M)|\le f\log\zeta(M)\le f.             \tag{QW.3}
\]
For the last bound, the integral test gives
$\zeta(M)\le\zeta(2)\le2$, and $\log2\le1$.
This establishes the source Euler-product existence and logarithmic
bound; it makes no assumption about zeros of an analytically continued
$L$-function.

\paragraph{QW2: count the two local level groups before taking their ratio.}
For a finite field of cardinality $q\ge3$, choosing successive independent
columns gives
\[
 |\mathrm{GL}_n(\mathbb F_q)|=q^{n^2}\theta_q(n),\qquad
 \theta_q(n)=\prod_{j=1}^n(1-q^{-j}),\qquad
                 \tfrac12\le\theta_q(n)\le1.                   \tag{QW.4}
\]
The lower bound follows from
$\prod(1-t_j)\ge1-\sum t_j$ for $t_j\in[0,1]$ and
$\sum_{j\ge1}q^{-j}=1/(q-1)\le1/2$.
The column-level unitriangular group has $k(k-1)/2$ independent entries.
For the row level put $s=m/2-1$, with the actual even $m\ge2$ and
$r=k+m-1$. Its free entries comprise the $k$-by-$k$ unitriangular
part and the two outer blocks, giving exactly
\[
       s(2k+1+2s)+\frac{k(k-1)}2
       =\frac{r^2-(k+1)^2-2s}2+\frac{k(k-1)}2.                  \tag{QW.5}
\]
Every choice of these entries gives an upper triangular matrix with
diagonal one. Conversely the defining zero pattern leaves precisely
these entries free. Its cardinality is therefore the corresponding
power of $q$; no quotient by an additional determinant character occurs.
Dividing (QW.4) by these two subgroup cardinalities proves
\[
 \tfrac12q^{k(k+1)/2}\le I_{\rm col}(q)\le q^{k(k+1)/2},\qquad
 \tfrac12q^{(r^2+3k+1+2s)/2}\le I_{\rm row}(q)
                   \le q^{(r^2+3k+1+2s)/2}.                   \tag{QW.6}
\]
The exponents are integers for the actual parities. Their difference
is $m(2k+m-1)/2$. At the two host primes $g_0,g_1$, each with $f$
copies, (QW.6) gives positive indices and
\[
 \frac{I_{\rm row}}{I_{\rm col}}
 \ge2^{-2f}(g_0g_1)^{fm(2k+m-1)/2}.                            \tag{QW.7}
\]
This is the original \texttt{idx\_\allowbreak place}/\texttt{idx\_\allowbreak ratio} statement
for the actual level groups, including its factor $2^{-2f}$.

\paragraph{QW3: telescope the stated mass hypothesis without strengthening it.}
Write $\mu_N$ for the genus mass of the standard rank-$N$ lattice.
IG9 proves its positivity directly from the finite positive mass sum.
The retained \texttt{SiegelMassRatioCM} hypothesis says, for every
unramified CM field and every $N\ge3$,
\[
 \frac{\mu_N}{\mu_{N-1}}
 =R_N:=\frac{\big((N-1)!(v/\pi)^N\big)^f}{\sqrt{(2v)^f}}L_E(N),
 \qquad v=\frac{|\operatorname{disc}E|^{1/(2f)}}2.               \tag{QW.8}
\]
For $k\ge2$ and $r=k+m-1$, multiply these identities, cancelling the
positive intermediate masses, to obtain
$\mu_r=\mu_k\prod_{N=k+1}^rR_N$. If $m=1$ the product is empty and
the identity is $\mu_k=\mu_k$. In the actual host, the identities
$f\mathrm{Deg}(E)=f$ and $v\mathrm{Inv}(E)=v$ follow from
$[E:\mathbb Q]=2f$ and total complexity, so (QW.8) has exactly the
normalization required by \texttt{mass\_\allowbreak tel}.

This paragraph proves the telescope from the displayed ratio law. The
ratio law itself is not a consequence of positivity. Its pinned
producer passes through \texttt{HermRepUnram} and \texttt{LIStat},
the representation and local-density mass identities. Those are
retained source targets.

\paragraph{QW4: an elementary Vandermonde comparison for interlacing columns.}
For a decreasing $N$-tuple $z$ put
$V_N(z)=\prod_{0\le i<j<N}(z_i-z_j)$ and
$d_N(z)=V_N(z)/\prod_{i=0}^{N-1}i!$.
Suppose $k\ge2$, $v_i\ge0$, $v_i-v_{i+1}=4\Delta$ with
$\Delta\ge1$, and
$u_i>v_i>u_{i+1}$. Let $b_0<\cdots<b_{m-1}$ be positive with
successive gaps at least one, and append $-b_0,\ldots,-b_{m-1}$ to $v$.
Every cross factor $v_i+b_j$ is at least $b_j$ and every bottom factor
$b_j-b_i$ is at least one. Hence
\[
 V_r(v,-b)\ge V_{k-1}(v)\prod_{j=0}^{m-1}b_j^{k-1}.             \tag{QW.9}
\]
If every gap in $u$ is at most $D\ge1$, the $2k-3$ pairs involving
its first or last entry contribute at most $D^{2k-3}$. Every remaining
pair satisfies
$u_i-u_j\le v_{i-1}-v_j$ for $1\le i<j\le k-2$.
These are exactly the nonadjacent pairs of $v$. The omitted adjacent
factors of $V_{k-1}(v)$ are at least one, so
\[
                      V_k(u)\le D^{2k-3}V_{k-1}(v).             \tag{QW.10}
\]
The proof covers $k=2$ by empty products. Thus this comparison loses
only $2k-3$ endpoint factors, rather than the full number of pairs.

For the literal \texttt{PlaceData}, $1\le x$, $a=xB$, and
$v_i\le4k\Delta$. Either $u_0\le2a$, $u_{k-1}\ge-a$, or
$u_0\le12k\Delta$, $u_{k-1}=v_{k-2}-1/2\ge-1/2$.
The former gives total gap at most $3a$. In the latter use
$4k\Delta\le b_0/x$ and $b_0<w_0\le w_{m-1}\le2B$ to bound it by
$6B+1/2\le8a$. Taking $D=8a$ in (QW.9)--(QW.10) proves, in particular,
\[
 \log V_r(v,-b)-\log V_k(u)
 \ge(k-1)\sum_j\log b_j-2k\log(8a)-2k(1+\log k).              \tag{QW.11}
\]
All Vandermondes are positive. The last extra nonpositive term merely
weakens our bound to the form used by the pinned source.

\paragraph{QW5: pass to the middle weights with an explicit constant.}
The remaining \texttt{PlaceData} clauses say
$b_j<w_j<b_{j+1}$ where defined, $w_{m-1}\le2B$, and
$4k\Delta\le b_0/x$. Thus $b_0\ge1$ and
\[
 \sum_j\log(w_j/b_j)
 \le\log(w_{m-1}/b_0)\le1+L,\qquad
                     \sum_j\log b_j\le m(1+L).               \tag{QW.12}
\]
The first inequality telescopes, including the case $m=1$.
For $\epsilon>0$, $L=x^\epsilon\ge1$ and $k\le x^{C_0}$,
\[
 \log x\le L/\epsilon,\quad
 \log k\le C_0L/\epsilon,\quad
 \log(8a)\le7+L/\epsilon+L.
\]
Insert (QW.12) into (QW.11). The loss is at most
\[
 k(1+L)+m(1+L)+2k(7+L/\epsilon+L)
                 +2k(1+C_0L/\epsilon).
\]
Since $L\ge1$, this proves the exact useful one-place bound
\[
 k\sum_j\log w_j\le\log V_r(v,-b)-\log V_k(u)
       +kL\left(20+\frac{2+2C_0}{\epsilon}\right)+2mL.          \tag{QW.13}
\]
No asymptotic estimate for individual gamma factors is hidden here.

\paragraph{QW6: factorial cancellation gives the global logarithmic ratio.}
Suppose $f,m\ge1$, $k\ge2$, and every one of the $f$ places satisfies
the same \texttt{PlaceData}. Assume $v\ge\pi$, $\log v\le L$,
$g_0,g_1\ge1$, the positive masses and indices, (QW.3), (QW.7),
and the mass telescope (QW.8). Retain the exact hypotheses
\[
 d_* =\mu_k I_{\rm col}\prod_\rho d_k(u_\rho),\qquad
 D_*\ge(2/3)^f\mu_r I_{\rm row}\prod_\rho d_r(v_\rho,-b_\rho).
                                                                    \tag{QW.14}
\]
These already give $d_*,D_*>0$. In the logarithm of the ratio, the
factorial terms cancel exactly, because
\[
 \sum_{N=k+1}^r\log((N-1)!)+\sum_{i=0}^{k-1}\log(i!)
                         =\sum_{i=0}^{r-1}\log(i!).            \tag{QW.15}
\]
Put $\ell=\log(v/\pi)\in[0,L]$, $G=\log g_0+\log g_1$,
$T=\sum_{\rho,j}\log w_{\rho j}$, and $K=(2+2C_0)/\epsilon$.
The remaining mass contribution is bounded below by
\[
 fk(m-1)\ell-\tfrac f2(m-1)(1+L)-f(m-1),
\]
since $\sum_{N=k+1}^rN\ge k(m-1)$, $\log2\le1$, and (QW.3).
The index contribution is at least $fkmG-2f\log2$; the Vandermonde
contribution is at least $kT-fkL(20+K)-2fmL$ by (QW.13).
Finally $f\log(2/3)\ge-f/2$.
Relative to the desired main term $kT+fkm\ell+fkmG$, the total loss
divided by $f$ is no more than
\[
 kL(20+K)+2mL+k\ell+\tfrac{m-1}{2}(1+L)
                     +(m-1)+2\log2+\tfrac12.
\]
Using $L\ge1$, $m\ge1$, $k\ge2$, this is at most
$(30+K)(kL+m^2L)$. Consequently the literal
\texttt{LogRatioLowerBound} holds with the explicit choice
\[
 \boxed{C_\epsilon=30+(2+2C_0)/\epsilon},                       \tag{QW.16}
\]
and conclusion
\[
 \frac{\log(D_*/d_*)}{fk}\ge
 \frac1f\sum_{\rho,j}\log(w_{\rho j}v/\pi)
 +m(\log g_0+\log g_1)-C_\epsilon\left(L+\frac{m^2L}{k}\right).
                                                                    \tag{QW.17}
\]
The restrictions $0<\epsilon<1/52$ and $C_0\ge0$ of the source are
preserved, although this calculation uses only $\epsilon>0$ and the
displayed size assumptions.

\paragraph{QW7: verify PlaceData on the actual top, thick and thin places.}
The scale and rank estimates RF11--RF14 were proved with thresholds
independent of the eventual datum. Fix $c_D,c_h\ge0$ and set
$\Delta=\lceil x^{c_D}\rceil$, $H_0=2\lceil x^{c_h}\rceil$.
They give, uniformly in the actual rank record,
\[
 x(4k\Delta+H_0)\le w''\le B',\quad 3B'\le B,\quad
 mH_0+8f^2+2\le B,\quad k\le x^5.                             \tag{QW.18}
\]
For completeness, $\Delta\le2x^{c_D}$, $H_0\le4x^{c_h}$,
$k\le x^5$, $f\le x^3$, $m\le x$ bound each polynomial on the left.
RF11 makes $w''$ and $B$ exceed any prescribed such polynomial,
uniformly for $1\le n\le x^3$, and gives $B'/B=x/X'\to0$.
The actual bottom-window length is at most
$2n x^{-\epsilon/8}L\Lambda B'/k\le B'$. Indeed RF14 gives
$\Lambda\le C L(\log x)^2$. Use also the sharper original rank
inequalities $n\le2x^{2+8\epsilon}$ and
$k\ge x^{9/2+10\epsilon}$; the coefficient is then at most
$C' x^{-5/2-\epsilon/8}(\log x)^2$, which tends to zero.

Write the actual bottom start at a place as $\beta$. The \texttt{HCVal}
equalities are
\[
 v_i=(k-2-i)4\Delta,\qquad b_j=\beta+jH_0,\qquad
                         w_j=b_j+H_0/2.                       \tag{QW.19}
\]
At a top place, $\beta=b_T+\lfloor B'/x\rfloor$ and the original
window plus the preceding length bound puts $B'\le b_T\le2B'$.
At a thick place, $\lceil B\rceil\le\beta<\lceil B\rceil+
|\mu(E)|$. At a thin place, $\beta=\lceil w''\rceil$.
The source root-of-unity bound is
$|\mu(E)|\le2[E:\mathbb Q]^2=8f^2$: if its order is $q$, then
$\varphi(q)\le[E:\mathbb Q]$ and $q\le2\varphi(q)^2$.
The latter follows from the prime-power formula for $\varphi$:
every prime-power contribution to $\varphi(q)^2/q$ is at least one
except a single possible factor $1/2$ from $2^1$.
These three cases therefore give
\[
                   w''\le\beta\le B+8f^2+2.                  \tag{QW.20}
\]
Together (QW.18)--(QW.20) imply
$x(4k\Delta+H_0)\le\beta$ and $\beta+mH_0\le2B$.

The existing \texttt{HCParam} gaps give the strict interlacing of $u,v$,
$-b_0<u_{k-1}$, $b_j<w_j<b_{j+1}$ and $b_{j+1}-b_j\ge1$.
Equation (QW.19) gives $|w_j-b_j|=H_0/2\le b_j/x$,
$4k\Delta\le b_0/x$, $w_{m-1}\le2B$, and $w_j\ge1$.
At a top place the prescribed endpoints give $u_0\le2a$ and
$u_{k-1}\ge-a$; the small endpoint correction is at most
$a/(k x^{A/2})\le a/2$. At another place the exact endpoints are
$u_0=v_0+4k\Delta+1/2$, $u_{k-1}=-1/2=v_{k-2}-1/2$.
Since $4k\Delta\ge1$, they meet the alternative endpoint inequalities.
This checks all nine original \texttt{PlaceData} clauses and $w_j\ge1$
on the same datum, at every word and every place.

\paragraph{QW8: the host constants and rank losses are uniformly small.}
The actual host discriminant inequalities give
\[
 \operatorname{rd}(E)\ge
 2^{e+1-2^{-e}}\sqrt d\ge4\sqrt3\ge32/5,
 \qquad v=\operatorname{rd}(E)/2\ge16/5>\pi,
\]
using $e\ge2$ and $d\ge3$. The other host inequality and its supplied
prime-window bounds give
\[
 \operatorname{rd}(E)\le(Q_S/2)\sqrt{q d}\le Q_S qd,
 \quad \log(Q_S qd)\le2x^{\epsilon/5}(\log x)^2\le L
\]
eventually. The last step is RF11 with positive exponent gap
$4\epsilon/5$. Hence $\log v\le L$.
Also the same \texttt{HostVal} has $g_i\ge(kB)^8$, so
$\log g_i\ge8L$ and $g_i\ge3$ for large $x$.

For any fixed real $C,C_\Gamma$, RF12 gives, eventually,
\[
      2|C|\le m,\qquad m^2\le k,\qquad 2C_\Gamma m^4\le k,
      \qquad L\ge1.                                           \tag{QW.21}
\]
One explicit way to see the margins is
$m\ge x^{1/4}$, $m\le x$, $k\ge x^{9/2+10\epsilon}$.
After $2|C_\Gamma|\le x^{1/4}$, the third left side is at most
$x^{17/4}\le k$; the other inequalities follow at once.
All constants were fixed before this final threshold. Neither the word,
the host, nor the chosen coordinates affect it.

\paragraph{QW9: prove the original quantitative window with its quantifier order.}
Fix \texttt{SiegelMassRatioCM}, $0<\epsilon<1/52$, a standard constant
record $\mathcal C$, and $c_D\ge0$. Choose $c_2=0$. For each
$c_h\ge c_2$, take a common $x_0$ for QW7--QW8, with
$C=C_\epsilon=30+12/\epsilon$ in (QW.21). Now let any scale $s$,
any original \texttt{PreCh} $P$, and the remaining $q_1,a,b,c,du,D_s$
be supplied with $s.\epsilon=\epsilon$, $x\ge x_0$, the literal
\texttt{ChOK}, and the two exact size hypotheses (QW.14).
For each word, QW1--QW8 apply with $C_0=5$.
The rows' $v,b$ coordinates agree with those of the base word by the
original \texttt{hcW\_\allowbreak v}/\texttt{hcW\_\allowbreak b} identities; thus the single
given lower bound for $D_s$ is exactly the one needed for every word.
Enumerating $\mathrm{Tau}$ replaces its product by the $f$ place product
without changing any factor.

The average logarithm in (QW.17) is nonnegative, since $w_j\ge1$ and
$v\ge\pi$. Moreover
\[
 m(\log g_0+\log g_1)\ge16mL,\qquad
 C_\epsilon(L+m^2L/k)\le2C_\epsilon L\le mL.
\]
Consequently $\log(D_s/du)/(fk)\ge15mL$. Put
$\Gamma=\exp(C_\Gamma m^5fL)$, the exact source \texttt{GamS}.
Positivity from (QW.14) justifies the real powers and logarithms, and
\[
 \begin{split}
 \log\!\left((D_s/du)^{1/(fk)}\Gamma^{-2/(fk)}\right)
 &=\frac{\log(D_s/du)}{fk}-\frac{2C_\Gamma m^4}{k}mL\\
 &\ge14mL\ge mL.
 \end{split}                                                    \tag{QW.22}
\]
This is \texttt{window\_\allowbreak lower}, with some slack preserved. The threshold
was chosen before $P,q_1,a,b,c,du,D_s$ and the word. This proves the
literal conditional theorem, not the existence of all its compatible
geometric inputs, and not the retained mass-ratio law itself.

\paragraph{QW10: remove the prime powers from the actual weighted sum.}
For the source family $\mathcal F=\texttt{LFam}$, retain its coefficient
laws $\lambda_g(n)\ge0$, $\sum_g\lambda_g(n)\le2\Lambda(n)$, and
the arithmetic support statement
\[
 p\text{ prime},\ \lambda_0(p)\ne0
                  \quad\Longrightarrow\quad\mathrm{NormOne}(d,M,p).
                                                                    \tag{QW.23}
\]
The latter is a condition on the same family, not merely an unrelated
prime represented by a norm. Split the sum over integers in $(X,Y]$
into primes and nonprimes. A prime outside the norm-prime set has
$\lambda_0(p)=0$ by (QW.23), whereas a prime inside it has
$\lambda_0(p)\le2\log p\le2\log Y$. At nonprimes the bound by
$2\Lambda(n)$ leaves only prime powers of exponent at least two.
For $Y\ge1$ this proves exactly
\[
 \sum_{\lfloor X\rfloor<n\le\lfloor Y\rfloor}\lambda_0(n)
 \le2\log Y\,N_{d,M}(X,Y)+2(\psi(Y)-\vartheta(Y)),               \tag{QW.24}
\]
where $N_{d,M}$ is the cardinality of the original norm-prime set.
The nonnegative remainder was enlarged to the full interval $(0,Y]$.
Elementary Chebyshev estimates give
$\psi(Y)-\vartheta(Y)\le C\sqrt Y$ for one absolute $C$.
For example the binomial coefficient estimate first gives
$\vartheta(t)\le C_1t$; then
$\sum_{j\ge2}\vartheta(Y^{1/j})\le C_1\sqrt Y+
C_1(\log_2Y)Y^{1/3}\le C_2\sqrt Y$ after enlarging the constant.
This uses no prime-number-theorem error term.

\paragraph{QW11: isolate exactly the analytic input that supplies ShortNorm.}
Retain the source arithmetic construction of an $\mathcal F$ satisfying
(QW.23) with conductor bound $\mathcal B=4dM^2$, for
$d$ prime, $d\equiv3\pmod4$, $d\ge100$, $M\ge3$.
Retain also its uniform weighted estimate: there are $A,c_0,e>0$ such
that, for every such family and $X\ge1$,
\[
 \log\mathcal B\le c_0\sqrt{\log X},\quad
 e^{-c_0\sqrt{\log X}}\le h\le1
 \quad\Longrightarrow\quad
 \sum_{X<n\le X(1+h)}\lambda_0(n)\ge\frac{e hX}{\mathcal B^A}.
                                                                    \tag{QW.25}
\]
Here and below integer endpoints mean the same floor convention as
(QW.24). Set $K=8(|C|+1)/e$ and choose
\[
 c=\min\left(c_0,\frac1{2(1+A)+2+K}\right)>0.
\]
For $s=\sqrt{\log X}$, $\mathcal B\ge4$ implies
$1\le\log\mathcal B\le cs$, so $s\ge2(1+A)+2+K$.
Consequently $c\le1$, $\mathcal B^A\le e^{As}$,
$h\ge e^{-s}$, and
\[
 K\le e^{s^2/2-(1+A)s},\qquad
                       K\sqrt X\le hX/\mathcal B^A.           \tag{QW.26}
\]
The first inequality follows from $e^t\ge1+t$ and the displayed
lower bound for $s$; multiplying the three exponential bounds proves
the second. In particular $\log X\ge1$.

For $Y=X(1+h)\le2X$, one has
$\log Y\le2\log X$ and $\sqrt Y\le2\sqrt X$.
Equations (QW.24)--(QW.26) show that the prime-power remainder consumes
at most half the right side of (QW.25). Dividing by $4\log X$ gives
\[
       N_{d,M}(X,X(1+h))\ge
          \frac{(e/8)hX}{\mathcal B^A\log X}.                  \tag{QW.27}
\]
This reconstructs the complete reduction to \texttt{ShortNorm}, with
$c'=e/8$. The arithmetic \texttt{ExistsLFamRayStatement} and analytic
\texttt{TrivialClassShortIntervalStatement} suppliers of (QW.23)--(QW.25)
remain; their names do not replace proofs of those assertions.

\paragraph{QW12: derive the source level-prime window uniformly in its endpoint.}
Assume \texttt{ShortNorm}, namely (QW.27) with some fixed $A,c,c'>0$
under its displayed hypotheses. The original \texttt{ModAt} says
\[
 d\text{ prime},\ d\equiv3\pmod4,\ d\ge100,\quad
 N>0,\ d\mid N,\quad \log N\le3x^{\epsilon/5}(\log x)^2.       \tag{QW.28}
\]
Fix $0<\epsilon<1/52$, $c_s,C_Y>0$. Suppose further
$g_i\ge1$, $\log g_i\le C_Yx^\epsilon$, $m\ge x^{1/3}$,
$1\le V\le x^{12}$, and $T_s>0$ with
$\log T_s\ge c_s m x^\epsilon$.
Set $M=N(g_0g_1)^2$ and $h=1/V$. Then $M\ge N\ge d\ge100$, and
\[
 \log(4dM^2)\le\log4+3\log N+8C_Yx^\epsilon,
 \qquad \sqrt{\log T_s}\ge\sqrt{c_s}x^{1/6+\epsilon/2}.
\]
Both $x^{\epsilon/5}(\log x)^2$ and $x^\epsilon$ are little-oh of
$x^{1/6+\epsilon/2}$; the second uses $\epsilon<1/3$.
Also $\log V\le12\log x$ is little-oh of that same power.
RF11 therefore supplies a single threshold, depending only on the fixed
constants, for which
$\log(4dM^2)\le c\sqrt{\log T_s}$ and
$h\ge e^{-c\sqrt{\log T_s}}$.
The positive right side of (QW.27) makes the finite norm-prime set
nonempty. Every chosen member satisfies exactly
\[
 T_s<T\le(1+1/V)T_s,\qquad
 T\text{ prime},\quad \mathrm{NormOne}(d,N(g_0g_1)^2,T).        \tag{QW.29}
\]
The floor endpoints give the strict lower inequality in (QW.29).
This proves \texttt{levelPrimeN} with the source quantifier order:
$T_s,V,m,N,d,g_i$ are chosen after the common threshold. Arbitrarily
large $T_s$ only improve the two tested inequalities. For the actual
host one may insert $C_Y=9$, $V=fk\le x^{12}$ and the original endpoint
lower bound to obtain the corresponding conditional \texttt{LevelPrime}.
The unconditional producer still needs QW11's retained suppliers.

\paragraph{QW13: the same input gives the literal nine-tenths dyadic count.}
Assume \texttt{ShortNorm} and \texttt{ModAt} again, now with any fixed
$\epsilon>0$. Put $t=x^{\epsilon/5}$, $\mathcal B=4dN^2$ and
$y=\log Y\ge8x^\epsilon=8t^5$. For sufficiently large $x$,
\[
 \log\mathcal B\le\log4+3\log N
 \le\log4+9t(\log x)^2\le10t^2.                               \tag{QW.30}
\]
Thus $\log\mathcal B/\sqrt y\le(10/\sqrt8)t^{-1/2}\le c$.
Apply (QW.27) with $h=1$ to obtain
$N_{d,N}(Y,2Y)\ge c'Y/(\mathcal B^A y)$.
The inequality needed for this to exceed $Y^{9/10}$ is
\[
                A\log\mathcal B+\log y-\log c'\le y/10.       \tag{QW.31}
\]
For $y\ge10$, $y/10-\log y$ is increasing. Hence it is enough to
check $y=8t^5$. At that value (QW.30) reduces (QW.31) to
$10At^2+\log(8t^5)-\log c'\le(4/5)t^5$, true eventually.
This proves, uniformly over every $Y$ with the original lower bound,
\[
                         N_{d,N}(Y,2Y)\ge Y^{9/10}.             \tag{QW.32}
\]
The exponent $9/10$ and the condition $\log Y\ge8x^\epsilon$ are
the literal \texttt{normPrimesN} output. They are conditional on the
quantitative analytic input, not consequences of IG's qualitative
unboundedness of ray primes.

\paragraph{QW14: a nonnegative character product proves nonvanishing at one.}
For completeness we supply the classical nonvanishing used by the
finite-modulus source lemma. Fix a modulus $N\ge1$. Use the standard
analytic continuation of Dirichlet $L$-functions: a nonprincipal
$L(s,\chi)$ is entire, and the principal factor is
$\zeta(s)\prod_{p\mid N}(1-p^{-s})$, with a single simple pole at one
and a zero at $s=-2$. These are the classical continuation and trivial
zero facts, independent of RH.\footnote{For the stated continuation,
principal pole and trivial-zero formulas, see NIST DLMF,
\url{https://dlmf.nist.gov/25.15}. The argument below supplies the
nonvanishing endpoint used by the pinned source's \texttt{noExc\_\allowbreak one}.}
Suppose some nonprincipal factor vanished at one. Then
\[
                  F(s)=\prod_{\chi\bmod N}L(s,\chi)
\]
would be entire, and $F(-2)=0$. In $\Re s>1$ its Dirichlet coefficients
are nonnegative: at a prime $p\nmid N$, if its class has order $d_p$
in $G=(\mathbb Z/N\mathbb Z)^\times$, character orthogonality gives
\[
 \prod_\chi(1-\chi(p)T)^{-1}=(1-T^{d_p})^{-|G|/d_p}.
\]
Every coefficient is nonnegative and the constant coefficient is one.
Write $F(s)=\sum_{n\ge1}a_n n^{-s}$, $a_n\ge0$, $a_1=1$.
For a fixed real $\sigma>1$, termwise differentiation is justified by
absolute convergence in a smaller right half-plane, and
\[
 (-1)^jF^{(j)}(\sigma)=\sum_{n\ge1}a_n n^{-\sigma}(\log n)^j\ge0.
\]
The Taylor series of an entire function has infinite radius. Evaluating
at $\sigma-t$ with $t=\sigma+2>0$ gives a sum of nonnegative terms,
whose constant term is $F(\sigma)\ge1$. Thus $F(-2)\ge1$, contradicting
$F(-2)=0$. Therefore every nonprincipal $L(1,\chi)$ is nonzero.
The argument establishes this classical endpoint without any assertion
about nontrivial zeros elsewhere in the critical strip.

\paragraph{QW15: finite-conductor real-zero margins preserve the final quantifier.}
For each fixed nonprincipal $\chi$, continuity and QW14 provide
$c_\chi>0$ such that $L(\beta,\chi)\ne0$ whenever
$1-c_\chi<\beta<1$ is real. Assign the harmless value $c_\chi=1$
to the principal character, for which no assertion is required.
There are finitely many characters modulo each positive $N$, and
finitely many $N\le N_0$. Their minimum is positive. It follows that
\[
 \forall N_0\ \exists c_0>0\ \forall\,1\le N\le N_0\
 \forall\chi\ne1\ \forall\beta\in(1-c_0,1):\quad
                              L(\beta,\chi)\ne0.               \tag{QW.33}
\]
For $N_0=0$ the assertion is vacuous and $c_0=1$ works. This is exactly
\texttt{QH.small\_\allowbreak zeros} and its literal real-variable \texttt{NoExc}.
The constant can depend on $N_0$; the proof provides no effective rate,
no fixed constant for all conductors, and no complex zero-free box.
In particular it cannot replace the uniform estimate (QW.25), the
independent simultaneous-prime selection, or the final RH/GRH argument.


\section{Normalized adic characters and their actual coefficient rings}
\label{sec:adic-realization}

This section closes the one-dimensional adic-character supplier left
explicit in AR. The input is a specified global reciprocity map
\[
 \operatorname{Art}:\mathbb A_E^\times\longrightarrow G_E^{\mathrm{ab}},
 \qquad
 \ker\operatorname{Art}
   =\overline{E^\times E_{\infty,+}^\times},                 \tag{GA.1}
\]
with the open quotient, inertia, decomposition and
\emph{inverse arithmetic Frobenius} clauses of RC3. Global class field
theory itself is retained. All conclusions below use this same map.
The algebraic closure of $\mathbb Q_\ell$, denoted
$\overline{\mathbb Q}_\ell$, carries its extended $\ell$-adic norm;
it is not silently replaced by its completion. The field isomorphism
$\iota:\overline{\mathbb Q}_\ell\simeq\mathbb C$ is algebraic and is
never assumed continuous.

Our coefficient datum consists of a DVR $O$, finite as a
$\mathbb Z_\ell$-module, its actual maximal-ideal-adic topology, and
an embedding
$\kappa:\operatorname{Frac}(O)\hookrightarrow\overline{\mathbb Q}_\ell$
that agrees with the standard map on $\mathbb Z_\ell$.
Write $g=\kappa|_O$. A Hecke character is a continuous
$\chi:\mathbb A_E^\times\to\mathbb C^\times$ trivial on $E^\times$.
For $n:\operatorname{Hom}(E,\mathbb C)\to\mathbb Z$,
\texttt{IsAlgPos} means
$\chi_\infty(a)=\prod_\tau\tau(a)^{n_\tau}$ for totally positive
$a\in E^\times$. A normalized adic realization is a continuous
$r:G_E\to O^\times$ satisfying
\[
 \kappa(r(\sigma))
 =\iota^{-1}\!\left(\chi(x)F_{\infty,n}(x)^{-1}\right)
       F_{\ell,n}(x)
 \quad\text{if }\operatorname{Art}(x)=\bar\sigma.            \tag{GA.2}
\]
Here the two finite products $F_{\infty,n}$ and $F_{\ell,n}$ are
defined in GA4. This is the literal \texttt{IsGal} identity,
including its coefficient embeddings.

\paragraph{GA1: the norm detects the actual DVR and its units.}
For the coefficient datum above, every $a\in O$ satisfies
$|g(a)|_\ell\leq1$, every unit has norm $1$, and
\[
 x\in O\ \Longleftrightarrow\ |\kappa(x)|_\ell\leq1,
 \qquad
 x\in O^\times\ \Longleftrightarrow\ |\kappa(x)|_\ell=1.      \tag{GA.3}
\]
We identify $O$ with its image in its fraction field in (GA.3).
Moreover $\ell\in\mathfrak m_O$.

Finiteness over $\mathbb Z_\ell$ makes $a$ integral. In a monic
equation for $a$, if $|g(a)|_\ell>1$, the leading term has larger
norm than every lower term, since their coefficients have norm at
most $1$. The ultrametric inequality rules out their sum being zero.
Applying this to a unit and its inverse gives norm $1$.
The element $\ell$ has norm $\ell^{-1}<1$ and is nonzero, so it is
a nonunit. Choose a uniformizer $\pi$ and write
$\ell=u\pi^e$, where $u\in O^\times$ and $e\geq1$.
Then
$0<|g(\pi)|_\ell=\ell^{-1/e}<1$.
Every nonzero fraction is uniquely a unit times an integral power of
$\pi$. Its norm is therefore $|g(\pi)|_\ell^a$, and it belongs to
$O$, respectively $O^\times$, exactly when $a\geq0$,
respectively $a=0$. This proves both directions, rather than just
boundedness of integral elements.

\paragraph{GA2: the coefficient fraction field has closed image.}
The image $L=\kappa(\operatorname{Frac}(O))$ is a finite-dimensional
$\mathbb Q_\ell$-subspace of $\overline{\mathbb Q}_\ell$, hence is
closed in that normed field.
Choose $\mathbb Z_\ell$-module generators $s_1,\ldots,s_d$ of $O$.
Every nonzero denominator $b=u_b\pi^a$ divides a power of $\ell$:
with $\ell=u\pi^e$, $a\geq0$,
\[
 b\bigl(u_b^{-1}u^a\pi^{(e-1)a}\bigr)=\ell^a.
\]
Consequently $\kappa(a_0/b)$ is a $\mathbb Q_\ell$-linear combination
of the $g(s_j)$. Conversely every such combination belongs to $L$:
a $\mathbb Q_\ell$ scalar is a fraction of elements of
$\mathbb Z_\ell$, and $\kappa$ respects their maps. Thus
$L=\sum_j\mathbb Q_\ell g(s_j)$.

For completeness, finite-dimensional subspaces over
$\mathbb Q_\ell$ are closed even when the ambient normed space is
incomplete. On a basis of such a subspace the coordinate sup norm
and the induced norm are equivalent. The lower estimate follows
by taking the minimum of the induced norm on the compact coordinate
unit sphere, where linear independence prevents zero; the upper
estimate follows from the triangle inequality. Cauchy sequences
therefore have convergent coordinates in complete $\mathbb Q_\ell$.
Their limits stay in the subspace, which proves closedness.
This is the closed-image fact needed for the quotient argument in
GA8; a merely algebraic subfield of
$\overline{\mathbb Q}_\ell$ would not suffice.

\paragraph{GA3: the embedding reflects topology and coefficient maps are continuous.}
The injection $g:O\to\overline{\mathbb Q}_\ell$ is a topological
embedding for the specified adic topology. Put
$c=|g(\pi)|_\ell\in(0,1)$. The inclusions
\[
 g(\pi^aO)\subseteq\{z:|z|_\ell\leq c^a\},
 \qquad
 |g(x)|_\ell<c^a\Longrightarrow x\in\pi^aO                \tag{GA.4}
\]
compare neighborhood bases. For the second, apply GA1 to
$x/\pi^a$. Since $c^a\to0$, they prove the assertion.
Continuity of an $O$-valued map can thus be tested after $g$.
For a homomorphism from a topological group into $O^\times$,
continuity of its underlying $O$-valued map also gives continuity
in the unit topology: the inverse coordinate is its value at
the inverse group element.

More generally, let $O'$ be an adic topological local ring and
$t:O\to O'$ a unital ring homomorphism, with
$t(\ell)\in\mathfrak m_{O'}$. This map is continuous.
Indeed
$t(\pi)^{ea}=t(u)^{-a}t(\ell)^a\in\mathfrak m_{O'}^a$;
the image of $\pi^{ea}O$ lies in $\mathfrak m_{O'}^a$.
This uses the actual adic topologies, and does not assume $t$
is a local map beforehand. In particular it applies to the
coefficient extensions used in GA19 and GA21.

\paragraph{GA4: the archimedean and adic algebraic products.}
For an idele $x$ define
\[
 F_{\infty,n}(x)=
 \prod_{v\mid\infty}\ \prod_{\substack{j:E_v\to\mathbb C\\j\ {\rm continuous}}}
       j(x_v)^{\,n_{j|_E}},
 \qquad
 F_{\ell,n}(x)=
 \prod_{v\mid\ell}\ \prod_{\substack{j:E_v\to\overline{\mathbb Q}_\ell\\
                                  j\ {\rm continuous}}}
       j(x_v)^{\,n_{\iota j|_E}}.                         \tag{GA.5}
\]
They are continuous multiplicative maps into the corresponding
unit groups. Each indexing set is finite: restriction to the dense
subfield $E$ is injective, and a number field has finitely many
embeddings. There are no continuous embeddings from $E_v$ into
$\overline{\mathbb Q}_\ell$ when $v\nmid\ell$.
For otherwise choose an integer divisible by the residue
characteristic of $v$ and prime to $\ell$; its powers tend to zero
in $E_v$ while their images have $\ell$-adic norm $1$.

Restriction identifies the union of the archimedean embedding sets
with all embeddings $E\to\mathbb C$. Likewise the union of the
$\ell$-adic embedding sets identifies, after $\iota$, with those
same embeddings. To see the latter without assuming continuity of
$\iota$, begin with $\tau:E\to\overline{\mathbb Q}_\ell$.
The valuation of $\tau$ gives a place $v\mid\ell$; equivalence of
the discrete valuations shows continuity on $E$ for $v$.
The image lies in the finite extension
$\mathbb Q_\ell(\tau(E))$, which is complete by the argument of
GA2, so $\tau$ extends uniquely to $E_v$. Conversely a continuous
extension determines its place from the induced valuation, and
density gives uniqueness. At infinity the same assertion is
the extension to $\mathbb R$ or $\mathbb C$, with complex
conjugation giving the second embedding of a complex completion.
Consequently, for $a\in E^\times$,
\[
 F_{\infty,n}(a)=\prod_\tau\tau(a)^{n_\tau},\qquad
 F_{\ell,n}(a)=\iota^{-1}\!\left(\prod_\tau\tau(a)^{n_\tau}\right).
                                                                  \tag{GA.6}
\]
The adic product is $1$ on purely infinite ideles; the infinite
product is $1$ on purely finite local ideles.

\paragraph{GA5: removing the infinite type gives an open kernel.}
Suppose $\chi$ satisfies \texttt{IsAlgPos} for $n$ and set
$\theta=\chi F_{\infty,n}^{-1}$. Then $\theta$ is continuous and
trivial on $E_{\infty,+}^\times$, and its kernel in the idele group
is open.

First, the connected component of $1$ in
$E_\infty^\times$ is the product of positive real factors and
complex multiplicative factors; it is open. Weak approximation
makes $E^\times$ dense in $E_\infty^\times$. Its intersection with
that component consists of totally positive elements and is dense
there. The defining \texttt{IsAlgPos} equality and continuity
therefore give $\theta=1$ on the whole component.

Every continuous homomorphism from the finite idele group to
$\mathbb C^\times$ has open kernel. Here is the needed elementary
argument. There is a neighborhood of $1$ on which all image values
lie in $|z-1|<1/2$. Inside it choose an open subgroup of the form
\[
 \{u:u-1\in N\widehat{\mathcal O}_E,\quad
           u^{-1}-1\in N\widehat{\mathcal O}_E\},\quad N\geq1.
                                                                  \tag{GA.7}
\]
Such subgroups form a neighborhood basis by the restricted-product
topology. If $z$ and all its powers satisfy $|z^a-1|<1/2$,
then $z=1$: the geometric sum $S_a=\sum_{h<a}z^h$ satisfies
$|S_a-a|\leq a/2$, so $|S_a|\geq a/2$, while
$|(z-1)S_a|=|z^a-1|<1/2$, a contradiction for large $a$ if
$z\ne1$. All powers of an element of (GA.7) stay in that
subgroup, proving it is killed. Combining this finite open
subgroup with $E_{\infty,+}^\times$ proves the claim.
It follows that $\iota^{-1}\theta$ is continuous: an arbitrary
homomorphism with open kernel is locally constant, irrespective
of the topology of $\iota$.

\paragraph{GA6: construction and uniqueness of the normalized adic character.}
Let $\chi,n$ satisfy \texttt{IsAlgPos}. Define a continuous
idele character
\[
 \Psi(x)=\iota^{-1}\!\left(\chi(x)F_{\infty,n}(x)^{-1}\right)
                          F_{\ell,n}(x).                 \tag{GA.8}
\]
Continuity follows from GA4 and GA5. On a principal idele,
$\chi(a)=1$ and (GA.6) cancels the other two factors.
On $E_{\infty,+}^\times$ the first factor is $1$ by GA5 and
the second is $1$ by GA4. The closed kernel of $\Psi$
therefore contains the kernel (GA.1). It descends uniquely to
a continuous character of $G_E^{\mathrm{ab}}$, since
$\operatorname{Art}$ is an open quotient map. Composing with
$G_E\to G_E^{\mathrm{ab}}$ gives a continuous
$r_\chi:G_E\to\overline{\mathbb Q}_\ell^\times$ satisfying
(GA.2).

If another character satisfies (GA.2), choose an idele mapping
to each $\bar\sigma$ by surjectivity of $\operatorname{Art}$.
The displayed equality gives the same value on $\sigma$.
This proves uniqueness. Thus the source's published
\texttt{HeckeLAdicCharacter} statement follows from the retained
\texttt{GlobalArtinReciprocity}; no second character-existence
hypothesis is required.

\paragraph{GA7: local reciprocity, inertia, and the Frobenius inverse.}
If $\operatorname{Art}(x_v)=\bar\sigma$, then
\[
 r_\chi(\sigma)=\iota^{-1}(\chi_v(x_v))
       \prod_jj(x_v)^{n_{\iota j|_E}}.                    \tag{GA.9}
\]
This is (GA.2) for a local finite idele. The product is $1$
when $v\nmid\ell$, or when $n=0$. Under either condition,
the inertia-image clause in (GA.1) gives exactly
\[
 r_\chi(I_v)=\iota^{-1}(\chi_v(\mathcal O_v^\times)).
                                                                  \tag{GA.10}
\]
In particular $r_\chi$ is unramified at $v$ if and only if
$\chi$ is, and the two inertia images have the same exponent
and, when finite, the same order. The decomposition clause,
continuity and closure similarly show that $r_\chi$ is
trivial on the whole decomposition group if and only if
$\chi_v$ is the trivial local character.

For arithmetic Frobenius $\sigma$ at any valuation above an
unramified $v$, the specified Artin convention reads a
uniformizer as $\bar\sigma^{-1}$. Consequently
\[
 \iota(r_\chi(\sigma))=\chi(v)^{-1}
       \quad(v\nmid\ell\text{ or }n=0).                 \tag{GA.11}
\]
These assertions also hold for $O$-valued realizations after
applying $g$. Multiplication and inversion of normalized
characters add and negate $n$, because the two products in
(GA.5) have exactly those laws. No inverse is added to the
local unit formula (GA.9).

\paragraph{GA8: descent into the specified coefficient DVR.}
Suppose a field $K$ has embeddings
$\epsilon:K\hookrightarrow\mathbb C$ and
$f:K\hookrightarrow\operatorname{Frac}(O)$ satisfying
$\iota\kappa f=\epsilon$, and that $\chi(v)\in\epsilon(K)$
for all but finitely many finite places. Then the character
of GA6 takes its values in $g(O^\times)$ and gives the exact
\texttt{exists\_\allowbreak isGal} conclusion.

Let $L$ be the closed field image from GA2. The subgroup
$L^\times\subseteq\overline{\mathbb Q}_\ell^\times$ is
closed, since it is the inverse image of $L$ under the
continuous underlying-value map. The quotient by it is a
Hausdorff commutative topological group. Compose $r_\chi$
with this quotient. Outside the finitely many excluded,
ramified and $\ell$-adic places, (GA.11) belongs to
$\kappa f(K^\times)$, so the quotient character kills every
Frobenius. RC4 proves that the quotient character is
identically $1$. Thus every value of $r_\chi$ belongs to $L$.

The absolute Galois group is compact. A continuous character
from a compact group into a normed field has values of norm
$1$: its image is bounded, and a value with norm greater
than $1$ would have unbounded positive powers; apply the
same argument to its inverse. By GA1 the unique preimage
of each value lies in $O^\times$. Uniqueness makes these
preimages multiplicative. GA3 makes the resulting $O$-valued
map continuous, and the unit inverse coordinate is continuous
as well. Its embedded character is precisely $r_\chi$, so it
satisfies (GA.2) for the original $\chi,n,\kappa,\iota$.
This constructs the normalized character itself; it does not
replace it by a finite-order character.

\paragraph{GA9: finite-order characters and exact recovery.}
If $\chi^m=1$ with $m>0$, then it satisfies
\texttt{IsAlgPos} with $n=0$. Indeed every totally positive
infinite idele has an $m$th root in the product of positive
real and nonzero complex factors; evaluating $\chi$ on
that power gives $1$. Formula (GA.2) then gives
$r_\chi^m=1$. Its image is a finite set of roots of unity,
so its kernel is open: the complement of $1$ in that
finite set is closed, and continuity pulls its complement
back to the kernel.

Conversely, let $r:G_E\to\overline{\mathbb Q}_\ell^\times$
be continuous with $r^m=1$. It descends through topological
abelianization, since its target is commutative and its
kernel is closed. Set
$\chi(x)=\iota(\bar r(\operatorname{Art}(x)))$.
This is continuous even though $\iota$ need not be:
$r$ has open kernel and the resulting character is locally
constant. It kills principal ideles, has order dividing
$m$, and satisfies (GA.2) with $n=0$.
Surjectivity onto the abelianization proves uniqueness.
The local identities of GA7 hold at \emph{all} finite places
in this case, including those above $\ell$.

\paragraph{GA10: anticyclotomic symmetry follows from the actual character.}
Suppose $E/F$ is the host quadratic CM extension with conjugation
$c$, and $\widetilde c$ is the supplied automorphism of the
algebraic closure extending it. Retain the host's finite-place
unramified equality
$v_P(a)=v_{P\cap F}(a)$ for $a\in F^\times$, and let
$\chi$ be trivial on the ideles of $F$. Every normalized
realization satisfies
\[
 r_\chi(\widetilde c\sigma\widetilde c^{-1})
      =r_\chi(\sigma)^{-1}.                             \tag{GA.12}
\]
The map $\sigma\mapsto\widetilde c\sigma\widetilde c^{-1}$
is continuous for the Krull topology: fixing the conjugates
of a finite field basis forces the conjugate automorphism
to fix the original finite field.

At an unramified pair $P,\bar P$ outside the $\ell$-adic
places, restriction of $\chi$ to a uniformizer of the
underlying place of $F$ gives
$\chi(P)\chi(\bar P)=1$. If $P=\bar P$, the same restriction
and the retained valuation equality give $\chi(P)=1$,
which is enough for this product identity. If they are
different, the two primes are all primes over that place,
each with exponent one. Conjugating a valuation and its
arithmetic Frobenius preserves the residue-field cardinality
and gives Frobenius at $\bar P$. Equation (GA.11) therefore
makes the continuous character
$\sigma\mapsto r_\chi(\widetilde c\sigma\widetilde c^{-1})
r_\chi(\sigma)$ trivial on every Frobenius outside a finite
conjugation-stable set. RC4 proves (GA.12). This proves the
source's \texttt{AntiHom} property with its supplied
conjugation, rather than choosing a new involution.

\paragraph{GA11: the CM parity calculation gives integral infinite types.}
Use the host CM type $\Phi=\{\phi_\tau\}$ of AH and HF.
Every complex embedding of $E$ occurs exactly once among
$\phi_\tau,\bar\phi_\tau$: equality of their infinite places
first gives equality of $\tau$, and total complexity excludes
an embedding fixed by conjugation. For $z\ne0$ and $a\in\mathbb Z$,
\[
 (z/|z|)^{2a}=z^a\bar z^{-a}.                            \tag{GA.13}
\]
Thus a character with argument exponent $2a_\tau$ has
algebraic exponents $(a_\tau,-a_\tau)$ on this pair.
In particular the HF character $\mu_{0,j}$, whose argument
exponent is $-2e_{j,\tau}$, is algebraic when the prescribed
$e_{j,\tau}$ are integers.

For an odd integer $r=2h+1$ and the HF character $\eta$ with
argument exponent $1$, the infinite contribution of
$\eta^{-1}|\cdot|^{-r/2}$ at a complex place is
\[
 (z/|z|)^{-1}|z|^{-r}
       =z^{-h-1}\bar z^{-h}.                            \tag{GA.14}
\]
Here the local absolute-value factor at a complex place
counts both embeddings. Hence this character is algebraic
with integral exponents $(-h-1,-h)$.
The source has $r=k+m-1$, with $k,m$ even, so $r$ is odd;
$t$ is separately odd. The same argument applies to $t$.
It does not identify the global chain length $m$ with any
local polynomial degree.

\paragraph{GA12: odd normalization preserves the specified field of values.}
Write
\[
 B_r(\mu,\eta)=\mu\,\eta^{-1}|\cdot|^{-r/2}.
 \qquad
 B_r(\mu,\eta)(P)=\mu(P)\eta(P)^{-1}N(P)^{r/2}.           \tag{GA.15}
\]
The second equality uses $|\varpi_P|=N(P)^{-1}$.
Suppose $\mu(P)\in\epsilon(K)$ cofinitely often, and that
$\eta(P)N(P)^{-1/2}\in\epsilon(K)$ at every unramified
place, the literal \texttt{ValuesIn} normalization for
$\eta$. If $r=2h+1$, choose $y,z\in K^\times$ with
$\epsilon(y)=\mu(P)$ and
$\epsilon(z)=\eta(P)N(P)^{-1/2}$. Then
\[
 B_r(\mu,\eta)(P)=\epsilon(yz^{-1}N(P)^h).                \tag{GA.16}
\]
The absolute norm is a positive integer, so its integral
power belongs to the prime field. This gives
\texttt{AlmostIn} for the \emph{same} field $K$.
Products and inverses preserve this cofinite condition
by intersecting the finite exceptional sets.
On a local unit, if $\eta$ is unramified,
\[
 B_r(\mu,\eta)_P(u)=\mu_P(u).                            \tag{GA.17}
\]
The norm factor and $\eta_P(u)$ are both $1$.

\paragraph{GA13: construction of the exact shift and block characters.}
Fix the actual $D$, its host character frame $C$, an index
$i$ of the supplied $\mathrm{NTw}$, a side $b$, and a
\texttt{Reading} $(\kappa,\iota)$ for its ring $O_0(i)$.
Thus
\[
 \iota\kappa\,\operatorname{embSide}_{i,b}(x)
       =\operatorname{emb}_{LL_0}(x)
 \quad(x\in LL_0),                                     \tag{GA.18}
\]
and $\kappa$ has the specified $\mathbb Z_\ell$ compatibility.
Retain \texttt{ValLL0} for this actual frame; HF16 supplies
it when its value field is contained in this actual $LL_0$.
There exist normalized $O_0(i)^\times$-valued characters
for
\[
 \operatorname{shiftChar}(C)=B_t(1,\eta),\qquad
 \operatorname{blockChar}(C,\ell,j)
        =B_r(\mu_{0,j}\delta_\ell,\eta),                 \tag{GA.19}
\]
where $\delta_\ell=\delta'$ if $\ell=p_\delta$, and
$\delta_\ell=\delta$ otherwise.

GA11 gives the integral infinite type of the shift.
The auxiliary character is finite order, hence algebraic
of type zero by GA9; GA11 then gives the block's type.
The odd $t$ and odd $r$ calculations of GA12 put their
values cofinitely in the actual $LL_0$.
Apply GA8 with $K=LL_0$, $f=\operatorname{embSide}_{i,b}$
and (GA.18). This proves the literal
\texttt{Triple.exists\_\allowbreak shift} and
\texttt{Triple.exists\_\allowbreak blockChar} endpoints. The character
and its integral type are chosen after the actual reading;
the construction does not select a coefficient place or
produce a twist index.

\paragraph{GA14: changing a reading does not change the integral character.}
Let $\rho,\rho'$ be two readings of the same
$(i,b,O_0(i),\operatorname{embSide}_{i,b})$. If
$y\in O_0(i)$ and
$\rho_{\mathbb C}(y)=\operatorname{emb}_{LL_0}(x)$, then
injectivity of $\rho.\iota$ and $\rho.\kappa$ gives the
fraction-field equality
$y=\operatorname{embSide}_{i,b}(x)$.
Applying $\rho'$ gives the same complex value.

Suppose an integral character $r$ realizes the block for
$\rho$. Construct $r'$ for $\rho'$ by GA13. At cofinally
many unramified places the inverse block value lies in
the common $LL_0$ by GA12. Equation (GA.11), the preceding
fraction-field equality and injectivity show
$r(\mathrm{Frob})=r'(\mathrm{Frob})$.
Both are continuous in the same adic ring; RC4 gives
$r=r'$. It follows that $r$ also realizes the block for
$\rho'$, with its appropriate infinite type.
The tame-generator equality in \texttt{IsChars} is
independent of the reading too: lift the specified
inertia generator to a local unit using the same Artin
map and apply (GA.9) to both readings. This proves
\texttt{isGal\_\allowbreak of\_\allowbreak reading} and \texttt{isChars\_\allowbreak of\_\allowbreak reading}
without treating arbitrary complex embeddings as continuous.

\paragraph{GA15: identification with the already installed first character.}
Retain the literal \texttt{NormalizedCharacterCompatibility}
record for $D$ and its $NChiData$, and the three identities
identifying its $\eta$, $\mu_0(j+1)$ and $\delta_i$ with
the same frame $C$. For every existing $i,b,j$ and every
reading $\rho$, the \emph{already installed}
$D.\mathrm{ntw}.\chi_1(i,j+1,b)$ realizes the exact block
in (GA.19).

Indeed GA13 constructs a comparison character $r$.
The retained \texttt{chi1} reading says, away from a finite
set, that for some $x_P\in LL_0$ and unramified Hecke value
$\alpha_P$,
\[
 \operatorname{emb}_{LL_0}(x_P)
       =\alpha_P^{-1}\sqrt{N(P)}^{-r},\qquad
 \chi_1(\sigma)=\operatorname{embSide}_{i,b}(x_P)           \tag{GA.20}
\]
in the fraction field at every Frobenius above the specified
valuation. The three character identities make
$\alpha_P=(\mu_{0,j}\delta_\ell\eta^{-1})(P)$.
Since $N(P)>0$, (GA.15) shows that the first expression
in (GA.20) is exactly the inverse block value.
Equations (GA.11) and (GA.18) therefore identify the two
embedded Frobenius values.

The source initially specifies one valuation above $P$.
Any other is its Galois conjugate; conjugating the
Frobenius equation gives Frobenius at the specified one,
and homomorphisms into commutative groups are unchanged
by this conjugation. Thus equality holds at all valuations
required by RC4. The installed character is continuous
by its $\mathrm{NTw}$ field, so RC4 and injectivity give
$r=\chi_1$. This proves
\texttt{firstTwistCharacter\_\allowbreak eq\_\allowbreak ladic\_\allowbreak heckeCharacter}.
The proof consumes the normalized reading record; existence
of a newly constructed $r$ alone would not identify an
arbitrarily preinstalled $\chi_1$.

\paragraph{GA16: exact auxiliary inertia and both side families.}
Let the actual HF frame and the integral $e_{j,\tau}$ be
fixed. Assume the two type primes split completely as in
the source, and retain both coefficient embeddings
$f,f':K\to\operatorname{Frac}(O)$ with
\[
 \iota\kappa f=\epsilon,\qquad
 \iota\kappa f'=\overline{\epsilon}.
\]
Put $\iota_+=\iota$, $\iota_-=\overline{\iota}$.
Then $\iota_b\kappa f_b=\epsilon$ for $f_+=f,f_-=f'$.
Choose the auxiliary prime
$p=p_\delta'$ when $\ell=p_\delta$, and $p=p_\delta$
otherwise. It is prime, distinct from $\ell$ and both
type primes, and $p\equiv1\pmod{15}$.
GA8, GA9 and GA13 construct $\Delta$ for $\delta_\ell$
and $\chi_{j,b}$ for each block with reading $\iota_b$.

The kernel of $\Delta$ is open, and GA10 makes it
anticyclotomic. GA7 and the HF12 exact local exponent
give
\[
 \Delta\text{ ramifies exactly above }p,\qquad
 \#\Delta(I_P)=15\quad(P\mid p).                        \tag{GA.21}
\]
For the order assertion the inertia image is a subgroup
of the $15$th roots in a field, hence finite cyclic.
Its exponent is $15$ because the local exponents divide
$a$ exactly when $15\mid a$; a finite cyclic group's
exponent equals its order. Also $N(P)=p^{f_P}$, so
$N(P)\equiv1\pmod{15}$.

At $P\mid p$, (GA.17) kills the $\mu_0$ factor and
leaves $\delta_\ell$. Since finite-order complex values
have conjugate equal to inverse, (GA.9) gives
\[
 \chi_{j,+}|_{I_P}=\Delta|_{I_P},\qquad
 \chi_{j,-}|_{I_P}=\Delta^{-1}|_{I_P}.                   \tag{GA.22}
\]
Away from $\ell$, the two type primes and $p$, these
characters are unramified. At a type prime $g_k\ne\ell$,
their inertia values have order dividing $g_k-1$, by
tameness and complete splitting, as detailed in GA17.
To define the literal family for every natural $j$,
use the genuine finite family at
$\min(j-1,m-1)$; $m\geq2$ makes this index legitimate,
and it recovers each prescribed $j+1$ for $j\in\operatorname{Fin}(m)$.
This proves every clause of
\texttt{HeckeTwistCharactersDeltaChi}.

\paragraph{GA17: tame and wild local unit orders.}
At a finite place with residue cardinality $q$, reduction
gives $u^{q-1}\equiv1$ for every local unit $u$.
A tame Hecke character therefore satisfies
$\chi_P(u)^{q-1}=1$.
Conversely, suppose a continuous local character kills
all $(q-1)$st roots of unity in the local field.
Its restriction to local units has open kernel: choose
a sufficiently small congruence subgroup mapped into
$|z-1|<1/2$ and apply the power argument of GA5.
Hensel's lemma applied to $X^{q-1}-1$, whose derivative
is a unit, gives a root $\zeta$ with the same residue as
$u$. Then $v=u/\zeta$ belongs to $1+\mathfrak m_P$.
Writing $q$ in the maximal ideal and using
\[
 v^q-1=(v-1)\sum_{a<q}v^a,\qquad
 \sum_{a<q}v^a=q+\sum_{a<q}(v^a-1)
 \in\mathfrak m_P,
\]
shows that each $q$th power raises the additive depth
by at least one. Hence $v^{q^a}$ eventually lies in
the open kernel of the character, and
$\chi_P(u)^{q^a}=1$ for some $a$.
This proves the two exact \texttt{orders\_\allowbreak on\_\allowbreak units}
statements, including dyadic places.

We will also need the Galois version. Let a continuous
commutative character have inertia image of exponent
$m$ prime to the residue characteristic, and let $\gamma$
be a specified tame generator. Its image kills wild
inertia, by the finite ramification filtration argument
of AR3. Continuity and the finiteness of the $m$th roots
allow this argument in a finite Galois quotient: choose
an open normal subgroup whose image avoids every
nonidentity $m$th root. By the definition of a tame
generator, inertia is the closure of the group generated
by $\gamma$ and wild inertia. The preimage of the finite
cyclic group generated by the character value at $\gamma$
is closed and contains those generators. Thus every
inertial value is a power of that value. This proves
\texttt{mem\_\allowbreak zpowers\_\allowbreak of\_\allowbreak tameGen} for the continuous
characters used below.

\paragraph{GA18: a uniform calibration of the specified tame generator.}
Fix a type prime $p=g_k$, a place $P$ in its chosen half,
a valuation above it and its specified tame generator
$\gamma_P$. Complete splitting gives residue field
$\mathbb F_p$. Lift $\gamma_P$ by the Artin inertia
clause to a local unit $u$, and denote its residue by
$a\in\mathbb F_p^\times$. The HF local law, first on
rational residue lifts and then on all units by tameness,
reads
\[
 \mu_{0,j,P}(u)=\omega_k(a)^{-2(j+1)}.                  \tag{GA.23}
\]
We claim that $a^2$ generates the subgroup of squares.
Use the row $j=0$ and an adic realization of $\mu_{0,0}$
at the distinct auxiliary prime $p_\delta$; GA6 supplies
it, without any coefficient-field restriction. Its
inertia image at $P$ has exponent dividing $p-1$.
For every $b\in\mathbb F_p^\times$, lift $b$ to a rational
local unit and then to inertia. GA17 makes its character
value a power of the value at $\gamma_P$. Equation
(GA.23) and injectivity of $\omega_k$ imply
$b^2\in\langle a^2\rangle$.

In a cyclic group of order $p-1$, this implies there is a
generator $a_0$ with $a_0^2=a^2$. To check the possible
factor of two, write the group as $\langle z\rangle$,
$a=z^d$. The condition is
$\gcd(2d,p-1)=\gcd(2,p-1)$.
If $p-1$ is odd, $a$ already generates.
If it is even and $\gcd(d,p-1)=1$, the same is true.
Otherwise $\gcd(d,p-1)=2$ and $(p-1)/2$ is odd;
$a z^{(p-1)/2}$ generates and has the same square.

The two injective characters $\omega_k$ and the specified
Teichmuller map of $\mathbb F_p^\times$ have the same
image, namely all $(p-1)$st roots of unity. Transport
$a_0$ through these injections, obtaining a generator
$a_P$ read by the specified Teichmuller map. For
\emph{every} $\ell\ne p$, every allowed coefficient ring
and reading, every $j$, and every $\xi$ unramified at $P$,
any realization of $B_r(\mu_{0,j}\xi,\eta)$ satisfies
\[
 \iota\kappa(\chi(\gamma_P))
   =\operatorname{Teich}\bigl(a_P^{-2(j+1)}\bigr).        \tag{GA.24}
\]
This follows directly from (GA.9), (GA.17), (GA.23)
and $a_0^2=a^2$. Crucially $a_P$ is chosen before all
these coefficient data and $j$. This proves
\texttt{TameInertiaGeneratorCharacterSpec}; setting
$c(k,P,j)=a_P^{2(j+1)}$ gives the actual
\texttt{exists\_\allowbreak tameTie}, for a new such array.
It does not identify that array with an arbitrary
preinstalled $D.\mathrm{nd}.\chi_T$.

\paragraph{GA19: the toric quotient character and its integral Frobenius values.}
Let $\vartheta$ be a Hecke character trivial on $F$,
with even integral argument exponents. Assume the
specified $K$ contains the values of $\mu_{0,j}$,
$\delta_\ell$ and $\vartheta$.
GA8 and GA11 construct the normalized realization
$\lambda_0$ of
$\mu_{0,j}\delta_\ell\vartheta^{-1}$, using the selected
side reading. For any coefficient map $O\to O'$ satisfying
GA3, let $\lambda$ be its image in $O'^\times$.
It is continuous and anticyclotomic by GA10.
If $S_0$ contains the $\ell$-adic, type and auxiliary
places and all ramification of $\vartheta$, then
$\lambda$ is unramified outside $S_0$ by GA7.

Suppose at the auxiliary places $\vartheta_P$ kills
the $(N(P)-1)$st roots of unity. On local units the
quotient of $\lambda$ by the mapped
$\Delta^{\epsilon_b}$ is read as $\vartheta_P^{-1}$,
where $\epsilon_+=1,\epsilon_-=-1$.
GA17 therefore proves, for every $\sigma\in I_P$,
\[
 \bigl(\lambda(\sigma)\,
          \Delta'(\sigma)^{-\epsilon_b}\bigr)^{N(P)^a}
        =1\quad\text{for some }a\geq0.                 \tag{GA.25}
\]
This is the exact power-of-residue-cardinality conclusion,
not a claim that the quotient is trivial.
Define $X(P)$ as the image of $\lambda_0$ at the source's
chosen Frobenius at $P$. If $P\notin S_0$, equation
(GA.11) shows the value is independent of every choice
of valuation and Frobenius. Thus $X(P)$ lies in $O'$,
with the literal equality to $\lambda(\sigma)$ for
all such choices. These constructions prove
\texttt{ToricHeckeQuotientCharacter}.

\paragraph{GA20: localization characters with nontrivial inertia on their support.}
Fix an actual Hecke character $\psi$ satisfying
$\texttt{IsPsi}(\ell,Q)$: it has $\ell$-power order,
is trivial on $F$, is tame, ramifies exactly on $Q$,
and its entire local character is trivial above $\ell$.
Retain the prescribed separation of $Q$ from the type
and selected auxiliary primes and the common field of
values for $\mu_{0,j},\eta,\delta_\ell,\psi$.
GA8 constructs $\psi_\ell$ and a character $\chi_{\rm base}$ realizing
$B_r(\mu_{0,j}\delta_\ell\psi,\eta)$. With
$\varepsilon_\ell:G_E\to\mathbb Z_\ell^\times$
the actual cyclotomic character, define
\[
 \chi_c=\chi_{\rm base}\,\varepsilon_\ell^{-(c-1)}
       \quad(c\in\mathbb N),\qquad \chi_1=\chi_{\rm base}, \tag{GA.26}
\]
mapping $\varepsilon_\ell$ into $O^\times$. GA3 and
compatibility on $\mathbb Z_\ell$ prove continuity.
The integer exponent in (GA.26) also defines $c=0$;
later source uses may restrict to $1\leq c\leq t$.

GA7--GA10 show that $\psi_\ell$ has open kernel,
$\ell$-power order, anticyclotomic symmetry, exact
ramification set $Q$, and trivial action of every
decomposition group above $\ell$.
For each $P\in Q$, $P\nmid\ell$ follows already from
the local triviality of $\psi$ above $\ell$.
On its units the other factors of the block are trivial.
If $\chi_{\rm base}$ killed the whole inertia group, (GA.10) would
force $\psi_P$ to be trivial on units, contradicting
$P\in Q$. Hence $\chi_1$ is genuinely ramified at
every place in $Q$.

Away from $\ell$, the type primes, the selected auxiliary
prime and $Q$, put
$\vartheta=\mu_{0,j}\delta_\ell\psi\eta^{-1}$. The exact
Frobenius identity is
\[
 \iota\kappa(\chi_c(\sigma))
     =\vartheta(P)^{-1}N(P)^{-r/2-(c-1)}.               \tag{GA.27}
\]
Indeed (GA.11) and (GA.15) prove it for $c=1$.
Arithmetic Frobenius acts on every $\ell$-power root
of unity by its $N(P)$th power: reduction has that
action, and reduction is injective on these roots
when $P\nmid\ell$, since $X^{\ell^a}-1$ has unit
derivative at each root. This proves its cyclotomic
value is $N(P)$.
The source exponent at $c=1$ is
$((t-1)/2)+(1-n_V)/2=(t-n_V)/2=-r/2$,
since $n_V=r+t$. This proves every clause of
\texttt{LocalisationCharacterSpec}, conditional on
the actual $\psi,Q$; it does not construct
\texttt{IsPsi} for an arbitrary nonempty $Q$.

\paragraph{GA21: coefficient extension commutes with the normalized twist.}
Let $O_0,O$ have the same residue prime $\ell$, actual
readings $(\kappa_0,\iota_0)$ and $(\kappa,\iota)$,
and maps $f_0:K\to\operatorname{Frac}(O_0)$,
$f:K\to\operatorname{Frac}(O)$ reading the same
$\epsilon:K\to\mathbb C$. Suppose
$t:O_0\to O$ has the precise compatibility
\[
 y=f_0(x)\text{ in }\operatorname{Frac}(O_0)
       \quad\Longrightarrow\quad t(y)=f(x)
            \text{ in }\operatorname{Frac}(O)
 \quad(y\in O_0,\ x\in K).                             \tag{GA.28}
\]
Let $\chi_1,\chi,\psi_\ell$ be normalized realizations
of respectively $B_r(\mu_{0,j}\delta_\ell,\eta)$,
$B_r(\mu_{0,j}\delta_\ell\psi,\eta)$ and $\psi$.
Then
\[
                 \chi=t_*(\chi_1)\,\psi_\ell.           \tag{GA.29}
\]
Here $t$ is continuous by GA3 and both sides are continuous
characters into the same Hausdorff unit group $O^\times$.
For cofinally many places the inverse first-block value
is $\epsilon(x)$ for some $x\in K$, by GA12.
Equation (GA.11) identifies $\chi_1(\sigma)$ with
$f_0(x)$ in the fraction field, so (GA.28) identifies
its image under $t$ with $f(x)$.
Multiplying the two Frobenius values now gives the value
of $\chi$, since the two Hecke blocks differ by $\psi$.
RC4 proves (GA.29). This proves
\texttt{LocalisationTwistCharacterCompat} even when
the two algebraic complex readings differ away from $K$.
It does not assume that they commute on the entire
fraction field.

\paragraph{GA22: installed local characters, DeltaRead, and IsGalOf.}
The coefficient descent has a useful local form. Suppose
a supplied continuous $r_0:G_E\to O_0^\times$ has
cofinite normalized Frobenius readings in $K$ for a
specific algebraic Hecke character $\chi$. Let
$t:O_0\to O$ and the local field embedding
$f:K\to\operatorname{Frac}(O)$ satisfy the literal
integral-value compatibility (GA.28).
The target $\kappa,\iota$ must read $f$ as $\epsilon$.
Then $t_*r_0$ is the normalized realization of $\chi$.
To prove this, its cofinite readings put $\chi(P)$
cofinitely in $\epsilon(K)$; GA8 constructs a target
realization $r$. GA3 proves continuity of $t_*r_0$.
At the specified Frobenii its values and those of $r$
agree by the actual reading and coefficient equation;
conjugation transfers equality to every valuation above
the place. RC4 gives $r=t_*r_0$.

Apply this with the \emph{same} $D$, normalized
\texttt{SideB.chi.chi1} reading, local index, map
$T.\mathrm{twist}_2$, and $\operatorname{Sp.iq}$ of
the source \texttt{chi1\_\allowbreak loc\_\allowbreak isGal}. The retained
$\operatorname{AutSpec}$ identity
$Y.\mathrm{iotaQ}=\operatorname{Sp.iq}$ and its $LL$
reading give the target commuting square.
The equality $D.\mathrm{ntw}.\ell(T.\mathrm{twOf}(a))
=X.\mathrm{nd}.\ell(a)$ gives the same residue prime.
The three actual Hecke-character equalities of GA15
give exactly
$B_r(\mu_{0,j}\delta_{\ell(a)},\eta)$.
Thus this proves that literal endpoint, with all its
other premises preserved, although several are unused
by this character argument.

For any such local realization, GA7 immediately gives
\texttt{IsGalOf} on every set of places where its
specified Hecke factor is unramified and $P\nmid\ell$,
with the exponent in (GA.27). For a finite-order
installed $\Delta$ having the actual cofinite normalized
reading and \texttt{ValuesIn} in $LL_0$, the same
argument first identifies it with the realization of
$\delta_\ell$. At every unramified $P$, choose
$x_P\in LL_0$ reading $\delta_\ell(P)^{-1}$.
Injectivity gives the exact fraction-field equality for
$\Delta(\sigma)$, proving the full \texttt{DeltaRead}
clause, rather than only a cofinite version.
Transport to a local coefficient ring gives $\delta_\ell$
on the true side and $\delta_\ell^{-1}$ on the false
side: the supplied conjugate embedding conjugates its
root-of-unity value, and conjugation equals inversion.
Likewise an installed $\psi_\ell$ identified with the
same chosen $\psi$ gives the literal local
\texttt{NGalChrTied.psiL} equality. These implications
retain the equality of the chosen $\psi$ with the tuple
component; neither a tuple nor a local index is created
by them.

\paragraph{GA23: what can be installed, and the remaining simultaneous construction.}
For every already supplied twist index $i$, if the actual
$LL_0$ is star-stable and contains the HF values, GA16
constructs the entire source record
$\texttt{Triple.Model.Chars}(D,C,\operatorname{Art},i)$.
Choose a reading of its existing coefficient ring, as
in NC and RC8; star stability and the supplied equation
for the two coefficient embeddings give the conjugate
reading required in GA16. The resulting record has its
actual $\Delta$, both $\chi_1$ families, normalized
realization, open kernel, anticyclotomic law, exact
ramification and inertia order $15$, continuity and
the three inertia clauses (GA.21)--(GA.22).
GA14 removes dependence of these realization statements
on that initially chosen reading.

The literal \texttt{exists\_\allowbreak isTripleOf} can consequently
replace these character fields on a supplied twist
skeleton when $\mathrm{nd.Loc}$ is empty and the
specified tame array already satisfies
\texttt{TameTie}. It preserves the supplied set
$\mathrm{Tw}$, its coefficient rings, embeddings and
\texttt{NonzeroIndexPlace}. But its construction sets
all $\mathrm{Sig}$ and $\mathrm{abelian}_2$ to empty
and all form spaces to the zero module. The remaining
triple representation and occurrence assertions then
have empty domains. This is an accurately scoped
conditional endpoint, not a compatible arithmetic
realization of the form spaces used in RC or NC.

The unresolved construction is therefore precise:
one must supply nonempty twist and local index data on
the same $D$, cover every actual coefficient place by
the stipulated side embedding, install the normalized
character readings, and simultaneously satisfy the
form, Hecke, local and tuple identifications.
In particular the literal \texttt{SideB.cover} asserts
$\forall w\,\exists i$ with matching residue
characteristic and true-side valuation.
An empty $\mathrm{Tw}$ cannot satisfy it when a
coefficient place exists, and it cannot supply
$\mathrm{NChoice.twI}$.
GA13 constructs characters \emph{after} $i$ and its
reading have been supplied; it never reverses those
quantifiers. The local character constructions of
GA20--GA22 remove that part of the supplier once its
compatible inputs are present. Global reciprocity,
the nonempty indexing and field-containment/reading
conditions, the genuine representation suppliers and
the same-object $\mathrm{MtReal}/\mathrm{NExpTied}/
\mathrm{ColW}/\mathrm{Pre2}/\mathrm{St}$ joins remain
explicit. No conclusion of RH or GRH is asserted by
this section.


\section{Finite Weil arithmetic and the actual normalized theta periods}
\label{sec:finite-weil-arithmetic}

This section supplies the qualitative finite-Weil assertion retained in
TT13 and TT16.  It treats the actual representation and its actual
normalized vacuum.  In particular, the argument determines an algebraic
phase; preservation of a norm alone would determine only its absolute
square.  The consequences are the actual normalized jet field, and the
actual normalized definite-period field for every supplied bottom
triangle.  The existence of additional full adelic Weil representations
for new bottom spaces is separated at WA17--WA18.

Here are the common hypotheses for WA1--WA11.  Fix the source's CM host
$D$, standard additive character $\psi$ of sign $-1$, its specified
infinity embeddings $\theta$, Hermitian spaces $S,T$, self-dual full
$\mathcal O_E$-lattices $\Lambda_S,\Lambda_T$, and their specified Cartan
involutions at infinity.  Let $v$ be the actual unit vacuum
\texttt{IsVacO} for these data, and let $\omega$ be the actual
\texttt{IsWeilDatumO}, with its given splitting Hecke characters
$\chi_S,\chi_T$ satisfying \texttt{ResEta} at the respective ranks.
Thus Heisenberg covariance, continuity, norm preservation, rational
automorphy and the positive-scalar split normalization are hypotheses
on this particular representation.  We do not infer its existence
from finite arithmetic.  Write $\mathcal W=\operatorname{Hom}_E(S,T)$,
write $b$ for the source's oriented symplectic pairing, and put
\[
 U(x)f(w)=\psi\bigl(b(w,x)/2\bigr)f(w+x),\qquad
 \Lambda=\widehat{\mathfrak a}_D\,
          \operatorname{Hom}_{\widehat{\mathcal O}_E}
                 (\widehat\Lambda_S,\widehat\Lambda_T),
             \qquad\mathfrak a_D=D.\mathrm{host}.\mathrm{AA}.
                                                               \tag{WA.1}
\]
This is exactly the source's \texttt{latO}, including its fractional
ideal factor; the Hom lattice alone would generally be different.
Elements of $\Lambda$ have zero infinite component.  The identity
$\xi_0\mathfrak a_D\overline{\mathfrak a_D}
=\mathfrak D_{E/\mathbb Q}^{-1}$ and the perfect Hermitian
pairings on the two self-dual lattices show, by local dual bases,
that the finite annihilator of $\Lambda$ under $\psi\circ b$ is
exactly $\Lambda$.  This is the trace-dual calculation of VI2,
with the actual \texttt{latO} multiplier and present frames.
Throughout, $\overline{\mathbb Q}$ denotes the actual algebraic closure
of $\mathbb Q$ inside $\mathbb C$, and
\[
 \mathcal A_v=\operatorname{span}_{\overline{\mathbb Q}}
                     \{U(x)v:x_\infty=0\}.                  \tag{WA.2}
\]
Every span expression in this section is a finite sum.

\paragraph{WA1: finite Heisenberg phases and the algebraic span.}
If $x_\infty=y_\infty=0$, choose positive integers $a,b$ with
$ax,by\in\Lambda$.  Such integers exist because a finite adele is
integral away from finitely many places and each remaining lattice
denominator can be cleared by an integer.  Self-duality gives
$\psi(b(ax,by))=1$.  Hence $\psi(b(x,y))$ has finite order dividing
$ab$, and $\psi(b(x,y)/2)$ has finite order dividing $2ab$.
These numbers are algebraic, nonzero and of absolute value one.
The actual operator identity is
\[
 U(x)U(y)=\psi\bigl(b(x,y)/2\bigr)U(x+y),\qquad
 U(-x)U(x)=1.                                                 \tag{WA.3}
\]
Consequently $\mathcal A_v$ contains $v$, is closed under algebraic
linear combinations, and is preserved by every $U(x)$ with finite
$x$.  This uses the half-pairing in (WA.1); dropping its factor of
$1/2$ would change the phases at primes over $2$.

The literal lattice-vacuum law gives one uniquely determined sign
$\varepsilon(\lambda)\in\{1,-1\}$ with
$U(\lambda)v=\varepsilon(\lambda)v$.  Uniqueness uses $v\ne0$,
which follows from its norm one.  The sign is one on $2\Lambda$
and on either isotropic half of $\Lambda$ at a split place.
To obtain the sign for all of $\Lambda$ from these literal clauses,
use that multiplication by $2$ is invertible at every odd place,
while the host makes every place over $2$ split.  The two local
idempotent halves there, together with $2\Lambda$, generate
$\Lambda$.  Their individual operators fix $v$; the product
phases have square one by the integral self-pairing.  This proves
the claimed sign and shows that it factors through $\Lambda/2\Lambda$.
Equation (WA.3) gives the exact semicharacter relation for
$\varepsilon$.  No simultaneous sign choice for unrelated vacua
is used.

\paragraph{WA2: algebraic pairings and rigidity of a scalar.}
Use the source's covolume-one theta inner product, conjugate-linear
in its first variable.  It is positive on nonzero theta functions
and satisfies $\langle v,v\rangle=1$.  We have
\[
 \langle v,U(x)v\rangle=
 \begin{cases}\varepsilon(x),&x\in\Lambda,\\0,&x\notin\Lambda,
 \end{cases}\qquad x_\infty=0.                              \tag{WA.4}
\]
The first case is the vacuum law.  For the second, self-duality
provides $\lambda\in\Lambda$ with $\psi(b(\lambda,x))\ne1$.
Apply $U(\lambda)$ to both entries of the inner product and use
(WA.3).  The two signs cancel, leaving the original inner product
multiplied by $\psi(b(\lambda,x))$; it must therefore vanish.
Heisenberg invariance now expresses
$\langle U(x)v,U(y)v\rangle$ as a phase from WA1 times a value
of (WA.4).  Expansion of two finite sums proves
\[
               F,F'\in\mathcal A_v
       \quad\Longrightarrow\quad
               \langle F,F'\rangle\in\overline{\mathbb Q}.
                                                               \tag{WA.5}
\]
If $F\ne0$ and $F'=cF$, positivity and (WA.5) give
$c=\langle F,F'\rangle/\langle F,F\rangle\in\overline{\mathbb Q}$.
In particular an equality of actual vectors, rather than equality
of their norms, detects the scalar.

\paragraph{WA3: a finite Weil move has an algebraic projective expansion.}
For a finite pair $g=(g_S,g_T)$, there are $F\in\mathcal A_v$ and
$\kappa\ne0$ such that
\[
                         \omega(g)v=\kappa F,
                              \qquad F\ne0.                \tag{WA.6}
\]
At this stage no algebraicity of $\kappa$ is asserted.  We give
the finite-dimensional argument behind the source's signed-average
proof.

First recall the analytic uniqueness used here.  For a self-dual
finite lattice, impose its specified semicharacter law and the
Fock annihilation equations for a positive Cartan complex structure.
Choose a symplectic integral basis of the corresponding rational
trace lattice.  Restriction to a fundamental finite coset then
gives the usual entire theta function with that positive complex
structure.  Fourier expansion along one integral half determines
all coefficients from one coefficient: translating by the opposite
half gives the theta recurrence.  Positive definiteness makes the
resulting Gaussian series normally convergent.  Thus this space has
dimension at most one, just as in VI6.  The change of real coordinates
is used only for this analytic uniqueness; its coefficients need
not be algebraic.  A known nonzero vacuum makes the dimension one.

For completeness, this uniqueness also gives a finite algebraic
model at every deeper level.  Let $V_M$ be the Fock theta functions
fixed by $U(2M\lambda)$ for all $\lambda\in\Lambda$, where $M>0$.
On $V_M$ the operators
$T(\lambda)=\varepsilon(\lambda)U(\lambda)$ form a representation
of the finite group $\Lambda/2M\Lambda$.  The semicharacter identity
proves the representation law, and the factor $2$ makes its kernel
contain $2M\Lambda$.  The finite character projectors are
\[
  P_\alpha=\frac1{|\Lambda/2M\Lambda|}
       \sum_{\lambda\bmod 2M\Lambda}
                          \alpha(\lambda)^{-1}T(\lambda).   \tag{WA.7}
\]
They sum to the identity.  Perfect finite-lattice duality represents
each $\alpha$ by pairing with some
$w\in(2M)^{-1}\Lambda/\Lambda$.  Translating an $\alpha$-eigenvector
by $U(-w)$ gives the original lattice-vacuum law.  Analytic uniqueness
therefore says that each eigenspace is the line generated by the
corresponding $U(w)v$.  This proves that $V_M$ has a finite basis of
such translates and that all Heisenberg operators preserving this
level have algebraic matrices in that basis: they permute these
lines with the roots of unity of WA1 and the lattice signs.

The finite pair $g$ preserves the Fock equations.  One may justify
differentiation here without assuming a linear map on all functions
is continuous: norm preservation makes $\omega(g)$ an isometry on
the theta $L^2$ space; on a fixed smooth finite level, archimedean
Heisenberg difference quotients converge in $L^2$.  Covariance moves
these quotients through that isometry.  The smooth output has the
resulting derivatives, giving the same annihilation equations.
Choose $M$ so that $\omega(g)v\in V_M$ and
$g\Lambda\subset(2M)^{-1}\Lambda$; increase $M$ if needed to have
$M\Lambda\subset g\Lambda$.  Commensurability gives all these
conditions.  The equations
\[
             U(g\lambda)f=\varepsilon(\lambda)f
                         \quad(\lambda\in\Lambda)          \tag{WA.8}
\]
are algebraic linear equations in the finite basis of $V_M$.
Their solution space is the one-dimensional vacuum line for the
moved self-dual lattice, by the same uniqueness argument.  An
intersection of kernels in a finite-dimensional space is determined
by finitely many of its equations.  Row reduction over
$\overline{\mathbb Q}$ therefore supplies a nonzero algebraic
coordinate vector $F$ spanning it.  This is (WA.6).

Equivalently, the source sums $\varepsilon(u)U(u)\phi$ over
$\Lambda/2M\Lambda$.  Its inner product with $v$ is the number of
summands times $\langle v,\phi\rangle$.  A nonzero character
projection supplies a translate with nonzero overlap, and applying
$\omega(g)$ to that average gives (WA.6).  Both proofs retain the
same unknown scalar.

\paragraph{WA4: a principal finite congruence fixes the actual vacuum.}
The simultaneous lattice stabilizer preserves the vacuum line:
covariance preserves the lattice law and its signs after a
sufficiently deep congruence, and a finite move preserves the
Fock equations as in WA3.  Shrink first to a compact open subgroup
on which the signs are preserved.  Thus
$\omega(k)v=c(k)v$ there.  The map $c$ is a character, has absolute
value one, and is continuous: evaluate at any point where $v$ is
nonzero and use the literal continuity clause of
\texttt{IsWeilDatumO}.

A continuous image of a profinite compact group in the circle is
a compact quotient of a profinite group, hence a profinite closed
subgroup of the circle.  The closed subgroups of the circle are
the finite cyclic groups and the whole circle; the latter is not
profinite.  Therefore $c$ has finite image and open kernel.
Principal finite congruences form a neighbourhood basis in this
lattice stabilizer.  Some positive integer $N$ consequently gives
\[
 \omega(k_S,k_T)v=v
       \quad\hbox{whenever }k_S,k_T\hbox{ satisfy the actual
                                      level-}N\hbox{ congruence}.
                                                               \tag{WA.9}
\]
Taking a multiple of $N$ accommodates both spaces and all lattice
denominators.  This proves the generic
\texttt{congruent\_\allowbreak pairs\_\allowbreak fix\_\allowbreak vacuum\_\allowbreak weilDatum}; it is stronger
than merely asserting that such pairs preserve a line.

\paragraph{WA5: a nonzero split-half coinvariant functional.}
Let $e$ be the actual split-place idempotent, put $\bar e=c(e)$,
and let $Y=\bar e\mathcal W(\mathbb A_f)$.  There exists a complex
linear functional $\ell$ on the ambient function space such that
\[
 \ell(v)=1,\qquad
 \ell(U(y)\phi)=\ell(\phi)quad(y\in Y, \phi\hbox{ theta}).
                                                               \tag{WA.10}
\]
No continuity of $\ell$ is needed or claimed.

Here is an algebraic proof that the required coinvariant is nonzero.
For $y\in Y\setminus\Lambda$, self-duality supplies
$\lambda\in e\Lambda$ with
$\zeta=\psi(b(\lambda,y))\ne1$.  Since $U(\lambda)v=v$,
the operator $U(\lambda)-\zeta$ kills $U(y)v$ and sends $v$ to
$(1-\zeta)v$.  Operators obtained in this way commute, since
$e\mathcal W$ is isotropic.  Multiplying one for each member of a
finite list shows that $v$ is outside the span of the translates
with $y\notin\Lambda$.  Linear extension gives a functional $q$
with
\[
                 q(U(y)v)=1_{Y\cap\Lambda}(y).              \tag{WA.11}
\]

Let $D_Y$ be the span of $U(y)\phi-\phi$ for theta functions
$\phi$ and $y\in Y$.  For each of a finite list of such differences,
smoothness at finite places gives an integer $n>0$ with
$U(ny)\phi=\phi$.  The sum
$A_y=\sum_{a=0}^{n-1}U(ay)$ kills that difference by telescoping.
All these sums commute.  Their product kills the finite list,
whereas (WA.11) evaluates its value on $v$ as a positive integer:
every summand is zero or one, and the all-zero summand is one.
Thus $v\notin D_Y$.  Take a linear functional on the quotient
which is one at $v$, and extend it to the ambient space.  This
proves (WA.10), including its invariance on every theta function.

\paragraph{WA6: make that functional algebraic on a prescribed finite list.}
For any finite list $x_1,\ldots,x_r$ with zero infinite component,
the functional in WA5 can be chosen so that
\[
                         \ell(U(x_i)v)\in\overline{\mathbb Q}
                                      \quad(1\leq i\leq r). \tag{WA.12}
\]
Write $q_0=1-e-\bar e$ and decompose
$x=ex+\bar ex+q_0x$.  If $q_0x\in\Lambda$, the invariant
$\bar e$-part and the lattice $q_0$-part can be removed using
(WA.3), with algebraic phases and a lattice sign.  For the
remaining $ex$, the value is one when $ex\in\Lambda$.  Otherwise
some $\lambda\in\bar e\Lambda$ has nontrivial commutator with
$ex$, and (WA.10) forces the value to be zero.  Thus (WA.12)
already holds in this case.

If $q_0x\notin\Lambda$, choose $u\in q_0\Lambda$ with
$\zeta=\psi(b(u,x))\ne1$, and put
$\epsilon=\varepsilon(u)$.  Replace $\ell$ by
\[
       \ell'=
       \frac{\ell\circ U(u)-\epsilon\zeta\ell}
                          {\epsilon-\epsilon\zeta}.        \tag{WA.13}
\]
The denominator is nonzero.  The replacement is still one on
$v$, and still invariant under $Y$, since $u$ is orthogonal to
$Y$.  It is zero on $U(x)v$.  For every finite $z$ its value on
$U(z)v$ is the previous value multiplied by
\[
           \frac{\epsilon\psi(b(u,z))-\epsilon\zeta}
                          {\epsilon-\epsilon\zeta}
                         \in\overline{\mathbb Q}.           \tag{WA.14}
\]
Thus already algebraic values stay algebraic.  Applying this
operation successively to the finite list proves the assertion.
The functional depends on the list, as permitted by the source;
no single algebraic functional on an infinite-dimensional space
is being assumed.

\paragraph{WA7: the splitting characters have algebraic finite values.}
For either actual splitting character $\chi$, every idele $a$
with $a_\infty=1$ satisfies $\chi(a)\in\overline{\mathbb Q}$.
The restriction condition \texttt{ResEta} makes the infinity
character trivial on positive real scalars.  Continuous circle
characters then give integers $m_\tau$ such that its infinity
part is $\prod_\tau(z_\tau/|z_\tau|)^{m_\tau}$.  At an embedded
algebraic number each factor is algebraic, because
$|z|^2=z\bar z$ is algebraic and $|z|$ is its positive square
root.  The finite integral idelic units are profinite compact;
the same circle argument as WA4 makes the character's image on
them finite, hence roots of unity.

The fractional ideal of $a$ has finite order in the ideal class
group.  Raising $a$ to that positive order and dividing by a
principal idele leaves an integral finite unit and the inverse
archimedean part of the principal idele.  Triviality on principal
ideles and the preceding two facts show $\chi(a)^h$ algebraic
for some $h>0$.  Therefore $\chi(a)$ is algebraic.  This is a
pointwise statement.  It does not put the raw values of a
half-normalized unitary character at infinitely many primes in
one number field; the source's normalized \texttt{ValuesIn}
statements retain their precise norm powers.

\paragraph{WA8: the split normalization determines the phase.}
Let $g=(g_S,g_T)$ be supported at the split pair represented by
$e+\bar e$.  Choose the exact determinant lifts $x,y$ in
\texttt{IsDetAt}.  Their infinite components are one.  The
\texttt{om\_\allowbreak split} clause of this actual representation says that
there is $r>0$ such that every functional as in (WA.10) satisfies
\[
 \ell(\omega(g)\phi)=r z\ell(\phi),\qquad
                   z=\chi_T(x)^{-1}\chi_S(y).               \tag{WA.15}
\]
WA7 gives $z\in\overline{\mathbb Q}$, and $\bar z z=1$.
Write $\omega(g)v=\kappa F$ as in WA3.  Unitarity and WA2 give
\[
            \bar\kappa\kappa=
                 \langle v,v\rangle/\langle F,F\rangle
                                  \in\overline{\mathbb Q}.  \tag{WA.16}
\]
Choose the functional in WA6 for the finite translates in $F$,
so $a=\ell(F)$ is algebraic.  Applying (WA.15) to $v$ gives
$\kappa a=rz$.  In particular $a\ne0$.  Taking absolute squares
and using $\bar z z=1$ yields
\[
                  r^2=(\bar\kappa\kappa)(\bar a a).
                                                               \tag{WA.17}
\]
Thus $r$ is algebraic, and $\kappa=rz/a$ is algebraic.  This
proves $\omega(g)v\in\mathcal A_v$.  The positive real scalar
and the exact character factor in (WA.15) are essential: (WA.16)
alone gives no such conclusion.

\paragraph{WA9: the algebraic-action subgroup is closed under positive roots.}
Let $H$ be the set of finite pairs $g$ for which
$\omega(g)v\in\mathcal A_v$.  If
$\omega(g)v=\kappa F$ as in WA3, covariance and WA1 give
\[
                    \omega(g)\mathcal A_v\subset\kappa\mathcal A_v.
                                                               \tag{WA.18}
\]
Indeed $\omega(g)U(x)v=U(gx)\omega(g)v$, and $gx$ is finite.
The identity and products belong to $H$.  For an inverse write
$\omega(g^{-1})v=\kappa F$.  If $g\in H$, then
$v=\kappa\omega(g)F$ with both $v$ and the nonzero
$\omega(g)F$ in $\mathcal A_v$.  WA2 makes $\kappa$ algebraic,
so $g^{-1}\in H$.

Iterating (WA.18) gives
$\omega(g^m)v=\kappa^mF_m$ with $F_m\in\mathcal A_v$ nonzero.
If $m>0$ and $g^m\in H$, WA2 makes $\kappa^m$ algebraic;
then $\kappa$ is algebraic, since it is a root of a polynomial
over a finite extension of $\mathbb Q$.  Hence
\[
                      m>0,\quad g^m\in H
                                  \quad\Longrightarrow\quad g\in H.
                                                               \tag{WA.19}
\]
WA4 puts a simultaneous principal congruence inside $H$, and
WA8 puts every split-supported finite pair inside $H$.

\paragraph{WA10: local quasi-reflections generate the remaining finite points.}
We prove directly the group statement used by the source's
\texttt{finiteUnitaryPoints\_\allowbreak mem\_\allowbreak of\_\allowbreak level\_\allowbreak and\_\allowbreak split\_\allowbreak place}.
Let $H_0$ be a subgroup of the finite unitary points which contains
a principal congruence of positive integer level $N$, contains
every split-place supported point, and satisfies positive-power
closure (WA.19).  Then $H_0$ is the whole finite unitary group.

At a nonsplit finite place, $E_w/F_v$ is a quadratic field
extension.  For a nonisotropic vector $a$, the Hermitian projection
\[
            P_a(z)=\frac{h(z,a)}{h(a,a)}a
\]
satisfies $P_a^2=P_a=P_a^*$.  If $u\bar u=1$, then
$R_{a,u}=1+(u-1)P_a$ is unitary and
\[
                       R_{a,u}^{m}=1+(u^m-1)P_a.           \tag{WA.20}
\]
These quasi-reflections generate the local unitary group.  To
verify this, take $x,z$ with the same nonzero norm.  If
$a=z-x$ is nonisotropic, put $A=h(z,a)$.  Then
$h(a,a)=A+\bar A\ne0$, so $A\ne0$, and
$u=-\bar A/A$ has norm one.  Substitution gives
$R_{a,u}z=x$.  If $z-x$ is isotropic, $z+x$ has norm
$4h(x,x)\ne0$.  The preceding construction sends $z$ to $-x$,
and $R_{x,-1}$ sends $-x$ to $x$.  A nondegenerate Hermitian
space has an orthogonal basis of nonisotropic vectors, by the
usual elimination of cross terms.  Fix its first vector using
at most two such reflections and induct on its orthogonal
complement.  This proves generation, including rank one; rank
zero is trivial.

Because conjugation preserves the local valuation,
$u\bar u=1$ gives $v_w(u)=0$.  Every finite quotient
$\mathcal O_{E_w}^{\times}/(1+\mathfrak p_w^b)$ is finite.
Some positive power $u^m$ is therefore congruent to one to any
prescribed depth.  Choose that depth large enough to clear the
matrix denominators of $P_a$ relative to the fixed local lattice
and to meet level $N$.  Formula (WA.20) then puts $R_{a,u}^m$
in the prescribed principal congruence; its inverse is integral
and has the same congruence by the convergent geometric-series
inverse of a sufficiently small integral matrix perturbation.
Power closure puts $R_{a,u}$ itself in $H_0$.  No condition that
$m$ be prime to the residue characteristic is needed.  This
argument therefore includes ramified and dyadic places.

Finally any finite adelic point stabilizes the fixed lattice
away from a finite set of places.  Enlarge that set by the places
dividing $N$.  Removing its finitely many single-place factors
leaves a member of the principal congruence, since $N$ is a
local unit everywhere else.  Split factors are already in
$H_0$; nonsplit factors are products of the reflections just
treated.  Their product proves the assertion.  This local proof
replaces the source's longer rational/split/congruence
factorization, without changing its group-theoretic conclusion.

\paragraph{WA11: actual finite Weil operators preserve the algebraic span.}
Apply WA10 to the inverse image of the subgroup in WA9 under
$g_S\mapsto(g_S,1)$ and then under $g_T\mapsto(1,g_T)$.
All its hypotheses were proved in WA4, WA8 and WA9.  Thus, for
every actual finite pair,
\[
     \omega(g_S,g_T)v\in\mathcal A_v,
          \qquad \omega(g_S,g_T)\mathcal A_v=\mathcal A_v.   \tag{WA.21}
\]
The equality follows from covariance and the same assertion for
the inverse.  This is the literal conclusion of
\texttt{BotExpand.om\_\allowbreak vac\_\allowbreak algComb}.  The source's
\texttt{moved\_\allowbreak translate\_\allowbreak algComb} follows by applying it to
$U(x)v$ and using (WA.3), with every resulting translation still
having infinite component zero.

The proof uses the displayed generic data; it does not use the
nominal \texttt{FunPr} or universal \texttt{NormSum} arguments
carried by the bottom wrappers.  It proves qualitative
algebraicity of finite coefficients.  It does not prove that an
arbitrary distinguished dictionary scalar has absolute value
one under every algebraic embedding, or is a unit at specified
finite primes.  Those stronger assertions in
\texttt{ExistsUnitModulusScalarLocal} remain distinct.

\paragraph{WA12: both actual bottom datum functions are algebraic combinations.}
Fix the finite half-space and the completed half-lattice in a
literal \texttt{IsBottomTri}.  Each of its two standard finite
functions is algebraic-valued.  Its determinant character is
evaluated on an integral finite unit, hence has finite order by
the profinite argument of WA7; its remaining factors are
indicators.  On a suitable bounded fractional half-lattice the
function is periodic by $k$ times the integral half-lattice,
for one positive integer $k$.  It is consequently a finite
linear combination of small-coset indicators.

For one such indicator, finite character orthogonality gives
\[
  1_{a+L_0}(x)=\frac1{|L/L_0|}
        \sum_{\alpha\in(L/L_0)^\vee}\alpha(x-a)
                                      \quad (x\in L),       \tag{WA.22}
\]
where $L$ contains the bounded support and $L_0$ is a common
period lattice.  The dual half-lattice realizes these characters
by modulation.  Translation and modulation are the two
isotropic Heisenberg actions; converting their product to one
$U(w)$ introduces exactly the algebraic phase of WA1.
The actual theta map is linear, sends the prescribed basic
half-lattice indicator to the given vacuum, and intertwines
these translations and modulations.  Thus (WA.22) proves
\[
             X.f_t\in\mathcal A_{X.\mathrm{vac}},\qquad
             X.f'_t\in\mathcal A_{X.\mathrm{vac}'}.          \tag{WA.23}
\]
On the second side use the source's conjugate idempotent and its
own actual character and vacuum.  No Fourier transform is
silently substituted for that second standard function.

WA11 applies separately to the two actual oriented Weil data,
so (WA.23) remains true after every finite pair acts.  Finite
Heisenberg operators commute with archimedean polynomial
creation, by covariance at disjoint places and the $L^2$
difference-quotient justification in WA3.  Therefore every
monomial bottom kernel is a finite algebraic sum of finite
translate jets of its own vacuum.  This proves the actual
\texttt{theta\_\allowbreak kernel\_\allowbreak algebraic\_\allowbreak expansion}, its dual, and
\texttt{ker'\_\allowbreak eq\_\allowbreak sum\_\allowbreak jets} with their original orientations.

\paragraph{WA13: the actual normalized grid jets are algebraic.}
Fix one literal \texttt{Choice}, its grid Weil representations,
normal datum and rational frames.  Every edge used in TT13 has
the self-dual lattices and actual unit vacuum of TT12, with the
specified Cartan structures.  WA11 supplies (TT.19) for this
very representation.  For each pure coset datum, the same
finite Fourier calculation as WA12 supplies its vector as an
algebraic finite sum of translates of that vacuum.

The period model of that normal datum is the one already used
in TA10 and VI9--VI11.  Finite additive approximation identifies
a finite-adelic translate with a rational archimedean translate,
up to the lattice law and a root of unity from WA1.  Thus its
monomial jet has the exact form
\[
 c_{0,p}\prod_\tau
     \operatorname{opp}_p(\operatorname{period}(\delta,\tau))
                         ^{d_\tau}\,a,
                      \qquad a\in\overline{\mathbb Q}.      \tag{WA.24}
\]
Here $c_{0,p}$ is the actual nonzero comparison scalar of the normal datum;
it is not asserted algebraic.  The half-point, $\delta$, vacuum
and frame in (WA.24) are unchanged.  TT13 proves that the
homogeneous degree of the actual differential operator factors
out exactly the period powers occurring in
$Z.\mathrm{scal}(1,p)$, whose definition is their product with
$c_{0,p}$ inverted.  They are nonzero and cancel.  TT9's rational
polynomial and bideterminant operations then give
\[
                  Z.\mathrm{kk}(d,x,y)_{ab}
                                \in\overline{\mathbb Q}     \tag{WA.25}
\]
for every actual required entry.  This supplies the previously
retained arithmetic premise of TT13.  It uses the direct
rational-translate germ theorem, not a claim that arbitrary
algebraic entire functions are stable under translation.

\paragraph{WA14: the literal normalized jet field now exists.}
For a fixed choice, TT6--TT8 give finite support labels and
finite class sets.  At the least depth for each label, TT7
shows that the support of every allowed expansion is contained
in one finite set.  Taking all twists, top and row edges and
matrix indices therefore gives a genuinely finite set
$E_{\rm jet}$ of all entries demanded by \texttt{JetsIn}.
WA13 makes every member algebraic.  Consequently
\[
                   M_{\rm jet}=\mathbb Q(E_{\rm jet})
\]
is a number field, and its inclusion into $\mathbb C$ satisfies
the literal \texttt{JetCl} for the original choice.  Combining
this with TT8 gives \texttt{NFJet} with no new size threshold:
the two additional thresholds may be taken to be zero.  This
is a qualitative field-existence result; it gives no uniform
degree, height or denominator bound.  In particular it does
not replace the quantitative clauses of \texttt{DictAt}.

\paragraph{WA15: one fixed-degree period for the actual bottom vacua.}
Let $X$ be a literal bottom triangle and let $d$ be a fixed
place-degree vector.  Its positive spaces have standard Gram
matrices, integral lattices and identity orthonormal frames.
The exact square-rank identities make the total target rank
$\operatorname{rank}WW=\operatorname{botBig}+\operatorname{botSmall}$
even.  Hence both products
$\operatorname{rank}A\operatorname{rank}WW$ and
$\operatorname{rank}A'\operatorname{rank}WW$ are even.  We use
this product parity, not an assertion that both source ranks
are even.

Conjugating the negative source vacuum gives a positive vacuum
for $A,WW$.  Choose the same standard CM curve datum and half-point
as in the earlier theta construction.  The positive period-model
function and this actual vacuum satisfy the same lattice and Fock
laws and are both nonzero.  Analytic uniqueness from WA3 therefore
gives one constant $c\in\mathbb C^{\times}$ relating them.  The
period-normalized entire function need not have unit $L^2$ norm,
so no absolute-value normalization of $c$ is asserted.
VI9--VI11 and the exact Fock derivative formula then give the
common nonzero period
\[
          \Omega_d=\overline{c\prod_\tau
                      \operatorname{period}(\delta,\tau)^{d_\tau}}
                                                               \tag{WA.26}
\]
for all finite translates and all monomials of that place degree.
The negative-space identity conjugates the whole positive jet,
including $c$ and the period product.  On the positive source
$A'$ use its own nonzero comparison scalar $c'$ and the unconjugated
product.  All periods are nonzero; degree zero and rank zero
cause no division by a zero quantity.

WA12 expresses every finite-point kernel as an algebraic sum
of precisely these translate jets.  Division by (WA.26) therefore
makes every monomial kernel of degree $d$ algebraic.  This proves
the actual \texttt{period\_\allowbreak algebraic\_\allowbreak fixed\_\allowbreak degree} and both
finite-point kernel arithmetic statements without an independent
\texttt{NormSum.alg} or finite-Weil arithmetic premise.  The
comparison scalars are absorbed in the periods; no algebraicity
of $c,c'$ is needed.

\paragraph{WA16: the actual definite integrals have a common number field.}
For every supplied bottom triangle $X$, there is a number field
$M\subset\mathbb C$ and nonzero periods $\Pi,\Pi'$ such that
the literal \texttt{SumsIn} holds.  We now discharge TT16's
remaining finite-Weil premise by WA11--WA12 and its fixed-degree
period input by WA15.

The argument retains the actual measures.  A common compact open
finite level fixes the finitely many standard theta data, using
WA4 and the finite Heisenberg expansions.  Polynomial creation
commutes with that level, so it is independent of the monomial.
The definite adelic quotient is compact and has finitely many
classes at this level.  Its exact integral decomposition is
\[
 I_t(P,h)=v_0\sum_l|\Gamma_l|^{-1}
       \int_C\operatorname{ker}_t(P,h)(z_lk)\,dk,
                 \qquad v_0>0,                             \tag{WA.27}
\]
with probability Haar measure on the compact real group $C$.
The constant $v_0$ is the original compact-open volume; it need
not be rational.  The continuous kernels are integrable on this
compact quotient.  TT15 identifies the compact projector over
$\overline{\mathbb Q}$ and the single surviving degree $d_0$,
using the exact determinant weight.  WA15 then proves
$I_t(P,h)/(v_0\Omega_{d_0})\in\overline{\mathbb Q}$ for all
finite $h$ and monomials $P$.  The statement also holds when all
these integrals are zero; the chosen period remains nonzero.

To obtain one finite field, rather than infinitely many algebraic
values, use the finite generation proved in TT16.  A common
finite level and rational invariance reduce the finite points to
finitely many class representatives.  The compensating infinite
matrices are embeddings of rational $E$-matrices.  Fock
substitution and integral determinant powers have coefficients
in one fixed finite normal extension of $E$.  The central
character equation permits only finitely many degree vectors
and finitely many monomials in them.  Thus every integral is a
linear combination over one number field $M_0$ of a finite list
$I_t(P_b,h_a)$.  Adjoin their normalized values, already known
algebraic.  Repeat the argument for $I'_t$ with its own measure
and period.  The compositum is finite and supplies exactly
\texttt{SumsIn}, with the same $X$ on both sides of the join.

\paragraph{WA17: the standard bottom spaces and the precise remaining existence.}
For each square $(i,j)$ the source defines the two nonnegative
ranks \texttt{botBig}, \texttt{botSmall}.  Take $A=E^{\rm botBig}$,
$A'=E^{\rm botSmall}$ with Gram matrix one and lattice
$\mathcal O_E^n$, and take $WW=A\oplus A'$ with those same
standard choices.  Testing the Hermitian dual against the
coordinate vectors shows that each standard lattice is
self-dual.  The Cartan involutions and orthonormal frames are
identities; the direct-sum rank identity and the product parity
in WA15 follow from the literal square-rank formulas.

The bottom characters are the source's parity corrections of
$\mu_t(j+1)$ and $\chi_{m-i}$.  Multiplication by the prescribed
power $\eta^{\operatorname{rank}V\bmod2}$ gives their exact
\texttt{ResEta} ranks; this is addition of the integral
restriction exponents modulo two.  It does not permit changing
the finite-order twist or choosing a new $\eta$ independently.
Given the two full oriented Weil representations for
$(\operatorname{neg}A,WW)$ and $(A',WW)$ with these characters,
choose the normalized vacua, the original finite split
idempotent, the actual conductor exponents, and the two theta
maps for the standard finite functions.  Choose Haar measures
and fundamental domains for these actual groups.  These choices
give the literal bottom triangle: its two rigs, target rig,
representations, vacua, idempotent, conductors, finite functions,
measures and domains have exactly the stipulated identities.
WA16 then gives its period field.  A finite compositum over all
squares supplies \texttt{NFBot} under this exact existence input.

The new representations here are not the already supplied grid
representations.  Their full existence is not proved in
WA1--WA16.  The next paragraph narrows it to a canonical
splitting-character case by an explicit construction.

\paragraph{WA18: change splitting characters by a determinant-root character.}
Suppose a full oriented Weil datum exists for the same $S,T$
with characters $\eta^{\operatorname{rank}S}$ and
$\eta^{\operatorname{rank}T}$, where $\eta$ is the same
unramified splitting character.  Then a full datum exists for
any of the prescribed characters $\chi_S,\chi_T$ with the
required rank restrictions.  Put
\[
 \nu_S=\eta^{-\operatorname{rank}S}\chi_S,\qquad
 \nu_T=\eta^{-\operatorname{rank}T}\chi_T.
                                                               \tag{WA.28}
\]
They are Hecke characters trivial on the ideles of $F$.
For a unitary element $g$, choose a determinant root $x$ with
$x/\bar x=\det g$.  Hilbert 90 supplies such roots locally.
Outside a finite set the extension and $\nu$ are unramified.
There a norm-one unit has a unit root: in the quadratic field
case rescale any root by the appropriate power of a base-field
uniformizer, whose extension valuation is one; in the split
case use the explicit root $(t,1)$.  These roots assemble to
an idele.  Moreover $\widetilde\nu$ is trivial on the integral
unitary group at every such place, because $\nu$ is trivial on
the local integral units.
Two choices differ by an $F$-idele, so $\nu(x)$ is independent
of the choice.  This defines a continuous unitary character
$\widetilde\nu(g)$.  Multiplicativity follows by using the
product of two roots; rational unitary elements have rational
roots by the field form of Hilbert 90, so this character is
trivial on rational points.  Continuity on the restricted product
now reduces to the finitely many remaining factors.  There it
follows from local
sections of $x\mapsto x/\bar x$: near one solve the equation
by the nonzero expression $1+t$ for $t\bar t=1$, with a
fixed alternative chart at $t=-1$.

On the same theta functions set
\[
 \omega'(g,h)=\widetilde\nu_T(g)^{-1}
                          \widetilde\nu_S(h)\,\omega(g,h).
                                                               \tag{WA.29}
\]
The two characters are deliberately crossed.  Their scalars
commute with Heisenberg operators, have absolute value one,
are continuous, and are trivial at rational pairs.  Therefore
representation, covariance, automorphy, continuity and norm
preservation all hold.  At a split pair the determinant roots
can be the literal \texttt{IsDetAt} lifts, so the scalar in
(WA.29) changes the split factor from
$\eta(x)^{-\operatorname{rank}T}\eta(y)^{\operatorname{rank}S}$
to $\chi_T(x)^{-1}\chi_S(y)$, with the same positive real
multiplier.  This verifies the last defining clause, rather than
merely asserting that a twist is harmless.

Thus the remaining full-Weil supplier can be taken to be
\texttt{exists\_\allowbreak weilDatum\_\allowbreak etaTilde}: the actual construction
of its finite and real actions, their commutation and rational
theta-distribution automorphy.  The source constructs these;
their reconstruction is a separate target.  Nothing in this
section treats that endpoint as an axiom-free consequence of
the character adjustment.

\paragraph{WA19: the same-choice field consequence and its exact boundary.}
For a fixed literal choice, WA14 now supplies its jet field.
For each supplied bottom triangle, WA16 supplies its two-period
field.  Taking their finite compositum with the character,
root and phase fields of TT11 gives these portions of
\texttt{NumbersIn} on the same choice and the same embedded
complex realization.  If bottom Weil data are supplied through
WA17--WA18 and the canonical case retained there, the finite
compositum covers all squares as demanded by \texttt{NFBot}.

Critical toric $L$-value ratios are not finite Heisenberg
coefficients.  Their remaining source input is
\texttt{ToricRatioConjugates} with its actual CM field data
and reference twist.  A field containing those ratios can be
joined with the fields above, but its existence is not inferred
here from (WA.21).  Likewise, membership in a subfield of
$\mathbb C$ is distinct from membership in the range of the
already chosen embedding $\sigma:K\hookrightarrow\mathbb C$.
The literal \texttt{Choice.FieldOK} still requires its same
normalization field, frame, Hida setting, family and specified
embedding to agree.  TT19's compatibility argument applies
once those actual joins and the remaining ratio field are
supplied.  No replacement choice or independent collection of
embeddings is silently substituted for them.

The substantive new endpoints are consequently (WA.21),
\texttt{NFJet}, and \texttt{SumsIn} for the actual supplied
bottom data.  The quantitative all-embedding and finite-unit
dictionary estimates, the canonical full bottom-Weil existence,
the critical ratio field, and the remaining compatible-choice
join retain their stated source scopes.


\section{The unramified Hermitian mass ratio}
\label{sec:mass-formula}

\status{MF1--MF20 reconstruct the actual local densities, the local
orbit identity, the singular-series average, and the rank induction
behind \texttt{SiegelMassRatioCM}. The only named external theorem in
the global representation step is the classical Hasse--Minkowski
local-to-global theorem for quadratic forms over number fields.
Its statement and exact use are recorded in MF4; its proof is
not claimed here. No mass formula, Tamagawa-number theorem,
representation-average formula, or quantitative prime estimate is
assumed. The mass premise retained in QW3 and QW9 is reduced to this
explicitly identified classical theorem. The separate geometric
volume and compatible arithmetic-data premises remain.}

Write \(E\) for the source field \(K\), \(F=E^{\bar{\ }}\) for its
fixed field, \(f=[F:\mathbb Q]\), and
\(v=|\operatorname{disc}E|^{1/(2f)}/2\).
All Hermitian forms are conjugate linear in the first variable.
An integral lattice is a full \(\mathcal O_E\)-lattice on which
all pairings are integral; it need not be free globally. The mass is
the literal sum \(\sum_{[L']}|\operatorname{Aut}(L')|^{-1}\) over the
genus in the given rational Hermitian space, as in IG9.
At a finite place \(\mathfrak p\) of \(F\), put
\[
 R=\widehat{\mathcal O}_{F,\mathfrak p},\quad
 B=\mathcal O_E\otimes_{\mathcal O_F}R,\quad
 q=N\mathfrak p,\quad \eta_{\mathfrak p}=
 \begin{cases}1,&B=R\times R,\\-1,&B/R\text{ is unramified quadratic}.
 \end{cases}
\]
Here \(\pi\) is an actual local uniformizer and
\(\operatorname{ord}_{\mathfrak p}\pi=1\).
The source's congruences and localizations are over the fixed-field
prime \(\mathfrak p\); a split pair of primes of \(E\) is treated
together. No rational prime is identified with \(\pi\).

\paragraph{MF1: the fixed field, the different, and the normalizations.}
The literal \texttt{IsCMUnram} hypothesis gives complex conjugation
under every embedding and, at every prime \(P\) of \(E\), an integral
\(\theta\) with \(\theta-\bar\theta\notin P\).
Choose a nonzero \(w\) with \(\bar w=-w\). Such a \(w\) exists because
otherwise every embedded element of \(E\) would be real.
Every \(z\in E\) has the unique decomposition
\[
 z=\frac{z+\bar z}{2}
       +w\frac{z-\bar z}{2w},\qquad E=F\oplus Fw.       \tag{MF.1}
\]
Thus \(E/F\) is quadratic, \(F\) is totally real, and \(E\) is totally
imaginary. At a fixed finite prime the element
\(\theta-\bar\theta\) is a unit at both conjugate primes.
The coefficient formulas in DR1, applied to \(1,\theta\), show that
the integral quadratic algebra over \(R\) has this basis. Its
discriminant \((\theta-\bar\theta)^2\) is a unit. Hence it is either
split or unramified quadratic, also over primes above \(2\).
Its trace pairing is perfect and its relative different is the unit
ideal. Taking determinants of the successive trace pairings gives
\[
 |\operatorname{disc}E|=(\operatorname{disc}F)^2,\quad
 \operatorname{disc}F=(2v)^f,\quad
 \operatorname{covol}(\mathcal O_E)=v^f,\quad
 \operatorname{covol}(\mathcal O_F)=\sqrt{(2v)^f}.       \tag{MF.2}
\]
The complex covolume uses ordinary Lebesgue measure on one copy of
\(\mathbb C\) per conjugate pair: it is
\(2^{-f}\sqrt{|\operatorname{disc}E|}\). The real covolume is
\(\sqrt{\operatorname{disc}F}\).

Set
\[
 A_m(\mathfrak p)=\prod_{j=1}^{m}(1-\eta_{\mathfrak p}^{\,j}q^{-j}),
 \quad e_m(\mathfrak p)=1-\eta_{\mathfrak p}^{\,m}q^{-m},
 \quad L_m=\prod_{\mathfrak p}e_m(\mathfrak p)^{-1}.     \tag{MF.3}
\]
For \(m\ge2\), this product is positive and absolutely convergent
after taking logarithms, by QW1. It is exactly
\(\texttt{lValEta}(E,m)\): an inert prime of \(F\) has one prime of
\(E\) of norm \(q^2\), while a split prime gives two factors with
exponent \(-1/2\). Their product is (MF.3), with no residual
square-root factor.

\paragraph{MF2: local Hermitian diagonalization, including dyadic places.}
The norm \(B^\times\to R^\times\) is onto.
In the nonsplit case first solve the norm equation in
\(\mathbb F_{q^2}^\times\). If a unit has the desired norm modulo
\(\pi^j\), multiply it by \(1+\pi^j t\); the norm changes modulo
\(\pi^{j+1}\) by \(1+\pi^j\operatorname{Tr}(t)\).
The residue trace is onto, even in characteristic \(2\), so successive
corrections converge. In the split case the norm is \((a,b)\mapsto ab\).
This is the norm-lifting argument of DR1; it never divides by \(2\).

Every integral nondegenerate Hermitian lattice over nonsplit \(B\)
has an orthogonal basis with Gram matrix
\(\operatorname{diag}(\pi^{a_1},\ldots,\pi^{a_m})\), \(a_i\ge0\).
Choose a pairing of least valuation \(a\).
If no diagonal entry has that valuation, replacing a vector by
\(x+t y\) produces one: the leading cross term is a nonzero residue
trace for a suitable \(t\). Projection along this vector has integral
coefficients, since all pairings are divisible by \(\pi^a\), so
induction splits off the line. Unit-norm surjectivity removes each
diagonal unit. For split \(B\), the two components of a lattice are
free \(R\)-modules paired by a matrix; Smith reduction gives the same
description, with independent bases in the two components.
The multiset of \(a_i\) is intrinsic, as the elementary divisors of
the map to the dual lattice.

Over the fraction algebra, a nonsplit form has a basis with all
exponents \(0\), except possibly one exponent \(1\).
First reduce exponents modulo \(2\). Two entries \(\pi,\pi\)
can be replaced by \(1,1\): choose a unit \(t\) with
\(N(t)=\pi-1\); then
\(\pi N(\pi^{-1})+\pi N(\pi^{-1}t)=1\).
Split off this norm-one vector. The complementary coefficient has
even valuation, hence is a norm and can also be normalized to \(1\).
In the split case every fraction-field form is standard.
Consequently an integral lattice of determinant valuation at most
one is determined up to integral isometry by its rank and determinant
valuation. If a fixed vector has norm one, keep it as the first
basis vector and choose such a minimal lattice in its orthogonal
complement. This proves the local construction used in
\texttt{exists\_\allowbreak lattice\_\allowbreak vd\_\allowbreak le\_\allowbreak one}.

\paragraph{MF3: exact norm orbits and rational local approximation.}
For vectors \(x,y\) of the same nonzero norm, put \(u=x-y\).
If \(c=h(u,x)\ne0\), the map
\[
             z\longmapsto z-u\,c^{-1}h(u,z)            \tag{MF.4}
\]
is an isometry taking \(x\) to \(y\). Indeed
\(c+\bar c=h(u,u)\), which cancels the three extra terms on
expanding the pairing. If \(c=0\), pass through \(-y\);
the corresponding denominator is \(2h(x,x)\ne0\), and multiplication
by \(-1\) is an isometry. Thus the unitary group is transitive on
each nonzero-norm fiber. Over a positive definite global space,
the first denominator can vanish only when \(x=y\).

The rational unitary group is dense in its group over any one
completion, and simultaneously over any prescribed finite set.
On skew-adjoint matrices \(Z^*=-Z\), where here the star is the
adjoint for the fixed Hermitian form, the Cayley transform
\[
             Z\longmapsto (1-Z)(1+Z)^{-1}              \tag{MF.5}
\]
is rational, with inverse \((1-g)(1+g)^{-1}\) when
\(\det(1+g)\ne0\). The skew-adjoint matrices form an \(F\)-vector
space, so ordinary weak approximation applies to their coordinates.
To enter this chart near an arbitrary \(g\), multiply \(g\) by a
norm-one scalar arbitrarily close to \(1\) avoiding the finitely
many forbidden eigenvalues. Such scalars are supplied by the
one-dimensional Cayley transform on \(Fw\); weak approximation on
\(Fw\) chooses one rational scalar for all the specified places.
Apply weak approximation
to the resulting chart coordinates and then undo the scalar.
This proves the assertion without strong approximation.

Lattice stabilizers are open. A completed isometry carrying one
rational localized lattice to another can therefore be replaced
by a rational isometry with exactly the same lattice action.
The same argument applies to the orthogonal complement in (MF.4).
Thus the completed orbits below are exactly the rational local
orbits in \texttt{fibLoc}, \texttt{IsLocReps}, and
\texttt{SameGenus}. Those definitions do not merely ask for completed
isometries.

\paragraph{MF4: the binary representation theorem and its classical input.}
We use the following named classical theorem:
\emph{a nondegenerate quadratic form over a number field has a nonzero
isotropic vector if it does over every completion}.
This is Hasse--Minkowski in its isotropy formulation; the
primary references are Kim, \emph{Hasse--Minkowski Theorem},
Theorem~1.1, and Hanke, \emph{Quadratic Forms and Automorphic Forms},
Theorem~1.9.1 and Remark~1.9.2 (printed p.~20).\footnote{Author-written
texts: \href{https://www.csun.edu/~sungjin/HasseMinkowski.pdf}{Kim's
statement} and \href{https://swc-math.github.io/aws/2009/09HankeNotes.pdf}{Hanke's
expanded Arizona Winter School notes}. The report distinguishes
their stated general theorem from proofs specialized to \(\mathbb Q\).}
Its global proof is not included here.

Let \(a,b,n\in F\) be totally positive. At a finite place, the binary
Hermitian form \(aN(s)+bN(t)\) represents every element of
\(F_{\mathfrak p}\).
For split \(E_{\mathfrak p}\), one norm already represents every
element. In the unramified field case, remove norm factors from
\(a,b\), leaving exponents \(0\) or \(1\).
If their parities differ, one coefficient has the target's valuation
parity, and unit-norm surjectivity gives a representation with that
coordinate. If their parities agree, the two-dimensional norm sum
is isotropic because \(-1\) is a unit norm, and represents every
target: after absorbing that norm, use
\(N((c+1)/2)-N((c-1)/2)=c\).
This is a fraction-field identity, with no integral division by \(2\).

Choose \(w\) as in (MF.1), with \(\delta=w^2\in F\).
At every real embedding \(\delta<0\).
The five-variable form
\[
 a(X_1^2-\delta X_2^2)+b(X_3^2-\delta X_4^2)-nX_5^2   \tag{MF.6}
\]
is locally isotropic at finite places by the explicit representations,
and at real places because its signature is \((4,1)\).
Hasse--Minkowski gives a nonzero global zero. Its last coordinate
cannot vanish: positivity at a real embedding would then force the
other four coordinates to vanish. Set
\(s=(X_1+wX_2)/X_5,\ t=(X_3+wX_4)/X_5\).
Then \(a\,s\bar s+b\,t\bar t=n\).
This proves precisely \texttt{HermRepUnram}.
Gram--Schmidt now shows that every totally positive definite
Hermitian space of rank at least two represents every totally
positive element of \(F\). Only this global step and its descendants
use the named theorem.

\paragraph{MF5: the matrix-density measure and the local ratio law.}
For a local integral lattice \(L\) of rank \(m\), let \(H\) be its
Gram matrix and \(d=\operatorname{ord}_{\mathfrak p}\det H\).
The source's \(d=\texttt{vd}\) is exactly this integer: determinants
of all Gram matrices of \(m\)-tuples generate \((\det H)\).
Write
\[
 \alpha_L=\lim_{k\to\infty}
 \frac{\#\{X\bmod\pi^k:X^*HX\equiv H\bmod\pi^k\}}{q^{km^2}},
 \qquad \nu(L)=q^{-md}\alpha_L.                       \tag{MF.7}
\]
These are the literal \texttt{akCount}/\texttt{alphaL}
normalizations; the numerator initially counts endomorphisms.
For \(k>d\), determinants show that all such endomorphisms are
invertible. In the split algebra, their two nonnegative determinant
valuations have sum zero.

On the full matrix space give the integral matrix lattice volume
one; do the same on the \(F_{\mathfrak p}\)-space of Hermitian
matrices. The map \(X\mapsto X^*HX-H\) has surjective derivative at
every unitary \(X\). At \(1\) it is \(Z\mapsto Z^*H+HZ\).
To solve it, write a Hermitian target as \(Y+Y^*\):
use upper-triangular entries and an integral trace-one element on
the diagonal, and then multiply by \(H^{-1}\).
This exhibits bounded denominators, including at \(2\).
On a sufficiently small ball a linear coordinate change puts the
map into projection form. Successive corrections prove this:
after removing the linear error the quadratic error has strictly
larger valuation, and the corrections converge.
The quotient of ambient measure by target measure is then the
measure on the fiber. Finitely many such balls cover the compact
unitary stabilizer; compactness puts all sufficiently accurate
approximate solutions in their union. Hence (MF.7) stabilizes,
equals its fiber measure, and is positive.

If the basis changes by \(C\), conjugation on the full matrix space
has Jacobian one, whereas \(Y\mapsto C^*YC\) on Hermitian matrices
has Jacobian \(|N_{E/F}\det C|^m\).
For diagonal \(C\), the diagonal and off-diagonal coordinates give
this directly; elementary matrices have determinant one, so
elimination proves the general case.
The determinant valuation changes by
\(\operatorname{ord}_{\mathfrak p}N_{E/F}\det C\).
Therefore \(q^{-md}\) times the fiber measure is basis independent.
Unitary left and right translations preserve it since
\(N_{E/F}\det g=1\). We have constructed a Haar measure satisfying
\[
                \operatorname{vol}(K_L)=\nu(L).       \tag{MF.8}
\]
For two lattices, take the volume of the common finite-index
stabilizer to obtain
\[
 \nu(A)\,\operatorname{orb}(B,A)
          =\nu(B)\,\operatorname{orb}(A,B),             \tag{MF.9}
\]
where \(\operatorname{orb}(A,B)=[K_A:K_A\cap K_B]\).
This is exactly \texttt{LocalDensityRatio} after multiplication by
\(q^{m(d_A+d_B)}\), with the source's orbit orientation.

\paragraph{MF6: the unimodular factors and the full good-prime norm density.}
For the standard form over the residue algebra, the number of vectors
of norm \(a\ne0\) is
\[
                 q^{2m-1}-\eta_{\mathfrak p}^{\,m}q^{m-1}.
                                                               \tag{MF.10}
\]
In the split case choose a nonzero first component and a second
component with prescribed scalar product.
In the nonsplit case the one-coordinate norm has one preimage of
zero and \(q+1\) preimages of every nonzero value.
Convolution gives by induction the zero count
\(q^{2m-1}+(-1)^m(q-1)q^{m-1}\), and (MF.10).
Every nonzero residue vector is a smooth point of the norm map,
since its derivative is a nonzero trace pairing. Each lifting
step contributes \(q^{2m-1}\) choices.

Counting orthonormal columns gives
\[
 |U_m(\mathbb F_q)|=q^{m^2}\prod_{j=1}^m(1-(-1)^jq^{-j})
\]
in the nonsplit case. In the split case the group is
\(\mathrm{GL}_m(\mathbb F_q)\), with order
\(q^{m^2}\prod_{j=1}^m(1-q^{-j})\).
The correction \(Z+Z^*=Y\) lifts every residue isometry, with
\(q^{m^2}\) choices at each step. Thus, for a unimodular lattice,
\[
                         \alpha_L=A_m(\mathfrak p).   \tag{MF.11}
\]

For a nonzero local integer \(n\), put
\(e=\operatorname{ord}_{\mathfrak p}n\).
The normalized norm-vector density \(\alpha(n)\) is \(e_m\) when
\(e=0\). The primitive zero-residue vectors have density
\[
 1+\eta_{\mathfrak p}^{\,m}(q-1)q^{-m}-q^{1-2m}
        =e_m(1+\eta_{\mathfrak p}^{\,m}q^{1-m}).
\]
Divisible vectors contribute \(q^{2-2m}\alpha(n/\pi^2)\) if \(e\ge2\),
and zero if \(e=1\). Consequently
\[
 \frac{\alpha(n)}{e_m}
       =\sum_{j=0}^{e}
              (\eta_{\mathfrak p}^{\,m}q^{1-m})^j.     \tag{MF.12}
\]
This includes dyadic unramified places. For \(m\ge2\) it is
nonnegative and uniformly bounded, since the geometric ratio has
absolute value at most \(1/2\).
This is the actual \texttt{alphaN\_\allowbreak closed} supplier, stronger than
the later difference estimate used by the source.

\paragraph{MF7: arbitrary integral norm densities are bounded and have mean one.}
For any integral nondegenerate \(H\), put
\[
 \alpha_{L,k}(n)=q^{-k(2m-1)}
       \#\{x\in L/\pi^kL:h(x,x)\equiv n\bmod\pi^k\}.   \tag{MF.13}
\]
For \(m\ge1,\ n\ne0\), this stabilizes at sufficiently large \(k\).
For \(m\ge2\) it is bounded uniformly in \(n,k\), and its limit is
bounded, nonnegative, locally constant away from zero, and satisfies
\[
                         \int_R\alpha_L(n)\,dn=1,     \tag{MF.14}
\]
where \(R\) has volume one.

For a primitive vector \(x\), the norm derivative is
\(z\mapsto\operatorname{Tr}h(x,z)\).
If every coefficient were divisible by \(\pi^{d+1}\), perfectness of
the trace pairing and the adjugate identity would make \(x\)
divisible by \(\pi\). Its image is therefore \(\pi^tR\) with
\(0\le t\le d\). Partition primitive vectors into balls modulo
\(\pi^{d+1}L\). On each ball choose a coordinate with derivative
valuation \(t\). Successive correction in that coordinate gives a
uniform map onto a coset modulo \(\pi^{d+1+t}\); the other
\(2m-1\) coordinates are free. There are finitely many balls,
so their pushforward has a bounded locally constant density \(P(n)\).

Writing a nonzero vector as \(\pi^r x\) with \(x\) primitive gives
\[
 \alpha_L(n)=\sum_{\substack{r\ge0\\n\in\pi^{2r}R}}
       q^{-(2m-2)r}P(n/\pi^{2r}).                    \tag{MF.15}
\]
For nonzero \(n\) this is finite. For \(m\ge2\) it converges uniformly
even at zero. The same bound applies to finite-level densities:
the last ball with \(2r\ge k\) contributes at most
\(q^{k-2m\lceil k/2\rceil}\le1\).
These observations prove convergence, boundedness and local
constancy away from zero. The source's sufficient stabilization
bound \(k\ge2d+2\operatorname{ord}_{\mathfrak p}(n)+2\) follows
from these same balls and their scaled versions.
Every vector modulo \(\pi^k\) has one norm residue, whence
\(q^{-k}\sum_{n\bmod\pi^k}\alpha_{L,k}(n)=1\).
Dominated convergence, equivalently the geometric tail in (MF.15),
proves (MF.14). No interchange over a singular fiber is assumed.

\paragraph{MF8: the finite local fibers and their weighted orbit identity.}
Fix a nonzero vector \(x_0\) of norm \(n\in R\setminus\{0\}\), and
write \(V=Ex_0\perp W\), \(r=\dim_EW=m-1\).
For a lattice \(T\) in \(W\), let \(f_L(T)\) count the lattices
\(\Lambda\) locally isometric to \(L\), containing \(x_0\), with
\(\Lambda\cap W=T\). These are exactly the source \texttt{fibLoc}
sets after MF3. Each fiber is finite, because
\[
 Bx_0\oplus T\ \subseteq\ \Lambda\
                  \subseteq\ n^{-1}Bx_0\oplus T^\#.   \tag{MF.16}
\]
The first inclusion is immediate. For the second, projection onto
\(x_0\) has coefficient \(h(x_0,z)/n\in n^{-1}B\), and the
orthogonal component pairs integrally with \(T\).
Both endpoints in (MF.16) are full lattices, with finite quotient.

There are only finitely many local isometry classes of possible
\(T\). Let \(e=\operatorname{ord}_{\mathfrak p}n\) and
\(\ell=\log_q[\Lambda:Bx_0\oplus T]\), with the quotient read
in the containing direction. Projection identifies this quotient
with a subquotient of \(n^{-1}B/B\), so \(0\le\ell\le2e\).
Determinants give \(e+d_T=d_L+\ell\), whence \(d_T\le d_L+e\).
The finite list of Jordan exponents in MF2 therefore supplies the
required finite \texttt{IsLocReps}. We retain actual lattices in
\(W\), not abstract rational isometry classes.

Let \(G=U(V)\), \(H=U(W)\), with the Haar measures of MF5.
The measure on \(V\) which in an \(L\)-basis is
\(|\det H_L|\,dx\) is independent of that basis, and gives \(L\)
volume \(q^{-d_L}\).
Divide this measure by \(dn\) under the norm map to obtain an
invariant measure \(\sigma_n\) on the smooth fiber \(h(x,x)=n\).
Its value on the vectors in \(L\) is
\[
                      \sigma_n(L)=q^{-d_L}\alpha_L(n).
\]
By MF3 the fiber is a single homogeneous space \(H\backslash G\).
The quotient of the two Haar measures is therefore
\(\kappa\,\sigma_n\) for a positive scalar \(\kappa\) independent
of \(L\). This uniqueness can be proved by integrating a compactly
supported function along \(H\): it identifies invariant integration
on the quotient with Haar integration on \(G\).

A right \(K_L\)-orbit in \(H\backslash G\) has measure
\(\nu(L)/\operatorname{vol}(H\cap gK_Lg^{-1})\).
Regroup these orbits by their cut lattice \(T\). The group \(K_T\)
acts on the finite fiber in (MF.16), so finite orbit--stabilizer
gives
\[
 \sum_{[\Lambda]\text{ over }[T]}
 \frac{1}{\operatorname{vol}(\operatorname{Stab}_H\Lambda)}
                  =\frac{f_L(T)}{\nu(T)}.
\]
Consequently, with the literal definitions
\[
 \lambda_{\mathfrak p}(T)=\frac{A_rq^{r d_T}}{\alpha_T}
                         =\frac{A_r}{\nu(T)},\qquad
 \Psi_{\mathfrak p}(L,n)=\sum_{[T]}f_L(T)\lambda_{\mathfrak p}(T),
\]
we obtain
\[
 \Psi_{\mathfrak p}(L,n)
   =A_r\kappa\,\frac{q^{(m-1)d_L}\alpha_L(n)}{\alpha_L}.
                                                               \tag{MF.17}
\]
This proves \texttt{LIRatio} for any two integral lattices.
It also covers empty representation fibers: both sides are zero.
No cancellation by a possibly zero representation density was used.

\paragraph{MF9: pin the scalar, scale the form, and obtain literal LIStat.}
First suppose \(n=1\). Choose \(L_0=Bx_0\perp T_0\) with
\(d_{L_0}=d_{T_0}\le1\), as in MF2.
Every integral lattice of this local type containing \(x_0\)
splits orthogonally as \(Bx_0\perp T\): the projection coefficient
\(h(x_0,z)\) is integral. Its cut has the same rank and determinant
valuation as \(T_0\), so MF2 makes it integrally isometric to \(T_0\).
For a fixed cut, the lattice is uniquely determined. Hence
\[
                       \Psi_{\mathfrak p}(L_0,1)
                              =\lambda_{\mathfrak p}(T_0).       \tag{MF.18}
\]

The finite isometry group modulo \(\pi^k\), for \(k>d_{L_0}\),
acts transitively on all vectors of norm \(1\) modulo \(\pi^k\).
Such a vector lifts to norm exactly \(1\): its norm derivative is
surjective since its norm is a unit, so correction preserves its
residue modulo \(\pi^k\). Its integral orthogonal complement is
again of the unique minimal type just described.
The stabilizer of the first vector consists exactly of the
isometries of the cut modulo \(\pi^k\); the first row and column
are forced by its norm being one. Thus
\[
 a_{L_0}(k)=n_{L_0}(k,1)a_{T_0}(k),\qquad
                    \alpha_{L_0}=\alpha_{L_0}(1)\alpha_{T_0}.
                                                               \tag{MF.19}
\]
The powers of \(q\) match because
\(m^2=(2m-1)+(m-1)^2\). Positivity of the matrix densities shows
\(\alpha_{L_0}(1)>0\). Equations (MF.17)--(MF.19) therefore give
\(\kappa=1\) at norm one.

For general \(n\), use \(h_0=n^{-1}h\), for which \(x_0\) has norm
one. Write \(e=\operatorname{ord}_{\mathfrak p}n\).
Apply the following scaling comparison on a reference lattice
integral for \(h_0\), constructed below. The originally supplied
\(h\)-integral \(L\) need not be integral for \(h_0\); the
lattice-independent scalar \(\kappa\) in MF8, equivalently
\texttt{LIRatio} with the nonzero reference density, then transfers
the computed constant to that original \(L\).
Multiplication of a rank-\(s\) Gram matrix by \(n\) changes the
determinant valuation by \(se\), and comparison of the congruences
at levels \(k\) and \(k-e\) gives
\[
 \alpha_{n h_0,L}=q^{s^2e}\alpha_{h_0,L},\qquad
 \alpha_{n h_0,L}(n)=q^e\alpha_{h_0,L}(1).             \tag{MF.20}
\]
For the first equality, the \(s^2\) Hermitian equations lose \(e\)
digits; for the second, the single norm equation loses \(e\)
digits. More explicitly the free lifts of matrices contribute
\(q^{2s^2e}\), and those of vectors contribute \(q^{2se}\);
division by their respective normalizing powers gives (MF.20).
The isometry groups, cuts and fibers do not change when the form
is scaled. Also \(\lambda_{\mathfrak p}\) is unchanged, since its
two scaling factors \(q^{s^2e}\) cancel.
The ratio in (MF.17) changes by \(q^{-(m-1)e}\).
Its pinned scalar is thus \(q^{(m-1)e}\).

The locally chosen reference lattice is available in the literal
global formulation too. Approximate a local basis of the
complement by rational vectors, preserving its integral Gram
elementary divisors, and multiply these vectors by an integer of
\(F\) which is a unit at \(\mathfrak p\) and clears all other
pairing denominators. Keep \(x_0\) itself unchanged.
This gives a globally integral \(h_0\)-lattice with the specified
localization, and hence a globally integral \(h=n h_0\)-lattice.
We have proved the exact formula
\[
 \boxed{\displaystyle
 \Psi_{\mathfrak p}(L,n)=
 q^{(m-1)(\operatorname{ord}_{\mathfrak p}n+d_L)}
       A_{m-1}(\mathfrak p)\frac{\alpha_L(n)}{\alpha_L}.}
                                                               \tag{MF.21}
\]
This is every clause of \texttt{LIStat}, with the original
\texttt{VecSplit}, \texttt{fibLoc}, and density definitions.
For \(m=1\) the cut is zero dimensional, its automorphism group,
density and empty product all equal one. The argument still applies.

\paragraph{MF10: the rank-one mass constant is independent of the lattice.}
On a positive Hermitian line write \(h(x,y)=a\bar x y\), \(a\in
F_{>0}\), and identify a full lattice with a fractional ideal \(I\).
The rational unitary group is always
\(E^1=\{z:z\bar z=1\}\), independently of \(a,I\).
At each finite place its local stabilizer is also independent of
\(a,I\): it is the norm-one unit subgroup.
A lattice \(J\) is in the genus of \(I\) exactly when the relative
fractional ideal \(JI^{-1}\) has exponent zero at inert primes and
opposite exponents at a split pair. MF3 identifies this completed
description with the source rational-local description.
Hence the set of genus classes is the same quotient of the group
of these balanced ideals by \(\{z\mathcal O_E:z\bar z=1\}\)
for every \(a,I\). Its automorphism group is always
\(\mathcal O_E^\times\cap E^1\), a finite group by the bounded-image
argument of CR4. Finiteness of the genus and positivity of its mass
are IG9. Thus all positive integral Hermitian lines have the same
positive mass \(c_1\).

For rank one with determinant valuation \(d\), the matrix congruence
is \(N(z)\equiv1\bmod\pi^{k-d}\) when \(k>d\).
Counting the \(2d\) free digits above a norm-one residue gives
\(\alpha_L=q^d(1-\eta_{\mathfrak p}q^{-1})\).
Therefore \(\lambda_{\mathfrak p}(L)=1\) at every place.
We have proved precisely
\[
                  \operatorname{MassConst}(E,F,1,c_1),\qquad c_1>0.
                                                               \tag{MF.22}
\]
Its numerical value is unnecessary: it will cancel in the adjacent
rank ratio.

\paragraph{MF11: the real lattice-point constant, without a mass theorem.}
For a rank-\(m\) integral Hermitian lattice \(L\), let
\(\mathfrak d(L)\) be its determinant ideal and put
\(D(L)=\sqrt{N_E\mathfrak d(L)}>0\).
Hermitian symmetry makes the two split determinant valuations
equal, and an inert prime has norm \(q^2\); hence
\[
                         D(L)=\prod_{\mathfrak p}q^{d_L}.        \tag{MF.23}
\]
Only finitely many factors differ from one.

Choose an orthogonal \(E\)-basis \(b_i\) inside \(L\), with
\(A=\bigoplus_i\mathcal O_E b_i\subset L\) of finite index \(I\).
In its complex coordinates, the region
\(h_\tau(z,z)\le T\) for every \(\tau:F\hookrightarrow\mathbb R\)
is a product of ellipsoids. Its volume is
\[
 T^{fm}\left(\frac{\pi^m}{m!}\right)^f
                 \Big/\prod_\tau\det H_\tau.
\]
The covolume of \(A\) in these coordinates is \(v^{fm}\), and
that of \(L\) is \(v^{fm}/I\).
Local determinants, or the determinant of a finite-index basis
change, give
\[
                  D(L)I=\prod_\tau\det H_\tau.        \tag{MF.24}
\]
This also proves that \(D(L)\) is constant on a genus.
The number of lattice points in an expanding bounded ellipsoid,
divided by its scale to the ambient dimension, tends to volume
divided by covolume. To see this, tile by a fixed fundamental
parallelotope. Tiles intersecting the boundary lie in a bounded-width
boundary neighborhood; after rescaling, its volume tends to zero.
The same argument works for this finite product of ellipsoids.
Consequently
\[
 \lim_{T\to\infty}T^{-fm}
 \#\{x\in L:h_\tau(x,x)\le T\text{ for every }\tau\}
       =\frac{(\pi^m/m!)^f}{v^{fm}D(L)}.              \tag{MF.25}
\]
This is the source \texttt{tendsto\_\allowbreak ballCount}, including its powers
of \(2,\pi,v\). Along \(T=j^2\) it has exactly the source's exponent
\(j^{2fm}\).

\paragraph{MF12: the global weighted representation count is a cut-lattice count.}
For \(n\in\mathcal O_F\) totally positive, define
\[
 R_L(n)=\sum_{[M]\in\operatorname{genus}(L)}
       \frac{\#\{x\in M:h(x,x)=n\}}{|\operatorname{Aut}(M)|}.
                                                               \tag{MF.26}
\]
The sums and all vector sets are finite by IG9 and positivity.
Suppose \(m\ge2\). MF4 supplies a rational \(x_0\) of norm \(n\);
let \(W=x_0^\perp\).
Apply finite orbit--stabilizer to the action of
\(\operatorname{Aut}(M)\) on its norm-\(n\) vectors, then apply
MF3 to move each rational vector to \(x_0\).
One obtains the sum of reciprocal stabilizer orders over lattices
\(\Lambda\) in the genus containing \(x_0\), modulo \(U(W)(F)\).
The cut \(T=\Lambda\cap W\) is a full integral lattice.
For a fixed cut \(T\), its possible extensions \(\Lambda\) lie
between the two global lattices in (MF.16), so form a finite set
\(\mathcal F_T\). The finite group \(\operatorname{Aut}(T)\) acts on
it, and its stabilizer is exactly
\(\operatorname{Aut}(\Lambda,x_0)\): the extension of an isometry
of \(T\) fixes \(x_0\), and conversely.
Thus the second orbit--stabilizer calculation gives
\[
                       R_L(n)=\sum_{[T]}
                    \frac{|\mathcal F_T|}{|\operatorname{Aut}(T)|}.
                                                               \tag{MF.27}
\]
In particular this identity has the reciprocal automorphism weights;
an unweighted count of complements would be wrong.

\paragraph{MF13: patch the actual fibers and regroup by complement genera.}
For a fixed rational cut lattice \(T\),
\[
                    |\mathcal F_T|
                          =\prod_{\mathfrak p}f_L(T_{\mathfrak p}).
                                                               \tag{MF.28}
\]
A global lattice determines its localizations. Conversely the finite
local choices in (MF.16) patch by intersection to one global
\(\mathcal O_E\)-lattice with exactly those localizations.
The finite lower and upper bounds in (MF.16) allow the ordinary
Chinese-remainder proof of this assertion, as in LP1--LP3.
Containing \(x_0\), integrality, and the exact cut are all local
conditions, so are preserved by this patching. Rational local
isometry is recovered by MF3.
Outside a finite set containing the primes dividing \(n\), the
determinant of \(L\), and that of one reference cut, both spaces
are unimodular and \(n\) is a unit. Projection then splits the
lattice, and the fiber is a singleton. Thus (MF.28) is a genuinely
finite product. If a local fiber is empty, both sides are zero.

Assume, as an induction hypothesis, that in rank \(m-1\)
\[
                  \operatorname{mass}(T)
                    =c_{m-1}\prod_{\mathfrak p}\lambda_{\mathfrak p}(T).
                                                               \tag{MF.29}
\]
Group (MF.27) by the genus of \(T\). The local fiber cardinalities
and local \(\lambda\)'s are constant on each such genus.
Every finite tuple of allowed local cut types occurs globally:
choose actual representatives in the fixed \(W_{\mathfrak p}\)
and patch them to a lattice in the same rational space \(W\).
This assertion patches lattices; it does not claim the existence
of a new rational Hermitian space with arbitrary invariants.
The finite sum over choices therefore distributes into a product:
\[
                       R_L(n)
                   =c_{m-1}\prod_{\mathfrak p}\Psi_{\mathfrak p}(L,n).
                                                               \tag{MF.30}
\]
The case where a local choice is impossible is again the zero
identity. This proves \texttt{repMass\_\allowbreak eq\_\allowbreak prod\_\allowbreak psi} from its actual
global fiber sets, not from a presumed representation formula.

\paragraph{MF14: the factor identity and the finite singular series.}
Define the literal local and global mass factors and singular series:
\[
 \lambda_{\mathfrak p}(L)=
       \frac{A_m(\mathfrak p)q^{m d_L}}{\alpha_L},\qquad
 \lambda(L)=\prod_{\mathfrak p}\lambda_{\mathfrak p}(L),\qquad
 \mathfrak S_L(n)=\prod_{\mathfrak p}
                        \frac{\alpha_L(n)}{e_m(\mathfrak p)}.
                                                               \tag{MF.31}
\]
The first product is finite because (MF.11) gives
\(\lambda_{\mathfrak p}=1\) when \(d_L=0\).
For nonzero \(n\), the second is also finite: outside the primes
dividing \(n\mathfrak d(L)\), (MF.12) gives \(\alpha_L(n)=e_m\).

Using \(A_m=A_{m-1}e_m\), the local identity (MF.21) becomes
\[
 \Psi_{\mathfrak p}(L,n)=
 q^{(m-1)\operatorname{ord}_{\mathfrak p}n}
       \lambda_{\mathfrak p}(L)\,q^{-d_L}
                   \frac{\alpha_L(n)}{e_m(\mathfrak p)}.
                                                               \tag{MF.32}
\]
Multiplication, with (MF.23) and unique ideal factorization, gives
\[
 R_L(n)=c_{m-1}\lambda(L)D(L)^{-1}
                  |N_{F/\mathbb Q}n|^{m-1}\mathfrak S_L(n).
                                                               \tag{MF.33}
\]
This is the source \texttt{repMass\_\allowbreak eq}. Every product in this
identity has finite support; no analytic Euler-product interchange
has yet been used.

\paragraph{MF15: weighted ideal boxes and finite-place equidistribution.}
Let
\[
 \mathcal B_T=\{n\in\mathcal O_F:0<\tau(n)\le T
                                  \text{ for every }\tau:F\hookrightarrow\mathbb R\}.
\]
For a nonzero ideal \(\mathfrak a\), a residue class modulo a fixed
ideal coprime to \(\mathfrak a\) is a translate of the product ideal
inside the real Minkowski lattice. Tiling by a fundamental
parallelotope and using Riemann sums gives the weighted asymptotic
\[
 T^{-fm}\sum_{\substack{n\in\mathcal B_T\cap\mathfrak a\\
                         n\equiv b\bmod\mathfrak c}}
       |Nn|^{m-1}
       \longrightarrow
       \frac{m^{-f}}{\sqrt{\operatorname{disc}F}\,
                                  N\mathfrak a\,N\mathfrak c}.
                                                               \tag{MF.34}
\]
Indeed the weight rescales to
\(\prod_\tau x_\tau^{m-1}\) on \((0,1]^f\), whose integral is
\(m^{-f}\); the lattice covolume is
\(\sqrt{\operatorname{disc}F}\,N\mathfrak a\,N\mathfrak c\).
The boundary has measure zero, so no uniform error for varying
ideals is needed for this fixed-ideal limit.

We do need a uniform domination. Distinct points of \(\mathfrak a\)
cannot lie in the same half-open cube of side
\((N\mathfrak a)^{1/f}\): their nonzero difference would have norm
strictly smaller than \(N\mathfrak a\), although its principal ideal
is contained in \(\mathfrak a\).
If \(\mathcal B_T\cap\mathfrak a\ne\varnothing\), then
\(N\mathfrak a\le T^f\). Subdividing the box therefore proves
\[
 \#(\mathcal B_T\cap\mathfrak a)
      \le\frac{2^fT^f}{N\mathfrak a},\qquad
 T^{-fm}\sum_{n\in\mathcal B_T\cap\mathfrak a}|Nn|^{m-1}
      \le\frac{2^f}{N\mathfrak a}.                    \tag{MF.35}
\]

Fix a finite set \(S\) of places and put
\(g(n)=\prod_{\mathfrak p\in S}\alpha_{L,\mathfrak p}(n)/
e_m(\mathfrak p)\). For \(m\ge2\), MF7 makes \(g\) bounded.
Away from \(n\in\mathfrak p^j\), each factor is constant on
sufficiently fine residue classes. Replacing it there by the
corresponding finite step function introduces an error bounded
by a constant times \(1_{n\in\mathfrak p^j}\).
For \((\mathfrak a,\mathfrak p)=1\), (MF.35) bounds its normalized
weighted sum over \(\mathfrak a\) by
\(C/(N\mathfrak a\,q^j)\).
Apply (MF.34), the Chinese remainder theorem, and then let \(j\)
increase. The mean-one identity (MF.14) gives
\[
 T^{-fm}\sum_{n\in\mathcal B_T\cap\mathfrak a}
        |Nn|^{m-1}g(n)
 \longrightarrow
 \frac{m^{-f}}{\sqrt{\operatorname{disc}F}\,N\mathfrak a}
                   \prod_{\mathfrak p\in S}e_m(\mathfrak p)^{-1}
                                                               \tag{MF.36}
\]
whenever \(\mathfrak a\) is prime to \(S\).
The same sums are bounded in absolute value by \(C/N\mathfrak a\),
uniformly in \(T\) and \(\mathfrak a\). This proves the actual
finite-place weighted averaging statement, including the treatment
of the exceptional zero residue.

\paragraph{MF16: the singular-series mean by an absolutely dominated divisor sum.}
Choose \(S\) containing the primes with \(d_L\ne0\).
For ideals prime to \(S\), define
\(\chi_m(\mathfrak a)=
 \prod_{\mathfrak p\notin S}\eta_{\mathfrak p}^{\,m\operatorname{ord}_{\mathfrak p}\mathfrak a}\).
The exact good-prime formula (MF.12) gives, for \(n\ne0\),
\[
 \mathfrak S_L(n)=g(n)
      \sum_{\substack{\mathfrak a\mid(n)\\(\mathfrak a,S)=1}}
                \frac{\chi_m(\mathfrak a)}{(N\mathfrak a)^{m-1}}.
                                                               \tag{MF.37}
\]
The normalized sum of a fixed divisor's contribution is bounded
by \(C/(N\mathfrak a)^m\), using (MF.35).
This is summable for \(m\ge2\): there are at most \(f\) prime ideals
over a rational prime, each of norm at least that prime, so
\(\sum_{\mathfrak p}(N\mathfrak p)^{-m}<\infty\).
Expanding the resulting absolutely convergent Euler product proves
\(\sum_{\mathfrak a}(N\mathfrak a)^{-m}<\infty\).
We may therefore pass the limit through the divisor sum.
Using (MF.36) and then its absolutely convergent Euler product gives
\[
 \begin{split}
 &\lim_{T\to\infty}T^{-fm}
       \sum_{n\in\mathcal B_T}|Nn|^{m-1}\mathfrak S_L(n)\\
 &\hspace{12mm}=
 \frac{m^{-f}}{\sqrt{\operatorname{disc}F}}\,
       \prod_{\mathfrak p\in S}e_m(\mathfrak p)^{-1}
       \sum_{(\mathfrak a,S)=1}\frac{\chi_m(\mathfrak a)}{(N\mathfrak a)^m}
   =\frac{m^{-f}}{\sqrt{(2v)^f}}\,L_m.                
 \end{split}
 \tag{MF.38}
\]
This proves \texttt{tendsto\_\allowbreak singS\_\allowbreak sum}, with the source's actual
constant. The infinite product was not interchanged with an
uncontrolled average: (MF.35) and the summable divisor majorant
justify the passage, including at the induction's endpoint \(m=2\).

\paragraph{MF17: compare two counts of the same vectors.}
Let \(s\) be one finite set of genus representatives for \(L\).
A nonzero vector in a totally positive definite space has a norm
in \(\mathcal B_T\) exactly when all its archimedean norms are
at most \(T\); the zero vector is the only vector of norm zero.
Hence the following identity holds before taking any limit:
\[
 \sum_{M\in s}\frac{\#\{x\in M:h_\tau(x,x)\le T\ \forall\tau\}}
                       {|\operatorname{Aut}(M)|}
              =\operatorname{mass}(L)+\sum_{n\in\mathcal B_T}R_L(n).
                                                               \tag{MF.39}
\]
The left side, divided by \(T^{fm}\), tends by MF11 to
\[
            \operatorname{mass}(L)
                        \frac{(\pi^m/m!)^f}{v^{fm}D(L)}.
\]
Under the rank-\((m-1)\) induction hypothesis, MF14 and MF16 make
the right side tend to
\[
            c_{m-1}\lambda(L)D(L)^{-1}
                             \frac{m^{-f}}{\sqrt{(2v)^f}}L_m.
\]
The separate zero-vector mass tends to zero because \(fm>0\).
All quantities cancelled below are positive: \(D(L),v,\pi,m!\)
are positive, as is \(c_{m-1}\).
Equality of the two limits and \(m!=m(m-1)!\) therefore prove
\[
 \operatorname{mass}(L)=c_{m-1}R_m\lambda(L),\qquad
 R_m=\frac{\big((m-1)!(v/\pi)^m\big)^f}{\sqrt{(2v)^f}}\,L_m .
                                                               \tag{MF.40}
\]
This is \texttt{massConst\_\allowbreak succ}.
The factor \(m^{-f}\) cancels the extra \(m^f\) in the ball
factorial. The discriminant contribution remains exactly
\(\sqrt{(2v)^f}\), with no unidentified measure constant.

\paragraph{MF18: the genus mass formula by induction.}
Start with the positive constant \(c_1\) in MF10 and define
\[
                   c_m=c_1\prod_{j=2}^m R_j
                   \quad(m\ge1).                    \tag{MF.41}
\]
The empty product is one. MF1 gives \(R_j>0\) for all \(j\ge2\),
so each \(c_m>0\).
Induct on \(m\), using MF17 for the successor step.
For every totally positive definite rank-\(m\) Hermitian space
over \(E\), and every full integral lattice \(L\) in it, we obtain
\[
                    \operatorname{mass}(L)=c_m\lambda(L).
                                                               \tag{MF.42}
\]
This is the source \texttt{massConst\_\allowbreak all}.
It is uniform over the rational Hermitian space and over
projective lattices. The local type at bad places is carried by
\(\lambda(L)\); it has not been silently replaced by the
unimodular type. The proof uses only the preceding local counts,
finite genus sums, lattice patching, and the global representation
theorem MF4. In particular it is not an invocation of the desired
Siegel mass formula under another name.

\paragraph{MF19: the literal standard and model mass-ratio endpoints.}
For an \texttt{IsOrthoLattice} of rank \(m\), the supplied global
orthonormal basis identifies the lattice with
\(\mathcal O_E^m\) and the form with the identity matrix.
Thus \(d_L=0\) at every place and (MF.11) gives
\(\lambda(L)=1\). Its mass is \(c_m\).
For two such lattices of ranks \(M\) and \(M-1\), positivity allows
division of (MF.41), giving
\[
 \boxed{\displaystyle
 \frac{\operatorname{mass}(L_M)}{\operatorname{mass}(L_{M-1})}
    =\frac{\big((M-1)!(v/\pi)^M\big)^f}{\sqrt{(2v)^f}}\,
          \operatorname{lValEta}(E,M).}              \tag{MF.43}
\]
We need only \(M\ge3\), exactly the source quantifier.
The argument internally also treats the step from rank one to
rank two. No formula for the absolute rank-one constant is required.

This proves the literal \texttt{SiegelMassRatio} for arbitrary
supplied orthonormal models. Taking the standard coordinate form
and lattice gives \texttt{StandardLattice.SiegelMassRatioCM},
hence the mathematical content of
\texttt{massRatio\_\allowbreak all\_\allowbreak unramified\_\allowbreak cmFields}.
Conversely every orthonormal model is carried to the standard one
by its displayed basis; transporting lattice classes and stabilizers
preserves mass. The model and standard interfaces therefore agree
with exactly the source normalization. All these global assertions
are consequences of the explicitly retained classical
Hasse--Minkowski theorem in MF4.

\paragraph{MF20: the precise consequence for the quantitative window.}
Insert (MF.43) into QW3. For \(k\ge2\) and \(r\ge k\),
\[
        \operatorname{mass}(L_r)=\operatorname{mass}(L_k)
                         \prod_{M=k+1}^{r}R_M.        \tag{MF.44}
\]
The common positive rank-one constant has already cancelled.
Thus QW9's complete quantitative \texttt{window\_\allowbreak lower} now uses
the proved mass ratio rather than a separate unexplained mass-law
premise. Its original scale order, \(0<\epsilon<1/52\), standard
constants, same-\texttt{PreCh} arithmetic data, and the exact
\(d_u,D_s\) hypotheses remain unchanged.
In particular QW6's constant
\(30+(2+2C_0)/\epsilon\) and QW9's choice \(C_0=5\) require no
alteration.

There are three distinct boundaries. First, Hasse--Minkowski is
a named established theorem used here, whose proof has not been
reconstructed. Second, the source's geometric
\texttt{VolumeEqMass2}, \texttt{VolTie}, indefinite volume
normalization, and compatible \texttt{Core}/arithmetic data are
not consequences merely of (MF.43). Third, this section does not
supply the unrelated short-interval prime estimate or finish any
of the RH/GRH global branches. The completed local density law and
global mass ratio must not be relabeled as any of those stronger
statements.





\section{An actual ray-class family and its uniform supply of norm primes}
\label{sec:ray-prime-estimates}

\status{QR1--QR9 construct the literal arithmetic \texttt{LFam} left open
in QW11. QR10--QR24 prove its analytic short-interval supplier, with a
polynomial real-zero gap and constants uniform over every such family.
QR25 combines these results with the already proved QW reductions.
These are proofs of specific arithmetic and analytic suppliers. The
simultaneous geometric choices in the four global branches remain
separate. No Lean compilation or kernel replay is asserted.}

We write $\mathcal B$ for the conductor bound, to distinguish it from
the scale $B=e^{x^\epsilon}$ in QW1--QW9. Groups of ray classes are
written multiplicatively in the arithmetic construction and additively
in \texttt{LFam}; their identity is therefore denoted $1$ or $0$,
respectively. All implied constants below are absolute unless their
dependence on the fixed zero-free width $c_0$ is specified. Existence
of such constants is sufficient for the literal source statements.
The estimates do not use RH or GRH.

\paragraph{QR1: remove the unit ambiguity in the actual ray coset.}
Let $d\ge100$ be prime with $d\equiv3\pmod4$, let $M\ge3$, and set
$K=\mathbb Q(\sqrt{-d})$, $\omega=(1+\sqrt{-d})/2$.
Then $\mathcal O_K=\mathbb Z[\omega]$ and $\operatorname{disc}K=-d$.
Use the actual ray group modulo $M\mathcal O_K$. The construction in
IG1 chooses, for each ray class $g$, an integral ideal $\mathfrak a$
prime to $M$ in the inverse class and $\alpha_0\in\mathfrak a$ with
$\alpha_0\equiv1\pmod {M\mathcal O_K}$. It identifies the ideals in
class $g$ with elements of the coset
$\alpha_0+M\mathfrak a$, modulo units congruent to one modulo $M$,
by $\alpha\mapsto(\alpha)\mathfrak a^{-1}$. The congruence conditions
on $\alpha$ ensure that this quotient is an integral ideal prime to
$M$. The inverse follows from the definition of the ray relation.

In this case the unit quotient is trivial. If a unit is
$u=1+Mt\ne1$, then $|u|=1$ under the complex embedding, whereas the
nonzero algebraic integer $t$ has $|t|^2=N_{K/\mathbb Q}(t)\ge1$.
Thus $M\le M|t|=|u-1|\le2$, a contradiction. Zero is not in the
coset, since every element is congruent to one modulo $M$.
Consequently the preceding map is a bijection, and
\[
 N\big((\alpha)\mathfrak a^{-1}\big)
       =\frac{|\alpha|^2}{N\mathfrak a},\qquad
 \operatorname{covol}(M\mathfrak a)
       =\frac{M^2N\mathfrak a\sqrt d}{2}.                 \tag{QR.1}
\]
These are the original ray class, norm and Euclidean measure; there is
no unit multiplicity to absorb in a later leading constant.

\paragraph{QR2: obtain a reduced binary form with the required bounds.}
Put $N=N\mathfrak a$ and $\Delta=\operatorname{covol}(M\mathfrak a)$.
A shortest nonzero lattice vector $v$ is primitive: otherwise dividing
by an integer of absolute value at least two gives a shorter vector.
It therefore extends to a $\mathbb Z$-basis $v,w$. Subtract an integer
multiple of $v$ from $w$ so its projection onto $v$ has absolute value
at most $|v|/2$. Since $|w|\ge|v|$, its perpendicular height is at
least $\sqrt3|v|/2$. Hence
\[
 |v|^2\le2\Delta/\sqrt3\le2\Delta.
\]
Also $v=Mz$ with $z\in\mathfrak a\setminus\{0\}$, so
$(z)\subseteq\mathfrak a$ and $|v|^2\ge M^2N$.
Writing the squared norm in this basis gives exactly
\[
 Q(m,n)=\frac{|\alpha_0+mv+nw|^2}{N}
       =a(m+\theta_1+\ell n)^2+b(n+\theta_2)^2,           \tag{QR.2}
\]
where
\[
 a=|v|^2/N,\quad b=\Delta^2/(|v|^2N),\quad
 \ell=\frac{\operatorname{Re}(\bar v w)}{|v|^2},\quad
 \theta_1=\frac{\operatorname{Re}(\bar v\alpha_0)}{|v|^2},\quad
 \theta_2=\frac{\operatorname{Im}(\bar v\alpha_0)}
                  {\operatorname{Im}(\bar v w)}.
\]
Orient the basis so that the last denominator is positive; its sign
would not affect (QR.2). The bounds above imply
\[
 a\ge1,\qquad b\ge\frac{\Delta}{2N}
                 =\frac{M^2\sqrt d}{4}\ge1,\qquad
 ab=\frac{dM^4}{4}.                                      \tag{QR.3}
\]
By QR1, every $Q(m,n)$ is a positive integer, and the number of its
representations of $j$ is exactly $a_g(j)$, the number of integral
ideals of norm $j$, prime to $M$, in the original class $g$.

\paragraph{QR3: the ellipse count is uniform in its shifts and shear.}
For $x\ge1$ and fixed $n$, the condition $Q(m,n)\le x$ is an interval
in $m$ of length
\[
 f(n+\theta_2),\qquad
 f(t)=\frac2{\sqrt a}\sqrt{(x-bt^2)_+}.
\]
Its integer count differs from its length by at most one. There are
at most $2\sqrt{x/b}+3$ rows meeting the ellipse. For any even,
nonnegative function decreasing on $[0,\infty)$ with compact support,
\[
 \left|\sum_{n\in\mathbb Z}f(n+\theta)-\int_{\mathbb R}f(t)\,dt\right|
 \le f(0).                                                \tag{QR.4}
\]
Indeed write $f$ as the integral of the indicators of its superlevel
sets. Each nonempty superlevel set is a centered interval, and its
shifted integer count differs from its length by at most one.
Integrating that discrepancy over the levels proves (QR.4); endpoints
of levels do not affect the integral. The semicircle integral is
$\pi x/\sqrt{ab}$. Thus (QR.3) gives
\[
 \left|\#\{(m,n):Q(m,n)\le x\}-\frac{\pi x}{\sqrt{ab}}\right|
 \le2\sqrt{x/b}+3+2\sqrt{x/a}\le10\sqrt x.                 \tag{QR.5}
\]
The row intervals also prove finiteness directly. Neither $\ell$ nor
the two shifts changes this estimate.

\paragraph{QR4: a uniform shifted Gaussian estimate.}
For $s>0$, periodize $e^{-s t^2}$ on $\mathbb R/\mathbb Z$.
Its Fourier coefficients are
$\sqrt{\pi/s}\,e^{-\pi^2j^2/s}$. The Gaussian and the resulting
Fourier series converge absolutely and uniformly on a period, so
integration against the characters and Fourier inversion give
\[
 \sum_{m\in\mathbb Z}e^{-s(m+\theta)^2}
 =\sqrt{\pi/s}\sum_{j\in\mathbb Z}
                     e^{-\pi^2j^2/s}e^{2\pi i j\theta}.   \tag{QR.6}
\]
The Gaussian Fourier transform follows by completing the square, or
by its differential equation and its value at zero. The periodization
argument is the same elementary Poisson argument used in IG1.
Since $e^{-y}\le2y^{-2}$ for $y>0$ and
$\sum_{j\ge1}j^{-4}\le2$,
\[
 2\sum_{j\ge1}e^{-\pi^2j^2/s}
 \le\frac{4s^2}{\pi^4}\sum_{j\ge1}j^{-4}\le s^2.
\]
Subtracting the term $j=0$ in (QR.6) therefore proves, for every shift,
\[
 \left|\sum_m e^{-s(m+\theta)^2}-\sqrt{\pi/s}\right|
                  \le\sqrt{\pi/s}\,s^2.                 \tag{QR.7}
\]
All sums involved are absolutely convergent, so the same estimate
can be used inside a second Gaussian sum.

\paragraph{QR5: verify both actual ray counting laws.}
Apply QR4 to the inner sum in $m$ in $e^{-tQ(m,n)}$; its shift
$\theta_1+\ell n$ depends on $n$, but the error bound does not.
Apply it again to the remaining sum in $n$. Bounding that outer
sum by its Gaussian main term times $1+(bt)^2$ gives
\[
 \left|\sum_{m,n}e^{-tQ(m,n)}-\frac{\pi}{\sqrt{ab}\,t}\right|
 \le\frac{\pi}{\sqrt{ab}\,t}
           \big((at)^2+(bt)^2+(at)^2(bt)^2\big).            \tag{QR.8}
\]
For $0<t\le1$, (QR.3) bounds this by $3\pi(ab)^2t$.
The equality between the double sum and
$\sum_{j\ge0}a_g(j)e^{-tj}$ follows by partition into the finite
norm fibers of QR2; nonnegativity also permits Tonelli's theorem.
Set
\[
 \mathcal B=4dM^2,\qquad
 \kappa=\frac{2\pi}{\sqrt d\,M^2}=\frac{\pi}{\sqrt{ab}}.
                                                               \tag{QR.9}
\]
The elementary inequalities $d\ge100$, $M\ge3$ and $\pi<4$
give $3\pi(ab)^2\le\mathcal B^{10}$. Combining (QR.5), (QR.8)
and the integral values of $Q$ proves exactly
\[
 \left|\sum_{1\le j\le\lfloor x\rfloor}a_g(j)-\kappa x\right|
 \le10\sqrt x,\qquad
 \left|\sum_{j\ge0}a_g(j)e^{-tj}-\kappa/t\right|
 \le\mathcal B^{10}t.                                    \tag{QR.10}
\]
The first estimate holds for every $x\ge1$; the second for every
$0<t\le1$, uniformly over all ray classes of this same field and
modulus.

\paragraph{QR6: ideal factorization gives the literal convolution.}
Define $\lambda(\mathfrak p^r)=\log N\mathfrak p$ for a nonzero
prime ideal $\mathfrak p$ and $r\ge1$, and define it to be zero on
other nonzero ideals. Unique ideal factorization gives
\[
 \sum_{\mathfrak j\mid\mathfrak i}\lambda(\mathfrak j)
       =\sum_{\mathfrak p}v_{\mathfrak p}(\mathfrak i)
                                \log N\mathfrak p
       =\log N\mathfrak i.                               \tag{QR.11}
\]
Let $\lambda_g(n)$ be the sum of these weights over the ideals counted
by $a_g(n)$. The norm fibers and the ideal divisor sets are finite.
For an ideal prime to $M$, both factors in a factorization
$\mathfrak i=\mathfrak j\mathfrak k$ are prime to $M$; conversely
the product of two such ideals is prime to $M$. Their norms multiply
and their ray classes add. Sum (QR.11) over the fiber of $(g,n)$ and
use this factorization bijection to obtain
\[
 a_g(n)\log n
  =\sum_{h\in G}\ \sum_{uv=n}\lambda_h(u)a_{g-h}(v).       \tag{QR.12}
\]
For $n=0$ the two sides are zero under the source convention, and
for $n=1$ every ideal weight is zero. No rearrangement of a conditionally
convergent series occurs here.

\paragraph{QR7: the degree-two weight bound retains norm-one support.}
All the weights are nonnegative. A nonzero weight at an integer norm
is supported on $n=p^r$. If $N\mathfrak p=p^f$, the contribution
of $\mathfrak p^k$ at that norm occurs only when $fk=r$ and equals
$f\log p$. For fixed $p$, the identity
$\sum_{\mathfrak p\mid p}e_{\mathfrak p}f_{\mathfrak p}=2$
implies $\sum f_{\mathfrak p}\le2$. Removing ideals not prime to
$M$ can only decrease the sum. Hence
\[
             \sum_g\lambda_g(n)\le2\Lambda(n).            \tag{QR.13}
\]
If $p$ is prime and $\lambda_0(p)\ne0$, there is an integral ideal
$\mathfrak i$ of norm $p$ in the identity ray class. The defining
ray relation has $(a)\mathfrak i=(b)$ with nonzero $a$ and
$a,b\equiv1\pmod M$. Since $(b)\subseteq(a)$, write $b=a\pi$
with $\pi\in\mathcal O_K$. Cancellation gives $\mathfrak i=(\pi)$,
and
$\pi-1=(a\pi-1)-\pi(a-1)\in M\mathcal O_K$.
Write $\pi=1+M(r+s\omega)$, $r,s\in\mathbb Z$, and put
$A=2(1+Mr)+Ms$, $D=Ms$. The exact quadratic norm identity gives
\[
       4p=A^2+dD^2,\qquad M\mid D,\qquad
                          2M\mid A-D-2.                  \tag{QR.14}
\]
These are precisely the original \texttt{NormOne} conditions, with
the same $d,M,p$ as the family. In particular no unsupported
replacement of a ray condition by ordinary principalness has been made.

\paragraph{QR8: real-character convolution is nonnegative and counts squares.}
Let $\chi^2=1$ be a ray character, so $\chi$ takes values $\pm1$.
For an ideal prime to $M$ put
\[
 r_\chi(\mathfrak i)=\sum_{\mathfrak j\mid\mathfrak i}
                  \chi([\mathfrak j])
 =\prod_{\mathfrak p^e\Vert\mathfrak i}
       (1+\chi([\mathfrak p])+\cdots+\chi([\mathfrak p])^e).
                                                               \tag{QR.15}
\]
If the sign is $1$, the factor is $e+1$; if it is $-1$, the factor
is $0$ for odd $e$ and $1$ for even $e$. All factors are nonnegative,
and every factor of $r_\chi(\mathfrak j^2)$ is at least one.
Grouping ideal factorizations by their product shows that the ordinary
Dirichlet convolution coefficient
\[
 c_\chi(n)=\sum_{uv=n}\Big(\sum_g a_g(u)\Big)
                              \Big(\sum_g\chi(g)a_g(v)\Big)
       =\sum_{N\mathfrak i=n}r_\chi(\mathfrak i)           \tag{QR.16}
\]
is nonnegative; on the right, all ideals are prime to $M$.
Squaring such ideals is injective by ideal factorization. Keeping just
the terms $\mathfrak i=\mathfrak j^2$ in $c_\chi(n^2)$ proves
$c_\chi(n^2)\ge\sum_g a_g(n)$. The zero norm case is also zero.
This proves both clauses of the source \texttt{sq} field, not merely
positivity of a few initial coefficients.

\paragraph{QR9: assemble every field of the actual arithmetic family.}
The ideal of norm one is $\mathcal O_K$, in the identity class;
there is no nonzero ideal of norm zero. Thus
$a_0(1)=1$, $a_g(1)=0$ for $g\ne0$, and $a_g(0)=0$.
The ray exact sequence bounds
$|G|\le h_K\,|\mathcal O_K/M\mathcal O_K|=h_KM^2$.
Here the short-vector argument of TL14 has the following explicit
two-dimensional specialization. An integral ideal $I$ has covolume
$\Delta_I=N(I)\sqrt d/2$. A disk of radius $\sqrt{\Delta_I}$ has
area $\pi\Delta_I>\Delta_I$. Averaging its lattice-point count over
a fundamental domain therefore gives a translate containing two
distinct points of $I$. Their nonzero difference $a$ satisfies
$|a|^2\le4\Delta_I=2\sqrt d\,N(I)$. Thus the integral ideal
$(a)I^{-1}$ has norm at most $2\sqrt d$; choosing $I$ in the inverse
class supplies such a representative in every class.
The number of degree-two ideals of
norm $n$ is at most the divisor function $\tau(n)$, by considering
split, inert and ramified prime powers. Therefore the number with
norm at most $2\sqrt d$ is at most
$\lfloor2\sqrt d\rfloor^2\le4d$. Hence $|G|\le\mathcal B$.
Formula (QR.9) gives $\mathcal B\ge4$ and
$1/\mathcal B\le\kappa\le4$ directly.

Take the actual finite additive ray group, $\mathcal B,\kappa$, and
the actual coefficients $a_g,\lambda_g$. Their nonnegativity and
norm-zero and norm-one conditions have just been proved. QR5 gives
both counting fields; QR6 gives convolution; QR7 gives the total
von Mangoldt bound; QR8 gives both real-character clauses. These
are all fourteen laws in the \texttt{LFam} record. QR7 simultaneously
supplies its norm-one support. This constructs
\texttt{exists\_\allowbreak lFam\_\allowbreak ray} and its public statement with the exact
conductor $4dM^2$; no arithmetic \texttt{LFam} supplier remains
as a hypothesis of QW11.

\paragraph{QR10: continue the whole family by its actual counting remainder.}
For the rest of this section let $\mathcal F$ be any \texttt{LFam},
with group $G$ of order $h$. Put
\[
 a_\chi(n)=\sum_g\chi(g)a_g(n),\quad
 \lambda_\chi(n)=\sum_g\chi(g)\lambda_g(n),\quad
 \rho_\chi=\begin{cases}\kappa h,&\chi=1,\\0,&\chi\ne1.\end{cases}
\]
Finite character orthogonality and the count field imply
\[
 R_\chi(x)=\sum_{n\le x}a_\chi(n)-\rho_\chi x,
       \qquad |R_\chi(x)|\le10\mathcal B\sqrt x
                      \quad(x\ge1).                       \tag{QR.17}
\]
Consequently
$m_\chi(s)=\int_1^\infty R_\chi(x)x^{-s-1}\,dx$
is holomorphic for $\operatorname{Re}s>1/2$, with
$|m_\chi(s)|\le10\mathcal B/(\operatorname{Re}s-1/2)$.
On every compact sub-half-plane its differentiated integrands are
dominated by $C x^{-1-\eta}(1+\log x)^j$; this proves both convergence
and holomorphy by differentiation under the integral.
Abel summation, justified first for $\operatorname{Re}s>1$, gives
\[
 L_\chi(s):=\sum_{n\ge1}a_\chi(n)n^{-s}
       =\frac{\rho_\chi s}{s-1}+s m_\chi(s).              \tag{QR.18}
\]
Thus the literal regularized function is
\[
 F_\chi(s)=\begin{cases}
 \rho_\chi s+(s-1)s m_\chi(s),&\chi=1,\\
 s m_\chi(s),&\chi\ne1.
 \end{cases}                                               \tag{QR.19}
\]
It is holomorphic for $\operatorname{Re}s>1/2$, equals
$(s-1)L_1(s)$ or $L_\chi(s)$ in the absolutely convergent region,
and satisfies
\[
 |F_\chi(s)|\le300\mathcal B(|s|+1)^2
                    \quad(\operatorname{Re}s\ge11/20).    \tag{QR.20}
\]
Indeed $|\rho_\chi|\le4\mathcal B$ and the preceding Mellin bound
is at most $200\mathcal B$ there. Real characters give
$F_\chi(\bar s)=\overline{F_\chi(s)}$ by conjugating the integral.

\paragraph{QR11: prove nonvanishing and the logarithmic derivative to the right.}
The nonnegative weights and their total bound give
$|\lambda_\chi(n)|\le2\Lambda(n)$ and $\lambda_\chi(1)=0$.
Take the finite character transform of (QR.12). Absolute convergence
for $\sigma>1$ permits multiplication of the resulting Dirichlet
series and gives
\[
 L'_\chi(s)=-G_\chi(s)L_\chi(s),\qquad
 G_\chi(s)=\sum_{n\ge2}\lambda_\chi(n)n^{-s}.             \tag{QR.21}
\]
Absolute convergence of $L_\chi$ in this region follows from its
partial-sum bound $\sum_{n\le x}|a_\chi(n)|\le C\mathcal Bx$;
termwise derivatives follow on smaller half-planes. Define
\[
 H_\chi(s)=\sum_{n\ge2}\frac{\lambda_\chi(n)}{\log n}\,n^{-s}.
\]
It has derivative $-G_\chi(s)$. Therefore
$L_\chi(s)e^{-H_\chi(s)}$ has derivative zero on the connected
half-plane. Its limit along the positive real axis is one, since
$a_\chi(1)=1$ and all remaining absolutely convergent terms tend to
zero. It is identically one. In particular
\[
 L_\chi=e^{H_\chi}\ne0,\qquad
 G_\chi(s)=\frac{\mathbf1_{\chi=1}}{s-1}
                        -\frac{F'_\chi(s)}{F_\chi(s)}.    \tag{QR.22}
\]
Unique integer factorization yields the absolutely convergent identity
$\sum_{n\ge2}\Lambda(n)n^{-\sigma}/\log n=\log\zeta(\sigma)$.
The integral test gives $\zeta(\sigma)\le\sigma/(\sigma-1)$.
Hence
\[
 |L_\chi(\sigma+it)|\ge e^{-2\log\zeta(\sigma)}
                       \ge((\sigma-1)/\sigma)^2.          \tag{QR.23}
\]
This proves all the required right-half-plane assertions without
presupposing an Euler product for the abstract family.

\paragraph{QR12: turn the theta field into an integrable nonprincipal heat sum.}
Put $\Theta_\chi(t)=\sum_n a_\chi(n)e^{-tn}$ for $t>0$.
From the count bound and nonnegativity,
$a_g(n)\le4n+10\sqrt n\le14n$ for $n\ge1$; consequently
$|a_\chi(n)|\le14\mathcal B n$. These bounds prove absolute
convergence of every heat sum and justify the finite character sum.
For $0<t\le1$, the theta field and orthogonality give
\[
       |\Theta_\chi(t)-\rho_\chi/t|\le\mathcal B^{11}t.
\]
For $t\ge1$, summing $n e^{-tn}$ bounds $|\Theta_\chi(t)|$
by $C\mathcal B e^{-t}$. After increasing an absolute $C\ge1$,
these statements imply
\[
 |\Theta_\chi(t)-\rho_\chi/t|
      \le C\mathcal B^{11}\min(t,t^{-1})\quad(t>0).       \tag{QR.24}
\]
If $\chi\ne1$, the small-$t$ bound and the stronger exponential
large-$t$ bound make $\Theta_\chi$ integrable. Its integral is
exactly $F_\chi(1)$. To see this without exchanging a possibly
nonabsolute series at $s=1$, use Abel summation in the form
\[
 \Theta_\chi(t)=t\int_1^\infty A_\chi(x)e^{-tx}\,dx,
      \qquad A_\chi(x)=\sum_{n\le x}a_\chi(n).
\]
For nonprincipal $\chi$, (QR.17) bounds $|A_\chi(x)|$ by
$10\mathcal B\sqrt x$ and
$\int_1^\infty|A_\chi(x)|x^{-2}\,dx<\infty$.
Fubini and $\int_0^\infty t e^{-tx}\,dt=x^{-2}$ now give
\[
                  \int_0^\infty\Theta_\chi(t)\,dt
                          =m_\chi(1)=F_\chi(1).           \tag{QR.25}
\]

\paragraph{QR13: a positive Bessel integral gives a polynomial lower bound at one.}
Let $\chi\ne1$ be real. Write $P=\Theta_1$, $Q=\Theta_\chi$,
$\rho=\kappa h$, $E=C\mathcal B^{11}$. The \texttt{sq} field
says the Dirichlet convolution coefficients $c_\chi(n)$ of
$a_1$ and $a_\chi$ are nonnegative. The norm-one coefficients give
$c_\chi(1)=1$. Define the positive kernel
\[
 W(y)=\int_0^\infty e^{-y/w-w}\,\frac{dw}{w}\quad(y>0),
       \qquad w_0=W(1)>0.
\]
The inequality $e^{-y/w}\le6w^3/y^3$ proves $W(y)\le12y^{-3}$.
Together with $|a_1(n)|,|a_\chi(n)|\le14\mathcal B n$, this
shows absolute integrability of the expanded double series below:
its absolute integrals are bounded by a constant times
$\mathcal B^2v^{-3}\sum_{m,n\ge1}m^{-2}n^{-2}$.
Substitute $w\mapsto w/m$ in each term to obtain
\[
 I(v):=\int_0^\infty P(v/w)Q(w)\,\frac{dw}{w}
       =\sum_n c_\chi(n)W(nv)\ge W(v)\ge w_0
                                      \quad(0<v\le1).     \tag{QR.26}
\]
Since $\min(t,t^{-1})\le\sqrt t$, (QR.24) gives
\[
 \left|I(v)-\frac\rho v\int_0^\infty Q(w)\,dw\right|
 \le E^2\sqrt v\,K_0,\qquad
 K_0=\int_0^\infty\frac{\min(w,w^{-1})}{w\sqrt w}\,dw
                    =2+\frac23=\frac83.                  \tag{QR.27}
\]
The integrable majorant in (QR.27) justifies subtracting the main
term. Choose once and for all
\[
 \delta=\min\left(1,\frac{w_0}{2(C^2K_0+1)}\right),\quad
 s=\delta\mathcal B^{-22},\quad v=s^2,\quad
 c_* =\min(1,\delta^2w_0/8)>0.
\]
The error in (QR.27) is at most $w_0/2$, so (QR.25)--(QR.27) and
$\rho\le4\mathcal B$ imply
$|F_\chi(1)|\ge w_0s^2/(8\mathcal B)\ge c_*\mathcal B^{-45}$.
For the principal character, (QR.19) gives
$F_1(1)=\rho\ge1/\mathcal B$, which satisfies the same lower bound.
Thus for every real character, including the principal one,
\[
                      |F_\chi(1)|\ge c_*\mathcal B^{-45}. \tag{QR.28}
\]
Every constant was chosen independently of the family. In particular
this step does not invoke an ineffective exceptional-zero theorem.

\paragraph{QR14: the real-zero gap follows from a local derivative bound.}
If $|z-1|\le1/4$, the disk $|w-z|\le1/10$ lies in
$\operatorname{Re}w\ge13/20$ and $|w|\le27/20$.
The growth bound (QR.20) is therefore at most $1800\mathcal B$
on that disk. Cauchy's derivative estimate gives
\[
                   |F'_\chi(z)|\le18000\mathcal B.        \tag{QR.29}
\]
For a real character and a real zero $\beta\ge3/4$, QR11 rules
out $\beta>1$ and QR13 rules out $\beta=1$. Integrating the derivative
along $[\beta,1]$ and using (QR.28) yields
\[
 c_*\mathcal B^{-45}\le|F_\chi(1)-F_\chi(\beta)|
           \le18000\mathcal B(1-\beta),\qquad
 1-\beta\ge c_9\mathcal B^{-46},\quad
 c_9=c_*/18000\in(0,1].                                  \tag{QR.30}
\]
This is the literal uniform real-zero separation needed below. It
is a polynomial separation for these families, with exponent $46$;
it is not a fixed-width exclusion of real zeros for unbounded
conductors.

\paragraph{QR15: a finite Blaschke product controls the local zero sum.}
Here is the complex-analysis estimate used in the zero-free argument.
Fix $0<r<R<1$. Suppose $f$ is holomorphic on a neighborhood of the
closed unit disk, $f(0)\ne0$, and $|f|\le M$ there. Its zeros
$a_j$ in $|z|\le R$, with multiplicities $m_j$, form a finite set:
an accumulation point would force the analytic function to vanish
identically. Divide by the finite Blaschke product
\[
 B_f(z)=\prod_j\left(\frac{z-a_j}{1-\bar a_jz}\right)^{m_j},
                    \qquad g=f/B_f.
\]
Removable singularities at the $a_j$ make $g$ holomorphic on the
unit disk and its boundary. On the unit circle $|B_f|=1$, so the
maximum principle gives $|g|\le M$. At zero,
$|g(0)|=|f(0)|/\prod_j|a_j|^{m_j}$. Thus, with
$L=\log(M/|f(0)|)\ge0$ and $N=\sum_jm_j$,
\[
                     N\log(1/R)\le L.                    \tag{QR.31}
\]
The function $g$ has no zeros on the disk of radius $R$, so an
analytic logarithm exists there. The real part of
$H=\log M-\log g$ is nonnegative, and
$\operatorname{Re}H(0)\le L$. The Poisson formula for this real
part, differentiated in the smaller disk, gives
\[
       |H'(z)|\le\frac{2R}{(R-r)^2}\operatorname{Re}H(0)
                                        \quad(|z|\le r).
\]
One may obtain the displayed bound by integrating the derivative of
$(Re^{i\theta}+z)/(Re^{i\theta}-z)$ against the nonnegative boundary
values; its absolute value is at most $2R/(R-r)^2$, and their mean
is $\operatorname{Re}H(0)$. A radius tending down to $R$ handles
any boundary issue. Finally
\[
 \frac{f'}f-\sum_j\frac{m_j}{z-a_j}
       =\frac{g'}g+\sum_j\frac{m_j\bar a_j}{1-\bar a_jz}.
\]
The last sum is at most $N/(1-r)$ in absolute value. With (QR.31)
this proves, for a constant depending only on $r,R$,
\[
 N\le C_{r,R}(1+L),\qquad
 \left|\frac{f'(z)}{f(z)}-\sum_j\frac{m_j}{z-a_j}\right|
       \le C_{r,R}(1+L)\quad(|z|\le r,\ f(z)\ne0).        \tag{QR.32}
\]
This is a direct proof of the local logarithmic-derivative lemma;
no unproved global zero-count estimate is substituted for it.

\paragraph{QR16: specialize that estimate with uniform family constants.}
Put $\ell(t)=\log(\mathcal B(|t|+2))>2$ and center a disk at
$z_t=5/4+it$. Formula (QR.23) gives
$|F_\chi(z_t)|\ge1/100$, including the principal factor $s-1$.
On the disk of radius $7/10$, (QR.20) gives the upper bound
$300\mathcal B(|t|+3)^2$. Rescale this disk to the unit disk
and apply QR15 with $r=5/7$, $R=6/7$. The logarithm of the
upper-to-lower ratio, plus one, is at most a fixed multiple of
$\ell(t)$. There is therefore one absolute $C>0$ such that
the zeros in $|\rho-z_t|\le3/5$, counted with multiplicity,
satisfy
\[
 \sum_\rho m_\rho\le C\ell(t),\qquad
 \left|\frac{F'_\chi(s)}{F_\chi(s)}
                   -\sum_\rho\frac{m_\rho}{s-\rho}\right|
          \le C\ell(t)\quad(|s-z_t|\le1/2,\ F_\chi(s)\ne0).
                                                               \tag{QR.33}
\]
All zeros have real part at most one by QR11. Consequently, at
$s=1+a+it$, $0<a\le1/2$, every term
$\operatorname{Re}(s-\rho)^{-1}$ is nonnegative. Combining
(QR.22) and (QR.33), and keeping any selected zeros in that disk,
gives
\[
 \operatorname{Re}G_\chi(1+a+it)
 \le\mathbf1_{\chi=1}\frac{a}{a^2+t^2}+C\ell(t)
         -\sum_{\rho\text{ selected}}m_\rho
                                      \operatorname{Re}\frac1{s-\rho}.
                                                               \tag{QR.34}
\]
The same $C$ can be increased once to cover all subsequent uses.
The elementary comparisons $\ell(0)\le\ell(t)$ and
$\ell(2t)\le2\ell(t)$ will be used below.

\paragraph{QR17: positivity gives a zero-free region outside real characters.}
For $|z|=1$,
$3+4\operatorname{Re}z+\operatorname{Re}z^2
 =2(1+\operatorname{Re}z)^2\ge0$.
Apply this with $z=\chi(g)n^{-it}$, multiply by
$\lambda_g(n)n^{-\sigma}\ge0$, and sum absolutely for $\sigma>1$.
This proves
\[
 3G_1(\sigma)+4\operatorname{Re}G_\chi(\sigma+it)
              +\operatorname{Re}G_{\chi^2}(\sigma+2it)\ge0.\tag{QR.35}
\]
Also $G_1(\sigma)\ge0$; if $\chi$ is real, then
$G_1(\sigma)+G_\chi(\sigma)\ge0$. For two real characters,
\[
 G_1(\sigma)+G_\chi(\sigma)+G_\psi(\sigma)
                                +G_{\chi\psi}(\sigma)\ge0,\tag{QR.36}
\]
because the corresponding coefficient is
$(1+\chi(g))(1+\psi(g))\ge0$.

We give the numerical separation step to fix the quantifiers.
Suppose $F_\chi(\rho)=0$, $\rho=1-\delta+i\gamma$,
$0\le\delta\le\varepsilon\le1/400$, and either $\chi^2\ne1$
or $\gamma\ne0$. If $\chi$ is real and
$0<|\gamma|<30\varepsilon$, its conjugate is a distinct zero.
Both are in the disk of (QR.33) centered at $5/4$.
At $\sigma=1+60\varepsilon$, their combined contribution is
at least
\[
 \frac{2(60\varepsilon)}{(61\varepsilon)^2+(30\varepsilon)^2}
                         =\frac{120}{4621\varepsilon}.
\]
Use $G_1\ge0$ if $\chi=1$, and $G_1+G_\chi\ge0$ otherwise.
Equation (QR.34), with the possible principal pole kept, yields
\[
 \frac{120}{4621\varepsilon}
          \le\frac1{60\varepsilon}+2C\ell(0),
       \qquad \varepsilon^{-1}\le300C\ell(\gamma).        \tag{QR.37}
\]
In all other cases use (QR.35) at $\sigma=1+6\varepsilon$,
$t=\gamma$, keeping the zero $\rho$ in the middle term.
Its contribution is at least $4/(7\varepsilon)$. If
$\chi^2\ne1$, the shifted pole terms are absent. Otherwise
$|\gamma|\ge30\varepsilon$, and each possible shifted pole
term is at most $1/(156\varepsilon)$ before its coefficient.
The principal term contributes $3/(6\varepsilon)$, and the
total errors are at most $9C\ell(\gamma)$. Hence
\[
 \frac4{7\varepsilon}
   \le\frac1{2\varepsilon}+\frac5{156\varepsilon}
                                  +9C\ell(\gamma),
       \qquad \varepsilon^{-1}\le300C\ell(\gamma).        \tag{QR.38}
\]
Both final inequalities are direct rational comparisons.
Put $K=300C$ and $c_{\rm zf}=1/(4K+400)>0$.
If a zero satisfies
$\operatorname{Re}\rho\ge1-c_{\rm zf}/\ell(\operatorname{Im}\rho)$,
choose $\varepsilon=1/((2K+200)\ell(\gamma))$.
Then $0\le\delta\le\varepsilon/2$, $\varepsilon\le1/400$,
and (QR.37) or (QR.38) is impossible. Every zero in this region
is therefore real and belongs to a real character. This proves
the stated family zero-free region, allowing the precise exception
that the next steps control.

\paragraph{QR18: at most one exceptional zero occurs in the entire family.}
Increase the constant $C$ of (QR.34) if needed, and choose
$c_{\rm un}=1/(100(C+1))$. Put $\ell_0=\ell(0)$ and
$\sigma=1+2c_{\rm un}/\ell_0$.
Any real zero $\beta\ge1-c_{\rm un}/\ell_0$ lies in the
disk used in (QR.34) at $t=0$ and contributes at least
\[
             \frac{m_\beta}{\sigma-\beta}
                    \ge\frac{m_\beta\ell_0}{3c_{\rm un}}.\tag{QR.39}
\]
Suppose that two such zero occurrences exist, counting multiplicity,
for real characters. If they belong to the same nonprincipal
character, use $G_1+G_\chi\ge0$; for the principal character
use $G_1\ge0$. If they belong to distinct nonprincipal real
characters $\chi,\psi$, use (QR.36). The four characters
$1,\chi,\psi,\chi\psi$ are then distinct. If one character
is principal, the corresponding two-term inequality suffices.
In every case exactly one principal pole is present, there are at
most four logarithmic-derivative errors, and (QR.34),(QR.39) force
\[
       0\le\frac{\ell_0}{2c_{\rm un}}+4C\ell_0
                                 -\frac{2\ell_0}{3c_{\rm un}}<0.
\]
The strict inequality follows from $1/(6c_{\rm un})>4C$.
This contradiction proves uniqueness across all characters and
simplicity of the zero if it exists. It also excludes two zeros
of one character; a count of exceptional characters alone would
not have supplied that assertion.

\paragraph{QR19: cancel the exception before estimating the derivative.}
Fix $0<c_0\le1/20$ and a finite set $Z$ satisfying the literal
$\operatorname{ExcZ}(\chi,c_0,Z)$: it contains at most one
real simple zero in $[1-c_0/\ell(0),1)$ and lists every zero
in $\operatorname{Re}s\ge1-c_0/\ell(\operatorname{Im}s)$.
For
\[
 1-\frac{c_0}{2\ell(t)}\le\operatorname{Re}s\le3,
             \quad t=\operatorname{Im}s,\quad s\notin Z,
\]
we claim
\[
 \left|\frac{F'_\chi(s)}{F_\chi(s)}
                         -\sum_{\beta\in Z}\frac1{s-\beta}\right|
                     \le C(c_0)\ell(t)^2.                 \tag{QR.40}
\]
For $\operatorname{Re}s\ge7/4$, use the absolutely convergent
series (QR.22); its absolute bound is uniform, and the real pole
and exception denominators have size at least $3/4$.
Otherwise $s$ lies in the radius-$1/2$ disk of (QR.33).
A zero $\rho$ in the radius-$3/5$ disk has
$\ell(\operatorname{Im}\rho)\le(4/3)\ell(t)$, since its
height differs from $t$ by at most $3/5$. If $\rho\notin Z$,
the defining zero-free condition for $Z$ gives
\[
 \operatorname{Re}(s-\rho)
    \ge\frac{c_0}{\ell(\operatorname{Im}\rho)}
                        -\frac{c_0}{2\ell(t)}
    \ge\frac{c_0}{4\ell(t)}.
\]
There are at most $C\ell(t)$ such zeros with multiplicity;
their reciprocal sum is at most $4C\ell(t)^2/c_0$.
An exceptional zero inside the disk cancels exactly, because it
is simple. If it is outside, its distance from $s$ is at least
$3/5-1/2=1/10$, so its added reciprocal has size at most ten.
Combine these facts with (QR.33) to prove (QR.40). The argument
also proves $F_\chi(s)\ne0$ at the points under consideration.

\paragraph{QR20: truncated Perron with the same integer endpoints.}
Let $b(n)$ be any complex coefficients with $|b(n)|\le2\Lambda(n)$.
For $2\le T\le u$, put $a=1+1/\log u$ and
$D_b(s)=\sum_{n\ge1}b(n)n^{-s}$. There is one absolute $P>0$
such that
\[
 \left|\sum_{1\le n\le\lfloor u\rfloor}b(n)
 -\frac1{2\pi}\int_{-T}^T D_b(a+it)\frac{u^{a+it}}{a+it}\,dt\right|
       \le P\left(\frac{u\log^2u}{T}+\log u\right).       \tag{QR.41}
\]
Here are the kernel and endpoint details. For $y>0$,
\[
 K_T(y)=\frac1{2\pi}\int_{-T}^T\frac{y^{a+it}}{a+it}\,dt.
\]
Closing the infinite vertical integral on the appropriate side gives
$\mathbf1_{y>1}$ when $y\ne1$. Integration by parts in each
tail bounds the error by $C y^a/(T|\log y|)$. A uniform second
bound is $C y^a$: after pairing $t$ and $-t$, the cosine integral
has weight $a/(a^2+t^2)$ of bounded integral, and the sine integral
with weight $t/(a^2+t^2)$ is uniformly bounded by splitting at $a$
and applying the Dirichlet bound to $\sin(t\log y)/t$.
The difference of the latter weight and $1/t$ has bounded integral
on $[a,\infty)$. Thus for $y\ne1$ the error is bounded by
\[
             C y^a\min(1,(T|\log y|)^{-1}).               \tag{QR.42}
\]
At $y=1$, $K_T(1)=\arctan(T/a)/\pi$, so the error from the
full endpoint weight one is at most one.

Absolute convergence on $\operatorname{Re}s=a>1$ permits the
finite-height integral and the coefficient series to be interchanged.
Apply (QR.42) with $y=u/n$. For $u/2<n<2u$,
$|\log(u/n)|\ge|u-n|/(2u)$ and $(u/n)^a$ is bounded, since
$a\le3$. The at most two closest integers cost $O(\log u)$;
the remaining terms cost at most
$C\log u\,(u/T)\sum_{1\le j\le2u}j^{-1}
 =O(u\log^2u/T)$. Outside this interval,
$|\log(u/n)|\ge\log2$ and
\[
 \sum_{n\ge2}\frac{\log n}{n^a}
                   \le C(1+(a-1)^{-2})\le C'\log^2u.
\]
Since $u^a=eu$, this gives the same error for the far terms.
If $u$ is an integer, changing its half endpoint to the required
full weight costs at most $|b(u)|\le2\log u$, already included
in (QR.41). These bounds prove the stated formula for the exact
floor convention, without a hidden distance-to-integers restriction.

\paragraph{QR21: shift the rectangle while keeping the exceptional residue.}
Assume $\operatorname{ExcZ}(\chi,c_0,Z)$ as in QR19, let $u\ge2$
and $T\ge1$, and put
\[
 \tau=\frac{c_0}{8\ell(T)},\qquad
 \mathcal M=C(c_0)\ell(T)^2+2/\tau.
\]
One of the two numbers $1-2\tau$, $1-4\tau$ has distance at
least $\tau$ from the at most one member of $Z$. Choose that
number as $\sigma$. The rectangle from $\sigma-iT$ to $a+iT$,
$a=1+1/\log u\le3$, is contained in the strip of QR19, because
$\ell(t)\le\ell(T)$ for $|t|\le T$. On its left and horizontal
edges the exceptional reciprocal and the possible principal reciprocal
are each at most $1/\tau$ in absolute value. Thus the meromorphic
continuation
\[
        G_\chi(s)=\frac{\mathbf1_{\chi=1}}{s-1}
                                  -\frac{F'_\chi(s)}{F_\chi(s)}
\]
has absolute value at most $\mathcal M$ there. Its only poles in
the rectangle are the principal pole at one, of residue $1$ when
present, and the listed zeros $\beta>\sigma$, of residue $-1$.
The simplicity in \texttt{ExcZ} is essential for this last coefficient.

Apply the residue theorem to $G_\chi(s)u^s/s$. The two horizontal
edges have $|s|\ge T$, bounded length and $u^{\operatorname{Re}s}
\le u^a=eu\le3u$; their combined contribution is at most
$18\mathcal M u/T$. On the left edge $|s|\ge\sigma\ge1/2$,
giving at most $4T\mathcal M u^{1-2\tau}$.
If an exceptional zero lies to the left of this edge, inserting its
term in the final formula costs at most $2u^{1-2\tau}$, since
$\beta\ge1/2$ and $\beta\le\sigma$.
Consequently the Perron integral $J$ in (QR.41), with
$b=\lambda_\chi$, satisfies exactly the convenient source bound
\[
 \left|J-\mathbf1_{\chi=1}u+\sum_{\beta\in Z}\frac{u^\beta}{\beta}\right|
 \le18\mathcal M u/T+4T\mathcal M u^{1-2\tau}
                                      +2u^{1-2\tau}.       \tag{QR.43}
\]
The residue $1$ of the principal logarithmic derivative explains
the main term $u$; the residue $\rho$ of $L_1$ itself does not
multiply that main term.

\paragraph{QR22: choose a height uniformly in the conductor.}
Suppose $\log\mathcal B\le\sqrt{\log u}$ and put
$W=\sqrt{\log u}$, $\eta=c_0/32$. For $\eta W>1$, choose
$T=e^{\eta W}$. Then $2\le T\le u$, $W\ge1$, and
$\ell(T)\le3W$. Hence
\[
 \mathcal M\le K W^4,\quad K=9C(c_0)+48/c_0,
       \qquad u^{1-2\tau}\le u/T^2.                      \tag{QR.44}
\]
The last inequality follows from
$2\tau\log u=c_0W^2/(4\ell(T))\ge c_0W/12
\ge2\eta W$.
Adding (QR.41) and (QR.43) therefore bounds the total error by
$(2P+22K+2)W^4u/T$. The elementary exponential-series inequality
\[
 W^4e^{-\eta W}\le24(2/\eta)^4e^{-\eta W/2}
\]
absorbs this polynomial factor.

For the complementary range $\eta W\le1$, the elementary
Chebyshev estimate proved by the binomial argument in QW10 gives
$\sum_{n\le u}\Lambda(n)\le C_\psi u$ with an absolute
$C_\psi$. Thus $|\psi_\chi(u)|\le2C_\psi u$; the principal
term is at most $u$ and the exceptional term at most $2u$.
Since $e^{-\eta W/2}\ge1/2$, the same error estimate follows
after increasing its constant. One permissible uniform choice is
\[
 c_7=c_0/64,\qquad
 C_7=2(2C_\psi+3)+(2P+22K+2)24(2/\eta)^4.
\]
We have proved, for every family, character and admissible $Z$,
\[
 \left|\psi_\chi(u)-\mathbf1_{\chi=1}u
                         +\sum_{\beta\in Z}\frac{u^\beta}{\beta}\right|
       \le C_7u e^{-c_7\sqrt{\log u}},
       \quad u\ge2,\quad\log\mathcal B\le\sqrt{\log u}.  \tag{QR.45}
\]
This supplies the literal existential uniform explicit formula. The
human proof uses a coarse absolute Chebyshev constant in the small
range; the source's sharper value is not needed for its conclusion.

\paragraph{QR23: select a single family-wide exception and bound its interval loss.}
Choose once
\[
 0<c_0\le\min(1/20,c_{\rm zf},c_{\rm un}).
\]
For each character let $Z_\chi$ consist of its real zeros in
$[1-c_0/\ell(0),1)$. QR17 and QR18 say these sets satisfy
\texttt{ExcZ}, contain simple zeros, and have total cardinality
at most one across the entire character group. Each such character
is real. QR14 therefore supplies the common separation
\[
       1-\beta\ge\Delta_{\mathcal B}:=c_9\mathcal B^{-46}
           \quad(\beta\in Z_\chi),\qquad0<\Delta_{\mathcal B}\le1.
                                                               \tag{QR.46}
\]
Every such $\beta$ is at least $3/4$.
For $X\ge e$, $0\le h\le1$ and $Y=X(1+h)$, the exact difference
of its two residue terms is
\[
 \frac{Y^\beta-X^\beta}{\beta}
       =\int_X^Y t^{\beta-1}\,dt
       \le hX X^{\beta-1}
       \le hX e^{-\Delta_{\mathcal B}}
       \le hX(1-\Delta_{\mathcal B}/2).                   \tag{QR.47}
\]
The last inequality follows, for $0\le v\le1$, from convexity:
$e^{-v}\le1-(1-e^{-1})v\le1-v/2$.
If no exception exists the loss is zero, also bounded by the
nonnegative right side of (QR.47). Thus one common main-term
margin remains after summing over all characters; summing a separate
worst-case exceptional loss for every character would have destroyed
this argument.

\paragraph{QR24: obtain the full weighted short-interval lower bound.}
Take $A_0=46$ and the constants $c_9,c_7,C_7$ just proved. Choose
\[
 c=\min\left(1,\frac{c_7}{2(A_0+2)},
                          \frac{c_7c_9}{24C_7}\right)>0.  \tag{QR.48}
\]
Suppose $X\ge1$, $\log\mathcal B\le c\sqrt{\log X}$ and
$e^{-c\sqrt{\log X}}\le h\le1$. Write $W=\sqrt{\log X}$.
Since $\mathcal B\ge4$, one has $cW\ge\log\mathcal B\ge1$,
so $X\ge e$, and (QR.45) applies at both $X$ and $Y=X(1+h)$.
Character orthogonality gives the exact identity
\[
       \sum_\chi\psi_\chi(u)=|G|\sum_{n\le u}\lambda_0(n).
                                                               \tag{QR.49}
\]
The sum of the two errors in (QR.45), over the $|G|\le\mathcal B$
characters, is at most $3\mathcal B C_7X e^{-c_7W}$.
It is bounded by $hX\Delta_{\mathcal B}/4$. Indeed the ratio
of these two positive quantities is at most
\[
 \frac{12C_7}{c_9}
       e^{((A_0+2)c-c_7)W}
 \le\frac{12C_7}{c_9}e^{-c_7W/2}\le1.                   \tag{QR.50}
\]
The first step uses $\mathcal B\le e^{cW}$ and $h\ge e^{-cW}$;
the second uses (QR.48). For the last, $W\ge1/c$ and
$c\le c_7c_9/(24C_7)$ imply
$c_7W/2\ge12C_7/c_9$, and $e^t\ge t$ completes the bound.

Subtract (QR.49) at $X$ from its value at $Y$, use the single
exception estimate (QR.47), and then (QR.50). For the actual floor
interval sum $S=\sum_{\lfloor X\rfloor<n\le\lfloor Y\rfloor}
\lambda_0(n)\ge0$ this proves
\[
 |G|S\ge hX\Delta_{\mathcal B}/4,\qquad
 S\ge\frac{c_9hX}{4\mathcal B^{47}}.                     \tag{QR.51}
\]
This is the complete \texttt{trivialClass\_\allowbreak shortInterval} supplier
for every \texttt{LFam}, with the permissible concrete exponent
$47$. Its constants are chosen before the family, $X$ and $h$.
The lower bound retains the source restrictions on the conductor
and the relative interval length; it makes no assertion for arbitrary
short intervals.

\paragraph{QR25: close the arithmetic prime suppliers used by the reconstruction.}
Insert QR9's actual family, with $\mathcal B=4dM^2$ and norm-one
prime support, into (QR.51). Both inputs formerly retained in QW11
are now supplied. The complete prime-power removal in QW10--QW11
then gives a fixed $c_{\rm SN}>0$, chosen no larger than $c$, such
that
\[
 \#\{p\in(X,X(1+h)]:p\text{ prime},\ \mathrm{NormOne}(d,M,p)\}
       \ge\frac{c_9hX}{32(4dM^2)^{47}\log X},             \tag{QR.52}
\]
whenever $d$ is prime, $d\equiv3\pmod4$, $d\ge100$, $M\ge3$,
$X\ge1$, $\log(4dM^2)\le c_{\rm SN}\sqrt{\log X}$ and
$e^{-c_{\rm SN}\sqrt{\log X}}\le h\le1$.
For example use QW11 with its exponent $A=47$ and coefficient
$e=c_9/4$ to choose $c_{\rm SN}$ explicitly in terms of the
absolute Chebyshev remainder constant. This proves the exact
\texttt{ShortNorm} assertion, including the floor endpoints.

QW12's uniform scale comparison now proves the literal
\texttt{levelPrimeN} and its actual-host specialization: after the
single threshold depending only on the fixed source constants, one
finds a prime $T$ with
$T_s<T\le(1+1/V)T_s$ and
$\mathrm{NormOne}(d,N(g_0g_1)^2,T)$ for every admissible later
endpoint and datum. QW13 similarly gives the original
$Y^{9/10}$ lower bound in the dyadic norm-prime interval.
These conclusions no longer retain \texttt{ExistsLFamRayStatement}
or \texttt{TrivialClassShortIntervalStatement} as unproved inputs.
Their original \texttt{ModAt}, host and rank inequalities remain
the explicit hypotheses of the geometric applications. This supply
of norm primes does not itself choose compatible nonempty local
stages, simultaneous \texttt{PairEx} data, or the centered prime-sum
objects required later in the RH and Dirichlet-GRH branches.


\section{Nonempty coefficient families and integral arithmetic realizations}
\label{sec:nonempty-realization}

\status{NR1--NR18 construct the actual all-prime coefficient family
\texttt{TwR/RingSpec}, its normalized characters, finite universal
anticyclotomic quotients, and readings at every prescribed coefficient
place. NR19--NR24 prove finite-field descent and integralization for
compact representations and the source's bounded family of row
representations, conditional on its stated automorphic Galois inputs.
The original number fields, rows, complex embeddings and host are
fixed throughout. This supplies nonempty twist indices and the exact
coefficient part of \texttt{SideB.cover}; it does not supply the
remaining geometric and automorphic form realizations.}

Fix the source datum $D$, its host CM field $E$, an algebraic closure
$\Omega/E$, its prescribed places on $\Omega$, and its character frame
$C$. Put
\[
 K=D.\mathrm{LL0}\subset L=D.\mathrm{what}_2.\mathrm{LL},\qquad
 \epsilon:L\hookrightarrow\mathbb C .                  \tag{NR.1}
\]
Both $K$ and $L$ are number fields. Whenever conjugate coefficient
embeddings are used, retain the literal hypothesis $\bar K=K$, with
conjugation induced from the fixed $L$ and $\epsilon$.
For a rational prime $\ell$, write $\Omega_\ell=\overline{\mathbb Q}_\ell$
and $C_\ell=\widehat{\Omega_\ell}$, with $|\ell|=\ell^{-1}$.
The algebraic closure $\Omega_\ell$ is not complete. A finite extension
always means a finite subextension of $\Omega_\ell/\mathbb Q_\ell$.
Its integer ring is the closed unit ball. We use the standard local
field facts that finite extensions are complete and their extended
absolute value is unique. The integer-ring properties needed here
are verified in NR3. The algebraic closedness of $C_\ell$, elementary
finite-dimensional nonarchimedean norm equivalence, the Baire theorem,
and finiteness of number-field ideal class groups are stated classical
inputs. Global Artin reciprocity, where used, has exactly the source's
normalization; it is not inferred from these local constructions.

\paragraph{NR1: embeddings realizing each actual finite prime.}
Let $q$ be a nonzero prime of $\mathcal O_K$, and let $\ell$ be the
characteristic of its finite residue field. Thus
$q\cap\mathbb Z=(\ell)$, with $\ell$ a rational prime. The completion
$K_q$ is a finite extension of $\mathbb Q_\ell$. Choose an embedding
$K_q\hookrightarrow\Omega_\ell$ over $\mathbb Q_\ell$ and restrict it
to $K$. Uniqueness of the extended absolute value gives an embedding
$J_q:K\hookrightarrow\Omega_\ell$ satisfying
\[
 a\in q\quad\Longleftrightarrow\quad |J_q(a)|<1
       \qquad(a\in\mathcal O_K).                      \tag{NR.2}
\]
This is the literal \texttt{exists\_\allowbreak embAt} condition, so the prescribed
family $J$ in \texttt{exists\_\allowbreak ringSpec} can be supplied, rather than
left as a new existence premise.

More generally, let $A$ be a valuation ring on $\Omega$ above a prime
of $E$ of residue characteristic $\ell$. There is an embedding
$j:\Omega\hookrightarrow\Omega_\ell$ such that
\[
                    x\in A\quad\Longleftrightarrow\quad |j(x)|\le1.
                                                               \tag{NR.3}
\]
One proof first chooses any embedding of the algebraic closure of
$\mathbb Q$ into $\Omega_\ell$. Its pulled-back valuation and $A$
restrict to the same $\ell$-adic valuation ring on $\mathbb Q$.
Extensions of a valuation to an algebraic normal extension are
conjugate: in each finite normal subextension this is transitivity on
primes above the given prime, and compactness of the inverse system
of the nonempty finite sets of conjugating automorphisms gives one
automorphism on the algebraic closure. Precomposing the first
embedding by its inverse proves (NR.3). This argument also explains
why a choice of prime must accompany a global algebraic embedding.

\paragraph{NR2: density and uniqueness of the normalized local norm.}
For every embedding $j:\Omega\hookrightarrow\Omega_\ell$, its image
is dense. Indeed, let $a\in\Omega_\ell$, choose a monic polynomial
$f\in\mathbb Q_\ell[X]$ with $f(a)=0$, and approximate its finitely
many coefficients by rational numbers. For any $\varepsilon>0$ this
gives a monic $g\in\mathbb Q[X]$ of the same degree $d$ with
$|g(a)|<\varepsilon^d$. All roots $b_1,\ldots,b_d$ of $g$ belong to
$j(\Omega)$, since this is an algebraically closed field containing
$\mathbb Q$. The equality
$|g(a)|=\prod_i|a-b_i|$ implies $|a-b_i|<\varepsilon$ for some $i$.

If $j,j':\Omega\hookrightarrow\Omega_\ell$ pull back the same
valuation ring, then $|j(x)|=|j'(x)|$ for every $x$. For $x\ne0$, test
membership of $\ell^{-m}x^n$ for every integer $m$ and positive integer
$n$. This compares $|j(x)|$ and $|j'(x)|$ with all numbers
$\ell^{m/n}$; these comparisons determine a positive real number.
Thus the map $j(x)\mapsto j'(x)$ is an isometry. Extend it and its
inverse to the complete field $C_\ell$. The extensions are inverse
field isometries, fix $\mathbb Q_\ell$ by density of $\mathbb Q$, and
preserve the elements algebraic over $\mathbb Q_\ell$. Restriction
therefore gives a unique
\[
             \tau\in\operatorname{Aut}_{\mathbb Q_\ell}(\Omega_\ell),
             \qquad \tau j=j'.                         \tag{NR.4}
\]
The completion is used in this proof; an extension theorem into the
incomplete field $\Omega_\ell$ has not been assumed.

\paragraph{NR3: the actual local decomposition bridge.}
For $j$ as in (NR.3), put $K_v=\mathbb Q_\ell(j(E))$.
A rational basis of $E$ gives finitely many algebraic generators,
so $K_v/\mathbb Q_\ell$ is finite. Its closure is the closure of
$j(E)$: the latter contains the closure $\mathbb Q_\ell$ of
$\mathbb Q$, and is a field containing the finitely many generators;
conversely the finite field $K_v$ is closed.
An element $\sigma$ of the decomposition group $D_A\subset G_E$
preserves $A$. NR2 extends $j\sigma j^{-1}$ uniquely to an
automorphism of $\Omega_\ell$ fixing $K_v$. Conversely an automorphism
fixing $K_v$ preserves the set $j(\Omega)$ of elements algebraic over
$\mathbb Q$ and the $\ell$-adic norm. Its restriction gives an element
of $D_A$. These constructions are inverse and yield
\[
       D_A\simeq\operatorname{Gal}(\Omega_\ell/K_v),
       \qquad j(\sigma x)=\tau_\sigma(jx).              \tag{NR.5}
\]
The inverse map is continuous: the condition of fixing finitely many
elements of $\Omega$ pulls back to the open stabilizer of finitely
many algebraic elements of $\Omega_\ell$. Both groups are compact
Hausdorff profinite groups, so a continuous bijection is a
homeomorphism. This proves the required topology in both directions.

Choose a finite normal extension $L_v/\mathbb Q_\ell$ containing
$K_v$. Every $\mathbb Q_\ell$-embedding of $K_v$ lands in $L_v$.
For every finite extension $K_0\supset L_v$, the ring
$\mathcal O_{K_0}$, the map $j$, and (NR.5) therefore give the
literal \texttt{LocBridge}, including all its coefficient embeddings.
The equivalence (NR.3) is its \texttt{br\_\allowbreak ind} clause. This is a
bridge for the prescribed place $A$, not a replacement place.

Here are the needed integer-ring facts, to make the base-ring
construction explicit as well. If $K_0/\mathbb Q_\ell$ has degree
$d$, the field norm gives $|x|^d\in\ell^{\mathbb Z}$ for $x\ne0$.
Choose an element of least positive value. Every nonzero element
of the unit ball is a unit times a nonnegative power of it, by
repeated division, so the unit ball $O_0$ is a DVR. Relative to
a $\mathbb Q_\ell$-basis of $K_0$, norm equivalence bounds all
coordinates of its unit ball. Thus $O_0$ is a $\mathbb Z_\ell$-submodule
of a finite fractional $\mathbb Z_\ell$-lattice; it is finite and
torsion-free, hence free. It is consequently integral over
$\mathbb Z_\ell$. Conversely a monic integral equation forbids
norm greater than one, so this unit ball is exactly the integral
closure used in the source. Multiplication by a sufficiently
large power of $\ell$ puts any element of $K_0$ into $O_0$, proving
$\operatorname{Frac}(O_0)=K_0$. Its residue field is a quotient of
the finite ring $O_0/\ell O_0$, and is finite. Powers of a
uniformizer are precisely successively smaller closed norm
balls, so the topology is adic, and closedness in the complete
field gives adic completeness. In the bridge the coefficient
isomorphism is the identity of $O_0$, hence is continuous in
both directions and respects $\mathbb Z_\ell$.

\paragraph{NR4: the value group of the unramified algebraic field.}
Let $U_\ell\subset\Omega_\ell$ be the field generated by all roots of
unity of order prime to $\ell$, the source's maximal unramified
extension. For $(m,\ell)=1$, set
$R_m=\mathbb Z_\ell[\mu_m]$. Products of $m$th roots are again $m$th
roots, so $R_m$ is their finite $\mathbb Z_\ell$-span. It is a compact,
and hence closed, subset of $\Omega_\ell$, and all its elements have
norm at most one. Put $N=\ell^{\varphi(m)}>1$. Euler's congruence
gives $\zeta^N=\zeta$ for $\zeta^m=1$, and reduction modulo $\ell$
gives $a^N\equiv a$ for $a\in\mathbb Z_\ell$. Applying Frobenius in
the quotient $R_m/\ell R_m$ therefore gives
\[
                   y^N-y\in\ell R_m\qquad(y\in R_m).  \tag{NR.6}
\]
If $|y|<1$, iteration shows
$\ell^{-1}(y-y^{N^s})\in R_m$ for every $s$. These elements tend
to $y/\ell$; closedness gives $y/\ell\in R_m$.
For $y\ne0$, repeat this division while the norm is less than one.
The process must terminate, since the norm is multiplied by $\ell$
at each step and elements of $R_m$ have norm at most one. Consequently
$|y|=\ell^{-a}$ for an integer $a\ge0$.

Every element of $U_\ell$ is a quotient of two elements of some
$R_m$. To check this without a finiteness assumption on $U_\ell$,
take the finitely many roots occurring in a rational expression,
replace their orders by their least common multiple, and clear the
finitely many $\mathbb Q_\ell$ denominators. Hence
\[
                        |U_\ell^\times|=\ell^{\mathbb Z}.
                                                               \tag{NR.7}
\]
This proves the discrete norm assertion used below directly from
the prime-to-$\ell$ roots.

\paragraph{NR5: completion preserves a discrete valuation ring.}
Let $L_0/\mathbb Q_\ell$ be finite and put $F=L_0U_\ell$.
The extension $F/U_\ell$ has finite degree $d>0$. Every embedding
of $F$ over $U_\ell$ preserves the $\ell$-adic norm, because it fixes
$\mathbb Q_\ell$ and the extended norm on an algebraic extension is
unique. The field norm thus satisfies
\[
 |N_{F/U_\ell}(x)|=|x|^d,\qquad
 |x|^d\in\ell^{\mathbb Z}\quad(x\in F^\times).         \tag{NR.8}
\]
The positive additive values of $F$ lie in $d^{-1}\mathbb Z$ and
form a nonempty set, since it contains the value of $\ell$.
Choose $\pi\in F$ having its least positive value. Equivalently,
$0<|\pi|<1$ is the largest norm of an element of norm below one.

Let $\widehat F$ be the closure of $F$ in $C_\ell$ and put
$O=\{x\in\widehat F:|x|\le1\}$. The norm gap survives completion:
the closed set
$\{x:|x|\le|\pi|\}\cup\{x:|x|\ge1\}$ contains $F$, hence
contains its closure. Dividing repeatedly by $\pi$ now shows that
every nonzero $x\in O$ has a unique form $u\pi^a$ with $u\in O^\times$
and $a\ge0$. For an ideal, the least value of a nonzero element
generates it. Thus $O$ is a DVR with maximal ideal $(\pi)$, and
\[
 x\in(\pi)^a\quad\Longleftrightarrow\quad |x|\le|\pi|^a.
                                                               \tag{NR.9}
\]
Its induced topology is exactly its maximal-ideal adic topology.
The field $\widehat F$ and its closed unit ball are complete, so $O$
is also adically complete.

Its residue field is algebraically closed. The residue field of
the algebraically closed valued field $C_\ell$ is algebraically
closed: lift a monic residue polynomial, take a root in $C_\ell$,
and use its monic equation to see that every root is integral.
The map $O/(\pi)\to\overline{\mathbb F}_\ell$ is injective.
It is onto because every nonzero residue is represented by a
prime-to-$\ell$ root of unity in $U_\ell$. Explicitly, approximate
a norm-one element of $C_\ell$ to distance less than one by an
algebraic element; in a finite residue field its nonzero reduction
has order prime to $\ell$, and the simple root of
$X^{\ell^f-1}-1$ lifts to a root of unity with that reduction
(equivalently, use the residue-root lifting of LS8).
Finally $O$ is the closure of the unit ball of $F$: sufficiently
close approximants to an element of norm at most one also have
norm at most one. These are all the assertions of
\texttt{exists\_\allowbreak integers\_\allowbreak completion}; $O$ is not the full,
nondiscrete valuation ring of $C_\ell$.

\paragraph{NR6: bounded coefficients after one algebraic adjunction.}
Let $F\subset\Omega_\ell$ contain $\mathbb Q_\ell$, let
$\widehat F\subset C_\ell$ be its closure, and let $z\in\Omega_\ell$
be algebraic over $\mathbb Q_\ell$. If a monic polynomial of degree
$n>0$ over $\mathbb Q_\ell$ kills $z$, the image of
\[
 T:\widehat F^n\longrightarrow C_\ell,
       \qquad (a_0,\ldots,a_{n-1})\longmapsto
                         \sum_{a<n}a_a z^a             \tag{NR.10}
\]
is $\widehat F[z]=\widehat F(z)$. Indeed, polynomial division spans
the algebra by these powers, and a finite-dimensional domain over a
field is a field. This image is a finite-dimensional normed space
over the complete field $\widehat F$, so it is complete and closed.
It contains $F(z)$, and every displayed linear combination is a
limit of combinations with coefficients in $F$. It follows that
\[
                     \overline{F(z)}=\widehat F(z).
                                                               \tag{NR.11}
\]
Choose a vector-space complement to $\ker T$. On that complement,
$T$ is an isomorphism of finite-dimensional normed spaces, whose
inverse is bounded. There is consequently a constant $C>0$ such
that every $x\in\overline{F(z)}$ with $|x|\le1$ admits coefficients
in (NR.10) satisfying $|a_a|\le C$. Powers of $z$ need not remain
linearly independent over $\widehat F$; the complement is what
justifies the bound when the degree drops after completion.

\paragraph{NR7: the finite cyclotomic coefficient tower.}
For finite $K_0,L_0\subset\Omega_\ell$, let
\[
 F=K_0L_0U_\ell,\quad O=\mathcal O_{\widehat F},\quad
 F_k=F(\zeta_{\ell^k}),\quad O_k=\mathcal O_{\widehat{F_k}},
                                                               \tag{NR.12}
\]
where $\zeta_{\ell^k}$ is primitive and $\zeta_1=1$.
NR5 applies to both $F$ and $F_k$, so $O,O_k$ are complete DVRs
with algebraically closed residue fields and their exact adic
topologies. Their inclusions are injective and $\ell$ is in their
maximal ideals.

Apply NR6 with $z=\zeta_{\ell^k}$. Choose $N$ so large that
$|\ell|^N C\le1$. For every $x\in O_k$,
\[
              \ell^N x\in
                O+O\zeta_{\ell^k}+\cdots+O\zeta_{\ell^k}^{n-1}.
                                                               \tag{NR.13}
\]
Thus $O_k$ is an $O$-submodule of a finite fractional $O$-module.
Since the DVR $O$ is Noetherian, $O_k$ is a finite $O$-module.
For any $x\in\operatorname{Frac}(O_k)$, multiply first by a
sufficient power of $\ell$ to make it integral and then use
(NR.13). This proves the exact generation statement
\[
 \exists d\in O\setminus\{0\}\ \exists p\in O[X],
                         \qquad dx=p(\zeta_{\ell^k}).  \tag{NR.14}
\]
The root has norm one and is a unit of exact order $\ell^k$.
For $k=0$ the two rings agree. These statements supply
\texttt{Ok\_\allowbreak fin}, \texttt{zk\_\allowbreak prim}, and \texttt{Ok\_\allowbreak gen}; the
stronger, unnecessary assertion $O_k=O[\zeta_{\ell^k}]$ has not
been used.

\paragraph{NR8: principal units cannot carry prime-to-residue torsion.}
Let $v$ be a finite place of $E$ of residue characteristic $p\ne\ell$,
and let $\chi:E_v^\times\to B$ be a continuous character with open
kernel, where every element of the finite group $B$ has
$\ell$-power order. If $u\in1+\mathfrak m_v$, then
$u^{p^a}\to1$. For completeness, if $u=1+x$ with
$\operatorname{ord}_v(x)\ge1$, the binomial theorem gives
\[
 \operatorname{ord}_v((1+x)^p-1)
          \ge \operatorname{ord}_v(x)+1.              \tag{NR.15}
\]
The middle binomial coefficients are divisible by $p$, and the
last term has value $p\operatorname{ord}_v(x)$, so this also covers
$p=2$. Iteration proves convergence. The open kernel eventually
contains $u^{p^a}$; hence $\chi(u)^{p^a}=1$. Its order is both an
$\ell$-power and a divisor of $p^a$, so $\chi(u)=1$.

On an infinite complex component every homomorphism to a finite
group is trivial: if the target has exponent $M$, write any
nonzero complex number as an $M$th power. Thus an idele character
as above, unramified outside a finite set $Q$ and also at every
prime over $\ell$, kills
\[
 U_Q=E_\infty^\times\times
       \prod_{v\in Q}(1+\mathfrak m_v)
                      \times\prod_{v\notin Q}\mathcal O_{E_v}^\times.
                                                               \tag{NR.16}
\]
At $v\in Q$ over $\ell$, its assumed unramifiedness kills all units;
at the other primes of $Q$, (NR.15) applies. To pass from local
components to the full product, approximate a unit idele by the
net of its finite-support components. The character kernel is
closed because it is an open subgroup, so it contains the limit.

\paragraph{NR9: a finite common idelic quotient.}
Let $U=E_\infty^\times\prod_v\mathcal O_{E_v}^\times$.
The quotient $U/U_Q$ injects into
$\prod_{v\in Q}\kappa(v)^\times$, and is finite. The valuation
map from ideles to fractional ideals is surjective, has kernel
$U$, and carries the diagonal $E^\times$ to principal ideals.
Consequently
\[
 \mathbb I_E/(E^\times U)\simeq\operatorname{Cl}(E),\qquad
                  [\mathbb I_E:E^\times U_Q]<\infty.  \tag{NR.17}
\]
Here finiteness of $\operatorname{Cl}(E)$ is the classical ideal
class theorem. No uniform bound for its order is required.
Equivalently, any subgroup of $\mathbb I_E$ containing $E^\times$
and $U_Q$ has finite index: first use the finite residue quotient
inside $U$, then the ideal class quotient. This is exactly the
common finiteness statement needed for all the characters in
NR10; their conductor exponents do not have to be bounded separately.

\paragraph{NR10: the universal kernel is finite index and open.}
Fix the prescribed conjugation $c$ of $\Omega$, a rational prime
$\ell$, and a finite set $Q$ of primes of $E$. Let $\mathcal K_Q$
be the kernels of all homomorphisms from $G_E$ to finite abelian
$\ell$-primary groups which have open kernel, satisfy
$\phi(cgc^{-1})=\phi(g)^{-1}$, and are unramified at every prime
outside $Q$ \emph{and at every prime over $\ell$}. Define
\[
                         H_Q=\bigcap_{H\in\mathcal K_Q}H.
                                                               \tag{NR.18}
\]
The trivial character ensures that this is an intersection of a
nonempty family. Each kernel is closed and contains the closure
of the commutator subgroup. Thus $G_E/H_Q$ is an abelian quotient
of $G_E^{\mathrm{ab}}$.

Use the supplied Global Artin reciprocity map
$a_E:\mathbb I_E\to G_E^{\mathrm{ab}}$. Its exact source quotient
map is surjective and kills diagonal elements. Its local unit
law identifies the image of local units with inertia. By NR8,
the composite of $a_E$ with every character defining
$\mathcal K_Q$ kills $U_Q$. Hence the composite
$\mathbb I_E\to G_E/H_Q$ kills $E^\times U_Q$ and is surjective.
NR9 makes $G_E/H_Q$ finite.

There are now only finitely many distinct subgroups $H$ in
$\mathcal K_Q$, because each contains $H_Q$ and corresponds to a
subgroup of the finite quotient. A finite intersection of such
kernels is again of the required kind: use the product of the
finitely many characters, whose target remains finite abelian
and $\ell$-primary. Therefore $H_Q$ itself is one of the kernels.
In particular it is open, and its quotient is $\ell$-primary.
This last argument is essential: an arbitrary infinite
intersection of open subgroups need not be open.

\paragraph{NR11: universal anticyclotomic quotients and restriction.}
Set
\[
     \operatorname{RCG}_{\ell,Q}=G_E/H_Q,
             \qquad a_{\ell,Q}:G_E\twoheadrightarrow
                                      \operatorname{RCG}_{\ell,Q}.
                                                               \tag{NR.19}
\]
NR10 proves finiteness, an open kernel and $\ell$-power order for
every element. The definitions prove the anticyclotomic law and
unramifiedness outside $Q$, as well as unramifiedness at primes
over $\ell$. If $Q$ contains no prime over $\ell$, every finite
abelian $\ell$-primary anticyclotomic character with open kernel
and ramification contained in $Q$ occurs in (NR.18). It therefore
factors uniquely through $a_{\ell,Q}$. If $Q'\subset Q$, the family
of permitted characters for $Q'$ is smaller, so
$H_Q\subset H_{Q'}$. The quotient map gives
\[
 \operatorname{res}_{Q,Q'}:\operatorname{RCG}_{\ell,Q}
      \longrightarrow\operatorname{RCG}_{\ell,Q'},\qquad
 \operatorname{res}_{Q,Q'}a_{\ell,Q}=a_{\ell,Q'}.        \tag{NR.20}
\]
These maps are compatible under further restriction by surjectivity
of $a_{\ell,Q}$.
This proves the full \texttt{exists\_\allowbreak ringClass\_\allowbreak host} supplier.
The restriction that $Q$ avoid $\ell$ in \texttt{art\_\allowbreak univ} is
retained exactly. The quotient exists also for other $Q$, but it
still kills inertia at $\ell$; it is not asserted to classify
arbitrarily ramified $\ell$-adic characters. Nor does universality
assert that this finite quotient is nontrivial or that a character
ramifies at every prime in $Q$.

\paragraph{NR12: the lower bound is chosen before the coefficient rings.}
For every $(\ell,j)$, where $j:K\hookrightarrow\Omega_\ell$, define
$c_0(\ell,j)$ to be the compositum of
\[
 \mathbb Q_\ell(j(K)),\qquad
 \mathbb Q_\ell(\mu_\ell),\qquad
                 L_v\quad(v\mid\ell\text{ in }E),    \tag{NR.21}
\]
where the $L_v$ are the finite normal fields from NR3 for the
prescribed places on $\Omega$. There are finitely many primes
above $\ell$, so this is a finite extension. It is chosen before
any of the fields denoted $b_0,b_1$ in the source. For every
$K_0\supset c_0(\ell,j)$ it contains $j(K)$, the $\ell$th roots,
and the coefficients of every required local bridge. This is
the literal \texttt{exists\_\allowbreak coefficientField\_\allowbreak lower\_\allowbreak bounds},
with its order of quantifiers unchanged.

We also need compatible complex readings. Each of $\Omega_\ell$
and $C_\ell$ is an algebraically closed field of characteristic
zero and cardinality continuum. The lower cardinal bound comes
from $\mathbb Q_\ell$; the upper bound for the algebraic closure
comes from counting polynomials, and for its metric completion
from counting sequences. An algebraically closed field of this
uncountable cardinality has transcendence degree continuum over
the countable field $K$. Extend the prescribed isomorphism
$j(K)\to\epsilon(K)$ to their algebraic closures, choose a
bijection of transcendence bases, and extend over the resulting
algebraic closures. Thus there are abstract field isomorphisms
\[
 \iota:\Omega_\ell\simeq\mathbb C,\qquad
 \iota':C_\ell\simeq\mathbb C,\qquad
           \iota j=\epsilon|_K,\quad\iota'j=\epsilon|_K.
                                                               \tag{NR.22}
\]
They need not agree off $j(K)$ and are not claimed continuous.
Only the explicitly displayed agreements are used.

\paragraph{NR13: the normalized characters and the sign of their reading.}
Retain \texttt{ValLL0} for the \emph{fixed} embedding
$\epsilon|_K$, the integrality of every exponent $e(j,g)$, and
the actual character frame $C$. HF13 constructs that frame;
HF16 supplies \texttt{ValLL0} only when its common value-field
generator is contained in the specified $K$. No enlargement of
an already fixed $D$ is implicit here.
Write $\delta_\ell=\delta'$ when $\ell=p_\delta$ and
$\delta_\ell=\delta$ otherwise. Its ramification set is the
primes over $p_\delta'$ or $p_\delta$, respectively, so it always
avoids $\ell$. For $1\le j\le m$ put
\[
 \beta_{\ell,j}
     =\mu_{0,j}\delta_\ell\eta^{-1}|\cdot|^{-r/2},
 \qquad
 \beta_{\ell,j}(P)
     =\mu_{0,j}(P)\delta_\ell(P)\eta(P)^{-1}N(P)^{r/2}.
                                                               \tag{NR.23}
\]
The actual $r$ is odd. The half-integral value law for $\eta$
and the integral value laws for $\mu_{0,j},\delta_\ell$ therefore
put these unramified values in $\epsilon(K)$, by GA11--GA12.

Take $K_0\supset c_0(\ell,j_0)$, its integer ring $O_0$, the
inclusion $\kappa:\operatorname{Frac}(O_0)=K_0\hookrightarrow
\Omega_\ell$, and $\iota$ from (NR.22). GA16 constructs the
actual continuous characters $\Delta:G_E\to O_0^\times$ and
$\chi_{1,j,b}:G_E\to O_0^\times$. The first is the finite
avatar of $\delta_\ell$; the second is the normalized avatar of
$\beta_{\ell,j}$, using $\iota$ for $b=\mathrm{true}$ and
$\overline{\phantom z}\circ\iota$ for $b=\mathrm{false}$.
Their coefficient embeddings are respectively $j_0$ and
$j_0\circ\bar{\phantom x}$. Conjugating twice shows that both
read the same specified complex embedding of $K$.

For the source's arithmetic Frobenius $\sigma_P$,
\[
 \iota\kappa(\Delta(\sigma_P))=\delta_\ell(P)^{-1},
 \qquad
 \iota_b\kappa(\chi_{1,j,b}(\sigma_P))
                                   =\beta_{\ell,j}(P)^{-1}
                                                               \tag{NR.24}
\]
at the permitted unramified primes. Since
\[
 (\mu_{0,j}\delta_\ell\eta^{-1})(P)
       =\beta_{\ell,j}(P)(\sqrt{N(P)})^{-r},            \tag{NR.25}
\]
the literal \texttt{NC.Read} exponent is $-r$, while that for
$\Delta$ is zero. To verify the record, discard the finite
ramification set and the primes over $\ell$, choose
$y\in K$ with $\epsilon(y)=\beta_{\ell,j}(P)^{-1}$, and apply
injectivity of $\iota_b\kappa$ to (NR.24). This gives the required
equality in $\operatorname{Frac}(O_0)$ for every chosen Frobenius
lift, with (NR.25) giving its complex value clause.
GA16 also proves the exact auxiliary inertia order $15$,
anticyclotomic law, off-set unramifiedness and prescribed tame
type-prime characters. They are properties of these normalized
avatars, not of freely chosen finite-order complex characters.
For the source's family indexed by every $j\in\mathbb N$, use
its actual reindexing $a(j)=\min(j-1,m-1)\in\{0,\ldots,m-1\}$,
where subtraction in this formula is truncated natural
subtraction and the rank data give $m\ge1$. Define the $j$th
character using row $a(j)$. This agrees with row $j-1$ on
$1\le j\le m$; all continuity, ramification and tame laws hold
also outside that range because they hold for every valid row.
The \texttt{Read} assertion itself is required only in the
literal range $1\le j\le m$.

\paragraph{NR14: construction of a genuinely nonempty twist-ring family.}
Fix once and for all $J_q$ from NR1. Given arbitrary families
$b_0,b_1$ of finite extensions with $b_0(\ell,j)\supset c_0(\ell,j)$,
and a finite set $S_x$ of primes of $E$, make the following choices
for every $q\in\operatorname{Spec}_1\mathcal O_K$:
\[
 \begin{split}
 \mathrm{Tw}&=\operatorname{Spec}_1\mathcal O_K,\qquad
 \ell(q)=\operatorname{char}\kappa(q),\qquad j_q=J_q,\\
 K_{0,q}&=b_0(\ell(q),j_q),\qquad
 L_q=b_1(\ell(q),j_q),\qquad O_{0,q}=\mathcal O_{K_{0,q}},\\
 O_q&=\mathcal O_{\widehat{K_{0,q}L_qU_{\ell(q)}}},\qquad
 O_{q,k}=\mathcal O_{\widehat{K_{0,q}L_qU_{\ell(q)}
                                      (\zeta_{\ell(q)^k})}}.
 \end{split}                                                   \tag{NR.26}
\]
The embedding $\mathrm{emb}_0:K\to\operatorname{Frac}(O_{0,q})$
is the codomain restriction of $J_q$, and
$\mathrm{emb}_{0b}(x)=\mathrm{emb}_0(\bar x)$. All algebra maps
between the displayed rings are their actual inclusions.
Use NR13 for $\Delta,\chi_1$ and NR11 for
$\operatorname{RCG}_{\ell(q),Q}$, its Artin quotient and its
restriction maps. Finally set
\[
 S_0(q)=\{v:v\mid\ell(q)\}\cup
       \bigcup_{a=0,1}\{v:v\mid g_a\}\cup
       \{v:v\mid p_\delta\}\cup\{v:v\mid p_\delta'\}\cup S_x.
                                                               \tag{NR.27}
\]
This set is finite because each nonzero rational integer has
only finitely many prime divisors in a number field.

These formulas supply every field of the source's \texttt{TwR}:
the finite $\mathbb Z_\ell$ coefficient DVR and finite residue
field are $O_{0,q}$; NR5--NR7 supply the complete DVR $O_q$,
algebraically closed residue field and finite cyclotomic tower;
NR13 supplies the actual normalized characters and exact
ramification clauses; NR11 supplies all the universal finite
groups and transition laws. To choose the conjugation of
$\Omega$, embed this algebraic closure of the number field into
$\mathbb C$. Its image is the field of algebraic numbers and
is stable under complex conjugation. Pulling conjugation back
gives an involution; on $E$ it is the host CM conjugation, since
every complex embedding of a CM field respects that conjugation.
This is the exact \texttt{IsCMConjOm} input used in NR10 and GA16.
Every rational prime has at least one prime above it in $K$:
the integral extension $\mathbb Z\subset\mathcal O_K$ has lying
over, or equivalently the finite ring $\mathcal O_K/\ell$ has a
maximal ideal. Hence this \texttt{Tw} is nonempty, indeed infinite.
For every one of its indices, $\Delta$ has a nontrivial inertia
image of order $15$.

\paragraph{NR15: all RingSpec compatibilities on the same fields.}
For the family (NR.26), use the identity equivalence
$\mathrm{Tw}\simeq\operatorname{Spec}_1\mathcal O_K$, the literal
identifications of $O_0,O,O_k$ with their integer rings, and the
bridges from NR3. Equation (NR.2) identifies the center of
$\mathrm{emb}_0$ with $q$ and proves \texttt{pl\_\allowbreak emb}.
The bounds in \texttt{RingSpec} hold with equality
\[
        K_{0,q}=b_0(\ell(q),j_q),\qquad
                   L_q=b_1(\ell(q),j_q).              \tag{NR.28}
\]
The diagrams for $\mathbb Z_\ell$, the base ring, and its finite
cyclotomic extensions commute because all maps are inclusions.
For example, for $z\ne0$, the two equations
$\mathrm{emb}_0(x)z=y$ and $J_q(x)z=y$ are equivalent after
the same injective fraction-field inclusion. The condition
\texttt{mu\_\allowbreak g} follows from the $\mu_\ell$ summand in (NR.21),
and the excluded set is exactly (NR.27).

Embed $\operatorname{Frac}(O_q)\subset C_{\ell(q)}$ and compose
with $\iota'$ from (NR.22) to define \texttt{iCq}. Define
$\mathrm{embO}(q,b)$ by composing $\mathrm{emb}_{0,b}$ with the
fraction-field inclusion. Then
\[
 \mathrm{iCq}_q\mathrm{embO}(q,\mathrm{true})=\epsilon|_K,\qquad
 \mathrm{embO}(q,\mathrm{false})(x)
                   =\mathrm{embO}(q,\mathrm{true})(\bar x).
                                                               \tag{NR.29}
\]
If a value of the base embedding is integral, this same inclusion
gives its required value in $O_q$, proving the sided fixed-field
clause without choosing another embedding.

Take $\mathrm{mu0H}(j)=\mu_{0,j}$ in its literal row range, and
the trivial character outside that range; take
$\mathrm{etaH}=\eta$, $\mathrm{delH}(q)=\delta_{\ell(q)}$.
NR13 proves the two \texttt{Read} clauses and their precise
exponents. HF11 gives the unramifiedness and the specified
unit-argument infinity law of $\eta$. The value equalities with
the original character frame hold by definition.
If two indices have the same auxiliary ramification set, their
$\mathrm{delH}$ agree. The only two candidate sets are the
nonempty sets of primes above the distinct rational primes
$p_\delta,p_\delta'$; these sets are disjoint, since a prime ideal
cannot contain two coprime rational primes. Thus equal sets
choose the same character. Finally the stronger
\texttt{deltaRead} assertion follows from the first equality
of (NR.24) at every permitted unramified prime, not merely
cofinitely many primes.

All fields of \texttt{RingSpec} have now been checked. We have
proved \texttt{exists\_\allowbreak ringSpec} with its exact order
\[
       \exists c_0\ \forall b_0\supset c_0\ \forall b_1\ \forall S_x
                   \ \exists(T,\mathrm{RingSpec}),     \tag{NR.30}
\]
and with the prescribed $J_q$ unchanged. Neither the rows nor
$K,L,\epsilon$ were replaced in this construction.

\paragraph{NR16: equality of centers determines the valuation ring.}
Let $v$ be a nontrivial additive valuation on a number field $K$,
with values in $\mathbb Z\cup\{\infty\}$. Every algebraic integer
has nonnegative value: it is integral over $\mathbb Z$, hence
over the valuation ring, which is integrally closed. Its center
\[
                   q_v=\{a\in\mathcal O_K:v(a)>0\}    \tag{NR.31}
\]
is a prime ideal. Additivity and the ultrametric inequality
give the ideal laws, and nonnegative values on integers show
that if $v(ab)>0$, one of $v(a),v(b)$ is positive.

Suppose $q_v=q\ne0$ and $e:K\hookrightarrow\operatorname{Frac}(O)$
is an embedding into a DVR whose maximal-ideal center is also $q$.
Algebraic integers map into $O$ by integral closedness. An element
$s\in\mathcal O_K\setminus q$ has $v(s)=0$ and $e(s)\in O^\times$.
Since $(\mathcal O_K)_q$ is a DVR, for every $x\in K$ either
$x=a/s$ or $x=s/a$, with $a,s\in\mathcal O_K$ and $s\notin q$.
In the first case $v(x)\ge0$ and $e(x)\in O$. In the second,
if $a\notin q$ both $x$ and $e(x)$ are units; if $a\in q$ then
$v(x)<0$ and $e(x)$ cannot be integral, because the product of
an integral element and $e(a)\in\mathfrak m_O$ cannot be the
unit $e(s)$. Therefore
\[
                 v(x)\ge0\quad\Longleftrightarrow\quad e(x)\in O
                               \qquad(x\in K).         \tag{NR.32}
\]
This upgrades a condition on integral elements of $K$ to the
condition on every element of $K$ required by \texttt{SidePl}.

\paragraph{NR17: exact all-place coefficient coverage.}
Let $w$ be any of the fixed places of $L$ in $D$. Restrict its
valuation to $K$ and take the center $q$ from (NR.31). It is
nonzero because the residue characteristic $p$ of $w$ has positive
value. NR15 gives the index $i=q$, whose coefficient embedding
has exactly this center. Both $p$ and $\ell(i)$ lie in $q$;
if they were different primes, a Bezout relation would put
$1$ in $q$. Hence $\ell(i)=p$. Applying NR16 yields
\[
 \forall x\in K,\qquad
   v_w(x)\ge0\quad\Longleftrightarrow\quad
                 \mathrm{emb}_0(i)(x)\in O_{0,i}.     \tag{NR.33}
\]
This is \texttt{exists\_\allowbreak setting\_\allowbreak of\_\allowbreak place}, and exactly the
$b=\mathrm{true}$ coefficient assertion in \texttt{SideB.cover}.
Taking positive values on $\mathcal O_K$ also proves the literal
\texttt{TwCover} center condition. The proof applies to every
$w$, irrespective of whether a row happens to be nonzero at $w$.
It does not assert the other fields of \texttt{SideB}.

\paragraph{NR18: a reading inducing a prescribed full-field place.}
The preceding coverage can be strengthened to the full
\texttt{Reading.Induces} condition, which concerns $L$, not just
$K$. Fix $w$ and the index $i$ from NR17. Choose
$k_w:L\hookrightarrow\Omega_{\ell(i)}$ inducing $w$ by NR1.
The restrictions $k_w|_K$ and $j_i$ induce the same valuation ring
on $K$. The rational-power comparison in NR2 proves that their
normalized absolute values agree. Their map on $K$ consequently
extends to an isomorphism of the finite local completions
\[
       \mathbb Q_\ell(k_w(K))\simeq\mathbb Q_\ell(j_i(K)),
                              \qquad k_w(x)\longmapsto j_i(x).
                                                               \tag{NR.34}
\]
An embedding of a finite separable extension into an algebraic
closure extends to an automorphism of that algebraic closure.
Choose such a $\mathbb Q_\ell$-automorphism $\tau$ and replace
$k_w$ by $\tau k_w$. Its norm is unchanged, so it still induces
$w$, and its restriction to $K$ is now exactly $j_i$.

Apply the transcendence-basis argument of NR12 to the number
field $L$ and the embeddings $k_w,\epsilon$. It gives an abstract
isomorphism $\iota_w:\Omega_\ell\simeq\mathbb C$ satisfying
$\iota_w k_w=\epsilon$ on all of $L$. Keep the canonical inclusion
$\kappa_i:\operatorname{Frac}(O_{0,i})\hookrightarrow\Omega_\ell$.
Then
\[
 \iota_w\kappa_i\mathrm{emb}_0(i)=\epsilon|_K,\qquad
 v_w(x)\ge0\ \Longleftrightarrow\
             |\iota_w^{-1}\epsilon(x)|\le1\quad(x\in L).
                                                               \tag{NR.35}
\]
The first equation and the unchanged $\mathbb Z_\ell$ inclusion
make $(\kappa_i,\iota_w)$ a literal \texttt{Reading}; the second
is its \texttt{Induces} clause. For a general already supplied
index and side, the identical argument applies whenever its
\texttt{SidePl} condition for $w$ holds, using
$\kappa_i\mathrm{emb}_{0,b}$ in place of $j_i$.
After any installation of this same ring family into an
$\mathrm{NTw}$, the assertion refers to that installation's
unchanged rings and embeddings. It supplies no signature,
residual representation or form comparison on its own.

\paragraph{NR19: a countable cofinal family of finite local fields.}
Fix an embedding $j:\overline{\mathbb Q}\hookrightarrow\Omega_\ell$.
For each $x\in\Omega_\ell$, choose $y\in j(\overline{\mathbb Q})$
so close to $x$ that
\[
 |x-y|<|x-x'|\quad\text{for every distinct conjugate }x'
                                  \text{ of }x\text{ over }\mathbb Q_\ell.
                                                               \tag{NR.36}
\]
There are finitely many such conjugates and NR2 gives density.
If $x\in\mathbb Q_\ell$, the conclusion below already holds.
Otherwise, every automorphism $\tau$ fixing $\mathbb Q_\ell(y)$
preserves the norm and satisfies
$|\tau x-x|\le\max(|\tau x-y|,|y-x|)=|x-y|$.
Equation (NR.36) forces $\tau x=x$. The fixed-field theorem in
the separable algebraic closure gives
\[
                              x\in\mathbb Q_\ell(y).  \tag{NR.37}
\]
This is the Krasner argument needed here, with its strict
inequality made explicit.

The field $\overline{\mathbb Q}$ is countable, and so is its set
of finite subsets. The fields
$K_S=\mathbb Q_\ell(j(S))$ for these finite subsets are finite
extensions. Every finite subset of $\Omega_\ell$ is contained
in one $K_S$, by applying (NR.37) to each of its elements.
Thus this is a countable cofinal family of finite subextensions,
although neither $\mathbb Q_\ell$ nor $\Omega_\ell$ is countable.

\paragraph{NR20: a compact representation has one finite coefficient field.}
Let $G$ be a compact Hausdorff topological group and
$\rho:G\to\operatorname{GL}_n(\Omega_\ell)$ a homomorphism whose
matrix entries are continuous. For the countable family of NR19
define
\[
 H_S=\{g:\text{all entries of }\rho(g)\text{ and }\rho(g^{-1})
                                                \text{ lie in }K_S\}.
                                                               \tag{NR.38}
\]
These are subgroups, by matrix multiplication and inversion.
They are closed: a finite extension of $\mathbb Q_\ell$ is
complete and therefore closed in $\Omega_\ell$, and there are
only finitely many continuous entry conditions. They cover
$G$, because a matrix and its inverse have finitely many entries.
The compact Hausdorff Baire theorem gives an $H_S$ with nonempty
interior. A subgroup with nonempty interior is open, and an open
subgroup of a compact group has finite index.

Choose representatives $g_1,\ldots,g_a$ for its finitely many
cosets. Adjoin their finitely many matrix entries to $K_S$,
obtaining a finite extension $K'$. For $g=g_jh$ with $h\in H_S$,
the equality $\rho(g)=\rho(g_j)\rho(h)$ puts every entry in $K'$.
The same assertion applied to $g^{-1}$ gives the inverse matrix
over $K'$. Hence the entire representation is defined over one
finite extension. If an initial finite field $K_{\mathrm{given}}$
must be retained, replace $K'$ by its compositum with that field.
This proves \texttt{exists\_\allowbreak finExt\_\allowbreak entries\_\allowbreak mem} with its arbitrary
lower-bound field. It does not assert that the image of $\rho$
is finite; compactness is used through Baire and boundedness,
not through such a claim.

\paragraph{NR21: a stable integral lattice and exact conjugation.}
With the finite field $K'$ from NR20, put $O'=\mathcal O_{K'}$
and choose a uniformizer $\pi$. The finitely many coordinate
functions of $\rho$ have bounded image on the compact group.
For some $N$,
$\rho(g)O'^n\subset\pi^{-N}O'^n$ for every $g$. Define
\[
                         \Lambda=\sum_{g\in G}\rho(g)O'^n,
              \qquad O'^n\subset\Lambda\subset\pi^{-N}O'^n.
                                                               \tag{NR.39}
\]
The sum denotes finite sums, so it is an $O'$-submodule.
The right-hand module is finite over the Noetherian DVR $O'$;
its submodule $\Lambda$ is therefore finite. It is torsion-free,
so the elementary divisor theorem over a DVR makes it free.
The two inclusions show that its rank is $n$. For every $h\in G$,
left multiplication permutes the summands, giving
$\rho(h)\Lambda=\Lambda$.

Choose an $O'$-basis of $\Lambda$ and let $P\in\operatorname{GL}_n(K')$
have those columns. Then
\[
                    P^{-1}\rho(g)P\in\operatorname{GL}_n(O')
                                  \quad\text{for all }g.       \tag{NR.40}
\]
The inverse is integral as well, since (NR.40) also holds for
$g^{-1}$. Matrix multiplication by the fixed $P,P^{-1}$ preserves
continuity. The induced topology on $O'$ is its adic topology,
so the resulting integral representation is continuous with
that topology. Conjugation preserves characteristic polynomials
and carries invariant direct-sum decompositions bijectively,
so it also preserves semisimplicity. This proves the actual
\texttt{exists\_\allowbreak finExt\_\allowbreak conj\_\allowbreak intOf} endpoint, including rank
zero if it occurs, without asserting that its integral model
has been identified with any pre-existing residual model.

\paragraph{NR22: the row twist has the exact Galois normalization.}
Now retain the literal \texttt{SideA} and \texttt{GalEx} of $D$.
The latter asserts, for every rational prime $\ell$, every complex
reading $\iota$, and every conjugate-self-dual regular algebraic
cuspidal $\pi$ of positive rank $a$, existence of a continuous
semisimple representation realizing $\pi$ at shift $(1-a)/2$.
This is a substantive automorphic Galois existence input, and
is not proved by NR20--NR21.

For $a=r$, let $\rho_1$ be the representation supplied by
\texttt{GalEx}. The actual odd integer $t$ and the half-integral
infinity type of $\eta$ make
\[
                     \gamma=\eta^{-1}|\cdot|^{-t/2}   \tag{NR.41}
\]
an algebraic Hecke character, as in GA11--GA13. Global Artin
reciprocity and GA6 give its continuous one-dimensional avatar
$\nu:G_E\to\Omega_\ell^\times$ with
$\iota(\nu(\sigma_P))=\gamma(P)^{-1}$ at the specified unramified
Frobenius elements. Define $\rho=\nu\rho_1$ by scalar multiplication.
It is continuous. A subspace is invariant for $\rho$ precisely
when it is invariant for $\rho_1$, so semisimplicity is preserved.

The \texttt{SideA} character reading identifies the local
one-dimensional representation $\eta_T^{-1}$ with $\eta^{-1}$
in precisely the source convention. The scalar
$u_P=\gamma(P)^{-1}$ is therefore the sole Frobenius eigenvalue
of that local character with the extra shift $-t/2$. If
$\alpha_1,\ldots,\alpha_r$ are the normalized eigenvalues of
$\rho_1(\sigma_P)$, then
\[
 \det(XI-u_P\,\iota(\rho_1(\sigma_P)))
                              =\prod_{a=1}^r(X-u_P\alpha_a),
 \qquad {1-r\over2}-{t\over2}={1-n_V\over2},
                          \quad n_V=r+t.              \tag{NR.42}
\]
Thus $\rho$ realizes $\pi\otimes\eta_T^{-1}$ at exactly
$(1-n_V)/2$. The finitely many ramified primes and primes over
$\ell$ can be removed together; the equality holds for every
Frobenius lift at each remaining prime. This proves
\texttt{exists\_\allowbreak isRepOf\_\allowbreak tw}, not a new assertion of \texttt{GalEx}.

\paragraph{NR23: keeping the representation while changing its reading.}
There are three precise invariances used in the bounded family.
First, suppose $\iota_0,\iota_1:\Omega_\ell\simeq\mathbb C$
both satisfy $\iota_bj=\epsilon|_K$. Suppose a representation
$\rho$ realizes $\pi$ under $\iota_0$, and that cofinitely many
of the required normalized Frobenius polynomials are
$p_P^\epsilon$ for polynomials $p_P\in K[X]$. Injectivity of
$\iota_0$ forces
\[
                \det(XI-\rho(\sigma_P))=p_P^j.
                                                               \tag{NR.43}
\]
Applying $\iota_1$ gives the same required complex polynomial.
Continuity and semisimplicity are properties of $\rho$ and
remain unchanged. No continuity of either $\iota_b$ is needed.

Second, if $\tau\in\operatorname{Aut}_{\mathbb Q_\ell}(\Omega_\ell)$,
it is an isometry, and applying it entrywise preserves continuity,
semisimplicity and matrix characteristic polynomials. Thus a
representation read with $\iota\tau$ becomes one read with
$\iota$ after applying $\tau$ to its coefficients.

Third, retain the literal \texttt{ChebSplit}: the specified
Frobenius lifts at split primes outside any finite set are dense
in $G_E$. If two continuous representations have the same
normalized Frobenius polynomials at those primes, every
coefficient of their characteristic polynomials agrees on a
dense subset. The coefficient functions are continuous into
the Hausdorff field $\Omega_\ell$, so
\[
                \det(XI-\rho(g))=\det(XI-\rho'(g))
                                      \quad(g\in G_E).           \tag{NR.44}
\]
In particular, if $\rho'$ realizes $\pi'$ and the local
representations of $\pi,\pi'$ agree at split primes, the original
$\rho$ also has the required Frobenius polynomials for $\pi'$.
This proves the source's \texttt{isRepOf\_\allowbreak of\_\allowbreak split}; it need not
choose an isomorphism between the two representations, and it
does not assert that split agreement constructs either one.

\paragraph{NR24: a uniform finite bound for all occurring row representations.}
In addition to the data of NR22, retain the actual
\texttt{RowCusp} and \texttt{ChebSplit}. Write
\[
 \mathrm{GoodOcc}(\sigma)\ \Longleftrightarrow\
 D.\mathrm{goodRow}(\sigma)\ \text{and}\
 \exists j\in[1,m],\ D.\mathrm{isot}(j,\sigma)\ne0.
                                                               \tag{NR.45}
\]
For each such $\sigma$, \texttt{RowCusp} supplies one actual
cuspidal $\pi_\sigma$, conjugate-self-dual and regular algebraic,
whose local representations agree with the given row at every
split prime. It also supplies, cofinitely in $P$, the
$K$-polynomials for $\pi_\sigma\otimes\eta_T^{-1}$ at the shift
$(1-n_V)/2$. This latter clause is what permits NR23; it is not
an automatic consequence of a name for the row.

For each $(\ell,j:K\hookrightarrow\Omega_\ell)$ and each good
occurring row, choose a complex reading extending $j$, apply
NR22, and apply NR21. Obtain an integral representation
$\rho_{\ell,j,\sigma}$ over a finite field $L_{\ell,j,\sigma}$.
NR23 shows that this \emph{same} representation works for every
complex reading extending $j$ and every conjugate-self-dual
regular algebraic $\pi$ with the same split local row. In the
latter comparison, \texttt{GalEx} first supplies a representation
for $\pi$; (NR.44) then transfers its Frobenius identities to the
already chosen representation.

There are finitely many row types $\sigma$ in the source. Put
$j_{\mathrm{true}}=j$ and $j_{\mathrm{false}}=j\circ\bar{\phantom x}$,
and define the finite extension
\[
 c_1(\ell,j)=\bigvee_{\sigma,\,b\in\{\mathrm{false},\mathrm{true}\}}
                               L_{\ell,j_b,\sigma}.   \tag{NR.46}
\]
Here $\bigvee$ denotes the compositum of this finite family;
for a row failing \texttt{GoodOcc}, use the base field and the
trivial representation, with no assertion of an automorphic
realization for that row. Equivalently, (NR.46) is the smallest
subfield of $\Omega_\ell$ containing all the displayed fields.
The bound $c_1$ is chosen before any $b_0,b_1,S_x,T$ or
$\mathrm{RingSpec}$.

For every supplied \texttt{RingSpec} with $b_0\supset c_1$,
its actual $O_{0,i}\simeq\mathcal O_{K_{0,i}}$ contains the
matrix entries of $\rho_{\ell(i),j_{i,b},\sigma}$. Restrict the
representation, including its inverse matrices, to this ring
and transport by that specified ring isomorphism. The result
is a continuous
\[
       \rho_{i,\sigma,b}:G_E\longrightarrow
                                      \operatorname{GL}_r(O_{0,i}).
                                                               \tag{NR.47}
\]
Continuity follows from the adic topology and its embedding
into $\Omega_{\ell(i)}$, as in GA1--GA3.

This family has the full universal reading clause of
\texttt{rho\_\allowbreak rows}. Namely, let
$\phi:\operatorname{Frac}(O_{0,i})\hookrightarrow\Omega_{\ell(i)}$
respect the specified $\mathbb Z_{\ell(i)}$ map, and let
$\iota\phi\mathrm{emb}_{0,b}=\epsilon|_K$.
The two embeddings $\phi$ and the canonical $\kappa_i$ of the
finite local field are conjugate over $\mathbb Q_{\ell(i)}$:
extend their isomorphism of finite subfields to obtain
$\phi=\tau\kappa_i$. Now
$\iota\tau j_{i,b}=\epsilon|_K$. The universal statement for the
chosen representation and the second part of NR23 show that
$\phi(\rho_{i,\sigma,b})$ realizes every permitted
$\pi\otimes\eta_T^{-1}$ at $(1-n_V)/2$. The equality is about
these matrices under these actual embeddings, not an unrelated
representation having the same dimension.

Finally choose $b_0(\ell,j)=c_0(\ell,j)c_1(\ell,j)$, or any
larger finite field, and any $b_1$. NR15 constructs the all-prime
ring family for these choices; the preceding construction then
gives its row representations. NR18 supplies readings inducing
each specified full-field place, and the universal reading
clause just proved applies to them as well. This is a coherent
simultaneous choice on the fixed $D$. Its only new coefficient
enlargements are the explicitly permitted local fields $b_0,b_1$;
the global fields in (NR.1) have not changed.

\medskip
\noindent\textbf{Exact remaining realization boundary.}
The complete source \texttt{SideB} contains more than coefficient
coverage and bounded Galois matrices. Its producer
\texttt{FormsStep.second\_\allowbreak step} also assumes the actual rational
row datum and integral Hecke lattice, $\mathrm{MtFam}$ and
$\mathrm{MtInput}$, the same level and tuple data, and the
automorphic inputs retained in NR22--NR24. A member
$\mathrm{MtReal}(\mu,j)$ is an injective map from the specified
space $D.\mathrm{Mt}(\mu,j)$ onto an entire adelic type space
$\mathrm{TypeSp}(\mathrm{wt},J,\chi)$, with an actual open level;
an arbitrary abstract vector-space isomorphism does not give
this assertion. The further $\mathrm{MtAd}$ ties that map
pointwise to the fixed $Y.\mathrm{IL}.jMt$, the same row frames,
the rational group and weight action, and the exact finite
matrix coordinates. None of these identities follows from
the coefficient constructions above.

Likewise $\mathrm{IsTripleWith}$ requires the actual occurring
signature and cuspidal parameter with its infinity and level
labels, a Galois representation for that parameter, the exact
scalar-twist equation for the already installed $\rho_S$, its
Hecke eigenvalue equations, and the specified residual
characteristic polynomials and Teichmuller reduction. NR18
supplies its \texttt{Induces} clause and NR22--NR24 supply the
stated conditional representation primitives, not these
remaining simultaneous identifications. $\mathrm{IsTripleFor}$
also requires $\mathrm{IsChars}$ for both readings, including
the already fixed tame-generator labels, and spectral coverage
of every relevant eigenspace. A compatible installation into
$\mathrm{NTw}$ still needs its actual form part $\mathrm{TwF}$
and toric part $\mathrm{TwL}$; nonempty local indices further
need $\mathrm{LocsCover}$ and the local form/Galois ties. The
bare stage with empty local indices does not furnish them.
Thus the nonempty all-prime family proved here closes concrete
arithmetic suppliers while leaving these exact same-object
geometric and automorphic interfaces explicit.


\section{Constructing the canonical finite and real Weil actions}
\label{sec:canonical-weil}

We now construct the full representation retained in WA18.  Throughout,
the arithmetic datum is the literal $D:\texttt{D7}$, $E/F$ is its CM
extension, $\psi$ is its standard additive character of sign $-1$, and
$S,T$ are the specified Hermitian spaces.  Their ranks may be zero and
need not be even.  The given full lattices $\Lambda_S,\Lambda_T$ are
self-dual.  We retain the original infinity embeddings $\theta$, the
host ideal $\mathfrak A=\texttt{host.AA}$, and the original half
symplectic pairing $b=\texttt{symplOH}$; put $B=2b=\texttt{symplO}$.
Thus, on $X=\operatorname{Hom}_E(S,T)(\mathbb A_F)$,
\[
 H_xf(y)=\psi(b(y,x))f(y+x),\qquad
 r(g,h)x=hxg^{-1}.                                      \tag{CW.1}
\]
The theta space $\Theta$ consists of the source's continuous,
real-smooth, finite-level functions with its rational quasi-periodicity.
Statements about operators below concern this space; extending an
operator to every function on $X$ does not assert analytic properties
on that larger algebraic space.

We use ordinary finite-dimensional Fourier inversion, Haar integration,
the spectral theorem for Hermitian matrices, and elementary number-field
lattice theory.  Every infinite integral introduced below has an
explicit convergence argument.  CA3--CA4 supply the already reconstructed
finite norm Gauss sum and signed lattice-intertwiner composition.  Their
generic statements apply here without CA5's separate even-rank condition.
No full Weil representation is an input to CW1--CW22.

\paragraph{CW1: the actual finite module, integral half, and self-dual lattice.}
Let $X_f$ be the finite part of $X$, and let
\[
 A=\widehat{\mathfrak A\operatorname{Hom}_{\mathcal O_E}
                  (\Lambda_S,\Lambda_T)}\subset X_f .     \tag{CW.2}
\]
Here completion and the embedding into $X$ are the original ones; the
infinite component is zero.  WA1's integral trace-dual calculation gives
$A=A^{\perp}$ for $\gamma(x,y)=\psi(B(x,y))$.  More explicitly, local
dual bases for the two Hermitian lattices make the matrix units dual
for the trace pairing.  The factor
$\xi_0\mathfrak A\overline{\mathfrak A}$ is the host's codifferent,
so testing all integral multiples of each matrix unit gives precisely
(CW.2).  This argument works with projective lattices and does not
require a global orthonormal integral basis.

Choose $w\in\mathcal O_E$ with $w+\bar w=1$.  Such a choice is possible
because the relative trace ideal is $\mathcal O_F$: away from $2$ the
trace of $1$ is a unit locally, and at the primes over $2$ the host's
unramified quadratic algebra has surjective integral trace.  Localizing
the trace ideal proves the global assertion.  Define
\[
 W_x=\psi(-b(wx,x))H_x\quad(x\in X_f).
 \qquad W_xW_y=\gamma(wx,y)W_{x+y}.                     \tag{CW.3}
\]
To obtain the identity, expand $b(w(x+y),x+y)$, use
$b(ax,y)=b(x,\bar a y)$, skew-symmetry, and $w+\bar w=1$.
It follows that $W_xW_y=\gamma(x,y)W_yW_x$.
All these scalars have absolute value one.  On a self-dual
$\mathcal O_E$-submodule the cocycle in (CW.3) is one.

Every $r(p)A$ is self-dual and commensurable with $A$.  Indeed $r(p)$
preserves the form; its matrix and inverse are integral outside a finite
set, and at each remaining place powers of a local uniformizer bound
both images.  One rational integer clears these finitely many bounds.
The same argument shows that a sufficiently large integer multiple of
$A$ lies in any given finite open neighborhood of zero.  All ensuing
quotients of two such lattices by their intersection are finite.

\paragraph{CW2: the signed normalization is available for these lattices.}
Use the multiplicative function $\lambda:\mathbb N\to\mathbb C$
constructed in CA4.  It is nonzero at positive integers.  For each
residue field of size $q_P$, $\lambda(q_P)^2=q_P$; when conjugation
fixes $P$, its residue action is nontrivial and
$\lambda(q_P)=-\sqrt{q_P}$.  This last sign is essential at odd as
well as dyadic residue characteristic.  It is supplied by CA3's norm
Gauss sum, not chosen independently at successive lattice moves.

For clarity, the host permits a single such multiplicative function.
Its abelian Galois extension $E/\mathbb Q$ has a common residue degree
$f_p$ above a rational prime $p$, and conjugation fixes either all or
none of the primes above $p$.  Set
$\lambda(p)=\sqrt p\exp(\pi i/f_p)$ in the former case and
$\lambda(p)=\sqrt p$ in the latter, and extend by prime factorization.
The unramified CM condition makes the residue involution nontrivial
at the fixed primes.  Hence these choices have exactly the residue
values just stated.  If no rational inert prime divides $n>0$, then
$\lambda(n)=\sqrt n$.

For commensurable self-dual lattices write
\[
 n(A,C)=[C:A\cap C],\qquad
 J_{A,C}e=\lambda(n(A,C))^{-1}
       \sum_{x\in C/(A\cap C)}W_xe\quad(e\in\Theta^A).
                                                               \tag{CW.4}
\]
The sum is representative-independent because $(A\cap C)$ fixes $e$
and $C$ is isotropic.  It belongs to $\Theta^C$.  CA4 proves
$J_{C,L}J_{A,C}=J_{A,L}$ and $J_{A,A}=1$ by finite character
orthogonality and the residue norm Gauss sum.  In particular these are
isomorphisms.  The same proof, or inverse composition on a nonzero
fixed vector, gives $\lambda(n(A,C))^2=n(A,C)$ whenever that fixed
space is nonzero.  Consequently its absolute-value square is
$n(A,C)$.  This signed normalization is the one used throughout CW;
it is not replaced by positive square roots at nonsplit primes.

\paragraph{CW3: finite character projectors span the actual theta space.}
Every $f\in\Theta$ is fixed by $W_x$ for $x$ in a sufficiently small
open subgroup of $A$.  The raw finite-level condition first gives
fixing under $H_x$; intersect with the kernel of the phase in (CW.3).
By CW1 it contains $NA$ for some integer $N>0$.
Put $L=N^{-1}A$ and
$U=A\cap L^\perp=NA$.  The pairing $\gamma$ induces a perfect pairing
of the finite groups $A/U$ and $L/(A\cap L)$.  Injectivity on either
side follows from $A=A^\perp$; finite character orthogonality then
gives equality of their cardinalities and surjectivity to the dual.

For $t\in L$ form
\[
 e_t=\sum_{a\in A/U}W_aW_tf.
                                                               \tag{CW.5}
\]
The vector $W_tf$ is $U$-fixed by the commutator formula, so (CW.5)
is well-defined and fixed by $A$.  Commuting $W_{-t}$ through this
sum expresses a nonzero phase multiple of $W_{-t}e_t$ as
$\sum_{a\in A/U}\gamma(a,t)W_af$.
Sum this expression over the dual group $L/(A\cap L)$ with the
opposite sign of $t$.  Orthogonality kills every term except $a=0$,
and the remaining term is $|A/U|f$.  Division by this positive integer
proves
\[
             \Theta=\operatorname{span}_{\mathbb C}
                   \{W_xe:x\in X_f, e\in\Theta^A\}.    \tag{CW.6}
\]
Only a finite sum is used for each $f$.  The construction neither
asserts that $\Theta^A$ is one-dimensional nor requires a topological
closure of the span.

\paragraph{CW4: defining a genuine representation from the fixed vectors.}
For $p=(g,h)$ define an operator on the generators in (CW.6) by
\[
       \Omega_f(p)(W_xe)=W_{r(p)x}J_{A,r(p)A}e .          \tag{CW.7}
\]
Here is the required well-definedness argument.  An $A$-fixed vector
$e$ gives an $A$-eigenvector $W_xe$ with eigencharacter
$a\mapsto\gamma(a,x)$.  The characters of $x,x'$ are equal exactly
when $x-x'\in A$.  A finite relation between such vectors therefore
splits into its separate $X_f/A$ cosets: for finitely many distinct
characters, polynomials in the commuting $W_a$ isolate each
eigenspace.  Within one coset, move every representative to the same
$x$ using (CW.3), leaving a relation among $A$-fixed vectors.
$J_{A,r(p)A}$ is linear and preserves that relation.  The transformed
representatives have exactly the same scalar correction, since
$r(p)$ preserves $\gamma$ and commutes with $w$.  Thus (CW.7)
annihilates every relation and defines a linear operator on $\Theta$.

It intertwines $W_x$ with $W_{r(p)x}$ by direct multiplication.
Reindexing the finite sum proves
\[
 \Omega_f(p)J_{A,C}e
       =J_{r(p)A,r(p)C}\Omega_f(p)e .                  \tag{CW.8}
\]
Applying (CW.8) with $C=r(q)A$ and using CW2 proves
$\Omega_f(p)\Omega_f(q)e=J_{A,r(pq)A}e$ on $\Theta^A$,
hence the group law on all generators.  Also
$\Omega_f(1)=1$ and the inverse is $\Omega_f(p^{-1})$.
This proves the literal lattice-representation construction; it is
not an appeal to a projective representation whose multiplier remains
to be removed.

\paragraph{CW5: continuity, arbitrary commutants, and the Haar norm.}
If a linear map $L:\Theta\to\Theta$ commutes with all finite $W_x$,
then it maps $\Theta^A$ into itself and commutes with the finite sum
(CW.4).  Equations (CW.6)--(CW.7) therefore give
$L\Omega_f(p)=\Omega_f(p)L$.  No continuity assumption on $L$ is
needed for this algebraic assertion.  Real Heisenberg translations
commute with finite ones, since their cross pairings vanish; hence
they are such commutants.  The phase in (CW.3) is invariant under
$r(p)$, so covariance also holds for the untwisted $H_x$.

For each fixed $f$, $\Omega_f(p)f=f$ on a neighborhood of $p=1$.
For a generator $W_xe$, require $r(p)A=A$ and
$r(p)x-x\in NA$, with $N$ chosen to annihilate
$a\mapsto\gamma(wx,a)$ on $NA$.  These are open conditions:
use finite generators for the completed lattice and the finitely
many nonintegral local matrix entries.  Then $J_{A,r(p)A}e=e$ and
(CW.3) gives $W_{r(p)x}e=W_xe$.  Intersect finitely many such
neighborhoods for a presentation of $f$.  The group law transports
this assertion to every $p$, proving local constancy of this
unadjusted finite action on each vector.

Fix any additive Haar measure $\mu$ and any rational-lattice
fundamental domain $\mathcal F$.  The theta products
$f\overline{f'}$ are periodic and bounded on the compact quotient,
so their integrals define a finite Heisenberg-invariant Hermitian
form.  Distinct translates $W_xe,W_ye'$ of fixed vectors are
orthogonal unless $x-y\in A$: applying an $A$-element with unequal
eigencharacters proves the assertion.  Consequently (CW.4) multiplies
the pairing of two fixed vectors by
$n(A,C)/|\lambda(n(A,C))|^2=1$.  Covariance, (CW.6), and the same
orthogonality prove the equality for arbitrary pairs.  Thus
\[
 \int_{\mathcal F}|\Omega_f(p)f|^2\,d\mu
                     =\int_{\mathcal F}|f|^2\,d\mu .   \tag{CW.9}
\]
This retains the given Haar measure; no unrecorded volume rescaling
has been made.  A zero fixed space causes no division problem, since
all its assertions are then zero and (CW.6) handles the whole space.

\paragraph{CW6: the exact positive split scalar and nonsplit fixing.}
Let $e$ be a split-place idempotent, $P x=ex$, $Qx=\bar e x$, and
suppose $r(p)$ is the identity off $P+Q$.  Both projections preserve
$A$, satisfy $PQ=QP=0$, and are adjoint for $\gamma$.
Let $\ell$ be any linear functional invariant under $H_{Qx}$ on
theta functions.  The phase in (CW.3) is one on this half-space,
so $\ell$ is also invariant under $W_{Qx}$.  With representatives
$b_i$ for $r(p)A/(A\cap r(p)A)$ put
\[
 n_e=\#\{i:Pb_i\in A\}>0.
 \qquad
 \ell(\Omega_f(p)f)
       ={n_e\over\lambda(n(A,r(p)A))}\ell(f).          \tag{CW.10}
\]
To prove it first take $f$ fixed by $A$.  Remove the $Qb_i$ factor
using $\ell$-invariance, and the component off $P+Q$ using
$A$-fixing.  The remaining term is $\ell(W_{Pb_i}f)$.
It equals $\ell(f)$ when $Pb_i\in A$.  Otherwise self-duality
provides $a\in A$ with $\gamma(Pb_i,a)\ne1$; its $Qa$ component
forces that term to vanish.  The zero quotient class gives $n_e>0$.
For $W_xf$, decompose $x$ into $Px,Qx,(1-P-Q)x$ and apply the
same argument to $\ell W_{(1-P-Q)x}$; covariance gives the identical
constant.  CW3 completes the proof.

The index in (CW.10) is supported at this split pair.  A rational
prime inert in the host cannot occur in that support: conjugation
has the same splitting behavior at every prime above it, because
the host extension is abelian.  CW2 therefore makes its denominator
the positive square root of the index.  Thus the multiplier in
(CW.10) is strictly positive, independent of both $f$ and $\ell$.

At a nonsplit finite place $v$, let $P_v$ be its idempotent.
The host makes every dyadic place split, so $P_v=2b_v$ for an
integral finite adele $b_v$.  If the pair $p$ is supported at $v$
and preserves $P_vA$, it preserves all of $A$, since it is the
identity on $(1-P_v)A$.  A vector fixed by the raw $H_x$ for
$x\in P_vA$ is fixed by $W_x$ there: write $x=2b_va$ and use
isotropy of $A$ to obtain $\psi(-b(wx,x))=1$.  The elementary
local-projector argument now gives $\Omega_f(p)f=f$ on these
vectors.  One proof is to apply the finite projectors of CW3 to
such an $f$: components with nontrivial character on $P_vA$ vanish;
in each remaining component the $P_v$ translation lies in $A$ by
self-duality and can be absorbed into the fixed vector.  The
remaining translation is outside $v$, where $r(p)$ is the identity.
This is the literal local-vacuum condition, with the same completed
Hom lattice, rather than merely invariance under a smaller subgroup.

\paragraph{CW7: the crossed determinant character constructs IsFinWeil.}
Let $\chi_S,\chi_T$ be the given unitary Hecke characters with
the exact restrictions of ranks $\dim_E S,\dim_E T$.  Choose one
member of every split pair and let $\epsilon$ be its finite adele
idempotent, zero at all nonsplit and infinite places.  The choices
form an adele because all entries are integral zeros or ones.  For
a unit $u$ write
\[
 \operatorname{cut}_\epsilon(u)=\epsilon u+1-\epsilon,
 \qquad
 z(g,h)=\chi_T(\operatorname{cut}_\epsilon\det g)^{-1}
         \chi_S(\operatorname{cut}_\epsilon\det h).
                                                               \tag{CW.11}
\]
The cut is a continuous group homomorphism, its inverse being the
same cut of $u^{-1}$.  The rank restriction implies that each
$\chi$ kills $u\bar u$: its restriction to $F$ is a power of the
quadratic norm character, which is one on norms.  Thus selecting
the conjugate member of a supported split pair gives the same
character value, because a unitary determinant satisfies
$\det g\,\overline{\det g}=1$.  At a real or nonsplit-supported
element the cut determinant is one.

Set $\omega_f(p)=z(p)\Omega_f(p)$ on $\Theta$.  Equations
(CW.7)--(CW.11) prove its group law, pointwise continuity, norm
preservation, real triviality, covariance, and nonsplit vacuum
fixing.  At a split pair (CW.10) becomes the exact positive
multiple of
$\chi_T(x)^{-1}\chi_S(y)=\texttt{splitFacO}(\chi_S,\chi_T,x,y)$,
where $x,y$ are the original determinant lifts.  The characters are
deliberately crossed.

To obtain the literal endomorphism of the whole function space,
choose an algebraic complement $\mathcal C$ to $\Theta$ and act as
the identity on $\mathcal C$.  Equivalently, for a projection
$\pi$ onto $\Theta$, the extension is
$1+\iota(\omega_f(p)-1)\pi$.  Its multiplication law follows from
$\pi\iota=1$.  All required assertions are on $\Theta$ and are
unchanged.  This constructs every clause of \texttt{IsFinWeil}.

\paragraph{CW8: Gaussian projectors and their actual smooth domain.}
Let $V$ be a finite-dimensional Euclidean real space with an
orthogonal complex structure $J$, and let $\rho$ be a strongly
continuous isometric Heisenberg action on a complex Banach space
$\mathcal B$:
\[
 \rho(u)\rho(v)=e^{\pi i\langle u,Jv\rangle}\rho(u+v).
 \quad G(u)=e^{-\pi\|u\|^2/2},\quad
 P f=\int_VG(u)\rho(u)f\,du .                           \tag{CW.12}
\]
Lebesgue measure is the Euclidean normalization.  The integral is
absolutely convergent since $\|\rho(u)f\|=\|f\|$.
The one-variable Gaussian Fourier formula, applied in orthogonal
coordinates, gives $P^2=P$ and
$P\rho(v)P=G(v)P$.  For example the coefficient of $\rho(x)$
in $P^2$ is
$\int G(u)G(x-u)e^{\pi i\langle u,Jx\rangle}\,du=G(x)$;
completion of the square gives the last equality with precisely
the normalization in (CW.12).

Put $\mathcal B^\infty=\{f:v\mapsto\rho(v)f\text{ is }C^\infty
\text{ in Banach norm}\}$.  For $f$ in this space,
$F_f(u)=G(u)\rho(u)f$ is a Banach-valued Schwartz function.
Indeed
\[
 F_f(t+v)=G(v)\rho(v)
     e^{-\pi\langle v,t\rangle-\pi i\langle v,Jt\rangle}F_f(t).
                                                               \tag{CW.13}
\]
Differentiating at $t=0$ bounds each derivative at $v$ by a
polynomial in $\|v\|$ times $G(v)$ and finitely many derivative
norms at zero.  Every additional polynomial factor is bounded
against the Gaussian.  Fourier inversion is therefore legitimate
in Banach norm.  The Heisenberg law gives
\[
 P\rho(-v)f=\rho(-v)\widehat F_f(Jv),\qquad
       f=\int_V\rho(v)P\rho(-v)f\,dv .                 \tag{CW.14}
\]
In particular all polynomial moments of
$\|P\rho(-v)f\|$ are integrable.

Conversely these moment bounds imply $f\in\mathcal B^\infty$:
(CW.14)'s first identity makes all polynomial moments of
$\widehat F_f$ integrable, so differentiation under Fourier
inversion makes $F_f$ smooth.  Divide by the nonvanishing smooth
Gaussian.  This equivalence is the source's \texttt{Sm}/\texttt{Good}
criterion.  It supplies both the convergence and the smooth domain
for the operators in the next paragraphs.
\paragraph{CW9: evaluating the two Gaussian scalars and their phases.}
Choose complex Euclidean coordinates in which $Jz=is z$, where
$s=s^*=s^{-1}$.  For $g^*sg=s$ put
\[
 p={1+s\over2},\quad q=1-p,\quad
 c(g)=\det(pg+q),\quad
 L(g)=\int G(u)G(gu)\,du .                             \tag{CW.15}
\]
For two such matrices set
\[
 M(g_1,g_2)=\int\!\int G(u)G(v)
       e^{\pi i\langle g_1u,Jg_2v\rangle}G(g_1u+g_2v)\,du\,dv.
                                                               \tag{CW.16}
\]
The exact evaluations are
\[
 \begin{aligned}
 c(g^{-1})&=\overline{c(g)},\qquad
 L(g)&=[c(g)c(g^{-1})]^{-1},\\
 M(g_1,g_2)&=[c(g_1^{-1})c(g_2)c(g_2^{-1}g_1)]^{-1}.
 \end{aligned}                                                   \tag{CW.17}
\]
Here $c(g)\ne0$.  These formulas determine a phase, as well as
an absolute value.

We give the matrix calculation behind them.  The exponents in
(CW.15)--(CW.16) are $-\pi z^*Az$ and $-\pi(z,z')^*N(z,z')$,
where
\[
 A={1+g^*g\over2},\qquad
 N=\begin{pmatrix}
 (1+g_1^*g_1)/2&g_1^*pg_2\\
 g_2^*qg_1&(1+g_2^*g_2)/2
 \end{pmatrix}.                                                \tag{CW.18}
\]
$A$ is positive definite and
$N+N^*=1+[g_1\ g_2]^*[g_1\ g_2]$ is positive definite.
For any complex matrix with positive Hermitian part,
\[
              \int_{\mathbb C^n}e^{-\pi z^*Az}\,dz=(\det A)^{-1}.
                                                               \tag{CW.19}
\]
For positive Hermitian $A$, diagonalize and use the elementary
complex one-dimensional Gaussian integral.  Both sides are
holomorphic in the matrix entries on the connected positive-real-part
domain; local Gaussian domination permits differentiation under the
integral.  Agreement on its open positive Hermitian real slice
implies agreement in a complex neighborhood by the power series
identity theorem, and then throughout that domain.  This proves
(CW.19), including absolute integrability, without choosing a square
root of a determinant.

For completeness the determinant reductions in (CW.18) use only
$p^2=p$, $pq=0$, and $g^*=sg^{-1}s$.  If $h=g^{-1}$, then
\[
 A=sh(pgp-qgq),\qquad
 \det(pgp-qgq)=\det(s)c(g)c(h)\det(g).
\]
Thus $\det A=c(g)c(h)$; taking adjoints gives
$c(h)=\overline{c(g)}$.  For the second matrix factor out
$\operatorname{diag}(sg_1^{-1},sg_2^{-1})$.  The remaining matrix is
\[
 K=\begin{pmatrix}pg_1p-qg_1q&pg_2\\-qg_1&pg_2p-qg_2q\end{pmatrix}.
\]
Conjugate by $\begin{pmatrix}p&q\\q&p\end{pmatrix}$, whose square
is one.  This gives an upper triangular block matrix with diagonal
blocks $sg_2(g_2^{-1}g_1p+q)$ and $pg_2p-qg_1q$.
Taking their determinants gives
$\det N=c(g_1^{-1})c(g_2)c(g_2^{-1}g_1)$.
Equation (CW.19) now gives all of (CW.17).

\paragraph{CW10: an actual oscillator action, rather than a projective one.}
For a group acting through the matrices in CW9, let $P_g$ be
(CW.12)'s projector for $u\mapsto\rho(gu)$ and define on
$\mathcal B^\infty$
\[
 T_g=c(g)P_gP,\qquad
 \Omega_\infty(g)f=\int_V\rho(gv)T_g\rho(-v)f\,dv .      \tag{CW.20}
\]
This integral converges absolutely because
$\|T_g\rho(-v)f\|\le |c(g)|\|P_g\|\|P\rho(-v)f\|$.
There is an exact coherent-coefficient identity, stronger than a
kernel estimate:
\[
 P_g\rho(-gv)\Omega_\infty(g)f=T_g\rho(-v)f.
\]
Indeed insert (CW.20), use $P_gT_g=T_g$, and apply
$P_g\rho(g(w-v))P_g=G(w-v)P_g$ with the Heisenberg phase.
The resulting integral is also obtained by applying
$T_gP\rho(-v)$ to the resolution (CW.14), since $T_gP=T_g$.
Both operations commute with the integral by the preceding absolute
bound.  For every integer $k\ge0$ this gives the explicit estimate
\[
 \int\|v\|^k\|P_g\rho(-gv)\Omega_\infty(g)f\|\,dv
 \le \|T_g\|\int\|v\|^k\|P\rho(-v)f\|\,dv<\infty.
\]
Thus the result satisfies CW8's moment criterion for $\rho\circ g$.
It is smooth for that action, and composition with $g^{-1}$ proves
smoothness for $\rho$.  All later linearity and covariance operations
are therefore on the genuine smooth domain.

Here are the identities that verify the group law.  Completing
the square in the absolutely convergent double and triple
Gaussian integrals gives
\[
 PP_gP=L(g)P,\qquad
 PP_{g_1}P_{g_2}P=M(g_1,g_2)P.                         \tag{CW.21}
\]
Also (CW.14), followed by translation of the integration variable,
gives
\[
 P_g\Omega_\infty(g)f=T_gf,\quad
 \Omega_\infty(g)\rho(v)f=\rho(gv)\Omega_\infty(g)f,
 \quad \Omega_\infty(g)P_h=P_{gh}\Omega_\infty(g).       \tag{CW.22}
\]
CW9 says exactly
\[
 |c(g)|^2L(g)=1,\qquad
 c(g)c(h)M(gh,g)=c(gh)L(gh).                           \tag{CW.23}
\]
Applying $P$ and using (CW.21)--(CW.23) shows that the difference
between $c(g)c(h)P_{gh}P_gP$ and $c(gh)P_{gh}P$ is killed by $P$.
It is in the range of $P_{gh}$; applying
$P_{gh}PP_{gh}=L((gh)^{-1})P_{gh}$ shows that it is zero.
Use this equality in (CW.22), and test after every translated
$P_{gh}$.  Formula (CW.14) says these tests separate smooth vectors.
The result is
\[
 \Omega_\infty(g)\Omega_\infty(h)=\Omega_\infty(gh),
                   \qquad\Omega_\infty(1)=1.           \tag{CW.24}
\]

Every bounded linear map commuting with $\rho$ commutes with
$P_g,T_g$ and (CW.20).  Every bounded linear map into a Hilbert space
whose norm is Heisenberg invariant is preserved by (CW.20).
For the latter assertion $P,P_g$ are self-adjoint for the resulting
sesquilinear form; (CW.21) and (CW.23) show $T_g$ is isometric on
the range of $P$.  The resolution (CW.14) gives
\[
 \langle f,f'\rangle
   =\int\langle P\rho(-v)f,P\rho(-v)f'\rangle\,dv.
\]
Use (CW.22) and the analogous resolution for $\rho\circ g$ to
deduce equality of the inner products after $\Omega_\infty(g)$.
All double integrals here are justified by CW8's moment bounds
and Gaussian domination.  Finally, (CW.20) depends continuously
on $g$ for each fixed smooth $f$: on a compact neighborhood of
$g_0$, the projector norms are uniformly bounded, and the same
integrable function $C\|P\rho(-v)f\|$ dominates the integrands.

\paragraph{CW11: the Cartan-vacuum determinant and change of coordinates.}
The vacuum normalization in CW10 can be checked directly.  If
\[
 d\rho(z)f+i\,d\rho(Jz)f=0\quad\text{for every }z,
                                                               \tag{CW.25}
\]
then $Pf=f$.  Indeed, for each $u$, the function
$\zeta\mapsto e^{\pi|\zeta|^2\|u\|^2/2}
 \rho((\Re\zeta)u+(\Im\zeta)Ju)f$
is holomorphic: differentiating and using (CW.25) is the
Cauchy--Riemann equation.  Its circle mean is its value $f$ at
zero.  Therefore the circle mean of the unweighted orbit on
$|\zeta|=1$ is $G(u)f$.  Average the integral (CW.12) over those
rotations; Fubini is bounded by $G(u)\|f\|$.  It becomes
$\bigl(\int e^{-\pi\|u\|^2}\,du\bigr)f=f$.

If $g$ is Euclidean unitary, $P_g=P$.  Equations (CW.20)--(CW.22)
then give $\Omega_\infty(g)f=c(g)f$ on this vacuum space.
For a different Cartan involution $s'$ with $Hs'$ positive,
choose an $H$-unitary change of coordinates taking $s$ to $s'$.
This exists by choosing orthonormal bases on the positive and
negative eigenspaces for the two Cartan decompositions.  Conjugate
the preceding assertion and use (CW.24).  Since determinant is
unchanged by conjugation, an element commuting with $s'$ acts
on its vacuum space by
\[
             d_P(s',g)=\det\left({1+s'\over2}g
                                      +{1-s'\over2}\right).
                                                               \tag{CW.26}
\]
It is this determinant itself, not its reciprocal.

For an arbitrary nonsingular Hermitian matrix $H$, choose an
invertible complex matrix reducing it by congruence to a signature
involution.  Pull $\rho$, the real coordinates, the group matrices
and the vacuum equations through that same linear equivalence.
CW8--CW10 are preserved, and (CW.26) is unchanged because its
matrix is conjugated.  Thus no definiteness of $S$ or $T$ is
needed for the real construction.

\paragraph{CW12: theta smoothness really is Banach smoothness here.}
Let $\mathcal B$ be the closed subspace of bounded continuous
functions on $X$ satisfying the original rational quasi-periodicity,
with supremum norm.  Compactness of $X/X(E)$ bounds every continuous
theta function; the quasi-periodic factors have modulus one.
The action $H_x$ is isometric and strongly continuous on this space.
For the latter assertion use uniform continuity on a compact set
meeting all rational cosets and the continuous phase in (CW.1).
Point evaluations are bounded and norming.

Use the actual inverse infinity chart of $\theta$ to write the
real Heisenberg action as $\rho(v)$.  A theta function has pointwise
directional derivatives, and each derivative is again a theta
function: differentiate the quasi-periodic formula and finite-level
identities, and then the defining real smoothness condition.
Denote the derivative at zero by $D_fz$.  Evaluations show this
is real linear in $z$; its finite-dimensional domain makes it
continuous.  The Heisenberg law gives the candidate Banach derivative
\[
 d(\rho(\cdot)f)_v(z)=\rho(v)D_fz
                         -\pi i\langle v,Jz\rangle\rho(v)f.
                                                               \tag{CW.27}
\]
It is continuous in $v$ in operator norm: this follows from strong
continuity on finitely many basis derivatives $D_fe_i$ and $f$.
Apply the scalar mean-value inequality to every evaluation of the
remainder, uniformly bounded by the modulus of continuity of
(CW.27).  Taking the supremum proves the Fr\'echet derivative in
Banach norm.  Iterating on the derivatives proves
$f\in\mathcal B^\infty$.  This is the missing uniform argument;
pointwise differentiability by itself would not suffice.

Conversely a Banach-smooth vector in $\mathcal B$ that is fixed at
some finite open level is a literal theta function.  Evaluate its
smooth real orbit and cancel the explicit smooth phase in (CW.1)
to get the source's real-line smoothness.  Quasi-periodicity and
continuity were already part of $\mathcal B$.  Thus the operators
of CW10 may be applied to precisely the functions needed here.

\paragraph{CW13: constructing IsInfWeil on the original theta functions.}
Vectorizing the actual real Hom matrices gives the complex index
set $((\operatorname{Fin}\dim T\times\operatorname{Fin}\dim S)
\times\mathrm{Tau})$.  In these coordinates the matrix of the
unitary pair is
\[
 m(g,h)=\operatorname{diag}_\tau
   \bigl(h_\tau\otimes(g_\tau^{-1})^{\mathsf t}\bigr).
                                                               \tag{CW.28}
\]
The real alternating form is the imaginary part of the original
nondegenerate Hermitian Hom form, including its prescribed host
scale and the sign of $\psi$.  Expanding the trace pairing proves
$m^*Hm=H$.  Apply CW10--CW11 to these matrices and to CW12's
Banach space.

Finite Heisenberg operators commute with the real action: their
cross pairings are zero, so they commute with every real $\rho(v)$
and hence with the Gaussian operator of CW10.  A finite open level
fixing $f$ therefore fixes its image.  CW12 makes that image a theta
function.  Restrict the resulting linear action to $\Theta$, then
extend by the identity on an algebraic complement, exactly as in
CW7.  Call the result $\omega_\infty$.

It is trivial on the finite subgroup, since (CW.28) is then the
identity.  Covariance for real shifts follows from (CW.22), and
for finite shifts from the commutant assertion.  Splitting an
arbitrary shift into its two orthogonal parts proves full
Heisenberg covariance for a real group element.  Pointwise
continuity follows from CW10 followed by evaluation.  For every
Haar measure and rational fundamental domain, the map
$\mathcal B\to L^2(\mathcal F,\mu)$ is bounded because
$\mu(\mathcal F)<\infty$; its norm is Heisenberg invariant by
translation on the compact quotient.  CW10 therefore proves
the exact source norm equality.

Finally the Cartan involution on the Hom coordinates is
$s=\operatorname{diag}_\tau(s_{T,\tau}\otimes
s_{S,\tau}^{\mathsf t})$.  The source's Fock equations are
exactly (CW.25), as evaluation of (CW.27) at zero shows.
For a pair commuting with these Cartans, (CW.26) equals
\[
 \prod_\tau\det\left[
 {1+s_{T,\tau}\otimes s_{S,\tau}^{\mathsf t}\over2}
       (h_\tau\otimes(g_\tau^{-1})^{\mathsf t})
 +{1-s_{T,\tau}\otimes s_{S,\tau}^{\mathsf t}\over2}\right].
                                                               \tag{CW.29}
\]
This is literally \texttt{VacCharPlus}.  All clauses of
\texttt{IsInfWeil}, on the original functions and original charts,
have now been constructed.

\paragraph{CW14: commutation and continuity of the combined action.}
For our actual finite action,
$\omega_f(p)=z(p)\Omega_f(p)$, CW5 immediately implies
\[
            \omega_f(p)\omega_\infty(q)f
             =\omega_\infty(q)\omega_f(p)f\quad(f\in\Theta),
                                                               \tag{CW.30}
\]
because the real action commutes with every finite $W_x$.
Hence $\omega(p)=\omega_f(p)\omega_\infty(p)$ is a genuine
representation on $\Theta$.  Extend this product representation
itself to all functions by one complement.  The arbitrary
extensions of the two factors need not commute outside $\Theta$.

There is a useful direct proof of the required pointwise continuity
which avoids an unjustified joint continuity assertion.  Near
$p_0$, write $p=p_0q$.  For this fixed $f$, CW5 gives
$\Omega_f(q)f=f$ when $q$ is sufficiently near $1$.  Using (CW.30)
and commuting only fixed operators gives
\[
 \omega(p)f=z(q)\,\omega_\infty(p_0q)\omega_f(p_0)f.
                                                               \tag{CW.31}
\]
The right side is pointwise continuous by CW7 and CW13.
Each factor preserves the exact Haar norm, so their product does
as well.  The real and finite subgroup restrictions are exactly
$\omega_\infty$ and $\omega_f$ respectively.  Splitting both a
group element and a Heisenberg shift into real and finite parts
proves full covariance under $r(p)$.  At this stage rational
automorphy remains to be proved; it is not a consequence of
(CW.30) or of the norm equality.

\paragraph{CW15: the scalar lemma with its essential norm hypothesis.}
Suppose $M$ preserves $\Theta$, commutes with every $H_x$, and
preserves its Haar $L^2$ norm.  Then on $\Theta$ it is multiplication
by one scalar of absolute value one.  Polarization first proves
preservation of the integrated pairings.  For $a\in X(E)$ the
quasi-periodicity gives
\[
                    H_af(x)=\psi(B(x,a))f(x).           \tag{CW.32}
\]
Consequently every Fourier coefficient of the continuous periodic
function $(Mf)\overline{Mf'}-f\overline{f'}$ is zero.
The characters in (CW.32) separate points of $X/X(E)$: the
annihilator of the rational lattice for the standard trace character
is itself.  Here is an elementary check of the underlying assertion.
For $\mathbb Q$, subtract a rational number to make every finite
component integral.  Pairing with the integers then makes the real
component an integer; subtract it.  Pairing successively with
$p^{-k}$ forces the remaining $p$-component into every $p^k\mathbb Z_p$,
hence to zero.  A rational basis and its trace-dual basis reduce a
number field, and then the Hom space, to this calculation.
Character orthogonality and the density of their
self-conjugate algebra in continuous functions on this compact
quotient show that the function is zero.  One can use uniform
approximation followed by integration against its complex conjugate;
this needs no assumption about the spectrum of $M$.

Apply this identity to $f$ and $H_{y-x}f'$ and cancel the nonzero
phase.  It follows that
$(Mf)(x)\overline{(Mf')(y)}=f(x)\overline{f'(y)}$ for all $x,y$.
If $f_0(y)\ne0$, this determines the same scalar for every $f,x$,
and the case $f=f'=f_0,x=y$ gives absolute value one.  If the theta
space is zero, take the scalar $1$.

For a rational unitary pair $p$ let
$R_pf(x)=f(r(p^{-1})x)$.  It preserves quasi-periodicity, finite
level and real smoothness.  Its rational linear change of variables
preserves the lattice quotient and Haar measure: the adelic
modulus of its rational determinant is one by the product formula.
It has the same covariance as $\omega(p)$.  Apply the scalar lemma
to $\omega(p)R_{p^{-1}}$.  Thus
\[
                     \omega(p)f=c_pR_pf,
                              \qquad |c_p|=1.          \tag{CW.33}
\]
The next calculations determine $c_p$; no assertion that it is
algebraic or equal to one has been used.

\paragraph{CW16: the finite Hermitian Gauss sum with its signed value.}
Let $M$ be a finite $\mathcal O_E$-module with a nondegenerate,
symmetric bicharacter $\beta$, satisfying
$\beta(ax,y)=\beta(x,\bar a y)$.  Put
$Q(x)=\beta(wx,x)$.  The integral half in CW1 gives
\[
 Q(x+y)=Q(x)Q(y)\beta(x,y),\qquad
       \sum_{x\in M}Q(x)=\lambda(|M|),\quad
                        \lambda(|M|)^2=|M|.           \tag{CW.34}
\]
Here the same $\lambda$ as in CW2 is used.

We prove the two sum assertions together by induction on $|M|$.
The zero module gives $1$.  Otherwise choose a simple submodule
$D\simeq\mathcal O_E/P$, of order $q$.
If $\beta|_{D\times D}=1$, then $Q|_D=1$.
Summing $Q(x+d)$ over $d\in D$ kills $x\notin D^\perp$.
On $D^\perp$ the function descends to the nondegenerate module
$D^\perp/D$, whose order is $|M|/q^2$ by perfect character duality.
Thus the sum is $q$ times its Gauss sum.  Induction and
$\lambda(q)^2=q$ give both assertions in (CW.34).

If the restriction is nontrivial, its radical is an
$\mathcal O_E$-submodule of the simple module $D$, hence zero.
Perfect duality gives $M=D\oplus D^\perp$.  Adjointness forces
$\bar P=P$: otherwise $P+\bar P=\mathcal O_E$ would make the
restriction zero.  The residue involution is nontrivial by the
unramified host condition.  On $D$, the quadratic function is
$a\mapsto\chi(w a\bar a)$ for a nontrivial additive character
with the required conjugation invariance.  CA3's explicit norm
sum is $-\sqrt q=\lambda(q)$, also in characteristic two.
The sums on the two orthogonal factors multiply, proving (CW.34).
This also proves the version for a finite quotient $B/C$ when
$C$ is exactly the radical of $\beta|_B$: the formula descends,
and a transversal sums it once.  Thus no semisimplicity of the
whole finite module, and no residue-characteristic restriction,
is hidden in the argument.

\paragraph{CW17: paired rational lattice indices and the actual finite sum.}
Let $V=\operatorname{Hom}_E(S,T)$ and
$L=\{x\in V:x_f\in A\}$, the source's \texttt{latQ}.
It is a full finitely generated $\mathbb Z$-lattice.  Its finite
closure is $A$: choosing local bases, finite approximation with
integral denominators gives density, and a common integer bounds
the finitely many nonstandard local lattices.  Consequently
$L$ is self-dual for
$\gamma_Q(x,y)=\gamma(x_f,y_f)$.  The relative index of two such
rational lattices equals the relative index of their completions,
since a finite lattice quotient is unchanged by finite completion.

Fix a rational unitary pair $p=(g,h)$.  Write $a$ for the rational
matrix of $g$ and $b_0$ for that of $h^{-1}$.  Throughout this
calculation
\[
 u(x)=b_0xa=r(p^{-1})x,\quad
 U=b_0\otimes a^{\mathsf t},\quad
 m=\dim_E V,\quad d=\det(U-1),\quad\Delta=\det U.
                                                               \tag{CW.35}
\]
Assume $d\ne0$; $\Delta\ne0$ follows from unitarity.
Put
\[
 B_1=(u-1)^{-1}L,\quad B_2=u^{-1}L,\quad
 C=L\cap B_1=L\cap B_2,\quad n=[B_1:C],\quad i=[B_2:C].
                                                               \tag{CW.36}
\]
The equality of intersections is simply
$ux-x\in L\iff ux\in L$ when $x\in L$.
Moreover $i=[r(p)A:A\cap r(p)A]$, with the orientation used in
(CW.4).  Choose an integer $N>0$ with $NuL\subset L$.
Then $NL\subset C$; let $c_0=[C:NL]>0$.
Applying $u-1$ and $u$ to the corresponding lattice quotients
and taking determinants over $\mathbb Q$ gives
\[
 c_0n=|N_{E/\mathbb Q}(N^m d)|,\qquad
 c_0i=|N_{E/\mathbb Q}(N^m\Delta)|.                   \tag{CW.37}
\]
Both equalities have the same $N,c_0,m$.  This common denominator
is needed for the later phase cancellation.

For $x,y\in B_1$ define
\[
 \beta(y,x)=\gamma_Q(\bar w(uy-y),x)
             \gamma_Q(w(uy-y),ux)
          =\gamma_Q(uy-y,wx+\bar w ux).
\]
Expansion, skew-symmetry and $\gamma_Q(ux,uy)=\gamma_Q(x,y)$
show that this bicharacter is symmetric and conjugate-adjoint.
Its radical is exactly $C$.  Indeed
$wx+\bar w ux=x+\bar w(ux-x)$; since $u-1$ maps $B_1$
onto $L$, orthogonality against all of $B_1$ is equivalent,
by self-duality of $L$, to $x+\bar w(ux-x)\in L$,
and hence to $x\in L$.  Apply CW16 to $B_1/C$.
Using $w+\bar w=1$ and $u$-invariance once more simplifies its
quadratic function to
$\beta(wt,t)=\gamma_Q(\bar w(ut-t),t)$.
Thus, for any transversal $t_1,\ldots,t_n$,
\[
       \sum_{k=1}^n\gamma_Q(\bar w(ut_k-t_k),t_k)
                       =\lambda(n),\qquad\lambda(n)^2=n.
                                                               \tag{CW.38}
\]
All these are integral quotient cardinalities, not dimensions of
rational spans.

\paragraph{CW18: the finite and infinite phases of the same eta character.}
Now take the supplied character $\eta$ satisfying the literal
\texttt{IsEtaTilde}: it is Hecke, unramified at every finite place,
has \texttt{ResEta} with exponent $1$, and infinity exponent
$1/2$ in the source convention.  The last convention means the
phase exponent is $2(1/2)=1$.

The local signs can be obtained without a new class-field theorem.
HF9 constructs the character on $\mathbb A_F^\times$ having
local factors $(-1)^{\operatorname{ord}_v}$ at inert primes,
$1$ at split primes, and the sign at each real place.
It is continuous; HF9 proves that it kills principal ideles.
It kills norms because unramified quadratic norms have even
order, split norms are unrestricted, and complex norms are positive.
It is nontrivial on an idele negative at exactly one real place.
It therefore satisfies the literal \texttt{IsEta}.  The
\texttt{AllSome} clause in \texttt{ResEta} permits this particular
character.  At an inert place an $F_v$-uniformizer is also an
$E_P$-uniformizer, and consequently $\eta$ takes the value $-1$
on it.  It takes the value $1$ on all finite unit ideles by
unramifiedness.  This determines its entire inert factor.

Retain the same half selector $\varepsilon$ as in CW7.  For
$x\in E^\times$ let $\kappa(x)$ be its principal idele cut to
$\varepsilon$, and $\kappa'(x)$ the cut to
$\varepsilon+\bar\varepsilon$.
Write $\eta_{\rm in}(x)$ for the product over the inert places.
The Hecke property and the specified infinity type give
\[
 \eta(\kappa'(x))\eta_{\rm in}(x)
               \prod_\tau\frac{\tau(x)}{|\tau(x)|}=1.
                                                               \tag{CW.39}
\]
For integral $x\ne0$, factor its principal ideal.  At a fixed
prime $P$, CW2 gives $\lambda(NP)=-\sqrt{NP}$ and the preceding
local calculation gives $\eta_P(x)=(-1)^{\operatorname{ord}_P x}$.
At every moved prime the relevant rational primes are noninert,
so the corresponding $\lambda$ is the positive square root.
Multiplication proves
\[
 \lambda(|N_{E/\mathbb Q}x|)
       =\eta_{\rm in}(x)\sqrt{|N_{E/\mathbb Q}x|}.
\]
For arbitrary $x$ with $k=|N_{E/\mathbb Q}x|\in\mathbb N$, write
$x=u_0/z_0$ with $u_0,z_0$ nonzero algebraic integers.
The integer identity $k|N z_0|=|N u_0|$ and nonvanishing of
$\lambda(|Nz_0|)$ allow cancellation of the two integral formulas.
Combining with (CW.39) yields the exact positivity statement
\[
 \eta(\kappa'(x))\lambda(k)\prod_\tau\tau(x)
        =\sqrt{k}\prod_\tau|\tau(x)|>0.                \tag{CW.40}
\]
This uses the original $\eta,\lambda,\varepsilon$ and CM-type
embeddings; independent square-root choices would not give it.

\paragraph{CW19: explicit Gaussian theta tests and their zero coefficient.}
Fix Cartan involutions on $S,T$, and let $J$ be their real Hom
complex structure.  Write $\omega=\texttt{omO}$ and
$q(v)=\omega(v,Jv)>0$ for $v\ne0$.
Let $\iota$ be the inverse infinity chart and $\pi$ its projection.
For an adelic matrix $x$ define
\[
 \chi_A(x)=\mathbf1_A(x_f)\psi(b(wx_f,x_f)).
\]
For each real $a$ form the actual series
\[
 F_a(x)=\sum_{y\in V(E)}\psi(b(y,x))\,
    \psi(b(\iota a,y+x))\,\chi_A(y+x)
    \exp\bigl(-\tfrac\pi2 q(y_\infty+a+\pi x)\bigr).
                                                               \tag{CW.41}
\]
The rational vectors in this formula are diagonally embedded.
For fixed finite $x_f$, the support in $y$ is either empty or a
coset of $L$.  On a sufficiently small finite neighborhood the
same support and finite phases apply.  Real lattice-point counts
are polynomial, whereas the Gaussian and each derivative are a
polynomial times a positive Gaussian.  Therefore the series and
all real derivatives converge locally uniformly and absolutely.
This proves continuity and real smoothness, including exchange of
each derivative with the sum.

Reindexing $y$ proves the exact rational quasi-periodicity.
Expansion of the quadratic phase gives
\[
 \psi(-b(wt,t))\psi(b(x,t))\chi_A(x+t)=\chi_A(x)
                       \qquad(t\in A).
\]
It proves $W_tF_a=F_a$ term by term; real and finite components
have zero cross pairing.  On $2A$ the correcting phase is one,
so the raw finite Heisenberg condition holds on an open level.
The order $b(\iota a,y+x)$ in the second phase is essential:
it is the right magnetic translate of the Gaussian, which commutes
with the left Heisenberg derivatives; reversing its arguments
would not preserve the same vacuum equations.
Finally differentiating (CW.41) in directions $v,Jv$ cancels the
linear Gaussian derivative with the two Heisenberg phases, giving
$dH(v)F_a+i\,dH(Jv)F_a=0$.  Thus these are genuine theta functions
and simultaneous finite and Cartan vacua.

The family is not zero.  Setting $a=-v,x=\iota v$ expresses
$F_{-v}(\iota v)$ as an absolutely convergent Fourier series in
$v$, with frequency $y\mapsto-\omega(y_\infty,\cdot)$ and constant
coefficient $1$ from $y=0$.  All other frequencies are nonzero.
Self-duality makes them integral on $L$, so integration over its
real fundamental parallelepiped gives its positive volume.
In particular at least one $F_a$ is nonzero.  This supplies the
nonzero theta vector needed to make CW15's scalar unique.

Here is the precise trace calculation using these tests.  Use the
pair $p,u,C,n,i$ of CW17, and suppose that on every simultaneous
$A$-fixed Cartan vacuum one has
\[
 R_pf=\kappa\sum_{j=1}^{i}W_{b_j}f,\qquad
       b_j\in r(p)A/(A\cap r(p)A).                    \tag{CW.42}
\]
Define
\[
 \begin{split}
 Q_u(y)&=\mathbf1_{B_1}(y)\,
                \gamma_Q(\bar w(uy-y),y),\\
 G_u(v)&=e^{-\pi i\omega(uv,v)}
                     e^{-\pi q((u-1)v)/2}.
 \end{split}                                                   \tag{CW.43}
\]
Insert $F_{-v}$ in (CW.42) and evaluate at $\iota v$.
On the left, reindex by the bijection $y\mapsto(u-1)y$.
The phase equality is exactly
\[
 \chi_A((u-1)y)=Q_u(y)e^{-\pi i\omega(uy_\infty,y_\infty)},
\]
where the finite and real character values on a rational pairing
multiply to one.  Expanding the two remaining real phases then
gives $\sum_y Q_u(y)G_u(y_\infty+v)$.
On the right, the $j$th summand has Fourier coefficient
\[
 a_j(y)=\psi(-b(wb_j,b_j))\psi(b(y,b_j))
                    \chi_A(y+b_j)e^{-\pi q(y_\infty)/2}.
                                                               \tag{CW.44}
\]
These coefficients are absolutely summable over finitely many
lattice cosets.  The frequency is again
$-\omega(y_\infty,\cdot)$ and is integral on $C$ whenever
$a_j(y)\ne0$: test the self-dual lattices $A$ and $r(p)A$
against the finite images of $c$ and $uc$, for $c\in C$.
Only the zero frequency survives integration over a fundamental
parallelepiped for $C$.  Its total coefficient is
$\sum_j a_j(0)=1$, because exactly one representative belongs
to $A$, and there the two quadratic phases cancel.

On the left $Q_u$ is $C$-periodic, supported on the $n$ cosets
$t_k+C$.  Unfold each coset into the whole real vector space.
If $G_u$ is integrable, absolute integrability justifies this
operation and the right-hand Fourier integration.  CW20 proves
that integrability under the same condition $d\ne0$.
Writing $r_1>0$ for the real covolume of $C$, (CW.38) gives
\[
             \lambda(n)I_u=r_1\kappa,
                  \qquad I_u=\int G_u(v)\,dv.          \tag{CW.45}
\]
This proves the needed trace identity from the series, rather
than assuming it as a vacuum normalization axiom.

\paragraph{CW20: evaluating the real integral with its actual phase.}
Assume in addition that the rational pair commutes with the chosen
Cartans.  On the complex vectorization of the real Hom space let
$H$ be its original Hermitian matrix and $s$ its Cartan involution.
Thus $Hs$ is positive definite, $U^*HU=H$, and $Us=sU$;
the matrix $U$ here includes all CM-type embeddings of (CW.35).
The source convention is
$\omega(v_1,v_2)=-\Im(z_2^*Hz_1)$, so
$q(v)=z^*Hs z$.  Direct substitution in (CW.43) gives
$G_u(v)=e^{-\pi z^*E_u z}$ with
\[
 E_u=\tfrac12\bigl((U-1)^*Hs(U-1)-HU+(HU)^*\bigr).
                                                               \tag{CW.46}
\]
Its Hermitian part is
$\tfrac12(U-1)^*Hs(U-1)>0$, since $d\ne0$.
This proves absolute integrability.  CW9's complex Gaussian
formula therefore applies with no square-root ambiguity.
Put $P_+=(1+s)/2$, $P_-=(1-s)/2$.
Multiplying out (CW.46), using $U^*H=HU^{-1}$ and $Us=sU$,
gives the useful exact identity
\[
 E_u\,U(P_+U^{-1}+P_-)=-H(U-1).                       \tag{CW.47}
\]
Taking determinants and retaining the positive Jacobian between
the original real coordinates and these complex coordinates gives
\[
 I_u P_0=r_2\epsilon_0 Q_0\delta,\qquad r_2>0,\qquad
 P_0=\prod_\tau\tau(d),\quad Q_0=\prod_\tau\tau(\Delta),
                                                               \tag{CW.48}
\]
where $\epsilon_0=\det(-s)$ and
$\delta=\det(P_+U^{-1}+P_-)$.
Indeed $r_2$ is the coordinate Jacobian times $\det(Hs)^{-1}$,
which is positive real, and
$\det(-H)\det(-s)=\det(Hs)$ proves (CW.48).
The scalar $\delta$ is exactly CW13's real vacuum scalar for
$p$: $u$ was the inverse pair action.  It is nonzero by CW9,
and $\epsilon_0\ne0$ is fixed once the Cartans are fixed.
In particular no unknown complex factor is absorbed into $r_2$.

\paragraph{CW21: cancellation of all phases, including the finite index.}
Use the canonical characters
$\chi_S=\eta^{\dim S}$ and $\chi_T=\eta^{\dim T}$ in CW7.
For a rational pair commuting with the Cartans and satisfying
$d\ne0$, CW15, CW4 and CW13 give (CW.42) with
\[
       \kappa=c_p^{-1}\delta z\lambda(i)^{-1}.
                                                               \tag{CW.49}
\]
No uniqueness theorem for a merely unitary projective action is
being used: all three operators in this equation were constructed
on the same functions.  Thus CW19 gives
$\lambda(n)I_u=r_1c_p^{-1}\delta z\lambda(i)^{-1}$.
CW18 applied to the two numbers in (CW.37) gives positive reals
$r_3,r_4$ with
\[
 \begin{split}
 y_1\lambda(c_0)\lambda(n)ZP_0&=r_3,\\
 y_2\lambda(c_0)\lambda(i)ZQ_0&=r_4,
 \end{split}\quad
 \begin{gathered}
 y_1=\eta(\kappa'(N^m d)),\quad
 y_2=\eta(\kappa'(N^m\Delta)),\\
 Z=\prod_\tau\tau(N^m).
 \end{gathered}                                                \tag{CW.50}
\]
We check the remaining character identity explicitly.
Unitarity gives
\[
 \Delta\bar\Delta=1,\qquad
            \Delta\bar d=(-1)^m d.                    \tag{CW.51}
\]
The first is the determinant of $U^*HU=H$; the second follows
by taking determinants in
$U\bigl(\overline{U-1}\bigr)^{\mathsf t}$ after conjugation
by $H$, using $U^{-1}-1=-U^{-1}(U-1)$.
The cut identities and norm-triviality of $\eta$ give
\[
 \eta(\kappa'(x))\eta(\kappa(\bar x))=\eta(\kappa(x)),
 \qquad z\eta(\kappa(\Delta))=1.
\]
The last identity uses
$\Delta=\det(b_0)^{\dim S}\det(a)^{\dim T}$ and the
crossed character in (CW.11), including $b_0=h^{-1}$.
Applying these identities to (CW.51), and cancelling the common
$N^m$ factor, proves
\[
              zy_2=s_0y_1,\qquad
                 s_0=\eta(\kappa((-1)^m))\ne0.         \tag{CW.52}
\]
Inverting this order-two factor would give the same value;
no unspecified phase remains.

All factors cancelled below are nonzero: $d,\Delta,N$, the
character values, $c_p,\delta$, and positive-index values of
$\lambda$ have already been checked.  Multiply (CW.45) by
$c_p\lambda(i)$, substitute (CW.48), then eliminate $P_0,Q_0$
using (CW.50).  Finally use (CW.38) and (CW.52).  The result is
\[
             c_p\epsilon_0=r s_0,\qquad
                    r={r_1r_3\over r_2r_4 n}>0.       \tag{CW.53}
\]
This is where the same two rational lattice indices and the
signed Gauss sum are essential.  Absolute-value identities alone
would leave $c_p$ undetermined.

\paragraph{CW22: removing the no-fixed-vector restriction and proving automorphy.}
Keep any rational pair $p$ commuting with fixed Cartans, without
assuming $d\ne0$.  The norm-one elements of $E$ are infinite:
if $\xi\ne0$ is purely imaginary, then
$(a+\xi)/(a-\xi)$, $a\in\mathbb Q$, are distinct and have norm
one.  Choose such a $t$ outside the roots of
$\det(tU-1)$ and $\det(tI-1)$, two nonzero polynomials.
Let $k=tI$ on $T$ and $q=(1,k^{-1})$.
Then the inverse Hom actions for $q$ and $pq$ are $tI$ and $tU$.
They commute with the same Cartans and satisfy the no-fixed-vector
condition.  CW19 makes $\Theta$ nonzero, so the scalar in (CW.33)
is unique and the group law gives $c_{pq}=c_pc_q$.
Apply (CW.53) to $q$ and $pq$ and divide.  The same
$\epsilon_0,s_0$ cancel, proving $c_p$ is positive real.
Since $|c_p|=1$, necessarily $c_p=1$.
This argument also covers rank zero: the empty determinants and
zero-dimensional integrals are one.

Finally rational unitary groups are generated by the
quasi-reflections of WA10.  A quasi-reflection acts on a
nonisotropic line and fixes its orthogonal complement.
At each infinity embedding choose a Cartan decomposition preserving
that line, and choose positive and negative orthonormal bases on
its complement.  The resulting Cartan involution commutes with
the reflection.  For a reflection in one factor the identity in
the other factor commutes with any Cartan.  The preceding paragraph
therefore proves $c_p=1$ for each generator in either factor.
The actual group law extends it to all rational pairs.  We obtain
on the original theta space the full equality
\[
       \omega(\gamma,\gamma')f(x)
                =f\bigl(r(\gamma^{-1},\gamma'^{-1})x\bigr)
       \quad(\gamma\in G_S(E),\ \gamma'\in G_T(E)).     \tag{CW.54}
\]
Evaluating at zero gives rational theta-distribution automorphy
in both factors at once.  In particular this proves the literal
source and target automorphy clauses without a second representation
on transposed spaces.

\paragraph{CW23: the literal canonical full Weil datum.}
For the specified $D,\psi,S,T,\theta,\Lambda_S,\Lambda_T$ and
any supplied $\eta:\texttt{IsEtaTilde}$, the construction above
proves the actual conclusion of
\texttt{exists\_\allowbreak weilDatum\_\allowbreak etaTilde}.  We spell out its join.
The output is the group homomorphism on all functions obtained
by extending CW14's action on $\Theta$ once by an algebraic
complement.  Its characters are exactly
$\eta^{\dim S},\eta^{\dim T}$.  Powers preserve the Hecke
property, and raising the supplied restriction identity to a
power gives their two literal rank restrictions.

These are the first five fields of \texttt{IsWeilDatumO}, with
its standard-character field unchanged.  The remaining six are:
(CW4, CW13) preservation of the actual theta space;
(CW14) covariance with the full adelic Heisenberg action;
(CW22) rational evaluation-at-zero automorphy;
(CW14) pointwise continuity on each theta function;
(CW5, CW13) equality of the original Haar $L^2$ integrals;
and (CW6--CW7) the prescribed split factor times a positive
real scalar.  The norm equality holds for every additive Haar
measure and every rational fundamental domain in the definition,
so it also gives the literal nonnegative extended integral.

A split-supported group element has trivial real component:
apply the infinite projection to its support idempotent.
Thus its full action restricts to the finite action of CW7,
and the positive split-scalar clause is unchanged.  The same
argument for a nonsplit place shows that CW6 supplies
\texttt{FixesVacAt} for the original $\Lambda_S,\Lambda_T$.
At a real-supported pair the finite cut determinant is one and
its finite lattice action is the identity; hence CW13 supplies
\texttt{VacCharPlus}, with exactly (CW.29).
This proves both extra conclusions of the source theorem as
well as all eleven fields of its full datum.  It does not
replace them by a qualitative statement about algebraic jets.

\paragraph{CW24: prescribed splitting characters without an extra uniqueness premise.}
For arbitrary specified Hermitian spaces $S,T$ and specified
Hecke characters $\chi_S,\chi_T$ with the literal restrictions
$\texttt{ResEta}(\chi_S,\dim S)$ and
$\texttt{ResEta}(\chi_T,\dim T)$, a full oriented Weil datum
exists on these same spaces and functions.  We give a route
which also constructs the auxiliary character needed here,
without importing a new global norm-index theorem.

First use the actual standard lattice of each Hermitian space.
Its image under the given frame is
$\Lambda_S=\operatorname{frame}_S(S.L_0)$.
Finite generation and fullness are preserved by this linear
isomorphism.  The coordinate identity for the Hermitian form
shows that its dual is the image of the dual of $S.L_0$, hence
itself.  This proves the source's
\texttt{isSelfDualC\_\allowbreak latFrame}, and likewise for $T$;
no new Gram matrix is chosen.

HF11 supplies a unitary Hecke character $\eta_0$ of the actual
host, unramified at finite primes, with phase exponent one and
restriction equal to HF9's explicit quadratic character $\eta_F$.
Transport it to $(\mathbb A_F\otimes_F E)^\times$ by the
continuous tensor homomorphism
\[
 a\otimes x\longmapsto \operatorname{baseAdele}(a)
                                 \operatorname{diag}(x).
                                                               \tag{CW.55}
\]
On a finite component this is the completion extension, and at
infinity it is the specified embedding of $E$.  Agreement on
dense rational elements identifies each completion extension;
the finite-dimensional tensor presentation proves continuity.
It sends principal ideles to the same principal ideles, finite
integral units to finite integral units, and each local support
idempotent to its corresponding component.  Composing the HF11
character with (CW.55) therefore gives the same Hecke character
readings, infinity phases and inert values on the source's domain.
In particular its restriction is exactly $\eta_F$, and it kills
$x\bar x$ because $\eta_F$ kills norms.

There is a small logical distinction at this point.
\texttt{IsEtaTilde} contains a universal \texttt{AllSome} restriction
against every character satisfying \texttt{IsEta}.  Proving that
all those characters equal $\eta_F$ would be the separate global
norm-index-two uniqueness theorem.  We do not need that theorem.
Inspecting CW1--CW22, the construction of the operators uses only
that $\eta_0$ is Hecke, is unramified, has phase exponent one,
kills norms, and has the stated inert signs.  CW7 uses norm
triviality to define its cut character; CW18 uses precisely the
explicit local readings just proved.  Hence the same formulas
construct an action $\omega_0$ with covariance, continuity,
Haar norm, rational automorphy, and the positive split factor
for $\eta_0^{\dim S},\eta_0^{\dim T}$.  Only the two universal
restriction fields of a full record have not been asserted for
these auxiliary characters.

The supplied restrictions on $\chi_S,\chi_T$ can be tested
against the particular $\eta_F:\texttt{IsEta}$ from CW18.
Thus
\[
       \nu_S=\eta_0^{-\dim S}\chi_S,\qquad
       \nu_T=\eta_0^{-\dim T}\chi_T
\]
are Hecke characters trivial on \emph{all} ideles of $F$.
Apply the explicit determinant-root construction of WA18 to
these two characters and the action $\omega_0$:
\[
 \omega'(g,h)=\widetilde\nu_T(g)^{-1}
                            \widetilde\nu_S(h)\omega_0(g,h).
                                                               \tag{CW.56}
\]
WA18 proves root independence, restricted-product continuity,
rational triviality, and the exact crossed split-factor identity.
Those proofs use the operator properties just established, not
the missing universal restriction of $\eta_0$.
Consequently (CW.56) satisfies all operator clauses for the
prescribed $\chi_S,\chi_T$.  Its Hecke and two universal rank
restriction fields are the original supplied hypotheses.
This completes the full datum with no independent Weil-existence
or quadratic-character-uniqueness premise.  General character
twisting need not preserve the canonical extra vacuum determinant
or nonsplit fixing clauses; the full datum asserted here has
exactly its eleven stated fields.

\paragraph{CW25: actual bottom Weil families and the same bottom period field.}
Let $L$ be the supplied level-character datum with $L.\mathrm{Is}(\mu_0)$,
let $G$ be its supplied grid and $W:G.\mathrm{Weil}$ its supplied
realization.  For every square $(i,j)$, use WA17's actual standard
positive spaces $A,A',WW$ of ranks
$\mathrm{botBig},\mathrm{botSmall},\mathrm{botBig}+\mathrm{botSmall}$,
with their original identity frames and standard self-dual lattices.
Retain the prescribed bottom characters, including the exact
parity correction by $L.\eta^{\operatorname{rank}V\bmod2}$.
The rank identities in WA17 show that their restrictions have
exactly the required ranks.  CW24 now constructs, for every $t$,
the full representations on $(\operatorname{neg}A,WW)$ and
$(A',WW)$ with these characters and their prescribed character
quotients.  Negating the Hermitian form leaves the dual of the
same standard lattice unchanged.  The auxiliary $\eta_0$ of
CW24 does not replace $L.\eta$ in any bottom character.

This discharges the representation input explicitly retained in
WA17.  Its remaining construction chooses the normalized vacua,
finite idempotents, actual conductor exponents, standard finite
theta functions, Haar measures and fundamental domains on these
same spaces.  All identities refer to the same two
representations just constructed.  Thus the literal
\texttt{exists\_\allowbreak bottomTri} holds, including its actual
integrability clauses, by WA17 and TT12's definite quotient
argument.  No identification with the already supplied grid
representation $W$ is made.

If the twist set $T$ is finite, WA16 gives for each of these
actual bottom triangles $X_{ij}$ a finite number field
$M_{ij}\subset\mathbb C$ and nonzero periods for its two exact
integrals, satisfying \texttt{SumsIn}.  The compositum of the
finitely many $M_{ij}$ is finite.  Inclusion of fields preserves
all range-membership clauses without changing $X_{ij}$, its
measures or its periods.  This proves the literal bottom-field
portion \texttt{NFBot} for every supplied admissible realization;
its hypotheses provide this finiteness of $T$.

The canonical full bottom-Weil constructor is therefore no
longer a retained primitive.  The quantitative all-embedding
finite-unit dictionary estimates, the critical toric ratio field,
and the remaining same-choice \texttt{Choice.FieldOK} embedding
join retain the separate scopes recorded in WA19.  None is a
consequence of (CW.54) alone, and the supplied arithmetic host,
level data and grid realization have not been constructed anew.


\section{Actual adelic volumes and their compact-dual normalization}
\label{sec:geometric-volumes}

\status{GV1--GV9 prove the finite and real measure conversions, with
every determinant, dyadic, discriminant, and factorial factor specified.
GV11 proves the geometric definite quotient-to-mass identity.
GV10 and GV13--GV16 use one explicitly retained classical input: the
Tamagawa-number theorem for unitary groups, with the measure convention
stated in GV10. This gives genuine signature independence and the
literal \texttt{VolumeEqMass\textsubscript{2}} and \texttt{VolTieG}
suppliers on their stated data. GV14 also uses the named
Hasse--Minkowski theorem already retained in MF4. The final stage joins
keep their exact objects. Discrete-series existence and the formal-degree
formula, as well as the remaining nonempty compatible \texttt{V4.Core}
data, are not inferred from a volume formula. No Lean build or kernel
verification is asserted.}

\paragraph{Conventions and the new classical boundary.}
Let $E/F$ be a CM quadratic extension, with $F$ totally real of degree
$f$. Unless the assertion explicitly allows ramification, assume the
literal \texttt{CMData.cmUnram} condition. DR1 and MF1 identify this
condition with an unramified quadratic algebra at every finite place,
including places above $2$. Write $\epsilon_v=1$ at a split finite
place and $\epsilon_v=-1$ at a nonsplit place; $q_v$ is the residue
cardinality of $F_v$. Absolute values on $F$ are normalized by
$|\varpi_v|_v=q_v^{-1}$ and the usual absolute value at real places.
The quadratic Hecke character of $E/F$ is denoted $\epsilon$.
For $j\ge2$ put
\[
 L_F(j,\epsilon^j)=\prod_{v<\infty}
                  (1-\epsilon_v^j q_v^{-j})^{-1}.
 \tag{GV.1}
\]
These products converge absolutely and are positive. The positive
finite value $L_F(1,\epsilon)$ and its occurrence in Tamagawa measure
are part of the precise classical theorem retained in GV10; no
conditionally convergent Euler product at $1$ is manipulated below.

A supplied \texttt{HSpace} $W$ has a nondegenerate Hermitian form,
a rational frame with Gram matrix $J$, and an actual full self-dual
lattice $L=W.L_0$. No global integral basis of $L$ is assumed. Put
$n=\dim_E W$, $G_J=U(J)$, and $\Gamma=G_J(F)$, embedded diagonally.
Thus the source's $\texttt{HA}(\mathbb A_F)$ and
$\texttt{HFinIn}(\mathbb A_F)$ are respectively
$G_J(\mathbb A_F)$ and $\Gamma$. A CM-type chooses one complex
embedding over each real place. We use the actual supplied
$\texttt{CDHaar}$, rather than an independently chosen real Haar
measure. Set
\[
 a_n=\frac{n(n+1)}2,\qquad b_n=\frac{n(n-1)}2,\qquad
 v_n=\frac{2^n\pi^{a_n}}{\prod_{j=0}^{n-1}j!},\qquad
 A_n=2^{b_n}v_n=\frac{(2\pi)^{a_n}}{\prod_{j=0}^{n-1}j!}.
 \tag{GV.2}
\]
Empty products are one, so $v_0=A_0=1$.

\paragraph{GV1: the actual adelic group and unimodularity.}
Choose an $F$-basis of $E$. Expanding each matrix entry in this basis
identifies the source tensor algebra with the usual restricted product
of $E\otimes_F F_v$, and separates the finite and real components:
\[
 G_J(\mathbb A_F)\simeq G_J(\mathbb A_{F,f})
                         \times\prod_{\tau:F\hookrightarrow\mathbb R}
                                      U(p_\tau,q_\tau).
 \tag{GV.3}
\]
The equations $g^*Jg=J$ and $g^{-1}=J^{-1}g^*J$ prove that these are
topological group identifications, not merely bijections of sets.
The finitely many exceptional coefficient denominators do not change
the restricted-product topology. Finite local matrix spaces and the
real matrix spaces are locally compact and second countable; using the
countably many finite sets of places proves the same for (GV.3).

These groups are unimodular. Indeed, after extending a local base field
to split $E/F$, the Lie algebra becomes $M_n$ and the adjoint action is
$X\mapsto gXg^{-1}$. Its determinant is
$(\det g)^n(\det g^{-1})^n=1$. This determinant identity descends to
the original field. A nonzero invariant top differential therefore
gives a Haar measure invariant under right translation as well as
left translation. At real places the same conclusion follows from
the absolute determinant on the real Lie algebra. Local products and
the restricted product are consequently also unimodular. In
particular every inner conjugation preserves Haar measure.

\paragraph{GV2: the source completion is the completed lattice.}
The symbol \texttt{AddSubgroup.closure} in \texttt{latHat} means the
algebraically generated additive subgroup. It is not a topological
closure. Nevertheless, for a full lattice $L$ it is exactly
\[
 \widehat L=\prod_{v<\infty}L_v,
 \qquad L_v=\mathcal O_{F_v}\otimes_{\mathcal O_F}L.
 \tag{GV.4}
\]
To see this, choose finitely many $\mathcal O_F$-generators $x_i$ of
$L$. At each $v$, every element of $L_v$ is $\sum_i a_{iv}x_i$ with
$a_{iv}\in\mathcal O_{F_v}$. The coefficients $a_i=(a_{iv})_v$ lie
in $\widehat{\mathcal O_F}$ and give the finite sum of pure tensors
allowed by \texttt{latHat}. Conversely every such generator belongs
to every $L_v$. This proves (GV.4) in the source's fixed rational frame.

The lattice is compact, as the image of a finite power of
$\widehat{\mathcal O_F}$. It is open: a common nonzero denominator
gives
$c\widehat{\mathcal O_E}^{n}\subseteq\widehat L
 \subseteq c^{-1}\widehat{\mathcal O_E}^{n}$.
Its stabilizer
\[
 K_L=\{g\in G_J(\mathbb A_{F,f}):g\widehat L=\widehat L\}
      =\prod_{v<\infty}K_{L_v}
 \tag{GV.5}
\]
is compact and open. In local bases a lattice stabilizer in the general
linear group is $\mathrm{GL}_n(\mathcal O_{E_v})$ (the corresponding
product at a split place); intersect with the closed unitary equations.
These descriptions agree with the standard integral model outside a
finite set. They show compactness and openness directly, and also show
that both implications in the source's \texttt{StabHat} are required.

\paragraph{GV3: existence and uniqueness of the prescribed Haar measure.}
Give $G_J(\mathbb A_{F,f})$ its unique Haar measure with
$K_L$ of mass one. At every real place choose a matrix
$A_\tau$ with
$A_\tau^*J_{p_\tau,q_\tau}A_\tau=\tau J$ and transport
$\texttt{cd.cd}(p_\tau,q_\tau)$ by $g\mapsto A_\tau gA_\tau^{-1}$.
The product through (GV.3) is a Haar measure $\mu_{W,L}$ satisfying
\[
 \mu_{W,L}\left(K_L\times\prod_\tau A_\tau^{-1}B_\tau A_\tau\right)
                    =\prod_\tau\texttt{cd.cd}(p_\tau,q_\tau)(B_\tau)
 \tag{GV.6}
\]
for all measurable $B_\tau$.

Every other signature matrix is $u_\tau A_\tau$ with
$u_\tau\in U(p_\tau,q_\tau)$. GV1 makes (GV.6) independent of these
choices. The archimedean algebra map extending a chosen embedding of
$E$ is unique: on $a\otimes x$ it is $a_\tau\tau(x)$, and the
real-coordinate idempotents isolate the prescribed $\tau$.
Thus both existential and universal parts of the source's
\texttt{AllSome} clauses hold. Conversely two Haar measures satisfying
(GV.6) are proportional by Haar uniqueness; evaluate a rectangle with
each real factor compact of positive finite measure to see that the
proportionality constant is one. This proves literal
\texttt{IsVolMeasF} existence and uniqueness, including nonconstant
signature functions, and its constant-signature \texttt{IsVolMeas}
specialization.

\paragraph{GV4: rational discreteness and genuine fundamental domains.}
The diagonal $F$ in $\mathbb A_F$ is discrete. For example, force the
finite coordinates into $\widehat{\mathcal O_F}$ and take a sufficiently
small neighborhood of zero in all real coordinates; the Minkowski
image of $\mathcal O_F$ has no nonzero point in that neighborhood.
Apply this entrywise after fixing an $F$-basis of $E$. A neighborhood
of the identity in the adelic matrix group then contains no rational
matrix except the identity. Hence $\Gamma$ is a discrete countable
subgroup of the group in GV1.

For completeness one can construct a Borel fundamental domain without
presupposing finite volume. Cover the group by countably many open
sets $U_i$ such that $\gamma U_i\cap U_i$ is empty for
$\gamma\ne1$. Such sets exist by discreteness. Put
$D_i=U_i\setminus\bigcup_{j<i,\gamma\in\Gamma}\gamma U_j$.
The union $D=\bigcup_iD_i$ meets each orbit exactly once. It is Borel,
and its translates partition the group. For any two measurable
fundamental domains $D,D'$, Tonelli and left invariance give
\[
 \mu(D)=\int_D\sum_{\gamma\in\Gamma}1_{D'}(\gamma g)\,d\mu(g)
       =\mu(D').
 \tag{GV.7}
\]
The common value is positive: otherwise countably many null translates
would cover a group carrying nonzero Haar measure. These facts also
hold for the source's almost-everywhere fundamental domains by the
same integral argument. Finiteness has not been assumed here.

\paragraph{GV5: the exact real normalization, including the powers of two.}
First use the following real coordinates on $\mathfrak u(n)$:
diagonal entries $it_i$, and, for $i<j$, entries
$X_{ij}=a_{ij}+ib_{ij}$, $X_{ji}=-a_{ij}+ib_{ij}$.
Let $dX$ be coordinate Lebesgue measure. For the map
$(X,Y)\mapsto X+iY$, a diagonal pair has real determinant of absolute
value one. An off-diagonal pair has coordinates
\[
 (a,b,c,d)\longmapsto(a-d,b+c,-a-d,b-c)
 \tag{GV.8}
\]
and determinant of absolute value four. The unique measure satisfying
the source's \texttt{IsOmegaLeb} condition is therefore
$\lambda_n=2^{b_n}dX$: its product with itself pushes to entrywise
complex Lebesgue measure exactly.

The invariant top density $dX$ has total volume $v_n$ on $U(n)$.
Here is a direct normalization proof. The first-column fibration
$U(n)\to S^{2n-1}$ has fiber $U(n-1)$. At the identity the horizontal
coordinates $t_1,a_{1j},b_{1j}$ map, with absolute determinant one,
to the ordinary tangent coordinates of the unit sphere. Splitting
the remaining coordinates gives precisely the same $dX$ on the
fiber. Invariant-measure disintegration thus yields
\[
 v_n=\operatorname{area}(S^{2n-1})v_{n-1}
     =\frac{2\pi^n}{(n-1)!}v_{n-1},\qquad v_0=1,
 \tag{GV.9}
\]
which is (GV.2). This argument specifies the quotient density; it
does not replace it by an unrelated quotient Riemannian metric.

For $J_{p,q}=\operatorname{diag}(1_p,-1_q)$, multiplication
$X\mapsto J_{p,q}X$ carries its unitary Lie algebra to
$\mathfrak u(n)$. Its complexification is multiplication by
$J_{p,q}$ on $M_n(\mathbb C)$, of complex determinant of modulus
one. It therefore preserves the exact Omega measure. An invariant
smooth Haar density is determined by its density at the identity;
the exponential map has derivative the identity there. The small-ball
limit in \texttt{OmegaDensity} is exactly that density relative to
$\lambda_n$, for the source's entry-sum balls as well as Euclidean
balls. The \texttt{CDHaar.prob} and \texttt{CDHaar.trick} fields now
force
\[
 \texttt{cd.cd}(p,q)\text{ to have Lie density }
                  A_n^{-1}\lambda_{p,q}.
 \tag{GV.10}
\]
In particular this holds for every supplied \texttt{CDHaar}, with
the same explicit constant $A_n^{-1}$ for every signature.

\paragraph{GV6: a rational gauge and its determinant under isometries.}
Choose $0\ne\delta\in E$ with $\bar\delta=-\delta$.
Use the following $F$-basis $\mathcal E$ of skew-Hermitian matrices:
$\delta E_{ii}$, and $E_{ij}-E_{ji}$,
$\delta(E_{ij}+E_{ji})$ for $i<j$.
On $\mathfrak u_J$ define a rational top form $\omega_J$ by declaring
the linear map $X\mapsto JX$ to identify its tangent density with
the coordinate form in $\mathcal E$. Translate this form on the group.
This tangent identification is not asserted to be a Lie homomorphism.

For any $B$ over a local quadratic algebra,
\[
 \det_F(Y\mapsto B^*YB)=N_{E/F}(\det B)^n
 \quad\text{on skew-Hermitian matrices}.
 \tag{GV.11}
\]
One proof is to split the quadratic algebra. The skew-Hermitian
matrices then identify with $M_n$, and the map is left and right
multiplication by the two component matrices of $B$, whose determinant
is the displayed norm power. Equality descends because these are
polynomial determinant identities. This proof covers a split place
and a dyadic place without dividing by two.

If $B^*JB=I$, the conjugation $g\mapsto B^{-1}gB$ transforms
$Y=JX$ into $B^*YB$. Its tangent determinant is $(\det J)^{-n}$.
Consequently, if it carries $K_{L_v}$ to $K_{\mathrm{std},v}$,
\[
 |\omega_J|_v(K_{L_v})
       =|\det J|_v^n\,|\omega_I|_v(K_{\mathrm{std},v}).
 \tag{GV.12}
\]
At a real place with $A^*J_{p,q}A=\tau J$, the corresponding map on
$Y$ is $Y\mapsto A^{-*}YA^{-1}$, again with absolute determinant
$|\det\tau J|^{-n}$. If $t_\tau=|\tau\delta|$, exactly $a_n$ of
the coordinates in $\mathcal E$ acquire a factor $t_\tau$.
By GV5,
\[
 |\omega_J|_\tau
   =|\det\tau J|^n\,t_\tau^{-a_n}v_n\,
                         \mu_{W,L,\tau}.
 \tag{GV.13}
\]
Here $\mu_{W,L,\tau}$ is the very real factor used in GV3.

\paragraph{GV7: the integral Lie lattice at dyadic as well as odd places.}
Let $R=\mathcal O_{F_v}$ and $S=\mathcal O_{E\otimes F_v}$.
An unramified quadratic algebra has an anti-invariant unit
$\kappa\in S$ with $\bar\kappa=-\kappa$. In the split case take
$(1,-1)$. In the nonsplit case lift an element whose residue
$z-\bar z$ is nonzero; such a residue exists also in characteristic
two because the trace map of the quadratic finite-field extension
is onto. Thus $S^- =R\kappa$.

The integral Lie lattice is the direct sum of $n$ copies of $S^-$
on the diagonal and $b_n$ copies of $S$ off the diagonal. The diagonal
coordinate determinant relative to $\delta$ has absolute value
$|\delta^2|_v^{-1/2}$. An integral basis of $S$ has unit
discriminant, whereas the discriminant of $(1,\delta)$ is
$4\delta^2$. The off-diagonal determinant relative to that basis
therefore has absolute value $|2|_v^{-1}|\delta^2|_v^{-1/2}$.
The determinant of the whole integral Lie lattice relative to
$\mathcal E$ is exactly
\[
 d_v=|2|_v^{-b_n}|\delta^2|_v^{-a_n/2}.
 \tag{GV.14}
\]
The factor $|2|_v^{-b_n}$ is essential; omitting it would give an
incorrect archimedean-to-finite normalization.

For clarity, the integral unitary group is smooth at every such
place by a direct lifting calculation. If a matrix is unitary modulo
$\varpi^r$, multiplication by $1+\varpi^r X$ changes its Hermitian
error modulo $\varpi^{r+1}$ by $X+X^*$. On off-diagonal entries this
map is onto by freely choosing one entry of each pair; on the
diagonal it is the surjective residue trace $S\to R$. Its kernel
has residue dimension $n^2$. Successive corrections converge in
the complete DVR. In particular reduction onto the finite unitary
group is surjective, and each subsequent congruence fiber has
$q_v^{n^2}$ elements. Equivalently, the same equations give local
analytic coordinates with unit Jacobian and $n^2$ free variables.
An invariant integral top density gives each first congruence
neighborhood mass $q_v^{-n^2}$. This supplies the measure assertion
in addition to a finite-set lifting assertion.

\paragraph{GV8: the exact finite normalized factors.}
For the standard self-dual lattice, GV7 gives
\[
 |\omega_I|_v(K_{\mathrm{std},v})
       =d_vq_v^{-n^2}|U_n(k_v)|
       =d_v\prod_{j=1}^n(1-\epsilon_v^j q_v^{-j}).
 \tag{GV.15}
\]
The notation $U_n(k_v)$ means $\mathrm{GL}_n(k_v)$ at a split
place. There its order follows by choosing independent columns.
At a nonsplit place, the number of norm-one vectors in $k_{v^2}^n$
is $q_v^{n-1}(q_v^n-(-1)^n)$: separate the last coordinate and use
the $q_v+1$ preimages of every nonzero norm, as in MF6. The unitary
group acts transitively on them, with stabilizer $U_{n-1}(k_v)$;
induction yields (GV.15). The residue-trace argument in GV7
establishes the same formula at $v\mid2$.

For a general self-dual $L_v$, DR1 gives an integral orthonormal
basis in the nonsplit case. At a split place choose a basis in one
factor and the conjugate inverse-transpose basis in the other.
Together these provide $B^*JB=I$ carrying the standard lattice to
$L_v$. Combining (GV.12) and (GV.15), relative to Haar measure
$dg_v$ with $K_{L_v}$ of mass one, gives
\[
 |\omega_J|_v
  =|\det J|_v^n d_v
                  \prod_{j=1}^n(1-\epsilon_v^j q_v^{-j})\,dg_v.
 \tag{GV.16}
\]
All determinants here are local coordinate determinants; no
globally free basis for $L$ has been introduced.

\paragraph{GV9: the global conversion constant, with no unspecified scalar.}
For $n>0$ define the adelic gauge normalization
\[
 d\tau_J=|\operatorname{disc}F|^{-n^2/2}L_F(1,\epsilon)^{-1}
        \prod_{\tau\mid\infty}|\omega_J|_\tau
        \prod_{v<\infty}L_v(1,\epsilon)|\omega_J|_v,
 \quad L_v(1,\epsilon)=(1-\epsilon_vq_v^{-1})^{-1}.
 \tag{GV.17}
\]
This is the convention of the classical theorem in GV10.
Its local products converge: the factor $L_v(1,\epsilon)$ cancels
the $j=1$ term of (GV.16), leaving $1+O_n(q_v^{-2})$; the
coefficient and determinant exceptions have finite support.

Multiply (GV.13) and (GV.16). First
$\prod_v|\det J|_v^n=1$ by the product formula in $F$.
Next $\delta^2\in F^\times$, $|\tau(\delta^2)|=t_\tau^2$, and
\[
 \left(\prod_{v<\infty}d_v\right)
                  \left(\prod_{\tau\mid\infty}t_\tau^{-a_n}v_n\right)
     =2^{b_nf}v_n^f=A_n^f.
 \tag{GV.18}
\]
Indeed the $\delta^2$ powers cancel by the product formula, and
$\prod_{v<\infty}|2|_v=2^{-f}$. The remaining finite Euler products
are absolutely convergent, so rearranging the finitely many $j$-products
is legitimate. We obtain the exact equality of measures
\[
 \boxed{\quad d\tau_J=C_n\,d\mu_{W,L},\qquad
 C_n=\frac{|\operatorname{disc}F|^{-n^2/2}A_n^f}
                   {\displaystyle\prod_{j=1}^nL_F(j,\epsilon^j)}.
       \quad}                                      \tag{GV.19}
\]
This is independent of $J$, the self-dual lattice, the rational
frame, the signature at every real place, and the chosen
$\texttt{CDHaar}$. Each independence follows from the calculated
factors. In particular (GV.19) does not assume signature independence
of quotient volumes or a mass formula.

\paragraph{GV10: the one retained Tamagawa theorem and finite covolume.}
We use the following established theorem as an explicit external input:
for every CM quadratic $E/F$ and every nonzero nondegenerate Hermitian
space, the measure (GV.17), using its usual local gauges and quadratic
$L$-value, has
\[
                    \tau_J\bigl(G_J(F)\backslash G_J(\mathbb A_F)\bigr)=2.
 \tag{GV.20}
\]
This includes finiteness. A precise published normalization and statement
are Harris--Li--Skinner, \emph{p-Adic L-Functions for Unitary Shimura
Varieties I}, Documenta Mathematica, Extra Volume Coates (2006),
\S0.2, p.~400, equation (0.2.1) and the following Tamagawa-number
statement; equation (0.2.2) gives its hyperspecial local convention.
The classical treatment is Weil, \emph{Adeles and Algebraic Groups},
Chapter IV, \S4.4.
\footnote{Primary publisher text:
\url{https://ems.press/content/book-chapter-files/27474};
DOI \url{https://doi.org/10.4171/DMS/4/12}.}
The theorem itself is not reconstructed here. It is stronger than
an assertion that an unspecified Haar measure has finite quotient
volume: its normalization is precisely (GV.17).

For an unramified extension with a self-dual lattice, GV9 therefore
gives, for every fundamental domain $D$, the positive finite value
\[
 \mu_{W,L}(D)=M_n,\qquad
 M_n=\frac{2|\operatorname{disc}F|^{n^2/2}}
                   {A_n^f}\prod_{j=1}^nL_F(j,\epsilon^j)
                   \quad(n>0).                     \tag{GV.21}
\]
For $n=0$ the group and rational subgroup are trivial and the
normalization in GV3 is the point mass one; put $M_0=1$.
The more general source statement \texttt{FiniteVolume\textsubscript{2}}
also follows from the same classical theorem, even when the CM
extension is ramified or no self-dual lattice is specified: any Haar
measure on this group is a positive finite scalar multiple of its
Tamagawa measure. GV4 supplies its fundamental domains and the common
mass. This general finiteness assertion does not extend (GV.19) to
ramified local normalizations.

\paragraph{GV11: the definite geometric quotient really is the genus mass.}
Assume $W$ is totally positive definite, and set
$K=G_J(F\otimes\mathbb R)K_L$. Its $\mu_{W,L}$-mass is one by
\texttt{CDHaar.prob} and (GV.6). The map
\[
 g\longmapsto L_g=W\cap g_f\widehat L
 \tag{GV.22}
\]
identifies $\Gamma\backslash G_J(\mathbb A_F)/K$ with rational
isometry classes in the genus of $L$. Here intersection and localization
are the actual operations of LP1. A tuple of local isometries is
integral outside finitely many places and so gives an adele;
conversely its intersection has exactly the prescribed localizations. At each of the finitely many exceptional
places, approximate a completed basis by rational vectors closely
enough to preserve its integral span; invertibility and the lattice
stabilizer are open. LP1 then applies to actual rational local modules.
The source genus uses rational local actions. MF3's Cayley weak
approximation places a rational unitary matrix in each needed open
local transporter, so the completed and rational-local formulations
of this genus agree. This does not assert strong approximation.

The positive trace-form proof of TL1 applies to this definite form as
well as to the coordinate identity form: self-duality bounds the
integral trace Gram determinant, positivity bounds multiplication
operators by $\max_\tau|\tau(a)|$, and the same short-vector argument
gives finitely many classes and finite automorphism groups. For a
representative $g_i$, its rational stabilizer in $g_iKg_i^{-1}$ is
exactly $\operatorname{Aut}(L_{g_i})$: preserve all completions and
intersect, or complete a global lattice automorphism. Denote its
order by $s_i$.

The normalization can now be checked without a volume theorem.
The nonnegative function
\[
                  \Phi(g)=\sum_i s_i^{-1}1_{g_iK}(g)
 \tag{GV.23}
\]
has $\sum_{\gamma\in\Gamma}\Phi(\gamma g)=1$ for every $g$.
Indeed exactly one double coset occurs, and in it the number of
translates meeting $g_iK$ is $s_i$. Integrate this equality on a
fundamental domain and use Tonelli as in GV4. Since each $g_iK$ has
measure one,
\[
 \boxed{\ \mu_{W,L}(D)=\sum_i\frac1{s_i}
                    =\operatorname{mass}(W,L).\ }   \tag{GV.24}
\]
The source's fallback definition of \texttt{mass} is harmless here:
the actual finite class representatives have just been supplied.
GV11, including finite covolume in the definite case, is independent
of the Tamagawa input in GV10.

\paragraph{GV12: exact lattice and rational-isometry transport.}
For two self-dual lattices in the same Hermitian space, DR1 and the
split-pair construction in GV8 give a finite adelic unitary element
$u$ carrying one completion to the other; take $u_v=1$ outside the
finite set where the lattices differ. Thus their compact open
stabilizers are conjugate. By GV1 a fixed finite Haar measure gives
them equal mass, so the normalized measure in GV3 is the same for
both lattices. This proves the literal \texttt{LatticeIndep}.

If a rational isometry $e:W_1\to W_2$ carries the specified lattices,
let $T$ be its matrix in the two specified rational frames.
Conjugation $\varphi(g)=TgT^{-1}$ is a topological group isomorphism
carrying the rational subgroup, the completed stabilizer, and the
real rectangle data to their counterparts. In particular
$\varphi(g)T=Tg$, exactly \texttt{IsConjBy}. Uniqueness in GV3 gives
$\varphi_*\mu_1=\mu_2$, and the partition definition of a fundamental
domain proves both directions of its transport clause. If $e$ does
not carry the specified lattices, first replace $L_2$ by $eL_1$ and
use lattice independence. Under the literal supplied \texttt{VolTie},
the equal domain masses then give $D.\mathrm{vol}(W_1)=D.\mathrm{vol}(W_2)$.
This proves all three clauses of \texttt{VolIndep} without Tamagawa
or a global isometry between different signatures.

\paragraph{GV13: actual signature independence and \texttt{VolumeEqMass\textsubscript{2}}.}
Now use the named input in GV10. Formula (GV.21) gives equal normalized
fundamental-domain masses for any two self-dual Hermitian spaces of
the same rank. This includes varying signatures at different real
places, and in particular all quantifiers in the source's
\texttt{VolSigIndep}. Apply GV11 to the actual positive standard
space $(E^n,I,\mathcal O_E^n)$. Its mass is the source's
$\mathrm{massI}(D,n)$, and hence
\[
          \mathrm{massI}(D,n)=M_n,\qquad
          \mu_{W,L}(D_W)=\mathrm{massI}(D,n).
 \tag{GV.25}
\]
Thus every actual positive self-dual genus has the same absolute mass;
no prior identification of rational Hermitian spaces was used.

For a supplied record $D$ with \texttt{VolTie}, its own normalized
domain mass is $D.\mathrm{vol}(W)$. Positivity permits passing from
$\mathrm{ofReal}$ to real numbers, so
\[
              D.\mathrm{vol}(W)=\mathrm{massI}(D,\operatorname{rank}W)
 \tag{GV.26}
\]
for every constant-signature $W$. The signature sum gives
$\operatorname{rank}W=P+Q$. This proves both literal clauses of
\texttt{VolumeEqMass\textsubscript{2}}, including the positive clause
in rank zero and the indefinite clause with $2\le P$.
The proof actually gives the indefinite equality without that lower
bound. It does not assign this volume to an arbitrary preexisting
positive function called $D.\mathrm{vol}$; its \texttt{VolTie}
hypothesis is essential.

\paragraph{GV14: a norm-sign vector and a self-dual orthogonal complement.}
This paragraph uses the named Hasse--Minkowski theorem retained in MF4.
Let $W$ have rank at least three, constant signature $(P,Q)$, and
choose $s=1$ with $P\ge1$ or $s=-1$ with $Q\ge1$.
At each finite place a local orthonormal basis exists by GV8 and the
unit norm map is onto by DR1, so $h(x,x)=s$ has a local solution.
At real places it has a solution by the stipulated signature.
On the underlying $F$-space let $q(x)=h(x,x)$.
The quadratic form $q\perp\langle-s\rangle$ is locally isotropic
everywhere. Hasse--Minkowski gives a global isotropic pair $(x,t)$.
If $t\ne0$, replace $x$ by $x/t$. If $t=0$, $q$ has a nonzero
isotropic vector; nondegeneracy produces a hyperbolic plane, which
represents every scalar, in particular $s$. Thus there is a rational
$x$ with $h(x,x)=s$.

Let $W'=x^\perp$. The formula
$y\mapsto y-s^{-1}x h(x,y)$ is the rational orthogonal projection,
so $W=Ex\perp W'$ and $W'$ has the required decremented signature.
We must supply its lattice, rather than assume that $L\cap W'$
is self-dual at every place. At a finite place choose a unit of norm
$s$. MF3's equal-norm reflection transports its first coordinate
vector to $x$, leaving a standard self-dual lattice on the
orthogonal complement. At a split place the same construction is
elementary linear algebra: choose a first vector and a complementary
kernel of its paired covector, then use the inverse-transpose dual
basis in the other factor. At all but finitely many places $x\in L_v$;
there the projection is integral, so
$L_v=\mathcal O_{E_v}x\perp(L_v\cap W'_v)$ and its complement is
already self-dual. Approximate their finitely many exceptional completed bases by rational
bases inside the same open lattice transporter, then patch the resulting
rational local complement lattices by LP1. The result is a full global self-dual lattice in
$W'$. Choose any rational frame for $W'$ and restrict the form.
Its inclusion into $W$ has range $\ker h(x,-)$ and is isometric.
These are precisely all the \texttt{UnitSplit} fields. A global
integral orthonormal basis and a global Landherr normal form were
not required.

\paragraph{GV15: the geometric successor ratio and \texttt{Conclusion2}.}
For $n\ge3$, divide (GV.21) in ranks $n$ and $n-1$:
\[
 \frac{M_n}{M_{n-1}}
  =|\operatorname{disc}F|^{n-1/2}
       \left(\frac{(n-1)!}{(2\pi)^n}\right)^f
                     L_F(n,\epsilon^n).
 \tag{GV.27}
\]
MF1 proves $f=\mathrm{fDeg}(E)$,
$|\operatorname{disc}F|=(2\,\mathrm{vInv}(E))^f$, and
$L_F(n,\epsilon^n)=\mathrm{lValEta}(E,n)$, with the two split-prime
half powers paired exactly. Substitution turns (GV.27) into
\[
 \frac{\big((n-1)!(\mathrm{vInv}(E)/\pi)^n\big)^f}
            {\sqrt{(2\,\mathrm{vInv}(E))^f}}
                       \mathrm{lValEta}(E,n),         \tag{GV.28}
\]
the literal \texttt{massRatioRightSide}. All divisors are positive.
For the actual $x,W',\iota$ of GV14, GV13 and \texttt{VolTie}
identify the left side with $D.\mathrm{vol}(W)/D.\mathrm{vol}(W')$.
This proves every clause of \texttt{Conclusion2}, and, before using
\texttt{VolTie}, its actual-normalized-domain version
\texttt{VolRatioTrue}. It is consistent with MF19, but (GV.27) was
derived from the geometric quotient, not by relabeling MF19's
mass law as a volume law. Its classical premises are exactly those
stated in GV10 and GV14.

\paragraph{GV16: the source's chosen \texttt{trueVol} has its intended value.}
For every $W$ satisfying \texttt{SigOK}, GV3 gives at least one
normalized measure and proves the property for every such measure.
GV4 gives at least one fundamental domain and proves equality for
every such domain. GV10 gives their common positive finite real
value $M_n$. Consequently
\[
 \mathrm{VolIs}(\mathrm{cm},\tau,\mathrm{cd},W,M_n),\qquad
 \mathrm{trueVol}(\mathrm{cm},\tau,\mathrm{cd},W)=M_n.
 \tag{GV.29}
\]
For the second equality, the defining existential branch is now
inhabited; its chosen value has the same domain mass and is positive,
so uniqueness forces equality. The fallback value $1$ is never used
on these valid spaces. This proves the actual
\texttt{volIs\_\allowbreak trueVol} supplier and hence
\[
 [\,D.\mathrm{vol}(W)=\mathrm{trueVol}(D.\mathrm{sp.cm},
              \tau_D,D.\mathrm{cd},W)\text{ for every }W\,]
                 \ \Longrightarrow\ \mathrm{VolTieG}(D).
 \tag{GV.30}
\]
The host CM-type is the one in the original record; its bijection
with infinite places is part of the host interface, not a change
of embeddings. Constant-signature specialization also gives
\texttt{VolTie}. For malformed signature fields outside
\texttt{SigOK}, no assertion that a chosen fallback has geometric
meaning is made.

\paragraph{GV17: the exact compact-dual scale in every sign ordering.}
Fix a Boolean sign function $e:\{1,\ldots,n\}\to\{+,-\}$ and
a sorting bijection $\rho:\{1,\ldots,P\}\sqcup\{1,\ldots,Q\}
\to\{1,\ldots,n\}$. Reindexing matrices is a topological group
isomorphism $U(P,Q)\to U(J_e)$. Its pushforward of
$\texttt{cd.cd}(P,Q)$ and the source's fixed
$\mathrm{REAL.muG}(e)$ are Haar measures, so there is a unique
positive finite scalar $c_e$ relating them. If a different sorting
is chosen, the two maps differ by conjugation with a permutation
unitary matrix; GV1 makes the scalar unchanged. Thus the same $c_e$
works for every sorting and every measurable set, exactly as demanded
by \texttt{CdScaleIs}. Positivity and uniqueness can be checked on
any compact set with nonempty interior.

Choosing these $c_e$ for all $n,e$ inhabits the existential branch
in \texttt{trueScale}; uniqueness proves
\[
               \mathrm{CdScaleIs}(\mathrm{cd},
                                  \mathrm{trueScale}(\mathrm{cd})).
 \tag{GV.31}
\]
Here no numerical normalization of \texttt{REAL.muG} is silently
assumed: $c_e$ is its exact positive Haar ratio, and the defining
identity on all measurable sets is proved. The explicit numerical
normalization of \texttt{cd.cd} itself was already computed in GV5.
This proves \texttt{cdScaleTie} whenever the same record's scale
function equals \texttt{trueScale}. It is independent of the
Tamagawa theorem and of any discrete-series existence statement.

\paragraph{GV18: the actual quotient probability measure and level changes.}
Push the restriction of $\mu_{W,L}$ to a fundamental domain onto
$\Gamma\backslash G_J(\mathbb A_F)$. For a nonnegative measurable
function on this quotient, its pullback is $\Gamma$-invariant;
the argument of (GV.7) shows that the integral is independent of
the selected domain. Thus the quotient measure, not only its total
mass, is canonical. Its mass is $M_n$, and its probability
normalization is precisely
\[
                       \mu_Q^1=M_n^{-1}\mu_Q.
 \tag{GV.32}
\]
If the same data have \texttt{VolTieG}, replace $M_n$ by
$D.\mathrm{vol}(W)$ in (GV.32). This proves the source's
\texttt{VolPairEx}, \texttt{FinVol},
\texttt{normalized\_\allowbreak quotient\_\allowbreak measure\_\allowbreak eq}, and
\texttt{probability\_\allowbreak measure\_\allowbreak eq\_\allowbreak pair}; none of their fallback
branches is used on a valid space.

There is also an exact finite-level check. If $K'\subseteq K_L$
is an open subgroup of finite index $r$, then
$\mu_f(K')=1/r$. Renormalizing the finite Haar measure to give
$K'$ mass one multiplies the adelic quotient mass by $r$ and leaves
the normalized probability measure unchanged. More generally the
ratio for two compact open levels is the ratio of their indices
over their common intersection. This follows from their finite
coset decompositions, and fixes the direction of every level factor.
It does not identify an unspecified model level with $K_L$.

\paragraph{GV19: the original base-volume stage, on the same base and columns.}
IG12 constructs, for the actual $B,C$ and supplied
$\mathrm{cd}$, a stage $G$ with $G.\mathrm{cd}=\mathrm{cd}$ and
\[
 G.\mathrm{vol}(W)=\mathrm{baseVol}(B,G.\mathrm{cd},W)
  =\mathrm{trueVol}(B.\mathrm{sp.cm},
          (\tau\mapsto B.\mathrm{host.emb}(\tau)),G.\mathrm{cd},W).
 \tag{GV.33}
\]
GV16 now upgrades this exact original stage from positivity to
actual \texttt{VolTieG}. For every choice of the other stages
$W_0,V,S,N$ over this same $B,C$, the assembled record has the
same host, CM-type, Hermitian spaces, Haar family and volume
function. Therefore (GV.30) holds for every one of those assembled
records. This is the universal volume clause used in
\texttt{GIs2}; it does not require changing the row or local stages.

The same argument applies to any $G.\mathrm{withTrue}$: by its
literal definition only \texttt{vol}, its positivity proof, and
\texttt{gx} are reset, while $B,C,\mathrm{cd}$ and the jet data
remain fixed. A further reset of \texttt{gat} and its size proof
also leaves (GV.33) untouched. Thus the volume clause in
\texttt{exists\_\allowbreak trueVolume\_\allowbreak stage\_\allowbreak of\_\allowbreak any\_\allowbreak column\_\allowbreak choice} is
now supplied. Its other \texttt{ColumnDataSplitPrimes} and
$\Omega$-compatibility clauses still require their exact standard
column, local-family and coefficient data. A dependent layout or
Core on a previous record is not automatically transported across
\texttt{withTrue}. Empty jets in IG12 remain empty, and no nonempty
arithmetic twist is created by this upgrade.

\paragraph{GV20: precisely which real-data fields have been supplied.}
The literal \texttt{RealTied}(D) has four fields:
$\mathrm{cdScale}$, $\mathrm{dsig}$, $\mathrm{dsig\_all}$,
and $\mathrm{vol}$. For the same record $D$, assume
\[
 D.s.\mathrm{cdScale}=\mathrm{trueScale}(D.\mathrm{cd}),\qquad
 D.\mathrm{dsig}(\tau)=\mathrm{trueDsig}(D.R,D.\mathrm{hc}(\tau)).
 \tag{GV.34}
\]
GV17 supplies the first field, and GV16 or GV19 supplies the last
field when $D$ has the actual base-volume normalization.

For the middle fields we retain exactly the source's
\texttt{DiscreteSeriesExist} and \texttt{FormalDegreeFormula\textsubscript{2}}.
The former supplies a discrete-series representation for each
regular Harish--Chandra parameter; the latter identifies its
formal degree, for the Haar measure in (GV.31), with the displayed
Harish--Chandra product. With the supplied parameter laws,
\texttt{trueDsig} is precisely that product. One chosen representation
therefore gives \texttt{IsDsig}; applying the same formula to every
representation and every nonzero vector gives \texttt{dsigTie}.
This proves the exact conditional \texttt{RealOf} join, with
(GV.34) imposed on this $D$. It does not derive representation
existence or a formal-degree formula from the Tamagawa theorem,
and it does not treat the larger \texttt{retie} operation as the
identity: \texttt{retie} also changes \texttt{dsig}, scale and
\texttt{gx}, with their dependent consequences.

\paragraph{GV21: the honest same-object \texttt{V4.Core} consequence.}
Fix the actual $K,X,f,T$, put $D=X.\mathrm{toD7}$, and suppose
there is a layout $Y:\mathrm{Lay}(D)$ with the hypotheses of
\texttt{coreOfLay}. Impose (GV.33)--(GV.34) on this very $X$.
Then GV20 supplies its \texttt{RealTied} argument, and GV13 supplies
the geometric \texttt{VolumeEqMass\textsubscript{2}} consequence
used by the size and formal-degree joins. The constructor keeps
$O=\mathrm{layObj}(Y)$ and all its levels, lattices, comparisons,
theta objects and arithmetic choice. Its fields are verified as
follows: \texttt{RecOK} supplies the column bounds and the
\texttt{Del}/\texttt{etaF} identities; \texttt{LedOK}, together with
$\mathrm{MuIsW.off}$, supplies \texttt{fval} and the actual Gy
bounds; $Y$ supplies the remaining column, theta, gauge and
$\mathrm{NChTie}$ fields; the exact untwisted vanishing clause
completes \texttt{RowLay}; and GV20 supplies \texttt{real}.

The other arguments remain explicit: \texttt{Over}, the same-map
$\xi$ and \texttt{MuIsW}, the forced-prime inclusion,
\texttt{GatNum}, untwisted row vanishing, \texttt{MtStableAt},
the primitive type-character powers, and
\[
   Y.\mathrm{nch}.Q(w)=T.Q(w),\qquad
   Y.\mathrm{nch}.k'(w)=T.k(w)\qquad(w\in W_t).
 \tag{GV.35}
\]
These last identities are conductor-support and exponent identities,
not merely a choice of a prime or a side. This is a strengthened
conditional Core supplier with the geometric volume and scale
arguments proved; it is not a proof that an arbitrary original
stage admits all the other arguments. In particular the forced-prime
field contains the distinguished prime $g_0$, so an empty twist-prime
set cannot satisfy it. The bare empty models of earlier stage
existence results still do not produce a \texttt{V4.Core} or
\texttt{ExTrue}. The conditional Tamagawa and formal-degree inputs
are visible at the final endpoint, and no RH or GRH conclusion is
drawn from these geometric identities alone.


\section{Uniform thinning of the actual centered character sums}
\label{sec:centered-character-sums}

\status{CC1--CC20 reconstruct the Dirichlet analytic inputs and the
uniform tail argument behind \texttt{QH.thin} and \texttt{ThinH}.
The real-character zero exclusion is an explicit hypothesis of the
uniform theorem. For a fixed bounded conductor it follows from QW15;
for a host character involving either selected prime it is supplied
by the original \texttt{RealBox} condition. No assertion that all
such boxes exist is made. The simultaneous-pair zero-density and
relative progression estimates remain separate.}

Put $\delta_\chi=1$ for a principal character and $0$ otherwise,
and write
\[
 L_\chi(s)=L(s,\chi),\quad
 F_\chi(s)=(s-1)^{\delta_\chi}L_\chi(s),\quad
 \psi_0(\chi,u)=\sum_{1\le n\le\lfloor u\rfloor}\Lambda(n)\chi(n)
                                      -\delta_\chi u .       \tag{CC.1}
\]
Characters have their usual zero extension at nonunits. Unless stated
otherwise they are primitive of conductor $N\ge1$. A primitive
principal character has $N=1$. All zero sums count analytic
multiplicity. Constants called absolute are independent of the
modulus, character, height and summation endpoint.

\paragraph{CC1: the primitive functional equation and its normalization.}
Let $a\in\{0,1\}$ satisfy $\chi(-1)=(-1)^a$ and put
\[
 \tau(\chi)=\sum_{r\bmod N}\chi(r)e^{2\pi i r/N},\qquad
 \Theta_\chi(t)=\sum_{n\in\mathbb Z}\chi(n)n^a e^{-\pi n^2t/N}.
\]
For primitive nonprincipal $\chi$ the finite Fourier transform is
$\sum_r\chi(r)e^{2\pi i kr/N}=\overline{\chi(k)}\tau(\chi)$.
For a unit $k$ this follows by substitution. For a nonunit $k$,
primitivity supplies a unit $b$ in the kernel of reduction modulo
$N/(k,N)$ with $\chi(b)\ne1$; otherwise the character factors
through that smaller modulus. Multiplication by $b$ leaves the
additive character fixed and forces the Fourier coefficient to vanish.
The unit-reduction map used here is onto by the Chinese remainder
theorem. Parseval applied to this finite transform gives
$|\tau(\chi)|^2=N$.

Periodize the Schwartz function $x^a e^{-\pi t x^2/N}$ on $N\mathbb Z$.
Its Fourier series and all differentiated series converge absolutely.
The Gaussian Fourier transform, and its first derivative when $a=1$,
therefore give the exact identity
\[
 \Theta_\chi(t)=\epsilon_\chi t^{-a-1/2}\Theta_{\bar\chi}(1/t),
 \qquad \epsilon_\chi=\frac{\tau(\chi)}{i^a\sqrt N},
 \qquad |\epsilon_\chi|=1.                              \tag{CC.2}
\]
For $\Re s>1$, termwise Mellin integration gives
\[
 \mathcal L_\chi(s):=
 (N/\pi)^{(s+a)/2}\Gamma((s+a)/2)L_\chi(s)
 =\frac12\int_0^\infty\Theta_\chi(t)t^{(s+a)/2}\frac{dt}{t}.
                                                               \tag{CC.3}
\]
The theta series decays exponentially at infinity; (CC.2) gives
the corresponding decay at zero. Splitting the integral at one
therefore continues it to an entire function and proves
$\mathcal L_\chi(s)=\epsilon_\chi\mathcal L_{\bar\chi}(1-s)$.
For the primitive principal character use the ordinary theta series
and separate its constant term. The two elementary integrals give
the simple poles at $0,1$ and the usual functional equation for
$\pi^{-s/2}\Gamma(s/2)\zeta(s)$.
Thus the continuation and reflection used below follow from the
actual Gaussian and Gauss-sum normalizations.\footnote{These
normalizations agree with NIST DLMF, Section 25.15,
\url{https://dlmf.nist.gov/25.15}. The argument here supplies the
Poisson--Mellin derivation used in this section.}

\paragraph{CC2: polynomial growth on fixed-width strips and a lower bound on the right.}
There are absolute $C,A>0$ such that, for every character modulo
$N$, not necessarily primitive,
\[
 |F_\chi(s)|\le C(N+1)^A(|s|+2)^A
                         \quad(-2\le\Re s\le6),\qquad
 |F_\chi(2+it)|\ge\zeta(2)^{-1}.                       \tag{CC.4}
\]
Here are sufficient details for the growth assertion. For a
nonprincipal character the partial-sum step function $A_0(x)$ is
$N$-periodic and bounded by $N$. Subtract its mean, take a
periodic primitive, subtract the new mean and repeat. After four
integrations by parts, the initial formula
$L(s)=s\int_1^\infty A_0(x)x^{-s-1}\,dx$ is a polynomial in $s$
with bounded boundary coefficients, plus an integral of the form
\[
 s(s+1)\cdots(s+4)\int_1^\infty P_4(x)x^{-s-5}\,dx,
                  \qquad |P_4(x)|\le C_4N^5.
\]
The removed means contribute polynomial terms because the factor
$s+j$ cancels the denominator from the corresponding constant
integral. The remainder converges absolutely for $\Re s>-4$,
and is bounded by $C_4N^5/(\Re s+4)$. This proves (CC.4), after
increasing $A$. The same integration applied to $\lfloor x\rfloor-x$
proves the bound for $(s-1)\zeta(s)$. A principal character modulo
$N$ adds $\prod_{p\mid N}(1-p^{-s})$, whose absolute value on this
strip is at most $\prod_{p\mid N}2p^2\le N^3$ for $N\ge2$.

In $\Re s>1$, the absolutely convergent Euler product gives
\[
 -\frac{L'_\chi}{L_\chi}(s)
       =\sum_{n\ge1}\frac{\Lambda(n)\chi(n)}{n^s},
 \qquad |L_\chi(2+it)|\ge\prod_p(1+p^{-2})^{-1}
                                                   \ge\zeta(2)^{-1}.
                                                               \tag{CC.5}
\]
The principal factor $|1+it|$ is at least one. These bounds also
establish nonvanishing on the open right half-plane.

\paragraph{CC3: local zero counts and logarithmic derivatives.}
Put $\lambda_N=\max(1,\log N)$ and
$\ell_N(t)=\lambda_N+\log(|t|+3)$. Apply Jensen's formula to
$F_\chi$ in the disk of radius four centered at $2+it$, using
(CC.4). The number of zeros in the concentric radius-$13/4$ disk,
with multiplicity, is at most $C\ell_N(t)$.
Removing those zeros by the finite Blaschke construction of QR15
gives, uniformly for $|s-(2+it)|\le5/2$ away from zeros,
\[
 \frac{F'_\chi}{F_\chi}(s)
       =\sum_{|\rho-(2+it)|\le13/4}\frac{m_\rho}{s-\rho}
                                             +O(\ell_N(t)).    \tag{CC.6}
\]
The regular term is bounded by differentiating the logarithm of
the zero-free quotient on the smaller disk. The gaps between the
three radii are fixed, so the implied constant is absolute.
The proof is exactly the rescaled disk argument of QR15; no
global product for the completed function is needed here.

In particular, the number of nontrivial zeros in any height interval
$[t-1,t+1]$ is at most $C\ell_N(t)$: their real parts lie in
$[0,1]$, so their distances from the disk center are at most
$\sqrt5<5/2$. This proves the actual height-radius-one
\texttt{ZerosWindow} estimate, after enlarging the constant to
replace $\ell_N(t)$ by $\log(N(|t|+2))$.
For the near-one arguments we can instead use the center
$5/4+it$ and the radii $7/10,3/5,1/2$ of QR16. The corresponding
growth and Euler-product lower bounds follow from the same
partial-sum formula on $\Re s\ge11/20$. Thus every local input
of QR17--QR18 is available with $\log B$ replaced by a fixed
multiple of $\lambda_N$.

\paragraph{CC4: the boundary line and the real principal character.}
QW14 proves $L_\chi(1)\ne0$ for every nonprincipal character.
For $0<\sigma<1$, the alternating series
$\eta(\sigma)=\sum_{n\ge1}(-1)^{n-1}n^{-\sigma}$ is positive:
pair consecutive terms and then pass to the convergent limit.
Continuation of
$\eta(s)=(1-2^{1-s})\zeta(s)$ gives $\zeta(\sigma)<0$.
The finite principal Euler factors are positive there. Hence
a principal character has no real zero in $(0,1)$.

There are also no zeros at $1+it$, $t\ne0$. On the right of
that line use the coefficientwise nonnegativity
\[
 3+4\cos\theta+\cos2\theta=2(1+\cos\theta)^2\ge0.
                                                               \tag{CC.7}
\]
For primes dividing the modulus all three character coefficients
are zero; on units (CC.7) applies. If $L_\chi$ had a zero of
order $m\ge1$ at $1+it$, the resulting logarithmic-derivative
inequality at $1+\delta$ would have right-hand singular part
at most $(3-4m)/\delta$. A zero of $L_{\chi^2}$ at $1+2it$
only decreases that upper bound; it has no pole there because
$t\ne0$. All remaining terms are bounded for this fixed
modulus and height, by local factorization. Letting $\delta$
decrease to zero is impossible. Together with (CC.5) this puts
every nontrivial zero in the open critical strip.
For a primitive character, (CC.2)--(CC.3) and the absence of
zeros at one also identify the possible zero at zero: it is
simple precisely for nonprincipal even characters. There are
no other zeros in $-1/2\le\Re s\le0$.

\paragraph{CC5: the uniform zero-free region with its exact exception.}
For all characters modulo $N$ there is an absolute $b>0$ such
that a zero in
\[
           \Re s>1-b/\ell_N(\operatorname{Im}s)           \tag{CC.8}
\]
is real, belongs to a real nonprincipal character, and is simple;
across that modulus there is at most one such zero.
Here the analytic inequalities from QR17--QR18 apply without
the arithmetic-family axioms. We spell out their inputs so
that this reuse does not assume a Dirichlet instance of \texttt{LFam}.
Formula (CC.5) supplies the logarithmic derivative. Formula
(CC.6) supplies the local zero expansion and its error.
All nontrivial zeros have real part at most one by CC4,
so omitted zero terms have nonnegative real part to the right
of one. The principal pole contributes exactly $1/(s-1)$.

For a complex character, (CC.7) gives the three-character
inequality used in QR17. For a real character, conjugation
pairs its nonreal zeros and gives the two small-height terms;
at larger heights the principal reciprocal terms are bounded
as in that paragraph. The two numerical contradictions are
the same as (QR.35)--(QR.38), with a new common error constant
$C$. Choosing $b$ smaller than the resulting fixed reciprocal
multiple of $C$ proves the assertion except uniqueness.
For distinct real characters use
$(1+\chi(n))(1+\chi'(n))\ge0$ on units: its four terms
contain only one principal character. Two zeros inside the
chosen width contradict the same near-one estimate. The
one-character two-term inequality excludes multiplicity two.
These are the computations of QR18; their constants depend
only on the local error constant. The real principal case
was excluded separately in CC4. To match the source's closed strip
and denominator $\log(N(|t|+2))$, shrink $b$ by a fixed factor:
$\ell_N(t)$ and this logarithm are comparable uniformly for $N\ge1$,
and the narrower closed strip lies inside (CC.8).

Fix $c_0>0$. If a real nonprincipal $\chi$ has no real zero in
$(1-c_0,1)$, (CC.8) therefore gives, for some $c=c(c_0)>0$,
\[
 0<\beta<1,\quad L_\chi(\beta+i\gamma)=0
 \quad\Longrightarrow\quad
       \beta\le1-\frac{c}{\lambda+\log(|\gamma|+3)}
 \qquad(\lambda\ge1,\ \log N\le\lambda).               \tag{CC.9}
\]
Take, for example, $c\le\min(b,c_0,1)$ after the fixed denominator
comparison. For a nonreal or principal character the exclusion
hypothesis is vacuous. Zeros with $\Re\rho\le0$ satisfy the same
upper bound automatically because $c\le1$ and $\lambda\ge1$;
CC4 excludes zeros with $\Re\rho\ge1$. Thus this supplies the
literal every-zero form of \texttt{pt\_\allowbreak zero\_\allowbreak free}, with the
exceptional-zero hypothesis retained.

\paragraph{CC6: reflect zeros to control small denominators.}
Primitivity is preserved by inversion: factoring $\chi^{-1}$
through a smaller modulus would factor $\chi$ through it.
If $0<\Re\rho<1$ and $L_\chi(\rho)=0$, the gamma factors in
(CC.3) are finite and nonzero there, so
\[
                     L_{\chi^{-1}}(1-\rho)=0.           \tag{CC.10}
\]
The real-character exclusion passes to $\chi^{-1}$, because
a character of square one equals its inverse. Applying (CC.9)
to (CC.10), for $|\Im\rho|\le2$, gives
\[
          \Re\rho\ge\frac{c}{\lambda+\log5}
                         \ge\frac{c}{3\lambda}.         \tag{CC.11}
\]
Thus a zero near zero is controlled by the same exclusion as
a zero near one. This reflection is essential for the
$1/|\rho|$ terms in the explicit formula.

\paragraph{CC7: sum the reciprocal zero weights.}
Let $Z_\chi(T)=\{\rho:0<\Re\rho<1,\ |\Im\rho|<T,
L_\chi(\rho)=0\}$. It is finite by analyticity and CC4.
Put $\kappa=c/3>0$ in (CC.11).
The unit-height count of CC3 then implies
\[
 \sum_{\rho\in Z_\chi(T)}\frac{m_\rho}{|\rho|}
          \le C_2(\lambda+\log T)^2\qquad(T\ge1),       \tag{CC.12}
\]
where $C_2$ depends only on $c_0$.
Indeed the zeros with $|\gamma|\le2$ lie in a fixed number
of unit intervals; their total multiplicity is $O(\lambda)$,
and $1/|\rho|\le\lambda/\kappa$. For $j\le|\gamma|<j+1$,
$j\ge2$, the weight is at most $1/j$ and the count is
$O(\lambda+\log(j+3))$. Summation gives
$O((1+\kappa^{-1})(\lambda+\log T)^2)$, by the elementary
harmonic-sum bounds. One can choose an absolute multiple of
$C(1+\kappa^{-1})$ for $C_2$. This coarser constant is sufficient
for the existential \texttt{psi\_\allowbreak thin} endpoint; no particular
numerical value from the longer source estimate is asserted.

\paragraph{CC8: the literal explicit formula, including the pole at zero.}
There is an absolute $C_7>0$ such that for every primitive
$\chi$, $2\le T\le u$,
\[
 \left|\psi_0(\chi,u)+
       \sum_{\rho\in Z_\chi(T)}m_\rho\frac{u^\rho}{\rho}\right|
 \le C_7\left\{
 \frac{u\log^2(NuT)}T+\log(Nu)
       +\sum_{\rho\in Z_\chi(1)}\frac{m_\rho}{|\rho|}\right\}.
                                                               \tag{CC.13}
\]
We give the contour details behind this source interface.
Apply QR20 to $\Lambda(n)\chi(n)$ on
$\Re s=1+1/\log u$. Its endpoint error is
$O(u\log^2u/T+\log u)$, including integer $u$.

Choose upper and lower heights separately in $[T,T+1]$.
The CC3 count in a fixed-width neighborhood of either height
is $O(\ell_N(T))$. Deleting intervals of radius
$c/\ell_N(T)$ about their ordinates, for a sufficiently small
absolute $c$, removes less than half of $[T,T+1]$.
At a remaining height, (CC.6) bounds $L'/L$ by
$O(\ell_N(T)^2)$ throughout $-1/2\le\Re s\le3$.
Extending the original right-edge integral by at most one unit
at each end costs $O(u\log^2u/T)$, from (CC.5) and
$\sum_{n\ge2}(\log n)n^{-1-1/\log u}=O(\log^2u)$.
Both horizontal edges therefore cost
$O(u\log^2(NuT)/T)$.

On $\Re s=-1/2$, logarithmic differentiation of (CC.3) bounds
$L'/L$ by $O(\log(N(|t|+2)))$: the reflected Dirichlet series
has real part $3/2$, and the two gamma logarithmic derivatives
are $O(\log(|t|+2))$. For completeness that gamma bound follows
from
\[
 \frac{\Gamma'}\Gamma(z)=-\gamma+
       \sum_{j\ge0}\left(\frac1{j+1}-\frac1{j+z}\right).
\]
On a right half-plane, split the series at $2|z|+2$; the initial
part is $O(\log(|z|+2))$ and the tail is $O(1)$.
The recurrence shifts the two fixed real parts here into that
half-plane, a bounded distance from poles. The left edge costs
$O(u^{-1/2}\log^2(N(T+2)))$, absorbed in the displayed error
because $T\le u$.

The principal pole contributes $u$. A nontrivial zero contributes
$-m_\rho u^\rho/\rho$. At zero, write
$-L'/L(s)=-e_0/s+b_0+O(s)$, where $e_0=1$ for a nonprincipal
even character and $0$ otherwise, by CC4. Its residue after
multiplication by $u^s/s$ is $b_0-e_0\log u$.
Formula (CC.6) at zero, after removing the zero there, gives
\[
 |b_0|\le C\log(N+2)+
                  \sum_{\rho\in Z_\chi(1)}m_\rho/|\rho|.
\]
Zeros of larger height in that fixed disk have reciprocal
absolute value at most one and total multiplicity $O(\log(N+2))$;
the possible negative trivial zero in the disk is also a bounded
distance from zero. Thus no unbounded constant at zero is omitted.
Finally replace the good heights by the strict cutoff $|\gamma|<T$.
The added or removed zeros lie in unit intervals near $\pm T$;
each costs at most $m_\rho u/T$, so their total is absorbed.
The residue theorem proves (CC.13), with exactly the source
\texttt{ExplicitFormulaWith} convention.

\paragraph{CC9: choose the height before the conductor.}
Assume the real exclusion in CC5, $\lambda\ge1$,
$\log N\le\lambda$, and $u\ge e$. Put
\[
 v=\log u,\quad s_0=\sqrt v,\quad T=e^{s_0},\quad
 c'=\min(c/3,1),\quad
 E=\exp\left(-\frac{c'v}{\lambda+\sqrt v}\right).
\]
Then $2\le T\le u$, $T^{-1}\le E$, and $uE\ge1$.
For a zero below $T$,
$\lambda+\log(|\gamma|+3)\le\lambda+\sqrt v+2
\le3(\lambda+\sqrt v)$.
Hence (CC.9) gives $u^{\Re\rho}\le uE$.
Using (CC.12), its complete zero sum is bounded by
$C_2uE(\lambda+\sqrt v)^2$.

In (CC.13), $\log(NuT)\le2(\lambda+v)$.
The first error is at most $4uE(\lambda+v)^2$;
the logarithm is at most $uE(\lambda+v)^2$; and the low zero
sum is at most $C_2uE(\lambda+v)^2$.
Consequently
\[
 |\psi_0(\chi,u)|\le
 [C_7(5+C_2)+C_2]u(\log u+\lambda)^2
             \exp\left(-\frac{c'\log u}
                              {\lambda+\sqrt{\log u}}\right).
                                                               \tag{CC.14}
\]
This proves the actual conductor-sensitive estimate by a fixed
height choice. There is no premise of a zero-free critical strip.

\paragraph{CC10: the uniform primitive \texttt{psi\_\allowbreak thin} theorem.}
Collect CC5--CC9. For every $c_0>0$ there are $C,c>0$ such that
for all primitive $\chi$ of conductor $N$, all
$\lambda\ge\max(1,\log N)$, and all $u\ge e$,
\[
 |\psi_0(\chi,u)|\le Cu(\log u+\lambda)^2
                  e^{-c\log u/(\lambda+\sqrt{\log u})},        \tag{CC.15}
\]
provided
\[
 \chi^2=1,\ \chi\ne1\quad\Longrightarrow\quad
     L_\chi(\beta)\ne0\quad(1-c_0<\beta<1).             \tag{CC.16}
\]
The constants are chosen before $N,\chi,\lambda,u$.
This supplies \texttt{QH.psi\_\allowbreak thin}, including its principal
primitive case, and its explicit-formula, zero-window and
zero-free inputs have been established above.

\paragraph{CC11: fixed-character convergence needs no uniform Siegel bound.}
For any fixed modulus, QW15 supplies a positive real-zero
margin for all its real nonprincipal characters. Apply CC10
to the finitely many primitive characters of conductor dividing
that modulus. For each fixed one this gives
\[
                \psi_0(\chi,u)=O_\chi(u/\log^2u).       \tag{CC.17}
\]
Indeed, once $\sqrt{\log u}\ge\lambda$, the exponential
in (CC.15) is at most $e^{-(c/2)\sqrt{\log u}}$, which absorbs
any prescribed power of $\log u$. An imprimitive character
differs from its primitive inducer only at powers of primes
dividing the larger modulus. Their total contribution is
$O_N(\log u)$, so (CC.17) holds for it as well.
The function $\psi_0$ is a locally bounded step function minus
a linear function. Thus
\[
                   \int_1^\infty
                       |\psi_0(\chi,u)|u^{-2}\,du<\infty.
                                                               \tag{CC.18}
\]
This is a fixed-modulus convergence assertion. It makes no
uniform lower-bound assertion as the modulus increases.

\paragraph{CC12: identify the actual nonprincipal centering constant.}
For a nonprincipal character put $a_n=\Lambda(n)\chi(n)$
and $A(u)=\sum_{n\le u}a_n$. Abel summation gives
\[
 \sum_{n\le U}\frac{a_n}{n}
             =\frac{A(U)}U+\int_1^U\frac{A(u)}{u^2}\,du.
                                                               \tag{CC.19}
\]
Equation (CC.17) makes the first term tend to zero and makes
the integral absolutely convergent. For real $\sigma>1$,
the same calculation gives
$-L'/L(\sigma)=\sigma\int_1^\infty A(u)u^{-\sigma-1}\,du$.
Dominated convergence as $\sigma\downarrow1$, justified by
(CC.18), and QW14 nonvanishing identify the limit:
\[
 S_\chi:=\sum_{n\ge1}\frac{\Lambda(n)\chi(n)}n
                         =-L'_\chi(1)/L_\chi(1).        \tag{CC.20}
\]
The series is asserted convergent, not absolutely convergent.
Its value is the original logarithmic-derivative centering.
Subtract (CC.19) at $y$ from its limit and rearrange to obtain,
for every real $y\ge1$,
\[
 G_\chi(y):=\sum_{n\le y}a_n(y/n-1)-yS_\chi
                       =-y\int_y^\infty A(u)u^{-2}\,du.
                                                               \tag{CC.21}
\]
This also verifies the floor convention: if $y$ is an integer,
the newly included summand has coefficient zero.

\paragraph{CC13: the principal centering is exactly $-1-\gamma$.}
Let $R(u)=\sum_{n\le u}\Lambda(n)-u$. From (CC.18) for
the primitive principal character, $R(u)/u^2$ is integrable.
For $\sigma>1$,
\[
 \sigma\int_1^\infty R(u)u^{-\sigma-1}\,du
        =-\frac{\zeta'}\zeta(\sigma)-\frac{\sigma}{\sigma-1}.
\]
The Laurent expansion
$\zeta(s)=1/(s-1)+\gamma+O(s-1)$ gives the limit $-1-\gamma$.
This constant also follows directly from the sum-minus-integral
definition $\gamma=\lim_n(\sum_{j=1}^n1/j-\log n)$ in the
Euler summation formula for $\zeta$. Dominated convergence yields
\[
                      \int_1^\infty R(u)u^{-2}\,du=-1-\gamma.
                                                               \tag{CC.22}
\]
The finite integral is
$\int_1^y R(u)u^{-2}\,du
=\sum_{n\le y}\Lambda(n)(1/n-1/y)-\log y$.
Subtract it from (CC.22) to get
\[
 G_1(y):=\sum_{n\le y}\Lambda(n)(y/n-1)
                    -y(\log y-1-\gamma)
                       =-y\int_y^\infty R(u)u^{-2}\,du.
                                                               \tag{CC.23}
\]
Thus the principal and nonprincipal constants have both been
identified, without borrowing a uniformly quantified
Siegel--Walfisz estimate to justify convergence.

\paragraph{CC14: match the literal centered functions.}
The source $G_0(\chi,y)$ is $G_1(y)$ when $\chi$ is principal
and $G_\chi(y)$ otherwise. For primitive characters (CC.21)
and (CC.23) therefore give exactly
\[
             G_0(\chi,y)=-y\int_y^\infty
                           \psi_0(\chi,u)u^{-2}\,du.    \tag{CC.24}
\]
The integral exists absolutely. Primitivity in the principal
case is essential: it makes $N=1$, so its von Mangoldt sum
has no deleted Euler factors.
The host function \texttt{An.gChi} uses
$\max(y/n-1,0)$ on $1\le n\le\lfloor y\rfloor$.
On this finite interval $n\le y$, so the maximum is $y/n-1$.
Its nonprincipal centering has the sign $+yL'/L(1)$ from
(CC.20), and its principal term is precisely (CC.23).
Hence it equals $G_0$ on these actual arguments. This supplies
\texttt{QH.g\_\allowbreak integral} and the function-identification clause
used by \texttt{thinH}.

\paragraph{CC15: a uniform integral estimate above the exponential cutoff.}
Fix $\epsilon,c,C,A,c'>0$. Suppose
\[
 y\ge\tfrac13e^{x^{\epsilon/3}},\qquad
 1\le\lambda\le Ax^{\epsilon/5}\log^2x,
\]
and a measurable function $\phi$ satisfies the bound on the
right side of (CC.15) for every $u\ge y$.
For all sufficiently large $x$, independently of $\lambda,y,\phi$,
\[
 \left|y\int_y^\infty\phi(u)u^{-2}\,du\right|
                          \le y e^{-c'x^{\epsilon/8}}.   \tag{CC.25}
\]
Here is a short proof with the uniform quantifiers visible.
Put $V=\log y\ge\tfrac12x^{\epsilon/3}$.
Since $61/300>1/5$, the power $x^{\epsilon/300}$ eventually
dominates $\log^2x$ and the fixed constant $A$.
Consequently, for every $v\ge V$,
$\lambda\le v^{61/100}$ and $v\ge1$. Therefore
\[
 (v+\lambda)^2e^{-cv/(\lambda+\sqrt v)}
                    \le4v^2e^{-(c/2)v^{39/100}}.        \tag{CC.26}
\]
The substitution $u=e^v$ in the norm majorant reduces its
integral to the integral of the left side of (CC.26).
Factor half the exponential from the right side at $V$:
\[
 \int_V^\infty4v^2e^{-(c/2)v^{39/100}}\,dv
 \le K_c e^{-(c/4)V^{39/100}},\qquad
 K_c=4\int_0^\infty v^2e^{-(c/4)v^{39/100}}\,dv<\infty.
\]
The last integral is finite by $w=v^{39/100}$, or directly
by exponential decay. The exponent at $V$ is at least a
positive constant times $x^{13\epsilon/100}$.
Since $13/100-1/8=1/200>0$, this absorbs $CK_c$ and any
prescribed $c'x^{\epsilon/8}$. This proves (CC.25).
Measurability and the displayed integrable majorant also prove
the claimed absolute integrability. The source's more general
totalized integral convention assigns zero to a nonintegrable
integrand; that unused case satisfies its numerical inequality
trivially. For $\phi=\psi_0$, measurability was proved in CC11.

\paragraph{CC16: the exact uniform thinning theorem.}
Fix $\epsilon,A,c',c_0>0$. There is $x_0$ such that for every
$x\ge x_0$, every primitive character $\chi$ with
\[
       \log N\le Ax^{\epsilon/5}\log^2x,
\]
satisfying (CC.16), and every
$y\ge\tfrac13e^{x^{\epsilon/3}}$, one has
\[
                        |G_0(\chi,y)|
                                \le y e^{-c'x^{\epsilon/8}}.
                                                               \tag{CC.27}
\]
Choose $\lambda=\max(1,\log N)$. For large $x$ it is at most
$(A+1)x^{\epsilon/5}\log^2x$. Apply CC10 with its constants
chosen only from $c_0$, then CC15 and the exact identity CC14.
Increasing $x_0$ ensures $y\ge e$ and all displayed thresholds.
The resulting threshold depends only on the four fixed
parameters. This is \texttt{QH.thin}, including its universal
quantifier over every $y$ above the cutoff.

\paragraph{CC17: a primitive real character has restricted conductor exponents.}
If a primitive character has square one, every odd prime
occurs in its conductor to exponent at most one, and the
exponent at two is at most three.
For an odd prime $p$, the kernel of
$(\mathbb Z/p^k\mathbb Z)^\times\to(\mathbb Z/p\mathbb Z)^\times$
is a finite $p$-group. A homomorphism from it to $\{\pm1\}$
is trivial, since its image order divides both a power of $p$
and two. At two, every unit congruent to one modulo eight
is a square modulo $2^k$. To prove the latter, lift an odd
square root from modulus $2^j$ to $2^{j+1}$ by replacing
$z$ with $z+2^{j-1}$ if necessary; for $j\ge3$ this flips
exactly the next binary digit of $z^2$.
A square-one character kills these squares. Chinese remainders
then make it factor through $8\operatorname{rad}(N_{\rm odd})$.
Primitivity proves the asserted exponent bounds.

Consequently, if $8\mid M$, $q,d$ are primes, and every prime
dividing $N$ divides $M$ or equals $q$ or $d$, then
\[
                                  N\mid Mqd.             \tag{CC.28}
\]
This is the actual \texttt{real\_\allowbreak cond} assertion, with its
prime-support hypotheses. If $q,d\nmid N$, the stronger
$N\mid M$ follows immediately. No exponent bound is imposed
on a general complex character.

\paragraph{CC18: bounded conductors have one common real-zero exclusion interval.}
For each fixed integer $B_0\ge0$, there is $c_s>0$ such that
every nonprincipal character modulo $N\le B_0$
has no real zero in $(1-c_s,1)$.
For each such character QW14 makes its analytic function
nonzero at one; continuity gives a neighborhood of one free
of zeros. There are finitely many moduli and characters, so
take the minimum of finitely many positive radii. When the
set is empty, take any positive radius. This is the same
finite-minimum argument as QW15, now for all nonprincipal
characters with the precise \texttt{small\_\allowbreak zeros} quantifier;
primitivity and reality are not needed in this bounded-modulus
statement.
Its dependence on $B_0$ is allowed; it is not a uniform
Siegel-type bound for unbounded conductors.

\paragraph{CC19: install the actual host thinning statement.}
Fix $B\in\mathbb R$ and $0<\epsilon<1/52$, as in the
literal \texttt{ThinH} quantifier. Let $H$ be a supplied host,
and write $Q_S$ for its fixed
cyclotomic factor and $q,d$ for its two prime factors.
The actual conductor is $Q_Sqd$. A primitive host character
$\chi$ has $N\mid Q_Sqd$: it is the primitive inducer of
the character associated to its host index. This is the
same conductor compatibility already used in HC9.
Assume
\[
 Q_{S,\rm odd}\le B,\qquad
 \log(Q_Sqd)\le2x^{\epsilon/5}\log^2x,
\]
and the source \texttt{RealBox}: each real primitive host
character involving $q$ or $d$ is nonzero for real
$\beta>1-1/400$.

Apply CC18 with $B_0=8\lfloor B\rfloor$ and put
$c_0=\min(c_s,1/400)>0$.
For a real nonprincipal host character, if $q\mid N$ or
$d\mid N$, the given \texttt{RealBox} supplies (CC.16).
Otherwise CC17 and the prime support of $Q_Sqd$ give
$N\mid8Q_{S,\rm odd}$ and hence $N\le8\lfloor B\rfloor$;
CC18 supplies (CC.16). Nonreal and principal characters
need no such exclusion. Every host character also satisfies
$\log N\le\log(Q_Sqd)$.
Now apply CC16 with $A=2$ and $c'=1$, and identify its
function with \texttt{An.gChi} by CC14. This proves the
literal \texttt{ThinH}, with its outer constant chosen as one.
The threshold depends on $B,\epsilon$, and the statement is
uniform over all supplied hosts, primitive host characters
and $y$ above the original exponential cutoff.
Thus the thinning premise of HC9 is supplied. The construction
of a host satisfying \texttt{RealBox}, simultaneous
\texttt{PairEx}, the relative progression theorem and the
cohomological windows are not inferred from this implication.

\paragraph{CC20: a quantitative derivative consequence for the same centered sums.}
The proof also gives a useful strengthening above the same
exponential cutoff. Under the hypotheses of CC16, for any
fixed $b>0$ and all sufficiently large $x$,
\[
 |G_0(\chi,y_2)-G_0(\chi,y_1)|
       \le (y_2-y_1)e^{-b x^{\epsilon/8}}
 \quad\left(\tfrac13e^{x^{\epsilon/3}}\le y_1\le y_2\right).
                                                               \tag{CC.29}
\]
Away from integers, differentiate (CC.24):
\[
           G'_0(\chi,y)=\psi_0(\chi,y)/y+G_0(\chi,y)/y.
\]
The same identity holds for the right derivative at integers.
Each ramp is continuous there, so $G_0$ is locally absolutely
continuous and piecewise continuously differentiable.
The argument of CC15 applied pointwise to (CC.15) bounds
$|\psi_0(\chi,y)|/y$ by $e^{-2b x^{\epsilon/8}}$, uniformly
over all $y$ above the cutoff, after increasing the threshold.
Use CC16 with $c'=2b$ for the second term. For large $x$,
their sum is at most $e^{-b x^{\epsilon/8}}$.
Integration on $[y_1,y_2]$ proves (CC.29), including intervals
crossing integers. This is an additional human consequence of
the proved bounds; it does not substitute for the separate
square-root estimates on the earlier polynomial-height window.



\section{Actual form and signature spaces, with the initial theta fields explicit}
\label{sec:automorphic-realization}

\status{MR1--MR11 construct the actual definite adelic realization,
including its full range, prime-local integral Hecke action and the
pointwise \texttt{MtAd} comparison on a supplied compatible geometric
state. MR12--MR22 construct the integral form and signature portion of
\texttt{SideB} over the already chosen coefficient rings. The source's
\texttt{mkF/mkEx} initializes separate theta fields as zero and still
requires a theta witness; MR22 states this boundary literally.
MR23--MR26 prove four local Hecke-algebra and relation-lattice
primitives on the original objects. The good-row spectral input,
automorphic Galois input and compatible geometric state remain
explicit. No compatible nonempty theta/local realization or global
branch conclusion follows from these initial fields.}

Write $L=D.\mathrm{what}_2.\mathrm{LL}$, $K=D.\mathrm{LL0}\subset L$,
and $\epsilon:L\hookrightarrow\mathbb C$ for the fixed coefficient
embedding. A rational form datum has a number field $K_s\subset K$,
an $L$-space $W_L$ and a $K_s$-form $W_{K_s}\subset W_L$; ``form'' here
includes the assertion $L\otimes_{K_s}W_{K_s}\simeq W_L$.
The letter $E$ denotes the host CM field and $F=E^+$.
All the old choices of host, rows, $L$, $K$, embeddings, tuple maps and
level characters remain fixed unless an explicitly specified
\texttt{setTw} operation is made.

We use the preceding finite-class and coefficient constructions
CR1--CR6, FC/TX, AR5--AR9 and AR15--AR18 where their hypotheses hold.
For coefficient representations we retain the standard characteristic
zero highest-weight theorem in its precise polynomial form: the
Schur module of a dominant weight is irreducible, its indicated
highest-weight line is one-dimensional, and its dual has the opposite
contragredient highest weight. Tensor products at the distinct,
nonconjugate embeddings of a CM type have the corresponding product
property. This is a representation-theoretic input, not an assertion
about the existence of automorphic forms or Galois representations.
The literal source interfaces are indexed in the accompanying ledger.

\paragraph{MR1: finite class values reconstruct all adelic forms.}
Let $G$ be the finite adelic group of a positive definite unitary
space, $\Gamma$ its rational group, $J$ a compact open level, and
$\chi:J\to L^\times$ the actual finite-order level character. Fix
representatives $y_c$ of the finite set $\Gamma\backslash G/J$.
Let $V_L$ be the coefficient module, with complex realization $V$ and
rational action $\rho$. For each $c$, let $V_{L,c}$ consist of the
vectors satisfying
\[
 \rho(\gamma)v=\chi(k)v
 \quad\text{whenever }\gamma y_c=y_c k,
 \qquad \mathcal C_L=\bigoplus_c V_{L,c}.             \tag{MR.1}
\]
The sign in (MR.1) follows from the convention
$F(\gamma gk)=\rho(\gamma)F(g)\chi(k)$ and the displayed
stabilizer equation.
For $v=(v_c)$ define
\[
       \operatorname{Ext}(v)(\gamma y_c k)
          =\rho(\gamma)\epsilon(\chi(k))v_c.          \tag{MR.2}
\]
If two decompositions are used, their quotient is one of the
stabilizer pairs in (MR.1), so their values agree. Conversely, every
equivariant adelic function has precisely these stabilizer relations
at $y_c$, and (MR.2) recovers it. Evaluation and extension are inverse.
The same argument over $\mathbb C$ proves surjectivity onto the full
adelic form space, not merely onto the span of a selected collection
of forms. Local constancy follows from the open kernel of $\chi$.

The rational stabilizers are finite: they preserve a lattice, and
their matrices form a discrete subset of the product of the compact
unitary groups at all real places. Their fixed-vector equations are therefore finite linear
equations over $L$. Row reduction, or SE2, shows that their kernels
commute with $L\hookrightarrow\mathbb C$. Thus a basis of
$\mathcal C_L$ extends to a complex basis of all adelic forms.
With the actual coefficient lattices $\nabla(\Lambda_c)$, integrality
is exactly $v_c\in\nabla(\Lambda_c)$ at every class. This supplies the
evaluation, extension and scalar-extension facts used from
\texttt{ClassPres}; finiteness and the specified geometric and
coefficient-lattice presentation are still required inputs.

\paragraph{MR2: the finite level projector and its rational basis.}
Suppose the actual larger level acts on $\mathcal C_L$ through a
finite group $G_0$, with representation $a$, and let $P_0\le G_0$.
For a character $\psi:P_0\to L^\times$, put
\[
 e_\psi=\frac1{|P_0|}\sum_{p\in P_0}\psi(p)^{-1}a(p).
                                                               \tag{MR.3}
\]
Changing the variable $p$ to $q^{-1}p$ gives
$a(q)e_\psi=\psi(q)e_\psi$. If $a(q)v=\psi(q)v$ for all $q$, then
$e_\psi v=v$. Hence $e_\psi^2=e_\psi$ and its image is exactly the
specified character subspace. Choose independent columns of its
matrix and assemble them into $E_\psi:L^d\to\mathcal C_L$.
Gaussian elimination gives a left inverse over $L$ and proves that
the columns are a basis of its image.

Apply \texttt{Ext} to these columns. They remain independent after
applying $\epsilon$, because the left inverse also extends to
$\mathbb C$. Every complex form satisfying the extra level character
is fixed by the corresponding averaging operator. Average its
expansion in the full basis from MR1; the result belongs to the span
of the columns of $E_\psi$. Thus these columns span the entire cut
space. If its values at the $y_c$ belong to the rational coefficient
module, MR1 first descends the form to $\mathcal C_L$, and (MR.3)
then descends it to $E_\psi(L^d)$. Finally,
\[
 x\in\mathcal N_\psi
 \quad\Longleftrightarrow\quad
       (E_\psi x)_c\in\nabla(\Lambda_c)\quad\text{for every }c
                                                               \tag{MR.4}
\]
is the actual level lattice. No division by $|P_0|$ is asserted to
preserve this lattice; the lattice is the inverse image in (MR.4).

\paragraph{MR3: the standard-frame type space and the inverse character.}
For the source's identity Gram matrix and a unit frame $P$, changing
the tensor order gives an isomorphism $t_P:G^{\rm std}(\mathbb A_f)
\simeq G$. On rational points it is the same matrix over $E$:
injectivity of $E\to E\otimes_F\mathbb A_{F,f}$ proves both
directions of the rational-point assertion. There is therefore a
compatible rational-group isomorphism, and no approximation of a
rational element is needed.

Let $Q_0$ be the larger parabolic level, let
$r:Q_0\twoheadrightarrow G_0$ be the prescribed reduction, and let
$e_G:G_0\simeq G_0(P_g,r)$ be its actual identification with the
product residue group. The standard level is precisely
\[
 J_{\rm std}=t_P^{-1}\{q\in Q_0:r(q)\in P_0\},\qquad
 \chi_{\rm std}(g)=\operatorname{ch}(t_Pg)
                       \epsilon(\psi_0(r(t_Pg))^{-1}).\tag{MR.5}
\]
Indeed, membership says that the reduced matrix preserves the
specified flags and that its upper-left blocks lie in $e_G(P_0)$.
The source uses the conjugate level character in \texttt{levFn},
which is the inverse in (MR.5). Substitution in the two equivariance
laws now proves, in both directions,
\[
 F\in\operatorname{cutSp}(\psi_0^{-1})
 \quad\Longleftrightarrow\quad
 F\circ t_P\in\operatorname{TypeSp}
                      (\mathrm{wtRep},J_{\rm std},\chi_{\rm std}).
                                                               \tag{MR.6}
\]
The smaller original level lies in $\ker r$, so its character is
also recovered by (MR.5). Thus no additional equivariance condition
is lost in the reverse implication. MR2 is applied with
$\psi=\psi_0^{-1}$, exactly as in \texttt{SideIn.fn}.

\paragraph{MR4: genuine product levels and the local descriptions.}
The standard level in MR3 is bounded by the integral maximal compact.
Take $N$ to be the product of the norms of all primes in $P_g\cup Q$.
It is positive and belongs to each of those prime ideals. A matrix
congruent to $1$ modulo $N$ reduces to $1$ at every level prime.
It therefore belongs to $J_{\rm std}$ and its character is $1$.
This proves the literal open-level condition. At a real-field prime
below none of $P_g\cup Q$, the local group is the full integral
unitary group and its level character is trivial; this is
\texttt{HypOutside}.

For completeness, the product assertion can be proved inside the
adelic tensor algebra. Let $e_v$ be $1$ at the place $v$ and $0$
elsewhere, and put
\[
             g_v=e_vg+(1-e_v)1.                      \tag{MR.7}
\]
Since $e_v^2=e_v=\bar e_v$, the unitary equation for $g$ implies
the unitary equation for $g_v$; its inverse is obtained from $g^{-1}$
by the same formula. Both integrality and each flag condition are
tested componentwise. The two complementary cuts multiply to $g$,
and their characters multiply to the character of $g$.
At a conductor prime the reduction condition is the actual
$(r,n-r)$ parabolic, with character the prescribed component on the
determinant of the complementary block. At a type prime the
upper-left block is the specified unipotent group, and the
superdiagonal sum is read by the actual nontrivial residue additive
character. Its completion character is obtained by extending the
residue character from local integers, as proved in AR6; AR7--AR8
give the exact tuple component and sign. AR5 gives the conductor
character. These identifications require the literal half-selection,
split-place and generic-character hypotheses of \texttt{LevCh.Is};
productness alone does not supply them.

\paragraph{MR5: the actual algebraic coefficient representation.}
Let $c_0:VV_L(\lambda)\to V$ be the given semilinear coefficient
map, with complex-independent basis images spanning $V$ and with
the actual rational-action identity. Define
$T:VV_{\mathbb C}(\lambda)\simeq V$ by sending the standard basis
to those images. On this basis, and hence everywhere,
\[
       T\,\rho_{VV}(g)=\rho_V(g)T.                   \tag{MR.8}
\]
The inverse-transpose matrix convention is preserved: scalar
extension commutes with matrix inverse because the two inverse
identities remain true after applying a field homomorphism.
The polynomial matrices of $VV$ therefore give polynomial matrices
of $\rho_V$ in the transported basis. Irreducibility and the dual
highest-weight statement follow from the precise highest-weight
input stated above, transported through $T$ and the nonzero dual
pairing. Distinct nonconjugate embeddings are essential here.

The passage from rational unitary points to the product of general
linear groups has an elementary density justification. On the
rational vector space of skew-Hermitian matrices the Cayley map
$X\mapsto(1-X)(1+X)^{-1}$ is inverse to the same formula on the
open subset where $1+g$ is invertible. A polynomial vanishing on
all rational points of a vector space over $\mathbb Q$ is zero,
by induction in its coordinates. Clearing the Cayley denominators
therefore proves Zariski density of the rational unitary points in
the restriction-of-scalars group. After algebraic scalar extension
this group is a product of general linear groups, one for each
conjugate embedding pair. Projection to distinct, nonconjugate
embeddings is surjective. An invariant subspace is defined by
polynomial matrix equations in a basis and a quotient map; density
makes it invariant under that entire product. The retained
polynomial highest-weight theorem consequently applies to the
actual rational-group representation as well.

For the actual row $n=R.n_V$ and
$\lambda(\tau)=\mathrm{lamLdef}(n,j,\tau)$ the dual highest
weight is exactly $\lambda(\tau)$; substituting its entries in
$\operatorname{infOf}$ gives
\[
 \operatorname{infOf}(\lambda(\tau))
       =\operatorname{paramV}(R,h_c(\tau),j-1,1/2).   \tag{MR.9}
\]
Thus the weight in \texttt{MtReal} is the coefficient representation
already used in the geometric forms. It is not a newly chosen
irreducible representation with an equal dimension.

\paragraph{MR6: rational Hecke operators with their actual normalization.}
At a split prime outside the level, the Hecke double coset of the
specified diagonal matrix is a finite union of right $J_{\rm std}$
cosets. To see finiteness, $J\cap aJa^{-1}$ is open in compact $J$,
so its cosets have a finite subcover. Choose the actual representatives
$a_t$ supported at that prime, where the level character is trivial.
They give the source operator
\[
       T_{s,k}F(g)=q_s^{-k(n-1)}\sum_t F(ga_t).        \tag{MR.10}
\]
Changing a representative by a local level element has no effect,
since its character is trivial. Left rational equivariance commutes
with the sum. For right equivariance, the local component permutes
these cosets; the complementary level component commutes with
their supported representatives and contributes exactly its original
character. MR4 supplies this product decomposition. Thus
$T_{s,k}$ preserves the full type space.
Its normalization in (MR.10) is the literal \texttt{NG.tN}, rather
than an unnormalized double-coset sum.

Every value in the sum is in the image of the rational coefficient
module: MR1 and the actual coefficient map give this at an arbitrary
adelic argument, not just at the selected representatives.
The normalization is rational. By MR2 and type-space preservation,
the Hecke transform
has a unique coordinate vector $x'\in L^d$. Consequently
$T_{s,k}\Phi(x)=\Phi(x')$ for the precise basis map $\Phi$ of MR2--MR3.
This proves the rational clause without any assertion that its
matrix is integral at the residue characteristic of $s$.

\paragraph{MR7: prime-local Hecke integrality and one common denominator.}
Let $\mathfrak w$ be a prime of $\mathcal O_L$ whose residue
characteristic differs from that of $s$. Each representative in
(MR.10), together with its inverse, has denominators supported at
the rational prime below $s$. This follows directly for the
diagonal Hecke element, and multiplication by integral level
matrices preserves it. Thus the lattices attached to $g$ and
$ga_t$ are commensurable by an integer invertible at
$\mathfrak w$. They have equal localizations there.

The same is true of their actual $\nabla$ coefficient lattices.
To justify this last step, choose one basis of each localized
underlying lattice. The scalar-extended lattice, its dual, and all
polynomial coefficient modules defined from it use that same basis.
Their presentations therefore agree over $\mathcal O_{L,\mathfrak w}$.
This argument uses the actual dual-Weyl lattice $\nabla$; it makes
no assertion that integral $\Delta$ and $\nabla$ are equal.
Since $q_s$ is a unit at $\mathfrak w$, (MR.10) preserves the
localized coefficient lattice. If $x\in\mathcal N_{\psi}$ and $x'$
is its rational Hecke transform, each finite class $c$ admits
$d_c\in\mathcal O_L\setminus\mathfrak w$ with
$d_c(E_\psi x')_c\in\nabla(\Lambda_c)$. Put $d=\prod_c d_c$.
Then $d\notin\mathfrak w$ and every coordinate of $dx'$ is integral:
\[
                    dx'\in\mathcal N_\psi.          \tag{MR.11}
\]
This proves the exact denominator assertion of
\texttt{heckeOperator\_\allowbreak keeps\_\allowbreak rational\_\allowbreak integral\_\allowbreak levelForms}.

\paragraph{MR8: extending the fixed rational basis to the full MtReal map.}
Now use the actual tuple comparison
$e:L^d\simeq D.\mathrm{MtL}(\mu,j)$, and suppose, as in
\texttt{TupIn}, the vectors
$b_a=j_{\rm Mt}(e(e_a))$ are complex-independent and span the fixed
$D.\mathrm{Mt}(\mu,j)$. They are a basis. Define
\[
 \mathrm{real}\Bigl(\sum_a z_a b_a\Bigr)
             =\sum_a z_a\Phi(e_a).                  \tag{MR.12}
\]
MR2 proves injectivity and MR3 proves that the range is the entire
specified type space. Semilinearity gives
$\mathrm{real}(j_{\rm Mt}(e(x)))=\Phi(x)$ for every $x\in L^d$.
MR4 supplies the open level and MR5 the actual algebraic weight.
These are all seven non-instance fields of \texttt{MtReal}:
weight, level, character, openness, map, injectivity and full range.
This construction proves \texttt{side\_\allowbreak read} and \texttt{tup\_\allowbreak read}
under their literal geometric/basis hypotheses. Their exact maps
are specified by (MR.12); dimension counting is not used to invent
a map on an unrelated space.

\paragraph{MR9: all eight Good clauses and the exact MtInput lattice.}
The product, parahoric, unipotent and hyperspecial clauses are MR4;
the weight clause is MR5; rational Hecke stability is MR6.
For integrality take a generator $y$ of the fixed
$\mathrm{MtLat}(\mu,j)$. The given \texttt{OLatTie} writes
$y=j_{\rm Mt}(e(x))$ with $x\in\mathcal N_{\psi}$.
Use MR7, then write the integral vector $dx'$ in the given finite
generators as
\[
 j_{\rm Mt}(e(dx'))=\sum_g\epsilon(c_g)g,
       \qquad c_g\in\mathcal O_L.
\]
Because $d\notin\mathfrak w$, each $d^{-1}c_g$ has nonnegative
additive valuation at the specified place, and
\[
      T_{s,k}\mathrm{real}(y)
        =\mathrm{real}\Bigl(\sum_g\epsilon(d^{-1}c_g)g\Bigr).
                                                               \tag{MR.13}
\]
The generators already lie in the range of $j_{\rm Mt}$, and their
complex span is full by the same lattice tie. This is the eighth
\texttt{Good.full} clause. The exclusion of the residue prime in
MR7 is exactly the exclusion in \texttt{Good.int}.

Choose these realizations for every tuple and every $1\le j\le m$.
The result is \texttt{MtFam}. The source's generator-place condition
\texttt{GenAt}, together with the fixed conductor support and half
selection, places every Hecke generator outside both the level and
the relevant coefficient residue prime. Applying (MR.13) gives all
eight fields of \texttt{MtInput}, including the literal exponent
range $k\le r$ and the twisted-place qualification of its integral
and full-lattice clauses.

\paragraph{MR10: pointwise MtAd on used and unused tuples.}
Let $Y$ be the supplied \texttt{R3.Read} on the same fixed state.
For a used tuple its \texttt{eMt}, \texttt{formsL} and
\texttt{latL} are exactly the basis and lattice data in MR8--MR9.
The reading equality first holds on these basis vectors, by the
definition of \texttt{Y.mtRead} and the geometric coefficient map.
Evaluation at any fixed $g$ is complex-linear. Equality on the
basis therefore proves, for every $x$ and every $g$,
\[
 e_W\bigl(\mathrm{real}(x)(g)\bigr)
       =(Y.\mathrm{IL}.j_{\rm Mt}(\mu,h,i)(x))(P_i\psi(g)).
                                                               \tag{MR.14}
\]
Here $P_i\psi$ denotes the actual frame map $(PV_i).\psi$.
Its standardness, rational-point equivalence, coefficient-action
intertwining and entry-by-entry identity-frame equality follow from
the matrix construction in MR3 and MR5. These, (MR.14), and $e_W$
are precisely the seven \texttt{MtAd} fields.

Unused tuples cannot simply be omitted from \texttt{MtFam}.
The source passes to the all-tuples level, whose level is constructed
from the same half selection. The literal matrix equality
$m_L(\mu,j)=\mathrm{lMat}(\mu,j)$ transports the already supplied
\texttt{Rw.Tie1}: its complex basis, full span and integral-lattice
identity supply MR8 again. No reading equality is requested for an
unused tuple. Combining the used and unused constructions proves
\texttt{inMt\_\allowbreak state} and \texttt{sMt} on an actual supplied
\texttt{St}/\texttt{ColPin}. It does not construct that state or
the missing same-object ties required to obtain it.

\paragraph{MR11: what the rational input really supplies.}
For a rational form datum, a finite full lattice $M_{K_s}$ in
$W_{K_s}$ gives the literal \texttt{RatIn.fd} and
\texttt{RatIn.mk\_\allowbreak fg}. Choose a finite set of module generators
$t_1,\ldots,t_a$. Each lies in $M_{K_s}$, every lattice vector is an
$\mathcal O_{K_s}$-linear combination of them, and every rational
vector is a $K_s$-linear combination. A basis of $W_{K_s}$ which
remains a basis over $L$ gives \texttt{RatIn.ks}. The finite group
action and its character descend by their actual coefficient
identities; this supplies \texttt{rho} and \texttt{chi}.

The remaining assertions are substantial. With $H_{k,P}$ the
actual rational Hecke operators, $e$ the prescribed central
idempotent, and $U$ the prescribed finite unipotent group, put
\[
 W^{\rm good}=\ker(e-1)\cap
      \bigcap_{u\in U}\ker(a(u)-\psi_U(u)),           \tag{MR.15}
\]
where $a$ is the character-adjusted action. One must supply the
literal good-row decomposition
\[
 W^{\rm good}\subseteq
  \sum_{\substack{\sigma:\ \mathrm{isot}(j,\sigma)\ne0\\
                              \mathrm{goodRow}(\sigma)}}
    \left(W^{\rm good}\cap
        \bigcap_{\substack{1\le k\le r\\P\ {\rm split},\ P\notin S_b}}
            \ker(H_{k,P}-e_{j,\sigma,k,P})\right),   \tag{MR.16}
\]
with $e_{j,\sigma,k,P}\in K$, and the power-denominator bound
\[
    (\mathrm N P)^a H_{k,P}(M_L\cap W^{\rm good})
                       \subseteq M_L               \tag{MR.17}
\]
for some $a$, where $M_L=\mathcal O_L M_{K_s}$.
Together with the rational stability of $H$, $H_{0,P}=1$, the
specified nontrivial $\theta$ and the actual Hecke/eigenvalue
readings, these construct the fourteen fields of \texttt{RatIn}
and the five fields of \texttt{RatAd}. This is exactly the
\texttt{sRat\_\allowbreak of\_\allowbreak isotInternal} assembly. Its source derives
(MR.16)--(MR.17) from genuine internal row decomposition and
Hecke-transform/reading/integrality inputs. They are retained here;
finite class dimension alone does not prove either assertion.

\paragraph{MR12: an integral additive reading without enlarging the DVR.}
Let $O$ be a DVR, $H=\operatorname{Frac}(O)$, and
$\iota:K\hookrightarrow H$ the specified reading. There exists an
additive map $\lambda:L\to H$ such that
\[
 \lambda(1)=1,\qquad \lambda(ax)=\iota(a)\lambda(x)
 \ (a\in K),\qquad \lambda(\mathcal O_L)\subset O.    \tag{MR.18}
\]
Here is a direct construction. Extend $\iota$ to an embedding
$L\hookrightarrow\overline H$, and let $N$ be the $O$-span of the
images of an integral $\mathbb Z$-basis of $\mathcal O_L$. It is a
finite torsion-free $O$-module and contains $1$. Every element of $N$
is integral over $O$. Since $O$ is integrally closed,
\[
                      N\cap H\cdot1=O\cdot1.        \tag{MR.19}
\]
It follows that $N/O\cdot1$ is torsion-free: if a nonzero element
of $O$ annihilates a class, its representative lies in $H\cdot1$,
and (MR.19) applies. A finite torsion-free module over a DVR is free.
Lift a basis of this quotient, and define an $O$-linear functional
$N\to O$ to be $1$ on $1$ and $0$ on those lifts. Extend first
$H$-linearly to $HN$ and then to $\overline H$ by a vector-space
complement. Compose with $L\hookrightarrow\overline H$.
The resulting map has all three properties in (MR.18).

It is only additive and $K$-semilinear; multiplicativity is neither
claimed nor needed. This is why the argument works even when the
whole field $L$ does not embed into the already fixed $H$.
The source's \texttt{exists\_\allowbreak integral\_\allowbreak functional} also proves the
underlying primitive-vector splitting without assuming a full
lattice, by taking a basis of the image in $V/Hv$; the same proof
above applies to that image and extends the resulting functional.

\paragraph{MR13: finite integral equations descend to the original field.}
Choose a $K_s$-basis $b_1,\ldots,b_d$ of $W_{K_s}$ which is also
an $L$-basis of $W_L$, and define
\[
 \Phi\Bigl(\sum_a x_a b_a\Bigr)
           =\sum_a\lambda(x_a)(1\otimes b_a)
       \in H\otimes_{K_s}W_{K_s}.                  \tag{MR.20}
\]
Then $\Phi$ is $K$-semilinear, fixes rational vectors in the sense
$\Phi(w)=1\otimes w$, and commutes with every $L$-linear operator
whose restriction preserves $W_{K_s}$. This follows by writing its
matrix in the basis $b_a$: all its entries lie in $K_s\subset K$.
Also $\Phi(aw)=\lambda(a)(1\otimes w)$ for
$a\in L$, $w\in W_{K_s}$, so integral $a$ gives an $O$-multiple.

The further assertion needed for lattice stability is integral,
not merely vector-space descent. Suppose finitely many vectors
$v_{t,a}$ have all their $b$-coordinates in $K$, and
$o_t\in O$ satisfy
\[
                    \sum_t o_t\Phi(v_{t,a})=0
                                      \quad\text{for every }a.
                                                               \tag{MR.21}
\]
Clear one common denominator in the finitely many coordinates,
giving a matrix $A$ over $\mathcal O_K$. The embedding
$\mathcal O_K\to O$ is well-defined because images of global
integers are integral over the integrally closed ring $O$.
As an $\mathcal O_K$-module $O$ is torsion-free; since
$\mathcal O_K$ is Dedekind it is flat. For clarity, this flatness
can be checked locally: a localization of a Dedekind domain at a
nonzero prime is a DVR, and a torsion-free module over a DVR is a
filtered union of finite free modules, hence flat. Flatness is local.
Tensor the exact sequence
$0\to\ker A\to\mathcal O_K^t\to\mathcal O_K^a$ with $O$.
Thus every vector $(o_t)$ in (MR.21) is an $O$-linear combination
of images of solutions $(c_t)\in\mathcal O_K^t$ satisfying
$\sum_t c_tv_{t,a}=0$. Cancelling the nonzero common denominator
recovers the original equations. This proves the exact
finite-equation assertion of \texttt{exists\_\allowbreak integral\_\allowbreak descent}.

\paragraph{MR14: the full local lattice and its genuine good intersection.}
Fix a twist index $i$ from the already constructed \texttt{RingSpec}
and a side $b$. Put $O=O_i$, $H=\operatorname{Frac}(O_i)$, with
the side-dependent $K_s$-algebra structure obtained from
$\mathrm{read}_{i,b}$. For every actual row $1\le j\le m$, define
\[
 \mathcal F=H\otimes_{K_s}W_{K_s},\qquad
 B^{\rm sup}=O\langle1\otimes M_{K_s}\rangle.        \tag{MR.22}
\]
If $t_1,\ldots,t_a$ generate $M_{K_s}$, their images generate
$B^{\rm sup}$, since all images of $\mathcal O_{K_s}$ belong to
$O$. They span $\mathcal F$ over $H$. Consequently $B^{\rm sup}$
is finite, torsion-free and free over $O$.

Let $V_0\subset\mathcal F$ be the actual central-idempotent and
generic-character subspace, and put
\[
                       B=B^{\rm sup}\cap V_0.       \tag{MR.23}
\]
It is a submodule of a finite module over a noetherian ring, hence
finite, and is torsion-free, hence free. Moreover $HB=V_0$.
Indeed, express $v\in V_0$ in a finite generating set of
$B^{\rm sup}$ and clear the finitely many denominators to obtain
$0\ne d\in O$ with $dv\in B^{\rm sup}\cap V_0$.
The natural map $H\otimes_O B\to\mathcal F$ is injective:
a localized vector is zero only when a nonzero denominator kills
its numerator, impossible in a vector space. Its image is exactly
$V_0$. Thus (MR.23) is the source's \texttt{BI}, with its actual
rank, not a free module selected solely from that rank.

\paragraph{MR15: the actual idempotent, generic character and good space.}
The specified coefficients $e(g)\in K$ map under $\epsilon$ to
the central good-type projector $e_0$ in $\mathbb C[G_0]$.
Coefficientwise injectivity of $K[G_0]\to\mathbb C[G_0]$ reflects
$e_0^2=e_0$ and centrality. Thus $e^2=e$ and $e$ is central over
$K$, and the same polynomial identities hold after every actual
$\mathrm{read}_{i,b}$. The character-adjusted action is
\[
       a(g)=\chi(g)^{-1}\rho(g),\qquad
       a_{i,b}(g)=\mathrm{read}_{i,b}(\chi(g))^{-1}
                                      (1\otimes\rho(g)).\tag{MR.24}
\]
No ambiguity is introduced by changing the side. Each value of the
finite additive character $\theta_k$ has order dividing the prime
$g_k$, hence its complex conjugate is its inverse. The literal
conjugate-reading identity of \texttt{TwR} gives
$\mathrm{read}_{i,\mathrm{false}}(\zeta)
=\mathrm{read}_{i,\mathrm{true}}(\zeta)^{-1}$ for these values.
Injectivity preserves $\theta_k\ne1$. This proves the exact
\texttt{psiU\_\allowbreak side}, genericity and nonzero-coordinate conditions.

Let $p_U=|U|^{-1}\sum_u\psi_U(u)^{-1}a(u)$. The calculation of
MR2 makes it the projector onto the generic-character subspace;
centrality makes $e$ commute with it. Therefore
$V_0=e\operatorname{im}(p_U)$ is precisely the common kernel of
$e-1$ and $a(u)-\psi_U(u)$. These matrices are defined over $K$.
Finite-kernel descent implies that their common kernel over $H$
is spanned by the images under $\Phi$ of vectors in
$W^{\rm good}\subset W_L$. Apply (MR.16) to these vectors and
use the intertwining of MR13. This proves that $V_0$ is spanned by
the simultaneous eigenspaces with the actual row labels and
eigenvalues
\[
           c_{i,b,P}^{\,k}\,
                  \mathrm{read}_{i,b}(e_{j,\sigma,k,P}).\tag{MR.25}
\]
In passing between the $K$-kernel and the $L$-kernel, one may apply
the $K$-linear coefficient projection $\lambda$ to the $L$
coordinates; it is the identity on $K$. This proves the spanning
claim without extending $L\hookrightarrow H$ or changing $H$.

\paragraph{MR16: stability of the actual intersection lattice.}
Write $H'=c_{i,b,P}^{\,k}(1\otimes H_{k,P})$, where
$c_{i,b,P}\in O_i^\times$ is the prescribed delta correction.
Take $x\in B=B^{\rm sup}\cap V_0$ and write it as
$x=\sum_t o_t\Phi(w_t)$ with $w_t\in M_{K_s}$.
The equations defining $V_0$ say that $(o_t)$ satisfies a finite
system as in (MR.21): use the images of
$(e-1)w_t$ and $(a(u)-\psi_U(u))w_t$, whose coordinates lie in $K$.
MR13 expresses $x$ as an $O$-linear combination of $\Phi(y)$ with
\[
          y\in\mathcal O_L M_{K_s}\cap W^{\rm good}.\tag{MR.26}
\]
For such $y$, (MR.17) gives
$(\mathrm N P)^aH_{k,P}y\in\mathcal O_LM_{K_s}$.
MR12--MR13 send this lattice into $B^{\rm sup}$. Because
$P\notin S_0(i)$, it is not above the residue characteristic
$\ell_i$; thus $\mathrm N P$ is a unit of $O_i$. Cancelling its
power proves $\Phi(H_{k,P}y)\in B^{\rm sup}$.
The spectral expansion (MR.16) also proves
$H_{k,P}W^{\rm good}\subset W^{\rm good}$.
Hence $H'\Phi(y)\in V_0$, and multiplication by the unit
$c_{i,b,P}^{\,k}$ preserves both conditions. Extend over the
finite $O$-linear combination to obtain
\[
                              H'B\subseteq B.        \tag{MR.27}
\]
This is \texttt{HkStab} on the original coefficient ring. It is
proved for exactly $1\le j\le m$, $1\le k\le r$, split $P$
outside $S_0(i)$. The source defines the lattice operator to be
the identity outside this range, in particular at $k=0$; it does
not claim (MR.27) for a ramified or excluded Hecke operator.

\paragraph{MR17: integral eigenvalues and finite, separated signatures.}
An endomorphism of a finite free $O$-module satisfies its monic
characteristic polynomial. If its scalar extension has a nonzero
eigenvector of eigenvalue $a\in H$, applying that polynomial to
the eigenvector gives $p(a)=0$. Since $O$ is integrally closed,
$a\in O$. The same conclusion for a finite module follows from
the determinant trick; freeness from MR14 suffices here.

For fixed $i$, take triples $(j,b,\sigma)$ with
$1\le j\le m$, $\mathrm{isot}(j,\sigma)\ne0$ and
$\mathrm{goodRow}(\sigma)$ for which every admissible eigenvalue
in (MR.25) lies in $O_i$. Store that integral eigenvalue at the
admissible pairs $(k,P)$ and store $1$ elsewhere. Define
$\mathrm{Sig}(i)$ as the image of these triples under
\[
       (j,b,\sigma)\longmapsto
                (j,b,(k,P\mapsto h_{j,b,\sigma}(k,P))).\tag{MR.28}
\]
This is a finite set even though the eigenvalue function has
infinitely many arguments, because its domain of triples is finite.
Select an original row label for each element of this image.
Two signatures with the same row number, side and all admissible
eigenvalues are equal: outside that range both functions are $1$.
This proves the literal injectivity clause and $h(0,P)=1$.

MR15 spans $V_0$ by simultaneous eigenvectors. Every nonzero
simultaneous eigenvector used in that span has integral eigenvalues
by MR16 and the preceding
monic-polynomial argument. Its triple is therefore included in
(MR.28); zero vectors contribute nothing. Via MR14, these
signature eigenspaces span $H\otimes_O B$. The selected original
row is good and has nonzero prescribed isotypic space, and
(MR.25) is unchanged. These are all five conclusions of
\texttt{exists\_\allowbreak sigData\_\allowbreak hecke\_\allowbreak eigenvalue\_\allowbreak formula}.

\paragraph{MR18: characteristic polynomials and the positive TwF fields.}
For an admissible individual Hecke operator, MR17 supplies a basis
of each of its eigenspaces and their direct sum spans the fraction
space. Distinct eigenvalues have independent eigenspaces: applying
$\prod_{a\ne a_0}(H'-a)$ isolates each summand. The matrix in the
collected basis is diagonal, so its characteristic polynomial is
the product of the corresponding linear factors. The matrix of
the integral restriction in an $O$-basis maps to the matrix of
its scalar extension. Determinants commute with that map; hence
the literal integral characteristic polynomial splits over $H$.
For a non-admissible operator the defined matrix is the identity,
and the polynomial is $(X-1)^{\operatorname{rank}_O B}$.

This constructs the positive form part of \texttt{TwF}: the exact
finite signature set, its row number and side, the module $B$ of
(MR.23), its finite freeness, the restricted Hecke operators, and
the integral eigenvalues. The signature Galois representation is
the actual previously constructed $\rho_{i,\sigma,b}$ of NR24
at the selected original row and side. Its continuity and its
universal reading clause follow from NR24, under its explicit
\texttt{GalEx}, \texttt{RowCusp}, \texttt{ChebSplit} and
\texttt{SideA} premises. Signature selection does not weaken
those premises or change that representation.

\paragraph{MR19: explicit weights over the same fraction fields.}
The row parameter is genuinely realizable by a dominant integral
weight. Let $n=R.n_V$ and let $S=\operatorname{paramV}(R,h,j,s)$.
Its entries are the multiset
\[
 \{v_a\}_{a<k-1}\ \sqcup\ \{-b_a\}_{a<m}\ \sqcup\
 \{-e_j+(t-1-2c)/2\}_{c<t},
          \quad\text{then shifted by }-s.           \tag{MR.29}
\]
The source's adjacent gap conditions make each block strictly
decreasing and separate the three blocks: the third lies between
the appropriate consecutive $-b$ entries and below the $v$ entries.
Thus $S$ has $n$ distinct entries. Integrality of $v,b,e$, oddness
of $t$, and evenness of $n$ show that, when $s\in\mathbb Z+1/2$,
each entry is $z+(n-1)/2$ with $z\in\mathbb Z$.
List the negatives of these integers increasingly as
$a_0<\cdots<a_{n-1}$ and set
\[
                         \lambda_i=-a_i+i.          \tag{MR.30}
\]
Since $a_{i+1}\ge a_i+1$, $\lambda$ is weakly decreasing.
Its infinitesimal-character entries are
$\lambda_i+(n-1)/2-i=-a_i+(n-1)/2$, exactly $S$.
The actual character frame gives $s=1/2$ for the one-dimensional
$\eta_T$ parameter; more generally the proof only needs the stated
half-integrality. This proves the literal combinatorial weight
supplier rather than leaving it as a record field.

For every $i,j,b$, use the side-dependent embeddings
$E\to\operatorname{Frac}(O_i)$ prescribed by \texttt{RingSpec}
and take the dual of the tensor polynomial module $VV(\lambda)$.
Its finite basis and matrices are over that very field. The
highest-weight input of MR5 proves irreducibility and the exact
dual highest weights; extension of scalars preserves this tensor
polynomial realization and its dual, hence preserves
irreducibility. The comparison is explicit on the polynomial
basis, including the inverse determinant powers. Therefore the
weight bound $c_2$ in \texttt{wt\_\allowbreak weights} can be the bottom field:
no additional local field is needed for this weight construction.
Only the source's harmless out-of-row padding uses a default
weight; every row $1\le j\le m$ uses (MR.29)--(MR.30).

\paragraph{MR20: actual local levels and the whole coefficient-field type space.}
Fix an admissible character $\psi$ and conductor set $Q$ at
$(i,k)$. Select one prime in each conjugate pair using the prescribed
$P_q$; write $Q_P=\operatorname{qpOf}(P_q,Q)$ and
$D_P=\operatorname{qpOf}(P_q,p_D(i))$. The split hypotheses and
conjugation stability show that each pair is met once. The three
families of real places, those from $Q_P$, $D_P$ and the type
primes, are disjoint. At these places define $J$ inside the
integral maximal compact by, respectively, the $(r,n-r)$
parabolic, the same parabolic, and the type-prime unipotent level.

These are subgroups by reduction modulo the local maximal ideal:
block upper-triangular matrices form a subgroup, and requiring the
upper-left block to be upper unitriangular is also a subgroup.
The determinant of the complementary block is a unit and
multiplicative in the residue field. The superdiagonal sum of an
upper unitriangular block is additive under multiplication.
Consequently the following product is a character on $J$:
the actual generic additive/type determinant factor at every type
prime, the tame character of $\psi$ on each complementary
determinant in $Q_P$, and the tame character of
$\delta_i^{\varepsilon_b}$ on the corresponding determinants in
$D_P$, where $\varepsilon_{\rm true}=1$ and
$\varepsilon_{\rm false}=-1$. Tame residue values are well-defined
and multiplicative by GA17 and AR3; wild inertia is killed since
the character orders and these residue characteristics are
coprime. At the type primes the literal \texttt{tauT} is the
residue character of the already normalized \texttt{chi1}, not
an arbitrary finite-order replacement.

The product of the norms of these finitely many real places gives
a positive integer $N$ for which the congruence subgroup is in
$\ker\chi$. Outside those places $J$ is maximal compact with
trivial character. These prove all local, open and hyperspecial
clauses of \texttt{LevLawsT}, with the literal sets $Q_P,D_P$.

Finally, the whole space
$\operatorname{TypeSpF}(\rho_{i,j,b,k},J,\chi)$ is finite-dimensional.
Choose finitely many classes modulo the maximal compact, then
finitely many coset representatives of its level-$N$ congruence
subgroup. Evaluation at their products is injective: every group
element is rational times one such product times an element of
the congruence kernel, and equivariance recovers its value.
The target is a finite product of the finite-dimensional coefficient
space. The literal \texttt{MtSubT} chooses this actual type space
when these level laws exist, and its comparison map is the identity
under that equality. Take the $O_{i,k}$-span of any fraction-field
basis as \texttt{MtLatT}; it is finite free and spans the space.
Its comparison to the original global \texttt{MtLat} is a separate
arithmetic compatibility problem, not a consequence of this choice
of basis.

\paragraph{MR21: the initial NExpTied fields and exact normalized comparisons.}
MR14--MR20 construct the coefficient-field form modules, the
generic action, central idempotent, weight modules, and full
admissible type spaces used in the source's \texttt{mkEx}.
The central coefficients remain the original $e(g)\in K$.
The universal anticyclotomic characters and their support,
continuity and conjugation laws are NR8--NR11. Their inertia
surjectivity is only asserted under literal \texttt{AdmQ}:
that condition supplies, at each conductor pair, a character
surjective on the chosen inertia. Its conjugate is surjective on
the other prime by inversion. The universal character's coordinate
is the selected two-prime character, so it is surjective there and
zero on inertia at a different, nonconjugate conductor prime.
Thus no stronger inertia-surjectivity assumption has been inferred
from support alone. The source signs $\mathrm{sgnT}$ and
$\mathrm{sgnG}$ are both $1$; MR15 and MR20 check their actual
character equalities with this convention.

For admissible rows the two tensor descriptions
\texttt{FmT} and \texttt{Fm S} are identical after replacing
$\mathrm{WKr}(j)$ by $W_{K_s}$ and identifying their identical
$K_s$-algebra structures. Send $a\otimes w$ to $a\otimes w$.
This gives the literal $e_F$, carries $B^{\rm sup}$ to the
specified lattice, and intertwines the character-adjusted action.
The Hecke comparison is
\[
 e_F(H'_{k,P}x)=c_{i,b,P}^{\,k}
                         (1\otimes H_{k,P})e_F(x).  \tag{MR.31}
\]
The unit $c_{i,b,P}$ is the value of the already fixed delta
representation at the selected Frobenius, inverted on the false
side. Any two allowed Frobenius representatives differ by inertia,
on which delta is trivial outside $S_0(i)$, so the value is
independent of the selected representative. This is the literal
\texttt{cHk\_\allowbreak del}, with the existing \texttt{IsFrob} convention.
The lattice injection is the subtype map $B\hookrightarrow
\mathcal F$. Its Hecke equality is (MR.27) for $k>0$ and the
identity equality for $k=0$. All these comparisons use the original
\texttt{RingSpec} and the same two side embeddings.

\paragraph{MR22: exact SideB conclusion and its initial-stage boundary.}
Retain the full input list of \texttt{second\_\allowbreak step}: the fixed
$D$, character frame, $j_E$, \texttt{SideA}, actual rational datum
and \texttt{RatIn}, a finite low-prime set, $P_q$ lying above every
real place, an actual \texttt{MtFam}/\texttt{MtInput}, global Artin
reciprocity, \texttt{GalEx}, \texttt{RowCusp}, \texttt{ChebSplit},
conjugation stability of $K$, integrality of $e\mathrm{Exp}$, and
splitting at both auxiliary primes. MR8--MR10 can supply the
Mt-fields when their compatible geometric state is given; MR11
specifies exactly the remaining rational spectral inputs.

NR24 chooses a finite field bound $c_1(\ell,q)$ for all genuine
row representations before the later parameters. MR19 chooses
$c_2=\bot$. Put $c_0=c_1\vee c_2$. Thus the quantifier order is
\[
 \exists c_0\ \forall b_0,b_1,S_x,T,\mathrm{Sp},\qquad
 c_0\le b_0,\quad S_b\subseteq S_x,\quad
 \bar S_x=S_x,\quad
 \{P:P\mid p,\ p\in\mathrm{low}\}\subseteq S_x.    \tag{MR.32}
\]
For every such \texttt{RingSpec}, MR14--MR21 give the actual
form/signature fields and the initial \texttt{NExpTied}. The
remaining fields of \texttt{SideB} have explicit sources:
the coefficient frame, normalized characters and all embedding
squares are the same \texttt{Sp}; cover at every prescribed
coefficient place is NR17; rational action is the same
\texttt{RatIn}; $e_F$, the lattice and Hecke equalities are MR21;
signature span, separation and eigenvalues are MR17--MR18; and
\texttt{rhoS} is NR24 at the selected good row. This proves every
\texttt{SideB} field under the source's exact qualification below.

That qualification is material. Literal \texttt{mkF/mkEx} initialize
\texttt{Theta} and \texttt{ThetaPsi} as zero submodules,
\texttt{Dat=PEmpty}, \texttt{thP=0}, \texttt{twDel=0}, and
$e_H=1$; the toric family is \texttt{TwL.empty} and
\texttt{SideA.bare} says that the old local index set is empty.
These fields are initialization choices, not the genuine theta
image, toric realization or nonempty local family. The source
conclusion is exactly
\[
 \exists\mathrm{TwF},\mathrm{ex},h_0,\quad
  [\text{the full-lattice properties of }B^{\rm sup}]
  \ \land\ \forall\mathrm{th}:NTh(\mathrm{new\ data}),\quad
        \operatorname{Nonempty}(\mathrm{SideB}(D_{\rm new})).
                                                               \tag{MR.33}
\]
It does not assert that such a \texttt{th} has been constructed.
$D_{\rm new}$ is the specified \texttt{setTw} datum: it preserves
the original host, row spaces, $L$, $K$, normalized frame and
rational operators, and replaces its twist/form expansion fields
by exactly the objects above. Its twist index set can be the
genuinely nonempty all-prime family from NR14--NR18. Nonvanishing
of a good space is faithfully preserved by (MR.22)--(MR.23) and
MR15, but is not asserted for every row if the original space is
zero. No new nonempty local index, genuine theta operator or
\texttt{IsTripleFor} is supplied by (MR.33).

\paragraph{MR23: one finite local field splits the entire commuting Hecke family.}
Fix the original row $j$, coefficient place $w$, embedding
$\iota_w:L\hookrightarrow\overline{\mathbb Q}_\ell$, and actual
finite-dimensional space
$V_L=D.\mathrm{what}_2.\mathrm{MtL}(\mu_w,j)$.
Let $H_\alpha\in\operatorname{End}_L(V_L)$ be the source's
\texttt{hLC} family, indexed by all pairs of finite places and
Hecke indices, with its supplied pairwise commutation law.
The $L$-linear span of this possibly infinite family is finite
dimensional. Choose finitely many members $H_1,\ldots,H_s$
spanning it, and put
\[
             p(X)=\prod_{a=1}^{s}\det(X-H_a)\in L[X].
                                                               \tag{MR.34}
\]
This is a nonzero monic polynomial and annihilates every selected
$H_a$ by Cayley--Hamilton. The empty product is $1$ when the
operator span is zero. Adjoin to $\mathbb Q_\ell$ all roots of
$\iota_w(p)$ in its algebraic closure. There are finitely many
roots, each algebraic over $\mathbb Q_\ell$, so the resulting
$F_0/\mathbb Q_\ell$ is finite.

Now let $H$ be any later coefficient field whose embedding in
$\overline{\mathbb Q}_\ell$ contains $F_0$ and extends the same
$\iota_w$. Every root of the polynomial obtained from $p$ over
$H$ therefore lies in $H$, so each selected operator has split
minimal polynomial on $H\otimes_LV_L$. Commuting split operators
are simultaneously triangularizable over $H$: a nonzero
eigenspace of the first operator is invariant under the others;
successively taking nonzero eigenspaces gives a common
eigenvector, and induction on the dimension applied to the
quotient gives an invariant full flag. Restrictions and quotients
remain split because their minimal polynomials divide the
original ones. Every $H$-linear combination of the selected
operators preserves this flag. Hence every original $H_\alpha$
has split characteristic and minimal polynomials over $H$.

For a split minimal polynomial
$\prod_c(X-c)^{a_c}$, the factors for distinct $c$ are coprime.
B\'ezout polynomials give the direct-sum decomposition into
$\ker(H_\alpha-c)^{a_c}$. Consequently
\[
 \sum_{c\in H}\operatorname{GenEig}(1\otimes H_\alpha,c)
          =H\otimes_LV_L.                              \tag{MR.35}
\]
This is the literal \texttt{exists\_\allowbreak finExt\_\allowbreak split} conclusion
for all admissible $\alpha$, with the stronger elementary
argument valid for every member of the family. No semisimplicity
has been assumed or concluded. The quantifier is one $F_0$ for
the fixed $(j,w)$, before every later compatible \texttt{CoefIn};
it is not one finite global field for all coefficient places.

\paragraph{MR24: a nonempty residual Hecke block has a genuinely local algebra.}
Let $O$ be a local ring with maximal ideal $\mathfrak m$, let
$X$ be a finite nonempty set, and let
$T\subseteq O^X$ be an $O$-subalgebra whose coordinate values
are congruent modulo $\mathfrak m$:
\[
             f(x)-f(y)\in\mathfrak m
       \quad(f\in T,\ x,y\in X).                       \tag{MR.36}
\]
Fix $x_0\in X$. If $f(x_0)$ is a unit, (MR.36) makes every
$f(x)$ a unit. The polynomial
$P_f(Z)=\prod_{x\in X}(Z-f(x))$ annihilates $f$ in $O^X$;
its coefficients belong to $O$, hence are constants in $T$,
and its constant coefficient is a unit. Writing
$P_f(Z)=a_0+ZQ(Z)$ gives
$f[-a_0^{-1}Q(f)]=1$ in $T$. Thus $f$ is a unit of $T$
if and only if $f(x_0)$ is a unit of $O$. The inverse implication
also follows directly by applying coordinate evaluation.
The nonunits are therefore the kernel of the surjective map
\[
               T\longrightarrow O/\mathfrak m,
                   \qquad f\longmapsto f(x_0)\bmod\mathfrak m.
                                                               \tag{MR.37}
\]
Surjectivity follows from the constant functions. This kernel
is the unique maximal ideal, proving that $T$ is local and that
its residue field is the original $O/\mathfrak m$.

For the source's residual block, points are $O$-algebra
characters $\psi$ of the same Hecke algebra with a common
kernel after residue reduction. If $\psi,\phi$ are two such
characters, apply equality of their residue kernels to
$a-\phi(a)\cdot1$. Its $\phi$-value is zero, so
$\psi(a)-\phi(a)\in\mathfrak m$. This proves (MR.36) on the
Hecke generators; sums, products and constants preserve it,
so it holds on their entire generated algebra. The actual
\texttt{HeckeStd0}, finite-dimensionality and existing
\texttt{LocT} hypotheses give the finite nonempty point set
\texttt{XOf}. Applying the preceding argument to its actual
\texttt{TOf} proves the literal local-ring assertion. This is
locality of an existing nonzero-support block, not an existence
claim for a block at an arbitrary place.

\paragraph{MR25: every existing supported Hecke point has a nonzero eigenspace.}
Let $V$ be finite dimensional over a field $H$, let $S$ be a
commuting family of endomorphisms, and let $0\ne W\subseteq V$
be invariant under $S$. Suppose for every $a\in S$ that
$W$ lies in the generalized $\chi(a)$-eigenspace of $a$.
Choose a nonzero invariant subspace $U\subseteq W$ of smallest
positive dimension. For fixed $a\in S$ and $0\ne u\in U$,
choose the least $r>0$ with $(a-\chi(a))^ru=0$.
Then $(a-\chi(a))^{r-1}u$ is a nonzero element of
$U\cap\ker(a-\chi(a))$. This intersection is invariant under
all of $S$, because the operators commute. Its dimension cannot
be smaller than that of $U$, so it equals $U$. The argument is
valid for every $a$, without assuming that $S$ is finite.
Every nonzero $u\in U$ is therefore a simultaneous eigenvector.

In the actual localization construction, a point
$x\in\texttt{XOf}(i)$ includes the nonzero generalized
character space $W_x$. Apply the preceding argument to
$W_x$ and the admissible normalized Hecke operators. It gives
$u\ne0$ with their exact point values. Retain the actual
\texttt{MtLoc}, \texttt{ActEx}, and \texttt{ActHK} comparison:
its injective map $\Psi_i$ has range the prescribed residual
block, and it intertwines the specified Hecke action on the
localized relation lattice with those operators. Thus
$u=\Psi_i(z)$ for some $z\ne0$, and injectivity gives
\[
       a\,z=\operatorname{ev}_x(a)z
              \quad\text{for each Hecke generator }a. \tag{MR.38}
\]
The equality propagates to constants, sums and products, hence
to the entire generated algebra $T_i$. Therefore its actual
character eigenspace contains $z\ne0$; this is
\texttt{heckeEig\_\allowbreak ne\_\allowbreak bot}. The comparison $\Psi_i$ and
the point's nonzero support are essential retained premises.
Neither is inferred merely from MR23's generalized splitting
or from MR24's local-ring structure.

\paragraph{MR26: the original relation lattice stays free over the chosen DVR.}
Let $A$ be the original integral coefficient ring. The literal
\texttt{OsConcrete} condition says its given injection into
$L$ has image
$\{x\in F:x\text{ is integral over }\mathbb Z\}$ for a
number subfield $F\subseteq L$. Thus it identifies $A$ with
$\mathcal O_F$, so $A$ is Dedekind. It does not permit
replacement by an unrelated coefficient ring.

Let $v_1,\ldots,v_d$ be the specified complex generators of
the original lattice and let
\[
 N=\{(a_1,\ldots,a_d)\in A^d:
                 \sum_t\epsilon(a_t)v_t=0\}.          \tag{MR.39}
\]
If $0\ne c\in A$ and $ca\in N$, then
$\epsilon(c)\sum_t\epsilon(a_t)v_t=0$; the fixed coefficient
embedding is injective, so $a\in N$. Consequently $A^d/N$
is torsion free. It is finite, hence projective over the Dedekind
ring $A$: its localization at each nonzero prime is a finite
torsion-free module over a DVR and is free; finite presentation
and these local bases give projectivity. One may see the last
step by choosing finitely many denominator neighborhoods on
which bases exist and applying the finite locally free
criterion. The quotient map therefore has a section, giving
an $A$-linear projection $e:A^d\to N$ with $e|_N=1$.

For the prescribed ring homomorphism $s:A\to O$ to the chosen
DVR, apply $s$ to the matrix of $e$, obtaining $e_O:O^d\to O^d$.
Its image is exactly
$N_O=O\langle s(N)\rangle$: it maps into this span because
its columns are images of members of $N$, and it fixes the
span because it fixes every $s(n)$. Thus $e_O$ is an
idempotent and
\[
                 O^d/N_O\simeq\ker e_O.              \tag{MR.40}
\]
The right side is a finite torsion-free $O$-module, so it is
free. This proves \texttt{bjC\_\allowbreak free} for the literal quotient
of the original generator relations, for every existing
localization index and its actual coefficient map. It does not
identify that quotient with a different lattice by a dimension
comparison.

The remaining nonempty-realization problem now has a sharper
boundary. MR23--MR26 supply the finite splitting-field, local
block, nonzero-eigenvector and relation-freeness primitives.
The full \texttt{NdSpec}/\texttt{LocAll} constructor still
requires the compatible original \texttt{NChoice},
\texttt{ChoiceSide}, all coefficient/twist embedding squares,
the actual $\Psi$ comparison and local Galois/automorphic
specification. The initial object in MR22 does not supply these
requirements, nor a compatible nonplaceholder theta witness.




\section{The quantitative Weil dictionary on the fixed choice}
\label{sec:quantitative-weil-dictionary}

This section strengthens WA11--WA16 and CW23--CW25 in a different
direction: it proves the all-conjugate scalar and local-unit assertions,
and the exact denominator and coefficient bounds in the source's
\texttt{DictAt}.  The input is the \emph{same supplied} arithmetic host,
\texttt{Choice} $\mathfrak c$ and proof $\mathfrak c.\mathrm{Is}$.  In
the literal source endpoints we retain their arguments
$\texttt{FunPr.Hypotheses}_2$ and $\forall D,\texttt{NormSum.Inp}(D)$;
we do not assert the construction of those whole packages here.
The substantive ingredients of these endpoints are proved below,
using the actual Weil action and the already reconstructed theta and
Fock comparisons.  In particular there is no independent assumption
of the desired scalar, dictionary, or simultaneous local-unit property.

Write $E/F$ for the CM extension, $f=[F:\mathbb Q]=D.R.f$, and
$p$ for a literal \texttt{DefPair}.  Thus $p$ is either a top edge or
a rightmost $\chi$ edge; the $\mu$ and core cases are excluded.
Let $a,b$ be the ranks of its oriented source and target.  For finite
points $x,y$, let $L_s(x,y),L_t(x,y)$ be the source's
\texttt{srcLat}, \texttt{tgtLat}, in its fixed rational orthonormal
frames.  For a $\chi$ edge the vertex order is reversed and the first
Hermitian space is negated, exactly as in \texttt{toPair}.
The involution $\operatorname{opp}_p$ is complex conjugation on that
edge and the identity on a top edge.  Put
\[
 A_0=\operatorname{La}_0(D,\operatorname{tor}_0(D,\mathfrak c_K,\delta)),
 \qquad A_{x,y}=\operatorname{La}(L_s(x,y),L_t(x,y)),
 \qquad G_{x,y}=I_{A_0,A_{x,y}}e_0.                         \tag{QD.1}
\]
Here $I$ is the source's \emph{positively normalized} finite
Heisenberg intertwiner, $e_0=\mathfrak c.S.Z.e_0(p)$,
$\delta=\mathfrak c.S.Z.\delta$, and $c_0=\mathfrak c.S.Z.c_0(p)$.
Neither $e_0$ nor $c_0$ is rescaled.  In particular $c_0\ne0$ need
not have absolute value one.  Denote by $J_{x,y}(\alpha)$ the actual
Fock monomial jet of
$\omega_p(\operatorname{toPair}_p(xu_s,yu_t))v_p$, and put
$P_\alpha=\prod_\nu\operatorname{period}(\delta,\nu)^{\alpha_\nu}$.
The statement \texttt{JetKap} is precisely
\[
 J_{x,y}(\alpha)=c_0\operatorname{opp}_p
       \left(P_\alpha\,\partial^\alpha(\kappa G_{x,y})(0)\right)
                         \quad\hbox{for every }\alpha.     \tag{QD.2}
\]
The notation $\partial^\alpha$ means the source's
\texttt{applyAt0} on the monomial; it includes the multi-factorial
relative to the Taylor coefficient.

\paragraph{QD1: extracting arithmetic from a positive power.}
Suppose $z^m=u$, $m>0$, and $u\in\overline{\mathbb Q}$.
Then $z\in\overline{\mathbb Q}$.  If every complex embedding of
$u$ has modulus one, so does every complex embedding of $z$.
For a fixed rational prime $q$, if every embedding of $u$ into
$\mathbb C_q$ has norm one, the same is true of $z$.
Indeed $X^m-u$ proves algebraicity over the algebraic numbers,
and under any of these embeddings the equality gives
$|z|^m=|u|=1$.  The nonnegative real number $|z|$ is therefore
one.  Products, inverses, and complex conjugates preserve both
assertions, with the latter interpreted by precomposition with
the conjugation automorphism of $\overline{\mathbb Q}$.
Roots of unity satisfy all of them and are algebraic integers.
We use $\mathrm{Mod}_1(z)$ for the all-complex-embedding assertion,
including algebraicity, and $\mathrm{Unit}_q(z)$ for the assertion
at every $\mathbb C_q$ embedding.  These are the source's
\texttt{Mod1All} and \texttt{UnitAt}.  A single complex norm
identity is not an alternative proof of either statement.

\paragraph{QD2: finite Hecke values have the required conjugate norms.}
Let $\chi$ be either actual splitting Hecke character and let $z$
be a finite idele, with infinite component one.  Then
$\mathrm{Mod}_1(\chi(z))$.  Moreover,
\[
 z_P=1\quad\hbox{for every }P\mid q
                  \quad\Longrightarrow\quad
                       \mathrm{Unit}_q(\chi(z)).          \tag{QD.3}
\]
Here and below a circle value is included into $\mathbb C$.

The supplied \texttt{ResEta} may be tested against HF9's actual
quadratic character.  Its norm triviality makes $\chi$ trivial
on positive real scalars at infinity.  On each complex factor,
polar decomposition and the elementary classification of circle
characters give an integer $m_\tau$; its infinity value is
$\prod_\tau(z_\tau/|z_\tau|)^{m_\tau}$.  One way to prove that
classification is to lift the continuous map along
$t\mapsto e^{2\pi i t}$ near zero, extend using the homomorphism
law, and observe that its additive lift has integer value at $1$.

The finite integral idelic units form a profinite compact group.
A sufficiently small open subgroup maps into an arc of the
circle containing no nontrivial subgroup, so it maps to one.
The character on the whole compact group consequently has finite
image.  Finite ideal-class number now gives $h>0$, $b\in E^\times$
and an integral finite unit idele $w$ with $z^h$ equal, at finite
places, to $w b$.  If $z_P=1$, taking valuations gives $v_P(b)=0$.
Raising once more to the order $M$ of $\chi(w)$, and using
triviality on principal ideles, gives
\[
 \chi(z)^{hM}\prod_\tau
       \left({\tau b\over|\tau b|}\right)^{M m_\tau}=1,
 \qquad
 \chi(z)^{2hM}=
       \prod_\tau\left({\overline{\tau b}\over\tau b}
                  \right)^{M m_\tau}.                    \tag{QD.4}
\]
Every embedding of the CM field sends its involution to complex
conjugation.  Thus every complex conjugate of the right side has
modulus one.  If $v_P(b)=0$ at all $P\mid q$, both $b$ and $b^{-1}$
are integral over $\mathbb Z_{(q)}$.  Explicitly each has an integer
denominator prime to $q$, obtained by clearing its finitely many
bad primes outside $q$.  Their images in every $\mathbb C_q$ have
norm at most one, hence exactly one.  This applies to all the
factors in (QD.4).  QD1 proves the assertions.  No continuity of
an arbitrary field embedding, and no universal uniqueness theorem
for \texttt{IsEta}, has been used.

\paragraph{QD3: rational unitary determinants and exact local lattices.}
For $U\in U_n(E)$ in the fixed standard Hermitian coordinates,
$\mathrm{Mod}_1(\tau\det U)$ for every chosen embedding $\tau$.
If full $\mathcal O_E$-lattices $M,UM$ are both
\texttt{EqAt} $q$ to $\mathcal O_E^n$, then also
$\mathrm{Unit}_q(\tau\det U)$.

The first assertion follows by taking determinants in $U^*U=1$
and then applying each embedding of the CM field.  For the
second, \texttt{EqAt} supplies integer multipliers prime to $q$
for the inclusions in both directions.  Combining the inclusions
gives integers $s,t$ prime to $q$ for which $sU$ and $tU^{-1}$
have integral entries.  Under every embedding into $\mathbb C_q$,
their determinants have norm at most one, whereas $s,t$ have
norm one.  Hence both $\det U$ and its inverse have norm at most
one.  This proves equality.  The argument needs the two actual
integral lattices, not just equality of their dimensions or their
rational spans.  Conversely, equality after localization at all
$P\mid q$ gives the source's integer-denominator witnesses:
clear the finitely many other prime denominators on generators
in both directions by integers prime to $q$.

\paragraph{QD4: a finite quadratic Gauss sum has finite-order phase.}
Let $B$ be a finite abelian group and $q:B\to\mathbb C^\times$
take root-of-unity values, with $q(0)=1$ and
$b(x,y)=q(x+y)/(q(x)q(y))$ a bicharacter.  Then
$S=\sum_{x\in B}q(x)$ is either zero or a positive real number
times a root of unity.  This holds at $2$ as well as at odd primes.

On the radical $R$ of $b$, $q$ is a character.  If that character
is nontrivial, translation by an element of $R$ proves $S=0$.
Otherwise $q$ descends to $B/R$ and $S=|R|\overline S$.
It is enough to consider the nondegenerate case, of order $N$.
On $\mathbb C^B$ put
\[
 F_{x,y}=N^{-1/2}b(x,y),\qquad T_{x,y}=1_{x=y}q(x).
\]
Character orthogonality gives $F^2f(x)=f(-x)$ and $F^4=1$.
Translation of the sum gives
$\sum_yq(y)b(y,t)=S/q(t)$, so direct matrix multiplication gives
\[
 FTF={S\over\sqrt N}\,T^{-1}F^{-1}T^{-1},
              \qquad (FT)^3={S\over\sqrt N}\,1.           \tag{QD.5}
\]
The matrices $F,T$ are unitary.  Thus $|S|=\sqrt N$.
Their determinants are roots of unity, since $F^4=1$ and the
diagonal entries of $T$ have finite order.  Taking determinants
in (QD.5) shows that $(S/\sqrt N)^N$, and hence $S/\sqrt N$,
is a root of unity.  A nontrivial constant $q(0)$, when present,
can first be divided out.  This proof needs no prime-power
classification or choice of a square root of an algebraic phase.

\paragraph{QD5: phase control for the actual positive intertwiners.}
For three commensurable self-dual period lattices with their
actual compatible semicharacters, the scalar in composing their
positively normalized intertwiners is a root of unity.  Only this
conclusion, rather than a specified order dividing eight, is
needed below.

Choose a common integral lattice $D$ small enough that all three
lattices lie in $D^\#$ and all their semicharacters agree on
$D$.  Multiplication by an even integer achieves the last
condition.  The calculation is then in the finite Heisenberg
module $D^\#/D$.  The normalized intertwiner is a positive
normalization of the sum of the actual operators
$\alpha(x)U(x)$ over the target lattice modulo the intersection.
Their covariance identifies its image with the target theta line;
VI5 proves it is an isometry and is nonzero.  On that line the
composition therefore differs by a scalar of absolute value one.

For completeness, its phase is a finite quadratic Gauss phase.
Expand the trace of the three finite lattice projectors.
Heisenberg orthogonality kills every summand except
$x+y+z=0$.  On this kernel the remaining summand is
\[
 \alpha_A(x)\alpha_B(y)\alpha_C(z)
          e(x,y)e(x+y,z),                                \tag{QD.6}
\]
where $e(v,w)$ is the phase in the \emph{actual} law
$U(v)U(w)=e(v,w)U(v+w)$, with the same convention as VI2
and CA1.  All its values are roots of unity: the vectors are
rational modulo a full integral lattice and their denominators
are bounded.  The semicharacter law shows directly that the
polarization of (QD.6) is a bicharacter.  The omitted trace,
index and normalization factors are positive real numbers.
QD4 therefore proves that the nonzero composition scalar has
finite-order phase.  Its absolute value is one, so it is itself
a root of unity.  Passing through the finite quotient is legitimate
because every projector sum and every trace in this calculation
is finite.  This supplies the phase assertion used here without
assuming equality of positive and signed intertwiners in odd rank.

\paragraph{QD6: the determinant controls a change of theta product basis.}
Let $b=(b_i)$ and $b'=(b'_i)$ be rational orthonormal bases of
the same positive Hermitian $K$-space used in the source's
\texttt{pullB}, and put
$G_{ij}=\phi_0(h(b'_j,b_i))$.  For the same normalized one-variable
theta function, period and un-killed half-point, write $\theta_b$
and $\theta_{b'}$ for the two product theta functions.  There are
$c\ne0$ and a root of unity $\zeta$ such that
\[
 \theta_{b'}=c I_{A_b,A_{b'}}\theta_b,
                         \qquad c^2=\zeta\det G.          \tag{QD.7}
\]

The theta-line theorem gives the first identity.  Both products
have the same positive mean square, the product of the fixed
one-variable mean squares; VI5's isometry therefore gives
$|c|=1$.  First suppose $d=\det(1-G^*)\ne0$.  The coordinate
unitary map $u$ has $\det(1-u)=d$.  Take a common sublattice
$\Lambda'\subset\Lambda$ on which $u\Lambda'\subset\Lambda$
and the two semicharacters agree.  Expand the reproducing theta
kernel, and unfold its integral on a fundamental domain of
$\Lambda'$ by the cosets of $(1-u)\Lambda'$ in $\Lambda$.
For $B=(1-u)^{-1}$ the term belonging to $x$ has integral
\[
 {C_H\over\det(1-u)}\,
       \alpha(x)\exp\{-\pi i\,\operatorname{Im}H(Bx,x)\},
                         \qquad C_H>0.                   \tag{QD.8}
\]
Here $C_H$ is the positive Haar scale of the Gaussian integral.
To verify (QD.8), translate the Gaussian by $Bx$ and complete
the square.  Its real quadratic part is positive definite since
$\operatorname{Re}H((1-u)z,z)=\|z-uz\|_H^2/2$.
Unitary diagonalization reduces the remaining integral to
$\int_{\mathbb C}e^{-\pi(1-\lambda)|z|^2}dz=(1-\lambda)^{-1}$.
Gaussian domination justifies both integration and unfolding.
The remaining finite sum is quadratic: its polarization follows
from
$\operatorname{Im}H(Bx,y)-\operatorname{Im}H(By,x)
=\operatorname{Im}H(x,y)$.  QD4 gives finite-order phase.

There is also an exact way to read this same integral on the
theta line.  Substitute the first identity of (QD.7) into the
kernel.  Orthogonality on the common fundamental domain kills
all nontrivial translate cosets; the zero coset contributes its
positive volume.  Thus $cd$ has finite-order phase.  Unitarity
of $G$ gives
\[
       \overline d=(-1)^r\det(G)d,
       \qquad |d|^2=d\overline d.                         \tag{QD.9}
\]
If $\omega=cd/|d|$ is its root-of-unity phase, (QD.9) gives
$c^2=(-1)^r\omega^2\det G$, proving (QD.7).

Finally choose a rational norm-one $z\in K$ avoiding the finitely
many forbidden eigenvalues, and insert $b''=zb$.
The values $(t+\sqrt{-d})/(t-\sqrt{-d})$, $t\in\mathbb Q$,
give infinitely many choices.  Both consecutive changes now
satisfy the determinant condition.  Compose their identities
and use QD5; their determinants multiply to $\det G$ and all
extra composition scalars are roots of unity.  This proves
(QD.7) with no fixed-vector exception.  The zero-dimensional
case has $c=\zeta=1$.

\paragraph{QD7: rational changes of the period lattice preserve arithmetic phase.}
For rational unitary $U\in U_a(E)$ and $U'\in U_b(E)$ and
self-dual $M,M'$, there is $\epsilon'$ with
$\mathrm{Mod}_1(\epsilon')$ such that
\[
 I_{A_0,\operatorname{La}(M,M')}e_0
          \bigl(\operatorname{tPull}(U,U')z\bigr)
 =\epsilon' I_{A_0,\operatorname{La}(UM,U'M')}e_0(z).      \tag{QD.10}
\]
If $M,M',UM,U'M'$ are all \texttt{EqAt} $q$ to their standard
lattices, then $\mathrm{Unit}_q(\epsilon')$.

Use the source's product-theta presentation of this same $e_0$,
$e_0=\kappa_b I_{A_b,A_0}\theta_b$.  Apply the linear isometry
$X\mapsto U'XU^*$ to $b$ and the lattice.  Finite sums show
that pullback commutes with each intertwiner; the quotient index
and its positive square root are unchanged.  QD6 supplies the
scalar for $b'_i=U'b_iU^*$, and QD5 supplies the three composition
phases.  The prefactor $\kappa_b$ occurs on both sides and cancels.
In particular $\epsilon'$ is the product of that $c$ with roots
of unity and their inverses, exactly the source's
$c\mu_1\mu_3\mu_2^{-1}$ (where $\mu_2^{-1}=\mu_2^7$
when its original eighth-power assertion is used).

The determinant occurring in (QD.7) is
\[
 \phi_0\det\bigl(h(U'b_jU^*,b_i)\bigr)
   =\prod_{\tau:E\hookrightarrow\mathbb C,\ \tau|_K=\phi_0}
          \tau(\det U')^a\tau(\det U^*)^b.                \tag{QD.11}
\]
This follows by vectorizing $X\mapsto U'XU^*$ over $E$ and
taking its determinant after restriction of scalars to $K$;
an orthonormal basis introduces no further Gram determinant.
QD3 controls (QD.11) under every required embedding, and QD1
extracts its square root.  This proves both assertions about
$\epsilon'$.  No evenness of each rank separately is used;
in particular CA5's more restrictive identification of signed
and positive intertwiners has not been invoked.

\paragraph{QD8: real vacuum characters and rational finite moves.}
For the actual supplied Weil action on definite spaces, the
scalar of its real compact group on a fixed nonzero Fock/lattice
vacuum is an integer determinant character.  In orthonormal
coordinates it has the form
\[
 c(g,h)=\prod_\tau\det(g_\tau)^{m_\tau}
                         \det(h_\tau)^{n_\tau},
                         \qquad m_\tau,n_\tau\in\mathbb Z. \tag{QD.12}
\]
The exponents are the winding numbers of the \emph{given}
action; we do not substitute a new action or character for it.

The real compact action preserves the Fock equations, commutes
with finite Heisenberg operators and preserves the lattice
conditions.  Theta-line uniqueness therefore gives a scalar on
this line.  Its Haar norm and continuity make it a continuous
circle-valued character.  Restrict a factor $U(n)$ to the diagonal
torus.  The circle calculation in QD2 gives integer exponents.
Conjugation by permutation matrices makes them all equal.
Every unitary matrix is unitarily diagonalizable and a scalar
character is conjugation invariant, so the character is the
corresponding power of the determinant.  Apply this to each
factor; $U(0)$ contributes one.  This proof applies also to the
negative definite source of a $\chi$ edge.  It avoids the source
route's auxiliary universal \texttt{IsEtaTilde} uniqueness issue.

Let $r,r'$ be the finite components of rational unitary elements,
with their actual coordinate matrices $U,U'$.  Rational automorphy
of the same Weil action and its finite/real commutation give the
finite action as rational pullback followed by the inverse real
vacuum character.  Differentiating this identity in the fixed
frames gives the source's \texttt{JetTh} substitution
\[
 \Theta(z)\longmapsto
       \epsilon\,\Theta(\operatorname{tPull}(U_s,U_t)z).
                                                               \tag{QD.13}
\]
Here $(U_s,U_t)=(U,U')$ for a top edge and $(U',U)$ for a
$\chi$ edge, with its stipulated negation and conjugation.
If $c$ is the scalar before the finite/real rearrangement,
$\epsilon=c^{-1}$ in the top case and $\epsilon=c$ in the
$\chi$ case, since $\overline c^{-1}=c$.  Formula (QD.12) and
QD3 prove $\mathrm{Mod}_1(\epsilon)$ and
$\mathrm{Unit}_q(\epsilon)$ when the four before/after vertex
lattices are standard at $q$.  The actual periods and $c_0$
are unchanged.  Applying QD7 to (QD.13) proves the source's
rational propagation of \texttt{JetKap}, with extra scalar
$\epsilon\epsilon'$ and precisely these four lattice premises.

\paragraph{QD9: the standard-lattice starting scalar is a root of unity.}
Suppose both actual moved vertex lattices are globally standard
in the fixed frames.  Then (QD.2) holds with a root of unity
$\kappa_0$.  Indeed the finite elements $xu_s,yu_t$ preserve
the original full self-dual lattices.  Their stabilizers are
compact open subgroups: after choosing a full rational lattice
frame, this is the compact integral matrix group, with the
closed Hermitian equations imposed.  The same Fock/lattice
vacuum line is preserved.  Its scalar character is continuous
and has values in the circle.  A small open subgroup maps to
an arc containing no nontrivial subgroup; compactness makes
the resulting quotient finite.  The scalar of each element
is consequently a root of unity.

Equivalently one may use the actual smooth-vector principal
level: some positive powers of the two stabilizer elements lie
in it, so a common positive power fixes the nonzero vacuum.
For standard lattices $I_{A_0,A_0}e_0=e_0$.  The fixed
\texttt{kernelNorm\_\allowbreak constant\_\allowbreak spec} then gives (QD.2), with
$\kappa_0$ the scalar or its conjugate according to $p$.
QD1 gives all its complex and local norms.  This argument
does not change the period-normalized $e_0$ into a unit vector.

\paragraph{QD10: the distinguished scalar norm is exactly one.}
Any scalar satisfying (QD.2) obeys $|\kappa|=1$.
Here is the normalization calculation which will remove the
positive factor in a split move.  Put
$N_0=\|e_0\|_{A_0,\mathrm{mean}}^2>0$, with the exact mean-square
normalization of VI5.  The original unit vacuum and its moved
unit vacuum have Gaussian entire restrictions.  Their Fock jet
identities identify those restrictions with
$c_0 e_0$ and $c_0\kappa G_{x,y}$ after the same period scaling.
For a $\chi$ edge conjugate both identities first.

Partition a rational adelic fundamental domain by the finite
lattice cosets and apply its rational quasi-periodicity.  The
finite coset count multiplies both volume and integral, so it
cancels in the normalized mean.  The real restriction is
Gaussian times its entire factor, and the substitution by
$\delta$ changes the Hermitian Gaussian and covolume by the
same coordinate Jacobian.  Thus the Haar unit-norm assertion
is exactly the period-lattice mean-square assertion.  This
can also be checked on one common real fundamental domain:
translate all pieces there; the absolute square of each phase
is one, and the normalized averages agree.  VI5 gives
$\|G_{x,y}\|_{A_{x,y},\mathrm{mean}}^2=N_0$.  Consequently
\[
                   |c_0|^2N_0=1,
                   \qquad |c_0\kappa|^2N_0=1.             \tag{QD.14}
\]
Cancel the first positive quantity to obtain $|\kappa|^2=1$.
All derivatives determine the same entire function: if needed,
their Taylor series agree near zero, then on every complex line
by the identity theorem.  No equality of the two norms is
inferred from jets without this entire-function step.  Formula
(QD.14) is only a distinguished-complex assertion; QD2, QD3
and QD6 supply the arithmetic information independently.

\paragraph{QD11: reading a split move as the same lattice intertwiner.}
Let $e$ be a split-place idempotent and let $s,s'$ be actual
finite unitary elements supported on $e+\bar e$.  Choose their
literal determinant lifts $x_0,y_0$, supported on $e$ (or on
$\bar e$ after the $\chi$-edge conjugation).  If (QD.2) holds
at $x,y$, it holds at $sx,s'y$ with scalar
\[
       r\zeta\,\chi_T(x_0)^{-1}\chi_S(y_0)\kappa,
                 \qquad r>0,\quad \zeta\text{ a root of unity}.
                                                               \tag{QD.15}
\]
This is an identity for the actual action, not merely the
existence of an unrelated vector on the same complex line.

Here is the finite calculation behind it.  Let $\Lambda$ and
$\Lambda'$ be the old and new self-dual finite Hom lattices.
They agree outside $e+\bar e$.  Choose representatives $R,R'$
for their two changed half-lattices modulo their intersections
with the old ones.  The two halves are isotropic.  Character
orthogonality therefore makes
\[
                  F=\sum_{r\in R}\sum_{r'\in R'}
                                      U(r)U(r')v          \tag{QD.16}
\]
a nonzero vector fixed by the new lattice, with the unchanged
real Fock equations.  To check nonvanishing and phase, use the
coinvariant functional of WA5: it is one at $v$, invariant on
the $\bar e$ half, and takes value $1_{\Lambda}(r)$ on the
$e$ half.  The operators in (QD.16) commute on these two new
lattice halves, since their full symplectic pairings have
trivial character.  Choose zero as the representative of the
zero coset.  The functional then gives $|R'|>0$ on (QD.16).
The split clause of the given Weil action gives the positive
multiple of $\chi_T(x_0)^{-1}\chi_S(y_0)$ on this functional.
Vacuum-line uniqueness identifies $\omega(s,s')v$ with (QD.16)
with exactly that phase.  Its norm fixes only the positive
factor, which is kept explicit at this point.

To read (QD.16) in the period model, approximate each
$r+r'$ by a rational Hom vector $\gamma$ modulo
$2(\Lambda\cap\Lambda')$.  Such approximation follows from
the finite Chinese remainder theorem on the finitely many
denominator primes; the rational vector may be chosen in
$\Lambda'$ at every finite place.  The remaining double
lattice vector fixes the old vacuum.  The Heisenberg law
turns its coefficient into the actual semicharacter of the
period vector $P\gamma$.  There is no undetermined phase here:
apply the same calculation to the nonzero new vacuum; since
both half-operators fix it, its period semicharacter must be
the coefficient just obtained.  The pairs $(r,r')$ index
exactly the cosets of the new period lattice modulo its
intersection with the old one.  Injectivity follows by
projecting a putative equality onto the two halves; surjectivity
follows by decomposing a new vector into $e,\bar e$ and the
unchanged complement.  Thus the sum reads as the positive
intertwiner $I_{A_{x,y},A_{sx,s'y}}$, up to a positive real
factor.  Compose it with $I_{A_0,A_{x,y}}$ and use QD5.  This
proves (QD.15).  On a $\chi$ edge perform this calculation on
the conjugated positive model, then conjugate back; the two
characters and the reversed source/target order are exactly
those displayed in (QD.15).

\paragraph{QD12: split propagation is arithmetic at every avoided prime.}
In (QD.15), $r=1$.  Both before and after scalars satisfy QD10;
the root of unity and the two actual circle-character values
also have complex norm one.  Taking norms in (QD.15) therefore
gives $r=1$.

If the supporting base prime does not lie over a rational prime
$q$, the determinant lifts $x_0,y_0$ are one at every $P\mid q$.
QD2 shows that the extra scalar
\[
                       \zeta\chi_T(x_0)^{-1}\chi_S(y_0)
                                                               \tag{QD.17}
\]
has $\mathrm{Mod}_1$ and $\mathrm{Unit}_q$.  The involution on
idempotents preserves this property.  Inverses of a supported
unitary element have the same support, so the conclusion
propagates through any finite word of supported elements and
their inverses.  Here \texttt{splitOff} is the algebraic subgroup
generated by such elements, not a topological closure; no limit
of algebraic numbers is being declared algebraic.

\paragraph{QD13: norm-one corrections with controlled nonsplit denominators.}
Fix a finite set $S$ of base-field primes, a norm-one $u\in E$,
and any required finite precision.  There exists a rational
norm-one $u'$ approximating $u$ at $S$ and approximating $1$
to any prescribed precision at another finite disjoint set.
At every conjugation-fixed prime $P$, $v_P(u')=0$.

For the approximation assertion use the elementary Hilbert--90
formula $u=\alpha/\bar\alpha$.  It follows directly by taking
$\alpha=\theta+u\bar\theta$ for a choice of $\theta$ making
this expression nonzero; if all such expressions vanished,
conjugation would be the identity.  The split local case is the
same explicit ratio of two components.  At the finitely many
places choose nonzero local $\alpha_v$ giving the desired
norm-one value.  Field weak approximation gives one
$\alpha\in E$ sufficiently close to all of them; the map
$\alpha\mapsto\alpha/\bar\alpha$ is continuous there.
The valuation assertion follows from
$v_P(\bar\alpha)=v_P(\alpha)$, including at ramified primes.

This gives the source's level correction for a unitary
quasireflection.  Write
\[
 R_{v,u}=1+(u-1)P_v,\qquad
              P_v(w)=v\,{h(w,v)\over h(v,v)}              \tag{QD.18}
\]
with the linear-variable convention adjusted to the actual
Hermitian form.  Its inverse is $R_{v,u^{-1}}$.  The projector
$P_v$ has integral lattice entries outside finitely many primes.
Choose $u'$ close to $u$ at the specified level primes and close
to $1$ at every other nonsplit prime where $P_v$ has a denominator.
At the latter places both $(u'-1)P_v$ and
$(u'^{-1}-1)P_v$ are integral; at every other nonsplit place
this follows from $v_P(u')=0$.  Thus the corrected reflection
and its inverse preserve the lattice outside the level primes
and split primes.  Its difference from the original reflection
is as small at the level primes as prescribed.  The finite
local inclusions can be expressed with the source's global
integer multipliers prime to those primes, by clearing
denominators outside them.  This proves the support control
in \texttt{normOne\_\allowbreak approx} and in the reflection correction
on the actual definite spaces, without a prime-existence or
strong-approximation assumption.

\paragraph{QD14: rational, split and stabilizer factorization.}
Let $g$ be a finite unitary element of an actual definite vertex,
let $\Lambda$ be its fixed full lattice, and let $Q$ be a finite
set of rational primes.  There is an exact factorization
\[
                 g=r s k,                               \tag{QD.19}
\]
where $r$ is the finite component of a rational unitary element,
$s$ is a finite product of split-supported elements away from
$Q$, and $k$ stabilizes $\Lambda$ everywhere.  At primes over
$Q$, $k$ can satisfy any prescribed principal congruence.

First recall why rational unitary weak approximation at finitely
many finite places is available here.  Multiply the desired local
unitaries by one rational norm-one scalar avoiding their finitely
many forbidden eigenvalues.  On this open set the Cayley map
\[
      X\longmapsto(1-X)(1+X)^{-1},\qquad X^*=-X,           \tag{QD.20}
\]
is a rational chart for the unitary group.  Its inverse has the
same form.  The skew-Hermitian matrices form an $F$-vector space;
approximate its coordinates by elements of $F$.  The nonvanishing
determinant conditions remain open, and the Cayley map preserves
the required precision.  Undoing the scalar gives a rational
$\gamma_0$ with $\gamma_0 g$ arbitrarily close to one at the
chosen places.  This is just finite field weak approximation.

Choose those places to include all primes over $Q$ and every
place where $g$ or $g^{-1}$ fails to preserve $\Lambda$.
Factor $\gamma_0$ into finitely many unitary quasireflections,
as in WA10.  Apply QD13 to their norm-one parameters, with
enough precision that the product $\gamma$ still has
$\gamma g$ integral and invertible at the chosen places and
has the requested congruence over $Q$.  At other nonsplit
places every corrected factor and its inverse preserve
$\Lambda$.  Thus $\gamma g$ can fail to preserve $\Lambda$
only at a finite set of split base primes away from $Q$.

Let $s$ agree with $\gamma g$ at precisely those split primes,
and be one elsewhere; let $k$ be one there and agree with
$\gamma g$ elsewhere.  Each support cut is fixed by conjugation,
so both matrices are unitary; they commute and multiply to
$\gamma g$.  The matrix $k$ and its inverse preserve the
lattice at every place.  The rational lattice is the intersection
of its localizations, so $k$ is in the actual adelic stabilizer.
The element $s$ is the finite product of its single split-pair
pieces.  Taking $r=\gamma^{-1}$ proves (QD.19).
This argument proves the definite-vertex instance needed here;
it does not label every more general source approximation
theorem as reconstructed.

\paragraph{QD15: propagating one scalar while protecting a finite prime set.}
Fix finite $Q$ such that both original vertex lattices at $x,y$
are \texttt{EqAt} standard at every $q\in Q$.  There exists
$\kappa_Q$ satisfying (QD.2), with $\mathrm{Mod}_1(\kappa_Q)$
and $\mathrm{Unit}_q(\kappa_Q)$ for all $q\in Q$.

Apply QD14 separately to $xu_s$ and $yu_t$ and write
$x=r s x_0$, $y=r's'y_0$, with
$x_0=k u_s^{-1}$ and $y_0=k'u_t^{-1}$.
The actual fixed-frame lattices of $x_0,y_0$ are globally standard.
QD9 supplies their root-of-unity scalar.  Apply QD12 to the
finite words $s,s'$, using one factor on one side at a time.
All intervening words act trivially at $Q$, so the two resulting
lattices remain standard there.  All scalars acquired thus far
have the required units at $Q$.  Finally apply QD8 and QD7
to $r,r'$.  The before lattices are standard at $Q$ by construction,
and the after lattices are the original ones, standard at $Q$
by hypothesis.  The four lattice premises of QD3 are therefore
present.  The product of the scalars proves the assertion.
All operations take place on the fixed original frames and on
the exact actual lattice $A_{x,y}$.

\paragraph{QD16: uniqueness gives the simultaneous all-prime scalar.}
The scalar in (QD.2) is unique.  The function $G_{x,y}$ is entire
and nonzero by theta-line uniqueness and VI5's invertibility of
$I_{A_0,A_{x,y}}$.  Hence at least one Taylor coefficient is
nonzero: otherwise its Taylor series is zero near the origin,
and restriction to each complex line makes it zero everywhere.
Choose such a multi-index $\alpha$.  Its multi-factorial,
$P_\alpha$, and $c_0$ are nonzero.  Comparing (QD.2) for two
scalars, cancel these factors and use injectivity of
$\operatorname{opp}_p$ to get equality.

Take $\kappa=\kappa_{\varnothing}$ from QD15.  For each rational
prime $q$ where both original vertex lattices are standard,
apply QD15 with $Q=\{q\}$.  Its scalar satisfies the same
all-jet identity for the same $G_{x,y}$, so it equals $\kappa$.
Consequently one and the same scalar has
\[
 \mathrm{Mod}_1(\kappa),\qquad
 \mathrm{Unit}_q(\kappa)\quad
  \hbox{at every such }q,\qquad \texttt{JetKap}(x,y,\kappa).
                                                               \tag{QD.21}
\]
This proves the actual conclusion of \texttt{exists\_\allowbreak kappa}.
There is no compactness argument or infinite simultaneous choice
of unrelated scalars.  In particular the selected scalar is
independent of the finite decomposition used to obtain it.

\paragraph{QD17: the literal moved-vacuum dictionary and its eight clauses.}
Choose $L_s=L_s(x,y)$ and $L_t=L_t(x,y)$, and take $\kappa$
from QD16.  These choices prove \texttt{movedVacuum\_\allowbreak dict},
including its geometric and analytic clauses.

The actual moved rational lattices are recovered by intersecting
the finite adelic lattice with the rational vector space in the
unchanged frame.  Duality commutes with this operation: test on
finite generators locally, use the Hermitian inverse of the
unitary element, and intersect the integral local conditions.
They are therefore full self-dual lattices.  Acting on rank-one
Hom generators sends the old Hom lattice to the Hom lattice of
these two moved lattices; linearity and closure give the exact
orbit-lattice inclusion required in the source.  The actual Weil
covariance then gives the moved finite lattice invariance.  Its
finite components commute with real Fock operators and preserve
the fixed real Cartans.  The exact Haar norm is preserved.  Thus
the actual moved vector satisfies the source's full
\texttt{IsVacO}, including its norm condition, on these lattices.

VI2 gives their period-lattice and semicharacter identifications;
VI5 gives the nonzero entire function $G_{x,y}$.  QD16 supplies
its all-conjugate scalar, all local-unit clauses and all jets.
For a top edge the two \texttt{EqAt} conclusions are the two
given vertex conclusions; for a $\chi$ edge they are interchanged,
not replaced by a new frame.  These verify every conclusion of
the literal record.  The source's \texttt{Big} wrapper may take
both extra size thresholds to be zero, because this construction
uses only the supplied valid choice and has no further large-size
condition.  This does not assert that such a choice exists for
arbitrary raw arithmetic data.

\paragraph{QD18: exact size of the split-half quotient.}
Let $\Lambda$ be the actual full self-dual finite Hom lattice
and let $e$ be the orthogonal sum of the split idempotents in
the finite set $P=\operatorname{locPl}(\mathfrak c.G,t)$.
For $k>0$ put $M=\Lambda\cap\bar e\mathcal W_f$ and
$Y=k^{-1}\Lambda\cap\bar e\mathcal W_f$.  Then $Y/M$ is
finite, of order $n>0$, and
\[
 n=\left(\prod_{v\in P_F}|\mathcal O_{F_v}/k\mathcal O_{F_v}|
                                                \right)^{ab},
                           \qquad n\mid k^{f a b}.        \tag{QD.22}
\]
Here $P_F$ contains each supporting base prime once.  The source's
split-place disjointness condition excludes selecting both
conjugate halves, so the passage from $P$ to $P_F$ is injective.

At a split prime $E_v=F_v\times F_v$, projection to either half
identifies the Hom lattice with an integral matrix lattice of
rank $ab$ over $\mathcal O_{F_v}$.  The two original full
lattices are free over that DVR; changing their local bases
does not change the quotient cardinality.  Orthogonal support
idempotents give a direct product over $P_F$.  Multiplication
by $k$ identifies $Y/M$ with $M/kM$, proving the equality.
Finally the local ideal norm product gives
\[
       \prod_{v\text{ finite}}|\mathcal O_{F_v}/k\mathcal O_{F_v}|
            =|\mathcal O_F/k\mathcal O_F|=k^{[F:\mathbb Q]}.
                                                               \tag{QD.23}
\]
Every omitted factor is a positive integer and equals one away
from $k$.  Thus the partial product divides the full product.
This proves the precise exponent $f a b$, including all
ramification multiplicities; replacing it by a guessed residue
depth or by $[E:\mathbb Q]ab$ would lose the source's bound.
For $k=1$, an empty set of places, or $ab=0$, the quotient is
trivial and $n=1$.

\paragraph{QD19: the actual coset indicator has bounded Fourier coefficients.}
Let $d$ be the given \texttt{FinDatum} with
$\texttt{IsCosetDatum}(k,d)$, $k>0$, $d.z=1$ and
$d.P=\operatorname{locPl}(\mathfrak c.G,t)$.
Its vector is a finite Heisenberg combination of the actual
vacuum, with $n$ as in QD18:
\[
       \mathfrak c.S.Z.\operatorname{vec}(d)
                   =\sum_{i=1}^n c_i U(w_i)v_p,
 \quad kw_i\in\Lambda,\quad(w_i)_\infty=0,
 \quad c_i={\zeta_i\over n},\quad\zeta_i^{2k^2}=1.        \tag{QD.24}
\]
Consequently
\[
 k^{f a b}c_i\text{ is an algebraic integer},\qquad
       \sum_i|\rho(c_i)|=1
       \quad\text{for every }\rho:\overline{\mathbb Q}\to\mathbb C.
                                                               \tag{QD.25}
\]

To verify the formula, write the actual coset as $x_0+k\Lambda_e$,
with $kx_0\in\Lambda_e$, and choose representatives $r_i$ of
$k^{-1}\Lambda_{\bar e}/\Lambda_{\bar e}$.  Self-duality makes
the pairing between the two finite half-quotients perfect.
The sum of all its characters is $n$ on the zero coset and
zero elsewhere, by translation by a character taking a nontrivial
value on a nonzero coset.  Thus
\[
 1_{k\Lambda_e}(t)=\frac1n\sum_i
              \psi(\operatorname{sympl}(t,r_i))1_{\Lambda_e}(t).
                                                               \tag{QD.26}
\]
Apply the actual normalized half-Schwartz-model map, which sends
$1_{\Lambda_e}$ to $v_p$, and then translate by $-x_0$ in its
specified convention.  The Heisenberg multiplication law gives
$w_i=-x_0+r_i$ and
$\zeta_i=\psi(\operatorname{symplOH}(-x_0,r_i))$.
Since $2\operatorname{symplOH}=\operatorname{sympl}$ and
$kx_0,kr_i\in\Lambda$, integrality of the full symplectic
pairing gives $\zeta_i^{2k^2}=1$.  This proves (QD.24).
The normalizing map and $d.z=1$ are essential: an additional
Weil translate of the datum is not silently dropped.  QD18
and integrality of roots of unity prove the first part of
(QD.25); every conjugate of $\zeta_i$ has modulus one, so the
second is the sum of $n$ terms $1/n$.

\paragraph{QD20: finite torsion Heisenberg translates have the same period jets.}
For the lattices and scalar of QD17, let $w$ satisfy $kw\in\Lambda$
in the original finite Hom lattice.  There are a root of unity
$\zeta\in\overline{\mathbb Q}$ and a period vector $s$ such that
$ks\in A_{x,y}$ and the actual jets of $U(w)v_p$ at $xu_s,yu_t$
are the jets of
\[
                         \kappa\zeta U(s)G_{x,y},         \tag{QD.27}
\]
with exactly the outside factor $c_0\operatorname{opp}_p(P_\alpha\cdot)$.

First use the actual Weil covariance to move $w$ to
$w'=\operatorname{act}(xu_s,yu_t)w$.  The Hom-lattice inclusion
of QD17 gives $kw'$ in the moved lattice.  Finite rational
approximation supplies a rational Hom vector $M$ with
$u=w'-M_f$ in that lattice.  Then $kM_f=kw'-ku$ is in it too.
The actual lattice-vacuum law gives $U(u)F=\epsilon F$ with
$\epsilon\in\{1,-1\}$.  The two Heisenberg factors therefore
give
$U(w')F=\zeta U(M_f)F$, where
$\zeta=\epsilon\psi(-\operatorname{symplOH}(M_f,u))$.
Since $kM_f,u$ are integral and
$2\operatorname{symplOH}=\operatorname{sympl}$,
$\zeta^{2k}=1$.  This is the finite-rational decomposition in
the source, not a replacement of the adelic vector by an
arbitrary real vector.

Rational quasi-periodicity replaces $M_f$ by its negative real
component.  In the fixed orthonormal coordinates put
\[
 X_\tau=-\,\operatorname{fr}_{T,\tau}^{-1}M_{\infty,\tau}
                         \operatorname{fr}_{S,\tau},
                              \qquad s=\delta X.           \tag{QD.28}
\]
The rational-lattice/period-lattice identification gives
$ks\in A_{x,y}$, with this sign.  Write the real restriction
$F=\mathrm{Gauss}\,\Phi$.  The exact Fock monomial formula and
(QD.2) give $\Phi(Z)=c_0\kappa G_{x,y}(\delta Z)$ after
conjugation in the $\chi$ case.  A real Heisenberg translate
is
\[
 \Phi(Z+X)\exp\{-\pi(2\xi_R\langle Z,X\rangle
                                      +\xi_R\langle X,X\rangle)\}.
                                                               \tag{QD.29}
\]
The source's scaled torus Hermitian form turns (QD.29) exactly
into $U(\delta X)G_{x,y}$; its periods satisfy
$\operatorname{period}(\delta,\tau)
=\delta/\sqrt{2\pi\xi_R}$.  Differentiating gives (QD.27).
For the negative source conjugate the positive calculation and
conjugate back, retaining $\operatorname{opp}_p$ outside the
formula.  All functions differentiated here are entire; the
Heisenberg exponential and translation preserve that property.

\paragraph{QD21: the full quantitative dictionary on the original vector.}
Use QD19's $c_i,w_i$ and apply QD20 to each $w_i$, obtaining
$\zeta_i',s_i$.  Define
\[
             r_i=\operatorname{opp}_{p,\overline{\mathbb Q}}(c_i)
                                         \zeta_i'.         \tag{QD.30}
\]
Then the same $L_s,L_t,\kappa$ of QD17 and these $r_i,s_i$
satisfy every clause of \texttt{DictAt}.  Explicitly,
\[
 \begin{gathered}
  k s_i\in A_{x,y},\qquad
  k^{f a b}r_i\text{ is integral over }\mathbb Z,\qquad
  \sum_i|\rho(r_i)|\le1\quad(\rho:\overline{\mathbb Q}\to\mathbb C),\\
 \operatorname{jetTab}(\operatorname{vec}(d),xu_s,yu_t;\mathrm{mon})
  =c_0\operatorname{opp}_p\left(
        P_\alpha\,\partial^\alpha
          \left[\kappa\sum_i r_iU(s_i)G_{x,y}\right](0)\right),
 \qquad \alpha=\mathrm{mon}.\operatorname{mapDomain}
                                      (\operatorname{fvar}_p).
 \end{gathered}                                            \tag{QD.31}
\]
The all-prime local-unit assertion for $\kappa$ and the two
local lattice assertions are exactly those of QD17.
In fact this uncollected presentation has $r_i=\nu_i/n$ with
$\nu_i^{2k^2}=1$: combine QD19's order $2k^2$ with QD20's
order $2k$.  Thus its coefficients lie in the fixed cyclotomic
field $\mathbb Q(\mu_{2k^2})$, even as the adelic points vary;
this extra observation does not assert such a field for $\kappa$.

Indeed multiplication by a root of unity and the algebraic
involution preserve the integrality in (QD.25).  For a complex
embedding $\rho$, the norm of each $\rho(\zeta_i')$ is one,
so the sum of norms in (QD.31) equals that in (QD.25), evaluated
at $\rho\circ\operatorname{opp}_{p,\overline{\mathbb Q}}$.
Linearity of the given Weil action and the Fock differential
operators applies to (QD.24).  Finite Heisenberg translates
with zero infinite component commute with these real differential
operators.  Apply QD20 term by term.  Move the finite entire
sum through $\partial^\alpha$; the two occurrences of
$\operatorname{opp}_p$ on $c_i$ cancel.  This proves the exact
jet identity.  Omitting the conjugation in (QD.30) would give
the wrong $\chi$-edge coefficients.  No integral, infinite sum
or arbitrary field embedding is interchanged with differentiation.

\paragraph{QD22: the actual Big endpoint and the field boundary.}
QD21 proves \texttt{DictionarySizeBound.dictAt\_\allowbreak big} under
its original literal input hypotheses.  More precisely, it is
valid for every supplied $\mathfrak c.\mathrm{Is}$ in that
source context; QD17 and its finite-translation extension use
both additional size thresholds equal to zero.  Thus the
\texttt{Big} quantifier is satisfied by those same witnesses,
with no growth estimate hidden in the word ``big.''  The
quantitative content is instead the exact torsion level $k$,
denominator $k^{f a b}$, all-conjugate coefficient sum at most
one, and the simultaneous local-unit conclusion.  The finite
presentation even has the positive number of terms $n\mid
k^{f a b}$ obtained in QD18, though the literal \texttt{DictAt}
does not request a separate bound on $n$.

For a fixed finite collection of such dictionaries the finitely
many algebraic coefficients and scalars can of course be put
in one finite extension of the already embedded coefficient
field: adjoin them inside the original algebraic closure, then
use the tower law.  This observation does not produce a single
field for infinitely many arbitrary adelic points, nor does it
put the critical toric ratios in the originally prescribed field.
The surviving \texttt{ToricRatioConjugates}/\texttt{NFRatio}
and original-embedding \texttt{Choice.FieldOK} suppliers remain
separate targets.  The universal uniqueness of the auxiliary
quadratic character is also not asserted.  The real-character
argument of QD8 and CW24's prescribed-character construction
allow the present endpoint to avoid that unproved uniqueness
route while retaining the exact supplied character restrictions.


\section{The actual large-gap discrete series and the real-stage join}
\label{sec:discrete-series-core}

This section constructs the discrete series needed by the later row parameters.
It does not assert the source's universal \texttt{DiscreteSeriesExist} or
\texttt{FormalDegreeFormula\(_2\)}.  The construction below uses finite compact
representation theory, integration and Hilbert-space arguments.  The degree
calculation in DC18--DC22 has one additional, explicitly unreconstructed
classical premise: Harish--Chandra's formal-degree polynomial theorem, stated
precisely in DC18.  Its Haar constant is calculated here.  The volume part of
the final join retains GV's named Tamagawa theorem.  No statement below
constructs an entire \texttt{Core} from a volume or a representation alone.

Number coordinates by $0,\ldots,n-1$.  Write
\[
 G_\epsilon=\{g:g^*J_\epsilon g=J_\epsilon\},\qquad
 K_\epsilon=G_\epsilon\cap U(n),\qquad
 p=\#\epsilon^{-1}(+),\quad q=\#\epsilon^{-1}(-),\quad q_0=\min(p,q).
\]
The positive roots are $e_a-e_b$ with $a<b$.  Separate them into compact and
noncompact roots according as $\epsilon_a=\epsilon_b$ or not, and let
$\rho_c,\rho_n$ be their half-sums.  Thus
$\rho=\rho_c+\rho_n$, $\rho_i=(n-1)/2-i$.  Inner products on Hilbert spaces
are conjugate linear in the first argument.  A parameter is \emph{regular
integral} here precisely when
\[
 \lambda_0>\cdots>\lambda_{n-1},\qquad
 \lambda_i-(n-1)/2\in\mathbb Z.                                  \tag{DC.1}
\]
Put $\Lambda=\lambda+\rho_n-\rho_c$,
$c_\lambda=\|\lambda\|^2-\|\rho\|^2$ and
$c_K(\nu)=\|\nu+\rho_c\|^2-\|\rho_c\|^2$.  A noncompact gap $h$ means
$|\lambda_a-\lambda_b|\ge h$ whenever $\epsilon_a\ne\epsilon_b$.
For the constructive existence assertions we require
\[
 h>H_\epsilon:=\max\{npq+1,\ 16nq_0+8\sqrt{2q_0}\}.             \tag{DC.2}
\]
The smaller bound $h>npq+1$ suffices for the rigidity assertions once a
model has been constructed.  Compact signatures are treated separately;
there is no division by $q_0$ in that case.

\paragraph{DC1: the exact compact type and its integral construction.}
\label{claim:DC1}
The weight $\Lambda$ is integral and dominant on each sign block.  Indeed
$\lambda+\rho$ is integral and
$\Lambda=\lambda+\rho-2\rho_c$.  If $a<b$ are consecutive indices in one
sign block, and $r$ opposite-sign indices lie between them, then
\[
 \Lambda_a-\Lambda_b=\lambda_a-\lambda_b+r-1\ge0.
\]
The last inequality follows from (DC.1): consecutive distinct entries of
$\lambda$ differ by a positive integer.  For an integral dominant weight
$\alpha$ of $U(d)$, subtract its last entry to get a partition.  The Young
symmetrizer applied to the corresponding tensor power of $\mathbb C^d$,
followed by the determinant power $\det^{\alpha_{d-1}}$, gives the compact
irreducible representation of highest weight $\alpha$.  Averaging its
positive inner product over $U(d)$ makes it unitary.  The usual elementary
highest-weight argument proves irreducibility: a highest vector is a scalar
multiple of the row-symmetrized, column-alternating tableau vector, and
lowering matrix units generate every tableau.  Apply this independently
to the two blocks.  We obtain a finite-dimensional unitary $K_\epsilon$-module
$T$, with a unit highest vector $w$, the exact object called \texttt{KType}
in the source.  This uses no noncompact representation-existence theorem.

\paragraph{DC2: the principal-minor function and the normalized average.}
\label{claim:DC2}
Set $\mu_i=\lambda_i+\rho_i\in\mathbb Z$, put $\mu_n=0$, and let
$\Delta_j(z)$ be the leading $j$ by $j$ minor.  On the open set where these
minors are nonzero define
\[
 \Phi_\mu(z)=\prod_{j=1}^n\Delta_j(z)^{-(\mu_{j-1}-\mu_j)},
 \qquad
 P(z)=d_T\int_K\Phi_\mu(k^{-1}zk)\,T(k)P_wT(k)^{-1}\,dk,         \tag{DC.3}
\]
where $d_T=\dim T$, $dk$ is probability Haar and $P_w v=w\langle w,v\rangle$.
Every leading minor of a positive definite matrix is positive.  Compactness
of $K$ therefore gives an open complex neighborhood of the positive cone
on which the second integral is holomorphic; all its derivatives may be
taken under the integral on a smaller neighborhood.  Triangular
multiplication gives
\[
 \Phi_\mu(lzu)=\prod_i(l_{ii}u_{ii})^{-\mu_i}\Phi_\mu(z)
                                                                  \tag{DC.4}
\]
for lower triangular $l$ and upper triangular $u$.  Schur averaging gives
$d_T\int_K T(k)P_wT(k)^{-1}dk=1$, hence $P(1)=1$.
For $X=X^*$ with $JX=-XJ$, one has $e^{-X}=Je^XJ$; diagonal conjugation
by $J$ fixes every leading minor.  Consequently $P(e^{-X})=P(e^X)$.
The source formula with an inverse matrix in the argument of $\Phi$ is
this identity, not a reciprocal of $\Phi$.

\paragraph{DC3: the compact differential identities and the global kernel.}
\label{claim:DC3}
For every block-diagonal complex matrix $Y$, the average satisfies
\[
 dP_z(zY)=-dT(Y)P(z),\qquad dP_z(Yz)=-P(z)dT(Y).                \tag{DC.5}
\]
Here is the integration-by-parts calculation that supplies the essential
weight shift.  Write $B=k^{-1}Yk$, and decompose $B$ into its strictly lower,
diagonal and strictly upper parts inside the two blocks.  Differentiating
(DC.4) gives the diagonal contribution $-2\sum_i\mu_i B_{ii}$ to the
variation in $Yz+zY$.  Moving the lower-minus-upper variation onto
$T(k)P_wT(k)^{-1}$ differentiates the two representation factors.
The integral of a compact left-invariant derivative is zero.  Differentiating
its coefficient $B_{ab}$ contributes
\[
 \sum_{a\ne b,\,\epsilon_a=\epsilon_b}
       \operatorname{sgn}(a-b)(B_{aa}-B_{bb})
       =-4\sum_a(\rho_c)_aB_{aa}.
\]
The remaining diagonal action is exactly $-2\sum_i\Lambda_iB_{ii}$,
because $\mu=\Lambda+2\rho_c$.  Thus
$dP_z(Yz+zY)=-dT(Y)P(z)-P(z)dT(Y)$ for Hermitian block $Y$.
Differentiating conjugation equivariance supplies the difference identity;
adding and subtracting gives (DC.5) for skew-Hermitian block $Y$.
Complex linearity and holomorphic continuation from the real positive
cone give the asserted complex identities.

The polar decomposition $g=k e^X$, $k\in K$, $X=X^*$, $JX=-XJ$, is unique.
Define
\[
             \Psi(g)=T(k)P(e^X).                               \tag{DC.6}
\]
Then $\Psi(k_1gk_2)=T(k_1)\Psi(g)T(k_2)$ and $\Psi(1)=1$.
For a fixed $g_0=k_0e^{X_0}$, a local logarithm of $k(g)^{-1}k_0$ and
(DC.5) show that, for $g$ sufficiently close to $g_0$,
\[
                  \Psi(g)=T(k_0)P(g^{-1}k_0).                  \tag{DC.7}
\]
Explicitly, the derivative of $T(e^{tY})P(ze^{tY})$ is zero.
Integrate this identity along the short compact logarithm path and use
$P(e^{-X})=P(e^X)$.  The right side of (DC.7) is a holomorphic function
of an ambient invertible complex matrix near $g_0$.  This proves the
literal local ambient smoothness required by the source.

\paragraph{DC4: central differential operators have the actual parameter.}
\label{claim:DC4}
Let $E_{ab}$ denote the complexified right derivatives in the source
convention.  The kernel has
\[
 \sum_{a,b}E_{ab}E_{ba}\Psi=c_\lambda\Psi,\qquad
 \mathcal C\bigl(Y+(n-1)/2\bigr)\Psi
       =\prod_i(Y+\lambda_i)\Psi.                              \tag{DC.8}
\]
In the second formula the literal column-determinant convention is
\[
 \mathcal C(u)=\sum_{\sigma\in S_n}\operatorname{sgn}(\sigma)
   \prod_{j=0}^{n-1}\bigl(E_{\sigma(j),j}
                      +\delta_{\sigma(j),j}(u-j)\bigr),
\]
with factors composed in increasing $j$ from left to right (the
rightmost factor acts first).  Thus in rank one it is $Y+E_{00}$;
for the trivial compact representation in rank two, $\lambda=\rho$
and $u=Y+1/2$, it is $(Y+1/2)(Y-1/2)$.
These are actual scalar differential identities, not a classification
assumption.  To check them, use (DC.7), differentiate under the compact
integral, and move conjugations through the invariant contracted matrix
units.  On $\Phi_\mu$, strictly lower left derivatives vanish and diagonal
left derivatives are $-\mu_i$ by (DC.4).  Commuting an upper derivative
past a lower one uses
$[E_{ab},E_{cd}]=\delta_{bc}E_{ad}-\delta_{ad}E_{cb}$.
For the quadratic sum the diagonal result is
$\sum_i\mu_i^2-2\sum_i\rho_i\mu_i
=\|\lambda\|^2-\|\rho\|^2$ after the inversion in (DC.7).
For the determinant, expand by columns and commute each strictly lower
factor to the right.  Every off-diagonal commutator cancels the term of
the permutation differing by that transposition; its diagonal commutator
in column $i$ contributes the shift $n-1-i$.  Only the diagonal product
survives.  Substituting $\mu_i=\lambda_i+(n-1)/2-i$ and reversing left
derivatives through inversion yields the displayed product.  The same
calculation with the first $r$ columns proves the identity inductively,
so it is valid in every rank.  These coefficients generate the center:
their highest PBW symbols are the coefficients of the ordinary
characteristic polynomial, which generate conjugation-invariant
polynomials on matrices; subtracting that leading symbol and inducting
on the filtration proves the assertion.  In particular the full infinitesimal
parameter of this kernel is the unordered list $\{\lambda_i\}$; the
quadratic eigenvalue alone would not establish that fact.

\paragraph{DC5: a logarithmic minor estimate with explicit constants.}
\label{claim:DC5}
If $w_j\ge0$, $\sum w_j=1$, $\sum w_jy_j=0$ and $|y_j|\le R$, then
\[
 \log\sum_jw_je^{y_j}\ge
       \tfrac14\min(1,R^{-1})\sum_jw_jy_j^2.                   \tag{DC.9}
\]
Take $c=\min(1,R^{-1})$, use $e^u\ge1+u+u^2/3$ for $|u|\le1$,
then $\log(1+t)\ge3t/4$ for $0\le t\le1/3$.  Jensen's inequality
for the concave power $x^c$ passes from $cy_j$ to $y_j$ and proves
(DC.9).  If $R=0$ both sides vanish.

Let $Y=k^{-1}Xk$ with $X$ as in DC3 and put $r^2=\operatorname{tr}(X^2)$.
Diagonalize $Y=U\operatorname{diag}(y_i)U^*$.  Cauchy--Binet gives
\[
 \Delta_m(e^Y)=\sum_{|I|=m}|\det U_{[m],I}|^2
                              \exp\bigl(\sum_{i\in I}y_i\bigr).
\]
The coefficients sum to one.  The weighted mean is
$\operatorname{tr}(Y_{[m],[m]})=0$; its variance is
$\sum_{b<m\le a}|Y_{ab}|^2$.  The mean and variance identities follow by
expanding respectively the first and second derivatives at zero of
$\det((e^{tY})_{[m],[m]})$; the second derivative is
$\operatorname{tr}((Y^2)_{[m],[m]})-\operatorname{tr}(Y_{[m],[m]}^2)$.
The rank of $Y$ is at most $2q_0$, so
$|\sum_{i\in I}y_i|\le\sqrt{2q_0}\,r$ by Cauchy--Schwarz.  Hence
\[
 \log\Delta_m(e^Y)\ge\tfrac14
  \min\{1,(\sqrt{2q_0}r)^{-1}\}\sum_{b<m\le a}|Y_{ab}|^2.       \tag{DC.10}
\]
At $X=0$ this is read by continuity.

\paragraph{DC6: decay, square integrability and the exact large-gap bound.}
\label{claim:DC6}
For $q_0>0$, (DC.10) telescopes in the exponents of (DC.3).
The determinant term is one because $\operatorname{tr}X=0$.  Moreover
\[
 \sum_{m<n}(\mu_{m-1}-\mu_m)
                 \sum_{b<m\le a}|Y_{ab}|^2
 =\sum_{b<a}(\mu_b-\mu_a)|Y_{ab}|^2\ge\tfrac h2 r^2.
\]
Only opposite-sign entries of $Y$ are nonzero.  Since the integrand
projectors are positive and their normalized average is the identity,
there is no extraneous $\dim T$ in the operator bound:
\[
 \|\Psi(ke^X)\|\le
 \exp\left[-\tfrac h8\min\{r^2,r/\sqrt{2q_0}\}\right].        \tag{DC.11}
\]
For completeness the polar Jacobian is the determinant of
$\sinh(\operatorname{ad}X)/\operatorname{ad}X$ on the noncompact tangent
space.  Differentiating $k e^X$ proves this formula, including its
continuous value at zero.  Diagonalizing $X$, whose nonzero eigenvalues
occur in opposite pairs, and using $\sinh(t)/t\le e^{|t|}$ bounds it by
$\exp(n\sqrt{2q_0}\,r)$.  Polar Lebesgue measure is ordinary measure
on $\mathbb C^{pq}$; its shell factor is polynomial in $r$.
Thus (DC.11) is integrable for $h>16nq_0$, is square integrable for
$h>8nq_0$, and its product with $e^{sr}$ is square integrable when
$h>8nq_0+8s\sqrt{2q_0}$ (a sufficient bound).  In particular (DC.2)
supplies all unweighted integrability needed here.  Any other Haar
measure is a positive scalar multiple and has the same finiteness
properties.

\paragraph{DC7: the kernel entries are genuinely smooth vectors in $L^2$.}
\label{claim:DC7}
Fix $v_0\ne0$ in $T$ and define
$j(v)(g)=\langle v_0,\Psi(g)v\rangle$.
The finite-dimensional space $j(T)$ consists of smooth $L^2$ functions,
is stable under the compact right derivatives, and has Casimir
$c_\lambda$ by DC3--DC6.  These facts imply that \emph{every} right
derivative word of $j(v)$ lies in $L^2$.

Here is the required global regularity argument.  Choose orthonormal real
compact and noncompact Lie bases $Y_i,X_j$ for the trace form.  On smooth
functions the positive elliptic operator is
$L=-\sum_iY_i^2-\sum_jX_j^2=2C_K-C_G$, with the conventional compact
$C_K=-\sum_iY_i^2$.  Use compactly supported cutoffs $\chi_R$ on the
complete invariant metric, equal to one on the radius-$R$ ball, with
$|d\chi_R|\le C/R$.  Integration by parts for $\chi_R f$, followed by
Cauchy--Schwarz and absorption of half the derivative term, bounds
$\sum\|\chi_R X_jf\|_2^2$ by a fixed multiple of
$\sum\|Y_if\|_2^2+(1+|c_\lambda|)\|f\|_2^2$ and a vanishing cutoff
error.  No derivative integrability was assumed in this compactly
supported calculation.  Fatou gives the first noncompact derivatives
in $L^2$.  The space of words of length at most $m$ applied to $j(T)$
is finite-dimensional and compact-stable, and the central Casimir
still has eigenvalue $c_\lambda$.  Repeating the same argument proves
the assertion by induction on $m$.
Finally the fundamental theorem of calculus along a one-parameter
subgroup, now in $L^2$, identifies these derivatives with the generators
of the right regular representation.  Iterating proves smoothness of
the orbit map.  Evaluation at the identity and continuity show $j$ is
nonzero.  Its kernel is compact-invariant, so compact irreducibility
makes $j$ injective.

\paragraph{DC8: the finite tensor calculation that eliminates lower types.}
\label{claim:DC8}
For a unitary infinitesimal representation the matrix units satisfy
$E_{ab}^*=E_{ba}$ on compact roots and $E_{ab}^*=-E_{ba}$ on
noncompact roots.  If $u$ is a compact highest vector of weight $\nu$
and $C_Gu=c_\lambda u$, expansion of the Casimir gives
\[
 (c_\lambda+\|\rho_c-\rho_n\|^2)\|u\|^2
 =\|\nu+\rho_c-\rho_n\|^2\|u\|^2
       -2\sum_{a<b,\,\epsilon_a\ne\epsilon_b}\|E_{ba}u\|^2.    \tag{DC.12}
\]
The proof pairs the $(a,b)$ and $(b,a)$ summands, uses the adjoint signs,
and replaces their commutator by $E_{aa}-E_{bb}$.  In particular no
compact highest vector of weight $\Lambda-(e_a-e_b)$ occurs if
$\lambda_a-\lambda_b>1+2\langle\rho_c,\rho_n\rangle$.
Counting ordered sign pairs gives
$2\langle\rho_c,\rho_n\rangle\le npq$, so the first inequality in
(DC.2) suffices.

There is a useful finite-dimensional version of this assertion.  On
$\mathfrak p_{\mathbb C}\otimes T$ put
\[
 A=\sum_{\epsilon_a=\epsilon_b}
                \operatorname{ad}(E_{ab})\otimes dT(E_{ba}).
\]
The operator is self-adjoint.  The elementary standard/dual tensor rule
in each compact block shows that every summand has highest weight
$\Lambda+e_r-e_s$, $\epsilon_r\ne\epsilon_s$, when that weight is
dominant.  This rule follows by locating a largest nonzero coordinate
of a highest vector in $\mathbb C^d\otimes V$ and solving the upper
matrix-unit equations; the dual rule is identical with indices reversed.
Since the compact Casimir on $\mathfrak p_{\mathbb C}$ is $n$, its
eigenvalue on this summand is
\[
 \theta_{rs}=\tfrac12\{c_K(\Lambda+e_r-e_s)-c_K(\Lambda)-n\}
 =\lambda_r+(\rho_n)_r-\lambda_s-(\rho_n)_s+1-n/2.
\]
For $r<s$ it is at least $h-n+1>0$; for $r>s$ it is at most $-h-1<0$.
Let $W_-,W_+$ be the negative and positive spectral subspaces.  The
multiplication map $X\otimes v\mapsto E(X)v$ kills $W_-$ by (DC.12).
This finite set of first-order relations is the Schmid relation used
below; its vanishing has been proved, not assumed.

\paragraph{DC9: coercive PBW layers and the minimum-norm condition.}
\label{claim:DC9}
Let $F_j$ be the span of words of at most $j$ noncompact derivatives
applied to $T$.  PBW and compact stability show that these exhaust the
generated Lie module.  If a nonzero highest vector of weight $\nu$
occurs in $F_j/F_{j-1}$, $j\ge1$, then
\[
 c_K(\nu)\ge c_K(\Lambda)+2j(h-n+1).                           \tag{DC.13}
\]
This conclusion requires only the first-order relation $W_-\mapsto0$.
Indeed choose the minimal-norm lift in
$\mathfrak p_{\mathbb C}^{\otimes j}\otimes T$ orthogonal to the
preimage of $F_{j-1}$.  Moving a negative slot to the last position
introduces commutators of smaller word length, so every negative-slot
tensor maps into $F_{j-1}$.  Every slice of the chosen lift is therefore
in $W_+$.  Expanding the total compact Casimir gives
\[
 C_K^{\mathrm{total}}=C_K^{\mathrm{tensor}}+c_K(\Lambda)
                                      +2\sum_{i=1}^j A_i.
\]
The first term is nonnegative, and DC8 bounds each $A_i$ below by
$h-n+1$.  The minimal lift is compact-equivariant, proving (DC.13).

In the tensor rule $\nu-\Lambda$ is a sum of $j$ noncompact roots.
Each has absolute pairing at most $(n-2)/2$ with $\rho_c$.
Consequently
\[
 \|\nu+2\rho_c\|^2-\|\Lambda+2\rho_c\|^2
 \ge j(2h-3n+4)>0.                                           \tag{DC.14}
\]
The last inequality follows from $p,q\ge1$, $pq\ge n-1$ and
$h>npq+1$.  Every fixed compact highest weight can occur only in
layers
$j\le[c_K(\nu)-c_K(\Lambda)]/[2(h-n+1)]$.
The original compact type occurs only in $F_0$ and with multiplicity
one.  These are finite-dimensional inequalities; no general
Blattner formula or discrete-series classification has been used.

\paragraph{DC10: closing the Lie module does not introduce extra compact types.}
\label{claim:DC10}
In the right regular $L^2$ space, close the module generated by $j(T)$
and call the resulting Hilbert space $H_\lambda$.  It is invariant under
$G_\epsilon$.  To justify passage from Lie operators to the group, choose
$A\ge0$ bounding the finitely many weights of $T$ and the Casimir scalar.
Every highest weight in $F_j$ has norm at most $A+\sqrt2 j$; the energy
identity used in DC7 therefore implies, for $X\in\mathfrak p$ and
$v\in F_k$,
\[
 \|E(X)^m v\|\le \|X\|_{\rm HS}^{m}
       \prod_{r=0}^{m-1}(A+\sqrt2(k+r))\,\|v\|.               \tag{DC.15}
\]
The Taylor series converges whenever
$|t|\sqrt2\|X\|_{\rm HS}<1$.  The integral remainder estimate, using
unitarity at every base point of the orbit, identifies its sum with
the actual group orbit.  Subdividing a finite interval gives invariance
under $e^X$; compact invariance and polar decomposition give invariance
under $G$.  Invariance under inverses gives equality.

For a compact irreducible type $\nu$, its compact character projector
is a bounded operator.  Apply it to the dense algebraic module.  DC9
places its range in one fixed finite-dimensional PBW layer, which is
closed.  The same is therefore true of its range on $H_\lambda$.
This proves admissibility, smoothness of all compact-finite vectors,
and persistence of the unique copy $j(T)$.  In particular completing
the module has not silently enlarged its lowest-type multiplicity.

\paragraph{DC11: irreducibility and identification of the actual coefficient.}
\label{claim:DC11}
Let $R$ be the restriction of the right regular representation to
$H_\lambda$.  If $S$ is a closed invariant subspace, its orthogonal
projection commutes with $G$ and hence with $K$.  By DC10 it maps
$j(T)$ into itself, and Schur's lemma makes its restriction a scalar.
Since it is a projection, that scalar is zero or one.  The space $j(T)$
generates $H_\lambda$, so $S=0$ or $H_\lambda$.  Thus $R$ is irreducible.

The invariant inner product on $j(T)$ is a positive scalar multiple of
that of $T$.  If $w$ is unit highest, put $a=\|j(w)\|_2^2>0$.
The exact coefficient is
\[
           \langle j(v'),R(g)j(v)\rangle
                  =a\langle v',\Psi(g)v\rangle.               \tag{DC.16}
\]
Here is a direct identification, avoiding an assumed positive-definite
kernel.  Project $R(g)j(v)$ onto $j(T)$, say to $j(u)$, and let the
orthogonal remainder be $f$.  Its compact matrix-coefficient averages
give compact maps $T\to H_\lambda$ orthogonal to $j(T)$; by uniqueness
they are zero.  The functions being averaged have continuous
representatives.  Their zero $L^2$ classes are therefore zero pointwise,
since Haar measure is positive on nonempty open sets.  Evaluate at $1$,
using $\Psi(kg)=T(k)\Psi(g)$ and $\Psi(k)=T(k)$.  The relevant compact
Schur integral is
$\|v_0\|^2\|\Psi(g)v-u\|^2/d_T$, so $u=\Psi(g)v$.
The inner product with $j(v')$ proves (DC.16).

\paragraph{DC12: all coefficients, Schur orthogonality and the source degree.}
\label{claim:DC12}
Equation (DC.16) and DC6 give a nonzero square-integrable coefficient
of $R$.  This implies every coefficient is square integrable and proves
the full four-vector orthogonality identity.  The functional-analytic
step can be seen explicitly.  For fixed $v$ define the closed operator
\[
 T_vu(g)=\langle v,R(g)u\rangle
\]
on the vectors for which this function lies in $L^2(G)$.  For an initial
nonzero square-integrable coefficient its domain contains the dense
span of translates of a nonzero vector.  Closedness follows by an
$L^2$ subsequence and pointwise continuity of the coefficient in $u$.
The positive self-adjoint operator $T_v^*T_v$ commutes with $R$.
All its spectral projections do too, so irreducibility makes it a
finite nonzero scalar times the identity.  Hence $T_v$ is bounded
and defined on the entire Hilbert space.  Inversion and complex
conjugation exchange the two coefficient slots; repeating the same
argument and polarizing in both slots gives
\[
 \int_G\langle w,R(g)v\rangle
       \overline{\langle w',R(g)v'\rangle}\,dg
       ={\langle v',v\rangle\langle w,w'\rangle\over d_R},
       \qquad d_R>0.                                         \tag{DC.17}
\]
Products are integrable by Cauchy--Schwarz.  This is precisely
\texttt{IsFdeg}.  Setting all four vectors equal to a nonzero $w$
shows that the source quotient \texttt{fdegOf} equals $d_R$ and is
independent of $w$.  Replacing $dg$ by $c\,dg$, $c>0$, replaces
$d_R$ by $d_R/c$.  DC9 gives the minimum Vogan norm required in the
literal \texttt{IsDS}; completeness and nontriviality are already
established.  We have constructed that entire predicate under (DC.2).

\paragraph{DC13: rigidity for the source's lowest-type definition.}
\label{claim:DC13}
An arbitrary representation satisfying the source's \texttt{IsDS}
for the same $(\epsilon,\lambda)$ is unitarily equivalent to the
constructed model under (DC.2).  This assertion is needed because
\texttt{IsDS} specifies a minimum compact type rather than a full
infinitesimal character.

First smooth a nonzero copy of the prescribed compact type by a
compact-conjugation-invariant approximate identity.  On that finite
copy the convolution is arbitrarily close to the identity, hence
injective, compact-equivariant and smooth.  It supplies a smooth
copy $V$ of $T$.  A first-layer type $\Lambda-e_a+e_b$, $a<b$,
would have
\[
 \|\Lambda-e_a+e_b+2\rho_c\|^2
       -\|\Lambda+2\rho_c\|^2
       =2-2(\lambda_a-\lambda_b+b-a)<0,
\]
contrary to the literal minimum-norm condition.  Thus the same
negative tensor space $W_-$ maps to zero.  The Schmid version of DC9
applies even before a Casimir scalar is known.

To recover that scalar without a classification theorem, take the
diagonal copy of $T$ in the direct sum of $V$ and the constructed
model's initial type.  Both projections kill $W_-$, so DC9 implies
that the generated diagonal module contains just one highest line
of weight $\Lambda$.  The central quadratic Casimir maps that line
to itself.  Projection to the constructed model determines its
scalar to be $c_\lambda$; projection to $V$ gives the same scalar
there.  Centrality extends it to the whole generated module.
DC10 closes this diagonal module to an invariant Hilbert subspace,
and DC11 makes it irreducible.  Its two coordinate projections are
nonzero intertwiners with dense range onto the two irreducible
representations.  The positive operators obtained by composing a
projection with its adjoint are scalar by Schur's lemma, hence each
projection is a positive multiple of a unitary equivalence.
Their ratio is the desired unitary intertwiner.  Consequently all
the central eigenvalues in (DC.8), as well as the formal degree for
any fixed Haar measure, hold for every literal \texttt{IsDS} at
these large-gap parameters.  This proof uses neither a general
lowest-type classification nor a universal existence theorem.

\paragraph{DC14: the later row weights satisfy the gap in the literal order.}
\label{claim:DC14}
The actual row rank is $n=R.r=k+m-1$.  Its ordered parameter consists
of $v_0,\ldots,v_{k-2}$, followed by $-b_0,\ldots,-b_{m-1}$;
the head has sign $+$ and the tail has alternating signs starting
with $+$.  Interlacing gives $v_{k-2}>u_{k-1}>-b_0$.
Therefore every opposite-sign head/tail difference is larger than
$b_1-b_0$, and every opposite-sign tail difference is at least the
least successive $b$-gap.  Head/head pairs are compact.
The source parity assumptions make $k,m$ even, $n$ odd and every
entry of this row parameter integral, so (DC.1) holds.  Strict
decrease follows from the same interlacing and increasing $b$'s.

The context has $b_j=b_0+jh_0$.  Its actual eventual quantifier
supplies, for arbitrary fixed $A\ge0,N,M$, simultaneously
\[
                   M\le m,\qquad A n^N\le h_0.                \tag{DC.18}
\]
For clarity, the source witnesses may be taken as $\epsilon_0=1$,
zero transported alternating data, $c_{D,0}=0$, $c_{h,0}=6N+1$ and
$x_0=\max(\max(4,A),M^4)$.  Its rank estimates give $n\le x^6$
and $m\ge x^{1/4}$, while $h_0\ge x^{c_h}$.  Thus these witnesses
prove (DC.18) in the original nested order, uniformly for every
later coefficient field and every same-predata object.
Take $A=16,N=19,M=1801$.  Then $n\ge1801$ and
$h_0\ge16n^{19}>H_\epsilon$, since $p,q,q_0\le n$.
This is the literal \texttt{GapHyp} route, and it preserves $B$, its
host, $R$, every $\mathrm{hc}_\tau$ and the supplied stage objects.

\paragraph{DC15: actual row existence, including the endpoint clauses.}
\label{claim:DC15}
For every $D$ satisfying \texttt{GapHyp}, and every real place $\tau$,
DC1--DC14 give a complete nonzero Hilbert representation
\[
             \operatorname{IsDS}(\epsilon_{SH},
                                  \lambda_{SH}(R,\mathrm{hc}_\tau),\sigma_\tau).
                                                                  \tag{DC.19}
\]
The choices are over the same finite set of real places; they do not
alter $D$.  The construction also gives finite compact isotypic
spaces, smooth compact-finite vectors, a unique lowest compact type,
the exact central parameter, and uniqueness among the source's
\texttt{IsDS} representations at these parameters.  For compact
signatures the representation of DC1 itself has these properties;
ordinary compact Schur orthogonality gives its degree.  In ranks
zero and one the empty products and the one-dimensional characters
give degree one for probability Haar.  Thus no rank-zero or
compact-signature exceptional branch is hidden in (DC.19).

\paragraph{DC16: the precise compact-dual quotient density.}
\label{claim:DC16}
Write $W_d=\prod_{j=0}^{d-1}j!$, with $W_0=1$.  GV5 proved that
compact probability Haar in ordinary skew-Hermitian coordinates
has density $1/v_d$, where $v_d=2^d\pi^{d(d+1)/2}/W_d$.
The source's compact-dual convention carries this same tangent
density to $U(p,q)$.  In block coordinates the compact factor has
probability density $1/(v_pv_q)$, so quotient measure on
$G/K$ has origin density
\[
             C_{pq}={v_pv_q\over v_n}
                    =\pi^{-pq}{W_n\over W_pW_q}.              \tag{DC.20}
\]
Identify $G/K$ with
$\mathcal D_{pq}=\{Z\in M_{p,q}(\mathbb C):1-Z^*Z>0\}$.
The derivative at the origin sends the off-diagonal Lie block
to $Z$ with Jacobian one.  The derivative of
$Z\mapsto(aZ+b)(cZ+d)^{-1}$ and the identity
\[
 1-(gZ)^*(gZ)=(cZ+d)^{-*}(1-Z^*Z)(cZ+d)^{-1}
\]
give the invariant quotient density
$C_{pq}\det(1-Z^*Z)^{-n}dZ$.
In particular (DC.20) concerns the actual noncompact quotient
measure induced from the compact-dual tangent convention, not
just the total volume of the compact group.

\paragraph{DC17: a scalar holomorphic model pins the normalization.}
\label{claim:DC17}
For an integer $s>n-1$, use holomorphic functions on $\mathcal D_{pq}$
with norm proportional to
\[
          \int_{\mathcal D_{pq}}|f(Z)|^2
                            \det(1-Z^*Z)^{s-n}\,dZ.           \tag{DC.21}
\]
The action
$f(Z)\mapsto\det(c_{g^{-1}}Z+d_{g^{-1}})^{-s}f(g^{-1}Z)$
is unitary by the preceding determinant and Jacobian identities.
Polynomials are dense: dilations $f(rZ)$ converge in the norm,
and their Taylor series converge on the smaller closed domains.
The central circle of $K$ separates homogeneous degrees; the
constant type occurs only in degree zero.  Translates of the
constant span densely.  Indeed averaging on that circle computes
the value at zero, and after a group translation computes the
value at any point; a function orthogonal to all such translates
vanishes everywhere.  The unique constant compact type and this
cyclicity prove irreducibility by the projection argument of DC11.

The beta integral in (DC.21) is elementary.  Remove the last row
$z$ of $Z$ and put $A=1-Z'^*Z'$.  The substitution $z=uA^{1/2}$
has real Jacobian $\det A$, and
$\det(1-Z^*Z)=\det A(1-\|u\|^2)$.  Radial integration in
$\mathbb C^q$ yields
\[
 I_{p,q}(a):=\int_{\mathcal D_{pq}}\det(1-Z^*Z)^a\,dZ
 =\pi^{pq}\prod_{r=0}^{p-1}{\Gamma(a+r+1)\over\Gamma(a+q+r+1)},
 \quad a>-1.                                                  \tag{DC.22}
\]
This follows by $p$ iterations of
$I_{p,q}(a)=\pi^q\Gamma(a+1)I_{p-1,q}(a+1)/\Gamma(a+q+1)$.
For a unit constant vector the squared coefficient is
$\det(1-Z^*Z)^s$.  Equations (DC.20)--(DC.22) and DC12 therefore give
\[
 d_s={1\over C_{pq}I_{p,q}(s-n)}
       ={W_pW_q\over W_n}
          \prod_{i=0}^{p-1}\prod_{j=0}^{q-1}(s-n+1+i+j).       \tag{DC.23}
\]
The central parameter can be checked on the constant vector by
triangular matrix units.  Order the negative block before the
positive block.  The two lists are
\[
 s+(q-p-1)/2-j\quad(0\le j<q),\qquad
 (n-1)/2-i\quad(0\le i<p).
\]
They are strictly ordered for $s>n-1$, have the required parity,
and their within-block Vandermondes are $W_q,W_p$.
Their cross differences have product equal to the double product
in (DC.23).  Hence this explicit model has
$d_s=|\prod_{a<b}(\lambda_a-\lambda_b)|/W_n$ for exactly the
compact-dual Haar measure.  No formal-degree theorem is used to
compute this calibration example.

\paragraph{DC18: the named degree theorem and its fully determined constant.}
\label{claim:DC18}
\emph{Explicit unreconstructed classical input HC-Deg.}
For a connected semisimple real group with finite center and a
compact Cartan, the formal degrees of its discrete series are a
positive group-and-Haar constant times the absolute product of
positive-root pairings of their regular infinitesimal parameters;
the constant is independent of the parameter and of the packet.
We use Harish--Chandra, \emph{Discrete series for semisimple Lie
groups. II: Explicit determination of the characters}, Acta
Mathematica 116 (1966), 1--111, \S40 Corollary 2 (p.~94) and
\S41 Theorem 16 (p.~96), DOI \texttt{10.1007/BF02392813}.
The publisher-version copy is available at
\url{https://repository.ias.ac.in/30833/1/336.pdf}.
The statement concerns the regular infinitesimal parameter; a
quadratic Casimir or a positive real number is not a substitute.

The reduction to $U(p,q)$ adds no unknown normalization depending
on the representation.  The finite covering
$U(1)\times SU(p,q)\to U(p,q)$ pulls an irreducible representation
back to a central character tensor an irreducible representation
of $SU(p,q)$.  To see irreducibility of the latter, a subspace
stable under $SU(p,q)$ is stable under the scalar central action
and hence under $U(p,q)$.  Probability measure on $U(1)$ contributes
one to the squared coefficient integral.  Finite-cover Haar
conversion is a fixed scalar independent of the character.
Root pairings depend only on differences of the $\lambda_i$.
HC-Deg therefore gives
$d_\lambda=A_{pq}|\prod_{a<b}(\lambda_a-\lambda_b)|$ for one
positive constant $A_{pq}$.  DC17 computes one such representation
and forces $A_{pq}=1/W_n$.  Compact groups give the same result
directly by the Weyl dimension formula and probability Schur
orthogonality.  Thus, \emph{conditionally on HC-Deg},
\[
               d_\lambda=
            \prod_{a<b}{|\lambda_a-\lambda_b|\over b-a}
                    =\operatorname{fdegHC}(\lambda).          \tag{DC.24}
\]
The equality $\prod_{a<b}(b-a)=W_n$ follows by collecting factors
with fixed $b$.  Every use of (DC.24) below retains HC-Deg.

\paragraph{DC19: the exact degree for every literal large-gap representation.}
\label{claim:DC19}
Assume HC-Deg, (DC.1) and (DC.2).  DC4 identifies the actual
infinitesimal parameter of the constructed model, so DC18 applies
to it.  DC13 transfers the result to \emph{every} representation
satisfying the literal \texttt{IsDS} at this parameter.
If $\mathrm{CdScaleIs}(cd,sc)$ holds, its scaled source measure
is the same compact-dual Haar measure under every allowed sign
sorting, by GV17 and uniqueness of Haar measure.  Consequently,
for each nonzero vector $w$,
\[
 \operatorname{fdegOf}(sc(n,\epsilon)\mu_G,\sigma,w)
             =\operatorname{fdegHC}(\lambda).                 \tag{DC.25}
\]
Conversely any $d$ satisfying \texttt{IsFdeg} for that same
representation and measure equals this number: insert $w$ in
all four slots of (DC.17) and divide by $\|w\|^4>0$.
Thus both the existential and the universal degree clauses
needed at the actual row parameters are supplied.  This is a
restricted supplier with all the source's vector and representation
quantifiers; it is not the all-parameter \texttt{FormalDegreeFormula\(_2\)}.

\paragraph{DC20: all four fields of the actual real-stage record.}
\label{claim:DC20}
Assume HC-Deg.  Let the \emph{same} $D$ satisfy
\texttt{GapHyp}, $\mathrm{VolTieG}(D)$,
\[
 D.s.cdScale=\mathrm{trueScale}(D.cd),\qquad
 D.dsig(\tau)=\mathrm{trueDsig}(D.R,D.\mathrm{hc}_\tau)
                    \quad\hbox{for every }\tau.               \tag{DC.26}
\]
Then $\mathrm{RealTied}(D)$ holds.  The field \texttt{cdScale}
is GV17.  The field \texttt{dsig} is DC15 together with
(DC.24), under the literal measure in (DC.26).  The field
\texttt{dsig\_\allowbreak all} is DC19, with its universal quantifier over
all Hilbert representations and all their asserted degrees.
The field \texttt{vol} is the supplied \texttt{VolTieG}.
This establishes the precise large-gap replacement for
\texttt{LevelThreeToFour.realOf}.  It neither assumes nor proves
the universal \texttt{RealOf} record, whose input $D$ has no
gap restriction.

When $D.vol$ is the actual \texttt{baseVol}/\texttt{trueVol}
on the same host, spaces and $cd$, GV18--GV19 supply the last
field under their explicitly retained Tamagawa-number theorem.
In this use DC20 is conditional on both HC-Deg and that theorem;
the two premises do different work.  An arbitrary stored
positive volume does not satisfy the field merely by positivity.

\paragraph{DC21: the same-stage Core join and the fields still needed.}
\label{claim:DC21}
Fix the literal $K,X,f,T$ of \texttt{coreOfLay}, a layout
$Y:\mathrm{Lay}(X.\mathrm{toD7})$, and its exact
\texttt{Over}, \texttt{RecOK}, \texttt{LedOK} witnesses.
Assume the following remaining inputs on these very objects:
the given $\xi$ and twist set $Wt$, the same-map \texttt{MuIsW},
forced inclusion of $g_0,g_1,d$ and the specified small primes,
\texttt{GatNum}, the untwisted row vanishing, \texttt{MtStableAt},
the primitive type-character power identities, and
$NChoice.Q=T.Q$, $NChoice.k'=T.k$.  Suppose also that $X.\mathrm{toD7}$
satisfies the gap and (DC.26), with actual volume tied as in DC20.
Then the literal \texttt{Core} on $X,f,T$ exists, conditionally
on HC-Deg and on any named theorem used to supply that volume.

Here is the field audit, rather than an implicit reassembly.
The data fields \texttt{xi}, \texttt{Wt} and \texttt{mu} are the
ones supplied; \texttt{isOver}, \texttt{forced}, \texttt{gatNum},
\texttt{mtStable} and \texttt{chiT} are the corresponding assumptions.
The object field is literally $O=\texttt{layObj}(Y)$, with the same
layout $Y$ throughout these assignments.
\texttt{fval} is \texttt{LedOK.1} with \texttt{MuIsW.off};
\texttt{col} is the same layout's prime, embedding, Minkowski and
type-place fields together with \texttt{RecOK};
\texttt{gy} is its ramification/filling fields with \texttt{LedOK.2}.
The nine \texttt{row} fields are the literal row layout, with only
\texttt{untw} supplied separately.  The theta/scalar fields
\texttt{th\_\allowbreak io}, \texttt{gauge\_\allowbreak def}, \texttt{del}, \texttt{etaF}
are read from that same layout and \texttt{RecOK}.
\texttt{real} is DC20; \texttt{dsig\_\allowbreak eq} is (DC.26).
\texttt{nch\_\allowbreak tie} is the layout's actual local-character tie and
\texttt{nch\_\allowbreak Q} is the asserted conductor-support/exponent pair.
This accounts for all twenty-one fields.  In particular it does
not turn a bare empty-twist stage into a geometric core: the
\texttt{forced} field already requires the actual distinguished
prime to occur.  Changing a stored spectral or volume component
requires a new layout or an explicit dependent transport theorem;
this argument never reuses a layout on a different $X$ by notation.

\paragraph{DC22: a restricted level-three-to-four replacement and its frontier.}
\label{claim:DC22}
Assume HC-Deg and retain the literal source hypotheses
\texttt{LevelThreeCondition}, \texttt{SatNew}, \texttt{PackEx},
the actual volume-to-mass implication, and the real idempotent
tie.  For the source's \texttt{SExt.none} branch, the
\texttt{step\_\allowbreak of}/\texttt{lv4\_\allowbreak of} argument works with DC20 in
place of the universal \texttt{RealOf} premise.  The replacement
is justified at \emph{both} uses of that premise.  First DC14
supplies the gap in the original eventual order.  The explicit
\texttt{exists\_\allowbreak sStage\_\allowbreak true} sets the spectral scalar,
\texttt{trueScale} and \texttt{trueDsig}, and supplies the ordinary
local tie on the same $B,C,W$; one uses the stipulated
\texttt{LevelThreeCondition} to move the level-three data to
this new stage.  DC20 supplies its real record.  \texttt{SatNew}
then returns $Z_s$ and its universal \texttt{SameL} comparison
property.  The source constructs the replacement stage using
the actual old divisibility lattices and heights, transports
columns through \texttt{ColRe.exists\_\allowbreak colsTiedSym}, and applies
the returned \texttt{SameL}/carrier isometries.  The row rank and
every $\mathrm{hc}_\tau$ are unchanged, so DC14's gap and
(DC.26) still hold at the second real-record use; the new volume
tie is furnished by the new stage's actual \texttt{GIs2} witness.
DC20 applies again.  This proves the restricted transition with
all remaining hypotheses displayed.

This conclusion removes the two universal discrete-series wrappers
from this \emph{particular} later-stage route.  It does not prove
\texttt{SatNew}: its spectral construction still requires the
actual cusp finite-dimensionality, rapid-decay, infinitesimal
character and uniqueness inputs in their exact scopes, or new
proofs of the restricted variants used there.  It does not prove
\texttt{PackEx}, a new compatible nonempty motive/twist layout,
or unconditional \texttt{ExTrue}.  The all-parameter existence
and degree wrappers, and the source's optional internal
translation/Laplace proof of degree, are retained as distinct
targets.  No claim in this section is a proof of RH or GRH.


\section{Relative progression errors and simultaneous prime pairs}
\label{sec:simultaneous-prime-pairs}

This section continues CC1--CC20.  It supplies the analytic prime-selection
endpoints \texttt{PrimesRel}, \texttt{twistK}, \texttt{meanSqK},
\texttt{pair\_\allowbreak exists}, and \texttt{PairEx}, and the real-zero exclusion on
an already specified compatible host.  Constants denoted by $C,c,A$ are
positive and absolute unless their dependence is explicitly indicated.
Zeros are counted with analytic multiplicity.  No assertion about zeros
on the critical line is used here.

Three classical inputs are retained, with their mathematical content
specified here rather than hidden inside a source theorem name.
\begin{enumerate}
\item[(CN)] The analytic class-number formula for quadratic fields, with
the usual ideal class number $h$, number of roots of unity $w$, and
regulator $R$: the residue is
$2^{r_1}(2\pi)^{r_2}hR/(w\sqrt{|D|})$.
In particular a real primitive nonprincipal Dirichlet character of
conductor $f$ has $L(1,\chi)\ge c/\sqrt f$.
The convention for a real quadratic field is
$L(1,\chi_D)=2h\log\epsilon_D/\sqrt D$, with $\epsilon_D>1$
generating its units modulo torsion.
\item[(BD)] Bombieri's log-free density theorem, including its exceptional
factor.  There are absolute $A_B,C_B,c_B>0$ such that the total number
$N_R(\sigma)$ of zeros of primitive characters of conductor at most $R$,
with $\sigma\le\Re\rho<1$ and $|\Im\rho|\le R$, is at most
$C_BR^{A_B(1-\sigma)}$, for $R\ge2$ and $1/2\le\sigma\le1$.
If a real primitive character in this family has a real zero
$\beta=1-\delta>1-c_B/\log R$, omit that one simple zero; the resulting
count is at most
$C_B\delta\log R\,R^{A_B(1-\sigma)}$.
\item[(HD)] A classical hybrid density estimate: for each fixed $h>0$,
\[
 \sum_{f\le Q}\ \sum_{\chi\bmod f}^{*}N(\sigma,T,\chi)
       \le C_h(Q^2T)^{2(1-\sigma)+h},
 \quad Q\ge2, T\ge1, 11/14\le\sigma\le1.
\]
Here the asterisk restricts to primitive characters, and both signs of
the imaginary part are counted.  This is the form of Heath-Brown's
1979 Theorem 1 that will be used.
\end{enumerate}
These are named classical premises, not proofs reconstructed in this
section.  Primary references are Bombieri, \emph{Le grand crible dans la
th\'eorie analytique des nombres}, Ast\'erisque 18, \S6, Theorem 14
(the 1987 second edition preserves that statement), and Heath-Brown,
\emph{The density of zeros of Dirichlet's $L$-functions},
Canad. J. Math. 31 (1979), 231--240, Theorem 1,
\url{https://doi.org/10.4153/CJM-1979-024-0}.
For the normalization in (CN), see Sutherland's MIT 18.785 notes,
Lecture 19 (2021), on the analytic class-number formula.
The source's particular large numerical density constants and its
fixed-$1/8$ exceptional-pair wrapper are not asserted by these substitutions.

Write $\psi_\chi(u)=\sum_{n\le u}\Lambda(n)\chi(n)$ and
$\operatorname{Li}(u)=\int_2^u dt/\log t$.  Primitive principal characters
occur only at conductor one.  The tail estimate is first proved for
nonprincipal primitive characters and then extended to every nonprincipal
character at its original level by an explicit deleted-Euler-factor bound.

\paragraph{QP1: an effective real-zero gap with the right normalization.}
Suppose $f\le m$, $m\ge3$, and $\chi$ is a real primitive nonprincipal
character of conductor $f$.  The two quadratic cases of (CN) give the
stated uniform lower bound: for $D<0$, $w\le6$ and $h\ge1$; for $D>0$,
$h\ge1$ and $\epsilon_D\ge(1+\sqrt5)/2$.  The latter inequality follows
by taking the integer trace of a quadratic unit of norm $1$ or $-1$.

For $1-1/\log m\le s\le1$, periodic cancellation gives
$|\sum_{n\le u}\chi(n)|\le f$.  Split $L(s,\chi)$ at $m$ and differentiate
the absolutely convergent partial-summation integral for its tail.
The finite part is bounded by
$e\sum_{n\le m}(\log n)/n\le C\log^2m$.
The tail and its derivative are bounded by
$Cf m^{-s}(1+\log m)\le C(1+\log m)$, after decreasing the fixed
neighborhood of one if necessary.  Thus $|L'(s,\chi)|\le C\log^2m$.
If $L(\beta,\chi)=0$ in this neighborhood, the real mean-value theorem
implies
\[
 1-\beta\ge \frac{c}{\sqrt f\log^2m}
       \ge\frac{c'}{m^2\log m}\ge\frac{c'}{m^3}.       \tag{QP.1}
\]
The second inequality is uniform for $m\ge3$, since
$\log m\le m^{3/2}$.  The constants do not depend on the character.
This proves the effective lower-gap clause needed by
\texttt{PrimesRel}; it does not assert Siegel's ineffective lower bound.

\paragraph{QP2: one exceptional primitive character across bounded conductors.}
There is $c_P>0$ such that among primitive conductors at most $R$, at most
one real character has a zero $\beta>1-c_P/\log R$; that zero is simple.
Here is the positivity argument, including the across-conductor point.
For two distinct real primitive characters $\chi_1,\chi_2$, inflate them
to a common modulus.  The coefficients of
\[
 -\frac d{ds}\log\{\zeta(s)L(s,\chi_1)L(s,\chi_2)
                                      L(s,\chi_1\chi_2)\}
\]
are $\Lambda(n)(1+\chi_1(n))(1+\chi_2(n))\ge0$ on $s>1$.
The product character is nonprincipal: otherwise the primitive characters
would agree.  Removing Euler factors at the common modulus changes the
local logarithmic-derivative estimates by $O(\log R)$.
The real part of every critical-zero term at a real $s>1$ is positive.
The local expansion of CC3, or its Hadamard-product limit, therefore gives
\[
 \frac1{s-\beta_1}+\frac1{s-\beta_2}
                  \le\frac1{s-1}+C\log R.             \tag{QP.2}
\]
Set $s=1+a/\log R$.  First choose $a>0$ with $Ca<1/4$, and then
$c_P<a/10$.  The inequality
$2/(a+c_P)\le1/a+C$ is impossible.  Applying the same argument to
$\zeta(s)L(s,\chi)$, counting multiplicity, excludes two nearby real
zeros of one character and excludes a multiple nearby zero.
The standard trigonometric positivity argument of CC5 excludes the
nonreal alternatives in a sufficiently small region
$\Re s>1-c/\log(f(|\Im s|+3))$.
Consequently at a fixed modulus $m$ one can select a unique exceptional
pair, or select no pair, using a sufficiently small fixed $c_1\le1/8$.
Its primitive conductor divides $m$, and induction back to modulus $m$
is unique.  QP1 applies to this pair.

\paragraph{QP3: a weighted zero sum with the exceptional factor retained.}
Fix $m\ge3$, $T\ge2$, put $H=\log(m(T+3))$, and suppose
$V\ge C_0H$ for a sufficiently large absolute $C_0$.
For characters modulo $m$, omit the selected exceptional zero if present.
Their remaining zeros satisfy $\Re\rho\le1-c/H$ by QP2 and CC5.
Let $D=1$ when no exceptional pair is selected, and let
$D=\min(1,\delta H)$ when the selected zero is $1-\delta$.
Then (BD) gives, uniformly for $1/2\le\sigma\le1$,
\[
 N'(\sigma;m,T)\le C D\exp\{A_BH(1-\sigma)\}.          \tag{QP.3}
\]
Indeed use the primitive family with cutoff $R=m(T+3)$.  If
$\delta H<c_B$, the selected zero is exceptional in that larger family,
so the second part of (BD) applies.  If $\delta H\ge c_B$, its first part
suffices after multiplying by $\max(1,c_B^{-1})$.  Distinct characters
modulo $m$ have distinct primitive inducers.  Extra Euler factors have
zeros only on $\Re s=0$, outside the region being counted.

For $b=\Re\rho\ge1/2$, use
\[
 e^{-(1-b)V}=e^{-V/2}+V\int_{1/2}^{b}e^{-(1-\sigma)V}\,d\sigma.
\]
Summation, followed by (QP.3), bounds the sum by an initial term
$CD e^{-(V-A_BH)/2}$ and
$CDV\int_{c/H}^{1/2}e^{-(V-A_BH)t}\,dt$.
Take $C_0\ge2A_B+4$.  Both terms are at most
$CD e^{-c'V/H}$.  Since $|\rho|\ge1/2$, this proves
\[
 \sum_{\chi\bmod m}\ \sum_{\substack{\Re\rho\ge1/2, |\Im\rho|\le T\cr
                                    \rho\ {\rm not\ exceptional}}}
       \frac{e^{(Re\rho-1)V}}{|\rho|}
             \le C D e^{-c'V/H}.                       \tag{QP.4}
\]
The factor $D$ is essential when the main term is small.  Replacing it
by $1$ would not prove a relative progression estimate.

\paragraph{QP4: the right-half explicit formula and relative Chebyshev error.}
Let $u=e^V$, $m\ge3$, and choose
$T=m^8e^{\sqrt V}$, with $u\ge m^{C_2}$ and $C_2$ sufficiently large.
Then $T\le u$ and $H\le10\log m+\sqrt V$, while $V\ge C_0H$.
Character orthogonality and the explicit formula yield
\[
 \varphi(m)\psi(u;m,a)
 =u-\theta\frac{u^\beta}{\beta}
   -\sum_{\chi\bmod m}\overline{\chi(a)}
      \sum_{\rho\ne\beta}\frac{u^\rho}{\rho}+E,         \tag{QP.5}
\]
where the displayed zero sum can be restricted to
$\Re\rho\ge1/2$, $|\Im\rho|\le T$, and
\[
 |E|\le C m\left\{u\log^2(muT)/T+
                              u^{1/2}\log^2(muT)\right\}.\tag{QP.6}
\]
For this right-half form there is no small-zero denominator left in the
error.  In the notation of CC8, the zero-removed local logarithmic
derivative gives the signed identity
\[
 b_0=\sum_{\rho\in Z_\chi(1)}\frac{m_\rho}{\rho}
                         +O(\log(N+2)).
\]
Indeed, each local term $-m_\rho/(s-\rho)$ has constant term
$m_\rho/\rho$ at zero; the omitted zeros at height at least one and
the regular gamma terms cost $O(\log(N+2))$ by CC3.
The right-half summands in this constant have $|\rho|\ge1/2$ and
are harmless.  For the left-half low zeros, combining the constant
with the explicit-formula sum gives exactly
$-m_\rho(u^\rho-1)/\rho$.  If $|\rho|$ is small and $\Re\rho\le1/2$, its
absolute value is at most $m_\rho u^{1/2}\log u$ using
$(u^\rho-1)/\rho=\int_1^u t^{\rho-1}dt$.
The CC3 unit-height counts then give the second term of (QP.6).
Zeros away from zero are bounded directly.  The removal of Euler factors
at primes dividing $m$ costs at most $C\log m\log u$, also absorbed there.
The CC8 floor-endpoint convention remains in force.

Here $\theta=0$ if no pair was selected; otherwise
$\theta=\chi_1(a)\in\{1,-1\}$ and $7/8\le\beta<1$.
For $V\ge8$, put $w=1$ in the former case and
$w=\min(1,(1-\beta)V)$ in the latter.  Elementary exponentiation gives
\[
 A(u):=u-\theta u^\beta/\beta\ge uw/4.                \tag{QP.7}
\]
For example $e^{-\delta V}/(1-\delta)\le e^{-\delta V/2}$
when $V\ge8$ and $\delta\le1/8$; this proves (QP.7) for the only delicate
case $\theta=1$.
Because $H\le V$, $D\le w$, so (QP.4) bounds the zero term relative to
$A(u)$ by $Ce^{-cV/H}$.
For the error (QP.6), (QP.1) gives $w\ge\min(1,c m^{-3}V)$.
Consequently its first relative part is at most
$C m^{-4}(V+H)^2e^{-\sqrt V}$, and its second is at most
$C m^4(V+H)^2e^{-V/2}$.  Increasing $C_2$ and decreasing $c$ absorbs
the polynomial factors uniformly.  Finally
$V/(10\log m+\sqrt V)\ge\frac1{20}\min(V/\log m,\sqrt V)$.
We have proved
\[
 |\varphi(m)\psi(u;m,a)-A(u)|
 \le C\{e^{-cV/\log m}+e^{-c\sqrt V}\}A(u).           \tag{QP.8}
\]
All constants are chosen before $m,a,u$ and before the exceptional pair.

\paragraph{QP5: partial summation preserves the relative error.}
Set $B(u)=\operatorname{Li}(u)-\theta\operatorname{Li}(u^\beta)$.
The identity $B'(u)=(1-\theta u^{\beta-1})/\log u$ exactly matches
$A'(u)/\log u$; in particular there is no extra factor $1/\beta$ in
the prime-counting main term.  Prime powers of exponent at least two
contribute $O(\sqrt u\log^2u)$ to $\psi$.
Apply partial summation on $[\sqrt u,u]$, using (QP.8) there after
increasing the starting power of $m$.  On this interval the error factor
is at most $C(e^{-cV/(2\log m)}+e^{-c\sqrt{V/2}})$.
The positive weight gives an error at most that factor times
\[
 A(u)/\log u+\int_{\sqrt u}^{u}A(t)/(t\log^2t)\,dt
                                  \le 2A(u)/\log u.
\]
Here $A(t)$ is increasing on the interval.  The portion below $\sqrt u$
and the prime-power correction cost $O(m\sqrt u\log^2u)$.

For $\theta=1$, integration over $[u^\beta,u]$ shows
$B(u)\ge c\,u\min(1,\delta\log u)/\log u$;
also $0<A(u)\le u\min(1,\delta\log u)$.
For $\theta=0$ or $-1$, both comparison statements hold with the minimum
replaced by $1$ and with changed absolute constants.  Thus the preceding
error is relative to $B(u)$, and (QP.1) absorbs the initial segment.
We obtain, for $u\ge m^{C_3}$,
\[
 |\varphi(m)\pi(u;m,a)-B(u)|
 \le C\{e^{-c\log u/\log m}+e^{-c\sqrt{\log u}}\}B(u).
                                                               \tag{QP.9}
\]
Taking $c_1$ from QP2, the zero-gap constant from QP1, and changing $c$
to match $e^{-\sqrt{c_3\log u}}$, gives the literal existential constants
of \texttt{GlobalTheory.Lin.PrimesRel} and its \texttt{Emain} convention.
This obtains that endpoint under (CN) and (BD).  It does not identify the
source's fixed-$1/8$ \texttt{PNTAPRelW} record; the argument below uses
(QP.9) directly and needs no such identification.

\paragraph{QP6: many primes in every admitted progression.}
Call $(Y,U,m,a)$ admitted at exponent $\epsilon>0$ when
$m>0$, $8\mid m$, $(a,m)=1$, and
\[
 \tfrac14\log^2Y\log m\le\log U,\qquad
 Y^{\epsilon/6}\le\log U\le Y^{\epsilon/5}\log^2Y.     \tag{QP.10}
\]
Let $P=\{q:\ q\text{ prime},\ \sqrt U<q\le U, q\equiv a\pmod m\}$.
For all sufficiently large $Y$, uniformly in the admitted data,
\[
                         |P|\ge U^{9/10}.              \tag{QP.11}
\]
Indeed (QP.10) gives $\log m=o(\log U)$ uniformly, so (QP.9) applies
and its relative error tends to zero.  Even in the exceptional case,
(QP.1) and QP5 give
$B(U)/\varphi(m)\ge cU/(m^4\log U)$.
Since $\log m\le4\log U/\log^2Y$ and
$\log\log U=O_\epsilon(\log Y)$, this lower bound is eventually
at least $2U^{9/10}$.  Removing at most $\sqrt U$ small primes proves
(QP.11).  This uses the effective lower gap, not an ineffective
constant depending on a moving modulus.

\paragraph{QP7: cancellation outside one squarefree exceptional index.}
For any fixed $A,B,\epsilon>0$ and all sufficiently large $Y$, admitted
data have an index $s_*\ge0$ such that for every odd squarefree
$1<s\le Y^A$, $(s,m)=1$, $s\ne s_*$,
\[
 \left|\sum_{q\in P}\left(\frac qs\right)\right|
                                        \le Y^{-B}|P|.\tag{QP.12}
\]
Put $R=mY^A$.  On the reduced residue classes modulo $ms$ lying over
$a\pmod m$, the function $(b/s)$ is a nontrivial character of the
kernel of reduction to $m$.  Its sum is zero.  If an exceptional
character is present, its contribution to this weighted sum also
vanishes unless its restriction to that kernel is exactly $(\cdot/s)$.
Call such a character aligned.  By QP2 at most one aligned index can
have $1-\beta<c_P/\log R$: two such characters must have the same
primitive inducer, and restricting to the kernel at any prime dividing
one of the two squarefree indices but not the other contradicts equality.
The Chinese remainder theorem provides the required test unit.
Take this index for $s_*$, or $s_*=0$ if none exists.

For every other index an aligned exceptional term is at most
$r\operatorname{Li}(U)$, with
$r\le4\log U\,e^{-c_P\log U/\log R}$.
The relative errors from (QP.9), summed over the fiber, are at most
$k$ times the sum of its positive main terms, where
$k=C(e^{-c\log U/\log(ms)}+e^{-c\sqrt{\log U}})$.
If alignment occurs the unweighted exceptional character sums to zero
on the fiber, so its total main term is the ordinary one; if it does
not occur the weighted main term is zero.  In either case division by
the total count costs at most $(1-k)^{-1}$, not $\varphi(s)$.
Removing primes below $\sqrt U$ changes the weighted sum by at most
$\sqrt U$; primes dividing $s$ have Jacobi value zero.

Uniformly for (QP.10),
$\log U/(\log m+A\log Y)\gg_{A,\epsilon}\log^2Y$.
Thus $r+k=o(Y^{-B})$, while (QP.11) gives
$\sqrt U/|P|\le U^{-2/5}=o(Y^{-B})$.
After increasing the single threshold these bounds prove (QP.12),
the literal \texttt{twistK} conclusion with arbitrary fixed $A,B$.

\paragraph{QP8: the coefficient estimates have absolute constants.}
For $y\ge3$ put $c_y(n)=\Lambda(n)\min(1,y/n)$ and, for $N\ge1$,
\[
 b_p=\sum_{p^j\le N}c_y(p^j),\quad
 R_m=\sum_{\substack{p^j\le N\cr j\ {\rm even}\ {\rm or}\ p\mid m}}c_y(p^j).
\]
The elementary binomial-coefficient proof of Chebyshev's bound gives
$\sum_{p\le t}\log p\le t\log4$ and $\psi(t)\le Ct$.
For $p\le y$, splitting the geometric series at
$\lfloor\log y/\log p\rfloor$ gives $b_p\le3\log y$.
For $p>y$, it gives $b_p\le2y\log p/p\le2\log y$.
The same Chebyshev bound implies $\pi(y)\le Cy/\log y$ by separating
the primes below $\sqrt y$.  Partial summation of
$\sum_{p>y}\log^2p/p^2$ then proves
\[
 \sup_p b_p\le C\log y,\qquad \sum_p b_p^2\le Cy\log y.\tag{QP.13}
\]
The even powers in $R_m$ are bounded by
$\sum_{n\le N}\Lambda(n)\min(1,y/n^2)\le C\sqrt y$,
using $\psi(t)\le Ct$ and the decreasing kernel.
If every prime divisor of $m$ that is at most $N$ is at most $w$, the
other powers contribute at most $Cw\log y$.  Hence
\[
 R_m\le C(\sqrt y+w\log y),\qquad
 \sum_{n\le N}c_y(n)\le Cy\log(2N).                 \tag{QP.14}
\]
Finally the identity
$\log(\lfloor z\rfloor!)=\sum_{n\le z}\Lambda(n)\lfloor\lfloor z\rfloor/n\rfloor$,
Stirling's elementary integral bounds, and $\psi(z)\le Cz$ show
\[
                    \sum_{n\le z}\Lambda(n)/n\le\log z+C.\tag{QP.15}
\]
Only its leading coefficient $1$, not the numerical value of $C$, is
needed below.  These estimates are uniform in the cutoff $N$.

\paragraph{QP9: expand the mean square without assuming multiplicativity.}
Let $P$ consist of primes $q>N^2$ in one reduced class modulo $m$,
where $8\mid m$.  Let $|\kappa(n)|\le1$; no multiplicativity is imposed.
Put
\[
 S_N(y,q,\kappa)=\sum_{n\le N}c_y(n)\kappa(n)(n/q).
\]
Quadratic reciprocity on this fixed class groups its terms as
$S_N=b_0+\sum_{p\nmid m}b'_p(q/p)$, where the latter indices are odd
primes, $|b_0|\le R_m$, and $|b'_p|\le b_p$.
Even powers and powers of primes dividing $m$ enter $b_0$; the signs
from reciprocity and the prime $2$ are fixed by the residue class.
Suppose all the squarefree correlations in (QP.12), up to $N^2$, are
bounded by $\eta|P|$ except possibly $s_*$.
The diagonal of the square expansion is bounded by
$(R_m^2+\sum_pb_p^2)|P|$.
An exceptional correlation can occur only at $s_*=p$, contributing
at most $2R_m\sup b_p$, or at $s_*=pp'$ for distinct primes,
contributing at most $2(\sup b_p)^2$.
All remaining cross terms total at most
$\eta(\sum_{n\le N}c_y(n))^2|P|$, by the triangle inequality.
Therefore
\[
 \sum_{q\in P}|S_N(y,q,\kappa)|^2\le |P|\left\{
 \sum_pb_p^2+R_m^2+2B^2+2BR_m+
                        \eta\left(\sum_{n\le N}c_y(n)\right)^2\right\},
 \quad B=\sup_pb_p.                                    \tag{QP.16}
\]
This is the source's \texttt{meanSq\_\allowbreak abs} inequality, obtained by an
explicit prime-power partition.  It explains why one exceptional
squarefree index does not destroy the desired mean square.

\paragraph{QP10: the actual polynomial-window mean square.}
Fix $0<\epsilon\le1$, $K\ge1$, set $Z=(Y^K)^{400}$ and
$N=\lfloor Z\rfloor$.  In addition to (QP.10), suppose every prime
dividing $m$ is at most $Y^{\epsilon/5}$ or exceeds $Z$.
QP7 with $A=800K$, $B=K+1$ gives the correlations for (QP.16), with
$\eta=Y^{-(K+1)}$.  For sufficiently large $Y$,
$\sqrt U>N^2$, so its prime-size condition also holds.
For $Y\le y\le Y^K$, take $w=Y^{\epsilon/5}$ in QP8.
Uniformly on this interval, $w\log y\le\sqrt y$ eventually.
Thus $R_m^2\le Cy$, and (QP.13) bounds the diagonal by $Cy\log y$.
The two exceptional terms are at most $C(\log^2y+\sqrt y\log y)$.
The final term divided by $y\log y$ is at most
$CY^{-1}(400K\log Y)^2/\log Y$, which tends to zero.
We have proved that one absolute $C>0$ satisfies
\[
 \forall y\in[Y,Y^K]\ \forall\kappa,quad |\kappa(n)|\le1
 \Longrightarrow\quad
 \sum_{q\in P}|S_N(y,q,\kappa)|^2\le Cy\log y\,|P|.   \tag{QP.17}
\]
The threshold depends on $\epsilon,K$, but $C$ does not.
This is the literal stronger \texttt{meanSqK} endpoint, including its
quantification over all bounded coefficient sequences.

\paragraph{QP11: a maximal estimate for a continuum of arguments.}
For fixed $q$, the function $S_N(y,q,\kappa)$ is continuous and piecewise
linear, with right derivative
\[
 D_q(v)=\sum_{v<n\le N}\Lambda(n)\kappa(n)(n/q)/n.
\]
The exact identity
$vD_q(v)=S_N(v,q,\kappa\,1_{n>v})$
allows (QP.17) to be applied to this derivative, because the truncated
sequence is still bounded by one.  Hence
$\sum_q|D_q(v)|^2\le C(\log v)|P|/v$.
On a block $[a,b]\subseteq[Y,Y^K]$ with $b\le2a$, the fundamental
theorem of calculus and Cauchy--Schwarz give, for every $y_q\in[a,b]$,
\[
 |S_N(y_q,q,\kappa)|^2
 \le2|S_N(a,q,\kappa)|^2+2a\int_a^b|D_q(v)|^2dv.
\]
Summing and integrating the derivative bound proves
\[
 \sum_{q\in P}|S_N(y_q,q,\kappa)|^2
                                  \le4Ca\log(2a)|P|.\tag{QP.18}
\]
All sums are finite and $D_q$ is a step function; thus no measurability
or exchange-of-limit issue occurs.  In particular each prime may choose
its own witness $y_q$.  This is the step that upgrades pointwise control
to control throughout a real interval.

\paragraph{QP12: discard only a small fraction of primes.}
Let $\tau>0$.  If a prime is bad on $[a,b]\subseteq[a,2a]$, choose a
witness with $|S_N(y_q,q,\kappa)|>y_q^{1/2+\tau}/2$.
By (QP.18), the proportion of such primes is at most
$16C\log(2a)a^{-2\tau}\le(32C/\tau)a^{-\tau}$ for $a\ge2$.
Use blocks $[2^iY,\min(2^{i+1}Y,Y^K)]$.
Every bad prime belongs to at least one block, so the geometric series
gives
\[
 \frac{\#\{q\in P:\ \exists y\in[Y,Y^K],\
                    |S_N(y,q,\kappa)|>y^{1/2+\tau}/2\}}{|P|}
 \le \frac{32C}{\tau(1-2^{-\tau})}Y^{-\tau}.          \tag{QP.19}
\]
This tends to zero uniformly in $P,N,\kappa$ whenever (QP.17) holds.
For a fixed finite family of at most $R$ coefficient sequences, choose
the threshold so that the right side is at most $\eta/(2(R+1))$ and
take their union.  This is the precise \texttt{bad\_\allowbreak share} supplier.

\paragraph{QP13: extend the explicit formula to the needed cutoff heights.}
Suppose $\chi$ is primitive nonprincipal, its conductor is at most $Q$,
and it has no zero with $\Re s>399/400$, $|\Im s|\le Q$.
Reflection by CC1 gives $1/400\le\Re\rho\le399/400$ for its nontrivial
zeros below height $Q$ (apply conjugation as well as reflection).
Take $T=Q/2$, with $Q$ sufficiently large.
The contour proof in CC8 also works when $T>u$, provided one retains
the left-edge term $Cu^{-1/2}\log^2(N(T+2))$ instead of absorbing it.
To check this extension, the Perron kernel bounds the term nearest $u$
by $C\log(2u)$; the remaining terms sum to
$Cu\log^2(2u)/T$ irrespective of the order of $u$ and $T$.
Choose the good horizontal edges in $[T,T+1]\subseteq[0,Q]$ as in CC8.
Their cost is $Cu\log^2(Qu)/Q$, and the left edge is the term just stated.
The pole-at-zero constant is $O(\log^2Q)$ by the reflected lower bound
and CC3.  The reciprocal zero weights sum to $C\log^2Q$ by CC7,
now with the fixed lower bound $1/400$ at small height.
It follows that for every $u\ge3$,
\[
 |\psi_\chi(u)|\le C\{u^{399/400}\log^2Q
       +(u/Q)\log^2(Qu)+\log(Qu)+\log^2Q
                              +u^{-1/2}\log^2Q\}.     \tag{QP.20}
\]
This inequality concerns this zero-free box only; it asserts no
zero-free strip beyond height $Q$.

\paragraph{QP14: a two-range tail estimate supplies polynomial truncation.}
Put $\lambda=\log Q$, and let $V=\exp(B\lambda^2)$, where $B$ is a
sufficiently large fixed absolute constant.  The real-zero exclusion
required in CC10 holds with $c_0=1/400$: real zeros have imaginary part
zero and are therefore in the box at every height.
CC11--CC12 justify convergence and partial summation of
\[
 T_\chi(Z):=-L'/L(1,\chi)-\sum_{n\le Z}\Lambda(n)\chi(n)/n
       =-\psi_\chi(Z)/Z+\int_Z^\infty\psi_\chi(u)/u^2\,du.\tag{QP.21}
\]
For $3\le Z\le V$, integrate (QP.20) up to $V$.
The zero term contributes at most $C Z^{-1/400}\lambda^2$;
the $1/Q$ term contributes at most $C\lambda^6/Q$;
the remaining terms contribute at most
$CZ^{-1}\log^2(QZ)$, including the endpoint in (QP.21).
Above $V$, CC10 gives an integral bounded by
\[
 C\int_{B\lambda^2}^{\infty}(v+\lambda)^2
             e^{-cv/(\lambda+\sqrt v)}dv.
\]
For $v\ge B\lambda^2$, $v/(\lambda+\sqrt v)\ge\sqrt v/2$.
Substitute $r=\sqrt v$ and integrate a fixed polynomial times
$e^{-cr/2}$.  Choose $B$ so large that this is at most $C\lambda^6/Q^2$.
We obtain the useful uniform estimate
\[
 |T_\chi(Z)|\le C\{Z^{-1/400}\lambda^2+
                  \lambda^6/Q+Z^{-1}\log^2(QZ)\}.     \tag{QP.22}
\]
This estimate extends to a nonprincipal character $\chi$ at any level
$N\le Q$.  Its primitive inducer $\chi^*$ is nonprincipal, and the deleted
Euler factors are nonzero on $\Re s>0$, so $\chi^*$ has the same required
zero-free box.  The difference of the two tails is a sum over the finitely
many deleted primes.  At each such prime,
\[
 \sum_{p^j>Z}\frac{\log p}{p^j}\le\frac{2\log p}{Z}.
\]
Summing gives $|T_\chi(Z)-T_{\chi^*}(Z)|\le2\log N/Z$, already
absorbed in (QP.22).  Thus (QP.22) holds under the literal
\texttt{GoodC} hypothesis, which does not require primitivity.

Let $x=Y^K$, $Z=x^{400}$ and
$Y^{\epsilon/6}\le\lambda\le Y^{\epsilon/4}$, where
$0<\epsilon\le1$, $K\ge1$.  Then $3\le Z\le V$ eventually and
\[
 xZ^{-1/400}\lambda^2=\lambda^2\le Y^{\epsilon/2},
 \qquad x\lambda^6/Q\longrightarrow0,
 \qquad xZ^{-1}\log^2(QZ)\longrightarrow0.
\]
The last two limits are uniform in the stated range of $\lambda$.
Thus $x|T_\chi(Z)|\le Y^\epsilon$ eventually, exactly \texttt{tailK}.
For $y\le x$, the exact centered identity is
\[
 G_\chi(y)=-\sum_{n\le Z}\Lambda(n)\chi(n)\min(1,y/n)
                                                 -yT_\chi(Z).\tag{QP.23}
\]
It follows by expanding the finite sum and (QP.21), with the same
$+yL'/L(1,\chi)$ centering as CC12.  Taking $K=1$ and using the
coarser form of (QP.15) also supplies \texttt{kappaC}, with an absolute
coefficient of $\log Y$.

\paragraph{QP15: the exact coefficient four follows from a log-log bound.}
In (QP.22) instead choose $Z=\lambda^{1000}$.
The first term is $\lambda^{-1/2}$, and the other two are bounded
absolutely for sufficiently large $\lambda$.
By (QP.15),
\[
 |L'/L(1,\chi)|\le1000\log\log Q+C.                  \tag{QP.24}
\]
This proves the source's \texttt{kappa4\_\allowbreak loglog} conclusion by the
two-range explicit formula, without the longer differenced-contour route.
If $0<\epsilon<1/52$, $A>0$, and
$Y^{\epsilon/6}\le\log Q\le A Y^{\epsilon/5}\log^2Y$, then
\[
 1000\log\log Q+C
 \le200\epsilon\log Y+2000\log\log Y+1000\log A+C
                                                    \le4\log Y
                                                               \tag{QP.25}
\]
for all sufficiently large $Y$.  The strict gap is
$4-200/52=2/13>0$.  Thresholds may depend on $\epsilon,A$, but not on
$Q,N,\chi$.  This proves the literal \texttt{kappa4} bound needed in
\texttt{PairGood}; a bound with an unspecified larger coefficient would
not have sufficed.

\paragraph{QP16: a hybrid density theorem counts bad primitive boxes.}
Apply (HD) at $\sigma=399/400$, $T=Q$, $h=1/100$.
The exponent is
$3(2/400+1/100)=9/200<1/4$.  Each primitive character of conductor
at most $Q$ whose box fails has at least one zero in this rectangle,
so the number of such nonprincipal characters is at most
\[
                             C Q^{9/200}\le C Q^{1/4}.\tag{QP.26}
\]
The one primitive principal character may be added separately.  In
particular a totalized value assigned to its pole at one is not counted
as an analytic zero in (HD); it costs at most one character and is
absorbed by changing $C$.  This is the existential
\texttt{ZeroDensityBox} endpoint actually used
by \texttt{count\_\allowbreak bad\_\allowbreak core}.  No explicit enormous numerical constant
from another density interface is claimed.
Moreover, if each bad prime $q>F$ is assigned a primitive character of
conductor $fq\le Q$ with $f\le F$, this assignment is injective.
Indeed $fq=f'q'$ implies $q\mid f'$ or $q=q'$; the first case is
impossible since $0<f'\le F<q$.  Thus the same $CQ^{1/4}$ bound applies
to the number of bad primes.  The argument permits the character to
depend on the prime in any manner.

\paragraph{QP17: primitive twists and induced-character quantifiers.}
Let $\xi$ be primitive of conductor $f$, and let $q$ be an odd prime
with $q>f$.  The quadratic character $(\cdot/q)$ is primitive modulo
$q$.  The Chinese remainder decomposition of unit groups shows that
$\xi(\cdot)(\cdot/q)$ is primitive of conductor $fq$: factoring through
a proper divisor would factor one of the two primitive components.
It is nonprincipal because its $q$ component is nontrivial.

There is also a useful uniqueness assertion in the source's
\texttt{Ind} convention.  If two primitive characters agree on all
integers coprime to a positive ambient modulus $M$, then they have the same
conductor and are the same primitive character.  To prove this, compare
them at a unit modulo the product of their conductors and use the
Chinese remainder theorem to choose a congruent representative avoiding
the primes of $M$ not in that product.  They induce the same character
at the product level; uniqueness of the conductor proves the assertion.
In the uses here the prescribed twist is a character at a known common
level, so this argument identifies every primitive character quantified
by \texttt{Ind}, not just a preferred choice of representative.
Removing an unused extra prime from the ambient level is justified by
the same argument.  This supplies the \texttt{ind\_\allowbreak drop} step when the
second quadratic exponent is zero.

\paragraph{QP18: the generic good-prime stage.}
Fix $0<\epsilon\le1/4$, $\tau>0$, $K\ge1$, $\eta>0$, and an integer
$R\ge0$.  In admitted data satisfying the prime-divisor condition of
QP10, let $1\le F\le\sqrt U$ and let $\Xi$ be any family of at most
$R$ primitive characters, each of conductor at most $F$.
For $q\in P$, QP17 gives nonprincipal primitive twists of conductor at
most $FU$; QP16 shows that at most $C(FU)^{1/4}\le CU^{3/8}$ primes
can have a bad box for one of these twists.
One can assign a single witness to each such prime, so this count does
not acquire a factor $R$.
By (QP.11), this is less than $\eta|P|/2$ eventually.

For each $\xi\in\Xi$, use QP12 with $\kappa(n)=\xi(n)$.
Outside another set of at most $\eta|P|/2$ primes, all its truncated
sums satisfy $|S_N(y,q,\xi)|\le y^{1/2+\tau}/2$ throughout
$[Y,Y^K]$.  Here $N=\lfloor(Y^K)^{400}\rfloor$.
For $Q=FU$, (QP.10) and $F\le\sqrt U$ give
$Y^{\epsilon/6}\le\log Q\le Y^{\epsilon/4}$ eventually.
QP14 and $Y^\epsilon\le y^{1/2+\tau}/2$ for $\epsilon\le1/4$
then give $|G_\chi(y)|\le y^{1/2+\tau}$.
QP14's \texttt{kappaC} supplies one absolute logarithmic-derivative
coefficient, and QP17 transfers these conclusions to every primitive
inducer at every permitted ambient level.
Thus all but $\eta|P|$ primes satisfy the full \texttt{BoxAll} conditions
for every base character.  For $\eta<1$ at least one such prime exists.
The threshold depends only on the fixed displayed parameters, including
$R$, and is uniform in the moving family and in $F$.

\paragraph{QP19: the first prime can be chosen uniformly over all base characters.}
Fix $N_0>0$, $0<\epsilon<1/52$, $\tau>0$, $K\ge1$.
Set $z=Y^{\epsilon/5}$ and
$P_1=\{q:\sqrt{e^z}<q\le e^z,\ q\equiv3\pmod8,\ q\text{ prime}\}$.
Use QP18 with $m=8$, $a=3$, $U=e^z$, $F=N_0$, and the primitive
inducers of all characters modulo $N_0$.
Their number is at most $\varphi(N_0)$; $F\le\sqrt U$ eventually.
The admitted inequalities and the prime-divisor condition hold for
large $Y$.  Thus $P_1$ is nonempty and all but any prescribed small
fraction satisfy \texttt{GoodQ} for every character modulo $N_0$.
Also $q>N_0$ and $q>(Y^K)^{400}$ uniformly for these primes eventually,
because $\log q>z/2$ dominates every fixed multiple of $\log Y$.
The constant in the \texttt{Box} records is independent of $N_0$.

\paragraph{QP20: a second progression with one large modulus prime.}
For any $q\in P_1$ put
$M_o=q\prod_{3\le p\le z, p\ {\rm prime}}p$.
The conditions $d\equiv7\pmod{16}$ and $d\equiv-1\pmod{M_o}$
give a single reduced class $a_0$ modulo $m_0=16M_o$ by CRT.
Let $U=e^{z'}$, where
\[
 z\log^2Y\le z'\le K^2 z\log^2Y.                    \tag{QP.27}
\]
The elementary primorial bound gives
$\log m_0\le\log16+z+z\log4\le4z$ eventually.
Apply QP18 with the enlarged exponent $\epsilon'=6\epsilon/5$.
Its lower admitted bound is $Y^{\epsilon'/6}=z$;
its upper one follows from $K^2\le Y^{\epsilon/25}$ eventually.
The remaining admitted inequality follows from $\log m_0\le4z$.
Every prime divisor of $m_0$ is at most $Y^{\epsilon'/5}$ except $q$,
which is larger than $(Y^K)^{400}$ by QP19.  This is precisely why the
large-prime alternative in QP10 was retained.

Take $F=N_0e^z$ and the primitive inducers of
$\xi(\cdot)(\cdot/q)^i$, $\xi\bmod N_0$, $i=0,1$.
There are at most $2\varphi(N_0)$ of them, their conductors are at most
$N_0q\le F$, and $F\le\sqrt U$ by (QP.27) for large $Y$.
Thus the second prime set is nonempty and almost all its elements
satisfy \texttt{GoodD} simultaneously for these base characters.
All thresholds are uniform over $q\in P_1$ and $z'$ in (QP.27).
The congruences also imply $q<d$ and that every odd prime $p\le z$,
and $p=q$, divides $d+1$.

\paragraph{QP21: positive proportions give the actual simultaneous pair.}
Take the good proportion at both stages to be at least $1/2$.
QP19 supplies a nonempty good subset of $P_1$; for every prime in that
subset, QP20 supplies a nonempty good subset of its second progression.
Choose $q$ and then $d$ from these subsets.  Use the maximum of the two
absolute \texttt{Box} constants, and the maximum of their finitely many
thresholds.  The \texttt{Box} conditions are monotone in that constant
and in a decrease of the upper endpoint $X$.
We obtain one absolute $C>0$ such that, for every fixed
$\epsilon,\tau,K,N_0$ as above and all sufficiently large $Y$, every
$X\le Y^K$ and every $z'$ in (QP.27) have primes $q,d$ satisfying
\texttt{Rows}$(Y^{\epsilon/5},z',q,d)$, $N_0<q$,
\texttt{GoodQ} and \texttt{GoodD}.  In particular both windows,
both prime congruences, their ordering, and every divisibility
condition hold for the same pair.
This proves the literal \texttt{pair\_\allowbreak share} and \texttt{pair\_\allowbreak exists}
conclusions under the named classical inputs above; positivity of the
proportions is used together with positivity of the actual prime counts.

\paragraph{QP22: install the complete PairGood record.}
Given $M>0$, $0<\epsilon<1/52$, $A\ge1$, apply QP21 with
$N_0=8M$, $\tau=\epsilon/2$, $K=A$, $X=Y^A$, and
$z'=Y^{\epsilon/5}\log^2Y$.
The resulting primes are coprime to $M$ and satisfy the literal
\texttt{WinPair} congruences and logarithmic windows.
For a character $\xi\bmod M$ and $(i,j)\in\{0,1\}^2\setminus\{(0,0)\}$,
inflate $\xi$ to modulus $8M$ and consider
$\xi(n)(n/q)^i(n/d)^j$ on integers coprime to $2Mqd$.
If $j=0$ then $i=1$; QP17 removes the unused $d$ and \texttt{GoodQ}
applies.  If $j=1$, \texttt{GoodD} applies for either value of $i$.
Every primitive inducer is nonprincipal and is covered by this argument.

Its height parameter is exactly
\[
 Q_j=8M\exp(Y^{\epsilon/5})
            \exp(Y^{\epsilon/5}\log^2Y)^j,\qquad
 Y^{\epsilon/6}\le\log Q_j\le3Y^{\epsilon/5}\log^2Y.  \tag{QP.28}
\]
Thus (QP.25) gives $|L'/L(1,\chi)|\le4\log Y\le4\log y$.
The centered bound is $|G_\chi(y)|\le y^{1/2+\epsilon/2}$
for every real $Y\le y\le Y^A$; on the finite sum $n\le y$,
$\max(y/n-1,0)=y/n-1$, so this is exactly \texttt{An.gChi}.
Applying the same argument to every real primitive base character with
conductor dividing $8M$ gives the required \texttt{ZeroBox} at height
$Q_j$.  Therefore the same $q,d$ satisfy all four clauses of
\texttt{PairGood}, with its universal induced-character quantifiers.
This proves \texttt{UpperTheory.Top.PairEx}.

\paragraph{QP23: real-character factorization gives the actual host exclusion.}
Let an already specified host $H$ have distinguished primes $q,d$ from
QP22, with $(H.QS,qd)=1$, every prime of $H.QS$ dividing $2M$, and
$\operatorname{ordCompl}_2(H.QS)\mid M$.
For a real primitive character belonging to $H$, its conductor divides
$H.QS\,qd$.  CRT decomposes it into its prime-power components.
At an odd prime the subgroup $1+p\mathbb Z/p^r\mathbb Z$ has odd
order, so a quadratic character is trivial there; its conductor
exponent is at most one.  At $2$, every unit congruent to $1\pmod8$
is a square (lift its odd square root one binary digit at a time), so
the conductor exponent is at most three.
The $q$ and $d$ components are therefore either trivial or the unique
quadratic characters modulo those primes.  Removing them leaves a
real primitive character $\rho$ with
$\rho.N\mid8\operatorname{ordCompl}_2(H.QS)$ and
\[
 \chi.N=\rho.N\,q^i d^j,\qquad
 \chi(n)=\rho(n)(n/q)^i(n/d)^j\quad((n,\chi.N)=1).    \tag{QP.29}
\]
If $q\mid\chi.N$ or $d\mid\chi.N$, then $(i,j)\ne(0,0)$.
QP22 gives the box for $\chi$; its height is positive and hence includes
every real zero.  This proves \texttt{RealBox H}:
$L(\beta,\chi)\ne0$ for every $\beta>399/400$ of the specified real
characters.  Moreover (QP.29) and the upper prime windows give the
literal conductor bound
\[
 \chi.N\le8\operatorname{ordCompl}_2(H.QS)
          e^{Y^{\epsilon/5}}
          e^{Y^{\epsilon/5}\log^2Y\,1_{d\mid\chi.N}}. \tag{QP.30}
\]
This is the same-host \texttt{realBox} endpoint; no replacement host
or enlargement of its prescribed fields has occurred.

\paragraph{QP24: close the exceptional-zero premise in the thin host estimate.}
Apply CC19--CC20 to the very host in QP23.
Whenever its other displayed hypotheses hold,
$\operatorname{ordCompl}_2(H.QS)\le B$ and
$\log H.\mathrm{condN}\le2x^{\epsilon/5}\log^2x$,
QP23 supplies the \texttt{RealBox H} premise previously retained there.
The conclusion is the actual \texttt{H.ChebyshevCharacterBound}
with the same centering, the same lower range
$y\ge\frac13\exp(x^{\epsilon/3})$, and one uniform positive decay
constant.  Thus the CC thin estimate now has a concrete prime-choice
supplier for its exceptional-zero exclusion whenever the chosen host
uses the pair of QP22 and satisfies QP23's level conditions.
This assertion does not construct an arbitrary host, prove a missing
geometric window, or supply the global automorphic input.  Those
requirements are independent of the real-character factorization.

\paragraph{QP25: the remaining window premise is not lost in a composition.}
The source's \texttt{winPairs} conclusion follows from QP22 and an
independently supplied \texttt{VWinAt}$(\psi,401)$, with the exact order
of quantifiers preserved.  To check the rescaling, choose the integer
page function with
$u^{2+8\epsilon}\le2^{e(u)}\le2u^{2+8\epsilon}$.
Let $J\ge1$ and its initial threshold come from \texttt{VWinAt}.
For a desired factor $c\ge1$, put $k=c^{1/(2+8\epsilon)}$.
Choose $M$ above the PairEx threshold and above $k^3,k^{400}$, and take
$X\ge M^J$.  QP22 gives a good pair at $Y=X^{1/J}$.
The window premise supplies $W\ge Y$, $W^2\le X$ and its geometric
window for that same pair.  Put $Z=kW$.
For $Z\le x\le Z^{3/2}$, one has
$W\le x/k\le x\le W^2$, because $W\ge k^3$.
Use the page $e(x/k)$ and multiply the preceding page inequalities by
$c$.  The geometric lower cutoff transfers because
$k^{400}\le x/k$, which gives
$x^{400}\le(x/k)^{401}$; its upper cutoff is monotone.
PairGood transfers from $X$ to $x$ by monotonicity.
This proves the literal \texttt{WinDatum} conclusion on
$[Z,Z^{3/2}]$ for arbitrarily large $Z$, conditional on the original
\texttt{VWinAt} input.  It neither proves that input nor advances any
global RH/GRH branch solely by reusing this implication.


\section{A closed coefficient-representation chain by Gaussian elimination}
\label{sec:highest-weight-closure}

\status{HP1--HP6 remove the retained characteristic-zero polynomial
highest-weight inputs from MR5/MR19 and AR13--AR14.  The proof applies to
the literal polynomial module used in the source, over every field of
characteristic zero.  It proves the two highest-fixed-vector interfaces,
absolute tensor irreducibility, classification, and the precise
contragredient and scalar-extension statements, then reaches the full
\texttt{wt\_\allowbreak weights} endpoint on its original data.  The argument uses
finite polynomial equations, Gaussian elimination, and a factorial
pairing.  It does not assume a Schur-module irreducibility theorem,
highest-weight classification, or a noncompact representation theorem.
The compatible geometric/theta/local state and the separate automorphic
Galois inputs remain the frontier of the larger form construction.}

We use indices $1,\ldots,n$ in this section.  In the source they are
$0,\ldots,n-1$, with the evident shift.  If $\lambda_1\ge\cdots\ge\lambda_n$
are integers, put $d=\lambda_n$ and $\nu_i=\lambda_i-d$; for $n=0$ use
$d=0$.  Let $X=(x_{ij})$ be the matrix of polynomial variables.  Write
$U_+$ and $U_-$ for the upper and lower unitriangular groups and $T$ for
the diagonal group.  Define
\[
 \mathcal N_\nu(K)=\{f\in K[x_{ij}]:
 f(Xt)=t_1^{\nu_1}\cdots t_n^{\nu_n}f(X),\quad
 f(Xu)=f(X)\ (u\in U_+)\}.
 \tag{HP.1}
\]
The first equality means column multihomogeneity.  Both equalities can
be read as polynomial identities.  Over an infinite field, equality at
all indicated points is equivalent to this reading.

This is exactly the source's \texttt{dualWeyl K n n nu}, hence
\texttt{nablaSq K n lambda} for $\nu=\lambda.\mathrm{natW}$.
The action, including its transpose and determinant, is
\[
       (\rho_\lambda(g)f)(X)=(\det g)^d f(g^tX).
 \tag{HP.2}
\]
It is \texttt{sqAct}, not substitution by $gX$ or $g^{-1}X$.
For a finite index set $I$ and weights $\lambda_\tau$, use separate
matrix variables $X_\tau$ and take the image of the tensor product
$\bigotimes_{\tau\in I}\mathcal N_{\nu_\tau}(K)$ under multiplication.
The polynomial monomial bases in disjoint variables show that this
map is injective.  Its image is the literal \texttt{VV K lambda};
(HP.2) factor by factor is its \texttt{rhoK} action.  Empty $I$ or
$n=0$ gives $K$, with highest vector $1$.

\paragraph{HP1: the literal fixed line and its rational structure.}
\label{claim:HP1}
Every $\mathcal N_\nu(K)$ is the solution of a finite system of linear
equations with integer coefficients inside a finite polynomial space.
Indeed, column multidegree $\nu$ forces total degree $N=\sum\nu_i$.
For each $i<j$, substitute column $j+t$ column $i$ for column $j$ and
set all coefficients of $t^a x^\alpha$ in the difference to zero.
There are finitely many equations because its degree is at most $N$.
The elementary column operations generate $U_+$.  These are the same
equations as the generic \texttt{colAddGen} kernels in (HP.1).
Thus, in characteristic zero, row reduction over $\mathbb Q$ gives
\[
 K\otimes_{\mathbb Q}\mathcal N_\nu(\mathbb Q)
       \simeq\mathcal N_\nu(K).
 \tag{HP.3}
\]
No integral freeness or equality between the integral Weyl and dual
Weyl lattices is needed or asserted.  The same argument applies to
the tensor polynomial space; it also identifies its scalar extension
with the same polynomial subspace over the new field.  Moreover, the
joint equations in separate blocks describe exactly that tensor space:
expand in a basis of the other blocks, apply one block's equations to
each coefficient, and repeat.  Equivalently, choose a complement to
each one-block kernel in its finite polynomial space; the tensor
decomposition then shows that a term with any complementary factor
cannot satisfy all the block equations.

Let $\Delta_r(X)$ be the leading $r$ by $r$ minor and set
$\nu_{n+1}=0$.  Then
\[
             h_\nu(X)=\prod_{r=1}^n
                   \Delta_r(X)^{\nu_r-\nu_{r+1}}
 \tag{HP.4}
\]
belongs to $\mathcal N_\nu(K)$, is nonzero, and is fixed by
$\rho_\lambda(U_+)$.  Each principal minor is unchanged by lower
unitriangular row operations and by upper unitriangular column
operations.  The column weights telescope to $\nu$, and
$h_\nu(1)=1$.  The determinant twist is $1$ on $U_+$.
The diagonal action on this vector is the character $\lambda$:
\[
 \rho_\lambda(\operatorname{diag}t)h_\nu
                      =\prod_i t_i^{\lambda_i}h_\nu.
 \tag{HP.5}
\]

Conversely, suppose $f\in\mathcal N_\nu(K)$ is fixed by
$\rho_\lambda(U_+)$.  This says $f(lX)=f(X)$ for $l\in U_-$.
On the open set where every $\Delta_r(X)$ is nonzero, Gaussian
elimination gives $X=l a u$ with $l\in U_-$, $a$ diagonal and
$u\in U_+$, and $a_i=\Delta_i/\Delta_{i-1}$, where $\Delta_0=1$.
Consequently
\[
 f(X)=f(a)=f(1)\prod_i a_i^{\nu_i}
                =f(1)h_\nu(X).                    \tag{HP.6}
\]
This open set is dense over any infinite field: a polynomial vanishing
there has product with a suitable power of $\prod_r\Delta_r$ equal
to zero, and the polynomial ring is a domain.  Equivalently, its
Gaussian coordinates give an open rational parametrization and the
usual polynomial identity principle.  Hence (HP.6) holds identically.
For several blocks apply Gaussian elimination in every block.  Thus
\[
 (\mathrm{VV}_K(\lambda))^{\prod U_+}
       =K h,\qquad h=\prod_{\tau\in I}h_{\nu_\tau}(X_\tau).
 \tag{HP.7}
\]
This is the exact source \texttt{HwFixedLine}; its distinguished vector
is \texttt{mtwHw}, not a highest vector chosen after an isomorphism.

We will also use the same argument before dominance is known.
If $\nu\in\mathbb N^n$ is arbitrary and a nonzero
$\mathcal N_\nu(K)$ contains a nonzero left-$U_-$-fixed polynomial,
(HP.6) holds as an equality of rational functions and forces
$\nu_r\ge\nu_{r+1}$.  For $r<n$, take the identity matrix except
for the rows/columns $r,r+1$, where the block is
$\left(\begin{smallmatrix}t&-1\\1&0\end{smallmatrix}\right)$.
All its leading minors are $1$, except $\Delta_r=t$.
Then (HP.6) reads $f(X_t)=f(1)t^{\nu_r-\nu_{r+1}}$ for $t\ne0$.
A polynomial cannot equal a negative power times a nonzero constant.
Moreover $f(1)\ne0$ by the identity principle.  This proves the claimed
dominance without invoking a classification theorem.

\paragraph{HP2: fixed-vector existence and absolute irreducibility.}
\label{claim:HP2}
Let $S$ be a nonzero subspace of $\mathrm{VV}_K(\lambda)$ stable under
$G=\prod_{\tau\in I}\mathrm{GL}_n(K)$, with $\operatorname{char}K=0$.
The diagonal action has finitely many weights.  For the polynomial
model these are the row multidegrees, shifted by $d_\tau$ in every
coordinate of block $\tau$.  The weight projections preserve $S$:
choose a diagonal element at which these finitely many different
Laurent characters have distinct values, and interpolate its spectral
projections by polynomials.  Such an element exists because a finite
product of nonzero Laurent polynomials cannot vanish on the whole
torus over an infinite field.

Choose a nonzero vector $v\in S$ of weight $\mu$ maximizing
$H(\mu)=\sum_{\tau,i}(n-i)\mu_{\tau,i}$ among weights occurring in $S$.
For an upper elementary matrix $g=1+tE_{ij}$, $i<j$, the polynomial
$\rho_\lambda(g)v$ has coefficients $v_a$ of weights
$\mu+a(e_{\tau,i}-e_{\tau,j})$.  This follows directly from replacing
row $j$ of $X_\tau$ by row $j+t$ row $i$; the determinant twist is $1$.
All $v_a$ belong to $S$, since coefficients are linear combinations of
values at sufficiently many distinct $t\in K$.  For $a>0$ their heights
are strictly greater than $H(\mu)$, so they vanish.  Hence $v$ is fixed
by every elementary upper unipotent and therefore by $\prod U_+$.
Together with HP1 this proves
\[
             S\ne0\quad\Longrightarrow\quad h\in S.
 \tag{HP.8}
\]
In particular this proves the exact \texttt{HwFixedEx} interface,
even over any characteristic-zero field.  The proof also works for
$\mathcal N_\nu$ with arbitrary nonnegative $\nu$; with HP1 it shows
that a nonzero such space necessarily has dominant $\nu$.

To show that (HP.8) implies simplicity here, the required duality
must be justified.  On the polynomial ring define the bilinear
factorial pairing
\[
        B(f,g)=\sum_\alpha \alpha! f_\alpha g_\alpha,
       \qquad \alpha!=\prod_j\alpha_j! .             \tag{HP.9}
\]
It has no complex conjugation.  Its restriction to
$\mathrm{VV}_{\mathbb Q}(\lambda)$ is nondegenerate: every nonzero
rational polynomial $f$ has $B(f,f)>0$ when read in $\mathbb R$.
A rational basis therefore has a positive-definite rational Gram
matrix, whose nonzero rational determinant remains nonzero in every
characteristic-zero field.  Equation (HP.3) proves nondegeneracy on
the entire $\mathrm{VV}_K(\lambda)$, not merely on an arbitrary
subspace of polynomials.  In particular $B(h,h)$ is a strictly
positive integer and remains nonzero in $K$.

The substitution identity for this pairing is
\[
              B(\rho_\lambda(g)f,z)
                       =B(f,\rho_\lambda(g^t)z).     \tag{HP.10}
\]
One direct verification writes
$B(f,z)=f(\partial)z|_{X=0}$ and uses the linear chain rule; checking
monomials gives the same multinomial coefficients $\alpha!$.
Both determinant factors are $(\det g)^{d_\tau}$, so the identity
holds for negative determinant twists as well as nonnegative ones.
It holds block by block for the tensor product.

If $S$ is proper, its annihilator $S^{\perp_B}$ is nonzero by
nondegeneracy.  It is stable under $G$, because for $x\in S^{\perp_B}$
and $y\in S$,
$B(\rho(g)x,y)=B(x,\rho(g^t)y)=0$.
Equation (HP.8) now puts $h$ in both $S$ and $S^{\perp_B}$,
contradicting $B(h,h)\ne0$.  Thus $S$ is the whole space.  This proves
simplicity over every characteristic-zero field and, by (HP.3), after
every extension of that field.  It closes the dual-Weyl, tensor and
highest-fixed-line chain with finite-dimensional polynomial arguments.

\paragraph{HP3: the actual rational unitary group and every coefficient field.}
\label{claim:HP3}
Let $E$ be a CM field, $c$ its specified involution, and let
$\sigma_\tau:E\hookrightarrow K$ be pairwise distinct embeddings
with $\sigma_\tau\ne\sigma_{\tau'}c$ for all $\tau,\tau'$.
Then the representation
\[
 \gamma\longmapsto
   \rho_K((\sigma_\tau(\gamma))_\tau)
       \quad(\gamma\in U_n(E/E^+)(\mathbb Q))          \tag{HP.11}
\]
on the same $\mathrm{VV}_K(\lambda)$ is absolutely irreducible.
Here the rational notation denotes the rational points of the
restriction of scalars from $E^+$ to $\mathbb Q$, exactly the source's
\texttt{GQ}, with its injective finite-adelic realization forgotten.

For completeness the density assertion needed in this passage has a
short proof.  The rational vector space
$\mathfrak s=\{X\in M_n(E):X+{}^t c(X)=0\}$ has the Cayley map
\[
                 X\longmapsto(1-X)(1+X)^{-1}.        \tag{HP.12}
\]
It is a rational isomorphism between nonempty open subsets of
$\mathfrak s$ and the rational unitary group, with the same inverse
formula.  In rational coordinates a polynomial over any extension
field which vanishes on all $\mathbb Q$-points is zero: induct on the
number of coordinates and use infinitely many rational values in one
coordinate.  Removing the finitely many Cayley denominators preserves
this argument.  Hence the rational points are Zariski dense, including
after extension to an algebraic closure $\overline K$.

Over $\overline K$, the pairs of conjugate embeddings of $E$ give
\[
 (\operatorname{Res}_{E^+/\mathbb Q}U_n)_{\overline K}
                \simeq\prod_{\{\sigma,\sigma c\}}\mathrm{GL}_n.
 \tag{HP.13}
\]
Indeed the unitary equation reads $A_{\sigma c}=A_\sigma^{-t}$,
so one matrix in each pair is arbitrary.  Projection to the selected
$\sigma_\tau$ is surjective by the distinct-nonconjugate hypothesis.
A stable subspace is defined by the equations obtained by composing
its inclusion, the action matrix, and its quotient map.  These are
regular equations in the group entries; density makes them hold on
all of (HP.13).  Its scalar extension is consequently stable under
all selected general linear groups and is zero or full by HP2.
Faithfulness of scalar extension gives the same conclusion over $K$.
The argument applies after every $K\hookrightarrow K'$, since
injectivity preserves the distinctness conditions.  It proves the
literal \texttt{tenRepF\_\allowbreak isIrreducible} and its base-change clauses;
no embedding of all of $K$ into $\mathbb C$ is assumed.

The action matrices are algebraic in exactly the source's sense.
The polynomial substitution part is polynomial in matrix entries.
A negative power of $\det\sigma_\tau(\gamma)$ is replaced, on the
unitary group, by the corresponding nonnegative power of
$\det\sigma_\tau c(\gamma)$.  Express every entry of $\gamma$ in a
fixed $\mathbb Q$-basis of $E$; each embedding is a linear polynomial
in those coordinates.  This gives the finite polynomial matrix
presentation in \texttt{IsAlgWeightF}.  In the dual representation
replace $\gamma^{-1}$ by ${}^t c(\gamma)$ and transpose its action
matrix.  Its coefficients are again polynomials in the same rational
coordinates, over the original $K$.

\paragraph{HP4: classification and the source's dual conventions.}
\label{claim:HP4}
We first derive the classification used in AR13, rather than leaving
it as an additional representation theorem.  An algebraic matrix
representation of the rational unitary points extends to the split
product (HP.13): its identity and multiplication equations hold on
dense rational points, hence identically.  Its matrix entries are
regular functions of the general linear matrices, thus polynomials
in their entries and inverse determinants.

For an arbitrary finite-dimensional rational representation of a
product of general linear groups over $\mathbb C$, diagonal torus
restriction has a finite direct sum of integral weights.  This can
be seen without simultaneous triangularization: expand its matrix
as $R(t)=\sum_\mu t^\mu A_\mu$.  The equations
$R(st)=R(s)R(t)$ and $R(1)=1$ give
$A_\mu A_\eta=0$ for $\mu\ne\eta$, $A_\mu^2=A_\mu$, and
$\sum A_\mu=1$.  Choose a maximal-height weight as in HP2.
Restriction to $1+tE_{ij}$ is polynomial in $t$, since its determinant
is $1$; conjugation by the diagonal torus shows that its coefficient
of $t^a$ raises that weight by $a(e_i-e_j)$.  Hence a maximal-height
weight vector $v$ is fixed by the upper unipotents and transforms
under the diagonal by some integral weight tuple $\lambda$.

Suppose now that the representation $R$ on $W$ is irreducible.
Its orbit of $v$ spans $W$.  For sufficiently large nonnegative
integers $m_\tau$, every matrix coefficient
\[
 f_\phi((X_\tau))=
       \prod_\tau\det(X_\tau)^{m_\tau}\,
                  \phi(R((X_\tau))v),\qquad \phi\in W^*,
 \tag{HP.14}
\]
is polynomial.  These polynomials have right-upper invariance and
column weights $\mu_{\tau,i}=\lambda_{\tau,i}+m_\tau$.
They form a nonzero invariant subspace of the tensor polynomial
space for these column degrees; the map $\phi\mapsto f_\phi$ is
injective because the orbit of $v$ spans $W$.
The column weights are nonnegative because the functions are
polynomial.  HP1--HP2, also in their stated pre-dominance version,
show that $\mu_\tau$ is dominant and that this entire polynomial
space is irreducible.  Thus the coefficient map is onto it.

Its action is exactly
\[
 f_\phi(g^tX)=\prod_\tau\det(g_\tau)^{m_\tau}
                                f_{\phi\circ R(g^t)}(X).
 \tag{HP.15}
\]
Let $\iota(g)=g^{-t}$.  Equation (HP.15) identifies
$(R^*\circ\iota)\otimes\prod\det^{m_\tau}$ with this polynomial
space.  Multiplication by
$\prod\det(X_\tau)^{\mu_{\tau,n}}$ is a nonzero equivariant map
from the canonical tensor module with column degrees
$\mu_{\tau,i}-\mu_{\tau,n}$, with the corresponding determinant
twist, onto that same space.  It is an isomorphism by HP2.
After canceling $\prod\det^{m_\tau}$ the resulting canonical weight
is precisely $\lambda$.  Finally (HP.10) gives the explicit
isomorphism
\[
      V_\lambda\simeq V_\lambda^*\circ\iota,
             \qquad x\longmapsto B(x,-).
 \tag{HP.16}
\]
The operation $R\mapsto R^*\circ\iota$ is involutive, so
$R\simeq V_\lambda$.  The fixed line (HP.7) proves uniqueness of
$\lambda$.  This establishes the required complex classification
and dominance by a matrix-coefficient argument.

The distinction between the two coefficient conventions matters:
\[
\begin{array}{c|c|c}
\text{source representation}&\text{action on rational }\gamma
       &\text{representation whose highest weight is }\lambda\\ \hline
\texttt{tenRepF}(\sigma,\lambda)&\rho_K(\sigma\gamma)
       &\texttt{tenRepF}\ \text{itself}\\
\texttt{wtRepF}(\sigma,\lambda)&\texttt{tenRepF}(\sigma,\lambda)^*
       &\texttt{wtRepF}^{*}\simeq\texttt{tenRepF}\\
\texttt{rhoVr}/\texttt{mtRep}(\sigma,\lambda)
       &\rho_K((\sigma\gamma)^{-t})
       &\texttt{mtRep}^{*}\simeq\texttt{tenRepF}.
\end{array}                                                   \tag{HP.17}
\]
The last equivalence uses the bilinear form (HP.10), and the middle
one is the natural double-dual evaluation map.  For unitary $\gamma$,
$(\sigma\gamma)^{-t}=\sigma c(\gamma)$, so the last line is also
\texttt{tenRep} at the conjugate CM type.  Taking a single dual at
an unchanged embedding would give the usual reversed negative
weight; (HP.17) explains exactly why the source's \emph{dual} weight
is the displayed $\lambda$, with no reversed-negative substitution.

This gives both existence and uniqueness for the literal
\texttt{HWAtF} predicate, including its polynomial-presentation
quantifier.  For existence use the matrices just constructed and the
coordinates of $h$.  Evaluating them at \texttt{evPtF(c,sigma,B)}
puts $B$ in the selected factor, $B^{-t}$ in its conjugate, and the
identity in the others; (HP.5) gives the required Borel eigenvalue.
For uniqueness, any competing polynomial presentation agrees with
the actual one on rational unitary points.  HP3 extends this agreement
to the split group, so it agrees at those exact evaluation points.
A nonzero vector satisfying its Borel equations in one factor may
have arbitrary components in other factors.  Apply a linear functional
to the other tensor factors that does not kill it.  The result is a
nonzero vector in that one polynomial factor, fixed by $U_+$ and of
the proposed diagonal weight.  HP1 and (HP.5) force that weight to be
$\lambda_\tau$.  Distinct Laurent characters are different over an
infinite field.  Thus no change of basis or alternative presentation
can manufacture a different highest weight at the same embedding.

\paragraph{HP5: coefficient descent and the closed CellWt branch.}
\label{claim:HP5}
Let $e:K\hookrightarrow\mathbb C$ and fix a complete CM type
$\sigma:E\hookrightarrow K$ as in AR13.  Let $W$ be an irreducible
finite-dimensional $K$-representation of the rational unitary group,
with the polynomial presentation obtained from its dual's
\texttt{HWAtF} data.  Complexify and choose a nonzero invariant subspace
of minimal dimension.  Apply the classification just proved to its
dual, then dualize back.  Using (HP.16)--(HP.17), this produces
a nonzero complex intertwiner
from a canonical inverse-transpose coefficient module
$V_{\lambda,\mathbb C}$ into $W_\mathbb C$.

The canonical module and its matrices are defined over $K$, by
(HP.3) and the original $\sigma$.  To descend that intertwiner,
choose a $K$-linear functional $\ell:\mathbb C\to K$ not vanishing
on one of its nonzero entries and apply $\ell$ entry by entry.
Every intertwining equation has coefficients in $K$, so the result
is a nonzero $K$-intertwiner.  This argument respects all rational
group elements simultaneously; it is not an unjustified intersection
of infinitely many kernels with scalar extension.
The canonical module is irreducible by HP3, so the descended map
is injective; $W$ is irreducible, so it is surjective.  Hence
\[
       W\simeq_K V_{\lambda,K}
       \quad\text{with its original inverse-transpose action}.
 \tag{HP.18}
\]
The prescribed dual highest weights transfer through this map;
HP4's uniqueness identifies them with this same $\lambda$.
This supplies both representation inputs retained in AR13.
It does not assume that an arbitrary irreducible $K$-representation
has irreducible complexification; in this case that conclusion
follows from (HP.18).

AR14 now completes its original \texttt{CellWt} comparison without
an extra highest-weight premise.  The source's infinitesimal-character
entries $\lambda_i+(n-1)/2-(i-1)$ are strictly decreasing, so equality
of their multisets identifies the dominant weights term by term.
The original \texttt{SideRead} field identifies the coefficient
embeddings, and the canonical polynomial bases commute with $e$.
Conjugating by the actual cell basis-change matrices gives equality
of the representation matrices in those specified bases, exactly
as in AR14.  All other AR14 hypotheses remain: its supplied actual
\texttt{MtReal}, actual \texttt{CharSide} infinitesimal parameters,
coefficient-field maps, and cell base-change witnesses.  The closed
branch is the representation classification and descent used within
that comparison, not the existence of those geometric data.

\paragraph{HP6: the full original-field weight supplier and the new frontier.}
\label{claim:HP6}
The source's \texttt{wt\_\allowbreak weights} theorem is now reconstructed with
no additional representation-theoretic input.  Retain its literal
hypotheses: \texttt{GlobalArtinReciprocity}, $D$, its frame, $j_E$,
\texttt{SideA}, star stability of $D.\mathrm{LL0}$, integral
\texttt{eExp}, and the subsequently quantified original
$\texttt{TwR}/\texttt{RingSpec}$.  The half-integrality needed here is already forced by these records:
$\texttt{cf.isEta.inf}$ says that the framed character has infinite
parameter $1/2$; $\texttt{A.chr.etaT}$ identifies it with the actual
$\eta_T$, and $\texttt{A.chr.infG}$ transports that parameter to
$D.naut$.  Thus a singleton $\{sh\}$ in the antecedent has
$sh=1/2$ (so the source's integer in $sh=z+1/2$ can be $z=0$).
The row-gap and parity calculation, already proved in MR19, gives dominant
$\lambda_j(\tau)$ with
\[
 \operatorname{infOf}(\lambda_j(\tau))
       =\operatorname{paramV}(D.R,D.hc_\tau,j-1,sh)
 \quad\text{when }\operatorname{infG}(\eta_T,\tau)=\{sh\}.
 \tag{HP.19}
\]
The pointwise $sh$ is uniquely determined whenever it exists; where
none exists the implication is empty and the zero weight supplies
the padding.  Use zero weights also outside $1\le j\le m$.

For each original coefficient ring $O_i$, put
\[
 K_i=\operatorname{Frac}(O_i),\qquad
 W_{i,j,b}=\mathrm{VV}_{K_i}(\lambda_j)^*,\qquad
 \rho_{i,j,b}=\texttt{wtRepF}(\texttt{kEmbT}_{i,b},\lambda_j).
 \tag{HP.20}
\]
The exact original side-dependent embeddings are distinct and
nonconjugate.  Explicitly, each is the same injective field reading
composed with the host automorphism labelled by the corresponding
CM-type element.  Equality of two such embeddings would give equality
of the labels; equality with a conjugate would put that same label
in both halves of the CM type.  Both are excluded by the actual frame.
Characteristic zero also follows from the coefficient records: the
finite algebra map $\mathbb Z_\ell\to O_{0,i}$ into a domain DVR is
injective.  Otherwise its prime kernel is $(\ell)$, making $O_{0,i}$
a finite-dimensional domain over $\mathbb F_\ell$, hence a field,
contrary to its DVR nonfield clause.  Composing with the stipulated
injection $O_{0,i}\to O_i$ and the fraction-field injection proves
$\operatorname{char}K_i=0$.  No larger coefficient field is selected.  The five required clauses
now have a single closed supply chain:
\begin{enumerate}
\item HP3 makes the tensor representation irreducible.  Its finite
      dual is irreducible, since the annihilator of a proper nonzero
      dual-stable subspace would be a proper nonzero original-stable
      subspace.  This gives \texttt{rhoW.IsIrreducible}.
\item HP3's explicit polynomial matrices give
      \texttt{IsAlgWeightF} on that very $W_{i,j,b}$.
\item Equation (HP.3), the natural dual-basis map, and injectivity of
      the original coefficient extension identify every
      \texttt{rhoKT} with the same construction over
      $K_{i,k}=\operatorname{Frac}(O_{i,k})$.
      HP3 therefore gives all its irreducibility clauses.
\item The natural double dual and HP4 supply \texttt{HWAtF} for
      $\rho_{i,j,b}^{*}$ at every prescribed embedding, with weights
      $\lambda_j(\tau)$.
\item HP4 forces every weight admitted by that literal predicate to
      equal $\lambda_j(\tau)$; (HP.19) therefore gives the final
      universally quantified infinitesimal-character identity.
\end{enumerate}
There is no representation-dependent extension to put into the lower
bound: one may take $c_2(\ell,j)=\bot$, exactly as in the source.
This closes the original-field, both-side, all-row, all-base-change
\texttt{WtData} supplier, and removes the retained highest-weight
premise from MR19 and from MR5's actual coefficient representation.

The closure is noncircular.  HP1--HP4 use only polynomial linear
algebra, Gaussian elimination, the explicitly proved factorial
pairing, and the Cayley parametrization of the given rational
unitary group.  HP5 then uses the already proved coefficient-extraction
descent and the supplied geometric comparisons.  HP6 uses the already
proved row-weight arithmetic only after the representation theorem is
available.  No \texttt{Core}, theta nonvanishing, automorphic Galois
representation, prime-window estimate, RH or GRH is used in proving
the representation leaf.

For the surrounding construction the frontier has genuinely changed:
\texttt{HwFixedEx}, \texttt{HwFixedLine}, canonical tensor
irreducibility/classification, the conjugate-CM-type dual highest
weight, and weight base-change irreducibility are no longer open
inputs.  \texttt{St/ColPin/R3.Read}, the original tuple and lattice
ties, the good-row spectrum, the supplied \texttt{SideA}, and the
compatible nonzero theta/toric/local installation remain required
where MR/AR use them.  This section supplies none of those by an
unrelated representation of the same dimension.


% Version 0.23: a complete downstream cohomological assembly.
% Four consolidated endpoints VW1--VW4. This fragment introduces no new macros.
\section{Closing the cohomological window assembly}\label{sec:cohomological-window-closure}

\status{The full implication from the two precise source interfaces
\texttt{JointWin} and \texttt{ThmBc} to \texttt{VWinAt} is proved below,
including the actual fields, all finite prime sums, and uniform thresholds.
The two input theorems themselves remain unreconstructed.  With QP22 this
closes the subsequent fixed-pair window construction.  It does not assert
an unconditional cohomological window or a complete RH/GRH proof.}

The usefulness of this assembly is that it removes a collection of
separate field, parity, conductor, scale and cutoff obligations from the
continuation frontier.  All the bounds concern the \emph{same} number
field $E$ specified by \texttt{IsChiField}; no new compatible-form or
coefficient-field witness is manufactured by an isomorphism.

\subsection{The measured groups and the exact three-part sum}

Fix a number field $E$, a rational prime $p$, an algebraic closure and
choices of places over the places of $E$ above $p$.  These are the
source's \texttt{LData}.  Put
\[
 \mathscr S_p(E,j)=\ker\!\left(
 H^1(G_{E,S_p},(\mathbb Q_p/\mathbb Z_p)(-j))
 \longrightarrow\prod_{v\mid p}
 H^1(E_v,(\mathbb Q_p/\mathbb Z_p)(-j))\right),
\]\[
 \ell_p(E,j)=\operatorname{length}_{\mathbb Z_p}\mathscr S_p(E,j).
\]
Here $G_{E,S_p}$ is the Galois group unramified outside the finite places
above $p$, with the usual allowance at infinity.  The lengths are initially
in $\mathbb N\cup\{\infty\}$.  This is exactly the source's negative twist:
\texttt{qzTw D j} uses the cyclotomic character to power $-j$.
We retain all the choices in each argument.  Every bound in this section
is required for \emph{every} algebraic closure and every such system of
local choices, as in \texttt{VVle}; no independence theorem is needed to
replace this quantifier.

For a CM field $E$, complex conjugation acts on $\mathscr S_p(E,j)$.
Write $\ell_p^+(E,j)$ and $\ell_p^-(E,j)$ for the lengths of the
eigenspaces with eigenvalues $(-1)^j$ and $(-1)^{j+1}$, respectively.
Thus the labels $+$ and $-$ include the twist parity; they are not
always the eigenvalues $+1$ and $-1$ of the untwisted action.

For an integer $t\ge1$, use the following notation for the literal source
conditions.  Each condition includes finiteness of every length being
summed, for $1\le j<t$:
\begin{align*}
 \mathcal B_2(E,t;B_2)&:
       \sum_{1\le j<t}\ell_2(E,j)\le B_2,\\
 \mathcal B_\pm(E,t;B_\pm)&:
       \sum_{p\in S}\log p\sum_{1\le j<t}\ell_p^\pm(E,j)\le B_\pm
       \quad\text{for every finite set $S$ of odd primes},\\
 \mathcal B(E,t;B)&:
       \sum_{p\in S}\log p\sum_{1\le j<t}\ell_p(E,j)\le B
       \quad\text{for every finite set $S$ of primes}.
\end{align*}
These are \texttt{BoundedHeightTwoPlaceCondition},
\texttt{VVplusLe}/\texttt{VVminusLe}, and \texttt{VVle}, respectively.
In particular, $\mathcal B_2$ has no factor $\log2$ in its definition.

\begin{proposition}[VW1: exact cohomological assembly]\label{prop:vw-parts}
For every CM field $E$, every integer $t\ge1$, and every specified
system of the preceding data, the three conditions
$\mathcal B_2(E,t;B_2)$, $\mathcal B_-(E,t;B_-)$ and
$\mathcal B_+(E,t;B_+)$ imply
\[
                 \mathcal B(E,t;B_2\log2+B_-+B_+).
                 \tag{VW.1}
\]
This implication preserves all the universal local-data quantifiers.
\end{proposition}

\begin{proof}
Here are the cohomological details needed for the splitting.  Lift an
automorphism $\tau$ of $E$ to an automorphism $\sigma$ of its algebraic
closure, and let $\chi_p$ be the cyclotomic character.  On a continuous
$1$-cocycle $z$ with values in the twist $-j$, the induced map is
\[
 (T_\sigma z)(g)=\chi_p(\sigma)^{-j}z(\sigma^{-1}g\sigma).
 \tag{VW.2}
\]
Conjugation preserves the subgroup generated by inertia outside $p$,
since it permutes places above each rational prime, so this formula
descends to $G_{E,S_p}$.  The cocycle identity shows that it takes
cocycles to cocycles and coboundaries to coboundaries.  If $\sigma$ fixes
$E$, the same identity gives
\[
 T_\sigma z(g)-z(g)=(\rho(g)-1)z(\sigma),
 \qquad \rho=\chi_p^{-j},
\]
which is a coboundary.  Consequently the action on cohomology is
independent of the lift and is multiplicative in $\tau$.  Local
restriction commutes with (VW.2), with the places permuted.  It therefore
preserves the kernel $\mathscr S_p(E,j)$.  For complex conjugation $c$,
this proves $c^2=1$ on that kernel, for the original local data.

When $p$ is odd, $2$ is a unit in $\mathbb Z_p$.  The two endomorphisms
\[
 P_+=\frac{1+(-1)^jc}{2},\qquad
 P_-=\frac{1-(-1)^jc}{2}
\]
are complementary idempotents.  Thus the module itself is the direct
sum of the two eigenspaces, and additivity of module length gives
\[
                \ell_p(E,j)=\ell_p^+(E,j)+\ell_p^-(E,j).
                \tag{VW.3}
\]
This equality holds for extended lengths as well.  The assumed finiteness
of the two summands makes every full odd-primary length finite.
The two-primary finiteness is supplied separately by $\mathcal B_2$.

Now split any finite prime set into its possible member $2$ and its odd
members.  Apply (VW.3) term by term on the odd part, and bound the
$2$-term by $B_2\log2$.  This proves (VW.1), including the finiteness
conjunct in \texttt{VVle}.  The argument applies separately to every
algebraic closure and every local-data system, so none is changed or
selected anew.  No duality or automorphy theorem is used in this assembly.
\end{proof}

\subsection{The actual field and the uniform numerical budget}

Fix a primitive target character $\psi$ of conductor $M>0$.
The condition \texttt{IsChiField}$(\psi,e,q,d;E)$ means that $E$ embeds
in a cyclotomic field and that its fixing automorphisms are exactly those
whose exponents $a$ satisfy
\[
 a\equiv\pm1\pmod{2^{e+2}},\quad \psi(a)=1,\quad
 (a/q)=(a/d)=1.
 \tag{VW.4}
\]
For prime $q\equiv3\pmod8$ and prime $d\equiv7\pmod{16}$ coprime to
$M$, this identifies its image with the compositum
\[
 E_e=\mathbb Q(\zeta_{2^{e+2}})^+
          \mathbb Q(\zeta_M)^{\ker\psi}(\sqrt{-q},\sqrt{-d}).
 \tag{VW.5}
\]
Let $N(E)$ be the least positive $N$ admitting a rational embedding
$E\hookrightarrow\mathbb Q(\zeta_N)$; this is the literal
\texttt{condF E}.  The existence of such $N$ is part of (VW.4).

For $0<\epsilon<1/52$, call $e(x)$ a page function if
\[
 x^{2+8\epsilon}\le2^{e(x)}\le2x^{2+8\epsilon}
                  \quad(x\ge1).
 \tag{VW.6}
\]
The explicit choice
$e(x)=\lceil(2+8\epsilon)\log_2x\rceil$ works for $x\ge1$.
A \emph{window pair at $w$} means the full \texttt{WinPair} condition:
$q,d$ are primes with the congruences just stated,
\begin{gather*}
 \tfrac12w^{\epsilon/5}<\log q\le w^{\epsilon/5},\qquad
 \tfrac12w^{\epsilon/5}(\log w)^2<\log d
                   \le w^{\epsilon/5}(\log w)^2,\\
 q\mid d+1,\qquad
 p\mid d+1\quad\text{for every odd prime }p\le w^{\epsilon/5}.
 \tag{VW.7}
\end{gather*}
The divisibility clauses are retained even though the estimates immediately
below only need the displayed size bounds.

\begin{proposition}[VW2: same-field and uniform numerical lemmas]
\label{prop:vw-budget}
The following assertions hold, with constants depending only on the
parameters declared fixed.
\begin{enumerate}
\item Every field satisfying (VW.4), with $q$ prime and $q\equiv3\pmod4$,
is a CM field and satisfies
\[
 1\le N(E)\le2^{e+2}Mqd,\qquad [E:\mathbb Q]\le N(E).
 \tag{VW.8}
\]
For the host assumptions in (VW.5) and $e\ge2$, its conductor is exactly
\[
                 N(E)=\operatorname{lcm}(2^{e+2},M)qd.
                 \tag{VW.9}
\]
\item Fix $C_t\ge1$.  A page function and $t(x)=\lfloor x^{C_t}\rfloor$
satisfy, for every $x\ge2$,
\[
 e(x)\ge2,\qquad 2^{e(x)}\le x^4,\qquad
             2\le t(x)\le x^{C_t}.
 \tag{VW.10}
\]
\item Uniformly over every window pair at $w$, for sufficiently large
$w$ one has $\log(qd)\le w^{\epsilon/4}$.  For fixed $J'>0,C_t,P$,
all window pairs at $w=X^{1/J'}$, once $X$ is sufficiently large, satisfy
\[
 q,d>2X^{C_t}+1,\qquad q,d>P.
 \tag{VW.11}
\]
\item Uniformly over $E,e,q,d,t$ satisfying (VW.6), (VW.8),
$1\le t\le x^{C_t}$, and $\log(qd)\le x^{\epsilon/4}$, one has
\[
                    \log(N(E)t)\le2x^{\epsilon/4}
                    \tag{VW.12}
\]
for all sufficiently large $x$.
\end{enumerate}
\end{proposition}

\begin{proof}
The fixing condition in (VW.4) and finite Galois correspondence identify
the image of $E$ with the compositum (VW.5), inside the \emph{same}
cyclotomic ambient field.  For the assertion of being CM, it suffices to
see $\sqrt{-q}$ in this image.  With $\zeta_q$ a primitive $q$th root,
the quadratic Gauss sum
$g=\sum_{a\bmod q}(a/q)\zeta_q^a$ satisfies
$\tau_b(g)=(b/q)g$.  Moreover $|g|^2=q$ by additive-character
orthogonality: writing $a=tb$ in the double sum gives $q-1$ for $t=1$
and $-1$ otherwise, and the sum of the nontrivial quadratic character
is zero.  Since $\overline g=(-1/q)g=-g$, it follows that $g^2=-q$.
Every automorphism fixing $E$ fixes $g$ by (VW.4), so $g$ lies in the
image of $E$.  Therefore $E$ has no real embedding.  It is an abelian
subextension of a cyclotomic extension.  Complex conjugation is then a
central nontrivial involution, its fixed field is totally real, and
$E$ is quadratic over that fixed field.  This proves the CM assertion.

Put $L=\operatorname{lcm}(2^{e+2},M,q,d)$.  Every automorphism trivial
modulo $L$ satisfies (VW.4); Galois correspondence gives an embedding
$E\hookrightarrow\mathbb Q(\zeta_L)$.  Hence $N(E)\le L\le2^{e+2}Mqd$.
At the minimal conductor, dimension of the embedding gives
$[E:\mathbb Q]\le\varphi(N(E))\le N(E)$, proving (VW.8).

For (VW.9), the character group of the compositum contains $\psi$, the
primitive quadratic characters of conductors $q,d$, and a primitive
even character of conductor $2^{e+2}$ from the real cyclotomic factor.
Such an even character is obtained by sending $5$ to a primitive
$2^e$th root and $-1$ to $1$ in
$(\mathbb Z/2^{e+2}\mathbb Z)^\times
 =\langle-1\rangle\times\langle5\rangle$.
The identity $v_2(5^{2^r}-1)=r+2$, proved by induction by squaring,
shows that this character does not factor modulo $2^{e+1}$.
Every cyclotomic modulus for the compositum must therefore be divisible
by $M,q,d,2^{e+2}$.  Conversely their least common multiple is a
cyclotomic modulus by the preceding paragraph.  Since $q,d$ are distinct
and prime to $2M$, this proves (VW.9).  This argument is invariant under
the given rational embedding of $E$ and proves the equality for the
original field, not just for a separately chosen model.

For (VW.10), $2+8\epsilon<3$.  At $x\ge2$, (VW.6) gives
$2^{e(x)}\ge x^2\ge4$ and
$2^{e(x)}\le2x^3\le x^4$.
The floor inequalities follow from $C_t\ge1$.

The upper pair bounds give
$\log(qd)\le w^{\epsilon/5}(1+\log^2w)$.
Since $1+\log^2w\le w^{\epsilon/20}$ eventually, the first estimate
in part 3 follows uniformly in the pair.  Set
$\delta=\epsilon/(5J')>0$.  For $w=X^{1/J'}$ and $\log w\ge1$,
both lower pair bounds give $q,d>\exp(X^\delta/2)$.
That expression eventually exceeds $2X^{C_t}+1$ and any fixed $P$.
This proves (VW.11) uniformly over all allowed pairs.

Finally (VW.6) and (VW.8) imply
\[
 \log(N(E)t)\le\log(8M)+(2+8\epsilon+C_t)\log x+\log(qd).
\]
For large $x$, the first two terms are at most $x^{\epsilon/4}$,
independently of $E,e,q,d,t$; the assumed bound on the last term proves
(VW.12).  All ``sufficiently large'' thresholds in this proof involve
only fixed real parameters and $M$, as required.
\end{proof}

\paragraph{The exact odd-minus expression.}
For $o,n,N,t\in\mathbb N$ and $C_b\in\mathbb R$, write
\[
 \mathfrak M(C_b;o,n,N,t)
 =3nt^2\log(tN)+C_b\sum_{\substack{\ell\mid o\\\ell\text{ odd prime}}}
          \ell n^2t^2\log(t\ell N)^2.
 \tag{VW.13}
\]
This is the source's \texttt{minusBound}.  If $1\le N,t$, $n\le N$,
and $\log(Nt)\le2W$, then
\[
 \mathfrak M(C_b;o,n,N,t)\le K\exp(8W),\qquad
 K=3+|C_b|\sum_{\substack{\ell\mid o\\\ell\text{ odd prime}}}\ell^3.
 \tag{VW.14}
\]
Indeed put $u=Nt\ge1$.  The first term is at most $3u^4$ because
$nt\le u$, $t\le u$ and $\log u\le u$.
For the summand at $\ell$, use
$\log(t\ell N)\le\ell u$ and $nt\le u$, obtaining at most
$\ell^3u^4$.  Replace $C_b$ by $|C_b|$ and use $u^4\le e^{8W}$.
The crucial point is that $o$ and $C_b$, hence $K$, are fixed before
$e,q,d,x$ vary.

\paragraph{A uniform final budget.}
Fix real $C$ and $K\ge0$.  If $u\ge1$ and
$\log u\le2x^{\epsilon/4}$, then eventually
\[
 e^{x^{4\epsilon/5}}\log2
       +Ke^{8x^{\epsilon/4}}+u^C+Ce^{x^{4\epsilon/5}}
                     \le e^{x^\epsilon/2}.
 \tag{VW.15}
\]
For a fully explicit threshold set
$A=\log2+|C|+K+1$, $D=9+2|C|$, and require
\[
 x\ge\max\{1,[2(D+\log A+1)]^{5/\epsilon}\}.
\]
Then $b=x^{4\epsilon/5}\ge1$ and $x^{\epsilon/4}\le b$.
Each exponential on the left of (VW.15), and also $u^C$, is bounded
by $e^{Db}$ after replacing coefficients by their absolute values.
The sum is at most $Ae^{Db}$.  The threshold gives
$x^\epsilon/2\ge b(D+\log A+1)\ge Db+\log A$, proving (VW.15).

\subsection{Two precise input interfaces and their closed consequence}

We now state the input interfaces in ordinary mathematical notation.
Their names refer to literal source predicates, rather than additional
assumptions on the final theorem.

Fix $C_t\ge1$.  A pair of functions $(e,t)$ is \emph{admitted with
exponent $A_1$} when, for all $x\ge2$,
\[
 e(x)\ge2,\quad 2^{e(x)}\le x^{A_1},\quad
                   2\le t(x)\le x^{C_t}.
\]
Let $M_\psi=\mathbb Q(\zeta_M)^{\ker\psi}$ and
$o_\psi=[M_\psi:\mathbb Q]$; this equals the order of $\psi$.

\begin{description}[style=nextline,leftmargin=0pt,labelwidth=0pt,labelsep=0pt]
\item[Joint window input $\mathsf J(C_t)$.]
For each primitive $\psi$, each $0<\epsilon<1/52$, each real $A_1$,
and each admitted $(e,t)$, there are constants
$J>0,C,P,X_J$ such that the following holds for every $X\ge X_J$.
Whenever $q,d$ are primes, $q\equiv3\pmod8$, $d\equiv7\pmod{16}$,
$\gcd(qd,M)=1$,
\[
 q,d>2X^{C_t}+1,\quad q,d>P,\quad
                  \log(qd)\le X^{\epsilon/(4J)},
\]
there is $Y$ with $X^{1/J}\le Y$ and $Y^2\le X$ such that,
for every $x\in[Y,Y^2]$ and every field $E$ satisfying
\texttt{IsChiField}$(\psi,e(x),q,d;E)$,
\[
 \begin{split}
 \mathcal B_2(E,t(x);e^{x^{4\epsilon/5}}),\
 \mathcal B_+(E,t(x);(N(E)t(x))^C+Ce^{x^{4\epsilon/5}})
 \end{split}
 \tag{VW.16}
\]
hold with their full local-data quantifiers.  This is
\texttt{JointWin}$(C_t)$: its \texttt{BaAt} and \texttt{HVAt} concern
the same $E$, using
$\texttt{IsCompField}(\texttt{AbF.ofChar}\,\psi)\leftrightarrow
\texttt{IsChiField}\,\psi$.

\item[Odd-minus input $\mathsf M(C_t)$.]
The literal \texttt{ThmBc}$(C_t)$ supplies a constant $C_b$,
uniform over its abelian base fields and admitted functions, such that
for the same fields and parameters
\[
 \mathcal B_-(E,t(x);
       \mathfrak M(C_b;o_\psi,[E:\mathbb Q],N(E),t(x)))
 \tag{VW.17}
\]
holds whenever $x\ge2$ and the prime congruence and coprimality
conditions hold.  Only its specialization to the fixed base $M_\psi$
is needed here; the universal source statement is stronger.
\end{description}

\begin{theorem}[VW3: complete $\mathsf J+\mathsf M$ to $\mathtt{VWinAt}$ chain]
\label{thm:vw-closed}
Let $C_t\ge1$.  Assuming the two interfaces $\mathsf J(C_t)$ and
$\mathsf M(C_t)$ just stated, the literal \texttt{VWinAt}$(\psi,C_t)$
holds for every primitive $\psi$.
More precisely, for every $0<\epsilon<1/52$ and every page function
$e$, there exist $J'\ge1$ and $X_0$ such that for every $X\ge X_0$
and \emph{every} window pair $q,d$ at $w=X^{1/J'}$, there is $Y$ with
\[
 w\le Y,\qquad Y^2\le X,
\]
and for all $x\in[Y,Y^2]$, every field $E$ satisfying
\texttt{IsChiField}$(\psi,e(x),q,d;E)$ satisfies
\[
       \mathcal B(E,\lfloor x^{C_t}\rfloor;e^{x^\epsilon/2}).
       \tag{VW.18}
\]
This includes finiteness of every individual length and bounds every
finite prime subsum, for every algebraic closure and every local-data
system.  The same pair $q,d$ is used throughout the entire square window.
\end{theorem}

\begin{proof}
Fix $\psi,\epsilon,e$.  Put $t(x)=\lfloor x^{C_t}\rfloor$.
By (VW.10), $(e,t)$ is admitted with $A_1=4$.
Apply $\mathsf J(C_t)$ to obtain $J,C,P,X_J$, and let $C_b$ be
provided by $\mathsf M(C_t)$.  Fix
\[
 J'=\max(J,1),\qquad P'=\max(P,M),\qquad
 K=3+|C_b|\sum_{\substack{\ell\mid o_\psi\\\ell\text{ odd prime}}}\ell^3.
\]
These constants are all fixed before $X,q,d,Y,x,E$ are chosen.
Choose $L\ge2$ above the thresholds for the first pair estimate in
VW2, (VW.12), and (VW.15).  Increase $X_0$ to exceed
$X_J,L^{J'}$ and the threshold in (VW.11) for $J',C_t,P'$.

Now let $X\ge X_0$ and let $q,d$ be \emph{any} window pair at
$w=X^{1/J'}$.  Then $w\ge L$ and
\[
 \log(qd)\le w^{\epsilon/4}
       =X^{\epsilon/(4J')}\le X^{\epsilon/(4J)}.
\]
By (VW.11), $q,d>2X^{C_t}+1$ and $q,d>P'$.  Since $M>0$ and these
primes are greater than $M$, they are coprime to $M$.  Thus every
hypothesis of $\mathsf J(C_t)$ holds for this same pair.  It supplies
$Y$ with $X^{1/J}\le Y$, $Y^2\le X$, and (VW.16) on $[Y,Y^2]$.
Because $J'\ge J>0$ and $X\ge1$,
$w=X^{1/J'}\le X^{1/J}\le Y$.

Fix any $x\in[Y,Y^2]$ and any field $E$ in (VW.18).
VW2 proves directly from its fixing condition that this $E$ is CM.
It also gives $1\le N(E)$ and $[E:\mathbb Q]\le N(E)$.
Moreover
\[
 \log(qd)\le w^{\epsilon/4}\le x^{\epsilon/4},\qquad x\ge L.
\]
Hence $u=N(E)t(x)\ge1$ satisfies (VW.12).
The odd-minus input (VW.17), followed by (VW.14) with
$W=x^{\epsilon/4}$, gives
$\mathcal B_-(E,t(x);Ke^{8x^{\epsilon/4}})$.
Use (VW.16) for the other two parts and combine them by VW1.
The resulting upper bound is exactly the left side of (VW.15), so
(VW.18) follows.  No stage chooses a different field, pair, algebraic
closure or local-data system.  Since all thresholds were chosen before
the arbitrary pair, the universal \texttt{WinPair} quantifier is preserved.
\end{proof}

\subsection{The minimal supplier boundary and the fixed-pair endpoint}

The new theorem removes \texttt{VWinAt} as a separate unresolved
\emph{assembly} once $\mathsf J$ and $\mathsf M$ are available.  It does
not construct the cohomological estimates (VW.16)--(VW.17).
For continuation, the source provides a more structural choice of input
boundary:
\[
 \boxed{\quad
 \mathtt{ThmBab}(K,C_t)\ +\ \mathtt{PointHV}\ +\ \mathtt{ThmBcFree}
 \quad}\ .
 \tag{VW.19}
\]
Here \texttt{ThmBab} supplies, on one common square window, the
$2$-primary bound and the finitely many specified odd-primary layer
bounds; \texttt{PointHV} converts those layer bounds into the
odd-plus estimate.  \texttt{ThmBcFree} is precisely the odd-minus
bound (VW.13) for all integer $e,t\ge2$, before admitting scale functions.
These are substantive inputs, not consequences of the arguments above.

\begin{proposition}[VW4: structural supplier and final fixed-pair closure]
\label{prop:vw-fixed-pair}
For $C_t\ge1$, the literal interfaces in (VW.19) imply
$\mathsf J(C_t)$ and $\mathsf M(C_t)$, hence VW3.
For $C_t=401$, combining this with QP22 gives, for every primitive
$\psi$, every $0<\epsilon<1/52$, and every $c\ge1$, a $J_0\ge1$
such that arbitrarily large $Z$ admit a \emph{single} $Y_0,q,d$ with
\[
 \begin{gathered}
 x^{1/J_0}\le Y_0\le x,\qquad
       \mathtt{PairGood}(M,\epsilon,Y_0,x;q,d),\\
 x^{2+8\epsilon}\le c2^{e(x/k)}\le2x^{2+8\epsilon},\qquad
 \mathcal B(E,\lfloor x^{400}\rfloor;e^{x^\epsilon/2})
 \end{gathered}
 \tag{VW.20}
\]
for every $x\in[Z,Z^{3/2}]$ and every
$E$ satisfying \texttt{IsChiField}$(\psi,e(x/k),q,d;E)$, where
$k=c^{1/(2+8\epsilon)}$.  All clauses of \texttt{PairGood}, including
its zero-free boxes and all primitive-inducer clauses, are retained.
When QP22 is used, its named classical inputs CN, BD and HD remain the
analytic leaves of this fixed-pair chain.
\end{proposition}

\begin{proof}
We first check the structural supplier, rather than merely quote an
arrow between predicate names.  For fixed $\psi$, put
$B'=\mathtt{Bchi}\,\psi$ and
$L_\ell=\mathbb Q(\zeta_{\ell M})$ for $\ell\in B'$.
Explicitly, $B'$ consists of $3$ and the odd prime divisors of
$o_\psi M p_Kh_*p_1$.  Here $p_K$ is the least prime at least $7$,
congruent to $3$ modulo $4$ and coprime to $M$; $h_*$ is the class
number of $\mathbb Q(\sqrt{-p_K})$; and $p_1$ is the least prime
congruent to $1$ modulo the odd part of $o_\psi$ and coprime to
$2Mp_K$.  These fixed constants are the ones in GH7 and
\texttt{StmtFields}.
Every $\ell\in B'$ is an odd prime and $L_\ell$ contains $\mu_\ell$.
These are the exact finite set and fields in the source.
The record $K$ specifies nonnegative integers $R_0,R_\ell$ and $C_K>0$.
Put
\[
 R=R_0(M_\psi)+\sum_{\ell\in B'}R_\ell(L_\ell),\qquad
 J_K=2^{R+2}(2C_KC_t/\epsilon)^R,
 \qquad j_\ell=\lfloor\log_\ell R_\ell(L_\ell)\rfloor+3,
 \tag{VW.21}
\]
where the source's natural-number logarithm is $0$ when its argument is
$0$.  The integer $j_\ell$ is therefore defined even when
$R_\ell(L_\ell)=0$.

For the literal \texttt{PointHV} hypothesis, first fix $\psi$ and the
function $\ell\mapsto j_\ell$.  It supplies integers $C,r_1,r_2$ with
$r_1,r_2$ prime.  For every host
$h=(\psi,e,q,d)$ with $e\ge2$ and the prime congruence/coprimality
conditions, it asserts the following.  If $q,d$ exceed $2t+1$,
$\psi.N,\operatorname{ord}(\psi),p_K,h_*,p_1,r_1,r_2$, and if the
specified $\ell$-primary length sum on every prescribed cyclotomic
layer of degree $\ell^{j_\ell}$ over $L_\ell$ is bounded by the same
$B\ge0$ for each $\ell\in B'$, then
\[
       \mathcal B_+(E,t;(N_h t)^C+CB)
       \quad\text{for every corresponding }E,
 \qquad N_h=\operatorname{lcm}(2^{e+2},M)qd.
 \tag{VW.22}
\]
Here the layer quantifier includes every field model $L_0$ of $L_\ell$
and every finite extension $L_1/L_0$ generated by roots of unity of
$\ell$-power order with $[L_1:L_0]=\ell^{j_\ell}$, and retains the
universal algebraic-closure and local-data quantifiers of
$\mathcal B_\ell$ (the unweighted $\ell$-primary length sum).
None of $p_K,h_*,p_1,r_1,r_2$ depends on the varying host.

Apply \texttt{ThmBab} to the abelian base $M_\psi$, this $B'$, and these
$L_\ell$.  It supplies a square window with exponent $J_K$, the
$2$-primary bound in (VW.16), and all the layer bounds just described,
with $B=e^{x^{4\epsilon/5}}$, on the \emph{same} window and for the
\emph{same} pair.  Choose the fixed lower prime cutoff $P$ above the
seven reference numbers in the preceding paragraph.  Increase the
threshold for $X$ so $X^{1/J_K}\ge2$ and the thresholds of both
inputs hold.  Because $t(x)\le x^{C_t}\le X^{C_t}$ on the window,
$q,d>2X^{C_t}+1$ ensures the host inequalities $q,d>2t(x)+1$.
The remaining host inequalities follow from $q,d>P$.
Thus (VW.22) applies at every $x$ in this very window.  The exact
conductor identity (VW.9) gives $N_h=N(E)$ for each corresponding
original field $E$.  Consequently (VW.22) is the odd-plus bound in
(VW.16), with no loss in exponent or conductor.  This proves
$\mathsf J(C_t)$, including its dependence order.  The same field is
used on both sides of the equality; neither a cohomology group nor a
coefficient system has been transported to a newly chosen host.
The free odd-minus input is applied to the integer values $e(x),t(x)$;
(VW.10), or the definition of an admitted pair, proves their lower
bounds.  This gives $\mathsf M(C_t)$ directly.

For (VW.20), use VW3 with $C_t=401$ to fix $J_0$ and its threshold
for the page function.  Apply QP22 with polynomial-window exponent
$J_0$.  Given a lower bound for $Z$, choose $X$ large enough that
$Y_0=X^{1/J_0}$ exceeds that bound, the QP22 selection threshold,
$k^3$ and $k^{400}$.  QP22 gives one pair $q,d$ satisfying the full
\texttt{PairGood} condition on $[Y_0,X]$.  Its \texttt{WinPair} part
allows VW3 to supply $W\ge Y_0$, $W^2\le X$, for this same pair.
Set $Z=kW$.

For $Z\le x\le Z^{3/2}$, the condition $W\ge k^3$ gives
$x\le W^2$, and with $u=x/k$ one has $W\le u\le x\le W^2$.
Restrict the good pair's upper endpoint from $X$ to $x$.
The inequalities $x\le X$ and $Y_0\le W\le x$ give
$x^{1/J_0}\le Y_0\le x$.
Multiplying (VW.6) at $u$ by $c=k^{2+8\epsilon}$ proves the degree
bounds in (VW.20).  Since $u\ge W\ge k^{400}$,
\[
                   x^{400}=k^{400}u^{400}\le u^{401}.
\]
VW3 at $u$ therefore bounds a longer twist range than needed at $x$;
all lengths and prime weights are nonnegative, so the bound persists
on decreasing the cutoff to $\lfloor x^{400}\rfloor$.
Finally $u^\epsilon\le x^\epsilon$ increases the allowed upper bound.
This proves (VW.20) with $Y_0,q,d$ fixed throughout the window.
For ordinary RH, $c=1$ permits retaining the whole square window
$[W,W^2]$, as in OH3.
\end{proof}

\paragraph{The closed chain and its open leaves.}
The chain established in this section is
\[
 \begin{gathered}
 (\mathtt{ThmBab},\mathtt{PointHV},\mathtt{ThmBcFree})
 \Longrightarrow(\mathtt{JointWin},\mathtt{ThmBc})
 \Longrightarrow\mathtt{VWinAt},\\
 \mathtt{VWinAt}+\mathrm{QP22}
 \Longrightarrow\text{one fixed good pair and cohomological window}.
 \end{gathered}
\]
It leaves exactly the deep supplier interfaces on the left (or the
coarser two-interface boundary in VW3).  The source proves them using
its Iwasawa, cohomological counting, congruence and automorphic machinery;
that work has not been reconstructed by this assembly.  In particular,
\texttt{LinkN}/\texttt{LinkV}, the required counting statements and the
actual odd-minus estimate cannot be retired merely because VW3 is now
proved.  Likewise (VW.20) is an input to the later arithmetic geometry;
it does not itself supply the lattice-gain inequality needed to bound
$H_\psi$.


\section{Closing the character-to-window dependency chain}
\label{sec:window-bridge}

\status{WB1--WB8 give one complete implication, from the explicitly
stated arithmetic gap and size data below to the fixed positive
character window and then to RH and Dirichlet GRH. The finite Fourier
identity, ramified-prime terms, subgroup averaging, moving cutoffs,
primitive-character reading, and uniform quantifier order are proved.
The arithmetic gap and the existence of one family satisfying its
premises remain inputs. This section closes the character/analytic
bridge left open in SC2; it does not assert those arithmetic inputs.}

The point of this consolidation is to exhibit the whole chain in one
place. Its leaves are concrete inequalities on one set of data, rather
than a collection of mutually unrelated existence statements. The
source route is \texttt{ChBalance}, \texttt{ChCore}, \texttt{GmAEq},
\texttt{SigALe}, \texttt{CoreGm}, and \texttt{TopWin}.

% BEGIN WB1
\subsection{WB1. The actual characters, including ramification}

Let $A$ be a finite abelian group of order $n$, and write
\[
 G=(\mathbb Z/2\mathbb Z)^2\times A,\quad
 \Phi=\{(0,b,h)\},\quad f=|\Phi|=2n,\quad
 c_0=(1,0,0),\quad c_\infty=(1,1,0).
\]
Use multiplicative notation for characters of the additive group $G$;
the trivial character, written $1$ here, is written $0$ in the Lean
\texttt{AddChar} notation. Put $\chi_d(a,b,h)=(-1)^a$.
An arithmetic datum assigns an inertia subgroup $I_p\le G$ and a
Frobenius class $F_p\in G$ to every prime. Define, for $j\ge1$,
\[
 v_\xi(p^j)=\frac1{|I_p|}\sum_{h\in I_p}\xi(jF_p+h),
 \qquad
 v_\xi(r)=\prod_{p^j\parallel r}v_\xi(p^j)\quad(r\ge1).
 \tag{WB.1}
\]
Thus $|v_\xi(r)|\le1$ and $v_1(r)=1$. Finite-character
orthogonality gives the exact ramified formula
\[
 \frac1f\sum_{\xi\in\widehat G}\overline{\xi(g)}v_\xi(p^j)
 =\frac{2}{|I_p|}\mathbf1_{g-jF_p\in I_p}.
 \tag{WB.2}
\]
Indeed, expand (WB.1) and interchange the two finite sums. The inner
character sum is $|G|=2f$ when $jF_p+h=g$, and zero otherwise.
There is at most one such $h\in I_p$.

Here is why these are the values of the \emph{primitive} Dirichlet
characters on an actual host. Suppose its positive modulus is $N_0$
and its specified reciprocity map
\[
 r:(\mathbb Z/N_0\mathbb Z)^\times\longrightarrow G
\]
is onto. At a prime $p\mid N_0$, use CRT to split the unit group
into its $p$-primary factor and its prime-to-$p$ factor. Let $I_p$
be the image of the first factor and let $F_p$ be the image of a unit
which is $p$ modulo the prime-to-$p$ part. Its value at the
$p$-primary part is immaterial modulo $I_p$. For $p\nmid N_0$,
take $I_p=0$ and $F_p=r(p)$.

Let $\xi^*$ be the primitive inducer of $\xi\circ r$. If
$\xi|_{I_p}$ is nontrivial, character averaging makes (WB.1) zero;
the primitive conductor is then divisible by $p$, and
$\xi^*(p)=0$ as well. If $\xi|_{I_p}=1$, the induced character
factors through the prime-to-$p$ part. Its primitive conductor is
prime to $p$, and its value at $p^j$ is exactly $\xi(jF_p)$.
This proves
\[
 v_\xi(r)=\xi^*(r)\qquad(r\ge1),
 \tag{WB.3}
\]
also at ramified primes, not only on integers prime to $N_0$.
Uniqueness of primitive induction and surjectivity of $r$ imply
\[
 \xi^*=\zeta^*\Longrightarrow\xi=\zeta,
 \qquad \xi^*=1\Longleftrightarrow\xi=1.
 \tag{WB.4}
\]
For the first implication, induce both primitive characters back to
$N_0$, and then use surjectivity. In particular no growing host
introduces repeated copies of one primitive target.

Extend the already centered principal ramp to every $y>0$ by the
same finite-sum formula
\[
 G(y)=\sum_{r\le y}\Lambda(r)(y/r-1)-y\log y+(1+\gamma)y.
\]
For $0<y<1$ the sum is empty. This extension matters in the exact
balance identity, where $a/w$ need not be at least $1$.
Fix numbers $\ell_\xi\in\mathbb C$ for $\xi\ne1$ and set, for
$y>0$,
\[
 \mathcal G_1(y)=G(y),\qquad
 \mathcal G_\xi(y)=\sum_{r\le y}\Lambda(r)v_\xi(r)(y/r-1)
                         +y\ell_\xi\quad(\xi\ne1).
 \tag{WB.5}
\]
On the original host require the actual equality
$\ell_\xi=L'/L(1,\xi^*)$. Equations (WB.3)--(WB.5) then identify
$\mathcal G_\xi$ with the already centered primitive ramp
$G_{\xi^*}$. This is the literal content of the
\texttt{HostAddCharCompatible}/\texttt{prim\_\allowbreak read} join.
The arbitrary value $\ell_1$ is never used.
% END WB1

% BEGIN WB2
\subsection{WB2. Exact balance, with the exceptional prime retained}

Assume that a finite prime set $S$ and a prime $d\notin S$ satisfy
\[
 I_p=0\ (p\notin S\cup\{d\}),\quad
 I_p\subset\Phi,\ F_p\in\Phi\ (p\in S),\quad
 I_d=\{0,c_0\}.
 \tag{WB.6}
\]
These are the substantive fields of \texttt{HostFactsS}. The source
also includes $2,q\in S$ to identify its inert-prime indexing.
For $t\in A$, the map
\[
 g\longmapsto\rho_t(g):=c_\infty+g+(0,0,t)
\]
is a bijection $G\setminus\Phi\to\Phi$. Given $a>0$ and
$w:\Phi\to\mathbb R_{>0}$, choose an integer $N$ at least every
$a/w(\rho)$. Write $(z)_+=\max(z,0)$ and define the actual asked
prime sum
\begin{align*}
 \mathcal A(a,t,w)={}&
 2\sum_{\substack{p\notin S\cup\{d\}\\F_p\notin\Phi}}
    \log p\sum_{\substack{j\ge1\\jF_p\notin\Phi}}
       (a/p^j-w(\rho_t(jF_p)))_+\\
 &+\log d\sum_{j\ge1}
       \sum_{\substack{g\notin\Phi\\g-jF_d\in I_d}}
          (a/d^j-w(\rho_t(g)))_+ .
 \tag{WB.7}
\end{align*}
All sums are finite: terms with $p^j>N$ vanish. This is the source's
\texttt{ASK}; when $d>N$ its last summand vanishes, rather than
requiring a separate assumption on $d$.

Define
\begin{align*}
 \mathcal M(a,w)&=a\left(\log a-1-\gamma+\ell_{\chi_d}
                       -\frac1f\sum_{\rho\in\Phi}\log w(\rho)\right),\\
 \mathcal D(a,t,w)&=\frac1f\sum_{g\notin\Phi}
       \sum_{\xi\in\widehat G}\overline{\xi(g)}
       w(\rho_t(g))\mathcal G_\xi(a/w(\rho_t(g))).
\end{align*}
Then the precise complex identity is
\[
                 \boxed{\ \mathcal A-\mathcal M=\mathcal D.\ }
 \tag{WB.8}
\]
To prove it, apply (WB.2) to each prime power. At an unramified prime
the coefficient is $2$ at $g=jF_p$ and zero elsewhere. At $p\in S$
the coset $jF_p+I_p$ lies in $\Phi$, so it contributes nothing.
At $d$ the coefficient is $1$ on its inertia coset; that coset has
exactly one member outside $\Phi$, because $c_0\notin\Phi$.
Consequently (WB.7) equals
\[
 \frac1f\sum_{g\notin\Phi}\sum_\xi\overline{\xi(g)}
       \sum_{r\ge1}\Lambda(r)v_\xi(r)(a/r-w(\rho_t(g)))_+ .
\]
For $W>0$, the innermost sum is $W\mathcal G_\xi(a/W)-a\ell_\xi$
if $\xi\ne1$. For $\xi=1$ it is
\[
 WG(a/W)+a\log a-a\log W-(1+\gamma)a.
\]
Finally orthogonality on the index-two subgroup $\Phi$ gives
\[
 \sum_{g\notin\Phi}\overline{\xi(g)}=
 \begin{cases} f,&\xi=1,\\-f,&\xi=\chi_d,\\0,&\text{otherwise}.
 \end{cases}
\]
The $-a\ell_{\chi_d}$ correction is therefore multiplied by $-f$:
this proves the \emph{plus} sign in $a\ell_{\chi_d}$ in
$\mathcal M$. The bijection $\rho_t$ gives the logarithmic sum.
This proves (WB.8) without dropping any ramified contribution.
% END WB2

% BEGIN WB3
\subsection{WB3. Uniform regularity of the centered ramps}

The following estimates use only $|v_\xi(r)|\le1$ and elementary
Chebyshev bounds; their constants are independent of the host.
First,
\[
 \Psi(y):=\sum_{r\le y}\Lambda(r)\le3y\quad(y\ge0),\qquad
 \log k-1\le\sum_{r\le k}\frac{\Lambda(r)}r\le\log k+3
 \quad(k\in\mathbb Z_{\ge1}).
 \tag{WB.9}
\]
For completeness, central binomial coefficients give
$\vartheta(y)=\sum_{p\le y}\log p\le y\log4$ by induction on
$\lfloor y\rfloor$. At $2m$ use
$\vartheta(2m)-\vartheta(m)\le\log\binom{2m}{m}\le2m\log2$;
at $2m+1$ use
$\vartheta(2m+1)-\vartheta(m+1)\le
\log\binom{2m+1}{m}\le2m\log2$.
The induction starts at $0,1,2$.
Summing prime powers gives, for $y\ge1$,
\[
 \Psi(y)\le y\log4+2\sqrt y\log y
 \le(\log4+4/e)y<3y;
\]
here $\log t/t\le1/e$ with $t=\sqrt y$. For the second bound use
\[
 \log(k!)=\sum_{r\le k}\Lambda(r)\lfloor k/r\rfloor,
 \qquad k\log k-k\le\log(k!)\le k\log k.
\]
The two inequalities $k/r-1\le\lfloor k/r\rfloor\le k/r$ and
the first part of (WB.9) give the result.

There is an absolute $C_0>0$ such that, for $y,z\ge1$,
\begin{align*}
 |G(y)-G(z)|&\le C_0|y-z|,\\
 |\mathcal G_\xi(y)-\mathcal G_\xi(z)|
 &\le |y-z|\bigl(\log\max(y,z)+C_0+|\ell_\xi|\bigr)
                 \quad(\xi\ne1),\\
 |\mathcal G_\xi(x)|&\le x(\log x+C_0+|\ell_\xi|)
                 \quad(x\ge1,\ \xi\ne1).
 \tag{WB.10}
\end{align*}
Between consecutive integers, differentiate the finite sums. For $G$
the derivative is $\sum_{r\le y}\Lambda(r)/r-\log y+\gamma$,
bounded by (WB.9); replacing $\log\lfloor y\rfloor$ by $\log y$
costs at most $\log2$. For a nontrivial character, bound the
derivative termwise using (WB.9). Each ramp is continuous at the
integers because the newly appearing term is zero there; integrate
the derivative bounds interval by interval. The final estimate uses
$\mathcal G_\xi(1)=\ell_\xi$. One may take $C_0=6+|\gamma|$.

For every character, $a>0$ and $w,w'>0$, we also have the exact
constant needed by the source:
\[
 \left|w\mathcal G_\xi(a/w)-w'\mathcal G_\xi(a/w')\right|
 \le\frac{3a}{\min(w,w')}|w-w'|.
 \tag{WB.11}
\]
For $\xi\ne1$, rewrite the left expression using
\[
 w\mathcal G_\xi(a/w)
 =\sum_{r\ge1}\Lambda(r)v_\xi(r)(a/r-w)_++a\ell_\xi.
\]
The last term cancels; the positive-part function is $1$-Lipschitz,
and (WB.9) bounds the remaining coefficient sum by
$3a/\min(w,w')$. For the trivial character suppose $w\le w'$.
The difference of the two uncentered ramp sums lies between $0$ and
$3a(w'-w)/w$. The centering subtracts
$a\log(w'/w)$, which lies between $0$ and $a(w'-w)/w$.
The difference of these two nonnegative numbers has absolute value
at most $3a(w'-w)/w$. Interchanging $w,w'$ proves (WB.11).
% END WB3

% BEGIN WB4
\subsection{WB4. Subgroup averaging produces the scalar coefficient}

Choose a subgroup $A_0\le A$, sets $T\subset A$ and
$T_2\subset A_0$, and an integer $m\ge1$. Write
\[
 n_T=|T|>0,\quad n_2=|T_2|>0,\quad
 d_2=n_2/n,\quad Q=n_Tmb,\quad
 \mathcal B=\{\xi\in\widehat G:\xi(0,0,A_0)=1\}.
\]
Let $x\ge3$, $b,b',X'>0$ and $\eta,\theta\ge0$. Supply
$a_t>0$ for $t\in T$ and positive weights
$w_{u,j}(\rho)$ for $u\in A_0$, $1\le j\le m$, $\rho\in\Phi$,
with
\begin{gather*}
 |a_t/b-x|\le1,\qquad b'X'\le xb,\\
 b\le w_{u,j}(0,1,s+u)\le b+\eta+1\quad(s\in T_2),\\
 w_{u,j}(\rho)\le2b'\quad
       \bigl(\rho\notin\{(0,1,s+u):s\in T_2\}\bigr),\\
 |\mathcal G_\xi(y)|\le y\theta
       \quad(\xi\in\widehat G,\ y\ge X'/3),\\
 2f(x+1)\theta\le1/x,\qquad
 6f(x+1)(\eta+1)\le b/x.
 \tag{WB.12}
\end{gather*}
The last two inequalities are precisely \texttt{CostWindows}.
The bound for every $y\ge X'/3$ is the supplied
\texttt{ThinScales}, not an estimate inferred from Fourier
orthogonality. Choose a common finite cutoff $N$ for these weights.

The one remaining arithmetic inequality at this point is
\[
 \sum_{t\in T}\sum_{j=1}^m
   \bigl(\mathcal A(a_t,t,w_{u,j})
                     -\Re\mathcal M(a_t,w_{u,j})\bigr)
 \le E\qquad\text{for every }u\in A_0.
 \tag{WB.13}\label{eq:wb-arithmetic-gap}
\]
It must hold for these exact data. It is not supplied by the following
argument.

Put $e_1(t)=(1,0,t)$ and define the ideal thick expression
\[
 \mathcal T_u(t,y)=\frac b f\sum_{s\in T_2}\sum_\xi
             \xi(-e_1(s+u-t))\mathcal G_\xi(y).
\]
The preimage under $\rho_t$ of $(0,1,s+u)$ is
$e_1(s+u-t)$. At such a thick point, (WB.11) bounds the change
from $w\mathcal G_\xi(a_t/w)$ to
$b\mathcal G_\xi(a_t/b)$ by $3(x+1)(\eta+1)$.
At a thin point,
\[
 \frac{a_t}{w}\ge\frac{(x-1)b}{2b'}\ge X'/3
\]
because $x\ge3$ and $b'X'\le xb$. Thus its discarded term has
absolute value at most $a_t\theta\le(x+1)b\theta$.
There are $f$ outer points and $2f$ characters; after division by
$f$, the total error is at most
\[
 2f\{3(x+1)(\eta+1)+(x+1)b\theta\}\le2b/x.
\]
Use (WB.8) and sum over $t,j$ in (WB.13) to obtain
\[
 \sum_{t,j}\Re\mathcal T_u(t,a_t/b)\le E+2Q/x.
 \tag{WB.14}
\]
Now average over $u\in A_0$. If $\xi\notin\mathcal B$ its
subgroup sum vanishes. If $\xi\in\mathcal B$, the factors from
$s$ and $u$ are one, leaving $\xi(e_1(t))$. Consequently
\[
 \frac1{|A_0|}\sum_u\mathcal T_u(t,y)
 =\frac{bd_2}{2}\sum_{\xi\in\mathcal B}
                         \xi(e_1(t))\mathcal G_\xi(y).
 \tag{WB.15}
\]
Combining (WB.14)--(WB.15) proves the complete averaging inequality
\[
 \frac{mbd_2}{2}\sum_{t\in T}
     \Re\sum_{\xi\in\mathcal B}\xi(e_1(t))
                                     \mathcal G_\xi(a_t/b)
 \le E+2Q/x.
 \tag{WB.16}
\]
In particular the factor $d_2/2$ is forced by
$f=2|A|$; there is no extra factor depending on $|A_0|$.
% END WB4

% BEGIN WB5
\subsection{WB5. The biased set gives a fixed positive signal}

Fix $\varepsilon>0$. Let $K\subset\mathcal B$ contain $1$, with $\xi(c_0)=1$ for
$\xi\in K$. Let $0<c\le1$ and define
\[
 \omega_\xi=\begin{cases}1,&\xi(0,0,h)=1\text{ for all }h\in A,\\
                         c,&\text{otherwise}.
              \end{cases}
\]
Assume that $M,K_0,C_E\ge0$ are fixed independently of the growing
host and that
\begin{gather*}
 |\mathcal B|\le M,\qquad
 \left|\sum_{t\in T}\xi(0,0,t)-\omega_\xi n_T\right|
       \le Mx^{-1+\varepsilon}n_T\quad(\xi\in K),\\
 |\ell_\xi|\le K_0\quad(\xi\in K\setminus\{1\}),\qquad
 |\ell_\xi|\le C_E\log x\quad(\xi\in\mathcal B\setminus K).
 \tag{WB.17}
\end{gather*}
For characters trivial on $A$ the first error is exactly zero.
Write
\[
 V(x)=\sum_{\xi\in K}\omega_\xi\Re\mathcal G_\xi(x)
                    -\sum_{\xi\in\mathcal B\setminus K}
                                      |\mathcal G_\xi(x)|.
 \tag{WB.18}
\]
We claim that, for all sufficiently large $x$ with a threshold depending
only on $M,K_0,C_E,\varepsilon$,
\[
 \frac1{n_T}\sum_{t\in T}\Re\sum_{\xi\in\mathcal B}
       \xi(e_1(t))\mathcal G_\xi(a_t/b)
 \ge V(x)-4x^{1/2+\varepsilon}.
 \tag{WB.19}
\]
Here are all the error terms. Put $y_t=a_t/b$. The principal character
costs at most $C_0$ by (WB.10). For a retained nonprincipal character,
$\xi(e_1(t))=\xi(0,0,t)$, and the difference from
$\omega_\xi n_T\mathcal G_\xi(x)$ is
\[
 \sum_t\xi(0,0,t)(\mathcal G_\xi(y_t)-\mathcal G_\xi(x))
 +\left(\sum_t\xi(0,0,t)-\omega_\xi n_T\right)
                                                   \mathcal G_\xi(x).
\]
After division by $n_T$ its norm is at most
\[
 \log(x+1)+C_0+K_0
       +Mx^\varepsilon(\log x+C_0+K_0).
\]
For an omitted character, use absolute values; its lower bound is
$-|\mathcal G_\xi(x)|-\log(x+1)-C_0-C_E\log x$.
Thus the total loss is bounded by $MR(x)$, where $R$ is the maximum
of these three displayed bounds. Since all their coefficients are
fixed,
\[
 \frac{MR(x)}{x^{1/2+\varepsilon}}\longrightarrow0.
\]
This proves (WB.19), including uniformity in $n$, $|A_0|$, the prime
pair, $m$, the individual weights, and the number of all characters
of $G$. Only the bounded quotient family $\mathcal B$ is counted.

The actual biased-set hypothesis supplies the Fourier estimate rather
than assuming it separately. Let $[A:A_0]=o$ and
$c=(3o+1)^{-1}$. If for every coset $r+A_0$
\[
 \left|\#(T\cap(r+A_0))-
       \frac{3+\mathbf1_{r\in A_0}}{3o+1}\,n_T\right|
       \le x^{-1+\varepsilon}n_T,
 \tag{WB.20}
\]
then every nontrivial character of $A/A_0$ has sum over $T$ within
$ox^{-1+\varepsilon}n_T$ of $cn_T$: the constant part vanishes
by orthogonality and the distinguished coset contributes $cn_T$.
The triangle inequality bounds the $o$ errors. The quotient
$G/(0,0,A_0)$ has order $4o$, so
$|\mathcal B|=4o$. Taking $M=4o$ proves (WB.17). This is the
mathematical reason for the source's particular $3:4$ bias.
% END WB5

% BEGIN WB6
\subsection{WB6. Closing the scalar and primitive-character joins}

Combining (WB.16) and (WB.19) gives
\[
 \frac{Qd_2}{2}\bigl(V(x)-4x^{1/2+\varepsilon}\bigr)
                   \le E+2Q/x.
 \tag{WB.21}
\]
Assume now the exact scalar sizes from SC1:
\begin{gather*}
 \frac{n_T}{n}\ge\tfrac12x^{-\varepsilon/8},\qquad
 x^{-\varepsilon}\le d_2\le1,\\
 x^{1/2-\varepsilon}\le m\le x^{1/2+\varepsilon/2},\qquad
 xm\le k\le n^2x^{1/2-\varepsilon},\qquad
 0\le a_E\le2xb,
\end{gather*}
and put
\[
 B=n_Ta_E(1+m^2/k)
       +m\left(2nmbx^\varepsilon+
                \frac{n_2kbx^\varepsilon}{n(nx^{-\varepsilon})}\right).
 \tag{WB.22}
\]
For some fixed $C\ge0$, the remaining gap bound must give
$E\le x^{C\varepsilon}B$. The complete proof of SC1 applied
to (WB.21) gives, once $26\le x^{3\varepsilon/8}$,
\[
                    V(x)\le x^{1/2+(C+3)\varepsilon}.
 \tag{WB.23}
\]
The elementary proof can be read without a dependency search: divide
$B$ by $Q$. Its four terms are at most $2Z,2Z,4Z,2Z$, respectively,
where $Z=x^{1/2+13\varepsilon/8}$. Divide (WB.21) by $Q/2$,
use $d_2\ge x^{-\varepsilon}$, and obtain
$V\le26x^{1/2+(C+21/8)\varepsilon}$. Absorb $26$ as stated.

For the fixed primitive target $\psi$ of order $o$, retain exactly
the primitive inducers of its powers, read via (WB.3). Require this
reading to be onto that fixed family. Its injectivity follows from
(WB.4). The source's cyclic level has \texttt{psK=none}; its
\texttt{EAtH.ps} clause excludes the auxiliary character
$\chi_q(a,b,h)=(-1)^b$ from $K$.
Any character trivial on $A$ and taking value $1$ at $c_0$ is
either $1$ or $\chi_q$. Thus only the trivial character in $K$ has
weight $1$, and all other retained characters have weight $c$.
Consequently the retained term in (WB.18) is exactly
\[
 H_\psi(x)=G(x)+\frac1{3o+1}
       \sum_{\substack{\chi\text{ primitive inducer of }\psi^j\\\chi\ne1}}
                                             \Re G_\chi(x).
 \tag{WB.24}
\]
This equality uses the actual primitive values and the actual linear
terms at $1$. A mere equality of the number of characters would not
suffice. A general source level can retain $\chi_q$ with weight
$1$ through its \texttt{psK} field; no such extra term occurs for
the cyclic levels used here.

If every omitted member of $\mathcal B$ obeys the supplied estimate
$|\mathcal G_\xi(x)|\le x^{1/2+\varepsilon}$, then
\[
 \Sigma(x):=\sum_{\mathcal B\setminus K}|\mathcal G_\xi(x)|
                  \le4o\,x^{1/2+\varepsilon}.
\]
The pointwise conclusion is therefore
\[
 \boxed{\ H_\psi(x)\le x^{1/2+(C+4)\varepsilon}\ }
 \quad\text{when }x^\varepsilon\ge1+4o.
 \tag{WB.25}
\]
In the source $C=e_D+e_L+3\ge0$, giving exactly
$1/2+(e_D+e_L+7)\varepsilon$. All thresholds used to reach
(WB.25) depend on fixed target constants and $\varepsilon$, before
the varying host and its data are considered.
% END WB6

% BEGIN WB7
\subsection{WB7. The complete input contract at the arithmetic boundary}

The preceding implication has the following explicit contract. For
each fixed primitive target and each sufficiently small
$\varepsilon>0$, supply arbitrarily late windows
$[Y,Y^{3/2}]$. For every $x$ in one such window, supply
\emph{one common datum} satisfying all of the following:
\begin{enumerate}
\item An onto reciprocity reading, its actual inertia and Frobenius,
the split/ramified conditions (WB.6), and the actual values
$\ell_\xi=L'/L(1,\xi^*)$.
\item A subgroup of fixed index $o$, the fixed primitive-power
retained family, its auxiliary-character exclusion, the biased
counts (WB.20), and fixed bounds $K_0,C_E$ in (WB.17).
\item The same $T,T_2,a_t,w_{u,j},m,b,b',X',\eta,\theta,N$
satisfying (WB.12), together with the omitted-ramp estimates in WB6.
\item The scalar size inequalities of WB6 and the actual arithmetic
gap (\ref{eq:wb-arithmetic-gap}) with
$E\le x^{C\varepsilon}B$, for a fixed $C\ge0$.
\end{enumerate}
Here $C$ is chosen before $\varepsilon$. Bounds such as $K_0,C_E$
may be chosen for the fixed target (or even for a fixed
$\varepsilon$), but must be uniform as the windows move to infinity.
The common datum may depend on $x$; the target family and $C$ may not.

Every implication from this contract to (WB.25) is proved in
WB1--WB6. There are no remaining character-orthogonality,
ramified-value, moving-cutoff, or scalar-comparison inputs in this
part of the chain. The unproved task is to construct this contract
from the arithmetic geometry on the same maps and fields.

For precise navigation, the source packages items 1--3 in
\texttt{HostAx}, \texttt{EAtH}, the reading part of
\texttt{CoreAtH}, \texttt{SizesA}, \texttt{ThinScales}, and
\texttt{CostWindows}. Item 4 is the conclusion of \texttt{ch\_\allowbreak top}
and the remaining fields of \texttt{CoreAtH}. Its gap is obtained
there from the eleven-field \texttt{SizeData2} package, uniform
error windows, and the two bounds on $B_0,N_c$. Those are not proved
merely by unpacking \texttt{CoreAtH}.

The algebraic part of that last source join is short and worth
recording. \texttt{IdentU} supplies the exact bijections from the
original top indices to $T$ and from the place indices to $\Phi$,
the equalities of $a,w,\gamma$, the equality
$\mathrm{LpL}=\Re\ell_{\chi_d}$, and the equality of the asked
sums. Substitution into the definition of the ledger main term gives
\[
 \mathrm{Main}(u)=\sum_{t,j}\Re\mathcal M(a_t,w_{\mathrm{tr}(u),j}),
 \qquad
 \mathrm{ASK}(u)-\mathrm{Main}(u)
    =\operatorname{gap}(\mathrm{tr}(u)).
 \tag{WB.26}
\]
The source's coverage clause $A_0\subset\mathrm{range}(\mathrm{tr})$
then transfers its bound to every $u\in A_0$ needed in (WB.13).
The $B_0,N_c$ inequalities replace its error bracket by (WB.22),
termwise and with nonnegative coefficients. Finally
$e_{\rm Div}\le e_D$, $e_{\rm Loss}\le e_L$ bound its error
exponent by $C=e_D+e_L+3$. These are exact substitutions on one
supplied ledger. The independent divisibility, metric, dimension,
theta-norm, and spectral estimates that establish that ledger are
still the arithmetic frontier.
% END WB7

% BEGIN WB8
\subsection{WB8. A closed terminal theorem with the quantifiers visible}

\begin{theorem}[Arithmetic contract implies the fixed target's zero-free half-plane]
Fix a primitive Dirichlet character $\psi$. Suppose there exist
$C\ge0$ and $\varepsilon_0>0$ such that for every
$0<\varepsilon<\varepsilon_0$, the contract in WB7 can be supplied
on arbitrarily late windows $[Y,Y^{3/2}]$. Then every primitive
inducer of a power of $\psi$ has no zero in $\Re s>1/2$,
away from the principal pole. If the contract is available for every
primitive $\psi$, RH and GRH for all Dirichlet $L$-functions follow.
\end{theorem}

\begin{proof}
Choose $0<\delta<1/2$ and then
\[
 0<\varepsilon<\min\{\varepsilon_0,\delta/(C+4)\}.
\]
All the pointwise thresholds in WB3--WB6 are now fixed. Choose one
of the contract's windows whose left endpoint exceeds them and any
prescribed $Y_0$. Equation (WB.25) holds at \emph{every} real $x$
in this window and gives $H_\psi(x)\le x^{1/2+\delta}$.
The positive-family one-sided criterion AN3 now rules out any
constituent zero with real part greater than $1/2+\delta$.
Given a proposed zero with real part $\beta>1/2$, choose
$0<\delta<\min\{1/2,\beta-1/2\}$; this is a contradiction. Positivity of
all coefficients in (WB.24), including the residue at a possible
real zero, is the point that allows a one-sided upper bound.

For an arbitrary character induced by primitive $\chi^*$,
\[
 L(s,\chi)=L(s,\chi^*)
       \prod_{\substack{p\mid N\\p\nmid\mathrm{cond}(\chi^*)}}
                       (1-\chi^*(p)p^{-s}).
\]
Every additional factor is nonzero for $\Re s>0$. The primitive
functional equation maps a zero left of $1/2$ in the open critical
strip to a zero right of $1/2$ of the conjugate character; its gamma,
conductor, and root-number factors are nonzero there. This proves
the all-character strip statement. The ordinary target $\psi=1$
has $o=1$ and $H_\psi=G$, giving RH by AN2, including its standard
functional-equation treatment outside the open strip. At $s=1$,
nonprincipal characters use their classical nonvanishing theorem;
the principal pole is excluded as in the exported statement.
\end{proof}

The retained analytic leaves are exactly those stated in AN2--AN3
and GH6: ordinary analytic continuation, the functional equations,
Euler products and classical nonvanishing on $\Re s=1$, together
with elementary complex analysis. This theorem adds no cohomology
or representation hypothesis beyond WB7's explicitly open
arithmetic contract.

The source's \texttt{tiedWin\_\allowbreak of\_\allowbreak ex} supplies the matching
quantifier rule: first choose the common bounds demanded by all
links, and only then use a family-existence theorem valid for
\emph{every} such bound. The dependency audit below spells this out.
On the literal \texttt{ShapeWin}\,$\to$\,\texttt{topDesWin}
route the chosen values are $e_D=2,e_L=0$, so its conditional
terminal exponent is $1/2+9\varepsilon$. This is an identification
of the constants in that route, not a new construction of its
required simultaneous family.
% END WB8


% A dependency audit and two explicit conditional joining lemmas.
% This fragment uses only the macros and theorem environments of the main note.
\section{A dependency audit organized by complete chains}
\label{sec:dependency-audit}

\status{This is a human mathematical audit of the reconstruction and a
targeted reading of the pinned source interfaces.  It is not an extraction
of Lean proof-term dependencies.  Its two joining arguments below are
complete implications from the stated premises.  The arithmetic suppliers
of those premises remain separate obligations.}

There are three different meanings of ``complete'' that matter here.
A local theorem can have a complete proof for every datum satisfying its
hypotheses.  A chain can be complete relative to explicitly named earlier
theorems.  A global application is complete only when all those hypotheses
hold on one actual family, with the required order of choices.  The local
Type-I and Type-II results are examples of the first meaning; QP's
\texttt{PairEx}, relative to CN/BD/HD, is an example of the second.
The remaining work toward the arithmetic window is predominantly the third.

The historical claim graph was designed to record mathematical references,
not to express every retained premise as a graph vertex.  For example,
AN2 has no predecessor vertex but assumes prime windows; AN3 and SC2 have
empty \texttt{remaining\_\allowbreak inputs} arrays although their prose retains
substantive hypotheses.  QP25 records QP22 as a predecessor and also
assumes \texttt{VWinAt}.  None of these is a mathematical contradiction.
It means that graph acyclicity, or the number of recorded claims, cannot
certify completion.  The separate audit data makes such obligations
first-class objects and leaves all historical edges unchanged.

\subsection{The exact top-level sufficient interfaces}

Fix a primitive target $\psi$ and let $H_\psi$ be its fixed positive
character family.  The downstream chain is
\[
 \begin{gathered}
 \text{one compatible family on arbitrarily late windows}\\
 +\ \text{actual arithmetic gap and numerical bounds}\\
 \Longrightarrow H_\psi(x)\le x^{1/2+A_\psi\epsilon}
 \Longrightarrow \text{zero-freeness for the target}.
 \end{gathered}                                                    \tag{DA.1}
\]
Here $A_\psi$ and a positive permitted range for $\epsilon$ are fixed
\emph{before} $\epsilon$ is chosen.  Taking every primitive target and
using GH6 gives the Dirichlet conclusions; taking only the trivial target
gives ordinary RH.  The full positive-family analytic implication is
AN2--AN3, not an additional open analytic problem.

The literal source organizes the upstream cut using
\texttt{GChi}$(\psi)$, which requires, on the same host and selected set,
\[
 \texttt{EAtH}(\operatorname{Level.cyc}\psi,H,T_T,x,\epsilon,4)
 \quad\text{and}\quad
 \texttt{VVle}(E,t,\sqrt b).                               \tag{DA.2}
\]
Its family class also requires arbitrarily late windows.  The first
condition supplies the character estimates and the second the
cohomological budget.  The source's
\texttt{HostExW} supplies the actual ranks, \texttt{HostGood} data and
this target requirement on a scale family.  A description of the abstract
field $E$, or the existence of a pair $q,d$, does not yet provide that
whole record.

The sufficient source cut is
\[
 \begin{gathered}
 \texttt{HostExW}(3/2,\texttt{GChi}(\psi).\texttt{req})\quad\text{and}\quad
 \text{the design/link construction}\\
 \Longrightarrow
 \texttt{FamilyClassConnected}(\texttt{GChi}(\psi).\texttt{cls}).
 \end{gathered}
                                                               \tag{DA.3}
\]
The last predicate supplies fixed $e_D,e_L$ before $\epsilon$ and a
family with \texttt{ChainTied} at every sufficiently late scale in its
window set.  The final character/scalar chain then gives (DA.1).
This is a sufficient cut through the source argument, not a claim of
logical necessity or a minimal set of possible mathematical ideas.

\subsection{How one family can satisfy finitely many closed conditions}

The source has an explicit solution to the order-of-choices problem.
Write $\mathcal C$ for a class of families and
$\operatorname{Large}(F,K,H)$ for its monotone threshold condition.
Thresholds have a finite supremum satisfying
\[
 \operatorname{Large}(F,K\vee K',H)
 \quad\Longleftrightarrow\quad
 \operatorname{Large}(F,K,H)\ \text{and}\
 \operatorname{Large}(F,K',H).                             \tag{DA.4}
\]
The literal definition of \texttt{Closed} is
\[
 \begin{split}
 \operatorname{Closed}_{\mathcal C}(P):\quad
 &\exists\epsilon_P>0\ \forall\epsilon\in(0,\epsilon_P]
 \ \exists K_P\ \forall F,H,\\
 &F\in\mathcal C,\ \operatorname{Large}(F,K_P,H)
 \ \Longrightarrow P(F,H).
 \end{split}                                                   \tag{DA.5}
\]
In contrast, \texttt{DesignExists} quantifies over the threshold first:
for each sufficiently small $\epsilon$ and \emph{every} $K$, it supplies
one $F,H$ in the required class, large for $K$, with
\texttt{ChosenE}, \texttt{TuSize} and \texttt{DGamSize}.

\begin{proposition}[DA1: finite uniform joining on one family]
\label{prop:da-closed-join}
If $P_1,\ldots,P_r$ satisfy (DA.5) on the same family class, then
their conjunction does too.  If \texttt{DesignExists} supplies families
in that class, one chosen family satisfies all $P_i$ and all the design
conditions simultaneously.  The positive range for $\epsilon$ is chosen
before $\epsilon$, and the joined threshold is chosen before the family.
\end{proposition}

\begin{proof}
Include the design supplier's range in the finite minimum of all
positive ranges.  For an $\epsilon$ in that range, take
the threshold $K_i$ furnished by each closed statement and put
$K=\bigvee_iK_i$.  Design existence, which holds for every $K$, now
chooses $F,H$ large for this $K$.  Repeated use of (DA.4) makes it large
for each $K_i$, so every $P_i(F,H)$ follows.  For no conditions there is
nothing to join; choose any design threshold.  All conclusions use the
one family just selected.  This is the entire mathematical content of
the source's \texttt{Closed.and} and its use in \texttt{topDesWin}.
It does not prove \texttt{DesignExists} or any of the closed premises.
\end{proof}

This quantifier structure is stronger than obtaining an individual
witness for each desired property.  It is also more precise than
invoking compatibility after all witnesses have already been chosen.
Every closure must refer to the same class; equal dimensions, isomorphic
fields, or unrelated enlargements of coefficient fields do not supply
the dependent identifications in (DA.5).

\subsection{The complete logical passage through the global links}

Here is the finite conjunction used by \texttt{topDesWin}.  In the table,
``closed'' means the exact quantifiers (DA.5), on the same design class.
Conditional clauses retain their dependence on the design's
\texttt{ChosenE}, \texttt{TuSize} and \texttt{DGamSize} fields.

\begin{center}\small
\begin{tabular}{p{0.24\linewidth}p{0.67\linewidth}}
\toprule
Input & Required content on the actual family\\
\midrule
\texttt{DesignExists} & One target-compatible family after every
finite threshold choice, with the selected parameters and size fields.\\
\texttt{LinkN} & Congruence-denominator cost bounded by the actual
cohomological sum and polynomial budget.\\
\texttt{LinkEta} & Coverage of the actual twisted local data.\\
\texttt{LinkIRR} & Nonnegative retained-row defect and its stated
exponentially small upper bound.\\
\texttt{SmallLevelRankEq} & The actual equality $d_{Z,0}(j)=D_0$
for every lift used.\\
\texttt{LinkDIVt} & Given \texttt{LinkEta} and \texttt{DIVinN}, the
common period target, integral/height bounds and \texttt{LinkDIV}.\\
Fullness and rank & \texttt{ChosenE} implies \texttt{NuFullF} and
positive retained ranks \texttt{J.MuRank}.\\
\texttt{LinkG} & The local divisibility/gain theorem on these actual
lattices and served orbits.\\
\texttt{LinkS} & The actual normalized off-diagonal and diagonal Gram
bounds and total-target dimension bound.\\
\texttt{LinkGLOBAL} & The same-family slope comparison turns
\texttt{ChosenE} and the five global links into \texttt{ChainTied}.\\
\bottomrule
\end{tabular}
\end{center}

The target class supplies \texttt{HypV}, namely
$\texttt{VVle}(E,t,\sqrt b)$ at every sufficiently late scale.
The exact metric transport used by the source is additionally visible:
\[
 \texttt{LinkDIV}+\texttt{NuFullF}+\texttt{J.MuRank}
 +\texttt{J.MuTu}+\texttt{J.MuDGam}
 \Longrightarrow\texttt{LinkMu}.                          \tag{DA.6}
\]
Its final wrapper is \texttt{J.linkMu\_\allowbreak of\_\allowbreak linkDIV}; it relies on the
actual \texttt{NprimeMuR} lattice/metric comparison.  That comparison
is retained here.  In particular, (DA.6) is not being deduced merely
from the ranks of unrelated spaces.

\begin{proposition}[DA2: the conditional link chain closes]
\label{prop:da-topdeswin}
Assume the exact common-class design and closed premises in the table,
the target-class inclusion and \texttt{HypV}, the size consequences
\texttt{TuSize}$\Rightarrow$\texttt{J.MuTu} and
\texttt{DGamSize}$\Rightarrow$\texttt{J.MuDGam}, and the actual
metric transport (DA.6).  Then they supply one family with
\texttt{FamilyClassConnected}, with the constants $e_D,e_L$ in
\texttt{LinkGLOBAL} fixed before $\epsilon$.
\end{proposition}

\begin{proof}
Apply DA1 to the nine closed link/rank conditions in the table and to
the design supplier, after transporting along the stated class-inclusion
maps.  All statements now concern one $F,H$.  It remains to check the
one numerical arrow that turns \texttt{HypV} into the needed
congruence budget; the other arrows are the explicit hypotheses.

At a sufficiently late scale, write $V$ for the actual weighted sum of
the finite cohomological lengths.  The finite-subsums clause in
\texttt{VVle}, together with nonnegativity, shows that this series is
summable and
\[
                     0\le V\le\sqrt b\le b\qquad(b\ge1).
                                                               \tag{DA.7}
\]
The finiteness clause and the literal \texttt{NData.sha\_\allowbreak eq} identify
these lengths with the stored lengths in the same $D$.  Thus
\texttt{ShaFinOdd} holds and \texttt{LinkN} applies.
Put $y=x^\epsilon$, so $y\ge1$, and let
$B_0=fmb\,y$, with $fm\ge1$ and $N_c\ge0$.  Eventually
$\log^2x\le y$.  The \texttt{LinkN} inequality, with
$d_{Z,0}(j)=D_0$ from \texttt{SmallLevelRankEq}, is
\[
 \eta_j\le D_0\{C(fmb\,y^2+\log^2x\,N_c)+V\}.
                                                               \tag{DA.8}
\]
The expression multiplying $C$ is nonnegative and is at most
$y(B_0+N_c)$.  Also $V\le b\le B_0\le y(B_0+N_c)$.
Replacing $C$ by $|C|$ proves
\[
           \eta_j\le D_0(|C|+1)y(B_0+N_c),                \tag{DA.9}
\]
which is precisely \texttt{DIVinN}.  There are only finitely many
eventual statements in this argument, each already uniform in all
family indices and used lifts.  Taking the maximum of their thresholds
therefore preserves that uniformity; it does not assume that the index
sets themselves are finite.

Now \texttt{LinkEta} and \texttt{LinkDIVt} give \texttt{LinkDIV}.
The design size fields give the two size consequences in (DA.6);
fullness and positive rank are supplied on this very family, so (DA.6)
gives \texttt{LinkMu}.  The actual package
\[
 \texttt{GlobalTheoremInputs}
 =\texttt{LinkG}\wedge\texttt{LinkS}\wedge\texttt{LinkIRR}
     \wedge\texttt{LinkDIV}\wedge\texttt{LinkMu}
\]
is therefore available.  Apply \texttt{LinkGLOBAL} with
\texttt{ChosenE}.  This supplies the remaining constants and
\texttt{ChainTied} at every sufficiently late scale in the same
family's window set.  The definition of
\texttt{FamilyClassConnected} is exactly this conclusion, with
$e_D,e_L$ already fixed.  This proves the asserted conditional chain.
\end{proof}

DA2 proves a complete assembly; it does not turn any listed closed
supplier into an existence theorem.  The source's particular
\texttt{shapeWin} call uses $C_{\rm abs}=2,C_0=3,e_D=2,e_L=0$.
Thus its literal final coefficient is $e_D+e_L+7=9$ if that entire
route and all its hypotheses are supplied.  The reconstruction need
not fix this numerical value: any finite coefficient fixed before
$\epsilon$ suffices for the analytic contradiction.

\subsection{The remaining cut and the alternative routes}

The audit uses an AND gate for simultaneous requirements and an OR gate
only when either complete route suffices.  The following are the main
branches; none can be retired solely because its downstream assembly
has been written.

The strongest new local leaf is HP1--HP4: the characteristic-zero
canonical polynomial highest-weight vector, fixed line, and tensor
irreducibility theorem are now proved.  HP5--HP6 identify these with
the existing weight/form interfaces while retaining the exact supplied
field, frame and row data.  Thus the Schur-theoretic premise formerly
left in MR5/MR19 can be discharged.  The geometric realization of that
representation and the other fields of the actual \texttt{Core} are
different obligations.  WB1--WB8 and VW1--VW4 close two complete
downstream interfaces, as specified below; their arithmetic and
cohomological suppliers remain explicit.

\begin{enumerate}
\item \textbf{Cohomological windows.} The complete downstream route is
\texttt{JointWin}+\texttt{ThmBc}$\Rightarrow$\texttt{VWinAt}, then QP22
supplies the fixed-pair windows.  An upstream sufficient alternative
is \texttt{ThmBab}, \texttt{PointHV}, and \texttt{ThmBcFree}.
The actual two-primary and odd-plus/odd-minus bounds remain the leaves.
Separate bounds on separate intervals cannot replace the joint-window
premise.
\item \textbf{Actual geometric and arithmetic family.} Full
\texttt{Core}, the dependent layout, original coefficient embeddings,
\texttt{St}/\texttt{ColPin}/\texttt{R3.Read}, rational good-row input
and genuine theta/local installation remain.  MR22's initial zero
theta and empty toric family are exactly its stated initialization;
they do not supply a later nonplaceholder \texttt{NTh} or local family.
\item \textbf{Actual theta and global analytic operators.} Critical
toric ratios and \texttt{Choice.FieldOK}; convergence, separation,
nonvanishing and operator identities; actual chain norms, normalized
spectral words, positive kernel/mass identities and exceptional terms
are distinct obligations.  The finite Weil dictionary does not imply
these identities for the global operators.
\item \textbf{Integral/metric and congruence comparison.} The true
source/target lattices and metrics, discarded directions, retained
ranks and \texttt{LinkN} denominator bounds must remain on the same
objects.  Rank identities and abstract isometries on different
spaces do not prove the required transport.
\end{enumerate}

There is a useful OR gate within the real-stage branch.  For the actual
large-gap rows, DC20 supplies the four \texttt{RealTied} fields using
the named Harish--Chandra degree theorem and GV's Tamagawa input with
the exact normalizations.  The stronger universal route uses
\texttt{DiscreteSeriesExist} and \texttt{FormalDegreeFormula2}.
It is an alternative and need not also be reconstructed if every row used
is covered by the first route.  In either case the other fields of
\texttt{Core} remain required.

Named established theorems, such as CN, BD, HD, Harish--Chandra formal
degree, Tamagawa and exact global Artin reciprocity, are distinguished
from these bespoke suppliers.  They can be accepted as explicit citation
leaves in a readable proof.  Expanding all classical background is a
different goal from closing the new argument's dependency chains.
An interface merely bearing the word ``literature'' is not such a
citation leaf unless its full mathematical statement has been matched
to an established theorem.

The machine-readable companion records the gates, source anchors,
quantifier obligations and binding keys.  It reports no completion
percentage and does not infer a global proof from a count of local claims.
Future work should retire a leaf only with its exact supplier or a fully
proved replacement, then propagate that closure through the gates above.


% Version 0.24.  Classical finite duality, with all source conventions matched.
\section{Finite duality and unramified reciprocity on the original fields}
\label{sec:local-duality-closure}

\status{This section supplies the literal finite-module
\texttt{Unramified\allowbreak Orthogonality} interface and matches
\texttt{Published.\allowbreak Poitou\allowbreak Tate\allowbreak Duality} to the classical finite-module
Poitou--Tate theorem.  The elementary procyclic, Euler-characteristic,
annihilator and reciprocity deductions are proved.  Local Tate duality,
local class field theory and finite-module Poitou--Tate duality are
explicit classical citation leaves.  No main conjecture, special-value
estimate, or arithmetic prime window is assumed or proved here.}

The relevant coefficients are finite modules, not necessarily free
modules.  Fix a prime $\ell$, an integer $r>0$, and
$R=\mathbb Z/\ell^r\mathbb Z$.  For a number field $F$ with its specified
algebraic closure $\Omega$, write $R(1)=\mu_{\ell^r}(\Omega)$, in additive
notation.  A primitive root $\zeta$ gives the identification
$a\mapsto\zeta^a$ with the source's rank-one representation
\texttt{char\allowbreak Rep(cyc\allowbreak Mod)}.  It is equivariant because
$g\zeta=\zeta^{\chi_{\ell,r}(g)}$.  The character
$\chi_{\ell,r}$ is independent of the chosen primitive root.
For a finite continuous $R$-representation $M$, put
\[
 M^D=\operatorname{Hom}_R(M,R(1)),\qquad
 (g\varphi)(m)=\chi_{\ell,r}(g)\varphi(g^{-1}m).
 \tag{LD.1}
\]
Thus $M^D$ is exactly \texttt{tw\allowbreak Dual}, including its inverse in the
argument of $\varphi$.

At a finite place $v$, keep the chosen valuation subring $A\subset\Omega$.
Its decomposition group $D_A$ and inertia group $I_A$ identify with
$G_{F_v}$ and its inertia.  One can make this identification by extending
the chosen valuation to an embedding in an algebraic closure of $F_v$:
finite extensions of the henselization and of the completion have the
same finite separable extensions, with the same ramification and residue
extensions.  Passing to inverse limits gives the asserted group
identification.  A different embedding conjugates both groups and every
restriction map.  We use the given $A$ throughout; this identification
does not replace $F$, $\Omega$, or its local data by a new arithmetic
witness.

\subsection{Procyclic cohomology and the Euler calculation}

\begin{proposition}[LD1: finite procyclic and tame Euler calculation]
\label{prop:ld-procyclic}
Let $\Gamma=\widehat{\mathbb Z}$ with its specified topological generator
$\phi$, acting continuously on a finite $\ell$-primary module $T$.  Then
\[
 H^0(\Gamma,T)=\ker(\phi-1),\qquad
 H^1(\Gamma,T)=T/(\phi-1)T,\qquad
 H^i(\Gamma,T)=0\quad(i\ge2).
 \tag{LD.2}
\]
In particular $|H^0(\Gamma,T)|=|H^1(\Gamma,T)|$.
If $K$ is a nonarchimedean local field of characteristic zero and
residue characteristic $p\ne\ell$, and $M$ is any finite continuous
$\ell$-primary $G_K$-module, then its cohomology is finite and
\[
 |H^1(K,M)|=|H^0(K,M)|\,|H^2(K,M)|,
 \qquad H^i(K,M)=0\quad(i\ge3).
 \tag{LD.3}
\]
No unramifiedness assumption on $M$ is needed in (LD.3).
\end{proposition}

\begin{proof}
A continuous crossed homomorphism on $\Gamma$ is determined by its
value at $\phi$, since the integer powers of $\phi$ are dense.  Every
$t\in T$ occurs: if $\phi^d$ acts trivially and $\ell^aT=0$, the rule
\[
 f(\phi^j)=\sum_{i=0}^{j-1}\phi^it\quad(j\ge0)
\]
factors through the cyclic quotient of order $d\ell^a$.  Indeed its
value at that order is $\ell^a\sum_{i=0}^{d-1}\phi^it=0$.
Coboundaries take at $\phi$ the values $(\phi-1)t$, proving the
degree-one formula.  Invariants give degree zero.

For completeness, the higher-degree calculation is the continuous
version of the cyclic resolution with alternating operators
$\phi-1$ and $1+\phi+\cdots+\phi^{n-1}$.  A continuous cochain with
finite target factors through a finite cyclic quotient; hence continuous
cohomology is the direct limit over such quotients.  Inflation from
order $n$ to order $kn$, where both actions factor through order $n$,
multiplies the even-degree norm-class representative by $k$ in degree
two.  In every degree $i\ge2$ the comparison of the periodic
resolutions contains a factor $k$ (and possibly a higher power of $k$).
Choosing $k$ divisible by the exponent of $T$ kills that class.
This proves the vanishing in (LD.2).  Alternatively the same calculation
is the two-term continuous resolution for the pro-$\ell$ factor,
after taking invariants under the prime-to-$\ell$ factor.  Finally the
kernel and cokernel of an endomorphism of a finite group have equal
orders, which proves the cardinality assertion.

For (LD.3), let $I$ be inertia and $P$ wild inertia.  The group $P$ is
pro-$p$, so its higher cohomology on $\ell$-primary modules vanishes:
on each finite quotient the usual averaging homotopy divides by a
$p$-power, which is invertible on the module.  The tame quotient is
\[
 I/P\simeq\prod_{q\ne p}\mathbb Z_q(1).
\]
The prime-to-$\ell$ factor again has exact invariants.  Applying the
procyclic calculation to its remaining $\mathbb Z_\ell(1)$ factor
shows that $A=M^I$ and $B=H^1(I,M)$ are finite and $H^i(I,M)=0$ for
$i\ge2$.  The residue Galois group $G_K/I$ is
$\widehat{\mathbb Z}$.  The Hochschild--Serre cochain filtration,
with only its two inertia rows and two residue columns remaining, gives
\[
 H^0(K,M)=A^\phi,\qquad
 0\longrightarrow A/(\phi-1)A\longrightarrow H^1(K,M)
   \longrightarrow B^\phi\longrightarrow0,
 \tag{LD.4}
\]
and $H^2(K,M)=B/(\phi-1)B$, with no higher groups.  All possible
differentials leave these two columns, so are zero.  Applying the
kernel--cokernel order equality separately to $A$ and $B$ gives
(LD.3).  This also proves all required finiteness in this case.
\end{proof}

\subsection{The exact annihilator of unramified classes}

\begin{proposition}[LD2: the literal unramified orthogonality supplier]
\label{prop:ld-unramified}
Let $K=F_v$ have residue characteristic different from $\ell$, and
let $M$ be a finite continuous $R$-representation on which inertia acts
trivially.  Define
\[
 H^1_{\rm ur}(K,M)=
 \ker\bigl(H^1(K,M)\longrightarrow H^1(I,M)\bigr).
\]
In the finite local Tate pairing
\[
 H^1(K,M^D)\times H^1(K,M)
 \xrightarrow{\ \cup\ }H^2(K,R(1))\simeq R,
 \tag{LD.5}
\]
the two unramified subgroups are exact annihilators in both directions.
Consequently \texttt{Reciprocity\allowbreak Theory.\allowbreak Gal.\allowbreak Unramified\allowbreak Orthogonality}
holds for every original number field, algebraic closure, place and
finite coefficient module occurring in that interface.
\end{proposition}

\begin{proof}
Prime-to-$p$ roots of unity are unramified over $K$.  To see this,
choose a finite residue extension containing the roots of
$X^{\ell^r}-1$; its derivative is invertible at every root.  Hensel
lifting in the corresponding unramified local extension gives all
$\ell^r$-th roots.  Thus $R(1)$, and hence $M^D$, are unramified.

Inflation--restriction and LD1 identify
\[
 H^1_{\rm ur}(K,M)=H^1(G_K/I,M)=M/(\phi-1)M,
 \qquad |H^1_{\rm ur}(K,M)|=|H^0(K,M)|.
 \tag{LD.6}
\]
The same holds for $M^D$.  The product of two unramified classes is
inflated from $H^2(G_K/I,R(1))=0$.  Each unramified subgroup therefore
annihilates the other.

The one classical input in this argument is finite local Tate duality:
cup product followed by the local invariant is a perfect pairing
$H^i(K,M)\times H^{2-i}(K,M^D)\to\mathbb Q/\mathbb Z$,
for $i=0,1,2$.  See Milne, \emph{Arithmetic Duality Theorems},
second edition, I, Corollary~2.3.  Since these groups are killed by
$\ell^r$, the target is its $\ell^r$-torsion, identified with $R$.
In particular
\[
 |H^2(K,M)|=|H^0(K,M^D)|.
\]
Combine this with LD1 and (LD.6):
\[
 |H^1(K,M)|=
 |H^1_{\rm ur}(K,M)|\,|H^1_{\rm ur}(K,M^D)|.
 \tag{LD.7}
\]
For a perfect pairing of finite abelian groups $X$ and $Y$, the
annihilator in $X$ of $Y_0\subset Y$ has order $|X|/|Y_0|$.
Indeed every character of $Y_0$ extends to $Y$, since
$\mathbb Q/\mathbb Z$ is divisible, so the restriction map from
$Y^\vee$ onto $Y_0^\vee$ is surjective.  Equation (LD.7) makes the
already proved inclusion an equality.  Interchanging $M$ and $M^D$
proves the other equality.

There is no hidden freeness condition.  A finite $R$-module is a direct
sum of $R/(\ell^a)$ with $0<a\le r$.  Its $R$-dual on that summand is
$\ell^{r-a}R$, of the same order, and evaluation identifies it with
its double dual.  Thus the $R$-linear duality used in the source agrees
with the finite abelian-group duality above.  The source places its
dual class first in (LD.5); reversing degree-one factors only changes
a sign, which does not change a zero or an annihilator.  Finally the
source requires vanishing in $H^2(K,R(1))$, before applying an
invariant.  The invariant is an isomorphism, so this is exactly the
same condition.  All identifications were made on the originally
chosen decomposition and inertia groups.
\end{proof}

\subsection{One normalized family of local invariants}

\begin{proposition}[LD3: exact finite Poitou--Tate specialization]
\label{prop:ld-poitou-tate}
For every number field $F$, specified algebraic closure $\Omega$,
prime $\ell$, integer $r>0$, and primitive $\ell^r$-th root
$\zeta\in\Omega$, assume $\ell\ne2$ or $F$ is totally complex.
The classical finite-module Poitou--Tate theorem supplies the exact
interface \texttt{Global\allowbreak Theory.\allowbreak Published.\allowbreak Poitou\allowbreak Tate\allowbreak Duality}.
In particular its local invariant family can be fixed \emph{before}
the finite set $S$, all choices of places over $S$, the finite module,
and its continuous action are chosen.
\end{proposition}

\begin{proof}
We first match the normalization, rather than choosing unrelated
isomorphisms $H^2\simeq R$ at different places.  For each finite place
of $F$, Kummer theory embeds $H^2(F_v,\mu_{\ell^r})$ in the Brauer
group with image its $\ell^r$-torsion.  The local invariant
$\operatorname{inv}^{\rm Br}_v:\operatorname{Br}(F_v)\simeq
\mathbb Q/\mathbb Z$ identifies this image with
$\ell^{-r}\mathbb Z/\mathbb Z$.  Let
$\iota_r:R\to\mathbb Q/\mathbb Z$ be
$a\mapsto a/\ell^r$.

Use the common sign for which the Kummer class of a uniformizer,
cupped on the left with the unramified character taking arithmetic
Frobenius to $1$, has invariant $1\in R$.  More explicitly, with the
arithmetic local reciprocity convention in which
\[
 \operatorname{inv}^{\rm Br}_v(\chi\cup\kappa(\varpi))
       =\iota_r(\chi(\operatorname{Frob}_v)),
\]
the source's order $\kappa(\varpi)\cup\chi$ requires
$j_v=-\iota_r^{-1}\operatorname{inv}^{\rm Br}_v$.
The sign follows from graded commutativity.  This elementary cyclic
algebra normalization is part of local class field theory; it is
independent of $v$, $S$, and $M$.

For a chosen valuation subring $A$ above $v$, transport this invariant
along its decomposition-group identification and the map
$a\mapsto\zeta^a$.  The formula is a bijective $R$-linear map
$j_A:H^2(D_A,R(1))\to R$.  Its uniformizer normalization is literally
the source's \texttt{Is\allowbreak Normalised}: a Kummer cocycle $\kappa$ and an
unramified additive character $\chi$ have the displayed representative
$\kappa(g)\chi(h)$, and arithmetic Frobenius is the residue power
map used in \texttt{Is\allowbreak Frob}.  At valuation subrings which do not lie
over finite places, extend the family by the zero map; the interface
places no condition there.

This family is also \texttt{Conj\allowbreak Compatible}.  If $\sigma$ transports
$A$ to $A'$ and $\sigma\zeta=\zeta^c$, the coefficient description
of a transported cocycle is $z'=cz$.  Under the root-of-unity
identification, this says precisely
$\zeta^{z'}=\sigma(\zeta^z)$.  The Brauer class is therefore
transported by the induced field isomorphism, which preserves the
local invariant.  Hence $j_{A'}([z'])=j_A([z])$, the exact source
identity, including the coefficient $c$.  All choices in the family
have now been made.

Let $S$ be an arbitrary finite set of finite places containing all
places above $\ell$.  Enlarge it by the archimedean places only for
the statement of the classical theorem, and write $G_{F,S}$ for the
Galois group unramified outside its finite places.  The source's
quotient by the closed normal subgroup generated by the omitted
inertia groups is precisely this group: a finite quotient factors
through it if and only if its finite extension is unramified at
every omitted finite place.  The coefficient $R(1)$ is unramified
there by the root-of-unity argument in LD2, and hence descends to
$G_{F,S}$.  If $\mu$ is any character whose pullback is the original
cyclotomic character, it is this same descended character, since
the quotient map is onto.  In particular the source's
\texttt{to\allowbreak Mu\allowbreak Loc} is induced by the identity on $R$, not by an
uncontrolled new coefficient identification.

We now apply Milne, \emph{Arithmetic Duality Theorems}, second
edition, I, Theorem~4.10, to an arbitrary finite continuous
$R[G_{F,S}]$-module $M$.  The order of $M$ is a power of $\ell$,
so it is a unit in $\mathcal O_{F,S}$, as required by that theorem.
Let $P^i_S(M)=\prod_{v\in S}H^i(F_v,M)$, for finite $S$, and let
$T^\vee=\operatorname{Hom}_R(T,R)$.  The resulting nine-term sequence
is
\[
\begin{aligned}
0\longrightarrow H^0(G_{F,S},M)&\longrightarrow P^0_S(M)
 \longrightarrow H^2(G_{F,S},M^D)^\vee\\
 &\longrightarrow H^1(G_{F,S},M)
 \longrightarrow P^1_S(M)
 \longrightarrow H^1(G_{F,S},M^D)^\vee\\
 &\longrightarrow H^2(G_{F,S},M)
 \longrightarrow P^2_S(M)
 \longrightarrow H^0(G_{F,S},M^D)^\vee\longrightarrow0.
\end{aligned}
 \tag{LD.8}
\]
Here the maps from the local terms are the sums of the local
cup-product pairings followed by the \emph{same} $j_A$ just fixed.
The source's middle pairing reverses the degree-one factors, which
multiplies that map by $-1$ relative to the opposite convention.
Changing the common invariant sign also multiplies the transgression
maps by $-1$.  Such changes preserve their images and kernels, and
the connecting maps may be given the corresponding signs.  Thus
the exact sequence has precisely the source's
\texttt{loc0,loc1,loc2} and \texttt{tr0,tr1,tr2}; it asserts the
existence of its two $R$-linear connecting maps.  This is
\texttt{Nine\allowbreak Term\allowbreak Exact}, not merely equality of orders.

The classical local terms at infinity use Tate cohomology, including
degree zero.  At a complex place they vanish.  At a real place they
vanish on an odd-primary module because averaging over the group of
order two is exact and its norm is onto the invariants.  Under the
assumption in the proposition, all these terms therefore vanish
and (LD.8) has exactly the finite-place products in the source.
Ordinary $H^0$ at a real place must not be inserted here.

For $i=1,2$ let $\operatorname{Sha}^i_S(M)$ be the kernel of localization to the
finite places in $S$.  Part (a) of the same classical theorem gives
perfect pairings of finite groups
\[
 \operatorname{Sha}^1_S(M)\times\operatorname{Sha}^2_S(M^D)\longrightarrow\mathbb Q/\mathbb Z,
 \qquad
 \operatorname{Sha}^1_S(M^D)\times\operatorname{Sha}^2_S(M)\longrightarrow\mathbb Q/\mathbb Z.
 \tag{LD.9}
\]
They take values in $\ell^{-r}\mathbb Z/\mathbb Z$ and are
$R$-bilinear.  The finite dual-module calculation in LD2 converts
them into the two $R$-linear equivalences and equal cardinalities
specified by \texttt{Sha\allowbreak Dual}.  Part (c) identifies higher global
cohomology with the real-place groups; those vanish by the same
argument, so $H^n(G_{F,S},M)=0$ for all $n\ge3$.
This verifies every clause and the order of every quantifier in
\texttt{Poitou\allowbreak Tate\allowbreak Duality}.
\end{proof}

\subsection{The finite reciprocity consequence actually used}

\begin{proposition}[LD4: global cup reciprocity and one exceptional place]
\label{prop:ld-reciprocity}
Retain LD3's original field, finite set $S$ and local invariant
family.  For every $z\in H^2(G_{F,S},R(1))$,
\[
                    \sum_{v\in S}j_{A_v}(z|_{D_{A_v}})=0.
                    \tag{LD.10}
\]
Consequently, if all but one localization of $z$ vanish, the remaining
one vanishes too.  In particular this applies to the cup product of
two global finite cocycles with equivariant pairing into $R(1)$.
At a place away from $\ell$ where the coefficient representations
are unramified, two unramified local classes contribute zero; and
orthogonality to every unramified class in the Tate-dual module
forces a class itself to be unramified.
\end{proposition}

\begin{proof}
Apply (LD.8) to $M=R(1)$.  Its dual is the trivial rank-one module
$R$, and the identity functional is the global invariant $1$ in
$H^0(G_{F,S},R)$.  Exactness at $P^2_S(R(1))$ says that the image of
any global $z$ is killed by the last transgression.  Evaluating that
functional at $1$ gives exactly (LD.10).  The termwise description
uses evaluation, so no trace factor or degree appears.  If only
one term remains, its local invariant is zero and its class is
zero because $j_{A_v}$ is bijective.  Naturality of cup products
under restriction gives the assertion for global cocycles.  A local
coboundary contributes zero to cohomology; a pair of unramified
classes contributes zero by LD2.  The last assertion is the reverse
annihilator implication of LD2.

These arguments can be made at any finite coefficient level before
taking a limit.  A continuous action on a finite module and finitely
many continuous cocycles factor through a finite Galois quotient:
local constancy on the compact profinite domain supplies a common
open normal subgroup.  Only finitely many finite places ramify in
that extension.  Adding them and the places above $\ell$ to $S$
therefore puts all these objects in the same $G_{F,S}$.  Enlarging
$S$ does not change the already chosen invariants and merely adds
zero unramified contributions at the new places where appropriate.
This is the global/single-exception input used by the source's
\texttt{Over\allowbreak Iff.\allowbreak glob} and its finite reciprocity construction.
\end{proof}

\paragraph{Classical references and the exact closed boundary.}
The primary reference is J.~S.~Milne,
\href{https://jmilne.org/math/Books/ADTnot.pdf}{\emph{Arithmetic Duality
Theorems}, second edition (2006)}: I, Corollary~2.3 for finite local
Tate duality, I, Theorem~2.6 for the classical unramified-annihilator
statement, I, Lemma~2.9 for the prime-to-residue-characteristic Euler
identity, and I, Theorem~4.10 for finite Poitou--Tate duality.
LD1--LD2 provide direct proofs of the latter two local deductions
from the stated local duality theorem.  The Kummer/unramified
cyclic-algebra invariant is the standard local class field
normalization; the order of its factors and its sign were fixed
explicitly above.  For this normalization see Milne,
\href{https://www.jmilne.org/math/CourseNotes/CFT.pdf}{\emph{Class Field
Theory}, version 4.03 (2020)}, III, Proposition~3.6, Lemma~3.7,
and the cup-product construction in III, Section~4, Steps~1--3.
The author's
\href{https://www.jmilne.org/math/Books/add/ADT2006.pdf}{published
errata} and the exact finite-module hypotheses were checked.
These citations are established general duality results, not
unexpanded bespoke special-value bounds.

The supplied source endpoints are \texttt{Unramified\allowbreak Orthogonality}
and \texttt{Published.\allowbreak Poitou\allowbreak Tate\allowbreak Duality}, with their full original-field,
finite-module and local-choice quantifiers.  LD4 records their finite
reciprocity consequence.  The Jacobi-sum/Kolyvagin and generalized
Bernoulli argument needed for \texttt{Psi\allowbreak Bound} is a separate
arithmetic step; none of (LD.2)--(LD.10) asserts that special-value
bound on its own.


% Version 0.24. Complete finite supplier for the literal PsiBound interface.
% Source: commit 7cbf7a0c32fca1badf05931879a37c3f8a2b03b4.
\section{The finite arithmetic supplier for the odd-minus bound}
\label{sec:odd-minus-supplier}

\status{This section proves the literal \texttt{Psi\allowbreak Bound} interface by a
finite Jacobi-sum and Kolyvagin argument.  It includes the construction of
the auxiliary classes, the selection of auxiliary primes, the finite
length descent, the Bernoulli normalization, the residual exceptional
case, and descent to the original character eigenspace.  The global
reciprocity input is the finite duality statement reconstructed in
LD2--LD4.  The remaining classical arithmetic inputs are stated explicitly
below.  No Iwasawa main conjecture or unproved Euler-system existence
statement is used in this supplier.}

\subsection{Conventions and classical inputs}

Fix an odd prime $p$, a finite extension of $\mathbb Q_p$ with ring of
integers $O$, fraction field $K$, uniformizer $\varpi$, and residue field
$k$.  The specified embedding $O\subset\overline{\mathbb Q}_p$ is kept
throughout.  Write $A=K/O$ and normalize $v=v_O$ by $v(\varpi)=1$,
with $v(0)=\infty$.  Thus $A[a]\simeq O/(a)$ and
$\operatorname{length}_O A[a]=v(a)$ for $a\ne0$.
For the source's finite DVR subalgebra $O$, this is literally its
coefficient module \texttt{QO}: the localization
$\mathbb Q_p\otimes_{\mathbb Z_p}O$ is the fraction field $K$,
and the specified quotient is $K/O$.  Indeed $O$ is a finite
torsion-free $\mathbb Z_p$-algebra and a domain, so this localization
is a finite-dimensional domain over $\mathbb Q_p$, hence a field.
Lengths of arbitrary modules initially belong to
$\mathbb N\cup\{\infty\}$.  Let $G=\operatorname{Gal}(\Omega/\mathbb Q)$
for the specified algebraic closure, and let $\chi_p:G\to O^\times$ be
the cyclotomic character.  Decomposition and inertia groups at a chosen
place above a rational prime $q$ are denoted by $D_q$ and $I_q$; whenever
several choices matter, the assertion is required at every such place.

For a continuous character $\alpha:G\to O^\times$, put
\[
 \mathscr S(\alpha)=\ker\!\left(
 H^1(G,A(\alpha))\longrightarrow
 H^1(D_p,A(\alpha))\times\prod_q H^1(I_q,A(\alpha))\right).
 \tag{OS.1}
\]
Here $A(\alpha)$ has action $g x=\alpha(g)x$.  The product notation
expresses simultaneous local triviality, and does not assert that only
finitely many local groups are nonzero.  This is the source's
$\operatorname{ShaU}$ over $\mathbb Q$.

We use the following classical theorems in their ordinary finite forms.
They are foundational inputs, rather than newly reconstructed claims.
Cyclotomic theory identifies $\operatorname{Gal}(\mathbb Q(\zeta_n)/\mathbb Q)$
with $(\mathbb Z/n)^\times$, identifies the integers with
$\mathbb Z[\zeta_n]$, gives ramification only at primes dividing $n$,
and identifies Frobenius at $q\nmid n$ with $\zeta_n\mapsto\zeta_n^q$.
At a prime $q$, roots of unity of order prime to $q$ reduce injectively;
adjoining an $M$th root of a local unit is unramified when $q\nmid M$.
The latter follows also directly from Hensel's lemma over the maximal
unramified extension, whose residue field is algebraically closed.
Chebotarev supplies Frobenius elements in every nonempty open conjugacy
class condition, outside a prescribed finite set.  Together with the
Hilbert class field and its Artin isomorphism, it supplies a degree-one
prime in every ideal class of a number field, outside a prescribed finite
set.  Indeed degree $\ge2$ primes have Dirichlet density zero: their
norm-reciprocal series converges at $s=1$, being bounded by a fixed
multiple of $\sum_r r^{-2}$ over rational primes.  The desired Artin
class has positive Dirichlet density.  Finally, Kummer theory gives
$H^1(\mathbb Q,\mu_p)=\mathbb Q^\times/\mathbb Q^{\times p}$, with its
usual compatible local identifications; this is the Kummer exact
sequence and Hilbert's theorem 90.  The exact global/local finite
duality input is isolated in OS5 and supplied by LD2--LD4.

The source map at the end of this component records pinned declarations
for each use.  In particular, these named classical inputs do not stand
for the stronger assertion that the auxiliary classes constructed below
already exist or satisfy the required recurrence.

\subsection{Jacobi elements and their two distribution laws}

Write $K_n=\mathbb Q(\zeta_n)\subset\Omega$ and
$\sigma_t(\zeta_n)=\zeta_n^t$.  For $1\le b<n$ define
\[
 \theta_{n,b}(t)=
 \left\lfloor\frac{b\langle t\rangle_n}{n}\right\rfloor-
 \left\lfloor\frac{b\langle-t\rangle_n}{n}\right\rfloor,
 \qquad t\in(\mathbb Z/n)^\times.
 \tag{OS.2}
\]
Thus $\theta_{n,b}(-t)=-\theta_{n,b}(t)$.
For an ideal $\mathfrak a\ne0$ of $K_n$, call $X$ an associated
Jacobi element if
\[
 X\overline X=1,\qquad
 (X)=\prod_{t\in(\mathbb Z/n)^\times}
       \sigma_t^{-1}\mathfrak a^{\,\theta_{n,b}(t)}.
 \tag{OS.3}
\]
The inverse in this formula fixes the convention used in all subsequent
Frobenius and character formulas.

% BEGIN OS1
\begin{proposition}[OS1: existence, norm, and congruence of Jacobi elements]
For every nonzero ideal $\mathfrak a$ an element satisfying (OS.3)
exists, and any two differ by a root of unity whose order divides $2n$.
Let $q\nmid n$ be prime, $b<q$, and let $X'$ and $X$ satisfy (OS.3)
at levels $nq$ and $n$, for $\mathfrak a'$ and
$N_{K_{nq}/K_n}\mathfrak a'$, respectively.  Then
\[
 N_{K_{nq}/K_n}X'
   =\zeta\,X/\sigma_q^{-1}X,\qquad \zeta^{2n}=1.
 \tag{OS.4}
\]
If $\mathfrak a'$ is prime to $q$, there is a single root of unity
$\xi$, with $\xi^{2nq}=1$, such that, at every prime above $q$,
\[
                  X'\equiv\xi\,\sigma_q^{-1}X\pmod q.
 \tag{OS.5}
\]
Congruences here mean equality in the local residue field; all the
elements being compared are local units.
\end{proposition}
\begin{proof}
First suppose $\mathfrak a=P$ has prime norm $r\nmid n$.
Then $r\equiv1\pmod n$, and reduction at $P$ identifies
$\mu_n$ with the subgroup of order $n$ in $\mathbb F_r^\times$.
Let $d=(r-1)/n$ and let $\eta:\mathbb F_r^\times\to\mu_n$ be
the character reducing to $x\mapsto x^d$, extended by zero at zero.
Set
\[
 J=\prod_{j=1}^{b-1}J(\eta,\eta^j),\quad
 J(\eta,\eta^j)=\sum_{x\in\mathbb F_r}\eta(x)\eta^j(1-x),
 \qquad X=\overline J/J.
 \tag{OS.6}
\]
The standard finite-sum change of variables gives
$J(\eta,\eta^j)J(\eta^{-1},\eta^{-j})=r$ when
$\eta,\eta^j,\eta^{j+1}$ are nontrivial.  One way to check the
identity is to introduce a nontrivial additive character, expand the
two Gauss sums, and substitute $x+y=u$; this gives
$g(\eta)g(\eta^j)=J(\eta,\eta^j)g(\eta^{j+1})$ and
$g(\lambda)g(\lambda^{-1})=\lambda(-1)r$.
Consequently each Jacobi sum is nonzero and its integral ideal divides
$(r)$.

Put $u=\langle t\rangle_n$ and $v=\langle jt\rangle_n$.
Reduction at $\sigma_t^{-1}P$ gives
\[
 J(\eta,\eta^j)\equiv
       \sum_{x\in\mathbb F_r}x^{du}(1-x)^{dv}.
\]
For $0<i,j<r-1$, the sum $\sum_x x^i(1-x)^j$ vanishes exactly
when $i+j<r-1$.  To see this, expand $(1-x)^j$; the only possible
nonzero power sum has exponent $r-1$, and its coefficient is
$-(-1)^{r-1-i}\binom{j}{r-1-i}$, which is nonzero modulo $r$ when
$0\le r-1-i\le j<r$.  Thus the exponent of
$\sigma_t^{-1}P$ in the Jacobi sum is
$1-\operatorname{carry}(u,v)$, where the carry is $0$ or $1$
according as $u+v<n$ or $u+v\ge n$.
The prime $r$ splits without ramification, so the exponents really are
zero or one.  Telescoping carries gives
\[
 \operatorname{ord}_{\sigma_t^{-1}P}J
    =b-1-\left\lfloor\frac{b\langle t\rangle_n}{n}\right\rfloor.
\]
Subtracting this formula from its conjugate proves (OS.3) for
$\overline J/J$.

For a general $\mathfrak a$, choose a degree-one prime $P$, prime to
$n$, whose ideal class is $[\mathfrak a]^{-1}$, and write
$\mathfrak aP=(y)$.  The explicit element
\[
                 Y=\prod_t\sigma_t^{-1}(y)^{\theta_{n,b}(t)}
 \tag{OS.7}
\]
satisfies (OS.3) for $(y)$: the ideal identity is immediate and the
conjugate factors cancel in pairs.  Dividing $Y$ by the element (OS.6)
for $P$ proves existence for $\mathfrak a$.
The quotient of two choices is a unit all of whose complex conjugates
have absolute value one.  It is a root of unity: the monic polynomials
whose roots are the conjugates of its powers have integral, uniformly
bounded coefficients, hence only finitely many possibilities; some two
powers agree.  If a root of order $h$ lies in $K_n$, then
$K_{\operatorname{lcm}(n,h)}=K_n$, so their Euler totients are equal.
The formula for the totient permits no increase in a prime-power
factor, except adjoining a factor $2$ when $n$ is odd.  Thus
$h\mid2n$.  This proves
the stated uniqueness.

The norm identity rests on the following integer identity:
\[
 \sum_{\substack{t'\in(\mathbb Z/nq)^\times\\t'\equiv t\pmod n}}
 \theta_{nq,b}(t')
       =\theta_{n,b}(t)-\theta_{n,b}(q^{-1}t).
 \tag{OS.8}
\]
To verify it, first sum over the $q$ representatives $u+in$ of $t$,
$0\le i<q$.  Since $0\le b<q$, the usual floor-sum identity gives
the floor at level $n$, plus the constant
$\sum_{i=0}^{q-1}\lfloor ib/q\rfloor$.
Remove the unique representative divisible by $q$; its contribution
is the floor at $q^{-1}t$.  Subtract the identical computation for
$-t$; the constants cancel.  Apply (OS.8) to ideals and use norm
compatibility with conjugation.  The two sides of (OS.4), with the
root of unity omitted, then have the same ideal and conjugate product
one.  The uniqueness just proved supplies $\zeta$.

For (OS.5), first take $\mathfrak a'$ of degree one and norm $r$.
The characters $\eta'$ at level $nq$ and $\eta$ at level $n$ in
(OS.6) satisfy $(\eta')^q=\eta$ after inclusion in the larger field:
both sides have the same reduction at the prime above $r$, and
reduction is injective on $\mu_{nq}$.  Also
$(\sigma_q^{-1}\eta)^q=\eta$.  In characteristic $q$, the Frobenius
map is injective, so $\eta'$ and $\sigma_q^{-1}\eta$ have the same
reduction at every prime above $q$.  The Jacobi sums, their products,
and their conjugates therefore have the same reductions.  Their
denominators are units there, since they divide powers of $r\ne q$.
This proves (OS.5) for the canonical choices, and uniqueness inserts
one root of unity for any other choices.

For a principal ideal $(y)$ prime to $q$, use (OS.7).
Reduction of $\sigma_{t'}^{-1}y$ depends only on $t'\bmod n$,
because $\zeta_q$ reduces to one.  Call this value $f(t)$.
Frobenius gives $f(t)=f(qt)^q$.  Combining this with (OS.8), the
reduction of the upper-level product is
\[
 \prod_t f(t)^{\theta(t)-\theta(q^{-1}t)}
   =\prod_t f(qt)^{(q-1)\theta(t)}.
\]
Every one of the $q-1$ relative conjugates of $y$ has the same
reduction.  The last expression is consequently the reduction of
the lower-level product attached to $N(y)$, acted on by
$\sigma_q^{-1}$.  This proves (OS.5) for principal ideals.
Finally choose a degree-one prime $P'$ in $[\mathfrak a']^{-1}$,
outside $nq$, and apply the principal calculation to
$P'\mathfrak a'=(y)$.  Divide by the already proved calculation
for $P'$.  All factors are units at $q$, and the root of unity from
the degree-one calculation is global, so a single $\xi$ works at
every prime above $q$.  This proves (OS.5) in full.
\end{proof}
% END OS1

\subsection{Removing the roots of unity}

For the ensuing finite construction let $M=p^m$, $m\ge1$, and
$F=N_0M$ with $(N_0,p)=1$.  Suppose
\[
 \alpha|_{U}\equiv1\pmod M,\quad
 \alpha(j)=-1,\quad
 \alpha(g_0)-\chi_p(g_0)\in O^\times,
 \qquad U=\operatorname{Gal}(\Omega/K_F),
 \tag{OS.9}
\]
where $j$ is a complex conjugation.  Put
$\beta=\chi_p\alpha^{-1}$.  Thus $\beta(j)=1$,
$\beta|_U\equiv1\pmod M$, and some $\beta(g)-1$ is a unit.
Choose an integer $a$ such that $g_0$ acts on $\mu_M$ by the
$a$th power, and put
\[
                   P(y)=g_0(y)/y^a.
 \tag{OS.10}
\]
Then $\alpha(g_0)-a$ is a unit, since $M$ is in the maximal ideal.

Choose a prime $r\equiv1\pmod F$ and one place $A_0$ above $r$.
Call a prime $q$ admissible if
\[
 q\equiv1\pmod F,\qquad q\ne r,\qquad
              p\nmid\operatorname{ord}_{\mathbb F_q^\times}(r).
 \tag{OS.11}
\]
An admissible list has distinct admissible entries, and $d$ denotes
their product.  All cyclotomic fields below are the actual subfields
of the chosen $\Omega$.

% BEGIN OS2
\begin{proposition}[OS2: a strict finite Euler system]
For $1\le b<F$ there are nonzero elements $x_d\in K_{Fd}$,
units outside $r$, and positive integers $s(d,q)$ prime to $p$,
such that for every admissible extension $d\mapsto dq$,
\begin{align*}
 N_{K_{Fdq}/K_{Fd}}x_{dq}
       &=(x_d/\sigma_q^{-1}x_d)^{s(d,q)},\\
 x_{dq}&\equiv(\sigma_q^{-1}x_d)^{s(d,q)}\pmod q.
\end{align*}
The second equality holds at every place above $q$.
The elements can be chosen in the form
$x_d=P(X_d)^{2N_0d}$, where $X_1$ has
\[
 \operatorname{ord}_{A_0}(\sigma_tX_1)
                 =\theta_{F,b}(t)\operatorname{ord}_{A_0}(r).
 \tag{OS.12}
\]
\end{proposition}
\begin{proof}
At each level take the prime below the fixed place $A_0$, and use OS1
to choose its Jacobi element $X_d$.  These elements are units outside
$r$.  Relative norm of the upper prime is a power of the lower
prime, with exponent $s_0(d,q)$ dividing the order of $r$ modulo $q$.
Indeed residue degrees in cyclotomic fields are orders of $r$ modulo
their levels, and their relative ratio divides that order.  The
admissibility condition makes $s_0(d,q)$ prime to $p$.
OS1 therefore gives the displayed norm and congruence for $X_d$
with exponent $s_0(d,q)$ and possible root-of-unity factors of order
dividing $2Fdq$.  At level $F$, $r$ splits completely, so the ideal
identity (OS.3) gives precisely (OS.12).

The operator $P$ commutes with cyclotomic conjugations and with their
norms.  If $\zeta^{2Fdq}=1$, then
$\zeta^{2N_0dq}\in\mu_M$, and consequently
\[
                     P(\zeta)^{2N_0dq}=1.
\]
Apply $P$ and then the power $2N_0dq$ to both lax identities.
The root-of-unity factors disappear.  On the lower side this power
is $q$ times the exponent used in $x_d$, so the new exponent is
$s(d,q)=q s_0(d,q)$, still prime to $p$.
Congruences remain valid after applying $g_0$, because the original
congruence held at every prime above $q$, and after division and
integer powers because all relevant elements are local units.
This proves both strict identities simultaneously.
\end{proof}
% END OS2

\subsection{Derivative classes with exact local recurrence}

For every admissible $q$, choose $\tau_q$ acting trivially on all
roots of unity of order prime to $q$, and generating the action on
$\mu_q$.  Such an element exists in inertia at $q$.
In multiplicative notation put
\[
 D_q y=\prod_{i=0}^{q-2}\tau_q^i(y)^i,
 \qquad D_L=\prod_{q\in L}D_q,
 \qquad y_L=D_Lx_d.
 \tag{OS.13}
\]
The last product means composition of commuting group-ring operators
on the indicated cyclotomic field.  Fix a primitive $M$th root
$\zeta\in\Omega$, representatives $\sigma_t\in G$ of
$t\in(\mathbb Z/F)^\times$, and $w\in A[M]$.

% BEGIN OS3
\begin{proposition}[OS3: construction of the auxiliary homomorphisms]
There are continuous additive homomorphisms $\nu_L:U\to A[M]$
such that
\[
 \nu_L(gug^{-1})=\beta(g)\nu_L(u),\qquad
 \nu_L(I_q)=0\quad(q\nmid Frd).
 \tag{OS.14}
\]
If $q$ is admissible and absent from $L$, then for every place above
$q$ and every Frobenius $\phi_q$ there are
$\tau\in I_q$ and $u_q\in O^\times$ with
\[
                  \nu_{qL}(\tau)=u_q\nu_L(\phi_q).
 \tag{OS.15}
\]
At the fixed place above $r$ there are $\tau_0\in I_r$ and
$u_0\in O^\times$ with
\[
 \nu_\varnothing(\tau_0)=u_0zw,\qquad
 z=\sum_{t\in(\mathbb Z/F)^\times}
          \theta_{F,b}(t)\alpha(\sigma_t)^{-1}.
 \tag{OS.16}
\]
These are the full source conditions \texttt{Nu\allowbreak Sys}, for the same
$M,F,r,\alpha,A_0,\sigma$ and $w$.
\end{proposition}
\begin{proof}
The group-ring identity
\[
                 (\tau_q-1)D_q=(q-1)-N_q
 \tag{OS.17}
\]
and the norm formula of OS2 imply, by induction on $L$, that
$g y_L/y_L$ is an $M$th power in $K_{Fd}$ for every $g\in U$.
Here are the induction details.  For a new $q$ write
$w'=D_Lx_{dq}$ and $y=D_Lx_d$.  Then
$N_qw'=(y/\phi_q^{-1}y)^s$, where $s=s(d,q)$.
By induction choose $z'$ in $K_{Fd}$ with
$(z')^M=\phi_q^{-1}y/y$.  Since $q-1=Mh$, (OS.17) gives
\[
 \frac{\tau_q(D_qw')}{D_qw'}
          =\frac{(w')^{q-1}}{N_qw'}
          =((w')^h(z')^s)^M.
\]
Repeat this calculation for each member of the enlarged list, using
the derivative for the list with that member removed; the induction
hypothesis applies to that shorter list.  For elements fixing
$K_{Fdq}$ the quotient is one.
These elements generate $\operatorname{Gal}(K_{Fdq}/K_F)$, by the
Chinese remainder theorem.  The property is closed under products
and inverses, proving the induction.

Choose $Y_L^M=y_L$.  For $g\in U$ choose $z_g\in K_{Fd}$ with
$z_g^M=g y_L/y_L$, and define $\kappa_L(g)\in\mathbb Z/M$ by
\[
 \zeta^{\kappa_L(g)}=
       \frac{g(PY_L)}{PY_L}\,\frac1{Pz_g}.
 \tag{OS.18}
\]
The $M$th power of the right side is one, because $g_0$ and $g$
commute on $K_{Fd}$.  Changing $z_g$ multiplies it by an element of
$\mu_M$, which $P$ kills; changing $Y_L$ has the same harmless
effect.  For $g,h\in U$, one can take $z_{gh}=g(z_h)z_g$.
Substitution in (OS.18), using the same commutation on $K_{Fd}$,
proves additivity.  Its kernel contains the open subgroup fixing
$K_{Fd}$ and $PY_L$, so $\kappa_L$ is continuous.
At a prime outside $Frd$, the field $K_{Fd}$ is unramified and
$y_L$ is a unit.  Its $M$th roots are unramified there, as are their
$g_0$-conjugates.  Inertia therefore gives zero in (OS.18).

Define
\[
 \nu_L(u)=\sum_t\kappa_L(\sigma_t^{-1}u\sigma_t)
                              \beta(\sigma_t)w.
 \tag{OS.19}
\]
Additivity, continuity, $M$-torsion and the stated unramifiedness are
immediate.  To check equivariance, reindex the cosets by $g$.
The extra factors are in $U$; an additive homomorphism on $U$ is
unchanged by conjugation inside $U$, and $\beta|_U$ acts trivially
on $A[M]$.  The remaining factor is $\beta(g)$, proving (OS.14).

For completeness the local recurrence in (OS.15) is an exact residue
calculation, not an additional Euler-system axiom.  Choose inertia
$\tau$ with the same action on $\mu_q$ as $\tau_q$.
Put $w'=D_Lx_{dq}$, $y=D_Lx_d$, $Y^M=D_qw'$ and $(Y')^M=y$.
For $\phi=\phi_q$ choose $z'$ with $(z')^M=\phi^{-1}y/y$.
The previous induction shows that $(w')^h(z')^s$ is an allowable
root for $\tau$, where $Mh=q-1$.
Substituting these roots in (OS.18) and cancelling factors gives
\[
 \frac{\zeta^{\kappa_{qL}(\tau)}}
      {(\zeta^{\kappa_L(\phi^{-1})})^s}
 =\frac{\tau(PY)}{PY}
   \left(\frac{PY'}{\phi^{-1}(PY')^q}\right)^s
   \left[P\!\left(\frac{w'}{(\phi^{-1}y)^s}\right)\right]^{-h}.
 \tag{OS.20}
\]
The first factor reduces to one because $\tau$ is inertia; the
second does so because $\phi$ is Frobenius; the third does so by
the congruence in OS2 and its images under $D_L$ and $P$.
All elements involved are units at $q$.
Both sides on the left belong to $\mu_M$, whose reduction is
injective since $q\ne p$.  Hence
$\kappa_{qL}(\tau)=-s\kappa_L(\phi)$ in $\mathbb Z/M$.
The same identity holds after conjugating the chosen place by each
$\sigma_t$: cyclotomic automorphisms commute, so the action of
the conjugated inertia element on $\mu_q$ is still the required
generator.  Sum as in (OS.19).  This gives (OS.15) with
$u_q=-s$, a unit.

For initialization, if $y\in K_F$ and
$\operatorname{ord}_{A_0}(\sigma_tPy)=k(t)
\operatorname{ord}_{A_0}(r)$, choose $\rho^M=r$ and tame inertia
$\tau_0$ with $\tau_0\rho=\zeta\rho$.
The quotient $\sigma_t(PY)/\rho^{k(t)}$ is a local unit whose
$M$th power is a unit, so it is fixed by inertia.  Formula (OS.18)
therefore gives
\[
       t\,\kappa(\sigma_t^{-1}\tau_0\sigma_t)=k(t)
                         \quad\hbox{in }\mathbb Z/M.
\]
Since $\chi_p(\sigma_t)\equiv t\pmod M$, its weighted sum is
$\sum_t k(t)\alpha(\sigma_t)^{-1}w$.
For $y=x_1=P(X_1)^{2N_0}$, write $d_0$ for the image of $g_0$
in $(\mathbb Z/F)^\times$.  By (OS.12),
\[
 k(t)=2N_0\bigl(\theta(td_0^2)-2a\theta(td_0)+a^2\theta(t)\bigr).
\]
Reindexing the sum, and using $\alpha|_U\equiv1\pmod M$, gives
\[
 \sum_t k(t)\alpha(\sigma_t)^{-1}w
       =2N_0(\alpha(g_0)-a)^2zw.
\]
The coefficient is a unit by (OS.9) and $p\ne2$.
This proves (OS.16) and the entire assertion.
\end{proof}
% END OS3

\subsection{Compatible auxiliary primes}

% BEGIN OS4
\begin{proposition}[OS4: one prime with all the required properties]
Assume (OS.9).  Let $c:U\to A[M]$ be a nonzero continuous
$\alpha$-equivariant homomorphism, and let
$\nu:U\to A[M]$ be a continuous $\beta$-equivariant homomorphism.
For every finite set of excluded primes, there are an admissible prime
$q$ outside that set, a place above $q$, and its Frobenius $\phi$, such
that
\[
       c(\phi)\ne0,\qquad
       O\nu(\phi)=\sum_{u\in U}O\nu(u).
 \tag{OS.21}
\]
The same prime can in addition be required to satisfy
$q\not\equiv1\pmod{Fp}$.
\end{proposition}
\begin{proof}
Choose $\rho^M=r$ and put $H=\operatorname{Gal}(\Omega/K_F(\rho))$.
The extension $K_F(\rho)/K_F$ is cyclic of degree $M$:
$X^M-r$ is Eisenstein at a prime above $r$, and $\mu_M\subset K_F$.
It is Galois over $\mathbb Q$, because every conjugate of $\rho$
is $\zeta^i\rho$.  On its relative cyclic group, conjugation by
$g\in G$ acts by the scalar $\chi_p(g)$ modulo $M$.

Let $B=\ker c$.  If $H\subset B$, then $c$ factors through this
cyclic Kummer group.  Comparing its $\alpha$-equivariance with
the preceding conjugation action gives
$(\alpha(g)-\chi_p(g))c(u)=0$ for every $g,u$.
At $g_0$ the coefficient is a unit, contrary to $c\ne0$.
Thus $H\not\subset B$.

The $O$-module $C=\sum O\nu(U)$ is finite cyclic: $\nu$ has finite
image by compactness, and every finite submodule of $A$ is one of
the modules $A[\varpi^a]$.
If $C\ne0$, let $B'=\{u:\nu(u)\in\varpi C\}$.
If $H\subset B'$, the nonzero homomorphism
$U\to C/\varpi C$ would again factor through the Kummer group.
Its two conjugation scalars would be $\beta(j)=1$ and
$\chi_p(j)=-1$.  Their difference is a unit because $p$ is odd,
a contradiction.  Thus $H\not\subset B'$.
If $C=0$, simply set $B'=B$; the second condition of (OS.21) is
then automatic.

Put $B''=\operatorname{Gal}(\Omega/K_{Fp})$, a subgroup of $U$.
We also have $H\not\subset B''$.
Otherwise $K_{Fp}\subset K_F(\rho)$.  The relative extension
$K_{Fp}/K_F$ has degree $p$ and is abelian over $\mathbb Q$.
Complex conjugation would therefore act trivially on its relative
Galois group.  But, as a quotient of the Kummer group, that action
is inversion.  A nontrivial group of odd order $p$ cannot have
both actions.  This proves the third noncontainment.

Here is the small group-theoretic compatibility argument.
A group cannot be the union of three proper subgroups $A_1,B_1,C_1$
when $x^M\in B_1$ for every $x$ and $M$ is odd.
Two proper subgroups cannot cover a group: the product of an element
in each exclusive part lies outside both.
If a three-subgroup cover existed, choose
$x\notin B_1\cup C_1$ and $y\notin A_1\cup C_1$ using this
two-subgroup fact.  Then $x\in A_1$, $y\in B_1$, and $xy\in C_1$.
The element $x^2y$ is in neither $A_1$ nor $C_1$, so belongs to
$B_1$.  Hence $x^2\in B_1$.  Combining this with $x^M\in B_1$
and oddness of $M$ gives $x\in B_1$, a contradiction.
Apply the argument inside $H$ to its intersections with $B',B,B''$.
The required power condition follows because $c(x^M)=Mc(x)=0$.
Hence there is an element
\[
                         h\in H\setminus(B\cup B'\cup B'').
\]
All three subgroups are open in $U$; thus this is a nonempty open
condition in $G$.  Pass to a finite Galois quotient small enough
to preserve it.  Chebotarev supplies a Frobenius conjugate in this
open set outside the prescribed primes.  Changing the place above
that rational prime selects the indicated conjugate itself.
This explains compatibility of all conditions in the compositum;
no independent, possibly conflicting congruence choices have been
made.

The selected Frobenius lies in $U\setminus B''$, so
$q\equiv1\pmod F$ but $q\not\equiv1\pmod{Fp}$.
Write $q-1=N_0M h_0$; then $p\nmid h_0$.
It fixes $\rho$, so its residue satisfies $\rho^{q-1}=1$,
and consequently $r^{N_0h_0}=1$ modulo $q$.
Therefore $\operatorname{ord}_{\mathbb F_q^\times}(r)$ divides
$N_0h_0$, which is prime to $p$.  Excluding $r$ makes $q$
admissible.  Avoidance of $B$ gives the first assertion of (OS.21),
and avoidance of $B'$ says that $\nu(\phi)$ is a generator of
the cyclic module $C$.  This proves every required condition.
\end{proof}
% END OS4

\subsection{The precise reciprocity consequence}

% BEGIN OS5
\begin{proposition}[OS5: finite reciprocity and the annihilator comparison]
The duality results LD2--LD4 imply the source condition
\texttt{Rec}.  In particular, let $a_0\mid a\mid a'$ with
$a'\mid M$, and let $c:G\to A$ be an $\alpha$-cocycle killed by
$a_0$.  Assume $c$ is zero on inertia away $F$, is a coboundary
on decomposition groups above $F$ with witnesses killed by $a_0$,
and is zero on the decomposition groups above the primes dividing
$n$ outside $F$.
Suppose $q\equiv1\pmod F$, $q\nmid Fn$, and
$\nu:U\to A[M]$ is $\beta$-equivariant, unramified outside $Fnq$,
with
\[
 \operatorname{Ann}_O\nu(U)=(a'),\qquad
                  \operatorname{Ann}_O\nu(\tau)=(a)
 \tag{OS.22}
\]
for some $\tau\in I_q$.
Then $a'=a d$ for some $d\in O$ and
\[
                         d\,c(\phi_q)=0.
 \tag{OS.23}
\]
\end{proposition}
\begin{proof}
We first spell out the finite coefficient reduction, so that the
reference to LD4 preserves the literal \texttt{Rec} quantifiers.
For any nonzero annihilator $a'$ choose $w_0\in A$ with
$\operatorname{Ann}(w_0)=(a')$ and identify $A[a']$ with
$B=O/(a')$ by $b\mapsto bw_0$.
Choose $p^h$ divisible by $a'$ and let
$\lambda:B\to\mathbb Z/p^h$ be any
$\mathbb Z/p^h$-linear map.
For an $\alpha$-cocycle $c$ and a
$\beta=\chi_p\alpha^{-1}$-cocycle $n$, the function
\[
             (g,h)\longmapsto
                 \lambda\bigl(n(g)\alpha(g)c(h)\bigr)
 \tag{OS.24}
\]
is their cup-product cocycle in $(\mathbb Z/p^h)(1)$.
Use the assumed continuity of the reduction of $\alpha$ modulo $a'$;
this is exactly the open-kernel hypothesis in \texttt{Rec}.
No continuity of an unreduced character is needed at this step.
At primes outside the finite ramification
set this cup class is unramified; at every specified prime where
$c$ is a coboundary, it is a coboundary.  For example if
$c(g)=(\alpha(g)-1)m$, a cobounding cochain is
$-\lambda(n(g)\alpha(g)m)$.
At the single remaining prime $q$, LD4 first makes the local cup
class zero.  The following argument converts that fact into the
elementwise assertion actually needed.  Let
$x_0:D_q/I_q\to R=\mathbb Z/p^h$ send arithmetic Frobenius to $1$.
Since $c$ is unramified and the characters are trivial on $B$,
$c(g)=x_0(g)c(\phi_q)$ in $B$.
The function $N(g)=\lambda(n(g)c(\phi_q))$ is a cocycle in $R(1)$:
on these values the action of $\beta$ agrees with that of $\chi_p$,
as both are trivial modulo $a'$ on $D_q$.
Thus the local cup (OS.24) is $N\cup x_0$.
Every unramified $R$-valued class is an $R$-multiple of $x_0$,
so $N$ pairs trivially with every such class.
LD2, applied to the trivial module $R$ and its dual $R(1)$, makes
$N$ unramified.  Inertia acts trivially on $R(1)$ at $q\ne p$,
so an unramified class has $N(\tau)=0$ as an actual cocycle value.
Multiplying by $x_0(g)$ now gives
\[
                  \lambda\bigl(n(\tau)c(g)\bigr)=0
                 \quad(\tau\in I_q,\ g\in D_q),
 \tag{OS.25}
\]
under exactly the stated trivial-character and unramified conditions.
This deduction uses the same finite modules and the same place choices
throughout.
All the maps $\lambda$ together separate points of $B$:
decompose its finite additive $p$-group into cyclic factors and
map a nonzero coordinate into $\mathbb Z/p^h$.
Thus (OS.25) holds before applying $\lambda$ as well.
Equivalently, $n(\tau)=yw_0$ implies $y c(g)=0$ in $A$.
This is exactly \texttt{Rec}; no trace pairing of potentially
nonunit discriminant has been substituted for these separating maps.

To apply it to $\nu$, choose $g_1$ with $\beta(g_1)-1$ a unit.
The quotient $G/U$ is abelian, so
\[
 n(g)=-(\beta(g_1)-1)^{-1}
                \nu(gg_1g^{-1}g_1^{-1})
 \tag{OS.26}
\]
is defined.  The commutator identity and $\beta$-equivariance give
$n(gh)=n(g)+\beta(g)n(h)$, and for $u\in U$ give $n(u)=\nu(u)$.
Its values are $O$-multiples of values of $\nu$, so $a'n=0$.
It is continuous, and it has exactly the unramifiedness needed
outside $Fnq$, since inertia there lies in $U$.
At $q$, $D_q\subset U$ because $q$ splits in $K_F$; hence
$\alpha,\beta\equiv1\pmod{a'}$.  The hypotheses on $c$
provide the other local conditions in \texttt{Rec}.
We obtain $y c(\phi_q)=0$ if $\nu(\tau)=yw_0$.

The inclusion $\nu(\tau)\in\nu(U)$ shows $a\mid a'$.
In a DVR,
\[
 v(a)=\max\{v(a')-v(y),0\},\qquad
 v(d)=v(a')-v(a)=\min\{v(y),v(a')\}.
\]
If $v(y)\le v(a')$, then $d$ is a multiple of $y$ up to a unit,
so it kills $c(\phi_q)$.  If $v(y)\ge v(a')$, then $a$ is a
unit, and $a_0\mid a$ makes $c=0$.
This proves (OS.23).
\end{proof}
% END OS5

\subsection{The finite descent actually bounds length}

For $\alpha(j)=-1$, let $\mathcal T(F,a_0)$ consist of continuous
$\alpha$-cocycles $c:G\to A$ with $c(j)=0$, killed by $a_0$,
zero on inertia away $F$, and locally coboundaries at primes
dividing $F$ with witnesses killed by $a_0$.
This is the exact normalized module \texttt{Sh\allowbreak T}.
Let $\mathcal T_n$ be its submodule whose cocycles vanish on every
decomposition group above a prime dividing $n$.

% BEGIN OS6
\begin{proposition}[OS6: the complete finite length descent]
Under (OS.9), and with $z$ as in (OS.16), if
$a_0z\mid M$, then
\[
              \operatorname{length}_O\mathcal T(F,a_0)\le v(z).
 \tag{OS.27}
\]
No finiteness assumption on the module on the left is required.
\end{proposition}
\begin{proof}
The divisibility implies $z\ne0$; write $M=k_0z$, so $a_0\mid k_0$.
Choose $w\in A$ with annihilator $(M)$, and construct $\nu_L$
by OS3.  Let $(a_L)=\operatorname{Ann}\nu_L(U)$.
Every $a_L$ divides $M$, and $a_L\ne0$.
Initialization gives $\operatorname{Ann}\nu_\varnothing(\tau_0)=(k_0)$,
so $a_\varnothing=k_0d_0$.

We will repeatedly use the following length comparison, which is
stronger than an elementwise annihilation statement.
If $\nu$ has annihilator $(a')$ and $\nu(\tau)$ has annihilator
$(a)$, with $a_0\mid a$ and $a'=ad$, OS5 shows that evaluation
\[
             \mathcal T_n\longrightarrow A,\qquad
                       c\longmapsto c(\phi_q)
\]
has image in $A[d]$, of length at most $v(d)$.
Its kernel is $\mathcal T_{nq}$.
Indeed inertia at $q$ acts trivially on the values and $c$ vanishes
there; the unramified quotient is procyclic, so a continuous
homomorphism zero at Frobenius vanishes on $D_q$.
The same conclusion holds at every conjugate place: the cocycle
conjugation formula changes the value only by
$(\alpha(g)-1)c(h)$ with $g\in D_q\subset U$, which is zero
because $a_0\mid M$ and $\alpha|_U\equiv1\pmod M$.
Thus additivity and monotonicity of extended module length give
\[
        \operatorname{length}\mathcal T_n
       \le\operatorname{length}\mathcal T_{nq}+v(d).
 \tag{OS.28}
\]
Apply this first with $n=1$, $q=r$, and the initialization data:
\[
 \operatorname{length}\mathcal T(F,a_0)
       \le\operatorname{length}\mathcal T_r+v(d_0).
 \tag{OS.29}
\]

For an admissible list $L$ with $a_0\mid a_L$, put
$S_L=\mathcal T_{r\prod L}$ and $J=v(M)$.
If $S_L\ne0$, choose a nonzero $c\in S_L$.
Its restriction to $U$ is an $\alpha$-equivariant homomorphism:
the action is trivial modulo $M$, and the cocycle conjugation
formula gives equivariance.  This restriction is nonzero.
For if it were zero, $c$ would factor through the abelian group
$G/U$; the commutativity relation for $j$ and $g$, together with
$c(j)=0$, would give $(\alpha(j)-1)c(g)=0$, hence $c=0$.

Use OS4 with this $c$ and $\nu_L$, excluding primes dividing
$Fr\prod L$.  The resulting $q$ has
$c(\phi_q)\ne0$ and $\nu_L(\phi_q)$ generates $\nu_L(U)$.
By OS3, $\nu_{qL}(\tau)$ therefore has annihilator $(a_L)$.
Write $a_{qL}=a_Ld$ and apply (OS.28):
\[
 \operatorname{length}S_L+v(a_L)
       \le\operatorname{length}S_{qL}+v(a_{qL}).
 \tag{OS.30}
\]
Moreover $v(d)>0$, because (OS.23) kills the nonzero value
$c(\phi_q)$; a unit cannot do so.  Thus $v(a_{qL})>v(a_L)$,
while both are at most the finite integer $J$.

Induction on $J-v(a_L)$ now proves
\[
                   \operatorname{length}S_L+v(a_L)\le J.
 \tag{OS.31}
\]
If $S_L=0$ this is immediate.  Otherwise (OS.30) reduces it to
a strictly smaller induction parameter.  This proof also establishes
finiteness; it has not subtracted an infinite length or presumed
that a nonzero element exhausts a module.
At $L=\varnothing$, substitute
$a_\varnothing=k_0d_0$ and $M=k_0z$ in (OS.31):
\[
                  \operatorname{length}\mathcal T_r+v(d_0)\le v(z).
\]
Only the finite integer $v(k_0)$ has been cancelled.
Combine this with (OS.29) to obtain (OS.27).
\end{proof}
% END OS6

\subsection{The Bernoulli value with the correct denominator and sign}

For a primitive Dirichlet character $\eta$ of conductor $f$, set
\[
 B_{k,\eta}=f^{k-1}\sum_{t=1}^{f}\eta(t)B_k(t/f),
 \qquad \operatorname{den}(N,k)=N\prod_{\substack{\ell\le k+1\\
                                    \ell\ {
m prime}}}\ell.
 \tag{OS.32}
\]
The generalized Bernoulli number in the source's
$\operatorname{Bprim}$ is always taken at the primitive conductor.
Thus for an induced character at level $N$ the source value is
$\operatorname{den}(N,k)B_{k,\eta}$; it has no additional Euler
factors at primes dividing $N/f$.

% BEGIN OS7
\begin{proposition}[OS7: finite theta evaluation]
Let $e\ge1$, $k=e+1$, and let $\eta$ be primitive of conductor
$f$, with values in $O$, satisfying
$\eta(-1)=(-1)^k$.  Put
\[
       \alpha=\chi_p^{-e}\eta^{-1},\qquad
       Y=B_{k,\eta}/k.
 \tag{OS.33}
\]
Suppose $Y\ne0$ and $\alpha(g_0)-\chi_p(g_0)$ is a unit for
some $g_0$.  For every $a_0\in O\setminus\{0\}$, there are
$m\ge1$, $F=f/p^{v_p(f)}\cdot p^m$, $1\le b<F$ and
representatives $\sigma_t$ such that (OS.9) holds and the element
$z$ in (OS.16) satisfies
\[
              v(z)=v(Y),\qquad a_0z\mid p^m.
 \tag{OS.34}
\]
In particular, if $f\mid N$ and
$a=\operatorname{den}(N,k)B_{k,\eta}\in O$, then $z\mid a$.
The field, the coefficient ring and the character are unchanged.
\end{proposition}
\begin{proof}
We use a simple denominator-clearing version of the source's
Bernoulli congruence.  Choose a nonzero integer $D=D_k$ clearing
the finitely many rational coefficients that occur below.
It depends only on $k$, and not on $F,b,t$ or the character.
Writing
\[
 b_{F,k}(t)=F^{k-1}B_k(t/F)
          =\frac{t^k}{F}-\frac{k}{2}t^{k-1}
             +\sum_{j=2}^k\binom{k}{j}B_jF^{j-1}t^{k-j},
\]
take $u=\langle t\rangle_F$, $y=\langle bt\rangle_F$,
and $h=\lfloor bu/F\rfloor$, with $b$ a unit modulo $F$.
Then $y=bu-Fh$ and
$\theta_{F,b}(t)=2h-(b-1)$.
Expand $y^k$ modulo $F^2$ and $y^{k-1}$ modulo $F$.
After the leading terms cancel, this gives
\[
 D\left[b^{k-1}\theta_{F,b}(t)u^{k-1}
       +2\left(\frac{b_{F,k}(y)}k
                   -b^k\frac{b_{F,k}(u)}k\right)\right]\in F\mathbb Z.
 \tag{OS.35}
\]
This calculation also explains why a denominator independent of
the increasing level suffices: all terms after the two displayed
leading Bernoulli terms contain a factor $F$.
The source uses the explicit choice
$D=2k\prod_{\ell\le k+1}\ell$; the argument here only needs a
fixed common denominator and therefore does not invoke a theorem
on Bernoulli denominators.

Let $\eta_F$ be $\eta$ induced to $F$ and put
\[
 S=\sum_{t\in(\mathbb Z/F)^\times}
          \theta_{F,b}(t)\langle t\rangle_F^{e}\eta_F(t).
\]
Multiply (OS.35) by $\eta_F(bt)$ and sum.
Reindexing the terms containing $y=bt$ gives
\[
 D\left[b^e\eta_F(b)S
       -2\bigl(\eta_F(b)b^k-1\bigr)B_{k,\eta_F}/k\right]
                                               \in p^m O.
 \tag{OS.36}
\]
The generalized Bernoulli numbers here are algebraic numbers in
$K$.  Adding only powers of $p$ to the conductor gives
\[
 B_{k,\eta_F}=
 \begin{cases}
 B_{k,\eta},&p\mid f,\\
 (1-\eta(p)p^{k-1})B_{k,\eta},&p\nmid f,
 \end{cases}
 \tag{OS.37}
\]
once $m\ge v_p(f)$ and $m\ge1$.
For a direct verification, use
$\sum_{i=0}^{q-1}B_k((x+i)/q)=q^{1-k}B_k(x)$,
which follows by summing the generating function
$te^{xt}/(e^t-1)$.  Summation over all lifts at level $qf$
preserves the Bernoulli value; if $q\nmid f$, removing the terms
divisible by $q$ subtracts $\eta(q)q^{k-1}B_{k,\eta}$.
Once $q$ divides the level there is no further removal.
This proves (OS.37).  Its displayed factor is a $p$-adic unit
because $k\ge2$.

Choose $m$ sufficiently large that $f\mid F$ and
\[
                  |p^m|_p<|a_0|_p\,|D|_p\,|Y|_p.
 \tag{OS.38}
\]
Choose $b$ to be the least positive representative of the action of
$g_0$ on $\mu_F$.  Then $b$ is a unit modulo $F$ and
$\eta_F(b)b^k-1$ is a unit in $O$:
modulo the maximal ideal it is associated to
$\alpha(g_0)-\chi_p(g_0)$, after multiplication by the unit
$-\chi_p(g_0)^e\eta(g_0)$.
The factors $b^e\eta_F(b)$ and $2$ are also units.
By (OS.36)--(OS.38) and the ultrametric inequality,
\[
                             |S|_p=|Y|_p.
\]
Indeed the error after division by $D$ is strictly smaller than
the norm of the nonzero term
$2(\eta_F(b)b^k-1)B_{k,\eta_F}/k$.

For any lift $\sigma_t$ of $t$ we have
$\alpha(\sigma_t)^{-1}=\chi_p(\sigma_t)^e\eta_F(t)$ and
$\chi_p(\sigma_t)\equiv\langle t\rangle_F\pmod{p^m}$.
Therefore $z-S\in p^mO$.
By (OS.38) again, $|z|_p=|S|_p=|Y|_p$.
In particular $Y$ has nonnegative valuation, as $z\in O$.
The same inequality implies $v(p^m)\ge v(a_0)+v(z)$, proving
the divisibility in (OS.34).
Also
\[
        a=\bigl(k\operatorname{den}(N,k)\bigr)Y
\]
has valuation at least $v(Y)=v(z)$, so $z\mid a$.
This last argument avoids assuming separately that the primitive
Bernoulli normalization is integral in $O$.

Finally $\alpha|_U\equiv1\pmod{p^m}$ since $f\mid F$,
and parity gives
$\alpha(j)=(-1)^{-e}(-1)^{-(e+1)}=-1$.
The remaining condition in (OS.9) was the hypothesis on $g_0$.
All constructions took place in the original $\Omega$ and $O$.
\end{proof}
% END OS7

\subsection{Normalized Selmer classes and the residual exception}

% BEGIN OS8
\begin{proposition}[OS8: the full cohomology module is controlled]
Let $\eta$ be a primitive Dirichlet character of conductor $f$,
of order prime to $p$, and put $\alpha=\chi_p^{-e}\eta^{-1}$,
with the parity and $e\ge1$ of OS7.
Every class of $\mathscr S(\alpha)$ has a unique cocycle representative
normalized by $c(j)=0$, and these representatives are the increasing
union of $\mathcal T(fp,\varpi^h)$, $h\ge0$.
Consequently a uniform bound $L$ for their lengths gives
$\operatorname{length}_O\mathscr S(\alpha)\le L$.
If
\[
                   \alpha(g)\equiv\chi_p(g)\pmod\varpi
                      \quad\hbox{for every }g\in G,
 \tag{OS.39}
\]
then $\mathscr S(\alpha)=0$.
\end{proposition}
\begin{proof}
Since $\alpha(j)-1=-2$ is a unit, subtracting the coboundary of
$(\alpha(j)-1)^{-1}c(j)$ gives a representative with $c(j)=0$.
Two such representatives differ by a coboundary
$(\alpha(g)-1)m$ vanishing at $j$; this forces $m=0$.
Normalization is therefore unique and $O$-linear.

We check carefully the local conditions in the asserted union.
The character $\alpha$ is unramified outside $fp$.
At each prime $q\mid f$, $q\ne p$, there is
$\tau\in I_q$ with $\alpha(\tau)-1$ a unit.
In fact primitivity means $\eta$ is nontrivial on the $q$-primary
factor of $(\mathbb Z/f)^\times$; otherwise it would factor
through a proper divisor of $f$.  Cyclotomic inertia surjects onto
that factor.  A nontrivial root of unity of order prime to $p$
cannot be congruent to one modulo $\varpi$: if $u^n=1$,
$(n,p)=1$ and $u\ne1$, the identity
$1+u+\cdots+u^{n-1}=0$ would reduce to $n=0$ in $k$.
As $\chi_p$ is unramified at $q\ne p$, this proves the unit claim.

Take a normalized cocycle representing a class in (OS.1).
Its finite image is killed by some $\varpi^h$, since a continuous
map from the compact group $G$ to the discrete torsion group $A$
has finite image.  Choose also one local coboundary witness at
$p$, and enlarge $h$ to kill it.
Conjugating that witness to every other place above $p$ replaces
it by $\alpha(g)m-c(g)$, so the same $\varpi^h$ kills all those
witnesses.
Away from $fp$ the inertia character is trivial; local cohomological
triviality therefore says that the actual cocycle is zero on inertia.
At $q\mid f$, $q\ne p$, let $m_q$ be an inertia coboundary
witness.  The unit value $\alpha(\tau)-1$ just found shows that
$\varpi^h m_q=0$.  It also promotes the witness from inertia to
the decomposition group: for $d\in D_q$, compare the cocycle
identities for $\tau$ and $d^{-1}\tau d$.
They give
\[
  (\alpha(\tau)-1)
       \bigl(c(d)-(\alpha(d)-1)m_q\bigr)=0,
\]
and the coefficient is a unit.  Thus the normalized cocycle belongs
to $\mathcal T(fp,\varpi^h)$.
Conversely the defining conditions of this module imply all the
local triviality conditions in (OS.1), proving the union assertion.

Module length respects this increasing union without a noetherian
assumption.  For any finite strictly increasing chain of submodules
of the union, choose one element witnessing each strict inclusion.
All the finitely many elements occur in one stage; intersecting the
chain with that stage leaves all its inclusions strict.
A common finite length bound for the stages therefore bounds every
such chain and hence the length of the union.

It remains to treat (OS.39), which the finite derivative argument
deliberately excluded.  At a prime $q\ne p$, cyclotomic inertia is
trivial.  If $\eta$ were nontrivial on that inertia group, the
root-of-unity unit argument above would contradict (OS.39).
Thus $\alpha$ is unramified outside $p$.
Moreover $A(\alpha)^G=0$, because $\alpha(j)-1=-2$ is a unit.
Also $A(\alpha)^{D_p}=0$: inertia at $p$ surjects onto
$\operatorname{Gal}(\mathbb Q(\zeta_p)/\mathbb Q)$, so it contains
an element with $\chi_p\equiv-1\pmod\varpi$; by (OS.39) its
$\alpha-1$ is a unit.

Apply cohomology to
$0\to A[\varpi](\alpha)\to A(\alpha)
\xrightarrow{\varpi}A(\alpha)\to0$.
The two vanishing invariant groups just established imply that a
$\varpi$-torsion class in $\mathscr S(\alpha)$ lifts uniquely to
$H^1(\mathbb Q,A[\varpi](\alpha))$ and is still locally zero at
$p$.  At inertia away from $p$, the action on $A$ is trivial and
multiplication by $\varpi$ on its invariants $A$ is surjective;
the same exact sequence makes the lifted class unramified there.
By (OS.39), $A[\varpi](\alpha)\simeq k(1)$.
Choose an $\mathbb F_p$-basis of $k$ and apply the Kummer
identification to each component.  Each component is represented
by $t\in\mathbb Q^\times$.
Unramifiedness at $q\ne p$ forces $v_q(t)\equiv0\pmod p$:
a $p$th root of a unit is unramified there, while a valuation not
divisible by $p$ forces ramification of order $p$.
Local triviality at $p$ makes $t$ a $p$th power in $\mathbb Q_p$,
so $v_p(t)$ is divisible by $p$ as well.
All rational prime exponents in $t$ are thus divisible by $p$,
and its sign is a $p$th power since $p$ is odd.
Hence $t\in\mathbb Q^{\times p}$ and every lifted component is zero.
We have proved $\mathscr S(\alpha)[\varpi]=0$.
Every cohomology class here is killed by a power of $\varpi$,
again because its cocycle has finite image.  Any nonzero such class
would yield nonzero $\varpi$-torsion by multiplying by its largest
nonvanishing power.  Therefore $\mathscr S(\alpha)=0$, as claimed.
\end{proof}
% END OS8

\subsection{Closing the original-field and eigenspace interface}

% BEGIN OS9
\begin{theorem}[OS9: the literal \texttt{Psi\allowbreak Bound} supplier]
Let $E/\mathbb Q$ be the original abelian CM field, with the
specified scalar towers into $\Omega$, and let
$\rho:(\mathbb Z/N)^\times\to\operatorname{Gal}(E/\mathbb Q)$,
$N>0$, have the specified cyclotomic compatibility
\texttt{Is\allowbreak Cyc\allowbreak Act}.  Let $p$ be odd with $p\nmid[E:\mathbb Q]$,
let $O\subset\overline{\mathbb Q}_p$ be the specified finite DVR,
and let $e\ge1$ and
$\psi:\operatorname{Gal}(E/\mathbb Q)\to O^\times$ satisfy
\[
                 (\psi\circ\rho)(-1)=(-1)^{e+1}.
\]
For every supplied $a\in O$ with
\[
           a=\operatorname{den}(N,e+1)
                 B_{e+1,(\psi\circ\rho)_{\rm prim}},
 \tag{OS.40}
\]
the original eigenspace length satisfies
\[
                         \operatorname{psiLen}(\Omega,e,\psi)
                                  \le v_O(a).
 \tag{OS.41}
\]
This is the universal source proposition \texttt{Psi\allowbreak Bound}, with
all of its original fields, coefficient rings and character readings.
\end{theorem}
\begin{proof}
Let $\eta_N$ be the Dirichlet character obtained from
$\psi\circ\rho$, and let $\eta$ be its primitive inducer of
conductor $f\mid N$.  Its values lie in the same $O$:
reduction of unit groups onto $(\mathbb Z/f)^\times$ is surjective,
and the induced values are independent of a lift.
Its order divides $|\operatorname{Gal}(E/\mathbb Q)|=[E:\mathbb Q]$,
so is prime to $p$.  Parity is unchanged.
The Galois character attached to $\eta_N$ is identical to that
attached to $\eta$, since their cyclotomic reductions commute.
Write this character as $\widetilde\eta$ and put
$\alpha=\chi_p^{-e}\widetilde\eta^{-1}$.

If $a=0$, (OS.41) is immediate because its right side is $\infty$.
Otherwise the Bernoulli value $Y=B_{e+1,\eta}/(e+1)$ is nonzero.
If (OS.39) holds, OS8 makes $\mathscr S(\alpha)=0$.
If it does not hold, there is a $g_0$ with
$\alpha(g_0)-\chi_p(g_0)$ a unit, since $O$ is local.
For each $a_0\ne0$, apply OS7, then OS6, to obtain
\[
 \operatorname{length}\mathcal T(F,a_0)
                    \le v(z)\le v(a).
\]
All these levels $F=f/p^{v_p(f)}\cdot p^m$ have exactly the
same prime divisors as $fp$.  The definition of $\mathcal T$
depends on this prime support, so
$\mathcal T(F,a_0)=\mathcal T(fp,a_0)$.
OS8 now gives, in both cases,
\[
                   \operatorname{length}_O\mathscr S(\alpha)
                                      \le v_O(a).
 \tag{OS.42}
\]
This also explains a useful simplification: Bernoulli nonvanishing
need not be imported into this supplier.  Its zero case was already
covered by the extended-valuation convention.  Separate nonvanishing
is used later when a finite quantitative bound is wanted.

We finally identify (OS.42) with the exact group in (OS.41).
Cyclotomic compatibility of $\rho$ says, for every $g\in G$,
\[
                  \rho(g|_{\mu_N})=g|_E.
 \tag{OS.43}
\]
Indeed both sides have the same action on every element of the
given copy of $E$ in $\Omega$, by \texttt{Is\allowbreak Cyc\allowbreak Act} and the
definition of restriction.  Hence
$\widetilde\eta=\psi\circ\operatorname{res}_E$ as actual
characters on $G$, not merely as isomorphic representations.
Let $U_E=\operatorname{Gal}(\Omega/E)$ and
$\Delta=G/U_E=\operatorname{Gal}(E/\mathbb Q)$.
Its order is a unit in $O$.

Restriction in cohomology therefore gives
\[
 H^1(G,A(\alpha))\simeq H^1(U_E,A(\chi_p^{-e}))^{\psi}.
 \tag{OS.44}
\]
Here is the elementary descent argument including the twist sign.
For a finite group of unit order, positive-degree cohomology
vanishes by the averaging homotopy; inflation--restriction is
therefore an isomorphism onto invariants.
On $U_E$ the finite character is trivial, so the coefficient
module really is the same $A(\chi_p^{-e})$.
The ordinary action on a restricted cocycle is
$z(u)\mapsto\chi_p(g)^{-e}z(g^{-1}ug)$.
The invariant action for the twisted coefficient $\alpha$ instead
uses $\alpha(g)=\chi_p(g)^{-e}\psi(g|_E)^{-1}$.
Multiplying by $\psi(g|_E)$ turns this invariant condition into
the $\psi$-eigencondition in (OS.44).

The isomorphism preserves the exact local kernels.
For each decomposition group and each inertia group of $G$,
its intersection with $U_E$ has quotient a subgroup of $\Delta$,
again of unit order.  The same inflation--restriction argument
makes local restriction injective.  Thus a global class is locally
trivial before restriction if and only if it is trivial on all the
corresponding local groups after restriction.
Primes above $p$ correspond to the specified decomposition groups;
all finite inertias correspond to the other part of (OS.1).
Conjugation simply permutes these actual places, so this argument
works for every original choice of them.
Consequently (OS.44) restricts to an $O$-linear isomorphism
\[
 \mathscr S(\alpha)\simeq
 \operatorname{Sha}(E,A(-e))^{\psi},
 \tag{OS.45}
\]
whose right side is precisely the module defining
$\operatorname{psiLen}(\Omega,e,\psi)$.
Length is invariant under this isomorphism.  Combining (OS.42)
and (OS.45) proves (OS.41), with every required quantifier retained.
\end{proof}
% END OS9

\paragraph{Closed dependency boundary.}
OS1--OS9 supply \texttt{Psi\allowbreak Bound}; they do not assume it at an
intermediate stage.  Their only duality boundary is LD2--LD4,
whose classical inputs are stated there.  Their remaining classical
arithmetic inputs are the finite cyclotomic/local facts, Hilbert
class field, Chebotarev and Kummer theorem stated at the start.
The Jacobi identities, roots-of-unity removal, derivative construction,
local transition, simultaneous prime selection, finite-length descent,
exceptional residual case and tame eigenspace transport are proved
here.  For the classical number-theoretic inputs see J.~S.~Milne,
\href{https://www.jmilne.org/math/CourseNotes/ANT.pdf}{\emph{Algebraic
Number Theory}, version 3.08}, Sections~6--7 and Theorem~8.31,
and
\href{https://www.jmilne.org/math/CourseNotes/CFT.pdf}{\emph{Class
Field Theory}, version 4.03}, V, Theorem~3.6, Corollary~3.7 and
Example~3.9, and II, Proposition~1.22 for Hilbert's theorem~90.
These references specify general classical theorem leaves; the
component provenance records the exact source declarations and pinned
file hashes for the reconstructed special-value argument.

The literal output (OS.41) is the input to the following odd-minus
quantitative argument.  Thus the former bespoke \texttt{Psi\allowbreak Bound}
frontier can be removed once the stated classical inputs and LD are
adopted; this does not remove the separate odd-plus and two-primary
cohomological suppliers or the geometric common-family obligations.


\section{The complete odd-minus estimate from the character-length bound}
\label{sec:odd-minus-bounds}

\status{OB1--OB8 reconstruct the full quantitative supplier
\texttt{Thm\allowbreak Bc\allowbreak Free} from OS1 and \texttt{Psi\allowbreak Bound}, with the explicit
constant $240$.  The proof includes the exceptional primes, the
original-field and local-data comparisons, the negative cyclotomic
twist, and all finite prime and twist sums.  The bespoke dependencies
are the literal \texttt{Psi\allowbreak Bound} proved in OS9 and the finite Jacobi
annihilation proved in OS1.  The classical inputs listed below are
standard theorems, not new project-specific hypotheses.  In particular,
\texttt{kill\allowbreak Pow}, \texttt{gen}, \texttt{prod\allowbreak Val}, \texttt{at\allowbreak Prime} and
\texttt{deg} are proved here in the scope needed by the endpoint, rather
than retained as names for unresolved arguments.}

Fix a CM number field $E$, its given algebraic closure $\Omega$, a
prime $\ell$, and the source's local-data record $D$.  Thus $D$ lists
exactly the primes of $E$ above $\ell$, chooses a place of $\Omega$
above each, and realizes the cyclotomic character on the quotient
$G_{E,S}$ unramified outside those primes.  Put
\[
 A_\ell(e)=(\mathbb Q_\ell/\mathbb Z_\ell)(-e),\qquad
 \mathcal S_\ell(E,e;D)
 =\ker\!\left(H^1(G_{E,S},A_\ell(e))
       \longrightarrow\prod_{v\mid\ell}H^1(E_v,A_\ell(e))\right).
 \tag{OB.1}
\]
All cochains are continuous and the coefficients are discrete.  This is
the literal \texttt{Sha D e}, since \texttt{qz\allowbreak Tw D e} has action
$\chi_{\rm cyc}^{-e}$.  Complex conjugation acts on this cohomology by
conjugation on the group and the same negative twist on coefficients.
For $e\ge1$ write
\[
 \mathcal S_\ell^-(E,e;D)
 =\{x\in\mathcal S_\ell(E,e;D):cx=(-1)^{e+1}x\},\qquad
 m_\ell(E,e;D)=\operatorname{length}_{\mathbb Z_\ell}
                       \mathcal S_\ell^-(E,e;D).
 \tag{OB.2}
\]
Initially lengths take values in $\mathbb N\cup\{\infty\}$; we prove
finiteness before treating them as real numbers.  The notation agrees
with \texttt{Sha\allowbreak Minus} and \texttt{minus\allowbreak Len}.  Write $n=[E:\mathbb Q]$,
and let $N\ge1$ be any modulus equipped with a rational embedding
$E\hookrightarrow\mathbb Q(\zeta_N)$.

\subsection{Classical inputs and the exact remaining interface}
\label{sec:ob-classical}

The following classical theorems are used in the precise forms stated.
They are recorded separately so that their use cannot be mistaken for
a reconstruction of a deeper project-specific supplier.
\begin{description}[style=nextline,leftmargin=0pt,labelwidth=0pt,labelsep=0pt]
\item[OB-CF: unramified class-field reciprocity.]
\label{ob-ref-classfield}
For a totally complex number field $K$, its maximal everywhere
unramified abelian extension has Galois group canonically isomorphic
to $\operatorname{Cl}(K)$ by the Artin map.  This isomorphism respects
automorphisms of $K$.  See J.~S.~Milne, \emph{Class Field Theory},
version 4.03, V, Theorem 3.6, Corollary 3.7 and Example 3.9,
pp.~158--159, \url{https://www.jmilne.org/math/CourseNotes/CFT.pdf}.
The equivariance also follows directly by conjugating Frobenius at
unramified primes.

\item[OB-KU: Kummer theory and the unit theorem.]
\label{ob-ref-kummer}
If $K$ contains $\mu_\ell$, then
$H^1(G_K,\mu_\ell)=K^\times/K^{\times\ell}$, with the usual
valuation criterion away from $\ell$.  The group of $S$-units has
rank $r_1+r_2-1+|S|$ and finite cyclic torsion.  See Milne,
\emph{Class Field Theory}, VII, Appendix A, Theorem A.3 and
Proposition A.5, pp.~226--227, and \emph{Algebraic Number Theory},
version 3.08, Theorem 5.11, p.~90,
\url{https://www.jmilne.org/math/CourseNotes/ANT.pdf}.
The Selmer cardinality bound used below is derived from these facts.

\item[OB-MD: ideal representatives and abelian discriminants.]
\label{ob-ref-minkowski}
Every ideal class of a number field of degree $r$ and discriminant
$\Delta$ has an integral representative of norm at most
$(4/\pi)^{r_2}r!r^{-r}\sqrt{|\Delta|}$ (Milne,
\emph{Algebraic Number Theory}, Theorem 4.3, p.~70).
For an abelian field, $|\Delta|$ is the product of the conductors of
its characters (R.~Sharifi, \emph{Iwasawa Theory}, equation (4.2.1),
\url{https://math.ucla.edu/~sharifi/notes/iwasawa-ch04.html}).

\item[OB-BC: Bernoulli denominators.]
\label{ob-ref-bernoulli}
For positive even $j$, von Staudt--Clausen implies that the
denominator of $B_j$ divides $\prod_{p-1\mid j}p$; also
$B_1=-1/2$ and odd $B_j$ vanish for $j>1$.
See T.~Ciurca, \emph{Arithmetic of cyclotomic fields},
Theorem 6.1.5, p.~111,
\url{https://sites.pitt.edu/~caw203/pdfs/Ciurca--Arithmetic_of_cyclotomic_fields--UROP2018.pdf}.

\item[OB-LF: finite extensions of local fields.]
\label{ob-ref-local}
The integers of a finite extension of $\mathbb Q_\ell$ form a finite
free $\mathbb Z_\ell$-module and a complete DVR, with the unique
extended absolute value.  See Milne, \emph{Algebraic Number Theory},
Theorem 7.38, Corollaries 7.39 and 7.42, and Remark 7.43,
pp.~123--124.  The length and norm formulas needed here are proved
explicitly below.
\end{description}

The main project-specific interface is the theorem supplied by OS9.  Let
$\rho:(\mathbb Z/N)^\times\twoheadrightarrow G=\operatorname{Gal}(E/\mathbb Q)$
be the actual cyclotomic restriction map.  For a finite DVR
$O\subset\overline{\mathbb Q}_\ell$, finite over $\mathbb Z_\ell$,
and a character $\psi:G\to O^\times$, put
$\chi=\psi\circ\rho$, extended by zero on nonunits.  For $k\ge2$ put
\[
 P_k=\prod_{p\le k+1}p,\qquad d(N,k)=NP_k,\qquad
 \alpha_{N,k}(\chi)=d(N,k)B_{k,\chi^*},\qquad
 B(N,k)=16N^{k+1}k!,                         \tag{OB.3}
\]
where $\chi^*$ is the primitive inducing character.  These are the
source's \texttt{den}, \texttt{al\allowbreak N}, \texttt{Bprim}, and \texttt{bnd}.
Let $\mathcal S_{O}(E,e)$ be the absolute-Galois version of (OB.1)
with coefficients $(\operatorname{Frac}O/O)(-e)$, and let
$\mathcal S_O(E,e)^\psi$ be its $\psi$-eigenspace for the actual
$G$-action.  This is the source's \texttt{sha\allowbreak Eig} in
\texttt{psi\allowbreak Len}; the comparison is proved in OB3.

The hypothesis \textbf{PB} is the literal \texttt{Psi\allowbreak Bound}:
for every such original $E,\Omega,N,\rho$, every odd prime
$\ell\nmid n$, every such $O$, every $e\ge1$, every $\psi$ with
$\chi(-1)=(-1)^{e+1}$, and every $a\in O$ whose image is
$\alpha_{N,e+1}(\chi)$,
\[
       \operatorname{length}_O\mathcal S_O(E,e)^\psi
                         \le v_O(a).                \tag{OB.4}
\]
Here $v_O$ is the DVR valuation with value $1$ on a uniformizer and
$v_O(0)=\infty$.  The interface is not replaced by an estimate over
a different field or an unspecified isomorphic coefficient module.

\paragraph{OB1: cyclotomic actions, character count, and the exceptional set.}
\label{claim:OB1}
The given embedding into a cyclotomic field produces a surjective
$\rho$ with
\[
 G\simeq(\mathbb Z/N)^\times/\ker\rho,
 \qquad |G|=[E:\mathbb Q]=n,
 \qquad \rho(-1)=c\ne1.                         \tag{OB.5}
\]
Its lift to a multiple of $N$ is obtained by reduction of units.
For any prescribed parity, the number of complex characters of
$(\mathbb Z/N)^\times$ trivial on $\ker\rho$ is at most $n$.
If $E$ satisfies \texttt{Is\allowbreak Comp\allowbreak Field}$(M,s,q,d;E)$, with
$o=M.\mathrm{deg}$ and the stated \texttt{Prime\allowbreak Pair} hypotheses,
then every odd prime divisor of $n$ divides $o$:
\[
                 \ell\mid n,\quad\ell\ne2
                       \quad\Longrightarrow\quad\ell\mid o.
 \tag{OB.6}
\]

\begin{proof}
The cyclotomic Galois group is $(\mathbb Z/N)^\times$.
All of its subgroups are normal, so the image of the given embedding
of $E$ is Galois over $\mathbb Q$, and restriction is surjective.
Choose an embedding of this cyclotomic overfield into the already
given $\Omega$ extending the original embedding of $E$.  An
automorphism of $\Omega$ acting as the $b$th power on $\mu_N$
therefore restricts to this very $\rho(b)$ on $E$.  This is
\texttt{Is\allowbreak Cyc\allowbreak Act}.  Restriction at a larger modulus commutes with
reduction of units.  Complex conjugation sends every root of unity
to its inverse, giving (OB.5); its restriction is nontrivial because
$E$ is CM.  Distinct characters of a finite abelian group are linearly
independent as functions, by the usual orthogonality sum, so any
subset of the character group has size at most $|G|$.

For (OB.6), work in the actual cyclotomic overfield supplied by
\texttt{Is\allowbreak Comp\allowbreak Field}.  Its fixing condition says that an exponent
$b$ fixes $E$ precisely when it is $\pm1$ modulo $2^{s+2}$,
belongs to the kernel specifying the abelian base $M$, and has
Jacobi symbol $1$ at $q$ and $d$.  Every cyclotomic automorphism
raised to $R=2^{s+1}o$ meets all these conditions: Euler's theorem
handles $2^{s+2}$, the finite base quotient has order $o$, and the
even exponent makes both Jacobi symbols $1$.  Thus the exponent of
$\operatorname{Gal}(E/\mathbb Q)$ divides $R$.  Cauchy's theorem
gives an element of order $\ell$ if $\ell\mid n$, so
$\ell\mid2^{s+1}o$; an odd $\ell$ must divide $o$.
This argument uses the original fixing subgroup and degree.
\end{proof}

\paragraph{OB2: the integral Bernoulli product and its prime valuation budget.}
\label{claim:OB2}
For $k\ge2$, fix a set $P$ of subgroups of $(\mathbb Z/N)^\times$
and let $\mathcal C$ consist of all characters with $\ker\chi\in P$
and $\chi(-1)=(-1)^k$.  There is a nonzero integer $J_{N,k,\mathcal C}$
such that, in every characteristic-zero splitting field,
\[
 \prod_{\chi\in\mathcal C}\alpha_{N,k}(\chi)
     =J_{N,k,\mathcal C},\qquad
 \log|J_{N,k,\mathcal C}|\le|\mathcal C|\log B(N,k).
 \tag{OB.7}
\]
Put $w_\ell=v_\ell(J_{N,k,\mathcal C})$.  For any finite set
$T_\ell$ of the corresponding $\ell$-adic characters and any finite
set $S$ of primes,
\[
 \sum_{\chi\in T_\ell}-\log|\alpha_{N,k}(\chi)|_\ell
       \le w_\ell\log\ell,\qquad
 \sum_{\ell\in S}w_\ell\log\ell
       \le |\mathcal C|\log B(N,k).                 \tag{OB.8}
\]

\begin{proof}
For a character $\eta$ of modulus $f$, the generalized Bernoulli
number is
\[
 B_{k,\eta}=f^{k-1}\sum_{a=1}^{f}\eta(a)B_k(a/f),\qquad
 B_k(X)=\sum_{j=0}^k\binom kj B_jX^{k-j}.          \tag{OB.9}
\]
OB-BC shows that $P_kB_j$ is integral for $0\le j\le k$.
Multiplication of (OB.9) by $fP_k$ consequently yields an integer
linear combination of the values of $\eta$.  If $f\mid N$, the
same is true after multiplication by $NP_k$.  The values of the
primitive inducer occur among those of the original character:
reduction $(\mathbb Z/N)^\times\to(\mathbb Z/f)^\times$ is
surjective, as follows by choosing lifts avoiding the additional
prime divisors of $N$.  Therefore $\alpha_{N,k}(\chi)$ belongs
to the ring generated by the values of $\chi$, and is an algebraic
integer.  In particular it lies in the actual $O$ whenever those
values lie in $O$.

Here are elementary bounds with the source's constant.  Integration
by parts in the Fourier coefficients of $B_k$, using
$B_k'=kB_{k-1}$ and $B_k(1)=B_k(0)$ for $k\ge2$, gives
\[
 B_k(x)=-\frac{k!}{(2\pi i)^k}
             \sum_{r\ne0}\frac{e^{2\pi irx}}{r^k}.
\]
For $k\ge2$ this converges absolutely and
$\zeta(k)\le1+\int_1^\infty u^{-k}du\le2$.
Thus $|B_{k,\eta}|\le4f^kk!/(2\pi)^k$.
The elementary inequality $\prod_{p\le h}p\le4^h$ follows by
induction: primes in $(m,2m]$ divide $\binom{2m}{m}$; primes in
$(m+1,2m+1]$ divide $\binom{2m+1}{m}\le4^m$, since the two
central binomial coefficients are equal.  Applying the induction
hypothesis to $m$ or $m+1$ proves the claim.  Since $4<2\pi$,
\[
 |\alpha_{N,k}(\chi)|
 \le N4^{k+1}\frac{4N^kk!}{(2\pi)^k}
 \le16N^{k+1}k!=B(N,k).                          \tag{OB.10}
\]

Nonvanishing here requires no result at $s=1$.  For primitive
$\eta$ of parity $(-1)^k$, insert the Fourier series into (OB.9).
The finite Gauss-sum identity gives
\[
 B_{k,\eta}
   =-\frac{2k!f^{k-1}\tau(\eta)}{(2\pi i)^k}
                              L(k,\overline\eta).
\]
The Gauss sum is nonzero: character orthogonality gives
$|\tau(\eta)|^2=f$; for the trivial conductor-one character its
value is $1$.  The absolutely convergent Euler product at $k>1$
is nonzero.  One can verify the Gauss identity for all integer
frequencies by substitution when the frequency is coprime to $f$;
otherwise primitivity supplies a unit that fixes the additive
character and has nontrivial multiplicative value, forcing its
sum to vanish.  Parseval on the finite cyclic group then gives
the stated squared norm.  Thus every factor in (OB.7) is nonzero.

Place the character values in one finite cyclotomic field.
Its Galois automorphisms raise these values to powers coprime to
their order, preserving their kernels and the parity.  They
therefore permute the factors of (OB.7).  The product is rational
and integral, hence an integer; it is nonzero by the preceding
paragraph.  The construction uses rational polynomials in the
roots of unity, so the same integer is obtained in every splitting
field.  Taking the product of (OB.10) gives its logarithmic bound.

In an $\ell$-adic splitting field each factor is integral, so its
absolute value is at most $1$.  The negative logarithms are
nonnegative, and their sum over the full set is
$-\log|J|_\ell=v_\ell(J)\log\ell$.  A partial sum is no larger.
Unique factorization of the nonzero integer $J$ gives
$\sum_{\ell\in S}v_\ell(J)\log\ell\le\log|J|$.
These prove (OB.8), including finiteness of the prime budget.
\end{proof}

\paragraph{OB3: the tame-prime length bound on the original local data.}
\label{claim:OB3}
Assume PB.  For every odd prime $\ell\nmid n$ and every $e\ge1$,
$m_\ell(E,e;D)$ is finite.  With
\[
 \mathcal C=\{\chi:\ker\rho\subseteq\ker\chi,
                          \ \chi(-1)=(-1)^{e+1}\},
\]
one has, for every finite set $S$ of odd primes not dividing $n$,
\[
 \sum_{\ell\in S}\log\ell\,m_\ell(E,e;D_\ell)
                \le n\log B(N,e+1).                 \tag{OB.11}
\]
The local-data system $D_\ell$ is arbitrary, and the same $E$ and
$\Omega$ occur in all terms.

\begin{proof}
First identify (OB.1) with the strict Selmer subgroup in absolute
Galois cohomology: its classes vanish on decomposition groups above
$\ell$ and on all inertia groups.  Away from $\ell$, inertia acts
trivially on $A_\ell(e)$; vanishing of its cohomology class is
therefore literal vanishing of a representative cocycle there.
Such a cocycle vanishes on the closed normal subgroup generated
by these inertia groups, so factors through $G_{E,S}$.
Conversely an inflated cocycle is unramified there.  The conditions
at $\ell$ are exactly the local coboundary conditions in (OB.1).
Changing the chosen place above a prime conjugates its decomposition
group and preserves this condition.  This proves the comparison
for the actual $D$, without choosing another host field.

Choose a finite extension of $\mathbb Q_\ell$ inside
$\overline{\mathbb Q}_\ell$ containing the $|G|$th roots of unity,
and let $O$ be its integer ring.  Write
$h=v_O(\ell)=\operatorname{length}_O(O/\ell O)$; this is a positive
finite integer.  Since $O$ is finite free over $\mathbb Z_\ell$,
tensoring continuous cochains with $O$ is the same as taking a
finite direct sum of their coordinate cochains.  It commutes with
cocycles, coboundaries and all local restriction kernels, including
the full family of inertia conditions.  Also
$O\otimes_{\mathbb Z_\ell}(\mathbb Q_\ell/\mathbb Z_\ell)
=\operatorname{Frac}(O)/O$.  Hence
\[
 O\otimes\mathcal S_\ell(E,e;D)=\mathcal S_O(E,e),
 \quad
 \operatorname{length}_O(O\otimes X)
             =h\operatorname{length}_{\mathbb Z_\ell}X.
 \tag{OB.12}
\]
For finite length the second formula follows by tensoring a
composition series and using the length of $O/\ell O$.
For arbitrary torsion modules it follows by taking the supremum
over finitely generated submodules.  All cohomology classes here
are $\ell$-power torsion, since a continuous cocycle from a compact
group into the discrete torsion module has finite image.

Conjugation and coefficient action commute with the comparison in
(OB.12).  Inner automorphisms act trivially on cohomology, as the
cocycle identity expresses their difference from the identity by
a coboundary.  Thus the action factors through the original
$G=\operatorname{Gal}(E/\mathbb Q)$, with exactly the sign in (OB.2).
Since $\ell\nmid |G|$, the elements
\[
 e_\psi=|G|^{-1}\sum_{g\in G}\psi(g)^{-1}[g]\in O[G]
\]
are mutually orthogonal projectors summing to $1$.
This follows by summing finite characters on $G$; the character
set has $|G|$ elements, by decomposing $G$ into cyclic factors.
The $(-1)^{e+1}$ eigenspace of $c$ is therefore the direct sum of
$\mathcal S_O(E,e)^\psi$ for
$\psi(c)=(-1)^{e+1}$.  These are exactly the \texttt{psi\allowbreak Len}
objects in PB, obtained from \texttt{sha\allowbreak Eig} by factoring out the
inner action, not different cohomology groups selected by dimension.

OB2 supplies $a_\psi\in O$ representing
$\alpha_{N,e+1}(\psi\circ\rho)$ and proves $a_\psi\ne0$.
PB and (OB.12) give
\[
 h\,m_\ell(E,e;D)
  =\sum_{\psi(c)=(-1)^{e+1}}
                  \operatorname{length}_O\mathcal S_O(E,e)^\psi
  \le\sum_{\psi(c)=(-1)^{e+1}}v_O(a_\psi)<\infty.
 \tag{OB.13}
\]
Writing $\ell=u\varpi^h$ and $a=v\varpi^{v_O(a)}$ shows
\[
            v_O(a)\log\ell=-h\log|a|_\ell.
\]
Indeed units have absolute value $1$, because both the unit and
its inverse are integral, and $|\ell|_\ell=\ell^{-1}$.
Multiplying (OB.13) by $\log\ell$ and canceling the positive $h$
gives a sum of the logarithms in (OB.8).  Surjectivity of $\rho$
makes $\psi\mapsto\psi\circ\rho$ injective, so no character is
counted twice.  Finally OB2 and $|\mathcal C|\le n$ from OB1
prove (OB.11).  Its integer product depends on $E,N,e$, not on
the varying prime or the choices of $D_\ell$ and $O$.
\end{proof}

\paragraph{OB4: a uniform bound on the number of cohomology generators.}
\label{claim:OB4}
For any number field $E\hookrightarrow\mathbb Q(\zeta_N)$,
any prime $\ell$, any integer $e$, and any original $D$, the
$\ell$-torsion of $H^1(G_{E,S},A_\ell(e))$ has finite length $g$ and
\[
                    g\le8\ell n\log(\ell N).
 \tag{OB.14}
\]

\begin{proof}
Set $K=E(\mu_\ell)$ inside the given $\Omega$.  Then
$[K:E]\mid\ell-1$, $K$ contains $\mu_\ell$, and the given
cyclotomic embedding extends to
$K\hookrightarrow\mathbb Q(\zeta_{\ell N})$.
Put $r=[K:\mathbb Q]\le\ell n$.
An $\ell$-torsion class has a representative cocycle taking values
in $A_\ell(e)[\ell]$: if $\ell f=\delta b$, choose $b'$ with
$\ell b'=b$ and replace $f$ by $f-\delta b'$.
The action on these values becomes trivial over $K$, so restriction
is an additive character $G_K\to\mathbb F_\ell$, unramified
outside the primes over $\ell$.

Choosing one such representative for each class gives an injective
map into this set of characters.  To see injectivity, if two
restrictions agree, the difference cocycle vanishes on $G_K$.
It descends to $\operatorname{Gal}(K/E)$ with values in the
$G_K$-invariants.  Its order is prime to $\ell$; averaging a
cocycle over this finite group writes it as a coboundary, since
division by the group order is available on every
$\mathbb Z_\ell$-module.  The original two classes are equal.
This argument asserts a set injection, and does not incorrectly
assume that the chosen representatives define a canonical linear map.

By Kummer theory, the target injects into
\[
 \operatorname{Sel}_S(K,\ell)
 =\{a\in K^\times:v_{\mathfrak p}(a)\in\ell\mathbb Z
                  \ (\mathfrak p\notin S)\}/K^{\times\ell}.
\]
The valuation condition follows by taking the valuation of
$\beta^\ell=a$ in an unramified local extension away from $\ell$.
Factoring the principal ideal outside $S$ gives the exact sequence
\[
 0\longrightarrow O_{K,S}^\times/(O_{K,S}^\times)^\ell
 \longrightarrow\operatorname{Sel}_S(K,\ell)
 \longrightarrow\operatorname{Cl}_S(K)[\ell]
 \longrightarrow0.
 \tag{OB.15}
\]
The last map sends $a$ to the class of the ideal whose $\ell$th
power is $(a)$ away from $S$.  Its kernel and surjectivity are
immediate from this construction.  There are at most $r$ primes
over $\ell$.  OB-KU therefore bounds the first term in (OB.15)
by $\ell^{2r}$, and the second class-group term by $h_K$.
Consequently the original $\ell$-torsion group has at most
$\ell^{2r}h_K$ elements.  Being killed by $\ell$, it is an
$\mathbb F_\ell$-vector space of dimension and length $g$, so
\[
                    g\log\ell\le2r\log\ell+\log h_K.
 \tag{OB.16}
\]

We record the class-number bound with its constants.  If
$K\hookrightarrow\mathbb Q(\zeta_L)$ and $L\ge2$, OB-MD gives
$|\Delta_K|\le L^r$ and $r\le\varphi(L)\le L$.
Minkowski's constant is at most $1$: for $r\ge2$, use
$r!\le((r+1)/2)^r$ and $(r+1)/(2r)\le3/4$; the degree-one
case is immediate.  Each class has a representative of norm at
most $X=\sqrt{|\Delta_K|}$.
The number of ideals of norm at most $X$ is at most
$X(1+\log X)^{r-1}$: at every rational prime the ideal-counting
Euler factor is coefficientwise at most $(1-T)^{-r}$, and the
sum of the $r$-fold divisor function up to $X$ is bounded by
$X(1+\log X)^{r-1}$ by induction using the harmonic sum.
It follows that
\[
 \log h_K\le\tfrac12\log|\Delta_K|
       +(r-1)\log(1+\tfrac12\log|\Delta_K|)
                       \le3r\log L.                 \tag{OB.17}
\]
For the last inequality use $\log|\Delta_K|\le r\log L$ and
$1+\tfrac12r\log L\le L^2$, which follows from
$r\le L$ and $\log L\le L-1$.
Apply (OB.17) with $L=\ell N$.  Then (OB.16) is at most
$5r\log(\ell N)$.  Since $\log\ell\ge\log2>5/8$,
$g\le8r\log(\ell N)\le8\ell n\log(\ell N)$, as claimed.
\end{proof}

\paragraph{OB5: an actual annihilator of the negatively twisted minus group.}
\label{claim:OB5}
For an odd prime $\ell$ and $e\ge1$, put $R=N\ell$,
$a=2R+1$, $k=e+1$, and let
$\pi:(\mathbb Z/R)^\times\twoheadrightarrow G$
be the lift of the original restriction map.  In $\mathbb Q[G]$
define
\[
 \begin{split}
 E_{L,k}(b)&=\frac{L^{k-1}}{k}B_k(b/L),\\
 y&=(1-a^k[\pi(a)])
                  \left(-\sum_{b\in(\mathbb Z/R)^\times}
                                  E_{R,k}(b)[\pi(b)]\right).
 \end{split}                                                   \tag{OB.18}
\]
Then $y$ has coefficients in $\mathbb Z_\ell$, and its resulting
element $z\in\mathbb Z_\ell[G]$ satisfies
\[
                 z\,\varphi(R)x=0
                 \qquad(x\in\mathcal S_\ell(E,e;D)).
 \tag{OB.19}
\]
The source writes this element as
\texttt{yEl (e+1) a (piQ rho)\^{}(-1)}: the inverse in that
argument compensates for the inverse used in \texttt{sum\allowbreak El},
leaving exactly the group elements $[\pi(b)]$ in (OB.18).

\begin{proof}
For $L=R\ell^m$, let $b$ denote a unit representative in $[0,L)$,
let $b'$ represent $a^{-1}b$ modulo $L$, and set
\[
 j_L(b)=\frac{ab'-b}{L}\in\mathbb Z,\qquad
 E_1(L,a,b)=\frac{a-1}{2}-j_L(b)\in\mathbb Z,
 \qquad
 z_m=-\sum_{b\in(\mathbb Z/L)^\times}
                  b^e E_1(L,a,b)[\pi(b)].             \tag{OB.20}
\]
Here $a$ is odd and coprime to every $L$.
We first prove $z_m\to y$ coefficientwise $\ell$-adically.
The first two terms of the Bernoulli polynomial give
\[
 E_{L,k}(b)=\frac{b^k}{kL}-\frac{b^{k-1}}2
                      +\frac{L}{kP_k}u_b,
                  \qquad u_b\in\mathbb Z.
\]
Indeed all remaining terms have a factor $L$ after their
denominators are multiplied by $kP_k$.
Since $ab'=b+Lj_L(b)$, the binomial theorem, first modulo $L^2$
for the $k$th powers and then modulo $L$ for the $(k-1)$st powers,
shows
\[
 E_{L,k}(b)-a^kE_{L,k}(b')
      =b^eE_1(L,a,b)+\frac{L}{2kP_k}v_b,
                         \qquad v_b\in\mathbb Z.
 \tag{OB.21}
\]
Thus the error has arbitrarily large $\ell$-adic valuation as
$m\to\infty$, uniformly over its finitely many coefficients.
The distribution identity
$\sum_{j=0}^{h-1}B_k((x+j)/h)=h^{1-k}B_k(x)$
follows at once by summing the defining generating function.
Since $\ell\mid R$, every lift of a unit modulo $R$ to $R\ell^m$
is a unit.  This identity proves that the higher-level element in
(OB.18), projected to $G$, is the same $y$.  Reindexing multiplication
by $a$ gives the coefficient difference in (OB.21).  Summing
(OB.21) therefore proves that $y-z_m$ belongs to
$\ell^r\mathbb Z_\ell[G]$ once
$m\ge r+v_\ell(2kP_k)$.  In particular $y$ is integral at $\ell$;
its $\ell$-adic realization is $z$.

We now prove the cohomological annihilation, including the twist.
Pull a cocycle $f$ representing $x\in\mathcal S_\ell(E,e;D)$
back to $G_E$.  Its image is finite, hence killed by some
$\ell^u$.  At each of the finitely many primes of $E$ over
$\ell$, its restriction is $\delta b_v$ for an element $b_v$
of $A_\ell(e)$.  Increase $u$ so that all $b_v$ are killed by
$\ell^u$.  For $m$ large enough, cyclotomic values on
$G_{K_m}$, where $K_m=\mathbb Q(\zeta_{R\ell^m})\subset\Omega$,
are $1$ modulo $\ell^u$.  Consequently
$f|_{G_{K_m}}$ is an additive character and vanishes on the
decomposition groups over $\ell$.  It is unramified at every
other finite prime because the original cocycle factors through
$G_{E,S}$.  These assertions hold at every conjugate place:
the cocycle conjugation identity adds only a term
$(1-\chi_{\rm cyc}(g)^{-e})f(h)$, which is zero on this level.
Since $K_m$ is totally complex, OB-CF makes this restriction an
ordinary class-group character.

We use the finite Jacobi annihilator proved in OS1, keeping its
group-ring orientation.  Increase $m$ so that $a<L$.  For a unit
$t$ modulo $L$, represented in $(0,L)$, its coefficient there is
\[
 \theta_{L,a}(t)
 =\left\lfloor\frac{at}{L}\right\rfloor
       -\left\lfloor\frac{a(L-t)}{L}\right\rfloor
 =2\left\lfloor\frac{at}{L}\right\rfloor-(a-1).
\]
The last equality holds because $at/L$ is not an integer.  Put
$b\equiv at\pmod L$, so $b'=t$ and
$j_L(b)=\lfloor at/L\rfloor$.  Reindexing gives the exact identity
\[
 2\sum_bE_1(L,a,b)\sigma_b^{-1}
       =-\sigma_a^{-1}
               \sum_t\theta_{L,a}(t)\sigma_t^{-1}.
\]
OS1 constructs a principal ideal for the right-hand sum applied to
every ideal, so that sum annihilates $\operatorname{Cl}(K_m)$.
The values of the class-group character $f|_{G_{K_m}}$ have
odd $\ell$-power order; multiplication by $2$ is invertible on
them.  Artin equivariance therefore gives, for $v\in G_{K_m}$
and lifts $\sigma_b$ of the cyclotomic automorphisms,
\[
              \sum_b E_1(L,a,b)f(\sigma_b^{-1}v\sigma_b)=0.
 \tag{OB.22}
\]
For comparison only, the source's classical Stickelberger route is
the integral identity
$\sum_bj_L(b)\sigma_b^{-1}=-\sigma_a^{-1}(\sigma_a-a)\theta_L$,
with $\theta_L=L^{-1}\sum_b b\sigma_b^{-1}$.
Its classical annihilation theorem is Sharifi,
\emph{Iwasawa Theory}, Definition 4.5.1, Lemma 4.5.2 and
Theorem 4.5.6.  It is not a retained input: the OS1 argument just
given proves exactly the odd-primary consequence needed here.

Let $U=G_{E(\mu_{\ell^\infty})}$ and let
$v\in U\cap G_{\mathbb Q(\mu_R)}$.  Such $v$ belongs to every
$G_{K_m}$.  On restriction to $U$, the action on a cohomology class is
\[
 \operatorname{res}(\pi(b)x)(v)
   =\chi_{\rm cyc}(\sigma_b)^{-e}
                       f(\sigma_b^{-1}v\sigma_b).
 \tag{OB.23}
\]
The lift satisfies $\chi_{\rm cyc}(\sigma_b)\equiv b\pmod{\ell^m}$.
Because its values are units and the displayed cocycle values are
killed by $\ell^u$, multiplication by
$b^e\chi_{\rm cyc}(\sigma_b)^{-e}$ acts as $1$ for $m\ge u$.
Equations (OB.20), (OB.22), and (OB.23) now give
$\operatorname{res}(z_mx)(v)=0$.  The coefficient congruence
$z-z_m\in\ell^u\mathbb Z_\ell[G]$ allows $z_m$ to be replaced
by $z$.  This explicitly accounts for the negative twist.

Finally restriction to $U$ is injective on the original $H^1$.
Indeed its kernel is represented by cocycles descending to
$\operatorname{Gal}(E(\mu_{\ell^\infty})/E)$, an infinite open
subgroup of $\mathbb Z_\ell^\times$.  For odd $\ell$ it is
procyclic: its torsion part has order prime to $\ell$ and its
principal-unit part is a copy of $\mathbb Z_\ell$.
If $\gamma$ is a topological generator, then
$\chi_{\rm cyc}(\gamma)^{-e}-1\ne0$, because $e\ge1$ and
$\gamma$ has infinite order.  Multiplication by this nonzero
$\ell$-adic integer is surjective on
$\mathbb Q_\ell/\mathbb Z_\ell$.  Subtract a coboundary agreeing
with the cocycle at $\gamma$; the remaining cocycle is zero on
the dense cyclic subgroup and hence everywhere by continuity.
This proves injectivity.  The homomorphism
$\operatorname{res}(zx):U\to A_\ell(e)$ vanishes on
$U\cap G_{\mathbb Q(\mu_R)}$.  Its quotient is a subgroup of
$(\mathbb Z/R)^\times$, so its exponent divides $\varphi(R)$.
Thus $\varphi(R)\operatorname{res}(zx)=0$, and injectivity gives
(OB.19).  No infinite-level finiteness assertion or implicit change
of the local-data system has been used.
\end{proof}

\paragraph{OB6: a determinant gives an explicit annihilating power.}
\label{claim:OB6}
For every odd $\ell$, every $e\ge1$, and every original $D$, there
is $r\ge0$ such that
\[
 \ell^r\mathcal S_\ell^-(E,e;D)=0,\qquad
      r\log\ell\le30n(e+1)\log(e\ell N).
 \tag{OB.24}
\]

\begin{proof}
Keep $R,a,k,G,z,y$ from OB5 and put $\varepsilon=(-1)^k$.
Because $c\in G$ is a nonidentity involution,
\[
       Q=\mathbb Z_\ell[G]/([c]-\varepsilon)
\]
is free of rank $n/2$: choose one representative from each pair
$\{g,cg\}$, replacing the other by $\varepsilon g$.
Let $\delta$ be the determinant of multiplication by $z$ on $Q$.
The adjugate identity, applied to $1\in Q$, gives
$\delta\in(z,[c]-\varepsilon)$ in $\mathbb Z_\ell[G]$.
Together with (OB.19) and $cx=\varepsilon x$, it follows that
$\delta\varphi(R)x=0$ for every $x$ in the minus group.

Put
\[
 D_{R,k}=Rk(k+1)!,\quad
 S_{R,a,k}=2a^kR^kk!,\quad
 X_{R,a,k}=D_{R,k}S_{R,a,k}.                    \tag{OB.25}
\]
Multiplying $y$ by $D_{R,k}$ gives an integral group-ring element:
the Bernoulli polynomial expansion and OB-BC clear every coefficient.
The determinant on the integral signed quotient is therefore a
rational integer $J$, and
\[
      J=D_{R,k}^{\,n/2}\delta,\qquad
      0<|J|\le X_{R,a,k}^{\,n/2}.                  \tag{OB.26}
\]
Here is a proof of nonvanishing and the bound.  Over $\mathbb C$ the
signed quotient diagonalizes into the $n/2$ characters $\eta$ of
$G$ with $\eta(c)=\varepsilon$.  The eigenvalue of $y$ is
\[
       (1-a^k\chi(a))\frac{-B_{k,\chi}}{k},
               \qquad \chi=\eta\circ\pi.
 \tag{OB.27}
\]
The smoothing factor is nonzero, since $a>1$ and $|\chi(a)|=1$.
For an imprimitive character of level $R$, the distribution formula
in (OB.9), excluding the extra nonunits prime by prime, gives
\[
 B_{k,\chi}=B_{k,\chi^*}
       \prod_{p\mid R,\ p\nmid\operatorname{cond}\chi}
                         (1-\chi^*(p)p^{k-1}).
\]
Each extra factor is nonzero for $k\ge2$, and the primitive factor
is nonzero by OB2.  Thus (OB.27) is nonzero.  The same Fourier bound
as in OB2 gives $|B_{k,\chi}|\le R^kk!$, because
$4/(2\pi)^k\le1$.  Also
$|1-a^k\chi(a)|\le2a^k$ and $k\ge1$.
This proves the bound $S_{R,a,k}$ for each eigenvalue, hence (OB.26).

Set $r_0=v_\ell(J)$ and $r_1=v_\ell(\varphi(R))$.
Since $\delta\mid J$ in $\mathbb Z_\ell$, it divides
$\ell^{r_0}$ there; similarly $\varphi(R)\mid\ell^{r_1}$.
Thus $r=r_0+r_1$ kills the minus group.  Moreover
\[
 r\log\ell\le n\log X_{R,a,k}+\log R
                    \le n\log(X_{R,a,k}R).          \tag{OB.28}
\]
The first inequality even permits $n/2$ in place of $n$ in its
first term; the weaker form is the one used by the source.

For completeness, its numerical constant is explicit.  Put
$U=e\ell N\ge2$ and $T=3U$.  Then
$R\le T$, $a=2R+1\le T^3$, $k=e+1\le T$,
$k+1\le T^2$, and $k\ge2$.  Hence
\[
 \begin{split}
 X_{R,a,k}R
 &=Rk(k+1)!\,2a^kR^kk!\,R\\
 &\le T^{7k+6}\le T^{10k}.
 \end{split}
\]
We used $(k+1)!\le(T^2)^{k+1}$ and $k!\le T^k$.
Since $3\le U^2$, $\log T\le3\log U$.
Substituting in (OB.28) proves (OB.24) with $30$.
The construction fixes $r$ before $x$ varies; it is not a
pointwise collection of unrelated annihilating exponents.
\end{proof}

\paragraph{OB7: exponent times generators gives the exceptional-prime bound.}
\label{claim:OB7}
For every odd prime $\ell$, every $e\ge1$, and every $D$,
\[
 m_\ell(E,e;D)<\infty,\qquad
 m_\ell(E,e;D)\log\ell
       \le240\ell n^2(e+1)\log(e\ell N)^2.
 \tag{OB.29}
\]

\begin{proof}
For an endomorphism $f$ of any module, the map induced by $f$ from
$\ker f^{j+1}$ to $\ker f^j$ has kernel $\ker f$.
The length inequality for this exact sequence proves inductively
that $\operatorname{length}\ker f^r\le
r\operatorname{length}\ker f$.
Apply this to multiplication by $\ell$ on the ambient $H^1$.
OB6 puts the minus group inside $\ker\ell^r$, while OB4 bounds
the length $g$ of $\ker\ell$.  Thus its length is finite and at
most $rg$.  Multiplication by $\log\ell$ and (OB.14), (OB.24)
give
\[
 m_\ell\log\ell
   \le240\ell n^2(e+1)\log(e\ell N)\log(\ell N)
   \le240\ell n^2(e+1)\log(e\ell N)^2.
\]
Both logarithms are nonnegative and $e\ge1$.  This is precisely
the odd-prime portion of \texttt{Stick\allowbreak Bound 240} needed below.
No claim about its unused $\ell=2$ extension is required.
\end{proof}

\paragraph{OB8: the full uniform \texttt{Thm\allowbreak Bc\allowbreak Free} supplier.}
\label{claim:OB8}
Use PB from OS9, the Jacobi annihilator from OS1, and the explicitly
listed classical theorems.  Let
$M=(m,H)$ be any \texttt{Ab\allowbreak F}, put $o=[(\mathbb Z/m)^\times:H]$,
let $s,t\ge2$, and let $q,d$ satisfy \texttt{Prime\allowbreak Pair}$(m,q,d)$.
For every original CM number field $E$ satisfying
\texttt{Is\allowbreak Comp\allowbreak Field}$(M,s,q,d;E)$, put
$n=[E:\mathbb Q]$ and $F=\operatorname{condF}(E)$.
Then for every given algebraic closure $\Omega$, every family of
local data $(D_\ell)$ on that same $E,\Omega$, and every finite
set $S$ of odd primes,
\[
 \begin{split}
 &m_\ell(E,e;D_\ell)<\infty
                 \qquad(\ell\text{ odd},\ 1\le e<t),\\
 &\sum_{\ell\in S}\log\ell\sum_{e=1}^{t-1}
                         m_\ell(E,e;D_\ell)\\
 &\qquad\le3nt^2\log(tF)
       +240\sum_{\substack{\ell\mid o\\\ell\text{ odd prime}}}
                       \ell n^2t^2\log(t\ell F)^2.
 \end{split}                                                   \tag{OB.30}
\]
Thus the literal \texttt{Thm\allowbreak Bc\allowbreak Free} holds with $C=240$, and it
supplies VW.13--VW.17 with that same constant.

\begin{proof}
The set defining $F$ is nonempty by \texttt{Is\allowbreak Comp\allowbreak Field}.  Its
least element is positive and itself supplies an embedding of the
original $E$ into $\mathbb Q(\zeta_F)$.  Apply OB3 and OB7 with
this $N=F$.  They prove all asserted individual finiteness
statements before any real-valued lengths are summed.

Split $S$ into $S_0=\{\ell\in S:\ell\nmid n\}$ and its complement.
For $1\le e<t$, (OB.11) bounds the sum over $S_0$ by
$n\log B(F,e+1)$.  The source's elementary numerical bound is
\[
        B(F,e+1)=16F^{e+2}(e+1)!
                       \le(tF)^{2t+5}.              \tag{OB.31}
\]
Indeed $16\le(tF)^4$, $F^{e+2}\le(tF)^{t+1}$, and
$(e+1)!\le t^t\le(tF)^t$.  Interchanging the two finite sums and
using $(t-1)(2t+5)\le3t^2$ for integers $t\ge2$ gives
\[
       \sum_{\ell\in S_0}\log\ell
                    \sum_{e=1}^{t-1}m_\ell
                        \le3nt^2\log(tF).           \tag{OB.32}
\]
For a prime in the complementary set, (OB.29) and
$e+1\le t$, $\log(e\ell F)\le\log(t\ell F)$ give
\[
      \log\ell\sum_{e=1}^{t-1}m_\ell
              \le240\ell n^2t^2\log(t\ell F)^2.
 \tag{OB.33}
\]
All summands are nonnegative.  OB1 puts every prime in this
complementary set among the odd prime divisors of the fixed
base degree $o$, so (OB.33) may be summed over that full set.
Combining it with (OB.32) gives (OB.30).  The same calculation works
for every $S$, with no need to assert beforehand that only finitely
many prime contributions are nonzero.  The resulting right side
is exactly \texttt{minus\allowbreak Bound 240 o n F t}, and the universal
quantifiers are exactly those of \texttt{VVminus\allowbreak Le} and
\texttt{Thm\allowbreak Bc\allowbreak Free}.
\end{proof}

\paragraph{Dependency consequence.}
OB1--OB8 close the full odd-minus supplier from the literal PB
interface proved in OS9, the Jacobi annihilation proved in OS1,
and the named classical theorems.  Together these replace
\texttt{Thm\allowbreak Bc\allowbreak Free} in the structural boundary of VW4; no additional
bespoke odd-minus arithmetic or quantitative interface remains
in this branch.  The resulting uniform
constant is independent of the abelian base, conductor, host primes,
height, twist range, algebraic closure and local-data choices.
This statement leaves the other VW4 suppliers, and the later
lattice-gain/character-balance inputs, as separate obligations.


\section{Congruence denominators, the good cut, and the original family}
\label{sec:congruence-chain}

\status{LN1--LN10 close the chain from an integral good-cut operator and
explicit local denominator estimates to \texttt{Link\allowbreak N} on the original
family.  Two useful simplifications are proved: a regular-representation
trace makes the central coefficients integral in the original DVR, and
a filtration by individual eigenspaces bounds a congruence module by its
coordinate congruence ideals.  The argument retains the actual local
lattices, theta ranks, place weights and cohomological sum.  The
per-prime Galois bounds and the compatibility of the actual theta/type
and local data are explicit inputs; their existence is not inferred
from this numerical and algebraic closure.}

Fix one source datum $D$, and write
\[
 \begin{gathered}
 N_G=D.\mathrm{what}_2.\mathrm{restRowsFiniteGroupCard}>0,\\
 Y_i=D.\mathrm{nd}.Z_{\rm good}(i)\subseteq
 Z_i=D.\mathrm{nd}.Z(i)\subseteq X_i.
 \end{gathered}
 \tag{LN.1}
\]
For a localisation $i$, put $\mathcal O=\mathcal O_i$, $K=\operatorname{Frac}
\mathcal O$, $M=M_i$ and $V=K\otimes_{\mathcal O}M$.  The source supplies
a finite set $X=X_i$, a separating $\mathcal O$-subalgebra
$T\subseteq\mathcal O^X$, the actual finite free $\mathcal O$-module
$M$, an integral action of $T$ on $M$, and the
eigenspace-spanning assertion
\[
 V=\sum_{x\in X}V_x,\qquad
 V_x=\{v:f v=f(x)v\text{ for every }f\in T\}.
 \tag{LN.2}
\]
The spanning assertion is a real hypothesis; separation of $T$ alone
does not imply it.  The $\mathcal O_i$ are the original DVRs, finite
over $\mathbb Z_{\ell_i}$ and of characteristic zero through their
specified number-field embeddings.  All lengths below are ordinary
finite module lengths, whose finiteness will be proved before division
or conversion from the source's extended-natural lengths.

\paragraph{LN1: central coefficients in the original coefficient ring.}
\label{claim:LN1}
Let $G$ be a finite group of order $N$, and let
$e\in\mathbb C[G]$ be a central idempotent.  If $\zeta$ is a primitive
$N$th root of unity, then
\[
                   N e(g)\in\mathbb Z[\zeta]\qquad(g\in G).
 \tag{LN.3}
\]
Here $e(g)$ denotes the coefficient of $g$.  To prove this without
character classification, set $y=g^{-1}e$.  Centrality, $e^2=e$ and
$g^N=1$ give $y^{N+1}=y$.  Left multiplication by $y$ on
$\mathbb C[G]$ therefore has eigenvalues in
$\{0\}\cup\{\zeta^a:0\le a<N\}$: its annihilating polynomial is
$t(t^N-1)$.  Its trace is a sum of these eigenvalues and belongs to
$\mathbb Z[\zeta]$.  In the basis $h\in G$, every diagonal entry of
left multiplication by $y$ is $y(1)=e(g)$, so the same trace is
$Ne(g)$.  This proves (LN.3), also for $N=1$.

Suppose the specified embedding $\iota:\mathcal O\hookrightarrow
\mathbb C$ contains a primitive $N$th root.  Lift the coefficients in
(LN.3) through $\iota$ to obtain $z\in\mathcal O[G]$ with
$\iota(z)=Ne$.  The coefficient lifts are unique.  For every $h\in G$,
$hz-zh$ maps to zero, so injectivity gives $hz=zh$.  Thus $z$ belongs
to the center of the original group algebra.  No extension of
$\mathcal O$, change of lattice or unramifiedness assumption is used.

In the actual datum,
\[
 G=\prod_{k=0,1}\prod_{v\in\mathrm{PlG}(k)}
       \mathrm{GL}_r(\mathbb F_{g_k}),\qquad |G|=N_G.
 \tag{LN.4}
\]
The equality is precisely the two source fields
\texttt{card\_\allowbreak Pl\allowbreak G} and \texttt{rest\allowbreak N\_\allowbreak card\_\allowbreak level\allowbreak Group\_\allowbreak eq}.
The chosen \texttt{Pend} supplies a primitive $N_G$th root in every
$\mathcal O_i$.  For every tuple of local labels, the product of the
central local idempotents, embedded in different factors of (LN.4),
is again a central idempotent.  Apply (LN.3) to this product.  This
proves the literal \texttt{Central\allowbreak Elt} coefficient-lifting interface.

The element called $e_0$ in the source is a finite sum over its
\emph{ordered root vectors}.  Lift each good summand as above and add
the lifts, with exactly the original multiplicities.  The resulting
central integral element $z_i$ satisfies
\[
                  \iota(z_i)=N_Ge_0.                 \tag{LN.5}
\]
This proof does not replace the ordered tuples by unordered
polynomials, and does not assert that the sum $e_0$ is an idempotent.
Its specific action on the relevant theta types is the separate
compatibility assertion used next.

\paragraph{LN2: construction of the literal good cut.}
\label{claim:LN2}
We spell out the source inputs used to turn (LN.5) into a cut.  For
each $i$, actual automorphic objects $\Pi$ map to points
$\operatorname{pt}_i(\Pi)\in X_i$ or to no point; every $x\in X_i$
has a representative.  Each such $\Pi$ has a nonzero type subspace.
On every type vector above $x\in Z_i$ the same $e_0$ acts as
\[
 e_0F=F\quad(x\in Y_i),\qquad
 e_0F=0\quad(x\in Z_i\setminus Y_i).                 \tag{LN.6}
\]
This is the exact \texttt{on\allowbreak Mt\_\allowbreak e0\_\allowbreak eq\_\allowbreak self\_\allowbreak of\_\allowbreak mem\_\allowbreak zgood\_\allowbreak and\_\allowbreak eq\_\allowbreak zero\_\allowbreak of\_\allowbreak not\allowbreak Mem}
interface.  It depends on the actual parameter/type dictionary and
the theta decomposition; LN1 does not prove it.  The original
\texttt{Pend} also gives the integral action $z\mapsto z_{\mathcal O}$,
its commutation with Hecke operators, and \texttt{Pass.\allowbreak Eig\allowbreak Span}: if the
complex action of $z$ is a given scalar on every type space above $x$,
then $(z_{\mathcal O})_K$ is that scalar on $V_x$.

Define, using these actual type spaces,
\[
 S_i=\{x:\exists\Pi,F,\ \operatorname{pt}_i(\Pi)=x,
          \ F\in\operatorname{type}_i(\Pi),\ e_0F\ne0\},
 \qquad \operatorname{keep}(i)\ \Longleftrightarrow\ Y_i\ne\varnothing.
 \tag{LN.7}
\]
Choose a representative and a nonzero type vector for $x\in Y_i$.
Equation (LN.6) proves $Y_i\subseteq S_i$.  Its zero clause proves
$S_i\cap Z_i=Y_i$, and hence $S_i$ is nonempty whenever $i$ is kept.
Let $\alpha_i$ be the action of the actual $z_i$ from (LN.5).
The scalar-compatibility input just stated gives
\[
 [\alpha_i,T_i]=0,\qquad
 (\alpha_i)_K|_{V_x}=N_G\quad(x\in Y_i),\qquad
 (\alpha_i)_K|_{V_x}=0\quad(x\notin S_i).
 \tag{LN.8}
\]
No scalar action on the other selected eigenspaces is required.

There is one further field in \texttt{Cut\allowbreak D.\allowbreak Good}: every actual
automorphic representative of a selected point must satisfy
\texttt{Nc} at both generator-prime places.  The precise remaining
type input says that an occurring local idempotent with label $f$
identifies the finite-field inertia-label polynomial with
$f(X)(X-u)^t$, and supplies monodromy zero.  Its degree is the
actual local dimension, at most $n_V$; the generator primes split
completely and satisfy the source's largeness bound, so the
\texttt{Sqf} multiplicity threshold is $2$ there.  For a
good tuple, $f$ is monic, has nonzero constant term, is squarefree,
and $f(u)\ne0$, $f(u^{-1})\ne0$.  Consequently its only possible repeated
label root is $u$ (or $u^{-1}$ at the conjugate place), a unit of
$\mathbb F_{g_k}$ with $(g_k-1)$st power $1$.
The actual complex inertia eigenvalues are read through the
\texttt{Pend} characters $\omega_a(\gamma)$, not by an embedding
of finite fields into $\mathbb C$.  These characters are locally
constant, trivial on wild inertia, separate distinct nonzero label
roots, and are multiplicative in the root $a$.  If two such
characters agree at a tame generator, they agree on its powers
and on wild inertia; density and local constancy give equality on
all inertia, hence equality of their labels.  Evaluation at any
tame generator is therefore injective on the nonzero labels and
preserves their multiplicities.  Also
$a^{g_k-1}=1$ implies $\omega_a(\gamma)^{g_k-1}=1$ by
multiplicativity.  The supplied inertia-characteristic-polynomial
formula now proves exactly the multiplicity-at-least-two clause
of \texttt{Nc}.  On wild inertia every character is $1$, which
proves its \texttt{Tm} clause; with the supplied zero monodromy
this proves \texttt{Nc}.  The argument applies to every tame
generator, as its definition requires.

For completeness, nonzero $e_0F$ supplies the needed occurring
labels.  A nonzero finite sum has a nonzero summand.  If a product of
commuting central local operators is nonzero on $F$, every factor
is nonzero on $F$.  Expanding such a factor as the sum of the
primitive local idempotents having its label produces an occurring
local type with that label.  These are purely finite-dimensional
arguments.  The definition of $S_i$ gives one representative;
the source's \texttt{provision\allowbreak Forty\allowbreak One} says that two representatives
of the same selected Hecke point have the same base-change
constituents up to permutation.  Their local parameter multisets
therefore agree.  This proves \texttt{Nc} for \emph{every}
representative, as required, rather than only for the witness in
(LN.7).  Equations (LN.7)--(LN.8) and this last argument construct
all five fields of \texttt{Cut\allowbreak D.\allowbreak Good}.

The input boundary is exact: (LN.6), integral/complex action
compatibility, the local parameter dictionary, and the
same-point permutation assertion are retained here.  This is a
complete construction from those fields, not an assertion that
arbitrary finite-group types produce the requisite automorphic data.

\paragraph{LN3: the actual spectral quotient and finite congruence module.}
\label{claim:LN3}
For each $x\in X$, separation gives, for each $y\ne x$, an
$f_{xy}\in T$ with $f_{xy}(x)\ne f_{xy}(y)$.  Then
\[
 h_x=\prod_{y\ne x}(f_{xy}-f_{xy}(y))\in T,
 \quad h_x(x)\ne0,\quad h_x(y)=0\ (y\ne x).          \tag{LN.9}
\]
The empty product is $1$.  The operators $h_x/h_x(x)$ project onto
$V_x$ in (LN.2), so that sum is direct.  Write
$V_A=\bigoplus_{x\in A}V_x$ and $\pi_A:V\to V_A$.
Because $M$ is free, it embeds in $V$; we use that embedding below.

For any $A\subseteq X$, set
\[
 U_A=\pi_A(M)=V_A\cap(M+V_{A^c}),\quad
 L_A=M\cap V_A,\quad C_A(M)=U_A/L_A.                 \tag{LN.10}
\]
The image $U_A$ is a finite free $\mathcal O$-module of rank
$\dim_KV_A$, because it is a finitely generated torsion-free
submodule of a finite-dimensional $K$-space.  Clearing the finitely
many denominators in the projections (LN.9) gives a nonzero
$d\in\mathcal O$ with $dU_A\subseteq L_A$.  Thus $L_A$ has full
rank in $U_A$ and $C_A(M)$ has finite length.  In particular (LN.10)
for $A=Y_i$ is precisely the source's
\texttt{Mt\allowbreak Up}/\texttt{Mt\allowbreak Low} quotient.

For a selected set $S$, the source defines \texttt{out\allowbreak M} as the
intersection of kernels of the integral operators $f$ which vanish
on $S^c$.  By (LN.9) and the direct sum,
\[
 \operatorname{outM}=M\cap V_{S^c},\qquad
 M'=M/(M\cap V_{S^c}),\qquad
 K\otimes M'\simeq V/V_{S^c}\simeq V_S.              \tag{LN.11}
\]
The quotient $M'$ is torsion-free: if $am\in V_{S^c}$ for a
nonzero $a\in\mathcal O$, then $m\in V_{S^c}$.  Hence $M'$ is
finite free.  Restricting $T$ to $S$ gives the actual \texttt{cut\allowbreak T};
its kernel kills $M'$ and the induced algebra still separates $S$.
The eigenspace at every retained $x$ is the image of $V_x$ under the
quotient, and that map is an isomorphism.  Thus multiplicities of
retained points agree exactly.  These identifications give the
literal source cut module, not a new freely chosen lattice on the
same rational representation.

\paragraph{LN4: a direct filtration bound for coordinate congruence ideals.}
\label{claim:LN4}
For $A\subseteq X$ and $x\in A$, define the nonzero ideal
\[
 I_x(A)=\{f(x):f\in T,\ f|_{A^c}=0\}\subseteq\mathcal O.
 \tag{LN.12}
\]
It is an ideal because $T$ contains constant functions.  It is
nonzero by (LN.9).  The following finite-length estimate is valid
with arbitrary eigenspace multiplicities:
\[
 \operatorname{length}_{\mathcal O} C_A(M)
 \le\sum_{x\in A}\dim_KV_x\,
             \operatorname{length}_{\mathcal O}(\mathcal O/I_x(A)).
 \tag{LN.13}
\]
Order $A=\{x_1,\ldots,x_s\}$ and put
$F_a=\bigoplus_{b\le a}V_{x_b}$ and $U_a=U_A\cap F_a$.
The images of $U_a$ filter $C_A(M)$.  Its $a$th graded quotient is
a quotient of the free module $U_a/U_{a-1}$, whose rank is at most
$\dim_KV_{x_a}$.  Freeness follows because this quotient embeds in
$F_a/F_{a-1}$.  If $f\in T$ vanishes on $A^c$, then
$f\pi_A(m)=fm\in L_A$ for $m\in M$.  On $F_a/F_{a-1}$ it acts by
$f(x_a)$.  More explicitly,
$(f-f(x_a))U_a\subseteq U_{a-1}$, since $f$ preserves $U_A$.
Thus $I_{x_a}(A)$ kills the graded quotient.  A free basis of
$U_a/U_{a-1}$ therefore gives a surjection
$(\mathcal O/I_{x_a}(A))^{\dim V_{x_a}}$ onto it, after adding
zero generators if necessary.  Length additivity along the finite
filtration proves (LN.13).  The case $A=\varnothing$ has both sides
zero.

This argument replaces a longer congruence-module decomposition.
It uses no Galois representation, semisimplicity theorem for an
integral Hecke algebra, or equality of different lattices.  The
ideals (LN.12) are the literal \texttt{cc\allowbreak At} images: operators
supported on the stated theta subset of the stated $X$.

\paragraph{LN5: cut defect and the original-row denominator inequality.}
\label{claim:LN5}
Suppose $Y\subseteq S$, $S\cap Z=Y$ and $\alpha\in
\operatorname{End}_{\mathcal O}(M)$ obeys (LN.8), with a nonzero
integer $N$ in place of $N_G$.  Let $q:V\to V/V_{S^c}$ be (LN.11).
It induces a map
$C_Y(M)\to C_Y(M')$: indeed $q(U_Y)\subseteq U'_Y$ and
$q(L_Y)\subseteq L'_Y$.  If a class represented by $u\in U_Y$
lies in its kernel, $q(u)=q(m)$ for some $m\in M$; hence
$u-m\in V_{S^c}$ and
\[
              Nu=\alpha u=\alpha m\in M\cap V_Y=L_Y.
 \tag{LN.14}
\]
The kernel is killed by $N$.  Its inverse image in $U_Y$ is a
submodule of a free module of rank $m_Y=\sum_{x\in Y}\dim V_x$,
and is therefore generated by at most $m_Y$ elements.  It follows
that the kernel is a quotient of $(\mathcal O/N\mathcal O)^{m_Y}$.
Length additivity for kernel and image, followed by inclusion of
the image in $C_Y(M')$, gives
\[
 \operatorname{length}C_Y(M)
 \le\operatorname{length}C_Y(M')+
               m_Y\operatorname{length}(\mathcal O/N\mathcal O).
 \tag{LN.15}
\]
Surjectivity onto the cut congruence module is neither needed nor
asserted.

For the normalisation, let $\ell=\ell_i$ and
$e_i=\operatorname{length}_{\mathcal O_i}(\mathcal O_i/\ell\mathcal O_i)$.
This is a positive finite integer.  Indeed $\ell\ne0$, and $\ell$
is not a unit: otherwise the finite nonzero
$\mathbb Z_\ell$-module $\mathcal O_i=\ell\mathcal O_i$ would
vanish by the elementary Nakayama argument (reduce a finite
generating set modulo the maximal ideal, or use its determinant
relation).  The other prime factors of an integer are units through
$\mathbb Z_\ell\to\mathcal O_i$.  Factorising $N$ and adding DVR
lengths gives the exact identity
\[
 \frac{\operatorname{length}(\mathcal O_i/N\mathcal O_i)}{e_i}
                   =v_\ell(N).                         \tag{LN.16}
\]
All involved lengths are finite by LN3 and $N\ne0$, so this is also
the source's \texttt{len\allowbreak Unit}, including its \texttt{to\allowbreak Nat}
convention.  Dividing (LN.15) by $e_i$ proves
\[
 \eta_i\le\eta'_i+
             \sum_{x\in Y_i}\operatorname{mult}_i(x)v_{\ell_i}(N_G).
 \tag{LN.17}
\]
If $Y_i$ is empty, $U_{Y_i}=0$, so $\eta_i=0$; this handles every
discarded localisation.

There is also a useful direct export to the original ideals.  If
$f\in T$ vanishes on $Z^c$, its restriction to $S$ vanishes on
$S\setminus Y$.  Consequently
\[
 I_x(Z)\subseteq I'_x(Y),\qquad
 \operatorname{length}(\mathcal O/I'_x(Y))
       \le\operatorname{length}(\mathcal O/I_x(Z))\quad(x\in Y).
 \tag{LN.18}
\]
Apply LN4 to the cut module and combine (LN.15)--(LN.18).  Define
$v\!c_i(x)=\operatorname{length}(\mathcal O_i/I_x(Z_i))/e_i$.
Then
\[
 \boxed{\quad\eta_i\le
       \sum_{x\in Y_i}\operatorname{mult}_i(x)
           \{v\!c_i(x)+v_{\ell_i}(N_G)\}.\quad}        \tag{LN.19}
\]
This is the actual original-row \texttt{Twisted\allowbreak Places.\allowbreak adm1}; the
proof works at every prime where the cut operator is supplied,
including more than its required non-generator range.  The local
denominator $v\!c_i(x)$ remains the original \texttt{vcc}, not the
denominator for an unrelated representation or quotient.

\paragraph{LN6: multiplicities are paid for by the same theta rank.}
\label{claim:LN6}
Fix an actual lift $j$ and a twisted place $w$.  The same complex
space $M_j$ has an internal direct sum decomposition
$M_j=\bigoplus_\sigma E_{j,\sigma}$.  Let $A_i$ be the finite set
of good rows occurring at localisation $i$.  The exact
\texttt{Conclusion\allowbreak Z} identification supplies
\[
 \sum_{x\in Y_i}\operatorname{mult}_i(x)
          \le\sum_{\sigma\in A_i}\dim E_{j,\sigma}
           \quad(i\in\operatorname{locs}(j,w)).          \tag{LN.20}
\]
For a nonempty $Y_i$ this is equality in the source.  For an empty
$Y_i$, the left side is zero and the right side is nonnegative.
\texttt{Localisation\allowbreak Uniqueness} says that the sets $A_i$, for this
fixed $j,w$, are pairwise disjoint.  Summing (LN.20) and using the
internal direct sum therefore gives
\[
 \sum_{i\in\operatorname{locs}(j,w)}
          \sum_{x\in Y_i}\operatorname{mult}_i(x)\le\dim M_j.
 \tag{LN.21}
\]

The required conversion from $\dim M_j$ to a theta-image rank has
a short complete proof.  Suppose $T:M_j\to H$ and $P:H\to M_j$
are the actual theta maps and, for each nonzero row summand,
$(PT)|_{E_{j,\sigma}}=\lambda_\sigma$ with
$\lambda_\sigma\ne0$.  If $Tv=0$, decompose
$v=\sum v_\sigma$.  Then $0=PTv=\sum\lambda_\sigma v_\sigma$;
independence of the summands forces every $v_\sigma=0$.  Hence
$T$ is injective and $\dim\operatorname{im}T=\dim M_j$.
This is the entire use of \texttt{Conclusion42} here.  It supplies
the actual scalars, not an unproved assertion that dimensions happen
to match.

Thus (LN.21) is \texttt{Mult\allowbreak Le}, with right side
$\operatorname{zeroSpaceRankAtPlace}(j,w)$.  The further equality
\[
 \operatorname{zeroSpaceRankAtPlace}(j,w)=d_j,
              \qquad d_j=L.dZ0(j),                     \tag{LN.22}
\]
is the source's \texttt{adm3}.  It follows, for example, by the same
injectivity argument for both this twist and the reference twist,
on the same $M_j$.  Where those scalar identities have not been
supplied, (LN.22) remains an explicit input.  The original
\texttt{GaloisPreconditions.what\textsubscript{2}.rank} supplies it
in the source join.  Neither the existence of a localisation nor
an identification of distinct theta spaces is silently assumed.

\paragraph{LN7: exact weighting, global cut loss, and denominator sum.}
\label{claim:LN7}
Let $\mathcal W$ be the actual finite support of twisted places and
put
\[
 a_w=\frac{e_wf_w\log\ell_w}{[F:\mathbb Q]}\ge0,
 \quad \mathcal E_j=\sum_{w\in\mathcal W}a_w\eta_J(j,w),
 \quad \mathcal V=\sum_{\ell\text{ prime}}\log\ell\,V_\ell.
 \tag{LN.23}
\]
Here $F=D.\mathrm{what}_2.\mathrm{LL}$, $e_w,f_w$ are its actual
ramification index and residue degree, and
$V_\ell=\sum_{1\le e<t}D.\mathrm{what}_2.\mathrm{shaLen}(\ell,e)$.
Thus $a_w$ is literally \texttt{lw\allowbreak P(w)*ram\allowbreak Idx(w)}, and
$\mathcal V$ is the source's \texttt{tsum vv\allowbreak Term}.  In particular
the ramification factor is present once.

The actual place identification is injective.  For each rational
prime $\ell$, the fundamental degree identity
$\sum_{w\mid\ell}e_wf_w=[F:\mathbb Q]$ gives
\[
              \sum_{w\in\mathcal W,\ w\mid\ell}a_w\le\log\ell.
 \tag{LN.24}
\]
One elementary proof of the degree identity is to compute the
dimension of $\mathcal O_F/\ell\mathcal O_F$ over $\mathbb F_\ell$.
It is $[F:\mathbb Q]$ because $\mathcal O_F$ is a free abelian group
of that rank.  Factorisation of $\ell\mathcal O_F$ and the Chinese
remainder theorem give local quotients of dimensions $e_wf_w$.
This requires no Galois hypothesis on $F$.

The actual tying equation is
$\eta_J(j,w)=\sum_{i\in\operatorname{locs}(j,w)}\eta_i$.
Apply LN5 and LN6, then (LN.24), and finally the prime factorisation
of the positive integer $N_G$.  This proves the exact global loss
\[
 \boxed{\quad\mathcal E_j(D)\le
             \mathcal E_j(\operatorname{cut}D)+d_j\log N_G.\quad}
 \tag{LN.25}
\]
Indeed the remaining scalar sum is bounded by
$\sum_\ell v_\ell(N_G)\log\ell=\log N_G$.  Discarded
localisations contribute zero by LN5.

We also give a direct local-to-global denominator version.  Suppose
there are nonnegative $\nu_\ell,\lambda_\ell$, a positive integer
$h_E$ (the source class number), and, off the two generator primes,
\[
 v\!c_i(x)\le\nu_{\ell_i}+\frac2{\ell_i^2}
                  +\lambda_{\ell_i}+V_{\ell_i}+v_{\ell_i}(h_E)
                       \quad(x\in Y_i).                \tag{LN.26}
\]
At a generator prime assume the actual weighted eta estimate
\[
 \sum_{\substack{w\in\mathcal W\\w\mid\ell}}a_w\eta_J(j,w)
 \le d_j\log\ell\left\{\frac2{\ell^2}+\lambda_\ell
                          +V_\ell+v_\ell(h_E)\right\}.
 \tag{LN.27}
\]
It need not arise from an individual congruence-ideal estimate.
Suppose $\sum_{\ell\in A}\nu_\ell\log\ell\le B_\nu$ for every
finite set of primes $A$, and that the nonnegative series
$\mathcal V$ is summable.  LN5, LN6 and (LN.24) yield
\[
 \boxed{\quad\mathcal E_j\le d_j\left\{
      4+B_\nu+\log h_E+\log N_G+
      \sum_{\ell\in\mathcal P}\log\ell\,\lambda_\ell
                       +\mathcal V\right\}.\quad}     \tag{LN.28}
\]
Here $\mathcal P$ is just the finite set of residue primes of
$\mathcal W$.  Nonnegativity permits extending finite valuation
sums to all prime factors and the finite $V_\ell$ sum to
$\mathcal V$.  The elementary bound
$\sum_{n\ge2}\log n/n^2\le2$ gives the constant $4$; the source
uses the harmless larger $6$.  This preserves coefficient exactly
$d_j$ on $\mathcal V$, an essential feature of the later scalar
join.  No bound on $\mathcal V$ or assertion of its summability is
being proved here.

The source's actual \texttt{Archimedean\allowbreak Signature\allowbreak Matching} in this
branch is the denominator inequality
\[
 v\!c_i(x)\le s_i(x)+\sum_{\lambda\mid\ell_i}b_{i,\lambda}
                 +V_{\ell_i}+v_{\ell_i}(h_E)+B_i+\ell_i^{-2}.
 \tag{LN.29}
\]
Together with $s_i(x)+\ell_i^{-2}\le\nu_{\ell_i}+2\ell_i^{-2}$
and an upper envelope for $\sum b_{i,\lambda}+B_i$, it is exactly
(LN.26).  Its name does not turn (LN.29) into a metric theorem.
The cyclotomic input used by the source is
$\nu_\ell\le v_\ell(QS)+\mathbf1_{\ell\mid QS}v_\ell(tr)$,
with zero at the distinguished excluded primes.  Factorisation then
gives $B_\nu=\log QS+\log(tr)$ directly.  The stated cyclotomic
image inequality, and (LN.29) itself, remain arithmetic inputs.
\paragraph{LN8: the complete numerical prime-envelope estimate.}
\label{claim:LN8}
We now prove the numerical export, including the small-order inert
primes.  This makes the input to LN9 a list of per-prime inequalities,
not a retained global summation theorem.  Use the actual parameter
record of PN16 and its notation from PN1.  Throughout this numerical
argument $n=D.R.n$, $f=D.R.f=2n$, $k=D.R.k$, $m=D.R.m$,
$r=D.R.r=k+m-1$, $t=D.R.t$, $e=D.R.e$ and $o=D.R.o$.
In particular $n$ is the source rank parameter, not a local degree
or the number of selected places.  Set
\[
 \mathcal L=x^\epsilon,\quad b=e^{\mathcal L},\quad
 b'=xb/e^{x^{\epsilon/3}},\quad
 w''=b/(n e^{x^{\epsilon/3}}),\quad S=L.\mathrm{Sc},\quad
 \eta=L.\mathrm{etaSz},\quad d_2=L.\mathrm{de2}.
 \tag{LN.30}
\]
The integer $N_m=k+m(t+1)$ used below is different from $N_G$.
Put
\[
 \begin{split}
 D_*&=nt^2(N_m+t)\mathcal L^2+kn^2(S+t+\mathcal L),\\
 z&=mb/D_*,\qquad \alpha=k\mathcal L/(md_2),\qquad
 H=b+S+\eta+t+2,\qquad \eta_3=\eta+t+2.
 \end{split}                                                   \tag{LN.31}
\]
These are respectively \texttt{Den}, \texttt{l\allowbreak Star}, \texttt{l\allowbreak A},
\texttt{h\allowbreak B}, and \texttt{eta\textsubscript{3}}.  The source's
$w_P$ and $G_P$ obey $w_P+\eta_3\le9b'$ and $G_P\le10b'$.
These follow from the actual window bounds:
$w_B\le3b'+\eta$, $S\le b'$, $t\le\eta\le b'$, with the
constant terms absorbed by $b'\ge1$.

Fix $C\ge1$, independent of $x$ and of every family index, which
bounds the constants and exponents in
\texttt{Parameters\allowbreak Polynomially\allowbreak Bounded} as well as the local table
constants.  Thus $k,n,t\le x^C$, $d_2\ge x^{-\epsilon}$,
$\log h_E\le Cx^C$, and $\log N_G\le Cx^C\mathcal L$.
The scale has $0<\epsilon<1$.  A sufficiently large threshold,
depending only on the fixed constants and $\epsilon$, gives
\[
 \begin{gathered}
 \mathcal L\ge4,\quad \log x\ge1,\quad
 1\le m\le k,\quad r\le2k,\quad f=2n,\quad n_V\le N_m+t,\\
 \mathcal L\le\alpha<z\le b/\mathcal L,\quad
 (3+\log\alpha)^2\le\mathcal L,\quad
 1\le H\le2b,\quad \mathcal L\le b',\\
 r/\alpha\le2m/\mathcal L,\qquad
 rb'\alpha\log x\le mb,\qquad
 r\eta_3\le4mw'',\qquad
 C(\eta_3+S+\mathcal L)\le w''/8.
 \end{gathered}                                                   \tag{LN.32}
\]
Here are the actual sources of these scalar inequalities, so their
uniformity is not hidden.  Positivity, $m\le k$, $r\le2k$ and
$n_V\le N_m+t$ are rank-record identities.  PN1, PN4--PN6 and PN15
give the scales and the parameter fields
\texttt{rb\allowbreak S}, \texttt{Neta}, \texttt{teta}, \texttt{k\allowbreak S\_\allowbreak le},
\texttt{logeta}, \texttt{log\allowbreak Del}, and \texttt{lstar1}.
The identity $zD_*=mb$, with $D_*\ge kn^2\mathcal L\ge m\mathcal L$,
gives $z\le b/\mathcal L$.  The formulas for $\alpha$ give
$\mathcal L\le\alpha\le x^{C+2\epsilon}$ and
$r/\alpha\le2m/\mathcal L$.
The \texttt{rb\allowbreak S} inequality is
$r(b'+S)\log x\,\alpha\le d_2mb$; discard $S$ and use
$d_2\le1$.  From $N_m\eta\le d_2mw''$, $r\le N_m$ and
$\eta_3\le4\eta$ obtain the penultimate bound.
For the last bound, \texttt{teta} gives $\eta\le w''/x$,
\texttt{k\allowbreak S\_\allowbreak le} and $k\ge\sqrt x$ give $S\le w''/\sqrt x$,
and $n\le x^C$ and exponential domination give
$\mathcal L\le w''/x$.  Thus its left side is at most
$6Cw''/\sqrt x$, which is at most $w''/8$ eventually.
All remaining comparisons in (LN.32) follow from
$\log x=o(x^a)$ for each fixed $a>0$.  For example
$\log\alpha=O(\log x)$, while $D_*$ is subexponential in
$x^{\epsilon/3}$ by the displayed bounds on $\log\eta$ and
$\log\Delta_{\rm cl}$.  Hence $\alpha<mb/D_*$ eventually.
This supplies finite, uniform thresholds, and requires no theorem
about a distribution of primes.

We use the following elementary integer sums, valid for $Y\ge2$
and $A\ge2$ (a finite subset has no larger sum):
\[
 \begin{aligned}
 \sum_{2\le n\le Y}\log n&\le Y\log Y,&
 \sum_{2\le n\le Y}\frac{\log n}{n}&\le(1+\log Y)^2,\\
 \sum_{2\le n\le Y}\frac{\log n}{n-1}&\le2(1+\log Y)^2,&
 \sum_{n\ge A}\frac{\log n}{n^2}&\le\frac{4(\log A+2)}A,\qquad
 \sum_{n\ge2}\frac{\log n}{n^2}\le2.
 \end{aligned}                                                  \tag{LN.33}
\]
For the second bound use $\log n\le\log Y$ and
$\sum_{n\le Y}1/n\le1+\log Y$; the third follows from
$(n-1)^{-1}\le2/n$.  The function $\log u/u^2$ decreases for
$u\ge2$.  Integral comparison with
$\int_A^\infty\log u/u^2\,du=(\log A+1)/A$ proves the fourth,
including a first term and replacing $A$ by its ceiling.  For the
last bound use the first term at $2$ plus the integral from $2$:
$\log2/4+(\log2+1)/2<2$.  The first bound just counts at most
$Y$ nonnegative terms, each at most $\log Y$.

Define
\[
 \begin{split}
 E(\ell)&=C\{1+H/\ell+\log H/\log\ell\},\\
 J(\ell)&=\max\{1,\lceil\log\alpha/\log(\ell-1)\rceil\},\\
 B(\ell)&=H/(\ell-1)^{J(\ell)+1}
                         +(G_P+1)(J(\ell)+1)+\mathcal L,\\
 F(\ell)&=(w_P+\eta_3)/\ell+S+\mathcal L+
                                                \eta_3 I(\ell),\\
 B_0&=Cfnt^2n_V\mathcal L^2.
 \end{split}                                                   \tag{LN.34}
\]
The indicator $I(\ell)$ is $1$ exactly when there is an actual
allowed band parameter $w''\le s\le w_B$ with
\[
 |\ell s-b|\le C(\ell\eta_3+\ell S+b/\ell+s+\ell\mathcal L).
 \tag{LN.35}
\]
This is the source's \texttt{band\allowbreak I}; replacing a particular
\texttt{band i lambda} by $I(\ell)$ only enlarges an upper bound.
For $\ell\ge2C$, (LN.32) and the upper half of (LN.35) give
\[
 \ell s\le b+\ell s/8+b/2+\ell s/2,
 \qquad w''\le s\le4b/\ell.
 \tag{LN.36}
\]
Consequently
\[
 r\eta_3 I(\ell)\log\ell
 \le16mb\frac{\log\ell}{\ell}
          +\mathbf1_{\ell\le2C}\,4mb\log\ell.       \tag{LN.37}
\]
For the small-prime alternative use $w''\le b'\le b$; if the
indicator is zero there is nothing to prove.

The raw local length inequalities in \texttt{Adm\allowbreak T} imply the
following envelope for an ordinary prime $\ell\ge5$, up to $z$
or in the exceptional set:
\[
 \Lambda(\ell)=B_0+
 \begin{cases}
 (2fm+f)E(\ell),&\ell\text{ split},\\
 4fCrB(\ell),&\ell\text{ inert},\ \ell<\alpha,\ o_\ell>4L_w,\\
 4fCr\{H/\ell^2+F(\ell)\}+2fCmE(\ell),
                    &\ell\text{ inert otherwise}.
 \end{cases}                                                   \tag{LN.38}
\]
Here $o_\ell$ is the least positive exponent with
$\ell^{o_\ell}\equiv\pm1\pmod{2^{e+2}}$.  The class alternatives
are the literal \texttt{Split\allowbreak K} and \texttt{Inert\allowbreak K} away from the
listed special primes.  To check (LN.38) from the individual place
inequalities, use
$\sum_{\lambda\mid\ell}n_\lambda=[E:\mathbb Q]=2f$ and
$\#\{\lambda\mid\ell\}\le2f$.  Thus
$\sum(n_\lambda+1)\le4f$.  The split estimate
$2\sum b_{i,\lambda}\le\sum(2n_\lambda m+1)E(\ell)$ gives
$(2fm+f)E(\ell)$ after division by $2$; the inert estimates give
the two other cases exactly.  This is where the original local
denominator estimates enter, and their truth is not deduced from
the scalar envelope.

We first sum (LN.38) for $5\le\ell\le z$.
The common terms are bounded by
\[
 \sum_{\ell\le z}\log\ell\, B_0\le Cfmb\mathcal L,
 \qquad \sum_{\ell\le z}\log\ell\, E(\ell)
                                      \le11Cb\mathcal L^2.
 \tag{LN.39}
\]
For the first, $znt^2n_V\mathcal L^2\le mb$ and
$\log z\le\mathcal L$.  For the second, expand
$\log\ell E(\ell)=C\{\log\ell+H\log\ell/\ell+\log H\}$.
Use (LN.33), $z\mathcal L\le b$, $H\le2b$ and
$\log H\le\mathcal L+1$.  The three sums are at most
$Cb$, $8Cb\mathcal L^2$ and $2Cb$, respectively.

For the small inert primes of large order, $\ell\ge5$ gives
$(\ell-1)^{J(\ell)}\ge\alpha$ and
$J(\ell)+1\le3+\log\alpha$.  Thus
\[
 B(\ell)\le\frac H{\alpha(\ell-1)}
                      +11b'(3+\log\alpha)+\mathcal L.
 \tag{LN.40}
\]
Multiplying by $r\log\ell$ and summing over $\ell<\alpha$,
the first term is at most
$2(r/\alpha)H(1+\log\alpha)^2\le8mb$.
The other terms are at most
$12rb'\alpha(3+\log\alpha)^2\le12mb\mathcal L$, using
$\mathcal L\le b'$ and $\log x\ge1$.
Multiplication by $4fC$ therefore gives at most
$80Cfmb\mathcal L^2$.
For the medium inert primes, (LN.33) and $r/\alpha\le2m/\mathcal L$
give
\[
 4fCrH\sum_{\ell\ge\alpha}\frac{\log\ell}{\ell^2}
      \le64Cfmb\le64Cfmb\mathcal L^2.                 \tag{LN.41}
\]

The small-order inert primes require a sparsity argument, which
we prove explicitly.  If $\ell$ is odd and
$\ell^o\equiv\pm1\pmod{2^a}$, write $o=2^vu$ with $u$ odd.
Odd powers do not change the valuations of $q-1$ and $q+1$ in
the respective odd-power factors.  If $v>0$, then
$\ell^{2^v}\equiv1\pmod8$, so a congruence to $-1$ modulo
$2^a$ with $a\ge3$ is impossible.  The factorisation into
successive differences of squares gives
\[
 v_2(\ell^{2^v}-1)=v_2(\ell-1)+v_2(\ell+1)+v-1.
 \tag{LN.42}
\]
Indeed the factors $\ell^{2^j}+1$ have valuation $1$ for
$j\ge1$, and exactly one of $\ell-1,\ell+1$ has valuation $1$.
It follows that $\ell\equiv\pm1\pmod{2^{a-v}}$ whenever that
modulus is nontrivial.  The cases $v=0$ or $a\le2$ follow directly
from the odd-power fact and the residues of odd integers.
Taking $a=e+2$, $o=o_\ell\le4L_w$ gives
\[
 \ell\equiv\pm1\pmod M,\qquad
 M=2^{e+2-\lfloor\log_2(4L_w)\rfloor}.               \tag{LN.43}
\]
At the thresholds in use the exponent is nonnegative.

Set $R=n/\mathcal L$, the literal \texttt{Level\allowbreak Star\allowbreak Exponents.\allowbreak R0}.
The field \texttt{radius\_\allowbreak ge\_\allowbreak sqrt\_\allowbreak ratio} gives
$\sqrt{k/m}\le R$, hence $r\le2k\le2mR^2$.
The field \texttt{k\_\allowbreak ge\_\allowbreak mx} gives $R\ge3$ eventually.
The frame field \texttt{ord\_\allowbreak le\_\allowbreak o\_\allowbreak mul\_\allowbreak two\_\allowbreak pow} supplies
$n\le o2^e$, and the size-datum field \texttt{o\_\allowbreak tie} supplies
$1\le o\le x^{\epsilon/100}$ (positivity of $o$ is a frame field).
Finally \texttt{Lw\_\allowbreak le}, together with $P.C_{Lw}\le C$, gives
$L_w\le C\log x$.
Eventually $4o^2\le\mathcal L$,
$2oL_w\le\mathcal L$ and
$(\log R+2)L_w\le\mathcal L$.  The first two and
$R\mathcal L\le o2^e$ imply
$R+1\le2^e/L_w$.  Since
$2^{\lfloor\log_2(4L_w)\rfloor}\le4L_w$, (LN.43) gives
$M\ge2^e/L_w\ge R+1$.  This proves both the claimed sign of
the exponent and the required spacing.  Each of the two progressions
$\pm1\pmod M$, after discarding $1$, starts at $u\ge M-1\ge R$.
For a decreasing $g(t)=\log t/t^2$ on $t\ge3$, interval comparison
gives
\[
 \sum_{h\ge0}g(u+hM)
 \le g(u)+\frac1M\int_u^\infty g(t)\,dt
 \le\frac{\log R}{R^2}+\frac{\log R+2}{RM}.
 \tag{LN.44}
\]
Both functions in the last bound decrease in $R\ge3$.
Combining the two progressions, $r\le2mR^2$ and
$n=R\mathcal L\le oML_w$, we obtain
\[
 r\sum_{o_\ell\le4L_w}\frac{\log\ell}{\ell^2}
 \le4m\log R+
       4mo\frac{(\log R+2)L_w}{\mathcal L}
 \le4m(\log R+o)\le4m\mathcal L^2.                 \tag{LN.45}
\]
Thus this part of (LN.38) costs at most
$32Cfmb\mathcal L^2$.  Counting these primes as an arbitrary
set of primes would lose the essential factor between $r$ and $m$.

For the flat terms, $rb'\le mb$ follows from (LN.32), and
$r(S+\mathcal L)z\le2mb$ follows from the second term of $D_*$.
Use (LN.33) and (LN.37).  Before the factor $4fC$, the four
contributions are at most
\[
 36mb\mathcal L^2,\quad2mb\mathcal L^2,\quad
 64mb\mathcal L^2,\quad4(2C+1)^2mb\mathcal L^2.
 \tag{LN.46}
\]
The last uses
$\sum_{2\le n\le2C}\log n\le(2C+1)^2$.
The extra $E(\ell)$ coefficient in (LN.38) is bounded by
$(3+2C)fm$, since $m\ge1$.  Combining (LN.39)--(LN.46), all
ordinary primes up to $z$ cost at most
\[
 \begin{split}
 &\{C+11C(3+2C)+80C+64C+32C\\
 &\hspace{35mm}+4C[102+4(2C+1)^2]\}\,fmb\mathcal L^2
           \le1000(C+1)^3fmb\mathcal L^2.
 \end{split}                                          \tag{LN.47}
\]
The last inequality follows by expanding the cubic polynomials;
it holds already for every $C\ge0$.

It remains to include the exceptional, small, ramified and generator
primes.  Suppose the actual table supplies the following raw bounds,
exactly the remaining clauses of \texttt{Adm\allowbreak T}:
\[
 \begin{array}{ll}
 \ell>z,\ \ell\notin\mathrm{Exc}:&
              \sum b_{i,\lambda}+B_i=0,\\
 \ell\in\mathrm{small}\text{ or }\ell=3:&
              \sum b_{i,\lambda}+B_i
                   \le Cfmb\log x+Cx^C\mathcal L^2,\\
 \ell=q\text{ or }d:&\sum b_{i,\lambda}+B_i\le2Cfmb,\\
 \ell=g_0\text{ or }g_1:&\lambda_\ell\le Cx^C\eta,\\
 |\mathrm{Exc}|\le(|TT|+2)(S+3\eta+t+8b/z+9),&
                    \ell\in\mathrm{Exc}\Longrightarrow\ell\le3b.
 \end{array}                                                    \tag{LN.48}
\]
The first line is used only off the generator primes and the
special cases in the following lines.  The generator line accompanies
the eta estimate (LN.27), not an assumption about each $v\!c_i(x)$.
In the source \texttt{small} is the set of prime divisors of $QS$,
with $QS\le n^2$ and $2\mid QS$; also $\log q,\log d\le\mathcal L$.
It follows that the small-prime logarithms sum to at most
$2C\log x+\log3$.  The second line therefore costs at most
$4C(C+1)fmb\mathcal L^2$ eventually: use
$\log^2x\le\mathcal L^2$ and absorb
$x^C\log x=o(b)$ in the second term.  The two ramified primes
cost at most $4Cfmb\mathcal L^2$.

For the remaining terms put
$U=x^{C+2}\exp(x^{\epsilon/4})$.  The parameter bounds give
\[
 \begin{gathered}
 S\le4U,\quad \eta,k,n,t,m,\mathcal L\le U,\quad
 f,r\le2U,\quad n_V,N_m+t\le4U^2,\\
 |TT|\le2U,\qquad D_*\le10U^7,\quad b/z\le10U^7.
 \end{gathered}                                                   \tag{LN.49}
\]
The first line follows from \texttt{logeta}, \texttt{log\allowbreak Del},
$S=4k\Delta_{\rm cl}$ and the rank identities; the second follows
by substitution into (LN.31) and the actual \texttt{TT\_\allowbreak card\_\allowbreak le}.
Thus (LN.48) gives $|\mathrm{Exc}|\le388U^8$.
If $\ell>z$, then $\ell>\alpha$, so the large-order small-inert
case of (LN.38) cannot occur.  Since $H\le2b$,
\[
 H/\ell\le20U^7,\quad
 H/\ell^2\le20U^7,\quad
 (w_P+\eta_3)/\ell\le90U^7,\quad
 \log H/\log\ell\le(\mathcal L+1)/\log2.
 \tag{LN.50}
\]
Together with (LN.49), these inequalities give
$\Lambda(\ell)\le K_1(C)U^9$ for an explicit polynomial
$K_1$, for instance $10^4(C+1)^2$.  Also
$\log\ell\le\log(3b)\le3U$ on the exceptional set.
Their whole contribution is at most
$1164K_1(C)U^{18}$.
The two generator primes satisfy
$\log g_k\le C\mathcal L+\log2$ by \texttt{g\_\allowbreak Y} and
$C_{YL}\le C$, so their total cost is at most
$2Cx^C\eta(C\mathcal L+\log2)=O_C(U^3)$.
Finally $4+B_\nu+\log h_E+\log N_G=O_C(U^3)$, using
$B_\nu\le2\log n+\log(tr)$.

For every fixed $a,A>0$, $A U^a/b\to0$: its logarithm is
$\log A+a(C+2)\log x+a x^{\epsilon/4}-x^\epsilon\to-\infty$.
The exceptional, generator and last scalar costs therefore sum to
at most $b\le fmb\mathcal L^2$ at one further uniform threshold.
Combining this with (LN.47)--(LN.48) proves, for \emph{every} finite
set of primes $\mathcal P$,
\[
 4+B_\nu+\log h_E+\log N_G+
        \sum_{\ell\in\mathcal P}\log\ell\,\lambda_\ell
       \le 2000(C+1)^3 fmb\mathcal L^2.               \tag{LN.51}
\]
Here $\lambda$ chooses the generator, small, ramified and ordinary
envelopes in that order, and is zero outside their supported cases.
If cases overlap, bounding by the sum of their nonnegative costs
proves the same estimate.  This completes the source's numerical
export with a harmless different threshold; it preserves the actual
constant $2000(C+1)^3$ and all original local denominator meanings.

\paragraph{LN9: restoration of the original family at one threshold.}
\label{claim:LN9}
Let $\mathrm{Fm}$ be the original \texttt{Fam\allowbreak P eps}.  For every
$x$ and $i\in\mathrm{Fm}.I_x$, choose the cut constructed in LN2
when $x$ is beyond the original threshold.  Below that threshold
use the all-points cut; no good-cut assertion is needed there.
The choice is pointwise, but all its subsequent estimates will
have constants independent of $x,i$.  A finite family of indices
is not assumed.

The actual \texttt{gd} construction keeps the family parameter,
index sets and original size constants, and replaces only the
local Hecke points/modules and their congruence losses as follows:
\[
 \begin{array}{c|c}
 \text{original data}&\text{data in the cut family}\\ \hline
 \mathrm{Loc},X_i,T_i,M_i&
 \{i:Y_i\ne\varnothing\},\ S_i,\ T_i|_{S_i},\ M_i/(M_i\cap V_{S_i^c})\\
 \eta_J(j,w)&\displaystyle\sum_{i\in\operatorname{locs}(j,w),\ Y_i\ne\varnothing}
                                                   \eta'_i\\
 F,\mathcal W,e_w,f_w,\mu_{\rm Tw},d_j&\text{the same entries}\\
 f,m,b,\mathcal L,N_c,N_G&\text{the same entries}\\
 \Omega,\mathrm{ldat}_\ell,\mathrm{shaLen},\mathrm{VV},\mathrm{vvTerm}
                                      &\text{the same entries}.
 \end{array}                                                    \tag{LN.52}
\]
The original global theta spaces and maps in the rank calculation
are retained.  This follows by reading the concrete
\texttt{nd\allowbreak Cut}, \texttt{lay\allowbreak Cut}, \texttt{set\allowbreak Raw}, \texttt{led\allowbreak Tr},
and \texttt{extra\allowbreak Tr} constructors: \texttt{nd\allowbreak Cut} takes the quotient
in (LN.11), while \texttt{set\allowbreak Raw} changes the eta entry and its ties.
The number-field embedding, the \texttt{ldat} and \texttt{sha\_\allowbreak eq}
entries are copied.  The norm choice is the transported original
\texttt{nch\allowbreak Tr}, and \texttt{extra\allowbreak Tr} transports its accompanying
identities.  Thus the finiteness and summability premises of
\texttt{Link\allowbreak N} are literally the same propositions in the two
families.  The twisted-place support is also the same, which
proves \texttt{tw\allowbreak Good\_\allowbreak gd} with the original threshold.

There is an especially simple bound for the cost of restoration.
Since $\mathrm{GL}_r(\mathbb F_g)$ is a subset of all $r$ by $r$
matrices, $|\mathrm{GL}_r(\mathbb F_g)|\le g^{r^2}$.
Equation (LN.4) gives
\[
       \log N_G\le fr^2(\log g_0+\log g_1).
 \tag{LN.53}
\]
If $k,n\le x^C$, $r\le2k$, $f=2n$ and
$\log g_k\le C\mathcal L+\log2$, then, for $\mathcal L\ge1$,
\[
 \log N_G\le16(C+1)x^{3C}\mathcal L
                     \le C_Nx^{C_N}\mathcal L,
 \quad C_N=\max\{1,3C,16(C+1)\}.
 \tag{LN.54}
\]
This proof needs no class-number bound.  The polynomial rank and
generator-prime logarithm bounds are uniform on the original
family, by its size and parameter records.  Since $\epsilon>0$,
$C_N\log x\le x^\epsilon$ eventually, so $x^{C_N}\le b$.
For every actual lift $1\le j\le m$, $fm\ge1$, and $\mathcal L\ge1$.
Therefore, with the source's exact budget
\[
 \mathcal B=fmb\mathcal L^2+(\log x)^2N_c,
          \qquad N_c\ge0,
 \quad \log N_G\le C_N\mathcal B.                    \tag{LN.55}
\]
Nonnegativity of $N_c$ follows directly from its product/quotient
definition with $x,n,d_2>0$, not from a congruence assertion.

Suppose the cut family has \texttt{Link\allowbreak N} constant $K$ and
threshold $x_c$, and suppose the multiplicity and rank inputs of
LN6 hold on the original family beyond $x_r$.  At
\[
 x\ge\max\{x_c,x_r,x_{\rm par},x_{\rm exp},
                                      \mathrm{Fm}.x_1+1\},
 \tag{LN.56}
\]
apply the cut estimate with the unchanged \texttt{Sha\allowbreak Fin\allowbreak Odd} and
summability premises, then (LN.25) and (LN.55).  For all actual
indices and all allowed $j$,
\[
 \mathcal E_j(D)
 \le d_j\{K\mathcal B+\mathcal V\}+d_j\log N_G
 \le d_j\{(K+C_N)\mathcal B+\mathcal V\}.
 \tag{LN.57}
\]
This proves the complete original-family
\texttt{link\allowbreak N\_\allowbreak of\_\allowbreak gd} implication.  The coefficient of
$\mathcal V$ stays exactly $d_j$; restoration costs only the
explicit polynomial budget, and never a new cohomological factor.

\paragraph{LN10: complete supplier and precise remaining frontier.}
\label{claim:LN10}
The preceding arguments give the following consolidated supplier.
On one fixed actual family, assume beyond one uniform threshold:
the source-compatible type and spectral data of LN2--LN3; the
row identifications, localisation uniqueness and theta scalar/rank
data of LN6; the parameter record of PN16 and the fixed polynomial
parameter bound in LN8.  Suppose, for every actual $x,i$ satisfying
\texttt{Sha\allowbreak Fin\allowbreak Odd}, there is a local denominator table with one
fixed constant $C$ satisfying (LN.29), its cyclotomic scalar bound,
the split/inert length bounds leading to (LN.38), the band condition
(LN.35), and all boundary/exceptional-prime bounds (LN.48).
The table is allowed to depend on $x,i$ and the finiteness premise;
its constant and threshold are not.

If these tables are on the original data, LN1--LN5 supply their
previously independent \texttt{adm1} requirement, and LN6 supplies
\texttt{adm2}.  Apply LN7 and LN8 to the finite set of actual
residue primes.  For every allowed $j$ and summable
\texttt{vv\allowbreak Term},
\[
 \mathcal E_j\le d_j\{2000(C+1)^3fmb\mathcal L^2+\mathcal V\}
       \le d_j\{2000(C+1)^3\mathcal B+\mathcal V\}.
 \tag{LN.58}
\]
This is \texttt{Link\allowbreak N} on the original family directly.  In the
source's preferred route the tables are instead on the actual
good-cut data, where the required local regularity is available.
There $Z'_i=Y'_i$, so LN4 alone bounds the cut congruence module
by its cut coordinate ideals; in particular its \texttt{adm1}
holds after adding the nonnegative $v_{\ell_i}(N_G)$ allowance.
The same numerical argument gives \texttt{Link\allowbreak N} on that cut
family, and LN9 gives (LN.58) on the original family with constant
$2000(C+1)^3+C_N$.

In both cases the order of quantifiers is exactly
\[
 \begin{gathered}
 \exists K,x_0\ \forall x\ge x_0\ \forall i\in\mathrm{Fm}.I_x:\\
 \mathrm{ShaFinOdd}(D_{x,i})\Longrightarrow
 \forall j\in\mathrm{lifts}(D_{x,i}),\\
 \mathrm{Summable}(\mathrm{vvTerm}_{D_{x,i}})
           \Longrightarrow\text{the original-family bound}.
 \end{gathered}
 \tag{LN.59}
\]
Choose $C$ first, increasing it to cover the finitely many original
constant bounds; the inequalities are monotone under this increase
for $x\ge1$.  Then choose the maximum of the finitely many
\emph{uniform} thresholds used above.  Neither step involves a
maximum over the family indices.  All cohomological terms remain
those on the original datum, with precisely the original finiteness
and summability guards.

Here is how this supplier meets the source's local-bound entry
point, without hiding a new analytic assertion in a record name.
A \texttt{Concl} supplies its actual representation family $\rho$,
the ordinary local lengths from its \texttt{Hub\allowbreak Data}, and a small
prime table $t_f$.  Its combined table has
$s_i(x)=\mathrm{Br.sPi}(\rho_i,x)$.  At primes in
\texttt{small}, set all length entries to zero and put their entire
supplied bound $\sum t_f.\mathrm{len}+t_f.BB$ in $B_i$.
At $3$ outside \texttt{small}, distribute that same total equally
among the finitely many places above $3$; this set is nonempty,
so summing recovers the original total exactly.  At $q,d$ put
the supplied \texttt{mid\allowbreak G} in $B_i$.  On the ordinary supported
range use the original Hub lengths and error bound, and put them
zero elsewhere.  Keep the Hub band and exceptional set and set
the generator cost to $Cx^C\eta$ and its scalar error to
$\ell^{-2}$.  This is the literal piecewise \texttt{tab}
construction.  The source's host disjointness and range conditions
make its cases agree with the preceding alternatives.

The pointwise denominator bounds survive this redistribution
because only their sum with $B_i$ occurs in (LN.29).
The split/inert estimates use the untouched ordinary entries;
the small-prime, $q,d$ and generator estimates are precisely their
corresponding supplied \texttt{Concl} clauses.  Large-prime
vanishing and the exceptional-set bounds are unchanged.  Thus a
uniform family of these actual local conclusions, with its fixed
parameter bounds, feeds (LN.59).  This expands the substantive
content of the source's table assembly, local-bound family export
and final \texttt{Link\allowbreak NFam} restoration.  Their exact declarations
are recorded in the provenance; none is used as an unexplained
numerical theorem.

The remaining leaves are narrower than \texttt{Link\allowbreak N} and remain
visible.  They are the actual denominator comparison (LN.29),
the cyclotomic scalar inequality, the four local length cases,
the band/exceptional bounds and the boundary-prime estimates,
all on compatible data; and the type/spectral/theta identification
inputs listed in LN2 and LN6.  In the source the first group is
supplied through \texttt{Norm\allowbreak Choice\allowbreak Facts}, local Galois and
integral $p$-adic comparison arguments, including the estimates
at $2$ and $3$.  The label \texttt{Hypotheses\textsubscript{2}}
bundles Fontaine, Artin reciprocity, Bloch--Kato, local duality,
local Euler characteristic, local finiteness and integral
comparison interfaces.  These names do not by themselves certify
that their exact quantified interfaces have already been matched
to classical theorems; any completed LD or other leaf can be
substituted separately in that input ledger.

The source's \texttt{Cut\allowbreak Gives\allowbreak Alt} has further requirements.
In addition to the good cut proved in LN2, it demands all norm-choice
facts and \texttt{Tie4} on the transported data.  These are the
literal \texttt{Norm\allowbreak Choice\allowbreak Facts} fields.  Those compatibility assertions
are not manufactured by the quotient construction.  LN1--LN10
close the integral selection, denominator accounting, prime
summation, rank and family-restoration chain from the explicit
local inputs just stated.  Their output is the original
\texttt{Link\allowbreak N} input in DA2; DA2 still separately requires the
cohomological bound, rank equality and actual metric comparison
specified there.  No claim of a completed RH/GRH proof or of an
unconditional local Galois construction follows from this chapter.


\section{Dependency closure after the odd-minus reconstruction}
\label{sec:dependency-audit-v24}

\paragraph{A complete supplier, with its classical leaves visible.}
The dependency chain is
\[
 \begin{aligned}
 \text{finite duality (LD)}&\longrightarrow
 \text{Jacobi/Kolyvagin construction (OS)}\\
 &\longrightarrow\mathtt{PsiBound}\longrightarrow
 \mathtt{ThmBcFree}(240)\quad\text{(OB)}.
 \end{aligned}
\]
The local and global finite-duality theorems, finite class field
theory, Chebotarev, Kummer and the other classical results stated in
those chapters remain identified citation leaves. Every bespoke
auxiliary-class, recurrence, length, special-value, coefficient,
negative-twist and quantitative step needed for this endpoint is
written out. OS1 also supplies the exact odd-primary class-group
annihilator used in OB5. OS9 is independent of OB2's Bernoulli
nonvanishing, so neither join introduces a circular dependency.

\paragraph{Immediate reduction of the cohomological frontier.}
Apply OB8 to the original field and local-data choices in VW3--VW4.
Its constant $240$ is independent of every varying field, prime,
twist and finite prime set allowed there. It therefore supplies
the \texttt{Thm\allowbreak Bc\allowbreak Free} quantifiers before the window and permitted
prime pair are chosen. The already proved VW argument now yields the
following statement. Fix a \texttt{BConst} record $K$ and $C_t\ge1$. Then
\[
 \mathtt{ThmBab}(K,C_t)+\mathtt{PointHV}
 \ \Longrightarrow\
 \mathtt{JointWin}(C_t)+\mathtt{ThmBc}(C_t)
 \ \Longrightarrow\
 \forall\text{ primitive }\psi,\ \mathtt{VWinAt}(\psi,C_t).
\]
With the unchanged QP22 prime-pair inputs and $C_t=401$, VW4 gives
the same fixed pair throughout the arbitrarily late three-halves
window with cutoff exponent $400$. This substitution removes one
deep supplier; it does not construct the two remaining ones or the
common arithmetic host. In particular, \texttt{Thm\allowbreak Bab} must control
all prescribed cyclotomic layers on one window, and
\texttt{Point\allowbreak HV} must use the original conductor and fields.

\paragraph{The exact congruence frontier.}
LN1--LN10 prove that the stated actual type/theta identities,
local Galois denominator tables and family scales imply
\texttt{Link\allowbreak N} on the original family. Its numerical summation
uses elementary estimates and a proved sparse-residue-class bound;
its cut restoration costs exactly $d_j\log|G|$ and preserves the
cohomological coefficient. Those parts are complete. Construction
of the actual local tables and compatible geometric state remains
open. The stronger \texttt{Cut\allowbreak Gives\allowbreak Alt} conclusion additionally
requires all \texttt{Norm\allowbreak Choice\allowbreak Facts} and \texttt{Tie4} fields;
the good-cut construction alone does not assert them.

\paragraph{What still separates this from the full proof.}
The remaining cohomological suppliers must be joined with the
actual common family, arithmetic gap and size/error packages,
genuine theta/local installation, global and spectral operator
comparisons, and exact lattice-metric transport. DA1 governs the
uniform threshold and same-family quantifiers, and DA2 retains
the exact metric premise. WB and the analytic endpoint are already
complete deductions from their stated contract. The updated
\texttt{dependency\_\allowbreak audit\_\allowbreak v0\_\allowbreak 24/obligation\_\allowbreak map.\allowbreak json}
records these AND/OR obligations separately from the historical
reference graph. Counts of claims or source links do not measure
how much of these substantive obligations has been discharged.

% Version 0.25. Exact native local-K interfaces, relative to classical finite local duality.
\section{Finite local duality on the specified local fields}
\label{sec:local-k-duality}

\status{This section supplies the three literal interfaces
\texttt{Local\allowbreak Finiteness\allowbreak K}, \texttt{Local\allowbreak Tate\allowbreak Duality\allowbreak K} and
\texttt{Local\allowbreak Euler\allowbreak Char\allowbreak K}.  The prime may be $2$, the coefficient modules
may be nonfree, and the coefficient prime equals the residue
characteristic.  Classical finite local duality and the local Euler
formula are accepted in their precise usual forms; all field,
cochain, coefficient, cyclotomic and cup-product identifications needed
to apply them to the source interfaces are written here.}

Fix a prime $\ell$, a finite intermediate field
\[
 \mathbb Q_\ell\subset K\subset\Omega=\overline{\mathbb Q}_\ell
       =\texttt{Padic\allowbreak Alg\allowbreak Cl}(\ell),\qquad
 d=[K:\mathbb Q_\ell],\qquad G=\operatorname{Aut}_K(\Omega).
 \tag{LK.1}
\]
The symbol $\Omega$ denotes the source's specified algebraic closure,
not its metric completion.  For $r>0$, put
$m=\ell^r$ and $R=\mathbb Z/m\mathbb Z$.
Let $M$ be any finite $R$-module, with the discrete topology and a
continuous $R$-linear action $\rho$ of $G$.  Its source cohomology is
denoted by $H^i(G,M)$ below.

\paragraph{Classical input boundary.}
The local-field fact used is uniqueness and completeness of the
extended absolute value on a finite extension of $\mathbb Q_\ell$;
see Milne, \emph{Algebraic Number Theory}, v3.08,
\href{https://www.jmilne.org/math/CourseNotes/ANT.pdf}{Theorem~7.38}.
The finite local cohomology input is finiteness, vanishing above degree
two, and perfect cup-product duality for finite discrete modules over
a characteristic-zero nonarchimedean local field.  Its dual coefficient
is $A^D=\operatorname{Hom}_{\mathbb Z}(A,\Omega^\times)$, with evaluation
and the usual local Brauer invariant.  See Milne,
\href{https://jmilne.org/math/Books/ADTnot.pdf}{\emph{Arithmetic Duality
Theorems}}, second edition, I, Theorem~2.1 and Corollary~2.3,
pp.~26--28.  Finally we use the local Euler formula, I, Theorem~2.8,
p.~31.  These statements apply when the module order is invertible in
$K$, which holds here even for $\ell$-power order.  They do not require
the order to be prime to the residue characteristic.  The cited
\href{https://www.jmilne.org/math/Books/add/ADT2006.pdf}{author's errata}
were checked; the correction concerning higher local fields is not a
hypothesis of the one-dimensional local-field statements used here.

\begin{proposition}[LK1: the actual field, cyclotomic action and cochains]
\label{prop:lk-identifications}
The field and group in (LK.1) are a characteristic-zero nonarchimedean
local field and its absolute Galois group, on the original underlying
objects.  Reduction of the source's $\ell$-adic cyclotomic character is
the character $\chi=\texttt{cyc\allowbreak K}(\ell,K,r)$ characterized by
\[
   g\zeta=\zeta^{\chi(g)}\qquad
   (g\in G,\ \zeta\in\mu_m(\Omega)).
   \tag{LK.2}
\]
For any primitive $m$th root $\zeta$, the map
$a\mapsto\zeta^a$ identifies the source representation
\texttt{char\allowbreak Rep}$(\chi)$ with $\mu_m(\Omega)$, in additive notation.
The source's $H^0,H^1,H^2$ and $H^n$ are the usual continuous
cohomology groups, with their original $R$-module structures.
\end{proposition}

\begin{proof}
The norm on $\Omega$ extends the norm of $\mathbb Q_\ell$; its
restriction to the finite extension $K$ is therefore the unique
extension in the cited local-field theorem.  It is complete and
discrete.  Here is also the lattice consequence that will be needed
for the Euler factor.  Choose a $\mathbb Q_\ell$-basis $b_1,\ldots,b_d$
and put $\Lambda=\sum_i\mathbb Z_\ell b_i$.  The coordinate maximum
norm and the restricted norm are equivalent.  Indeed, the upper bound
follows from the ultrametric inequality; for the lower bound minimize
$\|\sum_i x_i b_i\|$ on the compact set
$\{(x_i)\in\mathbb Z_\ell^d:\max_i|x_i|_\ell=1\}$.
The minimum is positive because the basis is independent.
Scaling a nonzero coordinate vector to this set proves the lower bound.
Consequently, for some integer $a\ge0$,
\[
 \ell^a\Lambda\subset\mathcal O_K\subset\ell^{-a}\Lambda.
 \tag{LK.3}
\]
The valuation ring $\mathcal O_K$ is a submodule of a finite free
$\mathbb Z_\ell$-module and contains a full lattice.  Since
$\mathbb Z_\ell$ is a principal ideal domain, it is free of rank $d$.
In particular its residue field is finite: it is a quotient of the
finite group $\mathcal O_K/\ell\mathcal O_K$.
Thus $K$ satisfies precisely the local-field hypotheses of the
classical theorems.  No completion changes its elements or its degree.

Every element of $\Omega$ is algebraic and separable over $K$, and
$\Omega$ is algebraically closed, so it is a separable closure of
this very $K$.  The source group $G$ is exactly
$\operatorname{Gal}(\Omega/K)$ with its Krull topology.  The subgroups
fixing finite extensions of $K$ give its neighborhood basis.  If one
describes the topology through the action on the valued algebraic
closure, the same basis results: each algebraic element has finitely
many conjugates, and a sufficiently small ball separates its chosen
value from the other conjugates.  This is also why every algebraic
$K$-automorphism is compatible with the local-field topology; the
unique extended norm is invariant under it.

Choose compatible primitive roots $\zeta_s$ of order $\ell^s$ with
$\zeta_{s+1}^{\ell}=\zeta_s$.  Write
$g\zeta_s=\zeta_s^{a_s(g)}$.  The residues $a_s(g)$ are units and
are compatible under reduction.  Their inverse limit is the
$\ell$-adic cyclotomic character.  A different compatible primitive
system multiplies its generators by a compatible system of units and
leaves every $a_s(g)$ unchanged.  Its reduction modulo $\ell^r$ is
therefore exactly the source definition of $\chi$, proving (LK.2)
for every $m$th root.  Its kernel fixes the finite extension $K(\mu_m)$,
so $\chi$ is continuous.  A primitive root exists in the specified
algebraic closure and its powers give a bijection $R\simeq\mu_m$;
(LK.2) proves equivariance.  Both finite coefficient groups are
discrete.  This does not assert that the roots lie in $K$.

For the given action, continuity means that every orbit map is
continuous.  Because $M$ is finite and discrete, the intersection of
the stabilizers of its elements is an open normal subgroup.  Thus the
action factors through a finite quotient, as required for a finite
discrete Galois module.

To check cohomology itself, the source uses the continuous
inhomogeneous cochains $C^n(G,M)=\operatorname{Cont}(G^n,M)$ and
the differential
\[
\begin{split}
 (\partial f)(g_1,\ldots,g_{n+1})={}&
 g_1f(g_2,\ldots,g_{n+1})\\
 &+\sum_{i=1}^{n}(-1)^i
 f(g_1,\ldots,g_i g_{i+1},\ldots,g_{n+1})\\
 &+(-1)^{n+1}f(g_1,\ldots,g_n).
\end{split}
\tag{LK.4}
\]
In degree zero this gives the invariant submodule.  In degrees one
and two the source spells out respectively
\[
 c(gh)=c(g)+g c(h),\qquad
 g z(h,k)-z(gh,k)+z(g,hk)-z(g,h)=0;
 \tag{LK.5}
\]
its boundaries are $g\mapsto gm-m$ and
$(g,h)\mapsto g b(h)-b(gh)+b(g)$, with $b$ continuous.
These are exactly (LK.4).  In every higher degree the source's
\texttt{Fin.\allowbreak contract\allowbreak Nth} terms multiply adjacent arguments, and the
last term deletes the final argument, giving the displayed formula.
Thus kernels and quotients agree, rather than merely having equal
cardinalities.  Forgetting $R$ to the underlying abelian group changes
neither the cochains nor these differentials.  Multiplication by
every element of $R$ is represented by multiplication by an integer,
so the identification preserves the $R$-module structures as well.
\end{proof}

\begin{proposition}[LK2: literal local finiteness]
\label{prop:lk-finiteness}
For all the original data above, $H^0(G,M)$, $H^1(G,M)$ and
$H^2(G,M)$ are finite, and $H^n(G,M)=0$ for $n\ge3$.
Consequently the source's \texttt{Local\allowbreak Finiteness\allowbreak K}
holds with its full quantifiers.
\end{proposition}

\begin{proof}
LK1 gives the actual finite local field, its specified separable
closure and a finite discrete continuous module.  Its order is a
positive integer and hence invertible in the characteristic-zero
field $K$.  The finite local cohomology theorem therefore gives
finiteness and vanishing for the continuous cohomology of these
objects.  LK1 identifies these groups with the source groups in each
degree, including the separate definition \texttt{Hn} in degrees
at least three.  A zero cohomology group has exactly one element,
which is the source's \texttt{Subsingleton} conclusion.  Neither
freeness over $R$ nor triviality of the action has been used.
\end{proof}

\begin{proposition}[LK3: invariant and all four literal cup bijections]
\label{prop:lk-tate}
\leavevmode\par
Put $P=H^2(G,\texttt{char\allowbreak Rep}(\chi))$ and
$M^D=\operatorname{Hom}_R(M,R)$ with action
\[
 (g\varphi)(x)=\chi(g)\varphi(g^{-1}x).
 \tag{LK.6}
\]
There is an $R$-linear isomorphism $j:P\simeq R$.  Each of the
following source maps is bijective:
\[
\begin{aligned}
 H^0(G,M)&\longrightarrow\operatorname{Hom}_R(H^2(G,M^D),P),\\
 H^1(G,M^D)&\longrightarrow\operatorname{Hom}_R(H^1(G,M),P),\\
 H^1(G,M)&\longrightarrow\operatorname{Hom}_R(H^1(G,M^D),P),\\
 H^2(G,M)&\longrightarrow\operatorname{Hom}_R(H^0(G,M^D),P).
\end{aligned}
\tag{LK.7}
\]
They are, in order, \texttt{cup\allowbreak H0H2Twisted\allowbreak Dual},
\texttt{cup\allowbreak H1H1Twisted\allowbreak Dual}, its \texttt{flip}, and
\texttt{cup20}, all with the source's identity group homomorphism.
Thus the literal \texttt{Local\allowbreak Tate\allowbreak Duality\allowbreak K} holds.
\end{proposition}

\begin{proof}
First consider coefficients carefully.  Every additive homomorphism
from $M$ to $\Omega^\times$ has image in $\mu_m$, because $mM=0$.
Via $a\mapsto\zeta^a$ from LK1 it becomes a unique map $M\to R$.
An additive map between $R$-modules is $R$-linear, since scalar
multiplication is integer multiplication modulo $m$.  Therefore
the classical Tate dual is the source's \texttt{TDual} with action
(LK.6), exactly \texttt{tw\allowbreak Dual}; the inverse on the argument is
essential.  Its action is continuous: the common open kernel of
$\rho$ and $\chi$ acts trivially on this finite dual.

No coefficient freeness is implicit.  Write a finite $R$-module as
a direct sum of $R/(\ell^{a_i})$, $1\le a_i\le r$; the zero module
corresponds to the empty sum.  On such a summand,
\[
 \operatorname{Hom}_R(R/(\ell^a),R)=\ell^{r-a}R.
 \tag{LK.8}
\]
The evaluation map into the double dual is injective: a nonzero
residue $x$ modulo $\ell^a$ is detected by the functional taking
$1$ to $\ell^{r-a}$.  Both modules have order $\ell^a$, so the
evaluation is an isomorphism.  Taking direct sums proves this for
all finite $R$-modules.  For the twisted duals it is equivariant,
since the two factors $\chi(g)$ and $\chi(g)^{-1}$ cancel in
evaluation.

The Kummer sequence and Hilbert~90 identify
$H^2(G,\mu_m)$ with the $m$-torsion of $\operatorname{Br}(K)$.
The local Brauer invariant identifies that torsion with
$m^{-1}\mathbb Z/\mathbb Z$.  Composing with the inverse of
\[
 \iota_m:R\longrightarrow m^{-1}\mathbb Z/\mathbb Z,
 \qquad a\longmapsto a/m,
 \tag{LK.9}
\]
and the LK1 coefficient identification gives $j:P\simeq R$.
This can be chosen before $M$ and $\rho$.  All its maps are additive
and hence $R$-linear.  The opposite common invariant sign, as in
LD's arithmetic-Frobenius normalization, multiplies $j$ by $-1$ and
leaves all the bijectivity conclusions unchanged.  The local-K
interface asks for an isomorphism, not an extra normalization.

For clarity the actual source cochain pairings in (LK.7) are as
follows.  If $x\in M^G$, $z_D$ is a dual-valued two-cocycle,
$f$ is a dual-valued one-cocycle, $c$ is an $M$-valued one-cocycle,
$z$ is an $M$-valued two-cocycle, and $\varphi\in(M^D)^G$, they are
\[
\begin{aligned}
  (x\cup z_D)(g,h)&=z_D(g,h)(x),\\
  (f\cup c)(g,h)&=f(g)\bigl(\rho(g)c(h)\bigr),\\
  (z\cup\varphi)(g,h)&=\varphi\bigl(z(g,h)\bigr).
\end{aligned}
\tag{LK.10}
\]
These follow directly from the source's evaluation pairing and its
continuous inhomogeneous cup products.  In the last formula the
usual action of $gh$ on the degree-zero second argument disappears
because $\varphi$ is invariant.  Thus these are the classical cup
products before the invariant, with the displayed order of the
arguments; the degree-one source map puts the dual first.

Finite local Tate duality, applied to $M$ and to $M^D$ with the
evaluation identification $M^{DD}\simeq M$, now says that composing
the first, second and fourth maps with the invariant gives perfect
pairings into $\mathbb Q/\mathbb Z$.  All cohomology groups here
are finite by LK2 and are killed by $m$, so every additive character
on them lands in $m^{-1}\mathbb Z/\mathbb Z$.  Equation (LK.9)
therefore identifies their character groups with their $R$-linear
duals.  This proves bijectivity after composing the target with $j$.
Postcomposition
\[
 \operatorname{Hom}_R(Y,P)\longrightarrow\operatorname{Hom}_R(Y,R),
 \qquad u\longmapsto j\circ u
 \tag{LK.11}
\]
is itself an isomorphism, with inverse $v\mapsto j^{-1}\circ v$.
Consequently the original maps into $P$, not only the scalar-valued
maps, are bijective.

Finally let $X=H^1(G,M^D)$ and $Y=H^1(G,M)$, and write
$b:X\to Y^*$ for the already proved scalar-valued isomorphism.
The flipped pairing is the composition
\[
 Y\xrightarrow{\ \mathrm{ev}\ }Y^{**}
   \xrightarrow{\ b^*\ }X^*,
 \qquad y\longmapsto\bigl(x\longmapsto b(x)(y)\bigr).
 \tag{LK.12}
\]
Both arrows are isomorphisms by finite biduality and bijectivity of
$b$.  Applying (LK.11) once more proves the exact \texttt{flip}
clause.  This uses the same bilinear form with its arguments exchanged;
no alternatingness or division by $2$ is involved.  It is valid also
when $\ell=2$.
\end{proof}

\begin{proposition}[LK4: literal residue-characteristic Euler factor]
\label{prop:lk-euler}
For the same original field and arbitrary finite coefficient module,
\[
 |H^1(G,M)|=|H^0(G,M)|\,|H^2(G,M)|\,|M|^{[K:\mathbb Q_\ell]}.
 \tag{LK.13}
\]
Hence the source's \texttt{Local\allowbreak Euler\allowbreak Char\allowbreak K} holds.
\end{proposition}

\begin{proof}
Put $q=|M|$.  From the finite module decomposition in (LK.8),
$q=\ell^s$ for some $s\ge0$.  LK1 proves that $\mathcal O_K$ is
free of rank $d$ over $\mathbb Z_\ell$.  Reduction of a basis gives
\[
 |\mathcal O_K/q\mathcal O_K|
   =|\mathbb Z_\ell/\ell^s\mathbb Z_\ell|^d
   =\ell^{sd}=q^d.
 \tag{LK.14}
\]
This computation uses the rank of the local field's ring of
integers; it does not identify $M$ with a free coefficient module.
It includes $s=0$, when the displayed quotient is the zero group
and has order one.

The classical local Euler formula for this finite discrete module is
\[
 \frac{|H^0(G,M)|\,|H^2(G,M)|}{|H^1(G,M)|}
     =|\mathcal O_K/q\mathcal O_K|^{-1}.
 \tag{LK.15}
\]
All three cohomology groups are finite by LK2, and each contains zero,
so their cardinalities are positive.  Substitute (LK.14) and clear
the positive denominators to obtain (LK.13).  LK1's identifications
preserve all cardinalities, and its $d$ is exactly the source's
\texttt{Module.\allowbreak finrank}$(\mathbb Q_\ell,K)$.  Thus (LK.13) is the
literal equality of natural cardinalities in \texttt{Local\allowbreak Euler\allowbreak Char\allowbreak K},
with the residue-characteristic factor retained.
\end{proof}

\paragraph{Closed interface and remaining scope.}
LK2, LK3 and LK4 can now be inserted directly as the three local-K
inputs wherever those literal predicates occur, including the
\texttt{Lemma\allowbreak K} and two-primary/reference-layer chains.  Their
quantifiers cover every finite intermediate $K$, every $r>0$, every
finite discrete $R$-module and every continuous action.  This does
not by itself prove the global integral-cohomology reduction, a
$K$-theoretic comparison, an arithmetic-group torsion bound, or the
odd-plus count/vanishing suppliers.  Those remain their own
mathematical chains.


% Version 0.25. Complete all-prime arithmetic-group torsion supplier.
\section{A complete torsion bound for the arithmetic homology group}
\label{sec:torsion-homology}

\status{TB1--TB10 prove the literal arithmetic-group supplier
\texttt{Tors\allowbreak Bound}$(49)$ and the explicit stronger specialization
\texttt{Tors\allowbreak Bound\allowbreak Pow}$(22)$.  The proof includes torsion at every prime.
It uses the source's ordered perfect-form nerve, but replaces the long
chain induction by a double-complex argument, the matrix normal-form
count by a row-module count, and the coordinate estimate by an inverse
Hermitian-matrix estimate.  No project-specific torsion, orbit-count,
acyclicity or stabilizer predicate is retained as an input.}

Write $T(A)$ for the torsion subgroup of an abelian group $A$.  In this
section homology has integral trivial coefficients unless specified.
All logarithms are natural.  The elementary background consists of
free resolutions, Smith normal form, finite-dimensional linear algebra,
and the integral lattice of a number field (for the latter, see
Milne, \href{https://www.jmilne.org/math/CourseNotes/ANT.pdf}{\emph{Algebraic
Number Theory}, Proposition~4.26}).  The exact constructions
that use this background are given below; in particular the acyclic
complex and the estimates entering the final bound are constructed.

\begin{proposition}[TB1: two integral torsion estimates]
\label{prop:tb-elementary-torsion}
Let $D:\mathbb Z^{(J)}\to\mathbb Z^a$ have columns of Euclidean norm
at most $b\ge1$, with $a<\infty$ and $J$ possibly infinite.  Then
$T(\operatorname{coker}D)$ is finite and has order at most $b^a$.
If $S$ is a finite group of order $s\ge1$, then, for every $q>0$,
\[
       |H_q(S,\mathbb Z)|\le s^{s^q}.
       \tag{TB.1}
\]
\end{proposition}

\begin{proof}
Let $L=\operatorname{im}D\subset\mathbb Z^a$ and let
$\overline L=(L\otimes\mathbb Q)\cap\mathbb Z^a$ be its saturation.
If its rank is zero the assertion is immediate.  Otherwise choose
$r=\operatorname{rank}L$ independent columns and let $L_0$ be their
span.  Choose $r$ coordinates on which these columns have nonzero
determinant.  Projection to these coordinates is injective on
$L\otimes\mathbb Q$, and consequently
\[
 [\overline L:L]\le[\overline L:L_0]
  =[\pi\overline L:\pi L_0]
  \le[\mathbb Z^r:\pi L_0]
  =|\det(\pi D_0)|\le b^r\le b^a.
\]
The penultimate bound is Hadamard's inequality, since coordinate
projection cannot increase a column norm.  The group
$\overline L/L$ is precisely the torsion of the cokernel.  This
also proves finiteness and does not require a bound on $|J|$.

For the second assertion use the inhomogeneous bar complex for $S$.
Its group in degree $q$ is free of rank $s^q$.  Its cycle subgroup is
therefore generated by at most $s^q$ elements, and so is its homology.
For completeness, the usual transfer to the trivial subgroup gives
\[
 H_q(S,\mathbb Z)\longrightarrow H_q(1,\mathbb Z)
       \longrightarrow H_q(S,\mathbb Z),
\]
whose composite is multiplication by $s$.  One constructs the first
map by summing all lifts of a simplex in the $s$-sheeted cover
$ES\to BS$ given by the bar construction; the boundary commutes
with this sum.  Projection sends each of those $s$ lifts back to
the original simplex, so the composite is multiplication by $s$.
For $q>0$ the intermediate homology is zero.  Hence $s$ annihilates
$H_q(S,\mathbb Z)$.  A group generated by $u$ elements and annihilated
by $s$ is a quotient of $(\mathbb Z/s)^u$, proving (TB.1), including
$s=1$.
\end{proof}

\begin{proposition}[TB2: torsion from an exact permutation complex]
\label{prop:tb-equivariant-bound}
Suppose a group $G$ acts on sets $X_p$, $p\ge0$, with equivariant
face maps deleting one of $p+1$ entries.  Suppose the augmented
complex
\[
 \cdots\longrightarrow\mathbb Z[X_1]\longrightarrow
 \mathbb Z[X_0]\longrightarrow\mathbb Z\longrightarrow0
 \tag{TB.2}
\]
is exact.  Fix $m\ge1$.  If there are at most $A$ orbits in each
$X_p$, $0\le p\le m$, and every stabilizer in those degrees is
finite of order at most $s\ge1$, then
\[
 |T(H_m(G,\mathbb Z))|
       \le (m+2)^A\,s^{m A s^m}.
 \tag{TB.3}
\]
In particular this torsion group is finite.
\end{proposition}

\begin{proof}
Take a free right $\mathbb ZG$-resolution $P_\bullet$ of $\mathbb Z$
and form the first-quadrant double complex
$P_q\otimes_{\mathbb ZG}\mathbb Z[X_p]$.  Exactness of (TB.2)
and freeness of each $P_q$ identify its total homology with
$H_*(G,\mathbb Z)$.  Computing the other direction first gives
\[
 E^1_{p,q}=\bigoplus_{[x]\in G\backslash X_p}
                    H_q(G_x,\mathbb Z)
       \ \Longrightarrow\ H_{p+q}(G,\mathbb Z).
 \tag{TB.4}
\]
Indeed, the permutation module of an orbit is
$\mathbb ZG\otimes_{\mathbb ZG_x}\mathbb Z$; restricting the free
resolution to $G_x$ proves the displayed identification directly.
Filtering the total complex by $p$ produces (TB.4); in each total
degree its filtration is finite, so the abutment has no convergence
qualification.

For $q>0$ and $p\le m$, TB1 bounds the corresponding group by
$s^{A s^q}$.  Its later-page subquotients have no larger order.
The row $q=0$ is the coinvariant complex
$\mathbb Z[G\backslash X_p]$.  In degree $m$ it has rank at most
$A$.  Every column of its incoming differential is a signed sum
of $m+2$ orbit basis elements, and hence has norm at most $m+2$.
TB1 bounds the torsion of its cokernel by $(m+2)^A$; the row's
$m$th homology is a subgroup of that cokernel, so it has the same
upper bound on torsion.  This is the required bound on
$T(E^2_{m,0})$.

For $r\ge2$ the homological differential has bidegree $(-r,r-1)$.
There is no incoming differential at $(m,0)$, since its potential
source has second index $1-r<0$.  Thus $E^\infty_{m,0}$ is a
subgroup of $E^2_{m,0}$, and the torsion bound persists.  The other
graded pieces on total degree $m$ are finite and bounded by their
$E^1$ groups.  In an exact sequence $0\to U\to V\to W\to0$ with
finite torsion subgroups, the map $T(V)\to T(W)$ has kernel in
$T(U)$, so $|T(V)|\le |T(U)|\,|T(W)|$.  Applying this to the
finite abutment filtration gives
\[
 |T(H_m(G,\mathbb Z))|
 \le (m+2)^A s^{A\sum_{q=1}^m s^q}
 \le (m+2)^A s^{m A s^m}.
\]
No stabilizer order has been inverted.  Thus this argument retains
small-prime torsion as well as all other torsion.
\end{proof}

Fix henceforth a totally complex number field $L$, with ring of
integers $\mathcal O$, degree $n=[L:\mathbb Q]$ and absolute
discriminant $D\ge1$.  Let $N\ge1$, $d=nN$ and
$\Gamma=\operatorname{GL}_N(\mathcal O)$.  Choose one embedding
$\sigma_w:L\hookrightarrow\mathbb C$ for each conjugate pair,
and put
\[
 V=\prod_w\mathbb C^N,\qquad
 E=\prod_w\operatorname{Herm}_N(\mathbb C),\qquad
 \Lambda=\mathcal O^N\subset V.
 \tag{TB.5}
\]
Thus $\dim_{\mathbb R}V=d$ and
$\dim_{\mathbb R}E=nN^2/2\le dN$.  For $h=(h_w)\in E$ set
\[
 q_h(y)=\sum_w \sigma_w(y)^*h_w\sigma_w(y),\qquad
 \mathcal R=\{h:q_h(y)\ge1\text{ for }0\ne y\in\Lambda\}.
 \tag{TB.6}
\]
Let $\mathcal P$ be the cone of tuples of positive semidefinite
matrices and $\mathcal P^\circ$ its positive definite interior.
Let $U\subset E^*$ be the cone of functionals strictly positive
on $\mathcal P\setminus\{0\}$.  It is open, nonempty and convex:
the sum of the matrix traces belongs to it, and positivity has a
positive minimum on the compact unit section of $\mathcal P$.

\begin{proposition}[TB3: the locally finite Hermitian polyhedron]
\label{prop:tb-hermitian-geometry}
The set $\mathcal R$ is nonempty, closed, convex and contained in
$\mathcal P^\circ$.  Its inequalities are locally finite there.
A point $v\in\mathcal R$ is an extreme point precisely when its
active functionals $h\mapsto q_h(y)$, with $q_v(y)=1$, span $E^*$.
For every $u\in U$, the minimum of $u$ on $\mathcal R$ is attained
at such a vertex.  Only finitely many vertices minimize a functional
in any fixed compact subset of $U$.
\end{proposition}

\begin{proof}
The embedding of $\mathcal O$ is a full lattice; its covolume is
$2^{-n/2}\sqrt D$.  This follows by taking an integral basis,
writing its complex embeddings as real and imaginary parts, and
comparing the square of the embedding determinant with the trace
discriminant.  The product embedding makes $\Lambda$ a full lattice
in $V$.  A sufficiently large scalar multiple of the identity form
therefore belongs to $\mathcal R$.  Closedness and convexity follow
from its defining half-spaces.

If $q_h$ is negative in a real direction, rational approximation in
a lattice basis and multiplication by a common denominator give a
nonzero lattice vector on which it is negative.  Thus $h\in\mathcal R$
is positive semidefinite.  It cannot be singular.  In fact, a
symmetric convex box which is arbitrarily long in its nullspace
and sufficiently thin in the complementary directions has
$q_h<1$ everywhere and arbitrarily large volume.  Minkowski's
convex-body argument then supplies a nonzero lattice point in it,
a contradiction.  The argument used here has the following short
form.  A measurable set of volume greater than the covolume of a
lattice has two distinct points with the same image in a fundamental
parallelepiped, by adding the measures of its translates in that
parallelepiped.  Apply this to half a symmetric convex body of
volume greater than $2^d$ times the covolume.  The difference of
the two points lies in the original body and is a nonzero lattice
point.  This proves the needed version of Minkowski's theorem.

Near a positive definite $h$, all forms have least eigenvalue at
least some $\delta>0$.  A defining inequality that can fail or become
an equality there must therefore have $\|y\|^2\le1/\delta$.
There are only finitely many such lattice vectors.  This proves
local finiteness.  If the active functionals at a point fail to
span, a nonzero direction annihilates all of them; sufficiently
small displacements in either sign remain feasible by local
finiteness.  The point is not extreme.  Conversely, if the active
functionals span and the point is the midpoint of two feasible
points, all its active equalities hold at both endpoints.  Their
difference is annihilated by a spanning set and hence is zero.

For $u\in U$ there is $c>0$ such that $u(h)\ge c\|h\|$ on
$\mathcal P$.  A nonempty sublevel of $u$ on $\mathcal R$ is
therefore compact.  Its minimum face is a nonempty compact convex
set.  Successively minimizing coordinates in a basis produces an
extreme point of that face, which is also extreme in $\mathcal R$.

Finally let $K\subset U$ be compact and fix $h_0\in\mathcal R$.
Uniformly on $u\in K$, one has $u(h)\ge c\|h\|$ on $\mathcal P$
and $u(h_0)\le M$, for some $c>0$ and $M<\infty$.  All minimizing
vertices lie in the compact set
$\mathcal R\cap\{\|h\|\le M/c\}$, which is contained in
$\mathcal P^\circ$.  On this set the least eigenvalue has a positive
minimum.  Thus all active lattice vectors belong to one finite set.
Each vertex is the unique solution to a full-rank selection of the
associated active equalities, so there are only finitely many.
\end{proof}

Let $\mathcal V$ be the vertex set from TB3.  Let $X_p$ consist of
ordered $(p+1)$-tuples of these vertices which simultaneously
minimize some $u\in U$.  Repeated entries are permitted.  These
are exactly the source sets \texttt{Sx}$(L,N,p)$: the forms,
active-constraint definition of a vertex, positive dual cone,
and minimization condition are the same.  Deleting one entry is
the face map.  The action is
\[
       (g\cdot h)_w=(\sigma_w(g)^{-1})^*h_w\sigma_w(g)^{-1}.
       \tag{TB.7}
\]
It preserves $\mathcal R$, vertices and common minimizing faces.

\begin{proposition}[TB4: acyclicity and a common minimal basis]
\label{prop:tb-nerve}
The augmented ordered-tuple complex (TB.2) for these $X_p$ is exact.
For every tuple in $X_p$, its common minimal vectors
\[
 \{y\in\Lambda:q_{v_j}(y)=1\text{ for all }j\}
 \tag{TB.8}
\]
span $L^N$ over $L$.
\end{proposition}

\begin{proof}
For each vertex $v$ let
$C_v=\{u\in U:u(v)\le u(h)\text{ for every }h\in\mathcal R\}$.
These are relatively closed convex subsets of $U$.  They cover $U$
and form a locally finite family by TB3.  A set of vertices is a
simplex of their nerve precisely when its members have a common
minimizing functional.

Here is a finite argument for the augmented homological acyclicity
of this nerve.  A finite cycle uses finitely many simplices.  Choose
a point in the intersection defining each of those simplices,
and let $K\subset U$ be a compact convex set containing these
finitely many points.  Only finitely many $C_v$ meet $K$.  Their
intersections $A_v=C_v\cap K$ are compact convex sets covering
$K$, and the original cycle belongs to their finite nerve.  For
all sufficiently small $\eta>0$, the open $\eta$-neighborhoods
of these sets have the same nerve.  Indeed, an empty intersection
of finitely many compact sets has a positive minimum of the
maximum distance to those sets; there are only finitely many
subfamilies to check.  The union of these neighborhoods is exactly
the $\eta$-neighborhood of $K$, which is convex, and every nonempty
intersection is convex.

The nerve of this finite open convex cover is homologically
acyclic.  One can see the assertion directly from the double
complex of singular chains on the intersections.  In one direction
the straight-line contractions of the nonempty intersections leave
exactly the augmented nerve complex.  In the other direction the
augmented cover complex is exact on chains small enough to lie in
one cover member: for each such singular simplex, its possible
indices form a nonempty simplex, contracted by adjoining any one
of its indices.  Repeated barycentric subdivision makes any finite
collection of chains this small, and the subdivision prism is a
chain homotopy.  The resulting homology is the homology of the
convex union, which vanishes after augmentation.  Thus the chosen
cycle bounds already in the finite nerve, and hence in the whole
nerve.  Nonemptiness follows from TB3.

To pass to ordered tuples, choose any order on the vertex set.
Map an ordered tuple with distinct entries to its increasing order
with the permutation sign, and map tuples with repetitions to
zero.  This is a chain map onto the oriented nerve complex, with
the increasing-tuples inclusion as a section.  On all tuples with
support in a fixed simplex, the augmented complex is contracted
by prepending one chosen vertex.  This proves the section followed
by the projection is homotopic to the identity: inductively in
degree, the difference on a basis tuple, after subtracting the
already constructed homotopy of its boundary, is a cycle supported
on that same simplex; fill it by the prepending contraction.
The ordered complex consequently has the same augmented homology,
and is exact.  The order and the contraction need not be equivariant;
the original ordered complex itself is a permutation complex with
the action (TB.7).

For the spanning assertion let $h$ be the average of the vertices
of the tuple.  Since every value $q_{v_j}(y)$ is at least one,
$q_h(y)=1$ holds exactly for the common minimal vectors (TB.8).
If these vectors did not span $L^N$, a nonzero $L$-linear functional
$a:L^N\to L$ would annihilate them.  At the complex places the
forms $|\sigma_w(a)y|^2$ define a nonzero $k\in\mathcal P$ with
$q_k(y)=0$ on every active constraint at $h$.  Local finiteness
then implies $h-\epsilon k\in\mathcal R$ for all sufficiently
small $\epsilon>0$.  More explicitly, choose a neighborhood with
only finitely many potentially active inequalities; the active
ones do not change, and each other one has a positive margin.
But the common minimizing functional $u$ has $u(k)>0$, and
$u(h-\epsilon k)<u(h)$ contradicts minimality.  This proves the
span assertion and permits a choice of $N$ independent common
minimal vectors in $\mathcal O^N$.
\end{proof}

Put
\[
       I_0=D^N d^d,\qquad R_0=N^N I_0^2,
       \qquad V_0=(2R_0+1)^d.
       \tag{TB.9}
\]
These are the source integers \texttt{I0}, \texttt{RR}, and the
near-vector bound.  They are positive.

\begin{proposition}[TB5: determinant and common-basis index bounds]
\label{prop:tb-determinant}
For $h\in\mathcal R$ one has
\[
                 1\le I_0\prod_w\det h_w.
                 \tag{TB.10}
\]
If $A\in\operatorname{Mat}_N(\mathcal O)$ is nonsingular and all
its columns are minimal for $h$, then
$\Delta=|N_{L/\mathbb Q}(\det A)|$ satisfies
$1\le\Delta\le I_0$.  Moreover the matrices
$H_w=\sigma_w(A)^*h_w\sigma_w(A)$ satisfy
\[
 H_w>0,\qquad (H_w)_{ii}\le1,\qquad
                 I_0^{-1}\le\det H_w\le1.
                 \tag{TB.11}
\]
\end{proposition}

\begin{proof}
Let $G_h$ be the real Gram matrix of $q_h$ in an integral basis of
$\Lambda$.  Its minimum on $\mathbb Z^d\setminus\{0\}$ is at
least one.  The inverse image under $G_h^{1/2}$ of the open cube
$(-d^{-1/2},d^{-1/2})^d$ contains no nonzero lattice point.
The convex-body argument in TB3 therefore bounds its volume by
$2^d$.  This gives $\det G_h\ge d^{-d}$.  The discriminant
calculation used in TB3 and the real determinant of a Hermitian
form give the exact identity
\[
        2^d\det G_h=D^N\left(\prod_w\det h_w\right)^2.
        \tag{TB.12}
\]
Consequently $\prod_w\det h_w\ge\sqrt{2^d/I_0}\ge I_0^{-1}$,
which proves (TB.10).

For a minimal column of $A$, the sum over $w$ of the corresponding
diagonal entries of $H_w$ is one.  These entries are nonnegative,
so each is at most one.  Hadamard's inequality gives
$0<\det H_w\le1$.  The product formula for the determinant is
\[
       \prod_w\det H_w
           =\Delta\prod_w\det h_w.
           \tag{TB.13}
\]
Here $\Delta$ is a positive integer.  Since the left side is at
most one, (TB.10) implies $\Delta\le I_0$.  Since all factors
on the left are at most one, each one is at least their product,
which by (TB.10) and $\Delta\ge1$ is at least $I_0^{-1}$.
This proves (TB.11).
\end{proof}

\begin{proposition}[TB6: a finite left-matrix quotient without normal forms]
\label{prop:tb-row-module}
For every integer $J\ge1$, the nonsingular matrices
$A\in\operatorname{Mat}_N(\mathcal O)$ with
$|N_{L/\mathbb Q}(\det A)|\le J$ have at most $J^{1+dN}$
orbits under left multiplication by $\Gamma$.
\end{proposition}

\begin{proof}
Let $M_A\subset\mathcal O^N$ be the module generated by the rows
of $A$.  Its index as an abelian subgroup is
$I=|N_{L/\mathbb Q}(\det A)|$.  This follows either by the real
embedding determinant, or by the determinant of the induced
$\mathbb Q$-linear map on $L^N$ in an integral basis.
Since the finite quotient has order $I$, it is annihilated by $I$,
so $I\mathcal O^N\subset M_A$.  The image of $M_A$ in
$(\mathcal O/I\mathcal O)^N$ is generated over $\mathcal O/I\mathcal O$
by its $N$ row residues.  There are at most
\[
             (I^{nN})^N=I^{dN}
\]
possible ordered lists of such residues, and each list determines
its generated submodule and hence its inverse image $M_A$.

Two nonsingular matrices have the same row module precisely when
their rows are two $\mathcal O$-bases of that module, and their
basis-change matrix belongs to $\operatorname{GL}_N(\mathcal O)$.
Thus they are in the same left orbit.  Summing the preceding bound
over $1\le I\le J$ gives
$\sum_{I=1}^J I^{dN}\le J^{1+dN}$.  No principal-ideal hypothesis
on $\mathcal O$ was used; the row modules here are free because
their displayed rows are independent over $L$.
\end{proof}

\begin{proposition}[TB7: near-vector and vertex counts]
\label{prop:tb-near-count}
Fix a nonsingular integral $A$ of determinant norm
$\Delta\le I_0$.  The union of all minimal vectors of all
$h\in\mathcal R$ for which the columns of $A$ are minimal has
cardinality at most $V_0$.  The number of such forms which are
vertices is at most $(V_0+1)^{dN}$.
\end{proposition}

\begin{proof}
The column lattice $A\mathcal O^N$ has index $\Delta$ in
$\mathcal O^N$.  Thus $\Delta\mathcal O^N\subset A\mathcal O^N$,
and there is an integral matrix $B$ with
$AB=BA=\Delta I$.  It is nonsingular.  Let $y$ be a vector in
the stated union and choose an associated form $h$.
At a complex place put $c=\sigma_w(A)^{-1}\sigma_w(y)$.
Then $c^*H_wc\le q_h(y)=1$.  The inverse-matrix Cauchy inequality
gives
\[
        |c_i|^2\le (H_w^{-1})_{ii}\,(c^*H_wc).
        \tag{TB.14}
\]
Indeed this is ordinary Cauchy--Schwarz for the inner product with
matrix $H_w$, applied to $c$ and $H_w^{-1}e_i$.
The diagonal cofactor of $H_w$ is at most one by Hadamard's
inequality and (TB.11), including the empty cofactor when $N=1$.
Therefore $(H_w^{-1})_{ii}\le I_0$ and
\[
       |\sigma_w((By)_i)|=\Delta |c_i|
                \le I_0^{3/2}\le R_0.
                \tag{TB.15}
\]
The same bound holds at conjugate embeddings.

There are at most $(2R_0+1)^n$ algebraic integers all of whose
conjugates have absolute value at most $R_0$.  To prove this,
reduce them modulo $Q\mathcal O$ where $Q=2R_0+1$.  If two
different such integers have difference $Qz$, the nonzero
integer $N_{L/\mathbb Q}(z)$ has absolute value at least one.
Some conjugate of $z$ consequently has absolute value at least
one, whereas the corresponding conjugate of the difference has
absolute value at most $2R_0<Q$, a contradiction.  Reduction is
injective, and $|\mathcal O/Q\mathcal O|=Q^n$ because
$\mathcal O$ is free of rank $n$ over $\mathbb Z$.
Applying this coordinatewise to $By$ and using injectivity of $B$
gives at most $Q^{nN}=V_0$ choices for $y$.

At a vertex the active functionals span $E^*$ by TB3.  Choose
at most $\dim E\le dN$ minimal vectors whose functionals form
a basis.  Their equations $q_h(y)=1$ determine the vertex
uniquely.  Each chosen vector lies in the preceding finite union.
Choose one ordered list for each vertex and pad to length $dN$
with a dummy symbol.  Two vertices with the same list agree by
the full-rank equalities.  Hence the number of vertices is at
most $(V_0+1)^{dN}$, as asserted.
\end{proof}

\begin{proposition}[TB8: all orbit and stabilizer budgets]
\label{prop:tb-orbits-stabilizers}
For every $p\ge0$, the $\Gamma$-orbits of $X_p$ are finite in
number and satisfy
\[
 \begin{split}
 |\Gamma\backslash X_p|
 &\le A_p:=I_0^{1+dN}(V_0+1)^{dN(p+1)}\\
 &\le (2Dd)^{12(p+1)d^3N}.
 \end{split}
 \tag{TB.16}
\]
Every tuple stabilizer is finite of order at most
\[
                         s=3^{dN}.
                         \tag{TB.17}
\]
\end{proposition}

\begin{proof}
TB4 supplies $N$ independent common minimal vectors for any tuple;
make them the columns of $A$.  TB5 gives $\Delta\le I_0$.
By TB6, a left translate $gA$ can be chosen from a list of at
most $I_0^{1+dN}$ matrices.  Under (TB.7) the minimal vectors
translate by $g$, so every vertex of the translated tuple belongs
to the near-vertex set for this representative.  TB7 gives at
most $(V_0+1)^{dN(p+1)}$ such ordered tuples.  This proves the
first bound, without requiring unique choices of common basis
or orbit representative.

For the numerical bound put $B_0=2Dd\ge2$.  Since $N\le d$,
$D,d,N\le B_0$, and
\[
 I_0\le B_0^{2d},\qquad R_0\le B_0^{5d},\qquad
 2R_0+1\le B_0^{5d+2},\qquad
 V_0+1\le B_0^{(5d+2)d+1}.
\]
The exponent in the resulting upper bound on $A_p$ is
\[
 2d(1+dN)+\bigl((5d+2)d+1\bigr)dN(p+1)
                    \le12(p+1)d^3N.
\]
For an explicit check, the terms $2d$, $2d^2N$,
$5d^3N(p+1)$, $2d^2N(p+1)$ and $dN(p+1)$ are bounded
respectively by $2,2,5,2,1$ times $(p+1)d^3N$.

A stabilizer of a tuple lies in the stabilizer of its first
positive definite form.  Under all complex embeddings that
stabilizer lies in a product of compact unitary groups.  Its
matrix entries lie in the lattice $\operatorname{Mat}_N(\mathcal O)$,
so it is finite.  A finite subgroup of $\Gamma$ injects under
reduction modulo $3\mathcal O$.  Here is the needed torsion-free
kernel argument.  If a nonidentity element of that kernel has
finite order, a power has prime order $q$; write it
$I+3^r C$, where $r\ge1$ is maximal and some integral entry of
$C$ is not divisible by $3$.  For $q\ne3$, the binomial expansion
of its $q$th power, reduced modulo $3^{r+1}$, has nonzero term
$q3^r C$.  For $q=3$, reduction modulo $3^{r+2}$ has nonzero
term $3^{r+1}C$; the remaining terms have factors
$3^{2r+1}$ and $3^{3r}$ and vanish in that quotient.  Both
contradict the claimed order.  The maximal $r$ exists because
$\mathcal O$ is a finite free abelian group.
Finally
$|\operatorname{Mat}_N(\mathcal O/3\mathcal O)|
=3^{nN^2}=3^{dN}$, which bounds the size of the finite subgroup
and proves (TB.17).
\end{proof}

\begin{theorem}[TB9: the literal arithmetic torsion bound]
\label{thm:tb-literal}
For every totally complex number field $L$ and every integer
$e\ge1$, set $m=2e$, $N=4e+4$ and $d=nN$.
The torsion subgroup of $H_m(\operatorname{GL}_N(\mathcal O),\mathbb Z)$
is finite, and
\[
 \log|T(H_m(\operatorname{GL}_N(\mathcal O),\mathbb Z))|
     \le \exp\!\bigl(49\,m d^3N\log(2Dd)\bigr).
 \tag{TB.18}
\]
This is exactly \texttt{Tors\allowbreak Bound}$(49)$.
\end{theorem}

\begin{proof}
Use the exact complex of TB4 in TB2.  The integers $A_p$ in
TB8 increase with $p$, so take $A=A_m$ and $s=3^{dN}$.
Put $X=m d^3N\log(2Dd)$.  Since $m\ge2$, (TB.16) gives
$A\le e^{24X}$.  Also $X\ge1$,
$\log(m+2)\le X$, and
$\alpha=m\log s=m dN\log3\le X$; these follow respectively
from $d\ge2$, $N\ge1$, $\log(m+2)\le m$ for $m\ge2$,
and $d^2\log(2Dd)\ge\log3$.
Taking logarithms in (TB.3) yields
\[
 \log|T(H_m)|
    \le A\bigl(\log(m+2)+m s^m\log s\bigr)
    \le e^{24X}(X+Xe^X)
    \le e^{27X}\le e^{49X}.
\]
In the penultimate step use $X\le e^X$ and
$2e^{2X}\le e^{3X}$ for $X\ge1$.
The source's \texttt{HGL} is this group homology with the
trivial integral representation, and \texttt{tors\allowbreak HGL} is its
full torsion subgroup.  Its \texttt{t\allowbreak Bound} is the right side
of (TB.18), with $n$ the original degree and $D$ the original
absolute discriminant.  The finiteness and inequality therefore
supply the literal predicate with all its field quantifiers.
\end{proof}

\begin{theorem}[TB10: an explicit power supplier for LemmaK]
\label{thm:tb-power}
For all the fields and integers in TB9,
\[
 \log|T(H_{2e}(\operatorname{GL}_{4e+4}(\mathcal O),\mathbb Z))|
        \le \exp\!\bigl((2en)^{22}\log(2D)\bigr).
        \tag{TB.19}
\]
Thus the source's \texttt{Tors\allowbreak Bound\allowbreak Pow}$(22)$ holds, with the
constant chosen before the number field and $e$.
\end{theorem}

\begin{proof}
Put $R=2en\ge2$.  Then $m\le R$,
$N=4e+4\le8e\le4R$, and $d=nN\le4R$.
For $R\ge2$,
$\log(4R)\le2R\log2$; the inequality holds at $R=2$ and the
derivative of the right side minus the left side is positive
thereafter.  Consequently
\[
 \log(2Dd)\le\log(2D)+\log(4R)
           \le3R\log(2D).
\]
It follows that
\[
 \begin{split}
 49m d^3N\log(2Dd)
 &\le 49R(4R)^3(4R)(3R)\log(2D)\\
 &=37632R^6\log(2D)
 \le R^{22}\log(2D),
 \end{split}
\]
since $R^{16}\ge2^{16}>37632$.
Insert this into TB9.  The right side is exactly
\texttt{k\allowbreak Bound}$(22,e,n,D)$, and the same torsion subgroup is
finite by TB9.  No field extension, replacement discriminant,
or excluded torsion prime is involved.
\end{proof}

\paragraph{Closed endpoint.}
TB10 supplies the complete torsion premise in the K-theoretic
reduction.  That reduction still needs its actual comparison
maps and the finite local/global cohomology argument; these
are supplied separately in the K-theory chapter.  The present
chapter alone makes no claim about a common arithmetic family,
the analytic gap, RH or Dirichlet GRH.


% Version 0.25.  The actual algebraic K-theory bridge and its quantitative join.
\section{From algebraic \texorpdfstring{$K$}{K}-theory to the uniform second-cohomology bound}
\label{sec:k-theory-bridge}

\status{KB1--KB4 construct the literal \texttt{KBridge} using classical
algebraic $K$-theory on the given number field.  KB5--KB8 close its
cardinality and constant chain: together with LD3, LK2--LK3 and TB10,
they prove the literal statement \texttt{Lemma\allowbreak K 31}.  Quillen finite
generation, Borel's rank theorem, Dwyer--Friedlander surjectivity,
Arlettaz's Hurewicz bound and integral homology stability are explicit
classical citation leaves with their hypotheses matched below.  None
is inferred merely from the name of a \texttt{Published} interface.}

Fix a prime $\ell$, a totally complex number field $L$, its specified
algebraic closure $\Omega$, and an arbitrary datum
$D:\texttt{LData}\ L\ \Omega\ \ell$.  Thus $S=D.S$ is exactly the set
of finite places above $\ell$, $D.A_v$ is the chosen valuation subring
above each $v\in S$, and $D.\chi:G_{L,S}\to\mathbb Z_\ell^\times$
pulls back to the cyclotomic character on $G_L$.
Assume \texttt{Has\allowbreak Mu\allowbreak Q}: $L$ contains a primitive $\ell$-th root if
$\ell$ is odd, and a primitive fourth root if $\ell=2$.
Throughout put
\[
\begin{gathered}
 n=[L:\mathbb Q],\quad \Delta=|\operatorname{disc}L|,\quad
 e\ge1,\quad m=2e,\quad N=4e+4=2m+4,\\
 X=H^2_{\rm cts}(G_{L,S},\mathbb Z_\ell(e+1)).
\end{gathered}
 \tag{KB.1}
\]
The cohomology in (KB.1) uses continuous inhomogeneous cochains and
the compact $\ell$-adic topology on $\mathbb Z_\ell$.
Write $h_L$ for the ordinary ideal class number, and
\[
 H=H_m(\operatorname{GL}_N(\mathcal O_L),\mathbb Z),\qquad
 T=H_{\rm tors},\qquad
 c(m)=\sum_{a=1}^{m}(a+2)=\frac{m(m+5)}2.
 \tag{KB.2}
\]
These are respectively the source's \texttt{HGL}, \texttt{tors\allowbreak HGL}
and \texttt{c\allowbreak One}; all group homology has trivial integral coefficients.

\subsection{Matching the original cochains and the arithmetic comparison}

\begin{proposition}[KB1: original datum and integral cochain comparison]
\label{claim:KB1}
\leavevmode\par
Let $U=\operatorname{Spec}\mathcal O_L[1/\ell]$ and let $j$ be an
integer.  With the character supplied by $D$, there are compatible
identifications
\[
 H^q_{\rm et}(U,\mathbb Z/\ell^r(j))
       \simeq H^q_{\rm cts}(G_{L,S},\mathbb Z/\ell^r(j)),
 \qquad
 H^2_{\rm cts}(G_{L,S},\mathbb Z_\ell(j))
       \simeq\varprojlim_r H^2_{\rm et}(U,\mathbb Z/\ell^r(j)).
 \tag{KB.3}
\]
The finite-level groups are finite; they vanish for $q\ge3$.
Moreover, directly on these cochains,
\[
 H^2_{\rm cts}(G_{L,S},\mathbb Z_\ell(j))/\ell
       \lhook\joinrel\longrightarrow
 H^2_{\rm cts}(G_{L,S},\mathbb F_\ell(j)).
 \tag{KB.4}
\]
No vanishing assertion about integral $H^3$ is needed for (KB.4).
\end{proposition}

\begin{proof}
The source's $G_{L,S}$ is the quotient of $G_L$ by the closed normal
subgroup generated by inertia at the omitted finite places.  Its
finite quotients are therefore exactly the finite Galois extensions
unramified outside $S$.  This identifies it with $\pi_1(U)$, using
the geometric point specified by $\Omega$.  Since $G_L\to G_{L,S}$
is onto, the identity $D.\chi\circ(G_L\to G_{L,S})=\chi_\ell$
determines $D.\chi$ uniquely.  It is continuous by the quotient
topology.  Consequently the action in \texttt{zl\allowbreak Tw D j} is precisely
$\chi_\ell^j$, including for negative $j$, and reduction gives the
usual finite Tate twist.  Restriction to each original $D.A_v$
commutes with these identifications.

The finite comparison is Milne, \emph{Arithmetic Duality Theorems},
II, Proposition~2.9, applied to the locally constant finite sheaf
$\mathbb Z/\ell^r(j)$; $\ell$ is invertible on $U$.
Finiteness is II, Theorem~2.13.  The higher-degree vanishing is the
finite global conclusion already matched in LD3, since $L$ has no
real places.\footnote{J.S. Milne, \emph{Arithmetic Duality Theorems},
2nd edition (2006), I, Theorem~4.10, and II, Proposition~2.9 and
Theorem~2.13; \url{https://jmilne.org/math/Books/ADTnot.pdf}.}

Here is the inverse-limit argument, to fix which integral cohomology
is meant.  Put $C_r^q=C_{\rm cts}(G_{L,S}^q,\mathbb Z/\ell^r(j))$.
The transition $C_{r+1}^q\to C_r^q$ is onto: any cochain has finitely
many values and can be lifted by choosing one residue lift for each
value.  The resulting lift is locally constant.  Also
\[
 C_{\rm cts}(G_{L,S}^q,\mathbb Z_\ell(j))=\varprojlim_r C_r^q;
 \tag{KB.5}
\]
a map into the inverse-limit topology is continuous exactly when
all its coordinate maps are continuous.  The short exact sequence
of complexes with product map
$(c_r)_r\mapsto(c_r-\operatorname{red}(c_{r+1}))_r$ gives
\[
 0\longrightarrow\varprojlim{}^1 H^{q-1}(C_r)
 \longrightarrow H^q(\varprojlim C_r)
 \longrightarrow\varprojlim H^q(C_r)\longrightarrow0.
 \tag{KB.6}
\]
For an inverse system of finite groups the product difference map
is onto: any finite set of its equations can be solved backwards
from its largest index, and compactness of the product of the finite
groups then solves all equations.  Thus its first derived limit
vanishes.  Apply this to the finite $H^1(C_r)$ to obtain (KB.3).

Finally the coefficient sequence
$0\to\mathbb Z_\ell(j)\xrightarrow{\ell}\mathbb Z_\ell(j)
\to\mathbb F_\ell(j)\to0$ is exact on continuous cochains.
Surjectivity uses the same finite-valued lift; division by $\ell$
on $\ell\mathbb Z_\ell$ is continuous, so exactness also holds in
the middle.  The cohomology exact sequence gives (KB.4).
In degree two its differential is
$dc(g,h)=g c(h)-c(gh)+c(g)$, exactly the source's differential and
continuous-coboundary quotient.  This matches the actual $X$, not
just an abstract group of the same order.
\end{proof}

\begin{proposition}[KB2: finite arithmetic K-group mapping onto the actual target]
\label{claim:KB2}
Let $A_m=K_m(\mathcal O_L)$ be the usual Quillen group, and set
$K_\ell=\mathbb Z_\ell\otimes_{\mathbb Z}A_m$.  Then $A_m$ and
$K_\ell$ are finite, and there is a surjective
$\mathbb Z_\ell$-linear map
\[
                    p:K_\ell\twoheadrightarrow X.
 \tag{KB.7}
\]
This includes $e=1$ and $\ell=2$.
\end{proposition}

\begin{proof}
Quillen's finite-generation theorem says that every
$K_i(\mathcal O_L)$ is finitely generated.  Borel's rank formula
gives rank zero for every positive even index.  Thus $A_m$ is a
finite abelian group, and $K_\ell$ is its $\ell$-primary part.
The groups here can be defined by
$K_i(\mathcal O_L)=\pi_i(B\operatorname{GL}_\infty(\mathcal O_L)^+)$
for $i\ge1$, the standard plus-construction model of Quillen
$K$-theory.\footnote{D. Quillen, ``Finite generation of the groups
$K_i$ of rings of algebraic integers'', \emph{Lecture Notes in
Mathematics} 341 (1973), 179--198, Theorem~1,
\url{https://doi.org/10.1007/BFb0067056}; original conference volume:
\url{https://www.sas.rochester.edu/mth/sites/doug-ravenel/otherpapers/bass-seattle.pdf}.
A. Borel, ``Stable real cohomology of arithmetic groups'',
\emph{Ann. Sci. ENS} 7 (1974), 235--272, Proposition~12.2,
\url{https://www.numdam.org/item/10.24033/asens.1269.pdf}.
The canonical plus model and its agreement with the exact-category
model are part of Quillen's $K$-theory; see also Arlettaz, cited in
KB3, Definition~3.12 and Remark~3.16.}

The arithmetic input is Dwyer--Friedlander's comparison.
Their Proposition~8.1 removes the primes over $\ell$ in finite
coefficients; Theorem~8.5 gives a surjection from algebraic to
\'{e}tale $K$-theory in degree $m\ge2$.  Their Proposition~8.2
also supplies the degree-two case directly.
For $\ell=2$ use the cofinal levels $r\ge2$ and the primitive
fourth root required by that theorem.  Our root is integral and
its square is $-1$.  Totally complex number fields have
$\ell$-cohomological dimension at most two, as required.
Their Proposition~5.2 supplies the finite-coefficient
\'{e}tale spectral sequence.\footnote{W.G. Dwyer and E.M. Friedlander,
``Algebraic and \'{e}tale $K$-theory'', \emph{Trans. Amer. Math. Soc.}
292 (1985), 247--280, Propositions~5.2, 8.1--8.2 and Theorem~8.5;
the completed conclusion is Theorem~8.7 and Remark~8.8.
\url{https://doi.org/10.1090/S0002-9947-1985-0805962-2}.
Author-uploaded full text:
\url{https://www.researchgate.net/publication/239378605_Algebraic_and_Etale_K-Theory}.}

We spell out the passage to compact coefficients.  The coefficient
exact sequence is
\[
 0\longrightarrow A_m/\ell^r A_m
 \longrightarrow K_m(\mathcal O_L;\mathbb Z/\ell^r)
 \longrightarrow K_{m-1}(\mathcal O_L)[\ell^r]
 \longrightarrow0.
 \tag{KB.8}
\]
The transition on its rightmost terms is multiplication by $\ell$,
as seen from the two-term free resolutions of $\mathbb Z/\ell^r$.
Since $K_{m-1}(\mathcal O_L)$ is finitely generated, its torsion has
bounded exponent, and that inverse limit is zero.  The leftmost
groups are finite and have vanishing first derived limit.  Hence
the inverse limit of the middle terms is $K_\ell$.

Write $E_r=K_m^{\rm et}(U;\mathbb Z/\ell^r)$.  The spectral sequence
has only cohomological columns $0,1,2$, by KB1, and only even
coefficient degrees.  In total degree $m=2e$ it therefore gives
\[
 0\longrightarrow H^2_{\rm et}(U,\mathbb Z/\ell^r(e+1))
 \longrightarrow E_r
 \longrightarrow H^0_{\rm et}(U,\mathbb Z/\ell^r(e))
 \longrightarrow0.
 \tag{KB.9}
\]
All these groups are finite.  The inverse limit of the last terms
is $H^0(G_{L,S},\mathbb Z_\ell(e))=0$.  Indeed the cyclotomic
image is infinite, since $[\mathbb Q(\mu_{\ell^r}):\mathbb Q]$
is unbounded whereas $[L:\mathbb Q]=n$ is fixed.  It contains a
unit $u$ with $u^e\ne1$; the torsion-free module
$\mathbb Z_\ell$ has no nonzero invariant under that unit.
By (KB.9) and KB1, $\varprojlim E_r\simeq X$.

The finite-level comparison maps in the preceding paragraph are
surjective and compatible.  A compatible target sequence has a
nonempty finite set of lifts at each level.  Every finite string
of compatibility conditions can be satisfied by lifting at its
largest level; compactness gives a compatible infinite string.
Their inverse-limit map is thus onto.  Combining it with (KB.8),
(KB.9) and (KB.3) constructs (KB.7).  All maps respect the scalar
actions, so the limit map is $\mathbb Z_\ell$-linear.
The cohomological-dimension assertion used for the fraction field
also follows from LD3: a finite-valued continuous cocycle on $G_L$
factors through an extension ramified at only finitely many places;
enlarge $S$ to contain those places and apply the higher-degree
vanishing there.  This uses no integral second-cohomology bound.
\end{proof}

\subsection{The homology target and the Hurewicz loss}

\begin{proposition}[KB3: integral stable Hurewicz map with the required kernel exponent]
\label{claim:KB3}
There is a $\mathbb Z_\ell$-linear map
\[
 h:K_\ell\longrightarrow\mathbb Z_\ell\otimes_{\mathbb Z}H
 \quad\hbox{such that}\quad
 h(x)=0\ \Longrightarrow\ \ell^{c(m)}x=0.
 \tag{KB.10}
\]
It is obtained from the canonical Hurewicz map and the actual block
stabilizations of the source.
\end{proposition}

\begin{proof}
We first check the ring hypothesis for stability.  For a Dedekind
domain $R$, let $(a_0,\ldots,a_t)$ be unimodular with $t\ge2$.
If all first $t$ entries vanish, $a_t$ is a unit, and adding it to
one of those entries shortens the row.  Otherwise choose $j<t$
with $a_j\ne0$, and choose $k<t$ different from $j$.  Only finitely
many maximal ideals contain $a_j$.  The Chinese remainder theorem
gives $b\in R$ which is $1$ modulo each such ideal containing
$a_k$, and $0$ modulo the remaining ones.  Replace $a_k$ by
$a_k+ba_t$, leaving the other first $t$ entries fixed.  A maximal
ideal containing the resulting row contains $a_j$.  If it does
not contain $a_k$, the choice $b=0$ is contradictory.  If it does,
the choice $b=1$ forces $a_t$ into it, contradicting the original
unimodularity.  This proves precisely \texttt{Stable\allowbreak Range\allowbreak Le R 2}.

Apply this to $R=\mathcal O_L$.  Van der Kallen's integral stability
theorem applies to block inclusions $\operatorname{GL}_a(R)
\hookrightarrow\operatorname{GL}_{a+1}(R)$.  In his Convention~2.2,
our stable range permits $\mathrm{sdim}=1$; his Theorem~4.11 gives
isomorphisms in degree $m$ for $a\ge2m+1$.
In particular all the maps starting at $N=2m+4$ are
isomorphisms.\footnote{W. van der Kallen, ``Homology stability for
linear groups'', \emph{Invent. Math.} 60 (1980), 269--295,
Convention~2.2 and Theorem~4.11 (p.~289).
The theorem states surjectivity for $a\ge2m+\max(1,\mathrm{sdim})-1$
and injectivity for $a\ge2m+\max(1,\mathrm{sdim})$.
Author's copy: \url{https://webspace.science.uu.nl/~kalle101/homstab1980.pdf}.}
The bar complex commutes with the union of these groups: each bar
chain involves finitely many matrices, all in some finite rank.
Filtered unions are exact for abelian groups, so
\[
 H\xrightarrow{\ \sim\ }
 H_m(\operatorname{GL}_\infty(\mathcal O_L),\mathbb Z).
 \tag{KB.11}
\]
The block maps are exactly \texttt{stab}, and the union is the
source's \texttt{GLinf}.  Singular homology of the nerve model
$B\operatorname{GL}_\infty$ agrees with this bar homology.

The canonical plus construction preserves integral homology.
Compose its degree-$m$ Hurewicz map with that identification to get
\[
 f:A_m\longrightarrow
 H_m(\operatorname{GL}_\infty(\mathcal O_L),\mathbb Z).
\]
Arlettaz's Corollary~7.19(a) states that $R_m\ker f=0$, with
$R_m=\prod_{a=2}^m L_a$ and every $L_a$ squarefree.\footnote{
D. Arlettaz, ``Algebraic $K$-theory of rings from a topological
viewpoint'', \emph{Publicacions Matem\`atiques} 44 (2000), 3--84,
Definitions~7.11 and 7.14 and Corollary~7.19(a), pp.~52--55.
\url{https://mat.uab.es/pubmat/fitxers/download/FileType:pdf/FolderName:v44(1)/FileName:44100_01.pdf}.}
Thus
\[
                   v_\ell(R_m)\le m-1\le c(m).
 \tag{KB.12}
\]
For an exact check of the source constant, $L_a$ is the product of
primes $p$ admitting a nonnegative integer sequence $(b_i)_{i\ge0}$
with sum $a$, $b_0\equiv0\pmod{2p-2}$,
$b_i\equiv0\text{ or }1\pmod{2p-2}$ for $i>0$, and
$p b_{i+1}\le b_i$.  These are \texttt{Arl\allowbreak Adm} and \texttt{arl\allowbreak L}.
The sequence is finitely supported because its sum is finite;
for $a>0$ one has $b_0\ge2p-2$, so the source's bound $p\le a$
omits no prime.  Also $L_1=1$, so its \texttt{arl\allowbreak R m} is this
same $R_m$.  Only squarefreeness of each factor is used in (KB.12).

The torsion-free abelian group $\mathbb Z_\ell$ is flat over
$\mathbb Z$.  Tensoring the kernel sequence for $f$ therefore shows
that $R_m$ annihilates the kernel of $1\otimes f$.  Write
$R_m=\ell^a u$, with $a=v_\ell(R_m)$ and $u$ an $\ell$-adic unit.
Multiplication by $u$ is invertible, so that kernel is killed by
$\ell^a$, hence by $\ell^{c(m)}$.  Compose $1\otimes f$ with
the tensor of the inverse of (KB.11).  This is (KB.10), including
all small primes; no prime is discarded from the homology group.
\end{proof}

\begin{proposition}[KB4: the actual KBridge supplier]
\label{claim:KB4}
The literal proposition \texttt{KBridge} holds.
For every original $\ell,L,\Omega,D,e$ satisfying its hypotheses,
one may take the finite $\mathbb Z_\ell$-module $K_\ell$ of KB2,
the surjection $p$ of (KB.7), and the map $h$ of (KB.10).
In particular $X$ is finite.
\end{proposition}

\begin{proof}
KB2 supplies the finite module and surjectivity.  KB3 supplies the
map (KB.10), with the specific homology group $H$ of (KB.2),
and the exponent $c(2e)$.
These are all the fields in \texttt{KBridge}.  Every coefficient
identification in KB1 used the original $D$ and $\Omega$, and no
new local places were selected.  Finiteness of $X$ follows by
surjectivity from the finite $K_\ell$.

The source files \texttt{KBridge\allowbreak Of\allowbreak KPublished} and
\texttt{KBridge\allowbreak Of\allowbreak KPublished\allowbreak W} specify packages of interfaces.
This proof constructs the existential bridge itself with canonical
$K$-theory.  It therefore does not need, or assert, a universal
comparison theorem for every arbitrary record satisfying the
source's definition of a \texttt{Plus} model.
\end{proof}

\subsection{Finite groups, strict local kernels, and the class number}

\begin{proposition}[KB5: the finite-module and tensor cardinality calculation]
\label{claim:KB5}
Suppose $K$ is a finite $\mathbb Z_\ell$-module,
$p:K\twoheadrightarrow X$ and $h:K\to W$ are linear, and
$\ell^c\ker h=0$.  Then
\[
                  |X|\le |X/\ell X|^c\,|\operatorname{im}h|.
 \tag{KB.13}
\]
If $W=\mathbb Z_\ell\otimes_{\mathbb Z}H$ and $T=H_{\rm tors}$
is finite, every finite $\mathbb Z_\ell$-submodule $Y\subset W$
satisfies $|Y|\le|T|$.  Finite generation of $H$ is not required.
\end{proposition}

\begin{proof}
Put $Y=p(\ker h)$.  Then $\ell^cY=0$ and $X/Y$ is a quotient
of $K/\ker h\simeq\operatorname{im}h$.  For a finite group $X$,
multiplication by $\ell^c$ has kernel and cokernel of equal order.
Moreover multiplication by $\ell^i$ maps $X/\ell X$ onto
$\ell^iX/\ell^{i+1}X$.  The filtration of $X/\ell^cX$ therefore
gives
\[
 |Y|\le|X[\ell^c]|=|X/\ell^cX|
                \le|X/\ell X|^c.
\]
Multiplication by $|X/Y|$ proves (KB.13), also when $c=0$.

For the tensor assertion, $H/T$ is torsion-free and hence flat over
$\mathbb Z$.  Tensor its flatness with the injection
$\mathbb Z_\ell\hookrightarrow\mathbb Q_\ell$ to see that
$\mathbb Z_\ell\otimes(H/T)$ embeds in
$\mathbb Q_\ell\otimes(H/T)$; it is thus torsion-free as an
abelian group.  Flatness of $\mathbb Z_\ell$ gives the exact sequence
\[
 0\longrightarrow\mathbb Z_\ell\otimes T
 \longrightarrow\mathbb Z_\ell\otimes H
 \longrightarrow\mathbb Z_\ell\otimes(H/T)\longrightarrow0.
\]
The image of finite $Y$ in the last group is zero.
Decomposing the finite abelian group $T$ into cyclic groups shows
that $\mathbb Z_\ell\otimes T$ is its $\ell$-primary part:
$\mathbb Z_\ell\otimes\mathbb Z/a\mathbb Z
\simeq\mathbb Z_\ell/a\mathbb Z_\ell$.
Its order is at most $|T|$, which proves the claim.
\end{proof}

\begin{proposition}[KB6: the generator bound on the same arithmetic datum]
\label{claim:KB6}
For every integer $j$, with the datum and roots of unity fixed above,
\[
 \left|H^2_{\rm cts}(G_{L,S},\mathbb Z_\ell(j))/\ell\right|
                         \le h_L\ell^n.
 \tag{KB.14}
\]
Here finiteness of the left side follows from KB1, independently of
the integral finiteness supplied by KB4.
\end{proposition}

\begin{proof}
Let $M=\mathbb F_\ell(j)$ and
$M^D=\operatorname{Hom}_{\mathbb F_\ell}(M,\mathbb F_\ell(1))$.
The dual action is $\chi_\ell(g)\varphi(g^{-1}x)$, so
$M^D=\mathbb F_\ell(1-j)$.  This is precisely the inverse-and-twist
convention in \texttt{zd}.  LD3 gives a perfect finite pairing
between the strict local kernels
$\operatorname{Sha}^2_S(M)$ and $\operatorname{Sha}^1_S(M^D)$;
their orders are equal.

We make the local-field join explicit.  Keep $A_v=D.A_v$ and let
$L_v^h\subset\Omega$ be the henselization selected by this
valuation, whose absolute Galois group is $\operatorname{Dec}(A_v)$.
Its completion is $L_v$.  Completion gives an equivalence of finite
separable extensions of $L_v^h$ and $L_v$, with the same embeddings
and degrees; this is the henselian completion comparison obtained
from Krasner's lemma.\footnote{This is the finite-extension
comparison already used in the setup of LD3.  For its precise
henselian hypotheses, see Milne, \emph{Arithmetic Duality Theorems},
I, Appendix A, Case~2, pp.~130--131: henselizations of local rings
of global fields are excellent, and completion preserves the
finite separable extensions and their Galois groups.}
Embed the finite $\mathbb Q_\ell$-extension $L_v$ into
$\texttt{Padic\allowbreak Alg\allowbreak Cl}(\ell)$ and denote its image by $K_v$.
Transporting the separable closures gives a continuous group
isomorphism, with continuous inverse,
\[
 \operatorname{Dec}(A_v)\simeq
 \operatorname{Gal}\bigl(\texttt{Padic\allowbreak Alg\allowbreak Cl}(\ell)/K_v\bigr).
\]
It preserves the action on every $\ell$-power root of unity;
therefore the restriction of the original cyclotomic character
becomes exactly \texttt{cyc\allowbreak K}.  Precomposition in each cochain
degree commutes with the differential, cup product and the
original restriction $\operatorname{Dec}(A_v)\to G_{L,S}$.
This is the source's
\texttt{decomposition\allowbreak Group\_\allowbreak mul\allowbreak Equiv\_\allowbreak local\allowbreak Galois\_\allowbreak cyclotomic}
join, on the chosen $A_v$.  LK1 identifies its local cochains,
and LK2--LK3, with $r=1$, consequently give
\[
 |H^2(L_v,M)|=|H^0(L_v,M^D)|\le\ell \quad(v\in S).
\]
The kernel of localization $H^2(G_{L,S},M)\to
\prod_{v\in S}H^2(L_v,M)$ is exactly
$\operatorname{Sha}^2_S(M)$.  Taking its image, which may be a
proper subgroup of the product, yields
\[
 |H^2(G_{L,S},M)|
       \le |\operatorname{Sha}^1_S(M^D)|\ell^{|S|}.
 \tag{KB.15}
\]
The sum of the products of ramification and residue degrees at
$\ell$ is $n$, and each product is at least one; hence $|S|\le n$.

The reduction of the cyclotomic character modulo $\ell$ is trivial:
all $\ell$-th roots already belong to $L$.  At $\ell=2$ this also
follows directly from $\mathbb F_2^\times=1$.
Thus $M^D$ is the trivial one-dimensional module.  Evaluation of a
dual linear functional at $1$ identifies it with $\mathbb F_\ell$
and commutes with all the original local restriction maps.
Its degree-one cocycles are continuous additive characters, and
its coboundaries are zero.  A strict local-kernel character is
zero on the chosen decomposition groups at every $v\mid\ell$.
Inflated to $G_L$, it is unramified outside $S$ by definition and
at $S$ by this zero restriction.  Conjugate decomposition groups
give the same conclusion, because its target is abelian.

Global unramified class field theory now factors this character
through $\operatorname{Cl}(\mathcal O_L)$.  There is no real-place
condition, since $L$ is totally complex.\footnote{The precise
classical input is the Artin identification of the maximal abelian
extension unramified at all places with the ordinary ideal class
group.  See J.S. Milne, \emph{Class Field Theory}, v4.03 (2020),
Theorem~0.4 and V, Example~3.9,
\url{https://www.jmilne.org/math/CourseNotes/CFT.pdf}.}
Consequently
$\operatorname{Sha}^1_S(M^D)$ injects into
$\operatorname{Hom}(\operatorname{Cl}(\mathcal O_L),\mathbb F_\ell)$.
Each cyclic factor of the class group contributes either $1$ or
$\ell$ to the latter order, never more than its own order.
Thus $|\operatorname{Sha}^1_S(M^D)|\le h_L$.
Combine this with (KB.15), $|S|\le n$, and the injection (KB.4).
This proves exactly the bound used by \texttt{card\_\allowbreak mod\_\allowbreak le},
with no switch to a larger number field or a different $D$.
\end{proof}

\subsection{The complete quantitative join}

\begin{proposition}[KB7: the literal reduction inequality]
\label{claim:KB7}
If $T$ of (KB.2) is finite, then $X$ is finite and
\[
 \log |X|\le
 4n\log(2\Delta)\,(2e)(2e+3)\log(n+1)+\log|T|.
 \tag{KB.16}
\]
All constants in this inequality are independent of
$\ell$, $\Omega$, and the local choices in $D$.
\end{proposition}

\begin{proof}
We include the elementary class-number estimate needed to control
the first factor in (KB.13).  Minkowski's theorem supplies an
integral ideal in each class of norm at most
$\frac{n!}{n^n}(4/\pi)^{r_2}\sqrt\Delta\le\sqrt\Delta$.
For the last inequality use $r_2\le n/2$ and, for $n\ge2$,
$n!\le((n+1)/2)^n$; the resulting factor is at most
$(3/(2\sqrt\pi))^n<1$.  For $n=1$ it equals one.\footnote{
Minkowski's ideal-class representative bound is stated with this
normalization in J.S. Milne, \emph{Algebraic Number Theory},
v3.08 (2020), Theorem~4.3,
\url{https://www.jmilne.org/math/CourseNotes/ANT.pdf}.}
At a rational prime the number of prime ideals is at most $n$ and
their residue degrees are positive.  Its ideal-counting Euler
factor is therefore coefficientwise bounded by $(1-t)^{-n}$.
Multiplication over rational primes shows that the number of
ideals of norm $a$ is at most the number $d_n(a)$ of ordered
factorizations $a=a_1\cdots a_n$ into positive integers.
For $B\ge1$, summing one variable at a time and using
$\sum_{a\le B}1/a\le1+\log B$ gives
\[
 \sum_{a\le B}d_n(a)\le B(1+\log B)^{n-1},\qquad
 \log h_L\le\frac12\log\Delta
 +(n-1)\log(1+\tfrac12\log\Delta)
 \le\frac n2\log\Delta.
 \tag{KB.17}
\]
Here $\Delta\ge1$ and $\log(1+t)\le t$ for $t\ge0$.

Put $a=|X/\ell X|$.  KB6 and (KB.17), with $\ell\ge2$, imply
\[
 \log a\le\log h_L+n\log\ell
                  \le4n\log(2\Delta)\log\ell.
 \tag{KB.18}
\]
For example, $\log\ell\ge\log2>1/2$ bounds
$\tfrac12\log\Delta$ by $\log\Delta\log\ell$, and
$4\log2>1$ supplies the remaining $\log\ell$ term.
Also $\ell\le n+1$.  For odd $\ell$, the primitive root gives
$\mathbb Q(\mu_\ell)\subset L$ and degree $\ell-1\le n$;
the degree formula follows from Eisenstein applied to
$\Phi_\ell(Y+1)$.  For $\ell=2$ use $n\ge1$.

KB4 and KB5 now give
$|X|\le a^{c(2e)}|T|$.  All group orders are positive, so
logarithms preserve the inequality.  Since
$c(m)=m(m+5)/2\le m(m+3)$ and
$\log\ell\le\log(n+1)$, (KB.18) gives (KB.16).
This is the numerical inequality in \texttt{Reduction.\allowbreak reduction},
including its homology rank $4e+4$ and cohomological degree $2e$.
\end{proof}

\begin{proposition}[KB8: one universal constant and the literal LemmaK endpoint]
\label{claim:KB8}
For any real $C$, \texttt{Tors\allowbreak Bound\allowbreak Pow C} implies
\[
              \texttt{Lemma\allowbreak K}\bigl(\max(C,0)+9\bigr).
 \tag{KB.19}
\]
In particular TB10, which supplies \texttt{Tors\allowbreak Bound\allowbreak Pow 22}, gives
$\texttt{Lemma\allowbreak K 31}$.  Explicitly, for every original arithmetic
datum in (KB.1),
\[
 X\text{ is finite},\qquad
 \log|X|\le
       \exp\!\left((2en)^{31}\log(2\Delta)\right).
 \tag{KB.20}
\]
Thus $\exists C_K>0\;\texttt{Lemma\allowbreak K}\ C_K$, with $C_K$ chosen
before $\ell,L,\Omega,D,e$.
\end{proposition}

\begin{proof}
Set $x=2en\ge2$, $A=\log(2\Delta)>0$, and
$p=x^{\max(C,0)}\ge1$.  The hypothesis gives finiteness of $T$
and $\log|T|\le\exp(x^C A)$.  The coefficient of $A$ in the
first term of (KB.16) satisfies
\[
 Q=4n(2e)(2e+3)\log(n+1)
 \le 4n(2e)(5e)n=10x^2\le x^9-1.
 \tag{KB.21}
\]
The first inequality uses $e\ge1$ and $\log(n+1)\le n$.
For the last, $x^7\ge2^7$ and $x^2\ge4$ give
$x^9-10x^2=x^2(x^7-10)\ge472>1$.
Therefore $QA\le B:=(x^9-1)pA$, while
$\exp(x^CA)\le\exp(pA)$.  Since $B\ge0$ and $\exp(pA)\ge1$,
\[
 \exp(pA)+B\le\exp(pA)(1+B)
                \le\exp(pA+B)
                =\exp\!\left(x^{\max(C,0)+9}A\right).
\]
Combining this with KB7 proves (KB.19).  The exponent is positive,
and it is a function only of the single constant $C$ in the
torsion theorem.  Substitution $C=22$ proves (KB.20) and gives the
literal quantified statement required downstream.
\end{proof}

\paragraph{Scope of the closure.}
KB4 is an actual supplier, rather than an implication whose
\texttt{KBridge} premise remains unfilled.  KB8 closes the whole
shared \texttt{Lemma\allowbreak K} branch once the separately proved TB10 is
joined: finite arithmetic $K$-theory, the correct continuous
cohomology, the finite-rank integral homology target, every kernel
exponent, the strict local conditions, and the universal constant
have all been connected.  The classical leaves just cited remain
citation leaves; no proof of those foundational theorems from
first principles is claimed.  This component asserts neither a
full RH/GRH reconstruction nor a replay of the Lean checker.


\section{One simultaneous window for the two-primary and reference-layer terms}
\label{sec:simultaneous-cohomological-windows}

\status{TW1--TW8 reconstruct the complete supplier
$\exists K\,\forall C_t\ge1,\ \texttt{Thm\allowbreak Bab}\ K\ C_t$.
The input from the quantitative cohomology branch is the literal
$\exists C_K,\ \texttt{Lemma\allowbreak K}\ C_K$, supplied by KB8 with $C_K=31$;
finite Poitou--Tate duality
is supplied by LD3 and its actual local-field hypotheses by LK2--LK3.
Ferrero--Washington is used in its classical $\mu=0$ form, and its
translation to the precise uniform class-rank statement in the source
is proved below.  Every tower, exceptional-twist, common-window and
field-model step is included.  This chapter does not assume
\texttt{Fixed\allowbreak In}, \texttt{Two\allowbreak In}, \texttt{ABound}, \texttt{First\allowbreak Bound}
or \texttt{Thm\allowbreak Bab} as an additional hypothesis.}

For a number field $F$, a prime $\ell$ and $k\ge1$, put
\[
 A_\ell(k)=(\mathbb Q_\ell/\mathbb Z_\ell)(-k),\qquad
 s_\ell(F,k)=\operatorname{length}_{\mathbb Z_\ell}
 \ker\!\left(H^1(G_{F,S_\ell},A_\ell(k))
       \longrightarrow\prod_{v\mid\ell}H^1(F_v,A_\ell(k))\right).
 \tag{TW.1}
\]
Here $S_\ell$ consists of exactly the primes above $\ell$.
All cochains with these torsion coefficients are continuous for the
discrete topology.  Thus $s_\ell(F,k)$ is the source's
\texttt{sha\allowbreak Inv F l k}, and also \texttt{sha\allowbreak Len D k} for every
algebraic closure and every admissible local-data record $D$ on the
given $F$.  At first, lengths may be infinite.  The compact dual
coefficient is $T_\ell(k)=\mathbb Z_\ell(k+1)$, and
\[
 V_\ell(F,k;D)=
 \operatorname{length}_{\mathbb Z_\ell}H^2(G_{F,S_\ell},T_\ell(k))
 \tag{TW.2}
\]
is \texttt{VLen D k}.

Write $\operatorname{VV}_\ell(F,t,B)$ for the assertion, for every
such closure and $D$, that all the lengths in (TW.1) with
$1\le k<t$ are finite and that their sum is at most $B$.
This is exactly \texttt{VVat}; at $\ell=2$ it is
\texttt{Bounded\allowbreak Height\allowbreak Two\allowbreak Place\allowbreak Condition}.
An \emph{$m$th cyclotomic layer} of $F$ is a field
$F(\mu_{\ell^a})$ of degree $\ell^m$ over $F$ for some $a$,
allowing any isomorphic field model equipped with its $F$-algebra
structure.  This is the literal \texttt{Is\allowbreak Layer l m F F'}.

\subsection{Named classical inputs}

Only the following classical inputs are retained in this chapter;
their finite-layer, coefficient and quantitative consequences are
proved in TW1--TW8.

\begin{description}[style=nextline,leftmargin=0pt,labelwidth=0pt,labelsep=0pt]
\item[TW-FW: Ferrero--Washington.]
For every finite abelian number field $F/\mathbb Q$ and every prime
$\ell$, the unramified Iwasawa module of its cyclotomic
$\mathbb Z_\ell$-extension has $\mu=0$.
This includes $\ell=2$.
See B.~Ferrero and L.~C.~Washington,
\emph{The Iwasawa invariant $\mu_p$ vanishes for abelian number fields},
Annals of Mathematics \textbf{109} (1979), 377--395,
\url{https://doi.org/10.2307/1971116},
and the precise all-prime formulation in R.~Sharifi,
\emph{Iwasawa Theory}, Theorem 6.2.1,
\url{https://math.ucla.edu/~sharifi/notes/iwasawa-ch06.html}.

\item[TW-IW: the ordinary Iwasawa module theorem.]
For a $\mathbb Z_\ell$-extension $F_\infty/F$, the inverse limit
$X=\varprojlim\operatorname{Cl}(F_n)\{\ell\}$ under norms is a
finitely generated torsion
$\Lambda=\mathbb Z_\ell[[\operatorname{Gal}(F_\infty/F)]]$-module.
Artin reciprocity identifies it with the Galois group of the maximal
unramified abelian pro-$\ell$ extension of $F_\infty$.
See Sharifi, Corollary 3.2.7 and Proposition 3.2.11,
\url{https://math.ucla.edu/~sharifi/notes/iwasawa-ch03.html}.
The structure theorem used with $\mu=0$ is Theorem 2.2.13,
together with Remark 2.2.11 and Lemma 2.2.12,
\url{https://math.ucla.edu/~sharifi/notes/iwasawa-ch02.html}.
It says that a torsion $\Lambda$-module has, up to finite kernel and
cokernel, distinguished-polynomial and $\ell$-power factors;
$\mu=0$ removes the latter.

\item[TW-CF: finite class-field reciprocity.]
The maximal everywhere unramified abelian extension of a totally
complex number field has Galois group its ideal class group.
Artin maps intertwine restriction of automorphisms with ideal norms.
See J.~S.~Milne, \emph{Class Field Theory}, version 4.03,
V, Theorem 3.6, Corollary 3.7 and Example 3.9, pp.~158--159,
\url{https://www.jmilne.org/math/CourseNotes/CFT.pdf}.
The norm compatibility can also be checked directly on prime ideals
using the residue degree in the Frobenius formula.

\item[TW-CD: abelian conductor--discriminant formula.]
For an abelian number field, its absolute discriminant is the product
of the conductors of its characters.  In particular,
$F\subseteq\mathbb Q(\zeta_N)$ implies
$|\Delta_F|\le N^{[F:\mathbb Q]}$.
See Sharifi, equation (4.2.1),
\url{https://math.ucla.edu/~sharifi/notes/iwasawa-ch04.html}.
This is the same classical discriminant input already recorded in
OB-MD.
\end{description}

Finite Poitou--Tate duality is the already reconstructed LD3, rather
than another unnamed classical assumption.  LK2 and LK3 realize its
finite local-cohomology and local-duality ingredients on the actual
finite extensions inside the specified $\ell$-adic algebraic closure.
Elementary finite abelian groups, compact duality, the two-term
cohomology complex of a procyclic pro-$\ell$ group, and the elementary
ramification structure of local fields are used explicitly below.

\paragraph{TW1: the exact Ferrero--Washington rank interface, including two.}
\label{claim:TW1}
Let $F/\mathbb Q$ be abelian and contain $\mu_\ell$ if $\ell$ is odd,
or $\mu_4$ if $\ell=2$.  There exists an integer $r\ge0$ such that
for every $a\ge0$ and every field model $F'$ of
$F(\mu_{\ell^a})$ over $F$,
\[
 \left|\operatorname{Cl}(F')/\ell\operatorname{Cl}(F')\right|
       \le\ell^r.
 \tag{TW.3}
\]
Moreover, the continuous everywhere-unramified characters
$G_{F(\mu_{\ell^\infty})}\to\mathbb F_\ell$ have dimension at most
one fixed finite integer.  Both statements are invariant under
changing the field model or the algebraic closure.

\begin{proof}
The cyclotomic image of $G_F$ is open in
$\mathbb Z_\ell^\times$.  The root hypothesis puts it inside
$1+\ell\mathbb Z_\ell$ for odd $\ell$, and inside
$1+4\mathbb Z_2$ at two.  The $\ell$-adic logarithm, or the binomial
valuation formula, identifies these groups with $\mathbb Z_\ell$.
Their open subgroups are $1+\ell^c\mathbb Z_\ell$, with $c\ge1$
and $c\ge2$ at two.  Consequently
$F_\infty=F(\mu_{\ell^\infty})$ is precisely the cyclotomic
$\mathbb Z_\ell$-extension of $F$; its root-adjoining fields are
the usual finite layers, possibly repeated at the beginning.
In particular, no extra order-two quotient is present at $\ell=2$.

By TW-IW and TW-FW, the inverse-limit class module $X$ is finitely
generated over $\mathbb Z_\ell$ itself.  Indeed its structure theorem
has no $\Lambda/(\ell^b)$ terms, and every remaining
$\Lambda/(f^a)$, with $f$ distinguished, is finite free over
$\mathbb Z_\ell$.  Finite kernel and cokernel preserve finite
generation over the noetherian ring $\mathbb Z_\ell$.

We spell out why this gives the actual finite-layer assertion
(TW.3), without assuming a stronger control theorem.
At a prime above $\ell$ the local cyclotomic extension is eventually
totally ramified.  To check this after base change, let $K/\mathbb
Q_\ell$ be the local base and let $L_b/\mathbb Q_\ell$ be a finite
cyclotomic root layer.  It is totally ramified, while
$[KL_b:K]\le[L_b:\mathbb Q_\ell]$ and
$e(KL_b/K)\ge e(L_b/\mathbb Q_\ell)/e(K/\mathbb Q_\ell)$.
Consequently the residue degrees $f(KL_b/K)$ are bounded by
$e(K/\mathbb Q_\ell)$; in the increasing tower they stabilize.
The local cyclotomic image is open in the global
$\mathbb Z_\ell$-image, so the number of primes also stabilizes.
Thus, after one layer $n_0$,
every subsequent extension of layers is totally ramified at a
chosen prime above $\ell$.  If $m\ge n\ge n_0$, no nontrivial
subextension of $F_m/F_n$ is unramified at that prime.
Let $H_n/F_n$ be the Hilbert $\ell$-class field.  Then
$H_n\cap F_m=F_n$, and $H_nF_m/F_m$ is everywhere unramified.
It is an abelian $\ell$-extension, so it lies in $H_m$.
Restriction therefore maps
$\operatorname{Gal}(H_m/F_m)$ onto
$\operatorname{Gal}(H_n/F_n)$.  By TW-CF this is the norm
$\operatorname{Cl}(F_m)\{\ell\}\to
\operatorname{Cl}(F_n)\{\ell\}$.

All these later norm maps are surjective.  Compactness of the
inverse limit of finite groups gives a surjection
$X\to\operatorname{Cl}(F_n)\{\ell\}$ for every $n\ge n_0$.
Hence their mod-$\ell$ dimensions are at most
$\dim_{\mathbb F_\ell}X/\ell X$.  Take the maximum with the
finitely many earlier dimensions.  The prime-to-$\ell$ part of a
finite class group has zero quotient by $\ell$, proving (TW.3).
This is the literal statement
\texttt{Published.\allowbreak Ferrero\allowbreak Washington}.

The class-field identification in TW-IW gives
\[
 \operatorname{Hom}_{\mathrm{cont,unr}}
 (G_{F_\infty},\mathbb F_\ell)
   =\operatorname{Hom}_{\mathrm{cont}}(X,\mathbb F_\ell),
 \tag{TW.4}
\]
so the same finite generator bound suffices at infinity.
Finally, a field isomorphism transports ideals, class groups, roots
of unity and their adjoining fields.  Extending it to algebraic
closures conjugates the corresponding absolute Galois groups and
their inertia subgroups.  All maps just used therefore preserve
both dimensions, for every field and closure model.
\end{proof}

\paragraph{TW2: one two-primary rank bound before the auxiliary primes are chosen.}
\label{claim:TW2}
Fix the abelian field datum $M$ with modulus $m>0$ and degree
$o=M.\deg$.  There is an integer $R_0(M)\ge1$, independent of
$q,d$, with the following property.  For every \texttt{Prime\allowbreak Pair}
$(q,d)$ there is an abelian field $F_0$ such that
\[
 \mu_4\subset F_0,\quad [F_0:\mathbb Q]\le8o,\quad
 F_0\hookrightarrow\mathbb Q(\zeta_{4mqd}),
 \tag{TW.5}
\]
and the compact Galois module of the maximal unramified abelian
pro-$2$ extension of $F_0(\mu_{2^\infty})$ in which primes above
$2$ split is generated by at most $R_0(M)$ elements over
$\mathbb Z_2$.  One may take
\[
                 R_0(M)=4(g+6o),                         \tag{TW.6}
\]
where $g$ is a fixed bound from TW1 for $M(i)$.
For every $s\ge2$ and every \emph{original} field $E$ satisfying
$\texttt{Is\allowbreak Comp\allowbreak Field}\ M\ s\ q\ d\ E$, some layer $F_a/F_0$ with
$a\le s+2$ satisfies, for every $k\ge1$,
\[
                    s_2(E,k)\le1+s_2(F_a,k).             \tag{TW.7}
\]
For each odd prime $\ell$ and abelian datum $A$ containing
$\mu_\ell$, an analogous fixed finite generator bound
$R_\ell(A)\ge1$ exists for every field model of $A$.

\begin{proof}
Choose a model $M_0$ of $M$ inside $\overline{\mathbb Q}$.
For an odd prime $p$, let $Q_p$ be the quadratic subfield of
$\mathbb Q(\zeta_p)$ fixed by the squares in
$(\mathbb Z/p)^\times$.  It is unramified outside $p$.
Set
\[
 H_\infty=M_0(\mu_{2^\infty}),\qquad
 F_0=M_0(i)Q_qQ_d,\qquad
 F_\infty=H_\infty Q_qQ_d.                              \tag{TW.8}
\]
These fields give (TW.5) immediately.  The root and degree data
are uniform even if some composita have smaller than their maximal
possible degree.

Let $\mathcal U$ be the $\mathbb F_2$-space of characters of
$G_{H_\infty}$ unramified outside primes above $q,d$.
Its everywhere-unramified subspace has dimension at most $g$
by TW1.  Restriction to the inertia groups at those primes has
that subspace as its kernel.  At an odd residue characteristic,
wild inertia has no nontrivial map to $\mathbb F_2$, and the
maximal pro-$2$ quotient of tame inertia is procyclic.
Thus each prime contributes at most one dimension.

There are at most $2o$ primes of $H_\infty$ over $q$ and at most
$4o$ over $d$.  Here is the precise congruence calculation.
For $q\equiv3\pmod8$, one has $v_2(q^2-1)=3$, so for $n\ge3$
the order of $q$ in $(\mathbb Z/2^n)^\times$ is $2^{n-2}$.
For $d\equiv7\pmod{16}$, one has $v_2(d^2-1)=4$, and for
$n\ge4$ its order is $2^{n-3}$.  These statements follow by
successively squaring $1+2^b u$ with $u$ odd; each squaring
increases its distance valuation from $1$ by exactly one.
The Frobenius indices in the respective cyclotomic fields are
therefore $2$ and $4$.  Passing from $\mathbb Q(\mu_{2^\infty})$
to its compositum with $M_0$ multiplies the number of primes by
at most $o$.  The bounded increasing finite-layer prime counts
stabilize, so these are also the counts at infinity.  We obtain
\[
                     \dim\mathcal U\le g+6o.            \tag{TW.9}
\]

For completeness, an index-two extension multiplies this bound
by at most two.  If $V\triangleleft U$ has quotient $C_2$, set
$W=\operatorname{Hom}_{\rm cont}(V,\mathbb F_2)$.
For the induced involution $\tau$, $(\tau-1)^2=0$, so
$\dim W\le2\dim W^{C_2}$: its image lies in its kernel.
The restriction--inflation sequence
\[
 0\longrightarrow H^1(C_2,\mathbb F_2)
 \longrightarrow H^1(U,\mathbb F_2)
 \longrightarrow W^{C_2}
 \longrightarrow H^2(C_2,\mathbb F_2)
 \tag{TW.10}
\]
and the fact that both outer groups have dimension one imply
$\dim W^{C_2}\le\dim H^1(U,\mathbb F_2)$.
This also proves finiteness before using the dimension formula.
Apply this twice to the extensions adjoining $Q_q$ and $Q_d$.
They are unramified outside the allowed primes, so their maximal
extensions unramified outside those primes can be taken inside
one common Galois closure; the indicated subgroups really have
index at most two.  A trivial extension only improves the bound.
The everywhere-unramified, $2$-split characters of $G_{F_\infty}$
form a subspace, so their dimension is at most (TW.6).

Let $X$ be the compact $\mathbb Z_2$-dual of the discrete group
of all such $\mathbb Q_2/\mathbb Z_2$-valued characters.
Then $(X/2X)^\vee$ is the space just bounded.  Compact Nakayama
gives a surjection $\mathbb Z_2^{R_0}\to X$.
One direct proof is to lift a basis of $X/2X$, approximate an
arbitrary element successively modulo $2,2^2,\ldots$, and take
the limit of its coefficients in $\mathbb Z_2$.
The image is compact and closed, and $\bigcap_n2^nX=0$:
an element of this intersection annihilates every finite-order
character in the discrete dual.  The subgroups $2^nX$ tend to
zero for the same reason.
The same argument applies to every prime $\ell$.

The fixing subgroup in \texttt{Is\allowbreak Comp\allowbreak Field} identifies the actual
image of $E$ with
$M_0\,\mathbb Q(\zeta_{2^{s+2}})^+ Q_qQ_d$ in its given
cyclotomic overfield.  Adjoining $i$ gives
$F_a=F_0(\mu_{2^{s+2}})$, whose degree over $F_0$ is $2^a$
for some $a\le s+2$.  The extension $F_a/E$ has degree at most
two.  Restriction of cohomology preserves every unramified and
strict local condition in (TW.1).  Its kernel embeds in
$H^1(C_2,A_2(k)^{G_{F_a}})$ when the degree is two.
The invariant coefficient group is either a finite cyclic
$2$-group or $\mathbb Q_2/\mathbb Z_2$.
For a cyclic coefficient subgroup, the formula
$H^1(C_2,B)=\ker(1+\tau)/(\tau-1)B$ shows that it is killed by
$2$ and has at most one generator; its order is at most $2$.
For the divisible group the only involutions are $1,-1$:
the same formula gives respectively $B[2]$ and $B/2B=0$.
Thus the kernel length is at most one, proving (TW.7).
Conjugating the embedding back transports this inequality to
the original $E$, not to a newly substituted host.

Finally, for a fixed odd $\ell$ and $A$, the condition
$A.\texttt{Has\allowbreak Mu}\ \ell$ means its fixing subgroup acts trivially
on $\mu_\ell$, so every model contains those roots.
TW1 bounds the mod-$\ell$
dimension of its unramified tower module.  Taking the
$\ell$-split quotient cannot increase the number of generators.
Enlarge the bound to at least one, and use the field-model
transport at the end of TW1.  This supplies $R_\ell(A)$.
\end{proof}

\paragraph{TW3: exact twisted tower descent and the reference-layer comparison.}
\label{claim:TW3}
Let $F_0$ contain $\mu_\ell$ if $\ell$ is odd and $\mu_4$ at two.
Put $F_\infty=F_0(\mu_{\ell^\infty})$, choose a topological
generator $\gamma$ of $\Gamma=\operatorname{Gal}(F_\infty/F_0)$,
and put $u=\chi_{\rm cyc}(\gamma)$.
Let $X$ be the compact module of unramified abelian pro-$\ell$
extensions of $F_\infty$ split at primes above $\ell$,
with its conjugation action $\rho(\gamma)$.
For every $k\ge1$ and every model of the $a$th layer,
\[
 s_\ell(F_a,k)=
 \operatorname{length}_{\mathbb Z_\ell}
 \frac{X}{((u^k\rho(\gamma))^{\ell^a}-1)X}.              \tag{TW.11}
\]
For every original $\Omega,D$ on any such layer $F$, LD3 and
LK2--LK3 give
\[
                         s_\ell(F,k)\le V_\ell(F,k;D).
 \tag{TW.12}
\]
Both statements keep the negative twist in (TW.1).

\begin{proof}
Over $F_\infty$ the twist is trivial.  Restrict a global cocycle
to $G_{F_\infty}$; its restriction is a continuous additive
character.  The strict local conditions say precisely that it
vanishes on every inertia group and on each decomposition group
above $\ell$.  Let $D_\infty=X^\vee$ be this character group.
Each finite-image character has an open stabilizer under conjugation:
continuity on the compact Galois group gives one open neighborhood
working uniformly on its finitely many fibers.  Thus this is a
discrete continuous $\Gamma$-module, and its compact dual action
is continuous.  Conjugation on it is
$f(h)\mapsto f(\gamma^{-1}h\gamma)$; the twisted action is
\[
         f(h)\longmapsto u^{-k}f(\gamma^{-1}h\gamma).
 \tag{TW.13}
\]
Restriction therefore lands in the invariants for (TW.13)
under $\gamma^{\ell^a}$.

This restriction is an isomorphism.  Indeed the continuous
cohomology complex of $\mathbb Z_\ell$ on a discrete
$\ell$-primary module is the two-term complex
$A\xrightarrow{\gamma-1}A$.  In the present module
$u^{-k\ell^a}-1\ne0$, and multiplication by this scalar on
$\mathbb Q_\ell/\mathbb Z_\ell$ is surjective.
Thus $H^1(\Gamma^{\ell^a},A_\ell(k))=0$ and
$H^2(\Gamma^{\ell^a},A_\ell(k))=0$.
The restriction--inflation sequence proves injectivity and
surjectivity onto invariant characters.
The local conditions survive the inverse construction.
Away from $\ell$ the cyclotomic tower is unramified, so the
whole inertia subgroup lies in $G_{F_\infty}$ and the lifted
cocycle vanishes there.  At a prime above $\ell$, the local
cyclotomic image is open in $\Gamma^{\ell^a}$.
A class whose restriction to the corresponding infinite-level
decomposition kernel is zero is inflated from this open
procyclic group; its $H^1$ vanishes by the same nonzero-scalar
argument.  Hence the lift is locally zero as required.

The dual of (TW.13) has positive twist
$u^k\rho(\gamma)$.  Duality changes invariants into
coinvariants.  If one first obtains the inverse operator,
its image minus the identity is the same image, since
$g^{-1}-1=-g^{-1}(g-1)$ for an automorphism $g$.
Compact/discrete duality preserves length, including the value
$\infty$, giving (TW.11).  There is a unique subgroup of index
$\ell^a$ in $\Gamma$; thus every field satisfying
\texttt{Is\allowbreak Layer} is carried to this same layer by an embedding
extending the specified $F_0$-algebra structure.

We give the limit argument for (TW.12), rather than treating
it as an extra duality hypothesis.
For $r\ge1$ let
$T_r=\mathbb Z/\ell^r(k+1)$ and
$A_r=A_\ell(k)[\ell^r]$.
The Cartier dual of $T_r$ is exactly $A_r$:
the roots-of-unity twist subtracts $k+1$ from $1$, giving
$-k$.  The finite Poitou--Tate pairing in LD3 identifies
the duals of the finite groups
\[
 \operatorname{Sha}^2(F,T_r)=
 \ker\!\left(H^2(G_{F,S_\ell},T_r)\to
          \prod_{v\mid\ell}H^2(F_v,T_r)\right),
 \qquad \operatorname{Sha}^1(F,A_r).
 \tag{TW.14}
\]
These are the natural cup-product pairings.  Reduction
$T_{r+1}\to T_r$ is paired with inclusion
$A_r\to A_{r+1}$: this follows directly by evaluating the
coefficient pairing and then the same local invariant maps.
The local transport proved in KB6 identifies the original
$\operatorname{Dec}(D.A_v)$ with the absolute Galois group of
the corresponding henselization, then of its completion $F_v$,
and finally of a finite field
$K_v\subset\texttt{Padic\allowbreak Alg\allowbreak Cl}(\ell)$.
These are continuous isomorphisms preserving every root of
unity and hence the original cyclotomic character.
Pullback on cochains preserves cup products and the given
localization map.  LK2--LK3 therefore supply those maps and
finite local cohomology on the prescribed local fields,
including at $\ell=2$ and at every coefficient level $r$.
There are no real-place terms, since the root hypothesis
implies that $F$ is totally complex.

All finite-level $H^1(G_{F,S_\ell},T_r)$ are finite, as also
established on these exact cochains in KB1:
their localization kernels are finite by finite duality,
applied with the Cartier-dual coefficient as necessary,
and their images lie in the finite product of the finite
local $H^1$ groups.  Their inverse systems therefore satisfy
the Mittag--Leffler condition.  Continuous cochains with
coefficients in $\mathbb Z_\ell(k+1)$ are the inverse limits
of the finite coefficient cochains; transition maps on cochains
are surjective by lifting a function on its finitely many
clopen fibers.  The usual kernel--cokernel calculation for
inverse-limit complexes, with the just established
$\varprojlim{}^1H^1=0$, gives
$H^2(F,T_\ell(k))=\varprojlim_r H^2(F,T_r)$, globally
and locally.  Inverse limits preserve kernels, so
$\operatorname{Sha}^2(F,T_\ell(k))
=\varprojlim_r\operatorname{Sha}^2(F,T_r)$.

On the torsion side continuous cohomology commutes with this
filtered union of coefficients: every continuous cochain on a
compact profinite group has finite image.  There are only
finitely many local terms, and filtered unions are exact.
Consequently
$\operatorname{Sha}^1(F,A_\ell(k))
=\varinjlim_r\operatorname{Sha}^1(F,A_r)$.
Taking the limit of the compatible finite pairings in (TW.14)
identifies its compact dual with $\operatorname{Sha}^2(F,T_\ell(k))$.
This latter group is a submodule of the actual global $H^2$
in (TW.2), which proves (TW.12).  All these constructions
commute with field isomorphisms and with conjugating a chosen
place in $\Omega$; therefore (TW.12) holds for every original
local-data record, not just a convenient representative.
\end{proof}

\paragraph{TW4: complete coinvariant bounds with at most \(R\) centers.}
\label{claim:TW4}
Suppose the module $X$ in TW3 is generated by $R\ge1$ elements
over $\mathbb Z_\ell$.  Set
\[
 j_0=\lfloor\log_\ell R\rfloor+3,\qquad
 c=v_\ell(u-1).
 \tag{TW.15}
\]
There are $s_1,\ldots,s_R\in\mathbb Z_\ell$ such that
for every $k\ge1$ and every $a\ge0$,
\[
 s_\ell(F_a,k)\le
   \sum_{i=1}^R\bigl(\max(a,j_0)+c+v_\ell(s_i-k)\bigr).
 \tag{TW.16}
\]
Here $v_\ell(0)=\infty$, and
$c\ge1$, with $c\ge2$ at two; moreover $F_0$ contains a
primitive $\ell^c$th root of unity.
For every model $F_{j_0}$ of the reference layer and every
local-data record $D$ on it,
\[
 s_\ell(F_a,k)\le
 V_\ell(F_{j_0},k;D)+R\max(a-j_0,0).                    \tag{TW.17}
\]
These are the two parts of the literal \texttt{ABound l R F0}.

\begin{proof}
The cyclotomic image is $1+\ell^c\mathbb Z_\ell$ by the
root calculation in TW1, and $c$ is the largest exponent for
which $\mu_{\ell^c}\subset F_0$.
The action of $\Gamma$ on the finite vector space $X/\ell X$
is continuous, hence factors through a finite $\ell$-group.
Thus $\rho(\gamma)-1$ is nilpotent modulo $\ell$, with
nilpotence degree at most $\dim X/\ell X\le R$.
For $g=u^k\rho(\gamma)$ the reduction is the same,
so this assertion is uniform in $k$.

We first prove the layer growth bound
\[
 \operatorname{length}\operatorname{coker}(g^{\ell^a}-1)
 \le
 \operatorname{length}\operatorname{coker}(g^{\ell^{j_0}}-1)
                +R\max(a-j_0,0).                       \tag{TW.18}
\]
Put $h=\lfloor\log_\ell R\rfloor+1$, so $\ell^h\ge R$.
In characteristic $\ell$,
$(g-1)^{\ell^h}=g^{\ell^h}-1$, hence
$(g^{\ell^h}-1)X\subset\ell X$.
One further $\ell$th power gives image in $\ell^2X$,
by the binomial theorem; this is valid at two as well.
It remains true after further $\ell$th powers.
If $(w-1)X\subset\ell^2X$, set
$\Phi_w=1+w+\cdots+w^{\ell-1}$.
For every $x$, $\Phi_wx=\ell x+\ell^2y$ for some $y$.
Thus $\ell X\subset\Phi_wX+\ell(\ell X)$.
Nakayama applied to
$(\ell X+\Phi_wX)/\Phi_wX$ gives
$\ell X\subset\Phi_wX$, so
$\operatorname{length}\operatorname{coker}\Phi_w\le R$.
Factor $w^\ell-1=\Phi_w(w-1)$ and use the exact sequence
of cokernels of a composite; the length increases by at most
$R$ at each later layer.  At earlier layers monotonicity follows
from the same factorization of $g^{\ell^{j_0}}-1$ by
$g^{\ell^a}-1$.  This proves (TW.18), also when a length is
infinite.

Choose a surjection $\pi:\mathbb Z_\ell^R\to X$, and lift
$\rho(\gamma)^{\ell^{j_0}}$ to a matrix $S$ on its generators.
Then $\pi S=\rho(\gamma)^{\ell^{j_0}}\pi$.
There is a surjection of the corresponding cokernels, so
Smith normal form gives
\[
 \operatorname{length}\operatorname{coker}
       (g^{\ell^{j_0}}-1)
 \le v_\ell\det\bigl((u^{\ell^{j_0}})^kS-I\bigr).
 \tag{TW.19}
\]
A zero determinant gives $\infty$, consistently with the
inequality.  In an algebraic extension of $\mathbb Q_\ell$
factor the determinant as
$\prod_{i=1}^R(a_k z_i-1)$, where
$a_k=(u^{\ell^{j_0}})^k$ and the $z_i$ are the eigenvalues
of $S$, with multiplicities.

For each $z=z_i$, we construct a center in $\mathbb Z_\ell$
without requiring $z$ to be a power of $u$.
For all nonnegative integers $b,k$ the binomial valuation
formula gives
\[
 |a_b-a_k|_\ell=\ell^{-j_0-c}|b-k|_\ell.
 \tag{TW.20}
\]
For any units $A,B$ and any $z$ in the normed extension,
\[
 |A-B|_\ell\le
       \max\{|Az-1|_\ell,\ |Bz-1|_\ell\}.
 \tag{TW.21}
\]
If $|z|_\ell\ge1$, this follows by subtracting the two
expressions and dividing by $|z|_\ell$.
If $|z|_\ell<1$, both norms on the right are one.
Therefore the closed balls in $\mathbb Z_\ell$ centered at
$k$ with radii $\ell^{j_0+c}|a_kz-1|_\ell$ intersect
pairwise.  Intersecting ultrametric balls are nested; their
finite intersections are nonempty.  Compactness of
$\mathbb Z_\ell$ gives a common point $s_i$.
It satisfies
$v_\ell(a_kz_i-1)\le j_0+c+v_\ell(s_i-k)$.
Summing these inequalities in (TW.19), then using (TW.18)
and TW3, proves (TW.16).

Finally, (TW.18) and (TW.11) compare the arbitrary layer with
the reference layer.  Apply (TW.12) to that reference layer
with its given $D$, giving (TW.17).
The centers may be regarded as elements of $\mathbb Q_\ell$,
which is an allowed normed extension in \texttt{ABound}.
Taking $R'=R$ in its first part and translating
$v_\ell(s_i-k)\le B$ to
$\ell^{-B}\le|s_i-k|_\ell$ matches its norm convention exactly.
\end{proof}

\paragraph{TW5: exceptional twists and a common square window.}
\label{claim:TW5}
Under TW4, for every integer $Z\ge1$ there is a set
$\mathcal E\subseteq\{1,\ldots,Z\}$ of size at most $R$
such that, simultaneously for every layer and every
$1\le k\le Z$ outside $\mathcal E$,
\[
 s_\ell(F_a,k)\le
 R\bigl(\max(a,j_0)+c+\lfloor\log_\ell Z\rfloor\bigr).
 \tag{TW.22}
\]
Independently, let $\mathcal E$ be any finite set of
nonnegative integers of size at most $R$, let
$0<\eta\le C_t$, and let $X>1$.  With
\[
                 J=2^{R+2}(C_t/\eta)^R,                \tag{TW.23}
\]
there is $Y$ such that
\[
 X^{1/J}\le Y,\qquad Y^2\le X,\qquad
 \left.
 \begin{array}{c}
 x\in[Y,Y^2],\ k\in\mathcal E,\ k\le x^{C_t}
 \end{array}\right\}
       \Longrightarrow k\le x^\eta .                  \tag{TW.24}
\]

\begin{proof}
For each center $s_i$ select the integers $1\le k\le Z$
with $|s_i-k|_\ell<\ell^{-\lfloor\log_\ell Z\rfloor}$.
There is at most one: two distinct choices would have
$\ell$-adic difference valuation at least
$\lfloor\log_\ell Z\rfloor+1$, so their ordinary integer
difference would be at least
$\ell^{\lfloor\log_\ell Z\rfloor+1}>Z$.
Take the union of these sets and apply (TW.16).
This proves the full \texttt{First\allowbreak Bound}, not merely a count
of exceptional twists.

For the second assertion put $B=C_t/\eta\ge1$ and
$U=\log X>0$.  We use the following elementary interval fact:
if a finite real set $T$ has at most $R$ elements, some $y$
satisfies
\[
 \frac{U}{2^{R+2}B^R}\le y,\quad 2y\le U,\quad
 a\in T,\ y\le v\le2y,\ a\le v\ \Longrightarrow\ Ba\le v.
 \tag{TW.25}
\]
Here is an induction proof, including its constants.
For $R=0$, take $y=U/2$.
If there is an $a_0\in T$ with $U/(2B)<a_0\le U$,
remove it and apply the induction hypothesis at
$U'=U/(2B)$ and size bound $R-1$.
Its lower bound is exactly the one in (TW.25), while
$v\le2y\le U'<a_0$ makes the deleted point irrelevant.
If no such point exists, take $y=U/2$.
Every $a\in T$ with $a\le v\le U$ then satisfies
$a\le U/(2B)$, so $Ba\le U/2\le v$.
The required lower bound is weaker than $U/2$.

Apply (TW.25) to
$T=\{\log k/C_t:k\in\mathcal E,\ k\ge1\}$ and take
$Y=\exp y$.  If $k\le x^{C_t}$, then
$a=\log k/C_t\le\log x$.  The conclusion
$Ba\le\log x$ is exactly $k\le x^\eta$.
The value $k=0$ causes no difficulty.  Exponentiating the
other inequalities proves (TW.24).  A union of exception
sets can be used here once, yielding the same $Y$ for all
of them.
\end{proof}

\paragraph{TW6: uniform reference-field sizes and exact model transfer.}
\label{claim:TW6}
Fix $\ell,j,n_1,N_0$ with $\ell$ prime and $N_0\ge1$.
There exist $n_0\ge1$ and $c_0\ge0$, depending only on these
four parameters, such that the following holds.
If $[F_0:\mathbb Q]\le n_1$,
$F_0\hookrightarrow\mathbb Q(\zeta_{N_0P})$, $P\ge1$,
$F_j/F_0$ is a $j$th layer, and
$1\le w$, $\log P\le w$, then
\[
 [F_j:\mathbb Q]\le n_0,\qquad
             \log(2|\Delta_{F_j}|)\le c_0w.             \tag{TW.26}
\]
For a fixed abelian datum $A$ and a fixed layer degree
$\ell^j$, all models of $A$ and all such layers are isomorphic
as towers.  The lengths (TW.1) are preserved by these
isomorphisms for every closure and local-data record.

\begin{proof}
Take $n_0=(n_1+1)\ell^j$.
Write $F_j=F_0(\mu_{\ell^b})$ for some $b$.
Since $F_j$ contains a primitive $\ell^b$th root,
$\varphi(\ell^b)\le[F_j:\mathbb Q]$.
For a prime power this implies
$\ell^b\le2[F_j:\mathbb Q]\le2n_0$, also for $b=0$.
Embedding the compositum in
$\mathbb Q(\zeta_{\ell^bN_0P})$ and using TW-CD gives
\[
 \log(2|\Delta_{F_j}|)
 \le \log2+n_0\log(2n_0N_0)+n_0\log P.
 \tag{TW.27}
\]
Thus one may take
$c_0=\log2+n_0\log(2n_0N_0)+n_0$.
The degree bound is the tower formula
$[F_j:\mathbb Q]=[F_0:\mathbb Q]\ell^j$.

A model of $A$ is the fixed field of its specified subgroup
of a cyclotomic Galois group, so two models have a rational
isomorphism.  Extend it to algebraic closures.  It transports
roots of unity and hence the root-adjoining towers.  These
fields are nested, so two with the same relative degree are
equal; under the hypothesis of TW1 this is also the unique
index-$\ell^j$ subgroup description.
The resulting isomorphism respects the
given base-algebra maps.  Conjugation transports absolute
Galois groups, the cyclotomic character, every inertia group
and every decomposition group.  Pullback on continuous
cocycles and its inverse commute with coboundaries and
localization.  Thus they give a linear isomorphism of the
strict kernels (TW.1).  Changing which place above a finite
prime represents a local term only conjugates that
restriction map, inducing the same kernel on cohomology.
This proves the stated model and local-data invariance.
\end{proof}

\paragraph{TW7: all twists on a good interval, from the literal \(\texttt{Lemma\allowbreak K}\).}
\label{claim:TW7}
Suppose $C_K\ge1$ satisfies the literal \texttt{Lemma\allowbreak K}:
for every prime $\ell$, every totally complex number field $F$
containing the roots required in TW1, every original $\Omega,D$,
and every $k\ge1$, the actual compact-coefficient group is
finite and
\[
 \log\left|H^2(G_{F,S_\ell},\mathbb Z_\ell(k+1))\right|
 \le
 \exp\!\left((2k[F:\mathbb Q])^{C_K}
                         \log(2|\Delta_F|)\right).       \tag{TW.28}
\]
Fix the prime $\ell$ and $R,c_1,n_0,c_0,J,C_t,\varepsilon,a$ with
$R\ge1$, $c_0,a,C_t\ge0$, $J,\varepsilon>0$.
There is $C\ge0$ depending only on these fixed parameters
and $C_K$ with the following property.
Suppose a tower as in TW4 has $c\le c_1$, its exception set
from TW5 is constructed up to $Z=\lfloor X^{C_t}\rfloor$,
and its reference layer satisfies
\[
 [F_{j_0}:\mathbb Q]\le n_0,\qquad
 \log(2|\Delta_{F_{j_0}}|)\le c_0x^{\varepsilon/4}.        \tag{TW.29}
\]
If $1\le x\le X$, $X^{1/J}\le x$, and
\[
 k\in\mathcal E,\ k\le x^{C_t}
              \Longrightarrow k\le x^{\varepsilon/(2C_K)},
 \tag{TW.30}
\]
then for every $1\le k\le x^{C_t}$ and every layer index
$b\le a(1+\log x)$ the length $s_\ell(F_b,k)$ is finite and
\[
 s_\ell(F_b,k)+1\le
       C(1+\log x)+2\exp(Cx^{3\varepsilon/4}).
 \tag{TW.31}
\]
For fixed $C,C_t,\varepsilon>0$, the right side times $x^{C_t}$
is at most $\exp(x^{4\varepsilon/5})$ for all sufficiently
large $x$.

\begin{proof}
A finite $\mathbb Z_\ell$-module of length $v$ has order
$\ell^v$.  Since $\log\ell\ge\log2>1/2$, (TW.28) gives
\[
 V_\ell(F,k;D)\le
 2\exp\!\left((2k[F:\mathbb Q])^{C_K}
                                  \log(2|\Delta_F|)\right).
 \tag{TW.32}
\]
If $k\in\mathcal E$, (TW.30) and (TW.29) imply
\[
 (2k[F_{j_0}:\mathbb Q])^{C_K}
       \log(2|\Delta_{F_{j_0}}|)
 \le (2n_0)^{C_K}c_0x^{3\varepsilon/4}.                 \tag{TW.33}
\]
Apply (TW.17) and (TW.32).  If $k\notin\mathcal E$, use
(TW.22); the inequality $k\le Z$ follows from
$k\le x^{C_t}\le X^{C_t}$.
Both cases prove finiteness and the common bound
\[
 s_\ell(F_b,k)\le
 R\bigl(b+j_0+c+\lfloor\log_\ell Z\rfloor\bigr)
       +2\exp((2n_0)^{C_K}c_0x^{3\varepsilon/4}).         \tag{TW.34}
\]
The logarithmic inequalities
$\log X\le J\log x$ and
$\log\ell>1/2$ give
$\lfloor\log_\ell Z\rfloor\le2C_tJ\log x$.
Thus (TW.31) holds with
\[
 C=R(a+j_0+c_1+2C_tJ)+1+(2n_0)^{C_K}c_0.
 \tag{TW.35}
\]
In particular the constant is uniform in the varying base
field whenever the displayed degree, discriminant and root
bounds are uniform.

For the final absorption, divide the logarithm of
$x^{C_t}[C(1+\log x)+2\exp(Cx^{3\varepsilon/4})]$
by $x^{4\varepsilon/5}$.
The exponential term contributes at most a constant times
$x^{-\varepsilon/20}$; the remaining terms are bounded by
constants times $\log x$ and $\log(2+\log x)$ divided by
$x^{4\varepsilon/5}$.  They all tend to zero.
The ratio is therefore at most one after a threshold depending
only on $C,C_t,\varepsilon$.  This proves the assertion.
\end{proof}

\paragraph{TW8: the literal simultaneous \(\texttt{Thm\allowbreak Bab}\), with its quantifiers.}
\label{claim:TW8}
Assume $\exists C_K,\ \texttt{Lemma\allowbreak K}\ C_K$.
Together with TW1--TW7, LD3 and LK2--LK3, this supplies
\[
                \exists K:\texttt{BConst},\
                \forall C_t\ge1,\ \texttt{Thm\allowbreak Bab}\ K\ C_t .
 \tag{TW.36}
\]
More explicitly, the single $K=(R_0,R_\ell,C_K)$ is chosen
before $C_t$, the abelian datum $M$, the finite set $B'$ of
odd primes, the family of abelian data $L(\ell)$ containing
$\mu_\ell$, $\varepsilon,A_1$, the admitted functions, the
prime pair and the eventual window.  For these data put
\[
 \begin{split}
 R&=R_0(M)+\sum_{\ell\in B'}R_\ell(L(\ell)),\\
 J&=2^{R+2}\left(\frac{2C_KC_t}{\varepsilon}\right)^R.
 \end{split}                                           \tag{TW.37}
\]
For $0<\varepsilon<1/52$ and any real $A_1$ there exists $X_0$
such that for every pair of functions $e,t:\mathbb R\to\mathbb N$
satisfying, for all $x\ge2$,
\[
 2\le e(x),\quad 2^{e(x)}\le x^{A_1},\quad
                  2\le t(x)\le x^{C_t},                 \tag{TW.38}
\]
every $X\ge X_0$, and every \texttt{Prime\allowbreak Pair} $(q,d)$ with
\[
                         \log(qd)\le X^{\varepsilon/(4J)},
 \tag{TW.39}
\]
one common $Y$ satisfies $X^{1/J}\le Y$, $Y^2\le X$ and the
following two assertions for every $x\in[Y,Y^2]$:
\[
 \begin{array}{ll}
 \operatorname{VV}_2(E,t(x),\exp(x^{4\varepsilon/5}))
 &\text{for every original }
 E\text{ with }\texttt{Is\allowbreak Comp\allowbreak Field}\ M\ e(x)\ q\ d\ E;\\[2mm]
 \operatorname{VV}_\ell(F,t(x),\exp(x^{4\varepsilon/5}))
 &\text{for every }\ell\in B',\ \text{every model }F_0
   \text{ of }L(\ell),\\
 &\text{and every }j_0(\ell,R_\ell(L(\ell)))\text{th layer }
 F/F_0 .
 \end{array}                                           \tag{TW.40}
\]
The second line quantifies separate reference-layer models;
it does not identify them with the field $E$ in the first.

\begin{proof}
First enlarge a witness for \texttt{Lemma\allowbreak K} to $C_K\ge1$.
This preserves (TW.28), since $2k[F:\mathbb Q]\ge2$ and
$\log(2|\Delta_F|)>0$.  Choose the function $R_0(M)$ by
TW2 for every abelian datum.  For each odd $\ell$ and
root-containing datum $A$, choose its fixed bound
$R_\ell(A)$ by TW2; assign arbitrary values, for example zero,
outside that domain.  These choices and $C_K$ define one
\texttt{BConst}.  They are made before any $C_t$.
Now fix the remaining parameters in the statement.

For the two-primary part, write $o=M.\deg$,
$R_2=R_0(M)$, and $j_2=j_0(2,R_2)$.
Every $F_0$ obtained in TW2 has degree at most $8o$ and is
contained in modulus $4mqd$.  It contains a primitive
$2^c$th root.  Since
$\varphi(2^c)\le[F_0:\mathbb Q]$, we have the coarse uniform
bound $c\le16o$.  Apply TW6 with
$\ell=2$, $j=j_2$, $n_1=8o$, $N_0=4m$, and $P=qd$.
It gives $n_0,c_0$ independent of $q,d$.
If $x\ge X^{1/J}$, (TW.39) implies
$\log(qd)\le x^{\varepsilon/4}$, so every reference layer
in this tower satisfies (TW.29).

By (TW.38),
$e(x)\log2\le A_1\log x$, hence
$e(x)+2\le(2|A_1|+2)(1+\log x)$.
For the layer used in (TW.7), the parameter $a$ in TW7 may
therefore be fixed as $2|A_1|+2$.
TW7, with $c_1=16o$ and the uniform $n_0,c_0$, gives one
constant $C_2$ and an absorption threshold $x_2$ for this
whole family, before the functions, $X$ and $(q,d)$ are
chosen.  It follows that, whenever the exception set of this
tower is good as in (TW.30), each relevant length on the
original $E$ is bounded by the right side of (TW.31):
the extra $1$ in that bound accounts exactly for (TW.7).
There are $t(x)-1\le x^{C_t}$ twists.  Their sum is bounded
by $\exp(x^{4\varepsilon/5})$ once $x\ge x_2$.

For each of the finitely many $\ell\in B'$, choose one model
of $L(\ell)$ and its $j_0(\ell,R_\ell(L(\ell)))$th layer.
Its degree, discriminant and maximal root exponent are now
fixed constants.  Apply TW7 with that root exponent as $c_1$,
the exact degree as $n_0$, and
$c_0=\log(2|\Delta|)$; inequality (TW.29) holds for every
$x\ge1$.  Take $a=j_0$, so the fixed layer index obeys
$j_0\le a(1+\log x)$.
This gives a constant $C_\ell$ and a threshold $x_\ell$,
again before $e,t,X,q,d$.  TW6 transports every conclusion
to every other model of this base and every prescribed layer
over it.  It also transports finiteness and the full family
of original local-data records, not just numerical degrees.

Choose $x_*\ge2$ larger than all the finitely many thresholds
$x_2,x_\ell$, and then set $X_0=\max(2,x_*^J)$.
These choices have precisely the permitted dependencies.
Now choose arbitrary admitted $e,t$, $X\ge X_0$ and a prime
pair satisfying (TW.39).  Use TW2 to select its two-primary
base.  For this tower and for each fixed odd reference tower,
use TW5 with the same cutoff $Z=\lfloor X^{C_t}\rfloor$.
Let their exception sets be $\mathcal E_2,\mathcal E_\ell$.
Their union has size at most the $R$ in (TW.37).
Use the square-window part of TW5 once, with
$\eta=\varepsilon/(2C_K)$ and $C_t/\eta=
2C_KC_t/\varepsilon$.
The hypothesis $\eta\le C_t$ follows from
$C_K,C_t\ge1$ and $\varepsilon<1/52$.
It supplies exactly the $J$ in (TW.37) and one common $Y$.

For every $x\in[Y,Y^2]$, one has
$x\ge X^{1/J}\ge x_*$ and $x\le X$, and every constituent
exception set satisfies (TW.30).  The two preceding uniform
estimates therefore apply to all integers $1\le k<t(x)$.
Summing them proves both lines of (TW.40), for all the
specified fields and all their local data, on this same
window and with this fixed pair $(q,d)$.
By the definitions of \texttt{Ba\allowbreak At}, \texttt{Bb\allowbreak At},
\texttt{Is\allowbreak BWindow} and \texttt{Thm\allowbreak Bab}, this is (TW.36).
\end{proof}

\subsection{Scope of the completed dependency chain}

TW1 supplies the source's exact
\texttt{Published.\allowbreak Ferrero\allowbreak Washington} interface.
TW2--TW4 supply \texttt{Two\allowbreak In} and \texttt{Fixed\allowbreak In}, including
the complete \texttt{Is\allowbreak Tower} to \texttt{ABound} step.
TW5--TW8 supply their simultaneous quantitative window, with one
$K$ before all $C_t\ge1$ and with all field models quantified.
The only project-specific incoming quantitative statement is
the literal \texttt{Lemma\allowbreak K}; KB8 supplies it with $C_K=31$.
The existing VW3--VW4 assembly can therefore replace its
\texttt{Thm\allowbreak Bab} input by TW8 once that supplier is joined.
The odd-minus supplier is already OB8 with constant $240$.
The separate \texttt{Point\allowbreak HV} supplier for the odd-plus comparison,
and the other arithmetic, lattice, host and representation inputs
outside this branch, remain exactly where the cumulative dependency
audit records them.  No full \texttt{Link\allowbreak N}, lattice gain, RH or
GRH conclusion is asserted by this chapter alone.


% Version 0.25. Complete exceptional odd-plus count supplier.
\section{Vanishing in the exceptional odd-plus case}
\label{sec:count-z-closure}

\status{CZ1--CZ6 prove the literal predicates \texttt{Prop\allowbreak Z} and
\texttt{Count\allowbreak Z}, with the uniform exponent $C=0$.  The proof works on
the original algebraic closure and coefficient ring.  It proves the
two residual-character vanishing arguments, the needed cyclotomic
control, and the quadratic class-number estimate.  The accepted
arithmetic input is finite global class field theory in its ray-class
and Hilbert-class-field forms.  No local-duality, K-theory, automorphy,
Bernoulli-unit or other count predicate is retained as a premise.}

Fix an odd prime $\ell$, an algebraic closure $\Omega$ of
$\mathbb Q$, and a discrete valuation ring $O$ finite as a
$\mathbb Z_\ell$-algebra.  Put
\[
 G=\operatorname{Aut}_{\mathbb Q}(\Omega),\qquad
 M=O\otimes_{\mathbb Z_\ell}(\mathbb Q_\ell/\mathbb Z_\ell),
 \qquad \alpha=\chi_\ell^{-e}:G\longrightarrow O^\times.
 \tag{CZ.1}
\]
Here $M$ has the discrete topology, and $\chi_\ell$ is the
source's actual cyclotomic character, followed by the structure map
to $O$.  The action on $M$ is scalar multiplication by $\alpha$.
For a valuation $A$ of $\Omega$ above a finite rational prime,
write $D_A$ and $I_A$ for its decomposition and inertia groups.
Define the strict unramified Selmer subgroup
\[
 \begin{split}
 \mathcal S_e(O)={}&
 \ker\!\left(H^1(G,M(\alpha))\longrightarrow
     \prod_{A\mid\ell}H^1(D_A,M(\alpha))\right)
 \\
 &\cap\bigcap_A\ker\!\left(H^1(G,M(\alpha))
                       \longrightarrow H^1(I_A,M(\alpha))\right).
 \end{split}
 \tag{CZ.2}
\]
The source imposes its inertia condition for all its valuation
subrings; all finite-place conditions displayed here are among them.
We prove vanishing already under (CZ.2), which therefore proves
vanishing for the source group as well.

\paragraph{Classical input boundary.}
We use the finite global reciprocity and existence theorem: a finite
abelian extension is a quotient of the ray class group for a modulus
supported on its ramified places, and the ray class field of modulus
$1$ has Galois group the ordinary ideal class group, with real places
remaining real.  The rational ray class field of modulus
$\ell^r\infty$ is $\mathbb Q(\mu_{\ell^r})$.
These are the statements in Milne,
\href{https://www.jmilne.org/math/CourseNotes/CFT.pdf}{\emph{Class
Field Theory}}, v4.03, V, Theorems~3.5--3.6, Corollary~3.7,
and Examples~3.9--3.10.  We explain below the transfer from each
actual cocycle to the relevant finite extension.  Basic integral
bases, ideal factorization, finite Galois ramification, Eisenstein's
criterion and Hensel lifting are used in their usual elementary forms;
see Milne,
\href{https://www.jmilne.org/math/CourseNotes/ANT.pdf}{\emph{Algebraic
Number Theory}}, v3.08, \S\S2--3 and 6--7.  In particular the
class-number inequality below is proved, not retained as a leaf.

\begin{proposition}[CZ1: the actual coefficients and small cocycles]
\label{prop:cz-coefficients}
The module $M$ in (CZ.1) is $\ell$-divisible and every element is
annihilated by a power of $\ell$.  Its $\ell$-torsion $M[\ell]$
is a finite-dimensional $\mathbb F_\ell$-vector space.  Every class
in $H^1(G,M(\alpha))$ is annihilated by a power of $\ell$.
Every class annihilated by $\ell$ has a continuous cocycle
representative taking values in $M[\ell]$.

The source's \texttt{Chr.\allowbreak V}$(\ell,e,1)$ is the $O$-length of
its subgroup of (CZ.2), on exactly these coefficients and this
negative cyclotomic action.  For a representative $c$, its local
conditions say that there is a possibly different $m_H\in M$ for
each required local group $H$, with
\[
                  c(g)=(\alpha(g)-1)m_H\quad(g\in H).
                  \tag{CZ.3}
\]
Consequently it suffices to show that the $\ell$-torsion of
$\mathcal S_e(O)$ is zero.
\end{proposition}

\begin{proof}
For the actual source coefficient $O=c.O$, the structure map
$\mathbb Z_\ell\to O$ is injective because $O$ is a subalgebra
of $\texttt{Padic\allowbreak Alg\allowbreak Cl}(\ell)$.  More generally it is injective
under the stated discrete valuation ring hypotheses: otherwise
its kernel is the nonzero prime ideal $(\ell)$; the finite algebra
$O$ would then be a finite integral domain, hence a field,
contrary to the discrete valuation ring hypothesis.  Since $O$ is
a domain, it is torsion-free and finite over the principal ideal
domain $\mathbb Z_\ell$, hence free of some finite rank $a$.
As a $\mathbb Z_\ell$-module, $M$ is therefore a direct sum of
$a$ copies of $\mathbb Q_\ell/\mathbb Z_\ell$.  This proves its
divisibility, primary torsion and the asserted description of
$M[\ell]$; no freeness over $O/\ell O$ is required.

The action on each $M[\ell^r]$ factors through the reduction of
the continuous cyclotomic character modulo $\ell^r$, so it is
continuous.  A continuous cocycle from the compact group $G$ to
the discrete module $M$ has finite image.  A single power of
$\ell$ kills all its values, and therefore also its cohomology
class.  If $\ell[c]=0$, there is $m\in M$ such that
$\ell c(g)=(\alpha(g)-1)m$.  Choose $b\in M$ with $\ell b=m$.
Then
\[
                c'(g)=c(g)-(\alpha(g)-1)b
                \tag{CZ.4}
\]
is a continuous cohomologous cocycle with $\ell c'=0$.
Subtracting this coboundary preserves every local condition (CZ.3).

In the source, \texttt{QO} is the tensor product in (CZ.1),
\texttt{rho}$(\ell,e,1)$ is scalar multiplication by
$\chi_\ell^{-e}$, and \texttt{Sha\allowbreak U} at the full Galois group is
the intersection of the restriction kernels for decomposition
groups above $\ell$ and all inertia groups.  Its cohomology uses
the cocycle equation
$c(gh)=c(g)+\alpha(g)c(h)$ and coboundaries as in (CZ.3), so
these identifications preserve both the original module and the
local restrictions.  Its length is exactly \texttt{Chr.\allowbreak V}.
Finally, if a nonzero primary-torsion group element has least
annihilating power $\ell^r$, multiplication by $\ell^{r-1}$
gives a nonzero element killed by $\ell$.  Thus the last
assertion follows without any finiteness assumption on the Selmer
group.
\end{proof}

\begin{proposition}[CZ2: cyclotomic inertia and rational control]
\label{prop:cz-cyclotomic-control}
Put $G_\infty=\ker(\chi_\ell:G\to\mathbb Z_\ell^\times)$.
For every $A\mid\ell$, the map
$I_A\to\mathbb Z_\ell^\times$ induced by $\chi_\ell$ is
surjective.  At primes different from $\ell$ the character is
trivial on inertia.  Every continuous additive homomorphism
$f:G\to M[\ell]$ unramified outside $\ell$ is zero on
$G_\infty$.

If a class in $\mathcal S_e(O)$ restricts to zero on $G_\infty$,
it is zero already on $G$.  This is the exact injectivity needed
here; it involves no change of algebraic closure or coefficient
ring.
\end{proposition}

\begin{proof}
For $r\ge1$, the polynomial
$\Phi_{\ell^r}(1+T)$ is Eisenstein at $\ell$: its constant
term is $\ell$, and modulo $\ell$ it is
$T^{\ell^{r-1}(\ell-1)}$.  Thus adjoining a primitive
$\ell^r$th root to $\mathbb Q_\ell$ gives a totally ramified
extension of degree $\varphi(\ell^r)$.  Globally the degree is
at most this number and at least this local degree, so the
cyclotomic Galois group is the full
$(\mathbb Z/\ell^r\mathbb Z)^\times$, and inertia at every
chosen $A\mid\ell$ maps onto it.  For a prescribed
$u\in\mathbb Z_\ell^\times$, its finite-level inverse images
in the compact group $I_A$ are nonempty nested closed sets.
Their intersection gives an element whose actual cyclotomic
character is $u$.  This supplies compatible infinite-level values,
not just unrelated finite-level surjections.

At a prime $q\ne\ell$, choose a finite residue extension containing
the $\ell^r$th roots of unity.  Their derivatives in
$X^{\ell^r}-1$ are units, so Hensel lifting places all these
roots in the unramified extension with that residue field.  Inertia
therefore acts trivially on every $\mu_{\ell^r}$, and hence on
$\chi_\ell$.  These finite ramification statements apply to the
decomposition and inertia groups of the original valuations of
$\Omega$.

The finite image of $f$ cuts out a finite abelian extension of
$\mathbb Q$ inside the original $\Omega$.  It is unramified
outside $\ell$ by hypothesis; complex conjugation has zero image
since its order divides $2$ and $M[\ell]$ has odd exponent.
Global class field theory consequently places this field in the
ray class field of some modulus $\ell^r\infty$, hence in
$\mathbb Q(\mu_{\ell^r})\subset\Omega$.
The subgroup $G_\infty$ fixes that field, proving $f|_{G_\infty}=0$.

For the last assertion, choose a representative $c$ for the given
class.  The action on $G_\infty$ is trivial, so every coboundary
on that subgroup is identically zero.  Trivial restriction of the
class therefore means $c|_{G_\infty}=0$.  The cocycle equation
shows that $c(g)$ depends only on $\chi_\ell(g)$.
Fix $A\mid\ell$ and choose $b\in M$ from its decomposition
condition (CZ.3).  For arbitrary $g\in G$, choose
$t\in I_A\subset D_A$ with $\chi_\ell(t)=\chi_\ell(g)$.
Then $c(g)=c(t)=(\alpha(t)-1)b=(\alpha(g)-1)b$.
This is one global coboundary, as required.
\end{proof}

\begin{proposition}[CZ3: the quadratic field and its class number]
\label{prop:cz-quadratic-class-number}
Suppose $\ell\ge5$ and $\ell\equiv1\pmod4$.  Inside the
original $\Omega$ there is a real quadratic field
$K=\mathbb Q(\sqrt\ell)$ of discriminant $\ell$ such that
\[
 G_K=\ker\bigl(\eta\circ\overline\chi_\ell\bigr),
 \qquad h_K<\ell,
 \tag{CZ.5}
\]
where $\eta$ is the quadratic character of
$\mathbb F_\ell^\times$ and $h_K$ is the ordinary ideal class
number.  Every continuous additive homomorphism
$G_K\to M[\ell]$ which is unramified at all finite places is
zero.
\end{proposition}

\begin{proof}
Choose a primitive $\ell$th root $\zeta\in\Omega$, extend
$\eta$ by $\eta(0)=0$, and set
$\tau=\sum_{a\in\mathbb F_\ell}\eta(a)\zeta^a$.
If $g\zeta=\zeta^u$, reindexing the sum gives
$g\tau=\eta(u)\tau$.  Let
$\tau'=\sum_b\eta(b)\zeta^{-b}=\eta(-1)\tau$.
The change of variables $a=tb$ for $b\ne0$ gives
\[
 \begin{split}
 \tau\tau'
 &=\sum_{t\in\mathbb F_\ell}\eta(t)
                  \sum_{b\ne0}\zeta^{(t-1)b}\\
 &=(\ell-1)-\sum_{t\ne1}\eta(t)=\ell.
 \end{split}
 \tag{CZ.6}
\]
Here the inner sum is $\ell-1$ for $t=1$ and $-1$ otherwise,
and $\sum_t\eta(t)=0$.  Since $\eta(-1)=1$, we have
$\tau^2=\ell$.  A rational square cannot have odd valuation at
the prime $\ell$, so $\tau\notin\mathbb Q$.  Thus
$K=\mathbb Q(\tau)\subset\Omega$ is quadratic and real, and
the transformation formula gives the asserted subgroup $G_K$.

The element $\omega=(1+\tau)/2$ is integral, satisfying
$X^2-X+(1-\ell)/4=0$.  The discriminant of the order
$\mathbb Z[\omega]$ is $\ell$.  The square of its index in
$\mathcal O_K$ divides this prime, by the change-of-integral-basis
formula for discriminants.  The index is therefore one, proving
$\operatorname{disc}(K)=\ell$.

Here is the class-number bound in full.  A nonzero integral ideal
$\mathfrak a$ embeds in $\mathbb R^2$ as a lattice of covolume
$\sqrt\ell\,N\mathfrak a$, by the discriminant formula and
the index $[\mathcal O_K:\mathfrak a]=N\mathfrak a$.
The symmetric diamond $|x_1|+|x_2|<t$ has area $2t^2$.
If this exceeds four times the covolume, it contains a nonzero
lattice point: half the diamond has area greater than the
covolume, so two of its points have the same image in a
fundamental parallelogram, and their difference is a nonzero
lattice point in the diamond.  This is the required
two-dimensional convex-body argument.  Consequently there is
$0\ne a\in\mathfrak a$ with
\[
 |N_{K/\mathbb Q}a|\le\tfrac12\sqrt\ell\,N\mathfrak a.
 \tag{CZ.7}
\]
Indeed $|x_1x_2|\le t^2/4$, and the bound at the limiting
value follows because nonzero absolute norms are positive integers.
The integral ideal $(a)\mathfrak a^{-1}$ represents the inverse
class and has norm at most $\sqrt\ell/2$.  Every ideal class
therefore has an integral representative of norm at most
$X=\sqrt\ell$.

For an integer $n\ge1$, the number of ideals of norm $n$ in a
quadratic field is at most the divisor number $d(n)$.
At a rational prime $q$, the split case has $r+1$ ideals of norm
$q^r$, the ramified case has one, and the inert case has at most
one.  Unique ideal factorization multiplies these bounds.
It follows that
\[
 \begin{split}
 h_K&\le\sum_{n\le X}d(n)
     =\sum_{a\le X}\lfloor X/a\rfloor
     \le X\sum_{a\le X}\frac1a\\
    &\le X(1+\log X)<X^2=\ell.
 \end{split}
 \tag{CZ.8}
\]
The harmonic-sum inequality follows by comparison with the integral
of $1/t$, and the strict inequality uses $\log X<X-1$ for
$X>1$.

Finally, the finite image of an unramified homomorphism as stated
defines a finite abelian extension of $K$.  At real places its
inertia has order at most two and therefore maps to zero in
$M[\ell]$.  The extension is consequently a quotient of the
ordinary Hilbert class field.  Its image is annihilated by $h_K$
as well as by $\ell$.  Since $0<h_K<\ell$, these two integers
are coprime, and the image is zero.
\end{proof}

\begin{proposition}[CZ4: residual-trivial vanishing]
\label{prop:cz-trivial-vanishing}
For $e\ge1$ with $\ell-1\mid e$, one has
$\mathcal S_e(O)=0$ and hence
$\texttt{Chr.\allowbreak V}(\ell,e,1)=0$.
\end{proposition}

\begin{proof}
Let $x\in\mathcal S_e(O)$ be annihilated by $\ell$ and use
CZ1 to represent it by an $M[\ell]$-valued cocycle $c$.
For every $g\in G$, the reduction of $\chi_\ell(g)^{-e}$
modulo $\ell$ is one because $\ell-1\mid e$.
Thus $\alpha(g)-1$ is a multiple of $\ell$ in $O$ and the
action on $M[\ell]$ is trivial.  The cocycle $c$ is therefore
an additive homomorphism.
At every $q\ne\ell$, CZ2 says $\alpha$ is trivial on
$I_A$; its local coboundary condition (CZ.3) then says
$c|_{I_A}=0$.  It is an unramified-outside-$\ell$ homomorphism,
so CZ2 gives $c|_{G_\infty}=0$.  The injectivity proved in
the same proposition gives $x=0$.  By CZ1 this kills the whole
Selmer group.  Its original $O$-length is consequently zero.
\end{proof}

\begin{proposition}[CZ5: residual-quadratic vanishing]
\label{prop:cz-quadratic-vanishing}
Suppose $\ell\ge5$, $e$ is even,
$\ell-1\mid2e$ and $\ell-1\nmid e$.  Then
$\mathcal S_e(O)=0$ and hence
$\texttt{Chr.\allowbreak V}(\ell,e,1)=0$.
\end{proposition}

\begin{proof}
Write $\ell-1=2a$.  The divisibility assumptions give
$e=ak$ with $k$ odd.  Since $e$ is even, $a$ is even and
$\ell\equiv1\pmod4$.  The residual character of
$\alpha=\chi_\ell^{-e}$ is exactly
$\eta\circ\overline\chi_\ell$, so it is trivial on the
actual $G_K$ in CZ3.  In particular $G_K$ acts trivially on
$M[\ell]$.

Choose $u\in\mathbb Z_\ell^\times$ with nonsquare reduction.
For any $g$ with $\chi_\ell(g)=u$, the scalar
$\alpha(g)-1$ is a unit of $O$: it is the image of the unit
$u^{-e}-1\in\mathbb Z_\ell$, whose reduction is $-2\ne0$.
Fix one place $A_0\mid\ell$ and, using CZ2, choose
$\delta\in I_{A_0}$ with cyclotomic value $u$.  Then
$\delta\notin G_K$.

Let $x\in\mathcal S_e(O)$ be killed by $\ell$ and choose an
$M[\ell]$-valued representative $c$ by CZ1.  Subtract the
coboundary of
$b=(\alpha(\delta)-1)^{-1}c(\delta)\in M[\ell]$.
The resulting representative, still denoted $c$, has
\[
                  c(\delta)=0,\qquad \ell c(g)=0
                  \quad(g\in G).
                  \tag{CZ.9}
\]
It retains all its local conditions.

Fix any $A\mid\ell$ and its local witness $m_A$ from (CZ.3)
on $D_A$.  Choose $t\in I_A$ with $\chi_\ell(t)=u$.
Then
\[
       0=\ell c(t)=(\alpha(t)-1)\ell m_A
       \quad\Longrightarrow\quad \ell m_A=0,
       \tag{CZ.10}
\]
because $\alpha(t)-1$ is a unit.  For $g\in D_A\cap G_K$
the action on $m_A\in M[\ell]$ is trivial; hence
$c(g)=0$.  At primes different from $\ell$, $\alpha$ is
trivial on inertia and (CZ.3) likewise gives $c|_{I_A}=0$.
Restriction of $c$ to $G_K$ is an additive homomorphism, since
the coefficient action is trivial there.  The preceding statements
make it unramified at every finite place of $K$ (and even zero
on its decomposition groups above $\ell$).  CZ3 therefore gives
$c|_{G_K}=0$.

The subgroup $G_K$ has index two and $\delta$ lies in the other
coset.  For any $g\notin G_K$ write $g=\delta h$ with
$h\in G_K$; then
$c(g)=c(\delta)+\alpha(\delta)c(h)=0$ by (CZ.9).
Thus $c=0$, so $x=0$.  Applying the primary-torsion argument
of CZ1 proves the full vanishing and the stated length equality.
\end{proof}

\begin{theorem}[CZ6: the literal exceptional count, uniformly with exponent zero]
\label{thm:cz-count-z}
The source predicates \texttt{Prop\allowbreak Z} and \texttt{Count\allowbreak Z} hold.
In \texttt{Count\allowbreak Z}, one may choose $C=0$ for every primitive
character, before choosing the host, twist cutoff, algebraic
closure, prime, coefficient ring, twist or character in the host.
In fact the counted length vanishes identically in this case.
\end{theorem}

\begin{proof}
For \texttt{Prop\allowbreak Z}, fix its original $\Omega,\ell,c,e$ with
$\ell\ge5$, $e\ge2$, $e$ even and $\ell-1\mid2e$;
here $O=c.O$.  If $\ell-1\mid e$, apply CZ4.  Otherwise
apply CZ5.  Both return the exact length
$\texttt{Chr.\allowbreak V}(\ell,e,1)=0$ on the original objects.

Now fix any data in the universal statement \texttt{Count\allowbreak Shape}
for the case \texttt{Case\allowbreak Z}.  By its definition this case says
\[
          \ell\notin B_\chi,\qquad
          \ell-1\mid2e,\qquad \psi=1.
          \tag{CZ.11}
\]
The exceptional prime set $B_\chi$ contains $3$; the count
statement already assumes the prime $\ell\ne2$.  Thus
$\ell\ge5$.  Membership $\psi\in h.X(c,e)$ includes
$\psi(-1)=(-1)^e$.  Because $\psi=1$ and $O$ has
characteristic zero, this forces $e$ even; the given $e>0$
then gives $e\ge2$.  The first part of the theorem applies.

The Galois character attached to the trivial Dirichlet character
is again trivial.  The source's \texttt{Vlen} in this case is
$\texttt{Chr.\allowbreak V'}(\ell,h.N,e,1)$, whose local conditions impose
decomposition restrictions above $\ell h.N$.  These include
all decomposition restrictions above $\ell$ and the same
inertia restrictions, so its Selmer group is a subgroup of the
zero group just proved.  Hence
\[
 \texttt{Vlen}(\Omega,c,e,1)=0,
 \qquad \texttt{Vn}(\Omega,c,e,1)=0.
 \tag{CZ.12}
\]
The first equality proves finiteness; the second follows from
the source's normalization of the length, whose numerator is
zero.  Therefore its actual predicate \texttt{VLe} holds with
left side $\texttt{Vn}\log\ell=0$ and right side
$((\ell:\mathbb R)h.N e)^0=1$.
Choosing $C=0$ before all the universally quantified data proves
\texttt{Count\allowbreak Z}.  No threshold, host, splitting field, coefficient
extension or discriminant is replaced in this last step.
\end{proof}

\paragraph{Closed dependency and remaining work.}
This closes one entire input of the source's implication from the
counts to \texttt{Point\allowbreak HV}.  The other count suppliers
\texttt{Count\allowbreak I}, \texttt{Count\allowbreak II}, \texttt{Count\allowbreak III} and
\texttt{Count\allowbreak V} retain their own obligations; in particular this
chapter does not supply the automorphy and host-rigidity inputs
used for \texttt{Count\allowbreak V}.  The current vanishing proof is independent
of those inputs and of \texttt{Lemma\allowbreak K}.


\section{Dependency closure after the shared cohomology and tower reconstruction}
\label{sec:dependency-audit-v25}

\paragraph{A complete quantitative supplier.}
The new chapters prove the chain
\[
 \begin{aligned}
 \text{Hermitian lattice and ordered nerve (TB)}
     &\longrightarrow\mathtt{TorsBoundPow}(22),\\
 \text{classical $K$-theory comparisons (KB)}
     &\longrightarrow\mathtt{KBridge},\\
 \mathtt{TorsBoundPow}(22)+\mathtt{KBridge}+\text{LD/LK}
     &\longrightarrow\mathtt{LemmaK}(31).
 \end{aligned}
\]
The torsion subgroup is the full integral torsion of the actual
$\operatorname{GL}_{4e+4}(\mathcal O_L)$ homology. The comparison
module, maps, Hurewicz kernel exponent and finite-rank homology
identification are constructed. The final compact cohomology is
that of the original $L,\Omega,D$, with the original discriminant
and degree. The constant $31$ is uniform before all of these choices.
The precise Quillen, Borel, Dwyer--Friedlander, Arlettaz, stability,
class-field and finite-duality results remain named classical leaves.

\paragraph{Closing the simultaneous window.}
TW1 proves that classical Ferrero--Washington $\mu=0$, together
with the ordinary Iwasawa module theorem, supplies exactly the
bounded class-group rank interface, including the prime $2$.
TW2--TW4 prove the uniform two-primary generator bound, strict
negative-twist coinvariant comparison, finite-layer growth and
the at-most-$R$ center bound. TW5 gives one square window for a
finite union of exceptional twists. TW6--TW8 preserve every field
model and the quantifier order. The resulting endpoint is
\[
             \exists K:\mathtt{BConst},\quad
             \forall C_t\ge1,\quad\mathtt{ThmBab}(K,C_t).
\]
The two-primary and odd reference-layer fields are separately
quantified on this same window and prime pair. Neither an isolated
window nor a convenient substitute field is used.

\paragraph{Immediate downstream substitution.}
Take the $K$ supplied by TW8. Together with OB8's constant-$240$
odd-minus theorem, the existing VW proof now gives
\[
 \mathtt{PointHV}\ \Longrightarrow\
 \forall C_t\ge1,\ \forall\text{ primitive }\psi,\
                          \mathtt{VWinAt}(\psi,C_t).
\]
This displayed implication retains the exact VW definitions and
its already stated classical prime-pair roles. With $C_t=401$,
the unchanged VW4/QP22 rescaling gives the same fixed pair
throughout the arbitrarily late three-halves window with cutoff
exponent $400$. Thus the selected cohomological frontier has one
remaining major supplier, \texttt{Point\allowbreak HV}. It does not follow
that the arithmetic host or the final gap is already available.

\paragraph{The exceptional odd-plus count is also supplied.}
CZ proves actual Selmer vanishing, first when the residual twist
is trivial and then when it is quadratic. The latter identifies
the original quadratic field by a finite Gauss sum and proves
$h(\mathbb Q(\sqrt\ell))<\ell$ directly. Local coboundary
normalization at every original $\ell$-adic place then gives
the full vanishing. The literal conclusion is
\texttt{Prop\allowbreak Z} and \texttt{Count\allowbreak Z}, with exponent $C=0$.
It is independent of the automorphy and rigidity inputs.

\paragraph{The exact remaining cohomological and global scope.}
For \texttt{Point\allowbreak HV}, LK, KB8 and CZ6 are available suppliers.
The actual \texttt{Count\allowbreak I/II/III/V}, their substantive arithmetic
inputs, and the remaining original-field case/layer and finite-sum
assembly remain explicit. The audit separates these tasks from
the completed \texttt{Thm\allowbreak Bab} and odd-minus chains.
Outside cohomology, the actual common family, genuine theta/local
state, local denominator tables, arithmetic gap and size/error
packages, global/spectral comparisons and exact lattice-metric
transport are retained. LN supplies its full conditional
denominator/restoration chain; WB and AN supply the full analytic
deduction from their arithmetic contract. No necessity or minimality
claim is made for this sufficient-condition map.

% Version 0.26.  Complete finite-character rank-one Selmer suppliers.
\section{Finite character descent and rank-one vector vanishing}
\label{sec:rank-one-cohomology}

\status{RV1--RV7 close the literal suppliers \texttt{So\allowbreak Fin} and
\texttt{So\allowbreak V} on every original algebraic closure and at every prime.
The proof uses the already reconstructed finite duality, local-K and
\texttt{Lemma\allowbreak K 31} suppliers.  It includes the finite-extension capture
of an algebraic-valued cocycle and the passage from a finite-length
torsion Selmer module to a vector coboundary.  RV1 also proves the
exact normalized-length transport needed when an auxiliary coefficient
field is used in the symplectic count.  The automorphy, rigidity,
counting and other cohomological packages in that count remain separate.}

Fix a prime $\ell$, a number field $k$, and its specified algebraic
closure $\Omega$.  Write $G_k=\operatorname{Gal}(\Omega/k)$ with its
profinite topology and $\chi=\chi_{k,\ell}:G_k\to\mathbb Z_\ell^\times$
for the actual cyclotomic character.  A coefficient ring $O$ will be
a discrete valuation ring, with its given $\mathbb Z_\ell$-algebra
structure, finite as a $\mathbb Z_\ell$-module.  As in the source,
a discrete valuation ring here is not a field.  Put
\[
 A_O=O\otimes_{\mathbb Z_\ell}(\mathbb Q_\ell/\mathbb Z_\ell),
 \qquad
 e_O=\operatorname{length}_O(O/\ell O).
 \tag{RV.1}
\]
Give $A_O$ the discrete topology.  For a character
$\psi:G_k\to O^\times$ with open kernel and an integer $e\ge1$, set
\[
 \alpha(g)=\psi(g)^{-1}\chi_O(g)^{-e},\qquad
 \rho(g)a=\alpha(g)a.
 \tag{RV.2}
\]
Here $\chi_O$ is obtained from the original $\chi$ by the specified
coefficient map; in particular the finite character is inverted.

For a valuation subring $A\subset\Omega$, let $D_A$ and $I_A$ denote
the actual decomposition and inertia subgroups of $G_k$.  Define
\[
 \begin{aligned}
 \mathscr D_\ell(k,\Omega)&=
 \{I_A:A\subset\Omega\text{ a valuation subring}\}
 \ \cup\
 \{D_A:A(\ell)<1\},\\
 \mathcal S(k,\Omega;\alpha)&=
 \bigcap_{U\in\mathscr D_\ell(k,\Omega)}
 \ker\!\left(H^1_{\rm cts}(G_k,A_O(\alpha))
                  \longrightarrow H^1_{\rm cts}(U,A_O(\alpha))\right).
 \end{aligned}
 \tag{RV.3}
\]
Thus a class in (RV.3) is represented by a continuous cocycle $z$
such that, for each indicated subgroup $U$, there is an $m_U\in A_O$
with $z(g)=(\alpha(g)-1)m_U$ for all $g\in U$.
These are precisely the source's
\texttt{Loc\allowbreak Sel rho top (sha\allowbreak Fam k Omega ell)} and
\texttt{Chr.\allowbreak V ell e psi}; the latter is the extended
$O$-length of (RV.3).
No finite set of valuation subrings is substituted in this definition.

\subsection{Coefficients, normalized lengths and the finite duality join}

\begin{proposition}[RV1: exact coefficient extension and normalized length]
\label{claim:RV1}
The map $\mathbb Z_\ell\to O$ is injective, $O$ is finite free over
$\mathbb Z_\ell$, and $e_O=v_O(\ell)$ is a positive finite integer.
With $K=\operatorname{Frac}O$, there is a natural identification
\[
             A_O\simeq K/O,\qquad
 \operatorname{length}_O A_O[a]=v_O(a)\quad(0\ne a\in O).
 \tag{RV.4}
\]
Let $O\hookrightarrow O'$ be an injective compatible map of such
coefficient rings.  Put $b=v_{O'}(\varpi)$ for a uniformizer
$\varpi$ of $O$.  For every profinite group, every discrete
$O$-module representation, and every fixed family of local
subgroups, extension of coefficients gives the natural isomorphism
\[
 O'\otimes_O\mathcal S_{\mathscr D}(G,M)
       \simeq\mathcal S_{\mathscr D}(G,O'\otimes_O M).
 \tag{RV.5}
\]
For the rank-one coefficients in (RV.2), the right side uses the
transported character and the same original Galois and local groups.
Their extended lengths satisfy
\[
 \operatorname{length}_{O'}\mathcal S'
       =b\,\operatorname{length}_O\mathcal S,\qquad
 e_{O'}=b e_O.
 \tag{RV.6}
\]
Consequently finiteness is equivalent on the two sides, and, whenever
the lengths are finite,
\[
       \frac{\operatorname{length}_{O'}\mathcal S'}{e_{O'}}
          =\frac{\operatorname{length}_O\mathcal S}{e_O}.
 \tag{RV.7}
\]
The assertions apply to the source's level-$\ell N$ module
\texttt{Chr.\allowbreak V'} as well as to \texttt{Chr.\allowbreak V}.
\end{proposition}

\begin{proof}
A nonzero kernel of $\mathbb Z_\ell\to O$ would be its maximal
ideal, because $O$ is a domain.  Then $O$ would be a finite-dimensional
domain over $\mathbb F_\ell$, hence a field, contrary to the
coefficient hypothesis.  The map is therefore injective.
The finite torsion-free module $O$ over the principal ideal ring
$\mathbb Z_\ell$ is free.  The element $\ell$ is not a unit in $O$:
otherwise $O=\ell O$ and finite-module Nakayama over $\mathbb Z_\ell$
would give $O=0$.  Thus $0<v_O(\ell)<\infty$.
Factoring $\ell$ into a unit times a power of a uniformizer proves
the formula for $e_O$.

The finite-dimensional $\mathbb Q_\ell$-algebra $O[1/\ell]$ is a
domain and therefore a field: multiplication by any nonzero element
is injective and hence surjective.  It is $K$.
Tensor the exact sequence
$0\to\mathbb Z_\ell\to\mathbb Q_\ell\to
\mathbb Q_\ell/\mathbb Z_\ell\to0$ with the free module $O$.
This gives the first identification in (RV.4).
The $a$-torsion is $a^{-1}O/O$, isomorphic to $O/aO$.
In particular it has the asserted length, and multiplication by any
nonzero $a$ on $A_O$ is surjective.
The residue field of $O$ is finite, since it is a quotient of the
finite additive group $O/\ell O$.

The extension $O'$ is finite over $O$: any finite
$\mathbb Z_\ell$-generating set of $O'$ also generates it over $O$.
It is torsion-free, hence finite free.  The map is local.
Indeed a maximal ideal in an integral extension contracts to a
maximal ideal, so the contraction of the maximal ideal of $O'$ is
the unique maximal ideal of $O$.  Therefore $b$ is positive.

Choose an $O$-basis of $O'$.  Tensoring a discrete module with $O'$
then identifies it with a finite direct sum of copies of that module.
Continuous functions into this sum are exactly tuples of continuous
functions, so in every degree
\[
 O'\otimes_O C^q_{\rm cts}(G,M)
        \simeq C^q_{\rm cts}(G,O'\otimes_O M).
 \tag{RV.8}
\]
The natural map in (RV.8) commutes with the cochain differential
and every restriction map.  Finite freeness makes it exact.
In the chosen basis all restrictions act componentwise, so their
simultaneous kernels commute with this base change even when
$\mathscr D$ is infinite.  This proves (RV.5).
Tensor associativity identifies $O'\otimes_O A_O$ with $A_{O'}$;
the scalar action is the given transported action, not an arbitrary
isomorphic representation.

A finite composition series over $O$ has factors $O/\varpi O$.
After the exact base change, each factor is
$O'/\varpi O'$ and has $O'$-length $b$, proving the length formula
when the length is finite.
For completeness, infinite length is also preserved.  The element
$1\in O'$ is primitive in the finite free $O$-module $O'$:
if it belonged to $\varpi O'$, then $\varpi$ would be a unit in $O'$.
It can therefore be included in an $O$-basis.  It follows that
the map $M\to O'\otimes_O M$, $m\mapsto1\otimes m$, is split
injective over $O$.  If the latter module had finite $O'$-length,
it would have finite $O$-length, since the finite extension of
residue fields makes every simple $O'$-factor have finite $O$-length.
Then $M$ would have finite length as well.  This proves the extended
version of the first equality in (RV.6).
Finally write $\ell=u\varpi^{e_O}$ and apply $v_{O'}$ to obtain
the second equality.  Equation (RV.7) follows with both positive
denominators in place.  Replacing the local family by
$\mathscr D_{\ell N}$ in the same proof gives the level assertion.
\end{proof}

\begin{proposition}[RV2: the pure negative-twist strict Selmer module is finite]
\label{claim:RV2}
Suppose a number field $L\subset\Omega$ contains a primitive
$\ell$-th root of unity for odd $\ell$, or a primitive fourth
root for $\ell=2$.  For every $e\ge1$ and every original
$D:\texttt{LData}\ L\ \Omega\ \ell$,
\[
 \operatorname{length}_{\mathbb Z_\ell}
 \mathcal S(L,\Omega;\chi_{L,\ell}^{-e})
 \le
 \operatorname{length}_{\mathbb Z_\ell}
 H^2_{\rm cts}(G_{L,S_\ell},\mathbb Z_\ell(e+1))
 <\infty .
 \tag{RV.9}
\]
Here the coefficient on the left is $\mathbb Q_\ell/\mathbb Z_\ell$,
and the group on the right is the original compact-cochain group
defined using $D$.
In particular, if $n=[L:\mathbb Q]$ and
$\Delta_L=|\operatorname{disc}L|$, its left side is at most
\[
       \frac{\exp\!\big((2en)^{31}\log(2\Delta_L)\big)}
                       {\log\ell}.
 \tag{RV.10}
\]
\end{proposition}

\begin{proof}
The roots-of-unity hypothesis implies that $L$ is totally complex,
including when $\ell=2$.
The cyclotomic character is unramified away from $\ell$.
For instance, at a residue characteristic prime to $\ell$, the
polynomial $X^{\ell^r}-1$ has invertible derivative at every root,
so these roots lie in the maximal unramified local extension.
This also constructs the descended character in an \texttt{LData}
record; choose the finitely many primes above $\ell$ and one
extension of each valuation to $\Omega$.

On an inertia group away from $\ell$ the pure negative twist is
trivial.  A local coboundary there is therefore identically zero.
A cocycle with this property descends through the closed normal
subgroup generated by those inertia groups to $G_{L,S_\ell}$.
To see the normal-closure assertion directly, if a cocycle vanishes
on an inertia element acting trivially on the coefficient, its
cocycle identity makes it vanish on every conjugate; sums, inverses
and continuity then give the closed normal subgroup.
Conversely an inflated cocycle has all these inertia restrictions
zero.  Inflation is injective on $H^1$, by lifting the same
coboundary witness through the surjective quotient map.
At primes above $\ell$, its remaining condition is exactly the
strict local kernel defining $\operatorname{Sha}^1(D,e)$.
Changing a chosen valuation to a conjugate does not change vanishing
of a restriction: if $z|_U=\delta m$, then on $tUt^{-1}$ the
witness is $\alpha(t)m-z(t)$.
Thus the all-valuation module in (RV.9) is the original
$\operatorname{Sha}^1(D,e)$.

We recall the finite-to-integral duality argument explicitly.
Put $T_r=\mathbb Z/\ell^r(e+1)$ and
$A_r=(\mathbb Q_\ell/\mathbb Z_\ell)(-e)[\ell^r]$.
The Cartier dual of $T_r$ is $A_r$, since its twist is
$1-(e+1)=-e$.
LD3 gives perfect finite pairings
\[
 \operatorname{Sha}^2(L,T_r)\times
       \operatorname{Sha}^1(L,A_r)\longrightarrow\mathbb Q/\mathbb Z.
 \tag{RV.11}
\]
Reduction on $T_r$ is paired with inclusion on $A_r$.
This compatibility follows from the coefficient evaluation followed
by the same local invariant maps, so these are not unrelated choices
of perfect pairings.
KB6 transports the original decomposition groups, through their
chosen henselizations and completions, to the finite native local
fields of LK2--LK3.  It preserves roots of unity, cup products and
the actual restriction maps at every $r$, including at two.

KB1 supplies finite global $H^1$ at every coefficient level.
The local counterparts are finite by LK2.
Thus the relevant inverse systems of $H^1$ are Mittag--Leffler.
Continuous integral cochains are the inverse limit of the finite
cochains; each transition on cochains is surjective by lifting a
function on its finitely many clopen fibers.
The resulting $\varprojlim^1 H^1$ obstruction vanishes, globally
and locally.  Hence
\[
 \operatorname{Sha}^2(L,\mathbb Z_\ell(e+1))
       =\varprojlim_r\operatorname{Sha}^2(L,T_r).
 \tag{RV.12}
\]
For discrete torsion coefficients, continuous cochains commute with
the filtered union: their domains are compact and their images
finite.  There are finitely many local places in this presentation,
and filtered colimits are exact.  Therefore
\[
 \operatorname{Sha}^1(L,(\mathbb Q_\ell/\mathbb Z_\ell)(-e))
       =\varinjlim_r\operatorname{Sha}^1(L,A_r).
 \tag{RV.13}
\]
The compatible finite pairings identify the compact dual of
(RV.13) with (RV.12), a submodule of the actual integral $H^2$.
This is also the finite duality argument underlying (TW.12);
no Iwasawa or Ferrero--Washington theorem is used here.
KB8 makes that $H^2$ finite and bounds the logarithm of its
cardinality by the numerator of (RV.10).
A finite $\mathbb Z_\ell$-module of length $h$ has cardinality
$\ell^h$.  Duality preserves that length, proving (RV.9)--(RV.10).
\end{proof}

\subsection{The actual finite-character supplier}

\begin{theorem}[RV3: finite-character descent and literal \texttt{So\allowbreak Fin}]
\label{claim:RV3}
For every original $k,\Omega,\ell,O,\psi$ as in (RV.1)--(RV.2)
and every $e\ge1$,
\[
                    \texttt{Chr.\allowbreak V}\ \ell\ e\ \psi<\infty .
 \tag{RV.14}
\]
No assumption that the order of $\psi$ is prime to $\ell$ is made.
In particular the source's literal \texttt{So\allowbreak Fin Omega ell} holds
for every algebraic closure $\Omega$ of $\mathbb Q$ and every prime.
\end{theorem}

\begin{proof}
Let $N=\ell$ for odd $\ell$ and $N=4$ at two.  Inside the original
$\Omega$ take the fixed field
\[
 L=\Omega^{\,\ker\psi\cap\ker(\chi\bmod N)}.
 \tag{RV.15}
\]
Both kernels are open, so $L/k$ is finite Galois.
Its Galois group embeds in the product of two abelian character
images, so it is abelian.  The first kernel makes $\psi|_{G_L}=1$.
The second fixes every $N$-th root of unity in $\Omega$; hence $L$
contains a primitive such root.  It is totally complex and satisfies
the exact root hypothesis of RV2.
This is an extension of the original $k$ inside the original
$\Omega$, with the original embeddings.

Set $H=G_L$, $\Gamma=G_k$, $\Delta=\Gamma/H$, and $d=|\Delta|$.
There is an $h\in H$ with $\chi(h)^e\ne1$.
Indeed the mod-$\ell^r$ cyclotomic image of $G_L$ has order
$[L(\mu_{\ell^r}):L]\ge
\varphi(\ell^r)/[L:\mathbb Q]$, which tends to infinity.
The degree formula used here has a short direct proof.
The primitive $\ell^r$-th roots are the roots of
$P_r(X)=\sum_{j=0}^{\ell-1}X^{j\ell^{r-1}}$.
Modulo $\ell$, the polynomial $P_r(1+T)$ is
$T^{\ell^{r-1}(\ell-1)}$, and its constant coefficient is
$\ell$, not divisible by $\ell^2$.
Eisenstein's criterion makes it irreducible over $\mathbb Q$,
of degree $\varphi(\ell^r)$.
The image of $\chi|_H$ is therefore infinite, whereas the equation
$u^e=1$ has at most $e$ roots in $\mathbb Q_\ell$.
Put
\[
                       a=\chi(h)^e-1\in\mathbb Z_\ell\setminus\{0\}.
 \tag{RV.16}
\]
On $H$ the coefficient action is the pure negative twist.
Consequently
$B=A_O^H\subset A_O[a]$: the scalar
$\chi(h)^{-e}-1$ differs from $a$ by a unit.
RV1 gives
\[
             \operatorname{length}_O B
                     \le e_O\,v_\ell(a)<\infty.
 \tag{RV.17}
\]
This supplies the needed $H^0$ control before using any restriction
kernel bound.

We give that bound without averaging over $\Delta$.
If a cocycle $z$ restricts to a coboundary on $H$, subtract that
coboundary so that $z|_H=0$.
The identity $hg=g(g^{-1}hg)$ shows
$\alpha(h)z(g)=z(g)$ for $h\in H$.
Also $z(gh)=z(g)$, so $z$ descends to a cocycle of $\Delta$
with values in $B$.
Two such normalized cocycles differ by a coboundary precisely when
their witnesses differ by an element of $B$.
This proves the exact kernel identification
\[
 \ker\!\left(H^1(\Gamma,A_O(\alpha))
                  \longrightarrow H^1(H,A_O(\alpha))\right)
                      \simeq H^1(\Delta,B).
 \tag{RV.18}
\]
A normalized finite-group cocycle is determined by its values on
the $d-1$ nonidentity elements.  Therefore
$\operatorname{length}_O H^1(\Delta,B)
\le(d-1)\operatorname{length}_O B$.
This remains valid when $\ell$ divides $d$.

Restriction sends the strict module (RV.3) into the strict module
over $L$.  For the same valuation subring $A\subset\Omega$, the
local groups over $L$ are $H\cap D_A$ and $H\cap I_A$; the
condition $A(\ell)<1$ is unchanged.  Thus every original local
coboundary witness restricts to the required witness.
Since $\psi|_H=1$, RV1 applied to the pure twist gives
\[
 \operatorname{length}_O\mathcal S(L,\Omega;\alpha|_H)
       =e_O\,
          \operatorname{length}_{\mathbb Z_\ell}
                  \mathcal S(L,\Omega;\chi_{L,\ell}^{-e}).
 \tag{RV.19}
\]
Its right side is finite by RV2.
The image and kernel of the restriction of (RV.3) therefore both
have finite length.  More explicitly, (RV.17)--(RV.19) yield
\[
 \frac{\operatorname{length}_O\mathcal S(k,\Omega;\alpha)}{e_O}
 \le
 \operatorname{length}_{\mathbb Z_\ell}
                  \mathcal S(L,\Omega;\chi_{L,\ell}^{-e})
                    +(d-1)v_\ell(a).
 \tag{RV.20}
\]
The numerator on the left is proved finite before this real-valued
normalization is used.
This proves (RV.14) for an arbitrary original coefficient ring.
Setting $k=\mathbb Q$ and $O=c.O$ proves the exact \texttt{So\allowbreak Fin}
predicate.  The estimate (RV.20) is a sufficient alternative to
the source's sharper generator-count restriction estimate; no
unstated prime-to-$\ell$ restriction or new finiteness hypothesis
has been introduced.
\end{proof}

\subsection{Capturing the coefficient field of a continuous cocycle}

\begin{proposition}[RV4: finite-extension capture from Baire and the cocycle equation]
\label{claim:RV4}
Let $\Gamma$ be a compact Hausdorff topological group,
$K_0\subset\overline{\mathbb Q}_\ell$ a finite extension of
$\mathbb Q_\ell$, and let
$\theta,c:\Gamma\to\overline{\mathbb Q}_\ell$ satisfy
\[
 \theta(g)\in K_0,\qquad \theta(1)=1,\qquad
 c\text{ continuous},\qquad
 c(gh)=c(g)+\theta(g)c(h).
 \tag{RV.21}
\]
Then a finite extension $E\subset\overline{\mathbb Q}_\ell$
contains $K_0$ and every value of $c$.
Here $\overline{\mathbb Q}_\ell$ is the algebraic closure with
its extended valuation, not its completion.
\end{proposition}

\begin{proof}
We first construct a countable family of finite subfields
containing every finite subset of $\overline{\mathbb Q}_\ell$.
Choose an embedded algebraic closure $\overline{\mathbb Q}$ of
$\mathbb Q$ in $\overline{\mathbb Q}_\ell$.
It is countable and dense.
For the density assertion, let $\alpha$ be algebraic over
$\mathbb Q_\ell$ with monic minimal polynomial $f$ of degree $n$.
Approximate the coefficients of $f$ by rational numbers, retaining
the leading coefficient $1$, so that the resulting
$g\in\mathbb Q[X]$ satisfies $|g(\alpha)|<\varepsilon^n$.
This is possible because $\mathbb Q$ is dense in $\mathbb Q_\ell$
and there are finitely many coefficients.
All roots $\beta_1,\ldots,\beta_n$ of $g$ lie in the embedded
$\overline{\mathbb Q}$, and
$\prod_i|\alpha-\beta_i|=|g(\alpha)|$.
At least one $\beta_i$ is within $\varepsilon$ of $\alpha$.

For $\alpha\notin\mathbb Q_\ell$, choose such a $\beta$ satisfying
\[
 |\alpha-\beta|<
       \min_{\alpha'\ne\alpha}
                         |\alpha-\alpha'|,
 \tag{RV.22}
\]
where $\alpha'$ ranges over its $\mathbb Q_\ell$-conjugates.
Every automorphism of $\overline{\mathbb Q}_\ell$ fixing
$\mathbb Q_\ell$ is an isometry, by uniqueness of the extended
absolute value.  If it also fixes $\beta$, then
$|\sigma\alpha-\alpha|\le|\alpha-\beta|$, so (RV.22) forces
$\sigma\alpha=\alpha$.
Finite Galois theory gives $\alpha\in\mathbb Q_\ell(\beta)$.
This proves the needed Krasner assertion directly.
If $\alpha\in\mathbb Q_\ell$, no adjoining is needed.
The fields generated over $\mathbb Q_\ell$ by finite subsets of
the countable $\overline{\mathbb Q}$ consequently form a
countable family $(K_i)$, cofinal for finite subsets of
$\overline{\mathbb Q}_\ell$.
Each $K_0K_i$ is finite over $\mathbb Q_\ell$, complete for
its extended absolute value, and hence closed in
$\overline{\mathbb Q}_\ell$.

The cocycle identity gives $c(1)=0$ and
$c(g^{-1})=-\theta(g^{-1})c(g)$.
Thus
\[
                 H_i=\{g\in\Gamma:c(g)\in K_0K_i\}
 \tag{RV.23}
\]
is a subgroup: products and inverses stay in the indicated field.
It is closed by continuity of $c$.
The $H_i$ cover $\Gamma$.
The Baire property of a compact Hausdorff space now implies that
some $H_i$ has nonempty interior.
One can see the needed Baire assertion directly: if a countable
closed cover had all members with empty interior, successive
nonempty open sets could be chosen with compact closures nested
inside the previous open set and avoiding the next member.
Compactness would give a point in all those closures, outside
the proposed cover.

A subgroup with nonempty interior is open, and an open subgroup
of a compact group has finite index.
Choose finitely many representatives $g_1,\ldots,g_s$ for its
left cosets and let
\[
                    E=K_0K_i(c(g_1),\ldots,c(g_s)).
 \tag{RV.24}
\]
This is a finite extension because all its generators are algebraic.
For $g=g_jh$ with $h\in H_i$, the identity (RV.21) gives
$c(g)=c(g_j)+\theta(g_j)c(h)\in E$.
This proves the assertion.
The subgroup structure in (RV.23) is essential: the argument
does not assert that every compact subset of
$\overline{\mathbb Q}_\ell$ is contained in a finite extension.
\end{proof}

\subsection{From torsion Selmer finiteness to an exact vector coboundary}

\begin{proposition}[RV5: a finite-length quotient controls vector cocycles]
\label{claim:RV5}
Let $E/\mathbb Q_\ell$ be finite, $O_E$ its valuation ring,
$\alpha:\Gamma\to O_E^\times$ a continuous scalar action, and
$\mathscr D$ any family of subgroups of a topological group $\Gamma$.
Suppose
$\mathcal S_{\mathscr D}(\Gamma,(E/O_E)(\alpha))$
has finite $O_E$-length.
If a continuous $c:\Gamma\to E$ is an $\alpha$-cocycle and is
a coboundary on every $U\in\mathscr D$, with witnesses in $E$,
then there is an $m\in E$ with
\[
                         c(g)=(\alpha(g)-1)m
                         \quad\text{for every }g.
 \tag{RV.25}
\]
\end{proposition}

\begin{proof}
The quotient map $\pi:E\to E/O_E$ is continuous for the discrete
quotient topology, since its kernel $O_E$ is open.
A finite-length module of length $h$ is killed by $\ell^h$:
each simple composition factor is killed by the maximal ideal,
which contains $\ell$, and induction along the series gives
the assertion.  Choose $N$ such that $\ell^N$ kills the strict
Selmer module in the hypothesis.

For each $n\ge0$, reduce
$\ell^{-(n+N)}c$ modulo $O_E$.
It is a continuous cocycle in that strict module; its local
witness on $U$ is the reduction of $\ell^{-(n+N)}m_U$.
Multiplication by $\ell^N$ kills its class, so for some $y_n\in E$
\[
            \ell^{-n}c(g)-(\alpha(g)-1)y_n\in O_E
                         \quad\text{for all }g.
 \tag{RV.26}
\]
Put $m_n=\ell^n y_n$ and $r=|\ell|=\ell^{-1}<1$.
Then
\[
                 |c(g)-(\alpha(g)-1)m_n|\le r^n
                         \quad\text{for all }g.
 \tag{RV.27}
\]

If $\alpha(g)=1$ for every $g$, (RV.27) directly gives $c=0$,
and $m=0$ works.
Otherwise choose $g_0$ with $d_0=\alpha(g_0)-1\ne0$ and set
$m=c(g_0)/d_0$.
For $d_g=\alpha(g)-1$, the exact identity
\[
 c(g)-d_gm
   =\big(c(g)-d_gm_n\big)
        -\frac{d_g}{d_0}\big(c(g_0)-d_0m_n\big)
 \tag{RV.28}
\]
bounds its absolute value by
$r^n\max(1,|d_g/d_0|)$ for every $n$.
Letting $n$ tend to infinity proves (RV.25).
The sequence of witnesses $m_n$ need not be asserted to converge:
the single explicit witness $m=c(g_0)/d_0$ gives the conclusion.
\end{proof}

\begin{proposition}[RV6: the vector theorem on an actual finite coefficient field]
\label{claim:RV6}
Let $E\subset\overline{\mathbb Q}_\ell$ be finite.
Let $\psi:G_k\to\overline{\mathbb Q}_\ell^\times$ have open
kernel and all its values in $E$, and let $e\ge1$.
Suppose $\theta,c:G_k\to\overline{\mathbb Q}_\ell$ satisfy
\[
 \theta(g)\psi(g)\chi(g)^e=1,\qquad
 c\text{ continuous},\qquad c(g)\in E,\qquad
 c(gh)=c(g)+\theta(g)c(h).
 \tag{RV.29}
\]
Suppose also that $c$ is a coboundary on every original inertia
group, and on every original decomposition group above $\ell$,
with witnesses initially allowed in $\overline{\mathbb Q}_\ell$.
Then a witness $m$ in this same $E$ makes $c$ a global coboundary.
This is the complete finite-field statement \texttt{so\allowbreak Vec}.
\end{proposition}

\begin{proof}
We record the coefficient-lattice facts needed here.
The finite extension $E$ has a unique extended absolute value,
is complete, and its valuation ring
$O_E=\{x\in E:|x|\le1\}$ is the integral closure of
$\mathbb Z_\ell$.
The last equality can also be checked from the minimal polynomial.
An integral element cannot have absolute value greater than one,
since then the leading term of a monic integral relation would
strictly dominate the others.
Conversely, if $|x|\le1$, all its conjugates have absolute value
at most one; the elementary symmetric functions show that its
monic minimal polynomial has coefficients in $\mathbb Z_\ell$.

For clarity, $O_E$ is finite free over $\mathbb Z_\ell$.
Choose a $\mathbb Q_\ell$-basis consisting of integral elements.
The traces of their products with any element of $O_E$ belong
to $\mathbb Z_\ell$.
Thus $O_E$ lies in the trace-dual lattice of that basis.
Nondegeneracy of the trace pairing in characteristic zero and
Noetherianity of $\mathbb Z_\ell$ make it a finite torsion-free,
hence free, module.
Its value group is discrete: the norm formula gives
$|x|^{[E:\mathbb Q_\ell]}
=|N_{E/\mathbb Q_\ell}(x)|\in\ell^{\mathbb Z}$.
The smallest positive additive valuation yields a uniformizer;
factoring by it shows that $O_E$ is a discrete valuation ring.
It is not a field because $0<|\ell|<1$.
It therefore has exactly the coefficient properties required in RV3,
and (RV.4) identifies $A_{O_E}$ with $E/O_E$.

Compactness of $G_k$ and openness of $\ker\psi$ make its quotient
finite.  Every value of $\psi$ is a root of unity and hence has
absolute value one.  Thus $\psi$ is a character
$\psi_E:G_k\to O_E^\times$ with exactly the same kernel.
The identity in (RV.29) uniquely determines
\[
                    \theta=\psi_E^{-1}\chi_{O_E}^{-e};
 \tag{RV.30}
\]
in particular $\theta$ takes values in $O_E^\times$ and is continuous.
The map $c$, with values in the subspace $E$, is continuous as
an $E$-valued map.
RV3 supplies finite length for the exact strict module
$\mathcal S(k,\Omega;(E/O_E)(\theta))$.

Each local witness can be taken in $E$ without enlarging it.
For a local subgroup $U$, if $\theta|_U=1$, its coboundary
condition already says $c|_U=0$, and the witness $0$ works.
Otherwise choose $g_0\in U$ with $\theta(g_0)\ne1$.
Any witness must equal
\[
                         \frac{c(g_0)}{\theta(g_0)-1}\in E.
 \tag{RV.31}
\]
It therefore gives the entire local coboundary identity in $E$.
This argument treats each original subgroup directly; there is no
attempt to adjoin infinitely many independently selected witnesses.
Apply RV5 with the unchanged family
$\mathscr D_\ell(k,\Omega)$ and the coefficient identification
above.  The resulting $m\in E$ proves the assertion.
\end{proof}

\begin{theorem}[RV7: the literal rank-one suppliers on every original closure]
\label{claim:RV7}
For every original algebraic closure $\Omega$ of $\mathbb Q$
and every prime $\ell$,
\[
       \texttt{So\allowbreak V}\ \Omega\ \ell
              \quad\text{and}\quad
       \texttt{So\allowbreak Fin}\ \Omega\ \ell
 \tag{RV.32}
\]
hold.  The vector statement in fact holds over every number
field $k$ in the full generality of \texttt{so\allowbreak Bar}.
Under an injective extension of coefficient rings, the original
\texttt{Vlen} is finite if and only if its transported version is
finite, and its normalized \texttt{Vn} is unchanged.
More generally, the same conclusion holds for two coefficient rings
inside the prescribed $\overline{\mathbb Q}_\ell$ whose Dirichlet
characters have equal values there.  This coefficient assertion
allows every $e\ge0$ and is the literal \texttt{coef\allowbreak Indep} statement.
\end{theorem}

\begin{proof}
RV3 gives \texttt{So\allowbreak Fin}.
For the vector statement, take exactly the original
$k,\Omega,\ell,\psi,e,\theta,c$ and its original local witnesses
in the definition of \texttt{so\allowbreak Bar}.
The finite image of $\psi$ generates a finite
$K_0\subset\overline{\mathbb Q}_\ell$.
Since $\theta(g)=(\psi(g)\chi(g)^e)^{-1}$, every value of
$\theta$ lies in $K_0$ and $\theta(1)=1$.
RV4 therefore supplies a finite $E\supset K_0$ containing all
values of $c$.
RV6 applies on this same $E$ and produces a global witness.
Regard that witness as an element of the originally prescribed
$\texttt{Padic\allowbreak Alg\allowbreak Cl}\ \ell$; this is the exact conclusion of
\texttt{so\allowbreak Bar}.  Specializing $k=\mathbb Q$ is precisely
\texttt{So\allowbreak V Omega ell}.

For the last statement, the Dirichlet character transported by an
injective coefficient map gives the transported Galois character
by composition with the same cyclotomic restriction map.
The source's \texttt{Vlen} is the length of the level-$\ell N$
module \texttt{Chr.\allowbreak V'}, and \texttt{Vn} is that length divided by
$e_O$.  Apply RV1 to that unchanged local family.
Both lengths are finite together, and (RV.7) is the required
exact normalized equality in the finite case.  If both lengths
are infinite, the source's convention $\infty.\texttt{to\allowbreak Nat}=0$
makes both displayed real normalizations zero; this convention
does not assert a finite cohomological bound.
For two possibly nonnested coefficient rings, take the compositum
of their fraction fields inside $\overline{\mathbb Q}_\ell$.
It is finite over $\mathbb Q_\ell$ by RV1.
Its valuation ring, constructed in the proof of RV6, contains both
original rings because their elements are integral over
$\mathbb Z_\ell$.  The two transported characters there are
identical by the given equality of their original values.
Apply RV1 to the two inclusions and compare through this common
ring.  This proves the full \texttt{coef\allowbreak Indep} assertion, without
changing the original $\Omega$, level or local family.
In particular this step returns a
bound from an auxiliary coefficient field to the original ring;
it does not merely produce an estimate on the auxiliary ring.
\end{proof}

\paragraph{Classical input and scope.}
The new local-field background used above is the uniqueness of the
extended absolute value on a finite extension of $\mathbb Q_\ell$,
its completeness and the norm formula.  See J.\,S.~Milne,
\emph{Algebraic Number Theory}, version 3.08 (2020),
Theorem 7.38 and Corollary 7.40, pp.~123--124,
\url{https://www.jmilne.org/math/CourseNotes/ANT.pdf}.
The comparison in (RV.22) is proved here; it is also Proposition
7.60 of those notes, pp.~131--132.
The elementary prime-power cyclotomic degree formula is proved
by Eisenstein's criterion in RV3; compare Proposition 6.2(a)
of those notes, p.~96.
The finite Galois correspondence, elementary cochain algebra,
finite-module structure over a discrete valuation ring, and the
Baire argument were used in their explicit forms above.
LD3, LK2--LK3 and KB1/KB6/KB8 retain their own previously recorded
classical theorem leaves.

No Iwasawa control theorem, arbitrary compact-set capture principle,
vanishing of an unconstructed automorphic representation, or
prime-to-$\ell$ hypothesis on the finite character is an input.
The two predicates in (RV.32) are now actual supplied leaves in
the symplectic count.  Its \texttt{Cor\allowbreak X}, other cohomology package,
finite algebra models, automorphy, rigidity and trace-counting
requirements are separate.  This chapter alone does not assert
\texttt{Count\allowbreak II}, \texttt{Count\allowbreak III}, \texttt{Point\allowbreak HV}, RH or GRH.


% Version 0.26. The literal Bernoulli-unit and residual CorX supplier.
% Pinned source: 7cbf7a0c32fca1badf05931879a37c3f8a2b03b4.
\section{Bernoulli units and residual unramified cocycles}
\label{sec:bernoulli-residual}

\status{This section proves the prescribed Bernoulli-unit theorem
\texttt{Lemma\allowbreak X}, the full residual unramified cocycle theorem,
and the literal \texttt{Cor\allowbreak X} interface.  The power-of-two conductor,
both inverse characters, the original coefficient and residue fields,
and the prime-to-$\ell$ relative-index local condition are retained.
The Bernoulli-unit argument is finite and elementary.  Its use in
cohomology rests on the Jacobi ideal relation already proved in OS1
and the finite Hilbert class field theorem stated below.}

\subsection{The exact statements and the classical boundary}

Put $G=\operatorname{Gal}(\Omega/\mathbb Q)$ for the specified algebraic
closure.  Write $I_A$ for inertia at a valuation subring $A$ of $\Omega$.
For a character $\gamma:G\to k^\times$, a cocycle means a locally
constant function $z:G\to k$ satisfying
\[
                  z(gh)=z(g)+\gamma(g)z(h).
                  \tag{BR.1}
\]
This is the source's \texttt{Is\allowbreak Coc}; no topology on the arbitrary field
$k$ has to be added to its statement.  Its condition
$\operatorname{UnrEv}(\ell,\gamma,z)$ says that for every $A$ there is
a subgroup $H_A\le G$ such that
\[
 \begin{gathered}
  K_A=I_A\cap\ker\gamma,\qquad
  \ell\nmid[K_A:K_A\cap H_A],\\
  z|_{K_A\cap H_A}=0.
 \end{gathered}
 \tag{BR.2}
\]
The natural-number relative index in the source is zero when the index
is infinite.  Thus the nondivisibility assertion in (BR.2) in
particular makes this index finite.  There is no requirement that
$H_A$ itself be contained in $K_A$.

Dirichlet characters are extended by zero off the units.  If $\varphi$
has level $f$ and $\xi$ has level $2^n$, let
\[
 \chi=\operatorname{prodX}(\varphi,\xi),\qquad
 \chi(a)=\varphi(a\bmod f)\xi(a\bmod 2^n).
 \tag{BR.3}
\]
The equality holds also at nonunits: a nonunit of the product level is
a nonunit at at least one factor.  We use the source normalization
\[
 B_{1,\chi}=\frac1M\sum_{a=0}^{M-1}a\chi(a)
 \quad\text{if }\chi\ne1\text{ has level }M.
 \tag{BR.4}
\]
Indeed the defining first Bernoulli polynomial is $T-\tfrac12$, and
$\sum_a\chi(a)=0$ for a nontrivial character.  Formula (BR.4) is
therefore the literal \texttt{gen\allowbreak B 1}, including for imprimitive
characters at their specified level.
When a Dirichlet character is evaluated on $g\in G$, its argument
means the cyclotomic residue $\operatorname{cls}_M(g)$ at its stated
level $M$.  Thus
$\operatorname{redChar}(\chi)(g)=\operatorname{red}(\chi(g))$.

The data $\operatorname{XData}(\ell,d,\varphi,\xi)$ mean:
$\ell$ is prime, $d>0$, $4\mid\ell^d-1$, $f$ is odd,
$\varphi$ and $\xi$ are primitive at their respective levels,
\[
 \operatorname{ord}(\varphi)\mid\ell^d-1,\qquad
 f\,4^{v_2(\ell^d-1)}\le2^n,\qquad
 \varphi(-1)\xi(-1)=-1.
 \tag{BR.5}
\]
In particular $f>0$.  The precise Bernoulli-unit assertion is
$|B_{1,\operatorname{prodX}(\varphi,\xi)}|_\ell=1$ in the specified
$\overline{\mathbb Q}_\ell$.

\paragraph{Accepted classical input.}
For a number field $F$, the finite Hilbert class field is its maximal
abelian extension unramified at all finite places and split at its
real places; its Artin group is the ordinary ideal class group.
We use this in the form of Milne,
\href{https://www.jmilne.org/math/CourseNotes/CFT.pdf}{\emph{Class
Field Theory}}, v4.03, V, Theorems~3.5--3.6 and Example~3.9.
We also use the elementary cyclotomic identification
$\operatorname{Gal}(\mathbb Q(\mu_M)/\mathbb Q)
\simeq(\mathbb Z/M)^\times$, with $\sigma_b(\zeta)=\zeta^b$,
and extension of a finite Galois automorphism to the specified
algebraic closure; see Milne,
\href{https://www.jmilne.org/math/CourseNotes/ANT.pdf}{\emph{Algebraic
Number Theory}}, v3.08, \S6.
All transfers from these finite theorems to the actual cocycle and its
original field are given in BR4.  The needed annihilator is derived
from OS1; no general Stickelberger annihilation theorem is retained
as an additional input.

\subsection{Dyadic powers and a finite Frobenius projection}

\begin{proposition}[BR1: dyadic generation and character-sum projection]
\label{prop:br-dyadic}
Let $Q=1+2^e a$ with $e\ge2$ and $a$ odd.  For $r\ge e$, the
powers of $Q$ modulo $2^r$ are exactly the units congruent to $1$
modulo $2^e$.

Let $k$ have odd characteristic $\ell$, let $Q=\ell^d$ with
$v_2(Q-1)=e\ge2$, and let $n\ge e$.  For any finite family
$z_i,g_i\in k$ such that
$z_i=0$ or $z_i^{2^n}=1$, and $g_i^Q=g_i$, one has
\[
 \sum_i z_i g_i=0
 \quad\Longrightarrow\quad
 \sum_{\{i:z_i^{2^e}=1\}}z_i g_i=0.
 \tag{BR.6}
\]
Finally, put $n=m+e$ with $m\ge e\ge2$.  Suppose
$\eta:\mathbb Z/2^n\to k$ is multiplicative, is $1$ at $1$,
vanishes on nonunits, and
\[
                    \eta(1+2^{n-1})=-1.
                    \tag{BR.7}
\]
If $\eta(u)^{2^e}=1$, then
$u\equiv1$ or $-1\pmod{2^m}$.
\end{proposition}

\begin{proof}
Repeated squaring gives
\[
 Q^{2^j}=1+2^{e+j}a_j,\qquad a_j\text{ odd},\qquad
 a_{j+1}=a_j+2^{e+j-1}a_j^2.
 \tag{BR.8}
\]
The last summand is even, which proves the assertion about $a_j$
inductively.  Thus $Q$ has order exactly $2^{r-e}$ modulo $2^r$
when $r>e$, and order one when $r=e$.  Its powers lie in the
subgroup $1+2^e\mathbb Z/2^r$, whose size is $2^{r-e}$.
They consequently exhaust that subgroup.

Raise the assumed zero sum in (BR.6) to $Q^j$.  Frobenius is a
ring homomorphism and $g_i^{Q^j}=g_i$, so
$\sum_i z_i^{Q^j}g_i=0$.  By the first assertion, every exponent
$1+2^e t$ occurs modulo $2^n$ among the $Q^j$.  Zero entries cause
no problem, since all these exponents are positive.  Summing for
$0\le t<2^{n-e}$ gives
\[
 \begin{split}
 0&=\sum_i g_i z_i
       \sum_{t=0}^{2^{n-e}-1}(z_i^{2^e})^t\\
  &=2^{n-e}\sum_{\{i:z_i^{2^e}=1\}}z_i g_i.
 \end{split}
 \tag{BR.9}
\]
For a nonzero $z_i$ not satisfying $z_i^{2^e}=1$, the geometric
sum vanishes because its ratio has $2^{n-e}$th power one.
For $z_i=0$ the prefactor makes the term zero.  The scalar
$2^{n-e}$ is invertible in $k$, proving (BR.6).

Apply (BR.8) to $5=1+4$.
It generates all units congruent to $1$ modulo $4$ at level $2^n$.
Every odd $u$ therefore has a unique expression
$u=(-1)^\epsilon5^j$, with $j$ modulo $2^{n-2}$.
Moreover
$5^{2^{n-3}}=1+2^{n-1}$ modulo $2^n$.
By (BR.7), $\eta(5)$ has order exactly $2^{n-2}$:
its $2^{n-2}$th power is one and its $2^{n-3}$rd power is $-1$.
Since $\eta(-1)^2=1$, the condition
$\eta(u)^{2^e}=1$ implies
$2^{n-2}\mid j2^e$, or $2^{m-2}\mid j$.
The powers of $5$ with this property are $1$ modulo $2^m$,
again by (BR.8).  This proves the last assertion.
Nonunits cannot occur because their $\eta$-value is zero.
\end{proof}

\subsection{An integral partial Bernoulli sum}

\begin{proposition}[BR2: partial sums without dividing by the odd conductor]
\label{prop:br-partial}
Let $F$ be a field of characteristic zero, let $(f,N)=1$ with
$f,N>0$, and let $w:\mathbb N\to F$ be periodic modulo $f$.
Suppose a subring $B\subset F$ contains all $w(a)$ and $N^{-1}$.
There are $g:\mathbb Z/N\to B$ and $c\in B$ such that, for
$ft\le N$,
\[
 g(u)-g(u-ft)=tc-
           \begin{cases}w(\langle u\rangle_N),&
                          \langle u\rangle_N<ft,\\0,&
                          \langle u\rangle_N\ge ft,
           \end{cases}
 \tag{BR.10}
\]
where $0\le\langle u\rangle_N<N$.
For every $\xi:\mathbb Z/N\to F$ of sum zero,
\[
 \sum_{a=0}^{fN-1}a\,w(a)\xi(a)
             =fN\sum_{u\bmod N}\xi(u)g(u).
 \tag{BR.11}
\]
If $w(a)=\varphi(a\bmod f)$, $\xi$ is a nontrivial character
of level $N$, and $\varphi(-1)\xi(-1)=-1$, then the last sum is
$B_{1,\operatorname{prodX}(\varphi,\xi)}$.
\end{proposition}

\begin{proof}
Define, initially in $F$,
\[
 \begin{gathered}
 S(u)=\sum_{\substack{0\le a<fN\\a\equiv u\ (N)}}w(a),
 \qquad
 A(u)=\sum_{\substack{0\le a<fN\\a\equiv u\ (N)}}a\,w(a),\\
 c=S(0)/N,\qquad g(u)=\frac{A(u)-A(0)}{fN}.
 \end{gathered}
 \tag{BR.12}
\]
The $f$ integers in each residue class $u$ modulo $N$ run through
all residues modulo $f$, so $S(u)$ is independent of $u$.
This also follows by shifting by $f$, since $f$ generates
$\mathbb Z/N$ additively.

Here is the shift identity, including the terms crossing the endpoints.
Put $M=fN$ and
$\psi_u(a)=\mathbf1_{a\equiv u\ (N)}w(a)$.
For any function $h$ and any nonnegative $s$,
\[
 \sum_{a=0}^{M-1}h(a+s)
 =\sum_{a=0}^{M-1}h(a)
   +\sum_{a=0}^{s-1}\bigl(h(M+a)-h(a)\bigr).
 \tag{BR.13}
\]
It follows by writing the sum from $0$ to $M+s-1$ in its two
possible ways.  With $s=ft$ and $h(a)=a\psi_u(a)$, periodicity gives
$\psi_u(a+s)=\psi_{u-s}(a)$ and $\psi_u(M+a)=\psi_u(a)$.
Consequently (BR.13) becomes
\[
 A(u-s)+sS(u-s)
          =A(u)+M\sum_{a=0}^{s-1}\psi_u(a).
 \tag{BR.14}
\]
Subtract and divide by $M$ to obtain, for every $t\ge0$,
\[
 g(u)-g(u-ft)=tc-
            \sum_{\substack{0\le a<ft\\a\equiv u\ (N)}}w(a).
 \tag{BR.15}
\]
Choose $t$ with $ft=u$ modulo $N$.  Since $g(0)=0$, this formula
shows that $g(u)$ belongs to $B$.  It is only a multiple of $c$
minus a sum of the $w(a)$; division by $f$ has disappeared.
Also $c\in B$ by its definition.  When $ft\le N$, the final sum
in (BR.15) has exactly the possible term described in (BR.10).

Finally, grouping the left side of (BR.11) by residues gives
$\sum_u\xi(u)A(u)$.
The constant $A(0)$ contributes zero because $\xi$ has sum zero;
this proves (BR.11).
For the character assertion, the product character is odd and
therefore nontrivial; (BR.3)--(BR.4) identify its Bernoulli number
with the left side divided by $fN$.  This proves the asserted
identity at the original product level.
\end{proof}

\begin{proposition}[BR3: the prescribed Bernoulli unit]
\label{prop:br-lemma-x}
For every $\operatorname{XData}(\ell,d,\varphi,\xi)$ with values in
the specified $\overline{\mathbb Q}_\ell$,
\[
             |B_{1,\operatorname{prodX}(\varphi,\xi)}|_\ell=1.
             \tag{BR.16}
\]
The bound in (BR.5) is used with no enlargement or subsequent choice
of $n$.  In particular this proves the entire source
\texttt{Lemma\allowbreak X} interface.
\end{proposition}

\begin{proof}
Put $Q=\ell^d$ and $e=v_2(Q-1)$.  The hypotheses give $e\ge2$
and $\ell\ne2$.  Write $n=m+e$.  Since
$f\,2^{2e}\le2^n$ and $f\ge1$, one has $m\ge e$ and
\[
                   s=f2^e,\qquad 1<s\le2^m.
                   \tag{BR.17}
\]
Let $\mathcal A=\{x:|x|_\ell\le1\}$ in
$\overline{\mathbb Q}_\ell$, and let
$\mathfrak m=\{x:|x|_\ell<1\}$.  The ultrametric inequality makes
$\mathcal A$ a ring and $\mathfrak m$ an ideal.  Each element of
norm one has its inverse in $\mathcal A$, so $\mathfrak m$ is
maximal.  The field $\kappa=\mathcal A/\mathfrak m$ has
characteristic $\ell$, since $|\ell|_\ell<1$.
This auxiliary residue field is used only to test (BR.16).
It does not replace the given field $k$ in BR4--BR7.

Every value of $\varphi$ is either zero or a root of unity of
order dividing $Q-1$.  Thus it is integral and its residue is
fixed by $x\mapsto x^Q$.  Define the subring
\[
 B=\{x\in\mathcal A:\overline{x}^{\,Q}=\overline{x}\}.
 \tag{BR.18}
\]
It is a subring because the $Q$th-power map on $\kappa$ is an
iterated Frobenius.  The element $(2^n)^{-1}$ belongs to $B$:
it has norm one and its residue is an invertible prime-field element.
Apply BR2 with $N=2^n$ and $w(a)=\varphi(a\bmod f)$.
It supplies $g(u),c\in B$ and
\[
 B_{1,\operatorname{prodX}(\varphi,\xi)}
                         =\sum_{u\bmod 2^n}\xi(u)g(u).
 \tag{BR.19}
\]
The character $\xi$ is nontrivial, as follows also from its
primitivity and $n\ge4$.

For completeness, primitivity gives the exact half-period value
\[
                         \xi(1+2^{n-1})=-1.
                         \tag{BR.20}
\]
Indeed $h=1+2^{n-1}$ has order two.  The kernel of
$(\mathbb Z/2^n)^\times\to(\mathbb Z/2^{n-1})^\times$
is exactly $\{1,h\}$, as is seen from the two representatives
congruent to $1$ modulo $2^{n-1}$.
If $\xi(h)=1$, the character factors through level $2^{n-1}$,
contradicting primitivity.  In a characteristic-zero field its
value is therefore $-1$.

All values of $\xi$ are integral, so (BR.19) is integral.
Write $\eta=\overline{\xi}$, $b(u)=\overline{g(u)}$,
$w_0(a)=\overline{\varphi(a)}$, and $C=2^e\overline c$.
Then $b(u)^Q=b(u)$ and (BR.10), with $t=2^e$, says
\[
 b(u)-b(u-s)
        =C-\mathbf1_{\langle u\rangle_{2^n}<s}
                              w_0(\langle u\rangle_{2^n}).
 \tag{BR.21}
\]
Also $w_0(1)=1$ and
$w_0(s-1)\eta(-1)=-1$, because $s-1=-1$ modulo $f$.
By (BR.20), $\eta(h)=-1$ in the odd characteristic field $\kappa$.

Suppose the residue of (BR.19) were zero.  Put
\[
 T=\{u\bmod2^n:\eta(u)^{2^e}=1\},\qquad
 \pi:\mathbb Z/2^n\longrightarrow\mathbb Z/2^m.
 \tag{BR.22}
\]
For every unit $y$, changing variables in the zero sum gives
$\sum_u\eta(u)b(uy)=0$: multiplication by $\eta(y)\ne0$
recovers the original sum.  BR1 therefore implies
\[
                         \sum_{u\in T}\eta(u)b(uy)=0.
                         \tag{BR.23}
\]
BR1 also gives $\pi(u)=1$ or $-1$ for $u\in T$.
Define $\epsilon(u)=1$ in the first case and $-1$ in the second.
Since $2^e\mid s$, these two cases give respectively
$u(1-s)=u-s$ and $u(1-s)=u+s$ modulo $2^n$.
The element $1-s$ is a unit.

Use (BR.21) to subtract the two instances $y=1$ and $y=1-s$
of (BR.23).  For each $u\in T$ the resulting summand is
\[
 \eta(u)b(u)-\eta(u)b(u(1-s))
     =\epsilon(u)\eta(u)C-\mathbf1_{u=1}-\mathbf1_{u=-1}.
 \tag{BR.24}
\]
Here are the endpoint checks.  In the case $\pi(u)=1$,
a representative less than $s\le2^m$ must equal $1$,
giving $-\eta(1)w_0(1)=-1$.
In the case $\pi(u)=-1$, apply (BR.21) at $u+s$ and change
the sign.  A representative of $u+s$ less than $s$ must
equal $s-1$, hence $u=-1$; its contribution is
$\eta(-1)w_0(s-1)=-1$.
The two points are distinct since $m\ge2$.

Multiplication by $h=1+2^{n-1}$ permutes $T$, fixes $\pi(u)$,
and multiplies $\eta(u)$ by $-1$.  The first two facts use
$n-1\ge m$ and the evenness of $2^e$.  Hence
\[
                         \sum_{u\in T}\epsilon(u)\eta(u)=0,
                         \tag{BR.25}
\]
since $2$ is invertible in $\kappa$.
Both $1$ and $-1$ belong to $T$.  Summing (BR.24), its left
side is zero by (BR.23), whereas its right side is $-2$ by
(BR.25), a contradiction.
The integral number (BR.19) consequently has nonzero residue,
which is exactly norm one.  This proves (BR.16) with all of
the source's prescribed data.
\end{proof}

\subsection{Unramified homomorphisms and the finite Jacobi relation}

For $N>2$, an odd positive integer $a$ prime to $N$, and
$b\in(\mathbb Z/N)^\times$, set
\[
 \begin{gathered}
 y_b=\langle a^{-1}b\rangle_N,\qquad
 j_{N,a}(b)=\frac{a y_b-\langle b\rangle_N}{N},\\
 E_{N,a}(b)=\frac{a-1}{2}-j_{N,a}(b).
 \end{gathered}
 \tag{BR.26}
\]
These are integers and are exactly the source's \texttt{y\allowbreak Cl},
\texttt{jv}, and \texttt{Eone}, evaluated on units.

\begin{proposition}[BR4: the original tame-index condition and Jacobi annihilator]
\label{prop:br-unramified}
Let $\ell$ be odd, let $k$ be any field of characteristic $\ell$,
and let $z$ satisfy (BR.1)--(BR.2).
Then $z$ vanishes on $I_A\cap\ker\gamma$ for every $A$.

Let $F=\mathbb Q(\mu_N)\subset\Omega$ for $N>2$, put
$U=\operatorname{Gal}(\Omega/F)$, and suppose $\gamma|_U=1$.
Choose lifts $\widetilde\sigma_b\in G$ of
$\sigma_b(\zeta_N)=\zeta_N^b$.
For every odd positive $a$ prime to $N$ and every $u\in U$,
\[
 \sum_{b\in(\mathbb Z/N)^\times}
 E_{N,a}(b)\,
 z(\widetilde\sigma_b^{-1}u\widetilde\sigma_b)=0.
 \tag{BR.27}
\]
The statement uses the original $\Omega$ and $k$ and does not require
$k$ finite, a surjective residue map, or a common choice of $H_A$.
\end{proposition}

\begin{proof}
On $K_A=I_A\cap\ker\gamma$, equation (BR.1) makes $z$ an
additive homomorphism.  Its kernel contains $K_A\cap H_A$.
Thus the finite index of that kernel divides
$[K_A:K_A\cap H_A]$, which is prime to $\ell$.
Its image is a finite additive subgroup of $k$ and therefore an
elementary abelian $\ell$-group.  Its order is both a power of
$\ell$ and prime to $\ell$, so it is one.  This proves the first
assertion, without assuming that $H_A$ is contained in inertia.

The restriction $\lambda=z|_U$ is now an additive, locally constant
homomorphism, zero on every $I_A\cap U$.  The group $U$ is compact,
so the locally constant function $\lambda$ has finite image.
Its kernel is open and normal in $U$, and its fixed field is an
actual finite abelian extension of $F$ inside the original $\Omega$.
For every place of $F$, inertia in $G_F$ is exactly
$I_A\cap U$ under this inclusion: both groups fix the chosen
valuation and act trivially on its residue field.
Thus this finite extension is unramified at all finite places.
The field $F$ is totally complex because $N>2$; there is no real
place condition to add.  The finite Hilbert class field theorem
therefore gives an additive map
\[
       \Lambda:\operatorname{Cl}(F)\longrightarrow k,\qquad
       \lambda=\Lambda\circ\operatorname{Art}_F.
       \tag{BR.28}
\]
Only its finite image is used.  The range remains a subgroup of $k$.

We derive the required ideal relation from OS1.
For $1\le r<N$, let
\[
 \theta_{N,r}(t)=
    \left\lfloor\frac{r\langle t\rangle_N}{N}\right\rfloor
    -\left\lfloor\frac{r(N-\langle t\rangle_N)}{N}\right\rfloor.
 \tag{BR.29}
\]
For every nonzero ideal $\mathfrak a$, OS1 constructs an element
whose principal ideal is
\[
                       \prod_t
                 \sigma_t^{-1}\mathfrak a^{\,\theta_{N,r}(t)}.
\]
Consequently the group-ring element
$\sum_t\theta_{N,r}(t)\sigma_t^{-1}$
annihilates the ideal class group.
This consequence uses only the ideal assertion in OS1.

There is no restriction $a<N$ in (BR.27).
Write $a=qN+r$ with $1\le r<N$; this is possible since $(a,N)=1$.
The same floor expression at the positive integer $a$ satisfies
\[
 \theta_{N,a}(t)=\theta_{N,r}(t)+q(2\langle t\rangle_N-N),
 \qquad
 \theta_{N,N-1}(t)=2\langle t\rangle_N-N.
 \tag{BR.30}
\]
Both follow by separating the integer part $q\langle t\rangle_N$.
The relation from OS1 for $r$ and for $N-1$ therefore proves
that $\Theta_a=\sum_t\theta_{N,a}(t)\sigma_t^{-1}$
annihilates $\operatorname{Cl}(F)$ for this arbitrary $a$.

Because $a\langle t\rangle_N$ is not divisible by $N$,
\[
 \theta_{N,a}(t)=
       2\left\lfloor\frac{a\langle t\rangle_N}{N}\right\rfloor
                  -(a-1).
 \tag{BR.31}
\]
Putting $b=at$ modulo $N$ in (BR.26) gives the exact group-ring
identity
\[
 2\sum_b E_{N,a}(b)\sigma_b^{-1}
                       =-\sigma_a^{-1}\Theta_a.
 \tag{BR.32}
\]
Thus the left side annihilates every class.
After applying the additive map $\Lambda$ to it, division by
$2$ is legitimate in the original odd-characteristic field $k$.

Finally Artin is compatible with conjugation:
\[
 \operatorname{Art}_F(\widetilde\sigma_b^{-1}u
                         \widetilde\sigma_b)
                   =\sigma_b^{-1}\operatorname{Art}_F(u).
 \tag{BR.33}
\]
One can check the convention on a Frobenius at a finite prime:
conjugating the automorphism conjugates that prime by
$\sigma_b^{-1}$, and Artin sends its Frobenius to its ideal class.
The resulting compatibility on all elements follows from the
Artin map, or from the density of Frobenius in each finite abelian
quotient.  This is finite reciprocity on the same field $F$;
the inverse agrees with the ideal action in OS1.
Apply $\Lambda$ to (BR.32) at
$\operatorname{Art}_F(u)$ and use (BR.28)--(BR.33).
This proves (BR.27).
\end{proof}

\subsection{The exact Bernoulli scalar on the original coefficient ring}

\begin{proposition}[BR5: the nonzero residual annihilator]
\label{prop:br-scalar}
Let $O=c.O\subset\overline{\mathbb Q}_\ell$ be the source's finite
discrete valuation ring over $\mathbb Z_\ell$, let $k$ be any field
of characteristic $\ell$, and let $\operatorname{red}:O\to k$
be the specified unital ring map.
Let $\chi$ be a nontrivial Dirichlet character of level $N>0$
and order prime to $\ell$, where $\ell$ is odd.
Write $\gamma=\operatorname{redChar}(\chi)$.
Suppose $\gamma$ differs somewhere from the reduction
$\overline{\chi_\ell}$ of the cyclotomic character, and suppose
\[
                         |B_{1,\chi^{-1}}|_\ell=1.
                         \tag{BR.34}
\]
Then $N>2$, $\gamma$ is nontrivial, and there is an odd positive
integer $a$ prime to $N$ such that
\[
 \operatorname{red}\left(
      \sum_{b\in(\mathbb Z/N)^\times}E_{N,a}(b)\chi^{-1}(b)
                    \right)\ne0.
 \tag{BR.35}
\]
The Bernoulli number in (BR.34) is evaluated through the original
inclusion $O\subset\overline{\mathbb Q}_\ell$.
\end{proposition}

\begin{proof}
We first record precisely what the given residue map implies.
The structure map $\mathbb Z_\ell\to O$ is injective, because $O$
is an actual subalgebra of $\overline{\mathbb Q}_\ell$.
The prime ideal $\ker(\operatorname{red})$ contains the nonzero
element $\ell$.  In a discrete valuation ring the only nonzero
prime ideal is the maximal ideal.  Hence this kernel is exactly
the maximal ideal, even though the map need not be surjective.
In particular $\operatorname{red}(x)\ne0$ if and only if $x\in O^\times$.

For use with the norm, all elements of $O$ have norm at most one:
they are integral over $\mathbb Z_\ell$ by module finiteness,
and a monic integral equation contradicts a norm greater than one
by the ultrametric inequality.  A unit and its inverse are integral,
so every unit has norm one.  Write
$\ell=u\varpi^h$ with $u$ a unit and $h>0$, which is possible
because $\ell$ is a nonzero nonunit.
Then $|\varpi|_\ell=\ell^{-1/h}<1$.
The discrete valuation factorization of every nonzero $x\in O$
shows conversely that $|x|_\ell=1$ implies $x$ is a unit.
This proves the norm-to-residue test on the actual ring $O$.

Reduction is injective on roots of unity of order prime to $\ell$.
Indeed, if $x^r=1$, $x\ne1$, and $\ell\nmid r$, then
$1+x+\cdots+x^{r-1}=0$.  Were $\operatorname{red}(x)=1$,
reduction would give $r=0$ in $k$, a contradiction.
Apply this with $r=\operatorname{ord}(\chi)$ to a unit residue
where $\chi$ is not one.  The cyclotomic restriction map
$G\to(\mathbb Z/N)^\times$ is surjective, so $\gamma\ne1$.
Also $N>2$, since the unit groups at levels $1$ and $2$ are trivial.

We choose $a$ so that
\[
                    \operatorname{red}(\chi(a))\ne \overline a.
                    \tag{BR.36}
\]
If $\ell\nmid N$, take $a=\ell$.
It is odd and prime to $N$; the left side is a nonzero character
value whereas the right side is zero.
If $\ell\mid N$, choose the given $g\in G$ where
$\gamma(g)\ne\overline{\chi_\ell}(g)$ and put
$b=\operatorname{cls}_N(g)$.
Choose a positive odd representative $a$ of $b$ modulo $N$.
This is possible: for even $N$ every unit representative is odd;
for odd $N$, adding $N$ changes parity.
The action of $g$ on $\mu_\ell\subset\mu_N$ is the $a$th-power
map.  Thus $\chi_\ell(g)-a\in\ell\mathbb Z_\ell$, and reduction
in the given $k$ gives
$\overline{\chi_\ell}(g)=\overline a$.
The chosen inequality is now (BR.36).

For any nontrivial character $\psi$ with values in a
characteristic-zero field, direct summation in (BR.26) gives
\[
 \sum_b E_{N,a}(b)\psi(b)
                    =(1-a\psi(a))B_{1,\psi}.
                    \tag{BR.37}
\]
Here is the calculation, including the inverse in $y_b$.
The constant term vanishes because $\sum_b\psi(b)=0$.
Multiplication by $a$ permutes the units, so
\[
 \sum_b y_b\psi(b)
      =\psi(a)\sum_t\langle t\rangle_N\psi(t).
 \tag{BR.38}
\]
Multiply the left side of (BR.37) by $N$ and use
$Nj_{N,a}(b)=a y_b-\langle b\rangle_N$.
Equations (BR.38) and (BR.4) give
$N(1-a\psi(a))B_{1,\psi}$, proving (BR.37) since $N\ne0$
in characteristic zero.

Apply (BR.37) to $\psi=\chi^{-1}$ in the specified
$\overline{\mathbb Q}_\ell$.
By (BR.36) the element
$v=1-a\chi^{-1}(a)\in O$ has nonzero reduction and hence is a unit.
Let $S=\sum_b E_{N,a}(b)\chi^{-1}(b)\in O$.
The element $v^{-1}S\in O$ maps to $B_{1,\chi^{-1}}$.
By (BR.34) and the norm criterion above it is a unit.
Thus $S$ is a unit, and its reduction is nonzero.
This proves (BR.35) without making an integrality or nonvanishing
assumption on a different character or coefficient ring.
\end{proof}

\begin{proposition}[BR6: the full residual unramified cocycle theorem]
\label{prop:br-residual}
Under exactly the hypotheses of BR5, every cocycle $z:G\to k$
for $\gamma=\operatorname{redChar}(\chi)$ satisfying
$\operatorname{UnrEv}(\ell,\gamma,z)$ has the form
\[
                     z(g)=(\gamma(g)-1)m\qquad(g\in G)
                     \tag{BR.39}
\]
for some $m$ in the original field $k$.
This is the literal source interface
\texttt{Unramified\allowbreak Cocycle\allowbreak Is\allowbreak Coboundary\allowbreak Of\allowbreak Bernoulli\allowbreak Unit}.
\end{proposition}

\begin{proof}
Put $F=\mathbb Q(\mu_N)\subset\Omega$ and $U=G_F$.
The character $\gamma$ is trivial on $U$.
For $u\in U$ and any lift $\widetilde\sigma_b$ as in BR4,
the cocycle identity gives
\[
 z(\widetilde\sigma_b^{-1}u\widetilde\sigma_b)
                    =\gamma(\widetilde\sigma_b)^{-1}z(u)
                    =\operatorname{red}(\chi^{-1}(b))z(u).
 \tag{BR.40}
\]
For example, compare $z(u\widetilde\sigma_b)$ with
$z(\widetilde\sigma_b
        (\widetilde\sigma_b^{-1}u\widetilde\sigma_b))$;
the terms $z(\widetilde\sigma_b)$ cancel because $\gamma(u)=1$.
Use in BR4 the integer $a$ supplied by BR5.
Equations (BR.27) and (BR.40) then say
\[
      \operatorname{red}(S)\,z(u)=0,\qquad u\in U.
      \tag{BR.41}
\]
The scalar $\operatorname{red}(S)$ is nonzero, so $z|_U=0$.

Choose $h\in G$ with $\gamma(h)\ne1$, which exists by BR5.
The quotient $G/U=\operatorname{Gal}(F/\mathbb Q)$ is abelian.
Hence $gh$ and $hg$ differ by an element of $U$, on which
both $\gamma-1$ and $z$ vanish.  Thus $z(gh)=z(hg)$.
Using (BR.1) on both products gives
\[
                (\gamma(h)-1)z(g)
                    =(\gamma(g)-1)z(h).
                \tag{BR.42}
\]
Set $m=z(h)/(\gamma(h)-1)\in k$.
Equation (BR.42) proves (BR.39) for all $g$.
In particular no division by the order of
$\operatorname{Gal}(F/\mathbb Q)$ is required; that order may
be divisible by $\ell$.  All local conditions, lifts and the
final coboundary remain on the original $\Omega$ and $k$.
\end{proof}

\subsection{Both characters and the literal CorX endpoint}

\begin{proposition}[BR7: the complete original CorX supplier]
\label{prop:br-cor-x}
Let $\Omega,\ell,c,O,k,\operatorname{red},d,f,n,\varphi,\xi$
be exactly as in the source's \texttt{Cor\allowbreak X}: $O=c.O$ is the
specified finite DVR subalgebra, $k$ is any field of characteristic
$\ell$, $\operatorname{red}:O\to k$ is any unital ring map,
and $\operatorname{XData}(\ell,d,\varphi,\xi)$ holds.
Put $\chi=\operatorname{prodX}(\varphi,\xi)$ and let
\[
        \gamma=\operatorname{redChar}(\chi)
          \quad\text{or}\quad
        \gamma=\operatorname{redChar}(\chi)^{-1}.
        \tag{BR.43}
\]
For every locally constant $\gamma$-cocycle $z$ satisfying the
exact condition (BR.2), there is $m\in k$ with
$z(g)=(\gamma(g)-1)m$ for every $g\in G$.
There are no remaining bespoke assumptions in this supplier.
\end{proposition}

\begin{proof}
Write $e=v_2(\ell^d-1)$.  As in BR3,
$\ell$ is odd and $n\ge2e\ge4$.
The order of $\xi$ divides a power of $2$, since it is a
character of the finite group $(\mathbb Z/2^n)^\times$.
The order of $\varphi$ divides $\ell^d-1$.
Therefore
\[
                \chi^{\,(\ell^d-1)2^n}=1,
                \qquad \ell\nmid\operatorname{ord}(\chi).
                \tag{BR.44}
\]
The same statements hold for $\chi^{-1}$.

We verify explicitly both nontriviality and the cyclotomic exclusion
required in BR5.  Put $b=1+2^{n-1}$ modulo $2^n$.
The primitive character calculation (BR.20) gives $\xi(b)=-1$
also in $O$, since its inclusion into
$\overline{\mathbb Q}_\ell$ is injective.
As $\ell f$ is odd, the Chinese remainder theorem gives a unit
$a$ modulo $\ell f2^n$ satisfying
\[
                   a\equiv1\pmod{\ell f},
                   \qquad a\equiv b\pmod{2^n}.
                   \tag{BR.45}
\]
Choose in the original $G$ a lift $g$ of this cyclotomic
automorphism.  It fixes $\mu_\ell$, so
$\overline{\chi_\ell}(g)=1$ in the original $k$.
On the other hand (BR.3) gives
\[
                 \chi(g)=-1,\qquad \chi^{-1}(g)=-1.
                 \tag{BR.46}
\]
Reduction preserves $-1\ne1$ because $\ell$ is odd.
Thus both residual characters are nontrivial and each differs
from $\overline{\chi_\ell}$ at this very same $g$.
In particular $\chi$ and $\chi^{-1}$ are nontrivial before reduction.

The data (BR.5) are preserved under simultaneous inversion:
the conductors and the character orders are unchanged,
$\varphi^{-1}(-1)=\varphi(-1)$ and
$\xi^{-1}(-1)=\xi(-1)$, and the numerical data are identical.
They are also preserved under the actual inclusion
$O\hookrightarrow\overline{\mathbb Q}_\ell$.
For conductor preservation, a character factors through a lower
level exactly when its unit-group kernel contains the corresponding
reduction kernel; an injective coefficient map does not change that
kernel.  Multiplication, inversion and the level changes in
(BR.3) commute with this inclusion.  Therefore BR3 applied to
both pairs gives
\[
                         |B_{1,\chi}|_\ell
                         =|B_{1,\chi^{-1}}|_\ell=1.
                         \tag{BR.47}
\]

For the first alternative of (BR.43), apply BR6 to $\chi$,
using the second equality in (BR.47).  For the second alternative,
apply BR6 to $\chi^{-1}$, using the first equality.
The identity
$\operatorname{redChar}(\chi^{-1})
 =\operatorname{redChar}(\chi)^{-1}$
follows by evaluating the unit-valued character and then applying
the same ring map $\operatorname{red}$.
Thus the input cocycle and its condition (BR.2) are literally
the ones required by BR6 in either case.
Its output lies in the originally quantified $k$.
This proves the full \texttt{Cor\allowbreak X} statement.
\end{proof}

\paragraph{Closure and the downstream boundary.}
The complete chain established here is the finite Frobenius projection
and integral partial sum, followed by the prescribed Bernoulli unit;
the original tame-index condition, finite Hilbert class-field descent
and OS1 Jacobi relation; and finally the nonzero scalar, cocycle
vanishing and both-character \texttt{Cor\allowbreak X} endpoint.
BR4 uses only the odd-characteristic specialization of the source's
general \texttt{hom\allowbreak Kill}; it does not assert the unrestricted
integral \texttt{cl\allowbreak Kill} or arbitrary-coefficient \texttt{hom\allowbreak Kill}
interfaces.  BR7 supplies the exact input to the paired
\texttt{Count\allowbreak II}/\texttt{Count\allowbreak III} argument.  Automorphy, rigidity,
congruence-number bounds and the other hypotheses of that argument
are separate obligations.


% Version 0.26.  Complete weighted intersection suppliers and the paired-count cut.
\section{Weighted intersection lengths and the symplectic paired-count frontier}
\label{sec:symplectic-count-closure}

\status{SP1 proves the literal \texttt{Lemma\allowbreak M'}.  SP2--SP6 prove the
literal \texttt{Lemma\allowbreak M}, including its finite orbit cover and its factor
$2$, by a three-variable Rees calculation.  These are complete algebra
suppliers for both the rank-three and rank-four routes.  The later
statements separate the finite trace-product arithmetic and coefficient
return from the still unproved construction and identification of the
symplectic deformation model.  Neither generic automorphy nor a package
named \texttt{Hypotheses\allowbreak H'} is accepted as a classical citation leaf.}

\subsection{The cotangent lattice and the additive-injection lemma}

Let $O$ be a discrete valuation ring, with uniformizer $\pi$, and let
$x:A\to O$ be an $O$-algebra augmentation.  Put
\[
 \mathfrak a=\ker x,\qquad
 \mathfrak a^{(2)}=\{f\in A:\ cf\in\mathfrak a^2
                         \text{ for some }0\ne c\in O\}.
 \tag{SP.1}
\]
For a grading $A=\bigoplus_{j\in X}A_j$, define
\[
 C_j=\operatorname{im}(A_j\cap\mathfrak a\longrightarrow
                          A/\mathfrak a^{(2)}),\qquad
 C_j(F)=C_j/\langle\overline F_i\rangle_O
 \quad(F_i\in A_j\cap\mathfrak a).
 \tag{SP.2}
\]
Thus $C_j$ and $\ell_O C_j(F)$ are exactly \texttt{cot\allowbreak Lat} and
\texttt{cot\allowbreak Len}.  Length is allowed to be $\infty$.
The source's \texttt{Spans\allowbreak Cot} means that the $F_i$ span $C_j$ after
inverting $\pi$; it does not mean that they span the integral lattice.

\begin{proposition}[SP1: the literal \texttt{Lemma\allowbreak M'}]
\label{claim:SP1}
Suppose the residue field of $O$ is finite, $X$ is any abelian grading
group, $x(A_j)=0$, and $0\ne c\in O$.  Let $M$ be an $O$-module and
let
\[
 D:M\longrightarrow\operatorname{Hom}_O(A,O/(c))
 \tag{SP.3}
\]
be an injective additive map.  Suppose $D\mu(F_i)=0$ for every $\mu,i$,
and every $D\mu$ is the reduction of an $O$-valued derivation at $x$
which vanishes on all weights other than $j$.  Then
\[
                    \ell_O M\le\ell_O C_j(F).
 \tag{SP.4}
\]
No $O$-linearity of the assignment $\mu\mapsto D\mu$ is assumed.
\end{proposition}

\begin{proof}
An $O$-valued derivation $\delta$ at $x$ kills $\mathfrak a^2$ by the
Leibniz rule.  If $cf\in\mathfrak a^2$ and $c\ne0$, then
$c\delta(f)=0$ in the domain $O$, so it also kills
$\mathfrak a^{(2)}$.  Restricting the reductions to $A_j$ therefore
gives an additive map
\[
 M\longrightarrow\operatorname{Hom}_O(C_j(F),O/(c)).
 \tag{SP.5}
\]
It is injective: equality on $A_j$ implies equality on every homogeneous
component, because both maps vanish on the other components; every
element of $A$ is a finite sum of homogeneous elements.  Injectivity
of (SP.3) then identifies the two elements of $M$.

If $C_j(F)$ has infinite length, (SP.4) is immediate.  Otherwise the
structure theorem over a DVR gives
$C_j(F)\simeq\bigoplus_iO/(\pi^{a_i})$.  Writing $b=v_O(c)$, the
target of (SP.5) has length $\sum_i\min(a_i,b)\le\sum_i a_i$ and
cardinality $q^{\sum_i\min(a_i,b)}$, where $q=|O/\pi|>1$.
The additive injection makes $M$ finite.  Every finite $O$-module has
cardinality $q^{\ell_O M}$, so comparison of cardinalities proves
(SP.4).  This also explains the finite-residue-field hypothesis in the
source: an additive injection alone is what is available.
\end{proof}

\subsection{The exact weighted intersection problem}

For SP2--SP6, $O$ is a complete DVR of characteristic zero,
$E=\operatorname{Frac}(O)$, and $K$ is an algebraically closed field
with a faithful $O$-algebra structure.  A valuation
$v:K^\times\to\mathbb R$ is normalized by $v(\pi)=1$ and
$v(c)=v_O(c)$ for $c\in O\setminus\{0\}$; these are exactly the
nonzero-input clauses of \texttt{Is\allowbreak NVal}.  Let
$A=\bigoplus_{j\in\mathbb Z}A_j$, $B=A_0$, where $B$ is complete
Noetherian local, $O\to B$ is local, and the induced residue-field map
is an isomorphism.  The last three conditions are \texttt{Is\allowbreak CO O B}.
Assume that $A$ is generated over $B$ by finitely many elements of
weights $-2,-1,1,2$, and that $x(A_j)=0$ for $j\ne0$.
These are precisely the remaining clauses of \texttt{Is\allowbreak WRing}.

Fix integers $r\ge1$, $s,r',s',m\ge0$ and equations
\[
\begin{array}{c|ccccc}
 \text{equations}&F_0,\ldots,F_r&G_1,\ldots,G_s&E_1,\ldots,E_{r'}
       &E'_1,\ldots,E'_{s'}&Z_1,\ldots,Z_m\\
 \text{weights}&-1&-2&1&2&0 .
\end{array}
 \tag{SP.6}
\]
All $x(Z_i)$ vanish, and $I=(F,G,E,E',Z)$.  Write $S(K)$ for the
$O$-algebra maps $\xi:A/I\to K$.  Such a point is \emph{two-sided}
if it is nonzero on some negative and on some positive weight.
The action is $(t\xi)(f_j)=t^j\xi(f_j)$.

We use the following exact five hypotheses, with their source names
recorded to make the endpoint checkable.
\begin{enumerate}
\item[(M1)] There is a weight-preserving formal chart
$\phi:A\to E[[z_i]]$ with constant term $x$, whose linear coordinate
functionals are independently realized by elements of $A$.  Their
common kernel on $\ker x$ is $\mathfrak a^{(2)}$.  The numbers of
coordinates of weights $-1,-2,1,2,0$ are respectively
$r,s,r'+1,s'+1,m$, and there are no other weights.  The equations
$F,G,Z$ span $C_{-1},C_{-2},C_0$ rationally.
\item[(M2)] The image of $B$ in $A/I$ is a finite $O$-module.
\item[(M3)] Over any fixed point $x_1\ne x$ of $S(K)$, the fibre of
restriction to $B$ cannot contain both a nontrivial point supported
only in nonnegative weights and a nontrivial point supported only in
nonpositive weights.
\item[(M4)] At every two-sided point of $S(K)$, any two $K$-valued
point derivations annihilating $I$ are linearly dependent.
\item[(M5)] The chosen $a\in B$ has $x(a)=0$ and $\xi(a)\ne0$ at
every two-sided point of $S(K)$.
\end{enumerate}
The statement called \texttt{Lemma\allowbreak M} is that there is a finite set
$T$ consisting entirely of two-sided points of $S(K)$ and covering
all their orbits, that
$\ell_OC_{-1}(F)<\infty$, and, for \emph{every} such cover,
\[
                    \ell_OC_{-1}(F)\le2\sum_{\xi\in T}v(\xi(a)).
 \tag{SP.7}
\]
A cover may contain more than one representative of an orbit.

We record two elementary finiteness observations used below.  Each
$A_j$ is a finite $B$-module: in a monomial in the chosen generators,
remove submonomials of total weight zero; a remaining monomial of
fixed weight has bounded length, since the only nonzero generator
weights are $\pm1,\pm2$.  The same grouping, separating a pair of
weight-one factors or one weight-two factor, gives
\[
                       A_{2n}=A_2^n,\qquad A_{-2n}=A_{-2}^n
                       \quad(n\ge0),
 \tag{SP.8}
\]
where products are products of $B$-submodules of $A$.
For example, a mixed monomial is reduced by removing opposite pairs
or triples $(2,-1,-1)$ and $(-2,1,1)$; an even one-sided remainder
partitions into blocks of weight $2$ or $-2$.
Also $B$ acts on every $C_j$ through $x:B\to O$, since
$(b-x(b))f\in\mathfrak a^2$ for $f\in\ker x$.  Thus $C_{-1}$ is
finite over $O$, torsion free by (SP.1), and has rank $r$ by the
injective and surjective linear-term clauses of the chart.

\begin{lemma}[SP2: the multigraded Euler calculation]
\label{claim:SP2}
Let $R$ be a Noetherian $\mathbb N^k$-graded $B$-algebra with
$R_0=B$, generated by finite $B$-modules in degrees
$\mathbf e_1,\ldots,\mathbf e_k$, and put
$R_{++}=\prod_i(R_{\mathbf e_i})R$.  Let $M$ be a finite graded
$R$-module, and let $\mathbf t=(t_1,\ldots,t_h)$ be homogeneous
sections of nonnegative multidegrees.  Suppose
\[
        (\mathfrak m_BR_{++})^N
                  \subset\operatorname{Ann}_R(M)+(\mathbf t)
                  \quad\text{for some }N.
 \tag{SP.9}
\]
For all sufficiently large multidegrees $d$, the Koszul homology
pieces have finite $B$-length.  Their alternating sum
$\chi(\mathbf t;M)_d$ is additive in graded exact sequences,
unchanged by permutation of the sections, and nonnegative.  It is
monotone under a graded quotient, and
\[
 \chi(\mathbf t,\mathbf u;M)_d
                 \le\chi(\mathbf u;M/(\mathbf t)M)_d.
 \tag{SP.10}
\]
Suppose in addition that $\pi$ is injective on $M$, that $B$ acts on
$M/\pi M$ through its residue field, and that $y_1,\ldots,y_n$ all
have degree $\alpha$.  If the section list
$(\pi^{b_i}y_i)_i,\mathbf u$ satisfies (SP.9), then
\[
 \chi((\pi^{b_i}y_i)_i,\mathbf u;M)_d
 \ \ge\ \Big(\sum_i b_i\Big)
       \Delta_\alpha^{n-1}\prod_{u\in\mathbf u}\Delta_{\deg u}
                  \ell_B(M/\pi M)_d .
 \tag{SP.11}
\]
Here $\Delta_\beta h(d)=h(d)-h(d-\beta)$ and every assertion is
eventual in \emph{all} coordinates of $d$.
\end{lemma}

\begin{proof}
Every Koszul homology module is finite graded and is annihilated by
$\operatorname{Ann}(M)+(\mathbf t)$.  If its homogeneous generators
have degrees $d_i$, then, once every coordinate of $d$ exceeds those
of the $d_i$ by $N$, its degree-$d$ part is generated using
$R_{++}^N$.  Consequently $\mathfrak m_B^N$ kills that part.  It is
finite over $B$, hence has finite length.  Exact-sequence additivity
follows from the homology long exact sequence.  Permuting a Koszul
basis proves invariance under permutations.

For completeness, positivity here does not invoke an unspecified
intersection-multiplicity theorem.  A finite graded module has a
finite filtration with factors $(R/\mathfrak p)(-b)$ for homogeneous
prime ideals $\mathfrak p$.  Indeed, choose a nonzero homogeneous
element with maximal annihilator; its annihilator is homogeneous
prime, and iterate in the quotient.  Homogeneous primality implies
primality by comparing highest terms in a translation-invariant
ordering of $\mathbb Z^k$; Noetherianity makes the iteration finite.
Additivity reduces positivity to a factor $R/\mathfrak p$.
If $R_{++}\subset\mathfrak p$, its large graded pieces vanish.
Otherwise there is a section outside $\mathfrak p$ of every
nonnegative multidegree: multiply suitable sections of degrees
$\mathbf e_i$ outside $\mathfrak p$.

Induct on the number of sections.  If the first section $t$ is
nonzero in $R/\mathfrak p$, it is a nonzerodivisor, and its Koszul
complex reduces to the quotient by $t$; apply the induction
hypothesis to the remaining list.  If $t=0$ there, its Euler
contribution is the difference between the Euler functions of
$R/\mathfrak p$ and its shift by $\deg t$.  Choose a section $s$ of
that degree outside $\mathfrak p$.  The exact sequence for
multiplication by $s$ identifies this difference with the Euler
function on $(R/\mathfrak p)/s$.  The remaining list is admissible
because $t$ was already zero.  Induction again gives nonnegativity.
The case of an empty list is simply a length.  This proves
positivity without an infinite-residue-field assumption.

Monotonicity follows from additivity and positivity on the kernel.
The recursion
\[
 \chi(t,\mathbf u;M)_d=
 \chi(\mathbf u;M/tM)_d-
 \chi(\mathbf u;M[t](-\deg t))_d
 \tag{SP.12}
\]
and positivity give (SP.10), one section at a time.

Finally the kernel--cokernel exact sequence for a product of two
endomorphisms, or the equivalent two mapping cones, gives
\[
 \chi(\pi^b y,\mathbf u;M)
       =b\,\chi(\pi,\mathbf u;M)+\chi(y,\mathbf u;M).
 \tag{SP.13}
\]
There is no grading shift in the first term because $\pi$ has degree
zero.  Injectivity of $\pi$ identifies that term with
$\chi(\mathbf u;M/\pi M)$; its graded terms have finite length, so
Euler--Poincar\'e identifies it with the indicated iterated difference
of their lengths.  Peel the factors $\pi^{b_i}$ successively,
permuting each $y_i$ to the end of the list after it is peeled.  The
degrees of the other $n-1$ sections remain $\alpha$.  Each contributes
$b_i$ times the same difference in (SP.11); the residual Euler
function is nonnegative.  Replacing a section by its divisor only
enlarges the section ideal, so admissibility holds at every step.
\end{proof}

\subsection{The Rees support argument and the two-sided fibres}

By Smith normal form, change the $F_i$ by an invertible $O$-linear
change of their $r+1$ generators, without changing their ideal or
their cotangent span, so that
\[
 F_i\equiv\pi^{c_i}y_i\pmod{\mathfrak a^{(2)}}\quad(0\le i<r),
 \qquad H:=F_r\in\mathfrak a^{(2)},\qquad
 \ell_OC_{-1}(F)=\sum_{i<r}c_i .
 \tag{SP.14}
\]
Here the classes of $y_i\in A_{-1}$ form an $O$-basis of $C_{-1}$.
The extra generator is obtained from the rank-one kernel of the map
$O^{r+1}\to\langle\overline F_i\rangle_O$; the target is free and
the map splits.  In particular this change really preserves the
integral span of the equations.

Introduce three finite $B$-submodules of $A$ and their multigraded
Rees algebra:
\[
 \begin{gathered}
 N_0=A_{-2},\quad N_1=A_2,\quad
 N_2=BH^2+A_{-2}^2A_2\subset A_{-2},\\
 R=\bigoplus_{d\in\mathbb N^3}L_dT^d\subset A[T_0,T_1,T_2],
 \qquad L_d=N_0^{d_0}N_1^{d_1}N_2^{d_2}.
 \end{gathered}
 \tag{SP.15}
\]
The section list is
\[
 \mathbf t=(F_i^2T_0)_{i<r},\ (G_iT_0)_i,\
                  (E_i^2T_1)_i,\ (E'_iT_1)_i,\ H^2T_2.
 \tag{SP.16}
\]
Put $D_0=(A_{-2}A_2)B\subset B$ and $I_0=I\cap B$.

\begin{lemma}[SP3: support and admissibility on this Rees algebra]
\label{claim:SP3}
Assume (M2), (M3), and rational spanning by $F$ and $G$.
Let $J'$ be any homogeneous ideal of $R$ containing $\mathbf t$
and the $Z_i$.  For all sufficiently large $d$,
$M_d=(R/J')_d$ is a finite $O$-module.  Every nonzero generic
localization of $M_d$ is at a prime
$\mathfrak p=\ker(\sigma:B\to K)$ with
\[
                     I_0\subset\mathfrak p,
                     \qquad D_0\not\subset\mathfrak p.
 \tag{SP.17}
\]
Such a $\sigma$ is the restriction of a two-sided point of $S(K)$.
If also (M5) holds, then for every homogeneous ideal $J\subset R$ containing
the $Z_i$, the list $(\mathbf t,a)$ satisfies (SP.9) on $R/J$.
\end{lemma}

\begin{proof}
We give the prime argument, including the use of (M3).  Let $P$ be a
prime of $R$ containing $\mathbf t,Z$ but not $R_{++}$.  Choose
$f\in N_0$, $g\in N_1$, $h\in N_2$ such that
$fT_0,gT_1,hT_2\notin P$.  The relation
\[
             (H^2T_0)(hT_2)=(hT_0)(H^2T_2)
 \tag{SP.18}
\]
shows that $H^2T_0\in P$, so all $F_i^2T_0$ belong to $P$.
Let $u=(fT_0)(gT_1)\notin P$.  If $cF_i$ has weight zero, then
\[
              (cF_i)^2u=(F_i^2T_0)(c^2fgT_1)\in P.
\]
For $cG_i$ use $(cG_i)u=(G_iT_0)(cfgT_1)$; the analogous identities
for $E_i^2T_1$ and $E'_iT_1$ interchange the two variables.  Terms
involving $Z_i$ already lie in $P$.  Taking weight-zero components
in an expression for $b\in I$ now shows $I_0\subset P\cap B$.

By (M2), $B/I_0$ is finite over $O$.  Its primes are either the
closed prime $\mathfrak m_B$ or horizontal primes whose residue
fraction fields are finite over $E$.  In the latter case choose an
embedding into $K$, obtaining $\sigma:B\to K$ with
$\ker\sigma=P\cap B$.  We show that $D_0\subset\ker\sigma$ is
impossible.

Here is an explicit construction of the two forbidden one-sided
points.  For $\epsilon\in\{-1,1\}$ use the Rees direction $T_0$
or $T_1$, respectively, and set
\[
 \iota_n(b)=\operatorname{pr}_{2\epsilon n}(b)T_\epsilon^n,
 \qquad
 J_\epsilon=\{b\in A:\ \iota_n(cb)\in P
                         \text{ for all }c\in A,n\ge0\}.
 \tag{SP.19}
\]
The notation $T_{-1}=T_0$, $T_1=T_1$ is used only in (SP.19).
By (SP.8), every displayed monomial lies in $R$.
The set $J_\epsilon$ is an ideal.  Choose $q\in A_{2\epsilon}$
with $qT_\epsilon\notin P$.  For homogeneous elements the identities
\[
 \iota_{n+n'}(bc)=\iota_n(b)\iota_{n'}(c),\qquad
 \iota_{2n}(b^2)=\iota_n(b)^2
 \tag{SP.20}
\]
give the following facts by primality: $b^2\in J_\epsilon$ implies
$b\in J_\epsilon$; $qb\in J_\epsilon$ implies
$b\in J_\epsilon$; and $bT_\epsilon\in P$ for
$b\in A_{2\epsilon}$ implies $b\in J_\epsilon$.
To verify the last assertion for a test factor of weight
$2\epsilon n-2\epsilon$, multiply it by $q$ and cancel
$qT_\epsilon$ in (SP.20).  Likewise
$P\cap B\subset J_\epsilon$ and
$bq^n\notin J_\epsilon$ if $b\notin P\cap B$.
Arbitrary test factors reduce to homogeneous ones by their finite
weight decomposition.

If $D_0\subset\ker\sigma$, then $qA_{-2\epsilon}$ lies in
$P\cap B$, so cancellation puts $A_{-2\epsilon}$ in
$J_\epsilon$.  Squaring and (SP.8) put every weight of the opposite
sign in that ideal.  The section equations in the chosen direction
put the equations of the same sign in it, using squares for weight
$\epsilon$; (SP.18) includes the last $F$.  Thus
$I\subset J_\epsilon$.
The ideal $J_\epsilon$ is disjoint from the multiplicative set
$\{bq^n:\sigma(b)\ne0\}$.  Put
$L=\operatorname{Frac}(B/\ker\sigma)$.
The algebra $((A/J_\epsilon)\otimes_B L)[1/q]$ is a nonzero
finite-type $L$-algebra: disjointness proves nonvanishing, and the
finite generation of $A$ over $B$ proves finite type.
Faithfully flat extension from $L$ to $K$ preserves both properties.
The Nullstellensatz applied to that algebra produces a point
$\xi_\epsilon$ over $\sigma$
annihilating $I$ and all weights of the opposite sign, with
$\xi_\epsilon(q)\ne0$.

There is also a fixed point over $\sigma$, namely
$\xi_0(b)=\sigma(\operatorname{pr}_0b)$.  It is multiplicative:
the square of a product of opposite nonzero weights is a product of
elements of $A_{-2}A_2$ by (SP.8), so its value under $\sigma$ is
zero.  The equations of nonzero weight and the $Z_i$ vanish at
$\xi_0$.  If $\xi_0\ne x$, the two constructed points contradict
(M3).  If $\xi_0=x$, use the negative point.  The functional
$b\mapsto\xi_{-1}(\operatorname{pr}_{-1}b)$ kills
$\mathfrak a^2$: in a product of two augmentation-zero homogeneous
elements of total weight $-1$, at least one has nonnegative weight
and evaluates to zero.  It kills the saturation as well, since
$O\to K$ is injective.  Rational spanning by $F$ forces its
vanishing on $A_{-1}$.  For total weight $-2$, the only additional
case is a product of two weight-$-1$ elements; it now also vanishes.
Rational spanning by $G$ forces vanishing on $A_{-2}$, and (SP.8)
forces vanishing on all negative weights.  This contradicts the
choice of $\xi_{-1}$.  We have proved that a relevant $P$ containing
$D_0$ must contain $\mathfrak m_B$.

These prime statements imply, by Noetherianity, power containments
\[
 (I_0R\,R_{++})^b\subset J',\qquad
 (\mathfrak m_BR\,R_{++})^{b'}\subset J'+D_0R
 \tag{SP.21}
\]
for some $b,b'$.  In large multidegrees they give
$I_0^bM_d=0$ and $\mathfrak m_B^{b'}M_d\subset D_0M_d$.
The first makes $M_d$ finite over $O$: $B/I_0^b$ is finite over
$O$, by its filtration with finite $B/I_0$-modules.  At a horizontal
prime containing $D_0$, $\pi$ is a unit, so the second containment
and Nakayama give $(M_d)_{\mathfrak p}=0$.  This proves (SP.17).

If (SP.17) holds, the degree-zero part of
$K\otimes_B(A/I)$ is $K$, so this finite-type $K$-algebra is
nonzero.  A $K$-point furnished by the Nullstellensatz restricts to
$\sigma$.  A mixed product in $D_0$ has nonzero value, hence this
point is two-sided.  Finally, a relevant prime containing
$J,\mathbf t,a$ and not $\mathfrak m_B$ would yield such a point
with $\xi(a)=0$, contradicting (M5).  Therefore
$\mathfrak m_BR\,R_{++}$ lies in the radical of $J+(\mathbf t,a)$;
a power lies in that ideal, which is (SP.9).
\end{proof}

\begin{lemma}[SP4: finite orbit cover and generic multiplicity one]
\label{claim:SP4}
Assume (M2) and (M4).  The two-sided points of $S(K)$ have finitely
many $K^\times$-orbits.  If $\sigma:B\to K$ and
$D_0\not\subset\ker\sigma$, then every weight piece of $A/I$,
localized at $\ker\sigma$, has $B_{\ker\sigma}$-length at most one.
Replacing all the $F_i$ and $E_i$ by their squares raises this bound
to at most $2^{r+1+r'}$.
\end{lemma}

\begin{proof}
Choose $f\in A_{-2}$ and $g\in A_2$.  Over $E$, the slice
\[
                   C=(A/I)[1/(\pi g)]/(f-1)
 \tag{SP.22}
\]
is a finite-type algebra.  A $K$-point of it is two-sided.  Its Euler
derivation on $A$ is $\mathcal E(b_j)=j\xi(b_j)$ and has
$\mathcal E(f)=-2$.  Any point derivation of $C$ kills $f$.
By (M4) it is linearly dependent with $\mathcal E$; evaluating that
dependence on $f$ forces the Euler coefficient to be zero and then
forces the other derivation to vanish.  After base change to $K$,
every closed point thus has zero cotangent space.  Nakayama makes
its local ring a field.  Nullstellensatz then shows that this
zero-dimensional reduced finite-type algebra is a finite product
of copies of $K$.  Consequently $C$ is a finite \emph{\'etale}
$E$-algebra.

Finite $B$-generating sets of $A_{-2}$ and $A_2$ supply finitely many
such pairs $(f,g)$.  Every two-sided point has both one of the $f$'s
and one of the $g$'s nonzero, by (SP.8).  A square root in $K$
normalizes $f$ to $1$ by the torus action.  The finite point sets of
the slices (SP.22) therefore give a finite orbit cover.

We also need the scheme multiplicity, not just the number of points.
On $U=(A/I)[1/(\pi fg)]$, $f$ is a unit.  Adjoin $z$ with
$z^2=f$.  There is an explicit isomorphism
\[
 U[z,z^{-1}]/(z^2-f)\ \simeq\ C[t,t^{-1}],
 \qquad b_j\longmapsto\overline b_jt^{-j},\quad z\longmapsto t.
 \tag{SP.23}
\]
Its inverse sends $\overline b_j$ to $b_jz^j$.
The left algebra is faithfully flat over $U$ (indeed finite free of
rank two), and the right algebra is reduced.  Thus $U$ is reduced.
This proves the required local reducedness directly, without
assuming that a finite orbit set was already a reduced scheme.

Let $\mathfrak p=\ker\sigma$.  If the localized quotient is zero
there is nothing to prove.  Otherwise the image of $B$ in $A/I$ is
finite over $O$, so its localization at this horizontal prime is an
Artinian local $E$-algebra.  Choose $fg\notin\mathfrak p$.
Its maximal ideal is nilpotent and acts by nilpotents on the reduced
ring in (SP.23), hence acts by zero.  The localized ring is therefore
an algebra over the residue field $L=\operatorname{Frac}(B/\mathfrak p)$,
and its degree-zero part is $L$.  Its even weight pieces have
dimension at most one, since $f$ is an invertible element of degree
$-2$.  If a piece of odd weight $j$ contains $0\ne h$, reducedness
gives $h^2\ne0$.  The even piece of weight $2j$ is one-dimensional,
so $h^2f^j\in L^\times$.  Thus $h$ is invertible; every other
element of weight $j$ is an $L$-multiple of $h$.  This proves the
length-one assertion for every weight.

For the final assertion filter a quotient by $u^2$ using
$uM/u^2M$.  That subquotient is an image of $M/uM$, with the
appropriate weight shift, and the other subquotient is $M/uM$.
Since the length-one bound holds at every weight, each replacement
of one equation by its square multiplies the bound by at most two.
There are $r+1+r'$ such replacements.
\end{proof}

\subsection{The leading form and the exact factor two}

\begin{lemma}[SP5: the lower Euler bound]
\label{claim:SP5}
Under (M1), make the Smith change (SP.14).  There is a
weight-preserving chart $\phi_0$ with only nonzero-weight variables,
with the same linear coordinates in those weights, and
$\phi_0(Z_i)=0$ for all $i$.  Let $J\subset R$ consist of the Rees
polynomials all of whose coefficients lie in $\ker\phi_0$.
If $x(a)=0$ and $(\mathbf t,a)$ is admissible, then for all
sufficiently large $d$,
\[
 2^{r+r'}\sum_{i<r}c_i
              \le\chi(\mathbf t;(R/J)/a(R/J))_d.
 \tag{SP.24}
\]
\end{lemma}

\begin{proof}
The $m$ equations $Z_i$ have linearly independent linear terms in
the $m$ zero-weight variables, because they rationally span that
cotangent space.  Use those equations as the new zero-weight
coordinates.  A formal substitution with invertible linear matrix
has an inverse: after solving the linear terms, solve its
homogeneous terms one total degree at a time.  Each correction has
higher total degree, so the coefficients stabilize.  Uniqueness
also makes the substitution weight preserving.  Set these new
zero-weight coordinates to zero.  This constructs $\phi_0$ and
preserves the linear terms of every remaining variable.

For a monomial $z^\nu$ in these variables let $n_-(\nu)$ and
$n_+(\nu)$ denote its negative and positive weighted degrees.
Every coefficient of a Rees element of degree $d$ has weight
$2(d_1-d_0-d_2)$, and its image under $\phi_0$ has
\[
                    n_-\ge 2d_0+4d_2.
 \tag{SP.25}
\]
For $N_0$ and $N_1$ this follows from the weight.  For $N_2$, the
term $A_{-2}^2A_2$ has this lower bound, while $H$ has no constant
or linear term; a weight-$-1$ monomial without a linear term has
$n_-\ge2$, so $H^2$ has $n_-\ge4$.  Products prove (SP.25).

Let $Q\subset R$ be the ideal obtained by requiring strict inequality
in (SP.25) in every Rees degree.  Leading coefficients show directly
that it is a homogeneous ideal, $J\subset Q$, $a\in Q$, and
$B$ acts on $R/Q$ through $x$.  It is $O$-torsion free: multiplying
a leading coefficient in $E$ by a nonzero scalar cannot kill it.
Moreover (SP.14) gives, in $R/Q$,
\[
                      F_i^2T_0=\pi^{2c_i}y_i^2T_0.
 \tag{SP.26}
\]
Indeed the difference $F_i-\pi^{c_i}y_i$ has no linear term and
has weight $-1$, so every term of the difference of squares has
negative weighted degree strictly greater than $2$.

The degree-$d$ part of $R/Q$ is a finite torsion-free $O$-module.
Its rank is
\[
 H(d)=\#\{\nu:\ n_-(\nu)=2d_0+4d_2,
                       \ n_+(\nu)=2d_1+2d_2\}.
 \tag{SP.27}
\]
To check equality of the rank, rather than only an upper bound,
use the independently realized linear coordinates in the chart.
The leading weight-$\pm1$ forms are independent linear forms; the
leading weight-$\pm2$ forms are independent linear forms in the
weight-two variables plus quadratics in the weight-one variables.
Subtract those quadratics over $E$.  These leading forms therefore
generate the entire polynomial algebra on each sign of variables.
Every monomial of even weighted degree is a product of
weighted-degree-two monomials.  Finally
\[
             A_{-2}^{d_0+2d_2}A_2^{d_1+d_2}\subset L_d
\]
realizes every monomial counted in (SP.27).  This proves full rank
after tensoring with $E$.  No claim that the integral lattice is
unimodular is being used.  Reduction of a finite free $O$-lattice
modulo $\pi$ has length its rank, so
$\ell_B((R/Q)/\pi)_d=H(d)$.

There is a graded surjection $(R/J)/a\to R/Q$.  Monotonicity in
SP2 and the scaling inequality (SP.11), applied to the $r$ sections
in (SP.26), give the lower bound
\[
 2\sum_i c_i\,
       \Delta_0^{r-1+s}\Delta_1^{r'+s'}\Delta_2H(d).
 \tag{SP.28}
\]
All needed admissibility conditions descend under the surjection.
The remaining numerical difference has a short exact calculation.
The generating function for monomials of bidegree $(a,b)$ is
\[
 \frac{1}{(1-X)^r(1-X^2)^s(1-Y)^{r'+1}(1-Y^2)^{s'+1}}.
\]
The three differences in (SP.28) are shifts by $(2,0),(0,2),(4,2)$.
After multiplying by their factors the generating function becomes
\[
 (1+X)^{r-1}(1+Y)^{r'}
               \frac{1-X^4Y^2}{(1-X)(1-Y)(1-Y^2)}.
 \tag{SP.29}
\]
For sufficiently large $a,b$, the coefficient in the last quotient
is
$(\lfloor b/2\rfloor+1)-(\lfloor(b-2)/2\rfloor+1)=1$.
The polynomial factor therefore makes the coefficient exactly
$2^{r-1+r'}$.  Substitution in (SP.28) proves (SP.24).
\end{proof}

\begin{theorem}[SP6: the literal \texttt{Lemma\allowbreak M}]
\label{claim:SP6}
Under exactly (M1)--(M5) and the ring and valuation hypotheses
above, there is a finite cover of the two-sided orbits,
$\ell_OC_{-1}(F)$ is finite, and (SP.7) holds for every finite orbit
cover.  Thus \texttt{Lemma\allowbreak M} has no remaining bespoke input.
\end{theorem}

\begin{proof}
The cotangent finiteness and (SP.14) already prove the finiteness of
the left side, and SP4 supplies a finite cover.  We first arrange
$\phi_0(a)\ne0$.  If it is zero, rational spanning by $F$ and the
nonzero weight-$-1$ linear coordinate imply that some $F_i$ has
nonzero chart image.  Choose $g\in A_1$ with nonzero chart image;
there is such a linear coordinate because its dimension is $r'+1$.
Replace $a$ by $a+gF_i$.  It still has weight zero and augmentation
zero, it agrees with $a$ on $S(K)$, and its chart image is a product
of two nonzero power series in a domain.  This replacement preserves
(M5) and the right side of (SP.7).

For the coefficientwise kernel $J$ of SP5, multiplication by $a$
is injective on $R/J$: its image is multiplication by the nonzero
series $\phi_0(a)$ on each coefficient.  SP3 supplies admissibility.
Apply SP2, first moving $a$ past the other sections, then passing to
their quotient.  Together with SP5 this gives
\[
 2^{r+r'}\sum_i c_i
 \le \chi(\mathbf t;(R/J)/a)_d
 =\chi(\mathbf t,a;R/J)_d
 \le\chi(a;R/(J+(\mathbf t)))_d .
 \tag{SP.30}
\]
We bound the final term on the actual $B$-module
$M_d=(R/(J+(\mathbf t)))_d$.

By SP3, $M_d$ is finite over $O$, and its generic support satisfies
(SP.17).  At such a prime choose $f\in A_{-2}$ and $g\in A_2$
with $fg\notin\mathfrak p$.  They become invertible after
localization, so $N_0,N_1,N_2$ all become free rank-one modules in
their indicated weights; $N_2=N_0$ there since
$f^2g=(fg)f\in A_{-2}^2A_2$.  In degree $d$ with $d_i\ge1$, the
section relations (SP.16) impose every homogeneous multiple of the
equations
\[
                    I^{[2]}=(F_i^2,G_i,E_i^2,E'_i,Z_i).
\]
Indeed one may take that multiple in degree $d-\mathbf e_i$ in
the corresponding Rees direction.  Consequently the localized
weight-$2(d_1-d_0-d_2)$ piece of $A/I^{[2]}$ surjects onto
$(M_d)_{\mathfrak p}$.  SP4 gives
\[
                    \ell_{B_{\mathfrak p}}(M_d)_{\mathfrak p}
                         \le2^{r+1+r'}.
 \tag{SP.31}
\]

Here is the length--valuation formula needed to sum these
multiplicities.  For any finitely generated $O$-module $M$ with a finite
$B$-action and any $b\in B$ invertible on $E\otimes_OM$,
\[
 \ell_B(M/bM)-\ell_B(M[b])
       =\sum_{\sigma}m_\sigma v(\sigma(b)),\qquad
 m_\sigma=\ell_{B_{\ker\sigma}}M_{\ker\sigma}.
 \tag{SP.32}
\]
The sum is over all embeddings of the generic support fields into
$K$.  To prove it, remove the $O$-torsion of $M$; its kernel and
cokernel under any endomorphism have the same finite length, so
contribute zero.  On the remaining free $O$-lattice the left side is
$v_O(\det(b\mid E\otimes_OM))$.  Decomposition of the finite
commutative $E$-algebra acting on this vector space into Artinian
local factors expresses this determinant as a product of field
norms, each raised to the length of the corresponding localized
module.  In characteristic zero each norm is the product over its
embeddings into $K$; applying $v$ gives (SP.32).  Finite $B$-length
equals finite $O$-length because their residue fields agree.

Every embedding in (SP.32) is the restriction of a two-sided point
by SP3, and hence of some representative in the given cover $T$.
By (M5), $a$ is nonzero at all such embeddings and is invertible on
$E\otimes_OM_d$.  Since the image of $B$ on $A/I$ is finite over
$O$, each $\xi(a)$ is integral over $O$.  A root of a monic
polynomial with coefficients of nonnegative valuation cannot have
negative valuation: its leading power would be the unique term
of smallest valuation.  Thus all $v(\xi(a))$ are nonnegative.
Repeated restrictions in $T$ can therefore only increase their
sum.  Equations (SP.31)--(SP.32) imply
\[
            \chi(a;M_d)\le
                    2^{r+1+r'}\sum_{\xi\in T}v(\xi(a)).
 \tag{SP.33}
\]
Combine (SP.30) and (SP.33), divide by $2^{r+r'}$, and use (SP.14).
This is exactly (SP.7), including the source's constant $2$.
\end{proof}

The classical background used in SP1--SP6 is the DVR structure
theorem, finite-dimensional linear algebra, the Nullstellensatz and
Nakayama's lemma.  The special multigraded positivity, prime support,
formal substitution, orbit reducedness and numerical difference
arguments have been proved above.  For the two algebraic-geometric
background statements, see the Stacks Project, Tags
\href{https://stacks.math.columbia.edu/tag/00FV}{00FV} and
\href{https://stacks.math.columbia.edu/tag/00DV}{00DV}; the finite
\'etale conclusion in SP4 is also Tag
\href{https://stacks.math.columbia.edu/tag/00U3}{00U3}.

\subsection{From actual rank-four cocycles to this cotangent lattice}

We now return to the source's original $\Omega$, odd prime $\ell$,
coefficient ring $O=c.O\subset\overline{\mathbb Q}_\ell$, and datum
$D:\texttt{Raw}\ \Omega\ c$.  The statement $D.\texttt{Good}$
includes $e=D.e\ge2$, an open-kernel finite-order character $D.\psi$,
the specified places $D.\mathrm{pl}(p)$, and
$\Psi=D.\psi\chi_\ell^e$.  It also includes
$\overline\Psi(g_0)\ne1$ and a local element at $\ell$ on which
$\chi_\ell$ has infinite order and $D.\psi$ is trivial.
Put
\[
 S_0=D.S,\quad
 T=\{\ell\}\cup\{p\in S_0:p\ne\ell,
                      \ \ell\mid|\psi(I_p)|\},\quad
 W=\operatorname{Frac}(O)/O(\Psi^{-1}).
 \tag{SP.34}
\]
The underlying torsion module in (SP.34) is the source's
\texttt{QO}; its twist, valuations and decomposition subgroups are
unchanged.  Let $\mathcal S$ be the strict Selmer module imposing
trivial restriction to every inertia group and to the decomposition
groups at primes in $S_0$.  Then
$D.\mathrm{sha}=\ell_O\mathcal S$ by definition.

\begin{proposition}[SP7: the actual \texttt{sha\allowbreak Coc} and \texttt{m\allowbreak Prime} comparison]
\label{claim:SP7}
Assume $D.\texttt{Good}$ and RV3.  Then $\mathcal S$ has finite
length.  There is a single $a\in O\setminus\{0\}$ and an additive
injection
\[
          Z:\mathcal S\lhook\joinrel\longrightarrow
                  \{z:G_{\mathbb Q}\to O/(a)\}
 \tag{SP.35}
\]
whose values are exactly of the source's type $D.\texttt{Sha\allowbreak Coc}$:
locally constant $\Psi^{-1}$-cocycles, zero at $g_0$, zero on the
residual-trivial inertia tested by $D.\texttt{Unr\allowbreak T}$, and
coboundaries over $O/(a)$ on every specified decomposition group
in $T$.
For every actual $\mathcal M:\texttt{Model}\ D$, equipped with the
source \texttt{Graded\allowbreak Algebra} structure on its first-weight-zero
algebra $A$ and its pieces \texttt{Model.\allowbreak gr}, one has
\[
       D.\mathrm{sha}\le
       \ell_OC_{-1}(\mathcal M.F).
 \tag{SP.36}
\]
Thus this particular incoming edge to the trace inequality is closed
once an actual model is supplied; no additional cocycle comparison
or choice-linearity hypothesis remains.
\end{proposition}

\begin{proof}
The strict Selmer module for $S_0$ is a submodule of the one imposing
only the decomposition condition at $\ell$, because $\ell\in S_0$.
RV3 therefore gives finite length.  Choose $0\ne b\in O$ killing
$\mathcal S$ and write $\theta=\Psi^{-1}$.
The element $u=\theta(g_0)-1$ is a unit, since its reduction is
nonzero.  Every cohomology class has the unique normalized cocycle
\[
                 N[z](g)=z(g)-(\theta(g)-1)u^{-1}z(g_0).
 \tag{SP.37}
\]
Changing $z$ by a coboundary leaves (SP.37) unchanged; it vanishes
identically exactly when the class is zero.  Hence $N$ is an
$O$-linear injection into locally constant cocycles and every
$N[z](g)$ is killed by $b$.

For each $v\in T$, choose $s_v\in D_v$ with
$\theta(s_v)\ne1$.  At $\ell$ use the infinite-order cyclotomic
element specified by $D.\texttt{Good}$.  At $v\ne\ell$ the
definition of $T$ makes $\psi(I_v)$ a nontrivial group of order
divisible by $\ell$, while $\chi_\ell$ is unramified there; an
element on which $\psi$ is nontrivial suffices.  Set
\[
                       a=b\prod_{v\in T}(\theta(s_v)-1)\ne0.
 \tag{SP.38}
\]
The normalized cocycle is a local coboundary
$N[z](g)=(\theta(g)-1)m_v$ on $D_v$.  Since
$b(\theta(s_v)-1)m_v=bN[z](s_v)=0$, (SP.38) kills $m_v$ as well
as all the values of $N[z]$.
The map $O/(a)\to\operatorname{Frac}(O)/O$ sending
$\overline y$ to $y/a+O$ identifies its source with the $a$-torsion.
Transport both the cocycle and these local parameters through this
map.  This constructs the additive injection (SP.35), with all its
local and normalization properties.

We check the potentially stronger-looking $\texttt{Unr\allowbreak T}$ clause.
Outside $S_0$ the finite character is unramified; at a prime of
$S_0\setminus T$ its inertia image has order prime to $\ell$.
The cyclotomic character is unramified away from $\ell$.
If the residual diagonal representation is trivial on an inertia
element, then $\theta$ is $1$ modulo the maximal ideal there.  A
root of unity of order prime to $\ell$ that is $1$ modulo that ideal
equals $1$: divide $X^n-1$ by $X-1$ and evaluate at the root, where
the second factor reduces to the unit $n$.
Thus $\theta(g)=1$ for exactly the inertia elements needed in
$\texttt{Unr\allowbreak T}$, and local coboundary gives $N[z](g)=0$.
Conjugating the chosen places gives the same assertion for every
valuation subring in that definition.  A valuation trivial on
$\mathbb Q$ is trivial on its algebraic closure and has trivial
inertia, so it adds no condition.  This verifies the literal
$D.\texttt{Sha\allowbreak Coc}$ statement.

It remains to use the actual model rather than postulate an
injection into cotangent duals.  Its record supplies a surjection
$q:R\to S$ onto the universal representation ring, a
$\mathbb Z^2$-weighted presentation on $R$, and, for every cocycle
$Z\mu$, an $O$-valued point derivation $\delta_\mu$ of weight
$-w_\Psi=(0,-1)$.  It kills $\ker q$ modulo $a$.  On lifts of the
two relevant universal matrix entries its reductions are
\[
 \delta_\mu(r(g)_{20})=D.\mathrm{th}(g)Z\mu(g),\qquad
 \delta_\mu(r(g)_{31})=\Psi(g)Z\mu(g)
                         \pmod a.
 \tag{SP.39}
\]
All other entries have different weights.  These are precisely the
\texttt{der}, \texttt{hom}, \texttt{at0} and \texttt{gen} fields of
$\texttt{Model}$, including the universal quantification over
possible lifts of an entry.

The choices $\delta_\mu$ need not be additive.  Their reductions
are additive.  First the augmentation descends through $q$:
if $q(f)=0$, choose a lift $b$ of the $(0,0)$ entry at the identity;
the \texttt{at0} field applied to $b+f$ and to $b$ gives
$x(b+f)=x(b)$, hence $x(f)=0$.
Now the derivation
$\delta_{\mu+\nu}-\delta_\mu-\delta_\nu$ descends modulo $a$ to
$S$, and is zero on every universal matrix entry by (SP.39) and
the weight restriction.  Associate to it the square-zero algebra
map $S\to O\oplus O/(a)$ with first coordinate the augmentation.
The model's generation field says that two maps to this complete
local $O$-algebra which agree on the universal representation are
equal.  Compare with the zero derivation to deduce vanishing on
all of $S$ and hence all of $R$.
Restriction to the first-weight-zero subalgebra $A\subset R$
therefore gives an additive map
$\mathcal S\to\operatorname{Hom}_O(A,O/(a))$.
It is injective: the $(2,0)$ entry lies in that subalgebra and
the coefficient $D.\mathrm{th}(g)$ in (SP.39) is a unit, so a zero
derivation forces $Z\mu=0$.
Each equation $\mathcal M.F_i$ lies in $\ker q$ and has weight
$-1$ in $A$.  The hypotheses of SP1 now all hold and yield
(SP.36).  This proves the actual source \texttt{m\allowbreak Prime} conclusion
without assuming that an arbitrary choice of lifts was linear.
\end{proof}

\subsection{Returning from auxiliary coefficients to the original character}

\begin{proposition}[SP8: exact character invariants and coefficient return]
\label{claim:SP8}
Let $f:O\hookrightarrow O'$ be an injective compatible inclusion of
coefficient DVRs in $\overline{\mathbb Q}_\ell$, and let
$\psi' = f\circ\psi$ be a Dirichlet character of the same modulus
$N$.  Then order, unit-character kernel, conductor, the parity
condition $\psi(-1)=(-1)^e$, and all the source's quantities
\[
          \operatorname{comp}_r\psi,\quad\theta_0,\quad\theta_1,
          \chi_P,\quad n_0,\quad c_2,\quad n_{\rm II},\quad n_{\rm III}
 \tag{SP.40}
\]
are preserved or commute with $f$ as appropriate.  In particular
membership in $h.X(c,e)$ and either case II or case III is preserved
on the same host and the same $\Omega$.
Together with RV1 this removes all auxiliary root-of-unity
assumptions from a bound expressed using the normalized Selmer
length: the identical bound holds over the original $O$.
\end{proposition}

\begin{proof}
Injectivity gives $f(\psi(u))=1$ if and only if $\psi(u)=1$,
for every unit $u$ modulo $N$, and $\psi'^n=1$ if and only if
$\psi^n=1$.  Thus the unit kernels and orders agree.
A character factors through $d\mid N$ exactly when its unit kernel
contains the kernel of $(\mathbb Z/N)^\times\to(\mathbb Z/d)^\times$.
Its complete set of possible conductors, and hence its conductor,
is unchanged.  The operations $\operatorname{comp}_r$ are pullbacks
by the fixed primary projection of $(\mathbb Z/N)^\times$; they
commute with any coefficient map.
The prime-to-$\ell$ part $\theta_0=\ell\mathrm{Prime}(\psi)$ is
defined by group powers determined by the order.  It too commutes
with $f$, as do
\[
 \theta_1=\theta_0(\operatorname{comp}_2\theta_0)^{-1},\qquad
 \chi_P=\theta_0(\operatorname{comp}_\ell\theta_0)^{-1}.
\]
The definitions
\[
 \begin{aligned}
 f_0&=\operatorname{ord}_{\operatorname{lcm}(4,\operatorname{ord}\theta_1)}(\ell),
 &n_0&=v_2(\ell^{f_0}-1),\\
 c_2&=v_2(\operatorname{cond}(\operatorname{comp}_2\psi))
 \end{aligned}
 \tag{SP.41}
\]
then prove their preservation.  So do the defining ceiling logarithms
for $n_{\rm II}$ and $n_{\rm III}$, written in (SP.48) below.
Injectivity also gives
$f(\psi(-1))=(-1)^e\Longleftrightarrow\psi(-1)=(-1)^e$.
Host character membership is a unit-kernel condition together with
this parity.  Case II uses $e\ge2$, the parity of
$\operatorname{ord}\chi_P$, and the fixed numerical condition on
$\ell,h,e$; case III uses that same fixed numerical condition,
odd order of $\psi$, and $\psi\ne1$.  All are therefore preserved.

For an individual $\psi$, adjoin a primitive
$2^{n_{\rm II}+n_{\rm III}}$-th root of unity to the finite
coefficient field and take its compositum with the original field.
Its ring of integers $O'$ is again an allowed coefficient ring,
contains the required roots for each separate exponent, and
contains $O$ injectively.  RV1 identifies the original strict
Selmer module after this extension with its scalar extension.
If $e_{O'/O}$ is the relative ramification index, then
\[
 \ell_{O'}\mathcal S'=e_{O'/O}\ell_O\mathcal S,
 \qquad v_{O'}(\ell)=e_{O'/O}v_O(\ell),\qquad
 \frac{\ell_{O'}\mathcal S'}{v_{O'}(\ell)}
       =\frac{\ell_O\mathcal S}{v_O(\ell)}.
 \tag{SP.42}
\]
Finiteness is equivalent in both directions.  Thus a bound proved
over $O'$ gives the same bound over $O$, with its original
character, conductor, parity, case, and normalization.  The chosen
coefficient extension may depend on $\psi$; a constant selected
before $\psi$ is not changed by this construction.
\end{proof}

\subsection{The exact arithmetic remaining after a finite trace catalogue}

The existence of the following catalogue is a retained input, not a
synonym for a general automorphy theorem.  Its quantifier order is
part of the input.  There must be integers $c_*,A_*\ge0$ such that,
for every $e\ge2$, $M\ge1$, and prime $p\ge5$ with $p\nmid M$,
there is a finite set $\mathcal D$ of complex multisets
$\beta_s$ of cardinality four with the following properties:
\begin{enumerate}
\item $|\mathcal D|\le c_*[M(e+2)]^{A_*}$, and every member of every
$\beta_s$ is an algebraic integer.
\item Every automorphism of $\mathbb C$ permutes the multisets
$\beta_s$, with multiplicity, and
\[
 \left|\sum\beta_s\right|
       \le (3\sqrt p+1)p^{(e+1)/2}.
 \tag{SP.43}
\]
\item For every odd prime $\ell\ne p$, every original algebraic
closure $\Omega/\mathbb Q$, every isomorphism
$\iota:\overline{\mathbb Q}_\ell\simeq\mathbb C$, and every
representation satisfying the literal \texttt{Adm\allowbreak RAlt} conditions
at level $M$, weights $(e+1,e,1,0)$ and duality weight $e+1$, there
is a label $s\in\mathcal D$.  Equal labels imply representation
isomorphism.  At every allowed place above $p$ and every arithmetic
Frobenius $\sigma$ there,
\[
                  \iota(\operatorname{tr}r(\sigma))=\sum\beta_s.
 \tag{SP.44}
\]
\end{enumerate}
Here \texttt{Adm\allowbreak RAlt} includes irreducibility, the literal
\texttt{Is\allowbreak Aut\allowbreak QA} condition, \texttt{Level}, the stated cyclotomic
dual twist, and \texttt{Safe\allowbreak RA}.  In the source the catalogue is
first expressed using geometric Frobenius; its inverse is used in
(SP.44).  Both this inverse convention and the classification of
all the actual representations are required.  The catalogue is
chosen before $\ell,\Omega,\iota$ and the representations.

\begin{proposition}[SP9: the integral product and its uniform trace bound]
\label{claim:SP9}
The preceding catalogue input supplies the literal
\texttt{Number\allowbreak Sp4Alt}, with
\[
                    C=c_*2^{A_*+1},\qquad A=A_*+1.
 \tag{SP.45}
\]
More precisely, after $e,M,p,k>0$ and three $k$-th roots of unity
$z_1,z_2,z_3$ have been fixed, there is a nonzero integer $N_*$,
chosen before $\ell,\Omega,\iota$ and any finite representation
family, such that
\[
 \begin{aligned}
 \log|N_*|&\le C\varphi(k)[M(e+1)]^A\log p,\\
 |N_*|_\ell&\le
 \prod_{r\in\mathcal X}
 \left|\operatorname{tr}r(\sigma)
       -\iota^{-1}\bigl(p^{e+1}z_1+p^ez_2+pz_3+1\bigr)\right|_\ell.
 \end{aligned}
 \tag{SP.46}
\]
Every displayed trace difference is nonzero.  The second inequality
holds for every finite set $\mathcal X$ of the specified actual
representations, with no two distinct members isomorphic, at every
place and Frobenius allowed in (SP.44).
\end{proposition}

\begin{proof}
Put $y=p^e$ and $x=p^{(e+1)/2}$.  Since $e\ge2$ and $p\ge5$,
$x\le y/2$ and $\sqrt p\,x\le y$, so (SP.43) is at most
$7y/2$.  For any three complex numbers of modulus one, the leading
term $p^{e+1}z_1$ has modulus $py$, whereas the sum of the trace and
the other three terms has modulus at most
\[
          \frac72y+y+p+1<py.
\]
The strict inequality follows from
$(p-9/2)p^e>p+1$, whose smallest case is $p=5,e=2$ and whose two
sides then separate monotonically.  Thus the difference is nonzero;
its modulus is less than $2py\le (py)^2=p^{2(e+1)}$.

For $u\in(\mathbb Z/k)^\times$ let
\[
 b(s,u)=\sum\beta_s-
          \bigl(p^{e+1}z_1^u+p^ez_2^u+pz_3^u+1\bigr),
 \qquad N_* =\prod_{s,u}b(s,u).
 \tag{SP.47}
\]
Every factor is an algebraic integer.  An automorphism of
$\mathbb C$ permutes the $s$'s and acts on all $k$-th roots by one
common unit modulo $k$; it therefore permutes the pairs $(s,u)$.
Consequently $N_*$ is fixed by every automorphism of its normal
number field over $\mathbb Q$, so is rational, hence an integer.
The nonvanishing just proved gives $N_*\ne0$.  Its logarithmic
size is at most
$2|\mathcal D|\varphi(k)(e+1)\log p$.
Use $e+2\le2(e+1)$ and $e+1\le M(e+1)$ to obtain (SP.45)--(SP.46).

For a fixed finite family $\mathcal X$, assign to $r$ its label
$s(r)$ and select the factor $b(s(r),1)$.  These selected factors
are distinct because equal labels imply isomorphism.  Equation
(SP.44) identifies their inverse $\iota$-images with exactly the
trace differences in (SP.46).  Every omitted factor is still an
algebraic integer and thus has $\ell$-adic norm at most one.
Taking the norm of the complete product proves the asserted
inequality.  The same argument with a singleton proves the
nonvanishing for every individual representation.  All the choices
defining $N_*$ preceded the auxiliary prime and the family, so no
uniformity has been exchanged in this argument.
\end{proof}

\subsection{The paired-count join and its remaining geometric input}

For precision, we state the outstanding input at the level at which
the preceding proofs can actually use it.  The following two
assertions remain to be reconstructed.
\begin{enumerate}
\item[(G1)] For every original host $h$, $h.\mathrm{Big}(t)$,
odd $\ell$, $0<e<t$ and $\psi\in h.X(c,e)$ in case II or III,
after adjoining the roots allowed in SP8 there is an actual
$D:\texttt{Raw}\,\Omega\,c$ with \texttt{Good} and
\texttt{Of\allowbreak Case}.  There are the source parameters $n',p_1,m_0$,
with $m_0\mid h.o$, $c_2\le n'$, $D.e=e$ and
$D.\psi=\operatorname{dirG}(\psi)$, for which
\[
 \begin{aligned}
 n_{
m II}&=\lceil\log_2\max
       (\ell h.N4^{n_0},\,2^{c_2+6}h.o\ell^3h.q^2)\rceil,\\
 n_{
m III}&=\lceil\log_2\max
       (\ell h.q\,p_{\!1}(\chi)4^{n_0},\,8\ell^2)\rceil,\\
 (n',p_1)&=(n_{
m II},1)\quad\hbox{or}\quad
                  (n_{
m III},p_{\!1}(\chi)),\\
 L&=\ell 2^{n'}h.Np_1,
 & k&=(\ell-1)2^{n'}(2h.o).
 \end{aligned}
 \tag{SP.48}
\]
Here $p_{\!1}(\chi)$ denotes the fixed source \texttt{p\allowbreak One}, not
the variable $p_1$.  The characters $D.\psi,D.\xi,D.\mathrm{om}$
have $k$-th power one and satisfy the source \texttt{Kills\allowbreak Mu}
conditions at level $L$; $D.S$ is the prime support of
$2\ell h.Np_1$.  With one common $\iota$, every actual irreducible
$D.\mathrm{EPt}$ representation satisfies the \texttt{Adm\allowbreak RAlt}
conditions used in SP9 at level $L$.
\item[(G2)] For this same $D$ there is an actual
$\mathcal M:\texttt{Model}\,D$ whose first-weight-zero algebra,
point, grading and equations satisfy (M1)--(M4) above.  Every
two-sided point of its equation scheme determines an irreducible
$D.\mathrm{EPt}$ representation.  The assignment is compatible
with the trace element $a_\sigma$:
\[
 x(a_\sigma)=0,\qquad
 \xi(a_\sigma)=\operatorname{tr}\rho_\xi(\sigma)
                       -\operatorname{tr}D.r_0(\sigma).
 \tag{SP.49}
\]
Isomorphic assigned representations come from the same
$\mathbb G_m$ orbit.  The trace element belongs to the degree-zero
algebra, and the valuation and coefficient map are those of the
original $D$, without a change of residual character or local
condition.
\end{enumerate}
These assertions explicitly include the construction of the
weighted model and its local-to-global comparison.  Naming
\texttt{Lemma\allowbreak P}, \texttt{Cohomology\allowbreak Hyp} or an automorphy package
does not prove them.

\begin{theorem}[SP10: the complete arithmetic and coefficient-return join]
\label{claim:SP10}
Assume (G1), (G2), and the finite catalogue preceding SP9.
Then SP1--SP9, RV3 and RV1 imply the literal
\texttt{Count\allowbreak II}$\wedge$\texttt{Count\allowbreak III}.  In particular, for
each primitive character $\chi$ there is one natural exponent
$C_\chi$, chosen before $h,t,\Omega,\ell,c,e,\psi$, such that
every original character in either case has finite strict Selmer
length and
\[
 \frac{\ell_O V(\Omega,c,e,\psi)}{v_O(\ell)}\log\ell
                         \le(\ell h.N e)^{C_\chi}.
 \tag{SP.50}
\]
This is a conditional endpoint.  The theorem closes the join after
the displayed geometric and catalogue inputs; it does not close
(G1), (G2), or the catalogue itself.
\end{theorem}

\begin{proof}
Choose $C,A$ from SP9, before any host or coefficient choices.
For the current $h,\ell,e,\psi$, put
$p=p_0(\ell h.Np_1)$, the least prime at least five not dividing
$\ell h.Np_1$.  Bertrand's theorem applied to $\ell h.Np_1+5$
gives
\[
                    5\le p\le2\ell h.Np_1+10.
 \tag{SP.51}
\]
In particular $p\nmid L$ and $p\notin D.S$.  Choose arithmetic
Frobenius $\sigma$ at $D$'s specified place over $p$.  Its
cyclotomic value is $p$.  Set
\[
 \begin{aligned}
 z_1&=\iota((D.\mathrm{om}^{-1}D.\xi D.\psi)(\sigma)),\quad
 z_2=\iota(D.\psi(\sigma)),\quad
 z_3=\iota((D.\mathrm{om}^{-1}D.\xi)(\sigma)).
 \end{aligned}
\]
They are $k$-th roots of unity, and direct substitution in the four
diagonal characters $D.r_0=(D.\mathrm{th},D.\Psi,D.\mathrm{al},1)$
shows that its trace is
$\iota^{-1}(p^{e+1}z_1+p^ez_2+pz_3+1)$.
SP9 thus makes every irreducible trace difference nonzero.
By (SP.49), $a_\sigma$ satisfies (M5).

Apply SP6 to obtain a finite cover $T$ of two-sided points and the
cotangent inequality.  Remove redundant representatives until
there is one per orbit.  Its assigned representations form a
finite set $\mathcal X$ with no two distinct members isomorphic:
this is exactly the last orbit assertion of (G2).  The value of
$a_\sigma$ is constant on each orbit.  SP7 and SP6 now give
\[
 \ell_O(D.\mathrm{Sha})
       \le\ell_OC_{-1}(F)
       \le2\sum_{r\in\mathcal X}
                 v_O(\operatorname{tr}r(\sigma)
                                  -\operatorname{tr}D.r_0(\sigma)).
 \tag{SP.52}
\]
For an algebraic $\ell$-adic number $z\ne0$ the extension of the
normalized $O$-valuation is
$v_O(z)=-v_O(\ell)\log|z|_\ell/\log\ell$.
Consequently (SP.52) is the source's actual $D.\mathrm{Ineq}(\sigma)$,
with its exact factor two and its finite-length assertion.
The product estimate (SP.46) then yields
\[
 \frac{\ell_O(D.\mathrm{Sha})}{v_O(\ell)}\log\ell
 \le -2\log|N_*|_\ell
 \le 2\log|N_*|
 \le 2Ck[L(e+1)]^A\log p.
 \tag{SP.53}
\]
The middle inequality is just
$v_\ell(N_*)\log\ell\le\log|N_*|$ for a nonzero integer;
the last also uses $\varphi(k)\le k$.  This explains why both the
integrality and the order of the quantifiers in SP9 are needed.

The original Selmer group injects into $D.\mathrm{Sha}$ even
though $D.S$ can contain additional primes.  Indeed, at a prime
$p\nmid\ell h.N$ the original character is unramified, and its
negative-twist Frobenius value is $\zeta p^{-e}$ for a root of
unity $\zeta$.  It is not one since $e>0$.  On $E/O$ multiplication
by $\zeta p^{-e}-1\ne0$ is surjective.  Thus the unramified
cohomology
$(E/O)/(\zeta p^{-e}-1)(E/O)$ vanishes, and restriction from
decomposition-group cohomology to inertia is injective.
Every original class already vanishes on inertia at every place;
it therefore vanishes on each added decomposition group as well.
At a prime dividing $\ell h.N$ this was already its defining
strict local condition.  This proves the exact inequality called
\texttt{sha\allowbreak Eq}, rather than an unwarranted equality, on the same
$\Omega$ and all its chosen valuations.

It remains to verify a single exponent uniformly in the host.
Write $o=\operatorname{ord}\chi$, $P=p_{\!1}(\chi)\ge1$,
$f=\varphi(4o)$, and $X=\ell h.Ne$.  The host's numerical and
character conditions give
\[
 \begin{gathered}
 X\ge16,\quad \ell,h.N,e,h.o,h.q\le X,\quad
 2^{c_2}\le h.N,\quad 2^{n_0}<\ell^f. 
 \end{gathered}
 \tag{SP.54}
\]
For clarity, the last bound uses no analytic theorem.  The order
of $\theta_1$ is prime to $\ell$ and divides $2o$.
Here is the finite-group reason for the latter statement.
Under the Chinese remainder decomposition, a unit supported on an
odd prime-power factor dividing $h.q$ has its $o$-th power in
$h.H$: it is trivial at the other primary factors and its image
under $\chi$ has order dividing $o$.  A unit supported on any
other odd prime-power factor has its square in $h.H$, by the
two square-subgroup conditions defining $h.H$.  Since
$h.H\subset\ker\psi$, every odd primary component of $\psi$
has order dividing $2o$.  Taking its prime-to-$\ell$ part preserves
that divisibility, and removing its $2$-primary-modulus component
leaves exactly $\theta_1$.  The product of the remaining primary
components therefore has order dividing $2o$.
It follows that
$\operatorname{lcm}(4,\operatorname{ord}\theta_1)$ divides $4o$;
the order $f_0$ of $\ell$ modulo this modulus divides its unit
group order, hence is at most $\varphi(4o)$.  Therefore
$2^{v_2(\ell^{f_0}-1)}\le\ell^{f_0}-1<\ell^f$.
The other clauses use $16\mid h.N$, $h.q\mid h.N$,
$h.o=o\le h.q$, and the conductor divisibility
$2^{c_2}\mid h.N$, all on the original host and character.

The inequality $2^{\lceil\log_2 m\rceil}\le2m$ for positive
integers $m$ gives, from (SP.48) and (SP.54),
\[
 2^{n_{
m II}}\le X^{2f+10},\qquad
 2^{n_{
m III}}\le X^{2f+P+3}.
 \tag{SP.55}
\]
For the first, the two arguments of the maximum are at most
$X^{2f+2}$ and $64X^7\le X^9$, and the extra factor two is at
most $X$.  For the second they are at most $X^{2f+P+2}$ and
$8X^2\le X^3$, using $P\le X^P$.  Set $a=2f+P+13$; either
choice of $n'$ then satisfies $2^{n'}\le X^a$.  Directly,
\[
 \begin{aligned}
 k&\le X^{a+3},& L(e+1)&\le X^{a+P+4},\\
 2C&\le X^{2C},& \log p&\le p\le X^{P+3}.
 \end{aligned}
 \tag{SP.56}
\]
For the last bound use (SP.51), $\ell h.Np_1\le X^{P+2}$
and $X\ge16$.  Hence one explicit, deliberately generous choice is
\[
             C_\chi=2C+a+3+A(a+P+4)+P+3.
 \tag{SP.57}
\]
It depends only on $\chi,C,A$.  Combining (SP.53), the original
Selmer inclusion, and (SP.56) proves (SP.50) over the enlarged
coefficient ring.  SP8 returns it to the original $O$ with the
same exponent, case, order, kernel, conductor, parity and
normalized length.  The same $C_\chi$ works for the union of the
two cases, so gives the two literal count predicates separately.
\end{proof}

\paragraph{Precisely what is still open in the shared count branch.}
The algebraic leaves \texttt{Lemma\allowbreak M'} and \texttt{Lemma\allowbreak M} have been
proved in SP1 and SP6.  SP7 also proves \texttt{sha\allowbreak Coc} and
\texttt{m\allowbreak Prime} for the actual model; RV7 supplies both
\texttt{So\allowbreak V} and \texttt{So\allowbreak Fin}, while BR7 supplies the literal
\texttt{Cor\allowbreak X}.  These remove complete suppliers from the source's
\texttt{sp4Main\_\allowbreak alt} input list.  The following are still
substantive inputs to the human reconstruction:
\begin{enumerate}
\item \texttt{Symplectic\allowbreak Four\allowbreak Image\allowbreak Dense}, including both the finite
weight-zero algebra and exclusion of the endoscopic locus, with
one embedding valid for every actual irreducible point.
\item The entire \texttt{Cohomology\allowbreak Hyp} package.  LK2--LK4 and LD3
cover identified finite local and finite duality pieces; they do
not by themselves identify the arbitrary period-ring records in
its Bloch--Kato local clause, or prove its remaining global
finiteness, Euler, orthogonality and Greenberg--Wiles clauses.
\item The application of the literal \texttt{Lemma\allowbreak P} to the actual
Galois data, with its framed weighted representation functor,
specified local conditions, relation weights and exact
cohomological dimension counts.  The universal algebraic
implication \texttt{Lemma\allowbreak P} is handled separately by the framed
deformation component; this chapter does not infer its actual
hypotheses or the resulting \texttt{ex\allowbreak Model} comparison from that
implication alone.  The cumulative ledger records the status and
supplier of that separate algebraic theorem.
\item The finite catalogue used in SP9, including its polynomial
size, algebraicity, conjugation, purity and classification
properties for the literal \texttt{Adm\allowbreak RAlt}.  The source derives
it from a larger automorphic and local-level package; its name
\texttt{Hypotheses'} is not a classical theorem citation.
\item \texttt{Rigidity\allowbreak Symplectic\allowbreak Four} at these actual points.  Its
source producer uses the specified unramified deformation inputs
and adjoint lifting theorem; no generic rigidity assertion has
been substituted.
\item The internal construction and comparison edges in (G1)--(G2):
\texttt{data\allowbreak Ex}, \texttt{ex\allowbreak Model}, the verification of (M1)--(M4),
the two-sided point and conjugacy-orbit comparison, and the
\texttt{level}/\texttt{safe4'} passage to the exact admissible
representation class.  Their required inputs and source ranges
are recorded separately in the provenance ledger.  They remain
open human proof obligations even when their larger source input
packages have been named.
\end{enumerate}
The source's exact top-level implication is recorded for navigation;
SP10 is the readable conditional join actually proved here.  Thus
this chapter does not assert that \texttt{Count\allowbreak II} or
\texttt{Count\allowbreak III}, the original-family \texttt{Point\allowbreak HV}, or RH/GRH
has now been proved without the remaining suppliers.

\paragraph{Classical reference for the prime choice.}
The only additional number-theoretic theorem in the last join is
Bertrand's theorem, in the form: for every integer $n>1$ a prime
lies strictly between $n$ and $2n$.  See P.~Erd\H{o}s,
\emph{Beweis eines Satzes von Tschebyschef}, Acta Sci. Math.
(Szeged) \textbf{5} (1932), 194--198,
\url{https://users.renyi.hu/~p_erdos/1932-01.pdf}.
It is applied only with $n=\ell h.Np_1+5$ to obtain (SP.51).



% Version 0.26. Full literal LemmaP supplier.
% Pinned Lean source: 7cbf7a0c32fca1badf05931879a37c3f8a2b03b4.
\section{The weighted framed deformation presentation}
\label{sec:framed-deformation-presentation}

\status{FP1--FP8 prove the entire general \texttt{Lemma\allowbreak P} statement.
The proof constructs the universal ring, its weighted relative
generators, and its relations, and proves the exact cohomological
counts.  The stated hypotheses on the original group, local
conditions and framed obstruction spaces remain hypotheses of this
general theorem.  This section does not assert that an arithmetic
datum satisfies them, and does not construct its later
\texttt{Model} or prove automorphy.}

\subsection{The original data and the exact output}

Let $O$ be a complete Noetherian local ring with finite residue field
$k$ of odd characteristic $\ell$; $O$ need not be a domain or a DVR.
Write $\mathcal C_O$ for the complete Noetherian local $O$-algebras
whose residue map identifies their residue field with the specified
$k$.  This is the source's \texttt{Is\allowbreak CO}.  All maps are $O$-algebra
maps; they are automatically local and continuous, as FP1 proves.
For a group with a topology, a representation over $A\in\mathcal C_O$
is called continuous if its entries modulo $\mathfrak m_A^n$ are
locally constant for every $n$.  This is precisely \texttt{Lift}.

The global group $\Gamma$ is compact, Hausdorff, totally disconnected
and a topological group.  Its given hypothesis is
\[
 \Phi_\ell(\Gamma):\quad
 \operatorname{Hom}_{\rm cts}(H,\mathbb F_\ell)
       \text{ is finite for every open subgroup }H\le\Gamma .
 \tag{FP.1}
\]
The target is additive.  Locally constant homomorphisms to this
finite group are exactly the source's open-kernel homomorphisms in
\texttt{Phi\allowbreak Cond}; no finite-generation hypothesis on $\Gamma$ is
added.

Let $\iota$ be a finite index set, $X$ a free abelian group of finite
rank, and $w:\iota\to X$.  The specified \texttt{Form\allowbreak Data} consists
of an invertible matrix $J$ over $O$, with $J^t=J$ or $J^t=-J$,
a block map $\beta:\iota\to\mathcal B$, an involution
$\varsigma$ on $\mathcal B$, and $c\in X$, such that
\[
 J_{ab}\ne0\ \Longrightarrow\
       \beta(b)=\varsigma\beta(a),\quad w(a)+w(b)=c .
 \tag{FP.2}
\]
For $A\in\mathcal C_O$, put
\[
 \begin{split}
 H(A)&=\{h:h^tJh=J,\ h_{ab}=0\text{ if }\beta(a)\ne\beta(b),
                         \ h\bmod\mathfrak m_A=1\},\\
 \mathfrak g(A)&=\{Y:Y^tJ+JY=0,\
                     Y_{ab}=0\text{ if }\beta(a)\ne\beta(b)\},\\
 \mathfrak g_j(k)&=\{Y\in\mathfrak g(k):
                       Y_{ab}=0\text{ if }w(a)-w(b)\ne j\}.
 \end{split}
 \tag{FP.3}
\]
The first set is a group: preservation of the invertible form gives
an inverse, and all block and residue conditions are preserved.

Fix an original residual representation $\bar r$ and a continuous
lift $r_0$ over $O$.  Both satisfy \texttt{Normalises}: their
matrices preserve $J$ up to a scalar, permute the prescribed blocks,
and have zero entries when $w(a)\ne w(b)$.  Conjugation by
$\bar r$ therefore preserves every $\mathfrak g_j(k)$; denote this
representation by $\operatorname{ad}_j$.
A direct check of the form condition is
$\rho^{-t}[Y^t(\rho^tJ\rho)+(\rho^tJ\rho)Y]\rho^{-1}=0$;
the block permutation and equal-weight support give the other
two conditions in (FP.3).  Thus no extra adjoint-stability input
is being used.
A lift $r$ satisfies \texttt{In\allowbreak H} if
$r(s)r_0(s)^{-1}\in H(A)$ for every $s\in\Gamma$.

The specified frame has finite sets $T,Q$, a distinguished $v_0\in T$,
groups $\Gamma_v,\Gamma_q$ with topologies, continuous homomorphisms
$\phi_v:\Gamma_v\to\Gamma$ and $\psi_q:\Gamma_q\to\Gamma$,
and subgroups $N_q\le\Gamma_q$ with
$\bar r(\psi_qN_q)=1$.  No topological-group hypothesis on the
individual local groups is inserted.
For each $q$ a condition $L_q$ on local lifts has the following
five properties, which are exactly \texttt{Is\allowbreak Loc\allowbreak Cond}.
\begin{enumerate}
\item It is cut out, on any fixed lift satisfying \texttt{In\allowbreak H}, by
an ideal of its coefficient ring, universally under base change.
\item It is invariant under conjugation by $H(A)$.
\item It is invariant under diagonal torus conjugation whose
characters all reduce to one.
\item It lifts across every surjection $A\to A'$ in $\mathcal C_O$
whose kernel is killed by $\mathfrak m_A$.
\item For two \texttt{In\allowbreak H} lifts across such a surjection with the
same reduction, if the first satisfies $L_q$, the second does so
exactly when the two representations agree on $N_q$.
\end{enumerate}
Also $L_q(r_0|_{\Gamma_q})$ holds.

The framed spaces used in the output are specified cochain spaces,
not an unspecified choice of Selmer complex.  For a $k$-representation
$V$ of $\Gamma$, let $H^1_F(V)$ consist of pairs
$(z,(y_v))$ such that $z$ is a locally constant cocycle,
\[
 z(\psi_q n)=0\ (n\in N_q),\qquad
 z(\phi_v s)=(\phi_v s)\cdot y_v-y_v,\qquad y_{v_0}=0,
 \tag{FP.4}
\]
where the dot denotes the given action on $V$.
This is \texttt{Frame.\allowbreak H1c}, without quotienting by coboundaries.
Write $d h(s,t)=h(s)+s h(t)-h(st)$ and $d^0y(s)=sy-y$.
Let $Z^2_F(V)$ consist of locally constant triples
$(z,(z_v),(z_q))$ satisfying
\[
 dz=0,\qquad z|_{\Gamma_v}=dz_v,\qquad
 z|_{\Gamma_q}=dz_q .
 \tag{FP.5}
\]
Let $B^2_F(V)$ be the subspace of triples of the form
\[
       (dh,\ (h\circ\phi_v-d^0y_v)_v,\
                         (h\circ\psi_q-e_q)_q),
 \tag{FP.6}
\]
where $h$ is locally constant and each $e_q$ is a locally constant
cocycle vanishing on $N_q$.  These triples form a subspace, so their
linear span, as written in the source, is the same set.
Set
$\operatorname{Ob}_F(V)=Z^2_F(V)/(Z^2_F(V)\cap B^2_F(V))$,
as in the source.  For the actual $V=\operatorname{ad}_j$ used
below, the action is locally constant and the local maps are
continuous.  Consequently every triple in (FP.6) is locally
constant and satisfies (FP.5), so $B^2_F(V)\subset Z^2_F(V)$
in these applications.
The hypothesis is that every $\operatorname{Ob}_F(\operatorname{ad}_j)$
is finite dimensional.

Finally $B\in\mathcal C_O$ has a given finite weighted presentation
$W_B:O[[x_s:s\in\Sigma_B]]\twoheadrightarrow B$ with homogeneous
kernel and all variable images in $\mathfrak m_B$.
For each $v$ fix an \texttt{In\allowbreak H} lift $\rho_v$ over this same $B$,
whose $(a,b)$ entry has weight $w(a)-w(b)$.
We will construct a universal tuple
\[
       (i:B\to R,\ r,\ (\alpha_v)_{v\in T}),\qquad
       \alpha_v\in H(R),\quad\alpha_{v_0}=1,
 \tag{FP.7}
\]
with $r$ satisfying \texttt{In\allowbreak H} and all $L_q$, and
\[
       r(\phi_v s)\alpha_v=\alpha_v\,i(\rho_v(s)).
 \tag{FP.8}
\]
Universality means unique maps preserving $i,r$ and every $\alpha_v$;
there is no quotient by change of frame.

\subsection{Finite coefficient rings and genuine representability}

% BEGIN FP1
\begin{proposition}[FP1: complete local algebra and finite tangent separation]
\label{claim:FP1}
Rings in $\mathcal C_O$ are compact Hausdorff for their adic topology;
their ideals are closed, and their Artinian quotients are finite.
An $O$-algebra map between them is local and continuous.  Such a map
$f:A\to A'$ is surjective if
\[
       \mathfrak m_{A'}=f(\mathfrak m_A)A'+\mathfrak m_{A'}^2.
 \tag{FP.9}
\]
Under (FP.1), finitely many elements of $\Gamma$ distinguish all
lifts of $\bar r$ to $k[\epsilon]/(\epsilon^2)$.
\end{proposition}

\begin{proof}
The quotients $A/\mathfrak m_A^n$ are finite: their successive
layers are finite-dimensional vector spaces over the finite field
$k$.  Completeness identifies $A$ with their inverse limit.
This proves compactness and Hausdorffness.  Every ideal has finitely
many generators; it is the image of a finite power of the compact
space $A$ under the continuous sum-of-products map, hence is closed.
Every proper quotient is consequently complete local with the same
residue field.

If $a\in\mathfrak m_A$ had nonzero residue after applying $f$,
choose $u\in O$ whose residue is its inverse.  Then $1-ua$ is a
unit, whereas its image has zero residue, a contradiction.  Thus
$f$ is local; it preserves powers of maximal ideals and is continuous.
The residue fields agree, so the converse membership also holds.
Nakayama applied to the finite $A'$-module $\mathfrak m_{A'}$
in (FP.9) gives finitely many $a_i\in\mathfrak m_A$ whose images
generate $\mathfrak m_{A'}$.  Approximate an element of $A'$
successively by polynomials over $O$ in the $f(a_i)$.
At stage $n$, reduce the coefficients of its degree-$n$ error
modulo $\mathfrak m_{A'}$ and lift them through $O\to k$.
The resulting correction is the image of an element of
$\mathfrak m_A^n$.  Their sum converges in $A$ and maps to the
desired element.  This proves surjectivity.
The same argument gives a power-series presentation of any
$A\in\mathcal C_O$ using lifts of a basis of
$\mathfrak m_A/(\mathfrak m_OA+\mathfrak m_A^2)$.

A lift of $\bar r$ to dual numbers has the unique form
$(1+\epsilon z(s))\bar r(s)$, where $z$ is a locally constant
cocycle for the full matrix adjoint action.
Let $H=\ker\bar r$.  It is open, since $r_0$ has locally constant
reduction, and has finite index by compactness.
The restriction $z|_H$ is an additive homomorphism into the finite
$\mathbb F_\ell$-vector space of matrices over $k$.
Its coordinates have finitely many possibilities by (FP.1).
The values on finitely many coset representatives also have
finitely many possibilities, and these data determine $z$.
There are therefore finitely many dual-number lifts.
For every distinct pair choose an element where they differ.
The finite collection of chosen elements distinguishes all pairs.
\end{proof}
% END FP1

% BEGIN FP2
\begin{proposition}[FP2: the universal continuous matrix lift]
\label{claim:FP2}
There is $R_{\rm rep}\in\mathcal C_O$ with a continuous lift
$r_{\rm rep}$ of $\bar r$ such that every continuous lift over any
$A\in\mathcal C_O$ is obtained by a unique $O$-algebra map
$R_{\rm rep}\to A$.  Its construction uses only (FP.1), not a
representability theorem or an unproved deformation-ring supplier.
\end{proposition}

\begin{proof}
Choose the finite set $S\subset\Gamma$ from FP1 and lifts in $O$ of
the residual matrix entries at its elements.  Let
$P=O[[x_{s,a,b}:s\in S,\ a,b\in\iota]]$.
For a lift $r$ over a finite Artinian $A$, let $C\subset A$ be the
$O$-subalgebra generated by all its matrix entries.
This is a finite local ring with residue field $k$: the entries
have the prescribed residues, and an element with nonzero residue
has an inverse given by a finite geometric series after division
by a suitable unit of $O$.

The map from $P$ evaluating its variables at the selected entry
differences surjects onto $C$.
Indeed, if those differences did not span
$\mathfrak m_C/(\mathfrak m_OC+\mathfrak m_C^2)$, a nonzero
$k$-linear functional vanishing on them would give two distinct
$O$-algebra maps $C\to k[\epsilon]$ with the same values on the
selected entries.  The associated representation lifts would be
equal by the definition of $S$, so the maps would agree on every
matrix entry and hence on $C$, a contradiction.
Here the functional defines a derivation by subtracting a constant
in $O$ with the same residue; the result is independent of that
constant modulo $\mathfrak m_OC+\mathfrak m_C^2$.
FP1, or its finite approximation argument, now gives the asserted
surjection.

Fix $n\ge1$.  Consider the kernels of all these evaluations with
$\mathfrak m_A^n=0$.  Each contains $\mathfrak m_P^n$, so there are
only finitely many possible kernels.
They are closed under finite intersections: take the $O$-subalgebra
of a product generated by the tuples of the representation entries.
It lies in the fibre product over $k$, is a finite local ring, has
$n$th maximal-ideal power zero, and again is a quotient of $P$ by
the selected-entry argument.  Its kernel is exactly the
intersection.  Let $J_n$ be the intersection of all kernels.
Thus $P/J_n$ itself carries a representation and dominates every
lift of this Artinian level.  The map is unique, since the selected
entry differences generate this ring.
Two putative representations on the same quotient with those
selected entries agree by applying the same product argument.

The rings $R_n=P/J_n$ have compatible surjections
$R_{n+1}\to R_n$ and representations.
Universality shows
$R_{n+1}/\mathfrak m_{R_{n+1}}^n\simeq R_n$:
both have the just-proved universal property among rings with
$n$th maximal-ideal power zero.
Put $J=\bigcap_nJ_n$ and $R_{\rm rep}=P/J$.
Compactness of $P$ gives
\[
          J_n=J+\mathfrak m_P^n,\qquad
          R_{\rm rep}\simeq\varprojlim_nR_n .
 \tag{FP.10}
\]
For the first equality, the compatibility at finite levels supplies
nonempty nested compact sets of lifts in the $J_m$ of any chosen
residue class in $J_n/\mathfrak m_P^n$; their intersection supplies
the required element of $J$.
The inverse-limit representation is continuous modulo every
maximal-ideal power.  Every lift over a complete $A$ gives unique,
compatible maps to its finite quotients, hence a map
$R_{\rm rep}\to A$.  Conversely this map recovers the original
entries.  Uniqueness follows on the selected variables and then
by continuity on the power series.  This proves the proposition.
\end{proof}
% END FP2

\subsection{Weights, small extensions and the local conditions}

% BEGIN FP3
\begin{proposition}[FP3: formal torus extraction and homogeneous presentations]
\label{claim:FP3}
For the free finite-rank group $X$, near-identity torus characters
detect all weights over every coefficient ring.
Consequently a finite power-series presentation stable under all
such weighted substitutions has homogeneous kernel.
An element which transforms by the character of weight $j$ belongs
to its actual weight-$j$ part.  Every homogeneous ideal has finitely
many homogeneous generators, and its quotient modulo the maximal
ideal times that ideal has a basis represented by homogeneous
elements.
\end{proposition}

\begin{proof}
Choose a basis $X\simeq\mathbb Z^d$ and put
\[
            \tau^j=\prod_{\nu=1}^d(1+u_\nu)^{j_\nu}
                    \in S[[u_1,\ldots,u_d]]^\times .
 \tag{FP.11}
\]
These characters all reduce to one.  Any finite family of them is
linearly independent over $S$, even in positive characteristic or
with zero divisors.  Multiply a purported relation by a sufficiently
large positive character; it becomes a polynomial relation in
the $1+u_\nu$.  The invertible polynomial substitution
$T_\nu=1+u_\nu$ shows every coefficient is zero.

Let $P=O[[x_s]]\twoheadrightarrow A$ be a presentation with variable
weights $d_s$.  In $A/\mathfrak m_A^n$ every evaluation of a power
series uses only finitely many monomials of total degree less than
$n$.  Substitute $x_s\mapsto\tau^{d_s}\overline x_s$ and group
them by weight.  If this substitution kills a relation, (FP.11)
says that it kills each separate weight component.
Letting $n$ increase and using separatedness proves homogeneity of
the kernel.  Comparing a transforming element with $\tau^j$ times
itself proves the assertion about weight parts by the same argument.

The weight projection on $P$ preserves the maximal ideal and is
continuous.  Each weight part is closed in the compact ring $P$,
and its image in a quotient is therefore closed.
If $I$ is homogeneous, so is $\mathfrak m_PI$.
The finite-dimensional space $I/\mathfrak m_PI$ is discrete:
$\mathfrak m_PI$ is a closed ideal of finite index in the compact
additive group $I$.  A power series is the limit of finite sums of
its weight components, so this quotient is the direct sum of its
weight subspaces, with only finitely many of them nonzero.
Choose a basis weight by weight and lift it homogeneously to $I$.
Nakayama says that these finitely many lifts generate $I$.
The same argument applies to a quotient of one homogeneous ideal
by another, as used below.
\end{proof}
% END FP3

% BEGIN FP4
\begin{proposition}[FP4: graded small-extension lifts and matrix corrections]
\label{claim:FP4}
Let $p:A\twoheadrightarrow R$ in $\mathcal C_O$ have homogeneous
kernel $I$ with $\mathfrak m_AI=0$, and give $A,R$ compatible
finite weighted presentations.  A graded continuous family of
matrices in $H(R)$ has a graded continuous lift in $H(A)$.
Moreover a graded \texttt{In\allowbreak H} representation over $R$ satisfying
one of the specified conditions $L_q$ has a graded lift over $A$
satisfying that same condition.
These assertions hold with the original topology on its local
group, without assuming it is a topological group.
\end{proposition}

\begin{proof}
The ideal $I$ is a finite-dimensional $k$-vector space.  It is
finite and discrete in $A$; choose $n$ with
$I\cap\mathfrak m_A^n=0$.
The map on each weight part is onto its indicated image.
It has continuous set sections: on the deep subgroup
$A_j\cap\mathfrak m_A^n$ it is injective, and its image is a closed
finite-index subgroup of the target weight part.
The inverse is continuous by compactness; finitely many coset
representatives complete a section.
These sections preserve adic continuity of families.
Equivalently, homogeneous projection of lifts shows that their
deep images are the corresponding deep target parts.

Lift the entries of $h\in H(R)$ in their required weights, setting
off-block entries to zero, to obtain $h_0$ with residue $1$.
The error $E=h_0^tJh_0-J$ has entries in $I$.
Since $\ell$ is odd, $2$ is a unit in $O$.
The correction
\[
                Z=-\tfrac12J^{-1}E,\qquad h=h_0+Z
 \tag{FP.12}
\]
has entries in $I$, has the required blocks and weights, and
satisfies $Z^tJ+JZ=-E$.
The support assertion follows from (FP.2), also for $J^{-1}$:
the projections onto a block and a weight commute through $J$
to their paired block and complementary weight, and the same
identity multiplied by $J^{-1}$ proves it for the inverse.
Because $(h_0-1)I=I^2=0$, expansion of $h^tJh$ leaves exactly
$h_0^tJh_0+Z^tJ+JZ=J$.  Thus (FP.12) is a lift in $H(A)$.
All operations are continuous and preserve the asserted weights.
The same expansion shows that adding any $I$-valued element of
$\mathfrak g(A)$ to a member of $H(A)$ stays in $H(A)$.

For the assertion about $L_q$, its lifting axiom first gives some
continuous lift $r_1$ satisfying $L_q$.
Choose homogeneous continuous lifts $y_{ab}(s)$ of the entries
of its graded reduction.  Then
$\delta_{ab}(s)=r_1(s)_{ab}-y_{ab}(s)\in I$ is locally constant:
its reduction modulo the chosen $\mathfrak m_A^n$ determines it.
Write $I=\bigoplus_c I_c$ using FP3 and decompose each $\delta$
accordingly.  Only finitely many weights $c$ occur.

Here is the point that ensures that imposing weights preserves
$L_q$, rather than merely the representation equations.
Let $S=A[[u_1,\ldots,u_d]]$ and let
$c_-:A\to S$ be the inverse weighted torus substitution:
$c_-(a_j)=\tau^{-j}a_j$.
For an entry of required weight $b=w(a)-w(b')$, torus conjugation
of $c_-(r_1)$ has entry
\[
                 y_{ab'}+\sum_c\tau^{\,b-c}\delta_{ab',c}.
 \tag{FP.13}
\]
Let $\mathcal E$ be the finite set of nonzero differences $b-c$
which occur.  Form the finite square-zero extension
\[
 A^\dagger=A\oplus\bigoplus_{e\in\mathcal E}I\,e_e,\qquad
 \iota^\dagger(a,(\eta_e))=a+\sum_e\tau^e\eta_e\in S.
 \tag{FP.14}
\]
Its extra summand has the usual $A$-module action and zero product.
It belongs to $\mathcal C_O$.
The map $\iota^\dagger$ is a ring injection: $I^2=0$ proves
multiplicativity and (FP.11) proves injectivity.
The entries (FP.13) lie in its image.  Hence the matrix equations
and the \texttt{In\allowbreak H} identities hold already over $A^\dagger$,
by injectivity.  Their entries are continuous: the $y$'s are
adic-continuous and the finitely many $\delta$'s are locally
constant.  Reduction gives the original residual representation.

Closedness of $L_q$ makes it functorial under base change.
Thus $c_-(r_1)$ satisfies it, and its torus conjugate does too.
Apply the same closedness axiom to the just-constructed
$A^\dagger$-representation: its defining ideal vanishes under the
injective map $\iota^\dagger$, so that ideal is zero.
It therefore still vanishes under the projection $A^\dagger\to A$.
The resulting representation satisfies $L_q$ and has entries
$y_{ab'}+\delta_{ab',b}$, of precisely the required weight.
It lifts the given representation over $R$.
Every step used pointwise matrix identities and locally constant
reductions, so no additional topology on a local group was used.
\end{proof}
% END FP4

\subsection{The universal framed object and its minimal coordinates}

% BEGIN FP5
\begin{proposition}[FP5: the actual universal framed ring and its weights]
\label{claim:FP5}
For the original $B,\rho_v,\Gamma,r_0,F,L_q$, there is a tuple
(FP.7)--(FP.8) with the literal \texttt{Is\allowbreak Univ\allowbreak Framed} property.
It has a finite weighted presentation, all entries of $r$ and
$\alpha_v$ have weights $w(a)-w(b)$, and the map $i$ preserves
the given weight parts of $B$.
\end{proposition}

\begin{proof}
Present $R_{\rm rep}$ and $B$ by finitely many power-series
variables and take the quotient of the power-series ring in
both sets by the two old relation ideals.
This completed coproduct represents pairs of maps out of the
two rings: evaluation at maximal-ideal elements is unique and
converges, as in FP1.
Adjoin finitely many variables for the entries of $\alpha_v-1$.
Impose their off-block entries, form equations and
$\alpha_{v_0}=1$, then impose the identities (FP.8).
Impose \texttt{In\allowbreak H} for the universal representation by its form,
block and residue equations, and impose each $L_q$ by its
universally defining ideal after this cut.
There may be infinitely many indexed equations, but their ideal
is finitely generated in the Noetherian ambient power-series ring.
Quotients and the evaluation construction give the asserted
universal property, with unique maps preserving every entry and
the base map.
The quotient is proper: the residual representation, all frames
equal to $1$, and the residue map $B\to k$ satisfy the equations.
For the $L_q$ equations this follows from the given condition at
$r_0$ by closedness and base change to $k$.
Thus the constructed ring lies in $\mathcal C_O$.

To construct its weights, use the coordinate family consisting of
the images of the variables of $W_B$, all entries
$r(s)_{ab}-r_0(s)_{ab}$, and all entries $(\alpha_v-1)_{ab}$.
All lie in the maximal ideal.  Give them respectively their old
weights and $w(a)-w(b)$.  This family distinguishes maps out of
$R$, by universality and the original presentation of $B$.
Consequently its classes span the relative cotangent space over
$O$: a missing linear functional would give two maps to dual
numbers agreeing on that family but not on $R$, exactly as in FP2.
Choose a finite spanning subfamily.  FP1 gives a surjective
power-series presentation using it.

For a near-identity character $t:X\to A^\times$ after any base
change, twist each coordinate of $B$ by its weight, and conjugate
the representation and the frames by the diagonal characters
$t^{w(a)}$.  The kernel of $W_B$ is homogeneous, so the changed
base map is well defined.  Relation (FP.2) makes diagonal
conjugation preserve $H$; $r_0$ is fixed because it has entries
only at equal weights.
The local $L_q$ conditions are preserved by their torus axiom.
The graded hypothesis on $\rho_v$ makes (FP.8) compatible with
the changed base map.  Thus universality gives the corresponding
map out of $R$.
It scales each member of the coordinate family by its declared
weight.
FP3 proves that the chosen presentation has homogeneous kernel,
and that every coordinate in the full family belongs to its
asserted weight part.
The same weighted substitution argument applied to power series
in $W_B$ proves that $i$ preserves all its weight parts.
This establishes the full weighted assertion on the original tuple.
\end{proof}
% END FP5

% BEGIN FP6
\begin{proposition}[FP6: framed tangent coordinates and minimal relative relations]
\label{claim:FP6}
The ring in FP5 has a weighted presentation
\[
 P=O[[x_s\ (s\in\Sigma_B),\,z_t\ (t\in\mathcal T)]]
       \twoheadrightarrow R
 \tag{FP.15}
\]
whose old variables map through $i$ and whose new variables are
actual entries $r(s)_{ab}-r_0(s)_{ab}$ or $(\alpha_v-1)_{ab}$.
For every weight $j$,
\[
 \#\{t:\operatorname{wt}(z_t)=j\}
            =\dim_k H^1_F(\operatorname{ad}_j).
 \tag{FP.16}
\]
If $K=\ker(P\to R)$ and $K_B$ is the extended old relation ideal,
there are homogeneous $f_u\in K$ whose residues form a basis of
$K/(K_B+\mathfrak m_PK)$, and
\[
                     K=K_B+(f_u)_u .
 \tag{FP.17}
\]
Put $A=P/(K_B+\mathfrak m_PK)$ and $I=\ker(A\to R)$.
Then $I$ has the homogeneous $k$-basis
$\epsilon_u=\overline f_u$, is killed by $\mathfrak m_A$, and
\[
                     I\subset\mathfrak m_BA+\mathfrak m_A^2.
 \tag{FP.18}
\]
\end{proposition}

\begin{proof}
Maps $R\to k[\epsilon]$ whose restriction to $B$ is its residue
map are the relative tangent vectors.
Write their representations and frames uniquely as
\[
 r(s)=(1+\epsilon z(s))\bar r(s),\qquad
                   \alpha_v=1-\epsilon y_v .
 \tag{FP.19}
\]
The group and form equations give a cocycle $z$ in
$\mathfrak g(k)$.
The framing equations become $z\circ\phi_v=d^0y_v$, with
$y_{v_0}=0$.
The local tangent axiom, applied to the constant lift of $r_0$
and to this lift over $k[\epsilon]\to k$, gives exactly
$z(\psi_qn)=0$ on $N_q$.
Conversely these equations give a framed tuple over dual numbers,
including the local conditions by the same tangent axiom.
Every function used is locally constant precisely when the
dual-number lift is continuous.
Thus the relative tangent space is exactly
$H^1_F(\mathfrak g(k))$ with these signs.

Conjugation by $\bar r$ respects weights by its normalization
hypothesis, so masking the matrix entries gives
\[
 H^1_F(\mathfrak g(k))
                 =\bigoplus_jH^1_F(\operatorname{ad}_j).
 \tag{FP.20}
\]
There are only finitely many relevant weights $w(a)-w(b)$.
On the other hand the relative tangent space is the dual of
$V=\mathfrak m_R/(\mathfrak m_BR+\mathfrak m_R^2)$.
This identification follows directly from the derivation
construction in FP2.
A functional supported on the weight-$j$ cotangent coordinates
corresponds in (FP.19) to matrices in $\mathfrak g_j(k)$:
multiplication by $\bar r(s)^{-1}$ preserves those entry weights.
Thus the weight index in (FP.16) is $j$, with no inverse-weight
convention hidden by the terminology of duality.
The coordinate entries in FP5 span $V$; the old coordinates are
zero there.  Choose a basis from the remaining entries separately
in each weight.  Dualizing (FP.20) gives (FP.16).
Keeping all the specified old variables and these chosen new ones
gives (FP.15) by FP1.
Torus stability and FP3 give its homogeneous kernel; the old
kernel maps into it, with the original rename of old variables.
Its weight parts coincide with those of FP5: by FP3, in either
presentation they are exactly the elements transforming by the
specified intrinsic formal-torus character.  Thus all entries
of $r$ and $\alpha_v$, and the original map from $B$, retain
their weights in this new minimal relative presentation.

The module $K/K_B$ is finite over $P$.
Its residue quotient has a homogeneous basis by FP3.
Lift this basis to homogeneous $f_u$.
Nakayama gives (FP.17).  It also shows directly that their images
form the asserted basis of $I$, which is killed by
$\mathfrak m_A$ by its definition.
There are no missing coefficient-ring relations: $K_B$ is the
extended kernel of the original presentation of $B$ throughout.

To prove (FP.18), let $g\in K$ and expand it modulo
$\mathfrak m_P^2$ into its constant and linear terms.
The constant term belongs to $\mathfrak m_O$, since $g$ maps to
zero and all variables map into $\mathfrak m_R$.
The old linear terms lie in $\mathfrak m_BA$.
In the relative cotangent space $V$, the coefficients of the new
linear terms must all have zero residue, by the selected basis.
Those terms therefore lie in $\mathfrak m_A^2$.
This proves (FP.18) for every element of the kernel after passage
to $A$.
\end{proof}
% END FP6

\subsection{The actual obstruction classes and their detection of relations}

% BEGIN FP7
\begin{proposition}[FP7: the framed small-extension obstruction, with its correction]
\label{claim:FP7}
For $A\to R$ and its homogeneous basis $\epsilon_u$ in FP6, there
are classes
\[
        o_u\in\operatorname{Ob}_F
                      (\operatorname{ad}_{\operatorname{wt}(f_u)}).
 \tag{FP.21}
\]
For any $k$-linear map $\lambda:I\to k$, homogeneous of weight $j$,
let $A_\lambda=A/\ker\lambda$ if $\lambda\ne0$.
If $\sum_{\operatorname{wt}(f_u)=j}\lambda(\epsilon_u)o_u=0$,
the universal framed tuple over $R$ lifts to a framed tuple over
$A_\lambda$, with the specified map from $B$.
\end{proposition}

\begin{proof}
Use FP4 to lift the global matrices
$r(s)r_0(s)^{-1}$ into $H(A)$, and multiply back by $r_0(s)$.
This gives continuous graded matrices $\widetilde r(s)$.
Lift every frame by FP4, setting $\widetilde\alpha_{v_0}=1$.
For each $q$, FP4 gives a genuine graded local representation
$r_q$ lifting $r|_{\Gamma_q}$ and satisfying $L_q$.
The base map $B\to A$ is the one supplied by the old variables.

Write $\lambda_u:I\to k$ for the coefficient of $\epsilon_u$.
The following defects have entries in $I$:
\[
 \begin{split}
 D(s,t)&=(\widetilde r(s)\widetilde r(t)-\widetilde r(st))
                                      \bar r(st)^{-1},\\
 E_v(s)&=(\widetilde r(\phi_vs)\widetilde\alpha_v
             -\widetilde\alpha_v\,i(\rho_v(s)))
                                      \bar r(\phi_vs)^{-1},\\
 F_q(s)&=(\widetilde r(\psi_qs)-r_q(s))
                                      \bar r(\psi_qs)^{-1}.
 \end{split}
 \tag{FP.22}
\]
Multiplication of an $I$-valued matrix by a residual matrix is
unambiguous: any lift of that residual matrix gives the same
product, because $\mathfrak m_AI=0$.
Apply $\lambda_u$ entry by entry to (FP.22).
The results lie in $\mathfrak g(k)$.
Indeed each difference compares two matrices in the same
form-preserving coset, and the first-order form equation is
$Y^tJ+JY=0$; the block conditions are identical.
Since the defects are graded and $\epsilon_u$ has its specified
weight, their only possible nonzero entries have
$w(a)-w(b)=\operatorname{wt}(f_u)$.
For this last assertion decompose the finite $k$-space $I$ into
its actual homogeneous subspaces; their sum is direct by FP3.

Associativity of three $\widetilde r$ matrices gives
$dD=0$ after taking coefficients.
Comparison with
$\widetilde\alpha_v i(\rho_v)\widetilde\alpha_v^{-1}$
gives $D|_{\Gamma_v}=dE_v$; comparison with $r_q$ gives
$D|_{\Gamma_q}=dF_q$.
In these expansions every product of two errors vanishes and
every other factor on an error reduces to $\bar r$.
For example the global equation is exactly
\[
 sD(t,u)+D(s,tu)=D(st,u)+D(s,t).
 \tag{FP.23}
\]
The local equations are
$D(\phi_vs,\phi_vt)=E_v(s)+\phi_vs\,E_v(t)-E_v(st)$
and the same formula with $\psi_q,F_q$.
All defects are locally constant: they are adic-continuous and
take values in the finite discrete kernel $I$.
They therefore define a member of $Z^2_F(\operatorname{ad}_j)$
and its class $o_u$.

Suppose the displayed combination of classes vanishes.
The ideal $\ker\lambda$ is legitimate because $A$ acts on $I$
through $k$.  In $A_\lambda$ its remaining kernel is the
one-dimensional space $k\eta$; write $\mu:k\to k\eta$ for its
specified identification.
The combined triple in (FP.22), now multiplied by $\eta$, has
the form (FP.6).  We may take $y_{v_0}=0$:
replace $h$ by $h-d^0y_{v_0}$, every $y_v$ by $y_v-y_{v_0}$,
and $e_q$ by $e_q-d^0y_{v_0}|_{\Gamma_q}$.
The last cocycle still vanishes on $N_q$ because
$\bar r(\psi_qN_q)=1$.
This changes none of the three defects.

The corrected matrices and frames are
\[
 \begin{split}
 r'(s)&=\widetilde r(s)-\mu(h(s)\bar r(s)),\\
 \alpha'_v&=\widetilde\alpha_v+\mu(y_v).
 \end{split}
 \tag{FP.24}
\]
Here and below tildes have been mapped into $A_\lambda$.
Expanding the first product shows that its error is
$\mu(D-dh)\bar r(st)=0$, so $r'$ is multiplicative.
It preserves \texttt{In\allowbreak H} by the last small Lie correction in
FP4; in particular its matrices are invertible, and multiplicativity
then gives $r'(1)=1$.
The framing error is
$\mu(E_v-h\circ\phi_v+d^0y_v)\bar r(\phi_vs)=0$.
Thus (FP.8) holds and the distinguished frame remains $1$.
The difference between $r'|_{\Gamma_q}$ and the image of $r_q$
is $-\mu(e_q(s)\bar r(\psi_qs))$.
It vanishes on $N_q$, so the local tangent axiom gives $L_q(r')$.
Continuity is preserved because all correcting cochains are
locally constant.
All matrices and frames return to the original ones over $R$.
This is the required actual framed lift, on the same base map.
\end{proof}
% END FP7

% BEGIN FP8
\begin{theorem}[FP8: the literal \texttt{Lemma\allowbreak P} and its exact weight counts]
\label{claim:FP8}
For all the data and hypotheses stated at the beginning of this
section, the tuple and presentation in FP5--FP6 can be chosen with
finite equation index set $\mathcal K$ such that, at every $j\in X$,
\[
 \begin{split}
 \#\{t:\operatorname{wt}(z_t)=j\}
      &=\dim_kH^1_F(\operatorname{ad}_j),\\
 \#\{u:\operatorname{wt}(f_u)=j\}
      &=\dim_k\operatorname{Ob}_F(\operatorname{ad}_j).
 \end{split}
 \tag{FP.25}
\]
The old variables, their weights and their original relations are
preserved; the new variables are the actual representation or frame
entries specified in FP6.  The kernel is exactly (FP.17), and every
entry of $r,\alpha_v$ has its required weight.
These are all fields of the source's \texttt{Is\allowbreak Univ\allowbreak Framed},
\texttt{Is\allowbreak Rel\allowbreak Pres} and literal \texttt{Lemma\allowbreak P} conclusion.
There is no remaining bespoke supplier in this general implication.
\end{theorem}

\begin{proof}
We prove that the classes $o_u$ of each weight in FP7 are linearly
independent.  A nontrivial linear dependence defines a nonzero
functional $\lambda:I\to k$, zero on every other weight, because
the $\epsilon_u$ form a basis.
FP7 gives a framed lift over $A_\lambda$.
Universality of $R$ supplies an $O$-algebra map
$s:R\to A_\lambda$ respecting the original $B$-map, the
representation and all frames.
Its composite with $p_\lambda:A_\lambda\to R$ preserves the
universal tuple, so uniqueness implies $p_\lambda s=1_R$.

The nonzero kernel $k\eta$ of $p_\lambda$ is killed by the maximal
ideal.  However (FP.18) implies
\[
       k\eta\subset\mathfrak m_BA_\lambda+\mathfrak m_{A_\lambda}^2.
 \tag{FP.26}
\]
For any $x\in\mathfrak m_{A_\lambda}$, the difference
$x-s(p_\lambda x)$ belongs to this kernel.
Since $s$ agrees with the specified map on $B$,
the first summand on the right of (FP.26) lies in
$s(\mathfrak m_R)A_\lambda$.
It follows that
\[
 \mathfrak m_{A_\lambda}
          =s(\mathfrak m_R)A_\lambda+\mathfrak m_{A_\lambda}^2.
 \tag{FP.27}
\]
FP1 makes $s$ surjective.  A surjective section of $p_\lambda$
makes $p_\lambda$ injective as well, contradicting its nonzero
kernel.  No nontrivial dependence was possible.  Hence
\[
   \#\{u:\operatorname{wt}(f_u)=j\}
        \le\dim_k\operatorname{Ob}_F(\operatorname{ad}_j).
 \tag{FP.28}
\]
This proves the actual relation bound; no dimension bound was
imported as a deformation-theoretic theorem.

Only the finite set $\{w(a)-w(b):a,b\in\iota\}$ can contribute:
outside it $\mathfrak g_j(k)=0$, and all its cochain spaces and
obstruction spaces are zero.
By hypothesis the finitely many remaining obstruction dimensions
are finite.  Adjoin enough equations equal to zero at each weight
to turn (FP.28) into equality.
Zero belongs to every weight part, and adjoining zero equations
does not change (FP.17).
The generator equality is already (FP.16).
The disjoint union $\Sigma_B\sqcup\mathcal T$ is literally the
old/new variable bijection required by \texttt{Is\allowbreak Rel\allowbreak Pres}; its
rename sends the old presentation to the original map $i$.
FP5 proves universality and all representation/frame weights.
FP6 identifies each new variable as an actual entry of the
required weight.  These observations verify every outgoing field
of \texttt{Lemma\allowbreak P}, including exact counts rather than inequalities.
\end{proof}
% END FP8

\paragraph{Classical boundary and downstream scope.}
The basic algebra used here is that finite formal power-series
rings over Noetherian rings are Noetherian, and Nakayama's lemma;
see the Stacks Project,
\href{https://stacks.math.columbia.edu/tag/0306}{Tag 0306}
and \href{https://stacks.math.columbia.edu/tag/00DV}{Tag 00DV}.
Finite-dimensional linear algebra, the universal property of a
power-series evaluation, and compactness of inverse limits of
finite spaces are used in their elementary forms.
The finite-level universal representation, closedness of the
required ideals, formal torus extraction, local graded lift,
matrix obstruction correction and relation-detection argument
have all been proved in this section.
No Mazur or Schlessinger representability assertion, general
deformation-obstruction bound, or arithmetic automorphy theorem
is a retained leaf.
The theorem still requires the original $\Phi_\ell(\Gamma)$,
all five local-condition axioms, the local base-point conditions,
the finite framed obstruction spaces, and the original graded
local data over $B$.  Supplying these for a later arithmetic
\texttt{Raw} datum, and obtaining its \texttt{Model}, remain
separate tasks.


\section{Original-field comparison and the complete PointHV assembly}
\label{sec:point-hv-assembly}

\status{HV1--HV10 prove the entire original-field assembly of
\texttt{Point\allowbreak HV}.  The two auxiliary count suppliers
\texttt{Count\allowbreak Vs} and \texttt{Count\allowbreak IV} are proved here; \texttt{Count\allowbreak Z}
is already supplied by CZ6.  The precise remaining count premise is
\texttt{Count\allowbreak I}$\wedge$\texttt{Count\allowbreak II}$\wedge$\texttt{Count\allowbreak III}
$\wedge$\texttt{Count\allowbreak V}.  No character comparison, layer descent,
case partition, finite-prime summation or uniformity-of-constants
premise remains in this implication.  In particular, the comparison
does not discard primes dividing the original field degree.
The arithmetic theorems that supply the remaining four counts are
separate obligations.}

\subsection{Notation and the exact interfaces}

Let $\chi$ be the primitive Dirichlet character of the source, put
$q=\operatorname{cond}(\chi)$ and $o=\operatorname{ord}(\chi)$,
and let $h$ be a host with exponent $s\ge2$ and prime pair $q_1,d$.
Thus $q_1\equiv3\pmod8$, $d\equiv7\pmod{16}$,
$(q_1d,q)=1$, and
\[
 N=\operatorname{lcm}(2^{s+2},q)q_1d.
 \tag{HV.1}
\]
Inside $(\mathbb Z/N)^\times$ let $H$ be the simultaneous kernel
of the four conditions: reduction modulo $2^{s+2}$ lies in
$\langle-1\rangle$; the $\chi$ value is $1$; and the reductions
modulo $q_1$ and $d$ are squares.  Write $n=[(\mathbb Z/N)^\times:H]$.
The finite exceptional set is the literal source set
\[
 \mathcal B_\chi=\{3\}\ \cup\
 \{\text{odd prime divisors of }o q p_K h_* p_1\},
 \tag{HV.2}
\]
where $p_K=p_K(\chi)$, $h_*=h_*(\chi)$ and $p_1=p_1(\chi)$
are the fixed quantities of \texttt{Stmt\allowbreak Fields}.
The condition \texttt{Big}$(t)$ says that $q_1,d>2t+1$ and that
each of $q,o,p_K,h_*,p_1$ is smaller than both $q_1$ and $d$.
All these constants have their original source meanings.

For an odd prime $\ell$, let $O$ be any source coefficient DVR
inside the fixed $\ell$-adic algebraic closure, finite over
$\mathbb Z_\ell$, and put
\[
 a_O=\operatorname{length}_O(O/\ell O)>0,
 \qquad A_O=O\otimes_{\mathbb Z_\ell}
                         (\mathbb Q_\ell/\mathbb Z_\ell).
 \tag{HV.3}
\]
For a character $\psi$ modulo $N$ trivial on $H$, the relevant
representation on $A_O$ is
\[
                 \rho_{e,\psi}=\psi^{-1}\chi_{\rm cyc}^{-e}.
 \tag{HV.4}
\]
Let $S'_\ell(\psi,e)$ be its continuous $H^1$ classes over
$\mathbb Q$ which are unramified outside $\ell N$ and locally
zero at every prime dividing $\ell N$.  Its length and normalized
length are exactly
\[
 L_\ell(\psi,e)=\mathtt{Vlen}(\psi,e),\qquad
 v_\ell(\psi,e)=L_\ell(\psi,e)/a_O=\mathtt{Vn}(\psi,e).
 \tag{HV.5}
\]
The quotient notation for the length is used only after finiteness
has been proved; before then we use extended lengths.  Let $X_e$
be the characters trivial on $H$ with $\psi(-1)=(-1)^e$.
The assertion \texttt{VLe}$(\psi,e,B)$ is precisely that
$L_\ell(\psi,e)$ is finite and $v_\ell(\psi,e)\log\ell\le B$.

For a number field $F$ we write $s_\ell(F,e)$ for the strict
negative-twist Selmer length (TW.1).  In particular
\texttt{VVat}$(\ell,F,t,B)$ bounds the sum of these Selmer
lengths for $1\le e<t$; it is not a different normalization
of the compact $H^2$ length.  The layer condition
\texttt{Layer\allowbreak Le}$(\chi,\ell,j,t,B)$ requires this bound on
every model of the $j$th cyclotomic $\mathbb Z_\ell$-layer of
$\mathbb Q(\mu_{\ell q})$.  This is the original universal
quantifier over \texttt{Is\allowbreak Ab\allowbreak Field} and \texttt{Is\allowbreak Layer}.

The only classical arithmetic inputs below are the elementary
cyclotomic local-field description, the abelian
conductor--discriminant formula already matched in TW-CD, and
the previously proved comparison TW3 and bound KB8.  All
restriction errors and character decompositions are proved here.
For the cyclotomic local description one can use the usual
Eisenstein polynomial $\Phi_{\ell^r}(1+X)$; the valuation argument
needed below is included in HV2.  The cohomological exact
sequence used below is restriction--inflation, with its
transgression; its cochain construction is recalled explicitly.
For the finite cyclotomic and local-field facts see Milne,
\emph{Algebraic Number Theory}, v3.08, Proposition 6.2 and
Proposition 7.55,
\url{https://www.jmilne.org/math/CourseNotes/ANT.pdf}.
The conductor--discriminant formula is equation (4.2.1) in
Sharifi, \emph{Iwasawa Theory}, Chapter 4,
\url{https://math.ucla.edu/~sharifi/notes/iwasawa-ch04.html}.

\paragraph{HV1: the original host field and its elementary character bounds.}
\label{claim:HV1}
Every original field $E$ satisfying \texttt{Is\allowbreak Chi\allowbreak Field} for $h$
has, in any algebraic closure of $\mathbb Q$, image
$\mathbb Q(\mu_N)^H$.  In particular it is abelian,
unramified outside $N$, and has degree $n$, with
\[
 n\mid2^{s+2}o,\qquad n\le N,\qquad16\mid N,\qquad
 1\le o\le q,\qquad |X_e|\le N.
 \tag{HV.6}
\]
Write $\psi_p$ for the CRT component of $\psi$ supported at $p$.
For odd $p$, $\operatorname{ord}(\psi_p)\mid2o$; if $p\nmid q$
then $\psi_p^2=1$.  If
$c_2=v_2(\operatorname{cond}(\psi_2))$, then
\[
                   \operatorname{ord}(\psi)\le2o\,2^{c_2}.
 \tag{HV.7}
\]
Write $\psi=\kappa\theta$ for its $\ell$-power and prime-to-$\ell$
order parts, and put
$\chi_P=\theta/\theta_\ell$ and $\theta_1=\theta/\theta_2$,
as in the source.  If $\ell$ is odd, $\ell\ne q_1,d$ and
$\operatorname{ord}(\chi_P)$ is odd, then
\[
          \psi_2=1,\qquad\operatorname{cond}(\psi)\mid q.
 \tag{HV.8}
\]
The same conclusions hold if $\operatorname{ord}(\psi)$ is odd.
Consequently the number of host characters satisfying (HV.8)
is at most $q$, in any coefficient ring.
Under \texttt{Big}$(t)$, every $\ell\in\mathcal B_\chi$ is
smaller than $q_1,d$ and hence smaller than $N$.  If $\ell$ is
odd and outside $\mathcal B_\chi$, then
\[
 \ell\ge5,\quad\ell\nmid oqn;
 \quad 1\le e<t, \ell-1\mid2e
       \ \Longrightarrow\ e\ge2,\ \ell\le2t,\ \ell\nmid N.
 \tag{HV.9}
\]

\begin{proof}
The defining embedding of \texttt{Is\allowbreak Chi\allowbreak Field} specifies exactly
the automorphisms which fix its image: the sign condition at
$2^{s+2}$, the kernel of $\chi$, and the two quadratic-residue
conditions.  For a unit modulo an odd prime, being a square is
equivalent to Legendre symbol $1$.  If the initially supplied
cyclotomic level is $N'$, its divisibilities give $N\mid N'$.
Its fixing group is the preimage of $H$ under reduction modulo
$N$, since all four defining conditions depend only on those
residues.  In particular it contains the kernel of the restriction
to $\mathbb Q(\mu_N)$, so the image of $E$ already lies in
that smaller field.  Cyclotomic restriction is surjective;
its fixing group there is exactly $H$.
The Galois correspondence gives the asserted field, with no
change to the original $E$.
Different embeddings are conjugate; since the cyclotomic Galois
group is abelian they have this same image.

The quotient by $H$ injects into the product of the quotient
of $(\mathbb Z/2^{s+2})^\times$ by $\langle-1\rangle$, the
image of $\chi$, and two groups of order two.  These have orders
$2^s,o,2,2$, proving the divisibility in (HV.6).
Also $n\le\varphi(N)\le N$, and $o\mid\varphi(q)$ gives
$o\le q$.  A finite abelian group of order $n$ has at most $n$
characters in a characteristic-zero domain: decompose it into
cyclic groups and use the bound of $r$ roots for $X^r-1$.
This proves the bound on $X_e$ and also on the union $X_0\cup X_1$.

On an odd CRT factor the first quotient is trivial, so its
exponent divides $2o$.  If $p\nmid q$, only the quadratic factors
remain.  The component $\psi_2$ factors through the units modulo
$2^{c_2}$, and its order divides $2^{c_2}$.  The product of all
odd components has order dividing $2o$, which proves (HV.7).
There are no components outside the prime divisors of $N$.

For (HV.8), the $2$-component has $2$-power order.  Since $\ell$
is odd, it is unchanged on removing the $\ell$-power order part,
and also on removing the component at $\ell$.  Its order
therefore divides the odd order of $\chi_P$, and it is trivial.
At $q_1,d$, which differ from $\ell$, the component of $\psi$
is a product of an odd-order component of $\chi_P$ and an
$\ell$-power order component of $\kappa$.  It has odd order;
it also has order dividing two, and is trivial.  At every other
odd prime the exponent of that prime in $N$ is exactly its
exponent in $q$.  This proves the conductor divisibility.
For an odd-order $\psi$ the same argument applies directly.
A character of conductor dividing $q$ factors uniquely modulo
$q$, giving at most $\varphi(q)\le q$ possibilities.

All factors defining (HV.2) are positive.  Every prime divisor
is no larger than one of these factors; \texttt{Big} therefore
bounds it by $q_1,d$.  The inserted prime $3$ also satisfies the
bound.  Outside this set the assertions about $o,q,n$ follow
from (HV.6), and $\ell\ge5$.  If $\ell-1\mid2e$, then
$\ell-1\le2e$, so $\ell\le2t$ and $e\ge2$.
The only possible prime divisors of $N$ are $2$, primes of $q$,
$q_1$ and $d$.  The size condition excludes the last two and
the earlier conclusions exclude the first two.  This proves
(HV.9).
\end{proof}

\paragraph{HV2: quantitative restriction without inverting the degree.}
\label{claim:HV2}
Let $F/\mathbb Q$ be finite Galois of degree $b$, and suppose a
finite-order character $\psi$ is trivial on $G_F$.
For $e\ge1$ and odd $\ell$, the restriction kernel on global
$H^1(\mathbb Q,A_O\rho_{e,\psi})$ has normalized length at most
$3b^2e$.  Whenever restriction carries a Selmer group $S$ into
the strict negative-twist group over $F$, one consequently has
\[
       \frac{\operatorname{length}_O S}{a_O}
                         \le s_\ell(F,e)+3b^2e.
 \tag{HV.10}
\]
Finiteness on the right implies finiteness on the left.
The assertion applies to the original algebraic closure and to
every original choice of local places.

\begin{proof}
For a finite extension $K/\mathbb Q_\ell$ of degree at most $b$,
the image $J$ of its cyclotomic character has finite index
$a\le b$ in $\mathbb Z_\ell^\times$.  Indeed the cyclotomic
character of $\mathbb Q_\ell$ is surjective, and the map from
the cosets of $G_K$ to the cosets of $J$ is surjective.
Surjectivity at the base follows from the Eisenstein
polynomials $\Phi_{\ell^r}(1+X)$, which give degree
$\varphi(\ell^r)$ for the local cyclotomic extensions.
The full image is compact and hence closed; its surjectivity
at every finite level therefore gives surjectivity onto
$\mathbb Z_\ell^\times$.
Since the quotient by $J$ has order $a$, $(1+\ell)^a\in J$.
The invariants of $A_O(-e)$ over $K$ are therefore killed by
$(1+\ell)^{ae}-1$.  The binomial formula, successively applied
to an exponent prime to $\ell$ and to powers of $\ell$, gives
\[
 \frac{\operatorname{length}_O A_O(-e)^{G_K}}{a_O}
 \le v_\ell((1+\ell)^{ae}-1)
 =1+v_\ell(a)+v_\ell(e)\le3be.
 \tag{HV.11}
\]
Here the $u$-torsion of $A_O$ has length $a_Ov_\ell(u)$
for $0\ne u\in\mathbb Z_\ell$.  Global invariants over $F$
are a submodule of these invariants at any one place above
$\ell$, so the same bound holds for
$W=A_O(-e)^{G_F}$.

Put $G=\operatorname{Gal}(F/\mathbb Q)$.
The restriction--inflation sequence identifies the restriction
kernel with $H^1(G,W)$, where the quotient action on $W$ is
the one inherited from $\rho_{e,\psi}$, including its
$\psi^{-1}$ factor.  Its underlying $O$-module is the invariant
module just bounded.  In inhomogeneous cochains the degree-one
group consists of $b$ copies of $W$, so
$\operatorname{length}H^1(G,W)\le b\operatorname{length}W$.
This proves the kernel bound, even when $\ell\mid b$.
Restriction of a cocycle preserves each local coboundary and
unramified-inertia condition.  Tensoring the strict group over
$F$ with $O$ multiplies its $\mathbb Z_\ell$-length by $a_O$:
the ring is finite free, so cochains, coboundaries and all local
kernels commute with this tensor product.  These facts prove
(HV.10), including finiteness.

Conjugating a selected place or transporting an algebraic closure
conjugates the cochains and their coboundary witnesses.  The
restriction diagram commutes with this operation.  Thus this
proof concerns the given $F$ and given local data, rather than
only one convenient model.
\end{proof}

\paragraph{HV3: integral decomposition into characters.}
\label{claim:HV3}
Let $G$ be finite abelian, let $O$ contain its character values,
and let $M$ be any $O[G]$-module, initially of possibly infinite
length.  For $c\in G$ of order two, odd residue characteristic,
and $\epsilon=\pm1$, put
$M_\psi=\{x:gx=\psi(g)x\text{ for every }g\in G\}$.
Then
\[
 \operatorname{length}_O M^{c=\epsilon}
 \le\sum_{\psi(c)=\epsilon}\operatorname{length}_O M_\psi.
 \tag{HV.12}
\]
In particular, finiteness of the displayed eigenspaces proves
finiteness of the left side.  No condition on divisibility of
$|G|$ by $\ell$ is imposed.

\begin{proof}
For commuting endomorphisms $f,g$, the map
$g:\ker(fg)\to\ker f$ has kernel $\ker g$.
Additivity and monotonicity of length give
$\operatorname{length}\ker(fg)\le
\operatorname{length}\ker f+\operatorname{length}\ker g$.
If $T^r=1$, factor
$T^r-1=\prod_{\zeta^r=1}(T-\zeta)$ and iterate this inequality.
Apply it successively to cyclic generators in a decomposition
$G=\prod_i C_{r_i}$, starting on $M^{c=\epsilon}$.
All the kernels are stable under the remaining generators.
Every final intersection is one of the simultaneous
eigenspaces $M_\psi\cap M^{c=\epsilon}$.
If $\psi(c)\ne\epsilon$, their difference is the unit $\pm2$,
so that intersection is zero.  This proves (HV.12).
The same successive kernel exact sequences prove finiteness
when all final terms are finite; no arithmetic with infinite
lengths is needed to infer it.
\end{proof}

\paragraph{HV4: the comparison on every original CM field.}
\label{claim:HV4}
Let $E$ be an original CM field of host $h$, and let $D$ be
any of its original local-data records at an odd prime $\ell$.
Fix $e\ge1$.
Choose a source coefficient ring $O$ containing the values of
all characters of $G=\operatorname{Gal}(E/\mathbb Q)$.
If all $L_\ell(\psi,e)$ for $\psi\in X_e$ are finite, then
$\mathtt{plusLen}_D(e)$ is finite and
\[
 \mathtt{plusLen}_D(e)
 \le\sum_{\psi\in X_e}
                 \bigl(v_\ell(\psi,e)+(Ne)^6\bigr).
 \tag{HV.13}
\]
If $\ell\nmid n$, the error can be omitted:
\[
             \mathtt{plusLen}_D(e)
                     \le\sum_{\psi\in X_e}v_\ell(\psi,e).
 \tag{HV.14}
\]
The source's normalization, the original field, and all its
local conditions are retained.

\begin{proof}
Use HV1 to embed the actual $E$ into the fixed rational
algebraic closure.  Let $B$ be its strict Selmer module with
coefficient $A_O(-e)$.  Normality gives the natural $G$-action
on $B$.  If $c$ is complex conjugation, the defining plus
condition, after this identification, is the eigenspace
$B^{c=(-1)^e}$.  To check the sign, the natural action on a
cocycle $z$ is $g\mapsto c\,z(c^{-1}gc)$, and $c$ acts on
the coefficient as $(-1)^e$.  Thus the untwisted conjugation
invariance defining the plus part is precisely the stated
eigenvalue.  Finite free coefficient extension and conjugacy
of the closures give the exact equality
\[
 a_O\mathtt{plusLen}_D(e)
                    =\operatorname{length}_O B^{c=(-1)^e}.
 \tag{HV.15}
\]
This is also the explicit \texttt{part\allowbreak Model} identification.
Changing a local place changes the cocycle and its local
coboundary witness by conjugacy; hence it changes no kernel
in this equality.
For the precise quotient used in $D$, a pure negative-twist
cocycle unramified away from $\ell$ vanishes on those inertia
groups: their action on the coefficient is trivial, so their
coboundaries are zero.  The cocycle equation then makes it
vanish on their conjugates and on the closed normal subgroup
they generate.  It therefore descends to the actual
$G_{E,S_\ell}$ quotient.  Inflation is its inverse on these
classes and preserves every selected local kernel.  This
identifies the absolute-cochain description with the one in $D$.

For a character $\psi$ of $G$, twist the coefficient action by
$\psi^{-1}$.  Over $E$ this does not change the coefficient,
but on $H^1(E,A_O(-e))$ it multiplies the $G$-action by
$\psi^{-1}$.  Thus $B_\psi$ becomes an invariant submodule.
Restriction--inflation gives an exact sequence
\[
 H^1(\mathbb Q,A_O\rho_{e,\psi})
 \xrightarrow{\mathrm{res}}H^1(E,A_O(-e))^G
 \longrightarrow H^2(G,W_\psi).
 \tag{HV.16}
\]
Here the action on the middle term is the twisted one just
specified, and $W_\psi$ has underlying module
$A_O(-e)^{G_E}$ with the quotient action inherited from
$\rho_{e,\psi}$.  In particular this action retains the
$\psi^{-1}$ factor; it need not be the untwisted action.
The sequence can be constructed directly: choose representatives
of $G_{\mathbb Q}/G_E$, express invariance of a cocycle by
coboundary witnesses on these representatives, and compare the
witness for a product with the two individual witnesses.
The resulting degree-two cocycle takes values in
$A_O(-e)^{G_E}$.  Altering the witnesses changes it by a
degree-two coboundary, and its being zero is exactly the
condition that the cocycle extend to $G_{\mathbb Q}$.
All quotient choices are finite, so this is the continuous
restriction--inflation sequence as well.

Let $P$ be the preimage of $B_\psi$ under restriction in
(HV.16).  At a prime $p\ne\ell$, the pure negative twist is
unramified over $E_w$, and
\[
 H^1_{\rm unr}(E_w,A_O(-e))
       =A_O/(p^{-fe}-1)A_O=0,
 \tag{HV.17}
\]
where $f$ is the residue degree.  The scalar is nonzero for
$e\ge1$ and acts surjectively on the divisible module $A_O$.
Thus every strict Selmer class in $B_\psi$ is locally zero
at every finite place of $E$.
If $p\nmid\ell N$, the extension $E_w/\mathbb Q_p$ is
unramified and its inertia is the same subgroup as rational
inertia.  Any lift in $P$ is consequently unramified there.
For $v\mid\ell N$, the localization of such a lift lies in
the restriction kernel
\[
 K_v=H^1(G_v,A_O(-e)^{G_{E_w}}),
       \qquad G_v=\operatorname{Gal}(E_w/\mathbb Q_v).
 \tag{HV.18}
\]
Here too the quotient action on the displayed invariant
module is the one inherited from $\rho_{e,\psi}$.
The kernel of $P\to\prod_{v\mid\ell N}K_v$ is exactly
$S'_\ell(\psi,e)$.  Conversely every class in that rational
Selmer group restricts into $B_\psi$, since local zero at
primes of $N$ gives the unramified condition over $E$ there.
This proves the equality of kernels, including every original
local condition.

The global invariant module in (HV.16) has normalized length
at most $3ne$ by (HV.11).  Degree-two inhomogeneous cochains
have $n^2$ copies, so the cokernel contribution is at most
$3n^3e$.  At $\ell$, (HV.11) and degree-one cochains bound
(HV.18) by $3n^2e$.  At a prime $p\mid N$, $p\ne\ell$,
the normalized invariant length is
\[
 v_\ell(p^{fe}-1)\le fe\log_\ell p\le neN,
 \tag{HV.19}
\]
since $f\le n$, $p\le N$ and $\ell\ge3$.
The corresponding degree-one cochains contribute at most
$n^2eN$.  There are at most $N$ such primes.  We have proved
\[
 \frac{\operatorname{length}_O B_\psi}{a_O}
 \le v_\ell(\psi,e)+3n^3e+3n^2e+Nn^2eN
 \le v_\ell(\psi,e)+7N^4e
 \le v_\ell(\psi,e)+(Ne)^6.
 \tag{HV.20}
\]
Finiteness follows before taking these ordinary lengths: the
local kernels and the degree-two obstruction are finite, and
$S'_\ell(\psi,e)$ is finite by hypothesis.  The last inequality
uses $N\ge16$ and $e\ge1$.
If $\ell\nmid n$, positive-degree cohomology of $G$ and all
$G_v$ is zero on these $\ell$-primary modules.  The averaging
homotopy divides by their orders, which are units in $O$;
equivalently restriction followed by transfer gives the same
vanishing.  Thus every error in (HV.20) is zero in this case.
Combine (HV.12), (HV.15) and (HV.20) to obtain (HV.13)--(HV.14).

There always exists such a coefficient ring: adjoin the finitely
many roots of unity of order dividing $n$ inside the specified
$\ell$-adic algebraic closure and take its ring of integers.
If an initial coefficient ring is prescribed, use the compositum.
For a finite extension $O\subset O'$ of DVRs the length and
$a_O$ are both multiplied by the relative ramification index.
Continuous cochains and all local kernels commute with this
finite free extension, so normalized lengths and their
finiteness are unchanged.  This also verifies compatibility
with the source's unrestricted coefficient quantifiers.
\end{proof}

\subsection{The two auxiliary counts are complete suppliers}

\paragraph{HV5: the complete small-conductor count \texttt{Count\allowbreak Vs}.}
\label{claim:HV5}
The literal assertion \texttt{Count\allowbreak Vs} holds.  Its exponent can
be chosen as one natural number depending only on $\chi$, before
the host, $t$, prime, coefficient ring, character and algebraic
closure are chosen.

\begin{proof}
In this case $\ell\in\mathcal B_\chi$, $e=1$, and
\[
 2^{c_2}<R_1(\ell,\psi)
      =16o(\ell q)^2 4^{n_0},\qquad
 n_0=v_2(\ell^{f_0}-1),\quad
 f_0=\operatorname{ord}_{\operatorname{lcm}(4,
                      \operatorname{ord}(\theta_1))}(\ell).
 \tag{HV.21}
\]
The order of $\theta_1$ divides $2o$ by HV1 and is prime to
$\ell$, by construction of the prime-to-$\ell$ part.
Consequently the modulus $m=\operatorname{lcm}(4,
\operatorname{ord}(\theta_1))$ divides $4o$, and $\ell$ is
a unit modulo $m$.  Thus $f_0$ is positive,
$f_0\mid\varphi(m)$ and $\varphi(m)\mid\varphi(4o)$.
The last divisibility follows directly from the prime-power
formula for $\varphi$ whenever $m\mid4o$.
It follows that
\[
 2^{n_0}\le\ell^{f_0}-1<\ell^{\varphi(4o)},\qquad
 R_1\le16o(\ell q)^2\ell^{2\varphi(4o)}.
 \tag{HV.22}
\]
Let $M_\chi=\max\mathcal B_\chi$; this is defined since
$3\in\mathcal B_\chi$.  Choose the fixed integer
\[
 b_* = 2o\,[16o(M_\chi q)^2M_\chi^{2\varphi(4o)}]M_\chi.
 \tag{HV.23}
\]
Equations (HV.7), (HV.21) and (HV.22) give
$\ell\operatorname{ord}(\psi)\le b_*$.

Inside the given algebraic closure, take the abelian field
fixed by the intersection of the kernel of $\psi$ and the
kernel of the mod-$\ell$ cyclotomic character.  Denote it by
$F$.  It contains $\mu_\ell$, has degree
$b\le(\ell-1)\operatorname{ord}(\psi)\le b_*$, and lies
in $\mathbb Q(\mu_{\ell N})$.
Restriction carries $S'_\ell(\psi,1)$ into the strict group
over $F$.  Indeed at primes dividing $N$ the rational class
is already locally zero, and at every other prime away from
$\ell$ it is unramified; these properties persist on restriction.
HV2 therefore gives
\[
                  v_\ell(\psi,1)\le s_\ell(F,1)+3b_*^2.
 \tag{HV.24}
\]
The field is totally complex, since $\ell$ is odd and
$\mu_\ell\subset F$.  Apply TW3's strict Selmer comparison
and the full original-field theorem KB8 to this $F$.
With $\Delta_F=|\operatorname{disc}(F)|$ they give finiteness
and
\[
 s_\ell(F,1)\log\ell
 \le\log\bigl|H^2(G_{F,S_\ell},\mathbb Z_\ell(2))\bigr|
 \le\exp\bigl((2b)^{31}\log(2\Delta_F)\bigr).
 \tag{HV.25}
\]
The named abelian conductor--discriminant formula gives
$\Delta_F\le(\ell N)^b$.  Put $x=\ell N\ge2$ and
$A=(2b_*)^{31}(b_*+1)$.  Then
$\log(2\Delta_F)\le(b_*+1)\log x$, so (HV.24)--(HV.25)
imply
\[
             v_\ell(\psi,1)\log\ell
                    \le x^A+3b_*^2\log\ell
                    \le x^A+3b_*^2x.
 \tag{HV.26}
\]
For every integer $b_*\ge1$ one has
$3b_*^2\le2^{2b_*+2}$; this follows, for example, from
$b_*\le2^{b_*}$.
Thus the last expression is at most $x^{A+2b_*+4}$.
Taking
\[
                     C_{\rm Vs}=A+2b_*+4
 \tag{HV.27}
\]
proves the source's \texttt{VLe}$(\ell Ne)^{C_{\rm Vs}}$
in every small-conductor case.  The proof established the
finiteness part first.  Every quantity in (HV.23), (HV.27)
depends only on $\chi$.  This proves the literal \texttt{Count\allowbreak Vs},
with a permissible larger exponent than the source construction.
\end{proof}

\paragraph{HV6: the complete layer comparison \texttt{Count\allowbreak IV}.}
\label{claim:HV6}
The literal assertion \texttt{Count\allowbreak IV} holds.  For fixed
$\chi$, $\ell\in\mathcal B_\chi$ and layer index $j$, set
$b_0=\ell q\ell^j$.  One may take
$C_4=2b_0+4$, before the host and all its remaining data.
In detail, for every $t\ge2$ with \texttt{Big}$(t)$ and every
real $B$ satisfying \texttt{Layer\allowbreak Le}$(\chi,\ell,j,t,B)$,
every host character with odd $\operatorname{ord}(\chi_P)$
has finite $L_\ell(\psi,e)$ for $1\le e<t$ and
\[
          \sum_{1\le e<t}v_\ell(\psi,e)
                                  \le B+(Nt)^{C_4}.
 \tag{HV.28}
\]
No parity condition on this character is required in this claim.

\begin{proof}
HV1 gives $\ell\ne q_1,d$ and
$\operatorname{cond}(\psi)\mid q$.
Construct $F_0=\mathbb Q(\mu_{\ell q})$ in the given closure.
It contains $\mu_\ell$ and has degree at most $\ell q$.
Its cyclotomic image is an open subgroup of
$1+\ell\mathbb Z_\ell$, so its root-adjoining tower is a
$\mathbb Z_\ell$-extension.  More explicitly, the valuation
formula from HV2 identifies this group with $\mathbb Z_\ell$;
its successive open subgroups have indices $\ell^j$.
Thus its $j$th layer $F_1$ is a root-adjoining field of exactly
that degree over $F_0$, and
$[F_1:\mathbb Q]\le b_0$.
This constructs an actual field satisfying the source's
\texttt{Is\allowbreak Ab\allowbreak Field} and \texttt{Is\allowbreak Layer} in the original closure.
The universal layer premise applies to it and gives
\[
       \sum_{1\le e<t}s_\ell(F_1,e)\le B,
       \qquad s_\ell(F_1,e)<\infty.
 \tag{HV.29}
\]
Since the conductor of $\psi$ divides $q$, it is trivial on
$G_{F_1}$.  The same local-condition argument as in HV5 shows
that rational restriction maps $S'_\ell(\psi,e)$ into the
strict Selmer group of $F_1$.  Applying HV2 and summing gives
\[
 \sum_{1\le e<t}v_\ell(\psi,e)
 \le B+3b_0^2\sum_{1\le e<t}e
 \le B+3b_0^2t^2.
 \tag{HV.30}
\]
All these groups are finite by the same restriction argument.
For $N\ge16$, $t\ge2$ and $b_0\ge1$,
$3b_0^2t^2\le(Nt)^{2b_0+4}$; the elementary estimate
$b_0\le2^{b_0}$ proves it immediately.  This is (HV.28).
We used the premise on an actual layer in the original closure;
any supplied layer model can be transported to it by its
$F_0$-isomorphism, which conjugates all cocycles and local
witnesses.  Thus no field-model restriction has entered the
conclusion.  This proves the full \texttt{Count\allowbreak IV}.
\end{proof}

\subsection{The complete partition and summation}

\paragraph{HV7: an exhaustive partition with the original predicates.}
\label{claim:HV7}
Let $\ell$ be odd, $1\le e<t$, assume \texttt{Big}$(t)$,
and let $\psi\in X_e$.  The source cases are exhaustive:
\[
\begin{array}{ll}
\mathrm{I}:&\ell\notin\mathcal B_\chi,\quad\ell-1\nmid2e;\\
\mathrm{Z}:&\ell\notin\mathcal B_\chi,\quad\ell-1\mid2e,
                                                        \quad\psi=1;\\
\mathrm{II}:&e\ge2,\quad\operatorname{ord}(\chi_P)\text{ even},
                 \quad(\ell\in\mathcal B_\chi\text{ or }\ell-1\mid2e);\\
\mathrm{III}:&\ell\notin\mathcal B_\chi,\quad\ell-1\mid2e,
                 \quad\operatorname{ord}(\psi)\text{ odd},\quad\psi\ne1;\\
\mathrm{IV}:&\ell\in\mathcal B_\chi,\quad e\ge2,
                                    \quad\operatorname{ord}(\chi_P)\text{ odd};\\
\mathrm{V}:&\ell\in\mathcal B_\chi,\quad e=1.
\end{array}
\tag{HV.31}
\]
Case V splits into the two disjoint inequalities
$R_1\le2^{c_2}$ and $2^{c_2}<R_1$, to which
\texttt{Count\allowbreak V} and the proved \texttt{Count\allowbreak Vs} apply.

\begin{proof}
Inside $\mathcal B_\chi$, first separate $e=1$; otherwise
$e\ge2$ and the positive order of $\chi_P$ is even or odd.
Outside $\mathcal B_\chi$, failure of divisibility is case I.
When divisibility holds, HV1 gives $e\ge2$, $\ell\nmid N$
and $\ell\nmid\operatorname{ord}(\psi)$: the latter follows
from $\operatorname{ord}(\psi)\mid n$ and (HV.6).
Hence $\psi$ already has prime-to-$\ell$ order, and its
$\ell$-component is trivial.  Thus $\chi_P=\psi$.
Separate the trivial character, and then even and odd order,
to obtain Z, II and III respectively.  The numerical split in
case V is the ordinary ordering of the two integers.
\end{proof}

\paragraph{HV8: every exceptional-prime sum, with the prescribed layer.}
\label{claim:HV8}
Assume the actual source statements \texttt{Count\allowbreak II} and
\texttt{Count\allowbreak V}.  Then the literal statement \texttt{Sum\allowbreak In}
holds.  Its constant for $(\chi,\ell,j)$ is chosen before
the host, the twist cutoff, the layer budget, the original
field and all its local data.

\begin{proof}
Fix $\chi$, $\ell\in\mathcal B_\chi$ and $j$.
Take the exponents $C_2,C_5$ provided by the two assumed
counts, $C_{\rm Vs}$ from HV5, and $C_4$ from HV6.
Put
\[
 C_m=\max(C_2,C_5,C_{\rm Vs}),\qquad
 C_{\rm in}=C_4+2C_m+12+\ell q.
 \tag{HV.32}
\]
Now take any subsequent host, $t\ge2$, $B\ge0$ and its
prescribed layer bound, and any original $E,\Omega,D$.
Put $M=Nt\ge32$ and choose a coefficient ring splitting
$\operatorname{Gal}(E/\mathbb Q)$, as in HV4.

Let $Y$ be all host characters for which
$\operatorname{ord}(\chi_P)$ is odd.  By HV1, $|Y|\le q$;
this is a single finite set independent of $e$.
For each $\psi\in Y$, HV6 gives
$\sum_{1\le e<t}v_\ell(\psi,e)\le B+M^{C_4}$.
For every $\psi\in X_e\setminus Y$, HV7 gives case II
when $e\ge2$, or one of the two parts of case V when $e=1$.
The assumed counts and HV5 therefore give
\[
          v_\ell(\psi,e)\log\ell
                \le(\ell Ne)^{C_m}\le M^{2C_m},
 \tag{HV.33}
\]
because $\ell<N$ and $e<t$.
Every character length required by HV4 is finite: use HV6
for $Y$ and these pointwise counts for its complement.
Using $|X_e|\le N$, nonnegativity, and interchanging only
finite sums yields
\[
 \log\ell\sum_{1\le e<t}\sum_{\psi\in X_e}v_\ell(\psi,e)
 \le q\log\ell\,(B+M^{C_4})+M^{2C_m+1}.
 \tag{HV.34}
\]
The total error in (HV.13) is bounded by
\[
 \log\ell\sum_{1\le e<t}|X_e|(Ne)^6
             \le\ell Nt\,M^6\le M^8.
 \tag{HV.35}
\]
Here $\log\ell\le\ell\le N$, and $q\le N$.
Combining (HV.13), (HV.34) and (HV.35) gives
\[
 \log\ell\sum_{1\le e<t}\mathtt{plusLen}_D(e)
 \le\ell q B+M^{C_4+2}+M^{2C_m+1}+M^8
 \le M^{C_{\rm in}}+C_{\rm in}B.
 \tag{HV.36}
\]
For the last inequality each of the three exponents is at
most $C_4+2C_m+8$, and $3M^a\le M^{a+2}$ for $M\ge2$;
also $C_{\rm in}\ge\ell q$.
All plus lengths are finite by HV4.
This is exactly \texttt{Plus\allowbreak Le} for the singleton $\{\ell\}$
and proves \texttt{Sum\allowbreak In}, with every quantifier in its
original order.
\end{proof}

\paragraph{HV9: the sum over all nonexceptional odd primes.}
\label{claim:HV9}
Assume the actual source statements \texttt{Count\allowbreak I},
\texttt{Count\allowbreak II} and \texttt{Count\allowbreak III}.  Then the literal
statement \texttt{Sum\allowbreak Out} holds.  Only the two auxiliary
primes supplied by \texttt{Count\allowbreak I} occur in its host threshold.

\begin{proof}
Fix $\chi$.  Choose $C_1,r_1,r_2$ from \texttt{Count\allowbreak I},
with $r_1,r_2$ prime, and choose the exponents $C_2,C_3$ of
the two other counts.  Put $C_Z=0$ by CZ6 and
\[
       K=\max(C_Z,C_2,C_3),\qquad
       C_{\rm out}=C_1+3K+4.
 \tag{HV.37}
\]
All choices precede the host and field.
Now take a permissible host above $r_1,r_2$, a cutoff $t\ge2$
with \texttt{Big}$(t)$, any original CM field and local data,
and an arbitrary finite set $S$ of odd primes outside
$\mathcal B_\chi$.  Set $M=Nt$.
Choose, for each prime, a coefficient ring splitting the
finite Galois group of this original field.

For a fixed $1\le e<t$, split $S=S_0\sqcup S_1$ according
to whether $\ell-1\nmid2e$ or $\ell-1\mid2e$.
On $S_0$, \texttt{Count\allowbreak I} applies to the whole finite set
and gives the joint bound $M^{C_1}$, as well as finiteness
of every required character length.
For $\ell\in S_1$, HV1 gives $\ell\le2t$.
The cases in HV7 are now Z, II or III; CZ6 and the two
assumed counts give, for every $\psi\in X_e$,
\[
                 v_\ell(\psi,e)\log\ell
                        \le(\ell Ne)^K\le M^{3K}.
 \tag{HV.38}
\]
There are at most $2t$ such primes and at most $N$ characters
at each prime.  Thus the combined character sum at this
weight is at most
\[
                         M^{C_1}+2M^{3K+1}.
 \tag{HV.39}
\]
Every prime of $S$ is coprime to $n$ by HV1.  Consequently
the error-free comparison (HV.14) applies, including at the
possibly very large primes of $S_0$; no sum of an uncontrolled
error over those primes occurs.
Summing (HV.39) over $e$ gives
\[
 \sum_{\ell\in S}\log\ell\sum_{1\le e<t}
                         \mathtt{plusLen}_{D_\ell}(e)
 \le t\bigl(M^{C_1}+2M^{3K+1}\bigr)
 \le M^{C_1+3K+4}.
 \tag{HV.40}
\]
For the last bound use $t\le M$, $M\ge2$, and bound both
terms by the larger displayed monomial, with one extra
power to absorb their sum.
Finiteness follows at every $(\ell,e)$ from the relevant
count and HV4.  Since $S$ was an arbitrary finite set, this
is precisely \texttt{Sum\allowbreak Out}, with the original uniform
prime-set quantifier.
\end{proof}

\paragraph{HV10: literal PointHV and immediate window propagation.}
\label{claim:HV10}
The completed suppliers above prove
\[
 \boxed{\quad
 \mathtt{CountI}\wedge\mathtt{CountII}\wedge
 \mathtt{CountIII}\wedge\mathtt{CountV}
                    \ \Longrightarrow\ \mathtt{PointHV}.
 \quad}
 \tag{HV.41}
\]
More explicitly, for each $\chi$ and each prescribed function
$j:\mathbb N\to\mathbb N$, there are natural numbers
$C,r_1,r_2$, with $r_1,r_2$ prime, before any host is chosen,
such that for every permissible host above those thresholds,
every $t\ge2$ satisfying \texttt{Big}$(t)$, and every $B\ge0$
for which
\[
 \mathtt{LayerLe}(\chi,\ell,j(\ell),t,B)
                      \quad(\ell\in\mathcal B_\chi)
 \tag{HV.42}
\]
holds, every original CM field $E$ of that host satisfies
\[
                       \mathtt{VVplusLe}
                               \bigl(E,t,(Nt)^C+CB\bigr).
 \tag{HV.43}
\]
With TW8, OB8 and the already reconstructed VW assembly,
the same four-count premise therefore implies the full
cohomological window conclusion
\[
 \forall C_t\ge1\ \forall\text{ primitive }\psi,
                         \mathtt{VWinAt}(\psi,C_t).
 \tag{HV.44}
\]
The previously stated prime-pair inputs of VW retain their
stated role.  In particular its $C_t=401$ specialization
retains the fixed pair and the cutoff exponent $400$.

\begin{proof}
By HV8 and HV9 the four-count premise supplies \texttt{Sum\allowbreak In}
and \texttt{Sum\allowbreak Out}; HV5, HV6 and CZ6 have already supplied
their other inputs.  Fix $\chi$ and the whole function $j$.
Choose $C_{\rm out},r_1,r_2$ from HV9 and, for each prime
of the finite set $\mathcal B_\chi$, choose the constant
$C_{\rm in}(\ell,j(\ell))$ from HV8.  Put
\[
 D=\sum_{\ell\in\mathcal B_\chi}C_{\rm in}(\ell,j(\ell)),
 \qquad k=|\mathcal B_\chi|,
 \qquad C=C_{\rm out}+D+k+1.
 \tag{HV.45}
\]
These are choices before the host, $t,B,E,\Omega$ and local
data.  For any subsequent choices satisfying (HV.42), put
$M=Nt\ge2$ and split an arbitrary finite set $S$ of odd primes
into its intersection with $\mathcal B_\chi$ and its complement.
The sum in (HV.43) over $S$ is bounded by
\[
 \begin{split}
 M^{C_{\rm out}}+
 \sum_{\ell\in S\cap\mathcal B_\chi}
       \bigl(M^{C_{\rm in}(\ell,j(\ell))}
                          +C_{\rm in}(\ell,j(\ell))B\bigr)
 &\le (k+1)M^{C_{\rm out}+D}+DB\\
 &\le M^C+CB.
 \end{split}
 \tag{HV.46}
\]
Here each $C_{\rm in}$ is positive, $B\ge0$, and
$k+1\le2^k\le M^k$ (also valid at $k=0$).
The two component results prove all individual finiteness
assertions.  The field and local data were arbitrary,
and (HV.46) is uniform over every finite $S$.
This is exactly \texttt{VVplus\allowbreak Le}, hence the literal
\texttt{Point\allowbreak HV}, with the full quantifier order in (HV.43).

For (HV.44), use the single constant record and universal
\texttt{Thm\allowbreak Bab} supplied by TW8, and the odd-minus
\texttt{Thm\allowbreak Bc\allowbreak Free} supplied by OB8.  Substitution of (HV.41)
into VW3--VW4 now fills their remaining odd-plus premise.
Those proofs already keep the same original field, all the
prescribed layer models, one common window and the fixed
prime pair.  Their rescaling with $C_t=401$ is unchanged.
This proves the asserted immediate propagation.
\end{proof}

\paragraph{Scope at the end of the chain.}
HV proves the two auxiliary counts and every step after the
four-count cut.  In particular the separate
\texttt{Point\allowbreak HV}-assembly task is complete, including bad
degree primes and the coefficient and field comparisons.
The full human reconstruction still needs the actual
\texttt{Count\allowbreak I}, \texttt{Count\allowbreak II}, \texttt{Count\allowbreak III} and
\texttt{Count\allowbreak V}, or complete sufficient suppliers for them.
The common arithmetic family, genuine theta/local installation,
denominator tables, arithmetic gap and exact metric/spectral
comparisons remain the independent global obligations already
identified in the audit.  Neither (HV.41) nor (HV.44) asserts
that those inputs have been constructed.


\section{Dependency closure after rank-one, algebra and height assembly}
\label{sec:dependency-audit-v26}

\paragraph{The rank-one branches now reach their literal endpoints.}
RV proves the complete chain from finite duality and KB8 through an
actual finite root-containing extension, a finite quotient
restriction kernel and arbitrary finite-character Selmer finiteness.
Its Baire argument uses closed subgroups defined by the cocycle;
it does not assert that an arbitrary compact subset of the algebraic
closure lies in a finite field.  One finite coefficient lattice and
an explicit limiting identity then give the actual vector coboundary.
The endpoints are \texttt{So\allowbreak Fin}, \texttt{So\allowbreak V} and exact coefficient
independence, including nonnested rings.

BR proves the prescribed \texttt{Lemma\allowbreak X} by integral partial sums
and finite Frobenius projection.  The precise conductor inequality
forces a two-point contradiction in odd characteristic.  It then
uses the original tame-index condition, finite class-field
reciprocity and the already proved OS1 Jacobi ideal relation to
derive the required scalar annihilator.  The resulting residual
coboundary lies in the original residue field.  The final
\texttt{Cor\allowbreak X} handles both characters with the exact source data.

\paragraph{Complete algebra suppliers, with arithmetic installation separate.}
SP1 proves \texttt{Lemma\allowbreak M'} even though its injection is only
additive.  SP2--SP6 construct the exact weighted intersection
argument: admissible multigraded Euler characteristics, the actual
Rees support, two-sided orbit slices, generic multiplicity one,
an explicit leading-form Hilbert difference and the length--norm
formula.  The lower and upper multiplicities differ by exactly
a factor two, giving the literal \texttt{Lemma\allowbreak M}.
FP supplies the complete universal \texttt{Lemma\allowbreak P}, including
representability and the actual relation/obstruction construction.
Its original profinite-group, local-condition and finite-obstruction
hypotheses remain part of that theorem; proving it does not
construct a model for an arbitrary arithmetic datum.

SP7 proves the actual \texttt{sha\allowbreak Coc}/\texttt{m\allowbreak Prime} edge on a
supplied \texttt{Model}, with additive reductions of derivation
choices justified.  SP8 gives the exact return from auxiliary
coefficients.  SP9 proves the integral trace-product bound from
the precisely quantified finite catalogue.  SP10 carries the
actual data and compatible-model hypotheses and that catalogue
through all trace, valuation, polynomial absorption and coefficient
steps to \texttt{Count\allowbreak II}\,$\wedge$\,\texttt{Count\allowbreak III}.
The data/model/catalogue requirements remain actual obligations;
generic automorphy or a named package cannot substitute for them.

\paragraph{The complete original-field height assembly.}
HV proves an integral character inequality at every odd prime.
Finite quotient cohomology gives a polynomial error when the prime
divides the original field degree, and gives zero error outside
the exceptional prime set.  The proof keeps the original
$E,\Omega,D$, the negative twist and the normalized coefficient
length.  HV5 proves \texttt{Count\allowbreak Vs} from a field whose degree is
bounded before the host; HV6 proves the exact unweighted
\texttt{Count\allowbreak IV} layer comparison.  The case partition and all
finite-prime sums are then complete, yielding
\[
 \begin{gathered}
 \mathtt{CountI}\wedge\mathtt{CountII}\wedge
 \mathtt{CountIII}\wedge\mathtt{CountV}\\
 \Longrightarrow\ \mathtt{PointHV}
 \Longrightarrow\ \forall C_t\ge1\ \forall\psi,
                              \mathtt{VWinAt}(\psi,C_t).
 \end{gathered}
\]
The second implication uses the already proved TW8 and OB8 in VW.
The fixed-pair $C_t=401$ rescaling and cutoff exponent $400$ are
unchanged.  No separate original-field assembly task remains
after the four actual counts.

\paragraph{The retained frontier.}
The actual \texttt{Count\allowbreak I/II/III/V}, their remaining same-object
deformation, local-global cohomology, automorphy, finite catalogue
and rigidity suppliers, and the global common family/gap, genuine
theta/local installation, denominator tables and exact
metric/spectral comparisons remain.  The current audit preserves
the historical nodes and target snapshots, records the newly
supplied interfaces and narrows the retained targets.  It is a
finite sufficient-condition map, without a claim of necessity,
minimality or a percentage of the full proof.

% Version 0.27.  Actual arithmetic data and the complete level comparison.
% Pinned Lean source: 7cbf7a0c32fca1badf05931879a37c3f8a2b03b4.
\section{The actual symplectic arithmetic data}
\label{sec:symplectic-arithmetic-data}

\status{AX1--AX10 close the arithmetic construction behind G1.
They prove the literal \texttt{data\allowbreak Ex}, all of \texttt{facts},
\texttt{lev\allowbreak Of}, the ordinary specialization theorem \texttt{sp\allowbreak Pt\allowbreak L},
and the level, duality and safety assertions on the actual points.
The final admissible class uses the explicitly retained
\texttt{Symplectic\allowbreak Four\allowbreak Image\allowbreak Dense} hypothesis to supply its common
automorphic embedding, \texttt{Fin0} and \texttt{No\allowbreak End}.  None of those
three assertions is inferred from the character construction.}

\subsection{The original objects and their conditions}

Fix the source algebraic closure $\Omega/\mathbb Q$ and put
$G=\operatorname{Gal}(\Omega/\mathbb Q)$.  Let $\ell$ be an odd prime,
let $O\subset\overline{\mathbb Q}_\ell$ be the original coefficient
DVR, finite over $\mathbb Z_\ell$, and let $k=O/\mathfrak m_O$.
Write $c:G\to\mathbb Z_\ell^\times\subset O^\times$ for the
cyclotomic character, $\omega$ for its Teichm\"uller part, and a bar
for reduction to $k$.  A Dirichlet character is evaluated on $G$
through the cyclotomic action at its specified level.  Characters of
different levels are multiplied after inflation to a common multiple;
all displayed identities involving them mean identities on units.

Keep the host of HV1: a primitive character $\chi$ of conductor $q$
and order $o$, an integer $s\ge2$, and the original prime pair $q_1,d$.
Put
\[
 N=\operatorname{lcm}(2^{s+2},q)q_1d,
 \qquad H\le(\mathbb Z/N)^\times
 \tag{AX.1}
\]
with exactly the four defining conditions in HV1.  A host character
$\psi$ is trivial on $H$.  Membership in $h.X(O,e)$ also requires
$\psi(-1)=(-1)^e$.  Throughout, $\operatorname{clog}_2(x)$ is the
least $n\ge0$ with $x\le2^n$ for the positive integers used below.

For a finite-order character, let $\operatorname{comp}_p$ denote its
CRT component supported at $p$.  Split
\[
 \psi=\psi_2\theta\kappa,
 \quad \psi_2=\operatorname{comp}_2\psi,
 \quad \kappa=\psi_{\ell\text{-primary}},
 \quad \theta=\psi_{\ell'}(\operatorname{comp}_2\psi_{\ell'})^{-1}.
 \tag{AX.2}
\]
Also put $\psi_P=\psi_{\ell'}
(\operatorname{comp}_\ell\psi_{\ell'})^{-1}$,
$c_2=v_2(\operatorname{cond}\psi_2)$,
$f_0=\operatorname{ord}_{\operatorname{lcm}(4,\operatorname{ord}\theta)}(\ell)$,
and $n_0=v_2(\ell^{f_0}-1)$.  These are the literal source
\texttt{psi2}, \texttt{th1}, \texttt{kap}, \texttt{chi\allowbreak P},
\texttt{c2}, \texttt{f0} and \texttt{n0}.
The two prescribed dyadic levels are
\[
\begin{split}
 n_{\rm II}&=\operatorname{clog}_2\max
  \{\ell N4^{n_0},\ 2^{c_2+6}o\ell^3q^2\},\\
 n_{\rm III}&=\operatorname{clog}_2\max
  \{\ell q p_{\!1}4^{n_0},\ 8\ell^2\}.
\end{split}\tag{AX.3}
\]
Here $p_{\!1}=\mathtt{pOne}(\chi)$ is the fixed prime congruent to
$1$ modulo the odd part of $o$ and coprime to $2qp_K$.  It does not
vary with the host.  The exceptional set is the same
$\mathcal B_\chi$ as in HV1.  In the two cases,
\[
\begin{array}{ll}
\mathrm{II}:&e\ge2,\quad \operatorname{ord}(\psi_P)\text{ even},
 \quad[\ell\in\mathcal B_\chi\ \text{or}\ \ell-1\mid2e],\\
\mathrm{III}:&\ell\notin\mathcal B_\chi,\quad\ell-1\mid2e,
 \quad\operatorname{ord}(\psi)\text{ odd},\quad\psi\ne1.
\end{array}\tag{AX.4}
\]

A source \texttt{Raw} datum $D$ consists of $S,e,\psi,\xi,\omega$,
a complex conjugation $\mathrm{cc}$, a place $A_p$ and an element
$\sigma_p$ for every rational prime $p$, and $\phi,g_0,s_0$.
The same symbol $\psi$ will denote its associated Galois character.
The derived characters and four weights are
\[
 \begin{gathered}
 \Psi=c^e\psi,\quad \alpha=c\omega^{-1}\xi,
 \quad\beta=\Psi\alpha^{-1},\quad\vartheta=\alpha\Psi,\\
 (p_0,p_1,p_2,p_3)=(\vartheta,\Psi,\alpha,1),
 \quad(w_0,w_1,w_2,w_3)=(e+1,e,1,0).
 \end{gathered}
 \tag{AX.5}
\]
Here the subscripts on $p_i$ denote characters, not rational primes.
Let $r_0=\operatorname{diag}(p_i)$ and $\bar r_0=\overline{r_0}$.
The source local sets are
\[
 S_b=\{p\in S:p\ne\ell,
       \ \ell\mid|\psi(I_{A_p})|\},\qquad T=\{\ell\}\cup S_b.
 \tag{AX.6}
\]
The condition $\operatorname{UnrT}(z,0)$ says that $z(g)=0$ for
$g\in I_A\cap\ker\bar r_0$ whenever $A$ is away from $T$.
The condition $X\!\operatorname{Van}(\gamma)$ says that a locally
constant $k(\bar\gamma)$-cocycle with this property and locally a
coboundary on each $G_{A_v}$, $v\in T$, is a global coboundary.

For clarity, the full \texttt{Good} condition has the following
content.  The set $S$ contains $2,\ell$ and consists of primes;
$e\ge2$; the selected $A_p$ lie over $p$; and $\mathrm{cc}$ is a
complex conjugation.  The character $\psi$ has open kernel, is
unramified outside $S$, and $\Psi(\mathrm{cc})=1$.
One has $\omega^{\ell-1}=1$, $\bar\omega=\bar c$.
The character $\xi$ has open kernel and order prime to $\ell$,
is unramified outside $S$ and at all $v\in T$, and is odd.
The elements $\sigma_v$ are arithmetic Frobenius lifts for $v\in T$,
$c(\sigma_\ell)=1$, and
\[
 \overline{\xi\Psi^i}(\sigma_v)\ne1
       \quad(v\in T,\ i\in\mathbb Z).
 \tag{AX.7}
\]
Both $X\!\operatorname{Van}(\gamma)$ and
$X\!\operatorname{Van}(\gamma^{-1})$ hold for
$\gamma=\alpha,\beta,\vartheta$.  Finally,
\[
 \begin{gathered}
 \sigma_\ell^{\ell^{n!}}\longrightarrow\phi,
 \qquad\bar\alpha(g_0)\ne1,\quad\bar\Psi(g_0)\ne1,\\
 \qquad s_0\in I_{A_\ell}\cap\ker\psi\cap\ker\omega,
 \quad c(s_0)\text{ has infinite order}.
 \end{gathered}
 \tag{AX.8}
\]
This list includes every field of the source record; no stronger
goodness assumption will be added later.

\paragraph{Classical boundary.}
We use the elementary finite cyclotomic and local-field facts in
Milne, \href{https://www.jmilne.org/math/CourseNotes/ANT.pdf}
{\emph{Algebraic Number Theory}}, v3.08, Sections 6--8: the
cyclotomic Galois group, extension and conjugacy of places, the
decomposition-group comparison with the local completion, and the
wild/tame inertia description.  The exact consequences used below
are proved from these finite statements.  The norm on a finite
extension of $\mathbb Q_\ell$ is its unique extension, as in RV1.
The one commutative-algebra fact used in AX8 is that a finite algebra
over a Henselian local ring is a product of Henselian local rings
(\href{https://stacks.math.columbia.edu/tag/04GH}{Stacks, Tag 04GH});
a domain is therefore local.  Completeness implies Henselianity by
successive lifting of simple roots
(\href{https://stacks.math.columbia.edu/tag/04GM}{Tag 04GM}).
BR7 supplies the full
residual vanishing theorem, with its already recorded classical
inputs.  There is no automorphy, crystalline-period, image-density,
or rigidity input before AX10.

\subsection{Finite characters with the prescribed conductors}

\begin{proposition}[AX1: exact host arithmetic and the two cases]
\label{prop:ax-host}
For a host character the decomposition (AX.2) holds.  Its dyadic
component has conductor $2^{c_2}$ and order dividing $2^{c_2}$;
$\kappa$ has $\ell$-power order; and $\theta$ has odd conductor,
prime-to-$\ell$ order dividing $2o$.
In particular $f_0>0$, $4\mid\ell^{f_0}-1$,
$\theta^{\ell^{f_0}-1}=1$, and $\ell-1\mid\ell^{f_0}-1$.

Put $T_b=\{\ell\}\cup\{p\mid N:p\ne\ell,
\operatorname{comp}_p\kappa\ne1\}$.  Every $p\in T_b\setminus\{\ell\}$
is odd, divides $q$, and is congruent to $1$ modulo $\ell$.
Every $v\in T_b$ divides $\ell q$.
For any places chosen over the rational primes, the actual set
$T$ in (AX.6), with $D.\psi=\psi$ and prime-valued $S$, lies in $T_b$.

In Case II, $\psi_P\ne1$ and $n_{\rm II}\ge c_2+6$.
In Case III, under the original $h.\mathtt{Big}(t)$ and
$0<e<t$, one has
\[
 \begin{gathered}
 e\ge2,\quad\ell\nmid oN,\quad\ell\ne p_{\!1},
 \quad\psi_2=\kappa=1,\quad\theta=\psi_P=\psi,\\
 \quad\operatorname{cond}\psi\mid q,\quad T_b=\{\ell\},
 \quad\operatorname{ord}\psi\mid\operatorname{oddpart}(o),
 \quad n_{\rm III}\ge3.
 \end{gathered}
 \tag{AX.9}
\]
\end{proposition}

\begin{proof}
The finite unit group modulo $N$ is the product of its prime-power
unit groups.  Projection to a factor gives
$\operatorname{comp}_p$, so the components commute, their product
is the original character, and their orders divide the original
order.  Splitting a finite cyclic character image into its
$\ell$-primary and prime-to-$\ell$ factors commutes with every
component.  The group of units modulo a power of $2$ is a $2$-group;
thus $\operatorname{comp}_2\kappa=1$.  Removing the remaining dyadic
component gives (AX.2), odd conductor for $\theta$, and the asserted
dyadic order bound by the size of the unit group.

For each unit $u$ modulo $N$, the power $u^{2^{s+1}o}$ lies in $H$:
it is $1$ modulo $2^{s+2}$, is killed by $\chi$, and is a square
at the two auxiliary primes.  Hence
$\operatorname{ord}\psi\mid2^{s+1}o$.  More precisely, for an odd
prime $p\mid q$ the $p$-supported unit raised to $o$ lies in $H$,
so $\operatorname{ord}(\operatorname{comp}_p\psi)\mid o$.
For an odd prime $p\nmid q$ its square lies in $H$, so this
component has order dividing $2$.  Consequently every odd
component, and then their product after taking prime-to-$\ell$
parts, has order dividing $2o$.  The unit $\ell$ modulo
$\operatorname{lcm}(4,\operatorname{ord}\theta)$ has finite positive
order $f_0$, which gives its stated divisibilities.  The last one
also follows directly from $\ell^f-1=(\ell-1)(1+\cdots+\ell^{f-1})$.

A nontrivial $\ell$-primary component at an odd prime $p\ne\ell$
cannot occur at either auxiliary prime: there the character has
order dividing $2$.  Thus $p\mid q$.  On
$(\mathbb Z/p^a)^\times$ the kernel of reduction modulo $p$ is
a $p$-group, so an $\ell$-primary character factors through
$(\mathbb Z/p)^\times$ and its nontrivial order divides $p-1$.
This proves the assertions about $T_b$.  Inertia on roots of unity
of order prime to $p$ is trivial.  Therefore the restriction of
$\psi$ to $I_{A_p}$ equals the restriction of its $p$-component.
If the former image has order divisible by $\ell$, so does the
latter.  This is equivalent to
$\operatorname{comp}_p\kappa\ne1$, proving $T\subset T_b$ without
changing the selected places.

Even order makes $\psi_P$ nontrivial.  The second term in the
maximum defining $n_{\rm II}$ is at least $2^{c_2+6}$, so its
claimed lower bound follows.  For Case III, an odd-order character
has trivial dyadic component and trivial components at the two
auxiliary primes, hence conductor dividing $q$ and order dividing
the odd part of $o$.  Since $\ell\notin\mathcal B_\chi$, the
prime $\ell$ is at least $5$, does not divide $oq$, and differs
from $p_{\!1}$.  Also $\ell-1\mid2e$ gives
$\ell\le2e+1<2t+1<q_1,d$.  Thus $\ell\nmid N$ and $e\ge2$.
The $\ell$-primary and $\ell$-local components of $\psi$ vanish,
which proves all identifications in (AX.9).  Finally $8\ell^2$
in (AX.3) gives $n_{\rm III}\ge3$.

Only the fixed defining property of $p_{\!1}$ is used here.
Its existence in the stated class can even be seen without an
analytic progression theorem.  For $m>1$, choose an integer $a$
divisible by $m$ and any prescribed finite excluded set, with
$a\ge3$.  Any prime divisor $p$ of $\Phi_m(a)$ is outside that
set and does not divide $m$, since $\Phi_m(a)\equiv1\pmod a$.
The reduction of $X^m-1$ is separable at $p$, so a root of
$\Phi_m$ has order exactly $m$, and $p\equiv1\pmod m$.
The inequality $|\Phi_m(a)|\ge(a-1)^{\varphi(m)}>1$ supplies a
prime divisor.  For $m=1$ use the usual Euclidean construction.
Taking $m=\operatorname{oddpart}(o)$ and excluding primes dividing
$2qp_K$ proves that the least such prime has precisely the
properties used above.  The earlier fixed definitions of $p_K$
and $\mathcal B_\chi$ retain their meanings.
\end{proof}

\begin{proposition}[AX2: primitive dyadic characters and quantitative separation]
\label{prop:ax-dyadic}
If $n\ge3$ and $O$ contains all $2^n$-th roots of unity, there is a
primitive odd character $\xi_*\pmod{2^n}$.
Twisting a primitive character modulo $2^n$ by a character of
conductor $2^c$ with $c<n$ preserves primitivity.
For every primitive $\xi_*\pmod{2^n}$, odd $v>1$, odd positive $u$,
and $a\ge0$,
\[
 \xi_*(v)^{2^a u}=1\quad\Longrightarrow\quad
                     2^n<v^2 2^a.
 \tag{AX.10}
\]

Suppose $p$ is odd (allow $p=1$), $d>0$,
$4\mid\ell^d-1$, and $\omega_0,\theta_0,\xi_0$ have respective
levels $\ell,N,p$, with odd conductor for $\theta_0$ and orders
dividing $\ell^d-1$.  Let $\chi_2$ have conductor $2^c$, $c<n$.
If the product $\omega_0\chi_2\theta_0\xi_0\xi_*$ is odd and
\[
       p\ell\operatorname{cond}(\theta_0)
                  4^{v_2(\ell^d-1)}\le2^n,
 \tag{AX.11}
\]
then it equals $\operatorname{prodX}(\varphi,\xi_1)$ on units
of the common level, for literal $X\!\operatorname{Data}
(\ell,d,\varphi,\xi_1)$, with
$\operatorname{cond}(\varphi)\mid p\ell N$.
\end{proposition}

\begin{proof}
The standard generators give
$(\mathbb Z/2^n)^\times=\langle-1\rangle\times\langle5\rangle$,
where $5$ has order $2^{n-2}$.  To check the order, repeated squaring
shows $v_2(5^{2^j}-1)=j+2$; the subgroup $1+4\mathbb Z/2^n$
has exactly $2^{n-2}$ elements.  Send $-1$ to $-1$ and $5$ to
a primitive $2^{n-2}$-th root in $O$.  The resulting character is
odd and is nontrivial on $1+2^{n-1}$, the kernel of reduction to
level $2^{n-1}$, so it is primitive.  Conversely every primitive
character must be nontrivial on this order-two kernel, and takes
the value $-1$ there.  A lower-conductor twist is trivial on that
kernel; this proves the twisting assertion.

For (AX.10), the assertion for $n\le1$ is immediate from $v\ge3$.
For $n\ge2$, put $b=v_2(v^2-1)\ge3$.  Odd powering preserves
$v_2(x-1)$ for $x\equiv1\pmod4$, and each subsequent squaring
increases it by one.  Indeed
$(1+2^b t)^2-1=2^{b+1}(t+2^{b-1}t^2)$ with the parenthesis odd;
the odd-power statement follows from the geometric sum modulo $2$.
Thus
$v_2(v^{2u2^a}-1)=b+a$.
If $b+a<n$, further squaring produces a power of $v^{2^a u}$
congruent to $1+2^{n-1}$ modulo $2^n$.  Its character value is
$1$ under the assumed equality, contradicting primitivity.
Therefore $n\le b+a$, and
$2^n\le2^{b+a}<v^2 2^a$, as $2^b\le v^2-1$.

Take the product of $\omega_0,\theta_0,\xi_0$ at a common level
and pass to its primitive inducing character $\varphi$.  Its
conductor $f$ divides $p\ell\operatorname{cond}(\theta_0)$,
is odd, and its order divides $\ell^d-1$.  These assertions follow
directly by quotienting the finite unit group by the character
kernel: removing a redundant level preserves its image and order.
Inflate $\chi_2$ to $2^n$ and multiply by $\xi_*$ to obtain the
primitive $\xi_1$.  The conductors, the original parity equality,
and $f4^{v_2(\ell^d-1)}\le2^n$ now give every clause of the
$X\!\operatorname{Data}$ record in BR7.  Evaluation on a Galois
element is compatible with this identity: choose its one
cyclotomic residue modulo $p\ell N2^n$, represented by an integer
coprime to that level, and reduce the same integer at every factor.
This also proves the source \texttt{dir\allowbreak X} and \texttt{dir\allowbreak Val}
compatibility at the original $\Omega$.
\end{proof}

\subsection{Choices on the original algebraic closure}

\begin{proposition}[AX3: Frobenius choices and the prime-to-$\ell$ limit]
\label{prop:ax-choices}
For each rational prime choose a place $A_p$ of the original
$\Omega$ and an arithmetic Frobenius $\sigma_p$.  They may be
chosen so that $c(\sigma_\ell)=1$.  For any finite character
$\psi:G\to O^\times$ with open kernel there is
$s_0\in I_{A_\ell}\cap\ker\psi\cap\ker\omega$
with $c(s_0)$ of infinite order.

If $G_0$ is any profinite group, $g\in G_0$ and $\ell$ is prime,
the sequence $g^{\ell^{n!}}$ converges to an element $g^{(\ell')}$.
For any adically continuous lift $\rho$ of a finite-residue
representation, of any finite matrix size, suppose
$\bar\rho(g)^m=1$ with $(m,\ell)=1$.  Then
\[
 \bar\rho(g^{(\ell')})=\bar\rho(g),
 \qquad \rho(g^{(\ell')})^m=1
 \tag{AX.12}
\]
in its source complete Noetherian local coefficient algebra.
The residual continuity needed here follows from the lift.
Finally, any two nontrivial group
characters to groups have an element at which both are nontrivial.
These choices supply the literal \texttt{raw\allowbreak Ex} output whenever
the two required residual characters are nontrivial.
\end{proposition}

\begin{proof}
Finite separable local extensions identify the chosen decomposition
group with the Galois group of $\mathbb Q_p$; the quotient by inertia
is $\widehat{\mathbb Z}$ with arithmetic Frobenius generator.  Choose
a lift.  At $\ell$, inertia maps onto
$\operatorname{Gal}(\mathbb Q_\ell(\mu_{\ell^r})/\mathbb Q_\ell)
\simeq(\mathbb Z/\ell^r)^\times$ for every $r$: the polynomial
$\Phi_{\ell^r}(1+X)$ is Eisenstein, so the finite cyclotomic
extension is totally ramified and has exactly that degree.
The nested closed fibers in the compact inertia group show that
$c(I_{A_\ell})=\mathbb Z_\ell^\times$, with the full inverse-limit
target, not only each separate quotient.  Multiply the chosen
Frobenius by an inertia element of cyclotomic value
$c(\sigma_\ell)^{-1}$.  This changes no residue Frobenius and gives
the asserted cyclotomic value.  At $p\ne\ell$, roots of unity of
order prime to $p$ lift from unramified residue extensions by
Hensel's lemma, so inertia fixes them and arithmetic Frobenius
acts by $\zeta\mapsto\zeta^p$.

Choose $u\in I_{A_\ell}$ with $c(u)=1+\ell$.
The open kernel of $\psi$ has finite index $b$; its image is
therefore finite of order $b$.  Then
$s_0=u^{b(\ell-1)}$ kills both $\psi$ and $\omega$.
The rational integer $1+\ell$ has infinite order in
$\mathbb Z_\ell^\times$ (injectivity of $\mathbb Z\to\mathbb Z_\ell$
would otherwise give $(1+\ell)^a=1$ in $\mathbb Z$ for some $a>0$).
Its indicated positive power still has infinite order.
Teichm\"uller lifting is the unique lifting of every element of
$\mathbb F_\ell^\times$ to a root of $X^{\ell-1}-1$ in
$\mathbb Z_\ell$; uniqueness proves multiplicativity.  This gives
$\omega^{\ell-1}=1$, $\bar\omega=\bar c$, and its factorization
through the cyclotomic action modulo $\ell$.
Complex conjugation extends from the algebraic real closure to
$\Omega$ and has cyclotomic value $-1$; all these choices are
made in the initially specified $G$.

In a finite quotient, write the order of $g$ as $\ell^a b$ with
$(b,\ell)=1$.  For all sufficiently large $n$, $\ell^{n!}$ is
$0$ modulo $\ell^a$ and $1$ modulo $b$, by Euler's theorem and
$\varphi(b)\mid n!$.  Thus the image of the sequence is eventually
constant in every finite quotient.  The constants are compatible,
so give an element of the profinite inverse limit and prove
convergence.  This proves the existence assertion without a
first-countability assumption on the group.

To prove the lift assertion let $\mathfrak m$ be the maximal ideal
of its coefficient algebra and $X=\rho(g)^m$.  Then
$X\equiv1\pmod{\mathfrak m}$ and $\ell\in\mathfrak m$.
The identity
\[
 X^\ell-1=(X-1)\left(\ell\,1+
                   \sum_{j=0}^{\ell-1}(X^j-1)\right)
 \tag{AX.13}
\]
shows that entrywise divisibility $X-1\in M_\nu(\mathfrak m^a)$,
$a\ge1$, implies $X^\ell-1\in M_\nu(\mathfrak m^{a+1})$
for any finite matrix size $\nu$.
Hence $X^{\ell^j}\equiv1\pmod{\mathfrak m^{j+1}}$.
Continuity modulo each $\mathfrak m^b$ and the convergence just
proved give $\rho(g^{(\ell')})^m\equiv1$ modulo every such power.
The ring is separated, proving (AX.12).  The residual identity
uses $\ell^{n!}\equiv1\pmod m$ and continuity of the finite
residual representation.

For the final assertion choose $x$ nontrivial under the first
character and $y$ nontrivial under the second.  If neither already
works, the first character is $1$ at $y$ and the second is $1$
at $x$, so $xy$ works.  Applied to $\bar\alpha,\bar\Psi$, this
chooses the same $g_0$ for both inequalities in (AX.8).
\end{proof}

\begin{proposition}[AX4: the residual local conditions and the generic good datum]
\label{prop:ax-generic}
Assume $e\ge2$, $\psi\in h.X(O,e)$ and $\psi_P\ne1$.
Let $p$ be odd, allowing $p=1$, and choose an even character
$\xi_0\pmod p$ and a primitive odd character $\xi_*\pmod{2^n}$.
Suppose
\[
\begin{gathered}
 v\nmid p\quad(v\in T_b),\qquad
 \xi_0^{\ell^{f_0}-1}=1,\qquad c_2<n,\\
 p\ell\operatorname{cond}(\theta)4^{n_0}\le2^n,
 \qquad u\text{ odd},\quad
 \psi^{2^a u}=\xi_0^{2^a u}=1,\qquad
 v^2 2^a\le2^n\quad(v\in T_b).
\end{gathered}\tag{AX.14}
\]
Then, using BR7, there is an actual $D$ on the original $\Omega,O$
with $D.\mathtt{Good}$ and
\[
 D.e=e,\quad D.\psi=\operatorname{dirG}\psi,\quad
 D.S=\operatorname{primeFactors}(2\ell Np),\quad
 D.\xi=\operatorname{dirG}\xi_0\operatorname{dirG}\xi_*.
 \tag{AX.15}
\]
It has the literal $\mathtt{OfChar}(e,\psi,n,p,
\operatorname{ord}\xi_0)$ witnesses.
\end{proposition}

\begin{proof}
Take the character data and the original-field choices of AX3.
For now postpone only $g_0$; AX1 already gives $T\subset T_b$.
Reduction kills $\kappa$, since $k^\times$ has order prime to
$\ell$.  With $\omega$ also read as the Dirichlet character
modulo $\ell$, one has
\[
 \bar\alpha=\overline{\xi_0\xi_*},\qquad
 \bar\Psi=\overline{\omega^e\psi_2\theta}.
 \tag{AX.16}
\]
For every $v\in T$, $\bar\omega(\sigma_v)=1$.
At $\ell$ this is AX3; at every other $v$ it follows from
$v\equiv1\pmod\ell$ and the Frobenius rule.
If $\overline{\xi\Psi^i}(\sigma_v)=1$, raise the equality to
$2^a u$.  The factors $\psi$ and $\xi_0$ disappear by (AX.14),
and the remaining cyclotomic residue is $1$.  Thus
$\overline{\xi_*(v)^{2^a u}}=1$.
Reduction is injective on roots of unity of order prime to $\ell$:
if $z\equiv1$ and $z^m=1$, $(m,\ell)=1$, then
$1+z+\cdots+z^{m-1}$ is a unit, so
$(z-1)(1+\cdots+z^{m-1})=0$ forces $z=1$.
The dyadic root in question therefore equals $1$ in $O$.
AX2 gives $2^n<v^2 2^a$, a contradiction.  This proves the full
integer-parameter condition (AX.7), including negative $i$.

We next verify the six residual vanishing conditions, with their
actual local subgroups.  The prime-to-$\ell$ residual diagonal
character group has order dividing $(|k|-1)^4$; consequently
$K=\ker\bar r_0$ is normal of finite index prime to $\ell$.
Let $z$ satisfy the defining conditions of $X\!\operatorname{Van}
(\gamma)$, where $\bar\gamma$ is a character covered by BR7.
At a place away from $T$, choose $H_A=K$ in BR7's
$\operatorname{UnrEv}$.  Its relative index in
$I_A\cap\ker\bar\gamma$ divides $[G:K]$, since $K$ is normal,
and the required vanishing is exactly $\operatorname{UnrT}$.
At a place above $v\in T$, conjugate inertia into the selected
$I_{A_v}$.  The local coboundary is zero on the kernel of
$\bar\gamma$.  For such an inertia element $x$, the cocycle
identity gives, explicitly,
\[
 z(g^{-1}xg)=\bar\gamma(g)^{-1}z(x)
       \quad\text{if }\bar\gamma(x)=1.
 \tag{AX.17}
\]
Thus $z$ vanishes on $I_A\cap\ker\bar\gamma$ and one may take
$H_A=G$.  This verifies the original relative-index condition at
every valuation subring.  BR7 gives the global coboundary in the
original field $k$.  Applying the same argument to the inverse
character gives its inverse case as well.

To apply it, use AX2 with the following three rows:
\[
\begin{array}{c|ccccc}
\gamma&\omega_0&\chi_2&\theta_0&\text{prime character}&
                 \text{dyadic character}\\ \hline
\alpha&1&1&1&\xi_0&\xi_*\\
\beta&\omega^e&\psi_2&\theta&\xi_0^{-1}&\xi_*^{-1}\\
\vartheta&\omega^e&\psi_2&\theta&\xi_0&\xi_*
\end{array}\tag{AX.18}
\]
Take $d=f_0$ in all rows.  AX1 and (AX.14) give the order and
conductor inequalities; for the first row use
$1\le\operatorname{cond}\theta$.  The parity of
$\omega^e\psi_2\theta$ is even, because
$\kappa(-1)=1$ and $\psi(-1)=(-1)^e$.
Each row is therefore odd.  Formula (AX.16), its inverse, and
the common-residue comparison in AX2 identify the residual
characters exactly.  All six applications use the same $D$, $K$
and originally chosen places; no new residual coefficient field
has appeared.

For $g_0$, (AX.7) with $v=\ell$ and $i=0$ proves
$\bar\alpha\ne1$.  Also $\bar\Psi\ne1$.  Here is the needed
argument rather than an assumption of that fact.  The prime-to-$\ell$
part of $\psi$ is the product of $\psi_P$, unramified at $\ell$,
and an $\ell$-local character.  The latter factors through
$(\mathbb Z/\ell)^\times$, since its restriction to the
principal-unit $\ell$-group is trivial.  It is some power
$\omega^b$.  If $\bar\Psi=1$, injectivity of reduction on
prime-to-$\ell$ roots gives $\omega^{e+b}\psi_P=1$ in $O$.
Restriction to $I_{A_\ell}$ and cyclotomic surjectivity force
$\omega^{e+b}=1$.  Hence $\psi_P=1$, contrary to hypothesis.
The finite cyclotomic action on the original $G$ is surjective,
so this last conclusion is also an equality of the original
Dirichlet characters.  AX3 now provides $g_0$.

It remains to check the other fields of \texttt{Good}.
The set in (AX.15) contains $2,\ell$ and only primes; its data
have the stated $e$, selected places and complex conjugation.
Finite Dirichlet characters have open kernel and are unramified
away from the primes in their levels, by AX3.
The order of $\xi_0\xi_*$ divides
$(\ell^{f_0}-1)\varphi(2^n)$, which is prime to $\ell$.
No prime of $T$ divides $p2^n$, so $\xi$ is also unramified
there.  Its parity is odd, whereas $\Psi(\mathrm{cc})=1$.
The Teichm\"uller, Frobenius, limit and $s_0$ clauses are exactly
AX3.  Together with (AX.7), (AX.17)--(AX.18) and the chosen $g_0$,
these are all the fields listed above.  Finally $\xi_*,\xi_0$
themselves witness \texttt{Of\allowbreak Char}, with precisely their original
levels and $\operatorname{ord}\xi_0$.
\end{proof}

\begin{proposition}[AX5: literal dataEx in both original host cases]
\label{prop:ax-dataex}
Let $h.\mathtt{Big}(t)$, $0<e<t$, $\ell\ne2$ and
$\psi\in h.X(O,e)$.  Suppose Case II or Case III holds and the
original $O$ contains all roots of unity of orders
$2^{n_{\rm II}}$ and $2^{n_{\rm III}}$.  Then there is
$D:\mathtt{Raw}(\Omega,O)$ with $D.\mathtt{Good}$ and the literal
$\mathtt{OfCase}(D,h,e,\psi)$.  More precisely, the parameters
$(n',p,m_0)$ in \texttt{Of\allowbreak Char} may be taken as
\[
 (n_{\rm II},1,1)\quad\text{in Case II},\qquad
 (n_{\rm III},p_{\!1},\operatorname{ord}\psi)
                         \quad\text{in Case III}.
 \tag{AX.19}
\]
\end{proposition}

\begin{proof}
In Case II choose the primitive odd dyadic character at the
prescribed $n=n_{\rm II}$ by AX2, set $p=1$, $\xi_0=1$, and write
$\kappa^{\ell^b}=1$.  Use
\[
 a=c_2+1+v_2(o),\qquad u=\operatorname{oddpart}(o)\ell^b.
 \tag{AX.20}
\]
The integer $u$ is odd and
$2^a u=2^{c_2}(2o)\ell^b$ kills every factor of (AX.2).
It also kills $\xi_0$.  The first maximum in (AX.3) gives
$\ell\operatorname{cond}(\theta)4^{n_0}\le2^n$.
For $v\in T_b$, AX1 gives $v\le\ell q$; hence
\[
 v^2 2^a\le\ell^2q^2\,2^{c_2+1}o
            \le2^{c_2+6}o\ell^3q^2\le2^n.
 \tag{AX.21}
\]
All other conditions of (AX.14) follow from AX1 or $p=1$.
AX4 supplies the first parameter choice in (AX.19).

In Case III, (AX.9) applies.  Put $p=p_{\!1}$ and
$m=\operatorname{ord}\psi$.  Since $m\mid\operatorname{oddpart}(o)$
and $p\equiv1\pmod{\operatorname{oddpart}(o)}$, choose a character
$\xi_0$ of $(\mathbb Z/p)^\times$ of order exactly $m$.
There is already a root of order $m$ in $O$: the finite image of
$\psi$ in the multiplicative group of a field is cyclic of order
$m$.  Send a generator of $(\mathbb Z/p)^\times$ to that root.
Thus no extra roots are required.  The order $m$ is odd, so
$\xi_0(-1)=1$.  Choose the primitive odd dyadic character at
$n=n_{\rm III}$; set $a=0,u=m$.  Then $T_b=\{\ell\}$,
$\ell\ne p$, and
\[
 p\ell\operatorname{cond}(\theta)4^{n_0}
 \le p\ell q4^{n_0}\le2^n,
 \qquad \ell^2\le8\ell^2\le2^n.
 \tag{AX.22}
\]
The order of $\xi_0$ divides $\ell^{f_0}-1$, because
$\theta=\psi$.  All conditions of AX4 now hold, giving the second
case with $m_0=m$.  This is the full quantified \texttt{data\allowbreak Ex}
implication with BR7 installed, on the original $\Omega,O$.
\end{proof}

\subsection{All local facts for any good datum}

\begin{proposition}[AX6: literal facts and the normalization unit]
\label{prop:ax-facts}
For every actual $D:\mathtt{Raw}(\Omega,O)$, $D.\mathtt{Good}$
implies all clauses of \texttt{facts} in \texttt{LMem\allowbreak Max}:
\begin{enumerate}
\item If $p\in S_b$, then $p$ is prime, $p\ne\ell$,
$\ell\mid p-1$, and $\omega|_{G_{A_p}}=1$.
\item For $v\in T$, $\sigma_v\in G_{A_v}$,
$\bar\omega(\sigma_v)=1$, and
$\bar p_i(\sigma_v)\ne\bar p_j(\sigma_v)$ whenever
$\lambda_i\ne\lambda_j$, where $\lambda=(1,0,1,0)$.
\item $\operatorname{UnrT}(p_i,1)$ holds for all $i$.  On
$I_{A_\ell}\cap\ker\psi\cap\ker\omega$ one has $p_i=c^{w_i}$.
For $p\in S_b$ and $u\in I_{A_p}$,
$\alpha(u)=1$ and $\Psi(u)=\vartheta(u)=\psi(u)$.
\item $p_i p_{3-i}=\vartheta$, the four values $p_i(s_0)$ are
distinct, and $\alpha(\mathrm{cc})=\beta(\mathrm{cc})=
\vartheta(\mathrm{cc})=-1$.
\item $\phi\in G_{A_\ell}$ and
$\bar r_0(\phi)=\bar r_0(\sigma_\ell)$.  Each $p_i$ is adically
continuous; $\ker\bar r_0$ and $\ker\omega$ are open.
Each $p_i$ is unramified outside $S$.
\end{enumerate}
In addition the characters $\alpha,\beta,\vartheta$ on the common
inertia subgroup in item 3 have positive weights $1,e-1,e+1$.
For every actual complete-local lift $\rho$ of $\bar r_0$, put
\[
 H(\rho)=\rho(g_0)+\vartheta(g_0)\rho(g_0^{-1}).
\]
Then $H(\rho)_{00}-H(\rho)_{22}$ is a unit, with exact residue
\[
 \overline{H(\rho)_{00}-H(\rho)_{22}}
       =(\bar\alpha(g_0)-1)(\bar\Psi(g_0)-1).
 \tag{AX.23}
\]
Every algebraic specialization of this unit is nonzero.
\end{proposition}

\begin{proof}
This proposition uses only \texttt{Good}, so also applies when
$D$ was not constructed by AX5.  First $\omega$ is the unique
Teichm\"uller character: its values and those of the character
in AX3 are prime-to-$\ell$ roots with identical reductions.
The injectivity argument in AX4 proves their equality.  In
particular it is unramified away from $\ell$ and has open kernel.

The open kernel of $\psi$ gives a finite abelian image.  In a
finite abelian extension of $\mathbb Q_p$, the prime-to-$p$ part
of inertia has exponent dividing $p-1$.  Here is the precise local
fact needed.  Wild inertia is pro-$p$.  In the tame quotient inertia
is generated topologically by compatible roots of a uniformizer;
arithmetic Frobenius conjugates a tame generator $\tau$ to
$\tau^p$.  This identity follows by lifting the action
$\zeta\mapsto\zeta^p$ on the prime-to-$p$ roots of unity used in
$\tau(\pi^{1/m})=\zeta_m\pi^{1/m}$.  After passing to an abelian
finite quotient it gives $\tau^{p-1}=1$.  Thus an $\ell$-divisor
of its inertia order with $p\ne\ell$ divides $p-1$.
Applying this to the actual image of $\psi|_{G_{A_p}}$ proves
$\ell\mid p-1$ for $p\in S_b$.
The Teichm\"uller character is trivial on $I_{A_p}$ and takes
the value $1$ at Frobenius because $p\equiv1\pmod\ell$.
Inertia and a Frobenius lift generate a dense subgroup of the
decomposition group, so $\omega=1$ on all of $G_{A_p}$.

The second assertion follows from the Frobenius clause of
\texttt{Good}, the preceding paragraph, and $c(\sigma_\ell)=1$.
Put $a=\bar\xi(\sigma_v)$ and $b=\bar\Psi(\sigma_v)$.
The four residual values are $(ab,b,a,1)$; condition (AX.7) says
$ab^j\ne1$ for every integer $j$.  The four cross-$\lambda$
comparisons are $ab\ne b$, $ab\ne1$, $a\ne b$ and $a\ne1$.
They follow respectively from $j=0,1,-1,0$.  This lists every
comparison, including their reverses.

For $\operatorname{UnrT}$, a place away from $T$ is away from
$\ell$.  There $c=\omega=1$ on inertia, so
$\alpha=\xi$ and $\Psi=\psi$.  The residual-kernel condition
gives $\bar\xi=\bar\psi=1$.  The order of $\xi$ is prime to
$\ell$, so $\xi=1$ by AX4.  If the place is away from $S$,
$\psi=1$ directly.  Otherwise it lies above some $p\in S$
with $p\notin S_b$; conjugacy of places identifies the inertia
image with the selected one, whose order is prime to $\ell$.
Thus $\psi=1$ there as well.  All $p_i$ equal $1$.
On the common inertia subgroup above $\ell$, both $\psi$ and
$\omega$ are $1$, and $\xi=1$ by \texttt{Good}; substitution in
(AX.5) gives the weights, including the three positive ones.
For $p\in S_b$, $\xi$ is unramified and the preceding
$c=\omega=1$ identities give item 3.

The pairing identities $p_i p_{3-i}=\vartheta$ follow from
(AX.5).  Distinctness at $s_0$ follows from the infinite order
of $c(s_0)$ and the four distinct weights for $e\ge2$.
Complex conjugation has $c=\omega=-1$, while $\xi=-1$ and
$\Psi=1$, proving the three oddness assertions.
The closed decomposition group contains all powers of
$\sigma_\ell$, hence their limit $\phi$.
Reduction modulo any power of $\mathfrak m_O$ of a cyclotomic
character is locally constant: a sufficiently deep cyclotomic
principal subgroup maps into that ideal power.  The finite
characters have open kernels.  Products and powers therefore
give the asserted adic continuity of all $p_i$, and reduction
of the diagonal representation has open kernel.
Its image at $\sigma_\ell$ has order dividing $|k|-1$.
AX3 proves $\bar r_0(\phi)=\bar r_0(\sigma_\ell)$.
Outside $S$, the finite characters and the cyclotomic character
are all unramified, since $\ell\in S$.  This finishes every
clause of the literal source \texttt{facts}.

Finally write $a=\bar\alpha(g_0)$ and $b=\bar\Psi(g_0)$.
The diagonal entries $0,2$ of the residual $H$ are respectively
$ab+1$ and $a+b$.  Their difference is $(a-1)(b-1)$, which is
nonzero by \texttt{Good}.  An element of a local ring with
nonzero residue is a unit.  Ring maps carry units to units,
so its image in $\overline{\mathbb Q}_\ell$ is nonzero.  This
proves (AX.23), with no extra normalization premise on $\rho$.
\end{proof}

\begin{proposition}[AX7: the literal common level and exponent]
\label{prop:ax-level-integer}
Suppose the same good $D$ has
$\mathtt{OfChar}(e,\psi,n',p,m_0)$, with $c_2\le n'$ and $m_0\mid o$.
Define
\[
             L=\ell 2^{n'}Np,
             \qquad k_0=(\ell-1)2^{n'}(2o).
             \tag{AX.24}
\]
Then $L,k_0>0$, each of $D.\psi,D.\xi,D.\omega$ has $k_0$-th
power one, and each satisfies the literal
$\operatorname{KillsMu}(L)$: it is $1$ on every $g$ fixing all
$L$-th roots of unity in $\Omega$.  For the data of AX5,
$D.S=\operatorname{primeFactors}(2\ell Np)$ has exactly the
prime support required by G1, and the parameters are (AX.19).
\end{proposition}

\begin{proof}
The levels $N,p,2^{n'}$ divide $L$.  Their cyclotomic actions
are therefore trivial whenever the action on $\mu_L$ is trivial.
The Teichm\"uller character factors through $\mu_\ell$, also
contained in $\mu_L$.  This proves all three \texttt{Kills\allowbreak Mu}
assertions on the same $G$.

The dyadic component of $\psi$ has order dividing $2^{c_2}$,
and every odd component has order dividing $2o$, by the CRT
argument of AX1.  Hence $\psi^{2^{n'}(2o)}=1$.
The prime factor $\xi_0$ has order $m_0\mid o$ and the dyadic
factor has order dividing $\varphi(2^{n'})\mid2^{n'}$.
Thus $\xi^{2^{n'}o}=1$.  Finally $\omega^{\ell-1}=1$.
Each exponent divides $k_0$, proving the literal \texttt{lev\allowbreak Of}
claims.  Positivity follows from $\ell>2$, $o,N,p>0$.
AX5 supplies $m_0=1$ in Case II and
$m_0=\operatorname{ord}\psi\mid\operatorname{oddpart}(o)$ in
Case III, together with the required $c_2\le n'$.
The assertion about $S$ is the exact equality (AX.15).
\end{proof}

\subsection{Ordinary flags on the component containing the diagonal point}

The source's local condition is a polynomial condition, whose
meaning matters under specialization.  For scalars $a_i$ put
\[
 P_i(X)=\prod_{j\ne i}(X-a_j1),\qquad
                  d_i=\prod_{j\ne i}(a_i-a_j).
 \tag{AX.25}
\]
The condition $\operatorname{Ord}(I,(p_i),s_0,r)$ says that
$r(s_0)$ has characteristic polynomial $\prod_i(T-p_i(s_0))$,
that $P_j(r(s_0))r(g)P_i(r(s_0))=0$ for $i<j$, and that
$P_i(r(s_0))(r(u)-p_i(u)1)P_i(r(s_0))=0$ for $u\in I$.
These are exactly the three source fields, with unnormalized
projectors, so they make sense even before $d_i$ is invertible.

An actual point of the ordinary component is a specialization
of a lift $\rho:G_{A_\ell}\to M_4(B)$, where $B$ is a source
complete Noetherian local $O$-algebra which is a domain,
$\rho$ satisfies the similitude and ordinary equations, and there
is an $O$-algebra map $x:B\to O$ carrying $\rho$ to
$r_0|_{G_{A_\ell}}$.  A second map
$y:B\to\overline{\mathbb Q}_\ell$ gives the point.
An intervening coefficient algebra in \texttt{Pt\allowbreak On}/\texttt{On\allowbreak Comp}
is absorbed by composing its map with $y$; the domain $B$ and
the base point $x$ are retained.

\begin{proposition}[AX8: the full ordinary specialization theorem]
\label{prop:ax-ordinary}
Let $D$ be good.  For every actual
$r:G_{A_\ell}\to M_4(\overline{\mathbb Q}_\ell)$ on the source
ordinary component, there is a basis making every $r(g)$ upper
triangular with diagonal continuous characters $\kappa_i$, such
that
\[
\begin{gathered}
 \kappa_i(u)=p_i(u)\quad(u\in I_{A_\ell}),\qquad
 |\kappa_i(g)-p_i(g)|_\ell<1\quad(g\in G_{A_\ell}),\\
 \kappa_0\kappa_3=\kappa_1\kappa_2=\vartheta.
\end{gathered}\tag{AX.26}
\]
For every open subgroup $U\le G_{A_\ell}$ containing
$\sigma_\ell$ and every $i\ne j$, there is $g\in U$ with
$\kappa_i(g)\ne c(g)\kappa_j(g)$.  The same nonidentity holds
on $G_{A_\ell}$ also for $i=j$.  These assertions supply the
literal \texttt{sp\allowbreak Pt\allowbreak L} interface, including its uniformity over
all such open subgroups.
\end{proposition}

\begin{proof}
We first prove the specialization estimate used in the argument.
Let $R$ be a local $O$-algebra with $\mathfrak m_O R\subset
\mathfrak m_R$, and let $y:R\to\overline{\mathbb Q}_\ell$ be an
$O$-algebra map.  Then
\[
 |y(a)|_\ell\le1\ (a\in R),\qquad
 |y(a)|_\ell<1\ (a\in\mathfrak m_R).
 \tag{AX.27}
\]
For the first assertion, if $|y(a)|>1$, the inverse of $y(a)$
has a monic minimal polynomial over $\operatorname{Frac}(O)$
whose nonleading coefficients lie in $\mathfrak m_O$.
Indeed all conjugates have the same norm less than $1$, by
uniqueness of the extension of the complete-field valuation;
their elementary symmetric functions lie in $O$ and have norm
less than $1$.  Multiplying the polynomial identity by a power
of $y(a)$ gives $1+b_1y(a)+\cdots+b_ny(a)^n=0$ with
$b_i\in\mathfrak m_O$.  The corresponding element
$1+b_1a+\cdots+b_na^n$ of $R$ is a unit and cannot map to zero.
If $a\in\mathfrak m_R$ and $|y(a)|=1$, instead use its monic
minimal polynomial.  All coefficients are in $O$ and the constant
coefficient is a unit.  Evaluating the polynomial at $a$ again
gives a unit of $R$ mapping to zero.  This proves (AX.27).
The arguments apply also to zero images, which already satisfy
the asserted inequalities.

If $R$ is Noetherian, choose generators $a_1,\ldots,a_b$ of
$\mathfrak m_R$ and put $a=\max_i|y(a_i)|<1$; the zero ideal
case is immediate.  Products of generators and (AX.27) give
$|y(z)|\le a^n$ for $z\in\mathfrak m_R^n$.  Thus an adically
continuous entry function becomes an ordinary continuous
$\overline{\mathbb Q}_\ell$-valued function.  In particular all
specialized matrix entries below are continuous.

We also need the following precise divided-eigenvalue argument.
Let $B$ be the domain above, $d\in O\setminus\{0\}$,
$t\in B$, $x(t)=0$, and let $v=t/d\in\operatorname{Frac}(B)$
be integral over $B$.  The base point $x$ is a retraction on
$O$, so $O\to B$ is injective and this denominator is nonzero.
The algebra $B[v]$ is finite over the
Henselian local domain $B$.  It is a domain and hence a local
ring by the stated finite-algebra theorem.  The maps $x,y$
extend to the localization $B[1/d]$, with targets
$\operatorname{Frac}(O)$ and $\overline{\mathbb Q}_\ell$.
Restrict them to $B[v]\subset B[1/d]$.  The first sends $v$ to
zero, so $v$ is not a unit and lies in the maximal ideal of
$B[v]$.  Its maximal ideal contracts to that of $B$, since the
extension is finite integral.  Applying (AX.27) to the second
map proves
\[
                      |y(t)/d|_\ell<1.
 \tag{AX.28}
\]
This is the point where the same domain component and its actual
diagonal base point are used; a bare ordinary specialization
would not justify the conclusion.

Now set $a_i=p_i(s_0)$.  AX6 makes them distinct.  Over either
$\operatorname{Frac}(B)$ or $\overline{\mathbb Q}_\ell$ the
operators $d_i^{-1}P_i(r(s_0))$ are pairwise orthogonal rank-one
idempotents summing to $1$.  This follows by evaluating the
Lagrange interpolation identities modulo
$\prod_i(T-a_i)$ and using the characteristic polynomial; each
eigenspace is one-dimensional because its root is simple.
Choose one vector from each eigenline, in index order.
The stable equations imply that the sums of the first $i+1$
lines are invariant under every $r(g)$.  The matrices are
therefore upper triangular.  Their diagonal entries are
characters $\kappa_i$, since multiplication of upper triangular
matrices multiplies the diagonal entries, and their inverse
values exist from $r(g^{-1})$.  The inertia equations give
$\kappa_i(u)=p_i(u)$ exactly.  Moreover
\[
 d_i\kappa_i(g)=\operatorname{tr}\big(P_i(r(s_0))r(g)\big).
 \tag{AX.29}
\]

Use the same argument first over $\operatorname{Frac}(B)$.
Its diagonal value $v_i(g)$ is integral over $B$, since it is
a root of the monic characteristic polynomial of $\rho(g)$.
Write $t_i(g)=\operatorname{tr}(P_i(\rho(s_0))\rho(g))$.
Then $d_i v_i(g)=t_i(g)$ and at the actual diagonal point
$x(t_i(g))=d_i p_i(g)$.
Apply (AX.28) to the integral element $v_i(g)-p_i(g)$,
with numerator $t_i(g)-d_i p_i(g)$.  Equation (AX.29) after
specialization identifies its image with
$\kappa_i(g)-p_i(g)$, and proves the strict estimate in (AX.26).
Continuity of $\kappa_i$ follows from continuity of the entries
of $r$ after a fixed conjugation, or directly from (AX.29).

The original alternating form is preserved up to $\vartheta$.
In the chosen eigenbasis for $r(s_0)$ its matrix can have a
nonzero entry at $(i,j)$ only if
$a_i a_j=\vartheta(s_0)$.
Since $a_i=c(s_0)^{w_i}$ with $c(s_0)$ of infinite order,
this means $w_i+w_j=e+1$.
For the four weights and $e\ge2$ this occurs precisely when
$j=3-i$.  Nondegeneracy makes all four antidiagonal entries
nonzero.  For an upper triangular $r(g)$, the $(i,3-i)$ entry
of $r(g)^tJr(g)=\vartheta(g)J$ then gives
$\kappa_i(g)\kappa_{3-i}(g)=\vartheta(g)$.
This proves the exact two pairing identities in (AX.26).

Finally suppose $\kappa_i=c\kappa_j$ on an open subgroup $U$
containing $\sigma_\ell$.  The subgroup has finite index in
the compact decomposition group, so $s_0^b\in U$ for some
positive $b$.  Evaluation there, using inertia and infinite
order, gives $w_i=w_j+1$.  This is impossible for $i=j$.
For distinct indices the only pairs allowed by the four weights
are $(0,1),(2,3)$ and, when $e=2$, $(1,2)$.
Every such pair has $\lambda_i\ne\lambda_j$.
Evaluation at $\sigma_\ell$ gives
$\kappa_i(\sigma_\ell)=\kappa_j(\sigma_\ell)$ because
$c(\sigma_\ell)=1$.  The strict estimates in (AX.26) imply
$\bar p_i(\sigma_\ell)=\bar p_j(\sigma_\ell)$: the difference
of their values has norm less than $1$, which is precisely
membership in $\mathfrak m_O$.  This contradicts AX6.
Thus the required element of $U$ exists for every pair.
\end{proof}

\subsection{The full admissible class on the same data}

\begin{proposition}[AX9: literal level, dual twist and safety]
\label{prop:ax-admissible}
Let $D$ be any good datum and $M\ge1$.  Suppose
$\operatorname{KillsMu}(M)$ holds for $\psi,\xi,\omega$.
For every actual $D.\mathtt{EPt}$ representation $r$,
\[
 \mathtt{Level}\bigl(\ell,M,(e+1,e,1,0),r\bigr),
 \qquad \mathtt{IsCycWt}(\ell,e+1,\vartheta),
 \qquad \mathtt{DualTw}(r,\vartheta)
 \tag{AX.30}
\]
hold, with the original coefficient embedding into
$\overline{\mathbb Q}_\ell$.  Also $\ell\mid M$ and
$r$ satisfies $\mathtt{UnlAll}$, hence the first alternative of
the literal $\mathtt{SafeRA}$ condition.  This proves
\texttt{level}, \texttt{sp\allowbreak Unl}, \texttt{teich\allowbreak Dvd}, and
\texttt{safe4'} on their actual quantified inputs.
\end{proposition}

\begin{proof}
An actual point is the image under $R\to\overline{\mathbb Q}_\ell$
of an adically continuous lift over a source complete local
$O$-algebra.  Formula (AX.27) and the finite-generator estimate
in AX8 make its matrix entries continuous.  Since
$r(g)^{-1}=r(g^{-1})$, this is also continuity as a
general-linear-group representation.  Its similitude equations,
the equality $r=1$ on the prescribed inertia kernel outside $T$,
and its equations at every $p\in S_b$ survive specialization:
they are matrix polynomial equalities.  For the \texttt{On\allowbreak Comp}
local conditions, compose with the actual domain-component map.
This explains the needed part of the source \texttt{pt\allowbreak In}
transport without assuming a representability theorem.

We prove the away-from-$\ell$ part of \texttt{Level} with global
characters, as its definition requires.  Let
$\varphi=(\psi,\xi,\omega):G\to A$ with $A$ its actual finite
abelian image.  Choose a place over a given rational prime
$p\ne\ell$.  If it is away from $T$, then for
$g\in I_A\cap\ker\varphi$ the cyclotomic character is $1$,
so $\bar r_0(g)=1$ and the point's unramified condition gives
$r(g)=1$.  If it lies above a member of $T$, that member is
some $p\in S_b$; choose the specified $A_p$ instead.
The point's actual $\mathtt{condP}$ equation gives
$r=1$ on $I_{A_p}\cap\ker\psi$, in particular on
$I_{A_p}\cap\ker\varphi$.

Thus in either case the inertia representation factors through
the finite abelian subgroup $\varphi(I_A)\le A$.
Every representing matrix has finite order and is diagonalizable
in characteristic zero, since its minimal polynomial divides a
separable polynomial $T^m-1$.  Commuting diagonalizable matrices
have a common eigenbasis: split into eigenspaces of one nonscalar
matrix and induct on their dimensions; if all are scalar the
claim is immediate.  The diagonal characters of $\varphi(I_A)$
extend to $A$.  Explicitly, when adjoining one generator $a$ to
a subgroup and its coset has order $b$, choose in the algebraic
closure a $b$-th root of the already specified value at $a^b$;
the rule on $ha^j$ is then well-defined and multiplicative.
Iterate through generators of the finite group.  Compose the
four extended characters with $\varphi$ to obtain four global
characters on the original $G$, diagonalizing inertia.  They
are all trivial on $\ker\varphi$, hence satisfy
$\operatorname{KillsMu}(M)$.  This is the complete
$\mathtt{LevelAway}$ assertion; it does not merely produce
local inertia characters.

At $\ell$, AX8 supplies the actual ordinary flag and characters
$\kappa_i$.  Define the following global finite characters:
\[
 (\chi_0,\chi_1,\chi_2,\chi_3)
       =(\omega^{-1}\xi\psi,\psi,\omega^{-1}\xi,1).
 \tag{AX.31}
\]
They all satisfy $\operatorname{KillsMu}(M)$, and
$p_i=c^{w_i}\chi_i$.  The inertia formula in (AX.26) is
therefore exactly the formula in $\mathtt{LevelAt}$.
The last part of AX8 gives
$\kappa_i\ne c\kappa_j$ for every pair, so we use the first
alternative of that definition.  No potentially crystalline
alternative is needed.  Over $\mathbb Q$ there is only one
finite place above $\ell$, and the chosen $A_\ell$ lies above
it; the comparison with the source prime ideal is just
$\mathcal O_{\mathbb Q}=\mathbb Z$ and its ideal $(\ell)$.
This proves $\mathtt{Level}$.

The multiplier is
$\vartheta=c^{e+1}(\omega^{-1}\xi\psi)$, whose finite part
satisfies $\operatorname{KillsMu}(M)$.  This is the actual
$\mathtt{IsCycWt}$ witness with the same $M$.
For the duality witness take the fixed matrix
\[
 J=\begin{pmatrix}0&0&0&1\\0&0&1&0\\0&-1&0&0\\-1&0&0&0\end{pmatrix},
 \qquad J^{-1}=-J.
 \tag{AX.32}
\]
The source adjoint is $J^{-1}X^tJ$; its similitude equation is
therefore exactly $r(g)^tJr(g)=\vartheta(g)J$.
This supplies the original $\mathtt{DualTw}$, with no change of
Frobenius or cyclotomic convention.

To see that $\ell\mid M$, suppose the contrary.  Inertia at
$A_\ell$ fixes $\mu_M$, so $\omega=1$ there by
$\operatorname{KillsMu}(M)$.  But AX3 supplies an inertia element
with $c=-1$, whose Teichm\"uller value is $-1$.  Since $\ell$
is odd this is a contradiction.  The noncritical ordinary flag
just proved holds for every place above $\ell$, not only the
selected one: if $A'=gA_\ell$, transport the flag by $r(g)$
and the characters by $u\mapsto g^{-1}ug$.
The global cyclotomic character is invariant under conjugation,
so all inequalities $\kappa_i\ne c\kappa_j$ are preserved.
The transport map on decomposition groups is an isomorphism,
so no nonidentity is lost.  This is precisely
$\mathtt{UnlAll}$.  Combining it with $\ell\mid M$ proves
the first clause of $\mathtt{SafeRA}$, and hence the literal
\texttt{safe4'} conclusion.
\end{proof}

\begin{proposition}[AX10: the complete G1 implication]
\label{prop:ax-g1}
Assume the image-density hypothesis, with the literal
\texttt{Symplectic\allowbreak Four\allowbreak Image\allowbreak Dense} statement.
For every original host and original Case II or Case III
character in AX5, adjoining the roots already allowed by SP8
produces an actual $D$ satisfying every clause of G1 in SP10.
Its parameters and level are exactly (AX.19) and (AX.24).
Moreover $D.\mathtt{Fin0}$ and $D.\mathtt{NoEnd}$ hold, and there
is one embedding
$\iota:\overline{\mathbb Q}_\ell\simeq\mathbb C$, chosen before
the varying point, such that every actual irreducible
$D.\mathtt{EPt}$ representation belongs to
\[
 \mathtt{AdmRAlt}\bigl(\iota,L,(e+1,e,1,0),e+1\bigr).
 \tag{AX.33}
\]
There is no remaining construction, character, residue-field,
level, duality, ordinary-specialization or safety premise in
this implication.  The actual image-density supplier itself
remains an arithmetic obligation.
\end{proposition}

\begin{proof}
First suppose the original $O$ already has the two required
sets of dyadic roots.  AX5 constructs $D$ on its original
$\Omega,O$ and the exact \texttt{Of\allowbreak Case}.  AX1 and AX5 provide
$c_2\le n'$ and $m_0\mid o$, so AX7 gives the common $L,k_0$
and all three $\mathtt{KillsMu}$ and power assertions.
AX9 applies to every actual point of this same $D$, giving
its level, multiplier of weight $e+1$, dual twist and safety.
Now apply \texttt{Symplectic\allowbreak Four\allowbreak Image\allowbreak Dense} to this $D$, with
the same host, $t,e,\psi$ and original \texttt{Good}/\texttt{Of\allowbreak Case}
witnesses.  Its conclusion is exactly
\[
 D.\mathtt{Fin0}\ \wedge\ D.\mathtt{NoEnd}\ \wedge\
 \exists\iota\ \forall r\in D.\mathtt{EPt},\quad
 \mathtt{Irr}(r)\Longrightarrow\mathtt{IsAutQA}(\iota,r).
 \tag{AX.34}
\]
Here $\mathtt{Irr}$ has the source meaning that the linear span
of the matrices $r(g)$ is all of $M_4$.  It implies the
$\mathtt{Af.IsIrr}$ field of the admissible record: an invariant
subspace would be invariant under every matrix, and a nonzero
vector can be carried to every vector by suitable matrices.
Thus the subspace is the whole vector space.  This is the exact
direction used for the actual points in SP10.
The quantifiers in (AX.34), rather than separately chosen
embeddings for each point, supply the common $\iota$.
Irreducibility, the stated automorphy assertion, and the four
clauses already proved are precisely the fields of
$\mathtt{AdmRAlt}$ in (AX.33).

For a coefficient DVR initially lacking the roots, use the
finite extension allowed in SP8 and RV1.  Its integral closure
inside the same $\overline{\mathbb Q}_\ell$ is a finite
coefficient DVR containing the two finite sets of roots.
The injective coefficient map preserves the original
character's kernel, order, conductor, CRT components, parity,
$f_0,n_0,c_2$, and hence both integers in (AX.3).
These facts can also be read directly: equality of two
character values is unchanged by an injective ring map,
so every finite quotient through which it factors is unchanged.
Construct $D$ over this allowed extension and perform the
preceding argument there.  This does not assert that the
new $D$ is a datum over the smaller coefficient ring.
The final normalized Selmer-length comparison is returned to
the original coefficient ring by the exact ramification-factor
cancellation already proved in RV1 and SP8.  The global field,
host, character case and integer bounds have stayed fixed.
Thus the conditional G1 statement used by SP10 is completely
supplied, with its original coefficient-return mechanism.
\end{proof}

\paragraph{The resulting boundary.}
AX5 and AX6 are unconditional arithmetic suppliers relative to
the explicitly listed classical facts and the reconstructed BR7.
AX7--AX9 are implications on any actual good datum, with no
image-density premise.  AX10 uses
\texttt{Symplectic\allowbreak Four\allowbreak Image\allowbreak Dense} exactly once for (AX.34).
It does not supply the automorphic catalogue, actual model
construction, geometric point comparison or rigidity.  Those
are the separately recorded G2 and catalogue chains.  In
particular the source formulas that later produce
\texttt{Symplectic\allowbreak Four\allowbreak Image\allowbreak Dense} from automorphic hypotheses
are continuation targets, not classical consequences of
the elementary construction above.


% Version 0.27. Actual symplectic local and cohomological installation.
% Pinned source: 7cbf7a0c32fca1badf05931879a37c3f8a2b03b4.
\section{Installing the arithmetic symplectic deformation data}
\label{sec:symplectic-local-installation}

\status{DX1--DX8 close the three actual arithmetic inputs to
\texttt{ex\allowbreak Model}: \texttt{Phi\allowbreak Cond}, \texttt{Nums}, and an original
\texttt{Loc\allowbreak Base}.  The ordinary and finite-inertia charts are proved
by continuous cohomology and explicit matrix deformations.
Their smoothness is at the characteristic-zero origin, precisely
the source's \texttt{Smooth\allowbreak Pt}; integral formal smoothness is not
asserted.  The full \texttt{Bloch\allowbreak Kato\allowbreak Local} package is not assumed
or promoted.  The original \texttt{Good} datum remains the input,
supplied by the separate arithmetic construction.}

\subsection{Original fields, characters and roots}

Fix the source's original algebraic closure $\Omega/\mathbb Q$,
odd prime $\ell$, coefficient DVR $O=c.O$, fraction field $E$,
and finite residue field $k$.  Fix a \texttt{Raw} datum $D$
satisfying \texttt{Good}, without replacing its places or
decomposition subgroups.  Write
\[
 \begin{aligned}
 \chi&=\chi_\ell,&
 \alpha&=\chi\omega^{-1}\xi,&
 \Psi&=\chi^e\psi,\\
 \beta&=\Psi\alpha^{-1},&
 \theta&=\alpha\Psi,&
 r_0&=\operatorname{diag}(\theta,\Psi,\alpha,1).
 \end{aligned}
 \tag{DX.1}
\]
Here $e\ge2$, $\psi,\omega,\xi$ have finite image, $\xi$ is
unramified at every member of $T=\{\ell\}\cup S_b$, and
$\omega$ is the specified residual cyclotomic lift.
All further arithmetic facts used below are the original
\texttt{facts(D,Good)} supplied by AX6.

Let $G_v=D.Gp(v)$ and $I_v=D.Ip(v)$, and retain the original
chosen $s_0\in I_\ell\cap\ker\psi\cap\ker\omega$.
Its cyclotomic value has infinite order.  Thus the four eigenvalues
of $r_0(s_0)$ are distinct, since their exponents are
$e+1,e,1,0$.  At every $v\in T$ the original $\sigma_v=D.sg(v)$
has the stated residual separation; at $\ell$,
$\chi(\sigma_\ell)=1$.

Use the weights
\[
 \begin{aligned}
 a&=(1,0),& b&=(-1,1),& p&=(0,1),& t&=(1,1),\\
 w_0&=(1,1),& w_1&=(0,1),& w_2&=(1,0),& w_3&=0 .
 \end{aligned}
 \tag{DX.2}
\]
The letter $p$ in this weight list denotes $w_{\rm Ps}$;
prime subscripts below are written $q$ when ambiguity is possible.
With the original alternating form
\[
 J=\begin{pmatrix}
 0&0&0&1\\0&0&1&0\\0&-1&0&0\\-1&0&0&0
 \end{pmatrix},
\]
the Lie algebra $\mathfrak{sp}_4$ has the following weight basis:
\[
\begin{array}{c|c|c}
\text{weight}&\text{basis matrix}&\text{adjoint character}\\ \hline
0&\operatorname{diag}(1,0,0,-1),\
       \operatorname{diag}(0,1,-1,0)&1\\
a&E_{01}-E_{23}&\alpha\\
p&E_{02}+E_{13}&\Psi\\
b&E_{12}&\beta\\
t&E_{03}&\theta\\
-a&E_{10}-E_{32}&\alpha^{-1}\\
-p&E_{20}+E_{31}&\Psi^{-1}\\
-b&E_{21}&\beta^{-1}\\
-t&E_{30}&\theta^{-1}
\end{array}
 \tag{DX.3}
\]
Direct multiplication verifies $Y^tJ+JY=0$ for these matrices;
that equation also determines all other entries from the displayed
ten.  The table is therefore valid over $O$, $E$, and $k$.
Put
\[
 d_P(j)=
 \begin{cases}2&j=0,\\1&j\in\{\pm a,\pm b,\pm p,\pm t\},\\0&\text{otherwise},
 \end{cases}
 \quad
 d_L(j)=
 \begin{cases}2&j\in\{0,a,b,p,t\},\\1&j\in\{-a,-b,-p,-t\},\\0&\text{otherwise}.
 \end{cases}
 \tag{DX.4}
\]
These are literally \texttt{dim\allowbreak P} and \texttt{dim\allowbreak L}.

\subsection{The exact finite and rational cohomology inputs}

% BEGIN DX1
\begin{proposition}[DX1: original local duality, global Euler formula and normalized finite duality]
\label{claim:DX1}
The native \texttt{Local\allowbreak Finiteness}, \texttt{Local\allowbreak Tate\allowbreak Duality},
\texttt{Local\allowbreak Euler\allowbreak Char}, \texttt{Global\allowbreak Finiteness} and
\texttt{Global\allowbreak Euler\allowbreak Char} interfaces follow from the finite local
and global theorems stated below, with their original coefficient
modules, places and cup products.  LD3 supplies the literal
\texttt{Poitou\allowbreak Tate\allowbreak Normalised} by forgetting its additional
conjugation-compatibility field.
These universally quantified interface assertions range over
every prime, number field and finite coefficient module
allowed by their definitions; they do not inherit an
odd-prime restriction from the later good-datum application.

In particular, for a finite $k[G_{\mathbb Q,S}]$-module $M$,
with $\ell\in S$, write $h^i=\dim_kH^i$ and let $c$ be the
original complex conjugation.  Then
\[
 h^1_{\rm glob}+ \dim_k M^{c=1}
       =h^0_{\rm glob}+h^2_{\rm glob}+\dim_k M,
 \tag{DX.5}
\]
and on the original $G_v$,
\[
 h^1_v=h^0_v+h^2_v+
                  {\bf1}_{v=\ell}\dim_kM .
 \tag{DX.6}
\]
For a continuous unit-valued character
$\kappa:G_\ell\to O^\times$, rational continuous cohomology is
finite dimensional and
\[
 \dim_EH^2(G_\ell,E(\kappa))
       =\dim_E H^0(G_\ell,E(\kappa^{-1}\chi)),\qquad
 \dim_EH^1=\dim_EH^0+\dim_EH^2+1 .
 \tag{DX.7}
\]
All coefficient twists in these formulas use the original
cyclotomic action.
\end{proposition}

\begin{proof}
The named classical inputs are finite local Tate duality and the
local Euler--Poincare formula, global finiteness, and the global
Euler--Poincare formula.  Precise primary references are given
at the end of the section.  We spell out their translation,
including the rational passage needed for the charts.

For a valuation subring $A$ over a finite place $v$ of an arbitrary
number field $F$, the original decomposition group is the absolute
Galois group of the corresponding henselization.  Finite separable
extensions of that henselian valued field and of its completion
correspond, preserving the embeddings into the chosen separable
closures.  This gives the topological Galois-group identification
with $G_{F_v}$ used in KB6.  Its inertia subgroup is the original
inertia subgroup.  It also preserves roots of unity: the action
on $\mu_{\ell^r}$ is exactly the pullback of
$\operatorname{cycMod}(F,\Omega,\ell,r)$.
Precomposition consequently gives isomorphisms of the actual
continuous cochain complexes, commuting with restriction and cup
products.  No unrelated coefficient or local field is substituted.

For $R=\mathbb Z/\ell^r$ and any finite, possibly nonfree, $R$-module
$M$, evaluation and a chosen primitive root identify the source's
twisted dual with $\operatorname{Hom}(M,\mu_{\ell^r})$.
The finite local theorem gives finiteness, vanishing above degree
two, and perfect pairings in degrees $(0,2),(1,1),(2,0)$.
Kummer theory identifies $H^2(R(1))$ with the $\ell^r$-torsion
in the local Brauer group.  Its invariant identifies that target
with $R$.  Thus perfection gives the source's maps into the
dual \emph{before} selecting the invariant as well.  The reversed
degree-one map is also bijective: graded commutativity changes
only a sign.  This proves every native local-duality clause.
It is the same finite-coefficient argument as LK1--LK3,
now transported to the actual decomposition groups.

The finite local Euler formula is
\[
 \frac{|H^0(F_v,M)|\,|H^2(F_v,M)|}{|H^1(F_v,M)|}
                   =|\mathcal O_{F_v}/|M|\mathcal O_{F_v}|^{-1}.
 \tag{DX.8}
\]
If $v\nmid\ell$ the right side is one.
If $v\mid\ell$ it is $|M|^{-[F_v:\mathbb Q_\ell]}$:
compute successively along the residue-field and ramification
filtrations, exactly as in LK4.  This proves the native
\texttt{loc\allowbreak Deg} exponent even when $M$ is not free.

For the global assertions, the quotient by the closed normal
subgroup generated by the omitted inertia groups is precisely
$G_{F,S}$: the assertion is tested on finite quotients, where it
is the definition of being unramified outside $S$.
The classical finite set includes all archimedean places; this
places no additional restriction on that quotient.
Because $S$ contains the places above $\ell$, $|M|$ is a unit
in the ring of $S$-integers.
The global finiteness theorem therefore applies.  Its
high-degree comparison is with the real-place groups, which
vanish for odd $\ell$, proving the exact extra clause in
\texttt{Global\allowbreak Finiteness}.

The global Euler formula reads
\[
 |H^0|\,|H^2|\,|M|^{[F:\mathbb Q]}
   =|H^1|
     \prod_{v\ {\rm real}}|M^{c_v=1}|\,
                       |M|^{r_2(F)}.
 \tag{DX.9}
\]
Indeed a real normalized absolute value sends $|M|$ to $|M|$,
and a complex one sends it to $|M|^2$.
This is exactly \texttt{arch\allowbreak H0Card}; it also covers $\ell=2$
in the stated Euler interface, without claiming high-degree
vanishing at real places in that case.
Milne explicitly defines the Euler expression immediately
before I.5.1 by the truncated ratio $|H^0||H^2|/|H^1|$,
including the real-place, $2$-primary case where higher
cohomology need not vanish.

The source's \texttt{Is\allowbreak Local\allowbreak Invariants} contains bijectivity and
the uniformizer normalization required by
\texttt{Poitou\allowbreak Tate\allowbreak Normalised}, together with conjugation
compatibility.  The rest of their quantified conclusions is
identical: the invariant family is chosen before $S$, the places,
the descended cyclotomic character and $M$.
Forget that one extra field in LD3.
This is a literal implication between the two definitions,
not an identification based on their names.
For a finite coefficient field $k$, restricting scalars to
$\mathbb F_\ell$ and using its nondegenerate trace identifies
the ordinary finite dual with the $k$-linear twisted dual.
All dimensions over $\mathbb F_\ell$ acquire the same factor
$[k:\mathbb F_\ell]$.  Taking cardinalities in (DX.8)--(DX.9)
gives (DX.5)--(DX.6), while the same trace identification
preserves the finite Poitou--Tate sequence.

For completeness, let $T=O(\kappa)$ and choose a uniformizer
$\varpi$ of $O$.  Continuous cochains with values in $T$ are
the inverse limit of those with values in $T/\varpi^nT$.
The cochain transition maps are onto: a locally constant map
into a finite quotient has a locally constant lift.
All finite-level cohomology groups are finite.  Their inverse
systems satisfy the Mittag--Leffler condition, so
\[
 H^i(G_\ell,T)=\varprojlim_nH^i(G_\ell,T/\varpi^nT).
 \tag{DX.10}
\]
The connecting sequence embeds
$H^i(T)/\varpi H^i(T)$ into the finite group $H^i(T/\varpi T)$.
Choose finitely many lifts of generators of this quotient.
Successive $\varpi$-adic approximation shows that they generate
$H^i(T)$: the corrections tend to zero in every finite-level
projection, and the image of a finite power of compact $O$
is closed.  Thus these compact cohomology groups are finite
$O$-modules.

A continuous cochain into $E(\kappa)$ has compact, hence bounded,
image on the compact source $G_\ell^i$; one power of $\varpi$
puts it into $T$.  Consequently rational cohomology is
$H^i(T)\otimes_OE$.
The compatible finite Tate pairings identify the Pontryagin dual
of $H^2(T)$ with the invariants of its discrete dual lattice.
Using the $\mathbb Z_\ell$-dual lattice rather than an
unjustified self-duality of $O$ introduces the codifferent,
but after tensoring with $E$ its character is still
$\kappa^{-1}\chi$.
If this character is not one, one nonzero scalar
$\kappa(g)^{-1}\chi(g)-1$ kills all its discrete invariants;
they are therefore finite.  If it is one, the invariants are
the entire divisible rank-one module.  The $E$-dimension of
$H^2$ is accordingly zero or one, giving the first formula
in (DX.7).

Finally the exact sequences
\[
 0\to H^i(T)/\varpi^nH^i(T)
 \to H^i(T/\varpi^nT)
 \to H^{i+1}(T)[\varpi^n]\to0
\]
show that the logarithm to base $|k|$ of the middle cardinality
is $n\operatorname{rank}_OH^i(T)+O(1)$.
Apply (DX.8), divide by $n$, and let $n$ increase.
The coefficient length of $T/\varpi^nT$ is $n$, regardless of
the ramification of $O/\mathbb Z_\ell$.
This proves the second formula in (DX.7) with the displayed
coefficient-field normalization.
\end{proof}
% END DX1

% BEGIN DX2
\begin{proposition}[DX2: actual global and local \texttt{Phi\allowbreak Cond}]
\label{claim:DX2}
For the original good datum,
$\Phi_\ell(D.Gam)$ holds.  It also holds for every original
decomposition group $D.Gp(v)$ at a finite prime $v$.
These conclusions require no Bloch--Kato theorem.
\end{proposition}

\begin{proof}
We use a concrete finite-cochain argument that avoids changing the
original group.  Let $\Gamma$ be a profinite group for which
$H^1(\Gamma,M)$ is finite for every finite continuous
$\mathbb F_\ell[\Gamma]$-module.  If $H\le\Gamma$ is open,
the coset set $X=\Gamma/H$ is finite and discrete.
Choose representatives $s(x)$, taking $s(H)=1$.
Give $P=\mathbb F_\ell^X$ the permutation action
$(g f)(x)=f(g^{-1}x)$.
For a continuous additive homomorphism $u:H\to\mathbb F_\ell$,
put
\[
       c_u(g)(x)=u\bigl(s(x)^{-1}g\,s(g^{-1}x)\bigr).
 \tag{DX.11}
\]
The argument belongs to $H$.  Multiplication of two such
arguments proves $c_u(gh)=c_u(g)+g c_u(h)$.
There are only finitely many cosets, so this is a continuous
cocycle.  Evaluation at $(h,H)$ gives $u(h)$, making
$u\mapsto c_u$ injective.
The cocycle group $Z^1(\Gamma,P)$ is finite: its quotient
$H^1(\Gamma,P)$ is finite and its coboundary subgroup is an
image of finite $P$.  Hence there are only finitely many $u$.
Open-kernel homomorphisms to the finite target are continuous,
so this is exactly the source's \texttt{Phi\allowbreak Cond}.

For $D.Gam=G_{\mathbb Q,S}$, the finite set $S$ consists of the
original rational primes and contains $\ell$.  Apply global
finiteness from DX1.  For $D.Gp(v)$ use the original
decomposition/completion identification in DX1 and finite
local cohomology, also for modules of order prime to $v$.
This proves both assertions with the stated original topologies.
\end{proof}
% END DX2

\subsection{The complete framed numerical calculation}

% BEGIN DX3
\begin{proposition}[DX3: finite Selmer duality and the actual framed Euler identity]
\label{claim:DX3}
Finite Poitou--Tate duality, with DX1's matched invariants, gives
the literal \texttt{Greenberg\allowbreak Wiles} identity for arbitrary finite
$\mathbb Z/\ell^r$-modules and arbitrary local submodules, when
$\ell$ is odd.

For the original rational field and a frame with
$S=T\sqcup Q$, $\ell\in T$, distinguished $v_0\in T$, put
$N_q=I_q\cap K$.  Suppose $K$ is normal, the action on $N_q$
is trivial, and $[I_q:N_q]$ is prime to $\ell$.
For a finite $k$-representation $M$ let $H^1_F(M)$ and
$\operatorname{Ob}_F(M)$ be the exact cochain spaces defined in
FP.4--FP.6, with the literal intersection quotient for
$\operatorname{Ob}$.  They are finite dimensional and
\[
 \dim\operatorname{Ob}_F(M)-\dim H^1_F(M)
       =\dim M^{c=1}+\dim H^0(G_{\mathbb Q,S},M^\vee(1)).
 \tag{DX.12}
\]
\end{proposition}

\begin{proof}
First fix an arbitrary number field and odd $\ell$, a finite
coefficient module $A$, and local submodules
$L_v\subset H^1(F_v,A)$.  Let $L_v^\perp$ be their exact
annihilators under the original local cup product.
Finite perfect duality and the middle and final terms of
LD3's nine-term sequence give the exact sequence
\[
\begin{aligned}
0\to\operatorname{Sel}_{L}(A)&\to H^1(G_{F,S},A)
 \to\bigoplus_{v\in S}H^1(F_v,A)/L_v\\
 &\to\operatorname{Sel}_{L^\perp}(A^\vee(1))^\vee
 \to H^2(G_{F,S},A)\\
 &\to\bigoplus_{v\in S}H^2(F_v,A)
 \to H^0(G_{F,S},A^\vee(1))^\vee\to0 .
\end{aligned}
 \tag{DX.13}
\]
To see the only modified middle term, divide the local
degree-one term in the Poitou--Tate sequence by $\bigoplus L_v$.
Its annihilator in the global dual cohomology is exactly the
inverse image of $\bigoplus L_v^\perp$.  Restriction of
functionals identifies the quotient of its dual by the
annihilator with the displayed Selmer dual.
Dualizing finite modules here is exact: the
$\mathbb Z/\ell^r$-linear dual is the group-character dual
under $a\mapsto a/\ell^r$, including nonfree modules.

Take cardinalities in (DX.13).  Apply (DX.8)--(DX.9) and
$\sum_{v\mid\ell}[F_v:\mathbb Q_\ell]=[F:\mathbb Q]$.
All powers of $|A|$ cancel, leaving
\[
\begin{aligned}
|\operatorname{Sel}_L(A)|\,|H^0(A^\vee(1))|
 \prod_{v\in S}|H^0(F_v,A)|\,\operatorname{archH0Card}(A)
={}&|\operatorname{Sel}_{L^\perp}(A^\vee(1))|\\
&{}\cdot |H^0(A)|\prod_{v\in S}|L_v|.
\end{aligned}
 \tag{DX.14}
\]
This is the full native \texttt{Greenberg\allowbreak Wiles} formula.
All local maps and annihilators were taken on the original
places; no alternate normalization is required.

We now prove the framed identity without treating the framed
obstruction space as an unnamed Selmer complex.
Write $U_q=H^1_{\rm ur}(G_q,M;N_q)$, and define
\[
 \begin{aligned}
 \operatorname{Sel}&=\ker\left[
 H^1(G_{\mathbb Q,S},M)\to
  \bigoplus_{v\in T}H^1(G_v,M)
       \oplus\bigoplus_{q\in Q}H^1(G_q,M)/U_q\right],\\
 \operatorname{Sha}^2&=\ker\left[
 H^2(G_{\mathbb Q,S},M)\to
                     \bigoplus_{v\in S}H^2(G_v,M)\right].
 \end{aligned}
 \tag{DX.15}
\]
The map from $H^1_F$ to $\operatorname{Sel}$ sends a cocycle
and its local primitives to its global cohomology class.
It is onto: choose local primitives at every $v\in T$ and
subtract the coboundary of the primitive at $v_0$.
At $q$, vanishing of the restricted class on $N_q$ is actual
vanishing of that cocycle, since $N_q$ acts trivially.
The kernel is parametrized by the local invariant vectors,
modulo one common global invariant vector.  Explicitly, for
a global coboundary $d^0m$, its primitives have the form
$m+b_v$, with $b_v\in M^{G_v}$; the normalization at $v_0$
determines $m=-b_{v_0}$, and the remaining redundancy is
$M^{G_{\mathbb Q,S}}$.  Thus
\[
 \dim H^1_F+h^0_{\rm glob}
             =\dim\operatorname{Sel}+\sum_{v\in T}h^0_v .
 \tag{DX.16}
\]

The map $\operatorname{Ob}_F\to H^2(G_{\mathbb Q,S},M)$
sends a triple to its first cocycle class.
Its image is exactly $\operatorname{Sha}^2$, since the local
components are precisely chosen primitives for its restrictions.
If that first class is zero, subtract its global primitive.
The remaining local entries are degree-one cocycles.
Quotienting by FP.6 divides the $T$-components by coboundaries,
the $Q$-components by $U_q$, and then divides by the image
of global $H^1$.  Conversely every such boundary gives
exactly those identifications.  Hence
\[
0\to\operatorname{coker}\left[
 H^1_{\rm glob}\to
 \bigoplus_T H^1_v\oplus\bigoplus_Q H^1_q/U_q\right]
 \to\operatorname{Ob}_F\to\operatorname{Sha}^2\to0 .
 \tag{DX.17}
\]
This proves finiteness as well as the dimension formula
\[
 \dim\operatorname{Ob}_F+h^1_{\rm glob}+\sum_Q\dim U_q
 =\dim\operatorname{Sha}^2+\sum_T h^1_v+\sum_Q h^1_q
                              +\dim\operatorname{Sel}.
 \tag{DX.18}
\]

Restriction from $I_q$ to $N_q$ has zero kernel on $H^1$:
inflation--restriction identifies that kernel with
$H^1(I_q/N_q,M)$, which is zero by averaging over the
finite group of order prime to $\ell$.
Thus $U_q$ is the usual unramified subgroup.
The quotient $G_q/I_q=\widehat{\mathbb Z}$ gives
$U_q=M^{I_q}/(\operatorname{Frob}_q-1)M^{I_q}$, so
$\dim U_q=h^0_q$ by elementary rank-nullity.
Since $q\ne\ell$, (DX.6) gives $h^1_q-\dim U_q=h^2_q$.
The sum of (DX.6) over $T$ contributes exactly one
$\dim M$, because $\ell\in T$.

Finally the last three nonzero terms of finite Poitou--Tate give
\[
 \dim\operatorname{Sha}^2+\sum_{v\in S}h^2_v
       =h^2_{\rm glob}+\dim H^0(M^\vee(1)).
 \tag{DX.19}
\]
Subtract (DX.16) from (DX.18), then use these local identities,
(DX.19) and (DX.5).  The result is exactly (DX.12).
\end{proof}
% END DX3

% BEGIN DX4
\begin{proposition}[DX4: the literal original \texttt{Nums}]
\label{claim:DX4}
Every original good $D$ satisfies $D.\texttt{Nums}$.
More explicitly, for every descent $\bar r_S$ of the specified
$\bar r$ to $D.Gam$, every supplied normalization proof, and
every $j\in\mathbb Z^2$, the actual framed obstruction space
is finite dimensional and
\[
 \begin{aligned}
 \dim_k\operatorname{Ob}_F(\operatorname{ad}_j)
   &=\dim_kH^1_F(\operatorname{ad}_j)+
      \begin{cases}2&j=0,\\1&j=p\text{ or }j=-p,\\0&\text{otherwise};
      \end{cases}\\
 H^1_F(\operatorname{ad}_j)&=0
                         \quad(j\notin\{0,p,-p\}).
 \end{aligned}
 \tag{DX.20}
\]
The frame, all inertia conditions, coefficient field and complex
conjugation are the original ones.
\end{proposition}

\begin{proof}
Let $K=\ker\bar r$.  AX6 makes it open and normal.
Its finite quotient embeds in $(k^\times)^4$, so its order
and every inertia index $[I_q:I_q\cap K]$ are prime to $\ell$.
The frame is exactly $T=\{\ell\}\cup S_b$, $Q=S\setminus T$,
with $N_q=I_q\cap K$ and $v_0=\ell$.
The set $T\sqcup Q$ is the original $S$.
Thus all hypotheses in DX3 hold, including the exact original
local triviality at $N_q$.

For the weight-$j$ module in (DX.3), the action is scalar by
\[
       \gamma_j=\overline{\alpha^{j_1}\Psi^{j_2}}.
 \tag{DX.21}
\]
Indeed an $(u,v)$ entry transforms by
$\overline{\operatorname{ps}_u/\operatorname{ps}_v}$.
The displayed basis therefore proves both the action and
dimension $d_P(j)$, independently of a choice of descent:
the quotient map to $D.Gam$ is onto.

The twisted dual has no global invariant unless its scalar
character is cyclotomic.  For $j=0,\pm p$ this cannot happen:
the original $c=D.cc$ acts by $1$ on $\gamma_j$ and by $-1$
on the residual cyclotomic character.
For $j\in\{\pm a,\pm b,\pm t\}$, use the original
$\sigma_\ell$.  The cyclotomic value is $1$, whereas
AX6's residual separation makes $\gamma_j(\sigma_\ell)\ne1$.
For all other weights the module is zero.
Consequently the last term of (DX.12) is always zero.

At complex conjugation,
$\gamma_j(c)=(-1)^{j_1}$.
Thus its fixed-space dimension is two at zero, one at
$\pm p$, and zero at the six remaining roots.
Formula (DX.12) proves the equality and finiteness in (DX.20).

It remains to prove the asserted vanishing of the actual framed
cocycles.  At a weight outside (DX.3) the coefficient space
is zero.  At an odd root it is one dimensional and the
appropriate \texttt{Good.\allowbreak xvan} applies to
$\alpha,\beta,\theta$ or its inverse.
Take a framed pair $(z,(y_v))$ and pull $z$ back to the
original $G_{\mathbb Q}$.
It vanishes on inertia outside $S$, since it factors through
$G_{\mathbb Q,S}$, and it vanishes on $I_q\cap K$ at
the original chosen place for every $q\in Q$.
The same holds at every conjugate place:
\[
 z(hnh^{-1})=\gamma_j(h)z(n)
                           +(1-\gamma_j(n))z(h)=0
       \quad(n\in K).
 \tag{DX.22}
\]
This is exactly the original \texttt{Unr\allowbreak T} condition,
with no loss of valuations.
The framed primitives give the required coboundaries at $T$.
By \texttt{xvan}, $z(g)=(\gamma_j(g)-1)m$ globally.
At $v_0=\ell$ the original primitive is zero.
Evaluate at $\sigma_\ell$, where $\gamma_j\ne1$, to obtain
$m=0$.  Then $z=0$.
At every other $v\in T$, evaluation at the original
$\sigma_v$ and its residual separation gives $y_v=0$.
Hence the framed space itself is zero, not only its
cohomology quotient.  This verifies every quantifier and
field of $D.\texttt{Nums}$.
\end{proof}
% END DX4

\subsection{From integral universal representations to actual weighted charts}

% BEGIN DX5
\begin{proposition}[DX5: the integral ring and the characteristic-zero chart]
\label{claim:DX5}
Let $\Gamma$ be an original decomposition group from DX2.
Let $\mathcal P$ be either $\operatorname{eqL}$ at $\ell$ or
$\operatorname{eqP}$ at $q\in S_b$.
There is an integral universal ring $R\in\mathcal C_O$ for the
actual continuous lifts satisfying $\mathcal P$, with a weighted
presentation $W$, all variable weights in
$\{0,\pm a,\pm b,\pm p,\pm t\}$, and an origin
$x:R\to O$ giving exactly $r_0|_\Gamma$.
Every universal $(i,j)$ matrix entry has weight $w_i-w_j$.

Suppose the functor of continuous $\mathcal P$-lifts of
$r_0|_\Gamma$ over augmented finite local $E$-algebras is
unobstructed for small extensions, and its weight-$j$ tangent
space has dimension $n(j)$.
Then the original $W.\!A$, with its actual origin and grading,
satisfies
\[
                \operatorname{SmoothPt}(W.\mathrm{gr},x,n).
 \tag{DX.23}
\]
More precisely, the resulting map
$\phi:R\to E[[t_1,\ldots,t_d]]$ has the source's
$\operatorname{IsCompl}(\ker\phi,x,\phi)$ property, where
$d=\sum_j n(j)$; some $y\in R$ with $x(y)\ne0$
annihilates $\ker\phi$.
No integral smoothness assertion is needed.
\end{proposition}

\begin{proof}
By DX2, FP2 constructs the universal continuous matrix lift.
The prescribed conditions are actual equations in its entries.
For $\operatorname{eqP}$ they are
\[
 \rho(g)^tJ\rho(g)=\theta(g)J,\qquad
 \rho(n)=1\quad(n\in I_q\cap\ker\psi).
 \tag{DX.24}
\]
For $\operatorname{eqL}$ add the specified characteristic
polynomial of $\rho(s_0)$ and the source's equations
\[
 \begin{aligned}
 P_j(\rho(s_0))\,\rho(g)\,P_i(\rho(s_0))&=0
                                      &&(i<j),\\
 P_i(\rho(s_0))(\rho(s)-\operatorname{ps}_i(s))P_i(\rho(s_0))
                                      &=0 &&(s\in I_\ell),
 \end{aligned}
 \tag{DX.25}
\]
where
$P_i(X)=\prod_{h\ne i}(X-\operatorname{ps}_h(s_0))$.
These are the unnormalized products in \texttt{Ord};
no eigenvalue difference has been inverted over $O$.
Quotient the universal ring by the ideal generated by all these
entries.  That ideal is closed and finitely generated by
Noetherianity, even though the list of equations was infinite.
The quotient is proper, because the original diagonal $r_0$
satisfies every equation.  Universality of the quotient follows
by applying the equations entry by entry under an arbitrary
$O$-algebra map.  Thus it is the actual integral
$\operatorname{Represents}(\mathcal P,\rho)$, for every
$\mathcal C_O$ target.

The two-dimensional diagonal torus acts by
$\rho(g)_{ij}\mapsto u^{w_i-w_j}\rho(g)_{ij}$.
The fixed multiplier is preserved, since
$w_i+w_{3-i}=(1,1)$.
Conjugation preserves the characteristic polynomial and all
products in (DX.25), as well as (DX.24).
Take finitely many universal entry differences
$\rho(g)_{ij}-(r_0(g))_{ij}$ generating the relative
cotangent space.  They generate the complete ring over $O$
by FP1.  The formal-torus argument in FP3 makes the kernel
of this presentation homogeneous and places every actual
matrix entry in its stated weight part.  This applies in
arbitrary residue characteristic; independence of formal
characters, rather than evaluation on the finite group
$k^\times$, is the relevant argument.
All selected differences vanish under evaluation at $r_0$.
Consequently the presentation has the actual origin $x$.
Every matrix-entry weight belongs to the displayed nine-element
set, including weight zero.

Put $\mathfrak a=\ker x$ and choose $\varpi\in O$.
The characteristic-zero completion at this origin is
\[
 \mathcal R=\varprojlim_m(R/\mathfrak a^m)\otimes_OE
        =\widehat{R[1/\varpi]}_{\mathfrak a R[1/\varpi]}.
 \tag{DX.26}
\]
The quotient $R/\mathfrak a^m$ is a finite $O$-module:
in the chosen presentation it is generated by monomials of
degree less than $m$ in the variables.
Thus $\mathcal R$ is a complete Noetherian local
$E$-algebra with residue field $E$.
Equivalently it is the quotient of $E[[z_1,\ldots,z_s]]$
by the extended presentation ideal.
These assertions are the Noetherian completion theorem,
applied to the actual presentation; in particular, completion
does not change any quotient by a power of the augmentation
ideal.

We verify the represented functor, since replacing integral
deformations by arbitrary characteristic-zero representations
without this step would be insufficient.
Let $B=E\oplus\mathfrak n_B$ be a finite local augmented
$E$-algebra, with $\mathfrak n_B^m=0$.
An $E$-algebra map $\mathcal R\to B$ produces the continuous
representation obtained from the integral universal one.
Indeed the variables land in $\mathfrak n_B$; the image of
$\mathfrak m_R^N$ lies in $\varpi^{N-m+1}$ times a fixed
finite $O$-module for $N\ge m-1$.  Hence the induced map is
continuous from the adic topology of $R$ to the
$\ell$-adic topology of finite-dimensional $B$.

Conversely let $\rho_B$ be a continuous lift of $r_0$ to $B$.
The differences of its entries from those of $r_0$ lie in
$\mathfrak n_B$.  They have compact image, because
$\Gamma$ is compact, and hence lie in some finite
$O$-lattice $L\subset\mathfrak n_B$.
The $O$-subalgebra $B_0=O[L]$ is finite over $O$:
products of $m$ elements of $L$ vanish, so only the finitely
many monomials of degree less than $m$ are required.
It is local with residue field $k$ and complete; its maximal
ideal is generated by $\varpi$ and its nilpotent augmentation
ideal.  Its induced topology from $B$ equals its adic topology.
Thus $\rho_B$ is an actual lift over $B_0\in\mathcal C_O$.
The polynomial equations hold in $B_0$ because $B_0\hookrightarrow B$
is injective.  Integral universality gives $R\to B_0\to B$.
The augmentation is $x$, by uniqueness on the matrix entries;
therefore this map factors uniquely through $\mathcal R$.
This argument uses only bounded subsets of the given
finite-dimensional $E$-algebra.  It makes no assertion that
a compact subset of an algebraic closure lies in a finite
extension.

Choose homogeneous elements $f_1,\ldots,f_d\in\mathfrak a$
whose images form an $E$-basis of
$\mathfrak m_{\mathcal R}/\mathfrak m_{\mathcal R}^2$.
They may be chosen in the original $R$: its homogeneous
variable images span this vector space, and denominators
in finite $E$-linear combinations can be cleared in $O$.
The dimension of each weight part is $n(j)$.
Here the weight agrees with the source's tangent convention:
a functional supported on weight $j$ gives precisely
an entry deformation in $\mathfrak{sp}_4(E)_j$.
No change from $j$ to $-j$ is made.

Successive approximation gives a surjection
\[
       E[[t_1,\ldots,t_d]]\longrightarrow\mathcal R,
                         \qquad t_i\longmapsto f_i .
 \tag{DX.27}
\]
Its kernel $J_0$ has no constant or linear term.
We prove it is zero, rather than inferring smoothness from
the dimension count.  If $J_0\ne0$, let $n\ge2$ be the
least order of a nonzero member and choose one of order $n$.
Writing $P=E[[t_1,\ldots,t_d]]$, the map
\[
   P/\mathfrak m_P^{n+1}\longrightarrow
                         P/(J_0+\mathfrak m_P^{n+1})
 \tag{DX.28}
\]
is a small extension: its kernel is contained in degree $n$
and is killed by the maximal ideal.
Unobstructedness would lift the natural $\mathcal R$-point
of the right side to the left side.
Under this lift each $t_i$ changes only in degree $n$.
Substitution therefore leaves the nonzero degree-$n$
initial form of the chosen relation unchanged, a
contradiction.  Hence (DX.27) is an isomorphism.
The chosen elements are homogeneous, so its inverse respects
the formal torus action and every weight part.
Composing $R\to\mathcal R$ with this inverse gives $\phi$.

Its constant coefficient is $x$.  Each $f_i$ maps to $t_i$,
so it gives the required separate nonzero linear coefficient.
For the injectivity clause, if $x(f)=0$ and $\phi(f)$ has
zero linear terms, the completion isomorphism modulo the
square shows that $f/1\in\mathfrak a^2R[1/\varpi]$.
Thus $cf\in\mathfrak a^2$ for some $0\ne c\in O$.
We must prove this in the source's graded subalgebra $W.\!A$,
not just in $R$.
For homogeneous $f$ of weight $j$, express
$cf=\sum_{u,v}a_{uv}z_uz_v$, where the $z_u$ are the
homogeneous variable images.
Project each $a_{uv}$ to weight $j-w_u-w_v$.
Homogeneity of the presentation makes this projection
well-defined in $R$, and leaves the equality unchanged.
All the resulting factors belong to $W.\!A$, and every
$z_u$ vanishes at $x$.
Therefore $cf$ belongs to
$\ker(x|_{W.A})^2$.
An arbitrary element of $W.\!A$ is a finite sum of homogeneous
parts.  Its vanishing constant and linear coefficients hold
part by part; multiplying the finitely many denominators
proves the same assertion for that element.
This is exactly membership in $\operatorname{satSq}(x|_{W.A})$.
The grading, constant, separate-linear-coefficient and
injectivity clauses of \texttt{Is\allowbreak Chart} are now all verified.
Their precise parameter counts give (DX.23).

The stronger completion statement follows from the same
argument modulo $\mathfrak a^N$ for each $N$.
A truncated power series has finitely many $E$-coefficients;
clearing their denominators gives the required density.
If $\phi(f)$ lies in the $N$th power of the parameter ideal,
then $cf\in\mathfrak a^N$ for some $0\ne c\in O$;
this is stronger than the stated
$cf\in\mathfrak a^N+\ker\phi$.
Finally localize $R$ at the prime $\mathfrak a$.
The completion in (DX.26) is also its maximal-ideal completion.
Krull intersection makes $R_{\mathfrak a}$ inject into that
completion.  Hence each $f\in\ker\phi$ is killed by some
$s\notin\mathfrak a$.  Choose finitely many generators of
this kernel and multiply the associated $s$'s.
The resulting $y$ has $x(y)\ne0$ and annihilates the entire
kernel, exactly as required.
If a different original ring and lift representing the same
condition are specified, their two universal maps are inverse:
each composite fixes every representation entry and hence
is the identity by uniqueness.  Transporting the constructed
$x,\phi,\ker\phi,y$ through this unique isomorphism gives
the same completion and matrix-entry statements on that
specified ring.  Thus the chart conclusions hold for every
universal ring in \texttt{sp\allowbreak Chart\allowbreak L} and \texttt{sp\allowbreak Chart\allowbreak P},
not just for the presentation used in the construction.
\end{proof}
% END DX5

% BEGIN DX6
\begin{proposition}[DX6: the ordinary fourteen-parameter chart and \texttt{tan\allowbreak Lz}]
\label{claim:DX6}
For the original good datum, the integral universal
$\operatorname{eqL}$ ring has the chart in DX5 with counts
$n(j)=d_L(j)$.  Thus its characteristic-zero germ is a formal
power-series ring in fourteen parameters with the exact weights
in (DX.4).

The following two assertions are also literal consequences
on the original $O$-valued tangent sets:
\[
\begin{aligned}
 z\in\operatorname{tanSet}_\ell(\operatorname{eqL},-p)
 &\ \Longrightarrow\
 \exists\,0\ne a_0\in O,\ m_0\in O:\\
 &\qquad a_0z(g)=(\Psi(g)^{-1}m_0-m_0)X_-
                        \quad(g\in G_\ell),\\
 z\in\operatorname{tanSet}_\ell(\operatorname{eqL},0)
 &\ \Longrightarrow\ z(s)=0\quad(s\in I_\ell).
\end{aligned}
 \tag{DX.29}
\]
Here $X_-=E_{20}+E_{31}$.
This is the complete \texttt{tan\allowbreak Lz} statement, including its
single integral denominator and all original local elements.
\end{proposition}

\begin{proof}
Let $B=E\oplus\mathfrak n$ be a finite local augmented
$E$-algebra and let $\rho$ reduce to $r_0$.
Put $\lambda_i=\operatorname{ps}_i(s_0)\in O$.
AX6 makes these pairwise distinct in $E$.
Consequently their nonzero differences are units in $B$.
The characteristic-polynomial condition in (DX.25) yields
orthogonal idempotents
\[
 e_i=
 \frac{\prod_{h\ne i}(\rho(s_0)-\lambda_h)}
                 {\prod_{h\ne i}(\lambda_i-\lambda_h)}.
 \tag{DX.30}
\]
They sum to one and their images are free lines, since
their reductions are the four coordinate projections.
The stability equations say that
$F_j=\bigoplus_{i\le j}\operatorname{im}e_i$ is an invariant
complete flag.  The inertia equations say that $I_\ell$
acts on its $i$th successive quotient by exactly
$\operatorname{ps}_i$.

This flag is isotropic and self-dual for $J$.
Indeed preservation of the form at $s_0$ implies
$(\lambda_i\lambda_j-\theta(s_0))\langle u_i,u_j\rangle=0$
for vectors in its eigenlines.
For $i+j\ne3$ the scalar is nonzero, and hence a unit in $B$:
the four exponents $e+1,e,1,0$ have sum $e+1$ precisely
in the two complementary pairs, since $e\ge2$.
This proves the asserted orthogonality and flag duality.

The flag reduces to the standard one.
There is a unique element $u$ of the opposite unipotent
subgroup of $\operatorname{Sp}_4(B)$, reducing to $1$,
which takes the standard flag to it.
This is the ordinary Gaussian-elimination chart on isotropic
flags: choose the unit pivots in successive columns and use
the pairing to determine the complementary columns.
Equivalently, multiply the four lower-root subgroups from
(DX.3) in a fixed root order; solving the four lower
coordinates successively gives the unique element.
It has one parameter of each weight $-a,-b,-p,-t$.
These operations use only unit pivots, so apply over every
such $B$ and commute with its quotients.
The representation $u^{-1}\rho u$ is upper triangular.

Conversely, take any near-identity point $u$ of that
opposite-unipotent chart and any upper triangular
symplectic-similitude representation $\tau$ with multiplier
$\theta$ whose diagonal characters restrict to
$\operatorname{ps}_i$ on $I_\ell$.
Then $\rho=u\tau u^{-1}$ satisfies (DX.25).
For an upper triangular matrix with distinct diagonal
$\lambda_i$, its ordered spectral flag is the standard flag;
the diagonal action on each successive quotient is exactly
the required inertia character.  Applying (DX.30) therefore
recovers precisely the chosen flag and its unique $u$.
This proves an equivalence of the actual deformation
functors with a product of the four-parameter flag chart
and the upper triangular deformation functor.
It does not choose an arbitrary full eigenbasis or introduce
extra upper triangular changes of basis.

We prove the upper triangular functor unobstructed.
Its four diagonal characters have the form
$\operatorname{ps}_i\delta_i$, where $\delta_i$ is unramified
and congruent to $1$ modulo $\mathfrak n$, and
\[
                 \delta_0\delta_3=\delta_1\delta_2=1.
 \tag{DX.31}
\]
There are exactly two independent such characters.
On $G_\ell/I_\ell=\widehat{\mathbb Z}$ each is determined
by an arbitrary Frobenius value in $1+\mathfrak n$.
The logarithm and exponential are finite polynomials on
this nilpotent ideal; the associated additive homomorphism
factors through the natural projection
$\widehat{\mathbb Z}\to\mathbb Z_\ell$.
This proves existence, continuity, and the lifting of each
diagonal character across an arbitrary Artinian quotient.

For each positive root character
\[
                    \kappa\in\{\alpha,\beta,\Psi,\theta\}
 \tag{DX.32}
\]
we have $\kappa\ne1$ and $\kappa\ne\chi$ on $G_\ell$.
Nontriviality follows at $s_0$: their respective exponents
are $1,e-1,e,e+1$, all positive.
For $\alpha,\beta,\theta$, residual separation at the
actual $\sigma_\ell$ gives
$\kappa(\sigma_\ell)\ne1=\chi(\sigma_\ell)$.
For $\Psi$, comparison at $s_0$ gives $\Psi\ne\chi$
because $e\ge2$ and $\chi(s_0)$ has infinite order.
Thus DX1 gives
\[
 H^0(G_\ell,E(\kappa))=H^2(G_\ell,E(\kappa))=0,\qquad
 \dim_EH^1(G_\ell,E(\kappa))=1.
 \tag{DX.33}
\]
In particular the continuous cocycle space has dimension
two: its coboundaries form a one-dimensional copy of
$E(\kappa)$.

Consider a small extension $B\to B'$ with kernel $I$,
where $\mathfrak n_B I=0$.
Lift the two unramified diagonal characters first.
The upper unipotent group has four polynomial root
coordinates and is an affine four-space.
Choose continuous $E$-linear lifts of the coordinate
functions of a given representation over $B'$.
Together with the lifted diagonal they form a continuous
matrix-valued lift, preserving the multiplier and the
inertia diagonals, though initially it need not multiply
correctly.
Its multiplicative defect is a continuous two-cocycle with
values in
\[
 I\otimes_E\mathfrak n^+
   =\bigoplus_{\kappa\ {\rm in}\ (DX.32)}
                              I\otimes_EE(\kappa).
 \tag{DX.34}
\]
Associativity gives the cocycle equation.  Since $I$ is
killed by the maximal ideal, the conjugation action here
is exactly that of $r_0$, not a varying deformed action.
Equation (DX.33) makes the defect a continuous coboundary.
Correct the lifted matrices by the corresponding
$1$-cochain in $I\otimes\mathfrak n^+$; the square-zero
multiplication formula then removes the defect.
The correction is upper unipotent, so preserves the
multiplier, diagonals and inertia characters.
This proves lifting across every small extension.
Factoring the kernel of any Artinian surjection by its
finite maximal-ideal filtration proves lifting for every
surjection.

The tangent computation is now explicit.
The diagonal directions are the two unramified characters,
both weight zero.
The four positive-root directions are the cocycles in
(DX.33), two of each of the weights $a,b,p,t$.
There is no quotient by coboundaries: these are actual
matrix representations.  Finally flag conjugation supplies
one direction at each negative root.
The total is $2+4\cdot2+4=14$, and the weight multiplicities
are exactly $d_L$.
The entry action of the diagonal torus preserves the
flag chart, the Borel subgroup and all the constructions;
its tangent characters are precisely the entry weights
in (DX.3).  DX5 therefore supplies the claimed weighted
formal chart, with the original origin $r_0$.

To obtain (DX.29), start with an actual $O$-valued tangent
$z$ and its dual-number lift from \texttt{tan\allowbreak Set}, and
extend coefficients injectively to $E[\epsilon]$.
At weight $-p$ the preceding product decomposition has
only a flag direction.  Infinitesimal conjugation by
$1+\epsilon mX_-$ changes $r_0(g)$ by the cocycle
$(1-\Psi(g)^{-1})mX_-$.
Changing the sign of $m$ yields
$z(g)=(\Psi(g)^{-1}-1)m_E X_-$ for one $m_E\in E$,
simultaneously for all $g$.
Choose $0\ne a_0\in O$ with $m_0=a_0m_E\in O$.
Injectivity of $O\hookrightarrow E$ proves the first
identity on the original $O$-valued matrices.
At weight zero the only directions are the two unramified
diagonal characters, so the tangent vanishes on every
element of the actual inertia subgroup.
Injectivity again returns that equality to $O$.
These are exactly the two quantified conclusions in
\texttt{Sp\allowbreak WApply.\allowbreak tan\allowbreak Lz}.
\end{proof}
% END DX6

% BEGIN DX7
\begin{proposition}[DX7: the original finite-inertia ten-parameter chart]
\label{claim:DX7}
For each original $q\in S_b$, the integral universal
$\operatorname{eqP}(q)$ ring has the chart in DX5 with counts
$n(j)=d_P(j)$.  Its characteristic-zero germ has ten parameters,
two of weight zero and one of each root weight.
The statement imposes the source's entire inertia-kernel
condition on every deformation.
\end{proposition}

\begin{proof}
Set $N=I_q\cap\ker\psi$.  This is closed and normal in
$G_q$, since $\psi$ is a character and inertia is normal.
Every $\operatorname{eqP}$ deformation is trivial on $N$,
by its defining equations (DX.24), even when the coefficient
ring has nilpotents.  Thus it factors through the actual
profinite quotient
\[
  1\longrightarrow C=I_q/N
    \longrightarrow G_q/N
    \longrightarrow\widehat{\mathbb Z}\longrightarrow1 .
 \tag{DX.35}
\]
The map induced by $\psi$ identifies $C$ with the finite
cyclic subgroup $\psi(I_q)\subset O^\times$.
Its finiteness follows from the open kernel in
\texttt{Good}; a finite subgroup of a field's multiplicative
group is cyclic.
It is nontrivial, since $\ell$ divides its order by the
definition of $S_b$.
Conjugation on $C$ is trivial: $\psi$ is abelian and is
injective on this quotient.
Choose a lift $\gamma$ of $1\in\widehat{\mathbb Z}$.
The continuous power map from $\widehat{\mathbb Z}$ to the
procyclic closure of $\langle\gamma\rangle$ composes to the
identity on $\widehat{\mathbb Z}$.
It therefore splits (DX.35), giving
$G_q/N=C\times\widehat{\mathbb Z}$.
This quotient argument uses the original $I_q$ and $N$,
not merely their residual images.

On inertia, the original cyclotomic and $\omega$ characters
are trivial and $\xi$ is unramified.  Hence
\[
               r_0|_C=\operatorname{diag}(\psi,\psi,1,1).
 \tag{DX.36}
\]
Let $B=E\oplus\mathfrak n$ be an augmented finite local
$E$-algebra.  Every deformation of (DX.36) with the fixed
multiplier is conjugate by a symplectic matrix congruent
to $1$ to (DX.36) itself.
Here averaging is legitimate over $E$, including when
$\ell$ divides $|C|$; it is not used over the integral
residue field.
Explicitly, set
\[
                 u_0=|C|^{-1}\sum_{c\in C}
                                      \rho(c)r_0(c)^{-1}.
 \tag{DX.37}
\]
This reduces to $1$, is invertible, and satisfies
$\rho(c)u_0=u_0r_0(c)$.
Put $J'=u_0^tJu_0$ and $H=J^{-1}J'$.
The two forms $J,J'$ transform under $r_0(C)$ by the
same multiplier $\psi$, so $H$ commutes with $r_0(C)$.
It is self-adjoint for $J$ and congruent to $1$.
The finite binomial series $H^{-1/2}$ is therefore
defined in the nilpotent ideal and commutes with $r_0(C)$.
Then
$u=u_0H^{-1/2}$ still intertwines, and
$u^tJu=J$.
This proves the claimed symplectic rigidity of the
finite inertia representation, with explicit formulas
compatible with all coefficient quotients.

The centralizer $H_C$ of (DX.36) in $\operatorname{Sp}_4$
is $\operatorname{GL}_2$.
To check this over every $B$, choose $c\in C$ with
$\psi(c)\ne1$.  This difference is a unit in $E$ and
hence in $B$.  A commuting matrix is block diagonal on
the two eigenspaces $\langle e_0,e_1\rangle$ and
$\langle e_2,e_3\rangle$.
Preservation of $J$ makes its second block the paired
inverse transpose of the first, which is otherwise an
arbitrary invertible $2$-by-$2$ matrix.
Its Lie algebra has the two zero-weight directions
and the directions $a,-a$ in (DX.3).

There is a six-parameter formal slice for
$\operatorname{Sp}_4/H_C$ at the identity, of weights
$\pm p,\pm b,\pm t$.
One can construct it without a smooth-quotient theorem.
Near the identity the Cayley transform
\[
               Y\longmapsto(1+Y/2)(1-Y/2)^{-1}
 \tag{DX.38}
\]
identifies the formal Lie algebra with the formal
symplectic group; the inverse is
$g\mapsto2(g-1)(g+1)^{-1}$.
Take the six indicated root spaces as a torus-stable
linear complement to $\operatorname{Lie}H_C$.
The product of their Cayley slice with $H_C$ has derivative
the identity on the direct sum of these two Lie spaces.
Solving successively in each power of a nilpotent
augmentation ideal gives a unique formal factorization.
Thus every conjugating $u$ in (DX.37) has a unique
representative in this six-dimensional slice modulo
$H_C$.
The construction is torus equivariant, since (DX.38)
and the direct-sum decomposition are.

After this unique slice conjugation, inertia is (DX.36).
Because $\gamma$ commutes with $C$, its matrix must
centralize (DX.36).
It has the prescribed multiplier $\theta(\gamma)$
and reduces to $r_0(\gamma)$.
It is therefore uniquely an element of
$r_0(\gamma)\widehat H_C$, where $\widehat H_C$
is the formal group of $H_C$ at one.
Conversely every such matrix defines a continuous
representation of $\widehat{\mathbb Z}$.
To verify continuity, its entries and inverse entries
belong to a finite local $O$-subalgebra of $B$:
adjoin the finitely many nilpotent entry differences to
$O$, using $r_0(\gamma),r_0(\gamma)^{-1}\in\operatorname{GL}_4(O)$.
Products of sufficiently many nilpotent generators vanish,
so this subalgebra is finite over $O$.
Its matrix unit group is compact and profinite.
The integer power map of the chosen matrix therefore
extends to $\widehat{\mathbb Z}$.
Its multiplier agrees with the original $\theta$
on the dense integer subgroup, hence everywhere.
Combining this homomorphism with (DX.36) gives the
desired continuous representation of the original quotient
(DX.35), and then of $G_q$.

We have thus constructed an actual equivalence of functors
\[
 \{\operatorname{eqP}\text{-lifts near }r_0\}
   \ \simeq\
 \{\text{six slice parameters}\}\times\widehat H_C .
 \tag{DX.39}
\]
Both factors lift across all finite Artinian surjections.
The four centralizer parameters have weights $0,0,a,-a$;
the six slice parameters have the remaining root weights.
Their total multiplicities are exactly $d_P$.
The equivariance and exact functor identification now
apply DX5, giving the asserted chart and all its
completion properties on the original integral ring.
\end{proof}
% END DX7

% BEGIN DX8
\begin{theorem}[DX8: the complete original \texttt{Loc\allowbreak Base} installation]
\label{claim:DX8}
For every original good datum $D$, there exists an actual
$\operatorname{LocBase}D$ with all its original coefficient,
decomposition and inertia data.
Together with DX2 and DX4, this installs precisely
\[
 \Phi_\ell(D.Gam),\qquad D.\operatorname{Nums},\qquad
                  \operatorname{Nonempty}(\operatorname{LocBase}D).
 \tag{DX.40}
\]
The local base has the exact smoothness multiplicities
\[
                       d_L(j)+|S_b|d_P(j)
                           \qquad(j\in\mathbb Z^2).
 \tag{DX.41}
\]
Thus these three actual arithmetic inputs in the literal
$\operatorname{exModel}$ constructor require no remaining
bespoke supplier once the original good datum is supplied.
\end{theorem}

\begin{proof}
For $v=\ell$ use DX5--DX6's integral universal local ring
$R_\ell$ with its actual $\operatorname{eqL}$ representation.
For each original $q\in S_b$ use DX5--DX7's $R_q$
with its $\operatorname{eqP}(q)$ representation.
The prime $\ell$ does not belong to $S_b$ by definition.
All these rings are over the same original $O$, and their
origins give exactly the corresponding restriction of
the same $r_0$.

Take their completed coproduct over $O$.
Explicitly, write their finite weighted presentations
$R_v=O[[z_{v,i}]]/J_v$, use the disjoint union of all
variable sets, and define
\[
       B=O[[z_{v,i}:v\in T]]/
                     \Bigl(\sum_{v\in T}J_v\Bigr).
 \tag{DX.42}
\]
There are finitely many variables and factors.
The ideal in parentheses is proper, because evaluation
of all variables at zero gives a map to $O$.
It is closed and homogeneous.  Thus $B\in\mathcal C_O$,
and its union presentation $W$ has all variable weights
in the required set and has an origin.
For any target $A\in\mathcal C_O$, maps $B\to A$
are exactly tuples of maps $R_v\to A$:
the variable images lie in $\mathfrak m_A$, so all
power-series substitutions converge, and the defining
ideals give exactly the required relations.
The integral universal properties of the $R_v$ identify
these with the arbitrary original local representation
tuples satisfying \texttt{Loc\allowbreak Eqs}.
Both existence and uniqueness are for actual tuples,
without quotienting by conjugacy.
This proves the \texttt{univ} field of \texttt{Loc\allowbreak Base}.

The images of the universal local representations give
its \texttt{rho} and \texttt{eqs} fields.
Their entries remain in the actual weight parts of $W$,
giving \texttt{hom}.
Its origin evaluates them at $r_0$, giving \texttt{at0}.
These assertions hold on every element of every original
$G_v$; no isomorphic substitute field or representation
is chosen.

Complete (DX.42) at its characteristic-zero origin.
The resulting ring is the completed tensor product over
$E$ of the corresponding completions of all $R_v$.
This follows directly from the disjoint-variable
presentation, or by reducing modulo successive powers
of its augmentation ideal.
By DX6--DX7 it is therefore
\[
                 E[[t_1,\ldots,t_{14+10|S_b|}]],
 \tag{DX.43}
\]
with the disjoint union of their parameter weights.
The multiplicity of each $j$ is (DX.41).
The final graded-subalgebra argument of DX5 applies to
this union presentation verbatim: constants agree with
the origin, each parameter has a separately detected
linear coefficient, and a zero tangent is in the
$O$-saturation of the actual square of the origin ideal
inside $W.\!A$.
Consequently (DX.43) supplies exactly
\[
 \operatorname{SmoothPt}
      \bigl(W.\mathrm{gr},\,W.\mathrm{origin}|_{W.A},
                          j\mapsto d_L(j)+|S_b|d_P(j)\bigr).
\]
This is the remaining \texttt{smooth} field.
All fields of the original record have now been constructed.
DX2 and DX4 provide the other two statements in (DX.40).
\end{proof}
% END DX8

\subsection{Scope and named classical leaves}

The full \texttt{Cohomology\allowbreak Hyp} record packages more than is
needed by this installation.  DX1 proves its native local
finiteness, local duality and local Euler clauses, global
finiteness and global Euler clauses, and the projection of
LD3 to its normalized Poitou--Tate clause.
LK2--LK4 supply the corresponding finite native-$K$ clauses,
LD2 supplies its unramified orthogonality clause, and DX3
proves its full odd-prime Greenberg--Wiles clause.
The entire \texttt{Bloch\allowbreak Kato\allowbreak Local} field, including its
arbitrary local period-module models, is not asserted here.
Accordingly this section does not promote the full
\texttt{Local\allowbreak Published\allowbreak Theorems} or
\texttt{Cohomology\allowbreak Hyp} records.
Instead DX5--DX8 prove the local universal rings and
the exact required ordinary and finite-inertia charts
directly.  Equations (DX.40)--(DX.41) are unconditional
relative to the original \texttt{Good} datum and the
following named classical inputs.

\begin{itemize}
\item \textbf{DX-LT and DX-LE.}
Finite local Tate duality and the local Euler--Poincare
formula, with the Brauer invariant normalization:
J.\ S.\ Milne, \emph{Arithmetic Duality Theorems},
second edition (2006), Chapter I, Theorem 2.1,
Corollary 2.3 and Theorem 2.8, printed pages 26--31.
The source-to-classical match for arbitrary finite,
possibly nonfree coefficients is the one already made
in LK1--LK4; DX1 supplies the original decomposition-group
transport and the rational inverse-limit argument.
\url{https://jmilne.org/math/Books/ADTnot.pdf}.
\item \textbf{DX-GF and DX-GE.}
Global finiteness for finite modules whose order is
invertible on the ring of $S$-integers, and the global
Euler--Poincare formula:
Milne, the same book, Chapter I, Corollary 4.15,
printed page 62, and Theorem 5.1, printed page 67.
The real-place factor and the range of high-degree
vanishing are retained exactly in DX1.
The finite Poitou--Tate input is the already proved LD3,
with its explicit Chapter I, Theorem 4.10 specialization
and compatible invariants.
\item \textbf{DX-CM.}
The Noetherian completion theorem, including exactness
for finite modules and agreement modulo powers of the
ideal; and Krull intersection.
The Stacks Project, Algebra, Sections 10.96--10.97
and Lemma 10.51.4.
Only the completion/quotient assertions and the injectivity
of a Noetherian local ring into its completion are used
in DX5.  The deformation representation, its
unobstructedness and the weighted chart are proved above,
not retained as a separate deformation-theoretic leaf.
\url{https://stacks.math.columbia.edu/tag/00M9};
\url{https://stacks.math.columbia.edu/tag/0BNH};
\url{https://stacks.math.columbia.edu/tag/00IP}.
\end{itemize}

The original decomposition-to-completion equivalence and
its roots-of-unity compatibility are the same precisely
matched classical input already recorded in KB6 and LK1.
All original arithmetic character, inertia and residual
separation facts are AX6.  The incoming good datum itself
is constructed in the separate arithmetic-data branch.
No automorphy theorem, rigidity predicate, prime-counting
supplier or final RH/GRH conclusion is claimed in this
local installation.


% Version 0.27. Actual symplectic Model construction on the original datum.
% Pinned Lean source: 7cbf7a0c32fca1badf05931879a37c3f8a2b03b4.
\section{Constructing the actual symplectic deformation model}
\label{sec:symplectic-model-construction}

\status{MX1--MX10 reconstruct every field of the source's
\texttt{Model D}.  The intrinsic constructor takes the original
\texttt{Good D}, \texttt{Phi\allowbreak Cond D.\allowbreak Gam}, \texttt{D.\allowbreak Nums} and an
actual \texttt{Loc\allowbreak Base D}; FP8 supplies its universal
\texttt{Lemma\allowbreak P}.  Substitution of the exact arithmetic suppliers
DX2, DX4 and DX8 gives a model for every original good datum.
The construction does not supply image density, automorphy,
rigidity, a finite trace catalogue or the later geometric
comparison of points.}

\subsection{The objects that must be kept fixed}

Fix an algebraic closure $\Omega$ of $\mathbb Q$, an odd prime
$\ell$, the specified coefficient DVR $O=c.O$, its fraction field
$E$ and finite residue field $k$, and an original
$D:\texttt{Raw}\,\Omega\,c$ satisfying \texttt{Good}.
No field, coefficient ring, character or decomposition group is
changed in this chapter.  Put $e=D.e\ge2$.
Write $G=\operatorname{Gal}(\Omega/\mathbb Q)$,
$G_p=D.Gp(p)$, $I_p=D.Ip(p)$ and
\[
 \Psi=\chi_\ell^{e}\psi,\qquad
 \alpha=\chi_\ell\omega^{-1}\xi,\qquad
 \vartheta=\alpha\Psi,\qquad
 r_0=\operatorname{diag}(\vartheta,\Psi,\alpha,1),
 \qquad \bar r=r_0\bmod\mathfrak m_O.
 \tag{MX.1}
\]
Here $\chi_\ell$ is the integral cyclotomic character.  Set
$S_b=D.Sb$, $T=\{\ell\}\cup S_b$, and $b=|S_b|$.
The four vector weights and relevant matrix weights are
\[
 \begin{gathered}
 w_0=(1,1),\quad w_1=(0,1),\quad
 w_2=(1,0),\quad w_3=(0,0),\qquad w_{ij}=w_i-w_j,\\
 A=(1,0),\quad B=(-1,1),\quad P=(0,1),\quad H=(1,1),\\
 \mathcal W=\{0,\pm A,\pm B,\pm P,\pm H\}.
 \end{gathered}
 \tag{MX.2}
\]
These are exactly \texttt{w4}, \texttt{wt}, \texttt{w\allowbreak Al},
\texttt{w\allowbreak Be}, \texttt{w\allowbreak Ps}, \texttt{w\allowbreak Th} and \texttt{wts}.
The letter $P$ in a weight means $(0,1)$; a ring denoted
$\mathcal P$ below is different.

Use the alternating form
\[
 J=\begin{pmatrix}0&0&0&1\\0&0&1&0\\0&-1&0&0\\-1&0&0&0\end{pmatrix},
 \qquad X^*=J^{-1}X^tJ,
 \qquad \mathfrak {sp}_4(O)=\{Y:Y^*=-Y\}.
 \tag{MX.3}
\]
Thus \texttt{Is\allowbreak GSp}$(\nu,g)$ means $g^*g=\nu I$.
In particular the form convention is the original one, including
the signs of the two entries of
$M=E_{20}+E_{31}\in\mathfrak {sp}_4(O)$, of weight $-P$.

Write $\mathcal C_O$ for the complete Noetherian local
$O$-algebras with the specified residue field $k$.
It is exactly \texttt{Is\allowbreak CO}; algebra maps between these rings are
automatically local and continuous by FP1.  A lift has locally
constant entries modulo every power of the maximal ideal.
For an augmented ring $x:R\to O$, an $O$-valued point derivation
satisfies
\[
 \delta(fg)=x(f)\delta(g)+x(g)\delta(f).
 \tag{MX.4}
\]
It has weight $j$ if it vanishes on every homogeneous piece of
weight different from $j$.  This convention evaluates on the
weight-$j$ piece; it does not reverse the sign of $j$.

We use a \emph{weighted presentation} in the exact sense of FP:
a surjection $O[[z_s:s\in\Sigma]]\to R$ with finite $\Sigma$
and homogeneous kernel.  Its weight pieces $R_j$ are images of
homogeneous formal series, and $R^{\mathrm{gr}}=\sum_jR_j$ consists
of their finite sums.  This subalgebra need not be the whole
complete ring.  An augmented weighted presentation always uses
the source's constant-coefficient origin: $xW$ is evaluation
at zero, so every presentation variable has origin zero.
The augmentation of $R^{\mathrm{gr}}$ is the restriction of $x$.

The input \texttt{Loc\allowbreak Base D} is an actual ring
$B_0\in\mathcal C_O$, with a weighted presentation, origin
$x_B:B_0\to O$, and a universal family $(\rho_v)_{v\in T}$
of the source's local equations.  At $\ell$ these are the
fixed multiplier and \texttt{Ord} projector equations; at
$p\in S_b$ they are the fixed multiplier and identity on
$I_p\cap\ker\psi$.  Universality ranges over all rings of
$\mathcal C_O$, on the original $G_p$.
Its entries have weight $w_{ij}$, their origins are
$r_0(g)_{ij}$, and its prescribed smooth-point dimensions are
\[
 n_B(j)=d_\ell(j)+b\,d_p(j),
 \tag{MX.5}
\]
where $d_\ell$ is $2$ on $0,A,B,P,H$, $1$ on the other
four nonzero weights, and $0$ elsewhere; $d_p$ is $2$ at $0$,
$1$ at every nonzero weight in $\mathcal W$, and $0$ elsewhere.
The precise meaning of this smooth-point assertion is explained
and preserved in MX4--MX5.

\subsection{The actual framing problem}

% BEGIN MX1
\begin{proposition}[MX1: diagonal base point and original quotient descent]
\label{claim:MX1}
Let $\Gamma=D.Gam$ be the quotient of $G$ by the closed normal
subgroup generated by all inertia outside $D.S$.
The original $r_0,\bar r$ descend to $\Gamma$, with the form and
weight normalizations needed by FP8.  On $\Gamma$, imposing
identity on
\[
 N_q=I_q\cap\ker\bar r\ \le G_q,
                       \qquad q\in D.S\setminus T,
 \tag{MX.6}
\]
is equivalent, after pullback to $G$, to the literal
\texttt{D.\allowbreak Unr\allowbreak T} condition.  This is an equivalence of actual
continuous lifts and preserves the prescribed multiplier.
The original diagonal local family satisfies all
\texttt{Loc\allowbreak Eqs} equations, and these equations, multiplier,
normalization and unramified conditions commute with coefficient
maps in $\mathcal C_O$.
\end{proposition}
\begin{proof}
AX6 supplies the literal \texttt{facts} consequences of
\texttt{Good} on this same datum.  In particular the four
characters in (MX.1) are adically continuous, $\ker\bar r$
is open, all four characters are trivial on inertia outside
$D.S$, and
\[
 g\in I_A,\ A\text{ away from }T,\ \bar r(g)=1
           \quad\Longrightarrow\quad r_0(g)=1.
 \tag{MX.7}
\]
Here $A$ ranges over the original valuation subrings of
$\Omega$, exactly as in \texttt{Unr\allowbreak T}.  There is no replacement
by a convenient finite list of completions.
The kernel of any continuous lift into a complete ring is
closed: it is the intersection of the clopen conditions
$r(g)\equiv1\pmod {\mathfrak m^n}$.  Thus $r_0,\bar r$
factor through the indicated closed normal quotient.
The quotient is profinite.  For example its points are separated
by the finite quotients of $G$ containing the closed normal
subgroup; compactness then gives its inverse-limit description.
Factorization preserves adic continuity because the quotient
map is a quotient map on each finite matrix-valued reduction.

We check the stronger correspondence for arbitrary lifts.
If a lift on $G$ satisfies \texttt{Unr\allowbreak T}, then it kills inertia
outside $D.S$: such inertia is away from $T\subset D.S$,
and its residual image is one by AX6.  Its closed kernel
therefore contains the defining subgroup of $\Gamma$.
Conversely, a lift on $\Gamma$ satisfying (MX.6) kills the
required residual-trivial inertia at every original valuation.
Outside $D.S$ this follows from the quotient definition.
Over $q\in D.S\setminus T$, conjugacy of extensions of the
$q$-adic valuation identifies its inertia with a conjugate of
the specified $I_q$; both the representation kernel and
$\ker\bar r$ are normal.  At the trivial valuation the inertia
is trivial.  A valuation cannot lie over two different rational
primes, by B\'ezout's identity.  These cases are exhaustive,
and prove the asserted equivalence and its inverse.

The diagonal identities
$\vartheta\cdot1=\Psi\alpha$ give
$r_0(g)^*r_0(g)=\vartheta(g)I$.
The matrices are diagonal, so the block and equal-weight clauses
of \texttt{Normalises} hold over both $O$ and $k$.
For any lift $r$ of $\bar r$, the condition
$r(g)r_0(g)^{-1}\in\operatorname{Sp}_4$ with reduction one
is precisely the same fixed multiplier condition.  This proves
the form comparison used in descent.

The origin of the given local base sends its actual equations
to the diagonal local tuple.  Thus $r_0|_{G_v}$ satisfies
\texttt{Loc\allowbreak Eqs} without an additional existence assumption.
All its defining equations are polynomial: characteristic
polynomials and the projectors
$\prod_{j\ne i}(X-\psi_j(s_0)I)$, their products in
\texttt{Ord}, the multiplier equations, and prescribed
matrix equalities on inertia.  Polynomial evaluation commutes
with every algebra map.  So do the two \texttt{h\allowbreak El}
equations and the block equations in \texttt{Nrm}.
Every equality in \texttt{Unr\allowbreak T} is preserved as well.
In particular an \texttt{On\allowbreak Comp} witness maps to the stated
local equation; the domain witness is only needed when the
stronger component condition itself is asserted.
\end{proof}
% END MX1

% BEGIN MX2
\begin{proposition}[MX2: the entire kill-inertia local condition]
\label{claim:MX2}
For every $q\in D.S\setminus T$, the condition that a local
lift kill $N_q$ satisfies all five \texttt{Is\allowbreak Loc\allowbreak Cond} axioms
used in FP8, on the original group $G_q$.  The original
diagonal local representation satisfies this condition.
\end{proposition}
\begin{proof}
Set $H=G_q$, $I=I_q$ and $N=I\cap\ker\bar r$.
These are closed normal subgroups where appropriate.
The finite group $I/N$ embeds by $\bar r$ into the diagonal
group $(k^\times)^4$, so its order is prime to $\ell$.
The original decomposition group satisfies
$H/I\simeq\widehat{\mathbb Z}$, with arithmetic Frobenius
mapping to $1$.  This is the finite unramified-extension
theorem, transported from the completion as in KB6; it concerns
the specified $D.pl(q)$, not a new local place.\footnote{J. S. Milne,
\emph{Algebraic Number Theory}, version 3.08, Proposition 7.50,
Corollary 7.52 and Theorem 7.58(a), with the finite-residue-field
discussion following Remark 7.53;
\url{https://www.jmilne.org/math/CourseNotes/ANT.pdf}.
Taking compatible finite unramified extensions gives the stated
inverse-limit group.  KB6 supplies the comparison with the
original decomposition group.}

For every finite continuous $k[H/N]$-module $V$ we have
\[
                    H^2_{\mathrm{cts}}(H/N,V)=0.
 \tag{MX.8}
\]
Here is the needed specialization in detail.  Averaging over
the finite normal subgroup $I/N$ is possible because its order
is invertible in $k$, and makes its positive-degree cohomology
zero.  The Hochschild--Serre cochain filtration therefore
identifies the group in (MX.8) with
$H^2_{\mathrm{cts}}(\widehat{\mathbb Z},V^{I/N})$.
The latter vanishes by a direct finite calculation.  A
continuous cocycle and its action factor through a finite
cyclic quotient $C_m$.  A normalized extension representing
its class is determined by a lift $u$ of the generator and
the element $u^m$ in the invariant coefficient group.
Pulling back to $C_{\ell m}$ replaces this element by
$u^{\ell m}=\ell u^m=0$.  The resulting extension splits,
so the class dies in the direct limit.  Compactness ensures
that all the cochains here factor through finite quotients;
the finite-group cochain filtration and its averaging
contraction pass to that filtered limit.  This proves (MX.8)
for exactly the locally constant cochains that the source uses.
The general filtration theorem is the classical theorem of
Hochschild--Serre; only its finite-normal-subgroup specialization
is used here.\footnote{G. Hochschild and J.-P. Serre,
\emph{Cohomology of group extensions}, Trans. Amer. Math. Soc.
74 (1953), 110--134, Chapter II, Theorems 1--2
(the finite cochain filtration and its second page);
\url{https://www.sas.rochester.edu/mth/sites/doug-ravenel/otherpapers/hochschild-serre.pdf}.
The passage to continuous finite cochains and the cyclic
vanishing calculation are given above.}

Four local-condition axioms are immediate with their full
quantifiers.  Closedness is given by the ideal generated by
the entries of $r(n)-I$, $n\in N$.  This ideal has precisely
the required universal base-change property.  Conjugation
and diagonal torus conjugation preserve identity on $N$.
If one of two lifts is one on $N$, the other is one there
exactly when their restrictions to $N$ agree; this is the
literal tangent axiom, including its reverse implication.
The diagonal base tuple kills $N$ by (MX.7).

For the remaining lifting axiom take an arbitrary surjection
$p:A\to A'$ in $\mathcal C_O$ with kernel $J_0$ satisfying
$\mathfrak m_AJ_0=0$.  Thus $J_0^2=0$ and $J_0$ is a finite
$k$-vector space.  Let $r'$ be a prescribed local lift of
the fixed multiplier which kills $N$.
Apply the form-preserving matrix lift construction of FP4 to
$r'(g)r_0(g)^{-1}$, as a function on $H/N$.
It gives continuous matrices $\widetilde r(g)$ of the same
multiplier, reducing to $r'(g)$ and constant on $N$-cosets;
choose $\widetilde r(1)=I$.
For clarity, continuity requires no Artinian hypothesis on
$A$.  The finite kernel $J_0$ has zero intersection with
$\mathfrak m_A^n$ for some $n$; matrix lifts can be chosen
consistently on the corresponding clopen residue classes,
as in FP4.  The first-order correction of a form error uses
$1/2\in O$ and preserves the fixed multiplier.

The multiplicative error
\[
 \kappa(g,h)=\widetilde r(g)\widetilde r(h)
                      \widetilde r(gh)^{-1}-I
       \in\mathfrak {sp}_4(k)\otimes_kJ_0
 \tag{MX.9}
\]
is a locally constant two-cocycle for the residual adjoint
action.  Associativity gives its cocycle identity, because
all products of two kernel entries vanish.  That action is
trivial on $N$ and continuous.  Equation (MX.8) gives a
locally constant $a:H/N\to\mathfrak {sp}_4(k)\otimes J_0$
with $da=\kappa$.  Since $\kappa$ is normalized, $a(1)=0$.
Then
\[
                   r(g)=(I-a(g))\widetilde r(g)
 \tag{MX.10}
\]
is multiplicative: its error is $\kappa-da$.
It is continuous, has the same multiplier, reduces to $r'$,
and is identity on $N$.  The linear form condition on $a$
ensures preservation of the form to first order, and
$J_0^2=0$ makes that preservation exact.  This is the
required lift for every small surjection and proves all five
axioms on the actual local group.
\end{proof}
% END MX2

% BEGIN MX3
\begin{proposition}[MX3: the actual universal framed ring and its origin]
\label{claim:MX3}
Assume \texttt{Phi\allowbreak Cond D.\allowbreak Gam} and the literal \texttt{D.\allowbreak Nums},
and fix the actual local base $B_0$ above.
There is a full source record \texttt{Fr D L}.
In particular there are $U\in\mathcal C_O$, $i:B_0\to U$,
a continuous global lift $r_U$ and matrices $\alpha_v$
which are universal for
\[
 \begin{gathered}
 r_U(g)^*r_U(g)=\vartheta(g)I,\qquad
 \texttt{Unr\allowbreak T}(r_U,1),\qquad
 \alpha_v^*\alpha_v=I,\quad \bar\alpha_v=I,
 \quad\alpha_\ell=I,\\
 r_U(g)\alpha_v=\alpha_v i(\rho_v(g))
                                      \quad(g\in G_v).
 \end{gathered}
 \tag{MX.11}
\]
Let $a_j$ be the number of new relative variables of weight
$j$.  The relative presentation has exactly
\[
 a_j+d_j\quad\hbox{relations of weight }j,
 \qquad
 d_0=2,\quad d_P=d_{-P}=1,\quad d_j=0\text{ otherwise},
 \tag{MX.12}
\]
and $a_j=0$ for $j\notin\{0,P,-P\}$.
Its origin sends $r_U$ to the original $r_0$, all
$\alpha_v$ to $I$, and $i$ to $x_B$.
\end{proposition}
\begin{proof}
Apply FP8 with the group $\Gamma$ and the descended base
representations of MX1, the four weights in (MX.2), the form
(MX.3), the original local maps $G_v\to\Gamma$, distinguished
place $\ell$, and the kill-$N_q$ conditions of MX2.
The one block and its trivial involution give exactly the
source's \texttt{form O}; $w_i+w_{3-i}=(1,1)$ wherever
$J_{i,3-i}\ne0$.  MX1 checks both normalizers and all
\texttt{In\allowbreak H}/multiplier comparisons.  MX2 supplies every
local-condition axiom and its value at the base tuple.
The given local base supplies the original $B_0$, the
homogeneous entries, their local equations and their
fixed multiplier.  Its presentation variables have origin
zero, hence belong to $\mathfrak m_{B_0}$.
All hypotheses of FP8 have now been installed explicitly.

The obstruction finiteness required there is the first
clause of \texttt{Nums}, applied to the descended residual
representation and its specified normalizer.  FP8 makes the
number of new variables of weight $j$ equal to the
dimension of the source's unquotiented framed space
$H^1_F(\operatorname{ad}_j)$, and the number of relations
equal to $\dim_k\operatorname{Ob}_F(\operatorname{ad}_j)$.
The other two \texttt{Nums} clauses give exactly (MX.12)
and the stated vanishing of $a_j$.  These are equalities,
not estimates with a possibly larger obstruction space.
Each new variable is an actual entry of $r_U-r_0$ or
$\alpha_v-I$, with its original weight $w_{ij}$.

Universality of FP8 descends and pulls back along the
surjection $G\to\Gamma$ by MX1.  Equalities of pulled-back
representations are equivalent to equalities on $\Gamma$;
the local maps and all frame matrices are unchanged.
This is exactly universality in (MX.11) for every target
ring in $\mathcal C_O$, and not only for finite targets.

The tuple $(x_B,r_0,(I)_v)$ satisfies (MX.11), so universality
gives $x_U:U\to O$ with those values.  It kills every old
presentation variable because $x_B$ does, and every new one
because it is a centered entry of this tuple.  The composite
with the formal-series presentation is therefore constant
coefficient, by continuity and finite Taylor approximation.
Consequently its kernel has zero constant term and
$x_U$ is the presentation's actual origin.
All asserted origin identities follow from that one
universal map.  This proves every field of \texttt{Fr}.
\end{proof}
% END MX3

\subsection{Charts, independent cuts and a genuine domain}

% BEGIN MX4
\begin{proposition}[MX4: finite jets, relative smooth covers and derivation lifts]
\label{claim:MX4}
Let $x:R\to O$ be the constant-coefficient origin of a finite
weighted presentation, and put
$I=\ker x$.  Then, with $I_{\mathrm{gr}}=I\cap R^{\mathrm{gr}}$,
\[
 R^{\mathrm{gr}}/I_{\mathrm{gr}}^n\ \simeq\ R/I^n
                         \qquad(n\ge1).
 \tag{MX.13}
\]
Charts and smooth-point dimensions therefore pass between
the homogeneous subalgebra and the complete ring.
For the actual framed relative presentation of MX3 there
is a cover
\[
 B_0\ \longrightarrow\ \mathcal P=B_0[[u_t:t\in\tau]]
             \mathop{\longrightarrow}^{p}\ U
 \tag{MX.14}
\]
with homogeneous kernel generated by the original relative
relations $g_\nu$, exact origin and weight lifting, and
smooth-point dimensions $n_B(j)+a_j$.
Moreover every quotient-valued point derivation on $U$
whose restriction to $B_0$ lifts to an $O$-valued derivation
of weight $j$, and which is zero on new variables of other
weights, lifts to an $O$-valued point derivation of weight
$j$ on $\mathcal P$.
\end{proposition}
\begin{proof}
First clarify the chart definition.  Put
\[
       V_x=(I/I^2)\otimes_O E.
 \tag{MX.15}
\]
A source \texttt{Is\allowbreak Chart} is a homogeneous $O$-algebra map
$\phi:R\to E[[t_1,\ldots,t_n]]$ whose constant coefficient
is $x$ and whose linear-coefficient map identifies $V_x$
with $E^n$.  Indeed its \texttt{inj} clause says precisely
that the kernel on $I$ is
\[
 \operatorname{sat}(I^2)=\{f:\exists\,0\ne d\in O, df\in I^2\},
 \tag{MX.16}
\]
and its \texttt{onto} clause supplies each coordinate
direction with a nonzero coefficient.  Conversely an
isomorphism on $V_x$ supplies those directions after
clearing finitely many denominators.  The variables $t_i$
have the specified weights, and their counts are
\texttt{Smooth\allowbreak Pt}.  This is a rational cotangent chart;
we make no integral formal-smoothness assertion.

We prove (MX.13), including the injectivity that is needed
for saturation.  The images of the finitely many formal
variables generate $I$.  A series modulo their $n$th power
is its finite Taylor polynomial, which belongs to
$R^{\mathrm{gr}}$; this proves surjectivity.
For injectivity, express any $f\in I^n\cap R^{\mathrm{gr}}$
as a finite sum
$f=\sum_{|a|=n}z^a h_a$, $h_a\in R$.
Let $J_f$ be the finite set of weights of $f$.
Homogeneity of the presentation kernel makes each weight
component of an element of $R$ well defined.  Taking the
$J_f$-components of this equation replaces $h_a$ by the
finite sum of its components of weights
$j-\operatorname{wt}(z^a)$, $j\in J_f$.
The new coefficients belong to $R^{\mathrm{gr}}$ and give
the same $f$.  Each $z^a$ is a product of $n$ elements
of $I_{\mathrm{gr}}$, proving the contraction equality.

A chart on the homogeneous subalgebra sends
$I_{\mathrm{gr}}^n$ into $(t)^n$.  The isomorphisms (MX.13)
therefore extend it coefficient by coefficient to $R$:
to calculate a coefficient of degree $d$, choose an
approximation modulo $I^{d+1}$.  The answer is independent
of that choice, and addition and multiplication follow
from the corresponding finite-jet operations.
The extension preserves weights and constant coefficients;
the $n=2$ isomorphism proves its tangent injectivity and
surjectivity.  Restricting a chart to $R^{\mathrm{gr}}$
uses the same argument in reverse.  This proves the first
assertion without identifying $R$ with its graded subalgebra.

Write the framed presentation as
\[
 U=O[[z_s,u_t]]/
       (J_B,\,f_\nu),\qquad
 B_0=O[[z_s]]/J_B.
 \tag{MX.17}
\]
Here $J_B$ is extended by the actual old-variable inclusion,
and $f_\nu$ are the relative homogeneous equations from FP8.
Define $\mathcal P=O[[z_s,u_t]]/(J_B)$, with $g_\nu$
the images of $f_\nu$.  This is $B_0[[u_t]]$: coefficientwise
lifts give surjectivity, and finite generation of $J_B$
identifies its coefficientwise extension with its generated
ideal.  The finite power-series ring and its proper quotient
are complete Noetherian local with residue field $k$.
These standard completeness statements are also used in
FP1; the quotient assertion is Stacks, Lemma 10.160.2,
tag 0325.  The origin is evaluation at $u=0$ followed by
$x_B$.  Thus $p$ is surjective, its kernel is exactly
$(g_\nu)$, it preserves origins, and every homogeneous
element of $U$ has a homogeneous preimage in $\mathcal P$.

Extend the given chart of $B_0^{\mathrm{gr}}$ to $B_0$
using (MX.13), and send each new variable to a new formal
variable of the same weight.  This gives a chart of
$\mathcal P$ with dimensions $n_B(j)+a_j$.
The cotangent assertion follows either directly from
coefficients or from
\[
 V_{x_{\mathcal P}}=V_{x_B}\oplus
                              \bigoplus_{t\in\tau} E[u_t].
 \tag{MX.18}
\]

For the derivation assertion let
$D_U:U\to O/(a)$ and $D_B:B_0\to O$ be the given point
derivations, with $D_Ui=D_B\bmod a$.
For new variables of weight $j$ choose any integral lifts
$d_t$ of $D_U(pu_t)$; put $d_t=0$ for all other weights.
There is a unique point derivation $\delta$ on
$B_0[[u_t]]$ restricting to $D_B$ and sending $u_t$ to
$d_t$.  Explicitly, if $F=b_0+\sum_t b_tu_t+$ terms of
degree at least two in $u$, then
\[
       \delta(F)=D_B(b_0)+\sum_t x_B(b_t)d_t.
 \tag{MX.19}
\]
This proves the Leibniz identity even for arbitrary formal
series.  The assigned values and the weight of $D_B$
make $\delta$ homogeneous of weight $j$.
Modulo $(a)$ it equals $D_Up$: two point derivations on a
finite formal-series presentation agreeing on its variables
agree everywhere, since both kill the square of the
augmentation ideal.  This is every field of the cover's
derivation-lifting clause.
\end{proof}
% END MX4

% BEGIN MX5
\begin{proposition}[MX5: independent homogeneous cuts with tangent preservation]
\label{claim:MX5}
Let an augmented weighted $\mathcal P\in\mathcal C_O$
have a source chart, and let $t_1,\ldots,t_h$ be homogeneous
elements of origin zero.  Suppose each $t_i$ has an
$O$-valued point derivation nonzero on $t_i$ and zero on
the other $t_j$.  There is a surjection
\[
        q:\mathcal P\twoheadrightarrow R
 \tag{MX.20}
\]
onto a domain in $\mathcal C_O$, killing every $t_i$,
preserving origins and homogeneous lifting, and carrying
a chart with one variable of weight $\operatorname{wt}(t_i)$
removed for each $i$.  Every $O$-valued point derivation on
$\mathcal P$ that kills all $t_i$ annihilates $\ker q$.
The final chart also restricts to the homogeneous subalgebra
of $R$, exactly as required by \texttt{Model.\allowbreak smooth}.
\end{proposition}
\begin{proof}
Every point derivation kills $I_{\mathcal P}^2$ and,
because $O$ is a domain, its saturation.  Applying the
separating derivations to a relation between the classes
of the $t_i$ shows those classes are $O$-linearly independent
modulo the saturated square.  Clearing denominators makes
them independent in $V_{x_{\mathcal P}}$ as well.
Their chart images therefore have independent linear terms
over $E$, with the specified weights.

Within each weight block perform Gaussian elimination on
these linear terms.  Choose $h$ pivot variables and
invert the corresponding matrix over $E$; no inversion in
$O$ is asserted.  The resulting equations have the form
\[
       x_i+\sum_j c_{ij}y_j+H_i(x,y)=0,
                   \qquad H_i\in(x,y)^2,
 \tag{MX.21}
\]
homogeneously in each original weight.  They can be solved
uniquely in $E[[y]]$ by successive homogeneous degree:
the degree-one part is $-\sum_jc_{ij}y_j$, and the
degree-$n$ part is determined by the already known lower
parts of $-H_i$.  Equivalently iteration
$x^{(n+1)}=-Cy-H(x^{(n)},y)$ converges in the variable-adic
topology.  Every iteration preserves weights.
Substitution into the original chart gives
$\chi:\mathcal P\to E[[y]]$ killing the $t_i$.

First put $Q=\mathcal P/(t_i)$.  Its rational cotangent
is $V_{x_{\mathcal P}}$ modulo the span of the classes of
the $t_i$.  The linear map induced by (MX.21) identifies
that quotient with the remaining coordinate space.
Hence the induced $\chi_Q:Q\to E[[y]]$ is still a source
chart.  Now put
\[
                  R=Q/\ker\chi_Q.
 \tag{MX.22}
\]
This is a domain, since it embeds as a subring of $E[[y]]$.
The kernel lies in the augmentation ideal, by the constant
coefficient identity, so the origin descends and the
quotient is nontrivial.  Proper quotients of complete
Noetherian local rings are complete Noetherian local;
their residue field remains $k$ here.  Thus $R\in\mathcal C_O$.

The entire kernel of $O[[z_s]]\to E[[y]]$ is homogeneous.
Indeed the images of the $z_s$ have zero constant term, so
an output coefficient of degree $n$ sees only finitely
many input monomials of degree at most $n$.  On this finite
sum the weight decomposition commutes with the map.
If its whole image is zero, the image of each weight
component is therefore zero coefficient by coefficient.
Consequently $R$ has the induced weighted presentation,
with no new variable weights, and its weight pieces are
exactly the images of those of $\mathcal P$.

To check the last, essential assertion, let $\delta$ kill
all $t_i$.  The Leibniz rule and $x(t_i)=0$ show that it
kills their entire ideal, so it descends to $Q$.
If $f\in\ker\chi_Q$, the chart property says
$df\in I_Q^2$ for some $0\ne d\in O$.
Then $d\delta(f)=0$, and $\delta(f)=0$ because $O$ is
a domain.  Thus it descends further through (MX.22).
The new chart retains its tangent isomorphism: taking
this last quotient removes only classes already zero in
the rational cotangent.  Finally apply (MX.13) for $n=2$
to restrict it to $R^{\mathrm{gr}}$.
This proves the domain, all exact dimension losses, and
the integral derivation preservation simultaneously.
\end{proof}
% END MX5

\subsection{Normalization and the final rings}

% BEGIN MX6
\begin{proposition}[MX6: independent six-root cuts and exact normalization]
\label{claim:MX6}
In the actual framed ring $U$ define the six positions
\[
 \mathcal P_6=((0,1),(0,3),(2,1),(1,0),(3,0),(1,2)).
 \tag{MX.23}
\]
For $v=\ell$ take the entries in these positions of
$r_U(D.ph)$; for $v\in S_b$ take those of $\alpha_v$.
Denote the resulting $6(b+1)$ elements by $t_{v,h}$.
They have origin zero, have the six weights
$A,H,-B,-A,-H,B$, and admit the separating integral
derivations needed in MX5.
After killing them and, for each $v\in S_b$, the entries
$(0,2),(2,0),(0,0),(1,1)$ of $\alpha_v-I$, every
$\alpha_v$ is exactly $I$.
The six cuts at $\ell$ make $r(D.ph)$ block diagonal in
the source's $\lambda$-blocks.
\end{proposition}
\begin{proof}
We first construct actual tangent directions.
For any $Y\in\mathfrak {sp}_4(O)$ and
$A_v\in\mathfrak {sp}_4(O)$ with $A_\ell=0$, conjugate
the diagonal global representation over $O[\epsilon]$
by $I+\epsilon Y$.  Its derivative is $[Y,r_0(g)]$.
Take $\alpha'_v=I+\epsilon A_v$ and conjugate the
diagonal local tuple by $I+\epsilon(Y-A_v)$.
The local equations hold: the symplectic form is preserved
because $Y^*=-Y$, and characteristic polynomials,
polynomial projectors, their products and inertia identities
are invariant under conjugation.  The diagonal local
equations were proved in MX1.
The universality of $B_0$ gives a map to $O[\epsilon]$
for this local tuple.  Direct multiplication with
$\epsilon^2=0$ gives
\[
 r'(g)(I+\epsilon A_v)
  =(I+\epsilon A_v)\rho'_v(g).
 \tag{MX.24}
\]
The global multiplier and \texttt{Unr\allowbreak T} conditions also
hold: conjugating a matrix equal to $I$ leaves it equal
to $I$.  Thus MX3 gives a map $U\to O[\epsilon]$.
Its constant part is $x_U$, by the two universality
statements, so its $\epsilon$-part is a point derivation
with values
\[
 \delta(r_U(g))=[Y,r_0(g)],\qquad
                         \delta(\alpha_v)=A_v.
 \tag{MX.25}
\]

For the six positions in (MX.23), respectively, take the
symplectic root matrices
\[
 E_{01}-E_{23},\ E_{03},\ E_{21},\
 E_{10}-E_{32},\ E_{30},\ E_{12}.
 \tag{MX.26}
\]
Their six chosen entries form the standard basis.
For a cut at $v\ne\ell$, use $Y=0$, that root matrix
for $A_v$, and zero for the other $A_{v'}$.
For a cut at $\ell$, use that root matrix for $Y$ and
all $A_v=0$.
The derivative of the selected entry of $r_U(D.ph)$ is
the difference of its two diagonal characters at $D.ph$,
and the other five selected entries have derivative zero.
AX6 gives $\bar r(D.ph)=\bar r(D.sg(\ell))$ and the
cross-$\lambda$ residual eigenvalue separation from
\texttt{Good.\allowbreak cond\allowbreak D}.  The difference is therefore a unit
of $O$, in particular nonzero.  Formula (MX.25) gives
the asserted independent directions for all $t_{v,h}$.
Their origins and weights follow from the actual entries
and the diagonal base tuple.

We record the matrix calculation giving exact
normalization; it is valid in any local ring $A$.
Let $g^*g=\nu I$ with $\nu$ a unit, all off-diagonal
entries in $\mathfrak m_A$, and the six entries (MX.23)
zero.  The $(0,0)$ and $(1,1)$ multiplier equations show
that $g_{00}g_{33}$ and $g_{11}g_{22}$ are units: each
is $\nu$ plus a product containing an off-diagonal
entry.  Thus $g_{00},g_{11}$ are units.
The $(0,1)$ and $(1,0)$ equations then give
$g_{23}=g_{32}=0$.  These, together with the original
six zeros, are exactly all cross-$\lambda$ entries.
If also $\nu=1$, $g_{02}=g_{20}=0$ and
$g_{00}=g_{11}=1$, the same equations give
$g_{22}=g_{33}=1$; the $(1,3)$ and $(2,0)$ equations
then give $g_{13}=g_{31}=0$.  Hence $g=I$.

Apply this calculation to each frame matrix, whose
reduction is $I$, and to $r_U(D.ph)$, whose reduction
is diagonal.  It proves both asserted normalizations
after every subsequent quotient, with no appeal to a
generic characteristic-zero point.
\end{proof}
% END MX6

% BEGIN MX7
\begin{proposition}[MX7: all rings, equations, dimensions and local components]
\label{claim:MX7}
From the data of MX3 one can construct rings $R,S\in\mathcal C_O$,
with $R$ a domain, a weighted presentation $W$ on $R$,
an origin $x_R$, and a surjection $\pi:R\to S$, having
all the \texttt{Model} fields except those concerning
generation and cocycle derivations.
Their dimensions are explicitly
\[
 r=1+b+a_{-P},\qquad r'=2+b+a_P,\qquad
                         m=2+2b+a_0.
 \tag{MX.27}
\]
The kernel of $\pi$ is generated by exactly $r+1$
equations $F$ of weight $-P$, $r'$ equations $E$ of
weight $P$, and $m$ equations $Z$ of weight zero and
origin zero.
\end{proposition}
\begin{proof}
Use the cover $p:\mathcal P\to U$ of MX4 and choose
homogeneous lifts $\widetilde t_{v,h}$ of the cuts in
MX6.  Their origins are zero.  Pulling the separating
derivations back along $p$ gives the hypotheses of MX5.
Apply it to obtain $q:\mathcal P\twoheadrightarrow R$,
with all its domain, origin, weight and tangent properties.

In $U$ also form the following homogeneous elements:
the four entries of $\alpha_v-I$ at positions
$(0,2),(2,0),(0,0),(1,1)$ for each $v\in S_b$,
and the two entries $(0,2),(2,0)$ of
\[
                 h_U=r_U(g_0)+\vartheta(g_0)r_U(g_0^{-1}).
 \tag{MX.28}
\]
Call this list $s_\eta$.  All its origins are zero;
its weights are $P,-P,0,0$ per $v$ and $P,-P$ for
the last two.  Choose homogeneous lifts
$\widetilde s_\eta\in\mathcal P$ and set
\[
 \begin{gathered}
 J_S=(q(g_\nu),q(\widetilde s_\eta))\subset R,
 \qquad S=R/J_S,\qquad\pi:R\twoheadrightarrow S,\\
 \sigma:U\twoheadrightarrow S,
 \qquad \sigma p=\pi q,
 \qquad i_R=q|_{B_0}.
 \end{gathered}
 \tag{MX.29}
\]
The map $\sigma$ exists because $\ker p=(g_\nu)$;
it is surjective because $p,q,\pi$ are.
Every generator of $J_S$ has origin zero, so this ideal
is proper and $S\in\mathcal C_O$.
The same origin descends to $S$ from both sides of
the square.  In particular
\[
 \pi f=\sigma u\quad\Longrightarrow\quad
                            x_R(f)=x_U(u).
 \tag{MX.30}
\]
Every homogeneous $u\in U_j$ has a lift $f\in R_j$
with $\pi f=\sigma u$: lift first through $p$, then
apply $q$.  All presentation-variable weights in $R$
already occurred in $\mathcal P$, hence in $U$ or
$B_0$, and therefore belong to $\mathcal W$.

The smooth-point dimensions before the cuts are
$d_\ell(j)+bd_p(j)+a_j$.
Each of the six nonzero weights $\pm A,\pm B,\pm H$
loses $b+1$ dimensions.  At $A,B,H$ the result is
$2+b-(b+1)=1$, and at their negatives it is zero.
No cut has weight $0,P,-P$, giving respectively
$m,r',r$ in (MX.27).  The resulting dimension function
is exactly \texttt{t\allowbreak Dims r r' m}; in particular $r\ge1$.
MX5 already gives this chart on the homogeneous subalgebra
of $R$, with its restricted origin, as the record demands.

The relative equations in (MX.12) occur only at
$0,P,-P$.  The additional list $s_\eta$ contributes
$b+1$ equations at each of $P,-P$ and $2b$ at $0$.
Thus their total counts are
\[
 \begin{array}{c|c|c}
 \text{weight}&\text{number of equations}&\text{dimension}\\ \hline
 -P&a_{-P}+1+b+1=r+1&r\\
 P&a_P+1+b+1=r'&r'\\
 0&a_0+2+2b=m&m.
 \end{array}
 \tag{MX.31}
\]
Enumerate these complete lists, including any zero
equations, as $F,E,Z$.  They generate exactly $J_S$,
so the prescribed kernel equality holds.  Their weights
and the origin statement for $Z$ have already been proved.

Define $r_S=\sigma_*r_U$.  It is a continuous lift of
the original $\bar r$.  Multiplier and \texttt{Unr\allowbreak T}
pass through $\sigma$.  MX6 gives $\sigma(\alpha_v)=I$
and the block condition at $D.ph$; (MX.28) gives the
other two normalization equations.  Hence
\[
      (r_S)|_{G_v}=\pi_* (i_R)_*\rho_v.
 \tag{MX.32}
\]
For each local condition, use this very domain $R$ as
the \texttt{On\allowbreak Comp} witness.  Its local lift
$(i_R)_*\rho_v$ satisfies the original polynomial
equations, and $x_Ri_R=x_B$ maps it to $r_0|_{G_v}$.
The map to $S$ is exactly $\pi$, so (MX.32) proves
both component clauses of \texttt{D.\allowbreak Is\allowbreak Pt}.  There is
no assumption that a solution automatically lies on the
desired component.
Finally homogeneous entry lifting follows from the
lifting statement before (MX.30), and for every preimage
of any matrix entry (MX.30) gives its prescribed value
at the origin.  These are \texttt{hom} and \texttt{at0},
with their universal preimage quantifiers.
\end{proof}
% END MX7

% BEGIN MX8
\begin{proposition}[MX8: the actual representation generates the quotient]
\label{claim:MX8}
The representation $r_S$ of MX7 satisfies the literal
\texttt{Lift.\allowbreak Generates}: for every $A\in\mathcal C_O$,
two $O$-algebra maps $f,f':S\to A$ with
$f_*r_S=f'_*r_S$ are equal.
\end{proposition}
\begin{proof}
By (MX.32), the maps
$f\sigma i,f'\sigma i:B_0\to A$ give the same local
tuple, namely the restrictions of the equal global
representations.  The polynomial local equations are
preserved, so the uniqueness clause of the actual local
base gives $f\sigma i=f'\sigma i$.
Both maps $f\sigma,f'\sigma:U\to A$ now give the
same base map, global representation and frame tuple:
every frame matrix maps to $I$ by MX6.
The uniqueness clause of the actual framed universal
ring gives $f\sigma=f'\sigma$.  Surjectivity of
$\sigma$ proves $f=f'$.  This verifies the source's
test against all rings in $\mathcal C_O$; an informal
statement about a dense algebra of entries is unnecessary.
\end{proof}
% END MX8

\subsection{The integral Selmer-cocycle derivation}

% BEGIN MX9
\begin{proposition}[MX9: every original Sha cocycle gives the required derivation]
\label{claim:MX9}
For every nonzero $a\in O$ and every literal
$z\in\texttt{D.\allowbreak Sha\allowbreak Coc}(a)$, the rings and representation
of MX7 admit an $O$-valued point derivation $\delta$ of
weight $-P$ such that
\[
 \begin{gathered}
 \pi f=0\ \Longrightarrow\ \delta(f)\in(a),\\
 \pi f=r_S(g)_{20}\ \Longrightarrow\
       \delta(f)\bmod a=\vartheta(g)z(g),\\
 \pi f=r_S(g)_{31}\ \Longrightarrow\
       \delta(f)\bmod a=\Psi(g)z(g).
 \end{gathered}
 \tag{MX.33}
\]
All local coboundaries and the normalization $z(g_0)=0$
are those in the original source definition.
\end{proposition}
\begin{proof}
Put $Q=O/(a)$ and $\gamma=\Psi^{-1}$.
The hypotheses say that $z:G\to Q$ is locally constant,
$z(gh)=z(g)+\gamma(g)z(h)$, satisfies
\texttt{Unr\allowbreak T} with value zero, vanishes at $g_0$, and
on each original $G_v$ is a coboundary.
Choose an integral lift $m_v\in O$ of its specified
local parameter, so
\[
             z(g)=\gamma(g)m_v-m_v\pmod a
                                      \quad(g\in G_v).
 \tag{MX.34}
\]
No division by $\gamma(g)-1$ is used.

Work first in the square-zero ring $O\oplus Q$.
The matrix-valued function
\[
 r_z(g)=r_0(g)+\epsilon M r_0(g)z(g),
                    \qquad M=E_{20}+E_{31},
 \tag{MX.35}
\]
is a continuous representation of the same multiplier.
Indeed
$r_0(g)M r_0(g)^{-1}=\gamma(g)M$,
$M^*=-M$, and the cocycle identity proves multiplication
when $\epsilon^2=0$.  Adic continuity follows from that
of $r_0$ and local constancy of $z$.
For residual-trivial inertia away from $T$, both
$r_0(g)=I$ and $z(g)=0$ by MX1 and \texttt{Sha\allowbreak Coc}.
Thus (MX.35) satisfies the actual global unramified
condition as well.

For each $v$ conjugate the original diagonal local
representation over $O[\epsilon]$ by
$I-\epsilon m_vM$.  This satisfies the local equations
by MX6's conjugation argument and has derivative
\[
 [-m_vM,r_0(g)]
       =(\gamma(g)m_v-m_v)Mr_0(g).
 \tag{MX.36}
\]
The local universality gives a map $y_B:B_0\to O[\epsilon]$
for the whole tuple.  Its constant part is $x_B$ by
uniqueness.  Its derivative $D_B$ has weight $-P$.
We justify this last statement on every homogeneous
element, not only on representation entries.

The local entries determine maps from $B_0$ to
$O[\epsilon]$, by the same local uniqueness.
For a unit character $s:\mathbb Z^2\to O^\times$,
the homogeneous presentation defines a torus automorphism
of $B_0$, multiplying weight-$j$ elements by $s(j)$.
Precompose $y_B$ with it and multiply the $\epsilon$
coordinate by $s(-P)^{-1}$.
This map agrees with $y_B$ on every local entry:
the constant part is zero at nonzero weight, and the
derivative in (MX.36) is supported only at weight $-P$.
Thus the two maps agree on all of $B_0$.
For $x\in(B_0)_j$ this gives
\[
           (s(j)/s(-P)-1)D_B(x)=0.
 \tag{MX.37}
\]
If $j\ne-P$, choose $s$ using powers of the
infinite-order unit $1+\ell\in O^\times$ in a coordinate
where $j$ and $-P$ differ.  The parenthesis is nonzero;
the domain property of $O$ forces $D_B(x)=0$.
This establishes the precise \texttt{Has\allowbreak Wt} clause.

Reduce the second coordinate of $y_B$ modulo $a$.
Equations (MX.34)--(MX.36) show that its local tuple
is exactly the restriction of $r_z$.  Therefore
$(y_B\bmod a,r_z,(I)_v)$ satisfies (MX.11).
The actual framed universality gives
$y_U:U\to O\oplus Q$.
Its constant part is $x_U$, again by uniqueness, so
its second coordinate is a point derivation $D_U$ with
\[
 \begin{split}
 D_Ui&=D_B\bmod a,\\
 D_U(r_U(g)_{ij})&=M_{ij}\psi_j(g)z(g),\\
 D_U((\alpha_v-I)_{ij})&=0,
 \end{split}
 \tag{MX.38}
\]
where $(\psi_0,\psi_1,\psi_2,\psi_3)
=(\vartheta,\Psi,\alpha,1)$.
Every new presentation variable of weight other than
$-P$ has derivative zero: by MX3 it is a centered
representation or frame entry, and (MX.38) proves the
claim with its original weight.

Apply the integral lifting assertion of MX4 to obtain
an $O$-valued point derivation $\widetilde\delta$ on
$\mathcal P$ of weight $-P$ reducing to $D_Up$.
All cuts $\widetilde t_{v,h}$ have one of the six other
nonzero weights.  Thus $\widetilde\delta$ kills them,
and MX5 makes it descend to an $O$-valued point
derivation $\delta$ on the \emph{actual domain} $R$,
still of weight $-P$.

It remains to check the later quotient $S$, where an
integral descent is not claimed.  The old relations have
$D_U$-value zero.  The frame equations in $s_\eta$
have value zero by (MX.38).  The cocycle identity gives
$z(1)=0$, and then $z(g_0)=0$ gives $z(g_0^{-1})=0$.
Both \texttt{h\allowbreak El} equations in (MX.28) consequently
have $D_U$-value zero as well.
Thus $\delta\bmod a$ kills every generator of $J_S$.
All these generators have origin zero, so the Leibniz
rule makes it kill their entire ideal.

More generally, for every $f\in R,u\in U$ with
$\pi f=\sigma u$, choose lifts $f=q(h)$ and $u=p(h')$
in $\mathcal P$.  Their difference $q(h-h')$ belongs
to $J_S$.  The preceding vanishing and the reduction
$\widetilde\delta\bmod a=D_Up$ give
\[
                  \delta(f)\bmod a=D_U(u).
 \tag{MX.39}
\]
This proves the first clause of (MX.33) by taking $u=0$.
For the other two take the actual $(2,0)$ and $(3,1)$
entries of $r_U(g)$ and use (MX.38): $M_{20}=M_{31}=1$,
$\psi_0=\vartheta$, and $\psi_1=\Psi$.
The equalities hold for every preimage $f$, as required.
This finishes the entire \texttt{Model.\allowbreak der} field.
\end{proof}
% END MX9

% BEGIN MX10
\begin{theorem}[MX10: the literal actual Model constructor]
\label{claim:MX10}
On every original $\Omega,\ell,c,D$ the following exact
implication holds, with the universal \texttt{Lemma\allowbreak P}
already supplied by FP8:
\[
 \begin{gathered}
 \texttt{Good D}\ \wedge\ \texttt{Phi\allowbreak Cond D.\allowbreak Gam}\
 \wedge\ \texttt{D.\allowbreak Nums}\ \wedge\
                         \texttt{Nonempty(Loc\allowbreak Base D)}\\
             \Longrightarrow\quad\texttt{Nonempty(Model D)}.
 \end{gathered}
 \tag{MX.40}
\]
Consequently DX2, DX4 and DX8, in their exact scopes on
the same good datum, give
\[
       \texttt{Good D}\quad\Longrightarrow\quad
                              \texttt{Nonempty(Model D)}.
 \tag{MX.41}
\]
This conclusion includes every field of the record,
including its domain, actual local components, generation
and integral cocycle derivations.
\end{theorem}
\begin{proof}
Choose the supplied local base once.  MX1--MX3 give its
actual framed universal deformation ring with its origin,
homogeneous variables, relative equations and exact counts.
MX4 constructs the relative smooth cover.  MX5--MX7
construct the domain $R$, the quotient $S$, the presentation,
origin, integers $(r,r',m)$, all three equation lists,
and the actual representation $r_S$ with its \texttt{Is\allowbreak Pt}
witnesses.  MX8 supplies generation and MX9 supplies
the complete derivation assertion.
The following table identifies every record field.
\[
 \begin{array}{c|l}
 \text{source fields}&\text{constructed supplier}\\ \hline
 R,\ \texttt{Is\allowbreak CO},\ \texttt{Is\allowbreak Domain},\ W,\ \texttt{org}
       &\text{MX4--MX7}\!\\
 \texttt{wt\_\allowbreak mem},\ r,r',m,\ \texttt{one\_\allowbreak le},\ \texttt{smooth}
       &\text{MX7, (MX.27)}\\
 F,E,Z,\ \texttt{h\allowbreak F},\texttt{h\allowbreak E},\texttt{h\allowbreak Z}
       &\text{MX7, (MX.31)}\\
 S,\ \texttt{Is\allowbreak CO},\ \pi,\ \texttt{surj},\ \texttt{ker}
       &\text{MX7, (MX.29)}\\
 \texttt{ru},\ \texttt{pt},\ \texttt{hom},\ \texttt{at0}
       &\text{MX7, (MX.30)--(MX.32)}\\
 \texttt{gen}&\text{MX8}\\
 \texttt{der}&\text{MX9, (MX.33).}
 \end{array}
 \tag{MX.42}
\]
All instances and coefficient maps belong to the chosen
rings; every local component is witnessed by the same
constructed domain $R$.  This is exactly the source
\texttt{ex\allowbreak Model} conclusion and proves (MX.40).
For (MX.41), substitute DX2 for \texttt{Phi\allowbreak Cond}, DX4
for \texttt{Nums}, and the actual local base of DX8.
No weakening of those predicates, change of original
coefficient data, or converse to a conditional theorem
is used in this substitution.
\end{proof}
% END MX10

\paragraph{Scope of the completed construction.}
The new result supplies an actual \texttt{Model} on the
original good datum.  Its construction no longer leaves
\texttt{Fr}, \texttt{Cover}, a generic smooth quotient,
normalization, generation or a cocycle-derivation package
as an unexplained supplier.  The local equations enter
through the actual \texttt{Loc\allowbreak Base}; the arithmetic
cohomology enters through the precise \texttt{Nums} and
\texttt{Phi\allowbreak Cond} interfaces, whose installations are DX2,
DX4 and DX8.  The first-weight-zero geometry and the
representation-point/orbit/trace comparison are addressed
separately by PX.  None of these constructions by itself
asserts automorphy, image density, rigidity, the finite
trace catalogue, an unconditional count statement, or the
complete RH/Dirichlet-GRH reconstruction.


\section{The actual model, its points, and the symplectic trace inequality}
\label{sec:symplectic-model-comparison}

\status{This chapter reconstructs the passage from the actual
\texttt{Model} to the weighted intersection theorem, including the
point, tangent, trace and orbit comparisons.  The model is the one
constructed in MX, on the unchanged arithmetic datum of AX.  The
arithmetic assertions \texttt{Fin0}, \texttt{No\allowbreak End} and
\texttt{All\allowbreak Irr Rigid\allowbreak At} are stated explicitly below and are not
deduced from a dimension count.  Together with these assertions,
the already reconstructed rank-one theorems and the stated
classical local duality theorem supply the full comparison.
Neither an actual automorphic image theorem nor a finite
automorphic catalogue is inferred from this chapter.}

\subsection{The objects that must agree}

Fix an odd prime $\ell$, the original coefficient DVR
$O\subset\overline{\mathbb Q}_\ell$, its fraction field $E$,
and a datum $D$ satisfying \texttt{Good}.  Put
$K=\overline{\mathbb Q}_\ell$ and $G=G_{\mathbb Q}$ on the
specified algebraic closure.  All inertia and decomposition
groups retain the places chosen in $D$.  The four diagonal
characters and their two auxiliary weights are
\[
 (\theta_0,\theta_1,\theta_2,\theta_3)
       =(\alpha\Psi,\Psi,\alpha,1),\qquad
 (\lambda_i)=(1,0,1,0),\quad(\mu_i)=(1,1,0,0).
 \tag{PX.1}
\]
Thus $w_{ij}=(\lambda_i-\lambda_j,\mu_i-\mu_j)$, and
\[
 w_\alpha=(1,0),\quad w_\beta=(-1,1),\quad
 w_\Psi=(0,1),\quad w_\theta=(1,1).
 \tag{PX.2}
\]
The torus $(t_1,t_2)$ acts on an entry of weight $w=(a,b)$
by $t_1^at_2^b$, equivalently by conjugation with
$\operatorname{diag}(t_1t_2,t_2,t_1,1)$.

Let $M$ be an actual \texttt{Model} $D$.  Its weighted
presentation is $W:O[[z_1,\ldots,z_d]]\twoheadrightarrow R$,
with homogeneous kernel, origin $x_R:R\to O$ and
$\mathbb Z^2$-pieces $R_w$.  Write
\[
 B=\bigoplus_{w\in\mathbb Z^2}R_w\ \subset R,\qquad
 A=\bigoplus_{n\in\mathbb Z}R_{(0,n)}\ \subset B.
 \tag{PX.3}
\]
These are the finite-sum graded subalgebras; a completed ring
is not silently being declared a direct sum.  The model also
has a quotient $\pi:R\twoheadrightarrow S$ whose kernel is
generated by
\[
 F_0,\ldots,F_r\in R_{-w_\Psi},\qquad
 E_1,\ldots,E_{r'}\in R_{w_\Psi},\qquad
 Z_1,\ldots,Z_m\in R_0\cap\ker x_R,
 \quad r\ge1.
 \tag{PX.4}
\]
Its universal representation $\rho^u:G\to\operatorname{GSp}_4(S)$
is an actual \texttt{Is\allowbreak Pt} lift, is generated by its matrix
entries in the source sense, and has homogeneous entry lifts
of weight $w_{ij}$ and origin $\theta_i\delta_{ij}$.
Let
\[
 \mathcal S=\bigoplus_w\pi(R_w)\subset S,
 \qquad \varphi:A\longrightarrow\mathcal S,
 \qquad I=(F_0,\ldots,F_r,E_1,\ldots,E_{r'},Z_1,\ldots,Z_m)_A.
 \tag{PX.5}
\]
We denote the induced grading of $A$ by $A_n$ and its origin
by $x_0$.  A point $\xi:\mathcal S\to K$ gives
$\rho_\xi(g)_{ij}=\xi(\rho^u(g)_{ij})$.

Here are the three retained arithmetic assertions, with their
literal quantifiers.  \texttt{Fin0} says that for every complete
Noetherian local $O$-algebra with the prescribed residue field,
and every actual \texttt{Is\allowbreak Pt} lift over it, the $O$-algebra
generated by all entry monomials of total weight $(0,0)$ is a
finite $O$-module.  \texttt{No\allowbreak End} says that an actual
algebraic-valued point which is block diagonal for the partition
$\{0,2\}\sqcup\{1,3\}$ cannot have its $\{0,2\}$ block
span all $2$-by-$2$ matrices.  \texttt{All\allowbreak Irr Rigid\allowbreak At} says that
for every actual point $\rho$ whose matrix span is $M_4(K)$,
every actual dual-number point reducing to $\rho$ has derivative
$[X,\rho]$ for some $X\in\mathfrak{sp}_4(K)$.
These assertions concern all the points on this same datum,
not a selected finite family or a generic fibre alone.

We use the finite and rational local Tate duality and Euler
formulas already stated in LD/LK and DX, including their exact
transport to the chosen decomposition groups.  All remaining
commutative-algebra citation leaves are Noetherian formal
power series, Nakayama, exact completion on finite modules,
and the weak Nullstellensatz.  Precise references are Stacks
Project tags \href{https://stacks.math.columbia.edu/tag/0306}{0306},
\href{https://stacks.math.columbia.edu/tag/00DV}{00DV},
\href{https://stacks.math.columbia.edu/tag/00MA}{00MA}, and
\href{https://stacks.math.columbia.edu/tag/00FV}{00FV}.
The substitutions and invariant-theory arguments needed here
are proved below rather than left as additional suppliers.

\paragraph{PX1: actual gradings and generation in degrees at most two.}
\label{claim:PX1}
The sums in (PX.3) are direct; $A=\bigoplus_nA_n$ is a
graded $O$-algebra.  Its degree-zero ring is complete
Noetherian local with the original residue field.  It is
generated over $A_0$ by finitely many homogeneous elements
of degrees $-2,-1,1,2$, and $x_0(A_n)=0$ for $n\ne0$.
This is the literal \texttt{a\allowbreak Graded}/\texttt{a\allowbreak W} conclusion.

\begin{proof}
For a formal series let $p_w$ retain exactly the monomials
of weight $w$.  The kernel of $W$ is stable under every
$p_w$.  If a finite sum of elements of different weights
vanishes in $R$, lift the summands and apply each $p_w$;
each summand vanishes.  This proves both direct-sum assertions.
Products add weights because the coefficient of every monomial
in a formal product is a finite sum.

We record a useful finite-monomial argument.  For any fixed
weight $w$, the set of exponent vectors in $\mathbb N^d$ of
weight $w$ has finitely many coordinatewise minimal elements.
Here the elementary finiteness statement follows by induction
on $d$: an infinite sequence has an infinite subsequence with
nondecreasing first coordinate, and repeat on the remaining
coordinates; an antichain is therefore finite.  Every exponent
of weight $w$ dominates one of these minimal exponents, with
difference of weight zero.  Assign each monomial to the first
such minimal exponent.  Consequently every weight-$w$ series
is a finite sum of these monomials times weight-zero series.
Similarly, the nonzero minimal exponent vectors of weight zero
generate that monoid: subtract one and induct on total degree.
There are only finitely many such generators.  Thus the full
weight-zero ring is a quotient of a formal power series ring
over $O$ on these nonconstant monomials.  Substitution is
well-defined because their degrees are positive.  A series
with unit constant coefficient has an inverse of weight zero,
obtained either by a geometric series or by projecting an
inverse.  This proves locality, completeness, Noetherianness
and the original residue field.  The same argument applies
after a homogeneous quotient.

Each original variable has first weight in $\{-1,0,1\}$
and second weight between $-1$ and $1$.  A monomial of first
weight zero is a product of variables of first weight zero
and pairs consisting of one variable of first weight $1$
and one of first weight $-1$.  Remove one such factor and
induct on degree.  Apply the preceding finite-monomial
decomposition to each full weight $(0,n)$.  It shows that
$A$ is generated over $A_0$ by the finitely many corresponding
single variables and pairs.  Their second weights lie between
$-2$ and $2$; those of weight zero are absorbed into $A_0$.
Finally a nonzero-weight series has constant coefficient zero,
which is exactly the origin assertion.  Notice that none of
these steps replaces a formal series by an unjustified
infinite sum in the algebraic direct sum.
\end{proof}

\paragraph{PX2: the invariant chart and its exact dimensions.}
\label{claim:PX2}
The source \texttt{Smooth\allowbreak Pt} condition on $B$ induces
$\texttt{Smooth\allowbreak Pt}(A_n,x_0,\mathtt{mDims}(r,0,r',0,m))$.
Explicitly the saturated cotangent ranks of $A$ are
\[
 \begin{array}{c|rrrrr}
 n&-2&-1&0&1&2\\\hline
 \operatorname{rank}_O C_n(A,x_0)&0&r&m&r'+1&1.
 \end{array}
 \tag{PX.6}
\]
The chart and the saturation in this statement are over $E$;
no integral formal smoothness is asserted.

\begin{proof}
Write $J=\ker(B\to O)$ and $J_A=\ker(A\to O)$.  The source
chart $\Phi:B\to E[[Y_i:i\in T]]$ respects both weights
and the origin.  Its linear coefficients identify
\[
 (J/J^2)\otimes_O E\ \simeq\ E^T.
 \tag{PX.7}
\]
Indeed its \texttt{inj} field says precisely that a zero
linear part lies in the $O$-saturation of $J^2$, and its
\texttt{onto} field supplies a diagonal basis of nonzero
linear parts.  Projection to the appropriate weight and
subtraction of the constant term make these basis lifts
homogeneous.  This also proves directly that (PX.7) respects
weights.

For clarity we justify the use of formal coordinates here.
The finite-monomial proof of PX1 makes $B$ Noetherian: it is
generated by the finitely many $z_i$ over the complete
Noetherian full weight-zero ring.  Complete $B\otimes_O E$
at the augmentation ideal.  Homogeneous lifts of the basis
in (PX.7) give a surjection from $E[[Y_i]]$ to this
completion, by successive approximation and Nakayama.
Composing with $\Phi$ gives a formal substitution with
invertible linear matrix.  Such a substitution is invertible:
after fixing the inverse through degree $q$, solve its
degree-$q+1$ error with the inverse linear matrix.  The
successive corrections converge in the variable-adic
topology.  It follows that the surjection and the completed
chart are isomorphisms.  The construction is equivariant
for both weights, since every correction may be projected
to its prescribed weight.

Passing to first-weight zero commutes with these completions
in this situation.  Here is the needed topology check,
which prevents an unwarranted use of formal invariants.
A first-weight-zero monomial in the original variables
splits into singles and opposite-weight pairs.  Any such
monomial of total degree at least $2q$ has at least $q$
nonconstant blocks, and lies in $J_A^q$.  Conversely
$J_A^q\subset J^q$.  On a homogeneous component an element
of $J^{2q}$ has a formal representative of order at least
$2q$: lift products of $2q$ augmentation elements, then
project to that weight.  Grouping its monomials by the
finitely many possible first $q$ blocks proves the same
inclusion for series.  Taking the finite sum of components
gives
\[
        J^{2q}\cap A\subset J_A^q\subset J^q\cap A.
 \tag{PX.8}
\]
The argument survives localization in $O\setminus\{0\}$.
Weight-zero projection is an additive retraction compatible
with each of these filtrations.  Thus compatible invariant
truncations lift compatibly, and (PX.8) identifies the
completion of $A\otimes E$ with the invariant part of the
completed $B\otimes E$.  This is also an explicit proof of
the required exactness of taking this torus's invariants.

The source dimension table has one variable $Y_\alpha$
of weight $(1,0)$, one $Y_\beta$ of weight $(-1,1)$,
one $Y_\theta$ of weight $(1,1)$, $r'$ variables of
weight $(0,1)$, $r$ of weight $(0,-1)$, and $m$ of
weight zero.  There is exactly one negative first-weight
variable, namely $Y_\beta$.  A monomial of first weight
zero is uniquely a monomial in
\[
 Y_i\quad(w_i^{(1)}=0),\qquad
 U_\alpha=Y_\alpha Y_\beta,\qquad
 U_\theta=Y_\theta Y_\beta.
 \tag{PX.9}
\]
The exponent of $Y_\beta$ is the sum of the exponents of
the two positive-first-weight variables.  This gives an
isomorphism of formal rings, coefficient by coefficient;
there are no relations among the displayed generators.
Their second weights are respectively unchanged, $1$, and
$2$.  Restrict the source chart and relabel these monomials.
The completed-ring identification proves its linear map is
an isomorphism on the rational cotangent space.  Clearing
the finitely many denominators supplies the source
\texttt{onto} witnesses in $A$, and zero linear part is
equivalent to membership in the saturation of $J_A^2$.
This proves every field of \texttt{Is\allowbreak Chart} and (PX.6).
\end{proof}

\paragraph{PX3: global tangent vanishing and all three spanning conditions.}
\label{claim:PX3}
Using RV7 and the ordinary tangent consequences in DX, the
equations (PX.4) satisfy
\[
 \mathtt{SpansCot}(-1,F),\qquad
 \mathtt{SpansCot}(-2,\varnothing),\qquad
 \mathtt{SpansCot}(0,Z).
 \tag{PX.10}
\]
Together with PX2 and $r\ge1$, this is the literal condition
$\mathtt{M1}$ for $M.\mathtt{eqns}$.

\begin{proof}
We first prove the actual global tangent assertion used in
the spanning test.  Suppose $\rho':G\to M_4(O[\epsilon])$
is an actual point, reduces to $\rho_0$, and its derivative
has only weight $w=0$ or $w=-w_\Psi$.  Set
\[
 z(g)_{ij}=\epsilon^{-1}(\rho'(g)-\rho_0(g))_{ij}
                                  \theta_j(g)^{-1}.
 \tag{PX.11}
\]
Expanding the representation identity shows that each entry
is a cocycle for $\theta_i\theta_j^{-1}$; the fixed
similitude identity places $z$ in $\mathfrak{sp}_4$.
The entries are continuous.  One direct check is that a
map from a complete local presentation to $O[\epsilon]$
sends the first component of its maximal ideal into
$\mathfrak m_O$, and the second component of its
$(q+1)$st power into $\mathfrak m_O^q$.  Thus the
given adic continuity implies continuity of the derivative.

For $w=0$ the entries are additive diagonal characters.
DX's weight-zero ordinary tangent calculation makes them
zero on inertia at $\ell$.  At $p\in S_b$ they vanish on
$I_p\cap\ker\psi$, and away from $T$ on
$I_p\cap\ker\bar\rho$.  The omitted subgroups have finite
index.  An additive characteristic-zero cocycle zero on a
finite-index subgroup is zero on all inertia, by taking a
power.  Transport by conjugation handles every place.
A continuous additive character of $G_{\mathbb Q}$
unramified everywhere is zero: each finite quotient of its
compact image cuts out an everywhere-unramified abelian
extension of $\mathbb Q$, which is trivial by the class
number of $\mathbb Q$ (equivalently Minkowski and class
field theory).  Passing through all finite quotients gives
the assertion for the original continuous character.

For $w=-w_\Psi$, only the entries $(2,0)$ and $(3,1)$
can occur; the symplectic identity relates them.  The first
entry is a cocycle for $\Psi^{-1}=\psi^{-1}\chi_{\rm cyc}^{-e}$.
DX's negative ordinary tangent calculation makes it a
coboundary on $G_\ell$ after passing to $E$.  At all other
inertia groups the same finite-index argument makes it a
coboundary.  Explicitly, if a cocycle vanishes on a
finite-index subgroup of inertia, it vanishes on the
kernel of its character by taking powers; choosing one
element with character different from $1$ then writes it
as $(\gamma-1)c$.  If the character is trivial it is
already zero.  RV7 applies to precisely this cocycle,
including its local condition at $\ell$, and gives
$z(g)_{20}=(\Psi(g)^{-1}-1)c$ globally.  The derivative
of the normalization $h(\rho')_{20}=0$ at $g_0$ gives
\[
             c(\alpha(g_0)-1)(\Psi(g_0)-1)=0.
 \tag{PX.12}
\]
Both factors are units by \texttt{Good}.  Hence $c=0$,
and the symplectic relation kills the other entry.  This
proves the full \texttt{tan\allowbreak Glob} assertion in both weights.

Now work first on $B$.  Its cotangent lattice of any
fixed weight is finite: a homogeneous series modulo the
square of the augmentation is a combination of the
finitely many original variables of that weight.  Suppose
the $F_i$ failed to span the rational cotangent space at
$-w_\Psi$.  A nonzero $E$-linear functional annihilating
them can be multiplied by one denominator to take values
in $O$ on this finite lattice.  It defines an $O$-valued
point derivation of the specified weight.  The same
argument applies to the $Z_i$ at weight zero.
Such a derivation on $B$ extends to $R$: on the formal
presentation use the finite linear part at the origin;
homogeneous kernel relations are killed, and the
quadratic terms give the Leibniz rule.  Its weight forces
it to kill the other two groups of equations in (PX.4).
Consequently $x_R+\epsilon\delta$ factors through $S$.
The universal representation gives an actual tangent
point of the kind just proved zero.  The model's
entry-generation property then makes the two maps
$S\to O[\epsilon]$ equal, so $\delta=0$, a contradiction.
Thus $F$ and $Z$ span on $B$ after rationalization.

Inclusion $A\subset B$ induces surjections of cotangent
spaces at weights $-1\mapsto-w_\Psi$ and $0\mapsto0$:
the homogeneous pieces themselves agree.  PX2 gives
their ranks at most $r,m$, while the original model
chart gives the target ranks exactly $r,m$.
They are therefore rational isomorphisms, transferring
the just-proved spanning statements to $A$.  The
degree-$-2$ cotangent has rational rank zero by PX2;
its defining saturation makes it torsion free, so it
vanishes and is spanned by the empty family.  Finally,
spanning after rationalization is exactly the displayed
\texttt{Spans\allowbreak Cot} predicate: for each integral class
clear its denominator, without asserting integral
generation or changing the integral cotangent quotient.
\end{proof}

\paragraph{PX4: the exact equation ideal and lifting every algebraic point.}
\label{claim:PX4}
The map $\varphi$ has kernel $I$, is onto each
first-weight-zero piece of $\mathcal S$, and preserves
its second weight.  Every $x:A\to K$ with $x(I)=0$
extends to $\xi:\mathcal S\to K$ with $x=\xi\varphi$.
These are the literal \texttt{ker\allowbreak A} and \texttt{pt\allowbreak Lift}
statements; no irreducibility hypothesis is used.

\begin{proof}
The ideal in (PX.4) is homogeneous, so $S$ has the
induced weighted presentation and its finite homogeneous
sum $\mathcal S$ has a genuine bigrading by the projection
argument of PX1.  If $f\in A$ lies in $\ker\pi$, write
it as a finite linear combination of the generators in
(PX.4) with coefficients in $R$.  For each of the
finitely many full weights of $f$, project this identity
to that weight.  Since every generator has first weight
zero, only the first-weight-zero components of its
coefficient survive, and these lie in $A$.  The resulting
finite identity puts $f$ in $I$.  The reverse inclusion
is immediate.  A homogeneous lift of a class in
$\pi(R_{(0,n)})$ proves the other two assertions about
$\varphi$.

Let $C=\varphi(A)=\mathcal S^{\mathbb G_m^{(1)}}$.
The point $x$ descends to $C$.  Projection
$P:\mathcal S\to C$ to first weight zero is $C$-linear
and has $P(1)=1$.  Hence the fibre
$\mathcal S\otimes_{C,x}K$ is nonzero: $xP$ induces
a $K$-linear map from this fibre to $K$ sending $1$ to
$1$.  It is of finite type over $K$.  Indeed the
finite-monomial decomposition in PX1 says that the full
graded algebra is generated over its full weight-zero
piece by the finitely many original variables; that
piece is contained in $C$.  Choose a maximal ideal of
the nonzero fibre.  The weak Nullstellensatz makes its
residue field $K$, yielding the desired extension.
Thus surjectivity on points is a proved Reynolds
projection and finite-type argument, not an assumption
that every point of a subalgebra extends automatically.
\end{proof}

\paragraph{PX5: algebraic entry generation, actual points and condition M2.}
\label{claim:PX5}
Assume \texttt{Fin0}.  Each homogeneous piece of
$\mathcal S$ is the $O$-linear span of the entry
monomials of that full weight.  Two algebra maps out
of $\mathcal S$ into any $O$-algebra agree if their
evaluated representations agree.  Every point
$\xi:\mathcal S\to K$ can be rescaled by the full
two-dimensional torus to the restriction of a map
$S\to K$, so its representation becomes an actual
$D.\mathtt{EPt}$.  Finally the degree-zero image of
$A/I$ is finite over $O$, the literal condition
$\mathtt{M2}$.

\begin{proof}
Let $C_0$ be the $O$-algebra generated by the entry
monomials of weight zero.  It is finite by
\texttt{Fin0}, applied to the actual universal lift.
Choose finitely many homogeneous entry differences
(subtract a scalar from diagonal entries) whose
relative tangent classes generate that of $S$.
They exist by the model's generation assertion: a
linear functional killing all these classes would
give two different maps into the residue-field dual
numbers with identical representations.  Nakayama
and successive approximation show that these finitely
many differences generate $S$ topologically.
The resulting new presentation has homogeneous kernel
for the original grading.  To see this directly, lift
each chosen difference to a homogeneous series in the
original $O[[z_s]]$.  Its constant coefficient can be
taken zero: the origin descends to $S$, since every
equation in (PX.4) has origin zero, and the differences
were centered at that origin.  Substitute these lifts
into a series in the new variables.  Their positive
formal order makes every output coefficient a finite
sum, so substitution commutes with each weight
projection.  The original presentation kernel is
homogeneous, hence so is the new kernel.  The same
identity shows that every old homogeneous element
has a new homogeneous preimage.  No assertion that
an arbitrary image of a weight piece is closed is
needed here.

The finite ring $C_0\subset S$ is complete over $O$.
It is local: the usual finite-ring decomposition over
a complete local ring would otherwise give a nontrivial
idempotent inside the local ring $S$.  Its maximal
ideal is $C_0\cap\mathfrak m_S$, and its topology is
equivalent to its $\mathfrak m_O$-adic topology.
The weight-zero monomials in the chosen differences
have positive order in $S$ and lie in that maximal
ideal.  Their formal series therefore converge in
$C_0$, and the inclusion in $S$ identifies their
limits with the original series.  For any weight $w$,
the finite-monomial argument of PX1 writes a
homogeneous series as a finite sum of entry monomials
times such weight-zero series.  Thus its value is
an $O$-linear combination of entry monomials of
weight $w$.  This proves the exact span assertion,
including the removal of a possible topological
closure.  It immediately proves uniqueness of an
algebra map from all its matrix entries.  Torus
rescaling defines an algebra map on $\mathcal S$,
because weights add.

Here is the point-extension argument with its
convergence condition included.  For presentation
variables $z_s$ put $a_s=v(\xi(z_s))$ when their
values are nonzero, and write $w_s\in\mathbb Z^2$
for their weights.  If $n_s\ge0$ are integers, not
all zero, with $\sum n_sw_s=0$, then
\[
                   \sum n_sa_s>0.
 \tag{PX.13}
\]
To prove this, the corresponding monomial $u$ has
weight zero and positive formal order.  If
$\xi(u)^{-1}$ were integral over $O$, a monic
polynomial annihilating it could be reversed to
give a polynomial in $u$ with constant coefficient
$1$, hence a unit in the weight-zero formal ring,
whose image under $\xi$ is zero.  This is impossible.
The valuation characterization of integrality in
an algebraic extension of a complete DVR proves
(PX.13).

The finite strict linear inequalities
$a_s+\langle h,w_s\rangle>0$ therefore have a
solution $h\in\mathbb Q^2$.  For completeness,
failure of these inequalities separates the affine
subspace $a+\{(\langle h,w_s\rangle)_s\}$ from
the positive orthant, producing a nonzero vector
$q_s\ge0$ with $\sum q_sw_s=0$ and
$\sum q_sa_s\le0$.  All data are rational, so
a rational such vector can be chosen by solving
the finite linear equalities on a supporting face.
Clearing denominators contradicts (PX.13).
A real feasible vector may likewise be approximated
by a rational one, because the inequalities are
strict.  Choose $t_1,t_2\in K^\times$ of valuations
the two coordinates of $h$; roots of a uniformizer
provide them.  All scaled variable values now have
positive valuation.  Variables whose values were
zero impose no condition.

The scaled values lie in one finite extension of
$E$.  Adjoin to $O$ these finitely many values,
generators of the finite module $\xi(C_0)$, and
$y_s^{N_s}/\varpi$ with $N_s$ large enough that
these last elements are integral.  This gives a
finite $O$-algebra $C\subset K$ in which all $y_s$
are in the radical of $\varpi C$.  It is complete,
so the power series may be evaluated in $C$.
This evaluation is initially defined on the formal
source, before descent to $S$.  Let $P_0$ be the
full-weight-zero subring of that source.  Its map
$W_0:P_0\to C_0=\mathcal S_0$ is local and continuous:
both rings are complete local with residue field $k$,
and every nonconstant monomial maps into
$C_0\cap\mathfrak m_S$.  The restriction $\xi|_{C_0}$
is continuous because it is $O$-linear on a finite
$O$-module.  Evaluation on $P_0$ is continuous as
well, since its finitely many nonconstant monomial
generators have positive valuation.  It agrees with
$\xi W_0$ on these generators, because their full
weight is zero.  Successive approximation therefore
proves agreement on all of $P_0$.
For every other weight the finite-monomial
decomposition gives agreement with $\xi W$ after
the indicated torus factor.  Consequently evaluation
kills each homogeneous component of a relation.
For each $n$, sum the complete weight components
whose weights occur in total degree less than $n$;
there are finitely many, and the remaining series
has order at least $n$.  Evaluation is continuous,
so letting $n$ grow shows that it kills the whole
relation.  Thus evaluation factors through $S$
without assuming that descent in advance.  The
resulting representation is an actual point by its
displayed \texttt{Is\allowbreak Pt} witness $\rho^u$.

Finally PX4 identifies $A/I$ with a subalgebra of
$\mathcal S$.  Its degree-zero piece maps into the
finite module $C_0$ by the span assertion.  A
submodule of a finite module over the Noetherian
ring $O$ is finite.  This proves $\mathtt{M2}$.
\end{proof}

\paragraph{PX6: normalized points and the three possible matrix patterns.}
\label{claim:PX6}
Every actual point has continuous entries in the integers
of one finite extension of $E$, strictly congruent to
$\rho_0$ in the $\ell$-adic norm.  Every diagonal matrix
unit $e_{ii}$ belongs to the span of its matrices, already
with integral coefficients in a common finite extension.
For an actual point $\rho$, write $i\rightsquigarrow j$
if $\rho(g)_{ij}\ne0$ for some $g$.  Exactly the following
exhaustive alternatives are needed:
\[
 \begin{split}
 &\text{the matrix span of $\rho$ is $M_4(K)$;}\\
 &i\ne j,\ i\rightsquigarrow j\ \Longrightarrow\
                                      j\not\rightsquigarrow i;\\
 &0\rightsquigarrow2,\ 2\rightsquigarrow0,
 \quad\text{and all cross-block entries vanish in one direction.}
 \end{split}
 \tag{PX.14}
\]
Here cross blocks refer to $\lambda=1$ and $\lambda=0$.
The second alternative is the source \texttt{Diag\allowbreak Type}.
This conclusion uses the residual vanishing already in
\texttt{Good} and finite local Tate duality and Euler at
the primes in $S_b$; no automorphy or irreducibility
criterion is an unstated input.

\begin{proof}
An algebra map $y:C\to K$ from an actual complete local
coefficient algebra is automatically bounded in the
needed sense.  If $c\in\mathfrak m_C$ and $|y(c)|\ge1$,
then $y(c)^{-1}$ is integral over $O$.  Reverse a monic
polynomial annihilating it.  Its value at $c$ is a unit,
since its constant coefficient is $1$, but its image is
zero, a contradiction.  Thus $|y(c)|<1$ on
$\mathfrak m_C$.  Every element is a scalar from $O$
plus an element of that ideal, so all values have norm
at most $1$.  Finitely many generators of
$\mathfrak m_C$ give a uniform bound $a<1$ on its
image and a bound $a^q$ on its $q$th power.  This
proves continuity, including continuity of the
specialized representation.  A finite formal
presentation now puts all values in the finite
extension generated by its finitely many variable
values; convergence takes place in that complete
field.  The residual identity gives
\[
                  |\rho(g)_{ij}-\rho_0(g)_{ij}|<1.
 \tag{PX.15}
\]

We next spell out why normalization supplies the
matrix projectors.  AX3 and AX6 give that
$\rho(\phi)$ has order prime to $\ell$, is block
diagonal, and has different residual eigenvalues
on the two blocks.  A matrix of prime-to-$\ell$
order which is scalar modulo the maximal ideal
is that scalar's unique finite-order lift:
factor $X^d-1$ at that root, whose complementary
factor is a unit.  Applying this to each block
shows that the two block projectors are
polynomials in $\rho(\phi)$ with integral
coefficients.  In the $\{0,2\}$ block the
self-adjoint element
\[
 H=\rho(g_0)+\theta_0(g_0)\rho(g_0^{-1})
 \tag{PX.16}
\]
has zero off-diagonal entries by normalization.
Its diagonal difference is a unit, by AX6:
its residue is
$(\bar\alpha(g_0)-1)(\bar\Psi(g_0)-1)$.
Linear interpolation in this block gives $e_{00}$
and $e_{22}$.  The symplectic adjoint gives
$e_{33}$ and $e_{11}$.  The span of the group
matrices is closed under multiplication and
adjoint, because products are group matrices
and the adjoint of $\rho(g)$ is its prescribed
similitude times $\rho(g^{-1})$.  This proves
the projector assertion over the coefficient
algebra and after every specialization.

It follows that $\rightsquigarrow$ is reflexive
and transitive, and is carried by symplectic
duality to $j^*\rightsquigarrow i^*$, where
$i^*=3-i$.  Indeed a nonzero entry and the
two projectors supply the matrix unit $e_{ij}$;
multiply matrix units to prove transitivity.
The full matrix span is equivalent to every
pair being linked.

The substantive step is that a two-element
strongly connected component cannot straddle
the two $\lambda$ blocks.  Suppose it consists
of $s,u$ with $\lambda_s\ne\lambda_u$.
The component gives a two-dimensional
subquotient; products of its matrices have no
contribution through an intermediate vertex
outside this component.  The ordinary flag
of AX8 gives this local subquotient a stable
line at $\ell$.  In one of the two coordinate
charts that line has slope $v$ with $|v|\le1$;
interchange $s,u$ if necessary.  Both links
exist, so the chosen off-diagonal coefficient
is nonzero.  In a common coefficient DVR
$O'$, let $b\ne0$ generate the ideal of all
$\rho(g)_{su}$ and write
$\rho(g)_{su}=b c(g)$ with $c(g)\in O'$.
The reductions of the $c(g)$ are not all zero.
The matrix-product identity gives the residual
cocycle
\[
 z(g)=\frac{\overline{c(g)}}{\bar\theta_u(g)},
 \qquad \gamma=\bar\theta_s\bar\theta_u^{-1}.
 \tag{PX.17}
\]
It cannot be a coboundary.  If
$\overline{c(g)}=a(\bar\theta_s(g)-\bar\theta_u(g))$,
apply the integral linear combination of group
matrices giving $e_{ss}$.  Its normalized
off-diagonal entry gives $0=a$, contradicting
that some $\overline{c(g)}$ is nonzero.

On the other hand this cocycle has all the
local conditions of $D.\mathtt{XVan}(\theta_s/\theta_u)$.
Outside $T$, the actual unramified condition
kills it on inertia intersected with
$\ker\bar\rho$.  At $p\in S_b$ the local
$H^1$ of its residual character is zero.
Here is the precise duality calculation.
AX6 provides $\sigma_p\in G_p$ for which
$\gamma(\sigma_p)\ne1$ and the residual
cyclotomic value is $1$.  Thus both $H^0(k(\gamma))$
and $H^0(k(\gamma^{-1}\bar\chi_{\rm cyc}))$
vanish.  Local Tate duality gives $H^2=0$;
the local Euler formula at $p\ne\ell$ then
gives $\dim H^1=0$.  This is valid for
the actual finite coefficient extension $k$.

At $\ell$, the stable-line identity is
\[
 \rho(g)_{ss}v+b c(g)
             =(\rho(g)_{us}v+\rho(g)_{uu})v.
 \tag{PX.18}
\]
Apply it also to $\sigma_\ell$ and eliminate
$b$ and the linear slope term.  Reduction
gives
\[
 \overline{c(g)}(\bar\theta_u(\sigma_\ell)
                      -\bar\theta_s(\sigma_\ell))
 =\overline{c(\sigma_\ell)}
                     (\bar\theta_u(g)-\bar\theta_s(g)).
 \tag{PX.19}
\]
The discarded terms are multiples of
$v\rho(g)_{us}$ or $v\rho(\sigma_\ell)_{us}$;
their norms are strictly less than $1$ by
(PX.15).  The first parenthesis is nonzero
by AX6, so (PX.19) makes (PX.17) a local
coboundary.  The ratio $\theta_s/\theta_u$
is one of $\alpha,\beta,\theta$ or their
inverses, because the first weights differ.
The residual vanishing in \texttt{Good}
applies.  It remains valid after extension
of the residue field: expand the cocycle
and its local coboundary constants in a
basis over the original residue field and
apply the original vanishing coordinate by
coordinate.  Thus (PX.17) is a global
coboundary, the required contradiction.

Finish by considering the strongly connected
components on four vertices.  A three-vertex
component and its dual overlap, hence coincide;
this is impossible since duality has no fixed
vertex and a stable subset has even size.
A four-vertex component is the first case of
(PX.14).  If all components are singletons
we have \texttt{Diag\allowbreak Type}.  Otherwise the
only permitted two-vertex component is one
of $\{0,2\}$ and $\{1,3\}$; duality gives
both.  Arrows between these two components
can point in at most one direction, or all
four vertices would be strongly connected.
This is exactly the third case in (PX.14).
\end{proof}

\paragraph{PX7: one-sided reducible points, two-sided irreducibility and M3.}
\label{claim:PX7}
Assume \texttt{Fin0} and \texttt{No\allowbreak End}.  Every two-sided
point of $A/I$ gives, through PX4 and PX5, an irreducible
actual representation.  Moreover $A,I,x_0$ satisfy the
literal $\mathtt{M3}$ condition.  The latter conclusion
uses \texttt{Fin0} but does not need \texttt{No\allowbreak End}.

\begin{proof}
First consider an actual representation of
\texttt{Diag\allowbreak Type}.  Transitivity and the absence
of a two-way link make its diagonal entries
characters $\kappa_i$.  Restrict to $G_\ell$.
The ordinary flag of AX8 and the equality of
traces identify their multiset with the four
ordinary diagonal characters.  To see that
the matching is unique, restrict to $s_0$:
its four character values are distinct by
AX6.  Independence of distinct characters
then matches them as characters, not only
as values at one element.  Residual separation
at $\sigma_\ell$ preserves $\lambda$, and
symplectic duality preserves $i\mapsto3-i$.
There are precisely two resulting permutations:
\[
                 m=\mathrm{id}\quad\hbox{or}\quad
                 m=(2,3,0,1).
 \tag{PX.20}
\]
This last enumeration follows by choosing the
image of $0$ in $\{0,2\}$; duality fixes
the image of $3$, and the remaining two
images are then fixed.  Thus on $I_\ell$
we have $\kappa_i=\theta_{m(i)}$.

Each $\kappa_i$ is
$\chi_{\rm cyc}^{h_{m(i)}}\delta_i$, with
$(h_0,h_1,h_2,h_3)=(e+1,e,1,0)$ and
$\delta_i$ of finite order.  Here is the
global finiteness check.  Divide by the
displayed cyclotomic power.  On inertia
this quotient is trivial on the intersection
with one fixed open subgroup
$\ker\bar\rho\cap\ker\psi\cap\ker\omega$,
using the local conditions and AX8 at
$\ell$.  Raising to its index makes the
character unramified at every finite prime.
An everywhere-unramified abelian character
of $G_{\mathbb Q}$ is trivial, as in PX3.
Thus a fixed power of $\delta_i$ is trivial;
continuity gives its open kernel.  This
argument does not assume in advance that
an arbitrary algebraic-valued character
has finite image.

We claim that a link $i\rightsquigarrow j$
between distinct vertices forces $m(i)<m(j)$.
It is enough to prove this for a cover link,
with no vertex strictly between the two:
any link in a finite partially ordered
set is a chain of cover links.  For such
a link put
\[
 z(g)=\rho(g)_{ij}/\kappa_j(g),
 \qquad \eta=\kappa_i/\kappa_j.
 \tag{PX.21}
\]
There is no intermediate contribution to
matrix multiplication, so this is a
continuous $\eta$-cocycle.  If $m(i)>m(j)$,
then $\eta$ is a finite character times
a strictly negative power of cyclotomic,
since $e+1>e>1>0$.  The ordinary local
flag splits this reversed extension at
$\ell$: its higher-character line is
complementary to the lower-character
line, and the distinct $s_0$ eigenvalues
make the complement a graph.  Equivalently
the $2$-by-$2$ quotient has
$\rho(g)_{ij}=a(\kappa_i(g)-\kappa_j(g))$
on $G_\ell$.  At every other inertia
group the actual unramified equations
and a finite-index argument give the
same coboundary condition.  RV7 now
makes (PX.21) a global coboundary.
Applying the linear combination giving
$e_{ii}$ forces its constant to be zero,
so the entry vanishes for every $g$,
contrary to the link.  The claim follows.

When $m=\mathrm{id}$, every negative
second-weight entry therefore vanishes,
and $\rho(s_0)_{00}=\theta_0(s_0)$.
When $m=(2,3,0,1)$, every positive
second-weight entry vanishes, and
$\rho(s_0)_{00}=\theta_2(s_0)$.
These two distinguished values differ.
By PX5 a negative-weight entry must
occur in every entry monomial of
negative total second weight; similarly
for positive weight.  Hence these
entry vanishings imply respectively
$\mathtt{NegZero}$ and $\mathtt{PosZero}$
on the restricted point of $A$.
Conversely $\mathtt{PosZero}$ kills
every $(0,2)$ entry and
$\mathtt{NegZero}$ every $(2,0)$ entry,
because their weights are $(0,1)$ and
$(0,-1)$ and PX4 lifts those pieces.

Apply PX6 to a torus-scaled actual
representative of a two-sided point.
The diagonal-type case has just been
excluded.  In the block-triangular
case, scale by the first torus parameter
and let it tend to zero in the direction
which kills the cross-block entries.
This is an algebraic specialization,
not an analytic limiting assertion:
entry monomials have only nonnegative
powers of the parameter wherever their
values are nonzero, by the one-sided
block condition and PX5.  Evaluation
at zero is an algebra map on
$\mathcal S$ retaining its first-weight-zero
pieces.  PX5 rescales its resulting
representation to an actual point.
Its $(0,2)$ and $(2,0)$ links survive;
the projectors and these links make its
$\{0,2\}$ block span $M_2(K)$.
This contradicts \texttt{No\allowbreak End}.  Only
the irreducible case remains.

For $\mathtt{M3}$ take a point $x_1$
with both signs zero, vanishing on $I$,
and satisfying the source's additional
inequality $x_1\ne x_0$.  Lift and
scale it.  The absent $(0,2)$ and
$(2,0)$ links exclude the first and
third cases of PX6, so its representation
has diagonal type.  Points with the
same degree-zero restriction have the
same diagonal entries: each is of
full weight zero.  They also share
diagonal type, since every product
$\rho(g)_{ij}\rho(h)_{ji}$ is a
degree-zero function and vanishes
for a diagonal-type representation.
They must have the same alternative
in (PX.20), because their $(0,0)$
entries at $s_0$ agree and the two
possible values are distinct.
Thus they all have the same sign
vanishing.  One of the two opposite
one-sided families in the definition
of $\mathtt{M3}$ is empty, exactly
as required.  The proof even gives
this conclusion without using the
extra inequality $x_1\ne x_0$.
\end{proof}

\paragraph{PX8: normalized conjugacy is torus conjugacy, with the correct orbit.}
\label{claim:PX8}
Two conjugate actual points, one of which is irreducible,
are conjugate by the torus in (PX.1).  Consequently, under
\texttt{Fin0} and \texttt{No\allowbreak End}, conjugate representations assigned to
two-sided points of $A/I$ come from the same orbit of the
second one-dimensional torus on $A$.

\begin{proof}
Let $U\rho(g)=\rho'(g)U$ with $U$ invertible.
The normalization and (PX.15) force $U$
to be block diagonal.  Here is an elementary
way to see the separation argument, also
valid without diagonalizing either matrix.
In an off-block entry of
$U\rho(\phi)=\rho'(\phi)U$, subtract the
common residual diagonal matrix.  Choose
an off-block entry of $U$ of maximum norm.
Its coefficient is a difference of two
residually distinct diagonal values and
has norm $1$.  Every error term is an
off-block entry of norm at most that
maximum, multiplied by an error of
norm strictly less than $1$.
The ultrametric inequality gives a
strict contradiction unless the maximum
is zero.

Intertwining also holds for the two
elements $H,H'$ in (PX.16).  Their
$(0,2)$ and $(2,0)$ entries are zero,
and the required opposite diagonal
values are distinct by the same
unit separation from AX6 and (PX.15).
The $(0,2)$ and $(2,0)$ entries of
$UH=H'U$ therefore force
$U_{02}=U_{20}=0$.
Moreover $U^*U$ commutes with every
$\rho(g)$, since both representations
have the same similitude.  Irreducibility
means their matrix span is all $M_4(K)$,
so $U^*U$ is scalar.  This follows
directly by commuting with the matrix
units; no version of Schur's lemma
over an unspecified endomorphism
field is needed.  Thus $U$ is a
symplectic similitude.  Its block
form and the two zero corner entries
then force the remaining off-diagonal
entries $U_{13},U_{31}$ to vanish.
Its diagonal entries $d_i$ obey
$d_0d_3=d_1d_2$ and are nonzero.
Dividing by $d_3$ gives exactly
\[
       (d_0,d_1,d_2,d_3)/d_3
             =(t_1t_2,t_2,t_1,1),
 \qquad t_1=d_2/d_3,\ t_2=d_1/d_3.
 \tag{PX.22}
\]
Scalar matrices act trivially, proving
the claimed torus conjugacy.

Now let $\xi,\xi'$ lift two points of
$A/I$.  PX5 scales them by torus
elements $t,t'$ to actual points.
If these are conjugate, (PX.22)
applies because PX7 makes the two-sided
representations irreducible under the
two stated hypotheses.  Composition of scalings gives a
torus element carrying the evaluated
representation of $\xi$ to that
of $\xi'$.  Entry generation in
PX5 makes the corresponding algebra
maps equal on all of $\mathcal S$.
Restrict to $A$.  Its first weights
are zero, so the first parameter
disappears and only the second
parameter acts.  This is the exact
\texttt{Same\allowbreak Orbit} in Lemma M;
quotienting by a larger conjugacy
group has not lost multiplicities.
\end{proof}

\paragraph{PX9: lifting derivations and the exact one-dimensional tangent bound.}
\label{claim:PX9}
Assume \texttt{Fin0}, \texttt{No\allowbreak End} and
\texttt{All\allowbreak Irr Rigid\allowbreak At}.  If $x:A\to K$ vanishes on
$I$, is two-sided, and $\delta$ is an $O$-linear
point derivation at $x$ killing $I$, then there is
$b\in K$ such that
\[
                   \delta(f)=b n x(f)
                       \quad(f\in A_n).
 \tag{PX.23}
\]
In particular the exact predicate $\mathtt{M4}$ holds.

\begin{proof}
Package the derivation as
$X=x+\epsilon\delta:A\to K[\epsilon]$.
PX4 first gives a lift $\xi_0:\mathcal S\to K$
of $x$.  PX7 makes its representation
irreducible.  Hence there are entry functions
$u\in\mathcal S_{(1,0)}$ and
$v\in\mathcal S_{(-1,0)}$ with
$\xi_0(u)\xi_0(v)\ne0$ (use entries
$(0,1)$ and $(1,0)$).  PX4 gives
$a_0\in A_0$ with $\varphi(a_0)=uv$.
The element $c=X(a_0)$ is a unit in
$K[\epsilon]$.

This supplies an explicit extension of
$X$ to $\mathcal S$.  After inverting
$uv$, the first-weight grading has
an invertible element $u$ of degree
$1$ and its degree-zero subring is
$\varphi(A)[(uv)^{-1}]$.  Every element
of first degree $a$ is uniquely $u^a$
times an element of that subring.
Send $u$ to $1$ and extend the map
$X$ on the degree-zero subring.
In nonlocalized notation, multiply
a homogeneous $f$ of first degree
$a$ by $u^{(-a)_+}v^{a_+}$, lift
this balanced element to $A$, apply
$X$, and divide by $c^{a_+}$.
The kernel assertion of PX4 proves
independence of the lift; balancing
two factors proves multiplicativity.
This gives
$\Xi:\mathcal S\to K[\epsilon]$
with $\Xi\varphi=X$.  Its reduction
need not equal the initially chosen
$\xi_0$; its restriction to $A$ is
$x$, which is exactly what is needed.

Scale the reduction of $\Xi$ to an
actual point using PX5, and use the
same constant torus element on $\Xi$.
The resulting representation over
$K[\epsilon]$ is itself an actual
point.  We include the reason for
this sometimes delicate passage.
The reduced entries satisfy (PX.15).
The span assertion of PX5 implies
that all reduced elements of
$\mathcal S$ have norm at most $1$,
and reduced formal variables have
norm strictly less than $1$; to
check the latter subtract their
original scalar residues and use
the same bound on entry polynomials.
In a dual-number algebra, an element
with integral first component is
integral over $O$: its nilpotent
part is integral, and integrality
is closed under addition.  Adjoin
the finitely many variable values,
the finite module of weight-zero
values, and sufficiently large
powers of the variables divided
by a uniformizer, just as in PX5.
The resulting finite $O$-algebra
is complete, including its nilpotent
ideal; the variables are in the
radical of the uniformizer ideal.
The identical finite-monomial and
truncation argument extends $\Xi$
to $S\to K[\epsilon]$.  Thus the
original universal \texttt{Is\allowbreak Pt}
lift is an actual witness for
$D.\mathtt{PtIn}$, as required by
the rigidity hypothesis.

Write this reduced actual point as
$\rho$; it is irreducible by PX7.
Rigidity gives the derivative
$[Y,\rho]$ with
$Y\in\mathfrak{sp}_4(K)$.
The differentiated block normalization
at $\phi$ and the maximum-norm
argument of PX8 force $Y$ to be
block diagonal.  Differentiating
the two zero corners of (PX.16)
then forces $Y_{02}=Y_{20}=0$,
because $H_{00}-H_{22}\ne0$.
The symplectic Lie-algebra equation
kills $Y_{13},Y_{31}$ as well.
Thus $Y$ is diagonal, and there
are $a,b\in K$ with
\[
 [Y,\rho(g)]_{ij}
   =(a(\lambda_i-\lambda_j)
                      +b(\mu_i-\mu_j))\rho(g)_{ij}.
 \tag{PX.24}
\]
This is the infinitesimal torus
$(1+\epsilon a,1+\epsilon b)$.
Its action on the constant point
has the same entries as the
dual-number point.  PX5 gives
equality of the maps on all of
$\mathcal S$.  On $A_n$, the
first-weight contribution is zero.
Undo the earlier constant scaling;
its nonzero factor cancels from
both sides, giving (PX.23).

For two derivations with parameters
$b,b'$, choose coefficients
$(b',-b)$ unless $b=0$, in which
case choose $(1,0)$.  They are
not both zero and their linear
combination vanishes on every
homogeneous piece by (PX.23),
hence on $A$.  This proves the
literal $\mathtt{M4}$, including
the actual-point restriction in
its use of rigidity.
\end{proof}

\paragraph{PX10: the trace element, orbit cover and actual Selmer inequality.}
\label{claim:PX10}
Under the same three arithmetic assertions, let
$g\in G$ satisfy
\[
 \operatorname{tr}\rho(g)-\operatorname{tr}\rho_0(g)\ne0
       \quad\text{for every irreducible actual point }\rho.
 \tag{PX.25}
\]
Then the literal source inequality $D.\mathtt{Ineq}(g)$
holds.  In particular its Selmer length is finite and
there is a finite family $\mathcal X$ of irreducible
actual representations, with no two distinct members
conjugate, such that
\[
 \operatorname{nrm}_O(D.\mathtt{sha})\log\ell
 \ \le\ 2\sum_{\rho\in\mathcal X}
 -\log\left|\operatorname{tr}\rho(g)
                         -\operatorname{tr}\rho_0(g)\right|.
 \tag{PX.26}
\]
This closes the entire internal model-to-comparison
part of SP's $G2$, with exactly \texttt{Fin0},
\texttt{No\allowbreak End} and \texttt{All\allowbreak Irr Rigid\allowbreak At} retained
as arithmetic input predicates.

\begin{proof}
Choose weight-zero lifts $f_i\in R_0$ of the
four diagonal entries of $\rho^u(g)$.  Set
\[
           a_g=\sum_{i=0}^3 f_i
                       -\operatorname{tr}\rho_0(g)\in A_0.
 \tag{PX.27}
\]
The model's origin identity gives $x_0(a_g)=0$;
at every point of $\mathcal S$ its value is
exactly the trace difference.  Torus
conjugation leaves this value unchanged.
PX7 and (PX.25) therefore prove $\mathtt{M5}$.

All the other hypotheses of the full
intersection theorem SP6 are now proved
on this same $A$: PX1 gives the weighted
ring, PX2--PX3 give $\mathtt{M1}$, PX5
gives $\mathtt{M2}$, PX7 gives
$\mathtt{M3}$, and PX9 gives
$\mathtt{M4}$.  Use the original
normalized valuation
\[
 v(z)=\frac{v_O(\ell)}{\log\ell}(-\log|z|),
                  \qquad |\ell|=\ell^{-1}.
 \tag{PX.28}
\]
Multiplicativity, the ultrametric inequality
and the DVR normalization check precisely
$\mathtt{IsNVal}$; no coefficient
extension is introduced by this step.
SP6 supplies a finite cover of all
two-sided orbits and the bound
\[
 \operatorname{length}_O C_{-1}(F)
                       \le2\sum_{x\in T}v(x(a_g)).
 \tag{PX.29}
\]
Choose such a cover of minimum cardinality.
It contains only two-sided points and
contains no two points in the same orbit,
since one of those could be deleted.
Lift and scale each point by PX4--PX5.
They are irreducible by PX7, and PX8
shows that distinct chosen points give
nonconjugate representations.  Thus
their trace values give exactly the
sum in (PX.26), without a multiplicity
loss or an enlargement of the set.

SP7 supplies the actual normalized
Selmer-cocycle injection into the
cotangent quotient on $M$, using
the model's original derivation
field and RV3's finiteness.  Together
with SP1 it bounds the original
Selmer length by the left side of
(PX.29), and proves that length is
finite before its real normalization.
Divide (PX.29) by $v_O(\ell)$ and
multiply by $\log\ell$ to obtain
(PX.26).  The factor is exactly
$2$.  Every representation, trace,
origin, equation, lattice and
coefficient normalization is the
one attached to the original $D$.
\end{proof}

\subsection{What the completed comparison leaves to the arithmetic count}

MX constructs the actual model from the local and
global installation in DX; AX constructs the original
$\mathtt{Good}$/$\mathtt{OfCase}$ datum and its exact
level and ordinary specializations.  Consequently
PX10 removes the internal weighted algebra, point
lifting, reducibility classification, tangent lifting
and orbit comparison from SP's remaining $G2$ input.
To use the paired-count conclusion SP10, the remaining
arithmetic theorem must still supply
$\mathtt{Fin0}$, $\mathtt{NoEnd}$, actual irreducible
rigidity and the common admissible representation
class with its uniform finite catalogue.  The image
theorem supplying the first two assertions is the
same one used in AX's final admissibility join.
This is a precise sufficient cut, not a proof of
the image theorem or a claim that an arbitrary
finite set of representations is an adequate
catalogue.  Once these inputs are supplied on
the specified objects, SP10 gives actual
$\mathtt{CountII}$ and $\mathtt{CountIII}$, and
HV10 already supplies their complete downstream
place in the original-field cohomological window
assembly.  The other actual counts and the
common-family, gap, theta, denominator and
metric/spectral branches keep their separate
roles in the full RH/Dirichlet-GRH reconstruction.


\section{Dependency closure after actual symplectic model construction}
\label{sec:dependency-audit-v27}

\paragraph{The original arithmetic datum is now constructed.}
AX carries the quantitative dyadic-character choice and the
previous residual vanishing theorem through both cases of the
literal \texttt{Raw/Good/Of\allowbreak Case} construction.  Required roots
are taken in the given coefficient ring, or in the permitted
finite extension with SP8/RV1 normalized coefficient return.
Its later arguments
prove the actual ordinary-point condition, level, duality and
safe unlinking fields.  One common complex embedding is chosen
before every irreducible point under exactly the retained
\texttt{Symplectic\allowbreak Four\allowbreak Image\allowbreak Dense} predicate.  This supplies SP10's
G1 in that stated scope.

\paragraph{The deformation theorem is installed on those objects.}
DX proves the exact finite local/global duality and Euler
specializations, the global and local \texttt{Phi\allowbreak Cond}, and the
literal framed \texttt{Nums}.  Explicit ordinary and finite-inertia
deformations give the required characteristic-zero charts and
the exact tangent restriction.  They supply the original
\texttt{Loc\allowbreak Base}; the full \texttt{Bloch\allowbreak Kato\allowbreak Local} package is not
an unexplained premise of this construction.

MX then closes the whole constructor: inertia local conditions,
universal framed presentation, homogeneous cuts, a domain quotient
with actual \texttt{On\allowbreak Comp} witnesses, representation generation,
and integral derivations attached to every actual Selmer cocycle.
The coefficient ring, cotangent lattice, matrix normalizations and
all preimage quantifiers are retained.  The result is an actual
\texttt{Model} for every original good symplectic datum.

\paragraph{The internal comparison reaches the arithmetic inequality.}
PX replaces the long invariant-chart calculation by a formal
weight-pairing argument and proves all required integral
generation and scaling statements.  It classifies the actual
points, identifies the two-sided irreducible locus, turns
conjugacy into the prescribed torus orbit, lifts tangent
derivations and applies the actual rigidity predicate.  The
previous weighted intersection theorem then gives the exact
factor-two trace/cotangent bound and the literal \texttt{D.\allowbreak Ineq}.
This last implication retains exactly the same datum's
\texttt{Fin0}, \texttt{No\allowbreak End} and \texttt{All\allowbreak Irr Rigid\allowbreak At}.

\paragraph{The remaining paired-count cut is explicit.}
A sufficient remaining input is the conjunction of
\texttt{Symplectic\allowbreak Four\allowbreak Image\allowbreak Dense}, actual irreducible rigidity on
each constructed datum, and SP9's exact uniform finite
automorphic catalogue.  Image-density supplies \texttt{Fin0},
\texttt{No\allowbreak End} and admissibility with a common embedding.  The
source \texttt{rig\allowbreak Sp4\_\allowbreak alt} implication, including any required
local Hodge/adjoint inputs, is still a continuation target; merely
naming \texttt{Rigidity\allowbreak Symplectic\allowbreak Four} does not close it.
Uniform catalogue bounds must precede the varying prime, original
field, embedding and representation family.

\texttt{Count\allowbreak II/III} remain open in the human reconstruction until
these inputs are supplied.  So do \texttt{Count\allowbreak I} and \texttt{Count\allowbreak V}.
HV10 and TW/OB/VW already close the entire downstream height/window
implication from the four-count cut.  The common geometric family,
arithmetic gap, genuine theta/local installation, actual denominator
tables and exact metric/global/spectral comparisons remain separate.
The machine-readable semantic map is a finite sufficient-condition
audit, not an exhaustive Lean dependency graph or a completion
percentage.  Every historical proof and original source record is
preserved, with additive current-scope annotations.

% Version 0.28. Literal local Bloch--Kato and full CohomologyHyp installation.
% Pinned source: 7cbf7a0c32fca1badf05931879a37c3f8a2b03b4.
\section{Closing the local Bloch--Kato package on the original period models}
\label{sec:bloch-kato-installation}

\status{BX1--BX8 supply the full source predicate
\texttt{Published.\allowbreak Bloch\allowbreak Kato\allowbreak Local}, including both integral annihilator
directions and the original cocycle representatives.  Together with
LK, LD and DX this supplies the entire \texttt{Cohomology\allowbreak Hyp} record.
The classical inputs are Fontaine's standard period-ring construction
and the local theorem of Bloch--Kato for arbitrary de Rham
representations.  Their use is matched to the literal source models;
the global Bloch--Kato conjecture is not an input.}

\subsection{The interface that must actually be supplied}

For this chapter $\ell$ is any prime, including $2$, and
$K/\mathbb Q_\ell$ is an arbitrary finite extension inside the source's
specified algebraic closure.  Write $G_K$ for its absolute Galois group.
The input is an arbitrary \texttt{Loc\allowbreak Prime} $P$, an arbitrary
\texttt{Loc\allowbreak Ainf} $AI$ for it, and a witness
$T:\texttt{Loc\allowbreak Tied}(\ell,P,AI)$.  Thus the argument must work for every
such presentation of the period rings.  It need not, and does not,
claim a Bloch--Kato theorem for an abstract \texttt{Period\allowbreak Data} that has
no \texttt{Loc\allowbreak Tied} witness.

Put $Q=\mathbb Q_\ell$, $O=\mathbb Z_\ell$, and $PD=P.PD(K)$.
For a continuous finite-dimensional $Q$-representation $V$ the
source uses
\[
 \begin{split}
 Z^1(V)&=\{a:G_K\longrightarrow V:
     a\text{ continuous},\
     a(gh)=a(g)+g a(h)\},\\
 Z_f^1(PD,V)&=\{a\in Z^1(V):
     \exists x\in V\otimes_Q P.R.B,\
     a(g)\otimes1=gx-x\quad(g\in G_K)\},\\
 B^1(V)&=\{g\longmapsto gv-v:v\in V\}.
 \end{split}                                                   \tag{BX.1}
\]
Its \texttt{Finite\allowbreak Cohomology} is $Z_f^1/B^1$.  If $L\subset V$ is a
stable $O$-lattice, its finite lattice cocycles are
$Z_f^1(PD,V)\cap Z^1(G_K,L)$.  For $A=V/L$, its \texttt{Hf\allowbreak D} is the
sum of the image of $Z_f^1(PD,V)$ in functions to $A$ and the actual
coboundaries in $A$.  These are statements about representatives,
not merely unnamed cohomology subspaces.

Only two classical results specific to period rings are used below.
Fontaine's construction identifies the usual
$A_{\mathrm{cris}}$, $B_{\mathrm{cris}}$, $B_{\mathrm{dR}}^+$
and their filtration; see \emph{Le corps des p\'eriodes
$p$-adiques}, Ast\'erisque 223 (1994), \S1.5 and
\S2.3.3--2.3.4.
Bloch--Kato, \emph{$L$-functions and Tamagawa numbers of motives},
\S3, Proposition 3.8 and Corollary 3.8.4, give the finite-part
dimension formula and rational local annihilator theorem.
Both apply to every finite $K/Q$, every continuous de Rham $V$,
and every prime $\ell$.  Finite local Tate duality is the already
matched LK3/DX1 input.  The following proofs supply the intervening
identifications, lattice passage and record assembly.

% BEGIN BX1
\begin{proposition}[BX1: the actual tied period models and twist]
\label{claim:BX1}
The maps in \texttt{Loc\allowbreak Tied} identify $P.R.B$ and $P.R.BdR$,
as $Q$-algebras with $G_K$-action, with the standard
$B_{\mathrm{cris}}$ and $B_{\mathrm{dR}}$.
They carry the source filtration to
$\operatorname{Fil}^iB_{\mathrm{dR}}=t^iB_{\mathrm{dR}}^+$.
The character $PD.\chi$ is the actual $\ell$-adic cyclotomic
character on $G_K$.
\end{proposition}
\begin{proof}
The source's \texttt{Per} rings are explicit constructions.
For $C=\mathbb C_\ell$ they start with
$R=\varprojlim_{x\mapsto x^\ell}O_C/\ell$ and
$A_{\inf}=W(R)$, and the usual map
$\theta:A_{\inf}\to O_C$.
The ring called \texttt{Acris0} is the $A_{\inf}$-subalgebra of
$A_{\inf}[1/\ell]$ generated by
\[
                   \frac{x^n}{n!}\quad
          (x\in\ker\theta,\ n\ge0).
                                                               \tag{BX.2}
\]
Its $\ell$-adic completion is \texttt{Acris}.
This is the standard integral crystalline period ring: if
$\ker\theta=(\xi)$, then $x=a\xi$ gives
$x^n/n!=a^n\xi^n/n!$, so the displayed algebra is exactly
$A_{\inf}[\xi^n/n!:n\ge0]$.
This also identifies its specified action, induced on the tilt,
Witt vectors, divided powers and completion.

The source first inverts $\ell$ and then all logarithmic periods
of primitive compatible systems of $\ell$-power roots of unity.
This last localization is the usual inversion of a single $t$.
Indeed $T_\ell(\mu)$ is a free rank-one $O$-module: fixing a
primitive compatible system identifies every other such system
with its $u$-th power for a unique $u\in O^\times$.
The period logarithm is $O$-linear on this module, hence the
associated period is $ut$.  Its linearity follows from the
formal identity $\log(ab)=\log a+\log b$ on the divided-power
ideal and passage to compatible $\ell$-adic approximations.
All of these $u$ are units already before localization.
The two localizations therefore have inverse maps by their
universal properties.  This proves the claimed identification
of \texttt{Per.\allowbreak Bcris}.

Similarly, \texttt{Per.\allowbreak Bd\allowbreak Rp} is the completion of
$A_{\inf}[1/\ell]$ for powers of $\ker\theta$, and
\texttt{Per.\allowbreak Bd\allowbreak R} is its fraction field.  The completed kernel
is generated by $t$; its powers give $t^nB_{\mathrm{dR}}^+$.
The source's definition of negative filtration degrees uses the
quotient of the corresponding fractional ideals and thus gives
$t^iB_{\mathrm{dR}}^+$ for every integer $i$.
These are precisely the standard ring and filtration in the
classical statements just cited.

The supplied ring isomorphisms $T.eB$ and $T.eD$ commute with
the specified $O$-scalar maps by \texttt{z\allowbreak B}, \texttt{z\allowbreak D}.
They consequently commute with every $Q$-scalar: write
$q=\ell^{-m}z$, with $z\in O$, and use preservation of inverses.
The fields \texttt{act\allowbreak B}, \texttt{act\allowbreak D}, together with
$P.\texttt{PD\_\allowbreak act\allowbreak B}$ and $P.\texttt{PD\_\allowbreak actd\allowbreak R}$, give
$G_K$-equivariance.  The field \texttt{fil} gives the filtration
identification.  For the present endpoint these facts suffice;
no extra compatibility of a newly chosen ring embedding is
being inferred.

Finally let $g\in G_K$, and choose a primitive $\ell^m$-th root
$\zeta$.  The field $P.\texttt{chi\allowbreak Q\_\allowbreak cyc}$ gives an integer $a$
such that $g\zeta=\zeta^a$ and
$|\chi_Q(g)-a|_\ell\le\ell^{-m}$.
The first equality determines $a$ modulo $\ell^m$ as the
cyclotomic character.  The second identifies $\chi_Q(g)$ with
that residue modulo $\ell^m$.  For all $m$ these residues
determine a unique $\ell$-adic number, so
$\chi_Q(g)=\chi_\ell(g)$.  The equality $P.\texttt{PD\_\allowbreak chi}$
gives the asserted identity for $PD.\chi$.
No step here requires $\ell$ to be odd.
\end{proof}
% END BX1

% BEGIN BX2
\begin{proposition}[BX2: exact finite cocycles and period invariants]
\label{claim:BX2}
For every continuous finite-dimensional $Q$-representation $V$,
the source's $Z_f^1(PD,V)$ is the inverse image in $Z^1(V)$ of
the standard Bloch--Kato subspace $H_f^1(K,V)$.
Its \texttt{Dd\allowbreak R}, intersection with \texttt{Fil\allowbreak V(0)}, and
\texttt{Finite\allowbreak Cohomology} identify respectively with
$D_{\mathrm{dR}}(V)$, $\operatorname{Fil}^0D_{\mathrm{dR}}(V)$,
and $H_f^1(K,V)$ as $Q$-vector spaces.
In particular, \texttt{Is\allowbreak DR} is precisely the usual de Rham
condition.
\end{proposition}
\begin{proof}
Tensor the two $Q$-linear isomorphisms in BX1 with the identity
of $V$.  The resulting maps carry the source action
$g(v\otimes b)=gv\otimes gb$ to the standard diagonal action.
They therefore identify its common kernels of $g-1$ with the
usual period invariants.  The submodule \texttt{Fil\allowbreak V(0)} is the
image of $V\otimes_Q PD.\operatorname{Fil}^0$ in
$V\otimes_Q P.R.BdR$; its image is exactly
$V\otimes_Q B_{\mathrm{dR}}^+$.  Intersecting with invariants
gives the stated filtered subspace.

For each $a\in Z^1(V)$, the witness $x$ in (BX.1) is carried
bijectively to a standard $B_{\mathrm{cris}}$-coboundary witness.
Consequently the image of $[a]$ in
$H^1(K,V\otimes_Q B_{\mathrm{cris}})$ vanishes exactly when
$a\in Z_f^1(PD,V)$.  This is the definition of $H_f^1(K,V)$.
Every rational coboundary is finite, by the witness $v\otimes1$.
Taking the quotient by those actual coboundaries therefore
gives the asserted isomorphism on \texttt{Finite\allowbreak Cohomology}.
Continuity of $a$ is unchanged, since its coefficient space $V$
and topology have not changed.

The standard invariant space $D_{\mathrm{dR}}(V)$ is a
finite-dimensional $K$-vector space.  Its $Q$-dimension is
$[K:Q]\dim_KD_{\mathrm{dR}}(V)$.  Thus the source equality
\[
 \dim_Q\texttt{Dd\allowbreak R}(PD,V)=[K:Q]\dim_QV
                                                               \tag{BX.3}
\]
is exactly $\dim_KD_{\mathrm{dR}}(V)=\dim_QV$.
This verifies the precise antecedent of the classical theorem,
including the source's restriction of scalars.
\end{proof}
% END BX2

% BEGIN BX3
\begin{proposition}[BX3: the actual compact--discrete Tate passage]
\label{claim:BX3}
Let $L\subset V$ and $L'\subset V'=V^*(1)$ be dual stable
$O$-lattices, and put $A'=V'/L'$.
Then $H^1(K,L)$ is a finitely generated compact $O$-module,
$H^1(K,L)\otimes_OQ=H^1(K,V)$, and finite local Tate duality
induces a perfect continuous pairing
\[
 H^1(K,L)\times H^1(K,A')\longrightarrow Q/O.                    \tag{BX.4}
\]
The kernel of $H^1(K,V')\to H^1(K,A')$ is the image of
$H^1(K,L')$, a full lattice in $H^1(K,V')$.
The rational cup pairing
$H^1(K,V)\times H^1(K,V')\to Q$ is perfect and compatible
with (BX.4) by reduction modulo $O$.
\end{proposition}
\begin{proof}
Continuous cochains into $L$ are the inverse limit of continuous
cochains into $L/\ell^mL$.  At every finite quotient a cochain
has finite image, so it can be lifted to the next quotient by
choosing representatives on its finitely many clopen fibres.
The finite cohomology groups satisfy the Mittag--Leffler
condition: in each fixed finite group the descending images
stabilize.  The cocycle and coboundary exact sequences therefore
give
\[
 H^i(K,L)=\varprojlim_m H^i(K,L/\ell^mL).
                                                               \tag{BX.5}
\]
More explicitly, the possible obstruction is
$\varprojlim{}^1H^{i-1}(K,L/\ell^mL)$; the stabilized-image
property just proved makes it zero.  The usual proof chooses
successive cochain corrections after restricting to those
stable images.  In particular $M=H^1(K,L)$ is compact and
pro-$\ell$.

The long exact sequence for multiplication by $\ell$ embeds
$M/\ell M$ in the finite group $H^1(K,L/\ell L)$.
Choose finitely many lifts of its generators, and let $M_0$
be the closed $O$-submodule that they generate.  Then
$M/M_0=\ell(M/M_0)$.  A nonzero profinite pro-$\ell$ group has
a nonzero finite $\ell$-group quotient, on which multiplication
by $\ell$ is not surjective.  Hence $M/M_0=0$.
This proves finite generation and also identifies its compact
topology with its usual $O$-module topology.

A continuous cocycle on the compact group $G_K$ with values in
$V$ has bounded image and thus takes values in $\ell^{-m}L$
for some $m$.  A rational coboundary has the same property.
Localization of the cocycle quotient proves
$M\otimes_OQ=H^1(K,V)$.  The same argument applies to $L'$.
Continuous cochains into $A'$ have finite image and land in
some $\ell^{-m}L'/L'$, so
\[
 H^i(K,A')=\varinjlim_m H^i(K,\ell^{-m}L'/L').
                                                               \tag{BX.6}
\]
The two finite modules $L/\ell^mL$ and
$\ell^{-m}L'/L'$ are exact Tate duals with values in
$\ell^{-m}O/O(1)$.  LK3 and its normalized invariant identify
their cohomology pairings at each level.  The transition on the
left is reduction; its dual transition on the right is
inclusion.  Passing to (BX.5)--(BX.6) gives (BX.4), with its
original cup-product normalization.

The exact cohomology sequence of
$0\to L'\to V'\to A'\to0$ identifies the stated kernel.
Its image is the quotient of the finitely generated
$H^1(K,L')$ by its torsion and spans $H^1(K,V')$, so it is a
full lattice.  To check rational nondegeneracy without
localizing a torsion target, first repeat the finite-generation
argument in degree two.  Thus $H^2(K,L')$ is finitely generated
and the kernel of $H^2(K,L')\to H^2(K,V')$ is finite.
The same long exact sequence shows that
$N_0=\operatorname{im}(H^1(K,V')\to H^1(K,A'))$ has finite
index in $N=H^1(K,A')$.
By (BX.4), its annihilator in $M$ is finite.  Every nonzero
class in $M\otimes Q$ has, after scaling, a nontorsion
representative $m\in M$, so some element of $N_0$ pairs
nontrivially with $m$.  Lift that element to $H^1(K,V')$.
Compatibility of the cup formulas and the target map
$Q(1)\to Q/O(1)$ says that the nonzero value just obtained
is the reduction of the rational cup value.  Hence the
rational pairing has zero left kernel.  Interchanging
$L$ and $L'$ gives zero right kernel.  Both rational spaces
are finite-dimensional by the same finite-generation argument,
so the pairing is perfect.  This also proves the asserted
compatibility and normalization.
\end{proof}
% END BX3

% BEGIN BX4
\begin{proposition}[BX4: the literal finite-part dimension formula]
\label{claim:BX4}
For every representation in \texttt{Bloch\allowbreak Kato.\allowbreak Dim\allowbreak Formula}
whose source hypothesis \texttt{Is\allowbreak DR} holds,
\texttt{Finite\allowbreak Cohomology} is finite-dimensional and
\[
 \dim_Q H_f^1(PD,V)
   =\dim_Q V^{G_K}
      +\dim_Q D_{\mathrm{dR}}(PD,V)
      -\dim_Q\bigl(D_{\mathrm{dR}}(PD,V)
                    \cap\operatorname{FilV}(0)\bigr).          \tag{BX.7}
\]
Thus the complete source predicate
\texttt{Bloch\allowbreak Kato.\allowbreak Dim\allowbreak Formula}(PD) holds.
\end{proposition}
\begin{proof}
By BX2 the actual input representation is de Rham, with its
specified finite-dimensional module topology and continuous
action.  The finite-part formula of Bloch--Kato,
Corollary 3.8.4, applies:
\[
 \dim_QH_f^1(K,V)
      =\dim_QH^0(K,V)
          +\dim_Q\bigl(D_{\mathrm{dR}}(V)/
                        \operatorname{Fil}^0D_{\mathrm{dR}}(V)\bigr).
                                                               \tag{BX.8}
\]
It also gives finite dimensionality.  BX2 identifies the left
side and both period subspaces.  The source's \texttt{H0V} is
the intersection of the kernels of $\rho(g)-1$, hence exactly
$V^{G_K}$.  Since the filtered subspace is contained in the
finite-dimensional invariant space, the dimension of the
quotient in (BX.8) is the difference of the two dimensions
in (BX.7).  Interpreting that equality in $\mathbb Z$ gives
the source statement, whose subtraction is integer
subtraction.  This proof is uniform in $P,AI,T,K$ and in every
representation and topology quantified by \texttt{Dim\allowbreak Formula}.
\end{proof}
% END BX4

% BEGIN BX5
\begin{proposition}[BX5: both rational annihilator implications]
\label{claim:BX5}
The full predicate \texttt{Bloch\allowbreak Kato.\allowbreak Rational\allowbreak Duality}(PD)
holds for every $P,AI,T,K$ above and every source pairing satisfying
\[
                     \texttt{Is\allowbreak Dual1}(PD,\rho,\rho',P_0).
\]
\end{proposition}
\begin{proof}
Keep the source's spaces $V,V'$ and its given bilinear form
$P_0:V\times V'\to Q$.  The map
\[
             V'\longrightarrow V^*(1),\qquad
             y\longmapsto\bigl(x\longmapsto P_0(x,y)\bigr)
                                                               \tag{BX.9}
\]
is injective by the second nondegeneracy condition and
surjective by the equality of dimensions in \texttt{Is\allowbreak Dual1}.
Its equivariance follows by substituting $g^{-1}x,y$ in
$P_0(gx,gy)=\chi_\ell(g)P_0(x,y)$, where BX1 has identified
the exact character.  It is a continuous linear isomorphism
for the specified finite-dimensional module topologies.
Thus this is an identification of the original pairing with
the evaluation pairing, not a replacement by an unrelated
duality.

For de Rham $V$, the rational local theorem of Bloch--Kato,
Proposition 3.8 with both local conditions equal to $f$,
says that $H_f^1(K,V)$ and $H_f^1(K,V^*(1))$ are exact
annihilators under the rational Tate cup pairing.
In particular it applies without a crystalline hypothesis.
The isomorphism (BX.9) preserves finite cocycles: tensor it
with the identity of $B_{\mathrm{cris}}$ and carry the
explicit coboundary witnesses in both directions.
BX2 therefore identifies the two subspaces with the source's
ones for $V$ and $V'$.

The source's cup cocycle is exactly
\[
              (g,h)\longmapsto P_0(a(g),g b(h)).
                                                               \tag{BX.10}
\]
Its \texttt{Cob\allowbreak L} predicate means that this cocycle is the
continuous coboundary of a $Q(1)$-valued $1$-cochain.
Consequently \texttt{Cob\allowbreak L} is equivalent to its class
vanishing in $H^2(K,Q(1))$.  The normalized Tate invariant is
an isomorphism from this group to $Q$, so vanishing is
equivalent to annihilation in the scalar pairing.
Every finite cohomology class has a finite cocycle
representative by BX2.  Testing all source finite cocycles
is therefore exactly testing all elements of the relevant
$H_f^1$.  Applying the two annihilator equalities gives
both implications and both orders in the literal definition
of \texttt{Rational\allowbreak Duality}.
\end{proof}
% END BX5

% BEGIN BX6
\begin{proposition}[BX6: integral annihilators, including torsion]
\label{claim:BX6}
For dual stable lattices $L,L'$ in de Rham $V,V'$,
define
\[
 \begin{split}
 H_f^1(K,L)&=\{m\in H^1(K,L):m_Q\in H_f^1(K,V)\},\\
 H_f^1(K,A')&=\operatorname{im}\bigl(
                H_f^1(K,V')\longrightarrow H^1(K,A')\bigr).
 \end{split}                                                   \tag{BX.11}
\]
These two subgroups are exact annihilators under (BX.4),
without a torsion-freeness assumption on $H^1(K,L)$.
\end{proposition}
\begin{proof}
Put $M=H^1(K,L)$, $U=M\otimes_OQ$, $W=H_f^1(K,V)$,
and $F=\{m\in M:m_Q\in W\}$.
The subgroup $F$ contains all torsion in $M$.  Moreover
$M/F$ is torsion-free: if $\ell^jm\in F$, then
$\ell^jm_Q\in W$, and $Q$-linearity of $W$ gives $m_Q\in W$.
By BX3 it follows that $M/F$ is a finite free $O$-module,
of rank $\dim_QU-\dim_QW$.
The subgroup $F$ is closed, since $W$ is closed in the
finite-dimensional $Q$-space $U$ and $M\to U$ is continuous.

The perfect pairing (BX.4) identifies its annihilator with
\[
       \operatorname{Hom}_{\mathrm{cont}}(M/F,Q/O)
            \simeq (Q/O)^{\dim U-\dim W}.                     \tag{BX.12}
\]
Let $W'=H_f^1(K,V')$, and let $B$ be its image in
$N=H^1(K,A')$.  BX5 gives $W'=W^\perp$ in the perfect
rational pairing.  Cup-product compatibility in BX3
therefore gives $B\subset F^\perp$.
The kernel of $W'\to N$ is the intersection of $W'$ with
the full lattice that is the image of $H^1(K,L')$.
This intersection is a full lattice in $W'$: multiply a
$Q$-basis of $W'$ by a sufficiently large common power of
$\ell$ to put it in the ambient lattice, and bound its
intersection above by a finite lattice in that basis.
Hence
\[
                     B\simeq (Q/O)^{\dim W'}.
                                                               \tag{BX.13}
\]
Rational duality gives
$\dim W'=\dim U-\dim W$.

For completeness, an inclusion of divisible $\ell$-primary
$O$-modules of the forms (BX.12)--(BX.13), with the same
finite corank, is equality.  A divisible module over the
PID $O$ is injective: extending a homomorphism from the
ideal $\ell^jO$ amounts to solving one equation
$\ell^jx=y$.  It follows that the inclusion splits.
Its complementary summand is divisible and $\ell$-primary
with zero corank.  Such a summand is zero, since a nonzero
element of order $\ell$ in a divisible module has
compatible preimages of arbitrarily high $\ell$-power
order, producing a copy of $Q/O$.  Thus $B=F^\perp$.

Finally the double annihilator of the closed subgroup $F$
under the perfect compact--discrete pairing is $F$.
Explicitly, if $m\notin F$, its nonzero image in the finite
free module $M/F$ survives in a finite quotient; a character
of that finite quotient into $Q/O$ detects it.
This gives $B^\perp=F$ and proves both directions, with the
torsion subgroup retained throughout.
\end{proof}
% END BX6

% BEGIN BX7
\begin{theorem}[BX7: the complete literal local Bloch--Kato predicate]
\label{claim:BX7}
For every prime $\ell$, every $P:\texttt{Loc\allowbreak Prime}(\ell)$,
every $AI:\texttt{Loc\allowbreak Ainf}(\ell,P)$ and every
$T:\texttt{Loc\allowbreak Tied}(\ell,P,AI)$, all three clauses
\[
 \texttt{Dim\allowbreak Formula}(PD),\qquad
 \texttt{Rational\allowbreak Duality}(PD),\qquad
 \texttt{Integral\allowbreak Duality}(PD)
                                                               \tag{BX.14}
\]
hold for every finite extension $K/Q$.
This supplies \texttt{Global\allowbreak Theory.\allowbreak Published.\allowbreak Bloch\allowbreak Kato\allowbreak Local}
with its original quantifiers and representatives.
\end{theorem}
\begin{proof}
BX4 and BX5 give the first two clauses.  To identify the
third, retain the given form $P_0$ and lattices $L,L'$ in
\texttt{Is\allowbreak Dual\allowbreak Lat1}.  The isomorphism (BX.9) sends $L'$
onto the exact $O$-dual of $L$: this is precisely the
source biconditional
\[
 y\in L'
   \quad\Longleftrightarrow\quad
    P_0(x,y)\in O\ \hbox{ for every }x\in L.                   \tag{BX.15}
\]
The target lattice $L_1$ in the scalar cyclotomic
representation has underlying submodule equal to the
image of $O\to Q$, by the explicit hypothesis in
\texttt{Integral\allowbreak Duality}.  Its quotient is therefore
the actual $Q/O(1)$, by BX1.

A lattice cocycle $d$ belongs to the source's
\texttt{lattice\allowbreak Finite\allowbreak Cocycles} exactly when its rational
class lies in $H_f^1(K,V)$, by BX2.
For a divisible cocycle $c:G_K\to V'/L'$, the definition
of \texttt{Hf\allowbreak D} says
\[
             c=q(a)+\delta\bar y
       \quad\text{for some }a\in Z_f^1(PD,V'),\
                                \bar y\in V'/L'.              \tag{BX.16}
\]
This is equivalent to $[c]$ lying in the image in (BX.11).
The forward implication is immediate.  For the reverse,
equality of the two classes says exactly that their
cocycle difference is an actual coboundary, yielding
(BX.16).  There is no need to lift $c$ to a rational
cocycle before testing this condition.

The source defines its integral cup using an arbitrary
choice of lift $\widetilde c(h)\in V'$:
\[
       (g,h)\longmapsto P_0(d(g),g\widetilde c(h))
                                  \pmod O.                    \tag{BX.17}
\]
Changing the lift changes this value by an element of $O$,
by stability of $L'$ and (BX.15).  Thus (BX.17) is the
ordinary cup product in $Q/O(1)$ and is independent of
the chosen section.  It is locally constant:
$c$ has finite image by compactness, hence is killed
by some $\ell^m$; the pairing then depends only on
$d$ modulo $\ell^mL$ and the action on the finite
module $\ell^{-m}L'/L'$.  Both are locally constant.
This argument also removes any continuity concern about
the arbitrary lift function itself.

The source's \texttt{Cob\allowbreak O} says exactly that (BX.17) is
in the group of locally constant $2$-coboundaries.
By the normalized invariant
$H^2(K,Q/O(1))\simeq Q/O$, this is exactly vanishing
of the Tate pairing in (BX.4).  BX6, applied to the
actual lattices and actual form, now says that $c$
satisfies (BX.16) if and only if it pairs trivially
with every finite lattice cocycle, and that a lattice
cocycle $d$ is finite if and only if it pairs trivially
with every such $c$.  These are exactly the two
conjuncts of \texttt{Integral\allowbreak Duality}.  All variables
were arbitrary in their original order, completing
the universally quantified predicate.
\end{proof}
% END BX7

% BEGIN BX8
\begin{theorem}[BX8: the full cohomology record]
\label{claim:BX8}
The complete predicates \texttt{Local\allowbreak Published\allowbreak Theorems}
and \texttt{Cohomology\allowbreak Hyp} are supplied by the reconstructed
local and global cohomology chains, with the classical
inputs stated in those chains.
\end{theorem}
\begin{proof}
The record \texttt{Local\allowbreak Published\allowbreak Theorems} has exactly
seven fields.  LK2, LK3 and LK4 give respectively
\texttt{lf\allowbreak K}, \texttt{lt\allowbreak K}, \texttt{le\allowbreak K}, on the original
finite local field and cyclotomic action.  DX1 gives
the native fields \texttt{lf}, \texttt{lt}, \texttt{le},
including nonfree finite modules and both degree-one
cup orders.  BX7 gives the entire field \texttt{bk}.
These seven statements have the same original
quantifiers, so they assemble the full local record.

The record \texttt{Cohomology\allowbreak Hyp} then has the following
complete field-by-field assembly:
\[
\begin{array}{c|l}
\text{field}&\text{supplier}\\ \hline
\texttt{loc}&\text{the seven-field record just assembled}\\
\texttt{pt}&\text{LD3, with the exact DX1 projection}\\
\texttt{gf}&\text{DX1: original global finite cohomology}\\
\texttt{ge}&\text{DX1: original truncated global Euler identity}\\
\texttt{uo}&\text{LD2: both unramified annihilator directions}\\
\texttt{gw}&\text{DX3: the complete Greenberg--Wiles formula}.
\end{array}                                                     \tag{BX.18}
\]
The normalized local invariant family is the one fixed
in LD3.  No choice of a different invariant is used to
combine the fields.  The global Euler field includes
the source's $2$-primary case through the truncated
degree-$0,1,2$ formula; its separate high-degree
vanishing clause and Greenberg--Wiles retain exactly
their source odd-prime antecedents.  The new local
Bloch--Kato field holds also at $\ell=2$.
Thus no field of the record is left as a new
bespoke input.
\end{proof}
% END BX8

\subsection{Consequence for the reconstruction frontier}

The older DX chapter correctly avoided assuming
\texttt{Bloch\allowbreak Kato\allowbreak Local} for its explicit local charts.
BX7 now supplies that larger predicate, and BX8
supplies the containing record for other consumers
that use it.  This is an additive completion of the
previous scope.  It does not by itself construct a
promodular support, an automorphic seed, or a uniform
automorphic catalogue.  Those source inputs have
their own arithmetic content and remain separately
accounted for in the dependency audit.


\section{Actual symplectic rigidity from adjoint Selmer vanishing}
\label{sec:actual-symplectic-rigidity}

This chapter completes the irreducible-point rigidity argument on the
original arithmetic datum.  The source passes through a substantial
ordinary automorphy-lifting construction.  A more direct route is available:
Thorne's adjoint Selmer theorem requires irreducibility over the cyclotomic
tower, and the ordinary characters of these particular points prove that
condition.  We give the restriction argument, the local crystalline
calculation, and the polarization comparison explicitly.  The resulting
statement supplies the literal \texttt{Rigidity\allowbreak Symplectic\allowbreak Four} interface.
It retains the automorphy premise in that interface; automorphy is not
inferred from the existence of a deformation ring.

\subsection{The unchanged datum and the classical inputs}

Fix the original algebraic closure $\Omega/\mathbb Q$, coefficient ring
$O\subset\overline{\mathbb Q}_\ell$, and a datum
$D:\texttt{Raw}(\Omega,O)$ with $D.\texttt{Good}$.  Thus $\ell$ is odd and
$e\ge2$.  Write $G=\operatorname{Gal}(\Omega/\mathbb Q)$,
$\chi:G\to\mathbb Z_\ell^\times$ for the cyclotomic character, and
$c=D.\mathrm{cc}$ for the chosen complex conjugation.  All decomposition
and inertia groups below are the original ones in $G$.  The characters are
\[
 \alpha=\chi\omega^{-1}\xi,\qquad
 \Psi=\chi^e\psi,\qquad \vartheta=\alpha\Psi,
 \qquad (p_0,p_1,p_2,p_3)=(\vartheta,\Psi,\alpha,1).
 \tag{RX.1}
\]
AX6 supplies their exact continuity, finite-order, inertia and separation
properties.  In particular $\vartheta(c)=-1$, and the four inertia weights
on the subgroup where $\psi$ and $\omega$ vanish are
\[
                    (w_0,w_1,w_2,w_3)=(e+1,e,1,0).
 \tag{RX.2}
\]
The fixed alternating matrix and its adjoint operation are
\[
 J=\begin{pmatrix}0&0&0&1\\0&0&1&0\\0&-1&0&0\\-1&0&0&0\end{pmatrix},
 \quad X^\dagger=J^{-1}X^tJ,\quad
 \mathfrak{sp}_4=\{X:X^\dagger=-X\},\quad
 P(X)=\tfrac12(X-X^\dagger).
 \tag{RX.3}
\]

An actual point $r$ is an $O$-algebra specialization of an actual
$\texttt{Is\allowbreak Pt}$ lift over a complete Noetherian local $O$-algebra $R$ with
the same residue field.  This is precisely $D.\texttt{Pt\allowbreak In}(A,r)$ for
$A=\overline{\mathbb Q}_\ell$ or
$\overline{\mathbb Q}_\ell[\epsilon]/(\epsilon^2)$.
The local condition is the original \texttt{On\allowbreak Comp} condition, including
its domain and diagonal base point.  In particular AX8 applies to every
such characteristic-zero point.  The source's $\texttt{Irr}(r)$ says that
the matrices $r(g)$ span $M_4(\overline{\mathbb Q}_\ell)$; it implies
absolute irreducibility because every invariant subspace is then invariant
under every matrix.

We use the following named classical inputs, with the specific
normalizations required here.

\begin{itemize}
\item Fontaine's category of crystalline representations is stable under
subquotients, tensor products and duals; its period functors are exact on
admissible objects.  An unramified representation is crystalline, and a
cyclotomic power $\chi^a$ has its de Rham filtration in degree $-a$.
These are the consequences of Proposition~1.5.2 and
\S\S5.1.3--5.1.7 of \cite{RX-Fontaine}.  BX1--BX2 identify these rings,
actions, filtrations and finite cocycles with the literal source models;
BX4 supplies the Bloch--Kato dimension formula.  DX1 supplies finite local
Tate duality and the finite Euler formula; RX2 below spells out the
rational passage for the possibly higher-dimensional representation.
\item Finite abelian class-field theory over $\mathbb Q$, including the
cyclotomic description used in AX3, attaches algebraic Hecke characters
to finite-order characters times integral cyclotomic powers.  We use
geometric Frobenius throughout.  The quadratic character of an imaginary
quadratic extension is trivial after composition with its idele norm.
\item Strong multiplicity one identifies two cuspidal
$\mathrm{GL}_n$ representations whose unramified local parameters agree
outside a finite set; see \cite[\S4]{RX-JS}.
The normalized Galois representation of a regular algebraic polarized
cuspidal representation is the one in
\cite[Theorem~2.1.1]{RX-BLGGT}.  Its geometric Frobenius eigenvalues are
$q^{(n-1)/2}$ times the unitary Satake parameters.  We use Chebotarev
density and the characteristic-zero character criterion for semisimple
representations to identify two such attached systems.
\item The substantive global vanishing input is
\cite[Theorem~6.1]{RX-Thorne}: if $F$ is CM, $n\ge2$, and
$(\pi,\lambda)$ is regular algebraic, cuspidal and polarized, then
irreducibility of $r_{\pi,\iota}$ on
$G_{F(\mu_{\ell^\infty})}$ implies
$H^1_f(F^+,\operatorname{ad}r_{\pi,\iota})=0$.
Here the adjoint is the polarized, or unitary, extension to $G_{F^+}$.
It will be identified by an explicit formula below.  The theorem has no
residual irreducibility or residual image-size premise.
\end{itemize}

The standard results are leaves of this reconstruction.  The class-field
and cyclotomic facts in the second item are the already matched AX3
inputs.  The enlarged
source packages \texttt{Unramified\allowbreak Inputs} and
\texttt{Adjoint\allowbreak Lifting\allowbreak Theorem} are not assumed in the argument below and
are not claimed as consequences of it.

\subsection{Actual coefficients and ordinary finite cocycles}

% BEGIN RX1
\begin{proposition}[RX1: finite coefficients and the actual tangent cocycle]
\label{prop:rx-coefficients}
Every actual $\overline{\mathbb Q}_\ell$ point is continuous and defined
over one finite extension of $\mathbb Q_\ell$.  The same holds for both
coordinates of every actual dual-number point.  If $r'$ is such a point
and its first coordinate is $r$, write
\[
 r'(g)=r(g)+\epsilon\dot r(g),\qquad
 z(g)=\dot r(g)r(g)^{-1}.
 \tag{RX.4}
\]
Then $z$ is a continuous adjoint cocycle, and every $z(g)$ belongs to the
actual $\mathfrak{sp}_4$ in (RX.3).  Moreover $r$ is unramified away from
a finite set of rational primes.
\end{proposition}

\begin{proof}
We first prove the coefficient assertion for the algebra map defining the
point.  Let $f:R\to\overline{\mathbb Q}_\ell$ be an $O$-algebra map.
For $a\in\mathfrak m_R$ one has $|f(a)|<1$.  Otherwise the inverse of
the nonzero algebraic number $f(a)$ is integral over $\mathbb Z_\ell$.
Its monic polynomial, multiplied by a suitable power of $f(a)$, gives
\[
              1+b_1f(a)+\cdots+b_df(a)^d=0,
                   \qquad b_i\in\mathbb Z_\ell.
 \tag{RX.5}
\]
The corresponding element of $R$ lies in $1+\mathfrak m_R$ and is a
unit, which cannot map to zero.  Every element of $R$ differs from an
element of $O$ by an element of $\mathfrak m_R$, so $|f(R)|\le1$.
Choose finitely many generators of $\mathfrak m_R$ and a real number
$0<a<1$ bounding their image norms.  The ultrametric inequality gives
\[
                 |f(x)|\le a^m\quad(x\in\mathfrak m_R^m).
 \tag{RX.6}
\]
Thus $f$ is continuous.  Lifting generators of
$\mathfrak m_R/(\mathfrak m_R^2+\mathfrak m_O R)$ and successive
adic approximation gives a surjection $O[[t_1,\ldots,t_d]]\to R$.
Choose a finite extension $E/\mathbb Q_\ell$ containing
$\operatorname{Frac}(O)$ and the finitely many $f(t_i)$.
Polynomial approximations converge after applying $f$ by (RX.6).
Since $E$ is complete and is closed in $\overline{\mathbb Q}_\ell$,
they show $f(R)\subset O_E$.

Now let $F=f+\epsilon\delta:R\to
\overline{\mathbb Q}_\ell[\epsilon]/(\epsilon^2)$ be an $O$-algebra
map.  The function $\delta$ is an $O$-linear $f$-derivation.
The ideal $I=\ker f$ has finitely many generators $a_j$.
The $O$-submodule $f(R)\subset O_E$ is finite; choose generators
$f(b_i)$.  Every $r\in R$ can therefore be written
\[
             r=\sum_i o_i b_i+\sum_j u_j a_j,
                      \qquad o_i\in O,\quad u_j\in R.
\]
Consequently
\[
             \delta(r)=\sum_i o_i\delta(b_i)
                         +\sum_j f(u_j)\delta(a_j).
 \tag{RX.7}
\]
Adjoining the finitely many displayed derivative values to $E$ gives one
finite coefficient field $E'$ containing all of $F(R)$, and (RX.7)
gives a uniform bound $|\delta(R)|\le C$.
Leibniz's rule and (RX.6) then give
$|\delta(x)|\le Ca^m$ on $\mathfrak m_R^{m+1}$.
The coefficients multiplying products of generators cause no problem:
their derivative term has one additional factor bounded by $a$.
Thus $\delta$ is also adically continuous.  This proves finite
coefficient capture and continuity for the original maps themselves;
it does not use any assertion that a compact subset of
$\overline{\mathbb Q}_\ell$ lies in a finite extension.

The actual lift has continuous entries modulo every power of
$\mathfrak m_R$.  Applying the preceding estimates proves continuity
of $r$ and $\dot r$ in their finite coefficient fields.  Multiplication
of the dual-number matrices gives
\[
 z(gh)=z(g)+r(g)z(h)r(g)^{-1}.
 \tag{RX.8}
\]
The multiplier is fixed by the $O$-algebra map, so the coefficient of
$\epsilon$ in
$r'(g)^tJr'(g)=\vartheta(g)J$ is
$z(g)^tJ+Jz(g)=0$.  This is exactly $z(g)^\dagger=-z(g)$.

Finally the residual representation has open kernel and finite image
by AX6.  Its finite fixed field is ramified at finitely many primes.
Outside these primes and $D.T$, its entire inertia group lies in
$\ker\bar r$, and the original \texttt{Unr\allowbreak T} equation says that the
actual lift, hence $r$, is trivial on that inertia group.
\end{proof}
% END RX1

% BEGIN RX2
\begin{proposition}[RX2: the ordinary crystalline calculation and descent]
\label{prop:rx-ordinary-crystalline}
Let $L/\mathbb Q_\ell$ and $E/\mathbb Q_\ell$ be finite, and let
$A=E$ or $E[\epsilon]/(\epsilon^2)$.  Suppose a continuous upper
triangular $A$-representation of $G_L$ has diagonal characters
$\kappa_1,\ldots,\kappa_n$ such that
\[
 \kappa_i|_{I_L}=\chi^{w_i},\qquad
 w_1>\cdots>w_n,\qquad
 \bar\kappa_i\ne\bar\kappa_j\chi\quad(i<j).
 \tag{RX.9}
\]
Here the bar means reduction modulo $\epsilon$ when $A$ has dual
numbers, and is the identity otherwise.  Its underlying
$\mathbb Q_\ell$ representation is crystalline.
In particular a dual-number self-extension of a representation that
satisfies (RX.9) has finite adjoint cocycle.  If such a cocycle becomes
finite on an open subgroup, it is finite on the original local group.
\end{proposition}

\begin{proof}
Each rank-one $A$-character $\kappa_i\chi^{-w_i}$ is unramified.
As an $E$-representation it may be a nontrivial unramified extension
of a character by itself, but is still crystalline: its restriction
to inertia is trivial, and Fontaine's unramified descent
\cite[\S5.1.5]{RX-Fontaine} gives crystalline admissibility.
Twisting by $\chi^{w_i}$ preserves it.

Induct on $n$.  The first $n-1$ coordinates give a crystalline
$A$-module $V$ by induction.  Twisting by $\kappa_n^{-1}$ gives a
crystalline underlying $E$-representation
$W=V\otimes_A A(\kappa_n^{-1})$.
Indeed this tensor product is a quotient of
$V\otimes_E A(\kappa_n^{-1})$, so the assertion follows from
tensor stability and subquotient stability of crystalline representations.
It has an $E$-character filtration with factors
$\bar\kappa_i/\bar\kappa_n$ for $i<n$, each occurring twice when
$A=E[\epsilon]$ and once when $A=E$.
Their inertia exponents are positive.  Exactness and strictness of
$D_{\mathrm{dR}}$ therefore give
\[
 \operatorname{Fil}^0D_{\mathrm{dR}}(W)=0,
 \qquad \dim_E t_W=[L:\mathbb Q_\ell]\dim_E W.
 \tag{RX.10}
\]
The constituents of $W^\vee(1)$ are all nontrivial by (RX.9).
Successive left exactness of invariants gives $H^0(L,W^\vee(1))=0$.
We record the needed higher-dimensional rational Euler passage.
Choose a stable $\mathbb Z_\ell$-lattice $T\subset W$; it exists
because the continuous image of the compact group is bounded, so
the sum of the translates of a lattice is contained in a fixed
larger lattice.  The finite-cochain limit and finite-generation
argument of BX3 applies in every degree at issue.  It gives
$H^i(L,T)\otimes\mathbb Q_\ell=H^i(L,W)$ and finite generation
of $H^i(L,T)$ for $i=0,1,2$, with $H^3(L,T)=0$.
The multiplication-by-$\ell^m$ exact sequence gives
\[
 0\longrightarrow H^i(L,T)/\ell^m
 \longrightarrow H^i(L,T/\ell^mT)
 \longrightarrow H^{i+1}(L,T)[\ell^m]\longrightarrow0.
\]
Its outside torsion groups have bounded size independently of $m$.
Thus the base-$\ell$ logarithm of the middle cardinality is
$m\operatorname{rank}_{\mathbb Z_\ell}H^i(L,T)+O(1)$.
Divide the finite Euler identity of DX1 by $m$ and let $m$ increase.
This proves the rational Euler identity with factor
$[L:\mathbb Q_\ell]\dim_{\mathbb Q_\ell}W$.
The degree-$(2,0)$ finite Tate pairings similarly identify
$H^2(L,T)$ with the compact dual of $H^0(L,W^\vee(1)/T^\vee(1))$.
If $H^0(L,W^\vee(1))=0$, the latter invariant group embeds into
the finite torsion of $H^1(L,T^\vee(1))$, by the lattice sequence.
Hence $H^2(L,W)=0$.
After dividing all dimensions by $[E:\mathbb Q_\ell]$, these
facts and BX4 yield
\[
 \begin{split}
 H^2(L,W)&=0,\\
 \dim_E H^1(L,W)&=\dim_E H^0(L,W)
                              +[L:\mathbb Q_\ell]\dim_E W\\
                    &=\dim_E H^1_f(L,W).
 \end{split}\tag{RX.11}
\]
These are finite-dimensional spaces; hence $H^1_f(L,W)=H^1(L,W)$.
The upper-right column of the representation, divided by
$\kappa_n$, is a cocycle in $W$.  Its image in
$\operatorname{Hom}_E(A(\kappa_n),V)$ is the class of the
underlying $E$-extension.  A $B_{\mathrm{cris}}$ primitive in $W$
maps to a primitive in this Hom space.  After tensoring with
$B_{\mathrm{cris}}$ the extension therefore splits.  Since its
subrepresentation and quotient are crystalline, their comparison
isomorphisms give one for the whole extension.  This completes the
induction.  In particular it includes all dual-number unramified
diagonal variations, rather than assuming that those variations vanish.

For a crystalline self-extension
$0\to V\to V'\to V\to0$, exactness of $D_{\mathrm{cris}}$
and the comparison isomorphisms give a $G_L$-equivariant splitting
after tensoring with $B_{\mathrm{cris}}$.
One can choose the splitting over the semisimple coefficient algebra
$L_0\otimes_{\mathbb Q_\ell}E$ before extending scalars.
Writing it in a chosen vector-space splitting says precisely that
the adjoint cocycle has a $B_{\mathrm{cris}}$ primitive.  BX1--BX2
identify this with the source's \texttt{finite\allowbreak Cocycles}, including
its specified tensor order and diagonal action.

Here is the descent argument at the cochain level.  Shrinking an open
subgroup to its normal core, suppose $H\lhd G_L$ is open and
$z(h)=hx-x$ for $h\in H$, with
$x\in V\otimes_{\mathbb Q_\ell}B_{\mathrm{cris}}$.
The cocycle $F(g)=z(g)-(gx-x)$ vanishes on $H$.  The cocycle identity
shows that it factors through $Q=G_L/H$ and takes values in the
$H$-invariants.  Put $S=\sum_{q\in Q}F(q)$.  Then
\[
     gS-S=-|Q|F(g),\qquad
     z(g)=g\bigl(x-|Q|^{-1}S\bigr)
                         -\bigl(x-|Q|^{-1}S\bigr).
 \tag{RX.12}
\]
The division is valid in characteristic zero, including when
$\ell\mid|Q|$.  This proves descent without an unramified-extension
restriction.  Equivariant coefficient maps and projections carry
the displayed primitive to the corresponding primitive, so the same
argument also proves the coefficient transports used below.
\end{proof}
% END RX2

% BEGIN RX3
\begin{proposition}[RX3: the full actual \texttt{tan\allowbreak BK} conclusion]
\label{prop:rx-tanbk}
Let $D$ be good, let $r$ be an actual point, and let $r'$ be any
actual dual-number point reducing to $r$.  The cocycle in (RX.4)
satisfies, on the original $G$,
\[
 z\in Z^1(G,\operatorname{Ad}r),\qquad
 \texttt{Af.\allowbreak Is\allowbreak BK}(\operatorname{Ad}r,z),\qquad
                   z(g)\in\mathfrak{sp}_4\quad(g\in G).
 \tag{RX.13}
\]
Thus the source's \texttt{tan\allowbreak BK} output is supplied with every
original local condition, coefficient map and placement retained.
\end{proposition}

\begin{proof}
RX1 gives the first and third assertions.  Enlarge its single finite
coefficient field if needed for the finitely many changes of basis
below; the original embeddings into $\overline{\mathbb Q}_\ell$ remain
fixed.

At $\ell$, the actual \texttt{On\allowbreak Comp} condition implies the
unnormalized ordinary polynomial equations.  The four prescribed
eigenvalues $p_i(s_0)$ are distinct by AX6, so their differences are
units after passage to $E[\epsilon]$.
The spectral projectors
\[
        \prod_{j\ne i}
           \frac{r'(s_0)-p_j(s_0)}{p_i(s_0)-p_j(s_0)}
 \tag{RX.14}
\]
decompose the free module into rank-one summands.  Their partial sums
are stable by the ordinary equations.  A basis of these summands
therefore makes all $r'(g)$ upper triangular, with continuous
diagonal characters $\kappa_i$ and
$\kappa_i|_{I_\ell}=p_i|_{I_\ell}$ exactly.
Reduction gives an ordinary flag for $r$.  Its ordered characters
are those of AX8: Jordan--H\"older uniqueness identifies the sets of
characters, and evaluation at the distinct prescribed eigenvalues
of $s_0$ identifies the indices.  No permutation of the ordinary
weights has been introduced.

Put $\sigma=D.\mathrm{sg}(\ell)$.
The finite character $\tau=(\psi,\omega)$ of the original local
group has an unramified companion $\tau_{\mathrm{ur}}$ taking
$\sigma$ to $\tau(\sigma)$.  Indeed
$G_\ell/I_\ell=\widehat{\mathbb Z}$ and a chosen Frobenius is its
topological generator, so a homomorphism to the finite abelian image
is determined by this value.  Therefore
\[
 U=\ker(\tau\tau_{\mathrm{ur}}^{-1})
 \tag{RX.15}
\]
is open and normal, contains $\sigma$, and kills both $\psi$ and
$\omega$ on $U\cap I_\ell$.
The inertia characters of $r'|_U$ are consequently the exact powers
$\chi^{e+1},\chi^e,\chi,1$.
For every $i<j$, AX8, applied to this very subgroup $U$, says that
$\bar\kappa_i|_U\ne\bar\kappa_j|_U\chi$.
The presence of $\sigma$ in $U$ is essential in this use of AX8;
distinct integer weights alone would not exclude the two adjacent
weight differences equal to one.

The actual decomposition group, its henselization and completion
identify $U$ with the absolute Galois group of a finite local
extension.  The inertia and cyclotomic identifications are those
already established in DX1 and transported to the literal periods
in BX1--BX2.  RX2 makes the underlying $r'$ and $r$ crystalline on
$U$, so $z|_U$ is finite.  Formula (RX.12) then returns the finite
cocycle to the entire original decomposition group.  The period
primitive is transported through the original graph subgroup
$\texttt{dec\allowbreak J}$, giving the exact \texttt{Is\allowbreak BK.\allowbreak above} clause.

For a prime $q\ne\ell$ in $D.S_b$, the original condition kills
$I_q\cap\ker\psi$.  Away from $D.T$, it kills
$I_q\cap\ker\bar r$.  Both are finite-index normal subgroups of
inertia, and $r'$ is the identity on the relevant subgroup.
Thus $z$ vanishes there and factors, on inertia, through a finite
quotient.  Averaging this finite quotient as in (RX.12), now with
coefficients in $M_4(E)$, shows that $z|_{I_q}$ is a coboundary.
For every rational finite place one may choose the original selected
valuation if $q\in D.S_b$, and any original valuation above it
otherwise.  Conjugating a placement transports a cocycle and its
primitive by the cocycle identity.  This is precisely the
existential original-valuation clause \texttt{Is\allowbreak BK.\allowbreak away}, at every
place, and proves (RX.13).
\end{proof}
% END RX3

\subsection{Cyclotomic irreducibility and an admissible CM field}

% BEGIN RX4
\begin{proposition}[RX4: irreducibility over the required infinite tower]
\label{prop:rx-cyclotomic-irreducibility}
If $D$ is good and $r$ is an irreducible actual point, then
$r|_{\ker\chi}$ is irreducible.  There is an imaginary quadratic
field $K\subset\Omega$ such that
\[
       r|_{G_{K(\mu_{\ell^\infty})}}\text{ is irreducible}.
 \tag{RX.16}
\]
In particular $r|_{G_K}$ is irreducible.  The field $K$ is chosen
inside the original $\Omega$; it does not replace the original
arithmetic datum or its coefficient field.
\end{proposition}

\begin{proof}
We need the following elementary profinite Clifford argument.
Let $\Gamma$ act continuously and irreducibly on a finite-dimensional
space over $\overline{\mathbb Q}_\ell$, let $H\lhd\Gamma$ be closed,
and suppose $\Gamma/H$ is procyclic.  If the restriction to $H$ is
reducible, there is a nontrivial finite-order continuous character
$\eta$ of $\Gamma/H$ such that
\[
                          r\otimes\eta\simeq r.
 \tag{RX.17}
\]
Here is a proof that also covers an infinite subgroup $H$.
Choose an $H$-invariant nonzero subspace of minimal dimension.  It
is simple.  Its $\Gamma$-translates are simple $H$-submodules, and
their sum is a nonzero $\Gamma$-invariant subspace, hence the whole
space.  A sum of simple modules is semisimple: successively add a
simple summand only when it is not already contained in the sum;
its intersection with that sum is then zero.  Finite dimension
ends the procedure.  Thus the restriction has finitely many
isotypic components, permuted transitively by $\Gamma$.

Each isotypic component is a closed vector subspace.  Its stabilizer
in $\Gamma$ is closed, since the condition that a continuous matrix
map preserve a given finite-dimensional subspace is closed.  Its
index is finite, so its complement is a finite union of closed
cosets and the stabilizer is open.  The action on the components
is consequently a continuous finite quotient of the procyclic
group.  If the number of components is $d>1$, a generator cycles
them.  Multiplication by $1,\zeta_d,\ldots,\zeta_d^{d-1}$ on
successive components gives an intertwiner for a nontrivial
character of this cyclic quotient, proving (RX.17).

It remains to exclude multiplicity without a component permutation.
In that case the restriction is $W\otimes M$, with $W$ simple for
$H$ and $H$ acting trivially on the multiplicity space $M$.
Choose $t\in\Gamma$ whose image topologically generates
$\Gamma/H$.  The conjugate of $W$ by $t$ is isomorphic to $W$;
choose an invertible intertwiner $A$ for this conjugation.
Schur's lemma gives $\operatorname{End}_H(W)=
\overline{\mathbb Q}_\ell$, so the action of $t$ is $A\otimes B$
for some invertible $B$ on $M$.
An eigenline of $B$ makes $W\otimes L$ invariant under $H$ and $t$.
The stabilizer of this subspace is closed and contains $H$ and a
dense generator modulo $H$, hence is all of $\Gamma$.
Irreducibility forces $\dim M=1$.  This proves the assertion, with
no finite-group Clifford theorem silently applied to an infinite
group.

Apply it to $\Gamma=G$ and $H=\ker\chi$.
For odd $\ell$,
$\mathbb Z_\ell^\times\simeq C_{\ell-1}\times\mathbb Z_\ell$
is procyclic, and the global cyclotomic character is surjective.
Suppose that (RX.17) holds.  At the original group $G_\ell$,
AX8 gives the ordinary characters $\kappa_0,\ldots,\kappa_3$.
The two isomorphic representations have the same semisimplification,
so their constituents satisfy
\[
               \{\kappa_i\eta:0\le i<4\}
                       =\{\kappa_i:0\le i<4\}.
 \tag{RX.18}
\]
Restrict further to an open subgroup of inertia killing $\eta$,
$\psi$ and $\omega$.  The four characters there are the distinct
powers (RX.2); the image of $\chi$ is infinite on every open
inertia subgroup.  Equality of two powers therefore forces equality
of their integer exponents.  Thus the permutation in (RX.18)
fixes every index, and $\eta|_{G_\ell}=1$.

The image of $I_\ell$ under $\chi$ is all of
$\mathbb Z_\ell^\times$.  One elementary verification is that
$\Phi_{\ell^m}(1+T)$ is Eisenstein at $\ell$; hence the local
cyclotomic extension is totally ramified of degree
$\varphi(\ell^m)$.  The original global inertia therefore surjects
onto every finite cyclotomic Galois group.  Compactness and the
inverse-limit description give the claimed full image.
A character of $G/H$ trivial on $I_\ell$ is consequently trivial.
This contradicts (RX.17), proving irreducibility on $H$.

By RX1 choose an odd prime $q\ne\ell$ at which $r$ is unramified,
and take $K=\mathbb Q(\sqrt{-q})$ inside $\Omega$.
Its quadratic character $\eta_K$ is nontrivial on $I_q$, since
$q$ has odd valuation in $-q$ and ramifies in $K$.
The cyclotomic character is unramified at $q$, so $I_q\subset H$.
Hence $H\cap G_K$ has index two in $H$.
If the restriction of the irreducible $r|_H$ to this subgroup were
reducible, the same Clifford argument for this cyclic quotient of
order two would give
$r|_H\simeq r|_H\otimes\eta_K|_H$.
On $I_q$ the first representation is the identity, while the second
contains the nontrivial scalar $-1$.  This is impossible.
Finally $H\cap G_K=G_{K(\mu_{\ell^\infty})}$, which proves
(RX.16).
\end{proof}
% END RX4

\subsection{The exact automorphic and adjoint comparison}

% BEGIN RX5
\begin{proposition}[RX5: the polarized adjoint of the actual automorphic point]
\label{prop:rx-adjoint-vanishing}
Suppose $D$ is good, $r$ is an irreducible actual point, and
$\texttt{Is\allowbreak Aut4Alt}(\iota,r)$ holds with its literal source meaning
$\texttt{Is\allowbreak Aut\allowbreak QA}$.  Let $K$ be as in RX4, and let $\eta_K$ be
its quadratic character.  On $M_4(\overline{\mathbb Q}_\ell)$ define
\[
 \mathcal R(g)X
   =r(g)P(X)r(g)^{-1}
       +\eta_K(g)r(g)(X-P(X))r(g)^{-1}.
 \tag{RX.19}
\]
This is the polarized adjoint representation of the automorphic
representation supplied by \texttt{Is\allowbreak Aut4Alt}, with the same
original $G$ and $r|_{G_K}$.  Its actual Bloch--Kato Selmer group
over $\mathbb Q$ vanishes.
\end{proposition}

\begin{proof}
RX1 supplies continuity, and RX4 supplies irreducibility of
$r|_{G_K}$, so the universal imaginary-quadratic clause of
$\texttt{Is\allowbreak Aut\allowbreak QA}$ applies to this $K\subset\Omega$.
It gives a regular algebraic cuspidal representation $\pi/K$
whose \texttt{Att\allowbreak C} system is $r|_{G_K}$.
The additional \texttt{Wit\allowbreak A} assertion in the literal definition
is not needed in this argument.

We must first verify polarization with the normalization of that
system.  Since $r(g)^tJr(g)=\vartheta(g)J$, the representation
is self-dual up to $\vartheta$; conjugation by $c$ also gives an
isomorphism from $r|_{G_K}$ to its conjugate.
At an unramified place $v$ of norm $Q_v$, write the unitary
Satake parameters of $\pi_v$ as $a_{v,i}$.
The definition of \texttt{Att\allowbreak C} uses geometric Frobenius $F_v$
and gives the eigenvalues $Q_v^{3/2}a_{v,i}$.
Because $\chi(F_v)=Q_v^{-1}$, the preceding self-duality and
conjugation identities yield the exact multiset relation
\[
 \{a_{v^c,i}\}_i
       =\{\iota(\vartheta(F_v))Q_v^{-3}a_{v,i}^{-1}\}_i.
 \tag{RX.20}
\]
The character $\vartheta\chi^{-(e+1)}$ has finite order and open
kernel.  Class-field theory therefore supplies an algebraic
Hecke character $\lambda_0$ of $\mathbb Q$ whose pullback to $K$
has $\ell$-adic character
\[
                    (\vartheta\chi^3)|_{G_K}.
 \tag{RX.21}
\]
Concretely this is the finite Dirichlet character supplied by the
cyclotomic description of the finite-order part, multiplied by
the integral norm power fixed by its geometric Frobenius values.
Thus its local factor in (RX.20) is exactly
$\iota(\vartheta(F_v))Q_v^{-3}$.
Strong multiplicity one, in the form of
\cite[Theorem~4.2]{RX-JS}, gives
\[
       \pi^c\simeq\pi^\vee\otimes
                    (\lambda_0\circ N_{K/\mathbb Q}\circ\det).
 \tag{RX.22}
\]
If needed, multiply $\lambda_0$ by the quadratic Hecke character
of $K/\mathbb Q$ so that its value at $-1$ at the real place is
$(-1)^4=1$.  That character has value $-1$ there and is trivial
after the idele norm from $K$.  Hence neither (RX.21) nor (RX.22)
changes.  This is the exact sign correction in the source's
\texttt{hecke\allowbreak Char\_\allowbreak sign\_\allowbreak fixed}; no equality between a Hecke
infinity sign and an uncorrected cyclotomic exponent is assumed.
The resulting pair $(\pi,\lambda)$ is polarized in the precise
sense of \cite[\S2.1]{RX-BLGGT}.

The normalized Galois representation of this pair has the same
geometric Frobenius characteristic polynomials as $r|_{G_K}$
outside a finite set, by (RX.20)--(RX.21) and
\cite[Theorem~2.1.1]{RX-BLGGT}.
Both representations are continuous and semisimple, so Chebotarev
\cite[Theorem~8.31]{RX-Milne}
and the character criterion identify them.  To see that the use
of $\overline{\mathbb Q}_\ell$ hides no change of topology,
choose finite coefficient fields for both representations, take
their compositum, and compare their continuous traces there.
The finite-dimensional character criterion can be seen by taking the
algebra generated by the two actions in
$\operatorname{End}(V\oplus V')$.  Semisimplicity and faithfulness
make this a semisimple algebra.  After scalar extension its simple
matrix factors have linearly independent trace characters, detected
by their central idempotents.  Equality of traces therefore forces
equality of the multiplicities of every simple module, and hence
an isomorphism.  In particular RX4 proves exactly the
cyclotomic irreducibility hypothesis in Thorne's theorem for this
attached representation.

It remains to identify its extension of the adjoint to $G$.
Put
\[
                            B=Jr(c)^{-1}.
 \tag{RX.23}
\]
Since $c^2=1$ and $\vartheta(c)=-1$, the matrix $B$ is symmetric.
For $s\in G_K$, direct multiplication gives
\[
         r(s)^t B r(csc^{-1})=\vartheta(s)B.
 \tag{RX.24}
\]
Thus the polarized adjoint acts by ordinary conjugation on $G_K$
and by
\[
                           X\longmapsto-B^{-1}X^tB
 \tag{RX.25}
\]
at $c$.  These are the usual Lie-algebra formulas for the polarized
extension to $\mathcal G_4$.  Any symmetric polarization pairing
with the same multiplier on $G_K$ differs from $B$ by a scalar:
the ratio gives an endomorphism of the irreducible $r|_{G_K}$.
Thus this formula agrees with the pairing used for the normalized
automorphic representation, and is independent of its scalar
normalization.

For completeness (RX.19) gives that same representation without
choosing a basis for $B$.  Conjugation by a symplectic similitude
commutes with $\dagger$, hence with the idempotent $P$.
The two summands $\operatorname{im}P$ and $\ker P$ are invariant;
(RX.19) is their ordinary and quadratic-twisted action, so is a
representation.  On $G_K$ it is ordinary conjugation.  At $c$,
\[
 P(X)-(X-P(X))=-X^\dagger
\]
makes (RX.19) equal to (RX.25), using (RX.23).
Since $G_K$ and $c$ generate $G$, the representations agree.
Now apply \cite[Theorem~6.1]{RX-Thorne}, with $F=K$, $F^+=\mathbb Q$
and $n=4$.  It gives $H^1_f(\mathbb Q,\mathcal R)=0$.
The local condition away from $\ell$ is restriction to inertia,
and the condition at $\ell$ is a $B_{\mathrm{cris}}$ primitive.
These are exactly \texttt{Af.\allowbreak Is\allowbreak BK}, by BX1--BX2 and the original
decomposition-group comparison.  At the real place positive
cohomology vanishes by averaging, since the coefficients have
characteristic zero.  There is no additional omitted archimedean
condition.

We also justify the literal algebraic-closure coefficient scope.
The action $\mathcal R$ is defined over a finite field $E_0$ by
RX1.  For any continuous $\mathcal R$-cocycle $a$, use the countable
cofinal family of finite fields constructed in RV4, enlarged to
contain $E_0$.  The sets
$H_i=\{g:a(g)\in M_4(E_i)\}$ are closed subgroups: the cocycle
and inverse formulas preserve these matrices because all action
coefficients lie in $E_0$.  They cover the compact group $G$.
Baire's theorem makes one $H_i$ open; adjoining the finitely many
matrix entries of $a$ on representatives of $G/H_i$ then captures
all its values in one finite extension $E$.
For a local primitive with coefficients in the algebraic closure,
its finitely many coefficient entries lie in some finite $E'/E$.
An $E$-linear retraction $E'\to E$ fixing $1$ commutes with the
action and returns that primitive to $M_4(E)$; at $\ell$ tensor
the retraction with $1$ on $B_{\mathrm{cris}}$.
This is done separately at each place and does not require a
single field containing infinitely many chosen local witnesses.
Thus every cocycle in the literal \texttt{Af.\allowbreak Is\allowbreak BK} condition is
already in the usual finite-coefficient Selmer group after one
finite extension.  Thorne's vanishing gives its actual coboundary
there and hence in the original algebraic closure.
\end{proof}
% END RX5

% BEGIN RX6
\begin{theorem}[RX6: literal actual irreducible rigidity]
\label{thm:rx-literal-rigidity}
For every good original datum $D$, every embedding
$\iota:\overline{\mathbb Q}_\ell\simeq\mathbb C$, and every
irreducible actual point $r$,
\[
          \texttt{Is\allowbreak Aut4Alt}(\iota,r)
                       \quad\Longrightarrow\quad D.\texttt{Rigid\allowbreak At}(r).
 \tag{RX.26}
\]
Consequently the universally quantified source proposition
\texttt{Rigidity\allowbreak Symplectic\allowbreak Four} is supplied.
\end{theorem}

\begin{proof}
Let $r'$ be any actual dual-number point reducing entrywise to $r$.
RX3 supplies the original cocycle $z$ and all its finite local
conditions.  It takes values in $\operatorname{im}P$, and
\[
 P\circ\operatorname{Ad}_r(g)
             =\mathcal R(g)\circ P,\qquad Pz=z.
 \tag{RX.27}
\]
Apply $P$ to every cochain identity and to every local primitive.
For the prime-$\ell$ condition use $P\otimes1$ on the actual
$B_{\mathrm{cris}}$ tensor module.  Formula (RX.27) proves directly
that the unchanged function $z$ is a cocycle for $\mathcal R$ and
satisfies its actual \texttt{Is\allowbreak BK} conditions.

RX5 gives a matrix $X$ with $z(g)=\mathcal R(g)X-X$.
Apply $P$ once more and multiply on the right by $r(g)$:
\[
    \dot r(g)=r(g)P(X)-P(X)r(g)
                   =Y r(g)-r(g)Y,\qquad Y=-P(X)\in\mathfrak{sp}_4.
 \tag{RX.28}
\]
This proves exactly the existential matrix and the entrywise
identity required by $D.\texttt{Rigid\allowbreak At}(r)$, for every allowed
$r'$; no normalization or choice of its local frames was added.

The literal \texttt{Rigidity\allowbreak Symplectic\allowbreak Four} statement assumes
$\texttt{Ctx2}(D)\lor\texttt{Ctx3}(D)$, a positive level $N$
killing the three finite characters, and then the displayed
actual-point and automorphy premises.  Either context includes
$D.\texttt{Good}$.  Applying (RX.26) therefore proves all its
quantifiers, with the same original $D$, $\Omega$, $O$, $\iota$
and $r$.  The unused extra level hypotheses need not be strengthened.
\end{proof}
% END RX6

% BEGIN RX7
\begin{corollary}[RX7: installation on the same selected datum]
\label{cor:rx-same-datum}
Fix an original host $h$, a threshold $t$ with $h.\texttt{Big}(t)$,
an odd prime $\ell$, the original coefficient ring $O$, an integer
$0<e<t$, and $\psi\in h.X(O,e)$.  Let $D$ be the original datum
with $D.\texttt{Good}$ and the literal
$\texttt{Of\allowbreak Case}(D,h,e,\psi)$ in either paired-count case.
The literal \texttt{Symplectic\allowbreak Four\allowbreak Image\allowbreak Dense} premise supplies an
embedding $\iota$ and
$D.\texttt{All\allowbreak Irr}(\texttt{Is\allowbreak Aut4Alt}(\iota,\cdot))$ on that
same datum.  It follows that
\[
                          D.\texttt{All\allowbreak Irr}(D.\texttt{Rigid\allowbreak At}).
 \tag{RX.29}
\]
In the already reconstructed paired-count chain, the independent
actual-rigidity premise can therefore be removed from the cut:
the remaining automorphic requirements are the same-datum image
density supplier and the exact finite catalogue with its trace
nonvanishing guard.
\end{corollary}

\begin{proof}
Choose the one embedding supplied by image density.  For every
$r$, its actual-point and irreducibility assumptions give the
corresponding \texttt{Is\allowbreak Aut4Alt} assertion at this embedding.
RX6 proves $D.\texttt{Rigid\allowbreak At}(r)$ pointwise, which is exactly
the quantifier order in \texttt{All\allowbreak Irr}.
This argument changes neither the selected coefficient ring nor
any local subgroup.  AX10, the actual model constructor MX10,
the model comparison PX1--PX10 and SP10 can therefore consume the
result on their existing common $D$.  Their coefficient-extension
and return steps, the separate \texttt{Fin0}/\texttt{No\allowbreak End}
consequences of image density, and the catalogue's guarded trace
quantifiers retain precisely their previous scope.
No assertion is made here that image density or that catalogue
holds without its own supplier, or that either count or the full
RH/GRH argument is thereby unconditional.
\end{proof}
% END RX7

\subsection{Source boundary and bibliography}

The source chain begins at \texttt{Rig\allowbreak Sp4.\allowbreak rig\allowbreak Sp4\_\allowbreak alt}.
Its auxiliary CM-field theorem obtains adjoint vanishing from
\texttt{Unramified\allowbreak Inputs} and
\texttt{Adjoint\allowbreak Lifting\allowbreak Theorem}.
The exact declaration and its proof range are recorded in the source ledger.
RX4--RX5 replace that intermediate supplier by the stated
cyclotomic irreducibility proof and Thorne's published theorem.
RX1--RX3 reconstruct the actual \texttt{Cop\allowbreak Norm\allowbreak Lt.\allowbreak co\allowbreak Pt},
\texttt{Tb\allowbreak Tri\allowbreak Diag.\allowbreak tan\allowbreak BK}, ordinary crystalline induction and
finite-index primitive descent.  RX6 uses exactly the projection
and sign in \texttt{Rig\allowbreak Sp4.\allowbreak rg\_\allowbreak alg}.  Source anchors distinguish
the literal endpoints from the intentionally bypassed larger
interfaces.  In particular this chapter does not promote the
whole Hida automorphy-lifting, unramified tower, or modularity
packages.

\begin{thebibliography}{9}
\bibitem{RX-Fontaine}
J.-M. Fontaine, \emph{Repr\'esentations $p$-adiques semi-stables},
Ast\'erisque \textbf{223} (1994), 113--184,
Proposition~1.5.2 (printed p.~124; PDF index 12),
\S\S5.1.3--5.1.7 (printed pp.~155--158; PDF indices 43--46).
\url{https://www.numdam.org/item/AST_1994__223__113_0.pdf}.

\bibitem{RX-JS}
H. Jacquet and J. A. Shalika,
\emph{On Euler Products and the Classification of Automorphic Forms II},
Amer. J. Math. \textbf{103} (1981), 777--815, Theorem~4.2,
printed pp.~806--807 (PDF indices 30--31); its linear independence
statement implies the strong multiplicity-one form used here.
\url{https://www.math.columbia.edu/~hj/On%20Euler%20products%20II.pdf}.

\bibitem{RX-BLGGT}
T. Barnet-Lamb, T. Gee, D. Geraghty and R. Taylor,
\emph{Potential automorphy and change of weight},
Ann. of Math. \textbf{179} (2014), 501--609,
\S2.1 and Theorem~2.1.1, printed pp.~535--538
(published PDF indices 34--37).
\url{https://annals.math.princeton.edu/wp-content/uploads/annals-v179-n2-p03-p.pdf}.

\bibitem{RX-Thorne}
J. A. Thorne, \emph{On the vanishing of adjoint Bloch--Kato Selmer
groups of irreducible automorphic Galois representations},
Pure Appl. Math. Q. \textbf{18} (2022), no.~5,
Theorem~6.1; author preprint printed p.~24 (PDF index 23).
The definition of the polarized adjoint is recalled on pp.~3 and 25.
\url{https://arxiv.org/pdf/2207.04925}.
DOI: \url{https://doi.org/10.4310/PAMQ.2022.v18.n5.a5}.

\bibitem{RX-Milne}
J. S. Milne, \emph{Algebraic Number Theory}, version 3.08 (2020),
Theorem~8.31, printed p.~147 (PDF index 148).
\url{https://www.jmilne.org/math/CourseNotes/ANT.pdf}.
\end{thebibliography}


\section{Closing the uniform symplectic Frobenius catalogue}
\label{sec:CX}

\status{CX1--CX9 supply the literal \texttt{Number\allowbreak Sp4Alt}.  The
proof uses the precisely stated classical results below and the
already reconstructed polynomial coefficient modules and SP9.  It
has no retained bespoke catalogue, automorphy-counting, or
level-comparison premise.  An individual representation must still
satisfy the original \texttt{Adm\allowbreak RAlt} hypotheses; constructing its
\texttt{Is\allowbreak Aut\allowbreak QA} witness is a different obligation.}

The useful target is a catalogue of four-dimensional
representations with weights
\[
 m=(e+1,e,1,0),\qquad w=e+1,\qquad e\ge2.                 \tag{CX.1}
\]
We prove exactly the catalogue used by SP9, with its labels chosen
before the varying coefficient prime, algebraic closure, complex
embedding, and representation family.  The main simplification is
to allow the auxiliary automorphic representation to have level
$M^7$.  Its original Galois representation continues to have level
$M$.  Since a fixed power of $M$ only changes a polynomial counting
exponent, the final source predicate permits this simplification.
We do not assert the source's stronger, unchanged-level
\texttt{Level\allowbreak PNAlt} or its all-rank \texttt{Count\allowbreak QAlt}.

\subsection{Conventions and the classical theorems being used}

For a number field $F$, $G_F$ denotes its absolute Galois group in
the specified algebraic closure.  At a finite place $v$, write
$q_v$ for the residue-field cardinality.  Reciprocity sends a
uniformizer to \emph{geometric} Frobenius.  Thus the cyclotomic
character has value $q_v^{-1}$ there.  The source's
\texttt{Frob\allowbreak Eig} uses this geometric convention, whereas
\texttt{Number\allowbreak Sp4Alt} evaluates an \emph{arithmetic} Frobenius.
If $B_v(\pi)$ is the usual unramified Satake multiset for
$\mathrm{GL}_4$, put
\[
 A_v(\pi)=q_v^{3/2}B_v(\pi),\qquad
 \beta_v(\pi)=\{a^{-1}:a\in A_v(\pi)\}.                \tag{CX.2}
\]
The first multiset is the geometric Galois multiset; the second is
the arithmetic multiset.  The factor $q_v^{3/2}$ is exactly the
source's \texttt{alg\allowbreak Sat}, equivalently local Langlands for
$\pi_v\otimes|\det|^{-3/2}$.

Here are the classical inputs, in the forms actually used.  Their
primary references and exact source-interface matching are also
recorded in the component's classical-input registry.

\begin{enumerate}
\item A regular algebraic cuspidal representation of
$\mathrm{GL}_n$ has nonzero relative Lie algebra cohomology with
the irreducible algebraic coefficient of the corresponding
infinitesimal character.  Its finite part has a number-field
model, and every automorphism of $\mathbb C$ produces a cuspidal
conjugate with the permuted infinity type.  These are Clozel's
cohomological and rationality theorems, specifically Lemma~3.14
and Theorem~3.13 of \emph{Motifs et formes automorphes}.
\item Cuspidal cohomology injects into ordinary cohomology with
the same coefficient system.  We use the arithmetic case of
Borel--Labesse--Schwermer, Theorem~5.2.  The Borel--Serre
compactification gives a finite CW model for each fixed
torsion-free arithmetic quotient.  These assertions alone do not
provide a uniform counting theorem; CX1 proves the required
uniform bound from the degree of a congruence covering.
\item For a regular algebraic cuspidal polarized pair
$(\pi,\chi)$ over a CM field, and every prime $a$ and isomorphism
$\overline{\mathbb Q}_a\simeq\mathbb C$, the associated
semisimple $a$-adic representation exists.  At places away from
$a$, it satisfies full Frobenius-semisimplified local-global
compatibility and purity.  Its Hodge--Tate numbers are
$\lambda_{\tau,i}+n-i$ when the coefficient has highest weight
$\lambda_\tau$; a spherical place over $a$ gives a crystalline
representation.  This is Barnet-Lamb--Gee--Geraghty--Taylor,
\emph{Potential automorphy and change of weight},
Theorem~2.1.1.  At an arbitrary place over $a$ we use only the
\emph{semisimplified} compatibility of their
\emph{Local-global compatibility for $l=p$, II}, Theorem~A.
No Shin-regularity or monodromy equality at $a$ is assumed.
\item Every de Rham representation is potentially semistable
(Berger, \emph{Repr\'esentations $p$-adiques et \'equations
diff\'erentielles}, Theorem~0.6, equivalently Theorem~5.19).
The potentially semistable comparison functor is exact
on potentially semistable representations; a crystalline
representation gives a weakly admissible filtered Frobenius
module, with the same Hodge--Tate numbers.  We use the classical
Fontaine comparison theorem in the form recalled by
Colmez--Fontaine, Proposition~4.2 and the corollary to
Theorem~4.3.  CX3 spells out the slope inequality needed here.
\item Local Langlands for general linear groups preserves
conductors, is compatible with characters and unramified
parameters, and every local component of a cuspidal
$\mathrm{GL}_n$ representation is generic.  The generic newvector
theorem of Jacquet--Piatetski-Shapiro--Shalika, Theorem~5.1,
then gives a vector fixed by the usual $K_1(\mathfrak p^a)$ at
conductor exponent $a$.  We use only the resulting invariance
under the smaller principal congruence subgroup.  Conductor
compatibility is the $\varepsilon$-factor part of local
Langlands; Scholze's Theorem~14.1 is a primary reference.
\item We use global and local class field theory, Chebotarev,
and strong multiplicity one for cuspidal $\mathrm{GL}_n$.
CX4 writes the cyclotomic character, polarization twist and
comparison of the original representations explicitly.  None of
these inputs asserts automorphy of an arbitrary Galois
representation.
\end{enumerate}

These theorems have their usual characteristic-zero hypotheses.
In particular, purity below applies only after a polarized
regular algebraic cuspidal representation has actually been
constructed.  The sole weak automorphy witness used for that
construction is already a field of \texttt{Adm\allowbreak RAlt}.

\subsection{A uniform CM catalogue, including its arithmetic properties}

\begin{proposition}[CX1: uniform counting by fixed cells and polynomial coefficients]
\label{claim:CX1}
For each fixed imaginary quadratic field $K$ there is an integer
$c_K\ge1$ such that, for every $e\ge2$ and $L\ge1$, the number
of isomorphism classes of cuspidal representations $\pi$ of
$\mathrm{GL}_4(\mathbb A_K)$ with
\[
 \operatorname{Inf}_\tau(\pi)=\{3/2+m_i:0\le i<4\}
 \quad(\tau:K\hookrightarrow\mathbb C),\qquad
 \pi_f^{K_K(L)}\ne0                                         \tag{CX.3}
\]
is at most $c_K L^{32}(e+2)^{32}$.  Here $K_K(L)$ is the
literal principal congruence subgroup modulo the integer $L$.
The same bound holds after restricting to polarized pairs.
\end{proposition}

\begin{proof}
The required coefficient has, at both embeddings, highest weight
\[
 \lambda=(-3,-3,-e-1,-e-1).
                                                               \tag{CX.4}
\]
Indeed its contragredient has infinitesimal character
$\{3/2+m_i\}$.  More explicitly, if indices are $1,\ldots,4$,
then $\lambda_i+4-i=-m_{5-i}$; reversing these entries gives
the prescribed character of the dual coefficient.

We can use the polynomial realization already proved in
HP1--HP4.  Set $H=e+1$.  The source's polynomial highest weight
is
\[
 \nu_i=i+H-m_{5-i}=(e+2,e+2,4,4)_i,
 \qquad \lambda=\nu-(e+5)(1,1,1,1).
                                                               \tag{CX.5}
\]
The polynomial module is contained in the polynomials of total
degree $2e+12$ in sixteen matrix entries.  Counting all possible
exponent tuples, even those of smaller degree, bounds its
dimension by $(2e+13)^{16}$.  Tensoring the two embedding factors
and applying the indicated determinant twists therefore gives a
coefficient system $E_e$ of rank
\[
             d_e\le(2e+13)^{32}\le5^{32}(e+2)^{32}.
                                                               \tag{CX.6}
\]
This argument uses the actual highest weight; it does not count
all representations of a merely comparable dimension.

Consider the adelic locally symmetric space with level $K_K(3)$,
with the positive real central factor removed in the usual
cohomological convention.  It has finitely many arithmetic
components.  For $\mathrm{GL}_4$ their finiteness also follows
directly by classifying the underlying rank-four
$\mathcal O_K$-lattices: each is isomorphic to
$\mathcal O_K^3\oplus\mathfrak a$, and the ideal class group is
finite.  Passing from level one to level three only introduces a
finite index.  The level-three groups are torsion free.  To see
the last assertion, use an integral basis of each lattice to
embed its stabilizer in $\mathrm{GL}_8(\mathbb Z)$; an element
congruent to $1$ modulo $3$ cannot have finite order unless it is
$1$.  If a nonidentity element of prime order were
$1+3^sA$, with $s\ge1$ maximal, the binomial expansion of its
prime power, divided by the first nonzero power of $3$, gives a
nonzero reduction modulo $3$, a contradiction.  The same
argument for a prime-order power excludes every finite order.

Choose finite CW models from Borel--Serre for all these fixed
components, and let $b_K$ be their total number of cells.  The
level $K_K(3L)$ space is a finite cover.  Its degree is bounded by
\[
 [K_K(1):K_K(3L)]
 \le\bigl|M_4(\mathcal O_K/3L\mathcal O_K)\bigr|
 =(3L)^{32}.                                                  \tag{CX.7}
\]
Indeed reduction modulo $3L$ identifies this quotient with a
subgroup of the invertible matrices in the displayed finite
matrix ring.  Each base cell has at most this many lifts.  A
rank-$d_e$ local system contributes $d_e$ dimensions per cell to
the cellular cochain complex, independently of its attaching
maps or monodromy.  Hence the sum of all cohomology dimensions
is at most $b_K(3L)^{32}d_e$.

Every $\pi$ in (CX.3) also has $K_K(3L)$-fixed vectors.
Clozel's cohomological theorem and cuspidal-cohomology injection
put a nonzero summand for $\pi_f^{K_K(3L)}$ in the cohomology
just bounded.  Distinct cuspidal isomorphism classes give
distinct summands; equivalently their full unramified Hecke
eigensystems are distinct by strong multiplicity one.
Consequently any finite set of distinct classes has at most
$b_K3^{32}5^{32}L^{32}(e+2)^{32}$ members.  The whole set is
therefore finite with that bound: otherwise one could select a
finite subset of larger cardinality.  Enlarging the fixed
integer constant proves the proposition.  In particular the
constant is chosen before $L$, $e$, or any Frobenius prime.
\end{proof}

\begin{proposition}[CX2: the full labels and their conjugation action]
\label{claim:CX2}
Let $\mathcal S_K(e,L)$ be the finite set of isomorphism classes
in (CX.3) that admit a polarized pair.  Every
$\gamma\in\operatorname{Aut}(\mathbb C)$ permutes this set, and
at every $v\nmid L$,
\[
 A_v({}^{\gamma}\pi)=\gamma A_v(\pi),\qquad
 \beta_v({}^{\gamma}\pi)=\gamma\beta_v(\pi).
                                                               \tag{CX.8}
\]
All these multisets consist of algebraic numbers.  The labels
remain the full global isomorphism classes; classes with equal
eigenvalues at one prime are not identified.
\end{proposition}

\begin{proof}
Clozel's conjugation theorem gives ${}^{\gamma}\pi$, whose
finite part is the scalar conjugate of $\pi_f$.  Fixed vectors
for the actual compact subgroup are carried to fixed vectors by
scalar conjugation.  The infinity type is unchanged because
(CX.3) is the same at both embeddings.  To check polarization,
write
\[
                 \pi^c\simeq\pi^\vee\otimes(\chi\circ N_{K/\mathbb Q}),
 \qquad \chi_\infty(-1)=1.                                   \tag{CX.9}
\]
Here $c$ is complex conjugation of $K$, and $\chi$ is an
algebraic Hecke character of $\mathbb Q$.  Such a character is a
finite Dirichlet character times an integral norm power.
Conjugate the finite character and keep the integral power.
This produces $\chi_\gamma$ with the same sign.  Scalar
conjugation commutes with the contragredient, the permutation of
places induced by $c$, and this character twist.  Thus (CX.9)
holds at almost every place for ${}^{\gamma}\pi$ and
$\chi_\gamma$.  Strong multiplicity one makes it a global
identity.  This proves that the polarized subset is stable.
Applying $\gamma^{-1}$ shows that the action is a permutation.

There is a useful normalization check at an unramified place.
Normalize Haar measure so that $\mathrm{GL}_4(\mathcal O_{K_v})$
has volume one, and let $T_{v,j}$ be the spherical Hecke
operator for the diagonal element with $j$ uniformizers.  Its
eigenvalue satisfies
\[
 e_j(A_v(\pi))=q_v^{j(j-1)/2}T_{v,j}(\pi),\qquad 0\le j\le4.
                                                               \tag{CX.10}
\]
The exponent on the right is an integer.  The Hecke operator is
defined over $\mathbb Q$, so its eigenvalue is conjugated by
$\gamma$.  Equality of the elementary symmetric polynomials
gives the first equality in (CX.8), including the potentially
misleading square-root factor in (CX.2).  Inversion gives the
second equality.  Clozel's number-field model for $\pi_f$
makes the eigenvalues of $T_{v,j}$ algebraic, so (CX.10) makes
the entries of $A_v$ and $\beta_v$ algebraic as well.
\end{proof}

\begin{proposition}[CX3: integral arithmetic eigenvalues and their size]
\label{claim:CX3}
Let $\pi\in\mathcal S_K(e,L)$, let $p\nmid L\operatorname{disc}K$
be prime, and let $v\mid p$, with $q=q_v$.  Every
$b\in\beta_v(\pi)$ is an algebraic integer, as is $q^{e+1}b^{-1}$.
Moreover
\[
                 |b|=q^{(e+1)/2}.                            \tag{CX.11}
\]
The assertions hold for every complex conjugate with the same
$K,e,L,p,v$.  In particular this supplies the source's full
\texttt{int} and \texttt{int\allowbreak Inv} conditions for the weights
(CX.1), and supplies its weaker \texttt{Win} condition whenever
$q>1$.
\end{proposition}

\begin{proof}
Apply the polarized automorphic Galois theorem with any
auxiliary prime $a$.  From (CX.4), at every embedding its
Hodge--Tate numbers are
\[
                     0,-1,-e,-e-1.                          \tag{CX.12}
\]
The weight of its pure Weil--Deligne representation at
$v\nmid a$ is therefore $-e-1$: the theorem pairs the Hodge
numbers at the two conjugate embeddings by $h\mapsto w_0-h$,
and (CX.12) forces $w_0=-e-1$.  Its geometric Frobenius
multiset is $A_v(\pi)$ by local-global compatibility.  Purity
gives $|a_v|=q^{-(e+1)/2}$ for those eigenvalues, which proves
(CX.11) after inversion.  CX2 gives the same argument for every
complex conjugate.

We check integrality at \emph{every} finite place of the number
field generated by $b$.  First choose a prime $a\ne p$ and an
embedding of that number field into
$\overline{\mathbb Q}_a$.  Extend its inverse to the complex
coefficient isomorphism used in the automorphic Galois theorem.
Its representation can be written over a finite extension of
$\mathbb Q_a$.  For completeness, the finite-coefficient
capture needed here follows from the same Baire argument as
RV4: enumerate the finite extensions of $\mathbb Q_a$ inside
its algebraic closure; the closed subgroups on which all matrix
entries lie in one such field cover the compact profinite
domain.  One has interior and is therefore open.  Adjoining
the finitely many matrices for its cosets gives one finite
coefficient field for the whole representation.

The compact image preserves a lattice.  Indeed the sum of all
translates of one lattice is bounded above by a fixed
fractional lattice, contains the original lattice, and is a
finitely generated module over the coefficient DVR.  It is
therefore a full stable lattice.  Frobenius and its inverse
act invertibly on it.  Their eigenvalues, including the chosen
image of $b$, are $a$-adic units.

For the prime $p$ itself, use the associated $p$-adic
representation.  The local factor $\pi_v$ is spherical, so
the representation is crystalline at $v$.  Let $f$ be the
residue degree, so $q=p^f$.  The eigenvalues of $\varphi^f$ in
its crystalline module are the images of $A_v(\pi)$.  Weak
admissibility and (CX.12) put every normalized Newton slope in
$[-e-1,0]$.  Here is the needed elementary implication.  A
lowest-slope Frobenius subobject, with its induced filtration,
has $t_N\ge t_H$, and every filtration jump is at least
$-e-1$; its slope cannot be below $-e-1$.  A highest-slope
quotient has $t_N\le t_H$, by subtracting the subobject
inequality from the equality for the whole module.  Every
filtration jump of the quotient is at most zero, so its slope
cannot exceed zero.  The slope subobject or quotient can be
formed after a finite coefficient extension by grouping the
generalized eigenvectors of $\varphi^f$ with the indicated
valuation, together with their cyclic Frobenius translates.
This proves the bounds without assuming a splitting of the
original crystalline representation.  The fact that $p$ is
unramified in $K$ makes $\varphi^f$ precisely the Frobenius
with residue cardinality $q$; the same argument is applied to
each coefficient embedding.

Consequently, for a geometric eigenvalue $a_v=b^{-1}$,
\[
             -f(e+1)\le v_p(a_v)\le0.
                                                               \tag{CX.13}
\]
Thus $b$ and $q^{e+1}b^{-1}$ have nonnegative valuation at
every place over $p$, and they have nonnegative valuation at
all other finite places by the preceding unit argument.
Since they are algebraic, they are algebraic integers.
For the source's \texttt{int} condition, $-a\le m_i$ for all
$i$ forces $a\ge0$, because $\min m_i=0$.  For
\texttt{int\allowbreak Inv}, $m_i\le a$ forces $a\ge e+1$.
Multiplication by a further nonnegative integral power of $q$
therefore gives exactly both universally quantified conditions.
Finally (CX.11) gives the correct product modulus and places
every entry strictly between $q^{(e+1)/2}/\sqrt q$ and
$\sqrt q\,q^{(e+1)/2}$.  This is \texttt{Win}.
\end{proof}
\subsection{Attaching the actual representation at a polynomial level}

\begin{proposition}[CX4: polarization and comparison in the original closure]
\label{claim:CX4}
Fix an odd prime $\ell$, an original algebraic closure
$\Omega/\mathbb Q$, an isomorphism
$\iota:\overline{\mathbb Q}_\ell\simeq\mathbb C$, and a
representation $r$ satisfying the literal
\texttt{Adm\allowbreak RAlt} at $(M,m,w)$ from (CX.1).  Suppose that an
imaginary quadratic subfield $K\subset\Omega$ has irreducible
restriction $r_K$.  The \texttt{Is\allowbreak Aut\allowbreak QA} field supplies a
regular algebraic cuspidal representation $\pi/K$ weakly
attached to $r_K$.  This same $\pi$ admits a polarized pair,
and its associated $\ell$-adic representation is isomorphic to
the actual $r_K$ in the same coefficient field and closure.
Its infinity type is (CX.3), and at every place away from
$\ell$ its full Weil--Deligne parameter is the one of $r_K$.
\end{proposition}

\begin{proof}
Continuity is a field of the original \texttt{Level}.  Thus
the literal \texttt{Is\allowbreak Aut\allowbreak QA} hypothesis applies to this very
quadratic restriction, with the given $\Omega$ and $\iota$.
It provides $\pi$ and equality of geometric Frobenius
polynomials at all but finitely many places.  It also provides
the additional \texttt{Wit\allowbreak A} assertion, which we need not
strengthen or use for this argument.

Write the original multiplier as
\[
 \theta=\varepsilon_\ell^{e+1}\eta,
 \qquad \eta\text{ factors through }
       \operatorname{Gal}(\mathbb Q(\mu_N)/\mathbb Q)
 \quad\text{for some }N\ge1.                                 \tag{CX.14}
\]
This is exactly \texttt{Is\allowbreak Cyc\allowbreak Wt}; factoring follows because its
\texttt{Kills\allowbreak Mu} condition kills the fixing subgroup of all
$N$-th roots of unity.  The finite character is continuous.
Class field theory makes
$\theta\varepsilon_\ell^3$ the $\ell$-adic character of an
algebraic Hecke character $\chi$ of $\mathbb Q$.  In particular
at a good rational place, with geometric Frobenius $F_v$,
\[
             \chi(v)=\iota\theta(F_v)\,q_v^{-3}.
                                                               \tag{CX.15}
\]
One can construct it directly from the finite Dirichlet
character in (CX.14) and the integral norm power $e+4$.
This is the source's shift by $n-1=3$; using only $\theta$
would give the wrong automorphic twist.

Choose a lift in $G_{\mathbb Q}$ of complex conjugation on
$K$.  Conjugating $r_K$ by its matrix under $r$ identifies
$r_K^c$ with $r_K$.  The original nondegenerate bilinear
form in \texttt{Dual\allowbreak Tw} identifies $r_K$ with
$r_K^\vee\otimes\theta|_{G_K}$.  At a common unramified
place, put $a=q_v^{3/2}b$, where $a$ is a geometric Galois
eigenvalue and $b$ a Satake eigenvalue.  The duality relation
sends $a$ to $\theta(F_v)/a$.  After dividing by $q_v^{3/2}$
it sends $b$ to $\theta(F_v)q_v^{-3}/b$.  By (CX.15) these
are precisely the Satake parameters of
$\pi^\vee\otimes(\chi\circ N_{K/\mathbb Q})$.
Conjugation of places gives those of $\pi^c$.
Strong multiplicity one proves (CX.9).

If $\chi_\infty(-1)$ is not $+1$, multiply $\chi$ by the
quadratic character $\eta_K$ of $K/\mathbb Q$.
This changes the real sign and leaves
$\chi\circ N_{K/\mathbb Q}$ unchanged: local reciprocity
sends a norm into the subgroup fixing $K$, or equivalently the
quadratic character is trivial on these norms.  It also leaves
the restriction of the associated $\ell$-adic character to
$G_K$ unchanged.  Thus
$\varepsilon_\ell^{-3}\chi_\ell|_{G_K}=\theta|_{G_K}$
after this correction; no equality of the corrected characters
on all of $G_{\mathbb Q}$ is asserted.  We now have the literal polarized-pair sign
for rank four.  No polarization theorem about an arbitrary
nonautomorphic representation has been invoked.

Apply the polarized automorphic Galois theorem to this pair.
Its representation $R_{\ell,\iota}(\pi)$ is semisimple and
continuous.  The original $r_K$ is irreducible and continuous.
Their geometric Frobenius polynomials agree at almost every
place, by weak attachment and the theorem's normalization
(CX.2).  Both images lie in a finite coefficient extension by
the Baire argument recorded in CX3.  Chebotarev and continuity
make their traces agree on every element of $G_K$.
Equal characters of semisimple representations in
characteristic zero give an isomorphism.  For an elementary
form of this last step, take the finite-dimensional algebra
spanned by the image on their direct sum.  Its radical acts
trivially on this semisimple faithful module, so the algebra
is a product of matrix algebras.  Equality of traces on that
algebra, and then on its central idempotents, gives equal
multiplicities of every simple factor.  This proves
\[
                      r_K\simeq R_{\ell,\iota}(\pi).
                                                               \tag{CX.16}
\]

The infinity type is also forced, not chosen.  The original
flag at $\ell$ has one-dimensional quotients $\kappa_i$ with
inertial characters $\varepsilon_\ell^{m_i}\eta_i$, where
$\eta_i$ is finite.  Each quotient differs from that character
by an unramified character and has Hodge--Tate number $-m_i$.
Restriction to $K$ preserves these numbers at every embedding.
Conjugating the original place transports the same flag to
any other place needed for this statement.  By (CX.16) the
whole representation is de Rham; exactness of its de Rham
filtration now makes its Hodge--Tate multiset precisely
$(0,-1,-e,-e-1)$.  We have not inferred that an arbitrary
extension of de Rham characters is de Rham.
The automorphic Galois theorem identifies these numbers with
$\lambda_i+4-i$.  Sorting them gives exactly (CX.4), at both
embeddings, and hence (CX.3).  Finally its away-$\ell$
local-global compatibility transports through the actual
isomorphism (CX.16).  The decomposition groups, inertia
groups, and coefficient embedding remain those with which
we started.
\end{proof}

\begin{proposition}[CX5: the actual conductor bound and the level $M^7$]
\label{claim:CX5}
In CX4, the same cuspidal representation $\pi$ has a nonzero
$K_K(M^7)$-fixed vector.  If $p\nmid M$ and $\ell\ne p$, then
at every place $v\mid p$ the original $r_K$ is unramified and
its geometric Frobenius polynomial is $A_v(\pi)$.  This holds
at every original place in $\Omega$ above $v$ and every
corresponding Frobenius.
\end{proposition}

\begin{proof}
The second disjunct of \texttt{Safe\allowbreak RA} says that no two
specified weights differ by one.  It is false because
$m_1=m_2+1$ in one-based indexing.  Thus its first disjunct
holds and
\[
                              \ell\mid M.                    \tag{CX.17}
\]
This elementary observation is what permits a polynomial
level enlargement at $\ell$.  There is no exceptional
uncontrolled $\ell$-factor in the catalogue.

For a place $v$ of $K$, write $a_v=v(M)$, using the valuation
with $v(\mathfrak p_v)=1$.  Every finite character obtained by
restricting one of the original \texttt{Kills\allowbreak Mu}$\,M$
characters to $G_{K_v}$ has conductor exponent at most $a_v$.
Here is the ramification check, including ramified $K_v$.
Local reciprocity and its norm compatibility identify the
cyclotomic action of a unit with the inverse of its norm in
the relevant cyclotomic quotient.  If $a_v>0$, a unit in
$1+M\mathcal O_{K_v}$ acts trivially, since
\[
 N_{K_v/\mathbb Q_b}(1+M x)\equiv1\pmod{M\mathbb Z_b},
                                                               \tag{CX.18}
\]
where $b$ is the residue characteristic.  To prove the
congruence, express multiplication by $1+Mx$ in an integral
$\mathbb Z_b$-basis of $\mathcal O_{K_v}$; its matrix is
$1+M A$, and its determinant is $1$ modulo $M$.  The ideal
$M\mathcal O_{K_v}$ is $\mathfrak p_v^{a_v}$.
If $a_v=0$, all $M$-th roots of unity lie in the maximal
unramified extension at $b$, by the separability of
$X^M-1$ in the residue field and Hensel lifting.  The character
is then unramified.  These arguments prove the stated
conductor bound in every case.

At $v\nmid\ell$, \texttt{Level\allowbreak Away} diagonalizes the
restriction of $r_K$ to inertia into four such finite
characters.  Its inertia image is finite, so its monodromy
operator is zero, and its Artin conductor is at most $4a_v$.
Full local-global compatibility in CX4 and conductor
compatibility of local Langlands give
$a(\pi_v)\le4a_v$.  In particular, when $a_v=0$, the
Weil--Deligne parameter is unramified with zero monodromy;
$\pi_v$ is spherical.

At $v\mid\ell$, (CX.16) is de Rham and potentially
semistable.  The original local flag and exactness of the
potentially semistable functor identify the semisimplified
Weil representation with four characters.  Their inertial
parts are exactly the finite characters just bounded;
cyclotomic and unramified twists contribute no inertia to
their Weil--Deligne characters.  Semisimplified local-global
compatibility identifies this semisimple Weil representation
with that of $\pi_v\otimes|\det|^{-3/2}$.  It does not
identify the monodromy operators, and we do not need it to.
For any four-dimensional Weil--Deligne parameter $(V,\rho,N)$,
\[
 \begin{aligned}
 a(\rho,N)
  &=a(\rho)+\dim V^I-\dim(\ker N)^I,\\
  &\le a(\rho)+\operatorname{rank}N
   \le a(\rho)+3.
 \end{aligned}                                                \tag{CX.19}
\]
The first equality is the definition of its conductor; the
last inequality uses nilpotence of $N$ in dimension four.
Conductor is unchanged by an unramified twist.  Since
$a_v\ge1$ at $v\mid\ell$ by (CX.17), we obtain
\[
                    a(\pi_v)\le4a_v+3\le7a_v.                \tag{CX.20}
\]
Thus (CX.20) holds at every finite place, with zero conductor
outside $M$.

Every $\pi_v$ is generic.  The local newvector theorem gives
a nonzero vector fixed by $K_1(\mathfrak p_v^{a(\pi_v)})$.
The principal subgroup congruent to $1$ modulo
$\mathfrak p_v^{7a_v}$ is contained in that group.  At
$a_v=0$ use the spherical vector.  Their restricted tensor
product is a nonzero vector fixed by the actual adelic
principal subgroup $K_K(M^7)$.  No character conductor or
coefficient prime has been inserted into this level after
the fact.

For $p\nmid M$ and $\ell\ne p$, the preceding away-$\ell$
argument gives unramifiedness and the exact polynomial
$A_v(\pi)$.  Different lifts of geometric Frobenius differ
by inertia and hence have identical matrices under $r_K$.
Different places of $\Omega$ above $v$ are conjugate by
$G_K$, and conjugating the place conjugates the Frobenius
matrix.  This proves all of the place and Frobenius
quantifiers in the proposition.
\end{proof}

\subsection{Fixed quadratic fields and the finite ambiguity of restriction}

\begin{proposition}[CX6: eighteen fixed fields and the two-extension bound]
\label{claim:CX6}
There are eighteen imaginary quadratic fields $K_i$, fixed
once and for all, with pairwise coprime discriminants.
For every prime $p$, at least seventeen are unramified at $p$.
For every irreducible four-dimensional representation $r$
over an algebraically closed characteristic-zero field, at
least one of those seventeen quadratic restrictions is
irreducible.  Once an irreducible restriction to a fixed
$K_i$ is specified, it has at most two extensions to
$G_{\mathbb Q}$ up to isomorphism.
\end{proposition}

\begin{proof}
Choose eighteen distinct primes $d_i\equiv3\pmod4$ and put
$K_i=\mathbb Q(\sqrt{-d_i})$.  There are enough such primes
by the elementary Euclidean argument: if all had occurred
in a finite list, $4\prod_i d_i-1$ would have a new prime
factor congruent to $3$ modulo $4$.  Indeed it is $3$ modulo
$4$, some prime factor with odd multiplicity has that
congruence, and no old $d_i$ divides it.  The quadratic
integer-basis calculation gives
$\operatorname{disc}K_i=-d_i$.  Thus a prime divides at
most one discriminant.  Each remaining residue degree
$f_i$ over $p$ is either one or two.

For a representation $V$ of dimension $n$ and every distinct
self-twist character $\delta$, choose an invertible
intertwiner $T_\delta:V\simeq V\otimes\delta$.
Under conjugation on $\operatorname{End}(V)$ these are
nonzero eigenvectors for distinct one-dimensional
characters.  They are linearly independent.  One direct
proof takes a nontrivial relation of minimal length and
applies a group element distinguishing two characters;
subtracting the appropriate scalar multiple of the
original relation gives a shorter relation, a contradiction.
Consequently there are at most $n^2$ self twists, hence at
most sixteen in dimension four.  The seventeen quadratic
characters belonging to the selected fields are distinct.
One, say $\delta_i$, is not a self twist.

Write $H=G_{K_i}\triangleleft G_{\mathbb Q}$ and choose
$t\notin H$.  If a nonzero proper $H$-stable subspace $W$
existed, then $W\cap tW$ and $W+tW$ would be stable under
the whole group, since $t^2\in H$.  Irreducibility forces
$W\cap tW=0$ and $W+tW=V$.  The operator which is $+1$
on $W$ and $-1$ on $tW$ commutes with $H$ and anticommutes
with the other coset.  It is therefore an invertible
intertwiner between $r$ and $r\otimes\delta_i$, contrary
to the choice of $i$.  Thus $r|_H$ is irreducible.  This
argument requires no unproved semisimplicity assertion
about an arbitrary restriction.

Finally suppose $r|_H\simeq r'|_H$ is irreducible.  Conjugate
$r'$ so that the two restrictions agree.  The operator
$r(t)^{-1}r'(t)$ commutes with $r(H)$ and is scalar by
Schur's lemma.  An elementary proof of Schur's lemma takes
an eigenvalue of this operator: its nonzero eigenspace is
$H$-stable and must be all of $V$.  Let the scalar be $a$.
Since $t^2\in H$, it satisfies $a^2=1$.  Hence $a=1$ or
$a=-1$, and the two possible extensions are $r$ and
$r\otimes\delta_i$.  The argument works in the original
algebraic closure after any embedding of the fixed $K_i$.
\end{proof}

\subsection{The common catalogue and the literal integer endpoint}

\begin{proposition}[CX7: complete fixed-field catalogues before the coefficient prime]
\label{claim:CX7}
Fix one imaginary quadratic field $K$.  For $e\ge2$, $M\ge1$,
a prime $p\nmid M\operatorname{disc}K$ and a place $v\mid p$,
there is a finite labelled catalogue $\mathcal D_K(e,M,v)$
whose four-element multisets satisfy the literal
\texttt{Eig\allowbreak Data} conditions at $(q_v,m,e+1)$.  Its cardinality
is at most
\[
                    c_K M^{224}(e+2)^{32}.                   \tag{CX.21}
\]
For every odd $\ell\ne p$, every original $\Omega$ and
$\iota$, and every irreducible restriction to $K$ of an
original \texttt{Adm\allowbreak RAlt} representation, there is a label
with its exact arithmetic Frobenius multiset at $v$.
Equal labels give isomorphic restrictions.  The catalogue
is independent of all those later choices.
\end{proposition}

\begin{proof}
Take the label set to be $\mathcal S_K(e,M^7)$ from CX2 and
attach $\beta_v(\pi)$ to the label $[\pi]$.  CX1 gives
(CX.21); CX2 gives the permutation under every automorphism
of $\mathbb C$; CX3 gives the four-entry cardinality,
\texttt{Win}, and both universally quantified integral
conditions.  Thus this is the literal finite data structure,
even though its proof used a stronger pure absolute-value
bound.  The level appearing in its label set is explicitly
$M^7$.

Given one of the later actual restrictions, CX4 and CX5
produce a label in that already chosen set.  CX5 identifies
its geometric Frobenius multiset with $A_v(\pi)$, so (CX.2)
identifies its arithmetic multiset with the assigned
$\beta_v(\pi)$.  Equal labels give the same automorphic
Galois representation and hence isomorphic restrictions by
(CX.16).  Equivalently, equality of full global Satake
systems would suffice, by Chebotarev and semisimplicity.
We keep the full labels and therefore never replace this
argument by an inference from one Frobenius polynomial.
The label set was formed using only $K,e,M,v$; no choice
made to classify an individual restriction changes it or
the fixed constant $c_K$.
\end{proof}

\begin{proposition}[CX8: the entire SP9 catalogue with uniform quantifiers]
\label{claim:CX8}
There is an integer $c_*\ge1$ such that, for every $e\ge2$,
$M\ge1$ and prime $p\ge5$ with $p\nmid M$, the full
catalogue required before SP9 exists with
\[
                   A_*=224,
 \qquad |\mathcal D|\le c_*[M(e+2)]^{224}.                    \tag{CX.22}
\]
It is chosen before $\ell$, $\Omega$, $\iota$, the places
and Frobenius elements in $\Omega$, and the representation
family.  It classifies every original \texttt{Adm\allowbreak RAlt}
representation with weights (CX.1), has algebraic-integer
entries, and is permuted with multiplicity by every
automorphism of $\mathbb C$.  In fact every entry has
absolute value $p^{(e+1)/2}$, so its trace satisfies SP9's
bound.
\end{proposition}

\begin{proof}
Use the eighteen fields fixed in CX6.  For this $p$, let
$I_p$ consist of the indices with $p\nmid\operatorname{disc}K_i$.
Choose one place $v_i$ of $K_i$ over $p$, and let
$f_i\in\{1,2\}$ be its residue degree.  Apply CX7 to get
a catalogue $\mathcal D_i$ at $q_i=p^{f_i}$.
All these choices precede $\ell$, $\Omega$ and $\iota$.

Define $\mathcal D$ to have the following labels:
\[
 (i,s,B,\epsilon),\quad
 i\in I_p,\ s\in\mathcal D_i,\quad
 B\text{ a four-element multiset with }B^{f_i}=\beta_{i,s},
 \quad \epsilon\in\{0,1\}.                                  \tag{CX.23}
\]
The multiset attached to this label is $B$ itself.
Here $B^{f_i}$ means raising every entry to that power,
with multiplicity.  Enumerate the four entries of
$\beta_{i,s}$ for the purpose of counting.  There are at
most $f_i^4\le16$ choices of their roots and therefore at
most sixteen possible root multisets.  Hence
\[
 \begin{aligned}
 |\mathcal D|
   &\le32\sum_{i\in I_p}|\mathcal D_i|\\
   &\le32\left(\sum_{i=1}^{18}c_{K_i}\right)
                   M^{224}(e+2)^{32}
    \le c_*[M(e+2)]^{224},
 \end{aligned}                                                \tag{CX.24}
\]
where $c_*=32\sum_i c_{K_i}$ is fixed once and for all.

For an entry $b\in B$, $b^{f_i}$ is an algebraic integer,
so $b$ is integral over the algebraic integers and hence
over $\mathbb Z$.  Moreover
$(p^{e+1}b^{-1})^{f_i}=q_i^{e+1}(b^{f_i})^{-1}$ is
integral by CX3, proving the inverse integral condition
as well.  CX3 also gives
$|b|^{f_i}=q_i^{(e+1)/2}$, so
$|b|=p^{(e+1)/2}$.  An automorphism $\gamma$ permutes
the labels by
\[
        (i,s,B,\epsilon)\longmapsto
          (i,\gamma s,\gamma B,\epsilon).                    \tag{CX.25}
\]
The power equation in (CX.23) is preserved by (CX.8), and
the inverse permutation uses $\gamma^{-1}$.  This proves
the permutation with all multiplicities retained, even
when distinct labels give the same multiset.  Finally
\[
       \left|\sum B\right|
          \le4p^{(e+1)/2}
          \le(3\sqrt p+1)p^{(e+1)/2}.                         \tag{CX.26}
\]

It remains to prove the uniform classification, which is
stronger than the finiteness of (CX.23).  Fix the later
$\ell,\Omega,\iota$.  Choose a place $B_0$ of $\Omega$
above $p$.  Embed each fixed $K_i$ in $\Omega$ so that
$B_0$ restricts to the selected $v_i$.  Such a choice is
possible: at an inert prime there is only one place, and
at a split prime the two embeddings interchange the two
places.  Changing this embedding only conjugates the
restriction of a representation of $G_{\mathbb Q}$ and
does not change any of the counting bounds.

For an actual admissible $r$, CX6 supplies an index
$i\in I_p$ with irreducible restriction.  Choose the
least such index.  CX7 attaches its restriction to a
label $s\in\mathcal D_i$.  At $p$ the original $r$ is
unramified by \texttt{Level\allowbreak Away}, since $p\nmid M$ and
$p\ne\ell$.  Let $B(r)$ be the complex eigenmultiset
of an arithmetic Frobenius $\sigma$ at $B_0$.  In the
unramified quadratic extension, $\sigma^{f_i}$ fixes
$K_i$ and is an arithmetic Frobenius over $v_i$.
Therefore
\[
                     B(r)^{f_i}=\beta_{i,s}.                  \tag{CX.27}
\]
This is an identity of characteristic-polynomial roots
with multiplicity, so no diagonalizability of the
Frobenius matrix is required.

For each fixed pair $(i,s)$, the admitted restrictions
are isomorphic by CX7.  By CX6 they arise from at most
two isomorphism classes of the original representation
of $G_{\mathbb Q}$.  Choose distinct bits for those
classes.  Sending $r$ to
$(i,s,B(r),\epsilon(r))$ now gives a classification map
into the fixed set (CX.23).  Equal labels imply the same
original isomorphism class.  The choice of bits is made
only for the classification map; it adds no labels and
does not change the already fixed catalogue.

Different arithmetic Frobenius lifts at $B_0$ differ by
inertia, which acts trivially.  Every other original
place above $p$ is a $G_{\mathbb Q}$-conjugate of $B_0$,
so its arithmetic Frobenius has a conjugate matrix.
Consequently the assigned label satisfies
\[
       \iota\bigl(\operatorname{tr}r(\sigma')\bigr)
                   =\sum B(r)                               \tag{CX.28}
\]
for every place and every Frobenius in the source
quantifiers, not just the auxiliary $B_0$ used in the
construction.  We have thus proved exactly all three
SP9 catalogue conditions, in their original order of
quantification.
\end{proof}

\begin{proposition}[CX9: literal \texttt{Number\allowbreak Sp4Alt}]
\label{claim:CX9}
The source predicate \texttt{Number\allowbreak Sp4Alt} holds
relative only to the classical theorems specified above.
One may take
\[
                  A=225,\qquad C=c_*2^{225}.                  \tag{CX.29}
\]
In particular, after $e\ge2$, $M\ge1$, $p\ge5$ prime
with $p\nmid M$, $k\ge1$, and three $k$-th roots of
unity $z_1,z_2,z_3$ are fixed, there is a nonzero
integer $N_*$ with
\[
 \log|N_*|\le C\varphi(k)[M(e+1)]^{225}\log p.                 \tag{CX.30}
\]
It is chosen before every odd $\ell\ne p$, every original
$\Omega,\iota$, every place and arithmetic Frobenius
$\sigma$ above $p$, and every finite pairwise
nonisomorphic family $\mathcal X$ of the original
\texttt{Adm\allowbreak RAlt} representations.  For all of them,
every difference below is nonzero and
\[
 |N_*|_\ell\le
 \prod_{r\in\mathcal X}
 \left|\operatorname{tr}r(\sigma)-
 \iota^{-1}\bigl(p^{e+1}z_1+p^e z_2+pz_3+1\bigr)\right|_\ell.
                                                               \tag{CX.31}
\]
\end{proposition}

\begin{proof}
CX8 supplies the exact antecedent of SP9, with $A_*=224$
and the same $c_*$ for all parameters.  Apply its proved
integral-product construction.  Explicitly its factors
are indexed by the labels in (CX.23) and by
$u\in(\mathbb Z/k)^\times$:
\[
 N_* =\prod_{s,u}
 \left(\sum\beta_s-
     p^{e+1}z_1^u-p^e z_2^u-pz_3^u-1\right).                 \tag{CX.32}
\]
CX8 gives integrality of every factor and a permutation
of the label multiset under every complex automorphism;
the same automorphism acts on the three roots of unity
by one common $u$.  The product is consequently a
rational algebraic integer.  The trace separation and
the bound on its nonzero factors are exactly the
elementary inequalities proved in SP9, with (CX.26)
as their hypothesis.  Its constants give (CX.29) and
(CX.30).  For a particular family, equal labels imply
isomorphism, so each member selects a different factor
of (CX.32).  All unselected factors have $\ell$-adic
norm at most one.  Equation (CX.28) identifies the
selected factors with exactly those in (CX.31), with
the original arithmetic Frobenius and coefficient
isomorphism.  This proves the normalized valuation
inequality and all the stated universal quantifiers.
The integer itself was formed before any of those
later choices.  This is the literal source endpoint.
\end{proof}

\subsection{What has closed, and what the closure does not assert}

CX1--CX8 remove the actual uniform catalogue premise of
SP9.  The bound is uniform because the eighteen number
fields, their cell counts and their coefficient bounds
are fixed before $e,M,p$; the finite catalogue is then
fixed before $\ell,\Omega,\iota$ or a representation
family.  CX4--CX5 use the same actual restriction of
the original Galois representation.  The passage from
$M$ to $M^7$ takes place only in the auxiliary
automorphic label set.  CX6--CX8 return to the original
representation, with its original level, through a
proved two-element ambiguity and the exact Frobenius
power identity.

The full source predicate \texttt{Hypotheses'} has not been promoted.
The general record \texttt{Literature\allowbreak Rank\allowbreak N\allowbreak Hypothesis} also remains
separate.  Neither have unchanged-level
\texttt{Level\allowbreak PNAlt} and all-rank \texttt{Count\allowbreak QAlt}.
In particular, this chapter
does not construct \texttt{Is\allowbreak Aut\allowbreak QA} for a new Galois
representation and does not prove
\texttt{Symplectic\allowbreak Four\allowbreak Image\allowbreak Dense}.  It proves the
literal universal counting theorem whose inputs are
already \texttt{Adm\allowbreak RAlt}, and closes precisely the
catalogue edge needed by the paired-count chain.

\paragraph{Primary references.}
The cohomological and rationality results are in
L.~Clozel, \emph{Motifs et formes automorphes}, in
\emph{Automorphic Forms, Shimura Varieties, and
$L$-functions}, vol.~I (1990), pp.~77--159,
Theorem~3.13 and Lemma~3.14;
\href{https://www.jmilne.org/math/Books/AA1988a.pdf}
{the editor's copy} is the primary locator.
For an explicit modern statement with the same
$C$-algebraic convention, see S.~Patrikis,
\href{https://people.math.osu.edu/patrikis.1/variationsrevision.pdf}
{\emph{Variations on a theorem of Tate}},
Theorem~2.4.5; the cohomological specialization is also
spelled out by S.~W.~Shin and N.~Templier,
\href{https://math.berkeley.edu/~swshin/FieldRat.pdf}
{\emph{On fields of rationality for automorphic
representations}}, Lemma~2.12 and Proposition~2.15.
The exact injection used here is A.~Borel,
J.-P.~Labesse and J.~Schwermer,
\href{https://www.numdam.org/item/CM_1996__102_1_1_0.pdf}
{\emph{On the cuspidal cohomology of $S$-arithmetic
subgroups}}, Theorem~5.2.
The finite compactification is A.~Borel and J.-P.~Serre,
\emph{Corners and arithmetic groups},
\emph{Comment.~Math.~Helv.} 48 (1973), 436--491,
Theorem~9.3 and the manifold-with-corners construction,
\href{https://doi.org/10.1007/BF02566134}{original article}.

For the Galois, Hodge and purity assertions see
T.~Barnet-Lamb, T.~Gee, D.~Geraghty and R.~Taylor,
\href{https://www.ma.imperial.ac.uk/~tsg/Index_files/potentialautomorphy.pdf}
{\emph{Potential automorphy and change of weight}},
Theorem~2.1.1, and
\href{https://arxiv.org/pdf/1105.2242}
{\emph{Local-global compatibility for $l=p$, II}},
Theorem~A.  The finite-coefficient and local
normalizations are checked in CX3--CX5; only
semisimplification at $\ell$ is used in (CX.19).
The filtered-module fact is P.~Colmez and J.-M.~Fontaine,
\href{https://webusers.imj-prg.fr/~pierre.colmez/CF.pdf}
{\emph{Construction des repr\'esentations $p$-adiques
semi-stables}}, Proposition~4.2 and the corollary to
Theorem~4.3.
For the implication from de Rham to potentially semistable see
L.~Berger,
\href{https://arxiv.org/pdf/math/0102179}
{\emph{Repr\'esentations $p$-adiques et \'equations
diff\'erentielles}}, Theorems~0.6 and~5.19.
For conductors see P.~Scholze,
\href{https://people.mpim-bonn.mpg.de/scholze/LocalLanglands.pdf}
{\emph{The local Langlands correspondence for
$\mathrm{GL}_n$ over $p$-adic fields}}, Theorem~14.1,
and H.~Jacquet, I.~I.~Piatetski-Shapiro and J.~Shalika,
\href{https://lesesvre.perso.math.cnrs.fr/newforms-references/jpss.pdf}
{\emph{Conducteur des repr\'esentations du groupe
lin\'eaire}}, Theorem~5.1.
Local norm compatibility and cyclotomic reciprocity are
used in their classical finite form; convenient primary
author references are J.~S.~Milne,
\href{https://www.jmilne.org/math/CourseNotes/CFT.pdf}
{\emph{Class Field Theory}}, local reciprocity and
global reciprocity chapters, and
\href{https://www.jmilne.org/math/CourseNotes/ANT.pdf}
{\emph{Algebraic Number Theory}}, Theorem~8.31
(Chebotarev).  Strong multiplicity one and genericity
are the standard general-linear cuspidal theorems;
J.~W.~Cogdell's
\href{https://people.math.osu.edu/cogdell.1/fields-www.pdf}
{\emph{Lectures on $L$-functions, converse theorems,
and functoriality for $\mathrm{GL}_n$}}, Theorem~4.3
and Corollary~4.1.2, gives the precise general-linear formulations.



% Version 0.28.  The actual valued-point image-density reduction.
% Pinned source: 7cbf7a0c32fca1badf05931879a37c3f8a2b03b4.
\section{Valued points, nonsplit blocks, and the image-density reduction}
\label{sec:symplectic-image-density}

\status{IX1--IX8 prove the entire valued-point reduction to
\texttt{Fin0}, \texttt{No\allowbreak End}, and automorphy with one common
complex embedding.  All of its local and global cohomological
inputs are supplied by AX, DX, BR, FP and LD; the full
\texttt{Cohomology\allowbreak Hyp} record is unnecessary here.
IX9--IX10 also prove the characteristic-$\ell$ trace obstruction
and the characteristic-zero transport from an actual
\texttt{Pro\allowbreak Mod\allowbreak Alt} witness.  These results do not construct
that witness.  The remaining seed, modular-support and
component-reaching assertions are recorded at the end with
their original quantifiers.  Consequently this chapter does
not assert that \texttt{Symplectic\allowbreak Four\allowbreak Image\allowbreak Dense} has already
been supplied.}

\subsection{Fixed arithmetic data and the three exact incoming assertions}

Fix the original algebraic closure $\Omega/\mathbb Q$, an odd
prime $\ell$, the specified coefficient DVR
$O\subset\overline{\mathbb Q}_\ell$, and an actual good datum
$D$.  Write $k=O/\mathfrak m_O$, $E=\operatorname{Frac}O$,
$G=\operatorname{Gal}(\Omega/\mathbb Q)$, and retain every
chosen $G_p,I_p$ in $D$.  As in MX and PX, put
\[
 \theta=(\alpha\Psi,\Psi,\alpha,1),\quad
 r_0=\operatorname{diag}(\theta_i),\quad
 \lambda=(1,0,1,0),\quad \mu=(1,1,0,0),\quad
 w_{ij}=(\lambda_i-\lambda_j,\mu_i-\mu_j).
 \tag{IX.1}
\]
The cyclotomic exponents of $\theta$ on
$I_\ell\cap\ker\psi\cap\ker\omega$ are
$(e+1,e,1,0)$, with $e\ge2$.  We use exactly the alternating
form $J$ of (MX.3), and $X^*=J^{-1}X^tJ$.
Write $V_A=\langle e_0,e_2\rangle$ and
$V_B=\langle e_1,e_3\rangle$.  A block matrix preserves both
these summands.  They are isotropic and are paired perfectly.

Here a \emph{valued ambient field} is the source's
$\mathfrak a:\texttt{Amb}\,\ell$: an algebraically closed field
$K$ with a valuation ring $V$, compatible local maps
$a_{O'}:O'\to V$ for all permitted coefficient DVRs $O'$, and
the source's residue-field algebraicity condition.  A
representation, together with finitely many auxiliary matrices,
is \emph{small} if it is defined over a complete DVR $A\subset V$
with finite residue field, is continuous modulo each power of
$\mathfrak m_A$, and $K$ is algebraic over
$\operatorname{Frac}A$.  This is \texttt{Amb.\allowbreak Small}, including
all clauses of \texttt{Is\allowbreak Dvr}; a merely bounded infinite set is
not being substituted for it.

The exact \texttt{KPair} consists of $r:G\to\operatorname{GSp}_4(V)$
and $u\in\operatorname{Sp}_4(V)$ such that:
the pair is small; $r$ reduces to $a_O r_0$;
its multiplier is $a_O\vartheta$; it satisfies
$D.\texttt{Unr\allowbreak T}$ and kills $I_p\cap\ker\psi$ for $p\in S_b$;
$u^{-1}r(G_\ell)u$ is upper triangular; its diagonal characters
reduce to the prescribed $\theta_i$ on $G_\ell$ and equal
$a_O\theta_i$ on $I_\ell$.
The exact \texttt{VPt} adds a witness consisting of an actual
$\texttt{Is\allowbreak Pt}$ lift $\rho$ over some $R\in\mathcal C_O$ and a
local coefficient-compatible map $R\to V$ giving $r$; before
choosing $u$, smallness refers to $(r,1)$.

Set $\operatorname{Irr}(r)$ to mean that the $K$-span of the
matrices $r(g)$ is $M_4(K)$.  This is the source definition.
Set $\operatorname{BlockExt}(r)$ to mean that $r$ is block
triangular with its $A\leftarrow B$ block zero, its $B$-block
span equal to $M_2(K)$, and no matrix $Y\in\operatorname{Hom}(V_A,V_B)$
turns every $(1-Y)r(g)(1+Y)$ into a block diagonal matrix.
Put $\operatorname{Star}=\operatorname{Irr}\lor\operatorname{BlockExt}$.
The direction of the off-diagonal block is part of this definition.

Fix a field isomorphism
$\iota:\overline{\mathbb Q}_\ell\simeq\mathbb C$.
The three incoming assertions are exactly these:
\begin{enumerate}
\item $\operatorname{Con}(D)$: over every ambient field of
characteristic $\ell$, every actual \texttt{KPair} has
$\neg\operatorname{Star}(r_K)$.
\item $\operatorname{PurityAtPoints}(D,\iota)$: every
\texttt{KPair} with $\operatorname{Star}$, in every ambient
field identified with $\overline{\mathbb Q}_\ell$ by the
source's coefficient- and valuation-preserving
\texttt{Is\allowbreak Zero} isomorphism, satisfies \texttt{Pure\allowbreak W}.
\item $\operatorname{AutmAlt}(D,\iota)$: every such
characteristic-zero \texttt{KPair} with full matrix span
satisfies the literal \texttt{Is\allowbreak Aut4Alt}$\,\iota$.
\end{enumerate}
In particular the second assertion concerns all these
\texttt{KPair}s, including the new nonsplit pair constructed
below; it is not just a purity statement for original
\texttt{EPt}s.  The source's \texttt{Pure\allowbreak W} is the following
strict inequality, and is weaker than equality of absolute values:
there are an open $H\le G$ and $M\ne0$ such that, at every
prime $p\nmid M$, for every original place over $p$ and every
arithmetic Frobenius $\sigma\in H$ of degree $f>0$,
\[
 |\iota x|<p^f|\iota y|
 \quad\text{for all roots }x,y\text{ of }
                 \det(X-r(\sigma)),\quad y\ne0.
 \tag{IX.2}
\]
We keep \texttt{Is\allowbreak Aut4Alt} as the literal \texttt{Is\allowbreak Aut\allowbreak QA}
of the source, including its quadratic CM extensions,
regular algebraic cuspidal attachment and \texttt{Wit\allowbreak A}
clauses.  A generic assertion of potential automorphy is
not substituted for it.

The additional classical facts used below are Noetherian
completion and Hensel's lemma, the weak Nullstellensatz,
properness of the symplectic flag variety, finite extensions
of complete discrete valuation fields, and Minkowski's theorem
that $\mathbb Q$ has no nontrivial finite extension unramified
at all finite primes.  The precise forms are recorded in the
classical-input registry.  For the geometric facts see Stacks
Project tags \href{https://stacks.math.columbia.edu/tag/04GH}{04GH},
\href{https://stacks.math.columbia.edu/tag/01WC}{01WC} and
\href{https://stacks.math.columbia.edu/tag/01KF}{01KF};
for the two number-theoretic facts see Milne,
\emph{Algebraic Number Theory}, Theorems 4.9 and 7.38--7.40,
\url{https://www.jmilne.org/math/CourseNotes/ANT.pdf}.
All finite local/global duality, finiteness and Euler formulas
are the already matched DX1 and LD3 inputs on the original
groups and finite coefficient modules.

% BEGIN IX1
\begin{proposition}[IX1: one universal ring for all actual points]
\label{claim:IX1}
There is a complete Noetherian local $O$-algebra $U$ with the
original residue field, an actual universal $\texttt{Is\allowbreak Pt}$ lift
$\rho^u$, its origin $x:U\to O$, and a finite weighted
presentation
\[
       O[[z_1,\ldots,z_d]]\twoheadrightarrow U
 \tag{IX.3}
\]
with homogeneous kernel.  The variables are actual matrix-entry
differences $\rho^u(g)_{ij}-\theta_i(g)\delta_{ij}$ of weight
$w_{ij}$.  Every actual $\texttt{Is\allowbreak Pt}$ lift over every
$R\in\mathcal C_O$ is induced by its unique $O$-algebra map
$U\to R$.
The full-weight-zero ring $B=U_0$ is complete Noetherian local
with residue field $k$.  If $B$ is finite over $O$, then
$D.\texttt{Fin0}$ holds.
\end{proposition}

\begin{proof}
DX2 and FP2 give the universal continuous matrix lift on
$D.Gam$.  MX1 identifies its pullbacks with precisely the
original global lifts satisfying the required inertia-killing
equations.  It remains to check that \texttt{Is\allowbreak Pt}, including
its \texttt{On\allowbreak Comp} clauses, is a closed condition on these
actual lifts.

For each original local group, use the integral universal
ring $U_v$ of DX5--DX7 for $\operatorname{condL}$ or
$\operatorname{condP}$.  Let $x_v:U_v\to O$ be its origin,
$\phi_v:U_v\to E[[t_1,\ldots,t_{d_v}]]$ its
characteristic-zero chart, and $P_v=\ker\phi_v$.
The chart is a domain, $P_v$ is prime, and DX5 supplies
$a_v\notin\ker x_v$ with $a_vP_v=0$.
For every $f:U_v\to R$ we claim
\[
 \operatorname{OnComp}(\rho^u_v\otimes_fR,r_{0,v})
                  \quad\Longleftrightarrow\quad f(P_v)=0.
 \tag{IX.4}
\]
If $f$ kills $P_v$, the domain $U_v/P_v$, its induced origin,
its universal local representation and its map to $R$ are the
required \texttt{On\allowbreak Comp} witness.  Conversely a domain witness
$C$ gives, by local universality, a map $h:U_v\to C$ with
its composite to $O$ equal to $x_v$.  Hence $h(a_v)\ne0$.
The equality $a_vP_v=0$ and the fact that $C$ is a domain
force $h(P_v)=0$, and then $f(P_v)=0$.
This proves (IX.4) even for nonreduced target rings.

On an arbitrary lift first impose the local polynomial
equations, obtaining its quotient ring and its unique map
from $U_v$.  Impose the image of $P_v$ in that quotient.
The inverse image of this ideal gives exactly its
\texttt{On\allowbreak Comp} condition by (IX.4).  Add the multiplier,
inertia and normalization equations.  There are finitely
many local groups, and the ideal generated by all equations
is finitely generated and closed.  Its quotient is nonzero
because the actual diagonal lift $r_0$ satisfies them, by
MX1.  The quotient therefore represents the full actual
point functor, not only its characteristic-zero tangent
functor or one preselected family of points.

Conjugation by the formal diagonal torus preserves this
condition.  The \texttt{On\allowbreak Comp} clause is preserved directly:
for a domain witness $C$, replace it by
$C[[T_1,T_2]]$, conjugate by
$\operatorname{diag}((1+T_1)(1+T_2),1+T_2,1+T_1,1)$,
and send $T_i$ to the desired unit minus $1$ in the target.
The origin sends the $T_i$ to zero and still gives $r_0$.
The other clauses are polynomial conjugation identities.
Take finitely many universal entry differences spanning the
relative cotangent space.  FP1 makes them topological
generators, and the formal-torus coefficient argument of FP3
makes their presentation kernel homogeneous, with the stated
entry weights and actual origin.

The finite-monomial proof in PX1 applies to any finite set of
weights: a fixed weight has finitely many coordinatewise
minimal exponent vectors; every exponent of that weight is
one of these plus a weight-zero exponent.  The nonzero minimal
weight-zero exponents generate their monoid.  Therefore $B$
is a quotient of a formal power series ring over $O$ on
finitely many positive-order invariant monomials.  It has the
stated completeness, locality and residue field.

Finally any entry monomial of total weight zero in any actual
lift is the image of an element of $B$.  Its generated
$O$-algebra lies in the image of $B$ under the universal map.
If $B$ is finite, that image is a finite $O$-module, and any
submodule is finite because $O$ is Noetherian.  This is
\texttt{Fin0} with its full quantification over $R$ and $\rho$.
\end{proof}
% END IX1

% BEGIN IX2
\begin{proposition}[IX2: ordinary flags on every small valued point]
\label{claim:IX2}
Every actual $\operatorname{VPt}(D,\mathfrak a,r)$ admits
$u\in\operatorname{Sp}_4(V)$ making it an actual
$\operatorname{KPair}(D,\mathfrak a,r,u)$.
Every original $D.\texttt{EPt}$ is such a valued point in the
standard ambient $\overline{\mathbb Q}_\ell$ with its valuation
ring.  These assertions do not require
$\operatorname{CohomologyHyp}$ or any automorphy input.
\end{proposition}

\begin{proof}
Choose the original $\texttt{Is\allowbreak Pt}$ witness and its local
domain witness $C$ at $\ell$.  The map $C\to O$ giving the
origin makes $O\hookrightarrow C$ injective.  Over
$F=\operatorname{Frac}C$ the four prescribed eigenvalues
$\theta_i(s_0)$ are distinct by AX6.  The unnormalized
projector equations in \texttt{Ord}, after these nonzero
differences are inverted, give their ordered stable full
flag; the multiplier pairs the $i$th and $(3-i)$th lines,
so it is a symplectic flag.  Its inertia quotient characters
are exactly $\theta_i$.

Take the schematic closure of this generic flag in the
symplectic flag variety over $C$.  The latter is projective:
embed a full flag in the product of the Grassmannians and
impose the closed incidence and isotropy equations.  Its
closed generic-flag closure is proper over $C$ and has image
all of $\operatorname{Spec}C$, since that image is closed and
contains the generic point.  All stability and fixed-inertia
equations hold on this closure: on each tautological
quotient they are regular equations vanishing at its generic
point.  After the given map $C\to V$ its generic fibre has a
$K$-point, because $K$ is algebraically closed.  Properness
extends that flag over the valuation ring $V$.
All its subquotients are free direct summands; successive
choice of a primitive vector and its symplectic dual gives
an adapted symplectic basis.  Let $u$ be its matrix.

The residual labels also need justification.  For $g\in G_\ell$
the $i$th generic quotient eigenvalue $b_i(g)$ belongs to
$F$: the ordered eigenspaces of $\rho(s_0)$ are defined over
$F$.  It is integral over $C$, because it is a root of the
monic characteristic polynomial of $\rho(g)$.
The finite domain $C[b_i(g)]$ is local because $C$ is
complete henselian.  Evaluating the projector expression
at the origin sends $b_i(g)$ to $\theta_i(g)$; this is
well-defined because its denominator is a product of the
nonzero scalars $\theta_i(s_0)-\theta_j(s_0)$.
Initially this evaluation takes values in $E$.  Its values
are integral over $O$, so they belong to $O$, which is
integrally closed.  The resulting map $C[b_i(g)]\to O$
is local: the inverse image of $\mathfrak m_O$ is a
maximal ideal above $\mathfrak m_C$, and the finite
henselian domain has only one such maximal ideal.
Thus $b_i(g)-\theta_i(g)$ belongs to the maximal ideal of
$C[b_i(g)]$.  Modulo $\mathfrak m_C$ it is nilpotent.
The determinant trick for this finite $C$-module supplies
a monic polynomial for it whose nonleading coefficients
lie in $\mathfrak m_C$.
The same polynomial equation holds on the generic-flag
closure and hence on the specialized quotient line.
In the valuation ring its root is in $\mathfrak m_V$.
This proves exactly the required residual character label
on every element of $G_\ell$, even when the characteristic
is $\ell$ and the four special eigenvalues coincide.

The remaining local conditions pass through the domain
witness by their polynomial equations, including identity
on $I_p\cap\ker\psi$ for $p\in S_b$.
Smallness is preserved: start with the DVR $A$ defining $r$;
the finitely many entries of $u$ and $u^{-1}$ are algebraic
over $\operatorname{Frac}A$.  Enlarge to their finite field
extension and its valuation ring $A'$.  It is a complete
DVR with finite residue field, and $u\in\operatorname{Sp}_4(A')$.
The induced maps are local, the entrywise adic continuity
of $r$ survives, and $K$ remains algebraic over its fraction
field.  This is the complete \texttt{Small} assertion.
The standard theorem for finite separable extensions is
Milne 7.38; in equal characteristic the purely inseparable
part is obtained by adjoining $\ell$-power roots of a
uniformizer.  Since the residue field is finite and perfect,
its valuation ring is again a finite complete DVR.

For an original $\texttt{EPt}$ with map $R\to
\overline{\mathbb Q}_\ell$, the boundedness argument of PX6
puts every element of $\mathfrak m_R$ in the open unit ball
and proves continuity.  A finite formal presentation shows
that its image lies in the finite extension of $E$ generated
by the variable values.  Its ring of integers gives the
required small local map and original residual congruence.
The identity ambient isomorphism fixes every coefficient
and its valuation, so satisfies the literal \texttt{Is\allowbreak Zero}.
This also explains why the auxiliary DVR does not replace
the original datum or its coefficient maps.
\end{proof}
% END IX2

% BEGIN IX3
\begin{proposition}[IX3: an actual nonconstant characteristic-$\ell$ curve]
\label{claim:IX3}
If $D.\texttt{Fin0}$ fails, there are an ambient field of
characteristic $\ell$, an actual valued point $r$, and an
entry monomial $x$ of total weight zero such that
\[
                   x^{\ell^n}\ne x\qquad(n>0).
 \tag{IX.5}
\]
In particular the conclusion is a small actual point on the
original functor, not an arbitrary point of a larger special
fibre.
\end{proposition}

\begin{proof}
Use $U$ and $B=U_0$ from IX1.  Then $B$ is not finite over
$O$.  If $B/\mathfrak m_OB$ were Artinian, some
$\mathfrak m_B^N$ would lie in $\mathfrak m_OB$; the two
adic topologies would be equivalent, and successive
approximation from a basis modulo $\mathfrak m_O$ would
make $B$ a finite $O$-module.  Thus choose a prime
$P\supset\mathfrak m_OB$, maximal among those different
from $\mathfrak m_B$.  Such a maximal prime exists by
Noetherianity.  For $t\in\mathfrak m_B\setminus P$, the
domain $C=B/P$ has only its maximal ideal above $(t)$,
so $\mathfrak m_C^N\subset(t)$ for some $N$.
It is complete, contains $k$, and the continuous map
$k[[T]]\to C$, $T\mapsto t$, is injective: a nonzero series
is $T^a$ times a unit.  Successive approximation modulo
$t$ makes $C$ finite over $k[[T]]$.
Embed its fraction field in an algebraic closure
$K$ of $k((T))$ and take the extended valuation ring $V$.
This gives a local $y:B\to V$ with $y(t)\ne0$ and
$y(\mathfrak m_OB)=0$; its values lie in one finite extension
of $k((T))$.

To make the original ambient coefficient system, identify
the residue field of the union of the original $\ell$-adic
coefficient DVRs with $\overline{\mathbb F}_\ell\subset K$,
extending the specified embedding of $k$.  Reduction gives
all the maps $a_{O'}$.  They are local and compatible, and
every residue element comes from a finite residue extension,
hence from an unramified coefficient extension.  Therefore
the full \texttt{Amb} condition holds, not just the choice
of $K$ and $V$.

We now lift the point of the invariant ring.  Write
$U^{\mathrm{gr}}=\bigoplus_wU_w$.  Projection to weight zero
is $B$-linear and sends $1$ to $1$, so the fibre
$U^{\mathrm{gr}}\otimes_{B,y}K$ is nonzero.  The
finite-monomial argument of PX1 makes it a finite-type
$K$-algebra: it is generated by the finitely many $z_i$.
The weak Nullstellensatz gives a map
$\xi:U^{\mathrm{gr}}\to K$ extending $y$.
For its nonzero variable values put $a_i=v(\xi(z_i))$.
Whenever $n_i\ge0$ are integers, not all zero, and
$\sum n_iw_i=0$, locality of $y$ gives
\[
                         \sum n_i a_i>0.
 \tag{IX.6}
\]
The strict linear inequalities
$a_i+\langle h,w_i\rangle>0$ consequently have a rational
solution $h$.  Indeed otherwise separation from the positive
orthant gives nonzero rational $q_i\ge0$ with
$\sum q_iw_i=0$ and $\sum q_i a_i\le0$; clearing denominators
contradicts (IX.6).  This is the same finite rational
separation argument proved in PX5.
Scale by a torus element with valuation vector $h$.
It fixes $B$ and makes every nonzero variable value have
positive valuation.

These finitely many values, together with the finite field
of definition of $y(B)$, lie in one finite extension of
$k((T))$.  Let $A'$ be its complete valuation ring.
Power series can now be evaluated there.  This evaluation
agrees with $y$ on the full weight-zero formal ring: both
maps are continuous and agree on its finite monomial
generators.  For another weight, the finite-monomial
decomposition gives agreement with the scaled $\xi$.
Thus evaluation kills each homogeneous part of a defining
relation of $U$.  To kill the whole relation, sum the
finitely many complete weight components which have a
monomial of degree less than $N$; the remaining series has
order at least $N$.  Continuity and $N\to\infty$ finish the
argument.  We have obtained an actual local map $U\to A'$,
preserving $y$ on $B$, and hence a small actual valued point.

The maximal ideal of $B$ is generated by
$\mathfrak m_O$ and finitely many positive-order invariant
monomials in the entry differences.  Since $y(t)\ne0$,
at least one of these monomials has nonzero image in
$\mathfrak m_V$.  Every such invariant monomial is an
$O$-linear combination of entry monomials of total weight
zero: subtracting the origin affects only diagonal entries,
whose weights are zero.  If all the latter values were
fixed by some positive power of Frobenius, their sums and
products and the finite coefficient field would have the
same property after taking a common multiple of the powers.
But a nonzero element of $\mathfrak m_V$ cannot satisfy
$b^{\ell^n}=b$: the factor $1-b^{\ell^n-1}$ is a unit.
This contradiction proves (IX.5).
\end{proof}
% END IX3

% BEGIN IX4
\begin{proposition}[IX4: all reducible patterns and their constant invariants]
\label{claim:IX4}
For any actual small valued point, in either characteristic,
the alternatives of (PX.14) hold: full matrix span,
\texttt{Diag\allowbreak Type}, or the two $2$-element blocks with all
cross-block entries zero in one direction and both $02,20$
entries present.  In the third case there is a block diagonal
\texttt{KPair} whose $A$-block span is $M_2(K)$.
In characteristic $\ell$, every weight-zero entry monomial
of a \texttt{Diag\allowbreak Type} point satisfies
$x^{\ell^n}=x$ for some $n>0$.
\end{proposition}

\begin{proof}
The polynomial projector proof of PX6 takes place over the
actual complete coefficient ring before specialization.
It supplies all $e_{ii}$ in the $V$-span of $r(G)$, and
remains valid in characteristic $\ell$.  Consequently the
relation $i\rightsquigarrow j$ defined by a nonzero entry
is reflexive and transitive, is reversed by
$(i,j)\mapsto(3-j,3-i)$, and has full matrix span exactly
when all pairs occur.

We spell out the arithmetic exclusion in the new valued
setting.  A two-element strongly connected component
$\{s,u\}$ crossing the $\lambda$ blocks gives a
two-dimensional subquotient and an ordinary stable local
line by IX2.  In one of its coordinate charts the slope
$v$ lies in $V$, interchanging $s,u$ if necessary.
In the DVR defining smallness, choose a least-valuation
nonzero coefficient $b=r(g_1)_{su}$ and write
$r(g)_{su}=b c(g)$ with $c(g)\in V$.
The reductions of the $c(g)$ are locally constant and
not all zero: continuity modulo the valuation of $b$
plus one proves the first assertion.  The component
matrix identity gives the residual cocycle
\[
 z(g)=\frac{\overline{c(g)}}{\bar\theta_u(g)},
 \qquad \gamma=\bar\theta_s\bar\theta_u^{-1}.
 \tag{IX.7}
\]
It cannot be a coboundary.  A coboundary would say
$\overline{c(g)}=a(\bar\theta_s(g)-\bar\theta_u(g))$.
Apply the integral linear combination of $r(g)$ giving
$e_{ss}$: its off-diagonal entry says $0=a$, whereas some
$\overline{c(g)}$ is nonzero.

Outside $T$ the original unramified condition kills the
cocycle on the inertia kernel specified by \texttt{XVan}.
At $p\in S_b$, AX6 supplies $\sigma_p$ with
$\bar\chi_\ell(\sigma_p)=1$ and
$\gamma(\sigma_p)\ne1$.  Hence both the invariants of
$k(\gamma)$ and of its cyclotomic dual vanish; DX1 local
duality and the Euler formula for $p\ne\ell$ give $H^1=0$.
The same holds after any finite extension of the residue
field.  At $\ell$, use the coordinate chart in which the
stable line is generated by $(v,1)$, and write the
two-dimensional matrix as
$\bigl(\begin{smallmatrix}a_g&bc(g)\\ \beta_g&d_g\end{smallmatrix}\bigr)$.
Stability says
$bc(g)=(d_g-a_g+\beta_gv)v$.
At $g=\sigma_\ell$ the factor in parentheses is a unit:
$d_g-a_g$ is a unit by AX6, whereas
$\beta_g\in\mathfrak m_V$ and $v\in V$.
It follows first that $v/b\in V$.
Now divide every stable-line equation by $b$.
The extra term is $\beta_gv(v/b)$, whose reduction is
zero because the opposite off-diagonal entry $\beta_g$
is in $\mathfrak m_V$.  Hence
$\overline{c(g)}=(\bar\theta_u(g)-\bar\theta_s(g))
\overline{v/b}$ for $g\in G_\ell$.
Eliminating $\overline{v/b}$ gives exactly
\[
 \overline{c(g)}(\bar\theta_u(\sigma_\ell)
                         -\bar\theta_s(\sigma_\ell))
 =\overline{c(\sigma_\ell)}(\bar\theta_u(g)
                         -\bar\theta_s(g)).
 \tag{IX.8}
\]
The denominator on the left is nonzero by AX6, so the local
cocycle is a coboundary also at $\ell$.
The residue field of an ambient is algebraic over
$\mathbb F_\ell$; a locally constant cocycle on the compact
group has finite image.  Its values therefore lie in one
finite extension of $k$.  Applying a $k$-basis coordinate
to this cocycle and all its local primitives reduces it to
the exact residual vanishing in $D.\texttt{Good}$, supplied
by BR7.  Recombining coordinates makes (IX.7) a coboundary,
the desired contradiction.

The remaining finite graph argument is elementary.
Strong components are permuted by $i\mapsto3-i$.
A component of size three would meet its reversed image,
so would be reversal-stable, impossible for this
fixed-point-free involution.  A component of size four
gives full span.  A component of size two cannot cross
$\lambda$ by the preceding paragraph, so the two components
are $\{0,2\}$ and $\{1,3\}$.  Their quotient graph cannot
have edges in both directions.  If there is no component
of size two, there is no directed cycle between distinct
vertices, which is \texttt{Diag\allowbreak Type}.

In the two-block case erase all cross-block entries.
Block projection is multiplicative on matrices triangular
in the same block direction and commutes with the
symplectic adjoint.  It therefore gives a representation
with the same multiplier, residual representation and
inertia-killing properties.  For the ordinary flag use
the following explicit specialization.  Capture the
original $r,u$ in its small DVR $A$.  Multiply its only
possible cross block by an indeterminate $t$.  The
resulting $r_t$ is a representation over $A[t]$ with the
same multiplier, and over $\operatorname{Frac}(A)(t)$
is conjugate to $r$ by a one-parameter diagonal torus.
Its ordered local diagonal characters are the original
elements of $A$, independent of $t$.
Take the closure of this generic flag in the projective
symplectic flag variety over $A[t]$, including the closed
equations for those fixed quotient characters.  Its
proper image contains the generic point and hence all
of $\operatorname{Spec}A[t]$.  At $t=0$ there is a flag
over $K$ with those same ordered characters, and
properness extends it over $V$.
Choose an integral adapted symplectic basis and capture
its finitely many entries in a finite DVR extension as
in IX2.  This proves every \texttt{KPair} clause for the
block projection.  The existing $e_{00},e_{22}$ and both
nonzero $02,20$ entries give all four $A$-block matrix
units, so that block has full span.

Finally suppose the characteristic is $\ell$ and the
pattern is \texttt{Diag\allowbreak Type}.  Each diagonal entry
$d_i(g)=r(g)_{ii}$ is a character, since no off-diagonal
two-step return contributes to a diagonal product.
Put $U_1=\ker\bar r\cap\ker\psi$ and let $m=[G:U_1]$.
For inertia at $\ell$, $r(s^m)$ is conjugate to an upper
unipotent matrix; its fourth nilpotent power is zero, so
$r(s^m)^{\ell^2}=1$.  At $p\in S_b$ it is already $1$,
and away from $T$ the \texttt{Unr\allowbreak T} condition gives the
same conclusion.  Conjugacy of places gives
$d_i(s)^{m\ell^2}=1$ on every inertia group.
Reduce the powered character $d_i^{m\ell^2}$ modulo each power of the small
DVR's maximal ideal.  Its finite continuous image then
defines an extension of $\mathbb Q$ unramified at every
finite prime, hence trivial by Minkowski's theorem.
Separatedness gives $d_i(g)^{m\ell^2}=1$ for all $g$.

There is an integral linear functional $L$ on the weight
lattice which is positive on every present nonzero
off-diagonal weight.  One explicit choice is obtained
from $a_i=|\{j:i\rightsquigarrow j\}|$,
$h_i=a_i-a_{3-i}$ and
$L(x,y)=(h_0-h_1)x+(h_0-h_2)y$.
For a present strict edge its value is $h_i-h_j>0$.
A nonzero monomial of total weight zero can therefore
use only diagonal entries.  Their values have finite
order, so in characteristic $\ell$ they belong to a
finite field.  The monomial, or zero if one factor
vanishes, is fixed by a positive Frobenius power.
\end{proof}
% END IX4

% BEGIN IX5
\begin{theorem}[IX5: finite Selmer sizes produce a nonzero integral cocycle]
\label{claim:IX5}
Let $A$ be a complete DVR with finite residue field of
characteristic $\ell\ne2$, uniformizer $\pi$, and let $T_A$
be a finite free $A$-module with an adically continuous
action $\rho$ of the original $G$.
Fix a finite prime set $S$ containing $\ell$, an original
complex conjugation $c$, the original decomposition group
$H=G_\ell$, and $s\in H$ such that
\[
 \rho(c)=-1,\qquad \chi_\ell(s)=1,\qquad
                       \rho(s)-1\text{ is an automorphism of }T_A.
 \tag{IX.9}
\]
Let $P\ne0$ be a saturated $A$-submodule of $T_A$ stable
under $H$.  Let $K_0\le\ker\rho$ contain inertia outside
$S$, and put
\[
 N=\{g\in K_0:g\text{ belongs to an inertia group away from }\ell\}.
 \tag{IX.10}
\]
Then there are an adically continuous nonzero cocycle
$z:G\to T_A$ and $y\in T_A$ satisfying
\[
 \begin{gathered}
 z(gh)=z(g)+\rho(g)z(h),\qquad z(c)=0,
                  \qquad z(N)=0,\\
 z(h)-(\rho(h)y-y)\in P\qquad(h\in H).
 \end{gathered}
 \tag{IX.11}
\]
This holds in mixed and equal characteristic.  Its only
arithmetic input is the exact finite duality, finiteness
and Euler package already installed in DX1 and LD3.
\end{theorem}

\begin{proof}
Put $M_n=T_A/\pi^nT_A$ and
$P_n=(P+\pi^nT_A)/\pi^nT_A$.  Saturation over a DVR makes
$T_A/P$ free, so
$|P_n|=q^{n d}$, where $q=|A/\pi|$ and $d=\operatorname{rank}P>0$.
The finite module $M_n$ is killed by $\ell^n$, since
$\ell\in\mathfrak m_A$; it is therefore a legitimate
finite module for the installed $\mathbb Z/\ell^{n+1}$
cohomology interfaces, even in equal characteristic.
No freeness over that coefficient ring is asserted.

On $G_\ell$ choose the local condition
\[
 L_\ell=\operatorname{im}
                  \bigl(H^1(G_\ell,P_n)\to H^1(G_\ell,M_n)\bigr).
 \tag{IX.12}
\]
The map in (IX.12) is injective.  Indeed $\rho(s)-1$
is bijective on $M_n$ and on its stable finite subgroup
$P_n$; if a $P_n$-valued cocycle is the coboundary of
$m\in M_n$, evaluating at $s$ puts $m$ in $P_n$.
The local Euler formula at $\mathbb Q_\ell$ gives
\[
 |H^1(G_\ell,P_n)|
       =|P_n|\,|H^0(G_\ell,P_n)|\,|H^2(G_\ell,P_n)|
       \ \ge\ |P_n|.
 \tag{IX.13}
\]
At every other prime in $S$ use the unramified local
condition.  Its cardinality is that of $H^0(G_p,M_n)$:
both are the kernel or cokernel of Frobenius minus $1$
on the same finite inertia-invariant group.

The Greenberg--Wiles cardinality identity obtained from
the normalized finite Poitou--Tate pairing now gives
\[
                         |\operatorname{Sel}(M_n)|\ge |L_\ell|.
 \tag{IX.14}
\]
Here are all the cancellations in that formula.
The global invariants and real-place invariants of $M_n$
are zero because $c$ acts by $-1$ and $2$ is a unit.
The local invariants at $\ell$ are zero by $s-1$.
The global invariants of the cyclotomic dual are also
zero: a fixed functional evaluated on
$(\rho(s)-1)m$ is zero, since $\chi_\ell(s)=1$,
and these elements exhaust $M_n$.
At the other finite places the unramified factors
$|L_p|/|H^0(G_p,M_n)|$ equal $1$.
The only factor left besides $|L_\ell|$ is the positive
cardinality of the dual Selmer group.  All groups in
this argument are finite by DX1, with the same local
decomposition-group transport as LD3 and KB6.

Each Selmer class has a unique representative with
$z(c)=0$: subtract a coboundary, using
$\rho(c)-1=-2$.  This representative vanishes on $N$.
At $p\ne\ell$ a locally unramified class is a
coboundary on inertia, and a coboundary vanishes on
inertia intersected with $K_0\le\ker\rho$.
Outside $S$ it is already unramified.
Condition (IX.12) says exactly that some $y_n\in M_n$
makes its restriction differ from $\delta y_n$ by
$P_n$-valued elements.  Consequently the finite set
\[
 Z_n=\{z:G\to M_n:\ z\text{ satisfies the finite version of
                                      (IX.11) for some }y_n\}
 \tag{IX.15}
\]
satisfies $|Z_n|\ge q^{nd}$.  Its finiteness follows
also directly from its injection into the finite global
$H^1$: normalized coboundaries are zero.

We must produce compatible nonzero elements; mere size
growth would not suffice.  Call $z\in Z_n$ primitive
if its reduction in $M_1$ is nonzero.  Suppose for a
fixed $n\ge1$ that no element of $Z_n$ is primitive.
The same holds at every higher level by reduction.
For a nonprimitive $z\in Z_{m+1}$, divide its values
by $\pi$ to obtain a unique $z'\colon G\to M_m$.
This is a continuous cocycle, with the same zero
conditions.  To check its local condition, let $y$
be a witness for $z$.  At $s$, reduction modulo
$P+\pi T_A$ gives $(\rho(s)-1)y=0$.
The endomorphism $\rho(s)-1$ preserves $P$ and is
bijective on $P$ and $T_A/P$: its matrix is block
triangular for a splitting of the saturated submodule,
and its determinant is a unit.  Thus
$y=p+\pi y'$ with $p\in P$.
Replace $y$ by $\pi y'$ and divide the local equation
by $\pi$; saturation of $P$ gives the required condition
for $z'$ modulo $\pi^m$.
This division map $Z_{m+1}\to Z_m$ is injective whenever
all its elements are nonprimitive.
The resulting bounded sequence of cardinalities for
$m\ge n$ contradicts $|Z_m|\ge q^{md}$.
Hence each $Z_n$ contains a primitive element.

For each level take all pairs $(z_n,y_n)$ with $z_n$
primitive and $y_n$ a local witness.  They form nonempty
finite sets; reduction maps preserve both the equations
and primitivity for levels at least one.
The finite-branching tree of such pairs has paths of
every length, so has an infinite path: successively
choose a vertex with arbitrarily long continuations.
Completeness of the finite free module $T_A$ then gives
$z(g)=\lim z_n(g)$ and $y=\lim y_n$.
All the equations pass to the limit.
The submodule $P$ is closed since $T_A/P$ is finite
free and separated.  Adic continuity follows at each
finite level, and primitivity at level one gives
$z\ne0$.  This proves (IX.11) with no inverse-limit
surjectivity assumption.
\end{proof}
% END IX5

% BEGIN IX6
\begin{theorem}[IX6: the actual nonsplit block constructor]
\label{claim:IX6}
Let $(r,u)$ be any actual \texttt{KPair} for the original
good datum.  Suppose $r$ is block diagonal and its
$A$-block span over $K$ is $M_2(K)$.  Then there is another
actual \texttt{KPair} $(r',u')$ in the same ambient field
such that
\[
 \operatorname{BlockExt}(r'_K),\qquad
          \det(X-r'(g))=\det(X-r(g))\quad(g\in G).
 \tag{IX.16}
\]
No assertion that $r'$ is an original \texttt{EPt} is
needed or made.
\end{theorem}

\begin{proof}
Capture $r,u$ in their defining complete DVR $A$.
All identities hold there by injectivity into $V$.
Let
\[
 \mathcal U=\operatorname{Hom}(V_A,V_B)
                         \cap\mathfrak {sp}_4(A)
   =A(E_{10}-E_{32})\oplus AE_{12}\oplus AE_{30}.
 \tag{IX.17}
\]
It is free of rank three, its elements multiply to zero,
and it is stable under $\rho(g)Y=r(g)Yr(g)^{-1}$.
Put
\[
             P=\{Y\in\mathcal U:u^{-1}Yu
                                      \text{ is strictly upper triangular}\}.
 \tag{IX.18}
\]
It is saturated because its defining entries are linear
equations over a domain, and is stable under $G_\ell$
because upper triangular matrices preserve strict upper
triangularity by conjugation.

We check that $P\ne0$, rather than assuming it from a
dimension count.  Write $X=r(\sigma_\ell)$ and
$U=u^{-1}Xu$.  Residual cross-block separation makes
the $B$-block of $X-U_{00}$ invertible, so the first
column of $u$ is $v=(v_0,0,v_2,0)$.
Likewise $X|_{V_A}-U_{11}$ and $U_{00}-U_{11}$ are
invertible.  The first two column equations then give
the $A$-part of the second column as $a v$, for some
$a\in A$.  Put $w=u_1-av=(0,w_1,0,w_3)$.
The columns are independent, so $v,w\ne0$.
Isotropy gives $v_0w_3=v_2w_1$.
For $\varphi=(v_2,0,-v_0,0)$ the nonzero outer product
$Y=w\varphi$ lies in (IX.17).  It kills the first two
flag steps and has image in the second; hence
$u^{-1}Yu$ is strictly upper triangular.  This proves
$P\ne0$ over the original $A$.

Modulo $\mathfrak m_A$, the operator
$\rho(\sigma_\ell)-1$ on (IX.17) has the nonzero
diagonal eigenvalues
$\bar\theta_i(\sigma_\ell)/\bar\theta_j(\sigma_\ell)-1$
for the indicated cross-block positions.
Its determinant is therefore a unit, so it is an
automorphism over $A$.  AX6 also gives
$\chi_\ell(\sigma_\ell)=1$.
Complex conjugation acts on the two blocks as $-1,+1$:
each block of the involution $r(c)$ is congruent to that
scalar, and the other factor in $(r(c)-1)(r(c)+1)=0$
is a unit.  Thus
\[
                  r(c)=\operatorname{diag}(-1,1,-1,1),
                   \qquad \rho(c)|_{\mathcal U}=-1.
 \tag{IX.19}
\]
The original \texttt{Unr\allowbreak T} condition and AX6 make $r$
unramified outside $D.S$.  Its kernel lies in $\ker\rho$.
Every hypothesis of IX5 is now verified for $T_A=\mathcal U$,
$P$ as in (IX.18), and $K_0=\ker r$.

Let $z,y$ be the resulting nonzero cocycle and primitive.
Multiply both by $\pi$; they still satisfy (IX.11),
$z$ remains nonzero, and all entries of $z(g)$ lie in
$\mathfrak m_A$.  Define
\[
                       r'(g)=(1+z(g))r(g),
                       \qquad u'=(1-y)u.
 \tag{IX.20}
\]
Because products of two members of $\mathcal U$ vanish,
the cocycle identity is exactly the multiplicativity of
$r'$.  Also $(1+z)^*(1+z)=1$ and
$(1-y)^*(1-y)=1$; hence its multiplier and the symplectic
condition on $u'$ are unchanged.
The representation remains small and adically continuous,
with the same reduction.  If $g$ is in inertia away from
$\ell$ and $r(g)=1$, (IX.11) gives $z(g)=0$, so $r'(g)=1$.
This verifies the full original \texttt{Unr\allowbreak T} and all
$I_p\cap\ker\psi$ conditions, including their specified
kernel restrictions.

For $h\in G_\ell$ set
$W_h=z(h)-(r(h)yr(h)^{-1}-y)\in P$.
Using the same zero-product identities gives
\[
 (u')^{-1}r'(h)u'
           =(1+u^{-1}W_hu)(u^{-1}r(h)u).
 \tag{IX.21}
\]
This is upper triangular with exactly the old diagonal,
because the first added term is strictly upper triangular.
Thus the integral flag, residual labels and prescribed
inertia characters are all unchanged.  Every field of
the new \texttt{KPair} has been checked.

The matrix $r'$ is block triangular with the same two
diagonal blocks as $r$, proving the characteristic
polynomial identity.  The $B$-block has full span:
the symplectic multiplier identifies it with the
inverse transpose of the $A$-block, multiplied for
each $g$ by its nonzero similitude.  As $g$ ranges over
$G$, inversion and transpose preserve the full-span
condition, and these individual nonzero scalar factors
do not change the span.
If a cross-block $Y$ split the representation, expansion
of $(1-Y)r'(g)(1+Y)$ would give
$z(g)=Y-r(g)Yr(g)^{-1}$.
At $c$, (IX.19) and $z(c)=0$ give $2Y=0$, so $Y=0$;
then $z=0$, a contradiction.  The same argument works
after passage from $A$ to $K$, so the nonsplitting is
exactly the required \texttt{Block\allowbreak Ext} over $K$.
\end{proof}
% END IX6

% BEGIN IX7
\begin{proposition}[IX7: a block diagonal pair violates the exact weak purity bound]
\label{claim:IX7}
For every characteristic-zero ambient identification
satisfying \texttt{Is\allowbreak Zero}, every block diagonal
\texttt{KPair} fails the literal \texttt{Pure\allowbreak W} condition
(IX.2), for every fixed $\iota$.
\end{proposition}

\begin{proof}
Transport the pair to $\overline{\mathbb Q}_\ell$ using
the specified identification.  Smallness makes all its
entries continuous: an element in the $n$th power of
the defining DVR's maximal ideal has norm at most
$a^n$ for a fixed $a<1$.  Thus its two block determinants
$\delta_A,\delta_B$ are continuous characters.

The ordered ordinary flag gives
\[
 \delta_A|_{I_\ell}=\theta_0\theta_2,
                  \qquad
 \delta_B|_{I_\ell}=\theta_1\theta_3.
 \tag{IX.22}
\]
To justify the assignment of the two characters, use the
two blocks of $r(\sigma_\ell)$.  Their characteristic
polynomials are coprime over the valuation ring, since
their residual eigenvalues are disjoint by AX6.
The polynomial block projectors therefore preserve the
ordinary flag.  Their diagonal values on the flag
quotients are $1,0,1,0$ and $0,1,0,1$, as determined
by the original residual labels.  Multiplying the
appropriate quotient characters proves (IX.22).

Each of $\delta_A\chi_\ell^{-(e+2)}$ and
$\delta_B\chi_\ell^{-e}$ has finite image.
Here is the full global argument.  Choose an integer
$m>0$ killing the finite quotients attached to
$\bar r,\psi,\omega$.  On inertia at $\ell$, (IX.22)
makes the $m$th power of the relevant ratio trivial.
At $p\in S_b$ this follows from identity on
$I_p\cap\ker\psi$, and away from $T$ from
\texttt{Unr\allowbreak T}; the cyclotomic character is unramified
away from $\ell$.
Conjugacy of places gives the assertion at every finite
prime.  The characteristic-zero small DVR is finite
over $\mathbb Z_\ell$: it is complete, its residue field
is finite, and $\ell$ is a nonzero power of its
uniformizer times a unit, so successive approximation
from a basis modulo $\ell$ proves finiteness.
In this finite coefficient field defining the
pair, reduce the powered ratio modulo every maximal
ideal power.  Minkowski's theorem makes each of these
everywhere-unramified finite characters trivial.
Separatedness shows the powered ratio itself is $1$.
A continuous character with image in the finite set of
$m$th roots of unity has open kernel.
Hence on some common open subgroup $H_0$,
\[
                       \delta_A=\chi_\ell^{e+2},
                       \qquad \delta_B=\chi_\ell^e.
 \tag{IX.23}
\]

Suppose (IX.2) held with $H,M$.  Choose a prime
$p\nmid \ell M$.  An arithmetic Frobenius in the
chosen decomposition group at $p$ has some positive
power $\sigma$ in $H\cap H_0$, since that subgroup
has finite index.  If the power is $f>0$, then
$\chi_\ell(\sigma)=p^f$ and $\sigma$ is an arithmetic
Frobenius of degree $f$ at that same place.
Let $a_1,a_2$ and $b_1,b_2$ be the roots of the two
block characteristic polynomials.  They are all
nonzero because both blocks are invertible.
Using (IX.2) for $(a_1,b_1)$ and $(a_2,b_2)$ gives
\[
 |\iota\delta_A(\sigma)|
                   <p^{2f}|\iota\delta_B(\sigma)|.
 \tag{IX.24}
\]
But (IX.23) makes the two sides respectively
$p^{f(e+2)}$ and $p^{2f}p^{fe}$, which are equal.
This is a strict-inequality contradiction.
It uses exactly \texttt{Pure\allowbreak W}, including its finite
exceptional set, open subgroup, all-place and
positive-Frobenius-degree quantifiers.
\end{proof}
% END IX7

% BEGIN IX8
\begin{theorem}[IX8: the complete actual image-density reduction]
\label{claim:IX8}
For the original good datum and one fixed $\iota$,
\[
 \begin{gathered}
 \operatorname{Con}(D),\quad
 \operatorname{PurityAtPoints}(D,\iota),\quad
 \operatorname{AutmAlt}(D,\iota)\\
 \Longrightarrow\quad
 D.\texttt{Fin0}\ \land\ D.\texttt{No\allowbreak End}\ \land\
 D.\texttt{All\allowbreak Irr}(\texttt{Is\allowbreak Aut4Alt}\,\iota).
 \end{gathered}
 \tag{IX.25}
\]
All cohomological, valued-point, curve and block-extension
steps in this implication are proved, including the
coefficient identifications.  The three displayed arithmetic inputs
are not themselves asserted by it.
\end{theorem}

\begin{proof}
If \texttt{Fin0} failed, IX3 would give a characteristic-$\ell$
actual valued point with a weight-zero monomial satisfying
(IX.5).  IX2 gives its \texttt{KPair}.
Apply IX4.  Full matrix span contradicts
$\operatorname{Con}(D)$.  In \texttt{Diag\allowbreak Type}, IX4 makes
the chosen monomial fixed by a positive Frobenius power,
contradicting (IX.5).  In the two-block case, IX4 first
gives a block diagonal pair with full $A$-block span;
IX6 turns it into a \texttt{Block\allowbreak Ext} pair.  This also
contradicts $\operatorname{Con}(D)$.
Thus \texttt{Fin0} holds on all original actual lifts.

To prove \texttt{No\allowbreak End}, suppose an original algebraic
point were block diagonal with full $A$-block span.
IX2 places it in the standard ambient as an actual
\texttt{KPair}.  IX6 gives a new \texttt{Block\allowbreak Ext} pair
with exactly the same characteristic polynomial at
every group element.  The stated
$\operatorname{PurityAtPoints}$ applies to this new pair.
Since (IX.2) concerns only those characteristic
polynomials, it transfers to the old block diagonal
pair.  IX7 rules this out.  This is why the scope of
$\operatorname{PurityAtPoints}$ was retained literally.

Finally take any original algebraic point with full
matrix span.  IX2 makes it a \texttt{KPair} in the same
standard ambient, so $\operatorname{AutmAlt}(D,\iota)$
applies.  A field isomorphism carries a full matrix
span onto a full matrix span: map a finite spanning
linear combination entrywise and map its scalar
coefficients by the same isomorphism.
Thus no restriction on the actual point has been lost.
The isomorphism $\iota$ was fixed before this arbitrary
point was chosen.  It is the one common embedding in
the conclusion, not a point-dependent choice.
\end{proof}
% END IX8

\subsection{The finite-weight-fibre obstruction and the actual promodular witness}

The next two results unpack the exact source witness rather
than treating ``promodular'' as a classical theorem.  A
\texttt{Pro\allowbreak Mod\allowbreak Alt} witness consists of a tower $T$, a
permitted coefficient extension $O\subset O'$, an actual
\texttt{Space} $X$ for $T$, an ideal $J\subset X.B$, a map
$y:X.A\to V$, and a global representation $r_1$ over $V$.
Here $X.B$ is complete Noetherian local with residue
$k'=O'/\mathfrak m_{O'}$, $X.B\hookrightarrow X.A$ is
injective, and $X.A$ is a finite-type $X.B$-algebra.
The universal framed pair $P$ has characteristic
polynomials $c_g(X)\in X.B[X]$ and local characters
$\psi_{vi}(s)\in X.B$.  Their residual values are the
specified four characters of $D$, with the tower's
permutation of labels.  The field $F=T.F$ is totally
real and abelian over $\mathbb Q$, $K_0=T.K$ is imaginary
quadratic, and $L=FK_0$.

The witness includes the full, unchanged
\texttt{T.\allowbreak Fit}, \texttt{T.\allowbreak Deep}, and
\texttt{Potential\allowbreak Modularity\allowbreak Tower}$\,\iota\,X\,J$.
It also includes all of the following equations:
\[
 \begin{gathered}
 y|_{O'}=a_{O'},\qquad
 y(\mathfrak m_{X.B})\subset\mathfrak m_V,\qquad y(J)=0,\\
 r_K\simeq (r_1)_K,\qquad
 r_1|_{G_F}=P\otimes_yV,\qquad
 r_1(g)^*r_1(g)=a_O\vartheta(g),\\
 y(\psi_{vi}(s))=a_O\theta_i(s)
       \quad(s\in I_{F,v}),
 \end{gathered}
 \tag{IX.26}
\]
where the last line uses the original label convention of
\texttt{Pro\allowbreak Mod\allowbreak Alt}; the residual characteristic-polynomial
formula uses the tower's possibly permuted labels.
The five clauses of the potential-modularity tower record are
a finite special weight fibre, local vanishing at the
specified characteristic-$\ell$ curve points, local
vanishing at actual seed points, weak purity for the
specified arithmetic points, and literal
\texttt{Is\allowbreak Aut4Alt} for their irreducible global descents.
Only the first, fourth and fifth clauses are used below.
The unused clauses remain part of the actual source
record; no existence of such a record is inferred from
their names.

% BEGIN IX9
\begin{proposition}[IX9: finite weight modules, residual traces, and the full obstruction]
\label{claim:IX9}
The finite-fibre clause used in the obstruction has a
complete algebraic supplier.  Namely, on an actual
\texttt{Fit}, \texttt{Deep} tower retain a
\texttt{Link}, its universal ordinary determinant ring
$R$, its universal inertial-character algebra $\Lambda$,
an ideal $J_0\subset R$, and the actual
\texttt{Is\allowbreak Phi}$\,0$ map $\phi:R\to B=X.B$.
If $R/J_0$ is finite over $\Lambda$, then for
$J=J_0B$ there is $N\ge0$ such that
\[
 \mathfrak m_B^N\subset
 I:=J+\mathfrak m_{O'}B+
       (\psi_{vi}(s)-1:s\in I_{F,v}).
 \tag{IX.27}
\]
Consequently, every actual characteristic-$\ell$
\texttt{KPair} admitting a \texttt{Pro\allowbreak Mod\allowbreak Alt} witness
has neither full matrix span nor \texttt{Block\allowbreak Ext}.
The conclusion uses no automorphic purity or
\texttt{Cohomology\allowbreak Hyp} assumption.
\end{proposition}

\begin{proof}
\emph{Finite module to finite special fibre.}
Take a prime $\mathfrak q\supset I$ and put
$C=B/\mathfrak q$.  It is a complete local domain with
residue $k'$, and $\mathfrak m_{O'}C=0$.
An element of $C$ integral over $k'$ must belong to
$k'$: subtract its residue in $k'$; if the difference
were nonzero, it would be algebraic over a field and
therefore invertible, although it lies in the maximal
ideal.  Write $\pi:B\to C$.

The map $\Lambda\to R\xrightarrow{\phi}B\to C$
sends every universal inertial character to a scalar
in $k'$.  Indeed the exact \texttt{Is\allowbreak Phi}$\,0$
equation is
\[
 \phi(\psi^R_{vi}(s))
        =\kappa(s)\psi_{e(v),i}(s|_F),
 \tag{IX.28}
\]
with $\kappa(s)\in(O')^\times$; the last factor is $1$
in $C$ for inertial $s$.  The \texttt{Link} identifies
the places and inertia groups in this formula.
These characters therefore equal those from the
residue map $\Lambda\to k'\to C$.
Both maps satisfy the same inertial lifting conditions:
they are obtained from actual local characters, their
residues agree, and their prescribed product and
multiplier equations are preserved.
The uniqueness clause of \texttt{Lam\allowbreak Obj.\allowbreak univ} makes
the two maps equal.  Thus the entire image of
$\Lambda$ is contained in $k'$.
Since $R/J_0$ is finite over $\Lambda$, every element
of the image of $R$ in $C$ is integral over $k'$.

For $g\in G_F$, its square belongs to $G_L$, since
$L/F$ is quadratic.  The characteristic-polynomial
part of \texttt{Is\allowbreak Phi}$\,0$ identifies the polynomial
of the universal determinant at a lift of $g^2$ with
$c_{g^2}$ after multiplying its roots by the coefficient
unit $\kappa(g^2)$.  Consequently every coefficient
of $\pi(c_{g^2})$ is integral over $k'$; multiplication
by an $O'$-unit does not affect integrality.
It follows that every coefficient of $\pi(c_g)$
is integral over $k'$ as well.  Here is a proof that
does not assume that $B\hookrightarrow X.A$ remains
injective after quotienting.  For formal roots
$z_1,\ldots,z_4$, each $z_i$ satisfies the monic
polynomial $\prod_j(T^2-z_j^2)$.  Thus each elementary
symmetric polynomial in the $z_i$ is integral over
the ring of elementary symmetric polynomials in the
$z_j^2$.  The resulting monic polynomial identities
are universal identities for the coefficients of
the characteristic polynomials of a matrix and its
square.  Apply them to $P(g)$ in $X.A$ and descend
the identities to $B$ using the original injection
$B\hookrightarrow X.A$; only then apply $\pi$.
Transitivity of integrality proves the assertion.

Every element of a chosen decomposition group of
$G_F$ at $\ell$ lifts to $G_L$: the splitting clause
of \texttt{T.\allowbreak Ok}, with $\ell\in D.S$, puts that
decomposition group in the subgroup fixing $K_0$.
Equation (IX.28), now for arbitrary decomposition-group
elements, shows that each $\pi(\psi_{vi}(s))$ is
integral over $k'$.  All characteristic-polynomial
coefficients and all these local character values
therefore equal their residues in $k'\subset C$.
The exact \texttt{Space.\allowbreak gen\allowbreak B} uniqueness clause says
that two maps $B\to C$ agreeing on these quantities
are equal.  Compare $\pi$ with the residue map
$B\to k'\to C$.  It follows that
$\pi(\mathfrak m_B)=0$, hence
$\mathfrak q=\mathfrak m_B$.
Every prime over $I$ is therefore the maximal ideal.
Since $B$ is Noetherian, some power of $\mathfrak m_B$
lies in $I$: take powers of finitely many generators
lying in $\sqrt I$ and expand a sufficiently large
power of their ideal.  This proves (IX.27).
It supplies exactly \texttt{TPm\allowbreak Fin}'s conclusion;
the source's additional \texttt{Unramified\allowbreak Inputs}
argument is not used in this algebraic proof.

\emph{Finite fibre to constant characteristic polynomials.}
Now take an actual \texttt{Pro\allowbreak Mod\allowbreak Alt} witness in
characteristic $\ell$.  Its \texttt{fin} clause is
(IX.27), possibly with a larger ideal $J$.
Every coefficient map $a_{O'}$ kills
$\mathfrak m_{O'}$: its prime kernel contains the
nonzero element $\ell$, and a nonzero prime of a DVR
is its maximal ideal.  Equation (IX.26) kills $J$.
The \texttt{T.\allowbreak Ok.\allowbreak ine\allowbreak L} condition says that the original
residual representation is trivial on $I_{F,v}$;
therefore (IX.26) also kills each $\psi_{vi}(s)-1$.
The map $B\to V$ thus kills $\mathfrak m_B^N$.
Its kernel is prime, so it kills $\mathfrak m_B$ and
factors through $k'$.
The actual \texttt{Space.\allowbreak cp\_\allowbreak red} and the conjugacy
in (IX.26) now give
\[
 \det(X-r(g))=\prod_{i=0}^3
       (X-a_O\theta_i(g))\qquad(g\in G_F).
 \tag{IX.29}
\]
The tower permutation of the four labels does not
change this product.

For any $g\in G$, choose $a>0$ such that $g^a\in G_F$.
The four elements $a_O\theta_i(g)$ have finite order,
since they belong to the finite residue field of $O$.
Equation (IX.29) for $g^a$ implies that every eigenvalue
$z$ of $r(g)$ is algebraic over $\mathbb F_\ell$:
its $a$th power is one of the algebraic elements
$(a_O\theta_i(g))^a$.  Its trace is therefore algebraic
over $\mathbb F_\ell$.  The difference
\[
              \operatorname{tr}r(g)
                         -\sum_i a_O\theta_i(g)
 \tag{IX.30}
\]
lies in $\mathfrak m_V$ by the original
\texttt{KPair.\allowbreak red}.  A nonzero element algebraic over
a finite field has finite multiplicative order and
is a unit in every valuation ring containing it.
Thus (IX.30) is zero.  Since the $\theta_i$ are
characters, we obtain
\[
 \operatorname{tr}r(xgh)=\operatorname{tr}r(xhg)
                    \qquad(x,g,h\in G).
 \tag{IX.31}
\]

\emph{The trace obstruction for both kinds of star point.}
Suppose a matrix representation $S$ of dimension
$d>1$ spans $M_d(K)$, and an invertible $Q$ satisfies
\[
 \operatorname{tr}(QS(xgh))
                         =\operatorname{tr}(QS(xhg))
                    \qquad(x,g,h\in G).
 \tag{IX.32}
\]
Linear extension in $x$ makes
$\operatorname{tr}(QM[S(g),S(h)])=0$ for every
$M\in M_d(K)$.  Taking
$M=Q^{-1}e_{ji}$ proves that every matrix entry of
the commutator is zero.  Linear extension in $g,h$
would make the entire matrix algebra commutative,
contrary to its two elementary matrix units.
Taking $Q=1$ in (IX.31) rules out full matrix span.

For \texttt{Block\allowbreak Ext}, write the diagonal blocks
as $A(g),B(g)$ in the order $V_A,V_B$, with $B$
spanning $M_2(K)$.  The distinguished $g_0\in G_F$
from \texttt{T.\allowbreak Ok.\allowbreak dist} has four distinct eigenvalues
by (IX.29).  Put
\[
 f_A(X)=X^2-\operatorname{tr}(A(g_0))X+\det A(g_0),
                 \qquad Q=f_A(B(g_0)).
 \tag{IX.33}
\]
The two block characteristic polynomials are coprime.
Bezout's identity and Cayley--Hamilton make $Q$
invertible.  Cayley--Hamilton on the $A$ block also gives
\[
 \begin{split}
 \operatorname{tr}(QB(y))={}&\operatorname{tr}r(g_0^2y)
     -\operatorname{tr}(A(g_0))\operatorname{tr}r(g_0y)\\
     &+\det A(g_0)\operatorname{tr}r(y).
 \end{split}
 \tag{IX.34}
\]
Apply (IX.31) to the three terms.  It gives (IX.32)
for the full $B$-block span, another contradiction.
This last argument uses only the triangular pattern
and the full block span; it therefore rules out the
source's stronger nonsplit \texttt{Block\allowbreak Ext} assertion
with no extra premise.  This proves the
finite-weight-fibre-to-star-obstruction chain.
\end{proof}
% END IX9

% BEGIN IX10
\begin{theorem}[IX10: exact characteristic-zero transport and the common embedding]
\label{claim:IX10}
An actual \texttt{Pro\allowbreak Mod\allowbreak Alt} witness, together with
the source's \texttt{Is\allowbreak Zero} identification, gives
the literal \texttt{Pure\allowbreak W} assertion for $r$.
If $r,u$ is also an actual \texttt{KPair} with full
matrix span, it gives the literal
\texttt{Is\allowbreak Aut4Alt}$\,\iota$ assertion for $r$.
Consequently, for one fixed original good datum and
one fixed $\iota$, the assertion
\[
 \begin{gathered}
 \text{every actual \texttt{KPair} with \texttt{Star},
 in every \texttt{Amb},}\\
 \text{admits the literal \texttt{Pro\allowbreak Mod\allowbreak Alt}$\,D\,\iota$
 witness}
 \end{gathered}
 \tag{IX.35}
\]
implies the entire conclusion of (IX.25), on the same
datum and with that one common $\iota$.
This theorem proves the reduction from (IX.35);
it does not assert (IX.35).
\end{theorem}

\begin{proof}
Let $f:V\to\overline{\mathbb Q}_\ell$ be the
valuation-preserving coefficient-compatible
identification, and compose $y$ with $f$.
It is an $O'$-algebra map, since
$f a_{O'}$ is the original inclusion of $O'$.
It kills $J$, and the transported $r_1$ restricts
to $P\otimes_{fy}\overline{\mathbb Q}_\ell$ on $G_F$.

The weight hypothesis of \texttt{Potential\allowbreak Modularity\allowbreak Tower}
must be checked with its actual normalization.  On the
open subgroup
\[
             G_F\cap\ker\psi\cap\ker\omega
 \tag{IX.36}
\]
and on inertia at every place over $\ell$, the last
line of (IX.26) becomes
$f y(\psi_{vi}(s))=\chi_\ell(s)^{w_i}$,
where $(w_0,w_1,w_2,w_3)=(e+1,e,1,0)$.
The original character formula at the chosen place
extends to every such place by conjugacy of places;
one-dimensional characters are unchanged by
conjugation.  The source's \texttt{T.\allowbreak Wts} takes
\emph{squares} of these local characters.
Its parameter is therefore precisely
$(2(e+1),2e,2,0)$, not $(e+1,e,1,0)$.
The \texttt{pur} clause now supplies (IX.2) for
$f r_1$.  Conjugacy in (IX.26) makes all characteristic
polynomials of $fr$ and $fr_1$ identical.  Every
quantifier of \texttt{Pure\allowbreak W}, including the same
open subgroup and excluded integer, transfers directly.

For automorphy there are two more hypotheses to prove.
First, smallness supplies continuity of $fr$: a
uniformizer of its defining DVR has image of norm
strictly less than $1$, and congruence modulo its
$n$th power bounds each matrix-entry difference by
that norm to the $n$th power.  Conjugacy gives
continuity of $fr_1$.  The multiplier equation of
(IX.26), with coefficient compatibility of $f$,
gives exactly $D.\texttt{Mult}(fr_1)$.

Second, full matrix span on $G$ must remain full
on $G_L$.  We give the whole argument.
Write $\rho=fr$ and $H=G_L$.  The extension $L/\mathbb Q$
is finite abelian, since $F/\mathbb Q$ is abelian and
$K_0/\mathbb Q$ is quadratic.  Thus $H$ is normal and
$G/H$ is finite abelian.
The \texttt{T.\allowbreak Ok.\allowbreak g3} clause supplies, for each
nonzero root difference, an element of $G_L$ on
which that residual root character is not $1$.
In the source's notation the root character satisfies
$\operatorname{rt}(w_{ij})\chi_j=\chi_i$;
this follows by substituting the two torus characters
and $\chi_0\chi_3=\chi_1\chi_2$.
After undoing the tower permutation, the four
$\bar\theta_i|_H$ are consequently distinct.

Full matrix span makes the four-dimensional
representation of $G$ simple.  Its restriction to
$H$ is semisimple: choose a simple $H$-submodule;
the sum of its finitely many $G/H$-translates is
a nonzero $G$-submodule and hence the whole space.
A sum of simple submodules is semisimple.
If this restriction were reducible, its commutant
$\operatorname{End}_H(\overline{\mathbb Q}_\ell^4)$
would have dimension greater than $1$.
The finite abelian quotient acts by conjugation
on this commutant.  In characteristic zero that
action diagonalizes.  Its fixed part is
$\operatorname{End}_G=\overline{\mathbb Q}_\ell$
by Schur's lemma, so there is a nontrivial character
$\eta$ of $G/H$ and a nonzero eigenvector $T$
intertwining $\rho$ with $\rho\otimes\eta$.
Its kernel and image are $G$-stable, so $T$ is invertible.
For $g$ with $\eta(g)\ne1$ and any $h\in H$, equality
of traces under this intertwiner gives
\[
                        \operatorname{tr}\rho(gh)=0.
 \tag{IX.37}
\]
The map $f$ is injective.  Thus (IX.37) is already
an equality for $r$ over $V$, which can be reduced
modulo $\mathfrak m_V$.  The original residual
condition gives
\[
       \sum_{i=0}^3\bar\theta_i(g)\bar\theta_i(h)=0
                                  \qquad(h\in H).
 \tag{IX.38}
\]
Distinct characters of any group into any field are
linearly independent: in a relation with the fewest
nonzero coefficients, translate by an element on
which two characters differ and subtract a suitable
scalar multiple to obtain a shorter relation.
The four characters here are distinct, whereas
every coefficient $\bar\theta_i(g)$ is a nonzero
unit.  Equation (IX.38) is impossible.
The restriction is therefore simple, and its matrix
span over the algebraically closed field is full.
For completeness, the latter implication is the
finite-dimensional density argument: a proper
matrix subalgebra acting simply has scalar
commutant, and the semisimple quotient acts through
a single full matrix block; its radical acts as
zero on the simple faithful module, so the algebra
itself is that full block.
This is the finite-dimensional density theorem over
an algebraically closed field; see Etingof et al.,
\emph{Introduction to Representation Theory},
Theorem~2.5 and Section~2.5,
\url{https://math.mit.edu/~etingof/replect.pdf}.

Conjugacy transfers full span on $G_L$ to $fr_1$.
All hypotheses of the actual \texttt{aut} clause
are now verified, so it gives literal
\texttt{Is\allowbreak Aut4Alt}$\,\iota$ for $fr_1$.
This predicate is invariant under isomorphism of
the Galois representation.  In its defining
\texttt{Is\allowbreak Aut\allowbreak QA}, continuity and irreducibility
are invariant under a fixed conjugation; for each
quadratic CM extension the same cuspidal
representation and regular algebraic weight remain,
and its attachment intertwiner is composed with
that conjugation.  The same operation transports
the intertwiners in \texttt{Wit\allowbreak A}.
Thus it gives exactly the stated predicate for $fr$.

Finally (IX.35) and IX9 give
$\operatorname{Con}(D)$.
The first part of this theorem gives
$\operatorname{PurityAtPoints}(D,\iota)$ on all
star pairs, including the pairs created by IX6.
The second part gives
$\operatorname{AutmAlt}(D,\iota)$ on all full-span
pairs.  IX8 now proves (IX.25).
Every construction used the original $D$, coefficient
maps and local groups.  The auxiliary $O'$ occurs
only inside the literal witness and is compatible
with the ambient maps; it does not change the
original quantified endpoint.
The single $\iota$ is fixed before any ambient or
point is chosen, exactly as required.
\end{proof}
% END IX10

\subsection{The remaining arithmetic frontier, with source quantifiers}

The completed results in this chapter retain exactly
the hypotheses printed in each statement: the
geometric/cohomological constructions IX1--IX7 and
the algebraic reductions IX9--IX10.  IX8 is the
complete implication from the three explicitly
stated arithmetic predicates.  No one of those
predicates, and no actual witness (IX.35), has been
manufactured by assuming that a Lean theorem named
``published'' is an applicable classical theorem.

The next source producer of (IX.35) separates
the original \texttt{Of\allowbreak Case} data into
\texttt{Ctx2} and \texttt{Ctx3}, with the same
$\Omega,O,D$ and one choice of $\iota$.
Its remaining inputs are the following.
\begin{enumerate}
\item \textbf{Actual modular support.}
The literal \texttt{PMEx\allowbreak Alt} says: for every tower
$T$ satisfying \texttt{Fit} and \texttt{Deep}, and
every $\iota$, there is a coefficient threshold
$O_1$ such that for every permitted $O'$
containing both $O$ and $O_1$, and for every actual
$\texttt{Space}\ X$ over $O'$, there is an ideal
$J\subset X.B$ satisfying all five clauses of
\texttt{Potential\allowbreak Modularity\allowbreak Tower}.
The threshold precedes both $O'$ and $X$.
The source constructs $J$ as a sum of images of
two modular ideals under two actual
\texttt{Link.\allowbreak Is\allowbreak Phi} maps.  IX9 closes its
finite-weight-module-to-\texttt{fin} step completely.
It does not construct the modular ideals, their
finite modules, the two links, or the other four
clauses.  In particular finite cohomological
windows from TW do not supply a modular support.

\item \textbf{Seed points with the prescribed residual entries.}
\texttt{Seed\allowbreak IIAlt} quantifies over every original
\texttt{Ctx2} datum and every
\texttt{Ok}, \texttt{Two}, \texttt{Deep} tower,
then over every $\iota$; it requires a coefficient
extension and an actual continuous residual
\texttt{Seed\allowbreak Pair\allowbreak Alt} in \texttt{In\allowbreak V}, with the
specified $-w_{\rm Al}$ alternative or both
$w_{\rm Be},-w_{\rm Th}$ entries.
For \texttt{Seed\allowbreak IVAlt}, fix an original
\texttt{Ctx3} datum and an \texttt{Ok}, \texttt{Four}
tower $T_0$.  For every $\iota$ and every ambient
field, one must choose a prime $w$ splitting in
$T_0$ and an actual residual pair $s$ for the
frame with the added prime $w$, supported in
$n_M$.  This choice is made before the following
universal quantifier: for every tower $T\ge T_0$
with \texttt{Ok} and \texttt{Dies}$\,\ell\,w\,T.F$,
there must be a coefficient extension and actual
\texttt{Seed\allowbreak Pair\allowbreak Alt} $Q$ whose reduction lies over
that same $s$, is in \texttt{In\allowbreak A}, and has both
the $-w_{\rm Be}$ and $-w_{\rm Th}$ entries.
These are exact seed constructions, including
regular algebraic cuspidal attachment after the
specified twist and the parity and gap conditions;
they are not supplied by generic potential
automorphy or by a theorem about already
automorphic points.

\item \textbf{Reaching the support from the actual star point.}
The case-two and case-three producers require
their actual tower and space constructions,
entry-pattern reductions, characteristic-$\ell$
curve connections and component chains.
The resulting point must satisfy every equation
of (IX.26), especially $y(J)=0$, its original
inertial-character equations and conjugacy of the
global representation.  Existence of a nonempty
modular support does not give these assertions.
The case-three source route includes
\texttt{Red\allowbreak A}, \texttt{YCoin}, \texttt{Fib\allowbreak Y},
\texttt{Chain4}, \texttt{Chain\allowbreak Desc4} and
the literal fourth-chain reduction.  This chapter records their
exact producer anchors but does not promote them
to reconstructed human proofs.
\end{enumerate}

The original top-level source obtains these
arithmetic objects through \texttt{Hypotheses\allowbreak H'}
and the unramified-input record, printed Hida and
adjoint-lifting constructions, and the rank-two
and rank-four seed producers.
\texttt{Hypotheses\allowbreak H'} contains exact literal
automorphy, cohomological, base-change, descent,
semisimplicity and quaternionic assertions; it
has not been treated as one named standard
theorem.  Even a complete supplier of
\texttt{Cohomology\allowbreak Hyp} would not by itself supply
the modular ideals or seed and component chains.
Conversely the entire lower reduction proved
here already uses only DX1, DX2, DX5--DX7,
AX6, BR7, FP2--FP3, MX1, PX1/PX5/PX6 and LD3.
It needs no still-open Bloch--Kato record.

The exact unclosed endpoint is therefore the
actual arithmetic realization (IX.35), uniformly
for the original \texttt{Good}/\texttt{Of\allowbreak Case}
data and with one $\iota$ before the point
quantifiers.  Once it is supplied, IX8--IX10
give literal \texttt{Symplectic\allowbreak Four\allowbreak Image\allowbreak Dense};
the later MX/PX model and point comparisons can
then use its \texttt{Fin0}, \texttt{No\allowbreak End} and
common-embedding automorphy on the very same
datum.  This is a closed downstream dependency
chain with an explicit arithmetic frontier,
not a claim that the full image-density
arithmetic input, RH, or GRH has been proved
in the human reconstruction.


% Version 0.28. Original paired counts with one remaining image-density input.
% Pinned source: 7cbf7a0c32fca1badf05931879a37c3f8a2b03b4.
\section{Joining the paired counts on the original arithmetic objects}
\label{sec:paired-count-installation}

\status{JX1 closes every join from the literal
\texttt{Symplectic\allowbreak Four\allowbreak Image\allowbreak Dense} predicate to the actual
\texttt{Count\allowbreak II} and \texttt{Count\allowbreak III}.  The catalogue and
rigidity are now supplied by CX and RX, while AX/DX/MX/PX
retain the same arithmetic datum and coefficient normalization.
Image-density itself remains an arithmetic input.
JX2 consequently reduces the selected height/window branch
to \texttt{Count\allowbreak I}, \texttt{Count\allowbreak V}, and that one predicate,
with the already stated prime-pair hypotheses retained.}

% BEGIN JX1
\begin{theorem}[JX1: one-input reconstruction of the actual paired counts]
\label{claim:JX1}
The reconstructed chains give the literal implication
\[
 \boxed{\quad
 \texttt{Symplectic\allowbreak Four\allowbreak Image\allowbreak Dense}
      \ \Longrightarrow\
              \texttt{Count\allowbreak II}\wedge\texttt{Count\allowbreak III}.
 \quad}                                                        \tag{JX.1}
\]
Precisely, for every primitive source character $\chi$ there
is a natural exponent $C_\chi$, chosen before
$h,t,\Omega,\ell,c,e,\psi$, for which both source cases have
finite strict Selmer length and
\[
 \frac{\ell_{c.O}V(\Omega,c,e,\psi)}{v_{c.O}(\ell)}
                  \log\ell
                     \le(\ell h.N e)^{C_\chi}.                 \tag{JX.2}
\]
Every host, character, local condition and coefficient
quantifier in \texttt{Count\allowbreak Shape} is retained.
\end{theorem}
\begin{proof}
Fix one witness $h_A$ of the source image-density predicate.
The catalogue theorem CX9 supplies the actual
\texttt{Number\allowbreak Sp4Alt}, with constants chosen before the
varying prime, algebraic closure, complex embedding and
finite family.  RX6 supplies the actual
\texttt{Rigidity\allowbreak Symplectic\allowbreak Four}, including its literal
ordinary and infinitesimal local conditions.  BX8
supplies the complete \texttt{Cohomology\allowbreak Hyp}; RV7, BR7,
FP8 and SP1--SP6 supply the rank-one, residual,
presentation and weighted-intersection inputs in their
already proved scopes.  We verify how these suppliers
meet on one original datum, rather than infer the result
from their names.

Fix $\chi$ and then an original host $h$ with $h.\chi=\chi$,
a permissible $t$, an original $\Omega$, an odd $\ell$,
an actual coefficient ring $O=c.O$, and
$\psi\in h.X(c,e)$ in case II or III, where $0<e<t$.
For the moment suppose $O$ contains the roots of unity
required by both $n_{\mathrm{II}}$ and
$n_{\mathrm{III}}$.  AX1--AX5 construct an actual
$D:\texttt{Raw}\,\Omega\,c$ with \texttt{Good} and
\texttt{Of\allowbreak Case}, $D.e=e$ and the original character.
In particular $e\ge2$ for these cases.  Apply $h_A$
to exactly these choices.  It gives
\[
 D.\texttt{Fin0},\qquad D.\texttt{No\allowbreak End},\qquad
 \exists\iota\ \forall r,\
    D.\texttt{EPt}(r)\wedge\operatorname{Irr}(r)
                  \Longrightarrow\texttt{Is\allowbreak Aut4Alt}(\iota,r).
                                                               \tag{JX.3}
\]
Fix this one $\iota$ before using any irreducible point.
AX6--AX10 give the exact ordinary, level, dual and safe
admissibility properties of every such point for this
same $\iota$.  The auxiliary imaginary quadratic fields
used in RX and CX may vary as their proofs specify;
they do not change $\iota$, $D$, or the original
representations over $\mathbb Q$.

Let $n',p_1,m_0$ be the source's case parameters.
AX7 supplies their precise level and root-order data
\[
      L=\ell\,2^{n'}h.Np_1,\qquad
      k=(\ell-1)\,2^{n'}(2h.o),
                                                               \tag{JX.4}
\]
with $m_0\mid h.o$, the required \texttt{Kills\allowbreak Mu} facts,
and the $k$-th power identities for
$D.\psi,D.\xi,D.\mathrm{om}$.  The alternatives for
$(n',p_1)$ are exactly those in (SP.48).
DX2, DX4 and DX8 supply the actual global/local
finiteness, framed dimensions and local base;
MX10 therefore constructs a \texttt{Model} over this
same $O$.  Apply RX7 to (JX.3), using the just
established level and case facts.  It gives
$D.\texttt{All\allowbreak Irr}\,D.\texttt{Rigid\allowbreak At}$ on this $D$,
not only rigidity of a separate automorphic
representation.

The point, conjugacy, orbit and tangent comparisons
PX1--PX9 now supply all of (G2) in SP10.
The hypotheses \texttt{Fin0}, \texttt{No\allowbreak End} and
actual irreducible rigidity used there have each been
supplied in (JX.3) and the preceding paragraph.
AX10 supplies all of (G1) under $h_A$.
CX8 gives exactly SP9's finite catalogue for the
original Galois level $L$.  Its use of an enlarged
auxiliary automorphic level is already included in
its fixed polynomial exponent; it changes neither
$L$ in (JX.4) nor a hypothesis about $D$.

For clarity, the numerical comparison reached by
these joins is the following.  Choose the source
prime $p=p_0(\ell h.Np_1)\ge5$ and arithmetic
Frobenius $\sigma$ at the original chosen place.
It is outside $D.S$ and prime to $L$.  The three
$k$-th roots obtained by evaluating the original
finite characters have Eisenstein trace
\[
 \operatorname{tr}D.r_0(\sigma)
   =\iota^{-1}(p^{e+1}z_1+p^ez_2+pz_3+1).
                                                               \tag{JX.5}
\]
CX9, or equivalently CX8 followed by SP9, gives
nonvanishing of every irreducible trace difference
and an integer $N_*\ne0$ which bounds the product
over every finite pairwise nonisomorphic family.
PX10 gives the actual $D.\texttt{Ineq}(\sigma)$.
Thus, with the original valuation and the exact
factor two,
\[
 \frac{\ell_O D.\mathrm{Sha}}{v_O(\ell)}\log\ell
 \le -2\log|N_*|_\ell
 \le 2\log|N_*|
 \le 2Ck[L(e+1)]^A\log p.                                     \tag{JX.6}
\]
The first inequality includes finiteness of the
left-hand module.  The original strict Selmer
module embeds in the one used for this bound:
at each added prime its unramified negative-twist
Frobenius scalar is a root of unity times $p^{-e}$,
hence differs from $1$, and multiplication by that
nonzero scalar difference on $E/O$ is surjective.
This is the precise SP10 \texttt{sha\allowbreak Eq} inequality.
It does not assert equality of the two modules.

SP10's arithmetic absorption of (JX.6) uses only
the displayed $L,k,p$, the case formulas for $n'$,
and the host's already stated inequalities.
The new catalogue exponent $A$ is fixed in advance,
so that argument applies unchanged and gives one
$C_\chi$ in (JX.2) for both cases.  This is a choice
before any of the objects fixed in the second
paragraph, since $C_\chi$ depends only on $\chi$
and the fixed catalogue constants.

Finally remove the extra root assumption on $O$.
Adjoin enough $2$-power roots for
$2^{n_{\mathrm{II}}+n_{\mathrm{III}}}$ and take a
common finite coefficient overring $O'$ containing
the original $O$.  SP8 and RV1 prove all required
character transports: kernels, orders, conductors,
the twists $\chi_P$, parity, both case tests and
both integers $n_{\mathrm{II}},n_{\mathrm{III}}$
are unchanged.  They also prove equality of the
normalized Selmer length on the original $\Omega$.
Apply the preceding bound over $O'$ and return
using that equality.  The exponent does not
depend on $O'$ or its ramification.  The resulting
\texttt{Count\allowbreak Shape} for the union of the two cases
restricts to each case, proving (JX.1) with the
exact original coefficient-ring quantifier.
\end{proof}
% END JX1

% BEGIN JX2
\begin{corollary}[JX2: the reduced height and window input]
\label{claim:JX2}
The selected cohomological branch has the sufficient input
\[
 \texttt{Count\allowbreak I}\wedge\texttt{Count\allowbreak V}
      \wedge\texttt{Symplectic\allowbreak Four\allowbreak Image\allowbreak Dense}
                     \ \Longrightarrow\ \texttt{Point\allowbreak HV}.
                                                               \tag{JX.7}
\]
With the already stated prime-pair inputs of VW,
the same premise gives its original-field cohomological
window conclusion for every $C_t\ge1$.  The
$C_t=401$ fixed-pair consequence retains the cutoff
exponent $400$.
\end{corollary}
\begin{proof}
JX1 supplies \texttt{Count\allowbreak II} and \texttt{Count\allowbreak III}
with the exact \texttt{Count\allowbreak Shape} quantifiers.
Together with the two stated counts these are
precisely the four premises of HV10.
The auxiliary \texttt{Count\allowbreak Vs}, \texttt{Count\allowbreak IV}
and \texttt{Count\allowbreak Z} have already been supplied
by HV5, HV6 and CZ6.  Applying HV10 gives the
literal \texttt{Point\allowbreak HV}.

The order of choices in that application is
unchanged.  For each $\chi$ and each prescribed
layer-depth function $j$, first choose the four
count exponents, then the finitely many
bad-prime layer constants and the prime
thresholds, and then the single $C$ in (HV.45).
Only afterwards vary the host, $t,B,E,\Omega$,
local models and finite prime sets.
Consequently its conclusion is the original
\[
  \texttt{VVplus\allowbreak Le}\bigl(E,t,(Nt)^C+CB\bigr)
                                                               \tag{JX.8}
\]
on every original CM field of every permissible
host, with all the same layer hypotheses.
No constant is chosen after observing the field
or finite prime set.

TW8 supplies the single constant record and the
simultaneous layer theorem, and OB8 supplies the
odd-minus bound with its constant $240$.
Substituting the newly obtained \texttt{Point\allowbreak HV}
in the already reconstructed VW3--VW4 proof gives
the claimed window and fixed-pair consequences,
with the same prime-pair inputs and rescaling.
This closes the propagation from (JX.7).
It does not assert that its three premises, or
the separate global family, arithmetic gap,
theta/local installation and metric/spectral
comparison, have been constructed.
\end{proof}
% END JX2

\paragraph{What has and has not been reduced.}
The original four-count statements remain valid
interfaces in all historical arguments.  The
current sufficient condition uses two actual
count suppliers and image-density because JX1
now supplies the other two counts from that
predicate alone.  IX traces a substantive part
of the image-density route and isolates its
remaining actual promodular realization and
automorphic seed inputs.  No logical equivalence
between that sufficient route and RH is claimed.


\section{Dependency audit after the complete rigidity and catalogue joins}
\label{sec:dependency-audit-v28}

The current audit records exact sufficient obligations.  It is
not an exhaustive transitive graph of Lean declarations and does
not assign a completion percentage.  Every version-0.27 semantic
node is retained; changed nodes include their full prior snapshots.
All historical claim fields, source arrays and claim-dependency
edges remain unchanged.  The new supplier conclusions are additive.

The important completed dependency chains are the following.
BX carries the actual tied period models through finite cocycles,
local duality and torsion-inclusive integral annihilators to the
whole \texttt{Cohomology\allowbreak Hyp}.  RX carries each actual dual-number
point through the ordinary crystalline condition, cyclotomic
irreducibility, precise automorphic polarization and the unitary
adjoint projection to \texttt{Rigid\allowbreak At}.  CX carries fixed-cell
cohomological bounds and normalized arithmetic Frobenius data
through an actual restriction of the original representation to
the full SP9 catalogue and \texttt{Number\allowbreak Sp4Alt}.

Consequently the current paired-count cut has one substantive
predicate:
\[
 \texttt{Symplectic\allowbreak Four\allowbreak Image\allowbreak Dense}
 \xRightarrow{\mathrm{JX1}}
       \texttt{Count\allowbreak II}\wedge\texttt{Count\allowbreak III}.
\]
The antecedent is not yet supplied.  IX proves the internal
\texttt{Con}/\texttt{Purity\allowbreak At\allowbreak Points}/\texttt{Autm\allowbreak Alt} endpoint
and actual algebraic/promodular transports, while retaining
their finite modular quotient, compatible link/support and
arithmetic seed producers.  Keeping these retained inputs on
the same datum is essential to using the implication.

The selected cohomological branch therefore has the sufficient
cut
\[
 \texttt{Count\allowbreak I}\wedge\texttt{Count\allowbreak V}
       \wedge\texttt{Symplectic\allowbreak Four\allowbreak Image\allowbreak Dense}
       \xRightarrow{\mathrm{JX2}}\texttt{Point\allowbreak HV}.
\]
TW8 and OB8 then feed the already proved VW window, with its
stated prime-pair inputs.  Its $C_t=401$ fixed-pair specialization
still has cutoff exponent $400$.  No constant is moved past the
original host, coefficient or field quantifiers.

The alternate RX and CX proofs supply their exact endpoints.
They do not assert the entire original \texttt{Unramified\allowbreak Inputs},
\texttt{Adjoint\allowbreak Lifting\allowbreak Theorem}, \texttt{Hypotheses\allowbreak H'},
\texttt{Literature\allowbreak Rank\allowbreak NHypothesis}, unchanged-level
\texttt{Level\allowbreak PNAlt}, or all-rank \texttt{Count\allowbreak QAlt} packages.
The separate actual common-family, arithmetic-gap, genuine
theta/local, denominator and exact metric/\allowbreak global/\allowbreak spectral
branches remain in the audit.  The milestone is closure of the
cohomology, rigidity, catalogue and conditional count chains;
it is not completion of the full RH/Dirichlet-GRH reconstruction.
% Version 0.29. Complete arithmetic suppliers for the ordinary universal rings.
% Pinned source: 7cbf7a0c32fca1badf05931879a37c3f8a2b03b4.
\section{The actual level, weight and ordinary determinant rings}
\label{sec:symplectic-universal-arithmetic}

\status{KX1--KX9 close the complete literal predicates
\texttt{Lev\allowbreak Ex}, \texttt{Loc\allowbreak CFTAll}, \texttt{Lam\allowbreak Ex},
\texttt{Phi\allowbreak All}, and \texttt{Pro\allowbreak Modular\allowbreak Ring\allowbreak Exists}.
The last statement retains its actual nonempty-pair antecedent.
The proof constructs the characters, universal maps and defining ideal;
it does not infer existence of a modular-support ideal from existence of
a deformation ring.  Classical local reciprocity and Chenevier's
determinant representability theorem are used in their precise scopes.}

Throughout this section, $\ell$ is a prime, $L$ is a CM number field,
$F=L^+$, and $\Omega$ is the original algebraic closure.
A source \texttt{Hida\allowbreak Setting} $s$ includes the chosen complex
conjugation, the places and valuations above $F$, the auxiliary place
$v_0$, and the splitting and unramifiedness hypotheses recorded in
\texttt{Continuous\allowbreak Res\allowbreak P}.  A level $V$ uses the quotient
$\Gamma=G_{F,S'}$ for its specified finite set $S'$;
$G_0=\ker\eta$ is the CM subgroup of this quotient.
For $v\mid\ell$ we always use the source group
$G_v=\operatorname{Dec}_{L}(s.A(v))$ and its actual inertia $I_v$.
The coefficient ring $O=c.O$ is the ring of integers of the given
finite extension of $\mathbb Q_\ell$, with residue field $k$.
Objects of $\mathcal C_O$ are exactly the source \texttt{Is\allowbreak CO}
algebras: complete Noetherian local $O$-algebras with local structure
map inducing the specified isomorphism of residue fields with $k$.
All occurrences of continuity below mean the maximal-ideal topology.

% BEGIN KX1
\begin{lemma}[KX1: morphisms and principal-unit coordinates]
\label{claim:KX1}
Every $O$-algebra homomorphism between objects of $\mathcal C_O$
is local, induces the identity on $k$, and is continuous.
For $R\in\mathcal C_O$, continuous characters
$\mathbb Z_\ell^d\longrightarrow 1+\mathfrak m_R$
are in bijection with $d$-tuples of elements of $\mathfrak m_R$,
by the images $1+x_j$ of the standard basis vectors.
Consequently $O[[X_1,\ldots,X_d]]$ represents this character functor
for all the algebra homomorphisms occurring in the source.
\end{lemma}
\begin{proof}
If $f:R\to R'$ is an $O$-algebra map, the composite
$R\to R'\to k$ is surjective, since its restriction to $O$ is the
residue map.  Its kernel is therefore a maximal ideal of $R$ and
equals $\mathfrak m_R$.  Thus $f$ is local, and its residue map is
the identity because $O\to k$ is surjective.  The inclusions
$f(\mathfrak m_R^a)\subseteq\mathfrak m_{R'}^a$ prove continuity.
This establishes the required assertion about \emph{all} algebra
maps; no extra continuous-map restriction is being inserted into
the universal properties.

The quotients $R/\mathfrak m_R^a$ are finite.  Indeed their
successive maximal-ideal quotients are finite dimensional over the
finite field $k$, by Noetherianity.  The successive quotients of
$1+\mathfrak m_R$ are subquotients of these additive $k$-spaces.
Hence $1+\mathfrak m_R$ is an abelian pro-$\ell$ group.
For later use there is an elementary uniform estimate: if
$u\in1+\mathfrak m_R$, then
\[
             u^{\ell^a}-1\in\mathfrak m_R^{a+1}.
                                                               \tag{KX.1}
\]
For the induction step factor $z^\ell-1$ as
$(z-1)(1+z+\cdots+z^{\ell-1})$; the second factor belongs to
$\mathfrak m_R$ whenever $z\equiv1\pmod {\mathfrak m_R}$.

If integers $b_m$ converge to $b\in\mathbb Z_\ell$, (KX.1)
shows that $u^{b_m}$ has a well-defined limit $u^b$ independent
of the approximating sequence.  The same estimate proves continuity
and the identity $u^{b+b'}=u^bu^{b'}$.  Thus the character belonging
to $(x_1,\ldots,x_d)$ is
\[
          (b_1,\ldots,b_d)\longmapsto
                       \prod_{j=1}^d(1+x_j)^{b_j}.
                                                               \tag{KX.2}
\]
Uniqueness follows from the density of $\mathbb Z^d$ and the
Hausdorff property of $R$.  Evaluation of a power series at elements
of $\mathfrak m_R$ converges, and defines the unique continuous
$O$-algebra map sending $X_j$ to $x_j$.  Conversely every algebra
map is local by the first paragraph and so has such images.
The argument includes $d=0$ and $\ell=2$.
\end{proof}
% END KX1

% BEGIN KX2
\begin{theorem}[KX2: the complete level existence predicate]
\label{claim:KX2}
For every source $s$ and every finite $S'\supseteq s.S$ satisfying
the source splitting condition at each of its places, there is
$V:s.\texttt{Lev}$ with $V.S'=S'$.  In particular the entire
predicate \texttt{Lev\allowbreak Ex} holds, with the original cyclotomic
character and the original complex conjugation.
\end{theorem}
\begin{proof}
Let $G=\operatorname{Gal}(\Omega/F)$ and let
$H=\operatorname{Gal}(\Omega/L)$, viewed through the source
restriction-of-scalars inclusion.  The fixed-field description gives
$H=\ker(G\to\operatorname{Gal}(L/F))$ and $[G:H]=2$.
This is a closed open normal subgroup.  Define
$\eta_0(g)=1$ on $H$ and $\eta_0(g)=-1$ outside $H$.
The quotient of order two shows that this is a continuous character.
The chosen $c$ acts as complex conjugation under the embedding in
the definition of $s$.  If $c\in H$, that embedding would make
$L$ real.  This contradicts $L$ being totally imaginary; hence
$\eta_0(c)=-1$.  The same embedding proves $c^2=1$.

Let $N_{S'}$ be the closed normal subgroup generated by the inertia
groups at finite places outside $S'$.  The field $L/F$ is unramified
at the finite places by $s.\texttt{unr}$, so every generator of
$N_{S'}$ belongs to $H$.  Since $H$ is closed and normal,
$N_{S'}\subseteq H$.  Also $S'$ contains all places above
$\ell$.  At a finite place of residue characteristic different from
$\ell$, inertia acts trivially on all $\ell$-power roots of unity:
such roots are obtained by lifting roots over the maximal unramified
extension, since their orders are units there.  Thus $N_{S'}$ lies
in the kernel of the original cyclotomic character
$\epsilon_0:G\to\mathbb Z_\ell^\times$.

Both characters therefore descend to $G/N_{S'}=G_{F,S'}$.
Write them as $\eta$ and $\epsilon$.  Their pullbacks are exactly
$\eta_0$ and $\epsilon_0$, so the source fields
\texttt{h\ensuremath{\eta}}, \texttt{\ensuremath{\eta}\_c}, and
\texttt{h\ensuremath{\epsilon}} hold.  The quotient topology gives
continuity of $\epsilon$ and openness of $\ker\eta$.
Together with the given finite set, inclusion and splitting proofs,
these are all the fields of \texttt{Lev}; the underlying $S'$ was
never altered.
\end{proof}
% END KX2

% BEGIN KX3
\begin{theorem}[KX3: local abelian coordinates on the actual decomposition groups]
\label{claim:KX3}
For each $v\mid\ell$ there is a number
$d_v=[F_v:\mathbb Q_\ell]$ and an actual
$\texttt{Loc\allowbreak Ab}(\ell,G_v,I_v,d_v)$.  Moreover
\[
                      \sum_{v\mid\ell}d_v=[F:\mathbb Q].
                                                               \tag{KX.3}
\]
These choices supply the whole predicate \texttt{Loc\allowbreak CFTAll}.
\end{theorem}
\begin{proof}
Completion at the chosen valuation identifies
$\operatorname{Dec}_L(s.A(v))$ with the
absolute Galois group of $K=L_{s.w(v)}$, carrying the original
inertia subgroup to $I_K$.  This identification can be made by
extending the embedding of $L$ to the chosen valued algebraic
closure; restriction and completion are inverse on finite
extensions.  The splitting assumption in $s$ identifies $K$ with
$F_v$.  It is essential here to use this chosen decomposition group,
rather than an unrelated copy of a local Galois group.

Local reciprocity identifies the topological abelianization
$G_K^{\mathrm{ab}}$ with the profinite completion of $K^\times$,
and the image of inertia with $\mathcal O_K^\times$.
We use the arithmetic Frobenius normalization, although reversing
reciprocity changes none of the existence or inertia assertions
below.\footnote{J. S. Milne, \emph{Class Field Theory},
Chapter I, Theorem 1.1 and the local existence theorem,
\url{https://www.jmilne.org/math/CourseNotes/CFT.pdf}.
The classical input is the profinite reciprocity isomorphism and
its identification of units with inertia.}

Here is the necessary group structure.  A uniformizer splits the
valuation sequence for $K^\times$.  The unit group is the product
of the prime-to-$\ell$ Teichm\"uller roots and its principal units.
On a sufficiently deep principal-unit subgroup the convergent
logarithm and exponential give an isomorphism with a nonzero ideal
of $\mathcal O_K$, which is free of rank $d_v$ over
$\mathbb Z_\ell$.  All principal units form a finitely generated
$\mathbb Z_\ell$-module: the deep subgroup has finite index, and
finitely many lifts of generators of the finite quotient suffice.
Its torsion is precisely the finite group $\mu_{\ell^\infty}(K)$.
The structure theorem over $\mathbb Z_\ell$ consequently gives
\[
 G_K^{\mathrm{ab}}\simeq
 \widehat{\mathbb Z}\ \times\mu_{q-1}(K)\ \times P\ \times
                                      \mathbb Z_\ell^{d_v},
 \qquad P=\mu_{\ell^\infty}(K),
                                                               \tag{KX.4}
\]
where $q$ is the residue-field cardinality.  Inertia is exactly the
last three factors.  The logarithm was used only sufficiently deep,
so the assertion also accounts for $2$-power roots of unity.

The $\ell$-primary torsion of the entire topological abelianization
is exactly $P$: $\widehat{\mathbb Z}$ is torsion free and
$\mu_{q-1}(K)$ has order prime to $\ell$.  Projecting (KX.4) gives
continuous maps
\[
 a:G_v\twoheadrightarrow\mathbb Z_\ell^{d_v}\times\mathbb Z_\ell,
       \qquad b:G_v\longrightarrow P.
                                                               \tag{KX.5}
\]
The last coordinate of $a$ is valuation, and
$a(I_v)=\mathbb Z_\ell^{d_v}\times\{0\}$.  The finite-valued
map $b$ has open kernel and restricts to the quotient map on
$\mu_v$, the inverse image of $P$ in $G_v$.
Every continuous map to a finite abelian $\ell$-group factors
through the maximal pro-$\ell$ quotient of (KX.4).  If it kills
$\mu_v$, it kills $P$, and hence factors through $a$.
This verifies, individually, the fields \texttt{cont},
\texttt{surj}, \texttt{inertia}, \texttt{univ},
\texttt{b\_\allowbreak open}, and \texttt{b\_\allowbreak mu} of \texttt{Loc\allowbreak Ab}.
Finally (KX.3) is the dimension identity in
$F\otimes_{\mathbb Q}\mathbb Q_\ell=\prod_{v\mid\ell}F_v$.
Thus the exact degree sum and every original local object required
by \texttt{Loc\allowbreak CFTAll} have been supplied.
\end{proof}
% END KX3

% BEGIN KX4
\begin{proposition}[KX4: prescribed torsion and every local character lift]
\label{claim:KX4}
For every actual residual datum $r:\texttt{Resid}$, every $v,i$
has a character $\psi^0_{v,i}:G_v\to O^\times$ satisfying the
entire predicate $r.\texttt{Is\allowbreak Psi}\,v\,i$.
For any $R\in\mathcal C_O$, all characters satisfying this
predicate are uniquely
\[
 \psi_{v,i}(g)=\psi^0_{v,i}(g)
                 \chi_{v,i}(a_v(g)),\qquad
 \chi_{v,i}:\mathbb Z_\ell^{d_v+1}\to1+\mathfrak m_R
                         \text{ continuous}.
                                                               \tag{KX.6}
\]
Here $\psi^0$ on the right is transported by the original
$O$-algebra structure map, and $a_v$ is the actual map in KX3.
\end{proposition}
\begin{proof}
Let $q=|k|$.  Hensel's lemma applied to $X^{q-1}-1$ gives a unique
multiplicative lift $[\ ]:k^\times\to O^\times$.
Its multiplicativity follows from uniqueness of the lifted root
with each residue.  The residual character $\bar\psi_{v,i}$ has
open kernel and therefore finite image.  It kills $\mu_v$:
it kills the closed commutator subgroup by continuity, and its
remaining image of the $\ell$-primary torsion has order dividing
both an $\ell$-power and $q-1$.

Write $P_v$ for the finite torsion group in KX3.  The map
$\mu_v\to P_v$ is onto, with kernel the closed commutator
subgroup.  The source assumption \texttt{th\_\allowbreak comm} therefore
descends the prescribed character $\theta_{v,i}$ to a character
$t_{v,i}:P_v\to O^\times$.  Every value of $t_{v,i}$ has
$\ell$-power order; its residue in $k^\times$ is consequently $1$.
Set
\[
            \psi^0_{v,i}(g)
              =[\bar\psi_{v,i}(g)]\,t_{v,i}(b_v(g)).
                                                               \tag{KX.7}
\]
Both factors have open kernels, so this is continuous.  Its
reduction is $\bar\psi_{v,i}$, and on $\mu_v$ its first factor is
$1$ while its second factor is exactly $\theta_{v,i}$ by
\texttt{b\_\allowbreak mu}.  These are precisely the three fields
\texttt{cont}, \texttt{red}, and \texttt{mu} of \texttt{Is\allowbreak Psi}.
In particular the prescribed torsion condition has not been
discarded in passing to free coordinates.

For another lift $\psi$ in $R$, put
$\delta=\psi/(\psi^0\otimes_OR)$.  It has values in
$1+\mathfrak m_R$ and kills $\mu_v$.  Reduce it modulo
$\mathfrak m_R^a$.  Its image is a finite abelian $\ell$-group,
by KX1, and its kernel is open.  The \texttt{univ} field of
KX3 shows that it kills $\ker a_v$.  This holds for every $a$;
separatedness of $R$ therefore shows that $\delta$ itself kills
$\ker a_v$.  It factors uniquely through the surjection $a_v$.
This factor is continuous: a continuous surjection from a compact
space to the Hausdorff space $\mathbb Z_\ell^{d_v+1}$ is a
quotient map, and the same is true after testing each finite
quotient of the target units.  Conversely any character $\chi$
in (KX.6) has residue $1$ and is trivial on $a_v(\mu_v)=0$.
The product in (KX.6) thus satisfies all three original conditions.
Surjectivity of $a_v$ proves the asserted uniqueness.
\end{proof}
% END KX4

% BEGIN KX5
\begin{theorem}[KX5: the full inertial weight algebra and its universal maps]
\label{claim:KX5}
For every source $s,V,c,n,r$, including $n=0$, there is a
$\texttt{Lam\allowbreak Obj}\,r$ with
\[
 \Lambda\simeq O[[X_{v,i,j}:v\mid\ell,\ 1\le i\le n,
                                      \ 1\le j\le d_v]]
       \simeq O[[X_1,\ldots,X_{n[F:\mathbb Q]}]].
                                                               \tag{KX.8}
\]
Its universal characters satisfy the exact \texttt{Is\allowbreak Inert}
predicate, and its universal property quantifies over all the
$O$-algebra homomorphisms in \texttt{Lam\allowbreak Obj}.
Consequently the entire source predicate \texttt{Lam\allowbreak Ex} holds.
\end{theorem}
\begin{proof}
Write $a_v=(a_v^{\mathrm{in}},a_v^{\mathrm{ur}})$ as in (KX.5)
and form the first power-series ring in (KX.8).  For $g\in I_v$
define
\[
 \kappa_{v,i}(g)=\psi^0_{v,i}(g)
                    \prod_{j=1}^{d_v}(1+X_{v,i,j})^{
                                                   a^{\mathrm{in}}_{v,j}(g)}.
                                                               \tag{KX.9}
\]
The expression in KX1 defines its powers.  The same formula on all
of $G_v$, with the last coordinate ignored, is an \texttt{Is\allowbreak Psi}
character by KX4.  Thus each $\kappa_{v,i}$ has an actual extension,
which is precisely \texttt{Is\allowbreak Inert}.  This also supplies the
universal character before invoking any universal property.

Let $\kappa'_{v,i}$ be an arbitrary \texttt{Is\allowbreak Inert} system in
$R$.  Choose its witnessing full character $\psi'_{v,i}$, and
factor it using (KX.6).  Restrict the factor to
$\mathbb Z_\ell^{d_v}\times\{0\}$ and let $u_{v,i,j}$ be its
value at the $j$th basis vector.  Then
$u_{v,i,j}-1\in\mathfrak m_R$, so KX1 gives an evaluation map
$X_{v,i,j}\mapsto u_{v,i,j}-1$.  Formula (KX.9) shows that this
map carries the universal inertial system to $\kappa'$.

For uniqueness, choose for each basis vector $e_j$ an element
$g_{v,j}\in I_v$ with $a_v(g_{v,j})=(e_j,0)$, using the exact
inertia image in KX3.  At this element (KX.9) reads
\[
       \kappa_{v,i}(g_{v,j})
                   =\psi^0_{v,i}(g_{v,j})(1+X_{v,i,j}).
                                                               \tag{KX.10}
\]
Every competing $O$-map preserves the first, invertible factor;
hence its image of each $X_{v,i,j}$ is forced.  KX1 proves equality
of the maps.  This proves independence of the initially chosen
extensions $\psi'$, as well as uniqueness in the literal source
category.  There are $n\sum_vd_v=n[F:\mathbb Q]$ variables by
KX3; renaming their finite index set gives the second isomorphism
in (KX.8).  For $n=0$ the system is empty and the ring is $O$.
Finally any other \texttt{Lam\allowbreak Obj} is uniquely isomorphic to this
one by its two universal maps; their composites fix all the
universal characters and are the identities by uniqueness.
\end{proof}
% END KX5

% BEGIN KX6
\begin{theorem}[KX6: the actual global and local finiteness predicates]
\label{claim:KX6}
The cohomology chain already supplied in BX8 gives, for every
source $s$ and $V$,
\[
 \Phi_\ell(\Gamma),\qquad\Phi_\ell(G_0),\qquad
                                \Phi_\ell(G_v)\quad(v\mid\ell).
                                                               \tag{KX.11}
\]
Here $\Phi_\ell(G)$ means that every open subgroup $H$ has only
finitely many continuous homomorphisms $H\to\mathbb F_\ell$,
exactly as in \texttt{Phi\allowbreak Cond}.  Thus the complete predicate
\texttt{Phi\allowbreak All} is supplied.
\end{theorem}
\begin{proof}
We first give the finite-cohomology implication on a compact
profinite group $G$.  For an open $H$, let $M$ be the finite
permutation module of $\mathbb F_\ell$-valued functions on the
left cosets $G/H$, with $(g f)(x)=f(g^{-1}x)$.
Choose representatives $s_x$, choosing $s_H=1$.
For $\varphi:H\to\mathbb F_\ell$ define
\[
       c_\varphi(g)(x)
             =\varphi(s_x^{-1}g s_{g^{-1}x}).
                                                               \tag{KX.12}
\]
The argument of $\varphi$ belongs to $H$.  Multiplication inside
that argument proves the cocycle identity.  A finite set of cosets
and the open kernel of $\varphi$ show that $c_\varphi$ is locally
constant: intersect the finitely many conjugates of this kernel
inside $H$ to obtain an open neighborhood on which all coordinates
are unchanged.  If $c_\varphi-c_{\varphi'}$ is a coboundary,
evaluate at $h\in H$ and the fixed coset $H$.  Every such
coboundary has value $f(h^{-1}H)-f(H)=0$, whereas (KX.12) gives
$\varphi(h)-\varphi'(h)$.  Thus
\[
           \operatorname{Hom}_{\mathrm{cont}}
                    (H,\mathbb F_\ell)\hookrightarrow H^1(G,M).
                                                               \tag{KX.13}
\]
Finite first cohomology for these finite continuous modules proves
$\Phi_\ell(G)$.

Apply (KX.13) first to the original quotient
$\Gamma=G_{F,S'}$.  The set $S'$ is finite and contains all places
over $\ell$, exactly the hypotheses of the global finiteness
field in BX8.  That field applies to the actual permutation
representation, so gives the first assertion of (KX.11).
The subgroup $G_0=\ker\eta$ is open by KX2 or by the given
level record.  Any open subgroup of $G_0$ is also open in
$\Gamma$, which proves the second assertion without a new
cohomology theorem.  Finally the local finiteness field in BX8,
transported to the original chosen decomposition group as in
KX3, applies to the same permutation construction on $G_v$.
This gives the last assertion for every $v$.
Only the two finite-cohomology fields of BX8 were used; no
automorphy or image-density hypothesis occurs in this argument.
\end{proof}
% END KX6

% BEGIN KX7
\begin{theorem}[KX7: universal determinants over the original coefficient ring]
\label{claim:KX7}
Let $G$ be one of the original profinite groups satisfying
$\Phi_\ell$, and let
$\bar D=\det(\bigoplus_{i=1}^n\bar\chi_i)$ be the actual
continuous residual determinant.  There is
$R_{\bar D}\in\mathcal C_O$ and a continuous determinant
$D^{\mathrm{univ}}$ reducing to $\bar D$ such that
\[
 \operatorname{Hom}_{O\text{-alg}}(R_{\bar D},R)
       \longleftrightarrow
 \{D\text{ continuous of degree }n\text{ on }G:
                                      D\bmod\mathfrak m_R=\bar D\}
                                                               \tag{KX.14}
\]
is a natural bijection.  This is the complete \texttt{Det\allowbreak Obj}
interface in the present arithmetic coefficient scope.
\end{theorem}
\begin{proof}
A source \texttt{GDet} is a multiplicative homogeneous
polynomial law on the group algebra, with its full base-change
naturality.  This is precisely the meaning of determinant in
Chenevier's representability theorem.  Its continuity is the
continuity of the characteristic-polynomial coefficients on
$G$.  In the source those coefficients are
$\lambda_i(g)=(-1)^i[X^{n-i}]D(X-[g])$.  For $0\le i\le n$
they are the same coefficients; the source's extra indices
$i>n$ give signed multiples of the constant coefficient and impose
no additional continuity requirement.

For $n\ge1$, Chenevier's Propositions 3.3 and 3.7, with Remark 3.5,
give a complete Noetherian local $W(k)$-algebra representing
continuous determinant deformations of $\bar D$.
The finiteness condition (F) of that paper is exactly
$\Phi_\ell(G)$, supplied by KX6; no irreducibility or inequality
$\ell>n$ is required.\footnote{G. Chenevier,
\emph{The $p$-adic analytic space of pseudocharacters of a profinite
group and pseudorepresentations over arbitrary rings},
\url{https://arxiv.org/abs/0809.0415}, version 2,
\S3.1--3.3, especially Propositions 3.3 and 3.7.
The polynomial-law definition, rather than a trace definition
requiring division by $n!$, is the one used here.}
The residual representation here is the given finite direct sum
of characters with open kernels.  Thus its continuity and its
semisimplicity are explicit, including when $\ell\le n$.

Let this universal $W(k)$-algebra be $A$.  The maximal unramified
subring of the given finite coefficient ring $O$ is canonically
$W(k)$.  Set
\[
                    R_{\bar D}=A\widehat\otimes_{W(k)}O.
                                                               \tag{KX.15}
\]
One may present $A$ as $W(k)[[Z_1,\ldots,Z_t]]/J$;
then (KX.15) is $O[[Z_1,\ldots,Z_t]]/JO[[Z_1,\ldots,Z_t]]$.
It is complete Noetherian local with residue $k$.
The universal law base changes to this ring.  If $R$ is in
$\mathcal C_O$, a determinant over $R$ first gives the unique
continuous $W(k)$-map $A\to R$, and the fixed structure map
$O\to R$ extends it uniquely across (KX.15).  Conversely that
extension pulls the universal law back to the given determinant.
For complete $R$ this assertion follows either directly from the
profinite version of the theorem or by reducing modulo
$\mathfrak m_R^a$ and taking the compatible inverse limit.
KX1 shows that the resulting bijection includes every algebra map
in the literal source interface, proving (KX.14).

If $n=0$, homogeneity applied with scalar $0$ shows that every
value of a degree-zero determinant equals its value at $0$;
multiplicativity and the value at $1$ make this constant $1$.
There is therefore just one degree-zero law and its universal
ring is $O$.  In subsequent equations we still use the source's
literal coefficient functions, which use truncated natural
subtraction $n-i$; no rank-zero convention for their function
\texttt{tr} is silently substituted.
\end{proof}
% END KX7

% BEGIN KX8
\begin{theorem}[KX8: the actual ordinary pair ring]
\label{claim:KX8}
For the original $s,V,c,n,r$, suppose that there is a
$\texttt{Pair}\,r\,R_0$ for some $R_0\in\mathcal C_O$.
Then there is an actual \texttt{Universal\allowbreak Ordinary\allowbreak Determinant\allowbreak Ring}
for $r$, with all its \texttt{Pair} fields and its unique-map
property.  Thus the complete predicate
\texttt{Pro\allowbreak Modular\allowbreak Ring\allowbreak Exists} holds in exactly its stated scope.
\end{theorem}
\begin{proof}
Use KX7 on $G_0$ and $\bar D=r.\texttt{Dbar}$, obtaining
$R_{\bar D}$ and $D^{\mathrm{univ}}$.
Use KX4 and KX1 to represent all the full local character lifts by
\[
 T=O[[Y_{v,i,j}:v\mid\ell,\ 1\le i\le n,
                                         \ 1\le j\le d_v+1]].
                                                               \tag{KX.16}
\]
The last variable for each $v,i$ records the unramified coordinate.
Its universal character is $\psi^0$ times the product (KX.2)
in all $d_v+1$ coordinates.  This is the full \texttt{Lam\allowbreak TObj}
universal property, distinct from the inertial ring of KX5.
The finite completed coproduct with the determinant ring is simply
\[
                 A=R_{\bar D}[[Y_{v,i,j}]].
                                                               \tag{KX.17}
\]
It is in $\mathcal C_O$.  The two constructions prove, directly,
that it represents pairs consisting of a continuous determinant
lifting $\bar D$ and all the prescribed local character lifts,
before imposing the ordinary equations.  Indeed a pair gives a
unique map on $R_{\bar D}$ and forced images in
$\mathfrak m_R$ for every $Y$; evaluation gives the unique map
from (KX.17).

Write $D,\psi$ for this preliminary universal system.  For a group
element $g$ write $[g]$ for its basis element in the group algebra.
Retain exactly the source anti-involution
\[
                 [g]^*=\epsilon(g)^{1-n}[cgc]^{-1}.
                                                               \tag{KX.18}
\]
For $\mathbf g=(g_1,\ldots,g_n)\in G_v^n$, let
\[
       o_v(\psi,\mathbf g)
             =\prod_{i=1}^n([\phi_v(g_i)]-\psi_{v,i}(g_i)[1]),
                                                               \tag{KX.19}
\]
with the product taken in the original order.  Also retain the
literal tame subgroup
$I_p(u)=I_u\cap\bigcap_i\ker\bar\chi_i$.
Let $\mathfrak a\subseteq A$ be the ideal generated by the
following elements, for all of the indicated indices:
\[
 \begin{aligned}
 \lambda_j([g]^*)-\lambda_j([g])
                                      &\quad(j\in\mathbb N, g\in G_0),\\
 \operatorname{tr}_D([h]o_v(\psi,\mathbf g))
                                      &\quad(v,\ \mathbf g\in G_v^n,\ h\in G_0),\\
 \operatorname{tr}_D([h]([g]-[1]))
                                      &\quad(u, g\in I_p(u), h\in G_0).
 \end{aligned}                                                 \tag{KX.20}
\]
In the middle line $\mathbf g\in G_v^n$ and $h\in G_0$;
$\operatorname{tr}_D$ is exactly $\lambda_1$ of the source.
The ideal is closed: every ideal of a complete Noetherian local
ring is closed for its maximal-ideal topology, as follows from
Krull intersection applied to the quotient.

The hypothesized pair over $R_0$ supplies a map $A\to R_0$.
Its \texttt{P}, \texttt{Ord}, and \texttt{I} fields kill,
respectively, the three lines of (KX.20).
Hence $\mathfrak a$ is proper, and lies in $\mathfrak m_A$.
Set
\[
                              B=A/\mathfrak a.
                                                               \tag{KX.21}
\]
It is again an object of $\mathcal C_O$.
Base change of $D,\psi$ gives the universal pair on $B$.
Their continuity, residual determinant and the three
\texttt{Is\allowbreak Psi} fields are preserved under the local quotient map.
The three remaining \texttt{Pair} fields are precisely (KX.20),
so hold by construction.  This verifies every field of the
universal pair rather than only its characteristic polynomial
at unramified Frobenius elements.

For an arbitrary pair $p$ over $R$, the preliminary universal
property gives a unique $f:A\to R$.  Naturality of a determinant
commutes with characteristic-polynomial coefficients; the star
operation and (KX.19) commute with an $O$-algebra map because
their scalars come from $O$ and the local characters.
Thus the three fields of $p$ kill each generator of
$\mathfrak a$, and $f$ factors uniquely through $B$.
Conversely such a factor carries the universal determinant and
every full local character to those of $p$.
If two factors do this, their composites with $A\to B$ coincide
by the preliminary unique-map property, and the factors coincide
by surjectivity.  These are exactly the existence and uniqueness
clauses of \texttt{Universal\allowbreak Ordinary\allowbreak Determinant\allowbreak Ring}.

Nothing in this argument requires $n>0$: KX7 supplied the
rank-zero preliminary law, and any possible degeneracy of the
literal trace equations is controlled by the actual nonempty-pair
antecedent.  Conversely the proof has made no claim that such
a pair exists for arbitrary residual data.  The source predicate
\texttt{Pro\allowbreak Modular\allowbreak Ring\allowbreak Exists} includes this same antecedent.
\end{proof}
% END KX8

% BEGIN KX9
\begin{corollary}[KX9: fixed universal maps for the subsequent modular argument]
\label{claim:KX9}
For every actual input in KX8, the inertial ring $\Lambda$ of KX5
has its prescribed map to the universal pair ring $B$, and for
every $p:\texttt{Pair}\,r\,R$ one has a commuting triangle
\[
                    \Lambda\longrightarrow B\longrightarrow R,
 \qquad \Lambda\longrightarrow R
                         =\Lambda.\texttt{to\allowbreak Pair}(p).
                                                               \tag{KX.22}
\]
All the level, local-character, finiteness and universal-ring
suppliers used in these constructions have now been proved on
the same original arithmetic objects.
\end{corollary}
\begin{proof}
Restrict the full universal characters of KX8 to each actual
$I_v$.  Their \texttt{Is\allowbreak Psi} fields supply the witnesses for
\texttt{Is\allowbreak Inert}; KX5 gives a unique map $\Lambda\to B$.
After composing with the unique map $B\to R$ of KX8, their
restrictions become the local characters of $p$.  The uniqueness
in KX5 therefore identifies the composite with the exact
$\Lambda.\texttt{to\allowbreak Pair}(p)$ of the source.  The same argument
proves compatibility with every further $O$-algebra map, and
identifies these maps with those for any originally chosen
universal rings via their unique universal isomorphisms.

Thus a later argument may fix its original $s,V,c,r,\Lambda,B$
once and use these actual maps.  KX2 supplies the entire
\texttt{Lev\allowbreak Ex}; KX3 the entire \texttt{Loc\allowbreak CFTAll}; KX5 the
entire \texttt{Lam\allowbreak Ex}; KX6 the entire \texttt{Phi\allowbreak All}; and KX8
the entire \texttt{Pro\allowbreak Modular\allowbreak Ring\allowbreak Exists}.
This does not supply the extra \texttt{Ring\allowbreak Props} conditions
on an ideal $J\subset B$, the automorphic support statements,
or the seed and component-reaching assertions of \texttt{PMEx\allowbreak Alt}.
In particular, (KX.20) is the precise \texttt{Pair} presentation;
full determinant annihilation and polarized identities over all
base rings in \texttt{Ring\allowbreak Props} require their own suppliers.
\end{proof}
% END KX9


% Version 0.29. Actual modular support and the two-link construction.
% Pinned source: 7cbf7a0c32fca1badf05931879a37c3f8a2b03b4.
\section{Adjoint vanishing, modular support, and the two-link construction}
\label{sec:symplectic-modular-support}

\status{This chapter proves the literal \texttt{Adjoint\allowbreak Lifting\allowbreak Theorem}
and the characteristic-zero fixed-weight support argument.  It also
constructs the two comparison links and their actual maps from universal
ordinary determinant rings.  Its final tower assertion retains the
precisely stated ordinary modular-quotient and characteristic-$\ell$
support inputs.  Those inputs are not consequences of adjoint Selmer
vanishing.  In particular no complete \texttt{Unramified\allowbreak Inputs},
\texttt{Hypotheses\allowbreak H'} or \texttt{Mod\allowbreak UAlt} package is inferred from the
adjoint theorem.}

\subsection{Conventions and the independent adjoint input}

We retain the original algebraic closure $\Omega$, coefficient subrings
$O\subset\overline{\mathbb Q}_\ell$, and their specified inclusions.
For a CM field $L\subset\Omega$ put $L^+=\texttt{Lp}\,L$.
A source level $V$ gives a quotient
\[
 G_{L^+}\twoheadrightarrow\Gamma,\qquad
 G_L\twoheadrightarrow\Gamma_0=\ker(\eta),\qquad
 \Gamma'=\ker(\varepsilon|_{\Gamma_0}),
 \quad \eta:\Gamma\longrightarrow\{\pm1\}.
 \tag{QX.1}
\]
The maps commute with the $\ell$-adic cyclotomic character $\chi$.
The chosen complex conjugation $c\in\Gamma$ satisfies $c^2=1$ and
$\eta(c)=-1$.  For a representation $\rho:\Gamma_0\to
\operatorname{GL}_n(\overline{\mathbb Q}_\ell)$ the source's
\texttt{Type\allowbreak I} is irreducibility of $\rho|_{\Gamma'}$, over the
algebraically closed coefficient field.  It is not a residual-image
hypothesis.

The source's \texttt{Is\allowbreak Aut\allowbreak Csd\allowbreak A}$\,\iota\rho$ includes an actual regular
algebraic conjugate-self-dual cuspidal representation $\pi$ of
$\operatorname{GL}_n(\mathbb A_L)$ and its \texttt{Att\allowbreak C} attachment.
At an unramified place of norm $q$, this attachment compares the
characteristic polynomial of geometric Frobenius with
$q^{(n-1)/2}$ times the unitary Satake roots.  We keep this convention
throughout.  The extra \texttt{Wit\allowbreak A} condition will matter in the
automorphic descent below, but is unnecessary in QX1.

We use Thorne, \emph{The vanishing of adjoint Bloch--Kato Selmer groups
of irreducible automorphic Galois representations}, Theorem 6.1,
\url{https://arxiv.org/abs/2207.04925}.  Its precise statement is:
for a CM field $L$, $n\ge2$, a regular algebraic cuspidal polarized
automorphic pair $(\pi,\xi)$, and an isomorphism
$\iota:\overline{\mathbb Q}_\ell\simeq\mathbb C$, irreducibility of
$r_{\pi,\iota}|_{G_{L(\mu_{\ell^\infty})}}$ implies
$H^1_f(L^+,\operatorname{ad}r_{\pi,\iota})=0$.
Here the adjoint is the polarized, or unitary, adjoint over $L^+$.
There is no condition on the residual image and no lower bound on
$\ell$.  The Galois-representation normalization is that of
Barnet-Lamb--Gee--Geraghty--Taylor, \emph{Potential automorphy and
change of weight}, Theorem 2.1.1,
\url{https://annals.math.princeton.edu/wp-content/uploads/annals-v179-n2-p03-p.pdf}.
Both are used with their hypotheses, rather than as an assertion that
all ordinary points are automorphic.

% BEGIN QX1
\begin{proposition}[QX1: the complete adjoint-lifting input]
\label{claim:QX1}
Let $n\ge2$, and let $\rho:\Gamma_0\to
\operatorname{GL}_n(\overline{\mathbb Q}_\ell)$ be continuous.
Suppose that $\rho\circ\texttt{pr}$ satisfies
\texttt{Is\allowbreak Aut\allowbreak Csd\allowbreak A}$\,\iota$ and that $\rho|_{\Gamma'}$ is
irreducible.  Then the literal \texttt{Ad\allowbreak Zero}$\,V\rho$ holds.
Consequently the source's entire \texttt{Adjoint\allowbreak Lifting\allowbreak Theorem}
is supplied, for every one of its original level, residual, point,
coefficient and embedding parameters.
\end{proposition}
\begin{proof}
First, the surjection $G_L\to\Gamma_0$ maps
$G_{L(\mu_{\ell^\infty})}$ \emph{onto} $\Gamma'$: lift an element
of $\Gamma'$, and use cyclotomic compatibility to see that its lift
has $\chi=1$.  Thus the irreducibility hypothesis is precisely the
one required by Thorne's theorem.  In particular $\rho$ and its
inflation to $G_L$ are absolutely irreducible.

Choose the actual $\pi$ in \texttt{Is\allowbreak Aut\allowbreak Csd\allowbreak A}.  The quadratic
Hecke character $\delta_{L/L^+}$ is trivial after pullback by the
norm from $L$, so $\pi$, together with
$\delta_{L/L^+}^{\,n}$, is a polarized automorphic pair.  At each
real place this character has value $(-1)^n$ on $-1$, the required
parity.  The Galois multiplier is
\[
       \nu=\delta_{L/L^+}^{\,n}\chi^{\,1-n};
       \qquad \nu(c)=-1.
       \tag{QX.2}
\]
The displayed exponent and sign are also stated in Thorne's
\S6.3.  The \texttt{Att\allowbreak C} Frobenius convention above is exactly
the normalization of $r_{\pi,\iota}$.  Chebotarev and the
semisimple-character criterion identify $\rho\circ\texttt{pr}$
with $r_{\pi,\iota}$: their characteristic polynomials agree at
all but finitely many unramified places, hence on the whole compact
Galois group by continuity.  Both representations are semisimple.

Now take \emph{any} $P$ and $R$ occurring in \texttt{Ad\allowbreak Zero}.
The exact form and action identities are
\[
 \rho(g)^tP\rho(cgc^{-1})=\chi(g)^{1-n}P,
 \quad R(g)X=\rho(g)X\rho(g)^{-1},
 \quad R(c)X=-P^{-1}X^tP.
 \tag{QX.3}
\]
The first identity is the source's \texttt{Is\allowbreak Form}, and the latter
two are \texttt{Is\allowbreak Ad\allowbreak Ext}.  A standard polarization of
$r_{\pi,\iota}$ gives another intertwiner $P_\pi$ in the first
identity.  Absolute irreducibility and Schur's lemma give
$P=aP_\pi$ for a scalar $a\ne0$.  This scalar cancels in the last
formula.  Therefore, after inflation to $G_{L^+}$, the given $R$
is exactly the polarized adjoint in Thorne's theorem, including
the minus sign at complex conjugation.  No smaller symplectic Lie
algebra or unpolarized adjoint is substituted for it.

All these assertions can be made over one finite extension
$E/\mathbb Q_\ell$.  Here is the coefficient point explicitly.
The finite subextensions of $\overline{\mathbb Q}_\ell$ have a
countable cofinal family.  The inverse images of
$\operatorname{GL}_N(E_j)$ under a continuous representation are
closed subgroups covering the compact source group.  Baire gives
one with interior; adjoining the finitely many matrices of coset
representatives then captures the entire image over a finite
extension.  Apply this to the representation and, for a continuous
cocycle $z$, to its affine block representation.  Adjoin the entries
of $P$ and of the chosen intertwiner.  The action $R$ is continuous
on the two cosets of $\Gamma_0$ by (QX.3).  This is the same finite
coefficient capture used in RV4, with no completion of the
algebraic closure assumed.

The hypothesis \texttt{Af.\allowbreak Is\allowbreak BK} on the inflated $z$ is exactly
membership in $H^1_f(L^+,R)$.  Away from $\ell$ its stated local
primitive is the vanishing of the inertia restriction.  Above
$\ell$, KB6 identifies the original decomposition group with the
absolute Galois group of the completion, and BX1--BX4 identify
\texttt{finite\allowbreak Cocycles} with the kernel after tensoring with
$B_{\rm cris}$.  These comparisons commute with the original
restriction maps and with finite coefficient extension.  Descent
back to $E$ is harmless: a finite extension that contains any
needed primitive admits an $E$-linear projection taking $1$ to $1$;
equivalently, the local finite-condition maps commute with finite
flat coefficient extension.  Thus the inflated class belongs to
the exact Selmer group in Thorne's theorem and is zero.

Choose a matrix $X$ giving its coboundary.  Surjectivity
$G_{L^+}\to\Gamma$ gives
$z(g)=R(g)X-X$ for every $g\in\Gamma$, which is precisely
\texttt{Ad\allowbreak Zero}.  The literal \texttt{Adjoint\allowbreak Lifting\allowbreak Theorem}
quantifies only over $2\le n\le4$; its point record already supplies
continuity, \texttt{Is\allowbreak Aut\allowbreak Csd\allowbreak A} and \texttt{Type\allowbreak I}.  Its additional
\texttt{Par}, \texttt{Free}, \texttt{N3}, residual and level
conditions need not be strengthened or used.  There are no rank-zero
or rank-one cases hidden in that literal statement.
\end{proof}
% END QX1

\subsection{Fixed-weight tangent vanishing and local support}

An ordinary determinant ring $B$ below is the actual
\texttt{Universal\allowbreak Ordinary\allowbreak Determinant\allowbreak Ring} of the source.  Its
universal datum consists of a degree-$n$ determinant law $D_B$ and
ordered local characters $\psi_{v,i}$.  It retains the residual
determinant, local residual characters and prescribed torsion.
Its \texttt{P} clause equates the characteristic coefficients of
each group element and its polarized star; its \texttt{Ord} and
\texttt{I} clauses are trace identities after left multiplication
by every group element.  These are the weaker universal
\texttt{Pair} relations of KX8.  They are not being replaced by
the full scalar-extended relations in \texttt{Ring\allowbreak Props}.  A weight object
$\Lambda$ represents exactly the allowed restrictions of the
$\psi_{v,i}$ to inertia.  KX supplies
\[
       \Lambda\simeq O[[T_1,\ldots,T_N]],
       \qquad N=n[L^+:\mathbb Q].
       \tag{QX.4}
\]
Thus fixing a homomorphism from $\Lambda$ fixes all those ordered
inertial characters, not just their product.

% BEGIN QX2
\begin{lemma}[QX2: a complete commutative-algebra support criterion]
\label{claim:QX2}
Let $B$ be a complete Noetherian local $O$-algebra with the original
finite residue field, $f:\Lambda=O[[T_1,\ldots,T_N]]\to B$, and
$J\subset B$.  Assume that $B/J$ is finite and torsion-free as a
$\Lambda$-module.  Let $x:B\to\overline{\mathbb Q}_\ell$ be an
$O$-algebra map killing $J$.  If every lift
$x':B\to\overline{\mathbb Q}_\ell[\epsilon]/(\epsilon^2)$ of $x$
whose derivative vanishes on $f(\Lambda)$ has zero derivative on
$B$, then
\[
       J\subset\mathfrak q
       \quad\hbox{for every prime }
       \mathfrak q\subset\ker x.
       \tag{QX.5}
\]
\end{lemma}
\begin{proof}
Put $\mathfrak p=\ker x$ and
$\mathfrak a=\ker(xf)$.  These characteristic-zero points have
residue fields finite over $\operatorname{Frac}O$.  To check the
point without a hidden continuity assumption, present $B$ as a
quotient of $O[[X_1,\ldots,X_s]]$.  Each $x(X_i)$ is algebraic and
has absolute value less than one: otherwise a polynomial relation
with unit constant term would give a unit in the maximal ideal of
the power-series ring.  A distinguished polynomial annihilates
each $x(X_i)$, and division by these polynomials reduces the entire
quotient to a finite $O$-module.  The map is consequently continuous
and has image in a finite extension.  Squaring these polynomial
relations proves the corresponding finite coefficient statement
for a dual-number lift.  This is the point comparison of DX5.

The regular local ring $\Lambda_{\mathfrak a}$ has dimension $N$.
The tangent hypothesis makes restriction injective on derivations
at the point.  Since its residue field is separable over
$\operatorname{Frac}O$, the tangent space of
$\Lambda_{\mathfrak a}$ has dimension $N$.  The cotangent space
of $B_{\mathfrak p}$ therefore has dimension at most $N$.

On the other hand $(B/J)_{\mathfrak p}$ has dimension at least
$N$.  Indeed every minimal prime of the finite
$\Lambda$-algebra $B/J$ contracts to zero: minimal primes are
associated primes, and a nonzero element of $\Lambda$ is a
nonzerodivisor by torsion-freeness.  On a component through the
point we therefore have a finite integral extension of the domain
$\Lambda$.  Going down over the integrally closed ring $\Lambda$
lifts a full prime chain under $\mathfrak a$ to one ending at the
chosen point.  Its length is $N$.

Dimension is bounded by embedding dimension.  Hence
$B_{\mathfrak p}$ is a regular local ring of dimension $N$, and
its quotient by $J_{\mathfrak p}$ has the same dimension.  A regular
local ring is a domain; every nonzero ideal in it lowers dimension.
Thus $J_{\mathfrak p}=0$.  For $a\in J$ choose
$s\notin\mathfrak p$ with $sa=0$.  If
$\mathfrak q\subset\mathfrak p$ is prime, $s\notin\mathfrak q$
and therefore $a\in\mathfrak q$.  This proves (QX.5), including
the case $N=0$.
\end{proof}
% END QX2

% BEGIN QX3
\begin{proposition}[QX3: the actual characteristic-zero support leaf]
\label{claim:QX3}
Let $B,\Lambda$ be actual ordinary objects for a level $V$ and
residual datum $\bar r$ with \texttt{V0ok}.  Let $J$ have the
finite and torsion-free properties in QX2.  Suppose that an actual
point $x:B\to\overline{\mathbb Q}_\ell$ kills $J$ and is represented
by a continuous $\rho:G_L\to\operatorname{GL}_n
(\overline{\mathbb Q}_\ell)$ with the ordered flags
$\psi_{v,i}=x(\psi^B_{v,i})$.  Assume $n\ge2$,
\texttt{Is\allowbreak Aut\allowbreak Csd\allowbreak A}$\,\iota\rho$, cyclotomic irreducibility,
unramifiedness at the auxiliary place, and integers $m_{v,i}$ and
$k>0$ such that
\[
 m_{v,i}\ge m_{v,j}+2\ (i<j),\qquad
       (\psi_{v,i}\chi^{-m_{v,i}})^k=1
                 \quad\hbox{on }I_v.
       \tag{QX.6}
\]
Then every prime below $\ker x$ contains $J$.
In particular this supplies the entire fixed-weight support
conclusion of the source's \texttt{eigenbasis\_\allowbreak rank\_\allowbreak four}, without
requiring its unused residual self-twist and large-order conditions.
\end{proposition}
\begin{proof}
We prove the tangent hypothesis of QX2.  Take a dual-number lift
$x'$ constant on $\Lambda$.  At the absolutely irreducible residual
point, its determinant law is the determinant of a matrix
representation $\rho'$ lifting $\rho$, unique up to conjugacy.
This is the split absolutely irreducible case of Chenevier,
\emph{The $p$-adic analytic space of pseudocharacters of a profinite
group and pseudorepresentations over arbitrary rings}, Theorem 2.22,
\url{https://gaetan.chenevier.perso.math.cnrs.fr/articles/determinants.pdf}:
first pass to the Cayley--Hamilton quotient, then lift matrix units
over the henselian ring $E[\epsilon]/\epsilon^2$.  The resulting
algebra is $M_n(E[\epsilon]/\epsilon^2)$.  Continuity is preserved:
choose finitely many group matrices forming a basis modulo
$\epsilon$; the trace pairing expresses the coordinates of every
$\rho'(g)$ through the continuous determinant traces at finitely
many products of group elements.  QX2 permits a single finite $E$
for this argument.

Write $\rho'(g)=(1+\epsilon y(g))\rho(g)$.  Burnside's theorem
makes the group matrices span $M_n(E)$ at the given irreducible
point.  Nakayama then makes their lifts span
$M_n(E[\epsilon]/\epsilon^2)$; finitely many already suffice.
The trace pairing on this matrix algebra is nondegenerate, as
the matrix units show.  Thus the actual universal equations
$\operatorname{tr}(\rho'(h)Q)=0$ for all $h$ force $Q=0$,
both for an ordinary product and for each prescribed inertial
matrix $\rho'(g)-1$.  These arguments derive the strong matrix
identities at this point from the weak \texttt{Pair} relations.
The ordinary identities give the ordered flags of $\rho'$
with characters $x'(\psi^B_{v,i})$.  For completeness,
on inertia kill the finitely many characters in (QX.6), and choose
an element whose cyclotomic value has infinite order.  Its ordered
eigenvalues are distinct.  Lift their spectral projectors to the
dual numbers.  In the resulting basis, the ordinary product
identities, with this element in all but one factor, force each
entry below the prescribed flag to vanish.  This gives an actual
flag over the dual numbers, not just a factorization of its
characteristic polynomial.

The \texttt{P} identities equate the characteristic polynomials,
and hence the characters, of $\rho'$ and its polarized conjugate
dual.  In the same matrix algebra the trace pairing makes equality
of their characters an intertwining isomorphism.  Lifting its
residual matrix gives an invertible form $P'$ for $\rho'$; this is
also the uniqueness assertion in the split Cayley--Hamilton
construction.  The
first derivative of this identity says that the class of $y$ is
fixed by the polarized complex-conjugation action.  Since $2$ is
invertible in $E$, inflation--restriction for
$\Gamma/\Gamma_0$ identifies these fixed classes with
$H^1(\Gamma,R)$, where $R$ has the action (QX.3).  Thus there is a
continuous cocycle $z$ extending the class of $y$.  Changing $y$
by a coboundary, if necessary, makes this an equality on
$\Gamma_0$.

We verify every local finite condition.  At places over $\ell$,
constancy on $\Lambda$ fixes all the inertial diagonal characters.
After a finite local extension the finite-order factors in
(QX.6) disappear.  The successive strictly upper-triangular
characters are $\chi^{m_{v,i}-m_{v,j}}$ times unramified
characters, with $m_{v,i}-m_{v,j}\ge2$.  Successive extensions of
the ordered characters make the base flag crystalline: at each
step its extension module has all these positive differences,
so local duality and the Bloch--Kato exponential give
$H^2=0$ and $H^1=H^1_f$.  This is an induction on the length of
the flag, beginning with the crystalline rank-one characters.

Here is the dimension argument for the tangent, which avoids
assuming surjectivity for a finite local condition.  Let
$\mathfrak n$ be the strictly upper-triangular adjoint module,
and $\mathfrak b$ the upper-triangular one.  They are crystalline
subquotients of the adjoint of this crystalline flag, and
\[
 0\longrightarrow\mathfrak n\longrightarrow\mathfrak b
       \longrightarrow E^n\longrightarrow0.
 \tag{QX.7}
\]
Every graded character of $\mathfrak n$ has exponent at least
two.  Consequently $H^0(\mathfrak n)=H^2(\mathfrak n)=0$ and
$H^1(\mathfrak n)=H^1_f(\mathfrak n)$, of dimension
$\dim_E t_{\mathfrak n}$, where
$t_M=D_{\rm dR}(M)/\operatorname{Fil}^0$.
The long exact sequence of (QX.7) shows that the inverse image
of $H^1_{\rm unr}(E^n)$ in $H^1(\mathfrak b)$ has dimension
\[
 \dim_E t_{\mathfrak n}
       -(n-h^0(\mathfrak b))+n
       =\dim_E t_{\mathfrak n}+h^0(\mathfrak b).
 \tag{QX.8}
\]
Functoriality puts $H^1_f(\mathfrak b)$ inside that inverse
image.  The Bloch--Kato dimension formula gives its dimension
as $\dim_E t_{\mathfrak b}+h^0(\mathfrak b)$, and
$t_{\mathfrak b}=t_{\mathfrak n}$ since the quotient has weight
zero.  Thus the two spaces coincide.  Constancy on $\Lambda$
puts the tangent class in this inverse image.  Its image in
the full adjoint is therefore finite.  This proves the local
claim with exactly the gaps in (QX.6).  Restriction followed by
corestriction descends it from the finite local extension.

At a non-$\ell$ place in the prescribed level, the universal
inertial relations make $\rho'$ trivial on the stated finite-index
subgroup $I_p\cap\ker\bar r$.  The cocycle is therefore zero there;
averaging over the remaining finite quotient makes its class
unramified.  Outside the level this is built into $\Gamma$.
At the auxiliary place $v_0$, the base representation is
unramified.  Choose a stable lattice over a sufficiently large
finite coefficient field.  Its Frobenius characteristic
polynomial reduces, by the determinant's residual identity,
to $(T-a)^n$ for one nonzero $a$; the residual matrix need not
itself be scalar.  All the residual eigenvalue ratios of its
adjoint are nevertheless $1$.  Moreover
$q_{v_0}\not\equiv1\pmod\ell$ by \texttt{V0ok}.  Hence
$\operatorname{Ad}(\rho(\operatorname{Frob}_{v_0}))-q_{v_0}$
is invertible: its determinant reduces to
$(1-q_{v_0})^{n^2}$.  The tame relation forces the infinitesimal
inertia value to be zero; the wild pro-$p$ subgroup has no
continuous additive image in an $\ell$-adic vector space for
$p\ne\ell$.  This proves the required unramified condition at
$v_0$ as well.  KB6 and BX transport all these checks to the
original valuations and \texttt{Af.\allowbreak Is\allowbreak BK} cochain convention.

QX1 now makes $z$, and hence $y$, a coboundary.  Thus $\rho'$ is
conjugate to the constant lift of $\rho$.  Its ordered diagonal
characters are constant too: the distinct inertia projectors
fix their order, so conjugacy cannot permute them.  The determinant
and every ordered local character of $x'$ agree with the constant
ones.  We do not apply the source's universal property to a
dual-number field, which is not an object of $\mathcal C_O$.
Instead KX7--KX8, or Chenevier's Remark 3.5 followed by the
ordinary relation quotient, shows that $B$ is topologically
generated by the determinant coefficients and the ordered
local character values.  Both point maps are continuous by
the finite coefficient argument in QX2.  Equality on these
generators makes $x'$ itself constant.
QX2 proves the result.  Notice that the proof required membership
$x(J)=0$; it does not create a modular quotient or an automorphic
point on one.
\end{proof}
% END QX3

\subsection{Two actual links, with one coefficient threshold}

Fix an actual $\texttt{Fit}+\texttt{Deep}$ tower $T$ for the
original $D$.  Thus $F/\mathbb Q$ is totally real and abelian,
$K/\mathbb Q$ is imaginary quadratic, $L=FK$, and $L/F$ is
unramified at finite places.  The primes in the original $S$,
including $2$ and $\ell$, split in $K$.  The residual characters
and the chosen $\ell$-adic decomposition groups are the ones
in $T$.  In particular \texttt{Deep} says that
$\psi|_{G_F}$ has finite order prime to $\ell$.
Write $\theta_i=\chi_i(D,T.\texttt{pr})$, so their labels include
the source's case-dependent permutation.

% BEGIN QX4
\begin{proposition}[QX4: construction of the two literal links]
\label{claim:QX4}
There is a coefficient ring $O_1$ such that, for every permitted
$O'\supset O,O_1$, there are two literal objects
$\mathcal L,\mathcal L':\texttt{Link}(T,O')$ with identical
characters $\kappa,\beta,t$ and with different auxiliary
rational primes.  They use the original tower, its original
places and its original residual-character labels.  The choice
of $O_1$ is independent of any later \texttt{Space}.
\end{proposition}
\begin{proof}
We give the character construction first, before choosing either
auxiliary prime.  Write $\omega$ for the Teichmuller part of
$\chi$.  The \texttt{Ctx2}/\texttt{Ctx3} and \texttt{Good}
identities make
\[
       \eta=(\omega^{e+3}\xi\psi_{\ell'})^{-1}
       \tag{QX.9}
\]
an even finite character of $G_{\mathbb Q}$ of order prime to
$\ell$, unramified outside $S$.  Here $\psi_{\ell'}$ is the
prime-to-$\ell$ part of $\psi$, obtained by an integer exponent
which is zero on its $\ell$-primary part and one on its
prime-to-$\ell$ part.  Deepness gives
$\psi_{\ell'}=\psi$ on $G_L$.

We need a finite character $\kappa_0$ of $G_K$ with
$\kappa_0\kappa_0^c=\eta|_{G_K}$ and the same permitted
ramification.  This has a direct ray-class construction.
Represent $\eta$ by an even Dirichlet character modulo $M$
supported on its ramified primes.  Choose, over each such prime,
one of the two primes of $K$; the splitting and Chinese
remainder theorem give $\mathcal O_K\to\mathbb Z/M\mathbb Z$.
For a principal ideal $(a)$ prime to $M$, evaluate the
Dirichlet character on this image of $a$.  The value is
independent of its generator: $2$ splits in $K$, so the units
of $\mathcal O_K$ are $\{\pm1\}$, and $\eta$ is even.
This character on principal ideals extends to the group of
ideals prime to $M$ with values in the divisible group of
prime-to-$\ell$ roots of unity.  It is trivial on principal
ideals congruent to $1$ modulo $M$, so it factors through the
finite ray class group.  Class field theory gives $\kappa_0$.
For every ideal $\mathfrak a$, the equality
$\mathfrak a\overline{\mathfrak a}=(N\mathfrak a)$ proves the
norm identity on Frobenius elements, hence everywhere by
Chebotarev.  Inverting the character if the opposite Artin
normalization is used gives the stated identity in the source's
Frobenius convention.

The unique square root in $1+\ell\mathbb Z_\ell$ gives the
continuous character
$\gamma=(\chi/\omega)^{1/2}$.  It is trivial on $\ker\chi$.
Next use the chosen prime of $K$ over $\ell$.  If $h_K$ is the
class number, a principal generator of $\mathfrak a^{h_K}$
for ideals away from $\ell$ gives, under the corresponding
embedding $K\hookrightarrow\mathbb Q_\ell$, the ray-class
character with infinity exponent $h_K(\ell-1)$ at that
embedding and zero at its conjugate.  Passing to the inverse
limit of the $\ell$-power ray class groups gives a continuous
$\mu:G_K\to1+\ell\mathbb Z_\ell$ and an integer $d>0$ with
\[
 \mu\mu^c=\chi^d,\qquad
 \mu|_{I_{\mathfrak l}}=\chi^d,\qquad
 \mu|_{I_{\bar{\mathfrak l}}}=1 .
 \tag{QX.10}
\]
One can take a positive multiple of $h_K(\ell-1)$ for $d$.
The ideal description proves these identities on principal
ideles; finite class field theory gives them on the inverse
limit.  This is a construction of the character, not an
assumption of a $\mathbb Z_\ell$-extension with additional
local behavior.

Set $\pi_0=\mu^2\chi^{-d}$.  Then $\pi_0^c=\pi_0^{-1}$, its
image lies in $1+\ell\mathbb Z_\ell$, and it is unramified
outside $\ell$.  There is a continuous additive
$t:G_K\to\mathbb Z_\ell$ with $\pi_0=(1+\ell)^t$.
Choose $a\in\overline{\mathbb Q}_\ell$ with
$a^{2d}=1+\ell$ and $|a-1|_\ell<1$, and put $\beta=a^t$.
The continuous $\mathbb Z_\ell$-power is defined as the limit
of integral powers; $a^{\ell^r}\to1$ proves that it exists.
It has image in a single finite coefficient field and satisfies
\[
 \beta^{2d}=\pi_0,\quad \beta^c=\beta^{-1},\quad
 \ker t\subset\ker\beta,\quad \bar\beta=1,\quad
 (\beta^2\chi^{-b_v})^d=1\ \hbox{on }I_v,\quad b_v\in\{1,-1\}.
 \tag{QX.11}
\]
In particular the exponents $b_v$ are odd.
Restrict these characters to $G_L$ and define
\[
        \kappa=\kappa_0\,\gamma^{-(e+4)}
                              \beta^{\,e\bmod2}.
        \tag{QX.12}
\]
Using (QX.9), $\gamma^2\omega=\chi$, and Deepness, direct
multiplication gives
\[
 \kappa(g)\kappa(cgc^{-1})\,\vartheta(g)\chi(g)^3=1,
 \quad \beta(g)\beta(cgc^{-1})=1,\quad
                         t(cgc^{-1})=-t(g).
 \tag{QX.13}
\]
At each place over $\ell$, (QX.11) shows that a positive power
of $\kappa\chi^{-a_v}$ is trivial for an \emph{integer} $a_v$:
the numerator $-(e+4)+b_v(e\bmod2)$ is even.
Away from $\ell$, $t$ and $\beta$ are unramified and
$\kappa$ has bounded prime-to-$\ell$ inertial order.
All three characters are trivial on the required
$\ell$-primary torsion of the local abelianization.  For
$\kappa_0$ this follows from its prime-to-$\ell$ order;
for $\gamma$ and $\beta$ it follows from their torsion-free
one-parameter images.  This also proves the exact
\texttt{mu} clause for the chosen decomposition groups.

Choose $O_1$ containing the finite character values, $a$, and
all embeddings of $\mathcal O_L$ into
$\overline{\mathbb Q}_\ell$.  The latter is a finite requirement
because $\mathcal O_L$ is finitely generated over $\mathbb Z$.
All character images then lie in $O_1^\times$.  Their
continuity is adic continuity over every $O'\supset O_1$.
This verifies the \texttt{large} field before choosing any
space or auxiliary prime.

We construct the auxiliary place.  The \texttt{g2} condition
in $\texttt{Fit}$ says that the residual cyclotomic character
is not in the subgroup generated by the two torus-root
characters.  Every ratio $\bar\theta_i/\bar\theta_j$ belongs
to that subgroup, using
$\theta_0\theta_3=\theta_1\theta_2$.
Duality for the finite abelian image of these characters gives
an element $g_0\in G_L$ with all $\bar\theta_i(g_0)$ equal
and $\bar\chi(g_0)\ne1$.  Indeed, failure of this conclusion
would put $\bar\chi$ in the annihilator of the common kernel,
which is exactly the subgroup generated by those ratios.

Choose an open normal subgroup on which these residual
characters are trivial.  Chebotarev, applied to its finite
Galois quotient and the class of $g_0$, gives a place $v_0$
of $F$ of residue degree one over a rational prime $p$,
split in $L$, with this Frobenius class.  The prime can avoid
any preassigned finite set.  Take that set to include the
original level, the discriminants, $2,\ell$ and the
ramification of the finite characters.  Then the residual
characters coincide on the whole decomposition group at
$v_0$, inertia is killed there, and $p\not\equiv1\pmod\ell$.
The last condition also proves $\mu_\ell\not\subset L$.

Take $S_L$ to be the places of $F$ over the original $S$
together with $v_0$, with the source's specified
$\ell$-adic places.  All are split in $L$; $[F:\mathbb Q]$
is even by Fit.  These are exactly the arithmetic conditions
of the \texttt{Hida\allowbreak Setting}.  KX2 supplies its level with
$V.S'=S_L$.  The characters
\[
 \bar r_i=\overline{\kappa\,\theta_i|_{G_L}},
 \qquad \bar\psi_{v,i}=\bar r_i|_{G_v},\qquad
                        \texttt{th}_{v,i}=1
 \tag{QX.14}
\]
descend to that quotient.  They have open kernels and remain
pairwise distinct by Fit; (QX.13) gives their determinant
polarization.  The previous paragraph is exactly
$\bar r.\texttt{V0ok}$.  All fields of \texttt{Link} are now
accounted for.

Repeat only the Chebotarev choice, adding $p$ to the excluded
set.  This gives $p'\ne p$ and hence two links with the
\emph{same} $\kappa,\beta,t$ and the same $O_1$.
No embedding, original residual label, or original
$\ell$-adic place has changed.
\end{proof}
% END QX4

% BEGIN QX5
\begin{proposition}[QX5: the actual determinant comparison maps]
\label{claim:QX5}
For either link $\mathcal L$ in QX4, every actual
$X:\texttt{Space}(T.\texttt{fr}\,\varnothing,O\to O')$ and every
integer $a\ge0$, there is an actual ordinary pair over $X.B$.
It supplies a universal ordinary determinant ring $B$ by KX9.
For \emph{every} such universal ring there is a unique comparison
map $\phi_a:B\to X.B$ satisfying the literal
$\mathcal L.\texttt{Is\allowbreak Phi}\,X\,a\,B\,\phi_a$:
\[
 \begin{split}
 \phi_a\bigl(\det(T-D_B(g))\bigr)
  &=\operatorname{scaleRoots}_{\kappa(g)\beta(g)^a}
                    \bigl(X.\operatorname{cp}(g|_F)\bigr),\\
 \phi_a(\psi^B_{v,i}(g))
  &=\kappa(g)\beta(g)^a\,X.\psi_{v,i}(g|_F).
 \end{split}
 \tag{QX.15}
\]
The first line denotes the characteristic polynomial of the
universal determinant at $g$, rather than a matrix realization
of $D_B$.  Both maps land in the original $X.B$, with its
specified injection into $X.A$.
\end{proposition}
\begin{proof}
In $X.A$ take the actual matrices
$\rho_a(g)=\kappa(g)\beta(g)^a X.P.r(g|_F)$.
Their characteristic coefficients lie in $X.B$ by the first
formula in (QX.15).  They are adically continuous and reduce
to (QX.14).  We first explain why this gives a determinant
\emph{law} over $X.B$, even in residue characteristic three.

For finitely many matrices $A_1,\ldots,A_s$, the integral
Amitsur identity expresses every coefficient of
$\det(1-\sum Z_iA_i)$ as an integer polynomial in the
characteristic coefficients of products of the $A_i$.
One form of the identity is the formal product of
$\det(1-Z_wA_w)$ over primitive cyclic words, in Lyndon
order; only finitely many words enter any fixed degree.
It follows by the unique nonincreasing Lyndon factorization
of a word and multiplicativity of determinant.  In the
present application each product $A_w$ is another group
matrix.  Consequently all these mixed coefficients belong
to $X.B$.  They define the determinant law there, and the
injection $X.B\hookrightarrow X.A$ verifies its homogeneity,
multiplicativity and all base-change identities.  This is
also the integral determinant descent in Chenevier,
Lemma 1.12 and Corollary 1.14.  It uses no division by $4!$.

Descent from $G_L$ to the level quotient needs a small
topological precaution.  The ring $X.A$ is only finite type
over $X.B$; its adic topology need not be separated, so we
do not infer that the kernel of the matrix representation
is closed.  Instead take the subgroup consisting of those
$u$ for which all characteristic coefficients at $gu$
equal those at $g$, for every $g$, and the corresponding
mixed coefficients do also.  The just-proved word identities
make this the normal kernel of the determinant law.
It is closed because its coefficients take values in the
separated ring $X.B$ and are adically continuous.  It
contains every inertia subgroup killed by $X.P.r$ and
$\kappa\beta^a$, so it contains their closed normal closure,
which is the kernel of $G_L\twoheadrightarrow\Gamma_0$
in KX2.  The determinant therefore descends to $\Gamma_0$.
Its continuity descends through this quotient map.

For the ordered local characters, the upper-triangular
matrices given by $X.P.n_v$ show directly that
$X.\psi_{v,i}(gh)=X.\psi_{v,i}(g)X.\psi_{v,i}(h)$ in $X.A$,
hence in $X.B$ by injectivity.  Their residual values are
nonzero, so they are units of the local ring $X.B$.
Multiplying by $\kappa\beta^a$ gives the second line of
(QX.15), including continuity, residual values and the
prescribed local torsion.

We check the remaining weak \texttt{Pair} relations.
The symplectic multiplier of $X.P.r$ and (QX.13) identify
the characteristic polynomial at a group element with
its polarized conjugate dual; $\beta\beta^c=1$ makes this
independent of $a$.  This is exactly \texttt{P}.
The actual ordered upper-triangular matrices make the
ordinary product zero, so multiplying by any group matrix
and taking its trace gives \texttt{Ord}.

For \texttt{I}, let $g$ be in one of the specified non-$\ell$
inertias and in all residual-character kernels (QX.14).
Choose an integer $m$ prime to $\ell$ killing
$\kappa$ on this inertia.  Here $\beta=1$.  The element
$g^m$ has original residual matrix $1$, so the defining
kernel condition of $T.\texttt{fr}\,\varnothing$ gives
$\rho_a(g)^m=1$.  Its characteristic polynomial, with
coefficients in $X.B$, reduces to $(T-1)^4$.
This already forces $\rho_a(g)=1$, including over $X.A$.
Indeed write $T^m-1=(T-1)Q(T)$, with $Q(1)=m$ a unit.
The resultant of $Q$ and that characteristic polynomial is
a unit of $X.B$, by reduction modulo its maximal ideal.
Bezout and Cayley--Hamilton make $Q(\rho_a(g))$ invertible;
multiplication by $(\rho_a(g)-1)$ then gives the result.
Its trace version is the required \texttt{I} relation.

We have constructed an actual \texttt{Pair} over the ring
$X.B\in\mathcal C_{O'}$.  Thus the nonempty-pair antecedent
in KX9 is genuinely satisfied, rather than discarded.
The universal property gives $\phi_a$, and evaluating
the determinant and characters gives precisely (QX.15).
Conversely those formulas determine their full polynomial
law by the word identity and determine every ordered
character.  Universality in the actual category
$\mathcal C_{O'}$ proves uniqueness.  Every construction
works for all $a\ge0$ and both links on the same space.
\end{proof}
% END QX5

\subsection{The ordinary ideal properties that are actually used}

For each link we now suppose an ideal $J_0\subset B$ has the
literal rank-four \texttt{Hida\allowbreak Props\allowbreak Alt} fields.  We list the
fields to keep the boundary explicit:
\begin{enumerate}
\item \texttt{ring}: $B/J_0$ is finite and torsion-free over
$\Lambda$, with the specified polarization, ordinary,
inertial-type and auxiliary-unramified identities.
\item \texttt{dense}: for every strictly increasing integral
weight, the source's good arithmetic points killing $J_0$
separate $B/J_0$.
\item \texttt{pt}: an actual continuous attached regular
algebraic conjugate-self-dual cuspidal representation with
\texttt{Wit\allowbreak A}, the stated auxiliary unramifiedness and
ordered local generic flags belongs to this quotient.
\item \texttt{end1}: every arithmetic point on the quotient
has the exact semisimple decomposition into consecutive
cyclotomic strings, and the stated strict Frobenius-root
bound when those strings all have length one.
\item \texttt{end2}: every irreducible arithmetic point is
actually conjugate-self-dual cuspidal and has
\texttt{Wit\allowbreak Off} away from the auxiliary place.
\item \texttt{twist}: the good arithmetic points are stable
under the specified finite $\ell$-power-order polarized
twists, trivial on the prescribed local kernels.
\end{enumerate}
These are assumptions on an actual ideal in an actual
universal ring; a citation to ordinary automorphic forms
without their exact local types does not supply them.
The separate characteristic-$\ell$ support assertion
\texttt{P1core} will be stated at QX8.  No \texttt{P3core}
field is used in this chapter.

% BEGIN QX6
\begin{proposition}[QX6: integral twist invariance of the image ideal]
\label{claim:QX6}
Assume the \texttt{dense} and \texttt{twist} fields above.
For the maps of QX5,
\[
              \phi_a(J_0)X.B=\phi_0(J_0)X.B
                    \qquad(a\ge0).
              \tag{QX.16}
\]
Here equality is equality of ideals, before evaluating any
point and over the original coefficient ring $O'$.
\end{proposition}
\begin{proof}
The character $\beta$ descends to $\Gamma_0$, since it kills
the closed inertia kernel of the level quotient.  It is
conjugate-inverse, is $1$ residually, and is trivial on the
specified non-$\ell$ inertia and the local torsion.
Twisting the universal determinant and every ordered
character by such a character $\gamma$ therefore produces
another actual \texttt{Pair}.  Universality gives a map
$\tau_\gamma:B\to B$ for coefficient-valued $\gamma$,
with inverse $\tau_{\gamma^{-1}}$.  The coefficient of
degree $4-j$ at $g$ is multiplied by $\gamma(g)^j$.

We prove that $\tau_\gamma(J_0)\subset J_0$ without any
incorrect assertion that finite-order characters converge
pointwise to $\gamma$.  Fix $r\ge1$.  The image of $\gamma$
modulo $\mathfrak m_B^r$ is a finite abelian $\ell$-group
$Q$: it lies in $1+\mathfrak m_B$, whose successive
quotients are finite additive groups of characteristic
$\ell$.  Its quotient map $q:\Gamma_0\to Q$ has open
kernel and has all the same polarization and local-kernel
identities.  The group algebra $B[Q]$ is a finite free,
complete local $O'$-algebra with residue field equal to
that of $B$.  Its maximal ideal is generated by
$\mathfrak m_B$ and $[u]-1$, since a finite $\ell$-group
has local group algebra in characteristic $\ell$.
Twist by the tautological character $g\mapsto[q(g)]$.
This gives an actual map $\tau_Q:B\to B[Q]$.

Take $b\in J_0$ and write $\tau_Q(b)=\sum_{u\in Q}b_u[u]$.
For any good arithmetic point $x$ killing $J_0$ and any
character $\lambda:Q\to\overline{\mathbb Q}_\ell^\times$,
the \texttt{twist} field says
\[
                  \sum_{u\in Q}x(b_u)\lambda(u)=0.
                  \tag{QX.17}
\]
These twists have finite $\ell$-power order, preserve the
ordered-character ratios, and satisfy every stipulated
local condition.  Character orthogonality over the
characteristic-zero field makes each $x(b_u)$ zero.
The \texttt{dense} field then gives $b_u\in J_0$ for every
$u$.  Evaluating $[q(g)]$ at $\gamma(g)$ modulo
$\mathfrak m_B^r$ and using uniqueness of the twisted pair
shows
$\tau_\gamma(b)\in J_0+\mathfrak m_B^r$.
Krull separation of the complete Noetherian quotient
$B/J_0$ gives $\tau_\gamma(b)\in J_0$.
The inverse twist gives equality.

Apply this with the coefficient-valued character $\beta^a$.
Formula (QX.15) and QX5 uniqueness give
$\phi_a=\phi_0\tau_{\beta^a}$.  Equality of image ideals
follows.  This argument also proves the universal
$B$-valued twist-stability statement used in the source;
the finite group algebra, rather than a quotient requiring
division by $\ell$, is what makes it integral.
\end{proof}
% END QX6

% BEGIN QX7
\begin{proposition}[QX7: full local support at an actual seed]
\label{claim:QX7}
Under the stated rank-four \texttt{Hida\allowbreak Props\allowbreak Alt} hypotheses,
put $J_{\mathcal L}=\phi_0(J_0)X.B$.  For every actual
$O'$-map $y:X.A\to\overline{\mathbb Q}_\ell$ for which
$T.\texttt{Is\allowbreak Seed\allowbreak Alt}\,\iota(X.P_y)$ holds, one has
\[
                  X.\texttt{Loc\allowbreak Zero}(J_{\mathcal L},y).
                  \tag{QX.18}
\]
This means that every prime of $X.B$ contained in
$\ker(y|_{X.B})$ contains the entire ideal
$J_{\mathcal L}$.
\end{proposition}
\begin{proof}
Unpack the seed as $(w,\kappa_s,\pi,a_s,k_s)$.
Its consecutive $w_i$ differ by even integers at least
four, $w_0+a_s$ is even, and
$(\kappa_s^2\chi^{-a_s})^{k_s}=1$ on every $\ell$-inertia.
Its automorphic representation is an actual regular
algebraic rank-four cusp form with \texttt{All\allowbreak Win}, attached
to $\kappa_s X.P_y.r|_{G_L}$.  The full-matrix-span condition
holds already on $G_L\cap\ker\chi$.

Choose $a\ge0$ with $w_0+a$ even and use $\phi_a$ from QX5.
At each $v\mid\ell$, (QX.11)--(QX.12) give an exponent
$e_v\equiv a\pmod2$ for
$(\kappa\beta^a)^2$.  All the integers
\[
               m_{v,i}=(w_i+e_v)/2
               \tag{QX.19}
\]
therefore exist, and their differences are
$(w_i-w_j)/2$.  A common positive power makes
$(x_a(\psi^B_{v,i})\chi^{-m_{v,i}})$ trivial on inertia,
where $x_a=y|_{X.B}\phi_a$.  To obtain a single power on
all inertia, take the index of the normal core of the
open subgroup in $T.\texttt{Wts}$, and multiply by the
finitely many local exponents.  Thus (QX.6) holds with
all its gaps at least two.

The ratio $\lambda=\kappa\beta^a\kappa_s^{-1}$ is a
continuous algebraic Galois character.  In fact
$e_v-a_s$ is even, and a positive power of
$\lambda\chi^{-(e_v-a_s)/2}$ is trivial on every local
inertia.  Local class field theory makes it locally
algebraic there; global class field theory gives the
corresponding algebraic Hecke character.  This is the
rank-one algebraic-character theorem, with the integer
exponents just checked.  Twisting $\pi$ by that Hecke
character attaches
\[
                 \rho_a=(\kappa\beta^a)X.P_y.r|_{G_L}.
                 \tag{QX.20}
\]
The polarization (QX.13) and strong multiplicity one
identify the twist's conjugate with its contragredient,
so the result is conjugate-self-dual.  The same comparison
uses geometric Frobenius on both sides; scalar twisting
does not introduce an additional power of $\chi$.

We also retain the \texttt{Wit\allowbreak A} part.  An algebraic
Hecke-character twist translates the center of every
strict window by its real norm exponent, preserving
\texttt{All\allowbreak Win}.  Every other regular algebraic cusp
form attached to (QX.20) is the same form by strong
multiplicity one.  Finally (QX.19) implies
\texttt{Pot\allowbreak Unl} at every place over $\ell$: an equality
between a positive power of two ordered characters
with one cyclotomic twist would, on an open inertia
subgroup, force $m_{v,i}-m_{v,j}=1$, which is impossible.
Thus both clauses of \texttt{Wit\allowbreak A} follow, without a
general assertion that arbitrary attached forms have it.

Continuity here follows from the actual point, not from an
unstated clause of the seed.  Absolute irreducibility and
Burnside give finitely many matrices $\rho_a(g_j)$ spanning
the matrix algebra.  The nondegenerate trace pairing recovers
every matrix entry of $\rho_a(g)$ from the functions
$\operatorname{tr}(\rho_a(gg_j))$.  These are continuous
by the characteristic-coefficient continuity of $X$ and
the continuous small-point comparison of QX2; scalar
twisting is continuous as well.  Thus $\rho_a$ is continuous.
Its ordered flags are the actual
flags of $X.P_y$, and its cyclotomic irreducibility is
unchanged by the scalar twist.  At the auxiliary prime
the residual characters coincide, and the same
prime-to-$\ell$ finite-order matrix argument as in QX5
gives the required unramifiedness.  The
\texttt{Point\allowbreak Theorem\allowbreak Base\allowbreak Object} field now says
$x_a(J_0)=0$: the ordinary weights are arithmetic by
local reciprocity, and their gaps exclude every forbidden
ratio $\psi_i=\psi_j\chi$.

Apply QX3 to $x_a$.  Given a prime
$\mathfrak q\subset\ker(y|_{X.B})$, its inverse image
under $\phi_a$ is a prime below $\ker x_a$.
It contains $J_0$, so $\mathfrak q$ contains
$\phi_a(J_0)X.B$.  QX6 identifies this ideal with
$J_{\mathcal L}$ and proves (QX.18).
In particular an actual seed supplies a nonzero point
on the appropriate support; it is not assumed at the
start of the fixed-weight argument.
\end{proof}
% END QX7

% BEGIN QX8
\begin{proposition}[QX8: the complete characteristic-$\ell$ support transport]
\label{claim:QX8}
Assume in addition the literal $B.\texttt{P1core}\,J_0$:
if an actual finite-field power-series point is
\texttt{Nice\allowbreak I}, matches the universal pair, and kills
$J_0$, every prime below its kernel contains $J_0$.
Then $J_{\mathcal L}$ satisfies the entire
\texttt{p1} clause of \texttt{Potential\allowbreak Modularity\allowbreak Tower}
on the original $X$.  No further residual-image or
generic-point hypothesis is required.
\end{proposition}
\begin{proof}
Take the exact test in that clause: a finite field $k$,
an $O'$-map $y:X.A\to k[[t]]$ whose constant term kills
$\mathfrak m_{X.B}$, with $y(J_{\mathcal L})=0$,
the prescribed local character separations on every
$T.U_1(v)$, and full matrix span over $k((t))$.
The constant-term condition makes $O'\to k$ local,
so the given residual field embeds in $k$.
Twisting by $\kappa$ and descending as in QX5 gives
an actual representation over $k[[t]]$ of $\Gamma_0$,
its symmetric polarization and its ordered characters.
Here the matrix descent itself is valid, since
$k[[t]]$ is separated.  Its scalar-extended determinant,
residual characters, inertia and flags give the actual
source \texttt{Point}; its map from $B$ is
$f=y|_{X.B}\phi_0$.

We verify all four \texttt{Nice\allowbreak I} fields rather than
assuming that full span already means \texttt{Nice\allowbreak I}.
First, restriction of the original irreducible
representation from $G_F$ to $G_L$ stays absolutely
irreducible.  Otherwise Clifford theory for this
index-two subgroup would make its trace zero on the
other coset.  Its reduction is the sum of the four
original residual characters, whose restrictions to
$G_L$ are distinct by (QX.14).  On any other coset,
linear independence of those four restricted
characters forbids an identically zero sum with their
nonzero coset values as coefficients.  This contradiction
works in every odd characteristic, including three.

The local separation condition gives, for each ordered
pair $i\ne j$, an element of $\Gamma'$ at which the
residual characters agree but their power-series
characters do not.  The ratio there lies in
$1+tk[[t]]$ and has infinite order; this group is
torsion-free, by the Frobenius map for the $\ell$-part
and the simple-root identity for the prime-to-$\ell$
part.  Hence the kernel of the ratio modulo its finite
constant torsion is closed and not open.
Avoiding finitely many such kernels in the compact
local group gives one element $u\in\Gamma'$ whose
four eigenvalue ratios all have infinite order.
Every positive power of $\rho(u)$ has four distinct
eigenvalues.  In particular local characters are
distinct and no ratio is $\chi$, proving \texttt{N1}.

Suppose restriction to $\Gamma'$ were reducible.
Clifford theory gives a sum of conjugate simple
constituents.  Their multiplicity is one, since a
positive power of $\rho(u)$ still has simple spectrum.
The number of constituents is therefore $2$ or $4$.
The quotient $\Gamma_0/\Gamma'$ is abelian; a
nontrivial quadratic character of the resulting
permutation quotient gives a quadratic self-twist.
Its kernel is open: stabilizers of subspaces are
closed by continuity, and they have finite index.
Because $\ell$ is odd, this quadratic character
of a closed subgroup of $\mathbb Z_\ell^\times$
is trivial on its pro-$\ell$ part, so factors through
the residual cyclotomic character.  Reduction of
the characteristic-polynomial identity permutes
the four residual characters.  This is excluded
by the exact \texttt{g3'} field of Fit, transported
through $\kappa$.  Thus \texttt{Type\allowbreak I} holds.

We give the rank-four regular-coset argument needed
for \texttt{N3}.  Work over an algebraic closure of
$k((t))$, and let $H$ be the Zariski closure of
the image of $\Gamma'$.  No open finite-index
subgroup of $\Gamma'$ can have commuting image.
If one did, its Zariski closure would be
commutative and contain $H^\circ$: finitely
many translates cover $H$, so one of them
contains its identity component.  The group
$H^\circ$ is normal in the Zariski closure of
$\Gamma_0$, since $\Gamma'$ is normal.
Its inverse image in $\Gamma'$ is closed
of finite index and contains a positive power
of $u$.  The four distinct eigenspaces of
that power are therefore permuted transitively
by $\Gamma_0$, making the representation
monomial.  On the kernel of the
permutation action, all four reduced diagonal
characters coincide.  To see this, reduce their
integral eigenvalues and compare with the four
given residual characters; the latter are
conjugation-invariant, whereas transitivity
permutes the former.  Independence of characters
gives the assertion.  Every residual ratio thus
factors through a transitive subgroup of $S_4$.
Such a subgroup is $C_4$, $C_2^2$, a dihedral
group of order eight, $A_4$, or $S_4$; its
one-dimensional characters have order at most
four.  This contradicts Fit's \texttt{ord} field.

The inverse image of $H^\circ$ is closed and of finite index,
hence open.  Every open subgroup is Zariski dense
in the relevant connected component.  A power
of $u$ lies there and generates a diagonal torus
whose four weights are distinct.  Choose a
cocharacter $z\mapsto\operatorname{diag}(z^{a_i})$
with distinct integers $a_i$; this avoids finitely
many proper hyperplanes in its cocharacter lattice.

Restriction to $H^\circ$ is semisimple by Clifford's
argument: conjugates of one simple submodule span
the whole irreducible representation.  The simple
summands have equal dimension and multiplicity
one by the simple spectrum just noted.  Dimension
one would make $H^\circ$ commuting, already
excluded.  Thus there is one four-dimensional
block or two two-dimensional blocks.  On any
component $gH^\circ$, a generic matrix $M$ therefore
has either all four diagonal entries and all
four diagonal cofactors nonzero, or the block
form
\[
                  M=\begin{pmatrix}0&A\\ B&0\end{pmatrix}
                  \tag{QX.21}
\]
with every entry of $A,B$ nonzero and
$\det A\det B\ne0$.  This follows from Burnside
on each simple block; nonvanishing is an open
condition on the irreducible variety $H^\circ$.

In the first case the characteristic polynomial of
$M\operatorname{diag}(z^{a_i})$ is separable over
the Laurent-series field.  Order the $a_i$ strictly.
The extreme linear and cubic coefficients have
nonzero leading terms, so its Newton polygon has
two outer edges of length one; any remaining
edge has length at most two.  A length-two edge
has no intermediate leading coefficient: the
possible pair sums jump strictly past its midpoint.
Its residual binomial is separable because
$\operatorname{char}k\ne2$.  In the block-swap
case the polynomial is
$T^4-c(z)T^2+d(z)$.
The largest term of $c(z)$ is the unique product
using the largest exponent in each block and
has nonzero coefficient.  Its exponent is
strictly more than half the exponent of $d(z)$.
The two roots in $T^2$ have distinct orders and
each has two distinct square roots.  This proves
separability in that case too.

The discriminant is therefore a nonzero regular
function on a suitable translate of that torus.
Powers of $u$ are Zariski dense in it.  Choosing
the generic $M$ inside any prescribed open
coset and then multiplying by a suitable power
of $u$ produces a separable characteristic
polynomial in that coset.  This proves the
source's \texttt{N3} on every nonempty open
intersection with $\Gamma'$.

Finally \texttt{N4} is the auxiliary-prime
calculation already used in QX3.  The
Frobenius characteristic polynomial reduces
to $(T-a)^4$ with $a\ne0$, while the residual
cyclotomic value is not $1$.  Hence the
operator expressing a cyclotomic-twisted
adjoint invariant has invertible determinant
over $k[[t]]$.  Clearing denominators proves
that its kernel over $k((t))$ is zero.

We have proved the full \texttt{Nice\allowbreak I} hypothesis
for this actual point.  Apply \texttt{P1core}
to $f$.  For a prime
$\mathfrak q\subset\ker(y|_{X.B})$, its inverse
image under $\phi_0$ is prime below $\ker f$.
It contains $J_0$, so $\mathfrak q$ contains
$J_{\mathcal L}$.  This is exactly the
required \texttt{Loc\allowbreak Zero}, with all original
finite-field and local-separation quantifiers.
\end{proof}
% END QX8

\subsection{Purity and the original-field automorphic conclusion}

We use the following precise classical results in the final transport.
Arthur--Clozel, \emph{Simple algebras, base change, and the advanced
theory of the trace formula}, Chapter III, Theorems 4.2 and 5.1,
give cyclic prime-degree cuspidal descent and base change, with
strong local compatibility; Chapter I, \S6 gives their local
tempered compatibility.  The primary text is
\url{https://www.claymath.org/library/cw/arthur/pdf/30.pdf}.
Harris--Lan--Taylor--Thorne, \emph{On the rigid cohomology of
certain Shimura varieties}, Theorem A (Corollary 7.14),
\url{https://virtualmath1.stanford.edu/~rltaylor/rigcoh.pdf},
attaches a continuous semisimple Galois representation, with the
geometric normalization above, to every regular algebraic cusp
form over a CM field.  No polarization is required for this
existence assertion.  We use it only to identify a finite-character
ambiguity in cyclic descent.

For a regular algebraic conjugate-self-dual cusp form over a CM
field, Caraiani's \emph{Local-global compatibility and the action
of monodromy on nearby cycles}, Theorems 1.1--1.2, and
\emph{Monodromy and local-global compatibility for $l=p$},
Theorem 1.1, give full Frobenius-semisimplified local compatibility
and temperedness at every finite place, for every rank $n\ge2$:
\url{https://www.ma.imperial.ac.uk/~acaraian/papers/lgc1.pdf} and
\url{https://www.ma.imperial.ac.uk/~acaraian/papers/lgc2.pdf}.
There is no Shin-regularity condition in these statements.
The word ``cohomological'' in Caraiani's statements matches
the actual regular algebraic cusp used here: Clozel's
\emph{Motifs et formes automorphes}, Lemma 3.14, identifies
its archimedean representation with one having nonzero
relative Lie algebra cohomology for the corresponding
algebraic coefficient, as already recorded in the CX
classical-input registry.
For the polarized forms obtained after descent, we use
BLGGT, Theorem 2.1.1(2)--(3), and \emph{Local-global
compatibility for $l=p$, II}, Theorem A,
\url{https://virtualmath1.stanford.edu/~rltaylor/lg2.pdf}.
The latter identifies the semisimplified Weil representations
at $\ell$ for every regular algebraic polarized cusp form;
we do not invoke its stronger, Shin-regular conclusion.
The needed vanishing of monodromy is proved below.

% BEGIN QX9
\begin{proposition}[QX9: purity and the complete automorphic return]
\label{claim:QX9}
Let $J_{\mathcal L}=\phi_0(J_0)X.B$ for one of the actual
links of QX4--QX5, and assume the six stated
\texttt{Hida\allowbreak Props\allowbreak Alt} fields for $J_0$.
For every $y:X.A\to\overline{\mathbb Q}_\ell$ killing
$J_{\mathcal L}$ and satisfying
$T.\texttt{Wts}(2D.\texttt{wts4})$, the entire
\texttt{pur} clause of \texttt{Potential\allowbreak Modularity\allowbreak Tower}
holds.  If, as in its \texttt{aut} clause, the original
$r:G_{\mathbb Q}\to\operatorname{GL}_4
(\overline{\mathbb Q}_\ell)$ is continuous, has
multiplier $D.\texttt{th}$, restricts to $X.P_y$ over
$F$, and is irreducible on $G_L$, then the literal
\texttt{Is\allowbreak Aut4Alt}$\,\iota r$ holds, with the same
$\iota$ and for every imaginary quadratic field in
its definition.
\end{proposition}
\begin{proof}
Put $x=y|_{X.B}\phi_0$.  The weight calculation
in QX7, now with the even weights $2D.\texttt{wts4}$,
makes $x$ an actual arithmetic point.  The local
integer weights are
$D.\texttt{wts4}(i)+e_v/2$; here the exponent $e_v$
of $\kappa^2$ is even, since $a=0$.
Local reciprocity identifies their restrictions to
units with the required algebraic weight.  This
establishes the source's \texttt{Arith} condition,
not merely a Hodge-weight multiset.

Apply \texttt{end1} to $x$.  It supplies a continuous
semisimple $\rho$ with the actual determinant at $x$
and a characteristic-polynomial decomposition
\[
 \rho^{\rm ss}\simeq
 \bigoplus_i\ \bigoplus_{j=0}^{m_i-1}
                  \rho'_i\chi^{-j},
 \qquad d_i=\dim\rho'_i\ge1,\quad m_i\ge1.
 \tag{QX.22}
\]
All the $m_i$ are one.  Indeed choose stable lattices
in the finitely many continuous constituents.  Their
semisimplified reductions together give the four
distinct residual characters of $\kappa X.P$.
If $m_i\ge2$, a character of the reduction of
$\rho'_i$ and its cyclotomic twist both occur.
Their ratio is $\bar\chi$ or its inverse.
The common factor $\bar\kappa$ cancels, contrary
to Fit's exact \texttt{g2} exclusion of
$\bar\chi$ from the subgroup generated by the
original residual root characters.  This argument
uses reductions of stable lattices, not an
unjustified specialization of characteristic-zero
irreducible constituents.

Consequently \texttt{end1} gives, outside a finite
set of places $w$ of $L$, the strict inequalities
\[
        |\iota\alpha|<\operatorname{N}w\,
                           |\iota\alpha'|
        \quad\hbox{for every pair of Frobenius roots.}
        \tag{QX.23}
\]
The semisimple $\rho$ is unramified outside the
finite level set.  To check this from its determinant,
the universal law is constant on each killed
inertial coset.  Hence
$\operatorname{tr}(\rho(g)(\rho(u)-1))=0$ for
every $g$ and every killed inertia element $u$.
The semisimple image algebra is a product of matrix
algebras, and its trace pairing is nondegenerate in
characteristic zero.  Thus $\rho(u)=1$.

Take $H=G_L\subset G_{\mathbb Q}$ and let $M\ne0$
be the product of $\ell$, the rational primes under
these exceptional places, and the finitely many
level and discriminant primes.  If
$\sigma\in H$ is an arithmetic Frobenius of degree
$f$ at a rational prime $p\nmid M$, its image at
$w\mid p$ is the $(f/f_w)$-th power of the local
arithmetic Frobenius, where
$\operatorname{N}w=p^{f_w}$ and $f_w\mid f$.
The source \texttt{Att\allowbreak C} convention uses geometric
Frobenius, but (QX.23) is preserved on replacing
all roots by their inverses and exchanging the
ordered pair.  Raising it to the positive power
$f/f_w$ gives the bound $p^f$.
The characteristic polynomial of $\rho(\sigma)$
is that of $\kappa(\sigma)r(\sigma)$ by QX5.
This scalar cancels from every ratio.  Thus
\[
       |\iota z|<p^f|\iota z'|
       \quad(z,z'\hbox{ roots of }r(\sigma),\ z'\ne0),
       \tag{QX.24}
\]
with exactly the open subgroup, nonzero $M$, original
valuation, positive degree and Frobenius quantifiers
of \texttt{Pure\allowbreak W}.  Neither irreducibility nor
continuity of the given extension $r$ was used here.

For the automorphic assertion, assume its additional
hypotheses.  The scalar twist
$\rho_L=\kappa r|_{G_L}$ is irreducible and
continuous.  The \texttt{end2} field at $x$ gives
an actual regular algebraic conjugate-self-dual
cuspidal $\pi_L$ attached to $\rho_L$.
In fact this field supplies \texttt{Aut\allowbreak Off\allowbreak N},
including a window assertion away from the auxiliary
prime.  We use its actual cusp and the matched
stronger local theorems above.  This is a deliberate
simplification: two auxiliary primes remain available
for the integral support construction, but the
automorphic return does not need an unmatched
\texttt{Win\allowbreak Of} or \texttt{Hypotheses\allowbreak H'} package.

The character $\kappa$ is algebraic by (QX.11)--(QX.12)
and the rank-one reciprocity argument of QX7.
Twist $\pi_L$ by its inverse Hecke character.
The resulting regular algebraic cusp $\pi_L^0$
attaches to the \emph{original} $r|_{G_L}$.
The twist also transports all local Langlands
parameters; it does not change the rank and so
introduces no extra Tate twist.

We spell out cyclic descent with the actual Galois
representation fixed.  Suppose $E/E_0$ is a cyclic
prime-degree step in $L/K$, and an actual cusp
$\pi_E$ attaches to $r|_{G_E}$.  This restriction
is irreducible, since its further restriction to
$G_L$ is.  For $\tau\in\operatorname{Gal}(E/E_0)$,
the attached representation is unchanged up to
isomorphism; strong multiplicity one gives
$\pi_E^\tau\simeq\pi_E$.
Arthur--Clozel III.4.2(d) supplies a cuspidal descent
$\pi_0$.  Its archimedean infinitesimal character
is the same at each extending complex embedding,
so it is regular algebraic.
Use HLTT to attach the continuous semisimple
$r_{\pi_0,\iota}$ without first presupposing
polarization of this chosen descent.
Its restriction to $G_E$ equals $r|_{G_E}$ by
Chebotarev and base change.  Both extensions of
this irreducible representation differ by a
character of $\operatorname{Gal}(E/E_0)$:
after choosing an intertwiner on $G_E$, Schur's
lemma makes the ratio of the two matrices a scalar
on every coset; these scalars form the required
finite continuous character.  Twist $\pi_0$ by
its finite Hecke character to obtain precisely
$r|_{G_{E_0}}$.  This proves descent without
silently changing the original Galois representation.
Iterating reaches the original imaginary quadratic
$K$, since $L/K$ is abelian.

The multiplier also descends correctly.  Set
$\theta=D.\texttt{th}$.  The symplectic pairing
and complex conjugation give
$r^c\simeq r^\vee\theta$ over every imaginary
quadratic field.  In our normalization the
Galois representation of $\pi^\vee$ is
$r_\pi^\vee\chi^{-3}$.  Thus strong multiplicity
one makes the descended cusp essentially
conjugate-self-dual with the algebraic Hecke
character corresponding to $\theta\chi^3$.
This character comes from $\mathbb Q$.  If its
real-place sign is not the required rank-four
sign $+1$, multiply it by the quadratic Hecke
character $\eta_{K'}$ of the imaginary quadratic
field in question.  That character changes the
real sign and is trivial after norm pullback
to $K'$; its Galois avatar is trivial on
$G_{K'}$.  Hence the automorphic polarization
and its Galois multiplier on $G_{K'}$ are
unchanged.  We do not identify a Hecke infinity
sign with the value of its uncorrected
cyclotomic avatar at complex conjugation, nor
assert global avatar equality after this
correction.  This proves all
the hypotheses of the polarized Galois and local
compatibility theorems used below.  The multiplier
is not replaced by its restriction or by a
character with a merely similar infinity type.

Now let $K'$ be any imaginary quadratic field
such that $r|_{G_{K'}}$ is irreducible.  The
representation $r$ is unramified outside a finite
set: outside the original residual ramification
and the discriminant of $F$, inertia lies in
$G_F\cap\ker\bar r$, where the actual
$X.\texttt{pair.\allowbreak kill}$ makes it trivial.
Choose another imaginary quadratic $K_0$ ramified
at an odd prime $q$ outside this set and outside
the discriminants of $K,K'$.
Then both $r|_{G_{KK_0}}$ and
$r|_{G_{K'K_0}}$ are irreducible.  Otherwise
index-two Clifford theory supplies the corresponding
quadratic self-twist of $r|_{G_K}$ or
$r|_{G_{K'}}$; on inertia at $q$ the original
representation is trivial but the twisting
character is nontrivial, a contradiction.
Cyclic base change from $K$ to $KK_0$ is
cuspidal: the exceptional alternative in
Arthur--Clozel III.4.2(b) would give exactly
this prohibited self-twist.
Descend to $K_0$, base change to $K'K_0$, and
descend to $K'$ by the same argument.  All
finite-character ambiguities are corrected as
above, and all uses retain the original $\iota$.
We have obtained an actual regular algebraic
polarized cusp attached to $r|_{G_{K'}}$ for
\emph{every} field in the source's quantifier.

It remains to prove the full \texttt{Wit\allowbreak A}, rather
than stopping at an automorphic attachment.
First note two local consequences of the actual
space.  At $p\ne\ell$, the open group
$G_F\cap\ker\bar r$ contains the commutator
subgroup of $G_{\mathbb Q}$: both $F/\mathbb Q$
and $\bar r$ are abelian.  The killing relation
therefore annihilates every $[d,u]$ with $d$ in
a decomposition group and $u$ in its inertia.
It also kills an open subgroup of inertia.
Thus the inertia image is finite and central
in the local image.  A lift of Frobenius and
this finite central inertia have dense generated
image, so the entire local image is commutative.
Its semisimplification is a sum of continuous
one-dimensional characters and its monodromy is
zero.  This proves the exact \texttt{Ine\allowbreak Ab}
condition on the original group and on every
restriction to $K'$.

At $\ell$, the flag of $X.P_y$ has the distinct
parallel weights
\[
                     e+1, e, 1, 0,\qquad e\ge2.
                     \tag{QX.25}
\]
Choose an open normal subgroup of the original
local decomposition group on which all its
diagonal characters are exactly the corresponding
powers of $\chi$; the finite-order errors in
$T.\texttt{Wts}$ disappear after shrinking.
Choose an element $u$ of its inertia on which
$\chi$ has infinite order.  Such an element
exists because the cyclotomic image of every
open local inertia subgroup is infinite.
The matrices of $u$ and $dud^{-1}$, in the
given flag basis, are upper triangular with
the same distinct diagonal entries
$\chi(u)^{e+1},\chi(u)^e,\chi(u),1$.
The equation
$r(d)r(u)=r(dud^{-1})r(d)$ forces $r(d)$
upper triangular: starting at the lower-left
corner, each entry below the diagonal is
multiplied by the nonzero difference of its
two diagonal eigenvalues, and induction makes
it zero.  Thus the ordered flag descends to
the original decomposition group.  Its diagonal
characters remain locally algebraic and parallel
with the weights (QX.25).  This proves the
actual \texttt{Flagged} condition, including
all original places; conjugating the chosen
place transports the construction.  The
restriction to any $K'$ has the same property.

Every cusp constructed above is essentially
tempered at every finite place.  The initial
$\pi_L$ is tempered by Caraiani's Theorem 1.2.
An algebraic Hecke character has absolute value
a single real power of the adelic norm, so the
inverse $\kappa$-twist preserves essential
temperedness with a single global normalization.
Strong local cyclic base change preserves and
reflects temperedness after this normalization:
in the local Langlands parameter, boundedness
on the Weil group before inserting its
$\mathrm{SL}_2$ factor is unchanged on a subgroup
of finite index.  The finitely many cosets give
the reverse implication as well.  Finite-character
twists are unitary.  This proves the assertion
through each of the descent and base-change
steps, without asserting the Ramanujan conjecture
for an unrelated unpolarized cusp form.

Here is the necessary local calculation at
$v\mid\ell$.  Write the characters in the
descended flag as
$\psi_i=\chi^{w_i}\eta_i$, where $\eta_i$
has finite inertia.  This follows from local
reciprocity and the parallel weights, or by
killing the finite inertia characters on an
open subgroup and descending back.  Continuity
makes the value of $\eta_i$ on Frobenius an
$\ell$-adic unit.  The Weil character of
$\chi$ at geometric Frobenius has value
$q^{-1}$.  Hence an equality
\[
 \operatorname{WD}(\psi_i)
   =\operatorname{WD}(\psi_j)\,|\cdot|
 \quad\Longrightarrow\quad
 w_i-w_j=1,\qquad \eta_i=\eta_j,\qquad
 \psi_i=\psi_j\chi.
 \tag{QX.26}
\]
Indeed the $\ell$-adic valuations determine the
weight difference, and the remaining Weil
characters of finite-inertia Galois characters
are faithful data.  The same implication holds
after exchanging $i,j$.  All one-dimensional
composition factors of a flag are intrinsic
as a multiset, so a different flag in
\texttt{Unl1} has the same test.

Suppose \texttt{Unl1} holds.  BLGGT's
semisimplified local compatibility identifies
the Weil constituents of the constructed cusp
with these four characters.  A nonzero
monodromy entry would give exactly a norm
ratio in (QX.26), contradicting \texttt{Unl1}.
Thus its local monodromy is zero.  Its
essential temperedness then makes all four
Weil-character absolute values equal after
one scalar normalization.  An equality
$\psi_i^m=(\psi_j\chi)^m$ with $m\ge1$
would instead give a Weil-character norm ratio
$q^{\pm m}$, a contradiction.  This proves
the first clause \texttt{Unl1}$\Rightarrow$
\texttt{Pot\allowbreak Unl}, with every positive power
and every original decomposition group.

At a place away from $\ell$, full polarized
local compatibility identifies the Weil--Deligne
parameter with the already proved sum of
characters and zero monodromy.  At a place
above $\ell$ where \texttt{Pot\allowbreak Unl} holds,
the preceding argument does the same.  Local
Langlands now gives a toral generic principal
series; all of its cuspidal support characters
are one-dimensional and have the same
uniformizer absolute value after the global
essential-unitary normalization.  In particular,
for one real number $t$ fixed by the global
central character, at every such place
\[
          q^{t-1/2}<|\xi(\varpi)|<q^{t+1/2}
          \quad\hbox{for every cuspidal-support character }\xi.
          \tag{QX.27}
\]
The inequalities are strict since $q>1$; the
middle value is $q^t$.  This includes the
source's \texttt{Tor\allowbreak At} assertion and is
exactly \texttt{Win\allowbreak At}, with one $t$ for
all places.  Any other regular algebraic cusp
attached to the same original $r|_{G_{K'}}$
is isomorphic to this one by strong multiplicity
one, so the universal quantifier over such
cusps in \texttt{Wit\allowbreak A} is also satisfied.

We have proved the actual automorphic attachment
and \texttt{Wit\allowbreak A} over every requested imaginary
quadratic field, using the same $\iota$ throughout.
This is the literal \texttt{Is\allowbreak Aut\allowbreak QA}, hence
\texttt{Is\allowbreak Aut4Alt}, for the original $r$.
\end{proof}
% END QX9

\subsection{The exact remaining support input and the five-clause join}

Call the following assumption the \emph{rank-four ordinary-support
input on the two constructed links}.  For each original Fit and
Deep tower and each fixed $\iota$, use the coefficient threshold
and the two links constructed in QX4.  Over every further
coefficient ring, for each of these links and every actual
universal ordinary ring $B$ and weight object $\Lambda$, there
exists an ideal $J_0$ satisfying all six \texttt{Hida\allowbreak Props\allowbreak Alt}
fields printed above and the literal $B.\texttt{P1core}\,J_0$.
This assumption is restricted to rank four and these actual
data.  It asserts existence of the ideal, including finite
torsion-free weight-module structure, classical points and
their two arithmetic endpoints, and characteristic-$\ell$
support.  It is a retained arithmetic input, not a consequence
of QX1 or a name for a theorem already proved in this chapter.

% BEGIN QX10
\begin{theorem}[QX10: the actual modular-support constructor]
\label{claim:QX10}
The stated rank-four ordinary-support input implies the entire
literal \texttt{PMEx\allowbreak Alt}.  More explicitly, for every original
$D,T$ with $T.\texttt{Fit}$ and $T.\texttt{Deep}$ and every
fixed $\iota$, there is a coefficient threshold $c_1$ such that,
for every $c'$ containing both $c$ and $c_1$ and every actual
$X:\texttt{Space}(T.\texttt{fr}\,\varnothing)$ over $c'$, an
ideal $J\subset X.B$ satisfies all five clauses of
$T.\texttt{Potential\allowbreak Modularity\allowbreak Tower}\,\iota XJ$.
\end{theorem}
\begin{proof}
Choose the common threshold of QX4; it depends on the
original tower and its character data, and may be enlarged
for $\iota$ if the retained support input requires that.
It is chosen before $c'$ and before $X$.
Fix any allowed $c'$ and any actual $X$.
For the first link, QX5 constructs an actual ordinary
\texttt{Pair} over $X.B$.  Thus the nonempty-pair
antecedent of KX9 is fulfilled, and KX9 supplies an
actual universal ring $B$.  KX5 supplies its actual
weight object $\Lambda$.
The retained input supplies $J_0$ for these data.
Perform the same construction for the second link,
obtaining $B',\Lambda',J'_0$ and the two maps
$\phi_0,\phi'_0$.  Define
\[
        J=\phi_0(J_0)X.B+\phi'_0(J'_0)X.B.
        \tag{QX.28}
\]
This is a concrete ideal in the original $X.B$;
no coefficient specialization or change of the
original space has been made.

For \texttt{fin}, use the first quotient's actual
finite weight-module field.  IX9 proves the full
algebraic finite-fibre argument: the universal
square-characteristic identities are formed
before quotienting; they make the original
characteristic coefficients integral over the
weight image.  After killing the coefficient
maximal ideal and all
$\psi_{v,i}(u)-1$ for original inertia elements,
the finite weight-module quotient is finite
dimensional over the residue field.  Its
maximal ideal is nilpotent.  Thus, for some
integer $N$,
\[
 \mathfrak m_{X.B}^{N}\subset
 J_{\mathcal L}+\mathfrak m_{c'.O}X.B+
 (\psi_{v,i}(u)-1:\ u\in I_v).
 \tag{QX.29}
\]
All hypotheses of IX9 are supplied by the actual
\texttt{Ring\allowbreak Props.\allowbreak fin}, QX5's \texttt{Is\allowbreak Phi},
and the unchanged \texttt{Space} generation
fields.  Enlarging $J_{\mathcal L}$ to $J$
gives exactly \texttt{fin}; no injectivity of
a map after quotienting is assumed.

For \texttt{p1}, take any of its finite-field
power-series tests killing $J$.  It kills both
summands.  QX8 says that every prime below its
$X.B$-kernel contains each summand, hence their
sum.  This is the required \texttt{Loc\allowbreak Zero}.
For \texttt{seed}, QX7 applies to the same seed
point for each link.  Again every prime below
its kernel contains both summands and therefore
$J$.  Its auxiliary twisting parameter may
be chosen separately for the two proofs; QX6
returns each to its original $a=0$ image ideal.

For \texttt{pur} and \texttt{aut}, a point
killing $J$ kills its first summand.  Apply QX9
with exactly the additional hypotheses printed
in the corresponding clause.  Its proof
returns the original $r$, every place and
imaginary quadratic field, and the same fixed
$\iota$.  Thus all five fields are supplied
on this same $X,J$.
The construction worked for arbitrary $c'$
above the common threshold and arbitrary $X$;
these universal quantifiers were never
interchanged with an $X$-dependent coefficient
choice.  This is precisely \texttt{PMEx\allowbreak Alt}.
\end{proof}
% END QX10

\paragraph{What this closes and what remains.}
QX1 supplies the literal \texttt{Adjoint\allowbreak Lifting\allowbreak Theorem}.
QX2--QX3 close its fixed-weight tangent-to-support use;
QX4--QX10 close every internal step from the actual
two ordinary-support ideals to the five-clause
\texttt{PMEx\allowbreak Alt} conclusion.  The original
\texttt{P3core} field is not used on this route.
The six-field ordinary modular quotient and its
\texttt{P1core} support remain to be constructed on
the actual rank-four residual data.  In particular
the residual-Schur and local-degree assumptions
of Allen--Newton--Thorne's residually reducible
automorphy-lifting theorem cannot be inferred from
the present Fit and Deep fields.  Neither this
chapter nor adjoint vanishing supplies those
hypotheses.  The exact Hida-control, Hecke-ideal,
arithmetic-point and characteristic-$\ell$ support
producer cut is recorded in the accompanying
frontier registry.  No unconditional \texttt{PMEx\allowbreak Alt},
full \texttt{Unramified\allowbreak Inputs}, paired-count result,
\texttt{Point\allowbreak HV}, RH or GRH conclusion is promoted.


\section{Constructing the mixed automorphic seed}
\label{sec:symplectic-automorphic-seeds}

This chapter constructs the seed used in case II of the symplectic
support argument.  In particular, an automorphic representation is
produced, rather than added to the hypotheses of a deformation point.
The construction has three parts: an Eisenstein congruence with different
ordinary refinements at two places; a change of weight inside an ordinary
Hecke algebra; and an explicit integral basis for an induced
four-dimensional representation.  The last part is essential: abstract
automorphy of an induction would not specify the residual matrix entries
required by the support argument.

The source endpoint is \texttt{Seed\allowbreak IIAlt}.  Its premise already includes
an arbitrary original tower satisfying \texttt{Ok}, \texttt{Two}, and
\texttt{Deep}.  We preserve that tower, its labels, its chosen element
$g_0$, all its specified valuations, and the previously fixed field
isomorphism $\iota:\overline{\mathbb Q}_\ell\simeq\mathbb C$.
The only choices made afterwards are the seed and a finite enlargement
of the coefficient field.  There is no assertion here that an arbitrary
tower satisfies these three premises.

\subsection{Conventions and precisely stated classical inputs}

Write $\chi$ for the cyclotomic character, with arithmetic Frobenius
sent to the residue cardinality.  The source's attachment convention
uses geometric Frobenius: if $B_v$ is the normalized Satake multiset
of an automorphic representation of rank $n$, its attached Galois
multiset is
\[
                 q_v^{(n-1)/2} B_v.                 \tag{SX.1}
\]
Thus twisting a Galois representation by $\chi$ corresponds to
twisting its automorphic representation by $|\det|$.  This convention
accounts for the extra $|\det|$ in the rank-two to rank-four induction
below.

We use the following established theorems.  The applications and the
remaining algebra are proved in the chapter.

\begin{itemize}
\item \textbf{Hilbert Eisenstein series (SX-E).}
For a totally real field and finite ray characters $\eta,\zeta$ of
the appropriate parity, the holomorphic series $E_j(\eta,\zeta)$
has positive ideal coefficient
$\sum_{\mathfrak a\mid\mathfrak m}
 \eta(\mathfrak m/\mathfrak a)\zeta(\mathfrak a)
 N\mathfrak a^{j-1}$.
The weight-one series $E_1(1,\zeta)$ exists when $\zeta$ is totally
odd and nontrivial.  Its standard-cusp constants and its degeneracy
maps are those of this normalized expansion.
Dasgupta--Darmon--Pollack, Proposition 2.1 and its proof, printed
pp.~458--462, give this precise weight-one statement.  We also use
the usual decomposition of holomorphic Hilbert forms of weight $>2$
into cuspidal and Eisenstein eigensystems, and the rational
$q$-expansion principle.

\item \textbf{An integral Hasse invariant (SX-H).}
For a Hilbert modular variety at a prime unramified in the totally
real field, the determinant Hasse invariant has parallel weight
$\ell-1$, extends to the minimal compactification, and has expansion
$1$ at every Tate cusp.  A sufficiently high power lifts to
characteristic zero.  Mok--Tan, \S2.6, especially Proposition 2.11
on printed p.~906, gives the lifting by Serre vanishing; we use its
genus-one specialization.  Finite level descent by a product, described
in SX2, gives exactly the level-one form needed here.

\item \textbf{Ordinary control and ordinary Galois representations
(SX-O).}
At tame level prime to $\ell$, the cuspidal nearly ordinary Hecke
algebra with fixed algebraic central character is finite and
torsion-free over the corresponding weight algebra.  Its arithmetic
specializations of dominant algebraic weight are the classical
ordinary Hecke algebras after inverting $\ell$.  The identification
preserves away-from-$\ell$ Hecke operators and the normalized
operators $U_v^{\circ}$.  The attached rank-two Galois representation
has the usual unramified characteristic polynomials and the ordered
ordinary local characters specified by $U_v^{\circ}$ and the two
weight coordinates.  This statement applies to residual Eisenstein
systems.  Hida's \emph{Control of nearly ordinary Hecke algebras},
Theorems 1.2 and 2.4, Corollary 2.5, and Theorem 3.1, pp.~3--8,
state precisely these facts.  Remark 3.1 requires residual absolute
irreducibility only for matrices over the integral big Hecke algebra;
we only use individual characteristic-zero representations.

\item \textbf{Cyclic base change and automorphic induction (SX-BC).}
Cyclic base change for cuspidal $\mathrm{GL}_n$ representations and
cyclic automorphic induction exist, commute with local parameters,
and have the standard cuspidality criteria: a prime-degree base change
fails to be cuspidal precisely in the self-twist case, and an induction
of a cuspidal representation in degree two fails to be cuspidal
precisely when its two conjugates agree.  We use successive cyclic
base changes for an abelian extension.  Arthur--Clozel, Chapter III,
Theorems 4.2 and 5.1 and \S6, Theorem 6.2, are the global inputs.
For all-place compatibility of induction we use Henniart,
\emph{Induction automorphe globale pour les corps de nombres},
Theorems 4--5 and \S3.7, printed pp.~6--7 and 16.
Local compatibility here only concerns principal series and unramified
induction, where it is also the explicit character/norm description.

\item \textbf{Galois compatibility after CM base change (SX-G).}
A regular algebraic polarized cuspidal representation over a CM field
has an attached semisimple Galois representation, with semisimplified
local--global compatibility away from $\ell$.
We use BLGGT, \emph{Potential automorphy and change of weight},
Theorem 2.1.1, solely for this rank-two conclusion.
The quadratic character of the CM extension may be multiplied into
the polarization character to fix its real-place sign; its norm
pullback is trivial.  The comparison with the already constructed
ordinary representation is by Chebotarev and semisimplicity.
No automorphy lifting theorem is used to construct the seed.

\item \textbf{Classical local and global background (SX-C).}
We use global and local class field theory, Chebotarev density,
Dirichlet's unit theorem, the local newvector conductor theorem,
and the principal-series part of the local Langlands correspondence.
The latter identifies the conductor and determinant with those of
the Weil--Deligne parameter.  The elementary conductor and unitary
estimates needed from it are proved below.  The finite coefficient
capture theorem RV4 is already reconstructed and is used with an
actual continuous group homomorphism, with the matrix argument
spelled out in SX7.  For four characters, the absence of ratios
$|\cdot|^{\pm1}$ implies irreducibility of normalized induction
by Bernstein--Zelevinsky, Theorem 4.2, printed p.~454.
\end{itemize}

\subsection{The original quadratic data}

Fix an original datum $D$ over $O_c$ and a tower $T$ with
\[
 D.\texttt{Ctx2},\qquad T.\texttt{Ok},\qquad
 T.\texttt{Two},\qquad T.\texttt{Deep}.             \tag{SX.2}
\]
Let $F=T.F$, $K=T.K$, and $L=FK=T.L$ inside the original closure
$\Omega$.  We identify their absolute Galois groups with subgroups
of $G_{\mathbb Q}=\operatorname{Gal}(\Omega/\mathbb Q)$.
Put $t_0=T.g_0\in G_F$, and let $p_i$ be the four characters
$\texttt{chi}(D,T.\mathrm{pr},i)$ with the original labeling.
We use $\xi,\psi,\omega,\vartheta$ for their scalar extensions to
$\overline{\mathbb Q}_\ell$.  Thus
\[
 \alpha=\chi\omega^{-1}\xi,
 \quad \Psi=\chi^e\psi,
 \quad \vartheta=\alpha\Psi,
 \quad p_i p_{3-i}=\vartheta.                         \tag{SX.3}
\]
An overline denotes reduction modulo the maximal ideal, after any
needed finite coefficient extension.

\begin{proposition}[SX1: the mixed quadratic input on the original tower]
\label{claim:SX1}
There is a totally real quadratic field $F_0\subset\Omega$, already
specified by $T.\texttt{Two}$, split at $2$ and $\ell$.  Its sign
character $\varphi:G_{\mathbb Q}\to\{\pm1\}$ has the following
properties.  On $H=G_F\cap G_{F_0}$,
\[
 (\bar p_0,\bar p_1,\bar p_2,\bar p_3)
       =(\bar\xi,1,\bar\xi,1),                    \tag{SX.4}
\]
while $t_0\notin G_{F_0}$ and
\[
 \bar p_1(t_0)=-\bar p_3(t_0),\quad
 \bar p_2(t_0)=-\bar p_0(t_0),\quad
 \bar p_3(t_0)\in\{1,-1\},\quad
 \bar\xi(t_0^2)\ne1.                              \tag{SX.5}
\]
The character $\xi|_{G_{F_0}}$ is totally odd, has finite
prime-to-$\ell$ order, is unramified outside the primes above $2$,
and is nontrivial on inertia at a prime of norm $2$.  Its reduction
has nontrivial square somewhere and is nontrivial on each
decomposition group above $\ell$.
For every $k$ divisible by $\ell-1$,
\[
 \mu:=\chi^k\xi\varphi
       \quad\hbox{satisfies}\quad
 \bar\mu|_{G_F}=\bar\vartheta|_{G_F}.              \tag{SX.6}
\]
\end{proposition}
\begin{proof}
The field and splitting assertions are literal clauses of
$T.\texttt{Two.\allowbreak quad}$.  Its kernel in $G_F$ is the kernel of
$T.\mathrm{rb}(w_{\Psi})$; the latter has order at most two.
The character taking $1$ on $G_{F_0}$ and $-1$ on its other coset
is multiplicative because this subgroup has index two, and continuous
because it is open.

The two possible labelings are
\[
 (\alpha\Psi,\Psi,\alpha,1),\qquad
 (\alpha,1,\alpha\Psi,\Psi).                        \tag{SX.7}
\]
As $\bar\omega=\bar\chi$, one has $\bar\alpha=\bar\xi$.
On $H$, $\bar\Psi=1$, proving (SX.4).  If $t_0$ were in $H$,
the second and fourth residual eigenvalues at $t_0$ would coincide,
contrary to $T.\texttt{Ok.\allowbreak dist}$.  Thus $\bar\Psi(t_0)=-1$.
Substitution in (SX.7) proves the first three assertions of (SX.5).
If $\bar\xi(t_0)^2=1$, then $\bar\xi(t_0)=\pm1$, and again two
of the four displayed eigenvalues coincide.  This proves the last
assertion.

The finite order and the ramification of $\xi$ are the clauses
$\texttt{Two.\allowbreak ord}$ and $\texttt{Two.\allowbreak xi}$; its order is a power of
two, whereas $\ell$ is odd.  Its oddness follows from $D.\texttt{Good}$
and conjugacy of complex conjugations.  At either prime above $2$,
$\texttt{Two.\allowbreak big2}$ supplies inertia on which
$\bar\xi^4\ne1$; the splitting of $F_0$ at $2$ puts that inertia
inside $G_{F_0}$.  The square of this element proves that
$\bar\xi^2$ is not the trivial character.  Similarly
$\texttt{Two.\allowbreak big\allowbreak L}$ and splitting at $\ell$ give the required
nontriviality on both decomposition groups above $\ell$.
In particular the conductor of $\xi|_{G_{F_0}}$ is nontrivial.

Finally $\bar\chi^k=1$, and on $G_F$ the characters $\bar\varphi$
and $\bar\Psi$ agree: both are $1$ on $H$ and $-1$ on its other
coset.  Hence $\bar\mu=\bar\xi\bar\Psi=\bar\vartheta$.
Every statement is on the original groups; no isomorphic copy of
the tower has been substituted.
\end{proof}

\subsection{An actual cuspidal ordinary congruence}

\begin{proposition}[SX2: mixed refinements force a cuspidal eigenlift]
\label{claim:SX2}
Let $F_0$ and $\xi$ be as in SX1, and fix $\iota$.
Choose one of the two primes above $\ell$, and call the singleton
set consisting of it $P$.  There exists an arbitrarily large
$k\ge2$, divisible by $\ell-1$, and a cuspidal ordinary Hilbert
eigenform of parallel weight $k+1$, nebentype corresponding to
$\xi$ in the geometric-Frobenius convention, and tame level its
conductor $\mathfrak c$, whose residual arithmetic Hecke system
is $1\oplus\bar\xi^{-1}$.  At the primes above $\ell$
the residual ordinary refinement is $\bar\xi^{-1}$ on $P$ and $1$
outside $P$.
\end{proposition}
\begin{proof}
Use class field theory and the fixed $\iota$ to write the finite
character as a complex narrow-ray character $\zeta$ with
$\zeta(v)=\iota(\xi(\operatorname{Frob}_{v,\mathrm{arith}})^{-1})$
at unramified places.  Its arithmetic Galois avatar is therefore
$\xi^{-1}$.  The series $E=E_1(1,\zeta)$
has first positive ideal coefficient $1$ and integral positive
coefficients.  At each $v\mid\ell$, its two stabilization roots
are $1$ and $\zeta(v)$.  Let $V_v$ be the degeneracy operator
whose normalized positive coefficient at $\mathfrak m$ is the
coefficient at $\mathfrak m/v$, or zero if $v\nmid\mathfrak m$.
Here is its constant-term normalization.  Let $u_v$ be the finite
idele equal to $\varpi_v^{-1}$ at $v$ and $1$ elsewhere.
For weight $j$, set
$V_v=q_v^{-j/2}R(\operatorname{diag}(u_v,1))$.
Substituting in the adelic Fourier expansion cancels the factor
$|u_v|^{j/2}=q_v^{j/2}$ and sends the normalized constant
function $c_0(y)$ to $c_0(yu_v)$.  The positive coefficient
indices shift by $v$, giving the stated formula.
For $E_1(1,\zeta)$, formula (50) of SX-E gives
$c_0(y)=2^{-[F_0:\mathbb Q]}L(\zeta,0)$, independently of
the narrow-class component.  Only that first constant term
occurs, since $\zeta$ has nontrivial conductor by SX1.
Apply
\[
 E_P=\prod_{v\in P}(1-V_v)
       \prod_{v\mid\ell,\,v\notin P}(1-\zeta(v)V_v)E.
                                                            \tag{SX.8}
\]
The identity $(U_v-1)(U_v-\zeta(v))E=0$, or the divisor-sum
coefficient formula, shows that $U_v E_P=\zeta(v)E_P$ on $P$
and $U_v E_P=E_P$ elsewhere.  At every standard cusp the constant
of $V_vE$ equals that of $E$ in this normalization.  Since $P$
is nonempty, the first product in (SX.8) kills these constants.
All positive coefficients remain integral, and the first one remains
$1$.  This assertion concerns the standard expansions on every
narrow-class component.  It does not assert that every cusp constant
of $E_P$ vanishes, and does not assert that $E_P$ is cuspidal.

Here is the required form congruent to $1$.  Choose an auxiliary
prime-to-$\ell$ fine level.  The Hilbert determinant Hasse invariant
on its minimal compactification has parallel weight $\ell-1$.
On a Tate object the Cartier operator fixes the logarithmic
differentials, so its determinant has expansion $1$.  For all large
even $a$, Serre vanishing applied to multiplication by $\ell$ in the
Hodge line bundle shows that its $a$th power lifts to an integral
characteristic-zero section $G_a$ of weight $a(\ell-1)$.
Make the auxiliary cover Galois and multiply all deck translates
of $G_a$.  This product descends to level one, has trivial
nebentype, and has expansion $1$ modulo the maximal ideal at
every cusp.  This descent uses a product, so no division by the
order of the deck group is involved.  On all narrow-class components
the same canonical Hasse section is used.  To include the adelic
central descent, take the product also over the finite central
ideal-class action on these components.  The normalized line-bundle
action factors through this finite class group: principal positive
ideles act by the prescribed parallel weight, and integral units
by the level-one trivial character.  Choose ideal-class representatives
prime to $\ell$.  Because the starting weight is a multiple of
$2(\ell-1)$, the rational norm factors in the normalized central
action are $\ell$-adic units congruent to $1$.
The Hasse section is fixed
by this action, by the Cartier functoriality under prime-to-$\ell$
isogenies (Mok--Tan, Lemma 2.10).  Every factor therefore still
has reduction $1$.  Taking a further common power if needed
trivializes the finite stabilizer action on the Hodge line.
Invariant sections descend on the finite quotient, and the
resulting adelic form has the trivial finite central character.
Increasing $a$ gives
arbitrarily large resulting weights $k$ divisible by $\ell-1$.

Multiply $E_P$ by such a level-one form $G$.  In a standard expansion
the positive coefficient of the product is the constant coefficient
of $G$ times that of $E_P$, plus a sum of products of positive
coefficients.  The possibly nonintegral constant of $E_P$ is zero.
Consequently every positive coefficient of $E_PG$ is integral,
and its reduction equals that of $E_P$.  The first coefficient is a
unit.  The explicit Hecke coefficient formulas then show that
$E_PG$ is an approximate simultaneous eigenvector modulo the
maximal ideal for all the away Hecke operators and all $U_v$.
At weight $k+1$ the factor $(Nv)^k$ occurring in these formulas
reduces to $1$ for $v\nmid\ell$, as $\ell-1\mid k$.

For completeness, the eigenlift is an algebraic step.  Work over a
finite coefficient extension containing $E_PG$, its character, and
the rational holomorphic space.  The module of forms with integral
positive coefficients is a lattice: finitely many coefficient
functionals already inject the finite-dimensional space into a
finite product of the coefficient field, and the $q$-expansion
principle supplies bounded denominators in the other direction.
It is Hecke stable by the coefficient formulas.  Let $\mathbb T$
be the finite commuting Hecke algebra acting on it.  The nonzero
residual approximate eigenvector gives a maximal ideal
$\mathfrak m\subset\mathbb T$ with the specified residue
eigenvalues.  The localized lattice is nonzero and torsion-free
over the discrete valuation ring.  Hence some prime of
$\mathbb T_{\mathfrak m}$ avoids $\ell$.  After passing to a
finite characteristic-zero residue field and a common eigenline,
one obtains a holomorphic eigenform with exactly this residual
system.  Its $U_v$ eigenvalues are units, so it is ordinary.

It remains to prove that this lift is cuspidal.  Suppose otherwise.
Its Eisenstein system has two finite characters $\eta_1,\eta_2$,
with away eigenvalues
$\eta_1(v)+(Nv)^k\eta_2(v)$ and product
$\eta_1\eta_2=\zeta$.  Reduction and Chebotarev give, for every
Galois element $g$,
\[
 \bar\eta_1(g)+\bar\eta_2(g)=1+\bar\xi(g)^{-1},
 \qquad\bar\eta_1(g)\bar\eta_2(g)=\bar\xi(g)^{-1}.
                                                            \tag{SX.9}
\]
Thus $(\bar\eta_1(g)-1)(\bar\eta_2(g)-1)=0$.
A group cannot be the union of two proper subgroups: if $a$ is
outside the first kernel and $b$ outside the second, each lies in
the other, whereas $ab$ lies in neither.  Hence one of the two
characters is identically $1$ and the other is $\bar\xi^{-1}$.

At $v\mid\ell$, both inducing characters are unramified.  Indeed
the nebentype is unramified there and the level has exponent one;
if either character were ramified, cancellation in the central
character would make both ramified and their conductor sum would
be at least two.  The two possible $U_v$ eigenvalues are
$\eta_1(v)$ and $(Nv)^k\eta_2(v)$.  Only the first is a unit.
Therefore the residual ordinary eigencharacter is the same
global character $\bar\eta_1$ at every $v\mid\ell$.
Our prescribed refinement is $\bar\xi^{-1}$ at the prime in $P$
and $1$ at the other prime, while $\bar\xi^{-1}$ is nontrivial at both.
Neither of the possibilities following (SX.9) is compatible with
these two refinements.  This contradiction proves cuspidality.
It is this argument, rather than multiplication by $G$, that kills
the Eisenstein alternative.
\end{proof}

\subsection{Changing weights with fixed determinant}

\begin{proposition}[SX3: the full mixed rank-two representation]
\label{claim:SX3}
On the same $F_0$, $\xi$, and fixed $\iota$, there are integers
$k,t\ge2$, each divisible by $\ell-1$, and a continuous
representation $r:G_{F_0}\to\mathrm{GL}_2(\overline{\mathbb Q}_\ell)$
with the following properties:
\[
 \det r=\chi^k\xi,
 \qquad\bar r^{\mathrm{ss}}=1\oplus\bar\xi.         \tag{SX.10}
\]
It is irreducible on every open subgroup.  Away from $\ell$,
its inertia is conjugate to $\operatorname{diag}(1,\xi)$.
At the two places above $\ell$ it has ordered ordinary flags
with inertia weights and residual diagonal characters
\[
 \begin{array}{c|c|c}
    &\text{inertia weights}&\text{residual ordered characters}\\ \hline
   v\in P &(k+t,-t)&(\bar\xi,1)\\
   v\notin P &(k,0)&(1,\bar\xi).
 \end{array}                                                   \tag{SX.11}
\]
There is an actual regular algebraic cuspidal representation
$\pi_0$ of $\mathrm{GL}_2(\mathbb A_{F_0})$ attached to $r$ in
(SX.1), with infinitesimal weights
$\frac12+(k+t,-t)$ and $\frac12+(k,0)$.
Every finite component of $\pi_0$ is a principal series with no
linked pair of inducing characters.  Thus the literal source
$\texttt{Mix\allowbreak Dat}$ and $\texttt{PR}$ conclusions are supplied.
\end{proposition}
\begin{proof}
Start with SX2 and its residual ordinary maximal ideal.
Fix the algebraic central character, including its norm power;
equivalently, fix the determinant $\chi^k\xi^{-1}$ of the usual
covariant Hilbert Galois representation.  In the notation of
Hida's control theorem the weight coordinates are initially
$(\kappa_1,\kappa_2)=(0,k-1)$ at both embeddings.  They may vary
with their sum fixed at $k-1$.  This parallel sum is exactly the
condition imposed by totally positive central units.
Here $F_0$ is real quadratic, so one may use the quaternion
algebra ramified at its two real places and split at every finite
place.  Jacquet--Langlands preserves the finite Hecke operators
and the prescribed tame level.  Its definite ordinary module
satisfies the finite projectivity and control statements of
SX-O (Hida, Lemmas 2.1--2.2); passing to sufficiently deep
$\ell$-level removes stabilizers.  The descent in Theorem 1.2
uses faithful flatness of weight algebras, so it does not require
dividing by a tame-level stabilizer order.
Localize the cuspidal
ordinary Hecke algebra at the maximal ideal of SX2 and take its
completed nonzero factor $B$.  The relevant component of the
weight algebra $\Lambda$ is a power-series domain.  By SX-O,
$B$ is finite and torsion-free over $\Lambda$.
In particular $\Lambda\to B$ is injective.

Choose $t\ge2$ divisible by $\ell-1$ and, if necessary, by an
arbitrarily high power of $\ell$.  Specialize the weights to
$(\kappa_1,\kappa_2)=(-t,k+t-1)$ at $P$ and $(0,k-1)$ outside $P$.
They have the required central sum and the same torsion character
modulo the maximal ideal.  Integral lying-over for
$\Lambda\hookrightarrow B$ gives a prime over this arithmetic
weight which avoids $\ell$; a finite extension of the specialization
field contains its residue field.  The control isogeny identifies
this characteristic-zero point with a classical cuspidal nearly
ordinary eigenform.  It has the original tame level and central
character.  The normalized $U_v^\circ$ and the away Hecke
operators retain the residual values from SX2.  This proves the
existence of the new form at the required nonparallel weights;
it is not an assumption about a hypothetical component through a
Galois deformation point.

Let $r_f$ be the usual covariant Galois representation of this
form.  Hida's Theorem 3.1 writes its ordinary flag as
$\left(\begin{smallmatrix}\beta&*\\0&\alpha\end{smallmatrix}\right)$,
where $\alpha$ has inertia $\chi^{\kappa_1}$ and $\beta$ has
inertia $\chi^{\kappa_2+1}$.  The normalized unit eigenvalue is
that of the \emph{quotient} $\alpha$ after removing its algebraic
inertia weight.  It has residual character $\bar\xi^{-1}$ on
$P$ and $1$ elsewhere.  These assertions hold on the entire
decomposition group: they agree on inertia and on an unramified
Frobenius lift, which together topologically generate it.
Define
\[
                         r=\chi^k r_f^\vee.
\]
The annihilator of the ordinary subline is the subline in the
dual.  Its character is $\alpha^{-1}$, so the upper character
of $r$ has weight $k+t$ and residual character $\bar\xi$ on
$P$, and weight $k$ and residual character $1$ elsewhere.
The lower characters are determined by the determinant.  Thus
the ordered flags are exactly (SX.11), with no reversal hidden
in the ordinary terminology.  Good-prime characteristic
polynomials and Chebotarev give
$\det r_f=\chi^k\xi^{-1}$ and
$\bar r_f^{\mathrm{ss}}=1\oplus\bar\xi^{-1}$.
Dualizing and twisting, with $\bar\chi^k=1$, proves (SX.10).
All characters are continuous; their finite-image parts have open
kernels.  The ordinary theorem is being used only at the
characteristic-zero specialization, so no assumption of residual
irreducibility is present.

Here are two elementary irreducibility checks.  If $r$ had a
global invariant line, its character would be unramified outside
finitely many places and would have one of the weights $k+t,-t$
at the first embedding and one of $k,0$ at the second.  A global
character of a totally real field with such algebraic local
inertia has parallel weights.  To see this directly, apply class
field theory to totally positive global units sufficiently close
to $1$ at the finitely many ramified places.  The product of their
embedding powers, with these integer weights, is $1$.  Apply the
fixed $\iota$ and the real logarithm.  Dirichlet's unit lattice
spans the hyperplane of vectors with coordinate sum zero, so the
weights must all be equal.  The two sets here are disjoint because
$k,t\ge2$.  This proves absolute irreducibility.

Suppose next that restriction to an open subgroup were reducible,
and replace that subgroup by its open normal core $N$.
Inertia above $\ell$ has open cyclotomic image, even after
intersection with $N$; hence $r(N)$ contains an element with two
distinct eigenvalues, using either nonzero weight difference in
(SX.11).  A line invariant under $N$ is therefore one of its two
eigenlines.  Normality makes the translates of that line again
$N$-invariant, so the full group permutes these two lines.
Irreducibility makes the permutation nontrivial.
Let $J$ be the index-two subgroup fixing each line.  Elements
outside $J$ have trace zero, hence reduction of (SX.10) gives
$\bar\xi=-1$ on that coset.  Multiplication by any element of
$J$ then gives $\bar\xi=1$ on $J$.  Thus $\bar\xi^2=1$
everywhere, contradicting SX1.  The restriction to every open
subgroup is irreducible.

We give the away local-type comparison carefully.  At a place
outside $\ell\mathfrak c$, the ordinary Galois theorem already
gives unramifiedness.  At a prime dividing $\mathfrak c$, the
automorphic conductor is at most $a(\xi)$, and its central
character has conductor exactly $a(\xi)$.  Thus both conductors
are equal.  A two-dimensional Weil--Deligne parameter with this
equality is a direct sum of two characters, exactly one
unramified, with zero monodromy.  Indeed an irreducible Weil
parameter has no inertia invariants: its tame conductor term is
two, whereas that of its determinant is at most one, and on
every higher ramification group determinant codimension is at
most vector codimension.  The conductor inequality is then
strict.  A special parameter also has strictly larger conductor
than its determinant.  Finally for two ramified characters the
conductor sum is strictly larger than the conductor of their
product.  These observations prove the assertion.  The remaining
ramified character restricts to $\xi$ or its inverse according
to the contravariant normalization.

To transfer this assertion to the actual Galois representation,
choose an imaginary quadratic extension of $\mathbb Q$ splitting
at the rational prime under the place being inspected.  Its
compositum $E/F_0$ is CM quadratic and splits at that place.
Such a quadratic field exists by the elementary congruence
conditions on a negative squarefree discriminant.
The cyclic base change of the Hilbert cusp is cuspidal: a
nontrivial self-twist would give a reducible restriction of $r$
to this open subgroup, just excluded.  Its central character
provides the polarization.  Multiplying the polarization character
by the odd quadratic character of $E/F_0$, if required, changes
its real-place sign and leaves its norm pullback unchanged.
SX-G applies.  Its representation equals $r|_{G_E}$ by good
Frobenius characteristic polynomials and semisimplicity.
The split local group is the original one.  Semisimplified local
compatibility therefore gives the inertial characters $1,\xi$.
They are distinct at a ramified conductor place; monodromy
commutes with inertia and hence is diagonal, so its nilpotence
forces it to vanish.  Thus the actual inertia representation,
not just a list of Frobenius traces, is
$\operatorname{diag}(1,\xi)$.  This argument avoids imposing
an unavailable finite discrete-series place in an even-degree
Hilbert local--global theorem.

At $v\mid\ell$ the classical form has Iwahori level and
unramified central character.  Its local component is either
unramified principal series or an unramified twist of Steinberg.
The Steinberg possibility has normalized ordinary slope
$(k-1)/2$ outside $P$ and $(k-1)/2+t$ on $P$, both positive.
It is excluded by the unit $U_v^\circ$ eigenvalue.  Thus it is
unramified principal series there as well.

To specify the attachment normalization, let $\pi_f$ be the
unitary cuspidal representation of this form and put
\[
              \pi_0=\pi_f\otimes
                         |\det|^{(k+1)/2}.          \tag{SX.12}
\]
If $\lambda_1,\lambda_2$ are the arithmetic eigenvalues of
$r_f$ at an unramified place, the Satake parameters of
$\pi_f$ are $q^{-k/2}\lambda_i$.  The geometric eigenvalues
of $r=\chi^k r_f^\vee$ are $q^{-k}\lambda_i$.
The parameters of (SX.12) are
$q^{-(k+1)/2}q^{-k/2}\lambda_i=q^{-k-1/2}\lambda_i$.
Multiplication by $q^{1/2}$ therefore gives exactly (SX.1).
The infinitesimal characters are $\frac12+(k+t,-t)$ and
$\frac12+(k,0)$, proving regular algebraicity.

We also record a local estimate needed in SX6.  For a unitary
irreducible principal series of $\mathrm{GL}_2$ with one
unramified and one ramified inducing character, both characters
are unitary: comparison with its conjugate dual either fixes
each character or exchanges them, and the exchange would force
their restrictions to units to agree.  At a spherical place write
the Satake parameters as $a,b$.  Conjugate duality and the unitary
central character imply either $|a|=|b|=1$, or
$(a,b)=(zu,z^{-1}u)$ with $z>1$ and $|u|=1$.
On the unit spherical vector, the normalized Hecke operator is
$q^{-1/2}$ times a sum of $q+1$ unitary translates.  Therefore
\[
                   |a+b|\le q^{1/2}+q^{-1/2}.
\]
It follows that $z\le q^{1/2}$.  Equality would force all these
translates to span the same line, by equality in the Hilbert-space
triangle inequality.  That line, fixed by the maximal compact
subgroup and stable under its uniformizer double coset, is stable
under $\mathrm{GL}_2$.  This would make the representation
one-dimensional, impossible for a cuspidal local component.
Thus $q^{-1/2}<|a|,|b|<q^{1/2}$.
After (SX.12), all finite inducing characters $\nu$ satisfy
\[
 q_v^{s-1/2}<|\nu(\varpi_v)|<q_v^{s+1/2},
                  \qquad s=-(k+1)/2.               \tag{SX.13}
\]
The same inequalities exclude a quotient $\nu_1/\nu_2=|\cdot|^{\pm1}$,
so these are the actual unlinked principal series required by
$\texttt{Plain\allowbreak R}$.
\end{proof}

\subsection{An integral induced representation with specified entries}

\begin{proposition}[SX4: the normalized Ribet lattice and mixed induction]
\label{claim:SX4}
After conjugating $r$, there is a continuous representation
$R:G_{\mathbb Q}\to\mathrm{GSp}_4(\overline{\mathbb Q}_\ell)$
with multiplier $\mu=\chi^k\xi\varphi$, abstractly
$\operatorname{Ind}_{G_{F_0}}^{G_{\mathbb Q}}r$, such that its
restriction to $G_F$ is integral, its diagonal reduces to
$(\bar p_i)$, and its off-diagonal residual support is contained
in
\[
        vV=\{w_\beta,-w_\alpha,-w_\vartheta\}.
                                                            \tag{SX.14}
\]
It has a unit entry of weight $-w_\alpha$, or it has both a unit
entry of weight $w_\beta$ and a unit entry of weight
$-w_\vartheta$.  Also $R(t_0)$ is diagonal.
\end{proposition}
\begin{proof}
Let $u=t_0^2\in H$.  Its two residual eigenvalues are distinct
by (SX.5).  Diagonalize $r(u)$ with its first eigenvalue reducing
to $1$ and its second to $\xi(u)$.  For $h\in H$, compare
$\operatorname{tr}r(h)$ and $\operatorname{tr}r(uh)$ with their
residual values.  The determinant of the resulting two-by-two
system is the unit difference of the two eigenvalues of $r(u)$.
It follows that
\[
       r(h)_{00}\equiv1,\qquad r(h)_{11}\equiv\xi(h).
                                                            \tag{SX.15}
\]
Comparing the $(0,0)$ entry of $r(hh')=r(h)r(h')$ then gives
$|r(h)_{01}r(h')_{10}|<1$ for every $h,h'\in H$.
Since $r|_H$ is irreducible, its upper-right entry is not
identically zero.  Compactness of $H$ and continuity give a
largest absolute value of that entry, attained at some $h_*$.
Conjugation by a diagonal matrix makes this maximum $1$.
The preceding product inequality now gives
\[
 r(h)_{00}\equiv1,\quad r(h)_{11}\equiv\xi(h),\quad
 |r(h)_{10}|<1,\quad |r(h)_{01}|\le1,
 \qquad |r(h_*)_{01}|=1.                            \tag{SX.16}
\]
The diagonalization of $r(u)$ is preserved.  This is an explicit
lattice normalization; it requires no assertion that an arbitrary
compact subset of $\overline{\mathbb Q}_\ell$ lies in a finite field.

Choose a square root $a$ of $r(u)_{00}$ reducing to $p_3(t_0)$,
and put $b=\mu(t_0)/a$.  The two choices of square root are
distinguished because $\ell\ne2$.  Equations (SX.5)--(SX.6)
give $\bar b=\bar p_0(t_0)$, and the determinant identity gives
$r(u)=\operatorname{diag}(a^2,b^2)$.
Both $a$ and $b$ are units.  Set $\gamma=b/a$.

In the usual induced basis, $h\in G_{F_0}$ acts by
$\operatorname{diag}(A,B)$, where
$A=r(h)$ and $B=r(t_0^{-1}ht_0)$, and $t_0$ acts by
$\left(\begin{smallmatrix}0&\operatorname{diag}(a^2,b^2)\\1&0\end{smallmatrix}\right)$.
Change basis by
\[
 P=\begin{pmatrix}
 0&a&0&a\\ b&0&-b&0\\0&-1&0&1\\1&0&1&0
 \end{pmatrix}.
                                                            \tag{SX.17}
\]
Direct multiplication gives $R(t_0)=\operatorname{diag}(b,-a,-b,a)$
and
\[
R(h)=\frac12\begin{pmatrix}
A_{11}+B_{11}&A_{10}/\gamma-B_{10}&B_{11}-A_{11}&A_{10}/\gamma+B_{10}\\
\gamma A_{01}-B_{01}&A_{00}+B_{00}&-\gamma A_{01}-B_{01}&A_{00}-B_{00}\\
B_{11}-A_{11}&-A_{10}/\gamma-B_{10}&A_{11}+B_{11}&-A_{10}/\gamma+B_{10}\\
\gamma A_{01}+B_{01}&A_{00}-B_{00}&-\gamma A_{01}+B_{01}&A_{00}+B_{00}
\end{pmatrix}.                                                \tag{SX.18}
\]
The group law follows from the induced representation and this
one fixed basis change.  For the multiplier, use the alternating
form $\operatorname{diag}(J_2,abJ_2)$ on the induced space, where
$J_2=\left(\begin{smallmatrix}0&1\\-1&0\end{smallmatrix}\right)$.
Its Gram matrix in (SX.17) is $-2abJ$, with $J$ the original
four-dimensional alternating matrix.  Both the subgroup and
$t_0$ act on it with multiplier $\mu$.  This proves the exact
symplectic identity, including the quadratic sign $\varphi$.

For $h\in H$, both $A$ and $B$ satisfy (SX.16), since $H$ is
normal in $G_F$ and $\xi$ is a global character.  Formula (SX.18)
therefore is integral, reduces on the diagonal to
$(\bar\xi,1,\bar\xi,1)$, and has possible nonzero off-diagonal
entries only at $(1,0),(1,2),(3,0),(3,2)$.
Their weights are respectively
$-w_\alpha,w_\beta,-w_\vartheta,-w_\alpha$.
All other entries reduce to zero by (SX.16).
At $h_*$, put $x=\gamma A_{01}$ and $y=B_{01}$.
The $(1,0)$ and $(3,0)$ entries are $(x-y)/2$ and $(x+y)/2$;
their sum is the unit $x$.  Hence one is a unit.
In the second case the $(1,2)$ entry is the negative of the
$(3,0)$ entry, so it too is a unit.
This proves precisely the asserted alternative.

Every remaining element of $G_F$ is an element of $H$ times
$t_0$.  Right multiplication by the diagonal unit matrix
$R(t_0)$ preserves the support and integrality.  Its residual
diagonal is correct by (SX.5) and the choices of $a,b$.
Continuity follows from the two open cosets and the explicit
formulas.  This proves all assertions on the original group.
\end{proof}

\subsection{The actual local flags and the correcting character}

\begin{proposition}[SX5: local flags, inertia kernel, and the exact weights]
\label{claim:SX5}
There is a continuous character $\delta:G_F\to1+\mathfrak m_{\bar{\mathbb Q}_\ell}$
and, for each original index $v\in T.\iota$, an integral lower
unipotent symplectic matrix $n_v$ such that
\[
        Q_{\bar{\mathbb Q}_\ell}(g)=\delta(g)R(g),
        \qquad\delta^2\mu=\vartheta.              \tag{SX.19}
\]
These matrices satisfy every global and local equation of
$(T.\mathrm{fr}(\varnothing)).\texttt{Is\allowbreak Pair}$.
Their residual local support is contained in
$qQ=\{-w_\alpha,-w_\Psi,-w_\vartheta\}$ and
$(n_v)_{20}$ is a unit.  Put
\[
 m=(k+t,k,0,-t),\qquad b_0=e+1-k,\qquad h_i=2m_i+b_0.
                                                            \tag{SX.20}
\]
Then the actual pair has source weights $h_i$, with
$m_i-m_{i+1}\ge2$.  There is one integer $N>0$ such that
\[
                (\delta(g)^2\chi(g)^{-b_0})^N=1
                         \quad(g\in G_F).          \tag{SX.21}
\]
\end{proposition}
\begin{proof}
The character $\vartheta/\mu$ is continuous and reduces to $1$
by SX1.  Squaring is an automorphism of the group of units
congruent to $1$ when $\ell$ is odd.  Its inverse is continuous:
on any finite coefficient field it is the unique Hensel root
congruent to $1$, or equivalently the convergent binomial series.
Uniqueness makes the pointwise root multiplicative.  This defines
$\delta$, proves (SX.19), and preserves every residual entry
from SX4.

Fix an original valuation above $\ell$.  It splits in $F_0$,
so its decomposition group lies in $G_{F_0}$; conjugation by
$t_0$ exchanges the two places.  At the one whose top residual
character is $\xi$, a common ordinary eigenline can be written
as $(u,1)$ with $|u|\le1$.  At the other it can be written as
$(1,w)$ with $|w|<1$.  To check these assertions, use an element
of the local group for which $\bar\xi\ne1$, supplied by SX1.
In the eigenvector equation its eigenvalue minus the opposite
diagonal entry is a unit.  Formula (SX.16) then bounds the ratio
of the two coordinates by the upper-right integral entry in the
first case and by the strictly small lower-left entry in the
second.  The required denominator coordinate cannot be zero.
The same line works for the whole local group because it is the
ordinary invariant line.

For $n(p,q,r,s,z)$ write
\[
 n(p,q,r,s,z)=
 \begin{pmatrix}1&0&0&0\\p&1&0&0\\q&s&1&0\\r&z&-p&1\end{pmatrix}.
                                                            \tag{SX.22}
\]
It is symplectic exactly when $z+ps-q=0$.
If the first block in (SX.18) has eigenline $(u,1)$, choose
\[
 n_v=n\left(\gamma u,-1,\gamma u,
       \frac{-2w}{1+\gamma uw},
       \frac{-(1-\gamma uw)}{1+\gamma uw}\right).     \tag{SX.23}
\]
If that eigenline occurs in the second block, choose instead
\[
 n_v=n\left(-u,1,u,
       \frac{-2w}{\gamma+uw},
       \frac{\gamma-uw}{\gamma+uw}\right).           \tag{SX.24}
\]
The denominators are units because $|\gamma|=1$ and $|uw|<1$.
These formulas satisfy $z+ps-q=0$, all their entries are
integral, their $(2,1)$ entry is strictly small, and their
$(2,0)$ entry is $-1$ or $1$.
This proves precisely the residual support and nonzero entry
assertions.

For clarity, triangularization in (SX.23) is an identity in six
scalars, not a dimension argument.  If $A(u,1)^t=a_1(u,1)^t$,
$\det A=a_1a_2$, then
\[
 A=\begin{pmatrix}
 a_2+ux&(a_1-a_2-ux)u\\x&a_1-xu
 \end{pmatrix}.
\]
If $B(1,w)^t=b_1(1,w)^t$, $\det B=b_1b_2$, then
\[
 B=\begin{pmatrix}
 b_1-yw&y\\(b_1-b_2-yw)w&b_2+yw
 \end{pmatrix}.
\]
Substitution into (SX.18) and (SX.23) makes
$n_v^{-1}R(g)n_v$ upper triangular.  Its diagonal is
$(a_1,b_1,b_2,a_2)$.
Interchanging the blocks gives (SX.24), whose diagonal is
$(b_1,a_1,a_2,b_2)$.
In both cases (SX.11) gives the inertia exponents
$m=(k+t,k,0,-t)$ and residual diagonal
$(\bar\xi,1,\bar\xi,1)$.
Thus the original local character order is respected.

We next verify the global inertia kernel in the frame.
Let $g\in G_F$ be an inertia element away from $\ell$ with
$D.\bar r(g)=1$.  The clause $\texttt{Two.\allowbreak unr}$ puts it in
$G_{F_0}$.  Away from $\ell$ one has $\chi(g)=\omega(g)=1$;
the residual diagonal then gives $\bar\xi(g)=1$.
The order of $\xi$ is prime to $\ell$, so $\xi(g)=1$.
Also $\texttt{Ok.\allowbreak kb}$ gives $\psi(g)=1$.
By SX3 both $r(g)$ and its $t_0$-conjugate are the identity.
Hence $R(g)=1$ by (SX.18), and (SX.19) gives $\delta(g)^2=1$.
The root congruent to $1$ is $1$, so $Q_{\bar{\mathbb Q}_\ell}(g)=1$.
The multiplier, diagonal-at-$g_0$, and local flag clauses were
already proved in (SX.19), SX4, and (SX.23)--(SX.24).

Dividing the two multipliers in (SX.19) gives the exact identity
\[
       \delta^2=\chi^{e+1-k}\omega^{-1}\psi\varphi^{-1}
                                  \quad\hbox{on }G_F.       \tag{SX.25}
\]
Choose $a>0$ prime to $\ell$ annihilating $\psi|_{G_F}$, as
provided by $\texttt{Deep}$.  We may take
$N=2a(\ell-1)$ in (SX.21).
On the open subgroup $U=\ker\omega\cap\ker\psi$ of $G_F$,
the inertia at every place above $\ell$ also kills $\varphi$
by splitting, so $\delta^2=\chi^{b_0}$ there.
The local diagonal characters $\delta\chi^{m_i}$ therefore
have square $\chi^{2m_i+b_0}$, exactly the source's
$\texttt{Wts}$ convention.

Finally the source's literal set in the local frame is
\[
 \operatorname{muOf}(D_v)=
 \{s\in D_v:\ s^{\ell^a}\in
       \overline{[D_v,D_v]}\text{ for some }a\ge0\}.
\]
It is contained in $I_v\cap\ker\chi$; equality is not needed.
Indeed $D_v/I_v\simeq\widehat{\mathbb Z}$ is abelian and has
no torsion, so the indicated power relation puts $s$ in $I_v$.
The continuous abelian character $\chi$ kills the closed
commutator subgroup, giving $\chi(s)^{\ell^a}=1$.
For odd $\ell$, $\mathbb Z_\ell^\times$ has no nontrivial
$\ell$-power torsion: its torsion has order dividing $\ell-1$
and its principal units are torsion-free, by the $\ell$-adic
logarithm.  Hence $\chi(s)=1$.
For every such $s$, one has $\chi(s)=\omega(s)=1$.  The clause
$\texttt{Ok.\allowbreak ine\allowbreak L}$ makes its residual diagonal trivial, hence
$\bar\psi=1$.  The prime-to-$\ell$ order supplied by
$\texttt{Deep}$ implies $\psi=1$ on this subgroup.
Equation (SX.25) gives $\delta=1$, and all diagonal local
characters are $1$.  This is the exact frame condition, not just
an equality on a smaller chosen inertia subgroup.
The adjacent differences of $m$ are $t,k,t$, all at least two.
\end{proof}

\subsection{Irreducibility and the actual automorphic witness}

\begin{proposition}[SX6: cyclotomic irreducibility and the local window]
\label{claim:SX6}
The representation $R$ is irreducible on
$G_L\cap\ker\chi$.  There is a cuspidal regular algebraic
representation $\Pi$ of $\mathrm{GL}_4(\mathbb A_L)$ attached
to $R|_{G_L}$ in convention (SX.1).  It satisfies the literal
\texttt{Imaginary\allowbreak Quadratic\allowbreak Eta\allowbreak Character.\allowbreak All\allowbreak Win} condition:
there is a single real number $s_4$ such that every finite local
component is toral and every character $\nu$ in its cuspidal
support satisfies
\[
 q_v^{s_4-1/2}<|\nu(\varpi_v)|<q_v^{s_4+1/2}.
                                                            \tag{SX.26}
\]
One can take $s_4=-(k+3)/2$.
\end{proposition}
\begin{proof}
Put $V=G_L$ and $U=G_L\cap\ker\chi$.
Write $\eta=\bar\Psi|_{G_F}$ for the quadratic character
cutting out $F_0$.  We prove directly that $\eta|_U\ne1$.
Suppose $\eta|_U=1$.  If $g\in V$ and $\omega(g)=1$,
then $\chi(g)\equiv1\pmod\ell$, and
$\chi(g^{\ell^n})\longrightarrow1$.
Compactness gives a cluster point $h\in V$ of the sequence
$g^{\ell^n}$.  The character $\eta$ is quadratic and $\ell$
is odd, so each term lies in the same closed $\eta$-coset
as $g$.  Thus $\chi(h)=1$ and $\eta(h)=\eta(g)$.
The assumption on $U$ now gives $\eta(g)=1$.
Consequently $\eta|_V$ is quadratic and trivial on
$V\cap\ker\omega$.  Multiplication by $\eta$ permutes the
four residual characters: in (SX.7) it exchanges indices
$0\leftrightarrow2$ and $1\leftrightarrow3$.
The precise no-self-twist clause $\texttt{Ok.\allowbreak g3'}$ therefore
forces $\eta|_V=1$.  But $\texttt{Ok.\allowbreak g3}$ applied to
$w_\Psi$ asserts that it is nontrivial on $V$.
This contradiction proves the required nontriviality.
Thus $U_0=U\cap G_{F_0}$ has index two in $U$.

The restriction $r|_{U_0}$ is irreducible.  Here is the precise
infinite-subgroup argument.  The quotient of the open group
$G_{F_0}\cap G_L$ by $U_0$ is a subgroup of
$\mathbb Z_\ell^\times$, hence procyclic since $\ell$ is odd.
The finite-dimensional algebra spanned by the restricted image
has radical stable under the whole open group.  Its radical
annihilates a nonzero subspace and hence, by irreducibility of
that open-group representation, annihilates the entire space.
The restriction is semisimple.  Its finitely many isotypic
summands have an open stabilizer; more than one summand would
contradict SX3.  A repeated single constituent has a projective
multiplicity representation of the procyclic quotient.
A lift of a topological generator has an eigenline; polynomial
identities and continuity make that line invariant under the
whole quotient.  Irreducibility forces multiplicity one.
This is the same finite-algebra/procyclic argument proved in RX4.

The two $U_0$ constituents $r$ and $r^{t_0}$ are not isomorphic.
Suppose they were.  Their determinants agree on the surrounding
open group, since $\det r$ is the restriction of a global
character.  Fix an intertwiner on $U_0$.  Schur's lemma shows
that, for every element $x$ of this surrounding group, the
comparison of the two extensions is scalar.  Its square is
$1$ by the determinant equality.  Consequently the two
representations have the same characteristic polynomial at
$x^2$.  At a place above $\ell$, take an inertia element in
this open group with cyclotomic value of infinite order.
Such elements exist because the local cyclotomic image has
finite index in $\mathbb Z_\ell^\times$.
The two polynomials, evaluated at a suitable positive even
power, have eigenvalues
\[
     \{\chi(x)^{k+t},\chi(x)^{-t}\},\qquad
     \{\chi(x)^k,1\}.
\]
The second contains $1$ and the first does not, since
$t,k+t>0$.  This is a contradiction.  The usual index-two
induction argument now proves $R|_U$ irreducible: its restriction
to $U_0$ has two nonisomorphic irreducible constituents exchanged
by the other coset, so no proper invariant sum is possible.

In particular $F_0$ is not contained in $L$.  Put
$E=F_0K$ and $M=F_0L$.  Then $E/F_0$ is CM quadratic,
$M/E$ is abelian, and $M/L$ is quadratic.
Apply cyclic base change successively to $\pi_0$ to obtain
$\pi_M$.  All intermediate base changes are cuspidal:
the self-twist obstruction would contradict irreducibility of
$r$ on the corresponding open subgroup.  The representations
$\pi_M$ and its conjugate over $L$ are distinct, since otherwise
their good Frobenius polynomials would identify $r$ and its
conjugate on $G_M$, contrary to the preceding argument on
$U_0$.  Automorphic induction therefore gives a cuspidal
$\operatorname{AI}_{M/L}(\pi_M)$.
Set
\[
             \Pi=\operatorname{AI}_{M/L}(\pi_M)
                                      \otimes|\det|.          \tag{SX.27}
\]
At a split prime the rank-two Satake multisets concatenate;
at an inert prime one takes both square roots of each
rank-two parameter.  Formula (SX.1) shows that before the
last twist this attaches $R\chi^{-1}$; the $|\det|$ twist
attaches $R$.  At infinity, the two extensions of an embedding
of $L$ restrict to the two different real embeddings of $F_0$.
The infinitesimal character in (SX.27) is
\[
               \tfrac32+\{k+t,k,0,-t\}.
\]
Its four entries are distinct, giving regular algebraicity.

We verify every finite place in (SX.26).  At a place split in
$M/L$ the local parameters are the two rank-two principal-series
parameters restricted to the local extension.  Pullback by a
local norm preserves the real exponents in (SX.13), because
$q_w=q_v^f$ and
$|\nu(N_{M_w/L_v}\varpi_w)|=|\nu(\varpi_v)|^f$.
At a nonsplit place, the quadratic field $F_0$ is unramified
over the relevant local field: $\texttt{Two.\allowbreak unr}$ says its
quadratic character kills every original inertia subgroup.
Moreover this place is not above $2$ or $\ell$, since both
primes split in $F_0$.  The rank-two representation and the
rank-two local automorphic representation are therefore
unramified there, by SX1 and SX3.  Inducing an unramified
character through an unramified quadratic extension gives
the two unramified characters whose uniformizer values are
the two square roots.  Their absolute exponents are unchanged.
Thus in both cases the four local inducing characters have
the strict inequalities (SX.13), with one common center.
No quotient of two of them can equal $|\cdot|^{\pm1}$:
its real exponent lies strictly between $-1$ and $1$.
The principal-series induction is consequently irreducible
and has exactly these characters as its cuspidal support.
Finally the twist in (SX.27) subtracts $1$ from the center,
giving $s_4=s-1=-(k+3)/2$.
This proves the literal toral and strict-window clauses at
all finite places, including those above $\ell$.
\end{proof}

\subsection{Return to a finite coefficient ring and close the endpoint}

\begin{proposition}[SX7: exact finite coefficient return]
\label{claim:SX7}
The actual pair $(\delta R,(n_v))$ of SX5 descends to a
continuous pair $Q$ over the ring $O_{c_1}$ of one finite
extension $c_1\ge c$.  All its pair equations hold over that
ring.  Its actual reduction, rather than an abstractly
isomorphic representation, satisfies $\texttt{In\allowbreak V}$ and the
root-entry alternative in SX4.
\end{proposition}
\begin{proof}
Here is the matrix version of the finite-field capture argument
in RV4.  There are only countably many finite subextensions of
$\overline{\mathbb Q}_\ell/\mathbb Q_\ell$: in each degree
Krasner's theorem gives finitely many isomorphism classes, each
with finitely many embeddings.  Enumerate them inside an
increasing sequence $E_j$.  Each is complete and hence closed
in $\overline{\mathbb Q}_\ell$.
For the actual continuous homomorphism
$\rho=\delta R$, the sets
\[
           H_j=\{g\in G_F:\rho(g)\in\mathrm{GL}_4(E_j)\}
\]
are closed subgroups and cover $G_F$; inversion preserves $E_j$
because it is given by adjugate divided by the nonzero determinant.
The compact Hausdorff Baire theorem makes some $H_j$ have
interior, hence makes it open.  Choose its finitely many coset
representatives and adjoin all their matrix entries to $E_j$.
The homomorphism law now puts the entries of \emph{every}
$\rho(g)$ in this single finite extension.  Thus the argument
uses a closed-subgroup cover, not compactness of an arbitrary
algebraic-valued set.
Adjoin also the original coefficient field and the entries of the
finitely many matrices $n_v$ to this field.
This yields one finite extension.  SX4--SX5 show that all
matrix entries have absolute value at most one, so they lie
in its valuation ring.  The continuous topology there is
the adic topology, giving the literal continuity required
by $\texttt{Seed\allowbreak Pair\allowbreak Alt}$.

Every pair equation is a polynomial identity with coefficients
in the original ring: the symplectic multiplier equation,
the killing of the frame inertia set, diagonalization at
$g_0$, the lower-unipotent equations for $n_v$, triangularity
on the original decomposition group, and triviality of its
diagonal on the original $\operatorname{muOf}(D_v)$ set.  Injectivity
of $O_{c_1}\hookrightarrow\overline{\mathbb Q}_\ell$
therefore descends them all.
Reduction also commutes with the symplectic adjoint and
conjugation.  An entry reduces to zero exactly when its
absolute value is less than one; a unit entry remains
nonzero.  Thus (SX.14), (SX.23)--(SX.24), and their unit
$(2,0)$ entries give literal $\texttt{In\allowbreak V}$.
The entries $(1,0)$, or $(1,2)$ and $(3,0)$, give respectively
$\operatorname{Ent}(-w_\alpha)$, or
$\operatorname{Ent}(w_\beta)$ and
$\operatorname{Ent}(-w_\vartheta)$.
The diagonal character congruences and the local ordered
character congruences give the original residual pair record.
All these assertions concern the descended matrices themselves.
\end{proof}

\begin{theorem}[SX8: the literal \texttt{Seed\allowbreak IIAlt} supplier]
\label{claim:SX8}
For every original $\Omega$, prime $\ell$, coefficient datum
$c$, raw datum $D$, and tower $T$ satisfying (SX.2), and for
every fixed $\iota:\overline{\mathbb Q}_\ell\simeq\mathbb C$,
there exist a finite coefficient extension $c_1\ge c$ and an
actual pair $Q$ over $O_{c_1}$ such that
\[
 \begin{gathered}
 \texttt{Seed\allowbreak Pair\allowbreak Alt}(\iota,c\le c_1,Q),\qquad
 \texttt{In\allowbreak V}(\bar Q),\\
 \operatorname{Ent}(-w_\alpha,\bar Q.r)\quad\text{or}\quad
 \bigl(\operatorname{Ent}(w_\beta,\bar Q.r) 
       \text{and}\ \operatorname{Ent}(-w_\vartheta,\bar Q.r)\bigr).
 \end{gathered}                                                \tag{SX.28}
\]
The only retained leaves of this construction are the precisely
stated classical theorems above and the already reconstructed
coefficient-capture theorem.  In particular no mixed-seed,
ordinary-weight-change, or automorphic-witness predicate remains
as an unproved bespoke supplier.
\end{theorem}
\begin{proof}
Take the pair of SX7.  It is a pair for the unchanged frame
$T.\mathrm{fr}(\varnothing)$, is continuous, and has its
literal residual pair and (SX.28), by SX4--SX7.
Twisting by a scalar character does not change invariant
subspaces, so SX6 gives irreducibility of $\delta R$ on
the original $T.G'=G_L\cap\ker\chi$.
Use the weights $h_i=2m_i+b_0$ of (SX.20) and put
\[
          \kappa=\delta^{-1}|_{G_L},\qquad a=-b_0.
                                                            \tag{SX.29}
\]
Then $h_i-h_{i+1}=2(m_i-m_{i+1})$ with each difference
$m_i-m_{i+1}\ge2$, and $a+h_0=2m_0$ is even.
The character $\kappa$ is continuous.  For the same $N>0$
from (SX.21),
\[
       (\kappa(s)^2\chi(s)^{-a})^N=1
\]
for every inertia element at every valuation above $\ell$;
in fact it holds for every $s\in G_L$.
Twisting $(\delta R)|_{G_L}$ by $\kappa$ gives exactly
$R|_{G_L}$, so SX6 supplies the actual $\Pi$, the attachment,
regular algebraicity, and the required $\texttt{All\allowbreak Win}$.
These are all clauses of $T.\texttt{Is\allowbreak Seed\allowbreak Alt}$.

The order of choices is the literal one: $D,T$ and all their
data are given first, then $\iota$ is given, and only then are
the Hilbert form, $r$, $R$, $\delta$, $\Pi$, and $c_1,Q$
chosen.  The argument applies to every such given tower.
It does not choose a new tower after seeing the seed, or
require a new complex embedding at a later stage.
\end{proof}

\subsection{Primary reference locations and scope of the closure}

The classical references used above are the following.  The accompanying
machine-readable registry records their individual applications and
the exact pinned Lean definitions replaced or matched by this chapter.
The papers themselves are research references and are not redistributed.

\begin{itemize}
\item S.~Dasgupta, H.~Darmon and R.~Pollack,
\emph{Hilbert modular forms and the Gross--Stark conjecture},
Annals of Mathematics 174 (2011), 439--484,
\S2.2, Proposition 2.1, especially formula (50), pp.~458--459;
\S2.3, p.~466, for the usual cuspidal/Eisenstein decomposition.
\href{https://annals.math.princeton.edu/wp-content/uploads/annals-v174-n1-p12-p.pdf}{Primary article (PDF)}.
The later Gross--Stark hypotheses are not used in the general
Eisenstein-series proposition.

\item C.~P.~Mok and F.~Tan,
\emph{Overconvergent Families of Siegel--Hilbert Modular Forms},
Canadian Journal of Mathematics 67 (2015), 893--922,
Lemma 2.10 and Proposition 2.11, pp.~905--906.
\url{https://doi.org/10.4153/CJM-2014-017-9}.
The genus-one determinant Hasse invariant and Serre lifting are
the only results used; no eigenvariety or overconvergent
automorphy assertion is imported.

\item H.~Hida, \emph{Control of nearly ordinary Hecke algebras},
Luminy lecture, 18 July 2007, pp.~3--8,
Theorem 1.2, Lemmas 2.1--2.2, Theorem 2.4,
Corollary 2.5 and Theorem 3.1 with Remark 3.1.
\url{https://www.math.ucla.edu/~hida/Lum07.pdf}.
The central character is fixed, $F_0$ is real quadratic,
$\ell$ splits, and only characteristic-zero classical
specializations are needed.

\item J.~Arthur and L.~Clozel,
\emph{Simple Algebras, Base Change, and the Advanced Theory
of the Trace Formula}, Annals of Mathematics Studies 120 (1989),
Chapter III, Theorems 4.2, 5.1 and 6.2,
pp.~202--203, 212--213 and 215--217.
\url{https://www.claymath.org/library/cw/arthur/pdf/30.pdf}.
For strong local compatibility also use G.~Henniart,
\emph{Induction automorphe globale pour les corps de nombres},
Bulletin de la Soci\'et\'e Math\'ematique de France 140 (2012),
1--17, Theorems 4--5 and \S3.7, pp.~6--7 and 16:
\url{https://www.numdam.org/article/BSMF_2012__140_1_1_0.pdf}.
Their unitary formulations apply after removing the displayed
common norm twist; base change, induction and restoring that
twist commute by their local character formulas.

\item T.~Barnet-Lamb, T.~Gee, D.~Geraghty and R.~Taylor,
\emph{Potential automorphy and change of weight},
Annals of Mathematics 179 (2014), 501--609,
\S2.1 and Theorem 2.1.1(1)--(2), pp.~535--538.
\href{https://annals.math.princeton.edu/wp-content/uploads/annals-v179-n2-p03-p.pdf}{Primary article (PDF)}.
Only the rank-two polarized CM Galois attachment and
away-from-$\ell$ compatibility are used.  The actual Hilbert
form has already been constructed before applying this theorem.

\item W.~Casselman, \emph{On some results of Atkin and Lehner},
Mathematische Annalen 201 (1973), 301--314,
Theorem 1 and its conductor computation, pp.~302--306.
\url{https://lesesvre.perso.math.cnrs.fr/newforms-references/casselman.pdf}.
The use is the exact equivalence between a prescribed
newvector level and the local conductor.  The principal-series
classification and its determinant/conductor normalization are
also in H.~Jacquet and R.~P.~Langlands,
\emph{Automorphic forms on GL(2)}, Chapter I, \S3:
\url{https://publications.ias.edu/sites/default/files/automorphic-forms-on-gl2_rpl_9.pdf}.

\item I.~N.~Bernstein and A.~V.~Zelevinsky,
\emph{Induced representations of reductive $p$-adic groups. I},
Annales scientifiques de l'\'Ecole Normale Sup\'erieure 10
(1977), 441--472, Theorem 4.2, p.~454.
\url{https://www.math.tau.ac.il/~bernstei/Publication_list/publication_texts/BZ-induced-1.pdf}.
We use its sufficient irreducibility criterion only for
four one-dimensional inducing characters.

\item J.~S.~Milne, \emph{Class Field Theory}, version 4.03,
Chapters I and VIII, for local and global reciprocity;
\emph{Algebraic Number Theory}, version 3.08,
the Dirichlet unit theorem and Theorem 8.31 (Chebotarev):
\url{https://www.jmilne.org/math/CourseNotes/CFT.pdf},
\url{https://www.jmilne.org/math/CourseNotes/ANT.pdf}.
The resulting elementary character and irreducibility arguments
are written out in SX1--SX6.
\end{itemize}

This closes the literal $\texttt{Seed\allowbreak IIAlt}$ supplier with its
original $\texttt{Ctx2}$, $\texttt{Ok}$, $\texttt{Two}$ and
$\texttt{Deep}$ hypotheses.  It does not claim the wider
$\texttt{Hypotheses\allowbreak H'}$ or $\texttt{Rank\allowbreak N}$ packages, general
rank-four ordinary support, $\texttt{Seed\allowbreak IVAlt}$, image density,
or the final counting theorems.  Those are separate dependency
edges, and the fixed $\iota$ in this endpoint remains the one
chosen before the seed.


% Version 0.29. Actual support realization, with the original valued point.
% Source pin: 7cbf7a0c32fca1badf05931879a37c3f8a2b03b4.
\section{From a case-II valued point to its promodular support}
\label{sec:symplectic-promodular-realization}

\status{This chapter constructs the actual case-II witness required by
\texttt{Pro\allowbreak Mod\allowbreak Alt}.  Its automorphic inputs are the literal
\texttt{Seed\allowbreak IIAlt} and \texttt{PMEx\allowbreak Alt}, with the same fixed embedding
$\iota$.  The cohomology, residual vanishing and Ferrero--Washington
interfaces used below have already been reconstructed in BX, BR, AX
and TW.  In particular, an automorphic residual seed is not being
identified with the original valued point: the component and
coefficient arguments transporting support between them are part of
the proof.  The case-III boundary is recorded separately.}

\subsection{Objects that must remain fixed}

Use the original $\Omega/\mathbb Q$, odd prime $\ell$, coefficient
DVR $O\subset\overline{\mathbb Q}_\ell$, actual datum $D$, and all
of its decomposition groups and inertia groups.  Put $G_F=
\operatorname{Gal}(\Omega/F)$.  A case-II datum means the exact
\texttt{Ctx2}: $D$ is good and is the datum
\texttt{Of\allowbreak Char} for an original host character in Case II, with
the prescribed parameters $(n_{\rm II},1,1)$ of AX5.  In
particular its $\xi$ is a primitive odd dyadic character, and
the original auxiliary factor $\xi_0$ equals $1$.

Fix once and for all $\iota:\overline{\mathbb Q}_\ell\simeq
\mathbb C$.  An ambient field $(K,V,\{a_{O'}\})$ and a small
pair $(r,u)$ have exactly the meaning in IX.  We suppose that
$(r,u)$ is an actual \texttt{KPair} and that $r_K$ is
\texttt{Star}, namely absolutely irreducible or the specified
nonsplit block extension.  Every later change of lattice is
conjugation over this same $K$.

Write
\[
 \alpha=(1,0),\quad \beta=(-1,1),\quad
 \tau=(1,1),\quad \sigma=(0,1),\qquad
 \mathcal R=\{\pm\alpha,\pm\beta,\pm\tau,\pm\sigma\}.
 \tag{UX.1}
\]
The names $\tau,\sigma$ in this chapter denote roots, not Galois
elements.  They are source \texttt{w\allowbreak Th,w\allowbreak Ps}.  The odd roots are
$\mathcal R_o=\{\pm\alpha,\pm\beta,\pm\tau\}$.  The matrix
weights, in the original basis $0,1,2,3$, are
\[
 (w_{ij})=
 \begin{pmatrix}
 0&\alpha&\sigma&\tau\\
 -\alpha&0&\beta&\sigma\\
 -\sigma&-\beta&0&\alpha\\
 -\tau&-\sigma&-\alpha&0
 \end{pmatrix}.
 \tag{UX.2}
\]
The three triangular patterns used below are
\[
 N_P=\{\beta,-\alpha,-\tau,-\sigma\},\quad
 N_M=\{-\alpha,-\beta,-\sigma,-\tau\},\quad
 Q=\{-\alpha,-\sigma,-\tau\},\quad
 V_0=\{\beta,-\alpha,-\tau\}.
 \tag{UX.3}
\]
Support in a root set means that every off-diagonal entry
whose weight is outside that set is zero.  An entry in a
root is \emph{present} if some such entry of some $r(g)$ is
nonzero.  This is the literal \texttt{Ent}; presence of an
entry of a local flag is stated separately.

A tower $T$ consists of the same label
$\mathrm{pr}\in\{0,1\}$, an imaginary quadratic field $K_0$,
a totally real abelian field $F$, their compositum $L=FK_0$,
a separating $g_0\in G_F$, and exactly one chosen place over
$\ell$ for each place of $F$, with the prescribed choice on
$K_0$.  Put $d=[F:\mathbb Q]$ and $s=\#\{v\mid\ell\}$.
The label permutes $(0,1,2,3)$ either trivially or by
$(2,3,0,1)$.  The relabelled characters are denoted $\chi_i$.
For a root $w=(a,b)$ put
\[
 \chi_w=(\chi_0/\chi_1)^a(\chi_0/\chi_2)^b.
 \tag{UX.4}
\]
The framed pair $P=(\rho,\{n_v\})$ has multiplier $D.\vartheta$,
$\rho(g_0)$ diagonal, and $n_v$ lower unipotent symplectic;
$n_v^{-1}\rho(G_{F_v})n_v$ is upper triangular.  Its local
characters are denoted $\psi_{v,i}$.  The frame kills the
specified inertia elements outside $\ell$ and a finite
auxiliary set $W$ whose original residual representation is
trivial.  It also imposes $\psi_{v,i}=1$ on the source's
$\mu_v$.  These are exactly $(T.\mathtt{fr}\,W).\mathtt{IsPair}$.
A residual pair has diagonal and ordinary characters
$\bar\chi_i$, is locally constant, and has either support
$N_P$ with flag support $Q$, or support $N_M$.

We retain all clauses of \texttt{Ok}, \texttt{Two},
\texttt{Fit}, \texttt{Deep} and \texttt{Le}.  In particular
$T\le T'$ means $K_0,\mathrm{pr}$ are unchanged,
$F\subset F'$, the quotient has $\ell$-power exponent, and
$g'_0$ restricts to an $\ell$-power of $g_0$.
\texttt{Deep} says $\psi|_{G_F}$ has order prime to $\ell$.
It does not merely assert that the layer index is large.

\subsection{The tower and the original lattice}

\begin{proposition}[UX1: the actual tame case-II tower and its extensions]
For every case-II $D$ and permitted label there exists an
actual tower $T_0$ satisfying \texttt{Ok}, \texttt{Two} and
\texttt{Tame\allowbreak L}.  For every finite set $W$ of source
\texttt{Splits} primes it has an extension $T_1$ with the
same properties and with \texttt{Dies}$\,(\ell,w,F_1)$ for
every $w\in W$.  From any such tower, and any prescribed
$B\in\mathbb N$, there is a cyclotomic layer $T\ge T_1$
which is deep and has $[F:\mathbb Q]\ge B$.
\end{proposition}

\begin{proof}
Separate the $\ell$-primary part of the Dirichlet character
$\psi$ into its factor at $\ell$ and the factor away from
$\ell$.  Let $\kappa$ be the Galois character of the latter.
It has $\ell$-power order and is unramified at $\ell$.  On
$\ker c_\ell\cap\ker\kappa$, the character $\psi$ has order
prime to $\ell$.  On inertia away from $\ell$, if
$\bar\Psi=1$ and $\kappa=1$, then $\psi=1$: the remaining
factor has order prime to $\ell$ and reduction is injective
on roots of that order.  This proves every clause of the
source's \texttt{Kap}, without changing $\psi$ or its field.

Choose three new primes and their quadratic characters
$\varphi,\eta,\eta'$, respectively even, odd, odd.  Each
character is ramified at its own new prime only; $\varphi$
splits the finite set $D.S$, $\eta$ splits $D.S$ and the
prime of $\varphi$, and $\eta'$ splits all these and the
prime of $\eta$.  For example choose the prime in the
required class modulo $8\prod_{q\in D.S}q$, with the extra
conditions for previously chosen primes.  Quadratic
reciprocity and the Chinese remainder theorem specify a
nonempty reduced residue class, and the prime-in-progression
theorem supplies a prime.  The character of
$\mathbb Q(\sqrt{\pm p})$, with its sign and congruence
chosen so the fundamental discriminant has only that prime,
has exactly these properties.  Equivalently this is the
quadratic, degree-one specialization of Chebotarev.

Inside the original $G_{\mathbb Q}$ define
\[
 H=\{g:\bar\Psi(g)=\varphi(g),\quad
                 \eta(g)=\eta'(g),\quad\kappa(g)=1\}.
 \tag{UX.5}
\]
It is open, normal and contains the commutator.  Let $F$ be
its fixed field and $K_0$ the fixed field of $\ker\eta$.
Every complex conjugation belongs to $H$, because $\Psi$ is
even, $\kappa$ has odd order, and both $\eta,\eta'$ are odd.
Thus $F$ is totally real abelian and $K_0$ is imaginary
quadratic.  At a ramified place of either $\eta$ or $\eta'$,
the equality between them on $H$ kills both inertia values;
else both are already unramified.  Consequently $FK_0/F$
is unramified everywhere at finite places.  The prescribed
splitting of $D.S$ in $K_0$ gives the tower's \texttt{split}
clause and permits the required places over $\ell$.

The three auxiliary characters can be independently adjusted
by inertia at their three distinct ramified primes, while
leaving $\Psi,\xi,\omega,\kappa$ unchanged.  This proves
that the conditions in (UX.5) have their stated index; in
particular the index is even.  It also permits the following
three elements of $H\cap\ker\eta$:
\[
\begin{array}{c|ccc}
 &\bar\xi&\bar\Psi&\bar\omega\\ \hline
 a&-1&1&-1\\
 u&z&1&1\\
 t&z'&-1&1
\end{array}
\qquad z^4\ne1,\quad (z')^2\ne1.
 \tag{UX.6}
\]
Here $a$ starts from an odd power of complex conjugation;
$u$ starts from inertia at $2$; and $t$ starts from the
nontrivial prime-to-$\ell$ component of the Case-II character.
The latter component is nontrivial by AX1.  If the first
choice of $t$ has $(z')^2=1$, replace it by $tu$.
All powers used to kill $\kappa$ are odd and so preserve
the indicated dyadic nontrivialities.

We give the quantitative point in the construction of $u$
and its local analogue.  Write $\psi^{2^a b}=1$ with $b$
odd, as in AX1.  The prescribed $n_{\rm II}$ satisfies
$2^{n_{\rm II}}\ge\ell^2 2^{a+2}$.  If the fourth power
of $\xi(\operatorname{Frob}_\ell)^{2^a b\ell^j}$ were $1$,
the primitive dyadic estimate AX2 would give the strict
reverse inequality.  The chosen Frobenius has $c_\ell=1$,
hence $\omega=1$.  It therefore gives, at every place over
$\ell$, an element with
$\omega=\bar\Psi=\kappa=1$ and $\bar\xi^4\ne1$.
At $2$, a generator of the dyadic inertia quotient has the
full primitive character order; the same odd-power argument
gives $\bar\xi^4\ne1$ after the required factors are killed.
These are the source's \texttt{big\allowbreak L} and \texttt{big2}, not
an unspecified large-image assertion.

On $G_F$ the residual $\sigma$-character is $\varphi$ (its
inverse is the same character), and the $\alpha$-character
is $\bar\xi$.  Thus it is quadratic, unramified everywhere,
trivial at $\ell$, and is cut out by the original real
quadratic field fixed by $\ker\varphi$.  That field splits
at $2$ and $\ell$.  These facts prove \texttt{Two}.
The rows of (UX.6) verify the remaining character clauses
of \texttt{Ok}: a root character cannot equal the
cyclotomic residual character (use $a$ for even first
coordinate and $u$ for odd first coordinate); every root
is nontrivial (use $u$ for odd roots and $t$ for
$\pm\sigma$); the $\alpha$-character has order greater
than four; and no nontrivial quadratic character trivial
on $\ker\omega$ permutes the four $\bar\chi_i$.
For the last assertion apply such a permutation to the
character $1$: evaluation at $u$ excludes $\xi$ and
$\xi\Psi$, and evaluation at $t$ excludes $\Psi$.
The four eigenvalues at $t$ are $(-z',-1,z',1)$, all
distinct, so choose $g_0$ above $t$.

For completeness the local \texttt{loc\allowbreak L} alternative is
already supplied by the Frobenius just constructed.  Square
it to satisfy both quadratic equalities in (UX.5).  Its
$\omega$ and $\sigma$ values remain $1$, whereas its
$\alpha$ value is not $1$.  At $\ell$, $\xi$ and
$\varphi$ are unramified; inertia in $H$ therefore kills
the whole residual representation.  The two assertions
about $\psi$ follow from the first paragraph.  Finally
inertia at $\ell$ has prime-to-$\ell$ image in $G/H$,
since $\kappa$ is unramified there and all remaining
characters in (UX.5) have order prime to $\ell$.
This is literal \texttt{Tame\allowbreak L}.

Now let $W$ be as stated.  Its primes are different from
$\ell$: complete splitting of $\bar r_0$ and $\omega$ at
$\ell$ would contradict AX6's separation at the prescribed
local Frobenius.  For each $w\in W$, take the compositum
of the cyclic degree-$\ell$ extensions of $\mathbb Q$
unramified outside $w$.  Global reciprocity identifies it
with the degree-at-most-$\ell$ subfield of
$\mathbb Q(\mu_w)$: the $w$-power principal units have
pro-$w$ order and contribute no map to $\mathbb F_\ell$.
This is a finite, totally real, abelian extension,
unramified at $\ell$.  Compositing with $F$ kills exactly
all the characters quantified in \texttt{Dies}.  Its
relative exponent divides $\ell$; take $g'_0$ above
$g_0^\ell$.  Distinct residual characters stay distinct
under this power.  Nontrivial prime-to-$\ell$ characters
and the other clauses of \texttt{Ok,Two} persist under an
$\ell$-extension: their restriction maps have
$\ell$-primary kernels and prime-to-$\ell$ images.
The fixed quadratic fields are unchanged.

For the final assertion set
\[
 G_{F_m}=G_{F_1}\cap
 \ker\bigl(c_\ell^{\ell-1}\bmod\ell^{m+1}\bigr).
 \tag{UX.7}
\]
These are exactly the source \texttt{Is\allowbreak Layer} fields.
The $\ell$-primary part of $\psi|_{G_{F_1}}$ vanishes
on $\ker c_\ell$, by \texttt{Ok.\allowbreak kc}, and hence factors
through the cyclotomic image.  It kills that image's
prime-to-$\ell$ Teichmuller subgroup.  Consequently
it is trivial wherever $c_\ell^{\ell-1}=1$.
By compactness it
vanishes on the intersection in (UX.7) for all sufficiently
large $m$.  More explicitly, the decreasing closed sets
on which that part is nontrivial and
$c_\ell^{\ell-1}\equiv1\pmod{\ell^{m+1}}$ have empty
intersection; one is already empty.  Thus these layers
are deep.  Tameness gives an element of $G_{F_1}$ whose
cyclotomic image has infinite order: raise the original
$s_0$ to the prime-to-$\ell$ power supplied by
\texttt{Tame\allowbreak L}.  Its orders in the finite cyclotomic
quotients are unbounded, so $[F_m:\mathbb Q]\to\infty$.
Use $g_0^{\ell^m}$, enlarge $m$ if necessary, and choose
the original-field place representatives.  This proves
all the assertions with the original $K_0$ and label.
\end{proof}

\begin{proposition}[UX2: a good residual lattice for the original star point]
There is a permitted label, chosen from the original
ordinary flag as below.  Fix an initial tower $T_0$
supplied by UX1 for that label and its separating element
$g_0$.  There is a pair $P_0$ on the original
$G_{\mathbb Q}$ and original $V$, with a conjugacy
$r_K\simeq(P_0)_K$, such that $P_0$ is a small
\texttt{QPair} with the original inertia characters and its
residue has literal \texttt{Lat\allowbreak Shape}, normalized at this
$g_0$.  For every compatible fit extension of $T_0$
(whose separating element is the prescribed $\ell$-power
of $g_0$), the restricted residual pair, with the
transported flags, is good: it either has both
$-\sigma,\beta$ present, or is in $V_0$ with all
$n_v(2,0)\ne0$ and with $-\alpha$ present.
\end{proposition}

\begin{proof}
At the original $G_\ell$, AX6 separates the two residual
$\lambda$-blocks.  The first ordinary flag line is therefore
contained in the span of $e_0,e_2$ after reduction.  One of
$u_{00},u_{20}$ is a unit.  Use the trivial label in the first
case and the symplectic permutation $(2,3,0,1)$ in the
second.  In the second case the first-column eigenvector
equation, for every $g\in G_\ell$, gives
$\bar\chi_0(g)=\bar\chi_2(g)$; hence
$\bar\Psi|_{G_\ell}=1$.  This is exactly
\texttt{Lab\allowbreak Ok}, including its implication when the label
is nontrivial.  The same residual flag calculation makes
the first two leading minors units.  Symplecticity supplies
the remaining minor.

Now choose $T_0$ and $g_0$ by UX1 for this label.
Work initially in a complete DVR
defining the small pair.  Hensel's lemma splits the
characteristic polynomial of $r(g_0)$ into its four simple
residual roots.  The corresponding spectral idempotents
give an eigenbasis congruent to the original one; rescale
the paired eigenvectors to make its matrix symplectic.
This uses only units and, if necessary, a finite extension
of this defining DVR.  The ambient $K,V$ are not changed.
Conjugate the representation by this matrix and perform
the symplectic lower--upper decomposition of the ordinary
flag matrix.  Its lower factor $n$ is unique with diagonal
$1$.  Its residue has all off-diagonal entries zero except
possibly the root $-\sigma$ entries.  Multiplication by
the upper factor changes neither the ordinary characters
nor the inertia equalities.  The resulting representation
still has diagonal residual representation and diagonal
$r(g_0)$.  These constructions are finite algebraic
operations on small data, so smallness and continuity
over a complete finite-residue DVR persist.

There are nonzero generic entries in $\sigma$ and
$-\sigma$, and in at least one of
$\beta,-\alpha,-\tau$.  For an absolutely irreducible
representation this follows because its matrix span is
all of $M_4(K)$.  In the nonsplit block case, the full
$2\times2$ block span gives the first two assertions.
If all three remaining entries were zero, the symplectic
relations would make the other off-block entries zero as
well.  The conjugating matrix preserves the block
parabolic: it is assembled from spectral idempotents of
an element already in that parabolic and a block
permutation.  A block-preserving conjugacy splitting the
representation would give a unipotent off-block splitting
of the original representation (factor the conjugating
matrix into its diagonal-block and unipotent parts).
That contradicts the specified \texttt{Block\allowbreak Ext}.

Here is the lattice selection, including the finite
calculation behind it.  Let $\pi$ be a uniformizer of a
complete DVR defining the pair, and measure nonzero
entries by their integral $\pi$-orders.  For each root
$w$ include the orders of all representation entries of
that weight and, for the four negative roots, the orders
of the independent flag entries $n_{10},n_{21},n_{20},n_{30}$.
All are positive except possibly the order at $-\sigma$.
The set of integral pairs $(A,B)$ satisfying
\[
 w_1A+w_2B+2k\ge0
 \quad\hbox{for every recorded order $k$ at root $w$}
 \tag{UX.8}
\]
is nonempty.  Maximize $B$, then $A$.  The nonzero
$-\sigma$ entry bounds $B$ above; one of
$\beta,-\alpha,-\tau$ bounds $A$ above at fixed $B$.
Because these are integer bounds the maxima are attained.
Adjoin a square root $q$ of $\pi$ in the fixed $K$ and
conjugate by the symplectic diagonal torus with parameters
$(q^A,q^B)$.  Equation (UX.8) says exactly that all new
representation and flag entries belong to $V$.

Call a root tight if equality holds in (UX.8), so its
residual representation entry or independent flag entry
is nonzero.  The failed moves $(1,0),(0,1),(-1,1),(1,1)$
at a lexicographic maximum force respectively a tight
root in the following four sets:
\[
 \{-\alpha,\beta,-\tau\},\quad
 \{-\sigma,-\beta,-\tau\},\quad
 \{-\sigma,\alpha,-\beta\},\quad
 \{-\sigma,-\alpha,-\tau\}.
 \tag{UX.9}
\]
For the last two moves note first that the first two
conditions give $A+B$ even.  Thus a drop by $2$ in an
inequality of that parity must start at equality.
For every tight root other than $-\sigma$, positivity
of its original order gives $w_1A+w_2B<0$; tightness
of $-\sigma$ gives $B\ge0$.  Substitution of the eight
roots in (UX.1) into these sign inequalities and (UX.9)
gives exactly the following six possibilities:
\[
\begin{array}{c|c}
 \hbox{two roots that are tight}&\hbox{all tight roots lie in}\\ \hline
 \beta,-\sigma&N_P\\
 -\alpha,-\sigma&Q\\
 -\tau,-\sigma&\{-\tau,-\sigma\}\\
 -\alpha,-\beta&N_M\\
 -\tau,-\beta&\{-\tau,-\sigma,-\beta\}\\
 -\tau,\alpha&\{-\tau,-\sigma,-\beta,\alpha\}
\end{array}
 \tag{UX.10}
\]
For example $\beta$ tight means $B<A$, while
$-\beta$ tight means $A<B$; $\pm\tau$ have the signs
of $\pm(A+B)$ and $\pm\alpha$ those of $\pm A$.
These inequalities rule out any additional root in each
row and give (UX.10) without a compactness assertion about
an infinite-dimensional building.

The residual vanishing supplied by BR7 and AX4 now removes
the last five rows.  Precisely, for any odd root, a
locally constant normalized cocycle, unramified at the
specified inertia and locally a coboundary at $\ell$,
is zero, and its local coboundary parameter is zero.
At the other places in $D.T$, the local $H^1$ vanishes
by the separated residual Frobenius and the exact finite
local duality supplied by BX.  Thus AX4's original
\texttt{Unr\allowbreak Ev}, also after any residue-field extension,
applies.  Its coboundary constant is zero because the
root character is not $1$ at $g_0$; the local parameter
is zero by AX6's local separation.

To see the two particular consequences needed here, write
the first-column ordinary identity as
\[
 r_{i0}(g)+r_{i1}(g)n_{10}+r_{i2}(g)n_{20}
       +r_{i3}(g)n_{30}=n_{i0}\bar\chi_0(g).
 \tag{UX.11}
\]
If $-\alpha$ is not tight, $r_{10}=r_{32}=n_{10}=0$;
the row $i=3$ makes $r_{30}/\bar\chi_0$ a locally
trivial odd-root cocycle.  Its vanishing gives absence
of $-\tau$.  If $\beta$ and $\sigma$ are not tight,
then $r_{12}=r_{13}=0$; the row $i=1$ similarly removes
$r_{10},n_{10}$ and symplecticity removes $r_{32}$.
Thus
\[
 -\alpha\notin\mathrm{Tight}\Rightarrow
 -\tau\notin\mathrm{Tight},\qquad
 \beta,\sigma\notin\mathrm{Tight}\Rightarrow
 -\alpha\notin\mathrm{Tight}.
 \tag{UX.12}
\]
Inspection of (UX.10) leaves its first row only.
Here $\beta$ can occur only in the representation, and
the flag has support $Q$.  If $-\sigma$ is absent from
the representation, tightness forces $n_{20}\ne0$.
Equation (UX.11) at $i=2$ then gives
$\bar\Psi|_{G_\ell}=1$, and the $\beta$ entry cannot
be locally a coboundary, by the same odd-root vanishing.
These are exactly all clauses of \texttt{Lat\allowbreak Shape}.

Finally restrict to $G_F$.  Nonzero scalar cocycles in
the indicated roots stay nonzero.  Indeed, if a
normalized scalar cocycle vanished on $G_F$, evaluate
its cocycle identity at $g g_0g^{-1}$: since the two
diagonal characters at $g_0$ are distinct, its value
at every $g$ would be zero.  This proves the assertion
directly, even when $[F:\mathbb Q]$ is divisible by
$\ell$.  The transported lower flags satisfy
\[
 n'_{20}\bar\chi_0(g)=r_{20}(g)+\bar\chi_2(g)n_{20},
 \quad
 n'_{10}\bar\chi_0(g)=r_{10}(g)+\bar\chi_1(g)n_{10}
                                      +r_{12}(g)n_{20}.
 \tag{UX.13}
\]
They also have $n'_{21}=0$.  These identities follow by
equating the first two columns in $r(g)n=n'b$ with $b$
upper triangular.  If $r_{20}$ is present, it and the
$\beta$ entry remain present.  If it is absent,
(UX.13) keeps every $n'_{20}$ nonzero.  Were the
$-\alpha$ entry also absent, its second equation would
make the original $\beta$ entry locally a coboundary,
contrary to the preceding paragraph.  This gives the
two asserted good alternatives and completes the proof.
\end{proof}

\subsection{Finite cohomology and the three residual lines}

\begin{proposition}[UX3: simultaneous nonzero cup extensions with removable ramification]
Let $T$ be fit, let $k$ be any field receiving the original
residue field, and let finitely many locally constant
scalar cocycles $y_j$ have character $\chi_{-\sigma}$,
vanish at $g_0$, and be locally coboundaries at the chosen
places over $\ell$.  Allow ramification at a finite set
$W$ of source \texttt{Splits} primes.  After adding finitely
many further such primes, there are a common normalized
$\chi_{-\alpha}$-cocycle $x$, local constants $m_v$, and
cochains $z_j$ such that
\[
\begin{aligned}
 x(gh)&=x(g)+\chi_{-\alpha}(g)x(h),\quad x(g_0)=0,\\
 x(s)&=m_v-\chi_{-\alpha}(s)m_v\quad(s\in G_{F_v}),\\
 z_j(gh)&=z_j(g)+\chi_{-\tau}(g)z_j(h)
                 +y_j(g)\chi_{-\alpha}(g)x(h),\\
 z_j(s)&=o_{jv}-\chi_{-\tau}(s)o_{jv}
                 -y_j(s)\chi_{-\alpha}(s)m_v,
 \qquad z_j(g_0)=0.
\end{aligned}
 \tag{UX.14}
\]
There is a $g$ with $\chi_{-\alpha}(g)=1$ and $x(g)\ne0$.
All these cochains kill the enlarged frame's forbidden
inertia.  They produce actual residual symplectic pairs
with support $Q$, a common $-\alpha$ entry and common
$n_{v,10}$, while retaining the given $y_j$ and
$n_{v,20}$.  A tower extension supplied by UX1 removes
all the added ramification, including when $\ell=3$.
\end{proposition}

\begin{proof}
First we prove the finite-module assertion used in this
construction.  Let $M$ be a finite $\mathbb F_\ell[G_F]$-
module, unramified outside a finite $S_0$ containing the
places over $\ell$, and choose $0\ne m_0\in M$.
Chebotarev supplies infinitely many rational primes
splitting completely in the finite Galois extension
cut out by $M$, $F$, $\bar r_0$ and $\mu_\ell$.
Use their degree-one places in $F$ and remove $S_0$.
These primes satisfy exactly \texttt{Splits}; at them
the module and the mod-$\ell$ cyclotomic character are
trivial on the whole decomposition group.

At an added place impose the following local condition:
the restriction of a cocycle to inertia lies in
$\mathbb F_\ell m_0$.  Unramified cocycles give a copy
of $M$; an independent tame cocycle with image in
$\mathbb F_\ell m_0$ gives another factor of $\ell$.
Indeed the tame local relation is
$\operatorname{Fr}\tau\operatorname{Fr}^{-1}=\tau^{q_v}$,
and $q_v\equiv1\pmod\ell$.  The two homomorphisms from
the pro-$\ell$ quotient, one on Frobenius and one on
tame inertia, are independent.  The local condition
therefore has size $\ell|M|\ge2|M|$.  At $S_0$ impose
the zero cohomology class.

The exact finite Greenberg--Wiles identity supplied by
DX3, or its direct Poitou--Tate cardinal proof, now gives
\[
 2^r\le |\operatorname{Sel}_{S_0\cup Q_r}(M)|\,
       |M^*(1)|\,|M|^{[F:\mathbb Q]+|S_0|}.
 \tag{UX.15}
\]
Here $Q_r$ is any set of $r$ of the new places.  To
check the coarsening, bound each local $H^0$ by $|M|$,
each real-place contribution by $|M|$, and the dual
global invariants by $|M^*(1)|$; the factors $|M|^r$
from the new local conditions cancel those new local
$H^0$ bounds.  No constant on the right depends on $r$.
The group of actual cocycles unramified outside $S_0$
is finite by BX8's finite global cohomology and the
finite group of coboundaries; call its cardinality $N$.
Put $C=|M^*(1)|\,|M|^{[F:\mathbb Q]+|S_0|}$ and take
$r=NC$.  If every Selmer cocycle had zero inertia at
the added places, choosing one representative per
class would inject this Selmer group into those $N$
cocycles.  Equation (UX.15) would imply $2^{NC}\le NC$,
a contradiction.  Thus some cocycle $c$ has a nonzero
inertia value in $\mathbb F_\ell m_0$; the acting
element has trivial action on $M$.  This cocycle is
locally a coboundary at every old place.

We apply this assertion to an explicit finite module,
so no finite-dimensionality assumption on $k$ is hidden.
The image of the coefficient residue map is a finite
field $k_0$.  Each locally constant $y_j$ has finite
image because $G_F$ is compact.  Let $V_0'$ be their
finite-dimensional $k_0$-span in $k$ and put
$M=k_0\oplus\bigoplus_j V_0'$.  With
$\rho=\chi_{-\alpha}$ and $\sigma'=\chi_{-\sigma}$,
define
\[
 g\cdot(a,(b_j))=
  \bigl(\rho(g)a,
    (\rho(g)\sigma'(g)b_j+\rho(g)a\,y_j(g))_j\bigr),
 \qquad m_0=(1,0).
 \tag{UX.16}
\]
The cocycle identities for $y_j$ show directly that
this is an action.  It factors through a finite quotient
and kills the original forbidden inertia.  Project the
cocycle $c$ just constructed to its coordinates.  Its
first coordinate is $x$, the remaining coordinates are
$z_j$, and the cocycle identity is exactly (UX.14).
Writing $c(s)=g m-m$ at an old local group gives the
local identities in (UX.14), with the negatives of
the coordinates of $m$.  On the nonzero inertia value
in $\mathbb F_\ell m_0$, the first coordinate is
nonzero; the first part of (UX.16) then forces
$\rho(g)=1$.

Normalize at $g_0$ without losing this nonzero value.
Put $b=x(g_0)/(1-\rho(g_0))$ and subtract the
coboundary $b-\rho b$ from $x$.  For each $j$ put
$c_j=z_j(g_0)/(1-\rho(g_0)\sigma'(g_0))$ and replace
\[
 z_j(g)\longmapsto z_j(g)-
 \{c_j-\rho(g)\sigma'(g)c_j-y_j(g)\rho(g)b\},
 \qquad m_v\longmapsto m_v-b.
 \tag{UX.17}
\]
The denominators are nonzero by the separating $g_0$.
Substitution into (UX.14) proves the normalized
identities.  Since $\rho(g)=1$ at the witnessed
nonzero value, that value of $x$ is unchanged.

Here are the resulting matrices.  Write
\[
 U(x,y,t)=
 \begin{pmatrix}
 1&0&0&0\\x&1&0&0\\y&0&1&0\\t&y&-x&1
 \end{pmatrix},\qquad
 U(x,y,t)U(x',y',t')
 =U(x+x',y+y',t+t'+yx'-xy').
 \tag{UX.18}
\]
For the $j$th datum define
\[
 r_j(g)=U(x(g),y_j(g),2z_j(g)-x(g)y_j(g))
                         \operatorname{diag}(\bar\chi_i(g)),
 \quad
 n_{jv}=U(m_v,n_{jv}^{\rm old},2o_{jv}-m_vn_{jv}^{\rm old}),
 \tag{UX.19}
\]
where $y_j(s)=n_{jv}^{\rm old}
(1-\sigma'(s))$.  The diagonal conjugation scales
$x,y,t$ by $\rho,\sigma',\rho\sigma'$.
Equations (UX.14) and (UX.18) show both that $r_j$ is
a representation and that $n_{jv}^{-1}r_j(s)n_{jv}$
is the prescribed diagonal matrix.  They also give
the symplectic multiplier, normalization at $g_0$,
support $Q$, and every required inertia equality.
The common $x,m_v$ give literal \texttt{Same\allowbreak Al}; the
unchanged $y_j,n_{jv}^{\rm old}$ give literal
\texttt{Same\allowbreak Ps}.  The nonzero $x$ gives presence of
$-\alpha$.  This verifies the simultaneous
\texttt{gysin\allowbreak Al} conclusion entry by entry.

Finally apply UX1 to the added rational primes.  At
one of them $w$, the residual diagonal characters are
trivial.  A matrix supported in $Q$ is unipotent of
nilpotence class at most two and has exponent $\ell$
for every odd $\ell$: (UX.18) gives $U^n=U(nx,ny,nt)$.
The pro-$\ell$ quotient of tame inertia is cyclic;
its unique nontrivial degree-$\ell$ character is the
one killed by the degree-$\ell$ cyclotomic subfield
used in \texttt{Dies}.  Wild inertia is pro-$w$ and
acts trivially.  Hence the restricted representation
is unramified there.  This proves the claim also at
$\ell=3$, without using the false assertion that
every $4\times4$ unipotent matrix in characteristic
three has exponent three.  Existing residual pairs
already unramified at $w$ need no such assertion.
\end{proof}

\begin{proposition}[UX4: the complete three-line case-II residual bridge]
For every tame \texttt{Ok,Two} tower $T_0$ and every
original \texttt{Lat\allowbreak Shape} residual pair, there is a
tame extension $T_1\ge T_0$ such that, for every deep
fit extension $T\ge T_1$, literal \texttt{Seed\allowbreak IIAlt}
at the fixed $\iota$ supplies an integral seed $Q$
whose residue is connected to every compatible
restriction of that original residual pair by a
chain of three polynomial residual pairs.  Every
endpoint in the chain is good in the literal sense.
\end{proposition}

\begin{proof}
Restrict the original pair and transport its ordinary
flags as in UX2; call the result $p$.  It has support
$N_P$ and flag support $Q$.  Its normalized entry
$y(g)=p_{20}(g)/\bar\chi_0(g)$ is a
$\chi_{-\sigma}$-cocycle, with local constant
$n_v=p_{v,20}$.  Apply UX3 (one may use the same
datum twice) to produce a residual pair $q$ supported
in $Q$, with $\mathtt{SamePs}(p,q)$ and a nonzero
$-\alpha$ entry.  Let $T_1$ kill its finitely many
new ramified primes.  This choice precedes $T$, the
seed and its coefficient ring.

Fix any later deep fit $T$ and use SeedIIAlt there.
It produces an actual integral seed whose residue $z$
is in $V_0$, with all $z_{v,20}\ne0$, and satisfies
\[
 \operatorname{Ent}_{-\alpha}(z)\quad\hbox{or}\quad
 \bigl(\operatorname{Ent}_{\beta}(z)\ \hbox{and}\
       \operatorname{Ent}_{-\tau}(z)\bigr).
 \tag{UX.20}
\]
Map its finite residue field into the given ambient
residue field; the compatible coefficient maps specify
this embedding.  Restrict $p,q$ to $T$, choosing their
flags over the tower's actual place representatives.
The identities (UX.13) transport \texttt{Same\allowbreak Ps},
and the nonzero normalized $-\alpha$ cocycle survives
restriction: its character at the appropriate
$\ell$-power of $g_0$ is still nontrivial.  Auxiliary
inertia is killed by UX3.  Thus the new pair, still
called $q$, is a genuine residual pair for $W=\varnothing$.

Project $q$ to the Levi root $-\alpha$: keep its
$r_{10},r_{32}$ entries and remove the other negative
root entries, and do the same to its local flags.
In (UX.18) this is $(x,y,t)\mapsto(x,0,0)$.
It preserves the representation, ordinary and
symplectic identities.  Replace its local $20$ flag
coordinate by the corresponding nonzero coordinate
 of $z$.  Explicitly, replace the Levi flag
 $U(m_v,0,0)$ by
 $U(m_v,0,0)U(0,z_{v,20},0)
   =U(m_v,z_{v,20},-m_vz_{v,20})$.
 The Levi flag has already diagonalized the local
 representation.  The second factor commutes with that
 diagonal matrix because \texttt{Two.\allowbreak at\allowbreak L} gives
 $\bar\chi_0=\bar\chi_2$ and symplecticity gives
 $\bar\chi_1=\bar\chi_3$ on every $G_{F_v}$.
 The resulting pair
$q'$ is in $V_0$, has all $q'_{v,20}\ne0$, and has
\texttt{Same\allowbreak Al} with $q$ and \texttt{Same\allowbreak Ps} with $z$.
The $-\alpha$ representation entry is still nonzero.

The links are actual affine matrix interpolation:
\[
 p\ \xrightarrow{\ P\ }\ q\
   \xrightarrow{\ M\ }\ q'\
   \xrightarrow{\ P\ }\ z,
 \qquad A(X)=(1-X)A(0)+XA(1).
 \tag{UX.21}
\]
Here is why interpolation respects multiplication;
this cannot be assumed for arbitrary representations.
For a $P$ link, the two matrices have support $N_P$,
the same diagonal and the same $20$ entry.  The
symplectic identity also makes their $31$ entries
equal.  Their difference is supported in the
$\lambda$-off-block square-zero ideal.  Products
of any two such differences are zero, and this
ideal is preserved on both sides by the common
block triangular algebra.  For an $M$ link, both
matrices are lower triangular, have the same
diagonal and the same $10$ entry; symplecticity
also makes their $32$ entries equal.  Their
difference lies in the lower left $2\times2$
square-zero block ideal.  In either case
\[
 A(X)B(X)=(1-X)A(0)B(0)+XA(1)B(1).
 \tag{UX.22}
\]
Adjoint preserves the same ideals.  Apply (UX.22)
first to the representation identity and then to
$n(X)^{-1}r(X)n(X)$ and the symplectic identity.
It proves all pair equations over $k[X]$, the
prescribed residual diagonals and local characters,
and the stated support.  Each coefficient function
is locally constant because both endpoint functions
are.  Evaluation at $0,1$ is exactly the indicated
endpoint, so these are literal \texttt{Linked}
witnesses.

It remains to check goodness, since connectedness
alone is insufficient below.  The starting pair is
good by UX2.  If $q$ has a $-\sigma$ entry, it is
good by that entry and its nonzero $-\alpha$ entry;
otherwise its $20$ flags are nonzero whenever
needed, and it is in $V_0$ with $-\alpha$ present.
More explicitly, if the original $p$ has a
$-\sigma$ entry then \texttt{Same\allowbreak Ps} preserves it;
if it does not, UX2 gives nonzero $20$ flags and
\texttt{Same\allowbreak Ps} preserves those.  The pair $q'$
is good by $V_0$ and $-\alpha$.  The seed is good
by (UX.20), since \texttt{Two.\allowbreak big\allowbreak L} makes the
$\alpha$-character nontrivial at every local group.
These are exactly the four alternatives defining
\texttt{Good}.  Thus (UX.21) proves the complete
\texttt{Chain}, uniformly for all later deep towers
and all compatible original restrictions.
\end{proof}

\subsection{Uniform bounds before choosing the deep layer}

For a residual coefficient map $a:O\to k$, write $Z(w)$ for
the locally constant scalar cocycles
\[
 z(gh)=z(g)+a(\chi_w(g))z(h)
 \tag{UX.23}
\]
which vanish on the prescribed forbidden inertia and satisfy
$z(g_0)=0$ when $w\ne0$.  Write $Z_v(w)$ for the corresponding
local cocycle space, without this normalization, and
$\operatorname{Str}(w)$ for the subspace of $Z(w)$ whose
restrictions to every $G_{F_v}$ are coboundaries.  These are
exactly \texttt{Zgl}, \texttt{Zloc} and \texttt{Str}.

\begin{proposition}[UX5: one bound for every later cyclotomic layer]
For the tame tower $T_1$ of UX4 there is an integer $C\ge1$,
chosen before the later layer or any seed coefficient field,
such that every cyclotomic layer $T\ge T_1$ satisfies
\[
 \begin{gathered}
 s\le C,\qquad
 \dim_k\operatorname{Str}(w)\le C\quad(w\in\mathcal R),\\
 \dim_k Z(w)\le C\quad(w=0,\sigma,-\sigma),\\
 \dim_k Z(w)\le d+C,\qquad
 \dim_k\prod_{v\mid\ell}Z_v(w)\le d+C
       \quad(w\in\mathcal R_o).
 \end{gathered}
 \tag{UX.24}
\]
This holds for every field $k$ and every $a:O\to k$ killing
the original maximal ideal.  One can choose the layer deep
enough that simultaneously
\[
 d\ge30(C+1)^2,\qquad d_v=[F_v:\mathbb Q_\ell]=d/s
                 \ge30(C+1).
 \tag{UX.25}
\]
Thus the literal \texttt{Bounds} used later has one
constant, independent of all later coefficient extensions.
\end{proposition}

\begin{proof}
First use the finite residue field $k_0=O/\mathfrak m_O$.
Put
\[
 H_1=G_{F_1}\cap\ker\bar r_D\cap\ker\bar c_\ell
       \ \subset G_{\mathbb Q}.
 \tag{UX.26}
\]
It is open and contains the commutator, since $F_1/\mathbb Q$
and all four residual characters are abelian.  Its fixed
field is therefore abelian and contains $\mu_\ell$.
For a layer of index $m$, the analogous subgroup is
$H_1\cap\{c_\ell\equiv1\bmod\ell^{m+1}\}$.  Indeed, on
$1+\ell\mathbb Z_\ell$, raising to $\ell-1$ is an
automorphism: in
$(u^{\ell-1}-1)/(u-1)=1+\cdots+u^{\ell-2}$ the second
factor is a unit.  This identifies the actual fixed fields
with the cyclotomic root layers, not just with abstract
fields of the same degree.  TW1 supplies a single bound $r$
for their mod-$\ell$ class ranks and hence for their
everywhere unramified additive characters.

There is also a single bound $g$ on the number of places
over $\prod_{p\in D.S}p$ in all these layers.  At $\ell$,
the local cyclotomic image is open, so the number of
places is eventually constant.  At $p\ne\ell$, the image
of Frobenius in the cyclotomic $\mathbb Z_\ell$-quotient
is open: its value is the nonzero $\ell$-adic logarithm
of $p^{\ell-1}$.  It cannot be zero, for a positive
rational prime is not an $\ell$-adic root of unity.
Multiplying these finite indices by the degree of the
fixed initial abelian field bounds the number of
decomposition-group orbits.  Sum over the finitely
many primes in $D.S$.  In particular $s\le g$.

Let $H$ be the subgroup (UX.26) in $G_F$.  Its index is
prime to $\ell$: the residual diagonal characters have
values in $k_0^\times$, and $\bar c_\ell$ has values in
$\mathbb F_\ell^\times$.  Every scalar root character
is trivial on $H$.  Restrict a strict cocycle to $H$.
It is an additive character and is unramified
everywhere.  At places over $\ell$ this follows from
the local coboundary formula, since the character is
$1$ on $H$; away from $\ell$ it follows from the
specified killed inertia.  Its coordinates in an
$\mathbb F_\ell$-basis of $k_0$ are therefore among
the $r$-dimensional space just bounded.
Restriction, together with one element of $k_0$,
determines the original cocycle.  Explicitly, if
$z|_H=0$ and $\sum_{x\in G_F/H}z(x)=0$, then
permuting the cosets by $h$ in the cocycle identity
gives $[G_F:H]z(h)=0$, hence $z(h)=0$.
Consequently
\[
 \dim_{k_0}\operatorname{Str}(w)\le r+1.
 \tag{UX.27}
\]

Here is the corresponding unrestricted global estimate.
Use DX3's finite Greenberg--Wiles equality with the full
$H^1$ local condition at every place above $D.S$.
Its dual condition is zero.  The dual strict Selmer
classes inject, after restriction to $H$, into the
everywhere unramified additive characters just bounded;
the preceding coset argument allows the same extra
factor $|k_0|$.  Local Tate duality and Euler
characteristic give, for a one-dimensional coefficient
module $M$ and a finite place $v$,
\[
 \frac{|H^1(F_v,M)|}{|H^0(F_v,M)|}
 \le |M|^{\,1+[F_v:\mathbb Q_\ell]}
       \quad(v\mid\ell),\qquad
 \frac{|H^1(F_v,M)|}{|H^0(F_v,M)|}\le |M|
       \quad(v\nmid\ell).
 \tag{UX.28}
\]
Use $|Z^1|\,|H^0|=|M|\,|H^1|$ in the global equality.
After cancelling the local $H^0$ factors and using
$|H^0(M^*(1))|\ge1$, it yields
\[
 |Z^1(G_{F,S},M)|\,
       \prod_{v\mid\infty}|H^0(F_v,M)|
       \le |M|^{r+2+g+d}.
 \tag{UX.29}
\]
This also proves finiteness; all local and global
finite-coefficient interfaces in this calculation are
the actual ones supplied by BX8.  The cocycles in $Z(w)$
inject into these $G_{F,S}$ cocycles.  For
$w=0,\pm\sigma$, complex conjugation acts trivially:
$\Psi$ is even, and the $\sigma$ root is $\Psi$ or its
inverse after relabelling.  Since $F$ is totally real,
the product in (UX.29) is then $|M|^d$.  Thus
\[
 \dim Z(w)\le r+2+g\ (w=0,\pm\sigma),\qquad
 \dim Z(w)\le d+r+2+g\ (\hbox{all }w).
 \tag{UX.30}
\]
Finally (UX.28), together with
$|Z^1|=|M||H^1|/|H^0|$, bounds each local cocycle
dimension by $2+[F_v:\mathbb Q_\ell]$.  Summing gives
$d+2s\le d+2g$.

Take $C=r+2+2g$.  These estimates prove (UX.24) over
$k_0$.  To pass to an arbitrary $k$ use its specified
embedding of $k_0$.  Every $k_0$-linear functional
$\lambda:k\to k_0$ carries each of the defining
cocycle, normalization and local-coboundary equations
to the same equation over $k_0$; in the strict case
apply $\lambda$ also to each local coboundary constant.
This preserves the bounds on dimension.  For
completeness, if $N+1$ functions over $k$ were linearly
independent, some $N+1$ evaluation vectors would have
nonzero determinant.  Expanding the finitely many
entries in a $k_0$-basis of their span and taking
coordinate functionals would put their span inside
the scalar extension of the at most $N$-dimensional
space over $k_0$, a contradiction.  The same argument
applies to functions on the disjoint union of local
groups.  No finiteness of $k$ is required.

Choose the deep layer by UX1 only now.  Every place of
the abelian field $F$ above $\ell$ has the same local
degree, so $\sum_{v\mid\ell}[F_v:\mathbb Q_\ell]=d$
gives $d_v=d/s$.  The first inequality of (UX.25)
and $1\le s\le C$ imply its second.  This order of
choices leaves (UX.24) valid after the seed and every
later coefficient enlargement.
\end{proof}

\subsection{An actual common space and its local deformation rings}

\begin{proposition}[UX6: the space, its lines, and its flat completions]
Fix a fit tower $T$ and any finite coefficient DVR
$\Lambda=O'\supset O$.  There exists the literal
$X:\mathtt{Space}(T.\mathtt{fr}\,\varnothing)(O\to\Lambda)$.
Thus $B=X.B$ is complete Noetherian local over $\Lambda$
with residue $k_\Lambda$, $A=X.A$ is a finite type
$B$-algebra, and $B\hookrightarrow A$.  The space carries
the actual pair $P$, its characteristic polynomials
$c_g(X)\in B[X]$, and its actual local characters
$\psi_{v,i}\in B$, with every clause of \texttt{Space}.
Every residual pair over $k_\Lambda[X]$ has an actual
map $A\to k_\Lambda[X]$ giving it and killing
$\mathfrak m_B$.  If $x:A\to k_\Lambda$ is a residual
pair, the completion
\[
 R_x=\widehat{A_{\ker x}},\qquad f_x:A\longrightarrow R_x
 \tag{UX.31}
\]
is flat over $A$ and is the universal complete local
ring of lifts of that exact residual pair, including
all its actual flags.
\end{proposition}

\begin{proof}
We construct all objects, since a universal pseudorepresentation
would not supply the local flags or the polynomial lines.
The residual root characters are locally constant and
have pairwise distinct values at $g_0$.  They satisfy
$\bar\chi_i\bar\chi_{i^{\rm rev}}=\bar\nu$,
$2\ne0$, and are trivial on the prescribed killed
inertia and on each $\mu_v$.  The image of $\mu_v$
in the local abelianization is $\ell$-primary torsion
and lies in abelian inertia; its image in the prime-to-$\ell$ group
$k_\Lambda^\times$ is trivial.
The quotient of $G_F$ by the closed normal subgroup
generated by the forbidden inertia is a quotient of
$G_{F,S}$ for the fixed finite set of places over
$D.S$.  It satisfies $\Phi_\ell$: additive characters
on any open subgroup inject into the finite $H^1$
over the corresponding number field, by BX8.
Each closed decomposition group satisfies the local
version by the same chapter.  These checks give
exactly \texttt{Std}, on the actual groups.

First disregard the flags and regard the $4s$ local
characters as independent.  FP1--FP2's finite-entry
construction applies to the split residual representation
$\operatorname{diag}(\bar\chi_i)$ on this quotient.
Concretely, finite $\Phi_\ell$ distinguishes its
dual-number lifts by finitely many group entries.
The finite-level entry algebras, with multiplication
and continuity imposed, have an inverse limit which
is a quotient of a power-series ring on those
tangent generators.  Impose the further closed
entry equations
\[
 \rho(g)^*\,\rho(g)=\nu(g)1,\quad
 \rho(n)=1\ (n\in N),\quad
 \rho(g_0)_{ij}=0\ (i\ne j).
 \tag{UX.32}
\]
Their ideal is closed because the power-series ring
is Noetherian and complete.  The split residual
point satisfies them, so this is a proper quotient.
Construct each independent local character by the
same argument in rank one and take the completed
coproduct of this finite list of rings.  The result
is a complete Noetherian local $\Lambda$-algebra
$R^{\rm sp}$ representing the actual matrices and
local characters with split residual matrices.
This is a direct application of the entry construction,
not an assertion that the extra hypotheses of
the full FP8 theorem hold here.

The two-dimensional diagonal symplectic torus acts
by conjugation on these matrices, fixing their
diagonal and all local characters.  It makes entry
$(i,j)$ homogeneous of weight $w_{ij}$.
Choose a finite collection of entry differences
and local-character differences spanning the
cotangent space.  They are homogeneous, vanish
at the split point, and give a surjection
$\Lambda[[z_1,\ldots,z_h]]\twoheadrightarrow R^{\rm sp}$
with homogeneous kernel; the formal torus argument
and successive lifting of homogeneous generators
are those proved in FP5.  Put
\[
 B_0=(R^{\rm sp})_0,\qquad
 A_0=\bigoplus_{w\in\mathbb Z^2}(R^{\rm sp})_w.
 \tag{UX.33}
\]
Here a homogeneous piece permits an infinite convergent
sum of monomials of its fixed weight, whereas the
direct sum in $A_0$ is a finite sum of such pieces.
The finite-monomial argument of PX1 applies here.
Every antichain in $\mathbb N^h$ is finite: from
an infinite sequence successively extract subsequences
nondecreasing in each coordinate.  Thus there are
finitely many componentwise minimal nonzero
zero-weight exponents.  They generate the whole
zero-weight monoid, by subtracting one such exponent
and inducting on total degree.  It follows that
$B_0$ is a quotient of a finite power-series ring
over $\Lambda$.  For any fixed weight, the same
antichain argument gives finitely many componentwise
minimal exponent vectors;
every other vector differs from one of them by
a zero-weight vector.  Therefore each piece is
generated over $B_0$ by finitely many monomials,
and $A_0=B_0[z_1,\ldots,z_h]$.
The degree-zero ring is complete local, with the
same residue field, and $A_0$ is finite type over it.
The same assertions descend through the homogeneous
kernel by taking its homogeneous components.

The characteristic polynomials and local characters
have degree zero.  We need their precise uniqueness
property for maps out of $B_0$.  Suppose
$u,u':B_0\to R'$ agree on these values, for a
complete local $\Lambda$-algebra $R'$ with the
specified residue field.  The diagonal eigenvalues
$d_i=\rho(g_0)_{ii}$ are determined individually by
the common polynomial $c_{g_0}$ and their distinct
residues.  Indeed
$\prod_j(u'(d_i)-u(d_j))=0$, and every factor
with $j\ne i$ is a unit; hence $u'(d_i)=u(d_i)$.
Their differences and their inverses are also
fixed by $u,u'$.

This determines every weight-zero product of
matrix entries.  Here is an explicit symplectic
invariant argument.  Let $\mathcal R$ be the algebra
generated by the matrices $\rho(g)$ over the
equalizer of $u,u'$.  It is closed under symplectic
adjoint, because $\rho(g)^*=\nu(g)\rho(g^{-1})$.
All its traces belong to that equalizer, and
\[
 e_{ii}=\prod_{j\ne i}
           \frac{\rho(g_0)-d_j1}{d_i-d_j}
                  \ \in\mathcal R .
 \tag{UX.34}
\]
In a zero-weight product, take a factor of nonzero
weight $w$.  Some other weight $w'$ has
$w\cdot w'<0$, since all weights sum to zero.
Inspection of the displayed matrix (UX.2) gives
\[
 w_{ij}\cdot w_{kl}<0\quad\Longrightarrow\quad
 j=k\ \hbox{or}\ l=i\ \hbox{or}\
 j=l^{\rm rev}\ \hbox{or}\ k^{\rm rev}=i .
 \tag{UX.35}
\]
In the first two cases multiply through the
intermediate projector $e_{jj}$ or $e_{ii}$;
in the other two first apply adjoint to one
factor.  Up to its explicit sign, the product
of the two selected entries becomes one entry
of a matrix in $\mathcal R$ of weight $w+w'$.
Induction on the number of factors reduces to
diagonal entries, which are
$\operatorname{tr}(e_{ii}M)$ and belong to the
equalizer.  This proves the assertion for every
zero-weight monomial.  The local-character
generators already belong to the equalizer.
Continuity and the finite power-series presentation
of $B_0$ now give $u=u'$.  This is the full
\texttt{gen\allowbreak B} property; a generic-point trace
identity alone would not suffice.

We next explain why $A_0$ describes non-split
triangular residual matrices too.  If the support
is $N_P$, use the functional
$\tau(a,b)=-2a-b$; if it is $N_M$, use
$\tau(a,b)=-a-2b$.  It is strictly positive on
the permitted nonzero residual entries.  Over an
Artinian local target $R'$, write
$\mathfrak m^N=0$, choose $c$ bounding the negative
values $\tau(w_{ij})$, and form the subring
\[
 S=R'[[t]][\,\mathfrak m\,t^{-c}\,]
                         \subset R'((t)).
 \tag{UX.36}
\]
It is a finite $R'[[t]]$-module because every
product of $N$ negative generators vanishes.
It is complete Noetherian local with residue
$k_\Lambda$ and maximal ideal generated by $t$
and the nilpotent negative generators.  Rescale
entry $(i,j)$ of the given matrix by
$t^{\tau(w_{ij})}$ and leave local characters
unchanged.  A negative exponent has coefficient
in $\mathfrak m$, so all entries lie in $S$;
every off-diagonal entry lies in its maximal
ideal.  The identities
$w_{ij}+w_{jk}=w_{ik}$ and
$w_{j^{\rm rev}i^{\rm rev}}=w_{ij}$ show that
matrix multiplication and symplectic adjoint
respect this rescaling.  It therefore gives a
split residual point of $R^{\rm sp}$ over $S$.

Its value on a homogeneous element of weight $w$
is a single Laurent monomial $t^{\tau(w)}b$.
For infinite homogeneous series, this follows
by truncation: the coefficient of a fixed
Laurent power vanishes on a sufficiently high
power of the maximal ideal of $S$.
Taking the coefficient $b$ on each homogeneous
piece is multiplicative, and hence defines
$A_0\to R'$ giving the original matrices and
characters.  For a general complete local
target, do this modulo each power of its
maximal ideal.  The maps agree on the chosen
entry generators and therefore are compatible;
their limit gives the required map.  Thus
no substitution $t=1$ in an arbitrary formal
power series has been used.

Adjoin to $A_0$ the four lower unipotent
symplectic coordinates at each place and impose
the equations that the flags triangularize
the representation with exactly the chosen
local characters, together with their $\mu_v$
conditions.  These are polynomial entry
equations; the resulting algebra is finite
type over $B_0$.  Replace $B_0$ by its image
$B$ in this algebra $A$ to ensure injectivity.
The split residual pair with identity flags
gives a point, so the image quotient is
proper and is still complete local.  The
matrix and flag entries generate $A$ over
$B$, and the preceding construction gives
the required universal map for every lift
of every permitted residual pair.  Its
entries modulo $\mathfrak m_B^n$ are locally
constant: homogeneous pieces are finite
modules over the complete ring $B$, and
this is inherited from the original
continuous formal matrices.  All maps,
residual polynomials and character functions
are consequently the ones required in
\texttt{Space}.

For the line property, embed $k_\Lambda[X]$
in $k_\Lambda[[X]]$ and apply this universality
to the given polynomial residual pair.  Its
characteristic polynomials are the fixed
split polynomials and its local characters
are the fixed $\bar\chi_i$, because its
permitted support is triangular.  The just
proved uniqueness for maps from $B$ implies
that $B$ maps to constant residues.  The
generators of $A$ map to the actual
polynomial entries and flags, so the map
lands in $k_\Lambda[X]$ and kills
$\mathfrak m_B$.  Its values reproduce
the given pair exactly.

Finally $A$ is Noetherian.  Localization
followed by completion at $\ker x$ is flat,
has residue $k_\Lambda$, and is complete
Noetherian local over $\Lambda$.  An actual
lift $Q$ of $x(P)$ first gives $g:A\to R'$
by the established universality.  Equality
of the residual pairs forces
$\bar g=x$: matrix and flag entries generate
$A$, and their characteristic polynomials
and local characters determine its base.
Thus $g$ is local at $\ker x$ and extends
uniquely to $R_x$.  Conversely such a map
gives $Q$.  Uniqueness again follows from
the same entry generation and base
uniqueness.  This proves the exact
\texttt{Is\allowbreak Def}, flatness and residue clauses
of \texttt{cmpl\allowbreak Ex}.
\end{proof}

\begin{proposition}[UX7: the exact presentation deficit]
Suppose in addition that $T$ is deep.  For every residual
pair $x$ over $k_\Lambda$, the same completed ring $R_x$
has a presentation
\[
 R_x=\Lambda[[z_1,\ldots,z_a]]/(f_1,\ldots,f_b),
 \qquad b+2+2d\le a+2s .
 \tag{UX.37}
\]
The assertion is uniform in the residual pair, in the
space $X$, and in every permitted coefficient extension.
\end{proposition}

\begin{proof}
Write $k=k_\Lambda$, $M=\mathfrak{sp}_4(k)$, and
$\rho=\operatorname{Ad}(\bar r_x)$.  For a local
flag $\bar n_v$, let $\mathfrak b_v$ be the conjugate
of the upper Borel Lie algebra by $\bar n_v$, and
let $\mathfrak n_v$ be its complementary lower
nilpotent subspace.  These have dimensions $6$ and
$4$, while $\dim M=10$ and the diagonal torus
has dimension $2$.  For example an upper
symplectic infinitesimal matrix is
\[
 \begin{pmatrix}
 a&c&d&e\\0&b&f&d\\0&0&-b&-c\\0&0&0&-a
 \end{pmatrix},
 \tag{UX.38}
\]
which proves the asserted Borel dimension directly.

Let $\Gamma$ be the quotient of $G_F$ by the
closed normal subgroup generated by
$I_w\cap\ker\bar r_D$ at every $w\nmid\ell$.
This is exactly the quotient specified by the
pair's forbidden inertia.  The representation
$\rho$ descends to it.  To record the cohomology
being used, put
\[
 C^\bullet_{\rm ord}
 =\operatorname{Cone}\left(
 C^\bullet(\Gamma,M)\longrightarrow
       \bigoplus_{v\mid\ell}
           C^\bullet(G_{F_v},M/\mathfrak b_v)
                  \right)[-1].
 \tag{UX.39}
\]
We use continuous finite cochains.  The groups in
degrees $0,1,2$ are the source groups
$\mathtt{coneH0}$, $\mathtt{H1C}$ and $\mathtt{coneH2}$,
respectively.  A degree-one class is represented by $(z,(Y_v))$
with $z|_{G_{F_v}}-dY_v\in\mathfrak b_v$,
modulo global coboundaries and changes
$Y_v\mapsto Y_v+B_v$ with $B_v\in\mathfrak b_v$.
In degree two use $(c,(c_v))$ with
$c|_{G_{F_v}}-dc_v\in\mathfrak b_v$.
Denote their dimensions by $h^0_o,h^1_o,h^2_o$.

We first justify the arithmetic Euler inequality
on this exact quotient:
\[
             h^2_o+h^0_o+6d\le h^1_o+4d .
 \tag{UX.40}
\]
The prime-to-$\ell$ index of
$K=\ker\bar r_D|_{G_F}$ is essential here;
the quotient is not silently replaced by
the group unramified outside $\ell$.
Choose a finite set $S'$ including the places
over $\ell$ and all ramification of $K$.
Then $\Gamma$ is a quotient of $G_{F,S'}$,
by the normal closure of the indicated
inertia intersections at $S'\setminus\{\ell\}$.
At such a place, the remaining inertia has
order prime to $\ell$.  Its higher cohomology
vanishes by averaging, and the residual
unramified quotient is procyclic.  In
particular, an inflated $H^2(\Gamma,M)$ class
has zero restriction at this place.
The local tame relation and the two-term
procyclic complex give
\[
 \dim_k\operatorname{Hom}_{G_{F_w}}
       (I_w\cap K,M)
       =\dim_k H^2(F_w,M)
            \quad(w\in S',\ w\nmid\ell);
 \tag{UX.41}
\]
the homomorphisms here are locally constant,
additive, and equivariant for conjugation.
One can see (UX.41) by removing the
pro-prime-to-$\ell$ wild inertia, then using
$\operatorname{Fr}\tau\operatorname{Fr}^{-1}
=\tau^{\#k_w}$ on the pro-$\ell$ tame
generator.  The equivariant homomorphisms
are the kernel of
$\operatorname{Fr}-\#k_w$ on $M^{I_w/K}$;
local Tate duality identifies its dimension
with $H^2(F_w,M)$.  Passing through the
finite prime-to-$\ell$ inertia index changes
neither side.

The inflation five-term sequence bounds the
kernel in degree two by the sum of the spaces
in (UX.41); restriction to the inertia
generators is injective because they normally
generate the kernel.  Combining this with
global Euler characteristic and the
Poitou--Tate tail for $G_{F,S'}$ cancels
these extra-place terms.  More explicitly,
write $h^2_{\rm str}$ for the kernel of
$H^2(\Gamma,M)\to\bigoplus_{v\mid\ell}H^2(F_v,M)$.
The resulting inequality is
\[
 \begin{split}
 h^2_{\rm str}+\sum_{v\mid\ell}h^2(F_v,M)
       +h^0(\Gamma,M)+d\dim M
 \ \le\ h^1(\Gamma,M)
       +\sum_{w\mid\infty}h^0(F_w,M)
       +h^0(F,M^*(1)).
 \end{split}
 \tag{UX.42}
\]
For clarity about the cancellation, the five-term
sequence gives
$\dim\ker(\mathrm{inf}:H^2(\Gamma,M)\to H^2(G_{F,S'},M))
=\dim H^1(\ker,\!M)^{\Gamma}
 -h^1(G_{F,S'},M)+h^1(\Gamma,M)$.
Its first term is at most the sum in (UX.41).
The strict kernel maps into the strict kernel
for all places of $S'$, by the local vanishing
above.  Apply the Poitou--Tate tail to that
kernel and insert the global Euler equality;
the same sum of $h^2(F_w,M)$ is subtracted
on the other side.  This proves (UX.42)
and finiteness, with no omitted local term.

The last term of (UX.42) vanishes.  The
triangular support of the residual pair
gives a filtration of the matrix adjoint
module with characters $\bar\chi_i/\bar\chi_j$.
No one of these is $\bar c_\ell$, by
$T.\mathtt{Ok.g2}$.  A homomorphism to
$k(1)$ must therefore vanish, successively
on the filtration.  This also applies
to $M$, since
$X\mapsto(X-X^*)/2$ is an equivariant
projection from all matrices onto $M$.
At a real place, the involution $\bar r_x(c)$
has multiplier $-1$; its $+1$ and $-1$
eigenspaces are paired perfectly by the
symplectic form and both have dimension $2$.
The fixed part of $M$ consists of an
endomorphism on the first eigenspace and
the forced negative transpose on the second,
and has dimension $4$.

Apply now the exact sequence of (UX.39).
Local Tate duality makes
$H^2(F_v,M)\to H^2(F_v,M/\mathfrak b_v)$
surjective, by dualizing the injective
map on $H^0$ of the corresponding duals.
Elementary kernel counting gives
\[
 \begin{split}
 &\dim\ker\left(H^2(\Gamma,M)\to
             \bigoplus_vH^2(F_v,M/\mathfrak b_v)\right)
            +\sum_vh^2(F_v,M/\mathfrak b_v)\\
 &\hspace{25mm}\le
 h^2_{\rm str}+\sum_vh^2(F_v,M).
 \end{split}
 \tag{UX.43}
\]
Local Euler characteristic gives
$h^1(F_v,Q_v)-h^0(F_v,Q_v)
 =h^2(F_v,Q_v)+d_v\dim Q_v$
for $Q_v=M/\mathfrak b_v$.  Subtract consecutive
ranks in the cone exact sequence and use
(UX.42)--(UX.43), $\dim Q_v=4$,
$\sum d_v=d$, and the vanishing dual
invariants.  The result is (UX.40).

Temporarily remove the $\mu_v$ conditions
from the deformation problem and call its
universal completed ring $R^0$.  Its
tangent dimension $a$ satisfies
\[
                     h^1_o+2\le a+h^0_o.
 \tag{UX.44}
\]
Indeed every cone class has a representative
whose local $Y_v$ lies in the complement
$\mathfrak n_v$.  Conjugation by $1+\epsilon X$
can make the off-diagonal value of $z(g_0)$
zero, since the residual eigenvalue
differences are units.  This is the actual
normalization imposed in (UX.32).
The remaining global conjugations are the
two diagonal directions; those acting
trivially are exactly the invariants lying
in every $\mathfrak b_v$, of dimension $h^0_o$.
This proves (UX.44), and in these separated
coordinates it in fact gives equality.

The space $H^2_o$ is a complete obstruction
space for $R^0$.  Here is its concrete
realization.  Across a small extension
with square-zero kernel $I$, lift the
global matrices and the local flags
preserving their symplectic form.  Such
lifts exist: the correction
$X\mapsto(X-X^*)/2$ solves the infinitesimal
form equation, and the lower unipotent
chart lifts its four coordinates freely.
The multiplication defects of the global
matrices form a continuous two-cocycle
in $M\otimes I$.  The failure of each
lifted flag to triangularize the local
matrices gives the local one-cochain
modulo $\mathfrak b_v\otimes I$.
Associativity and restriction give exactly
the cone equations above.  Changing
lifts changes this pair by a cone
coboundary.  If its class is zero,
correct the matrices by the corresponding
global one-cochain and correct the
flags by its local zero-cochains.
The remaining $g_0$ normalization is
restored by the separated-eigenvalue
conjugation, which preserves all other
conditions.  Thus vanishing is sufficient
as well as necessary for a lift.

The finite obstruction argument of FP7
now presents $R^0$ with $a$ variables
and at most $h^2_o$ relations.  Explicitly,
in a minimal surjection $\Lambda[[z_a]]
\twoheadrightarrow R^0$, dualizing the
relation module $J/\mathfrak mJ$ sends a
functional to the obstruction for the
associated one-dimensional small
extension.  If that obstruction vanishes,
the corrected universal lift gives a
section, so the functional on $J$ is zero.
This injects the dual relation module
into $H^2_o$ and proves the bound.

Finally impose $\mu_v$.  By local class
field theory, its image in the abelianization
of $G_{F_v}$ is exactly the $\ell$-power
torsion in the profinite completion of
$F_v^\times$, which is cyclic: it is
the group of $\ell$-power roots of unity.
Choose a lift $\zeta_v$ of a generator.
All characters kill the closed commutator,
so it suffices to impose their values
at $\zeta_v$.  Deepness gives
$\nu(\zeta_v)=1$: the cyclotomic character
has no $\ell$-power torsion for odd $\ell$,
the dyadic character has prime-to-$\ell$
order, and the remaining $\psi|_{G_F}$
does also.  The symplectic relation
$\psi_{v,i}\psi_{v,i^{\rm rev}}=\nu$
therefore reduces the four equations to
\[
             \psi_{v,0}(\zeta_v)=1,\qquad
             \psi_{v,1}(\zeta_v)=1 .
 \tag{UX.45}
\]
Thus $R_x$ is a quotient of $R^0$ by
at most $2s$ further equations.  If
$b\le h^2_o+2s$ is the resulting count,
(UX.40) and (UX.44) give
$b+2+2d\le a+2s$, as required.
\end{proof}

\subsection{The two loci that a connecting curve must avoid}

For reference, the nine source reducible patterns are the following.
The notation in each row lists all its permitted nonzero roots.
\[
\begin{array}{c|l}
0&\beta,-\beta,-\alpha,-\tau,-\sigma\\
1&\beta,-\beta,\alpha,\tau,\sigma\\
2&\tau,-\tau,\sigma,\beta,-\alpha\\
3&\tau,-\tau,-\sigma,-\beta,\alpha\\
4&\alpha,-\alpha,-\tau,-\sigma,-\beta\\
5&\alpha,-\alpha,\tau,\sigma,\beta\\
6&\sigma,-\sigma,\beta,-\alpha,-\tau\\
7&\sigma,-\sigma,-\beta,\alpha,\tau\\
8&\beta,-\beta,\tau,-\tau
\end{array}
\tag{UX.46}
\]
Write $I_S$ for the ideal generated by the representation
entries outside row $S$.  Over a field with separated
$g_0$, every non-absolutely-irreducible representation
lies in one of these loci.  To check the assertion,
the spectral projectors (UX.34) lie in the algebra
spanned by the representation.  Multiplying between
them shows that this algebra is spanned by its
nonzero matrix units.  If their directed graph is
strongly connected, multiplying along paths gives
all matrix units and the representation is absolutely
irreducible.  Otherwise a union of terminal components
gives a proper coordinate invariant subspace.
The algebra is closed under symplectic adjoint,
so an edge $(i,j)$ is present if and only if
$(j^{\rm rev},i^{\rm rev})$ is present.  List the
nonempty proper coordinate subsets and their
symplectic orthogonals.  The four one-dimensional
subsets give rows $0$--$3$, the four isotropic
two-dimensional subsets give rows $4$--$7$,
and the two complementary nonisotropic subsets
give the common block pattern in row $8$.
Three-dimensional subsets give the orthogonals
of the one-dimensional ones.  This exhausts
the list and proves the assertion without a
classification assumption.

\begin{proposition}[UX8: the complete reducible-locus bound]
At every good residual pair $x$ and for every row $S$
of (UX.46), the exact ring of UX6 satisfies
\[
                   \dim(R_x/I_S)\le d+20C+1 .
 \tag{UX.47}
\]
Here $C$ is the one constant of UX5; no bound on the
cardinality of the residue field is added.
\end{proposition}

\begin{proof}
There is nothing to prove if $I_S=R_x$.  Otherwise
the residual pair $c=(\bar r,\bar n)$ itself has
support in $S$.  A derivation at its residual
point gives the infinitesimal matrices
\[
 X(g)=\partial r(g)\bar r(g)^{-1},\qquad
 Y_v=\partial n_v\,\bar n_v^{-1}.
 \tag{UX.48}
\]
They lie in $\mathfrak{sp}_4(k)$, satisfy the
adjoint cocycle identity and vanish on the
forbidden inertia, and $X(g_0)$ is diagonal.
They also satisfy
\[
 \begin{gathered}
 (X(g)\bar r(g))_{ij}=0\quad
                       (i\ne j,\ w_{ij}\notin S),\\
 \bar n_v^{-1}Y_v\bar n_v\text{ is strictly lower},\\
 \bar n_v^{-1}\!
 \left(X(h)+\bar r(h)Y_v\bar r(h)^{-1}-Y_v\right)
 \bar n_v\text{ is upper}\quad(h\in G_{F_v}).
 \end{gathered}
 \tag{UX.49}
\]
Conversely these are exactly the linearized
equations needed for an upper bound.  The
map from derivations to $(X,Y)$ is injective:
if all entries in (UX.48) vanish, the two
dual-number maps give identical pairs;
UX6's universal uniqueness makes the maps
equal.  Thus it is enough to bound the
space defined by (UX.49).

We reduce this linear space to two split
calculations.  Conjugate matrices by
$\operatorname{diag}(t^{m_i})$ and take
the limit at $t=0$, using respectively
\[
 (m_i)=(0,2,1,3)\text{ for }N_P,\qquad
 (0,1,2,3)\text{ for }N_M.
 \tag{UX.50}
\]
The residual representation becomes split
and the flags become $1$.  If goodness
instead uses $V_0$ and its nonzero $20$
flag coordinates, use $(m_i)=(0,1,0,1)$.
The representation again becomes split,
but the flags become
$1+b_vX_{-\sigma}$, with $b_v=c.n_v{}_{20}\ne0$.
All these weights satisfy
$m_i+m_{i^{\rm rev}}$ constant, so they
preserve the symplectic equations.
The dimension of the original tangent
space is at most that of the limiting
space.  This assertion does not require
a finite-dimensional ambient space of
all functions: filter the function
coordinates by their finitely many
matrix weights.  The lowest nonzero
weight terms of the linear equations
(UX.49) give the limiting equations.
Each successive graded piece injects
into the corresponding weight piece
of their solution space.  Summing the
finitely many inequalities gives the
assertion.

For split $\bar r$, every scalar
$X_{ij}$ belongs to $Z(w_{ij})$.
When $\bar n_v=1$, every lower entry
is locally a coboundary, directly by
the last equation in (UX.49).
If the first or second goodness
alternative holds, only rows $0,4,6$
can contain its required entries.
In these three rows select, respectively,
the entries $12,01,12$.  Upon setting
the selected scalar cocycle equal to
zero, every other independent odd
entry is either lower or is forced
to be zero by the pattern.  Thus all
odd entries lie in their strict
spaces $\operatorname{Str}(w)$.

In the $V_0$ case, the local diagonal
characters satisfy $\bar\chi_0=\bar\chi_2$
and $\bar\chi_1=\bar\chi_3$.  The first
equality follows already from the $20$
entry of $\bar r(h)\bar n_v=
\bar n_v(\bar n_v^{-1}\bar r(h)\bar n_v)$,
using $b_v\ne0$; the second follows by
the multiplier.  In particular, on
this local group a root's character
depends only on its first coordinate.
Set
$Z=X(h)+\bar r(h)Y_v\bar r(h)^{-1}-Y_v$.
Multiplying out the limiting flag
$1+b_vX_{-\sigma}$ in (UX.49) gives
the three relations
\[
 Z_{10}=-b_vZ_{12},\qquad
 Z_{30}=-b_v^2Z_{12},\qquad
 Z_{21}=2b_vZ_{01}+b_v^2Z_{03}.
 \tag{UX.51}
\]
The possible rows are $0,2,4,5,6,8$;
choose the single exceptional entry as follows:
\[
 \begin{array}{c|rrrrrr}
 S&0&2&4&5&6&8\\ \hline
 (i,j)&12&12&01&01&12&03 .
 \end{array}
 \tag{UX.52}
\]
If the selected cocycle is zero, (UX.51)
and the entries excluded by that row
make all six independent odd entries
locally coboundaries.  For example the
first relation transfers this property
between $X_{12}$ and $X_{10}$, and
the second transfers it between
$X_{12}$ and $X_{30}$.  In the third,
any two known entries determine the
third.  Division by $b_v,b_v^2,2b_v$
is permitted.  These statements concern
coboundaries, not zero values of $Z$:
the same identities absorb the local
$Y_v$ entries into the required
coboundary constants.  This verifies
every row of (UX.52).

Project the whole tangent space to the
selected odd scalar cocycle.  Its image
has dimension at most $d+C$, by UX5.
In the kernel use the ten entries
$i\le j^{\rm rev}$ of a symplectic
infinitesimal matrix.  Each odd one
is strict, and each even one has
weight $0,\sigma$ or $-\sigma$;
each therefore contributes at most
$C$.  The local variables contribute
at most their six strictly lower
coordinates at each place, hence
at most $6s\le6C$.
These coordinates determine $(X,Y)$:
symplectic adjoint determines the
remaining matrix entries, and
conjugating back by $\bar n_v$
determines each $Y_v$.
Consequently
\[
 \dim\operatorname{Tan}_S(c)
        \le(d+C)+10C+6C
        \le d+20C.
 \tag{UX.53}
\]
The first or second goodness alternatives
use (UX.50); the third or fourth use
(UX.52).  Thus all good residual pairs
are covered.  A complete local algebra
over a DVR with relative cotangent
dimension at most $d+20C$ is a quotient
of a power-series ring in that many
variables over the DVR.  Its dimension
is therefore at most $d+20C+1$.
Apply this to $R_x/I_S$ and the
injectivity in (UX.48), proving (UX.47).
\end{proof}

\begin{proposition}[UX9: trace fibres and coincident local characters]
Let $J\subset B$ satisfy the literal five-field
promodular-support contract for $T,X,\iota$.
For a prime $\mathfrak q\subset R_x$ at which the
generic representation is absolutely irreducible,
put $\mathfrak q_B=\mathfrak q\cap B$.  Then
\[
       \dim R_x/\mathfrak q
             \le \dim B/\mathfrak q_B+2+4s .
 \tag{UX.54}
\]
If a prime $\mathfrak b\subset B$ contains
$J+\mathfrak m_\Lambda B$ and, for one $v$ and
$i\ne j$, all
$\psi_{v,i}(h)-\psi_{v,j}(h)$ with $h\in T.U_1(v)$,
then
\[
                    \dim B/\mathfrak b+d_v\le2d+1 .
 \tag{UX.55}
\]
Only the \texttt{fin} field of the support contract
is used in (UX.55).
\end{proposition}

\begin{proof}
First let $L$ be the fraction field of the image
of $A$ in $R_x/\mathfrak q$, and let $E$ be
the fraction field of $B/\mathfrak q_B$.
All representation traces and multipliers lie
in $E$.  Absolute irreducibility implies that
the algebra spanned by its matrices is $M_4(L)$.
Choose sixteen matrices forming a basis.  The
trace pairing on $M_4$ is nondegenerate, so
the inverse of their Gram matrix expresses
all multiplication structure constants and
all other representation matrices using
traces of products.  They lie in $E$.
Thus the representation algebra descends to
a central simple $E$-algebra of dimension $16$.
After a finite algebraic extension of $E$
it splits, and the representation matrices
are simultaneously defined over that
extension up to conjugation.

The unique alternating form with the
specified multiplier is defined there up
to scalar: solve its linear entry equations
over that field and use irreducibility
to see that their solution space is one
dimensional.  A symplectic basis and then
an ordered eigenbasis for $r(g_0)$ may be
chosen after another finite algebraic
extension.  The separating eigenvalues
are roots of $c_{g_0}\in E[X]$.  Both
the original and the new $g_0$ matrices
are diagonal with this same ordering,
so their conjugator is diagonal.
The fixed form forces
$h_0h_3=h_1h_2$ up to a scalar which
is algebraic over $E$ and can be
absorbed after adjoining its square
root.  Conjugation is insensitive to
a common scalar.  Thus only two
diagonal parameters are needed to
recover all original representation
entries over an algebraic extension
of $E$.  This argument can be performed
in an algebraic closure of $L$, which
does not change transcendence degrees.

A lower unipotent symplectic flag is
determined by its entries $10,20,30,21$:
its other nonconstant entries are
$n_{32}=-n_{10}$ and
$n_{31}=n_{20}-n_{21}n_{10}$.
Consequently
\[
                 \operatorname{trdeg}_E L\le2+4s .
 \tag{UX.56}
\]
Every entry of $A$ is generated over $B$
by these actual matrices and flags,
which is why (UX.56) bounds this field,
not just the field of traces.
The complete local ring $B$ is excellent
and universally catenary.  The dimension
formula for a finite type algebra at a
closed point with finite residue extension
gives
$\dim(A/\mathfrak q_A)_{\ker x}
 \le\dim B/\mathfrak q_B+
                 \operatorname{trdeg}_E L$.
Its completion is
$R_x/\mathfrak q_A R_x$ and has the
same dimension; its quotient
$R_x/\mathfrak q$ has no larger dimension.
This proves (UX.54), including the
completion step.

For (UX.55) put $\overline B=B/\mathfrak b$,
a complete local domain of characteristic
$\ell$.  The actual local characters
take inertia into $1+\mathfrak m_{\overline B}$
because $T.\mathtt{Ok.ineL}$ makes the
original residual representation trivial
on inertia above $\ell$.  Their products
in symplectically paired positions equal
$1$ there: the multiplier reduces to
its fixed coefficient residue.  The
torsion part of a pro-$\ell$ character
is zero over this domain, since
$u^{\ell^a}=1$ implies $(u-1)^{\ell^a}=0$.
Local reciprocity and
$\mathcal O_{F_v}^{\times}\{\text{pro-}\ell\}
 \simeq(\text{finite torsion})\times\mathbb Z_\ell^{d_v}$
therefore give $d_v$ inertia elements
which determine each such character.
The same fixed number $d_v=d/s$ works
at every place.

Also choose one inertia element $h_0$
whose cyclotomic value generates the
pro-$\ell$ part of the cyclotomic image.
Two characters into $1+\mathfrak m_{\overline B}$
which agree on $\ker c_\ell$ and on
$h_0$ agree on all inertia: their
ratio descends to a procyclic quotient,
and its prime-to-$\ell$ part is
automatically trivial.  Here
$I_v\cap\ker c_\ell=T.U_1(v)$,
since the residual kernel condition
is automatic on $I_v$.

Choose $a\in\{0,1,2,3\}$ so that
$\{a,a^{\rm rev},i,i^{\rm rev}\}$
is the whole set and either $a=j$
or $j=i^{\rm rev}$.  In $B/\mathfrak b$
use the following parameters:
the $d_v$ pinning values
$\psi_{u,a}(h)-1$ at each place $u$;
the $d_v$ values $\psi_{u,i}(h)-1$
at all $u\ne v$; and the single
value $\psi_{v,i}(h_0)-\psi_{v,j}(h_0)$.
Their number is
\[
               sd_v+(s-1)d_v+1=2d-d_v+1 .
 \tag{UX.57}
\]
They lie in the maximal ideal.  In
any prime quotient killing them,
all $\psi_{u,a}$ are $1$ on inertia,
and all $\psi_{u,i}$ with $u\ne v$
are $1$.  At $v$ the hypothesis on
$U_1(v)$ and the additional parameter
make $\psi_{v,i}=\psi_{v,j}$ on
inertia.  If $a=j$, this makes it
$1$; if $j=i^{\rm rev}$, the paired
product gives $\psi_{v,i}^2=1$,
and its residue $1$ with $2$ invertible
again makes it $1$.  The paired
characters now make all four local
characters trivial on all inertia.

The precise \texttt{fin} contract says
that some power of $\mathfrak m_B$
is contained in
\[
 J+\mathfrak m_\Lambda B+
   (\psi_{u,m}(h)-1:\ u,\ m,\ h\in I_u).
 \tag{UX.58}
\]
It therefore forces every such prime
to be the maximal ideal.  The parameters
in (UX.57) have maximal radical in
$B/\mathfrak b$.  The height theorem
bounds its dimension by their number,
proving (UX.55).  This argument neither
requires $B/J$ finite over $\Lambda$
nor replaces (UX.58) by mere generic
finiteness.
\end{proof}

\subsection{Reaching every component at a good point}

For a prime $\mathfrak p\subset A$, write
$\operatorname{SomeIn}(J,\mathfrak p)$ if some
minimal prime $\mathfrak q\subset\mathfrak p$
contains $JA$, and write
$\operatorname{AllIn}(J,\mathfrak p)$ if all
such minimal primes do.  For a map $y:A\to D$
to a domain, $\operatorname{LocZero}(J,y)$
means that every prime $\mathfrak b\subset B$
annihilated by $y|_B$ contains $J$.  These are
the exact source predicates, with their
quantifiers over primes retained.

\begin{proposition}[UX10: quantitative component propagation and residual lines]
Assume (UX.24)--(UX.25) and the literal support
contract for $T,X,J,\iota$.
At any good residual point $x:A\to k_\Lambda$,
\[
 \operatorname{SomeIn}(J,\ker x)
       \ \Longrightarrow\
 \operatorname{AllIn}(J,\ker x).
 \tag{UX.59}
\]
Moreover, if two residual points are the ends
of a finite chain of literal polynomial links
whose endpoints are good, support at one
endpoint in the sense \texttt{Some\allowbreak In} implies
\texttt{All\allowbreak In} at the other.  Only the
\texttt{fin} and \texttt{p1} fields of the
support contract are used in this proposition.
\end{proposition}

\begin{proof}
Put
\[
                    N=2d+1-2s .
 \tag{UX.60}
\]
It is a positive integer.  The numerical
choices imply the two strict inequalities
\[
 d+20C+1<N,\qquad
 2d+1+(2+4s)<N+d_v .
 \tag{UX.61}
\]
Indeed $s\le C$ reduces the first to
$d>22C$, and the second to $d_v>2+6s$;
both follow from (UX.25).

Complete at $x$, obtaining $R=R_x$.
By UX7 it is a quotient of the regular
complete local ring
$\Lambda[[z_1,\ldots,z_a]]$ by $b$
elements with $a-b\ge N+1$.
Grothendieck's connectedness theorem gives
a chain between any two minimal primes
of $R$ such that adjacent primes $P,Q$
satisfy
\[
                       \dim R/(P+Q)\ge N+1 .
 \tag{UX.62}
\]
Here is the exact specialization of the
classical theorem.  The power-series
domain has dimension $a+1$ and cannot
be disconnected by removing any proper
closed subset.  The local cohomological
dimension of a $b$-generated ideal is
at most $b$, by its \v Cech complex.
Stacks Project, Lemma 52.14.1
(\href{https://stacks.math.columbia.edu/tag/0ECR}{tag 0ECR}),
therefore says that its quotient remains
connected after removing a closed subset
of dimension less than $a-b$.
If the graph of minimal components
joined by intersections of dimension
at least $a-b$ were disconnected,
remove the finitely many cross
intersections between its two parts.
The resulting space would be disconnected,
contradicting the theorem.  This proves
(UX.62), without assuming that $R$
is a complete intersection.

We show that $JR\subset P$ propagates
across one such edge to $JR\subset Q$.
Let $\varpi$ be a uniformizer of
$\Lambda$.  The first family of bad
ideals is the nine $I_S$ of UX8.
The second consists, for every $v,i\ne j$,
of
\[
 B_{vij}=JR+(\varpi)+
     (\psi_{v,i}(h)-\psi_{v,j}(h):
                                      h\in U_1(v))R .
 \tag{UX.63}
\]
For every prime $\mathfrak q$ containing
one of these latter ideals but none of
the $I_S$, the generic representation
is absolutely irreducible by (UX.46).
UX9 and (UX.61) give
\[
        \dim R/\mathfrak q
        \le2d+1-d_v+(2+4s)<N .
 \tag{UX.64}
\]
Every quotient by an $I_S$ has
dimension less than $N$ by (UX.47)
and (UX.61).

Cut $R/(P+Q)$ by $\varpi$.
The principal ideal theorem leaves
a component of dimension at least $N$.
Choose its prime $\eta$.  It cannot
contain an $I_S$, by the preceding
bound.  Nor can it contain a $B_{vij}$,
by (UX.64).  For each bad ideal choose
one element outside $\eta$ and take
their product $f$, which remains
nonzero in the complete local domain
$R/\eta$.  There exists a
one-dimensional prime quotient of
this domain on which $f$ remains
nonzero.  For an explicit construction
in dimension $n>1$, extend $f$ to
a system of parameters
$f,t_2,\ldots,t_n$.  A minimal prime
over $(t_2,\ldots,t_n)$ avoiding $f$
exists by the height theorem; its
quotient has dimension one because
adding $f$ makes the ideal maximal.
For $n=1$ take the zero prime.
Let $p\subset R$ be its inverse
image.  Then
\[
 P+Q+(\varpi)\subset p,\qquad
 \dim R/p=1,\qquad
 I_S\not\subset p,\ B_{vij}\not\subset p .
 \tag{UX.65}
\]

The normalization of the complete
one-dimensional domain $R/p$ is
finite, since complete local rings
are Nagata.  It is local: a finite
algebra over a complete local ring
is a product of its local factors,
and a subring of a field has no
nontrivial idempotents.  Hence this
normalization is a complete DVR.
Its residue field $k'$ is a finite
extension of the finite field $k_\Lambda$.
It has characteristic $\ell$ because
$\varpi\in p$.  The equal-characteristic
complete DVR structure theorem
identifies it with $k'[[t]]$.
The coefficient-field embedding can
be chosen to extend the given
$k_\Lambda$: its finite field
elements are the unique lifts of
their simple equations
$X^{|k_\Lambda|}-X=0$.
Thus there is a genuine
$\Lambda$-algebra map
\[
 R\longrightarrow k'[[t]]
                  \quad\hbox{with kernel }p .
 \tag{UX.66}
\]
The normalization and structure
theorems used here are, precisely,
Stacks Project, Section 33.41
(\href{https://stacks.math.columbia.edu/tag/0C44}{tag 0C44})
and Lemma 10.160.10
(\href{https://stacks.math.columbia.edu/tag/0C0S}{tag 0C0S}).

Compose (UX.66) with $A\to R$.
This is one of the literal finite-field
power-series points in \texttt{p1}.
Its base maximal ideal has zero
constant coefficient, since both
maps are local.  It kills $J$ because
$JR\subset P\subset p$.
Its generic representation is
absolutely irreducible because
it avoids all $I_S$.  Every pair
of local characters is distinct
on some $U_1(v)$: otherwise all
their differences would lie in
$p$, contradicting (UX.65).
Apply \texttt{p1} to obtain
$\operatorname{LocZero}$ for this
actual point.  The prime
$Q\cap B$ is annihilated by it,
as $Q\subset p$; therefore
$J\subset Q\cap B$, as required.
No use of purity or characteristic-zero
automorphy was made in this step.

Faithful flatness of
$A_{\ker x}\to R$ lifts each
minimal prime of $A$ below
$\ker x$ to a minimal prime of
$R$, with that contraction.
To see the minimal assertion,
first use lying over for a
faithfully flat map and then
choose a minimal prime below
the lifted prime; going down
forces its contraction to be
the original minimal prime.
Support on one original
component therefore gives
support on one component
of $R$.  Propagation along
(UX.62) gives support on
every component of $R$,
and contraction gives
(UX.59).

Finally a polynomial link gives,
by UX6, a map $A\to k_\Lambda[X]$
whose evaluations are the two
actual endpoint maps.  Its
kernel is prime and is contained
in both endpoint maximal ideals.
Choose a minimal prime below
this common kernel.  If
\texttt{All\allowbreak In} holds at the
first endpoint, that minimal
prime contains $JA$, so
\texttt{Some\allowbreak In} holds at the
second.  Apply (UX.59) there.
Induction along a finite chain,
in either direction, proves
the final assertion.  This
uses goodness only at the
endpoints; it does not
assume all points of a
link are good or automorphic.
\end{proof}

\subsection{Returning to the original valued representation}

\begin{proposition}[UX11: integral restriction and exact coefficient descent]
Let the original global pair have the lattice shape of UX2,
and let $T$ be an extension of its initial tower as in UX1.
There is a small pair $P$ on $G_F$ whose representation is
the restriction of that same global representation,
whose residue is the prescribed residual restriction,
and whose local inertia characters satisfy
\[
 \psi_{v,i}(s)=a_O\bigl(D.\mathrm{ps}_i(s|_{\mathbb Q})\bigr)
 \qquad(s\in I_{F_v}).
 \tag{UX.67}
\]
For every small pair $P$ there is a coefficient threshold
$O_3$ such that, for every coefficient DVR
$\Lambda\supset O_3$, it descends along an injective map
$f:R\to V$, where $R$ is a complete Noetherian local
$\Lambda$-algebra with residue field exactly $k_\Lambda$,
and
\[
 f\circ a_{\Lambda,R}=a_\Lambda,
 \qquad \operatorname{res}_V\circ f
       =j_\Lambda\circ\operatorname{res}_R.
 \tag{UX.68}
\]
The descended pair is continuous for the maximal-ideal
topology of $R$.  Independently, every finite residual
chain in the ambient residue field descends, with all
its actual polynomial links and goodness conditions,
to $k_\Lambda$ for every $\Lambda$ above a single
coefficient threshold.
\end{proposition}

\begin{proof}
First choose a complete DVR $A$ and an injective local
map $A\to V$ witnessing smallness of the global pair.
For each of the finitely many places $v\mid\ell$, choose
the prescribed Galois element carrying the original
place over $\ell$ to the place selected in $T$.
Its representation matrix and the original ordinary flag
are matrices over $A$.  Their product has all leading
principal minors invertible.  This can be checked in
the residue field: support $N_P$ and flag support $Q$
make these minors, successively, the products of the
first one, two, three and four diagonal characters.
For instance the first three rows of their product have
the form
\[
 \begin{pmatrix}
  a&0&0&0\\ *&b&*&0\\ *&0&c&0
 \end{pmatrix},\qquad abc\ne0.
 \tag{UX.69}
\]
Gaussian elimination over $A$ therefore gives a unique
lower unipotent factor $n_v$ and an invertible upper
factor.  Both have the required symplectic similitude
properties: apply the adjoint relation to the
factorization and use uniqueness of lower-times-upper
factorization.  The lower factor has multiplier $1$.
This construction uses divisions by units only.

Conjugating the original local representation by this
factor gives an upper triangular representation.  Its
ordered diagonal characters are unchanged by the
remaining upper triangular conjugation.  Since the
original characters are abelian, conjugating the
decomposition group changes no character value.
This proves (UX.67) with the original indices $i$.
The tower label permutes residual global characters;
it does \emph{not} permute the integral inertia characters
in (UX.67).  The source's label condition is precisely
what makes the two residual descriptions agree on
these decomposition groups.

The same elimination in the residue field gives flag
support $Q$ and the residual restriction appearing in
UX4.  Restriction preserves normalization at $g_0$,
because the later $g_0$ is the specified $\ell$-power
of the earlier one.  The original unramified conditions,
together with the tower condition killing $\psi$ on
the remaining forbidden inertia, give the exact
\texttt{Is\allowbreak Pair} inertia condition outside $\ell$.
For $\mu_v$, the image in the local abelianization is
$\ell$-primary torsion.  Local class field theory places
it in inertia; the tower's cyclotomic and character
conditions make the original inertia characters
trivial there.  Thus the required $\mu_v$ condition
also holds.  The finitely many matrices $n_v$ were
constructed in $A$, and the restricted representation
is continuous there, so the new pair is still small.

We prove the coefficient assertion without assuming that
the ambient valued field has characteristic zero.
Write $f_A:A\to V$ for the small DVR map.  The ambient
field $K$ is algebraic over $\operatorname{Frac}A$ by
the definition of smallness.  A useful consequence is
\[
 \bigcap_{n\ge0}f_A(\pi)^nV=0,
 \tag{UX.70}
\]
where $\pi$ is a uniformizer of $A$.
Indeed, if nonzero $x$ belonged to the intersection,
an algebraic relation for $x$, after removing a power
of its variable and clearing denominators, would have
nonzero constant term $u\pi^e$ with $u\in A^\times$.
Writing $x=f_A(\pi)^{e+1}z$ in that relation would
put the unit $f_A(u)$ in $f_A(\pi)V$, a contradiction.
This also explains why no unspecified continuity of
the ambient valuation is being used.

For $t\in\mathbb Z_\ell$, choose ordinary integers
$t_n\equiv t\pmod{\ell^n}$.  Their images in $A$
form a Cauchy sequence, since $\ell\in\mathfrak m_A$;
let its limit be $t_A$.  In characteristic $\ell$
the sequence is eventually constant.  Compatibility
of the given coefficient maps with integers shows
that
$f_A(t_A)-a_O(t)$ belongs to every
$f_A(\pi)^nV$.  By (UX.70) it is zero.  Consequently
the image of $\mathbb Z_\ell$ under every given
coefficient map already lies in $f_A(A)$.

The residue field of $A$ is finite.  Choose one lift
of each of its elements.  The ambient residue
approximation property puts each of their residues
in the image of some coefficient residue field.
A finite compositum gives $O_3$ such that, for every
$\Lambda\supset O_3$, all these residues belong to
$j_\Lambda(k_\Lambda)$.
Choose finitely many $\mathbb Z_\ell$-module generators
of $\Lambda$.  Adjoin their prescribed images in $K$
to $\operatorname{Frac}A$, obtaining a finite field
extension $E$.  Then
\[
 A'=E\cap V
 \tag{UX.71}
\]
is a complete DVR, finite over $A$, with finite residue
field; it contains the image of $\Lambda$ and of $A$.
Here is an elementary finiteness argument suited to
the possibly positive-characteristic ambient field.
The ring $A'$ is a valuation ring of finite $A$-rank
at most $[E:\operatorname{Frac}A]$.  Independent
residue classes in $A'/\pi A'$ lift to $A$-linearly
independent elements: divide a putative relation by
the smallest valuation of a nonzero coefficient and
reduce modulo $\pi$.  Thus $A'/\pi A'$ is finite
over $A/\pi A$.  Lift a finite spanning family and
successively express each element modulo
$\pi,\pi^2,\ldots$.  Completeness of $A$ sums the
coefficient series; (UX.70) says the remaining error
is zero.  The family therefore generates $A'$ over
$A$.  Hence $A'$ is Noetherian; a nonfield
Noetherian valuation ring is a DVR.  Finiteness
over complete $A$ gives completeness, and its residue
field is finite.  This proves all claims in (UX.71).

Let $k'$ be the residue field of $A'$, and let
$k_\Lambda\hookrightarrow k'$ be the injection
induced by the actual coefficient map.  Define
\[
 R=\{x\in A':\bar x\in k_\Lambda\}.
 \tag{UX.72}
\]
Its maximal ideal is $\mathfrak m_{A'}$ as a set,
and its $n$th power is $\mathfrak m_{A'}^n$.
For the latter assertion write any element
$\pi'^na$ as the product of $n-1$ copies of
$\pi'\in R$ and $\pi'a\in R$.
The ring $R$ is closed and complete, has residue
field $k_\Lambda$, and has finite-dimensional
cotangent space: $k'$ is finite over $k_\Lambda$.
Lift a finite spanning family of that cotangent
space.  Successive approximation, with the residue
field supplied by $\Lambda$, gives a surjection
from a power-series ring in finitely many variables
over $\Lambda$ to $R$.  Thus $R$ is Noetherian.
The choice of $O_3$ puts all of $A$ inside $R$.
Its inclusion is local, hence carries
$\mathfrak m_A^n$ into $\mathfrak m_R^n$.
The original continuous pair therefore descends
continuously to $R$.  The inclusion $R\hookrightarrow V$
and the definition (UX.72) give both exact identities
(UX.68).  Notice that $R$ need not itself be a DVR;
its actual complete local residue-field property is
what the universal space requires.

Finally, a locally constant map from a compact Galois
group to the discrete ambient residue field has finite
image.  Hence every residual representation and its
finitely many flags use only finitely many residue
coefficients.  The same assertion holds for an actual
polynomial link: its matrix-entry maps into
$k_V[X]$ are locally constant, so only finitely many
polynomials, and therefore finitely many polynomial
coefficients, occur.  Apply the ambient residue
approximation property to these finite sets for
every point and every link of the finite chain.
Their finite compositum is the required threshold.
For every larger $\Lambda$, all objects descend
entry by entry to $k_\Lambda$ or $k_\Lambda[X]$.
Injectivity of $j_\Lambda$ descends multiplication,
symplectic, flag, inertia and support identities.
It also preserves nonzero-entry witnesses and the
nontrivial-character witnesses in goodness.  Since
the original maps have finite image, local constancy
descends as well.  Evaluation at $0$ and $1$ commutes
with the coefficient injection, so the actual links
have exactly their original endpoints.
\end{proof}

\begin{theorem}[UX12: the full case-II support witness]
Assume the literal \texttt{Seed\allowbreak IIAlt} and
\texttt{PMEx\allowbreak Alt}.  For every original case-II datum
$D$, every fixed $\iota$, every ambient valued field
and every actual small \texttt{KPair} $(r,u)$ with
$r_K$ \texttt{Star}, the literal
\texttt{Pro\allowbreak Mod\allowbreak Alt}$\,(D,\iota,V,r)$ holds.
The only further inputs are the cohomology and
Ferrero--Washington theorems already supplied by
BX and TW, and the residual vanishing supplied
by BR and AX.  In particular the original
representation, its multiplier, its ordered inertia
characters and the fixed $\iota$ are retained.
\end{theorem}

\begin{proof}
Choose the permitted label from the first-column
residual ordinary flag, as in the first paragraph
of the proof of UX2.  Apply UX1 to this label to
obtain the initial tame tower $T_0$ and its
separating element $g_0$.  Apply the remaining
construction of UX2 at this $g_0$ to put the
global representation in the normalized lattice
shape.  Let $r_1$ denote this integral representation
over $V$.  The change is conjugation over the
original $K$, and its multiplier is still
$a_O(D.\vartheta)$.  UX3--UX4 provide a later tame tower
$T_1$ such that every subsequent deep tower admits
an actual residual chain from the restriction of
this lattice to an actual integral seed.
All this is performed with the single initially
fixed $\iota$; the seed input has a universal
quantifier over $\iota$ and is applied at that
value only.

Apply UX5 at $T_1$ to choose its constant $C$.
Only afterwards choose a cyclotomic layer $T$
which is deep and satisfies
$d\ge30(C+1)^2$.  The bounds are uniform for
all residue coefficient fields and all later
layers, so they still hold at this $T$.
UX4 and the literal seed hypothesis now give
some $O_1\supset O$ and an actual integral
pair $Q$ over $O_1$ with
$T.\mathtt{SeedPairAlt}\,(\iota,Q)$.
UX11 gives the small restricted pair $P$ of
$r_1$, satisfying (UX.67).  The chain has the
form
\[
 \overline P\ \longleftrightarrow\ \cdots\
 \longleftrightarrow\ \overline Q
 \quad\hbox{in the one ambient residue field.}
 \tag{UX.73}
\]
It is the chain for this actual restriction and
these chosen places, not a chain to an isomorphic
representation on a newly chosen tower.

Take four coefficient thresholds: $O_1$ for the
integral seed, $O_2$ for the descent of (UX.73),
$O_3$ for the small-pair descent of UX11, and
$O_4$ from \texttt{PMEx\allowbreak Alt} at this $D,T,\iota$.
Choose a single coefficient DVR $\Lambda$
containing them all.  Such a finite compositum
is again an allowed coefficient DVR inside the
fixed $\overline{\mathbb Q}_\ell$.
The order of these choices matters:
\texttt{PMEx\allowbreak Alt} supplies $O_4$ before the
universal space is chosen, and promises its
conclusion for \emph{every} space over every
later coefficient DVR.  We therefore use UX6
to construct the actual space $X=(A,B,P_X)$
over $\Lambda$, and then obtain its ideal
$J\subset B$ with all five clauses of
$T.\mathtt{PotentialModularityTower}\,(\iota,X,J)$.
They include the stated finiteness, every
finite-field power-series point, every seed,
and both characteristic-zero conclusions.
No replacement of those clauses by a pointwise
or generic version has occurred.

Write $s,t$ for the descended endpoints of
(UX.73) in $k_\Lambda$.  UX11 gives an
injective local map $f:R\to V$ and a
continuous pair $P_R$ lifting $s$, with
$P_R\otimes_fV=P$.  The seed extends to
the integral pair $Q_\Lambda$ over
$\Lambda$, lifting $t$.
All equalities follow from the actual
coefficient inclusions and the injective
residue map $j_\Lambda$.  Universality of
this one space gives
\[
 f_0:A\longrightarrow R,
 \qquad f_1:A\longrightarrow\Lambda,
 \qquad P_X\otimes_{f_0}R=P_R,
 \quad P_X\otimes_{f_1}\Lambda=Q_\Lambda.
 \tag{UX.74}
\]
These are maps of $\Lambda$-algebras; their
reductions are precisely $s$ and $t$.

Compose $f_1$ with the given inclusion
$\Lambda\hookrightarrow\overline{\mathbb Q}_\ell$.
The resulting point is the original seed at
the fixed $\iota$, so the \texttt{seed}
clause gives \texttt{Loc\allowbreak Zero}.  Choose a
minimal prime $q$ of $A$ inside the kernel
of this characteristic-zero point.
Its contraction to $B$ is a prime annihilated
by the point; hence \texttt{Loc\allowbreak Zero} says
$J\subset q\cap B$.
The kernel is contained in the kernel of
the residue point $t$, since the map
$\Lambda\to\overline{\mathbb Q}_\ell$
is injective.  Thus
\[
 \mathtt{SomeIn}(J,\ker t).
 \tag{UX.75}
\]
This step does not assume that the seed
already annihilates $J$; support follows
from the stronger literal seed clause.

UX7--UX10 apply at the chosen uniform
bound and prove \texttt{All\allowbreak In} at every
good endpoint along (UX.73), in reverse
order starting from (UX.75).  In particular
every minimal prime $q\subset\ker s$
contains $JA$.
Set
\[
 y=f\circ f_0:A\longrightarrow V.
 \tag{UX.76}
\]
Its kernel is prime because $V$ is a
domain.  It lies in $\ker s$: injectivity
of $f$ says that $y(z)=0$ implies
$f_0(z)=0$, and then the reduction of
$f_0(z)$ is zero.  Choose a minimal prime
inside $\ker y$.  The just proved
\texttt{All\allowbreak In} statement gives
$JA\subset\ker y$.
Thus $y$ annihilates $J$ on the actual
valued point.

The remaining defining clauses can now
be checked as equalities rather than
existential compatibility assertions.
From (UX.68) and (UX.74),
\[
 \begin{gathered}
 y\circ a_{\Lambda,A}=a_\Lambda,
 \qquad y(\mathfrak m_B)\subset\mathfrak m_V,
 \qquad P_X\otimes_yV=P,\\
 r_1|_{G_F}=(P_X\otimes_yV).r,\qquad
 y\bigl(P_X.\psi_{v,i}(s)\bigr)
    =a_O\bigl(D.\mathrm{ps}_i(s|_{\mathbb Q})\bigr)
       \quad(s\in I_{F_v}).
 \end{gathered}
 \tag{UX.77}
\]
For the maximal-ideal assertion, the
restriction $B\to R$ is a map of
complete local $\Lambda$-algebras with
the same residue field; (UX.68) carries
its zero residues to zero residues in
$V$.  The last assertion is exactly
(UX.67), with the original indices.
The tower is fit and deep, its full
five-clause ideal was supplied above,
$r_1$ has the original multiplier, and
$r_K$ and $(r_1)_K$ are conjugate over
the same ambient $K$.  Together these
are every clause of the literal
\texttt{Pro\allowbreak Mod\allowbreak Alt} definition.

All universal choices in the conclusion
were arbitrary.  The existential tower
and coefficient DVR may depend on the
original point, as the source permits;
the residue bounds are uniform before
their use, and the \texttt{PMEx\allowbreak Alt}
coefficient threshold is uniform in
the subsequently chosen space.
No change of $\iota$ is made anywhere
in the construction.
\end{proof}

\subsection{The exact remaining boundary}

UX12 closes the original case-II implication
\texttt{pro\allowbreak Mod\_\allowbreak of\_\allowbreak pm\allowbreak Ex\_\allowbreak of\_\allowbreak ctx2\_\allowbreak of\_\allowbreak k\allowbreak Pair\_\allowbreak of\_\allowbreak star}.
It supplies the actual witness consumed by IX10,
including all coefficient, field, inertia and
support equalities.  Its two automorphic inputs
are the literal \texttt{Seed\allowbreak IIAlt} and
\texttt{PMEx\allowbreak Alt}; their independent producers
belong to SX and QX.  Applying any such producer
must preserve its own hypotheses.  In particular
this chapter does not infer the five-clause
support ideal from generic automorphic finiteness.

The case-III analogue, based on
\texttt{Seed\allowbreak IVAlt}, has additional tower and
residual-pattern quantifiers.  UX12 makes no
claim to that analogue or to universal
\texttt{Symplectic\allowbreak Four\allowbreak Image\allowbreak Dense}.  Consequently
the older continuation targets combining both
cases must retain their unsupplied case-III
scope even after this case-II chain is installed.



% Version 0.29. One complete conditional case-II count chain.
% Pinned source: 7cbf7a0c32fca1badf05931879a37c3f8a2b03b4.
\section{Closing the case-II chain through the original count}
\label{sec:case-two-installation}

\status{The case-II seed is supplied by SX8, and UX12 constructs
the complete actual case-II promodular witness from that seed and
\texttt{PMEx\allowbreak Alt}.  YX1 closes the resulting original-point
image-density statement in this case.  YX2 carries it through the
actual arithmetic datum, universal model, tangent comparison and
automorphic catalogue to the literal \texttt{Count\allowbreak II}.
The remaining input to this entire case-II chain is the actual
\texttt{PMEx\allowbreak Alt} producer, with its narrower sufficient support
input analyzed in QX.  Case III, its distinct seed and its
component-reaching argument remain separate.}

% BEGIN YX1
\begin{theorem}[YX1: case-II image density from the single tower input]
\label{claim:YX1}
Assume the literal \texttt{PMEx\allowbreak Alt}.  For every original
$D:\texttt{Raw}\,\Omega\,c$ with $\texttt{Ctx2}\,D$, and
every fixed $\iota:\overline{\mathbb Q}_\ell\simeq\mathbb C$,
\[
       D.\texttt{Fin0}\ \wedge\ D.\texttt{No\allowbreak End}\ \wedge\
              D.\texttt{All\allowbreak Irr}(\texttt{Is\allowbreak Aut4Alt}\,\iota).
                                                               \tag{YX.1}
\]
In particular the one embedding required for the case-II part of
the source image-density statement can be chosen before all its
varying points and coefficient witnesses.
\end{theorem}
\begin{proof}
Fix $D$ and $\iota$ first.  The exact \texttt{Ctx2} includes
$D.\texttt{Good}$ and an original host, exponent, character and
\texttt{Of\allowbreak Char} witness with the prescribed parameters
$(n_{\mathrm{II}},1,1)$.  BX8 supplies the whole cohomology
record, and TW1 supplies the Ferrero--Washington input in the
scope used by UX.  SX8 supplies the literal \texttt{Seed\allowbreak IIAlt}
for every one of these towers and the same chosen $\iota$.
The assumed \texttt{PMEx\allowbreak Alt} is its other arithmetic input.

Now take any original ambient $\mathfrak a$, any actual small
pair $(r,u)$ satisfying \texttt{KPair}, and suppose that the
representation in its ambient field is \texttt{Star}.
UX12 constructs an actual \texttt{Pro\allowbreak Mod\allowbreak Alt} witness for this
pair.  This includes its tower, coefficient inclusion, original
\texttt{Space}, ideal, specialization into the original
valuation ring, conjugate representation, restrictions and local
characters.  Its coefficient threshold is chosen in the order
required by \texttt{PMEx\allowbreak Alt}; none of these later choices changes
the already fixed $\iota$.
We have therefore supplied the universal star-point assertion
(IX.35) on this original $D$.

Here are its three uses in IX, to make the entire chain explicit.
For a characteristic-$\ell$ star pair, IX9 applies the finite
special-weight-fibre clause of its actual witness and obtains
constant residual characteristic polynomials.  The matrix
commutator obstruction excludes both irreducible full span and
the nonsplit \texttt{Block\allowbreak Ext} alternative.  This is the whole
$\operatorname{Con}(D)$ assertion.  In characteristic zero,
the same witness and IX10 give the precise weak Frobenius-root
ratio $\operatorname{PureW}$ on every star pair, including a
nonsplit extension constructed from a block diagonal point.
On full-span pairs IX10 also transports the actual automorphic
attachment and its \texttt{Wit\allowbreak A} data to
\texttt{Is\allowbreak Aut4Alt} at this same $\iota$.

IX8 now gives (YX.1).  Its proof includes the actual
characteristic-$\ell$ curve arising from failure of
\texttt{Fin0}, the integral nonsplit block construction, and
the weak-purity contradiction proving \texttt{No\allowbreak End}.
The ordinary flags used there come from IX2 on the actual
points.  Thus the conclusion covers every original lift and
every original irreducible algebraic point in the required
quantifier order.  It uses no case-III witness.
\end{proof}
% END YX1

% BEGIN YX2
\begin{theorem}[YX2: the complete case-II count from the single tower input]
\label{claim:YX2}
With all the completed suppliers above, one has the literal
implication
\[
                         \boxed{\ \texttt{PMEx\allowbreak Alt}
                                    \Longrightarrow\texttt{Count\allowbreak II}.\ }
                                                               \tag{YX.2}
\]
Precisely, for each primitive source character $\chi$ there is
$C_\chi\in\mathbb N$, chosen before
$h,t,\Omega,\ell,c,e,\psi$, such that every character in the
original Case II has finite strict Selmer length and
\[
        \frac{\ell_{c.O} V(\Omega,c,e,\psi)}{v_{c.O}(\ell)}
                    \log\ell\ \le (\ell h.N e)^{C_\chi}.
                                                               \tag{YX.3}
\]
Consequently \texttt{Count\allowbreak I}, \texttt{Count\allowbreak III},
\texttt{Count\allowbreak V}, and \texttt{PMEx\allowbreak Alt} suffice for the original
\texttt{Point\allowbreak HV}; the previously stated prime-pair hypotheses
then give the common cohomological-window consequences of
HV/TW/OB/VW.
\end{theorem}
\begin{proof}
Fix the tower input once.  CX8--CX9 supply the catalogue constants
$C,A$ before any host, coefficient ring or varying representation.
For a fixed primitive $\chi$, use the arithmetic absorption in
SP10 to choose the final exponent $C_\chi$ from these constants
and $\chi$.  The proof of that absorption uses only the source
host inequalities and the explicit level/exponent formulas below,
so this choice precedes all the remaining quantifiers.

Take an original host $h$ with $h.\chi=\chi$, an admissible
$t$, an odd $\ell$, the original $\Omega,c$, and
$\psi\in h.X(c,e)$ with $0<e<t$ in Case II.
First suppose that $O=c.O$ contains the finite dyadic root sets
required in AX5.  That construction gives an actual $D$ with
$D.\texttt{Good}$ and
\[
          \texttt{Of\allowbreak Char}(D,e,\psi,n',1,1),
                         \qquad n'=h.n_{\mathrm{II}}(\ell,\psi).
                                                               \tag{YX.4}
\]
Together with the data just fixed this is exactly
$\texttt{Ctx2}\,D$.  Choose one $\iota$ and apply YX1.
It supplies \texttt{Fin0}, \texttt{No\allowbreak End} and the automorphy
of every irreducible actual point on this same datum.

AX7 gives the original common level and character exponent
\[
           L=\ell 2^{n'}h.N,
       \qquad k=(\ell-1)2^{n'}(2h.o).
                                                               \tag{YX.5}
\]
All three finite characters of $D$ satisfy
\texttt{Kills\allowbreak Mu}(L) and have $k$th power one.
AX6, AX8 and AX9 give the actual continuity, ordinary weights,
level, dual twist and safety properties of every point.
Adding its automorphy from YX1 and its full matrix span gives
exactly the rank-four \texttt{Adm\allowbreak RAlt} record, for this one
$\iota$, the original $L$, and weights $(e+1,e,1,0)$.
This repeats the single pointwise use (AX.34) in AX10 and
completes G1 for Case II.

DX2, DX4 and DX8 give the original finiteness, dimensions and
local ring, so MX10 constructs the actual \texttt{Model}
over this same $O$.  Apply RX6 to each irreducible actual point,
using the automorphy supplied by YX1.  Its conclusion is
$D.\texttt{All\allowbreak Irr}(D.\texttt{Rigid\allowbreak At})$, with all the original
local tangent conditions.  PX1--PX9 now supply the point,
conjugacy, orbit and tangent comparisons on this model.
With \texttt{Fin0} and \texttt{No\allowbreak End} already available,
these are all of G2 used by SP10, and PX10 gives the actual
cotangent/trace inequality.  This supplies the geometric input
on the same original object as the arithmetic input.

For clarity, the resulting numerical step can be followed
without any additional image-density assumption.  Choose the
source prime $p=p_0(\ell h.N)\ge5$ and arithmetic Frobenius
$\sigma$ at the specified place.  Then $p\nmid L$ and
$p\notin D.S$.  The original diagonal representation has trace
\[
       \iota^{-1}(p^{e+1}z_1+p^e z_2+pz_3+1),\qquad z_j^k=1.
                                                               \tag{YX.6}
\]
CX9 gives a nonzero integer $N_*$ bounding the product of all
the relevant trace differences over any finite pairwise
nonisomorphic family.  Its full global labels, conjugacy
invariance and finite root multiplicities have already been
checked in CX; its enlarged auxiliary automorphic-label level
does not change the original Galois level $L$ in (YX.5).
PX10 and the same product inequality as in (JX.6) give
\[
 \frac{\ell_O D.\mathrm{Sha}}{v_O(\ell)}\log\ell
       \le -2\log|N_*|_\ell
       \le 2\log|N_*|
       \le 2Ck[L(e+1)]^A\log p.
                                                               \tag{YX.7}
\]
The finite-length assertion is part of the first step.
The factor two and the original valuation normalization
are retained.

The strict Selmer module in (YX.3) injects into the auxiliary
one used here by SP10's \texttt{sha\allowbreak Eq} comparison.
At every newly admitted prime its negative-weight Frobenius
scalar is a root of unity times $p^{-e}$ and differs from $1$;
the scalar difference acts surjectively on $E/O$, giving the
needed local comparison.  SP10's arithmetic estimates absorb
(YX.7), with $p_1=1$ and the exact $n_{\mathrm{II}}$ in
(YX.4), into (YX.3).  Only the Case-II part of that argument
is used, so no case-III datum or seed enters the proof.

Finally let the initial $O$ be arbitrary.  Adjoin the two finite
dyadic root sets allowed in SP8 and work in a common finite
coefficient overring $O'$.  SP8 and RV1 preserve the kernel,
order, conductor, parity, case test and integer $n_{\mathrm{II}}$
of the original character, on the original algebraic closure.
The preceding construction applies over $O'$.  Those same
claims give equality of normalized strict Selmer lengths:
the length and $v(\ell)$ acquire the same ramification factor.
Returning to $O$ proves (YX.3) for its original coefficient
quantifier.  Neither $C_\chi$ nor any bound is changed by the
chosen extension.  This is exactly the source \texttt{Count\allowbreak Shape}
with the Case-II predicate, hence the literal \texttt{Count\allowbreak II}.

For the last assertion combine this count with the three
stated counts and apply the original four-count assembly HV10.
The separately proved \texttt{Count\allowbreak Vs}, \texttt{Count\allowbreak IV}, and
\texttt{Count\allowbreak Z} have their existing suppliers HV5, HV6 and CZ6.
The established window propagation retains its original field,
layer and prime-pair hypotheses.  In particular it makes no
claim that \texttt{Count\allowbreak III} or the actual tower producer has
already been supplied.
\end{proof}
% END YX2


\section{Dependency audit after the complete case-II chain}
\label{sec:dependency-audit-v29}

This audit records exact sufficient obligations, preserves every
version-0.28 semantic node, and includes prior snapshots for changed
nodes.  It is not an exhaustive Lean graph or a completion percentage.
Historical claims, source arrays and dependency edges are unchanged.

The principal new closed chain is
\[
 \texttt{PMEx\allowbreak Alt}
 \xRightarrow[\mathrm{SX8}]{\mathrm{UX12}}
       \text{actual case-II promodularity}
 \xRightarrow{\mathrm{YX1}}\text{case-II image density}
 \xRightarrow{\mathrm{YX2}}\texttt{Count\allowbreak II}.
\]
SX8 supplies the seed at every given permissible tower and fixed
embedding.  UX12 supplies the actual valued point, coefficient
maps, restrictions, ordered inertia and support ideal, using three
explicit residual lines in one actual deformation space.  YX uses
the existing original-point, model, rigidity and catalogue machinery
with every constant and normalized coefficient return in its
original quantifier order.  No case-III input is used in this chain.

KX supplies the exact universal arithmetic packages.  QX supplies
the literal adjoint lifting theorem, the actual two Links and
\texttt{Is\allowbreak Phi} maps, and the complete five-clause tower join.
Its retained arithmetic input is the actual rank-four
\texttt{Hida\allowbreak Props\allowbreak Alt} ideals together with \texttt{P1core} on
those Links, uniformly after one coefficient threshold.
The six fields \texttt{ring}, \texttt{dense}, \texttt{pt},
\texttt{end1}, \texttt{end2} and \texttt{twist} belong to the
same ideal and the same fixed embedding.  A generic dense subset
or an unrelated quotient does not supply this record.
The stronger all-rank packages and unused \texttt{P3core} are
not promoted.

The added sufficient height cut is
\[
 \texttt{Count\allowbreak I}\wedge\texttt{Count\allowbreak III}\wedge\texttt{Count\allowbreak V}
       \wedge\texttt{PMEx\allowbreak Alt}
       \xRightarrow{\mathrm{YX2}}\texttt{Point\allowbreak HV}.
\]
The existing TW/OB/VW assembly retains its prime-pair hypotheses
and its $C_t=401$ cutoff exponent $400$.  The earlier JX cut
through full two-case image density remains valid as stated.

The historical seed target combines cases II and IV, while the
promodularity target combines cases II and III.  Supplying SX8
and UX12 narrows those targets but does not retire their entire
scope.  The actual ordinary Hecke/support producer, case-IV seed,
case-III construction, remaining counts, and independent actual
family/gap/theta/local/denominator/metric/spectral branches remain.
The milestone is a complete conditional case-II dependency chain,
not the full human-readable RH/Dirichlet-GRH proof.
\section{Installing the actual target host on the cohomological windows}
\label{sec:host-window-installation}

\status{HX1--HX4 close a previously unjoined dependency, rather than
claim a new estimate for Galois cohomology.  Relative to the named
classical inputs already recorded for QP and the least-prime theorem
specified below, the literal implication
\texttt{VWin\allowbreak At}$(\psi,401)\Rightarrow
\texttt{Host\allowbreak Ex\allowbreak W}(3/2,(\texttt{GChi}\,\psi).\texttt{req})$
is now supplied.  The host is the one constructed from AH18, on the
same prime pair and dyadic field layer that occur in the window input.
The estimate \texttt{VWin\allowbreak At} itself is not asserted unconditionally.}

\paragraph{The input and output, with their quantifiers.}
Fix a primitive Dirichlet character $\psi$ of conductor $M$ and
order $o$.  The notation $\mathrm{VVle}(E,t,B)$ means both finiteness
of all the relevant $\ell$-primary strict-Selmer lengths in twists
$1\le j<t$, and the bound
\[
 \sum_{\ell\in S}\log\ell
       \sum_{1\le j<t}\operatorname{length}_{\ell}(E,j)\le B
 \quad\text{for every finite set of primes }S.
 \tag{HX.1}
\]
The lengths and local conditions in (HX.1) are the original ones:
the source quantifies over every algebraic closure and every compatible
local-data system.  It is not enough to check one finite set $S$ or one
substitute coefficient model.  Its elementary monotonicity is
\[
 \mathrm{VVle}(E,t,B),\quad t'\le t,\quad B\le B'
       \quad\Longrightarrow\quad\mathrm{VVle}(E,t',B').
 \tag{HX.2}
\]
This follows by restricting the twists and using the nonnegativity of
each length and of $\log\ell$.

The hypothesis $\mathrm{VWinAt}(\psi,401)$ has the following order.
For every $0<\varepsilon<1/52$ and every dyadic page function $e(x)$
with
\[
 x^{2+8\varepsilon}\le2^{e(x)}\le2x^{2+8\varepsilon}
 \qquad(x\ge1),
 \tag{HX.3}
\]
one chooses $J\ge1$ and $X_0$ before every $X\ge X_0$ and every
admissible pair $q,d$ at base $X^{1/J}$.  One then obtains a window
$[W,W^2]\subseteq[X^{1/J},X]$ on which every original target field
$E$ at page $e(x)$ obeys
\[
 \mathrm{VVle}\bigl(E,\lfloor x^{401}\rfloor,
                         \sqrt{\exp(x^\varepsilon)}\bigr).
 \tag{HX.4}
\]
The precise admissible-pair congruences and bounds are the
\texttt{Win\allowbreak Pair} record; QP22 supplies the stronger \texttt{Pair\allowbreak Good}
record with those same clauses.

The conclusion $\mathrm{HostExW}(3/2,(\mathrm{GChi}\,\psi).\mathrm{req})$
says that, for each $\varepsilon$ in that range, there are constants
$C_{T_2},C_8,C_5,c_6,C_b,x_0$ and one scale set containing arbitrarily
late intervals $[Z,Z^{3/2}]$, together with one fixed $J\ge1$.
The membership condition on a scale $s$ includes $s.J=J$; also
$s.\varepsilon=\varepsilon$.  At every allowed scale $x=s.x\ge x_0$
there are ranks $R$, an actual host $H$, and its auxiliary prime $q_0$ such
that
\[
 \begin{gathered}
 \mathrm{RankVal}(s,R),\qquad H.C_{T_2}=C_{T_2},\qquad
 \mathrm{HostGood}(s,C_8,C_5,c_6,C_b,H,q_0),\\
 \mathrm{EAtH}(\mathrm{Level.cyc}\,\psi,H,TT,x,\varepsilon,4),\qquad
 \mathrm{VVle}(H.E,R.t,\sqrt{s.b}).
 \end{gathered}
 \tag{HX.5}
\]
The scale set and constants precede all of these scale-by-scale
choices.  The hosts may vary with the scale; no single field is
asserted to serve every dyadic page of the interval.

\paragraph{HX1: the exact small-prime input is a classical application.}
The source's \texttt{Small\allowbreak One} is
\[
 \exists a\ge1\ \forall N\ge2\ \exists p\text{ prime}:
                p\equiv1\pmod N,\qquad p\le N^a.
 \tag{HX.6}
\]
Xylouris, \emph{On the least prime in an arithmetic progression and
estimates for the zeros of Dirichlet $L$-functions}, Acta Arithmetica
150 (2011), Theorem~1.1, gives an absolute $C>0$ with
$P(a,q)\le Cq^{5.18}$ for coprime positive $a,q$.
Only existence of an absolute exponent is needed here.  Put
$C_+=\max(C,1)$ and
\[
 a_0=\max\left(1,5.18+\frac{\log C_+}{\log2}\right).
\]
For $N\ge2$ the hypotheses hold with residue $1$ and modulus $N$,
and $C_+\le N^{\log C_+/\log2}$.  Hence the supplied prime satisfies
(HX.6) with $a=a_0$.  If the published estimate is read with a
large-modulus threshold, enlarge $C$ over the finitely many smaller
moduli first; Dirichlet's theorem supplies those finitely many primes.
There is no use of RH, GRH, a uniform prime number theorem with a
moving modulus, or a numerical value of the implicit constant.

\paragraph{HX2: ranks and admission are already supplied on this field.}
Suppose AH18 supplies a frame $\phi$ and $H_0:\mathrm{HostAx}(\phi)$,
with
$x^{2+8\varepsilon}\le|\phi|\le2x^{2+8\varepsilon}$.
SZ13 constructs the entire rank record over this exact frame.
Explicitly it chooses even $m,k$ from the two prescribed real rank
scales, sets $f=2|\phi|$, $r=k+m-1$, and defines $t$ by the exact
ceiling expression (SZ.32).  Its proof gives
$R.\mathrm{toFrame}=\phi$ and $\mathrm{RankVal}(s,R)$; it does not
replace $\phi$ by an abstract group of equal cardinality.

For every permitted set of host choices, every
$\mathrm{FieldRows}(H_0,b)$ and $q,d>4tr+1$, HC4 copies the prescribed
CM type, splitting different ideal, square roots and choices into
\texttt{Host\allowbreak NF}.  Thus it gives the full predicate called
\texttt{admitted} in the source, with exact equality of the underlying
\texttt{Host\allowbreak Ax} and of the choices.  AH18 supplies the field facts used
in that copying argument.  This proves the complete rank/admission
input of \texttt{admissible\allowbreak Classes\allowbreak Twisted\allowbreak Char}, uniformly before the
frame and field are selected.  No arithmetic existence assumption is
being hidden in the word ``admitted.''

\paragraph{HX3: all clauses of the conditional host completion.}
Fix $\varepsilon,C_0$ with $0<\varepsilon<1/52$ and $C_0\ge1$.
Use QR25's actual \texttt{Short\allowbreak Norm} and HX1 in HC5.  The output has
an absolute exponent $C_1$ before $\varepsilon$ and the host.
For every later datum with
\[
 N=Qqd,\qquad d\mid N,\qquad d\ge100,\qquad
 \log N\le3x^{\varepsilon/5}(\log x)^2,
 \tag{HX.7}
\]
it gives the two distinct norm primes $g_0,g_1$, their disjoint
$BB(r,g_i)$ sets, and the three smaller primes $p,p',q_0$ of HC5.
Their norms, integral generators and congruences are (HC.7)--(HC.9),
not merely assertions that the primes split.

Here is the complete input matching for the HC9 host join.
\begin{center}
\begin{tabular}{p{.33\linewidth}p{.58\linewidth}}
Input required on the original host & Supplied construction\\\hline
\texttt{Short\allowbreak Norm}, \texttt{Small\allowbreak One}&QR25 and HX1; the conductor and
interval restrictions in QR25 remain in the HC5 application.\\
Actual frame and field, and admission&AH18 and HX2, with the exact
\texttt{Host\allowbreak Ax} equality.\\
Pair rows, large $q,d$ and modulus bound&HC7: the given
\texttt{Pair\allowbreak Good} and \texttt{Pair\allowbreak Win}, with $Q\le C_0n$.\\
Five compatible auxiliary primes&HC5, applied with the actual $N=Qqd$;
the small primes lie below $e^{x^\varepsilon}$ and the norm primes
above it.\\
Second set and its two cardinality bounds&HC6 at the actual $L_w$;
one may use $C_{T_2}=64$.\\
Biased top set, shares and bad-prime bound&HC1--HC2, with the original
Frobenius steps and a record update changing only $TT$.\\
Real-character boxes and thin-character estimate&QP23--QP24 and
HC8--HC9 on that same host; $Q_{\mathrm{odd}}\mid M$ follows from
$Q\mid\operatorname{lcm}(2^{e+2},M)$.\\
Logarithmic derivative&HC8 gives $|dlogH|\le4\log x$; enlarge the
threshold until $4\log x\le x^\varepsilon$.\\
\end{tabular}
\end{center}

Choose $C_8=\max(C_0,C_1)$, $C_5=4$, and the supplied positive
$c_6,C_b$.  Take the maximum of the thresholds for these finitely
many results.  This is before every later $s,R,H_0,M,Y$.
The exact equality $H.\mathrm{toHostAx}=H_0$ preserves $Q,q,d,E$
and all the character data.  Installing HC2's $TT$ changes none of
them.  The norm identities give the required integral elements of
$\mathcal O_E\cap K$, congruent to $1$ modulo $Qqd$.
The congruences $p,p'\equiv1\pmod{\operatorname{lcm}(15,Qqd)}$
and $q_0\equiv1\pmod{Qqd}$ give all the splitting clauses by the
actual Artin map.  The level set is exactly
$\{g_0,g_1,p,p',q_0\}$.

The table supplies every conjunct of \texttt{Host\allowbreak Good}: its size and
arithmetic records, both second-set inequalities, \texttt{QDq},
\texttt{Box\allowbreak W}, the logarithmic derivative, \texttt{Thin\allowbreak OK},
\texttt{Bad\allowbreak OK}, disjoint $BB$ sets, integral norm elements, the two
congruence primes, the extra prime and the precise level set.
It also supplies the separate shares conclusion.  This closes the
literal \texttt{host\allowbreak Of\allowbreak Field\allowbreak Of} implication with no remaining analytic
input beyond the matched classical leaves used by QP and HX1.
The supplied \texttt{Pair\allowbreak Good}, field, and rank inequalities remain
hypotheses of that implication.

\paragraph{HX4: install the full target request on the cohomological window.}
Assume $\mathrm{VWinAt}(\psi,401)$.  AH18 chooses integers
$c\ge1,e_0$ before the dyadic page and pair.  QP22 and QP25, applied
to that $c$ and the given window hypothesis, choose $J\ge1$ and
arbitrarily late intervals $[Z,Z^{3/2}]$ on which the literal
\texttt{Win\allowbreak Datum} exists.  Its construction uses a single $Y,q,d$
through each such interval, although the dyadic page $e=e(x)$ may
vary.  It has
\[
 \begin{gathered}
 x^{1/J}\le Y\le x,\qquad \mathrm{PairGood}(M,\varepsilon,Y,x,q,d),\\
 x^{2+8\varepsilon}\le c2^e\le2x^{2+8\varepsilon},\\
 \mathrm{VVle}\bigl(E,\lfloor x^{400}\rfloor,
                  \sqrt{\exp(x^\varepsilon)}\bigr)
 \quad\hbox{for every }\mathrm{IsChiField}(\psi,e,q,d,E).
 \end{gathered}
 \tag{HX.8}
\]
For precision, the passage from (HX.4) to (HX.8) uses
$\kappa=c^{1/(2+8\varepsilon)}$, $Z=\kappa W$ and
$x=\kappa x'$.  Taking
$W\ge\max(\kappa^3,\kappa^{400})$ gives
\[
 [Z,Z^{3/2}]\subseteq[\kappa W,W^2],\qquad
 W\le x'\le W^2,\qquad x^{400}\le(x')^{401}.
 \tag{HX.9}
\]
Also $x'\le x$, so (HX.2) preserves the bound
$\sqrt{\exp(x^\varepsilon)}$.  The exponent $400$ in (HX.8) is
therefore intentional; the later proof does not silently restore
$401$ after this rescaling.

Put
\[
 C_0=J+o+8M,\qquad K=5+4M,\qquad A=K^{5/\varepsilon}.
 \tag{HX.10}
\]
All are fixed before the scale.  Enlarge $x_0$ past $A^J$, the
rank and host thresholds, and $c2^{e_0}$.  Then
$Y\ge x^{1/J}\ge A$ gives $Y^{\varepsilon/5}\ge5+4M$.
The lower bound in (HX.8) forces $e\ge e_0$: otherwise
$c2^e\le c2^{e_0}<x\le x^{2+8\varepsilon}$.

The \texttt{Pair\allowbreak Good} congruences give the exact \texttt{Host\allowbreak Hyp},
$q\ne3$, coprimality with $M$, and the splitting conditions for
every prime at most $\lfloor Y^{\varepsilon/5}\rfloor$.
AH18 therefore supplies an actual frame $\phi,H_0$ with
\[
 \begin{gathered}
 |\phi|=c2^e,\quad\phi.o=o,\quad H_0.q=q,\quad H_0.d=d,\\
 M\mid Q\mid\operatorname{lcm}(2^{e+2},M),\qquad
 Q\le4M|\phi|,
 \end{gathered}
 \tag{HX.11}
\]
and all of its original \texttt{Field\allowbreak Rows}, kept-character,
exceptional-character, kernel and twist conclusions.  HX2 gives
the ranks over this very frame.  The rank window and (HX.11)
give $Q\le C_0R.n$ (the source uses the available factor $8M$),
while $J,o,M\le C_0$.  HX3 applies to these data, returning
$H.\mathrm{toHostAx}=H_0$ and the required $\mathrm{HostGood}$
and shares statements, with constants independent of this scale.

It remains to verify the two conjuncts of the target request;
neither is a new existence theorem.  For \texttt{EAt\allowbreak H}, its first
three fields are the equality $\phi.o=o$, kept-character statement
and exceptional-character statement from AH18.  Its shares field
is HX3.  For the last two fields we give the exact
\texttt{bounded\allowbreak Tw} argument, extending GH4's finite-character
identification.  Let $\theta$ be in the source's bounded set
$(X^+\setminus\mathrm{Kept})\cup X^-$, and let $\chi$ be any
of its primitive Dirichlet representatives.  Both branches are
trivial on the first frame factor, so AH18's \texttt{Twist\allowbreak Form}
gives a power $\xi$ of $\psi$ and $i,j\in\{0,1\}$ with
\[
 H.\mathrm{dirAt}(\theta,a)=\xi(a)(a/q)^i(a/d)^j
       \quad((a,Qqd)=1),\qquad
 \theta(\iota_d)=(-1)^j.
 \tag{HX.12}
\]
The pair $(i,j)$ cannot be $(0,0)$.  In the $X^+$ branch, that
would make the primitive representative of $\xi$ one of the
kept characters and identify it with $\theta$, using $M\mid Q$;
this contradicts the definition of that branch.  In the $X^-$
branch, the displayed sign would be $+1$ instead of $-1$.

Every prime dividing $Q$ divides $2M$, by (HX.11).
Consequently an integer coprime to $2Mqd$ is coprime to $Qqd$.
The identity (HX.12), together with the defining identity of the
primitive representative $\chi$, therefore says exactly that
$\chi$ induces the displayed twist at the source level $2Mqd$.
This verifies the full \texttt{An.\allowbreak Induces} premise of
\texttt{Pair\allowbreak Good.\allowbreak bounds}, not just an equality away from an
unspecified finite set of primes.  Apply that bound at $y=x$,
using $Y\le x$, to get
\[
 |G_\chi(x)|\le x^{1/2+\varepsilon/2}
                   \le x^{1/2+\varepsilon},\qquad
 |L'/L(1,\chi)|\le4\log x.
 \tag{HX.13}
\]
These are precisely the two remaining \texttt{EAt\allowbreak H} fields,
uniformly for every such primitive representative.

For cohomology, the same kernel identity identifies the actual
$H.E$ as $\mathrm{IsChiField}(\psi,e,q,d,H.E)$, so the last line of
(HX.8) applies to this field.  HC7 gives $R.t\le x^{400}$ after a
single further threshold; since $R.t$ is an integer,
$R.t\le\lfloor x^{400}\rfloor$.  Finally
$s.b=\exp(x^\varepsilon)$ by the scale definition.  Apply (HX.2)
to obtain exactly $\mathrm{VVle}(H.E,R.t,\sqrt{s.b})$.

All clauses of (HX.5) now hold.  Taking as scale set the union of
the \texttt{Win\allowbreak Datum} intervals proves the literal
$\mathrm{HostExW}(3/2,(\mathrm{GChi}\,\psi).\mathrm{req})$.
The constants may depend on the fixed character and $\varepsilon$;
they precede $x$, the page, the rank record, and the host.  The same
construction applies to $\psi=1$, with $M=o=c=1$; it does not
remove the cohomological premise (HX.4) from the RH specialization.

\paragraph{What is removed from the remaining dependency list.}
The actual target host is now an explicit consequence of the window
estimate, with the matching of all previous host inputs written out.
It is not an independent deep supplier.  In a connected proof graph
its incoming arithmetic edge is the complete quantified
\texttt{VWin\allowbreak At} estimate, together with the classical leaves of
the established prime-pair/host constructions.  Constructing the
compatible theta/form/lattice family above that host is a separate
task, represented by \texttt{Design\allowbreak Exists}; HX4 does not prove it.


\section{The actual lattice-to-metric transport is a complete elementary chain}
\label{sec:metric-transport-complete}

\status{MT1--MT7 prove the literal implication from
\texttt{Link\allowbreak DIV}, fullness, positive retained rank, \texttt{Tu\allowbreak Size} and
\texttt{DGam\allowbreak Size} to \texttt{Link\allowbreak Mu}.  Every lattice, basis, metric,
coefficient field and finite place is the one stored in the original
\texttt{Fam\allowbreak P}.  In particular, \texttt{Nprime\allowbreak Def} is a field of that
record, not an extra arithmetic hypothesis.  The actual family and the
other four inputs still have to be constructed.  This closes the separate
metric-transport leaf in the previous global audit.}

The useful simplification is that this comparison is an ideal-index
calculation followed by a rounding estimate.  It does not require a new
comparison theorem between independently constructed Hermitian spaces.
The source already keeps the two lattices inside the same metrized vector
space.  We give the entire argument, including the exact factor $2$ and
the original uniformity in the family index and word.

\subsection{The exact objects and the ideal-index identity}

% BEGIN MT1
\paragraph{MT1: the source really uses one fixed metric.}
\label{claim:MT1}
Fix a late member $(D,L)$ of an original family $\mathcal F$.  Write
$K=D.\texttt{what}_2.\texttt{LL}$, $d_K=[K:\mathbb Q]$, and let $u$
be a word.  The Hermitian-lattice record $D.\texttt{div.\allowbreak Nu}(u)$ consists
of a $K$-space $V_u$, a basis $b_u$, matrices $X_{u,v}$ at its infinite
places, and an $\mathcal O_K$-submodule $N_u$.  The comparison object is
\[
 L.\texttt{Nu\allowbreak H}(u)=
       \{D.\texttt{div.\allowbreak Nu}(u)\ \text{with lattice }M_u\},
 \qquad M_u=L.\texttt{Nu\allowbreak Lat}(u).                         \tag{MT.1}
\]
Thus its space, basis and all infinite-place matrices are literally those
of $N_u$.  The matrices are invertible by the original
$D.\texttt{div\_\allowbreak good}$ field.

For the actual finite-place labels $w$, put
\[
 \ell_w=\frac{\log N\mathfrak p_w}{d_K},\qquad
 e_w=e(\mathfrak p_w/p),\qquad s_{u,w}=L.\texttt{sg}(u,w),
 \qquad q_{u,w}=\lfloor e_ws_{u,w}\rfloor.
\]
Here $\mathfrak p_w=D.\texttt{pl\allowbreak Eq}(w)$ and $p$ is its actual residue
characteristic; $D.\texttt{ram\_\allowbreak eq}$ identifies the stored ramification
index with $e_w$.  Define the nonzero fractional ideal
\[
       I_u=\prod_{w\in L.\texttt{supp}}\mathfrak p_w^{-q_{u,w}}.
                                                                  \tag{MT.2}
\]
The source's \texttt{Nprime\allowbreak Def} is exactly
\[
                N_u=I_uM_u,                                      \tag{MT.3}
\]
where the product denotes the $\mathcal O_K$-span of products.
For $x>\mathcal F.x_1$, (MT.3) is supplied by the field
$\mathcal F.\texttt{nprime\allowbreak Def}(x,i)$ for every $i\in\mathcal F.I_x$.
No choice of a new lattice or an isometry is involved.
% END MT1

% BEGIN MT2
\paragraph{MT2: multiplication by an ideal changes the slope exactly.}
\label{claim:MT2}
Let $V$ be a finite-dimensional $K$-space with the fixed source metric,
$M\subset V$ a nonzero finitely generated $\mathcal O_K$-module, and
$I\ne0$ a fractional ideal.  If $r=\dim_K KM$, then
\[
       \boxed{\quad\mu(IM)=\mu(M)-\frac{2\log N(I)}{d_K}.\quad}     \tag{MT.4}
\]
This is valid for lattices in a proper subspace of $V$ as well as for
full lattices.

First let $I=\mathfrak a$ be an integral ideal.  Locally at each finite
prime, the torsion-free module $M$ is free of rank $r$ over the discrete
valuation ring.  Consequently the $\mathfrak p$-length of
$M/\mathfrak a M$ is $r\,v_{\mathfrak p}(\mathfrak a)$, and
\[
                  [M:\mathfrak a M]=N(\mathfrak a)^r.             \tag{MT.5}
\]
One may see this without a global free basis: decompose the finite
quotient into its prime-primary parts and use local bases only to count
each part.  The archimedean realization is injective, and
$\mathfrak a M$ and $M$ span the same real space, of dimension $d_Kr$.
A fundamental region for the smaller lattice consists of
$[M:\mathfrak a M]$ translates of one for the larger lattice.  Hence
its covolume is multiplied by $N(\mathfrak a)^r$.

The source normalization is
\[
 \mu(M)=-\frac{2}{d_Kr}\log\operatorname{covol}(M)
              -\log2+\frac{\log|\operatorname{disc}(K)|}{d_K}.
                                                                  \tag{MT.6}
\]
The field-dependent constant cancels, and (MT.5) gives (MT.4).
For a fractional $I$, choose $0\ne a\in\mathcal O_K$ with
$\mathfrak a=aI$ integral.  Then $a(IM)=\mathfrak a M$.
Apply the integral formula to $IM$ and $(a)$, and to $M$ and
$\mathfrak a$, and subtract.  Since
$N(I)=N(\mathfrak a)/N((a))$, this gives exactly (MT.4).

The hypotheses required in the source are slightly asymmetric: it
assumes that $N=IM$ is nonzero and finitely generated.  These imply the
same properties for $M$.  Indeed, choose $0\ne a\in I$.  Multiplication
by $a$ injects $M$ into the Noetherian module $N$, so $M$ is finitely
generated; $M=0$ would force $N=0$.  Conversely multiplication by a
nonzero fractional ideal preserves the $K$-span.  Thus the preceding
proof applies with exactly the source's stated hypotheses.
% END MT2

% BEGIN MT3
\paragraph{MT3: the entire metric discrepancy is a floor error.}
\label{claim:MT3}
Apply MT2 to (MT.3).  Because
$\log N(I_u)=-\sum_{w\in\mathrm{supp}}q_{u,w}\log N\mathfrak p_w$,
\[
 \mu(N_u)-\mu(M_u)
      =2\sum_{w\in\mathrm{supp}}\ell_w\lfloor e_ws_{u,w}\rfloor.
                                                                  \tag{MT.7}
\]
The literal local size term is
$\texttt{sz}(w,s)=2\ell_we_ws$.  It follows that
\[
 \begin{split}
 R_u&:=\mu(N_u)-\mu(M_u)
              -\sum_{w\in\mathrm{supp}}\texttt{sz}(w,s_{u,w})\\
    &=2\sum_{w\in\mathrm{supp}}\ell_w
          \bigl(\lfloor e_ws_{u,w}\rfloor-e_ws_{u,w}\bigr),
 \end{split}                                                     \tag{MT.8}
\]
and therefore
\[
      |R_u|\le2\sum_{\substack{w\in\mathrm{supp}\\s_{u,w}\ne0}}
                          \ell_w.                                \tag{MT.9}
\]
Indeed $\ell_w\ge0$, and each floor error has absolute value at most
one; it is zero if $s_{u,w}=0$.  The formula uses the ordinary integer
floor even when $s_{u,w}<0$.  The sign of $s_{u,w}$ is not a hypothesis.
In fact $R_u\le0$, although only the absolute bound is needed.
% END MT3

\subsection{The arithmetic support of the rounding error}

% BEGIN MT4
\paragraph{MT4: normalized prime-place sums and the exact support bound.}
\label{claim:MT4}
For any rational prime $p$, the degree formula
$\sum_{w\mid p}e_wf_w=d_K$ and $N\mathfrak p_w=p^{f_w}$ give
\[
                  \sum_{w\mid p}e_w\ell_w=\log p.                 \tag{MT.10}
\]
Thus for any subset of its finite places,
$\sum_{w\in S,\,w\mid p}\ell_w\le\log p$, since $e_w\ge1$.
This calculation does not assume that $K/\mathbb Q$ is Galois.  The
source's bijection $\texttt{pl\allowbreak Eq}$ and residue-characteristic identity
make (MT.10) apply to its actual labels, without duplicates.

Put
\[
 \begin{gathered}
 X=L.\texttt{sc}.X''(D.R.n),\quad
 T_u=(D.\texttt{mem}(u)).A.T,\\
 \Delta_u=(D.\texttt{mem}(u)).C.\texttt{DGam},\quad
 g_0=D.\texttt{host}.g(0),\quad g_1=D.\texttt{host}.g(1),
 \quad d=D.\texttt{host}.d.
 \end{gathered}
\]
The original \texttt{Zero\allowbreak Above} assertion says that above $X$, away
from $T_u,g_0,g_1,d$ and prime divisors of $\Delta_u$, both the gain
and the $E$-term vanish.  The $F$-term vanishes there by the field
$L.\texttt{F\_\allowbreak above}$.  Since
\[
 s_{u,w}=\sum_{j=1}^{m}
       \bigl(\texttt{gain\allowbreak D}(j,w)+F(j,w)-E(u,j,w)\bigr),
\]
we obtain
\[
 s_{u,w}\ne0\Longrightarrow
 p\le X\ \text{or}\ p\in\{T_u,g_0,g_1,d\}
                      \ \text{or}\ p\mid\Delta_u.                \tag{MT.11}
\]
The removal of dropped primes in \texttt{gain\allowbreak D} preserves the zero
gain implication.  Notice that $\Delta_u$ is not presumed prime.

For completeness, the elementary estimate
\[
       \sum_{p\le Y}\log p\le Y\log4\qquad(Y\ge0)                \tag{MT.12}
\]
can be proved by induction on $n=\lfloor Y\rfloor$.  Let
$P(n)=\prod_{p\le n}p$.  For odd $n=2a+1\ge3$, all primes in
$(a+1,2a+1]$ divide $\binom{2a+1}{a}$.  This binomial coefficient is
at most $2^{2a}=4^a$, because the two equal middle coefficients are
part of a row summing to $2^{2a+1}$.  Hence
$P(2a+1)\le P(a+1)4^a\le4^{2a+1}$.  For even $n\ge4$,
$n$ is composite, so $P(n)=P(n-1)\le4^{n-1}$; the cases $n\le2$
are immediate.  Taking logarithms proves (MT.12).  Also
$\prod_{p\mid N}p\le N$ for every positive integer $N$, so
$\sum_{p\mid N}\log p\le\log N$.

Group (MT.9) by rational residue characteristic, apply (MT.10), and
then (MT.11)--(MT.12).  Nonnegative terms allow overlaps in the union
of prime sets to be counted more than once.  This gives the full
pointwise estimate
\[
 \boxed{\quad
 |R_u|\le2\{X\log4+\log T_u+\log g_0+\log g_1+
                                      \log d+\log\Delta_u\}.\quad} \tag{MT.13}
\]
This is the source's \texttt{nprime\allowbreak Mu\_\allowbreak at}, including every exceptional
prime family and every normalization.
% END MT4

% BEGIN MT5
\paragraph{MT5: all support costs fit the original $X''$ scale.}
\label{claim:MT5}
Write $x=L.\texttt{sc}.x>1$, $\epsilon=L.\texttt{sc}.\texttt{eps}$,
$\theta=\epsilon/100$, $n=D.R.n$, and $N=D.R.N(m)$.  Recall
\[
 X'=\exp(x^{\epsilon/3}),\qquad X=X'xn,\qquad f=2n,
 \qquad m\le2n,\qquad 0<\epsilon<1/52.                            \tag{MT.14}
\]
The identity $X=X'xn$ follows directly from
$b'=a/X'$, $w''=b'/(xn)$ and $X''=a/w''$.  In particular $x\le X$.

The literal \texttt{Size\allowbreak G} and \texttt{Size\allowbreak QD} clauses yield
\[
 \log g_i\le x^{\theta+\epsilon}\le x\le X,
 \qquad
 \log d\le x^{\epsilon/5+\theta}\le x\le X.                       \tag{MT.15}
\]
\texttt{Size\allowbreak Tu} yields
\[
               \log T_u\le mx^{\theta+\epsilon}
                               \le2nx\le2X.                       \tag{MT.16}
\]
These exponent inequalities hold throughout the original permitted
$\epsilon$ range.

Here is the complete numerical argument for the remaining
$\Delta_u$ term.  The last \texttt{Size\allowbreak Nodes} clause gives
$Nm^{10}\le\eta/m$; since $m\ge2$, it implies $N\le\eta$.
\texttt{Size\allowbreak Eta} gives $\log\eta\le x^{\epsilon/4}$, so
\[
                   N\le\eta\le e^{x^{\epsilon/4}}\le X'.         \tag{MT.17}
\]
\texttt{Size\allowbreak N} gives $n\le x^{2+8\epsilon+\theta}\le x^3$.
Thus $\log n\le3\log x$, while $\log N\le x^{\epsilon/4}\le x$.
Using $\log x\le x-1$ and $\log2<1$ gives
\[
                  0\le\log(fN)=\log(2nN)\le4x.
\]
Together with (MT.17) this proves
\[
                    0\le fN\log(fN)\le8X\le10X.                  \tag{MT.18}
\]
If \texttt{Fact\allowbreak DGam} holds with constant $C_\Delta$, it says exactly
$\log\Delta_u\le C_\Delta fN\log(fN)$.  Hence
\[
                   \log\Delta_u\le10|C_\Delta|X.                 \tag{MT.19}
\]
The absolute value permits the constant in the original record to have
either sign.  Inserting (MT.15)--(MT.19) into (MT.13) proves
\[
                  |R_u|\le
          2\bigl(\log4+5+10|C_\Delta|\bigr)X.                    \tag{MT.20}
\]
More generally, if $\log T_u\le C_TX$ and
$\log\Delta_u\le C_\Gamma X$, the exact constant is
$2(\log4+3+C_T+C_\Gamma)$, as in \texttt{nprime\allowbreak Mu\allowbreak R\_\allowbreak at}.
% END MT5

\subsection{Uniform installation on the original family}

% BEGIN MT6
\paragraph{MT6: literal \texttt{Link\allowbreak Mu} from the original source inputs.}
\label{claim:MT6}
Let $\mathcal F:\texttt{Fam\allowbreak P}(\epsilon)$ with $\epsilon>0$.  Suppose
\[
 \texttt{Link\allowbreak DIV}(\mathcal F,C_{\rm abs},C_H),\qquad
 \texttt{Nu\allowbreak Full\allowbreak F}(\mathcal F),\qquad
 \texttt{J.\allowbreak Mu\allowbreak Rank}(\mathcal F),\qquad
 \texttt{Tu\allowbreak Size}(\mathcal F),\qquad
 \texttt{DGam\allowbreak Size}(\mathcal F).                                  \tag{MT.21}
\]
Then the literal \texttt{Link\allowbreak Mu} holds: there are one constant $C$
and one threshold $x_0$ such that for every real $x\ge x_0$, every
$i\in\mathcal F.I_x$ and every word $u$,
\[
 \left|\mu(N_u)-\mu(M_u)
       -\sum_{w\in\mathrm{supp}}\texttt{sz}(w,s_{u,w})\right|
                 \le C X''(D.R.n).                              \tag{MT.22}
\]

To prove this with precisely the recorded quantifiers, first extract the
threshold from each input in (MT.21).  \texttt{Link\allowbreak DIV} supplies
\texttt{Zero\allowbreak Above} for all late $x,i$.  \texttt{Nu\allowbreak Full\allowbreak F} supplies finite
generation of every $N_u$; \texttt{Mu\allowbreak Rank} gives its positive $K$-rank,
hence $N_u\ne0$.  The metric matrices remain invertible by
$D.\texttt{div\_\allowbreak good}$, and (MT.3) is already a field of
$\mathcal F$ once $x>\mathcal F.x_1$.

\texttt{Tu\allowbreak Size} initially bounds
$\log T_u\le\mathcal F.C_{Tu}\,m x^\epsilon$, with a family
constant $\mathcal F.C_{Tu}$.  Since $x^{\epsilon/100}\to\infty$,
one further threshold makes
$\mathcal F.C_{Tu}\le x^{\epsilon/100}$.  This gives the literal
\texttt{Size\allowbreak Tu} inequality, and hence (MT.16), uniformly in $i,u$.
\texttt{DGam\allowbreak Size} supplies one $C_\Delta$ before $x,i,u$ and the
literal \texttt{Fact\allowbreak DGam} bound.  Every size assertion used in MT5
belongs to $\mathcal F.\texttt{sizes}(x,i)$.

Take the maximum of these finitely many thresholds and
$\mathcal F.x_1+1$.  MT1--MT5 now apply simultaneously to every $i,u$.
Take $C=2(\log4+5+10|C_\Delta|)$.  This proves (MT.22), with no
finiteness assumption on the union of all index sets or on the family
of words as $x$ varies.  The thresholds were already uniform in those
variables before their maximum was taken.

If one starts instead with the source's \texttt{Mu\allowbreak Tu} and
\texttt{Mu\allowbreak DGam}, the same proof uses their supplied constants $C_T$
and $C_\Gamma$ directly and gives
$C=2(\log4+3+C_T+C_\Gamma)$.  This is exactly
\texttt{J.\allowbreak link\allowbreak Mu\_\allowbreak of\_\allowbreak link\allowbreak DIV}; the preceding two paragraphs also
supply its entire \texttt{Tu\allowbreak Size}/\texttt{DGam\allowbreak Size} installation.
% END MT6

% BEGIN MT7
\paragraph{MT7: a smaller complete global interface.}
\label{claim:MT7}
In the source's \texttt{top\allowbreak Des\allowbreak Win} theorem, \texttt{Design\allowbreak Exists}
chooses a single $\mathcal F,C_H$ after the finite threshold join.
It supplies \texttt{Chosen\allowbreak E}, \texttt{Tu\allowbreak Size} and \texttt{DGam\allowbreak Size}.
The closed rank/fullness input gives
\texttt{Nu\allowbreak Full\allowbreak F}$\wedge$\texttt{Mu\allowbreak Rank} on this same $\mathcal F$.
As in DA2, \texttt{Hyp\allowbreak V}, \texttt{Link\allowbreak N} and
\texttt{Small\allowbreak Level\allowbreak Rank\allowbreak Eq} give \texttt{DIVin\allowbreak N}; then
\texttt{Link\allowbreak Eta} and \texttt{Link\allowbreak DIVt} give \texttt{Link\allowbreak DIV}.
Apply MT6.  The resulting \texttt{Link\allowbreak Mu} is on exactly the family
already carrying the other four \texttt{Global\allowbreak Theorem\allowbreak Inputs}.
Consequently DA2's metric-transport premise is fully discharged.

The sufficient global interface is therefore the original design
existence, the nine original closed link/rank conditions, and the
original class inclusions and \texttt{Hyp\allowbreak V}.  No independent
\texttt{METRIC\_\allowbreak TRANSPORT} leaf is needed.  This is a genuine reduction
in the human-proof frontier, although it does not construct
\texttt{Design\allowbreak Exists}, \texttt{Link\allowbreak DIVt}, or the rank/fullness input.
In particular, actual spectral/Gram bounds in \texttt{Link\allowbreak S} and
height/operator bounds in \texttt{Link\allowbreak DIV} keep their independent roles.
% END MT7


\section{The exact finite-orbit reduction of CountI}
\label{sec:count-i-finite-orbit-audit}

\status{KA1--KA4 reconstruct the full implication
\texttt{Prop\allowbreak GK}$\Rightarrow$\texttt{Cor\allowbreak GK}$\Rightarrow$\texttt{Count\allowbreak I}.
The arithmetic construction of \texttt{Prop\allowbreak GK} remains an input.
In particular, the assertions below do not supply \texttt{One\allowbreak Char},
the rank-three congruence-number theorem, or the automorphy and
rigidity hypotheses needed to construct them.  This isolates the
deep producer from the finite aggregation at the end of the CountI
branch.  Every coefficient ring and finite set of primes in the
original predicate remains arbitrary.  KA5--KA6 prove that Case III
is empty for targets of power-of-two order, including the principal
target, and carry that specialization through the actual height and
cohomological-window consumers.}

\subsection{Interfaces and quantifier order}

Use the original host notation of HV1: a fixed primitive character
$\chi$, conductor $q$, order $o$, exponent $s\ge2$, and auxiliary
primes $q_1,d$.  Write
\[
 N=\operatorname{lcm}(2^{s+2},q)q_1d,
 \qquad G_N=(\mathbb Z/N\mathbb Z)^\times.
 \tag{KA.1}
\]
The host subgroup $H_h\le G_N$ is cut out by the sign condition at
$2^{s+2}$, the kernel of $\chi$, and the square conditions at $q_1,d$.
For any source coefficient DVR $O$ at an odd prime $\ell$, let
$X_h(O,e)$ denote the characters trivial on $H_h$ and satisfying
$\psi(-1)=(-1)^e$.  Their original extended Selmer length and
normalized length are denoted by $L_{\ell,O}(\psi,e)$ and
$v_{\ell,O}(\psi,e)$, exactly as in (HV.3)--(HV.5).  Before
finiteness is proved, only the extended length is used.  The
normalized quantity in the source is always nonnegative.

For $J\le G_N$, put
\[
 \mathcal O_O(J)=\{\psi:G_N\longrightarrow O^\times:
                         \ker\psi=J\}.
 \tag{KA.2}
\]
This is the literal source \texttt{orb}.  It can be empty; no common
coefficient field containing all roots of unity is assumed.

For a number field $K$, write $D_K=|\operatorname{disc}K|$.
The predicate \texttt{GD} consists of the following conditions:
$\ell$ is odd; $e\ge1$; $N$ is even; $\ell\nmid\operatorname{ord}\psi$;
the $\ell$-primary CRT component of $\psi$ has square one;
$\psi(-1)=(-1)^e$; $K$ is imaginary quadratic;
$\ell\nmid D_Kh_K$; $(D_K,N)=1$;
\[
 \ker\overline\chi_{\mathrm{cyc},\ell}
   \cap\ker\operatorname{dir}(\chi_P(\ell,\psi))
       \not\subseteq G_K;
 \qquad \ell-1\nmid2e.
 \tag{KA.3}
\]
Here $\chi_P$ is the source character obtained by deleting the
$\ell$-power-order part and the local $\ell$-component.  In
particular its conductor divides $N$.

The exact useful form of \texttt{Cor\allowbreak GK} is this.  For every fixed
number field $K$ there are integers $c_5,A\ge0$, chosen before
$\Omega,e,N,J,p,p'$, the coefficient rings, and the finite prime
set $\mathcal L$, such that if $p\ne p'$ are primes not dividing
$30ND_K$ and for each $\ell\in\mathcal L$ at least one character in
$\mathcal O_{O_\ell}(J)$ satisfies \texttt{GD}, then all lengths
in these orbits are finite and
\[
 \sum_{\ell\in\mathcal L}\log\ell
       \sum_{\psi\in\mathcal O_{O_\ell}(J)}v_{\ell,O_\ell}(\psi,e)
 \le 2c_5\varphi([G_N:J])\,[N(e+1)]^A
                      (\log p+\log p').
 \tag{KA.4}
\]
Only the existence of one admissible orbit member at each prime is
required to apply this exact input; its conclusion covers the whole
orbit in that coefficient ring.

The stronger input \texttt{Prop\allowbreak GK} supplies, for every such
$K,e,N,J$ and every prime $p\nmid30ND_K$, a nonzero integer $b_p$,
chosen before every later $\ell\ne p$ and coefficient ring, with
\[
 \log|b_p|\le c_5\varphi([G_N:J])\,[N(e+1)]^A\log p,
 \qquad
 \sum_{\psi\in\mathcal O_{O_\ell}(J)}v_{\ell,O_\ell}(\psi,e)
       \le 2\operatorname{ord}_{\ell}(b_p),
 \tag{KA.5}
\]
and finiteness of those lengths whenever the same \texttt{GD}
antecedent holds.  The second assertion is conditional on
$\ell\ne p$; this exception is why two auxiliary primes occur in
(KA.4).  The source \texttt{Cor\allowbreak G}, with fixed exponents and a
different constant, is not substituted for \texttt{Cor\allowbreak GK}.

\paragraph{KA1: total mass of character-kernel orbits.}
\label{claim:KA1}
Let $G$ be a finite abelian group, and let $\mathcal T$ be any
finite set of subgroups $J\le G$ such that $G/J$ is cyclic.  Then
\[
       \sum_{J\in\mathcal T}\varphi([G:J])\le |G|.
 \tag{KA.6}
\]
If $J$ is the kernel of a character with values in the units of a
characteristic-zero domain, then $G/J$ is cyclic.

\begin{proof}
A finite subgroup of the multiplicative group of a field is cyclic.
One proof is to write its exponent as $m$.  If the group had more
than $m$ elements, the polynomial $X^m-1$ would have too many roots;
the structure theorem for finite abelian groups then shows that a
group with cardinality equal to its exponent is cyclic.  Apply this
to the image of the character, using the fraction field of the
coefficient domain.

For a cyclic quotient $G/J$ of order $m=[G:J]$, exactly
$\varphi(m)$ complex characters are faithful: a generator may be
sent to any primitive $m$th root of unity.  Composition with the
quotient map identifies these with the complex characters of $G$
whose kernel is exactly $J$.  For different $J$ these sets are
disjoint.  A finite abelian group has exactly $|G|$ complex
characters, as follows by decomposing $G$ into cyclic factors and
choosing the image of one generator in each factor.  Counting the
disjoint subsets gives (KA.6).  This argument uses complex
characters only as a counting device; it does not enlarge or
identify the original coefficient DVRs.
\end{proof}

\paragraph{KA2: the exact Case-I host supplies GD.}
\label{claim:KA2}
Fix $\chi$ and let $p_K=p_K(\chi)$, the source prime satisfying
$p_K\equiv3\pmod4$, $p_K\ge7$, and $p_K\nmid q$.
Let $K_\chi=K_*(p_K)$ be the source splitting-field model of
$X^2+p_K$, and let $h_*=h(K_\chi)$.
For every original host $h$ with character $\chi$ and
$\texttt{Big}(t)$, every odd prime $\ell$ in Case I, every
$1\le e<t$, every original coefficient ring $O$, and every
$\psi\in X_h(O,e)$, the datum
$(\Omega,K_\chi,\ell,e,\psi)$ satisfies the literal \texttt{GD}.
The statement is valid for every embedding of $K_\chi$ into the
given algebraic closure $\Omega$.

\begin{proof}
There are infinitely many primes congruent to $3$ modulo $4$:
if a finite list included all of them, a prime factor congruent to
$3$ modulo $4$ of four times their product minus one would be a
new one.  Including $3$ and the prime divisors of $q$ among the
excluded finite list supplies an eligible $p_K$.  Thus the
source minimum is well defined.  Since $-p_K$ is squarefree and
congruent to $1$ modulo $4$, the ring of integers of
$\mathbb Q(\sqrt{-p_K})$ has basis
$1,(1+\sqrt{-p_K})/2$ and discriminant $-p_K$.
The two roots $\pm\sqrt{-p_K}$ show that its splitting field is
the literal $K_*(p_K)$ up to a rational-field
isomorphism.  Hence $K_\chi$ is imaginary quadratic and
$D_{K_\chi}=p_K$.

The source exceptional set is
\[
 \mathcal B_\chi=\{3\}\cup
    \{\text{odd prime divisors of }oqp_Kh_*p_1\}.
 \tag{KA.7}
\]
Case I means $\ell\notin\mathcal B_\chi$ and
$\ell-1\nmid2e$.  Therefore
$\ell\nmid oqp_Kh_*$ and $\ell\ne p_K$.
The assumptions $\texttt{Big}(t)$ give $q_1,d>p_K$;
using $p_K\nmid q$ and (KA.1), we obtain $p_K\nmid N$.
These facts prove the two discriminant/class-number conditions
in \texttt{GD}.

The host quotient $G_N/H_h$ injects into the product of the
quotient of $(\mathbb Z/2^{s+2})^\times$ by $\langle-1\rangle$,
the image of $\chi$, and the two order-two square quotients.
Their orders are $2^s,o,2,2$, respectively.  Thus every
character trivial on $H_h$ has order dividing $2^{s+2}o$;
the odd prime $\ell\nmid o$ cannot divide its order.
Since $\ell\nmid q$, its CRT factor in the host quotient is
trivial unless $\ell=q_1$ or $\ell=d$.  In either exceptional
possibility the quotient is a square quotient, of exponent two.
Thus the local $\ell$-component of $\psi$ has square one.
Also $16\mid N$, and membership in $X_h(O,e)$ supplies the
required parity.  The oddness of $\ell$, positivity of $e$,
and nondivisibility $\ell-1\nmid2e$ are already hypotheses.

It remains to prove the noncontainment in (KA.3), which cannot be
replaced by the assertion that the two characters are nontrivial.
At a place above $p_K$, the quadratic extension $K_\chi/\mathbb Q$
is ramified.  For example, $X^2+p_K$ is Eisenstein over
$\mathbb Q_{p_K}$ and its root has valuation $1/2$; the extension
is totally ramified.  Consequently the absolute inertia group
at this place acts nontrivially on $K_\chi$.
The mod-$\ell$ cyclotomic character is unramified there because
$p_K\ne\ell$.  The Dirichlet character $\chi_P(\ell,\psi)$ is
also unramified there because its conductor divides $N$ and
$p_K\nmid N$.  This inertia group lies in both kernels in
(KA.3), while one of its elements is outside $G_{K_\chi}$.
This proves the noncontainment.  All assertions concern the
given embedded fields and their decomposition/inertia groups;
conjugating the chosen embedding preserves them.
\end{proof}

\paragraph{KA3: one integer at each of two primes gives CorGK.}
\label{claim:KA3}
The literal \texttt{Prop\allowbreak GK} implies the literal \texttt{Cor\allowbreak GK},
with the same constants $c_5,A$.

\begin{proof}
If $\mathcal L$ is empty, both the finiteness assertion and the
sum bound are immediate; the right side of (KA.4) is nonnegative.
Otherwise the assumed \texttt{GD} witness for one prime implies
$e\ge1$, so \texttt{Prop\allowbreak GK} applies.  Use it at $p$ and $p'$,
obtaining $b_p,b_{p'}\ne0$ before inspecting any $\ell$ or $O_\ell$.
For each $\ell\in\mathcal L$, at least one of $\ell\ne p$ and
$\ell\ne p'$ holds.  Applying the corresponding integer bound
proves finiteness and, by nonnegativity of both valuations,
\[
  \sum_{\psi\in\mathcal O_{O_\ell}(J)}v_{\ell,O_\ell}(\psi,e)
    \le2\operatorname{ord}_\ell(b_p)
                      +2\operatorname{ord}_\ell(b_{p'}).
 \tag{KA.8}
\]
For every nonzero integer $b$ and finite set of distinct primes,
\[
 \prod_{\ell\in\mathcal L}\ell^{\operatorname{ord}_\ell(b)}
       \mid |b|,
 \qquad
 \sum_{\ell\in\mathcal L}\operatorname{ord}_\ell(b)\log\ell
       \le\log|b|.
 \tag{KA.9}
\]
The divisibility follows from unique factorization and the
logarithmic inequality follows by taking logarithms of positive
integers.  Multiplying (KA.8) by $\log\ell\ge0$, summing, and
applying (KA.9) to both integers gives
$2\log|b_p|+2\log|b_{p'}|$.  The first inequalities of (KA.5)
are exactly (KA.4).  This also shows why the same two integers
must work for all later coefficient rings and primes: integers
chosen after $\ell$ could not be combined using (KA.9).
\end{proof}

\paragraph{KA4: complete original-host aggregation to CountI.}
\label{claim:KA4}
The literal \texttt{Cor\allowbreak GK} implies the literal \texttt{Count\allowbreak I}.
Hence \texttt{Prop\allowbreak GK} implies \texttt{Count\allowbreak I} as well.

\begin{proof}
Fix the primitive target $\chi$, and apply \texttt{Cor\allowbreak GK} to
the fixed field $K_\chi=K_*(p_K(\chi))$.  Obtain $c_5,A\in\mathbb N$.
Choose distinct primes
\[
       r_1>30qp_K,\qquad r_2>r_1.
 \tag{KA.10}
\]
Euclid's theorem supplies them.  These choices depend only on
$\chi$.  Put $B=2c_5(\log r_1+\log r_2)\ge0$ and choose
$k\in\mathbb N$ so that $B\le16^k$.  Finally put
\[
                    C=k+1+A.
 \tag{KA.11}
\]
This has made every required existential choice before the host,
the cutoff $t$, the twist $e$, the algebraic closure, the family of
coefficient rings, and the finite set of primes.

Now take any original host of character $\chi$ with
$q_1,d>r_1,r_2$, any $\texttt{Big}(t)$, any integer
$1\le e<t$, any algebraic closure $\Omega$, any coefficient
family $(O_\ell)_\ell$, and any finite set $\mathcal L$ of odd
Case-I primes.  Embed $K_\chi$ into $\Omega$; the following
estimates are independent of that choice by KA2.
Every prime dividing $N$ is either $2$, a divisor of $q$,
$q_1$, or $d$.  Since each $r_i$ exceeds $30qp_K$ but is smaller
than $q_1,d$, neither $r_i$ divides $30Np_K$.  Thus the same two
primes are admissible in (KA.4) for every such host.

Let $\mathcal T$ be the union, over $\ell\in\mathcal L$, of
the kernels of all characters in $X_h(O_\ell,e)$.  For
$J\in\mathcal T$, let
\[
 \mathcal L_J=\{\ell\in\mathcal L:
     \text{some }\psi\in X_h(O_\ell,e)\text{ has kernel }J\}.
 \tag{KA.12}
\]
For each $\ell\in\mathcal L_J$ choose such a character only to
verify the antecedent of \texttt{Cor\allowbreak GK}; KA2 supplies its
\texttt{GD}.  Applying (KA.4) to $J$ and $\mathcal L_J$ gives
finiteness for every character in its full coefficient-ring orbit
and the estimate
\[
 \sum_{\ell\in\mathcal L_J}\log\ell
       \sum_{\psi\in\mathcal O_{O_\ell}(J)}v_{\ell,O_\ell}(\psi,e)
       \le B\,\varphi([G_N:J])\,[N(e+1)]^A.
 \tag{KA.13}
\]
Every host character lies in precisely one kernel fiber, which is
a subset of the corresponding full orbit.  The quantities being
summed are nonnegative; therefore enlarging each fiber to that
orbit and interchanging two finite sums is valid.  It gives
\begin{align*}
 &\sum_{\ell\in\mathcal L}\log\ell
                   \sum_{\psi\in X_h(O_\ell,e)}v_{\ell,O_\ell}(\psi,e)\\
 &\quad\le\sum_{J\in\mathcal T}
       \sum_{\ell\in\mathcal L_J}\log\ell
                   \sum_{\psi\in\mathcal O_{O_\ell}(J)}v_{\ell,O_\ell}(\psi,e)\\
 &\quad\le B[N(e+1)]^A
                    \sum_{J\in\mathcal T}\varphi([G_N:J])
       \le BN[N(e+1)]^A.
 \tag{KA.14}
\end{align*}
The last step is KA1 and $|G_N|=\varphi(N)\le N$.
Notice that $\mathcal T$ is a set of subgroups of the one fixed
group $G_N$; no identification between the coefficient rings is
used.  Finiteness for all original host characters was already
obtained from (KA.13).

Finally $16\mid N$ gives $N\ge16$, and $e<t$ with integer
$e,t$ gives $e+1\le t$.  Hence
\[
 BN[N(e+1)]^A
 \le N^kN(Nt)^A
 \le (Nt)^{k+1+A}=(Nt)^C.
 \tag{KA.15}
\]
This is exactly the numerical conclusion of \texttt{Count\allowbreak I},
together with its finiteness conclusion and the required order
of all quantifiers.  KA3 gives the stated further implication
from \texttt{Prop\allowbreak GK}.
\end{proof}

\paragraph{KA5: Case III vanishes for every target of power-of-two order.}
\label{claim:KA5}
Suppose $\operatorname{ord}(\chi)$ is a power of two, including
$\operatorname{ord}(\chi)=1$.  Then, on every original host of
character $\chi$, there is no $\psi\in X_h(O,e)$ satisfying
\texttt{Case\allowbreak III}, for any coefficient ring, prime or twist.
Consequently the fixed-$\chi$ specialization of \texttt{Count\allowbreak III}
holds with exponent zero.

\begin{proof}
The quotient calculation in KA2 shows
\[
  \operatorname{ord}(\psi)\mid2^{s+2}\operatorname{ord}(\chi).
 \tag{KA.16}
\]
(The source \texttt{host\allowbreak Comp} proves the stronger exponent
$2^{s+1}\operatorname{ord}(\chi)$, which is not needed.)
If the right side is a power of two, so is the positive integer
$\operatorname{ord}(\psi)$.  The literal \texttt{Case\allowbreak III}
requires this order to be odd and requires $\psi\ne1$.
An odd divisor of a power of two is one, whereas a character of
order one is the trivial character.  This is a contradiction.
The numerical and finiteness assertions in the fixed-target
\texttt{Count\allowbreak Shape} are therefore implications with false
antecedent.  Choosing $C=0$ satisfies all remaining quantifiers.
No assertion about \texttt{Count\allowbreak III} for targets with nontrivial
odd-order part has been made.
\end{proof}

\paragraph{KA6: the specialized height/window cut.}
\label{claim:KA6}
For any fixed target $\chi$ of power-of-two order, the following
are sufficient for the original \texttt{Point\allowbreak HV} assertion at
that target:
\[
 \operatorname{CorGK}_{K_\chi},\qquad
 \texttt{PMEx\allowbreak Alt},\qquad \operatorname{CountV}_{\chi}.
 \tag{KA.17}
\]
Here the subscript means the exact original predicate with only
its outer field or target variable fixed, leaving every later
quantifier unchanged.  One can replace the first premise by
$\operatorname{PropGK}_{K_\chi}$.  With the supplied TW/OB/VW
inputs and their unchanged prime-pair and window hypotheses, the
same premises supply $\texttt{VWin\allowbreak At}(\chi,401)$.  Thus no
Case-III seed or promodularity theorem is needed on this selected
cohomological route to the principal target.

\begin{proof}
KA4 uses only \texttt{Cor\allowbreak GK} at the single field $K_\chi$, and
therefore supplies $\operatorname{CountI}_\chi$ from that
specialization.  YX2 supplies the literal \texttt{Count\allowbreak II} from
\texttt{PMEx\allowbreak Alt}, so in particular its $\chi$ specialization.
KA5 supplies $\operatorname{CountIII}_\chi$ with exponent zero.
The previously reconstructed \texttt{Count\allowbreak Z}, \texttt{Count\allowbreak Vs}
and \texttt{Count\allowbreak IV} remain available in their original scopes.

To check that fixed-target counts really suffice, inspect the
quantifiers of the actual consumers, not merely their universal
wrapper.  The proof of \texttt{sum\allowbreak Out} fixes $\chi$ and calls
\texttt{Count\allowbreak I}, \texttt{Count\allowbreak Z}, \texttt{Count\allowbreak II} and
\texttt{Count\allowbreak III} only at this same $\chi$; it then performs
the finite prime, twist and character sums of HV9.  The proof of
\texttt{sum\allowbreak In} likewise uses \texttt{Count\allowbreak II}, \texttt{Count\allowbreak V},
\texttt{Count\allowbreak Vs} and the layer count only at this $\chi$.
In \texttt{point}, after fixing $(\chi,j)$, all exceptional
constants are chosen from \texttt{sum\allowbreak In} at
$(\chi,\ell,j(\ell))$ and the outside constant is chosen from
\texttt{sum\allowbreak Out} at $\chi$.  Thus the proofs of HV8--HV10 run
unchanged with these specialized predicates.  They give precisely
the same field-model, coefficient, layer-function and uniform
host-threshold quantifiers of $\operatorname{PointHV}_\chi$.

Finally \texttt{thm\allowbreak HV} uses \texttt{Point\allowbreak HV} only at the supplied
target $\chi$, with
$j(\ell)=\lfloor\log_\ell R_\ell(L_\chi)\rfloor+3$ in the
original integer-logarithm convention.  The source \texttt{joint\allowbreak Win\_\allowbreak of}
calls this height estimate at the same target, and
\texttt{vwin\allowbreak At\_\allowbreak of\_\allowbreak joint} uses the joint-window hypothesis only
at that same target.  The already supplied two-primary and
odd-minus inputs are global and may of course be specialized.
Consequently the full field/window chain remains valid at
$\chi$, including the original real continuum of scales and
the cutoff exponent $400$ when $C_t=401$.  None of these
arguments provides the still-separate common arithmetic family,
gap, theta, denominator or metric/spectral branches required
by the final theorem.
\end{proof}

\paragraph{Why this does not remove cases I, II or V for RH.}
The principal target has $p_K=7$, $p_1=3$, and
$h(\mathbb Q(\sqrt{-7}))=1$, hence $\mathcal B_1=\{3,7\}$.
The class-number identity has the elementary reduced-form check:
if $[a,b,c]$ is a reduced primitive positive binary quadratic form
of discriminant $-7$, then $|b|\le a\le c$ implies
$3a^2\le7$.  Thus $a=1$, $b=\pm1$, and $c=2$; the usual reduced
boundary convention selects $[1,1,2]$.  The standard
ideal-class/binary-form correspondence therefore gives one class.
Thus Case V specializes to $\ell\in\{3,7\}$ and $e=1$;
its dyadic conductor is still unbounded.  For a concrete
nonempty instance take $q_1=11$, $d=23$, $t=2$ and host exponent
$s=12$.  A primitive even character modulo $2^{14}$, multiplied
by the Legendre character modulo $11$ and the trivial character
modulo $23$, is an allowed odd host character.  Choose an
original $3$-adic coefficient ring containing its roots of unity.
For $\ell=3$ its source invariants satisfy
\[
  f_0=2,\qquad n_0=3,\qquad c_2=14,\qquad
  R_1=16\cdot3^2\cdot4^3=9216<2^{14}.
 \tag{KA.18}
\]
It is therefore in the literal \texttt{Count\allowbreak V} case.  The host
is \texttt{Big}(2), since $5,7,3<11,23$ and $q=o=h_*=1$.
A primitive even character of conductor $2^{14}$ exists by
$(\mathbb Z/2^{14})^\times\cong C_2\times C_{2^{12}}$:
send $-1$ to $1$ and $5$ to a primitive $2^{12}$th root.

Case I is nonempty already with $e=1$, $\ell=5$ and the odd
Legendre character modulo $11$, since $5\notin\{3,7\}$ and
$4\nmid2$.  Case II is nonempty with $e=2$, $t=3$,
$\ell=3$, and the product of the Legendre characters modulo
$11$ and $23$; its order is two and its parity is even.
These examples only prevent an erroneous empty-case elimination;
they assert no cohomological lower bound.  In particular the
principal route still encounters characteristic three, so a
support theorem restricted to $\ell>3$ cannot by itself supply
its whole modular-support obligation.

\subsection{What this closes and what it leaves}

For a general target, the remaining arithmetic input on the
Case-I branch may now be named precisely:
\texttt{Prop\allowbreak GK} for the original coefficient and host interfaces.
An even narrower sufficient input for a fixed target is its
specialization to the single quadratic field $K_*(p_K(\chi))$;
the argument above never varies $K$ after choosing $\chi$.
For a proof of every Dirichlet target, the field specialization
must be supplied for every $\chi$.

The source constructs \texttt{Prop\allowbreak GK} from \texttt{One\allowbreak Char} and
\texttt{Congruence\allowbreak Number\allowbreak Rank\allowbreak Three}.  That producer controls one
nonzero integer simultaneously for all later coefficient primes
and contains the rank-two/rank-three automorphic comparison.
It is not a consequence of the finite orbit count proved here.
The present reduction also does not identify rank-four RX/CX
statements with rank-three Kato or host-frame statements.
For ordinary RH and, more generally, power-of-two-order targets,
KA5--KA6 remove the separate Case-III producer from the selected
cohomological route.  They do not remove it from unrestricted
Dirichlet GRH, and they do not assert that every other arithmetic
branch is complete.

\paragraph{Exact source record.}
All source references below are at commit
\texttt{7cbf7a0c32fca1badf05931879a37c3f8a2b03b4} of
\texttt{Coco\allowbreak Sgt/\allowbreak riemann-hypothesis-lean}.  KA1 follows
\texttt{Height\allowbreak And\allowbreak Vanishing/\allowbreak Orb\allowbreak Sum.\allowbreak lean:\allowbreak orb\allowbreak Sum} and
\texttt{Orb\allowbreak Facts.\allowbreak lean:\allowbreak orb\allowbreak Facts}.  KA2 reconstructs
\texttt{Gd\allowbreak Of\allowbreak Host.\allowbreak lean:\allowbreak gd\allowbreak Of\allowbreak Host}, with the actual host definitions
in \texttt{Dirichlet\allowbreak Prime\allowbreak To\allowbreak LPart\allowbreak Mul\allowbreak LPower\allowbreak Part.\allowbreak lean},
\texttt{Bchi\allowbreak Facts.\allowbreak lean:\allowbreak bchi\allowbreak Facts},
\texttt{Host\allowbreak Index.\allowbreak lean:\allowbreak host\allowbreak Index}, and
\texttt{Splitting\allowbreak Field\allowbreak Iso\allowbreak Quadratic.\allowbreak lean:\allowbreak kstar}.
KA3 reconstructs \texttt{Unitary\allowbreak Three\allowbreak Groups/\allowbreak Cor\allowbreak GSum\allowbreak Log.\allowbreak lean:\allowbreak cor\allowbreak GK}.
KA4 reconstructs \texttt{Unitary\allowbreak Three\allowbreak Groups/\allowbreak Count\allowbreak I.\allowbreak lean:\allowbreak count\allowbreak I}.
The exact predicates are \texttt{Unitary\allowbreak Three\allowbreak Groups/\allowbreak Kato\allowbreak Use.\allowbreak lean}
and \texttt{Height\allowbreak And\allowbreak Vanishing/\allowbreak Counts.\allowbreak lean}.
The surviving producer is
\texttt{Unitary\allowbreak Three\allowbreak Groups/\allowbreak Orb\allowbreak Equiv.\allowbreak lean:\allowbreak prop\allowbreak GK}, with the
actual \texttt{One\allowbreak Char} supplier under
\texttt{Modularity\allowbreak Argument/\allowbreak Kato\allowbreak Quadratic\allowbreak Characters/\allowbreak One\allowbreak Char.\allowbreak lean}.
The accompanying provenance file records exact line ranges,
file hashes, Git blob identifiers and links for each claim.
KA5 additionally uses \texttt{Host\allowbreak Comp.\allowbreak lean:\allowbreak host\allowbreak Comp} and the
literal \texttt{Host.\allowbreak Case\allowbreak III} definition.  KA6 is checked against
\texttt{Sum\allowbreak Out.\allowbreak lean:\allowbreak sum\allowbreak Out}, \texttt{Final\allowbreak Le.\allowbreak lean:\allowbreak sum\allowbreak In},
\texttt{Point\allowbreak Arith.\allowbreak lean:\allowbreak point}, \texttt{Thm\allowbreak HV.\allowbreak lean:\allowbreak thm\allowbreak HV},
\texttt{Prime\allowbreak Set\allowbreak Bound\allowbreak Data\allowbreak Of\allowbreak Primitive\allowbreak Char.\allowbreak lean:\allowbreak joint\allowbreak Win\_\allowbreak of}, and
\texttt{CMField\allowbreak Exponential\allowbreak Bounds/\allowbreak Admissible\allowbreak Of\allowbreak Page\allowbreak N.\allowbreak lean:\allowbreak vwin\allowbreak At\_\allowbreak of\_\allowbreak joint}.


\section{The complete numerical passage from the original tied family to a target window}
\label{sec:tied-window-completion}

\status{ZW1--ZW4 close the numerical and source-identification step between
the literal \texttt{Chain\allowbreak Tied} supplied by the global theorem and WB's
positive character window.  The proof includes the weighted slope
calculation and every error absorption; these are not inferred merely
from the name of a source theorem.  The eleven arithmetic ledger
conditions remain fields of the supplied \texttt{Chain\allowbreak Tied}.  No new
arithmetic existence assertion is added after that record.}

There are two different meanings of ``gap'' in a dependency discussion.
Producing the simultaneous geometric estimates is substantial arithmetic.
Once those estimates occur as the eleven fields of the literal tied
ledger, turning them into the asked-minus-main inequality is a finite
calculation.  We give that calculation here.  A single uniform constant $200R^3$ simplifies the error absorption
without changing the source's final exponent.

% BEGIN ZW1
\paragraph{ZW1: unpacking the original family without changing the host.}
\label{claim:ZW1}
Fix $\mathcal F:\texttt{Fam\allowbreak P}(\epsilon)$ and suppose
\[
 \texttt{Chain\allowbreak Tied}(\mathcal F,e_D,e_L,K,c_{HH},e_K,C_L,x).
 \tag{ZW.1}
\]
Its witnesses include one original index $i_0$, the frame
$\phi=(\mathcal F.D(x,i_0)).R$, a map
$\mathrm{pt}:\phi.A\to\mathcal F.I_x$, and the ledger $D$ with word
sets $W_u=(\mathcal F.D(x,\mathrm{pt}(u))).\mathrm{Word}$.
They also include the top set $T$, second set $T_2$, coefficients
$a_t,w_{u,j}$, finite cutoff $N$, actual logarithmic derivatives
$\ell$, and all of the following on these same objects:
\[
 \begin{gathered}
 D.x=x,\quad D.\epsilon=\epsilon,\quad
 C.e_{\rm Div}\le e_D,\quad C.e_{\rm Loss}\le e_L,
 \quad C.c_{HH}\le c_{HH},\\
 \texttt{Ident},\quad \texttt{Size\allowbreak Data}_2,\quad
 \texttt{Err\allowbreak Windows},\quad\texttt{BNBounds},\\
 \texttt{Growth\allowbreak Inequality\allowbreak Sizes},\quad\texttt{Thin\allowbreak Scales},\quad
 \texttt{Cost\allowbreak Windows},\quad\texttt{Rest\allowbreak Size}.
 \end{gathered}                                                    \tag{ZW.2}
\]
Here \texttt{Size\allowbreak Data}$_2$ means the eleven predicates in
\texttt{Formalism.\allowbreak LFunction\allowbreak Chain\allowbreak Parameters.\allowbreak LChain}, not the different
same-named datum in local-ring cardinality.  Its eleven fields are
\[
 \begin{gathered}
 \texttt{Hida\allowbreak Setting},\quad \texttt{DIV},\quad
 \texttt{EComp5},\quad \texttt{HComp},\\
 \texttt{Count5},\quad \texttt{Slopes},\quad
 \texttt{Quotient\allowbreak Slope\allowbreak Bounds\allowbreak With},\\
 \texttt{Loss\allowbreak Bounds},\quad \texttt{Slope\allowbreak Ineq},\quad
 \texttt{HHBound},\quad \texttt{FParams}.
 \end{gathered}
\]
These are premises of the calculation below, not additional consequences
asserted here.

The source \texttt{Tie\allowbreak Arith} identifies membership in $T$ with membership
of the integer label in $\mathcal F.TTn(x)$.  The family field
\texttt{TT\_\allowbreak eq} gives exactly the same characterization for the top set
of the original host $H=(\mathcal F.D(x,i_0)).\mathrm{host}$.
Injectivity of the frame's integer labels therefore gives
$T=H.TT$, as sets rather than just equal cardinalities.  Its other
clause says that every element of $T_2$ is a digit of an original
second-set entry; such digits lie in $\phi.A_0$.  Finally
$\texttt{arith\allowbreak Of}(H)=\texttt{arith}(H)$ by their definitions.

Take the translate set to be $\phi.A$ and its coverage map to be the
identity.  It covers $\phi.A_0$.  Reindexing \texttt{Ident} by this map
gives precisely \texttt{Ident\allowbreak U}; \texttt{BNBounds} becomes
\texttt{BNBounds\allowbreak U}; restricting the node bounds to $\phi.A_0$ gives
\texttt{Sizes\allowbreak A}.  The word sets are finite and nonempty because their
original distinguished words are already fields of the datum.
\texttt{Tie\allowbreak Arith}'s last clause is the exact
\texttt{Host\allowbreak Add\allowbreak Char\allowbreak Compatible} equality at every primitive character,
including its linear centering term.  Consequently (ZW.1) supplies
\texttt{Core\allowbreak At\allowbreak H} on this original $H$ and its original $TT$.
The fields \texttt{Tie\allowbreak Global} and \texttt{Tie\allowbreak Local} are not replaced by
new witnesses; the already supplied ledger is simply retained.
% END ZW1

% BEGIN ZW2
\paragraph{ZW2: the finite weighted slope calculation.}
\label{claim:ZW2}
Fix a translate and suppress it from notation.  The finite word set $W$
is nonempty, its weights $d_w$ are positive, and $\langle Z\rangle$
denotes $\sum_wd_wZ_w/\sum_wd_w$.  Put
\[
 A=\sum_ta_t,\quad n_T=\#T,\quad
 s=\langle\mu_{N,w}-h_w\rangle,\quad
 u=\max_w(\mu_{N,w}^{\max}-h_w)-s.
\]
Write $G=\mathrm{Ga}$, $F=\mathrm{Fo}$, $E_w=\mathrm{Ee}_w$,
$\mu_C=\mathrm{muCC}$, and $t_w=\log T_w$.  The inequality
$\mu_{N,w}\le\mu_{N,w}^{\max}$ shows $u\ge0$ by taking a
positive weighted average.

All coefficients below are the original components of $C$.  Define
\[
 \begin{aligned}
 B_1&=c_Sn_TaL_a(1+m/k),& B_2&=c_Sn_TaL_a,\\
 P&=c_\ell mB_0,&
 L_P&=c_Pn_Tmb\,x^{e_{\rm Loss}\epsilon},\\
 E_D&=c_Dn_TaL_a+c_Dm(B_0+N_c)x^{e_{\rm Div}\epsilon},\\
 V&=2mc_EB_0+c_H(n_TaL_a+mB_0)+Ac_C\log k,\\
 L_w&=c_\ell\mathrm{IO}
       +c_\ell n_Tm(a/k+b')t_w+c_\ell X''+1,\\
 M_L&=c_\ell L_P+c_\ell n_Tm(a/k+b')(c_CmL)+c_\ell X''+1,\\
 K_3&=B_1+B_2+V+P,\qquad U=K_3+M_L.
 \end{aligned}                                                     \tag{ZW.3}
\]
The subscripts $D,E,H,C,S,\ell,P$ abbreviate, in order,
\texttt{c\allowbreak Div,c\allowbreak E,c\allowbreak H,c\allowbreak Count,c\allowbreak Slope,c\allowbreak Sl1,c\allowbreak Loss}; $L_a=D.\texttt{La}$.
Thus $L_w,M_L,U$ are exactly the source's \texttt{Lo5},
\texttt{MLo5}, and \texttt{UBar5}.

The quotient and ordinary slope fields give
\[
 \begin{split}
 \mu_{N,w}&\ge-B_1+G+F-E_w-L_w,\\
 \mu_{N,w}^{\max}&\le B_2+G+F-E_w+P.
 \end{split}                                                       \tag{ZW.4}
\]
The exponent and height comparisons, followed by the variation clause of
\texttt{Count5}, give for any $v,w$
\[
 (-E_v-h_v)-(-E_w-h_w)\le V.                                      \tag{ZW.5}
\]
Indeed the exponent difference costs $2mc_EB_0$; the height difference
costs $c_H(n_TaL_a+mB_0)$ plus $A(t_v-t_w)$; and the last term is
at most $Ac_C\log k$.  Subtract the first inequality of (ZW.4) for
$w$ from the second for $v$, and average in $w$.  This gives
$\mu_{N,v}^{\max}-h_v-s\le K_3+\langle L_w\rangle$.
Taking the maximum in $v$ and using
$\mathrm{IO}\le L_P$, $t_w\le c_CmL$, and nonnegative coefficients,
we obtain
\[
             0\le u\le K_3+\langle L_w\rangle\le U.              \tag{ZW.6}
\]

Let
\[
 \Lambda=\sum_{j=1}^m\frac1f\sum_{\rho=1}^f
                 \log(w_{j\rho}v/\pi)+m\,\mathrm{lg}
                       -c_C L_a(1+m^2/k).
\]
The count field gives $t_w\ge\Lambda$.
Writing $S_D=\sum_t m\{\operatorname{ph}(a_t)-a_t\log\Omega\}$,
the divisibility inequality gives
$F-h_w\ge\mu_C-S_D+At_w-E_D$.
Combine this with (ZW.4), $A\ge0$, and
$E_w\le m(A\,\mathrm{lg}+c_EB_0)$, then average:
\[
 s\ge G+\mu_C-K_5-\langle L_w\rangle,\quad
 K_5=B_1+m(A\,\mathrm{lg}+c_EB_0)+S_D-A\Lambda+E_D.
\]
Expanding $\operatorname{ph}(y)=y\log y-y$, with the continuous
convention $0\log0=0$, and
$\log\Omega=\gamma+\log(\pi/v)-\mathrm{LpL}$, and distributing
the finite sums, gives exactly
\[
 S_D-A\Lambda+mA\,\mathrm{lg}
       =\mathrm{Main}+Ac_C L_a(1+m^2/k).
\]
Thus, with
\[
              K_0=Ac_C L_a(1+m^2/k)+E_D+mc_EB_0+B_1,
\]
we have $G\le\mathrm{Main}+K_0+\langle L_w\rangle+s-\mu_C$.

It remains to insert the original slope inequality.  Its hypotheses are
supplied by \texttt{FParams}, \texttt{Hida\allowbreak Setting} and
\texttt{HHBound}.  Write $\alpha=C.\texttt{c\allowbreak HH}$ for the ledger's
height exponent, distinct from its fixed global upper bound $c_{HH}$.
Put $H=\mathrm{HH}>0$ and $X=x^\alpha$.
For $u\ge0$, $H\le Xb$ and $b>0$,
\[
                u\log(1+H/u)\le(U+b)\log(1+X).                  \tag{ZW.7}
\]
For $u>0$, concavity of the logarithm with weights
$u/(u+b)$ and $b/(u+b)$ first bounds the left side by
$(u+b)\log(1+H/(u+b))$; use $H\le Xb$ and $u\le U$.
At $u=0$ the expression has value zero by continuous extension,
which also agrees with the source's totalized expression after
multiplication by $u$.

The parameters obey $0\le\delta\le\bar\delta<1$,
$0\le\eta\le\bar\eta$ and $0\le e\le\bar e\le1/4$.
The error terms in the slope inequality are nonnegative.  Hence we may
increase its reciprocal to $(1-\bar\delta)^{-1}$, replace its other
parameters by their upper bounds, and apply (ZW.7).
Finally $\mathrm{ASK}-G\le L_P$.  We have proved the full pointwise
estimate
\[
\begin{split}
 \mathrm{ASK}-\mathrm{Main}\le{}&K_0+\log|W|+L_P+M_L\\
 &+\frac{4\log(1+x^\alpha)(U+b)
       +(\bar\eta+4\sqrt{\bar e})x^\alpha b}
                  {1-\bar\delta}.
\end{split}                                                       \tag{ZW.8}
\]
The right side is exactly \texttt{Err\allowbreak X5at}.  This proves its numerical
implication from all eleven original ledger fields, for every translate
separately, without a new mean or a choice of a better word.
% END ZW2

% BEGIN ZW3
\paragraph{ZW3: uniform absorption of every error.}
\label{claim:ZW3}
Assume the original \texttt{Err\allowbreak Windows} and (ZW.2).
Fix $\epsilon>0$ and the global bounds $K,c_{HH},e_K,e_D,e_L$ before
the data.  Put
\[
 R=1+|K|+|c_{HH}|+|e_K|,\quad z=x^\epsilon,\quad
 E=n_Ta(1+m^2/k)+m(B_0+N_c).
\]
For all sufficiently large $x$, uniformly in every such ledger and
translate,
\[
                 \mathrm{ErrX5at}\le
                     x^{(e_D+e_L+3)\epsilon}E.                   \tag{ZW.9}
\]
We need only enlarge the threshold until
$x\ge3$, $\log x\le z$, and $200R^3\le z$.
These are eventual inequalities for each fixed positive $\epsilon$.

Here are the complete bounds.  \texttt{Err\allowbreak Windows} says
$0<\epsilon\le1$, $n_T\ge1$, $b\ge1$, $xb\le a$,
$xm\le k\le a$, $1\le k\le x^{e_K}$, $kb'\le a$, $X''\le a$,
$0\le L,L_a\le2z$, $B_0,N_c\ge0$, and $0\le a_t\le2a$.
It also gives $m\le x^{1/2+\epsilon/2}\le x$ and
$\log|W|\le(k-3)\log2\le a$.
In particular
\[
 E\ge n_Ta\ge a\ge b\ge1,\quad E\ge mB_0,\quad
 A\le2n_Ta,\quad m/k\le1,\quad mb\le a.                         \tag{ZW.10}
\]
All seven multiplicative constants lie in $[0,K]\subset[0,R]$;
$0\le e_{\rm Div}\le e_D$ and $0\le e_{\rm Loss}\le e_L$.
The following table follows by direct substitution into (ZW.3):
\[
\begin{array}{c|l}
\text{quantity}&\text{upper bound}\\ \hline
 L_P&R E z^{e_L}\\
 V&7R^2 E z\\
 K_3&14R^2 E z\\
 M_L&7R^2 E z^{e_L+1}\\
 U&21R^2 E z^{e_L+1}\\
 K_0&11R E z+R E z^{e_D}.
\end{array}                                                       \tag{ZW.11}
\]
For clarity, the only potentially large mixed term is bounded by
\[
 c_\ell n_Tm(a/k+b')(c_CmL)
              \le4R^2 n_Ta(m^2/k)z\le4R^2Ez.
\]
Also $\log k\le e_K\log x\le Rz$, which bounds the final term
of $V$ by $2R^2Ez$.  The other two terms of $V$ cost at most
$5REz$.  The bounds for $K_3$ and $M_L$ then follow by addition;
$c_\ell X''+1\le(R+1)E\le2RE$.
These calculations prove every entry in (ZW.11), without a constant
depending on the number of words or on a changing field.

The remaining original error-window clause is
\[
 (\bar\eta+4\sqrt{\bar e})x^{C.c_{HH}}b
                         \le Kz\,n_Ta\le REz.                  \tag{ZW.12}
\]
Moreover $\bar\delta\le1/2$, so its reciprocal is at most $2$,
and, using $\alpha=C.\texttt{c\allowbreak HH}\le c_{HH}$,
$\log(1+x^\alpha)\le\log2+c_{HH}\log x\le Rz$.
Using (ZW.11)--(ZW.12), the last line of (ZW.8) is at most
\[
 8Rz\{21R^2Ez^{e_L+1}+E\}+2REz
                         \le178R^3Ez^{e_L+2}.
\]
Its first line is at most
$11REz+REz^{e_D}+E+REz^{e_L}+7R^2Ez^{e_L+1}$.
Because $z\ge1$ and $e_D,e_L\ge0$, the sum of both lines is
at most
\[
 200R^3 E z^{e_D+e_L+2}\le E z^{e_D+e_L+3},
\]
as required.  If the fixed bounds admit no ledger, the implication is
vacuous; otherwise their nonnegativity assertions follow from the
displayed original fields.  The proof is uniform in the original
indices, word sets, coefficients and host.
% END ZW3

% BEGIN ZW4
\paragraph{ZW4: the literal tied-family terminal theorem.}
\label{claim:ZW4}
Fix a cyclic target level $\Lambda=\mathrm{Level.cyc}(\psi)$,
$\epsilon>0$, and one family $\mathcal F:\texttt{Fam\allowbreak P}(\epsilon)$.
Suppose $\texttt{EFam\allowbreak C}(\Lambda,\mathcal F.\texttt{to\allowbreak Fam\allowbreak Core})$ holds
and there are fixed $e_D,e_L,K,c_{HH},e_K,C_L,x_0$ for which (ZW.1)
holds whenever $x\ge x_0$ and $x\in\mathcal F.\mathrm{Win}$.
Then, on all sufficiently late members of that family's windows,
\[
                    H_\psi(x)\le
                         x^{1/2+(e_D+e_L+7)\epsilon}.             \tag{ZW.13}
\]
Consequently the literal \texttt{Family\allowbreak Class\allowbreak Connected} conclusion for
$(\mathrm{GChi}\psi).\mathrm{cls}(3/2)$ suffices for the target's
zero-free half-plane.  No separate arithmetic-gap existence theorem is
needed after that conclusion.

Indeed ZW1 gives the exact \texttt{Core\allowbreak At\allowbreak H} ledger.  Its
\texttt{Ident\allowbreak U} field and the top/place bijections give the two
identities (WB.26), with the original $a_t,w_{u,j},\ell$ and $N$.
ZW2--ZW3 therefore bound the gap at every translate in the covered
subgroup.  The two \texttt{BNBounds\allowbreak U} inequalities increase $E$ to
the bracket $B$ of (WB.22), term by term with nonnegative
coefficients.  \texttt{Rest\allowbreak Size} supplies its lower size bounds;
the two upper bounds $m\le x^{1/2+\epsilon/2}$ and $xm\le k$
come from \texttt{Err\allowbreak Windows}.  The remaining scalar upper bounds
$a\le2xb$ and $k\le n^2x^{1/2-\epsilon}$ are also fields of
\texttt{Rest\allowbreak Size}, while $\#T_2/n\le1$ follows from $T_2\subset\phi.A$.
Its positive lower bounds imply $m\ge1$ and $\#T_2>0$, and
\texttt{Err\allowbreak Windows} gives $n_T\ge1$.  Thus every scalar hypothesis
of WB6 holds, with the fixed gap exponent $C=e_D+e_L+3$.

The original \texttt{Host\allowbreak Ax} supplies WB1--WB2's reciprocity and
ramification facts.  Its \texttt{EAt\allowbreak H} supplies the same primitive
retained characters, biased top-set counts, omitted-character ramp
bound, and nonretained logarithmic-derivative bound.  The finite sum
of the retained $|L'/L(1,\chi)|$ depends only on $\Lambda$, so gives
the uniform $K_0$ in WB5.  Enlarge the exceptional-character
constant to its maximum with zero, which preserves its bound for
$x\ge3$ and supplies WB5's nonnegativity convention.
The full \texttt{Sizes\allowbreak A},
\texttt{Thin\allowbreak Scales}, \texttt{Cost\allowbreak Windows} and $T_2\subset A_0$
are exactly the remaining WB4 premises.  WB6 consequently yields
(ZW.13), with its additional threshold $x^\epsilon\ge1+4o$.

The definition of \texttt{EFam\allowbreak C} gives one host threshold and one
exceptional-character constant before $x$ and the family index.
The tied-family condition gives one ledger threshold and constants
before $x$.  Join these thresholds with ZW3's and WB's by a finite
maximum.  At each $x$, use the index $i_0(x)$ supplied by ZW1
for that $x$; no single index or host is fixed across the window.
The estimate holds at every real $x$ in the late part of each family
window, because every threshold precedes this index choice.

Finally fix any proposed zero with real part $\beta>1/2$ and choose
$0<\delta<\min(1/2,\beta-1/2)$.
\texttt{Family\allowbreak Class\allowbreak Connected} first fixes $e_D,e_L$ and an upper
range for $\epsilon$.  Choose a permitted $\epsilon>0$ so small
that $(e_D+e_L+7)\epsilon\le\delta$, and only then choose its
family and a window beyond the finite common threshold.  Since the
windows have fixed stretch $3/2>1$ and occur arbitrarily late,
WB8/AN3 contradict the proposed zero.  With
$e_D=2,e_L=0$ on the selected \texttt{top\allowbreak Des\allowbreak Win} route, the
pointwise exponent is exactly $1/2+9\epsilon$.

The arithmetic work retained here is precisely the production of the
literal common family and its links.  The numerical slope argument is
now complete even though its geometric slope premise
\texttt{Slope\allowbreak Ineq} still has to be installed by the upstream global
producer.  This distinction is reflected by separate implication and
producer statuses in the dependency audit.
% END ZW4


% v0.30 candidate component. Independent human reconstruction.
% Source pin: 7cbf7a0c32fca1badf05931879a37c3f8a2b03b4.
\section{An exact weaker analytic input on the selected route}
\label{sec:analytic-envelope-audit}

The pointwise window bound in AN1 is sufficient, but is stronger than the
actual analytic interface in the pinned source.  Its
\texttt{Window\allowbreak Hyp\allowbreak U} requires a nonnegative envelope with a small weighted
integral.  The following proof reconstructs that interface independently.
The arithmetic construction of such envelopes remains an obligation.

\subsection{An exact positive-part formulation}

Fix one primitive target $\psi$ and its fixed positive character family,
and write $H=H_\psi$ as in AN1 and $g(u)=H(e^u)$.
For a real number $v$, put $v_+=\max(v,0)$.  Fix
$0<\theta<1$ and a fixed $\delta>0$.  The integral window condition is
\[
 \begin{gathered}
 \text{for every $\eta>0$ and $Y_0$, there exists
 $Y\ge\max(2,Y_0)$ such that}\\
 \int_Y^{Y^{1+\delta}} H(x)_+\,x^{-\theta-1}\,dx\le Y^\eta.
 \end{gathered}\tag{AQ.1}
\]
In logarithmic coordinates it is
\[
 \forall\eta>0\ \forall U_0\ \exists U\ge\max(1,U_0),\qquad
 \int_U^{(1+\delta)U}g(u)_+ e^{-\theta u}\,du\le e^{\eta U}.
 \tag{AQ.2}
\]
The harmless lower thresholds $2$ and $1$ are interchangeable because
the starting threshold is universally quantified.

\begin{proposition}[AQ1: the literal envelope interface]
\label{claim:AQ1}
For this $g$, condition (AQ.2) is equivalent to the source's
\texttt{Window\allowbreak Hyp\allowbreak U}$(g,\theta,\delta)$.
\end{proposition}
\begin{proof}
The source asks for a function $R$ on each selected interval satisfying
$R\ge0$, $g\le R$, integrability of $R(u)e^{-\theta u}$, and the
integral bound in (AQ.2) with $R$ in place of $g_+$.  Such an $R$
dominates $g_+$.  Conversely choose $R=g_+$.
The finite ramp sum defining $H$ is continuous, and its crude bound
$|H(x)|\le Mx^2$ for $x\ge1$ makes $g_+$ integrable on every bounded
interval.  This proves both implications.  Substitution $x=e^u$
proves equivalence with (AQ.1).  Notice that a bound on the signed
integral of $H$ would not imply (AQ.1).
\end{proof}

\subsection{Localization on every fixed multiplicative stretch}

\begin{lemma}[AQ2: a quantitative gamma localization]
\label{claim:AQ2}
Let $\delta>0$, put
\[
 d=\min(\delta,1),\quad m=1+\frac d2,\quad
 h=\frac{8m}{d},\quad \sigma=(h+2)^2+3,\quad \kappa=m\sigma.
\]
For $U\ge1$ let $k=\lceil\kappa U\rceil$.  There is a constant
$C_\delta$ independent of $U$ such that
\[
 \int_{(0,\infty)\setminus[U,(1+d)U]}
       \frac{u^k}{k!}e^{-(\sigma-2)u}\,du
       \le C_\delta\sigma^{-k}e^{-mU}.
 \tag{AQ.3}
\]
\end{lemma}
\begin{proof}
On the lower tail insert $e^{h(U-u)}\ge1$ and on the upper tail
insert $e^{h(u-(1+d)U)}\ge1$.  Extending both integrals gives
\[
 \frac{e^{hU}}{(\sigma+h-2)^{k+1}}
 +\frac{e^{-h(1+d)U}}{(\sigma-h-2)^{k+1}}.       \tag{AQ.4}
\]
The denominators are positive.  Multiplication by $\sigma^k$ makes
the first summand at most a fixed constant times
\[
 \exp\{[h-\kappa\log(1+(h-2)/\sigma)]U\},
\]
using $k\ge\kappa U$.  The second is at most a fixed constant times
\[
 \exp\{[-h(1+d)-\kappa\log(1-(h+2)/\sigma)]U\},
\]
using $k\le\kappa U+1$.  For $v\ge0$,
$\log(1+v)\ge v-v^2/2$; for $0\le v\le1/2$,
$-\log(1-v)\le v+v^2$.  Our choice of $\sigma$ allows the latter
inequality at $v=(h+2)/\sigma$.  Since
$(m-1)h=(1+d-m)h=4m$, the two exponents are bounded respectively by
\[
 -2m+\frac{m(h-2)^2}{2\sigma}\le-m,
 \qquad
 -2m+\frac{m(h+2)^2}{\sigma}\le-m.
\]
This proves (AQ.3), including the effect of the ceiling in $k$.
\end{proof}

\begin{theorem}[AQ3: positive integral windows exclude right-hand zeros]
\label{claim:AQ3}
If (AQ.1) holds for a fixed positive character family and fixed
$0<\theta<1$, $\delta>0$, then every constituent Dirichlet
$L$-function is nonzero on $\operatorname{Re}s>\theta$, apart from
the principal pole at $1$.
\end{theorem}
\begin{proof}
The Mellin computations AN2--AN3 give the meromorphic continuation
of the Laplace transform
\[
 \Phi(s)=\int_0^\infty g(u)e^{-su}\,du
 \qquad(\operatorname{Re}s>2)
\]
to $\operatorname{Re}s>0$.  Its poles there are exactly the zeros of
the constituent functions, after combining coincident poles, and a
zero $\rho$ has residue
\[
 \frac{b_\rho}{\rho(1-\rho)},\qquad b_\rho>0.       \tag{AQ.5}
\]
The apparent pole at $1$ has been canceled by the centering; the
nonprincipal values at $1$ are nonzero by QW14.  These poles are
locally finite, lie to the left of $\operatorname{Re}s=\sigma$,
and the residue at every real pole in $(0,1)$ is positive.  The
coefficient $b_\rho$ is a positive weighted sum of multiplicities
for the characters and their conjugates; hence there is no
cancellation even when different constituent functions vanish at
the same point.  The elementary bound $|g(u)|\le M e^{2u}$ is
uniform for $u\ge0$, with $M$ depending on the fixed family.

Use $\sigma,\kappa$ from AQ2 and put
\[
 A_k(z)=\frac{(-1)^k}{k!}\Phi^{(k)}(z)
       =\int_0^\infty g(u)\frac{u^k}{k!}e^{-zu}\,du.
\]
For every real $t$, reality of $g$ gives
\[
 |A_k(\sigma+it)|+A_k(\sigma)
 \le 2\int_0^\infty g(u)_+\frac{u^k}{k!}e^{-\sigma u}\,du.
 \tag{AQ.6}
\]
Restrict the good interval in (AQ.2) to $[U,(1+d)U]$; this only
decreases its nonnegative integral.  On this interval
\[
 \frac{u^k}{k!}e^{-(\sigma-\theta)u}
          \le(\sigma-\theta)^{-k},
\]
by the exponential series.  Outside it apply AQ2 and the crude
growth bound.  Therefore the selected $U$ and
$k=\lceil\kappa U\rceil$ satisfy
\[
 |A_k(\sigma+it)|+A_k(\sigma)
 \le2e^{\eta U}(\sigma-\theta)^{-k}+O(\sigma^{-k}).
 \tag{AQ.7}
\]
The implied constant is independent of the selected window and of
$t$.  Since $U\le k/\kappa$, the first term is at most
$2B^{-k}$ where $B=(\sigma-\theta)e^{-\eta/\kappa}$.

Suppose some pole lies to the right of $\theta$.  Set
\[
 D=\min\{\sigma-\theta,\ \inf_\rho|\sigma-\rho|\}.
\]
First suppose that some pole has horizontal distance from
$\operatorname{Re}s=\sigma$ strictly less than $D$.  Choose a real
$t$ for which $z=\sigma+it$ has a unique nearest pole $\rho$ at
distance $a<D$.  This is possible near the imaginary part of that
pole: only finitely many poles can compete inside the relevant
compact disks, and each pair with different imaginary parts
forbids only one $t$.  Distinct poles with the same imaginary part
have different distances because they lie to the left of $\sigma$.
Choose $a<r<D$ so that the closed $r$-disk about $z$ contains no
other pole.  The $r$-disk about $\sigma$ contains none.  Cauchy's
estimate, after subtracting the principal part at $\rho$, gives
\[
 A_k(z)=\frac{c}{(z-\rho)^{k+1}}+O(r^{-k}),\qquad
 A_k(\sigma)=O(r^{-k}),\quad c\ne0.
\]
Choose $\eta>0$ sufficiently small that
$B=(\sigma-\theta)e^{-\eta/\kappa}>r$.
This choice precedes the selection of arbitrarily late windows in
(AQ.2).  Inequality (AQ.7) now implies
\[
 |c|a^{-k-1}\le O(r^{-k})+2B^{-k},
\]
which is impossible as $U$, hence $k$, tends to infinity.

In the remaining case all poles have horizontal distance at least
$D$.  The existence of a pole to the right of $\theta$ forces
$D<\sigma-\theta$.  A nearest pole to $\sigma$ exists by local
finiteness in a compact disk contained in
$\operatorname{Re}s>\theta$.  It must be real, since its horizontal
distance is at least $D$ and its total distance equals $D$.
It is the unique pole at that distance, located at $\beta=\sigma-D$,
and its residue $c$ is positive.  Choose
$D<r<\sigma-\theta$ with no other pole in the closed $r$-disk.
Then
\[
 A_k(\sigma)=cD^{-k-1}+O(r^{-k}).
\]
Choose $\eta>0$ so that $B>r$ and use (AQ.7) at $t=0$.
For large $k$ its left side is $2A_k(\sigma)>0$, giving the same
impossible comparison of exponential rates.  Thus no such pole
exists in either case.  The Mellin residue identification (AQ.5)
proves the asserted nonvanishing.
\end{proof}

\begin{corollary}[AQ4: a minimal sufficient terminal contract]
\label{claim:AQ4}
For every fixed primitive $\psi$ and every
$\theta\in(1/2,1)$, suppose there is a $\delta_{\psi,\theta}>0$,
fixed before $\eta$ and the window, such that (AQ.1) holds.
It suffices to impose this for a sequence of such $\theta$
decreasing to $1/2$.  Then Dirichlet GRH follows.  The ordinary RH
case requires only $\psi=1$.
\end{corollary}
\begin{proof}
AQ3 excludes any proposed zero with real part $\beta>1/2$ by
choosing a member of the sequence below $\beta$.
The primitive functional equation reflects zeros in the open
strip on the left to right-hand zeros of the conjugate character.
For an imprimitive character the additional Euler factors
$1-\chi^*(p)p^{-s}$ are nonzero for $\operatorname{Re}s>0$.
These are the already written GH6 transitions.  This conclusion
concerns Dirichlet $L$-functions; it does not extend the scope to
arbitrary automorphic $L$-functions.
\end{proof}

\paragraph{What the contract saves, and what it does not supply.}
A pointwise bound $H(x)\le C_\theta x^\theta$ gives an integral
at most $C_\theta\delta\log Y$, which is at most $Y^\eta$ on
sufficiently late windows for every $\eta>0$.  Thus the previous
pointwise route is retained.  Conversely (AQ.1) permits short
positive spikes and does not imply a pointwise bound.  It is a
potentially useful alternate arithmetic target, not a new supplier
of the host, common family, theta data, arithmetic gap or operators.
The character family and $\delta$ are fixed before $\eta$ and the
window; constants or weights chosen anew with each $Y$ would not
satisfy the theorem as stated.

\subsection{A small discrete simplification}

\begin{lemma}[AQ5: interpolation of the original centered ramps]
\label{claim:AQ5}
For $n\ge1$ and $x\in[n,n+1]$, let $L_n(x)$ be the affine
interpolant of $H(n)$ and $H(n+1)$.  For the original target with
principal coefficient $1$,
\[
 0\le H(x)-L_n(x)
       \le\frac{(x-n)(n+1-x)}{2n}\le\frac1{8n}.
 \tag{AQ.8}
\]
Consequently for $0<\theta<1$ and $C\ge0$, bounds
$H(n)\le Cn^\theta$, $H(n+1)\le C(n+1)^\theta$ imply
$H(x)\le Cx^\theta+1/(8n)$ on that interval.
\end{lemma}
\begin{proof}
Every nonprincipal centered ramp is affine between consecutive
integers, as is the finite von Mangoldt part of the principal
ramp.  The only nonaffine term is $-x\log x$; therefore
$H''(x)=-1/x$ on $(n,n+1)$.  The difference $H-L_n$ is concave
with zero boundary values, giving its lower bound.  Subtract
$(x-n)(n+1-x)/(2n)$.  The resulting function has nonnegative
second derivative and zero boundary values, so it is at most
zero.  This proves (AQ.8).  Finally $x^\theta$ is concave, so its
affine interpolant is at most its value at $x$.
\end{proof}

In particular, if for a fixed $\lambda>1$ and each
$\theta\in(1/2,1)$ there are a fixed $C_\theta\ge0$ and arbitrarily
late windows with $H(n)\le C_\theta n^\theta$ for every integer
$n\in[Y,Y^\lambda]$, the same GRH conclusion follows.  Choose
$1<\mu<\lambda$ and $Z=\lceil Y\rceil$.  For sufficiently large
$Y$, $Z^\mu+1\le Y^\lambda$, so for every $x\in[Z,Z^\mu]$,
the two integers $\lfloor x\rfloor$ and $\lfloor x\rfloor+1$
lie in the supplied window.
AQ5 gives $H(x)\le(C_\theta+1)x^\theta$ throughout that smaller
window.  AQ3 applies with $\delta=\mu-1$.

\paragraph{Pinned-source scope.}
The source defines the original \texttt{Window\allowbreak Hyp} and its logarithmic
version \texttt{Window\allowbreak Hyp\allowbreak U}.  Its \texttt{landau\_\allowbreak window} proves the
logarithmic pole criterion and \texttt{lemma\allowbreak L} its character-family
zero-free consequence;
the source's power-window theorem derives the
pointwise special case for any fixed $\delta>0$.
The present gamma estimate is an independent human proof of the
same analytic implication, not a claim about a new arithmetic
window construction or a fresh Lean replay.


\section{Compatible design choices and the remaining existence obligation}
\label{sec:design-existence-v031}

\status{DEX1--DEX2 install the literal family thinning condition from the
already reconstructed centered-character estimates. DEX3 matches the
published Minkowski theorem, with its exact integer, and DEX4 proves the
required logarithmic bound. DEX5 proves the numerical proviso directly
from the original rank formula. DEX6--DEX11 reconstruct the common-threshold,
common-selector and actual-family joins. The final existence statement
remains conditional on the specified original recursive objects, rank
formula and level-prime size supplier. None of these joins asserts that
independently existing geometric or theta objects are compatible.}

Throughout, a family means the original \texttt{Fam\allowbreak P}, with its actual
fields, maps, word labels, lattices and metrics. All statements concern the
pinned source and the fixed target character used in the global audit.
The standard background of the preceding chapters is retained. In
particular, CC10--CC18 are used in their proved scopes, including the
classical zero-free-region input already recorded there.

\paragraph{DEX1: the conductor and odd-part bounds on one actual family.}
\label{claim:DEX1}
Let $D$ be an original datum and let $\chi$ satisfy
\texttt{Is\allowbreak Char\allowbreak ED}. Thus $\chi$ is primitive modulo $N$ and, on
integers prime to $M=Q_Sqd$, agrees with a character of the original
finite Artin group, pulled back by the actual Artin map. Then
\[
                         N\mid M.                         \tag{DEX.1}
\]
To see the precise conductor assertion, inflate both characters to
$\operatorname{lcm}(N,M)$. Every unit there has an integer representative
prime to both moduli, so the two inflations agree. The right-hand
character factors through $M$. By the prime-by-prime definition of the
conductor, each local conductor exponent of the primitive $\chi$ is at
most the corresponding exponent of $M$. This proves (DEX.1), including
$N=1$. Equivalently, the local kernel killed by reduction to $M$ is
killed by $\chi$; Chinese remainders separate these local assertions.

For an original family $F$, the late members satisfy
\[
 o\le C_o,\qquad Q_S\le C_o n,\qquad
 n\le o2^e,\qquad 2^e\mid Q_S .                           \tag{DEX.2}
\]
The first two are literal family fields and the last two are the
original frame and host fields. Write $Q_S=2^vQ_{S,\mathrm{odd}}$.
Since $v\ge e$, (DEX.2) gives
\[
 2^eQ_{S,\mathrm{odd}}\le Q_S
       \le C_on\le C_o^2 2^e,
 \qquad Q_{S,\mathrm{odd}}\le C_o^2.                      \tag{DEX.3}
\]
Here $C_o\ge0$ follows from $o\le C_o$ and $o\ge0$.
In particular $B=\lceil C_o^2\rceil$ bounds the odd part uniformly
over the late members. This does not bound the power of two in $Q_S$.

Finally, \texttt{Box\allowbreak E} at any supplied scale $Y$ implies the following
real-zero exclusion: every real primitive host character whose conductor
is divisible by $q$ or $d$ has no real zero $\beta>1-1/400$.
Indeed, if $d\mid N$, use the first \texttt{Good\allowbreak For} clause. Otherwise
$q\mid N$ and the second clause applies. In both cases the height of the
real argument is zero, which is at most the nonnegative box height.
No comparison between $Y$ and the large variable of the character sum is
needed for this real restriction.

\paragraph{DEX2: close \texttt{Chosen\allowbreak Thin}, with coefficient one.}
\label{claim:DEX2}
For every original family $F$ and $\epsilon>0$, we prove
\[
 \mathrm{BoxesF}(F)\ \wedge\ \mathrm{QDRowsF}(F)
               \quad\Longrightarrow\quad
                         \mathrm{ChosenThin}(F,1).       \tag{DEX.4}
\]
In particular the extra \texttt{Chosen\allowbreak QD} premise appearing in one
source wrapper is unnecessary for this implication.

Fix $F$. Beyond one threshold, \texttt{Boxes\allowbreak F} supplies a scale $Y'$
and \texttt{Box\allowbreak E} for each actual $D=F.D(x,i)$. The other hypothesis
supplies a possibly different scale $Y$, with
$x^{1/J}\le Y\le x$, and its conductor bound
\[
       \log(Q_Sqd)\le2Y^{\epsilon/5}\log^2Y.
\]
The family has $J\ge1$. After imposing $x>1$, we have $Y>1$;
both $Y^{\epsilon/5}$ and $\log^2Y$ increase on this range. Hence
\[
                     \log(Q_Sqd)
                         \le2x^{\epsilon/5}\log^2x.      \tag{DEX.5}
\]
The two scales need not be equal: the box is used only through its
real-zero exclusion, and the conductor inequality only through (DEX.5).

Take $B$ from DEX1. CC18 supplies $c_s>0$ excluding real zeros in
$(1-c_s,1)$ for every nonprincipal character of conductor at most $8B$.
Set $c_0=\min(c_s,1/400)>0$. For any real primitive nonprincipal
\texttt{Is\allowbreak Char\allowbreak ED} character, there are two cases. If $q\mid N$ or
$d\mid N$, DEX1 gives the required exclusion. Otherwise (DEX.1) and
the primality of $q,d$ give $N\mid Q_S$. CC17 gives two-adic conductor
exponent at most three, so $N\mid8Q_{S,\mathrm{odd}}$ and $N\le8B$.
The finite-conductor exclusion then applies. Thus every such character
satisfies the precise \texttt{No\allowbreak Exc} premise of CC16 with this common
$c_0$. Nonreal and primitive principal characters require no additional
real-zero exclusion in that premise.

By (DEX.1) and (DEX.5), every character also has
$\log N\le2x^{\epsilon/5}\log^2x$. Apply CC16 with $A=2$ and $c'=1$.
CC14 identifies its centered function with the literal
\texttt{Global\allowbreak Theory\allowbreak Theory.\allowbreak g\allowbreak Chi}. We obtain, simultaneously for all
late $x$, all $i$, all primitive host characters, and every
$y\ge e^{x^{\epsilon/3}}/3$,
\[
             |g_\chi(y)|\le y\exp(-x^{\epsilon/8}).        \tag{DEX.6}
\]
The threshold may depend on $F$ through $C_o$, as the definition of
\texttt{Chosen\allowbreak Thin} permits. The coefficient $1$ does not. Taking the
maximum with $F.x_1+1$ and the two supplied family thresholds proves
(DEX.4) with the exact universal quantifiers. A family itself supplies
$\epsilon>0$ through any index at $x=0$, so the pointwise implication
also yields \texttt{Closed} on any class: choose outer bound $1$ and
the zero transported threshold. Consequently no separate literal
\texttt{Lit\allowbreak Thin} package is needed at this consumer; the analytic inputs
of CC retain their previously stated roles.

\paragraph{DEX3: the exact published Minkowski application.}
\label{claim:DEX3}
For $n\ge0$ put
\[
 a_p(n)=\sum_{j\ge0}\left\lfloor\frac{n}{p^j(p-1)}\right\rfloor,
 \qquad M(n)=\prod_{p\le n+1}p^{a_p(n)}.                  \tag{DEX.7}
\]
These are precisely the source's \texttt{minkowski\allowbreak Exp} and
\texttt{mink\allowbreak M}. Terms with $j\ge n+1$ vanish: $p^j\ge2^j>n$.
Primes larger than $n+1$ contribute zero. Thus the finite truncations
in the source do not change either expression.

The published input is Minkowski's theorem in Serre,
\emph{Bounds for the orders of the finite subgroups of $G(k)$},
Lecture I, \S1.1, Theorem 1(i): for a finite subgroup
$H\subset\mathrm{GL}_n(\mathbb Q)$, the $p$-adic valuation of $|H|$
is at most $a_p(n)$. See
\url{https://www.college-de-france.fr/media/jean-pierre-serre/UPL3821667391778701726_6___Bounds_for_the_orders.pdf},
printed p.~4. The source \texttt{Mink\allowbreak Mult} has exactly this rational
general linear group and assumes the subgroup finite. Applying the
theorem at every prime gives $|H|\mid M(n)$. For $n=0$ the group is
trivial and $M(0)=1$. This checks every rank, including the case not
covered by the published statement $n\ge1$, and proves the literal
\texttt{Mink\allowbreak Div}. No estimate over a larger coefficient field is
substituted for this rational theorem.

\paragraph{DEX4: the exact denominator-size constant twelve.}
\label{claim:DEX4}
The numerical consequence of (DEX.7) has an independent elementary
proof. Legendre's formula and $p\le2(p-1)$ give, term by term,
\[
 \left\lfloor\frac{n}{p^j(p-1)}\right\rfloor
 \le \left\lfloor\frac{2n}{p^{j+1}}\right\rfloor,
 \qquad M(n)\mid(2n)! .                                 \tag{DEX.8}
\]
Legendre's formula itself counts, among $1,\ldots,2n$, one contribution
for each multiple of $p$, a second for each multiple of $p^2$, and so
on. Since the prime-power factors are coprime for distinct primes,
the valuation inequalities prove the displayed divisibility. For
$n\ge1$, it follows that
\[
                       \log M(n)\le2n\log(2n).
\]
Write $z=fN(m)$ for the actual rank product. The original ranks have
$z\ge2$. Therefore
\[
 \log M(2z)\le4z\log(4z)
       =4z(2\log2+\log z)\le12z\log z.                 \tag{DEX.9}
\]
If \texttt{DGam\allowbreak Dvd\allowbreak F} holds, then eventually every original positive
integer $D_\Gamma(u)$ divides this same $M(2fN(m))$. Divisibility by
a positive integer ensures $D_\Gamma(u)>0$ and bounds it above by that
integer. Taking logarithms in (DEX.9) proves \texttt{DGam\allowbreak Size} with
the literal constant $12$, uniformly over all words and family indices.
The original exceptional-prime support is unchanged: all prime divisors
of $M(n)$ are at most $n+1$ by (DEX.7). This is a consequence of the
supplied divisibility, not a new construction of the original lattice.

\paragraph{DEX5: the numerical proviso from the actual rank formula.}
\label{claim:DEX5}
Let $D,L$ be original data and assume the source \texttt{Rank\allowbreak M} identity
\[
 d_M(u)=d_*(u)(T_u^k-1)^f.                              \tag{DEX.10}
\]
Here $d_M(u)=D.duS(u)$, and $d_*(u)$ is the exact standard-lattice
quantity \texttt{du\allowbreak Star\allowbreak U}. Retain the original equalities
\texttt{Fact\allowbreak Du\allowbreak Star} and \texttt{Fact\allowbreak Dstar}. They identify the positive
point quantities with $d_*(u)$ and $D_*$, without an isomorphism choice.
Put $\Gamma=\mathrm{GamL}(L)$ and
\[
 Y=\left(\frac{D_*}{d_*(u)}\right)^{1/(fk)}
                  \Gamma^{-2/(fk)} .                    \tag{DEX.11}
\]
The size datum's \texttt{Gam\_\allowbreak eq} identifies this $\Gamma$ with the
point's $\mathrm{GamP}$; thus $Y$ is the actual $T_*$.
Its original prime window says
$Y\le T_u\le(1+1/(fk))Y$, and $T_u\ge2$.
The rank fields give $1\le f\le k$ and hence
$2f\le2k\le2^k$; the last inequality follows by induction for $k\ge1$.
Using $N=fk$ and positive quantities throughout, (DEX.10)--(DEX.11)
give the exact ratio
\[
 \frac{d_M(u)\Gamma^2}{D_*}
       =\left(\frac{T_u}{Y}\right)^N(1-T_u^{-k})^f.       \tag{DEX.12}
\]
Bernoulli's inequality gives
$(1-T_u^{-k})^f\ge1-fT_u^{-k}\ge1/2$, while
$(1+1/N)^N\le e$ follows from $\log(1+t)\le t$.
Consequently
\[
                  \frac12\le\frac{d_M(u)\Gamma^2}{D_*}
                                      \le e<3.           \tag{DEX.13}
\]
For completeness, $e<3$ follows from its series:
$n!\ge2^{n-1}$ for $n\ge2$, with strict inequality at $n=3$, so
$\sum_{n\ge2}1/n!<1$. Thus (DEX.13) implies the exact
\texttt{Ch\allowbreak Proviso\_\allowbreak at} with $C_0=3$.

For a family, \texttt{Fact\allowbreak Du\allowbreak Star} and \texttt{Fact\allowbreak Dstar} are fields of
its actual late extra record. Hence an eventual \texttt{Rank\allowbreak M} on
that same family proves \texttt{Chosen\allowbreak Proviso}(F,3), after taking
the maximum with $F.x_1+1$. This proof does not require an independent
formal-degree estimate. Producing (DEX.10), with its exact class-count
and coefficient-space meanings, remains a separate task.

\paragraph{DEX6: combining closed predicates before choosing the family.}
\label{claim:DEX6}
A transported lower-bound record has components
$K_{\mathrm{pre}},K_t,K_{\mathrm{mid}},K_{\mathrm{last}}$;
the last two are functions of earlier parameter lists. Define $K\vee K'$
by the pointwise maximum of corresponding components. By inspection
of the original \texttt{Large} inequalities,
\[
 F.\mathrm{Large}(K\vee K',C_H)
 \quad\Longleftrightarrow\quad
 F.\mathrm{Large}(K,C_H)\wedge F.\mathrm{Large}(K',C_H).   \tag{DEX.14}
\]
This includes the dependent arguments of the middle and last bounds:
both are evaluated on the same lists of the same $F,C_H$.

Suppose finitely many predicates are \texttt{Closed} on the same
class. Take the minimum of their positive epsilon bounds. Once
$\epsilon$ has been fixed in that range, each produces a transported
threshold; take their finite join. Any one family subsequently chosen
large for that join satisfies every predicate. Restricting the class
preserves the implication, since the original class membership is
then available. These observations prove the source operations
\texttt{Closed.\allowbreak and}, \texttt{Closed.\allowbreak sub} and implication transport.
They do not combine witnesses from different families: the uniform
closed statements precede the single existential family choice.

\paragraph{DEX7: a smaller sufficient interface for \texttt{Design\allowbreak Exists}.}
\label{claim:DEX7}
Fix one original class $\mathcal C$ and one function $c_T$.
The following three inputs suffice:
\[
\begin{split}
 &\mathrm{DesignChoices}(\mathcal C,3,c_T),\\
 &\mathrm{Closed}_{\mathcal C}
    [c_T(\epsilon,F.C_{YL})\le F.C_{Tu}\Rightarrow\mathrm{TuSize}(F)],\\
 &\mathrm{Closed}_{\mathcal C}
    [\exists x_0\ \forall x\ge x_0\ \forall i,
                                   \mathrm{RankM}(F.D(x,i))].
\end{split}
                                                               \tag{DEX.15}
\]
They imply the literal $\mathrm{DesignExists}(\mathcal C,3)$.
In the last line \texttt{Rank\allowbreak M} is (DEX.10) on the original datum;
it is not a dimension match on a separately constructed model.

Use DEX5 and DEX6 to turn the last two inputs into one closed statement
giving the prime-size implication and \texttt{Chosen\allowbreak Proviso}.
Take the minimum of its epsilon bound and that of
\texttt{Design\allowbreak Choices}. For the requested $K$, first choose the
closed-statement threshold $K'$, then apply \texttt{Design\allowbreak Choices}
to $K\vee K'$. It gives one $F,C_H,C_5,C_b$, with $C_H\ge0$,
$C_5,C_b>0$, the exact class membership, the required large-parameter
property, and
\[
 \begin{gathered}
 \mathrm{ThickSetEventually},\ \mathrm{ChosenTT},\
 \mathrm{ChosenBad}(C_b),\ \mathrm{ChosenQD}(C_5),\\
 \mathrm{ChosenCode}(3,C_H),\ \mathrm{BoxesF},\ \mathrm{QDRowsF},\\
 \mathrm{DGamDvdF},\ c_T(\epsilon,C_{YL})\le C_{Tu}.
 \end{gathered}
\]
The source also supplies tame distinctness, $C_F=C_H$,
\texttt{Code\allowbreak Up\allowbreak F} and $C_{YL}\le9$; these remain attached to the same
family. DEX2 supplies \texttt{Chosen\allowbreak Thin}(1), DEX4 supplies
\texttt{DGam\allowbreak Size}, and the closed statements supply
\texttt{Chosen\allowbreak Proviso}(3) and \texttt{Tu\allowbreak Size}.
The seven chosen clauses, with positive constants $(C_5,1,C_b)$,
are exactly $\texttt{Chosen\allowbreak E}(F,3,C_H)$. Equation (DEX.14) retains
$\texttt{Large}(K,C_H)$. This proves every field of
\texttt{Design\allowbreak Exists}, in its original order of quantifiers.

The same proof works with a supplied closed \texttt{Chosen\allowbreak Proviso}
instead of the last line of (DEX.15). The source obtains that proviso
through \texttt{Class\allowbreak Number\allowbreak Formula}$\,{}_2$ and
\texttt{Rank\allowbreak Ties\allowbreak F}; DEX5 records the smaller numerical input actually
used. We have not asserted that either source-specific producer exists.

\paragraph{DEX8: one selector for all eventual row conditions.}
\label{claim:DEX8}
We reconstruct the finite logical step behind \texttt{class\allowbreak Sel}.
An eventual parameter assertion has the original order
\[
 \begin{gathered}
 \exists\epsilon_0\ \forall\epsilon,C_8,C_5,c_6,C_b\
 \exists K_0\ \forall C\\
 \exists c_1\ \forall c_D\ge c_1\
 \exists c_2\ \forall c_h\ge c_2 .
 \end{gathered}                                           \tag{DEX.16}
\]
Here the source positivity, standardness and largeness antecedents
are retained. An \texttt{Ev\allowbreak P} assertion then has
$\kappa.\mathrm{Pos}\Rightarrow\exists x_0\forall K$ with
$\kappa.\mathrm{Of}(K)$ and $K.s.x\ge x_0$.
Two such assertions combine by a minimum for $\epsilon_0$, a join
for $K_0$, and maxima for $c_1$, then $c_2$, then $x_0$.
Each maximum is taken only after all of its permitted dependencies
have been fixed. Induction gives the finite-conjunction rule.

Write $b\in\mathrm{Bd}(D,O,c)$ for an actual bundle, and
$\mathrm{Lk}(K,\beta,c,b)$ for the source link to the original
constructed point $\beta$. The crucial distinction in
\texttt{Row\allowbreak E} is its universal quantifier: eventually, for
\emph{every} such linked $b$, the stated predicate $E(\kappa,D,O,c,b)$
holds. The separate \texttt{Ex\allowbreak B} supplies one linked $b$ for each
$c$. Therefore, after the common threshold, that one $b$ satisfies
the conjunction of all finitely many row predicates. No simultaneous
choice theorem for unrelated existential bundles is being presumed.

Define $\mathrm{Built}(Q,\mathrm{Lk},\kappa,D,O,c,b)$ to mean that
there are $K,\beta,b'$ with the source context and core hypotheses,
the same link, and the literal dependent equality
\[
                 (D,O,b)=(\beta.D,\beta.O,b').            \tag{DEX.17}
\]
Let $Q_O(\kappa,D,O)$ say that, whenever $\kappa$ is positive,
for every $c$ there is such a built $b$ satisfying the combined $E$.
Take the class of original families admitting one $\kappa$ with
$\texttt{Link}(\kappa,F)$ and the source seven-duty object property
$Q_O$. At each late index choose a witness at $c=F.C_{sz}$.
This produces one selector $\sigma$. The defining witness gives
\texttt{Is\allowbreak Sel}, including (DEX.17), and every projected row condition
holds for this same $\sigma$.

The source's additional large-base slots cause no compatibility problem.
For finitely many \texttt{Row\allowbreak L} conditions and one \texttt{Ex\allowbreak BL},
replace their slots $R_j$ by $R=\bigwedge_jR_j$. The large-base
eventual statements combine by (DEX.16); at each fixed context the
finitely many eventual sets of sufficiently large bases are intersected.
Each is a member of the same large-base filter, so their finite
intersection is a member too. Restriction from $Q'\wedge R$
to each $Q'\wedge R_j$ changes no bundle or link. The preceding
construction applies, and forgetting $R$ in (DEX.17) gives the required
selector over $Q'$. Thus the common class and selector follow from the
literal eventual row and existence premises. Their actual arithmetic
suppliers, and recursive existence in this class, are still required.

\paragraph{DEX9: construct the family on all real indices from actual kits.}
\label{claim:DEX9}
Suppose $S$ is the supplied unbounded set of allowed scales, with
fixed $J$, and $x_1>1$. For every allowed scale $s$ with
$s.x\ge x_1$, suppose an original nonempty kit $K_s$ and its
\texttt{Num\allowbreak At} records are supplied. Choose one allowed $y\ge x_1$.
For each real $x$, use scale $x$ if $x>x_1$ and $x\in S$;
otherwise use the fixed fallback scale $y$. Denote this scale by $s_x$.
Define
\[
 I_x=\{i\in K_{s_x}.I:\ x\le x_1\ \text{or}\ x\in S\}.    \tag{DEX.18}
\]
It is nonempty on $S$, and at $x=0$, as the family definition requires.
For late $x$ outside $S$ it is empty. Consequently every actual late
member has $s_x.x=x$; the fallback creates no false statement about
an unallowed late scale.

For $i\in I_x$, copy the point's $D,L,G_x,G_y$, extra record,
local objects and ties. Copy the common numerical constants from $C$.
The kit's equality clauses identify ranks, Artin maps, primes, root
discriminants, canonical infima, exponents, selected weights and
period data at this fixed $x$. Its word maps give the same word codes,
and its coverage gives the family coverage by retaining the proof in
(DEX.18). The chosen \texttt{Num\allowbreak At} records give all parameter, size,
top-bound and constant identities. In particular $b_{\rm top}$ is
independent of the selected kit index; the original
$\mathrm{LaG}\ge x^\epsilon$ inequality gives the required lower
bound for it. These are exactly the fields of \texttt{fam\allowbreak Of\allowbreak Kits}.

Every supplied point property therefore holds at every actual member
by substitution, with no new transport of a lattice or metric.
The common lists are $F.c_{\rm pre}=C.c_{\rm pre}$,
$F.c_{\rm mid}=C.c_{\rm mid}$ and $F.c_{\rm last}=C.c_{\rm last}$.
For the output list its index 15 is $\max(C_{pde},J)$.
Any fallback kit supplies $J\le C_{pde}$ through its host-size
condition, so this is $C_{pde}$ and the output list also agrees.
Thus the family has the exact parameter lists needed for
\texttt{Large}, not merely comparable lists. The assertions
$F.\mathrm{Win}=S$ and $F.J=J$ follow directly from the definition.

\paragraph{DEX10: every field of the compatible \texttt{Design\allowbreak Choices} join.}
\label{claim:DEX10}
Assume the literal \texttt{Kit\allowbreak CPW} premise of \texttt{Design\allowbreak S}, for
the point predicate
$((\mathrm{withDuty}(\mathrm{withObj}(Q_\kappa)))\wedge
\mathrm{WallPt})\wedge G.\mathrm{pt}$, and fix $\lambda>1$.
We prove its \texttt{Design\allowbreak Choices} conclusion, using DEX3 for
\texttt{Mink\allowbreak Div}. All kit and point conditions here are on the
same original objects.

Take $\epsilon_0$ to be the minimum of the kit bound and $1/53$.
For $0<\epsilon\le\epsilon_0$, the kit gives constants
$C_{T2},C_8,C_5,c_6,C_b$ and threshold $K_0$, before any family
is chosen. For the requested $K$, first replace it by
$\mathrm{thrTu}(K,c_T(\epsilon,9))$: this replaces entry 13 in the
argument lists of its middle and last bounds by $c_T(\epsilon,9)$.
It is the exact threshold transformation for the subsequent
\texttt{set\allowbreak Tu} operation.

There is an elementary standard parameter record large for this
threshold and $K_0$. Explicitly choose its pre-parameter $p$ at least
$1000$ and both pre-bounds. Choose its middle parameter $a$ at least
$1$ and both middle bounds on the already fixed output list. Then
choose its last parameter $b$ at least $1$ and both last bounds on
that list followed by the middle list. In the source's
\texttt{std\allowbreak Const}, every pre-entry is $p$, every middle-entry is $a$,
and every last-entry is $b$; its fixed entries are
$C_t=401$, $C_c=2$, $C_{YL}=9$, $c_w=9/10$, $e_P=12$,
$C_r=C_{dg}=C_{ccK}=1$, and the stated fixed gap constants.
Its other positive gap entries are $a,b,b^{-1}$, in their prescribed
positions. Substitution verifies \texttt{Std}, both largeness
conditions and nonnegativity of the common gap record. No later
choice enters an earlier bound.

For this $C$, take the maximum of the kit's $c_D$ threshold and
SZ9's numerical threshold, then the maximum of their $c_h$
thresholds, then $x_1=\max(2,x_{\rm kit},x_{\rm num})$.
Choose one supplied kit at each allowed late scale. SZ9 gives its
\texttt{Num\allowbreak At}, \texttt{Ch\allowbreak TT} and \texttt{Ch\allowbreak Code}; DEX9 constructs
one family $F$. Its common choice is
\[
       \kappa=(\epsilon,C,c_D,c_h,C_5,c_6,C_b),
       \qquad C_H=C.P.C_F\ge1.                          \tag{DEX.19}
\]
The copied point conditions give the original class property and
$G.\mathrm{pt}$. The same set of windows supplies $G.\mathrm{cls}$.
The \texttt{Num\allowbreak At} identities fix the page constants and windows,
in particular $C_{YL}=9$. They also give
\texttt{Thick\allowbreak Set\allowbreak Eventually}, \texttt{Chosen\allowbreak TT},
$\texttt{Chosen\allowbreak Bad}(C_b)$, $\texttt{Chosen\allowbreak QD}(C_5)$ and
$\texttt{Chosen\allowbreak Code}(3,C_H)$. Their translations to the family
variables use $L.sc.x=x$ at late indices and $L.sc.\epsilon=\epsilon$;
the threshold $F.x_1+1$ makes these substitutions legitimate.

The \texttt{Wall\allowbreak Pt} and original duty clauses give
\texttt{Boxes\allowbreak F} and \texttt{QDRows\allowbreak F}, as well as the chosen wall
and twist properties, by their definitions. The \texttt{Mink\allowbreak Pt}
clause, applied to the precise rational divisor of DEX3, gives
\texttt{DGam\allowbreak Dvd\allowbreak F}; its tame clause gives \texttt{Tame\allowbreak Distinct\allowbreak F}.
The middle-list equality identifies $F.C_F=C_H$. The word count
is literally
\[
                 |\mathrm{Word}|=\lceil3\Gamma^2x^{C_H}\rceil,
\]
so the ceiling inequality yields
$|\mathrm{Word}|\le3\Gamma^2x^{C_H}+1$, proving
\texttt{Code\allowbreak Up\allowbreak F} with its original additive one.

Finally set $F'=\mathrm{setTu}(F,c_T(\epsilon,9))$.
This changes only the scalar field $C_{Tu}$ of the family record.
All the actual data and point predicates just verified are unchanged.
The definition of \texttt{thr\allowbreak Tu} gives exactly
$F'.\mathrm{Large}(K,C_H)$, and $C_{YL}=9$ gives
$c_T(\epsilon,F'.C_{YL})\le F'.C_{Tu}$ with equality.
For the source $\texttt{Link}(\kappa,F')$, its output clause first
sets $C_{Tu}$ back to $C.CTu$; hence the earlier exact output-list
identity verifies that clause as well. This proves all fields of
the literal \texttt{Design\allowbreak Choices} conclusion, on the original
target-compatible class and with one witness family.

\paragraph{DEX11: the recursive input that still produces the compatible class.}
\label{claim:DEX11}
Let $Q_O(\kappa,D,O)$ be the original object predicate chosen in
DEX8, or any supplied predicate of that type. Assume its actual
$\mathrm{RecExP}(\mathrm{withQD}(Q_O))$ and an actual
$\mathrm{HostExW}(\lambda,G.\mathrm{req})$, with $\lambda>1$.
The reconstructed RF10 point join applies: AC supplies its
Harish--Chandra and level constructors, SZ its initial sizes, QR25
the required level-prime assertion, and the assumed recursive
object supplies the actual $D,Y$ after every level choice. It gives
\texttt{Obj\allowbreak CP} for precisely $Q_O$, including the unchanged
\texttt{with\allowbreak QD} object and all original host, code and canonical
instance equalities.

Here the quantifier ``after every level choice'' is essential:
the positive integers $d_u$, positive $D_*$ and logarithmic window
are selected before the universal $Lv$. The recursive premise
then works for each $Lv$ with those same values. RF10 can therefore
use the actual prime levels obtained from QR25. An existential
recursive datum for one unrelated level would not suffice.

Forget $Q_O$ temporarily. HX1 supplies the required small prime,
and HW1--HW2, using GW5, construct the degree-selected host from
this point constructor and the supplied \texttt{Host\allowbreak Ex\allowbreak W}.
The injective word code can also be chosen uniformly, before the
host and all its positions. Put $H=m^5f x^\epsilon$ and
$M=\lceil3\exp(2C_\Gamma H)x^{C_F}\rceil_+$. The literal rank
windows give
\[
 \begin{gathered}
 m\le2x^{1/2-\epsilon/16},\qquad
 f\le4x^{2+8\epsilon+\epsilon/100},\\
 H\le128x^{9/2+9\epsilon},\qquad
 k\ge x^{9/2+10\epsilon}.
 \end{gathered}
                                                               \tag{DEX.21}
\]
For fixed positive $\epsilon$ and nonnegative $C_\Gamma,C_F$,
the ceiling estimate $M<2(3\exp(2C_\Gamma H)x^{C_F})$ gives
\[
 \log M\le\log6+256C_\Gamma x^{9/2+9\epsilon}
                         +C_F\log x=o(x^{9/2+10\epsilon}).
                                                               \tag{DEX.22}
\]
Choose one threshold depending only on these three parameters so
that $k\ge6$ and the right side is at most $k/4$. Since
$\log2>1/2$, we have $k/4\le(k-3)\log2$, hence
$0<M\le2^{k-3}\le2^{k-2}$. Enumerating the Boolean cube therefore
gives an injection $\operatorname{Fin}M\hookrightarrow
\{0,1\}^{\operatorname{Fin}(k-2)}$. This threshold is fixed before
the scale, the host and every $u_0$; it supplies the literal
\texttt{exists\_\allowbreak injective\_\allowbreak words} input of \texttt{kit\allowbreak S}.

This also corrects the sharper intermediate estimate in the historical
RF.45--RF.46: the rank window for $f$ contains $\epsilon/100$,
which those displayed exponents omitted. Use (DEX.21) there, with
positive exponent gap $\epsilon$, and choose
$x^\epsilon\ge2048C_\Gamma$ for its first eighth-bound.
The logarithmic eighth-bound is unchanged. The argument above still
proves the stated RF.48 conclusion $\#\mathrm{Word}\le2^{k-3}$.
Only that historical intermediate estimate is superseded; the
word bound and its uniform quantifier order are preserved.

Fix that host, its degree pair and the resulting injective word code, then
invoke the original point constructor once for each $u_0<n$.
The common host and code identify their word values; the source
canonical infimum and filling formulas give the exact kit equalities
as proved in HW2. The point indexed by $u_0$ supplies kit coverage
at that position. HW1 gives the wall condition for every point,
and the copied target-host property gives $G.\mathrm{pt}$.
Combining thresholds in the order (DEX.16) yields the literal
\texttt{Kit\allowbreak CPW} used in DEX10. That claim gives
\texttt{Design\allowbreak Choices} for
\[
 \mathrm{WithTw}\{F:\ \mathrm{clsP}(Q_O)(F)
                                  \wedge G.\mathrm{cls}(\lambda,F)\}.
                                                               \tag{DEX.20}
\]
This is the full conditional \texttt{rec\allowbreak S} route, with the original
object predicate and target class, not just a nonempty family.

The remaining recursive supply must be read against the latest
chapters. QW9 already proves its actual logarithmic window, and
MF20 discharges the mass-ratio premise on the same data. Therefore
IG14's historical window condition is closed. What remains in
that join is the actual \texttt{Rec\allowbreak At\allowbreak G} core and the eventual
$Q_O$ property, with their original objects and dependent equalities.
The immediate recursive output has predicate
$\kappa.\mathrm{Pos}\Rightarrow Q_O$; for the particular
$Q_O=\mathrm{QOf}$ of DEX8 this is equivalent to $Q_O$, since
its definition already begins with that same positivity guard.
Thus this simplification neither discards a hypothesis nor strengthens
the recursive producer.
DEX8 supplies the finite conjunction and selection step once its
bundle-existence and row premises are proved. DEX7 then supplies
\texttt{Design\allowbreak Exists} once the same class also has the required
\texttt{Tu\allowbreak Size} implication and actual rank formula. These are
the residual construction obligations; DEX1--DEX11 do not assume
them away or identify independently chosen bundles.


\section{The original local, theta, and operator dictionary}
\label{sec:local-theta-dictionary}

\status{LTD1--LTD8 reconstruct exact transports on the selected objects and
the finite Hecke argument producing the full local rank coverage.  In
particular, an actual \texttt{NBuilt} certificate implies the literal
\texttt{Locs\allowbreak Cover\allowbreak Tw}, and eventual such certificates imply
\texttt{Link\allowbreak Eta}.  No additional published Hecke theorem is needed for
this implication: its precise hypotheses occur in \texttt{Nd\allowbreak Spec.\allowbreak std}.
The existence of the compatible original geometric, theta, spectral,
and local certificates remains a producer obligation.}

The same field names occur in several different records in this part of
the source.  We distinguish the rational row operators, acting on
\(S.W_L(j)\), from the ambient operators, acting on
\(D.\mathrm{MtL}(\mu_w,j)\).  The former have Hecke indices through
\(r\), and the latter through \(n_V\).  Their comparison must use
the prescribed maps and normalizing scalars.  Merely having the same
name \texttt{h\allowbreak L} does not identify them.

All assertions below concern the pinned original source and the exact
data in their premises.  ``B'' classifies the complete human proof of
the stated implication; it does not classify the unconstructed premise
as elementary.  No new Lean replay is claimed.

% BEGIN LTD1
\paragraph{LTD1: all global selections come from one literal bundle.}
\label{claim:LTD1}
Let \(\mathcal A\) be a family class and let \(\sigma\) be a
\texttt{Raw\allowbreak Sel Bundle} on that class.  Thus, whenever
\(\mathcal F\in\mathcal A\), \(x>\mathcal F.x_1\), and
\(i\in\mathcal F.I_x\), there is a specified bundle
\[
 b=\sigma(\mathcal F,x,i)
  =(R,S,P,\mathrm{pin},DH,C,B)
\]
on the original \(D=\mathcal F.D(x,i)\), its original object \(O\),
and the same size constant \(c=\mathcal F.C_{\rm sz}\).
Here \(R\) is the tied record, \(S\) its spectral realization,
\(P\) its row family, \(C\) the compatible columns, and
\[
 B:\texttt{Real\allowbreak TB}
   (R,P,C.\mathrm{face},C.\mathrm{Sp},
                 \mathrm{zsc}(P,DH,c),S).
 \tag{LTD.1}
\]
The source's selections are exactly
\[
 \begin{split}
 \mathrm{sel}(\mathcal F)&=
   \bigl\langle\mathcal F.x_1,
           (x,i)\longmapsto b.R\bigr\rangle,\\
 \mathrm{selS}&=b.S,\quad \mathrm{selP}=(b.P,b.\mathrm{pin}),\quad
 \mathrm{selDH}=b.DH,\quad \mathrm{selC}=b.C,\quad
 \mathrm{selB}=b.B.
 \end{split}                                                    \tag{LTD.2}
\]
Indeed, substituting these projections in the definition of
\texttt{Zsc\allowbreak At} gives the very scale in (LTD.1): its last component
is set to zero and its previous component is changed by
\(-c\,z_{6,u}\).  All other terms use the same \(P,DH,D,O.L\).
Consequently the dependent type of \(b.B\) is precisely the type
required by \texttt{Sel\allowbreak B}; no new compatible realization is chosen.

In particular, the indefinite realization selected by
\texttt{sel\allowbreak IOf} is exactly \(b.B.Q\), and its coefficient field
selected by \texttt{sel\allowbreak KOf} is exactly \(b.B.T.Q.c.F\).
The source's \texttt{wire\_\allowbreak sel} uses this one \(Q=b.B.Q\) for its
all-tuples comparison, its point comparison and its base-rank
comparison.  Likewise \texttt{link\_\allowbreak alt} takes its spectral
instance and its instance tie from \(b.S\), its real volume tie
from \(b.R.T.\mathrm{theta.real}\), and its scale tie from the
original \(\mathcal F\).  These are common-object identities, not
assertions inferred from equal dimensions.

The condition \(x>\mathcal F.x_1\) already implies the weak
threshold requirement \(\mathcal F.x_1\le x\) in (LTD.2).
No further eventual threshold is needed to align these projections.
This proves the compatibility of the selection once \(\sigma\)
is supplied; it does not construct the bundles.
% END LTD1

% BEGIN LTD2
\paragraph{LTD2: the local installation preserves the original theta and form objects.}
\label{claim:LTD2}
Fix \(D\) and \(L:\texttt{Loc\allowbreak All}(\texttt{lay0}(D))\), and put
\(D_L=\texttt{with\allowbreak Loc}(L)\).  The local installation replaces the
local index set, local rings, modules, projections, automorphic
specializations and associated local comparisons by the fields of
\(L\).  It also sets \(Z,Z_{\rm good}\) to the specified actual
subsets and replaces the row number \(\eta_J\), for \(1\le j\le m\),
by
\[
       \eta_{J,L}(j,w)=\sum_{a\in\mathrm{locs}_L(j,w)}\eta_{\rm loc}(a).
       \tag{LTD.3}
\]
Outside that range it preserves the previous \(\eta_J\).
It preserves the host, its coefficient fields and their embeddings,
all row spaces \(\mathrm{MtL},\mathrm{Mt},\mathrm{Mj}\), their
original realization maps and lattices, and their theta maps.

Here is a direct check of the claim.  \texttt{Loc\allowbreak Nd.\allowbreak nd} changes
exactly the local-indexed fields of \texttt{NData}.  In
\texttt{Loc\allowbreak Tw.\allowbreak ntw}, only the maps from these new local indices to
the old twist indices, their local coefficient maps, and the local
characters are replaced.  The twist indices, their rings \(O_0,O,O_k\),
the two side embeddings, the universal Artin maps, the twist
characters, signatures and toric theta values are copied.  In
\texttt{Loc\allowbreak All.\allowbreak lay}, \texttt{nexp} and \texttt{nth} are record
copies.  The subsequent \texttt{fix\allowbreak Z} changes only the two finite
sets; \texttt{re\allowbreak Eta} changes only \(\eta_J\) in the underlying
row datum.  Therefore every claimed carrier and every claimed map
has the same defining expression after the necessary dependent
type identifications.  In particular, the coefficient-valued
\(\mathrm{th},\mathrm{thPsi},\mathrm{thU},\mathrm{thP}\), the
form action and Hecke operators, and the actual evaluation maps
of \texttt{NTh} are the originals.

For an original form family \(M\), normalized choice \(C\),
character frame \(cf\), and tuple tie \(tp\), the source defines
\[
 M_L=\texttt{mt\allowbreak L}(L,M),\qquad C_L=\texttt{ch\allowbreak L}(L,C),\qquad
 cf_L=\texttt{cf\allowbreak L}(L,cf),\qquad tp_L=\texttt{tp\allowbreak L}(L,tp).
 \tag{LTD.4}
\]
These copy the old fields, including \(C.\mathrm{twI},C.\mathrm{sd},
C.Q,C.k',C.\psi',C.\mathrm{embW},C.e_{Mt}\), and the tuple
characters.  Thus the exact semilinear basis and lattice equalities
in \texttt{NChoice} retain their maps, not merely their numerical
ranks.  This is the source's \texttt{Kept} comparison at this step.

Preservation is an implication on an existing object.  It does not
turn an initialized zero theta map into a nonzero one.  In particular,
the separate \texttt{NTh.\allowbreak Th\allowbreak Ev} and \texttt{NTh.\allowbreak Emb\allowbreak BRange}
propositions are not fields proved merely by presenting an
\texttt{NTh} record.  Whenever such a proposition has been proved
for the original maps, the equality just checked transports it
unchanged; its mathematical content has not been manufactured by
the installation.
% END LTD2

% BEGIN LTD3
\paragraph{LTD3: integral simultaneous generalized eigenspaces.}
\label{claim:LTD3}
Let \(O\) be a discrete valuation ring, \(K\) its fraction field,
and \(V\) a finite-dimensional \(K\)-vector space.  Let
\(\mathcal S\subseteq\operatorname{End}_K(V)\) be a possibly
infinite pairwise commuting family.  Suppose each member splits
over \(K\), in the sense that its generalized eigenspaces span
\(V\), and every eigenvalue having a nonzero generalized
eigenspace lies in \(O\).  Let \(T=O[\mathcal S]\).  For an
\(O\)-algebra homomorphism \(\psi:T\to O\), define
\[
 W_\psi=\bigcap_{t\in T}\operatorname{GEig}(t,\psi(t)),
 \qquad
 \mathfrak m_\psi=\ker(T\xrightarrow{\psi}O\longrightarrow O/\mathfrak m_O).
 \tag{LTD.5}
\]
The scalar inclusion \(O\hookrightarrow K\) is implicit in
\(\operatorname{GEig}\).  There are only finitely many nonzero
\(W_\psi\), and
\[
                     V=\bigoplus_{W_\psi\ne0}W_\psi.             \tag{LTD.6}
\]

We prove this without a finiteness assumption on \(\mathcal S\).
Begin with \(V\).  If a nonzero current block is split into two
or more generalized eigenspaces by some \(s\in\mathcal S\),
perform that refinement.  Every new block is invariant under
all of \(\mathcal S\), since commuting operators preserve the
kernels of powers of \(s-\lambda\).  The restriction of another
split operator to an invariant subspace still splits: its
generalized spectral projections are polynomials in that
operator, hence preserve the invariant subspace.  Each strict
refinement increases the number of nonzero summands, which is
at most \(\dim_KV\).  We therefore reach a finite direct sum
\(V=\bigoplus_aV_a\) such that every \(s\in\mathcal S\) has on
each nonzero \(V_a\) exactly one generalized eigenvalue
\(\lambda_a(s)\in O\).

On \(V_a\), write \(s=\lambda_a(s)I+N_s\).  Each \(N_s\)
is nilpotent, and all these nilpotents commute.  A polynomial
with zero constant term in finitely many commuting nilpotents
is nilpotent: in a sufficiently high power, every monomial has
an exponent at least the nilpotence exponent of one factor.
Consequently every polynomial expression \(t\in T\) acts as
\(\lambda_a(t)I+N_t\), where \(N_t\) is nilpotent and the scalar
is the same polynomial in the \(\lambda_a(s)\) and its
\(O\)-coefficients.  This scalar is independent of the expression
for \(t\).  Indeed, a relation would make \(cI\) nilpotent for
the difference \(c\) of two putative scalars; on the nonzero
space \(V_a\), that implies \(c=0\).  Thus \(\lambda_a\) is an
\(O\)-algebra homomorphism \(T\to O\).

Two distinct blocks were separated at a refinement and hence
have distinct values on some generator.  Their characters
\(\lambda_a\) are distinct.  Their generalized character spaces
are exactly the blocks: on any other block, a generator with
a different value has invertible scalar-plus-nilpotent difference.
The same observation proves that any nonzero \(W_\psi\) is one
of these blocks.  This proves (LTD.6).  For \(V=0\), all spaces
in (LTD.5) are zero and the conclusion is immediate.

For a maximal ideal \(\mathfrak m\subset T\), put
\[
 V_{\mathfrak m}=\sum_{\mathfrak m_\psi=\mathfrak m}W_\psi,
 \qquad V^{\rm off}_{\mathfrak m}
            =\sum_{\mathfrak m_\psi\ne\mathfrak m}W_\psi.        \tag{LTD.7}
\]
Each \(\mathfrak m_\psi\) is maximal: residue evaluation is
surjective because it contains the images of all constant
elements of \(O\).  Therefore, for any finite set \(\mathcal M\)
containing every maximal ideal of \(T\), grouping (LTD.6) by
these residue kernels gives
\[
 \sum_{\mathfrak m\in\mathcal M}\dim_KV_{\mathfrak m}=\dim_KV,
 \qquad
 \bigcap_{\mathfrak m\in\mathcal M}V^{\rm off}_{\mathfrak m}=0.
 \tag{LTD.8}
\]
Ideals with no supported character contribute a zero summand.
The argument does not assert diagonalizability, and does not
need to prove that \(T\) is finite as an \(O\)-module.
% END LTD3

% BEGIN LTD4
\paragraph{LTD4: the actual local specialization gives complete rank coverage.}
\label{claim:LTD4}
Retain \(D,L,M,C\) from LTD2 and a specified
\[
 \mathrm{Sp}:\texttt{Nd\allowbreak Spec}(\mathrm{Pl},M,C,\mathrm{big},L.X,L.T).
 \tag{LTD.9}
\]
For every \(1\le j\le m\) and every original place \(w\) with
\(D.\mu_{\rm At}(w)\ne D.\mu_0\), one has
\[
 \sum_{a\in D_L.\mathrm{locs}(j,w)}
          \operatorname{rank}_{O_a}M_a
       =\dim_{D.L}D.\mathrm{MtL}(D.\mu_{\rm Tw}(w),j).
 \tag{LTD.10}
\]
Moreover the original projections \(\mathrm{pr}_{Mt,a}\), after
the explicit index casts, are jointly injective on the right-hand
space.  Equation (LTD.10), for these exact quantifiers, is
\texttt{Locs\allowbreak Cover\allowbreak Tw} on \(D_L\).

Fix \(j,w\).  Take \(O=\mathrm{Sp}.B.O_c(j,w)\), its fraction
field \(K\), \(V=\mathrm{Sp}.B.V(j,w)\), and the family
\(H_{v,k}=\mathrm{Sp}.B.h_K(j,w,v,k)\) with the literal
\texttt{Gen\allowbreak At} restrictions: \(v\) is split, away from the
coefficient prime, the specified auxiliary set and \(Q_c(j,w)\),
and \(1\le k\le n_V\).  The field \texttt{Sp.\allowbreak std} gives
pairwise commutation and integrality of every occurring
generalized eigenvalue.  The field \texttt{Sp.\allowbreak B.\allowbreak split} gives
the splitting hypothesis of LTD3.

The original row-basis axiom first makes
\(D.\mathrm{MtL}(\mu_w,j)\) finite dimensional: its basis maps
to a linearly independent family in the finite-dimensional
original complex row space.  The field \texttt{i\allowbreak V\_\allowbreak basis} maps
that same basis to a \(K\)-basis of \(V\).  Thus
\[
 \dim_KV=\dim_{D.L}D.\mathrm{MtL}(\mu_w,j),\qquad
                 i_V\text{ is injective}.                         \tag{LTD.11}
\]
For the second assertion, write a vector in the original basis.
A zero image gives a zero combination in the image basis, so
each coefficient has zero image under the field embedding and
is zero.

The fields \texttt{mx\_\allowbreak max}, \texttt{mx\_\allowbreak inj}, and
\texttt{mx\_\allowbreak all} identify the finite set
\(D_L.\mathrm{locs}(j,w)\) bijectively with the maximal ideals
used in (LTD.8).  For each index \(a\), the specified map
\[
 \Psi_a:K_a\otimes_{O_a}M_a\longrightarrow V
\]
is semilinear for the field isomorphism \(\mathrm{phK}_a:K_a\simeq K\),
is injective, and has image exactly \(V_{\mathrm{mx}(a)}\).
Transporting a basis along this isomorphism gives
\[
 \operatorname{rank}_{O_a}M_a
       =\dim_KV_{\mathrm{mx}(a)}.                                 \tag{LTD.12}
\]
Here \(M_a\) is finite free over the given DVR, so extending a
free basis to its fraction field preserves its length.
Summing (LTD.12), using injectivity and exhaustiveness of
\(\mathrm{mx}\), then (LTD.8) and (LTD.11), proves (LTD.10).

Finally suppose every original local projection of \(y\) is
zero.  Membership \(a\in\mathrm{locs}(j,w)\) supplies the
identities of the row number, place and twist index; these give
the source casts \texttt{cast\allowbreak Mt\allowbreak Inv}.  The actual
\texttt{Psi\_\allowbreak pr} equality says
\[
 \Psi_a(\mathrm{pr}_{Mt,a}(y))-i_V(y)
                    \in V^{\rm off}_{\mathrm{mx}(a)}.
 \tag{LTD.13}
\]
Its first term is zero.  By (LTD.8), \(i_V(y)=0\), and (LTD.11)
gives \(y=0\).  This proves joint injectivity on the original
space, not on an arbitrarily chosen space of the same dimension.

No \texttt{Aut\allowbreak Spec} hypothesis is needed for this proof.
If one wants the entire \texttt{Normalized\allowbreak Local}\allowbreak\texttt{Compatibility}
record, \texttt{Sp} supplies all its displayed algebraic fields;
the two additional equalities \texttt{Aut\allowbreak Spec.\allowbreak iota\allowbreak Q} and
\texttt{iota\allowbreak QC} turn \texttt{Sp.\allowbreak iota\allowbreak C\_\allowbreak coh} into the required
complex embedding square.  This is exactly the \texttt{cover}
construction in \texttt{hecke\allowbreak Open\_\allowbreak with\allowbreak Loc\_\allowbreak alt}.
% END LTD4

% BEGIN LTD5
\paragraph{LTD5: the exact original certificate implies the terminal coverage link.}
\label{claim:LTD5}
For every original datum \(D'\), every fixed local Langlands
input \(h_{LL}\), and every supplied
\(N:\texttt{NBuilt}(h_{LL},D')\), one has
\[
                    \texttt{Locs\allowbreak Cover\allowbreak Tw}(D').                    \tag{LTD.14}
\]
Indeed, \(N\) contains \(D,L,M,C,\mathrm{Sp}\) as in LTD4 and
the equality \(D'=\texttt{with\allowbreak Loc}(L)\).  Substitute this
equality and apply LTD4.  The proof never applies the theorem
denoted by \(h_{LL}\); it uses the already supplied local
certificate.  The role of local Langlands in producing that
certificate is separate.

The stronger \texttt{NTrue} record additionally identifies the
entire dependent tuple
\[
 (D',cf',C',tp')
       =(D_L,cf_L,C_L,tp_L).                                     \tag{LTD.15}
\]
Thus (LTD.14) is compatible with the exact characters and choice
used later.  In \texttt{Rows\allowbreak At\allowbreak Start.\allowbreak }\allowbreak\texttt{n\allowbreak True\allowbreak At\allowbreak Weight}, both the datum
equality and (LTD.15) are reflexive: the constructor packages
these very objects.  No isomorphic replacement is inserted.

Consequently, for a fixed \(\mathcal F:\texttt{Fam\allowbreak P}(\epsilon)\),
the sufficient family condition is precisely
\[
 \exists x_0\ \forall x\ge x_0\ \forall i\in\mathcal F.I_x,
       \exists N:\texttt{NBuilt}(h_{LL},\mathcal F.D(x,i)).
 \tag{LTD.16}
\]
Applying (LTD.14) separately at each \((x,i)\), without changing
\(x_0\), gives the source's \(\texttt{Link\allowbreak Eta}(\mathcal F)\).
When the producer requires \(x>\mathcal F.x_1\), replace
\(x_0\) by \(\max(x_0,\mathcal F.x_1+1)\) before applying it.
If (LTD.16) is provided for every member of the actual selected
class, the same assertion is \texttt{Closed} on that class,
and on its \(D3\) subclass by the defining inclusion.

This is a strict reduction of the coverage consumer to the
original \texttt{NBuilt}/\texttt{Nd\allowbreak Spec} producer.  It does not
assert that an arbitrary \texttt{TSel} carries an \texttt{NBuilt}
certificate.  The generic source theorem \texttt{link\_\allowbreak eta}
uses only the selected tied records' \texttt{Mt\allowbreak Fam},
\texttt{Normalized\allowbreak Local\allowbreak Compatibility} and \texttt{Hecke\allowbreak Tied};
its \texttt{Duty} argument is unused.  That generic route calls
\texttt{hecke\allowbreak Std}.  The concrete route (LTD.16) instead retains
the \texttt{std} proof already present in the actual local
construction, and therefore introduces no second independent
Hecke integrality supplier.
% END LTD5

% BEGIN LTD6
\paragraph{LTD6: the original row spectrum and its normalized Hecke action.}
\label{claim:LTD6}
Fix one source geometric state \(x:\texttt{St}(P)\) and a
compatible \(l:\texttt{LC}(h_{LL},x)\).  This is the source's
\texttt{Strict\allowbreak Henselian}\allowbreak\texttt{Acyclic\allowbreak Level} in the expanded naming.
Its rational form datum is
\[
 S=\texttt{rat\allowbreak Dat\allowbreak R}(x.m.z.R,x.m.z.cG,l.\mathrm{rat}.k_{K_s},
                                                l.\mathrm{rat}.hk).
\]
Under the installation \(L\), the source chooses
\[
 D_L=\texttt{DL}(x,L)=\texttt{with\allowbreak Loc}(L),\qquad
 \texttt{dat\allowbreak L}=\texttt{Dat.\allowbreak transfer}(x.m.z),
\]
using the identity finite-group comparison.  Consequently
\[
 S_L=\texttt{rat\allowbreak L}(L,S),\quad H^{\rm row}_{L,j,k,v}
              =l.\mathrm{rat}.\mathrm{In}.h_L(j,k,v),\quad
 e_{L,j,\tau,k,v}=l.\mathrm{rat}.\mathrm{In}.\mathrm{eig}(j,\tau,k,v).
 \tag{LTD.17}
\]
These are equalities of the operators on the transported original
spaces.  Their proof is substitution in \texttt{SL},
\texttt{ys\allowbreak L}, \texttt{hk\allowbreak H}, \texttt{hk\allowbreak HL}, and \texttt{eig\allowbreak D}:
\texttt{SL\_\allowbreak heq} and \texttt{h\allowbreak L\_\allowbreak eq} reduce to reflexive
equalities.  A final row reading \(Y\) is related by the specified
\texttt{YRel}, which keeps its actual representation and its
\(e_W,j_{Mt}\) readings.  \texttt{normalized\allowbreak Row\allowbreak Datum} uses that
relation to identify its representation with
\(\tau\mapsto\texttt{row\allowbreak Rep\allowbreak R}(R.\theta,Y.e_{FT}(\tau))\).

To check that this identification is useful for the spectral
consumer, write \(G_j\) for \texttt{good\allowbreak L}: the intersection of
the eigenspace of the original finite-group projector with the
specified generic-character eigenspaces.  All operators and
characters defining this intersection are preserved in LTD2 and
(LTD.17), so its transported subspace is literally the same
\(G_j\).  The original \texttt{Rat\allowbreak In.\allowbreak span} assertion is
\[
 G_j\ \subseteq\sum_{\substack{\tau:\,\mathrm{isot}(j,\tau)\ne0\\
                                      \mathrm{goodRow}(\tau)}}
 \left(G_j\cap\bigcap_{\substack{1\le k\le r\\
                      v\text{ split},\ v\notin S_b}}
       \ker(H^{\rm row}_{j,k,v}-e_{j,\tau,k,v})\right).
 \tag{LTD.18}
\]
Substitution proves exactly (LTD.18) for the final original row
model.  Likewise, if \(\Lambda_j\) is the
\(\mathcal O_{D.L}\)-span of its prescribed integral coefficient
module, the original \texttt{h\allowbreak L\_\allowbreak int} becomes, with the same
quantifiers and the same \(S_b\),
\[
 \forall j,k,v\ \exists n\quad
     (\operatorname{N}v)^n H^{\rm row}_{j,k,v}
                    (\Lambda_j\cap G_j)\subseteq\Lambda_j,
 \tag{LTD.19}
\]
for \(1\le j\le m\), split \(v\notin S_b\), and \(k\le r\).
No uniform denominator exponent is inferred: \(n\) may depend
on \(j,k,v\), exactly as in the original field.

The coefficient-valued Hecke action is related to this row
action by the actual map \(e_F\) and actual unit \(u_{i,b,v}\):
\[
 e_F H^{\rm exp}_{i,j,b,k,v}
       =u_{i,b,v}^{\,k}(1\otimes H^{\rm row}_{j,k,v})e_F,
       \qquad 0\le k\le r,\quad v\notin S_0(i),\ v\text{ split}.
 \tag{LTD.20}
\]
This follows from \texttt{Side\allowbreak B.\allowbreak hk\allowbreak Eq} and the equality
\(l.b.YS.e_F=l.b.Bs.e_F\); \texttt{hk\allowbreak H} defines the scalar as
the \(k\)-th power of the embedded \texttt{c\allowbreak Hk} unit.  For
\(k=0\) it is one.  For any scalar \(a\), induction on \(n\)
in (LTD.20) gives
\[
 e_F(H^{\rm exp}-u^ka)^n
       =u^{kn}(1\otimes H^{\rm row}-a)^n e_F.                    \tag{LTD.21}
\]
Thus \(e_F\) carries both ordinary and generalized eigenspaces
to those with the prescribed rescaled eigenvalue.  This proves
the exact spectral transport; it does not assert the initial
spanning or integrality in (LTD.18)--(LTD.19) without
\texttt{Rat\allowbreak In}.

There is a second family \(\mathrm{Sp}.B.h_L/h_K\), used in
LTD4, on \(\mathrm{MtL}(\mu_w,j)\) with \(k\le n_V\).
Its operator square is \texttt{NLoc\allowbreak Base.\allowbreak h\allowbreak L\_\allowbreak def/\allowbreak h\allowbreak K\_\allowbreak def}, and
its local action square is \texttt{Nd\allowbreak Spec.\allowbreak act\_\allowbreak h\allowbreak K}.  These
are not (LTD.17).  In particular,
\texttt{Choice\allowbreak Side.\allowbreak e\allowbreak Mt\_\allowbreak level} only gives an extending complex
embedding, equality of the level subgroup, and equality of its
character after \texttt{k\allowbreak Map\allowbreak E}.  It contains no pointwise
equality for \(e_{Mt}\) and no Hecke square.  The missing
comparison cannot be supplied by silently identifying these
two operator families.
% END LTD6

% BEGIN LTD7
\paragraph{LTD7: localization transports the existing theta evaluation and export series.}
\label{claim:LTD7}
Continue with the same \(x,l,L\) as in LTD6.  The source's
\texttt{Moduli\allowbreak Points\allowbreak State} contained in \(l\) supplies an actual
\texttt{Geo\allowbreak Tie.\allowbreak Theta\allowbreak Tie} between its original expansion and
\(D.\mathrm{nth}\).  Copy its points, covering maps, evaluation
identities, unramified forms, Hecke identities and unramified
extension condition.  By LTD2 the coefficient fields,
evaluation maps and Hecke operators appearing on both sides
are unchanged, and by LTD6 their real model has been transported
by the identity comparison.  Every copied equality therefore
has exactly the required target and source.  This constructs
the final \texttt{Theta\allowbreak Tie}; it is the source's
\texttt{geo\allowbreak L}.  In particular, it uses the existing evaluation
identity, rather than deducing one from the existence of a
theta-named record.

There is also an exact transport for the exported series.  Let
\(X'\) be any \texttt{NExp\allowbreak Tied} on \(D_L\), and suppose
\[
      h:\texttt{Exp\allowbreak Rel}(L,l.b.Bs.\mathrm{exp},X').                 \tag{LTD.22}
\]
This comparison has precisely three fields.  It transports
\texttt{Tau\allowbreak True}; it preserves the ordered pair of bad-place
sets \((QP,DP)\) of every admissible level; and it transports
every \texttt{Reads} factorization.  The latter says that a
specified pointwise reading \(r(f,g)\) factors through the
actual realization by an injective semilinear coefficient map:
\[
 r(f,g)=\omega((\mathrm{real}(f))(g)),\qquad\omega\text{ injective}.
 \tag{LTD.23}
\]
No equality of the two realizations themselves is presumed.

Construct the localized export by retaining the original
characters, finite class sets and representatives, measures,
Hermitian frames, integral lattices, additive character and
normalizing scalars.  Only their ambient datum is relabeled.
The twist maps in the two directions are inverse identity
readings of \((Q,\chi,\mathrm{res})\); thus all original
admissibility and finite-class conditions have the same
meaning.  The \texttt{setup\allowbreak L} construction has the following
pointwise equalities under this relabeling:
\[
 \theta_L=\theta,\qquad \mathrm{gens}_L=\mathrm{gens},\qquad
 \xi_L=\xi,\qquad \mathrm{classSum}_L=\mathrm{classSum},
 \tag{LTD.24}
\]
and its predicates \texttt{Is\allowbreak Datum} and \texttt{Is\allowbreak Hecke\allowbreak Datum}
are equivalent to the originals.  To see this, expand a theta
or class-sum vector into its original frame, place map, finite
representatives and vector function.  Each of these has just
been retained.  The datum predicates assert splitting,
avoidance, finiteness, support and membership for those same
arguments; the Hecke predicate asserts the same chosen element,
representatives and determinant.  Each clause is unchanged.

An original \texttt{Cell\allowbreak Tie} can now be copied field by field.
Its support, class, polynomial coefficient and \(\mathrm{rdV}\)
maps are kept.  Its datum-surjectivity and theta equalities use
the first equality in (LTD.24), its twisted-generator equality
uses the second, its universal coefficient identities and
off-support vanishing use the third, and its theta-projection
identity uses the fourth.  The equality of integral lattices
and injectivity of \(\mathrm{rdV}\) refer to the unchanged maps.
The finite-group and level-character equations are unchanged
as well.  This checks every field of \texttt{Cell\allowbreak Tie},
including its conditional antecedents.

Keep exactly the original normalization scalar \(\nu\ne0\).
The norm assertion \texttt{Is\allowbreak Norm} transports in both directions
because its actual family, measures and maps have been retained.
The three rationality conditions on the base theta values,
twisted generators and universal coefficients are equivalent
by (LTD.24).  Hence both conditional clauses in
\texttt{Is\allowbreak Scalar}, as well as \(\nu\ne0\), transport.  This
preserves the original scalar choice even when a conditional
rationality antecedent is false; it does not add an unconditional
nonvanishing theorem.  The copied Hecke tie follows from the
\texttt{Is\allowbreak Hecke\allowbreak Datum} equivalence, and the universal-character
tie is retained by LTD2.  We have constructed the localized
\texttt{Export\allowbreak Tie} with its original normalization.

To finish \texttt{Ser\allowbreak At}, retain its bad-place equality, its
selected-place equality, the bounded-lattice equivalence, the
off-range zero condition and its original lambda datum and
lambda ties.  Its only remaining field, \texttt{join}, is the
factorization (LTD.23) for the pointwise expression
\[
                 \mathrm{rdV}(f)(\mathrm{finHom}(\mathrm{PV}(g))).
\]
The old \texttt{join} supplies that reading; apply
\(h.\mathrm{reads}\) to it.  This produces its required
injective semilinear map for \(X'\), on the same \(f,g\).
We have therefore proved
\[
 \texttt{Exp\allowbreak Rel}(L,l.b.Bs.\mathrm{exp},X')
  \Longrightarrow
 \operatorname{Nonempty}\bigl(\texttt{Ser\allowbreak At}
     (\mathrm{ysL},\mathrm{datL}.\mathrm{geo}.v_0,
                     X',\mathrm{pvL},\mathrm{cfD},\mathrm{pqD})\bigr).
 \tag{LTD.25}
\]
The original \texttt{Theta\allowbreak Tie}, \texttt{Ser\allowbreak At} and the comparison
(LTD.22) remain explicit premises.  Their analytic rationality,
convergence and initial realization content is not replaced
by the transport proof.
% END LTD7

% BEGIN LTD8
\paragraph{LTD8: the coefficient field is the actual embedded field.}
\label{claim:LTD8}
Fix the source's real choice \(ch\).  Its predicate
\texttt{Field\allowbreak OK} means \texttt{Numbers\allowbreak In} for the particular
embedding \(\iota=D.\mathrm{embLL}:D.L\hookrightarrow\mathbb C\).
The required numbers comprise the specified character values,
eighth roots, the place-dependent roots of unity, phases,
square roots of place norms, finite expansion jets, critical
toric ratios and normalized bottom sums.  Thus, for any field
embedding \(\sigma:K\hookrightarrow\mathbb C\),
\[
 \texttt{Field\allowbreak In}_\sigma(ch)\quad\hbox{and}\quad
 \sigma(K)\subseteq\iota(D.L)
       \quad\Longrightarrow\quad\texttt{Field\allowbreak OK}(ch).            \tag{LTD.26}
\]
For every required number, compose its asserted membership in
\(\sigma(K)\) with the displayed containment.  In the bottom
sum clause retain the same triangles and the same nonzero
period scalars, and apply the containment to each normalized
sum.  All quantifiers over tuples, supported expansions and
places are untouched.  This proves every clause of
\texttt{Numbers\allowbreak In}, hence (LTD.26).

Conversely, \texttt{Field\allowbreak OK} implies \texttt{Has\allowbreak Number\allowbreak Field}:
take the image field \(\iota(D.L)\subseteq\mathbb C\).
The embedding identifies it with the original number field,
so it has the same finite dimension over \(\mathbb Q\);
every required number is in that image.  The converse needs
the containment in (LTD.26); existence of some number field
alone gives no such containment for the fixed datum.

The jets clause immediately makes the specified top and row
jets rational at their actual pairs of points, and the
character clause gives \texttt{Val\allowbreak LL}.  For the further
\texttt{Field\allowbreak Cond} consumer the source also requires the same
real additive character and theta datum, equality of the two
grid characters at every relevant index, equality of the
selected place set and its depths, nonvanishing of the
denominator Hecke values at the specified positive integers,
existing bottom triangles, and integrability of their two
kernel expressions.  Under those additional premises,
substitution gives the required roots and character values;
the critical-ratio clause supplies
\(\iota(\rho)=\Lambda(s)/\Lambda_0(s)\), and the nonzero
denominator gives \(\Lambda(s)=\iota(\rho)\Lambda_0(s)\).

For the last clause, the existing bottom-triangle witness
activates \texttt{Bottom\allowbreak In}.  This supplies a possibly newly
chosen triangle \(X\) satisfying the same \texttt{Is\allowbreak Bottom\allowbreak Tri}
conditions, and nonzero complex periods \(\Omega,\Omega'\).
Its \texttt{Sums\allowbreak In} certificate already asserts that every
normalized monomial integral \(I/\Omega\), respectively
\(I'/\Omega'\), lies in the actual field image, conditional on
integrability of its corresponding kernel.  The additional
\texttt{h\allowbreak I} hypothesis is quantified over every such triangle,
so it applies to this selected \(X\) and discharges both
integrability guards.  For every allowed \(t\), finite-adelic
point \(g\), and respective monomial indices \(k,k'\), there are
therefore
\(z,z'\in D.L\) with
\[
 \iota(z)=\frac{I(t,k,g)}{\Omega},\qquad
 \iota(z')=\frac{I'(t,k',g)}{\Omega'},\qquad
 I(t,k,g)=\Omega\iota(z),\quad
 I'(t,k',g)=\Omega'\iota(z').
 \tag{LTD.27}
\]
The last two equalities follow by multiplying through the
nonzero periods.  Together with the integrability statements
they are exactly \texttt{Sums\allowbreak In.\allowbreak rational\allowbreak Sums}; the remaining
triangle fields transport by the stated character, place and
rank equalities.  Rationality of these integrals is thus an
explicit part of the \texttt{Sums\allowbreak In} input.  Without these
premises, \texttt{Field\allowbreak OK} is not the whole period formula.
% END LTD8

\subsection*{The reduced producer boundary}

The following distinctions are consequences of the proofs above.
They specify remaining work on the original objects, rather than
introducing replacement models.

\begin{longtable}{p{0.22\textwidth}p{0.34\textwidth}p{0.36\textwidth}}
\textbf{Consumer}&\textbf{Completed implication}&\textbf{Retained producer}\\\hline
Common selected objects&LTD1: every selection is a projection of the
same bundle&The common bundle/compatible design existence theorem.\\
Local installation&LTD2: original forms, theta maps, embeddings,
choices and operators survive the installation&The original
\texttt{Loc\allowbreak All} and the earlier nonplaceholder theta realization.\\
Local rank coverage&LTD3--LTD5: actual \texttt{Nd\allowbreak Spec}/\texttt{NBuilt}
gives \texttt{Locs\allowbreak Cover\allowbreak Tw}; eventual certificates give
\texttt{Link\allowbreak Eta}&Actual local specialization on the common model,
including \(\Psi\), its range and its projection square.\\
Actual row spectrum&LTD6: the specified rational row operators,
eigenvalues, lattices and normalized coefficient action are transported
exactly&Original \texttt{Rat\allowbreak In.\allowbreak span}, \texttt{h\allowbreak L\_\allowbreak int}, row reading,
and the compatible geometric state.\\
Theta evaluation and series&LTD7: existing \texttt{Theta\allowbreak Tie} and
\texttt{Ser\allowbreak At} survive localization with their original scalar&The
original evaluation/series proofs and \texttt{Exp\allowbreak Rel}; the separate
\texttt{EMt\allowbreak Theta} comparison and actual row theta square.\\
Coefficient rationality&LTD8: every inclusion and field transport is
on the specified embedding&Actual toric-ratio/jet/period values,
their prescribed field containment and analytic hypotheses.\\
\end{longtable}

The next reconstruction should construct these retained inputs at
the state used by \texttt{Rows\allowbreak At\allowbreak Start}.  They must then be retained
through the bundle constructor.  In particular, the local coverage step no
longer warrants a separate deep reconstruction after the actual
\texttt{Nd\allowbreak Spec} has been supplied.  The shared local/theta dictionary
has been narrowed by exact transport theorems, while its substantive
existence content remains open in the human reconstruction.


\section{Small-level projection and identification of the base theta map}
\label{sec:small-level-rank}

\status{SLR1--SLR9 reconstruct the original level projection and the finite
rank argument.  The projection is onto on every supplied
\texttt{Bundle}, with an explicit right inverse and no additional size
threshold.  For the original selected row family, the remaining sufficient
input is its base row realization together with its theta comparison at
its base tuple.  From these inputs we prove the literal
\texttt{Small\allowbreak Level\allowbreak Rank\allowbreak Eq} conclusion, including the quantifiers of
\texttt{Closed}.  We do not claim to have constructed those actual row
and theta inputs.  This replaces the former opaque projection question
by an explicit, narrower producer boundary.}

The target is an equality of complex dimensions for the original record:
\[
 \exists x_0\in\mathbb R\ \forall x\ge x_0\ \forall i\in\mathcal F.I_x\
 \ \forall j\in\{1,\ldots,(\mathcal F.D_{x,i}).R.m\},\qquad
       dZ_0(j)=D_0.                                      \tag{SLR.1}
\]
Here $dZ_0(j)=\dim_{\mathbb C}Z_j(\mu_0)$ and
$Z_j(\mu)=\operatorname{im}\theta_{\mu,j}^{*}$, with
$\theta_{\mu,j}:M_t(\mu,j)\to M_j$ the actual map stored in
$D.\texttt{what}_2$.  The superscript $*$ in this expression means
Hilbert-space adjoint.  Matrix conjugate transpose will also be denoted
by $*$; their roles will always be specified.  The separate source
symbol \texttt{l\allowbreak Z0} denotes a coefficient matrix, not the complex
space $Z_j$.

All source references below use the pinned revision.  The accompanying
\texttt{provenance.\allowbreak json} gives exact paths, ranges and hashes.  The principal target references are S01--S03, S27
and S37--S38.  Existing CR, CS and TD arguments remain conditional on
their stated actual data; an abstract rank equality does not supply
that data.

\subsection{The coefficient projection has an explicit right inverse}

% BEGIN SLR1
\paragraph{SLR1: the actual level and good spaces are nested.}
\label{claim:SLR1}
Let $K$ be the coefficient number field, let $\sigma:K\hookrightarrow
\mathbb C$ be its selected embedding, and assume the coefficient
involution satisfies $\sigma(\bar a)=\overline{\sigma(a)}$.  Let a
finite group $G$ act on a finite-dimensional $K$-space.  For a subgroup
$P$ and a character $\psi:P\to K^{\times}$, set
\[
 e=\frac1{|P|}\sum_{h\in P}\psi(h^{-1})[h],\qquad
 C=\sum_{\chi\in X}\frac{\chi(1)}{|G|}
                  \sum_{g\in G}\chi(g^{-1})[g],
 \qquad a=Ce.                                           \tag{SLR.2}
\]
The finite set $X$ consists of characters of absolutely irreducible
representations, as in the source's \texttt{Is\allowbreak Abs\allowbreak Irr\allowbreak Char}.  Each such
character is constant on conjugacy classes: over the algebraic closure
it is a matrix trace, and conjugation replaces its representing matrix
by a similar matrix.  Injectivity of $K\to\overline K$ brings this
equality back to $K$.  Therefore $C$ is central in $K[G]$.

For $u\in P$, the coefficient of $[u]$ in $e^2$ is
\[
 \frac1{|P|^2}\sum_{h\in P}
          \psi(h^{-1})\psi((h^{-1}u)^{-1})
       =\frac1{|P|}\psi(u^{-1});
\]
there are no coefficients outside $P$.  Thus $e^2=e$, and centrality
gives $e(Ce)=Ce$.  Applying the group action shows
\[
           \operatorname{im}\rho(a)\subseteq
                         \operatorname{im}\rho(e).       \tag{SLR.3}
\]
This argument uses neither character orthogonality nor any equality of
integral lattices.  It also works when either space is zero.
\emph{Source: S04, S05, S07.}
% END SLR1

% BEGIN SLR2
\paragraph{SLR2: the matrices defining \texttt{l\allowbreak Prj} are onto.}
\label{claim:SLR2}
Use the original coefficient basis of the finite form space, of
dimension $n$, and let $H$ be its Hermitian matrix.  At $\sigma$ it is
positive definite.  Let $L\in M_{n,\ell}(K)$ have as columns the source's
chosen basis of $\operatorname{im}\rho(e)$, and let
$V\in M_{n,v}(K)$ have as columns its chosen basis of
$\operatorname{im}\rho(a)$.  For the original subspace $Z$ of forms,
let $U\in M_{v,d}(K)$ be the chosen basis matrix of the inverse image
of $Z$ under $V$.  These three matrices have full column rank by their
definitions as subspace basis matrices.  Define
\[
 \begin{aligned}
 L^+&=(L^*HL)^{-1}L^*H,&
 V^+&=(V^*HV)^{-1}V^*H,\\
 H_V&=|P|^{-1}V^*HV,&
 Q&=(U^*H_VU)^{-1}U^*H_V,\\
 E&=L^+VU,& \Pi&=QV^+L.
 \end{aligned}                                             \tag{SLR.4}
\]
These are precisely \texttt{lvl\allowbreak Pr}, \texttt{v\allowbreak Pr}, \texttt{HV},
\texttt{Pq}, \texttt{l\allowbreak E} and \texttt{l\allowbreak Pr} of the source.

For clarity, all inverses in (SLR.4) exist over $K$.  If a matrix $A$
has full column rank, some maximal minor is nonzero in $K$, hence also
at $\sigma$.  Thus $\sigma(A)$ is injective.  Positivity of $\sigma(H)$
makes $\sigma(A)^*\sigma(H)\sigma(A)$ positive definite, hence
nonsingular.  Its determinant is the image of
$\det(A^*HA)$, so this determinant is nonzero in $K$.  The same argument
applies to $H_V$, since $|P|^{-1}>0$ at $\sigma$.  The empty matrix
case has determinant $1$ and causes no exception.

Multiplying out gives $L^+L=1$, $V^+V=1$ and $QU=1$.
By (SLR.3), every column of $V$ lies in the image of $L$;
consequently $LL^+V=V$.  Therefore
\[
             LE=VU,\qquad
             \Pi E=QV^+LL^+VU=QV^+VU=QU=1.              \tag{SLR.5}
\]
The map $y\mapsto Ey$ is an explicit right inverse to $x\mapsto\Pi x$.
This proves surjectivity, not just equality of source and target ranks.

There is also an exact orthogonal projection identity.  Put
$H_L=|P|^{-1}L^*HL$.  From $LE=VU$ we get
\[
 E^*H_LE=U^*H_VU,
 \qquad
 E^*H_L=|P|^{-1}U^*V^*HL.
\]
After cancelling the nonzero factor $|P|^{-1}$ in (SLR.4), this yields
\[
             \Pi=(E^*H_LE)^{-1}E^*H_L.                  \tag{SLR.6}
\]
In particular $E^*H_L(x-E\Pi x)=0$ for every $x$.  Thus $E\Pi$ is
projection onto the columns of $E$ along their $H_L$-orthogonal
complement.  This verifies the normalization by $|P|$ as well as the
right inverse.
\emph{Source: S04--S08.}
% END SLR2

% BEGIN SLR3
\paragraph{SLR3: the precise lattice consequence and its limit.}
\label{claim:SLR3}
For an $\mathcal O_K$-submodule $N\subset K^n$, write
$\operatorname{cut}_A(N)=\{u:Au\in N\}$.  Let $N$ be the actual
full coefficient lattice of forms, and let $N'\subseteq N$ be the
source's smaller lattice.  With (SLR.4), the source's level and lift
lattices are
\[
 N_L=\operatorname{cut}_L(N),\qquad
 T=\Pi N_L,\qquad
 N_Z=\operatorname{cut}_{VU}(N),\qquad
 N'_Z=\operatorname{cut}_{VU}(N').                       \tag{SLR.7}
\]
The exact identity $LE=VU$ gives
$\operatorname{cut}_E(N_L)=N_Z$.  If $z\in N_Z$, then $Ez\in N_L$
and $\Pi Ez=z$.  Hence
\[
                         N'_Z\subseteq N_Z\subseteq T.   \tag{SLR.8}
\]

Here $N_L$, $N_Z$, $N'_Z$ (when $N'$ is full), and $T$ are full
lattices in their respective coefficient spaces.  One elementary
verification is as follows.  If $A:K^d\hookrightarrow K^n$ is
injective, choose a $K$-linear left inverse $B$.  Then
$\operatorname{cut}_A(N)\subseteq B(N)$; the latter is finitely
generated over $\mathcal O_K$, and its submodule is finitely generated
because $\mathcal O_K$ is Noetherian.  A common nonzero denominator
sends each vector $Ae_r$ into $N$, so the cut contains a nonzero scalar
multiple of every coordinate vector.  It spans $K^d$.
Images of finite generating sets generate $\Pi N_L$, and
$K(\Pi N_L)=\Pi(KN_L)=K^d$ by (SLR.5).  These observations prove the
assertions.

The surjectivity in SLR2 is over $K$.  We have not proved that $E$
preserves any independently specified integral lattice, or that
$N_Z=T$.  The map $N_L\to T$ is onto because $T$ is its image by
definition.  No stronger integral assertion enters (SLR.1).
\emph{Source: S05, S08.}
% END SLR3

\subsection{The finite type decomposition and the base dimension}

% BEGIN SLR4
\paragraph{SLR4: the actual types give the needed direct sums.}
\label{claim:SLR4}
We give the finite-dimensional argument because the source's type
spaces are sums of joint \emph{generalized} eigenspaces.  No
semisimplicity assumption is necessary.

Let $T_\alpha$ be a commuting family of endomorphisms of a finite
$K$-space $F$, each split over $K$.  Their generated unital algebra
is a finite-dimensional subspace of $\operatorname{End}_K(F)$.
Choose a basis of that algebra made from finitely many products of
operators; the finitely many operators occurring in those products
already generate the entire algebra.  For one such operator, its
minimal polynomial is a product of coprime powers $(X-a)^r$.
Bezout's identity, applied to these factors, supplies polynomial
idempotents that sum to $1$ and project onto their kernels.  All the
other operators preserve those kernels because they commute.  Applying
this construction successively to the finite generating family gives
\[
          F=\bigoplus_{\lambda\in\Lambda}F_\lambda,
          \qquad |\Lambda|<\infty,                      \tag{SLR.9}
\]
with polynomial projections $p_\lambda$ in the operators.
On a resulting block each generator is its scalar eigenvalue plus a
nilpotent operator.  The nilpotent parts commute; their generated ideal
is nilpotent, since a sufficiently long monomial contains a power at
least the nilpotence exponent of one generator.  It follows that every
operator of the original family has a single generalized eigenvalue
on that block.  Blocks with distinct scalar data are disjoint.  Thus
(SLR.9) is precisely the decomposition by the source's intersections
of maximal generalized eigenspaces.  If $F=0$, it is the empty direct
sum.

Every subspace stable under the operators is stable under the
$p_\lambda$.  If the operators commute with the $G$-action, these
projections also commute with $\rho(Ce)$.  Consequently
$Z\cap\operatorname{im}\rho(Ce)$ is the direct sum of its
intersections with $F_\lambda$, whenever $Z$ is stable.
Pulling back under the injective coefficient map $VU$ preserves this
direct-sum decomposition.

We now check the source's labels.  A type $p$ consists of residual
polynomials $p_1(P_0)$ and inert eigenvalues $p_2(I_0,y)$.  A joint
system $\lambda$ has that type if, at each actual Hecke place $P$,
\[
 H_P(\lambda;X)=p_1(\operatorname{pl}P;X)B_P(X),\qquad
 B_P(X)=\prod_{\nu\in\texttt{bl}}\prod_{c=0}^{t-1}
                      (1-\nu(P)q_P^cX),                 \tag{SLR.10}
\]
and its inert eigenvalues are the specified $p_2$ values.
The polynomial $B_P$ has constant term $1$, so it is nonzero.
Hence the residual polynomial, when it exists, is unique by cancellation
in $K[X]$.  The labels are normalized to $1$ at unused Hecke places
and to $0$ at unused inert places.  The place maps are injective, so
these assignments are well defined.  Two distinct normalized labels
therefore differ at an actual place, and cannot label the same
$\lambda$.  Since there are only finitely many occurring $\lambda$,
there are only finitely many occurring labels.  Thus the finite type
set \texttt{ix\allowbreak Fin3} really is that set; its conditional definition does
not select its fallback empty set here.

At the base tuple, $Z$ is contained in the prescribed block sum by
\texttt{Zs.\allowbreak In\allowbreak Block}.  Each nonzero primary component in
$Z\cap\operatorname{im}\rho(Ce)$ therefore satisfies the divisibility
in (SLR.10), and gives an occurring normalized type.  It also belongs
to the full good space used to define \texttt{ix\allowbreak Fin3}: the image of
$\rho(Ce)$ is contained in the image of $\rho(C)$, and every
$p_\lambda$ preserves both images.  The nonzero component witnesses
that its type occurs there.  The preceding projections consequently
give
\[
        K^{dQ(t_0)}=\bigoplus_{p\in I}Q_p(t_0).           \tag{SLR.11}
\]
For the lift spaces, distinct transformed labels are different at an
actual lift place by \texttt{Types\allowbreak Sep3}.  The same cancellation and
primary-space argument prove independence of their inverse-image
spaces $Z_{k,p}(t_0)$.  Their spanning is exactly the \texttt{types}
field of the supplied base \texttt{Indef\allowbreak Rows} record.  Thus
\[
        K^{dZ(k,t_0)}=\bigoplus_{p\in I}Z_{k,p}(t_0).      \tag{SLR.12}
\]
All hypotheses just used are fields of \texttt{Rows\allowbreak Hold0\_\allowbreak 3}: stable
and block-contained subspaces, some \texttt{fam3.\allowbreak Ok3 s}, the common
base \texttt{Indef\allowbreak Rows}, and type separation.  We have proved the
three pieces of \texttt{WF} needed for dimension transport; its metric
and lattice clauses need not be imported into this argument.
\emph{Source: S09--S14.}
% END SLR4

% BEGIN SLR5
\paragraph{SLR5: the rational lift rank equals the original $D_0$.}
\label{claim:SLR5}
Fix the original selected row family $P$, and a base row witness
$R_0$ with $R_0.\texttt{to\allowbreak Indef\allowbreak Data}=P.S_0$ as in
\texttt{Rows\allowbreak Hold0\_\allowbreak 3}.  Write $D_r=P.S_0.D$ and
$d_k=P.S_0.\texttt{d\allowbreak W}(k)$.  The \texttt{l1c} field gives bijections
\[
 \iota_0:K^{dQ(t_0)}\longrightarrow K^{D_r},\qquad
 \vartheta_{k,0}:K^{dZ(k,t_0)}\longrightarrow K^{d_k}.
                                                               \tag{SLR.13}
\]
The \texttt{l1a} field says these maps carry the summands in
(SLR.11)--(SLR.12) onto $E_p$ and $W_{k,p}$, respectively.
Therefore these latter spaces are direct sums of their indicated
summands.  The \texttt{l1d} field says
$\dim_KW_{k,p}=\dim_KE_p$ for every $p$.  Taking a basis of each
summand and concatenating these bases proves
\[
 d_k=\sum_{p\in I}\dim_KW_{k,p}
       =\sum_{p\in I}\dim_KE_p=D_r.                     \tag{SLR.14}
\]
There is no exchange of an infinite sum or choice of unrelated type
sets in this equality.

The equality $D_r=D_0$ follows directly from data already in the
original \texttt{Row\allowbreak Side\allowbreak Family}.  Its $P.I.\texttt{e\allowbreak RR0}$ is an
isomorphism
\[
                 K^{D_r}\simeq_K (\texttt{ex.\allowbreak p0.\allowbreak s.\allowbreak RR}).V.
\]
The latter rational space has a specified basis $b_a$, indexed by the
finite set $A_0$.  The field \texttt{ex.\allowbreak p0.\allowbreak s.\allowbreak j\allowbreak EE\_\allowbreak basis} states that
the images $j_{EE}(b_a)$ form a complex basis of the original
$EE_0$.  Finally, \texttt{onb0} is a complex orthonormal basis of
that same $EE_0$, indexed by $\{1,\ldots,D_0\}$.  Counting these
bases gives
\[
          D_r=\dim_K(\texttt{ex.\allowbreak p0.\allowbreak s.\allowbreak RR}).V
              =|A_0|=\dim_{\mathbb C}EE_0=D_0.           \tag{SLR.15}
\]
This uses no isometry between newly chosen spaces.  In particular the
separate \texttt{Rows\allowbreak Hold\allowbreak Lam3} predicate is unnecessary for this
numerical comparison.

The \texttt{li} and \texttt{span} fields of
$P.\texttt{tied.\allowbreak div.\allowbreak th}$ assert that the vectors
$j_{W,k}(e_r)$ form a complex basis of the original $M_{k+1}$,
where $r=1,\ldots,d_k$.  Thus (SLR.14)--(SLR.15) imply
\[
                      \dim_{\mathbb C}M_{k+1}=D_0.       \tag{SLR.16}
\]
\emph{Source: S13--S19.}
% END SLR5

\subsection{The exact theta comparison needed for the rank conclusion}

% BEGIN SLR6
\paragraph{SLR6: a base-tuple square makes the actual theta map onto.}
\label{claim:SLR6}
For this paragraph the additional input is the following equality on
the same selected $P$, at its single base tuple $t_0=P.A_3.A.\mu_0$:
\[
 \theta_{P.I.\mathrm{tup}(t_0),\,k+1}
       \bigl(j_{Mt}(e_{Mt,k,t_0}(x))\bigr)
     =j_{W,k}\bigl(\vartheta_{k,0}(\Pi_{k,t_0}x)\bigr)
       \quad\text{for every }k,x.                       \tag{SLR.17}
\]
Here $\Pi_{k,t_0}$ is the original \texttt{l\allowbreak Prj} matrix,
$e_{Mt,k,t_0}=P.I.\texttt{e\allowbreak Mt}(k,t_0)$, and the maps $j_{Mt}$ and
$j_{W,k}$ are the actual embeddings in the original record.  Equation
(SLR.17) is the base-tuple restriction of \texttt{th\allowbreak Theta}, not a
consequence of the basis assertions \texttt{li/\allowbreak span} alone.

For each coordinate vector $e_r\in K^{d_k}$ choose its unique
preimage $z_r=\vartheta_{k,0}^{-1}e_r$ under (SLR.13), and put
$x_r=E_{k,t_0}z_r$.  By SLR2,
$\Pi_{k,t_0}x_r=z_r$.  Substitution in (SLR.17) shows that the range
of the actual theta map contains every $j_{W,k}(e_r)$.  These vectors
are a complex basis by SLR5; its range is a complex subspace, so the
map is onto.  In fact (SLR.17) only needs to be checked at this finite
list of $x_r$ for each $k$.  We retain its all-$x$ base-tuple version
as a convenient sufficient contract.

For completeness, if a linear map $A:V\to W$ between finite
complex inner-product spaces is onto and $A^*w=0$, choose $v$ with
$Av=w$.  The adjoint identity gives
$\langle w,w\rangle=\langle Av,w\rangle
=\langle v,A^*w\rangle=0$, hence $w=0$.
Thus $A^*:W\to V$ is injective and its image has dimension
$\dim_{\mathbb C}W$.  Apply this to the actual theta map and use
(SLR.16).  Since $P.\texttt{tied.\allowbreak div.\allowbreak tup.\allowbreak mu0}$ identifies
$P.I.\mathrm{tup}(t_0)$ with the original $\mu_0$, we obtain
\[
       \dim_{\mathbb C}Z_{k+1}(\mu_0)=D_0.               \tag{SLR.18}
\]
If there are no lifts, the requested family of equalities is empty.
Otherwise every $1\le j\le m$ is uniquely $k+1$ with
$k=j-1\in\{0,\ldots,m-1\}$, so (SLR.18) covers exactly the source's
\texttt{Size\allowbreak Data}\textsubscript{2}\texttt{.\allowbreak lifts}.
\emph{Source: S01--S03, S17--S22.}
% END SLR6

% BEGIN SLR7
\paragraph{SLR7: what is still required to identify the actual square.}
\label{claim:SLR7}
The source separates \texttt{Div\allowbreak Def3}, already inside $P.\texttt{tied}$,
from \texttt{Div\allowbreak Rows3}.  Its \texttt{Th\allowbreak Def3} contains linear
independence, complex spanning and a lattice tie; it contains no
assertion (SLR.17).  The latter is a field of \texttt{Div\allowbreak Rows3}, and
its used-tuple version is a field of \texttt{Div\allowbreak Rows3G}.  Since
\[
 \texttt{Used0}(D,\mu)\ \Longleftrightarrow\
       \mu=\mu_0\ \text{or}\ \mu=\texttt{mu\allowbreak At}(w)\text{ for some }w,
                                                               \tag{SLR.19}
\]
the base tuple is used by its first alternative.  Therefore a supplied
\texttt{Div\allowbreak Rows3G} for this $P$ proves (SLR.17) directly.  There is no
need to prove \texttt{Tup\allowbreak All\allowbreak Used}, or to upgrade every other tuple to
\texttt{Div\allowbreak Rows3}, for the argument above.

A more detailed reading of the actual theta proof reveals a narrower
analytic boundary.  At an actual realization let $\Theta$ be its
analytic theta map on level forms and let $\epsilon_W:M_{k+1}\to H_W$
be the stored injective reading into the analytic model.  Write
$J(x)$ for the actual level form represented by the rational
coordinates $x$, and $g(y)$ for the actual good form with coordinates
$y$.  The needed comparisons are
\[
 \begin{split}
 \epsilon_W\theta\bigl(j_{Mt}(e_{Mt}x)\bigr)&=\Theta J(x),\\
 \epsilon_W j_W(y)&=g(y),\\
 \Theta J(Ey)&=g(\vartheta y),\\
 E^*H_Lx=0&\Longrightarrow\Theta J(x)=0.
 \end{split}                                             \tag{SLR.20}
\]
These formulas use the same realization, coefficient embedding and
level character as the matrices.  The source's phase is $1$ by its
\texttt{Pins.\allowbreak phs} condition; its possible conjugate scalar model is
transported back before the equalities in (SLR.20) are compared.

Once (SLR.20) is available the remaining implication is elementary.
By (SLR.6), $x-E\Pi x$ is orthogonal to every column of $E$, so the
last line of (SLR.20), together with additivity, gives
$\Theta J(x)=\Theta J(E\Pi x)=g(\vartheta\Pi x)$.
The first two lines and injectivity of $\epsilon_W$ then give
(SLR.17).  This proves the projection-factorization step completely.

We do not label all of (SLR.20) elementary.  In the pinned proof the
first line is the Hilbert-model/lowest-$K$-type reading, whose proof
uses restriction to lowest $K$-types and uniqueness there.  The last
line uses the actual rational adjoint-image and level compatibility.
The third line uses the actual rational theta matrix.  TD proves useful
finite and rational consequences from explicit analytic inputs, but
does not by itself supply all these identifications on the selected
record.  This is a precise remaining human reconstruction obligation.
The proof of \texttt{th\allowbreak Theta} is obtained before the extra pairing
and column fields of full \texttt{Div\allowbreak Rows3}; those extra fields are
not hypotheses of the rank implication established here.
\emph{Source: S19--S23, S30--S34.}
% END SLR7

\subsection{Specialization to the selected family and closure}

% BEGIN SLR8
\paragraph{SLR8: the projection certificate is pointwise on the literal bundle.}
\label{claim:SLR8}
Let $b$ be the original \texttt{Bundle} at $(D,O,c)$ and put $P=b.P$.
Its column-compatibility field $b.C.\texttt{ok1}$ includes
$P.A_3.\texttt{fam3.\allowbreak Ok3}(1)$.  Consequently the lift block at any
$k,t$ has: the stated coefficient involution at its selected embedding,
the irreducible character data of SLR1, and the positive Hermitian
matrix used in SLR2.  The exact lift level is
\[
 \begin{aligned}
 B&=P.A_3.A.\mathrm{fam}.L(k,t),\\
 \ell&=P.A_3.A.\mathrm{fam}.\mathrm{lvlA}(P.P_0,P.\psi_0)
          =(X,P.P_0,P.\psi_0^{-1}),\\
 Z&=P.Z_s.Z_l(k,t).
 \end{aligned}                                           \tag{SLR.21}
\]
By definition \texttt{l\allowbreak Z0}$=E$ and \texttt{l\allowbreak Prj}$=\Pi$ for these
objects.  Notice the inverse character on the lift level; replacing
it by the base-level character would change the objects.
SLR2 proves the actual projection surjective, in particular at $t_0$.
The appearance of \texttt{Forms\allowbreak Tie} in one source wrapper is not an
extra need for this proof: that wrapper uses only the family
\texttt{Ok} data in its proof body.

A raw bundle selector $\sigma$ gives
$\operatorname{selP}(\mathcal F,h_c,x,h_x,h,i).1
=(\sigma\mathcal F h_c x h_x i).P$ and likewise obtains the columns
and realization from the corresponding $C$ and $B$ fields of this
same bundle.  Its selection threshold is $\mathcal F.x_1$.
Thus all occurrences of $P$, $\Pi$, $E$, $j_W$, and $e_{Mt}$ above
refer to the same literal selection.  The realization used by
\texttt{sel\allowbreak IOf} is $(\sigma\mathcal F h_c x h_x i).B.Q$.
There is no second existential choice to reconcile.

For the original \texttt{H\_\allowbreak z\allowbreak L} predicate choose its additional
threshold $x_r=0$.  Choose $\varepsilon_1=1$ and all four components
of its \texttt{Closed} bound $K$ equal to zero.  For every allowed
$\varepsilon,\mathcal F,CH$ the pointwise argument just given proves
the required surjectivity; the largeness premise is unused.  This is
the exact closure of the projection obligation.  No period
admissibility input or new uniform arithmetic bound is required.
\emph{Source: S06, S07, S24--S29, S35--S36.}
% END SLR8

% BEGIN SLR9
\paragraph{SLR9: a narrower sufficient producer cut closes the original rank predicate.}
\label{claim:SLR9}
Fix a class $\mathcal A$ and the original compatible selectors.  It
suffices that, for each allowed family $\mathcal F\in\mathcal A$,
there is a real number $a_{\mathcal F}$ such that every selected
$P$ at $x\ge a_{\mathcal F}$ with $x>\mathcal F.x_1$ has both:
\begin{enumerate}
\item \texttt{Rows\allowbreak Hold0\_\allowbreak 3} for this $P$;
\item the base-tuple equality (SLR.17) for this same $P$.
\end{enumerate}
Here $a_{\mathcal F}$ is also taken at least the selector threshold;
if the two inputs initially have different eventual thresholds,
replace them by their maximum.  This is one common-family contract,
not two assertions that possibly use different row realizations.
The source obtains these inputs from the stronger hypotheses
\texttt{Prim\allowbreak All} and pointwise \texttt{Div\allowbreak Rows3G}, with their original
common selector.  SLR7 proves the restriction of the latter to the
base tuple.  On the row side, the primitive-row construction that turns
\texttt{Prim\allowbreak All} into an \texttt{Indef\allowbreak Rows} witness remains part of
the producer boundary here.  Once that witness has been supplied, its
transport uses the realization's actual equalities of $Z_s$ and $S_0$:
substitute the first equality in the type of the witness, then the
second in its underlying data.  The resulting witness has exactly
$P.S_0$ as underlying data.  No isomorphism of unrelated abstract row
spaces is substituted for these equalities.

Choose
\[
              x_0=\max\{a_{\mathcal F},\mathcal F.x_1+1\}.
                                                               \tag{SLR.22}
\]
If $x\ge x_0$, the strict condition $x>\mathcal F.x_1$ holds.
Every index $i\in\mathcal F.I_x$ has the same compatible objects
as in SLR8, the two producer inputs hold, and SLR4--SLR6 give
(SLR.18) for every $j\in\{1,\ldots,m\}$.  This proves (SLR.1) for
the original family, with one threshold for all $i,j$.

It also proves the source's
$\texttt{Closed}(D3\ 3\ \mathcal A)\,
(\mathcal F,CH\mapsto\texttt{Small\allowbreak Level\allowbreak Rank\allowbreak Eq}\,\mathcal F)$.
Indeed, choose $\varepsilon_1=1$ and the zero bound $K$ before
$\mathcal F,CH$ are supplied.  Membership in $D3\ 3\ \mathcal A$
forgets to membership in $\mathcal A$; apply the preceding argument.
The only family-dependent threshold is $x_0$ inside the existential
quantifier of \texttt{Small\allowbreak Level\allowbreak Rank\allowbreak Eq}, exactly where the definition
permits it.  No threshold is moved in front of a family quantifier.

\emph{Audit consequence.}  The logical implication from this common
producer cut to the original \texttt{GLOBAL\_\allowbreak SMALL\_\allowbreak RANK} target is
now completely reconstructed.  Its projection premise is discharged
by SLR1--SLR3 and SLR8, and its finite rank and threshold arguments by
SLR4--SLR6 and this paragraph.  The actual base row and theta producer
remains open in the human reconstruction.  Its arithmetic and analytic
construction is shared with the actual row/period program; such a
shared work relationship is not itself a logical proof of either
input.  The construction of the bundle/family also remains an input.
In particular, this chapter supplies the rank input to
\texttt{top\allowbreak Des\allowbreak Win} conditionally on the displayed actual producer;
it does not close every other input of that theorem.
\emph{Source: S01, S15, S20--S29, S35--S39.}
% END SLR9


\section{The original second layer and its integral ramified model}
\label{sec:second-layer-global}

\status{RSL1--RSL6 identify the exact mass, volume and finite-character
inputs on the original records.  RSL7--RSL10 prove the integral and
isometric transport from the actual level-zero forms to the augmentation
model; this supplies the previously retained transport hypothesis in
LC7.  RSL11--RSL12 give a smaller sufficient condition for actual
\texttt{Nu\allowbreak Full\allowbreak F} and \texttt{Mu\allowbreak Rank}.  RSL13 records the exact remaining
global slope, operator and assembly obligations.  A theorem about a
supplied second layer is not being counted as construction of the common
family.  All references to an unresolved obligation concern this human
reconstruction, under the standing assumption that the pinned formal
proof is correct.}

Throughout, $D$ is the original \texttt{D7} record, $E/F$ its CM
extension, $K=D.\texttt{what}_2.\texttt{LL}$ its coefficient field, and
$s_K:K\hookrightarrow\mathbb C$ its specified embedding.  The
normalizations of mass, Haar measure and squared Hermitian slope are
those already fixed in MF, GV and LC7.  In particular, a new numerical
function called a volume, an abstractly isomorphic form space, or a
finite-index sublattice cannot replace a field of $D$ in the arguments
below.

\subsection{Mass and volume on the very same Hermitian spaces}

% BEGIN RSL1
\paragraph{RSL1: the literal genus-mass input in every rank.}
\label{claim:RSL1}
For every original $D$ and every integer $n\ge0$, the following
statement is established:
\[
 \begin{gathered}
 W:\texttt{HSpace}(E),\quad W.\texttt{Sig\allowbreak OK},\quad
 W.\mathrm{sig}(\tau)=(W.\mathrm{rank},0)\quad\hbox{for every }\tau,
 \quad W.\mathrm{rank}=n\quad\Longrightarrow\\
 \operatorname{mass}(W.h_T,W.L_0)=\operatorname{massI}(D,n).
 \end{gathered}                                                   \tag{RSL.1}
\]
This is exactly \texttt{Genus\allowbreak Mass\allowbreak At D n}; in particular it supplies
the universally quantified input \texttt{h\allowbreak G} in \texttt{link\_\allowbreak gl}.

Here is the full transport, including the identities of the lattices.
The hypotheses \texttt{Sig\allowbreak OK} and the CM-unramified field of $D$
place the actual pair $(W.h_T,W.L_0)$ in the scope of GV8--GV11.
GV11 identifies the normalized adelic quotient with the genus of this
very lattice: a finite adelic representative $g$ gives
$W\cap g_f\widehat {W.L_0}$, and its rational stabilizer is its full
lattice automorphism group.  Consequently its quotient measure is
the sum of the reciprocal orders of those automorphism groups, which
is precisely the left side of (RSL.1).  It is not the mass of a
replacement lattice.  For $n>0$, GV9--GV10 give that quotient the
value $M_n$ of (GV.21).  Apply the identical calculation to the
standard positive pair $(E^n,I,\mathcal O_E^n)$.  The definition of
\texttt{mass\allowbreak I} is its actual genus mass, so it too is $M_n$.
For $n=0$, both spaces and lattices are zero, their automorphism group
is the one-element group, and both masses are one.  This proves
(RSL.1) in every rank.

The retained published input is exactly GV-TAM, with the measure
normalization already checked in GV9--GV10.  The use of that theorem
does not require a global isometry between $W$ and the standard
positive space.  Nor does (RSL.1) assume \texttt{Vol\allowbreak Tie} or change
$D.\mathrm{vol}$: this is a statement about the actual two genus
masses.
% END RSL1

% BEGIN RSL2
\paragraph{RSL2: the literal reciprocal volume ratio.}
\label{claim:RSL2}
For every original $D$, GV14--GV15 give
\[
 \texttt{Vol\allowbreak Tie}(D)\ \Longrightarrow\ \texttt{Volume\allowbreak Ratio}_2(D),
                                                                  \tag{RSL.2}
\]
which is exactly \texttt{Mass\allowbreak Ratio\allowbreak Under\allowbreak Volume D}, the input
\texttt{h\allowbreak U} of \texttt{link\_\allowbreak gl}.  We check both the witness and the
normalization rather than identifying two statements by their names.

Let $W$ have constant signature $(P,Q)$ and rank $n\ge3$, and let
$\epsilon\in\{1,-1\}$ be a sign represented by that signature.
GV14, with its retained Hasse--Minkowski input MF-HM, constructs a
rational $x\in W$ with $h_T(x,x)=\epsilon$.  It constructs a
self-dual lattice on the actual rational complement $W'=x^\perp$
and an isometric inclusion $\iota:W'\hookrightarrow W$ with range
$\ker h_T(x,-)$.  The signatures of $W'$ are $(P-1,Q)$ or
$(P,Q-1)$ according to the chosen sign.  These are exactly the
vector, complement, inclusion, range and signature witnesses that
occur in \texttt{Volume\allowbreak Ratio}\textsubscript{2}; none is reselected
when we pass to numerical volumes.

Choose normalized Haar measures and fundamental domains for $W$ and
$W'$ using the existence clauses of the original \texttt{Vol\allowbreak Tie}.
GV15 applies to every such pair of measures and domains.  The other
clauses of \texttt{Vol\allowbreak Tie} identify their measures with
$\operatorname{ofReal}(D.\mathrm{vol}(W))$ and
$\operatorname{ofReal}(D.\mathrm{vol}(W'))$.  The original
$D.\texttt{vol\_\allowbreak pos}$ fields justify taking real values.  Thus, with
$v=D.\mathrm{host}.v>0$, $f=[F:\mathbb Q]$ and
$L_n=\mathrm{lValEta}(E,n)>0$, GV15 yields
\[
 \frac{D.\mathrm{vol}(W)}{D.\mathrm{vol}(W')}
 =\frac{\big((n-1)!(v/\pi)^n\big)^f}{\sqrt{(2v)^f}}L_n.
                                                                  \tag{RSL.3}
\]
Taking reciprocals, using positive real powers, gives
\[
 \boxed{\quad
 \frac{D.\mathrm{vol}(W')}{D.\mathrm{vol}(W)}
   =\left(\frac{(2\pi)^n}{(n-1)!}\right)^f
        \big((2v)^f\big)^{1/2-n}L_n^{-1}.
 \quad}                                                          \tag{RSL.4}
\]
Indeed the ratio of the right sides of (RSL.4) and the reciprocal
of (RSL.3) is $2^{nf}(2v)^{-nf}v^{nf}=1$.
This is the precise coefficient in the source, including the
orientation of the ratio and its half power.  The proof works for
every $D$ and every admissible $W$, so the universal quantifier in
\texttt{h\allowbreak U} has not been weakened to a chosen family member.
Its retained inputs are GV-TAM and MF-HM, both already explicit in
GV14--GV15.
% END RSL2

% BEGIN RSL3
\paragraph{RSL3: real-place and base-volume transport without retieing.}
\label{claim:RSL3}
For the same $D$,
\[
 \texttt{Real\allowbreak Tied}(D)\Longrightarrow\texttt{Vol\allowbreak Tie}(D)
 \Longrightarrow
 D.\mathrm{vol}(W)=\operatorname{massI}(D,W.\mathrm{rank})
                                                                  \tag{RSL.5}
\]
for every supplied constant-signature $W$ satisfying
\texttt{Sig\allowbreak OK}.  The first implication simply specializes the
\texttt{vol} field of \texttt{Real\allowbreak Tied} to the constant signature
function.  Its measure and fundamental-domain quantifiers are
unchanged.  For the second, GV9--GV10 give the actual normalized
domain measure $M_n$, and RSL1 applied to the standard positive
space gives $M_n=\operatorname{massI}(D,n)$.  The actual
\texttt{Vol\allowbreak Tie} identifies that measure with
$\operatorname{ofReal}(D.\mathrm{vol}(W))$; positivity again removes
\texttt{of\allowbreak Real}.  This proves the equality on the same $D$.

Apply this to the original $D.\mathrm{sp}.Y(0)$.
Its stored signature and rank identities give rank $D.R.r$, so
(RSL.5) is the literal \texttt{Base\allowbreak Volume\allowbreak Equals\allowbreak Mass D}.  Alternatively
one can use RSL1 in rank two, the actual finite class decomposition
of GV11, and the successive complements of RSL2, exactly as in
\texttt{volume\_\allowbreak eq\_\allowbreak mass\_\allowbreak rank\_\allowbreak two}.  Both routes retain the
original base space and its volume.

If a family has an eventual \texttt{Real\allowbreak Tied} witness at threshold
$a$, this same $a$ supplies its eventual \texttt{Vol\allowbreak Tie} and
base-volume identities, with the already stipulated condition
$x>\mathcal F.x_1$ wherever the family predicate requests it.
The standard adjacent genus-mass ratio is separately supplied by
MF19.  Thus the fields \texttt{h2b} and \texttt{h2v} of
\texttt{Fam\allowbreak GLmin} have no remaining same-object transport gap.
This does not produce \texttt{Real\allowbreak Tied} for an arbitrary record:
DC20's normalization hypotheses, or an actual selected
$T.\mathrm{theta}.\mathrm{real}$, must still supply that record's
real-place witness.
% END RSL3

\subsection{Finite characters on the original coefficient field}

% BEGIN RSL4
\paragraph{RSL4: the required roots already belong to the chosen field.}
\label{claim:RSL4}
Fix the actual true-column data
\[
 I:\texttt{TC}(B,X),\quad I.\texttt{Ok},\quad
 l:I.\texttt{Laws},\quad Z:\texttt{Later}(B,I.\mathrm{cols}(l)),
 \qquad D=I.\mathrm{rec7}(l,Z).
                                                                  \tag{RSL.6}
\]
For $i=0,1$, put $p=g_i=D.\mathrm{host}.g(i)$.
The existing nontrivial additive character in $I.\texttt{Ok}$
supplies a primitive $p$th root of unity in the original $K$.
Here is the argument, which uses no coefficient-field enlargement.
The selected places over $g_i$ form a nonempty set: their cardinality
is $f=2\,\mathrm{ord}>0$.  At one of these places the original
residue map sends $g_i=\pi_i\overline\pi_i$ to zero.  As $g_i$ is
prime, the residue field has characteristic $p$.  The character
\[
                    \psi:(k_w,+)\longrightarrow K^\times
\]
stored in the character record is nontrivial, so choose $a$ with
$\zeta=\psi(a)\ne1$.  Additivity gives
$\zeta^p=\psi(pa)=\psi(0)=1$.  Its order is a nontrivial divisor
of the prime $p$, hence is exactly $p$.

Under the definition of \texttt{rec7}, this is literally the field
$D.\texttt{what}_2.\texttt{LL}$ and literally the same prime $g_i$.
The source also constructs roots of order $g_i-1$, but none is
needed in the zero-level argument below: there are no multiplicative
blocks at that level.  Keeping this distinction avoids imposing an
unnecessary cyclotomic extension or a stronger character condition.
% END RSL4

% BEGIN RSL5
\paragraph{RSL5: one explicit exponent works for any supplied symmetric forms.}
\label{claim:RSL5}
Suppose the original coefficient field contains primitive
$g_0$th and $g_1$th roots.  For every word $u$ and every actual
\texttt{Forms\allowbreak Tied\allowbreak Sym} object $T$ on $D,u$, its level-zero character
satisfies
\[
 \chi_T(\kappa)^{g_0g_1}=1,
 \qquad \chi_T(\kappa)\in s_K(K),
 \qquad \kappa\in T.\mathrm{cls}(0).\mathrm{lev}.U.
                                                                  \tag{RSL.7}
\]
In particular this proves the exact predicate \texttt{Chi\allowbreak Fin0 T},
with the explicit positive exponent $g_0g_1$.

To prove it, use the \texttt{glev0} field of this very $T$.
It expresses $\chi_T(\kappa)$ as the finite product, over the two
auxiliary primes and their selected places, of the inverses of the
local unipotent characters.  At level zero their multiplicative
block count is zero.  Each local term is therefore the value of
an additive character $\psi$ in residue characteristic $g_i$.
For every residue element $a$,
$\psi(a)^{g_i}=\psi(g_i a)=1$; the same holds for its inverse.
Each local term consequently has order dividing $g_0g_1$, and so
does their finite product.

For rationality, let $\zeta_i\in K$ be a primitive $g_i$th root.
The numbers $1,s_K(\zeta_i),\ldots,s_K(\zeta_i)^{g_i-1}$ are
distinct and are roots of $X^{g_i}-1$.  A polynomial of degree
$g_i$ has at most $g_i$ roots, so every complex $g_i$th root is
one of these powers.  Thus every local character value and its
inverse lie in $s_K(K)$; multiply their preimages in $K$ to obtain
the required preimage of $\chi_T(\kappa)$.  This proves both
clauses of (RSL.7), with no restriction on the particular supplied
$T$ beyond its original fields.

RSL4 gives the needed roots for (RSL.6).  TG12 supplies
$C:\texttt{Cols\allowbreak Tied\allowbreak Sym}(I.\mathrm{rec7}(l,Z))$, with its original
standard-column and standard-level ties, from the supplied
$I.\texttt{Ok}$, pins, laws and later-stage record; DR3 supplies
the initial realizer used there.  Applying (RSL.7) to each
$C.\mathrm{forms}(u)$ proves the complete finite-character
endpoint in the source.
More strongly, (RSL.7) applies to any previously selected
$C$ on this same $D$.  Thus the older FC19 terminal-geometry
boundary is superseded in this scope by TG12, and we need not
replace a selector's columns by newly chosen ones.
% END RSL5

% BEGIN RSL6
\paragraph{RSL6: the same-selector and family quantifiers.}
\label{claim:RSL6}
Fix an original family $\mathcal F$ and a threshold selection
$s:\texttt{TSel\allowbreak F}(\mathcal F)$, with threshold $a_s$.
For $x>\mathcal F.x_1$, $a_s\le x$ and $i\in\mathcal F.I_x$,
write
\[
       T_{x,i}=s(x,x>\mathcal F.x_1,a_s\le x,i).T.
\]
Assume that at a threshold $a_\chi$ the roots of RSL4 have been
identified in the original field of $\mathcal F.D(x,i)$, for
every such $x,i$.  For example this follows when those same
records are the true-column \texttt{rec7} records of (RSL.6),
with an equality of records or an explicit transport of the
entire dependent tuple.  The conclusion is
\[
             \texttt{Real\allowbreak Place\allowbreak Second\allowbreak Layer\allowbreak Conditions}(\mathcal F).
                                                                  \tag{RSL.8}
\]
Take
\[
             a=\max\{a_\chi,a_s,\mathcal F.x_1+1\}.
                                                                  \tag{RSL.9}
\]
For every real $x\ge a$ and every $i\in\mathcal F.I_x$, retain
the original pair
\[
       T_{x,i}.\mathrm{theta}.\mathrm{real},\qquad
       T_{x,i}.\mathrm{cols}.
                                                                  \tag{RSL.10}
\]
RSL5 supplies \texttt{Chi\allowbreak Fin0} for every word of these very
columns.  Equations (RSL.9)--(RSL.10) are exactly the witnesses
of (RSL.8), whose threshold precedes both $i$ and $u$.
The choice of the exponent $g_0g_1$ may depend on the family
member, as the predicate allows; no uniform field, prime or
character exponent is being assumed across members.

There is also a direct bundle formulation.  In
\texttt{Weil\allowbreak Unitary\allowbreak Entries.\allowbreak pt\allowbreak E\_\allowbreak of\_\allowbreak wit}, a witness $q$ supplies
the real-place field $q.x.\mathrm{real}$ and the true-column
data $q.w.I,q.w.\mathrm{ok},q.w.\mathrm{pins},q.w.l,q.Z$.
RSL5 supplies its columns on precisely $q.w.I.\mathrm{rec7}$.
The linked-bundle equality transports the whole dependent tuple
$(D,O,\mathrm{Bundle})$, so it transports their conjunction to
the original $\mathcal F.D(x,i)$.  If these witnesses exist for
every $x>\mathcal F.x_1$ and every $i$, threshold
$\mathcal F.x_1+1$ suffices.  An equality of ranks, an equality
of abstract fields, or separately chosen real and column models
would not justify this step.

Finally, applying RSL3 member by member gives eventual
\texttt{Vol\allowbreak Tie}.  Applying RSL7 below to
$(T_{x,i}.\mathrm{cols}.\mathrm{forms}(u)).\mathrm{to0}$ gives
the actual trivial-stabilizer condition.  Together with RSL5
this supplies \texttt{t\allowbreak CVTies} and
\texttt{Fam\allowbreak U Concrete\allowbreak Forms\allowbreak Level\allowbreak Zero}; threshold $a$ already
forces $x>\mathcal F.x_1$.  Thus no additional eventual
stabilizer theorem is needed on this selected route.  For a
class of families, the outer quantifiers are those of
\texttt{Closed}: there is $\varepsilon_1>0$ such that for
$0<\varepsilon\le\varepsilon_1$ there is a transported bound
$\mathcal K$ valid for every $\mathcal F,C_H$ in the class
with $\mathcal F.\mathrm{Large}(\mathcal K,C_H)$.
The maximum in (RSL.9) is taken inside those family quantifiers,
and before every member index and word.  Thus the proof lifts
an already closed supplier without exchanging its quantifiers.
% END RSL6

\subsection{The actual integral and metrized induction model}

For this subsection fix $D,u$ and an actual
$T_0:\texttt{Forms0}(D,u)$.  In the preceding construction $T_0$
is the literal \texttt{to0} of the supplied symmetric forms.  Write
\[
 \begin{gathered}
 B=(D.\mathrm{fm}(u)).\mathrm{lat}(0),\qquad
 \mathscr C=(D.\mathrm{mem}(u)).C.\mathrm{cls}(0),\\
 \mathscr G=\prod_{v\in A.v_T} k_v^\times,
 \qquad A=(D.\mathrm{mem}(u)).A.
 \end{gathered}
                                                                  \tag{RSL.11}
\]
Thus $B$ is a Hermitian $\mathcal O_K$-lattice in its original
$K$-space $B.V$, with its original basis and metrics.  Each $k_v$
has cardinality $T=A.T$ and $|A.v_T|=f=D.R.f$.  The form-valued
realization $j:B.V\longrightarrow\{
\mathscr C\to\hbox{polynomial functions}\}$ is the original
semilinear embedding $(D.\mathrm{fm}(u)).\mathrm{emb}(0)$.
The fields \texttt{val0.\allowbreak rat}, \texttt{val0.\allowbreak lat} and
\texttt{norm0} in $T_0$ identify rational forms, the integral
lattice and the metrics, respectively.  None of these is inferred
merely from a dimension count.

% BEGIN RSL7
\paragraph{RSL7: two auxiliary primes kill every actual stabilizer.}
\label{claim:RSL7}
Every stabilizer of an original class $c\in\mathscr C$ is trivial.
This is the exact predicate \texttt{adm\allowbreak S D u}; it follows already
from $T_0$, without a further period-admissibility or asymptotic
size hypothesis.

First consider a rational unitary matrix $\gamma$ such that
$y^{-1}\gamma y\in U$, where $y=T_0.\mathrm{cls0}.y(c)$ and
$U=T_0.\mathrm{cls0}.\mathrm{lev}.U$.  The level conditions imply
that this conjugate preserves the prescribed local lattices at
every finite place.  Their intersection is a full global
$\mathcal O_E$-lattice by LP1, and $\gamma$ preserves it.  The
level-zero Hermitian space in this construction is totally
positive definite at all archimedean embeddings.  Hence its
rational unitary lattice automorphism group is finite: after
restriction of scalars, the positive trace form bounds the images
of a fixed integral basis; a discrete lattice has only finitely
many points in each of these bounded balls.  In particular
$\gamma$ has finite order.

At either auxiliary prime $p=g_i$, the original
\texttt{glev0} conditions say that in the supplied split frame
the local matrix is integral and has upper-unipotent reduction.
Here the level-zero rank is $D.R.k$; the exact
\texttt{rank\_\allowbreak eq\_\allowbreak first\_\allowbreak level} identity gives this, so no
lower block with unrestricted eigenvalues remains.  The kernel
of reduction in $\mathrm{GL}_k$ over a complete discrete
valuation ring has no nonidentity torsion of order prime to
$p$.  To check this directly, suppose $M^n=1$, $p\nmid n$ and
$M=1+\varpi^a C$ with $a\ge1$ minimal and $C\not\equiv0$
modulo $\varpi$.  The binomial expansion modulo
$\varpi^{a+1}$ is
\[
                    M^n\equiv1+n\varpi^a C,
\]
which is not one, a contradiction.  If $M=1$ to every order,
separatedness of the complete ring gives $M=1$ directly.
An upper-unipotent matrix over a finite field of characteristic
$p$ has $p$-power order, since sufficiently large $p$-powers
kill its nilpotent part.  Combining these two facts shows that
every finite-order lift of such a matrix has $p$-power order.

Apply this at both $g_0$ and $g_1$.  The rational matrix maps
injectively into each split local factor, and conjugation does
not change its order.  That order is therefore both a
$g_0$-power and a $g_1$-power.  The distinctness field of the
original host gives $g_0\ne g_1$, so the order is one.
Finally \texttt{cls0.\allowbreak stab\allowbreak I} injects the abstract stabilizer
stored in $D$ into precisely this rational group, with exactly
this range.  Thus it too is the one-element group.  The class
weights in the original norm formula are consequently all one.
% END RSL7

% BEGIN RSL8
\paragraph{RSL8: actual diamonds and a free action on the actual classes.}
\label{claim:RSL8}
The group $\mathscr G$ in (RSL.11) acts freely on $\mathscr C$
by the original diamond dilations.  We give the local existence
and global freeness checks.

First the selected $T$-places are disjoint from the auxiliary
places.  The original \texttt{at\allowbreak T.\allowbreak pi\allowbreak U\_\allowbreak cong} gives
$\pi\equiv1$ modulo
$N=Q_Sqd(g_0g_1)^2$, and \texttt{pi\allowbreak U\_\allowbreak norm} gives
$\pi\overline\pi=T$.  At any prime above $g_i$, both
$\pi$ and $\overline\pi$ reduce to one, so $T$ reduces
to one.  It cannot equal $g_i$, whose residue is zero.
This derives $T\ne g_0,g_1$ from the supplied fields,
rather than treating it as an additional host condition.
The original tower-divisibility field also gives
$4\mid Q_S$, hence $4\mid N$.  The same norm congruence
therefore gives $T\equiv1\pmod4$.  Since $T$ is prime,
it follows in particular that $T\ge3$, pointwise and
without an asymptotic size argument.

At a selected place over $T$, the element $\pi$ in
$T_0.\mathrm{atT}$ has residue zero and $\overline\pi$ has
nonzero residue.  Put $\alpha=\pi-\overline\pi$; then
$\overline\alpha=-\alpha$ and $a=\alpha^2$ lies in $F$.
Choose an integral $b\in F$ lifting the nonzero residue of
$\alpha$ at this place.  The residue field is the prime field
of order $T$, and $T\ne2$, so
$b^2-a\in\mathfrak p$ and $2b$ is a unit.  Successive Newton
corrections for $X^2-a$ converge in the completed discrete
valuation ring to a unit $r$ with $r^2=a$.  The explicit
idempotents
\[
            e=\frac{1+\alpha/r}{2},\qquad
            1-e=\overline e
                                                                  \tag{RSL.12}
\]
are integral in the completed CM algebra and split its two
factors.

The level functional has an additive quotient of prime order
$T$.  An idempotent acts on this quotient by zero or identity:
its image is an additive subgroup, hence is zero or the whole
group, and a surjective idempotent is the identity.  Replacing
$e$ by $1-e$ if needed, take the one acting as identity.  For
$a_v\in k_v^\times$, choose an integer $k$ representing it,
with $T\nmid k$.  The local scalar
\[
                   \mu=k^{-1}e+k(1-e)
                                                                  \tag{RSL.13}
\]
is an integral unit of norm one.  It preserves the original
local lattices and scales the level functional by $k^{-1}$,
which is the prescribed diamond operation.  Put this scalar
at that place and the identity elsewhere.  This gives a
genuine finite adelic unitary element, not just a residue-class
operator.  It normalizes $U$, because scaling a level functional
by a unit preserves its kernel.  Elements at different places
commute, and its $(T-1)$st power belongs to $U$ by
$a_v^{T-1}=1$.  The quotient action is independent of the
integer lift.  The \texttt{fv0} and level fields of $T_0$
identify this action with the stored dilation of the original
class set and with the stored operators on forms.

To prove freeness, let a multidiamond $\delta$ representing
$g\in\mathscr G$ fix the class of $y$.  Thus
$y=\gamma y\delta\kappa$ for a rational unitary $\gamma$
and $\kappa\in U$, with the harmless inverse convention
absorbed into $g$.  Since $\delta$ normalizes $U$ and
$\delta^{T-1}\in U$, iteration gives
$y^{-1}\gamma^{T-1}y\in U$.  RSL7 implies
$\gamma^{T-1}=1$.  At the auxiliary primes $g_0,g_1$ the
diamond is the identity.  Therefore $y^{-1}\gamma y$ there
has the same upper-unipotent reduction as a level element.
The two-prime argument in RSL7 applies to this finite-order
$\gamma$ and gives $\gamma=1$.  It follows that
$\delta\in U$.  At each selected place the level functional
is nonzero on its quotient of order $T$; invariance under
its multiplier forces that multiplier to be one.  Thus every
coordinate of $g$ is one, which proves freeness.

All local fields, idempotents and functional quotients in this
argument belong to $T_0.\mathrm{atT}$ and
$T_0.\mathrm{cls0}.\mathrm{lev}$.  In particular the action is
the original one occurring in the definitions of
$\mathrm{ramSub}$ and \texttt{Free\allowbreak T}.
% END RSL8

% BEGIN RSL9
\paragraph{RSL9: exact integral orbit slices and isometry at every embedding.}
\label{claim:RSL9}
Assume in addition \texttt{Chi\allowbreak Fin0} for $T_0$.  There exist a
nonzero good Hermitian $\mathcal O_K$-lattice $W$ and an
injective $K$-linear map
\[
 \Psi:W.V^{\mathscr G}\longrightarrow B.V
                                                                  \tag{RSL.14}
\]
such that
\[
 \begin{split}
 B.L&=\Psi\bigl((W.L)^{\mathscr G}\bigr),\\
 \langle\Psi F,\Psi F'\rangle_\sigma
   &=\sum_{g\in\mathscr G}\langle F_g,F'_g\rangle_{W,\sigma}
       \quad\hbox{at every infinite place }\sigma,\\
 (D.\mathrm{fm}(u)).\mathrm{ramSub}
   &=\Psi\left(\bigcap_{v\in A.v_T}\ker\mathrm{Nm}_v\right),
 \end{split}                                                     \tag{RSL.15}
\]
Here $\mathrm{Nm}_v$ sums along the $k_v^\times$ coordinate.
The source calls these three properties
\texttt{Induced\allowbreak Lattice}.

\emph{Rationality and the slice decomposition.}
Each original diamond operator preserves $B.L$, by the original
\texttt{rho\allowbreak T\_\allowbreak lat} field, and therefore preserves $B.V$ because
$B.L$ spans it.  Injectivity of $j$ descends it uniquely to a
$K$-linear automorphism $\rho(g)$ of $B.V$; the group relations
descend by the same injectivity.  Thus the group action is on
the original rational space, not just its complexification.

Choose one representative in each of the finite free orbits of
RSL8, and let $W.V\subset B.V$ consist of forms supported on
these representatives.  Restricting a rational form to any set
of classes is still rational.  Indeed \texttt{val0.\allowbreak rat}
characterizes the image of $j$ by classwise rational polynomial
representatives; retain that polynomial on the chosen classes
and take zero elsewhere.  Trivial stabilizers, proved in RSL7,
make these classwise data legitimate automorphic forms, because
the only consistency condition inside a rational double coset
is its stabilizer condition.  Extending by the level character
and rational unitary action gives the required form, and its
classwise rationality puts it back in the image of $j$.

Define
\[
                       \Psi(F)=\sum_{g\in\mathscr G}\rho(g)F_g.
                                                                  \tag{RSL.16}
\]
At any class $c$, precisely one summand can be nonzero: freeness
gives exactly one $g$ carrying the appropriate class to the
chosen representative of its orbit.  Thus (RSL.16) is injective.
For a given form $b$, restrict it to each of the translated
representative sets and apply the inverse translate; these
pieces belong to $W.V$ and sum back to $b$.  This proves
surjectivity as well and describes the inverse explicitly.

\emph{Equality of the integral lattices.}
Put $W.L=B.L\cap W.V$.  The \texttt{val0.\allowbreak lat} field says that a
rational form lies in $B.L$ exactly when, at every class and
every coefficient place, its polynomial value lies in the
specified local integral polynomial module, in every allowed
local basis.  This is a classwise condition; zero satisfies it.
Consequently restriction to a set of classes preserves $B.L$.
Every $\rho(g)$ and $\rho(g^{-1})$ also preserves $B.L$.
The explicit inverse to (RSL.16) therefore preserves the
lattice coordinate by coordinate, and proves
\[
       \Psi(F)\in B.L\quad\Longleftrightarrow\quad
                    F_g\in W.L\ \hbox{for every }g.
                                                                  \tag{RSL.17}
\]
This is equality of integral modules, rather than commensurability.
The slice projector is $K$-linear and preserves $B.L$, so its
image $W.L$ is finitely generated and spans $W.V$; restricting
the positive definite metrics makes it a good Hermitian lattice.

It is nonzero.  The original class set contains the identity
double coset.  The level-zero coefficient representation has a
nonzero highest-weight function.  Explicitly, for a dominant
integer weight $\lambda_1\ge\cdots\ge\lambda_k$, take
\[
 (\det Z)^{\lambda_k}
       \prod_{j=1}^{k-1}\Delta_j(Z)^{\lambda_j-\lambda_{j+1}},
\]
where $\Delta_j$ is the leading $j\times j$ minor, and take
the product over the stored embeddings.  The minor exponents
are nonnegative, the determinant factor is defined on the
invertible frames even when $\lambda_k<0$, and its value on
the coordinate identity frame is one.  This is the extreme
weight element of the original coefficient module, expressed
in the designated frame coordinates.  It therefore gives a
nonzero value function of exactly the required weight.
Placing it on one class and zero elsewhere gives a nonzero
complex form, since the stabilizer is trivial.  The original
\texttt{emb\_\allowbreak basis} spans this complex form space; hence
$B.V\ne0$.  Since $B.L$ spans $B.V$ and (RSL.16) is a
bijection, $W.V$ and $W.L$ are nonzero as claimed.

\emph{The original metrics.}
The field \texttt{norm0} expresses each squared norm as the
sum over classes of the specified Fock norms of its polynomial
values, weighted by reciprocal stabilizer orders.  RSL7 makes
these weights one.  Forms supported on disjoint sets of classes
are orthogonal at every embedding.  One can check the mixed
term from the norm identity: apply it to $x+y$ and $x+a y$,
where $a\in K$ has nonreal image at the embedding in question.
Such an $a$ exists because $K$ is totally complex.  The two
equalities force both real and imaginary parts of the Hermitian
cross term to be zero.

It remains to prove that each $\rho(g)$ is an isometry; a
permutation of supports alone would not suffice.  For a single
diamond and each class, the original double-coset equality and
\texttt{val0} express its new polynomial value as the old value
after substitution by a rational unitary matrix $\gamma^{-1}$,
multiplied by a level-character value $\chi(\kappa)$, and
evaluated at the dilated class.  The character hypothesis gives
$r\in K$ and $N>0$ with
$s_K(r)=\chi(\kappa)$ and $s_K(r)^N=1$.
Injectivity gives $r^N=1$ in $K$, so
\[
                         |\sigma(r)|=1
                   \quad\hbox{for every embedding }\sigma.
                                                                  \tag{RSL.18}
\]
The rational unitary matrix has its entries in the original
$E$.  Mapping its entries to $K$ through the stored embeddings
commutes coefficient by coefficient with polynomial
substitution and with every subsequent embedding $\sigma$.
The source value function also contains its determinant shift:
the corresponding polynomial transport has the factor
$\prod_\tau\det(\gamma_\tau)^{\mathrm{shift}_\tau}$.
Each $\gamma_\tau$ is unitary for its Hermitian form, so
$|\det(\gamma_\tau)|=1$ at every embedding.  This factor
therefore changes no squared norm, including for negative
integer shifts.

The same unitary identity carries an orthonormal frame at
$\sigma$ to another orthonormal frame.  Apply \texttt{norm0}
to the rational restriction supported on just one class,
whose existence was proved above.  Its sum has only that
class's contribution, so this contribution is independent of
the chosen orthonormal frame.  This justifies changing frames
separately for each class even though $\gamma$ can depend on
that class.  Changing frame thus preserves the evaluated
polynomial norm, while (RSL.18) preserves the character
factor.  Summing over the permuted classes
shows that the original squared norm is invariant under the
diamond at every $\sigma$.  Polarization as above gives
invariance of the full Hermitian pairing.  Products of diamonds
are therefore isometries too.

Distinct translates of the slice have disjoint supports.  The
preceding two paragraphs now prove the pairing identity in
(RSL.15), with the restricted metric on $W$.

\emph{The actual ramified subspace.}
For a single coordinate $v$ and $a\in k_v^\times$, (RSL.16)
intertwines $\rho(a)$ with translation of that coordinate of
the function $F$.  Summing over $a$ intertwines the original
norm operator with coordinate summation.  The definition of
the original $\mathrm{ramSub}$ is the joint kernel of exactly
these norm operators; the original rational embedding is
injective.  This proves the last equality of (RSL.15).
Taking its intersection with (RSL.17) gives the exact original
ramified lattice, including its induced metric.
% END RSL9

% BEGIN RSL10
\paragraph{RSL10: the actual rank, mean and maximal-slope formulas.}
\label{claim:RSL10}
The source-facing conclusion retains
\texttt{Hida\allowbreak Setting D u}.  The argument below needs only
$T\ge3$, which is already pointwise true by RSL8.
Put $q=T-1$, and set
\[
 U=\mathcal O_K^{q},\qquad
 A_0=\ker\bigl(U\xrightarrow{\sum}\mathcal O_K\bigr).
\]
Give $U$ its standard metric.  RSL9 identifies the original
full lattice with $W\otimes U^{\otimes f}$ and the original
ramified lattice with $W\otimes A_0^{\otimes f}$,
integrally and isometrically.  The second identification uses
the integral splitting
$U=A_0\oplus\mathcal O_K e_q$: its tensor product shows that
the joint coordinate-sum kernel is precisely the indicated
tensor product, without a saturation or index correction.

If $h=\operatorname{rk}W$, their ranks are $h q^f$ and
$h(q-1)^f$.  Define the actual source integer
\[
 d_u=(D.\mathrm{mem}(u)).A.du
      =\dim_{\mathbb C}(D.\mathrm{mem}(u)).A.M_{\mathrm{ram}}.
\]
The original $\texttt{emb\_\allowbreak basis}$ identifies
$B.V\otimes_{K,s_K}\mathbb C$ with $A.M_{u,0}$,
so $D.duS(u)=h q^f$.  Each diamond and norm operator
descends to $K$, as proved in RSL9.  Their joint kernel
commutes with scalar extension: stack their finitely many
matrices into one matrix over $K$; its rank is unchanged
over $\mathbb C$, because rank is characterized by its
nonzero minors.  Under the original embedding, their
complex joint kernel is exactly $A.M_{\mathrm{ram}}$.
Therefore the original rational \texttt{ram\allowbreak Sub} has
$K$-dimension $d_u=h(q-1)^f$.  The injectivity and range
of the original $D.\mathrm{nuIn}$ identify the division
vector space with this very rational kernel.  Thus
\[
 \frac{d_u}{D.duS(u)}
       =\left(1-\frac1{T-1}\right)^f.                 \tag{RSL.19}
\]
The lattice $W$ is nonzero and $q-1>0$, so both ranks are
positive.  The basis $e_j-e_q$ of $A_0$ has Gram matrix
$I+J$ and determinant $q$.  The Kronecker determinant
calculation in LC7, now applied through the exact isometry of
RSL9, gives
\[
 \boxed{\quad
 (D.\mathrm{fm}(u)).\mathrm{ramLat}.\mu
       =B.\mu-\frac{f\log(T-1)}{T-2}.
 \quad}                                                          \tag{RSL.20}
\]
For clarity, the tensor Gram determinant is
$q^{f(q-1)^{f-1}}$ when $f>0$; dividing its logarithm by
$(q-1)^f$ gives the loss in (RSL.20).  Finite indices in a
free sublattice of the projective $W.L$ occur with the same
normalized multiplicity on both sides, so no freeness
assumption on $W.L$ is required.  The squared-metric slope
convention gives exactly (RSL.20), with no extra factor two.
When $f=0$ the tensor product is $W$, the rank factor is one
and the loss is zero, so this boundary case is consistent too.

The original ramified lattice is the intersection of $B.L$
with a rational subspace, with the restricted original
metrics.  Every lattice in this subspace used to compute its
maximal slope is also an admissible lattice in $B$.  Taking
the supremum gives
\[
        \mu_{\max}\bigl((D.\mathrm{fm}(u)).\mathrm{ramLat}\bigr)
                           \le\mu_{\max}(B).                     \tag{RSL.21}
\]
This is also valid if the convention restricts to saturated
submodules: the ramified intersection is saturated in $B.L$,
and saturation inside it remains saturated in $B.L$.

Finally choose a complex basis of $W.V\otimes_{K,s_K}\mathbb C$.
Through (RSL.16), the full form space becomes functions on the
finite free $\mathscr G$-set consisting of one copy of
$\mathscr G$ for each basis vector.  The single-coordinate
actions intertwine the stored $\rho_T$ operators.  This is
the literal \texttt{Free\allowbreak T D u} record.  Equation (RSL.20) is
the literal \texttt{Slope\allowbreak Line D u}; (RSL.19)--(RSL.21) are
the three fields of \texttt{family\_\allowbreak property\_\allowbreak and} on the
original data.  RSL8 alone already gives the corresponding
complex free model, without the character hypothesis; the
finite-character condition is used to obtain the isometry at
every coefficient embedding.

Applied to RSL6 at one common threshold, these results supply
the fields \texttt{h\allowbreak W} and $h_\mu$ of
\texttt{Fam\allowbreak GLmin}.  This closes LC7's actual integral and
isometric transport boundary.  It supplies no assertion
about the Hecke maps between the two different local
integral lattices appearing in RSL13.
% END RSL10

\subsection{A smaller fullness cut and the residual global ledger}

% BEGIN RSL11
\paragraph{RSL11: actual fullness needs only the ramified local inclusion.}
\label{claim:RSL11}
For an original family $\mathcal F$, suppose there is one
threshold $a_N$ such that, for every $x\ge a_N$, every
$i\in\mathcal F.I_x$, every word $u$ and every coefficient
place $w$ of residue characteristic $T_u$,
\[
                N_\lambda(D,u,w)\subset S_{\mathrm{ram}}(D,u,w),
       \qquad D=\mathcal F.D(x,i).                    \tag{RSL.22}
\]
Then \texttt{Nu\allowbreak Full\allowbreak F Fm} follows from the original fields of
\texttt{Fam\allowbreak P}.  In particular it does not require the entire
congruence length bound, a mean-slope theorem, or a new
arithmetic existence assertion.

Fix a member beyond $\max\{a_N,\mathcal F.x_1+1\}$, and write
$M=L.\mathrm{NuLat}(u)$ and
$R=\mathrm{MramL}(D,u)$, both in the original vector space
$(D.\mathrm{div}.\mathrm{Nu}(u)).V$.  The field
$(\mathcal F.\mathrm{ex}(x,i)).\mathrm{lat}.\mathrm{atT}$
identifies the localization of $M$ at $w\mid T_u$, through
the original $D.\mathrm{nuIn}(u)$, with $N_\lambda(D,u,w)$.
The first component of $S_{\mathrm{ram}}$ is the original
local full lattice $S(0,w)$.  The field $D.\mathrm{locSpan}$
clears its local denominators.  Thus (RSL.22) implies
$M_w\subset R_w$ at every $w\mid T_u$.
At the other places, the original
$\mathcal F.\mathrm{nuDef}(x,i)$ supplies this inclusion.
Consequently $M\subset R$ globally.  For completeness, if
$z\in M$ did not belong to $R$, its annihilator ideal
$\{a\in\mathcal O_K:az\in R\}$ would be proper and
contained in a maximal ideal $\mathfrak p$.  Local inclusion
at $\mathfrak p$ would supply $s\notin\mathfrak p$ with
$sz\in R$, a contradiction.

The original $R$ is finitely generated: its injective map
$D.\mathrm{nuIn}$ identifies it with a submodule of the
finitely generated lattice $B.L$, and $\mathcal O_K$ is
Noetherian.  Since
$\mathcal O_K$ is Noetherian, its submodule $M$ is finitely
generated too.  The original \texttt{Nu\allowbreak Def} also gives
$M_w=R_w$ at a suitable place away from $T_u$.  There is
always an eligible rational prime: use $2$ when $T_u\ne2$,
because $2$ belongs to the host's designated small-prime set;
when $T_u=2$, use the host prime $q$, which is distinct from
$2$ and has the designated non-inert class.  Choose a
coefficient place above this prime.  For finitely many
generators $r_1,\ldots,r_b$ of $R$, local equality gives
$s_j\notin\mathfrak p$ with $s_jr_j\in M$.
Their product $c=\prod_j s_j$ is nonzero and satisfies
\[
                              cR\subset M\subset R.             \tag{RSL.23}
\]
This is exactly \texttt{Nu\allowbreak Commens D L}.

The original $R$ spans the original division vector space:
apply $D.\mathrm{nuIn}$ to a vector, express it in the
$K$-span of $B.L$, clear the finitely many coefficients,
and pull back.  Thus (RSL.23) proves that $M$ spans it too.
Finally the actual \texttt{nprime\allowbreak Def} field of the same
$\mathcal F$ gives
\[
                 N=(D.\mathrm{div}.\mathrm{Nu}(u)).L=I_u M,
                                                                  \tag{RSL.24}
\]
with the same nonzero fractional ideal $I_u$ as in MT1.
Choose $0\ne d\in\mathcal O_K$ with $dI_u\subset\mathcal O_K$.
Then $dN\subset M$, so Noetherianity and injectivity of
multiplication by $d$ imply finite generation of $N$.
Choose a nonzero integral element $a\in I_u$; (RSL.23)
gives $acR\subset N$, so $N$ spans the division vector
space.  These are both fields of $D.\texttt{Nu\allowbreak Full}$.
The threshold was chosen before $i,u,w$, hence proves
\texttt{Nu\allowbreak Full\allowbreak F} with its original uniformity.  Neither
\texttt{nu\allowbreak Def} nor \texttt{nprime\allowbreak Def} has been added as a
new datum; both are fields of the original family.
% END RSL11

% BEGIN RSL12
\paragraph{RSL12: actual positive rank and the sufficient operator cut.}
\label{claim:RSL12}
Suppose an original family has the actual second layer of
RSL6, \texttt{Tu\allowbreak Size}, and the eventual inclusion (RSL.22).
Then
\[
                   \texttt{Nu\allowbreak Full\allowbreak F}\ \wedge\ \texttt{J.\allowbreak Mu\allowbreak Rank}
                                                                  \tag{RSL.25}
\]
holds for that same family.  The second conjunct says that
each original division lattice has strictly positive rank.

Indeed RSL8 gives $T_u\equiv1\pmod4$ and hence $T_u>2$
on each of the original records.  Take one maximum of
the second-layer threshold, $a_N$, and
$\mathcal F.x_1+1$.  RSL11 gives fullness of the original
division lattice.  Its full $K$-rank equals the dimension
$d_u=(D.\mathrm{mem}(u)).A.du$ of the original ramified form space, since
$D.\mathrm{nuIn}$ is injective with precisely that range.
RSL9--RSL10 give $D.duS(u)>0$, $T_u>2$ and (RSL.19),
so $d_u>0$.  This proves (RSL.25) at one threshold.
The nonzero slice constructed in RSL9 makes the positivity
argument slightly stronger than the generic source lemma
\texttt{mu\allowbreak Rank\_\allowbreak min}, which instead obtains $D.duS(u)>0$
from \texttt{Chosen\allowbreak Proviso}.  On the selected final route
\texttt{Chosen\allowbreak E Fm 3 CH} already contains that proviso;
either argument applies, and neither changes the family.
The retained \texttt{Tu\allowbreak Size} hypothesis is consequently
unused in this stronger actual-forms proof.  In particular
we have not appealed here to the longer numerical supplier
\texttt{Tu\allowbreak Size}$\Rightarrow$\texttt{Setting\allowbreak Eventually}.

We also isolate exactly what is needed to supply (RSL.22).
At a seen coefficient place above $T_u$, suppose the actual
level-zero operators satisfy
\[
 U_v(B_w(v))\subset S_w,\qquad
 U_v(M_{\mathrm{ram}})\subset M_{\mathrm{ram}},\qquad
 T_u^{b_v}S_w\subset B_w(v),
                                                                  \tag{RSL.26}
\]
where $b_v=\mathrm{bTop}(\tau_w(v))\ge0$ for the selected
places, as supplied by \texttt{Hida\allowbreak Setting}.  The normalized
operator $G_v=T_u^{b_v}U_v$ then preserves
$S_{\mathrm{ram},w}=S_w\cap M_{\mathrm{ram},w}$.
Rationality of $U_v$ follows from its integral map between
two full local lattices: choose a basis, clear denominators
in the full source lattice, and descend by the injective
original rational embedding.  The exact source definition
of $N_\lambda$ is the image of $S_{\mathrm{ram},w}$ under
the $m$th power of the product of these $G_v$, in the
source's ordered list of selected places.  Distributing the
central scalar factors gives the exponent
$m\sum_v b_v=\texttt{exp\allowbreak T}$ in that original definition;
no interchange of two Hecke operators is used here.
Successive preservation proves (RSL.22).  At an unseen
place the selected set is empty, its product is the
identity, and $N_\lambda=S_{\mathrm{ram},w}$ directly.

Thus (RSL.26) is sufficient for fullness and positive rank.
The full three contracts are
\[
 \begin{gathered}
 \texttt{Hecke\allowbreak Integral\allowbreak Level\allowbreak Zero},\qquad
 \texttt{Hecke\allowbreak Ramified\allowbreak Level\allowbreak Zero},\\
 \texttt{Second\allowbreak Branch\allowbreak Inclusion\allowbreak Length}.
 \end{gathered}
\]
They imply (RSL.26), but
their reverse-operator identity and exact length formula
are stronger than this use requires.  Those stronger
parts are still needed in the selected global error
ledger.  We have proved the implication from (RSL.26),
not silently asserted those operator and length contracts.
% END RSL12

% BEGIN RSL13
\paragraph{RSL13: the exact global cut and the remaining proof tasks.}
\label{claim:RSL13}
For one fixed class $\mathcal C$ the selected source interface
is the eleven-field conjunction \texttt{Fam\allowbreak GLmin C 4}.
The following ledger is literal; every field concerns the
same original family, member, coefficient field and word.
\begin{center}
\begin{tabular}{p{0.11\textwidth}p{0.47\textwidth}p{0.30\textwidth}}
\hline
Field & Contract & Human status in this chapter \\
\hline
$h_{2b}$ & Eventual standard adjacent genus-mass ratio & MF19, closed. \\
$h_{2v}$ & Eventual original base-volume equals mass & RSL1--RSL3, closed from the supplied real tie. \\
$h_3$ & \texttt{Sigma\allowbreak Degree\allowbreak Two\allowbreak Sided\allowbreak Bound\allowbreak Of\allowbreak Gap} & B: explicit finite-product bounds and original-row installation; deep discrete-series producer covered by DC1, DC10, DC19. \\
$h_{Za}$ & \texttt{Family\allowbreak Slope\allowbreak Lower\allowbreak Open\allowbreak Part} & B: exact column-lattice slope and standard-weight estimate. \\
$h_{Zb}$ & \texttt{Family\allowbreak Slope\allowbreak Max\allowbreak Open\allowbreak Part Fm 4} & B: original class-map degree and hook bound. \\
$tCV$ & Eventual \texttt{Nonempty (Col\allowbreak Val D u)} & B: actual polynomial/local-basis/archimedean dictionary. \\
$h_W$ & \texttt{Fam\allowbreak U (Nonempty Free\allowbreak T)} & RSL7--RSL10, closed from the second layer. \\
$h_\mu$ & \texttt{Fam\allowbreak U Slope\allowbreak Line} & RSL9--RSL10, closed from the second layer. \\
$h_I$ & \texttt{Fam\allowbreak U Hecke\allowbreak Integral\allowbreak Level\allowbreak Zero} & B: actual two integral Hecke maps. \\
$h_R$ & \texttt{Fam\allowbreak U Hecke\allowbreak Ramified\allowbreak Level\allowbreak Zero} & B: actual ramified stability and inverse composite. \\
$h_B$ & \texttt{Fam\allowbreak U Second\allowbreak Branch\allowbreak Inclusion\allowbreak Length} & B: actual local inclusion and exact length. \\
\hline
\end{tabular}
\end{center}

There is a separate B installation gate when these conditional
\texttt{Fam\allowbreak U} fields are used on actual late members:
\[
       \texttt{Tu\allowbreak Size}\ \Longrightarrow\
       \texttt{Setting\allowbreak Eventually},
                                                                  \tag{RSL.28}
\]
where the latter asserts the full \texttt{Hida\allowbreak Setting} at
one threshold before $i,u$.  The source supplier is
\texttt{Places.\allowbreak setting\allowbreak F\_\allowbreak of}, through
\texttt{nums\allowbreak L\_\allowbreak wts\allowbreak L\_\allowbreak fam} and
\texttt{level\allowbreak Prime\_\allowbreak setting\_\allowbreak of\_\allowbreak sizes}.
Its numerical size and weight comparisons still need a human
reconstruction.  In particular, the class projection supplying
\texttt{Tu\allowbreak Size} does not itself prove all the alcove and
$b_{\rm top}\ge0$ inequalities.  The strong Hecke/length route
to (RSL.26), and its use in the final global ledger, must retain
this gate.  The weaker direct fullness input (RSL.22) does not
need it, and RSL12 has avoided it for positivity by the original
congruence $T\equiv1\pmod4$.

The smallest \emph{literal GL4 cut at this interface} consists
of the three fields $h_{Za},h_{Zb},tCV$.  Its sufficiency for
these three slots is by their definitions; no assertion of
absolute mathematical minimality is intended.  More concrete
producer targets are:
\begingroup\raggedright
\begin{itemize}
\item \texttt{column\_\allowbreak lattice\_\allowbreak exact\_\allowbreak slope}
in \path{GL4/ColumnLatticeExactSlope.lean}, followed by
the standard Schur-lattice slope and hook-length estimates
used in \path{GL4/LowerSlopeInequalityFamily.lean}.
The equality must hold for the actual \texttt{Col\allowbreak Val.\allowbreak Q}
with its exact base-change lattice and metric.
\item \texttt{class\_\allowbreak map\_\allowbreak archimedean}
in \path{GL4/ClassMapArchimedean.lean}, followed by the
class-map slope inequality and the constant-four hook bound
in \path{GL4/PStmtFree.lean}.  The local class bases and
the original archimedean matrices must be used together.
\item \texttt{Col\allowbreak Val\allowbreak S.\allowbreak col\allowbreak Val\allowbreak Of\allowbreak Tied\allowbreak S\_\allowbreak holds} in
\path{GlobalTheory/TransposeBaseMatrixInfinity.lean}.
Its useful smaller dictionary target is
\texttt{Col\allowbreak Tie\allowbreak W}: an injective polynomial realization of
the original rational form space, the exact local integral
membership in bases of the actual $Q_{c,\tau}$, and the
exact sum-of-polynomial-norms formula for the original
metric.  RSL9 establishes the rational restriction and
metric transport mechanisms, but does not by itself
identify these $Q_{c,\tau}$ or write the complete
\texttt{Col\allowbreak Tie\allowbreak W} constructor.
\end{itemize}
\endgroup

For the three GL6 operator slots, the exact source producers are:
\begingroup\raggedright
\begin{itemize}
\item \texttt{integral\_\allowbreak at\_\allowbreak t0\_\allowbreak lattice} in
\path{OrbitalBoundsConclusion/IntegralAtT0Lattice.lean};
\item \texttt{hecke\_\allowbreak ramified\_\allowbreak zero\_\allowbreak all} in
\path{GlobalTheory/UnramifiedTrace/HeckeRamifiedZeroAll.lean},
which installs the same-place matching and unramified
composite calculations;
\item \texttt{second\_\allowbreak branch\_\allowbreak all\_\allowbreak members} in
\path{GlobalTheory/CoordinateLengths/SecondBranchAllMembers.lean}.
\end{itemize}
\endgroup
Their mathematical content is a finite local coset/integrality
calculation and a discrete-valuation length calculation;
their exact installations remain B reconstruction tasks.
The finite product of operators in RSL12 proves their
weaker fullness consequence once (RSL.26) is supplied.
For the full ledger one must also write the finite-index
length estimates for the composite and its $m$th power;
the source declaration is
\texttt{congruence\_\allowbreak no\_\allowbreak kernel\_\allowbreak minimal}.

The closed transport claims do not reinstate older
boundaries.  TG12 supplies the full columns on true-column
data; DC15 and DC19 supply the large-gap discrete-series
existence, uniqueness and degree under their named
normalizations; DC20 matches all four real-place fields
when those normalizations are supplied.  FR1 and GF8
prove the finite-dimensional frame and normalized
Harder--Narasimhan consequences.  They do not identify
the actual GL4 lattices or the two Hecke lattices, and
they do not check by themselves every input of the final
\texttt{link\allowbreak GLOBAL7\_\allowbreak items} assembly.  That last numerical
and lattice assembly remains an explicit B proof-of-implication
task, rather than being counted as closed because the
formal source calls it.

For precision, finite conjunction of already proved
\texttt{Closed} contracts is elementary.  Choose the
minimum of their finitely many positive $\varepsilon$
bounds.  For each allowed $\varepsilon$, take the supremum
of their finitely many transported bounds.  The definition
of \texttt{Large} at that supremum implies each original
\texttt{Large} condition.  Thus every contract holds for
one and the same $\mathcal F,C_H$; within that family,
take the maximum of the finitely many eventual thresholds.
This proves the finite joining step without changing
any outer quantifier.

On the selected final route put $\mathcal C=D3(3,\mathcal A)$.
The precise terminal source contract is
\[
 \begin{gathered}
 \texttt{Duty}\hbox{ on }\mathcal A,\quad
 \texttt{Closed}(\mathcal C,\texttt{Real\allowbreak Place\allowbreak Second\allowbreak Layer\allowbreak Conditions}),\\
 \forall D\,\texttt{Mass\allowbreak Ratio\allowbreak Under\allowbreak Volume}(D),\quad
 \forall D\,\texttt{Genus\allowbreak Mass\allowbreak At}(D,2)\\
 \Longrightarrow\quad
 \texttt{Closed}\bigl(\mathcal C,
  \texttt{Chosen\allowbreak E}(3,C_H)\Rightarrow
                  \texttt{Nu\allowbreak Full\allowbreak F}\wedge\texttt{Mu\allowbreak Rank}\bigr),\\
 \texttt{Closed}\bigl(\mathcal C,
   \texttt{Tu\allowbreak Size}\Rightarrow\texttt{DGam\allowbreak Size}\Rightarrow
                  \texttt{Link\allowbreak GLOBAL}(2,C_H,3,2,0)\bigr).
 \end{gathered}                                                   \tag{RSL.27}
\]
RSL1--RSL3 close both universal mass/volume inputs.
RSL4--RSL6 close the finite-character and real/column
join once the same-record selection has been installed.
The class projection $D3(3,\mathcal A)\to\mathcal A$
then supplies the original \texttt{Chosen\allowbreak Wall} from
\texttt{Duty}; it does not choose a new family.
Membership in this same $D3$ class already includes
\texttt{Tu\allowbreak Size}.  Consequently RSL12's \texttt{Tu\allowbreak Size}
hypothesis is a class projection on the selected
fullness/rank route, not a new producer to construct.
The numerical parameters are exactly $C_{\rm abs}=2$,
$C_0=3$, $e_D=2$, $e_L=0$; $C_5=4$ belongs only to
the intermediate slope estimate.  The external
\texttt{DGam\allowbreak Size} argument is retained in (RSL.27)
although the \texttt{link8} wrapper does not use it.
To count the full human implication (RSL.27) as finished,
the B entries, the size-installation gate (RSL.28), and the
final assembly just identified
must still be written.  This chapter does not replace
those proofs by an invocation of \texttt{link\_\allowbreak gl}.
% END RSL13


\section{Actual formal degrees and a direct construction of the level-prime size bound}
\label{sec:degree-prime-size}

\status{DSZ1--DSZ3 close the original two-sided degree input $h_3$ on
every supplied family, using the already matched Harish--Chandra input
of DC18. DSZ4--DSZ9 prove the literal \texttt{Tu\allowbreak Size} from the original
eventual volume identity. DSZ10 supplies that identity on the original
\texttt{Built} construction class and proves its \texttt{Closed} size
obligation. The common class and its actual families still have to be
constructed. The proof does not require the stronger Vandermonde
variation or prime-counting package.}

Write $W_n=\prod_{j=0}^{n-1}j!$ and
$V_n(z)=\prod_{a<b}(z_a-z_b)$. All ranks, coordinates, levels and
volumes below belong to the same original datum. In the family
calculation $L=x^\epsilon$, $B=e^L$, $a=xB$, $r=k+m-1$, and
$f=[E:\mathbb Q]/2$. The polynomial dimension is $d_n(z)=V_n(z)/W_n$.
The letter $\delta$ in DSZ2 is an error tolerance, not the family
exponent $\epsilon$ and not a Harish--Chandra coordinate gap.

% BEGIN DSZ1
\paragraph{DSZ1: the finite product comparison with the exact compact type.}
\label{claim:DSZ1}
Let $\lambda_0>\cdots>\lambda_{n-1}$ be regular integral in the sense
of (DC.1), and let $s$ be the original sign vector. Put
$p=\#s^{-1}(+)$, $q=\#s^{-1}(-)$, and
$A=\max_{a,\allowbreak{} b}|\lambda_a-\lambda_b|$, with $A=0$ for an empty
index set. Define
\[
 d_\tau(s,\lambda)=
 \frac{\displaystyle\prod_{a<b,\ s_a=s_b}
       (\lambda_a+\rho_{n,\allowbreak{} a}-\lambda_b-\rho_{n,\allowbreak{} b})}{W_pW_q}.
                                                               \tag{DSZ.1}
\]
Here $\rho_n$ is the half-sum of the noncompact positive roots,
with exactly the convention of DC1. If $a<b$ and $s_a=s_b$, direct
counting in its definition gives
\[
 \rho_{n,\allowbreak{} a}-\rho_{n,\allowbreak{} b}
       =\#\{c:a<c<b,\ s_c\ne s_a\}\ge0.
\]
Thus each compact numerator factor in (DSZ.1) dominates the
corresponding factor of $V_n(\lambda)$. There are exactly $pq$
unordered pairs of opposite signs; choosing their increasing order
identifies them with the noncompact positive roots. Each of their
factors is at most $A$. Finally
\[
 W_pW_q\le W_{p+q},
\]
because $W_{p+q}/W_p=\prod_{j=0}^{q-1}(p+j)!\ge\prod_{j=0}^{q-1}j!$.
Multiplying and dividing these positive quantities proves
\[
             0<d_n(\lambda)\le d_\tau(s,\lambda)A^{pq}.
                                                               \tag{DSZ.2}
\]
For $n=0$ both products are empty, $p q=0$ and the assertion is
$1\le1$; for $n=1$ the same convention applies. For $n\ge2$,
regular integrality makes every positive difference an integer,
so $A\ge1$. DC1 identifies (DSZ.1) with the dimension of the
specified compact type, and DC10 identifies the actual lowest
compact-type subspace with that type in the constructed large-gap
representation. DC13 transfers this identification to every literal
\texttt{Is\allowbreak DS} representation in the same large-gap range.
\emph{Source: Z01--Z03.}
% END DSZ1

% BEGIN DSZ2
\paragraph{DSZ2: the source gap reaches the reconstructed existence range.}
\label{claim:DSZ2}
Suppose $p,\allowbreak{} q\ge1$, $n=p+q$, $e^{-n}\le\delta<1/2$, and
\[
 h\ge8192\min(p,\allowbreak{} q)(pq)^3\frac{\log n}{\delta^2}.
                                                               \tag{DSZ.3}
\]
Then $\delta>0$, $n\ge2$, $pq\ge1$, $n\le2pq$, and
$\log n\ge\log2>1/2$. It follows that
$h>16384(pq)^3>32npq+2$: indeed
$32npq+2\le64(pq)^2+2\le66(pq)^3$.
Writing $q_0=\min(p,\allowbreak{} q)$, we have
$\sqrt{2q_0}\le2q_0$, so
\[
 16nq_0+8\sqrt{2q_0}\le24nq_0\le24npq<32npq+2.
\]
The other threshold $npq+1$ is smaller too. Consequently (DSZ.3)
implies the actual existence threshold (DC.2), and also the smaller
rigidity threshold used in DC13. This checks the needed hypothesis
against the reconstructed DC argument, rather than only against the
formal source's potentially smaller existence bound.

The general construction of DC1--DC12 therefore supplies a genuine
representation at this parameter, with its entire literal
\texttt{Is\allowbreak DS} predicate as proved at the end of DC12.
DC18--DC19 identify its degree, for the exact compact-dual Haar
measure, with $d_n(\lambda)$. This uses the already stated published
input HC-Deg, whose group-and-Haar constant DC16--DC18 determine.
The input applies to the connected finite-center group $SU(p,\allowbreak{} q)$;
DC18 proves the passage through the compact central factor and its
finite covering to $U(p,\allowbreak{} q)$. Thus no compact-center or measure
normalization hypothesis is suppressed here.

For every literal representation at this parameter and every nonzero
vector, its formal degree is this same number. The lower bound
$(1-\delta)d_n(\lambda)\le d_n(\lambda)$ is immediate. For the upper
bound use (DSZ.2), $3pq\ge1$, and the lowest-type dimension just
identified. This gives exactly
\[
 (1-\delta)d_n(\lambda)\le d(\sigma)
 \le4\dim(\sigma_{\rm lowK})(3pqA)^{pq}.                       \tag{DSZ.4}
\]
\emph{Source: Z04--Z06.}
% END DSZ2

% BEGIN DSZ3
\paragraph{DSZ3: identify the stored degree and close the actual $h_3$ slot.}
\label{claim:DSZ3}
For the original $D$, the row has length
$r=2r_S+k+1=(k-1)+m$, signature
$(p,\allowbreak{} q)=(k-1+m/2,\allowbreak{} m/2)$, sign vector
$\texttt{row\allowbreak Sgn}$ and parameter
\[
 \lambda_\rho=(v_0,\ldots,\allowbreak{} v_{k-2},-b_0,\ldots,-b_{m-1}).
                                                               \tag{DSZ.5}
\]
The original finite reindexing identifies these with
$\texttt{eps\allowbreak SH}$ and $\texttt{lam\allowbreak SH}$; it changes neither a root
difference nor a product. The \texttt{HCParam} clauses give precisely
regular integrality. Also $p,\allowbreak{} q\ge1$ because $m\ge2$ is even.

Assume the two original ties \texttt{cd\allowbreak Scale\allowbreak Tie} and
\texttt{dsig\allowbreak Tie}. If (DSZ.5) has a noncompact gap satisfying
(DSZ.3), DSZ2 constructs a representation and proves its
\texttt{Is\allowbreak Fdeg} identity for the same scaled measure. Apply
\texttt{dsig\allowbreak Tie} to that representation and that degree. It yields
the exact, stronger equality
\[
                    D.\mathrm{dsig}(\rho)=d_r(\lambda_\rho).
                                                               \tag{DSZ.6}
\]
This step is essential: a positive stored number or a degree of an
unrelated representation would not give (DSZ.6).

At the source threshold \texttt{spectral\allowbreak Gap\allowbreak Lower\allowbreak Bound}, (DSZ.3)
is exactly the displayed hypothesis. Substituting (DSZ.6) and
DSZ1 into (DSZ.4) proves the original
\texttt{Sigma\allowbreak Degree\allowbreak Two\allowbreak Sided\allowbreak Bound\allowbreak Of\allowbreak Gap}, including its quantifiers
over every $\delta$ and every place. In a \texttt{Fam\allowbreak P}, both ties
are fields of $\mathcal F.\texttt{ties}(x,\allowbreak{} i)$ for every $x,\allowbreak{} i$.
Hence this predicate holds at every member, with eventual threshold
$0$. Its \texttt{Closed} version on any class has
$\epsilon_1=1$ and the zero transported bound $K$. These bounds
are chosen before a family. No new existential row selector is
required for this slot. This closes \texttt{TASK\_\allowbreak GL3\_\allowbreak GAP\_\allowbreak DEGREE}
relative to the already recorded HC-Deg application.
\emph{Source: Z06--Z09.}
% END DSZ3

% BEGIN DSZ4
\paragraph{DSZ4: read the needed inequalities from the actual family.}
\label{claim:DSZ4}
Fix a supplied family, so $0<\epsilon<1/52$, and put
$C_0=\max\{5,\allowbreak{} C_{YL}+1\}$,
$\Delta=D.s.\mathrm{Del}$,
$S=4k\Delta=\mathcal L.\mathrm{Sc}$,
$\eta=\mathcal L.\mathrm{etaSz}$, and
$w_{j\rho}=\texttt{wj}(D,\allowbreak{} j+1,\rho)$.
Here $\mathcal L=\mathcal F.L(x,\allowbreak{} i)$; this notation distinguishes
that record from the scalar $L=x^\epsilon$.
The following checks use its actual \texttt{params},
\texttt{sizes}, and \texttt{hc\allowbreak W} identities. They do not replace its
coordinates by a new equally spaced model.

The rank windows give, uniformly in $i$ and eventually in $x$,
\[
 x^{1/4}\le m\le x,\quad f,\allowbreak{} k\le x^5,\quad k\ge3x,
 \quad m\le k,\quad m^2L\le k,\quad L\ge1.                    \tag{DSZ.7}
\]
For example the lower exponent for $m$ is
$1/2-\epsilon/16-\epsilon/100>1/4$; its upper exponent is
$1/2-\epsilon/16+\epsilon/100<1$. The actual upper exponent for
$k$ is $9/2+10\epsilon+\epsilon/100<5$, and its lower exponent
exceeds $2$. The $n$ bound and $f=2n$ give $f\le x^5$ for
$x\ge3$. Increasing the threshold so that $5\log x\le L$ also
gives $k\le B$. The inequality $m^2L\le k$ is a stated size clause.
These comparisons retain the $\epsilon/100$ margins.

The host discriminant lower bound, as in QW8, gives $v>\pi$.
For the upper bound, put $Q=Q_Sqd\ge1$. The same host bounds
give $v\le Q$, while \texttt{Size\allowbreak Cond} says
$Q^{36}t^{150}\le\sqrt B$ with $t\ge1$.
Thus $v\le B$ and $\log v\le L$.
The auxiliary-prime parameters give one $Y$ with
$(kB)^8\le Y\le g_i\le2Y$ and $\log Y\le C_{YL}L$.
Therefore
\[
                  8L\le\log g_i\le C_0L.                    \tag{DSZ.8}
\]
The size bound $N(m)m^{10}\le\Delta$ gives $\Delta\ge1$.

The exact cloud identities give, for the indicated indices,
\[
 0\le v_i\le4k\Delta,\quad v_i-v_{i+1}=4\Delta,
 \qquad u_l=v_{l-1}-\tfrac12-\xi_l\Delta,
 \quad \xi_l\in\{0,1\},\quad1\le l\le k-2.                 \tag{DSZ.9}
\]
At a non-top place the actual endpoint clauses give
$4k\Delta\le u_0-v_0\le8k\Delta$ and
$u_{k-1}=v_{k-2}-1/2$. At a top place the literal endpoint formulas
are $u_0=a_T+(k-1)/2$ and $u_{k-1}=-b_T-(k-1)/2$.
The window $a\le a_T\le a+a/(kx^{26/2})$ gives
$a_T\le4a/3$. The bound $b_T\le5B'/4$ established below,
together with $k\le B$, $B'\le B$ and $3B\le a$, therefore gives
$u_0\le4a/3+B/2\le2a$ and
$u_{k-1}\ge-5B/4-B/2\ge-a$.

Here are the bounds which justify those window estimates and the
remaining weight clauses. Put
$X'=\exp(x^{\epsilon/3})$, $B'=xB/X'$ and $w''=B'/(xn)$.
For fixed $\epsilon>0$, eventually $X'\ge2x$; hence
$0<w''\le B'\le B/2$.
The actual size inequalities include
\[
 \eta\le B/(t^3fk^3)\le B/27,\qquad
 t\eta\le w''/x,\qquad kS\le w'',\qquad t\ge3.
                                                               \tag{DSZ.10}
\]
Consequently $(x+1)\eta\le w''$ and $xS\le w''/3$.
The top window has $B'\le b_T\le5B'/4$.
For its upper bound, the source's lower-window width is at most
$B'/4$ by \texttt{Size\allowbreak Rho}; all quantities discarded from that
size inequality are nonnegative. At top, thick and thin places
the respective upper bounds on $w_j$ are
$b_T+B'/x+\eta$, $B+\eta+1$, and $w''+\eta$.
Their lower bounds are respectively $B'$, $B$, and $w''$.
Since $B\ge e$, the additional $1$ in the thick window is also
absorbed by $2B$. Thus in every case
\[
 w''\le w_j\le2B,\quad b_j<w_j,\quad0<w_j-b_j\le\eta,
 \quad |w_j-b_j|\le b_j/x.                                    \tag{DSZ.11}
\]
Indeed $b_j=w_j-(w_j-b_j)\ge w''-\eta\ge x\eta$.
Also $xS+\eta\le w''$ gives $4k\Delta\le b_0/x$.
The original $e_j,\allowbreak{} b_{j+1}$ separation gives $w_j<b_{j+1}$ and
$b_{j+1}-b_j\ge1$. Original interlacing gives
$u_i>v_i>u_{i+1}$ and $u_{k-1}>-b_0$.
Every $u_i-u_j$ with $i<j$ is at least one, by the original
integral column weight, not just by strict positivity.

Finally \texttt{nodes\_\allowbreak gap} and \texttt{Size\allowbreak Nodes} give
\[
 b_{j+1}-b_j\ge\eta/m
       \ge x^{\epsilon/100}r^{19}f^5L
       \ge r^{19}f^5L/C_0.                                   \tag{DSZ.12}
\]
Equations (DSZ.7)--(DSZ.12), with the displayed cloud and endpoint
identities, check every clause of the actual \texttt{Data\allowbreak B} needed
below. In particular they check QW4--QW6's \texttt{Place\allowbreak Data} on
every word and place; $w_j\ge1$ follows from $b_0\ge4k\Delta x$.
All choices made were thresholds for this given family, uniform in
$i$, words and places. No bound on a newly chosen $\Delta$ is assumed.
\emph{Source: Z10--Z15.}
% END DSZ4

% BEGIN DSZ5
\paragraph{DSZ5: the exact formal-degree polynomial on every late member.}
\label{claim:DSZ5}
In (DSZ.5) all cloud coordinates have positive sign. The signs of
the negative bottom coordinates alternate. Two opposite-sign bottom
coordinates are separated by at least one adjacent difference in
(DSZ.12). A cloud coordinate and a negative-sign bottom coordinate
are separated by at least the first bottom difference: the first
negative-sign bottom index is $1$, $v_i\ge0$, and $b_0>0$.
It follows that the actual row's noncompact gap is at least
$r^{19}f^5L/C_0$.

The threshold (DSZ.3) for $\delta=1/3$ is at most
$9\cdot8192\,\allowbreak{} r^9$: use $p,\allowbreak{} q,\min(p,\allowbreak{} q)\le r$ and
$\log r\le r$ (this estimate includes harmless slack).
Once $L\ge9\cdot8192 C_0$, the gap just proved exceeds this
threshold, since $r\ge1$ and $f\ge1$.
Also $r\ge2$ gives $e^{-r}\le1/3$.
DSZ3 therefore gives (DSZ.6) simultaneously for every original
place and index of every sufficiently late member of the family.
In particular the definition of $D_*$ becomes exactly
\[
 D_* =\mu_r I_{\rm row}\prod_\rho d_r(v_\rho,-b_\rho)
                                                               \tag{DSZ.13}
\]
whenever its original volume satisfies $\operatorname{vol}(Y_0)=\mu_r$.
The corresponding $d_*(u)=\mu_k I_{\rm col}\prod_\rho d_k(u_\rho)$
is already the original definition. The \texttt{Fact\allowbreak Dstar} and
\texttt{Fact\allowbreak Du\allowbreak Star} family fields identify these with the stored
positive $D_*,\allowbreak{} d_*(u)$ in the actual prime window.
\emph{Source: Z09, Z14--Z17.}
% END DSZ5

% BEGIN DSZ6
\paragraph{DSZ6: the original mass, index and factorial identity.}
\label{claim:DSZ6}
Assume the actual volume identity in DSZ5 and the mass-ratio theorem
already reconstructed in MF19 and RSL1--RSL3, with their named
classical inputs retained. Put
\[
 \begin{gathered}
 X_u=\frac1{fk}\log\frac{D_*}{d_*(u)},\quad
 \ell=\log(v/\pi),\quad G=\log(g_0g_1),\\
 A_m=(m-1)k+\frac{m(m-1)}2,\quad
 B_m=\frac{m(2k+m-1)}2.
 \end{gathered}
\]
By QW2 the exact index ratio, before weakening any Euler factors, is
\[
 \frac{I_{\rm row}}{I_{\rm col}}
  =(g_0g_1)^{fB_m}
       \prod_{i=0}^1\prod_{j=k+1}^{r}(1-g_i^{-j})^f.
\]
The equality follows by dividing the two exact finite-group
cardinalities in equations (QW.4)--(QW.6), not just their inequalities.
QW3 telescopes the actual positive masses, and its factorial
factors cancel those in (DSZ.13) and $d_*(u)$ exactly as in
(QW.15). Thus
\begin{align*}
 fkX_u={}&fA_m\ell-\frac{f(m-1)}2\log(2v)
     +\sum_{M=k+1}^{r}\log L_E(M)+fB_mG\\
 &+f\sum_{i=0}^1\sum_{j=k+1}^{r}\log(1-g_i^{-j})
     +\sum_\rho\log\frac{V_r(v_\rho,-b_\rho)}{V_k(u_\rho)}.
                                                               \tag{DSZ.14}
\end{align*}
Every logarithm is of a positive number; this was proved for the
Euler product in QW1, the masses in IG9/MF, the indices in QW2,
and the Vandermondes by the actual interlacing.
This is the source's \texttt{Lemma\allowbreak Five\allowbreak One\allowbreak Identity} after the
stronger degree equality (DSZ.6). The total factorial contribution
is zero, rather than an error estimated at size $k^2$.
\emph{Source: Z16--Z19.}
% END DSZ6

% BEGIN DSZ7
\paragraph{DSZ7: a short upper Vandermonde comparison on the actual cloud.}
\label{claim:DSZ7}
Fix one actual word and place satisfying DSZ4. Let $n=k-2$ and
compare the interior column coordinates $u_1,\ldots,\allowbreak{} u_{k-2}$
with $v_0,\ldots,\allowbreak{} v_{k-3}$. From (DSZ.9), for $1\le i<j\le k-2$,
\[
 u_i-u_j\ge\Delta(4(j-i)-1),\qquad
 v_{i-1}-v_{j-1}=4\Delta(j-i).
\]
Therefore, with empty products interpreted as one,
\[
 \log\frac{V_n(v_0,\ldots,\allowbreak{} v_{k-3})}
               {V_n(u_1,\ldots,\allowbreak{} u_{k-2})}
 \le\sum_{d=1}^{n-1}(n-d)\log\frac{4d}{4d-1}
 \le\frac{k}{3}(1+\log k).                                  \tag{DSZ.15}
\]
For the last inequality, use
$\log(1+t)\le t$, $4d-1\ge3d$, and
$\sum_{d=1}^{n-1}d^{-1}\le1+\log k$.
This harmonic sum is the reason the comparison costs $k\log k$
rather than $k^2$.

The remaining $k-2$ factors in $V_{k-1}(v)$ involve its last
coordinate and are at most $4k\Delta\le2B$.
The $(k-1)m$ cloud-to-bottom factors are at most $4B$, since
$0\le v_i\le4k\Delta\le2B$ and $0<b_j<2B$.
The $m(m-1)/2$ bottom factors are at most $2B$.
Every omitted endpoint factor in $V_k(u)$ is at least one by
DSZ4. Combining these facts gives
\begin{align*}
 \log\frac{V_r(v,-b)}{V_k(u)}
 \le{}&\frac{k}{3}(1+\log k)+(k-2)\log(2B)\\
     &+(k-1)m\log(4B)+\frac{m(m-1)}2\log(2B).
                                                               \tag{DSZ.16}
\end{align*}
Eventually $L\ge3$ and $\log k\le5\log x\le L$.
Then $\log(2B),\log(4B)\le2L$; using $2\le m\le k$,
the four terms in (DSZ.16) are bounded respectively by
$kmL/3$, $kmL$, $2kmL$, and $kmL$. In particular
\[
                 \log\frac{V_r(v,-b)}{V_k(u)}\le5kmL.
                                                               \tag{DSZ.17}
\]
This is an independent proof for the literal coordinates. It is
stronger than needed for the size bound and replaces only that
use of the source's more elaborate one-place upper estimate.
\emph{Source: Z12, Z14, Z20--Z21.}
% END DSZ7

% BEGIN DSZ8
\paragraph{DSZ8: two-sided logarithmic bounds with one family threshold.}
\label{claim:DSZ8}
Under the hypotheses of DSZ6, apply (DSZ.14).
We have $0\le\ell\le L$, $\log(2v)>0$, $G\le2C_0L$,
and $\sum_{M=k+1}^r\log L_E(M)\le fm$ by QW1.
The local-factor logarithms are nonpositive.
Also $A_m\le3km/2$, $B_m\le3km/2$, and (DSZ.17) bounds
the final place sum by $5fkmL$. Since $k\ge5$ and $L\ge1$,
$fm\le fkmL/2$. Dividing by $fk>0$ gives
\[
                         X_u\le(7+3C_0)mL.                   \tag{DSZ.18}
\]

For the lower bound use QW6 with the same actual coordinates,
$C_0$, and the exact (DSZ.13), which implies its weaker
$(2/3)^f$ hypothesis. All its remaining hypotheses were checked
in DSZ4. Put
$C_\epsilon=30+(2+2C_0)/\epsilon$.
The average of $\log(w_jv/\pi)$ is nonnegative, and (DSZ.8)
gives $mG\ge16mL$. Since $m^2L\le k$ and $L\ge1$,
\[
 X_u\ge16mL-C_\epsilon(L+m^2L/k)
       \ge16mL-2C_\epsilon L.
\]
Choose the family threshold also so that $m\ge2C_\epsilon$,
which follows from $m\ge x^{1/4}$. We obtain, simultaneously
for every original index and word,
\[
                    15mL\le X_u\le(7+3C_0)mL.                \tag{DSZ.19}
\]
The family-dependent constants $C_0,\allowbreak{} C_\epsilon$ precede this
eventual threshold, as the definition of \texttt{Tu\allowbreak Size} permits.
The original volume identity may itself have a family threshold;
take its maximum with the finitely many thresholds above.
The argument does not require a uniform value of this final
threshold for all possible families.
\emph{Source: Z17--Z22; QW1--QW6.}
% END DSZ8

% BEGIN DSZ9
\paragraph{DSZ9: obtain the literal prime-size predicate with the original constant.}
\label{claim:DSZ9}
Let $\Gamma_u$ be the actual $\texttt{Gam\allowbreak P}$.
The original \texttt{Gam\_\allowbreak eq} and \texttt{Size\allowbreak KRel} inequality
give $\log\Gamma_u\le fk$; also $\Gamma_u>1$ is an original
field. Thus the actual prime-window center satisfies
\[
 T_u^*=(D_*/d_*(u))^{1/(fk)}\Gamma_u^{-2/(fk)},\qquad
 X_u-2\le\log T_u^*\le X_u.                                  \tag{DSZ.20}
\]
The source selects its original prime in
$T_u^*\le T_u\le(1+1/(fk))T_u^*$.
Since $fk\ge1$, taking logarithms and using
$\log(1+1/(fk))\le1$ gives
\[
 X_u-2\le\log T_u\le X_u+1.
\]
By (DSZ.19), $m\ge2$ and $L\ge1$, the lower bound is at least
$mL/2$. The upper bound is at most $(8+3C_0)mL$.
Consequently
\[
 e^{mL/2}=B^{m/2}\le T_u,
 \qquad \log T_u\le(8+3C_0)mL.                               \tag{DSZ.21}
\]
These are the precise lower predicate \texttt{Tu\allowbreak Ge} and upper
inequality in \texttt{Tu\allowbreak Size}. In particular they prove the
source conclusion whenever
\[
 c_T(C_{YL})=\max\{11,16+3\max(5,\allowbreak{} C_{YL}+1)\}\le C_{Tu},       \tag{DSZ.22}
\]
using its original, larger constant without changing any chosen
prime or any family record. Only a threshold in $x$ was enlarged.

We have proved: every original supplied family with eventual
$\operatorname{vol}(Y_0)=\mu_r$ satisfies
$c_T(C_{YL})\le C_{Tu}\Longrightarrow\texttt{Tu\allowbreak Size}$.
All other premises were original family fields or the already
matched mass and formal-degree inputs. In particular no unresolved
Vandermonde variation, rank formula, local Hecke estimate, or
small-level rank theorem is used in this proof.
\emph{Source: Z15, Z17, Z23--Z25.}
% END DSZ9

% BEGIN DSZ10
\paragraph{DSZ10: the design-class row supplier and its quantifiers.}
\label{claim:DSZ10}
Take exactly the original construction class
$\mathcal A=\texttt{cls}(Q,\mathrm{Lk},\allowbreak{} E)$ and its common selector.
A \texttt{Built} witness is a context $K$, a point $\beta$, a
linked bundle $b'$ and an equality of the entire dependent triple
$(D,\allowbreak{} O,\allowbreak{} b)=(\beta.D,\beta.O,\allowbreak{} b')$.
At $\beta$, the stored core carries the actual
\texttt{Real\allowbreak Tied} witness. The subsequent \texttt{set\allowbreak Div} operation
changes only the division lattices and their prescribed heights;
it preserves the host, Hermitian spaces, degree function, compact-dual
measure and volume. Its four \texttt{Real\allowbreak Tied} clauses therefore
remain the same. Transport along the displayed dependent equality
gives \texttt{Real\allowbreak Tied} for the original $D$ of the selected bundle.
This is the content of \texttt{built\_\allowbreak val}; it is a projection and
transport of the supplied construction witness, not a new
existence assertion about a core.

Every member of $\mathcal A$ supplies such a \texttt{Built}
witness at every $x>\mathcal F.x_1$ and every actual $i$, by its
literal \texttt{QOf} and \texttt{Link} clauses. The parameter
positivity guard is discharged by that same Link's positivity
field. RSL1--RSL3 apply to this \texttt{Real\allowbreak Tied} and give
$\operatorname{vol}(Y_0)=\mu_r$ for those very objects, with
the already recorded GV Tamagawa and MF Hasse--Minkowski
applications retained. Thus DSZ9 applies on this class and on
every subclass obtained by adjoining \texttt{With\allowbreak Ch\allowbreak N},
\texttt{With\allowbreak Sm}, \texttt{With\allowbreak Tw}, or the target point condition.

For the source's \texttt{Closed} assertion choose
$\epsilon_1=1$ and the zero transported bound $K$, before the
family and $CH$ are given. Once a family in the specified class
is supplied, forget the added class conditions, project the
\texttt{Built} witness, and apply DSZ9 with its permitted
family-dependent eventual threshold. The \texttt{Large}
premise is unused. This proves the exact
\[
 \texttt{Closed}\bigl(\mathcal C_{\rm des},\,
       c_T(C_{YL})\le C_{Tu}\Longrightarrow\texttt{Tu\allowbreak Size}\bigr)
                                                               \tag{DSZ.23}
\]
for the one construction class and the original function
$c_T(\epsilon,\allowbreak{} y)=c_T(y)$ in (DSZ.22).
The same proof also works for an arbitrary class $\mathcal A$ and
selector $\sigma$ supplied with
$\texttt{Is\allowbreak Sel}(Q,\mathrm{Lk},\mathrm{True},\mathcal A,\sigma)$:
the definition of \texttt{Is\allowbreak Sel} supplies the required \texttt{Link}
and the \texttt{Built} witness for each literal selected bundle.
The proof above then applies without using the special definition
of $\texttt{cls}$. This establishes the universal premise of
\texttt{Row.\allowbreak of\_\allowbreak sel}, and hence the original
\texttt{row\_\allowbreak h\allowbreak Tu\_\allowbreak c\allowbreak Tmin} row with the trivial extra bundle predicate
$E=\mathrm{True}$.
The finite conjunction with other rows therefore imposes no new
bundle duty here. DEX6--DEX8 allow this proved row to be joined
before one family is selected.

\emph{Audit consequence.} The former \texttt{DESIGN\_\allowbreak TUSIZE\_\allowbreak CLOSED}
leaf is discharged on the selected construction-class route.
The independent compatible \texttt{Design\allowbreak Choices} producer remains;
(DSZ.23) does not construct a member of that class. The result also
does not assert a volume identity for arbitrary \texttt{Fam\allowbreak P}
records outside the construction class. These restrictions are
part of the theorem just proved, not new assumptions to be
silently supplied by positivity.
\emph{Source: Z26--Z29.}
% END DSZ10


\section{The original class number and the column-rank producer}
\label{sec:rank-producer}

\status{RMP1--RMP11 prove the literal \texttt{Class\allowbreak Number\allowbreak Formula}$\,{}_2$
for every supplied original $D:\texttt{D7}$ and
$\mathrm{Lv}:\texttt{Level\allowbreak Data}(D)$. They then prove the original
\texttt{Rank\allowbreak M} from these records and level-zero column forms on the
same $D$. The source's already supplied \texttt{TSel}, or its original
\texttt{Raw\allowbreak Sel Bundle}, supplies those forms with one common threshold.
For the selected construction class, its membership predicate itself
supplies this Bundle through \texttt{cls\_\allowbreak at} and \texttt{pick}.
Thus the separate rank-formula producer is discharged there without
a further selector premise. This does not construct a member of the
common arithmetic family class.
The mass input is the previously reconstructed MF18--MF19; the genus
dictionary is GV11. No new mass theorem or Tamagawa theorem is invoked.}

\paragraph{Notation and exact target.}
Fix the original datum $D$ and a word $u$. Write
\[
 \begin{gathered}
 E=D.\mathrm{host}.E,\quad F=D.\mathrm{host}.F,\quad
 k=D.R.k,\quad f=D.R.f=[F:\mathbb Q],\\
 T=(D.\mathrm{mem}(u)).A.T,\qquad g_i=D.\mathrm{host}.g(i).
 \end{gathered}
\]
The Hermitian form is conjugate linear in its first variable.
In the original coordinates $\mathrm{crd}_0$ it has Gram matrix
$J=D.\mathrm{sp}.\mathrm{gram}_0$. Put
\[
 \begin{gathered}
 G=\mathrm U(J)(\mathbb A_{F,\allowbreak{} f}),\qquad
 \Gamma=\mathrm U(J)(F),\qquad L=\mathrm{Lv}(u).\mathrm{Lb},\\
 U=\mathrm{Lv}(u).U,\qquad C_u=(D.\mathrm{mem}(u)).C.\mathrm{cls}(0).
 \end{gathered}
                                                               \tag{RMP.1}
\]
Here $\mathrm U(J)(F)$ means the rational unitary matrices with entries
in $E$; it is the source's \texttt{rat\allowbreak Pts} inside \texttt{UPts}.
The notation $L^{\#}$ denotes the actual Hermitian dual
\texttt{dual\allowbreak Set}, which is already an $\mathcal O_E$-module and
therefore equals its span \texttt{dual\allowbreak Lb}.
The target class number is
\[
 |C_u|=\mathrm{massI}(D,\allowbreak{} k)\,
       \mathrm{idxPi}(D,\allowbreak{} g_0,\allowbreak{} g_1)(T^k-1)^f.              \tag{RMP.2}
\]
The rank target is the literal definition
\[
 \mathrm{RankM}(D):\quad
 D.\mathrm{duS}(u)=\mathrm{duStarU}(D,\allowbreak{} u)(T^k-1)^f
       \quad\hbox{for every }u,                        \tag{RMP.3}
\]
where $D.\mathrm{duS}(u)=\dim_{\mathbb C}(D.\mathrm{mem}(u)).A.\mathrm{Mu}_0$.
All equalities involving the source's real-valued factors use the
natural embedding of the integer cardinalities into $\mathbb R$.

% BEGIN RMP1
\paragraph{RMP1: the exact base, primes and level conditions.}
\label{claim:RMP1}
The original $\mathrm{ColLevel}$ gives a finitely generated full lattice $L$ and
\[
 L\subset L^{\#},\qquad T L^{\#}\subset L,\qquad
 |L^{\#}/L|=T^{2f}.                                    \tag{RMP.4}
\]
It gives the local groups $U_v$, with $U_v\subset\operatorname{Stab}(L_v)$,
and identifies $U$ with the unitary adeles having components in all
$U_v$. Outside the primes above $g_0,\allowbreak{} g_1,\allowbreak{} T$ it gives
$U_v=\operatorname{Stab}(L_v)$. Above $g_i$, an actual split projection
and an actual frame identify $U_v$ with the inverse image of the
upper-unitriangular group under reduction
$\mathrm{GL}_k(\mathcal O_{F_v})\to\mathrm{GL}_k(k_v)$.
Above $T$ it gives an $\mathcal O_{E,\allowbreak{} v}$-module $M_v$ of additive
cardinality $T$ and a nonzero linear functional
\[
 t_v:L_v^{\#}\longrightarrow M_v,\qquad
 t_v(L_v)=0,\qquad t_v(TL_v^{\#})=0,
\]
and, literally,
\[
 U_v=\operatorname{Stab}(L_v)\cap\operatorname{Stab}(t_v).
                                                               \tag{RMP.5}
\]
The last stabilizer fixes the functional itself, not just its kernel.
These are stored fields, not proposed replacements for the level.
[Sources R01--R03.]

The source's $g_0,\allowbreak{} g_1$ are distinct rational primes. They split
completely in $E$: the field \texttt{nd.\allowbreak g\_\allowbreak split} says that the
number of primes is $[E:\mathbb Q]$. The degree formula
$\sum_{P\mid g_i}e_P f_P=[E:\mathbb Q]$, with positive integral
summands, forces every $e_P=f_P=1$. The ramification and residue-degree
multiplication formulas show the same in $F$; hence there are exactly
$f$ places of $F$ over $g_i$, each of residue cardinality $g_i$.

For $T$ the source stores $\pi\bar\pi=T$ and
$\pi\equiv1\pmod{N\mathcal O_E}$, where
$N=D.\mathrm{host}.\mathrm{QS}\,\allowbreak{} q\,\allowbreak{} d\,(g_0g_1)^2$.
Conjugation preserves this congruence, so $T-1\in N\mathcal O_E$.
A rational algebraic integer is an integer; thus $N\mid T-1$ in
$\mathbb Z$. In particular $T\ne g_0,\allowbreak{} g_1$ and
$T\equiv1\pmod{\mathrm{condN}}$.
The original host's Frobenius/Artin fields identify the Frobenius at
this unramified prime with the Artin image of its residue class.
That class is $1$, so Frobenius is $1$; the stored unramifiedness
makes inertia trivial. The decomposition group, generated modulo
inertia by Frobenius, is trivial. Thus $T$ splits completely in $E$
and in $F$, proving the source's \texttt{Splits\allowbreak F} statements without
a prime-existence assertion. Equivalently, the later argument can
retain exactly the two \texttt{Splits\allowbreak F} hypotheses of (RMP.2).
[Sources R04--R06.]

Finally, the base space $D.\mathrm{sp}.X(0)$ is totally positive
definite, not just definite of an unspecified sign. Its rank is $k$;
its supplied definite signature is either $(k,0)$ or $(0,\allowbreak{} k)$.
It embeds isometrically into $Y(0)$, whose signature is
$(k+r_S,\allowbreak{} r_S+1)$ with $r_S+1=m/2<k$. A negative-definite subspace of
rank $k$ would exceed that negative index. This proves the positive
alternative at every embedding by Sylvester's law of inertia.
The source's unramified CM condition is exactly the condition treated
in MF1 and DR1, including the dyadic places. [Sources R02, R07.]
% END RMP1

% BEGIN RMP2
\paragraph{RMP2: the actual local discriminant module has length two.}
\label{claim:RMP2}
Fix $v\mid T$, put $R=\mathcal O_{F_v}$, let $\varpi$ be its
uniformizer and let $k_v=\mathbb F_T$. Write
$B=\mathcal O_E\otimes_{\mathcal O_F}R$.
The unramified quadratic algebra $B$ is either $R\times R$ or an
unramified quadratic DVR. The action on the cyclic additive group
$M_v$ of prime order defines a unital ring map
$B\to\operatorname{End}_{\mathbb Z}(M_v)=\mathbb F_T$.
Its restriction to $R$ is reduction: its kernel is the unique maximal
ideal. An unramified quadratic residue field cannot embed in
$\mathbb F_T$. Thus $B=R\times R$, and the action on $M_v$ selects
one factor. This gives the exact ordering of the split factors
used by the source's \texttt{exists\_\allowbreak local}; it is not an arbitrary
choice subsequently imposed on $t_v$.

Let $e_+,\allowbreak{} e_-$ be the corresponding integral idempotents, with
$e_+M_v=M_v$ and $e_-M_v=0$. Under the two projections write
\[
 L_v=A_v\oplus B_v,\qquad
 L_v^{\#}=B_v^{\vee}\oplus A_v^{\vee}.                 \tag{RMP.6}
\]
The superscript $\vee$ uses the nondegenerate cross-pairing given
by $J$. The two same-factor pairings vanish because conjugation
exchanges $e_+$ and $e_-$. The two projected modules are free
$R$-lattices of rank $k$, and the two quotient lengths
\[
 d_v=\operatorname{length}_R(B_v^{\vee}/A_v)
     =\operatorname{length}_R(A_v^{\vee}/B_v)           \tag{RMP.7}
\]
are equal. To check this directly, choose a basis of $B_v$, use its
dual basis for $B_v^{\vee}$, and diagonalize the full sublattice
$A_v$ by elementary divisors. If the diagonal exponents are
$a_1,\ldots,\allowbreak{} a_k\ge0$, its dual enlarges $B_v$ by
$\varpi^{-a_1},\ldots,\varpi^{-a_k}$. Both lengths are
$\sum_j a_j$. This also proves that duality commutes with the
localization used in (RMP.6).

Because $t_v\ne0$, its image is all of the prime-order module
$M_v$. It factors through the first quotient in (RMP.6), so $d_v\ge1$.
Away from $T$, (RMP.4) makes $L_v^{\#}=L_v$. Localization and the
Chinese remainder theorem for the finite module $L^{\#}/L$ give
\[
 T^{2f}=|L^{\#}/L|=\prod_{v\mid T}T^{2d_v}.
                                                               \tag{RMP.8}
\]
There are exactly $f$ such places, and all $d_v\ge1$. Taking
integer exponents in (RMP.8) gives $d_v=1$ at every one of them.
Thus $t_v$ induces an isomorphism
$B_v^{\vee}/A_v\simeq M_v$ and kills the entire second component.
In particular its first-component kernel is exactly $A_v$.
The two length-one quotients also give the exact $B$-linear
discriminant-module isomorphism to the product of the two simple
residue modules in \texttt{Local\allowbreak At\allowbreak T}.
This is the kernel equality for which the source uses its
\texttt{mem\_\allowbreak of\_\allowbreak count} argument. [Sources R01, R08--R10.]
% END RMP2

% BEGIN RMP3
\paragraph{RMP3: construct the self-dual enlargement on the original base.}
\label{claim:RMP3}
At $v\mid T$ set
\[
 N_v=B_v^{\vee}\oplus B_v
     =e_+L_v^{\#}+e_-L_v;
 \qquad N_v=L_v\quad(v\nmid T).                       \tag{RMP.9}
\]
These local lattices are self-dual by the cross-pairing in (RMP.6).
Define the global lattice by the literal intersection formula
\[
 N=\{x\in L^{\#}:e_-x_v\in L_v\text{ for every }v\mid T\}.
                                                               \tag{RMP.10}
\]
This is exactly \texttt{over\allowbreak Lat} with the chosen negative-factor
idempotents. It contains $L$ and lies in $L^{\#}$. Multiplication by
$T$ embeds $L^{\#}$ into the finitely generated $\mathcal O_E$-module $L$;
since $\mathcal O_E$ is Noetherian, $L^{\#}$ and its submodule $N$
are finitely generated. Their spans are full because
they contain $L$.

The completion of (RMP.10) is precisely (RMP.9). One way to see both
inclusions without any approximation assumption is to use
$Q=L^{\#}/L$. It is a finite module supported above $T$.
Its canonical Chinese-remainder decomposition is the direct sum of
its localizations $Q_v$. Choose in each $Q_v$ the submodule
$e_+Q_v$, and take the inverse image of their direct sum in
$L^{\#}$. This inverse image is (RMP.10), and its localizations
are (RMP.9); completing these finite-length quotients changes
nothing. This is also the finite-lattice patching argument of LP1.

Membership of an element of $E$ in $\mathcal O_E$ is tested at every
finite place. Applying that fact to all pairings, and using
$N_v^{\#}=N_v$, gives $N^{\#}=N$. Thus $N$ is an actual full
self-dual lattice on the original Hermitian space.

At $v\mid T$, a unitary map preserving $L_v$ preserves its dual
and its two factors, hence preserves $N_v$. Conversely a unitary
map preserving $N_v$ is determined by its action on $B_v^{\vee}$,
with the adjoint inverse action on $B_v$. Let
\[
 f_v:B_v^{\vee}\longrightarrow k_v
                                                               \tag{RMP.11}
\]
be the original $t_v$ transported through any additive
$M_v\simeq k_v$ respecting the first-factor action. Its kernel is
$A_v$. Fixing $f_v$ therefore preserves $A_v$, and duality also
preserves $A_v^{\vee}$. Consequently (RMP.5) is exactly
\[
 U_v=\{a\in\operatorname{Stab}(N_v):f_v\circ a_+=f_v\}.
                                                               \tag{RMP.12}
\]
Away from $T$, $N_v=L_v$ and the original inclusion
$U_v\subset\operatorname{Stab}(L_v)$ suffices. Hence
$U\subset K_N:=\operatorname{Stab}(\widehat N)$.
All assertions of \texttt{Is\allowbreak Over\allowbreak Lattice}, the projection and
functional clauses of \texttt{Local\allowbreak At\allowbreak T}, and the inclusions in
\texttt{Has\allowbreak Over\allowbreak Lattice} now hold on this very $N$.
The unique unitary lifting of a first-factor matrix is given by
its adjoint inverse on the second factor; a basis of
$B_v^{\vee}$ gives the required frame. [Sources R08--R11.]
% END RMP3

% BEGIN RMP4
\paragraph{RMP4: the local index above $T$ is $T^k-1$.}
\label{claim:RMP4}
In a basis of $B_v^{\vee}$ the first split projection identifies
$K_{N,\allowbreak{} v}$ with $\mathrm{GL}_k(R)$. It identifies $f_v$ with a
nonzero row $r_v\in k_v^k$. Indeed $f_v$ is surjective and its
kernel has quotient of cardinality $T$, so it factors through
reduction and remains nonzero there. Equation (RMP.12) identifies
$U_v$ with the inverse image of
\[
 \{a\in\mathrm{GL}_k(k_v):r_va=r_v\}.                \tag{RMP.13}
\]
The reduction map from $\mathrm{GL}_k(R)$ is onto: lift any
invertible residue matrix entrywise; its determinant is a unit.
The group $\mathrm{GL}_k(k_v)$ acts transitively on the nonzero
rows, since any such row can be extended to a dual basis. Its
orbit therefore has $T^k-1$ elements. The ordinary orbit-stabilizer
bijection now gives
\[
                       [K_{N,\allowbreak{} v}:U_v]=T^k-1.           \tag{RMP.14}
\]
There is no division by $T-1$: the level fixes a specified nonzero
functional, not its projective line. This proves the exact
\texttt{Index\allowbreak At\allowbreak T} and local-index conclusion used in the original
class-number calculation. [Sources R11--R12.]
% END RMP4

% BEGIN RMP5
\paragraph{RMP5: auxiliary indices and their exact adelic product.}
\label{claim:RMP5}
At a place above $g_i$, the actual split projection and frame give
$U_v$ as the inverse image of the upper-unitriangular subgroup
$N_k(\mathbb F_{g_i})$. This is a subgroup: multiplication preserves
zero entries below the diagonal and diagonal entries equal to one,
and the inverse of $1+X$ is the finite series
$1-X+\cdots+(-X)^{k-1}$ for strictly upper-triangular $X$.
It has $g_i^{k(k-1)/2}$ elements. Relative to the maximal compact
in that frame, reduction therefore gives the index
\[
 \frac{|\mathrm{GL}_k(\mathbb F_{g_i})|}
      {|N_k(\mathbb F_{g_i})|}
       =\mathrm{idxColLoc}(g_i,\allowbreak{} k).                    \tag{RMP.15}
\]
That maximal compact and $K_{N,\allowbreak{} v}$ need not have been specified by
the same frame. Both contain $U_v$, and they are conjugate in the
local unitary group: the split first component is simply
$\mathrm{GL}_k(F_v)$, and any two full lattice bases are related
by an invertible matrix. GV1 proves this group unimodular, so its
Haar measure is invariant under conjugation. The two compact groups
have equal measure. Since $U_v$ is compact open, dividing by its
positive measure proves that (RMP.15) is also $[K_{N,\allowbreak{} v}:U_v]$.

Away from $g_0,\allowbreak{} g_1,\allowbreak{} T$, (RMP.9) and the original
\texttt{Uloc\_\allowbreak hyp} give $U_v=K_{N,\allowbreak{} v}$.
The three sets of exceptional places are disjoint by RMP1.
The map from $K_N/U$ to the product of local right-coset sets is
bijective. Equality is tested locally by the definition of $U$.
For surjectivity choose representatives at the finitely many
exceptional places and put the identity elsewhere; that is an
actual finite unitary adele preserving $\widehat N$.
There are $f$ places over each of the three rational primes.
Thus
\[
 \boxed{[K_N:U]=\mathrm{idxPi}(D,\allowbreak{} g_0,\allowbreak{} g_1)(T^k-1)^f.}   \tag{RMP.16}
\]
Both factors are positive and the index is finite.

This also proves the source's stronger
\texttt{Stabiliser\allowbreak Level\allowbreak Index\allowbreak Formula}, which quantifies over
\emph{any} full self-dual $N'$ with $U\subset K_{N'}$.
DR1 and the split dual-basis construction give a local unitary
isometry $N_v\to N'_v$ at every finite place. The lattices agree
outside a finite set, so these isometries assemble to an adele
conjugating $K_N$ to $K_{N'}$. GV1 and GV12 give equal Haar
measures. Dividing by the measure of the same compact open $U$
shows $[K_{N'}:U]=[K_N:U]$ and preserves (RMP.16).
No global integral orthonormal basis of $N$ or $N'$ is assumed.
[Sources R12--R14.]
% END RMP5

% BEGIN RMP6
\paragraph{RMP6: the chosen self-dual genus has exactly the required mass.}
\label{claim:RMP6}
Transport $N$ through the original coordinate isomorphism:
\[
              M=\mathrm{crd}_0^{-1}(N)
                  \subset(D.\mathrm{sp}.X(0)).V.       \tag{RMP.17}
\]
The displayed identity
$h(x,\allowbreak{} y)=\operatorname{hform}(J,\mathrm{crd}_0x,\mathrm{crd}_0y)$
shows that finite generation, full span and self-duality are
preserved. Thus (RMP.17) is precisely an admissible $M$ in the
source first-space mass identity, with $\mathrm{crdLat}_0(M)=N$. We have not asserted that the original
standard lattice $X(0).L_0$ is fixed by $U$; (RMP.17) supplies the
self-dual lattice for which that property was proved.

MF18 gives, for every positive definite rank-$k$ integral lattice
in this same unramified CM extension,
$\operatorname{mass}(M)=c_k\lambda(M)$.
Self-duality makes the local discriminant length $d_M$ zero at
every place. MF6 and MF14 then give
$\lambda_v(M)=1$, including the dyadic places; hence
$\lambda(M)=1$. The actual standard coordinate form and lattice
in the definition of $\mathrm{massI}(D,\allowbreak{} k)$ have the same property.
MF19 identifies their mass with $c_k$. Therefore
\[
             \operatorname{mass}(M)=\mathrm{massI}(D,\allowbreak{} k).
                                                               \tag{RMP.18}
\]
This argument uses the already proved uniformity in the rational
Hermitian space in MF18. It does not posit a rational or integral
isometry from the supplied positive definite space to the standard
one. Its inherited named classical input is exactly MF4's
Hasse--Minkowski theorem; no additional absolute mass formula is
needed. [Sources R02, R07, R15--R16; inherited MF6, MF14,
MF18--MF19.]
% END RMP6

% BEGIN RMP7
\paragraph{RMP7: the actual double-coset dictionary and finite stabilizers.}
\label{claim:RMP7}
For $a\in G$, intersect the adelic lattice $a\widehat N$ with
$E^k$. LP1 shows that its localizations are the prescribed ones.
The map
\[
 [a]\longmapsto[\,\allowbreak{} E^k\cap a\widehat N\,]
                                                               \tag{RMP.19}
\]
is the exact bijection from $\Gamma\backslash G/K_N$ to the
genus classes of $N$. Changing $a$ on the right by $K_N$ changes
no lattice; changing it on the left by $\Gamma$ gives a rational
isometry. Conversely equality up to a rational isometry is tested
at every completion and gives the required double-coset equality.
Every genus lattice supplies local isometries, equal to the identity
outside finitely many places, so it supplies an adele. If one uses
the source's rational-local definition of genus, MF3's Cayley
weak-approximation argument meets the relevant open local
transporter and gives the same local lattice. This is the precise
dictionary proved in GV11, not a strong-approximation assertion.

The rational stabilizer
$\Gamma_a=\Gamma\cap aK_Na^{-1}$ is exactly the automorphism
group of $E^k\cap a\widehat N$. It is finite. Indeed this is a
full integer lattice, and the sum over embeddings of its positive
Hermitian norms is a positive-definite real quadratic form.
An automorphism takes each member of an integer basis into a
discrete lattice with its norm fixed. A bounded set contains only
finitely many lattice points, so there are only finitely many
possible images of that basis.
The quotient $\Gamma\backslash G/K_N$ is also finite: it is the
surjective image of the original finite $C_u\simeq\Gamma\backslash G/U$
under enlargement of the right subgroup. Thus the genus mass is
literally the finite sum
\[
 \sum_{[a]\in\Gamma\backslash G/K_N}|\Gamma_a|^{-1}
             =\mathrm{massI}(D,\allowbreak{} k),                    \tag{RMP.20}
\]
using (RMP.18). The source's fallback convention for mass is not
entered, because actual finite representatives have been supplied.
[Sources R01, R15, R17--R18; inherited GV11, LP1, MF3.]

We also need the smaller level stabilizers to be trivial.
For $b\in G$, the group $\Gamma\cap bUb^{-1}$ preserves the
full lattice $E^k\cap b\widehat L$, so the same argument makes
it finite. At a place above $g_i$, conjugating by the original
frame puts every element into an integral matrix group whose
reduction is upper unitriangular.
If an element there has finite order $g_i^a c$ with
$(c,\allowbreak{} g_i)=1$, its $g_i^a$th power has prime-to-$g_i$ order and
upper-unipotent reduction. That reduction must be the identity.
Write the power as $1+X$, $X\in\varpi M_k(R)$. If
$X\in\varpi^rM_k(R)$, its $c$th power is
$1+cX$ modulo $\varpi^{r+1}$. Since $c$ is a unit, induction
puts $X$ in every power of $\varpi$, whence $X=0$.
Its order is consequently a power of $g_i$. The projection of a
rational matrix into either split completion is injective, so its
rational order is a power of both distinct $g_0,\allowbreak{} g_1$ and equals one.
This proves
\[
                    \Gamma\cap bUb^{-1}=\{1\}
                         \quad(b\in G).               \tag{RMP.21}
\]
The original \texttt{stab\allowbreak I\_\allowbreak inj} and \texttt{stab\allowbreak I\_\allowbreak range} fields
identify the stored $C_u$ stabilizers with these groups; hence
$(D.\mathrm{mem}(u)).C.\mathrm{StabZero}$ holds.
This is CR4 on the actual \texttt{Col\allowbreak Level}. It needs both given
auxiliary primes and definiteness, but no lower bound called
\texttt{GBig}. [Sources R01, R19; inherited CR4 and RSL7.]
% END RMP7

% BEGIN RMP8
\paragraph{RMP8: count the fibers and obtain \texttt{Class\allowbreak Number\allowbreak Formula}$\,{}_2$.}
\label{claim:RMP8}
For a representative $a$ of a class in $\Gamma\backslash G/K_N$,
its fiber in $\Gamma\backslash G/U$ is the orbit set of $K_N/U$
under $a^{-1}\Gamma_a a$. The action is free: an element fixing
$kU$ belongs after conjugation to $\Gamma\cap akU(ak)^{-1}$,
which is trivial by (RMP.21). Every orbit has $|\Gamma_a|$
elements, so the fiber cardinality is $[K_N:U]/|\Gamma_a|$.
Summing these finite cardinalities gives
\[
 \begin{split}
 |\Gamma\backslash G/U|
   &=[K_N:U]\sum_{[a]\in\Gamma\backslash G/K_N}|\Gamma_a|^{-1}\\
   &=\mathrm{idxPi}(D,\allowbreak{} g_0,\allowbreak{} g_1)(T^k-1)^f\,
                                      \mathrm{massI}(D,\allowbreak{} k).
 \end{split}                                                     \tag{RMP.22}
\]
The bijection here is the original \texttt{Lv(u).\allowbreak y\_\allowbreak bij} from
$C_u$; its representatives and stabilizer injections were used in
RMP7. Thus (RMP.22) proves exactly (RMP.2).

Introducing the hypotheses of the source definition therefore
proves its literal class-number predicate
$\texttt{Class\allowbreak Number\allowbreak Formula}_2(D,\mathrm{Lv})$.
Its initial \texttt{GBig} implication is harmless: the proof
has established a stronger assertion without it. All splitting
hypotheses have either been retained during the local count or
checked from $D$ in RMP1. In particular, for every original family,
every $x>\mathcal F.x_1$ and every actual $i$, this proves the
class-number formula for
$(\mathcal F.D(x,\allowbreak{} i),(\mathcal F.\mathrm{ex}(x,\allowbreak{} i)).\mathrm{lev})$.
This is the original \texttt{hn4} slot of \texttt{link\_\allowbreak ex},
with no eventual arithmetic producer left inside that slot.
[Sources R03, R15, R20--R21.]
% END RMP8

% BEGIN RMP9
\paragraph{RMP9: the actual level-zero values have the stated Weyl dimension.}
\label{claim:RMP9}
Supply $\mathcal T_u:\texttt{Forms0}(D,\allowbreak{} u)$ on this same $D$.
Its \texttt{val0.\allowbreak e\allowbreak Y} identifies the original value-variable set
with the original frame space, and \texttt{val0.\allowbreak Fi}, together
with \texttt{Fi\_\allowbreak zero}, identifies the original $\mathrm{Mu}_0$
with the restriction of adelic forms having polynomial coefficient
values of weights
$\lambda_\tau=(D.\mathrm{fm}(u)).\mathrm{lam}(0,\tau)$.
This is an equality of the specified subspaces after that displayed
linear change of variables; a separate representation of the same
abstract dimension would not suffice. [Sources R22--R24.]

RMP7 made the stored stabilizers $C_u.\mathrm{stab}(0,\allowbreak{} c)$ trivial.
The \texttt{cls0.\allowbreak stab\allowbreak I\_\allowbreak range} field of $\mathcal T_u$ therefore
makes its own rational level stabilizers trivial as well. Its
\texttt{cls0.\allowbreak y\_\allowbreak bij} is indexed by precisely the same $C_u$.
Consequently arbitrary coefficient values at these representatives
extend uniquely to an adelic form. Explicitly, write
$x=\gamma y_c\kappa$ and set
\[
       F(x,\allowbreak{} z)=\chi(\kappa)\,\allowbreak{} v_c(\gamma^{-1}z).
                                                               \tag{RMP.23}
\]
The class determines $c$, and the trivial stabilizer makes
$\gamma,\kappa$ unique. The character law and the coefficient
action give the required equivariance; the invertible coefficient
action preserves the polynomial coefficient space. Restriction
and this extension are inverse linear maps. This is the argument
of CR2 and CR5--CR6, now applied to the literal \texttt{Mu0\_\allowbreak iff}
and the literal $C_u$.

At each embedding the coefficient space is the characteristic-zero
Weyl space with dimension
\[
 \prod_{0\le a<b<k}
         \frac{\lambda_\tau(a)-\lambda_\tau(b)+b-a}{b-a}.
                                                               \tag{RMP.24}
\]
CB2--CB3 and IW3/CR1 give this dimension for the actual polynomial
span, including determinant shifts. Restricting its polynomial
functions to invertible frames is injective because those frames
form a Zariski-open dense set; multiplying by the nonzero determinant
power preserves injectivity. Tensoring the coefficient spaces over
$\tau$ multiplies their dimensions.
The original weight equation is
\[
 \lambda_\tau(a)=u_\tau(a)-\left(\frac{k-1}{2}-a\right),
       \qquad u_\tau=\mathrm{uColR}(D,\allowbreak{} u,\tau).
                                                               \tag{RMP.25}
\]
Thus the numerator in (RMP.24) is exactly $u_\tau(a)-u_\tau(b)$;
its denominator product is $\prod_{j=0}^{k-1}j!$.
These are precisely the source's \texttt{vand} and
\texttt{fac\allowbreak Prod}. We obtain
\[
 \boxed{\dim_{\mathbb C}\mathrm{Mu}_0
            =|C_u|\prod_\tau\mathrm{dimV}(\mathrm{uColR}(D,\allowbreak{} u,\tau)).}
                                                               \tag{RMP.26}
\]
This proves the original \texttt{Mu0Tie}. For this dimension
comparison the two supplied level descriptions need only be
indexed by the same stored $C_u$ with its stored stabilizers and
weights; no additional integral or metric equality is used.
The original \texttt{Cols\allowbreak Tied\allowbreak Sym} supplies stronger same-carrier
identities when they are needed elsewhere. [Sources R22--R26.]
% END RMP9

% BEGIN RMP10
\paragraph{RMP10: the literal \texttt{Rank\allowbreak M} and the stored $d_\star$.}
\label{claim:RMP10}
Supply $\mathrm{Lv}:\texttt{Level\allowbreak Data}(D)$ and
$\mathcal T_u:\texttt{Forms0}(D,\allowbreak{} u)$ for every word, all on the
original $D$. Multiply (RMP.2) by the product in (RMP.26).
The source defines
\[
 \mathrm{duStarU}(D,\allowbreak{} u)
   =\mathrm{massI}(D,\allowbreak{} k)\,\mathrm{idxPi}(D,\allowbreak{} g_0,\allowbreak{} g_1)
        \prod_\tau\mathrm{dimV}(\mathrm{uColR}(D,\allowbreak{} u,\tau)).
                                                               \tag{RMP.27}
\]
Together with $D.\mathrm{duS}(u)=\dim_{\mathbb C}\mathrm{Mu}_0$,
this gives (RMP.3) for every word. This is a pointwise proof of the
original \texttt{Rank\allowbreak M}, not an asymptotic estimate and not the
ramified dimension $(D.\mathrm{mem}(u)).A.\mathrm{du}$ of RSL10.

For an original family member with $x>\mathcal F.x_1$, its actual
extra record supplies
\[
  ((D.\mathrm{mem}(u)).A.\mathrm{duStar}:\mathbb R)
       =\mathrm{duStarU}(D,\allowbreak{} u).                         \tag{RMP.28}
\]
Consequently (RMP.3) is also the literal
$d_M(u)=d_\star(u)(T_u^k-1)^f$ used in DEX5.
No alternative standard lattice, different mass normalization,
or independently chosen number called $d_\star$ has been substituted.
For any supplied $C:\texttt{Cols\allowbreak Tied\allowbreak Sym}(D)$, take
$\mathcal T_u=(C.\mathrm{forms}(u)).\mathrm{to0}$.
TG12 supplies this $C$ on its original true-column input with its
stated laws and later-stage record. RMP10 adds no existence demand
beyond those inputs. [Sources R03, R26--R28.]
% END RMP10

% BEGIN RMP11
\paragraph{RMP11: common-selector thresholds and the closed producer.}
\label{claim:RMP11}
Let $\mathcal A$ be the original source class and retain its
specified $s:\texttt{TSel}(\mathcal A)$. For a family
$\mathcal F\in\mathcal A$, write $a_s=(s\mathcal F).1$.
The definition supplies
\[
 R_{x,\allowbreak{} i}:\texttt{Rec\allowbreak T}(\mathcal F.\mathrm{ex}(x,\allowbreak{} i))
 \quad\text{for all }x>\mathcal F.x_1,\ a_s\le x,
                         \ i\in\mathcal F.I_x.         \tag{RMP.29}
\]
Its actual field $R_{x,\allowbreak{} i}.T.\mathrm{cols}$ is a
$\texttt{Cols\allowbreak Tied\allowbreak Sym}(\mathcal F.D(x,\allowbreak{} i))$. Use these very columns,
and the level data from that very $\mathcal F.\mathrm{ex}(x,\allowbreak{} i)$,
in RMP10. Thus with
\[
                   a=\max\{a_s,\mathcal F.x_1+1\}
                                                               \tag{RMP.30}
\]
we have, in the original quantifier order,
\[
 \forall x\ge a\ \forall i\in\mathcal F.I_x\
                      \mathrm{RankM}(\mathcal F.D(x,\allowbreak{} i)).
                                                               \tag{RMP.31}
\]
The threshold precedes $i$ and every word $u$. No \texttt{Duty},
finite-character assertion, theta square, metric transport, or
additional size hypothesis is needed for this rank consumer.
For \texttt{Raw\allowbreak Sel Bundle}, the source defines $s$ with threshold
$\mathcal F.x_1$ and selects exactly the Bundle's $R$ field;
then $a=\mathcal F.x_1+1$ suffices. [Sources R29--R31.]

Fix the target condition $\mathcal G=G.\mathrm{cls}(\lambda)$.
The selected design classes are
\[
 \mathcal C^{\prime}=\texttt{With\allowbreak Ch\allowbreak N}(\mathcal A\cap\mathcal G),
 \qquad
 \mathcal C_{\mathrm{des}}=\texttt{With\allowbreak Tw}(\texttt{With\allowbreak Sm}
                 (\mathcal C^{\prime})).
\]
They map to $\mathcal A$ by their literal conjunction projections. Apply (RMP.31) to the same family
under this projection. In the definition of \texttt{Closed}, choose
$\varepsilon_1=1$ and the transported bound with all components zero.
For every $0<\varepsilon\le1$, every family and every $C_H$,
class membership permits evaluation of the fixed $s$; the argument works regardless of
its additional \texttt{Large} proof. This proves
\[
 \texttt{Closed}\bigl(\mathcal C_{\mathrm{des}},\
   \mathcal F\mapsto
     \exists a\ \forall x\ge a\ \forall i,
                     \mathrm{RankM}(\mathcal F.D(x,\allowbreak{} i))\bigr).
                                                               \tag{RMP.32}
\]
The all-member class-number assertion in RMP8 is itself closed
with the same choices, even before $s$ is supplied.
If one instead has only a \emph{closed} eventual Forms0 supplier,
its own $\varepsilon$ and transported bound must be retained;
joining that supplier with the finite implication preserves them
and only replaces the inner threshold by its maximum with
$\mathcal F.x_1+1$. They cannot be reset to zero in that weaker scope.

For the actual construction class there is no additional existence
premise for this selector. Put
$\mathcal A=\texttt{cls}(Q,\allowbreak{} Lk,\allowbreak{} E_{\mathrm{row}})$ with
$Bd=\texttt{Bundle}$. A proof $h:\mathcal A(\mathcal F)$
contains a parameter $\kappa$, a \texttt{Link} from $\kappa$ to
$\mathcal F$, and the pointwise \texttt{QOf} property.
The same \texttt{Link} contains $\kappa.\mathrm{Pos}$, so the
positivity guard in \texttt{QOf} is discharged. For each
$x>\mathcal F.x_1$ and $i$, evaluate \texttt{QOf} at
$c=\mathcal F.C_{\mathrm{sz}}$. It supplies a Bundle on the
literal dependent pair
$(\mathcal F.D(x,\allowbreak{} i),\texttt{obj\allowbreak At}(\mathcal F,\allowbreak{} x,\allowbreak{} i))$.
Choose these witnesses pointwise. This is precisely
\texttt{cls\_\allowbreak at} followed by the source's \texttt{pick}, giving
$\sigma:\texttt{Raw\allowbreak Sel Bundle}(\mathcal A)$.
The \texttt{Class\allowbreak Sel} construction uses this very class and this
very pick (with its actual combined slot in place of $Q$).
The selector of RMP11 is then the \texttt{sel} of $\sigma$:
its selected record is exactly $(\sigma(\mathcal F,\allowbreak{} h,\allowbreak{} x,\allowbreak{} i)).R$.
Consequently its columns, extra record, level groups and stored
$d_\star$ are the ones used throughout RMP8--RMP10.
This verifies (RMP.32) on the selected construction class and its
actual target intersection, without an independent selector premise.

This discharges the v0.31 leaf \texttt{DESIGN\_\allowbreak RANKM\_\allowbreak CLOSED}.
The shared \texttt{DESIGN\_\allowbreak CHOICES\_\allowbreak COMMON} obligation still has
to construct a family satisfying this class together with the
other required common conditions. The argument above proves a
property of every member of the class; it does not assert that
the class has a member. The design-choice and level-size tasks
keep their separate scopes. There is no remaining class-number,
local-index or Weyl-dimension obligation hidden in this rank leaf.
[Sources R20--R21, R29--R33; inherited DEX5--DEX7 and TG12.]
% END RMP11


\section{The actual base theta map: convergence, descent, and the remaining extension}
\label{sec:base-theta-v032}

\status{BTH1--BTH15 give complete proofs of the stated implications on
the original base tuple.  They produce its convergent lowest-type
integral, reconstruct the smooth-vector uniqueness and finite
orthogonal-kernel arguments, and reduce the original base theta square
to an explicit common-object producer.  They do not yet construct the
original smooth cuspidal extension, its geometric reading, or its
simultaneous rational tests.  Those producer obligations remain C.
No new Lean replay is claimed.}

The base-rank consumer in SLR8 needs \(\mathrm{RowsHold0}_3\) and a
theta square at \(t_0=A_3.A.\mu_0\).  Its need for that square is
different from the separate \(\mathrm{RankM}\) branch.  Everything
below concerns the former consumer and one fixed original bundle.
There is no change to its coefficient field, Hilbert metric, adjoint,
finite group action, or chosen theta map.

Here is the notation used throughout.  Put
\[
 K=D.\mathrm{what}_2.\mathrm{LL},\qquad
 \iota=D.\mathrm{what}_2.\mathrm{embLL},\qquad
 \sigma=A_3.A.v_0.\mathrm{embedding}.
 \tag{BTH.1}
\]
Let \(R,\allowbreak{} F,\allowbreak{} Y,\allowbreak{} Z\) be the original real datum, face, \texttt{Read},
and \texttt{Read\allowbreak L}; their identification with the supplied bundle is
an explicit premise.  For \(i<m\), put
\[
 \begin{aligned}
 H_i&=(R.HV(i,\allowbreak{} t_0)).H,&
 W_i&=(R.HW(i,\mathrm{sg}_0)).H,\\
 L_i&=R.\mathrm{extV}(i,\allowbreak{} t_0),&
 U_i&=R.\mathrm{thF}(i,\allowbreak{} t_0),&
 T_i&=R.\mathrm{thG}(i,\allowbreak{} t_0).
 \end{aligned}                                                  \tag{BTH.2}
\]
The realization \(e_i:W_i\to\mathrm{MMap}(\mathcal D_i)\) is injective
and has image \(\mathrm{Wall}_i\);
\(\mathcal D_i=R.MW(i,\mathrm{sg}_0)\) is the prescribed tensor of
discrete-series models.  These are
fields of \texttt{Is\allowbreak Models}.  The space \(\mathrm{MMap}(\mathcal D_i)\)
consists of multilinear maps from the smooth-vector factors to
automorphic functions.  The subspace \(\mathrm{Wall}_i\) is the
span of maps with the source's automorphy, cuspidality, level,
equivariance, and Sobolev-continuity properties.

The full level coordinate space will be denoted \(K^{d_i}\).
Its map to rational forms is
\[
 f_i(x)=\mathrm{liftForm}(i,\allowbreak{} t_0,\allowbreak{} x).
 \tag{BTH.3}
\]
The original good-cut space is \(K^{s_i}\), and its inclusion and
normalized projection are
\(E_i=lZ0(i,\allowbreak{} t_0)\) and \(\Pi_i=lPrj(i,\allowbreak{} t_0)\).
The original target-coordinate map is \(J_i=Z.jW(i)\), and its
Hilbert reading is \(q_i=Y.mj.eW(i)\).
In particular \(q_iJ_i=\mathrm{goodForm}_i\).  SLR2 proves the
normalization
\[
 \Pi_iE_i=1,\qquad
 x-E_i\Pi_i x\perp_{lH_i}\operatorname{im}E_i.                    \tag{BTH.4}
\]
The tuple identification is literal:
\(Y.\mathrm{tup}(t_0)=D.\mathrm{what}_2.\mu_0\), which is
always \(\mathrm{Used0}\).  No assertion about a nonbase tuple
is needed.

% BEGIN BTH1
\paragraph{BTH1: smooth vectors of the lowest compact type are nonzero and cyclic in the required topology.}
\label{claim:BTH1}
Let \(\pi\) be one of the source's nonzero irreducible unitary
representations of a real unitary Lie group \(G_\infty\), with
maximal compact subgroup \(K_\infty\).  Write \(H^\infty\) for its
smooth vectors with seminorms
\[
 p_{\mathcal X}(v)=\|d\pi(X_1)\cdots d\pi(X_s)v\|,
 \quad X_j\in\mathfrak g,\quad s\ge0.                            \tag{BTH.5}
\]
Its nonzero lowest compact isotypic space contains a nonzero
smooth vector.  Moreover, if \(0\ne u\in H^\infty\), every
continuous linear functional for the topology (BTH.5) that
vanishes on all \(\pi(g)u\) is zero.

\emph{Proof.}
Choose a nonnegative smooth bump \(\phi\) supported sufficiently
near \(1\), and normalize its Haar integral to one.  Averaging
\(\phi(k g k^{-1})\) over \(K_\infty\) makes it conjugation
invariant under \(K_\infty\), with arbitrarily small support.
Conjugation preserves Haar measure here: the modular function
has compact image on \(K_\infty\), so that image is \(\{1\}\).
The smoothing operator
\[
       S_\phi v=\int_{G_\infty}\phi(g)\pi(g)v\,\allowbreak{} dg
 \tag{BTH.6}
\]
commutes with \(K_\infty\).  Differentiating the compactly
supported scalar kernel after a change of variable gives
\(S_\phi H\subset H^\infty\).  More precisely, for each
derivative word \(\mathcal X\) there is a compactly supported
smooth derivative \(D_{\mathcal X}\phi\), with the appropriate
signs, such that
\[
 p_{\mathcal X}(S_\phi v)
       \le \|D_{\mathcal X}\phi\|_{L^1}\|v\|.                    \tag{BTH.7}
\]
This follows first for one derivative by differentiating
\(\phi(\exp(-tX)g)\), and then by induction.  Compact support
justifies every differentiation under the integral.
For a nonzero vector \(v\) in the lowest compact isotypic space,
strong continuity permits the support to be chosen so that
\(\|\pi(g)v-v\|<\|v\|/2\) there.  Hence
\(\|S_\phi v-v\|<\|v\|/2\); the smoothed vector is nonzero
and remains in the same isotypic space.

Now suppose \(\ell\) has the stated continuity and vanishing
properties.  Its continuity gives finitely many words and a
constant \(C\) with
\[
       |\ell(v)|\le C\sum_{\mathcal X}p_{\mathcal X}(v).
 \tag{BTH.8}
\]
By (BTH.7), \(\ell S_\phi\) extends to a bounded functional on
the Hilbert space \(H\).  For every \(g\),
\[
  \ell S_\phi(\pi(g)u)
      =\int\phi(h)\ell(\pi(hg)u)\,\allowbreak{} dh=0.                          \tag{BTH.9}
\]
The interchange is valid in the smooth topology: for a fixed
smooth \(u\), all derivative seminorms of the orbit map are
continuous and bounded on this compact support.
The closed Hilbert span of \(\{\pi(g)u:g\in G_\infty\}\) is a
nonzero invariant subspace, hence all of \(H\) by irreducibility.
Thus \(\ell S_\phi=0\).

Take bumps forming an approximate identity.  For every
\(v\in H^\infty\), \(S_\phi v\to v\) in every seminorm (BTH.5).
To see this directly, move a derivative across \(\pi(g)\) using
\[d\pi(X)\pi(g)=\pi(g)d\pi(\operatorname{Ad}(g^{-1})X);\]
iteration gives a finite sum whose coefficients and vectors
converge as \(g\to1\).  Integrating over shrinking supports gives
the convergence in each fixed seminorm.  Consequently
\(\ell(v)=\lim\ell(S_\phi v)=0\).

The source's \texttt{DSModel.\allowbreak ds} supplies the irreducibility
and unitarity used here, and its lowest compact type is nonzero
before smoothing.  Its frames conjugate the standard unitary
group to the original Gram-matrix group, preserving the argument.
This proves the smooth-topology statement; Hilbert density alone
would not have justified applying a Sobolev-continuous functional.
The corresponding source declarations are
\texttt{exists\_\allowbreak smooth\allowbreak Vec\_\allowbreak low\allowbreak KSub}, \texttt{low\allowbreak T\_\allowbreak ne\_\allowbreak bot},
\texttt{eq\_\allowbreak zero\_\allowbreak of\_\allowbreak translates}, and \texttt{low\allowbreak Generates}.
% END BTH1

% BEGIN BTH2
\paragraph{BTH2: a multiplicity map is determined by its lowest-type values.}
\label{claim:BTH2}
For the original \(\mathcal D_i\), suppose \(A,\allowbreak{} B\in\mathrm{Wall}_i\)
have equal values on all pure tensors of lowest-type vectors.
Then \(A=B\).

\emph{Proof.}
Put \(C=A-B\).  A scalar evaluation of \(C\), with the other
archimedean variables fixed, is Sobolev-continuous in the
remaining variable.  For a generator of \(\mathrm{Wall}_i\)
this is part of its defining continuity condition.  It remains
true under finite sums and scalar multiples, which is exactly
how the span defining \(\mathrm{Wall}_i\) is formed.

Order the finite set of archimedean factors.  Initially \(C\)
vanishes when every factor is of lowest type.  Suppose it
vanishes with the first \(s\) factors arbitrary and the rest
of lowest type.  Fix all but the next factor, and fix the
adelic evaluation point.  Equivariance in that next factor
shows that the resulting functional vanishes on every
translate of a nonzero smooth lowest-type vector: the translate
only changes the adelic evaluation point in a function that
already vanishes identically.  BTH1 makes this functional
zero on the entire smooth factor.  Induction removes all the
lowest-type restrictions.  Multilinearity and the spanning
property of pure tensors give \(C=0\).

The same proof applies to the matched finite target lifts used
below.  Those lifts normalize the prescribed principal flag
level and preserve its character, by the original level-lift
conditions.  Their right translations therefore preserve this
\(\mathrm{Wall}_i\), including its automorphic and Sobolev
conditions.  Thus an identity proved by the theta
integral on lowest types extends to all smooth-vector
arguments.  This is the exact content needed from
\texttt{eq\_\allowbreak zero\_\allowbreak of\_\allowbreak low\allowbreak T} and
\texttt{theta\allowbreak Lift\_\allowbreak intertwines\_\allowbreak lifts\_\allowbreak smooth\allowbreak Vectors}.
It is a uniqueness theorem, not an existence theorem for the
theta extension.
% END BTH2

% BEGIN BTH3
\paragraph{BTH3: the original base lowest-type theta integral converges.}
\label{claim:BTH3}
Assume the original model fields and weak theta fields
\(\omega P,\allowbreak{} vP,\allowbreak{} pt\_\allowbreak theta\).  For every \(y\in H_i\), every
lowest-type tensor \(a\), and every target adelic point \(g\),
the original integral
\[
 {\cal I}_i(y,\allowbreak{} a)(g)=
 \int_{\mathcal F_{V,\allowbreak{} i}}
  \theta KO(\omega_i)(h,\allowbreak{} g)
       \bigl(pt_i(io_i(a,(L_i y)(h)))\bigr)\,\allowbreak{} d\nu_i(h)            \tag{BTH.10}
\]
is absolutely convergent.  It is complex linear in \(y\)
and in \(a\).  Neither an \texttt{Is\allowbreak Io\allowbreak W} law nor an existing
cuspidal extension is needed for this conclusion.

\emph{Proof.}
The source space \(V_c(D,\allowbreak{} i)\) is positive definite at every
archimedean place: \texttt{models.\allowbreak vpos}, the original
\texttt{V\_\allowbreak sig}, and a change of frame give positive definite
Gram matrices.  A real unitary group of such a matrix is compact.
Indeed write its Gram matrix as \(B^*B\).  Conjugation by \(B\)
identifies the group with \(U(n)\), whose entries have absolute
value at most one and whose defining equations are closed.

For the finite part use the double-coset representatives
stored in \((R.B.L(i,\allowbreak{} t_0)).pres\); no new class-number
finiteness theorem is required.  The finite level is a closed
subgroup of a conjugate of \(GL_n(\widehat{\mathcal O}_E)\)
cut out by finitely many congruences.  It is compact: each
integral local matrix group is compact, the product is compact,
and conjugation and the closed unitary equations preserve this
property.  Its inverse image in the full adelic group is compact
because its real factor is compact as well.  Lifting the finitely
many stored representatives and taking their translates gives
a compact set \(C\) with
\[
       G_V(\mathbb A)=G_V(F)\,\allowbreak{} C.                                \tag{BTH.11}
\]
The finite projection used here is the source's range
projection, so each stored representative has an adelic lift.

The prescribed measure \(\nu_i\) is Haar and its prescribed
\(\mathcal F_{V,\allowbreak{} i}\) is a measurable fundamental domain.  It
has finite measure.  For completeness, the countable rational
group and (BTH.11) give
\(\sum_{\gamma\in G_V(F)}1_C(\gamma h)\ge1\).
Integrate over that fundamental domain and unfold its
translates.  Countable additivity for the nonnegative sum and
left invariance give
\[
 \nu_i(\mathcal F_{V,\allowbreak{} i})
 \le\int_{\mathcal F_{V,\allowbreak{} i}}\sum_\gamma1_C(\gamma h)\,\allowbreak{} d\nu_i(h)
 =\nu_i(C)<\infty.                                               \tag{BTH.12}
\]
The fundamental-domain equality is understood almost everywhere,
as in its definition.

The function \(F_y=L_i y\) is continuous and rational-left
invariant.  The first assertion follows from the continuous
coefficient representation and a sufficiently small principal
open subgroup on which the finite-level character is one.
The second follows by cancellation in
\[
     (c_V.\mathrm{ext}\,\allowbreak{} f)(h)
       =\rho_V(h^{-1})f(\mathrm{fin}(h)):
\]
the rational automorphy of \(f=(R.HV(i,\allowbreak{} t_0)).e(y)\) cancels
\(\rho_V(\gamma)\).

Choose a basis \(b_\alpha\) of the finite-dimensional
coefficient space of \(c_V\).  The integrand in (BTH.10) is
the finite sum
\[
 \sum_\alpha b_\alpha^*(F_y(h))\,
    \omega_i(h,\allowbreak{} g)\bigl(pt_i(io_i(a,\allowbreak{} b_\alpha))\bigr)(0).            \tag{BTH.13}
\]
Its first factors are continuous, and its second factors
are continuous by the \texttt{om\_\allowbreak cont} field of the actual
\texttt{Is\allowbreak Weil\allowbreak Datum\allowbreak O}.  The \texttt{pt\_\allowbreak theta} field checks
the exact theta-function hypothesis for each vector appearing
in (BTH.13).  Rational-left invariance follows from
\texttt{om\_\allowbreak aut}, \texttt{om\_\allowbreak pres}, and the invariance of
\(F_y\).  It is therefore bounded on all of the quotient by
its bound on \(C\).  Equations (BTH.11)--(BTH.12) prove absolute
integrability.  Bilinearity of \(io_i\), linearity of the
Weil action and \(pt_i\), and additivity of integrals of
integrable functions prove the two linearity assertions.

Thus the \texttt{Is\allowbreak Theta\allowbreak P} integrability guard is discharged
on the actual base data.  The substantially stronger assertion
that (BTH.10) extends to a smooth cuspidal multiplicity map is
not a consequence of compactness of the source quotient.
% END BTH3

% BEGIN BTH4
\paragraph{BTH4: finite extension data activate the original weak choices and construct the good-projected map.}
\label{claim:BTH4}
Fix a finite basis \(h_1,\ldots,\allowbreak{} h_n\) of the original \(H_i\).
Assume that for each \(h_\alpha\) there is
\(A_\alpha\in\mathrm{Wall}_i\) whose lowest-type restriction
is the integral (BTH.10).  Then the source's actual stored
\(U_i=\mathrm{thF}(i,\allowbreak{} t_0)\) satisfies \(\mathrm{PThF}(i,\allowbreak{} t_0)\).
The actual stored \(T_i=\mathrm{thG}(i,\allowbreak{} t_0)\) then satisfies
\[
       e_i T_i=\mathrm{postOn}(eWf_i)\,\allowbreak{} U_i.                      \tag{BTH.14}
\]
No independent \(\mathrm{PThG}\)-existence premise is needed.

\emph{Proof.}
Define \(U(\sum c_\alpha h_\alpha)=\sum c_\alpha A_\alpha\).
It is linear and lies in \(\mathrm{Wall}_i\).  BTH3 and the
basis expansion identify its lowest-type restriction with
\({\cal I}_i\) for every \(y\).  On \(\operatorname{im}L_i\),
use \(L_i^{-1}\), available from \(Y.\mathrm{hinj}\), to define
the low-type theta map.  Formula (BTH.10) and its convergence
give exactly the two fields of \texttt{Is\allowbreak Theta\allowbreak P}.
This proves \(\mathrm{PThF}(i,\allowbreak{} t_0,\allowbreak{} U)\).
BTH2 also proves uniqueness of such \(U\).

The weak-choice rule is
\[
 \mathrm{OrZero}(P,\allowbreak{} z)
 =((\exists w,\ P(w))\Rightarrow P(z))
   \ \wedge\ ((\nexists w,\ P(w))\Rightarrow z=0).                \tag{BTH.15}
\]
The first implication in \texttt{tw\_\allowbreak th\allowbreak F} therefore gives
\(\mathrm{PThF}(i,\allowbreak{} t_0,\allowbreak{} U_i)\) for the stored \(U_i\), not
for a substituted map.

Write the prescribed central projector as
\(eWf_i=\sum_{g\in G}c_g\,\mathrm{actL}(g)\).
The model identity for \(gW\) gives
\[
 e_i\!\left(\sum_g c_g\,\allowbreak{} gW_i(g)w\right)
       =\mathrm{postOn}(eWf_i)(e_iw).                           \tag{BTH.16}
\]
Since \(\operatorname{im}e_i=\mathrm{Wall}_i\), the map on
the right of (BTH.14) takes its values in
\(\operatorname{im}e_i\).  Pulling it back through the linear
isomorphism \(e_i:W_i\to\operatorname{im}e_i\) constructs a
linear map \(T\) satisfying (BTH.14).  Injectivity of \(e_i\)
proves its uniqueness.  The first implication in
\texttt{tw\_\allowbreak th\allowbreak G} now proves (BTH.14) for the stored \(T_i\).

This is the pointwise-at-\(t_0\) form of
\texttt{exists\_\allowbreak linear\_\allowbreak lift} and \texttt{exists\_\allowbreak factor\_\allowbreak hmod}.
The global eight-field \texttt{Ex\allowbreak At} package, including its
chain lifts and other tuples, has not been used.
% END BTH4

% BEGIN BTH5
\paragraph{BTH5: the actual finite character gives theta covariance and the inverse target level.}
\label{claim:BTH5}
Assume the original weak Weil, polynomial, group, and level
fields and the base extension from BTH4.  For matching source
and target finite-level lifts \(\kappa,\kappa'\), put
\(c=\mathrm{levChi}(\mathrm{flagL}(\mathrm{datV}))(\kappa)\).
Then
\[
 \omega_i(\kappa,\kappa')vP_i=c^{-1}vP_i.                        \tag{BTH.17}
\]
If
\[
        c\,\allowbreak{} L_i y'(h)=L_i y(h\kappa),
\]
the full multiplicity maps satisfy
\[
        U_i y'(v)(g)=U_i y(v)(g\kappa')                         \tag{BTH.18}
\]
for every smooth-vector tuple \(v\).  In particular a source
level character is inverted on the target level, with the
source's exact fixed embedding.

\emph{Proof.}
We record the finite normalization because a scalar ambiguity
would change the asserted map.  At a prime in \(P_g\), write
\(a,\allowbreak{} d\) for the determinant characters of the source's upper
and lower blocks and \(a'\) for the target character.  The
matched upper blocks give \(a'=a\).  At a prime in \(Q\), write
\(x_a,\allowbreak{} x_d,\allowbreak{} x_b\) for the corresponding prescribed \(\xi\)-values.
The source's vacuum and half-datum transformations are
\[
 \begin{aligned}
 m_V&=\prod_{P_g}(ad)^{-1}\prod_Q x_ax_d,&
 m_W&=\prod_{P_g}(a')^{-1}\prod_Q x_b,\\
 d_{\rm half}&=\prod_Q x_a^{-1}x_b,&
 c&=\prod_{P_g}d^{-1}\prod_Q x_d .
 \end{aligned}                                                  \tag{BTH.19}
\]
Thus \(m_V^{-1}m_W=d_{\rm half}c^{-1}\).
The half-support is preserved: its reduced adjoint frame
changes by \(A\mapsto GAH^{-1}\); the zero lower block
remains zero, the identity upper block at \(P_g\) remains
the identity, and the invertible upper block at \(Q\)
has inverse-determinant-character multiplier
\(x_a^{-1}x_b\).  Reindexing the finite Heisenberg sum therefore
gives
\(d_{\rm half}\omega_i(\kappa,\kappa')vP_i
   =(m_V^{-1}m_W)vP_i\).
All factors are units, so cancellation proves (BTH.17).
The vacuum multipliers are precisely the actual-action
stabilizer character proved in BTH15 below, with the recorded
determinants.  BTH13--BTH14 justify the identification with a
canonical action wherever that comparison is used.  Their
frame hypotheses are the self-dual
\(\beta V,\beta W\) fields in the original level data.
Their unramified-away hypothesis is the original
\(\mu\)-character condition away from \(P_g\cup Q\);
the second Weil character is \(\mu\eta^{-1}\).
The required extension character is the original \(\eta\):
the two restriction laws and
\(\operatorname{rank}V=r+t,\operatorname{rank}W=r\),
with \(t\) odd, give its required odd-parity restriction.
Explicitly, for every quadratic \(\eta_F\) appearing in
\texttt{Res\allowbreak Eta}, the two restrictions give
\(\mu|_F=\eta_F^{r+t}\) and
\((\mu\eta^{-1})|_F=\eta_F^r\).  Since
\(\eta_F^2=1\), division gives
\(\eta|_F=\eta_F^{2r+t}=\eta_F\).
This holds for every such \(\eta_F\), retaining the universal
restriction in \texttt{Is\allowbreak Eta\allowbreak Tilde}.
The level record's \texttt{Chr\allowbreak Eq}/\texttt{Is\allowbreak Eta\allowbreak T} gives the
prescribed infinity exponent and triviality on each finite
local unit group.  Products with finite support are dense in
the compact product of these groups, so continuity gives
triviality on every finite integral unit idele.
These are the checks carried out in
\texttt{om\_\allowbreak lifts} before its finite determinant cancellation.
No arbitrary Weil scalar is introduced.

The finite actions commute with the polynomial creation
operators at infinity.  This differentiation does not
presuppose continuity of an ambient endomorphism on all
functions: the actual Haar-norm identity makes it an isometry
on theta \(L^2\); on a fixed smooth finite level the real
Heisenberg difference quotients converge in \(L^2\), and
covariance carries those quotients through the isometry.
The output is again a smooth theta function, so the resulting
\(L^2\) derivative is its classical derivative.  Iteration is
exactly the difference-quotient argument in WA3.
Applying those operators to
(BTH.17), one at a time, gives the same scalar \(c^{-1}\)
on every \(pt_i(P)\).  Equivalently this is the pure-tensor
finite-action calculation in \texttt{om\_\allowbreak pure\allowbreak Tensor\allowbreak O\_\allowbreak of\_\allowbreak finite}.
The homomorphism law for \(\omega_i\) then gives the kernel
identity
\[
 \theta(g\kappa',\allowbreak{} h\kappa,\allowbreak{} pt_i(P))
       =c^{-1}\theta(g,\allowbreak{} h,\allowbreak{} pt_i(P)).                              \tag{BTH.20}
\]
Every level lift lies in a compact subgroup of the source
adelic group.  Its modular character is therefore one, so
right translation by \(\kappa\) preserves Haar measure.
The integral of an integrable rational-left-invariant
function over a fundamental domain is unchanged by that
translation: translate the domain, use the fundamental
domain partition, and cancel the left rational translates.
BTH3 checks the integrability for the present functions.
Substitution in (BTH.10), followed by (BTH.20), proves
(BTH.18) on lowest-type tensors.  BTH2 proves it on all
smooth tensors.

To check the inverse character explicitly, let \(p\) be a
target level element.  Its matched source lift follows
from the already supplied \texttt{Bridge.\allowbreak Act\allowbreak Tie.\allowbreak red\_\allowbreak surj}:
choose \(q\) above \(eG^{-1}(g_0)\), lift its finite
component with identity infinite component, and use
\texttt{grp\_\allowbreak red\_\allowbreak L} to identify its reduction with the
literal \(rTup\) of \(g_0\).  Thus no separate
parahoric-surjectivity producer is needed here.
The original level relation says that the source character
on this lift is \(c\,\psi(p)^{-1}\).  Hence one can take
\(y'=\psi(p)^{-1}y\) in (BTH.18), which gives
\(U_i y(v)(gp)=\psi(p)^{-1}U_i y(v)(g)\).
The source declaration
\texttt{lift\allowbreak Level\_\allowbreak fixed\allowbreak Embedding\_\allowbreak eq\_\allowbreak conj\_\allowbreak row\allowbreak Level}
identifies this with \(R.\mathrm{levL}(F,\mathrm{false},
\mathrm{sg}_0)\).  This verifies the full finite-level
membership needed below, not just principal-level membership.
% END BTH5

% BEGIN BTH6
\paragraph{BTH6: the stored target finite action is unitary and the base theta map intertwines it.}
\label{claim:BTH6}
Under the same assumptions, the target action \(gW_i\)
preserves the original Hilbert inner product, and
\[
 T_i jV_i(r_i(g)f)=gW_i(g)\,\allowbreak{} T_i jV_i(f)
       \qquad(g\in G,\ f\text{ an original rational form}).      \tag{BTH.21}
\]
The central projector acts as the identity on
\(\operatorname{im}T_i\).  Consequently the lift portions of
the source's \texttt{PGUnit}, \texttt{PGEquiv}, and
\texttt{e\allowbreak H\_\allowbreak laws} introduce no additional base producer.

\emph{Proof.}
The Hilbert realization identifies the inner product with
the prescribed multiplicity inner product \(ipAt_i\).
For a finite group element with a target adelic lift, its
action is a right translate multiplied by a level character
of absolute value one.  Right translation preserves Haar
measure because the lift belongs to a compact subgroup.
The defining integral for each multiplicity pairing is
unchanged by the translation and the two unit-modulus
factors cancel.  This also preserves the assertion that a
scalar is the multiplicity pairing.  Its uniqueness is
explicit: choose a nonzero smooth vector in each discrete-series
factor using BTH1, take the two tensor arguments equal, and
divide by the nonzero product of their squared norms.
Translation by the inverse gives the converse existence
implication.  Thus the translation preserves the source's
chosen \(ipAt_i\), including its conditional-choice branch.
Norms, and then inner
products by complex polarization, are unchanged.

If an element has no target lift, the source definition
sets its function action to zero.  The model identity
\(e_i gW_i(g)=\mathrm{postOn}(\mathrm{actL}(g))e_i\)
and injectivity of \(e_i\) would then make \(gW_i(g)\)
zero.  It is an invertible group action, so \(W_i=0\);
all the asserted identities in this case follow.
This deals with the no-lift branch rather than assuming
surjectivity of target lifts.

For a rational form \(f\), the original
\texttt{act\_\allowbreak tie.\allowbreak red\_\allowbreak surj} and \texttt{grp\_\allowbreak red\_\allowbreak L},
by the construction just given in BTH5, supply a source lift \(\kappa\)
with
\[
 (HV.e)(jV_i(r_i(g)f))(h)
 =c^{-1}(HV.e)(jV_i f)(h\,\mathrm{fin}(\kappa)).
 \tag{BTH.22}
\]
The coefficient extension converts (BTH.22) into the
hypothesis of BTH5.  Choose a target lift if one exists.
Equation (BTH.18) and the trivial principal-level target
character identify the resulting function with
\(\mathrm{actL}(g)\,\allowbreak{} U_i(jV_i f)\).

The good projector commutes with these finite actions.
Indeed its group-algebra element is the original central
character idempotent
\[
C=\sum_{\chi\in X}\frac{\chi(1)}{|G|}
                 \sum_{g\in G}\chi(g^{-1})g ;
 \tag{BTH.23}
\]
conjugation invariance of every \(\chi\in X\) proves centrality by
reindexing, and the irreducible character orthogonality
used in the source's \texttt{Ok} makes these distinct character
idempotents mutually orthogonal and gives \(C^2=C\).
The identification of its coefficients with \(eWf_i\)
is \texttt{grp.\allowbreak record}; the embeddings, including their
complex conjugation, are the ones in (BTH.1).
Apply the good projector to (BTH.18), use (BTH.14), and
then use injectivity of \(e_i\).  This proves (BTH.21).

Finally (BTH.16) writes \(T_i y=E_i^{\rm cen}w\) for
some \(w\), where \(E_i^{\rm cen}=\sum c_g gW_i(g)\).
Since \((E_i^{\rm cen})^2=E_i^{\rm cen}\), this projector
fixes \(T_i y\).  The same statements on the twisted
spaces follow by applying the exact \(\mathrm{Tw}\) functor.
The rational source action is unitary for its original
coordinate Hermitian form by the supplied \texttt{Ok}
block data.  Thus both unitary actions and their actual
intertwiner required in BTH9 are available.
% END BTH6

% BEGIN BTH7
\paragraph{BTH7: the exact original Hilbert reading, with an explicit witness.}
\label{claim:BTH7}
Assume that \(Y,\allowbreak{} Z\) are the actual readings on the original
bundle, and the base extension hypotheses of BTH4 hold.
For every \(x\in K^{d_i}\),
\[
 q_i\!\left(
 D.\theta P\bigl(\mu_0,\allowbreak{} i+1,\allowbreak{} D.jMt(\mu_0,\allowbreak{} i+1)(Y.eMt(i,\allowbreak{} t_0)x)\bigr)
       \right)=T_i jV_i(f_i(x)).                               \tag{BTH.24}
\]

\emph{Proof.}
Let \(x_M=D.jMt(\mu_0,\allowbreak{} i+1)(Y.eMt(i,\allowbreak{} t_0)x)\) and
\(z=jV_i(f_i(x))\).  The \texttt{mt\allowbreak Read} field says exactly
\[
       (Y.IL.jMt(\mu_0,\allowbreak{} i,\allowbreak{} x_M)).1=L_i z.                         \tag{BTH.25}
\]
By the range and level condition in the codomain of
\(IL.jMt\), this function is at the prescribed source
base level.  BTH5 therefore puts \(U_i z\) in the
corresponding full target multiplicity space.
The central projector preserves that target level:
its action commutes with the level action and preserves
the prescribed character.  Its idempotence puts the
projected map in the original good target space \(W_c\).

Use the explicit tensor-linear map
\[
                  w=\operatorname{lift}(U_i z).                \tag{BTH.26}
\]
The definition of \texttt{to\allowbreak Lift\allowbreak Real.\allowbreak th\allowbreak C} first inverts
\(L_i\) on its image and then restricts this tensor map to
lowest types.  Equation (BTH.25), \(L_i\)'s injectivity,
and \(Y.\mathrm{tupT\_\allowbreak tup}\) identify that inverse with
this same \(z\).  Consequently \(w\) has precisely the
required lowest-type restriction.  Restricting the
convergent integral from \(\operatorname{im}L_i\) to its
base-level subspace proves \texttt{Is\allowbreak Theta\allowbreak P} for that
restriction: equality is checked on pure tensors and
integrability is inherited from BTH3.

The universally quantified \texttt{Y.\allowbreak tied.\allowbreak lift.\allowbreak theta}
may now be applied to this particular \(w\).  It gives
\[
      IL.jMj(D.\theta P(\mu_0,\allowbreak{} i+1,\allowbreak{} x_M))
             =eWf_i\circ w.                                    \tag{BTH.27}
\]
The \texttt{mj\allowbreak Read} field changes the left side into
\(\operatorname{lift}(e_i q_i(D.\theta P(\cdots)))\).
On the right we have
\[
 \operatorname{lift}(\mathrm{postOn}(eWf_i)U_i z)
     =eWf_i\circ\operatorname{lift}(U_i z).
\]
By (BTH.14) this is
\(\operatorname{lift}(e_i T_i z)\).
Injectivity of the tensor-linear equivalence and of \(e_i\)
proves (BTH.24).

This direct choice avoids obtaining an unspecified witness
from \texttt{Theta\allowbreak Fire\allowbreak L} and then identifying it afterwards.
BTH2 is still used in BTH5 to establish its full-level
membership.  The geometric \texttt{Lift\allowbreak Tie.\allowbreak theta} premise
is still part of the actual reading; this argument does
not construct that premise from an arbitrary bundle.
% END BTH7

% BEGIN BTH8
\paragraph{BTH8: one base rational test family supplies the rational adjoint image.}
\label{claim:BTH8}
Use the source's twisted spaces
\[
 V_i=\mathrm{Tw}_{\epsilon}(H_i),\qquad
 \mathcal W_i=\mathrm{Tw}_{\epsilon}(W_i),\qquad
 \Theta_i=(R.\mathrm{toCx3}(F)).\theta'(i,\allowbreak{} t_0),                  \tag{BTH.28}
\]
where \(\epsilon=\mathrm{twAt}(\iota,\allowbreak{} v_0)\).  The phase is
one by \texttt{Pins.\allowbreak phs}, so that
\(\mathrm{Tw.out}\,\Theta_i(\mathrm{Tw.mk}\,\allowbreak{} y)=T_i y\).
The original rational-form realization
\(j_i:F_i\to_\sigma V_i\) spans \(V_i\) and is isometric
for its original \(K\)-Hermitian form.

Suppose there are finitely many complex-linear functionals
\(\lambda_1,\ldots,\lambda_N\) on \(\mathcal W_i\) such that
\[
 \begin{split}
  &\bigl(\lambda_a(w)=0\text{ for all }a\bigr)\Longrightarrow w=0,\\
  &\lambda_a(\Theta_i j_i f)\in\sigma(K)
                   \quad(f\in F_i),\\
  &\lambda_a(w_{0,\allowbreak{} i}e_k)\in\sigma(K)
                   \quad(k\text{ an original target basis index}).
 \end{split}                                                    \tag{BTH.29}
\]
Here \(w_{0,\allowbreak{} i}=(R.\mathrm{toCx3}(F)).w0C(\mathrm{false},\allowbreak{} i)\).
Then, writing
\(Z_i=\operatorname{im}\Theta_i^*\) and
\(Z_{K,\allowbreak{} i}=j_i^{-1}(Z_i)\),
\[
                   \operatorname{span}_{\mathbb C}j_i(Z_{K,\allowbreak{} i})
                         =Z_i.                                 \tag{BTH.30}
\]

\emph{Proof.}
Choose a \(K\)-basis of \(F_i\).  Its Gram matrix for the
original nondegenerate \(K\)-Hermitian form is invertible
and remains invertible under \(\sigma\).  By the isometric
realization, this is the Gram matrix of its image vectors,
so they are complex linearly independent.  The full-span
property makes them a complex basis of \(V_i\).
Applying all the tests to its
theta images gives a matrix with entries in \(\sigma(K)\).
Because the tests separate points, its kernel is exactly
\(\ker\Theta_i\) in these coordinates.  Gaussian elimination
over \(K\) gives a \(K\)-basis of that kernel, and its
complex span is the whole complex kernel: the nonzero
minors determining rank remain nonzero under \(\sigma\).

The matrix of the source Hermitian form also has entries
in \(K\), and \(\sigma(\bar a)=\overline{\sigma(a)}\).
Taking the orthogonal complement of the rational kernel
is therefore another system of linear equations over
\(K\), with the same rank after complexification.
Its complex span is
\((\ker\Theta_i)^\perp=\operatorname{im}\Theta_i^*\).
This is (BTH.30), with the original adjoint and metric.

Only the first two lines of (BTH.29) were used here;
the third will give rational target coordinates in
BTH11.  Thus one base-only \texttt{Lift\allowbreak Test\allowbreak B} family
does both jobs.  No tests for another tuple and no
independent \texttt{Lift\allowbreak Test} family are needed.
Existence of the simultaneous separating tests with
these rationality properties remains a substantive
producer; finite-dimensionality alone does not supply
them.
% END BTH8

% BEGIN BTH9
\paragraph{BTH9: the finite unitary idempotent kernel argument.}
\label{claim:BTH9}
Let \(F_K\) be a rational model of a finite-dimensional
complex Hilbert space \(V\), with full realization \(j\).
Let finite-group representations \(\rho,\tau\) on \(V,\allowbreak{} W\)
be unitary, let \(\Theta:V\to W\) intertwine them, and
assume
\(\operatorname{span}_{\mathbb C}j(Z_K)
  =Z=\operatorname{im}\Theta^*\).
Let \(e,\allowbreak{} a\in K[G]\), assume \(\rho(e)^2=\rho(e)\), and
assume \(\tau(a)\) fixes \(\operatorname{im}\Theta\).
For \(f\in F_K\), suppose
\[
 ef=af,\qquad
       \langle jz,\allowbreak{} jf\rangle=0
           \quad(z\in Z_K\cap\operatorname{im}e).
 \tag{BTH.31}
\]
Then \(\Theta(jf)=0\).

\emph{Proof.}
The kernel of \(\Theta\) is group invariant.  Since
the representation is unitary, its orthogonal complement
\(Z\) is also invariant.  Thus the orthogonal projection
\(P_Z\) commutes with every \(\rho(g)\) and consequently
with \(\rho(e)\).  It follows that
\[
 u=P_Z\rho(e)jf\in Z\cap\operatorname{im}\rho(e),
 \qquad
 \Theta u=\Theta\rho(e)jf=\tau(a)\Theta jf=\Theta jf.
 \tag{BTH.32}
\]
The rational equality
\[
 Z\cap\operatorname{im}\rho(e)
       =\operatorname{span}_{\mathbb C}
                     j(Z_K\cap\operatorname{im}e)               \tag{BTH.33}
\]
does not require \(e\) to be self-adjoint.  Indeed the
left side equals \(\rho(e)Z\): idempotence makes \(e\)
the identity on its image.  Apply \(\rho(e)\) to the
rational spanning expression for \(Z\); its rational
coefficients and preservation of \(Z\) give the right
side.  The reverse inclusion is immediate.
Equations (BTH.31)--(BTH.33) yield
\(\langle u,\allowbreak{} jf\rangle=0\).

Write \(u=\Theta^*w\).  Equation (BTH.32) gives
\(\Theta(u-jf)=0\), so
\[
 0=\langle\Theta^*w,\allowbreak{} u-jf\rangle
   =\langle u,\allowbreak{} u\rangle-\langle u,\allowbreak{} jf\rangle
   =\|u\|^2.                                                    \tag{BTH.34}
\]
Hence \(u=0\), and (BTH.32) finishes the proof.
This is a complete proof of the finite assertion in
\texttt{theta\_\allowbreak eq\_\allowbreak zero\_\allowbreak of\_\allowbreak orth}.  In particular the
source's group-algebra level idempotent has not been
silently replaced by an orthogonal projector.
% END BTH9

% BEGIN BTH10
\paragraph{BTH10: the original orthogonal complement is killed, with its exact level normalization.}
\label{claim:BTH10}
Assume (BTH.29), the original \texttt{Ok} block data, and
the hypotheses of BTH4--BTH7.  Then
\[
  \bigl(hform(lH_i)(E_i z,\allowbreak{} x)=0\text{ for every }z\bigr)
       \quad\Longrightarrow\quad T_i jV_i(f_i(x))=0.
 \tag{BTH.35}
\]

\emph{Proof.}
The source's \(\mathrm{lvlA}(P_0,\psi_0)\) uses
the character \(\psi_0^{-1}\).  Write \(P=P_0\),
\(\psi=\psi_0^{-1}\), and
\[
 p_\psi=\frac1{|P|}\sum_{p\in P}\psi(p^{-1})p,
            \qquad e=Cp_\psi,\qquad a=C.                       \tag{BTH.36}
\]
The central idempotent \(C\) in (BTH.23) commutes with
\(p_\psi\).  The latter is idempotent by reindexing the
double sum, so \(e^2=e\).  If \(f\) is a level-\(\psi\)
form, each summand in \(p_\psi f\) is \(|P|^{-1}f\);
thus \(p_\psi f=f\) and \(ef=af\).

The definition of \(\mathrm{zsOf}\) is the original
rational preimage of the adjoint image.  Consequently
the original good-cut form inclusion has range
\[
       f_i(E_iK^{s_i})
            =Z_{K,\allowbreak{} i}\cap\operatorname{im}(Cp_\psi).             \tag{BTH.37}
\]
This is a definition-level equality (\texttt{good\allowbreak R\_\allowbreak eq}
and \texttt{range\_\allowbreak to\allowbreak F}), not an equality deduced from
dimensions.  Notice the exact metric factor:
\[
 \mathrm{csum}(f_i(v),\allowbreak{} f_i(x))
              =|P|\,\allowbreak{} hform(lH_i)(v,\allowbreak{} x).                         \tag{BTH.38}
\]
The factor is \(|P|\), because \(lH_i\) was defined with
\(|P|^{-1}\); it is not its reciprocal in (BTH.38).
The model's \texttt{JIso\allowbreak L} identifies the Hilbert pairing
with the \(\sigma\)-image of this rational \(\mathrm{csum}\).
Thus the hypothesis of (BTH.35), together with (BTH.37),
is exactly the orthogonality premise (BTH.31).

BTH6 supplies unitarity, intertwining, and the fact that
\(C\) fixes the theta image.  BTH8 supplies (BTH.30).
BTH9 therefore proves \(\Theta_i(j_i f_i(x))=0\).
Applying \(\mathrm{Tw.out}\) and the actual phase-one
identity gives (BTH.35).
No positive theta-norm formula, full common-image
identity \(\mathrm H5\), or nonbase primitive row is
used in this kernel argument.
% END BTH10

% BEGIN BTH11
\paragraph{BTH11: the stored target basis gives containment and the original rational theta matrix.}
\label{claim:BTH11}
Assume in addition that the supplied bundle's
\texttt{Th\allowbreak Def3} basis is the same \(J_i\) in the reading.
Then the actual good-cut theta image lies in the image of
\(w_{0,\allowbreak{} i}\), and its original relation matrix is rational
over \(K\).  If
\[
             M_i=(A_3.\mathrm{indefDown}(s0Of(R,\allowbreak{} F))).\theta(i,\allowbreak{} t_0),
 \tag{BTH.39}
\]
then, for every \(z\in K^{s_i}\),
\[
                T_i jV_i(f_i(E_i z))=q_iJ_i(M_i z).
 \tag{BTH.40}
\]
There is no additional \texttt{PLev\allowbreak W}, \texttt{LUnif},
or full \(\mathrm H5\) hypothesis.

\emph{Proof.}
The \texttt{Th\allowbreak Def3.\allowbreak li/\allowbreak span} fields say that the complex
vectors \(J_i(e_k)\) are a basis of the original
\(D.Mj(i+1)\).  The map \(q_i\) is injective:
its \texttt{Mj\allowbreak Tie.\allowbreak e\allowbreak W\_\allowbreak norm} formula equates its norm
to a strictly positive fixed multiple of the original
norm, since \(\mathrm{nP}(\mathrm{false})>0\).
The exact source reading is
\[
       w_{0,\allowbreak{} i}(\sigma(z))
                =\mathrm{Tw.mk}_\epsilon(q_iJ_i z).             \tag{BTH.41}
\]
In particular \(w_{0,\allowbreak{} i}\) is injective.

Apply BTH7 to \(E_i e_\alpha\).
Its left side belongs to \(q_i(D.Mj(i+1))\), which is
spanned over \(\mathbb C\) by the \(q_iJ_i(e_k)\).
Equation (BTH.41) puts
\(\Theta_i j_i f_i(E_i e_\alpha)\) in
\(\operatorname{im}w_{0,\allowbreak{} i}\).  Linearity gives this
containment for the entire complexified good-cut
space.  This argument uses the already supplied
\texttt{Th\allowbreak Def3} basis; it does not assume the missing
theta square.

For each column, solve the unique coordinate equation
\[
       w_{0,\allowbreak{} i}c_\alpha=\Theta_i j_i f_i(E_i e_\alpha).
 \tag{BTH.42}
\]
Applying the tests in (BTH.29) gives \(Ac_\alpha=b_\alpha\),
where every entry of \(A\) and \(b_\alpha\) lies in
\(\sigma(K)\).  The test separation and injectivity of
\(w_{0,\allowbreak{} i}\) imply that \(A\) has full column rank.
Select a nonzero maximal minor.  Cramer's rule on that
minor gives \(c_\alpha\in\sigma(K)^{r_i}\), including
the empty-dimensional case in the usual unique way.

By definition the source's
\(\mathrm{liftMatC}=\mathrm{relMat}(w_{0,\allowbreak{} i},\allowbreak{} tC_i)\).
Its relation is uniquely solvable by (BTH.42), so this
is that matrix of columns \(c_\alpha\).  The operation
\(\mathrm{kMat}_\sigma\) takes their unique preimages
in \(K\).  The matrix in (BTH.39) is exactly that
stored \(\mathrm{liftMat}\), and
\(\mathrm{cr.toL}(z)=f_i(E_i z)\) by the original
coordinate definitions.  Applying (BTH.41) to the
rational relation equation and then undoing the twist
proves (BTH.40).

The source proves a stronger target-space equality
elsewhere using \texttt{PLev\allowbreak W}; the containment above
is sufficient for this base matrix and follows
already from BTH7 and the common basis.
% END BTH11

% BEGIN BTH12
\paragraph{BTH12: the base theta square and its exact remaining producer.}
\label{claim:BTH12}
Under the preceding hypotheses, the original base square is
\[
 \begin{split}
 &D.\theta P\bigl(\mu_0,\allowbreak{} i+1,
            D.jMt(\mu_0,\allowbreak{} i+1)(Y.eMt(i,\allowbreak{} t_0)x)\bigr)\\
 &\hspace{30mm}=J_i(M_i\Pi_i x)
       \qquad(x\in K^{d_i},\ 0\le i<m).
 \end{split}                                                    \tag{BTH.43}
\]
Together with \(\mathrm{RowsHold0}_3\) on this same
bundle, it gives the small-level rank conclusion of SLR8.

\emph{Proof.}
Put \(r=x-E_i\Pi_i x\).
By (BTH.4), \(r\) is \(lH_i\)-orthogonal to
\(\operatorname{im}E_i\).  BTH10 gives
\[
       T_i jV_i(f_i(x))
             =T_i jV_i(f_i(E_i\Pi_i x)).
\]
BTH11 identifies the right side with
\(q_iJ_i(M_i\Pi_i x)\).  BTH7 identifies the left
side with \(q_i\) of the left side of (BTH.43).
The injectivity of \(q_i\) proves (BTH.43).
The identification of \(A_3,\mathrm{zsOf},\allowbreak{} s0Of\) and of
the maps \(E_i,\Pi_i,\allowbreak{} J_i,\allowbreak{} M_i\) with the bundle's stored fields is
required, and follows by substitution once the
same original \texttt{Tw\allowbreak In}/bundle reading is supplied.

Here is a sufficient remaining producer, with the
quantifiers exposed.  For every family
\(\mathcal F\in\mathcal A\), choose one eventual
threshold \(a_\mathcal F\).  For all
\[
  x\ge a_\mathcal F,\qquad x>\mathcal F.x_1,\qquad
           \alpha\in\mathcal F.I_x,                             \tag{BTH.44}
\]
use the \emph{same} original bundle
\(b=\sigma(\mathcal F,\allowbreak{} x,\alpha)\), where \(\sigma\)
is the common compatible selector.  On it supply:

\begin{enumerate}
 \item the actual \(R,\allowbreak{} F,\allowbreak{} Y,\allowbreak{} Z\) reading, its weak
       characterization and phase-one pins, and its
       identification with \(b.P.A_3,\allowbreak{} b.P.Zs,\allowbreak{} b.P.S_0\)
       and the stored \texttt{Th\allowbreak Def3} maps;
 \item for each \(i<m\), a smooth cuspidal
       multiplicity extension of (BTH.10) for the
       finitely many vectors of one basis of the
       original \(H_i\);
 \item for each \(i<m\), one finite separating
       test family (BTH.29), simultaneously rational
       on the full rational source and on the
       original target-coordinate basis;
 \item \(\mathrm{RowsHold0}_3\) on this same bundle.
\end{enumerate}

Items 1--3 give (BTH.43); item 4 and SLR8 give
the base-rank equality.  The elementary passage
from the family threshold to the original
\(\mathrm{SmallLevelRankEq}\) uses the threshold
\(\max(a_\mathcal F,\mathcal F.x_1+1)\).
If this producer is supplied only in the source's
weaker \texttt{Closed} form, its own \(\varepsilon\)
and \(K\) bounds must be retained.  Nothing here
converts such a weaker producer into a universal
one with an arbitrary new bound.
% END BTH12

\subsection*{Removing the finite-action normalization ambiguity}

The next three arguments supply the actual-action assertion used
in BTH5.  The canonical action of CW23 alone would not prove a
statement about a different, already stored
\texttt{Is\allowbreak Weil\allowbreak Datum\allowbreak O}.  The comparison below is on the original
theta functions, with the same additive and splitting characters.

% BEGIN BTH13
\paragraph{BTH13: a continuous adelic unitary character is determined by rational and split points.}
\label{claim:BTH13}
Let \(E/F\) be the original CM extension, \(V\) any
nondegenerate Hermitian \(E\)-space, and
\(\xi:U(V)(\mathbb A_F)\to S^1\) a continuous character.
Suppose that \(\xi\) is trivial on \(U(V)(F)\) and on every
finite element supported at a split base-field place.  Then
\(\xi=1\).  No definiteness hypothesis on \(V\) is required.

\emph{Proof.}
Fix a full lattice \(L\).  For a nonisotropic rational vector
\(v\), its Hermitian projector \(P_v\) is rational, idempotent,
and self-adjoint.  The quasi-reflection
\[
 R_{v,\allowbreak{} u}=1+(u-1)P_v,\qquad u\bar u=1,
 \quad R_{v,\allowbreak{} u}^{-1}=R_{v,\allowbreak{} u^{-1}}                             \tag{BTH.45}
\]
is unitary.  Such reflections generate \(U(V)(F)\).
Here is the argument, including the indefinite case.
If \(x,\allowbreak{} z\) have the same nonzero norm and \(a=z-x\) is
nonisotropic, put \(A=h(z,\allowbreak{} a)\), using the convention linear
in the first variable.  Then \(h(a,\allowbreak{} a)=A+\bar A\ne0\), and
\(u=-\bar A/A\) gives \(R_{a,\allowbreak{} u}z=x\).
If \(z-x\) is isotropic, \(z+x\) has norm \(4h(x,\allowbreak{} x)\ne0\);
the preceding construction sends \(z\) to \(-x\), and
\(R_{x,-1}\) sends \(-x\) to \(x\).
A nondegenerate Hermitian space has an orthogonal basis of
nonisotropic vectors: otherwise polarization would make a
nonzero remaining form identically zero.  Fix its first
vector by at most two reflections and induct on its
orthogonal complement.  This is WA10's algebraic
generation proof, now over the rational field \(E/F\).

We first prove triviality on the infinite factor.
Write a rational norm-one \(u\) as
\(u=\alpha/\bar\alpha\) with \(0\ne\alpha\in\mathcal O_E\).
For completeness, choose \(b\) with \(b+u\bar b\ne0\)
and take \(\alpha=b+u\bar b\), then clear its rational
denominator.  Such \(b\) exists because otherwise \(b=1\)
would give \(u=-1\) and then conjugation would fix every \(b\).
For all sufficiently large integers \(n\), put
\[
 N_n=n!,\qquad
 u_n=\frac{1+N_n\alpha}{1+N_n\bar\alpha}.                     \tag{BTH.46}
\]
The numerator and denominator are nonzero, and \(u_n\bar u_n=1\).
At each infinite place \(u_n\to u\); at each fixed finite
place \(u_n\to1\).  At every nonsplit finite place it is a
unit, since conjugation preserves valuation and the two
valuations in (BTH.46) agree.

The projector \(P_v\) has lattice-integral entries outside
finitely many places.  At its finitely many nonsplit
denominator places, \(u_n-1\) and \(u_n^{-1}-1\) tend to zero.
Thus, for all sufficiently large \(n\), the finite component
of \(R_n=R_{v,\allowbreak{} u_n}\), and its inverse, preserve \(L\) at
every nonsplit place.  Remove its factors at precisely the
finitely many split places where lattice integrality or
inverse integrality fails.  Let \(s_n\) be their product and
let \(k_n\) be the remaining finite component.  Then
\[
 R_n=R_{n,\infty}\,\allowbreak{} s_n\,\allowbreak{} k_n,\qquad
 k_n\in K_L:=\prod_{w<\infty}U(L_w).                         \tag{BTH.47}
\]
Each \(U(L_w)\) is a closed subgroup of
\(\operatorname{GL}(L_w)\), hence compact.  The product
\(K_L\) has its product topology as a compact open
subgroup of the finite adelic group.
For any fixed finite place \(w\), \(R_n\to1\) there;
eventually its component is integral with integral inverse,
so no factor is removed there and \((k_n)_w\to1\).
Consequently \(k_n\to1\) inside this one fixed compact
product.  This is the required restricted-product
convergence, not merely componentwise convergence of
unbounded adelic elements.

Rational triviality gives \(\xi(R_n)=1\), and split
triviality gives \(\xi(s_n)=1\).  Continuity in (BTH.47)
therefore proves
\(\xi(R_{v,\allowbreak{} u,\infty})=1\).
Reflection generation proves the same for the infinite
component of every rational unitary matrix.

Those components are dense in the full infinite unitary
group.  More generally, rational unitary matrices satisfy
weak approximation at any finite set of places.  To see
this, multiply all the desired local matrices by one
rational norm-one scalar \(z\) so that \(1+zg_w\) is
invertible at each selected place.  There are only
finitely many forbidden scalars; the rational norm-one
scalars are infinite, for instance by the ratios
\((t+\beta)/(t+\bar\beta)\) with \(t\in F\) and
\(\beta\notin F\).
The Cayley transform
\[
 X=(1-zg)(1+zg)^{-1},\qquad
 X^*=-X,\qquad zg=(1-X)(1+X)^{-1}                            \tag{BTH.48}
\]
identifies this open set with an open subset of the
\(F\)-vector space of skew-Hermitian matrices.
Field weak approximation applied to its coordinates,
and openness of the nonzero determinant conditions,
give a rational approximation; undo \(z\).
This proof applies to indefinite forms as well.
Continuity now gives \(\xi=1\) on the entire infinite
factor.

For the finite factor we prove the required approximation
with a single fixed lattice.  Given a finite adelic
unitary \(g\) and an integer \(N>0\), choose a finite set
\(S\) containing all primes over \(N\) and every place
where \(g\) or \(g^{-1}\) fails to preserve \(L\).
By (BTH.48), choose rational \(\gamma_0\) so that
\(\gamma_0g\) is integral and invertible at \(S\) and is
congruent to \(1\) modulo \(NL_w\) at the primes over \(N\).
Factor \(\gamma_0\) into the rational reflections (BTH.45).

Correct each reflection parameter while preserving this
precision at \(S\).  The elementary ratio
\(\alpha/\bar\alpha\), together with field weak
approximation on nonzero local choices of \(\alpha\),
approximates any prescribed norm-one parameters at
finitely many places.  Choose the corrected parameter
close to its old value at \(S\), and close to \(1\) at
every other nonsplit denominator place of its projector.
Every rational norm-one parameter has valuation zero at
a nonsplit prime.  Hence the corrected reflection and
its inverse preserve \(L\) at every nonsplit place
outside \(S\); the closeness to \(1\) clears the finitely
many projector denominators.  Choose enough precision
at \(S\) that the corrected product \(\gamma\) retains
the original integrality and congruence there.
This is QD13's correction argument; all of its ingredients
have just been checked without a definiteness assumption.

The element \(\gamma g\) can now fail lattice
stabilization only at finitely many split places outside
\(S\).  Let \(s\) equal its components there and \(1\)
elsewhere; let \(k\) equal \(1\) there and its components
elsewhere.  These support cuts are unitary, commute, and
give the exact factorization
\[
 g=\gamma_f^{-1}s k,\qquad k\in K_L,\qquad
       (k-1)L_w\subset NL_w\quad(w\mid N).                   \tag{BTH.49}
\]
This proves the needed generalization of QD14.
Take \(N=n!\).  The resulting \(k_n\) converge to \(1\)
in \(K_L\).
The rational identity and the already proved infinite
triviality give \(\xi(\gamma_f)=1\);
split triviality gives \(\xi(s)=1\).
Thus \(\xi(g)=\xi(k_n)\to1\), proving finite triviality.
The infinite and finite factors together prove the claim.
These are the original quantifiers of
\texttt{unitary\_\allowbreak adelic\_\allowbreak character\_\allowbreak trivial\_\allowbreak rat\_\allowbreak and\_\allowbreak split}.
% END BTH13

% BEGIN BTH14
\paragraph{BTH14: full Weil data with the same characters agree on the original theta space.}
\label{claim:BTH14}
Let \(\omega,\omega'\) be two full
\texttt{Is\allowbreak Weil\allowbreak Datum\allowbreak O} actions for the same
\(D,\psi,\allowbreak{} S,\allowbreak{} T,\chi_S,\chi_T\), on the original functions.
Then for every theta function \(f\) and every adelic
pair \(p\),
\[
                         \omega(p)f=\omega'(p)f.             \tag{BTH.50}
\]
In particular this applies to the self-dual lattices
and the already stored action used in BTH5.

\emph{Proof.}
The operator \(\omega'(p^{-1})\omega(p)\) preserves
theta functions, commutes with every Heisenberg
translation by the two covariance laws, and preserves
the original Haar norm by the two \texttt{om\_\allowbreak unit}
laws.  CW15's scalar lemma, including its norm
hypothesis, therefore gives
\[
          \omega(p)f=c(p)\omega'(p)f,\qquad |c(p)|=1.          \tag{BTH.51}
\]
If the theta space is zero there is nothing to prove.
Otherwise fix a nonzero theta function \(f_0\).
The operator \(\omega'(p)\) is invertible on that
space, so \(c(p)\) is unique.  Applying the two group
laws to \(f_0\) proves that \(c\) is a character.
It is continuous: near any \(p_0\), choose \(w\) with
\((\omega'(p_0)f_0)(w)\ne0\).  The quotient
\[
       c(p)=\frac{(\omega(p)f_0)(w)}
                        {(\omega'(p)f_0)(w)}
\]
there is continuous by the literal pointwise
\texttt{om\_\allowbreak cont} laws.
Translate \(f_0\) by a Heisenberg operator if needed
so that \(f_0(0)\ne0\).  Both rational
\texttt{om\_\allowbreak aut} laws evaluated at \(0\) then give
\(c(\gamma,\gamma')=1\) for every rational pair.

For a split-supported pair \(p\), choose the same
determinant lifts \(x,\allowbreak{} y\) in both split clauses and
write \(z=\chi_T(x)^{-1}\chi_S(y)\).
They supply positive real numbers \(r,\allowbreak{} r'\) with
\[
 \ell(\omega(p)f)=rz\ell(f),\qquad
 \ell(\omega'(p)f)=r'z\ell(f)                                \tag{BTH.52}
\]
for every functional invariant under the specified
full split half.  There are such \(\ell,\allowbreak{} f\) with
\(\ell(f)\ne0\).
Here is why the earlier coinvariant proof applies.
The finite character projectors of CW3, on the
standard self-dual finite Hom lattice, express any
nonzero theta vector as a finite sum of translates
of lattice-fixed vectors, so at least one is nonzero.
The required self-dual lattices are the images of
the standard lattices already stored in \(S,\allowbreak{} T\)
under their frames; their duals transform by the
same isometry, as checked in CW24.
On each of the two split half-lattices the corrected
Heisenberg action in CW3 equals the raw action:
the phase is zero by isotropy of that half.
WA5's algebraic coinvariant argument applied to
this vector gives the required functional and
theta vector.  This construction concerns only
Heisenberg operators and lattices, not either
choice of Weil action.
Substitute (BTH.51) into (BTH.52) and cancel
the nonzero \(z\ell(f)\): \(c(p)r'=r\).
Since \(|c(p)|=1\) and \(r,\allowbreak{} r'>0\), it follows
that \(r=r'\) and \(c(p)=1\).

Restrict \(c\) separately to \(U(S)(\mathbb A_F)\)
and \(U(T)(\mathbb A_F)\).  Each restriction meets
all hypotheses of BTH13, so both are trivial.
The decomposition \(p=(p_S,1)(1,\allowbreak{} p_T)\) gives
\(c(p)=1\), proving (BTH.50).
The assertion is equality on theta functions;
no equality of arbitrary extensions to the ambient
space of all functions has been asserted.
This proves the source's
\texttt{weil\allowbreak Datum\_\allowbreak unique\_\allowbreak on\_\allowbreak theta\allowbreak Functions}.
% END BTH14

% BEGIN BTH15
\paragraph{BTH15: the original action has the exact finite lattice-stabilizer character.}
\label{claim:BTH15}
Use the original self-dual lattices \(L_S,\allowbreak{} L_T\),
their frames, a supplied \(\eta:\texttt{Is\allowbreak Eta\allowbreak Tilde}\),
and any full original
\(\omega:\texttt{Is\allowbreak Weil\allowbreak Datum\allowbreak O}(\psi,\allowbreak{} S,\allowbreak{} T,\chi_S,\chi_T)\).
Suppose \(\chi_S,\chi_T\) are unramified away from
the finitely many split pairs selected by an idempotent
\(e\).  If \(g,\allowbreak{} h\) are finite lattice stabilizers and
\(f\) is an original theta function satisfying the
original \texttt{Is\allowbreak Lat\allowbreak Vac} law for the Hom lattice
\(\Lambda\), then
\[
 \omega(g,\allowbreak{} h)f=
 \chi_T(\det\nolimits_e g)^{-1}
          \chi_S(\det\nolimits_e h)\,\allowbreak{} f.                       \tag{BTH.53}
\]
Here \(\det_e g=e\det g+1-e\), with the source's
idele interpretation; the two characters are crossed.

\emph{Proof.}
First let the pair be supported at one split place.
Its action preserves \(\Lambda\) in both directions
and commutes with its split-half projections.
Choose the literal determinant lifts \(x,\allowbreak{} y\), and
let \(r>0\) be the number supplied by
\(\texttt{om\_\allowbreak split}\), with
\(z=\chi_T(x)^{-1}\chi_S(y)\).
The theta function
\[
                         d=\omega(g,\allowbreak{} h)f-rzf
\]
is fixed by both complementary split half-lattices:
this holds for \(f\) by \texttt{Is\allowbreak Lat\allowbreak Vac}, and for
\(\omega(g,\allowbreak{} h)f\) by covariance and lattice
stabilization.  If \(d\ne0\), the WA5 coinvariant
proof applied to \(d\) produces a linear functional
\(\ell\), invariant under the full required half,
with \(\ell(d)\ne0\).
Self-duality supplies the separation used in that
proof; finite-level smoothness supplies the
telescoping sums.  The same
\(\texttt{om\_\allowbreak split}\) clause applied to \(f\)
gives \(\ell(d)=0\), a contradiction.
Thus \(\omega(g,\allowbreak{} h)f=rzf\).

If \(f=0\) the desired identity is immediate.
Otherwise its continuous norm-square descends to
the compact rational additive quotient of the
Hom adelic space.  Its Haar integral is finite
and positive: compactness bounds it, and a point
where it is nonzero has an open neighborhood of
positive Haar measure.  The norm identity
\(\texttt{om\_\allowbreak unit}\) now gives
\(r^2\|f\|_2^2=\|f\|_2^2\); since \(r>0\),
we obtain \(r=1\).  This proves the exact split
normalization for the supplied action itself.

Next let the pair be supported at a nonsplit place
\(v\), with both characters unramified there.
CW23 constructs a canonical full action \(\omega_0\)
on these same spaces and lattices, with characters
\(\eta^{\dim S},\eta^{\dim T}\) and its proved
local fixed-vacuum law.  Put
\[
 \nu_S=\eta^{-\dim S}\chi_S,\qquad
 \nu_T=\eta^{-\dim T}\chi_T.
\]
They are Hecke characters trivial on the ideles of
\(F\), by the exact rank restrictions.
The determinant-root construction WA18 gives
\[
 \omega_1(g,\allowbreak{} h)=\widetilde\nu_T(g)^{-1}
                  \widetilde\nu_S(h)\omega_0(g,\allowbreak{} h),             \tag{BTH.54}
\]
a full datum with the same characters as the
given \(\omega\).
BTH14 therefore identifies \(\omega\) and
\(\omega_1\) on every original theta function.

The scalars in (BTH.54) are one at this pair.
Indeed a norm-one determinant \(a\) at \(v\) is
a unit, since conjugation preserves valuation.
The original host's \texttt{cm\allowbreak Unram} condition
makes the involution nontrivial on the residue
field at every conjugation-fixed prime.
Choose an integral \(b\) for which
\[
                  x=b+a\bar b
\]
has nonzero residue.  Such \(b\) exists: if the
expression vanished for every residue \(b\), then
\(b=1\) would give \(a=-1\) and the remaining
identities would make the residue involution
trivial.  Thus \(x\) is a unit and
\(x=a\bar x\), so \(a=x/\bar x\).
Take this determinant root at \(v\) and \(1\)
elsewhere.  Both \(x\) and \(x^{-1}\) are
integral finite ideles, supported at \(v\).
The original characters and \(\eta\) are
unramified there, so their values on \(x\)
are one.  The same construction applies to
the other determinant, proving the assertion.

The canonical local law asks for fixing under
every Heisenberg vector in the local lattice.
The original \texttt{Is\allowbreak Lat\allowbreak Vac} first gives fixing
under \(2u\).  The host makes every dyadic
place split, so \(2\) is a unit at the present
nonsplit \(v\).  Divide any local lattice vector
by \(2\), still inside that lattice, to obtain
the required stronger local fixing.  Thus the
canonical \texttt{Fixes\allowbreak Vac\allowbreak At} applies, and
(BTH.54) gives \(\omega(g,\allowbreak{} h)f=f\) for this
original action.

For the general finite stabilizer pair, first
truncate to finitely many places.
At a selected split pair the first part gives
exactly its determinant factor in (BTH.53).
At an unselected split pair both determinant
lifts are integral units: a matrix and its
inverse preserve the framed lattice, so its
determinant and inverse determinant are integral.
The two unramified characters therefore give
factor one.  At each nonsplit place the preceding
argument gives factor one.  Multiplication proves
(BTH.53) for every finite-support truncation.
All these truncations lie in the same compact
product of framed lattice stabilizers and
converge to the given pair.  The original
\texttt{om\_\allowbreak cont}, the continuity of the
characters, and the continuous determinant cuts
let us pass to the limit at every function
argument.  This proves (BTH.53).

The precise source chain is recorded in the following files:
\begin{quote}\small
\texttt{Weil\allowbreak Datum\allowbreak Unique\allowbreak On\allowbreak Theta\allowbreak Functions};\\
\texttt{Weil\allowbreak Datum\allowbreak Fixes\allowbreak Vacuum\allowbreak Non\allowbreak Split\allowbreak Place\allowbreak Unramified};\\
\texttt{Weil\allowbreak Operator\allowbreak Lattice\allowbreak Stabiliser\allowbreak Fixes\allowbreak Vacuum\allowbreak Nonsplit};\\
\texttt{Split\allowbreak Place\allowbreak Weil\allowbreak Rep\allowbreak Lattice\allowbreak Vacuum\allowbreak Character};\\
\texttt{Weil\allowbreak Lattice\allowbreak Vacuum\allowbreak Stabilizer\allowbreak Splitting\allowbreak Characters}
(\texttt{stab\_\allowbreak vacuum}).
\end{quote}
The canonical-to-original comparison is now proved;
it is not an additional finite-covariance producer.
% END BTH15

\subsection*{What this closes, and what it leaves open}

The mathematical arguments BTH1--BTH15 are category B
proofs of their stated implications.  In particular,
BTH3 closes the actual integral-convergence guard
once the original models and weak Weil data are
supplied.  BTH4 constructs the good-projected Hilbert
lift, rather than leaving a second lift-existence
problem.  BTH6 reconstructs the required target
unitarity and base equivariance.  BTH10 is the full
finite orthogonal-kernel proof, and BTH11 replaces
the unnecessarily strong common-image package by
a consequence of the same stored basis.
BTH13--BTH15 remove the scalar ambiguity between
the canonical Weil construction and the original
action, proving its exact finite stabilizer character.

The remaining existence obligations are not classified
as B merely because the resulting formulas are short.
The precise residual table is:

\begin{center}
\begin{tabular}{p{0.18\linewidth}p{0.60\linewidth}p{0.10\linewidth}}
\hline
Obligation & Required original data & Class\\ \hline
Base reading &
Actual \texttt{Read}/\texttt{Read\allowbreak L}, their
\texttt{Lift\allowbreak Tie.\allowbreak theta}, \texttt{mt\allowbreak Read}, \texttt{mj\allowbreak Read},
weak models and pins, identified with the selected
bundle's row family and basis.  An unrelated model
with equal dimensions is insufficient. & C\\
Base extension &
For a basis of each original \(HV(i,\allowbreak{} t_0)\), extend
its now-convergent lowest-type integral to
\(\mathrm{Wall}_i\).  This includes the actual
smooth, equivariant, cuspidal multiplicity property. & C\\
Base rational tests &
The simultaneous finite separating tests (BTH.29)
for the exact field, embedding, theta map, and
target basis. & C\\
Base rows &
\(\mathrm{RowsHold0}_3\) on the same bundle;
its primitive row construction is a separate
producer. & C\\ \hline
\end{tabular}
\end{center}

The first three entries refine the theta part of the
previous joint same-bundle producer.  They are
shared mathematical existence tasks, not three
independent freely choosable witnesses: every
entry uses the same \(R,\allowbreak{} F,\allowbreak{} Y,\allowbreak{} Z,\omega_i,\allowbreak{} io_i,\allowbreak{} pt_i\),
measures, field, and stored weak choices.  There
is no new category-D instance ambiguity in the
displayed implication once these identifications
are supplied.  Their actual construction remains
unclosed, and the original joint producer must
not be marked complete.

For the remaining extension, the pinned source gives a
concrete next mathematical target.  In
\texttt{Exists\allowbreak Theta\allowbreak Lift\allowbreak Linear\allowbreak Map\allowbreak Cusp\allowbreak Form\allowbreak Multiplicity},
the declaration \texttt{fires\allowbreak At}
proves the pointwise extension by first establishing
smoothness, uniform moderate growth, zero transvection
constant terms, compact-type equivariance, and annihilation
by the lowering operators of the holomorphic discrete
series.  Its \texttt{exists\_\allowbreak extension} then realizes
these functions in the prescribed multiplicity model.
The exact suppliers include:
\begin{quote}\small
\texttt{theta\allowbreak Lift\_\allowbreak definite\_\allowbreak smooth\_\allowbreak moderate\_\allowbreak growth\_\allowbreak annihilated};\\
\texttt{theta\allowbreak Lift\_\allowbreak definite\_\allowbreak constant\allowbreak Term\_\allowbreak eq\_\allowbreak zero};\\
\texttt{fock\_\allowbreak vectors\_\allowbreak killed\_\allowbreak heisenberg\_\allowbreak average\_\allowbreak isotropic};\\
\texttt{bounded\allowbreak Holomorphic\allowbreak Form\_\allowbreak eq\_\allowbreak multiplicity\allowbreak Space\_\allowbreak value}.
\end{quote}
Their substantial analytic and representation-theoretic
proofs have not been replaced by naming a published
theorem.  The source's use of
\texttt{Published.\allowbreak Cusp\allowbreak Rapid\allowbreak Decay} in that route is
likewise not a new category-A certificate in this chapter.

The reduced base route does not need the remaining
fields of \texttt{Ex\allowbreak At}, the chain lift \(\mathrm{Th}\),
an independent \(\mathrm{thG}\) existence theorem,
a separate \texttt{Lift\allowbreak Test}, nonbase test families,
\texttt{PLev\allowbreak W}, \texttt{LUnif}, \(\mathrm H5\),
a positive theta-norm formula, \texttt{Rows\allowbreak Hold\allowbreak Lam3},
or the all-tuple \texttt{Tup\allowbreak All\allowbreak Used} comparison.
This is a reduction of the sufficient producer for
this consumer.  It does not establish existence
of the shared selector, the original analytic
reading, or the full RH/GRH proof.


\section{The original late-family setting and the local Hecke estimates}
\label{sec:global-local-estimates}

\status{This chapter reconstructs the numerical setting and the finite local
arguments on the original objects.  A supplied family, its original forms,
and their same-object ties are retained throughout.  The source is the
pinned commit recorded in the accompanying provenance ledger.}

\paragraph{Objects and conventions.}
Fix a source member $D$, a word $u$, and write
$K=D.\mathrm{what}_2.\mathrm{LL}$, $k=D.R.k=N(0)$,
$m=D.R.m$, $N=N(m)$, $f=D.R.f$, and $p=(D.\mathrm{mem}\,\allowbreak{} u).A.T$.
The prime $p$ is the original word prime.  At a coefficient place $w\mid p$,
write $\mathcal O_w$ for the original DVR and $e_w=v_w(p)$.
An original level-zero forms record is denoted $\mathcal F$; this avoids
using $T$ for both that record and the prime.  In the source its type is
\texttt{Forms0 D u}.  Its maps \texttt{at\allowbreak T}, \texttt{cls0}, \texttt{val0},
\texttt{lv0}, \texttt{up\_\allowbreak natural\_\allowbreak level\_\allowbreak zero},
\texttt{down\_\allowbreak natural\_\allowbreak level\_\allowbreak zero}, and \texttt{bw0} are the original
ones, obtained by projecting the supplied symmetric forms record.

\subsection{A direct size proof for the actual weights}

% BEGIN GLE1
\paragraph{GLE1: the original parameter windows give a common weight interval.}
\label{claim:GLE1}
Let $L:\texttt{Size\allowbreak Data}_2(D)$ and let $P:\texttt{Parameter\allowbreak Choice}'(D,\allowbreak{} L)$
be the supplied parameter record.  Put
\[
 \begin{gathered}
 x=L.\mathrm{sc}.x,\quad b=e^{x^\varepsilon},\quad a=xb,
 \quad X'=e^{x^{\varepsilon/3}},\quad b'=a/X',\\
 w''=b'/(xn),\quad S=L.\mathrm{Sc},\quad \eta=L.\mathrm{etaSz}.
 \end{gathered}
 \tag{GLE.1}
\]
Assume $x\ge4$ and $X'\ge x$.  Then, for every original word and every
embedding, every entry of every original lifted weight with $0\le j\le m$
satisfies
\[
 -2b\le\lambda_{j,\tau}(r)\le2xb.                 \tag{GLE.2}
\]
Moreover,
\[
 b\ge24x,\qquad b\ge3N,\qquad N\le b/4,
 \qquad a_{\max}\le2xb,
 \qquad b_{\mathrm{top},\tau}>0\quad(\tau\in\mathrm{tops}(D)).
 \tag{GLE.3}
\]

Here is the complete extraction from the parameter fields.  Since
$h_0=\eta/m$, the ninth inequality of \texttt{h0\_\allowbreak ge} gives
$Nm^{11}\le\eta$.  The field \texttt{mteta} gives
$m(t+1)\eta\le w''$, and \texttt{k\allowbreak S\_\allowbreak le} gives $kS\le w''$.
The rank laws give $m\ge2$, $k\ge5$, $n\ge1$, $N\ge1$, $t\ge1$;
\texttt{t\_\allowbreak le\_\allowbreak eta} gives $\eta\ge1$.  Consequently
\[
 Nm^{12}\le w'',\qquad 1\le\eta\le w'',\qquad
 0<S\le w'',\qquad b'=xnw'',\qquad 0<b'\le b.
 \tag{GLE.4}
\]
The last inequality follows directly from $b'=xb/X'$.
In particular $w''\ge4096N$, $b'\ge4096xN$, and
$w''\le b'/x\le b'/4$.  These imply the first three bounds in (GLE.3).

We check the first $k$ entries before the appended blocks.  Write
$u_r$ for the original half-integral Harish--Chandra entries and
\[
 \lambda_r=u_r-\bigl((k-1)/2-r\bigr).
 \tag{GLE.5}
\]
The intervening entries are integral and strictly interlace the $u_r$.
Thus $u_r-u_{r+1}\ge1$, and (GLE.5) is decreasing.  At a top place the
first and last entries are exactly $a_T$ and $-b_T$.  The original
window for $a_T$ gives $a\le a_T\le2a$.
For $b_T$, put $q=N_*x^\varepsilon L.\mathrm{LaG}/k$.
The fields \texttt{S\_\allowbreak le}, $S>0$, and $b'>0$ imply
$L.\mathrm{LaG}>0$.  The definition of \texttt{rho\allowbreak P} is
\[
 \rho_P=4q+C_\rho\log x\,\allowbreak{} L.\mathrm{LaG}/k.
\]
All terms on the right are nonnegative, and
$\rho_P\le x^{-2}\le1$.  Hence $q\le1/4$ and the actual
\texttt{b\allowbreak T\_\allowbreak le} and \texttt{b\allowbreak T\_\allowbreak ge} fields give
\[
 b'\le b_T\le\tfrac54 b'.                       \tag{GLE.6}
\]
Every middle top entry belongs to $[0,\allowbreak{} S]$.  Indeed $u_1\le S$ and
the shift $(k-1)/2-1$ is nonnegative, so $\lambda_1\le S$.
At the other endpoint $u_{k-2}\ge0$ and the shift
$(k-1)/2-(k-2)$ is nonpositive, so $\lambda_{k-2}\ge0$.
The monotonicity of (GLE.5) places all intermediate entries between
these two bounds.  Off the top set, \texttt{uk\_\allowbreak extreme} and the
last cloud bound give $\lambda_{k-1}\ge0$; the first cloud bound and
\texttt{cloud\_\allowbreak place\_\allowbreak bound} give $\lambda_0\le3S$.
Thus all off-top entries belong to $[0,3S]$.
Since $3S\le3b'/4\le b\le a$, these entries also satisfy (GLE.2).
The maximum defining $a_{\max}$ is the maximum of $0$ and the original
top entries $a_T$; hence $a_{\max}\le2a$.  The original definition
$b_{\mathrm{top},\tau}=-\lambda_{0,\tau}(k-1)$ identifies it with
$b_T$ at every top place, proving its positivity.

For the appended block indexed by $0\le r<m$, the original weight
formula is $-w_r$, where
\[
 w_r=w_{r+1}^{\mathrm{raw}}-\delta_r,
 \qquad
 \delta_r=\frac{k}{2}+\frac{(2r+1)(t+1)}2,
 \qquad 0\le\delta_r\le N.                       \tag{GLE.7}
\]
The three raw windows are
\[
 \begin{array}{c|c}
 \text{top}&[b_T+b'/x,\ b_T+b'/x+\eta]\\
 \text{thick}&[b,\ b+\eta+1]\\
 \text{thin}&[w'',\ w''+\eta].
 \end{array}                                      \tag{GLE.8}
\]
Every lower endpoint is at least $w''\ge N$.
The top upper endpoint is at most $7b'/4$, the thick one at most
$b+b'/2$, and the thin one at most $b'/2$; here $1\le w''\le b'/4$.
Subtracting (GLE.7) therefore gives $0\le w_r\le2b$.
The source's literal \texttt{lam\_\allowbreak eq} identifies the first $k$ entries
with (GLE.5) and every remaining entry with one of these $-w_r$.
This proves (GLE.2) for all $j$, including $j=0$ and $j=m$.
No general \texttt{Nums\allowbreak L} conclusion is being presumed in this proof.
% END GLE1

% BEGIN GLE2
\paragraph{GLE2: \texttt{Tu\allowbreak Size} supplies every clause of the eventual setting.}
\label{claim:GLE2}
For every original family $\mathrm{Fm}:\texttt{Fam\allowbreak P}\,\varepsilon$,
\[
 \texttt{Tu\allowbreak Size}(\mathrm{Fm})
 \quad\Longrightarrow\quad
 \texttt{Setting\allowbreak Eventually}(\mathrm{Fm}).          \tag{GLE.9}
\]
The conclusion means that there is one $A$ such that, for every
$x\ge A$ with $x>\mathrm{Fm}.x_1$, every $i\in\mathrm{Fm}.I_x$ and
every word $u$ of $D=\mathrm{Fm}.D(x,\allowbreak{} i)$ satisfy all four original
\texttt{Hida\allowbreak Setting} requirements.  The threshold is chosen before
$i$ and $u$.

To check its existence directly, put $\delta=\varepsilon/3>0$.
For $x\ge1$, $\log x\le x^{\delta/2}/(\delta/2)$, by
$\log y\le y$ with $y=x^{\delta/2}$.
Thus $x^{\delta/2}\ge2/\delta$ implies $x\le e^{x^\delta}=X'$.
Increase a single family threshold until this holds, $x\ge32$, and
the \texttt{Tu\allowbreak Size} threshold has been reached.  All original parameter
and size records are then supplied by \texttt{Fm.\allowbreak params} and
\texttt{Fm.\allowbreak sizes}; their scale satisfies $L.\mathrm{sc}.x=x$ and
$L.\mathrm{sc}.\varepsilon=\varepsilon$.

The lower $m$-size field says
\[
 m\ge x^{-\varepsilon/100}\sqrt{x\,\allowbreak{} x^{-\varepsilon/8}}
   =x^{1/2-29\varepsilon/400}\ge x^{2/5}\ge4.
 \tag{GLE.10}
\]
The middle inequality uses the original $0<\varepsilon<1/52$.
The lower prime bound in \texttt{Tu\allowbreak Size}, namely
$p\ge b^{m/2}$, now gives $p\ge b^2$.
The original rank law is $f\le N$, so (GLE.3) yields
\[
 2fN+1\le2N^2+1<9N^2\le b^2\le p,
 \qquad 4a_{\max}\le8xb<24xb\le b^2\le p.
 \tag{GLE.11}
\]
For every $j\le m$ one has $N(j)\le N$.  By (GLE.2),
\[
 \begin{split}
 2\bigl(\lambda_{j,\tau}(r)-\lambda_{j,\tau}(s)+N(j)\bigr)
 &\le4xb+4b+2N\\
 &\le6xb<24xb\le b^2\le p.
 \end{split}                                      \tag{GLE.12}
\]
Indeed $N\le b/4$ and $x\ge4$ give $4b+2N\le2xb$.
This is the literal strict \texttt{Alcove} inequality, for every ordered
pair of entries.  The last clause $b_{\mathrm{top},\tau}\ge0$ follows
from GLE1.  Equations (GLE.11)--(GLE.12) prove the original setting,
including its strict inequalities, on one common threshold.

This proof needs only the lower \texttt{Tu\allowbreak Ge} part of \texttt{Tu\allowbreak Size}.
The upper bound for $\log p$ and the logarithmic derivative estimate
are not extra premises.  They occur in the source's stronger general
$\texttt{Nums\allowbreak L}/\texttt{Wts\allowbreak L}$ supplier, whose remaining conclusions
are not claimed here.
% END GLE2

\subsection{The exact polynomial coordinate lattices}

% BEGIN GLE3
\paragraph{GLE3: integral and rational changes of polynomial coordinates.}
\label{claim:GLE3}
For a dominant integral weight $\lambda$ of rank $k$, put
$s=\lambda_{k-1}$ and $\nu_r=\lambda_r-s$.  Let
$\Delta_{\nu,\mathbb Z}$ be the original \texttt{del\allowbreak Sq} lattice: the
integer span of all divided-power root words applied to the product of
principal minors of weight $\nu$.  For a coefficient DVR $\mathcal O_w$
put
\[
 \Lambda_w=\bigotimes_\tau
       (\mathcal O_w\otimes_{\mathbb Z}\Delta_{\nu_\tau,\mathbb Z}),
 \qquad V_K=K\Lambda_w.
 \tag{GLE.13}
\]
The tensor variables for different $\tau$ are disjoint.  After the
stored embedding $K\hookrightarrow\mathbb C$, these are exactly
\texttt{Int\allowbreak Span D w (lam 0)} and \texttt{Rat\allowbreak Span D (lam 0)}.
Indeed the definitions are the additive spans of precisely the pure
tensors in (GLE.13), with respectively integral-at-$w$ and arbitrary
$K$ coefficients.  Negative determinant shifts belong to the value
function, not to this polynomial lattice.

For an invertible matrix $C$, use the source convention
\[
 R_\lambda(C)P(Z)=\det(C)^s P(CZ).                 \tag{GLE.14}
\]
If $C\in\mathrm{GL}_k(\mathcal O_w)$, then (GLE.14) preserves the
original integral lattice, with equality; if $C\in\mathrm{GL}_k(K)$,
it preserves its rational span.  To prove this without changing
lattices, a transvection $I+cE_{ab}$ acts by the finite Taylor expansion
$\sum_{r\ge0}c^r E_{ab}^{[r]}$.  Every divided-power root operator
preserves the span of root words.  A diagonal unit preserves each
integer row-weight summand, and its determinant shift is a unit even
when $s<0$.  Gaussian elimination over a local ring factors every
invertible matrix into transvections, permutations and diagonal units:
a unit determinant forces a unit entry in the first pivot column after
a row permutation; eliminate that column and induct.  Row permutations
are products of transvections and diagonal units.  This proves the
assertion and its inverse.  The same argument over $K$ proves the
rational assertion.  Tensoring proves both assertions for independent
matrices at the different $\tau$.

The lattice $\Delta_{\nu,\mathbb Z}$ has a basis homogeneous in all row
degrees.  Every generating root word has such a degree, it is a
submodule of the finite free integer module of polynomials of the
fixed total degree, and its distinct row-degree pieces have zero
intersection.  Choose an integer basis of each piece and take their
union.  If $r_\lambda$ is its rank and $d_j(h)$ is the $j$th row degree
of the basis element $h$, row permutations give equal total degrees
in all $k$ rows.  The total degree of every $h$ is $|\nu|$.  Hence
\[
 \sum_h(d_j(h)+s)=\frac{r_\lambda}{k}|\lambda|,
 \qquad d_j(h)+s\ge s.                            \tag{GLE.15}
\]
This elementary degree identity, including the shift $s$, will compute
the local length below.  It needs no semisimplicity in the residue
characteristic.

Finally, the polynomial-to-value map in the source is injective on
these coordinates.  Its frames range over all invertible matrices,
and the determinant factors in \texttt{fn\allowbreak Of} are nonzero.  A polynomial
vanishing on all such tuples is zero: for any fixed tuple of matrices
replace each component $Z$ by $Z+tI$.  Its determinant is nonzero for
all but finitely many $t$.  Here a separate variable $t_\tau$ is used
for each component.  The resulting polynomial vanishes on a product
of cofinite sets, hence is zero by induction on the number of
variables.  Setting all $t=0$ proves the assertion on every tuple.
The ordinary polynomial identity principle over the infinite field
$\mathbb C$ finishes the proof.  Therefore the displayed coordinate
transports are equalities of the actual forms, not merely maps on a
larger ambient polynomial space.
% END GLE3

% BEGIN GLE4
\paragraph{GLE4: exact coordinate descriptions of the original local lattices.}
\label{claim:GLE4}
Assume the original level-zero stabilizers are trivial, as in
\texttt{adm\allowbreak S}; RSL7 supplies this from the actual forms in the selected
application.  Let $\mathcal C$ be the original finite class set.
For each $c\in\mathcal C$ choose a $p$-local basis $B_c$ of the original
dual class lattice
$H_c=\texttt{Ld D u 0 c}$.  The prescribed frame coordinates give a
$K$-linear isomorphism
\[
 V= ((D.\mathrm{fm}\,\allowbreak{} u).\mathrm{lat}\,0).V
       \ \simeq\ \bigoplus_{c\in\mathcal C}V_K.
 \tag{GLE.16}
\]
There are no further classwise constraints: the original
\texttt{cls0.\allowbreak stab\allowbreak I\_\allowbreak range} identifies the compatibility group at $c$
with its trivial stabilizer.  Arbitrary values at the preferred
representatives therefore extend uniquely by the original adelic
equivariance rule.  The \texttt{val0.\allowbreak rat} equality identifies those
whose coordinates lie in the original rational polynomial span.
Injectivity is GLE3 and the original injective embedding of $V$.
Thus both directions of (GLE.16) use the original rational form.

Write $S_w=(D.\mathrm{fm}\,\allowbreak{} u).S(0,\allowbreak{} w)$.  In (GLE.16) one has exactly
\[
                         S_w=\bigoplus_{c\in\mathcal C}\Lambda_w.
 \tag{GLE.17}
\]
First, two allowed local bases differ, after each relevant embedding,
by a matrix in $\mathrm{GL}_k(\mathcal O_w)$.  GLE3 therefore converts
the source's condition ``in every local basis'' into membership in
$\Lambda_w$ in the chosen basis.  To check the localization itself,
choose a nonzero integral scalar $d$ with $d z$ in the original global
lattice.  If $z$ has integral coordinates at $w$, choose
$t\in\mathcal O_K\setminus w$ which is divisible, at every other prime
dividing $(d)$, by at least the corresponding valuation of $d$.
Such $t$ exists by the Chinese remainder theorem, with $t\equiv1$
modulo $w$.  At $w$, the coordinates of $tz$ remain integral; elsewhere
$t/d$ is locally integral.  The original \texttt{val0.\allowbreak lat} criterion
now puts $tz$ in the global lattice, so $z\in S_w$.  Conversely a
denominator outside $w$ is a unit in every coordinate in (GLE.17).
This proves the equality without a new lattice choice.

Fix a seen place $w\mid p$ and a selected place $v$ of the original
\texttt{At\allowbreak TTied} record.  Put $\tau_v=(D.\mathrm{fm}\,\allowbreak{} u).\mathrm{tauOf}(w,\allowbreak{} v)$.
Let
\[
 K_{v,\allowbreak{} c}=\texttt{Lk}\,\mathcal F.\mathrm{atT}\,0\,\allowbreak{} v\,\allowbreak{} c
        =\texttt{cls\allowbreak Crd}(c)+P_vH_c,
 \qquad P_v=\mathcal F.\mathrm{atT}.P_v.
 \tag{GLE.18}
\]
Choose a $p$-local basis $B'_{v,\allowbreak{} c}$ of $K_{v,\allowbreak{} c}$.  Such bases exist
because localization at the finitely many primes above $p$ is a
semilocal Dedekind ring, hence a PID: every ideal is generated by an
element with its finitely many prescribed valuations.  A finite
torsion-free module over this ring is free; expressing that basis in
the original $E$-space gives the required matrix.  The lattices in
(GLE.18) are full: $H_c$ is a Hermitian dual of a full lattice, and
$\texttt{cls\allowbreak Crd}(c)\subset K_{v,\allowbreak{} c}\subset H_c$.

Set $C_{v,\allowbreak{} c}=(B'_{v,\allowbreak{} c})^{-1}B_c$ at $\tau_v$, and set $C_{v,\allowbreak{} c}=1$
at the other embeddings.  In the same coordinates the original
neighbor lattice is exactly
\[
 B_{w,\allowbreak{} v}=(D.\mathrm{fm}\,\allowbreak{} u).\mathrm{Bw}(0,\allowbreak{} w,\allowbreak{} v)
       =\bigoplus_c R_{\lambda_{0,\tau_v}}(C_{v,\allowbreak{} c})\Lambda_w.
 \tag{GLE.19}
\]
The forward implication is the original \texttt{bw0} clause.
For the reverse implication, each right-hand coordinate is a rational
polynomial by GLE3, since the entries of the basis matrices are in $E$
and the stored embedding is $K$-valued.  Use (GLE.16) to realize those
coordinates in $V$.  Changing either local basis changes the mixed
coordinate matrix by integral invertible matrices at the appropriate
embeddings.  GLE3 preserves $\Lambda_w$ under those changes, so all
the bases quantified in \texttt{bw0} satisfy the required condition.
This proves surjectivity as well as inclusion in (GLE.19).

In particular, if $r$ is the rank of the tensor lattice $\Lambda_w$,
then the literal source dimension is
\[
                  d_S=D.\mathrm{duS}(u)=|\mathcal C|r.
 \tag{GLE.20}
\]
The original \texttt{emb\_\allowbreak basis} identifies scalar extension of $V$
with the complex forms; GLE3 and (GLE.16) calculate its dimension.
No identification of $d_S$ with the smaller ramified dimension is made.
% END GLE4

\subsection{Integral maps and the second branch on those lattices}

% BEGIN GLE5
\paragraph{GLE5: the two original Hecke operators are integral.}
\label{claim:GLE5}
Supply the same original $\mathcal F$, its finite level character with
values in $K$, and \texttt{adm\allowbreak S}.  For each seen $w\mid p$ and every
selected $v$, the original operators satisfy
\[
              U_v(B_{w,\allowbreak{} v})\subseteq S_w,
              D_v(S_w)\subseteq B_{w,\allowbreak{} v}.           \tag{GLE.21}
\]
They are $K$-rational operators on $V$.  This is the complete literal
\texttt{Hecke\allowbreak Integral\allowbreak Level\allowbreak Zero}; it in particular proves the
source's version with the additional \texttt{Hida\allowbreak Setting} hypothesis.

We verify the geometric and coefficient parts separately on the same
move.  Work at $P_v$ and use the original residue map.  We first check
the precise uniformizer assertion, following the elementary norm
argument of FC1.  There are $f$ distinct selected primes, each of
residue cardinality $p$, and every one divides $\pi_u$.  The original
host has degree $[E:\mathbb Q]=2f$ and
$\pi_u\bar\pi_u=p$.  Thus the absolute ideal norm of $(\pi_u)$ is
$p^f$: its square is the norm $p^{2f}$ of $(p)$, and conjugation
preserves ideal norm.  The $f$ displayed primes already account for
this whole norm.  Each therefore occurs in $(\pi_u)$ with exponent
one.  The original nonzero residue of $\bar\pi_u$ makes it a unit at
each of them.  Consequently $p$ has valuation one at $P_v$, and
$P_v\mathcal O_{E,\allowbreak{} P_v}=p\mathcal O_{E,\allowbreak{} P_v}$.
Write $H$ for the localized $H_c$, and let $z$ lift an up index with
$f(z)=1$.  The exact up naturality fields, already reconstructed in
NB2 and NB7, say that its rational unitary transport $\gamma$ identifies
\[
 \gamma\,\texttt{cls\allowbreak Crd}(c^+)=H,\qquad
 \gamma H_{c^+}=H+\mathcal O_{E,\allowbreak{} P_v}p^{-1}z.
 \tag{GLE.22}
\]
Consequently
\[
 \gamma (K_{v,\allowbreak{} c^+})_{P_v}
   =H+p\bigl(H+\mathcal O_{E,\allowbreak{} P_v}p^{-1}z\bigr)=H.
 \tag{GLE.23}
\]
For a down index $g$ with $g(t_e)=1$, NB3 and NB8 identify the original
transport by
\[
             \gamma(H_{c^-})_{P_v}
                     =(\texttt{cls\allowbreak Crd}(c))_{P_v}=(K_{v,\allowbreak{} c})_{P_v}.
 \tag{GLE.24}
\]
For the last equality, $pH_c\subseteq\texttt{cls\allowbreak Crd}(c)$ follows from
the original base-lattice dual bound.  At each selected place $v'\ne v$
the transported dual lattices are equal, because the adelic move is
the identity there and the residual level factor preserves the base
lattice and its Hermitian dual.  To descend that completed equality
to rational localization, clear the finitely many denominators at the
moved place with an element which is a unit at $P_{v'}$.  Equivalently,
the two rational lattices have identical localizations away from the
support of their finite quotient.  This proves the exact
\texttt{Away\allowbreak Lat0} assertion used by the source.

Let $B_{\rm in}$ denote the appropriate incoming base or neighbor
basis, and $B_{\rm out}$ the outgoing one.  Equations (GLE.23)--(GLE.24)
and their away-place versions say precisely that
\[
       B_{\rm in}^{-1}\gamma^{-1}B_{\rm out}
                \in\mathrm{GL}_k(\mathcal O_{E,\allowbreak{} P_{v'}})
 \tag{GLE.25}
\]
at the prime indexed by each $v'$.  Its inverse is integral too.
Here the place dictionary on the CM field deserves a check:
\texttt{tau\allowbreak Of} states a residue equality on the totally real field,
whereas these matrices have entries in $E$.  Let $\nu$ be the
$w$-valuation after the stored embedding $E\to K$ at
$\tau=(D.\mathrm{fm}\,\allowbreak{} u).\mathrm{tauOf}(w,\allowbreak{} v')$.  It is nonnegative on
$\mathcal O_E$.  The original \texttt{seen\allowbreak P} gives $\nu(\pi_u)>0$.
The element $\pi_u+\bar\pi_u$ lies in the totally real field and has
nonzero residue at $P_{v'}$.  Thus \texttt{tau\allowbreak Of} gives it valuation
zero, and hence $\nu(\bar\pi_u)=0$.  If $y\in P_{v'}$, then
$y\bar\pi_u+\bar y\pi_u$ is totally real and has zero residue.
Its valuation is positive, as is that of $\bar y\pi_u$; subtraction
shows $\nu(y\bar\pi_u)>0$, so $\nu(y)>0$.
If $b\notin P_{v'}$, choose $a\in\mathcal O_E$ with
$ab-1\in P_{v'}$.  Were $\nu(b)>0$, both $ab$ and $ab-1$ would have
positive valuation, a contradiction.  Thus $\nu(b)=0$.
It follows that $\mathcal O_{E,\allowbreak{} P_{v'}}$ maps to $\mathcal O_w$,
with every denominator outside $P_{v'}$ mapping to a unit.  Applying
this to the matrix and its inverse in (GLE.25) gives an integral
invertible coefficient matrix at every $\tau$.

The original function transport is
$\chi(\kappa)^{-1}$ times substitution by $\gamma^{-1}$ for the
source's factorization $\gamma y_{c'}=y_c\eta\kappa$, including
the determinant powers in (GLE.14).  Its character factor is a unit
at every coefficient place: if $\chi(\kappa)=\iota_K(r)$ and
$\chi(\kappa)^q=1$ with $q>0$, injectivity gives $r^q=1$;
both $r$ and $r^{-1}$ are algebraic integers.  GLE3 applied to
(GLE.25) therefore sends each integral incoming polynomial to an
integral outgoing polynomial.  This is valid for every choice of
local bases, so GLE4 identifies the resulting conditions with the
actual lattices in (GLE.21).

Finally the operators are the original finite sums
\[
 (U_v\phi)(c)=\sum_{z:f_v(z)=1}\operatorname{upT}_{v,\allowbreak{} c,\allowbreak{} z}
                         \phi(c^+_z),\qquad
 (D_v\phi)(c)=\sum_{g:g(t_e)=1}\operatorname{dnT}_{v,\allowbreak{} c,\allowbreak{} g}
                         \phi(c^-_g).             \tag{GLE.26}
\]
There is no averaging denominator.  Naturality identifies these sums
with the finite adelic translates, and changing a representative
permutes their residue indices.  They consequently remain in the
same complex forms space.  Every summand has a $K$-rational
polynomial value: $\gamma$ is over $E$, its stored embeddings take
values in $K$, and the character factor is in $K$.  The reverse
direction of the original \texttt{val0.\allowbreak rat} condition therefore places
the sum in the original $V$.  Summing the integral coordinate
conditions proves both inclusions in (GLE.21), with exactly the
source's \texttt{Into} preimages under its injective embedding.
% END GLE5

% BEGIN GLE6
\paragraph{GLE6: the exact second-branch inclusion and length.}
\label{claim:GLE6}
For the same original forms and trivial stabilizers, fix seen $w\mid p$
and $v$, and put
$b_v=(D.\mathrm{fm}\,\allowbreak{} u).bTop(\tau_v)=-\lambda_{0,\tau_v}(k-1)$.
If $b_v\ge0$, then
\[
 \begin{split}
 p^{b_v}S_w&\subseteq B_{w,\allowbreak{} v},\\
 \ell_{\mathcal O_w}(B_{w,\allowbreak{} v}/p^{b_v}S_w)
   &=\frac{e_wd_S}{k}
           \bigl(|\lambda_{0,\tau_v}|+kb_v\bigr).
 \end{split}                                      \tag{GLE.27}
\]
In particular the displayed rational number is a nonnegative integer.
This is the full original \texttt{Second\allowbreak Branch\allowbreak Inclusion\allowbreak Length}.

The original residue functional on $H_c$ is nonzero, has scalar
residue-field target of cardinality $p$, and has kernel $K_{v,\allowbreak{} c}$.
This is exactly \texttt{lv0.\allowbreak fv\_\allowbreak ker}; surjectivity follows by scaling
one nonzero value.  Thus, over $\mathcal O_{E,\allowbreak{} P_v}$, this is a
colength-one sublattice of $H_c$.  Choosing a lift $z_0$ with
$f(z_0)=1$ and lifting a basis of its kernel gives bases with relative
matrix $\operatorname{diag}(p,1,\ldots,1)$.  The prime $p$ has valuation
one here, as in GLE5.  Any chosen global $p$-local bases have the same
Smith form after localization.  Under the original embedding at
$\tau_v$, their relative matrix therefore has the form
\[
 B_c^{-1}B'_{v,\allowbreak{} c}
       =A_c\operatorname{diag}(p,1,\ldots,1)C_c,
       \qquad A_c,\allowbreak{} C_c\in\mathrm{GL}_k(\mathcal O_w).
 \tag{GLE.28}
\]
This description is made only in the coefficient DVR for computing
the existing lattices; it does not replace their global bases.
The other factors in the tensor representation are unchanged.

By GLE3, the unit matrices in (GLE.28) preserve $\Lambda_w$, and
also preserve its scalar multiple $p^{b_v}\Lambda_w$.
It suffices to calculate the quotient with the inverse diagonal
matrix in (GLE.19).  In a homogeneous row-degree basis, its entries
are $p^{-\beta_h}$, where
$\beta_h=d_1(h)+\lambda_{0,\tau_v}(k-1)$.
For every such basis element,
$b_v+\beta_h=d_1(h)\ge0$.  Thus the inclusion holds and its elementary
divisor has length $e_w(b_v+\beta_h)$.  Tensor the basis with the
bases of the other factors and use (GLE.15).  On one class the length
is
\[
                  e_wr\bigl(b_v+|\lambda_{0,\tau_v}|/k\bigr).
\]
Sum over the same finite class set and use the exact identity
$d_S=|\mathcal C|r$ from (GLE.20).  This proves (GLE.27), including
finiteness and the exact normalization by $e_w$ and $k$.
The determinant shift is essential: omitting it would give the wrong
index when the lowest weight is negative.
% END GLE6

\subsection{The finite local calculation on the original ramified forms}

% BEGIN GLE7
\paragraph{GLE7: diamond matchings and the exact local move dictionary.}
\label{claim:GLE7}
Fix a selected place $v$, write $\mathbb F=\mathbb F_p$ for its
original residue field, and localize the original lattices at $P_v$.
Write $H$ for the dual lattice, $L$ for the primal lattice, and
$f:H/pH\to\mathbb F$ for the original nonzero functional.  The
source's \texttt{fv\_\allowbreak ker} identifies $L/pH$ with $\ker f$, so
$pH\subset L\subset H$.  Put $V_0=L/pL$, of dimension $k$, and
choose $z_0\in H$ with $f(z_0)=1$.  Its original distinguished vector is
\[
                         t_e=pz_0\pmod {pL}\ne0.                 \tag{GLE.30}
\]
This is independent of $z_0$: two such lifts differ by an element of
$L$.  Nonvanishing follows because $pz_0\in pL$ would put $z_0\in L$.
The prescribed residue projections identify this $V_0$ and $t_e$
with \texttt{lk} and \texttt{te}, without a choice of a new residue
space.

The original local moves have the following precise description.
An up index is $z\in H/pH$ with $f(z)=1$.  After its original rational
transport the new dual lattice is $H+p^{-1}\mathcal O z$, the new
primal lattice is $H$, and the new functional is
\[
                         f^+(y)=f(py).
\]
An arbitrary down index is a nonzero functional $g\in V_0^*$.
After its original transport the new dual lattice is $L$, its new
primal lattice is the inverse image of $\ker g$ in $L$, and the new
functional is $g$.  The ordinary down operator restricts these
indices by $g(t_e)=1$; the auxiliary operator $E_v$ uses all
$g\ne0$.  These are the local equalities of NB2--NB3 and NB7--NB8.
At the conjugate prime the data are fixed by Hermitian duality;
at every other prime they are unchanged.

The original diamond indexed by $a\in\mathbb F^\times$ replaces
$f$ by $af$ and leaves $H,\allowbreak{} L$ unchanged.  It replaces $t_e$ by
$a^{-1}t_e$.  In an up sum the bijection $z\mapsto a^{-1}z$
preserves the enlarged lattice and multiplies the resulting
functional by $a$.  In a down sum the bijection $g\mapsto ag$
preserves the condition $g(t_e)=1$ after this change of $t_e$, and
again multiplies the resulting functional by $a$.
These give, on the original forms,
\[
 \rho_{v,\allowbreak{} a}U_v=U_v\rho_{v,\allowbreak{} a},\qquad
 \rho_{v,\allowbreak{} a}D_v=D_v\rho_{v,\allowbreak{} a},\qquad E_v\rho_{v,\allowbreak{} a}=E_v.          \tag{GLE.31}
\]
For the last equality, all nonzero functionals occur, and multiplication
by $a$ simply permutes them.  At two distinct selected places the
local moves have disjoint support.  Their finite-adelic matrices
commute componentwise, which proves the corresponding commutations
of $U_v,\allowbreak{} D_v$ with every $\rho_{v',\allowbreak{} a}$ for $v'\ne v$.

We check why these identifications are equalities of transported
values, rather than only equalities of lattice shapes.  Work at the
original preferred adelic representative and use moves supported at
the place in question.  If their final dual lattice, primal lattice
and functional agree, the ratio of the two move matrices stabilizes
these data.  The original \texttt{Uloc\allowbreak Alt} and \texttt{U\_\allowbreak eq}
identify this stabilizer with the original level subgroup.  Its
components away from the moved place are the identity.  The selected
place is disjoint from the two auxiliary bad primes: RSL8 proves this
from $\pi_u\equiv1$ modulo $(g_0g_1)^2$ and
$\pi_u\bar\pi_u=p$.  The original \texttt{chi\_\allowbreak triv} therefore
makes the level character of the ratio equal to one.  Adelic
equivariance gives equal values, with no residual scalar.
Factoring the moves through the prescribed class representatives
introduces the same natural transforms as in (GLE.26), so it does
not alter the equality.  This also proves preservation of the
original forms space by $U_v,\allowbreak{} D_v,\allowbreak{} E_v$ and the diamonds: changing a
level representative permutes the displayed finite indices and
uses this same equivariance.  No transitivity or approximation
assertion for a different class space is being invoked.
% END GLE7

% BEGIN GLE8
\paragraph{GLE8: the actual $DU$ formula by finite fibers.}
\label{claim:GLE8}
Continue with exactly the preceding local data.  First make a down
move indexed by $g\ne0$, then an up move from its new dual lattice
$L$.  The original down transport identifies the latter residue
space with $V_0=L/pL$.  Thus its up indices are in bijection with
\[
                         \{y\in V_0:g(y)=1\}.                  \tag{GLE.32}
\]
Both directions use the original residue projections: a lift of $y$
in $L$ pulls back to a lift in the down dual lattice, and two lifts
give the same index exactly when their difference lies in $pL$.
The projection is surjective and has this kernel by \texttt{prk}
and \texttt{prk\_\allowbreak ker}.  This verifies the whole bijection, including
its uniqueness.

For a nonzero $y\in V_0$ choose a lift $\widetilde y\in L$.
After the two transports the final dual and primal lattices and
functional are
\[
 H_y=L+p^{-1}\mathcal O\widetilde y,\qquad L_y=L,
 \qquad h_y(l+c p^{-1}\widetilde y)=\bar c.                    \tag{GLE.33}
\]
Here $\bar c$ denotes reduction to $\mathbb F$.
Indeed the up functional is $g(p\,\cdot)$, and $g(y)=1$;
its kernel in $H_y/pH_y$ is the image of $L$.
The expression is well defined: if $c p^{-1}\widetilde y\in L$,
then $\bar c\,\allowbreak{} y=0$, hence $\bar c=0$.
Changing the lift by $pL$ does not change any of these data.
In particular the endpoint depends only on $y$, not on $g$.
The conjugate local component is determined by its Hermitian dual,
and all other components are unchanged.  The character-one
comparison in GLE7 therefore proves equality of the actual
transported form values.  Denote that common value, at a fixed class
and frame, by $W(y)$.

For $y=t_e$, (GLE.33) gives the original $H,\allowbreak{} L,\allowbreak{} f$, so $W(t_e)=\phi$.
For $y=a t_e$ with $a\ne0$ it gives $H,\allowbreak{} L,\allowbreak{} a^{-1}f$.  Thus
$W(a t_e)=\rho_{v,\allowbreak{} a^{-1}}\phi$ at that class and frame.  In
particular
\[
             \sum_{a\in\mathbb F^\times} W(a t_e)
                       =\operatorname{Nm}_v\phi,\qquad
             \operatorname{Nm}_v=\sum_a\rho_{v,\allowbreak{} a}.             \tag{GLE.34}
\]
The inversion in this formula records the literal source diamond
convention.

It remains to count fibers, a calculation in the $k$-dimensional
space $V_0$.  If $y\ne0$, choosing a basis beginning with $y$
shows that exactly $p^{k-1}$ functionals satisfy $g(y)=1$;
none of them is zero.  For $y=0$ there are none.
Requiring in addition $g(t_e)=1$ gives
\[
 \#\{g:g(t_e)=g(y)=1\}=
 \begin{cases}
 p^{k-1},&y=t_e,\\
 0,&y\in\mathbb F t_e, y\ne t_e,\\
 p^{k-2},&y\notin\mathbb F t_e.
 \end{cases}                                                  \tag{GLE.35}
\]
The middle case follows by linearity.  In the last case extend
$t_e,\allowbreak{} y$ to a basis; two coordinates of $g$ are fixed and the other
$k-2$ are arbitrary.  The original ranks have $k\ge5$, so the
exponents are ordinary nonnegative integers.

Expanding the exact finite sums, reindexing by (GLE.32), and using
(GLE.35) now gives
\[
 \begin{split}
 D_vU_v\phi&=p^{k-1}\phi+
                   p^{k-2}\sum_{y\notin\mathbb F t_e}W(y),\\
 E_vU_v\phi&=p^{k-1}\sum_{y\ne0}W(y).
 \end{split}
\]
Subtract the line contribution (GLE.34).  With the original
normalization $\operatorname{trp}_v=p^{-(k-1)}E_vU_v$, the result is
\[
 D_vU_v\phi=p^{k-1}\phi-p^{k-2}\operatorname{Nm}_v\phi
                              +p^{k-2}\operatorname{trp}_v\phi. \tag{GLE.36}
\]
Every step holds at each original class and frame, so this is an
identity in the original forms space.  It is the complete literal
\texttt{DUFormula\allowbreak Level\allowbreak Zero}; its proof has supplied the move,
fiber, endpoint and scalar identifications needed for that formula.
% END GLE8

% BEGIN GLE9
\paragraph{GLE9: the inverse composite on the original ramified space.}
\label{claim:GLE9}
Let
\[
 M^{\rm ram}=\mathcal F.\mathrm{Fi}(0)\cap
                  \bigcap_{v'}\ker\operatorname{Nm}_{v'}
                    =\texttt{Mram\allowbreak I D u 0}.
\]
The commutations (GLE.31), including those at distinct places,
show that $U_v$ and $D_v$ preserve every displayed kernel and the
forms space.  Moreover (GLE.31) gives
\[
 \operatorname{trp}_v\rho_{v,\allowbreak{} a}
       =p^{-(k-1)}E_vU_v\rho_{v,\allowbreak{} a}
       =p^{-(k-1)}E_v\rho_{v,\allowbreak{} a}U_v
       =\operatorname{trp}_v.
\]
For $\phi\in M^{\rm ram}$ this implies
\[
 0=\operatorname{trp}_v(\operatorname{Nm}_v\phi)
                 =(p-1)\operatorname{trp}_v\phi.
\]
Since $p$ is a prime, $p-1\ne0$ in $\mathbb C$, and hence
$\operatorname{trp}_v\phi=0$.  Equation (GLE.36) reduces to
$D_vU_v\phi=p^{k-1}\phi$.

The original forms space is finite dimensional by its original
\texttt{emb\_\allowbreak basis}, so its subspace $M^{\rm ram}$ is too.
The last identity makes $U_v$ injective there.  A linear injective
endomorphism of a finite-dimensional space is surjective, and
therefore $D_v=p^{k-1}U_v^{-1}$ on this same space.  Consequently
\[
 U_v(M^{\rm ram})\subseteq M^{\rm ram},\qquad
 D_v(M^{\rm ram})\subseteq M^{\rm ram},\qquad
 U_vD_v\phi=p^{k-1}\phi\quad(\phi\in M^{\rm ram}).             \tag{GLE.37}
\]
This is exactly \texttt{Hecke\allowbreak Ramified\allowbreak Level\allowbreak Zero}, with $k=N(0)$.
No change of the ramified subspace, of the finite character, or of
the preferred representative system has occurred.  The proof uses
the original \texttt{Forms0} ties; neither a new abstract Hecke
module nor an additional matching predicate is required.
% END GLE9

\subsection{The original congruence lattice and its length}

% BEGIN GLE10
\paragraph{GLE10: a DVR estimate for one normalized original operator.}
\label{claim:GLE10}
Let $\mathcal O=\mathcal O_w$, let $e=e_w$ be the valuation of $p$,
and write $S=S_w$, $B=B_{w,\allowbreak{} v}$, $W=\mathcal F.\mathrm{ramSub}$.
The rationality in GLE5 descends the two original operators uniquely
to $K$-linear operators on $V$.  Their descent is unique because
the original embedding is injective.  GLE9 and the original
\texttt{ram\allowbreak Sub} equality make both descend to $W$, with
$UD=p^{k-1}I$ there.  Put
\[
 R=S\cap W,\qquad B_R=B\cap W,\qquad b=b_v\ge0,
 \qquad G=p^bU.
\]
Thus $R$ is the original \texttt{Sram}, not a substitute lattice.
The equalities $p^bR=(p^bS)\cap W$ and
$U(p^bR)=p^bU(R)$ follow because $W$ is a $K$-subspace.
GLE5--GLE6 and GLE9 give
\[
 p^bR\subseteq B_R,\qquad U(B_R)\subseteq R,\qquad
 D(R)\subseteq B_R,\qquad
 p^{k-1}R\subseteq U(B_R).
\]
In particular $G(R)\subseteq R$.

All these are finitely generated torsion-free modules over a DVR,
hence free.  Multiplication by $p$ has cokernel length $e$ on each
free generator.  The last inclusion gives
\[
 \ell(R/U(B_R))\le e(k-1)\operatorname{rank}_{\mathcal O}R
                                       \le e(k-1)d_S.
\]
The natural injection
$B_R/p^bR\hookrightarrow B/p^bS$ shows
$\ell(B_R/p^bR)\le\ell(B/p^bS)$.
Applying $U$ gives a surjection
$B_R/p^bR\twoheadrightarrow U(B_R)/U(p^bR)$, so the same upper
bound holds for the latter quotient.  The exact sequence of the
three nested modules
$U(p^bR)\subseteq U(B_R)\subseteq R$ therefore yields
\[
 \ell(R/G(R))\le
      \ell(B/p^bS)+e(k-1)d_S
   =e d_S\left(\frac{|\lambda_{0,\tau_v}|+kb_v}{k}+k-1\right).
                                                               \tag{GLE.38}
\]
The term on the left is finite; each step has just exhibited a
finite-length quotient that bounds it.  The rank inequality uses
only $\operatorname{rank}R\le\dim_K V=d_S$, established on the
original coordinates in GLE4.  This independently proves the
source's local \texttt{Lemma\allowbreak C\_\allowbreak core} estimate with the exact
original dimension and normalization.
% END GLE10

% BEGIN GLE11
\paragraph{GLE11: the ordered product, its power, and the complete length bound.}
\label{claim:GLE11}
Assume \texttt{Hida\allowbreak Setting D u} and fix $w\mid p$.
Let $\mathcal V_w$ be the original \texttt{VVT}: it consists of the
selected places whose $\tau_v$ is top if $w$ is seen, and is empty
otherwise.  On this set $b_v\ge0$ by the last clause of
\texttt{Hida\allowbreak Setting}.  For seen $w$ let $G_v=p^{b_v}U_v$ on the
same $R=\texttt{Sram D u w}$.  GLE10 applies to every $v$.

We spell out the finite-index argument to avoid any commutativity
assumption about the chosen order.  If $A(R),\allowbreak{} B(R)\subseteq R$, then
$A(R)/AB(R)$ is a quotient of $R/B(R)$ under $A$.  The exact
sequence with submodule $A(R)/AB(R)$ gives
\[
 \ell(R/AB(R))\le\ell(R/A(R))+\ell(R/B(R)).                    \tag{GLE.39}
\]
Induct over the original \texttt{to\allowbreak List.\allowbreak foldr} order and then over
the $m$ powers.  The resulting operator $G_*$ satisfies
\[
 G_*^m(R)\subseteq R,\qquad
 \ell(R/G_*^m(R))\le
 m e_wd_S\sum_{v\in\mathcal V_w}
       \left(\frac{|\lambda_{0,\tau_v}|+kb_v}{k}+k-1\right).
                                                               \tag{GLE.40}
\]
Scalars commute with all linear maps, so, in that very same order,
\[
 G_*^m=p^{m\sum_{v\in\mathcal V_w}b_v}
       \left(\prod_{v\in\mathcal V_w}^{\rm original\ order}U_v\right)^m.
\]
The exponent and product are literally \texttt{exp\allowbreak T} and
\texttt{Uprod}.  The source's definition of \texttt{img\allowbreak L}, together
with the unique rational descent in GLE10, therefore identifies
\[
              G_*^m(R)=\texttt{Nlam D u w}.                    \tag{GLE.41}
\]
This is equality of submodules of the original rational $V$.

For completeness, the weight simplification in the stated source
bound is also elementary.  At a top $\tau$, the first entry is
at most $a_{\max}$, and the absolute values of every later entry
are at most the source's $\operatorname{Wt}(\tau)$.  The last
entry has absolute value $b_v\le\operatorname{Wt}(\tau)$.
Thus
\[
 \frac{|\lambda_{0,\tau}|+kb_v}{k}
 \le \frac{a_{\max}}k+
       \frac{k-1}{k}\operatorname{Wt}(\tau)+b_v
 \le \frac{a_{\max}}k+2\operatorname{Wt}(\tau).
\]
Here $|\lambda|$ continues to denote the signed sum of its entries,
not their absolute-value sum.  Using $k-1\le k$ and $b_v\ge0$ in
(GLE.40) proves the exact, slightly looser source bound
\[
 \begin{split}
 \texttt{Nlam D u w}&\subseteq\texttt{Sram D u w},\\
 \ell(\texttt{Sram}/\texttt{Nlam})
 &\le e_w m d_S
   \sum_{v\in\mathcal V_w}
     \left(\frac{a_{\max}}k+2\operatorname{Wt}(\tau_v)+k+b_v\right).
 \end{split}                                                  \tag{GLE.42}
\]
The length is a natural number, as proved in (GLE.40).
If $w$ is unseen, the original list is empty, its product is the
identity, and its exponent is zero.  Hence \texttt{Nlam} equals
\texttt{Sram} and both sides of the required length inequality are
zero.  This handles that case without presupposing a seen place.
Equations (GLE.41)--(GLE.42) are the complete literal
\texttt{Congruence\allowbreak Lattice\allowbreak Length\allowbreak Bound D u}.
% END GLE11

% BEGIN GLE12
\paragraph{GLE12: uniform installation and the remaining common producer.}
\label{claim:GLE12}
The logical roles of the preceding statements are distinct.
GLE1--GLE2 prove
\[
             \texttt{Tu\allowbreak Size}(\mathcal D)
                 \Longrightarrow\texttt{Setting\allowbreak Eventually}(\mathcal D)
                                                               \tag{GLE.43}
\]
for the actual family with its original parameter witnesses.
GLE5, GLE9 and GLE6 respectively prove the literal pointwise
\texttt{Hecke\allowbreak Integral\allowbreak Level\allowbreak Zero}, \texttt{Hecke\allowbreak Ramified\allowbreak Level\allowbreak Zero},
and \texttt{Second\allowbreak Branch\allowbreak Inclusion\allowbreak Length} from the same original
\texttt{Forms0}, its finite $K$-valued level character and
\texttt{adm\allowbreak S}.  The second-branch statement retains its explicit
$b_v\ge0$ hypothesis.  GLE11 consumes these three statements and
\texttt{Hida\allowbreak Setting}, on the same $D,\allowbreak{} u$, to produce the complete
congruence lattice length bound.

In the selected application the original
\texttt{Real\allowbreak Place\allowbreak Second\allowbreak Layer\allowbreak Conditions} supplies, eventually, the
same \texttt{Forms\allowbreak Tied\allowbreak Sym}, its \texttt{Chi\allowbreak Fin} data, and the
required real-place ties.  RSL6 fixes the common original selector;
RSL7 supplies \texttt{adm\allowbreak S} on those forms.  Take the maximum of
that supplier threshold and the threshold in GLE2.  It is chosen
before $i$, before the word $u$, and before the places $w,\allowbreak{} v$.
The pointwise proofs above are uniform in all of these variables.
Thus the original guarded family predicate is satisfied for each
of the three operator conclusions:
\[
 \exists A\ \forall x\ge A\ \forall i\ \forall u,
 \quad \mathcal D.x_1<x\ \Longrightarrow
 \bigl(\texttt{Hida\allowbreak Setting}(D_{x,\allowbreak{} i},\allowbreak{} u)\Longrightarrow
          P(D_{x,\allowbreak{} i},\allowbreak{} u)\bigr).                                  \tag{GLE.44}
\]
This is the literal \texttt{Fam\allowbreak U} guard and quantifier order.
One may then discharge its guard using (GLE.43), on the same late
family.  Maxima of finitely many thresholds do not change the
family, its class constants, any parameter witness, or the
selector.  Accordingly this also transports to the source's
\texttt{Closed} predicates whenever those upstream predicates are
supplied on a class.  It does not choose a fresh family separately
for the three conclusions.

The five implication tasks
\texttt{TASK\_\allowbreak GL6\_\allowbreak SIZE\_\allowbreak SETTING},
\texttt{TASK\_\allowbreak GL6\_\allowbreak INTEGRAL},
\texttt{TASK\_\allowbreak GL6\_\allowbreak RAMIFIED},
\texttt{TASK\_\allowbreak GL6\_\allowbreak SECOND\_\allowbreak BRANCH}, and
\texttt{TASK\_\allowbreak GL6\_\allowbreak CONGRUENCE\_\allowbreak LENGTH} are thereby closed in this
stated scope.  The first now composes with the original construction
class's \texttt{Tu\allowbreak Size} producer; the other four retain the
original same-selector forms and finite-character input as explicit
upstream columns.  The compatible common selector remains a
producer obligation, and this chapter does not assert the GL4
$hZ_a,\allowbreak{} hZ_b,\allowbreak{} t_{CV}$ inputs or the entire eleven-field
\texttt{Fam\allowbreak GLmin} output.  Once their separate inputs are supplied,
the present conclusions are available on exactly the original
families required for that final assembly.  RSL11--RSL12 may also
reuse the inclusion in (GLE.42) to obtain their already proved
fullness conclusion without a separate congruence-lattice inclusion
producer.
% END GLE12




\section{The actual finite isotropic averaging operator}
\label{sec:finite-isotropic-average-v033}

\status{FIA1--FIA8 reconstruct the finite cancellation needed by the
original definite theta lift.  Under the five literal hypotheses in
the source's averaging theorem, they produce its vanishing sum on
every original polynomial kernel and every original adelic translate.
No additional character, finite model, or averaging hypothesis is
retained.  The hypotheses concern the supplied original real datum;
their existence and the later smooth cuspidal extension are separate
producer obligations.  These are category B proofs, with no new
published-theorem application and no new Lean replay.}

Fix the original host \(D\), with totally real field \(F\), CM field
\(E\), conjugation \(z\mapsto\bar z\), and nonzero purely imaginary
\(\xi_0=D.\mathrm{host}.\mathrm{xi0}\).  Write
\(\mathbb A_E=\mathbb A_F\otimes_F E\), with the source's
identification of its finite part with
\(E\otimes_F\mathbb A_{F,\allowbreak{} f}\).  Its integral finite ring is
\(\widehat{\mathcal O}_E\).  All inclusions, finite projections,
and residue maps below are the original
\texttt{inc\allowbreak A}, \texttt{to\allowbreak EAf}, and \texttt{res\allowbreak S}.

Let \(S,\allowbreak{} T\) be the original Hermitian coordinate spaces, of ranks
\(n,\allowbreak{} m\), with Gram matrices \(G_S,\allowbreak{} G_T\).  The oriented space is
\[
 X=\operatorname{Mat}_{m\times n}(\mathbb A_E),\qquad
 t^\dagger=G_S^{-1}\bar t^{\,\allowbreak{} t}G_T,\qquad
 B(x,\allowbreak{} y)=\operatorname{Tr}_{\mathbb A_E/\mathbb A_F}
                  \bigl(\xi_0\operatorname{tr}(y^\dagger x)\bigr).
 \tag{FIA.1}
\]
Thus \(B=\texttt{sympl\allowbreak O}\), whereas the source's raw Heisenberg
operator uses its half:
\[
 b=B/2,\qquad H_y f(x)=\psi(b(x,\allowbreak{} y))f(x+y).
 \tag{FIA.2}
\]
The distinction is essential: the opposite-half modulation law in
\texttt{Is\allowbreak Place\allowbreak Model} uses \(\psi(B(t,\allowbreak{} y))\), not
\(\psi(b(t,\allowbreak{} y))\).  We preserve both normalizations throughout.
For actual finite frames \(\beta_S,\beta_T\), put
\[
 {\cal F}(t)=\beta_S^{-1}(t^\dagger)_f\beta_T
       =\texttt{fr\allowbreak Map\allowbreak Adj}(S,\allowbreak{} T,\beta_S,\beta_T,\allowbreak{} t).
 \tag{FIA.3}
\]
The subscript \(f\) includes the source's tensor-factor interchange.
Rows and columns in all matrix identities below have these orders.

% BEGIN FIA1
\paragraph{FIA1: every local character is nontrivial and has an edge.}
\label{claim:FIA1}
Suppose \(\psi\) satisfies the literal \texttt{Is\allowbreak Char\allowbreak AF}:
it is a continuous additive character of \(\mathbb A_F\), has
absolute value one, is one on the diagonal \(F\), and is not
identically one.  For every finite place \(v\), its restriction
\(\psi_v\) to the single \(F_v\)-component is nontrivial.
Moreover there is \(x_v\in F_v\) such that
\[
 \psi_v(x_v)\ne1,\qquad
 |y|_v<|x_v|_v\ \Longrightarrow\ \psi_v(y)=1.
 \tag{FIA.4}
\]
Here \(|\cdot|_v\) denotes any multiplicative normalization of
the discrete valuation; only its ordering is used.

\emph{Proof of local nontriviality.}
We give the character argument because a nontrivial character of
an arbitrary product need not be nontrivial on every factor.
On \(F_\infty\simeq\mathbb R^{[F:\mathbb Q]}\), a continuous unitary
additive character has the form
\(\exp(2\pi iL(x))\), with \(L\) real linear.
For completeness, on a real line take the continuous argument in
a small arc around \(1\).  The additive law makes its argument
additive on a smaller interval, since the integer ambiguity is
continuous and zero at the origin.  Continuity and rational
approximation make it linear there; subdivision gives the result
on the whole line.  Multiplying the coordinate restrictions proves
the assertion in finitely many dimensions.

The finite restriction of \(\psi\) kills \(N\widehat{\mathcal O}_F\)
for some positive integer \(N\).  Indeed, continuity gives a
compact open additive subgroup whose image lies in
\(\{z\in\mathbb C:|z|=1,\ \operatorname{Re}z>0\}\).
A subgroup of the circle contained in this arc is trivial:
every nonidentity point has some integral power outside the arc.
A sufficiently divisible positive integer multiple of
\(\widehat{\mathcal O}_F\) lies in this open subgroup, by the
restricted-product topology.  Rational triviality now gives
\[
            L(Nr)\in\mathbb Z\qquad(r\in\mathcal O_F).
 \tag{FIA.5}
\]
An integral basis therefore makes \(L|_F\) rational valued.
The nondegenerate rational trace pairing produces \(a\in F\)
with \(L(r)=\operatorname{Tr}_{F/\mathbb Q}(ar)\) for \(r\in F\).
One can check nondegeneracy directly: the matrix of embeddings
of a rational basis is invertible for a separable number field,
and the trace matrix is its transpose times itself.  Since an
integral basis is a real basis of \(F_\infty\), the same trace
formula holds on \(F_\infty\).

Here is a concrete standard global character.  On \(\mathbb A_\mathbb Q\)
put
\[
 \Psi_\mathbb Q((u_\infty,\allowbreak{} u_p)_p)
   =\exp\!\left(2\pi i
             \left(u_\infty-\sum_p\{u_p\}_p\right)\right),
 \qquad
 \Psi_F=\Psi_\mathbb Q\circ
                  \operatorname{Tr}_{\mathbb A_F/\mathbb A_\mathbb Q}.
 \tag{FIA.6}
\]
The \(p\)-adic fractional part is represented in
\(\mathbb Z[1/p]\cap[0,1)\); its class modulo \(\mathbb Z\) is
additive, and only finitely many terms in the sum are nonzero.
For a rational number \(q\), the difference
\(q-\sum_p\{q\}_p\) is integral.  To see this, its denominator
has no prime divisor: at every \(p\), the \(p\)-fractional part
removes the entire negative-valuation part and the other
fractional parts are \(p\)-integral.  This proves rational
triviality, continuity, and the additive law of (FIA.6).
Applying the field trace proves the corresponding assertions
for \(\Psi_F\).

The characters \(\psi\) and \(z\mapsto\Psi_F(az)\) agree at
infinity and on \(F\).  Their quotient is consequently a
continuous character of \(\mathbb A_{F,\allowbreak{} f}\) trivial on \(F\).
The diagonal \(F\) is dense in \(\mathbb A_{F,\allowbreak{} f}\).  An explicit
approximation proves this last assertion.  Given finitely many
local targets and congruence precisions, choose a positive rational
integer \(M\) that clears their denominators.  At every prime
dividing \(M\) not in the target set, impose divisibility by \(M\);
at target primes prescribe the desired residues of \(M\) times
the target, to the requested precision after division by \(M\).
The Chinese remainder theorem in \(\mathcal O_F\) gives a
numerator \(r\).  Then \(r/M\) has all required approximations
and is integral at every other place.  Thus the quotient character
is one, and
\[
                    \psi(z)=\Psi_F(az).
 \tag{FIA.7}
\]
The original nontriviality implies \(a\ne0\).
For \(v\mid p\), the local restriction is the nontrivial standard
character of \(\mathbb Q_p\) applied to
\(\operatorname{Tr}_{F_v/\mathbb Q_p}(a\,\allowbreak{} z)\).
This trace map is onto: the trace of \(1\) is the positive degree,
which is invertible in \(\mathbb Q_p\).  Multiplication by
the nonzero image of \(a\) is onto as well.  Hence \(\psi_v\)
is nontrivial.  This establishes the precise local assertion
needed here without invoking the stronger conductor-product
formula.

\emph{Proof of the edge.}
The same small-arc argument makes the kernel of \(\psi_v\)
contain \(\varpi^d\mathcal O_v\) for some integer \(d\).
The integers \(r\) for which \(\psi_v\) is one on
\(\varpi^r\mathcal O_v\) form a nonempty upper ray, with a lower
bound because \(\psi_v\) is not identically one.  Let \(d_0\)
be its least element.  Choose
\(x_v\in\varpi^{d_0-1}\mathcal O_v\) with
\(\psi_v(x_v)\ne1\).  Its valuation is exactly \(d_0-1\).
Any \(y\) with \(|y|_v<|x_v|_v\) has valuation at least \(d_0\),
so (FIA.4) follows.  In particular \(x_v\ne0\).
This proves the original \texttt{Theta\allowbreak Support.\allowbreak exists\_\allowbreak edge}.
% END FIA1

% BEGIN FIA2
\paragraph{FIA2: split charts and residue precision.}
\label{claim:FIA2}
Let \(e\) satisfy the original
\(\texttt{Is\allowbreak Idem\allowbreak At}(e,\mathfrak P)\).  There is a finite place
\(v\) of \(F\) and a surjective ring homomorphism
\(\pi:\mathbb A_E\to F_v\) such that
\[
 \begin{gathered}
 \pi(z)=0\ \Longleftrightarrow\ ez=0,\qquad
 \pi(e)=1,\qquad e_\infty=0,\\
 \pi(\iota_F(y))=y_v,\qquad
 \operatorname{Tr}_{\mathbb A_E/\mathbb A_F}(ez)
       =\operatorname{single}_v(\pi(z)).
 \end{gathered}                                                   \tag{FIA.8}
\]
The single component on the last line has zero infinite part.
The chart factors through the original finite projection as
\(\pi=\phi\circ\texttt{to\allowbreak EAf}\), with
\(\phi:E\otimes_F\mathbb A_{F,\allowbreak{} f}\to F_v\).
For an integral finite element \(u\),
\[
 |\phi(u)|_v\le1,\qquad
 \operatorname{res}_{\mathfrak P}(u)=0
       \ \Longrightarrow\ |\phi(u)|_v<1.
 \tag{FIA.9}
\]

\emph{Proof.}
The literal split-idempotent conditions say that
\(e^2=e\), \(e\bar e=0\), and \(e+\bar e\) is the identity at
one finite \(F\)-place and zero elsewhere.  Multiplying the last
identity by \(e\) shows that \(e\) is zero at every other place,
including infinity.  At the indicated place the quadratic
algebra \(E\otimes_F F_v\) has the two nonzero orthogonal
idempotents \(e,\bar e\).  Its two factors have dimensions adding
to two, so each is one dimensional over \(F_v\).
Consequently this algebra is \(F_v\times F_v\), with
\(e=(1,0)\) after choosing the first factor.  Projection to this
factor is the source's chart \(\pi\); a single finite scalar
times \(e\) gives every value, proving surjectivity.
Multiplication, the kernel assertion, and the trace identity
are now coordinate identities.  In particular the source
\texttt{Split\allowbreak Chart.\allowbreak comp} is this projection: its defining
equation \(ez=\operatorname{single}_v(\pi(z))e\) determines it
uniquely.  It annihilates the infinite component and therefore
factors through the indicated finite tensor ring.

The integral ring of the split quadratic algebra is
\(\mathcal O_v\times\mathcal O_v\).  This is also the image of
the original completed integral ring: completion of the finite
\(\mathcal O_F\)-module \(\mathcal O_E\), and then projection to
the two primes over \(v\), gives precisely these two local
integer rings.  Equivalently, integral elements are characterized
by their monic equations over \(\mathcal O_v\) in each factor.
It follows that \(\phi\) sends integral elements into
\(\mathcal O_v\).

For a nonzero \(r\in\mathcal O_E\), the second clause of
\texttt{Is\allowbreak Idem\allowbreak At} says
\[
 r\in\mathfrak P
 \ \Longleftrightarrow\
       (er^{-1})_f\notin\widehat{\mathcal O}_E
 \ \Longleftrightarrow\ |\pi(r^{-1})|_v>1
 \ \Longleftrightarrow\ |\pi(r)|_v<1.
 \tag{FIA.10}
\]
The same conclusion is true for \(r=0\).
Thus the selected local prime is exactly \(\mathfrak P\),
not its conjugate.  The kernel of the original residue map is
\(\mathfrak P\widehat{\mathcal O}_E\): the source defines it by
the quotient by this extended ideal, identifying that quotient
with \(\mathcal O_E/\mathfrak P\).
This identification follows directly by completion and the
Chinese remainder theorem.  At every other prime
\(\mathfrak P\) becomes the unit ideal, while at \(\mathfrak P\)
completion preserves the residue field.
Every element of the extended ideal is a finite sum of
integral multiples of elements of \(\mathfrak P\).
Equation (FIA.10), integrality of \(\phi\), and the
nonarchimedean triangle inequality prove (FIA.9).
In particular
\[
       \psi\!\left(\operatorname{Tr}(ez)\right)=\psi_v(\pi(z)).
 \tag{FIA.11}
\]
No rescaling of the character, trace, or residue homomorphism
has occurred.
% END FIA2

% BEGIN FIA3
\paragraph{FIA3: the constant phase on the original top block.}
\label{claim:FIA3}
Assume \(m\le n\), \(\texttt{Is\allowbreak Char\allowbreak AF}(\psi)\), and
\(\texttt{Is\allowbreak Idem\allowbreak At}(e,\mathfrak P)\).
Let \(c\in\mathbb A_E^m\) with \(ec_j\ne0\) for some \(j\).
The frames \(\beta_S,\beta_T\) are the actual invertible finite
matrices; no change of them is made.  Put
\[
                         q_{\mathfrak P}
                            =\#(\mathcal O_E/\mathfrak P)>0.
 \tag{FIA.12}
\]
There exist \(l\in\mathbb A_E^n\) and \(\zeta\in\mathbb C\)
such that
\[
 el=l,\qquad \zeta\ne1,\qquad \zeta^{q_{\mathfrak P}}=1,
 \tag{FIA.13}
\]
and, for every \(t\in X\),
\[
 \begin{split}
 {\cal F}(t)\in\operatorname{Mat}_{n\times m}
                      (\widehat{\mathcal O}_E),\qquad
 \operatorname{top}_m(\operatorname{res}_{\mathfrak P}{\cal F}(t))=I_m\\
 \Longrightarrow\quad
            \psi\bigl(B(t,\allowbreak{} c\,\allowbreak{} l^t)\bigr)=\zeta .
 \end{split}                                                       \tag{FIA.14}
\]
This proves \texttt{Theta\allowbreak Support.\allowbreak exists\_\allowbreak const}, with its positive
period specified explicitly.

\emph{Proof.}
Take the chart \(\pi,\phi\) of FIA2 and the edge \(x=x_v\)
of FIA1.  Define
\[
 d=\phi(\beta_T^{-1})\,\pi(c),\qquad
 u=\pi(\xi_0).
 \tag{FIA.15}
\]
The vector \(\pi(c)\) is nonzero, and the finite matrix
\(\phi(\beta_T)\) is invertible, so \(d\ne0\).
Choose \(j<m\) with \(|d_j|_v\) maximal.  It is nonzero.
The composite \(E\to\mathbb A_E\to F_v\) is a unital
field homomorphism and hence injective, so \(u\ne0\).
Set
\[
 s=-\frac{x}{u\,\allowbreak{} d_j},\qquad
 \pi(l_{0,\allowbreak{} a})=s\,\phi((\beta_S^{-1})_{ja}),\qquad
 l_a=e\,\allowbreak{} l_{0,\allowbreak{} a},\qquad \zeta=\psi_v(x).
 \tag{FIA.16}
\]
Surjectivity of \(\pi\) supplies every \(l_{0,\allowbreak{} a}\);
the row \(j\) is valid in the \(n\)-row frame because \(m\le n\).
These choices depend on \(c\), the selected place, character,
and frames, and not on \(t\).

Let \(U={\cal F}(t)\) satisfy the two conditions in (FIA.14).
Its \(j\)-th row, after applying \(\phi\), satisfies
\[
          |\phi(U_{ja})-\delta_{ja}|_v<1
                             \qquad(0\le a<m),                  \tag{FIA.17}
\]
by (FIA.9).  Put \(z=\sum_{a<m}\phi(U_{ja})d_a\).
Every term of
\[
                  z-d_j=\sum_{a<m}
                         (\phi(U_{ja})-\delta_{ja})d_a
 \tag{FIA.18}
\]
has absolute value strictly smaller than \(|d_j|_v\);
the finite ultrametric triangle inequality therefore gives
\(|z-d_j|_v<|d_j|_v\).
The actual matrix identity
\[
       \phi(U)=\phi(\beta_S^{-1})\,\pi(t^\dagger)\,
                    \phi(\beta_T)
 \tag{FIA.19}
\]
and (FIA.15)--(FIA.16) yield
\(\pi(l_0^t t^\dagger c)=s z\).

We check the orientation and sign directly.
From (FIA.1),
\(\operatorname{tr}(t^\dagger c\,\allowbreak{} l_0^t)
=l_0^t t^\dagger c\).
Hermitian symmetry, \(\bar\xi_0=-\xi_0\), and invariance
of the quadratic trace under conjugation give
\[
 B(t,\allowbreak{} c(e l_0)^t)
   =\operatorname{Tr}
          \bigl(e(-\xi_0\,\allowbreak{} l_0^t t^\dagger c)\bigr).
 \tag{FIA.20}
\]
Applying (FIA.11), its character value is
\(\psi_v(-u\,\allowbreak{} s z)\).  The chosen negative sign in \(s\)
is exactly what makes
\[
             -u\,\allowbreak{} s z=x+x\frac{z-d_j}{d_j}.
 \tag{FIA.21}
\]
The error term has absolute value \(<|x|_v\), so FIA1
annihilates it.  Additivity proves (FIA.14).

Finally \(q_{\mathfrak P}\in\mathfrak P\): the additive group
of the finite residue field is killed by its cardinality.
Consequently \(|q_{\mathfrak P}|_v<1\) by (FIA.10), and
\[
 \zeta^{q_{\mathfrak P}}
   =\psi_v(q_{\mathfrak P}x)=1
 \tag{FIA.22}
\]
by the same edge inequality.  Nontriviality of \(\zeta\)
was built into the choice of \(x\).
Its order need only divide \(q_{\mathfrak P}\); no primitive-root
claim or new choice of residue character is required.
% END FIA3

% BEGIN FIA4
\paragraph{FIA4: the actual finite half-datum meets the model's guards.}
\label{claim:FIA4}
Let \(\Lambda_S,\Lambda_T\) be full
\(\mathcal O_E\)-lattices.  In the application they are the
original self-dual lattices.  Write
\[
 \Lambda=\texttt{lat\allowbreak O}(S,\allowbreak{} T,\Lambda_S,\Lambda_T,\mathrm{True}).
 \tag{FIA.23}
\]
It contains the actual fractional-ideal factor
\(\mathfrak A=D.\mathrm{host}.\mathrm{AA}\).  In the self-dual
case it is
\[
 \widehat{\mathfrak A}\,
 \operatorname{Hom}_{\widehat{\mathcal O}_E}
  (\widehat\Lambda_S,\widehat\Lambda_T).
\]
Replacing it by the unweighted Hom lattice would be incorrect.

Let \(P_g,\allowbreak{} Q\) be the finite sets in
\texttt{Is\allowbreak Th\allowbreak PVec\allowbreak Of}, and let \(e_\Sigma\) be its supplied
sum of split idempotents over \(P_g\cup Q\).
Set \(a=\bar e_\Sigma\).  The actual finite function is
\[
 \Phi(t)=
 \begin{cases}
 \Omega_K(P_g,\operatorname{res}{\cal F}(t))\,
 \Omega_Q(Q,\allowbreak{} m,\chi,-1,\operatorname{res}{\cal F}(t)),
       &at=t,\ {\cal F}(t)\text{ integral},\\
 0,&\text{otherwise}.
 \end{cases}                                                     \tag{FIA.24}
\]
Here \(\chi=\texttt{xi\allowbreak R}\) is the stored family of residue
characters.  The first factor is one exactly when, at each
\(\mathfrak P\in P_g\), the upper \(m\times m\) block is the
identity and all rows below it are zero; otherwise it is zero.
The second factor is zero unless its indicated lower rows
vanish and top determinants at \(Q\) are units, and otherwise
is the product of the stored inverse determinant characters.
Only its dependence on the residues at \(Q\) is used.

Then \(\Phi\) satisfies the exact \texttt{Is\allowbreak SBHalf} guard:
its support lies in \(aX\), and there is an integer \(K>0\)
such that
\[
 \begin{gathered}
 au=u,\ u\in\Lambda
       \quad\Longrightarrow\quad
       \Phi(t+Ku)=\Phi(t)\quad(t\in X),\\
 \Phi(t)\ne0\quad\Longrightarrow\quad Kt\in\Lambda .
 \end{gathered}                                                     \tag{FIA.25}
\]
\emph{Proof.}
We first justify uniform lattice denominators.
The source's lattice \(\Lambda\) is generated over the finite
integral \(F\)-adeles by the matrices
\[
           A\,\allowbreak{} v\,\bar w^{\,\allowbreak{} t}G_S,\qquad
 A\in\mathfrak A,\quad v\in\Lambda_T,\quad w\in\Lambda_S.
 \tag{FIA.26}
\]
The fractional ideal and both lattices are finitely generated
over their integer rings.  Their rational spans are \(E\)
and the two full coordinate spaces.  Since \(G_S\) is invertible,
these matrices span \(\operatorname{Mat}_{m\times n}(E)\).
Take finite \(\mathcal O_E\)-generating families in (FIA.26)
and multiply the resulting matrix generators by an integral
\(\mathbb Z\)-basis of \(\mathcal O_E\).
Their \(\mathbb Z\)-span is a full finitely generated
lattice.  Clearing the coefficients of its generators and of
expressions for the rational matrix units gives positive integer
inclusions in both directions between \(\Lambda\) and the
standard integral finite matrix lattice.
This also proves that every element of \(\Lambda\) has zero
infinite part.

On the finite matrix space, \({\cal F}\) is the composition of
conjugate transpose with left and right multiplication by
invertible finite-adelic matrices.  Its inverse is of the same
form.  Entries of those matrices and their inverses are integral
away from finitely many places; one positive rational integer
clears their finitely many remaining denominators.
Combined with the preceding lattice inclusions, this gives
positive integers \(M,\allowbreak{} J\) such that
\[
 M{\cal F}(\Lambda)\subset
                   \operatorname{Mat}(\widehat{\mathcal O}_E),
 \qquad
 t_\infty=0,\ {\cal F}(t)\text{ integral}
                       \ \Longrightarrow\ Jt\in\Lambda .
 \tag{FIA.27}
\]
This is the content of the source's
\texttt{fr\allowbreak Map\allowbreak Adj\_\allowbreak up/\allowbreak down} and lattice-denominator lemmas,
proved here for the actual fractional lattice.

Put \(q=\prod_{\mathfrak P\in P_g\cup Q}q_{\mathfrak P}\)
and \(K=JMq>0\); the empty product is \(1\).
For \(u\in\Lambda\), the increment
\({\cal F}(Ku)=Jq\,\allowbreak{} M{\cal F}(u)\) is integral and is zero
modulo every \(\mathfrak P\in P_g\cup Q\).
Thus integrality of \({\cal F}(t)\) is unchanged by this
increment, all relevant residue matrices are unchanged,
and the half condition is unchanged when \(au=u\).
Every branch of (FIA.24), including its zero branch,
is therefore invariant.  If \(\Phi(t)\ne0\), then \(at=t\).
Since \(a_\infty=0\), we have \(t_\infty=0\); the second
inclusion in (FIA.27) gives \(Jt\in\Lambda\), and hence
\(Kt\in\Lambda\).  This proves all guards in (FIA.25).
% END FIA4

% BEGIN FIA5
\paragraph{FIA5: a genuine eigencharacter of the stored finite vector.}
\label{claim:FIA5}
Suppose \(m\le n\), the original lattices are self-dual,
\(\mathfrak P\in P_g\), and
\[
 \texttt{Is\allowbreak Th\allowbreak PVec\allowbreak Of}
   (\psi,\theta,\allowbreak{} S,\allowbreak{} T,\allowbreak{} s_S,\allowbreak{} s_T,\Lambda_S,\Lambda_T,
         \beta_S,\beta_T,\allowbreak{} P_g,\allowbreak{} Q,\chi,\allowbreak{} f_0,\allowbreak{} f)
 \tag{FIA.28}
\]
holds for the actual vector \(f\).
Assume \(c\in\mathbb A_E^m\) has nonzero projection at every
split idempotent, exactly as in \texttt{exists\_\allowbreak eigen}.
There exist \(l,\zeta\) with \(\zeta\ne1\),
\(\zeta^{q_{\mathfrak P}}=1\), such that, for all integers
\(k\ge0\),
\[
        (c(kl)^t)_\infty=0,\qquad
        H_{c(kl)^t}f=\zeta^k f .
 \tag{FIA.29}
\]

\emph{Proof.}
Unpack (FIA.28).  It contains a standard vacuum \(f_0\),
a supplied sum \(e_\Sigma=\sum_{\mathfrak R\in P_g\cup Q}
e_{\mathfrak R}\), and the exact
\(\texttt{Is\allowbreak Theta\allowbreak At}\) assertion for the half \(a=\bar e_\Sigma\)
and finite function (FIA.24).
The single-prime idempotents are pairwise orthogonal.
Indeed FIA2 identifies each with the characteristic idempotent
of its \(\mathfrak R\)-component; distinct primes give disjoint
components.  Thus
\[
            e_{\mathfrak P}e_\Sigma=e_{\mathfrak P}.
 \tag{FIA.30}
\]
The vacuum's \texttt{Is\allowbreak Std\allowbreak Char} field implies the precise
\texttt{Is\allowbreak Char\allowbreak AF} hypothesis used in FIA1 and FIA3.
Apply FIA3 at \(e_{\mathfrak P}\) to obtain \(l,\zeta\).
It has \(e_{\mathfrak P}l=l\), so (FIA.30) gives
\(e_\Sigma c(kl)^t=c(kl)^t\).  This proves its zero
infinite component and makes it an element of the half
opposite to \(a=\bar e_\Sigma\).

The literal \texttt{Is\allowbreak Theta\allowbreak At} is an \texttt{All\allowbreak Some}:
at least one linear model \(T_a\) satisfies
\texttt{Is\allowbreak Place\allowbreak Model}, and \emph{every} such model sends
the finite function \(\Phi\) to the same supplied vector \(f\).
Choose one of those models.  FIA4 has verified the guard
needed for its modulation axiom, which gives
\[
 H_{c(kl)^t}f
   =T_a\!\left(t\longmapsto
           \psi(B(t,\allowbreak{} c(kl)^t))\,\Phi(t)\right).
 \tag{FIA.31}
\]
For \(\Phi(t)\ne0\), its \(\Omega_K\) factor forces the
top identity block at \(\mathfrak P\); its integrality guard
also holds.  FIA3 and additivity of \(B\) in its second
argument show that the character in (FIA.31) is \(\zeta^k\).
For \(\Phi(t)=0\), the equality of the two multiplied
functions holds automatically.  Linearity of the
\emph{chosen actual} model therefore gives
\(H_{c(kl)^t}f=T_a(\zeta^k\Phi)=\zeta^k f\).
This is an equality of the stored functions, not a replacement
of \(f\) by a newly normalized vector.  Neither a scalar
uniqueness theorem for Weil representations nor the existence
of a new model is needed.
% END FIA5

% BEGIN FIA6
\paragraph{FIA6: the original good-prime set is nonempty.}
\label{claim:FIA6}
If \(L:\texttt{Lev\allowbreak Ch}(D,\allowbreak{} T_0)\) and \(L.\texttt{Is}(\mu_0)\),
then \(L.P_g\ne\varnothing\).  Moreover the supplied
\(L.\beta_V(i)\), \(L.\beta_W(i)\) have the full self-dual
lattices required by FIA4--FIA5.

\emph{Proof.}
The original character frame contains
\(\texttt{Pi\allowbreak Bar.\allowbreak pb}(0)=\mathfrak p_0\), a nonzero prime of
\(\mathcal O_{K_0}\), where
\(K_0=D.\mathrm{host}.K\subset E\).
It also records that the positive rational prime
\(D.\mathrm{host}.g(0)\) belongs to \(\mathfrak p_0\).
There is a prime \(\mathfrak P\) of \(\mathcal O_E\) over
\(\mathfrak p_0\).  Here is a short proof of this particular
lying-over assertion.  Localize the finite torsion-free
\(\mathcal O_{K_0}\)-module \(\mathcal O_E\) at
\(\mathfrak p_0\).  Over the resulting discrete valuation
ring it is free of positive rank.  Its reduction modulo
the maximal ideal is consequently a nonzero finite-dimensional
algebra over \(k(\mathfrak p_0)\).  A maximal ideal of that
algebra pulls back to a prime of the localization and then
of \(\mathcal O_E\); its contraction is \(\mathfrak p_0\),
because the residue-field scalar map is injective.
The pulled-back prime is nonzero, as it contains
\(D.\mathrm{host}.g(0)\ne0\); in the Dedekind ring
\(\mathcal O_E\) it is a height-one prime.

By definition this says \(L.C.P.\texttt{Div}(0,\mathfrak P)\).
The literal field \(L.\texttt{Is}.P_g\) is the equivalence
\[
       \mathfrak P\in L.P_g
          \quad\Longleftrightarrow\quad
       \exists k\in\{0,1\},\ L.C.P.\texttt{Div}(k,\mathfrak P).
 \tag{FIA.32}
\]
It proves \(\mathfrak P\in L.P_g\).
This independent argument reaches the exact
\texttt{exists\_\allowbreak mem\_\allowbreak Pg} conclusion and avoids the source's
longer Galois-conjugacy detour.
Splitting at the selected \(\mathfrak P\) is supplied when
the actual \texttt{Is\allowbreak Th\allowbreak PVec\allowbreak Of} is unpacked; it has not
been inferred from lying over alone.

Finally, the two level fields give
\(\texttt{Is\allowbreak SDFrame}(V,\beta,\Lambda,1)\).
This means \(\texttt{Is\allowbreak Base\allowbreak Frame}(V,\beta,\Lambda,1^c)\)
for some \(c\).  Its two lattice inclusions identify
\(\Lambda\) with the coordinates of the original full
standard lattice \(V.L0\), and its duality field is
\(\Lambda^\vee=\Lambda\).  Hence \(\Lambda\) is finitely
generated, spans \(V\), and has precisely the
\texttt{Is\allowbreak Self\allowbreak Dual\allowbreak C} property used in FIA5.
All these facts concern the stored lattice and frame.
% END FIA6

% BEGIN FIA7
\paragraph{FIA7: finite Heisenberg shifts commute with the original creation words.}
\label{claim:FIA7}
Suppose \(\theta\) is the supplied \texttt{Is\allowbreak Inf\allowbreak Fam}, and
\(u\in X\) has zero infinite part.  For the original
\texttt{pure\allowbreak Tensor\allowbreak O} and any polynomial \(P\),
\[
 \operatorname{pureTensorO}(H_u f)(P)
       =H_u\bigl(\operatorname{pureTensorO}(f)(P)\bigr).
 \tag{FIA.33}
\]
All Cartan involutions, bases, and scaling factors in this
formula are the same on the two sides.  The operation is also
homogeneous in its seed:
\[
       \operatorname{pureTensorO}(z f)(P)
          =z\,\operatorname{pureTensorO}(f)(P)
                         \qquad(z\in\mathbb C).
 \tag{FIA.34}
\]
These statements apply to every actual polynomial kernel.

\emph{Proof.}
Let \(v(r)=\texttt{line\allowbreak At}(\theta,\tau,\allowbreak{} Z)(r)\) be any
one of the real lines used to define \(\texttt{d\allowbreak Rho}\).
The defining property of \(\texttt{Is\allowbreak Inf\allowbreak Fam}\) identifies
the infinite tensor algebra with its archimedean components;
the line has prescribed component \(rZ\) at \(\tau\), zero
at the other infinite places, and zero finite part.
Its existence and uniqueness follow by taking these
components and the zero finite component.  Thus it is the
source's chosen \texttt{line\allowbreak At}.
Since finite and infinite adele components multiply to zero,
\[
                B(u,\allowbreak{} v(r))=B(v(r),\allowbreak{} u)=0.
 \tag{FIA.35}
\]
The explicit formula (FIA.2) gives the Heisenberg composition
law
\[
           H_xH_y=\psi(b(x,\allowbreak{} y))H_{x+y}.
 \tag{FIA.36}
\]
Equation (FIA.35) therefore gives
\(H_{v(r)}H_u=H_uH_{v(r)}\).
For a fixed evaluation point \(w\), the right side is
\[
     \psi(b(w,\allowbreak{} u))\,(H_{v(r)}f)(w+u),
 \tag{FIA.37}
\]
whose prefactor is constant in the real variable \(r\).
Differentiating at \(r=0\) proves
\[
       \mathrm d\rho(Z)\,\allowbreak{} H_u f
           =H_u\,\mathrm d\rho(Z)f.
 \tag{FIA.38}
\]
This is pointwise differentiation of an explicit formula.
It does not pass a derivative through an arbitrary endomorphism
of the space of all functions.
It also agrees with the source's total \texttt{deriv} convention:
the phase in (FIA.37) is nonzero, so differentiability of one
side is equivalent to that of the other, and the zero value
assigned at a nondifferentiable point respects the equality.

Each creation operator is the prescribed scalar multiple of
\(\mathrm d\rho(Z)-i\,\mathrm d\rho(JZ)\).
Equation (FIA.38) and linearity of \(H_u\) give commutation
with each such operator.  Induction through the exact ordered
list defining \texttt{op\allowbreak Pow} gives commutation for every
monomial word.  Taking the finite linear combination of
monomial words defining the polynomial proves (FIA.33);
no reordering of creation operators is necessary.
Constant multiplication in the same one-variable derivatives
proves (FIA.34), first for each word and then for its finite
polynomial sum.
For the actual seed, the given
\texttt{th\allowbreak P\_\allowbreak pt\_\allowbreak theta} at \(P=1\) also identifies it as
the original theta function, since
\(\operatorname{pureTensorO}(f)(1)=f\).
% END FIA7

% BEGIN FIA8
\paragraph{FIA8: the uniform original averaging operator.}
\label{claim:FIA8}
Let \(R:\texttt{Real\allowbreak Dat}(D,\allowbreak{} A,\allowbreak{} K_s)\) be the actual supplied
datum.  Assume exactly
\[
 \begin{gathered}
 R.\texttt{to\allowbreak Lev\allowbreak Ch.\allowbreak Is}(A.\mu_0),\qquad R.\texttt{models\_\allowbreak inf},
 \qquad R.\texttt{th\allowbreak P\_\allowbreak }\omega_P,\\
 R.\texttt{th\allowbreak P\_\allowbreak v\allowbreak P},\qquad R.\texttt{th\allowbreak P\_\allowbreak pt\_\allowbreak theta}.
 \end{gathered}                                                    \tag{FIA.39}
\]
There is a single integer \(N_R>0\) such that for every
original index \(i\), every tuple \(t\), every
\(h\in G(V_c(i))(\mathbb A_F)\),
\(g\in G(W_c(i))(\mathbb A_F)\), and every nonzero
\(\xi\in E^{\operatorname{rank}W_c(i)}\), there is a list
\[
       a_0,\ldots,\allowbreak{} a_{N_R-1}\in
                   \mathbb A_E^{\operatorname{rank}V_c(i)}
 \tag{FIA.40}
\]
such that, simultaneously for every polynomial \(P\),
\[
 \sum_{k=0}^{N_R-1}
 H_{\operatorname{along}(\xi,\allowbreak{} a_k)}
       \left(R.\omega_P(i,\allowbreak{} t)(h,\allowbreak{} g)\,
                        R.\mathrm{ptP}(i,\allowbreak{} t)(P)\right)=0 .
 \tag{FIA.41}
\]
The equality is an equality of functions on the actual
oriented matrix space.  Here
\(\operatorname{along}(\xi,\allowbreak{} a)=\iota_E(\xi)a^t\)
is literally the source's \texttt{along}.

\emph{Proof.}
FIA6 chooses one \(\mathfrak P\in R.P_g\), once for the
whole datum; set \(N_R=q_{\mathfrak P}\).
Fix \(i,\allowbreak{} t,\allowbreak{} h,\allowbreak{} g,\xi\ne0\), and abbreviate
\(S=V_c(i)\), \(T=W_c(i)\).
The original rank formulas give
\(\operatorname{rank}T=r\le\operatorname{rank}S\);
the required self-dual lattice statements were checked in FIA6.
Let \(M_S(h),\allowbreak{} M_T(g)\) denote the actual coordinate matrices
\texttt{mt}.  Put
\[
                   c=M_T(g^{-1})\,\iota_E(\xi).
 \tag{FIA.42}
\]
For every split chart \(\pi\), the matrix \(\pi(M_T(g))\)
is invertible, while \(\pi(\iota_E(\xi))\ne0\):
the composite \(E\to F_v\) is injective.
Therefore \(\pi(c)\ne0\), exactly the nonvanishing condition
in FIA5.
Apply that claim to the \emph{stored}
\(f=R.vP(i,\allowbreak{} t)\) and its supplied \texttt{Is\allowbreak Th\allowbreak PVec\allowbreak Of}
to obtain \(l,\zeta\) with period \(N_R\).

Define the averaging row vector by the literal formula
\[
                         a_k^t=(kl)^t M_S(h^{-1})
                          \qquad(0\le k<N_R).
 \tag{FIA.43}
\]
Its construction precedes the choice of a polynomial.
Matrix multiplication gives
\[
 \begin{aligned}
 \operatorname{along}(\xi,\allowbreak{} a_k)
   &=M_T(g)\,\allowbreak{} c(kl)^tM_S(h^{-1})\\
   &=\operatorname{actO}(S,\allowbreak{} T,\allowbreak{} h,\allowbreak{} g)\bigl(c(kl)^t\bigr).
 \end{aligned}                                                     \tag{FIA.44}
\]
By FIA7 and (FIA.29),
\[
     H_{c(kl)^t}\,\allowbreak{} R.\mathrm{ptP}(i,\allowbreak{} t)(P)
                =\zeta^k R.\mathrm{ptP}(i,\allowbreak{} t)(P).
 \tag{FIA.45}
\]
The supplied \texttt{th\allowbreak P\_\allowbreak pt\_\allowbreak theta} makes each polynomial
vector a theta function, so the \emph{original action's}
\texttt{Is\allowbreak Weil\allowbreak Datum\allowbreak O.\allowbreak om\_\allowbreak heis} applies to precisely that
vector.  Its covariance and linearity, with (FIA.44), yield
\[
 \begin{split}
 H_{\operatorname{along}(\xi,\allowbreak{} a_k)}
   \bigl(R.\omega_P(i,\allowbreak{} t)(h,\allowbreak{} g)\,\allowbreak{} R.\mathrm{ptP}(i,\allowbreak{} t)(P)\bigr)
   =\zeta^k
       R.\omega_P(i,\allowbreak{} t)(h,\allowbreak{} g)\,\allowbreak{} R.\mathrm{ptP}(i,\allowbreak{} t)(P).
 \end{split}                                                       \tag{FIA.46}
\]
Summing this finite identity proves (FIA.41), because
\[
       \sum_{k=0}^{N_R-1}\zeta^k
                   =\frac{\zeta^{N_R}-1}{\zeta-1}=0 .
 \tag{FIA.47}
\]
There is no convergence or exchange-of-integrals issue in this
step.

The source's \(\texttt{Iso\allowbreak Vec}(T,\xi)\) is the conjunction
of \(\xi\ne0\) with its Hermitian norm being zero.
Only its first clause has been used.  Restricting (FIA.41)
to such vectors proves the exact
\texttt{Theta\allowbreak Support.\allowbreak sum\_\allowbreak heis\allowbreak O\_\allowbreak along\_\allowbreak eq\_\allowbreak zero} and its wrapper:\par
\noindent
\texttt{fock\_\allowbreak vectors\_\allowbreak killed\_\allowbreak
heisenberg\_\allowbreak average\_\allowbreak isotropic}.
Our quantifiers are slightly stronger: \(N_R\) is independent
of \(i,\allowbreak{} t,\allowbreak{} h,\allowbreak{} g,\xi,\allowbreak{} P\), and the list is independent of \(P\).
No independence of the list from \(i,\allowbreak{} t,\allowbreak{} h,\allowbreak{} g,\xi\) is asserted.
The case of a zero-dimensional target has no nonzero \(\xi\),
so it causes no exception.
% END FIA8

\paragraph{What has been supplied to the cuspidal-extension branch.}
FIA8 closes the finite \texttt{hsupp} premise on every original
tuple, including the base tuple.  It uses the original level
data, actual half-datum, original creation operators and original
Weil covariance, all coupled as in (FIA.39).
The same nonzero field embedding and the same residue kernel
are used in the frame computation and the transported-vector
computation.  The construction contains no unconstrained
auxiliary root of unity, replacement of a stored vector, or
new theta action.  BTH13--BTH15's scalar-normalization argument
is compatible with this calculation but is not a dependency:
only the already supplied covariance field is used.

The constant-term invariance under \(\operatorname{along}\),
its measure normalization, and the implication from those
constant terms to cuspidality are the next separate arguments.
The original smooth multiplicity-space extension still also
requires its analytic and representation-theoretic producers.
FIA8 makes no assertion of moderate growth, lowering
annihilation, bounded holomorphic realization, simultaneous
\texttt{io\allowbreak P} existence, or original geometric/rational reading.
In particular a published rapid-decay theorem is neither used
nor certified by this component.


\section{From the original isotropic averages to cuspidal theta integrals}
\label{sec:constant-terms-v33}

This chapter supplies the constant-term and cuspidality portion of
the base-theta extension in BTH12.  The finite cancellation itself is
proved in FIA1--FIA8.  We use the original full Weil datum, its actual
character, the original positive definite source, and the original
coefficient functions.  No new theta kernel is selected.  The final
holomorphic realization and the remaining smooth-growth and lowering
arguments are distinct obligations.

Write $E/F$ for the fixed CM extension, $\delta=\xi_0\in E^\times$
for its purely imaginary element, $S$ for the positive definite source,
and $T$ for the target.  Put $s=\operatorname{rank}S$,
$n=\operatorname{rank}T$, and
$W=M_{n,\allowbreak{} s}(\mathbb A_E)$.  The Hermitian Gram matrices in the original
frames are $G_S,\allowbreak{} G_T$.  We use the source's oriented form
$B=\operatorname{symplO}$, its half $B/2=\operatorname{symplOH}$,
and the actual operators
\[
 H_u f(w)=\psi\bigl(B(w,\allowbreak{} u)/2\bigr)f(w+u),\qquad
 \omega(h,\allowbreak{} g)H_u=H_{r(h,\allowbreak{} g)u}\omega(h,\allowbreak{} g).
 \tag{CTC.1}
\]
All operator identities below are on the original theta space
$\Theta$.  It consists of the smooth, finite-level, rationally
quasiperiodic functions with the prescribed character.  The Weil
datum preserves this space, preserves its Haar $L^2$ norm, has the
source's pointwise continuity, rational automorphy at zero and split
positivity.  These are the actual contracts already constructed in
CW23 and compared in BTH14; we do not assume joint orbit continuity
without proving what is needed.

% BEGIN CTC1
\paragraph{CTC1: the original transvection subgroup and its quotient.}
\label{claim:CTC1}
Let $\xi\in E^n$ be nonzero and isotropic.  Define
\[
 N_\xi=\xi\xi^*G_T,\qquad
 j(a)=\xi a\quad(a\in\mathbb A_E^s),\qquad
 R(w)=\xi^*G_Tw.
 \tag{CTC.2}
\]
Then $N_\xi^2=0$, $N_\xi^*G_T=G_TN_\xi$, and
$N_\xi\ne0$.  The last assertion follows because $G_T$ is invertible
and both factors of the rank-one matrix are nonzero.  Direct
multiplication gives
\[
 z(t)=1+(t\delta)N_\xi,\qquad
 z(t)z(t')=z(t+t'),\qquad t,\allowbreak{} t'\in\mathbb A_F.
 \tag{CTC.3}
\]
These are unitary.  Conversely, if $1+aN_\xi$ is unitary, its unitary
equation reads $(a+\bar a)G_TN_\xi=0$.  A nonzero rational entry of
$N_\xi$ is invertible in $\mathbb A_E$, so $a+\bar a=0$.
Writing $\mathbb A_E=\mathbb A_F\oplus\delta\mathbb A_F$ gives
$a=t\delta$ for a unique $t$.  Thus the generating set in the
definition of $Z_I(T,\xi)$ is already a group, and (CTC.3) is exactly
that subgroup, including the source's algebraic subgroup closure.

The map $t\mapsto z(t)$ and its inverse are continuous: recover $a$
by dividing one nonzero entry of $z-1$ by the corresponding entry of
$N_\xi$, then take its $\delta$-coordinate.  The same coordinate
description shows the image is closed.  Moreover
\[
 Z_I(T,\xi)\simeq(\mathbb A_F,+),\qquad
 Z_I(T,\xi)\cap U(T)(F)\simeq(F,+).
 \tag{CTC.4}
\]
For the rational assertion the recovered $a$ is in $E$ and its skew
condition makes $a/\delta\in F$.

Here and below the adelic vector quotients are compact.  To see this
directly, for $x\in\mathbb A_F$ choose a positive integer $q$ clearing
its finitely many finite denominators.  The Chinese remainder theorem
chooses $b\in\mathcal O_F$ with $qx-b\in q\widehat{\mathcal O}_F$.
Subtract $b/q$, then subtract an element of $\mathcal O_F$ to put the
infinite component in a fixed closed fundamental parallelepiped.
The resulting representative lies in a compact set.  The same
coordinates show that $F$ is discrete.  Products and an $F$-basis of
$E$ prove the assertion for all vector spaces used here.  Haar
measure on a quotient by a discrete cocompact subgroup has finite
positive total mass.  This can also be checked with a small open
set having disjoint rational translates and a finite cover of a
compact representative set.  Any measurable fundamental domain has
that same mass, by partitioning it and another domain by their
countably many rational translates.  In particular all the source's
transvection fundamental domains have finite, nonzero measure.

The matrix identities needed later are
\[
 R(j(a))=0,\qquad
 r(1,\allowbreak{} z(t))w=w+j(t\delta R(w)),\qquad
 r(1,\allowbreak{} z(t))j(a)=j(a),\qquad B(j(a),\allowbreak{} j(b))=0.
 \tag{CTC.5}
\]
They follow by multiplying (CTC.2), using $\xi^*G_T\xi=0$.
% END CTC1

% BEGIN CTC2
\paragraph{CTC2: a uniform pointwise bound on compact Weil orbits.}
\label{claim:CTC2}
For $f\in\Theta$ and a compact set $C$ of actual unitary pairs,
\[
                 \sup_{p\in C,\,\allowbreak{} w\in W}|\omega(p)f(w)|<\infty.
 \tag{CTC.6}
\]
Here is a proof from the source's pointwise contracts.  Take a real
basis $b_1,\ldots,\allowbreak{} b_d$ of $W_\infty$ and write $D_v$ for the
infinitesimal Heisenberg derivative.  The derivative is an ordinary
directional derivative plus a purely imaginary linear coefficient.
It preserves $\Theta$, is real-linear in $v$, and commutes with finite
Heisenberg operators.  Differentiated covariance gives
\[
 D_v\omega(p)g=\omega(p)D_{r(p)^{-1}v}g.
 \tag{CTC.7}
\]
The differentiation is justified by the $L^2$ difference-quotient
argument in WA3 and BTH5.  Iterating it uses only finitely many
derivative words of length at most $d$.

Choose a finite Heisenberg neighborhood fixing $f$.  It fixes all
these derivative words, since real and finite Heisenberg directions
commute.  Continuity of $(p,\allowbreak{} u)\mapsto r(p)^{-1}u$ and compactness of
$C$ give a single compact open finite neighborhood $K_f$ whose
translates fix every $\omega(p)D_{v_1}\cdots D_{v_a}f$, for $p\in C$
and $a\le d$.  In making this assertion, bounded $v_i$ are expanded
in the fixed basis, so only finitely many fixed words are involved.

We record the elementary estimate which converts a quotient norm
to a real-cube norm.  If $g$ has this finite level, $|g|^2$ is
constant on $K_f$-translates and periodic under $W(E)$.  For a fixed
real unit cube $Q$, integrate over $Q\times K_f$.  Its projection
to $W(E)\backslash W$ has bounded multiplicity: after choosing one
point of a fiber, all other rational differences lie in the finite
set $W(E)\cap((Q\times K_f)-(Q\times K_f))$.  The same bound holds
for every translate of this compact set.  Since $K_f$ has positive
measure, there is a constant $A$ independent of $g$ and $w$ such that
\[
 \int_{[0,1]^d}\left|g\left(w+\sum_i t_i b_i\right)\right|^2dt
             \le A\|g\|_2^2.
 \tag{CTC.8}
\]
Haar measure on $W$ factors into real Lebesgue measure and finite
Haar measure up to a fixed positive constant; that constant is
absorbed into $A$.  Translations of $|g|^2$ still have the same
rational periodicity and quotient integral, so (CTC.8) is uniform
in $w$.

For completeness the required Sobolev estimate uses only the
fundamental theorem of calculus on $[0,1]$.  Multiply a restriction
$t\mapsto g(w+tb_i)$ by a phase of absolute value one that integrates
its purely imaginary derivative coefficient.  Its ordinary
derivative then has absolute value $|D_{b_i}g(w+tb_i)|$.  Cauchy--Schwarz
and averaging the endpoint inequality give
\[
 |g(w)|^2\le 2\int_0^1|g(w+tb_i)|^2dt
                 +2\int_0^1|D_{b_i}g(w+tb_i)|^2dt.
 \tag{CTC.9}
\]
Iterate in the $d$ coordinate directions, applying the same inequality
to every derivative word which occurs.  If each cube integral of
every word of length at most $d$ is at most $M$, the result is
$|g(w)|^2\le4^dM$.  This includes $d=0$.

The vectors $r(p)^{-1}b_i$ have bounded coordinates for $p\in C$.
By (CTC.7), real multilinearity and unitarity of the original Weil
action, the $L^2$ norms of every required word applied to
$\omega(p)f$ are uniformly bounded by a finite sum of the norms
of fixed words applied to $f$.  Those norms are finite: each such
word is a continuous theta function and its absolute square descends
to the compact rational quotient.  Apply (CTC.8)--(CTC.9).  This proves
(CTC.6) without inferring joint continuity from pointwise continuity.
% END CTC2

% BEGIN CTC3
\paragraph{CTC3: local character duality and compact annihilators.}
\label{claim:CTC3}
Let $K=F_v$ at a finite place and let $\psi_v$ be the original local
character.  FIA1 proves it is nontrivial.  Every additive character
$\chi:K\to\mathbb C^\times$ trivial on an open subgroup has the form
\[
                         \chi(x)=\psi_v(ax)\quad(a\in K).
 \tag{CTC.10}
\]
We give the finite argument, rather than import a local duality
theorem.  Write $\mathfrak p^c$ for the largest fractional
$\mathcal O_K$-ideal on which $\psi_v$ is trivial.  Such $c$ exists:
continuity provides a ball on which all character values are close
to one, all integer multiples stay in that ball, and a complex
number all of whose powers stay sufficiently close to one is one.
Nontriviality excludes arbitrarily large fractional ideals from
the kernel.  Thus the kernel contains $\mathfrak p^c$ and does not
contain $\mathfrak p^{c-1}$.

Choose $m$ with $\chi(\mathfrak p^m)=1$.  For $j$ sufficiently large,
the pairing
\[
 \mathfrak p^{c-m}/\mathfrak p^{c+j}\ \times\
 \mathfrak p^{-j}/\mathfrak p^m\longrightarrow\mathbb C^\times,
                   \qquad(a,\allowbreak{} x)\longmapsto\psi_v(ax)
 \tag{CTC.11}
\]
is perfect.  Indeed, if a nonzero class of $a$ annihilates the
second factor, then $a\mathfrak p^{-j}\subset\ker\psi_v$, which
forces $v(a)-j\ge c$, a contradiction.  The two finite groups
have equal cardinalities.  Finite character orthogonality shows
their duals have those same cardinalities, proving perfection.
The set of $a\in\mathfrak p^{c-m}$ which realize $\chi$ on
$\mathfrak p^{-j}$ is therefore a nonempty compact coset of
$\mathfrak p^{c+j}$.  These cosets are nested.  Their intersection
is nonempty and realizes $\chi$ on all of $K$.  Coordinatewise the
same conclusion holds for $K^q$.

We shall also need the following consequence with its quantifiers.
For a nondegenerate $K$-bilinear pairing $b$ on two copies of $K^q$
and any neighborhood $V$ of zero in the first copy, there exists a
compact open additive subgroup $L$ of the second copy such that
\[
 \psi_v(b(\eta,\allowbreak{} L))=1\quad\Longrightarrow\quad\eta\in V.
 \tag{CTC.12}
\]
Choose linear coordinates for $b$, choose a small product of
fractional ideals contained in $V$, and take the reciprocal ideals
with conductor $\mathfrak p^c$.  Formula (CTC.11), or its valuation
argument before passing to a quotient, proves (CTC.12).  All changes
of coordinates are the original invertible Gram matrices.
% END CTC3

% BEGIN CTC4
\paragraph{CTC4: split-half coinvariant functionals separate theta vectors.}
\label{claim:CTC4}
Let $e,\bar e$ be the original pair of idempotents at a split finite
place.  Put $Y=\bar eW$ and $Y'=eW$.  If $g\in\Theta$ and
\[
 m(H_tg)=0\quad\hbox{for all }t\in Y'
 \quad\hbox{and all }m\hbox{ invariant under }H(Y),
 \tag{CTC.13}
\]
then $g=0$.  The functionals are algebraic; no continuity is imposed.
The form $B$ pairs $Y$ and $Y'$ nondegenerately, and each is isotropic.
In the original frames this follows by multiplying by $G_S^{-1}$,
$G_T$ and the nonzero $\delta$ in the trace formula for $B$.

Here is the needed algebraic separation argument.  Let a commuting
additive action $\rho$ on a complex vector space $V$ have the property
that for every $u,\allowbreak{} f$ some positive integer $N$ satisfies
$\rho(Nu)f=f$.  If $g\ne0$, there exist an additive character $\chi$
and a functional $\ell$ such that
\[
 \ell(g)\ne0,\qquad \ell(\rho(u)f)=\chi(u)\ell(f),
 \qquad \rho(u)g=g\Longrightarrow\chi(u)=1.
 \tag{CTC.14}
\]
To prove this, use the commutative group algebra
$A=\mathbb C[G]$, with basis elements $X_u$, and choose a maximal
ideal $\mathfrak m$ containing the annihilator of $g$.
Every $X_u$ has a positive power equal to one in $A/\mathfrak m$,
because it has such a power on $g$.  Thus this residue field is
algebraic over $\mathbb C$, so it is $\mathbb C$.  Write
$c:A\to\mathbb C$ for the resulting homomorphism and
$\chi(u)=c(X_u)$.

Let $V_0$ consist of those $f$ killed by some $a\in A$ with
$c(a)\ne0$.  Commutativity makes this a vector subspace, and
$g\notin V_0$.  Choose a functional on $V/V_0$ nonzero at $g$.
For given $u,\allowbreak{} f$, take a positive common multiple $N$ of a period of
$u$ on $f$ and a period on $g$.  Then $\chi(u)^N=1$.  The polynomial
\[
 q(X)=\sum_{i=0}^{N-1}X^i\chi(u)^{N-1-i}
 \quad\hbox{satisfies}\quad
 q(X)(X-\chi(u))=X^N-1,
 \quad q(\chi(u))=N\chi(u)^{N-1}\ne0.
\]
Consequently $\rho(u)f-\chi(u)f\in V_0$.  This proves every assertion
in (CTC.14), including the common eigenfunctional on all of $V$.

Apply this argument to $V=\Theta$ and $\rho(u)=H_u$, $u\in Y$.
Finite smoothness supplies the required periods: for a fixed local
vector $u$, multiplication by sufficiently large powers of the
residue characteristic puts it into any prescribed open stabilizer.
The stabilizer of $g$ is open, so (CTC.14) makes $\chi$ trivial on
an open subgroup.  CTC3 and the original nondegenerate pairing give
$t_0\in Y'$ with $\chi(u)=\psi(B(u,\allowbreak{} t_0))$.
The Heisenberg commutator is
\[
 H_uH_{t_0}=\psi(B(u,\allowbreak{} t_0))H_{t_0}H_u.
 \tag{CTC.15}
\]
It follows from (CTC.14)--(CTC.15) that
$m(f)=\ell(H_{t_0}f)$ is $H(Y)$-invariant.  But
$m(H_{-t_0}g)=\ell(g)\ne0$, contradicting (CTC.13).
An algebraic functional on $\Theta$ extends to the ambient function
space by a basis extension, exactly as required by the source.
% END CTC4

% BEGIN CTC5
\paragraph{CTC5: finite averaging transfers a small-row test to along coinvariants.}
\label{claim:CTC5}
Keep the split place of CTC4.  Suppose $g\in\Theta$ and $V$ is a
neighborhood of zero in $\mathbb A_E^s$ with the following property:
\[
 m(H_tg)=0\quad\text{whenever }t\in eW,\ R(t)\in V,
 \quad m\text{ is }H(\bar eW)\text{-invariant}.
 \tag{CTC.16}
\]
Then every functional $\ell$ invariant under all $H_{j(a)}$
satisfies $\ell(g)=0$.

The exact original trace calculation is
\[
 B\bigl(t,\allowbreak{} j(\bar s)\bigr)
   =\operatorname{Tr}_{\mathbb A_E/\mathbb A_F}
        \bigl(\delta\,(R(t)G_S^{-1})\cdot s\bigr).
 \tag{CTC.17}
\]
For $t\in eW$ this is the nondegenerate local pairing of CTC3
applied to the row coordinate $R(t)$; multiplication by
$\delta G_S^{-1}$ is invertible on that coordinate space.
Choose a compact open $L\subset e\mathbb A_E^s$ such that
triviality of (CTC.17) as a character of $L$ implies $R(t)\in V$.
Let $\alpha(s)=j(\bar s)$ for $s\in L$.  This lies in $\bar eW$
and in the along space, which is isotropic.

The continuous orbit $s\mapsto H_{\alpha(s)}g$ has an open
stabilizer in the compact group $L$, by the finite level of $g$.
That stabilizer has finite index.  For representatives $s_1,\ldots,\allowbreak{} s_q$,
where $q>0$, set $g_L=\sum_i H_{\alpha(s_i)}g$.
The sum is invariant under all $H_{\alpha(s)}$, $s\in L$.
For an $H(\bar eW)$-invariant $m$, the commutator gives
\[
 m(H_tH_{\alpha(s)}g)
          =\psi(B(t,\alpha(s)))m(H_tg).
 \tag{CTC.18}
\]
If this character of $L$ is trivial, (CTC.12), (CTC.16) and
(CTC.18) give $m(H_tg_L)=0$.  If it is nontrivial, apply (CTC.18)
to the invariant vector $g_L$ and to an $s$ where the character is
not one.  The same conclusion follows.  CTC4 now gives $g_L=0$.
Applying the along-invariant $\ell$ yields $q\ell(g)=0$, hence
$\ell(g)=0$.  This finite average is a proof about coinvariants;
it is separate from the arithmetic finite average FIA constructs
for the actual Fock vectors.
% END CTC5

% BEGIN CTC6
\paragraph{CTC6: positivity on the actual along coinvariants.}
\label{claim:CTC6}
For a transvection $z\in Z_I(T,\xi)$ supported at the split place
$e+\bar e$, there exists a single $r_z>0$ such that
\[
             \ell(\omega(1,\allowbreak{} z)f)=r_z\ell(f)
 \tag{CTC.19}
\]
for every $f\in\Theta$ and every along-invariant $\ell$.
First, the rank-one determinant formula gives
$\det(1+aN_\xi)=1+a\xi^*G_T\xi=1$.
In the original split-positivity axiom take the two determinant
lifts to be $1$.  Their splitting-character factors are exactly
$1$.  That axiom therefore supplies $r_z>0$ with (CTC.19) for
every functional invariant under the larger space $\bar eW$.
This is an application of the already constructed full Weil datum,
not an assumption that $r_z=1$.

Fix $f$ and put $g=\omega(1,\allowbreak{} z)f-r_zf$.  Write
$z^{-1}=1+aN_\xi$.  Its action on $t\in eW$ is
$r(1,\allowbreak{} z^{-1})t=t+j(aR(t))$.  If $R(t)$ is sufficiently small,
$j(aR(t))$ lies in the finite Heisenberg stabilizer of $f$.
It also lies in $eW$, as does $t$, so their Heisenberg half-phase
is zero.  Covariance consequently gives
\[
 H_t\omega(1,\allowbreak{} z)f=\omega(1,\allowbreak{} z)H_tf
                  \quad\text{when }R(t)\text{ is sufficiently small}.
 \tag{CTC.20}
\]
The split-positivity formula applied to $H_tf$ shows that $g$
satisfies (CTC.16).  CTC5 gives $\ell(g)=0$ for every along-invariant
$\ell$, proving (CTC.19).  The scalar $r_z$ came from the datum
before $f$ or $\ell$ was chosen, as required in the Fourier argument.
% END CTC6

% BEGIN CTC7
\paragraph{CTC7: a nontrivial eigenvalue for every nonzero rational row.}
\label{claim:CTC7}
Let $x\in M_{n,\allowbreak{} s}(E)$, embedded rationally in $W$, and suppose
$m=R(x)\ne0$.  Put
\[
                    N=\sum_{i,\allowbreak{} j}(G_S^{-1})_{ij}\bar m_jm_i.
 \tag{CTC.21}
\]
This is fixed by conjugation.  At any chosen complex embedding,
$G_S^{-1}$ is positive definite, so $N\ne0$.  Thus $N\in F^\times$
and $q=\delta^2N\in F^\times$.

The original host has a split place over $2$.  Use its original
split idempotent and choose a local $y\in F_v$ with
$\psi_v(y)\ne1$, by FIA1.  Let $t\in\mathbb A_F$ be supported
at this place with value $y/q$.  With $z=z(t)$ and
$b=t\delta m$, (CTC.5) gives $r(1,\allowbreak{} z)x=x+j(b)$.  The original
Hermitian trace formula, including its half, gives
\[
 h_O(j(b),\allowbreak{} x)=t\delta N,\qquad
 \tfrac12 B(j(b),\allowbreak{} x)
   =\tfrac12\operatorname{Tr}_{\mathbb A_E/\mathbb A_F}
                      (t\delta^2N)=tq.
 \tag{CTC.22}
\]
Consequently $\eta=\psi(B(j(b),\allowbreak{} x)/2)\ne1$.
No division by an integral unit at the dyadic place is asserted:
the calculation takes place in its field of fractions, and its
factor $2$ cancels before the value of the character is used.

For every along-invariant $\ell$ there is therefore a $c\ne1$
such that
\[
 \ell(H_x\omega(1,\allowbreak{} z)f)=c\ell(H_xf)
                         \qquad(f\in\Theta),\qquad c=\eta r_z.
 \tag{CTC.23}
\]
Indeed, covariance gives
$\omega(1,\allowbreak{} z)H_xf=H_{j(b)+x}\omega(1,\allowbreak{} z)f$.
Use $H_{j(b)}H_x=\eta H_{j(b)+x}$, invariance of $\ell$ under
$H_{j(b)}$, and (CTC.19) applied to $H_xf$.  This is precisely
(CTC.23).  If $c=1$, its absolute value and $r_z>0$ force
$r_z=1$, since $|\eta|=1$; then $\eta=1$, a contradiction.
The positivity of the source enters here through (CTC.21), and
the positivity of the Weil scalar enters through this last step.
% END CTC7

% BEGIN CTC8
\paragraph{CTC8: the compact-quotient Fourier averaging lemma.}
\label{claim:CTC8}
Let $Y/\Gamma$ and $Z/\Lambda$ be compact abelian quotients of
locally compact second countable groups, with finite positive Haar
measures.  Let $k(z,\allowbreak{} y)$ be bounded, separately continuous and
periodic in both variables.  Suppose a family of unitary characters
$\chi_a$ on $Y/\Gamma$ contains $1$, is closed under multiplication
and conjugation, and separates points.  Suppose also that, for each
nontrivial $\chi_a$, there are $z_a$ and $c_a\ne1$ such that
\[
 \int_{Y/\Gamma}\chi_a(y)k(z_a z,\allowbreak{} y)\,\allowbreak{} dy
       =c_a\int_{Y/\Gamma}\chi_a(y)k(z,\allowbreak{} y)\,\allowbreak{} dy
                      \quad\text{for every }z.
 \tag{CTC.24}
\]
Then $\Phi(y)=\int_{Z/\Lambda}k(z,\allowbreak{} y)\,\allowbreak{} dz$ is constant.

All integrals here are genuine integrals of bounded measurable
functions over finite measures.  Separate continuity on these
second countable spaces gives joint measurability: use countably
many shrinking neighborhoods and measurable choices of a nearby
point in one variable to express the function as a pointwise limit
of functions measurable in the product.  Dominated convergence
makes $\Phi$ continuous.  Fubini and invariance of quotient Haar
measure show that its $\chi_a$-coefficient $A_a$ satisfies
$A_a=c_aA_a$, hence $A_a=0$ for every nontrivial character in the
family.  Translation does not change the trivial coefficient.

For fixed $y_0$, the continuous function
$D(y)=\Phi(y+y_0)-\Phi(y)$ consequently has all these coefficients
zero.  To make the density step explicit, the complex linear span
of the characters is a conjugation-stable algebra containing the
constants and separating points.  Its real part is uniformly dense
in the real continuous functions by the usual elementary
Stone--Weierstrass argument: polynomial approximation to absolute
value makes its closure stable under maximum and minimum; functions
matching a prescribed continuous function at two specified points
are obtained from separation and constants; two finite-cover
arguments with maximum and minimum give a uniform approximation
to that function.  Thus its complex span uniformly approximates
$\bar D$.  Integrating $D$ against such approximations gives
$\int|D|^2=0$.  Since Haar measure is positive on nonempty open
sets and $D$ is continuous, $D=0$.  This proves the assertion for
every $y_0$, and proves constancy at every point, not merely almost
everywhere.  Normalizing either Haar measure is unnecessary.
% END CTC8

% BEGIN CTC9
\paragraph{CTC9: the actual transvection constant term is Heisenberg invariant.}
\label{claim:CTC9}
For every original Haar measure $\mu$ and rational fundamental
domain $\mathcal F$ of $Z_I(T,\xi)$, set
\[
 \mathcal C_\xi(f)=\int_{\mathcal F}\omega(1,\allowbreak{} z)f(0)\,\allowbreak{} d\mu(z).
\]
Then, for every $f\in\Theta$ and $a\in\mathbb A_E^s$,
\[
                   \mathcal C_\xi(H_{j(a)}f)=\mathcal C_\xi(f).
 \tag{CTC.25}
\]
This is the exact source theorem
\texttt{constant\allowbreak Term\_\allowbreak heis\allowbreak O\_\allowbreak along} on its original operators.

Let $Y=\mathbb A_E^s$ and $\Gamma=j^{-1}(W(E))$.  A nonzero
coordinate of $\xi$ shows $\Gamma=E^s$: divide a rational entry of
$j(a)$ in that row by its nonzero rational coefficient.  Thus
$Y/\Gamma$ is compact by CTC1.  Isotropy and quasiperiodicity make
$a\mapsto f(j(a))$ periodic modulo $\Gamma$.  Define the functional
\[
                         \ell(f)=\int_{Y/\Gamma}f(j(a))\,\allowbreak{} da.
 \tag{CTC.26}
\]
It is well-defined and linear on $\Theta$, because every theta
function is bounded in absolute value on the compact rational
quotient of $W$.  The phase in $H_{j(b)}f(j(a))$ is one by (CTC.5);
translation invariance of quotient Haar measure therefore makes
$\ell$ invariant under every $H_{j(b)}$.

For rational matrices $x\in W(E)$ use the characters
\[
             \chi_x(a)=\psi(B(j(a),\allowbreak{} x)).
 \tag{CTC.27}
\]
They are continuous, unitary, trivial on $\Gamma$, and satisfy
$\chi_{x+x'}=\chi_x\chi_{x'}$, $\chi_{-x}=\bar\chi_x$.
They separate $Y/\Gamma$: if all their values at $a$ are one,
the rational self-annihilator calculation of CW15 gives
$j(a)\in W(E)$, hence $a\in\Gamma$.
This uses the original oriented Gram form; its rational invertible
Gram matrix identifies its dual with the rational Hom space, as
in CW15, so no change of theta model is involved.

Take
\[
                         k(z,\allowbreak{} a)=\omega(1,\allowbreak{} z)f(j(a)).
 \tag{CTC.28}
\]
It is separately continuous by the original Weil and theta
contracts.  It is periodic in $a$ as above.  It is periodic in $z$
under rational transvections: these fix $j(a)$; move $H_{j(a)}$
past them by covariance, evaluate at zero, and use rational
automorphy of the original datum.  Choose a compact representative
set for $F\backslash\mathbb A_F$ in CTC1's identification
(CTC.4).  CTC2 bounds (CTC.28) there uniformly in $a$; rational
periodicity extends the bound to every $z$.  In particular all
integrals in (CTC.25)--(CTC.28) are absolutely convergent.

If $R(x)=0$, the original trace formula makes $\chi_x=1$.
Otherwise CTC7 applies to (CTC.26).  For rational $x$ the two
half-phases in the definition and the quasiperiodicity combine to
give the full modulation formula
\[
                             H_xg(w)=\psi(B(w,\allowbreak{} x))g(w).
 \tag{CTC.29}
\]
Therefore (CTC.23), applied to $\omega(1,\allowbreak{} z)f$, is exactly
(CTC.24) for the coefficient of (CTC.28).  CTC8 shows that
$\int_{\mathcal F}k(z,\allowbreak{} a)d\mu(z)$ is independent of $a$.
Finally covariance, (CTC.5), and $H_u g(0)=g(u)$ identify that
integral with the left side of (CTC.25); at $a=0$ it is the right
side.  This proves the claimed invariance with the source's
unscaled Haar measures and with no remaining analytic theorem
hidden in the Fourier step.
% END CTC9

% BEGIN CTC10
\paragraph{CTC10: finite cancellation gives zero constant terms of the original lift.}
\label{claim:CTC10}
Let $H=U(S)(\mathbb A)$ be the original definite source group.
Assume it has the compact rational representative set provided
by the original finite class presentation in BTH3.  Let $X$ be the
original finite-dimensional coefficient space, $L:X\to\Theta$
the original linear kernel map, and $F:H\to X$ the original
continuous rationally left-invariant coefficient function.  Let
$\nu,\mathcal H$ be its original Haar measure and fundamental
domain.  Suppose that for every $x,\allowbreak{} h,\allowbreak{} g$ and rational nonzero
isotropic $\xi$, there are $q>0$ and $a_1,\ldots,\allowbreak{} a_q$ with
\[
             \sum_{i=1}^q H_{j(a_i)}\omega(h,\allowbreak{} g)L(x)=0.
 \tag{CTC.30}
\]
Then the actual convergent theta integral
\[
                 \Theta_F(g)=\int_{\mathcal H}
                                 \omega(h,\allowbreak{} g)L(F(h))(0)\,\allowbreak{} d\nu(h)
 \tag{CTC.31}
\]
has $ZCusp$ in the literal source sense: every transvection
constant term is zero, for every Haar measure, rational fundamental
domain and right translate.

First fix $x,\allowbreak{} h,\allowbreak{} g,\xi$.  Apply the linear functional
$\mathcal C_\xi$ on the finite sum (CTC.30).  CTC9 makes all
$q$ terms equal.  Since $q>0$ in $\mathbb C$,
\[
          \int_{\mathcal F}\omega(h,\allowbreak{} zg)L(x)(0)\,\allowbreak{} d\mu(z)=0.
 \tag{CTC.32}
\]
Only ordinary linearity of an integral of a finite sum is being
used; all its terms are integrable by CTC9.

For the exchange with (CTC.31), put
$K(z,\allowbreak{} h)=\omega(h,\allowbreak{} zg)L(F(h))(0)$.  Choose a basis $x_b$ of $X$.
The expansion $L(F(h))=\sum_b F_b(h)L(x_b)$ is finite.
Each $F_b$ is continuous and each fixed-vector orbit evaluation
is continuous by the Weil contract.  Thus $K$ is jointly
continuous, even though no joint-continuity axiom for the entire
infinite-dimensional theta space was assumed.  It is rationally
periodic in each group variable by the original coefficient
invariance and rational Weil automorphy.  Compact rational
representatives for $H$ and CTC1 for $Z_I$ make it bounded
globally.  Both quotient measures are finite, so its product
integral is absolutely convergent.  Fubini and (CTC.32), with
$x=F(h)$, now give
\[
 \int_{\mathcal F}\Theta_F(zg)\,\allowbreak{} d\mu(z)
 =\int_{\mathcal H}\int_{\mathcal F}
           \omega(h,\allowbreak{} zg)L(F(h))(0)\,\allowbreak{} d\mu(z)d\nu(h)=0.
 \tag{CTC.33}
\]
This proves the original definite-source lift statement in exactly
the class-presentation scope needed here.  It does not assert that
a separate unrestricted class-number theorem for every possible
definite group has been reconstructed.
% END CTC10

% BEGIN CTC11
\paragraph{CTC11: the full unipotent quotient is compact and contains the central transvections.}
\label{claim:CTC11}
Let $C\in M_{n,\allowbreak{} r}(E)$ be an original isotropic frame: $r>0$,
its columns are independent, and $C^*G_TC=0$.  The subgroup
$N_I(T,\allowbreak{} C)$ in the source is the unipotent group fixing the
isotropic space $U=\operatorname{im}C$ and acting trivially on
$U^\perp/U$ and on $T/U^\perp$.  We check this directly from
its defining equations, and prove its rational quotient compact.

Choose a rational matrix $C'$ with $C^*G_TC'=1$; this exists
because $G_T$ is nondegenerate and $C$ has independent columns.
Replacing $C'$ by $C'-\tfrac12C(C'^*G_TC')$ makes it isotropic
without changing that pairing.  Complete the resulting hyperbolic
pair to a rational orthogonal decomposition
$T=U\oplus T_0\oplus U'$.  In this basis the Gram matrix is
\[
       J=\begin{pmatrix}0&0&1\\0&J_0&0\\1&0&0\end{pmatrix},
                         \qquad J_0\text{ nondegenerate}.
 \tag{CTC.34}
\]
The source's first two equations are $(n-1)C=0$ and
$C^*G_T(n-1)=0$.  The third says that $(n-1)x\in U(\mathbb A_E)$
for every rational $x\in U^\perp$.  A rational basis of $U^\perp$
spans its adelic extension, so the third equation holds on that
whole extension.  Together with unitarity these equations say
exactly that, for an adelic matrix $A$ and skew-Hermitian $B$,
\[
 n(A,\allowbreak{} B)=
 \begin{pmatrix}
  1&-A^*J_0&B-\tfrac12 A^*J_0A\\
  0&1&A\\0&0&1
 \end{pmatrix},\qquad B^*=-B.
 \tag{CTC.35}
\]
For example the first two equations fix the first block column
and last block row; the third fixes the middle diagonal block;
the remaining unitary equations determine the upper-middle
block and the Hermitian part of the upper-right block.  This
derives (CTC.35), including when $T_0$ has dimension zero.

Direct multiplication gives
\[
 n(A,\allowbreak{} B)n(A',\allowbreak{} B')=
 n\bigl(A+A',\allowbreak{} B+B'+\tfrac12(A'^*J_0A-A^*J_0A')\bigr).
 \tag{CTC.36}
\]
Thus these matrices already form a closed subgroup.  The source's
subgroup closure adds nothing.  The kernel of $n(A,\allowbreak{} B)\mapsto A$
is the central additive space of skew-Hermitian $B$; the quotient
is the additive $E$-matrix space of $A$.  Both are adelic vector
spaces defined over $F$, and their rational points are their
original rational subspaces.

For any $n(A,\allowbreak{} B)$, left multiply by a rational $n(a,0)$ to put
$A+a$ into a fixed compact vector representative set.  Then left
multiply by a rational central $n(0,\allowbreak{} b)$ to put its new central
coordinate into a fixed compact vector representative set.
CTC1 supplies both compact sets.  Their image under (CTC.35)
is compact and meets every left rational coset.  This proves
\[
                     N_I(T,\allowbreak{} C)(F)\backslash N_I(T,\allowbreak{} C)(\mathbb A)
                              \quad\text{is compact}.
 \tag{CTC.37}
\]
The rational subgroup is discrete in these coordinates.
Product additive Haar measure in $(A,\allowbreak{} B)$ is left and right
invariant by the shear law (CTC.36).  Every rational fundamental
domain therefore has finite positive measure, and every continuous
rationally invariant function has an absolutely convergent
integral over it.

Finally choose any column $\xi=C_j$.  It is nonzero and isotropic,
and
$N_\xi=C E_{jj}C^*G_T$.
For every defining matrix $n$ of $N_I$, the equations above give
$(n-1)N_\xi=N_\xi(n-1)=0$.
Consequently $Z_I(T,\xi)\subset N_I(T,\allowbreak{} C)$ and this subgroup is
central.  These are the exact source inclusions, not an appeal
to an unspecified parabolic conjugacy criterion.
% END CTC11

% BEGIN CTC12
\paragraph{CTC12: zero transvection constants imply genuine cuspidality.}
\label{claim:CTC12}
Let $\varphi$ be continuous and rationally left-invariant on the
original target group, and assume $ZCusp(T,\varphi)$.  Then
$IsCuspFun(T,\varphi)$ holds with genuinely convergent constant
terms.  The finite-level part of $IsAutFun$ is not needed for
this implication.

Fix an original isotropic frame $C$ and a right translate $g$.
Let $N=N_I(T,\allowbreak{} C)$, $\Gamma=N\cap U(T)(F)$ and
$F_g(n)=\varphi(ng)$.  CTC11 makes $F_g$ continuous on a compact
rational quotient, hence bounded and integrable for every
original Haar fundamental domain $\mathcal H_N$.  Choose a
column $\xi$ of $C$ and a rational fundamental domain
$\mathcal F_Z$ of the central subgroup $Z=Z_I(T,\xi)$.
Its measure $m_Z$ is positive and finite by CTC1.

For each $z\in Z$, centrality makes $n\mapsto zn$ descend to
the left quotient $\Gamma\backslash N$.  It preserves its
Haar measure.  One can verify this directly: $z^{-1}\mathcal H_N$
is another rational fundamental domain, and partitioning the two
domains into rational translates gives equality of their integrals.
Thus
\[
 \int_{\mathcal H_N}F_g(zn)\,\allowbreak{} dn
                   =\int_{\mathcal H_N}F_g(n)\,\allowbreak{} dn.
 \tag{CTC.38}
\]
The same equality for $|F_g|$ and finiteness of $m_Z$ justify
Fubini over $\mathcal F_Z\times\mathcal H_N$.  But $ZCusp$
says $\int_{\mathcal F_Z}F_g(zn)\,\allowbreak{} dz=0$ for every $n$.
Integrating and using (CTC.38) gives
\[
                 0=m_Z\int_{\mathcal H_N}\varphi(ng)\,\allowbreak{} dn.
 \tag{CTC.39}
\]
Since $m_Z>0$, the integral is zero.  The source quantifies
over all such frames, measures, domains and $g$; all have been
retained.  If no isotropic frame exists, the condition is vacuous.

The source's abstract measure lemma also covers nonintegrable
functions because its total Bochner integral is defined to be
zero there.  That convention is not needed in this human proof:
CTC11 proves integrability of every continuous automorphic
function on every unipotent quotient appearing in the definition.
% END CTC12

% BEGIN CTC13
\paragraph{CTC13: the exact base-kernel cuspidality output and its remaining interfaces.}
\label{claim:CTC13}
Fix one original base index, parameter and source coefficient
vector in BTH12.  Use its original polynomial kernel map $L$,
the same full Weil datum, and the original finite class presentation.
The finite-average endpoint of FIA applies to every translated
polynomial kernel in (CTC.30), with a positive period and one
averaging list independent of the polynomial.  Linearity therefore
also gives (CTC.30) on the entire original finite coefficient
space; equivalently one may use FIA's universal polynomial
statement directly for each $x$.  BTH3 gives convergence of
(CTC.31) and its compact source representative set.
CTC10 now gives $ZCusp$ of the actual scalar lift.

For clarity, the two hypotheses needed to apply CTC12 can also
be checked here without waiting for the full automorphic
membership theorem.  Rational target invariance follows from
the actual rational Weil identity evaluated at zero, with the
source variable unchanged.  For continuity near $g_0$, restrict
$g$ to a compact neighborhood, choose a finite basis of the
coefficient space, and restrict $h$ to the compact source
representative set.  The finite-basis expansion used in CTC10
is jointly continuous on this compact set, hence uniformly
bounded; it is rationally invariant in $h$.  Dominated convergence
for the finite source quotient measure makes $\Theta_F(g)$
continuous at $g_0$.  Therefore CTC12 proves $IsCuspFun$ of
the original lift.

This supplies the literal isotropic-average, zero-constant-term
and cuspidality portion of the source's
\texttt{Exists\allowbreak Theta\allowbreak Lift\allowbreak Linear\allowbreak Map\allowbreak Cusp\allowbreak Form\allowbreak Multiplicity} construction.
It applies for each original base parameter and uses no independently
selected replacement kernel or scalar normalization.  It does
\emph{not} by itself construct the original smooth cuspidal
multiplicity-space extension. That still needs the original
smoothness and uniform moderate-growth
derivative estimates, lowering and prescribed lowest-type conditions,
the holomorphic realization into the actual Hilbert multiplicity
model. Separately, the original source route still needs
$\mathrm{PIoP}$: one common $\mathrm{ioP}$ satisfying
$\mathrm{IsIoW}$ for every original parameter. This is not the
smooth extension itself. The same-Bundle readings and rational
test identities recorded in BTH12 also remain. Those obligations
must be supplied with their original universal parameter
quantifiers. The present result
removes the cuspidality subargument from that residual task; it
does not remove the entire producer leaf.
% END CTC13




\section{The original column values and the two global slope bounds}
\label{sec:global-slope}

\status{This chapter supplies the three GL4 entries retained in RSL13:
the actual column-value dictionary, the lower slope bound, and the
maximum slope bound with coefficient $4$.  The finite arguments below
are on the original lattices and metrics.  The one additional published
input is Banaszczyk's Euclidean transference theorem, with its hypotheses
checked in GLS8.  The common compatible family and its original second
layer remain inputs.}

\paragraph{Fixed objects and normalization.}
Fix the original record $D$ and a word $u$.  Write
\[
\begin{gathered}
 E=D.\mathrm{host}.E,\quad K=D.\mathrm{what}_2.\mathrm{LL},\quad
 k=N(0)=D.R.k,\quad d_E=[E:\mathbb Q]=2f,\\
 T=(D.\mathrm{mem}\,\allowbreak{} u).A.T,\qquad
 v=D.\mathrm{host}.v=\operatorname{rd}(E)/2.
\end{gathered}
 \tag{GLS.1}
\]
The class set is $\mathcal C=(D.\mathrm{mem}\,\allowbreak{} u).C.\mathrm{cls}(0)$.
Put $S=D.\mathrm{sp}.X(0)$, $L_c=D.\mathrm{clsLat}(u,0,\allowbreak{} c)$, and
$A=(D.\mathrm{fm}\,\allowbreak{} u).\mathrm{lat}(0)$.  All these objects, including
$A$'s integral lattice and all its infinite-place metrics, are fixed.
For the original dominant integral weight $\lambda_\tau$ set
\[
 s_\tau=\lambda_\tau(k-1),\quad
 \nu_{\tau,\allowbreak{} j}=\lambda_\tau(j)-s_\tau,\quad
 d_\tau=\sum_{j=0}^{k-1}\nu_{\tau,\allowbreak{} j},\quad
 a_\tau=\sum_{j=0}^{k-1}\lambda_\tau(j),\quad
 h_\tau=\nu_{\tau,0}.
 \tag{GLS.2}
\]
Thus $a_\tau,\allowbreak{} d_\tau,\allowbreak{} h_\tau$ are respectively the source's
\texttt{col\allowbreak Abs}, \texttt{col\allowbreak Abs\allowbreak P}, and \texttt{first\allowbreak Column\allowbreak Spread}.
The first of these is a signed sum.  The last two are nonnegative.
We use the source's squared Hermitian slope: for a rank-$r$ lattice
$M$ over a totally complex field $K$, a rational basis $e$, local basis
matrices $\Gamma_\mathfrak p$, and positive Gram matrices $H_\sigma$,
\[
 r\mu(M)=-\frac{2}{[K:\mathbb Q]}
       \sum_\mathfrak p\log N\mathfrak p\,
                      v_\mathfrak p(\det\Gamma_\mathfrak p)
       -\frac1{n_\infty(K)}\sum_\sigma\log\det H_\sigma.
 \tag{GLS.3}
\]
Here $n_\infty(K)=[K:\mathbb Q]/2$.  Only finitely many finite-place
terms are nonzero.  Formula (GLS.3) is the determinant/index formula
already used in MT2, with the same factor $2$.

\subsection{The same column-value record, at every place}

% BEGIN GLS1
\paragraph{GLS1: the original class lattices and their scalar extensions.}
\label{claim:GLS1}
Suppose an original \texttt{Cols\allowbreak Tied\allowbreak Sym D} record $\mathcal F$ is
supplied.  Its actual definite space $S$ is standard positive at every
embedding.  Its level-zero class set is nonempty: the identity adelic
class has a preimage under the stored class bijection.  The original
lattice $L_c$ is integral, and
\[
 L_c\subseteq L_c^\vee,\qquad
 L_c^\vee/L_c\simeq\mathcal O_E/T\mathcal O_E,
 \qquad [L_c^\vee:L_c]=T^{2f}.                 \tag{GLS.4}
\]
Here $L_c^\vee=\{y:h_S(L_c,\allowbreak{} y)\subseteq\mathcal O_E\}$ is exactly
the source's \texttt{dual\allowbreak Sub}.  To check the first assertion on the
original lattice, the stored local representative identifies its
completion with the adelic translate of the original integral base
lattice.  A unitary translate preserves inclusion in the Hermitian
dual.  Intersecting these local inclusions with $S$ gives the first
inclusion in (GLS.4).  The stored \texttt{at\allowbreak T.\allowbreak disc} identifies the
actual quotient with the displayed residue module, as in RMP2; its
cardinality is $T^{[E:\mathbb Q]}$.  This proves the whole literal
\texttt{Lat\allowbreak T} condition, including its index.

Choose the source embedding
\[
 \rho_\tau=\mathrm{embE}\circ\mathrm{conj}_E\circ
                   \sigma_{\mathrm{eTau}(\tau)}:E\longrightarrow K.
 \tag{GLS.5}
\]
For every $\sigma:K\hookrightarrow\mathbb C$ it satisfies
$\sigma\rho_\tau=\overline{\mathrm{embS}(D,\sigma,\tau)}$.
Let $G$ be the Hermitian Gram matrix of the original frame of $S$,
reindexed by $\mathrm{Fin}(k)$.  Define $Q_{c,\tau}$ on $K^k$ by
\[
 \operatorname{bc}_\tau(z)=\rho_\tau(\operatorname{coord}_S z),
 \qquad Q_{c,\tau}.L=
       \mathcal O_K\operatorname{-span}\operatorname{bc}_\tau(L_c),
 \qquad H_{c,\tau,\sigma}=\sigma\rho_\tau(G).
 \tag{GLS.6}
\]
Choose an invertible complex matrix $X_{c,\tau,\sigma}$ with
$X_{c,\tau,\sigma}^{-*}X_{c,\tau,\sigma}^{-1}
=H_{c,\tau,\sigma}$.  The positive Hermitian square root gives one.
The original full lattice $L_c$ has finitely many generators; their
images generate the lattice in (GLS.6).  Its $K$-span is $K^k$ because
the original rational frame is in the $E$-span of $L_c$.  Thus each $Q$
is good.  The coordinate map is additive and $\rho_\tau$-semilinear,
and at every infinite place
\[
 \|\operatorname{bc}_\tau(z)\|_{Q,\sigma}^{2}
       =\operatorname{Re}\sigma\rho_\tau(h_S(z,\allowbreak{} z)).
 \tag{GLS.7}
\]
This constructs all fields of an actual \texttt{Col\allowbreak Base D u} with
\texttt{rho = rho\allowbreak C D}.  In particular it uses the original $L_c$,
and (GLS.7) holds at every coefficient embedding.
% END GLS1

% BEGIN GLS2
\paragraph{GLS2: the coefficient map and the exact local criterion.}
\label{claim:GLS2}
Assume additionally \texttt{adm3 D}; its assertion at $u$ says that
all original level-zero stabilizers are trivial.  Project
$\mathcal F.\mathrm{forms}(u)$ to its original \texttt{Forms0} record.
The rational-value dictionary and polynomial injectivity give a
$K$-linear map
\[
 \operatorname{co}_0:A.V\longrightarrow
       \bigoplus_{c\in\mathcal C}W_K,\qquad
 W_K=\bigotimes_\tau \Delta_{\nu_\tau}\otimes_{\mathbb Z}K.
 \tag{GLS.8}
\]
The polynomial variables for different $\tau$ are disjoint.  The
integral lattice $\Delta_\nu$ is the divided-power Weyl lattice
generated by its actual highest vector; it is the literal
\texttt{del\allowbreak Sq}, rather than an arbitrary full polynomial lattice.
The value of $\operatorname{co}_0(x)_c$ on the original frame chart is
the original rational form $x$ on class $c$.  Vanishing of all these
polynomials implies vanishing of the form.  Conversely, any tuple in
$W_K$ gives an allowed function on each class: the stabilizer
equivariance condition is empty, and the stored rational-span equality
puts that tuple in $A.V$.  Polynomial uniqueness then gives (GLS.8)
as an isomorphism.  This is also the original rational surjectivity
used in GLE4 and RSL9.

Use exactly the source action
\[
 \mathcal R(g)P(Z)=\left(\prod_\tau\det(g_\tau)^{s_\tau}\right)
                              P((g_\tau^{t}Z_\tau)_\tau).
 \tag{GLS.9}
\]
Set $g_{c,\tau}=\rho_\tau(G)^{-1}$ for the coordinates (GLS.6), and
$\operatorname{co}(x)_c=\mathcal R(g_c)\operatorname{co}_0(x)_c$.
This is still an isomorphism onto the same polynomial space.  For a
general coordinate basis of $Q$ the formula is
$g_{c,\tau}=M_{c,\tau}\rho_\tau(G)^{-1}$, where $M$ is the matrix
of the original map $\operatorname{bc}$; this is precisely
\texttt{g\allowbreak C} in the source.

Fix an arbitrary coefficient place $w$, not just one over $T$.
Choose an $E$-rational basis matrix $Y_c$ for $L_c$ simultaneously at
the finitely many places over the rational residue characteristic of
$w$.  Such a basis exists by local freeness over the Dedekind domain
and the Chinese remainder theorem: prescribe a basis modulo each of
these primes, choose the columns in $L_c$, and use Nakayama at each
prime.  Put $\gamma_{c,\tau}=\rho_\tau(Y_c)$ in (GLS.6), or
$M_{c,\tau}\rho_\tau(Y_c)$ in the general coordinates.  These are
local bases of the exact scalar-extension lattice $Q_{c,\tau}.L$.
The corresponding original Hermitian dual basis has matrix
$(Y_c^*G)^{-1}$.  At the distinguished coefficient embedding,
\[
 \left(\mathrm{embLL}
             (\gamma_{c,\tau}^{-1}g_{c,\tau})\right)^t
       =\left(\mathrm{embT}_\tau(Y_c^*G)\right)^{-1}.
 \tag{GLS.10}
\]
Indeed $\gamma^{-1}g=\rho_\tau(GY_c)^{-1}$, and (GLS.5)
turns the transpose of its image into the inverse of $Y_c^*G$.
Substitute (GLS.10) in the original \texttt{fn\allowbreak Of} formula.
Its determinant power is exactly $s_\tau$, so both its polynomial
substitution and its scalar determinant factor agree with (GLS.9).
The original local integral-value criterion of GLE4 consequently says
\[
 x\in A.L_w\quad\Longleftrightarrow\quad
 \mathcal R(\gamma_c^{-1})\operatorname{co}(x)_c
       \in\bigotimes_\tau\Delta_{\nu_\tau,\mathcal O_w}
       \quad\hbox{for every }c.                  \tag{GLS.11}
\]
The right side is the source's \texttt{Int\allowbreak Span\allowbreak At}.  GLE3 proves its
invariance under every integral invertible change of basis, including
the signed determinant factor.  Any other local basis of the same
$Q_{c,\tau}.L$ differs from $\gamma_{c,\tau}$ by such a change.
Thus (GLS.11) holds for \emph{every} family of local bases.  This
checks the full quantifier in \texttt{Col\allowbreak Val.\allowbreak fin}.
% END GLS2

% BEGIN GLS3
\paragraph{GLS3: the original all-place norm, and a literal ColVal.}
\label{claim:GLS3}
At an infinite place $\sigma$ put
\[
 Z_{c,\tau}=X_{c,\tau,\sigma}^{-1}\,
                         \sigma(g_{c,\tau}).
 \tag{GLS.12}
\]
In general coordinates insert $M$ as in GLS2.  The identity
$(X^{-1}\sigma M)^*(X^{-1}\sigma M)=\sigma\rho_\tau(G)$,
together with (GLS.5), gives
\[
 (Z_{c,\tau}^{t})^*\,
      \mathrm{embS}(D,\sigma,\tau)(G)\,\allowbreak{} Z_{c,\tau}^{t}=1.
 \tag{GLS.13}
\]
This follows by multiplying out $Z=Y(\sigma\rho_\tau G)^{-1}$
with $Y^*Y=\sigma\rho_\tau G$ and transposing.  Thus $Z_c^t$ is
an orthonormal frame for the original definite form space at this
embedding.

Every class weight in the original norm formula is
$1/|\mathrm{stab}(0,\allowbreak{} c)|=1$.  To allow the orthonormal frame in
(GLS.13) to depend on $c$, use (GLS.8) to take the rational form
supported on that single class.  Its original \texttt{norm0} formula
shows that this class's contribution is independent of the chosen
orthonormal frame.  Apply this separately to each class and add.
The determinant term in \texttt{fn\allowbreak Of} is retained throughout;
$\det(Z^t)=\det Z$ and (GLS.9) includes its exact signed power.
One obtains
\[
 \|x\|_{A,\sigma}^{2}
 =\sum_{c\in\mathcal C}
  \frac{\left\|
    \mathcal R(X_{c,\sigma}^{-1})
                \sigma(\operatorname{co}(x)_c)\right\|_F^2}
       {\prod_\tau\|H_{\nu_\tau}\|_F^2}.        \tag{GLS.14}
\]
Here the Fock squared norm is $\|P\|_F^2=\sum_\alpha
|P_\alpha|^2\alpha!$, with the tensor convention on disjoint
blocks.  The equality follows from
$\mathcal R(X^{-1})\sigma(\operatorname{co}x)
=\mathcal R(Z)\sigma(\operatorname{co}_0x)$.

GLS1 supplies the nonempty class set, the exact \texttt{Lat\allowbreak T}, and
all scalar-extension fields; GLS2 supplies the rational map, its
range, and its local criterion; (GLS.14) supplies the last metric
field.  These are every field of the literal \texttt{Col\allowbreak Val D u}.
We have therefore proved, pointwise on the original record,
\[
 \texttt{Cols\allowbreak Tied\allowbreak Sym}(D)\ \wedge\ \texttt{adm3}(D)
       \Longrightarrow\ \forall u,\
                          \texttt{Nonempty} (\texttt{Col\allowbreak Val}(D,\allowbreak{} u)).
 \tag{GLS.15}
\]
This is the exact \texttt{Column\allowbreak Valuation\allowbreak From\allowbreak Twist\allowbreak Compatibility}
contract.  No column values, class representatives, or integral
lattice were reselected independently of the supplied forms.
% END GLS3

\subsection{Two finite polynomial estimates with their normalization}

% BEGIN GLS4
\paragraph{GLS4: the highest-vector norm and the hook bound.}
\label{claim:GLS4}
For a partition $\nu_1\ge\cdots\ge\nu_k\ge0$ use one-based indices
in this paragraph.  Put
$H_\nu=\prod_{j=1}^k\det(Z_{[j],[j]})^{\nu_j-\nu_{j+1}}$,
where $\nu_{k+1}=0$, and let $N_\nu=F(H_\nu,\allowbreak{} H_\nu)$ for the
bilinear rational Fock form.  Then
\[
 N_\nu=
 \frac{\prod_{i=1}^k(\nu_i+k-i)!}
      {\prod_{1\le i<j\le k}(\nu_i-\nu_j+j-i)}.
 \tag{GLS.16}
\]
We give the finite differential calculation behind the formula.
Multiplication by $z_{as}$ is Fock-adjoint to $\partial_{as}$, as
is seen on two monomials.  Set $E_{ab}=\sum_s z_{as}\partial_{bs}$.
The following normal-order identity holds, with factors in increasing
column order:
\[
 \sum_{\pi\in\mathfrak S_k}\operatorname{sgn}(\pi)
 \prod_{j=1}^k\bigl(E_{\pi(j),\allowbreak{} j}+(k-j)\delta_{\pi(j),\allowbreak{} j}\bigr)
           =\det Z\,\det\partial.              \tag{GLS.17}
\]
For completeness, this can be verified without a representation
theorem.  For lists $w,\allowbreak{} g$ of length $r$ let
\[
 P_r(w,\allowbreak{} g)=\sum_{s_1,\ldots,\allowbreak{} s_r}
 z_{w_1s_1}\cdots z_{w_rs_r}
 \partial_{g_1s_1}\cdots\partial_{g_rs_r},\qquad
 L_r(w,\allowbreak{} g)=\prod_{j=1}^r
       (E_{w_jg_j}+(r-j)\delta_{w_jg_j}).
\]
The relation $\partial_{bs}z_{at}=z_{at}\partial_{bs}
+\delta_{ab}\delta_{st}$ gives
\[
 P_r=E_{w_1g_1}P_{r-1}(w',\allowbreak{} g')
   -\sum_{j=2}^r\delta_{w_jg_1}
                   P_{r-1}(w'[j\leftarrow w_1],\allowbreak{} g').
 \tag{GLS.18}
\]
Subtract this recurrence from that for $L_r$.  Besides
$(E_{w_1g_1}+(r-1)\delta_{w_1g_1})(L_{r-1}-P_{r-1})$, the
difference has $(r-1)\delta_{w_1g_1}P_{r-1}$ and the displayed
sum with positive sign.  Antisymmetrize in $w$.  The first term
vanishes by induction.  For each $j$, the transposition of positions
$1,\allowbreak{} j$ changes the sign and identifies the $j$th term with
$\delta_{w_1g_1}P_{r-1}$; the remaining $r-1$ terms cancel it.
Thus the antisymmetrizations of $L_r$ and $P_r$ agree.  When
$w=g=(1,\ldots,\allowbreak{} k)$, terms with repeated $s_j$ vanish in the
determinant, and the remaining permutations give
$\det Z\det\partial$.  This proves (GLS.17).

On $H_{\nu+\mathbf1}$, $E_{ii}$ acts by $\nu_i+1$, and
$E_{ij}$ kills it when $i<j$; both statements follow directly from
the leading-minor product.  In a nonidentity permutation term of
(GLS.17), take the largest moved column $j$.  Then $\pi(j)<j$;
all factors to its right are diagonal, so that term kills the
highest vector.  Consequently (GLS.17), cancellation of the nonzero
polynomial $\det Z$, and Fock adjointness give
\[
 N_{\nu+\mathbf1}
          =\prod_{i=1}^k(\nu_i+k-i+1)\,\allowbreak{} N_\nu.
 \tag{GLS.19}
\]
If $\nu_k=0$, $H_\nu$ uses only the first $k-1$ rows and columns;
its Fock norm is unchanged on deleting the last row and column.
These two recurrences prove (GLS.16), by induction on $k$ and then
on $\nu_k$: the right side of (GLS.16) satisfies exactly the same
recurrences and is $1$ in dimension zero.

The coefficient of $\prod_i z_{ii}^{\nu_i}$ in $H_\nu$ is $1$.
Its Fock contribution alone is $\prod_i\nu_i!$.  Thus, with the
source normalization $c_\nu=N_\nu/\prod_i\nu_i!$,
\[
 1\le c_\nu\le(\nu_1+k)^{k(k-1)/2},\qquad
 \log c_\nu\le k^2\log(\nu_1+k).             \tag{GLS.20}
\]
For the upper bound, each positive integer factor in the denominator
of (GLS.16) is at least one, and
$(\nu_i+k-i)!/\nu_i!\le(\nu_1+k)^{k-i}$.
The last logarithmic inequality uses $\nu_1+k\ge1$.
This proves the literal \texttt{Hook\allowbreak Bound}, with its exact Fock
normalization.
% END GLS4

% BEGIN GLS5
\paragraph{GLS5: the normalized integral Weyl slope.}
\label{claim:GLS5}
If $k=0$, the only polynomial is a constant:
$\Delta=\nabla=\mathbb Z\cdot1$, its normalized Gram matrix is
$(1)$, $c_\nu=1$, and both the slope and the required lower bound
are zero.  Thus suppose $k\ge1$, put $d=\sum_j\nu_j$, and let
$r=\dim_\mathbb Q W_\nu>0$.
Let $\nabla_\nu$ be the integral coefficient lattice in the same
polynomial space.  IW5 gives, for the normalized partition, the
exact duality
\[
 \Delta_\nu^{\#}=\nabla_\nu,\qquad
 B_\nu=F/N_\nu.                              \tag{GLS.21}
\]
Here is the extension to every partition, including $\nu_k>0$.
Put $s=\nu_k$ and $\nu^0=\nu-s\mathbf1$. Multiplication by
$\det(Z)^s$ identifies the two rational highest-weight spaces,
and it maps $\Delta_{\nu^0}$ onto $\Delta_\nu$: every
off-diagonal root operator annihilates the determinant factor,
so it commutes with this multiplication, including its divided
powers.  It also maps $\nabla_{\nu^0}$ onto $\nabla_\nu$.
Indeed the rational spaces already match, and the polynomial
$\det(Z)^s$ is primitive; the DVR Gauss lemma shows that a rational
polynomial whose product with this factor has integral coefficients
itself has integral coefficients.  This also transports the
integral basis of CB2.

The pullback of $B_\nu$ has the same transpose-adjoint root
operators as $B_{\nu^0}$ and gives the highest vector norm $1$.
These properties determine the form: move each divided root word
in the first entry into the second by adjointness, and then take
the highest-weight component, which is a scalar multiple of the
highest vector.  Root words span the rational space, so all
pairings are determined.  Thus the pullback is exactly $B_{\nu^0}$,
and the two lattice identities transfer to the arbitrary partition.
This proves (GLS.21) in its full stated scope; it does not assume
that unnormalized Fock norms are unchanged by determinant
multiplication.

IW5 proves the normalized case at every DVR.  It holds over $\mathbb Z$ because
membership in a full rational lattice is equivalent to membership
at every prime.  In integral bases, the matrix of the pairing between
the two lattices is unimodular.  It follows that their positive Gram
determinants are reciprocal.

Here is the requisite bound on the coefficient lattice's determinant.
Every monomial $Z^\alpha$ occurring in it has column sums $\nu_j$.
The multinomial theorem gives
\[
 \frac{\prod_j\nu_j!}{\alpha!}\le k^d,
 \qquad \frac{\alpha!}{N_\nu}\ge
                        (k^d c_\nu)^{-1}=q_0.\tag{GLS.22}
\]
Choose a finite set of monomials containing all supports of an
integral basis of $\nabla_\nu$, and let $V$ be its integer
coefficient matrix.  Its normalized Gram matrix is
$V^t\operatorname{diag}(\alpha!/N_\nu)V$.
Cauchy--Binet expresses its determinant as the sum of squared
integer $r$-minors of $V$, weighted by products of $r$ diagonal
entries.  Some minor is nonzero, since the basis is independent.
Every weight is at least $q_0^r$, so the determinant is at least
$q_0^r$.  By (GLS.21), the Gram determinant on $\Delta_\nu$ is at
most $q_0^{-r}$.  Consequently
\[
 \mu_{\rm std}(\nu)
 :=-\frac1r\log\det\operatorname{Gram}_{B_\nu}(\Delta_\nu)
 \ \ge\ -d\log k-\log c_\nu.                 \tag{GLS.23}
\]
This is the source's \texttt{Std\allowbreak Slope\allowbreak Low}.  It also covers $d=0$:
the polynomial lattice is the one-dimensional lattice generated by
$1$, and both sides are zero.  All determinants here use the literal
divided-power lattice and the normalization of (GLS.14).
% END GLS5

\subsection{The lower slope on an arbitrary original ColVal}

% BEGIN GLS6
\paragraph{GLS6: exact slope of every column lattice in a supplied ColVal.}
\label{claim:GLS6}
Now take \emph{any} $\mathrm{CV}:\texttt{Col\allowbreak Val D u}$, without
assuming it was constructed by GLS1.  Its semilinear map extends to
a $K$-linear map $K\otimes_{E,\rho_\tau}S\to Q_{c,\tau}.V$.
The image contains the full lattice \texttt{CV.\allowbreak bc\_\allowbreak L}, so it is
surjective; both dimensions are $k$, so it is an isomorphism.
The metric condition \texttt{bc\_\allowbreak norm} identifies all diagonal
Hermitian values.  Polarization identifies the entire form: use
$z+z'$ for its real part and $z+az'$ for an $a\in E$ with nonreal
image to recover its imaginary part.  Such an $a$ exists because
$E$ is CM.  Every embedding of $E$ extends to $K$; hence positivity
of the good lattice $Q$ also proves positivity of $S$ at every
embedding.

The normalized slope is preserved by this scalar extension.  To
see every normalization, in a local basis of a full lattice a
determinant valuation at $\mathfrak p$ is multiplied by
$e(\mathfrak P/\mathfrak p)$ upstairs.  In (GLS.3) the sum
$\sum_{\mathfrak P\mid\mathfrak p}ef=[K:\rho_\tau(E)]$
cancels the increase of the field degree.  At infinity, each
complex place of $E$ has exactly $[K:E]$ places above it; a conjugate
choice has the same positive Gram determinant.  Thus the normalized
average in (GLS.3) is unchanged.  This argument applies to the
literal span \texttt{bc\_\allowbreak L}; no finite index is introduced.

For a positive Hermitian $E$-lattice $L$, the same determinant
formula gives $\mu(L^\vee)=-\mu(L)$: the dual ideal determinants
are inverse conjugates and the Gram matrices are inverse
transposes.  If $L\subseteq L^\vee$ has index $q$, the finite index
formula gives
\[
 \mu(L)=\mu(L^\vee)-\frac{2\log q}{d_E k},
 \qquad \mu(L)=-\frac{\log q}{d_E k}.          \tag{GLS.24}
\]
The \texttt{lat\allowbreak T} field of this very $\mathrm{CV}$ has
$q=T^{2f}$ and $d_E=2f$.  Scalar extension therefore gives
\[
             \boxed{\ \mu(Q_{c,\tau})=-\frac{\log T}{k}\ }
                  \qquad(c\in\mathcal C,\ \tau\in\mathrm{Tau}).
 \tag{GLS.25}
\]
This proves the full \texttt{QSlope} contract for an arbitrary
\texttt{Col\allowbreak Val}, including its given embeddings and metrics.
% END GLS6

% BEGIN GLS7
\paragraph{GLS7: the exact polynomial determinant and the lower bound.}
\label{claim:GLS7}
Let $r_\tau=\operatorname{rank}_{\mathbb Z}\Delta_{\nu_\tau}$ and
$r=\prod_\tau r_\tau$.  Choose the tensor integral basis $Z$ of
$\bigotimes_\tau\Delta_{\nu_\tau}$ and let $G_0$ be its normalized
Fock Gram matrix.  Disjoint variables make this matrix a Kronecker
product, so
\[
 -\frac1r\log\det G_0=\sum_\tau\mu_{\rm std}(\nu_\tau).
 \tag{GLS.26}
\]
For the representation (GLS.9), its determinant satisfies
\[
 \det\mathcal R(g)^k=\prod_\tau\det(g_\tau)^{r a_\tau}.
 \tag{GLS.27}
\]
Indeed the integral polynomial lattice is graded by row weights.
Permutation matrices permute the rows, so the sum of the weights
of any one row over a homogeneous basis is $rd_\tau/k$.
The determinant twist adds $rs_\tau$, making the exponent
$ra_\tau/k$.  This proves (GLS.27) on diagonal matrices.
Elementary unipotents act unipotently and have determinant one;
diagonal matrices and these elementary matrices generate
$\mathrm{GL}_k$.  Thus the identity holds for every invertible
matrix, and its integer-power form avoids any root choices.

Use the rational basis $\operatorname{co}^{-1}(\mathbf1_cZ_i)$
of the original $A.V$.  At a finite place its integral basis is
$\operatorname{co}^{-1}(\mathbf1_c\mathcal R(\gamma_c)Z_i)$ by
\texttt{CV.\allowbreak fin}; at infinity its Gram blocks are
$R_{c,\sigma}^*G_0R_{c,\sigma}$ by \texttt{CV.\allowbreak inf}, with
$R_{c,\sigma}=\mathcal R(X_{c,\sigma}^{-1})$ in the basis $Z$.
Substitute these determinants and (GLS.27) into (GLS.3).  The finite
and infinite contributions belonging to a column combine into
$ra_\tau\mu(Q_{c,\tau})$.  Therefore
\[
 r|\mathcal C|\mu(A)
 =-|\mathcal C|\log\det G_0+
                 \sum_{c,\tau}ra_\tau\mu(Q_{c,\tau}).
 \tag{GLS.28}
\]
Here $r>0$ and $|\mathcal C|>0$; the latter is a field of
\texttt{Col\allowbreak Val}.  Divide (GLS.28), and apply GLS6 and (GLS.26):
\[
 \mu(A)=-\sum_\tau\frac{a_\tau}{k}\log T
                       +\sum_\tau\mu_{\rm std}(\nu_\tau).
 \tag{GLS.29}
\]
Using GLS4--GLS5 now proves exactly
\[
 \boxed{\quad
 -\sum_\tau\frac{a_\tau}{k}\log T
 -\sum_\tau\bigl(d_\tau\log k+k^2\log(h_\tau+k)\bigr)
                  \ \le\ \mu(A).
 \quad}                                                       \tag{GLS.30}
\]
This is \texttt{Slope\allowbreak Low D u}, for every original $D,\allowbreak{} u$ and every
$\mathrm{CV}:\texttt{Col\allowbreak Val D u}$.  No size or eventuality
hypothesis was used.
% END GLS7

\subsection{The maximum slope with the exact coefficient four}

% BEGIN GLS8
\paragraph{GLS8: a published transference theorem on the exact trace lattice.}
\label{claim:GLS8}
For each original $L_c$ occurring in any \texttt{Col\allowbreak Val}, there are
$k$ $E$-linearly independent elements $y_1,\ldots,\allowbreak{} y_k\in L_c^\vee$
such that
\[
 \frac{\operatorname{Tr}_{E/\mathbb Q}h_S(y_a,\allowbreak{} y_a)}{2f}
                    \le k^2(2v)^2 \qquad(1\le a\le k).
 \tag{GLS.31}
\]
The published input is W.~Banaszczyk, \emph{New bounds in some
transference theorems in the geometry of numbers}, Math. Ann.
\textbf{296} (1993), 625--635, Theorem~2.1, p.~631,
\url{https://doi.org/10.1007/BF01445125}.
It states that a full Euclidean lattice $M\subset\mathbb R^n$
and its integer-pairing dual satisfy
$\lambda_i(M)\lambda_{n+1-i}(M^*)\le n$.
The case $n=1$ is equality.  This is the only result imported here
instead of being proved in this chapter.

Here is the complete matching of hypotheses and constants.
Use $M=L_c^\vee$, as a lattice of rank $n=d_E k$ in
$S\otimes_\mathbb Q\mathbb R$, with positive Euclidean form
\[
                  B(x,\allowbreak{} y)=\operatorname{Tr}_{E/\mathbb Q}h_S(x,\allowbreak{} y).
 \tag{GLS.32}
\]
It is real and symmetric because trace commutes with CM conjugation;
it is positive definite by GLS6.  The $\mathcal O_E$-dual of a full
projective lattice is full and finitely generated, so its Minkowski
image is discrete and spans this real vector space.  The integer
pairing dual in the theorem is exactly
\[
                  M^*=\mathfrak D_E^{-1}L_c.    \tag{GLS.33}
\]
Indeed $\operatorname{Tr}h(x,\allowbreak{} M)\subseteq\mathbb Z$ is equivalent
to $h(x,\allowbreak{} M)\subseteq\mathfrak D_E^{-1}$, since $M$ is
$\mathcal O_E$-stable.  Local duality in a basis, followed by
intersection over the primes, gives (GLS.33).

The original integrality $h(L_c,\allowbreak{} L_c)\subseteq\mathcal O_E$ and
$\overline{\mathfrak D_E}=\mathfrak D_E$ imply, for $0\ne z\in M^*$,
that $\beta=h_S(z,\allowbreak{} z)$ is totally positive and belongs to
$\mathfrak D_E^{-2}$.  Its principal fractional ideal multiplied by
$\mathfrak D_E^2$ is a nonzero integral ideal.  Hence
$|N_{E/\mathbb Q}\beta|\,|\operatorname{disc}E|^2\ge1$.
Arithmetic--geometric mean over the $d_E$ embeddings gives
\[
 B(z,\allowbreak{} z)\ge d_E/\operatorname{rd}(E)^2,
 \qquad \lambda_1(M^*)\ge\sqrt{d_E}/\operatorname{rd}(E).
 \tag{GLS.34}
\]
Apply the quoted theorem with $i=n$.  It gives $n$ independent
vectors in the original $M$ with squared norm at most
$d_Ek^2\operatorname{rd}(E)^2$; minima are attained because $M$
is discrete.  They span over $\mathbb Q$, hence over $E$.  Extract
an $E$-basis from them, preserving the bound, and divide it by
$d_E=2f$.  Since $\operatorname{rd}(E)=2v$, this is (GLS.31).
Thus the source's \texttt{Short\allowbreak At} witnesses are supplied in the
original dual lattice with the exact constant.
% END GLS8

% BEGIN GLS9
\paragraph{GLS9: coefficient ideals and the signed determinant correction.}
\label{claim:GLS9}
Fix a class $c$, its short dual basis from GLS8, and an arbitrary
$\mathrm{CV}$.  Work in the given rational coordinate basis of
$Q_{c,\tau}$.  The semilinear extension in GLS6 defines a unique
$K$-linear functional $\ell_{a,\tau}$ by
\[
 \ell_{a,\tau}(\operatorname{bc}_\tau z)
                         =\rho_\tau(h_S(y_a,\allowbreak{} z)).
 \tag{GLS.35}
\]
Let $g_\tau$ have these functionals as rows.  Nondegeneracy of $h_S$
and independence of the $y_a$ make $g_\tau$ invertible.  Every row
is integral on $Q_{c,\tau}.L$: this follows first on
$\operatorname{bc}_\tau(L_c)$ from $y_a\in L_c^\vee$ and
Hermitian symmetry, then on
its exact $\mathcal O_K$-span.  Define $\phi_c(x)$ to be the full
list of monomial coefficients of
\[
     P_c((g_\tau^{t}Z_\tau)_\tau),\qquad
                         P_c=\operatorname{co}(x)_c.
 \tag{GLS.36}
\]
All monomials have column multidegrees $\nu_\tau$; their finite
index set is exactly the source's \texttt{Mon\allowbreak I}.  Invertible
substitution and the range in (GLS.8) make this coefficient map
injective on the subspace supported at $c$.

Let $M_\tau$ be the actual coordinate lattice of $Q_{c,\tau}.L$,
and let $\det M_\tau$ denote its determinant fractional ideal in
this basis.  Set
\[
 I_c=\prod_\tau(\det M_\tau)^{s_\tau},\qquad
 D_{c,\sigma}=\prod_\tau|\det X_{c,\tau,\sigma}|^{2s_\tau}>0.
 \tag{GLS.37}
\]
Negative $s_\tau$ are allowed.  At a finite place choose local
matrices $\gamma_\tau$ for $M_\tau$.  The coefficients of
$\mathcal R(\gamma^{-1})P_c$ are integral by
\texttt{CV.\allowbreak fin}, and $g_\tau\gamma_\tau$ has integral entries
by (GLS.35).  Substitution of $(g\gamma)^tZ$ shows that
\[
 \left(\prod_\tau\det\gamma_\tau^{-s_\tau}\right)
                   \phi_c(x)_\alpha\in\mathcal O_{K,\allowbreak{} w}
                \quad(x\in A.L\hbox{ supported at }c).
 \tag{GLS.38}
\]
Intersecting at all finite places gives
$\phi_c(x)_\alpha\in I_c$.  This is exactly the source's
\texttt{tw\allowbreak Id}/\texttt{tw\allowbreak Ideal}: their definition intersects the
conditions (GLS.38) over every local basis.  Changes of local basis
have unit determinant and leave the same ideal.

Write $\mu_0(M_\tau)$ for the slope with the standard coordinate
metric.  Formula (GLS.3) yields both identities
\[
 \begin{split}
 -\frac{2}{[K:\mathbb Q]}\log N I_c
       &=\sum_\tau s_\tau k\mu_0(M_\tau),\\
 k\mu(Q_{c,\tau})
       &=k\mu_0(M_\tau)+\frac1{n_\infty(K)}
                            \sum_\sigma\log|\det X_{c,\tau,\sigma}|^2.
 \end{split}                                                   \tag{GLS.39}
\]
For the first equality the local exponent of $I_c$ is
$\sum_\tau s_\tau v_w(\det\gamma_\tau)$.
For the second, the coordinate Gram matrix is $X^{-*}X^{-1}$,
whose determinant is $|\det X|^{-2}$.
Combining these with GLS6 gives the exact cancellation
\[
 -\frac{2}{[K:\mathbb Q]}\log N I_c
 +\frac1{n_\infty(K)}\sum_\sigma\log D_{c,\sigma}
                       =-\sum_\tau s_\tau\log T.
 \tag{GLS.40}
\]
This computation keeps the original determinant shift and all
finite and infinite normalizations on the same $Q$.
% END GLS9

% BEGIN GLS10
\paragraph{GLS10: the class map's exact all-place norm estimate.}
\label{claim:GLS10}
Give the coefficient in monomial $\alpha$ the squared weight
\[
 q_\alpha=\frac{\alpha!}{\prod_\tau N_{\nu_\tau}},
 \qquad q_0=\left(k^{\sum_\tau d_\tau}
                                \prod_\tau c_{\nu_\tau}\right)^{-1}.
 \tag{GLS.41}
\]
The separate-column multinomial estimate of GLS5 proves
$q_\alpha\ge q_0>0$.  At an infinite place $\sigma$ let
\[
 G_{\tau,\sigma}=\sum_{a=1}^k
       \operatorname{Re}\sigma\rho_\tau(h_S(y_a,\allowbreak{} y_a))>0.
 \tag{GLS.42}
\]
The Riesz functional identity (GLS.35), in the orthonormal coordinates
given by $X_{c,\tau,\sigma}$, gives
\[
 \| (\sigma(g_\tau)X_{c,\tau,\sigma})^t\|_{\rm HS}^2
                                                   =G_{\tau,\sigma}.
 \tag{GLS.43}
\]
It follows by expanding $h_S$ in a basis and using the polarized
version of \texttt{bc\_\allowbreak norm}; the row of the dual functional is
the conjugate of the orthonormal coordinate vector of
$\operatorname{bc}_\tau(y_a)$.

For a polynomial of total degree $d_\tau$ in block $\tau$,
substitution by a matrix $U_\tau$ has Fock norm at most
$\|U_\tau\|_{\rm op}^{d_\tau}$ times the original norm.
One can check this directly by identifying a monomial of degree $d$
with its symmetric tensor: its Fock squared norm is $\alpha!$,
and the symmetric-tensor operator is the restriction of
$U^{\otimes d}$, whose norm is at most $\|U\|_{\rm op}^d$.
The normalization in this identification depends only on $d$ and
cancels.  Apply this separately to the blocks, and use
$\|U\|_{\rm op}\le\|U\|_{\rm HS}$.

In (GLS.36) write
$P=\mathcal R(X^{-1})\sigma(\operatorname{co}(x)_c)$.
Then (GLS.9) gives, including its signed determinant shift,
\[
 \sigma(\operatorname{co}(x)_c)((\sigma g)^tZ)
     =\left(\prod_\tau\det X_{c,\tau,\sigma}^{s_\tau}\right)
                     P(((\sigma g_\tau X_{c,\tau,\sigma})^tZ_\tau)_\tau).
\]
The $c$-summand of \texttt{CV.\allowbreak inf} is bounded by the whole norm.
Hence the coefficient map satisfies, on the original $A.V$,
\[
 \sum_\alpha q_\alpha|\sigma(\phi_c(x)_\alpha)|^2
 \le B_{c,\sigma}\|x\|_{A,\sigma}^2,
 \qquad B_{c,\sigma}=D_{c,\sigma}
                                \prod_\tau G_{\tau,\sigma}^{d_\tau}.
 \tag{GLS.44}
\]

All embeddings of $E$ extend equally often to $K$.  A conjugate
embedding leaves the real diagonal values in (GLS.42) unchanged.
Thus their average over $K$-places is the corresponding average
over $E$-places.  Since the full trace is twice the sum over the
complex places, (GLS.31) gives
\[
 \frac1{n_\infty(K)}\sum_\sigma G_{\tau,\sigma}
      =\sum_a\frac{\operatorname{Tr}h_S(y_a,\allowbreak{} y_a)}{2f}
      \le k^3(2v)^2.
\]
Concavity of $\log$ (or summing $\log t\le t-1$ after division
by the mean) gives
\[
 \frac1{n_\infty(K)}\sum_\sigma\log G_{\tau,\sigma}
                             \le\log(k^3(2v)^2).
 \tag{GLS.45}
\]
Equations (GLS.38), (GLS.40), (GLS.44), and (GLS.45) supply all
fields of the original \texttt{Class\allowbreak Map\allowbreak T} and \texttt{Class\allowbreak Map}.
% END GLS10

% BEGIN GLS11
\paragraph{GLS11: the maximum slope and the exact constant $4$.}
\label{claim:GLS11}
We first record the elementary slope comparison used for (GLS.44).
If a full lattice $N$ maps injectively to coefficient vectors in
$I^J$, with diagonal squared metric weights at least $q_0>0$, and
the squared norm of the map at $\sigma$ is at most $B_\sigma$, then
\[
 \mu_{\max}(N)\le-\log q_0
       -\frac{2\log NI}{[K:\mathbb Q]}
       +\frac1{n_\infty(K)}\sum_\sigma\log B_\sigma.
 \tag{GLS.46}
\]
Here is a determinant proof valid for every nonzero sublattice.
Write a rank-$r$ sublattice as $\bigoplus_j\mathfrak a_jz_j$ in
a pseudo-basis and choose a nonzero $r$-minor $\beta$ of its
coefficient vectors.  Integrality gives
$\beta\prod_j\mathfrak a_j\subseteq I^r$, whence
$|N\beta|\prod_jN\mathfrak a_j\ge(NI)^r$.
Cauchy--Binet bounds its coefficient Gram determinant below by
$q_0^r|\sigma\beta|^2$.  The product formula and (GLS.3) imply
that its coefficient-space slope is at most
$-\log q_0-2\log NI/[K:\mathbb Q]$.
Taking the $r$th exterior power in the norm bound gives the last
term in (GLS.46).  Taking the supremum over sublattices proves the
stated maximum-slope inequality.

The coordinate projection to a single class preserves $A.L$:
apply \texttt{CV.\allowbreak fin}, set the other class polynomials to zero,
and intersect the local integral lattices.  The projections sum to
the identity, and \texttt{CV.\allowbreak inf}, polarized, makes their ranges
orthogonal at every infinite place.  Thus $A.L$ is the actual
orthogonal direct sum of its class lattices $A_c$.
Its maximum slope is at most the maximum of theirs.  For example,
filter any sublattice by successive class projections; its successive
quotients inject into the individual $A_c$, with quotient norm at
least the projected norm.  Additivity of determinant degrees and
the same exterior-power comparison used for (GLS.46) give the
claimed bound.  The original class set is finite and nonempty.

Apply (GLS.46) to each $A_c$ with GLS9--GLS10.  Since
$-\log q_0=\sum_\tau(d_\tau\log k+\log c_{\nu_\tau})$,
the result before weakening the signed shift is
\[
 \mu_{\max}(A)\le
  \sum_\tau\left(d_\tau\log(4k^4v^2)
                    +\log c_{\nu_\tau}-s_\tau\log T\right).
 \tag{GLS.47}
\]
The original host bounds give $\operatorname{rd}(E)^2\ge2f$
by AH16.  Also $e\ge2$, $2^e\mid\operatorname{ord}$, and
$f=2\operatorname{ord}$, so $f\ge8$.  Positivity yields
$\operatorname{rd}(E)\ge4$, hence $v\ge2$.
Therefore $4k^4v^2\le(kv)^4$.  Using $d_\tau\ge0$,
$T\ge2$, and GLS4 in (GLS.47) gives exactly
\[
 \boxed{\quad
 \mu_{\max}(A)\le
   \sum_\tau\left(4d_\tau\log(kv)
          +k^2\log(h_\tau+k)+|s_\tau|\log T\right).
 \quad}                                                       \tag{GLS.48}
\]
This is \texttt{Slope\allowbreak Max D u 4}, on every original $D,\allowbreak{} u$ and
every given \texttt{Col\allowbreak Val D u}.  It uses no Hida size guard.
The transference input in GLS8 is an explicit, fully matched
published theorem; all subsequent arithmetic, ideal, and metric
comparisons have been proved here.
% END GLS11

\subsection{Uniform installation and the remaining final assembly}

% BEGIN GLS12
\paragraph{GLS12: all three literal GL4 family entries.}
\label{claim:GLS12}
For every original family $\mathrm{Fm}:\texttt{Fam\allowbreak P}\,\varepsilon$,
GLS7 and GLS11 apply to every $D=\mathrm{Fm}.D(x,\allowbreak{} i)$ and every
$u,\mathrm{CV}$.  Threshold zero consequently gives
\[
 \texttt{Family\allowbreak Slope\allowbreak Lower\allowbreak Open\allowbreak Part}\,\mathrm{Fm},\qquad
 \texttt{Family\allowbreak Slope\allowbreak Max\allowbreak Open\allowbreak Part}\,\mathrm{Fm}\,4.
 \tag{GLS.49}
\]
Their \texttt{fam0} quantifiers are
$\exists x_0\ \forall x\ge x_0\ \forall (h:\mathrm{Fm}.x_1<x)
\ \forall i$, followed by the pointwise predicate.  The argument
works for all such $x,\allowbreak{} i$, so these are universal \texttt{Closed}
statements on any family class, with no change of $\varepsilon$,
$C_H$, or the original member.
Explicitly, \texttt{Closed} means
\[
 \exists\varepsilon_1>0\ \forall\varepsilon\in(0,\varepsilon_1]
 \ \exists\mathcal K\ \forall\mathrm{Fm},\allowbreak{} C_H,\quad
 \mathcal A(\mathrm{Fm})\wedge\mathrm{Fm}.\mathrm{Large}(\mathcal K,\allowbreak{} C_H)
       \Longrightarrow P(\mathrm{Fm},\allowbreak{} C_H).
\]
For each assertion in (GLS.49) one may take
$\varepsilon_1=1$ and the original all-zero \texttt{Transported\allowbreak Alt}
tuple $\mathcal K$; the hypotheses of \texttt{Large} are unused.

Suppose the same family has \texttt{Real\allowbreak Place\allowbreak Second\allowbreak Layer\allowbreak Conditions}
at threshold $a$.  Thus, for every $x\ge a$ and every $i$, it
supplies the original real-place tie and one
$\mathcal F_{x,\allowbreak{} i}:\texttt{Cols\allowbreak Tied\allowbreak Sym}(\mathrm{Fm}.D(x,\allowbreak{} i))$,
whose forms have their supplied finite characters for every word.
RSL7 applied to each original $\mathcal F_{x,\allowbreak{} i}.\mathrm{forms}(u)$
already gives \texttt{adm\allowbreak S} pointwise, for every $u$; their conjunction
is \texttt{adm3}.  Thus one may take $b=a$.  The source's
\texttt{family\_\allowbreak conditions\_\allowbreak of\_\allowbreak layer\_\allowbreak two} also permits a separate
eventual stabilizer threshold $b$ with the guard
$\mathrm{Fm}.x_1<x$; the following common choice covers both routes:
\[
                    a_{\rm CV}=\max\{a,\allowbreak{} b,\mathrm{Fm}.x_1+1\}.
 \tag{GLS.50}
\]
For $x\ge a_{\rm CV}$, (GLS.15) applied to this very
$\mathcal F_{x,\allowbreak{} i}$ gives
\[
 \forall i:\mathrm{Fm}.Ix(x)\ \forall u,\quad
            \texttt{Nonempty}(\texttt{Col\allowbreak Val}(\mathrm{Fm}.D(x,\allowbreak{} i),\allowbreak{} u)).
 \tag{GLS.51}
\]
This is the literal unguarded \texttt{t\allowbreak CV} field: its threshold
is chosen before $i,\allowbreak{} u$, and (GLS.50) explicitly removes the strict
$x_1$ guard.  Any \texttt{Closed} second-layer input therefore
gives \texttt{Closed t\allowbreak CV} on the same class.

Equations (GLS.49)--(GLS.51) close three stated implication tasks:
the lower slope, maximum slope, and original column-value dictionary.
Their audit identifiers are respectively
\texttt{TASK\_\allowbreak GL4\_\allowbreak SLOPE\_\allowbreak LOWER},
\texttt{TASK\_\allowbreak GL4\_\allowbreak SLOPE\_\allowbreak MAX}, and
\texttt{TASK\_\allowbreak GL4\_\allowbreak COLUMN\_\allowbreak VALUES}.
The last retains the original same-selector second-layer producer.
The two slope inequalities are conditional on their given
\texttt{Col\allowbreak Val}, as their source predicates require.
% END GLS12

% BEGIN GLS13
\paragraph{GLS13: the complete eleven-field packet and its precise frontier.}
\label{claim:GLS13}
Let $\mathcal A$ be a family class and retain
\[
 \forall\mathrm{Fm}\in\mathcal A,\quad\texttt{Duty}(\mathrm{Fm}),
 \qquad
 \texttt{Closed}_{D_3(3,\mathcal A)}
                   (\texttt{Real\allowbreak Place\allowbreak Second\allowbreak Layer\allowbreak Conditions}).
 \tag{GLS.52}
\]
All the following applications concern that same class, family,
and second-layer witness.  MF19 supplies \texttt{h2b};
RSL1--RSL3, using the stated mass and volume inputs, supply
\texttt{h2v}; DSZ1--DSZ3 supply \texttt{h3}.
GLS12 supplies \texttt{h\allowbreak Za}, \texttt{h\allowbreak Zb} with $C_5=4$, and
\texttt{t\allowbreak CV}.  The original finite character from (GLS.52) and
RSL7--RSL10 supply \texttt{h\allowbreak W} and \texttt{h}$\mu$.
Finally GLE5, GLE9, and GLE6 supply \texttt{h\allowbreak I}, \texttt{h\allowbreak R}, and
\texttt{h\allowbreak B}, respectively, retaining the original \texttt{Fam\allowbreak U}
guard \texttt{Hida\allowbreak Setting}.  GLE12 proves their uniform installation.
These are exactly the eleven fields, in order, of
\[
                    \texttt{Fam\allowbreak GLmin}(D_3(3,\mathcal A),4).
 \tag{GLS.53}
\]
Each eventual intersection uses a maximum of finitely many existing
thresholds.  The maximum is taken before the member $i$, word $u$,
coefficient place, or level place is quantified.  There is no need
to choose a second form record or change any field of the family.
When combining the \texttt{Closed} assertions, take the minimum
of their positive $\varepsilon_1$ bounds and the componentwise
supremum of their finitely many \texttt{Transported\allowbreak Alt} tuples.
The original monotonicity of \texttt{Large} recovers each input
for the very same $\mathrm{Fm},\allowbreak{} C_H$.  These choices precede the
family and its members, as the displayed quantifier requires.
Notice also that (GLS.53), like the source's guarded fields, does
not itself require \texttt{Tu\allowbreak Size}.  That condition supplies
\texttt{Setting\allowbreak Eventually} when those guarded fields are consumed.

The smaller fullness/rank conclusion remains the already proved
RSL11--RSL12 and GLE route; it does not depend on GLS4--GLS11.
Membership in $D_3(3,\mathcal A)$ supplies the same family's
\texttt{Tu\allowbreak Size}.  The Duty clause supplies that family's
\texttt{Chosen\allowbreak Wall}.  These are exactly the two projections used
in the source's \texttt{link\_\allowbreak gl} after constructing (GLS.53).

We keep the final analytic assembly distinct.  The source's
\texttt{link\allowbreak GLOBAL7\_\allowbreak min} first transforms (GLS.53) into
\texttt{Part\allowbreak Five\allowbreak Family\allowbreak Bounds}, \texttt{Conclusion\allowbreak F5},
\texttt{Conclusion\allowbreak F6}, \texttt{family\_\allowbreak property\_\allowbreak and}, and
\texttt{Congruence\allowbreak Lattice\allowbreak Length\allowbreak Bound}.  Its
\texttt{link\allowbreak GLOBAL7\_\allowbreak items} then combines the original count,
two slope families, rational ranks, lattices, frames, loss windows,
characters, and size conditions in \texttt{link\allowbreak Core}.
Giving the names of those source maps is not a human proof of all
their numerical estimates and input identifications.  Accordingly
the remaining B obligation here is the actual numerical/identity
installation from this eleven-field packet to
\[
 \texttt{Tu\allowbreak Size}\to\texttt{DGam\allowbreak Size}\to
                    \texttt{Link\allowbreak GLOBAL}(\mathrm{Fm},2,\allowbreak{} C_H,3,2,0),
 \tag{GLS.54}
\]
through \texttt{link\allowbreak GLOBAL7\_\allowbreak min},
\texttt{link\allowbreak GLOBAL7\_\allowbreak items}, and \texttt{Fam\allowbreak GLmin.\allowbreak link8}.
The final \texttt{link8} projection retains (GLS.52)'s
\texttt{Chosen\allowbreak Wall}, not an independently chosen covering family.
The coefficient $C_5=4$ in (GLS.53), $C_{\rm abs}=2$, $C_0=3$,
$e_D=2$, and $e_L=0$ must remain the same in that task.
The common compatible selector and (GLS.52) remain upstream
producer obligations; neither is asserted to exist by this chapter.
% END GLS13


\section{The original joint Hecke ideal and its characteristic-prime support}
\label{sec:joint-support-producer}

\status{JSP1--JSP12 reconstruct the original ideal-kernel construction,
its full \texttt{Ring\allowbreak Props} and \texttt{Dense} consequences, the original
Taylor--Wiles base on the same module and ideal, and the single-point
\texttt{Nice\allowbreak I} implication from the explicitly retained patching inputs
to \texttt{P1core}. They do not construct the retained ordinary
arithmetic system, its four further Hida endpoints, or its all-level
automorphic modules. The joint producer \texttt{COUNT\_\allowbreak II\_\allowbreak SUPPORT}
therefore remains open. Every assertion below is an implication with
its mathematical antecedents printed; a source theorem name is not
used as its proof. The coefficient prime may be three.}

\subsection{The actual objects and the retained arithmetic data}

Fix one of the two original Links of QX4, at a permitted coefficient
object $c'$, and write $O=c'.O$, $k=O/\mathfrak m_O$ and $n=4$.
Retain its original CM field $L$, setting $s$, level $V$, residual
datum $r$, and the already fixed identification
$\iota:\overline{\mathbb Q}_\ell\simeq\mathbb C$.
Write $G_0=V.\mathrm{frame}.G_0$, $G_v$ and $I_v$ for its exact
local groups, and $I_u,\allowbreak{} I_0,\allowbreak{} D_0$ for its tame and auxiliary subgroups.
Let $B$ be any original universal ordinary determinant ring, with
pair $(D_B,\psi_B)$, and let $\Lambda$ be any original weight object.
Set $A=\Lambda.\Lambda$ and
\[
 \lambda=\Lambda.\mathrm{toPair}(B.u):A\longrightarrow B.
 \tag{JSP.1}
\]
KX1--KX9 already supply these universal objects and their maps, with
the nonempty-pair antecedent supplied here by QX5.

The ordinary construction uses the original data
\[
 Q=\mathrm{QData.empty}(V.\mathrm{hV}),\qquad
 R:\mathrm{UnitaryHeckeData}(Q,\iota,\Lambda).
\]
Such an $R$ includes
the integral residual type $\tau_0$, the character $\theta$ with its
\texttt{Is\allowbreak Theta} identity, and the residual Hecke eigenvalues $a$.
Its existence is retained. Put
\[
 H=\mathrm{hset}(Q,\allowbreak{} R.\tau_0),\qquad
 M=H.\mathrm{MK}(R.\theta,\allowbreak{} R.a,\allowbreak{} K),\qquad
 \mathbb T=H.\mathrm{TK}(R.\theta,\allowbreak{} R.a,\allowbreak{} K),                \tag{JSP.2}
\]
where $K$ is the original empty auxiliary level. The algebra
$\mathbb T$ is the subalgebra of $\operatorname{End}_A(M)$ generated
by the actual Hecke operators. The zero-module case is treated
separately and is never excluded by fiat.

For $M\ne0$ we retain the following ordinary arithmetic input,
on these precise objects.
\begin{enumerate}
\item The Hecke operators commute, $M$ is finite free over $A$,
and $\mathbb T$ is an object of $\mathcal C_O$. Its generators have
the original residual values. These are the relevant control and
integral module assertions, not just a generic existence of forms.
We retain the original $H.\mathrm{Ax}$, in particular its
prime-to-$\ell$ Hecke places and validity of its specified Artin maps.
\item There are an adically continuous determinant $D_{\mathbb T}$
and ordered local characters $\psi_{\mathbb T,\allowbreak{} v,\allowbreak{} i}$ with the
literal \texttt{Is\allowbreak Univ} Frobenius, unit and uniformizer identities.
Their reductions are $r.\mathrm{Dbar}$ and $r.\mathrm{psib}$;
the characters have the prescribed torsion restrictions, and
their restrictions to $I_v$ are the images of $\Lambda.\kappa_{v,\allowbreak{} i}$.
\item At each admitted weight, the original good arithmetic
homomorphisms $\mathbb T\to\overline{\mathbb Q}_\ell$ separate
$\mathbb T$. Their specialized determinants have the original
continuous representations with ordered flags, polarization and
the residual tame type. For the auxiliary assertion in JSP6,
the representations at one such separating weight are trivial
on $I_0$ when $r.\mathrm{V0ok}$ holds.
\end{enumerate}
These hypotheses are the precise parts of the ordinary control,
\texttt{core\allowbreak Det}, \texttt{Out}, residual-character and auxiliary
point producers that are used below. They are not asserted to have
been supplied by this chapter. In particular, separation at one
weight suffices for the determinant relations, whereas the full
\texttt{Dense} field needs separation at every weight used in JSP7.
For the literal weight dictionary we also retain the following exact
reciprocity matching: every map satisfying the source
\texttt{Is\allowbreak Local\allowbreak Artin} agrees with its specified $\operatorname{artOf}$,
and on inertia this map satisfies
$\epsilon(\sigma)=\prod_\tau\tau(\operatorname{Art}(\sigma))$.
These source-facing uniqueness and normalization assertions are
distinct from KX3's already proved abstract \texttt{Loc\allowbreak CFT} interface.
The source obtains them using global reciprocity, injectivity on
local ideles, and the cyclotomic norm formula. Their exact classical
matching is retained here, not silently attributed to KX3.
The exact source declarations are recorded in the accompanying
provenance file.

% BEGIN JSP1
\paragraph{JSP1: weight coordinates on the original inertia groups.}
\label{claim:JSP1}
Let $d=4[L^+:\mathbb Q]$. There are the source's
$\mathrm{WtVars}(\mathrm{locOf}(r),\allowbreak{} O,4,\allowbreak{} d)$ and an isomorphism
\[
 e:O[[X_1,\ldots,\allowbreak{} X_d]]\xrightarrow{\sim}A
\]
such that, for every coordinate $j$, its chosen
$\gamma_j\in I_{v_j}$, index $i_j$ and unit $c_j\in O^\times$
satisfy
\[
 \overline c_j\,\overline\psi_{v_j,\allowbreak{} i_j}(\gamma_j)=1,
 \qquad
 e(X_j)=c_j\Lambda.\kappa_{v_j,\allowbreak{} i_j}(\gamma_j)-1.       \tag{JSP.3}
\]
\begin{proof}
Use KX3's local abelian coordinates on the actual $G_v$ and its
actual inertia, and KX4's fixed character lift
$\psi^0_{v,\allowbreak{} i}:G_v\to O^\times$. For each basis vector of the
torsion-free inertia factor $\mathbb Z_\ell^{d_v}$, choose a
preimage $\gamma_{v,\allowbreak{} j}$ in $I_v$. The variable of the universal
principal-unit character at this basis vector is $X_{v,\allowbreak{} i,\allowbreak{} j}$.
The universal character therefore has value
$\psi^0_{v,\allowbreak{} i}(\gamma_{v,\allowbreak{} j})(1+X_{v,\allowbreak{} i,\allowbreak{} j})$ there.
KX5's universal isomorphism to the given $\Lambda$ gives
(JSP.3) with $c_j=\psi^0_{v,\allowbreak{} i}(\gamma_j)^{-1}$.
Reduction gives its first identity. The number of coordinates is
$\sum_v4d_v=4[L^+:\mathbb Q]$ by KX3.
Enumerate these finitely many triples to obtain the source's
$\mathrm{Fin}(d)$ coordinates; its $\mathrm{Empty}\sqcup\mathrm{Fin}(d)$
index set is the same set with an empty summand. No different local
group, torsion character, or coefficient ring has been substituted.
\end{proof}
% END JSP1

% BEGIN JSP2
\paragraph{JSP2: detecting full determinant identities by arithmetic fibers.}
\label{claim:JSP2}
Let $C$ be a commutative ring and $x_\alpha:C\to E_\alpha$ a
family of homomorphisms which jointly detects zero. Equalities of
determinant laws over $C$, and the source's full determinant-kernel
assertions, can be checked after every $x_\alpha$.
For a determinant law $D$ the latter assertion is literally
\[
 D\text{ kills }z
 \quad\Longleftrightarrow\quad
 D_S(1+h z_S)=1
 \quad\text{for every }C\text{-algebra }S\text{ and }h\in S[G_0].
 \tag{JSP.4}
\]
In a family of polarized flagged representations, the ordered
ordinary product and its star have zero image, and so are killed
in the sense of (JSP.4).
\begin{proof}
For a finite collection of group elements, introduce one independent
variable for each coefficient. Evaluate both determinant laws on
the resulting universal group-algebra element. Their values are
polynomials over $C$. A coefficient which maps to zero under all
$x_\alpha$ is zero; hence the polynomials agree. Naturality then
permits substitution of arbitrary coefficients in any $S$, proving
the asserted law identity. For (JSP.4), use the universal expression
$1+(\sum_j X_j[g_j])z$ and the same coefficient argument. This
does not assume that an injection into a product of fields remains
injective after arbitrary tensor products.

For the local ordinary relation, a common flag
$0=F_0\subset F_1\subset\cdots\subset F_4$ has successive
characters $\psi_0,\ldots,\psi_3$. The factor
$\rho(g_i)-\psi_i(g_i)$ preserves the flag and kills
$F_{i+1}/F_i$. Consequently the ordered product with indices
$0,1,2,3$ is zero: its rightmost factor first lowers $F_4$ to
$F_3$, and the next factors lower successively to zero.
This uses the source order and does not commute the factors.

If $P$ is the polarization matrix, the original form identity
gives
$\rho([g]^*)=P^{-1}\rho(g)^tP$, with
$[g]^*=\epsilon(g)^{1-4}[cgc^{-1}]^{-1}$.
Thus the representation of the star of a killed expression is
also zero, and determinant is unchanged by star. A zero matrix
image of $z$ gives determinant one for $1+h z$ over every
coefficient algebra. Apply the first paragraph to return all
these conclusions to $C$.
\end{proof}
% END JSP2

% BEGIN JSP3
\paragraph{JSP3: the actual ordinary pair and its weight-compatible map.}
\label{claim:JSP3}
Under the retained nonzero-module arithmetic input, there is a
unique $O$-algebra homomorphism
\[
 f:B\longrightarrow\mathbb T,\qquad
 f_*D_B=D_{\mathbb T},\quad f_*\psi_B=\psi_{\mathbb T},\quad
 f\lambda=\text{the structural map }A\longrightarrow\mathbb T.
 \tag{JSP.5}
\]
\begin{proof}
First verify the three relations of the original \texttt{Pair}.
JSP2 applied to the separating polarized arithmetic fibers gives
polarization and the ordered ordinary annihilation identity.
In particular their characteristic coefficients satisfy the
source's star relation, and the trace of $[h]o_v$ is zero.

For the tame relation, put $q=|k|$. Every nonzero residue class
has a unique lift $[a]\in O$ with $[a]^{q-1}=1$.
Indeed successive lifting of a root of $X^{q-1}-1$ works because
its derivative is a unit; completeness gives the limit. If two
lifts coincide residually, factor their difference of
$(q-1)$st powers. The remaining factor is a unit, since its
reduction is $(q-1)a^{q-2}\ne0$, proving uniqueness.
Uniqueness also proves multiplicativity. Thus the residual
characters have the specified Teichm\"uller lifts, and every such
lift takes the value one where its residual character does.
For $g\in I_p(u)=I_u\cap\bigcap_i\ker\overline\chi_i$,
the prescribed tame type makes every specialized $\rho(g)$
conjugate to the identity, hence equal to it. JSP2 gives
$D_{\mathbb T}$ kills $[g]-1$; the required traces are zero.

The remaining \texttt{Pair} fields are exactly the retained
continuity, residual determinant, local-character continuity,
reductions and torsion restrictions. We have therefore constructed
the literal pair over $\mathbb T\in\mathcal C_O$.
Universality of the given $B$ gives $f$ and its uniqueness.
Both $f\lambda$ and the structural map from $A$ carry every
universal inertial character to $\psi_{\mathbb T}|_{I_v}$.
The universal property of the given $\Lambda$ identifies them.
This proves the last identity without an extra choice of weight map.
The Teichm\"uller argument divides by $q-1$, a unit, and not by
$3$ or by $4!$.
\end{proof}
% END JSP3

% BEGIN JSP4
\paragraph{JSP4: recovery of every Hecke generator and surjectivity.}
\label{claim:JSP4}
The map $f$ of (JSP.5) is surjective as an ordinary ring map;
no completion of its image is needed.
\begin{proof}
The algebra $\mathbb T$ is, by its definition, generated over
$A$ by the actual spherical generators $T_{w,\allowbreak{} j}$
($1\le j\le4$) and ordinary generators $U_{v,\allowbreak{} j}$
($0\le j<4$). The structural copy of $A$ is in the image
by (JSP.5). Write $q_w=\operatorname{N}w$.
For the original inverse-Frobenius element $g_w$, the source
normalization is
\[
 \chi_{D_{\mathbb T}}([g_w];X)
   =\sum_{j=0}^{4}(-1)^j q_w^{j(j-1)/2}T_{w,\allowbreak{} j}X^{4-j},
 \qquad T_{w,0}=1.                                  \tag{JSP.6}
\]
Since $w\nmid\ell$, the integer $q_w$ is a unit in $O$.
Taking the coefficient of $X^{4-j}$ gives the explicit preimage
\[
 T_{w,\allowbreak{} j}=f\!\left((-1)^j q_w^{-j(j-1)/2}
                  [X^{4-j}]\chi_{D_B}([g_w];X)\right).
 \tag{JSP.7}
\]
The sign, inverse Frobenius, and exponent are those of the
original \texttt{hecke\allowbreak Pol}; none is changed to a different
normalization.

For each $v$, local reciprocity supplies a Weil element $g_v$
whose image under the source Artin map is its specified
uniformizer. The literal uniformizer identity is
\[
 U_{v,\allowbreak{} j}=\prod_{i=0}^{j}
       \psi_{\mathbb T,\allowbreak{} v,\allowbreak{} i}(g_v)\epsilon(g_v)^i
   =f\!\left(\prod_{i=0}^{j}\psi_{B,\allowbreak{} v,\allowbreak{} i}(g_v)\epsilon(g_v)^i\right).
 \tag{JSP.8}
\]
Thus every algebra generator is in $f(B)$. Polynomial
expressions in these generators and scalars from $A$ exhaust
$\mathbb T$, proving surjectivity. All inverses in (JSP.7)
are of norms prime to $\ell$, so this proof applies when
$\ell=3$.
\end{proof}
% END JSP4

% BEGIN JSP5
\paragraph{JSP5: the literal ideal and its annihilator are the same.}
\label{claim:JSP5}
If $M=0$, set $J=B$. Otherwise set $J=\ker f$ for (JSP.5).
This is the original \texttt{Is\allowbreak J} witness. In both cases the
original module has a $B$-action $\mu$ compatible with $\lambda$,
and
\[
               J=\ker\mu=\operatorname{Ann}_B(M).
 \tag{JSP.9}
\]
For $M\ne0$, $B/J\simeq\mathbb T$ as $A$-algebras.
\begin{proof}
In the nonzero case, let $\mu$ be $f$ followed by the inclusion
$\mathbb T\subset\operatorname{End}_A(M)$. This inclusion is
injective because the Hecke algebra was defined as a subalgebra
of endomorphisms. An element of $B$ therefore acts as zero on
$M$ exactly when its image under $f$ is zero. This proves
(JSP.9). The identity $f\lambda=A\to\mathbb T$ gives
$\mu(\lambda(a))m=am$. Surjectivity in JSP4 gives the
quotient isomorphism. The continuous determinant and character
identities in (JSP.5), together with their original
\texttt{Is\allowbreak Univ} formulas, are precisely \texttt{Is\allowbreak Map}.
These are all the nonzero-case \texttt{Is\allowbreak J} fields.

For $M=0$, its endomorphism ring is the zero ring. The unique
action has kernel $B$ and annihilator $B$; it satisfies the
scalar identity. The source explicitly permits this alternative
$M=0$, $J=B$. We have not inferred a nonzero ordinary module
from a residual datum.
\end{proof}
% END JSP5

% BEGIN JSP6
\paragraph{JSP6: all eight original ring properties.}
\label{claim:JSP6}
The ideal in JSP5 has the complete original
$B.\mathrm{RingProps}(\Lambda,\allowbreak{} J)$ under the retained arithmetic
input, including the auxiliary unramifiedness assertion stated
there.
\begin{proof}
If $M=0$, then $B/J=0$. It is finite over $A$, its torsion-free
condition is vacuous, and every determinant equality and kernel
identity holds in the zero ring. Suppose henceforth $M\ne0$.
By JSP1, $A$ is a Noetherian domain. Since $M$ is finite free,
$\operatorname{End}_A(M)$ is finite free over $A$.
Its submodule $\mathbb T$ is therefore finite and torsion-free
over $A$. By JSP5 these are exactly \texttt{fin} and
\texttt{tf} for $B/J$ and its specified weight map.

Transport the determinant and characters across
$B/J\simeq\mathbb T$. At a separating arithmetic fiber,
polarization gives the \texttt{polar} field by JSP2.
An upper triangular matrix has characteristic polynomial equal
to the product of its diagonal linear factors. The original
ordered flag therefore gives
\[
 \chi_{D_{B/J}}([\phi_v(g)];X)
       =\prod_{i=0}^{3}(X-\psi_{B/J,\allowbreak{} v,\allowbreak{} i}(g)),           \tag{JSP.10}
\]
which is \texttt{O1}. JSP2 gives both the ordinary-product
and star-product kernel assertions \texttt{O2}.

The Teichm\"uller uniqueness proof in JSP3 shows that every
lift allowed by \texttt{Is\allowbreak Teich} is the one already used there.
The tame type of each fiber thus identifies the whole restricted
determinant law on $I_u$ with the diagonal law of these lifts.
JSP2, applied coefficientwise on the group algebra of $I_u$,
proves the original \texttt{types} field, including its
quantifier over all permitted lifts. On $I_p(u)$ those
diagonal matrices are the identity; JSP2 proves both
\texttt{O3} kernel assertions. Finally the retained auxiliary
point statement gives identity matrices on $I_0$ whenever
$r.\mathrm{V0ok}$ holds. The same proof gives both
\texttt{unr0} assertions.

These are the eight source fields: finite generation,
torsion-freeness, polarization, \texttt{O1}, \texttt{O2},
\texttt{O3}, tame type and auxiliary unramifiedness.
The polynomial-law test in JSP2 verifies their assertions over
every coefficient algebra, rather than only at field-valued
points or at traces.
\end{proof}
% END JSP6

% BEGIN JSP7
\paragraph{JSP7: the original weight dictionary and full density.}
\label{claim:JSP7}
For every source strictly increasing weight $h$, the supplied
good Hecke fibers of the corresponding weight imply the exact
$B.\mathrm{Dense}(J)$. In this assertion the weight dictionary
is determined by the original $R.\mathrm{theta}$.
\begin{proof}
Use indices $0,\ldots,3$. At an embedding $\tau$ belonging to
the chosen place $v\mid\ell$, define the dominant weight by
\[
 \lambda_h(\tau)_a=h_{v,3-a,\tau}-(3-a).
 \tag{JSP.11}
\]
Because an increasing integer weight rises by at least one at
each step, these four entries are weakly decreasing.
The pairs $(v,\tau)$ restrict injectively to embeddings of
$L$: the valuation induced on $L$ determines the place, and
density of $L$ in its completion determines the continuous
embedding at that place. Hence (JSP.11) is consistent; assign
zero at unused embeddings. In the source notation
$\mathrm{htW}(\lambda_h,\allowbreak{} i)=\mathrm{lowW}(\lambda_h,\allowbreak{} i)+i$,
so it equals $h_{v,\allowbreak{} i,\tau}$ at the chosen embeddings.
The actual $H$ has $H.\mathrm{admitted}=\mathrm{Set.univ}$,
so every weight $\lambda_h$ just constructed is admitted.

Let $x:B\to\overline{\mathbb Q}_\ell$ and
$z=x\lambda:A\to\overline{\mathbb Q}_\ell$.
For a unit $y$ of $L_v$, let $F_{v,\allowbreak{} i}(y)$ be the source finite
part
\[
 F_{v,\allowbreak{} i}(y)=
 \frac{z(\theta(e_{v,\allowbreak{} i}(y)))}
      {\prod_\tau\tau(y)^{\mathrm{lowW}(\lambda_h,\allowbreak{} i)(\tau)}}.
\]
For $\sigma$ in the actual inertia group with
$\operatorname{Art}(\sigma)=y$, the \texttt{Is\allowbreak Theta}
identity and $\lambda$'s character identity give
\[
 x(\psi_{B,\allowbreak{} v,\allowbreak{} i}(\sigma))F_{v,\allowbreak{} i}(y)
                  =\prod_\tau\tau(y)^{-h_{v,\allowbreak{} i,\tau}}. \tag{JSP.12}
\]
Here the explicitly retained source reciprocity normalization has
$\epsilon(\sigma)=\prod_\tau\tau(y)$; its $i$th power
is the shift by $i$ in (JSP.11). This is the local norm formula
for the cyclotomic character in the specified reciprocity;
using the inverse reciprocity map requires inverting both
occurrences, not changing only this shift.

Equation (JSP.12) determines the finite weight-part character
on inertia. The retained uniqueness of the actual Artin map
makes the identity valid for \emph{every} Artin map quantified
in the source's \texttt{Is\allowbreak Wt\allowbreak Part}. It is locally constant
exactly when $F_{v,\allowbreak{} i}$ is
one on a sufficiently deep principal-unit group. Indeed a
locally constant character has open kernel, and reciprocity
identifies its finite quotient on units with a finite quotient
of inertia; conversely such a unit character pulls back to
an open-kernel character on inertia. There are finitely many
$(v,\allowbreak{} i)$, so the separate depths can be replaced by their maximum.
Thus $x.\mathrm{Arith}(h)$ is exactly the original
$H.\mathrm{Arith}(\theta,\allowbreak{} z,\lambda_h)$.
Surjectivity of inertia onto units under reciprocity also shows
that the characters for $i\ne j$ are distinct exactly when
$F_{v,\allowbreak{} i}(y)\ne F_{v,\allowbreak{} j}(y)$ for some unit $y$.
Consequently $x.\mathrm{Good}(h)$ is exactly
$H.\mathrm{Good}(\theta,\allowbreak{} z,\lambda_h)$.

Now suppose $b\in B$ vanishes at every good point of weight
$h$ killing $J$. Every good Hecke point
$t:\mathbb T\to\overline{\mathbb Q}_\ell$ yields such
a point $tf$ by (JSP.5) and the dictionary just proved.
It therefore kills $f(b)$. Separation at $\lambda_h$ gives
$f(b)=0$, or $b\in J$. If $M=0$, the conclusion
$b\in J=B$ holds without a point. This is precisely the full
\texttt{Dense} quantifier over every strictly increasing $h$.
\end{proof}
% END JSP7

% BEGIN JSP8
\paragraph{JSP8: the original Taylor--Wiles base on the same ideal.}
\label{claim:JSP8}
The module $M$ and ideal $J$ of JSP5 give the literal
\[
 b:\mathrm{TWBase}(\mathrm{pairOf}(B.u),\allowbreak{} J,
                         \mathrm{Empty},4[L^+:\mathbb Q]).
 \tag{JSP.13}
\]
Its module is $M$ itself, and its support agrees with $J$
by equality of annihilators, not only by an unproved comparison
of radicals.
\begin{proof}
Let $S=O[[X_1,\ldots,\allowbreak{} X_d]]$, use JSP1's $e:S\simeq A$,
and set $b.\mathrm{lam}=\lambda e$.
Equip $M$ with its $B$-action $\mu$ and its $S$-action
through $e$. The first identity of (JSP.3) gives
\texttt{wt\_\allowbreak res}; the second, composed with $\lambda$,
gives the exact source variable identity
\[
 b.\mathrm{lam}(X_j)=c_j\psi_{B,\allowbreak{} v_j,\allowbreak{} i_j}(\gamma_j)-1.
\]
The scalar identity of JSP5 says the two actions agree along
$b.\mathrm{lam}$. Transport an $A$-basis of $M$ along $e$
to obtain a finite free $S$-basis, including the empty basis
if $M=0$. In that case the zero action and empty basis suffice;
no nonzero Hecke algebra or nonzero-case arithmetic input is needed.
Finally (JSP.9) gives, for every prime $\mathfrak p$,
\[
 \operatorname{Ann}_B(M)\subseteq\mathfrak p
             \quad\Longleftrightarrow\quad J\subseteq\mathfrak p.
\]
These are all fields of \texttt{TWBase}. The replacement
\texttt{pair\allowbreak Of} retains the same underlying determinant and
ordered characters; it changes only the presentation of tame
inertia. The chosen $B$-action and its annihilator are unchanged.
This proves the base construction, not the existence of the
\texttt{TWRing}/\texttt{TWMaps} systems at all later levels.
\end{proof}
% END JSP8

\subsection{The single-point characteristic-prime support argument}

The remaining patching input is not replaced by characteristic-zero
density. We now prove the precise single-point implication used for
\texttt{P1core}, retaining its arithmetic and patching hypotheses.
This separates the elementary changes of rings and the discriminant
argument from the genuinely substantial support theorem.

% BEGIN JSP9
\paragraph{JSP9: topological generation by the original pair values.}
\label{claim:JSP9}
The original universal $B$ satisfies
$\mathrm{GenBy}(O,\mathrm{pairVals}(B.u))$: for every $N$ and
$b\in B$, a polynomial $a$ over $O$ in determinant coefficients
and ordered-character values satisfies
\[
                         b-a\in\mathfrak m_B^N.       \tag{JSP.14}
\]
\begin{proof}
Let $S$ be the set of source pair values. Two $O$-algebra maps
from $B$ to an object of $\mathcal C_O$ agreeing on $S$ give
the same characters and determinant. For the determinant this
uses the integral word identity already used in QX5: its
coefficients on finite sums of group elements are integral
polynomials in characteristic coefficients of group words.
For precision, work modulo any fixed total degree in formal
variables. The unique Lyndon factorization of words gives the
ordered product factorization of
$1-\sum_i Z_i[g_i]$ by factors $1-Z_w[g_w]$ indexed by
primitive cyclic words. Only finitely many factors contribute
at the chosen degree. Applying multiplicativity and homogeneity
of the determinant law expresses its coefficients through
the single-group-element factors. This argument applies to
the law itself and needs no matrix realization of $D_B$.
Thus equality on these coefficients gives equality of the full
laws, also in residue characteristic three. Universality of $B$
then makes the two maps equal.

Put $A_0=O[S]\subset B$ and
$I=(A_0\cap\mathfrak m_B)B$. Suppose
$\mathfrak m_B\ne I+\mathfrak m_B^2$.
The nonzero finite-dimensional $k$-space
$\mathfrak m_B/(I+\mathfrak m_B^2)$ has a nonzero linear
functional. It defines an $O$-derivation $\delta:B\to k$:
subtract from $b$ any lift in $O$ of its residue and apply the
functional. Different lifts differ by an element of
$\mathfrak m_O B\subset I$, so this is well defined.
Addition is immediate, and expanding a product after subtracting
its two scalar residues gives
$\delta(bb')=\bar b\delta(b')+\bar b'\delta(b)$, since
products of elements of $\mathfrak m_B$ are killed.
The two $O$-algebra maps
\[
 b\longmapsto\bar b+\delta(b)\varepsilon,
 \qquad b\longmapsto\bar b
 \quad\text{into }k[\varepsilon]/(\varepsilon^2)
\]
agree on $S$ but are different. This contradicts the preceding
unique-map property. Hence $\mathfrak m_B=I+\mathfrak m_B^2$.

Choose finitely many elements $y_j\in A_0\cap\mathfrak m_B$
whose residues span $\mathfrak m_B/\mathfrak m_B^2$.
Every degree-$r$ graded piece is generated over $k$ by the
degree-$r$ monomials in the $y_j$: express each factor of a
product in $\mathfrak m_B$ modulo $\mathfrak m_B^2$ and
expand. Scalars in $k$ lift to $O$. Starting with a scalar
approximation modulo $\mathfrak m_B$, correct each successive
degree using such a polynomial. After finitely many corrections
this gives (JSP.14). The ring $B$ and the set of character
values have not changed. This is exactly the hypothesis needed
by the source's \texttt{Raw} theorem.
\end{proof}
% END JSP9

% BEGIN JSP10
\paragraph{JSP10: finite coefficient extensions preserve kernels and give stable discriminants.}
\label{claim:JSP10}
Let $j:k_0[[t]]\to k_1[[t]]$ be an injective local homomorphism
between the power-series rings over finite fields occurring in
the source. Its restriction to constants comes from a field
embedding $k_0\to k_1$. In particular
\[
                      \ker(jf)=\ker f.               \tag{JSP.15}
\]
Let $z$ be a normalized source point \texttt{NP} with
$z.a=4$ and let $g:B\to z.K[[t]]$ carry $D_B$ to the
determinant of $z.\rho$. Then every ideal
$\mathfrak q\subseteq\ker g$ has the literal
$\mathrm{DiscSt}(D_B,\mathfrak q)$.
\begin{proof}
For $a\in k_0$, the series $j(a)$ satisfies $X^{|k_0|}=X$.
It is constant: otherwise its least positive nonzero degree
could not occur on the power side, whose corresponding degree
would have to come from a smaller positive degree. Thus $j$
restricts to a field embedding on constants. Locality makes
the constant coefficient of $j(t)$ zero. The induced residue
embedding therefore gives exactly the coefficient action needed
to regard a composite map $jf$ as an $O$-algebra map into the
new power-series ring. Equation (JSP.15) follows from
injectivity and does not require $j(t)=t$.

For the second statement put $F=z.K((t))$.
The normalized span field says that every integral matrix
becomes an $z.K[[t]]$-linear combination of $z.\rho(G')$
after multiplication by some power of $t$; here $G'$ is the
source cyclotomic kernel. Since $z.a=4$, its block algebra
is the full matrix algebra. Inverting $t$ gives
\[
          \operatorname{span}_F z.\rho(G')=M_4(F).
\]
Choose sixteen group images $\rho(g_i)$ forming a basis.
The trace pairing $(U,\allowbreak{} V)\mapsto\operatorname{tr}(UV)$ is
nondegenerate in every characteristic, since the matrix unit
$E_{ij}$ pairs with $E_{ji}$. Its Gram determinant in this
basis is therefore nonzero. Accordingly
\[
 \Delta=\det\bigl(\operatorname{tr}_{D_B}([g_i g_j])\bigr)_{i,\allowbreak{} j=1}^{16}
                 \quad\text{has }g(\Delta)\ne0.       \tag{JSP.16}
\]
The determinant specialization identifies each entry with the
matrix trace. If $\mathfrak q\subseteq\ker g$, then
$\Delta\notin\mathfrak q$. Taking the two source lists in
\texttt{Disc\allowbreak St} to be $(g_i)$ proves the assertion. No division
by the matrix size, a factorial, or the residue characteristic
is involved.
\end{proof}
% END JSP10

\paragraph{The retained support hypotheses.}
For JSP11 retain the source contracts \texttt{Norm},
\texttt{Choice}, \texttt{Ar\allowbreak Pl},
\texttt{First\allowbreak Order\allowbreak Dual\allowbreak Arithmetic\allowbreak Bound},
\texttt{Taylor\allowbreak Wiles\allowbreak Obstruction\allowbreak Bound}, and \texttt{Raw}, restricted
to the rank-four original data used here. Only their single-point
instances are needed. Their content is as follows.
\begin{enumerate}
\item \texttt{Norm} changes an actual \texttt{Nice\allowbreak I} point, by an
injective finite local extension of its power-series coefficient
ring, to a normalized point. It preserves the determinant and
ordered characters, gives the normalized full-matrix block
$a=4$, preserves \texttt{N3}, and retains \texttt{N4} for the
\emph{original} $D_0$.
\item \texttt{Choice} supplies finitely many elements of $G'$,
one further injective finite local extension, a source cocycle
frame \texttt{Kap}, and distinct eigenvalue lists at those
elements, with the exact \texttt{Detects} property. It
preserves the normalized representation, characters and block
size. The extra comparison of residue-field sizes in the
general source statement is automatic for a one-point family.
\item \texttt{Ar\allowbreak Pl}, from the given cohomological package
\texttt{Coh\allowbreak P}, realizes these elements by a group
$\mathcal T:\texttt{TWGrp}$ with its original
$\texttt{TWPl}$, including all level maps, prime norms,
Frobenius and tame inertia elements, and its subgroup $H$.
It also supplies the source finiteness condition used by
\texttt{Choice}.
\item The two arithmetic bound contracts give exactly
$\mathcal T.\texttt{B1}(k_z,4[L^+:\mathbb Q]+4\mathcal T.q)$
and $\mathcal T.\texttt{B2}(k_z)$ for this normalized point
and its chosen cocycle frame and eigenvalues.
\item \texttt{Raw} gives the source \texttt{Chain\allowbreak Concl} for
the actual pair, ideal and automorphic Taylor--Wiles data,
provided its \texttt{Adapted}, \texttt{Gen\allowbreak By} and point-bound
conditions hold. In the single-point instance $p=\ker g$,
the part used here says: if a prime
$\mathfrak q_1\subseteq p$ has stable discriminant and
contains $J$, then every prime $\mathfrak q\subseteq p$
having stable discriminant contains $J$.
\end{enumerate}
The last assertion is the substantial patched-support input.
The passage to it is not a proof of that input. The exact
quantified source predicates, rather than generic patching
terminology, are bound in the provenance file.

% BEGIN JSP11
\paragraph{JSP11: the complete one-point implication to \texttt{P1core}.}
\label{claim:JSP11}
Fix the same ideal $J\subseteq B$. Suppose that it has the
original $\mathrm{HasAut}(B,\allowbreak{} J)$, the six retained support
inputs above, and the source side conditions:
$\mathrm{CohP}$, $\mathrm{LocCFT}(s)$, $\mathrm{IsCoef}(\ell,\allowbreak{} O)$,
$V.S'=s.S$, and a nontrivial residual cyclotomic value on $G_0$.
Then $B.\mathrm{P1core}(J)$ holds. No \texttt{P3core}
assertion or \texttt{Nice\allowbreak E} point is used.
\begin{proof}
Take any original point $x$, map $f:B\to x.K[[t]]$, and
prime $\mathfrak q$ as in \texttt{P1core}: $x$ is
\texttt{Nice\allowbreak I}, $f$ matches the universal pair,
$J\subseteq\ker f$ and $\mathfrak q\subseteq\ker f$.
Use the unchanged arithmetic frame $\mathrm{globOf}(r,\allowbreak{} V)$.
Its tame inertia is $r.I_p(u)$: an element of the image of
the local inertia lies in the residual kernels precisely when
its local lift lies in their preimages. Consequently
$\mathrm{residOf}(r).I_p(u)=r.I_p(u)$ and this is already
the tame subgroup of the new frame. This verifies the exact
inertia premise for \texttt{Norm}; it also verifies the tame
part of \texttt{Adapted} later.

First take the \texttt{Phi\allowbreak Cond} component of \texttt{Ar\allowbreak Pl}.
Apply \texttt{Norm} to $x$, then the one-point
\texttt{Choice} construction. We obtain a normalized point
$z$ with $z.a=4$, an injective local composite
$j:x.K[[t]]\to z.K[[t]]$, a cocycle frame $k_z$, and the
finite detection and eigenvalue data. JSP10 supplies its
correct $O$-algebra structure and
\[
                         g=jf,\qquad \ker g=\ker f.   \tag{JSP.17}
\]
All determinant and ordered-character identities compose
along these maps, so $g$ still matches the original pair.
The original \texttt{N4} survives the second extension:
it is the vanishing of the kernel of a family of linear
equations in sixteen matrix entries over the fraction field.
A nonzero fraction-field solution could be multiplied by a
common denominator to give a nonzero integral solution, so
the integral \texttt{N4} from \texttt{Norm} gives this assertion.
A finite subsystem has full rank when that kernel is zero,
and its nonzero minor stays nonzero under the injective
field extension. This is the same original auxiliary subgroup,
not the trivial $D_0$ stored in an intermediate frame.

Use \texttt{Ar\allowbreak Pl} to realize the chosen elements by
$\mathcal T$ and its actual $\mathrm{TWPl}$.
Its identity for $H$ says in particular that every residual
character is one on $H$. Thus the residual characteristic
polynomial of $z.\rho(h)$ is $(X-1)^4$ for $h\in H$;
this follows from the residual determinant identity and the
constant-field embedding of JSP10. We verify the source
$\mathrm{PL}(\ell,\allowbreak{} H)$ condition used in the two bounds.
If $N=z.\rho(h)-1$, Cayley--Hamilton gives
$N^4\in tM_4(z.K[[t]])$. In characteristic $\ell$,
\[
 z.\rho(h)^{\ell^{4m}}-1=N^{\ell^{4m}}
                                      \in t^mM_4(z.K[[t]]).
 \tag{JSP.18}
\]
The exponent works since $\ell^{4m}\ge4m$.
At each depth $m$ take the open kernel of reduction of the
representation and its normal core in the original profinite
group. The original $\mathcal T.H$ is an open normal subgroup;
its residual-kernel description also makes that normality explicit.
Every conjugate of $h\in H$ is again in $H$, so
(JSP.18) puts a power of $\ell$ of every such $h$ in that
normal core. By the definition of $\mathrm{PL}(\ell,\allowbreak{} H)$
as the intersection of these eligible open normal subgroups,
the representation is one there modulo $t^m$ for every $m$.
Separation of the power-series ring proves it is exactly one.

The two retained arithmetic bound contracts now apply and
give the precise \texttt{B1} and \texttt{B2} fields printed
above. Along with $z,\allowbreak{} k_z$, the eigenvalues and $g$, these
construct the source point
$\mathrm{Pt}(\mathcal T,\mathrm{pairOf}(B.u),
4[L^+:\mathbb Q]+4\mathcal T.q)$.
The first \texttt{Adapted} field follows from the same
$H$ identity, and its second field was checked in the first
paragraph. JSP9 supplies \texttt{Gen\allowbreak By} for this exact pair.
The pair replacement preserves all coefficient and character
values, so no new generation argument is required.

The original \texttt{Has\allowbreak Aut}(B,\allowbreak{} J) applies to this very
$\mathcal T$ and its \texttt{TWPl}, supplying
\[
 \mathrm{TWAut}(\mathrm{pairOf}(B.u),\allowbreak{} J,\mathrm{Empty},
                         4[L^+:\mathbb Q],\mathcal T).
\]
In particular its base ideal
is still $J$, selected before $x$ and $\mathcal T$.
Apply the retained \texttt{Raw} theorem to the one-point
family. Its residue-cardinality comparison is reflexive.
By JSP10, $p=\ker g$ is a prime with stable discriminant,
and every prime below it also has stable discriminant.
Since $J\subseteq p$, choose $\mathfrak q_1=p$ in the
single-point support conclusion. It gives
$J\subseteq\mathfrak q$ for the prime originally chosen,
using (JSP.17). This proves the full original \texttt{P1core}
quantifier. There are no connecting thick primes in a
one-point family and no appeal to their separate
augmentation-kernel comparison.
\end{proof}
% END JSP11

% BEGIN JSP12
\paragraph{JSP12: exact same-ideal assembly and the two-Link quantifiers.}
\label{claim:JSP12}
The preceding proofs reduce the original joint producer to the
following remaining inputs on each of QX4's two actual Links:
the ordinary data in (JSP.2) and the three arithmetic items
following it, together with the printed exact reciprocity matching;
the four remaining Hida fields
\texttt{pt}, \texttt{end1}, \texttt{end2}, \texttt{twist}
on the ideal of JSP5; the all-level automorphic system on this
same ideal; and the retained support inputs of JSP11.
If these inputs are supplied with one common coefficient
threshold after the fixed $\iota$, they give
\texttt{COUNT\_\allowbreak II\_\allowbreak SUPPORT}, hence \texttt{PMEx\allowbreak Alt} and
\texttt{Count\allowbreak II} by the existing QX10 and YX2 implications.
\begin{proof}
Fix the original Fit/Deep tower and $\iota$ first.
Choose a coefficient threshold containing QX4's threshold
and any thresholds required by the retained inputs for the
two Links. A finite compositum of coefficient fields gives
one threshold in the original coefficient system. It is
chosen before the later $c'$, before an original Space $X$,
and before its universal data. At every allowed $c'$ use
both QX4 Links with their common $\kappa$ and distinct
auxiliary primes. They keep their own residual data and
universal rings. Work on each with arbitrary actual $B$
and $\Lambda$; no preferred universal representative is
being substituted for the one supplied.

For each Link, JSP3--JSP5 construct its one ideal $J$;
JSP6 and JSP7 give \texttt{ring} and \texttt{dense}.
Combine these with the retained four further fields on
that exact $J$ to obtain \texttt{Hida\allowbreak Props\allowbreak Alt}.
JSP8 supplies its original Taylor--Wiles base whenever
the retained all-level construction uses that base.
More explicitly, if for every realized $\mathcal T$
the all-level construction supplies the residual lifts
$u_{j,\allowbreak{} i}$, each $\mathrm{TWRing}_N$, and
$\mathrm{TWMaps}_N$ with one common exponent (the source
construction uses $4!=24$), these records form
$\mathrm{TWAut}$ and hence \texttt{Has\allowbreak Aut} on the same $J$.
This observation assembles supplied all-level data; it
does not produce their modules or prove their local operator
identities. If \texttt{Has\allowbreak Aut} itself is supplied, JSP11
uses its witness directly.

KX3 supplies the original \texttt{Loc\allowbreak CFT} interface.
The original coefficient object gives \texttt{Is\allowbreak Coef}:
it is a complete DVR with finite residue field of
characteristic $\ell$. The Link gives $V.S'=s.S$ and
\texttt{no\allowbreak Mu}. The latter means that $L$ does not contain
the $\ell$th roots of unity. Thus the action of $G_L$ on
them is nontrivial. That action factors through the level
quotient, since the omitted inertia is away from $\ell$,
so the required residual cyclotomic value already occurs
on the original $G_0$. Its cohomological input is the
explicit \texttt{Coh\allowbreak P} in JSP11; no unnamed strengthening
is added. JSP11 therefore gives \texttt{P1core} on the
very same $J$.

Doing this twice produces $J_0\subset B_0$ and
$J_1\subset B_1$, each with all seven required properties.
Their common threshold precedes all later choices.
QX10 forms
$\phi_0(J_0)X.B+\phi_1(J_1)X.B$ on the original $X$
and gives its five fields. YX2 then supplies the original
\texttt{Count\allowbreak II}. Neither finite composita nor any algebraic
step here imposes $\ell>3$.
\end{proof}
% END JSP12

\paragraph{The remaining boundary and its contribution.}
The kernel and support transport is now a written human argument.
What remains is the production of the ordinary arithmetic system,
the four specified Hida endpoints on that kernel, the actual
all-level automorphic modules, and the six explicitly listed
arithmetic/patching inputs in their single-point scope.
In particular, \texttt{Has\allowbreak Aut} is stronger than the
\texttt{TWBase} constructed in JSP8. The latter's finite freeness
and support identity do not construct the higher-level modules or
prove patched support. A rank-two control theorem, ordinary
characteristic-zero density, or QX1's adjoint vanishing does not
fill these missing rank-four inputs.

The frontier remains category C; the exact published/source
matching of a proposed support replacement in characteristic $\ell$
remains category D until its hypotheses and ideal are matched.
The completed JSP implications are category B human arguments,
not new independent existence premises. They are useful on the
RH route because KA5--KA6 remove only Case III for 2-primary
targets: Cases I, II and V still require their actual inputs.
For unrestricted Dirichlet GRH the same \texttt{PMEx\allowbreak Alt}
also feeds the selected Case-III route, but it supplies neither
the separate original \texttt{Ctx3} promodularity chain nor
\texttt{Seed\allowbreak IVAlt}. No RH, GRH, or unconditional Count II
conclusion is asserted here.


\section{Cross-word counting and the original ordinary column slopes}
\label{sec:global-count}

\status{GVC1--GVC7 prove the original logarithmic \texttt{Count\allowbreak F}
package, including comparison of different words on the same datum.
GVC8--GVC11 prove the full \texttt{Nums\allowbreak L}/\texttt{Wts\allowbreak L} supplier,
retaining the derivative hypothesis required by its literal statement.
GVC12--GVC17 install the original ordinary column slopes and
\texttt{Slopes\allowbreak MF}. GVC18 shows that the literal
\texttt{Fam\allowbreak P.\allowbreak sizes} field already supplies the eventual derivative
bound, so no additional derivative premise is needed on that family.
All results concern a supplied
original family; they do not construct its compatible selector or
second layer.}

Fix the original family $\mathcal F=\mathrm{Fm}$, with
$0<\epsilon<1/52$. At a member $x,\allowbreak{} i$ write $D=\mathcal F.D(x,\allowbreak{} i)$
and $\mathcal L=\mathcal F.L(x,\allowbreak{} i)$. The scalar and rank notation is
\[
 \begin{gathered}
 \ell=x^\epsilon,\quad b=e^\ell,\quad a=xb,\quad
 X'=\exp(x^{\epsilon/3}),\quad b'=a/X',\\
 w''=\frac{b'}{xn},\quad S=\mathcal L.\mathrm{Sc}=4k\Delta,\quad
 \eta=\mathcal L.\mathrm{etaSz},\\
 \mathcal A=\mathcal L.\mathrm{La}=\log(a\,\operatorname{rd}E).
 \end{gathered}\tag{GVC.1}
\]
Here $\Delta=D.s.\mathrm{Del}$, $f=[E:\mathbb Q]/2=2n$,
$N=k+m(t+1)$, and $n_T=\#D.\mathrm{host}.\mathrm{TT}$.
These are equalities in the original records. In particular the
parameter $\Delta$ is not chosen anew in this chapter.
A word is denoted $\omega$, to distinguish it from its
Harish--Chandra coordinates $u_{\omega,\tau,\allowbreak{} j}$.
All logarithms of arithmetic quantities below are of positive numbers.

The finite-family predicates use eventual thresholds before the member
$i$, the word $\omega$, the archimedean place $\tau$ and any
coefficient place. The auxiliary constant
$C_Y=\max\{5,\allowbreak{} C_{YL}+1\}$ below is unrelated to the fixed global
constant $C_0=3$. The exact GL4 maximum-slope coefficient remains
$C_5=4$; the free constants in \texttt{Conclusion5} and
\texttt{Conclusion6} will be denoted $C_-$ and $C_+$.

\subsection{The logarithmic Vandermonde comparison}

% BEGIN GVC1
\paragraph{GVC1: logarithmic perturbations and the harmonic bound.}
\label{claim:GVC1}
If $c>0$, $z,\allowbreak{} z'\ge c$ and $|z-z'|\le e$, then
\[
                  |\log z-\log z'|\le e/c.                 \tag{GVC.2}
\]
Indeed, if $z\ge z'$, integrate $1/s\le1/c$ over
$[z',\allowbreak{} z]$; interchange the two variables otherwise.
For an integer $q\ge1$,
\[
 H_q=\sum_{j=1}^{q}\frac1j
 \le1+\int_1^q\frac{ds}{s}=1+\log q.
\]
The integral comparison follows termwise from
$1/j\le\int_{j-1}^j ds/s$ for $j\ge2$.
Consequently, for $0\le r<k$,
\[
                  \sum_{s=r+1}^{k-1}\frac1{s-r}
                         \le1+\log k.                     \tag{GVC.3}
\]
An empty sum is zero, so this includes $r=k-1$ and $k=1$.
These prove the source's \texttt{abs\_\allowbreak log\_\allowbreak sub\_\allowbreak le\_\allowbreak div}
and \texttt{harmonic\_\allowbreak gap\_\allowbreak sum\_\allowbreak le} independently.
% END GVC1

% BEGIN GVC2
\paragraph{GVC2: separated coordinate intervals give the exact variation.}
\label{claim:GVC2}
Let $\Delta>0$ and let $u,\allowbreak{} u'$ have $k$ real coordinates. Suppose
that every pair $0<s<k$ satisfies
\[
 u_0-u_s,\ u'_0-u'_s\ge4k\Delta,\qquad
 |(u_0-u_s)-(u'_0-u'_s)|\le\Delta,
\]
and that every $1\le r<s<k$ satisfies
\[
 \begin{split}
 u_r-u_s,\ u'_r-u'_s&\ge3(s-r)\Delta,\\
 |(u_r-u_s)-(u'_r-u'_s)|&\le2\Delta.
 \end{split}
\]
Writing $V_k(u)=\prod_{r<s}(u_r-u_s)$, we have
\[
 \boxed{\quad
 |\log V_k(u)-\log V_k(u')|
 \le\frac14+\frac23 k(1+\log k).
 \quad}                                                     \tag{GVC.4}
\]
For $k=0$, both Vandermondes are the empty product $1$,
so the conclusion is immediate. For $k\ge1$ all their factors
are positive. Expand their logarithms and apply GVC1.
Each pair with $r=0$ costs at most $1/(4k)$, and their total is
at most $1/4$. For any fixed $r\ge1$ the cost is at most
$(2/3)\sum_{s>r}(s-r)^{-1}\le(2/3)(1+\log k)$.
There are at most $k$ such rows. The triangle inequality proves
(GVC.4), with exactly the source's coefficient $2/3$.
% END GVC2

% BEGIN GVC3
\paragraph{GVC3: check both endpoints and every interior interval on the original data.}
\label{claim:GVC3}
For the moment suppose $k\ge44$, $\Delta\ge1$, and at a given
place $\tau$ the actual coordinates satisfy
\[
 \begin{gathered}
 v_j-v_{j+1}=4\Delta,\qquad v_j\ge0,\\
 u_{\omega,\allowbreak{} j}=v_{j-1}-\tfrac12-\xi_{\omega,\allowbreak{} j}\Delta,
 \qquad \xi_{\omega,\allowbreak{} j}\in\{0,1\},\quad1\le j\le k-2,\\
 u_{\omega,0}-v_0\ge4k\Delta.
 \end{gathered}                                               \tag{GVC.5}
\]
At a non-top place assume
$u_{\omega,\allowbreak{} k-1}=v_{k-2}-1/2$; at a top place assume
$u_{\omega,\allowbreak{} k-1}\le-4k\Delta$.
Finally, require that both endpoints are independent of $\omega$.
These are exactly the parts of \texttt{Data\allowbreak A} used by the source
variation argument. They imply the hypotheses of GVC2 for any two
words on this same $D,\tau$.

Here is the complete pair check. The chain identity gives
$v_j=v_0-4j\Delta$. Each interior coordinate lies in the same
interval
\[
 [v_0-4(j-1)\Delta-\tfrac12-\Delta,\
   v_0-4(j-1)\Delta-\tfrac12]
\]
of length $\Delta$, for every word. Every coordinate after the
first is at most $v_0$, including the last in either place type.
The fixed first endpoint therefore gives the first-row lower bound
$4k\Delta$, while the difference of the two corresponding gaps
has absolute value at most $\Delta$.

For two interior coordinates at distance $d=s-r\ge1$, each gap
is at least $(4d-1)\Delta\ge3d\Delta$.
The two gaps differ by at most $2\Delta$.
If $s=k-1$ is non-top, substitute
$v_{k-2}-1/2$ for the last coordinate and obtain the identical
$(4d-1)\Delta$ estimate; its variation is at most $\Delta$.
If $s=k-1$ is top, nonnegativity of $v_{r-1}$ gives
$u_{\omega,\allowbreak{} r}\ge-\Delta-1/2$, and hence the gap is at least
$(4k-1)\Delta-1/2$. This dominates
$3(k-1-r)\Delta$: their difference is
$(k+2+3r)\Delta-1/2>0$.
The last endpoint is fixed, so here too the variation is at most
$\Delta$. This exhausts the pairs.

For the literal family, DSZ4 proves the cloud formulas and the
non-top endpoint formulas in (GVC.5) on the original
$\Delta=D.s.\mathrm{Del}$. The identities
$D.\mathrm{hcW\_\allowbreak first}$ and $D.\mathrm{hcW\_\allowbreak last}$ identify both
endpoints with the endpoints of the one original $D.\mathrm{hc}$,
so they are independent of the word, including at top places.
At a top place these endpoints are
\[
 u_{\omega,0}=a_T+(k-1)/2,\qquad
 u_{\omega,\allowbreak{} k-1}=-b_T-(k-1)/2.
\]
The original parameters give $kS\le w''\le b'\le b$ and
$a_T\ge a$, $b_T\ge b'$. Thus the last endpoint is at most
$-b'\le-S=-4k\Delta$. Also $v_0\le S$ and
$a=xb\ge2S$ once $x\ge2$, so $u_{\omega,0}-v_0\ge S$.
This verifies the top inequalities rather than assuming an
unproved full \texttt{Rows\allowbreak A} package.

DSZ4 also gives $\Delta\ge1$ and $m\ge x^{1/4}$,
with $k\ge m$. Enlarge one threshold so that $m\ge44$.
Every check above then holds for all $i,\omega,\tau$.
In particular (GVC.4) proves the literal
\texttt{vandermonde\_\allowbreak variation\_\allowbreak at\_\allowbreak place} whenever
\texttt{Data\allowbreak A} is supplied, and proves its needed actual-family
conclusion directly without claiming every stronger
\texttt{Data\allowbreak A}/\texttt{Family\allowbreak Part\allowbreak Five} clause.
% END GVC3

% BEGIN GVC4
\paragraph{GVC4: cancellation of the word-independent mass factors.}
\label{claim:GVC4}
The exact source definitions give
\[
 d_*^{\,\mathrm{form}}(\omega)
 =K_D\prod_{\tau}V_k(u_{\omega,\tau}),\qquad
 K_D=\frac{\mathrm{massI}(D,\allowbreak{} k)\,
          \mathrm{idxPi}(D,\allowbreak{} g_0,\allowbreak{} g_1)}
          {\prod_{\tau}\prod_{j=0}^{k-1}j!}.                 \tag{GVC.6}
\]
Both $K_D$ and $D_*^{\,\mathrm{form}}=\mathrm{DstarD}(D)$
belong to this same $D$ and are independent of $\omega$.
For the original
\[
 X_\omega=\frac1{fk}
          \log\frac{D_*^{\,\mathrm{form}}}
                    {d_*^{\,\mathrm{form}}(\omega)}
            =\mathrm{logXuD}(D,\omega)
\]
their contributions cancel in a difference. Since the original
place set has cardinality $f$, GVC3 gives
\[
 \boxed{\quad
 |X_\omega-X_{\omega'}|
 \le \frac1{4k}+\frac23(1+\log k)
 \le 1+\log k+\frac2k .
 \quad}                                                      \tag{GVC.7}
\]
Indeed expansion of the positive Vandermonde products bounds the
absolute value of their log difference by
$f[1/4+(2/3)k(1+\log k)]$; divide by $fk$.

This also proves the full generic
\texttt{Vandermonde\allowbreak Variation\allowbreak Explicit\allowbreak Bound} under \texttt{Data\allowbreak A}.
Its first requested bound is
$k(1+\log k)+2$, which is larger than (GVC.4).
For completeness, the formal generic expression allows a zero
common factor or a zero numerator. With the formal convention
$\log0=0$, and real logarithm $\log|z|$ for $z\ne0$, either zero
case makes both log-ratios zero. In the other cases
$\log|ab|=\log|a|+\log|b|$ gives the same cancellation.
No nonzero hypothesis is omitted in this generic statement.
For the actual family all these factors are positive by
QW1--QW2, the mass identity and the original positive degree fields,
so only the ordinary positive-logarithm calculation is used.

Finally the stored original equalities
$\mathcal F.\mathrm{ex}.\mathrm{duStar\_\allowbreak def}$ and
$\mathcal F.\mathrm{ex}.\mathrm{Dstar\_\allowbreak def}$ give
\[
 d_*^{\,\mathrm{actual}}(\omega)=d_*^{\,\mathrm{form}}(\omega),
 \qquad D_*^{\,\mathrm{actual}}=D_*^{\,\mathrm{form}}.
\]
Consequently $X_\omega=\log((D.\mathrm{mem}\,\omega).A.\mathrm{Xu})$.
This is the complete \texttt{jnt\_\allowbreak log\allowbreak Xu} dictionary: the positive
real power in \texttt{Xu} has exponent $1/(fk)$.
No new mass, row degree, coefficient field or column is substituted.
% END GVC4

% BEGIN GVC5
\paragraph{GVC5: the same prime center gives cross-word prime variation.}
\label{claim:GVC5}
The original field $D.\mathrm{GamP\_\allowbreak eq}$ identifies all
$\Gamma_\omega=(D.\mathrm{mem}\,\omega).A.\mathrm{GamP}$
with the same positive $\Gamma$. Thus
\[
 \log T_\omega^*
       =X_\omega-\frac2{fk}\log\Gamma .
\]
The stored window $T_\omega^*\le T_\omega
\le(1+1/(fk))T_\omega^*$ says that
\[
 \delta_\omega:=\log T_\omega-\log T_\omega^*
          \in[0,\log(1+1/(fk))]\subseteq[0,1].
\]
In particular $|\delta_\omega-\delta_{\omega'}|\le1$.
The common $\Gamma$ term cancels before estimating, and GVC4 yields
\[
 |\log T_\omega-\log T_{\omega'}|
 \le\frac1{4k}+\frac23(1+\log k)+1
 \le4\log k.                                                 \tag{GVC.8}
\]
For the last inequality $k\ge5$ gives $\log k\ge1$ and
$1/(4k)\le1/4$; the resulting constant term
$1/4+2/3+1=23/12$ is at most $(10/3)\log k$.
All primes in (GVC.8) are those stored for the original two words.
One cannot obtain this bound from the separate size estimates
$\log T_\omega=O(m\ell)$; the fixed center factor and the
Vandermonde cancellation are essential.
% END GVC5

\subsection{The three clauses of the original count}

% BEGIN GVC6
\paragraph{GVC6: assemble the literal pointwise \texttt{Count\_\allowbreak at}.}
\label{claim:GVC6}
Assume the same eventual mass-ratio and base-volume inputs as
DSZ6, together with the original degree ties of DSZ3.
Set
\[
 C_Y=\max\{5,\allowbreak{} C_{YL}+1\},\qquad
 C_\epsilon=30+\frac{2+2C_Y}{\epsilon},\qquad
 C_c=\max\{C_\epsilon+2,\ 8+3C_Y,\ 4\}.                      \tag{GVC.9}
\]
For every sufficiently late member, define the original lower main
term
\[
 H=\sum_{j=1}^{m}\frac1f\sum_\tau
       \log\frac{w_j(D,\tau)v}{\pi}
       +m\log(g_0g_1).
\]
The QW6 lower estimate, applied with the exact DSZ6 mass and
factorial identity and DSZ4's checked coordinates, says
$X_\omega\ge H-C_\epsilon(\ell+m^2\ell/k)$.
DSZ9 proves $X_\omega-2\le\log T_\omega$ and
$\log T_\omega\le(8+3C_Y)m\ell$ on these same data.
We retain the full QW6 lower expression here, before the
coarser positive lower bound used for \texttt{Tu\allowbreak Size}.

Since $\operatorname{rd}E\ge2$, $x>1$ and $\ell\ge1$,
\[
 \mathcal A=\log x+\ell+\log\operatorname{rd}E\ge\ell\ge1.
\]
Therefore
\[
 \begin{split}
 C_\epsilon(\ell+m^2\ell/k)+2
 &\le(C_\epsilon+2)(\mathcal A+m^2\mathcal A/k),\\
 H-C_c(\mathcal A+m^2\mathcal A/k)
 &\le\log T_\omega\le C_cm\ell .
 \end{split}\tag{GVC.10}
\]
Together with GVC5 and $C_c\ge4$ this is every clause of
$\mathrm{Count\_\allowbreak at}(\mathcal L,\allowbreak{} C_c)$.

The finite-index dictionary is exact. The source set
$\mathrm{SizeData}_2.\mathrm{lifts}(D)$ is $\{1,\ldots,\allowbreak{} m\}$;
QW6's zero-based $w_{j,\tau}$ is the original
$\mathrm{wj}(D,\allowbreak{} j+1,\tau)$. The two auxiliary primes are $g_0,\allowbreak{} g_1$,
and positivity gives
$\log g_0+\log g_1=\log(g_0g_1)=\mathrm{lgG}(D)$.
The real place set is the original $\mathrm{Tau}$ of cardinality
$f$, so its factor is exactly $1/f$, with no additional factor two.
These identifications turn (GVC.10) and (GVC.8) into the literal
source formula, rather than an estimate for a reindexed model.
% END GVC6

% BEGIN GVC7
\paragraph{GVC7: a complete uniform \texttt{Count\allowbreak F} on the same family.}
\label{claim:GVC7}
Suppose the original family is supplied with
\[
 \begin{split}
 &\mathrm{fam0}(\mathcal F,\mathrm{MassRatio}_2),\\
 &\mathrm{fam0}(\mathcal F,\mathrm{BaseVolumeEqualsMass}),\\
 &\mathrm{fam0}(\mathcal F,\mathrm{SigmaDegreeTwoSidedBoundOfGap}).
 \end{split}                                                 \tag{GVC.11}
\]
The third input is in fact available on every original family
by DSZ3, including its already recorded Harish--Chandra input.
The first two retain the same-record mass/volume hypotheses of
MF19 and RSL1--RSL3; no new realization is chosen here.
DSZ5--DSZ6 give the exact degree and mass identities after one
family threshold. DSZ4 and DSZ8--DSZ9 give all remaining
coordinate and prime-center bounds after finitely many further
thresholds. In particular, the row degree in the lower QW6
application is the one tied to this $D.\mathrm{dsig}$.

The number $C_{YL}$ is constant on this family: apply the
projection to the entry at index $19$ in the stored equality
$\mathrm{parConsts}(\mathcal F.\mathrm{params}(x,\allowbreak{} i))
 =\mathcal F.\mathrm{PK}$.
Thus choose the constants (GVC.9) first.
Then take a finite maximum of the mass/volume thresholds, all
DSZ/QW thresholds just listed, $2$, the threshold $m\ge44$,
and $\mathcal F.x_1+1$. For every $x$ beyond this maximum and
every $i\in\mathcal F.\mathrm{Ix}(x)$, the source equality
$\mathcal L.\mathrm{sc}.x=x$ is valid and GVC6 applies
simultaneously to all words. We have proved
\[
 \exists C_c\ge0\ \exists x_0\
 \forall x\ge x_0\ \forall i\in\mathcal F.\mathrm{Ix}(x),\
               \mathrm{Count\_\allowbreak at}(\mathcal F.L(x,\allowbreak{} i),\allowbreak{} C_c),
                                                               \tag{GVC.12}
\]
which is the original \texttt{Count\allowbreak F}.

This closes the count interface needed by the global ledger.
It does not assert the stronger
\texttt{Family\allowbreak Vandermonde\allowbreak Count\allowbreak Explicit} or
\texttt{Part\allowbreak Five\allowbreak Family\allowbreak Bounds} package: their unused assertions are
not needed for (GVC.12). The proof is a direct human proof of the
literal consumer, with all three clauses checked.
% END GVC7

\subsection{Full numerical and weight conditions on the original family}

% BEGIN GVC8
\paragraph{GVC8: the complete original \texttt{Wts} predicate.}
\label{claim:GVC8}
Let the original $\mathrm{ParameterChoice}'(D,\mathcal L)$ be
given, $x\ge4$, $x\le X'$, and $C_n\ge3$.
For every word, the full source predicate holds:
\[
              \mathrm{Wts}(D,\omega,\allowbreak{} C_n,\allowbreak{} x,\allowbreak{} b,\allowbreak{} b',\allowbreak{} S,\eta).
\]

We check the clauses, including the appended weights that were
not needed for the shorter Hida-setting proof GLE1.
Put
\[
 \lambda_{\tau,\allowbreak{} j}=u_{\omega,\tau,\allowbreak{} j}
                  -\bigl((k-1)/2-j\bigr),\qquad
 z_{\tau,\allowbreak{} j}=w_j^{\rm raw}
              -\delta_j,\quad
 \delta_j=\frac{k}{2}+\frac{(2j+1)(t+1)}2,\quad0\le j<m,
                                                               \tag{GVC.13}
\]
where $w_j^{\rm raw}=\mathrm{wj}(D,\allowbreak{} j+1,\tau)$.
The original identity $D.\mathrm{hcW\_\allowbreak e}$ shows that this
$z_{\tau,\allowbreak{} j}$ is exactly \texttt{w\allowbreak J}, with no word dependence
hidden in the raw appended coordinates. For $j<m$,
$0\le\delta_j\le N$ and
$\delta_{j+1}-\delta_j=t+1$.

The half-integrality of $u_{\omega,\tau,\allowbreak{} j}$ and evenness of $k$
make every $\lambda_{\tau,\allowbreak{} j}$ integral. Interlacing through the
integral coordinate $v_j$ gives
$u_{\omega,\tau,\allowbreak{} j}-u_{\omega,\tau,\allowbreak{} j+1}\ge1$:
writing the two half-integers as $r+1/2,\allowbreak{} s+1/2$, an integer
strictly between them forces $r\ge s+1$.
Hence $\lambda_{\tau,\allowbreak{} j}$ is nonincreasing.
GLE1, with its exact cloud and endpoint calculation, gives
\[
 \begin{array}{ll}
 \tau\text{ top}:&
 a\le\lambda_{\tau,0}\le2a,\quad
 b'\le-\lambda_{\tau,\allowbreak{} k-1}\le5b'/4<2b',\quad
 0\le\lambda_{\tau,\allowbreak{} j}\le S\ (0<j<k-1),\\
 \tau\text{ non-top}:&
 0\le\lambda_{\tau,\allowbreak{} j}\le3S\le C_nS\quad(0\le j<k).
 \end{array}                                                 \tag{GVC.14}
\]
That same parameter calculation proved
\[
 Nm^{12}\le w'',\quad S,\eta,\allowbreak{} mt\le w'',\quad
 \eta\ge1,\quad b'=xn w'',\quad
 N\le w''\le b'/x,\quad0<b'\le b.                            \tag{GVC.15}
\]
For clarity, the first inequality follows from the actual
$h_0$ bound $Nm^{11}\le\eta$ and
$m\eta\le w''$ supplied by the original \texttt{mteta}
inequality; the $S$ estimate follows from $kS\le w''$.
These are inequalities on the supplied parameters, not a new
choice of appended weights.

The two original Harish--Chandra gaps give
\[
 e_{j+1}-b_{j+1}\ge(t+3)/2,\qquad
 b_{j+1}-e_j\ge(t+3)/2.
\]
Thus $w_{j+1}^{\rm raw}-w_j^{\rm raw}\ge t+3$,
and subtracting $\delta_{j+1}-\delta_j=t+1$ proves
$z_{\tau,\allowbreak{} j+1}\ge z_{\tau,\allowbreak{} j}$.

At a top place the original window is
$[b_T+b'/x,\ b_T+b'/x+\eta]$.
Since $b_T\le5b'/4$ and both $b'/x,\eta\le b'/4$,
every $z_{\tau,\allowbreak{} j}\le7b'/4<2b'$.
At a thin place its raw window is $[w'',\allowbreak{} w''+\eta]$,
so $z_{\tau,\allowbreak{} j}\le2w''\le b'/2<2b'$.
These are all non-thick places, using the original disjoint
top/thick/thin partition.
At a thick place the raw window is $[b,\allowbreak{} b+\eta+1]$;
therefore
\[
                    b-N\le z_{\tau,\allowbreak{} j}\le b+\eta+1.
\]
For the first-weight clause, at top
$z_{\tau,0}\ge b_T+b'/x-N\ge b_T=-\lambda_{\tau,\allowbreak{} k-1}$.
Off top, the raw lower endpoint is at least $w''$,
so $z_{\tau,0}\ge w''-N\ge0\ge-\lambda_{\tau,\allowbreak{} k-1}$.

Finally all raw weights at a fixed place lie in one common
interval $[\mathrm{lo},\mathrm{lo}+\eta+1]$; choose
$\mathrm{lo}=b_T+b'/x$, $b$, or $w''$ in the three cases.
As $m\ge2$ and $t\ge1$,
\[
 z_{\tau,\allowbreak{} m-1}-z_{\tau,0}
 \le\eta+1-(m-1)(t+1)\le\eta.
\]
This proves every clause of the literal \texttt{Wts}, including
its last two inequalities, with the same $\omega,\tau$ and the
same original appended coordinates.
% END GVC8

% BEGIN GVC9
\paragraph{GVC9: the four large numerical inequalities.}
\label{claim:GVC9}
Let $N,\allowbreak{} m\ge1$, $1\le\ell\le m$, $c_L\ge3$,
$C_c\ge0$, $C_n\ge1$, and suppose
\[
 m\ge c_L+C_c+6C_n+9,\quad
 T_0\le C_cm\ell,\quad G_0\le2C_n\ell,\quad
 b'\ge Nm^4.
\]
Then
\[
 \begin{split}
 N^2(c_L\ell+T_0)&\le Nb',\\
 Nc_L\ell+T_0&\le b',\\
 m^2N^2(c_L\ell+T_0+G_0)&\le Nb'c_L\ell,\\
 (4C_n+9)m^3N&\le b'.
 \end{split}\tag{GVC.16}
\]
Indeed
$c_L\ell+C_cm\ell\le(c_L+C_c)m\ell
 \le m^2\ell\le m^3$.
The first inequality follows by multiplying by $N^2$ and using
$Nm^3\le Nm^4\le b'$.
For the second,
$Nc_L\ell+T_0\le N(c_L\ell+C_cm\ell)\le Nm^3\le b'$,
because $N\ge1$ and $C_cm\ell\ge0$.
Also $c_L+2C_n\le m$ and $C_c\le m$, so
\[
 c_L+C_cm+2C_n\le m+m^2\le2m^2\le c_Lm^2.
\]
Thus $c_L\ell+T_0+G_0\le c_Lm^2\ell$.
Multiplying by $m^2N^2$ and then using $Nm^4\le b'$
gives the third assertion.
Finally $4C_n+9\le m$, so the fourth follows from the same
$Nm^4$ bound. This proves the source's \texttt{nums\_\allowbreak alg}
with no suppressed asymptotic step.
% END GVC9

% BEGIN GVC10
\paragraph{GVC10: check all sixteen \texttt{Nums} clauses.}
\label{claim:GVC10}
For an original $\mathrm{ParameterChoice}'(D,\mathcal L)$
assume $x\ge4$, $x\le X'$, and
\[
 \begin{gathered}
 7\log x\le\ell,\quad
 n\ge3,\quad \log n\le3\log x,\quad
 \log q,\log d\le\ell,\quad
 \operatorname{rd}E\le2n^2\sqrt{qd},\\
 |\mathcal L.\mathrm{logDerivativeAtOne}|\le\ell,\quad
 \ell\le m,\quad m\ge c_L+C_c+6C_n+9,\\
 c_L\ge3,\quad C_c\ge0,\quad C_n\ge3,\quad
 C_c\le C_n,\quad C_{YL}+1\le C_n,\quad
 \log T_\omega\le C_cm\ell .
 \end{gathered}                                               \tag{GVC.17}
\]
Then the literal
$\mathrm{Nums}(D,\allowbreak{} T_\omega,\allowbreak{} g_0,\allowbreak{} g_1,\allowbreak{} C_n,\allowbreak{} x,\allowbreak{} b,\allowbreak{} c_L\ell,\allowbreak{} b',\allowbreak{} S,\eta)$
holds.

Here is a complete check in the original clause order.
First $b=e^\ell>0$.
Since $v=\operatorname{rd}E/2$ and the host discriminant bound
gives $v\ge1$,
\[
 \log v\le2\log n+\tfrac12(\log q+\log d)
          \le6\log x+\ell .
\]
It follows that
$\log(xbv)\le7\log x+2\ell\le3\ell\le c_L\ell$.
The third clause follows from
$\log T_\omega\le C_cm\ell\le C_nm(c_L\ell)$.
The original $g_Y$ window has $Y\le g_j\le2Y$ and
$\log Y\le C_{YL}\ell$, so
$\log g_j\le\log2+C_{YL}\ell\le(C_{YL}+1)\ell
 \le C_n(c_L\ell)$; this proves clauses four and five.

GVC8 gives $b'\le b=xb/x$, $\eta\ge1$, and
$S,\eta,\allowbreak{} mt\le w''$.
Also $S=4k\Delta\ge1$, since the original
$\Delta\ge Nm^{10}\ge1$ and $k\ge1$.
As $n\ge3$,
\[
              S+\eta+mt\le3w''\le n w''=b'/x .
\]
These prove clauses six through nine.
From $Nm^{12}\le w''\le b'$ and $m\ge1$ we have
$Nm^4\le b'$.
Set $T_0=\log T_\omega$ and
$G_0=\log(g_0g_1)\le2C_n\ell$ in GVC9.
Its first three bounds are exactly clauses ten through twelve.

The stored equality
$\mathcal L.\mathrm{logDerivative\_\allowbreak at\_\allowbreak one\_\allowbreak eq}$ identifies
the derivative in (GVC.17) with the original $\mathrm{dlogL}(D,1)$.
Hence
$|\mathrm{dlogL}(D,1)|\le\ell\le C_n(c_L\ell)$,
which is clause thirteen; this uses the supplied derivative
bound and does not prove a new analytic estimate.
Clauses fourteen and fifteen are $x\ge2$ and $v\ge1$.
The last GVC9 inequality is clause sixteen because
$1+2(2C_n+4)=4C_n+9$.
This completes the full conjunction, with the exact source
normalization $\widetilde L=c_L\ell$.
% END GVC10

% BEGIN GVC11
\paragraph{GVC11: one full \texttt{Nums\allowbreak L}/\texttt{Wts\allowbreak L} supplier, with constants before the threshold.}
\label{claim:GVC11}
Suppose the supplied family has an eventual uniform upper bound
\[
 \log T_\omega\le C_cm\ell
 \quad\text{and an eventual bound}\quad
 |\mathcal L.\mathrm{logDerivativeAtOne}|\le\ell .
                                                               \tag{GVC.18}
\]
The first is the second clause of \texttt{Count\allowbreak F};
\texttt{Tu\allowbreak Size} also suffices for this particular input.
Write $\widehat C_c=\max(C_c,0)$ and fix
\[
 C_{n,0}=
 \max\bigl\{\max(3,\widehat C_c),\
                    \mathcal F.\mathrm{PK}[19]+1\bigr\}.
                                                               \tag{GVC.19}
\]
Then the exact source statement is
\[
 \boxed{\begin{gathered}
 \forall C_n\ge C_{n,0}\ \forall c_L\ge3\
 \exists x_0\ \forall x\ge x_0\ \forall i\in\mathcal F.\mathrm{Ix}(x):\\
 \mathrm{NumsL}(D,\mathcal L,\allowbreak{} C_n,\allowbreak{} c_L)
 \ \wedge\ \mathrm{WtsL}(D,\mathcal L,\allowbreak{} C_n).
 \end{gathered}}                                             \tag{GVC.20}
\]
Here \texttt{Nums\allowbreak L} is GVC10's \texttt{Nums} for every original
word, and \texttt{Wts\allowbreak L} is GVC8's \texttt{Wts} for every word.

We verify all uniformity and growth inputs.
For fixed $\epsilon>0$, $x^{14}\le e^{x^{\epsilon/3}}=X'$
eventually. One elementary proof puts $y=\log x$ and notes
$e^{(\epsilon/3)y}/y\longrightarrow\infty$, which follows from
$e^{(\epsilon/3)y}\ge(\epsilon y/3)^2/2$ for $y\ge0$.
This gives $14\log x\le x^{\epsilon/3}$ and hence
$7\log x\le\ell$ as well as $x\le X'$.

The actual \texttt{Size\allowbreak M} exponent is
$1/2-\epsilon/16-\epsilon/100>2/5$; the source
\texttt{Size\allowbreak N} exponents are $2+8\epsilon\pm\epsilon/100$.
Thus for $x\ge4$ the stored size inequalities imply
\[
 x^{2/5}\le m,\quad
 \ell\le x^{2/5},\quad n\ge x^2\ge3,\quad
 \log n\le3\log x.                                         \tag{GVC.21}
\]
The upper \texttt{Size\allowbreak QD} exponents for $\log q$ and $\log d$
are $\epsilon/5+\epsilon/100<\epsilon$, so both are at most
$\ell$. The original discriminant formula gives
$\operatorname{rd}E\le2Q_S\sqrt{qd}\le2n^2\sqrt{qd}$;
its lower bound gives $\operatorname{rd}E\ge2$.
These are the already matched host bounds of AH16, on the same
host, together with $Q_S\le n^2$.

Now fix $C_n,\allowbreak{} c_L$ in the order in (GVC.20).
Take a threshold also making
\[
       x^{2/5}\ge c_L+\widehat C_c+6C_n+9.
\]
This is possible because $x^{2/5}$ tends to infinity.
The identity for the entry at index $19$ used in GVC7 gives
$C_{YL}+1\le C_n$ uniformly in $i$.
Enlarge the threshold to dominate the two bounds in (GVC.18),
the power-exponential threshold, $4$, and $\mathcal F.x_1+1$.
For every later $x,\allowbreak{} i$, GVC10 applies with $\widehat C_c$ and
GVC8 applies with $C_n$; both apply to all words at once.
This proves (GVC.20), including every clause of the literal
\texttt{nums\_\allowbreak wts\_\allowbreak of\_\allowbreak fam} and
\texttt{nums\allowbreak L\_\allowbreak of\_\allowbreak fam\_\allowbreak thr} interfaces.

In particular, choose $C_n=C_{n,0}$ and $c_L=3$ before the
last threshold. This gives the exact common numerical packet
used both by ordinary column slopes and by the adjacent level-prime
length calculation. This chapter does not replace $3$ by a
member-dependent constant.
% END GVC11

\subsection{Finite weight estimates and the original column slopes}

% BEGIN GVC12
\paragraph{GVC12: finite sums of the literal column weights.}
\label{claim:GVC12}
Let $D,\omega$ be arbitrary original data, and supply
\texttt{Nums} and \texttt{Wts} with parameters
$(C_n,\allowbreak{} x,\allowbreak{} b,\widetilde L,\allowbreak{} b',\allowbreak{} S,\eta)$.
Assume $x\ge K:=\max\{1,\allowbreak{} C_n\}$, and put $a=xb$.
For each original place define
\[
 s_\tau=\lambda_{\tau,\allowbreak{} k-1},\quad
 a_\tau=\sum_{j=0}^{k-1}\lambda_{\tau,\allowbreak{} j},\quad
 d_\tau=\sum_{j=0}^{k-1}(\lambda_{\tau,\allowbreak{} j}-s_\tau),\quad
 h_\tau=\lambda_{\tau,0}-s_\tau .
\]
These are respectively the last weight, \texttt{col\allowbreak Abs},
\texttt{col\allowbreak Abs\allowbreak P} and \texttt{first\allowbreak Column\allowbreak Spread}.
In particular $a_\tau$ is a signed sum, not a sum of absolute
values. By the integral-interlacing argument in GVC8 the weights
are nonincreasing, hence $d_\tau,\allowbreak{} h_\tau\ge0$.

The top place set has exactly $n_T$ members.
Explicitly, its bijection with $t\in D.\mathrm{host}.\mathrm{TT}$
is
$t\mapsto D.eTau^{-1}(\tau_1(t))$.
The original $\tau_1(t)$ lies in $\mathrm{Psi}$; injectivity
follows from the stored frame bijection, and membership in
\texttt{top\allowbreak Places}, defined as this image, gives surjectivity.
Thus this count is for the actual place set, not a doubled set
of complex embeddings.

The elementary consequences of \texttt{Nums} needed here are
\[
 b'\le b,\quad x\ge2,\quad v\ge1,\quad
 S,\eta\ge1,\quad S+\eta+mt\le b'/x,\quad
 N\widetilde L+\log T_\omega\le b'.                          \tag{GVC.22}
\]
The original ranks give $m\ge2,\allowbreak{} t\ge1,\allowbreak{} k\ge5,\allowbreak{} N\ge k$.
It follows that $b'\ge4x\ge8$, $a\ge16$, and
\[
 \widetilde L\ge\log(av)\ge\log a>1,\qquad
                  k\le b'\le a.
\]
Here $\log T_\omega\ge0$, since $T_\omega$ is the original
positive prime. Also $KS\le xS\le b'$.
Consequently the top spread is at most
$2a+2b'\le4a$, and every non-top spread is at most
$C_nS\le KS\le a$. Hence
\[
              \log(h_\tau+k)\le\log(5a)
                            \le2\widetilde L.              \tag{GVC.23}
\]
The last step uses $\log5\le\log a\le\widetilde L$.

At top, write $\lambda_{\tau,0}=a_T(\tau)$ and
$-\lambda_{\tau,\allowbreak{} k-1}=b_T(\tau)$.
The middle weights lie in $[0,\allowbreak{} S]$, so
\[
 a_\tau\le a_T(\tau)+kS,\qquad
 d_\tau\le2a+2kb'+kS,\qquad |s_\tau|\le2b'.
\]
Off top all entries lie in $[0,\allowbreak{} C_nS]$, giving
$a_\tau,\allowbreak{} d_\tau\le kKS$ and $|s_\tau|\le KS$.
Sum over the original $f$ places and use $K\ge1$:
\[
 \begin{split}
 \sum_\tau a_\tau
    &\le\sum_{\tau\ {\rm top}}a_T(\tau)+fkKS,\\
 \sum_\tau d_\tau
    &\le2n_Ta+fk(2b'+KS),\\
 \sum_\tau|s_\tau|&\le f(2b'+KS),\\
 \sum_\tau k^2\log(h_\tau+k)&\le2fk^2\widetilde L .
 \end{split}                                                 \tag{GVC.24}
\]
All summands multiplied by logarithms in the next step have
the signs explicitly checked here. This finite calculation
uses neither a square-mean estimate for arbitrary auxiliary
spaces nor the derivative clause of \texttt{Nums}.
% END GVC12

% BEGIN GVC13
\paragraph{GVC13: the exact \texttt{Conclusion5} and \texttt{Conclusion6} formulas.}
\label{claim:GVC13}
For the same $D$ suppose its original open lower and upper
column bounds hold, with arbitrary fixed upper coefficient $c_5$.
Set
\[
 K=\max\{1,\allowbreak{} C_n\},\qquad q_5=\max\{c_5,0\},\qquad
 C_-=K+2,\qquad C_+=(q_5+1)(K+4).                            \tag{GVC.25}
\]
For every $x\ge K$ these prove
$\mathrm{Conclusion5}(D,\allowbreak{} C_n,\allowbreak{} C_-,\allowbreak{} x)$ and
$\mathrm{Conclusion6}(D,\allowbreak{} C_n,\allowbreak{} C_+,\allowbreak{} x)$.

To prove this, take the arbitrary word, parameters, \texttt{Nums},
\texttt{Wts}, and original $\mathrm{CV}:\mathrm{ColVal}(D,\omega)$
quantified in those statements. Write
$M_\omega=(D.\mathrm{fm}\,\omega).\mathrm{lat}(0)$.
The lower raw inequality, proved for the actual model by GLS7, is
\[
 \mu(M_\omega)\ge
 -\frac{\log T_\omega}{k}\sum_\tau a_\tau
 -\log k\sum_\tau d_\tau
 -\sum_\tau k^2\log(h_\tau+k).
\]
Insert (GVC.24). Since $K+2$ dominates $K$, $2$ and the
coefficients of $2b'+KS$, this gives
\[
 \begin{split}
 \mu(M_\omega)\ge{}&
 -\sum_{\tau\ {\rm top}}\frac{a_T(\tau)}{k}\log T_\omega
 -C_-n_Ta\log k-C_-fk(b'+S)\log k\\
 &-C_-fS\log T_\omega-C_-fk^2\widetilde L .
 \end{split}                                                 \tag{GVC.26}
\]
This is the literal \texttt{Conclusion5}, including its signed
top sum and all four error terms.

The upper raw inequality is
\[
 \mu_{\max}(M_\omega)\le
 c_5\log(kv)\sum_\tau d_\tau
 +\sum_\tau k^2\log(h_\tau+k)
 +\log T_\omega\sum_\tau|s_\tau|.
\]
Here $\log(kv)\ge0$ because $k\ge5,\allowbreak{} v\ge1$.
We may therefore replace $c_5$ by $q_5$, and (GVC.24) yields
\[
 \begin{split}
 \mu_{\max}(M_\omega)\le{}&
 2q_5n_Ta\log(kv)
 +q_5fk(2b'+KS)\log(kv)\\
 &+f(2b'+KS)\log T_\omega
 +2fk^2\widetilde L .
 \end{split}
\]
The positive number $(q_5+1)(K+4)$ dominates
$2q_5$, $q_5(2+K)$, $2+K$ and $2$.
Thus
\[
 \begin{split}
 \mu_{\max}(M_\omega)\le{}&
 C_+n_Ta\log(kv)
 +C_+f(b'+S)\bigl(k\log(kv)+\log T_\omega\bigr)\\
 &+C_+fk^2\widetilde L ,
 \end{split}                                                 \tag{GVC.27}
\]
which is the literal \texttt{Conclusion6}.
For the actual GLS11 bound $c_5=4$, so $q_5=4$ and one may
take $C_+=5(K+4)$. The original coefficient $4$ is retained;
only the existential constant in the ensuing finite estimate
has been enlarged.

Equations (GVC.26)--(GVC.27) give independent human proofs of
the source's \texttt{ramified\allowbreak Lattice\_\allowbreak of\_\allowbreak chain\_\allowbreak estimates}
and \texttt{class\allowbreak Map\_\allowbreak estimate\_\allowbreak of\_\allowbreak partial} interfaces.
They do not claim the stronger arbitrary-space
\texttt{zeta\_\allowbreak proposition\_\allowbreak second} or
\texttt{zeta\_\allowbreak proposition\_\allowbreak third}; the simpler direct sum
argument proves exactly the two original consumers required here.
% END GVC13

% BEGIN GVC14
\paragraph{GVC14: the ordinary column statements with the exact family quantifiers.}
\label{claim:GVC14}
Suppose
$\mathrm{FamilySlopeLowerOpenPart}(\mathcal F)$ and
$\mathrm{FamilySlopeMaxOpenPart}(\mathcal F,\allowbreak{} c_5)$ hold.
Fix any real $C_n$ first, and use the constants (GVC.25).
Take the maximum of the two eventual open-part thresholds and
$K=\max(1,\allowbreak{} C_n)$. For every later $x>\mathcal F.x_1$ and every
$i$, the original $\mathcal L.\mathrm{sc}.x=x$;
GVC13 therefore applies to every word, every tuple
$(b,\widetilde L,\allowbreak{} b',\allowbreak{} S,\eta)$ satisfying its numerical predicates,
and every supplied original \texttt{Col\allowbreak Val}.
This proves exactly
\[
 \begin{split}
 &\forall C_n\ \exists C_-,\allowbreak{} x_-\ \forall x\ge x_-\
   (x>\mathcal F.x_1)\ \forall i:
             \mathrm{Conclusion5}(D,\allowbreak{} C_n,\allowbreak{} C_-,\allowbreak{} x),\\
 &\forall C_n\ \exists C_+,\allowbreak{} x_+\ \forall x\ge x_+\
   (x>\mathcal F.x_1)\ \forall i:
             \mathrm{Conclusion6}(D,\allowbreak{} C_n,\allowbreak{} C_+,\allowbreak{} x).
 \end{split}                                                 \tag{GVC.28}
\]
These are the literal \texttt{Conclusion\allowbreak F5} and
\texttt{Conclusion\allowbreak F6}, not just their values at one convenient
tuple of numerical parameters.
GLS7 and GLS11 supply the two open inequalities on every original
$D$, conditional on the same \texttt{Col\allowbreak Val}; GLS12 supplies
their family formulations with $c_5=4$.
Their geometric and published inputs retain their stated scopes.
This step does not supply \texttt{Col\allowbreak Val}; the eventual
\texttt{t\allowbreak CV} input is used explicitly below.
% END GVC14

\subsection{Uniform ordinary slopes in the global ledger}

% BEGIN GVC15
\paragraph{GVC15: the two exact size dominations.}
\label{claim:GVC15}
For the original $\mathrm{SizesAt}(\mathcal L)$, suppose
$x\ge2$ and $x^{14}\le X'$. Then
\[
 \log k\le5\log x,\qquad
 fk(b'+S)\le n_Ta,\qquad fk^2\le n_Ta.                       \tag{GVC.29}
\]
Indeed \texttt{Size\allowbreak N} and \texttt{Size\allowbreak K}, with
$0<\epsilon<1/52$, give $n\le x^3$ and $k\le x^5$.
The positive lower \texttt{Size\allowbreak TT} bound makes $n_T>0$;
as an integer it satisfies $n_T\ge1$.
The original \texttt{Size\allowbreak Del} clause gives
$kS\le b'/(xn)$, and $x,\allowbreak{} n,\allowbreak{} k\ge1$, so $S\le b'$.
Using $f=2n$,
\[
 \begin{split}
 2fk&\le4x^8\le x^{14}\le X',\\
 fk(b'+S)&\le2fkb'\le X'b'=a\le n_Ta.
 \end{split}
\]
The middle comparison holds because $x^6\ge64\ge4$.
Also
\[
 fk^2\le2x^{13}\le x^{14}\le X'\le b\le a\le n_Ta.
\]
Here $X'\le b$ follows from $x^{\epsilon/3}\le x^\epsilon$.
Finally taking logarithms of $k\le x^5$ gives the first bound
of (GVC.29). This proves \texttt{jnt\_\allowbreak sizes\_\allowbreak slopes} with its
literal ranks and does not assume a new bound on $\Delta$.
% END GVC15

% BEGIN GVC16
\paragraph{GVC16: finite substitution into the exact \texttt{Slopes\allowbreak M\_\allowbreak at}.}
\label{claim:GVC16}
Suppose on this original $D,\mathcal L$ that
\texttt{Conclusion5} and \texttt{Conclusion6} hold with constants
$C_-,\allowbreak{} C_+$, that \texttt{Nums\allowbreak L}, \texttt{Wts\allowbreak L} hold with
$C_n,\allowbreak{} c_L$ where $c_L\ge0$, and that every word has an original
\texttt{Col\allowbreak Val}. Suppose also
$\log T_\omega\le C_cm\ell$ with $C_c\ge0$ and that (GVC.29)
holds. Put
\[
 P=\max\{C_-,\allowbreak{} C_+,0\},\qquad
                  C_s=2C_c+P(12+C_c+c_L).                  \tag{GVC.30}
\]
Then, for every word,
\[
 \boxed{\quad
 -C_sn_Ta\mathcal A(1+m/k)\le\mu(M_\omega),\qquad
 \mu_{\max}(M_\omega)\le C_sn_Ta\mathcal A .
 \quad}                                                      \tag{GVC.31}
\]
These are the two clauses of the original
$\mathrm{SlopesM\_\allowbreak at}(\mathcal L,\allowbreak{} C_s)$.

We check every numerical cost. Put $Z=n_Ta\mathcal A\ge0$.
The host bound $v\ge1$ implies $\operatorname{rd}E\ge2$, so
\[
 \ell\le\mathcal A,\quad
 0\le\log k\le5\mathcal A,\quad
 0\le\log(kv)\le6\mathcal A,\quad
 0\le\log T_\omega\le C_cm\mathcal A .
\]
The third inequality uses
$\log v\le\log\operatorname{rd}E$ and
$\mathcal A=\log x+\ell+\log\operatorname{rd}E$.
The top sum in (GVC.26) has exactly $n_T$ terms and
$a_T(\tau)\le2a$; its contribution is at most
$2C_c(m/k)Z$.

All four unsigned error terms of (GVC.26) are nonnegative.
Replacing their coefficient by $P$ is therefore valid even if
one of the original constants is negative.
The first two cost at most $5PZ$ each, by (GVC.29).
Since $fkS\le fk(b'+S)\le n_Ta$,
the third costs at most $PC_c(m/k)Z$.
The fourth, with $\widetilde L=c_L\ell$, costs at most $Pc_LZ$.
Thus the total lower cost is at most
\[
 [P(10+c_L)+(2C_c+PC_c)m/k]Z .
\]
The constant $C_s$ dominates both $P(10+c_L)$ and
$2C_c+PC_c$. This proves the lower assertion in (GVC.31).

For (GVC.27), its first term costs at most $6PZ$ and
the $k\log(kv)$ part of the second costs at most another $6PZ$.
The $\log T_\omega$ part costs at most $PC_c(m/k)Z\le PC_cZ$,
because $m\le k$. The last term costs at most $Pc_LZ$.
Their sum is bounded by $P(12+C_c+c_L)Z\le C_sZ$.
This proves the upper assertion.
In particular the exact source constant (GVC.30), including
the number $12$, is justified by a visible term-by-term calculation.
% END GVC16

% BEGIN GVC17
\paragraph{GVC17: the literal \texttt{Slopes\allowbreak MF} installation.}
\label{claim:GVC17}
Suppose this same original family has
\texttt{Count\allowbreak F}, \texttt{Conclusion\allowbreak F5}, \texttt{Conclusion\allowbreak F6},
the eventual derivative bound in (GVC.18), and
\[
 \exists x_{\rm CV}\ \forall x\ge x_{\rm CV}\
       \forall i\ \forall\omega,\
              \mathrm{Nonempty}(\mathrm{ColVal}(D,\omega)).
                                                               \tag{GVC.32}
\]
Extract a nonnegative $C_c$ and its threshold from
\texttt{Count\allowbreak F}. GVC11 supplies $C_n=C_{n,0}$ and $c_L=3$,
then a threshold for the full numerical and weight packet.
Only after fixing $C_n$ extract the two constants $C_-,\allowbreak{} C_+$
and their thresholds from (GVC.28).
Define the single $C_s\ge0$ by (GVC.30).

Take the maximum of the count, numerical, open-column,
$\mathrm{ColVal}$ and $x^{14}\le X'$ thresholds, of $2$, and
of $\mathcal F.x_1+1$. For each $x$ beyond it and each original
$i$, GVC15 gives the exact size inequalities.
For each word use an inhabitant of (GVC.32), on this same
$D,\omega$, in GVC16. Its constants and threshold do not depend
on the choice of inhabitant, the word, or the member.
We obtain
\[
 \exists C_s\ge0\ \exists x_0\
   \forall x\ge x_0\ \forall i\in\mathcal F.\mathrm{Ix}(x),\
          \mathrm{SlopesM\_\allowbreak at}(\mathcal F.L(x,\allowbreak{} i),\allowbreak{} C_s).
                                                               \tag{GVC.33}
\]
This is the original \texttt{Slopes\allowbreak MF}, and proves the full
\texttt{slopes\allowbreak MF\_\allowbreak of\_\allowbreak count\_\allowbreak fam} implication.
No finite or infinite metric is transported at this step:
$M_\omega$ has throughout been exactly
$(D.\mathrm{fm}\,\omega).\mathrm{lat}(0)$, with the original
norms in \texttt{Col\allowbreak Val}. Those norms and the compatible
column supplier are the retained GLS/RSL hypotheses.
% END GVC17

% BEGIN GVC18
\paragraph{GVC18: remove a redundant derivative supplier by a literal family projection.}
\label{claim:GVC18}
In the actual \texttt{Fam\allowbreak P}, the separate eventual derivative
hypothesis of GVC11 and GVC17 is already supplied.
Indeed \texttt{Fm.\allowbreak sizes} gives
$\mathrm{SizesAt}(\mathcal F.L(x,\allowbreak{} i))$ for every
$x>\mathcal F.x_1$. The seventh conjunct of \texttt{Sizes\allowbreak At}
is \texttt{Size\allowbreak QD}, whose sixth conjunct is precisely
\[
 |\mathcal L.\mathrm{logDerivativeAtOne}|
          \le\mathcal L.\mathrm{sc}.x^{
                         \mathcal L.\mathrm{sc}.\epsilon}.
\]
For $x\ge\mathcal F.x_1+1$ the original equalities
$\mathcal L.\mathrm{sc}.x=x$ and
$\mathcal L.\mathrm{sc}.\epsilon=\epsilon$ identify its right
side with $\ell=x^\epsilon$.
This proves the exact eventual hypothesis (GVC.18), with one
threshold before every $i$. It uses neither a new derivative
estimate nor a different host.

Consequently the exact outputs on a supplied original family are
\[
 \begin{gathered}
 \mathrm{CountF}\Longrightarrow
 \bigl[\text{the complete supplier (GVC.20)}\bigr],\\
 \mathrm{CountF}\wedge\mathrm{LowF}\wedge
 \mathrm{MaxF}(c_5)\wedge\mathrm{CVF}
       \Longrightarrow\mathrm{SlopesMF}.
 \end{gathered}                                               \tag{GVC.34}
\]
Here $\mathrm{LowF}$ and $\mathrm{MaxF}(c_5)$ abbreviate the
two original family open-part predicates, and $\mathrm{CVF}$
abbreviates (GVC.32).
For the selected original route, GLS12 supplies the open bounds
with $c_5=4$, and its \texttt{t\allowbreak CV} supplies (GVC.32) under the
same compatible \texttt{Cols\allowbreak Tied\allowbreak Sym}/\texttt{adm3} data.
Combining those inputs with GVC7 therefore closes both count and
ordinary-slope entries of the one-family ledger.
The choice $c_L=3$ in the first line is also available to the
level-prime branch.

The word \emph{supplied} matters. This argument projects properties
from the existing \texttt{Fam\allowbreak P}; it does not construct that family
or prove the analytic derivative estimate independently of its
size fields. The common selector, real-place second layer,
actual forms and mass/volume interfaces keep their inherited
conditional scope. The complementary quotient slopes, frame
estimate, loss sum and terminal \texttt{Chain\allowbreak Tied} constructor
are separate remaining joins. No statement here closes the
whole RH or unrestricted Dirichlet-GRH human proof.
% END GVC18


\section{Installing the original quotient slopes and frame inequality}
\label{sec:global-assembly-v034}

\status{The following arguments install the finite-place and frame inputs
on the original supplied family.  They keep the arithmetic inputs
\texttt{Link\allowbreak G}, \texttt{Link\allowbreak S}, \texttt{Link\allowbreak IRR}, \texttt{Link\allowbreak DIV} and
\texttt{Link\allowbreak Mu} as inputs.  In particular, they do not manufacture the
common second layer or the spectral estimates.  The remaining final
assembly is specified at the end of the chapter.}

Fix an original family $\mathcal F:\texttt{Fam\allowbreak P}(\epsilon)$,
$0<\epsilon<1/52$, and its fixed constants.  A statement asserted
eventually means that a single threshold is chosen before the index
$i\in\mathcal F.I_x$, the word $u$, and every finite place.  We write
$D=\mathcal F.D(x,\allowbreak{} i)$ and $L=\mathcal F.L(x,\allowbreak{} i)$ only after that choice.
The scales retain their original meanings:
\[
 b=e^{x^\epsilon},\qquad a=xb,\qquad
 X'=e^{x^{\epsilon/3}},\qquad b'=a/X',\qquad
 X''=xnX'.                                                    \tag{GFA.1}
\]
For a word let $M_u$ be the original ramified reference lattice,
$N_u^0=L.\texttt{Nu\allowbreak Lat}(u)$, and
$N_u=D.\texttt{div.\allowbreak Nu}(u)=I_uN_u^0$, with the unchanged metrics of MT1.
Write $d_u=D.\texttt{du\allowbreak S}(u)$, $r_u=\operatorname{rank}N_u$ and
$T=T_u$.  The full source lattice is denoted $A_u=(D.\texttt{fm}(u)).
\texttt{lat}(0)$.  Thus $A_u$ and $M_u$ are not confused.

% BEGIN GFA1
\paragraph{GFA1: the exact type-prime cost, with constant $36$.}
\label{claim:GFA1}
Assume the original ramified-model identities and congruence-length
bound supplied by RSL9--RSL10 and GLE11, together with
\texttt{Tu\allowbreak Size}.  GLE2 activates the guarded statements simultaneously.
The full numerical and weight windows from GVC11's
\texttt{Nums\allowbreak L}/\texttt{Wts\allowbreak L} installation may be taken with $c_L=3$.
Here the needed logarithmic-derivative bound is already the
\texttt{Dl\allowbreak Size} clause of the original size record.

We first check the numerical factor in the congruence estimate.  The
actual top-coordinate identity gives $a_{\max}\le2a$.  Every other
entry of the lifted weight at a top has absolute value at most $2b'$:
the interior and last entries have this bound by \texttt{Wts\allowbreak L}; the
absolute value of each appended block entry lies between the
bottom weight and the largest non-thick $w_j\le2b'$.  The determinant shift
$i(t+1)/2$ is at most $N(m)\le b'$.  Consequently the original
shifted absolute weight $\operatorname{Wt}$ is at most $3b'$.
Also $k\le N(m)\le b'$ and $b_{\rm top}\le2b'$.
At any selected place over $T$, its summand therefore satisfies
\[
 \frac{a_{\max}}k+2\operatorname{Wt}+k+b_{\rm top}
 \le 2a/k+9b'\le9(a/k+b').                              \tag{GFA.2}
\]
The map from the selected local positions to the original tops is
injective (it is the restriction of the actual $\tau$-bijection).
There are at most $|TT|$ such positions.  Hence (GFA.2) proves
\texttt{Level\allowbreak Prime\allowbreak Size\allowbreak Bounds} with $C_1=9$ on the original data.

Let $\ell_w=\log N\mathfrak p_w/[K:\mathbb Q]$ and
$e_w=v_w(T)$ for $w\mid T$, where $K=D.\texttt{what}_2.LL$.
For a finite set $F$ of such places put
$B=C_1|TT|(a/k+b')$.  GLE11 bounds the actual local quotient length by
\[
 n_w\le e_w m d_u B.
\]
The source's \texttt{Nu\allowbreak Def} and \texttt{Len\allowbreak At\allowbreak T} identify this same
integer with $e_wL.\texttt{len\allowbreak MN}(u,\allowbreak{} w)$; no alternative quotient is
substituted.  The degree identity of LC3 now gives
\[
 \sum_{w\in F}2\ell_we_wL.\texttt{len\allowbreak MN}(u,\allowbreak{} w)
 \le2md_uB\sum_{w\in F}e_w\ell_w
 \le2md_uB\log T.                                      \tag{GFA.3}
\]
The setting supplies $T>2fN(m)+1\ge2f+1$.  By Bernoulli,
\[
 q_u:=\left(1-\frac1{T-1}\right)^f
 \ge1-\frac f{T-1}\ge\frac12.
\]
There is a rank identification to check before this division.
GLE11 gives the original inclusion $N_\lambda\subset S_{\rm ram}$.
On this same supplied family it activates RSL11, hence fullness
of $N_u^0$ and the actual \texttt{Nu\allowbreak Full} property of $N_u$.
The source identities
\texttt{Nu\allowbreak V\_\allowbreak finrank} and \texttt{Nu\allowbreak Full.\allowbreak rank} then identify the
rank of the actual $N_u$ with the ambient retained dimension
$(D.\texttt{mem}(u)).A.\texttt{du}$.  RSL10 computes that dimension
as $d_uq_u$.  Thus $r_u=d_uq_u$ for the very lattice used here.
Dividing (GFA.3) by $r_u$ yields
\[
 \boxed{\quad
 \frac1{r_u}\sum_{w\in F}2\ell_we_wL.\texttt{len\allowbreak MN}(u,\allowbreak{} w)
 \le36|TT|m(a/k+b')\log T.
 \quad}                                                 \tag{GFA.4}
\]
If $d_u=0$, the numerator and retained rank vanish; the formal
total-division convention gives zero on the left.  Every later use
has $r_u>0$, so uses the ordinary division just written.

The other four clauses of the original \texttt{adm26} are also
accounted for.  Its rank, mean and maximal-slope clauses are exactly
RSL10.  For the local inclusion at $T$, \texttt{Nu\allowbreak At\allowbreak TDef} maps
$(N_u^0)_w$ into the same congruence lattice $N_\lambda$.
GLE11 gives $N_\lambda\subset S_0$; the original localized class
lattice identity pulls this back to $(M_u)_w$.  Together with
(GFA.4), this proves \texttt{Level\allowbreak TFamily} with constant $36$.
Taking the maximum of the setting, numerical, weight, and
congruence thresholds occurs before $i,\allowbreak{} u,\allowbreak{} w$.
% END GFA1

% BEGIN GFA2
\paragraph{GFA2: the required total correction bound directly from its definitions.}
\label{claim:GFA2}
The quotient-slope argument needs only
\[
 \mathrm{Ta}:=\sum_{w\in L.\mathrm{supp}}2\ell_we_w\tau_w\le b^2.
                                                               \tag{GFA.5}
\]
We prove this on the actual family without assuming the stronger
\texttt{Family\allowbreak Weight\allowbreak Sum\allowbreak Bound}.  This is a deliberately coarse
estimate for this particular consumer.

Set $S=\exp(x^{\epsilon/2})$.  The original rank windows, GII2--GII3,
and the elementary comparison of fixed powers with exponentials give,
eventually and uniformly,
\[
\begin{gathered}
 S\ge40,\quad a\ge1,\quad
 f,\allowbreak{} m,\allowbreak{} n,\allowbreak{} k,\allowbreak{} N_1,\allowbreak{} r_G,\allowbreak{} x^{26},\allowbreak{} X',\log a/\log2\le S,\\
 \max(1,|C_e|,\text{each fixed local constant})\le S,
 \qquad 0\le c_G=4L.\mathrm{Sc}\le4b'/x\le5a,\quad b'\le a,\\
 \Lambda:=C_{\rm La}N_1^2\log a\le S^4,
 \qquad 0\le c_0=2b'+2c_G\le12a .
\end{gathered}                                               \tag{GFA.6}
\]
For every $p<X'$, GII2 gives
$G_p\le C_e x^{\epsilon/8}\log x$.  Therefore
\[
 \log(p^{G_p})\le C_e x^{11\epsilon/24}\log x
 \le x^{\epsilon/2},\qquad Q_p=p^{G_p}\le S.                \tag{GFA.7}
\]
The strict exponent difference is $\epsilon/24>0$; the threshold
depends only on the fixed $C_e,\epsilon$.  This is also valid for
the original natural ceilings, already included in GII2.

If the original set $VV(p)$ is nonempty, its definition implies that
$p$ is prime, $p<X'$, and it is inert.  The exact orbit identification
then gives $n_p=\operatorname{ord}(h_p)\le n\le S$.
There are at most $n$ representatives and at most $n$ tops in any
orbit.  Since $\delta_1\le1$,
$N^+=40n_p\delta_1\le S^2$.  We have
$\log_p a\le S$, $p^{-L_1}\le1$, $pk\ge1$, and
$0\le b_t\le2b'\le2a$ at every top.  These are the only
orbit and weight facts used below; no new good-orbit estimate is
assumed.

Here are upper bounds for every term in the literal definitions.
They keep the two branches of $D_p,\allowbreak{} t_v$ distinct:
\[
\begin{array}{c|c|l}
\text{quantity}&\text{upper bound}&\text{reason before rounding}\\\hline
 D^{\rm I}_p&S^{11}&
 C_{DI}N^+x^{26}(N_1^5+k^3)\log_pa\le2S^{10}\\
 D^{\rm II}_p&aS^{11}&
 3S^9+10aS^2+12aS^3\\
 t^{\rm II}_v&aS^8&
 S^3(S^4a+a+12a+S^2)+10aS^2\\
 c_6^{\rm I}&aS^{12}&
 8(S^{11}+1)+S^5+20aS^2+10aS\\
 t^{\rm I}_v&aS^{16}&
 2aS+2aS^7+2aS^{15}+10aS .
\end{array}                                                  \tag{GFA.8}
\]
For the first two rows use $\lceil z\rceil_+\le\max(z,0)+1$.
The inequalities displayed follow for $S\ge40,\allowbreak{} a\ge1$.
For the last row, the predecessor sum has at most $S$ terms and
each term is bounded by its bottom weight $2a$; if its distance
is zero the term is zero.  Its second summand is at most
$S^3(a+S^4a)$.  The third is at most
$S^2(S\cdot aS^{12}+2S^2+5a)$.  This proves the final row,
including the actual minimum and distance convention.

The stored ramification exponent is
$e_p=\lceil C_{\rm ram} f k^2\log_pa\rceil_+$, where
$C_{\rm ram}$ is the fixed original \texttt{c\allowbreak X.\allowbreak Ce} (distinct
from the orbit constant $C_e$).  Thus $e_p\le2S^5$.
The original definition
\[
 \tau_p=e_p+\sum_{v\in VV(p)}m(\lceil t_v\rceil_++D_p+1)
\]
and (GFA.8) imply $0\le\tau_p\le aS^{20}$.  When $VV(p)$
is empty the definition instead sets $\tau_p=0$.
In particular every $\tau_p$ is a nonnegative integer.
The original fill equality is $\tau_w=\tau_{p(w)}$.
For the finite support, LC3, nonnegativity, and (GFA.7) now give
\[
 \mathrm{Ta}\le2\sum_{p<X'}\log p\,\tau_p
 \le4aS^{22}\le aS^{24}.
\]
Indeed there are at most $\lceil X'\rceil\le2S$ eligible primes,
and $\log p\le S$.  Finally
\[
 \log(aS^{24})=x^\epsilon+\log x+24x^{\epsilon/2}
 \le2x^\epsilon=\log b^2
\]
eventually.  This proves (GFA.5) with one threshold for every
original member and word.  It neither proves nor needs the
sharper Mertens-weighted estimate of the alternative source route.
% END GFA2

% BEGIN GFA3
\paragraph{GFA3: exact mean and maximum slopes from all finite places.}
\label{claim:GFA3}
Let $M$ be a full positive-rank hermitian lattice over $K$, and
$N^0\subset M$ a commensurable lattice with the same metric.
For $n_w=\operatorname{length}_{\mathcal O_w}(M_w/N_w^0)$,
local elementary divisors and the norm of the finite quotient give
\[
 \mu(N^0)=\mu(M)-\frac1r\sum_w2\ell_wn_w.                  \tag{GFA.9}
\]
Only finitely many $n_w$ are nonzero: a common nonzero denominator
contains a multiple of $M$ in $N^0$, and its prime divisors contain
the support.  Multiplication by
$I=\prod_w\mathfrak p_w^{-\lfloor e_ws_w\rfloor}$ gives, by MT2,
\[
 \mu(IN^0)=\mu(M)+\sum_w2\ell_w\lfloor e_ws_w\rfloor
                     -\frac1r\sum_w2\ell_wn_w.              \tag{GFA.10}
\]

For completeness the corresponding maximal-slope inequality is
also local to global.  If $N_w^0\subset\mathfrak p_w^{k_w}M_w$
at a finite set, with ordinary inclusion elsewhere, then
$N^0\subset JM$ for $J=\prod_w\mathfrak p_w^{k_w}$.
To see this, localize the proposed quotient at any maximal ideal;
all other prime-ideal factors become units and the stated local
inclusion applies.  A nonzero quotient element would have a proper
annihilator contained in some maximal ideal, contradicting these
zero localizations.  Taking every nonzero rational subspace, its
induced sublattice in $N^0$ has no larger degree than its induced
sublattice in $JM$.  Taking suprema, and using MT2 on each such
subspace, proves
\[
 \mu_{\max}(IN^0)\le\mu_{\max}(M)
       +\sum_w2\ell_w(\lfloor e_ws_w\rfloor-k_w).             \tag{GFA.11}
\]
This argument also verifies that the scalar-ideal shift for maximum
slope uses the same normalized squared metrics as the mean shift.

On the original records, $\texttt{Nu\allowbreak Def}$ reads
$n_w=e_wL.\texttt{len\allowbreak MN}(u,\allowbreak{} w)$, and \texttt{Nprime\allowbreak Def} is the
literal equality $N_u=I_uN_u^0$.  RSL11 provides commensurability
and fullness; its proof preserves the ambient $\texttt{Nu.\allowbreak V}$ and
the reference $\texttt{Mram\allowbreak L}$.  Thus (GFA.10)--(GFA.11) apply
to these objects, with $r=r_u$, once the support and inclusions
checked next are supplied.
% END GFA3

% BEGIN GFA4
\paragraph{GFA4: the support, integer exponents, and each word's type prime.}
\label{claim:GFA4}
Assume \texttt{Link\allowbreak G}, \texttt{Tu\allowbreak Size}, and the actual family fill,
local-lattice and error-definition fields.  Since $m\to\infty$
in its size window, eventually $m\ge6$ and
\[
 T_u\ge b^{m/2}\ge b^3\ge X'.                              \tag{GFA.12}
\]
The definition of $VV(p)$ forces $p<X'$ when it is nonempty.
Consequently $\tau_{T_u}=0$ for every word, including when
different words have different $T_u$.
GFA2 proves the exact integrality $e_w\tau_w\in\mathbb Z_{\ge0}$,
so $\lfloor e_w\tau_w\rfloor_+=e_w\tau_w$.
The original $\mathrm{io}_p$ is nonnegative: its defining
$\mathrm{Sig}$ is a sum of nonnegative $\kappa$ values.
In the explicit minimum branch this follows from
$\mathcal F.\texttt{GK\_\allowbreak nonneg}$.  In the other branch the
original $\texttt{Extra.\allowbreak kap\allowbreak I\_\allowbreak def}$ makes $\kappa$ the infimum
of admissible nonnegative values.  That set is nonempty:
$\kappa=1$ permits the zero-dimensional choice, whose required
rank bound is $(1-\kappa)\dim=0$ and whose remaining indexed
conditions are empty.  Thus its infimum is also nonnegative.
Furthermore
$S_v\ge0$ because it is the ceiling of a sum of nonnegative
parts.  The empty-set branch is zero.  The fill equality therefore
gives $\mathrm{io}_w\ge0$ and $\mathrm{IOsum}\ge0$.

At $p(w)\ne T_u$, \texttt{Link\allowbreak G} supplies precisely
\[
 (N_u^0)_w\subset\mathfrak p_w^{e_w\tau_w}(M_u)_w,
 \qquad L.\mathrm{lenMN}(u,\allowbreak{} w)\le d_u(\tau_w+\mathrm{io}_w).
                                                               \tag{GFA.13}
\]
At $T_u$ use GFA1's inclusion, which has exponent zero, and
retain its separate cost (GFA.4).

It remains to justify deletion of every place outside the stored
finite support.  At a regular such place, support saturation
implies $p(w)>X''\ge X'$, whence $VV(p(w))$ is empty.
The actual \texttt{lats0} field then says
$L.\mathrm{lenMN}(u,\allowbreak{} w)=0$.  At a type place, use the error field
at the original lift $j=1$ (it exists since $m\ge2$).  If
\texttt{seen\allowbreak P} were true, its $\mathrm{EatT}$ formula would have
the strictly positive summand $fN(m)$ and only nonnegative other
summands.  This contradicts $L.E(u,1,\allowbreak{} w)=0$ outside support.
Thus \texttt{seen\allowbreak P} is false.  The actual type-place set is then
empty, its exponent is zero and its ordered operator product is
the identity.  Injectivity of the original scalar-extension map
gives $N_\lambda=S_{\rm ram}$.  Its quotient has length zero,
and \texttt{Nu\allowbreak At\allowbreak TDef} identifies it with the required local
length.  This proves the full off-support assertion used in
(GFA.10), at regular and exceptional places alike.
% END GFA4

% BEGIN GFA5
\paragraph{GFA5: the literal \texttt{Places\allowbreak F} and \texttt{Slopes\allowbreak NF} inputs.}
\label{claim:GFA5}
Assume, on the same family, \texttt{Level\allowbreak TFamily},
\texttt{Nu\allowbreak Commens\allowbreak F}, \texttt{Mu\allowbreak Rank}, \texttt{Tu\allowbreak Size},
\texttt{Link\allowbreak G} and \texttt{Link\allowbreak Mu}.  GFA1 supplies the first
with constant $C_{26}=36$ under its stated inputs; alternatively
retain any supplied $C_{26}$.  Write $C_\mu$ for the original
\texttt{Link\allowbreak Mu} constant and choose
$C=\max(C_{26},\allowbreak{} C_\mu,4)$ before the final threshold.

Put
\[
\begin{gathered}
\delta_u=1-r_u/d_u,\qquad
R_u=\sum_{w\in\mathrm{supp}}2\ell_w
                    (e_ws_{u,\allowbreak{} w}-\lfloor e_ws_{u,\allowbreak{} w}\rfloor),\\
B_u=\frac{\sum_{w\in\mathrm{supp},\ p(w)=T_u}
                 2\ell_we_wL.\mathrm{lenMN}(u,\allowbreak{} w)}{r_u}.
\end{gathered}
                                                               \tag{GFA.14}
\]
Here $r_u>0$, $d_u>0$, $0\le\delta_u<1$ and
$\delta_u\le2f/T_u$ by RSL10 and Bernoulli.  GFA1 gives
$B_u\le C|TT|m(a/k+b')\log T_u$.
The error $R_u$ is nonnegative, and it is exactly the negative of
the signed metric discrepancy in MT3, not merely another upper
bound for it.  Thus \texttt{Link\allowbreak Mu} gives $R_u\le CX''$.

The parameter conditions needed to absorb the retained-rank error
follow directly from (GFA.12).  Eventually $b\ge4$ and $16f\le b$.
Then $2f/T_u\le1/2$, and
\[
 \frac{4fb^2}{T_u}\le\frac{4f}{b}\le\frac14,
 \qquad
 \frac{f\log(T_u-1)}{T_u-2}
 \le\frac{4f}{\sqrt{T_u}}\le\frac14.                       \tag{GFA.15}
\]
For the second inequality use $\log t\le2\sqrt t$ for $t\ge1$,
$T_u-2\ge T_u/2$, and $\sqrt{T_u}\ge b$.
In particular the sum in (GFA.15) is at most $1$, with room to spare.

The exact definitions give
\[
 \sum_w2\ell_we_ws_{u,\allowbreak{} w}
       =\mathrm{Ga}+\mathrm{Fo}-\mathrm{Ee}_u+\mathrm{Ta}.
\]
At the regular places use (GFA.13), and at the type places
$\tau_w=0$.  Equation (GFA.10) and RSL10 therefore give
\[
\begin{split}
 \mu(N_u)\ge{}&\mu(A_u)+\mathrm{Ga}+\mathrm{Fo}-\mathrm{Ee}_u\\
 &-\frac{2\mathrm{IOsum}}{1-\delta_u}
   -\frac{\delta_u\mathrm{Ta}}{1-\delta_u}
   -\frac{f\log(T_u-1)}{T_u-2}-B_u-R_u\\
 \ge{}&\mu(A_u)+\mathrm{Ga}+\mathrm{Fo}-\mathrm{Ee}_u
             -4\mathrm{IOsum}-B_u-R_u-1 .
\end{split}                                                    \tag{GFA.16}
\]
GFA2 supplies $\mathrm{Ta}\le b^2$ in this calculation.
The factor $2$ before $\mathrm{IOsum}$ is the squared-norm
prime-place normalization, and the second factor $2$ comes from
$1/(1-\delta_u)\le2$.  Both factors are retained.
For the maximum slope take $k_w=e_w\tau_w$ in (GFA.11).
The floor is at most $e_ws_{u,\allowbreak{} w}$, and RSL10 bounds
$\mu_{\max}(M_u)$ by $\mu_{\max}(A_u)$.  Hence
\[
 \mu_{\max}(N_u)\le
 \mu_{\max}(A_u)+\mathrm{Ga}+\mathrm{Fo}-\mathrm{Ee}_u.
                                                               \tag{GFA.17}
\]

For clarity, the witnesses of the source's existential
\texttt{Places\_\allowbreak at} predicate are
\[
 s'_{u,\allowbreak{} w}=\frac{\lfloor e_ws_{u,\allowbreak{} w}\rfloor}{e_w},\quad
 \delta_u,\quad \mu(M_u),\quad\mu_{\max}(M_u),\quad
 P_T=L_T=B_u,\quad R_d=R_X=R_u.
\]
Equations (GFA.10), (GFA.13)--(GFA.16) give its two lower-bound
packages; (GFA.11) gives its upper-bound package.  The last two
inequalities are precisely the stated bounds for $B_u,\allowbreak{} R_u$.
This checks every witness on the original objects and proves
\texttt{Places\allowbreak F}.

With
\[
 \mathrm{Lo}_u=C\mathrm{IOsum}
    +C|TT|m(a/k+b')\log T_u+CX''+1,
\]
(GFA.16) proves the lower clause of \texttt{Slopes\allowbreak N\_\allowbreak at}.
Its upper clause follows from (GFA.17) by adding the nonnegative
allowance $CmB_0$.  This is the literal \texttt{Slopes\allowbreak NF}.
All estimates used finitely many uniform thresholds; their maximum
is chosen before any $i,\allowbreak{} u,\allowbreak{} w$.  In particular the shortcut in GFA2
does not change the family or any of these quantifiers.
% END GFA5

% BEGIN GFA6
\paragraph{GFA6: positive target rank and the exact lattice dual.}
\label{claim:GFA6}
Assume \texttt{Link\allowbreak S}, \texttt{Chosen\allowbreak Proviso} with $C_0>0$,
\texttt{Link\allowbreak IRR} and \texttt{Link\allowbreak DIV}.  The proviso gives
$d_u\ge D_*/(C_0\Gamma^2)>0$.  The spectral diagonal bound has
\[
 \|K_{uu}-1\|_{\rm HS}^2\le e_0d_u,
 \qquad 0\le e_0,\quad\sqrt{e_0}\le e^{-x^{\epsilon/8}}<1.
\]
If the full analytic target had dimension zero, its period maps
would vanish and $K_{uu}=0$.  The left side would be $d_u$,
contradicting $e_0<1$.  Thus $D_{\rm all}>0$.
The actual target defect obeys
$0\le\delta=D.\texttt{what}_2.\texttt{delp}
 \le e^{-x^{e_z/7}}<1$ by \texttt{Link\allowbreak IRR}.
The stored equality $D_{\rm rows}=(1-\delta)D_{\rm all}$ and
the fullness of $RR$ from \texttt{Link\allowbreak DIV} give
$\operatorname{rank}RR=D_{\rm rows}>0$.
This proves \texttt{RRRank\allowbreak F} on the same family.

The dual used below is the literal algebraic dual lattice
$RR^\vee=\{\lambda:\lambda(RR)\subset\mathcal O_K\}$ with its
stored dual metric.  A finite generating set and a nonzero
denominator show that this is finitely generated and spans the
full dual space.  There is no trace dual or different ideal in
this assertion.  For a saturated sublattice the determinant
identity for a lattice and its metric quotient proves
$\deg(Q^\vee)=-\deg(Q)$.  Annihilators identify saturated
subobjects with dual quotients, reversing their ranks and
slopes.  Taking the largest and smallest polygon slopes gives
\[
 \mu_{\min}(RR^\vee)=-\mu_{\max}(RR).
\]
For each full positive-rank $N_u$, its mean slope is the weighted
mean of its polygon slopes and is at most its maximum.
Together with \texttt{Nu\allowbreak Full\allowbreak F} and \texttt{Mu\allowbreak Rank}, these are
exactly the two fields of \texttt{Lat\_\allowbreak at}; a common maximum
threshold proves \texttt{Lat\allowbreak F}.
% END GFA6

% BEGIN GFA7
\paragraph{GFA7: putting the original maps in the common Gram calculation.}
\label{claim:GFA7}
Keep the hypotheses of GFA6, now with $C_0\ge1$, and supply
the original full and retained lattices.  The period normalization is
\[
 c_u=\frac{\operatorname{vol}(N_u^{\rm an})
                   \operatorname{vol}(H_u^{\rm an})}{\dim\Pi_u}.
\]
It is nonnegative by the source's positive volumes and dimension
formula.  If $c_u=0$, the source's normalized kernel is zero,
again contradicting its diagonal defect $e_0<1$ and $d_u>0$.
Thus $c_u>0$ for every word.  With $A_u^{\rm an}=c_u^{-1/2}P_u$,
the Gram blocks in GF1 are exactly
$(A_u^{\rm an})^*A_{u'}^{\rm an}$.

We spell out the coordinate identification needed to use GF8.
Choose an orthonormal basis of the original retained target
subspace and extend it to the full analytic target.  Its first
block has dimension $\operatorname{rank}RR$, by GFA6.
Choose orthonormal bases of every full source and every retained
source.  The inclusion matrix $V_u$ of the retained space has
$V_u^*V_u=1$.  The original \texttt{Bessel\allowbreak Id} in
\texttt{Link\allowbreak DIV} identifies the projected analytic map with the
scalar extension of the original rational map
\[
 B_u:N_u\longrightarrow RR^\vee.
\]
The scalar in this identity is nonzero, as explicitly required
there; absorb it and $c_u^{1/2}$ into the invertible source
coordinate matrix.  The target coordinate change is the inverse
transpose of the primal basis change, as required by the dual
basis.  Thus, at the actual designated complex embedding,
\[
 YB_u^\sigma=P_0 A_u^{\rm an}J_uX_u,
 \qquad J_u^*J_u=1,\quad X_u,\allowbreak{} Y\text{ invertible}.             \tag{GFA.18}
\]
If the representative infinite-place embedding is the conjugate
of the designated embedding, conjugate every matrix in this
equality.  Conjugation preserves inverses, products, Gram
matrices and squared norms, so the same equality holds there.
This verifies the source's possible identity-or-conjugation
choice, without choosing unrelated embeddings for different
words.

The remaining geometric requirements in GF8 are also literal
fields of the supplied records.  \texttt{One\allowbreak Map\allowbreak Into\allowbreak CC} gives
$B_u(N_u)\subset RR^\vee$.  The nonzero-map height bound of
\texttt{Link\allowbreak DIV} gives the normalized squared height
$h(B_u)\le D.\texttt{div.\allowbreak h}(u)$ for every nonzero $B_u$;
the zero maps impose no height condition.  The target is the
same $RR^\vee$ as in GFA6.  Its retained fraction is
$1-\delta$.  The isometries $J_u$ use the original retained
source ranks $r_u$, not the full $d_u$.  Finally
\texttt{Link\allowbreak S}, activated by the same \texttt{Chosen\allowbreak Proviso},
supplies all the mixed and diagonal Gram estimates and
$D_{\rm all}\le(1+\eta_1)D_*$.
Replacing $\eta_0,\eta_1$ by their nonnegative parts weakens
these inequalities and preserves their uniform upper bounds.
Consequently every object, normalization, and map hypothesis
in GF1--GF8 has now been matched to this one original family.
% END GFA7

% BEGIN GFA8
\paragraph{GFA8: the literal frame inequality for every height parameter.}
\label{claim:GFA8}
Let $\sigma=x^{-C_H}$, and suppose $0<\sigma\le1/20$,
$0\le d\le1/4$, and
$0\le\eta_0,\eta_1,\sqrt{e_0}\le\sigma/2$.
The proviso and code condition give
\[
 \sum_ud_u\ge\frac{D_*}{\sigma},
 \qquad \sum_u(d_u-r_u)\le d\sum_ud_u.                       \tag{GFA.19}
\]
The first follows by multiplying
$d_u\ge D_*/(C_0\Gamma^2)$ by
$\#\mathrm{Word}\ge C_0\Gamma^2x^{C_H}$.

Define the actual frame quantities
\[
 s=\frac{\sum_ur_u(\mu(N_u)-h_u)}{\sum_ur_u},\qquad
 u_* =\max_u(\mu_{\max}(N_u)-h_u)-s\ge0.                    \tag{GFA.20}
\]
The nonnegativity follows from GFA6 and positive ranks; the
denominator is positive because the word set is nonempty.
Every adjusted source polygon slope is at most $s+u_*$, and
their weighted average is exactly $s$.
For \emph{any} $H>0$ satisfying
$s-\mu_{\min}(RR^\vee)\le H$, all target slopes are at least
$s-H$.  Apply GF8 with the objects and heights in GFA7:
\[
 s-\frac{4u_*H/(u_*+H)+(5\sigma+4d+4\sqrt{e_0})H}{1-\delta}
       \le\mu(RR^\vee).                                    \tag{GFA.21}
\]
For $u_*>0$, integration of $(1+t)^{-1}$ on $[0,\allowbreak{} H/u_*]$
gives
$H/(u_*+H)\le\log(1+H/u_*)$.
If $u_*=0$, both multiplied expressions are zero; this also
agrees with the source's total real-division convention.
Since $4d\le3\sqrt d$ for $d\le1/4$ and
$\sqrt{e_0}\le\sqrt{2e_0}$, (GFA.21) implies
\[
 s-\log\#\mathrm{Word}
 -\frac{4u_*\log(1+H/u_*)+
        (5\sigma+3\sqrt d+4\sqrt{2e_0})H}{1-\delta}
 \le\mu(RR^\vee).                                         \tag{GFA.22}
\]
The extra logarithm is nonnegative because the word set is
nonempty.  This proves the exact \texttt{Frame\_\allowbreak at} predicate,
with witnesses $\eta=5\sigma+3\sqrt d$, $e=2e_0$.
The proof does not restrict $H$ to one later application or
assume a numerical upper bound for it.  In particular it supplies
the full universal quantifier required by that predicate.
% END GFA8

% BEGIN GFA9
\paragraph{GFA9: one pair of frame errors before the family member.}
\label{claim:GFA9}
Assume the preceding family inputs, \texttt{Tu\allowbreak Size}, and the
actual ramified rank formula.  Write
$E(x)=\exp(-x^{\epsilon/8})$.
Bernoulli and $T_u\ge b$ give
\[
 d_u-r_u\le(2f/b)d_u.
\]
The rank windows imply $f\le x^5$; hence, eventually,
$2f/b\le E(x)^2$, uniformly for every $i,\allowbreak{} u$.
Indeed
$\log(2x^5)-x^\epsilon\le-2x^{\epsilon/8}$ eventually.
Thus in GFA8 we can choose
\[
 d=E(x)^2,\qquad e_B=2E(x)^2,\qquad
 \eta_B=5x^{-C_H}+3E(x),                                  \tag{GFA.23}
\]
before choosing $i$.  The spectral errors satisfy
$\eta_0,\eta_1,\sqrt{e_0}\le E(x)$ on every member.
For fixed $C_H>0$, enlarge a single threshold until
\[
 x^{-C_H}\le1/20,\quad E(x)\le x^{-C_H}/2,
 \quad 2E(x)^2\le1/4,
 \quad (3+4\sqrt2)E(x)x^{C_H}\le1.
\]
All follow from exponential domination of a fixed power.
Now $\eta,\allowbreak{} e$ from GFA8 obey
$0\le\eta\le\eta_B$, $0\le e\le e_B$, and
\[
 (\eta_B+4\sqrt{e_B})x^{C_H}
 =5+(3+4\sqrt2)E(x)x^{C_H}\le6.                            \tag{GFA.24}
\]
Equations (GFA.22)--(GFA.24) are the complete literal
$\texttt{Frame\allowbreak F}(\mathcal F,\allowbreak{} C_H)$.
Its source discard $d$ and target discard $\delta$ remain
separate: only the latter occurs in the denominator
$(1-\delta)^{-1}$.  The fixed normalization constants
$C_{\rm abs}=2$, $C_0=3$, $C_5=4$, $e_D=2$, $e_L=0$
have not been altered.
% END GFA9

% BEGIN GFA10
\paragraph{GFA10: the original thick, top, and thin nodes.}
\label{claim:GFA10}
We prove \texttt{Nodes\allowbreak F} for any supplied original family.
Put $w''=b'/(xn)$ and $\eta=L.\mathrm{etaSz}$.
The size record gives $w''\ge\sqrt b\ge1$, so $b'\ge1$.
The parameter condition $t\eta\le w''/x$ and $t\ge1$ give
$\eta\le w''/x\le b'/3$ for $x\ge3$.
The nonnegative second summand of the original $\rho$ inequality
shows
\[
 \frac{N_*x^\epsilon\Lambda}{k}\le\frac{x^{-2}}4\le\frac1{36}.
\]
Consequently the original bottom window gives
$b_t\le(1+1/36)b'$.  It follows that
\[
 w''+\eta\le2b',\qquad
 b_t+b'/x+\eta\le\frac{61}{36}b'\le2b'.                     \tag{GFA.25}
\]

Here is the exact place dictionary in \texttt{Nodes\_\allowbreak at}.
For $\rho=((0,\rho_2),\allowbreak{} h)\in\Phi$, set
\[
 g(\rho)=((-1)^{\rho_2},1,\iota h)\in D.R.\Psi,
 \qquad \tau=D.eTau^{-1}(g(\rho)).
\]
Its code is exactly $\texttt{mcode}(\rho)$, so the stipulated
reading of $w(j,\rho)$ is the original $w_{j+1}(\tau)$.
The thick condition says $\rho=((0,1),\allowbreak{} t+u_0)$ with
$t\in T_2$.  The source hypotheses identify $T_2$ with the
digits of the original integer second set and identify $u_0$
by its injective integer label.  Therefore $g(\rho)$ is thick
exactly when $\rho$ is thick.  This follows by comparing the
two sign coordinates and the injective $\iota$ coordinate;
it does not follow merely from equality of cardinalities.

At a thick place the stored interval is
$b\le w_{j+1}(\tau)\le b+\eta+1$.
At a non-thick top it is
$b_t+b'/x\le w_{j+1}(\tau)\le b_t+b'/x+\eta$;
here $b_t\ge b'\ge1$ and (GFA.25) applies.
At a place that is neither thick nor top it is
$w''\le w_{j+1}(\tau)\le w''+\eta$, so again
$1\le w_{j+1}(\tau)\le2b'$.
These cases prove all three clauses of \texttt{Nodes\_\allowbreak at}:
all nodes in $\Phi$ are at least $1$, thick nodes have their
stated $b$-interval, and the remaining nodes are at most $2b'$.
No clause for an unrepresented code is needed on $\Phi$ because
the displayed inverse gives a representative there.  The
definition's extra convention $w=1$ elsewhere is retained.
The threshold $\max(3,\mathcal F.x_1+1)$ proves
\texttt{Nodes\allowbreak F}.
% END GFA10

% BEGIN GFA11
\paragraph{GFA11: all seventeen original size-window clauses.}
\label{claim:GFA11}
Supply the same family's \texttt{Thick\allowbreak Set\allowbreak Density\allowbreak Eventually}.
For every fixed $c_6>0$ we prove \texttt{Win\allowbreak F} with $e_K=5$.
The raw \texttt{Size\allowbreak M}, \texttt{Size\allowbreak K}, \texttt{Size\allowbreak N}
windows have exponents
\[
 m\in[x^{1/2-\epsilon/16-\epsilon/100},
       x^{1/2-\epsilon/16+\epsilon/100}],\quad
 k\in[x^{9/2+10\epsilon-\epsilon/100},
       x^{9/2+10\epsilon+\epsilon/100}],\quad n\le x^3.
\]
Since $0<\epsilon<1/52$, comparison of these exponents gives
\[
 x^{1/2-\epsilon}\le m\le x^{1/2+\epsilon/2},\qquad
 xm\le k\le x^5,
 \qquad k\ge x^4.                                         \tag{GFA.26}
\]
The separate parameter field $k_{\rm up}$ reads exactly
$k\le n^2x^{1/2-\epsilon}$.  This is retained as a different
inequality; it is not inferred from the loose windows in
(GFA.26).

Eventually $k\le x^5\le X'\le a$, whence $kb'\le a$.
GFA10 gives $w''\ge1$, so $X''=a/w''\le a$.
The top interval
$a\le a_t\le a+a/(kx^{13})$ proves both
$a_t\le2a$ and
$0\le a_t/b-x\le x/(kx^{13})\le1$.
The literal code-size upper bound
$\#\mathrm{Word}\le2^{k-3}$ and its nonempty word set give
$\log\#\mathrm{Word}\le(k-3)\log2$.

We also check $\mathcal L_a=L.\mathrm{La}=\log(a\,\mathrm{rd}_E)$,
which differs from $\Lambda=L.\mathrm{LaG}$.
The size record has $\mathrm{rd}_E\ge2$ and
\[
 \mathrm{rd}_E\le\frac{Q_S}{2}\sqrt{qd},\quad
 Q_S\le4n x^{\epsilon/100},\quad
 \log q,\log d\le x^{\epsilon/5+\epsilon/100}.
\]
Hence
\[
 0\le\mathcal L_a
 \le x^\epsilon+5\log x+x^{\epsilon/5+\epsilon/100}
 \le2x^\epsilon                                              \tag{GFA.27}
\]
eventually.  In the middle inequality use
$\log n+(\epsilon/100)\log x\le3\log x$ from the actual
\texttt{Size\allowbreak N} upper exponent and $\log2\le\log x$.
In the last use $5\log x\le x^\epsilon/2$ and
$2x^{\epsilon/5+\epsilon/100}\le x^\epsilon$.
The same record gives $v=\mathrm{rd}_E/2>0$.

There are two small-error clauses.  Since $f=2n\le2x^3$,
choose a threshold with $x^7\le e^{c_6x^{\epsilon/8}}$.
For $x\ge3$, $4(x+1)\le x^3$, and therefore
\[
 2f(x+1)e^{-c_6x^{\epsilon/8}}\le1/x.                       \tag{GFA.28}
\]
The source's different parameter inequality
$\eta\le b/(t^3fk^3)$, with $t\ge1,\allowbreak{} k\ge x^4$, gives
$f\eta x^{12}\le b$.  Since
$12x(x+1)\le x^{12}$ for $x\ge3$,
$12x(x+1)f\eta\le b$.
Also $12x(x+1)f\le x^9\le b$ eventually.
Adding and dividing by $2x$ proves
\[
 6f(x+1)(\eta+1)\le b/x.                                  \tag{GFA.29}
\]

The seventeen fields of $\texttt{Win\_\allowbreak at}(L,5,\allowbreak{} c_6)$, in their
source order, are now accounted for as follows:
\[
\begin{array}{c|l}
1,2,3 & m\text{ upper},\ xm\le k,\ k\le x^5\quad\text{by (GFA.26)}\\
4,5,6 & k\le a,\ 0\le\mathcal L_a,\ \mathcal L_a\le2x^\epsilon\\
7,8 & kb'\le a,\ X''\le a\\
9,10 & a_t\le2a,\ \log\#\mathrm{Word}\le(k-3)\log2\\
11,12 & \text{(GFA.28), (GFA.29)}\\
13,14 & m\text{ lower in (GFA.26)},\quad k\le n^2x^{1/2-\epsilon}\\
15 & x^{-\epsilon}\le |T_2|/n\quad\text{by the supplied thick density}\\
16,17 & |a_t/b-x|\le1,\quad v>0.
\end{array}
\]
All exponential comparisons involve fixed positive exponents
and fixed constants.  Their finite maximum, together with the
thick-density and original family thresholds, precedes $i$.
This proves the literal \texttt{Win\allowbreak F}, with the prescribed
order $\forall c_6>0\ \exists e_K,\allowbreak{} x_0\ \forall x\ge x_0\ \forall i$.
% END GFA11

% BEGIN GFA12
\paragraph{GFA12: the original primitive-character reading and \texttt{Char\allowbreak F}.}
\label{claim:GFA12}
There is an elementary uniqueness observation useful for the
exact source choices.  Suppose primitive Dirichlet characters
$\chi_1,\chi_2$ agree at every prime outside the divisors of
a fixed nonzero integer $M$.  For any unit residue class modulo
$N_1N_2$, CRT gives a positive representative $a$ avoiding
every prime dividing $M$: retain the given unit residue at
the primes dividing $N_1N_2$, and impose $a\equiv1$ at the
other prime divisors of $M$.  Every prime factor of $a$ is
outside $M$.  Complete multiplicativity gives
$\chi_1(a)=\chi_2(a)$.  Thus their inductions to $N_1N_2$
are equal.  Factoring that character through the gcd of its
two primitive moduli proves $N_1=N_2$ and equality of the
primitive characters, hence equality at \emph{all} positive
integers.  More explicitly, CRT decomposes the common character
into characters of the groups of units modulo prime powers.
If such a local character factors through both exponents,
it factors through their minimum: the kernel at that minimum
is the larger of the two nested reduction kernels.  Taking
the product proves the asserted gcd factorization and the
uniqueness of the minimal conductor.  No prime-in-progression existence theorem is needed
for this uniqueness argument.

On the actual host, AH9--AH10 identify the specified Galois
map and Artin map, including their inertia and Frobenius
cosets.  Let
\[
 t_E:D.R.GG\longrightarrow(\mathbb Z/2\mathbb Z)^2\times D.R.A
\]
be the source's exact multiplicative-to-additive isomorphism.
It sends the second multiplicative sign to the first additive
sign and the first multiplicative sign to the second, as in
\texttt{gcode}.  For an additive character $\phi$, set
$c(\phi)(g)=\phi(t_Eg)$; the inverse is precomposition by
$t_E^{-1}$.  These maps are inverse character isomorphisms,
and $c(\phi)=1$ exactly when $\phi$ is the trivial additive
character (written $0$ in the source).

The original reciprocity map is surjective.  WB1 applied to
$c(\phi)\circ\mathrm{art}$ supplies its primitive inducer
$\chi_\phi$, with the precise values
\[
 \chi_\phi(n)=\mathrm{primVal}(\mathrm{arithOf}(D.host),\phi,\allowbreak{} n)
 \qquad(n\ge1).                                           \tag{GFA.30}
\]
Its values at unramified primes are the required Frobenius
values by AH9.  Therefore it is an actual witness of
\texttt{Reads}, proving that the source's primitive-reading
choice has a witness.  Every possible such reading agrees
with this one off the finite discriminant set.  The preceding
uniqueness proof identifies the source's stored
$\mathrm{prim}(D,\allowbreak{} c(\phi))$ with this same primitive character,
including every ramified positive integer in (GFA.30).
Surjectivity of reciprocity also gives
$\chi_\phi=1$ if and only if $\phi=0$.

Define the witness required by \texttt{Char\_\allowbreak at} to be
\[
 \ell(\phi)=\mathrm{Lp}(D,\allowbreak{} c(\phi)).
\]
If any primitive $\chi$ is prescribed by all the positive
integer values in (GFA.30), the same uniqueness argument
identifies it with the stored character, so
$\ell(\phi)=L'/L(1,\chi)$ for $\phi\ne0$.
The centered ramps also agree exactly:
\[
 \mathrm{GPsi}(\mathrm{arithOf}(D.host),\ell,\phi,\allowbreak{} y)
       =\mathrm{Gc}(D,\allowbreak{} c(\phi),\allowbreak{} y).                           \tag{GFA.31}
\]
For a nontrivial character substitute (GFA.30) in the finite
von Mangoldt sum and the same $y\ell(\phi)$ term.  Its
positive-part factor equals $y/n-1$ on the actual summation
range $1\le n\le\lfloor y\rfloor$.  If that range is empty,
both sums are zero.  For the trivial character both sides
are the same centered principal ramp, independently of the
unused value $\ell(0)$.  This checks the centering, including
$0<y<1$.

The distinguished additive characters $\chi_C,\psi_C$ map
under $c$ to the second and first multiplicative signs,
respectively.  Their product maps to the product of these
signs.  The original host includes the field
\texttt{Host.\allowbreak fl\allowbreak PS}, which supplies its actual \texttt{Ps\allowbreak Q}
quadratic-root identities; AH9--AH10 identify their Artin
and Frobenius readings on every such host.  These identities
identify the corresponding primitive characters with the
inducers of $\operatorname{chDir}(d)$,
$\operatorname{psDir}(q)$ and their product, in the source's
modulus-$8qd$ reading.  For $d\equiv3\pmod4$, quadratic
reciprocity identifies $\operatorname{chDir}(d)$ with the
Legendre character modulo $d$, which is the actual
$\mathrm{chiK}(D)$.  Thus the definition of $\mathrm{dlogL}$
and the stored logarithmic-derivative field give
\[
 L.\mathrm{logDerivativeAtOne}=\operatorname{Re}\ell(\chi_C).
                                                               \tag{GFA.32}
\]
Only this real-part equality is required here.

Every $\chi_\phi$ has the source's \texttt{Is\allowbreak Char\allowbreak ED}
witness $c(\phi)$: at integers prime to the host modulus,
(GFA.30) equals its value on the actual Artin image.
Consequently \texttt{Ch\allowbreak Thin\_\allowbreak at} applied to these very
characters and (GFA.31) gives \texttt{Thin\allowbreak Scales}, with the
same $X'=e^{x^{\epsilon/3}}$, the same $c_6$, and the same
lower limit $y\ge X'/3$.
Likewise the three clauses of \texttt{Ch\allowbreak QD\_\allowbreak at}, using the
three just-identified quadratic readings, give
\texttt{Values\allowbreak At\allowbreak X} with its original $C_5$.
Equations (GFA.30)--(GFA.32) and these two implications are
all four fields of \texttt{Char\_\allowbreak at}.  They are pointwise
identities for every original record; thus \texttt{Char\allowbreak F}
has threshold zero for every supplied $c_6,\allowbreak{} C_5$.
The thin and quadratic estimates remain their stated
\texttt{Chosen\allowbreak Thin}/\texttt{Chosen\allowbreak QD} inputs.
% END GFA12

\paragraph{The remaining final-global work after these installations.}
On the same original family, GVC supplies \texttt{Count\allowbreak F} and
\texttt{Slopes\allowbreak MF}; GFA1--GFA5 supply \texttt{Places\allowbreak F} and
\texttt{Slopes\allowbreak NF}; GFA6--GFA9 supply \texttt{RRRank\allowbreak F},
\texttt{Lat\allowbreak F} and the literal \texttt{Frame\allowbreak F}; GFA10--GFA12
supply \texttt{Nodes\allowbreak F}, \texttt{Win\allowbreak F} and \texttt{Char\allowbreak F}.
These are proved implications from their stated original
inputs, not new producers for those inputs.

The final global obligation still has two substantive parts:
the original $\texttt{Loss\allowbreak FS}(\mathcal F,0)$ estimate from
the chosen small shares, and the exact common-family
\texttt{link\allowbreak Core}/\texttt{Chain\allowbreak Tied} installation with those
losses.  The latter must carry the same selected top and
bottom values, original host, characters, all eleven ledger
fields, and the fixed constants through the final
\texttt{Closed} wrapper.  ZW1--ZW4 already prove the numerical
and analytic consequence once that actual tied ledger is
available.  Therefore \texttt{TASK\_\allowbreak GL\_\allowbreak FINAL\_\allowbreak ASSEMBLY}
remains open B with a narrower, explicit contract; it is not
declared closed by listing the now-completed inputs.




\section{The original theta integral: real derivatives, uniform growth, and lowering}
\label{sec:theta-extension-v034}

\status{TAE1--TAE13 independently reconstruct the real analytic and
lowering block of the original base theta extension.  They apply to the
same full Weil datum, theta functions, frames, coefficient space, Haar
measure and fundamental domain already used in BTH, FIA and CTC.
The remaining common polynomial identification, lowest-type
multiplicity-one, boundedness and Hilbert realization are stated
separately at the end.  No new Lean replay or category-A application is
claimed.}

Fix the original number-field datum \(D\), standard character
\(\psi\) with \(\mathrm{IsStdChar}(D,\psi,-1)\), Hermitian spaces
\(S,\allowbreak{} T\), infinity chart \(\theta\), and full action
\[
 \omega:G_S(\mathbb A_F)\times G_T(\mathbb A_F)
       \longrightarrow\operatorname{End}_{\mathbb C}
                    \bigl(\mathbb C^{W}\bigr),
 \qquad W=\mathrm{WO}(S,\allowbreak{} T).
 \tag{TAE.1}
\]
The premise on this action is the literal
\(\mathrm{IsWeilDatumO}(\psi,\allowbreak{} S,\allowbreak{} T,\chi_S,\chi_T,\omega)\).
Write \(\Theta\) for its original space
\(\mathrm{IsThetaFunO}(\psi,\allowbreak{} S,\allowbreak{} T)\).  The real Hom space is
\[
 V_{\mathbb R}=\mathrm{RW}(T,\allowbreak{} S)
       =\prod_{\tau}\operatorname{Mat}_{r_T,\allowbreak{} r_S}(\mathbb C),
 \qquad N=\dim_{\mathbb R}V_{\mathbb R}.
 \tag{TAE.2}
\]
The inverse infinity chart gives mutually compatible real inclusion
\(\iota:V_{\mathbb R}\to W\) and projection \(\pi:W\to V_{\mathbb R}\).
These are the actual \(\mathrm{iotaW}\) and \(\mathrm{piW}\).
The real Heisenberg action is
\(\rho(v)=\mathrm{heisO}(\psi,\allowbreak{} S,\allowbreak{} T,\iota(v))\).

CW12 identifies every \(f\in\Theta\) with a smooth vector in the Banach
space \(\mathcal B\) of bounded continuous functions satisfying the
original rational quasi-periodicity, equipped with the supremum norm.
We use its smooth topology
\[
 p_{\mathbf v}(f)=
  \|D_{v_1}\cdots D_{v_k}f\|_{\infty},
 \qquad
 D_vf=\left.\frac{d}{dt}\rho(tv)f\right|_{t=0}.
 \tag{TAE.3}
\]
The empty word is included.  Every \(D_v\) preserves \(\Theta\), and
all real derivatives preserve any fixed finite Heisenberg level.
These facts concern the genuine smooth domain; they do not assert
continuity of the ambient endomorphism on all functions.

The theta kernel and the source integral will always mean
\[
 K_\omega(h,\allowbreak{} g;f)=(\omega(h,\allowbreak{} g)f)(0),\qquad
 \Phi_L(g)=\int_{\mathcal F_S}
        K_\omega(h,\allowbreak{} g;L(F(h)))\,\allowbreak{} d\nu(h).
 \tag{TAE.4}
\]
Here \(E\) is a finite-dimensional complex normed space,
\(L:E\to\Theta\) is complex linear, and
\(F:G_S(\mathbb A_F)\to E\) is continuous and left
\(G_S(F)\)-invariant.  Whenever the integral is used, its original
source quotient has a compact set of representatives \(C_S\).
For the supplied base bundle this is exactly the finite class
presentation established in BTH3, together with compactness of the
definite real group and of the finite level.  Thus no new general
class-number assertion is being imported.

\subsection*{The real action on its genuine smooth domain}

% BEGIN TAE1
\paragraph{TAE1: the explicit Gaussian oscillator action is smooth in all real parameters.}
\label{claim:TAE1}
Use CW8--CW11's real coordinates: \(V\) is Euclidean,
\(J^2=-1\) is orthogonal, and
\[
 \rho(u)\rho(v)=e^{\pi i\beta(u,\allowbreak{} v)}\rho(u+v),
 \qquad \beta(u,\allowbreak{} v)=\langle u,\allowbreak{} Jv\rangle .
 \tag{TAE.5}
\]
Let \(g\) vary in the unitary matrix group preserving the same
Hermitian form.  CW10's operator
\(\Omega(g)\) acts smoothly on each smooth vector, in every seminorm
(TAE.3).  Its derivatives preserve the smooth-vector space.  On a
compact parameter set each derivative is bounded by a finite sum of
seminorms of the input.

\emph{Proof.}
Put \(G(u)=e^{-\pi\|u\|^2/2}\),
\[
 P=\int_VG(u)\rho(u)\,\allowbreak{} du,\qquad
 a_f(v)=P\rho(-v)f .
 \tag{TAE.6}
\]
The Gaussian and the Euclidean Lebesgue measure are precisely those
of CW8.  That proof gives
\[
 a_f(v)=\rho(-v)\widehat{(G(\cdot)\rho(\cdot)f)}(Jv),
 \qquad
 \int_V(1+\|v\|)^q\|a_f(v)\|\,\allowbreak{} dv<\infty
 \quad(q\ge0).
 \tag{TAE.7}
\]
The integral in (TAE.7) is controlled by finitely many smooth
seminorms of \(f\).  Indeed differentiate
\(G(u)\rho(u)f\) in a fixed real basis.  Each derivative is a finite
sum of a polynomial times \(G(u)\rho(u)D_{\mathbf e}f\), by the
Heisenberg law.  Its \(L^1\) norm is therefore bounded by a finite
sum of the \(p_{\mathbf e}(f)\).  Fourier integration by parts,
using more than \(q+\dim V\) coordinate derivatives, bounds its
Fourier transform by a constant times
\((1+\|v\|)^{-q-\dim V-1}\).  This proves the claimed quantitative
version of (TAE.7).

We also need a uniform estimate on derivatives of \(P\), not just
pointwise differentiability.  Changing variables in
\[
 \rho(tz)Pf
  =\int_VG(u-tz)e^{\pi i\beta(tz,\allowbreak{} u)}\rho(u)f\,\allowbreak{} du
 \tag{TAE.8}
\]
shows that \(D_zP\) is a bounded operator on \(\mathcal B\):
its kernel is a degree-one polynomial times a Gaussian.
Repeating the argument proves that
\(D_{z_1}\cdots D_{z_k}P\) is bounded, with norm at most
\(C_k\prod_j\|z_j\|\).  All differentiations of (TAE.8) are
dominated by an integrable polynomial times a Gaussian, uniformly
for \(t\) in a compact interval.  Since \(P^2=P\), these estimates
apply to \(a_f(v)\), with its own norm on the right.

Expanding the two projectors in CW10 and using covariance gives
the following particularly useful version of its defining integral:
\[
 \Omega(g)f=c(g)\int_V\int_V
   G(u)e^{\pi i\beta(v,\allowbreak{} u)}
           \rho(g(v+u))a_f(v)\,\allowbreak{} du\,\allowbreak{} dv .
 \tag{TAE.9}
\]
Here \(c(g)=\det(p g+q)\) in the signature coordinates of CW9;
after a fixed Hermitian congruence it is the same determinant
polynomial of the conjugated matrix.  The phase in (TAE.9) is
independent of \(g\), since \(g\) preserves \(\beta\).

For a smooth curve \(y(t)\) and a smooth vector \(a\), the
Heisenberg law yields
\[
 \frac{d}{dt}\rho(y(t))a
 =\rho(y(t))
       \bigl(D_{y'(t)}a-\pi i\beta(y(t),\allowbreak{} y'(t))a\bigr).
 \tag{TAE.10}
\]
Apply this with \(y(t)=g(t)(v+u)\) and \(a=a_f(v)\).
All derivatives of \(g\) are bounded on a fixed compact parameter
set.  Iterating (TAE.10), and then applying any fixed word of
left Heisenberg derivatives, bounds the resulting integrand by
\[
 C(1+\|u\|+\|v\|)^q\,\allowbreak{} G(u)\,\|a_f(v)\|.
 \tag{TAE.11}
\]
The derivatives of \(a_f(v)\) which occur here are bounded using
(TAE.8), because \(a_f(v)=Pa_f(v)\).
The function in (TAE.11) is integrable by (TAE.7) and the Gaussian
moments in \(u\).  Dominated differentiation in the Banach norm
therefore applies at every order.  It also proves continuity of
each derivative in every seminorm (TAE.3).

For clarity, this argument is valid even though a Heisenberg
derivative is an unbounded operator on \(\mathcal B\).
One can apply a finite Heisenberg difference quotient to (TAE.9)
first, dominate it by (TAE.11), and then pass to the limit.  The
limit is the derivative asserted above.  Equivalently the
generator of an isometric one-parameter group is closed, so the
same dominated estimates identify its value on the integral.
Thus no exchange of an unbounded operator and an integral is
assumed.

In the original adelic model, a fixed finite Heisenberg
translation commutes with every real \(\rho(u)\).  It therefore
commutes with \(P\), with (TAE.9), and with all the derivatives
just obtained.  A finite open level fixing \(f\) fixes all those
derivatives.  CW12's converse now makes them original theta
functions.  This proves the assertion on exactly the space that
will be used below.
% END TAE1

% BEGIN TAE2
\paragraph{TAE2: the continuous scalar twists have genuine real derivatives.}
\label{claim:TAE2}
A continuous character \(s:\mathbb R\to\mathbb C^\times\) with
\(s(t+u)=s(t)s(u)\) is smooth and has the form
\(s(t)=e^{a t}\).  A continuous character of a finite-dimensional
matrix Lie group is smooth on that group.  In particular the
determinant-root characters in CW24 are smooth on the real
factor of the original adelic group.

\emph{Proof.}
Let \(A(t)=\int_0^t s(u)\,\allowbreak{} du\).  Then \(A'=s\), and
\(A(\epsilon)\ne0\) for some \(\epsilon\), since
\(A'(0)=s(0)=1\).  Translation of this ordinary integral gives
\[
 s(t)=\frac{A(t+\epsilon)-A(t)}{A(\epsilon)}.
 \tag{TAE.12}
\]
The right side is continuously differentiable.  Differentiating
the character law at its second variable zero gives
\(s'(t)=s(t)s'(0)\).  Hence
\(e^{-s'(0)t}s(t)\) has zero derivative and value one at zero.
This proves the formula and all higher differentiability.
If the character is unitary, \(a\) is purely imaginary, but this
extra fact is not needed for differentiability.

Choose a real basis \(X_1,\ldots,\allowbreak{} X_d\) of a matrix Lie algebra.
The product
\(\exp(t_1X_1)\cdots\exp(t_dX_d)\) is a smooth local chart at the
identity: its differential there is the basis isomorphism, so the
inverse function theorem applies.  A character on this chart
is the product of the smooth one-variable characters just proved.
Translation gives smoothness near every point.  CW24 proves that
its determinant-root twists are continuous characters on the
full original adelic groups.  Restriction to the real factors
therefore satisfies precisely these hypotheses.
% END TAE2

% BEGIN TAE3
\paragraph{TAE3: the original action has the required generator, commutators, and translated point derivatives.}
\label{claim:TAE3}
For every real place \(\tau\) and
\[
 X^*H_{T,\tau}+H_{T,\tau}X=0
 \tag{TAE.13}
\]
there is a complex-linear map \(\mathcal G_{\tau,\allowbreak{} X}:\Theta\to\Theta\)
such that, for every \(f\in\Theta\), every adelic pair \(p\), and
every \(w\in W\),
\[
 \left.\frac{d}{dt}
  \bigl(\omega(p(1,\mathrm{oneP}_\tau(X,\allowbreak{} t)))f\bigr)(w)
       \right|_{t=0}
       =(\omega(p)\mathcal G_{\tau,\allowbreak{} X}f)(w).
 \tag{TAE.14}
\]
Writing \(D_{\tau',\allowbreak{} u}\) for the original \(\mathrm{dRho}\), one has
\[
 [\mathcal G_{\tau,\allowbreak{} X},\allowbreak{} D_{\tau,\allowbreak{} u}]f=D_{\tau,\allowbreak{} Xu}f,
 \qquad
 [\mathcal G_{\tau,\allowbreak{} X},\allowbreak{} D_{\tau',\allowbreak{} u}]f=0\quad(\tau'\ne\tau).
 \tag{TAE.15}
\]
Extending \(\mathcal G_{\tau,\allowbreak{} X}\) linearly off \(\Theta\) gives
the literal ambient endomorphism in the original
\texttt{exists\_\allowbreak generator} statement.

\emph{Proof.}
CW24 constructs a full datum with the same
\(D,\psi,\allowbreak{} S,\allowbreak{} T,\chi_S,\chi_T\).  BTH14 identifies it with the
specified \(\omega\) on every theta function.  On the real
factor this constructed action is
\[
                  \omega(\gamma)f=\zeta(\gamma)\Omega(m(\gamma))f,
 \tag{TAE.16}
\]
where \(\zeta\) is the continuous unitary determinant-root
character of CW24 and
\[
 m(g,\allowbreak{} h)=\operatorname{diag}_{\tau'}
       \bigl(h_{\tau'}\otimes(g_{\tau'}^{-1})^{\mathsf t}\bigr)
 \tag{TAE.17}
\]
is the original action on the real Hom space.  Neither a new
character nor a different action is chosen in (TAE.16).
The finite factor of a real-supported element is the identity.
TAE1 and TAE2 therefore make (TAE.16) smooth in the topology
(TAE.3).  Its derivative along \((1,\mathrm{oneP}_\tau(X,\allowbreak{} t))\)
defines \(\mathcal G_{\tau,\allowbreak{} X}\).  Linearity in \(f\) follows from
linearity of the action.  Preservation of the real smooth
domain and of a finite level follows from TAE1; scalar
differentiation does not change a finite level.  Thus its
values are in \(\Theta\).  It is also real linear in \(X\),
being the differential of the smooth real group action.

The group covariance law gives, on \(\Theta\),
\[
 \omega(1,\mathrm{oneP}_\tau(X,\allowbreak{} t))D_{\tau',\allowbreak{} u}
 =D_{\tau',\,\allowbreak{} e^{tX}u}
       \omega(1,\mathrm{oneP}_\tau(X,\allowbreak{} t))
       \quad(\tau'=\tau);
 \tag{TAE.18}
\]
if \(\tau'\ne\tau\), its right-hand direction is just \(u\).
To justify (TAE.18), first differentiate the Heisenberg
covariance law in the real shift parameter.  Both sides are
smooth in the topology (TAE.3) by TAE1, so the differentiation
is legitimate.  Differentiating once more in \(t\) at zero
gives exactly (TAE.15), including its plus sign.

There is a separate point concerning the arbitrary adelic
left factor \(p\) in (TAE.14).  We do not assume that
\(\omega(p)\) is bounded in the supremum norm.
Set
\[
 u(t)=\omega(\gamma(t))f,\quad
 A=\mathcal G_{\tau,\allowbreak{} X}f,\quad
 v(t)=\omega(p\gamma(t))f,
 \qquad \gamma(t)=(1,\mathrm{oneP}_\tau(X,\allowbreak{} t)).
 \tag{TAE.19}
\]
The first curve is differentiable in \(\mathcal B\).
The real subgroup is normal in the adelic group, and
\(p\gamma(t)p^{-1}\) is a smooth real matrix curve.
Applying (TAE.16) to this curve and to \(\omega(p)f\in\Theta\)
shows that \(v(t)\) is also differentiable in \(\mathcal B\);
write its derivative \(B\).

Choose an additive Haar measure on \(W\), normalized so that
the compact rational quotient has volume one, and a rational
fundamental domain \(\mathcal F_W\).  The map
\(\mathcal B\to L^2(\mathcal F_W)\) is bounded.  The original
\(\mathrm{om\_\allowbreak unit}\) gives
\[
 \|v(t)-v(0)-t\,\omega(p)A\|_{L^2}
       =\|u(t)-u(0)-tA\|_{L^2}=o(|t|).
 \tag{TAE.20}
\]
All vectors inside this equality are theta functions, so that
the stated isometry applies.  Uniqueness of an \(L^2\)
derivative gives \(B=\omega(p)A\) in \(L^2\).
Their difference is continuous and its norm is rationally
periodic.  A nonzero value would give a positive-measure open
set of positive norm in the quotient; hence equality in
\(L^2\) implies equality at every point of \(W\).
Point evaluation on \(\mathcal B\) is bounded, so (TAE.14)
follows.

Finally choose a complex vector-space complement to \(\Theta\)
in the space of all functions and let the generator be zero
on that complement.  This gives an ambient complex-linear
endomorphism.  All assertions involve theta functions and
their Heisenberg derivatives, so the extension changes none
of them.  This proves exactly the original generator
contract, without leaving a real oscillator differentiability
theorem as an additional input.
% END TAE3

\subsection*{One growth exponent for every derivative word}

% BEGIN TAE4
\paragraph{TAE4: the finite-level real-cube bound is uniform along the real group.}
\label{claim:TAE4}
Fix a basis \(e_1,\ldots,\allowbreak{} e_N\) of \(V_{\mathbb R}\), an open
finite Heisenberg level \(U\), a target point \(x\), and a compact
source set \(C\).  There is \(b_0\ge0\), independent of
\(h\in C\), \(g\in G_{T,\infty}\), and \(\varphi\in\Theta\)
fixed at level \(U\), such that
\[
 \int_{[0,1]^N}
   \left|\,\omega(h,\allowbreak{} xg)\varphi
          \left(\iota\left(\sum_j t_je_j\right)\right)\right|^2dt
             \le b_0\|\varphi\|_\infty^2.
 \tag{TAE.21}
\]
The same constant applies to every real Heisenberg derivative
of every such \(\varphi\).

\emph{Proof.}
The action of \((h,\allowbreak{} xg)^{-1}\) on a purely finite shift depends
only on \(h,\allowbreak{} x\); the real-supported \(g\) has trivial finite
component.  Continuity of this linear action and compactness
of \(C\) provide a common compact open neighborhood
\(U_0\) in the finite Hom space whose image lies in \(U\) for
every \(h\in C\).  Here one first chooses a neighborhood for
each \(h\), covers \(C\) by finitely many corresponding
neighborhoods, and intersects the finitely many finite
neighborhoods.  The full covariance law consequently fixes
\(\omega(h,\allowbreak{} xg)\varphi\) under the Heisenberg shifts in \(U_0\).
Their phases have modulus one.  Thus the squared norm at
\(\iota(v)+u\), \(u\in U_0\), equals its squared norm at
\(\iota(v)\).

The continuous splitting \(W=V_{\mathbb R}\oplus W_f\)
identifies Haar measure with a positive scalar multiple of
the product Haar measures.  The compact set
\[
 K_0=\left\{\iota\left(\sum_jt_je_j\right)+u:
                         t\in[0,1]^N,\ u\in U_0\right\}
 \tag{TAE.22}
\]
therefore converts the left side of (TAE.21) to a fixed
positive multiple of the integral of
\(|\omega(h,\allowbreak{} xg)\varphi|^2\) over \(K_0\).
This is the explicit real-slice measure comparison; restricting
an adelic integral to a measure-zero real slice would not give
it without the common finite level.

Let \(\Gamma=W(F)\), a discrete cocompact additive lattice.
Each coset meets \(K_0\) with multiplicity at most
\[
 M_0=\#\bigl(\Gamma\cap(K_0-K_0)\bigr)<\infty.
 \tag{TAE.23}
\]
Indeed, relative to one such intersection point, every other
rational translate lies in the displayed finite set.
Unfolding a nonnegative rationally periodic function over a
fundamental domain consequently bounds its \(K_0\)-integral
by \(M_0\) times its quotient integral.
The norm of a theta function is rationally periodic, and
\(\mathrm{om\_\allowbreak unit}\) identifies the quotient integral of
\(|\omega(h,\allowbreak{} xg)\varphi|^2\) with that of \(|\varphi|^2\).
Normalize the latter quotient measure to have volume one;
it is then at most \(\|\varphi\|_\infty^2\).
Combining these bounds gives (TAE.21).

Real and finite Heisenberg shifts commute.  Differentiating
that equality proves that every real derivative of
\(\varphi\) is fixed at the same \(U\).  Nothing in the
construction of \(U_0,\allowbreak{} K_0,\allowbreak{} M_0,\allowbreak{} b_0\) depends on the particular
vector, proving the final assertion.
% END TAE4

% BEGIN TAE5
\paragraph{TAE5: the original theta kernel has exponent \(N=\dim_{\mathbb R}\mathrm{RW}(T,\allowbreak{} S)\).}
\label{claim:TAE5}
For every \(f\in\Theta\), target point \(x\), and compact
source set \(C\), there is \(c_f\ge0\) such that
\[
 |K_\omega(h,\allowbreak{} xg;f)|
       \le c_f\,\allowbreak{} H_T(g)^N
       \quad(h\in C,\ g\in G_{T,\infty}),
 \qquad H_T(g)=\mathrm{normInf}(\theta,\allowbreak{} T,\allowbreak{} g).
 \tag{TAE.24}
\]
The exponent is independent of \(f,\allowbreak{} x,\allowbreak{} C\).

\emph{Proof.}
Use the supremum norm on the coordinates of \(V_{\mathbb R}\).
The original height is exactly
\[
 H_T(g)=\sum_{\tau,\allowbreak{} a,\allowbreak{} b}
       \left(|(g_\tau)_{ab}|+|(g_\tau^{-1})_{ab}|\right).
 \tag{TAE.25}
\]
For the original real Hom action write
\[
 M(h,\allowbreak{} k)u=(k_\tau u_\tau h_\tau^{-1})_\tau .
 \tag{TAE.26}
\]
The full Heisenberg covariance, differentiated as in TAE3,
gives
\[
 D_u\,\omega(h,\allowbreak{} k)f
        =\omega(h,\allowbreak{} k)\,\allowbreak{} D_{M(h^{-1},\allowbreak{} k^{-1})u}f .
 \tag{TAE.27}
\]
The same \(L^2\) difference-quotient argument as (TAE.20)
justifies this formula for an arbitrary adelic pair; all
Heisenberg derivatives exist in \(\mathcal B\) by CW12.

Matrix multiplication and the triangle inequality give the
precise height comparison
\[
 \|M(h,(xg)^{-1})u\|
       \le \|M(h,\allowbreak{} x^{-1})u\|\,\allowbreak{} H_T(g).
 \tag{TAE.28}
\]
To verify it, the left matrix is
\(g_\tau^{-1}(x_\tau^{-1}u_\tau h_\tau^{-1})\).
Each row sum of \(|g_\tau^{-1}|\) is bounded by (TAE.25),
and each entry of the other factor is bounded by its
supremum norm.  No changed height or hidden exponent is used.
Compactness of \(C\) gives a single \(\kappa\ge0\) with
\[
 \|M(h^{-1},\allowbreak{} x^{-1})e_j\|\le\kappa
                       \quad(h\in C,\ 1\le j\le N).
 \tag{TAE.29}
\]

We record an elementary word estimate.  For each fixed
\(f\in\Theta\) there is \(A_f\ge0\) such that
\[
 \|D_{u_k}\cdots D_{u_1}f\|_\infty
        \le A_f\rho^N
 \quad\left(k\le N,\ \rho\ge1,\ \|u_j\|\le\rho\right).
 \tag{TAE.30}
\]
Expand every \(u_j\) in the fixed real basis.  The coordinate
functionals have bounded norm, and \(D_u\) is real linear in
\(u\).  The expansion is a finite sum of the finitely many
basis words of lengths at most \(N\), with coefficient bounds
of the form \(C^k\rho^k\).  Each basis word is a bounded theta
function by CW12.  Summing their bounds and using
\(\rho^k\le\rho^N\) proves (TAE.30), including the empty word.

For a real direction \(e_j\), the one-dimensional fundamental
theorem of calculus and Cauchy--Schwarz imply
\[
 |q(0)|^2\le
       2\int_0^1|q(t)|^2\,\allowbreak{} dt+
       2\int_0^1|q'(t)|^2\,\allowbreak{} dt .
 \tag{TAE.31}
\]
Indeed \(q(0)=q(t)-\int_0^tq'(r)\,\allowbreak{} dr\); square, integrate
over \(t\), and enlarge the derivative integral to \([0,1]\).
Apply this to a Heisenberg orbit evaluated at a point.
The unit-modulus phase shows that its two terms are the
ordinary squared norms of the vector and its covariant
derivative on the translated real segment.  Iterating along
the \(N\) coordinate directions gives \(4^N\) times a common
upper bound for all real-cube integrals of derivative words
of length at most \(N\).  This is also an immediate induction:
at each step there are two terms, each with coefficient two.

Apply this inequality to \(\omega(h,\allowbreak{} xg)f\) at zero.
By (TAE.27)--(TAE.29), every derivative that appears is
\(\omega(h,\allowbreak{} xg)\varphi\), where \(\varphi\) is a word on \(f\)
of length at most \(N\) whose directions have norms at most
\(\kappa H_T(g)\).  All such words have the same finite
level.  TAE4 and (TAE.30), with
\(\rho=\max(1,\kappa H_T(g))\), therefore give
\[
 |K_\omega(h,\allowbreak{} xg;f)|^2
        \le 4^N b_0 A_f^2\rho^{2N}.
 \tag{TAE.32}
\]
If \(N>0\), then \(H_T(g)\ge1\).  To see this, select one
nonempty real matrix block and an index \(a\); the identity
\(g_\tau g_\tau^{-1}=1\) gives
\(1\le\sum_b |(g_\tau)_{ab}||(g_\tau^{-1})_{ba}|
 \le H_T(g)^2\).
Hence
\(\rho^N\le\max(1,\kappa)^N H_T(g)^N\).
Taking square roots in (TAE.32) proves (TAE.24), with
\[
       c_f=2^N\sqrt{b_0}\,\allowbreak{} A_f\max(1,\kappa)^N .
 \tag{TAE.33}
\]
If \(N=0\), the cube has one point, all words are empty, and
\(H_T(g)^0=1\); (TAE.21) proves the same assertion directly.
Thus the degenerate case needs no positive-height assumption.
% END TAE5

\subsection*{Differentiating the original definite-source integral}

% BEGIN TAE6
\paragraph{TAE6: compact representatives give dominated differentiation of the original integral.}
\label{claim:TAE6}
For the data of (TAE.4), \(\Phi_L\) is continuous and left
\(G_T(F)\)-invariant.  For every \(\tau,\allowbreak{} X\) satisfying
(TAE.13), put \(L'=\mathcal G_{\tau,\allowbreak{} X}\circ L\).  Then
\[
 R_{\tau,\allowbreak{} X}\Phi_L=\Phi_{L'},\qquad
 \left.\frac{d}{dt}
     K_\omega(h,\allowbreak{} g\,\mathrm{oneP}_\tau(X,\allowbreak{} t);Lz)
       \right|_{t=0}
     =K_\omega(h,\allowbreak{} g;L'z).
 \tag{TAE.34}
\]
Both integrals are genuinely absolutely convergent on the
original \((\nu,\mathcal F_S)\).

\emph{Proof.}
Choose a basis \(b_1,\ldots,\allowbreak{} b_d\) of \(E\), and write
\(F(h)=\sum_jF_j(h)b_j\).  For either \(L\) or \(L'\), the
kernel is the finite sum
\[
 K_\omega(h,\allowbreak{} g;L(F(h)))
          =\sum_j F_j(h)\,(\omega(h,\allowbreak{} g)Lb_j)(0).
 \tag{TAE.35}
\]
Each summand is jointly continuous by the actual
\(\mathrm{om\_\allowbreak cont}\) clause and continuity of \(F_j\).
Rational source invariance follows from
\(\mathrm{om\_\allowbreak aut}\) and the source invariance of \(F\);
rational target invariance follows from the other
\(\mathrm{om\_\allowbreak aut}\) clause.

The source quotient is covered by the compact set \(C_S\).
It has finite Haar volume, and every continuous source
invariant kernel is bounded by its maximum on \(C_S\).
This proves absolute convergence.  For a compact neighborhood
of a target point, (TAE.35) has a uniform bound on its product
with \(C_S\); ordinary dominated convergence proves
continuity of \(\Phi_L\).

The point derivative in (TAE.34) is TAE3 with \(p=(h,\allowbreak{} g)\)
and \(w=0\).  The one-parameter group law translates it
from zero to every \(t\).  Its derivative is therefore
\[
 K_\omega\bigl(h,\allowbreak{} g\,\mathrm{oneP}_\tau(X,\allowbreak{} t);L'(F(h))\bigr).
 \tag{TAE.36}
\]
This is jointly continuous in \((t,\allowbreak{} h)\), by (TAE.35) for
\(L'\).  On a compact interval about zero and \(C_S\) it has
a finite bound \(M\).  Rational source invariance extends
this same bound to every \(h\).
The difference quotient is dominated by \(M\), by the
one-variable integral mean-value estimate, and the constant
function \(M\) is integrable on \(\mathcal F_S\).
Dominated convergence therefore moves the derivative through
the integral.  This proves the first equality in (TAE.34).
This argument keeps the original measure and fundamental
domain; compact representatives are used only to produce a
dominating function.
% END TAE6

% BEGIN TAE7
\paragraph{TAE7: the lift is smooth and has one exponent for all real derivative words.}
\label{claim:TAE7}
The function \(\Phi_L\) satisfies the original
\(\mathrm{IsSmoothFn}(T,\theta)\) and
\(\mathrm{IsUnifModGrowth}(T,\theta)\).  In the latter
predicate one can take the same \(N\) of (TAE.2):
\[
 \forall\mathbf X,\ \forall x,\ \exists C_{\mathbf X,\allowbreak{} x}\ge0,\
 \forall g\in G_{T,\infty},\
 \bigl|R_{\mathbf X}\Phi_L(xg)\bigr|
                 \le C_{\mathbf X,\allowbreak{} x}H_T(g)^N .
 \tag{TAE.37}
\]
The derivative word consists of real Lie directions at any
of the original places, in its original order.
If a complex lowering derivative of every kernel
\(g\mapsto K_\omega(h,\allowbreak{} g;Lz)\) is zero, the corresponding
complex lowering derivative of \(\Phi_L\) is zero.

\emph{Proof.}
Repeatedly apply TAE6.  If
\(\mathbf X=((\tau_1,\allowbreak{} X_1),\ldots,(\tau_k,\allowbreak{} X_k))\) denotes
the original outer-to-inner order of \(\mathrm{foldr}\),
then
\[
 R_{\mathbf X}\Phi_L
   =\Phi_{\mathcal G_{\tau_1,\allowbreak{} X_1}\cdots
                     \mathcal G_{\tau_k,\allowbreak{} X_k}L}.
 \tag{TAE.38}
\]
Every map on the right is linear into \(\Theta\), and every
function so obtained is continuous by TAE6.

We explain why this gives the precise smoothness assertion.
On the product of the real unitary groups choose a Lie algebra
basis.  Products of its exponentials are local coordinates by
the inverse function theorem, as in TAE2.  A coordinate
derivative of a function is a finite linear combination of
its right Lie derivatives, with smooth coefficients: move
the differentiated exponential through the factors to its
right and conjugate its Lie direction by their inverse
product.  In a fixed Lie basis these adjoint coordinates
are smooth matrix functions.  TAE6 supplies the resulting
first derivatives, continuously; induction using (TAE.38)
supplies all coordinate derivatives.  Translation covers
the real group.  Composition with any finite product of
the specified one-parameter groups proves exactly the
\(\mathrm{ContDiff}\) requirement in \(\mathrm{IsSmoothFn}\),
at every adelic left translate.

For any linear \(L_1:E\to\Theta\), apply TAE5 to each
of the finitely many \(L_1b_j\), with \(C=C_S\).
Let \(M_j=\sup_{h\in C_S}|F_j(h)|\), and let \(c_j\)
be the corresponding constants in (TAE.24).
Source invariance then gives, for every \(h\),
\[
 |K_\omega(h,\allowbreak{} xg;L_1(F(h)))|
       \le\left(\sum_j M_jc_j\right)H_T(g)^N.
 \tag{TAE.39}
\]
Integration multiplies the constant by
\(\nu(\mathcal F_S)<\infty\).  Take
\(L_1=\mathcal G_{\tau_1,\allowbreak{} X_1}\cdots
\mathcal G_{\tau_k,\allowbreak{} X_k}L\).  Its constants may change
with the word, while the exponent \(N\) in TAE5 does not.
This proves the quantifier order in (TAE.37).

Finally, for a complex matrix \(Z\), the source defines
\[
 X=\tfrac12(Z+\overline Z^{\,\allowbreak{} H}),\qquad
 Y=\tfrac1{2i}(Z-\overline Z^{\,\allowbreak{} H}),\qquad
 \overline Z^{\,\allowbreak{} H}=-H_T^{-1}Z^*H_T,\qquad
 R_Z^{\mathbb C}=R_X+iR_Y .
 \tag{TAE.40}
\]
Both \(X,\allowbreak{} Y\) satisfy (TAE.13).  Apply TAE6 separately
to them.  If the sum of their pointwise kernel
derivatives is zero, its integral is zero by ordinary
linearity of the two absolutely convergent derivative
integrals.  This proves the final assertion.
% END TAE7

\subsection*{The exact lowering operator on the original Fock polynomials}

The source's real phase convention matters in the next calculation.
At a fixed place put \(\xi_\tau=\tau_E(D.\mathrm{host}.\xi_0)\).
It is nonzero and purely imaginary, with the prescribed positive
imaginary part.  On real matrix directions define
\[
 \begin{split}
 \beta_\tau(u,\allowbreak{} v)
   &=2\operatorname{Re}\!
       \left(\xi_\tau\operatorname{tr}
                 (H_S^{-1}v^*H_Tu)\right),\\
 s_\tau&=2\sqrt{2\pi|\xi_\tau|}>0.
 \end{split}
 \tag{TAE.41}
\]
This \(\beta_\tau\) is literally \(\mathrm{omDir}\) at the
same place; directions at distinct places pair to zero.
The character sign \(-1\) and the half in the original
Heisenberg law give
\[
 [D_u,\allowbreak{} D_v]=-2\pi i\,\beta_\tau(u,\allowbreak{} v)\,1 .
 \tag{TAE.42}
\]
For example, this follows by differentiating the group
commutator
\(\rho(tu)\rho(rv)\rho(-tu)\rho(-rv)
       =e^{-2\pi itr\beta_\tau(u,\allowbreak{} v)}1\).
This specifies the sign relative to the abstract coordinates
of (TAE.5); those coordinates already absorbed the original
alternating form.

In the rest of this subsection \(S\) is positive definite at
infinity.  Use its original positive basis \(B_\tau\) and the
target model's original basis \(B'_\tau=M.\mathrm{fr}_\tau^{-1}\).
They satisfy
\[
 B_\tau^*H_SB_\tau=1,\quad
 (B'_\tau)^*H_TB'_\tau
       =\operatorname{diag}(\epsilon_a),\quad
 \epsilon_a=\begin{cases}1&a<p_\tau,\\-1&a\ge p_\tau.\end{cases}
 \tag{TAE.43}
\]
The Cartan involution is
\(s=M.\mathrm{cart}_\tau
  =B'_\tau\operatorname{diag}(\epsilon_a)(B'_\tau)^{-1}\).
Thus \(J_\tau u=i s u\) is the original
\(\mathrm{cplxO}(1,\allowbreak{} s)\).  Define
\[
 A_\tau(u)=D_u+iD_{J_\tau u},\quad
 a_\tau^+(u)=s_\tau^{-1}(D_u-iD_{J_\tau u}),\quad
 u_{ac}=B'_\tau E_{ac}B_\tau^{-1}.
 \tag{TAE.44}
\]
The original polynomial vector \(\mathrm{pt}(P)\) is
\(P(a^+)f\), with one creation operator for each original
variable \((\tau,\allowbreak{} a,\allowbreak{} c)\), applied to the original vacuum \(f\).
The equality with the source's \(\mathrm{pureTensorO}\)
will be checked in TAE10.

% BEGIN TAE8
\paragraph{TAE8: a quadratic commutator determines the full polynomial action.}
\label{claim:TAE8}
Let \(a_j\) be commuting linear operators on a common invariant
complex vector space, let \(G,\allowbreak{} b_j\) be linear operators there,
and suppose
\[
 [G,\allowbreak{} a_j]=b_j,\qquad
 [b_j,\allowbreak{} a_k]=(\mu_{jk}+\mu_{kj})1,\qquad
 b_jf=0,\qquad Gf=0 .
 \tag{TAE.45}
\]
Then for every polynomial \(P\) in the finitely many variables,
\[
 GP(a)f=
    \left(\sum_{j,\allowbreak{} k}\mu_{jk}\partial_j\partial_kP\right)(a)f .
 \tag{TAE.46}
\]
No symmetry of \(\mu\), injectivity of \(P\mapsto P(a)f\), or
completion of the polynomial space is required.

\emph{Proof.}
First commute \(b_j\) across a monomial, one factor at a time.
Its final term kills \(f\), and its other terms give
\[
 b_jP(a)f=\sum_k(\mu_{jk}+\mu_{kj})(\partial_kP)(a)f .
 \tag{TAE.47}
\]
Linearity extends this from monomials to all polynomials.
Now prove (TAE.46) by induction on monomial degree.
It holds for constants because \(Gf=0\).
If it holds for \(P\), then
\[
 G(a_jP(a)f)=a_jGP(a)f+b_jP(a)f.
 \tag{TAE.48}
\]
On the polynomial side, applying
\(\sum_{k,\allowbreak{} l}\mu_{kl}\partial_k\partial_l\) to \(x_jP\)
gives \(x_j\) times its value on \(P\), plus
\(\sum_k(\mu_{jk}+\mu_{kj})\partial_kP\).
This is exactly (TAE.48), by (TAE.47).
The commuting \(a_j\)'s make every monomial evaluation
independent of its chosen order.  Addition completes the proof.
% END TAE8

% BEGIN TAE9
\paragraph{TAE9: the original vacuum is killed by the complex lowering generator.}
\label{claim:TAE9}
Fix \(\tau\).  Let \(X\) satisfy (TAE.13) and
\(sX=-Xs\), and put \(Y=isX\).  Then \(Y\) also satisfies
(TAE.13).  If \(f\in\Theta\) is an original Fock vacuum at
every real place, then
\[
       \mathcal Gf=0,\qquad
       \mathcal G=\mathcal G_{\tau,\allowbreak{} X}
                           +i\mathcal G_{\tau,\allowbreak{} Y}.
 \tag{TAE.49}
\]

\emph{Proof.}
Since \(s^*H_T=H_Ts\), the conditions on \(X\) give
\[
 Y^*H_T+H_TY
   =-iX^*H_Ts+iH_TsX
   =iH_T(Xs+sX)=0.
 \tag{TAE.50}
\]
The commutators of TAE3 imply
\[
 [\mathcal G,\allowbreak{} D_{\tau,\allowbreak{} u}]=A_\tau(Xu);
 \tag{TAE.51}
\]
the commutator is zero at every other place.
Consequently \(\mathcal Gf\) is again annihilated by
every \(A_{\tau'}(u)\): commute \(\mathcal G\) through the
two terms of \(A_{\tau'}(u)f=0\), and use
\(A_\tau(Xu)f=A_\tau(XJ_{\tau'}u)f=0\) when the places agree.
TAE3 also gives \(\mathcal Gf\in\Theta\).

Let \(k=(1,\mathrm{atTau}_\tau(s))\).
CW11 states that the canonical oscillator action of this
Cartan element is one fixed nonzero scalar on the entire
vacuum space.  Multiplication by CW24's determinant-root
character changes that scalar by another nonzero scalar,
independent of the vacuum vector.  BTH14 transfers this
statement to the actual \(\omega\).  Thus there is a single
\(c\ne0\) such that \(\omega(k)v=cv\) for every original
vacuum \(v\).  This applies both to \(f\) and to
\(\mathcal Gf\).

On the other hand,
\(s e^{tX}=e^{-tX}s\), directly from the exponential
series and \(sX=-Xs\).  TAE3 permits differentiation of
this group identity after applying \(\omega\) to \(f\):
\[
       \omega(k)\mathcal G_{\tau,\allowbreak{} X}f
                           =-c\mathcal G_{\tau,\allowbreak{} X}f .
 \tag{TAE.52}
\]
The same reasoning applies to \(Y\), since \(sY=-Ys\).
Hence \(\omega(k)\mathcal Gf=-c\mathcal Gf\).
But \(\mathcal Gf\) is a vacuum, so its left side is
\(c\mathcal Gf\).  As \(2c\ne0\), (TAE.49) follows.
This argument needs the common vacuum scalar, not a choice
of scalar depending on the tested vacuum.  CW11 and the
full-action comparison supply exactly that uniform statement.

There is also an original-data passage needed when \(f\)
is the stored \(\mathrm{IsThPVecOf}\) vector, rather than
its initial lattice vacuum \(f_0\).  We give it explicitly.
That record writes \(f\) as an \(\mathrm{IsThetaAt}\)
value of a finite split place model \(M_0\), with
vacuum \(f_0\).  Its \(\mathrm{AllSome}\) clause says that
every place model with those same data has that same
value on the prescribed finite function \(\Phi\).
The two half directions in this model have zero
infinite component.

Fix a real place and a direction \(z\).
The functional \(q\mapsto q'(0)\) is complex linear
on the subspace of functions \(\mathbb R\to\mathbb C\)
differentiable at zero.  Extend it algebraically
to a linear functional \(\ell\) on all functions.
Apply \(\ell\), pointwise, to the two real Heisenberg
curves defining \(A_\tau(z)\).  This gives a linear
operator \(\widetilde A\) on all functions on \(W\)
which agrees with \(A_\tau(z)\) on \(\Theta\).
It commutes with every purely finite Heisenberg
operator: the real and finite operators commute
before differentiation, and their fixed scalar
factors pass through \(\ell\).
Thus
\[
                  M_1=M_0+\widetilde A\,\allowbreak{} M_0
 \tag{TAE.52a}
\]
satisfies both of the same finite place-model
intertwining laws.  Its vacuum value is unchanged,
because \(\widetilde A f_0=A_\tau(z)f_0=0\).
The universal part of \(\mathrm{IsThetaAt}\) gives
\(f=M_0\Phi=M_1\Phi=f+\widetilde A f\).
Since the stored \(f\) is a theta function, this
says \(A_\tau(z)f=0\).
Apply the argument at every place and every \(z\).
It proves the literal
\(\mathrm{isFockVacAt\_\allowbreak thPVec}\) passage used below,
without assuming smoothness of arbitrary place-model
outputs or differentiating an arbitrary linear map.
% END TAE9

% BEGIN TAE10
\paragraph{TAE10: the actual lowering derivative is the original mixed Laplacian.}
\label{claim:TAE10}
Assume (TAE.43) at every place and let \(f\in\Theta\) be the
original vacuum at every place.  For a polynomial \(P\), suppose
at one place \(\tau\) that
\[
       \sum_{c=0}^{r_S-1}
           \partial_{\tau,\allowbreak{} a,\allowbreak{} c}\partial_{\tau,\allowbreak{} b,\allowbreak{} c}P=0
           \quad(a\ge p_\tau,\ b<p_\tau).
 \tag{TAE.53}
\]
Then for every original \(Z\in\mathfrak p^-_\tau(M)\)
and every \(h\in G_S(\mathbb A_F)\),
\[
 R_{\tau,\allowbreak{} Z}^{\mathbb C}
       \bigl(g\mapsto K_\omega(h,\allowbreak{} g;\mathrm{pt}(P))\bigr)=0 .
 \tag{TAE.54}
\]

\emph{Proof.}
The definition of \(\mathfrak p^-_\tau(M)\) says that
\(M.\mathrm{fr}_\tau ZM.\mathrm{fr}_\tau^{-1}\) is supported
only in the block with negative row and positive column.
Equivalently
\[
                         sZ=-Z,\qquad Zs=Z .
 \tag{TAE.55}
\]
Use the real parts \(X,\allowbreak{} Y\) of \(Z\) in (TAE.40).
Conjugating (TAE.55) by the Hermitian form shows
\[
        sX=-Xs,\qquad Y=isX .
 \tag{TAE.56}
\]
For example \(sH_T^{-1}=H_T^{-1}s^*\), whence
\[
 \begin{aligned}
 s(H_T^{-1}Z^*H_T)&=H_T^{-1}Z^*H_T,\\
 (H_T^{-1}Z^*H_T)s&=-H_T^{-1}Z^*H_T.
 \end{aligned}
\]
Substitution in (TAE.40) gives (TAE.56).
TAE9 now supplies \(Gf=0\) for
\(G=\mathcal G_{\tau,\allowbreak{} X}+i\mathcal G_{\tau,\allowbreak{} Y}\).

We compute the other commutators rather than identifying
them only up to a scalar.  The form in (TAE.41) satisfies
\[
 \beta_\tau(Ju,\allowbreak{} Jv)=\beta_\tau(u,\allowbreak{} v),\qquad
 \beta_\tau(u,\allowbreak{} Jv)=-\beta_\tau(Ju,\allowbreak{} v).
 \tag{TAE.57}
\]
These follow from \(s^*H_Ts=H_T\) and \(J^2=-1\).
Expanding (TAE.44) with (TAE.42) and (TAE.57) gives
\[
 \begin{split}
 [a_\tau^+(u),\allowbreak{} a_\tau^+(v)]&=0,\\
 [A_\tau(u),\allowbreak{} a_\tau^+(v)]
   &=-\frac{4\pi i}{s_\tau}
      \bigl(\beta_\tau(u,\allowbreak{} v)+i\beta_\tau(Ju,\allowbreak{} v)\bigr)\,1 .
 \end{split}
 \tag{TAE.58}
\]
At different places all these operators commute.
In particular the source's \(\mathrm{opPow}\) ordering
gives precisely \(P(a^+)f\), because the creation
operators commute.  The definition by the monomial basis
therefore agrees with \(\mathrm{pureTensorO}\) on every
polynomial.

For \(j=(\rho,\allowbreak{} a,\allowbreak{} c)\), set \(a_j=a_\rho^+(u_{ac})\).
By (TAE.51),
\[
 [G,\allowbreak{} a_j]=b_j,\qquad
 b_j=\begin{cases}
  s_\tau^{-1}\bigl(A_\tau(Xu_{ac})
                   -iA_\tau(XJu_{ac})\bigr),&\rho=\tau,\\
 0,&\rho\ne\tau.
 \end{cases}
 \tag{TAE.59}
\]
Every \(b_j\) kills \(f\).  To compute \([b_j,\allowbreak{} a_k]\),
put
\[
 o_W(a,\allowbreak{} b)=
   2\operatorname{Re}\!
       \left(\xi_\tau
              \bigl((B'_\tau)^*H_TWB'_\tau\bigr)_{ab}\right).
 \tag{TAE.60}
\]
For a real Lie element \(W\), the matrix
\((B'_\tau)^*H_TWB'_\tau\) is skew Hermitian.
Because \(\xi_\tau\) is purely imaginary,
\(o_W(a,\allowbreak{} b)=o_W(b,\allowbreak{} a)\).  If \(sW=-Ws\), its entries with
equal signs \(\epsilon_a=\epsilon_b\) vanish:
multiplication by the two diagonal sign matrices gives
\((\epsilon_a+\epsilon_b)\) times that entry equal to zero.

The original source frame produces the crucial column delta.
Since \(H_S^{-1}=B_\tau B_\tau^*\), cyclicity of trace and
the matrix units give
\[
 \beta_\tau(Wu_{ac},\allowbreak{} u_{bd})
                         =\delta_{cd}\,\allowbreak{} o_W(b,\allowbreak{} a).
 \tag{TAE.61}
\]
More explicitly the trace in (TAE.41), before taking its
real part, reduces to
\[
 \operatorname{tr}
 (E_{bd}^*(B'_\tau)^*H_TWB'_\tau E_{ac});
\]
it is zero unless \(c=d\), and otherwise is the
\((b,\allowbreak{} a)\) entry.  Thus the two original bases and the source
Hermitian inverse have all been retained.

The identities \(JXu=Yu\), \(XJu=-Yu\), and
\(JXJu=Xu\), which follow from (TAE.56), now turn
(TAE.58)--(TAE.61) into the following exact coefficients:
\[
 \begin{split}
 \nu_{ab}&=
 \begin{cases}
   2\bigl(o_X(a,\allowbreak{} b)+i\,\allowbreak{} o_Y(a,\allowbreak{} b)\bigr),
                 &a\ge p_\tau,\ b<p_\tau,\\
   0,&\text{otherwise},
 \end{cases}\\
 \mu_{(\rho,\allowbreak{} a,\allowbreak{} c),(\rho',\allowbreak{} b,\allowbreak{} d)}
   &=\begin{cases}
    -4\pi i\,\allowbreak{} s_\tau^{-2}\delta_{cd}\nu_{ab},
                         &\rho=\rho'=\tau,\\
    0,&\text{otherwise}.
   \end{cases}
 \end{split}
 \tag{TAE.62}
\]
Indeed the scalar before the factor
\(-4\pi i\,\allowbreak{} s_\tau^{-2}\delta_{cd}\) is
\[
 o_X(b,\allowbreak{} a)+i o_Y(b,\allowbreak{} a)
       -i\bigl(-o_Y(b,\allowbreak{} a)+i o_X(b,\allowbreak{} a)\bigr)
                   =2\bigl(o_X(b,\allowbreak{} a)+i o_Y(b,\allowbreak{} a)\bigr).
 \tag{TAE.63}
\]
When the two signs agree this is zero.  When they differ
it is \(\nu_{ab}+\nu_{ba}\), using symmetry of \(o_X,\allowbreak{} o_Y\).
Consequently \([b_j,\allowbreak{} a_k]=(\mu_{jk}+\mu_{kj})1\).
All hypotheses of TAE8 hold on \(\Theta\), and it proves
the exact identity
\[
 G\,\mathrm{pt}(P)=
 \mathrm{pt}\left(
   -4\pi i\,\allowbreak{} s_\tau^{-2}
       \sum_{\substack{a\ge p_\tau\\b<p_\tau}}
           \nu_{ab}\sum_c
             \partial_{\tau,\allowbreak{} a,\allowbreak{} c}\partial_{\tau,\allowbreak{} b,\allowbreak{} c}P
              \right).
 \tag{TAE.64}
\]
The inverse powers of \(s_\tau\) are legitimate by (TAE.41).
In particular (TAE.53) makes this expression zero.
Finally TAE3, followed by evaluation at zero, identifies
\(R_{\tau,\allowbreak{} Z}^{\mathbb C}K_\omega(h,\cdot;\mathrm{pt}(P))\)
with \(K_\omega(h,\cdot;G\,\mathrm{pt}(P))\).
This is zero and proves (TAE.54).
% END TAE10

% BEGIN TAE11
\paragraph{TAE11: the stored product of minors satisfies every required mixed equation.}
\label{claim:TAE11}
Let \(m=r_T,\allowbreak{} n=r_S\), and let the original integer parameters
satisfy \(p_\tau+q_\tau\le n\).
With the source's literal polynomial
\[
 \begin{split}
 \Delta_\tau={}&
 \prod_{k=1}^{p_\tau}
   \det(z_{\tau,\beta,\alpha})_{0\le\alpha,\beta<k}
                   ^{\,\allowbreak{} a^+_{\tau,\allowbreak{} k-1}-a^+_{\tau,\allowbreak{} k}}\\
 &{}\cdot
 \prod_{k=1}^{q_\tau}
   \det(z_{\tau,\allowbreak{} p_\tau+\beta,\allowbreak{} n-k+\alpha})
                      _{0\le\alpha,\beta<k}
                   ^{\,\allowbreak{} a^-_{\tau,\allowbreak{} k-1}-a^-_{\tau,\allowbreak{} k}},
 \qquad \Delta=\prod_\tau\Delta_\tau ,
 \end{split}
 \tag{TAE.65}
\]
the terminal \(a^+_{\tau,\allowbreak{} p_\tau}\) and
\(a^-_{\tau,\allowbreak{} q_\tau}\) mean zero, and the exponents are
the original natural-number differences.
Out-of-range variables have the source value zero.
For every \(a\ge p_\tau,\allowbreak{} b<p_\tau\),
\[
       \sum_c
          \partial_{\tau,\allowbreak{} a,\allowbreak{} c}\partial_{\tau,\allowbreak{} b,\allowbreak{} c}\Delta=0.
 \tag{TAE.66}
\]

\emph{Proof.}
Every variable occurring in a positive minor at \(\tau\)
has row \(<p_\tau\) and column \(<p_\tau\).
Every variable occurring in a negative minor has row
\(\ge p_\tau\) and column at least
\(n-q_\tau\ge p_\tau\).
Thus \(\Delta\) is independent of every variable whose
row and column lie on opposite sides of \(p_\tau\).
This statement survives determinants, products and all
the displayed nonnegative powers.  Variables at other
places are distinct.

For a given column \(c<p_\tau\), the derivative
\(\partial_{\tau,\allowbreak{} a,\allowbreak{} c}\Delta\) is zero.
For \(c\ge p_\tau\), the derivative
\(\partial_{\tau,\allowbreak{} b,\allowbreak{} c}\Delta\) is zero.
Polynomial partial derivatives commute, as is verified
on monomials, so each summand of (TAE.66) is zero.
Summing gives the result.  No monotonicity or irreducibility
assertion about the polynomial representation is needed
for this harmonicity statement.

For the original bundle \(p_\tau+q_\tau=r_T\le r_S\).
Consequently TAE10 applies to the actual
\(\mathrm{RealDat.Delta}(i)\), with its original
Fock normalization, at every parameter \(t\).
% END TAE11

\subsection*{The exact interface to lowest-type and Hilbert realization}

% BEGIN TAE12
\paragraph{TAE12: lowering propagates to all coefficient values under the stated lowest-type hypotheses.}
\label{claim:TAE12}
Let \(c:E\) be the original finite-dimensional coefficient
representation with \(c.\mathrm{Is}\).  Assume explicitly
\(\mathrm{LowSm}(M)\) and \(\mathrm{LowIrr}(M)\), and let
\[
 io:\mathrm{LowT}(M)\longrightarrow
              \operatorname{Hom}_{\mathbb C}(E,\mathcal P)
 \tag{TAE.67}
\]
satisfy the original \(\mathrm{IsIoW}\) with this
\(\omega,\mathrm{pt},\Delta\).
If all kernel lowering derivatives of
\(\mathrm{pt}(\Delta)\) vanish, then for every pure
lowest-type tensor \(v\), every \(w\in E\), and every source
point \(h\), all lowering derivatives of
\[
             g\longmapsto
             K_\omega(h,\allowbreak{} g;\mathrm{pt}(io(v)w))
 \tag{TAE.68}
\]
vanish.  Its integral has the same property.

\emph{Proof.}
For precision, \(\mathrm{IsIoW}\) supplies three things.
First, for every real source element \(u\),
\[
       \omega(u,1)\mathrm{pt}(io(v)w)
                           =\mathrm{pt}(io(v)c(u)w).
 \tag{TAE.69}
\]
Second, a compact target element acting in one model
factor replaces that factor of \(v\), with its literal
frame-conjugated adelic lift on the \(\omega\) side.
Third, there are unit vectors \(v_{0,\tau}\) and
\(w_0\) with \(io(\otimes_\tau v_{0,\tau})w_0=\Delta\).
We use all three clauses.

Let \(\mathcal N\subset\mathcal P\) consist of polynomials
whose kernels are killed by every lowering derivative,
at every source point.  It is a complex vector subspace.
Indeed TAE3 and the two-real-part formula (TAE.40) express
any lowering derivative as a fixed linear operator on
the theta vector followed by \(\omega(h,\allowbreak{} g)\) and
evaluation.  Thus addition and scalar multiplication
preserve the vanishing.

For the fixed tensor \(v_0=\otimes_\tau v_{0,\tau}\),
the subspace
\(\{w:io(v_0)w\in\mathcal N\}\) contains \(w_0\ne0\).
It is invariant under every \(c(u)\).  For a real source
element, (TAE.69) replaces the kernel source variable
\(h\) by \(hu\), and this does not affect a condition
imposed for all \(h\).  A finite source element acts
trivially by \(c.\mathrm{Is}\).
Every adelic element is a product of its real-supported
and finite-supported parts; the original infinity chart
identifies its real part with the same real group.
The irreducibility clause in \(c.\mathrm{Is}\) therefore
makes this subspace all of \(E\).

Now fix \(w\).  The inverse image of \(\mathcal N\)
under \(v\mapsto io(v)w\) contains \(v_0\) and is stable
under each compact target factor.
Here is the derivative check in that assertion.
For a real target point \(k\),
\[
 R_Z^{\mathbb C}(r_k\varphi)
        =r_k R_{\operatorname{Ad}(k^{-1})Z}^{\mathbb C}\varphi,
       \qquad (r_k\varphi)(g)=\varphi(gk).
 \tag{TAE.70}
\]
It follows by differentiating
\(g e^{tX}k=gk e^{t k^{-1}Xk}\), separately in the
two real parts.  Smoothness needed here was proved
by TAE3 and TAE7.  In the model frame a compact
element commutes with \(\operatorname{diag}(\epsilon_a)\),
so it is block diagonal.  Conjugation by it preserves
the negative-row, positive-column block
\(\mathfrak p^-_\tau(M)\).  Elements at other places
leave that block unchanged.  The target clause of
\(\mathrm{IsIoW}\) and (TAE.70) prove the claimed stability.

By \(\mathrm{LowSm}\), each full lowest compact type lies
inside the model's smooth space.  By \(\mathrm{LowIrr}\),
it is a finite-dimensional irreducible compact
representation.  The algebraic span of the orbit of
each nonzero \(v_{0,\tau}\) is therefore its whole
lowest type.  Apply this fact in one tensor factor
at a time.  Multilinearity proves that the entire
algebraic tensor product maps into \(\mathcal N\).
This gives (TAE.68) for every stated \(v,\allowbreak{} w,\allowbreak{} h\).
Finally apply the last assertion of TAE7 to pass
the vanishing through the original source integral.

The two assumptions \(\mathrm{LowSm}\) and
\(\mathrm{LowIrr}\) have been used explicitly.
\(\mathrm{DSModel.ds}\) alone supplies a lowest
irreducible compact subspace; it does not definitionally
say that the supremum of all such subspaces is one copy.
This proof does not claim to have supplied that
multiplicity-one argument.
% END TAE12

% BEGIN TAE13
\paragraph{TAE13: the original lowest-type integral now has the complete analytic and cusp packet, with its remaining producers visible.}
\label{claim:TAE13}
Fix one original \(\mathrm{RealDat}\ R\) on the same bundle.
Retain the literal hypotheses
\(R.\mathrm{toLevCh.Is}(A.\mu_0)\),
\(R.\mathrm{IsModels}\),
\(R.\mathrm{thP\_\allowbreak \omega P}\),
\(R.\mathrm{thP\_\allowbreak vP}\), and
\(R.\mathrm{thP\_\allowbreak pt\_\allowbreak theta}\).
These keep its original full \(\omega_P(i,\allowbreak{} t)\),
vacuum \(v_P(i,\allowbreak{} t)\), polynomial map and source models.
For a fixed \(i\), additionally
suppose that one map \(io_i\), independent of \(t\), satisfies
\[
 \forall t,\quad
 \mathrm{IsIoW}\bigl(
     R.\theta,\ R.MW(i,\mathrm{sg}_0),\ R.cV(i),\
     R.\omega_P(i,\allowbreak{} t),\ R.\mathrm{ptP}(i,\allowbreak{} t),\
     \mathrm{Delta}(i),\ io_i\bigr),
 \tag{TAE.71}
\]
and retain \(\mathrm{LowSm}\) and \(\mathrm{LowIrr}\)
for that same target model.
Put \(F_{i,\allowbreak{} t,\allowbreak{} x}=R.\mathrm{extV}(i,\allowbreak{} t)x\) and
\(L_v=R.\mathrm{ptP}(i,\allowbreak{} t)\circ io_i(\otimes_\tau v_\tau)\).
Then, for every \(t\) and original source Hilbert vector \(x\),
its literal lowest-type integral
\[
 U_{i,\allowbreak{} t,\allowbreak{} x}(v)(g)=
  \int_{R.\mathrm{fdV}(i,\allowbreak{} t)}
   K_{R.\omega_P(i,\allowbreak{} t)}\bigl(h,\allowbreak{} g;L_v(F_{i,\allowbreak{} t,\allowbreak{} x}(h))\bigr)
                       \,\allowbreak{} dR.\nu_V(i,\allowbreak{} t)(h)
 \tag{TAE.72}
\]
satisfies the exact \(\mathrm{IsHolCuspKMap}\) analytic
packet at the original principal target level.
It is linear in \(x\) and multilinear in the lowest-type
vectors.  Its uniform growth exponent can be chosen as
\[
       N_i=2\,\allowbreak{} r_{V_c(i)}r_{W_c(i)}\,\#\mathrm{Tau},
 \tag{TAE.73}
\]
independently of \(t,\allowbreak{} x,\allowbreak{} v\) and of the real derivative word.
This assertion does not construct an element of
\(\mathrm{MMap}(R.MW(i,\mathrm{sg}_0))\).

\emph{Proof.}
The source in (TAE.72) is the same finite-dimensional
coefficient-valued extension used in BTH3.  Its continuity
and rational invariance are the original coefficient-model
identities, and its source quotient has BTH3's compact
representatives.  The map \(L\) in (TAE.4) is, for fixed
\(v\), the literal
\[
             L=R.\mathrm{ptP}(i,\allowbreak{} t)\circ
                          io_i(\otimes_\tau v_\tau).
 \tag{TAE.74}
\]
TAE6 proves absolute convergence, continuity and rational
target invariance.  TAE7 proves smoothness and uniform
moderate growth with (TAE.73), since the real Hom space
has exactly that dimension.  This exponent is selected
before the coefficient map, the vector and the derivative
word; only their bounding constants change.

The original finite covariance and principal-level fixing
were proved in BTH5 and BTH15, with the same action and
the same frames.  Only the finite vacuum calculation in
those proofs is used here; its subsequent extension to a
full multiplicity map is not used.  For a principal
target element \(k\), take its matched source element
to be \(1\).  Every reduced block is the identity, so
all determinant and half-support multipliers in
BTH.19 are one.  Thus \(\omega(1,\allowbreak{} k)v_P=v_P\).
Finite action commutes with every real creation
operator by the \(L^2\) derivative transport of TAE5,
so it fixes every original polynomial vector.
This calculation is pointwise in \(t\), uses only
the stated level and vacuum fields, and applies
unchanged at each \(t\).  The original principal
target character is one.  Inserting this identity under
the integral gives the required right level law.
The compact target clause of (TAE.71), inserted under
the same integral, gives its exact model-frame compact
equivariance.  Multilinearity and linearity in \(x\)
follow from linearity of the kernel and of the original
coefficient extension, using the absolute convergence
already established.

The original rank inequality \(r_W\le r_V\) and signature
order give the hypotheses of TAE11 for
\(\mathrm{Delta}(i)\).  The actual \(\mathrm{IsThPVecOf}\)
and polynomial-theta fields give the Fock-vacuum equations
for \(v_P(i,\allowbreak{} t)\), by the explicit
\(\mathrm{AllSome}\) argument at the end of TAE9.
It changes neither its scale nor its polynomial map.
TAE10 and TAE11 kill the lowering kernel of
\(\mathrm{ptP}(i,\allowbreak{} t)(\mathrm{Delta}(i))\).
TAE12 propagates this to every value of (TAE.74), and
TAE7 passes it through the integral.

FIA8 gives the finite isotropic average for every
polynomial and every source and target translate.
CTC9--CTC13 therefore give the zero transvection
constant terms and genuine full cuspidality of
(TAE.72).  Their unipotent quotient integrals converge;
no undefined-integral convention is used here.
Together these are exactly the eight fields
of the source's \(\mathrm{IsHolCuspKMap}\):
continuity, rational left invariance, level law,
smoothness, uniform growth, cuspidality, compact
equivariance and lowering.

Finally (TAE.71) is the original \(\mathrm{PIoP}(i,\allowbreak{} io_i)\).
It is not replaced by the weaker quantifier order
\(\forall t\,\exists io_{i,\allowbreak{} t}\).
If the original \(R.\mathrm{ioP}(i)\) is its stored
\(\mathrm{OrZero}\) choice, existence of a witness
with (TAE.71) activates that choice by the definition
of \(\mathrm{OrZero}\).  One may then use the actual
stored map in (TAE.72).
Nothing in this implication supplies such a witness,
or supplies \(\mathrm{LowSm}/\mathrm{LowIrr}\).
Nothing in it extends \(U_{i,\allowbreak{} t,\allowbreak{} x}\) from the lowest
compact type to all smooth vectors of the target
discrete-series model.  Those are the remaining
interfaces, stated next.
% END TAE13

\subsection*{The exact remaining producer}

The analytic part just reconstructed is pointwise on one
supplied original \(R\).  It introduces no new family
threshold, coefficient field, Hermitian adjoint, or theta
matrix.  At the base tuple it uses the same
\(t_0=A_3.A.\mu_0\) as BTH12 and SLR8.  Its consequences
can also be applied at all \(t\), with the common
polynomial map quantified as in (TAE.71).

The residual common-object obligations are as follows.
\begin{enumerate}
\item Produce one \(io_i\) satisfying (TAE.71), with both
original intertwining laws and the unit-vector
normalization to the literal \(\mathrm{Delta}(i)\).
The source supplier is
\texttt{Joint\allowbreak Harm.\allowbreak exists\_\allowbreak io\_\allowbreak poly}, followed by
\texttt{Joint\allowbreak Harm.\allowbreak exists\_\allowbreak io\allowbreak P}.
Its finite polynomial lowest-weight transfer has not
been independently reconstructed in this chapter.
\item Prove \(\mathrm{LowSm}\) and
\(\mathrm{LowIrr}\) for the full original
\(\mathrm{lowKSub}\).
The source supplier is
\begin{quote}\small
\texttt{Hol\allowbreak Low\allowbreak Once.\allowbreak low\allowbreak Sm\_\allowbreak and\_\allowbreak low\allowbreak Irr}.
\end{quote}
In particular its multiplicity-one step uses a
holomorphic coefficient uniqueness argument.
The earlier nonzero smooth-vector and cyclicity
argument in BTH1 is not a substitute for that step.
\item Supply boundedness for the actual smooth cuspidal
uniform-moderate-growth lowest-type form, and the
finite-volume target quotient needed by the Hilbert
realization.  The source route uses
\(\mathrm{Published.CuspRapidDecay}\) through
\texttt{bounded\_\allowbreak of\_\allowbreak smooth\_\allowbreak
cuspidal\_\allowbreak uniform\_\allowbreak moderate}.
Naming this predicate does not certify a published
application: its hypotheses must still be matched,
or the needed instance proved directly.
\item Reconstruct the actual bounded holomorphic
coefficient identity and smooth multiplicity-map
realization.  The immediate suppliers are
\begin{quote}\small
\texttt{inner\_\allowbreak reg\_\allowbreak eq\_\allowbreak coef\_\allowbreak extremal},\\
\texttt{bounded\allowbreak Holomorphic\allowbreak Form\_\allowbreak eq\_\allowbreak multiplicity\allowbreak Space\_\allowbreak value}.
\end{quote}
The remaining conclusion is an actual
\[
 \Theta_{i,\allowbreak{} t,\allowbreak{} x}\in\mathrm{MMap}(R.MW(i,\mathrm{sg}_0))
\]
satisfying \(\mathrm{IsMultMap}\), whose lowest-type
restriction is exactly (TAE.72), with its original
Sobolev continuity, finite level and central-finiteness.
\end{enumerate}
These are retained C producer subcontracts in the present
audit, not four new independent top-level obligations.
The proved implications TAE1--TAE13 are B.
No internal D translation issue is being concealed by
calling any of these producers automatic.
The geometric Hilbert reading, the simultaneous rational
tests and the original base-row producer in BTH12/SLR8
remain separate from this analytic extension.


\section{Exact reciprocity and the original irreducible-point normalization}
\label{sec:count-ii-producer}

\status{CIP1--CIP12 close two specified inputs of the original CountII
route.  They supply the exact reciprocity matching retained in JSP7,
and the complete \texttt{Nice\allowbreak I} instance of \texttt{Norm} used in
JSP11.  The latter statement is universal in the original point and
includes its integral representation, ordered characters, original
$D_0$, and determinant law.  The general \texttt{Nice\allowbreak E} instance is
outside this conclusion.  The ordinary arithmetic system, four further
Hida fields, all-level \texttt{Has\allowbreak Aut}, and five other single-point
support inputs remain explicit.  In particular
\texttt{COUNT\_\allowbreak II\_\allowbreak SUPPORT} is not asserted to be closed.}

\subsection{The reciprocity convention on the actual local groups}

Retain the original CM field $L$, its maximal real subfield $L^+$,
the chosen algebraic closure $\Omega$, and the original
$s:\mathrm{HidaSetting}(\ell,\allowbreak{} L,\Omega)$.  For $v\mid\ell$ put
$w=s.w(v)$ and $K=L_w$.  The valuation on $\Omega$ is the given
$s.A(v)$, not a newly selected place.  KX3 identifies its decomposition,
inertia and Weil groups with those of $K$.  A Weil element of degree
one is an arithmetic Frobenius.  The source instead requires
\[
 \operatorname{ord}_K(a(\sigma))=-\deg(\sigma).
 \tag{CIP.1}
\]
Write $\operatorname{rec}^{\mathrm{ar}}_K$ for reciprocity sending a
uniformizer to arithmetic Frobenius.  Define its geometric counterpart by
\[
 \operatorname{rec}^{\mathrm{geo}}_K(u)
       =\operatorname{rec}^{\mathrm{ar}}_K(u)^{-1}.
\]
Thus the desired map from the Weil group is the inverse of geometric
reciprocity, not the inverse of arithmetic reciprocity.

The published input in this subsection is classical local and global
class field theory: the reciprocity isomorphisms, local--global and
norm compatibility, and the global kernel and surjectivity statements.
Their precise primary reference is J.~S.~Milne,
\emph{Class Field Theory}, version 4.03:
Introduction, p.~15; I.1.1 and I.1.10--1.11, pp.~20, 23--24;
III.3.2, p.~107; V.5.2--5.5 and Remark 5.7, pp.~178--180.
The inverse convention is imposed explicitly below.
\footnote{\url{https://www.jmilne.org/math/CourseNotes/CFT.pdf}.
The page numbers are the printed numbers; the PDF indices are larger
by eight.  V.4.10, pp.~175--176, also checks the inverse exponent on
rational local units.  The accompanying application certificate
separates these published inputs from the elementary matching proofs.}

% BEGIN CIP1
\paragraph{CIP1: the source-compatible global reciprocity map.}
\label{claim:CIP1}
For every number field $E$ and the specified algebraic closure
$\Omega$, inverse arithmetic reciprocity gives a map
\[
 \alpha_E:\mathbb A_E^\times\longrightarrow G_E^{\mathrm{ab}}
\]
satisfying all five fields of the source \texttt{Is\allowbreak Global\allowbreak Artin}.
\begin{proof}
Apply global reciprocity to $E$ and identify its abelian Galois group
with the topological abelianization of $\operatorname{Gal}(\Omega/E)$
by restriction to its maximal abelian subextension.  The theorem
applies to a number field, with its usual finite and infinite
completions; these are precisely the original idele factors.
Invert the arithmetic Artin map.  Inversion is a continuous
automorphism of this abelian target, so it preserves surjectivity and
the kernel
\[
 \ker\alpha_E=
 \overline{E^\times\, E_\infty^{\times,0}}.          \tag{CIP.2}
\]
Here $E_\infty^{\times,0}$ is exactly the source's
\texttt{inf\allowbreak Conn}: its real coordinates are positive and its complex
coordinates unrestricted.  This proves the \texttt{ker} field.

The continuous surjection is open.  For completeness, the idele group
is locally compact and a countable union of compact sets: at the
finite places the compact full-unit subgroup has countably many
valuation cosets, and the archimedean factors are countable unions
of compact annuli.  If $U$ is a neighborhood of one, choose a compact
neighborhood $C$ with $CC^{-1}\subset U$.  Countably many translates
of $C$ cover the domain.  Their images are closed compact sets
covering the compact Hausdorff target.  The Baire theorem gives
nonempty interior to one such image; consequently
$\alpha_E(C)\alpha_E(C)^{-1}$ contains a neighborhood of one.
It lies in $\alpha_E(U)$.  Applying this argument inside each
translate of an open set proves openness.  This gives the exact
\texttt{quot} field, including continuity and the quotient property.

Local--global compatibility identifies the image of local units with
the image of local inertia.  The unit group is compact, so its image
is closed, and the local reciprocity identification is onto this
inertia image.  This is the two-way \texttt{ine} assertion.  A
uniformizer under $\alpha_E$ has the image of an inverse arithmetic
Frobenius.  Multiplication by units runs over the whole inertia
coset, which is precisely the set of images of all such inverse
Frobenius elements.  This is the two-way \texttt{frob} assertion,
with the source's inverse sign.  Finally the closure of the local
multiplicative image is the image of the local decomposition group:
local reciprocity has dense image in its abelianization, and the
decomposition group is compact.  This is \texttt{dec}.  The local
identifications preserve the specified valuation on $\Omega$, as in
KX3, so none of these fields uses an unrelated local subgroup.
\end{proof}
% END CIP1

% BEGIN CIP2
\paragraph{CIP2: uniqueness under the literal global predicate.}
\label{claim:CIP2}
If $E$ is CM, every map satisfying the source
\texttt{Is\allowbreak Global\allowbreak Artin} equals $\alpha_E$ of CIP1.
\begin{proof}
Let $\alpha,\beta$ be two such maps, and let $q$ be any finite
continuous quotient of $G_E^{\mathrm{ab}}$.  The characters
$q\alpha,\allowbreak{} q\beta$ are continuous.  Choose a finite set $S$ of finite
places such that both kill the full product of unit groups away
from $S$.  This follows from an open neighborhood of one in their
kernels and the definition of the restricted-product topology.
Their ratio $\delta=(q\alpha)(q\beta)^{-1}$ kills diagonal $E^\times$
and all infinite ideles: all infinite places of a CM field are
complex and hence connected.

At a finite place $u\notin S$, choose a uniformizer $\pi_u$ and a
valuation over $u$ on $\Omega$.  The \texttt{frob} field for
$\alpha$ expresses $\alpha(j_u\pi_u)$ as the image of an inverse
arithmetic Frobenius $\sigma^{-1}$.  The reverse implication of the
\texttt{frob} field for $\beta$, using this same $\sigma$, supplies
another uniformizer $y_u$ with the same global image.  Since
$y_u/\pi_u$ is a unit killed by $q\beta$, the two maps agree at
$\pi_u$ after $q$.  They already agree on local units.  Therefore
$\delta$ is trivial on every $E_u^\times$ for $u\notin S$.

Fix $u\in S$ and $z\in E_u^\times$.  Weak approximation supplies
$a\in E^\times$ sufficiently close to $z$ at $u$ and to one at the
other places in $S$ that the corresponding local values of the
finite-valued continuous character $\delta$ agree.  On a diagonal
idele its value is one.  The factors away from $S$ and at infinity
are already trivial, so $\delta(j_u z)=1$.  Thus all local factors
are trivial.  An idele is a product of finitely many nonunit local
factors and an element of the full unit product; the latter is
handled by the same choice of $S$.  Hence $\delta=1$ on all ideles.
Finite continuous quotients separate points of the profinite group
$G_E^{\mathrm{ab}}$, giving $\alpha=\beta$ before taking quotients.
This argument uses the actual \texttt{frob}, \texttt{quot} and
\texttt{ker} fields, rather than assuming uniqueness as an extra
property of that predicate.
\end{proof}
% END CIP2

% BEGIN CIP3
\paragraph{CIP3: local-unit injectivity at the actual split CM places.}
\label{claim:CIP3}
Let $E$ be CM with complex conjugation $c$, and let $w\ne c(w)$
be a finite place.  If $u\in\mathcal O_{E_w}^{\times}$ and
$\alpha_E(j_wu)=1$, then $u=1$.  Every original $w=s.w(v)$ for
$v\mid\ell$ has this distinctness property.
\begin{proof}
Project (CIP.2) to the finite ideles.  The idele $j_wu$ belongs to
the closure of diagonal $E^\times$.  The full-unit subgroup
$U_f=\prod_z\mathcal O_{E_z}^{\times}$ is open, and
$E^\times\cap U_f=\mathcal O_E^\times$.  Therefore $j_wu$ lies in
the closure, inside $U_f$, of the diagonal global units.

For a global unit $e$, the element $e/c(e)$ is an algebraic integer
unit all of whose complex conjugates have absolute value one.
It is a root of unity.  One elementary proof is to take its powers:
their monic minimal polynomials have bounded degree and bounded
integer coefficients, since all their roots have modulus one.
There are only finitely many such polynomials and roots, so two
powers coincide.  The same argument bounds the set of all roots
of unity in the fixed field $E$, proving that $\mu(E)$ is finite.

The continuous map $y\mapsto y/c(y)$ on $U_f$ consequently carries
$j_wu$ into the finite, hence closed, diagonal copy of $\mu(E)$.
At any finite place outside $\{w,\allowbreak{} c(w)\}$ this idele has value one.
The embedding of $E$ in that completion is injective, so the root
of unity is one.  Thus $j_wu=c(j_wu)$.  Their $w$-coordinates are
$u$ and one because $w\ne c(w)$, proving $u=1$.

For the original data, $s.\mathrm{split}$ says that the
decomposition subgroup at $v$ in $G_{L^+}$ fixes $L$.
The image of this subgroup in $\operatorname{Gal}(L/L^+)$ is the
decomposition subgroup of $w/v$: surjectivity follows by extending
a place-preserving automorphism and adjusting its extension by an
element of $G_L$.  That image is trivial.  Since $L/L^+$ is
quadratic, $v$ has two distinct places $w,\allowbreak{} c(w)$ in $L$.
Equivalently, the completion in KX3 is $L_w\simeq L_v^+$.
Thus the hypothesis just used is supplied by the actual setting
at every place required by JSP7.  No assertion about nonsplit
places or a generic local-idele injectivity theorem is needed.
\end{proof}
% END CIP3

% BEGIN CIP4
\paragraph{CIP4: validity and uniqueness of the chosen \texttt{art\allowbreak Of}.}
\label{claim:CIP4}
For each original $v\mid\ell$, the source's
$\operatorname{artOf}(s.hw(v))$ satisfies \texttt{Is\allowbreak Local\allowbreak Artin}.
Every map satisfying that predicate equals this map on the entire
original Weil group.
\begin{proof}
On $K=L_w$ use the local class field theory isomorphism
$K^\times\simeq W_K^{\mathrm{ab}}$, with its usual Weil topology,
and put
\[
 a_0(\sigma)
   =(\operatorname{rec}^{\mathrm{ar}}_K)^{-1}
            ([\sigma]^{-1})
   =(\operatorname{rec}^{\mathrm{geo}}_K)^{-1}([\sigma]).
 \tag{CIP.3}
\]
KX3's completion identification puts this map on the original
\texttt{Weil}(s.hw(v)).  It is onto, continuous on the original
inertia, and satisfies (CIP.1).  These give
\texttt{surj}, \texttt{cont}, and \texttt{ord\_\allowbreak eq}.

A smooth complex character of the Weil group has open kernel on
inertia, exactly as in the source definition.  Its kernel is open
in the Weil topology because inertia is open; the character is
therefore continuous.  It factors through the topological
abelianization, hence through $a_0$.  The quotient map to the
abelianization is open, so the resulting character of $K^\times$
has open kernel.  This proves the
literal \texttt{factor} field.  Local--global compatibility gives
\[
 \alpha_L(j_w a_0(\sigma))=\operatorname{toAb}(\sigma).
 \tag{CIP.4}
\]
CIP2 identifies every source \texttt{Is\allowbreak Global\allowbreak Artin} map with
$\alpha_L$, so (CIP.4) proves \texttt{glob} with its universal
quantifier, not just for one preferred global map.  Thus $a_0$
satisfies the full predicate.

If $a$ is any other map satisfying that predicate, then
$a(\sigma)/a_0(\sigma)$ is a unit by \texttt{ord\_\allowbreak eq}.
The \texttt{glob} field and (CIP.4) put its local idele in
$\ker\alpha_L$.  CIP3 forces this ratio to be one, for every
Weil element $\sigma$.  Thus $a=a_0$.  The source
\texttt{art\allowbreak Of} is the Hilbert-choice element of this nonempty
predicate, so it is valid and equals $a_0$ too.  In particular the
universal quantifier over local Artin maps in \texttt{Is\allowbreak Wt\allowbreak Part}
can be replaced by this one map with no change of meaning.
Finally (CIP.1) and surjectivity show that the restriction from
inertia is onto $\mathcal O_K^\times$: a preimage of a unit has
degree zero, which is exactly inertia.
\end{proof}
% END CIP4

% BEGIN CIP5
\paragraph{CIP5: the positive cyclotomic norm and the exact JSP7 input.}
\label{claim:CIP5}
For every original inertia element $\sigma$ and every source
\texttt{Is\allowbreak Local\allowbreak Artin} map $a$, one has, in
$\overline{\mathbb Q}_\ell$,
\[
 \epsilon(\sigma)=N_{K/\mathbb Q_\ell}(a(\sigma))
       =\prod_{\tau\in\mathrm{Emb}(\ell,\allowbreak{} w)}\tau(a(\sigma)).
 \tag{CIP.5}
\]
This supplies all of \texttt{JSP\_\allowbreak EXACT\_\allowbreak RECIPROCITY\_\allowbreak MATCH}.
\begin{proof}
First fix the sign over $\mathbb Q_\ell$.  Arithmetic reciprocity
on a unit $z\in\mathbb Z_\ell^\times$ acts on the
$\ell^m$th roots of unity by $z^{-1}$ for every $m$.
Here is a direct check consistent with the published example.
Choose a positive integer $b$, prime to $\ell$, sufficiently close
to $z^{-1}$ that $zb$ acts trivially on these roots under local
reciprocity, and also $b\equiv z^{-1}\pmod{\ell^m}$.
Continuity permits this choice.  Global arithmetic reciprocity
kills the principal idele $b$, and its real positive component.
At a prime $p\ne\ell$, the cyclotomic extension is unramified and
arithmetic Frobenius acts as the $p$th power.  Indeed its roots of
unity lie in unramified extensions by Hensel lifting of the
separable polynomial $X^{\ell^m}-1$ over finite residue fields.
The product of these away-from-$\ell$ actions is raising to
$\prod_{p\ne\ell}p^{v_p(b)}=b$.  The action of the $\ell$-component
$b$ is consequently raising to $b^{-1}$, and that of $z$ is
raising to $b$, hence to $z^{-1}$ modulo $\ell^m$.

Norm compatibility for local reciprocity now says that
$\operatorname{rec}^{\mathrm{ar}}_K(u)$ acts on all
$\ell$-power roots by $N_{K/\mathbb Q_\ell}(u)^{-1}$ for
$u\in\mathcal O_K^\times$.  CIP4 identifies
$\sigma$ in the local abelianization with
$\operatorname{rec}^{\mathrm{geo}}_K(a(\sigma))$.
Inverting the arithmetic action gives the positive norm in
(CIP.5).  This proves the identity for every admissible $a$.

The chosen valuation lies over both $v$ and $w$.  Since $v\mid\ell$,
its restriction places $\ell$ in $w$, so $K=L_w$ is a finite
extension of $\mathbb Q_\ell$ as required for the norm comparison.
The source set $\mathrm{Emb}(\ell,\allowbreak{} w)$ consists of continuous ring
embeddings $K\to\overline{\mathbb Q}_\ell$.  Such a map fixes
$\mathbb Q_\ell$ because it fixes $\mathbb Q$ and is continuous.
Conversely every $\mathbb Q_\ell$-embedding of the finite extension
$K$ is continuous by uniqueness of its extended valuation.
The extension is separable, so the norm is the product of exactly
these $[K:\mathbb Q_\ell]$ embeddings.  Thus the source's finite
product has precisely the factors in (CIP.5).

The definition of \texttt{eps\allowbreak V} is the original cyclotomic character
followed by the structural coefficient maps.  The field
$V.h\epsilon$ makes the same identification for the level frame.
Consequently (CIP.5) gives exactly the cyclotomic factor in
\texttt{Is\allowbreak Theta} and \texttt{Is\allowbreak Univ}, including its exponent $i$.
It is also \texttt{Is\allowbreak Wt\allowbreak Part} with weight $-1$ and finite character
one: its displayed exponent is $-(-1)=1$.  Together with CIP4 this
is every validity, every-map uniqueness, and normalization clause
retained in JSP7.  Its proof of the increasing-weight and finite-part
dictionary can therefore be used without that additional input.
Separation of actual good arithmetic fibers remains a distinct
antecedent of JSP7; the reciprocity calculation does not produce it.
\end{proof}
% END CIP5

\subsection{The complete normalization needed at an original NiceI point}

Fix an odd prime $\ell$, the original $F:\mathrm{GFrame}(\ell)$,
a local $\mathbb Z_\ell$-algebra $O$, $n>0$, a residual datum $r$,
and an actual $x:\mathrm{Point}(r)$ satisfying \texttt{Nice\allowbreak I}.
Put $A=x.K[[t]]$, $F_A=x.K((t))$, and
$G'=\ker(\epsilon|_{G_0})$.  Here $G'$ uses the full
$\mathbb Z_\ell$-valued cyclotomic character of the frame.
The original representation and its form are
\[
 \rho:G_0\longrightarrow\mathrm{GL}_n(A),\qquad
 P=P^t\in\mathrm{GL}_n(A),\qquad
 \rho(g)^tP\rho(cgc^{-1})=\epsilon(g)^{1-n}P.       \tag{CIP.6}
\]
The invertibility of $P$ over $A$, not just $F_A$, is part of
the original \texttt{Point}.  Its ordered local characters are
$\psi_{v,0},\ldots,\psi_{v,\allowbreak{} n-1}$, and its flags are over $F_A$.

% BEGIN CIP6
\paragraph{CIP6: a finite constant extension standardizes the integral form.}
\label{claim:CIP6}
There is a finite extension $k'/x.K$, a coefficientwise map
$j:A\to A'=k'[[t]]$, and $T\in\mathrm{GL}_n(A')$ such that
\[
  T^t j(P)T=1.                                      \tag{CIP.7}
\]
The map $j$ is injective, local and finite free, sends $t$ to $t$,
and commutes with the original $\mathbb Z_\ell$-structure maps.
\begin{proof}
The composite $\mathbb Z_\ell\to O\to x.K$ makes $x.K$ a finite
field of characteristic $\ell$.  If its characteristic were a
different prime $p$, the unit $p\in\mathbb Z_\ell$ would map to
zero, which is impossible.  Moreover every map from
$\mathbb Z_\ell$ to this characteristic-$\ell$ ring factors
through the prime field $\mathbb F_\ell$; this also proves the
asserted compatibility after extension.  In particular two is a
unit.

Diagonalize $P$ integrally.  Modulo $t$, some vector has nonzero
self-pairing: otherwise the polarization identity
$2B(u,\allowbreak{} v)=B(u+v,\allowbreak{} u+v)-B(u,\allowbreak{} u)-B(v,\allowbreak{} v)$ would make the nondegenerate
bilinear form zero.  Lift this vector to a primitive vector in
$A^n$; its pairing with itself is a unit.  A primitive vector
extends to a basis over the local ring, and subtracting its
orthogonal projections from the other basis vectors gives a direct
orthogonal complement.  The complementary Gram matrix is again
invertible.  Induction gives
\[
 U^tPU=\operatorname{diag}(d_1,\ldots,\allowbreak{} d_n),\qquad
 U\in\mathrm{GL}_n(A),\quad d_i\in A^\times.
\]
Choose a finite extension $k'/x.K$ containing square roots of the
finitely many nonzero constant coefficients $d_i(0)$; the splitting
field of their quadratic polynomials is a finite field extension.
Each $j(d_i)$ has a square root $b_i\in A'^\times$.  Explicitly,
after choosing $b_i(0)$, the coefficient of $t^m$ in
$b_i^2=j(d_i)$ determines the new coefficient uniquely by division
by $2b_i(0)$.  Set $T=j(U)\operatorname{diag}(b_i^{-1})$.
This proves (CIP.7) with both $T$ and $T^{-1}$ integral.

Coefficientwise extension is injective and preserves the constant
coefficient, so it preserves and reflects units; this is the
required local-homomorphism property.  If $e_1,\ldots,\allowbreak{} e_d$ is an
$x.K$-basis of $k'$, decomposing every coefficient of a power series
in this basis gives a unique expression
$\sum_{a=1}^d e_a f_a(t)$ with $f_a\in A$.
Hence $A'$ is free of rank $d$ over $A$.  This proves ring-map
finiteness, rather than only finiteness of the residue extension.
\end{proof}
% END CIP6

% BEGIN CIP7
\paragraph{CIP7: the full block and its integral span.}
\label{claim:CIP7}
Define
\[
 \rho'(g)=T^{-1}j(\rho(g))T,\qquad
 \psi'_{v,\allowbreak{} i}=j\psi_{v,\allowbreak{} i}.                            \tag{CIP.8}
\]
The $F_A'=k'((t))$-span of $\rho'(G')$ is
$M_n(F_A')$.  Consequently for every $Z\in M_n(A')$ there is
$s\ge0$ such that
\[
 t^s Z\in\operatorname{span}_{A'}\{\rho'(g):g\in G'\}.
 \tag{CIP.9}
\]
Thus the source's block and span fields hold with $a=n$.
\begin{proof}
The original \texttt{Type\allowbreak I} says exactly that $\rho|_{G'}$ is
irreducible after extension to an algebraic closure of $F_A$.
We give the finite-dimensional density argument it requires.
Over an algebraically closed field $C$, let $R$ be the $C$-span
of an irreducible group action on $V=C^n$.  It is a unital algebra,
since products of group matrices are group matrices.  Every
$R$-linear endomorphism of $V$ is scalar: it has an eigenvalue,
and the kernel of its difference with that scalar is a nonzero
invariant subspace, hence all of $V$.

For linearly independent $v_1,\ldots,\allowbreak{} v_r$, prove by induction that
$R\to V^r$, $a\mapsto(av_1,\ldots,\allowbreak{} av_r)$, is onto.  For $r=1$
this is irreducibility.  At the induction step its image $W$ is
an $R$-submodule whose projection on the first $r-1$ factors is
onto.  The kernel in the last factor is either zero or $V$.
In the latter case $W=V^r$.  In the former case $W$ is the graph
of an $R$-linear map $V^{r-1}\to V$.  Its component maps are scalar
by the preceding paragraph.  Applying this graph identity to
$a=1$ would express $v_r$ as a linear combination of its
predecessors, a contradiction.  At $r=n$ it follows that $R$
contains every endomorphism of $V$.

Apply this to the given absolute irreducibility.  The span over
$F_A$ must already have dimension $n^2$, since extension of scalars
preserves the dimension of a subspace and its matrix inclusion.
Thus finitely many matrices $\rho(g)$ form an $F_A$-basis of
$M_n(F_A)$.  After the coefficient extension and conjugation by
$T$ they still form a basis, proving the claimed full span.
Express a given integral $Z$ in this basis.  Each of its finitely
many coefficients lies in $k'((t))$; multiplying by a sufficiently
large power of $t$ makes all coefficients integral and proves
(CIP.9).  For $a=n$ every index of $\mathrm{Fin}(n)$ belongs to
the first block, so the literal space
$\mathrm{Ev}(A',\allowbreak{} A',\allowbreak{} n,\allowbreak{} n)$ is all of $M_n(A')$.
No factorial, trace denominator or characteristic-zero argument
has been used.
\end{proof}
% END CIP7

% BEGIN CIP8
\paragraph{CIP8: the original characters, torsion and integral flags.}
\label{claim:CIP8}
For every $v$, the characters in (CIP.8) are trivial on the
original subgroup $F.\mu(v)$ and admit an integral flag basis
$w_v\in\mathrm{GL}_n(A')$ in their original order:
\[
 w_v^{-1}\rho'(\phi_v(g))w_v
 \text{ is upper triangular with diagonal }
 (\psi'_{v,0}(g),\ldots,\psi'_{v,\allowbreak{} n-1}(g)).           \tag{CIP.10}
\]
Their complete original \texttt{N1} inequalities are preserved.
\begin{proof}
First prove the torsion assertion for the original $\psi_{v,\allowbreak{} i}$.
In the original rational flag basis, its value is a diagonal
entry of a conjugate of $\rho(\phi_v(g))$; it is therefore a
fixed finite $F_A$-linear combination of these matrix entries.
Clear all denominators by a power $t^N$.  The equation
$\psi_{v,\allowbreak{} i}(g)=1$ is then a finite integral linear equation in
these entries.  Modulo $t^m$ its solution set is closed, indeed
clopen, by the original entrywise \texttt{Adic\allowbreak Cont} and
continuity of $\phi_v$.  Intersecting over $m$, and using
$\bigcap_m t^mA=0$, shows that $\ker\psi_{v,\allowbreak{} i}$ is closed.
It contains the commutator subgroup because the target is abelian,
so it contains its topological closure.

By definition $F.\mu(v)$ is the inverse image of the
$\ell$-primary torsion in the topological abelianization of $G_v$.
For $g$ in this subgroup, some $g^{\ell^b}$ belongs to that closed
commutator subgroup.  Therefore $\psi_{v,\allowbreak{} i}(g)^{\ell^b}=1$.
In the domain $A$ of characteristic $\ell$ this implies
$(\psi_{v,\allowbreak{} i}(g)-1)^{\ell^b}=0$, hence
$\psi_{v,\allowbreak{} i}(g)=1$.  Applying $j$ proves the same assertion for
$\psi'_{v,\allowbreak{} i}$.  This proof does not assume an additional continuity
field for the original characters.

Extend the given rational flag coefficientwise and conjugate it
by $T^{-1}$.  It is a full flag
$0=V_0\subset\cdots\subset V_n=(F_A')^n$ for $\rho'$ with
successive characters $\psi'_{v,0},\ldots,\psi'_{v,\allowbreak{} n-1}$.
Set $L_i=V_i\cap(A')^n$.  It is a finite free module of rank $i$:
it is a submodule of the finite free module over the DVR, and
clearing denominators of a basis of $V_i$ gives its rational rank.
It is saturated, since $az\in V_i$ for nonzero $a$ implies
$z\in V_i$.  The successive quotients $L_i/L_{i-1}$ are therefore
finite torsion-free modules of rank one, hence free.  Choose their
bases successively and lift them.  As $L_n=(A')^n$, the resulting
adapted basis is integral and invertible; let its matrix be $w_v$.
The action and its inverse preserve each $L_i$.  Its quotient
character agrees over the fraction field with the original ordered
$\psi'_{v,\allowbreak{} i}$, hence agrees integrally.  This proves (CIP.10).
Choose $\mathrm{ty}_v(j)=\mathrm{true}$ for every index.  Since
$a=n$, the antecedent in the original \texttt{w\_\allowbreak blk} field is
always false, so that field holds as well.

Finally injectivity of $j$ preserves inequalities of characters.
CIP6 identifies the two coefficient images of $\epsilon$.
An equality $\psi'_{v,\allowbreak{} i}=\psi'_{v,\allowbreak{} j}\epsilon\circ\phi_v$
would therefore descend to the forbidden original equality.
This proves both parts of \texttt{N1}, with exactly the original
index order $i<j$ in the second part.
\end{proof}
% END CIP8

% BEGIN CIP9
\paragraph{CIP9: the full determinant law and both nondegeneracy fields.}
\label{claim:CIP9}
The representation (CIP.8) has the image under $j$ of the entire
original determinant law, satisfies the source's normalized
\texttt{N3}, and satisfies \texttt{N4} for the original subgroup
$F.D_0$.
\begin{proof}
Let $C$ be any commutative $A'$-algebra and let
$y=\sum_g a_g[g]\in C[G_0]$.  Naturality of matrix extension and
(CIP.8) give
\[
 \sum_g a_g\rho'_C(g)
       =T_C^{-1}\left(\sum_g a_g\rho_C(g)\right)T_C.
\]
Determinants are unchanged by this conjugation.  This proves
\[
 \mathrm{GDet.ofRep}(\rho')
             =j_*\mathrm{GDet.ofRep}(\rho)           \tag{CIP.11}
\]
as polynomial laws on every coefficient algebra, rather than only
as characteristic polynomials of individual group elements.

For \texttt{N3}, take any original open set $U$ meeting $G'$.
The given $x.N3$ supplies $\sigma\in U\cap G'$ whose characteristic
polynomial over $F_A$ is separable.  After the field extension and
conjugation its characteristic polynomial is the coefficient image
of the same polynomial.  A Bezout identity with its derivative
survives the field extension, so it remains separable.  This is
exactly the normalized \texttt{N3} field, on the same $U$ and
the same $\sigma$.

For \texttt{N4}, over $F_A$ the original condition is that the
linear equations
\[
 \epsilon(g)\rho(g)X-X\rho(g)=0
                      \qquad(g\in F.D_0)            \tag{CIP.12}
\]
have only the zero solution in $n^2$ matrix entries.  Their
coefficient functionals consequently span the full dual of
$M_n(F_A)$; otherwise a nonzero annihilator would solve all of
them.  Choose a finite basis of these functionals.  The resulting
square coefficient matrix has nonzero determinant, which remains
nonzero over $F_A'$.  Thus (CIP.12) has no nonzero solution even
over the extended field.  If an integral matrix $Z\in M_n(A')$
satisfies the normalized equations, conjugate it by $T$ and view
it in $M_n(F_A')$.  By CIP6 the cyclotomic scalars agree, so
$TZT^{-1}$ satisfies the extended (CIP.12).  It is zero; hence
$Z=0$.  Throughout this argument the subgroup is the original
$F.D_0$.  It has not been replaced by the bottom subgroup stored
in \texttt{re\allowbreak Iu}.
\end{proof}
% END CIP9

% BEGIN CIP10
\paragraph{CIP10: all fields of the exact NiceI normalization.}
\label{claim:CIP10}
For every original $I_u\le r.I_p(u)$ there exist
\[
 x':\mathrm{NP}(\mathrm{reIu}(F,\allowbreak{} I),\allowbreak{} n),\qquad
 j:x.A\longrightarrow x'.A
\]
with $j$ injective, local and finite, satisfying
\[
 \begin{gathered}
 \mathrm{GDet.ofRep}(\mathrm{mat}\,\allowbreak{} x'.\rho)
       =(\mathrm{GDet.ofRep}(\mathrm{mat}\,\allowbreak{} x.\rho)).\mathrm{map}(j),
 \qquad x'.\psi_{v,\allowbreak{} i}=\mathrm{pushCh}(j,\allowbreak{} x.\psi_{v,\allowbreak{} i}),\\
 x'.N3,\qquad x'.N4(F.D_0),\qquad x'.a=n.
 \end{gathered}                                      \tag{CIP.13}
\]
These are the complete conclusions of the original \texttt{Norm}
contract when its point hypothesis is \texttt{Nice\allowbreak I}.
\begin{proof}
Take $k',\allowbreak{} j,\allowbreak{} T$ from CIP6 and the representation and characters of
(CIP.8).  Give $k'$ the composite
$\mathbb Z_\ell\to O\to x.K\to k'$ algebra structure.
The coefficient field is finite.  Every entry of $\rho'$ is a
finite $A'$-linear combination of coefficient images of entries
of $\rho$ with fixed integral coefficients from $T,\allowbreak{} T^{-1}$.
The map $j$ preserves the $t$-adic filtration.  Reducing modulo
$t^m$ therefore proves the exact entrywise \texttt{Adic\allowbreak Cont}
required for every $m$.

Applying $j$ to (CIP.6) and conjugating by $T$ gives
\[
 \rho'(g)^t\rho'(cgc^{-1})=\epsilon(g)^{1-n}1.
\]
Indeed $T^t j(P)T=1$ implies
$(T^{-1})^tT^{-1}=j(P)$; substituting this in the left side
gives the displayed identity.  CIP6 gives the exact
$\mathbb Z_\ell$-coefficient convention.  Thus \texttt{form}
has the literal matrix one.  Set $a=n$; $1\le a\le n$
follows from $n>0$, and CIP7 gives \texttt{blk} and \texttt{span}.

If $g\in I_u\le r.I_p(u)$, the original point has $\rho(g)=1$.
Coefficient extension and conjugation preserve this equality,
giving \texttt{ine} for the exact new tame subgroup.
CIP8 provides \texttt{psi\_\allowbreak mu}, \texttt{w}, \texttt{ty},
\texttt{w\_\allowbreak blk}, \texttt{tri}, and both \texttt{N1} clauses.
This has checked every field in the original \texttt{NP} record.
The changes $D_0=I_0=\bot$ inside \texttt{re\allowbreak Iu} do not change
$\Gamma,\allowbreak{} G_0,\allowbreak{} G',\epsilon,\allowbreak{} G_v,\phi_v$ or $\mu_v$.
CIP9 gives (CIP.11), \texttt{N3}, and the separately requested
\texttt{N4} for the old $D_0$.  The character identities are the
definition (CIP.8).  CIP6 gives injectivity, locality and finiteness
of the one map $j$.  Finally $a=n$ holds outright, and therefore
also under the last implication $x.\mathrm{TypeI}\Rightarrow a=n$
printed in \texttt{Norm}.

This constructs the required object; it does not invoke the source
theorem named \texttt{norm}.  It proves the universal restriction
with the original \texttt{Nice\allowbreak I} hypothesis.  The disjunctive
\texttt{Nice\allowbreak I}\,$\vee$\,\texttt{Nice\allowbreak E} statement as a whole is
not needed for the following consumer and is not asserted here.
\end{proof}
% END CIP10

% BEGIN CIP11
\paragraph{CIP11: the original h=0 consumer keeps its point kernel and ideal.}
\label{claim:CIP11}
CIP10 supplies the entire \texttt{JSP\_\allowbreak SINGLE\_\allowbreak POINT\_\allowbreak NORM}
obligation in JSP11.  On the original rank-four route it has
$x'.a=4$, is valid at $\ell=3$, and introduces no change to
the selected ideal $J$ or to the kernel of an original point map.
\begin{proof}
Fix an arbitrary point $x$ and map $f:B\to x.A$ in the literal
\texttt{P1core} quantifier.  By definition $x$ is \texttt{Nice\allowbreak I}.
The $h=0$ instance of the original support chain uses the singleton
family $\mathrm{Fin}(1)$ with this very point; it never calls
\texttt{Norm} at a \texttt{Nice\allowbreak E} point.  The exact subgroup
identification in JSP11 is
\[
 I_u=(\mathrm{globOf}(r,\allowbreak{} V).\mathrm{frame}).I_u
                         =r.I_p(u).
\]
Thus CIP10 applies, with the original frame and its original $D_0$.
It returns precisely the normalized point required before the
remaining one-point \texttt{Choice} call, including both
nondegeneracy fields, all flags and all characters.

Let $O$ act on $k'$ through $O\to x.K\to k'$.
Coefficientwise construction in CIP6 gives
$j\circ(O\to A)=O\to A'$ and
\[
 \operatorname{const}_{k'}\circ j
              =(x.K\hookrightarrow k')\circ\operatorname{const}_{x.K}.
 \tag{CIP.14}
\]
Consequently $jf$ is an $O$-algebra map and retains the actual
residual determinant and ordered residual characters after this
one constant-field embedding.  Equation (CIP.11) is the full
determinant matching required by the original pair.  Injectivity
gives
\[
 \ker(jf)=\ker f,
 \qquad J\subseteq\ker f\ \Longleftrightarrow\ J\subseteq\ker(jf).
 \tag{CIP.15}
\]
No new ideal has been chosen.  A further injective extension supplied
by \texttt{Choice} preserves this equality and the displayed laws
by JSP10.  The original $N4$ transport used there applies to the
$N4$ just proved in CIP9.

The field extension and Gram--Schmidt argument only require that
two is invertible.  The full matrix span and determinant arguments
work in every characteristic.  Thus all conclusions apply to the
original $n=4,\ell=3$ case.  The construction depends on the point
$x$, as the source existential allows; it neither enlarges the
external coefficient object $c'$ nor alters the already selected
$B,\Lambda,\allowbreak{} J$.
\end{proof}
% END CIP11

% BEGIN CIP12
\paragraph{CIP12: the reduced joint-support frontier and its quantifiers.}
\label{claim:CIP12}
In the sufficient route of JSP12, replace the retained exact
reciprocity matching by CIP1--CIP5, and the retained one-point
\texttt{Norm} input by CIP6--CIP11.  All other producer inputs
remain on their original objects.  This reduction applies to both
actual rank-four Links, with the original fixed-$\iota$ and common
coefficient-threshold order.
\begin{proof}
First fix the original tower and the identification
$\iota:\overline{\mathbb Q}_\ell\simeq\mathbb C$.
The selected joint-support route must then provide one coefficient
threshold, before all later $c'$, original Spaces, and their
universal data.  At each allowed $c'$ it uses both Links with the
same $\kappa$ and their distinct auxiliary primes.  CIP1--CIP5
apply to every original $s$ at its actual $\ell$-adic places,
independently of $\iota$ and without enlarging $c'$.
CIP6--CIP11 apply to every later \texttt{Nice\allowbreak I} point over every
actual $B$, using only a pointwise finite constant extension already
allowed by the original \texttt{Norm} conclusion.  These operations
do not exchange any quantifiers or add a threshold.

On each Link the still-required ordinary system produces its own
specific $R,\allowbreak{} H,\allowbreak{} M,\mathbb T$ and the one ideal
$J=\ker(B\to\mathbb T)=\operatorname{Ann}_B(M)$ of JSP5.
JSP6 and JSP7 then give \texttt{Ring\allowbreak Props} and \texttt{Dense},
with the reciprocity antecedent now supplied.  Its four further
Hida endpoints must still hold on that exact $J$:
\texttt{Point\allowbreak Theorem\allowbreak Base\allowbreak Object},
\texttt{arithmetic\allowbreak Point\allowbreak Purity\allowbreak Dense},
\texttt{arithmetic\allowbreak Point\allowbreak Classicality\allowbreak Dense}, and \texttt{Twist\allowbreak Ok}.
The original all-level \texttt{Has\allowbreak Aut}(B,\allowbreak{} J) is also still needed.
JSP8 supplies only its already reconstructed base, not its higher
level modules or the common-exponent identities.

Of JSP11's six support inputs, the required first input is now
proved.  The remaining five are the literal one-point
\texttt{Choice}, the actual \texttt{Ar\allowbreak Pl}, the exact first-order
bound, the exact obstruction bound, and the one-point \texttt{Raw}
support theorem.  The original \texttt{Coh\allowbreak P}, \texttt{Loc\allowbreak CFT},
coefficient, level and nontrivial residual-cyclotomic side
conditions retain their exact scopes in JSP11.  Supplying these
remaining inputs yields \texttt{P1core} on the already selected
$J$ by the same proof, with CIP11 at its normalization step.
This is a sufficient conditional implication, not an assertion
that those producers have been supplied.

The two ideals belong to the two different original universal
rings; within each ring all Hida and support properties refer to
the same ideal.  Once the remaining data satisfy the common
threshold, JSP12, QX10 and YX2 still assemble the original CountII
conclusion.  For RH and two-primary target characters only the
separate CaseIII branch is vacuous by KA5--KA6; CountII, as well as
Cases I and V, remains substantive.  Thus the present closures
contribute to both RH and unrestricted GRH.  They do not resolve
the extra CaseIII arithmetic work needed for general characters.
\end{proof}
% END CIP12


\section{The original simultaneous loss estimate}
\label{sec:global-loss-v035}

\status{LFS1--LFS17 prove the literal
$\texttt{Loss\allowbreak FS}(\mathcal F,0)$ from the same original
\texttt{Chosen\allowbreak TT}, \texttt{Chosen\allowbreak Badp}, \texttt{Link\allowbreak G}, and
\texttt{Chosen\allowbreak WS}.  The proof includes both the prime-logarithmic
\texttt{IOsum} bound and the universally read \texttt{ASK}-minus-
\texttt{Ga} bound.  It preserves the chosen integer functions and
their realizing member.  Existence of the compatible family and its
local gain estimates remains an upstream input.}

Throughout the final application, $\mathcal F:\texttt{Fam\allowbreak P}(\epsilon)$
is the original supplied family and $0<\epsilon<1/52$.  Write
$D=\mathcal F.D(x,\allowbreak{} i)$ and $L=\mathcal F.L(x,\allowbreak{} i)$ only after fixing the
common threshold.  To keep the source's two small parameters distinct,
write
\[
 z=x^\epsilon,\quad u=x^{\epsilon/3},\quad
 b=e^z,\quad a=xb,\quad X=e^u=X',\quad
 b'=xb/X,\quad w''=b/(nX),\quad X''=xnX.                 \tag{LFS.1}
\]
The relative top-window width is denoted $h$; the exceptional-share
cutoff is $\theta=x^{-\epsilon}$.  Thus $h$ is the source field
\texttt{Par.\allowbreak th}, whereas $\theta$ is \texttt{Par.\allowbreak theta}.
All finite sums below are over the displayed sets; repeated indices
never denote independently chosen arithmetic data.

% BEGIN LFS1
\paragraph{LFS1: elementary prime-logarithmic estimates with their constants.}
\label{claim:LFS1}
For every real $Y\ge1$ and every finite set of primes $S\subset[1,\allowbreak{} Y]$,
\[
 \sum_{p\in S}\log p\le (\log4)Y,\qquad
 \sum_{p\in S}\frac{\log p}{p}\le\log Y+\log4.          \tag{LFS.2}
\]
Here is a proof of the exact constant needed later.  Let
$P(N)=\prod_{p\le N}p$.  Strong induction gives $P(N)\le4^N$.
The cases $N\le2$ are direct.  If $N=2r\ge4$, $N$ is composite and
$P(N)=P(N-1)$.  If $N=2r+1\ge3$, the product of the primes in
$(r+1,2r+1]$ divides $\binom{2r+1}{r}$: each occurs once in the
numerator factorial and not in either denominator factorial.
The two equal central binomial coefficients are among terms summing
to $2^{2r+1}$, so $\binom{2r+1}{r}\le4^r$.  Hence
$P(N)\le P(r+1)4^r\le4^{2r+1}$.  Taking logarithms and then
$N=\lfloor Y\rfloor$ proves the first estimate.

For the second, each $p^{\lfloor N/p\rfloor}$ divides $N!$, and
these factors for different primes are coprime.  Consequently
\[
 \sum_{p\le N}\lfloor N/p\rfloor\log p\le\log(N!)
       \le N\log N.
\]
Multiply the desired prime sum by $N$, use
$N/p\le\lfloor N/p\rfloor+1$, and apply the first estimate.
Division by $N>0$ proves (LFS.2) at integers; restriction to
$\lfloor Y\rfloor$ gives the real version.

We will also use, for an integer $p_0\ge3$ and any finite
$S\subset\{p_0,\allowbreak{} p_0+1,\ldots\}$,
\[
 \sum_{n\in S}\frac{\log n}{n^2}
 \le \frac{2(\log(p_0-1)+1)}{p_0-1}.                    \tag{LFS.3}
\]
Indeed $f(t)=\log t/t^2$ is decreasing for $t\ge2$, because
$f'(t)=(1-2\log t)/t^3\le0$; the inequality
$\log2\ge1/2$ follows by integrating $1/t\ge1/2$ on $[1,2]$.
Thus the full tail is at most
$\int_{p_0-1}^\infty f(t)\,\allowbreak{} dt=(\log(p_0-1)+1)/(p_0-1)$,
which is stronger than (LFS.3).  These arguments establish all three
estimates without a prime number theorem or a distribution theorem
for residue classes.
% END LFS1

% BEGIN LFS2
\paragraph{LFS2: the dyadic small-order prime bound.}
\label{claim:LFS2}
Let $e\ge1$, $J\ge1$, and $Y\ge1$.  For an odd prime $p$, let
$\bar p$ be its class in
$(\mathbb Z/2^{e+2}\mathbb Z)^\times/\{\pm1\}$.
Then every finite set $S$ of such primes with $p\le Y$ and
$\operatorname{ord}(\bar p)\le J$ satisfies
\[
 \sum_{p\in S}\frac{\log p}{p}
 \le\frac{12J\log^2(2^eY)}{2^e}.                       \tag{LFS.4}
\]
We give the residue calculation, rather than treating the order
condition as a prime-density assertion.  If $v$ is odd, one of
$v-1,\allowbreak{} v+1$ has 2-adic valuation exactly one.  Factoring
$v^2-1=(v-1)(v+1)$ therefore gives, for $r\ge1$,
\[
 2^{r+1}\mid v^2-1
 \quad\Longleftrightarrow\quad
 2^r\mid v-1\ \hbox{or}\ 2^r\mid v+1.
\]
Also $v^2+1\equiv2\pmod4$.  Induction on $j$, using these two
facts at each squaring, gives for $0\le j\le e$
\[
 v^{2^j}\equiv\pm1\pmod{2^{e+2}}
 \quad\Longleftrightarrow\quad
 v\equiv\pm1\pmod{2^{e+2-j}}.                          \tag{LFS.5}
\]
In particular the quotient group has exponent dividing $2^e$.
An element of order at most $J$ thus has order $2^r$, where
$r\le j_0:=\min(e,\lfloor\log_2J\rfloor)$.  Equation (LFS.5)
puts $p$ in one of the two classes $\pm1$ modulo
$M=2^{e+2-j_0}$.  We have $M\ge4$ and $2^e\le JM$.

Every integer $p\ge3$ in these two classes is $qM+1$ or $qM-1$
with $q\ge1$.  If $p\le Y$, then $q\le\lfloor Y\rfloor$ and
$p\ge qM/2$.  Extending the nonnegative sum along the two
progressions and using
$\sum_{q=1}^Nq^{-1}\le1+\log N$ proves
\[
 \sum_{p\in S}\frac{\log p}{p}
 \le\frac{4\log Y(1+\log Y)}M.
\]
Put $\lambda=\log(2^eY)$.  Since $e\log2\ge1/2$,
$\log Y(1+\log Y)\le3\lambda^2$.
The inequality $M^{-1}\le J/2^e$ now proves (LFS.4).
Thus the source's universal order constant may be taken to be
$C_{\rm ord}=12$.
% END LFS2

% BEGIN LFS3
\paragraph{LFS3: first-kind primes, including the narrow power bands.}
\label{claim:LFS3}
Let $X\ge2$, $0\le\rho\le1/2$, and $p_0\ge8$.
If a finite set of integers lies in
$[P(1-\rho),\allowbreak{} P(1+\rho)]\cap[p_0,\infty)$, with $1\le P\le X$,
then its cardinality is at most $2\rho P+1$.  Each reciprocal is
at most $\min(2/P,1/p_0)$ and each logarithm at most $\log(2P)$.
It follows that
\[
 \sum \frac{\log p}{p}
 \le (4\rho+1/p_0)\log(2X).                            \tag{LFS.6}
\]
For $p\ge p_0$ and $i\ge1$, the condition
$X(1-\rho)\le p^i\le X(1+\rho)$ implies
\[
 X^{1/i}(1-\rho)\le p\le X^{1/i}(1+\rho).
\]
Moreover $i\log8\le\log(2X)\le2\log X$, and
$\log8\ge2$.  One proof of the last inequality is
$\log2=2\int_0^{1/3}(1-t^2)^{-1}dt\ge2/3$.
Therefore $i\le\lfloor\log X\rfloor$.
Sum (LFS.6) over those exponents.  Overlaps can only increase this
upper bound.

Combining with LFS2 proves the following exact first-class bound.
If each prime of $S$ either has dyadic quotient order at most $J$
and lies below $Y$, or lies in one of the displayed bands, then
\[
 \sum_{p\in S}\frac{\log p}{p}
 \le\frac{12J\log^2(2^eY)}{2^e}
  +\lfloor\log X\rfloor(4\rho+1/p_0)\log(2X).          \tag{LFS.7}
\]
For the source parameters use
$J=\lfloor x^{\epsilon/4}\rfloor$,
$Y=\max(P.Y,\allowbreak{} X'')$, and $X=P.X'$.
An integer order strictly below $x^{\epsilon/4}$ is at most $J$.
Since $\log X'=x^{\epsilon/3}$, the number of bands is exactly
the source's $I_b=\lfloor x^{\epsilon/3}\rfloor$.
% END LFS3

% BEGIN LFS4
\paragraph{LFS4: the finite requested sums and the high-prime boundary.}
\label{claim:LFS4}
For a finite set $I\subset\mathbb N$ of exponents, write
\[
 A(I,\allowbreak{} A_0,\allowbreak{} p,\allowbreak{} w)=\sum_{r\in I}\max(A_0/p^r-w_r,0).
\]
If $p\ge2$, $0\le A_0\le3a/2$, all $r\in I$ are odd and
$w_r\ge0$, then the geometric series over positive odd exponents
gives
\[
 A(I,\allowbreak{} A_0,\allowbreak{} p,\allowbreak{} w)\le A_0\frac{p^{-1}}{1-p^{-2}}
 \le\frac{2a}{p},\qquad
 2\log p\,\allowbreak{} A(I,\allowbreak{} A_0,\allowbreak{} p,\allowbreak{} w)\le4a\frac{\log p}{p}.           \tag{LFS.8}
\]
For $d\ge4$, all positive exponents instead give
$A(I,\allowbreak{} A_0,\allowbreak{} d,\allowbreak{} w)\le A_0/(d-1)\le2a/d$ and hence the
discriminant contribution is at most $2a\log d/d$.

The sharper high-prime estimate uses the actual positive lower
weight $w''$, with $X''w''=a$, $X''\ge2$, and
$A_0\le a(1+h)$, $0\le h\le1$.
For $p>X''$, every term with exponent at least two is zero:
$p^2\ge2X''$ and $A_0/p^r\le w''$.  The exponent-one term
also vanishes for $p\ge X''(1+h)$; in the remaining interval it
is at most $w''h$.  There are at most $X''h+1$ integers there,
and $\log p\le\log(2X'')$.  Therefore
\[
 \sum_{p>X''}2\log p\,\allowbreak{} A(I_p,\allowbreak{} A_0,\allowbreak{} p,\allowbreak{} w_p)
 \le2\log(2X'')\bigl(ah^2+w''h\bigr).                 \tag{LFS.9}
\]
The endpoint allowance $+1$ is retained in the $w''h$ term.
All these estimates apply to finite subsets without any assumption
that every exponent or every prime is present.
% END LFS4

% BEGIN LFS5
\paragraph{LFS5: exact deficit splitting into four prime classes.}
\label{claim:LFS5}
Fix finite sets of lifts $J$, tops $T$, and primes $\mathcal P$;
put $m=|J|$ and $q=|T|$.  For each prime write $A_{j,\allowbreak{} t}(p)$ for
the finite requested sum in LFS4.  Let $g_j(p)$ be its retained
gain, $\tau(p)$ its correction, and $T_p\subset T$ its retained
tops.  The finite-place unit rule says that all places over $p$
have these same quantities and their normalized weights sum to
$\log p$, with zero contribution off $\mathcal P$.  Hence
\[
 \mathrm{ASK}-\mathrm{Ga}
 =\sum_{p\in\mathcal P}2\log p\left[
       \sum_{j,\allowbreak{} t}A_{j,\allowbreak{} t}(p)-\left(\sum_jg_j(p)-\tau(p)\right)
                                 \right]+\sum_{j,\allowbreak{} t}D_{j,\allowbreak{} t}.    \tag{LFS.10}
\]
This follows by interchanging finite sums; $D_{j,\allowbreak{} t}$ denotes the
discriminant term and the factor $2$ occurs on both sides.

At a served prime suppose the actual ledger inequality is
\[
 \sum_j\sum_{t\in T_p}(A_{j,\allowbreak{} t}(p)-\ell(p))-mf(p)
        \le\sum_jg_j(p)-\tau(p).
\]
Subtract and split $T=T_p\sqcup(T\setminus T_p)$.  The bracket
in (LFS.10) is at most
\[
 \sum_j\sum_{t\in T\setminus T_p}A_{j,\allowbreak{} t}(p)
                  +m|T_p|\ell(p)+mf(p).                       \tag{LFS.11}
\]
At an unserved prime the stipulated gain and correction are zero,
so the bracket is the entire requested sum.  Partition
$\mathcal P$ into served, first, second, and high primes, multiply
these inequalities by $2\log p\ge0$, and add them.  This proves
the source's deficit splitting with no positivity assumption on
the net gain and no change in the sets $T_p$.
% END LFS5

% BEGIN LFS6
\paragraph{LFS6: the four served-prime error bounds.}
\label{claim:LFS6}
Assume $C,\allowbreak{} f,\mathcal S,\allowbreak{} Q_1,\allowbreak{} Q_2,\Lambda\ge0$, $a\ge1$,
$k,\allowbreak{} K_0>0$, $p_0\ge3$, and $x,\allowbreak{} X,\allowbreak{} X''\ge1$.
The original envelopes are
\begin{align*}
 F(p)&=CfN^2\log_p a+Cf(\log_pa)^2,\\
 C_p&=C\mathcal S\log_p(2a),\\
 P_p&=\begin{cases}Q_1&p<p_0,\\Q_2&p\ge p_0,\end{cases}\\
 O_p&=C\log x\left(\frac{\Lambda a}{pk}
  +\begin{cases}ap^{-L_1(p)}&p<p_0,\\a/p^2&p\ge p_0.\end{cases}\right).
\end{align*}
Suppose served primes are at most $X''$, the $P_p,\allowbreak{} O_p$ terms
occur only for $p<X$, and $K_0\le p^{L_1(p)-1}$ for $p<p_0$.
Then their sums with the weight $2\log p$ satisfy
\begin{align*}
 \sum 2\log p\,\allowbreak{} F(p)
 &\le2CfX''\bigl(N^2\log a+2(\log a)^2\bigr),\\
 \sum 2\log p\,\allowbreak{} mqC_p
 &\le2C\mathcal S\log(2a)\,\allowbreak{} mqX'',\\
 \sum_{p<X}2\log p\,\allowbreak{} P_p
 &\le2\log4\,(p_0Q_1+XQ_2),\\
 \sum_{p<X}2\log p\,\allowbreak{} O_p
 &\le2C\log x\,\allowbreak{} a\left[
    (\Lambda/k+1/K_0)(\log X+\log4)
    +\frac{2(\log(p_0-1)+1)}{p_0-1}\right].             \tag{LFS.12}
\end{align*}
For the first line cancel $\log p$ in $\log_pa$, and use
$1/\log p\le2$ in the square term.  The number of served
primes is at most $\lfloor X''\rfloor\le X''$.
The second line is exact cancellation of $\log p$ followed by
the same cardinality bound.  For the third line, bound the low
prime sum by $\log4\,\allowbreak{} p_0$ and the whole sum below $X$ by
$\log4\,\allowbreak{} X$, using LFS1; counting the high contribution also
at low primes only increases this upper bound.
Finally $p^{-L_1(p)}\le1/(K_0p)$ below $p_0$.
Apply the reciprocal-prime bound in LFS1 to both terms with
$1/p$, and its square tail bound to the terms with $1/p^2$.
These steps prove every line of (LFS.12) with the stated constants.
% END LFS6

% BEGIN LFS7
\paragraph{LFS7: the complete eleven-term loss decomposition.}
\label{claim:LFS7}
Use the exact $\texttt{Par}$ quantities
\begin{equation}
\begin{aligned}
 f&=2n,\qquad\mathcal S=S_c+k+\eta,\qquad
 Q_1=Cr_Gb'\log x,\\
 Q_2&=C\left(b'/x+\eta+S_c+(w''+t)\log x/\delta_1\right)+b_o,
 \qquad\delta_1=x^{-\epsilon/8}.
\end{aligned}                                             \tag{LFS.13}
\end{equation}
Assume the eleven clauses of the finite ledger:
the unit rule in LFS5; $0\le a_t\le a(1+h)$;
odd prime exponents with weights at least $w''$;
positive discriminant exponents with nonnegative weights;
served primes at most $X''$; high primes greater than $X''$;
the served ledger in LFS5; zero unserved gains and corrections;
$f(p)\le F(p)$; $\ell(p)\le C_p+\mathbf1_{p<X}(P_p+O_p)$;
and $K_0\le p^{L_1(p)-1}$ below $p_0$.
Assume also the first-prime alternative of LFS3.

Take $h\le1/2$, so $a_t\le3a/2$.  LFS5, LFS4, and LFS6
give the following terms.  The first-class contribution is split
by LFS3 into its order and band parts.  At a served prime, the
envelope of $\ell(p)$ is nonnegative, so
$m|T_p|\ell(p)\le mq(C_p+\mathbf1_{p<X}(P_p+O_p))$ even
if $\ell(p)$ is negative.  Thus
\begin{equation}
\begin{aligned}
 \mathrm{ASK}-\mathrm{Ga}
 &\le E_F+E_C+E_{P,<}+E_{P,\ge}+E_O+E_{\rm out}\\
 &\qquad{}+E_{\rm ord}+E_{\rm band}+E_H+E_d+E_{\rm second}.
\end{aligned}                                               \tag{LFS.14}
\end{equation}
where the eleven terms, in order, are
\begin{align*}
 E_F&=m\,2CfX''(N^2\log a+2\log^2a),\\
 E_C&=2C\mathcal S\log(2a)mqX'',\\
 E_{P,<}&=mq\,2\log4\,\allowbreak{} p_0Q_1,\\
 E_{P,\ge}&=mq\,2\log4\,\allowbreak{} XQ_2,\\
 E_O&=mq\,2C\log x\,\allowbreak{} a
 \left[(\Lambda/k+1/K_0)(\log X+\log4)
       +\frac{2(\log(p_0-1)+1)}{p_0-1}\right],\\
 E_{\rm out}&=\sum_{p\ {\rm served}}4ma\frac{\log p}{p}|T\setminus T_p|,\\
 E_{\rm ord}&=4mqa\frac{C_{\rm ord}J_o
                       \log^2(2^{e_2}\max(Y,\allowbreak{} X''))}{2^{e_2}},\\
 E_{\rm band}&=4mqaI_b(4\rho+1/p_0)\log(2X),\\
 E_H&=mq\,2\log(2X'')(ah^2+w''h),\\
 E_d&=mq\,2a\log d/d,\\
 E_{\rm second}&=4mqa\sum_{p\ {\rm second}}\frac{\log p}{p}.
\end{align*}
Here $J_o=\lfloor x^{\epsilon/4}\rfloor$,
$I_b=\lfloor x^{\epsilon/3}\rfloor$, and $C_{\rm ord}=12$.
The terms correspond literally to
\texttt{t\allowbreak Floors}, \texttt{t\allowbreak Cp}, \texttt{t\allowbreak Pr\allowbreak Low},
\texttt{t\allowbreak Pr\allowbreak High}, \texttt{t\allowbreak Or}, the omitted-top sum,
\texttt{t\allowbreak Orders}, \texttt{t\allowbreak Band}, \texttt{t\allowbreak High},
\texttt{t\allowbreak D}, and the second-class sum.  No first-class or
discriminant term has been silently absorbed into a different
arithmetic object.
% END LFS7

% BEGIN LFS8
\paragraph{LFS8: the same small-share inequality pays three errors.}
\label{claim:LFS8}
Let $\mathcal C$ be the candidate set and $\Sigma_p\ge0$.
Suppose the original chosen-weight inequality is
\[
 W:=\sum_{p\in\mathcal C}\frac{\log p}{p}
    \left(2\mathbf1_{\{\Sigma_p>\theta\}}+4\Sigma_p\right)
       \le x^{-9}.                                      \tag{LFS.15}
\]
Suppose every second-class prime belongs to $\mathcal C$ and has
$\Sigma_p>\theta$.  Then
\[
 E_{\rm second}\le2amqW\le2amq x^{-9}.                  \tag{LFS.16}
\]
Indeed the indicator part of $2amqW$ is exactly the required
$4amq\log p/p$, and its remaining summands are nonnegative.
If, for each candidate, the original error satisfies
$io(p)\le2\Sigma_p(2mqa/p)$, then the $4\Sigma_p$ part of
$amqW$ proves simultaneously
\[
 \sum_{p\in\mathcal C}\log p\,\allowbreak{} io(p)
       \le amqW\le amq x^{-9}.                          \tag{LFS.17}
\]

Let $\mathcal R$ be the original large-order residue-prime set.
Suppose an omitted served top is bad at that prime, and the prime
lies in $\mathcal R$.  Inclusion of these top sets and nonnegativity
of their weights show
\[
 E_{\rm out}\le4am\sum_{p\in\mathcal R}\frac{\log p}{p}
                         |\{t\in T:\mathrm{bad}(p,\allowbreak{} t)\}|.
\]
The original chosen-top inequality bounds this by
\[
                      E_{\rm out}\le4amq x^{-2+\epsilon}.
                                                               \tag{LFS.18}
\]
If a served prime is absent from $\mathcal R$, the omission
condition says $T\setminus T_p$ is empty, so this extension of
the sum is valid also at those primes.  Equations
(LFS.16)--(LFS.18) use precisely the same $\Sigma$, bad predicate,
and degree functions; a second independent choice is unnecessary.
% END LFS8

% BEGIN LFS9
\paragraph{LFS9: uniform absorption with the original leading constant and thirteen errors.}
\label{claim:LFS9}
Fix $0<\epsilon<1/4$, $e\ge0$, $K\ge1$ and
$C_{\rm ord}=12\le K$.  Assume the original
\texttt{Windows\allowbreak D} inequalities:
\begin{gather*}
 x^2\le2^{e_2}\le n\le x^e,\qquad x^2\le k\le x^e,
 \qquad 0\le N,\allowbreak{} t,\allowbreak{} r_G\le x^e,\\
 x^4+1\le p_0\le x^e,\quad0\le\Lambda\le x^{2\epsilon},
 \quad0\le S_c,\eta\le e^{x^{\epsilon/4}},\\
 0\le\rho\le x^{-2},\quad0\le h\le x^{-1},\quad
 x^2\le K_0,\allowbreak{} d,\quad0\le C\le K,\quad
 1\le Y\le2X,\quad0\le b_o\le\sqrt b.                 \tag{LFS.19}
\end{gather*}
There is one threshold depending only on $\epsilon,\allowbreak{} e,\allowbreak{} K$ such that
every such parameter record and every finite ledger in LFS7--LFS8
satisfy, with $U=qmb$ and $\sigma=100Kx^{4\epsilon}/x$,
\[
 \mathrm{ASK}-\mathrm{Ga}
       \le(2\log4\,\allowbreak{} C+13\sigma)U,
 \qquad \sum_{p\in\mathcal C}\log p\,\allowbreak{} io(p)\le\sigma U,
 \qquad0\le\sigma\le1.                                \tag{LFS.20}
\]
Here $q\ge1$ and $m\ge0$.  We prove the uniform size step
explicitly; we do not infer it from the later, different
\texttt{Win\_\allowbreak at} predicate.

Put $E=x^e$, $S=e^{x^{\epsilon/4}}$, $A=\log x+z$, and
\[
 L_1=\log(2xEX),\qquad L_2=\log(2xE^2X).
\]
For sufficiently large $x$, substitution of (LFS.19) into the
eleven displayed terms gives the following upper coefficients
after factoring out $U$.  For $E_{P,\ge}$ the listed expression
is the coefficient after subtracting $2\log4\,\allowbreak{} C$.
\begin{center}\small
\begin{tabular}{p{0.20\textwidth}p{0.72\textwidth}}
Term & Upper coefficient\\\hline
$E_F$ & $4KxE^2(X/b)(E^2A+2A^2)$\\
$E_C$ & $2K(2S+E)(A+\log2)xE(X/b)$\\
$E_{P,<}$ & $2\log4\,\allowbreak{} KE^2x\log x/X$\\
High precision
 & $\begin{aligned}
 &2\log4\{K[2XS/b\\
 &\quad{}+(x^{-2}+XE/b)x^{\epsilon/8}\log x]+X/\sqrt b\}
 \end{aligned}$\\
$E_O$ & $\begin{aligned}
 &2Kx\log x[(z^2+1)x^{-2}(u+\log4)\\
 &\qquad{}+2(e\log x+1)x^{-4}]
 \end{aligned}$\\
$E_{\rm out}$ & $4x^{-1+\epsilon}$\\
$E_{\rm ord}$ & $4Kx^{-1+\epsilon/4}L_2^2$\\
$E_{\rm band}$ & $4xu(4x^{-2}+x^{-4})(u+\log2)$\\
$E_H$ & $2L_1(x^{-1}+x^{-3}/X)$\\
$E_d$ & $4\log x/x$\\
$E_{\rm second}$ & $2x^{-8}$\\
Prime-logarithmic $io$ sum & $x^{-8}$
\end{tabular}
\end{center}
For the floor term $1/q\le1$ was used.  The high-precision
term is worth checking separately: from $b'=xb/X$ and
$w''=b/(nX)$,
\[
 \frac{XQ_2}{b}
 =C\left[1+\frac{X(\eta+S_c)}b
       +\left(\frac1n+\frac{Xt}b\right)
                         x^{\epsilon/8}\log x\right]
        +\frac{Xb_o}b.                                    \tag{LFS.21}
\]
This is the unique leading term in the table.
For $E_O$ use $p_0-1\ge x^4$ and
$\log(p_0-1)\le e\log x$.  For the order term use
$2^{e_2}\max(Y,\allowbreak{} X'')\le2xE^2X$.  For the discriminant term,
$\log t/t$ is decreasing for $t\ge e^1$ and $d\ge x^2$,
so $2x\log d/d\le4\log x/x$; no upper bound on $d$ is
assumed.  These justify every row of the table.

Every coefficient in the table is $o(\sigma)$, uniformly in
all the parameter records satisfying (LFS.19).
For completeness the only non-exponential comparisons needed
are as follows.  Since $L_1,\allowbreak{} L_2=O_{e,\epsilon}(u)$, the ratios
to $\sigma$ of the order, band, and high rows are respectively
\[
 O(x^{-37\epsilon/12}),\qquad
 O(x^{-10\epsilon/3}),\qquad O(x^{-11\epsilon/3}).
\]
The principal part of the orbit row has ratio
$O(\log x\,\allowbreak{} x^{-5\epsilon/3})$; its other part is smaller.
The omitted-top row has ratio $O(x^{-3\epsilon})$.
The high-precision remainder without an exponential has ratio
$O(\log x\,\allowbreak{} x^{-1-31\epsilon/8})$.
The discriminant, second-class and $io$ rows have ratios
$O(\log x/x^{4\epsilon})$ and $O(x^{-7-4\epsilon})$.
All remaining ratios are polynomial-logarithmic factors times
one of $X/b$, $XS/b$, $X/\sqrt b$, or $1/X$, and tend to zero:
their logarithms tend to $-\infty$ because
$x^{\epsilon/3}+x^{\epsilon/4}=o(x^\epsilon)$ and
$\log x=o(x^r)$ for $r>0$.
These last limits follow, for example, by integrating
$1/t\le t^{r/2-1}$ and then dividing by $x^r$.

Choose a common threshold making each of the ten non-leading
loss terms at most $\sigma U$, the high-precision remainder at
most $3\sigma U$, and the $io$ sum at most $\sigma U$.
There are ten terms besides $E_{P,\ge}$ in (LFS.14), giving
exactly $10+3=13$.  Enlarge the threshold until
$100Kx^{4\epsilon-1}\le1$.  This proves (LFS.20), including
$m=0$ by multiplication of the coefficient inequalities rather
than division by $U$.

If desired the same threshold can be enlarged to satisfy all
three quantitative \texttt{Large\allowbreak X} comparisons:
$100Kx^{4\epsilon}\le x$,
$100Kx^{4e+4}\le e^u$, and
$100Kx^{4e+4}e^{u+x^{\epsilon/4}}\le e^{z/2}$.
The proof uses a possibly larger cutoff than the source's
particular sufficient one, while retaining its exact
\texttt{joint\_\allowbreak core} constants and quantifier order.
% END LFS9

% BEGIN LFS10
\paragraph{LFS10: every original loss-window clause, uniformly on the family.}
\label{claim:LFS10}
Fix the constant $C$ and function $\delta(x)\to0$ supplied by
\texttt{Link\allowbreak G}.  At each member put $b_o=b^{\delta(x)}$ and take
the literal parameter record \texttt{Loss.\allowbreak par}, with
$C_b=0$ and $C_{\rm ord}=12$.
The complete coordinate identification is
\[
\begin{array}{c|c@{\qquad}c|c}
 P.x&x&P.\epsilon&\epsilon\\
 P.e_2&D.R.e&P.n& D.R.n\\
 P.k&D.R.k&P.N& k+m(t+1)\\
 P.t&D.R.t&P.p_0&\lceil k^2x^{26}\rceil_+\\
 P.\Lambda&L.\mathrm{LaG}&P.r_G&r_G(\epsilon)\\
 P.S_c&L.\mathrm{Sc}=4kD.s.\mathrm{Del}&P.\eta&L.\mathrm{etaSz}\\
 P.\rho&G.\mathrm{rhoB}(D,\mathcal F.GK,\allowbreak{} L.sc)&P.h&1/(kx^{13})\\
 P.K_0&k^2x^{78}&P.d&D.host.d\\
 P.C&\max(C,0)&P.C_{\rm ord}&12\\
 P.Y&X'(1+\mathcal F.C_r\rho)&P.b_o&b^{\delta(x)}\\
 P.\beta&0&P.s_0&m|TT|x^{-12}.
\end{array}                                                   \tag{LFS.22}
\]
Choose
\[
                 e=\max(40,\mathcal F.C_t+3),\qquad
                 K=1+|C|+12.                                \tag{LFS.23}
\]
One threshold, independent of $i$, makes the entire original
\texttt{Windows}(P,\allowbreak{} e,\allowbreak{} K) true, and makes $k^2x^{26}\le X'$.
In particular its projection \texttt{Windows\allowbreak D} is true.

Here are all the checks.  The original sizes in SZ1--SZ8 give,
for $x\ge3$,
\[
 n\le x^3,\quad m\le x,\quad x^4\le k\le x^5,
 \quad t\le x^{\mathcal F.C_t+\epsilon/100}
                   \le x^{e-2}.                              \tag{LFS.24}
\]
The source also has the sharper stored $t$-window; the weaker
bound displayed here suffices.  Its rank frame has
$2^{e_2}\le n\le o2^{e_2}$ and $o\le x^{\epsilon/100}$.
The lower $n$-window is at least $x^{2+\epsilon/100}$;
divide by this bound for $o$ to obtain $x^2\le2^{e_2}$.
Also $N=k+m(t+1)\le3x^{e-1}\le x^e$.
The ceiling in (LFS.22) satisfies
\[
 x^4+1\le p_0\le x^{36}+1\le x^{40}\le x^e.
\]
The lower inequality follows already from $k\ge x^4$ and
$x^{26}\ge2$; the upper one uses $k\le x^5$ and retains the
ceiling's $+1$.

The remaining scalar checks are
\begin{align*}
 0\le\Lambda&\le x^{\epsilon+\epsilon/100}\le x^{2\epsilon},\\
 S_c=4kD.s.\mathrm{Del}
 &\le4x^5e^{x^{\epsilon/4}/2}\le e^{x^{\epsilon/4}},
 &0<\eta&\le e^{x^{\epsilon/4}},\\
 0\le\rho&\le x^{-2},
 &0<h&=1/(kx^{13})\le x^{-1},\\
 x^2&\le k^2x^{78},
 &x^2&\le d.                                                \tag{LFS.25}
\end{align*}
For $S_c$ use \texttt{Size\allowbreak Del}'s
$\log(D.s.\mathrm{Del})\le x^{\epsilon/4}/2$, and absorb
$4x^5$ into the other half exponential.  The $\eta$ inequality
is exactly its logarithmic size clause.  Nonnegativity of
$\rho$ follows by expanding its original expression into
nonnegative terms, using $\mathcal F.GK.\mathrm{Nonneg}$;
its upper bound is the literal \texttt{Fam\allowbreak P.\allowbreak rho\_\allowbreak le} field.
For $d$ use \texttt{Size\allowbreak QD}'s
$\log d\ge x^{19\mathcal F.ez/100}$.
Here $\mathcal F.ez=\epsilon/\mathcal F.J>0$ is fixed before
the threshold, so its exponential eventually exceeds $x^2$.
The integer $r_G(\epsilon)$ is fixed and is at most $x\le x^e$
eventually.

Since $\mathcal F.C_r\ge1$ is fixed, take $x\ge\mathcal F.C_r$.
Then $0\le\mathcal F.C_r\rho\le x^{-1}\le1$, proving
$1\le Y\le2X'$.  Since eventually $\delta(x)\le1/2$,
$0<b^{\delta(x)}\le e^{z/2}$; no sign assumption on
$\delta(x)$ is needed.  Clearly $0\le P.C\le K$,
$0\le P.C_{\rm ord}\le K$, and $0\le P.\beta=0\le Kx^{-2}$.
The top bound gives $|TT|\le2n\delta_1\le2x^3\le x^4$,
so $0\le s_0\le x^{-7}\le x^{-2}$.  These include the two
clauses discarded by \texttt{Windows\allowbreak D}; that projection is
equivalently the full window for the record with
$\beta=s_0=0$.
Finally $k^2x^{26}\le x^{36}\le e^{x^{\epsilon/3}}$ eventually.
Every threshold used here depends only on the original fixed
family constants, $C$, and the eventual bound for $\delta$.
It precedes $x$ and every $i\in\mathcal F.I_x$.
% END LFS10

% BEGIN LFS11
\paragraph{LFS11: exact lift, top, gain, and prime-logarithmic dictionaries.}
\label{claim:LFS11}
Fix one original $D,\allowbreak{} L$ and $X_G:G.\mathrm{Inst}(D,\allowbreak{} L)$.
Let $J=\{1,\ldots,\allowbreak{} m\}$, so $|J|=m$, and let
$T=\mathrm{TTv}(D)=\{\operatorname{val}(t):t\in TT\}$.
The frame's integer-label map is injective and
$\mathrm{ofNat}(\operatorname{val}(t))=t$.  Thus
\[
 |T|=|TT|,\quad
 \sum_{r\in T}F(r)=\sum_{t\in TT}F(\operatorname{val}(t)).     \tag{LFS.26}
\]
The same bijection restricts to the bad subsets defined by
\[
 \mathrm{badAt}(p,\allowbreak{} r)=\neg G.\mathrm{Good}
 (X_G,\allowbreak{} p,\allowbreak{} X_G.bs(G.\mathrm{stepP}(p),\mathrm{ofNat}(r))).
\]
Their cardinalities are exactly \texttt{bad\allowbreak Tops}, not merely
bounded by those cardinalities.

Define the original finite prime universe
\[
 B=1+\max\bigl(\lfloor a+a/(kx^{13})\rfloor_+,
                                  \lfloor X''\rfloor_+\bigr),
 \qquad\mathcal P=\{p<B:p\text{ prime and inert in }K_d\}.    \tag{LFS.27}
\]
The candidate set is literally
$\mathcal C=\texttt{cands}(L,\mathcal F.GK)
 =\texttt{s\allowbreak At}(L,\mathcal F.GK)$: each candidate has $p<X'$,
so the condition $p<\lceil X'\rceil_+$ adds nothing.
The source's rational-prime ledger \texttt{ld\allowbreak At} has
$W=\mathcal P$, $\mathrm{below}=\mathrm{id}$ and weight
$\log p$.  Its unit rule is therefore a singleton-sum identity.
Its gain is
\[
 g_j(p)=\begin{cases}0&p\in L.\mathrm{drop},\\
 G.\mathrm{gainP}(X_G,\mathcal F.GK,\allowbreak{} a_L,\allowbreak{} b_L,\allowbreak{} p,\allowbreak{} j)&p\notin L.\mathrm{drop},
 \end{cases}
\]
and its correction is the original $G.\mathrm{tauP}$.

Let $\ell_w=\log N\mathfrak p_w/[E_0:\mathbb Q]$ for the
original coefficient field $E_0=D.\mathrm{what}_2.LL$, and
$e_w=v_w(p)$ for $w\mid p$.  The weight called
\texttt{Lnk.\allowbreak l\allowbreak E} is $\ell_we_w$.
The exact finite-place degree formula of LC3 gives
\[
                       \sum_{w\mid p}\ell_we_w=\log p.       \tag{LFS.28}
\]
The family support contains the full fiber over each $p\le X''$.
The original \texttt{Parameters\allowbreak Fill} identities say, on that
support, that the gains and corrections are precisely the
rational-prime functions above.  At non-inert primes they
vanish by $L.\mathrm{not\_\allowbreak inert}$.  At $p>X''$, candidates
are absent, so $VV$ is empty, $\tau_P=io_P=0$, and the
explicit high-prime branch of $\mathrm{gainP}$ is zero.
At any further support places the summand is therefore zero.
Grouping finite sums by the full fibers in (LFS.28) proves
\[
 \boxed{\quad
 (\mathrm{ldAt}) .\mathrm{Ga}(J)=L.\mathrm{Ga},\qquad
 \sum_{p\in\mathcal C}\log p\,\allowbreak{} io_P(p)=\mathrm{IOsum}(L).
 \quad}                                                       \tag{LFS.29}
\]
The first identity uses
$\mathrm{sz}(w,\allowbreak{} y)=2\ell_we_wy$, hence retains the factor $2$.
The second uses the same $L.io(w)=io_P(p)$ from
\texttt{Parameters\allowbreak Fill}; no factor $2$ is inserted into it.
Neither identity replaces the support by an infinite sum of
possibly nonzero terms.
% END LFS11

% BEGIN LFS12
\paragraph{LFS12: the actual four prime classes satisfy the order and omission alternatives.}
\label{claim:LFS12}
The source classifies a prime as served if $p\notin L.\mathrm{drop}$;
otherwise as high if $p>X''$; otherwise as second if it is a
candidate with $\Sigma_p>\theta$; otherwise as first.
We check these branches on the original universe (LFS.27).

First no type prime $T_u$ is inert.  The original norm witness
$\pi\in\mathcal O_E$ has $\pi\bar\pi=T_u$ and
$\pi\equiv1\pmod{M\mathcal O_E}$, where
$M=Q_Sqd(g_0g_1)^2$.
Therefore $T_u-1\in M\mathcal O_E\cap\mathbb Z=M\mathbb Z$.
For the last equality, $(T_u-1)/M$ is a rational algebraic
integer, hence an integer.  Every divisor of $M$ is consequently
coprime to $T_u$.  Since $16\mid Q_S$, also $T_u\equiv1\pmod{4d}$.
Quadratic reciprocity gives $(-d/T_u)=1$: the $-1$ sign is $1$
and $T_u\equiv1\pmod d$ makes the other symbol $1$.
The literal host prime classification then says $T_u$ is split.
The prime $2$ is the original ramified-two class, so is not in
$\mathcal P$ either.  This is the same norm/congruence datum
used in RSL8 and the same arithmetic host of AH9--AH12.

The \texttt{Fam\allowbreak P.\allowbreak drop\_\allowbreak iff} clause now simplifies, on
$\mathcal P$, to
\[
             p\in L.\mathrm{drop}\ \Longleftrightarrow\
                         \mathrm{NoGain}(X_G,\mathcal F.GK,\allowbreak{} C_r,\allowbreak{} p),
                                                               \tag{LFS.30}
\]
because its additional exclusions $p\ne T_u$ have just been
verified.  The definition of \texttt{No\allowbreak Gain} lists: a small
order, a narrow power band, a candidate with large share,
a prime above $X''$, a non-inert class, or a nonzero odd-order
Frobenius step in the inert-residue class.
For a first prime the large-share and high-prime branches are
excluded by its definition, and the non-inert branch is
excluded by (LFS.27).  The small-order branch gives the first
alternative in LFS3, with bound $p<Y$.
The power-band branch gives its second alternative; for an
integer prime $p$, the condition $p\ge k^2x^{26}$ is exactly
$p\ge\lceil k^2x^{26}\rceil_+$, including the endpoint.

Here is the exact order dictionary in the other two cases.
The unramified reciprocity calculation of AH9--AH10 identifies
the dyadic projection of the actual Frobenius step $h_p$ with
$p$ in $(\mathbb Z/2^{e_2+2}\mathbb Z)^\times/\{\pm1\}$.
Projection can only divide its order.  If $h_p$ has odd order,
its projection has order dividing both that odd integer and
$2^{e_2}$; hence the projection is zero and has order $1$.
This is strictly smaller than $x^{\epsilon/4}$ for $x>1$.
Such a first prime is at most $X''$, so it again satisfies
LFS3 with bound $\max(Y,\allowbreak{} X'')$.  These arguments supply every
branch of the literal \texttt{First\allowbreak Hyp}.

They also show that \texttt{No\allowbreak Gain} forces $VV=\varnothing$.
The first five listed reasons directly contradict its defining
conditions (a band is exactly the candidate's excluded closed
band).  In the odd-order case, the frame implies
$\operatorname{ord}(h_p)\mid o$; since
$o\le x^{\epsilon/100}<x^{\epsilon/4}$, this contradicts
the candidate's lower-order bound.  Thus an unserved prime has
both gain $g_j(p)=0$ and $\tau_P(p)=0$, as required by LFS5.
Also a served prime must satisfy $p\le X''$ by (LFS.30).

Finally suppose a served top is omitted from $T_p$.
The literal definition of \texttt{TTP} retains all tops at a
signal prime and at $p\ge X'$.  Hence this omission forces
$p<X'$ in the inert-residue class.  The failure of
\texttt{No\allowbreak Gain} gives the candidate conditions and
$\Sigma_p\le\theta$, so $VV$ is exactly the good orbit
representatives.  The omitted top's actual representative
$X_G.bs(\mathrm{stepP}(p),\allowbreak{} t)$ is therefore not good.
Moreover $p$ lies in the original large-order residue-prime
set.  This proves the exact bad-top inclusion used in LFS8.
The second-class implication is immediate from its defining
branch: it is the same candidate and the same $\Sigma_p>\theta$.
% END LFS12

% BEGIN LFS13
\paragraph{LFS13: install all eleven finite-ledger clauses on the original record.}
\label{claim:LFS13}
Assume the original \texttt{Link\allowbreak G\_\allowbreak at} at $D,\allowbreak{} L$, with the same
$X_G,\mathcal F.GK,\allowbreak{} C,\allowbreak{} b_o$, and use the parameter record
(LFS.22).  LFS11 provides the unit rule and LFS12 the served,
high, and zero-unserved clauses.  We verify the others directly.

The top window says
$a\le a_t\le a+a/(kx^{13})=a(1+h)$, so in particular $a_t\ge0$.
The literal prime exponents in \texttt{ld\allowbreak At} are
\[
 I_{j,\allowbreak{} r,\allowbreak{} p}=\{s<\log_p^{\mathbb N}((a_t)_+)+2:s\text{ odd}\},
 \qquad r=\operatorname{val}(t).                             \tag{LFS.31}
\]
Their weights are the unchanged
$G.\mathrm{fill}(X_G,\allowbreak{} b_L,\allowbreak{} j,\mathrm{partner}(p,\allowbreak{} t,\allowbreak{} s))$.
At a top this fill is at least $b'/x$, at a thick place at
least $b$, and elsewhere at least $w''$.
Since $w''=b'/(xn)\le b'/x$ and $w''\le b$, every such
weight is at least $w''$.
The discriminant exponent set is
$[1,\log_d^{\mathbb N}((a_t)_+)+1]$.
Its weight is the original $w_{\rm At}$ at the source's
ramified representative; it is nonnegative and, when that
representative lies in $\Psi$, at least one by GFA10's exact
\texttt{Num\allowbreak N\_\allowbreak at} proof.  Outside $\Psi$ its defining
convention is zero.

At a served prime choose any original coefficient place over it.
For existence, the integral basis of the coefficient field makes
$\mathcal O_{E_0}/p\mathcal O_{E_0}$ a nonzero finite ring of
$p^{[E_0:\mathbb Q]}$ elements.  A maximal proper ideal exists
in this finite ring; its inverse image is a nonzero prime over
$(p)$.  Its residue characteristic is $p$,
and it belongs to the saturated support since $p\le X''$.
The \texttt{Parameters\allowbreak Fill} equalities identify its selected
tops with the image defining $T_p$, its $\mathrm{gain}$ with
$g_j(p)$, and its correction with $\tau_P(p)$.
They identify each $\mathrm{wPart}$ with the exact fill in
(LFS.31).  Reindex by the injective top labels and by
$J=\{1,\ldots,\allowbreak{} m\}$.  The third clause of
\texttt{Link\allowbreak G\_\allowbreak at} is then precisely
\[
 \sum_j\sum_{r\in T_p}(A_{j,\allowbreak{} r}(p)-\mathrm{loG}(L,\allowbreak{} C,\allowbreak{} b_o,\allowbreak{} p))
       -m\,\mathrm{flG}(L,\allowbreak{} C,\allowbreak{} p)
          \le\sum_jg_j(p)-\tau_P(p).                       \tag{LFS.32}
\]
In particular $T_p\subset T$ is the same actual subset, not a
new subset chosen to make the estimate hold.

For the floor envelope replace $C$ by $P.C=\max(C,0)$;
all factors multiplying $C$ are nonnegative for $a\ge1$ and
$p\ge2$.  The resulting formula is exactly $F(p)$ in LFS6,
because $P.f=2n=D.R.f$ and $P.N=D.R.N(m)$.
For the loss envelope put
\[
 A_1=r_Gb'\log x,\qquad
 A_2=b'/x+\eta+S_c+(w''+t)\log x/\delta_1,
\]
and
\[
 B_p=\Lambda a/(pk)+
       \begin{cases}ap^{-L_1(p)}&p<k^2x^{26},\\a/p^2&p\ge k^2x^{26}.
       \end{cases}
\]
All are nonnegative.  The actual \texttt{lo\allowbreak G} is its common
$C(S_c+k+\eta)\log_p(2a)$ term, plus a class-restricted
precision term $CA_1$ or $CA_2+b_o$, plus the class-restricted
orbit term $C\log x\,\allowbreak{} B_p$.  Replace $C$ by its nonnegative
maximum and remove class restrictions; this can only increase
each term.  The low-prime branch lies below $X'$ by LFS10,
and $p<P.p_0$ is exactly $p<k^2x^{26}$ for integer $p$.
Hence this is exactly
$C_p+\mathbf1_{p<X'}(P_p+O_p)$ with the source's $Q_1,\allowbreak{} Q_2$,
including the single additive $b_o$ in $Q_2$.

Finally the actual
\[
 L_1(p)=\lceil\log_p(k^2x^{78})\rceil_++1
\]
satisfies $p^{L_1(p)-1}\ge k^2x^{78}=K_0$ by monotonicity
of $p^y$ and the ceiling inequality.  This proves the last
clause and completes the literal \texttt{Parameters\allowbreak Hyp}.
The finite equation (LFS.32) uses a supplied local gain estimate;
the present argument reconstructs its exact installation rather
than asserting an independent producer for \texttt{Link\allowbreak G}.
% END LFS13

% BEGIN LFS14
\paragraph{LFS14: the universal original arithmetic reading of the asked sum.}
\label{claim:LFS14}
For every sufficiently late member of the original family,
every tuple $(e',\allowbreak{} m',\allowbreak{} AD,\allowbreak{} N,\allowbreak{} T',\allowbreak{} a',\allowbreak{} w')$ satisfying the literal
\texttt{Reads}, the heterogeneous equality
$AD=\mathrm{arithOf}(D.host)$, and the original cutoff condition,
we have
\[
 \sum_{t\in T'}\sum_{j\in\mathrm{Fin}(m')}
     \mathrm{ASK}(AD,\allowbreak{} N,\allowbreak{} a'_t,\allowbreak{} t,\allowbreak{} w'_j)
 =\mathrm{ASK}_{\rm original}(L)
 =(\mathrm{ldAt}).\mathrm{ASK}(d,\mathcal P,\allowbreak{} J,\allowbreak{} T).             \tag{LFS.33}
\]
We prove both equalities, including the ramified term.

The equalities in \texttt{Reads} first identify $e'$ with
the actual frame and $m'$ with $m$.  Its top-membership
equivalence and the injective integer labels identify $T'$
bijectively with the actual tops.  Its degree clause then gives
$a'_t=a_t$.  Its weight clause identifies, on the entire
arithmetic half $\Phi$, $w'_j$ with the original
$w_{\rm At}(j+1)$ under the original group equivalence
$E:D.R.GG\longrightarrow (\mathbb Z/2\mathbb Z)^2\times A$.
Explicitly, for a point of $\Phi$, its inverse lies in
$\Psi$; use $D.eTau^{-1}$ on that point to get exactly the
code requested by \texttt{Reads.\allowbreak hw}.  Thus this is equality
of weights at each point, not equality of distributions.
GFA10 supplies $w_{\rm At}\ge1$ there.

On the same arithmetic host AH9--AH11 identify the unramified
Frobenius and its two sign coordinates.  At an inert prime
its second multiplicative sign is $-1$.  Its $s$th power
lies outside $\Psi$ exactly when $s$ is odd.  For odd $s$,
the actual partner is
\[
 (\mathbf1_{\mathrm{signal}},\allowbreak{} t+s h_p),
\]
and multiplication of the two sign coordinates and the
additive frame coordinate shows that its image is exactly
the source's $\rho(\tau_1(t),\mathrm{Frob}_p^s)$.
The fill identities therefore make its weight exactly
$L.\mathrm{wPart}(j,\allowbreak{} t,\allowbreak{} p,\allowbreak{} s)$.
In additive notation this is the partner appearing in
$\mathrm{ASK}(AD,\allowbreak{} N,\cdot)$, since $E$ carries products to
sums and powers to multiples.  The host's quadratic sign
identifies its inert predicate with the original signal-or-
inert-residue predicate.  Thus every unramified prime summand
agrees term by term.

At $d$ the actual inertia group is
$\{1,\iota_d\}$ by AH11, with $\iota_d\notin\Psi$.
Every Frobenius-power coset has exactly one representative
outside $\Psi$: use $\mathrm{Frob}_d^s$ if it is outside,
and $\mathrm{Frob}_d^s\iota_d$ otherwise.
This is precisely the source's $g_D(s)$.  It lies outside
$\Psi$, and the partner map $\rho(\tau_1(t),\cdot)$ carries
that complement into $\Psi$.  Thus its resulting partner has
the actual weight at least one from GFA10.  The zero-weight
convention outside $\Psi$ mentioned in LFS13 is never used
for this tail argument.
The arithmetic sum over outside-half inertia representatives
therefore has one term, with the same partner and weight,
and the same coefficient $\log d$.  In particular it has
neither an extra factor $2$ nor a missing inertia average.

The cutoff causes no alteration.  For $W>0$ and
$a_t/W\le N$, every integer $M>N$ has
$\max(a_t/M-W,0)=0$.
This proves vanishing of all terms with exponent $s>N$,
because $p^s>s$, and of all primes $p>N$.
Independently, since the original weights on the partner
half are at least one, all terms with $p^s>(a_t)_+$ vanish.
Hence the infinite original prime-power sums can be cut
at exactly the sets (LFS.31) and the discriminant interval
in LFS13.  Primes above $(a_t)_+$ give zero, so extending
their finite universe to (LFS.27) gives zero extra terms.
The top window ensures $(a_t)_+<B$.

Finally replace $j\in\mathrm{Fin}(m)$ by $j+1\in J$
and use (LFS.26).  These finite bijections prove (LFS.33).
An arbitrary weight outside $\Phi$ is immaterial: every
partner in the defining asked sums belongs to $\Phi$.
The threshold needed here is only
$\max(3,\mathcal F.x_1+1)$, which activates the original
fill identities and GFA10's positivity proof.  In particular
it is chosen before the reading, its cutoff $N$, and all its
auxiliary types and functions.
% END LFS14

% BEGIN LFS15
\paragraph{LFS15: the actual same-record prime-class loss setting.}
\label{claim:LFS15}
Suppose $TT\ne\varnothing$, \texttt{Ch\allowbreak Badp\_\allowbreak at} holds,
\texttt{Ch\allowbreak W\_\allowbreak at} holds, and the original \texttt{Link\allowbreak G\_\allowbreak at}
holds on the record considered in LFS11--LFS14.
Then every hypothesis used by LFS7--LFS9 is satisfied with
\[
 \begin{split}
 \mathcal P&=\mathrm{PsAt}(L),\quad
 \mathcal C=\mathrm{cands}(L,\mathcal F.GK),\quad
 \mathcal R=\mathrm{largeOrderResiduePrimes}(L),\\
 J&=\mathrm{lifts}(D),\quad T=\mathrm{TTv}(D),\quad
 \Sigma_p=G.\mathrm{Sig}(X_G,\mathcal F.GK,\allowbreak{} a_L,\allowbreak{} b_L,\allowbreak{} p),\\
 io(p)&=G.\mathrm{ioP}(X_G,\mathcal F.GK,\allowbreak{} a_L,\allowbreak{} b_L,\allowbreak{} p),\quad
 D_{\rm fin}=\mathrm{ldAt}(X_G,\mathcal F.GK,\allowbreak{} C,\allowbreak{} b_o).
 \end{split}                                                   \tag{LFS.34}
\]
We verify the share clauses that have not yet been checked.
Each small-prime $\kappa$ is a minimum of $1$ and a
nonnegative expression.  Each large-prime $\kappa$ is the
infimum of a set of nonnegative feasible values, using the
literal \texttt{Extra.\allowbreak kap\allowbreak I\_\allowbreak def}; its infimum is nonnegative,
including the empty-set convention.  Thus $\Sigma_p\ge0$,
as also established in GFA4.  The \texttt{Ch\allowbreak W\_\allowbreak at} expression
is exactly (LFS.15), because the candidate sets agree and
$x^{-9}=(x^9)^{-1}$.  It has no hidden expectation remaining.

If $VV(p)$ is empty, $io(p)=0$.  Otherwise its literal value is
\[
             io(p)=2\Sigma_p\max_{v\in VV(p)}S_v.
\]
The second clause of \texttt{Link\allowbreak G\_\allowbreak at} gives
$S_v\le2mN_va/p$.  The original $N_v$ counts a subset of
the tops, so $N_v\le|TT|=q$; this is just the cardinality
bound for the actual filtered top set.
Multiplication by $2\Sigma_p\ge0$ proves the exact
$io(p)\le2\Sigma_p(2mqa/p)$ of LFS8.

By LFS11 the bad subsets have the exact original cardinalities.
Thus \texttt{Ch\allowbreak Badp\_\allowbreak at} is exactly the inequality used in
(LFS.18).  LFS12 supplies the first-class, second-class,
and omitted-top implications.  LFS13 supplies every
finite-ledger clause.  Every prime in $\mathcal P$ is odd
by LFS12, and $q\ge1$ by nonemptiness and (LFS.26).
This proves the full literal \texttt{Prime\allowbreak Class\allowbreak Loss\allowbreak Hypotheses}
at (LFS.34), or its source's \texttt{Loss\allowbreak Setting\allowbreak Family}
instance with the harmless external index type \texttt{Unit}.
Neither the selected top set nor the local gains are enlarged
to establish these hypotheses.
% END LFS15

% BEGIN LFS16
\paragraph{LFS16: both original losses, with one constant and one threshold.}
\label{claim:LFS16}
Assume, on the same original family,
\[
 \texttt{Chosen\allowbreak TT}(\mathcal F),\qquad
 \texttt{Chosen\allowbreak Badp}(\mathcal F),\qquad
 \texttt{Link\allowbreak G}(\mathcal F).                              \tag{LFS.35}
\]
Then the complete literal \texttt{Losses\allowbreak Small\allowbreak Shares} holds:
there are $C_L\ge0$ and $x_L$ such that, for all $x\ge x_L$
and every $i\in\mathcal F.I_x$,
\[
 \texttt{Ch\allowbreak W\_\allowbreak at}(\mathcal F.Gx(x,\allowbreak{} i),\mathcal F.GK)
       \Longrightarrow\texttt{Loss\_\allowbreak at}(\mathcal F.L(x,\allowbreak{} i),\allowbreak{} C_L,0).
                                                               \tag{LFS.36}
\]
To prove it, take once and for all the witnesses
$C,\allowbreak{} x_G,\delta$ of \texttt{Link\allowbreak G}.
Take $e,\allowbreak{} K$ from (LFS.23), $C_{\rm ord}=12$, and
\[
                              C_L=3|K|+13.                   \tag{LFS.37}
\]
Choose $x_L$ to be the maximum of the \texttt{Chosen\allowbreak TT} and
\texttt{Chosen\allowbreak Badp} thresholds, $x_G$, LFS10's window
threshold, LFS9's uniform numerical threshold, and LFS14's
reading threshold.  All are chosen before $i$ and before
any arithmetic reading.

Fix $x\ge x_L$ and $i$, and suppose the premise of (LFS.36).
The first chosen-top bound is
$|TT|\ge\delta_1 n/2>0$, so $TT$ is nonempty.
LFS10 installs the original parameter record with
$b_o=b^{\delta(x)}$, LFS15 installs the finite loss setting,
and LFS9 gives (LFS.20).  The cardinality identities of
LFS11 identify its $U$ exactly with $|TT|mb$.
Since $2\log4\le3$, $0\le P.C\le K\le|K|$, and
$\sigma\le1$,
\[
 \mathrm{IOsum}(L)\le U\le C_LU,
 \qquad\mathrm{ASK}_{\rm original}(L)-L.\mathrm{Ga}
                  \le(3P.C+13)U\le C_LU.                 \tag{LFS.38}
\]
For the elementary numerical inequality, $\log2<3/4$
follows from $e^{3/4}>1+3/4+(3/4)^2/2>2$, so
$2\log4=4\log2<3$.

Now take an arbitrary \texttt{Reads} tuple, arbitrary
$AD$ heterogeneously equal to the original host arithmetic,
and any cutoff satisfying its specified inequalities.
LFS14 replaces its asked sum by precisely the middle
quantity in (LFS.38).  LFS11 replaces its finite gain and
prime-logarithmic sum by the original $L.\mathrm{Ga}$ and
$\mathrm{IOsum}(L)$.  The source factor
$L.sc.x^{0\cdot L.sc.\epsilon}$ is exactly $1$.
Thus both clauses, with all their universal quantifiers,
are exactly $\texttt{Loss\_\allowbreak at}(L,\allowbreak{} C_L,0)$.
This proves (LFS.36), not merely a loss estimate for one
preferred reading or one cutoff.
% END LFS16

% BEGIN LFS17
\paragraph{LFS17: preserve the chosen realization and obtain literal \texttt{Loss\allowbreak FS}.}
\label{claim:LFS17}
In addition to (LFS.35), assume the original
$\texttt{Chosen\allowbreak WS}(\mathcal F)$, with its own threshold $x_W$.
Let $C_L,\allowbreak{} x_L$ be supplied by LFS16 and put
$x_0=\max(x_L,\allowbreak{} x_W)$.  For any $x\ge x_0$ in
$\mathcal F.\mathrm{Win}$, take the exact witnesses of
\texttt{Chosen\allowbreak WS}:
\[
 i_0\in\mathcal F.I_x,\qquad
 a_v,\allowbreak{} b_v:\mathbb N\longrightarrow\mathbb Z,
 \qquad i_1\in\mathcal F.I_x.                              \tag{LFS.39}
\]
Their stated windows are retained verbatim:
\[
 \forall r\in\mathcal F.\mathrm{TTn}(x),\quad
 \mathrm{aWin}(\mathcal F.L(x,\allowbreak{} i_0),\allowbreak{} a_v(r)),\quad
 (\mathcal F.L(x,\allowbreak{} i_0)).sc.b'\le b_v(r)\le\mathcal F.bTop(x).
                                                               \tag{LFS.40}
\]
The same $i_1$ realizes them at every one of its actual tops:
\[
 (\mathcal F.L(x,\allowbreak{} i_1)).aT(t)=a_v(\operatorname{val}(t)),\qquad
 (\mathcal F.L(x,\allowbreak{} i_1)).bT(t)=b_v(\operatorname{val}(t)).         \tag{LFS.41}
\]
For an arbitrary $i\in\mathcal F.I_x$ satisfying these same
two equalities at all its tops, \texttt{Chosen\allowbreak WS} supplies
$\texttt{Ch\allowbreak W\_\allowbreak at}$ for that $i$.  Apply (LFS.36), since
$x\ge x_L$, to obtain \texttt{Loss\_\allowbreak at} with the unchanged
$C_L$ and exponent $0$.
The tuple $C_L,\allowbreak{} x_0,\allowbreak{} i_0,\allowbreak{} a_v,\allowbreak{} b_v,\allowbreak{} i_1$, (LFS.40), (LFS.41),
and this universal implication is exactly
\[
 \boxed{\quad
 \texttt{Chosen\allowbreak TT}\ \wedge\ \texttt{Chosen\allowbreak Badp}\ \wedge\
 \texttt{Link\allowbreak G}\ \wedge\ \texttt{Chosen\allowbreak WS}
       \ \Longrightarrow\ \texttt{Loss\allowbreak FS}(\mathcal F,0).
 \quad}                                                       \tag{LFS.42}
\]
HW3 supplies the active preselected-pair form of
\texttt{Chosen\allowbreak WS} from its already stated constructor
hypotheses.  The present implication requires only that
literal predicate and does not select a new pair or realize
an arbitrary admissible pair.  In particular the distinct
\texttt{Cover\allowbreak All} property is absent from (LFS.42).
The source's later \texttt{cover\allowbreak S\_\allowbreak carry} use, starting from
the specific member $i_1$, belongs to the subsequent
same-family ledger assembly.
% END LFS17

\paragraph{Scope at the endpoint.}
The complete finite implication (LFS.42) closes the original
loss installation task on a supplied compatible family.
It leaves the original local arithmetic producer
\texttt{Link\allowbreak G}, the chosen-top and chosen-bad inequalities,
and the common chosen-weight realization in their exact
stated upstream scopes.  It neither replaces those inputs
by abstract numerically convenient data nor asserts that
the full human proof of RH or unrestricted Dirichlet GRH
has thereby been completed.


\section{The common global ledger and its final uniform installation}
\label{sec:global-final-installation-v035}

\status{This chapter constructs the actual tied ledger from the nine
original numerical and lattice predicates.  Together with LFS it
completes the remaining global assembly implication.  The common
family, compatible second layer, duties, and arithmetic and spectral
producer inputs retain their original scope.  No new family or
unrelated model is chosen to satisfy a missing condition.}

Throughout, $\mathcal F:\texttt{Fam\allowbreak P}(\epsilon)$ is the supplied
original family, with $0<\epsilon<1/52$.  At a scale $x>\mathcal F.x_1$
write $D_i=\mathcal F.D(x,\allowbreak{} i)$ and $L_i=\mathcal F.L(x,\allowbreak{} i)$.
We distinguish the original record $L_i:\texttt{Size\allowbreak Data}_2(D_i)$
from the eleven-predicate \texttt{LChain.\allowbreak Size\allowbreak Data}$_2$ on a numerical
ledger.  The latter is the input to ZW1--ZW4.  All eventual thresholds
below are chosen before the member, base point, word, or finite place.

% BEGIN GFI1
\paragraph{GFI1: the common family dictionary includes arithmetic data.}
\label{claim:GFI1}
The original fields \texttt{R\_\allowbreak eq}, \texttt{L\_\allowbreak x},
\texttt{L\_\allowbreak eps} and \texttt{L\_\allowbreak J} imply that all $D_i$ have the
same ranks and frame, and all $L_i$ have the same scale.  Indeed, a
scale is specified by $x,\epsilon,\allowbreak{} J$; its remaining components are
functions of those values.  Hence $a,\allowbreak{} b,\allowbreak{} b',\allowbreak{} k,\allowbreak{} n,\allowbreak{} m,\allowbreak{} f,\allowbreak{} L$ agree, where
$L=x^\epsilon$, $b=e^L$, $a=xb$ and $f=2n$.
The separate original family field \texttt{eta\_\allowbreak eq} gives the
common value of $L_i.\mathrm{etaSz}$; it is not inferred from
the scale alone.  This is the equality used by \texttt{Tie\allowbreak Global}
and the remaining-size comparison in GFI12.

The map from the original top set to integer labels is injective and
has image $T_n=\mathcal F.TTn(x)$.  It therefore gives a bijection
\[
 D_i.\mathrm{host}.TT\simeq T_n,
 \qquad \#D_i.\mathrm{host}.TT=\#T_n.                 \tag{GFI.1}
\]
The second-set integer labels agree by
\texttt{realised\allowbreak Chain\_\allowbreak torus\_\allowbreak eq}.  The family also fixes $q,\allowbreak{} d,\allowbreak{} QS$,
the root discriminant, $g$, and $C_\Gamma$.  Consequently it fixes
$v$, $\log(g_0g_1)$, $L_a$, $B_0$, $N_c$ and $\Gamma$.
For $N_c$ use the actual second-set equality before taking its
cardinality.  For $\Gamma$ use
$\exp(C_\Gamma m^5 fL)$.  The quadratic logarithmic derivative
is fixed because its character is defined by the same prime $d$;
this identifies the actual value at one, including its real part.

There is also equality of the arithmetic records after transporting
along the frame equality.  Their modulus is the same $QSqd$.
The field \texttt{art\_\allowbreak eq} gives the same signed reciprocity value
and cyclic integer label at every integer coprime to that modulus.
Every residue unit is represented by its least nonnegative integer,
which is coprime to the modulus.  Thus the two reciprocity
homomorphisms agree on every unit.  Injectivity of the cyclic labels
and of the two sign codes recovers equality of the actual values.
The definition of \texttt{arith\allowbreak Of} then gives equality of the
arithmetic records, with the prescribed interchange of the two sign
coordinates.  This is equality on the source objects, rather than a
choice of a second reciprocity model.
% END GFI1

% BEGIN GFI2
\paragraph{GFI2: carry one realized degree pair to every base point.}
\label{claim:GFI2}
Suppose the original \texttt{Loss\allowbreak FS} conclusion is available at $x$.
It supplies $i_0$, integer functions $a_v,\allowbreak{} b_v$, their admissible
windows, and an index $i_1$ realizing those functions on every top.
It also supplies \texttt{Loss\_\allowbreak at} for every member realizing them.
Let $H=D_{i_0}.R.A$.  For each $h\in H$, its integer label is below
$n$.  The family's \texttt{cover\allowbreak S} property applied to $i_1$ gives
an index $\mathrm{pt}(h)$ whose distinguished base point has that
label and whose top and bottom degrees agree with those at $i_1$
whenever the top labels agree.

Every top at $\mathrm{pt}(h)$ has a counterpart at $i_1$, by
(GFI.1).  Their degrees are therefore exactly
$a_v(\mathrm{val}(t))$ and $b_v(\mathrm{val}(t))$.  Choose these
indices for the finite set $H$.  We have constructed the literal
\texttt{Mod} record
\[
 M=(i_0,\allowbreak{} a_v,\allowbreak{} b_v,\mathrm{pt}),                            \tag{GFI.2}
\]
including both its base-point equality and its degree-carrying
condition.  All members $\mathrm{pt}(h)$ consequently satisfy the
same loss estimate.  Only the actual realized pair is transported;
no assertion that every admissible pair is realized is used.
% END GFI2

% BEGIN GFI3
\paragraph{GFI3: top sets, place labels, and weights on that model.}
\label{claim:GFI3}
Use the original frame $e=D_{i_0}.R.\mathrm{toFrame}$ and put
\[
 T=\{h\in H:\mathrm{val}(h)\in T_n\},\qquad
 T_2=e.\mathrm{dig}(D_{i_0}.\mathrm{host}.T2),\qquad
 a_h=a_v(\mathrm{val}(h)).                              \tag{GFI.3}
\]
The inverse of the integer-label map gives an equivalence
$e_T:T_n\simeq T$.  Hence $\#T=\#T_n$.  The original second-set
entries lie below the digit map's injectivity bound: its stored
range condition is $4L_w s<n_0$, and $L_w\ge2$.  Thus
\[
 \#T_2=\#D_{i_0}.\mathrm{host}.T2.                     \tag{GFI.4}
\]
The same set describes the second-set image for every $D_{\mathrm{pt}(h)}$
by GFI1.

The original place coding maps its $\tau$-set bijectively to
$\Phi\subset E_0\times H$.  Injectivity follows from injectivity
of both sign codes, the cyclic labels, and the original $e_\tau$;
surjectivity follows from membership and the common cardinality
$\#\Phi=2n=f$.  Choose the source's finite enumeration
$e_P:\{0,\ldots,\allowbreak{} f-1\}\simeq\Phi$.

For $h\in H$, $0\le j<m$ and $\rho\in E_0\times H$, define
$w_{h,\allowbreak{} j}(\rho)$ to be the original $w_{j+1}(\tau)$ if $\rho$
is the code of a place of $D_{\mathrm{pt}(h)}$, and to be $1$
otherwise.  The representing $\tau$ is unique, so this is the
literal source weight.  Reindexing gives, for any function $F$,
\[
 \sum_{r=0}^{f-1}F(w_{h,\allowbreak{} j}(e_P(r)))
 =\sum_\tau F(w_{j+1}(\tau)).                           \tag{GFI.5}
\]
Similarly the lifted indices are exactly $1,\ldots,\allowbreak{} m$, and
\[
 \sum_{t\in D_{\mathrm{pt}(h)}.\mathrm{host}.TT}F(a_T(t))
 =\sum_{t\in T_n}F(a_v(t)).                              \tag{GFI.6}
\]
These are bijective reindexings of the actual sums.  They preserve
their multiplicities and all metric-dependent quantities.
% END GFI3

% BEGIN GFI4
\paragraph{GFI4: construct every entry of the original numerical ledger.}
\label{claim:GFI4}
Let $W_h$ be the original word set of $D_{\mathrm{pt}(h)}$.
It is finite and nonempty, with its stored distinguished word.
For a cutoff $N$, constants $C_H,\bar\eta,\bar e$, and error
functions $\eta_h,\allowbreak{} e_h$, define the ledger $\mathscr D$ as follows.
Its scalar entries $x,\epsilon,\allowbreak{} a,\allowbreak{} b,\allowbreak{} b',\allowbreak{} k,\allowbreak{} L,\allowbreak{} L_a,\allowbreak{} B_0,\allowbreak{} N_c,\allowbreak{} v,
\log(g_0g_1),\log\Gamma$ are those at $i_0$; its Euler constant is
$\gamma$ and its logarithmic derivative is $L_{i_0}$'s actual one.
Its top and place weights are (GFI.3) and (GFI.5).

For $u\in W_h$, put
\[
\begin{aligned}
 d_{h,\allowbreak{} u}&=\operatorname{rank}(D_{\mathrm{pt}(h)}.\mathrm{div.Nu}(u)),&
 T_{h,\allowbreak{} u}&=(D_{\mathrm{pt}(h)}.\mathrm{mem}(u)).A.T,\\
 \mu_{N,\allowbreak{} h,\allowbreak{} u}&=\mu(D_{\mathrm{pt}(h)}.\mathrm{div.Nu}(u)),&
 \mu^{\max}_{N,\allowbreak{} h,\allowbreak{} u}&=\mu_{\max}(D_{\mathrm{pt}(h)}.\mathrm{div.Nu}(u)),\\
 \mu_{M,\allowbreak{} h,\allowbreak{} u}&=\mu((D_{\mathrm{pt}(h)}.\mathrm{fm}(u)).\mathrm{lat}(0)),&
 \mu^{\max}_{M,\allowbreak{} h,\allowbreak{} u}&=\mu_{\max}((D_{\mathrm{pt}(h)}.\mathrm{fm}(u)).\mathrm{lat}(0)),\\
 h_{h,\allowbreak{} u}&=D_{\mathrm{pt}(h)}.\mathrm{div.h}(u),&
 E_{h,\allowbreak{} u,\allowbreak{} j}&=\sum_{w\in\mathrm{supp}}\mathrm{sz}(w,\allowbreak{} E(u,\allowbreak{} j+1,\allowbreak{} w)).
\end{aligned}                                                     \tag{GFI.7}
\]
The last sum uses $L_{\mathrm{pt}(h)}$'s original support and
valuation normalization.  Set $\mu_{C,\allowbreak{} h}=\mu(RR^\vee)$ and
$\mu^{\min}_{C,\allowbreak{} h}=\mu_{\min}(RR^\vee)$ for that same member.
Copy its $\mathrm{Fo}$ and $\mathrm{Ga}$ and set
\[
 \mathrm{Ta}_h=\sum_{w\in\mathrm{supp}}\mathrm{sz}(w,\tau_w),
 \qquad
 \mathrm{IO}_h=\sum_{w\in\mathrm{supp}}\ell_we_w\,\mathrm{io}_w.
                                                               \tag{GFI.8}
\]
Thus $\mathrm{IO}$ has the source's single $\ell_we_w$ factor;
it is not confused with the doubled squared-metric sum.
Using the common arithmetic record $A$ from GFI1, set
\[
 \mathrm{ASK}_h=\sum_{t\in T}\sum_{j=0}^{m-1}
       \mathrm{ASK}(A,\allowbreak{} N,\allowbreak{} a_t,\allowbreak{} t,\allowbreak{} w_{h,\allowbreak{} j}),\qquad HH_h=x^{C_H}b.
                                                               \tag{GFI.9}
\]
Finally copy the actual retained defect $\delta_h$, put
$\bar\delta=e^{-x^{\epsilon/(7J)}}$, and use the chosen
$\eta_h,\allowbreak{} e_h,\bar\eta,\bar e$.

These formulas specify every field of the source's \texttt{Data}.
GFI1 proves \texttt{Tie\allowbreak Global} at every $h$, including the common
$X''$, $\eta_{\mathrm{Sz}}$ and $X'$.  Formulae (GFI.7)--(GFI.8)
prove every \texttt{Tie\allowbreak Local} field by definition.  GFI2 and the
common ranks and primes prove \texttt{Tie\allowbreak Pts}.  The source and target
lattices, rational maps, metrics and word sets are retained without
any change of model.
% END GFI4

% BEGIN GFI5
\paragraph{GFI5: one character functional and an admissible cutoff.}
\label{claim:GFI5}
Take the original character functional $\ell$ supplied by GFA12
at $i_0$, for the common arithmetic record $A$.  It agrees with the
logarithmic derivative of every primitive character realizing each
nonzero additive character.  It also satisfies
$L_{i_0}.\mathrm{logDerivativeAtOne}=\operatorname{Re}\ell(\mathrm{chC})$.
Keep this single $\ell$ for all $h$.

GFI3, the unchanged second set, and that character identity prove
\texttt{Tie\allowbreak Arith}.  The equalities of top weights, place weights,
Euler constant, quadratic logarithmic derivative, and the ASK sum
in (GFI.9) are exactly the five fields of \texttt{Ident}.

The original node estimate of GFA10 gives $w_{h,\allowbreak{} j}(\rho)\ge1$
for $\rho\in\Phi$.  Each common top label can be read at the
single member $\mathrm{pt}(0)$; its original top window and GFA11
give
\[
 a\le a_t\le2a,\qquad |a_t/b-x|\le1.                \tag{GFI.10}
\]
Choose $N=\lceil2a\rceil_{\mathbb N}$.  Then
$a_t/w_{h,\allowbreak{} j}(\rho)\le a_t\le N$ for all permitted $h,\allowbreak{} j,\allowbreak{} t,\rho$.
Moreover the five original \texttt{Reads} conditions hold:
frame and rank equality are GFI1, top membership is (GFI.1), top
degree equality is GFI2, and place-weight equality is GFI3.
The arithmetic-record equality is the one proved in GFI1.  Thus
the loss statement can actually be applied to (GFI.9).
% END GFI5

% BEGIN GFI6
\paragraph{GFI6: a common threshold, constants, and positivity record.}
\label{claim:GFI6}
Assume the nine original family predicates
\[
\begin{gathered}
 \texttt{Count\allowbreak F},\ \texttt{Slopes\allowbreak MF},\ \texttt{Slopes\allowbreak NF},\
 \texttt{Lat\allowbreak F},\ \texttt{Nodes\allowbreak F},\ \texttt{Frame\allowbreak F}(C_H),\\
 \texttt{Loss\allowbreak FS}(e_L),\ \texttt{Char\allowbreak F},\ \texttt{Win\allowbreak F},
\end{gathered}                                                     \tag{GFI.11}
\]
together with the original \texttt{Chosen\allowbreak TT},
\texttt{Chosen\allowbreak QD}$(C'_5)$, \texttt{Chosen\allowbreak Thin}$(c_6)$,
\texttt{Link\allowbreak IRR} and \texttt{Link\allowbreak DIV}$(C_{\rm abs},\allowbreak{} C_H)$.
Assume $C_{\rm abs},\allowbreak{} C_H,\allowbreak{} e_L\ge0$ and $c_6>0$.
Extract their constants and thresholds first, using
\texttt{Char\allowbreak F} at $(c_6,\allowbreak{} C'_5)$ and \texttt{Win\allowbreak F} at $c_6$.
Write $C_c,\allowbreak{} C_s,\allowbreak{} C_1,\allowbreak{} C_\ell\ge0$ for the counting, two slope,
and loss constants, and $C_d$ for the divisibility constant.
Put
\[
 C_d^+=\max(C_d,0),\quad
 K=C_d^++C_c+C_s+C_1+C_\ell+6,\quad
 c_{HH}=C_H,\quad C_L=C'_5.                              \tag{GFI.12}
\]
The ledger constant record has, in source order,
\[
 (c_D,\allowbreak{} e_D,\allowbreak{} c_E,\allowbreak{} c_H,\allowbreak{} c_C,\allowbreak{} c_S,\allowbreak{} c_{\mathrm{sl}},\allowbreak{} c_P,\allowbreak{} e_P,\allowbreak{} c_{HH})
 =(C_d^+,\allowbreak{} C_{\rm abs},\allowbreak{} C_d^+,\allowbreak{} C_d^+,\allowbreak{} C_c,\allowbreak{} C_s,\allowbreak{} C_1,\allowbreak{} C_\ell,\allowbreak{} e_L,\allowbreak{} C_H).
                                                               \tag{GFI.13}
\]
Here $c_H$ in the record is a height-comparison coefficient,
distinct from the exponent $C_H$.

Take one threshold greater than $\mathcal F.x_1$ and all extracted
thresholds.  For example, $\mathcal F.x_1+1$ plus the sum of
their absolute values works.  Only then fix $x$ in the original
window, choose (GFI.2), and use \texttt{Frame\allowbreak F} to choose
$\bar\eta,\bar e$ before any member.  For every $h$ choose its
frame errors $\eta_h,\allowbreak{} e_h$ within these same bounds.  They satisfy
\[
 0\le\bar e\le\tfrac14,\quad 0\le\bar\eta,\qquad
 (\bar\eta+4\sqrt{\bar e})x^{C_H}\le6.                 \tag{GFI.14}
\]

The ledger now satisfies all seventeen positivity conditions in
\texttt{Hida\allowbreak Setting}.  Indeed $x,\allowbreak{} b,\allowbreak{} a,\allowbreak{} b',\allowbreak{} k,\allowbreak{} v,\allowbreak{} f$ have their
required signs; $a_t\ge0$ follows from (GFI.10), and place weights
are positive by GFI5.  Actual ranks are positive by \texttt{Lat\allowbreak F},
and actual level primes are positive.  The same predicate gives
$\mu_N\le\mu_N^{\max}$.  Also $C_1\ge0$ and
$\bar\delta=e^{-x^{\epsilon/(7J)}}<1$, since $J\ge1$ and
$\epsilon>0$.  The remaining three error conditions are (GFI.14).
No positivity of an artificial rank or replacement form is used.
% END GFI6

% BEGIN GFI7
\paragraph{GFI7: divisibility and the two comparison fields.}
\label{claim:GFI7}
The last three estimates in the original \texttt{Link\allowbreak DIV} record
give \texttt{DIV}, \texttt{EComp5}, and \texttt{HComp} for
$\mathscr D$.  Here is the complete transport and constant check.
First $L_a,\allowbreak{} B_0,\allowbreak{} N_c\ge0$ by the original windows and their
definitions.  All the quantities multiplying $C_d$ in those three
estimates are therefore nonnegative.  Replacing $C_d$ by $C_d^+$
is consequently valid even if the supplied $C_d$ was negative.

For \texttt{DIV}, GFI1 identifies the original
$\log\Omega=\gamma+\log(\pi/v)-L'/L(1)$ at every member.
Using (GFI.6), its main term becomes exactly
\[
 \sum_{t\in T_n}\bigl[m(\phi(a_v(t))-a_v(t)\log\Omega)
                       -a_v(t)\log T_{h,\allowbreak{} u}\bigr],
 \qquad \phi(y)=y\log y-y.                            \tag{GFI.15}
\]
The other terms are $C_d^+n_TaL_a$ and
$C_d^+m(B_0+N_c)x^{C_{\rm abs}\epsilon}$ by GFI1.
This is precisely the ledger's divisibility inequality.

For the place-error comparison the source gives, for each lift,
\[
 |E_{h,\allowbreak{} u,\allowbreak{} j}-A_T\log(g_0g_1)|\le C_d^+B_0,
 \qquad A_T=\sum_{t\in T_n}a_v(t).
\]
Sum over the $m$ lifts and apply the triangle inequality.  The
result is
\[
 \left|\sum_{j=0}^{m-1}E_{h,\allowbreak{} u,\allowbreak{} j}
                -mA_T\log(g_0g_1)\right|\le C_d^+mB_0,
                                                               \tag{GFI.16}
\]
the full \texttt{EComp5}.  Finally the source compares
$h_{h,\allowbreak{} u}+A_T\log T_{h,\allowbreak{} u}$ for any two words, with upper bound
$C_d^+(n_TaL_a+mB_0)$.  Pulling the word-dependent logarithm
inside the finite top sum gives the literal \texttt{HComp}.
% END GFI7

% BEGIN GFI8
\paragraph{GFI8: counting and both slope fields retain their normalization.}
\label{claim:GFI8}
Apply the supplied \texttt{Count\_\allowbreak at} to $D_{\mathrm{pt}(h)}$.
Replace its lift sum by $j=0,\ldots,\allowbreak{} m-1$, and its place sum by
$e_P$ using (GFI.5).  The factor $1/f$ remains unchanged; GFI1
identifies $m,\allowbreak{} k,\allowbreak{} v,\allowbreak{} L,\allowbreak{} L_a$ and $\log(g_0g_1)$.  The three resulting
inequalities are exactly \texttt{Count5}: the lower bound for
$\log T_{h,\allowbreak{} u}$, the upper bound $C_cmL$, and the cross-word
variation $C_c\log k$.

The two \texttt{Slopes\allowbreak M\_\allowbreak at} inequalities become
\[
 -C_sn_TaL_a(1+m/k)\le\mu_{M,\allowbreak{} h,\allowbreak{} u},\qquad
 \mu^{\max}_{M,\allowbreak{} h,\allowbreak{} u}\le C_sn_TaL_a,                    \tag{GFI.17}
\]
by (GFI.1) and GFI1.  They are the ledger's \texttt{Slopes}.
For the quotient slopes, reindexing the lift sum gives equality
of the original $\mathrm{Ee}$ with $\sum_jE_{h,\allowbreak{} u,\allowbreak{} j}$.
The upper and lower quotient bounds then have the same correction
$\mathrm{Ga}_h+\mathrm{Fo}_h-\mathrm{Ee}_{h,\allowbreak{} u}$ and the same
lower loss
\[
 C_1\mathrm{IO}_h+C_1n_Tm(a/k+b')\log T_{h,\allowbreak{} u}+C_1X''+1.
                                                               \tag{GFI.18}
\]
This supplies \texttt{Quotient\allowbreak Slope\allowbreak Bounds\allowbreak With}, with the original
constant $1$, the original $X''$, and the original lattice metrics.
% END GFI8

% BEGIN GFI9
\paragraph{GFI9: the loss witness supplies both ledger inequalities.}
\label{claim:GFI9}
Every $\mathrm{pt}(h)$ carries the pair in GFI2, so it satisfies
the same \texttt{Loss\_\allowbreak at} with constant $C_\ell$.  Its first
clause is, by GFI1 and (GFI.8),
\[
 \mathrm{IO}_h\le C_\ell n_Tmb\,\allowbreak{} x^{e_L\epsilon}.
\]
Its second clause quantifies over every permitted reading and
cutoff.  Apply it to the actual $A,\allowbreak{} T,\allowbreak{} a,\allowbreak{} w$ of GFI3 and the cutoff
$\lceil2a\rceil$; GFI5 checked every hypothesis of that application.
It gives
\[
 \mathrm{ASK}_h-\mathrm{Ga}_h
                  \le C_\ell n_Tmb\,\allowbreak{} x^{e_L\epsilon}.             \tag{GFI.19}
\]
The right-hand side is exactly $\mathscr D.\mathrm{LoP}(C)$.
These are both fields of \texttt{Loss\allowbreak Bounds}, simultaneously
for every $h$.  In particular no selected reading has been substituted
for the universally quantified statement before its hypotheses were
checked.
% END GFI9

% BEGIN GFI10
\paragraph{GFI10: frame, height, and defect on the same rational map.}
\label{claim:GFI10}
The definitions (GFI.7) show that the ledger's weighted mean $s_h$
is the source $s_{\rm Frame}$ and its excess $u_h$ is the source
$u_{\rm Frame}$.  The finite rank sum in the denominator is unchanged.
Thus the full all-height estimate in \texttt{Frame\allowbreak F}, specialized
to $HH_h=x^{C_H}b$, is precisely \texttt{Slope\allowbreak Ineq} whenever its
height hypotheses hold.  Those hypotheses do hold: \texttt{Link\allowbreak DIV}
gives $s_h+\mu_{\max}(RR)\le x^{C_H}b$, and \texttt{Lat\allowbreak F}
gives $\mu_{\min}(RR^\vee)=-\mu_{\max}(RR)$.
Hence
\[
 s_h-\mu^{\min}_{C,\allowbreak{} h}\le HH_h,
 \qquad 0<HH_h=x^{C_H}b.                               \tag{GFI.20}
\]
The upper height bound in \texttt{HHBound} is equality because
the record has $c_{HH}=C_H$.

The source's \texttt{Link\allowbreak IRR} supplies
$0\le\delta_h\le\bar\delta$ for these actual retained rows.
The choices made after (GFI.14) supply
$0\le\eta_h\le\bar\eta$ and $0\le e_h\le\bar e$.
These six inequalities are \texttt{FParams}.  The frame estimates
and the height bound concern the same $RR^\vee$ and the same
original rational maps; no further dual or normalization is introduced.
% END GFI10

% BEGIN GFI11
\paragraph{GFI11: every error-window field and the trace-error budget.}
\label{claim:GFI11}
The chosen top density implies $n_T\ge1$: its lower bound is
$x^{-\epsilon/8}n/2>0$, and $n_T$ is an integer.  We have $x\ge1$,
$0<\epsilon\le1$, $b\ge1$, $a=xb$, and $k\ge1$.
The original windows of GFA11 give
\[
\begin{gathered}
 m\le x^{1/2+\epsilon/2},\quad xm\le k\le x^{e_K},\quad
 k\le a,\quad 0\le L_a\le2x^\epsilon,\\
 kb'\le a,\quad X''\le a,\quad
 \log\#W_h\le(k-3)\log2,
 \qquad e_K\ge0.                                      \tag{GFI.21}
\end{gathered}
\]
Also $0\le L=x^\epsilon\le2x^\epsilon$, $b'>0$, $X''\ge0$,
and $B_0,\allowbreak{} N_c\ge0$.  Top nonnegativity and the upper bound $2a$
are (GFI.10).  Since $x^{\epsilon/(7J)}\ge1$,
$\bar\delta\le e^{-1}\le1/2$.  Every nonnegative coefficient
in (GFI.13) is at most $K$, and the three exponent parameters
are nonnegative.  Together with (GFI.14) these prove every field
of \texttt{Err\allowbreak Windows0}; in particular the word-cardinality bound
is at every actual member.

For the extra field of \texttt{Err\allowbreak Windows}, (GFI.14) gives
\[
 (\bar\eta+4\sqrt{\bar e})x^{C_H}b\le6b
 \le Kx^\epsilon n_Ta.                                \tag{GFI.22}
\]
The second inequality uses only $K\ge6$, $x^\epsilon\ge1$,
$n_T\ge1$, and $a=xb\ge b$.  It proves the full trace-error
budget with the very same $K$ and $C_H$.  No weaker window
statement is substituted at this step.
% END GFI11

% BEGIN GFI12
\paragraph{GFI12: the remaining growth, cost, and size records.}
\label{claim:GFI12}
The actual definitions of the two budgets give equalities
\[
 B_0=2nmbx^\epsilon,\qquad
 N_c=\frac{\#T_2\,\allowbreak{} kb x^\epsilon}{n(nx^{-\epsilon})}.      \tag{GFI.23}
\]
For the first use $f=2n$; for the second use (GFI.4).
Thus the two \texttt{BNBounds} inequalities hold with equality.

Put $\eta=\eta_{\rm Sz}$, $X'=e^{x^{\epsilon/3}}$ and
$\theta=e^{-c_6x^{\epsilon/8}}$.  These have the required signs,
and $b'X'=a=xb$.  The near-$x$ top condition is (GFI.10).
The node conditions follow from GFA10 after transporting the same
second set and the same $\eta_{\rm Sz}$ using GFI1.  GFI5 supplies
the common cutoff condition.  These are all fields of
\texttt{Growth\allowbreak Inequality\allowbreak Sizes}.  Applying the last two clauses
of the single character record from GFI5 to the supplied
\texttt{Chosen\allowbreak Thin} and \texttt{Chosen\allowbreak QD} gives
\texttt{Thin\allowbreak Scales} and \texttt{Values\allowbreak At\allowbreak X} with this $\theta$
and $C_L=C'_5$.  Therefore the full remaining growth record holds.

The same windows give
\[
 2f(x+1)\theta\le1/x,\qquad
 6f(x+1)(\eta+1)\le b/x,                                \tag{GFI.24}
\]
and the place cardinality gives $f_{\rm Pl}(H)=f$, proving
\texttt{Cost\allowbreak Windows}.  Finally the five \texttt{Rest\allowbreak Size}
conditions are exactly
\[
 a\le2xb,\quad x^{-\epsilon/8}n/2\le n_T,\quad
 x^{-\epsilon}\le\#T_2/n,\quad
 x^{1/2-\epsilon}\le m,\quad k\le n^2x^{1/2-\epsilon}.
                                                               \tag{GFI.25}
\]
They follow, respectively, from $a=xb$, the original top density,
the second-set window and (GFI.4), and the final two rank windows.
All these conclusions concern the same ledger and cutoff.
% END GFI12

% BEGIN GFI13
\paragraph{GFI13: the literal \texttt{Chain\allowbreak Tied} exists eventually.}
\label{claim:GFI13}
Under the hypotheses (GFI.11) and GFI6, there are constants
$K,\allowbreak{} c_{HH},\allowbreak{} e_K,\allowbreak{} C_L$ and a single threshold $x_0$ such that
\[
 x\ge x_0,\quad x\in\mathcal F.\mathrm{Win}
 \quad\Longrightarrow\quad
 \texttt{Chain\allowbreak Tied}(\mathcal F,\allowbreak{} C_{\rm abs},\allowbreak{} e_L,
                              K,\allowbreak{} c_{HH},\allowbreak{} e_K,\allowbreak{} C_L,\allowbreak{} x).               \tag{GFI.26}
\]
Indeed choose the constants and threshold as in GFI6.  For each such
$x$ use the common model of GFI2.  Its original frame, rank and
word sets, top type $T_n$, constants (GFI.13), ledger (GFI.7)--(GFI.9),
sets (GFI.3), equivalences $e_T,\allowbreak{} e_P$, character functional and cutoff
from GFI5 are the witnesses in \texttt{Chain\allowbreak Tied}.

The four tie predicates are GFI4--GFI5.  The five scalar/constant
identities and inequalities are definitions: $\mathscr D.x=x$,
$\mathscr D.\epsilon=\epsilon$, $e_D=C_{\rm abs}$,
$e_P=e_L$ and $c_{HH}=C_H$.  The \texttt{Ident} record is GFI5.
The eleven ledger fields, in their source order, are GFI6;
the three results of GFI7; the three results of GFI8; GFI9;
and the three results of GFI10.  GFI11 supplies
\texttt{Err\allowbreak Windows}; GFI12 supplies \texttt{BNBounds}, the
remaining growth hypotheses, \texttt{Cost\allowbreak Windows}, and
\texttt{Rest\allowbreak Size}.  This checks every field of the original
existential record and proves (GFI.26).

In particular, the constants are chosen before $x$, although the
realized degree pair, base index and error bounds may depend on $x$
as the source permits.  A scalar bound established only after
varying the base point would not give this conclusion.
% END GFI13

% BEGIN GFI14
\paragraph{GFI14: the original chosen inputs supply the small-share inputs.}
\label{claim:GFI14}
Suppose $\texttt{Chosen\allowbreak E}(\mathcal F,\allowbreak{} C_0,\allowbreak{} C_H)$ and the same
family's \texttt{Chosen\allowbreak Wall} are supplied.  Unpack the original
positive constants $C'_5,\allowbreak{} c_6,\allowbreak{} C_b$.  The top, quadratic, thin,
proviso and code predicates are already present.  We verify the
three other fields of \texttt{Weight\allowbreak Vector\allowbreak Small\allowbreak Shares\allowbreak Seven}.

First the original second-set lower bound is
\[
 \#T_2^{\rm raw}\ge
 \frac{n_0}{L_w}
 \exp\!\left[-C_{T2}\sqrt{\log n_0\log L_w}\right],
 \qquad n=o n_0.                                      \tag{GFI.27}
\]
The family fixes $C_{T2}$, bounds $o\le C_o$, and has
$1\le n_0\le n\le x^{9/4}$ and $2\le L_w\le(\log x)^2$
eventually, by its original size record.  Set
$C=\max(C_{T2},0)$ and $O=\max(C_o,1)$.  For $x\ge e$,
the logarithm of the inverse lower density in (GFI.27) is at most
\[
 \log O+2\log\log x+
       C\sqrt{\tfrac92\log x\log\log x}.              \tag{GFI.28}
\]
Choose a single threshold making each of the three terms at most
$\epsilon\log x/3$.  This is possible because
$1/\log x\to0$ and $\log\log x/\log x\to0$; for the square-root
term the required condition is
$(9/2)C^2\log\log x/\log x\le\epsilon^2/9$.
Dividing (GFI.27) by $n$ and exponentiating gives
$x^{-\epsilon}\le\#T_2^{\rm raw}/n$ uniformly.  This is exactly
\texttt{Thick\allowbreak Set\allowbreak Density\allowbreak Eventually}.

Next the source \texttt{Chosen\allowbreak Bad} gives its original bad-top sum
at most
$4C_bx^{-2}\delta_1n\sum_{p\in\mathcal P}\log p/p$, where
$\mathcal P$ is the original large-order residue-prime set.
Since $\delta_1n\le2n_T$, it suffices to bound that prime sum.
An elementary bound is enough.  For $X\ge e$, putting
$M=\lceil X\rceil$ and using $\sum_{r=1}^M1/r\le1+\log M$ gives
\[
 \sum_{\substack{p<M\\p\ {
m prime}}}\frac{\log p}{p}
 \le \log M(1+\log M)\le8(\log X)^2.                  \tag{GFI.29}
\]
Here $M\le2X$, so $\log M\le2\log X$ and
$1+\log M\le3\log X$.  Apply (GFI.29) to $X=X'$.
The bad-top sum is at most
$64C_bn_Tx^{-2+2\epsilon/3}$.  After imposing
$64C_b\le x^{\epsilon/3}$, this is at most
$n_Tx^{-2+\epsilon}$, the literal \texttt{Chosen\allowbreak Badp} bound.
No distribution theorem for primes is needed for this conversion.

Finally, at a late allowed scale choose an existing index $i_0$.
Its actual degrees define integer functions on $T_n$, extended by
zero elsewhere; injectivity of the labels makes them well-defined.
The original top window and $b_T\le b_{\rm top}$ give admissibility.
The same $i_0$ is a realizing index.  Every other carrier satisfies
\texttt{Ch\allowbreak W\_\allowbreak at} by \texttt{Chosen\allowbreak Wall}.  This is
\texttt{Chosen\allowbreak WS}, also as in HW3.  Combining these fields with the
unchanged five chosen inputs proves the entire original Seven record.
Its constants and thresholds retain the required order.
% END GFI14

% BEGIN GFI15
\paragraph{GFI15: from the eleven original family fields to \texttt{Link\allowbreak GLOBAL}.}
\label{claim:GFI15}
Fix $C_{\rm abs}\ge0$, $C_H>0$, $C_0\ge3$, and suppose the
eleven original \texttt{Fam\allowbreak GLmin} fields with upper-column constant
$4$ hold for this family.  Supply \texttt{Tu\allowbreak Size},
$\texttt{Weight\allowbreak Vector\allowbreak Small\allowbreak Shares\allowbreak Seven}(\mathcal F,\allowbreak{} C_0,\allowbreak{} C_H)$ and
\[
 \texttt{Link\allowbreak G},\quad\texttt{Link\allowbreak S}(C_0),\quad
 \texttt{Link\allowbreak IRR},\quad\texttt{Link\allowbreak DIV}(C_{\rm abs},\allowbreak{} C_H),\quad
 \texttt{Link\allowbreak Mu}.                                       \tag{GFI.30}
\]
Then the literal \texttt{Link\allowbreak GLOBAL7} conclusion holds with
$e_D=C_{\rm abs}$ and $e_L=0$.

Here is the full dependency installation.  The original guarded
local fields are activated by \texttt{Tu\allowbreak Size} through GLE2.
We first recover the three ramified-model conclusions directly
from the supplied $h_W,\allowbreak{} h_\mu$; the stronger second-layer hypotheses
of the RSL construction are not required for this generic consumer.
Fix an original word, put $T=(D.\mathrm{mem}(u)).A.T$, and take
its supplied \texttt{Free\allowbreak T} record.  Its finite set $X$ has a free
action of
\[
 \mathscr G=\prod_{v\mid T}k_v^\times,
 \qquad |k_v|=T,\qquad |\{v\mid T\}|=f.
\]
Choose one representative of each orbit.  If there are $r$ orbits,
then $X$ is a disjoint union of $r$ regular $\mathscr G$-sets.
The stored equivariant linear equivalence from the original full
form space to $\mathbb C^X$ carries each original norm operator
to summation in its coordinate.  On one orbit the function space
is $\bigotimes_{v\mid T}\mathbb C^{T-1}$.  Split each factor as
the kernel of coordinate summation plus the line spanned by one
coordinate vector of sum one.  Expanding the tensor product shows
that the joint kernel is the tensor product of those kernels,
whose dimension is $(T-2)^f$.  Thus, with the original dimensions,
\[
 D.duS(u)=r(T-1)^f,\qquad
 d_u=(D.\mathrm{mem}(u)).A.du=r(T-2)^f,
 \qquad d_u=D.duS(u)\left(1-\frac1{T-1}\right)^f.
                                                               \tag{GFI.30a}
\]
The identifications with these dimensions use the original
\texttt{emb\_\allowbreak basis}, the actual joint-kernel definition of the
ramified forms, and the range of the original $D.\mathrm{nuIn}$.
No integral or isometric version of the free model is needed for
this dimension computation.

The supplied $h_\mu$ is already the exact mean identity
\[
 \mu(\mathrm{ramLat})=\mu(\mathrm{lat}(0))
                         -\frac{f\log(T-1)}{T-2}.
                                                               \tag{GFI.30b}
\]
The original \texttt{Chosen\allowbreak Proviso}, contained in the Seven packet,
forces $D.duS(u)>0$: if it were zero, its normalized lower bound
would read $C_0^{-1}\le0$, whereas $C_0>0$.
The setting gives $T>2fN(m)+1>2$ since $f,\allowbreak{} N(m)\ge1$.
Consequently (GFI.30a) gives $d_u>0$.
The actual ramified lattice is the intersection of the full original
lattice with its rational ramified subspace, with restricted metrics.
Every positive-rank lattice used in its maximal-slope supremum is
therefore admissible for the full lattice's supremum.  If saturation
is part of the convention, this remains true because the ramified
intersection is saturated and saturation inside it is saturated in
the full lattice.  This proves
$\mu_{\max}(\mathrm{ramLat})\le\mu_{\max}(\mathrm{lat}(0))$
by the same inclusion argument written in RSL10.  Together with
(GFI.30a)--(GFI.30b), these are all three literal
\texttt{family\_\allowbreak property\_\allowbreak and} clauses.

Next $h_I,\allowbreak{} h_R,\allowbreak{} h_B$ give the original integral and ramified operator
identities and second-branch inclusion/length.  The finite-DVR and
ordered-product arguments of GLE10--GLE11 consume exactly these
clauses and the setting, giving the full original congruence length
bound and $N_\lambda\subset S_{\rm ram}$ at every place over $T$.
RSL11 then gives the original commensurability and fullness, using
the already stored \texttt{nu\allowbreak Def} and \texttt{nprime\allowbreak Def}.
Their ranks equal $d_u$ by the original rational embedding, so
(GFI.30a) and the proviso give \texttt{J.\allowbreak Mu\allowbreak Rank} as well as the
ramified rank identity.  All thresholds used in this paragraph are
maxima of the supplied family thresholds and precede every word
and place.  In particular no additional second-layer selector
has entered this implication.

GVC1--GVC18 use the mass, volume, numerical and column fields to
give \texttt{Count\allowbreak F} and \texttt{Slopes\allowbreak MF}; their derivative
window is the existing family-size projection.  This direct route
does not assume the stronger unused \texttt{Part\allowbreak Five\allowbreak Family\allowbreak Bounds}.
GFA1--GFA5, using (GFI.30), give \texttt{Slopes\allowbreak NF}.
GFA6 gives the positive target rank and \texttt{Lat\allowbreak F}.
GFA7--GFA9 use the same Bessel map, proviso, code, local ranks,
and links to give the full \texttt{Frame\allowbreak F}$(C_H)$.
GFA10--GFA12 give \texttt{Nodes\allowbreak F}, \texttt{Win\allowbreak F} and
\texttt{Char\allowbreak F}, using the supplied thick density for the window.
LFS17, with the supplied \texttt{Chosen\allowbreak TT},
\texttt{Chosen\allowbreak Badp}, \texttt{Chosen\allowbreak WS} and \texttt{Link\allowbreak G}, gives
the original \texttt{Loss\allowbreak FS}$(0)$.

These are exactly (GFI.11).  GFI13, with the unchanged chosen
quadratic and thin constants $C'_5,\allowbreak{} c_6$, now proves
\texttt{Link\allowbreak GLOBAL7}.  Notice that $C'_5=C_L$ is the quadratic
constant; it need not equal the upper-column constant $4$.

If instead \texttt{Chosen\allowbreak E} and \texttt{Chosen\allowbreak Wall} are supplied,
GFI14 constructs the required original Seven record on this same
family.  Applying the preceding argument to every such
\texttt{Chosen\allowbreak E} and every (GFI.30) proves the literal
\texttt{Link\allowbreak GLOBAL}$(\mathcal F,\allowbreak{} C_{\rm abs},\allowbreak{} C_H,\allowbreak{} C_0,\allowbreak{} C_{\rm abs},0)$.
The retained \texttt{DGam\allowbreak Size} premise in the final exported
signature is allowed; it is used upstream when MT6 constructs
\texttt{Link\allowbreak Mu}, and is not consumed a second time here.
% END GFI15

% BEGIN GFI16
\paragraph{GFI16: the \texttt{Closed} quantifiers are preserved.}
\label{claim:GFI16}
Let $\mathcal C$ be the same original family class.  Suppose
$\texttt{Fam\allowbreak GLmin}(\mathcal C,4)$ holds and every family in
$\mathcal C$ has its original \texttt{Chosen\allowbreak Wall}.  Then
\[
 \texttt{Closed}_{\mathcal C}\left(
 \texttt{Tu\allowbreak Size}\Longrightarrow\texttt{DGam\allowbreak Size}
 \Longrightarrow
 \texttt{Link\allowbreak GLOBAL}(C_{\rm abs},\allowbreak{} C_H,\allowbreak{} C_0,\allowbreak{} C_{\rm abs},0)
 \right)                                                \tag{GFI.31}
\]
for $C_{\rm abs}\ge0$ and $C_0\ge3$.

To check the exact order, write each of the eleven \texttt{Closed}
inputs as a positive cutoff $\epsilon_j$ followed, for every
permitted $\epsilon$, by a transported tuple $\mathcal K_j$.
Take the minimum of these finitely many cutoffs and a fixed
positive cutoff below $1/52$.  For that $\epsilon$ take the
componentwise supremum of the tuples.  Its \texttt{pre} component
is a maximum, and its \texttt{t}, \texttt{mid} and \texttt{last}
components are pointwise maxima of functions of the original lists.
Thus no choice here depends on the later family or on $C_H$.
The definition of \texttt{Large} and monotonicity recover each
original largeness input for the same family and lists.

Join also the tuple with \texttt{pre}$=0$, \texttt{t}$=0$,
\texttt{mid}$=1$ and \texttt{last}$=0$.  Because $C_H$ is one
of the original middle constants, its largeness condition forces
$C_H\ge1$.  This proves the strictly positive height exponent
needed in GFI15 before any frame estimate is used.

Now take any supplied $\mathcal F\in\mathcal C$, $C_H$ satisfying
that joined largeness condition.  All eleven fields and
\texttt{Chosen\allowbreak Wall} concern this same family.  Apply GFI15 to
any supplied \texttt{Tu\allowbreak Size}, \texttt{DGam\allowbreak Size}, chosen data and
global links.  Its eventual $x$-threshold may depend on those
objects, exactly as \texttt{Link\allowbreak GLOBAL} permits, but precedes
every index and word.  This is (GFI.31) with the literal
\texttt{Closed} quantifier order.
% END GFI16

% BEGIN GFI17
\paragraph{GFI17: completion of the selected global assembly obligation.}
\label{claim:GFI17}
Retain the original class $\mathcal A$, its \texttt{Duty} on every
member, and the compatible second-layer assertion on
$\mathcal C=D_3(3,\mathcal A)$ from GLS13.  That theorem gives
the full eleven-field $\texttt{Fam\allowbreak GLmin}(\mathcal C,4)$.
Membership in $D_3(3,\mathcal A)$ projects to membership in
$\mathcal A$, and the original Duty record projects to
\texttt{Chosen\allowbreak Wall} for that same family.  GFI16, with
$C_{\rm abs}=2$ and $C_0=3$, therefore proves
\[
 \boxed{\quad
 \texttt{Closed}_{D_3(3,\mathcal A)}
  \bigl(\texttt{Tu\allowbreak Size}\to\texttt{DGam\allowbreak Size}\to
          \texttt{Link\allowbreak GLOBAL}(2,\allowbreak{} C_H,3,2,0)\bigr).
 \quad}                                                   \tag{GFI.32}
\]
This is the full remaining numerical/lattice consumer in
\texttt{TASK\_\allowbreak GL\_\allowbreak FINAL\_\allowbreak ASSEMBLY}.  All its loss and ledger
conditions now have human proofs, and $e_D=2,\allowbreak{} e_L=0$ are unchanged.
When the final common-family inputs are supplied, ZW1--ZW4 apply
to this very \texttt{Chain\allowbreak Tied} and retain the final exponent
$1/2+9\epsilon$.  No numerical absorption needs to be added
between (GFI.32) and that existing endpoint.

The conclusion is conditional on the same common construction,
compatible selector and duties as GLS13.  It supplies none of
their existence, nor the actual denominator, retained-defect,
period-height, local-gain or spectral producers.  Those remain
explicit in the final dependency audit.  In particular (GFI.32)
closes the assembly obligation without asserting a complete human
proof of either final theorem.
% END GFI17


\section{Holomorphic coefficient uniqueness on the original frame domain}
\label{sec:holomorphic-uniqueness}

This chapter reconstructs the negative-frame uniqueness argument used by
the original theta extension, together with its representation-theoretic
consumers.  All source identifications refer to commit
\texttt{7cbf7a0c32fca1badf05931879a37c3f8a2b03b4}.
The arguments are on the supplied original objects.  An existential
auxiliary frame factor is constructed where the source asks for one;
the original Hermitian forms, discrete-series models and coefficient
maps are not changed.

\subsection*{The exact finite-dimensional dictionary}

Let \(I\) be a finite set.  At each \(\tau\in I\) fix an \(n\)-by-\(n\)
complex matrix \(H_\tau\), an integer \(0\le q_\tau\le n\), and a matrix
\(x_{0,\tau}\) satisfying the source's \(\mathrm{OrthFr}\) condition
\[
 x_{0,\tau}^{*}H_\tau x_{0,\tau}=J_\tau,\qquad
 J_\tau=-S_\tau,\qquad
 S_\tau=\operatorname{diag}(s_{\tau,\allowbreak{} a}),\quad
 s_{\tau,\allowbreak{} a}=
 \begin{cases}1&a<q_\tau,\\-1&a\ge q_\tau.\end{cases}
 \tag{HCU.1}
\]
Indices start at zero.  For a matrix \(A\), let \(\Delta_k(A)\) be the
determinant of its leading \(k\)-by-\(k\) submatrix, with
\(\Delta_0(A)=1\), and put
\(s_{\tau,[k]}=\prod_{a<k}s_{\tau,\allowbreak{} a}\).
The actual source domain is
\[
 \mathcal D_\tau=\{x:\det x\ne0,\quad
 s_{\tau,[k]}\operatorname{Re}\Delta_k(-x^*H_\tau x)>0
 \ (0\le k\le n)\},\qquad
 \mathcal D=\prod_{\tau\in I}\mathcal D_\tau .
 \tag{HCU.2}
\]
Let \(B\) be the product of the invertible upper-triangular groups,
\(G=\prod_\tau U(H_\tau)\), and
\[
 \chi_w(b)=\prod_{\tau,\allowbreak{} a}b_{\tau,\allowbreak{} aa}^{\,\allowbreak{} w_{\tau,\allowbreak{} a}},
 \quad w_{\tau,\allowbreak{} a}\in\mathbb Z,\qquad
 D(s)_\tau=\operatorname{diag}(e^{is_{\tau,\allowbreak{} a}}),\quad
 t(s)=x_0D(s)x_0^{-1}.
 \tag{HCU.3}
\]
Here and below \(s\) in \(D(s)\) is a tuple of real angles, distinct from
the fixed signs in (HCU.1).  Write
\(\chi_w(s)=\prod_{\tau,\allowbreak{} a}e^{iw_{\tau,\allowbreak{} a}s_{\tau,\allowbreak{} a}}\).
The source's \(\mathrm{HolFn}(H,\allowbreak{} q,\allowbreak{} x_0,\allowbreak{} w,\allowbreak{} f)\) means exactly that there
is a holomorphic function \(P\) on \(\mathcal D\) such that
\[
 P(xb)=\chi_w(b)P(x),\qquad f(u)=P(ux_0)\quad(u\in G).
 \tag{HCU.4}
\]
Holomorphy means complex differentiability on this open subset of the
finite-dimensional matrix space.  No extension to all matrices is
required.  Empty products and zero-dimensional matrix spaces are allowed.

% BEGIN HCU1
\paragraph{HCU1: a smooth triangular factor on the original domain.}
\label{claim:HCU1}
The matrices \(H_\tau,\allowbreak{} x_{0,\tau}\) are nonsingular, and \(H_\tau\) is
Hermitian.  There is a real-smooth map
\(\beta_\tau:\mathcal D_\tau\to B_n\) such that
\[
 u_\tau(x)=x\beta_\tau(x)^{-1}x_{0,\tau}^{-1}\in U(H_\tau),
 \qquad x=u_\tau(x)x_{0,\tau}\beta_\tau(x).
 \tag{HCU.5}
\]
This is a witness for the literal
\texttt{smooth\_\allowbreak gram\allowbreak Schmidt} existential statement.

\emph{Proof.}
Taking determinants in (HCU.1) proves nonsingularity, and solving that
equation for \(H_\tau\) proves Hermitian symmetry.  Omit \(\tau\).
For \(A=-x^*Hx\), all leading minors are nonzero.  Gaussian elimination
without row exchanges gives
\[
 A=L\operatorname{diag}(d_0,\ldots,\allowbreak{} d_{n-1})L^*,
 \quad L\text{ lower unitriangular},\qquad
 d_a=\frac{\Delta_{a+1}(A)}{\Delta_a(A)}\in\mathbb R^\times .
 \tag{HCU.6}
\]
For completeness, isolate the nonzero first pivot \(d=A_{00}\).
Subtracting the outer product of the first column and row, divided by
\(d\), leaves its Hermitian Schur complement.  Its leading \(k\)-minor
is \(\Delta_{k+1}(A)/d\).  Induction constructs (HCU.6), and comparison
of leading determinants gives its displayed pivot formula.  Every
entry of \(L\) and every \(d_a\) is a rational expression in entries of
\(A\), with the already nonzero pivots as denominators.

The signs in (HCU.2) give \(s_a d_a>0\).  Set
\[
 \beta(x)=\operatorname{diag}(\sqrt{|d_a|})\,\allowbreak{} L^* .
 \tag{HCU.7}
\]
Then \(\beta\) is upper triangular with positive diagonal, and
\(A=\beta^*S\beta\).  Consequently
\[
 (x\beta^{-1})^*H(x\beta^{-1})=-S=x_0^*Hx_0,
\]
which proves (HCU.5).  The rational formulas, complex conjugation and square root on
the positive real axis show real smoothness on the entire domain.
For \(n=0\) all matrices are the unique empty matrix and all formulas
have their empty-product meaning.  The source requires an existential
factor with (HCU.5); it does not prescribe the particular diagonal
sign convention of that factor.  Thus (HCU.7) changes no stored datum.
% END HCU1

% BEGIN HCU2
\paragraph{HCU2: openness, stability and connectedness of the actual domain.}
\label{claim:HCU2}
The set \(\mathcal D\) is open, nonempty and path connected.  It contains
\(x_0\) and is stable under \(x\mapsto uxb\), \(u\in G,\allowbreak{} b\in B\).
In particular it has exactly the properties used in
\texttt{fr\allowbreak Doms\_\allowbreak open\_\allowbreak preconnected\_\allowbreak stable}.

\emph{Proof.}
The finitely many determinant inequalities in (HCU.2) are open
conditions.  For upper-triangular \(b\), its first \(k\) columns involve
only the first \(k\) rows, so
\[
 \Delta_k\bigl(b^*(-x^*Hx)b\bigr)
   =|\det b_{[k]}|^2\Delta_k(-x^*Hx).
 \tag{HCU.8}
\]
The multiplying scalar is positive.  Left multiplication by \(U(H)\)
does not change the Gram matrix.  These observations prove stability;
(HCU.1) proves \(x_0\in\mathcal D\).  HCU1 gives the surjection
\((u,\allowbreak{} b)\mapsto ux_0b\) from \(G\times B\) to \(\mathcal D\).

Here is a direct path-connectedness proof for its two factors.
Conjugation by \(x_0\) identifies \(U(H)\) with \(U(J)\).
For \(g\in U(J)\), put \(A=g^*g\), \(p=A^{1/2}\), \(k=gp^{-1}\).
The positive square root is defined by unitary diagonalization of the
positive Hermitian matrix \(A\).  From \(g^{-1}=Jg^*J\) one obtains
\(JAJ=A^{-1}\), hence \(JpJ=p^{-1}\).  Thus \(k\) is both ordinarily
unitary and \(J\)-unitary, and therefore commutes with \(J\).
Choose a skew-Hermitian logarithm \(Y\) of \(k\), separately on its
two \(J\)-eigenspaces, and let \(X=\log p\), defined by the same
positive spectral calculus.  Then
\[
 JXJ=-X,\qquad
 [Y,\allowbreak{} J]=0,\qquad
 [0,1]\ni t\longmapsto e^{tY}e^{tX}\in U(J)
 \tag{HCU.9}
\]
joins \(1\) to \(g\).  These assertions follow on each eigenspace:
positive eigenvalues have real logarithms, unit eigenvalues have
imaginary logarithms, and conjugation by \(J\) takes the positive
spectral calculus of \(A\) to that of \(A^{-1}\).  Only the existence
of a path for each fixed \(g\) is used; no continuous global choice of
logarithms is being asserted.

An upper-triangular invertible matrix has the form \(d(1+N)\), with
\(d\) diagonal invertible and \(N\) strictly upper triangular.
For fixed logarithms \(\ell_a\) of its diagonal entries,
\(\operatorname{diag}(e^{t\ell_a})(1+tN)\) is a path from \(1\) to it.
Finite products and continuous images preserve path connectedness.
This proves the claim, also when a signature block, \(I\), or \(n\)
is empty.
% END HCU2

% BEGIN HCU3
\paragraph{HCU3: the local zero lemma with its exact hypotheses.}
\label{claim:HCU3}
Let \(\Omega\) be a neighborhood of the tuple \(1\) in
\(\prod_I M_n(\mathbb C)\), and let \(Q\) be a function on that
ambient space which is complex differentiable on \(\Omega\).
Suppose:
\begin{enumerate}
\item if \(y\in\Omega\), \(Q(y)=0\) and \(b\in B\), then \(Q(yb)=0\);
\item whenever \(y,\allowbreak{} y'\in\Omega\) and \(y'D(s)=D(s)y\), one has
\(Q(y')=Q(y)\);
\item \(Q(1)=0\).
\end{enumerate}
Then \(Q\) vanishes on a neighborhood of \(1\).  Only its restriction
to a small open ball inside \(\Omega\) is used analytically;
the first hypothesis is the literal ambient zero-set condition.

\emph{Proof.}
Use the maximum-entry norm and choose \(r>0\) with
\(B(1,\allowbreak{} r)\subset\Omega\).  For a tuple \(z\) put
\[
 L_\zeta(z)_{\tau,\allowbreak{} ac}=
 \begin{cases}\zeta^{a-c}z_{\tau,\allowbreak{} ac}&c<a,\\0&c\ge a,\end{cases}
 \qquad
 U(z)_{\tau,\allowbreak{} ac}=
 \begin{cases}z_{\tau,\allowbreak{} ac}&a\le c,\\0&a>c.\end{cases}
 \tag{HCU.10}
\]
If \(\|z\|<r\) and \(|\zeta|\le1\), then
\(\|L_\zeta(z)\|\le\|z\|\).
Thus \(h(\zeta)=Q(1+L_\zeta(z))\) is holomorphic on a neighborhood of
the closed unit disc.  For \(\zeta=e^{it}\),
\(D_\zeta=\operatorname{diag}(\zeta^a)\) gives
\[
 (1+L_\zeta(z))D_\zeta=D_\zeta(1+L_1(z)).
 \tag{HCU.11}
\]
Hypothesis 2 makes \(h\) constant on the unit circle.
The maximum-modulus principle applied to \(h-h(1)\), or the Cauchy
integral formula, makes it constant throughout the disc.
Since all exponents \(a-c\) in the lower part are positive,
\(L_0(z)=0\), and \(h(1)=h(0)=Q(1)=0\).

Now the polynomial map
\[
 \mathcal L(z)=(1+L_1(z))(1+U(z))
       =1+z+L_1(z)U(z)
 \tag{HCU.12}
\]
has derivative the identity at zero.  The finite-dimensional inverse
function theorem therefore maps a sufficiently small neighborhood of
zero onto a neighborhood of \(1\).  Shrink it to \(\|z\|<\min(r,1)\).
The diagonal entries \(1+z_{\tau,\allowbreak{} aa}\) are nonzero, so \(1+U(z)\in B\).
The first part of the proof and hypothesis 1 give
\(Q(\mathcal L(z))=0\) throughout that neighborhood.  This proves
exactly \texttt{Uniq.\allowbreak local\_\allowbreak zero}.  In a zero-dimensional ambient space,
the conclusion is already \(Q(1)=0\).
% END HCU3

% BEGIN HCU4
\paragraph{HCU4: global uniqueness with the original torus law.}
\label{claim:HCU4}
Assume \(\mathrm{HolFn}(H,\allowbreak{} q,\allowbreak{} x_0,\allowbreak{} w,\allowbreak{} f)\), together with the separate
left torus law
\[
 f(t(s)u)=\chi_w(s)f(u)\quad(u\in G),\qquad f(1)=0.
 \tag{HCU.13}
\]
Then \(f=0\) on \(G\).  This is the full original
\texttt{Uniq.\allowbreak eq\_\allowbreak zero\_\allowbreak of\_\allowbreak torus\_\allowbreak law} contract.

\emph{Proof.}
Let \(P\) be its section in (HCU.4).  For \(x=ux_0b\), HCU1 and
(HCU.13) give
\[
 P(t(s)x)=\chi_w(b)f(t(s)u)=\chi_w(s)P(x).
 \tag{HCU.14}
\]
Consider \(Q(y)=P(x_0y)\) on
\(\Omega=\{y:x_0y\in\mathcal D\}\), an open neighborhood of \(1\).
Its zero set is stable under right multiplication by \(B\).
If \(y'=D(s)yD(s)^{-1}\), then
\[
 Q(y')=P(t(s)x_0yD(s)^{-1})
      =\chi_w(s)\chi_w(D(s)^{-1})Q(y)=Q(y).
 \tag{HCU.15}
\]
The integer powers are well defined and the two characters cancel
exactly.  Also \(Q(1)=f(1)=0\).  HCU3 makes \(P\) vanish on an open
neighborhood of \(x_0\).

A holomorphic function on a connected open set that vanishes on a
nonempty open subset vanishes everywhere.  To see the needed
several-variable statement directly, cover a path by overlapping balls.
On a ball, restrict to a complex line through a point of an open zero
set; the one-variable power-series identity theorem gives vanishing
along that line wherever the line stays in the ball.  It follows that
the whole ball is zero, and finitely many overlapping balls propagate
the conclusion along the path.  HCU2 supplies such paths inside the
actual domain.  Therefore \(P=0\) on \(\mathcal D\), whence \(f=0\)
on \(G\).
% END HCU4

% BEGIN HCU5
\paragraph{HCU5: the full smooth Borel--holomorphic dictionary.}
\label{claim:HCU5}
Let \(\mathfrak g_\tau=\{X:X^*H_\tau+H_\tau X=0\}\).
Use the source right derivatives \(R_X\) and their complexification
\[
 c_H(Z)=-H^{-1}Z^*H,\qquad
 R_Z^{\mathbb C}
 =R_{(Z+c_H(Z))/2}
    +iR_{(Z-c_H(Z))/(2i)}.
 \tag{HCU.16}
\]
A function \(f\) on \(G\) has \(\mathrm{HolFn}(H,\allowbreak{} q,\allowbreak{} x_0,\allowbreak{} w,\allowbreak{} f)\) if and
only if it has the source \(\mathrm{SmFn}\) smoothness and
\[
 R^{\mathbb C}_{x_{0,\tau}E_{ac}x_{0,\tau}^{-1}}f(u)
   =\begin{cases}w_{\tau,\allowbreak{} a}f(u)&a=c,\\0&a<c\end{cases}
 \quad(a\le c,\ u\in G).
 \tag{HCU.17}
\]
Thus both \texttt{hol\allowbreak Fn\_\allowbreak of\_\allowbreak killed} and
\texttt{sm\allowbreak Fn\_\allowbreak and\_\allowbreak killed\allowbreak Fn\_\allowbreak of\_\allowbreak hol\allowbreak Fn} hold on the literal source domain.

\emph{Proof.}
First, the source smoothness condition, which tests finite products of
real one-parameter exponentials at every \(u\), is ordinary smoothness
on \(G\).  Here is the chart verification.  On the \(H\)-skew matrices,
the Cayley map
\((1+X/2)(1-X/2)^{-1}\) is a smooth chart into \(U(H)\) near \(1\);
its inverse is \(2(g-1)(g+1)^{-1}\), and its derivative at zero is the
identity.  Given a real basis \(X_1,\ldots,\allowbreak{} X_d\) of \(\mathfrak g\),
the product \(\prod_j e^{t_jX_j}\) has the same derivative in this chart.
The inverse function theorem makes these product exponentials local
coordinates.  The source tests their left translates, so it supplies
all smooth coordinate restrictions.  Conversely a smooth group function
satisfies every such test.

The right differential of a smooth function is real-linear in \(X\).
The involution \(c_H\) has fixed space \(\mathfrak g\), and every
complex matrix is uniquely \(X+iY\) with \(X,\allowbreak{} Y\in\mathfrak g\).
Consequently (HCU.16) is its complex-linear extension.
For a fixed real angle tuple \(s\), equation (HCU.17) on diagonal
units gives the scalar ordinary differential equation along
\(u\,\allowbreak{} t(ts)\), hence
\[
 f(u\,\allowbreak{} t(s))=\chi_w(s)f(u).
 \tag{HCU.18}
\]
The integer weights ensure that this is the prescribed character of
the diagonal torus.

Use HCU1 to define
\(P(x)=\chi_w(\beta(x))f(u(x))\).
This is real smooth.  It has the stronger formula
\[
 P(ux_0b)=\chi_w(b)f(u)\qquad(u\in G,\ b\in B).
 \tag{HCU.19}
\]
Indeed, two such decompositions differ by an element of
\(U(J_\tau)\cap B_n\).  An upper-triangular \(J\)-unitary matrix is
diagonal with entries of modulus one: its first column has only its
first entry, its norm determines that entry's modulus, and its
orthogonality to the other columns kills the remainder of the first
row; induction finishes.  Thus the two factors differ by \(D(s)\),
and (HCU.18) makes their two expressions in (HCU.19) identical.

It remains to prove complex differentiability; real smoothness alone
is insufficient.  At \(x=ux_0b\), write a tangent as
\(\dot x=ux_0Zb\).  Every complex matrix \(Z\), considered as a real
vector, is \(X+B_0\), where \(X\in\mathfrak u(J)\) and \(B_0\)
is upper triangular.  An explicit choice sets
\(X_{ac}=Z_{ac}\) for \(a>c\),
\(X_{ca}=-J_{aa}J_{cc}\overline{Z_{ac}}\), and \(X_{aa}=0\);
then \(B_0=Z-X\) is upper triangular.  Differentiating the factorized
curve
\(u e^{t x_0Xx_0^{-1}}x_0 e^{tB_0}b\) in (HCU.19) gives
\[
 dP_x(ux_0Zb)=\chi_w(b)
 \left(R_{x_0Xx_0^{-1}}f(u)
            +\sum_{\tau,\allowbreak{} a}w_{\tau,\allowbreak{} a}(B_0)_{\tau,\allowbreak{} aa}f(u)\right)
 =\chi_w(b)R^{\mathbb C}_{x_0Zx_0^{-1}}f(u).
 \tag{HCU.20}
\]
The last equality uses (HCU.17) for the whole upper-triangular
\(B_0\), by complex linearity.  Its right side is complex-linear in
\(Z\); multiplication by the fixed invertible frames identifies \(Z\)
complex-linearly with \(\dot x\).  Therefore the real derivative of
\(P\) is complex-linear at every point.  This is complex
differentiability, and proves the forward construction of (HCU.4).

Conversely let \(P\) satisfy (HCU.4).
The restrictions \(f(u)=P(ux_0)\) are smooth.
At \(ux_0\), its complex derivative in the tangent
\(ux_0E_{ac}\) is obtained from the curve \(ux_0e^{zE_{ac}}\).
For \(a<c\), its upper-triangular character is identically one; for
\(a=c\), it is \(e^{zw_{\tau,\allowbreak{} a}}\).
Differentiating gives (HCU.17), with precisely the complexification
(HCU.16).  This also proves (HCU.18) directly from (HCU.4).
% END HCU5

% BEGIN HCU6
\paragraph{HCU6: bounded integration preserves the original holomorphic type.}
\label{claim:HCU6}
Let \((\Omega,\nu)\) be a finite measure space, and suppose
\(f_o:G\to\mathbb C\) has the same \(\mathrm{HolFn}(H,\allowbreak{} q,\allowbreak{} x_0,\allowbreak{} w)\)
for every \(o\).  Suppose
\(|f_o(u)|\le C\) for all \(o,\allowbreak{} u\), and \(o\mapsto f_o(u)\) is
almost-everywhere strongly measurable for every fixed \(u\).
Then
\[
 u\longmapsto\int_\Omega f_o(u)\,\allowbreak{} d\nu(o)
 \tag{HCU.21}
\]
has the same \(\mathrm{HolFn}\).  This is the full
\texttt{hol\allowbreak Fn\_\allowbreak integral} statement.

\emph{Proof.}
Choose any section \(P_o\) for \(f_o\).  Its right transformation law
and HCU1 force the pointwise identity
\[
 P_o(x)=\chi_w(\beta(x))f_o(u(x)).
 \tag{HCU.22}
\]
In particular every section value is measurable in \(o\), regardless
of how the sections were chosen.  Near each \(x\in\mathcal D\),
the continuous factor \(\chi_w(\beta(y))\) is bounded: the diagonal
entries never vanish, including when some weights are negative.
Thus on a sufficiently small closed polydisc in \(\mathcal D\)
all \(P_o\) are bounded by a common finite constant.  Its constant
function on \(\Omega\) is integrable.

For clarity, holomorphy under this integral follows directly from
Cauchy's formula on a slightly larger polydisc.  Its coefficient
bounds give a uniformly summable majorant for the power series on
each smaller polydisc.  Integrate each coefficient; domination
justifies exchange of integral and series and gives a convergent
holomorphic power series for \(\int P_o\).
Measurability on the boundary products used in Cauchy's formula
follows by approximating their continuous parameter dependence with
finite grids: separate measurability in \(o\) and continuity in the
finite-dimensional compact parameter give measurable step
approximations with pointwise limits.  Hence Fubini is applicable.
The same argument with \(\max(C,0)\) covers the zero-measure or empty
index cases.  The right character law passes through the integral,
and evaluation at \(ux_0\) gives (HCU.21), proving the assertion.
% END HCU6

% BEGIN HCU7
\paragraph{HCU7: proportionality and a cyclic coefficient test.}
\label{claim:HCU7}
Two functions of the same \(\mathrm{HolFn}\) type, both satisfying
the left torus law in (HCU.13), and with values \(c,1\) at \(1\),
satisfy \(f=cg\) on \(G\).

For a single place, let \(\sigma\) be a strongly continuous unitary
irreducible representation of \(U(H)\) on a complex Hilbert space.
Let \(v\ne0\) be smooth and satisfy the original Borel equations
\[
 d\sigma^{\mathbb C}(x_0E_{ac}x_0^{-1})v
   =\delta_{ac}w_av\quad(a\le c).
 \tag{HCU.23}
\]
If \(l\) has torus character \(\chi_w\) and \(\langle l,\allowbreak{} v\rangle=0\),
then \(l=0\).  This closes the entire
\texttt{Hol\allowbreak Low\allowbreak Once.\allowbreak }\allowbreak\texttt{eq\_\allowbreak zero\_\allowbreak of\_\allowbreak orth} implication.

\emph{Proof.}
The difference \(f-cg\) satisfies HCU4; this proves proportionality.
For the second assertion, use the convention that the Hilbert inner
product is conjugate-linear in its first argument.  The smooth
coefficient \(F(u)=\langle l,\sigma(u)v\rangle\) satisfies (HCU.17)
by differentiating in the second argument, so HCU5 supplies its
literal \(\mathrm{HolFn}\).  Since \(\sigma(t(s))l=\chi_w(s)l\)
and \(|\chi_w(s)|=1\), unitarity gives
\[
 \langle l,\sigma(t(s)u)v\rangle
       =\chi_w(s)\langle l,\sigma(u)v\rangle .
 \tag{HCU.24}
\]
HCU4 makes this coefficient zero.  The closed span of
\(\{\sigma(u)v:u\in U(H)\}\) is a nonzero closed invariant subspace,
hence the whole Hilbert space.  Therefore \(l\) is orthogonal to
the whole space and is zero.  No assertion that all lowest compact
types were already irreducible was used in this test.
% END HCU7

\subsection*{The full original lowest compact type}

The next arguments concern the actual representation, not a replacement
model.  For a sign pattern \(\epsilon\), let
\(J_\epsilon=\operatorname{diag}(1_{\epsilon(a)}-1_{\neg\epsilon(a)})\),
\(K_\epsilon=U(n)\cap U(J_\epsilon)\), and write \(E_{ac}\) for matrix
units.  The source's \(\mathrm{IsKTypeOf}(\epsilon,\sigma,\Lambda,\allowbreak{} S)\)
says that \(S\) is finite dimensional and \(K_\epsilon\)-stable,
irreducible under this compact group, contains a nonzero torus vector
of weight \(\Lambda\), and every weight \(\mu\) in \(S\) satisfies
\[
 \Lambda-\mu
   =\sum_{\substack{a<c\\\epsilon(a)=\epsilon(c)}}m_{ac}(e_a-e_c),
 \qquad m_{ac}\in\mathbb N .
 \tag{HCU.25}
\]
This is the literal relation \(\mathrm{CDom}(\epsilon,\mu,\Lambda)\).
For \(\Lambda=\mathrm{lowK}(\epsilon,\lambda)\), the original
\(\mathrm{lowKSub}\) is the algebraic supremum of \emph{all} such
subspaces \(S\).  It is not initially assumed finite dimensional or
irreducible.

% BEGIN HCU8
\paragraph{HCU8: smooth vectors, derived actions and a smooth compact-type seed.}
\label{claim:HCU8}
Let \(\sigma\) be a strongly continuous unitary representation of \(U(H)\)
on a complete complex Hilbert space.  Its source smooth vectors form a
complex vector subspace, are preserved by the group and by all
complexified derived operators, and satisfy on that subspace
\[
 \begin{split}
 d\sigma^{\mathbb C}([Z,\allowbreak{} W])&=
     [d\sigma^{\mathbb C}(Z),\allowbreak{} d\sigma^{\mathbb C}(W)],\\
 \sigma(g)d\sigma^{\mathbb C}(Z)v
   &=d\sigma^{\mathbb C}(gZg^{-1})\sigma(g)v,\\
 \langle d\sigma^{\mathbb C}(Z)v,\allowbreak{} w\rangle
   &=-\langle v,\allowbreak{} d\sigma^{\mathbb C}(c_H(Z))w\rangle .
 \end{split}
 \tag{HCU.26}
\]
Every finite-dimensional irreducible subspace for a compact subgroup
\(K\) has an isomorphic \(K\)-subspace consisting of smooth vectors.
For \(K=K_\epsilon\), this preserves its exact
\(\mathrm{IsKTypeOf}\) weight.  In particular every original DSModel
has a smooth compact-type seed of its specified lowest weight.

\emph{Proof.}
The product-exponential charts established in HCU5 apply equally to
Hilbert-valued functions.  A smooth orbit map remains smooth after
addition, scalar multiplication and translation.  If \(v\) is smooth,
the function \((g,\allowbreak{} t)\mapsto\sigma(g e^{tX})v\) is smooth.
Differentiating in \(t\) at zero gives the smooth orbit map of
\(d\sigma(X)v\).  This proves stability under differentiation without
an assumption that an unbounded derived operator is continuous on the
whole Hilbert space.  The ordinary chain rule in these charts proves
real linearity in \(X\) and complex linearity in \(v\).  Complexification
gives linearity in every \(Z\).

Differentiate
\(g e^{tX}g^{-1}=e^{t gXg^{-1}}\) for covariance.
The mixed derivatives of
\(e^{sX}e^{tY}e^{-sX}e^{-tY}=1+st[X,\allowbreak{} Y]+o(st)\), or the
corresponding right-invariant vector fields in a chart, give the
commutator identity.  All derivatives here act on the same smooth
orbit map, so their interchange is legitimate.
Differentiating
\(\langle\sigma(e^{tX})v,\sigma(e^{tX})w\rangle=\langle v,\allowbreak{} w\rangle\)
gives skew adjointness for real \(X\); decomposing \(Z=X+iY\)
and remembering conjugate linearity in the first argument gives the
last formula in (HCU.26).

Here is the smoothing argument with its exact compact type preserved.
Let \(S\) be the given finite-dimensional irreducible \(K\)-subspace
and \(v\in S\setminus\{0\}\).
The real part of \(\langle v,\sigma(g)v\rangle\) is positive on a
neighborhood of \(1\).  Compactness of \(K\) gives a smaller neighborhood
all of whose \(K\)-conjugates lie in it: continuity of conjugation on
\(K\times\{1\}\) and a finite cover of \(K\) prove this assertion.
Choose a nonnegative smooth bump supported in that smaller neighborhood,
positive at \(1\), and average it under compact conjugation to obtain
\(a\).  The average is smooth, compactly supported, nonnegative and
\(K\)-conjugation invariant, and its support remains in the positivity
neighborhood.  These bumps exist directly in a finite-dimensional
real chart by the usual function
\(e^{-1/(1-\|x\|^2)}\) inside a ball, extended by zero.

Conjugation by \(K\) preserves left Haar measure on \(U(H)\).
Indeed its scaling factor is a continuous homomorphism
\(K\to\mathbb R_{>0}\), whose compact image must be \(\{1\}\);
taking logarithms proves the last assertion.
Define
\[
 A v'=\int_{U(H)}a(g)\sigma(g)v'\,\allowbreak{} dg\quad(v'\in S).
 \tag{HCU.27}
\]
Strong continuity and compact support give its Hilbert-valued integral.
The real part of \(\langle v,\allowbreak{} Av\rangle\) is positive, so \(A\ne0\).
Changing variables by compact conjugation shows
\(A\sigma(k)=\sigma(k)A\).  Its kernel in \(S\) is invariant;
irreducibility therefore makes \(A\) injective.  Its image is finite
dimensional and isomorphic to \(S\) as a compact representation, so all
parts of \(\mathrm{IsKTypeOf}\), including (HCU.25), are unchanged.

Finally
\(\sigma(x)Av'=\int a(x^{-1}g)\sigma(g)v'\,\allowbreak{} dg\).
On a compact neighborhood of any \(x\), this integrand and all its
coordinate derivatives have a common compact support in \(g\);
their norms are bounded by a scalar smooth-bump derivative times
\(\|v'\|\).  Dominated differentiation of the Bochner integral gives
a smooth orbit map for \(Av'\).  The argument is unchanged after
conjugating the matrix group by the original \(M.\mathrm{fr}\).
Since the source \(M.\mathrm{sm}\) is the span of smooth vectors,
the already proved linear-subspace property identifies it with the
smooth vectors themselves.  This proves the complete seed contract
\texttt{exists\_\allowbreak low\allowbreak K\_\allowbreak le\_\allowbreak sm}.
% END HCU8

% BEGIN HCU9
\paragraph{HCU9: the finite compact-weight facts actually needed.}
\label{claim:HCU9}
The following facts hold for a unitary continuous representation of
\(K_\epsilon\).
\begin{enumerate}
\item Every nonzero finite-dimensional invariant subspace contains an
\(\mathrm{IsKTypeOf}(\epsilon,\sigma,\mu,\allowbreak{} S')\).
Every such highest weight \(\mu\) is integral and weakly decreasing
along each increasing-index sign block.
\item If \(L\) is the supremum of all compact subspaces of fixed type
\(\Lambda\), then every vector of \(L\) belongs to a finite-dimensional
compact-stable subspace of \(L\).
If \(y\) has torus weight \(\nu\) and
\(\langle w,\allowbreak{} y\rangle\ne0\) for some \(w\in L\), then
\(\mathrm{CDom}(\epsilon,\nu,\Lambda)\).
\item If a smooth \(v\in L\) has weight \(\Lambda\), then
\(d\sigma^{\mathbb C}(E_{ac})v=0\) for
\(a<c,\epsilon(a)=\epsilon(c)\), and
\(d\sigma^{\mathbb C}(E_{aa})v=\Lambda_a v\).
\end{enumerate}

\emph{Proof.}
A nonzero invariant subspace of smallest positive dimension is
irreducible.  Commuting unitary operators of the diagonal torus on it
are simultaneously diagonalizable: diagonalize one nonscalar operator,
restrict the remaining commuting operators to its eigenspaces, and
induct on dimension.  The resulting simultaneous characters are
products \(e^{i\sum\mu_a s_a}\), with \(\mu_a\in\mathbb Z\).
One elementary verification is to smooth the finite representation by
the averaging argument of HCU8, now on the compact group itself; on an
irreducible finite summand its nonzero intertwiner is invertible.
Each character is therefore smooth, its differential is a real
linear form in the angles, and the ordinary differential equation
along a line gives the exponential formula.
Period \(2\pi\) in each coordinate makes every \(\mu_a\) integral.

Choose a weight maximizing \(\sum_a(-a)\mu_a\) among the finite weight
set, and a nonzero vector of that weight.  Every matrix unit \(E_{ac}\)
with \(a<c\) in the same sign block raises that quantity; hence its
complexified derivative kills this vector.
The span of words in the opposite, lower matrix units applied to it
is invariant under the whole complexified compact Lie algebra.
Here is the algebraic check: move an upper matrix unit to the right
through a word using
\([E_{ab},\allowbreak{} E_{cd}]=\delta_{bc}E_{ad}-\delta_{da}E_{cb}\);
each commutator shortens the word, and induction reorders it into
lower units, diagonal units and upper units.  Diagonal units act by
the current weight, and upper units at the far right kill the vector.
Thus all resulting weights are the chosen weight minus nonnegative
integer sums of compact positive roots.
The compact group is a product of unitary groups and is connected;
each of its elements is an exponential of a skew-Hermitian block
matrix.  A finite-dimensional Lie-stable subspace is stable under
these exponentials by the matrix exponential series.  Irreducibility
therefore makes this lower-word span the whole summand.
This proves precisely (HCU.25) and all fields of
\(\mathrm{IsKTypeOf}\), without an appeal to a classification theorem.

For any \(i<j\) in a sign block, permutation of these coordinates is
in \(K_\epsilon\).  It carries a vector of weight \(\mu\) to one of
weight \(\mu\circ(ij)\).
Apply (HCU.25) and pair its positive-root sum with the strictly
decreasing vector \((-a)_a\).  The resulting inequality is
\((\mu_i-\mu_j)(j-i)\ge0\), proving dominance.

A vector in an algebraic supremum is a finite sum of vectors from the
defining subspaces; their finite sum is the required finite invariant
space.  The orthogonal projection onto any one defining compact
subspace commutes with the compact action, since its orthogonal
complement is invariant too.  If \(\nu\) fails (HCU.25), the projection
of \(y\) onto every such subspace is zero, by its weight condition.
Then \(y\) is orthogonal to their algebraic sum \(L\), which proves
assertion 2 by contraposition.

For assertion 3, choose a finite compact-stable subspace of \(L\)
containing \(v\).  Differentiation along compact real directions
preserves it because it is closed.  Their complex span contains
\(E_{ac}\) within each sign block.  Covariance in (HCU.26) shows that
\(d\sigma^{\mathbb C}(E_{ac})v\) has weight
\(\Lambda+e_a-e_c\).  This weight cannot satisfy (HCU.25): pairing
\(-e_a+e_c\) with \((-j)_j\) is negative.  Assertion 2, applied to
the derivative vector itself, makes it zero.  Differentiation of the
diagonal torus eigencharacter proves the diagonal formula.
% END HCU9

% BEGIN HCU10
\paragraph{HCU10: the precise lowest-norm inequality.}
\label{claim:HCU10}
Suppose \(\epsilon(a)\) is true exactly when \(a<p\), with
\(0<p<n\), and \(\lambda\) is strictly decreasing.
Write the literal source weights and norm as
\[
 \Lambda=\lambda+\rho_N-\rho_C,\quad
 V(\eta)=\|\eta+2\rho_C\|^2,\quad
 \rho_N(a)+\rho_C(a)=\frac{n-1}{2}-a.
 \tag{HCU.28}
\]
The last equality follows by counting all indices before and after \(a\).
Assume \(\Lambda,\mu\) are dominant on each sign block and
\[
 \mathrm{CDom}(\epsilon,\mu+e_a-e_c,\Lambda),
 \qquad a<p\le c .
 \tag{HCU.29}
\]
Then \(V(\mu)<V(\Lambda)\).

\emph{Proof.}
Put \(L=\Lambda+2\rho_C\), \(M=\mu+2\rho_C\), and
\(\gamma=\Lambda-\mu-e_a+e_c\), a nonnegative sum of compact positive
roots.  The vectors \(\rho_C,\allowbreak{} L,\allowbreak{} M\) are weakly decreasing within each
block, by their definitions and dominance.
For \(a_0=p-1,\allowbreak{} c_0=p\), (HCU.28) gives
\[
 L_{a_0}-L_{c_0}
       =\lambda_{a_0}-\lambda_{c_0}+1>1.
 \tag{HCU.30}
\]
Every compact positive root has nonpositive \(a_0\)-coordinate and
nonnegative \(c_0\)-coordinate.  Therefore
\(\gamma_{a_0}-\gamma_{c_0}\le0\).
Since \(L-M=e_a-e_c+\gamma\), it follows that
\(M_{a_0}-M_{c_0}\ge L_{a_0}-L_{c_0}-2\).
Pairing the root expansion with the block-decreasing vector \(L+M\)
now gives
\[
 \begin{split}
 \|L\|^2-\|M\|^2
 &= (L_a+M_a)-(L_c+M_c)+\langle\gamma,\allowbreak{} L+M\rangle\\
 &\ge (L_{a_0}+M_{a_0})-(L_{c_0}+M_{c_0})\\
 &\ge 2(L_{a_0}-L_{c_0})-2>0.
 \end{split}
 \tag{HCU.31}
\]
This is the exact \texttt{Lowering.\allowbreak v\allowbreak Norm\allowbreak Sq\_\allowbreak lt} argument needed below.
For an empty sign block there is no pair \(a<p\le c\); the associated
annihilation statement is vacuous.
% END HCU10

% BEGIN HCU11
\paragraph{HCU11: lowering on the full lowest subspace.}
\label{claim:HCU11}
Suppose \(\mathrm{IsDS}(\epsilon,\lambda,\sigma)\),
\(\mathrm{IsHCParam}(\lambda)\), and the ordered sign pattern above.
Let
\(L=\mathrm{lowKSub}(\epsilon,\lambda,\sigma)\).
For every smooth \(v\in L\),
\[
 d\sigma^{\mathbb C}(E_{ca})v=0\quad(a<p\le c).
 \tag{HCU.32}
\]
This proves the original \texttt{low\allowbreak K\_\allowbreak killed}, including its use on
the whole supremum \(L\).

\emph{Proof.}
The source \(\mathrm{IsDS}\) supplies a complete unitary representation,
a compact type of weight \(\Lambda=\mathrm{lowK}(\epsilon,\lambda)\),
and the minimality condition
\[
 V(\Lambda)\le V(\mu)
 \quad\hbox{for every actual compact type of weight }\mu.
 \tag{HCU.33}
\]
Choose, by HCU9, a finite compact-stable space \(W\subset L\) containing
\(v\), and replace it by its intersection with the smooth vectors.
This intersection is a vector subspace by HCU8, is finite dimensional,
is compact-stable, and still contains \(v\).
Let \(S\) be the span of
\(d\sigma^{\mathbb C}(Z)w\), where \(w\in W\) and \(Z\) is a linear
combination of \(E_{ca}\) with \(a<p\le c\).
It is finite dimensional, being the image of the finite-dimensional
space of such \(Z\)'s tensor \(W\).
It consists of smooth vectors and is compact-stable by covariance,
since compact block conjugation preserves these off-diagonal units.

If one expression in (HCU.32) is nonzero, then \(S\ne0\).
HCU9 gives an irreducible compact subspace \(S'\subset S\) of highest
weight \(\mu\), with a nonzero weight vector \(u\).
Some generator has
\(\langle d\sigma^{\mathbb C}(E_{ca})w,\allowbreak{} u\rangle\ne0\); otherwise
\(u\), which belongs to their span, would be orthogonal to itself.
Since \(c_{J_\epsilon}(E_{ca})=E_{ac}\), the last identity in
(HCU.26) makes
\[
 \langle w,\allowbreak{} d\sigma^{\mathbb C}(E_{ac})u\rangle\ne0.
 \tag{HCU.34}
\]
(The adjoint identity has a minus sign, which does not affect
nonvanishing.)  Covariance under the diagonal torus gives the latter
vector weight \(\mu+e_a-e_c\).
Because \(w\in L\), HCU9 implies (HCU.29).
Both \(\mu\) and \(\Lambda\) are compact dominant by HCU9.
HCU10 contradicts (HCU.33).  Thus \(S=0\), proving (HCU.32).
No prior finite-dimensionality or multiplicity-one assertion for \(L\)
has entered the proof.
% END HCU11

% BEGIN HCU12
\paragraph{HCU12: original frames, extremal vectors and Borel equations.}
\label{claim:HCU12}
Let \(M\) be any original DSModel with \(\mathrm{HolShape}(M)\).
At each real place put \(q=n-p\).
Let \(r=\mathrm{revNeg}(\epsilon)\), the permutation reversing the
negative sign block and fixing the positive block, and let
\[
 a_q(j)=
 \begin{cases}n-1-j&j<q,\\j-q&j\ge q.\end{cases}
 \quad
 x_0=M.\mathrm{fr}^{-1}P_{a_q},\quad
 w_j=\Lambda_{\,\allowbreak{} r(a_q(j))}.
 \tag{HCU.35}
\]
These are the source's \(\mathrm{arr}\), \(\mathrm{x0Mat}\) and
\(\mathrm{WtIs}\) formulas.  If \(v\) is a smooth vector in the full
lowest subspace and has the extremal torus weight
\(\Lambda\circ r\), then the original \(M.\mathrm{rep}\) satisfies
\[
 d(M.\mathrm{rep})^{\mathbb C}(x_0E_{ac}x_0^{-1})v
       =\delta_{ac}w_av\qquad(a\le c).
 \tag{HCU.36}
\]
In particular every supplied source \(\mathrm{IsExtVec}\) satisfies
the full \texttt{borel\_\allowbreak ext\allowbreak Vec} conclusion.

\emph{Proof.}
The stored identity
\(M.\mathrm{fr}^*J_\epsilon M.\mathrm{fr}=H\) gives
\(x_0^*Hx_0=\operatorname{diag}(-1_q,1_p)\), exactly (HCU.1).
Conjugating the representation by \(M.\mathrm{fr}\) preserves
unitarity, irreducibility and smoothness on their original Hilbert
space.  Direct multiplication gives
\[
 x_0E_{ac}x_0^{-1}
    =M.\mathrm{fr}^{-1}E_{a_q(a),\allowbreak{} a_q(c)}M.\mathrm{fr}.
 \tag{HCU.37}
\]
The same identity for derived operators follows by differentiating
\(F e^{tX}F^{-1}=e^{tFXF^{-1}}\); its complexification respects
the original Hermitian involutions.

For a diagonal matrix unit, HCU9 gives the required extremal weight.
For a strict upper unit in the two negative-first coordinates of
(HCU.35), there are three cases.
If both indices lie in the negative block, their images under
\(r\circ a_q\) are \(p+a,\allowbreak{} p+c\), in increasing order.
If both lie in the positive block, these images are \(a-q,\allowbreak{} c-q\),
again in increasing order.  The permutation \(r\) is compact;
conjugating \(v\) by its permutation matrix turns it into a vector
of weight \(\Lambda\).  HCU9's compact raising annihilation applies
in both cases.  In the remaining case \(a<q\le c\), the matrix unit
in (HCU.37) has its row in the negative block and column in the
positive block, so HCU11 applies.
This proves (HCU.36) with the actual order and signs.

For use in constructing an extremal vector, HCU8 supplies a smooth
compact subspace of highest weight \(\Lambda\); applying the compact
permutation \(r\) to a normalized highest vector gives a unit vector
of weight \(\Lambda\circ r\) in that same subspace.
Its torus weights are integral by HCU9, so the \(w\)'s in (HCU.35)
are integers.  This produces the required weight tuple before
any appeal to lowest-type multiplicity one.
% END HCU12

% BEGIN HCU13
\paragraph{HCU13: full LowSm and LowIrr, with multiplicity one proved.}
\label{claim:HCU13}
For every original DSModel \(M\) with \(\mathrm{HolShape}(M)\),
\[
 \mathrm{LowSm}(M)\ \wedge\ \mathrm{LowIrr}(M).
 \tag{HCU.38}
\]
Thus the entire original
\texttt{Hol\allowbreak Low\allowbreak Once.\allowbreak low\allowbreak Sm\_\allowbreak and\_\allowbreak low\allowbreak Irr} producer is reconstructed.

\emph{Proof.}
Work at one real place.  HCU8 supplies a smooth irreducible compact
subspace \(S_0\) of weight \(\Lambda\).
Choose its unit extremal vector \(v\) by HCU12.
Let \(T\) be any other subspace satisfying
\(\mathrm{IsKTypeOf}(\epsilon,\sigma,\Lambda,\allowbreak{} T)\).
Its definition supplies a nonzero highest vector, and the compact
permutation \(r\), followed by normalization, gives a unit extremal
\(v'\in T\) of weight \(\Lambda\circ r\).
This \(v'\) is not assumed smooth for the full noncompact group.

Set \(l=v'-\langle v,\allowbreak{} v'\rangle v\).
It has the same torus weight, and \(\langle l,\allowbreak{} v\rangle=0\), by
conjugate symmetry and \(\|v\|=1\).
The representation on the original Hermitian coordinates is unitary
and irreducible by its stored DSModel identities.
HCU12 gives the Borel equations for the smooth vector \(v\), and
(HCU.35) identifies the torus character of both \(v\) and \(v'\)
with the exact \(\chi_w\).
HCU7 therefore gives \(l=0\), so \(v'\in S_0\).

The intersection \(T\cap S_0\) is nonzero and compact-stable.
Irreducibility first gives \(T\subset S_0\), and then \(T=S_0\).
Taking the supremum over \emph{all} defining subspaces proves
\(\mathrm{lowKSub}=S_0\).
It is consequently contained in the smooth vectors and is itself
an \(\mathrm{IsKTypeOf}\) of the required weight.  These are exactly
\(\mathrm{LowSm}\) and \(\mathrm{LowIrr}\), place by place.
The equality was proved on the original Hilbert space and with the
original compact action, not by selecting a different DSModel.
% END HCU13

\subsection*{Actual real lifts, holomorphic forms and the Hilbert coefficient}

% BEGIN HCU14
\paragraph{HCU14: the original archimedean lift, with no changed coordinates.}
\label{claim:HCU14}
For the original \(D,\theta,\allowbreak{} V\) with \(\mathrm{IsInfFam}(\theta)\),
the map
\[
 G_{\infty}(V)\longrightarrow
 \prod_\tau U(\mathrm{gramAt}(D,\allowbreak{} V,\tau)),\qquad
 g\longmapsto(\mathrm{archM}(V,\theta_\tau,\allowbreak{} g))_\tau
 \tag{HCU.39}
\]
is a continuous group isomorphism with continuous inverse \(L\).
The inverse of a tuple supported at \(\tau\) is the original
\(\mathrm{atTau}(\theta,\allowbreak{} V,\tau,\allowbreak{} g_\tau)\).
Every adelic \(g\) decomposes as \(g=au\) with
\(a\in G_f(V)\), \(u\in G_\infty(V)\), and its archimedean
evaluations belong to the stated unitary groups.
The finite and infinite adelic subgroups commute, and each original
\(\mathrm{gramAt}\) is Hermitian and nonsingular.
These are all the literal \texttt{exists\_\allowbreak real\_\allowbreak lift} fields.

\emph{Proof.}
Write \(E/F\) for the original CM quadratic extension.
The stored \(\mathrm{tauE}\) indices give exactly one member of each
conjugate pair of embeddings of \(E\).  More explicitly, the source
maps its index set bijectively to the chosen set \(\Psi\) of embeddings,
and the defining CM-type condition is
\(\sigma\in\Psi\) if and only if \(\overline{\sigma}\notin\Psi\).
Two embeddings of \(E\) over a real embedding of \(F\) are such a
pair.  Thus these indices enumerate the real places of \(F\) once.

Here the infinitary homomorphisms are forced by their stated
\(\mathrm{IsInfEmb}\) conditions, not chosen afresh.
The diagonal copy of \(F\) is dense in \(F_\infty=\prod_{v\mid\infty}F_v\).
One direct proof chooses a rational basis of \(F\); its real embedding
matrix is invertible by separability, as is seen from the Vandermonde
matrix of a primitive element and change of basis.
The image of the dense set \(\mathbb Q^{[F:\mathbb Q]}\) under this
real linear isomorphism is the required diagonal copy.
A homomorphism satisfying \(\mathrm{IsInfEmb}(\tau,\theta_\tau)\)
therefore agrees on \(F_\infty\) with evaluation at the real place
under \(\tau\), by continuity and its specified values on \(F\).
For a pure tensor it must consequently satisfy
\[
 \theta_\tau(a\otimes e)=a_{v(\tau)}\,\tau(e).
 \tag{HCU.40}
\]

To verify that these evaluations are a bijection, take
\(\delta=y-\overline y\ne0\) in \(E\).
The vectors \(1,\delta\) form an \(F\)-basis, since conjugation fixes
the first and negates the second.  At every real place,
\(\tau(\delta)\) is nonzero purely imaginary.
Thus the two real coordinates of \(a+b\delta\) give all of
\(\mathbb C\) at that place, uniquely and continuously.
Place by place, (HCU.40) is the continuous ring isomorphism
\[
 F_\infty\otimes_F E\simeq\prod_\tau\mathbb C
 \tag{HCU.41}
\]
that takes the original CM involution to complex conjugation.
The formula is the given \(\theta\), not just an isomorphic alternative.

Apply this isomorphism to each matrix entry.  The original Gram matrix
evaluates to \(\mathrm{gramAt}\); the unitary equation and inverse
matrices are preserved in both directions.
The original \(HSp\) identity that its Gram matrix equals its
CM-conjugate transpose evaluates to ordinary Hermitian symmetry.
Its nonzero determinant stays nonzero under the field embedding;
thus each evaluated Gram matrix is nonsingular.
In particular evaluating any adelic unitary \(g\), without assuming
that its finite component is the identity, gives an element of
each stated \(U(\mathrm{gramAt})\).
The adelic ring is the product of its finite and infinite parts.
Gluing the tuple in (HCU.41) to the identity finite matrix therefore
defines the unique element \(L(u)\in G_\infty(V)\).
It is continuous and multiplicative, and both inverse identities in
(HCU.39) follow entry by entry.
The stored \(\mathrm{atTau}\) has exactly the same finite identity
and single-place evaluations, so uniqueness identifies it with
\(L(\mathrm{single}_\tau g_\tau)\).
Finally matrices supported in the two different product factors
commute.  For an arbitrary adelic unitary \(g\), retain its finite
component and replace its infinite component by the identity to
obtain \(a\); perform the opposite replacement to obtain \(u\).
Both satisfy the original Gram equation componentwise, and \(g=au\).
This supplies the asserted finite--infinite decomposition.
These verifications retain the original \(F,\allowbreak{} E,\theta\),
Hermitian equation and all evaluation maps.
% END HCU14

% BEGIN HCU15
\paragraph{HCU15: the actual extremal model coefficient is holomorphic.}
\label{claim:HCU15}
Let \(M\) be an original DSModel, and assume its literal
\(\mathrm{HolShape}\), integer-weight \(\mathrm{WtIs}(M,\allowbreak{} w)\)
and \(\mathrm{IsExtVec}(M,\allowbreak{} v)\).
Then
\[
 B(u)=\prod_\tau
 \langle v_\tau,\allowbreak{} M.\mathrm{rep}_\tau(u_\tau)v_\tau\rangle
 \tag{HCU.42}
\]
has \(\mathrm{HolFn}(\mathrm{gramAt},\allowbreak{} q,\mathrm{x0Mat},\allowbreak{} w)\),
satisfies the left torus law of HCU4, and \(B(1)=1\).
This proves the whole \texttt{hol\allowbreak Fn\_\allowbreak coef\_\allowbreak ext\allowbreak Vec} output used
by the original L2 argument.

\emph{Proof.}
HCU8 identifies each \(v_\tau\in M.\mathrm{sm}_\tau\) with a genuine
smooth vector.  Differentiation of a coefficient gives
\[
 R_Z^{\mathbb C}
   \langle v_\tau,\allowbreak{} M.\mathrm{rep}_\tau(u)v_\tau\rangle
 =\langle v_\tau,\allowbreak{} M.\mathrm{rep}_\tau(u)
                 d(M.\mathrm{rep}_\tau)^{\mathbb C}(Z)v_\tau\rangle.
 \tag{HCU.43}
\]
This follows first for real \(Z\) by bounded linearity of the
fixed group operator and the inner product functional, and then
by the complexification (HCU.16).
HCU12 supplies the exact Borel equations for (HCU.43).
The finite product (HCU.42) thus has the source \(\mathrm{SmFn}\)
and \(\mathrm{KilledFn}\); HCU5 makes it a \(\mathrm{HolFn}\).

The original extremal weight and the permutation in (HCU.35) give
\[
 M.\mathrm{rep}_\tau(t(s)_\tau)v_\tau
 =\left(\prod_a e^{iw_{\tau,\allowbreak{} a}s_{\tau,\allowbreak{} a}}\right)v_\tau.
\]
Unitarity, as in (HCU.24), gives the left torus law for each
coefficient and hence their product.  The normalization in
\(\mathrm{IsExtVec}\) is \(\|v_\tau\|=1\), so \(B(1)=1\).
Every factor in (HCU.42) is the original model coefficient.
% END HCU15

% BEGIN HCU16
\paragraph{HCU16: original finite-volume representatives.}
\label{claim:HCU16}
Fix the original \(\mathrm{FinPair}(V,\mu,\mathcal F)\) and
\(\mathrm{IsReg}(V,\mu,\mathcal F,\rho)\), and put
\(\nu_Q=(\mu|_{\mathcal F})_*\mathrm{autMk}\).
A continuous bounded rational-left-invariant function \(\phi\)
has an \(L^2(\nu_Q)\) representative \(F\) satisfying the source
\(\mathrm{Repr}(F,\phi)\).
Representatives obey linearity and right translation, and
\[
 \langle F,\allowbreak{} F'\rangle
   =\int_{\mathcal F}\overline{\phi(x)}\phi'(x)\,\allowbreak{} d\mu(x)
 \tag{HCU.44}
\]
whenever \(\mathrm{Repr}(F,\phi)\) and
\(\mathrm{Repr}(F',\phi')\).
The integrand is integrable.  The \(L^2\) representative is unique;
a continuous representative is also unique pointwise.
In particular
\[
 \phi(xk)=c\phi(x)\ (\forall x)
 \quad\Longrightarrow\quad \rho(k)F=cF .
 \tag{HCU.45}
\]

\emph{Proof.}
The invariant \(\phi\) descends to a measurable function on the
quotient with its quotient measurable structure.  The pushed-forward
measure is finite, since \(\mu(\mathcal F)<\infty\).
Boundedness makes the descended function square integrable.
Its chosen \(L^2\) representative equals it almost everywhere on the
quotient, hence almost everywhere on \(\mathcal F\) after pullback.
The exceptional set in the full group is rational-left-invariant.
Countably many rational translates of the fundamental domain cover
the group up to a null set; the rational group is countable since
it is a set of matrices over the number field.
Left invariance of \(\mu\) therefore makes that exceptional set
globally null.  This is exactly the source's globally quantified
\(\mathrm{Repr}\), not merely an equality on the fundamental domain.

Conversely equality of two representatives globally almost everywhere
implies their equality for \(\mu|_{\mathcal F}\) and then for \(\nu_Q\);
this proves \(L^2\) uniqueness.
If two continuous functions agree almost everywhere for Haar measure,
their difference is zero everywhere: a nonzero value would give a
nonempty open set on which its modulus is bounded below, and Haar
measure is positive on every nonempty open set.
The usual almost-everywhere identities for \(L^2\) addition and scalar
multiplication lift globally by the same invariant-null-set argument.
The \(\mathrm{IsReg.val}\) identity and the right invariance in
\(\mathrm{FinPair}\) give
\(\mathrm{Repr}(\rho(k)F,\phi(\,\cdot\,\allowbreak{} k))\).
Finally Cauchy--Schwarz gives integrability of the product of the two
\(L^2\) functions.  The definition of the pushforward measure and
the representative identities transform its inner product into
(HCU.44).  These facts prove linearity, translation and (HCU.45),
including zero vectors and zero scalar eigenvalues where applicable.
No existence of a finite-volume target quotient is asserted here:
the same \(\mathrm{FinPair}\) remains an explicit input.
% END HCU16

% BEGIN HCU17
\paragraph{HCU17: the original extremal coefficient identity.}
\label{claim:HCU17}
On the original \(D,\theta,\allowbreak{} V,\allowbreak{} M,\allowbreak{} w,\allowbreak{} v\), assume
\(\mathrm{IsInfFam}\), \(\mathrm{HolShape}\), \(\mathrm{WtIs}\),
\(\mathrm{IsExtVec}\), the same \(\mathrm{FinPair}\) and
\(\mathrm{IsReg}\).
Assume also \(\mathrm{HolForm}(\theta,\allowbreak{} M,\allowbreak{} w,\allowbreak{} J,\chi,\phi)\)
and \(\mathrm{Repr}(F,\phi)\).
Then for every \(g\in G_\infty(V)\),
\[
 \langle F,\rho(g)F\rangle
 =\|F\|^2\prod_\tau
  \langle v_\tau,\allowbreak{} M.\mathrm{rep}_\tau
          (\mathrm{archM}(V,\theta_\tau,\allowbreak{} g))v_\tau\rangle .
 \tag{HCU.46}
\]
This is the entire literal
\texttt{inner\_\allowbreak reg\_\allowbreak eq\_\allowbreak coef\_\allowbreak extremal} statement.

\emph{Proof.}
Use the original \(L\) of HCU14.
The \(\mathrm{HolAd}\) field of \(\mathrm{HolForm}\) says that, for
each fixed \(x\),
\(u\mapsto\phi(xL(u))\) has the exact \(\mathrm{HolFn}\) type.
If \(|\phi|\le C\), then
\[
 (x,\allowbreak{} u)\longmapsto\overline{\phi(x)}\phi(xL(u))
 \tag{HCU.47}
\]
is bounded in modulus by \(\max(C,0)^2\).
It is measurable in \(x\), since \(L(u)\) is fixed when this
measurability is required and \(\phi\) is continuous.
Apply HCU6 with the finite measure \(\mu|_{\mathcal F}\).
By HCU16 the resulting holomorphic-type function is
\(A(u)=\langle F,\rho(L(u))F\rangle\).

The right upper-triangular law (HCU.4), evaluated at the diagonal
torus, gives directly
\[
 \phi(xL(t(s)))=\chi_w(s)\phi(x).
\]
Thus \(\rho(L(t(s)))F=\chi_w(s)F\) by (HCU.45).
Unitarity of \(\rho\) gives
\(A(t(s)u)=\chi_w(s)A(u)\), with the same character orientation as
(HCU.24), and \(A(1)=\|F\|^2\).
HCU15 gives the function \(B\) in (HCU.42), with the identical
holomorphic type and left torus law, and \(B(1)=1\).
HCU7 therefore yields \(A=\|F\|^2B\).
Use both inverse identities of HCU14 to replace \(L(u)\) by the
given \(g\in G_\infty(V)\); this gives (HCU.46).
The conclusion is valid when \(F=0\).
The boundedness of \(\phi\) and finite target measure were used as
explicit inputs, not inferred from source compactness or from the
name of a rapid-decay theorem.
% END HCU17

% BEGIN HCU18
\paragraph{HCU18: the exact Hilbert multilinear realization.}
\label{claim:HCU18}
Let \(G_\tau=U(H_\tau)\) act unitarily and irreducibly on complete
Hilbert spaces \(\mathcal H_\tau\), with unit vectors \(v_\tau\).
Let \(\rho\) be a unitary representation of
\(G=\prod_\tau G_\tau\) on a complete Hilbert space \(\mathcal K\).
Suppose \(k\in\mathcal K\) satisfies
\[
 \langle k,\rho(u)k\rangle
  =\|k\|^2\prod_\tau\langle v_\tau,\sigma_\tau(u_\tau)v_\tau\rangle .
 \tag{HCU.48}
\]
Then there is a continuous multilinear map
\(\Phi:\prod_\tau\mathcal H_\tau\to\mathcal K\) with
\[
 \begin{aligned}
 \Phi(v)&=k,\\
 \Phi(y_1,\ldots,\sigma_\tau(g)y_\tau,\ldots)
      &=\rho(\mathrm{single}_\tau g)\Phi(y),\\
 \|\Phi(y)\|&=\|k\|\prod_\tau\|y_\tau\|.
 \end{aligned}
 \tag{HCU.49}
\]
and whose image lies in the closed span of \(\rho(G)k\).

This proves the full source
\texttt{exists\_\allowbreak equivariant\_\allowbreak }\allowbreak
\texttt{multilinear\_\allowbreak of\_\allowbreak }\allowbreak\texttt{matrix\allowbreak Coeff\_\allowbreak eq}
contract.

For an adelic \(\rho\), HCU14 transports this construction to the
original one-place equivariance with \(\mathrm{atTau}\).
Every finite-place eigenlaw \(\rho(a)k=c k\), \(a\in G_f(V)\),
then holds for every \(\Phi(y)\).  Hence, with the inputs of HCU17,
the whole original \texttt{Back\allowbreak Form.\allowbreak exists\_\allowbreak map} conclusion holds
for the same \(L^2\) space and same initial representative \(F\).

\emph{Proof.}
Irreducibility makes the span
\(E_\tau=\operatorname{span}\{\sigma_\tau(g)v_\tau:g\in G_\tau\}\)
dense in \(\mathcal H_\tau\), because its closure is nonzero and
invariant.  The inverse in each group and unitarity transform
(HCU.48) into the orbit Gram identity
\[
 \langle \rho(a)k,\rho(b)k\rangle
   =\|k\|^2\prod_\tau
       \langle\sigma_\tau(a_\tau)v_\tau,
              \sigma_\tau(b_\tau)v_\tau\rangle .
 \tag{HCU.50}
\]
For finite expressions
\(y_\tau=\sum_j c_{\tau,\allowbreak{} j}\sigma_\tau(g_{\tau,\allowbreak{} j})v_\tau\),
define
\[
 \Phi_0(y)=
  \sum_{(j_\tau)_\tau}
       \left(\prod_\tau c_{\tau,\allowbreak{} j_\tau}\right)
       \rho((g_{\tau,\allowbreak{} j_\tau})_\tau)k .
 \tag{HCU.51}
\]
Expanding (HCU.50) gives
\[
 \langle\Phi_0(y),\Phi_0(z)\rangle
       =\|k\|^2\prod_\tau\langle y_\tau,\allowbreak{} z_\tau\rangle .
 \tag{HCU.52}
\]
If two expressions represent the same tuple, expansion of the squared
norm of the difference using (HCU.52) gives zero.  Thus (HCU.51)
is well defined and multilinear, and its norm is the one in
(HCU.49).  Approximate an arbitrary tuple by tuples from the dense
\(E_\tau\)'s.  A telescoping multilinear difference, bounded by the
product norm formula, proves that the images are Cauchy.
Completeness gives a limit independent of all approximations and
yields the continuous multilinear \(\Phi\).
All stated identities pass from the finite sums to these limits.
This includes the image in the closed orbit span.  If \(k=0\), the
construction gives the zero map; if the index set is empty it gives
the constant value \(k\), the correct zero-variable multilinear map.

For the adelic conclusion, compose the actual \(\rho\) with the
real lift \(L\) and use (HCU.46).
The original model representations are unitary, irreducible and
complete by the DSModel fields and their fixed frame conjugations.
HCU14 identifies the one-place action with \(\mathrm{atTau}\).
Since a finite-place \(a\) commutes with every \(L(u)\), its
\(c\)-eigenlaw holds on every real orbit vector \(\rho(L(u))k\)
and then on their closed span by boundedness of \(\rho(a)-c\).
HCU16 transfers each pointwise finite eigenlaw of the original
\(\phi\) to its \(L^2\) representative \(F\).
These are precisely the four outputs of
\texttt{Back\allowbreak Form.\allowbreak exists\_\allowbreak map}, including its original finite
characters and its larger full-adelic closed-orbit containment.
This is a Hilbert-valued map; conversion to pointwise smooth
automorphic functions is a separate remaining step.
% END HCU18

% BEGIN HCU19
\paragraph{HCU19: the original lowest-type cusp packet is holomorphic at its extremal.}
\label{claim:HCU19}
Let \(M,\theta,\allowbreak{} w\) satisfy the original
\(\mathrm{HolShape}\), \(\mathrm{IsInfFam}\) and \(\mathrm{WtIs}\).
Let \(U\) be an original lowest-type multilinear map whose values
are smooth, whose compact equivariance is the source
\(\mathrm{KEquiv}\), and which is killed by every original
\(\mathfrak p^-\) direction.
If \(e\) is a tuple with
\(\mathrm{IsExtVec}(M,(e_\tau)_\tau)\), then
\[
 \mathrm{HolAd}(\theta,\allowbreak{} M,\allowbreak{} w,\allowbreak{} U(e)).
 \tag{HCU.53}
\]
Thus the complete \texttt{hol\allowbreak Ad\_\allowbreak of\_\allowbreak k\allowbreak Equiv} consumer is reconstructed
with LowSm and LowIrr supplied by HCU13.

\emph{Proof.}
The lowest compact subspaces are finite dimensional by HCU13.
For any fixed output point, evaluation of the multilinear map \(U\)
is therefore a continuous multilinear form on them.
Differentiate its compact equivariance along a real compact
one-parameter group.  Ordinary finite-dimensional differentiation
gives its derivative by replacing the corresponding input by the
derived compact action.  Complexify within the compact blocks.
No continuity of \(U\) on the full discrete-series Hilbert spaces
is being assumed.

Fix \(g\) in the target adelic group and consider
\(f_g(u)=U(e)(gL(u))\), with the original \(L\) of HCU14.
It satisfies the source \(\mathrm{SmFn}\).
For the diagonal and same-sign upper units in the \(x_0\) coordinates,
the preceding differentiation and the compact part of HCU12 give
exactly \(w_a f_g\) and zero, respectively.
The remaining upper units are precisely the original
\(\mathfrak p^-\) directions after (HCU.37), and vanish by the
assumed lowering equations.
Thus \(f_g\) satisfies (HCU.17).  HCU5 gives its literal
\(\mathrm{HolFn}\).
Since \(L\) inverts the original archimedean evaluation, this is
exactly the \(\mathrm{HolAd}\) requirement for every \(g\).
% END HCU19

% BEGIN HCU20
\paragraph{HCU20: installation on the same original base model.}
\label{claim:HCU20}
Let \(R:\mathrm{RealDat}(D,\allowbreak{} A,\allowbreak{} K_s)\) be the original datum and retain
its exact \(\mathrm{models\_\allowbreak order}\) field.  For every original \(i\),
put \(M_i=R.MW(i,\mathrm{sg}_0)\).  Then
\[
 \mathrm{HolShape}(M_i),\qquad
 \exists w_i,\ \mathrm{WtIs}(M_i,\allowbreak{} w_i),\qquad
 \mathrm{LowSm}(M_i)\ \wedge\ \mathrm{LowIrr}(M_i).
 \tag{HCU.54}
\]
There is also a unit original extremal tuple, with the same integer
weight \(w_i\).  This supplies the entire
\texttt{hol\allowbreak Shape\_\allowbreak fixed} installation and the full-lowest-type inputs
retained in TAE12--TAE13.  It makes no assertion about an arbitrary
other coefficient embedding.

\emph{Proof.}
The field \(R.hv0\) says that \(A.v_0\) is the infinite place of
\(D.\mathrm{what}_2.\mathrm{embLL}\).
By the definition of \(\mathrm{sgAt}\),
\(\mathrm{sg}_0=\mathrm{sgAt}(\mathrm{embLL},\allowbreak{} A.v_0)=\mathrm{embLL}\).
The stored identity \(D.\mathrm{embE}_{\phi_0}\) makes
\[
 \mathrm{labE}(\mathrm{embLL},\tau)
   =\phi_0\circ\mathrm{sigmaOf}(D.eTau(\tau))
   =\mathrm{tauE}(D,\tau).
 \tag{HCU.55}
\]
The chosen CM-type embeddings are distinct and contain no conjugate
pair, as checked in HCU14.  Therefore any index selected in
\(\mathrm{lab}(\mathrm{embLL},\rho)\) must equal \(\rho\);
the conjugate alternative in its definition is impossible.
Unfolding \(\mathrm{kapAt}\) and \(\mathrm{parAt}\) now gives
\(\mathrm{parAt}(W,\mathrm{sg}_0,\rho)
=\mathrm{holPar}(p,\allowbreak{} q,\kappa)\), with the original
\(p,\allowbreak{} q=\mathrm{sig}_W(\mathrm{tauE}(D,\rho))\),
\(p+q=\mathrm{rank}(W)=n\), and
\(\kappa=\mathrm{wtOf}(D,\allowbreak{} W,\rho)\in\mathbb Z^{p+q}\).
No relabeling or conjugation of the chosen model has occurred.

The model's stored \(\mathrm{enum}\) already asserts that
\(\lambda=M_i.\mathrm{lam}(\rho)\) is strictly decreasing and
enumerates this signed finite parameter set.
Its entries are consequently among
\[
 \begin{aligned}
  \kappa_c^+-\frac q2+\frac{p-1}{2}-c
      &=\frac{n-1}{2}+(\kappa_c^+-q-c)
        &&(0\le c<p),\\
  \kappa_c^-+\frac p2+\frac{q-1}{2}-c
      &=\frac{n-1}{2}+(\kappa_c^--c)
        &&(0\le c<q).
 \end{aligned}
 \tag{HCU.56}
\]
Thus for each index \(a\) there is an integer \(m_a\) with
\(\lambda_a=(n-1)/2+m_a\).  Together with strict decrease this is
exactly \(\mathrm{IsHCParam}(\lambda)\).
The original \(\mathrm{models\_\allowbreak order}\) states
\(\epsilon_a=\mathbf 1_{a<p}\); since \(p\le n\), the remaining
\(\mathrm{HolShape}\) field follows.

For an explicit integral weight, let
\(C_a^+=\#\{b>a:\epsilon_b=\epsilon_a\}\) and
\(C_a^-=\#\{b<a:\epsilon_b=\epsilon_a\}\).
The elementary partition of indices above and below \(a\) gives
\(\rho_N(a)+\rho_C(a)=(n-1)/2-a\).  Hence
\[
 \Lambda_a=\lambda_a+\rho_N(a)-\rho_C(a)
       =m_a+n-1-a-C_a^++C_a^-\in\mathbb Z .
 \tag{HCU.57}
\]
Take \(w_i(\rho,\allowbreak{} j)=\Lambda_{r(a_q(j))}\), with the actual
\(\mathrm{revNeg}\) and \(\mathrm{arr}\) of HCU12.
This is exactly \(\mathrm{WtIs}\), including empty signature blocks.
HCU13 now gives both full-lowest-type assertions on this same \(M_i\).
HCU8 and HCU12 provide its unit extremal vector at each real place;
their finite tuple uses the just computed weight.
All conclusions hold for every \(i\) before any choice of the
parameter \(t\).  They use \(R.\mathrm{models\_\allowbreak order}\), which is a
field of the already required \(\mathrm{IsModels}\), and do not
construct a new datum \(R\).
% END HCU20

\subsection*{What this closes, and what still must be produced}

The proved endpoints are the actual local and global coefficient
uniqueness, the complete smooth Borel--holomorphic dictionary,
bounded holomorphic integration, full lowest-type smoothness and
multiplicity one, the original extremal coefficient identity,
and its Hilbert multilinear realization with original finite
eigenlaws.  In particular the full-lowest-type producer retained
in TAE is now supplied under its exact \(\mathrm{HolShape}\)
hypothesis; it is not replaced by an existentially selected
irreducible summand.

For the base original model \(M_i=R.MW(i,\mathrm{sg}_0)\), the required
\(\mathrm{HolShape}\) and integer weight are supplied by its ordered
holomorphic parameter formulas, reconstructed with the
\(\mathrm{sg}_0=D.\mathrm{what}_2.\mathrm{embLL}\) identification
in HCU20.  For any other coefficient embedding, the
same HCU implications apply when \emph{that model's}
\(\mathrm{HolShape}\) is checked; this chapter does not assert that
every embedding has that shape.

The original TAE13 cusp packet, together with HCU13 and HCU19,
therefore reaches a holomorphic extremal function.
If the actual boundedness and finite target quotient are supplied,
HCU16--HCU18 give the literal coefficient and Hilbert-map stages.
They do not manufacture those analytic inputs.
The source's predicate \(\mathrm{Published.CuspRapidDecay}\)
has not been turned into a checked published application.
A finite-volume \(\mathrm{FinPair}\), a regular realization
\(\mathrm{IsReg}\), and the original scalar function are kept
on the same target quotient throughout.

To finish \texttt{exists\_\allowbreak mult\allowbreak Map}, one still needs the original
smooth representatives and all their derivative identities in that
Hilbert image, point-evaluation Sobolev bounds and resulting
uniform growth, central finiteness, and preservation of cusp
integrals under closed orbit span.  HCU8 supplies smooth-vector
calculus, but it does not prove those global automorphic
point-evaluation estimates.
The remaining exact sources are:
\begin{itemize}
\item \texttt{Smooth\allowbreak Vector\allowbreak Automorphic\allowbreak Derivatives};
\item \texttt{Sobolev\allowbreak Bound\allowbreak For\allowbreak Smooth\allowbreak Vectors};
\item \texttt{Gelfand\allowbreak Invariants\allowbreak Act\allowbreak By\allowbreak Scalars\allowbreak Highest\allowbreak Weight};
\item \texttt{Cuspidality\allowbreak Preserved\allowbreak In\allowbreak Closure}.
\end{itemize}
No full \(\mathrm{MMap}\), \(\mathrm{firesAt}\) or
\(\mathrm{exists\_\allowbreak thetaMap}\) conclusion is claimed yet.

Independently, one must still construct, on the same original \(R\),
one \(io_i\) with
\[
 \exists io_i\ \forall t,\quad
 \mathrm{IsIoW}\bigl(R.\theta,\allowbreak{} M_i,\allowbreak{} R.cV(i),
       R.\omega P(i,\allowbreak{} t),\allowbreak{} R.ptP(i,\allowbreak{} t),\mathrm{Delta}(i),\allowbreak{} io_i\bigr).
 \tag{HCU.58}
\]
The two intertwining laws, exact polynomial normalization and the
quantifier order are unchanged.  Only such a simultaneous witness
activates the stored \(\mathrm{OrZero}(\mathrm{PIoP}(i))\).
No argument above supplies this polynomial map.
The same-Bundle weak-data producer, geometric Hilbert reading,
rational adjoints and tests, and original \(\mathrm{RowsHold0_3}\)
remain separate.  The present proofs are pointwise and introduce
no new family threshold.  Formal correctness is assumed as in the
project; no new Lean replay is claimed.







\section{The original auxiliary-prime supplier and residual Hecke values}
\label{sec:count-auxiliary-primes}

\status{CAP1--CAP14 prove the full original \texttt{Ar\allowbreak Pl} contract,
conditional on its unchanged \texttt{Coh\allowbreak P} hypothesis.  This supplies
\texttt{JSP\_\allowbreak ACTUAL\_\allowbreak ARPL}, including \texttt{Phi\allowbreak Cond}, the exact residual
subgroup, a Taylor--Wiles system for every prescribed finite list, its
actual arithmetic realization, and one character bound for all levels.
The only new published input is the unconditional Chebotarev theorem,
with its hypotheses and Frobenius convention checked below.  Neither
ordinary Hecke data nor \texttt{Has\allowbreak Aut}, \texttt{Choice}, or the two
cohomological bounds are inferred from this construction.  CAP15--CAP17
also supply the full original \texttt{Is\allowbreak Res} values for any supplied
integral type, at every original \texttt{QData} level.  The type,
\texttt{TTheta\allowbreak Ex}/\texttt{neat\allowbreak U}, and ordinary-control producers remain
explicit.}

Retain the original prime $\ell$, CM number field $L$, maximal real
subfield $L^+$, algebraic closure $\Omega$, and
$s:\mathrm{HidaSetting}(\ell,\allowbreak{} L,\Omega)$.  Put $S=s.S$.  The initial
level $V_0:s.\mathrm{Lev}$ satisfies $V_0.S'=S$.  Thus
\[
 \Gamma=G_{L^+,\allowbreak{} S}
  =G_{L^+}/N_S,\qquad G_0=\ker\eta,\qquad
 G'=\ker(\varepsilon|_{G_0}),
 \tag{CAP.1}
\]
where $N_S$ is the closed normal subgroup generated by all inertia
groups outside $S$.  The map $\rho:G_L\to\Gamma$ is the original
$V_0.\mathrm{pr}$ followed by the inclusion of $G_0$.
The fixed local datum is $\mathrm{locOf}(r)$; in particular its local
groups are the decomposition groups at the original valuations
$s.A(v)$ and $s.A(u)$, including their original inertia conditions.
Changing the global level below will leave this local datum fixed.

Let $O$ be any coefficient DVR satisfying the original
$\mathrm{IsCoef}(\ell,\allowbreak{} O)$, with its specified $\mathbb Z_\ell$-algebra
structure, and let $r:\mathrm{Resid}(V_0.\mathrm{frame},\allowbreak{} O,\allowbreak{} n)$ for
\emph{any} $n\geq0$.  Write $k=O/\mathfrak m_O$ and $\bar\chi_i$
for its residual characters.  The original hypothesis implies
$\operatorname{char}k=\ell$.  We use additive notation for
\[
 X_\ell(K)=\operatorname{Hom}_{\mathrm{cont}}(K,\mathbb F_\ell).
 \tag{CAP.2}
\]
Since the target is finite discrete, continuity is equivalent to an
open kernel.  Thus this is exactly the source's $\mathrm{HomL}$,
written additively.  The assertion $\Phi_\ell(\Gamma)$ means that
$X_\ell(K)$ is finite for \emph{every} open subgroup $K\leq\Gamma$.

The source anchors for the contracts are:
\begin{quote}\small
\nolinkurl{TaylorWilesObstructionBound.lean}:23--31,37--71,88--124;\\
\nolinkurl{ThetaCoefficientMapEq.lean}:12--40;\\
\nolinkurl{PairedDualFirstOrderDualMatrixSpace.lean}:78--106.
\end{quote}
All source paths and exact byte ranges are recorded in the accompanying
provenance; the proofs below establish the mathematical implications.

\subsection{Finiteness and the exact residual subgroup}

% BEGIN CAP1
\paragraph{CAP1. Finite first cohomology implies the full $\Phi$ condition.}
\label{claim:CAP1}
Let $\Pi$ be a profinite group.  Suppose $H^1_{\mathrm{cont}}(\Pi,\allowbreak{} M)$
is finite for every finite discrete $\mathbb F_\ell$-module $M$ with
continuous $\Pi$-action.  Then $\Phi_\ell(\Pi)$ holds.

\emph{Proof.}
Fix an open subgroup $K$.  The left-coset space $X=\Pi/K$ is finite:
its disjoint open cosets cover the compact group.  Choose representatives
$a_x$, with $a_K=1$, and let $M=\operatorname{Map}(X,\mathbb F_\ell)$,
with $(g m)(x)=m(g^{-1}x)$.  This is a finite continuous module.  For
$\chi\in X_\ell(K)$ define
\[
 z_\chi(g)(x)
   =\chi\bigl(a_x^{-1}g a_{g^{-1}x}\bigr).
 \tag{CAP.3}
\]
The argument belongs to $K$.  Inserting $a_{g^{-1}x}a_{g^{-1}x}^{-1}$
gives
\[
 z_\chi(gh)(x)=z_\chi(g)(x)+z_\chi(h)(g^{-1}x),
\]
so $z_\chi$ is a cocycle.  For each $x$, the representative
$a_{g^{-1}x}$ is locally constant in $g$; on each of its finitely many
clopen pieces (CAP.3) is a continuous map into $K$, followed by
$\chi$.  Hence $z_\chi$ is a continuous cocycle.

The map $\chi\mapsto[z_\chi]$ is injective.  Indeed, for $h\in K$,
$z_\chi(h)(K)=\chi(h)$, whereas every coboundary $h m-m$ has value
zero at the coset $K$.  If two cocycles define the same class, apply
this observation to their difference.  Thus $X_\ell(K)$ injects into
the finite group $H^1_{\mathrm{cont}}(\Pi,\allowbreak{} M)$.  This proves the result
for every open $K$, without requiring $K$ to be normal.\hfill$\square$

This is the written induced-cocycle argument underlying
\nolinkurl{ChebotarevFamilies/PhiGS.lean}:47--141 and
\nolinkurl{PhiDec.lean}:19--69; no separate Shapiro theorem is needed.
% END CAP1

% BEGIN CAP2
\paragraph{CAP2. The actual global and local $\Phi$ conditions.}
\label{claim:CAP2}
The original $\mathrm{CohP}$ supplies $\Phi_\ell(\Gamma)$ and
$\Phi_\ell(D_{s.A(v),\allowbreak{} L})$, $\Phi_\ell(D_{s.A(u),\allowbreak{} L})$ for all the
original local indices.

\emph{Proof.}
For the global assertion apply $\mathrm{CohP.gf}$ to the number field
$L^+$, the given $\Omega$, the finite set $V_0.S'$, and exponent $1$.
This set contains every prime above $\ell$, since $s.S\subseteq V_0.S'$.
The module in CAP1 is finite, discrete, and its action is continuous,
so it satisfies every hypothesis of
$\mathrm{GlobalFinitenessAt}$.  Its $H^1$ is therefore finite.  CAP1
applies to $G_{L^+,\allowbreak{} V_0.S'}$.

For the local assertion use $\mathrm{CohP.lf}$, again at exponent $1$,
with the number field $L$, the same $\Omega$, the actual valuation
$s.A(v)$, and its recorded prime $s.w(v)$.  The field extension and
place-over hypotheses are exactly those stored in $s$.  A decomposition
group is a closed subgroup of the absolute Galois group, hence
profinite.  CAP1 applies to it; the argument for $u\in s.St$ is
identical with its own recorded valuation.  These are the precise groups
in $\mathrm{locOf}(r)$, rather than newly chosen isomorphic models.
The other fields of $\mathrm{CohP}$ remain part of the original
hypothesis; none is replaced or asserted to follow here.\hfill$\square$

The exact finiteness predicates are
\nolinkurl{ReciprocityTheory/Hypotheses/GalCoh.lean}:17--20,151--155,
195--199,219--224.  The corresponding original applications are
\nolinkurl{PhiGS.lean}:142--160 and \nolinkurl{PhiDec.lean}:70--79.
% END CAP2

% BEGIN CAP3
\paragraph{CAP3. The original residual kernel is open and normal.}
\label{claim:CAP3}
There is an open normal subgroup $H\leq\Gamma$ with the exact formula
\[
 g\in H\quad\Longleftrightarrow\quad
 \eta(g)=1,\quad\bar\chi_i(g)=1\ (\forall i),\quad
 \bar\varepsilon(g)=1,
 \tag{CAP.4}
\]
where $\bar\varepsilon$ is reduction of $\varepsilon$ in
$\mathbb F_\ell^\times$.  In particular $H\subseteq G_0$ and both
the source $H_\eta$ and $H_\varepsilon$ clauses hold.

\emph{Proof.}
The displayed set is a subgroup.  It is open because $G_0$ is open,
the finitely many $\bar\chi_i$ have open kernels, and the residual
cyclotomic character is continuous with finite discrete target.
It is stable under conjugation by $G_0$, since all its conditions are
kernels of characters on $G_0$.  For conjugation by the specified
complex conjugation $c$, the residual determinant polarization is
essential.

Fix $x\in H$.  The image of $\varepsilon(x)$ in $k$ is $1$:
$\bar\varepsilon(x)=1$ says $\varepsilon(x)-1\in\ell\mathbb Z_\ell$,
which maps to zero in $k$.  The source involution consequently sends
$t[x]-1$ to $t[cxc^{-1}]^{-1}-1$.  Applying the full residual polarized
determinant law gives, for every $t\in k$,
\[
 (t-1)^n=
   \prod_{i=1}^{n}\bigl(t\bar\chi_i(cxc^{-1})^{-1}-1\bigr).
 \tag{CAP.5}
\]
For a particular index $j$ take $t=\bar\chi_j(cxc^{-1})$.
The right side is zero, so $(t-1)^n=0$.  Existence of $j$ implies
$n>0$; since $k$ is a field, $t=1$.  When $n=0$ there are no such
conditions.  The cyclotomic and $\eta$ conditions are preserved by
every conjugation.  Thus $cHc^{-1}=H$.  As
$\eta(\Gamma)=\{1,-1\}$ and $\eta(c)=-1$, every element of $\Gamma$
belongs to $G_0$ or $cG_0$, proving normality.

This retains exactly (CAP.4).  Taking an arbitrary normal core would
give an inclusion but would not establish the required equality
$\mathrm{TWPl.H\_\allowbreak eq}$.  No division by $n!$ has occurred; the argument
includes $\ell=3$ in rank four.\hfill$\square$

Source polarization:
\begin{quote}\small
\nolinkurl{UniversalOrdinaryDeterminantRing.lean}:80--99,109--117.\\
Residual-kernel consumer: \nolinkurl{ChebotarevFamilies/StarSingle.lean}:20--152.
\end{quote}
% END CAP3

\subsection{The precise relative pro-\texorpdfstring{$\ell$}{ell} subgroup and its tower}

For a subgroup $K\leq\Pi$ put
\[
 P_\ell^\Pi(K)=
 \bigcap_{\substack{U\triangleleft_{\mathrm{open}}\Pi\\
  \forall x\in K\ \exists a\geq0:\ x^{\ell^a}\in U}}U.
 \tag{CAP.6}
\]
This is the literal source $\mathrm{PL}(\ell,\allowbreak{} K)$.  Because $\Pi/U$
is finite, its displayed condition says precisely that the image of
$K$ is an $\ell$-group: every element has $\ell$-power order, and
Cauchy's theorem excludes every other prime divisor of its order.
Each term is closed and normal, so their intersection is closed and
normal.  The ambient group in the superscript matters.

% BEGIN CAP4
\paragraph{CAP4. Character kernels and local pro-$\ell$ kernels.}
\label{claim:CAP4}
If $K$ is open normal in a profinite group $\Pi$, every
$\chi\in X_\ell(K)$ vanishes on $P_\ell^\Pi(K)$.  Also, if
$f:I\to\Pi$ is continuous and $f(I)\subseteq K$, then
\[
 f\bigl(\operatorname{proKer}_\ell(I)\bigr)
       \subseteq P_\ell^\Pi(K),
 \tag{CAP.7}
\]
where $\operatorname{proKer}_\ell(I)$ is the intersection of all open
normal subgroups of $I$ of $\ell$-power index.

\emph{Proof.}
For the first assertion $A=\ker\chi$ is open in $\Pi$ and has finite
index.  Its normal core $U=\bigcap_{g\in\Pi}gAg^{-1}$ is an intersection
of finitely many conjugates, hence is open normal.  For $x\in K$ and
every $g\in\Pi$, normality of $K$ gives $g^{-1}xg\in K$, and
$\chi((g^{-1}xg)^\ell)=0$.  Hence $x^\ell\in U$.
Thus $U$ is one of the terms in (CAP.6), so
$P_\ell^\Pi(K)\subseteq U\subseteq\ker\chi$.

For (CAP.7), fix any term $U$ of (CAP.6).  The image of
$I\xrightarrow{f}\Pi\to\Pi/U$ is a finite $\ell$-group, as it lies
in the image of $K$.  Its kernel is open normal in $I$ and has
$\ell$-power index.  It therefore contains
$\operatorname{proKer}_\ell(I)$.  Intersecting over $U$ proves (CAP.7).
In the local application $I$ is the actual inertia group, and its
image $\operatorname{proKer}$ is exactly the source
$\mathrm{ineKer}$.\hfill$\square$

Compare \nolinkurl{ChebotarevFamilies/HomBound.lean}:5--78,
\nolinkurl{ProkerPL.lean}:5--32, and the definitions in
\texttt{GaloisCohomology\textsubscript{2}/ContinuousCycMod.lean}:10--36.
% END CAP4

% BEGIN CAP5
\paragraph{CAP5. A normal Frattini tower cofinal above the exact $P_\ell$.}
\label{claim:CAP5}
If $\Phi_\ell(\Gamma)$ and $H$ is open normal, there are descending
open normal subgroups $W_N\subseteq H$, $N\geq0$, such that
\[
 P_\ell^\Gamma(H)\subseteq V\triangleleft_{\mathrm{open}}\Gamma
 \quad\Longrightarrow\quad
 \exists N_0\ \forall N\geq N_0:\ W_N\subseteq V.
 \tag{CAP.8}
\]

\emph{Proof.}
For an open normal $K$ define
$\operatorname{fr}(K)=\bigcap_{\chi\in X_\ell(K)}\ker\chi$.
This is a finite intersection by $\Phi_\ell(\Gamma)$, and hence is
open.  Conjugation by $\Gamma$ permutes $X_\ell(K)$, so it is normal
in $\Gamma$.  Set $W_0=H$ and $W_{N+1}=\operatorname{fr}(W_N)$.

Fix $V$ as in (CAP.8).  The family of open normal subgroups in (CAP.6)
is closed under finite intersections: in the quotient by an
intersection, the image of $H$ embeds into a product of finite
$\ell$-groups.  Compactness shows that some such subgroup $U$ lies
in $V$.  Indeed, otherwise their intersections with the compact set
$\Gamma\setminus V$ would have the finite intersection property and
a common point, contradicting $P_\ell^\Gamma(H)\subseteq V$.
Consequently the image of $H$ in $\Gamma/V$ is a finite $\ell$-group.

Every nontrivial finite $\ell$-group $Q$ has a nonzero homomorphism to
the additive group $\mathbb F_\ell$.  To see this without a structural
theorem, the class equation shows that its center has order divisible
by $\ell$.  Choose a central subgroup of order $\ell$ by Cauchy's
theorem.  If the quotient by it is nontrivial, induction on $|Q|$
gives the desired homomorphism through that quotient; if it is trivial,
$Q$ itself is cyclic of order $\ell$.

Apply this observation to the image of $W_N$ in $\Gamma/V$.  Whenever
that image is nontrivial, a nonzero character of it pulls back to a
continuous character of $W_N$.  Hence the image of
$\operatorname{fr}(W_N)$ lies in a proper kernel and has strictly
smaller order.  After finitely many steps the image is trivial.
The descending property then gives (CAP.8) for every later $N$.
This is cofinality relative to $P_\ell^\Gamma(H)$, which is exactly
the condition used by $\mathrm{TWGrp.ph\_\allowbreak lim}$.\hfill$\square$

The source construction is \nolinkurl{ChebotarevFamilies/FrSNormal.lean}:6--207.
% END CAP5

\subsection{Chebotarev with the required places and sign}

% BEGIN CAP6
\paragraph{CAP6. A prescribed coset has degree-one arithmetic Frobenius.}
\label{claim:CAP6}
Let $E$ be a number field, $\Omega$ an algebraic closure, $U$ an open
subgroup of $G_E$, $g\in G_E$, and $A$ a finite set of finite primes of
$E$.  There are a prime $w\notin A$ of absolute norm a rational prime
$p$, a valuation $B$ of $\Omega$ above $w$, and
$\sigma\in D_{B,\allowbreak{} E}$ with
\[
 \sigma\text{ arithmetic Frobenius at }B,\qquad
 \sigma g^{-1}\in U.
 \tag{CAP.9}
\]

\emph{Published input and hypotheses.}
We use Chebotarev for a finite Galois extension of number fields:
each conjugacy class occurs as an unramified Frobenius class at a
positive density of primes.  The primary reference is J.~S.~Milne,
\emph{Algebraic Number Theory}, v3.08 (19 July 2020), Theorem 8.31,
printed p.~147, \url{https://www.jmilne.org/math/CourseNotes/ANT.pdf}.
Its Frobenius convention is the residue-field power map, as in
the discussion preceding 8.14.  We apply the theorem below to the
explicit finite Galois extension $E'/\mathbb Q$; no effective estimate
or RH hypothesis is involved.

\emph{Proof of the matching.}
Choose an open normal subgroup $U_0\subseteq U$; the normal core of
an open subgroup suffices.  Its fixed field $F/E$ is finite Galois.
Choose a finite Galois extension $E'/\mathbb Q$ inside $\Omega$
containing $F$, and put $h=g|_{E'}$.  Since $h$ fixes $E$, it belongs
to $\operatorname{Gal}(E'/E)$.  Chebotarev supplies infinitely many
rational primes unramified in $E'$ with Frobenius conjugate to $h$.
Exclude the finitely many rational primes below $A$.  Conjugating a
prime of $E'$ if necessary, choose $\mathfrak P$ whose Frobenius is
\emph{exactly} $h$.  Put $w=\mathfrak P\cap O_E$.
For every $a\in O_E$, the equality $h(a)=a$ implies
$a^p\equiv a\pmod w$.  Thus the residue field of $w$ is
$\mathbb F_p$ and $\operatorname{Norm}w=p$.

Extend $\mathfrak P$ to a valuation $B$ of $\Omega$ and lift the
arithmetic residue Frobenius to $D_{B,\mathbb Q}$.
For completeness, the finite valuation facts and this limiting passage
can be proved as follows.  In a finite Galois extension $F/K$, suppose
primes $P,\allowbreak{} Q$ above the same prime were in different Galois orbits.
The Chinese remainder theorem gives an integer $b$ that is zero
modulo $Q$ and one modulo every conjugate of $P$.  Its norm belongs
to $Q\cap O_K$ and hence to $P$.  Since the norm is a product of
conjugates of $b$, primality of $P$ contradicts those congruences.
Thus the action on primes is transitive.

The decomposition group at $P$ induces the full residue Galois
group.  To verify surjectivity explicitly, choose an integer $z$
whose residue at $P$ generates the multiplicative group of the finite
residue field, and whose residue at every other prime above the base
prime is zero.  Put $f(T)=\prod_{g\in\operatorname{Gal}(F/K)}(T-gz)$.
Its coefficients lie in $O_K$.  If the base residue cardinality is
$q$, then $f(z^q)=0$ modulo $P$.  Some $g$ therefore satisfies
$gz\equiv z^q\pmod P$.  Such a $g$ must stabilize $P$, since otherwise
$gz\equiv0\pmod P$.  Its residue action is the $q$-power map, which
generates the residue Galois group.  The kernel of this residue action
is exactly inertia by definition.

Finite integral lying-over and passage through finite composita give
a valuation $B$ extending $\mathfrak P$.  Fix this valuation.  For
each finite Galois subextension of $\Omega/\mathbb Q$ containing
$E'$, require an automorphism to stabilize its restriction of $B$
and induce the $p$-power residue map.  These are nonempty closed
conditions on the compact absolute Galois group by the preceding
finite argument.  Any finitely many are simultaneously satisfied
in their compositum.  Compactness supplies an automorphism satisfying
all of them; it stabilizes $B$ and induces the $p$-power map on every
residue.  The same compactness argument applied to finite prime
transitivity gives transitivity on valuations over a fixed prime.
Because $E'/\mathbb Q$ is unramified at $\mathfrak P$, this lift
$\sigma$ restricts to $h$ on $E'$.  It fixes $E$, belongs to
$D_{B,\allowbreak{} E}$, and induces the $\operatorname{Norm}w$ power map on
residues.  Finally $\sigma g^{-1}$ fixes $F$, so it belongs to
$U_0\subseteq U$.\hfill$\square$

The source target is \nolinkurl{ChebotarevFamilies/ChebOne.lean}:54--108;
the degree-one input is traced in
\nolinkurl{Proofs/InertiaEigenvalues/ChebotarevDegreeOne.lean}:17--69.
The argument above derives the required degree-one conclusion directly
from the stated finite Galois theorem.
% END CAP6

% BEGIN CAP7
\paragraph{CAP7. Split CM places and every prescribed finite list.}
\label{claim:CAP7}
In the original CM field $L$, fix an open subgroup $U\leq G_L$,
a finite set $A$ of primes of $L^+$, and a list
$g_j\in G_L$, $j\in\operatorname{Fin}(q)$.  One can choose
$v_j,\allowbreak{} w_j,\allowbreak{} B_j,\sigma_j$ such that the $v_j$ are distinct and outside
$A$, $B_j$ lies over $w_j$ and $v_j$,
\[
 D_{B_j,\allowbreak{} L^+}\subseteq G_L,
 \qquad \operatorname{Norm}w_j=p_j\text{ is prime},\qquad
 \sigma_jg_j^{-1}\in U,
 \tag{CAP.10}
\]
and $\sigma_j$ is an arithmetic Frobenius at $B_j$ over $L$.

\emph{Proof.}
Choose an algebraic integer $a\in O_L$ with $a\ne\bar a$; a
nonreal element of $L$ becomes integral after multiplication by a
nonzero rational integer.  Enlarge the finite forbidden set of primes
of $L$ to include every prime above $A$ and every prime dividing
$a-\bar a$.  Finiteness follows from finite extension degree and the
factorization of the nonzero principal ideal $(a-\bar a)$.
Apply CAP6.  If $B$ is the resulting valuation and $w$ its prime,
an element $d\in D_{B,\allowbreak{} L^+}$ restricts on $L$ to either $1$ or CM
conjugation.  It stabilizes $w$ and acts trivially on its residue field
$\mathbb F_p$.  The second possibility would imply
$a-\bar a\in w$, which was excluded.  Thus it fixes $L$, proving
the decomposition-group inclusion.  The prime $v$ below $w$ is
outside $A$.

For the list choose its entries in order, adding previously chosen
$v_i$ to $A$ at each step.  This gives injective $j\mapsto v_j$ and
retains all coset conditions.  When $q=0$ use the empty family.

We will also use the following consequence of the valuation facts
in CAP6.  If $B$ and $B'$ lie over the same prime of $L^+$, they are
conjugate by $G_{L^+}$: the finite-level transitivity statements have
a simultaneous solution by compactness.  The subgroup $G_L$ is
normal in $G_{L^+}$ because $L/L^+$ is quadratic Galois.  Therefore
$D_{B,\allowbreak{} L^+}\subseteq G_L$ implies
$D_{B',\allowbreak{} L^+}\subseteq G_L$.  This supplies splitting at the original
stored valuation $s.A(v)$ even when the prime was selected using a
different valuation $B$.\hfill$\square$

These are the exact conditions of \nolinkurl{ChebOne.lean}:6--108 and
\nolinkurl{ChebFam.lean}:6--35.
% END CAP7

\subsection{The actual level maps and the uniform bound}

% BEGIN CAP8
\paragraph{CAP8. Enlarging the original level with its fixed local datum.}
\label{claim:CAP8}
Let $Q$ be a finite set of primes of $L^+$, each with a valuation
$B_v$ over $v$ satisfying $D_{B_v,\allowbreak{} L^+}\subseteq G_L$.  There is an
actual level $V_1:s.\mathrm{Lev}$ with
$S_1=V_1.S'=V_0.S'\cup Q$ and a surjective continuous map
\[
 p:G_{L^+,\allowbreak{} S_1}\longrightarrow G_{L^+,\allowbreak{} V_0.S'},\qquad
 p(\operatorname{toGS}_{S_1}g)=\operatorname{toGS}_{V_0.S'}g,
 \tag{CAP.11}
\]
which is a $\mathrm{Glob.Hom}$ from
$\mathrm{globOf}(r,\allowbreak{} V_1)$ to $\mathrm{globOf}(r,\allowbreak{} V_0)$.

\emph{Proof.}
Since $V_0.S'\subseteq S_1$, the generators of $N_{S_1}$ are among
those of $N_{V_0.S'}$, so $N_{S_1}\subseteq N_{V_0.S'}$.
The quotient map gives (CAP.11), including surjectivity and continuity.
Set $\eta_1=\eta_0p$ and $\varepsilon_1=\varepsilon_0p$.
The original $\eta$ characterization, value at $c$, open kernel,
cyclotomic identity, and continuity pull back along $p$.
The enlarged set is finite and contains $s.S$.  At old primes the
stored splitting condition is unchanged.  At new primes CAP7
transports the supplied splitting condition to $s.A(v)$.
These verify every field of $V_1$.

The global datum uses the same $\mathrm{locOf}(r)$ on both sides.
Its local maps are restrictions of the quotient maps followed by
$G_L\to G_{L^+}$.  Consequently (CAP.11) commutes with every
original $\varphi_v$ and $\varphi_u$, as well as $\eta$, $c$, and
$\varepsilon$.  These are precisely all fields of
$\mathrm{Glob.Hom}$.  In particular its induced map on $G_0$
satisfies $p_0(V_1.\mathrm{pr}(g))=V_0.\mathrm{pr}(g)$ for
every $g\in G_L$.\hfill$\square$

Source level construction and contracts:
\begin{quote}\small
\nolinkurl{ChebotarevFamilies/LevUp.lean}:5--53;\\
\nolinkurl{ContinuousResP.lean}:106--162;\\
\nolinkurl{PairedDualFirstOrderDualMatrixSpace.lean}:78--99.
\end{quote}
% END CAP8

% BEGIN CAP9
\paragraph{CAP9. The whole level kernel is generated modulo the exact $P_\ell$.}
\label{claim:CAP9}
Suppose $Q=\{v_j:j\in\operatorname{Fin}(q)\}$ is disjoint from
$V_0.S'$, with the valuations $B_j$ of CAP7.  Let $p$ be (CAP.11),
$K=p^{-1}H$, and let $\tau_j\in G_L$ satisfy
$\mathrm{IsProGen}(L,\allowbreak{} B_j,\ell,\tau_j)$.  Put
$t_j=V_1.\mathrm{pr}(\tau_j)$, regarded also as an element of
$\Gamma_1=G_{L^+,\allowbreak{} S_1}$.  Then $p(t_j)=1$ and
\[
 \ker p\ \subseteq\
 \overline{\langle\!\langle t_j:j\in\operatorname{Fin}(q)\rangle\!\rangle
       \ P_\ell^{\Gamma_1}(K)}.
 \tag{CAP.12}
\]
The bar means topological closure of the subgroup generated by the two
displayed normal subgroups.  Moreover the image of every
$\mathrm{ineKer}(L,\allowbreak{} B_j,\ell)$ lies in $P_\ell^{\Gamma_1}(K)$.

\emph{Proof.}
Every inertia element at $v_j$ is killed by the old quotient, because
$v_j\notin V_0.S'$.  Hence its image in $\Gamma_1$ is in $\ker p$
and therefore in $K$.  Apply CAP4 to the continuous map from its
actual inertia group to $\Gamma_1$: its pro-$\ell$ kernel maps into
$P_\ell^{\Gamma_1}(K)$.  The definition of $\mathrm{IsProGen}$ says
that the inertia group is contained in the closure of the subgroup
generated by $\tau_j$ and its $\mathrm{ineKer}$.  Its whole image
therefore lies in the right side $C$ of (CAP.12); continuity suffices
because $C$ is closed.  In particular $p(t_j)=1$.

The group $C$ is closed normal in $\Gamma_1$.  To show that it kills
the complete old ramification kernel, take an inertia group at a
prime $v\notin V_0.S'$, with any valuation $A$ above $v$.
If $v\notin S_1$ it is killed already by $\Gamma_1$.
Otherwise $v=v_j$.  By valuation transitivity choose a conjugation
carrying $A$ to $B_j$.  At $B_j$ the decomposition group over $L^+$
is contained in $G_L$, so its inertia is exactly the inertia over
$L$, with the same residue action.  The previous paragraph puts its
image in $C$, and normality transports this conclusion back to $A$.
Thus the inverse image of $C$ in $G_{L^+}$ contains every inertia
generator of $N_{V_0.S'}$.  It is closed normal, so contains all of
$N_{V_0.S'}$.  Every element of $\ker p$ lifts to an element of this
last group, which proves (CAP.12).\hfill$\square$

This accounts for the full, closed kernel in
\nolinkurl{ChebotarevFamilies/PLNormal.lean}:37--101, including arbitrary
valuations over each new prime.
% END CAP9

% BEGIN CAP10
\paragraph{CAP10. A character bound uniform in the enlarged level.}
\label{claim:CAP10}
Let $p:\Pi\twoheadrightarrow\Gamma$ be continuous and surjective
between profinite groups, let $H\triangleleft_{\mathrm{open}}\Gamma$,
and set $K=p^{-1}H$.  Suppose $X_\ell(H)$ is finite and $t_1,\ldots,\allowbreak{} t_q$
lie in $\ker p$, with (CAP.12) for this $p$ and $K$.  Then
\[
 |X_\ell(K)|\leq |X_\ell(H)|\,
       \ell^{\,\allowbreak{} q[\Gamma:H]}.
 \tag{CAP.13}
\]
In particular $X_\ell(K)$ is finite.

\emph{Proof.}
The quotient map identifies $\Pi/K$ with $\Gamma/H$.  Choose
representatives $a_1,\ldots,\allowbreak{} a_d$ for these cosets, where
$d=[\Gamma:H]$.  Evaluate a character on the $qd$ elements
$a_i t_j a_i^{-1}\in K$.  There are at most $\ell^{qd}$ possible
arrays of values.

If two characters have the same array, their difference $\chi$
vanishes on these elements and hence on every conjugate of every
$t_j$.  Indeed, write any $a\in\Pi$ as $k a_i$ with $k\in K$;
conjugation by $k$ does not affect an abelian character of $K$.
It vanishes on $P_\ell^\Pi(K)$ by CAP4.  Its kernel is closed in
$\Pi$, since $K$ is open and closed and the character is continuous.
The kernel therefore contains the right side of (CAP.12), so
$\chi$ kills $\ker p$.

It descends uniquely to a character $\chi_0:H\to\mathbb F_\ell$.
The restricted surjection $p:K\to H$ is a quotient map: a continuous
surjection from a compact space onto a Hausdorff space has this
property.  Consequently $\chi_0$ is continuous, since its composite
with $p$ is $\chi$.  In every nonempty evaluation fiber choose one
character; subtracting it injects that fiber into $X_\ell(H)$.
There are at most $\ell^{qd}$ fibers, proving both finiteness and
(CAP.13).  The proof includes $q=0$.

Only $q,\allowbreak{} H,\ell$ occur in the bound.  In particular one may use the
same integer
$h_b=|X_\ell(H)|\ell^{q[\Gamma:H]}$ at every level.\hfill$\square$

This proves the actual conclusion of
\nolinkurl{ChebotarevFamilies/HomBound.lean}:82--157.
% END CAP10

\subsection{Construction of every field of the original records}

% BEGIN CAP11
\paragraph{CAP11. Simultaneous congruence, approximation, and inverse Frobenius.}
\label{claim:CAP11}
Fix an arbitrary list $s_j\in G'$, $j\in\operatorname{Fin}(q)$.
There are, for every $N\geq0$, distinct new split primes $v_{Nj}$,
primes $w_{Nj}$ of $L$, valuations $B_{Nj}$, arithmetic Frobenius
elements $\sigma_{Nj}$, and tame generators $\tau_{Nj}$, such that
\[
 \begin{gathered}
 v_{Nj}\notin S,\qquad
 Q_{Nj}:=\operatorname{Norm}w_{Nj}\equiv1\pmod{\ell^N},\\
 \rho(\sigma_{Nj})^{-1}s_j^{-1}\in W_N.
 \end{gathered}
 \tag{CAP.14}
\]
With the actual level $V_N$ attached to
$S_N=S\cup\{v_{Nj}\}_j$, write $p_N$ for its map to the original
level, $K_N=p_N^{-1}H$, and
\[
 \phi_{Nj}=V_N.\mathrm{pr}(\sigma_{Nj})^{-1},\qquad
 t_{Nj}=V_N.\mathrm{pr}(\tau_{Nj}).
 \tag{CAP.15}
\]
Then $p_N(t_{Nj})=1$, (CAP.12) holds, and
\[
 (t_{Nj}^{Q_{Nj}})^{-1}
    (\phi_{Nj}^{-1}t_{Nj}\phi_{Nj})
       \in P_\ell^{\Gamma_N}(K_N).
 \tag{CAP.16}
\]

\emph{Proof.}
First, the original $V_0.\mathrm{pr}:G_L\to G_0$ is surjective.
Lift $s_j\in G_0$ through the quotient $G_{L^+}\to\Gamma$; its
$\eta$ value is $1$, so the exact $V_0.h\eta$ identifies the lift
as an element of the image of $G_L$.  Choose $g_j\in G_L$ with
$V_0.\mathrm{pr}(g_j)=s_j$.  Since $s_j\in G'$ and the original
$V_0.h\varepsilon$ is the full cyclotomic identity,
$\varepsilon_\ell(g_j)=1$ in $\mathbb Z_\ell^\times$, not merely
after residual reduction.

Set
\[
 U_N=\rho^{-1}(W_N)\cap\ker(\varepsilon_\ell\bmod\ell^N).
 \tag{CAP.17}
\]
This is open: $W_N$ is open, and the second character is continuous
with finite discrete target.  Apply CAP7 to $U_N$, the forbidden
set $S$, and the list $g_j^{-1}$.  The result satisfies
$\sigma_{Nj}g_j\in U_N$.  In particular
$\rho(\sigma_{Nj})s_j\in W_N$.  Normality of $W_N$, followed by
inversion, gives
\[
 \rho(\sigma_{Nj})^{-1}s_j^{-1}
  =\bigl(s_j(\rho(\sigma_{Nj})s_j)s_j^{-1}\bigr)^{-1}\in W_N,
\]
with the precise order of factors required in (CAP.14).

The prime $w_{Nj}$ is away from $\ell$, because otherwise its prime
$v_{Nj}$ below it would belong to $S$.  The cyclotomic image of an
arithmetic Frobenius there is $Q_{Nj}$ modulo $\ell^N$.
Here is the normalization check.  An $\ell^N$-th root of unity has
valuation one.  Arithmetic Frobenius sends its residue to its
$Q_{Nj}$-th power.  Reduction is injective on these roots: if
$u^{\ell^N}=1$ and $u\equiv1$ in the residue field, then
$(u-1)(1+u+\cdots+u^{\ell^N-1})=0$, and the second factor reduces
to the unit $\ell^N$, forcing $u=1$.  Thus Frobenius acts on the
roots themselves by the $Q_{Nj}$-th power, which is the claimed
cyclotomic value.  On the other hand
$\varepsilon_\ell(g_j)=1$ and $\sigma_{Nj}g_j\in U_N$ imply that
this value is $1$.  This proves the congruence, including the trivial
case $N=0$.

Now apply the retained $\mathrm{CohP.tame.1}$ to the number field
$L$, the given algebraic closure, $B_{Nj}$ over $w_{Nj}$, and the
prime $\ell$ which does not lie in $w_{Nj}$.  Its
$\mathrm{TameQuotient}$ clause supplies
$\tau_{Nj}$ with $\mathrm{IsProGen}$ and the exact relation
\[
 (\tau_{Nj}^{Q_{Nj}})^{-1}
   \sigma_{Nj}\tau_{Nj}\sigma_{Nj}^{-1}
       \in\mathrm{ineKer}(L,\allowbreak{} B_{Nj},\ell).
 \tag{CAP.18}
\]
CAP8 supplies $V_N$ and $p_N$.  CAP9 supplies its complete kernel
bound, $p_N(t_{Nj})=1$, and the inclusion of $\mathrm{ineKer}$ into
$P_\ell^{\Gamma_N}(K_N)$.  Mapping (CAP.18) to that quotient and
using $\phi_{Nj}^{-1}=V_N.\mathrm{pr}(\sigma_{Nj})$ gives exactly
(CAP.16).  The hypothesis $\mathrm{CohP}$ was retained throughout;
the tame relation is an explicit use of its original clause.\hfill$\square$

Source assembly and normalization:
\begin{quote}\small
\nolinkurl{AutomorphicDetection/ArPl.lean}:21--119;\\
\nolinkurl{ReciprocityTheory/GaloisValuations/Global.lean}:199--229;\\
tame contract: \nolinkurl{ContinuousCycMod.lean}:116--134.
\end{quote}
% END CAP11

% BEGIN CAP12
\paragraph{CAP12. The complete original $\mathrm{TWGrp}$ at all levels.}
\label{claim:CAP12}
For the fixed arbitrary list of CAP11, define $T$ by taking its
length to be $q$, its subgroup to be the exact $H$ of CAP3, its
listed elements to be $s_j$, and its level datum to be
$G_N=\mathrm{globOf}(r,\allowbreak{} V_N)$, with $h_N=p_N$ and the data
(CAP.14)--(CAP.15).  Give it the single bound
\[
 h_b=|X_\ell(H)|\ell^{q[\Gamma:H]}.
 \tag{CAP.19}
\]
These data satisfy every field of the original $\mathrm{TWGrp}$.

\emph{Proof.}
CAP3 provides $H$ open normal, $H_\eta$, and $H_\varepsilon$.
CAP8 provides every required $\mathrm{Glob.Hom}$ $h_N$ with the
same original local datum.  The list belongs to $G'$ by assumption.
The norm congruence is CAP11.  Its last relation in (CAP.14), together
with $W_N\subseteq H$, gives $\mathrm{ph\_\allowbreak H}$ at every $N$.
For any open normal $V\supseteq P_\ell^\Gamma(H)$ and any $j$,
CAP5 gives $W_N\subseteq V$ for all sufficiently large $N$; hence
$h_N(\phi_{Nj})s_j^{-1}\in V$ eventually.  This is the exact
$\mathrm{ph\_\allowbreak lim}$ quantifier, with $V$ and $j$ fixed before the
eventual bound.  In fact the constructed bound works for all $j$
at once.

The conditions $\mathrm{ta\_\allowbreak ker}$, $\mathrm{tame}$, and
$\mathrm{ker\_\allowbreak le}$ are CAP9 and CAP11.  Finiteness of $X_\ell(H)$
follows from CAP2.  CAP10 applied to every $p_N$ gives finiteness
of $X_\ell(K_N)$ and its bound by the same (CAP.19), which proves
both $\mathrm{fin0}$ and $\mathrm{finN}$.  Finally CAP2 supplies
$\mathrm{phiV}$ and $\mathrm{phiU}$ for all the original local
decomposition groups.  This exhausts the record fields in
\nolinkurl{ThetaCoefficientMapEq.lean}:15--40.

All $N\in\mathbb N$ are included, with primes and generators
chosen for each $N$.  The required compatibility is the map from
each level to the same original level, together with the displayed
uniform bound and eventual approximation.  No nestedness of the
sets of new primes is used.\hfill$\square$
% END CAP12

% BEGIN CAP13
\paragraph{CAP13. Actual arithmetic realization and the full original $\mathrm{ArPl}$.}
\label{claim:CAP13}
The system $T$ of CAP12 has an actual $\mathrm{TWPl}(r,\allowbreak{} T)$,
with its original length and list unchanged.  Consequently the entire
source predicate $\mathrm{ArPl}$ holds.

\emph{Proof.}
For every $N$ take the arithmetic level to be the already constructed
$V_N$.  Because $T.G_N=\mathrm{globOf}(r,\allowbreak{} V_N)$, take its comparison
$e_N$ to be the identity $\mathrm{Glob.Hom}$.  It is injective, and
its compatibility with $T.h_N$ is exactly (CAP.11), for every
$g\in G_{L^+}$.
Use the same $v_{Nj},\allowbreak{} w_{Nj},\allowbreak{} B_{Nj}$ of CAP11.  Their injectivity,
avoidance of $S$, actual place-over relations, and decomposition-group
splitting were proved in CAP7.  By construction
$V_N.S'=S\cup\operatorname{range}(v_N)$ and
$T.Nq_{Nj}=\operatorname{Norm}w_{Nj}$.

Set the record's Frobenius element to
$\mathrm{fr}_{Nj}=\sigma_{Nj}^{-1}$.  It belongs to the decomposition
group, and its \emph{inverse} is the arithmetic Frobenius.
This is precisely $\mathrm{fr\_\allowbreak frob}$ in the source.  Use the same
$\tau_{Nj}$ from CAP11, with its proved $\mathrm{IsProGen}$.
The comparisons for $\mathrm{ph}$ and $\mathrm{ta}$ are the
literal identities (CAP.15), since $e_N$ is the identity.
Finally $\mathrm{H\_\allowbreak eq}$ is (CAP.4), using the original residual
characters and the original cyclotomic character.  These checks cover
every field of \nolinkurl{TaylorWilesObstructionBound.lean}:88--111.

Before choosing the list, CAP2 already proved $\Phi_\ell(\Gamma)$.
For every $q$ and every list in $G'$, CAP12 and this construction now
give $T$ and its nonempty arithmetic realization.  The length equality
$T.q=q$ is definitional and its transport of the listed elements is
the identity.  This is exactly the conjunction and quantifier order
in the original \texttt{Ar\allowbreak Pl} definition at lines 119--124.
The argument permits $n=0$, $q=0$, and $N=0$; in the CountII
application it therefore applies in particular to $n=4$, $\ell=3$.
Besides the explicitly retained coefficient, level and residual data,
its original arithmetic hypothesis is the unchanged
$\mathrm{CohP}$.\hfill$\square$
% END CAP13

% BEGIN CAP14
\paragraph{CAP14. Contribution to the same-ideal CountII producer.}
\label{claim:CAP14}
The complete residual input \texttt{JSP\_\allowbreak ACTUAL\_\allowbreak ARPL} in CIP12 is
supplied by CAP1--CAP13.  The remaining sufficient same-ideal support
rule is JSP11 with CIP10--CIP11 and CAP13 inserted, retaining the
original \texttt{Choice}, first-order bound, obstruction bound,
\texttt{Raw}, and same-$J$ \texttt{Has\allowbreak Aut}, and all its original
\texttt{Coh\allowbreak P}/\texttt{Loc\allowbreak CFT}/coefficient/level/cyclotomic hypotheses.

\emph{Proof and scope.}
The original consumer first uses the $\Phi$ conjunct of $\mathrm{ArPl}$
to call $\mathrm{Choice}$; after Choice produces its finite list, it
uses the second conjunct to realize precisely that list at all levels.
CAP2 is independent of the list, and CAP13 is universal in it, so
this substitution introduces no circular dependence and no additional
restriction on the detector chosen by Choice.

ArPl has no ideal argument.  Therefore it works for every already
fixed original $B,\allowbreak{} r,\allowbreak{} J$, without replacing $J$ when a point, list, or
level is chosen.  The same-$J$ HasAut premise is still the original
universal assertion for every realized $\mathrm{TWGrp}$ and
$\mathrm{TWPl}$; constructing these records does not construct their
automorphic modules.  The original ordinary/control and separating
arithmetic-fiber data, the four further Hida fields, and this HasAut
producer remain separate obligations.

CAP13 is universal over the original coefficient ring, residual datum,
and Hida setting.  It adds no coefficient enlargement and no threshold.
Thus it can be inserted for each of both original rank-four Links
after the fixed $\iota$ and their one common coefficient threshold,
for every later coefficient/Space/universal datum.  Within each Link
the ideal stays the original $J$ of JSP1--JSP12; the two different
universal rings are not identified.  The common $\kappa$, the actual
distinct auxiliary primes required by LinkTwo, and $\ell=3$ are
unchanged.  QX4/QX10 and YX2 therefore retain exactly their original
conditional support-to-CountII application.  CountII remains needed
for RH and unrestricted Dirichlet GRH.  KA5--KA6 remove only CaseIII
for targets of two-power order; Cases I, II, and V retain their
arithmetic producers.\hfill$\square$

The exact consumer order is
\nolinkurl{ChebotarevFamilies/PFromAutOf.lean}:30--49,104--126;
the joint producer contracts are
\nolinkurl{DenseThresholdCases/HidaU.lean}:12--16,
\nolinkurl{PointTheoremBaseObject.lean}:26--48,
\nolinkurl{LinkTwo.lean}:13--51, and
\nolinkurl{PotentialModularityTower.lean}:38--66.
% END CAP14

\subsection{The complete residual Hecke-value subpacket}

For this subsection retain an arbitrary original level $V:s.\mathrm{Lev}$
and $Q:\mathrm{QData}(V)$, with its original finite auxiliary set.
Here no equality $V.S'=s.S$ is imposed.  Fix $n\geq1$, a coefficient
field datum $c$, and \emph{any supplied} representation
$\tau_0:K_T(n,\mathrm{wt}(s))\to\mathrm{GL}_m(c.O)$.
Let $r:\mathrm{Resid}(V.\mathrm{frame},\allowbreak{} c.O,\allowbreak{} n)$ and put
$k=c.O/\mathfrak m$, a finite field of characteristic $\ell$.
The original $\mathcal H=\mathrm{hset}(Q,\tau_0)$ has operator indices
\[
 \mathcal H.\mathrm{Op}
   =(\mathcal H.W_p\times\operatorname{Fin}(n))
      \sqcup(V.\mathrm{frame}.V_\ell\times\operatorname{Fin}(n)),
 \quad
 \mathcal H.W_p=
 \{w:w|_{L^+}\notin V.S',\ w\text{ split}\}.
 \tag{CAP.20}
\]
We construct one function $a:\mathcal H.\mathrm{Op}\to c.O$ with
both original fields of $\mathrm{IsRes}(Q,\allowbreak{} r,\tau_0,\allowbreak{} a)$.
The validity of the \emph{actual} $\mathrm{artOf}$ on each original
chosen local group is the already proved CIP4.

% BEGIN CAP15
\paragraph{CAP15. Every residual local character is constant on Artin fibers.}
\label{claim:CAP15}
Let $D$ be the decomposition group at a number-field prime, $W$ its
original Weil subgroup, and $b:W\twoheadrightarrow K^\times$ any map
satisfying the full original $\mathrm{IsLocalArtin}$.  If
$\psi:D\to M$ is an abelian-valued character with open kernel, then
\[
 b(x)=b(y)\quad\Longrightarrow\quad\psi(x)=\psi(y)
 \qquad(x,\allowbreak{} y\in W).
 \tag{CAP.21}
\]
This applies to all the original $r.\bar\psi_{v,\allowbreak{} i}$ and to the
residual character $\mathcal H.\varepsilon_v$.

\emph{Proof.}
The image $B=\psi(D)$ is finite, because its kernel is open in the
compact group $D$.  Characters $B\to\mathbb C^\times$ separate its
elements.  An elementary proof of the latter assertion is useful:
for $1\ne z\in B$, define a character on $\langle z\rangle$ by
sending $z$ to a primitive root of its order.  Extend a character
from a subgroup $C$ to $\langle C,\allowbreak{} t\rangle$ as follows.  If $e$ is
the least positive integer with $t^e\in C$, choose
$u\in\mathbb C^\times$ with $u^e=\chi(t^e)$ and set
$\chi'(ct^j)=\chi(c)u^j$.  The minimality of $e$ and this equation
prove independence of the expression and multiplicativity.  Iterating
over the finite group extends the original character to $B$.

For any such complex character $\lambda$, the composite
$\lambda\psi|_W$ is smooth: its kernel on inertia contains the
open kernel of $\psi$ restricted to inertia.  The literal
$\mathrm{IsLocalArtin.factor}$ therefore writes it as $\xi\circ b$
for a smooth character $\xi$ of $K^\times$.  If $b(x)=b(y)$, it
follows that every $\lambda$ is $1$ on $\psi(xy^{-1})$.
Character separation gives $\psi(x)=\psi(y)$.

The kernels of $r.\bar\psi_{v,\allowbreak{} i}$ are open by the original Resid
record.  For $\mathcal H.\varepsilon_v$, the continuous cyclotomic
character modulo $\ell$ has open kernel, and that kernel maps to $1$
in $k^\times$: the specified map $\mathbb Z_\ell\to k$ kills
$\ell\mathbb Z_\ell$.  Hence this residual character also has open
kernel.  CIP4 supplies every required field of the actual
$b=\mathrm{artOf}(s.hw(v))$, so (CAP.21) applies without replacing it
by a preferred reciprocity map.\hfill$\square$

The original local arguments are \nolinkurl{DenseThresholdCases/TResEx.lean}:65--104;
the factorization predicate is
\texttt{AutomorphicForms\textsubscript{2}/HasExponentZeroIff.lean}:145--152.
% END CAP15

% BEGIN CAP16
\paragraph{CAP16. The exact residual Hecke polynomial, independent of every Frobenius choice.}
\label{claim:CAP16}
For each $w\in\mathcal H.W_p$ there are elements
$b_{w,\allowbreak{} i}\in c.O$, $i\in\operatorname{Fin}(n)$, such that for
\emph{every} valuation $B$ over $w$ and every arithmetic Frobenius
$\sigma$ at $B$, the original residual polynomial satisfies
\[
 \bar D_r\bigl(X-[V.\mathrm{pr}(\sigma^{-1})]\bigr)
   =\sum_{j=0}^{n}
       (-1)^j q_w^{j(j-1)/2}\,\beta_{w,\allowbreak{} j} X^{n-j},
 \qquad
 \beta_{w,0}=1,\quad \beta_{w,\allowbreak{} j}=\bar b_{w,\allowbreak{} j-1}\ (1\leq j\leq n),
 \tag{CAP.22}
\]
where $q_w=\operatorname{Norm}w$.  The left side denotes the exact
source $\bar D_r.\chi([V.\mathrm{pr}(\sigma^{-1})])$.

\emph{Proof.}
Every character $\bar\chi_i\circ V.\mathrm{pr}$ is unramified at
$w$: since the prime below $w$ is outside $V.S'$, the original
$N_{V.S'}$ kills inertia at \emph{every} valuation over $w$.
For two valuations $B,\allowbreak{} B'$ over $w$, choose $g\in G_L$ carrying
$B$ to $B'$ by valuation transitivity.  Conjugate a Frobenius at
$B'$ by $g^{-1}$ to get a Frobenius at $B$.  Its ratio with a fixed
Frobenius there is an inertia element, as both have the same residue
power action.  An unramified abelian character kills this ratio and
is invariant under conjugation.  Thus the value of each
$\bar\chi_i(V.\mathrm{pr}(\sigma^{-1}))$ is independent of both
$B$ and $\sigma$, with the inverse exactly as displayed.

Choose one such Frobenius $\sigma_w$ and form the monic degree-$n$
polynomial
\[
 P_w(X)=\prod_{i\in\operatorname{Fin}(n)}
     \bigl(X-\bar\chi_i(V.\mathrm{pr}(\sigma_w^{-1}))\bigr).
\]
It is the required determinant polynomial because the residual
representation is diagonal.  The prime $w$ is away from $\ell$:
every prime of $L^+$ above $\ell$ belongs to $s.S\subseteq V.S'$.
Thus $q_w$ is an $\ell$-adic unit and is nonzero in $k$.
For $1\leq j\leq n$ define
\[
 \beta_{w,\allowbreak{} j}=(-1)^j q_w^{-j(j-1)/2}
                         [X^{n-j}]P_w(X).
 \tag{CAP.23}
\]
Choose $b_{w,\allowbreak{} j-1}\in c.O$ lifting this element by surjectivity of
the residue map.  Substitution into the right side of (CAP.22)
recovers every coefficient of $P_w$: the leading coefficient is $1$,
and the two sign factors and two powers of $q_w$ cancel for every
$j\geq1$.  The Frobenius independence already proved gives (CAP.22)
for all choices.  This is the literal $\mathrm{heckePol}$
normalization, including its factor $q_w^{j(j-1)/2}$ and its
operator-index shift $j\mapsto j-1$.  No factorial or inverse of
$\ell$ is used.\hfill$\square$

The polynomial is defined in \nolinkurl{HeckeAlgebras/Distinct.lean}:46--47;
the corresponding source construction and independence arguments are
\nolinkurl{DenseThresholdCases/TResEx.lean}:16--64,109--140.
% END CAP16

% BEGIN CAP17
\paragraph{CAP17. Ordered uniformizer values and the entire original $\mathrm{IsRes}$.}
\label{claim:CAP17}
For every original $V,\allowbreak{} Q,\allowbreak{} c,\allowbreak{} m,\tau_0,\allowbreak{} r$ in the scope of (CAP.20), there
exists one $a:\mathcal H.\mathrm{Op}\to c.O$ satisfying the full
original $\mathrm{IsRes}(Q,\allowbreak{} r,\tau_0,\allowbreak{} a)$.

\emph{Proof.}
At each original $v\mid\ell$, CIP4 gives surjectivity of the
actual $\mathrm{artOf}$ onto $L_{s.w(v)}^\times$.  Choose a Weil
element $y_v$ whose image is the original chosen uniformizer and
regard it as an element of the original decomposition group.
For $j\in\operatorname{Fin}(n)$ set
\[
 u_{v,\allowbreak{} j}=\prod_{\substack{i\in\operatorname{Fin}(n)\\i\leq j}}
       r.\bar\psi_{v,\allowbreak{} i}(y_v)\,
               \mathcal H.\varepsilon_v(y_v)^i\ \in k^\times.
 \tag{CAP.24}
\]
CAP15 shows that every factor is unchanged if $y_v$ is replaced by
\emph{any} Weil element having that same Artin image.  Hence (CAP.24)
is independent of this choice.  Choose $d_{v,\allowbreak{} j}\in c.O$ whose residue
is the underlying field element of $u_{v,\allowbreak{} j}$.

Define the single function on the original disjoint union of operator
indices by
\[
 a(\mathrm{inl}(w,\allowbreak{} i))=b_{w,\allowbreak{} i},\qquad
 a(\mathrm{inr}(v,\allowbreak{} j))=d_{v,\allowbreak{} j}.
 \tag{CAP.25}
\]
For the first field of $\mathrm{IsRes}$, CAP16 proves precisely its
polynomial equation for every $w,\allowbreak{} B,\sigma$.  For the second, let
$x$ be any Weil element and $g$ any element of the recorded
decomposition group having the same underlying Galois automorphism,
with $\mathrm{artOf}(x)$ equal to the chosen uniformizer.  Then $g$
is the image of $x$ under the original Weil inclusion.  CAP15 replaces
every factor at $g$ by the corresponding factor at $y_v$.  Equation
(CAP.24), followed by the choice of $d_{v,\allowbreak{} j}$, is exactly the
source's ordered product over $i\leq j$, with exponent $i$, and
proves its \texttt{unif} field.  These are both fields of
\nolinkurl{HeckeAlgebras/IsTheta.lean}:161--168.

The choices of lifts are made simultaneously in the same original
ring $c.O$; no extension of coefficients is needed.  This proves
the complete conclusion of the original \texttt{TRes\allowbreak Ex} supplier,
at every original level and for every supplied $\tau_0$.  Its source
use of a literature hypothesis merely to obtain $\mathrm{artOf}$
is replaced here by CIP4 with its checked class field theory input.
In constructing $R:\mathrm{UnitaryHeckeData}$ this supplies its
literal fields $a$ and $\mathrm{res}$ after $\tau_0$ is supplied.
The chosen $a$ is then held fixed for the ordinary module and its
future kernel ideal.  No equality of modules or ideals associated
with different choices of $a$ is asserted.

The integral type with $\mathrm{TypeForAlt}$, the torus character
with both $\mathrm{IsTheta}$ and the full $\mathrm{IsWt}$ (including
its actual neat-level condition), and the remaining original
ordinary control, determinant, and arithmetic-fiber producers remain open at their
previous scopes.  In particular CAP17 does not supply all of $R$
or close $\mathrm{COUNT\_\allowbreak II\_\allowbreak SUPPORT}$.\hfill$\square$

The source conclusion is \nolinkurl{DenseThresholdCases/TResEx.lean}:105--165;
its precise insertion is \nolinkurl{UnitaryHeckeData.lean}:59--70 and
\nolinkurl{HdEx.lean}:20--27.  The retained torus-character and full
weight supplier is \nolinkurl{TThetaEx.lean}:160--180.
% END CAP17


\section{Why simultaneous local exceptions have product size}\label{sec:joint-rank}

\status{The passage from single-place tests to the joint exceptional-generator bound is proved here. The polynomial lattice comparisons are proved in Section~\ref{sec:edge-lattices}; Sections~\ref{sec:symbols}--\ref{sec:digit-tests} now prove the generator, lattice-comparison and pointwise Gram/test constructions; incoming word closure is proved in IN/WT, and RK supplies the selected-symbol dimension. The local Type-II input is now proved in Corollary~\ref{cor:typeII}; the remaining global inputs are listed in Section~\ref{sec:remaining}. This section does not repeat the conversion of elementary-divisor tails to index.}

Work over a DVR $\mathcal O$ with fraction field $K$, maximal ideal $\mathfrak m$, and a nonzero nonunit $p$; in the source $w^2=p$. Fix an integer $D\ge0$. The base and dual coefficient lattices are finite free of constant rank $d_v$. At a target lattice $L$, the relevant dual-slot lattice is
\[
 P_v(L)=N_{\mathrm{tgt},v}(L)^\vee
   =\{x:\ell(x)\in\mathcal O\text{ for every }\ell\in N_{\mathrm{tgt},v}(L)\},
 \qquad \operatorname{rank}P_v(L)=d_v.
\]
The local input supplies finitely many incoming edges $e_i=(L_i^-,L)$ and vectors $F_i$ in the source lattice at $L_i^-$. Put $f_i=p^{A_v}F_i$, where $A_v$ is the source's sum of the first $r$ weights. The integral pairing statements imply $f_i\in P_v(L)$. Besides vanishing to order $w^{c_*}$ at every other target of $L_i^-$, the input supplies $n_v\ge(1-\kappa_v)d_v$, an $\mathcal O$-linear map $x\mapsto\ell_x$ into $N_{\mathrm{tgt},v}(L)$, and an integral matrix $C$ such that
\begin{equation}\label{eq:jr-tests}
 \sum_i C_{ji}\ell_x(f_i)=p^D x_j
 \qquad(x\in\mathcal O^{n_v},\ 1\le j\le n_v).
\end{equation}
Here $\kappa_v\ge0$ and $c_*\ge4(D+1)$. This is the concrete rank content of \texttt{LPprime}, rather than an assumed independence of exceptional subspaces. \textsf{[JR1, JR12]}

\begin{lemma}[Extracting good integral directions]\label{lem:jr-smith}
Under \eqref{eq:jr-tests}, $P_v(L)$ has a basis $b_{v,a}$ with at most $\kappa_vd_v$ bad indices, such that
\[
 p^D b_{v,a}\in\sum_i\mathcal O f_i
 \quad\text{for every good index }a.
\]
\end{lemma}
\begin{proof}
Put $q=p^D$ and apply Smith normal form to $\sum_i\mathcal O f_i\subseteq P_v(L)$, with nonzero elementary divisors $a_i$. The integral functionals $\ell_x$ and \eqref{eq:jr-tests} give an integral matrix identity $C'\operatorname{diag}(a_i)X=qI_{n_v}$. Call $i$ good if $a_i\mid q$. For every other $i$, valuation comparability gives $a_i=q\mu_i$ with $\mu_i\in\mathfrak m$. Moving these terms to the right gives
\[
 \sum_{i\ \mathrm{good}} C'_i a_i X_i=q(I-Y),
 \qquad Y\in M_{n_v}(\mathfrak m).
\]
The right side has rank $n_v$ over $K$, since $I-Y$ has unit determinant. The left side is a sum of at most one rank-one matrix per good index. Hence there are at least $n_v$ good indices. For their ambient Smith basis vectors, $q b_{v,a}$ lies in the generated sublattice. The other at most $d_v-n_v\le\kappa_vd_v$ indices are bad. \textsf{[JR4]}
\end{proof}

To globalize, use the finite set $\Omega$ of orbits of decorated lattice tuples. The acting groups have trivial stabilizers and preserve the coefficient lattices. Consequently a value $\xi$ at a representative $L$ extends uniquely to an equivariant section $\delta_L(\xi)$ supported on that orbit. Let $M_B$ use the dual coefficient lattice at coordinates $v\in B$ and the base coefficient lattice elsewhere. Let $J_B\subseteq M_B$ be the sum, over $v\in B$, of the neighbor-operator images with the source lattice in coordinate $v$. The neighbor operator itself includes the factor $p^{A_v}$; $A_v$ may be negative.

For a good index at $v$, tensor its relation in the lemma with arbitrary integral vectors in the other coordinates. Applying the neighbor operator to the corresponding section at $L_i^-$ produces the desired section at $L$ and terms at the other neighbors. The desired term occurs once, because stabilizers are trivial. All other terms lie in $w^{c_*}M_B\subseteq p^{D+1}M_B$. Equivariance gives the same assertion at every translate. Therefore every elementary tensor with at least one good selected coordinate satisfies
\begin{equation}\label{eq:jr-one-good}
 p^D\delta_L(b_{1,a_1}\otimes\cdots\otimes b_{s,a_s})
       \in J_B+p^{D+1}M_B.
\end{equation}
This cancellation uses the integral neighbor-map inclusion and equivariance proved in Section~\ref{sec:edge-lattices}, discharging the polynomial component formerly retained in \textsf{[JR12]}. The cancellation itself is proved by the preceding evaluation argument. \textsf{[JR5]}

\begin{proposition}[Joint exceptional generators]\label{prop:jr-joint}
Suppose these local tests and lattice properties hold for every selected coordinate, with common $D,c_*$. Write $d=\operatorname{rank}M_{\mathrm{na}}$ for the full base function lattice. For every nonempty set $B$ of selected coordinates, there is a finite set $\mathcal C_B\subseteq M_B$ with
\begin{equation}\label{eq:jr-joint}
 |\mathcal C_B|\le d\prod_{v\in B}\kappa_v,
 \qquad
 p^{D+1}M_B\subseteq J_B+\operatorname{span}_{\mathcal O}\mathcal C_B.
\end{equation}
\end{proposition}
\begin{proof}
At each orbit representative, retain precisely the elementary tensors whose index is bad at \emph{every} coordinate of $B$. Their number is at most
\[
 \left(\prod_{v\in B}\kappa_vd_v\right)
 \left(\prod_{v\notin B}d_v\right)
   =\left(\prod_{v\in B}\kappa_v\right)\prod_v d_v.
\]
Every other elementary tensor satisfies \eqref{eq:jr-one-good}. Delta-extend the retained tensors and take their union over $\Omega$ to obtain $\mathcal C_B$. The base coefficient tensors, similarly extended, are linearly independent: evaluate a relation at each representative. Thus $|\Omega|\prod_vd_v\le d$, proving the cardinality bound. \textsf{[JR6, JR7]}

With $T=J_B+\operatorname{span}_{\mathcal O}\mathcal C_B$, expansion in the tensor bases gives $p^DM_B\subseteq T+p^{D+1}M_B$. The module
\[
 Q=(p^DM_B+T)/T
\]
is finitely generated and satisfies $Q=pQ$. Nakayama's lemma gives $Q=0$. Multiplying the resulting inclusion by $p$ gives exactly \eqref{eq:jr-joint}, the exponent advertised by the source. \textsf{[JR7]}
\end{proof}

The product appears because a tensor survives only when \emph{all} selected coordinates are bad. Arbitrary subspace intersections do not satisfy such a product bound. Summing the displayed products over $|B|=k$ gives the elementary symmetric polynomial $e_k(\kappa)$; the subsequent lattice walk supplies the additional multiplicities and the exponent of $p$.

\paragraph{What remains at one place.}
The supplementary report proves how primitive pairings isolated at different incoming edges combine to give \eqref{eq:jr-tests}, and constructs $q_0^{r(k-r)}$ distinct edges from graph planes over $\mathbf F_{q_0}$. \textsf{[JR2, JR3, JR8]} For $r\ge2$, FG/OR/CP/EA now prove the full local Type-II input, including frames, selected-symbol rank, outgoing realization, primitive dual tests and all numerical compatibility. The separate $r=1$ treatment is now completed by RU/RD/HM/NP/TC, with the full $\texttt{Conclusion2}$ under literal $\texttt{HypI}$. GI/GII now supply the uniform numerical and good-orbit hypotheses for a supplied family, and BC/GJ give the finite-place application. Constructing all the global data together remains separate. Neither local conclusion is inferred merely from counting planes. \textsf{[JR9, JR10]}

\section{Local divisibility with its exact lattice costs}\label{sec:local-accounting}

\status{The joint-rank assembly in Section~\ref{sec:joint-rank}, the lattice walk below, and the tail-to-index calculation form a complete conditional chain. The actual lattice compatibility is now proved in Section~\ref{sec:edge-lattices}; the remaining quantitative test inputs are specified in Section~\ref{sec:remaining}, after the pointwise constructions of Sections~\ref{sec:path-generators}--\ref{sec:digit-tests}.}

\subsection{Why the local exceptional ranks give Smith tails}\label{sec:lattice-walk}

The step between a joint rank estimate and Proposition~\ref{prop:local-index} is a finite lattice argument. The following is the source's \texttt{WalkSharp}, specialized to its actual parameters; the proof explains both the exponent $m^{\beta+1}$ and the depth $\tau+\beta S$. \textsf{[LW1--LW3]}

Work over a DVR $\mathcal O$ containing $w$ with $w^2=p$. Order the local indices $1,\ldots,V$. There is a compatible lattice $L_\sigma$ for each choice of a base, ball, or dual slot at every index. Write $M=L_{\rm base}$ and $L_B$ for dual slots on $B$, base slots elsewhere. With the other slots fixed, assume
\[
 L_{\rm dual}\subseteq L_{\rm base}\subseteq L_{\rm ball},\qquad
 w^{t_v}L_{\rm ball}\subseteq L_{\rm base},\qquad
 w^{t_v}L_{\rm base}\subseteq L_{\rm dual},\qquad
 y_vL_{\rm ball}\subseteq L_{\rm dual}.
\]
The $y_v$ commute. In particular they preserve $M$. Suppose, for every nonempty $B$, that a finite set $G_B\subseteq L_B$ satisfies
\begin{equation}\label{eq:walk-joint}
 p^{D_0+1}L_B\subseteq
 \sum_{v\in B}y_vL_{B;v=\rm ball}+\langle G_B\rangle,
 \qquad |G_B|\le d\prod_{v\in B}\kappa_v,
\end{equation}
where $d=\operatorname{rank}M$ and $\kappa_v\ge0$. The preceding joint-rank argument supplies this estimate from the single-place inputs; it is not a statement about arbitrary intersecting subspaces.

For integers $m\ge1$ and $t_v,D_0,S_v,S,\beta\ge0$, put $h_v=t_v+D_0+1$, $\tau=m\sum_vh_v$, and
\[
 N=\sum_{A\subseteq\{1,\ldots,V\}}
 p^{\sum_{v\notin A}S_v}\Bigl(\prod_{v\in A}y_v^m\Bigr)M,
 \qquad 0\le S_v\le S.
\]
Then $N\subseteq M$, $p^{\sum_vS_v}M\subseteq N$, and for every $\beta\ge0$,
\begin{equation}\label{eq:walk-tail}
 p^{\tau+\beta S}M\subseteq N+\langle C_\beta\rangle,
 \qquad |C_\beta|\le d\,m^{\beta+1}e_{\beta+1}(\kappa).
\end{equation}
For an integral count one may take the floor of the right-hand side. This is precisely the exceptional-direction hypothesis used below.

\begin{proof}
Set $q=\beta+1$. If $q>V$, then $\sum S_v\le\beta S$ and the assertion follows with $C_\beta=\varnothing$. Otherwise keep $q$ active indices. A state consists of a set $B$ of size $q$ and counters $0\le i_v<m$ on $B$. The indices less than $\max B$ and outside $B$ are completed; larger indices are waiting. Define
\[
 Y_{B,i}=\prod_{u\,\rm completed}y_u^m\prod_{v\in B}y_v^{i_v},
 \quad
 A_{B,i}=\sum_{u\,\rm completed}2mh_u+
                 \sum_{v\in B}(t_v+2i_vh_v).
\]
Use as $C_\beta$ the union of all $Y_{B,i}G_B$. Every vector belongs to $M$, and each $B$ has exactly $m^q$ choices of counters. Hence its cardinality is at most
$m^q\sum_{|B|=q}|G_B|\le d m^q e_q(\kappa)$.

Let $E=2(\tau+\beta S)$. We prove the state invariant
\[
 w^{E-A_{B,i}}Y_{B,i}L_B\subseteq N+\langle C_\beta\rangle
\]
backwards in the sum of the operator exponents, which strictly increases at every transition and is bounded by $mV$. Apply \eqref{eq:walk-joint}, spending $2(D_0+1)$ powers of $w$. Its exceptional term is already in $\langle C_\beta\rangle$. Each other term has the form $Y_{B,i}y_vz$ with $z\in L_{B;v=\rm ball}$.

If $i_v+1<m$, multiply $z$ by $w^{2t_v}$ to return from ball to dual, and increase $i_v$ by one. The budget increases by $2h_v$, exactly the sum of these two costs. If $i_v+1=m$ and an index is waiting, return the completed $v$ from ball to base at cost $t_v$, enter the next waiting index $u$ in its dual slot at cost $t_u$, and replace $v$ by $u$ in $B$. The budget increases by
$2(D_0+1)+t_v+t_u$, again exactly the amount spent.

Finally suppose $i_v+1=m$ and no index is waiting. All but the remaining $\beta$ active indices now carry $y_u^m$. After paying $t_v$ to return to base, pay
$2\sum_{u\in B\setminus\{v\}}S_u\le2\beta S$ to use the corresponding generator of $N$. Each completed index has spent $2mh_u$; an active index has spent $t_u+2i_uh_u\le2mh_u$, with enough unused budget for the next decomposition or final exit. Thus all the displayed powers of $w$ are nonnegative, and the backwards induction is valid.

Start with $B=\{1,\ldots,q\}$ and all counters zero. Entering its dual slots costs $\sum_{v\in B}t_v=A_{B,0}$, so the invariant gives $w^EM\subseteq N+\langle C_\beta\rangle$. Since $w^2=p$, this is \eqref{eq:walk-tail}.
\end{proof}

The three slots themselves have a simple lattice interpretation. Locally the ball is
$w^{-t}H\cap\bigcap_eR_e$, and the dual slot is the integral polar of
$w^{-t}H^\vee\cap\bigcap_eR'_e$. Once $H\subseteq R_e$ and $H^\vee\subseteq R'_e$ are known, reflexivity gives the four displayed lattice inclusions. Tensoring and taking integral equivariant functions preserves them. Section~\ref{sec:edge-lattices} proves the actual edge inclusions, fullness, covariance, and compatibility of the adjacency operators with the polar pairing. These are no longer retained coefficient-lattice inputs. \textsf{[LW4]}

\subsection{From exceptional directions to a lattice index}

Let $\mathcal O$ be a discrete valuation ring with uniformizer $\pi$. Let $M$ be free of rank $r>0$, and let $N\subseteq M$ be a full lattice. For $a\ne0$ in the maximal ideal put $e=v_\pi(a)$. Say that $N$ has \emph{at most $g$ exceptional directions at depth $E$} if there are $g$ vectors in $M$ whose span, together with $N$, contains $a^E M$. This is the source's \texttt{AllButLE}; using fewer than $g$ vectors makes no difference.\cite{local-selection,local-smith}

Choose a Smith basis, so $M=\bigoplus_i\mathcal O e_i$ and $N=\bigoplus_i\pi^{c_i}\mathcal O e_i$. The condition is equivalent to
\[
\#\{i:c_i>eE\}\le g.
\]
For one direction, retain the corresponding basis vectors. Conversely, the summands with $c_i>eE$ require that many generators after reduction at depth $eE+1$; a submodule of a $g$-generated module over a DVR has at most $g$ generators. This is the elementary-divisor content of the formal statement. \textsf{[LC1]}

\begin{proposition}\label{prop:local-index}
Suppose $\tau,S\ge0$ are integers and $\kappa_v\ge0$ is a finite family. Write
\[
e_j(\kappa)=\sum_{|B|=j}\prod_{v\in B}\kappa_v,
\qquad \Sigma=m\sum_v\kappa_v,
\]
where $m\ge1$ is an integer. Suppose that for every $\beta\ge0$ the number of exceptional directions at depth $\tau+\beta S$ is at most
$r m^{\beta+1}e_{\beta+1}(\kappa)$. Then, for $N_*:=N\cap a^\tau M$,
\begin{equation}\label{eq:local-index-exact}
\operatorname{length}_{\mathcal O}(M/N_*)
\le er\left[\tau+S\left\{\prod_v(1+m\kappa_v)-1\right\}\right]
\le er\{\tau+S(e^\Sigma-1)\}.
\end{equation}
In particular, if $\Sigma\le1/2$, then the length is at most $er(\tau+2S\Sigma)$.
\end{proposition}

\begin{proof}
The elementary divisors of $N_*$ have exponents $\max(c_i,e\tau)$, so the length is
\[
er\tau+\sum_i(c_i-e\tau)_+.
\]
For $S>0$, divide the tail of each exponent into intervals of length $eS$. On the $\beta$th interval, the number of surviving exponents is at most $r m^{\beta+1}e_{\beta+1}(\kappa)$. Thus the sum of the tails is at most
\[
eSr\sum_{j\ge1}m^j e_j(\kappa)
=eSr\left\{\prod_v(1+m\kappa_v)-1\right\}.
\]
The sums are finite, since $e_j(\kappa)=0$ beyond the number of indices. If $S=0$, this eventual zero bound forces every $c_i\le e\tau$, giving the same conclusion. Finally $1+y\le e^y$, and for $0\le\Sigma\le1/2$,
$e^\Sigma-1\le\Sigma/(1-\Sigma)\le2\Sigma$.
\end{proof}

The last inequality is the source's index allowance. The product and exponential bounds in \eqref{eq:local-index-exact} retain information that its proof discards; these are separately derived improvements. There is also a useful saturated refinement. If $M'\subseteq M$ is saturated of rank $r'$, then
\[
\operatorname{length}\bigl(M'/(M'\cap N_*)\bigr)
\le er'\tau+erS\left\{\prod_v(1+m\kappa_v)-1\right\}.
\]
The baseline $er'\tau$ is exact; the remaining quotient embeds in $a^\tau M/(N\cap a^\tau M)$, whose length is the tail sum just estimated. Section~\ref{sec:saturation} checks the reference-lattice transports on the actual route and applies this refinement. It also absorbs the remaining rank ratio using the sharper tail bound. \textsf{[LC2]}

In the actual local selection argument, the candidate lattice is generated by
\[
a^{\sum_{v\notin B}S_v}\Bigl(\prod_{v\in B}y_v^m\Bigr)M,
\qquad B\subseteq\mathcal V,
\]
where the $y_v$ commute, $S_v\le S$, and
$\tau=\sum_{v\in\mathcal V}m(t_v+D_0+1)$. If necessary replace the generated lattice by its intersection with $M$; this leaves both the exceptional-direction hypothesis and $N_*$ unchanged. Sections~\ref{sec:joint-rank} and \ref{sec:lattice-walk} prove that hypothesis from the single-place tests and the compatible lattice system. The present proposition supplies its elementary-divisor consequence. The integral path-generator, Weyl comparison and pointwise constrained Gram/test constructions are proved in Sections~\ref{sec:path-generators}--\ref{sec:digit-tests}; IN/WT prove incoming membership, RK proves the selected-symbol rank, and FG/OR/CP/EA supply the entire local Type-II theorem. TC now supplies the separate Type-I theorem. The global arithmetic obligations remain.\cite{local-systems}

\subsection{Normalization and rounding over a number field}

Let $K$ be a number field of degree $k_0$, and for a finite place $w$ write
\[
\ell_w=\frac{\log Nw}{k_0},\qquad e_w=v_w(p)\quad(w\mid p),
\qquad \operatorname{sz}_w(y)=2\ell_w e_w y.
\]
The fundamental identity $\sum_{w\mid p}e_w f_w=k_0$ gives
\begin{equation}\label{eq:prime-normalization}
\sum_{w\mid p}\operatorname{sz}_w(y)=2y\log p.
\end{equation}
No Galois hypothesis is needed for this identity. It explains how local normalized gains become rational-prime sums without an extra factor from the degree of the coefficient field. \textsf{[LC3]}\cite{local-normalization}

For the slope statements take $K$ totally complex and use the source's squared-norm convention. For finitely supported real shifts $s_u(w)$, set
\[
I_u=\prod_w\mathfrak p_w^{-\lfloor e_ws_u(w)\rfloor},\qquad
\mathcal R_u=2\sum_w\ell_w\{e_ws_u(w)\},
\]
where $\{z\}=z-\lfloor z\rfloor\in[0,1)$, also for negative $z$. The slope identity, with the same archimedean metrics, is exact:
\begin{equation}\label{eq:round-exact}
\mu(I_uN_u)-\mu(N_u)
=\sum_w\operatorname{sz}_w(s_u(w))-\mathcal R_u.
\end{equation}
Indeed multiplication by $I_u$ multiplies covolume by $N(I_u)^{\operatorname{rank}N_u}$; take its normalized logarithm. Thus rounding is a nonnegative loss, not an unsigned unspecified error. The floor is taken \emph{after} multiplying by $e_w$.\cite{local-div-defs}

If the nonzero rounding terms lie over primes $p\le X$ and at most three exceptional primes $g_1,g_2,T_u$, then
\[
\mathcal R_u\le2\left(\vartheta(X)+\log g_1+\log g_2+\log T_u\right),
\qquad \vartheta(X)=\sum_{p\le X}\log p.
\]
Here $\sum_{w\mid p}\ell_w\le\log p$, since $e_w\ge1$. The elementary bound $\vartheta(X)\le X\log4$ follows by induction from the central binomial coefficients: use $\binom{2n}{n}\le2^{2n}$ for even arguments and $\binom{2n+1}{n}\le2^{2n}$ for odd ones, with the prime intervals $(n,2n]$ and $(n+1,2n+1]$, respectively. Therefore if each exceptional logarithm is at most $X$,
\begin{equation}\label{eq:round-source}
0\le\mathcal R_u\le(2\log4+6)X.
\end{equation}
This reconstructs the source's explicit rounding allowance. Its required support and type-prime bounds remain hypotheses until supplied by the chosen family. \textsf{[LC4]}\cite{local-rounding}

\subsection{The precise net-gain identity}

The source defines the real shift by
\[
s_u(w)=\sum_j\{g_j(w)+F_j(w)-E_{u,j}(w)\},
\]
where $g_j$ already incorporates the discarded-prime convention. Define
\begin{align*}
\mathrm{Ga}&=\sum_w\operatorname{sz}_w\!\left(\sum_jg_j(w)-\tau_w\right),&
\mathrm{Fo}&=\sum_w\operatorname{sz}_w\!\left(\sum_jF_j(w)\right),\\
\mathrm{Ee}_u&=\sum_w\operatorname{sz}_w\!\left(\sum_jE_{u,j}(w)\right),&
\mathrm{IO}&=\sum_w\operatorname{sz}_w(\mathrm{io}_w).
\end{align*}
Also put $\mathrm{Ta}=\sum_w\operatorname{sz}_w(\tau_w)$. Let $M_u$ be the retained reference source lattice of rank $r_u>0$, and let $d_u\ge r_u$ be the full column rank. Put $\rho_u=d_u/r_u$. The source bounds the local restriction index using the \emph{full} rank:
\[
\operatorname{length}((M_u)_w/(N_u)_w)
\le e_w d_u(\tau_w+\mathrm{io}_w).
\]
Assume this at the regular places, and write $\mathcal B_u$ for the remaining nonnegative bound on the excess signed relative index at exceptional places. The determinant formula for lattice index, followed by \eqref{eq:round-exact}, gives
\begin{equation}\label{eq:net-gain-exact}
\begin{split}
\mu(I_uN_u)\ge{}&\mu(M_u)+\mathrm{Ga}+\mathrm{Fo}-\mathrm{Ee}_u\\
&-\rho_u\mathrm{IO}-(\rho_u-1)\mathrm{Ta}-\mathcal R_u-\mathcal B_u.
\end{split}
\end{equation}
This is the source's coarser allowance. The ratio is required by that full-rank estimate alone; Section~\ref{sec:saturation} proves a stronger actual-lattice estimate that removes both displayed rank inflations. If $r_u\ge(1-d)d_u$, $0\le d<1$, and $\tau_w,\mathrm{io}_w\ge0$, the two index terms are bounded by
\[
\rho_u\mathrm{IO}+(\rho_u-1)\mathrm{Ta}
\le\frac{\mathrm{IO}}{1-d}+\frac{d\,\mathrm{Ta}}{1-d}.
\]
This is exactly the structure of the source's \texttt{ClaimLow}. At an exceptional place use the signed relative index
\[
\operatorname{length}(M_u/(M_u\cap N_u))_w
-\operatorname{length}(N_u/(M_u\cap N_u))_w;
\]
its nonnegative excess beyond $e_wd_u(\tau_w+\mathrm{io}_w)$, divided by $r_u$ and weighted by $2\ell_w$, defines $\mathcal B_u$. This convention also covers a local enlargement instead of a restriction. \textsf{[LC5]}

In the released construction the word's own type prime $T_u$ is explicitly excluded from the regular-place restriction hypothesis. Its contribution must therefore be bounded separately. The source packages this and other permitted losses in \texttt{Lo5}; the present argument does not set them to zero. Equations \eqref{eq:round-exact}--\eqref{eq:net-gain-exact} give the precise bookkeeping behind the schematic net gain of version 0.1.\cite{local-div-defs,local-costs}

\paragraph{Absorbing the ramified rank and mean corrections.}
For an integer $f\ge0$, the source also supplies the separate identities
\[
\frac{r_u}{d_u}=\left(1-\frac1{T_u-1}\right)^f,
\qquad
\mu(M_u)=\mu(M_u^{\rm full})-
\frac{f\log(T_u-1)}{T_u-2},\qquad T_u>2.
\]
The first gives $\delta_u:=1-r_u/d_u\le f/(T_u-1)\le2f/T_u$ by Bernoulli's inequality. The second is an additional metric input, not a consequence of the rank identity. Assume $\mathrm{IO}\ge0$, $0\le\mathrm{Ta}\le b^2$, $2f/T_u\le1/2$, and
\[
\frac{4fb^2}{T_u}+\frac{f\log(T_u-1)}{T_u-2}\le1.
\]
Then $\mathrm{IO}/(1-\delta_u)\le2\mathrm{IO}$, while
$\delta_u\mathrm{Ta}/(1-\delta_u)\le4fb^2/T_u$.
Thus \eqref{eq:net-gain-exact} implies
\[
\mu(I_uN_u)\ge\mu(M_u^{\rm full})+\mathrm{Ga}+\mathrm{Fo}-\mathrm{Ee}_u
-\bigl(2\mathrm{IO}+\mathcal B_u+\mathcal R_u+1\bigr).
\]
When $\mathrm{io}_w$ depends only on the rational prime $p$ below $w$, put $\mathrm{IOp}=\sum_p\mathrm{io}_p\log p$. Equation~\eqref{eq:prime-normalization} gives $\mathrm{IO}=2\mathrm{IOp}$, yielding the numerical $4\mathrm{IOp}$ allowance in \texttt{ClaimLow2}. This accounts for the ramified corrections explicitly without asserting their underlying metric theorem. \textsf{[LC6]}

% Inclusion fragment. Exact source records: augmentation_provenance.json.
\paragraph{LC7. The augmentation model for the retained rank and mean.}
Let $K$ be totally complex, let $W$ be a nonzero hermitian $\mathcal O_K$-lattice
of rank $h$, and let $T\ge3$, $f\ge0$ be integers.  Give
$U=\mathcal O_K^{T-1}$ the standard orthonormal metric at every infinite
place, and give tensor products their product hermitian metrics.  Set
\[
 A=\ker\bigl(U\xrightarrow{\sum}\mathcal O_K\bigr),\qquad
 M_{\rm full}=W\otimes U^{\otimes f},\qquad
 M_{\rm ret}=W\otimes A^{\otimes f}.
\]
The integral decomposition $U=A\oplus\mathcal O_K e_{T-1}$ shows that
$M_{\rm ret}$ is exactly the joint kernel of coordinate summation in
$M_{\rm full}$.  It follows that
\[
 \frac{\operatorname{rk}M_{\rm ret}}{\operatorname{rk}M_{\rm full}}
 =\left(\frac{T-2}{T-1}\right)^f,
 \qquad
 \delta=1-\left(1-\frac1{T-1}\right)^f
 \le\frac f{T-1}\le\frac{2f}{T}.
\]
The first inequality is Bernoulli's inequality.

Use the source's normalized squared-metric slope: the standard
$\mathcal O_K$ has slope zero, and a free lattice of rank $r$ with
hermitian Gram matrices $G_v$ in an integral basis has
$\mu=-\frac1{r\,d}\sum_{v\mid\infty}\log\det G_v$, where
$d=[K:\mathbb Q]/2$ and the sum has one place per conjugate pair.
The basis $e_i-e_{T-1}$ of $A$ has Gram matrix $I+J$ of size $T-2$,
where every entry of $J$ is one.  Its eigenvalues are $1$ and $T-1$,
with multiplicities $T-3$ and one, so $\det(I+J)=T-1$ and
$\mu(A)=-\log(T-1)/(T-2)$.
The identity
$\det(G\otimes H)=(\det G)^{\operatorname{rk}H}
 (\det H)^{\operatorname{rk}G}$ proves tensor additivity of this
normalized slope.  For a projective lattice $W$, choose a full free
sublattice; its finite index is raised to the number of copies in each
tensor product, giving the same identity after division by rank.
Since $\mu(U)=0$, we obtain
\[
 \boxed{\displaystyle
 \mu(M_{\rm ret})=\mu(M_{\rm full})
       -\frac{f\log(T-1)}{T-2}.}
\]
For $f=0$ both lattices are $W$, the rank ratio is one, and the loss is
zero.  No square-root or additional factor of two occurs in this metric
convention.  Applying this calculation to the actual ramified source
requires an integral isomorphism with this induced model and an isometry
at every infinite place; a rank formula alone does not supply that
identification.  The source's precise transport theorem is
\texttt{slope\_ramified\_induced\_prime\_field}; its actual column
specialization is \texttt{MramH\_mu}.


\section{Saturation removes an avoidable rank loss}\label{sec:saturation}

\status{This strengthens the local accounting for the actual retained reference lattice. The joint exceptional-direction estimates remain inputs. The reference-lattice comparisons and descent used below are explicit in the pinned source; the elementary deductions are proved here.}

Fix a word $u$ and a served finite place $w\mid p$, so $p\ne T_u$.
Write $K$ for the coefficient field, $A=\mathcal O_{K,w}$, and
$\mathcal O\subset F$ for the source's comparison DVR and its fraction field.
Let $L_u$ be the full integral column lattice, of rank $d_u$, and let
$U_u$ be its rational ramified subspace. Its retained reference lattice is
$R_u=L_u\cap U_u$, of rank $r_u>0$.
The source's coordinate inclusion identifies this with \texttt{DV.MramL}.
At $w$, the comparison map $\phi$ identifies
\begin{equation}\label{eq:sr-reference}
 M=\mathcal O\phi(L_u),\qquad
 U_F=F\phi(U_u),\qquad
 M_{\rm ram}=M\cap U_F=\mathcal O\phi(R_u).
\end{equation}
Here $M$ is \texttt{G.M}, not \texttt{G.Mna}. Since $U_F$ is a vector
subspace, $ax\in U_F$ and $a\ne0$ imply $x\in U_F$. Thus
$M_{\rm ram}$ is saturated in $M$, and
\[
 M_{\rm ram}\cap p^tM=p^tM_{\rm ram}\qquad(t\ge0).
\]
The same argument proves saturation of $R_u$ in $L_u$ before comparison.
No tensor-model or metric identification is needed for this fact.
The selected finite-index lattice need not itself be saturated. \textsf{[SR1]}

\paragraph{The actual comparison and descent.}
The source specifies an embedding $\sigma:K\hookrightarrow F$ with
$\sigma(x)\in\mathcal O$ if and only if $x\in A$, and requires images
of a $K$-basis under $\phi$ to remain $F$-linearly independent.
Consequently $A\to\mathcal O$ is an injective local map of DVRs.
It is faithfully flat, and scalar extension commutes with intersections
of lattices and rational subspaces. This proves \eqref{eq:sr-reference}.
The actual local selected lattice is
\[
 P=J\cap p^\tau M\cap U_F,
 \quad
 J=\sum_{B\subseteq\mathcal V}
 p^{\sum_{v\notin B}S_v}\Bigl(\prod_{v\in B}y_v^m\Bigr)M.
\]
The recorded rationality of the $y_v$ makes $J$ an extension of an
$A$-lattice; allowing a unit factor for each $y_v$ gives the same conclusion.
Flatness then shows that $P$ is an extension of its rational preimage.
The source's served-place identity identifies that preimage with
$(N_u)_w$, where $N_u=\texttt{NuLat}\,u$.
This is the content of \texttt{M\_inf\_Vram}, \texttt{NP\_eq\_span}, and
\texttt{length\_quotient\_of\_nonempty}, with rationality supplied through
\texttt{NPDesc\_of}. It is not an assumption that two commensurable
reference lattices have the same index.

Put $e_w=v_w(p)$ and $e_{\mathcal O}=v_{\mathcal O}(p)$.
If $a=v_{\mathcal O}(\pi_w)$, then extension multiplies each Smith
exponent by $a$, and $e_{\mathcal O}=ae_w$. Hence
\begin{equation}\label{eq:sr-length-transport}
 \frac{\operatorname{length}_{A}((R_u)_w/(N_u)_w)}{e_w}
 =\frac{\operatorname{length}_{\mathcal O}(M_{\rm ram}/P)}{e_{\mathcal O}}.
\end{equation}
This also preserves the ranks $d_u,r_u$. \textsf{[SR2]}

\paragraph{Why the smaller full lattice causes only a depth shift.}
The comparison lattice $M_{\rm na}=\texttt{G.Mna}$ satisfies
$p^eM_{\rm na}\subseteq M\subseteq M_{\rm na}$.
Both are preserved by the $y_v$. If $J_{\rm na}$ is defined as $J$ using
$M_{\rm na}$, then $p^eJ_{\rm na}\subseteq J$.
Thus, from
$p^EM_{\rm na}\subseteq J_{\rm na}+\sum_{i=1}^g\mathcal O z_i$,
multiplication by $p^e$ gives
\[
 p^{E+e}M\subseteq J+\sum_{i=1}^g\mathcal O(p^ez_i),
 \qquad p^ez_i\in M.
\]
The exceptional count is unchanged, and both full lattices have the same
rank. This is precisely the source's \texttt{inverse\_shift} step.
For a nonempty served set the resulting depth is the integer
\[
 \tau=e+\sum_{v\in\mathcal V}m(t_v+D_0+1).
\]
The real-valued \texttt{tauP} is this integer, and \texttt{Sv} is already
a nonnegative integer. Their displayed ceilings introduce no additional
rounding loss. For an empty served set, $\tau=\mathrm{io}=0$ and
$P=M_{\rm ram}$. At unserved regular places the recorded
\texttt{NuUnserved} identity gives $(N_u)_w=(R_u)_w$, hence zero index.
\textsf{[SR3]}

\begin{proposition}\label{prop:sr-actual-index}
Assume the source's joint exceptional-direction bounds: at depth
$\tau+\beta S$, $J\subseteq M$ has at most
$d_u m^{\beta+1}e_{\beta+1}(\kappa)$ exceptional directions for every
$\beta\ge0$, with $S\ge S_v$, $\kappa_v\ge0$ and
$\Sigma=m\sum_v\kappa_v\le1/2$.
Set $\Theta=\prod_v(1+m\kappa_v)-1$ and $\mathrm{io}=2S\Sigma$.
Then the actual selected lattice satisfies
\begin{equation}\label{eq:sr-index}
 \operatorname{length}_{A}((R_u)_w/(N_u)_w)
 \le e_w\{r_u\tau+d_uS\Theta\}
 \le e_w\{r_u\tau+d_u\mathrm{io}\}.
\end{equation}
\end{proposition}

\begin{proof}
Put $Q=J\cap p^\tau M$. Saturation gives an exact sequence
\[
 0\longrightarrow p^\tau M_{\rm ram}/P
 \longrightarrow M_{\rm ram}/P
 \longrightarrow M_{\rm ram}/p^\tau M_{\rm ram}\longrightarrow0.
\]
The last quotient has length $e_{\mathcal O}r_u\tau$.
The first embeds in $p^\tau M/Q$, whose length is at most
$e_{\mathcal O}d_uS\Theta$ by the elementary-divisor tail sum in LC2.
Apply \eqref{eq:sr-length-transport}, then $\Theta\le2\Sigma$.
\end{proof}

Writing $\rho_u=d_u/r_u$, the slope conclusion of LC5 therefore improves to
\begin{equation}\label{eq:sr-slope}
 \mu(I_uN_u)\ge\mu(R_u)+\mathrm{Ga}+\mathrm{Fo}-\mathrm{Ee}_u
 -\rho_u\mathrm{IO}-\mathcal R_u-\mathcal B_u.
\end{equation}
The entire $(\rho_u-1)\mathrm{Ta}$ term is removed. The type prime
$T_u$ still has its separately estimated cost $\mathcal B_u$.
The actual source sizes give $T_u\ge X'$, so its served set is empty and
$\tau=\mathrm{io}=0$ there. Thus the exceptional baseline in LC5 is also
zero, and $\mathcal B_u$ retains its previous meaning. If any additional
exceptional places were introduced with nonzero baseline, their excess
would have to be estimated relative to the new baseline. \textsf{[SR4]}

\paragraph{The source's rank sizes absorb the remaining ratio.}
One can retain more of the same tail estimate:
\[
 \Theta\le\sum_{j\ge1}\frac{\Sigma^j}{j!}
 \le\frac{\Sigma}{1-\Sigma/2}\le\frac43\Sigma,
\]
because $j!\ge2^{j-1}$. If $r_u/d_u\ge2/3$, then
$d_uS\Theta\le r_u(2S\Sigma)=r_u\mathrm{io}$.
The source's actual rank formula and size hypotheses give this with room
to spare:
\[
 \frac{r_u}{d_u}=\left(1-\frac1{T_u-1}\right)^f,
 \quad16f\le b,\quad4\le b,\quad b^3\le T_u
 \ \Longrightarrow\ 
 1-\frac{r_u}{d_u}\le\frac{2f}{T_u}
 \le\frac1{8b^2}\le\frac1{128}.
\]
These numerical hypotheses already occur in \texttt{PlacesP.sizes\_at}
and its family assembly. Thus, under the same retained arithmetic and
comparison inputs, \eqref{eq:sr-slope} strengthens to
\begin{equation}\label{eq:sr-final}
 \boxed{\mu(I_uN_u)\ge\mu(R_u)+\mathrm{Ga}+\mathrm{Fo}-\mathrm{Ee}_u
 -\mathrm{IO}-\mathcal R_u-\mathcal B_u.}
\end{equation}
If the actual integral and isometric transport required in LC7 is also
supplied, its separate metric identity is
$\mu(R_u)=\mu(L_u)-f\log(T_u-1)/(T_u-2)$.
The mean correction, the finite-place rounding, and the type-prime cost
remain explicit. In the rational-prime notation $\mathrm{IO}=2\mathrm{IOp}$,
the regular-place allowance is now $2\mathrm{IOp}$ instead of the source's
$4\mathrm{IOp}$ after its coarser absorption. \textsf{[SR5]}

\section{What produces the Gram estimates}
\label{sec:actual-gram}

The following reconstruction goes below the abstract frame argument.
Its conclusions remain conditional on the explicitly identified
spectral and arithmetic inputs below.  In particular, the estimates
are not merely bounds for functions at infinity: finite-place
density and mass identities are essential to their normalization.
Source identifiers AG1--AG10 refer to the accompanying provenance ledger.\cite{actual-gram-kernel,actual-gram-words}

\paragraph{The kernel and its four parts (AG1--AG3).}
Write \(P_u:\mathcal V_u\to\mathcal H\) for the full period map,
\(n_u=\dim\mathcal V_u\), and
\[
 c_u=\frac{\operatorname{vol}N_u\operatorname{vol}H_u}
                 {\dim\pi_u}>0,\qquad
 A_u=c_u^{-1/2}P_u,\qquad K_{uv}=A_u^*A_v .
\]
These are the full source spaces, before the column restriction.
The absolutely convergent geometric expansion separates as
\[
 K_{uv}=\mathbf1_{u=v}I+Z_{uv}+N_{uv}+G_{uv}.                 \tag{AG.1}
\]
Here \(Z\) consists of the nonidentity terms with zero Gram
determinant; \(N\) and \(G\) are respectively the nongeneric and
generic terms with nonzero determinant.  The identity-cell
normalization gives the displayed \(I\).  For \(u\ne v\), the
zero-determinant amplitude intertwines distinct compact types
(after an inner conjugation), so Schur's lemma gives \(Z_{uv}=0\).
The geometry establishing this intertwining, and the absolute
convergence needed for (AG.1), are retained inputs.

The generic squared norm is controlled by a scalar pair sum and a
small remainder.  Set
\[
 X_{uv}=\frac{n_un_v}{D_*},\quad
 F=\prod_{\mathfrak p\ {\rm typed}}(1+3/N\mathfrak p),\quad
 E=e^{-c_Lkf\log k}.
\]
The source's pair-trace identity gives
\[
 \|G_{uv}\|_{\rm HS}^2\le FE X_{uv}+S_{uv},                 \tag{AG.2}
\]
where \(S_{uv}=F\gamma_G\operatorname{Re}\sum_{\rm scalar}p\).
The omitted pairs comprise defect \(s_0\ge2\), and the nonscalar
generic pairs of defect \(s_0=1\).  Here \(s_0=k-\dim_E W_p\)
for the matching subspace attached to the pair.  Their arithmetic counting bound is the
input \(E X_{uv}\); it is not a consequence of Schur orthogonality.

\paragraph{The actual archimedean functions (AG3--AG5).}
At a complex-place representative \(\nu\), the Bessel amplitude is
the double equivariant average
\[
 J_\nu F(g)=
 \int \psi_\nu(n)\overline{\psi'_\nu(n')}\,
 \pi_\nu(w)F(w^{-1}n^{-1}gn'w')\pi'_\nu(w'^{-1})
 \,dw\,dn\,dw'\,dn'.
                                                               \tag{AG.3}
\]
The displayed integral uses the source's fixed Haar measures
and the required convergence hypotheses.  In its big-cell coordinates
\((z_1,\ldots,z_r,z)\), where \(r=r_S\), the energy is
\[
 \mathcal J=\kappa^{-(r+1)}
 \prod_\nu\int_{\mathbb C^{r+1}}
 \left[\prod_{a=1}^{r}|z_a|^{\,2(2(r-a)+k+1)}\right]
 (|z|^2-T_uT_v)_+^{\,k-1}
 \|J_\nu F_\nu(z_1,\ldots,z_r,z)\|_{\rm HS}^2\,dz .
                                                               \tag{AG.4}
\]
Here Lebesgue measure is used at each complex coordinate, and
\(\kappa>0\) is the common covolume of the scaled arithmetic
lattices in \(\mathbb C^f\); the source chooses the block-dependent
scalings so these covolumes agree.  Boundary values of the
positive part are immaterial to the integral.  The scalar pair identity rewrites
\(\gamma_G\operatorname{Re}\sum_{\rm scalar}p\) as
\((n_u/\dim\pi_u)\) times a positive quadratic form in tensor
products of these amplitudes.  Its coefficient kernel is bounded,
in the positive-semidefinite order, by an arithmetic lattice
kernel.  This is the precise input to the next argument.

The words indexing \(u\) are binary.  A differing letter at position
\(\ell\) gives, at every \(\nu\), the two compact parameters
\[
 v_{\nu\ell}-\tfrac12-\Delta,\qquad v_{\nu\ell}-\tfrac12 .
\]
Thus
\[
 I_\nu=[\,\tfrac52-v_{\nu\ell},\ \tfrac{\Delta}{2}-v_{\nu\ell}\,]
                                                               \tag{AG.5}
\]
lies strictly between their negatives and stays at least \(2\)
from both endpoints when \(\Delta\ge5\).  Only one differing
letter is needed; this step uses no lower bound for Hamming distance.
The large width of \(I_\nu\), together with the remaining
Harish--Chandra gaps, permits repeated spectral differentiation.
The normalization of a differential word \(w\) divides by
\(\prod_{\nu,Y\in w_\nu}|P_{\lambda_\nu}(Y)|\), with
\(P_\lambda(Y)=\prod_j(Y+\lambda_j)\).

Here is the local analytic input that still requires the
representation-theoretic proof.  At each level its Fourier transform
has zero mode equal to the preceding level's energy and satisfies
\[
 |\widehat{\mathcal N}_w(\xi)|
 \le\prod_\nu\min\{1,(x_{i\nu}/|\xi_\nu|)^K\}
             \sum_{w'}\theta_\xi(w')\mathcal E_{w'},\qquad
 \theta_\xi\ge0,\quad\sum_{w'}\theta_\xi(w')=1.                 \tag{AG.6}
\]
The extensions \(w'\) use at most \(K\) additional letters.
A factor with \(\xi_\nu=0\) is interpreted as \(1\).
There are corresponding off-diagonal estimates with
\(\xi_\nu=t_\nu^{-1}-(t'_\nu)^{-1}\), and estimates for blocks of
size \(s\ge2\).  The source takes \(K=5\).  These estimates,
positivity of the intermediate forms, and their common limiting
energy are substantive inputs, not definitions of a proof.

\paragraph{Why the error has a power \(f\) (AG6--AG7).}
Let \(E\) be the auxiliary CM field of degree \(2f\), with exactly
one embedding from each conjugate pair.  If \(\mathfrak a\ne0\)
is a fractional ideal and
\(\prod_\nu x_\nu\le \sqrt{N\mathfrak a}/\Theta^f\), the elementary
packing argument in the appendix proves
\[
 \sum_{0\ne z\in\mathfrak a}
       \prod_\nu\min\{1,(x_\nu/|\nu z|)^K\}
 \le \bigl(384\,\Theta^{-(K-1/2)}\bigr)^f
 \quad(K\ge3,\ \Theta\ge1).                                  \tag{AG.7}
\]
This is a new sufficient replacement for the weaker exponent
\(K-5/2\) used in the archive.  No improvement to the final RH
exponent is claimed.  The source defines \(\Theta\) from the
maximum of all relevant spectral-size products divided by square
roots of fractional-ideal norms, so the same \(\Theta\) controls
both Fourier and off-diagonal sums.

Poisson summation now separates the zero frequency from the
nonzero frequencies in (AG.6).  The former contributes exactly the
preceding energy; the latter contributes the small quantity in
(AG.7).  Splitting the final block of a positive Gram form and
using its triangle inequality gives the following concrete
recurrence for square roots of level energies:
\[
 x_i\le(1+\varepsilon)\sqrt{1+e_i}\,x_{i-1}
       +\sum_{s=2}^{i}a^{s-1}x_{i-s},\qquad
 a=(C/\Theta)^f,\quad 0\le\varepsilon\le a.                   \tag{AG.8}
\]
The \(e_i\ge0\) are the arithmetic kernel's density losses.
For the finitely many place patterns, the source closes the
finite supports of these probability distributions under the
bounded number of permitted extensions.  Taking a maximum over
that finite set bounds every average without a factor equal to
the number of words.  Finiteness follows because every nontrivial
extension strictly increases an integer total length, while the
length budget is bounded; the resulting tree has finite depth
and finite branching.

For completeness, suppose \(a\le1/2\) and \(na\le1\), where
\(n=r+1\).  Divide \(x_i\) by
\(\prod_{j\le i}\sqrt{1+e_j}\) and call the result \(y_i\).
Since the removed factors are at least one,
\[
 y_i\le(1+\varepsilon)y_{i-1}+
              \sum_{s=2}^{i}a^{s-1}y_{i-s}.
\]
If \(y_0=0\), induction gives \(y_i=0\) for all \(i\).
Otherwise induction gives \(y_i\le(1+3a)^iy_0\): the terms with \(s\ge2\),
after division by \((1+3a)^{i-1}y_0\), sum to at most
\(a/(1-a)\le2a\).  Consequently
\[
 x_n^2\le
 [1+(e^6-1)na]\prod_{j=1}^{n}(1+e_j)\,x_0^2.                \tag{AG.9}
\]
Indeed \((1+3a)^{2n}\le e^{6na}\le1+(e^6-1)na\).
This proves the leading constant \(1\) and the stated small
relative error.  Increasing the fixed constant \(C\) absorbs
the harmless absolute constants.  Taking finite sums first,
then the smooth-cutoff limit and increasing finite sets, is
justified by the stated positivity and absolute summability.

\paragraph{Returning to the period maps (AG8--AG10).}
The mass and formal-degree normalization, a separate arithmetic
input, bounds the base energy after multiplication by
\(n_uc_W/\dim\pi_u\) by \(e^{8q_S}X_{uv}\), where \(c_W>0\)
is the scalar in the arithmetic weight identity and
\(q_S=\sum_{\mathfrak p}(N\mathfrak p)^{-k}\).
The density factors and \(e^{8q_S}\) are absorbed into
\(1+e_{\rm loc}\).  Equations (AG.6)--(AG.9) therefore yield
\[
 S_{uv}\le
 [1+C(r+1)(C/\Theta)^f](1+e_{\rm loc})F X_{uv}
 \quad(u\ne v).                                             \tag{AG.10}
\]
The source's parameter comparison supplies
\((r+1)(C/\Theta)^f\le b^{-7}\), with \(b=e^L\).
It also gives, for \(\delta=b^{-7}\),
\[
 F\le1+\delta,\quad E\le\delta,\quad e_{\rm loc}\le\delta,
 \quad \|N_{uv}\|_{\rm HS}^2\le\delta^3X_{uv}.
\]
Use
\(\|N+G\|^2\le(1+\delta)\|G\|^2+
(1+\delta^{-1})\|N\|^2\) in (AG.1)--(AG.2).
The result is at most
\([1+(13+8C)\delta]X_{uv}\), hence
\[
 \|K_{uv}\|_{\rm HS}^2\le(1+b^{-6})\,n_un_v/D_*
\]
once \(b\ge13+8C\).
On the diagonal the scalar sum needs only a coarse
\(e^{C_E m^5fL}n_u^2/D_*\) bound.  The same decomposition gives
\[
 \|K_{uu}-Z_{uu}-I\|_{\rm HS}^2\le\Gamma n_u^2/D_*,
 \qquad
 \|Z_{uu}\|_{\rm HS}^2\le n_u e^{-kf\log k}.
\]
Since \(n_u\Gamma^2/D_*\le C_0\), it follows that
\[
 \|K_{uu}-I\|_{\rm HS}^2\le
 \bigl(2C_0/\Gamma+2e^{-kf\log k}\bigr)n_u .                 \tag{AG.11}
\]
A separate automorphic dimension estimate supplies
\[
 D_{\rm all}\le(1+\eta_1)D_*,
 \qquad \eta_1=e^{-c_7Rf\log R},\quad c_7>0,
\]
where \(R\) is the source's total rank, not \(r=r_S\).
This counting estimate is a retained input; it does not follow
from the period-kernel calculation.
These inequalities give the Gram hypotheses used below.
The retained obligations are now more specific:
the geometric expansion and trace identity, the local spectral
identities underlying (AG.6), the arithmetic kernel and mass
normalization, and the loop, nongeneric defect-one, and
zero-determinant diagonal estimates.

\section{The common-lattice slope comparison}\label{sec:frame}

The preceding constructions produce integral maps $B_i:N_i\to C$ into one common hermitian lattice $C=R^\vee$. Here $N_i$ denotes the \emph{shifted} source lattice; its rank is the retained source rank. For a hermitian lattice, $\mu$ is the normalized Arakelov degree divided by rank, with logarithms of squared norms. Its Harder--Narasimhan slopes describe the degrees of the successive semistable parts; their rank-weighted mean is $\mu$. We use the standard nonzero-morphism inequality
\[
\mu_{\min}(E)\le\mu_{\max}(F)+h(f)
\]
for an integral nonzero map $f:E\to F$, where $h(f)$ is the normalized mean logarithmic squared operator norm. This fixes both the sign and the factor of two in the following comparison.\cite{frame,division}

\subsection{A complete finite-dimensional frame lemma}

\begin{lemma}\label{lem:frame}
Let $A_i:V_i\to H$ be complex linear maps between finite-dimensional Hilbert spaces. Put $d_i=\dim V_i$, $M=\sum_i d_i>0$, $D=\dim H>0$, and $Q=\sum_iA_iA_i^*$. Suppose
\begin{gather*}
\sum_i\|A_i^*A_i-I\|_\HS^2\le eM,\qquad 0\le e\le1/4,\\
\tr(Q^2)\le(1+\eta)\frac{(\tr Q)^2}{D},\qquad\eta\ge0.
\end{gather*}
If subspaces $V_i'\subseteq V_i$ have total codimension at most $\alpha M$, then
\[
\codim_H\sum_i A_iV_i'\le(\eta+4\alpha+4\sqrt e)D.
\]
\end{lemma}

\begin{proof}
Let $t=\tr Q$ and let $Z=(\sum_iA_iV_i')^\perp$, of dimension $w$. Cauchy--Schwarz gives $|t-M|\le M\sqrt e$, so $t\ge M/2$. Let $P_Z$ be the orthogonal projection onto $Z$. Since $A_i^*Z\perp V_i'$, the omitted directions give
\[
\tr(P_ZQ)\le\alpha M+M\sqrt{\alpha e}.
\]
Indeed, project onto $(V_i')^\perp$, sum its ranks, and apply Hilbert--Schmidt Cauchy--Schwarz to the errors $A_i^*A_i-I$.

Put $X=Q-(t/D)I$. The second hypothesis gives $\|X\|_\HS^2\le\eta t^2/D$, hence
\[
\tr(P_ZQ)\ge tw/D-\sqrt w\,\|X\|_\HS.
\]
For $r=w/D$, combine the inequalities and $M/t\le2$:
\[
r\le\sqrt{\eta r}+2\alpha+2\sqrt{\alpha e}.
\]
If $\alpha\ge1/4$, the asserted bound is automatic. Otherwise use $\sqrt{\eta r}\le(\eta+r)/2$ and $\sqrt{\alpha e}\le\sqrt e$. Rearranging proves the lemma.
\end{proof}

This is a compression of the abstract frame theorem in the source. \textsf{[FR1]}\cite{frame} Its content is concrete: near-isometry of individual columns and near-isotropy of their aggregate Gram operator force their retained images to span almost all of the target.

% Replacement for the old "From spanning to a mean-slope bound" subsection.
% Assumes the preceding full-family frame lemma has label lem:frame.
% Source identifiers GF1--GF8 are resolved in gram_frame_provenance.json.
\subsection{From the actual Gram bounds to a mean slope}\label{sec:gram-to-slope}

Here is the finite-dimensional bridge from the released pairwise estimates to
the common-lattice comparison.  Section~\ref{sec:actual-gram} reconstructs their kernel reduction and several quantitative deductions. The spectral, arithmetic-kernel, and normalization inputs identified there remain explicit.  Write $\mathcal V_i$ for the full source spaces, $\mathcal H$ for the
full target, $n_i=\dim\mathcal V_i$, $D=\dim\mathcal H$, and $M=\sum_i n_i$.
The source normalizes its period operators by
\[
A_i=c_i^{-1/2}P_i,\qquad
c_i=\frac{\operatorname{vol}(N_i^{\rm an})
                  \operatorname{vol}(H_i^{\rm an})}{\dim\Pi_i}>0.
\]
Thus its kernel $K_{ij}$ is exactly $A_i^*A_j$.
The positivity of $c_i$ follows from the nonnegative normalization and the
diagonal estimate with $e_0<1$ and $n_i>0$.  The operators here act on the full
analytic spaces, before either source or target directions are discarded.
\textsf{[GF1, GF6]}

The pairwise input is
\begin{gather*}
\|A_i^*A_j\|_{\HS}^2\le(1+\eta_0)n_in_j/D_*\quad(i\ne j),\qquad
\|A_i^*A_i-I\|_{\HS}^2\le e_0n_i,\\
D\le(1+\eta_1)D_*.
\end{gather*}
We can enlarge $\eta_0,\eta_1$ to be nonnegative.  Set $a=\sqrt{e_0}$,
$Q=\sum_iA_iA_i^*$, and $t=\tr Q$.  The trace identities and
Hilbert--Schmidt Cauchy--Schwarz give the source's \texttt{CorBounds}:
\begin{equation}\label{eq:full-cor-bounds}
\begin{split}
t&\ge(1-a)M,\\
\tr(Q^2)&\le(1+\eta_0)M^2/D_*+(1+a)^2M,\\
\sum_i\|A_i^*A_i-I\|_{\HS}^2&\le e_0M.
\end{split}
\end{equation}
Indeed, $\tr(Q^2)=\sum_{i,j}\|A_i^*A_j\|_{\HS}^2$, and each diagonal
term is at most $(1+a)^2n_i$.  The code-size and dimension-balance conditions
$\#I\ge C_0\Gamma^2/\tau$ and
$n_i\ge D_*/(C_0\Gamma^2)$ give $M\ge D_*/\tau$.
\textsf{[GF2--GF4]}

\begin{proposition}[Direct Gram-to-slope bridge]\label{prop:direct-gram-slope}
Assume \eqref{eq:full-cor-bounds}, $D,D_*,M>0$, $e_0\ge0$,
$D\le(1+\eta_1)D_*$, $M\ge D_*/\tau$,
$0<\tau\le1/4$, and $0\le\eta_0,\eta_1,a\le\tau/2$.
For every $i$, let $J_i:V_i\hookrightarrow\mathcal V_i$ be an isometry.
Write $r_{\rm tot}=\sum_i\dim V_i>0$, and assume $M-r_{\rm tot}\le dM$, where $0\le d<1$.

Let $N_i,C$ be full hermitian lattices over one totally complex number field,
and let $B_i:N_i\to C$ be integral maps whose nonzero maps have normalized
logarithmic squared operator heights at most $h_i$.
At one complex embedding, suppose invertible identifications
$X_i:(N_i)_\sigma\to V_i$ and $Y:C_\sigma\to\mathcal H_0\subseteq\mathcal H$
satisfy
\[
YB_i^\sigma=P_0A_iJ_iX_i,\qquad
\dim\mathcal H_0\ge(1-\delta)D,\quad 0\le\delta<1,
\]
where $P_0$ is orthogonal projection.  Suppose the mean over the $r_{\rm tot}$
adjusted source slopes $\lambda_{ij}=\mu_j(N_i)-h_i$, with each
Harder--Narasimhan slope repeated according to its rank, is at least $s$;
every adjusted source slope is at most $s+u$; and every target slope is
at least $s-H$, where $u,H\ge0$.  Then
\begin{equation}\label{eq:direct-gram-slope}
\mu(C)\ge s-
\frac{4uH/(u+H)+(5\tau+4d+4\sqrt{e_0})H}{1-\delta}.
\end{equation}
The fraction $uH/(u+H)$ means zero when $u=H=0$.
\end{proposition}

\begin{proof}
First the full family satisfies the preceding frame lemma.  Put
$N_\tau=(1+\tau/2)[1+\tau/2+\tau(1+\tau/2)^2]$.  From
\eqref{eq:full-cor-bounds},
\[
\frac{D\tr(Q^2)}{t^2}
\le\frac{(1+\eta_1)[1+\eta_0+\tau(1+a)^2]}{(1-a)^2}
\le\frac{N_\tau}{(1-\tau/2)^2}\le1+5\tau.
\]
The last inequality follows by expansion: the numerator of the difference is
$\tau(16-52\tau+4\tau^2-\tau^3)/8\ge0$ for $0\le\tau\le1/4$.
Also $e_0\le1/64<1/4$.  Hence Lemma~\ref{lem:frame} applies to the
\emph{full} spaces with parameters $\eta=5\tau$, $e=e_0$.

Let $q(v)$ be the fraction of the $r_{\rm tot}$ adjusted source slopes below $s-v$,
and $p(v)$ the fraction of target slopes below $s-v$.  The usual nonzero-map
slope inequality implies that the source steps of adjusted slope at least
$s-v$ map into the corresponding target step.  Apply it to the quotient
by that target step: a nonzero induced map would contradict the strict
slope separation.  The transported source steps have total codimension
$M-r_{\rm tot}+r_{\rm tot}q(v)$ in the full source spaces, whose fraction is at most $d+q(v)$.
Their full images therefore miss at most
$(5\tau+4d+4q(v)+4\sqrt{e_0})D$ target dimensions.
Projection onto $\mathcal H_0$ cannot increase this codimension.
Consequently
\begin{equation}\label{eq:direct-tail}
p(v)\le\frac{5\tau+4d+4\sqrt{e_0}+4q(v)}{1-\delta}.
\end{equation}
Here $d$ counts discarded \emph{source} directions, while $\delta$ counts
discarded \emph{target} directions.

For every adjusted source slope $\lambda\le s+u$,
\[
\min\{H,(s-\lambda)_+\}
\le\frac{H(s+u-\lambda)}{H+u}.
\]
Averaging gives $\int_0^Hq(v)\,dv\le uH/(u+H)$.
Integrating \eqref{eq:direct-tail} bounds the target's mean positive deficit
from $s$, proving \eqref{eq:direct-gram-slope}.  If $H=0$ the conclusion
follows directly from the target's minimum-slope bound.
\end{proof}

\paragraph{Uniform parameters and a simpler error term.}
Fix $\varepsilon,C_H>0$.  The source supplies $\tau=x^{-C_H}$ and
$\eta_0,\eta_1,\sqrt{e_0}\le E(x)=\exp(-x^{\varepsilon/8})$ uniformly
over every auxiliary datum at that $x$.  Its ramified source ranks satisfy
$r_{\rm tot}\ge(1-d)M$ with $d\le2f/b\le E(x)^2$, where
$b=\exp(x^\varepsilon)$ and $f\le x^5$.
The discarded target fraction is independently bounded by
$\delta\le\exp(-x^{e_z/7})$, with $e_z>0$ fixed.
For sufficiently large $x$, uniformly in all those data,
$E\le\min(1,\tau/8)$ and $\delta\le1/7$.
Thus $5\tau+4d+4\sqrt{e_0}\le6\tau$.
If $H\le x^{C_H}b=b/\tau$, the proposition yields
\begin{equation}\label{eq:direct-frame-final}
s\le\mu(C)+\frac{4u+6b}{1-\delta}
\le\mu(C)+5u+7b.
\end{equation}
The common-coordinate realization comes from \texttt{BesselId}; nonzero
normalizing scalars are absorbed into the invertible $X_i$.
\textsf{[GF5--GF7]}

This direct proof does not charge a $\log(\#I)$ loss, and replaces the
later multiplier $\log(1+x^{C_H})$ on the source-width error by a constant.
The bound on $H$ is still needed to control the small frame-defect term.
The direct treatment of source loss and the clipped-deficit integration are
derived alternatives here; the archive takes a longer compression route.
They simplify the comparison without proving the remaining arithmetic
inputs or changing the final exponent by themselves. \textsf{[GF8]}

The clipped-deficit improvement was also identified in the agent-authored reading guide; no novelty claim is made. The complete argument above is included so that using it does not depend on that guide.\cite{guide}

\subsection{The cancellation that matters}

Let $s$ be the source-rank-weighted mean of the shifted source slopes minus the respective map heights. The source's arithmetic comparison has the form
\[
G_{\rm net}+\mu(C)-E_{\rm arith}\le s.
\]
The direct frame comparison gives $s\le\mu(C)+5u+7b$ under its displayed uniform hypotheses. Eliminating the same target slope therefore gives
\begin{equation}\label{eq:cancel}
G_{\rm net}\le E_{\rm arith}+5u+7b.
\end{equation}
The detailed companion calculation translates this into the source's \texttt{Ga\_bound5} notation and replaces its later logarithmic frame allowance by a constant multiple of the source-width error. The local height, congruence, and prime-sum errors are still required. The source identities and averaging then lead to Section~\ref{sec:scalar}. \textsf{[GL1, GF8]}\cite{congruence,chtop}

The use of one common target is essential: unrelated target lattices would leave an uncontrolled difference of slopes. This is why the descent, local gluing, finite-place estimates, and archimedean realization all have to refer to the same maps.

% Human reconstruction; source ledger: ordinary_host_provenance.json.
% This is an inclusion fragment, with no preamble or document environment.
\section{The ordinary host and a fixed pair on a whole window}\label{sec:ordinary-host}

Here $\mathbf1$ is the primitive character of conductor one.  The following
steps isolate the ordinary specialization without asserting the remaining
cohomological estimate.  All character bounds retain their centering at $s=1$.

\paragraph{OH1. The field.}
Let $q\equiv3\pmod8$ and $d\equiv7\pmod{16}$ be distinct primes, and put
\[
 T_e=\mathbb Q(\zeta_{2^{e+2}})^+,
 \qquad E_e=T_e(\sqrt{-q},\sqrt{-d})\quad(e\ge1).
\]
Then
\[
 \operatorname{Gal}(E_e/\mathbb Q)=C_{2^e}\times C_2\times C_2,
 \quad [E_e:\mathbb Q]=4\cdot2^e,
 \quad E_e^+=T_e(\sqrt{qd}).
\]
Indeed $T_e$ is cyclic of degree $2^e$ and ramifies only at $2$.
The three quadratic subfields of
$\mathbb Q(\sqrt{-q},\sqrt{-d})$ ramify at $q$ or $d$, so its intersection
with $T_e$ is $\mathbb Q$.  Complex conjugation is $(1,-1,-1)$ in the
displayed product.  In $\mathbb Q(\zeta_{2^{e+2}qd})$, an automorphism with
cyclotomic exponent $a$ fixes $E_e$ precisely when
\[
 a\equiv\pm1\pmod{2^{e+2}},\qquad
 J_q(a)=J_d(a)=1,\qquad J_r(a)=\Bigl(\frac ar\Bigr).
\]
The last assertion follows either from quadratic Gauss sums or the quadratic
subfield theorem for cyclotomic fields: the discriminants of
$\mathbb Q(\sqrt{-q})$ and $\mathbb Q(\sqrt{-d})$ are $-q$ and $-d$.
Thus this is exactly the field specified by \texttt{IsChiField} for
$\mathbf1$.  Its conductor is $2^{e+2}qd$.
The source's order parameter is $n=2^e$, not the field degree $4n$:
\texttt{ordAny} gives $1\le c\le\operatorname{ord}(\mathbf1)=1$, so its
general formula $n=c2^e$ has $c=1$.

\paragraph{OH2. The two inputs, with their quantifiers.}
Fix $0<\varepsilon<1/52$ and define
\[
 e(x)=\left\lceil(2+8\varepsilon)\log_2x\right\rceil\quad(x\ge1),
 \qquad x^{2+8\varepsilon}\le2^{e(x)}\le2x^{2+8\varepsilon}.
\]
Write $\operatorname{Good}(\varepsilon,Y,X;q,d)$ for the complete source
condition \texttt{PairGood 1}: besides its zero-free boxes, it includes
\[
 \tfrac12Y^{\varepsilon/5}<\log q\le Y^{\varepsilon/5},\quad
 \tfrac12Y^{\varepsilon/5}(\log Y)^2<\log d
       \le Y^{\varepsilon/5}(\log Y)^2,
\]
$q\mid d+1$, and $p\mid d+1$ for every odd prime $p\le Y^{\varepsilon/5}$.
Its three nontrivial twists $\chi=J_q,J_d,J_qJ_d$ satisfy
$|G_\chi(u)|\le u^{1/2+\varepsilon/2}$ and
$|L'/L(1,\chi)|\le4\log u$ for every $u\in[Y,X]$.
The pair-selection input says that for every fixed $A\ge1$ this condition
holds for some pair at every sufficiently large $Y$, with $X=Y^A$.

For a number field $E$, let $S_p$ consist of its places above $p$ and its
infinite places, and let $G_{E,S_p}$ be the Galois group of its maximal
extension unramified outside $S_p$. Put
$D_{p,j}=(\mathbb Q_p/\mathbb Z_p)(-j)$, the usual Tate-twisted module, and define
\[
\ell_p(E,j)=\operatorname{length}_{\mathbb Z_p}\ker\left(
H^1(G_{E,S_p},D_{p,j})\longrightarrow\prod_{w\mid p}H^1(E_w,D_{p,j})\right).
\]
The cohomological size is
\[
 \mathcal V(E,t)=\sup_{S\text{ finite primes}}
       \sum_{p\in S}\log p\sum_{1\le j<t}\ell_p(E,j).
\]
The cohomological input also asserts finiteness of every displayed length.
The specialization of \texttt{VWinAt} for $\mathbf1$ and $401$ used here says:
for this choice of $e(x)$, there are $J\ge1,X_0$ such that, for every
$X\ge X_0$ and every pair satisfying the preceding congruences and size
conditions at $Y=X^{1/J}$, there is $W$ with
\[
 X^{1/J}\le W,\quad W^2\le X,\quad
 \mathcal V(E_{e(x)},\lfloor x^{401}\rfloor)
       \le \exp(\tfrac12x^\varepsilon)\quad(W\le x\le W^2).
\]
This input is the unexpanded arithmetic theorem in this subsection.

\paragraph{OH3. Fixed-pair assembly.}
Assuming the two inputs, there is $J\ge1$ such that for every lower
threshold $B$ there are $W\ge B$, $Y$, and a \emph{single} pair $q,d$ for
which, simultaneously for every $x\in[W,W^2]$,
\[
 \begin{gathered}
 x^{1/J}\le Y\le x,\qquad \operatorname{Good}(\varepsilon,Y,x;q,d),\\
 x^{2+8\varepsilon}\le n(x)=2^{e(x)}\le2x^{2+8\varepsilon},\qquad
 \mathcal V(E_{e(x)},\lfloor x^{400}\rfloor)
       \le\exp(\tfrac12x^\varepsilon).
 \end{gathered}
\]
To prove this, first fix $J,X_0$ from the cohomological input.  Apply pair
selection with $A=J$.  Choose $X\ge X_0$ so large that
$Y=X^{1/J}\ge\max(1,B)$ and exceeds its selection threshold.  Choose
$q,d$ good on $[Y,Y^J]=[Y,X]$; then choose $W$ from the cohomological
input for this very pair.  For $W\le x\le W^2\le X$, restrict the good
pair's upper endpoint to $x$ and decrease the twist cutoff from
$\lfloor x^{401}\rfloor$ to $\lfloor x^{400}\rfloor$.
All summands are nonnegative, so the bound persists.  Only $e(x)$ changes
with $x$.  The general source proof rescales by $c^{1/(2+8\varepsilon)}$
and retains $[W,W^{3/2}]$; $c=1$ makes that loss unnecessary here.

\paragraph{OH4. The four characters and the sign.}
The characters trivial on $C_{2^e}$ are exactly
$1,J_q,J_d,J_qJ_d$.  Weighting them by $+,+,-,-$ gives
\[
 S_{q,d}(x)=G_{\mathbf1}(x)+G_{J_q}(x)-G_{J_d}(x)-G_{J_qJ_d}(x),
 \qquad |S_{q,d}(x)-G_{\mathbf1}(x)|\le3x^{1/2+\varepsilon/2}.
\]
At a prime power $p^k$, the corresponding von Mangoldt coefficient,
in units of $\log p$, is
$(1+J_q(p)^k)(1-J_d(p)^k)$.
For $p\nmid qd$, it is $4$ exactly when $k$ is odd,
$J_q(p)=1$, and $J_d(p)=-1$; otherwise it is zero.
Since $d\equiv-1\pmod q$, reciprocity gives $J_q(d)=-1$ and
$J_d(q)=1$.  Thus all terms at $q$ vanish; at $d$, the even powers have
coefficient $2$ and the odd powers vanish.  Also $J_d(2)=1$, and the
small-prime congruences give $J_d(p)=1$ for all odd
$p\le Y^{\varepsilon/5}$.  This identifies the prime class exactly.
It does not replace the centered ramp bounds, zero-free boxes, or the
cohomological input by a bound on a finite uncentered prime sum.

\paragraph{OH5. Why a square window can avoid finitely many exceptional twists.}
One elementary part of the cohomological argument can also be isolated. Suppose a set $T$ contains at most $R$ exceptional nonnegative integers. Fix $0<\eta\le C$ and $X>1$, and put
\[
A=2C/\eta,\qquad W_j=X^{1/(2A^j)}\quad(0\le j\le R).
\]
The $R+1$ bands $(W_j^\eta,W_j^{2C}]$ are pairwise disjoint, since the upper endpoint of one is the excluded lower endpoint of the preceding band. At least one band is empty of exceptions. For its $W=W_j$, every $x\in[W,W^2]$ and $r\in T$ with $r\le x^C$ satisfies $r\le x^\eta$. Moreover
\[
W\ge X^{1/J_*},\quad W^2\le X,\qquad J_*=2(2C/\eta)^R.
\]
This proves the finite-exception window mechanism, with a slightly smaller constant than the source. Bounding the number of exceptions and their contributions remains arithmetic work; small exceptional twists have not disappeared. The companion notes also prove the ultrametric deduction from a finite set of spectral centers. \textsf{[OH5]}



% Inclusion fragment. Exact declaration ledger: grh_provenance.json.
% New elementary deductions are distinguished from retained arithmetic inputs.
\section{What changes for Dirichlet GRH}\label{sec:grh}

\paragraph{GH1. Scope and fixed target.}
The exported general statement concerns every Dirichlet character, of every
positive modulus: $L(s,\chi)\ne0$ for $\Re s>1/2$, with the principal
character's pole at $s=1$ excluded.  It also asserts that every zero in
$0<\Re s<1$ lies on the critical line.  Unrestricted Dedekind or
automorphic GRH is not asserted by these declarations.\cite{grh-statements,grh-host}

Fix a primitive character $\psi$ of conductor $M$ and order $o$.
Let $\psi_j$ be the primitive inducer of $\psi^j$, $0\le j<o$.
These $o$ characters are distinct: equal primitive inducers induce the
same character back to modulus $M$.  They are closed under inversion,
$\psi_0=\mathbf1$, and $\psi$ occurs.  Thus the source's
\texttt{Level.cyc} is exactly the fixed positive family
\[
 H_\psi(x)=G(x)+\frac1{3o+1}\sum_{j=1}^{o-1}\Re G_{\psi_j}(x).
 \tag{GH2}
\]
Here the centered sums $G_\chi$ are those already used in AN3.  No target
character changes as the auxiliary fields grow.  The arithmetic goal is
the same window bound as before, with constants allowed to depend on this
fixed $\psi$; AN3 already proves that it excludes every constituent zero
to the right of $1/2$.

\paragraph{GH3. An explicit host and the constant $c$.}
Write $M=2^aM_0$, $M_0$ odd, and use the Chinese remainder theorem to write
$\psi=\psi_2\psi_{\rm odd}$.  Put
\[
 K_\psi=\mathbb Q(\zeta_M)^{\ker\psi},\quad
 T_e=\mathbb Q(\zeta_{2^{e+2}})^+,\quad
 c=\operatorname{lcm}\bigl(\operatorname{ord}\psi_{\rm odd},
                  \operatorname{ord}(\psi_2(-1))\bigr).
\]
For $e\ge1$ and $e+2\ge a$, and distinct primes
$q\equiv3\pmod8$, $d\equiv7\pmod{16}$ with $\gcd(M,qd)=1$,
\[
 E_e=T_eK_\psi(\sqrt{-q},\sqrt{-d}),\qquad
 [K_\psi\cap T_e:\mathbb Q]=o/c,\quad
 [E_e:\mathbb Q]=4c2^e,\quad n=c2^e.                 \tag{GH3}
\]
Indeed the character group of $K_\psi$ is $\langle\psi\rangle$.
Its member $\psi^j$ belongs to the character group of $T_e$ precisely
when its odd-conductor component is trivial and its $2$-component is
even: $\operatorname{ord}(\psi_{\rm odd})\mid j$ and
$\operatorname{ord}(\psi_2(-1))\mid j$.  Their intersection therefore
has $o/c$ elements.  The two new quadratic characters are independent
of this base group by CRT at $q$ and $d$, giving the factor four.
In the ambient cyclotomic field, fixing $E_e$ means exactly
\[
 u\equiv\pm1\pmod{2^{e+2}},\quad \psi(u)=1,
 \quad J_q(u)=J_d(u)=1.
\]
This is \texttt{IsChiField}; it identifies the actual source field up to
$\mathbb Q$-isomorphism.  The source's auxiliary modulus is also explicit:
$Q_S=\operatorname{lcm}(2^{e+2},M)$, since its definition gives both
$2^{e+2}\mid Q_S$ and $Q_S\mid\operatorname{lcm}(2^{e+2},M)$,
while primitivity gives $M\mid Q_S$.
The proof explains the eventual constant in \texttt{ordAny} without
an unspecified stabilization stage.  In particular $1\le c\le o$.
For $M$ odd, $c=o$; for a purely $2$-primary even character, $c=1$.
Total evenness alone does not give $c=1$: $\psi=\chi_{-4}\chi_{-3}$
is even but has $c=2$.

\paragraph{GH4. The finite character calculation.}
The quotient characters used by the source are
$\psi^jJ_q^iJ_d^k$ ($0\le j<o$, $i,k\in\{0,1\}$).
To see completeness, a character trivial on the first source factor
descends to $(\mathbb Z/M\mathbb Z)^\times/\ker\psi$; every character
of this cyclic group is a power of $\psi$.  The remaining two sign
factors supply $i,k$.  On integers coprime to $2Mqd$,
\[
 \sum_{j=0}^{o-1}\sum_{i,k=0}^1
   (-1)^k\psi(n)^jJ_q(n)^iJ_d(n)^k
 =o\,\mathbf1_{\psi(n)=1}(1+J_q(n))(1-J_d(n)).       \tag{GH4}
\]
This follows by multiplying three finite geometric sums.
It is an exact support identity, not a square-root estimate.
At ramified primes the primitive representatives can have different
values from the characters extended by zero modulo $M$; the displayed
coprimality condition must be retained.
The omitted source characters have $(i,k)\ne(0,0)$, exactly the twists
controlled by \texttt{PairGood}.  Its bounds, and the cohomological
window theorem, remain arithmetic inputs here.

\paragraph{GH5. A fixed pair after the degree rescaling.}
Assume the full \texttt{PairEx} and \texttt{VWinAt} inputs for $\psi$.
Fix $0<\varepsilon<1/52$, put $\alpha=2+8\varepsilon$ and
$\kappa=c^{1/\alpha}$.  There are $J\ge1$ and arbitrarily large $Z$
with a \emph{single} $Y,q,d$ such that every $x\in[Z,Z^{3/2}]$ has
\[
 \begin{gathered}
 x^{1/J}\le Y\le x,\quad \operatorname{PairGood}(M,\varepsilon,Y,x;q,d),\\
 x^\alpha\le c2^{e(x/\kappa)}\le2x^\alpha,\qquad
 \mathcal V(E_{e(x/\kappa)},\lfloor x^{400}\rfloor)
       \le\exp(\tfrac12x^\varepsilon),
 \end{gathered}                                                    \tag{GH5}
\]
where $e(t)=\lceil\alpha\log_2t\rceil$ for $t\ge1$.
For a proof, obtain $J$ from the cohomological input, then select
$q,d$ good on $[Y,X]$ with $Y=X^{1/J}$.  That input supplies
$Y\le W$, $W^2\le X$ and its bound through degree $\lfloor t^{401}\rfloor$
for all $W\le t\le W^2$.  Choose $X$ large enough that
$W\ge\max(\kappa^3,\kappa^{400},1)$, and set $Z=\kappa W$.
Then $x\le Z^{3/2}\le W^2\le X$, $t=x/\kappa\in[W,W^2]$, and
$x^{400}=\kappa^{400}t^{400}\le t^{401}$.
Also $c\,t^\alpha=x^\alpha$ and $t^\varepsilon\le x^\varepsilon$.
These are all the required comparisons.  Any fixed lower bound on $e$
can be met by increasing $X$.  The quantifiers fix $q,d$ before $x$;
only the tower exponent varies.

\paragraph{GH6. Primitive to arbitrary characters, including the boundary.}
Let $\chi$ modulo $N$ be induced by a primitive $\chi^*$ of conductor
$f\mid N$.  Away from a principal pole,
\[
 L(s,\chi)=L(s,\chi^*)
       \prod_{\substack{p\mid N\\p\nmid f}}
             (1-\chi^*(p)p^{-s}).                         \tag{GH6}
\]
For $\Re s>0$, every displayed factor is nonzero, because
$|\chi^*(p)p^{-s}|=p^{-\Re s}<1$.
Thus AN3 applied to GH2 for every primitive target yields nonvanishing
for all characters to the right of $1/2$.  At $s=1$ a nonprincipal
$\chi$ has nonprincipal $\chi^*$, so the same identity applies and
the classical theorem \texttt{LFunction\_ne\_zero\_of\_one\_le\_re}
for nonprincipal characters covers this point; a principal character
has a pole and is excluded.  Within $0<\Re s<1$ the primitive functional
equation reflects a zero to $1-s$ for $(\chi^*)^{-1}$: the conductor,
root-number, and gamma factors there are nonzero.  A zero left of
$1/2$ would contradict right-half-plane nonvanishing.  Added Euler
factors can have zeros on $\Re s=0$, which is why the all-character
zero statement is about the open strip.

\paragraph{GH7. Shared work and the remaining arithmetic.}
The local polynomial, period, Gram, integral-index, slope, and positive-family
analytic arguments are shared.  Ordinary RH sets $M=o=c=1$ and
$H_\psi=G$.  Dirichlet GRH additionally requires the full fixed-$M$
pair-selection theorem, the host $K_\psi$, and cohomological windows
uniform across its growing tower.  These latter inputs do not follow
from the elementary character and degree calculations above; their
precise extra parameters are recorded in the companion report.


\section{The arithmetic still to be written out}\label{sec:remaining}

Section~\ref{sec:dependency-audit-v35} gives every one of the 34
remaining selected contracts. The global assembly implication is
complete; the shared producer hypotheses remain exact.

\begin{description}[style=nextline]
\item[The common original construction.]
Construct the recursive core, one compatible Bundle and its duties,
the original geometric readings, actual base rows and simultaneous
rational tests, and the literal local NBuilt/NdSpec projection square.
The original fields, embeddings and coefficient maps must agree.
The construction-class size and rank implications remain complete.
\item[The original theta extension.]
Produce one ioP valid for every parameter on the same original datum.
Full LowSm/LowIrr and multiplicity one on the actual ordered base model
are now proved. The extremal coefficient identity and Hilbert map are
also proved under boundedness and the same finite-volume regular target.
Supply those target inputs, then the pointwise smooth representatives,
all derivative identities, Sobolev evaluation and uniform growth,
central finiteness and cusp closure needed for the full original MMap.
\item[The remaining global producers.]
The eleven-field FamGLmin package and all its numerical, lattice and
loss consumers now give the exact same-family ChainTied and Closed
assembly. Construct the compatible second layer and duties, together
with the original denominator, retained-defect, period-height,
local-gain and spectral estimates. Their hypotheses have not been
discharged by completing the consumer.
\item[The original arithmetic counts.]
RH and two-primary targets still need I/II/V. CountII now has exact
reciprocity, required universal NiceI normalization, full original ArPl
under CohP, and residual values for any supplied integral type.
Its ordinary/control, type/neatU/TThetaEx, four further Hida endpoints,
same-J all-level HasAut, Choice, first-order/obstruction bounds and Raw
remain. Both original Links retain the threshold at fixed iota.
Unrestricted Dirichlet targets also retain III, SeedIV and chains.
\end{description}

The next work should be selected from these producer contracts by its
contribution to the final ten-input interface. The bounded global
assembly need not be revisited. The precise next source locations and
couplings are recorded in the version-0.35 continuation notes. The
optional AQ route and the DEX11 correction of historical RF16 remain.

\section{Claim ledger and reading order}

\begin{longtable}{>{\raggedright\arraybackslash}p{0.13\linewidth}>{\raggedright\arraybackslash}p{0.30\linewidth}>{\raggedright\arraybackslash}p{0.47\linewidth}}
\toprule
Tags & Content & Present status\\
\midrule
\endfirsthead
\toprule
Tags & Content & Present status\\
\midrule
\endhead
AN1--3 & Analytic endpoint & Complete one-sided-window implication for the fixed positive family.\\[3pt]
SC1--2 & Square-root balance & Complete scalar deduction; arithmetic supply remains.\\[3pt]
PL1--20 & Periods and common target & Descent, gluing, integrality and norm deductions; theta identities retained.\\[3pt]
EL1--10 & Actual coefficient lattices & Direct polynomial weight bound, full lattices, signed normalization, polars, covariance and neighbor maps proved.\\[3pt]
SY1--8 & Single-edge polynomial lift & Localizer, keys, finite correction and degree budget proved; former path-generator inputs now supplied by IP. RK supplies dimension; FG/OR/CP/EA complete local Type-II compatibility.\\[3pt]
CL1--12 & Complement and dual tests & Complementary representation and incoming-test inputs supplied by CB/IN/WT; RK supplies symbol rank; CP proves the actual cloud with constants $10,54$.\\[3pt]
IP1--11 & Integral path generators & Unimodular path basis, all-prefix filtration, exact span and Frobenius key image proved.\\[3pt]
IW1--7 & Integral duality and projectors & Finite averaging, denominator, normalized Fock duality, tensor/edge transport and isotypic estimate proved.\\[3pt]
PG1--11 & Constrained local Gram test & Explicit projection and derivative construction, factorial cancellation, source constants and zero remainder proved.\\[3pt]
DT1--11 & Digit corrections and test assembly & Factorial theorem, normal test, selection and outgoing membership proved; incoming \texttt{MemAt} now supplied by IN/WT.\\[3pt]
CB1--7 & Integral complementary basis & Monic basis, exact dimension, parity-selected count and integral Cartan multiplication proved.\\[3pt]
IN1--9 & Initial symbol membership & Exact blocked route, normal projections, zero chain loss and link loss $s/8$ proved.\\[3pt]
WT1--9 & Operator-word membership & Exact digit/carry calculus, signed weights and numerical word budget prove incoming \texttt{MemAt}.\\[3pt]
SD1--8 & Selection and dimension ingredients & Primitive selection, multiplier, degree and elementary counts proved; former density frontier closed by CG/LA/GP/RS/RK.\\[3pt]
CG1--9 & Coarse generator congruence & Actual path-lattice saturation, prime-field Frobenius, label splitting and individual-generator lift proved.\\[3pt]
LA1--6 & Actual segment assembly & Exact scale formula, residue bases, mixed-radix injection and full multidegree matching proved.\\[3pt]
GP1--12 & Good paths and segment rank & Interlacing/Abel tail estimate, actual membership and independence, and terminal residue count proved.\\[3pt]
RS1--8 & Complete residue summation & All degree margins, label sums, binomial comparisons and loss factors proved for the actual data.\\[3pt]
RK1--3 & Quantitative selected-symbol theorem & Every local rank, primitive-basis, complement and degree clause proved with constants $1088,1,10$.\\[3pt]
FG1--8 & Actual periodic frames & Complete FramesOK, cyclic wrap, strong gaps, span and counts proved from OrbitHyp.\\[3pt]
OR1--8 & Outgoing integral realization & Determinant extraction and exact monPh lift proved under the actual hypotheses for $C\ge100$.\\[3pt]
CP1--8 & Actual chain and cloud parameters & Floor-correct margins, signed-chain constructor, $50L$ flat and $54L$ cloud costs, and height constant $10$ proved.\\[3pt]
EA1--9 & Full local Type-II theorem & All compatibility, realization, primitive edge tests and assembly proved; final threshold $C\ge4129$.\\[3pt]
TI1--9 & Rank-one polynomial mechanism & Exact gamma and averaging proofs, now incorporated into the complete local Type-I theorem TC6.\\[3pt]
GG1--8 & Actual graded Gram lattice & Integral image lifting, highest projection and exact determinant constant proved.\\[3pt]
TP1--11 & Simultaneous tie partners & Sharp gamma integrality, compatible integral partners and literal weighted mean with $C_M=1$ proved.\\[3pt]
SW1--12 & Segment forms and weights & Harmonic block comparison and weighted mean proved. RU supplies coarse rank, NP the uniform hypotheses and final height.\\[3pt]
RL1--11 & Actual cleaned paired lifts & Factorial correction, raw tails, finite cleanup and source/dual transport proved at height $\lfloor t_0\rfloor$.\\[3pt]
RU1--11 & Coarse upper rank and HH assembly & Actual path injection, digit sum, weighted-lattice Smith selection and integral paired certificate proved.\\[3pt]
RD1--15 & Coarse lower rank and density & Weighted bad-tail count, kept injection, finite harmonic density and exact exceptional mean proved.\\[3pt]
HM1--12 & Actual harmonic and cloud maps & Integral projections, all edge normalizations, complementary contraction, tensor and coefficient-ring transport proved.\\[3pt]
NP1--9 & Uniform numerical supply & Literal Std, rounded windows, depth, conductors and shifted final height proved with explicit thresholds.\\[3pt]
TC1--6 & Full local Type-I theorem & Actual general-weight and edge assembly; Conclusion2 for $C_0\ge880$, $C_H\ge2800001$, $C=302114$.\\[3pt]
GII1--10 & Actual Type-II family supply & Good-orbit count, all numerical fields, local application and uniform probability; BO/BS/HC now supply the biased host selection.\\[3pt]
GI1--8 & Actual Type-I family supply & All fourteen numerical fields, exact band/window/depth translation, cyclic time bound and canonical averaged deficit.\\[3pt]
BC1--8 & Finite places and coefficient extension & Concrete standard ring, prescribed number-field place, integral lattice/dual/clean transport and exact LPBaseChange; recC gives StdSubAt.\\[3pt]
GJ1--5 & Family orbit application & Universal infimum, explicit scale bounds, exact orbit identifications and OrbitLatticeBounds with unchanged height/depth/loss.\\[3pt]
GW1--6 & One global degree pair & Exact marginals, total prime-loss estimate, MeanEx/WReq and ChosenW; optional CoverAll supplies a carrier.\\[3pt]
HW1--3 & Actual degree realization & Exact point/host WReq--WOK transfer, fixed-pair rebuilding and ChosenWall--ChosenWS; HC now expands the conditional host assembly.\\[3pt]
BO1--5 & Biased orbit probability & Forced-point conditioning, empty and parity-marked cells, uniform $C_bx^{-2}$ badness bound and actual orbit transport.\\[3pt]
BS1--7 & One balanced set of tops & Simultaneous variance/bad-mass selection, exact $4:3$ shares, prime-weighted bound and Fourier explanation.\\[3pt]
HC1--10 & Actual host and family join & Literal tops, no-relation set with $C_{T_2}=64$, elementary prime/field joins and conditional host/window assembly; analytic and field inputs retained.\\[3pt]
LI1--8 & Actual local identifications & All local places, completions, parity-correct positions, local lattice frames and integral away character; global comparison algebra now in SF/CR/CS.\\[3pt]
SF1--7 & One global frame & Exact standard matrix, actual diagonal companion lattice and simultaneous local frame replacement; full standard-forms construction retained.\\[3pt]
CR1--6 & Actual class-rank equality & Weyl fiber ranks, class bijection, two-prime stabilizer triviality and literal pack dimension equality; no metric conclusion.\\[3pt]
CS1--9 & Exact integral comparison & Valuation descent, odd-position Weyl/Gram compatibility and exact integral span; VE now supplies its polynomial-value premise.\\[3pt]
LP1--5 & Patch the selected lattices & Exact finite-place patch and ramification-normalized length; RF now supplies the canonical refill and conditional point installation.\\[3pt]
AD1--4 & Adelic constructions & Explicit diagonal base point, supported conjugate-pair splice, global-lattice recovery and exact rational frames.\\[3pt]
DR1--3 & Decorated realization & Local unramified Hermitian normal form including dyadic places, actual unitary adelic realizer, and every residue and rational frame condition.\\[3pt]
FC1--19 & Actual forms and carriers & Finite charts, normalized residues, Fock/period identities, common characters and levels, distinguished-block and dilation ties, direct Forms0; DR supplies realization.\\[3pt]
CV1--10 & Exact chart values & Actual complex member spaces, rational preimages and all-prime integral Weyl criteria; terminal values conditional on its row chart.\\[3pt]
NB1--8 & Actual neighbors & Paired-prime duality, literal up/down lattices and functionals, coupling labels, projections and complete naturality via the supported splice.\\[3pt]
VE1--11 & Concrete comparison & Explicit signed block values, universal ValEx and injectivity, complete comparison hypotheses and actual integral image span with $C_e=1$; separate realizability premise discharged.\\[3pt]
RF1--16 & Refill, point and parameters & Canonical refill, conditional point installation, full ParamsT and seven families; PN/SZ/AC now supply its remaining numerical and conditional constructor inputs.\\[3pt]
TL1--21 & Column laws and canonical reductions & Full Ok-to-Laws; finite character fields, tuple readings, orientation, principalization and conditional jet reductions; primitive local/analytic inputs retained.\\[3pt]
TG1--12 & Full reference symmetric columns & Exact terminal geometry, all forms and Hermitian carriers, residue normalization and complete ColsTiedSym; actual host/Ok/Pins/Later retained.\\[3pt]
PN1--16 & Complete parameter choice & All twenty-one families and every field of ParameterChoice, with the same data and uniform threshold order.\\[3pt]
SZ1--13 & Sizes, initial data and ranks & Complete SizesAt/NumAt, all initial size fields and Over identities, prescribed and cyclic frame rank constructors.\\[3pt]
AC1--10 & Arithmetic constructors & Exact root congruences, coded columns and ChGood; integral norm elements and full LevCh from the retained analytic LevelPrime input.\\[3pt]
AH1--18 & Character field and host & Exact mixed frame, cyclotomic fixed field, all derived arithmetic and half-different fields; conditional prescribed-character fieldChi.\\[3pt]
PA1--15 & Ambient period arithmetic & Actual normalized periods, rational/local/complex isometries, AmbIn.Ok, finite correction, CM-j and FieldOK; exact torsion/even transfer.\\[3pt]
TA1--14 & Normalized theta arithmetic & Pure analytic normalization, formal ODE and CM rigidity, all rational-translate algebraicity and exact arbitrary-field even-theta conjugacy.\\[3pt]
CI1--19 & Canonical local indices & Finite cosets/normal forms, projected characters, finite-field nonvanishing and literal initial/canonical index suppliers.\\[3pt]
TX1--16 & Complete all-twist comparison & Exact terminal and common levels, integral Law0, chain/form/level/row data and ten-field standard join; NChoice/TupTied retained.\\[3pt]
VI1--11 & Actual vacuum and algebraic germs & Trace-dual lattices, StdOK, Fourier inversion, normalized four-copy root, all rational translates and actual indicator jets.\\[3pt]
CA1--13 & Exact covariance and canonical algebraicity & Finite indicator bounds, signed intertwiners, Gaussian phase, all equivariance laws and same-data FinAlg/EvLg.\\[3pt]
NC1--17 & Normalized arithmetic choices & Complete-DVR readings, local lattice comparisons, all choice/character-tie/tuple fields; RankTie and exact support interfaces retained.\\[3pt]
SE1--16 & Restricted lattices and stage existence & Common-kernel descent, DivStage, row/spectral/Haar/jet models, conditional Later and pins normal form; NStage and geometric input retained.\\[3pt]
RC1--21 & Actual arithmetic compatibility & Scalar extension, character descent, choice/tuple joins, local level rigidity and finite integral/avoidance choices; realization and twisting inputs retained.\\[3pt]
LS1--22 & Actual periods and bare Galois stages & Witt/theta/de Rham/crystalline constructions, invariants and fundamental intersection; exact empty-index NStage and Later.\\[3pt]
IG1--14 & Original sized geometry & Qualitative ray primes, explicit host/factors/spaces/letters, integer mass index, original PreX and chosen-volume jets; quantitative core/windows retained.\\[3pt]
TD1--16 & Rational theta targets and metrics & Adjoint/level descent, conditional common image and exact matrices/metric, common test normalization; primitive theta identities retained.\\[3pt]
AR1--AR22 & Actual local arithmetic realization & Local character/level and Joint proofs; conditional weight descent, harmonic maps and installed series reading.\\[3pt]
HF1--HF20 & Actual host-character frame & Character extensions, exact local restrictions, determinant package and same-base initial joins; full core retained.\\[3pt]
TT1--TT19 & Primitive theta-test inputs & Finite/local rationality and support calculations; same-field analytic and period/operator inputs retained.\\[3pt]
QW1--QW15 & Quantitative windows & Complete conditional mass window, prime-count and scale transports, classical finite-conductor real-zero margin.\\[3pt]
GA1--GA23 & Normalized adic realization & Exact character construction and local laws; installed readings and nonempty same-object joins retained.\\[3pt]
WA1--WA19 & Finite Weil arithmetic & Same-vacuum algebraic span, normalized jet field and supplied-bottom periods; full existence and field joins retained.\\[3pt]
MF1--MF20 & Exact mass ratio & Local densities, global weighted averages and rank induction; classical Hasse--Minkowski named explicitly.\\[3pt]
QR1--QR25 & Uniform norm-prime supply & Actual ray family, exceptional-zero estimates and weighted interval theorem; ShortNorm and level/dyadic consequences.\\[3pt]
NR1--NR24 & Nonempty coefficient families & All-prime rings and inducing readings; bounded integral row families under explicit automorphic inputs.\\[3pt]
CW1--CW25 & Full Weil construction & Finite and real actions, exact rational phase cancellation, canonical and prescribed-character bottom data.\\[3pt]
GV1--GV21 & Actual geometric volumes & Exact gauge normalization and true-volume stages; classical Tamagawa and separate formal-degree inputs explicit.\\[3pt]
CC1--CC20 & Exact centered thinning & Dirichlet analytic inputs, exact integral and uniform late-cutoff bounds conditional on real-zero exclusion.\\[3pt]
MR1--MR26 & Actual form and local primitives & Fixed-state full-range forms, original-DVR good lattices/signatures; initial theta boundary explicit.\\[3pt]
QD1--QD22 & Quantitative Weil dictionary & Common scalar units, all-conjugate control and the exact totally-real-degree denominator exponent.\\[3pt]
DC1--DC22 & Large-gap discrete series & Constructive sufficient range, calibrated formal degree with named input, restricted RealTied.\\[3pt]
QP1--QP25 & Relative progressions and prime pairs & Continuum bounds, all-inducer PairEx and compatible-host RealBox under named inputs.\\[3pt]
HP1--HP6 & Closed coefficient representations & Principal minors, absolute irreducibility, classification, original-field descent and full weight supplier.\\[3pt]
VW1--VW4 & Cohomological window assembly & Exact same-field parity/conductor/budget chain and fixed-pair rescaling; deep suppliers retained.\\[3pt]
WB1--WB8 & Character-to-window closure & Ramified Fourier balance, biased averaging, scalar exponent and full conditional analytic endpoint.\\[3pt]
DA1--DA2 & Dependency and uniform-family audit & Finite threshold join, conditional global-link assembly and explicit semantic obligations.\\[3pt]
LD1--LD4 & Finite duality and reciprocity & Exact local/global conventions, unramified annihilators and one-place reciprocity.\\[3pt]
OS1--OS9 & Complete finite arithmetic supplier & Jacobi elements, Kolyvagin classes, length descent and literal original-field PsiBound.\\[3pt]
OB1--OB8 & Complete quantitative odd-minus bound & Finiteness, valuation budget, negative-twist annihilator and ThmBcFree with constant 240.\\[3pt]
LN1--LN10 & Conditional original-family LinkN & Integral cuts, congruence filtration, prime sum and exact restoration; local arithmetic inputs retained.\\[3pt]
LK1--LK4 & Exact native local-K interfaces & Finite nonfree duality, original cyclotomic cochains and the residue-characteristic Euler factor.\\[3pt]
TB1--TB10 & Complete all-prime torsion supplier & Ordered perfect-form nerve, orbits and stabilizers; TorsBound49 and TorsBoundPow22.\\[3pt]
KB1--KB8 & Actual KBridge and LemmaK31 & Classical K-theory maps, original local transport, finite cardinalities and one universal bound.\\[3pt]
TW1--TW8 & Complete simultaneous ThmBab & Class-rank translation, twisted towers, finite exceptions, common window and all field models.\\[3pt]
CZ1--CZ6 & Complete exceptional CountZ & Actual Selmer vanishing in the residual-trivial and residual-quadratic cases; uniform exponent zero.\\[3pt]
RV1--RV7 & Full rank-one finiteness and vector vanishing & Arbitrary finite characters, original local conditions, Baire capture and exact coefficient independence.\\[3pt]
BR1--BR7 & Full Bernoulli-unit and CorX supplier & Exact dyadic threshold, finite Frobenius projection, OS1 ideal relation and original residual field.\\[3pt]
SP1--SP6 & Complete LemmaMprime and LemmaM & Additive-injection cardinalities, Rees support, orbit multiplicity and the exact factor two.\\[3pt]
SP7--SP10 & Actual cocycle and paired-count join & Exact model/cotangent edge, trace product and coefficient return; actual model and catalogue retained.\\[3pt]
FP & Full universal LemmaP supplier & Framed representability and weighted obstruction presentation under the original group and local hypotheses.\\[3pt]
HV1--HV10 & CountVs, CountIV and full PointHV assembly & Integral character comparison at every odd prime, original fields and exact four-count/window quantifiers.\\[3pt]
AX1--AX10 & Actual symplectic data and G1 & Original ring, both cases, ordinary/level/dual/unlinking properties; image-density retained.\\[3pt]
DX1--DX8 & Original arithmetic local installation & Exact finite duality, framed dimensions, ordinary and finite-inertia charts and tangent conditions.\\[3pt]
MX1--MX10 & Complete actual Model constructor & All local axioms, homogeneous integral cuts, domain chart and cocycle-to-derivation fields.\\[3pt]
PX1--PX10 & Complete internal model comparison & Invariant chart, actual points and orbits, tangent rigidity and exact trace bound; actual arithmetic predicates retained.\\[3pt]
BX1--BX8 & Complete actual cohomology records & All tied periods, local Bloch--Kato, dual lattices and cohomological torsion.\\[3pt]
RX1--RX7 & Complete actual rigidity & Ordinary finite cocycles, cyclotomic irreducibility and the polarized adjoint join.\\[3pt]
CX1--CX9 & Complete rank-four catalogue & Original Galois level, all embedding/Frobenius quantifiers and a uniform integer trace bound.\\[3pt]
IX1--IX10 & Actual image-density subchains & Valued points, nonsplit blocks, purity and promodular transports; arithmetic producers retained.\\[3pt]
JX1--JX2 & Complete conditional count joins & Image density gives both paired counts; CountI/CountV/image density suffice for PointHV.\\[3pt]
KX1--KX9 & Complete universal arithmetic & Actual levels, reciprocity, weights, Phi, determinants and the universal Pair ring.\\[3pt]
QX & Adjoint and actual support joins & Two Links and the five-clause tower join; the precise rank-four Hecke producer remains.\\[3pt]
SX1--SX8 & Complete case-II seed & Actual cuspidality, fixed-central weights, induced lattice, local flags and automorphy.\\[3pt]
UX1--UX12 & Complete case-II promodularity & Actual towers, residual lines, spaces, component propagation and coefficient descent.\\[3pt]
YX1--YX2 & Complete conditional case-II count & PMExAlt gives CountII with all original point, length and coefficient quantifiers.\\[3pt]
HX1--HX4 & Original host installation & VWinAt(401) gives HostExW on the same field and pair, with cutoff400.\\[3pt]
MT1--MT7 & Exact metric transport & Fractional-ideal index, common metrics, complete floor support and uniform LinkMu.\\[3pt]
KA1--KA6 & CountI and target specialization & Complete PropGK-to-CountI aggregation; no Case III for power-of-two targets.\\[3pt]
ZW1--ZW4 & Literal tied-family numerical endpoint & Complete weighted ledger and uniform error bound; exponent $1/2+9\varepsilon$ on the selected route.\\[3pt]
AQ1--AQ5 & Positive-part analytic envelope & Fixed-stretch integral endpoint and integer interpolation; arithmetic supplier retained.\\[3pt]
JR1--12 & Joint local rank & Smith extraction, tensor exceptions, orbit cancellation and Nakayama; former coefficient-lattice input now covered by EL.\\[3pt]
LW1--4 & Three-slot lattice walk & Complete finite-state proof from the joint test datum; slots now constructed in EL.\\[3pt]
LC1--7, SR1--6 & Local integral costs & Tail integration, actual saturated transport and sharper index; metric transport remains separate.\\[3pt]
AG1--10 & Actual Gram estimates & Kernel reduction, recurrence and ideal-lattice decay; spectral/arithmetic inputs retained.\\[3pt]
FR1, GF1--9, GL1 & Frame and slopes & Full-family frame and clipped-deficit arguments, with the replacement inequality.\\[3pt]
OH1--8 & Ordinary host and windows & Fixed ordinary field and pair, elementary window deduction; selection and cohomology retained.\\[3pt]
GH1--7 & Dirichlet-GRH branch & Exact fixed character family, explicit $c$ and host, Fourier and rescaling proofs, primitive-to-all-character endpoint.\\
\bottomrule

DEX1--DEX11 & Compatible design assembly & Three same-class inputs; exact thinning, denominator, rank proviso, selector and kit joins.\\[3pt]
LTD1--LTD8 & Original local and theta dictionary & Actual specialization gives finite coverage; exact maps, operators and embedded fields retained.\\[3pt]
SLR1--SLR9 & Small-level rank consumer & Exact projection and base-only theta square; actual row/theta producer retained.\\[3pt]
RSL1--RSL13 & Second layer and ramified lattice & Original mass, finite characters, integral isometry and smaller fullness cut; remaining global slots explicit.\\[3pt]

DSZ1--DSZ10 & Actual degree and prime size & Exact original h3, harmonic upper comparison and construction-class TuSize.\\[3pt]
RMP1--RMP11 & Original class number and rank & Self-dual enlargement, actual local indices, mass, Mu0Tie and construction-class RankM.\\[3pt]
BTH1--BTH15 & Original base theta reduction & Integral convergence, actual Weil normalization and exact Hilbert/finite/rational consumer, with common producer inputs explicit.\\[3pt]
GLE1--GLE12 & Original local Hecke and length & Late setting, actual integral and ramified maps, exact second branch and composite bound.\\[3pt]

FIA1--FIA8 & Actual finite isotropic average & Original split/residue chart, half-datum guards, phase and uniform polynomial annihilation.\\[3pt]
CTC1--CTC13 & Original theta cuspidality & Compact quotients, orbit bounds, coinvariants, Fourier constant terms and genuine unipotent cuspidality.\\[3pt]
GLS1--GLS13 & Original column slopes & Actual ColVal, exact column slope, hook/Fock and constant-four bounds, same-selector FamGLmin.\\[3pt]
JSP1--JSP12 & Joint Hecke ideal and support & Kernel and annihilator, conditional RingProps/Dense, same-ideal TWBase and single-point P1 transport.\\[3pt]

GVC1--GVC18 & Original global counts and source slopes & Exact degree/Vandermonde comparison, numerical and weight clauses, CountF and SlopesMF.\\[3pt]
GFA1--GFA12 & Original quotient and frame consumers & Type-prime cost, actual corrections and support, all-height frame, exact nodes/windows/characters.\\[3pt]
TAE1--TAE13 & Original real theta extension & Smooth generators, all-word growth, dominated differentiation, stored vacuum and exact lowering.\\[3pt]
CIP1--CIP12 & Count-II reciprocity and normalization & Every source local map at split CM places, cyclotomic sign and full required NiceI Norm.\\[3pt]

LFS1--LFS17 & Original simultaneous losses & Exact prime classes, eleven terms, every-reading dictionary and the same realized degree pair.\\[3pt]
GFI1--GFI17 & Complete original global assembly & Every tied-ledger field, original chosen-data conversion and final Closed quantifiers.\\[3pt]
HCU1--HCU20 & Original lowest type and Hilbert coefficient & Full multiplicity one, holomorphic uniqueness, exact base model and conditional Hilbert realization.\\[3pt]
CAP1--CAP17 & Original auxiliary primes and residual values & Full ArPl under CohP, all lists/levels and uniform bound; supplied-type IsRes.\\[3pt]
\end{longtable}

Start with the version-0.35 audit in
Section~\ref{sec:dependency-audit-v35}. Read LFS and GFI for the
complete original loss and final-global assembly consumer; HCU for
the full original lowest type, coefficient uniqueness and Hilbert
map; CAP for the original auxiliary-prime supplier and residual
Hecke values. Their immediate inputs are GVC/GFA/GLS/RSL/GLE,
TAE/CW/BTH/FIA/CTC, and CIP/JSP/QX. Earlier chapters retain the
host, metric, counting and final zero-free arguments. Historical
frontiers are governed by the current audit and additive scope
updates, including the preserved RF16 correction.


GJ4--GJ5 in Section~\ref{sec:family-join} apply the local theorems;
GW/HW in Sections~\ref{sec:global-weight-choice}--\ref{sec:host-weight-join}
select and realize the global degrees. GI/GII give their uniform
numerical inputs, and BC the coefficient rings and finite extension.
The complete local Type-I theorem is TC6 in
Section~\ref{sec:type-one-closure}; the complete local Type-II theorem
remains Corollary~\ref{cor:typeII}. JR, LW and LC/SR convert local
tests to integral costs. The ordinary and general-character host
sections distinguish their arithmetic specializations. Historical
reports retain their wording; current annotations and the continuation
checkpoint explain which earlier boundaries have been discharged.

Every earlier claim ID is preserved, with version-0.35 updates where an input is now discharged or narrowed. Earlier component reports are preserved as historical records; the current main note and merged ledger state the updated boundaries. A retained source statement is not counted as a proof in this text. A derived simplification is marked separately from a source translation. No Lean source was modified, and no full Lean build or independent kernel replay was run for this reconstruction.

\section{Appendix: an elementary ideal-lattice decay lemma}
\label{app:ag-lattice}

This argument is a derived alternative to the archive's
lattice-sum step; it uses only the algebraic norm and elementary
packing.  It does not prove the representation-theoretic Fourier
estimates to which the lemma is applied.

\paragraph{Statement.}
Let \(E\) be a totally complex number field of degree \(2f\),
\(f\ge1\), and choose exactly one embedding
\(\sigma_\nu:E\hookrightarrow\mathbb C\) from each complex-conjugate
pair.  Let \(\mathfrak a\ne0\) be a fractional ideal.  Its absolute
norm \(N\mathfrak a\) is the positive rational norm, extended
multiplicatively from integral ideals.  For \(x_\nu>0\), \(K\ge3\)
real, and \(\Theta\ge1\), assume
\[
 \prod_{\nu=1}^{f}x_\nu\le
        \frac{\sqrt{N\mathfrak a}}{\Theta^f}.
\]
Then
\[
 \sum_{\substack{z\in\mathfrak a\\z\ne0}}
 \prod_{\nu=1}^{f}
       \min\{1,(x_\nu/|\sigma_\nu z|)^K\}
 \le \bigl(384\,\Theta^{-(K-1/2)}\bigr)^f.                 \tag{AG.A1}
\]
The bound is formally valid at \(\Theta=1\); its smallness is
useful only when \(\Theta\) is sufficiently large compared with
the fixed constant.

\paragraph{Proof.}
For \(0\ne z\in\mathfrak a\), the ideal
\((z)\mathfrak a^{-1}\) is integral.  Its norm is a positive
integer, so
\[
 |N_{E/\mathbb Q}z|\ge N\mathfrak a,\qquad
 \prod_{\nu=1}^{f}|\sigma_\nu z|\ge\sqrt{N\mathfrak a}.
\]
The square root is essential: each chosen complex embedding
represents two embeddings in the field norm.
Write \(Z_\nu(z)=\sigma_\nu(z)/x_\nu\).
For distinct \(z,z'\in\mathfrak a\),
\[
 \prod_\nu|Z_\nu(z)-Z_\nu(z')|\ge\Theta^f.                 \tag{AG.A2}
\]

For radii \(R_\nu\ge1\), put \(P=\prod_\nu R_\nu\) and let
\(M(R)\) count the nonzero \(z\in\mathfrak a\) with
\(|Z_\nu(z)|\le R_\nu\) for every \(\nu\).
If \(P<\Theta^f\), then \(M(R)=0\).
Otherwise put \(q=P^{1/f}/\Theta\ge1\) and
\(h=\lceil2q\rceil\).
At each complex place partition the square
\([-R_\nu,R_\nu]^2\) into half-open squares of side \(R_\nu/h\),
assigning the outer boundaries consistently.
At most \((2h+2)^2\le(8q)^2\) squares are needed.
Two points lying in the same product of these squares would
satisfy
\[
 \prod_\nu|Z_\nu(z)-Z_\nu(z')|
 \le(\sqrt2/h)^fP
 \le 2^{-f/2}\Theta^f<\Theta^f,
\]
contradicting (AG.A2).  Thus each product box contains at most
one point, and the packing is in \(2f\) real dimensions:
\[
 M(R)\le64^f\,P^2/\Theta^{2f}
       \quad\hbox{if }P\ge\Theta^f.                        \tag{AG.A3}
\]

First take \(K=3\).  The nonnegative layer-cake identity
\[
 \min\{1,t^{-3}\}=\int_1^\infty 3R^{-4}\,
                                      \mathbf1_{t\le R}\,dR
 \qquad(t>0)
\]
and Tonelli's theorem give, using (AG.A3),
\begin{align*}
 \sum_{z\ne0}\prod_\nu\min\{1,|Z_\nu(z)|^{-3}\}
 &\le 192^f\Theta^{-2f}
 \int_{\substack{R_\nu\ge1\\\prod R_\nu\ge\Theta^f}}
                         \prod_\nu R_\nu^{-2}\,dR_\nu\\
 &=192^f\Theta^{-2f}
 \int_{\substack{t_\nu\ge0\\\sum t_\nu\ge f\log\Theta}}
                            e^{-\sum t_\nu}\,dt\\
 &\le192^f\Theta^{-5f/2}
       \int_{t_\nu\ge0}e^{-\frac12\sum t_\nu}\,dt\\
 &=384^f\Theta^{-5f/2}.
\end{align*}
The substitution in the second line is \(R_\nu=e^{t_\nu}\).
For general \(K\ge3\), every nonzero lattice point has
\(\prod_\nu\max\{1,|Z_\nu(z)|\}\ge\Theta^f\).  Consequently
\[
 \prod_\nu\min\{1,|Z_\nu(z)|^{-K}\}
 \le \Theta^{-(K-3)f}
                   \prod_\nu\min\{1,|Z_\nu(z)|^{-3}\}.
\]
Multiplying the preceding bound proves (AG.A1).

\paragraph{Relation to the archived estimate.}
The source's \texttt{latSum\_of\_core} proves the sufficient
bound \((C\Theta^{-(K-5/2)})^f\), for integer \(K\ge3\)
and \(\Theta>2\).  Formula (AG.A1) implies this form directly
with \(C=384\).  Thus it can replace that single lattice-sum
step in the human argument while all downstream estimates are
left at their archived strength.  Its validity does not depend
on the root discriminant, on a chosen integral basis, or on the
shape of the embedded fractional ideal.


\begin{thebibliography}{99}\small\raggedright
\bibitem{announcement} L. Alp\"oge, announcement and remarks, 20 November 2026. \href{https://alpo.ge/remark.pdf}{Original four-page PDF}.
\bibitem{verification} A. Sutherland, \emph{Verification files for the Lean proof of the Riemann Hypothesis}. Snapshot inspected: 20 November 2026, 22:05 UTC. \url{https://math.mit.edu/~drew/grh/}.
\bibitem{main} Released source mirror, \href{https://github.com/CocoSgt/riemann-hypothesis-lean/tree/7cbf7a0c32fca1badf05931879a37c3f8a2b03b4}{pinned revision}. Entry points: \src{Main.lean}{Main.lean}; \src{Statements.lean}{Statements.lean}.
\bibitem{charfamily} \src{Development/AutomorphicArithmetic/DedekindLFunctions/CharFamily.lean}{CharFamily.lean}; \src{Formalism/SharpRestriction/RamificationFrobenius/TotientGaloisClasses/Analytic.lean}{Analytic.lean}. Definitions of the centered ramps and positive families.
\bibitem{level} \src{Formalism/SharpRestriction/RamificationFrobenius/TotientGaloisClasses/ExplicitFormulaErrorOfHostConditions.lean}{ExplicitFormulaErrorOfHostConditions.lean}. Cyclic family, weights, and level definitions.
\bibitem{zerofree} \src{Development/AutomorphicArithmetic/DedekindLFunctions/LandauTheorems/ZeroFreeOfRpowWindow.lean}{ZeroFreeOfRpowWindow.lean}; \src{Development/AutomorphicArithmetic/DedekindLFunctions/LandauTheorems/lemmaL.lean}{lemmaL.lean}.
\bibitem{landau} \src{Development/AutomorphicArithmetic/DedekindLFunctions/LandauTheorems/landauWindow.lean}{landauWindow.lean}; \src{Development/AutomorphicArithmetic/DedekindLFunctions/LandauTheorems/landauWindowIneq.lean}{landauWindowIneq.lean}; \src{Development/AutomorphicArithmetic/DedekindLFunctions/LandauTheorems/landauTails.lean}{landauTails.lean}. Nearest-pole and gamma-kernel argument.
\bibitem{poles} \src{Development/AutomorphicArithmetic/DedekindLFunctions/LandauTheorems/poleData.lean}{poleData.lean}; \src{Development/AutomorphicArithmetic/DedekindLFunctions/LandauTheorems/mellinGchiOfNeOne.lean}{mellinGchiOfNeOne.lean}; \src{Development/AutomorphicArithmetic/DedekindLFunctions/LandauTheorems/mellinGchiModOne.lean}{mellinGchiModOne.lean}.
\bibitem{critline} \src{Formalism/GammaFactorNonvanishing/CritLine.lean}{CritLine.lean}; \src{Development/AutomorphicArithmetic/CharacterBoxes/ClaimStrip.lean}{ClaimStrip.lean}.
\bibitem{size} \src{Formalism/SizeWindowBrackets/GrowthBoundGivesRiemannHypothesis.lean}{GrowthBoundGivesRiemannHypothesis.lean}; \src{Formalism/SizeWindowBrackets/ChainSlopeWindows/LipschitzWeightAndMinimum.lean}{LipschitzWeightAndMinimum.lean}. Parameter hypotheses and scalar estimate.
\bibitem{balance} \src{Development/AutomorphicArithmetic/DedekindLFunctions/LandauTheorems/ChBalance.lean}{ChBalance.lean}. Exact arithmetic/character identity.
\bibitem{chcore} \src{Development/AutomorphicArithmetic/DedekindLFunctions/LandauTheorems/ChCore.lean}{ChCore.lean}. Averaging and scalar bridge.
\bibitem{coregm} \src{Development/AutomorphicArithmetic/PrimitiveCosets/CoreGm.lean}{CoreGm.lean}. Final exponent $1/2+(e_D+e_L+7)\varepsilon$.
\bibitem{gm} \src{Development/AutomorphicArithmetic/PrimitiveCosets/GmAEq.lean}{GmAEq.lean}. Identification with the fixed family.
\bibitem{sigma} \src{Development/AutomorphicArithmetic/PrimitiveCosets/SigALe.lean}{SigALe.lean}. Bound for omitted characters.
\bibitem{topwin} \src{Formalism/SharpRestriction/RamificationFrobenius/TotientGaloisClasses/TopWin.lean}{TopWin.lean}; \src{Development/AutomorphicArithmetic/PrimitiveCosets/LevelGmLeOnFamilyScales.lean}{LevelGmLeOnFamilyScales.lean}. Window quantifiers and passage to zero-freeness.
\bibitem{frame} \src{Formalism/HarderNarasimhan/StatementFrameF.lean}{StatementFrameF.lean}; \src{Formalism/HarderNarasimhan/HarderNarasimhanProofs/FrameLinearAlgebraFirst.lean}{FrameLinearAlgebraFirst.lean}; \src{Formalism/HarderNarasimhan/HarderNarasimhanProofs/TheoremFSharp.lean}{TheoremFSharp.lean}.
\bibitem{division} \src{Formalism/LocalToGlobal/Divisibility/DivisionHeightBoundStatement.lean}{DivisionHeightBoundStatement.lean}; \src{Formalism/LocalToGlobal/Divisibility/ConstantGeTwentyEight.lean}{ConstantGeTwentyEight.lean}.
\bibitem{congruence} \src{Formalism/SizeWindowBrackets/ChainSlopeWindows/LevelOfCongruence.lean}{LevelOfCongruence.lean}. Common-target slope cancellation.
\bibitem{chtop} \src{Development/AutomorphicArithmetic/DedekindLFunctions/LandauTheorems/ChTop.lean}{ChTop.lean}. Output as an arithmetic gap bound.
\bibitem{inputs} \src{Formalism/SharpRestriction/GlobalTheoremInputs.lean}{GlobalTheoremInputs.lean}. Quantitative arithmetic interfaces.
\bibitem{host} \src{Development/AutomorphicArithmetic/CharacterBoxes/ExistsPageN.lean}{ExistsPageN.lean}. Ordinary host, auxiliary primes, and window data.
\bibitem{conclusion} \src{Development/AutomorphicArithmetic/CharacterBoxes/MainConclusion.lean}{MainConclusion.lean}; \src{Development/AutomorphicArithmetic/CharacterBoxes/ZeroFreeAllOfTopG.lean}{ZeroFreeAllOfTopG.lean}. Assembly of the unconditional source theorem.
\bibitem{tao} T. Tao, announcement and discussion, 20 November 2026; see M. Jaswal's code-reading comments. \href{https://terrytao.wordpress.com/2026/11/20/unforced-euler-singularity-unforced-navier-stokes-singularity-and-the-generalized-riemann-hypothesis/comment-page-1/}{Tao's blog discussion}.
\bibitem{guide} \texttt{bramble\_trace}, \emph{Reading Alp\"oge's proposed proof of RH from theta lifts to prime sums}, rStation, revisions 33--34. \href{https://rstation.ai/r/riemann_hypothesis/posts/100}{Reading guide}; \href{https://rstation.ai/api/v1/posts/uw5cfb_J_Cj7XJO6/attachment?revision=34}{supporting bundle}. Agent-authored source map; \href{https://rstation.ai/about}{platform description}.

\bibitem{local-selection} \src{Formalism/SharpRestriction/GrowthBounds/SelectionLemmas.lean}{SelectionLemmas.lean}. Candidate lattice, depth allowance, and selection conclusion.
\bibitem{local-smith} \src{Formalism/SubspaceCounting/Statement.lean}{SubspaceCounting/Statement.lean}; \src{Formalism/SubspaceCounting/SubspaceCountingProofs/AllButLESmith.lean}{AllButLESmith.lean}. Exceptional directions and Smith exponents.
\bibitem{local-systems} \src{Formalism/SharpRestriction/GrowthBounds/LatticeSystems.lean}{LatticeSystems.lean}. Elementary-symmetric tail bound, index estimate, and theta-kernel expansion.
\bibitem{local-normalization} \src{Formalism/LocalRingCardinality/DataArith.lean}{DataArith.lean}; \src{Formalism/LocalRingCardinality/RankBoundMonotoneInDimension.lean}{RankBoundMonotoneInDimension.lean}. Finite-place normalization.
\bibitem{local-div-defs} \src{Formalism/LocalRingCardinality/DataDIVDefs.lean}{DataDIVDefs.lean}; \src{Formalism/LocalRingCardinality/DataParams₂.lean}{DataParams2.lean}. Shifts, reference lattices, and gain definitions.
\bibitem{local-rounding} \src{Formalism/SizeWindowBrackets/ChainSlopeWindows/AnalyticFirstToLatticeDiscriminant.lean}{AnalyticFirstToLatticeDiscriminant.lean}. The floor loss and its explicit allowance.
\bibitem{local-costs} \src{Formalism/SizeWindowBrackets/ClaimGaFree.lean}{ClaimGaFree.lean}; \src{Formalism/LocalToGlobal/GlobalTheory/FamRec.lean}{FamRec.lean}. Full versus retained ranks, exceptional-place costs, and ramified mean corrections.

\bibitem{walksharp} \src{Formalism/SubspaceCounting/SubspaceCountingProofs/WalkSharp.lean}{WalkSharp.lean}. Finite-state proof of the lattice walk.
\bibitem{saturated-transport} \src{Formalism/LocalRingCardinality/Lemmas.lean}{LocalRingCardinality/Lemmas.lean}. Actual ramified reference, selected lattice, and normalized length comparison.
\bibitem{actual-selected-lattice} \src{Development/AbelianVarietyTheory/Proofs/GlobalInterpolation/ExistsLedPart.lean}{ExistsLedPart.lean}. Concrete construction of the selected global lattice.
\bibitem{joint-tests} \src{Formalism/SharpRestriction/GrowthBounds/ScaleEdges/GoodRankOfPrim.lean}{GoodRankOfPrim.lean}; \src{Formalism/SharpRestriction/GrowthBounds/ScaleEdges/ColSpanMap.lean}{ColSpanMap.lean}. Primitive pairing tests and graph-plane edge count.
\bibitem{actual-gram-kernel} \src{Formalism/SharpRestriction/CuspSpanNormBounds/AdmissibleBoundsFromSource.lean}{AdmissibleBoundsFromSource.lean}; \src{Formalism/SharpRestriction/CuspSpanNormBounds/StandingHypothesesHoldFamily.lean}{StandingHypothesesHoldFamily.lean}. Actual diagonal, mixed, and dimension bounds.
\bibitem{actual-gram-words} \src{Formalism/SharpRestriction/CuspSpanNormBounds/BlockEstimateWords.lean}{BlockEstimateWords.lean}; \src{Formalism/SharpRestriction/CuspSpanNormBounds/OscillatoryIntegralBounds.lean}{OscillatoryIntegralBounds.lean}. Finite word closure, Poisson reduction, and archived lattice-sum estimate.

\bibitem{edge-coefficients} \src{Formalism/GradedWeylDivisibility/StatementA.lean}{StatementA.lean}; \src{Formalism/SharpRestriction/GrowthBounds/UnramifiedQuadraticSetting.lean}{UnramifiedQuadraticSetting.lean}. Literal dual-Weyl polynomial model and coefficient lattices; precise declarations are in EL.
\bibitem{edge-operators} \src{Formalism/SharpRestriction/GrowthBounds/LatticeSystems.lean}{LatticeSystems.lean}. Coefficient comparisons, equivariance and normalized neighbor operators.
\bibitem{symbol-source} \src{Formalism/SharpRestriction/GrowthBounds/SourceCleaning.lean}{SourceCleaning.lean}; \src{Formalism/SharpRestriction/GrowthBounds/WittFrameLocaliserOrders.lean}{WittFrameLocaliserOrders.lean}. Explicit localizer, determinant and pairing keys, correction, and symbol realization.
\bibitem{symbol-elementary} \src{Formalism/PolynomialDimension/DetMulEqSumAdjugate.lean}{DetMulEqSumAdjugate.lean}. Finite-field telescoping and finite correction induction.
\bibitem{cloud-construction} \src{Formalism/SharpRestriction/GrowthBounds/DualFockLinear/TheoremIIClosedOfTestBound.lean}{TheoremIIClosedOfTestBound.lean}. Complementary family and primitive-product argument.
\bibitem{cloud-flat} \src{Formalism/SharpRestriction/GrowthBounds/DiagonalSupport/Proposition5OfThreeParts.lean}{Proposition5OfThreeParts.lean}. Flat valuation argument; the former pointwise projection/Gram inputs are expanded in PG/DT.
\bibitem{grh-statements} \src{Statements.lean}{Statements.lean}; \src{Supplements/Kernels/MainTheorems.lean}{MainTheorems.lean}. Exact all-Dirichlet scope and exported endpoints.
\bibitem{grh-host} \src{Development/AutomorphicArithmetic/CharacterBoxes/OrdAny.lean}{OrdAny.lean}; \src{Development/AutomorphicArithmetic/CharacterBoxes/IsChiField.lean}{IsChiField.lean}. Character field and degree stabilization; GH gives a direct explicit formula.

\bibitem{path-basis} \src{Formalism/SharpRestriction/GrowthBounds/ObjPol.lean}{ObjPol.lean}. Signed column paths and their polynomial model; IP records the exact basis and prefix-filtration declarations.
\bibitem{path-symbol} \src{Formalism/SharpRestriction/GrowthBounds/PathWeight/LetterWeightPos.lean}{LetterWeightPos.lean}. Actual path-generator symbol image; IP gives a shorter route through split columns and their Frobenius keys.
\bibitem{weyl-integral} \src{Formalism/ThetaLogCoordinates/TraceIsometryWeights/SchurTraceProofs/InvInt.lean}{InvInt.lean}; \src{Formalism/ThetaLogCoordinates/TraceIsometryWeights/SchurTraceProofs/TypeInt.lean}{TypeInt.lean}. Integral finite averaging and isotypic projector denominator.
\bibitem{weyl-form} \src{Formalism/SharpRestriction/GrowthBounds/DualFockLinear/FormCoreSKXOf.lean}{FormCoreSKXOf.lean}. Exact normalized Fock lattice comparison; IW traces the defining square polynomial model and opposite-edge transports.
\bibitem{gram-construction} \src{Formalism/SharpRestriction/GrowthBounds/MonomialOrderSegments/LiftBoundOfHypotheses.lean}{LiftBoundOfHypotheses.lean}. Constrained projection--derivative construction; PG records the complete source route and its zero remainder.
\bibitem{digit-approximation} \src{Formalism/SharpRestriction/GrowthBounds/DiagonalSupport/BinomCombFOnFinset.lean}{BinomCombFOnFinset.lean}. Source factorial-unit interfaces; DT supplies the alternative degree-$2(c-1)$ proof.
\bibitem{test-membership} \src{Formalism/SharpRestriction/GrowthBounds/DiagonalSupport/PickCore.lean}{PickCore.lean}; \src{Formalism/SharpRestriction/GrowthBounds/DiagonalSupport/ExistsMapOfCoefGe.lean}{ExistsMapOfCoefGe.lean}; \src{Formalism/SharpRestriction/GrowthBounds/TwistedSubstitution/MemRed.lean}{MemRed.lean}. Selection and outgoing assembly; the distinct incoming word closure is now proved in IN/WT.

\bibitem{complement-basis} \src{Formalism/GradedWeylDivisibility/FinalC.lean}{FinalC.lean}; \src{Formalism/GradedWeylDivisibility/FinalD.lean}{FinalD.lean}. Integral column-invariant basis. CB expands the row-peeling and leading-term proof and derives its count by a finite determinant identity.
\bibitem{complement-cartan} \src{Formalism/LocalToGlobal/GroupLattices/FockModels/Cartan.lean}{Cartan.lean}; \src{Formalism/SharpRestriction/GrowthBounds/DualFockLinear/RestCloudBOfTwo.lean}{RestCloudBOfTwo.lean}. Integral multiplication and actual edge transport; exact locators are in CB6--CB7.
\bibitem{initial-blocks} \src{Formalism/SharpRestriction/GrowthBounds/PlaneCoefficients/PlaneCoeffOneTTens.lean}{PlaneCoeffOneTTens.lean}. The actual blocked initial-membership route. IN records every defining model and the coefficient, frame and degree calculations.
\bibitem{incoming-words} \src{Formalism/SharpRestriction/GrowthBounds/PlaneCoefficients/RestOneOpOf.lean}{RestOneOpOf.lean}; \src{Formalism/SharpRestriction/GrowthBounds/PlaneCoefficients/WBound.lean}{WBound.lean}; \src{Formalism/SharpRestriction/GrowthBounds/TwistedSubstitution/MemRed.lean}{MemRed.lean}. Single operators, exact word budget and incoming closure; all three are expanded in WT.
\bibitem{symbol-selection} \src{Formalism/SharpRestriction/GrowthBounds/SegmentDegree/SelIsComplRestrictSupport.lean}{SelIsComplRestrictSupport.lean}; \src{Formalism/SharpRestriction/GrowthBounds/SegmentDegree/StrictMonoIntExists.lean}{StrictMonoIntExists.lean}. Primitive generator selection and normal-degree bound. SD gives the shorter coefficient-minor proof and retains a sharper intermediate degree bound.
\bibitem{symbol-density} \src{Formalism/SharpRestriction/GrowthBounds/SegmentDegree.lean}{SegmentDegree.lean}; \src{Formalism/SharpRestriction/GrowthBounds/SegmentDegree/RestCountOf.lean}{RestCountOf.lean}. Exact numerical witnesses and selected-symbol rank assembly. CG/LA/GP/RS now expand its substantive inputs; RK proves the complete supplied-data conclusion.

\bibitem{coarse-lift} \src{Formalism/SharpRestriction/GrowthBounds/SegmentDegree/ColumnWordReductionOfFour.lean}{ColumnWordReductionOfFour.lean}; \src{Formalism/SharpRestriction/GrowthBounds/SegmentDegree/DiscriminantNeZeroOfResidueNeZero.lean}{DiscriminantNeZeroOfResidueNeZero.lean}. Actual congruence and homogeneous lift, reconstructed in CG.
\bibitem{segment-assembly} \src{Formalism/SharpRestriction/GrowthBounds/SegmentDegree/RestLamOfRests.lean}{RestLamOfRests.lean}; \src{Formalism/SharpRestriction/GrowthBounds/SegmentDegree/LamImgApply.lean}{LamImgApply.lean}. Actual residue-space injection and degree assembly, reconstructed in LA.
\bibitem{good-paths} \src{Formalism/SharpRestriction/GrowthBounds/SegmentDegree/GoodPathWeight.lean}{GoodPathWeight.lean}; \src{Formalism/ThetaLogCoordinates/ResidueClassEstimates/FinalT.lean}{FinalT.lean}. Weighted paths and their tail estimate, reconstructed and shortened in GP.
\bibitem{residue-sum} \src{Formalism/SharpRestriction/GrowthBounds/SegmentDegree/FrameZeroStrictMono.lean}{FrameZeroStrictMono.lean}; \src{Formalism/SharpRestriction/GrowthBounds/SegmentDegree/SumTotalLoss.lean}{SumTotalLoss.lean}. Exact parameter bounds and full summation, reconstructed in RS.
\bibitem{rank-closure} \src{Formalism/SharpRestriction/GrowthBounds/ValuationLowestDegreeRename.lean}{ValuationLowestDegreeRename.lean}; \src{Formalism/SharpRestriction/GrowthBounds/SegmentDegree/NonemptyDvrSettingOfStandard.lean}{NonemptyDvrSettingOfStandard.lean}. Literal Conclusion2 and assembly of the constants, completed in RK.


\bibitem{frames-actual} \src{Formalism/SharpRestriction/GrowthBounds/WittFrameLocaliserOrders.lean}{WittFrameLocaliserOrders.lean}. Full periodic frames, strict spans, gaps and counts, reconstructed in FG.
\bibitem{outgoing-actual} \src{Formalism/SharpRestriction/GrowthBounds/SourceCleaning.lean}{SourceCleaning.lean}; \src{Formalism/SharpRestriction/GrowthBounds/PathWeight/PhProdShift.lean}{PhProdShift.lean}; \src{Formalism/SharpRestriction/GrowthBounds/PathWeight/FixedSymbolFactorIntegral.lean}{FixedSymbolFactorIntegral.lean}. Actual outgoing determinant extraction and one-sided realization, reconstructed in OR.
\bibitem{cloud-actual} \src{Formalism/SharpRestriction/GrowthBounds/DiagonalSupport/LeOfOneLeMul.lean}{LeOfOneLeMul.lean}; \src{Formalism/SharpRestriction/GrowthBounds/DualFockLinear/TestSideExplicitIdentification.lean}{TestSideExplicitIdentification.lean}. Actual numerical chain data, normalized flat tests and uniform cloud, reconstructed in CP.
\bibitem{edges-actual} \src{Formalism/SharpRestriction/GrowthBounds/ScaleEdges/ManyEdges.lean}{ManyEdges.lean}; \src{Formalism/SharpRestriction/GrowthBounds/ScaleEdges/InitialLevelMonoWeights.lean}{InitialLevelMonoWeights.lean}. Literal Type-II hypotheses and final uniform edge assembly, reconstructed in EA.
\bibitem{type-one-actual} \src{Formalism/SharpRestriction/GrowthBounds/ObjSeg.lean}{ObjSeg.lean}; \src{Formalism/MahlerMeasure/MahlerMeasureProofs/GamDeg.lean}{GamDeg.lean}; \src{Formalism/MahlerMeasure/MahlerMeasureProofs/GamLead.lean}{GamLead.lean}; \src{Formalism/MahlerMeasure/MahlerMeasureProofs/WindowAP.lean}{WindowAP.lean}; \src{Formalism/SharpRestriction/GrowthBounds/ProjectionLemmaBaseLevel.lean}{ProjectionLemmaBaseLevel.lean}. Actual Type-I segment, contraction polynomials, matrix averaging and final local interfaces.


\bibitem{graded-gram-actual} \src{Formalism/MahlerMeasure/StatementT.lean}{StatementT.lean}; \src{Formalism/ThetaLogCoordinates/PartnerMeanCoordinates/PartnerMeanProofs/GramHypothesisOfLogBound.lean}{GramHypothesisOfLogBound.lean}. The actual restricted graded Gram hypothesis and its partner-mean interface; GG reconstructs the coefficient-image lattice argument.
\bibitem{tie-partners-actual} \src{Formalism/ThetaLogCoordinates/PartnerMeanCoordinates/PartnerMeanProofs/GamIntC.lean}{GamIntC.lean}; \src{Formalism/ThetaLogCoordinates/PartnerMeanCoordinates/PartnerMeanProofs/GammaZeroIntegral.lean}{GammaZeroIntegral.lean}. Sharp gamma integrality, compatible tie partners and the exact multiple-tie mean, reconstructed in TP.
\bibitem{segment-weighted-actual} \src{Formalism/MahlerMeasure/MahlerMeasureProofs/SSDirect.lean}{SSDirect.lean}; \src{Formalism/SharpRestriction/GrowthBounds/MonomialFactorials/BlockBoundOfDualFamilies.lean}{BlockBoundOfDualFamilies.lean}; \src{Formalism/SharpRestriction/GrowthBounds/MonomialFactorials/SegmentWeightedCountingBoundOfAllRemainingBounds.lean}{SegmentWeightedCountingBoundOfAllRemainingBounds.lean}. Actual block lattices, good dual families and dimension-weighted summation, reconstructed in SW with its coarse rank input explicit.
\bibitem{segment-numerics-actual} \src{Formalism/SharpRestriction/GrowthBounds/CoarseMultiplicity/PolyUpperBoundDelta.lean}{PolyUpperBoundDelta.lean}; \src{Formalism/SharpRestriction/GrowthBounds/BinomialRatio/TopDegreeLeTwoB.lean}{TopDegreeLeTwoB.lean}. Kept-block, terminal-degree and total-order estimates from the literal standard hypotheses.
\bibitem{paired-lifts-actual} \src{Formalism/SharpRestriction/GrowthBounds/CoarseMultiplicity/RestBlockForm.lean}{RestBlockForm.lean}; \src{Formalism/SharpRestriction/GrowthBounds/CoarseMultiplicity/PairedLiftAtOfIdentity.lean}{PairedLiftAtOfIdentity.lean}; \src{Formalism/SharpRestriction/GrowthBounds/BinomialRatio/PairLift.lean}{PairLift.lean}. Actual factorial correction, raw polynomial tails, finite cleanup and source/dual edge pairing, reconstructed in RL.


\bibitem{coarse-rank-upper} \src{Formalism/SharpRestriction/GrowthBounds/BinomialRatio/RankOnDigitsCast.lean}{RankOnDigitsCast.lean}; \src{Formalism/SharpRestriction/GrowthBounds/BinomialRatio/ExtendedStatementOfParts.lean}{ExtendedStatementOfParts.lean}. RU proves the actual upper count, digit-to-HH dimension comparison and integral paired assembly.
\bibitem{coarse-rank-lower} \src{Formalism/SharpRestriction/GrowthBounds/ColumnSpanDensity/ThresholdDensityFacts.lean}{ThresholdDensityFacts.lean}; \src{Formalism/SharpRestriction/GrowthBounds/ColumnSpanDensity/LocalHypothesisOfGlobal.lean}{LocalHypothesisOfGlobal.lean}. RD records the complete lower-rank source route and proves the density and final mean, using a shorter finite telescoping argument.
\bibitem{harmonic-cloud-actual} \src{Formalism/SharpRestriction/GrowthBounds/ColumnSpanDensity/CloudPropOfMapsRetractionBasis.lean}{CloudPropOfMapsRetractionBasis.lean}; \src{Formalism/SharpRestriction/GrowthBounds/ColumnSpanDensity/RestHLTransport.lean}{RestHLTransport.lean}. HM constructs the actual harmonic and complementary-middle maps, with exact integral normalization and coefficient-ring transport.
\bibitem{type-one-numerical} \src{Formalism/SharpRestriction/GrowthBounds/BinomialRatio/ThresholdLowerBounds.lean}{ThresholdLowerBounds.lean}; \src{Formalism/SharpRestriction/GrowthBounds/ColumnSpanDensity/ThresholdFixedWeightFacts.lean}{ThresholdFixedWeightFacts.lean}; \src{Formalism/SharpRestriction/GrowthBounds/ColumnSpanDensity/SegStartSpec.lean}{SegStartSpec.lean}. NP supplies the literal numerical fields and the final shifted height.
\bibitem{type-one-complete} \src{Formalism/SharpRestriction/GrowthBounds/ColumnSpanDensity/CloudCount.lean}{CloudCount.lean}; \src{Formalism/SharpRestriction/GrowthBounds/ColumnSpanDensity/LatticeCountPrimeMonoLevel.lean}{LatticeCountPrimeMonoLevel.lean}; \src{Formalism/SharpRestriction/GrowthBounds/ProjectionLemmaBaseLevel.lean}{ProjectionLemmaBaseLevel.lean}. TC proves the actual local Conclusion2 with explicit sufficient constants, bypassing unused bundled admissibility fields.


\bibitem{type-two-family} \src{Formalism/SharpRestriction/GrowthBounds/TopChoiceBounds.lean}{TopChoiceBounds.lean}; \src{Formalism/SharpRestriction/GrowthBounds/SizesAtRecord.lean}{SizesAtRecord.lean}; \src{Formalism/SharpRestriction/GrowthBounds/Joints.lean}{Joints.lean}. GII gives the actual good-orbit counts, uniform numerical supply, conditional bad-orbit probability and simultaneous tops selection.
\bibitem{type-one-family} \src{Formalism/SharpRestriction/GrowthBounds/LogarithmicSums/MeanNum.lean}{MeanNum.lean}; \src{Formalism/SharpRestriction/GrowthBounds/Joints.lean}{Joints.lean}. GI proves every actual Type-I family inequality, the exact band and window conversion, time comparison and canonical mean.
\bibitem{finite-base-change} \src{Formalism/LocalToGlobal/GroupLattices/FockModels/LPBase.lean}{LPBase.lean}; \src{Formalism/SharpRestriction/GrowthBounds/Joints.lean}{Joints.lean}. BC constructs the standard coefficient ring and prescribed-place extension, proves exact lattice and dual transport, and identifies the recC application. Its ledger follows the concrete construction to the individual ring declarations.
\bibitem{family-orbit-join} \src{Formalism/SharpRestriction/GrowthBounds/DefsInst.lean}{DefsInst.lean}; \src{Formalism/SharpRestriction/GrowthBounds/PlaneCountTheoremWithConstants.lean}{PlaneCountTheoremWithConstants.lean}; \src{Formalism/SharpRestriction/GrowthBounds/Joints.lean}{Joints.lean}. GJ proves the universal canonical infimum, uniform scales and actual finite-place join.
\bibitem{global-degree-choice} \src{Formalism/SharpRestriction/GrowthBounds/LogarithmicSums/SlopeEstimateOfLemmaBSlopeLine.lean}{SlopeEstimateOfLemmaBSlopeLine.lean}; \src{Formalism/LocalToGlobal/GlobalTheory/LossRequirementsOfSmallShares.lean}{LossRequirementsOfSmallShares.lean}; \src{Formalism/SharpRestriction/FamCore.lean}{FamCore.lean}. GW proves exact finite marginal averages and the single degree pair; it preserves the separate full-coverage obligation.
\bibitem{host-degree-realization} \src{Development/AutomorphicArithmetic/CharacterBoxes/HostWW.lean}{HostWW.lean}; \src{Development/AutomorphicArithmetic/CharacterBoxes/KitS.lean}{KitS.lean}; \src{Development/AbelianVarietyTheory/Proofs/GlobalInterpolation/ExistsFamPOfKits.lean}{ExistsFamPOfKits.lean}. HW proves the active host preselection, fixed-pair rebuilding and actual family transport, keeping the upstream construction premises explicit.


\bibitem{biased-orbit-probability} \src{Development/AutomorphicArithmetic/ThickFrames/MarkedLow.lean}{MarkedLow.lean}; \src{Development/AutomorphicArithmetic/ThickFrames/orbitTerms.lean}{orbitTerms.lean}; \src{Development/AutomorphicArithmetic/ThickFrames/orbitSum.lean}{orbitSum.lean}. BO proves the literal biased conditional estimate, marked-tile lower bound and transport to the prescribed based orbit.
\bibitem{balanced-tops} \src{Development/AutomorphicArithmetic/ThickFrames/select.lean}{select.lean}; \src{Development/AutomorphicArithmetic/ThickFrames/SharesOfCounts.lean}{SharesOfCounts.lean}; \src{Development/AutomorphicArithmetic/ThickFrames/topsGroup.lean}{topsGroup.lean}. BS gives one simultaneous choice with the exact host coset shares, cardinality and weighted bad mass; its Fourier calculation explains the character weights.
\bibitem{actual-host-tops} \src{Development/AutomorphicArithmetic/ThickFrames/topsChoiceH.lean}{topsChoiceH.lean}; \src{Development/AutomorphicArithmetic/CoprimeRanks/SecondSetNoRelation.lean}{SecondSetNoRelation.lean}; \src{Development/AutomorphicArithmetic/ShortNorms/hostOfFieldOf.lean}{hostOfFieldOf.lean}; \src{Development/AutomorphicArithmetic/ShortNorms/supplyOfAdmOf.lean}{supplyOfAdmOf.lean}. HC constructs the actual tops and no-relation set and expands the finite/algebraic host assembly; its field and analytic inputs are stated explicitly.
\bibitem{actual-local-places} \src{Development/AbelianVarietyTheory/Proofs/SatakeParameters/ExistsPlTieOfInertCore.lean}{ExistsPlTieOfInertCore.lean}; \src{Development/AbelianVarietyTheory/Proofs/IwasawaDecomposition/ExistsCompletionEquiv.lean}{ExistsCompletionEquiv.lean}; \src{Development/AbelianVarietyTheory/Proofs/IwasawaDecomposition/ExistsLocIdOfCompletions.lean}{ExistsLocIdOfCompletions.lean}; \src{Development/AbelianVarietyTheory/Proofs/IwasawaDecomposition/AddCharIsIntegral.lean}{AddCharIsIntegral.lean}. LI gives the exact local-place bijection, completions, splitting, frame, position compatibility and integral inverse character for the supplied data.

\end{thebibliography}
\end{document}
